% Punto de entrada aditivo del sucesor; el main.tex rector queda inmutable.
\RequirePackage{fix-cm}
\documentclass[11pt,a4paper,twoside,openany]{book}

\usepackage[T1]{fontenc}
\usepackage[utf8]{inputenc}
\DeclareUnicodeCharacter{03BC}{\ensuremath{\mu}}
\DeclareUnicodeCharacter{03C4}{\ensuremath{\tau}}
\DeclareUnicodeCharacter{03C3}{\ensuremath{\sigma}}
\DeclareUnicodeCharacter{03C0}{\ensuremath{\pi}}
\DeclareUnicodeCharacter{03B3}{\ensuremath{\gamma}}
\DeclareUnicodeCharacter{03A9}{\ensuremath{\Omega}}
\DeclareUnicodeCharacter{03C1}{\ensuremath{\rho}}
\DeclareUnicodeCharacter{03C6}{\ensuremath{\varphi}}
\DeclareUnicodeCharacter{03C8}{\ensuremath{\psi}}
\DeclareUnicodeCharacter{03A5}{\ensuremath{\Upsilon}}
\DeclareUnicodeCharacter{03B7}{\ensuremath{\eta}}
\DeclareUnicodeCharacter{03B1}{\ensuremath{\alpha}}
\DeclareUnicodeCharacter{03B5}{\ensuremath{\varepsilon}}
\DeclareUnicodeCharacter{03B6}{\ensuremath{\zeta}}
\DeclareUnicodeCharacter{03B8}{\ensuremath{\theta}}
\DeclareUnicodeCharacter{03BB}{\ensuremath{\lambda}}
\DeclareUnicodeCharacter{00B1}{\ensuremath{\pm}}
\DeclareUnicodeCharacter{2192}{\ensuremath{\longrightarrow}}
\DeclareUnicodeCharacter{2194}{\ensuremath{\leftrightarrow}}
\DeclareUnicodeCharacter{21A6}{\ensuremath{\longmapsto}}
\DeclareUnicodeCharacter{21D2}{\ensuremath{\Rightarrow}}
\DeclareUnicodeCharacter{2202}{\ensuremath{\partial}}
% lcdf-typetools glyphtounicode.tex, Version 2.95
% Contents: Glyph mapping information for pdftex, used for PDF searching
% Generated from:
% - glyphlist.txt, Version 2.0
% - texglyphlist.txt, Version 2.95
% - texglyphlist-g2u.txt, Version 2.95
\pdfglyphtounicode{A}{0041}
\pdfglyphtounicode{AE}{00C6}
\pdfglyphtounicode{AEacute}{01FC}
\pdfglyphtounicode{AEmacron}{01E2}
\pdfglyphtounicode{AEsmall}{00E6}
\pdfglyphtounicode{Aacute}{00C1}
\pdfglyphtounicode{Aacutesmall}{00E1}
\pdfglyphtounicode{Abreve}{0102}
\pdfglyphtounicode{Abreveacute}{1EAE}
\pdfglyphtounicode{Abrevecyrillic}{04D0}
\pdfglyphtounicode{Abrevedotbelow}{1EB6}
\pdfglyphtounicode{Abrevegrave}{1EB0}
\pdfglyphtounicode{Abrevehookabove}{1EB2}
\pdfglyphtounicode{Abrevetilde}{1EB4}
\pdfglyphtounicode{Acaron}{01CD}
\pdfglyphtounicode{Acircle}{24B6}
\pdfglyphtounicode{Acircumflex}{00C2}
\pdfglyphtounicode{Acircumflexacute}{1EA4}
\pdfglyphtounicode{Acircumflexdotbelow}{1EAC}
\pdfglyphtounicode{Acircumflexgrave}{1EA6}
\pdfglyphtounicode{Acircumflexhookabove}{1EA8}
\pdfglyphtounicode{Acircumflexsmall}{00E2}
\pdfglyphtounicode{Acircumflextilde}{1EAA}
\pdfglyphtounicode{Acute}{00B4}
\pdfglyphtounicode{Acutesmall}{00B4}
\pdfglyphtounicode{Acyrillic}{0410}
\pdfglyphtounicode{Adblgrave}{0200}
\pdfglyphtounicode{Adieresis}{00C4}
\pdfglyphtounicode{Adieresiscyrillic}{04D2}
\pdfglyphtounicode{Adieresismacron}{01DE}
\pdfglyphtounicode{Adieresissmall}{00E4}
\pdfglyphtounicode{Adotbelow}{1EA0}
\pdfglyphtounicode{Adotmacron}{01E0}
\pdfglyphtounicode{Agrave}{00C0}
\pdfglyphtounicode{Agravesmall}{00E0}
\pdfglyphtounicode{Ahookabove}{1EA2}
\pdfglyphtounicode{Aiecyrillic}{04D4}
\pdfglyphtounicode{Ainvertedbreve}{0202}
\pdfglyphtounicode{Alpha}{0391}
\pdfglyphtounicode{Alphatonos}{0386}
\pdfglyphtounicode{Amacron}{0100}
\pdfglyphtounicode{Amonospace}{FF21}
\pdfglyphtounicode{Aogonek}{0104}
\pdfglyphtounicode{Aring}{00C5}
\pdfglyphtounicode{Aringacute}{01FA}
\pdfglyphtounicode{Aringbelow}{1E00}
\pdfglyphtounicode{Aringsmall}{00E5}
\pdfglyphtounicode{Asmall}{0061}
\pdfglyphtounicode{Atilde}{00C3}
\pdfglyphtounicode{Atildesmall}{00E3}
\pdfglyphtounicode{Aybarmenian}{0531}
\pdfglyphtounicode{B}{0042}
\pdfglyphtounicode{Bcircle}{24B7}
\pdfglyphtounicode{Bdotaccent}{1E02}
\pdfglyphtounicode{Bdotbelow}{1E04}
\pdfglyphtounicode{Becyrillic}{0411}
\pdfglyphtounicode{Benarmenian}{0532}
\pdfglyphtounicode{Beta}{0392}
\pdfglyphtounicode{Bhook}{0181}
\pdfglyphtounicode{Blinebelow}{1E06}
\pdfglyphtounicode{Bmonospace}{FF22}
\pdfglyphtounicode{Brevesmall}{02D8}
\pdfglyphtounicode{Bsmall}{0062}
\pdfglyphtounicode{Btopbar}{0182}
\pdfglyphtounicode{C}{0043}
\pdfglyphtounicode{Caarmenian}{053E}
\pdfglyphtounicode{Cacute}{0106}
\pdfglyphtounicode{Caron}{02C7}
\pdfglyphtounicode{Caronsmall}{02C7}
\pdfglyphtounicode{Ccaron}{010C}
\pdfglyphtounicode{Ccedilla}{00C7}
\pdfglyphtounicode{Ccedillaacute}{1E08}
\pdfglyphtounicode{Ccedillasmall}{00E7}
\pdfglyphtounicode{Ccircle}{24B8}
\pdfglyphtounicode{Ccircumflex}{0108}
\pdfglyphtounicode{Cdot}{010A}
\pdfglyphtounicode{Cdotaccent}{010A}
\pdfglyphtounicode{Cedillasmall}{00B8}
\pdfglyphtounicode{Chaarmenian}{0549}
\pdfglyphtounicode{Cheabkhasiancyrillic}{04BC}
\pdfglyphtounicode{Checyrillic}{0427}
\pdfglyphtounicode{Chedescenderabkhasiancyrillic}{04BE}
\pdfglyphtounicode{Chedescendercyrillic}{04B6}
\pdfglyphtounicode{Chedieresiscyrillic}{04F4}
\pdfglyphtounicode{Cheharmenian}{0543}
\pdfglyphtounicode{Chekhakassiancyrillic}{04CB}
\pdfglyphtounicode{Cheverticalstrokecyrillic}{04B8}
\pdfglyphtounicode{Chi}{03A7}
\pdfglyphtounicode{Chook}{0187}
\pdfglyphtounicode{Circumflexsmall}{02C6}
\pdfglyphtounicode{Cmonospace}{FF23}
\pdfglyphtounicode{Coarmenian}{0551}
\pdfglyphtounicode{Csmall}{0063}
\pdfglyphtounicode{D}{0044}
\pdfglyphtounicode{DZ}{01F1}
\pdfglyphtounicode{DZcaron}{01C4}
\pdfglyphtounicode{Daarmenian}{0534}
\pdfglyphtounicode{Dafrican}{0189}
\pdfglyphtounicode{Dbar}{0110}
\pdfglyphtounicode{Dcaron}{010E}
\pdfglyphtounicode{Dcedilla}{1E10}
\pdfglyphtounicode{Dcircle}{24B9}
\pdfglyphtounicode{Dcircumflexbelow}{1E12}
\pdfglyphtounicode{Dcroat}{0110}
\pdfglyphtounicode{Ddotaccent}{1E0A}
\pdfglyphtounicode{Ddotbelow}{1E0C}
\pdfglyphtounicode{Decyrillic}{0414}
\pdfglyphtounicode{Deicoptic}{03EE}
\pdfglyphtounicode{Delta}{2206}
\pdfglyphtounicode{Deltagreek}{0394}
\pdfglyphtounicode{Dhook}{018A}
\pdfglyphtounicode{Dieresis}{00A8}
\pdfglyphtounicode{DieresisAcute}{F6CC}
\pdfglyphtounicode{DieresisGrave}{F6CD}
\pdfglyphtounicode{Dieresissmall}{00A8}
\pdfglyphtounicode{Digamma}{D875 DFCB}
\pdfglyphtounicode{Digammagreek}{03DC}
\pdfglyphtounicode{Djecyrillic}{0402}
\pdfglyphtounicode{Dlinebelow}{1E0E}
\pdfglyphtounicode{Dmonospace}{FF24}
\pdfglyphtounicode{Dotaccentsmall}{02D9}
\pdfglyphtounicode{Dslash}{0110}
\pdfglyphtounicode{Dsmall}{0064}
\pdfglyphtounicode{Dtopbar}{018B}
\pdfglyphtounicode{Dz}{01F2}
\pdfglyphtounicode{Dzcaron}{01C5}
\pdfglyphtounicode{Dzeabkhasiancyrillic}{04E0}
\pdfglyphtounicode{Dzecyrillic}{0405}
\pdfglyphtounicode{Dzhecyrillic}{040F}
\pdfglyphtounicode{E}{0045}
\pdfglyphtounicode{Eacute}{00C9}
\pdfglyphtounicode{Eacutesmall}{00E9}
\pdfglyphtounicode{Ebreve}{0114}
\pdfglyphtounicode{Ecaron}{011A}
\pdfglyphtounicode{Ecedillabreve}{1E1C}
\pdfglyphtounicode{Echarmenian}{0535}
\pdfglyphtounicode{Ecircle}{24BA}
\pdfglyphtounicode{Ecircumflex}{00CA}
\pdfglyphtounicode{Ecircumflexacute}{1EBE}
\pdfglyphtounicode{Ecircumflexbelow}{1E18}
\pdfglyphtounicode{Ecircumflexdotbelow}{1EC6}
\pdfglyphtounicode{Ecircumflexgrave}{1EC0}
\pdfglyphtounicode{Ecircumflexhookabove}{1EC2}
\pdfglyphtounicode{Ecircumflexsmall}{00EA}
\pdfglyphtounicode{Ecircumflextilde}{1EC4}
\pdfglyphtounicode{Ecyrillic}{0404}
\pdfglyphtounicode{Edblgrave}{0204}
\pdfglyphtounicode{Edieresis}{00CB}
\pdfglyphtounicode{Edieresissmall}{00EB}
\pdfglyphtounicode{Edot}{0116}
\pdfglyphtounicode{Edotaccent}{0116}
\pdfglyphtounicode{Edotbelow}{1EB8}
\pdfglyphtounicode{Efcyrillic}{0424}
\pdfglyphtounicode{Egrave}{00C8}
\pdfglyphtounicode{Egravesmall}{00E8}
\pdfglyphtounicode{Eharmenian}{0537}
\pdfglyphtounicode{Ehookabove}{1EBA}
\pdfglyphtounicode{Eightroman}{2167}
\pdfglyphtounicode{Einvertedbreve}{0206}
\pdfglyphtounicode{Eiotifiedcyrillic}{0464}
\pdfglyphtounicode{Elcyrillic}{041B}
\pdfglyphtounicode{Elevenroman}{216A}
\pdfglyphtounicode{Emacron}{0112}
\pdfglyphtounicode{Emacronacute}{1E16}
\pdfglyphtounicode{Emacrongrave}{1E14}
\pdfglyphtounicode{Emcyrillic}{041C}
\pdfglyphtounicode{Emonospace}{FF25}
\pdfglyphtounicode{Encyrillic}{041D}
\pdfglyphtounicode{Endescendercyrillic}{04A2}
\pdfglyphtounicode{Eng}{014A}
\pdfglyphtounicode{Enghecyrillic}{04A4}
\pdfglyphtounicode{Enhookcyrillic}{04C7}
\pdfglyphtounicode{Eogonek}{0118}
\pdfglyphtounicode{Eopen}{0190}
\pdfglyphtounicode{Epsilon}{0395}
\pdfglyphtounicode{Epsilontonos}{0388}
\pdfglyphtounicode{Ercyrillic}{0420}
\pdfglyphtounicode{Ereversed}{018E}
\pdfglyphtounicode{Ereversedcyrillic}{042D}
\pdfglyphtounicode{Escyrillic}{0421}
\pdfglyphtounicode{Esdescendercyrillic}{04AA}
\pdfglyphtounicode{Esh}{01A9}
\pdfglyphtounicode{Esmall}{0065}
\pdfglyphtounicode{Eta}{0397}
\pdfglyphtounicode{Etarmenian}{0538}
\pdfglyphtounicode{Etatonos}{0389}
\pdfglyphtounicode{Eth}{00D0}
\pdfglyphtounicode{Ethsmall}{00F0}
\pdfglyphtounicode{Etilde}{1EBC}
\pdfglyphtounicode{Etildebelow}{1E1A}
\pdfglyphtounicode{Euro}{20AC}
\pdfglyphtounicode{Ezh}{01B7}
\pdfglyphtounicode{Ezhcaron}{01EE}
\pdfglyphtounicode{Ezhreversed}{01B8}
\pdfglyphtounicode{F}{0046}
\pdfglyphtounicode{FFIsmall}{0066 0066 0069}
\pdfglyphtounicode{FFLsmall}{0066 0066 006C}
\pdfglyphtounicode{FFsmall}{0066 0066}
\pdfglyphtounicode{FIsmall}{0066 0069}
\pdfglyphtounicode{FLsmall}{0066 006C}
\pdfglyphtounicode{Fcircle}{24BB}
\pdfglyphtounicode{Fdotaccent}{1E1E}
\pdfglyphtounicode{Feharmenian}{0556}
\pdfglyphtounicode{Feicoptic}{03E4}
\pdfglyphtounicode{Fhook}{0191}
\pdfglyphtounicode{Finv}{2132}
\pdfglyphtounicode{Fitacyrillic}{0472}
\pdfglyphtounicode{Fiveroman}{2164}
\pdfglyphtounicode{Fmonospace}{FF26}
\pdfglyphtounicode{Fourroman}{2163}
\pdfglyphtounicode{Fsmall}{0066}
\pdfglyphtounicode{G}{0047}
\pdfglyphtounicode{GBsquare}{3387}
\pdfglyphtounicode{Gacute}{01F4}
\pdfglyphtounicode{Gamma}{0393}
\pdfglyphtounicode{Gammaafrican}{0194}
\pdfglyphtounicode{Gangiacoptic}{03EA}
\pdfglyphtounicode{Gbreve}{011E}
\pdfglyphtounicode{Gcaron}{01E6}
\pdfglyphtounicode{Gcedilla}{0122}
\pdfglyphtounicode{Gcircle}{24BC}
\pdfglyphtounicode{Gcircumflex}{011C}
\pdfglyphtounicode{Gcommaaccent}{0122}
\pdfglyphtounicode{Gdot}{0120}
\pdfglyphtounicode{Gdotaccent}{0120}
\pdfglyphtounicode{Gecyrillic}{0413}
\pdfglyphtounicode{Germandbls}{0053 0053}
\pdfglyphtounicode{Germandblssmall}{0073 0073}
\pdfglyphtounicode{Ghadarmenian}{0542}
\pdfglyphtounicode{Ghemiddlehookcyrillic}{0494}
\pdfglyphtounicode{Ghestrokecyrillic}{0492}
\pdfglyphtounicode{Gheupturncyrillic}{0490}
\pdfglyphtounicode{Ghook}{0193}
\pdfglyphtounicode{Gimarmenian}{0533}
\pdfglyphtounicode{Gjecyrillic}{0403}
\pdfglyphtounicode{Gmacron}{1E20}
\pdfglyphtounicode{Gmir}{2141}
\pdfglyphtounicode{Gmonospace}{FF27}
\pdfglyphtounicode{Grave}{0060}
\pdfglyphtounicode{Gravesmall}{0060}
\pdfglyphtounicode{Gsmall}{0067}
\pdfglyphtounicode{Gsmallhook}{029B}
\pdfglyphtounicode{Gstroke}{01E4}
\pdfglyphtounicode{H}{0048}
\pdfglyphtounicode{H18533}{25CF}
\pdfglyphtounicode{H18543}{25AA}
\pdfglyphtounicode{H18551}{25AB}
\pdfglyphtounicode{H22073}{25A1}
\pdfglyphtounicode{HPsquare}{33CB}
\pdfglyphtounicode{Haabkhasiancyrillic}{04A8}
\pdfglyphtounicode{Hadescendercyrillic}{04B2}
\pdfglyphtounicode{Hardsigncyrillic}{042A}
\pdfglyphtounicode{Hbar}{0126}
\pdfglyphtounicode{Hbrevebelow}{1E2A}
\pdfglyphtounicode{Hcedilla}{1E28}
\pdfglyphtounicode{Hcircle}{24BD}
\pdfglyphtounicode{Hcircumflex}{0124}
\pdfglyphtounicode{Hdieresis}{1E26}
\pdfglyphtounicode{Hdotaccent}{1E22}
\pdfglyphtounicode{Hdotbelow}{1E24}
\pdfglyphtounicode{Hmonospace}{FF28}
\pdfglyphtounicode{Hoarmenian}{0540}
\pdfglyphtounicode{Horicoptic}{03E8}
\pdfglyphtounicode{Hsmall}{0068}
\pdfglyphtounicode{Hungarumlaut}{02DD}
\pdfglyphtounicode{Hungarumlautsmall}{02DD}
\pdfglyphtounicode{Hzsquare}{3390}
\pdfglyphtounicode{I}{0049}
\pdfglyphtounicode{IAcyrillic}{042F}
\pdfglyphtounicode{IJ}{0132}
\pdfglyphtounicode{IUcyrillic}{042E}
\pdfglyphtounicode{Iacute}{00CD}
\pdfglyphtounicode{Iacutesmall}{00ED}
\pdfglyphtounicode{Ibreve}{012C}
\pdfglyphtounicode{Icaron}{01CF}
\pdfglyphtounicode{Icircle}{24BE}
\pdfglyphtounicode{Icircumflex}{00CE}
\pdfglyphtounicode{Icircumflexsmall}{00EE}
\pdfglyphtounicode{Icyrillic}{0406}
\pdfglyphtounicode{Idblgrave}{0208}
\pdfglyphtounicode{Idieresis}{00CF}
\pdfglyphtounicode{Idieresisacute}{1E2E}
\pdfglyphtounicode{Idieresiscyrillic}{04E4}
\pdfglyphtounicode{Idieresissmall}{00EF}
\pdfglyphtounicode{Idot}{0130}
\pdfglyphtounicode{Idotaccent}{0130}
\pdfglyphtounicode{Idotbelow}{1ECA}
\pdfglyphtounicode{Iebrevecyrillic}{04D6}
\pdfglyphtounicode{Iecyrillic}{0415}
\pdfglyphtounicode{Ifractur}{2111}
\pdfglyphtounicode{Ifraktur}{2111}
\pdfglyphtounicode{Igrave}{00CC}
\pdfglyphtounicode{Igravesmall}{00EC}
\pdfglyphtounicode{Ihookabove}{1EC8}
\pdfglyphtounicode{Iicyrillic}{0418}
\pdfglyphtounicode{Iinvertedbreve}{020A}
\pdfglyphtounicode{Iishortcyrillic}{0419}
\pdfglyphtounicode{Imacron}{012A}
\pdfglyphtounicode{Imacroncyrillic}{04E2}
\pdfglyphtounicode{Imonospace}{FF29}
\pdfglyphtounicode{Iniarmenian}{053B}
\pdfglyphtounicode{Iocyrillic}{0401}
\pdfglyphtounicode{Iogonek}{012E}
\pdfglyphtounicode{Iota}{0399}
\pdfglyphtounicode{Iotaafrican}{0196}
\pdfglyphtounicode{Iotadieresis}{03AA}
\pdfglyphtounicode{Iotatonos}{038A}
\pdfglyphtounicode{Ismall}{0069}
\pdfglyphtounicode{Istroke}{0197}
\pdfglyphtounicode{Itilde}{0128}
\pdfglyphtounicode{Itildebelow}{1E2C}
\pdfglyphtounicode{Izhitsacyrillic}{0474}
\pdfglyphtounicode{Izhitsadblgravecyrillic}{0476}
\pdfglyphtounicode{J}{004A}
\pdfglyphtounicode{Jaarmenian}{0541}
\pdfglyphtounicode{Jcircle}{24BF}
\pdfglyphtounicode{Jcircumflex}{0134}
\pdfglyphtounicode{Jecyrillic}{0408}
\pdfglyphtounicode{Jheharmenian}{054B}
\pdfglyphtounicode{Jmonospace}{FF2A}
\pdfglyphtounicode{Jsmall}{006A}
\pdfglyphtounicode{K}{004B}
\pdfglyphtounicode{KBsquare}{3385}
\pdfglyphtounicode{KKsquare}{33CD}
\pdfglyphtounicode{Kabashkircyrillic}{04A0}
\pdfglyphtounicode{Kacute}{1E30}
\pdfglyphtounicode{Kacyrillic}{041A}
\pdfglyphtounicode{Kadescendercyrillic}{049A}
\pdfglyphtounicode{Kahookcyrillic}{04C3}
\pdfglyphtounicode{Kappa}{039A}
\pdfglyphtounicode{Kastrokecyrillic}{049E}
\pdfglyphtounicode{Kaverticalstrokecyrillic}{049C}
\pdfglyphtounicode{Kcaron}{01E8}
\pdfglyphtounicode{Kcedilla}{0136}
\pdfglyphtounicode{Kcircle}{24C0}
\pdfglyphtounicode{Kcommaaccent}{0136}
\pdfglyphtounicode{Kdotbelow}{1E32}
\pdfglyphtounicode{Keharmenian}{0554}
\pdfglyphtounicode{Kenarmenian}{053F}
\pdfglyphtounicode{Khacyrillic}{0425}
\pdfglyphtounicode{Kheicoptic}{03E6}
\pdfglyphtounicode{Khook}{0198}
\pdfglyphtounicode{Kjecyrillic}{040C}
\pdfglyphtounicode{Klinebelow}{1E34}
\pdfglyphtounicode{Kmonospace}{FF2B}
\pdfglyphtounicode{Koppacyrillic}{0480}
\pdfglyphtounicode{Koppagreek}{03DE}
\pdfglyphtounicode{Ksicyrillic}{046E}
\pdfglyphtounicode{Ksmall}{006B}
\pdfglyphtounicode{L}{004C}
\pdfglyphtounicode{LJ}{01C7}
\pdfglyphtounicode{LL}{004C 004C}
\pdfglyphtounicode{Lacute}{0139}
\pdfglyphtounicode{Lambda}{039B}
\pdfglyphtounicode{Lcaron}{013D}
\pdfglyphtounicode{Lcedilla}{013B}
\pdfglyphtounicode{Lcircle}{24C1}
\pdfglyphtounicode{Lcircumflexbelow}{1E3C}
\pdfglyphtounicode{Lcommaaccent}{013B}
\pdfglyphtounicode{Ldot}{013F}
\pdfglyphtounicode{Ldotaccent}{013F}
\pdfglyphtounicode{Ldotbelow}{1E36}
\pdfglyphtounicode{Ldotbelowmacron}{1E38}
\pdfglyphtounicode{Liwnarmenian}{053C}
\pdfglyphtounicode{Lj}{01C8}
\pdfglyphtounicode{Ljecyrillic}{0409}
\pdfglyphtounicode{Llinebelow}{1E3A}
\pdfglyphtounicode{Lmonospace}{FF2C}
\pdfglyphtounicode{Lslash}{0141}
\pdfglyphtounicode{Lslashsmall}{0142}
\pdfglyphtounicode{Lsmall}{006C}
\pdfglyphtounicode{M}{004D}
\pdfglyphtounicode{MBsquare}{3386}
\pdfglyphtounicode{Macron}{00AF}
\pdfglyphtounicode{Macronsmall}{00AF}
\pdfglyphtounicode{Macute}{1E3E}
\pdfglyphtounicode{Mcircle}{24C2}
\pdfglyphtounicode{Mdotaccent}{1E40}
\pdfglyphtounicode{Mdotbelow}{1E42}
\pdfglyphtounicode{Menarmenian}{0544}
\pdfglyphtounicode{Mmonospace}{FF2D}
\pdfglyphtounicode{Msmall}{006D}
\pdfglyphtounicode{Mturned}{019C}
\pdfglyphtounicode{Mu}{039C}
\pdfglyphtounicode{N}{004E}
\pdfglyphtounicode{NJ}{01CA}
\pdfglyphtounicode{Nacute}{0143}
\pdfglyphtounicode{Ncaron}{0147}
\pdfglyphtounicode{Ncedilla}{0145}
\pdfglyphtounicode{Ncircle}{24C3}
\pdfglyphtounicode{Ncircumflexbelow}{1E4A}
\pdfglyphtounicode{Ncommaaccent}{0145}
\pdfglyphtounicode{Ndotaccent}{1E44}
\pdfglyphtounicode{Ndotbelow}{1E46}
\pdfglyphtounicode{Ng}{014A}
\pdfglyphtounicode{Nhookleft}{019D}
\pdfglyphtounicode{Nineroman}{2168}
\pdfglyphtounicode{Nj}{01CB}
\pdfglyphtounicode{Njecyrillic}{040A}
\pdfglyphtounicode{Nlinebelow}{1E48}
\pdfglyphtounicode{Nmonospace}{FF2E}
\pdfglyphtounicode{Nowarmenian}{0546}
\pdfglyphtounicode{Nsmall}{006E}
\pdfglyphtounicode{Ntilde}{00D1}
\pdfglyphtounicode{Ntildesmall}{00F1}
\pdfglyphtounicode{Nu}{039D}
\pdfglyphtounicode{O}{004F}
\pdfglyphtounicode{OE}{0152}
\pdfglyphtounicode{OEsmall}{0153}
\pdfglyphtounicode{Oacute}{00D3}
\pdfglyphtounicode{Oacutesmall}{00F3}
\pdfglyphtounicode{Obarredcyrillic}{04E8}
\pdfglyphtounicode{Obarreddieresiscyrillic}{04EA}
\pdfglyphtounicode{Obreve}{014E}
\pdfglyphtounicode{Ocaron}{01D1}
\pdfglyphtounicode{Ocenteredtilde}{019F}
\pdfglyphtounicode{Ocircle}{24C4}
\pdfglyphtounicode{Ocircumflex}{00D4}
\pdfglyphtounicode{Ocircumflexacute}{1ED0}
\pdfglyphtounicode{Ocircumflexdotbelow}{1ED8}
\pdfglyphtounicode{Ocircumflexgrave}{1ED2}
\pdfglyphtounicode{Ocircumflexhookabove}{1ED4}
\pdfglyphtounicode{Ocircumflexsmall}{00F4}
\pdfglyphtounicode{Ocircumflextilde}{1ED6}
\pdfglyphtounicode{Ocyrillic}{041E}
\pdfglyphtounicode{Odblacute}{0150}
\pdfglyphtounicode{Odblgrave}{020C}
\pdfglyphtounicode{Odieresis}{00D6}
\pdfglyphtounicode{Odieresiscyrillic}{04E6}
\pdfglyphtounicode{Odieresissmall}{00F6}
\pdfglyphtounicode{Odotbelow}{1ECC}
\pdfglyphtounicode{Ogoneksmall}{02DB}
\pdfglyphtounicode{Ograve}{00D2}
\pdfglyphtounicode{Ogravesmall}{00F2}
\pdfglyphtounicode{Oharmenian}{0555}
\pdfglyphtounicode{Ohm}{2126}
\pdfglyphtounicode{Ohookabove}{1ECE}
\pdfglyphtounicode{Ohorn}{01A0}
\pdfglyphtounicode{Ohornacute}{1EDA}
\pdfglyphtounicode{Ohorndotbelow}{1EE2}
\pdfglyphtounicode{Ohorngrave}{1EDC}
\pdfglyphtounicode{Ohornhookabove}{1EDE}
\pdfglyphtounicode{Ohorntilde}{1EE0}
\pdfglyphtounicode{Ohungarumlaut}{0150}
\pdfglyphtounicode{Oi}{01A2}
\pdfglyphtounicode{Oinvertedbreve}{020E}
\pdfglyphtounicode{Omacron}{014C}
\pdfglyphtounicode{Omacronacute}{1E52}
\pdfglyphtounicode{Omacrongrave}{1E50}
\pdfglyphtounicode{Omega}{2126}
\pdfglyphtounicode{Omegacyrillic}{0460}
\pdfglyphtounicode{Omegagreek}{03A9}
\pdfglyphtounicode{Omegainv}{2127}
\pdfglyphtounicode{Omegaroundcyrillic}{047A}
\pdfglyphtounicode{Omegatitlocyrillic}{047C}
\pdfglyphtounicode{Omegatonos}{038F}
\pdfglyphtounicode{Omicron}{039F}
\pdfglyphtounicode{Omicrontonos}{038C}
\pdfglyphtounicode{Omonospace}{FF2F}
\pdfglyphtounicode{Oneroman}{2160}
\pdfglyphtounicode{Oogonek}{01EA}
\pdfglyphtounicode{Oogonekmacron}{01EC}
\pdfglyphtounicode{Oopen}{0186}
\pdfglyphtounicode{Oslash}{00D8}
\pdfglyphtounicode{Oslashacute}{01FE}
\pdfglyphtounicode{Oslashsmall}{00F8}
\pdfglyphtounicode{Osmall}{006F}
\pdfglyphtounicode{Ostrokeacute}{01FE}
\pdfglyphtounicode{Otcyrillic}{047E}
\pdfglyphtounicode{Otilde}{00D5}
\pdfglyphtounicode{Otildeacute}{1E4C}
\pdfglyphtounicode{Otildedieresis}{1E4E}
\pdfglyphtounicode{Otildesmall}{00F5}
\pdfglyphtounicode{P}{0050}
\pdfglyphtounicode{Pacute}{1E54}
\pdfglyphtounicode{Pcircle}{24C5}
\pdfglyphtounicode{Pdotaccent}{1E56}
\pdfglyphtounicode{Pecyrillic}{041F}
\pdfglyphtounicode{Peharmenian}{054A}
\pdfglyphtounicode{Pemiddlehookcyrillic}{04A6}
\pdfglyphtounicode{Phi}{03A6}
\pdfglyphtounicode{Phook}{01A4}
\pdfglyphtounicode{Pi}{03A0}
\pdfglyphtounicode{Piwrarmenian}{0553}
\pdfglyphtounicode{Pmonospace}{FF30}
\pdfglyphtounicode{Psi}{03A8}
\pdfglyphtounicode{Psicyrillic}{0470}
\pdfglyphtounicode{Psmall}{0070}
\pdfglyphtounicode{Q}{0051}
\pdfglyphtounicode{Qcircle}{24C6}
\pdfglyphtounicode{Qmonospace}{FF31}
\pdfglyphtounicode{Qsmall}{0071}
\pdfglyphtounicode{R}{0052}
\pdfglyphtounicode{Raarmenian}{054C}
\pdfglyphtounicode{Racute}{0154}
\pdfglyphtounicode{Rcaron}{0158}
\pdfglyphtounicode{Rcedilla}{0156}
\pdfglyphtounicode{Rcircle}{24C7}
\pdfglyphtounicode{Rcommaaccent}{0156}
\pdfglyphtounicode{Rdblgrave}{0210}
\pdfglyphtounicode{Rdotaccent}{1E58}
\pdfglyphtounicode{Rdotbelow}{1E5A}
\pdfglyphtounicode{Rdotbelowmacron}{1E5C}
\pdfglyphtounicode{Reharmenian}{0550}
\pdfglyphtounicode{Rfractur}{211C}
\pdfglyphtounicode{Rfraktur}{211C}
\pdfglyphtounicode{Rho}{03A1}
\pdfglyphtounicode{Ringsmall}{02DA}
\pdfglyphtounicode{Rinvertedbreve}{0212}
\pdfglyphtounicode{Rlinebelow}{1E5E}
\pdfglyphtounicode{Rmonospace}{FF32}
\pdfglyphtounicode{Rsmall}{0072}
\pdfglyphtounicode{Rsmallinverted}{0281}
\pdfglyphtounicode{Rsmallinvertedsuperior}{02B6}
\pdfglyphtounicode{S}{0053}
\pdfglyphtounicode{SF010000}{250C}
\pdfglyphtounicode{SF020000}{2514}
\pdfglyphtounicode{SF030000}{2510}
\pdfglyphtounicode{SF040000}{2518}
\pdfglyphtounicode{SF050000}{253C}
\pdfglyphtounicode{SF060000}{252C}
\pdfglyphtounicode{SF070000}{2534}
\pdfglyphtounicode{SF080000}{251C}
\pdfglyphtounicode{SF090000}{2524}
\pdfglyphtounicode{SF100000}{2500}
\pdfglyphtounicode{SF110000}{2502}
\pdfglyphtounicode{SF190000}{2561}
\pdfglyphtounicode{SF200000}{2562}
\pdfglyphtounicode{SF210000}{2556}
\pdfglyphtounicode{SF220000}{2555}
\pdfglyphtounicode{SF230000}{2563}
\pdfglyphtounicode{SF240000}{2551}
\pdfglyphtounicode{SF250000}{2557}
\pdfglyphtounicode{SF260000}{255D}
\pdfglyphtounicode{SF270000}{255C}
\pdfglyphtounicode{SF280000}{255B}
\pdfglyphtounicode{SF360000}{255E}
\pdfglyphtounicode{SF370000}{255F}
\pdfglyphtounicode{SF380000}{255A}
\pdfglyphtounicode{SF390000}{2554}
\pdfglyphtounicode{SF400000}{2569}
\pdfglyphtounicode{SF410000}{2566}
\pdfglyphtounicode{SF420000}{2560}
\pdfglyphtounicode{SF430000}{2550}
\pdfglyphtounicode{SF440000}{256C}
\pdfglyphtounicode{SF450000}{2567}
\pdfglyphtounicode{SF460000}{2568}
\pdfglyphtounicode{SF470000}{2564}
\pdfglyphtounicode{SF480000}{2565}
\pdfglyphtounicode{SF490000}{2559}
\pdfglyphtounicode{SF500000}{2558}
\pdfglyphtounicode{SF510000}{2552}
\pdfglyphtounicode{SF520000}{2553}
\pdfglyphtounicode{SF530000}{256B}
\pdfglyphtounicode{SF540000}{256A}
\pdfglyphtounicode{SS}{0053 0053}
\pdfglyphtounicode{SSsmall}{0073 0073}
\pdfglyphtounicode{Sacute}{015A}
\pdfglyphtounicode{Sacutedotaccent}{1E64}
\pdfglyphtounicode{Sampigreek}{03E0}
\pdfglyphtounicode{Scaron}{0160}
\pdfglyphtounicode{Scarondotaccent}{1E66}
\pdfglyphtounicode{Scaronsmall}{0161}
\pdfglyphtounicode{Scedilla}{015E}
\pdfglyphtounicode{Schwa}{018F}
\pdfglyphtounicode{Schwacyrillic}{04D8}
\pdfglyphtounicode{Schwadieresiscyrillic}{04DA}
\pdfglyphtounicode{Scircle}{24C8}
\pdfglyphtounicode{Scircumflex}{015C}
\pdfglyphtounicode{Scommaaccent}{0218}
\pdfglyphtounicode{Sdotaccent}{1E60}
\pdfglyphtounicode{Sdotbelow}{1E62}
\pdfglyphtounicode{Sdotbelowdotaccent}{1E68}
\pdfglyphtounicode{Seharmenian}{054D}
\pdfglyphtounicode{Sevenroman}{2166}
\pdfglyphtounicode{Shaarmenian}{0547}
\pdfglyphtounicode{Shacyrillic}{0428}
\pdfglyphtounicode{Shchacyrillic}{0429}
\pdfglyphtounicode{Sheicoptic}{03E2}
\pdfglyphtounicode{Shhacyrillic}{04BA}
\pdfglyphtounicode{Shimacoptic}{03EC}
\pdfglyphtounicode{Sigma}{03A3}
\pdfglyphtounicode{Sixroman}{2165}
\pdfglyphtounicode{Smonospace}{FF33}
\pdfglyphtounicode{Softsigncyrillic}{042C}
\pdfglyphtounicode{Ssmall}{0073}
\pdfglyphtounicode{Stigmagreek}{03DA}
\pdfglyphtounicode{T}{0054}
\pdfglyphtounicode{Tau}{03A4}
\pdfglyphtounicode{Tbar}{0166}
\pdfglyphtounicode{Tcaron}{0164}
\pdfglyphtounicode{Tcedilla}{0162}
\pdfglyphtounicode{Tcircle}{24C9}
\pdfglyphtounicode{Tcircumflexbelow}{1E70}
\pdfglyphtounicode{Tcommaaccent}{0162}
\pdfglyphtounicode{Tdotaccent}{1E6A}
\pdfglyphtounicode{Tdotbelow}{1E6C}
\pdfglyphtounicode{Tecyrillic}{0422}
\pdfglyphtounicode{Tedescendercyrillic}{04AC}
\pdfglyphtounicode{Tenroman}{2169}
\pdfglyphtounicode{Tetsecyrillic}{04B4}
\pdfglyphtounicode{Theta}{0398}
\pdfglyphtounicode{Thook}{01AC}
\pdfglyphtounicode{Thorn}{00DE}
\pdfglyphtounicode{Thornsmall}{00FE}
\pdfglyphtounicode{Threeroman}{2162}
\pdfglyphtounicode{Tildesmall}{02DC}
\pdfglyphtounicode{Tiwnarmenian}{054F}
\pdfglyphtounicode{Tlinebelow}{1E6E}
\pdfglyphtounicode{Tmonospace}{FF34}
\pdfglyphtounicode{Toarmenian}{0539}
\pdfglyphtounicode{Tonefive}{01BC}
\pdfglyphtounicode{Tonesix}{0184}
\pdfglyphtounicode{Tonetwo}{01A7}
\pdfglyphtounicode{Tretroflexhook}{01AE}
\pdfglyphtounicode{Tsecyrillic}{0426}
\pdfglyphtounicode{Tshecyrillic}{040B}
\pdfglyphtounicode{Tsmall}{0074}
\pdfglyphtounicode{Twelveroman}{216B}
\pdfglyphtounicode{Tworoman}{2161}
\pdfglyphtounicode{U}{0055}
\pdfglyphtounicode{Uacute}{00DA}
\pdfglyphtounicode{Uacutesmall}{00FA}
\pdfglyphtounicode{Ubreve}{016C}
\pdfglyphtounicode{Ucaron}{01D3}
\pdfglyphtounicode{Ucircle}{24CA}
\pdfglyphtounicode{Ucircumflex}{00DB}
\pdfglyphtounicode{Ucircumflexbelow}{1E76}
\pdfglyphtounicode{Ucircumflexsmall}{00FB}
\pdfglyphtounicode{Ucyrillic}{0423}
\pdfglyphtounicode{Udblacute}{0170}
\pdfglyphtounicode{Udblgrave}{0214}
\pdfglyphtounicode{Udieresis}{00DC}
\pdfglyphtounicode{Udieresisacute}{01D7}
\pdfglyphtounicode{Udieresisbelow}{1E72}
\pdfglyphtounicode{Udieresiscaron}{01D9}
\pdfglyphtounicode{Udieresiscyrillic}{04F0}
\pdfglyphtounicode{Udieresisgrave}{01DB}
\pdfglyphtounicode{Udieresismacron}{01D5}
\pdfglyphtounicode{Udieresissmall}{00FC}
\pdfglyphtounicode{Udotbelow}{1EE4}
\pdfglyphtounicode{Ugrave}{00D9}
\pdfglyphtounicode{Ugravesmall}{00F9}
\pdfglyphtounicode{Uhookabove}{1EE6}
\pdfglyphtounicode{Uhorn}{01AF}
\pdfglyphtounicode{Uhornacute}{1EE8}
\pdfglyphtounicode{Uhorndotbelow}{1EF0}
\pdfglyphtounicode{Uhorngrave}{1EEA}
\pdfglyphtounicode{Uhornhookabove}{1EEC}
\pdfglyphtounicode{Uhorntilde}{1EEE}
\pdfglyphtounicode{Uhungarumlaut}{0170}
\pdfglyphtounicode{Uhungarumlautcyrillic}{04F2}
\pdfglyphtounicode{Uinvertedbreve}{0216}
\pdfglyphtounicode{Ukcyrillic}{0478}
\pdfglyphtounicode{Umacron}{016A}
\pdfglyphtounicode{Umacroncyrillic}{04EE}
\pdfglyphtounicode{Umacrondieresis}{1E7A}
\pdfglyphtounicode{Umonospace}{FF35}
\pdfglyphtounicode{Uogonek}{0172}
\pdfglyphtounicode{Upsilon}{03A5}
\pdfglyphtounicode{Upsilon1}{03D2}
\pdfglyphtounicode{Upsilonacutehooksymbolgreek}{03D3}
\pdfglyphtounicode{Upsilonafrican}{01B1}
\pdfglyphtounicode{Upsilondieresis}{03AB}
\pdfglyphtounicode{Upsilondieresishooksymbolgreek}{03D4}
\pdfglyphtounicode{Upsilonhooksymbol}{03D2}
\pdfglyphtounicode{Upsilontonos}{038E}
\pdfglyphtounicode{Uring}{016E}
\pdfglyphtounicode{Ushortcyrillic}{040E}
\pdfglyphtounicode{Usmall}{0075}
\pdfglyphtounicode{Ustraightcyrillic}{04AE}
\pdfglyphtounicode{Ustraightstrokecyrillic}{04B0}
\pdfglyphtounicode{Utilde}{0168}
\pdfglyphtounicode{Utildeacute}{1E78}
\pdfglyphtounicode{Utildebelow}{1E74}
\pdfglyphtounicode{V}{0056}
\pdfglyphtounicode{Vcircle}{24CB}
\pdfglyphtounicode{Vdotbelow}{1E7E}
\pdfglyphtounicode{Vecyrillic}{0412}
\pdfglyphtounicode{Vewarmenian}{054E}
\pdfglyphtounicode{Vhook}{01B2}
\pdfglyphtounicode{Vmonospace}{FF36}
\pdfglyphtounicode{Voarmenian}{0548}
\pdfglyphtounicode{Vsmall}{0076}
\pdfglyphtounicode{Vtilde}{1E7C}
\pdfglyphtounicode{W}{0057}
\pdfglyphtounicode{Wacute}{1E82}
\pdfglyphtounicode{Wcircle}{24CC}
\pdfglyphtounicode{Wcircumflex}{0174}
\pdfglyphtounicode{Wdieresis}{1E84}
\pdfglyphtounicode{Wdotaccent}{1E86}
\pdfglyphtounicode{Wdotbelow}{1E88}
\pdfglyphtounicode{Wgrave}{1E80}
\pdfglyphtounicode{Wmonospace}{FF37}
\pdfglyphtounicode{Wsmall}{0077}
\pdfglyphtounicode{X}{0058}
\pdfglyphtounicode{Xcircle}{24CD}
\pdfglyphtounicode{Xdieresis}{1E8C}
\pdfglyphtounicode{Xdotaccent}{1E8A}
\pdfglyphtounicode{Xeharmenian}{053D}
\pdfglyphtounicode{Xi}{039E}
\pdfglyphtounicode{Xmonospace}{FF38}
\pdfglyphtounicode{Xsmall}{0078}
\pdfglyphtounicode{Y}{0059}
\pdfglyphtounicode{Yacute}{00DD}
\pdfglyphtounicode{Yacutesmall}{00FD}
\pdfglyphtounicode{Yatcyrillic}{0462}
\pdfglyphtounicode{Ycircle}{24CE}
\pdfglyphtounicode{Ycircumflex}{0176}
\pdfglyphtounicode{Ydieresis}{0178}
\pdfglyphtounicode{Ydieresissmall}{00FF}
\pdfglyphtounicode{Ydotaccent}{1E8E}
\pdfglyphtounicode{Ydotbelow}{1EF4}
\pdfglyphtounicode{Yen}{00A5}
\pdfglyphtounicode{Yericyrillic}{042B}
\pdfglyphtounicode{Yerudieresiscyrillic}{04F8}
\pdfglyphtounicode{Ygrave}{1EF2}
\pdfglyphtounicode{Yhook}{01B3}
\pdfglyphtounicode{Yhookabove}{1EF6}
\pdfglyphtounicode{Yiarmenian}{0545}
\pdfglyphtounicode{Yicyrillic}{0407}
\pdfglyphtounicode{Yiwnarmenian}{0552}
\pdfglyphtounicode{Ymonospace}{FF39}
\pdfglyphtounicode{Ysmall}{0079}
\pdfglyphtounicode{Ytilde}{1EF8}
\pdfglyphtounicode{Yusbigcyrillic}{046A}
\pdfglyphtounicode{Yusbigiotifiedcyrillic}{046C}
\pdfglyphtounicode{Yuslittlecyrillic}{0466}
\pdfglyphtounicode{Yuslittleiotifiedcyrillic}{0468}
\pdfglyphtounicode{Z}{005A}
\pdfglyphtounicode{Zaarmenian}{0536}
\pdfglyphtounicode{Zacute}{0179}
\pdfglyphtounicode{Zcaron}{017D}
\pdfglyphtounicode{Zcaronsmall}{017E}
\pdfglyphtounicode{Zcircle}{24CF}
\pdfglyphtounicode{Zcircumflex}{1E90}
\pdfglyphtounicode{Zdot}{017B}
\pdfglyphtounicode{Zdotaccent}{017B}
\pdfglyphtounicode{Zdotbelow}{1E92}
\pdfglyphtounicode{Zecyrillic}{0417}
\pdfglyphtounicode{Zedescendercyrillic}{0498}
\pdfglyphtounicode{Zedieresiscyrillic}{04DE}
\pdfglyphtounicode{Zeta}{0396}
\pdfglyphtounicode{Zhearmenian}{053A}
\pdfglyphtounicode{Zhebrevecyrillic}{04C1}
\pdfglyphtounicode{Zhecyrillic}{0416}
\pdfglyphtounicode{Zhedescendercyrillic}{0496}
\pdfglyphtounicode{Zhedieresiscyrillic}{04DC}
\pdfglyphtounicode{Zlinebelow}{1E94}
\pdfglyphtounicode{Zmonospace}{FF3A}
\pdfglyphtounicode{Zsmall}{007A}
\pdfglyphtounicode{Zstroke}{01B5}
\pdfglyphtounicode{a}{0061}
\pdfglyphtounicode{aabengali}{0986}
\pdfglyphtounicode{aacute}{00E1}
\pdfglyphtounicode{aadeva}{0906}
\pdfglyphtounicode{aagujarati}{0A86}
\pdfglyphtounicode{aagurmukhi}{0A06}
\pdfglyphtounicode{aamatragurmukhi}{0A3E}
\pdfglyphtounicode{aarusquare}{3303}
\pdfglyphtounicode{aavowelsignbengali}{09BE}
\pdfglyphtounicode{aavowelsigndeva}{093E}
\pdfglyphtounicode{aavowelsigngujarati}{0ABE}
\pdfglyphtounicode{abbreviationmarkarmenian}{055F}
\pdfglyphtounicode{abbreviationsigndeva}{0970}
\pdfglyphtounicode{abengali}{0985}
\pdfglyphtounicode{abopomofo}{311A}
\pdfglyphtounicode{abreve}{0103}
\pdfglyphtounicode{abreveacute}{1EAF}
\pdfglyphtounicode{abrevecyrillic}{04D1}
\pdfglyphtounicode{abrevedotbelow}{1EB7}
\pdfglyphtounicode{abrevegrave}{1EB1}
\pdfglyphtounicode{abrevehookabove}{1EB3}
\pdfglyphtounicode{abrevetilde}{1EB5}
\pdfglyphtounicode{acaron}{01CE}
\pdfglyphtounicode{acircle}{24D0}
\pdfglyphtounicode{acircumflex}{00E2}
\pdfglyphtounicode{acircumflexacute}{1EA5}
\pdfglyphtounicode{acircumflexdotbelow}{1EAD}
\pdfglyphtounicode{acircumflexgrave}{1EA7}
\pdfglyphtounicode{acircumflexhookabove}{1EA9}
\pdfglyphtounicode{acircumflextilde}{1EAB}
\pdfglyphtounicode{acute}{00B4}
\pdfglyphtounicode{acutebelowcmb}{0317}
\pdfglyphtounicode{acutecmb}{0301}
\pdfglyphtounicode{acutecomb}{0301}
\pdfglyphtounicode{acutedeva}{0954}
\pdfglyphtounicode{acutelowmod}{02CF}
\pdfglyphtounicode{acutetonecmb}{0341}
\pdfglyphtounicode{acyrillic}{0430}
\pdfglyphtounicode{adblgrave}{0201}
\pdfglyphtounicode{addakgurmukhi}{0A71}
\pdfglyphtounicode{adeva}{0905}
\pdfglyphtounicode{adieresis}{00E4}
\pdfglyphtounicode{adieresiscyrillic}{04D3}
\pdfglyphtounicode{adieresismacron}{01DF}
\pdfglyphtounicode{adotbelow}{1EA1}
\pdfglyphtounicode{adotmacron}{01E1}
\pdfglyphtounicode{ae}{00E6}
\pdfglyphtounicode{aeacute}{01FD}
\pdfglyphtounicode{aekorean}{3150}
\pdfglyphtounicode{aemacron}{01E3}
\pdfglyphtounicode{afii00208}{2015}
\pdfglyphtounicode{afii08941}{20A4}
\pdfglyphtounicode{afii10017}{0410}
\pdfglyphtounicode{afii10018}{0411}
\pdfglyphtounicode{afii10019}{0412}
\pdfglyphtounicode{afii10020}{0413}
\pdfglyphtounicode{afii10021}{0414}
\pdfglyphtounicode{afii10022}{0415}
\pdfglyphtounicode{afii10023}{0401}
\pdfglyphtounicode{afii10024}{0416}
\pdfglyphtounicode{afii10025}{0417}
\pdfglyphtounicode{afii10026}{0418}
\pdfglyphtounicode{afii10027}{0419}
\pdfglyphtounicode{afii10028}{041A}
\pdfglyphtounicode{afii10029}{041B}
\pdfglyphtounicode{afii10030}{041C}
\pdfglyphtounicode{afii10031}{041D}
\pdfglyphtounicode{afii10032}{041E}
\pdfglyphtounicode{afii10033}{041F}
\pdfglyphtounicode{afii10034}{0420}
\pdfglyphtounicode{afii10035}{0421}
\pdfglyphtounicode{afii10036}{0422}
\pdfglyphtounicode{afii10037}{0423}
\pdfglyphtounicode{afii10038}{0424}
\pdfglyphtounicode{afii10039}{0425}
\pdfglyphtounicode{afii10040}{0426}
\pdfglyphtounicode{afii10041}{0427}
\pdfglyphtounicode{afii10042}{0428}
\pdfglyphtounicode{afii10043}{0429}
\pdfglyphtounicode{afii10044}{042A}
\pdfglyphtounicode{afii10045}{042B}
\pdfglyphtounicode{afii10046}{042C}
\pdfglyphtounicode{afii10047}{042D}
\pdfglyphtounicode{afii10048}{042E}
\pdfglyphtounicode{afii10049}{042F}
\pdfglyphtounicode{afii10050}{0490}
\pdfglyphtounicode{afii10051}{0402}
\pdfglyphtounicode{afii10052}{0403}
\pdfglyphtounicode{afii10053}{0404}
\pdfglyphtounicode{afii10054}{0405}
\pdfglyphtounicode{afii10055}{0406}
\pdfglyphtounicode{afii10056}{0407}
\pdfglyphtounicode{afii10057}{0408}
\pdfglyphtounicode{afii10058}{0409}
\pdfglyphtounicode{afii10059}{040A}
\pdfglyphtounicode{afii10060}{040B}
\pdfglyphtounicode{afii10061}{040C}
\pdfglyphtounicode{afii10062}{040E}
\pdfglyphtounicode{afii10063}{F6C4}
\pdfglyphtounicode{afii10064}{F6C5}
\pdfglyphtounicode{afii10065}{0430}
\pdfglyphtounicode{afii10066}{0431}
\pdfglyphtounicode{afii10067}{0432}
\pdfglyphtounicode{afii10068}{0433}
\pdfglyphtounicode{afii10069}{0434}
\pdfglyphtounicode{afii10070}{0435}
\pdfglyphtounicode{afii10071}{0451}
\pdfglyphtounicode{afii10072}{0436}
\pdfglyphtounicode{afii10073}{0437}
\pdfglyphtounicode{afii10074}{0438}
\pdfglyphtounicode{afii10075}{0439}
\pdfglyphtounicode{afii10076}{043A}
\pdfglyphtounicode{afii10077}{043B}
\pdfglyphtounicode{afii10078}{043C}
\pdfglyphtounicode{afii10079}{043D}
\pdfglyphtounicode{afii10080}{043E}
\pdfglyphtounicode{afii10081}{043F}
\pdfglyphtounicode{afii10082}{0440}
\pdfglyphtounicode{afii10083}{0441}
\pdfglyphtounicode{afii10084}{0442}
\pdfglyphtounicode{afii10085}{0443}
\pdfglyphtounicode{afii10086}{0444}
\pdfglyphtounicode{afii10087}{0445}
\pdfglyphtounicode{afii10088}{0446}
\pdfglyphtounicode{afii10089}{0447}
\pdfglyphtounicode{afii10090}{0448}
\pdfglyphtounicode{afii10091}{0449}
\pdfglyphtounicode{afii10092}{044A}
\pdfglyphtounicode{afii10093}{044B}
\pdfglyphtounicode{afii10094}{044C}
\pdfglyphtounicode{afii10095}{044D}
\pdfglyphtounicode{afii10096}{044E}
\pdfglyphtounicode{afii10097}{044F}
\pdfglyphtounicode{afii10098}{0491}
\pdfglyphtounicode{afii10099}{0452}
\pdfglyphtounicode{afii10100}{0453}
\pdfglyphtounicode{afii10101}{0454}
\pdfglyphtounicode{afii10102}{0455}
\pdfglyphtounicode{afii10103}{0456}
\pdfglyphtounicode{afii10104}{0457}
\pdfglyphtounicode{afii10105}{0458}
\pdfglyphtounicode{afii10106}{0459}
\pdfglyphtounicode{afii10107}{045A}
\pdfglyphtounicode{afii10108}{045B}
\pdfglyphtounicode{afii10109}{045C}
\pdfglyphtounicode{afii10110}{045E}
\pdfglyphtounicode{afii10145}{040F}
\pdfglyphtounicode{afii10146}{0462}
\pdfglyphtounicode{afii10147}{0472}
\pdfglyphtounicode{afii10148}{0474}
\pdfglyphtounicode{afii10192}{F6C6}
\pdfglyphtounicode{afii10193}{045F}
\pdfglyphtounicode{afii10194}{0463}
\pdfglyphtounicode{afii10195}{0473}
\pdfglyphtounicode{afii10196}{0475}
\pdfglyphtounicode{afii10831}{F6C7}
\pdfglyphtounicode{afii10832}{F6C8}
\pdfglyphtounicode{afii10846}{04D9}
\pdfglyphtounicode{afii299}{200E}
\pdfglyphtounicode{afii300}{200F}
\pdfglyphtounicode{afii301}{200D}
\pdfglyphtounicode{afii57381}{066A}
\pdfglyphtounicode{afii57388}{060C}
\pdfglyphtounicode{afii57392}{0660}
\pdfglyphtounicode{afii57393}{0661}
\pdfglyphtounicode{afii57394}{0662}
\pdfglyphtounicode{afii57395}{0663}
\pdfglyphtounicode{afii57396}{0664}
\pdfglyphtounicode{afii57397}{0665}
\pdfglyphtounicode{afii57398}{0666}
\pdfglyphtounicode{afii57399}{0667}
\pdfglyphtounicode{afii57400}{0668}
\pdfglyphtounicode{afii57401}{0669}
\pdfglyphtounicode{afii57403}{061B}
\pdfglyphtounicode{afii57407}{061F}
\pdfglyphtounicode{afii57409}{0621}
\pdfglyphtounicode{afii57410}{0622}
\pdfglyphtounicode{afii57411}{0623}
\pdfglyphtounicode{afii57412}{0624}
\pdfglyphtounicode{afii57413}{0625}
\pdfglyphtounicode{afii57414}{0626}
\pdfglyphtounicode{afii57415}{0627}
\pdfglyphtounicode{afii57416}{0628}
\pdfglyphtounicode{afii57417}{0629}
\pdfglyphtounicode{afii57418}{062A}
\pdfglyphtounicode{afii57419}{062B}
\pdfglyphtounicode{afii57420}{062C}
\pdfglyphtounicode{afii57421}{062D}
\pdfglyphtounicode{afii57422}{062E}
\pdfglyphtounicode{afii57423}{062F}
\pdfglyphtounicode{afii57424}{0630}
\pdfglyphtounicode{afii57425}{0631}
\pdfglyphtounicode{afii57426}{0632}
\pdfglyphtounicode{afii57427}{0633}
\pdfglyphtounicode{afii57428}{0634}
\pdfglyphtounicode{afii57429}{0635}
\pdfglyphtounicode{afii57430}{0636}
\pdfglyphtounicode{afii57431}{0637}
\pdfglyphtounicode{afii57432}{0638}
\pdfglyphtounicode{afii57433}{0639}
\pdfglyphtounicode{afii57434}{063A}
\pdfglyphtounicode{afii57440}{0640}
\pdfglyphtounicode{afii57441}{0641}
\pdfglyphtounicode{afii57442}{0642}
\pdfglyphtounicode{afii57443}{0643}
\pdfglyphtounicode{afii57444}{0644}
\pdfglyphtounicode{afii57445}{0645}
\pdfglyphtounicode{afii57446}{0646}
\pdfglyphtounicode{afii57448}{0648}
\pdfglyphtounicode{afii57449}{0649}
\pdfglyphtounicode{afii57450}{064A}
\pdfglyphtounicode{afii57451}{064B}
\pdfglyphtounicode{afii57452}{064C}
\pdfglyphtounicode{afii57453}{064D}
\pdfglyphtounicode{afii57454}{064E}
\pdfglyphtounicode{afii57455}{064F}
\pdfglyphtounicode{afii57456}{0650}
\pdfglyphtounicode{afii57457}{0651}
\pdfglyphtounicode{afii57458}{0652}
\pdfglyphtounicode{afii57470}{0647}
\pdfglyphtounicode{afii57505}{06A4}
\pdfglyphtounicode{afii57506}{067E}
\pdfglyphtounicode{afii57507}{0686}
\pdfglyphtounicode{afii57508}{0698}
\pdfglyphtounicode{afii57509}{06AF}
\pdfglyphtounicode{afii57511}{0679}
\pdfglyphtounicode{afii57512}{0688}
\pdfglyphtounicode{afii57513}{0691}
\pdfglyphtounicode{afii57514}{06BA}
\pdfglyphtounicode{afii57519}{06D2}
\pdfglyphtounicode{afii57534}{06D5}
\pdfglyphtounicode{afii57636}{20AA}
\pdfglyphtounicode{afii57645}{05BE}
\pdfglyphtounicode{afii57658}{05C3}
\pdfglyphtounicode{afii57664}{05D0}
\pdfglyphtounicode{afii57665}{05D1}
\pdfglyphtounicode{afii57666}{05D2}
\pdfglyphtounicode{afii57667}{05D3}
\pdfglyphtounicode{afii57668}{05D4}
\pdfglyphtounicode{afii57669}{05D5}
\pdfglyphtounicode{afii57670}{05D6}
\pdfglyphtounicode{afii57671}{05D7}
\pdfglyphtounicode{afii57672}{05D8}
\pdfglyphtounicode{afii57673}{05D9}
\pdfglyphtounicode{afii57674}{05DA}
\pdfglyphtounicode{afii57675}{05DB}
\pdfglyphtounicode{afii57676}{05DC}
\pdfglyphtounicode{afii57677}{05DD}
\pdfglyphtounicode{afii57678}{05DE}
\pdfglyphtounicode{afii57679}{05DF}
\pdfglyphtounicode{afii57680}{05E0}
\pdfglyphtounicode{afii57681}{05E1}
\pdfglyphtounicode{afii57682}{05E2}
\pdfglyphtounicode{afii57683}{05E3}
\pdfglyphtounicode{afii57684}{05E4}
\pdfglyphtounicode{afii57685}{05E5}
\pdfglyphtounicode{afii57686}{05E6}
\pdfglyphtounicode{afii57687}{05E7}
\pdfglyphtounicode{afii57688}{05E8}
\pdfglyphtounicode{afii57689}{05E9}
\pdfglyphtounicode{afii57690}{05EA}
\pdfglyphtounicode{afii57694}{FB2A}
\pdfglyphtounicode{afii57695}{FB2B}
\pdfglyphtounicode{afii57700}{FB4B}
\pdfglyphtounicode{afii57705}{FB1F}
\pdfglyphtounicode{afii57716}{05F0}
\pdfglyphtounicode{afii57717}{05F1}
\pdfglyphtounicode{afii57718}{05F2}
\pdfglyphtounicode{afii57723}{FB35}
\pdfglyphtounicode{afii57793}{05B4}
\pdfglyphtounicode{afii57794}{05B5}
\pdfglyphtounicode{afii57795}{05B6}
\pdfglyphtounicode{afii57796}{05BB}
\pdfglyphtounicode{afii57797}{05B8}
\pdfglyphtounicode{afii57798}{05B7}
\pdfglyphtounicode{afii57799}{05B0}
\pdfglyphtounicode{afii57800}{05B2}
\pdfglyphtounicode{afii57801}{05B1}
\pdfglyphtounicode{afii57802}{05B3}
\pdfglyphtounicode{afii57803}{05C2}
\pdfglyphtounicode{afii57804}{05C1}
\pdfglyphtounicode{afii57806}{05B9}
\pdfglyphtounicode{afii57807}{05BC}
\pdfglyphtounicode{afii57839}{05BD}
\pdfglyphtounicode{afii57841}{05BF}
\pdfglyphtounicode{afii57842}{05C0}
\pdfglyphtounicode{afii57929}{02BC}
\pdfglyphtounicode{afii61248}{2105}
\pdfglyphtounicode{afii61289}{2113}
\pdfglyphtounicode{afii61352}{2116}
\pdfglyphtounicode{afii61573}{202C}
\pdfglyphtounicode{afii61574}{202D}
\pdfglyphtounicode{afii61575}{202E}
\pdfglyphtounicode{afii61664}{200C}
\pdfglyphtounicode{afii63167}{066D}
\pdfglyphtounicode{afii64937}{02BD}
\pdfglyphtounicode{agrave}{00E0}
\pdfglyphtounicode{agujarati}{0A85}
\pdfglyphtounicode{agurmukhi}{0A05}
\pdfglyphtounicode{ahiragana}{3042}
\pdfglyphtounicode{ahookabove}{1EA3}
\pdfglyphtounicode{aibengali}{0990}
\pdfglyphtounicode{aibopomofo}{311E}
\pdfglyphtounicode{aideva}{0910}
\pdfglyphtounicode{aiecyrillic}{04D5}
\pdfglyphtounicode{aigujarati}{0A90}
\pdfglyphtounicode{aigurmukhi}{0A10}
\pdfglyphtounicode{aimatragurmukhi}{0A48}
\pdfglyphtounicode{ainarabic}{0639}
\pdfglyphtounicode{ainfinalarabic}{FECA}
\pdfglyphtounicode{aininitialarabic}{FECB}
\pdfglyphtounicode{ainmedialarabic}{FECC}
\pdfglyphtounicode{ainvertedbreve}{0203}
\pdfglyphtounicode{aivowelsignbengali}{09C8}
\pdfglyphtounicode{aivowelsigndeva}{0948}
\pdfglyphtounicode{aivowelsigngujarati}{0AC8}
\pdfglyphtounicode{akatakana}{30A2}
\pdfglyphtounicode{akatakanahalfwidth}{FF71}
\pdfglyphtounicode{akorean}{314F}
\pdfglyphtounicode{alef}{05D0}
\pdfglyphtounicode{alefarabic}{0627}
\pdfglyphtounicode{alefdageshhebrew}{FB30}
\pdfglyphtounicode{aleffinalarabic}{FE8E}
\pdfglyphtounicode{alefhamzaabovearabic}{0623}
\pdfglyphtounicode{alefhamzaabovefinalarabic}{FE84}
\pdfglyphtounicode{alefhamzabelowarabic}{0625}
\pdfglyphtounicode{alefhamzabelowfinalarabic}{FE88}
\pdfglyphtounicode{alefhebrew}{05D0}
\pdfglyphtounicode{aleflamedhebrew}{FB4F}
\pdfglyphtounicode{alefmaddaabovearabic}{0622}
\pdfglyphtounicode{alefmaddaabovefinalarabic}{FE82}
\pdfglyphtounicode{alefmaksuraarabic}{0649}
\pdfglyphtounicode{alefmaksurafinalarabic}{FEF0}
\pdfglyphtounicode{alefmaksurainitialarabic}{FEF3}
\pdfglyphtounicode{alefmaksuramedialarabic}{FEF4}
\pdfglyphtounicode{alefpatahhebrew}{FB2E}
\pdfglyphtounicode{alefqamatshebrew}{FB2F}
\pdfglyphtounicode{aleph}{2135}
\pdfglyphtounicode{allequal}{224C}
\pdfglyphtounicode{alpha}{03B1}
\pdfglyphtounicode{alphatonos}{03AC}
\pdfglyphtounicode{amacron}{0101}
\pdfglyphtounicode{amonospace}{FF41}
\pdfglyphtounicode{ampersand}{0026}
\pdfglyphtounicode{ampersandmonospace}{FF06}
\pdfglyphtounicode{ampersandsmall}{0026}
\pdfglyphtounicode{amsquare}{33C2}
\pdfglyphtounicode{anbopomofo}{3122}
\pdfglyphtounicode{angbopomofo}{3124}
\pdfglyphtounicode{angbracketleft}{27E8}
\pdfglyphtounicode{angbracketright}{27E9}
\pdfglyphtounicode{angkhankhuthai}{0E5A}
\pdfglyphtounicode{angle}{2220}
\pdfglyphtounicode{anglebracketleft}{3008}
\pdfglyphtounicode{anglebracketleftvertical}{FE3F}
\pdfglyphtounicode{anglebracketright}{3009}
\pdfglyphtounicode{anglebracketrightvertical}{FE40}
\pdfglyphtounicode{angleleft}{2329}
\pdfglyphtounicode{angleright}{232A}
\pdfglyphtounicode{angstrom}{212B}
\pdfglyphtounicode{anoteleia}{0387}
\pdfglyphtounicode{anticlockwise}{27F2}
\pdfglyphtounicode{anudattadeva}{0952}
\pdfglyphtounicode{anusvarabengali}{0982}
\pdfglyphtounicode{anusvaradeva}{0902}
\pdfglyphtounicode{anusvaragujarati}{0A82}
\pdfglyphtounicode{aogonek}{0105}
\pdfglyphtounicode{apaatosquare}{3300}
\pdfglyphtounicode{aparen}{249C}
\pdfglyphtounicode{apostrophearmenian}{055A}
\pdfglyphtounicode{apostrophemod}{02BC}
\pdfglyphtounicode{apple}{F8FF}
\pdfglyphtounicode{approaches}{2250}
\pdfglyphtounicode{approxequal}{2248}
\pdfglyphtounicode{approxequalorimage}{2252}
\pdfglyphtounicode{approximatelyequal}{2245}
\pdfglyphtounicode{approxorequal}{224A}
\pdfglyphtounicode{araeaekorean}{318E}
\pdfglyphtounicode{araeakorean}{318D}
\pdfglyphtounicode{arc}{2312}
\pdfglyphtounicode{archleftdown}{21B6}
\pdfglyphtounicode{archrightdown}{21B7}
\pdfglyphtounicode{arighthalfring}{1E9A}
\pdfglyphtounicode{aring}{00E5}
\pdfglyphtounicode{aringacute}{01FB}
\pdfglyphtounicode{aringbelow}{1E01}
\pdfglyphtounicode{arrowboth}{2194}
\pdfglyphtounicode{arrowbothv}{2195}
\pdfglyphtounicode{arrowdashdown}{21E3}
\pdfglyphtounicode{arrowdashleft}{21E0}
\pdfglyphtounicode{arrowdashright}{21E2}
\pdfglyphtounicode{arrowdashup}{21E1}
\pdfglyphtounicode{arrowdblboth}{21D4}
\pdfglyphtounicode{arrowdblbothv}{21D5}
\pdfglyphtounicode{arrowdbldown}{21D3}
\pdfglyphtounicode{arrowdblleft}{21D0}
\pdfglyphtounicode{arrowdblright}{21D2}
\pdfglyphtounicode{arrowdblup}{21D1}
\pdfglyphtounicode{arrowdown}{2193}
\pdfglyphtounicode{arrowdownleft}{2199}
\pdfglyphtounicode{arrowdownright}{2198}
\pdfglyphtounicode{arrowdownwhite}{21E9}
\pdfglyphtounicode{arrowheaddownmod}{02C5}
\pdfglyphtounicode{arrowheadleftmod}{02C2}
\pdfglyphtounicode{arrowheadrightmod}{02C3}
\pdfglyphtounicode{arrowheadupmod}{02C4}
\pdfglyphtounicode{arrowhorizex}{F8E7}
\pdfglyphtounicode{arrowleft}{2190}
\pdfglyphtounicode{arrowleftbothalf}{21BD}
\pdfglyphtounicode{arrowleftdbl}{21D0}
\pdfglyphtounicode{arrowleftdblstroke}{21CD}
\pdfglyphtounicode{arrowleftoverright}{21C6}
\pdfglyphtounicode{arrowlefttophalf}{21BC}
\pdfglyphtounicode{arrowleftwhite}{21E6}
\pdfglyphtounicode{arrownortheast}{2197}
\pdfglyphtounicode{arrownorthwest}{2196}
\pdfglyphtounicode{arrowparrleftright}{21C6}
\pdfglyphtounicode{arrowparrrightleft}{21C4}
\pdfglyphtounicode{arrowright}{2192}
\pdfglyphtounicode{arrowrightbothalf}{21C1}
\pdfglyphtounicode{arrowrightdblstroke}{21CF}
\pdfglyphtounicode{arrowrightheavy}{279E}
\pdfglyphtounicode{arrowrightoverleft}{21C4}
\pdfglyphtounicode{arrowrighttophalf}{21C0}
\pdfglyphtounicode{arrowrightwhite}{21E8}
\pdfglyphtounicode{arrowsoutheast}{2198}
\pdfglyphtounicode{arrowsouthwest}{2199}
\pdfglyphtounicode{arrowtableft}{21E4}
\pdfglyphtounicode{arrowtabright}{21E5}
\pdfglyphtounicode{arrowtailleft}{21A2}
\pdfglyphtounicode{arrowtailright}{21A3}
\pdfglyphtounicode{arrowtripleleft}{21DA}
\pdfglyphtounicode{arrowtripleright}{21DB}
\pdfglyphtounicode{arrowup}{2191}
\pdfglyphtounicode{arrowupdn}{2195}
\pdfglyphtounicode{arrowupdnbse}{21A8}
\pdfglyphtounicode{arrowupdownbase}{21A8}
\pdfglyphtounicode{arrowupleft}{2196}
\pdfglyphtounicode{arrowupleftofdown}{21C5}
\pdfglyphtounicode{arrowupright}{2197}
\pdfglyphtounicode{arrowupwhite}{21E7}
\pdfglyphtounicode{arrowvertex}{F8E6}
\pdfglyphtounicode{asciicircum}{005E}
\pdfglyphtounicode{asciicircummonospace}{FF3E}
\pdfglyphtounicode{asciitilde}{007E}
\pdfglyphtounicode{asciitildemonospace}{FF5E}
\pdfglyphtounicode{ascript}{0251}
\pdfglyphtounicode{ascriptturned}{0252}
\pdfglyphtounicode{asmallhiragana}{3041}
\pdfglyphtounicode{asmallkatakana}{30A1}
\pdfglyphtounicode{asmallkatakanahalfwidth}{FF67}
\pdfglyphtounicode{asterisk}{002A}
\pdfglyphtounicode{asteriskaltonearabic}{066D}
\pdfglyphtounicode{asteriskarabic}{066D}
\pdfglyphtounicode{asteriskcentered}{2217}
\pdfglyphtounicode{asteriskmath}{2217}
\pdfglyphtounicode{asteriskmonospace}{FF0A}
\pdfglyphtounicode{asterisksmall}{FE61}
\pdfglyphtounicode{asterism}{2042}
\pdfglyphtounicode{asuperior}{0061}
\pdfglyphtounicode{asymptoticallyequal}{2243}
\pdfglyphtounicode{at}{0040}
\pdfglyphtounicode{atilde}{00E3}
\pdfglyphtounicode{atmonospace}{FF20}
\pdfglyphtounicode{atsmall}{FE6B}
\pdfglyphtounicode{aturned}{0250}
\pdfglyphtounicode{aubengali}{0994}
\pdfglyphtounicode{aubopomofo}{3120}
\pdfglyphtounicode{audeva}{0914}
\pdfglyphtounicode{augujarati}{0A94}
\pdfglyphtounicode{augurmukhi}{0A14}
\pdfglyphtounicode{aulengthmarkbengali}{09D7}
\pdfglyphtounicode{aumatragurmukhi}{0A4C}
\pdfglyphtounicode{auvowelsignbengali}{09CC}
\pdfglyphtounicode{auvowelsigndeva}{094C}
\pdfglyphtounicode{auvowelsigngujarati}{0ACC}
\pdfglyphtounicode{avagrahadeva}{093D}
\pdfglyphtounicode{aybarmenian}{0561}
\pdfglyphtounicode{ayin}{05E2}
\pdfglyphtounicode{ayinaltonehebrew}{FB20}
\pdfglyphtounicode{ayinhebrew}{05E2}
\pdfglyphtounicode{b}{0062}
\pdfglyphtounicode{babengali}{09AC}
\pdfglyphtounicode{backslash}{005C}
\pdfglyphtounicode{backslashmonospace}{FF3C}
\pdfglyphtounicode{badeva}{092C}
\pdfglyphtounicode{bagujarati}{0AAC}
\pdfglyphtounicode{bagurmukhi}{0A2C}
\pdfglyphtounicode{bahiragana}{3070}
\pdfglyphtounicode{bahtthai}{0E3F}
\pdfglyphtounicode{bakatakana}{30D0}
\pdfglyphtounicode{bar}{007C}
\pdfglyphtounicode{bardbl}{2225}
\pdfglyphtounicode{barmonospace}{FF5C}
\pdfglyphtounicode{bbopomofo}{3105}
\pdfglyphtounicode{bcircle}{24D1}
\pdfglyphtounicode{bdotaccent}{1E03}
\pdfglyphtounicode{bdotbelow}{1E05}
\pdfglyphtounicode{beamedsixteenthnotes}{266C}
\pdfglyphtounicode{because}{2235}
\pdfglyphtounicode{becyrillic}{0431}
\pdfglyphtounicode{beharabic}{0628}
\pdfglyphtounicode{behfinalarabic}{FE90}
\pdfglyphtounicode{behinitialarabic}{FE91}
\pdfglyphtounicode{behiragana}{3079}
\pdfglyphtounicode{behmedialarabic}{FE92}
\pdfglyphtounicode{behmeeminitialarabic}{FC9F}
\pdfglyphtounicode{behmeemisolatedarabic}{FC08}
\pdfglyphtounicode{behnoonfinalarabic}{FC6D}
\pdfglyphtounicode{bekatakana}{30D9}
\pdfglyphtounicode{benarmenian}{0562}
\pdfglyphtounicode{bet}{05D1}
\pdfglyphtounicode{beta}{03B2}
\pdfglyphtounicode{betasymbolgreek}{03D0}
\pdfglyphtounicode{betdagesh}{FB31}
\pdfglyphtounicode{betdageshhebrew}{FB31}
\pdfglyphtounicode{beth}{2136}
\pdfglyphtounicode{bethebrew}{05D1}
\pdfglyphtounicode{betrafehebrew}{FB4C}
\pdfglyphtounicode{between}{226C}
\pdfglyphtounicode{bhabengali}{09AD}
\pdfglyphtounicode{bhadeva}{092D}
\pdfglyphtounicode{bhagujarati}{0AAD}
\pdfglyphtounicode{bhagurmukhi}{0A2D}
\pdfglyphtounicode{bhook}{0253}
\pdfglyphtounicode{bihiragana}{3073}
\pdfglyphtounicode{bikatakana}{30D3}
\pdfglyphtounicode{bilabialclick}{0298}
\pdfglyphtounicode{bindigurmukhi}{0A02}
\pdfglyphtounicode{birusquare}{3331}
\pdfglyphtounicode{blackcircle}{25CF}
\pdfglyphtounicode{blackdiamond}{25C6}
\pdfglyphtounicode{blackdownpointingtriangle}{25BC}
\pdfglyphtounicode{blackleftpointingpointer}{25C4}
\pdfglyphtounicode{blackleftpointingtriangle}{25C0}
\pdfglyphtounicode{blacklenticularbracketleft}{3010}
\pdfglyphtounicode{blacklenticularbracketleftvertical}{FE3B}
\pdfglyphtounicode{blacklenticularbracketright}{3011}
\pdfglyphtounicode{blacklenticularbracketrightvertical}{FE3C}
\pdfglyphtounicode{blacklowerlefttriangle}{25E3}
\pdfglyphtounicode{blacklowerrighttriangle}{25E2}
\pdfglyphtounicode{blackrectangle}{25AC}
\pdfglyphtounicode{blackrightpointingpointer}{25BA}
\pdfglyphtounicode{blackrightpointingtriangle}{25B6}
\pdfglyphtounicode{blacksmallsquare}{25AA}
\pdfglyphtounicode{blacksmilingface}{263B}
\pdfglyphtounicode{blacksquare}{25A0}
\pdfglyphtounicode{blackstar}{2605}
\pdfglyphtounicode{blackupperlefttriangle}{25E4}
\pdfglyphtounicode{blackupperrighttriangle}{25E5}
\pdfglyphtounicode{blackuppointingsmalltriangle}{25B4}
\pdfglyphtounicode{blackuppointingtriangle}{25B2}
\pdfglyphtounicode{blank}{2423}
\pdfglyphtounicode{blinebelow}{1E07}
\pdfglyphtounicode{block}{2588}
\pdfglyphtounicode{bmonospace}{FF42}
\pdfglyphtounicode{bobaimaithai}{0E1A}
\pdfglyphtounicode{bohiragana}{307C}
\pdfglyphtounicode{bokatakana}{30DC}
\pdfglyphtounicode{bparen}{249D}
\pdfglyphtounicode{bqsquare}{33C3}
\pdfglyphtounicode{braceex}{F8F4}
\pdfglyphtounicode{braceleft}{007B}
\pdfglyphtounicode{braceleftbt}{F8F3}
\pdfglyphtounicode{braceleftmid}{F8F2}
\pdfglyphtounicode{braceleftmonospace}{FF5B}
\pdfglyphtounicode{braceleftsmall}{FE5B}
\pdfglyphtounicode{bracelefttp}{F8F1}
\pdfglyphtounicode{braceleftvertical}{FE37}
\pdfglyphtounicode{braceright}{007D}
\pdfglyphtounicode{bracerightbt}{F8FE}
\pdfglyphtounicode{bracerightmid}{F8FD}
\pdfglyphtounicode{bracerightmonospace}{FF5D}
\pdfglyphtounicode{bracerightsmall}{FE5C}
\pdfglyphtounicode{bracerighttp}{F8FC}
\pdfglyphtounicode{bracerightvertical}{FE38}
\pdfglyphtounicode{bracketleft}{005B}
\pdfglyphtounicode{bracketleftbt}{F8F0}
\pdfglyphtounicode{bracketleftex}{F8EF}
\pdfglyphtounicode{bracketleftmonospace}{FF3B}
\pdfglyphtounicode{bracketlefttp}{F8EE}
\pdfglyphtounicode{bracketright}{005D}
\pdfglyphtounicode{bracketrightbt}{F8FB}
\pdfglyphtounicode{bracketrightex}{F8FA}
\pdfglyphtounicode{bracketrightmonospace}{FF3D}
\pdfglyphtounicode{bracketrighttp}{F8F9}
\pdfglyphtounicode{breve}{02D8}
\pdfglyphtounicode{brevebelowcmb}{032E}
\pdfglyphtounicode{brevecmb}{0306}
\pdfglyphtounicode{breveinvertedbelowcmb}{032F}
\pdfglyphtounicode{breveinvertedcmb}{0311}
\pdfglyphtounicode{breveinverteddoublecmb}{0361}
\pdfglyphtounicode{bridgebelowcmb}{032A}
\pdfglyphtounicode{bridgeinvertedbelowcmb}{033A}
\pdfglyphtounicode{brokenbar}{00A6}
\pdfglyphtounicode{bstroke}{0180}
\pdfglyphtounicode{bsuperior}{0062}
\pdfglyphtounicode{btopbar}{0183}
\pdfglyphtounicode{buhiragana}{3076}
\pdfglyphtounicode{bukatakana}{30D6}
\pdfglyphtounicode{bullet}{2022}
\pdfglyphtounicode{bulletinverse}{25D8}
\pdfglyphtounicode{bulletoperator}{2219}
\pdfglyphtounicode{bullseye}{25CE}
\pdfglyphtounicode{c}{0063}
\pdfglyphtounicode{caarmenian}{056E}
\pdfglyphtounicode{cabengali}{099A}
\pdfglyphtounicode{cacute}{0107}
\pdfglyphtounicode{cadeva}{091A}
\pdfglyphtounicode{cagujarati}{0A9A}
\pdfglyphtounicode{cagurmukhi}{0A1A}
\pdfglyphtounicode{calsquare}{3388}
\pdfglyphtounicode{candrabindubengali}{0981}
\pdfglyphtounicode{candrabinducmb}{0310}
\pdfglyphtounicode{candrabindudeva}{0901}
\pdfglyphtounicode{candrabindugujarati}{0A81}
\pdfglyphtounicode{capslock}{21EA}
\pdfglyphtounicode{careof}{2105}
\pdfglyphtounicode{caron}{02C7}
\pdfglyphtounicode{caronbelowcmb}{032C}
\pdfglyphtounicode{caroncmb}{030C}
\pdfglyphtounicode{carriagereturn}{21B5}
\pdfglyphtounicode{cbopomofo}{3118}
\pdfglyphtounicode{ccaron}{010D}
\pdfglyphtounicode{ccedilla}{00E7}
\pdfglyphtounicode{ccedillaacute}{1E09}
\pdfglyphtounicode{ccircle}{24D2}
\pdfglyphtounicode{ccircumflex}{0109}
\pdfglyphtounicode{ccurl}{0255}
\pdfglyphtounicode{cdot}{010B}
\pdfglyphtounicode{cdotaccent}{010B}
\pdfglyphtounicode{cdsquare}{33C5}
\pdfglyphtounicode{cedilla}{00B8}
\pdfglyphtounicode{cedillacmb}{0327}
\pdfglyphtounicode{ceilingleft}{2308}
\pdfglyphtounicode{ceilingright}{2309}
\pdfglyphtounicode{cent}{00A2}
\pdfglyphtounicode{centigrade}{2103}
\pdfglyphtounicode{centinferior}{00A2}
\pdfglyphtounicode{centmonospace}{FFE0}
\pdfglyphtounicode{centoldstyle}{00A2}
\pdfglyphtounicode{centsuperior}{00A2}
\pdfglyphtounicode{chaarmenian}{0579}
\pdfglyphtounicode{chabengali}{099B}
\pdfglyphtounicode{chadeva}{091B}
\pdfglyphtounicode{chagujarati}{0A9B}
\pdfglyphtounicode{chagurmukhi}{0A1B}
\pdfglyphtounicode{chbopomofo}{3114}
\pdfglyphtounicode{cheabkhasiancyrillic}{04BD}
\pdfglyphtounicode{check}{2713}
\pdfglyphtounicode{checkmark}{2713}
\pdfglyphtounicode{checyrillic}{0447}
\pdfglyphtounicode{chedescenderabkhasiancyrillic}{04BF}
\pdfglyphtounicode{chedescendercyrillic}{04B7}
\pdfglyphtounicode{chedieresiscyrillic}{04F5}
\pdfglyphtounicode{cheharmenian}{0573}
\pdfglyphtounicode{chekhakassiancyrillic}{04CC}
\pdfglyphtounicode{cheverticalstrokecyrillic}{04B9}
\pdfglyphtounicode{chi}{03C7}
\pdfglyphtounicode{chieuchacirclekorean}{3277}
\pdfglyphtounicode{chieuchaparenkorean}{3217}
\pdfglyphtounicode{chieuchcirclekorean}{3269}
\pdfglyphtounicode{chieuchkorean}{314A}
\pdfglyphtounicode{chieuchparenkorean}{3209}
\pdfglyphtounicode{chochangthai}{0E0A}
\pdfglyphtounicode{chochanthai}{0E08}
\pdfglyphtounicode{chochingthai}{0E09}
\pdfglyphtounicode{chochoethai}{0E0C}
\pdfglyphtounicode{chook}{0188}
\pdfglyphtounicode{cieucacirclekorean}{3276}
\pdfglyphtounicode{cieucaparenkorean}{3216}
\pdfglyphtounicode{cieuccirclekorean}{3268}
\pdfglyphtounicode{cieuckorean}{3148}
\pdfglyphtounicode{cieucparenkorean}{3208}
\pdfglyphtounicode{cieucuparenkorean}{321C}
\pdfglyphtounicode{circle}{25CB}
\pdfglyphtounicode{circleR}{00AE}
\pdfglyphtounicode{circleS}{24C8}
\pdfglyphtounicode{circleasterisk}{229B}
\pdfglyphtounicode{circlecopyrt}{20DD}
\pdfglyphtounicode{circledivide}{2298}
\pdfglyphtounicode{circledot}{2299}
\pdfglyphtounicode{circleequal}{229C}
\pdfglyphtounicode{circleminus}{2296}
\pdfglyphtounicode{circlemultiply}{2297}
\pdfglyphtounicode{circleot}{2299}
\pdfglyphtounicode{circleplus}{2295}
\pdfglyphtounicode{circlepostalmark}{3036}
\pdfglyphtounicode{circlering}{229A}
\pdfglyphtounicode{circlewithlefthalfblack}{25D0}
\pdfglyphtounicode{circlewithrighthalfblack}{25D1}
\pdfglyphtounicode{circumflex}{02C6}
\pdfglyphtounicode{circumflexbelowcmb}{032D}
\pdfglyphtounicode{circumflexcmb}{0302}
\pdfglyphtounicode{clear}{2327}
\pdfglyphtounicode{clickalveolar}{01C2}
\pdfglyphtounicode{clickdental}{01C0}
\pdfglyphtounicode{clicklateral}{01C1}
\pdfglyphtounicode{clickretroflex}{01C3}
\pdfglyphtounicode{clockwise}{27F3}
\pdfglyphtounicode{club}{2663}
\pdfglyphtounicode{clubsuitblack}{2663}
\pdfglyphtounicode{clubsuitwhite}{2667}
\pdfglyphtounicode{cmcubedsquare}{33A4}
\pdfglyphtounicode{cmonospace}{FF43}
\pdfglyphtounicode{cmsquaredsquare}{33A0}
\pdfglyphtounicode{coarmenian}{0581}
\pdfglyphtounicode{colon}{003A}
\pdfglyphtounicode{colonmonetary}{20A1}
\pdfglyphtounicode{colonmonospace}{FF1A}
\pdfglyphtounicode{colonsign}{20A1}
\pdfglyphtounicode{colonsmall}{FE55}
\pdfglyphtounicode{colontriangularhalfmod}{02D1}
\pdfglyphtounicode{colontriangularmod}{02D0}
\pdfglyphtounicode{comma}{002C}
\pdfglyphtounicode{commaabovecmb}{0313}
\pdfglyphtounicode{commaaboverightcmb}{0315}
\pdfglyphtounicode{commaaccent}{F6C3}
\pdfglyphtounicode{commaarabic}{060C}
\pdfglyphtounicode{commaarmenian}{055D}
\pdfglyphtounicode{commainferior}{002C}
\pdfglyphtounicode{commamonospace}{FF0C}
\pdfglyphtounicode{commareversedabovecmb}{0314}
\pdfglyphtounicode{commareversedmod}{02BD}
\pdfglyphtounicode{commasmall}{FE50}
\pdfglyphtounicode{commasuperior}{002C}
\pdfglyphtounicode{commaturnedabovecmb}{0312}
\pdfglyphtounicode{commaturnedmod}{02BB}
\pdfglyphtounicode{compass}{263C}
\pdfglyphtounicode{complement}{2201}
\pdfglyphtounicode{compwordmark}{200C}
\pdfglyphtounicode{congruent}{2245}
\pdfglyphtounicode{contourintegral}{222E}
\pdfglyphtounicode{control}{2303}
\pdfglyphtounicode{controlACK}{0006}
\pdfglyphtounicode{controlBEL}{0007}
\pdfglyphtounicode{controlBS}{0008}
\pdfglyphtounicode{controlCAN}{0018}
\pdfglyphtounicode{controlCR}{000D}
\pdfglyphtounicode{controlDC1}{0011}
\pdfglyphtounicode{controlDC2}{0012}
\pdfglyphtounicode{controlDC3}{0013}
\pdfglyphtounicode{controlDC4}{0014}
\pdfglyphtounicode{controlDEL}{007F}
\pdfglyphtounicode{controlDLE}{0010}
\pdfglyphtounicode{controlEM}{0019}
\pdfglyphtounicode{controlENQ}{0005}
\pdfglyphtounicode{controlEOT}{0004}
\pdfglyphtounicode{controlESC}{001B}
\pdfglyphtounicode{controlETB}{0017}
\pdfglyphtounicode{controlETX}{0003}
\pdfglyphtounicode{controlFF}{000C}
\pdfglyphtounicode{controlFS}{001C}
\pdfglyphtounicode{controlGS}{001D}
\pdfglyphtounicode{controlHT}{0009}
\pdfglyphtounicode{controlLF}{000A}
\pdfglyphtounicode{controlNAK}{0015}
\pdfglyphtounicode{controlRS}{001E}
\pdfglyphtounicode{controlSI}{000F}
\pdfglyphtounicode{controlSO}{000E}
\pdfglyphtounicode{controlSOT}{0002}
\pdfglyphtounicode{controlSTX}{0001}
\pdfglyphtounicode{controlSUB}{001A}
\pdfglyphtounicode{controlSYN}{0016}
\pdfglyphtounicode{controlUS}{001F}
\pdfglyphtounicode{controlVT}{000B}
\pdfglyphtounicode{coproduct}{2A3F}
\pdfglyphtounicode{copyright}{00A9}
\pdfglyphtounicode{copyrightsans}{00A9}
\pdfglyphtounicode{copyrightserif}{00A9}
\pdfglyphtounicode{cornerbracketleft}{300C}
\pdfglyphtounicode{cornerbracketlefthalfwidth}{FF62}
\pdfglyphtounicode{cornerbracketleftvertical}{FE41}
\pdfglyphtounicode{cornerbracketright}{300D}
\pdfglyphtounicode{cornerbracketrighthalfwidth}{FF63}
\pdfglyphtounicode{cornerbracketrightvertical}{FE42}
\pdfglyphtounicode{corporationsquare}{337F}
\pdfglyphtounicode{cosquare}{33C7}
\pdfglyphtounicode{coverkgsquare}{33C6}
\pdfglyphtounicode{cparen}{249E}
\pdfglyphtounicode{cruzeiro}{20A2}
\pdfglyphtounicode{cstretched}{0297}
\pdfglyphtounicode{ct}{0063 0074}
\pdfglyphtounicode{curlyand}{22CF}
\pdfglyphtounicode{curlyleft}{21AB}
\pdfglyphtounicode{curlyor}{22CE}
\pdfglyphtounicode{curlyright}{21AC}
\pdfglyphtounicode{currency}{00A4}
\pdfglyphtounicode{cwm}{200C}
\pdfglyphtounicode{cyrBreve}{02D8}
\pdfglyphtounicode{cyrFlex}{00A0 0311}
\pdfglyphtounicode{cyrbreve}{02D8}
\pdfglyphtounicode{cyrflex}{00A0 0311}
\pdfglyphtounicode{d}{0064}
\pdfglyphtounicode{daarmenian}{0564}
\pdfglyphtounicode{dabengali}{09A6}
\pdfglyphtounicode{dadarabic}{0636}
\pdfglyphtounicode{dadeva}{0926}
\pdfglyphtounicode{dadfinalarabic}{FEBE}
\pdfglyphtounicode{dadinitialarabic}{FEBF}
\pdfglyphtounicode{dadmedialarabic}{FEC0}
\pdfglyphtounicode{dagesh}{05BC}
\pdfglyphtounicode{dageshhebrew}{05BC}
\pdfglyphtounicode{dagger}{2020}
\pdfglyphtounicode{daggerdbl}{2021}
\pdfglyphtounicode{dagujarati}{0AA6}
\pdfglyphtounicode{dagurmukhi}{0A26}
\pdfglyphtounicode{dahiragana}{3060}
\pdfglyphtounicode{dakatakana}{30C0}
\pdfglyphtounicode{dalarabic}{062F}
\pdfglyphtounicode{dalet}{05D3}
\pdfglyphtounicode{daletdagesh}{FB33}
\pdfglyphtounicode{daletdageshhebrew}{FB33}
\pdfglyphtounicode{daleth}{2138}
\pdfglyphtounicode{dalethatafpatah}{05D3 05B2}
\pdfglyphtounicode{dalethatafpatahhebrew}{05D3 05B2}
\pdfglyphtounicode{dalethatafsegol}{05D3 05B1}
\pdfglyphtounicode{dalethatafsegolhebrew}{05D3 05B1}
\pdfglyphtounicode{dalethebrew}{05D3}
\pdfglyphtounicode{dalethiriq}{05D3 05B4}
\pdfglyphtounicode{dalethiriqhebrew}{05D3 05B4}
\pdfglyphtounicode{daletholam}{05D3 05B9}
\pdfglyphtounicode{daletholamhebrew}{05D3 05B9}
\pdfglyphtounicode{daletpatah}{05D3 05B7}
\pdfglyphtounicode{daletpatahhebrew}{05D3 05B7}
\pdfglyphtounicode{daletqamats}{05D3 05B8}
\pdfglyphtounicode{daletqamatshebrew}{05D3 05B8}
\pdfglyphtounicode{daletqubuts}{05D3 05BB}
\pdfglyphtounicode{daletqubutshebrew}{05D3 05BB}
\pdfglyphtounicode{daletsegol}{05D3 05B6}
\pdfglyphtounicode{daletsegolhebrew}{05D3 05B6}
\pdfglyphtounicode{daletsheva}{05D3 05B0}
\pdfglyphtounicode{daletshevahebrew}{05D3 05B0}
\pdfglyphtounicode{dalettsere}{05D3 05B5}
\pdfglyphtounicode{dalettserehebrew}{05D3 05B5}
\pdfglyphtounicode{dalfinalarabic}{FEAA}
\pdfglyphtounicode{dammaarabic}{064F}
\pdfglyphtounicode{dammalowarabic}{064F}
\pdfglyphtounicode{dammatanaltonearabic}{064C}
\pdfglyphtounicode{dammatanarabic}{064C}
\pdfglyphtounicode{danda}{0964}
\pdfglyphtounicode{dargahebrew}{05A7}
\pdfglyphtounicode{dargalefthebrew}{05A7}
\pdfglyphtounicode{dasiapneumatacyrilliccmb}{0485}
\pdfglyphtounicode{dbar}{0111}
\pdfglyphtounicode{dblGrave}{00A0 030F}
\pdfglyphtounicode{dblanglebracketleft}{300A}
\pdfglyphtounicode{dblanglebracketleftvertical}{FE3D}
\pdfglyphtounicode{dblanglebracketright}{300B}
\pdfglyphtounicode{dblanglebracketrightvertical}{FE3E}
\pdfglyphtounicode{dblarchinvertedbelowcmb}{032B}
\pdfglyphtounicode{dblarrowdwn}{21CA}
\pdfglyphtounicode{dblarrowheadleft}{219E}
\pdfglyphtounicode{dblarrowheadright}{21A0}
\pdfglyphtounicode{dblarrowleft}{21D4}
\pdfglyphtounicode{dblarrowright}{21D2}
\pdfglyphtounicode{dblarrowup}{21C8}
\pdfglyphtounicode{dblbracketleft}{27E6}
\pdfglyphtounicode{dblbracketright}{27E7}
\pdfglyphtounicode{dbldanda}{0965}
\pdfglyphtounicode{dblgrave}{00A0 030F}
\pdfglyphtounicode{dblgravecmb}{030F}
\pdfglyphtounicode{dblintegral}{222C}
\pdfglyphtounicode{dbllowline}{2017}
\pdfglyphtounicode{dbllowlinecmb}{0333}
\pdfglyphtounicode{dbloverlinecmb}{033F}
\pdfglyphtounicode{dblprimemod}{02BA}
\pdfglyphtounicode{dblverticalbar}{2016}
\pdfglyphtounicode{dblverticallineabovecmb}{030E}
\pdfglyphtounicode{dbopomofo}{3109}
\pdfglyphtounicode{dbsquare}{33C8}
\pdfglyphtounicode{dcaron}{010F}
\pdfglyphtounicode{dcedilla}{1E11}
\pdfglyphtounicode{dcircle}{24D3}
\pdfglyphtounicode{dcircumflexbelow}{1E13}
\pdfglyphtounicode{dcroat}{0111}
\pdfglyphtounicode{ddabengali}{09A1}
\pdfglyphtounicode{ddadeva}{0921}
\pdfglyphtounicode{ddagujarati}{0AA1}
\pdfglyphtounicode{ddagurmukhi}{0A21}
\pdfglyphtounicode{ddalarabic}{0688}
\pdfglyphtounicode{ddalfinalarabic}{FB89}
\pdfglyphtounicode{dddhadeva}{095C}
\pdfglyphtounicode{ddhabengali}{09A2}
\pdfglyphtounicode{ddhadeva}{0922}
\pdfglyphtounicode{ddhagujarati}{0AA2}
\pdfglyphtounicode{ddhagurmukhi}{0A22}
\pdfglyphtounicode{ddotaccent}{1E0B}
\pdfglyphtounicode{ddotbelow}{1E0D}
\pdfglyphtounicode{decimalseparatorarabic}{066B}
\pdfglyphtounicode{decimalseparatorpersian}{066B}
\pdfglyphtounicode{decyrillic}{0434}
\pdfglyphtounicode{defines}{225C}
\pdfglyphtounicode{degree}{00B0}
\pdfglyphtounicode{dehihebrew}{05AD}
\pdfglyphtounicode{dehiragana}{3067}
\pdfglyphtounicode{deicoptic}{03EF}
\pdfglyphtounicode{dekatakana}{30C7}
\pdfglyphtounicode{deleteleft}{232B}
\pdfglyphtounicode{deleteright}{2326}
\pdfglyphtounicode{delta}{03B4}
\pdfglyphtounicode{deltaturned}{018D}
\pdfglyphtounicode{denominatorminusonenumeratorbengali}{09F8}
\pdfglyphtounicode{dezh}{02A4}
\pdfglyphtounicode{dhabengali}{09A7}
\pdfglyphtounicode{dhadeva}{0927}
\pdfglyphtounicode{dhagujarati}{0AA7}
\pdfglyphtounicode{dhagurmukhi}{0A27}
\pdfglyphtounicode{dhook}{0257}
\pdfglyphtounicode{dialytikatonos}{0385}
\pdfglyphtounicode{dialytikatonoscmb}{0344}
\pdfglyphtounicode{diamond}{2666}
\pdfglyphtounicode{diamondmath}{22C4}
\pdfglyphtounicode{diamondsolid}{2666}
\pdfglyphtounicode{diamondsuitwhite}{2662}
\pdfglyphtounicode{dieresis}{00A8}
\pdfglyphtounicode{dieresisacute}{00A0 0308 0301}
\pdfglyphtounicode{dieresisbelowcmb}{0324}
\pdfglyphtounicode{dieresiscmb}{0308}
\pdfglyphtounicode{dieresisgrave}{00A0 0308 0300}
\pdfglyphtounicode{dieresistonos}{0385}
\pdfglyphtounicode{difference}{224F}
\pdfglyphtounicode{dihiragana}{3062}
\pdfglyphtounicode{dikatakana}{30C2}
\pdfglyphtounicode{dittomark}{3003}
\pdfglyphtounicode{divide}{00F7}
\pdfglyphtounicode{dividemultiply}{22C7}
\pdfglyphtounicode{divides}{2223}
\pdfglyphtounicode{divisionslash}{2215}
\pdfglyphtounicode{djecyrillic}{0452}
\pdfglyphtounicode{dkshade}{2593}
\pdfglyphtounicode{dlinebelow}{1E0F}
\pdfglyphtounicode{dlsquare}{3397}
\pdfglyphtounicode{dmacron}{0111}
\pdfglyphtounicode{dmonospace}{FF44}
\pdfglyphtounicode{dnblock}{2584}
\pdfglyphtounicode{dochadathai}{0E0E}
\pdfglyphtounicode{dodekthai}{0E14}
\pdfglyphtounicode{dohiragana}{3069}
\pdfglyphtounicode{dokatakana}{30C9}
\pdfglyphtounicode{dollar}{0024}
\pdfglyphtounicode{dollarinferior}{0024}
\pdfglyphtounicode{dollarmonospace}{FF04}
\pdfglyphtounicode{dollaroldstyle}{0024}
\pdfglyphtounicode{dollarsmall}{FE69}
\pdfglyphtounicode{dollarsuperior}{0024}
\pdfglyphtounicode{dong}{20AB}
\pdfglyphtounicode{dorusquare}{3326}
\pdfglyphtounicode{dotaccent}{02D9}
\pdfglyphtounicode{dotaccentcmb}{0307}
\pdfglyphtounicode{dotbelowcmb}{0323}
\pdfglyphtounicode{dotbelowcomb}{0323}
\pdfglyphtounicode{dotkatakana}{30FB}
\pdfglyphtounicode{dotlessi}{0131}
\pdfglyphtounicode{dotlessj}{0237}
\pdfglyphtounicode{dotlessjstrokehook}{0284}
\pdfglyphtounicode{dotmath}{22C5}
\pdfglyphtounicode{dotplus}{2214}
\pdfglyphtounicode{dottedcircle}{25CC}
\pdfglyphtounicode{doubleyodpatah}{FB1F}
\pdfglyphtounicode{doubleyodpatahhebrew}{FB1F}
\pdfglyphtounicode{downfall}{22CE}
\pdfglyphtounicode{downslope}{29F9}
\pdfglyphtounicode{downtackbelowcmb}{031E}
\pdfglyphtounicode{downtackmod}{02D5}
\pdfglyphtounicode{dparen}{249F}
\pdfglyphtounicode{dsuperior}{0064}
\pdfglyphtounicode{dtail}{0256}
\pdfglyphtounicode{dtopbar}{018C}
\pdfglyphtounicode{duhiragana}{3065}
\pdfglyphtounicode{dukatakana}{30C5}
\pdfglyphtounicode{dz}{01F3}
\pdfglyphtounicode{dzaltone}{02A3}
\pdfglyphtounicode{dzcaron}{01C6}
\pdfglyphtounicode{dzcurl}{02A5}
\pdfglyphtounicode{dzeabkhasiancyrillic}{04E1}
\pdfglyphtounicode{dzecyrillic}{0455}
\pdfglyphtounicode{dzhecyrillic}{045F}
\pdfglyphtounicode{e}{0065}
\pdfglyphtounicode{eacute}{00E9}
\pdfglyphtounicode{earth}{2641}
\pdfglyphtounicode{ebengali}{098F}
\pdfglyphtounicode{ebopomofo}{311C}
\pdfglyphtounicode{ebreve}{0115}
\pdfglyphtounicode{ecandradeva}{090D}
\pdfglyphtounicode{ecandragujarati}{0A8D}
\pdfglyphtounicode{ecandravowelsigndeva}{0945}
\pdfglyphtounicode{ecandravowelsigngujarati}{0AC5}
\pdfglyphtounicode{ecaron}{011B}
\pdfglyphtounicode{ecedillabreve}{1E1D}
\pdfglyphtounicode{echarmenian}{0565}
\pdfglyphtounicode{echyiwnarmenian}{0587}
\pdfglyphtounicode{ecircle}{24D4}
\pdfglyphtounicode{ecircumflex}{00EA}
\pdfglyphtounicode{ecircumflexacute}{1EBF}
\pdfglyphtounicode{ecircumflexbelow}{1E19}
\pdfglyphtounicode{ecircumflexdotbelow}{1EC7}
\pdfglyphtounicode{ecircumflexgrave}{1EC1}
\pdfglyphtounicode{ecircumflexhookabove}{1EC3}
\pdfglyphtounicode{ecircumflextilde}{1EC5}
\pdfglyphtounicode{ecyrillic}{0454}
\pdfglyphtounicode{edblgrave}{0205}
\pdfglyphtounicode{edeva}{090F}
\pdfglyphtounicode{edieresis}{00EB}
\pdfglyphtounicode{edot}{0117}
\pdfglyphtounicode{edotaccent}{0117}
\pdfglyphtounicode{edotbelow}{1EB9}
\pdfglyphtounicode{eegurmukhi}{0A0F}
\pdfglyphtounicode{eematragurmukhi}{0A47}
\pdfglyphtounicode{efcyrillic}{0444}
\pdfglyphtounicode{egrave}{00E8}
\pdfglyphtounicode{egujarati}{0A8F}
\pdfglyphtounicode{eharmenian}{0567}
\pdfglyphtounicode{ehbopomofo}{311D}
\pdfglyphtounicode{ehiragana}{3048}
\pdfglyphtounicode{ehookabove}{1EBB}
\pdfglyphtounicode{eibopomofo}{311F}
\pdfglyphtounicode{eight}{0038}
\pdfglyphtounicode{eightarabic}{0668}
\pdfglyphtounicode{eightbengali}{09EE}
\pdfglyphtounicode{eightcircle}{2467}
\pdfglyphtounicode{eightcircleinversesansserif}{2791}
\pdfglyphtounicode{eightdeva}{096E}
\pdfglyphtounicode{eighteencircle}{2471}
\pdfglyphtounicode{eighteenparen}{2485}
\pdfglyphtounicode{eighteenperiod}{2499}
\pdfglyphtounicode{eightgujarati}{0AEE}
\pdfglyphtounicode{eightgurmukhi}{0A6E}
\pdfglyphtounicode{eighthackarabic}{0668}
\pdfglyphtounicode{eighthangzhou}{3028}
\pdfglyphtounicode{eighthnotebeamed}{266B}
\pdfglyphtounicode{eightideographicparen}{3227}
\pdfglyphtounicode{eightinferior}{2088}
\pdfglyphtounicode{eightmonospace}{FF18}
\pdfglyphtounicode{eightoldstyle}{0038}
\pdfglyphtounicode{eightparen}{247B}
\pdfglyphtounicode{eightperiod}{248F}
\pdfglyphtounicode{eightpersian}{06F8}
\pdfglyphtounicode{eightroman}{2177}
\pdfglyphtounicode{eightsuperior}{2078}
\pdfglyphtounicode{eightthai}{0E58}
\pdfglyphtounicode{einvertedbreve}{0207}
\pdfglyphtounicode{eiotifiedcyrillic}{0465}
\pdfglyphtounicode{ekatakana}{30A8}
\pdfglyphtounicode{ekatakanahalfwidth}{FF74}
\pdfglyphtounicode{ekonkargurmukhi}{0A74}
\pdfglyphtounicode{ekorean}{3154}
\pdfglyphtounicode{elcyrillic}{043B}
\pdfglyphtounicode{element}{2208}
\pdfglyphtounicode{elevencircle}{246A}
\pdfglyphtounicode{elevenparen}{247E}
\pdfglyphtounicode{elevenperiod}{2492}
\pdfglyphtounicode{elevenroman}{217A}
\pdfglyphtounicode{ellipsis}{2026}
\pdfglyphtounicode{ellipsisvertical}{22EE}
\pdfglyphtounicode{emacron}{0113}
\pdfglyphtounicode{emacronacute}{1E17}
\pdfglyphtounicode{emacrongrave}{1E15}
\pdfglyphtounicode{emcyrillic}{043C}
\pdfglyphtounicode{emdash}{2014}
\pdfglyphtounicode{emdashvertical}{FE31}
\pdfglyphtounicode{emonospace}{FF45}
\pdfglyphtounicode{emphasismarkarmenian}{055B}
\pdfglyphtounicode{emptyset}{2205}
\pdfglyphtounicode{enbopomofo}{3123}
\pdfglyphtounicode{encyrillic}{043D}
\pdfglyphtounicode{endash}{2013}
\pdfglyphtounicode{endashvertical}{FE32}
\pdfglyphtounicode{endescendercyrillic}{04A3}
\pdfglyphtounicode{eng}{014B}
\pdfglyphtounicode{engbopomofo}{3125}
\pdfglyphtounicode{enghecyrillic}{04A5}
\pdfglyphtounicode{enhookcyrillic}{04C8}
\pdfglyphtounicode{enspace}{2002}
\pdfglyphtounicode{eogonek}{0119}
\pdfglyphtounicode{eokorean}{3153}
\pdfglyphtounicode{eopen}{025B}
\pdfglyphtounicode{eopenclosed}{029A}
\pdfglyphtounicode{eopenreversed}{025C}
\pdfglyphtounicode{eopenreversedclosed}{025E}
\pdfglyphtounicode{eopenreversedhook}{025D}
\pdfglyphtounicode{eparen}{24A0}
\pdfglyphtounicode{epsilon}{03B5}
\pdfglyphtounicode{epsilon1}{03F5}
\pdfglyphtounicode{epsiloninv}{03F6}
\pdfglyphtounicode{epsilontonos}{03AD}
\pdfglyphtounicode{equal}{003D}
\pdfglyphtounicode{equaldotleftright}{2252}
\pdfglyphtounicode{equaldotrightleft}{2253}
\pdfglyphtounicode{equalmonospace}{FF1D}
\pdfglyphtounicode{equalorfollows}{22DF}
\pdfglyphtounicode{equalorgreater}{2A96}
\pdfglyphtounicode{equalorless}{2A95}
\pdfglyphtounicode{equalorprecedes}{22DE}
\pdfglyphtounicode{equalorsimilar}{2242}
\pdfglyphtounicode{equalsdots}{2251}
\pdfglyphtounicode{equalsmall}{FE66}
\pdfglyphtounicode{equalsuperior}{207C}
\pdfglyphtounicode{equivalence}{2261}
\pdfglyphtounicode{equivasymptotic}{224D}
\pdfglyphtounicode{erbopomofo}{3126}
\pdfglyphtounicode{ercyrillic}{0440}
\pdfglyphtounicode{ereversed}{0258}
\pdfglyphtounicode{ereversedcyrillic}{044D}
\pdfglyphtounicode{escyrillic}{0441}
\pdfglyphtounicode{esdescendercyrillic}{04AB}
\pdfglyphtounicode{esh}{0283}
\pdfglyphtounicode{eshcurl}{0286}
\pdfglyphtounicode{eshortdeva}{090E}
\pdfglyphtounicode{eshortvowelsigndeva}{0946}
\pdfglyphtounicode{eshreversedloop}{01AA}
\pdfglyphtounicode{eshsquatreversed}{0285}
\pdfglyphtounicode{esmallhiragana}{3047}
\pdfglyphtounicode{esmallkatakana}{30A7}
\pdfglyphtounicode{esmallkatakanahalfwidth}{FF6A}
\pdfglyphtounicode{estimated}{212E}
\pdfglyphtounicode{esuperior}{0065}
\pdfglyphtounicode{eta}{03B7}
\pdfglyphtounicode{etarmenian}{0568}
\pdfglyphtounicode{etatonos}{03AE}
\pdfglyphtounicode{eth}{00F0}
\pdfglyphtounicode{etilde}{1EBD}
\pdfglyphtounicode{etildebelow}{1E1B}
\pdfglyphtounicode{etnahtafoukhhebrew}{0591}
\pdfglyphtounicode{etnahtafoukhlefthebrew}{0591}
\pdfglyphtounicode{etnahtahebrew}{0591}
\pdfglyphtounicode{etnahtalefthebrew}{0591}
\pdfglyphtounicode{eturned}{01DD}
\pdfglyphtounicode{eukorean}{3161}
\pdfglyphtounicode{euro}{20AC}
\pdfglyphtounicode{evowelsignbengali}{09C7}
\pdfglyphtounicode{evowelsigndeva}{0947}
\pdfglyphtounicode{evowelsigngujarati}{0AC7}
\pdfglyphtounicode{exclam}{0021}
\pdfglyphtounicode{exclamarmenian}{055C}
\pdfglyphtounicode{exclamdbl}{203C}
\pdfglyphtounicode{exclamdown}{00A1}
\pdfglyphtounicode{exclamdownsmall}{00A1}
\pdfglyphtounicode{exclammonospace}{FF01}
\pdfglyphtounicode{exclamsmall}{0021}
\pdfglyphtounicode{existential}{2203}
\pdfglyphtounicode{ezh}{0292}
\pdfglyphtounicode{ezhcaron}{01EF}
\pdfglyphtounicode{ezhcurl}{0293}
\pdfglyphtounicode{ezhreversed}{01B9}
\pdfglyphtounicode{ezhtail}{01BA}
\pdfglyphtounicode{f}{0066}
\pdfglyphtounicode{fadeva}{095E}
\pdfglyphtounicode{fagurmukhi}{0A5E}
\pdfglyphtounicode{fahrenheit}{2109}
\pdfglyphtounicode{fathaarabic}{064E}
\pdfglyphtounicode{fathalowarabic}{064E}
\pdfglyphtounicode{fathatanarabic}{064B}
\pdfglyphtounicode{fbopomofo}{3108}
\pdfglyphtounicode{fcircle}{24D5}
\pdfglyphtounicode{fdotaccent}{1E1F}
\pdfglyphtounicode{feharabic}{0641}
\pdfglyphtounicode{feharmenian}{0586}
\pdfglyphtounicode{fehfinalarabic}{FED2}
\pdfglyphtounicode{fehinitialarabic}{FED3}
\pdfglyphtounicode{fehmedialarabic}{FED4}
\pdfglyphtounicode{feicoptic}{03E5}
\pdfglyphtounicode{female}{2640}
\pdfglyphtounicode{ff}{0066 0066}
\pdfglyphtounicode{ffi}{0066 0066 0069}
\pdfglyphtounicode{ffl}{0066 0066 006C}
\pdfglyphtounicode{fi}{0066 0069}
\pdfglyphtounicode{fifteencircle}{246E}
\pdfglyphtounicode{fifteenparen}{2482}
\pdfglyphtounicode{fifteenperiod}{2496}
\pdfglyphtounicode{figuredash}{2012}
\pdfglyphtounicode{filledbox}{25A0}
\pdfglyphtounicode{filledrect}{25AC}
\pdfglyphtounicode{finalkaf}{05DA}
\pdfglyphtounicode{finalkafdagesh}{FB3A}
\pdfglyphtounicode{finalkafdageshhebrew}{FB3A}
\pdfglyphtounicode{finalkafhebrew}{05DA}
\pdfglyphtounicode{finalkafqamats}{05DA 05B8}
\pdfglyphtounicode{finalkafqamatshebrew}{05DA 05B8}
\pdfglyphtounicode{finalkafsheva}{05DA 05B0}
\pdfglyphtounicode{finalkafshevahebrew}{05DA 05B0}
\pdfglyphtounicode{finalmem}{05DD}
\pdfglyphtounicode{finalmemhebrew}{05DD}
\pdfglyphtounicode{finalnun}{05DF}
\pdfglyphtounicode{finalnunhebrew}{05DF}
\pdfglyphtounicode{finalpe}{05E3}
\pdfglyphtounicode{finalpehebrew}{05E3}
\pdfglyphtounicode{finaltsadi}{05E5}
\pdfglyphtounicode{finaltsadihebrew}{05E5}
\pdfglyphtounicode{firsttonechinese}{02C9}
\pdfglyphtounicode{fisheye}{25C9}
\pdfglyphtounicode{fitacyrillic}{0473}
\pdfglyphtounicode{five}{0035}
\pdfglyphtounicode{fivearabic}{0665}
\pdfglyphtounicode{fivebengali}{09EB}
\pdfglyphtounicode{fivecircle}{2464}
\pdfglyphtounicode{fivecircleinversesansserif}{278E}
\pdfglyphtounicode{fivedeva}{096B}
\pdfglyphtounicode{fiveeighths}{215D}
\pdfglyphtounicode{fivegujarati}{0AEB}
\pdfglyphtounicode{fivegurmukhi}{0A6B}
\pdfglyphtounicode{fivehackarabic}{0665}
\pdfglyphtounicode{fivehangzhou}{3025}
\pdfglyphtounicode{fiveideographicparen}{3224}
\pdfglyphtounicode{fiveinferior}{2085}
\pdfglyphtounicode{fivemonospace}{FF15}
\pdfglyphtounicode{fiveoldstyle}{0035}
\pdfglyphtounicode{fiveparen}{2478}
\pdfglyphtounicode{fiveperiod}{248C}
\pdfglyphtounicode{fivepersian}{06F5}
\pdfglyphtounicode{fiveroman}{2174}
\pdfglyphtounicode{fivesuperior}{2075}
\pdfglyphtounicode{fivethai}{0E55}
\pdfglyphtounicode{fl}{0066 006C}
\pdfglyphtounicode{flat}{266D}
\pdfglyphtounicode{floorleft}{230A}
\pdfglyphtounicode{floorright}{230B}
\pdfglyphtounicode{florin}{0192}
\pdfglyphtounicode{fmonospace}{FF46}
\pdfglyphtounicode{fmsquare}{3399}
\pdfglyphtounicode{fofanthai}{0E1F}
\pdfglyphtounicode{fofathai}{0E1D}
\pdfglyphtounicode{follownotdbleqv}{2ABA}
\pdfglyphtounicode{follownotslnteql}{2AB6}
\pdfglyphtounicode{followornoteqvlnt}{22E9}
\pdfglyphtounicode{follows}{227B}
\pdfglyphtounicode{followsequal}{2AB0}
\pdfglyphtounicode{followsorcurly}{227D}
\pdfglyphtounicode{followsorequal}{227F}
\pdfglyphtounicode{fongmanthai}{0E4F}
\pdfglyphtounicode{forall}{2200}
\pdfglyphtounicode{forces}{22A9}
\pdfglyphtounicode{forcesbar}{22AA}
\pdfglyphtounicode{fork}{22D4}
\pdfglyphtounicode{four}{0034}
\pdfglyphtounicode{fourarabic}{0664}
\pdfglyphtounicode{fourbengali}{09EA}
\pdfglyphtounicode{fourcircle}{2463}
\pdfglyphtounicode{fourcircleinversesansserif}{278D}
\pdfglyphtounicode{fourdeva}{096A}
\pdfglyphtounicode{fourgujarati}{0AEA}
\pdfglyphtounicode{fourgurmukhi}{0A6A}
\pdfglyphtounicode{fourhackarabic}{0664}
\pdfglyphtounicode{fourhangzhou}{3024}
\pdfglyphtounicode{fourideographicparen}{3223}
\pdfglyphtounicode{fourinferior}{2084}
\pdfglyphtounicode{fourmonospace}{FF14}
\pdfglyphtounicode{fournumeratorbengali}{09F7}
\pdfglyphtounicode{fouroldstyle}{0034}
\pdfglyphtounicode{fourparen}{2477}
\pdfglyphtounicode{fourperiod}{248B}
\pdfglyphtounicode{fourpersian}{06F4}
\pdfglyphtounicode{fourroman}{2173}
\pdfglyphtounicode{foursuperior}{2074}
\pdfglyphtounicode{fourteencircle}{246D}
\pdfglyphtounicode{fourteenparen}{2481}
\pdfglyphtounicode{fourteenperiod}{2495}
\pdfglyphtounicode{fourthai}{0E54}
\pdfglyphtounicode{fourthtonechinese}{02CB}
\pdfglyphtounicode{fparen}{24A1}
\pdfglyphtounicode{fraction}{2044}
\pdfglyphtounicode{franc}{20A3}
\pdfglyphtounicode{frown}{2322}
\pdfglyphtounicode{g}{0067}
\pdfglyphtounicode{gabengali}{0997}
\pdfglyphtounicode{gacute}{01F5}
\pdfglyphtounicode{gadeva}{0917}
\pdfglyphtounicode{gafarabic}{06AF}
\pdfglyphtounicode{gaffinalarabic}{FB93}
\pdfglyphtounicode{gafinitialarabic}{FB94}
\pdfglyphtounicode{gafmedialarabic}{FB95}
\pdfglyphtounicode{gagujarati}{0A97}
\pdfglyphtounicode{gagurmukhi}{0A17}
\pdfglyphtounicode{gahiragana}{304C}
\pdfglyphtounicode{gakatakana}{30AC}
\pdfglyphtounicode{gamma}{03B3}
\pdfglyphtounicode{gammalatinsmall}{0263}
\pdfglyphtounicode{gammasuperior}{02E0}
\pdfglyphtounicode{gangiacoptic}{03EB}
\pdfglyphtounicode{gbopomofo}{310D}
\pdfglyphtounicode{gbreve}{011F}
\pdfglyphtounicode{gcaron}{01E7}
\pdfglyphtounicode{gcedilla}{0123}
\pdfglyphtounicode{gcircle}{24D6}
\pdfglyphtounicode{gcircumflex}{011D}
\pdfglyphtounicode{gcommaaccent}{0123}
\pdfglyphtounicode{gdot}{0121}
\pdfglyphtounicode{gdotaccent}{0121}
\pdfglyphtounicode{gecyrillic}{0433}
\pdfglyphtounicode{gehiragana}{3052}
\pdfglyphtounicode{gekatakana}{30B2}
\pdfglyphtounicode{geomequivalent}{224E}
\pdfglyphtounicode{geometricallyequal}{2251}
\pdfglyphtounicode{gereshaccenthebrew}{059C}
\pdfglyphtounicode{gereshhebrew}{05F3}
\pdfglyphtounicode{gereshmuqdamhebrew}{059D}
\pdfglyphtounicode{germandbls}{00DF}
\pdfglyphtounicode{gershayimaccenthebrew}{059E}
\pdfglyphtounicode{gershayimhebrew}{05F4}
\pdfglyphtounicode{getamark}{3013}
\pdfglyphtounicode{ghabengali}{0998}
\pdfglyphtounicode{ghadarmenian}{0572}
\pdfglyphtounicode{ghadeva}{0918}
\pdfglyphtounicode{ghagujarati}{0A98}
\pdfglyphtounicode{ghagurmukhi}{0A18}
\pdfglyphtounicode{ghainarabic}{063A}
\pdfglyphtounicode{ghainfinalarabic}{FECE}
\pdfglyphtounicode{ghaininitialarabic}{FECF}
\pdfglyphtounicode{ghainmedialarabic}{FED0}
\pdfglyphtounicode{ghemiddlehookcyrillic}{0495}
\pdfglyphtounicode{ghestrokecyrillic}{0493}
\pdfglyphtounicode{gheupturncyrillic}{0491}
\pdfglyphtounicode{ghhadeva}{095A}
\pdfglyphtounicode{ghhagurmukhi}{0A5A}
\pdfglyphtounicode{ghook}{0260}
\pdfglyphtounicode{ghzsquare}{3393}
\pdfglyphtounicode{gihiragana}{304E}
\pdfglyphtounicode{gikatakana}{30AE}
\pdfglyphtounicode{gimarmenian}{0563}
\pdfglyphtounicode{gimel}{05D2}
\pdfglyphtounicode{gimeldagesh}{FB32}
\pdfglyphtounicode{gimeldageshhebrew}{FB32}
\pdfglyphtounicode{gimelhebrew}{05D2}
\pdfglyphtounicode{gjecyrillic}{0453}
\pdfglyphtounicode{glottalinvertedstroke}{01BE}
\pdfglyphtounicode{glottalstop}{0294}
\pdfglyphtounicode{glottalstopinverted}{0296}
\pdfglyphtounicode{glottalstopmod}{02C0}
\pdfglyphtounicode{glottalstopreversed}{0295}
\pdfglyphtounicode{glottalstopreversedmod}{02C1}
\pdfglyphtounicode{glottalstopreversedsuperior}{02E4}
\pdfglyphtounicode{glottalstopstroke}{02A1}
\pdfglyphtounicode{glottalstopstrokereversed}{02A2}
\pdfglyphtounicode{gmacron}{1E21}
\pdfglyphtounicode{gmonospace}{FF47}
\pdfglyphtounicode{gohiragana}{3054}
\pdfglyphtounicode{gokatakana}{30B4}
\pdfglyphtounicode{gparen}{24A2}
\pdfglyphtounicode{gpasquare}{33AC}
\pdfglyphtounicode{gradient}{2207}
\pdfglyphtounicode{grave}{0060}
\pdfglyphtounicode{gravebelowcmb}{0316}
\pdfglyphtounicode{gravecmb}{0300}
\pdfglyphtounicode{gravecomb}{0300}
\pdfglyphtounicode{gravedeva}{0953}
\pdfglyphtounicode{gravelowmod}{02CE}
\pdfglyphtounicode{gravemonospace}{FF40}
\pdfglyphtounicode{gravetonecmb}{0340}
\pdfglyphtounicode{greater}{003E}
\pdfglyphtounicode{greaterdbleqlless}{2A8C}
\pdfglyphtounicode{greaterdblequal}{2267}
\pdfglyphtounicode{greaterdot}{22D7}
\pdfglyphtounicode{greaterequal}{2265}
\pdfglyphtounicode{greaterequalorless}{22DB}
\pdfglyphtounicode{greaterlessequal}{22DB}
\pdfglyphtounicode{greatermonospace}{FF1E}
\pdfglyphtounicode{greatermuch}{226B}
\pdfglyphtounicode{greaternotdblequal}{2A8A}
\pdfglyphtounicode{greaternotequal}{2A88}
\pdfglyphtounicode{greaterorapproxeql}{2A86}
\pdfglyphtounicode{greaterorequalslant}{2A7E}
\pdfglyphtounicode{greaterorequivalent}{2273}
\pdfglyphtounicode{greaterorless}{2277}
\pdfglyphtounicode{greaterornotdbleql}{2269}
\pdfglyphtounicode{greaterornotequal}{2269}
\pdfglyphtounicode{greaterorsimilar}{2273}
\pdfglyphtounicode{greateroverequal}{2267}
\pdfglyphtounicode{greatersmall}{FE65}
\pdfglyphtounicode{gscript}{0261}
\pdfglyphtounicode{gstroke}{01E5}
\pdfglyphtounicode{guhiragana}{3050}
\pdfglyphtounicode{guillemotleft}{00AB}
\pdfglyphtounicode{guillemotright}{00BB}
\pdfglyphtounicode{guilsinglleft}{2039}
\pdfglyphtounicode{guilsinglright}{203A}
\pdfglyphtounicode{gukatakana}{30B0}
\pdfglyphtounicode{guramusquare}{3318}
\pdfglyphtounicode{gysquare}{33C9}
\pdfglyphtounicode{h}{0068}
\pdfglyphtounicode{haabkhasiancyrillic}{04A9}
\pdfglyphtounicode{haaltonearabic}{06C1}
\pdfglyphtounicode{habengali}{09B9}
\pdfglyphtounicode{hadescendercyrillic}{04B3}
\pdfglyphtounicode{hadeva}{0939}
\pdfglyphtounicode{hagujarati}{0AB9}
\pdfglyphtounicode{hagurmukhi}{0A39}
\pdfglyphtounicode{haharabic}{062D}
\pdfglyphtounicode{hahfinalarabic}{FEA2}
\pdfglyphtounicode{hahinitialarabic}{FEA3}
\pdfglyphtounicode{hahiragana}{306F}
\pdfglyphtounicode{hahmedialarabic}{FEA4}
\pdfglyphtounicode{haitusquare}{332A}
\pdfglyphtounicode{hakatakana}{30CF}
\pdfglyphtounicode{hakatakanahalfwidth}{FF8A}
\pdfglyphtounicode{halantgurmukhi}{0A4D}
\pdfglyphtounicode{hamzaarabic}{0621}
\pdfglyphtounicode{hamzadammaarabic}{0621 064F}
\pdfglyphtounicode{hamzadammatanarabic}{0621 064C}
\pdfglyphtounicode{hamzafathaarabic}{0621 064E}
\pdfglyphtounicode{hamzafathatanarabic}{0621 064B}
\pdfglyphtounicode{hamzalowarabic}{0621}
\pdfglyphtounicode{hamzalowkasraarabic}{0621 0650}
\pdfglyphtounicode{hamzalowkasratanarabic}{0621 064D}
\pdfglyphtounicode{hamzasukunarabic}{0621 0652}
\pdfglyphtounicode{hangulfiller}{3164}
\pdfglyphtounicode{hardsigncyrillic}{044A}
\pdfglyphtounicode{harpoondownleft}{21C3}
\pdfglyphtounicode{harpoondownright}{21C2}
\pdfglyphtounicode{harpoonleftbarbup}{21BC}
\pdfglyphtounicode{harpoonleftright}{21CC}
\pdfglyphtounicode{harpoonrightbarbup}{21C0}
\pdfglyphtounicode{harpoonrightleft}{21CB}
\pdfglyphtounicode{harpoonupleft}{21BF}
\pdfglyphtounicode{harpoonupright}{21BE}
\pdfglyphtounicode{hasquare}{33CA}
\pdfglyphtounicode{hatafpatah}{05B2}
\pdfglyphtounicode{hatafpatah16}{05B2}
\pdfglyphtounicode{hatafpatah23}{05B2}
\pdfglyphtounicode{hatafpatah2f}{05B2}
\pdfglyphtounicode{hatafpatahhebrew}{05B2}
\pdfglyphtounicode{hatafpatahnarrowhebrew}{05B2}
\pdfglyphtounicode{hatafpatahquarterhebrew}{05B2}
\pdfglyphtounicode{hatafpatahwidehebrew}{05B2}
\pdfglyphtounicode{hatafqamats}{05B3}
\pdfglyphtounicode{hatafqamats1b}{05B3}
\pdfglyphtounicode{hatafqamats28}{05B3}
\pdfglyphtounicode{hatafqamats34}{05B3}
\pdfglyphtounicode{hatafqamatshebrew}{05B3}
\pdfglyphtounicode{hatafqamatsnarrowhebrew}{05B3}
\pdfglyphtounicode{hatafqamatsquarterhebrew}{05B3}
\pdfglyphtounicode{hatafqamatswidehebrew}{05B3}
\pdfglyphtounicode{hatafsegol}{05B1}
\pdfglyphtounicode{hatafsegol17}{05B1}
\pdfglyphtounicode{hatafsegol24}{05B1}
\pdfglyphtounicode{hatafsegol30}{05B1}
\pdfglyphtounicode{hatafsegolhebrew}{05B1}
\pdfglyphtounicode{hatafsegolnarrowhebrew}{05B1}
\pdfglyphtounicode{hatafsegolquarterhebrew}{05B1}
\pdfglyphtounicode{hatafsegolwidehebrew}{05B1}
\pdfglyphtounicode{hbar}{0127}
\pdfglyphtounicode{hbopomofo}{310F}
\pdfglyphtounicode{hbrevebelow}{1E2B}
\pdfglyphtounicode{hcedilla}{1E29}
\pdfglyphtounicode{hcircle}{24D7}
\pdfglyphtounicode{hcircumflex}{0125}
\pdfglyphtounicode{hdieresis}{1E27}
\pdfglyphtounicode{hdotaccent}{1E23}
\pdfglyphtounicode{hdotbelow}{1E25}
\pdfglyphtounicode{he}{05D4}
\pdfglyphtounicode{heart}{2665}
\pdfglyphtounicode{heartsuitblack}{2665}
\pdfglyphtounicode{heartsuitwhite}{2661}
\pdfglyphtounicode{hedagesh}{FB34}
\pdfglyphtounicode{hedageshhebrew}{FB34}
\pdfglyphtounicode{hehaltonearabic}{06C1}
\pdfglyphtounicode{heharabic}{0647}
\pdfglyphtounicode{hehebrew}{05D4}
\pdfglyphtounicode{hehfinalaltonearabic}{FBA7}
\pdfglyphtounicode{hehfinalalttwoarabic}{FEEA}
\pdfglyphtounicode{hehfinalarabic}{FEEA}
\pdfglyphtounicode{hehhamzaabovefinalarabic}{FBA5}
\pdfglyphtounicode{hehhamzaaboveisolatedarabic}{FBA4}
\pdfglyphtounicode{hehinitialaltonearabic}{FBA8}
\pdfglyphtounicode{hehinitialarabic}{FEEB}
\pdfglyphtounicode{hehiragana}{3078}
\pdfglyphtounicode{hehmedialaltonearabic}{FBA9}
\pdfglyphtounicode{hehmedialarabic}{FEEC}
\pdfglyphtounicode{heiseierasquare}{337B}
\pdfglyphtounicode{hekatakana}{30D8}
\pdfglyphtounicode{hekatakanahalfwidth}{FF8D}
\pdfglyphtounicode{hekutaarusquare}{3336}
\pdfglyphtounicode{henghook}{0267}
\pdfglyphtounicode{herutusquare}{3339}
\pdfglyphtounicode{het}{05D7}
\pdfglyphtounicode{hethebrew}{05D7}
\pdfglyphtounicode{hhook}{0266}
\pdfglyphtounicode{hhooksuperior}{02B1}
\pdfglyphtounicode{hieuhacirclekorean}{327B}
\pdfglyphtounicode{hieuhaparenkorean}{321B}
\pdfglyphtounicode{hieuhcirclekorean}{326D}
\pdfglyphtounicode{hieuhkorean}{314E}
\pdfglyphtounicode{hieuhparenkorean}{320D}
\pdfglyphtounicode{hihiragana}{3072}
\pdfglyphtounicode{hikatakana}{30D2}
\pdfglyphtounicode{hikatakanahalfwidth}{FF8B}
\pdfglyphtounicode{hiriq}{05B4}
\pdfglyphtounicode{hiriq14}{05B4}
\pdfglyphtounicode{hiriq21}{05B4}
\pdfglyphtounicode{hiriq2d}{05B4}
\pdfglyphtounicode{hiriqhebrew}{05B4}
\pdfglyphtounicode{hiriqnarrowhebrew}{05B4}
\pdfglyphtounicode{hiriqquarterhebrew}{05B4}
\pdfglyphtounicode{hiriqwidehebrew}{05B4}
\pdfglyphtounicode{hlinebelow}{1E96}
\pdfglyphtounicode{hmonospace}{FF48}
\pdfglyphtounicode{hoarmenian}{0570}
\pdfglyphtounicode{hohipthai}{0E2B}
\pdfglyphtounicode{hohiragana}{307B}
\pdfglyphtounicode{hokatakana}{30DB}
\pdfglyphtounicode{hokatakanahalfwidth}{FF8E}
\pdfglyphtounicode{holam}{05B9}
\pdfglyphtounicode{holam19}{05B9}
\pdfglyphtounicode{holam26}{05B9}
\pdfglyphtounicode{holam32}{05B9}
\pdfglyphtounicode{holamhebrew}{05B9}
\pdfglyphtounicode{holamnarrowhebrew}{05B9}
\pdfglyphtounicode{holamquarterhebrew}{05B9}
\pdfglyphtounicode{holamwidehebrew}{05B9}
\pdfglyphtounicode{honokhukthai}{0E2E}
\pdfglyphtounicode{hookabovecomb}{0309}
\pdfglyphtounicode{hookcmb}{0309}
\pdfglyphtounicode{hookpalatalizedbelowcmb}{0321}
\pdfglyphtounicode{hookretroflexbelowcmb}{0322}
\pdfglyphtounicode{hoonsquare}{3342}
\pdfglyphtounicode{horicoptic}{03E9}
\pdfglyphtounicode{horizontalbar}{2015}
\pdfglyphtounicode{horncmb}{031B}
\pdfglyphtounicode{hotsprings}{2668}
\pdfglyphtounicode{house}{2302}
\pdfglyphtounicode{hparen}{24A3}
\pdfglyphtounicode{hsuperior}{02B0}
\pdfglyphtounicode{hturned}{0265}
\pdfglyphtounicode{huhiragana}{3075}
\pdfglyphtounicode{huiitosquare}{3333}
\pdfglyphtounicode{hukatakana}{30D5}
\pdfglyphtounicode{hukatakanahalfwidth}{FF8C}
\pdfglyphtounicode{hungarumlaut}{02DD}
\pdfglyphtounicode{hungarumlautcmb}{030B}
\pdfglyphtounicode{hv}{0195}
\pdfglyphtounicode{hyphen}{002D}
\pdfglyphtounicode{hyphenchar}{002D}
\pdfglyphtounicode{hypheninferior}{002D}
\pdfglyphtounicode{hyphenmonospace}{FF0D}
\pdfglyphtounicode{hyphensmall}{FE63}
\pdfglyphtounicode{hyphensuperior}{002D}
\pdfglyphtounicode{hyphentwo}{2010}
\pdfglyphtounicode{i}{0069}
\pdfglyphtounicode{iacute}{00ED}
\pdfglyphtounicode{iacyrillic}{044F}
\pdfglyphtounicode{ibengali}{0987}
\pdfglyphtounicode{ibopomofo}{3127}
\pdfglyphtounicode{ibreve}{012D}
\pdfglyphtounicode{icaron}{01D0}
\pdfglyphtounicode{icircle}{24D8}
\pdfglyphtounicode{icircumflex}{00EE}
\pdfglyphtounicode{icyrillic}{0456}
\pdfglyphtounicode{idblgrave}{0209}
\pdfglyphtounicode{ideographearthcircle}{328F}
\pdfglyphtounicode{ideographfirecircle}{328B}
\pdfglyphtounicode{ideographicallianceparen}{323F}
\pdfglyphtounicode{ideographiccallparen}{323A}
\pdfglyphtounicode{ideographiccentrecircle}{32A5}
\pdfglyphtounicode{ideographicclose}{3006}
\pdfglyphtounicode{ideographiccomma}{3001}
\pdfglyphtounicode{ideographiccommaleft}{FF64}
\pdfglyphtounicode{ideographiccongratulationparen}{3237}
\pdfglyphtounicode{ideographiccorrectcircle}{32A3}
\pdfglyphtounicode{ideographicearthparen}{322F}
\pdfglyphtounicode{ideographicenterpriseparen}{323D}
\pdfglyphtounicode{ideographicexcellentcircle}{329D}
\pdfglyphtounicode{ideographicfestivalparen}{3240}
\pdfglyphtounicode{ideographicfinancialcircle}{3296}
\pdfglyphtounicode{ideographicfinancialparen}{3236}
\pdfglyphtounicode{ideographicfireparen}{322B}
\pdfglyphtounicode{ideographichaveparen}{3232}
\pdfglyphtounicode{ideographichighcircle}{32A4}
\pdfglyphtounicode{ideographiciterationmark}{3005}
\pdfglyphtounicode{ideographiclaborcircle}{3298}
\pdfglyphtounicode{ideographiclaborparen}{3238}
\pdfglyphtounicode{ideographicleftcircle}{32A7}
\pdfglyphtounicode{ideographiclowcircle}{32A6}
\pdfglyphtounicode{ideographicmedicinecircle}{32A9}
\pdfglyphtounicode{ideographicmetalparen}{322E}
\pdfglyphtounicode{ideographicmoonparen}{322A}
\pdfglyphtounicode{ideographicnameparen}{3234}
\pdfglyphtounicode{ideographicperiod}{3002}
\pdfglyphtounicode{ideographicprintcircle}{329E}
\pdfglyphtounicode{ideographicreachparen}{3243}
\pdfglyphtounicode{ideographicrepresentparen}{3239}
\pdfglyphtounicode{ideographicresourceparen}{323E}
\pdfglyphtounicode{ideographicrightcircle}{32A8}
\pdfglyphtounicode{ideographicsecretcircle}{3299}
\pdfglyphtounicode{ideographicselfparen}{3242}
\pdfglyphtounicode{ideographicsocietyparen}{3233}
\pdfglyphtounicode{ideographicspace}{3000}
\pdfglyphtounicode{ideographicspecialparen}{3235}
\pdfglyphtounicode{ideographicstockparen}{3231}
\pdfglyphtounicode{ideographicstudyparen}{323B}
\pdfglyphtounicode{ideographicsunparen}{3230}
\pdfglyphtounicode{ideographicsuperviseparen}{323C}
\pdfglyphtounicode{ideographicwaterparen}{322C}
\pdfglyphtounicode{ideographicwoodparen}{322D}
\pdfglyphtounicode{ideographiczero}{3007}
\pdfglyphtounicode{ideographmetalcircle}{328E}
\pdfglyphtounicode{ideographmooncircle}{328A}
\pdfglyphtounicode{ideographnamecircle}{3294}
\pdfglyphtounicode{ideographsuncircle}{3290}
\pdfglyphtounicode{ideographwatercircle}{328C}
\pdfglyphtounicode{ideographwoodcircle}{328D}
\pdfglyphtounicode{ideva}{0907}
\pdfglyphtounicode{idieresis}{00EF}
\pdfglyphtounicode{idieresisacute}{1E2F}
\pdfglyphtounicode{idieresiscyrillic}{04E5}
\pdfglyphtounicode{idotbelow}{1ECB}
\pdfglyphtounicode{iebrevecyrillic}{04D7}
\pdfglyphtounicode{iecyrillic}{0435}
\pdfglyphtounicode{ieungacirclekorean}{3275}
\pdfglyphtounicode{ieungaparenkorean}{3215}
\pdfglyphtounicode{ieungcirclekorean}{3267}
\pdfglyphtounicode{ieungkorean}{3147}
\pdfglyphtounicode{ieungparenkorean}{3207}
\pdfglyphtounicode{igrave}{00EC}
\pdfglyphtounicode{igujarati}{0A87}
\pdfglyphtounicode{igurmukhi}{0A07}
\pdfglyphtounicode{ihiragana}{3044}
\pdfglyphtounicode{ihookabove}{1EC9}
\pdfglyphtounicode{iibengali}{0988}
\pdfglyphtounicode{iicyrillic}{0438}
\pdfglyphtounicode{iideva}{0908}
\pdfglyphtounicode{iigujarati}{0A88}
\pdfglyphtounicode{iigurmukhi}{0A08}
\pdfglyphtounicode{iimatragurmukhi}{0A40}
\pdfglyphtounicode{iinvertedbreve}{020B}
\pdfglyphtounicode{iishortcyrillic}{0439}
\pdfglyphtounicode{iivowelsignbengali}{09C0}
\pdfglyphtounicode{iivowelsigndeva}{0940}
\pdfglyphtounicode{iivowelsigngujarati}{0AC0}
\pdfglyphtounicode{ij}{0133}
\pdfglyphtounicode{ikatakana}{30A4}
\pdfglyphtounicode{ikatakanahalfwidth}{FF72}
\pdfglyphtounicode{ikorean}{3163}
\pdfglyphtounicode{ilde}{02DC}
\pdfglyphtounicode{iluyhebrew}{05AC}
\pdfglyphtounicode{imacron}{012B}
\pdfglyphtounicode{imacroncyrillic}{04E3}
\pdfglyphtounicode{imageorapproximatelyequal}{2253}
\pdfglyphtounicode{imatragurmukhi}{0A3F}
\pdfglyphtounicode{imonospace}{FF49}
\pdfglyphtounicode{increment}{2206}
\pdfglyphtounicode{infinity}{221E}
\pdfglyphtounicode{iniarmenian}{056B}
\pdfglyphtounicode{integerdivide}{2216}
\pdfglyphtounicode{integral}{222B}
\pdfglyphtounicode{integralbottom}{2321}
\pdfglyphtounicode{integralbt}{2321}
\pdfglyphtounicode{integralex}{F8F5}
\pdfglyphtounicode{integraltop}{2320}
\pdfglyphtounicode{integraltp}{2320}
\pdfglyphtounicode{intercal}{22BA}
\pdfglyphtounicode{interrobang}{203D}
\pdfglyphtounicode{interrobangdown}{2E18}
\pdfglyphtounicode{intersection}{2229}
\pdfglyphtounicode{intersectiondbl}{22D2}
\pdfglyphtounicode{intersectionsq}{2293}
\pdfglyphtounicode{intisquare}{3305}
\pdfglyphtounicode{invbullet}{25D8}
\pdfglyphtounicode{invcircle}{25D9}
\pdfglyphtounicode{invsmileface}{263B}
\pdfglyphtounicode{iocyrillic}{0451}
\pdfglyphtounicode{iogonek}{012F}
\pdfglyphtounicode{iota}{03B9}
\pdfglyphtounicode{iotadieresis}{03CA}
\pdfglyphtounicode{iotadieresistonos}{0390}
\pdfglyphtounicode{iotalatin}{0269}
\pdfglyphtounicode{iotatonos}{03AF}
\pdfglyphtounicode{iparen}{24A4}
\pdfglyphtounicode{irigurmukhi}{0A72}
\pdfglyphtounicode{ismallhiragana}{3043}
\pdfglyphtounicode{ismallkatakana}{30A3}
\pdfglyphtounicode{ismallkatakanahalfwidth}{FF68}
\pdfglyphtounicode{issharbengali}{09FA}
\pdfglyphtounicode{istroke}{0268}
\pdfglyphtounicode{isuperior}{0069}
\pdfglyphtounicode{iterationhiragana}{309D}
\pdfglyphtounicode{iterationkatakana}{30FD}
\pdfglyphtounicode{itilde}{0129}
\pdfglyphtounicode{itildebelow}{1E2D}
\pdfglyphtounicode{iubopomofo}{3129}
\pdfglyphtounicode{iucyrillic}{044E}
\pdfglyphtounicode{ivowelsignbengali}{09BF}
\pdfglyphtounicode{ivowelsigndeva}{093F}
\pdfglyphtounicode{ivowelsigngujarati}{0ABF}
\pdfglyphtounicode{izhitsacyrillic}{0475}
\pdfglyphtounicode{izhitsadblgravecyrillic}{0477}
\pdfglyphtounicode{j}{006A}
\pdfglyphtounicode{jaarmenian}{0571}
\pdfglyphtounicode{jabengali}{099C}
\pdfglyphtounicode{jadeva}{091C}
\pdfglyphtounicode{jagujarati}{0A9C}
\pdfglyphtounicode{jagurmukhi}{0A1C}
\pdfglyphtounicode{jbopomofo}{3110}
\pdfglyphtounicode{jcaron}{01F0}
\pdfglyphtounicode{jcircle}{24D9}
\pdfglyphtounicode{jcircumflex}{0135}
\pdfglyphtounicode{jcrossedtail}{029D}
\pdfglyphtounicode{jdotlessstroke}{025F}
\pdfglyphtounicode{jecyrillic}{0458}
\pdfglyphtounicode{jeemarabic}{062C}
\pdfglyphtounicode{jeemfinalarabic}{FE9E}
\pdfglyphtounicode{jeeminitialarabic}{FE9F}
\pdfglyphtounicode{jeemmedialarabic}{FEA0}
\pdfglyphtounicode{jeharabic}{0698}
\pdfglyphtounicode{jehfinalarabic}{FB8B}
\pdfglyphtounicode{jhabengali}{099D}
\pdfglyphtounicode{jhadeva}{091D}
\pdfglyphtounicode{jhagujarati}{0A9D}
\pdfglyphtounicode{jhagurmukhi}{0A1D}
\pdfglyphtounicode{jheharmenian}{057B}
\pdfglyphtounicode{jis}{3004}
\pdfglyphtounicode{jmonospace}{FF4A}
\pdfglyphtounicode{jparen}{24A5}
\pdfglyphtounicode{jsuperior}{02B2}
\pdfglyphtounicode{k}{006B}
\pdfglyphtounicode{kabashkircyrillic}{04A1}
\pdfglyphtounicode{kabengali}{0995}
\pdfglyphtounicode{kacute}{1E31}
\pdfglyphtounicode{kacyrillic}{043A}
\pdfglyphtounicode{kadescendercyrillic}{049B}
\pdfglyphtounicode{kadeva}{0915}
\pdfglyphtounicode{kaf}{05DB}
\pdfglyphtounicode{kafarabic}{0643}
\pdfglyphtounicode{kafdagesh}{FB3B}
\pdfglyphtounicode{kafdageshhebrew}{FB3B}
\pdfglyphtounicode{kaffinalarabic}{FEDA}
\pdfglyphtounicode{kafhebrew}{05DB}
\pdfglyphtounicode{kafinitialarabic}{FEDB}
\pdfglyphtounicode{kafmedialarabic}{FEDC}
\pdfglyphtounicode{kafrafehebrew}{FB4D}
\pdfglyphtounicode{kagujarati}{0A95}
\pdfglyphtounicode{kagurmukhi}{0A15}
\pdfglyphtounicode{kahiragana}{304B}
\pdfglyphtounicode{kahookcyrillic}{04C4}
\pdfglyphtounicode{kakatakana}{30AB}
\pdfglyphtounicode{kakatakanahalfwidth}{FF76}
\pdfglyphtounicode{kappa}{03BA}
\pdfglyphtounicode{kappasymbolgreek}{03F0}
\pdfglyphtounicode{kapyeounmieumkorean}{3171}
\pdfglyphtounicode{kapyeounphieuphkorean}{3184}
\pdfglyphtounicode{kapyeounpieupkorean}{3178}
\pdfglyphtounicode{kapyeounssangpieupkorean}{3179}
\pdfglyphtounicode{karoriisquare}{330D}
\pdfglyphtounicode{kashidaautoarabic}{0640}
\pdfglyphtounicode{kashidaautonosidebearingarabic}{0640}
\pdfglyphtounicode{kasmallkatakana}{30F5}
\pdfglyphtounicode{kasquare}{3384}
\pdfglyphtounicode{kasraarabic}{0650}
\pdfglyphtounicode{kasratanarabic}{064D}
\pdfglyphtounicode{kastrokecyrillic}{049F}
\pdfglyphtounicode{katahiraprolongmarkhalfwidth}{FF70}
\pdfglyphtounicode{kaverticalstrokecyrillic}{049D}
\pdfglyphtounicode{kbopomofo}{310E}
\pdfglyphtounicode{kcalsquare}{3389}
\pdfglyphtounicode{kcaron}{01E9}
\pdfglyphtounicode{kcedilla}{0137}
\pdfglyphtounicode{kcircle}{24DA}
\pdfglyphtounicode{kcommaaccent}{0137}
\pdfglyphtounicode{kdotbelow}{1E33}
\pdfglyphtounicode{keharmenian}{0584}
\pdfglyphtounicode{kehiragana}{3051}
\pdfglyphtounicode{kekatakana}{30B1}
\pdfglyphtounicode{kekatakanahalfwidth}{FF79}
\pdfglyphtounicode{kenarmenian}{056F}
\pdfglyphtounicode{kesmallkatakana}{30F6}
\pdfglyphtounicode{kgreenlandic}{0138}
\pdfglyphtounicode{khabengali}{0996}
\pdfglyphtounicode{khacyrillic}{0445}
\pdfglyphtounicode{khadeva}{0916}
\pdfglyphtounicode{khagujarati}{0A96}
\pdfglyphtounicode{khagurmukhi}{0A16}
\pdfglyphtounicode{khaharabic}{062E}
\pdfglyphtounicode{khahfinalarabic}{FEA6}
\pdfglyphtounicode{khahinitialarabic}{FEA7}
\pdfglyphtounicode{khahmedialarabic}{FEA8}
\pdfglyphtounicode{kheicoptic}{03E7}
\pdfglyphtounicode{khhadeva}{0959}
\pdfglyphtounicode{khhagurmukhi}{0A59}
\pdfglyphtounicode{khieukhacirclekorean}{3278}
\pdfglyphtounicode{khieukhaparenkorean}{3218}
\pdfglyphtounicode{khieukhcirclekorean}{326A}
\pdfglyphtounicode{khieukhkorean}{314B}
\pdfglyphtounicode{khieukhparenkorean}{320A}
\pdfglyphtounicode{khokhaithai}{0E02}
\pdfglyphtounicode{khokhonthai}{0E05}
\pdfglyphtounicode{khokhuatthai}{0E03}
\pdfglyphtounicode{khokhwaithai}{0E04}
\pdfglyphtounicode{khomutthai}{0E5B}
\pdfglyphtounicode{khook}{0199}
\pdfglyphtounicode{khorakhangthai}{0E06}
\pdfglyphtounicode{khzsquare}{3391}
\pdfglyphtounicode{kihiragana}{304D}
\pdfglyphtounicode{kikatakana}{30AD}
\pdfglyphtounicode{kikatakanahalfwidth}{FF77}
\pdfglyphtounicode{kiroguramusquare}{3315}
\pdfglyphtounicode{kiromeetorusquare}{3316}
\pdfglyphtounicode{kirosquare}{3314}
\pdfglyphtounicode{kiyeokacirclekorean}{326E}
\pdfglyphtounicode{kiyeokaparenkorean}{320E}
\pdfglyphtounicode{kiyeokcirclekorean}{3260}
\pdfglyphtounicode{kiyeokkorean}{3131}
\pdfglyphtounicode{kiyeokparenkorean}{3200}
\pdfglyphtounicode{kiyeoksioskorean}{3133}
\pdfglyphtounicode{kjecyrillic}{045C}
\pdfglyphtounicode{klinebelow}{1E35}
\pdfglyphtounicode{klsquare}{3398}
\pdfglyphtounicode{kmcubedsquare}{33A6}
\pdfglyphtounicode{kmonospace}{FF4B}
\pdfglyphtounicode{kmsquaredsquare}{33A2}
\pdfglyphtounicode{kohiragana}{3053}
\pdfglyphtounicode{kohmsquare}{33C0}
\pdfglyphtounicode{kokaithai}{0E01}
\pdfglyphtounicode{kokatakana}{30B3}
\pdfglyphtounicode{kokatakanahalfwidth}{FF7A}
\pdfglyphtounicode{kooposquare}{331E}
\pdfglyphtounicode{koppacyrillic}{0481}
\pdfglyphtounicode{koreanstandardsymbol}{327F}
\pdfglyphtounicode{koroniscmb}{0343}
\pdfglyphtounicode{kparen}{24A6}
\pdfglyphtounicode{kpasquare}{33AA}
\pdfglyphtounicode{ksicyrillic}{046F}
\pdfglyphtounicode{ktsquare}{33CF}
\pdfglyphtounicode{kturned}{029E}
\pdfglyphtounicode{kuhiragana}{304F}
\pdfglyphtounicode{kukatakana}{30AF}
\pdfglyphtounicode{kukatakanahalfwidth}{FF78}
\pdfglyphtounicode{kvsquare}{33B8}
\pdfglyphtounicode{kwsquare}{33BE}
\pdfglyphtounicode{l}{006C}
\pdfglyphtounicode{labengali}{09B2}
\pdfglyphtounicode{lacute}{013A}
\pdfglyphtounicode{ladeva}{0932}
\pdfglyphtounicode{lagujarati}{0AB2}
\pdfglyphtounicode{lagurmukhi}{0A32}
\pdfglyphtounicode{lakkhangyaothai}{0E45}
\pdfglyphtounicode{lamaleffinalarabic}{FEFC}
\pdfglyphtounicode{lamalefhamzaabovefinalarabic}{FEF8}
\pdfglyphtounicode{lamalefhamzaaboveisolatedarabic}{FEF7}
\pdfglyphtounicode{lamalefhamzabelowfinalarabic}{FEFA}
\pdfglyphtounicode{lamalefhamzabelowisolatedarabic}{FEF9}
\pdfglyphtounicode{lamalefisolatedarabic}{FEFB}
\pdfglyphtounicode{lamalefmaddaabovefinalarabic}{FEF6}
\pdfglyphtounicode{lamalefmaddaaboveisolatedarabic}{FEF5}
\pdfglyphtounicode{lamarabic}{0644}
\pdfglyphtounicode{lambda}{03BB}
\pdfglyphtounicode{lambdastroke}{019B}
\pdfglyphtounicode{lamed}{05DC}
\pdfglyphtounicode{lameddagesh}{FB3C}
\pdfglyphtounicode{lameddageshhebrew}{FB3C}
\pdfglyphtounicode{lamedhebrew}{05DC}
\pdfglyphtounicode{lamedholam}{05DC 05B9}
\pdfglyphtounicode{lamedholamdagesh}{05DC 05B9 05BC}
\pdfglyphtounicode{lamedholamdageshhebrew}{05DC 05B9 05BC}
\pdfglyphtounicode{lamedholamhebrew}{05DC 05B9}
\pdfglyphtounicode{lamfinalarabic}{FEDE}
\pdfglyphtounicode{lamhahinitialarabic}{FCCA}
\pdfglyphtounicode{laminitialarabic}{FEDF}
\pdfglyphtounicode{lamjeeminitialarabic}{FCC9}
\pdfglyphtounicode{lamkhahinitialarabic}{FCCB}
\pdfglyphtounicode{lamlamhehisolatedarabic}{FDF2}
\pdfglyphtounicode{lammedialarabic}{FEE0}
\pdfglyphtounicode{lammeemhahinitialarabic}{FD88}
\pdfglyphtounicode{lammeeminitialarabic}{FCCC}
\pdfglyphtounicode{lammeemjeeminitialarabic}{FEDF FEE4 FEA0}
\pdfglyphtounicode{lammeemkhahinitialarabic}{FEDF FEE4 FEA8}
\pdfglyphtounicode{largecircle}{25EF}
\pdfglyphtounicode{latticetop}{22A4}
\pdfglyphtounicode{lbar}{019A}
\pdfglyphtounicode{lbelt}{026C}
\pdfglyphtounicode{lbopomofo}{310C}
\pdfglyphtounicode{lcaron}{013E}
\pdfglyphtounicode{lcedilla}{013C}
\pdfglyphtounicode{lcircle}{24DB}
\pdfglyphtounicode{lcircumflexbelow}{1E3D}
\pdfglyphtounicode{lcommaaccent}{013C}
\pdfglyphtounicode{ldot}{0140}
\pdfglyphtounicode{ldotaccent}{0140}
\pdfglyphtounicode{ldotbelow}{1E37}
\pdfglyphtounicode{ldotbelowmacron}{1E39}
\pdfglyphtounicode{leftangleabovecmb}{031A}
\pdfglyphtounicode{lefttackbelowcmb}{0318}
\pdfglyphtounicode{less}{003C}
\pdfglyphtounicode{lessdbleqlgreater}{2A8B}
\pdfglyphtounicode{lessdblequal}{2266}
\pdfglyphtounicode{lessdot}{22D6}
\pdfglyphtounicode{lessequal}{2264}
\pdfglyphtounicode{lessequalgreater}{22DA}
\pdfglyphtounicode{lessequalorgreater}{22DA}
\pdfglyphtounicode{lessmonospace}{FF1C}
\pdfglyphtounicode{lessmuch}{226A}
\pdfglyphtounicode{lessnotdblequal}{2A89}
\pdfglyphtounicode{lessnotequal}{2A87}
\pdfglyphtounicode{lessorapproxeql}{2A85}
\pdfglyphtounicode{lessorequalslant}{2A7D}
\pdfglyphtounicode{lessorequivalent}{2272}
\pdfglyphtounicode{lessorgreater}{2276}
\pdfglyphtounicode{lessornotdbleql}{2268}
\pdfglyphtounicode{lessornotequal}{2268}
\pdfglyphtounicode{lessorsimilar}{2272}
\pdfglyphtounicode{lessoverequal}{2266}
\pdfglyphtounicode{lesssmall}{FE64}
\pdfglyphtounicode{lezh}{026E}
\pdfglyphtounicode{lfblock}{258C}
\pdfglyphtounicode{lhookretroflex}{026D}
\pdfglyphtounicode{lira}{20A4}
\pdfglyphtounicode{liwnarmenian}{056C}
\pdfglyphtounicode{lj}{01C9}
\pdfglyphtounicode{ljecyrillic}{0459}
\pdfglyphtounicode{ll}{006C 006C}
\pdfglyphtounicode{lladeva}{0933}
\pdfglyphtounicode{llagujarati}{0AB3}
\pdfglyphtounicode{llinebelow}{1E3B}
\pdfglyphtounicode{llladeva}{0934}
\pdfglyphtounicode{llvocalicbengali}{09E1}
\pdfglyphtounicode{llvocalicdeva}{0961}
\pdfglyphtounicode{llvocalicvowelsignbengali}{09E3}
\pdfglyphtounicode{llvocalicvowelsigndeva}{0963}
\pdfglyphtounicode{lmiddletilde}{026B}
\pdfglyphtounicode{lmonospace}{FF4C}
\pdfglyphtounicode{lmsquare}{33D0}
\pdfglyphtounicode{lochulathai}{0E2C}
\pdfglyphtounicode{logicaland}{2227}
\pdfglyphtounicode{logicalnot}{00AC}
\pdfglyphtounicode{logicalnotreversed}{2310}
\pdfglyphtounicode{logicalor}{2228}
\pdfglyphtounicode{lolingthai}{0E25}
\pdfglyphtounicode{longdbls}{017F 017F}
\pdfglyphtounicode{longs}{017F}
\pdfglyphtounicode{longsh}{017F 0068}
\pdfglyphtounicode{longsi}{017F 0069}
\pdfglyphtounicode{longsl}{017F 006C}
\pdfglyphtounicode{longst}{017F 0074}
\pdfglyphtounicode{lowlinecenterline}{FE4E}
\pdfglyphtounicode{lowlinecmb}{0332}
\pdfglyphtounicode{lowlinedashed}{FE4D}
\pdfglyphtounicode{lozenge}{25CA}
\pdfglyphtounicode{lparen}{24A7}
\pdfglyphtounicode{lscript}{2113}
\pdfglyphtounicode{lslash}{0142}
\pdfglyphtounicode{lsquare}{2113}
\pdfglyphtounicode{lsuperior}{006C}
\pdfglyphtounicode{ltshade}{2591}
\pdfglyphtounicode{luthai}{0E26}
\pdfglyphtounicode{lvocalicbengali}{098C}
\pdfglyphtounicode{lvocalicdeva}{090C}
\pdfglyphtounicode{lvocalicvowelsignbengali}{09E2}
\pdfglyphtounicode{lvocalicvowelsigndeva}{0962}
\pdfglyphtounicode{lxsquare}{33D3}
\pdfglyphtounicode{m}{006D}
\pdfglyphtounicode{mabengali}{09AE}
\pdfglyphtounicode{macron}{00AF}
\pdfglyphtounicode{macronbelowcmb}{0331}
\pdfglyphtounicode{macroncmb}{0304}
\pdfglyphtounicode{macronlowmod}{02CD}
\pdfglyphtounicode{macronmonospace}{FFE3}
\pdfglyphtounicode{macute}{1E3F}
\pdfglyphtounicode{madeva}{092E}
\pdfglyphtounicode{magujarati}{0AAE}
\pdfglyphtounicode{magurmukhi}{0A2E}
\pdfglyphtounicode{mahapakhhebrew}{05A4}
\pdfglyphtounicode{mahapakhlefthebrew}{05A4}
\pdfglyphtounicode{mahiragana}{307E}
\pdfglyphtounicode{maichattawalowleftthai}{F895}
\pdfglyphtounicode{maichattawalowrightthai}{F894}
\pdfglyphtounicode{maichattawathai}{0E4B}
\pdfglyphtounicode{maichattawaupperleftthai}{F893}
\pdfglyphtounicode{maieklowleftthai}{F88C}
\pdfglyphtounicode{maieklowrightthai}{F88B}
\pdfglyphtounicode{maiekthai}{0E48}
\pdfglyphtounicode{maiekupperleftthai}{F88A}
\pdfglyphtounicode{maihanakatleftthai}{F884}
\pdfglyphtounicode{maihanakatthai}{0E31}
\pdfglyphtounicode{maitaikhuleftthai}{F889}
\pdfglyphtounicode{maitaikhuthai}{0E47}
\pdfglyphtounicode{maitholowleftthai}{F88F}
\pdfglyphtounicode{maitholowrightthai}{F88E}
\pdfglyphtounicode{maithothai}{0E49}
\pdfglyphtounicode{maithoupperleftthai}{F88D}
\pdfglyphtounicode{maitrilowleftthai}{F892}
\pdfglyphtounicode{maitrilowrightthai}{F891}
\pdfglyphtounicode{maitrithai}{0E4A}
\pdfglyphtounicode{maitriupperleftthai}{F890}
\pdfglyphtounicode{maiyamokthai}{0E46}
\pdfglyphtounicode{makatakana}{30DE}
\pdfglyphtounicode{makatakanahalfwidth}{FF8F}
\pdfglyphtounicode{male}{2642}
\pdfglyphtounicode{maltesecross}{2720}
\pdfglyphtounicode{mansyonsquare}{3347}
\pdfglyphtounicode{maqafhebrew}{05BE}
\pdfglyphtounicode{mars}{2642}
\pdfglyphtounicode{masoracirclehebrew}{05AF}
\pdfglyphtounicode{masquare}{3383}
\pdfglyphtounicode{mbopomofo}{3107}
\pdfglyphtounicode{mbsquare}{33D4}
\pdfglyphtounicode{mcircle}{24DC}
\pdfglyphtounicode{mcubedsquare}{33A5}
\pdfglyphtounicode{mdotaccent}{1E41}
\pdfglyphtounicode{mdotbelow}{1E43}
\pdfglyphtounicode{measuredangle}{2221}
\pdfglyphtounicode{meemarabic}{0645}
\pdfglyphtounicode{meemfinalarabic}{FEE2}
\pdfglyphtounicode{meeminitialarabic}{FEE3}
\pdfglyphtounicode{meemmedialarabic}{FEE4}
\pdfglyphtounicode{meemmeeminitialarabic}{FCD1}
\pdfglyphtounicode{meemmeemisolatedarabic}{FC48}
\pdfglyphtounicode{meetorusquare}{334D}
\pdfglyphtounicode{mehiragana}{3081}
\pdfglyphtounicode{meizierasquare}{337E}
\pdfglyphtounicode{mekatakana}{30E1}
\pdfglyphtounicode{mekatakanahalfwidth}{FF92}
\pdfglyphtounicode{mem}{05DE}
\pdfglyphtounicode{memdagesh}{FB3E}
\pdfglyphtounicode{memdageshhebrew}{FB3E}
\pdfglyphtounicode{memhebrew}{05DE}
\pdfglyphtounicode{menarmenian}{0574}
\pdfglyphtounicode{merkhahebrew}{05A5}
\pdfglyphtounicode{merkhakefulahebrew}{05A6}
\pdfglyphtounicode{merkhakefulalefthebrew}{05A6}
\pdfglyphtounicode{merkhalefthebrew}{05A5}
\pdfglyphtounicode{mhook}{0271}
\pdfglyphtounicode{mhzsquare}{3392}
\pdfglyphtounicode{middledotkatakanahalfwidth}{FF65}
\pdfglyphtounicode{middot}{00B7}
\pdfglyphtounicode{mieumacirclekorean}{3272}
\pdfglyphtounicode{mieumaparenkorean}{3212}
\pdfglyphtounicode{mieumcirclekorean}{3264}
\pdfglyphtounicode{mieumkorean}{3141}
\pdfglyphtounicode{mieumpansioskorean}{3170}
\pdfglyphtounicode{mieumparenkorean}{3204}
\pdfglyphtounicode{mieumpieupkorean}{316E}
\pdfglyphtounicode{mieumsioskorean}{316F}
\pdfglyphtounicode{mihiragana}{307F}
\pdfglyphtounicode{mikatakana}{30DF}
\pdfglyphtounicode{mikatakanahalfwidth}{FF90}
\pdfglyphtounicode{minus}{2212}
\pdfglyphtounicode{minusbelowcmb}{0320}
\pdfglyphtounicode{minuscircle}{2296}
\pdfglyphtounicode{minusmod}{02D7}
\pdfglyphtounicode{minusplus}{2213}
\pdfglyphtounicode{minute}{2032}
\pdfglyphtounicode{miribaarusquare}{334A}
\pdfglyphtounicode{mirisquare}{3349}
\pdfglyphtounicode{mlonglegturned}{0270}
\pdfglyphtounicode{mlsquare}{3396}
\pdfglyphtounicode{mmcubedsquare}{33A3}
\pdfglyphtounicode{mmonospace}{FF4D}
\pdfglyphtounicode{mmsquaredsquare}{339F}
\pdfglyphtounicode{mohiragana}{3082}
\pdfglyphtounicode{mohmsquare}{33C1}
\pdfglyphtounicode{mokatakana}{30E2}
\pdfglyphtounicode{mokatakanahalfwidth}{FF93}
\pdfglyphtounicode{molsquare}{33D6}
\pdfglyphtounicode{momathai}{0E21}
\pdfglyphtounicode{moverssquare}{33A7}
\pdfglyphtounicode{moverssquaredsquare}{33A8}
\pdfglyphtounicode{mparen}{24A8}
\pdfglyphtounicode{mpasquare}{33AB}
\pdfglyphtounicode{mssquare}{33B3}
\pdfglyphtounicode{msuperior}{006D}
\pdfglyphtounicode{mturned}{026F}
\pdfglyphtounicode{mu}{00B5}
\pdfglyphtounicode{mu1}{00B5}
\pdfglyphtounicode{muasquare}{3382}
\pdfglyphtounicode{muchgreater}{226B}
\pdfglyphtounicode{muchless}{226A}
\pdfglyphtounicode{mufsquare}{338C}
\pdfglyphtounicode{mugreek}{03BC}
\pdfglyphtounicode{mugsquare}{338D}
\pdfglyphtounicode{muhiragana}{3080}
\pdfglyphtounicode{mukatakana}{30E0}
\pdfglyphtounicode{mukatakanahalfwidth}{FF91}
\pdfglyphtounicode{mulsquare}{3395}
\pdfglyphtounicode{multicloseleft}{22C9}
\pdfglyphtounicode{multicloseright}{22CA}
\pdfglyphtounicode{multimap}{22B8}
\pdfglyphtounicode{multiopenleft}{22CB}
\pdfglyphtounicode{multiopenright}{22CC}
\pdfglyphtounicode{multiply}{00D7}
\pdfglyphtounicode{mumsquare}{339B}
\pdfglyphtounicode{munahhebrew}{05A3}
\pdfglyphtounicode{munahlefthebrew}{05A3}
\pdfglyphtounicode{musicalnote}{266A}
\pdfglyphtounicode{musicalnotedbl}{266B}
\pdfglyphtounicode{musicflatsign}{266D}
\pdfglyphtounicode{musicsharpsign}{266F}
\pdfglyphtounicode{mussquare}{33B2}
\pdfglyphtounicode{muvsquare}{33B6}
\pdfglyphtounicode{muwsquare}{33BC}
\pdfglyphtounicode{mvmegasquare}{33B9}
\pdfglyphtounicode{mvsquare}{33B7}
\pdfglyphtounicode{mwmegasquare}{33BF}
\pdfglyphtounicode{mwsquare}{33BD}
\pdfglyphtounicode{n}{006E}
\pdfglyphtounicode{nabengali}{09A8}
\pdfglyphtounicode{nabla}{2207}
\pdfglyphtounicode{nacute}{0144}
\pdfglyphtounicode{nadeva}{0928}
\pdfglyphtounicode{nagujarati}{0AA8}
\pdfglyphtounicode{nagurmukhi}{0A28}
\pdfglyphtounicode{nahiragana}{306A}
\pdfglyphtounicode{nakatakana}{30CA}
\pdfglyphtounicode{nakatakanahalfwidth}{FF85}
\pdfglyphtounicode{nand}{22BC}
\pdfglyphtounicode{napostrophe}{0149}
\pdfglyphtounicode{nasquare}{3381}
\pdfglyphtounicode{natural}{266E}
\pdfglyphtounicode{nbopomofo}{310B}
\pdfglyphtounicode{nbspace}{00A0}
\pdfglyphtounicode{ncaron}{0148}
\pdfglyphtounicode{ncedilla}{0146}
\pdfglyphtounicode{ncircle}{24DD}
\pdfglyphtounicode{ncircumflexbelow}{1E4B}
\pdfglyphtounicode{ncommaaccent}{0146}
\pdfglyphtounicode{ndotaccent}{1E45}
\pdfglyphtounicode{ndotbelow}{1E47}
\pdfglyphtounicode{negationslash}{0338}
\pdfglyphtounicode{nehiragana}{306D}
\pdfglyphtounicode{nekatakana}{30CD}
\pdfglyphtounicode{nekatakanahalfwidth}{FF88}
\pdfglyphtounicode{newsheqelsign}{20AA}
\pdfglyphtounicode{nfsquare}{338B}
\pdfglyphtounicode{ng}{014B}
\pdfglyphtounicode{ngabengali}{0999}
\pdfglyphtounicode{ngadeva}{0919}
\pdfglyphtounicode{ngagujarati}{0A99}
\pdfglyphtounicode{ngagurmukhi}{0A19}
\pdfglyphtounicode{ngonguthai}{0E07}
\pdfglyphtounicode{nhiragana}{3093}
\pdfglyphtounicode{nhookleft}{0272}
\pdfglyphtounicode{nhookretroflex}{0273}
\pdfglyphtounicode{nieunacirclekorean}{326F}
\pdfglyphtounicode{nieunaparenkorean}{320F}
\pdfglyphtounicode{nieuncieuckorean}{3135}
\pdfglyphtounicode{nieuncirclekorean}{3261}
\pdfglyphtounicode{nieunhieuhkorean}{3136}
\pdfglyphtounicode{nieunkorean}{3134}
\pdfglyphtounicode{nieunpansioskorean}{3168}
\pdfglyphtounicode{nieunparenkorean}{3201}
\pdfglyphtounicode{nieunsioskorean}{3167}
\pdfglyphtounicode{nieuntikeutkorean}{3166}
\pdfglyphtounicode{nihiragana}{306B}
\pdfglyphtounicode{nikatakana}{30CB}
\pdfglyphtounicode{nikatakanahalfwidth}{FF86}
\pdfglyphtounicode{nikhahitleftthai}{F899}
\pdfglyphtounicode{nikhahitthai}{0E4D}
\pdfglyphtounicode{nine}{0039}
\pdfglyphtounicode{ninearabic}{0669}
\pdfglyphtounicode{ninebengali}{09EF}
\pdfglyphtounicode{ninecircle}{2468}
\pdfglyphtounicode{ninecircleinversesansserif}{2792}
\pdfglyphtounicode{ninedeva}{096F}
\pdfglyphtounicode{ninegujarati}{0AEF}
\pdfglyphtounicode{ninegurmukhi}{0A6F}
\pdfglyphtounicode{ninehackarabic}{0669}
\pdfglyphtounicode{ninehangzhou}{3029}
\pdfglyphtounicode{nineideographicparen}{3228}
\pdfglyphtounicode{nineinferior}{2089}
\pdfglyphtounicode{ninemonospace}{FF19}
\pdfglyphtounicode{nineoldstyle}{0039}
\pdfglyphtounicode{nineparen}{247C}
\pdfglyphtounicode{nineperiod}{2490}
\pdfglyphtounicode{ninepersian}{06F9}
\pdfglyphtounicode{nineroman}{2178}
\pdfglyphtounicode{ninesuperior}{2079}
\pdfglyphtounicode{nineteencircle}{2472}
\pdfglyphtounicode{nineteenparen}{2486}
\pdfglyphtounicode{nineteenperiod}{249A}
\pdfglyphtounicode{ninethai}{0E59}
\pdfglyphtounicode{nj}{01CC}
\pdfglyphtounicode{njecyrillic}{045A}
\pdfglyphtounicode{nkatakana}{30F3}
\pdfglyphtounicode{nkatakanahalfwidth}{FF9D}
\pdfglyphtounicode{nlegrightlong}{019E}
\pdfglyphtounicode{nlinebelow}{1E49}
\pdfglyphtounicode{nmonospace}{FF4E}
\pdfglyphtounicode{nmsquare}{339A}
\pdfglyphtounicode{nnabengali}{09A3}
\pdfglyphtounicode{nnadeva}{0923}
\pdfglyphtounicode{nnagujarati}{0AA3}
\pdfglyphtounicode{nnagurmukhi}{0A23}
\pdfglyphtounicode{nnnadeva}{0929}
\pdfglyphtounicode{nohiragana}{306E}
\pdfglyphtounicode{nokatakana}{30CE}
\pdfglyphtounicode{nokatakanahalfwidth}{FF89}
\pdfglyphtounicode{nonbreakingspace}{00A0}
\pdfglyphtounicode{nonenthai}{0E13}
\pdfglyphtounicode{nonuthai}{0E19}
\pdfglyphtounicode{noonarabic}{0646}
\pdfglyphtounicode{noonfinalarabic}{FEE6}
\pdfglyphtounicode{noonghunnaarabic}{06BA}
\pdfglyphtounicode{noonghunnafinalarabic}{FB9F}
\pdfglyphtounicode{noonhehinitialarabic}{FEE7 FEEC}
\pdfglyphtounicode{nooninitialarabic}{FEE7}
\pdfglyphtounicode{noonjeeminitialarabic}{FCD2}
\pdfglyphtounicode{noonjeemisolatedarabic}{FC4B}
\pdfglyphtounicode{noonmedialarabic}{FEE8}
\pdfglyphtounicode{noonmeeminitialarabic}{FCD5}
\pdfglyphtounicode{noonmeemisolatedarabic}{FC4E}
\pdfglyphtounicode{noonnoonfinalarabic}{FC8D}
\pdfglyphtounicode{notapproxequal}{2247}
\pdfglyphtounicode{notarrowboth}{21AE}
\pdfglyphtounicode{notarrowleft}{219A}
\pdfglyphtounicode{notarrowright}{219B}
\pdfglyphtounicode{notbar}{2224}
\pdfglyphtounicode{notcontains}{220C}
\pdfglyphtounicode{notdblarrowboth}{21CE}
\pdfglyphtounicode{notdblarrowleft}{21CD}
\pdfglyphtounicode{notdblarrowright}{21CF}
\pdfglyphtounicode{notelement}{2209}
\pdfglyphtounicode{notelementof}{2209}
\pdfglyphtounicode{notequal}{2260}
\pdfglyphtounicode{notexistential}{2204}
\pdfglyphtounicode{notfollows}{2281}
\pdfglyphtounicode{notfollowsoreql}{2AB0 0338}
\pdfglyphtounicode{notforces}{22AE}
\pdfglyphtounicode{notforcesextra}{22AF}
\pdfglyphtounicode{notgreater}{226F}
\pdfglyphtounicode{notgreaterdblequal}{2267 0338}
\pdfglyphtounicode{notgreaterequal}{2271}
\pdfglyphtounicode{notgreaternorequal}{2271}
\pdfglyphtounicode{notgreaternorless}{2279}
\pdfglyphtounicode{notgreaterorslnteql}{2A7E 0338}
\pdfglyphtounicode{notidentical}{2262}
\pdfglyphtounicode{notless}{226E}
\pdfglyphtounicode{notlessdblequal}{2266 0338}
\pdfglyphtounicode{notlessequal}{2270}
\pdfglyphtounicode{notlessnorequal}{2270}
\pdfglyphtounicode{notlessorslnteql}{2A7D 0338}
\pdfglyphtounicode{notparallel}{2226}
\pdfglyphtounicode{notprecedes}{2280}
\pdfglyphtounicode{notprecedesoreql}{2AAF 0338}
\pdfglyphtounicode{notsatisfies}{22AD}
\pdfglyphtounicode{notsimilar}{2241}
\pdfglyphtounicode{notsubset}{2284}
\pdfglyphtounicode{notsubseteql}{2288}
\pdfglyphtounicode{notsubsetordbleql}{2AC5 0338}
\pdfglyphtounicode{notsubsetoreql}{228A}
\pdfglyphtounicode{notsucceeds}{2281}
\pdfglyphtounicode{notsuperset}{2285}
\pdfglyphtounicode{notsuperseteql}{2289}
\pdfglyphtounicode{notsupersetordbleql}{2AC6 0338}
\pdfglyphtounicode{notsupersetoreql}{228B}
\pdfglyphtounicode{nottriangeqlleft}{22EC}
\pdfglyphtounicode{nottriangeqlright}{22ED}
\pdfglyphtounicode{nottriangleleft}{22EA}
\pdfglyphtounicode{nottriangleright}{22EB}
\pdfglyphtounicode{notturnstile}{22AC}
\pdfglyphtounicode{nowarmenian}{0576}
\pdfglyphtounicode{nparen}{24A9}
\pdfglyphtounicode{nssquare}{33B1}
\pdfglyphtounicode{nsuperior}{207F}
\pdfglyphtounicode{ntilde}{00F1}
\pdfglyphtounicode{nu}{03BD}
\pdfglyphtounicode{nuhiragana}{306C}
\pdfglyphtounicode{nukatakana}{30CC}
\pdfglyphtounicode{nukatakanahalfwidth}{FF87}
\pdfglyphtounicode{nuktabengali}{09BC}
\pdfglyphtounicode{nuktadeva}{093C}
\pdfglyphtounicode{nuktagujarati}{0ABC}
\pdfglyphtounicode{nuktagurmukhi}{0A3C}
\pdfglyphtounicode{numbersign}{0023}
\pdfglyphtounicode{numbersignmonospace}{FF03}
\pdfglyphtounicode{numbersignsmall}{FE5F}
\pdfglyphtounicode{numeralsigngreek}{0374}
\pdfglyphtounicode{numeralsignlowergreek}{0375}
\pdfglyphtounicode{numero}{2116}
\pdfglyphtounicode{nun}{05E0}
\pdfglyphtounicode{nundagesh}{FB40}
\pdfglyphtounicode{nundageshhebrew}{FB40}
\pdfglyphtounicode{nunhebrew}{05E0}
\pdfglyphtounicode{nvsquare}{33B5}
\pdfglyphtounicode{nwsquare}{33BB}
\pdfglyphtounicode{nyabengali}{099E}
\pdfglyphtounicode{nyadeva}{091E}
\pdfglyphtounicode{nyagujarati}{0A9E}
\pdfglyphtounicode{nyagurmukhi}{0A1E}
\pdfglyphtounicode{o}{006F}
\pdfglyphtounicode{oacute}{00F3}
\pdfglyphtounicode{oangthai}{0E2D}
\pdfglyphtounicode{obarred}{0275}
\pdfglyphtounicode{obarredcyrillic}{04E9}
\pdfglyphtounicode{obarreddieresiscyrillic}{04EB}
\pdfglyphtounicode{obengali}{0993}
\pdfglyphtounicode{obopomofo}{311B}
\pdfglyphtounicode{obreve}{014F}
\pdfglyphtounicode{ocandradeva}{0911}
\pdfglyphtounicode{ocandragujarati}{0A91}
\pdfglyphtounicode{ocandravowelsigndeva}{0949}
\pdfglyphtounicode{ocandravowelsigngujarati}{0AC9}
\pdfglyphtounicode{ocaron}{01D2}
\pdfglyphtounicode{ocircle}{24DE}
\pdfglyphtounicode{ocircumflex}{00F4}
\pdfglyphtounicode{ocircumflexacute}{1ED1}
\pdfglyphtounicode{ocircumflexdotbelow}{1ED9}
\pdfglyphtounicode{ocircumflexgrave}{1ED3}
\pdfglyphtounicode{ocircumflexhookabove}{1ED5}
\pdfglyphtounicode{ocircumflextilde}{1ED7}
\pdfglyphtounicode{ocyrillic}{043E}
\pdfglyphtounicode{odblacute}{0151}
\pdfglyphtounicode{odblgrave}{020D}
\pdfglyphtounicode{odeva}{0913}
\pdfglyphtounicode{odieresis}{00F6}
\pdfglyphtounicode{odieresiscyrillic}{04E7}
\pdfglyphtounicode{odotbelow}{1ECD}
\pdfglyphtounicode{oe}{0153}
\pdfglyphtounicode{oekorean}{315A}
\pdfglyphtounicode{ogonek}{02DB}
\pdfglyphtounicode{ogonekcmb}{0328}
\pdfglyphtounicode{ograve}{00F2}
\pdfglyphtounicode{ogujarati}{0A93}
\pdfglyphtounicode{oharmenian}{0585}
\pdfglyphtounicode{ohiragana}{304A}
\pdfglyphtounicode{ohookabove}{1ECF}
\pdfglyphtounicode{ohorn}{01A1}
\pdfglyphtounicode{ohornacute}{1EDB}
\pdfglyphtounicode{ohorndotbelow}{1EE3}
\pdfglyphtounicode{ohorngrave}{1EDD}
\pdfglyphtounicode{ohornhookabove}{1EDF}
\pdfglyphtounicode{ohorntilde}{1EE1}
\pdfglyphtounicode{ohungarumlaut}{0151}
\pdfglyphtounicode{oi}{01A3}
\pdfglyphtounicode{oinvertedbreve}{020F}
\pdfglyphtounicode{okatakana}{30AA}
\pdfglyphtounicode{okatakanahalfwidth}{FF75}
\pdfglyphtounicode{okorean}{3157}
\pdfglyphtounicode{olehebrew}{05AB}
\pdfglyphtounicode{omacron}{014D}
\pdfglyphtounicode{omacronacute}{1E53}
\pdfglyphtounicode{omacrongrave}{1E51}
\pdfglyphtounicode{omdeva}{0950}
\pdfglyphtounicode{omega}{03C9}
\pdfglyphtounicode{omega1}{03D6}
\pdfglyphtounicode{omegacyrillic}{0461}
\pdfglyphtounicode{omegalatinclosed}{0277}
\pdfglyphtounicode{omegaroundcyrillic}{047B}
\pdfglyphtounicode{omegatitlocyrillic}{047D}
\pdfglyphtounicode{omegatonos}{03CE}
\pdfglyphtounicode{omgujarati}{0AD0}
\pdfglyphtounicode{omicron}{03BF}
\pdfglyphtounicode{omicrontonos}{03CC}
\pdfglyphtounicode{omonospace}{FF4F}
\pdfglyphtounicode{one}{0031}
\pdfglyphtounicode{onearabic}{0661}
\pdfglyphtounicode{onebengali}{09E7}
\pdfglyphtounicode{onecircle}{2460}
\pdfglyphtounicode{onecircleinversesansserif}{278A}
\pdfglyphtounicode{onedeva}{0967}
\pdfglyphtounicode{onedotenleader}{2024}
\pdfglyphtounicode{oneeighth}{215B}
\pdfglyphtounicode{onefitted}{0031}
\pdfglyphtounicode{onegujarati}{0AE7}
\pdfglyphtounicode{onegurmukhi}{0A67}
\pdfglyphtounicode{onehackarabic}{0661}
\pdfglyphtounicode{onehalf}{00BD}
\pdfglyphtounicode{onehangzhou}{3021}
\pdfglyphtounicode{oneideographicparen}{3220}
\pdfglyphtounicode{oneinferior}{2081}
\pdfglyphtounicode{onemonospace}{FF11}
\pdfglyphtounicode{onenumeratorbengali}{09F4}
\pdfglyphtounicode{oneoldstyle}{0031}
\pdfglyphtounicode{oneparen}{2474}
\pdfglyphtounicode{oneperiod}{2488}
\pdfglyphtounicode{onepersian}{06F1}
\pdfglyphtounicode{onequarter}{00BC}
\pdfglyphtounicode{oneroman}{2170}
\pdfglyphtounicode{onesuperior}{00B9}
\pdfglyphtounicode{onethai}{0E51}
\pdfglyphtounicode{onethird}{2153}
\pdfglyphtounicode{oogonek}{01EB}
\pdfglyphtounicode{oogonekmacron}{01ED}
\pdfglyphtounicode{oogurmukhi}{0A13}
\pdfglyphtounicode{oomatragurmukhi}{0A4B}
\pdfglyphtounicode{oopen}{0254}
\pdfglyphtounicode{oparen}{24AA}
\pdfglyphtounicode{openbullet}{25E6}
\pdfglyphtounicode{option}{2325}
\pdfglyphtounicode{ordfeminine}{00AA}
\pdfglyphtounicode{ordmasculine}{00BA}
\pdfglyphtounicode{orthogonal}{221F}
\pdfglyphtounicode{orunderscore}{22BB}
\pdfglyphtounicode{oshortdeva}{0912}
\pdfglyphtounicode{oshortvowelsigndeva}{094A}
\pdfglyphtounicode{oslash}{00F8}
\pdfglyphtounicode{oslashacute}{01FF}
\pdfglyphtounicode{osmallhiragana}{3049}
\pdfglyphtounicode{osmallkatakana}{30A9}
\pdfglyphtounicode{osmallkatakanahalfwidth}{FF6B}
\pdfglyphtounicode{ostrokeacute}{01FF}
\pdfglyphtounicode{osuperior}{006F}
\pdfglyphtounicode{otcyrillic}{047F}
\pdfglyphtounicode{otilde}{00F5}
\pdfglyphtounicode{otildeacute}{1E4D}
\pdfglyphtounicode{otildedieresis}{1E4F}
\pdfglyphtounicode{oubopomofo}{3121}
\pdfglyphtounicode{overline}{203E}
\pdfglyphtounicode{overlinecenterline}{FE4A}
\pdfglyphtounicode{overlinecmb}{0305}
\pdfglyphtounicode{overlinedashed}{FE49}
\pdfglyphtounicode{overlinedblwavy}{FE4C}
\pdfglyphtounicode{overlinewavy}{FE4B}
\pdfglyphtounicode{overscore}{00AF}
\pdfglyphtounicode{ovowelsignbengali}{09CB}
\pdfglyphtounicode{ovowelsigndeva}{094B}
\pdfglyphtounicode{ovowelsigngujarati}{0ACB}
\pdfglyphtounicode{owner}{220B}
\pdfglyphtounicode{p}{0070}
\pdfglyphtounicode{paampssquare}{3380}
\pdfglyphtounicode{paasentosquare}{332B}
\pdfglyphtounicode{pabengali}{09AA}
\pdfglyphtounicode{pacute}{1E55}
\pdfglyphtounicode{padeva}{092A}
\pdfglyphtounicode{pagedown}{21DF}
\pdfglyphtounicode{pageup}{21DE}
\pdfglyphtounicode{pagujarati}{0AAA}
\pdfglyphtounicode{pagurmukhi}{0A2A}
\pdfglyphtounicode{pahiragana}{3071}
\pdfglyphtounicode{paiyannoithai}{0E2F}
\pdfglyphtounicode{pakatakana}{30D1}
\pdfglyphtounicode{palatalizationcyrilliccmb}{0484}
\pdfglyphtounicode{palochkacyrillic}{04C0}
\pdfglyphtounicode{pansioskorean}{317F}
\pdfglyphtounicode{paragraph}{00B6}
\pdfglyphtounicode{parallel}{2225}
\pdfglyphtounicode{parenleft}{0028}
\pdfglyphtounicode{parenleftaltonearabic}{FD3E}
\pdfglyphtounicode{parenleftbt}{F8ED}
\pdfglyphtounicode{parenleftex}{F8EC}
\pdfglyphtounicode{parenleftinferior}{208D}
\pdfglyphtounicode{parenleftmonospace}{FF08}
\pdfglyphtounicode{parenleftsmall}{FE59}
\pdfglyphtounicode{parenleftsuperior}{207D}
\pdfglyphtounicode{parenlefttp}{F8EB}
\pdfglyphtounicode{parenleftvertical}{FE35}
\pdfglyphtounicode{parenright}{0029}
\pdfglyphtounicode{parenrightaltonearabic}{FD3F}
\pdfglyphtounicode{parenrightbt}{F8F8}
\pdfglyphtounicode{parenrightex}{F8F7}
\pdfglyphtounicode{parenrightinferior}{208E}
\pdfglyphtounicode{parenrightmonospace}{FF09}
\pdfglyphtounicode{parenrightsmall}{FE5A}
\pdfglyphtounicode{parenrightsuperior}{207E}
\pdfglyphtounicode{parenrighttp}{F8F6}
\pdfglyphtounicode{parenrightvertical}{FE36}
\pdfglyphtounicode{partialdiff}{2202}
\pdfglyphtounicode{paseqhebrew}{05C0}
\pdfglyphtounicode{pashtahebrew}{0599}
\pdfglyphtounicode{pasquare}{33A9}
\pdfglyphtounicode{patah}{05B7}
\pdfglyphtounicode{patah11}{05B7}
\pdfglyphtounicode{patah1d}{05B7}
\pdfglyphtounicode{patah2a}{05B7}
\pdfglyphtounicode{patahhebrew}{05B7}
\pdfglyphtounicode{patahnarrowhebrew}{05B7}
\pdfglyphtounicode{patahquarterhebrew}{05B7}
\pdfglyphtounicode{patahwidehebrew}{05B7}
\pdfglyphtounicode{pazerhebrew}{05A1}
\pdfglyphtounicode{pbopomofo}{3106}
\pdfglyphtounicode{pcircle}{24DF}
\pdfglyphtounicode{pdotaccent}{1E57}
\pdfglyphtounicode{pe}{05E4}
\pdfglyphtounicode{pecyrillic}{043F}
\pdfglyphtounicode{pedagesh}{FB44}
\pdfglyphtounicode{pedageshhebrew}{FB44}
\pdfglyphtounicode{peezisquare}{333B}
\pdfglyphtounicode{pefinaldageshhebrew}{FB43}
\pdfglyphtounicode{peharabic}{067E}
\pdfglyphtounicode{peharmenian}{057A}
\pdfglyphtounicode{pehebrew}{05E4}
\pdfglyphtounicode{pehfinalarabic}{FB57}
\pdfglyphtounicode{pehinitialarabic}{FB58}
\pdfglyphtounicode{pehiragana}{307A}
\pdfglyphtounicode{pehmedialarabic}{FB59}
\pdfglyphtounicode{pekatakana}{30DA}
\pdfglyphtounicode{pemiddlehookcyrillic}{04A7}
\pdfglyphtounicode{perafehebrew}{FB4E}
\pdfglyphtounicode{percent}{0025}
\pdfglyphtounicode{percentarabic}{066A}
\pdfglyphtounicode{percentmonospace}{FF05}
\pdfglyphtounicode{percentsmall}{FE6A}
\pdfglyphtounicode{period}{002E}
\pdfglyphtounicode{periodarmenian}{0589}
\pdfglyphtounicode{periodcentered}{00B7}
\pdfglyphtounicode{periodhalfwidth}{FF61}
\pdfglyphtounicode{periodinferior}{002E}
\pdfglyphtounicode{periodmonospace}{FF0E}
\pdfglyphtounicode{periodsmall}{FE52}
\pdfglyphtounicode{periodsuperior}{002E}
\pdfglyphtounicode{perispomenigreekcmb}{0342}
\pdfglyphtounicode{perpcorrespond}{2A5E}
\pdfglyphtounicode{perpendicular}{22A5}
\pdfglyphtounicode{pertenthousand}{2031}
\pdfglyphtounicode{perthousand}{2030}
\pdfglyphtounicode{peseta}{20A7}
\pdfglyphtounicode{pfsquare}{338A}
\pdfglyphtounicode{phabengali}{09AB}
\pdfglyphtounicode{phadeva}{092B}
\pdfglyphtounicode{phagujarati}{0AAB}
\pdfglyphtounicode{phagurmukhi}{0A2B}
\pdfglyphtounicode{phi}{03C6}
\pdfglyphtounicode{phi1}{03D5}
\pdfglyphtounicode{phieuphacirclekorean}{327A}
\pdfglyphtounicode{phieuphaparenkorean}{321A}
\pdfglyphtounicode{phieuphcirclekorean}{326C}
\pdfglyphtounicode{phieuphkorean}{314D}
\pdfglyphtounicode{phieuphparenkorean}{320C}
\pdfglyphtounicode{philatin}{0278}
\pdfglyphtounicode{phinthuthai}{0E3A}
\pdfglyphtounicode{phisymbolgreek}{03D5}
\pdfglyphtounicode{phook}{01A5}
\pdfglyphtounicode{phophanthai}{0E1E}
\pdfglyphtounicode{phophungthai}{0E1C}
\pdfglyphtounicode{phosamphaothai}{0E20}
\pdfglyphtounicode{pi}{03C0}
\pdfglyphtounicode{pi1}{03D6}
\pdfglyphtounicode{pieupacirclekorean}{3273}
\pdfglyphtounicode{pieupaparenkorean}{3213}
\pdfglyphtounicode{pieupcieuckorean}{3176}
\pdfglyphtounicode{pieupcirclekorean}{3265}
\pdfglyphtounicode{pieupkiyeokkorean}{3172}
\pdfglyphtounicode{pieupkorean}{3142}
\pdfglyphtounicode{pieupparenkorean}{3205}
\pdfglyphtounicode{pieupsioskiyeokkorean}{3174}
\pdfglyphtounicode{pieupsioskorean}{3144}
\pdfglyphtounicode{pieupsiostikeutkorean}{3175}
\pdfglyphtounicode{pieupthieuthkorean}{3177}
\pdfglyphtounicode{pieuptikeutkorean}{3173}
\pdfglyphtounicode{pihiragana}{3074}
\pdfglyphtounicode{pikatakana}{30D4}
\pdfglyphtounicode{pisymbolgreek}{03D6}
\pdfglyphtounicode{piwrarmenian}{0583}
\pdfglyphtounicode{planckover2pi}{210F}
\pdfglyphtounicode{planckover2pi1}{210F}
\pdfglyphtounicode{plus}{002B}
\pdfglyphtounicode{plusbelowcmb}{031F}
\pdfglyphtounicode{pluscircle}{2295}
\pdfglyphtounicode{plusminus}{00B1}
\pdfglyphtounicode{plusmod}{02D6}
\pdfglyphtounicode{plusmonospace}{FF0B}
\pdfglyphtounicode{plussmall}{FE62}
\pdfglyphtounicode{plussuperior}{207A}
\pdfglyphtounicode{pmonospace}{FF50}
\pdfglyphtounicode{pmsquare}{33D8}
\pdfglyphtounicode{pohiragana}{307D}
\pdfglyphtounicode{pointingindexdownwhite}{261F}
\pdfglyphtounicode{pointingindexleftwhite}{261C}
\pdfglyphtounicode{pointingindexrightwhite}{261E}
\pdfglyphtounicode{pointingindexupwhite}{261D}
\pdfglyphtounicode{pokatakana}{30DD}
\pdfglyphtounicode{poplathai}{0E1B}
\pdfglyphtounicode{postalmark}{3012}
\pdfglyphtounicode{postalmarkface}{3020}
\pdfglyphtounicode{pparen}{24AB}
\pdfglyphtounicode{precedenotdbleqv}{2AB9}
\pdfglyphtounicode{precedenotslnteql}{2AB5}
\pdfglyphtounicode{precedeornoteqvlnt}{22E8}
\pdfglyphtounicode{precedes}{227A}
\pdfglyphtounicode{precedesequal}{2AAF}
\pdfglyphtounicode{precedesorcurly}{227C}
\pdfglyphtounicode{precedesorequal}{227E}
\pdfglyphtounicode{prescription}{211E}
\pdfglyphtounicode{prime}{2032}
\pdfglyphtounicode{primemod}{02B9}
\pdfglyphtounicode{primereverse}{2035}
\pdfglyphtounicode{primereversed}{2035}
\pdfglyphtounicode{product}{220F}
\pdfglyphtounicode{projective}{2305}
\pdfglyphtounicode{prolongedkana}{30FC}
\pdfglyphtounicode{propellor}{2318}
\pdfglyphtounicode{propersubset}{2282}
\pdfglyphtounicode{propersuperset}{2283}
\pdfglyphtounicode{proportion}{2237}
\pdfglyphtounicode{proportional}{221D}
\pdfglyphtounicode{psi}{03C8}
\pdfglyphtounicode{psicyrillic}{0471}
\pdfglyphtounicode{psilipneumatacyrilliccmb}{0486}
\pdfglyphtounicode{pssquare}{33B0}
\pdfglyphtounicode{puhiragana}{3077}
\pdfglyphtounicode{pukatakana}{30D7}
\pdfglyphtounicode{punctdash}{2014}
\pdfglyphtounicode{pvsquare}{33B4}
\pdfglyphtounicode{pwsquare}{33BA}
\pdfglyphtounicode{q}{0071}
\pdfglyphtounicode{qadeva}{0958}
\pdfglyphtounicode{qadmahebrew}{05A8}
\pdfglyphtounicode{qafarabic}{0642}
\pdfglyphtounicode{qaffinalarabic}{FED6}
\pdfglyphtounicode{qafinitialarabic}{FED7}
\pdfglyphtounicode{qafmedialarabic}{FED8}
\pdfglyphtounicode{qamats}{05B8}
\pdfglyphtounicode{qamats10}{05B8}
\pdfglyphtounicode{qamats1a}{05B8}
\pdfglyphtounicode{qamats1c}{05B8}
\pdfglyphtounicode{qamats27}{05B8}
\pdfglyphtounicode{qamats29}{05B8}
\pdfglyphtounicode{qamats33}{05B8}
\pdfglyphtounicode{qamatsde}{05B8}
\pdfglyphtounicode{qamatshebrew}{05B8}
\pdfglyphtounicode{qamatsnarrowhebrew}{05B8}
\pdfglyphtounicode{qamatsqatanhebrew}{05B8}
\pdfglyphtounicode{qamatsqatannarrowhebrew}{05B8}
\pdfglyphtounicode{qamatsqatanquarterhebrew}{05B8}
\pdfglyphtounicode{qamatsqatanwidehebrew}{05B8}
\pdfglyphtounicode{qamatsquarterhebrew}{05B8}
\pdfglyphtounicode{qamatswidehebrew}{05B8}
\pdfglyphtounicode{qarneyparahebrew}{059F}
\pdfglyphtounicode{qbopomofo}{3111}
\pdfglyphtounicode{qcircle}{24E0}
\pdfglyphtounicode{qhook}{02A0}
\pdfglyphtounicode{qmonospace}{FF51}
\pdfglyphtounicode{qof}{05E7}
\pdfglyphtounicode{qofdagesh}{FB47}
\pdfglyphtounicode{qofdageshhebrew}{FB47}
\pdfglyphtounicode{qofhatafpatah}{05E7 05B2}
\pdfglyphtounicode{qofhatafpatahhebrew}{05E7 05B2}
\pdfglyphtounicode{qofhatafsegol}{05E7 05B1}
\pdfglyphtounicode{qofhatafsegolhebrew}{05E7 05B1}
\pdfglyphtounicode{qofhebrew}{05E7}
\pdfglyphtounicode{qofhiriq}{05E7 05B4}
\pdfglyphtounicode{qofhiriqhebrew}{05E7 05B4}
\pdfglyphtounicode{qofholam}{05E7 05B9}
\pdfglyphtounicode{qofholamhebrew}{05E7 05B9}
\pdfglyphtounicode{qofpatah}{05E7 05B7}
\pdfglyphtounicode{qofpatahhebrew}{05E7 05B7}
\pdfglyphtounicode{qofqamats}{05E7 05B8}
\pdfglyphtounicode{qofqamatshebrew}{05E7 05B8}
\pdfglyphtounicode{qofqubuts}{05E7 05BB}
\pdfglyphtounicode{qofqubutshebrew}{05E7 05BB}
\pdfglyphtounicode{qofsegol}{05E7 05B6}
\pdfglyphtounicode{qofsegolhebrew}{05E7 05B6}
\pdfglyphtounicode{qofsheva}{05E7 05B0}
\pdfglyphtounicode{qofshevahebrew}{05E7 05B0}
\pdfglyphtounicode{qoftsere}{05E7 05B5}
\pdfglyphtounicode{qoftserehebrew}{05E7 05B5}
\pdfglyphtounicode{qparen}{24AC}
\pdfglyphtounicode{quarternote}{2669}
\pdfglyphtounicode{qubuts}{05BB}
\pdfglyphtounicode{qubuts18}{05BB}
\pdfglyphtounicode{qubuts25}{05BB}
\pdfglyphtounicode{qubuts31}{05BB}
\pdfglyphtounicode{qubutshebrew}{05BB}
\pdfglyphtounicode{qubutsnarrowhebrew}{05BB}
\pdfglyphtounicode{qubutsquarterhebrew}{05BB}
\pdfglyphtounicode{qubutswidehebrew}{05BB}
\pdfglyphtounicode{question}{003F}
\pdfglyphtounicode{questionarabic}{061F}
\pdfglyphtounicode{questionarmenian}{055E}
\pdfglyphtounicode{questiondown}{00BF}
\pdfglyphtounicode{questiondownsmall}{00BF}
\pdfglyphtounicode{questiongreek}{037E}
\pdfglyphtounicode{questionmonospace}{FF1F}
\pdfglyphtounicode{questionsmall}{003F}
\pdfglyphtounicode{quotedbl}{0022}
\pdfglyphtounicode{quotedblbase}{201E}
\pdfglyphtounicode{quotedblleft}{201C}
\pdfglyphtounicode{quotedblmonospace}{FF02}
\pdfglyphtounicode{quotedblprime}{301E}
\pdfglyphtounicode{quotedblprimereversed}{301D}
\pdfglyphtounicode{quotedblright}{201D}
\pdfglyphtounicode{quoteleft}{2018}
\pdfglyphtounicode{quoteleftreversed}{201B}
\pdfglyphtounicode{quotereversed}{201B}
\pdfglyphtounicode{quoteright}{2019}
\pdfglyphtounicode{quoterightn}{0149}
\pdfglyphtounicode{quotesinglbase}{201A}
\pdfglyphtounicode{quotesingle}{0027}
\pdfglyphtounicode{quotesinglemonospace}{FF07}
\pdfglyphtounicode{r}{0072}
\pdfglyphtounicode{raarmenian}{057C}
\pdfglyphtounicode{rabengali}{09B0}
\pdfglyphtounicode{racute}{0155}
\pdfglyphtounicode{radeva}{0930}
\pdfglyphtounicode{radical}{221A}
\pdfglyphtounicode{radicalex}{F8E5}
\pdfglyphtounicode{radoverssquare}{33AE}
\pdfglyphtounicode{radoverssquaredsquare}{33AF}
\pdfglyphtounicode{radsquare}{33AD}
\pdfglyphtounicode{rafe}{05BF}
\pdfglyphtounicode{rafehebrew}{05BF}
\pdfglyphtounicode{ragujarati}{0AB0}
\pdfglyphtounicode{ragurmukhi}{0A30}
\pdfglyphtounicode{rahiragana}{3089}
\pdfglyphtounicode{rakatakana}{30E9}
\pdfglyphtounicode{rakatakanahalfwidth}{FF97}
\pdfglyphtounicode{ralowerdiagonalbengali}{09F1}
\pdfglyphtounicode{ramiddlediagonalbengali}{09F0}
\pdfglyphtounicode{ramshorn}{0264}
\pdfglyphtounicode{rangedash}{2013}
\pdfglyphtounicode{ratio}{2236}
\pdfglyphtounicode{rbopomofo}{3116}
\pdfglyphtounicode{rcaron}{0159}
\pdfglyphtounicode{rcedilla}{0157}
\pdfglyphtounicode{rcircle}{24E1}
\pdfglyphtounicode{rcommaaccent}{0157}
\pdfglyphtounicode{rdblgrave}{0211}
\pdfglyphtounicode{rdotaccent}{1E59}
\pdfglyphtounicode{rdotbelow}{1E5B}
\pdfglyphtounicode{rdotbelowmacron}{1E5D}
\pdfglyphtounicode{referencemark}{203B}
\pdfglyphtounicode{reflexsubset}{2286}
\pdfglyphtounicode{reflexsuperset}{2287}
\pdfglyphtounicode{registered}{00AE}
\pdfglyphtounicode{registersans}{00AE}
\pdfglyphtounicode{registerserif}{00AE}
\pdfglyphtounicode{reharabic}{0631}
\pdfglyphtounicode{reharmenian}{0580}
\pdfglyphtounicode{rehfinalarabic}{FEAE}
\pdfglyphtounicode{rehiragana}{308C}
\pdfglyphtounicode{rehyehaleflamarabic}{0631 FEF3 FE8E 0644}
\pdfglyphtounicode{rekatakana}{30EC}
\pdfglyphtounicode{rekatakanahalfwidth}{FF9A}
\pdfglyphtounicode{resh}{05E8}
\pdfglyphtounicode{reshdageshhebrew}{FB48}
\pdfglyphtounicode{reshhatafpatah}{05E8 05B2}
\pdfglyphtounicode{reshhatafpatahhebrew}{05E8 05B2}
\pdfglyphtounicode{reshhatafsegol}{05E8 05B1}
\pdfglyphtounicode{reshhatafsegolhebrew}{05E8 05B1}
\pdfglyphtounicode{reshhebrew}{05E8}
\pdfglyphtounicode{reshhiriq}{05E8 05B4}
\pdfglyphtounicode{reshhiriqhebrew}{05E8 05B4}
\pdfglyphtounicode{reshholam}{05E8 05B9}
\pdfglyphtounicode{reshholamhebrew}{05E8 05B9}
\pdfglyphtounicode{reshpatah}{05E8 05B7}
\pdfglyphtounicode{reshpatahhebrew}{05E8 05B7}
\pdfglyphtounicode{reshqamats}{05E8 05B8}
\pdfglyphtounicode{reshqamatshebrew}{05E8 05B8}
\pdfglyphtounicode{reshqubuts}{05E8 05BB}
\pdfglyphtounicode{reshqubutshebrew}{05E8 05BB}
\pdfglyphtounicode{reshsegol}{05E8 05B6}
\pdfglyphtounicode{reshsegolhebrew}{05E8 05B6}
\pdfglyphtounicode{reshsheva}{05E8 05B0}
\pdfglyphtounicode{reshshevahebrew}{05E8 05B0}
\pdfglyphtounicode{reshtsere}{05E8 05B5}
\pdfglyphtounicode{reshtserehebrew}{05E8 05B5}
\pdfglyphtounicode{revasymptequal}{22CD}
\pdfglyphtounicode{reversedtilde}{223D}
\pdfglyphtounicode{reviahebrew}{0597}
\pdfglyphtounicode{reviamugrashhebrew}{0597}
\pdfglyphtounicode{revlogicalnot}{2310}
\pdfglyphtounicode{revsimilar}{223D}
\pdfglyphtounicode{rfishhook}{027E}
\pdfglyphtounicode{rfishhookreversed}{027F}
\pdfglyphtounicode{rhabengali}{09DD}
\pdfglyphtounicode{rhadeva}{095D}
\pdfglyphtounicode{rho}{03C1}
\pdfglyphtounicode{rho1}{03F1}
\pdfglyphtounicode{rhook}{027D}
\pdfglyphtounicode{rhookturned}{027B}
\pdfglyphtounicode{rhookturnedsuperior}{02B5}
\pdfglyphtounicode{rhosymbolgreek}{03F1}
\pdfglyphtounicode{rhotichookmod}{02DE}
\pdfglyphtounicode{rieulacirclekorean}{3271}
\pdfglyphtounicode{rieulaparenkorean}{3211}
\pdfglyphtounicode{rieulcirclekorean}{3263}
\pdfglyphtounicode{rieulhieuhkorean}{3140}
\pdfglyphtounicode{rieulkiyeokkorean}{313A}
\pdfglyphtounicode{rieulkiyeoksioskorean}{3169}
\pdfglyphtounicode{rieulkorean}{3139}
\pdfglyphtounicode{rieulmieumkorean}{313B}
\pdfglyphtounicode{rieulpansioskorean}{316C}
\pdfglyphtounicode{rieulparenkorean}{3203}
\pdfglyphtounicode{rieulphieuphkorean}{313F}
\pdfglyphtounicode{rieulpieupkorean}{313C}
\pdfglyphtounicode{rieulpieupsioskorean}{316B}
\pdfglyphtounicode{rieulsioskorean}{313D}
\pdfglyphtounicode{rieulthieuthkorean}{313E}
\pdfglyphtounicode{rieultikeutkorean}{316A}
\pdfglyphtounicode{rieulyeorinhieuhkorean}{316D}
\pdfglyphtounicode{rightangle}{221F}
\pdfglyphtounicode{rightanglene}{231D}
\pdfglyphtounicode{rightanglenw}{231C}
\pdfglyphtounicode{rightanglese}{231F}
\pdfglyphtounicode{rightanglesw}{231E}
\pdfglyphtounicode{righttackbelowcmb}{0319}
\pdfglyphtounicode{righttriangle}{22BF}
\pdfglyphtounicode{rihiragana}{308A}
\pdfglyphtounicode{rikatakana}{30EA}
\pdfglyphtounicode{rikatakanahalfwidth}{FF98}
\pdfglyphtounicode{ring}{02DA}
\pdfglyphtounicode{ringbelowcmb}{0325}
\pdfglyphtounicode{ringcmb}{030A}
\pdfglyphtounicode{ringhalfleft}{02BF}
\pdfglyphtounicode{ringhalfleftarmenian}{0559}
\pdfglyphtounicode{ringhalfleftbelowcmb}{031C}
\pdfglyphtounicode{ringhalfleftcentered}{02D3}
\pdfglyphtounicode{ringhalfright}{02BE}
\pdfglyphtounicode{ringhalfrightbelowcmb}{0339}
\pdfglyphtounicode{ringhalfrightcentered}{02D2}
\pdfglyphtounicode{ringinequal}{2256}
\pdfglyphtounicode{rinvertedbreve}{0213}
\pdfglyphtounicode{rittorusquare}{3351}
\pdfglyphtounicode{rlinebelow}{1E5F}
\pdfglyphtounicode{rlongleg}{027C}
\pdfglyphtounicode{rlonglegturned}{027A}
\pdfglyphtounicode{rmonospace}{FF52}
\pdfglyphtounicode{rohiragana}{308D}
\pdfglyphtounicode{rokatakana}{30ED}
\pdfglyphtounicode{rokatakanahalfwidth}{FF9B}
\pdfglyphtounicode{roruathai}{0E23}
\pdfglyphtounicode{rparen}{24AD}
\pdfglyphtounicode{rrabengali}{09DC}
\pdfglyphtounicode{rradeva}{0931}
\pdfglyphtounicode{rragurmukhi}{0A5C}
\pdfglyphtounicode{rreharabic}{0691}
\pdfglyphtounicode{rrehfinalarabic}{FB8D}
\pdfglyphtounicode{rrvocalicbengali}{09E0}
\pdfglyphtounicode{rrvocalicdeva}{0960}
\pdfglyphtounicode{rrvocalicgujarati}{0AE0}
\pdfglyphtounicode{rrvocalicvowelsignbengali}{09C4}
\pdfglyphtounicode{rrvocalicvowelsigndeva}{0944}
\pdfglyphtounicode{rrvocalicvowelsigngujarati}{0AC4}
\pdfglyphtounicode{rsuperior}{0072}
\pdfglyphtounicode{rtblock}{2590}
\pdfglyphtounicode{rturned}{0279}
\pdfglyphtounicode{rturnedsuperior}{02B4}
\pdfglyphtounicode{ruhiragana}{308B}
\pdfglyphtounicode{rukatakana}{30EB}
\pdfglyphtounicode{rukatakanahalfwidth}{FF99}
\pdfglyphtounicode{rupeemarkbengali}{09F2}
\pdfglyphtounicode{rupeesignbengali}{09F3}
\pdfglyphtounicode{rupiah}{20A8}
\pdfglyphtounicode{ruthai}{0E24}
\pdfglyphtounicode{rvocalicbengali}{098B}
\pdfglyphtounicode{rvocalicdeva}{090B}
\pdfglyphtounicode{rvocalicgujarati}{0A8B}
\pdfglyphtounicode{rvocalicvowelsignbengali}{09C3}
\pdfglyphtounicode{rvocalicvowelsigndeva}{0943}
\pdfglyphtounicode{rvocalicvowelsigngujarati}{0AC3}
\pdfglyphtounicode{s}{0073}
\pdfglyphtounicode{sabengali}{09B8}
\pdfglyphtounicode{sacute}{015B}
\pdfglyphtounicode{sacutedotaccent}{1E65}
\pdfglyphtounicode{sadarabic}{0635}
\pdfglyphtounicode{sadeva}{0938}
\pdfglyphtounicode{sadfinalarabic}{FEBA}
\pdfglyphtounicode{sadinitialarabic}{FEBB}
\pdfglyphtounicode{sadmedialarabic}{FEBC}
\pdfglyphtounicode{sagujarati}{0AB8}
\pdfglyphtounicode{sagurmukhi}{0A38}
\pdfglyphtounicode{sahiragana}{3055}
\pdfglyphtounicode{sakatakana}{30B5}
\pdfglyphtounicode{sakatakanahalfwidth}{FF7B}
\pdfglyphtounicode{sallallahoualayhewasallamarabic}{FDFA}
\pdfglyphtounicode{samekh}{05E1}
\pdfglyphtounicode{samekhdagesh}{FB41}
\pdfglyphtounicode{samekhdageshhebrew}{FB41}
\pdfglyphtounicode{samekhhebrew}{05E1}
\pdfglyphtounicode{saraaathai}{0E32}
\pdfglyphtounicode{saraaethai}{0E41}
\pdfglyphtounicode{saraaimaimalaithai}{0E44}
\pdfglyphtounicode{saraaimaimuanthai}{0E43}
\pdfglyphtounicode{saraamthai}{0E33}
\pdfglyphtounicode{saraathai}{0E30}
\pdfglyphtounicode{saraethai}{0E40}
\pdfglyphtounicode{saraiileftthai}{F886}
\pdfglyphtounicode{saraiithai}{0E35}
\pdfglyphtounicode{saraileftthai}{F885}
\pdfglyphtounicode{saraithai}{0E34}
\pdfglyphtounicode{saraothai}{0E42}
\pdfglyphtounicode{saraueeleftthai}{F888}
\pdfglyphtounicode{saraueethai}{0E37}
\pdfglyphtounicode{saraueleftthai}{F887}
\pdfglyphtounicode{sarauethai}{0E36}
\pdfglyphtounicode{sarauthai}{0E38}
\pdfglyphtounicode{sarauuthai}{0E39}
\pdfglyphtounicode{satisfies}{22A8}
\pdfglyphtounicode{sbopomofo}{3119}
\pdfglyphtounicode{scaron}{0161}
\pdfglyphtounicode{scarondotaccent}{1E67}
\pdfglyphtounicode{scedilla}{015F}
\pdfglyphtounicode{schwa}{0259}
\pdfglyphtounicode{schwacyrillic}{04D9}
\pdfglyphtounicode{schwadieresiscyrillic}{04DB}
\pdfglyphtounicode{schwahook}{025A}
\pdfglyphtounicode{scircle}{24E2}
\pdfglyphtounicode{scircumflex}{015D}
\pdfglyphtounicode{scommaaccent}{0219}
\pdfglyphtounicode{sdotaccent}{1E61}
\pdfglyphtounicode{sdotbelow}{1E63}
\pdfglyphtounicode{sdotbelowdotaccent}{1E69}
\pdfglyphtounicode{seagullbelowcmb}{033C}
\pdfglyphtounicode{second}{2033}
\pdfglyphtounicode{secondtonechinese}{02CA}
\pdfglyphtounicode{section}{00A7}
\pdfglyphtounicode{seenarabic}{0633}
\pdfglyphtounicode{seenfinalarabic}{FEB2}
\pdfglyphtounicode{seeninitialarabic}{FEB3}
\pdfglyphtounicode{seenmedialarabic}{FEB4}
\pdfglyphtounicode{segol}{05B6}
\pdfglyphtounicode{segol13}{05B6}
\pdfglyphtounicode{segol1f}{05B6}
\pdfglyphtounicode{segol2c}{05B6}
\pdfglyphtounicode{segolhebrew}{05B6}
\pdfglyphtounicode{segolnarrowhebrew}{05B6}
\pdfglyphtounicode{segolquarterhebrew}{05B6}
\pdfglyphtounicode{segoltahebrew}{0592}
\pdfglyphtounicode{segolwidehebrew}{05B6}
\pdfglyphtounicode{seharmenian}{057D}
\pdfglyphtounicode{sehiragana}{305B}
\pdfglyphtounicode{sekatakana}{30BB}
\pdfglyphtounicode{sekatakanahalfwidth}{FF7E}
\pdfglyphtounicode{semicolon}{003B}
\pdfglyphtounicode{semicolonarabic}{061B}
\pdfglyphtounicode{semicolonmonospace}{FF1B}
\pdfglyphtounicode{semicolonsmall}{FE54}
\pdfglyphtounicode{semivoicedmarkkana}{309C}
\pdfglyphtounicode{semivoicedmarkkanahalfwidth}{FF9F}
\pdfglyphtounicode{sentisquare}{3322}
\pdfglyphtounicode{sentosquare}{3323}
\pdfglyphtounicode{seven}{0037}
\pdfglyphtounicode{sevenarabic}{0667}
\pdfglyphtounicode{sevenbengali}{09ED}
\pdfglyphtounicode{sevencircle}{2466}
\pdfglyphtounicode{sevencircleinversesansserif}{2790}
\pdfglyphtounicode{sevendeva}{096D}
\pdfglyphtounicode{seveneighths}{215E}
\pdfglyphtounicode{sevengujarati}{0AED}
\pdfglyphtounicode{sevengurmukhi}{0A6D}
\pdfglyphtounicode{sevenhackarabic}{0667}
\pdfglyphtounicode{sevenhangzhou}{3027}
\pdfglyphtounicode{sevenideographicparen}{3226}
\pdfglyphtounicode{seveninferior}{2087}
\pdfglyphtounicode{sevenmonospace}{FF17}
\pdfglyphtounicode{sevenoldstyle}{0037}
\pdfglyphtounicode{sevenparen}{247A}
\pdfglyphtounicode{sevenperiod}{248E}
\pdfglyphtounicode{sevenpersian}{06F7}
\pdfglyphtounicode{sevenroman}{2176}
\pdfglyphtounicode{sevensuperior}{2077}
\pdfglyphtounicode{seventeencircle}{2470}
\pdfglyphtounicode{seventeenparen}{2484}
\pdfglyphtounicode{seventeenperiod}{2498}
\pdfglyphtounicode{seventhai}{0E57}
\pdfglyphtounicode{sfthyphen}{00AD}
\pdfglyphtounicode{shaarmenian}{0577}
\pdfglyphtounicode{shabengali}{09B6}
\pdfglyphtounicode{shacyrillic}{0448}
\pdfglyphtounicode{shaddaarabic}{0651}
\pdfglyphtounicode{shaddadammaarabic}{FC61}
\pdfglyphtounicode{shaddadammatanarabic}{FC5E}
\pdfglyphtounicode{shaddafathaarabic}{FC60}
\pdfglyphtounicode{shaddafathatanarabic}{0651 064B}
\pdfglyphtounicode{shaddakasraarabic}{FC62}
\pdfglyphtounicode{shaddakasratanarabic}{FC5F}
\pdfglyphtounicode{shade}{2592}
\pdfglyphtounicode{shadedark}{2593}
\pdfglyphtounicode{shadelight}{2591}
\pdfglyphtounicode{shademedium}{2592}
\pdfglyphtounicode{shadeva}{0936}
\pdfglyphtounicode{shagujarati}{0AB6}
\pdfglyphtounicode{shagurmukhi}{0A36}
\pdfglyphtounicode{shalshelethebrew}{0593}
\pdfglyphtounicode{sharp}{266F}
\pdfglyphtounicode{shbopomofo}{3115}
\pdfglyphtounicode{shchacyrillic}{0449}
\pdfglyphtounicode{sheenarabic}{0634}
\pdfglyphtounicode{sheenfinalarabic}{FEB6}
\pdfglyphtounicode{sheeninitialarabic}{FEB7}
\pdfglyphtounicode{sheenmedialarabic}{FEB8}
\pdfglyphtounicode{sheicoptic}{03E3}
\pdfglyphtounicode{sheqel}{20AA}
\pdfglyphtounicode{sheqelhebrew}{20AA}
\pdfglyphtounicode{sheva}{05B0}
\pdfglyphtounicode{sheva115}{05B0}
\pdfglyphtounicode{sheva15}{05B0}
\pdfglyphtounicode{sheva22}{05B0}
\pdfglyphtounicode{sheva2e}{05B0}
\pdfglyphtounicode{shevahebrew}{05B0}
\pdfglyphtounicode{shevanarrowhebrew}{05B0}
\pdfglyphtounicode{shevaquarterhebrew}{05B0}
\pdfglyphtounicode{shevawidehebrew}{05B0}
\pdfglyphtounicode{shhacyrillic}{04BB}
\pdfglyphtounicode{shiftleft}{21B0}
\pdfglyphtounicode{shiftright}{21B1}
\pdfglyphtounicode{shimacoptic}{03ED}
\pdfglyphtounicode{shin}{05E9}
\pdfglyphtounicode{shindagesh}{FB49}
\pdfglyphtounicode{shindageshhebrew}{FB49}
\pdfglyphtounicode{shindageshshindot}{FB2C}
\pdfglyphtounicode{shindageshshindothebrew}{FB2C}
\pdfglyphtounicode{shindageshsindot}{FB2D}
\pdfglyphtounicode{shindageshsindothebrew}{FB2D}
\pdfglyphtounicode{shindothebrew}{05C1}
\pdfglyphtounicode{shinhebrew}{05E9}
\pdfglyphtounicode{shinshindot}{FB2A}
\pdfglyphtounicode{shinshindothebrew}{FB2A}
\pdfglyphtounicode{shinsindot}{FB2B}
\pdfglyphtounicode{shinsindothebrew}{FB2B}
\pdfglyphtounicode{shook}{0282}
\pdfglyphtounicode{sigma}{03C3}
\pdfglyphtounicode{sigma1}{03C2}
\pdfglyphtounicode{sigmafinal}{03C2}
\pdfglyphtounicode{sigmalunatesymbolgreek}{03F2}
\pdfglyphtounicode{sihiragana}{3057}
\pdfglyphtounicode{sikatakana}{30B7}
\pdfglyphtounicode{sikatakanahalfwidth}{FF7C}
\pdfglyphtounicode{siluqhebrew}{05BD}
\pdfglyphtounicode{siluqlefthebrew}{05BD}
\pdfglyphtounicode{similar}{223C}
\pdfglyphtounicode{similarequal}{2243}
\pdfglyphtounicode{sindothebrew}{05C2}
\pdfglyphtounicode{siosacirclekorean}{3274}
\pdfglyphtounicode{siosaparenkorean}{3214}
\pdfglyphtounicode{sioscieuckorean}{317E}
\pdfglyphtounicode{sioscirclekorean}{3266}
\pdfglyphtounicode{sioskiyeokkorean}{317A}
\pdfglyphtounicode{sioskorean}{3145}
\pdfglyphtounicode{siosnieunkorean}{317B}
\pdfglyphtounicode{siosparenkorean}{3206}
\pdfglyphtounicode{siospieupkorean}{317D}
\pdfglyphtounicode{siostikeutkorean}{317C}
\pdfglyphtounicode{six}{0036}
\pdfglyphtounicode{sixarabic}{0666}
\pdfglyphtounicode{sixbengali}{09EC}
\pdfglyphtounicode{sixcircle}{2465}
\pdfglyphtounicode{sixcircleinversesansserif}{278F}
\pdfglyphtounicode{sixdeva}{096C}
\pdfglyphtounicode{sixgujarati}{0AEC}
\pdfglyphtounicode{sixgurmukhi}{0A6C}
\pdfglyphtounicode{sixhackarabic}{0666}
\pdfglyphtounicode{sixhangzhou}{3026}
\pdfglyphtounicode{sixideographicparen}{3225}
\pdfglyphtounicode{sixinferior}{2086}
\pdfglyphtounicode{sixmonospace}{FF16}
\pdfglyphtounicode{sixoldstyle}{0036}
\pdfglyphtounicode{sixparen}{2479}
\pdfglyphtounicode{sixperiod}{248D}
\pdfglyphtounicode{sixpersian}{06F6}
\pdfglyphtounicode{sixroman}{2175}
\pdfglyphtounicode{sixsuperior}{2076}
\pdfglyphtounicode{sixteencircle}{246F}
\pdfglyphtounicode{sixteencurrencydenominatorbengali}{09F9}
\pdfglyphtounicode{sixteenparen}{2483}
\pdfglyphtounicode{sixteenperiod}{2497}
\pdfglyphtounicode{sixthai}{0E56}
\pdfglyphtounicode{slash}{002F}
\pdfglyphtounicode{slashmonospace}{FF0F}
\pdfglyphtounicode{slong}{017F}
\pdfglyphtounicode{slongdotaccent}{1E9B}
\pdfglyphtounicode{slurabove}{2322}
\pdfglyphtounicode{slurbelow}{2323}
\pdfglyphtounicode{smile}{2323}
\pdfglyphtounicode{smileface}{263A}
\pdfglyphtounicode{smonospace}{FF53}
\pdfglyphtounicode{sofpasuqhebrew}{05C3}
\pdfglyphtounicode{softhyphen}{00AD}
\pdfglyphtounicode{softsigncyrillic}{044C}
\pdfglyphtounicode{sohiragana}{305D}
\pdfglyphtounicode{sokatakana}{30BD}
\pdfglyphtounicode{sokatakanahalfwidth}{FF7F}
\pdfglyphtounicode{soliduslongoverlaycmb}{0338}
\pdfglyphtounicode{solidusshortoverlaycmb}{0337}
\pdfglyphtounicode{sorusithai}{0E29}
\pdfglyphtounicode{sosalathai}{0E28}
\pdfglyphtounicode{sosothai}{0E0B}
\pdfglyphtounicode{sosuathai}{0E2A}
\pdfglyphtounicode{space}{0020}
\pdfglyphtounicode{spacehackarabic}{0020}
\pdfglyphtounicode{spade}{2660}
\pdfglyphtounicode{spadesuitblack}{2660}
\pdfglyphtounicode{spadesuitwhite}{2664}
\pdfglyphtounicode{sparen}{24AE}
\pdfglyphtounicode{sphericalangle}{2222}
\pdfglyphtounicode{square}{25A1}
\pdfglyphtounicode{squarebelowcmb}{033B}
\pdfglyphtounicode{squarecc}{33C4}
\pdfglyphtounicode{squarecm}{339D}
\pdfglyphtounicode{squarediagonalcrosshatchfill}{25A9}
\pdfglyphtounicode{squaredot}{22A1}
\pdfglyphtounicode{squarehorizontalfill}{25A4}
\pdfglyphtounicode{squareimage}{228F}
\pdfglyphtounicode{squarekg}{338F}
\pdfglyphtounicode{squarekm}{339E}
\pdfglyphtounicode{squarekmcapital}{33CE}
\pdfglyphtounicode{squareln}{33D1}
\pdfglyphtounicode{squarelog}{33D2}
\pdfglyphtounicode{squaremg}{338E}
\pdfglyphtounicode{squaremil}{33D5}
\pdfglyphtounicode{squareminus}{229F}
\pdfglyphtounicode{squaremm}{339C}
\pdfglyphtounicode{squaremsquared}{33A1}
\pdfglyphtounicode{squaremultiply}{22A0}
\pdfglyphtounicode{squareoriginal}{2290}
\pdfglyphtounicode{squareorthogonalcrosshatchfill}{25A6}
\pdfglyphtounicode{squareplus}{229E}
\pdfglyphtounicode{squaresolid}{25A0}
\pdfglyphtounicode{squareupperlefttolowerrightfill}{25A7}
\pdfglyphtounicode{squareupperrighttolowerleftfill}{25A8}
\pdfglyphtounicode{squareverticalfill}{25A5}
\pdfglyphtounicode{squarewhitewithsmallblack}{25A3}
\pdfglyphtounicode{squiggleleftright}{21AD}
\pdfglyphtounicode{squiggleright}{21DD}
\pdfglyphtounicode{srsquare}{33DB}
\pdfglyphtounicode{ssabengali}{09B7}
\pdfglyphtounicode{ssadeva}{0937}
\pdfglyphtounicode{ssagujarati}{0AB7}
\pdfglyphtounicode{ssangcieuckorean}{3149}
\pdfglyphtounicode{ssanghieuhkorean}{3185}
\pdfglyphtounicode{ssangieungkorean}{3180}
\pdfglyphtounicode{ssangkiyeokkorean}{3132}
\pdfglyphtounicode{ssangnieunkorean}{3165}
\pdfglyphtounicode{ssangpieupkorean}{3143}
\pdfglyphtounicode{ssangsioskorean}{3146}
\pdfglyphtounicode{ssangtikeutkorean}{3138}
\pdfglyphtounicode{ssuperior}{0073}
\pdfglyphtounicode{st}{0073 0074}
\pdfglyphtounicode{star}{22C6}
\pdfglyphtounicode{sterling}{00A3}
\pdfglyphtounicode{sterlingmonospace}{FFE1}
\pdfglyphtounicode{strokelongoverlaycmb}{0336}
\pdfglyphtounicode{strokeshortoverlaycmb}{0335}
\pdfglyphtounicode{subset}{2282}
\pdfglyphtounicode{subsetdbl}{22D0}
\pdfglyphtounicode{subsetdblequal}{2AC5}
\pdfglyphtounicode{subsetnoteql}{228A}
\pdfglyphtounicode{subsetnotequal}{228A}
\pdfglyphtounicode{subsetorequal}{2286}
\pdfglyphtounicode{subsetornotdbleql}{2ACB}
\pdfglyphtounicode{subsetsqequal}{2291}
\pdfglyphtounicode{succeeds}{227B}
\pdfglyphtounicode{suchthat}{220B}
\pdfglyphtounicode{suhiragana}{3059}
\pdfglyphtounicode{sukatakana}{30B9}
\pdfglyphtounicode{sukatakanahalfwidth}{FF7D}
\pdfglyphtounicode{sukunarabic}{0652}
\pdfglyphtounicode{summation}{2211}
\pdfglyphtounicode{sun}{263C}
\pdfglyphtounicode{superset}{2283}
\pdfglyphtounicode{supersetdbl}{22D1}
\pdfglyphtounicode{supersetdblequal}{2AC6}
\pdfglyphtounicode{supersetnoteql}{228B}
\pdfglyphtounicode{supersetnotequal}{228B}
\pdfglyphtounicode{supersetorequal}{2287}
\pdfglyphtounicode{supersetornotdbleql}{2ACC}
\pdfglyphtounicode{supersetsqequal}{2292}
\pdfglyphtounicode{svsquare}{33DC}
\pdfglyphtounicode{syouwaerasquare}{337C}
\pdfglyphtounicode{t}{0074}
\pdfglyphtounicode{tabengali}{09A4}
\pdfglyphtounicode{tackdown}{22A4}
\pdfglyphtounicode{tackleft}{22A3}
\pdfglyphtounicode{tadeva}{0924}
\pdfglyphtounicode{tagujarati}{0AA4}
\pdfglyphtounicode{tagurmukhi}{0A24}
\pdfglyphtounicode{taharabic}{0637}
\pdfglyphtounicode{tahfinalarabic}{FEC2}
\pdfglyphtounicode{tahinitialarabic}{FEC3}
\pdfglyphtounicode{tahiragana}{305F}
\pdfglyphtounicode{tahmedialarabic}{FEC4}
\pdfglyphtounicode{taisyouerasquare}{337D}
\pdfglyphtounicode{takatakana}{30BF}
\pdfglyphtounicode{takatakanahalfwidth}{FF80}
\pdfglyphtounicode{tatweelarabic}{0640}
\pdfglyphtounicode{tau}{03C4}
\pdfglyphtounicode{tav}{05EA}
\pdfglyphtounicode{tavdages}{FB4A}
\pdfglyphtounicode{tavdagesh}{FB4A}
\pdfglyphtounicode{tavdageshhebrew}{FB4A}
\pdfglyphtounicode{tavhebrew}{05EA}
\pdfglyphtounicode{tbar}{0167}
\pdfglyphtounicode{tbopomofo}{310A}
\pdfglyphtounicode{tcaron}{0165}
\pdfglyphtounicode{tccurl}{02A8}
\pdfglyphtounicode{tcedilla}{0163}
\pdfglyphtounicode{tcheharabic}{0686}
\pdfglyphtounicode{tchehfinalarabic}{FB7B}
\pdfglyphtounicode{tchehinitialarabic}{FB7C}
\pdfglyphtounicode{tchehmedialarabic}{FB7D}
\pdfglyphtounicode{tchehmeeminitialarabic}{FB7C FEE4}
\pdfglyphtounicode{tcircle}{24E3}
\pdfglyphtounicode{tcircumflexbelow}{1E71}
\pdfglyphtounicode{tcommaaccent}{0163}
\pdfglyphtounicode{tdieresis}{1E97}
\pdfglyphtounicode{tdotaccent}{1E6B}
\pdfglyphtounicode{tdotbelow}{1E6D}
\pdfglyphtounicode{tecyrillic}{0442}
\pdfglyphtounicode{tedescendercyrillic}{04AD}
\pdfglyphtounicode{teharabic}{062A}
\pdfglyphtounicode{tehfinalarabic}{FE96}
\pdfglyphtounicode{tehhahinitialarabic}{FCA2}
\pdfglyphtounicode{tehhahisolatedarabic}{FC0C}
\pdfglyphtounicode{tehinitialarabic}{FE97}
\pdfglyphtounicode{tehiragana}{3066}
\pdfglyphtounicode{tehjeeminitialarabic}{FCA1}
\pdfglyphtounicode{tehjeemisolatedarabic}{FC0B}
\pdfglyphtounicode{tehmarbutaarabic}{0629}
\pdfglyphtounicode{tehmarbutafinalarabic}{FE94}
\pdfglyphtounicode{tehmedialarabic}{FE98}
\pdfglyphtounicode{tehmeeminitialarabic}{FCA4}
\pdfglyphtounicode{tehmeemisolatedarabic}{FC0E}
\pdfglyphtounicode{tehnoonfinalarabic}{FC73}
\pdfglyphtounicode{tekatakana}{30C6}
\pdfglyphtounicode{tekatakanahalfwidth}{FF83}
\pdfglyphtounicode{telephone}{2121}
\pdfglyphtounicode{telephoneblack}{260E}
\pdfglyphtounicode{telishagedolahebrew}{05A0}
\pdfglyphtounicode{telishaqetanahebrew}{05A9}
\pdfglyphtounicode{tencircle}{2469}
\pdfglyphtounicode{tenideographicparen}{3229}
\pdfglyphtounicode{tenparen}{247D}
\pdfglyphtounicode{tenperiod}{2491}
\pdfglyphtounicode{tenroman}{2179}
\pdfglyphtounicode{tesh}{02A7}
\pdfglyphtounicode{tet}{05D8}
\pdfglyphtounicode{tetdagesh}{FB38}
\pdfglyphtounicode{tetdageshhebrew}{FB38}
\pdfglyphtounicode{tethebrew}{05D8}
\pdfglyphtounicode{tetsecyrillic}{04B5}
\pdfglyphtounicode{tevirhebrew}{059B}
\pdfglyphtounicode{tevirlefthebrew}{059B}
\pdfglyphtounicode{tfm:cmbsy10/diamond}{2662}
\pdfglyphtounicode{tfm:cmbsy10/heart}{2661}
\pdfglyphtounicode{tfm:cmbsy5/diamond}{2662}
\pdfglyphtounicode{tfm:cmbsy5/heart}{2661}
\pdfglyphtounicode{tfm:cmbsy6/diamond}{2662}
\pdfglyphtounicode{tfm:cmbsy6/heart}{2661}
\pdfglyphtounicode{tfm:cmbsy7/diamond}{2662}
\pdfglyphtounicode{tfm:cmbsy7/heart}{2661}
\pdfglyphtounicode{tfm:cmbsy8/diamond}{2662}
\pdfglyphtounicode{tfm:cmbsy8/heart}{2661}
\pdfglyphtounicode{tfm:cmbsy9/diamond}{2662}
\pdfglyphtounicode{tfm:cmbsy9/heart}{2661}
\pdfglyphtounicode{tfm:cmmi10/phi}{03D5}
\pdfglyphtounicode{tfm:cmmi10/phi1}{03C6}
\pdfglyphtounicode{tfm:cmmi12/phi}{03D5}
\pdfglyphtounicode{tfm:cmmi12/phi1}{03C6}
\pdfglyphtounicode{tfm:cmmi5/phi}{03D5}
\pdfglyphtounicode{tfm:cmmi5/phi1}{03C6}
\pdfglyphtounicode{tfm:cmmi6/phi}{03D5}
\pdfglyphtounicode{tfm:cmmi6/phi1}{03C6}
\pdfglyphtounicode{tfm:cmmi7/phi}{03D5}
\pdfglyphtounicode{tfm:cmmi7/phi1}{03C6}
\pdfglyphtounicode{tfm:cmmi8/phi}{03D5}
\pdfglyphtounicode{tfm:cmmi8/phi1}{03C6}
\pdfglyphtounicode{tfm:cmmi9/phi}{03D5}
\pdfglyphtounicode{tfm:cmmi9/phi1}{03C6}
\pdfglyphtounicode{tfm:cmmib10/phi}{03D5}
\pdfglyphtounicode{tfm:cmmib10/phi1}{03C6}
\pdfglyphtounicode{tfm:cmmib5/phi}{03D5}
\pdfglyphtounicode{tfm:cmmib5/phi1}{03C6}
\pdfglyphtounicode{tfm:cmmib6/phi}{03D5}
\pdfglyphtounicode{tfm:cmmib6/phi1}{03C6}
\pdfglyphtounicode{tfm:cmmib7/phi}{03D5}
\pdfglyphtounicode{tfm:cmmib7/phi1}{03C6}
\pdfglyphtounicode{tfm:cmmib8/phi}{03D5}
\pdfglyphtounicode{tfm:cmmib8/phi1}{03C6}
\pdfglyphtounicode{tfm:cmmib9/phi}{03D5}
\pdfglyphtounicode{tfm:cmmib9/phi1}{03C6}
\pdfglyphtounicode{tfm:cmsy10/diamond}{2662}
\pdfglyphtounicode{tfm:cmsy10/heart}{2661}
\pdfglyphtounicode{tfm:cmsy5/heart}{2661}
\pdfglyphtounicode{tfm:cmsy6/diamond}{2662}
\pdfglyphtounicode{tfm:cmsy6/heart}{2661}
\pdfglyphtounicode{tfm:cmsy7/diamond}{2662}
\pdfglyphtounicode{tfm:cmsy7/heart}{2661}
\pdfglyphtounicode{tfm:cmsy8/diamond}{2662}
\pdfglyphtounicode{tfm:cmsy8/heart}{2661}
\pdfglyphtounicode{tfm:cmsy9/diamond}{2662}
\pdfglyphtounicode{tfm:cmsy9/heart}{2661}
\pdfglyphtounicode{tfm:eurb10/phi}{03D5}
\pdfglyphtounicode{tfm:eurb10/phi1}{03C6}
\pdfglyphtounicode{tfm:eurb5/phi}{03D5}
\pdfglyphtounicode{tfm:eurb5/phi1}{03C6}
\pdfglyphtounicode{tfm:eurb6/phi}{03D5}
\pdfglyphtounicode{tfm:eurb6/phi1}{03C6}
\pdfglyphtounicode{tfm:eurb7/phi}{03D5}
\pdfglyphtounicode{tfm:eurb7/phi1}{03C6}
\pdfglyphtounicode{tfm:eurb8/phi}{03D5}
\pdfglyphtounicode{tfm:eurb8/phi1}{03C6}
\pdfglyphtounicode{tfm:eurb9/phi}{03D5}
\pdfglyphtounicode{tfm:eurb9/phi1}{03C6}
\pdfglyphtounicode{tfm:eurm10/phi}{03D5}
\pdfglyphtounicode{tfm:eurm10/phi1}{03C6}
\pdfglyphtounicode{tfm:eurm5/phi}{03D5}
\pdfglyphtounicode{tfm:eurm5/phi1}{03C6}
\pdfglyphtounicode{tfm:eurm6/phi}{03D5}
\pdfglyphtounicode{tfm:eurm6/phi1}{03C6}
\pdfglyphtounicode{tfm:eurm7/phi}{03D5}
\pdfglyphtounicode{tfm:eurm7/phi1}{03C6}
\pdfglyphtounicode{tfm:eurm8/phi}{03D5}
\pdfglyphtounicode{tfm:eurm8/phi1}{03C6}
\pdfglyphtounicode{tfm:eurm9/phi}{03D5}
\pdfglyphtounicode{tfm:eurm9/phi1}{03C6}
\pdfglyphtounicode{tfm:fplmbi/phi}{03D5}
\pdfglyphtounicode{tfm:fplmbi/phi1}{03C6}
\pdfglyphtounicode{tfm:fplmri/phi}{03D5}
\pdfglyphtounicode{tfm:fplmri/phi1}{03C6}
\pdfglyphtounicode{tfm:lmbsy10/diamond}{2662}
\pdfglyphtounicode{tfm:lmbsy10/heart}{2661}
\pdfglyphtounicode{tfm:lmbsy5/diamond}{2662}
\pdfglyphtounicode{tfm:lmbsy5/heart}{2661}
\pdfglyphtounicode{tfm:lmbsy7/diamond}{2662}
\pdfglyphtounicode{tfm:lmbsy7/heart}{2661}
\pdfglyphtounicode{tfm:lmmi10/phi}{03D5}
\pdfglyphtounicode{tfm:lmmi10/phi1}{03C6}
\pdfglyphtounicode{tfm:lmmi12/phi}{03D5}
\pdfglyphtounicode{tfm:lmmi12/phi1}{03C6}
\pdfglyphtounicode{tfm:lmmi5/phi}{03D5}
\pdfglyphtounicode{tfm:lmmi5/phi1}{03C6}
\pdfglyphtounicode{tfm:lmmi6/phi}{03D5}
\pdfglyphtounicode{tfm:lmmi6/phi1}{03C6}
\pdfglyphtounicode{tfm:lmmi7/phi}{03D5}
\pdfglyphtounicode{tfm:lmmi7/phi1}{03C6}
\pdfglyphtounicode{tfm:lmmi8/phi}{03D5}
\pdfglyphtounicode{tfm:lmmi8/phi1}{03C6}
\pdfglyphtounicode{tfm:lmmi9/phi}{03D5}
\pdfglyphtounicode{tfm:lmmi9/phi1}{03C6}
\pdfglyphtounicode{tfm:lmmib10/phi}{03D5}
\pdfglyphtounicode{tfm:lmmib10/phi1}{03C6}
\pdfglyphtounicode{tfm:lmmib5/phi}{03D5}
\pdfglyphtounicode{tfm:lmmib5/phi1}{03C6}
\pdfglyphtounicode{tfm:lmmib7/phi}{03D5}
\pdfglyphtounicode{tfm:lmmib7/phi1}{03C6}
\pdfglyphtounicode{tfm:lmsy10/diamond}{2662}
\pdfglyphtounicode{tfm:lmsy10/heart}{2661}
\pdfglyphtounicode{tfm:lmsy5/diamond}{2662}
\pdfglyphtounicode{tfm:lmsy5/heart}{2661}
\pdfglyphtounicode{tfm:lmsy6/diamond}{2662}
\pdfglyphtounicode{tfm:lmsy6/heart}{2661}
\pdfglyphtounicode{tfm:lmsy7/diamond}{2662}
\pdfglyphtounicode{tfm:lmsy7/heart}{2661}
\pdfglyphtounicode{tfm:lmsy8/diamond}{2662}
\pdfglyphtounicode{tfm:lmsy8/heart}{2661}
\pdfglyphtounicode{tfm:lmsy9/diamond}{2662}
\pdfglyphtounicode{tfm:lmsy9/heart}{2661}
\pdfglyphtounicode{tfm:msam10/diamond}{2662}
\pdfglyphtounicode{tfm:msam5/diamond}{2662}
\pdfglyphtounicode{tfm:msam6/diamond}{2662}
\pdfglyphtounicode{tfm:msam7/diamond}{2662}
\pdfglyphtounicode{tfm:msam8/diamond}{2662}
\pdfglyphtounicode{tfm:msam9/diamond}{2662}
\pdfglyphtounicode{tfm:pxbmia/phi}{03D5}
\pdfglyphtounicode{tfm:pxbmia/phi1}{03C6}
\pdfglyphtounicode{tfm:pxbsy/diamond}{2662}
\pdfglyphtounicode{tfm:pxbsy/heart}{2661}
\pdfglyphtounicode{tfm:pxbsya/diamond}{2662}
\pdfglyphtounicode{tfm:pxmia/phi}{03D5}
\pdfglyphtounicode{tfm:pxmia/phi1}{03C6}
\pdfglyphtounicode{tfm:pxsy/diamond}{2662}
\pdfglyphtounicode{tfm:pxsy/heart}{2661}
\pdfglyphtounicode{tfm:pxsya/diamond}{2662}
\pdfglyphtounicode{tfm:pzdr/a1}{2701}
\pdfglyphtounicode{tfm:pzdr/a10}{2721}
\pdfglyphtounicode{tfm:pzdr/a100}{275E}
\pdfglyphtounicode{tfm:pzdr/a101}{2761}
\pdfglyphtounicode{tfm:pzdr/a102}{2762}
\pdfglyphtounicode{tfm:pzdr/a103}{2763}
\pdfglyphtounicode{tfm:pzdr/a104}{2764}
\pdfglyphtounicode{tfm:pzdr/a105}{2710}
\pdfglyphtounicode{tfm:pzdr/a106}{2765}
\pdfglyphtounicode{tfm:pzdr/a107}{2766}
\pdfglyphtounicode{tfm:pzdr/a108}{2767}
\pdfglyphtounicode{tfm:pzdr/a109}{2660}
\pdfglyphtounicode{tfm:pzdr/a11}{261B}
\pdfglyphtounicode{tfm:pzdr/a110}{2665}
\pdfglyphtounicode{tfm:pzdr/a111}{2666}
\pdfglyphtounicode{tfm:pzdr/a112}{2663}
\pdfglyphtounicode{tfm:pzdr/a117}{2709}
\pdfglyphtounicode{tfm:pzdr/a118}{2708}
\pdfglyphtounicode{tfm:pzdr/a119}{2707}
\pdfglyphtounicode{tfm:pzdr/a12}{261E}
\pdfglyphtounicode{tfm:pzdr/a120}{2460}
\pdfglyphtounicode{tfm:pzdr/a121}{2461}
\pdfglyphtounicode{tfm:pzdr/a122}{2462}
\pdfglyphtounicode{tfm:pzdr/a123}{2463}
\pdfglyphtounicode{tfm:pzdr/a124}{2464}
\pdfglyphtounicode{tfm:pzdr/a125}{2465}
\pdfglyphtounicode{tfm:pzdr/a126}{2466}
\pdfglyphtounicode{tfm:pzdr/a127}{2467}
\pdfglyphtounicode{tfm:pzdr/a128}{2468}
\pdfglyphtounicode{tfm:pzdr/a129}{2469}
\pdfglyphtounicode{tfm:pzdr/a13}{270C}
\pdfglyphtounicode{tfm:pzdr/a130}{2776}
\pdfglyphtounicode{tfm:pzdr/a131}{2777}
\pdfglyphtounicode{tfm:pzdr/a132}{2778}
\pdfglyphtounicode{tfm:pzdr/a133}{2779}
\pdfglyphtounicode{tfm:pzdr/a134}{277A}
\pdfglyphtounicode{tfm:pzdr/a135}{277B}
\pdfglyphtounicode{tfm:pzdr/a136}{277C}
\pdfglyphtounicode{tfm:pzdr/a137}{277D}
\pdfglyphtounicode{tfm:pzdr/a138}{277E}
\pdfglyphtounicode{tfm:pzdr/a139}{277F}
\pdfglyphtounicode{tfm:pzdr/a14}{270D}
\pdfglyphtounicode{tfm:pzdr/a140}{2780}
\pdfglyphtounicode{tfm:pzdr/a141}{2781}
\pdfglyphtounicode{tfm:pzdr/a142}{2782}
\pdfglyphtounicode{tfm:pzdr/a143}{2783}
\pdfglyphtounicode{tfm:pzdr/a144}{2784}
\pdfglyphtounicode{tfm:pzdr/a145}{2785}
\pdfglyphtounicode{tfm:pzdr/a146}{2786}
\pdfglyphtounicode{tfm:pzdr/a147}{2787}
\pdfglyphtounicode{tfm:pzdr/a148}{2788}
\pdfglyphtounicode{tfm:pzdr/a149}{2789}
\pdfglyphtounicode{tfm:pzdr/a15}{270E}
\pdfglyphtounicode{tfm:pzdr/a150}{278A}
\pdfglyphtounicode{tfm:pzdr/a151}{278B}
\pdfglyphtounicode{tfm:pzdr/a152}{278C}
\pdfglyphtounicode{tfm:pzdr/a153}{278D}
\pdfglyphtounicode{tfm:pzdr/a154}{278E}
\pdfglyphtounicode{tfm:pzdr/a155}{278F}
\pdfglyphtounicode{tfm:pzdr/a156}{2790}
\pdfglyphtounicode{tfm:pzdr/a157}{2791}
\pdfglyphtounicode{tfm:pzdr/a158}{2792}
\pdfglyphtounicode{tfm:pzdr/a159}{2793}
\pdfglyphtounicode{tfm:pzdr/a16}{270F}
\pdfglyphtounicode{tfm:pzdr/a160}{2794}
\pdfglyphtounicode{tfm:pzdr/a161}{2192}
\pdfglyphtounicode{tfm:pzdr/a162}{27A3}
\pdfglyphtounicode{tfm:pzdr/a163}{2194}
\pdfglyphtounicode{tfm:pzdr/a164}{2195}
\pdfglyphtounicode{tfm:pzdr/a165}{2799}
\pdfglyphtounicode{tfm:pzdr/a166}{279B}
\pdfglyphtounicode{tfm:pzdr/a167}{279C}
\pdfglyphtounicode{tfm:pzdr/a168}{279D}
\pdfglyphtounicode{tfm:pzdr/a169}{279E}
\pdfglyphtounicode{tfm:pzdr/a17}{2711}
\pdfglyphtounicode{tfm:pzdr/a170}{279F}
\pdfglyphtounicode{tfm:pzdr/a171}{27A0}
\pdfglyphtounicode{tfm:pzdr/a172}{27A1}
\pdfglyphtounicode{tfm:pzdr/a173}{27A2}
\pdfglyphtounicode{tfm:pzdr/a174}{27A4}
\pdfglyphtounicode{tfm:pzdr/a175}{27A5}
\pdfglyphtounicode{tfm:pzdr/a176}{27A6}
\pdfglyphtounicode{tfm:pzdr/a177}{27A7}
\pdfglyphtounicode{tfm:pzdr/a178}{27A8}
\pdfglyphtounicode{tfm:pzdr/a179}{27A9}
\pdfglyphtounicode{tfm:pzdr/a18}{2712}
\pdfglyphtounicode{tfm:pzdr/a180}{27AB}
\pdfglyphtounicode{tfm:pzdr/a181}{27AD}
\pdfglyphtounicode{tfm:pzdr/a182}{27AF}
\pdfglyphtounicode{tfm:pzdr/a183}{27B2}
\pdfglyphtounicode{tfm:pzdr/a184}{27B3}
\pdfglyphtounicode{tfm:pzdr/a185}{27B5}
\pdfglyphtounicode{tfm:pzdr/a186}{27B8}
\pdfglyphtounicode{tfm:pzdr/a187}{27BA}
\pdfglyphtounicode{tfm:pzdr/a188}{27BB}
\pdfglyphtounicode{tfm:pzdr/a189}{27BC}
\pdfglyphtounicode{tfm:pzdr/a19}{2713}
\pdfglyphtounicode{tfm:pzdr/a190}{27BD}
\pdfglyphtounicode{tfm:pzdr/a191}{27BE}
\pdfglyphtounicode{tfm:pzdr/a192}{279A}
\pdfglyphtounicode{tfm:pzdr/a193}{27AA}
\pdfglyphtounicode{tfm:pzdr/a194}{27B6}
\pdfglyphtounicode{tfm:pzdr/a195}{27B9}
\pdfglyphtounicode{tfm:pzdr/a196}{2798}
\pdfglyphtounicode{tfm:pzdr/a197}{27B4}
\pdfglyphtounicode{tfm:pzdr/a198}{27B7}
\pdfglyphtounicode{tfm:pzdr/a199}{27AC}
\pdfglyphtounicode{tfm:pzdr/a2}{2702}
\pdfglyphtounicode{tfm:pzdr/a20}{2714}
\pdfglyphtounicode{tfm:pzdr/a200}{27AE}
\pdfglyphtounicode{tfm:pzdr/a201}{27B1}
\pdfglyphtounicode{tfm:pzdr/a202}{2703}
\pdfglyphtounicode{tfm:pzdr/a203}{2750}
\pdfglyphtounicode{tfm:pzdr/a204}{2752}
\pdfglyphtounicode{tfm:pzdr/a205}{276E}
\pdfglyphtounicode{tfm:pzdr/a206}{2770}
\pdfglyphtounicode{tfm:pzdr/a21}{2715}
\pdfglyphtounicode{tfm:pzdr/a22}{2716}
\pdfglyphtounicode{tfm:pzdr/a23}{2717}
\pdfglyphtounicode{tfm:pzdr/a24}{2718}
\pdfglyphtounicode{tfm:pzdr/a25}{2719}
\pdfglyphtounicode{tfm:pzdr/a26}{271A}
\pdfglyphtounicode{tfm:pzdr/a27}{271B}
\pdfglyphtounicode{tfm:pzdr/a28}{271C}
\pdfglyphtounicode{tfm:pzdr/a29}{2722}
\pdfglyphtounicode{tfm:pzdr/a3}{2704}
\pdfglyphtounicode{tfm:pzdr/a30}{2723}
\pdfglyphtounicode{tfm:pzdr/a31}{2724}
\pdfglyphtounicode{tfm:pzdr/a32}{2725}
\pdfglyphtounicode{tfm:pzdr/a33}{2726}
\pdfglyphtounicode{tfm:pzdr/a34}{2727}
\pdfglyphtounicode{tfm:pzdr/a35}{2605}
\pdfglyphtounicode{tfm:pzdr/a36}{2729}
\pdfglyphtounicode{tfm:pzdr/a37}{272A}
\pdfglyphtounicode{tfm:pzdr/a38}{272B}
\pdfglyphtounicode{tfm:pzdr/a39}{272C}
\pdfglyphtounicode{tfm:pzdr/a4}{260E}
\pdfglyphtounicode{tfm:pzdr/a40}{272D}
\pdfglyphtounicode{tfm:pzdr/a41}{272E}
\pdfglyphtounicode{tfm:pzdr/a42}{272F}
\pdfglyphtounicode{tfm:pzdr/a43}{2730}
\pdfglyphtounicode{tfm:pzdr/a44}{2731}
\pdfglyphtounicode{tfm:pzdr/a45}{2732}
\pdfglyphtounicode{tfm:pzdr/a46}{2733}
\pdfglyphtounicode{tfm:pzdr/a47}{2734}
\pdfglyphtounicode{tfm:pzdr/a48}{2735}
\pdfglyphtounicode{tfm:pzdr/a49}{2736}
\pdfglyphtounicode{tfm:pzdr/a5}{2706}
\pdfglyphtounicode{tfm:pzdr/a50}{2737}
\pdfglyphtounicode{tfm:pzdr/a51}{2738}
\pdfglyphtounicode{tfm:pzdr/a52}{2739}
\pdfglyphtounicode{tfm:pzdr/a53}{273A}
\pdfglyphtounicode{tfm:pzdr/a54}{273B}
\pdfglyphtounicode{tfm:pzdr/a55}{273C}
\pdfglyphtounicode{tfm:pzdr/a56}{273D}
\pdfglyphtounicode{tfm:pzdr/a57}{273E}
\pdfglyphtounicode{tfm:pzdr/a58}{273F}
\pdfglyphtounicode{tfm:pzdr/a59}{2740}
\pdfglyphtounicode{tfm:pzdr/a6}{271D}
\pdfglyphtounicode{tfm:pzdr/a60}{2741}
\pdfglyphtounicode{tfm:pzdr/a61}{2742}
\pdfglyphtounicode{tfm:pzdr/a62}{2743}
\pdfglyphtounicode{tfm:pzdr/a63}{2744}
\pdfglyphtounicode{tfm:pzdr/a64}{2745}
\pdfglyphtounicode{tfm:pzdr/a65}{2746}
\pdfglyphtounicode{tfm:pzdr/a66}{2747}
\pdfglyphtounicode{tfm:pzdr/a67}{2748}
\pdfglyphtounicode{tfm:pzdr/a68}{2749}
\pdfglyphtounicode{tfm:pzdr/a69}{274A}
\pdfglyphtounicode{tfm:pzdr/a7}{271E}
\pdfglyphtounicode{tfm:pzdr/a70}{274B}
\pdfglyphtounicode{tfm:pzdr/a71}{25CF}
\pdfglyphtounicode{tfm:pzdr/a72}{274D}
\pdfglyphtounicode{tfm:pzdr/a73}{25A0}
\pdfglyphtounicode{tfm:pzdr/a74}{274F}
\pdfglyphtounicode{tfm:pzdr/a75}{2751}
\pdfglyphtounicode{tfm:pzdr/a76}{25B2}
\pdfglyphtounicode{tfm:pzdr/a77}{25BC}
\pdfglyphtounicode{tfm:pzdr/a78}{25C6}
\pdfglyphtounicode{tfm:pzdr/a79}{2756}
\pdfglyphtounicode{tfm:pzdr/a8}{271F}
\pdfglyphtounicode{tfm:pzdr/a81}{25D7}
\pdfglyphtounicode{tfm:pzdr/a82}{2758}
\pdfglyphtounicode{tfm:pzdr/a83}{2759}
\pdfglyphtounicode{tfm:pzdr/a84}{275A}
\pdfglyphtounicode{tfm:pzdr/a85}{276F}
\pdfglyphtounicode{tfm:pzdr/a86}{2771}
\pdfglyphtounicode{tfm:pzdr/a87}{2772}
\pdfglyphtounicode{tfm:pzdr/a88}{2773}
\pdfglyphtounicode{tfm:pzdr/a89}{2768}
\pdfglyphtounicode{tfm:pzdr/a9}{2720}
\pdfglyphtounicode{tfm:pzdr/a90}{2769}
\pdfglyphtounicode{tfm:pzdr/a91}{276C}
\pdfglyphtounicode{tfm:pzdr/a92}{276D}
\pdfglyphtounicode{tfm:pzdr/a93}{276A}
\pdfglyphtounicode{tfm:pzdr/a94}{276B}
\pdfglyphtounicode{tfm:pzdr/a95}{2774}
\pdfglyphtounicode{tfm:pzdr/a96}{2775}
\pdfglyphtounicode{tfm:pzdr/a97}{275B}
\pdfglyphtounicode{tfm:pzdr/a98}{275C}
\pdfglyphtounicode{tfm:pzdr/a99}{275D}
\pdfglyphtounicode{tfm:rpxbmi/phi}{03D5}
\pdfglyphtounicode{tfm:rpxbmi/phi1}{03C6}
\pdfglyphtounicode{tfm:rpxmi/phi}{03D5}
\pdfglyphtounicode{tfm:rpxmi/phi1}{03C6}
\pdfglyphtounicode{tfm:rpzdr/a1}{2701}
\pdfglyphtounicode{tfm:rpzdr/a10}{2721}
\pdfglyphtounicode{tfm:rpzdr/a100}{275E}
\pdfglyphtounicode{tfm:rpzdr/a101}{2761}
\pdfglyphtounicode{tfm:rpzdr/a102}{2762}
\pdfglyphtounicode{tfm:rpzdr/a103}{2763}
\pdfglyphtounicode{tfm:rpzdr/a104}{2764}
\pdfglyphtounicode{tfm:rpzdr/a105}{2710}
\pdfglyphtounicode{tfm:rpzdr/a106}{2765}
\pdfglyphtounicode{tfm:rpzdr/a107}{2766}
\pdfglyphtounicode{tfm:rpzdr/a108}{2767}
\pdfglyphtounicode{tfm:rpzdr/a109}{2660}
\pdfglyphtounicode{tfm:rpzdr/a11}{261B}
\pdfglyphtounicode{tfm:rpzdr/a110}{2665}
\pdfglyphtounicode{tfm:rpzdr/a111}{2666}
\pdfglyphtounicode{tfm:rpzdr/a112}{2663}
\pdfglyphtounicode{tfm:rpzdr/a117}{2709}
\pdfglyphtounicode{tfm:rpzdr/a118}{2708}
\pdfglyphtounicode{tfm:rpzdr/a119}{2707}
\pdfglyphtounicode{tfm:rpzdr/a12}{261E}
\pdfglyphtounicode{tfm:rpzdr/a120}{2460}
\pdfglyphtounicode{tfm:rpzdr/a121}{2461}
\pdfglyphtounicode{tfm:rpzdr/a122}{2462}
\pdfglyphtounicode{tfm:rpzdr/a123}{2463}
\pdfglyphtounicode{tfm:rpzdr/a124}{2464}
\pdfglyphtounicode{tfm:rpzdr/a125}{2465}
\pdfglyphtounicode{tfm:rpzdr/a126}{2466}
\pdfglyphtounicode{tfm:rpzdr/a127}{2467}
\pdfglyphtounicode{tfm:rpzdr/a128}{2468}
\pdfglyphtounicode{tfm:rpzdr/a129}{2469}
\pdfglyphtounicode{tfm:rpzdr/a13}{270C}
\pdfglyphtounicode{tfm:rpzdr/a130}{2776}
\pdfglyphtounicode{tfm:rpzdr/a131}{2777}
\pdfglyphtounicode{tfm:rpzdr/a132}{2778}
\pdfglyphtounicode{tfm:rpzdr/a133}{2779}
\pdfglyphtounicode{tfm:rpzdr/a134}{277A}
\pdfglyphtounicode{tfm:rpzdr/a135}{277B}
\pdfglyphtounicode{tfm:rpzdr/a136}{277C}
\pdfglyphtounicode{tfm:rpzdr/a137}{277D}
\pdfglyphtounicode{tfm:rpzdr/a138}{277E}
\pdfglyphtounicode{tfm:rpzdr/a139}{277F}
\pdfglyphtounicode{tfm:rpzdr/a14}{270D}
\pdfglyphtounicode{tfm:rpzdr/a140}{2780}
\pdfglyphtounicode{tfm:rpzdr/a141}{2781}
\pdfglyphtounicode{tfm:rpzdr/a142}{2782}
\pdfglyphtounicode{tfm:rpzdr/a143}{2783}
\pdfglyphtounicode{tfm:rpzdr/a144}{2784}
\pdfglyphtounicode{tfm:rpzdr/a145}{2785}
\pdfglyphtounicode{tfm:rpzdr/a146}{2786}
\pdfglyphtounicode{tfm:rpzdr/a147}{2787}
\pdfglyphtounicode{tfm:rpzdr/a148}{2788}
\pdfglyphtounicode{tfm:rpzdr/a149}{2789}
\pdfglyphtounicode{tfm:rpzdr/a15}{270E}
\pdfglyphtounicode{tfm:rpzdr/a150}{278A}
\pdfglyphtounicode{tfm:rpzdr/a151}{278B}
\pdfglyphtounicode{tfm:rpzdr/a152}{278C}
\pdfglyphtounicode{tfm:rpzdr/a153}{278D}
\pdfglyphtounicode{tfm:rpzdr/a154}{278E}
\pdfglyphtounicode{tfm:rpzdr/a155}{278F}
\pdfglyphtounicode{tfm:rpzdr/a156}{2790}
\pdfglyphtounicode{tfm:rpzdr/a157}{2791}
\pdfglyphtounicode{tfm:rpzdr/a158}{2792}
\pdfglyphtounicode{tfm:rpzdr/a159}{2793}
\pdfglyphtounicode{tfm:rpzdr/a16}{270F}
\pdfglyphtounicode{tfm:rpzdr/a160}{2794}
\pdfglyphtounicode{tfm:rpzdr/a161}{2192}
\pdfglyphtounicode{tfm:rpzdr/a162}{27A3}
\pdfglyphtounicode{tfm:rpzdr/a163}{2194}
\pdfglyphtounicode{tfm:rpzdr/a164}{2195}
\pdfglyphtounicode{tfm:rpzdr/a165}{2799}
\pdfglyphtounicode{tfm:rpzdr/a166}{279B}
\pdfglyphtounicode{tfm:rpzdr/a167}{279C}
\pdfglyphtounicode{tfm:rpzdr/a168}{279D}
\pdfglyphtounicode{tfm:rpzdr/a169}{279E}
\pdfglyphtounicode{tfm:rpzdr/a17}{2711}
\pdfglyphtounicode{tfm:rpzdr/a170}{279F}
\pdfglyphtounicode{tfm:rpzdr/a171}{27A0}
\pdfglyphtounicode{tfm:rpzdr/a172}{27A1}
\pdfglyphtounicode{tfm:rpzdr/a173}{27A2}
\pdfglyphtounicode{tfm:rpzdr/a174}{27A4}
\pdfglyphtounicode{tfm:rpzdr/a175}{27A5}
\pdfglyphtounicode{tfm:rpzdr/a176}{27A6}
\pdfglyphtounicode{tfm:rpzdr/a177}{27A7}
\pdfglyphtounicode{tfm:rpzdr/a178}{27A8}
\pdfglyphtounicode{tfm:rpzdr/a179}{27A9}
\pdfglyphtounicode{tfm:rpzdr/a18}{2712}
\pdfglyphtounicode{tfm:rpzdr/a180}{27AB}
\pdfglyphtounicode{tfm:rpzdr/a181}{27AD}
\pdfglyphtounicode{tfm:rpzdr/a182}{27AF}
\pdfglyphtounicode{tfm:rpzdr/a183}{27B2}
\pdfglyphtounicode{tfm:rpzdr/a184}{27B3}
\pdfglyphtounicode{tfm:rpzdr/a185}{27B5}
\pdfglyphtounicode{tfm:rpzdr/a186}{27B8}
\pdfglyphtounicode{tfm:rpzdr/a187}{27BA}
\pdfglyphtounicode{tfm:rpzdr/a188}{27BB}
\pdfglyphtounicode{tfm:rpzdr/a189}{27BC}
\pdfglyphtounicode{tfm:rpzdr/a19}{2713}
\pdfglyphtounicode{tfm:rpzdr/a190}{27BD}
\pdfglyphtounicode{tfm:rpzdr/a191}{27BE}
\pdfglyphtounicode{tfm:rpzdr/a192}{279A}
\pdfglyphtounicode{tfm:rpzdr/a193}{27AA}
\pdfglyphtounicode{tfm:rpzdr/a194}{27B6}
\pdfglyphtounicode{tfm:rpzdr/a195}{27B9}
\pdfglyphtounicode{tfm:rpzdr/a196}{2798}
\pdfglyphtounicode{tfm:rpzdr/a197}{27B4}
\pdfglyphtounicode{tfm:rpzdr/a198}{27B7}
\pdfglyphtounicode{tfm:rpzdr/a199}{27AC}
\pdfglyphtounicode{tfm:rpzdr/a2}{2702}
\pdfglyphtounicode{tfm:rpzdr/a20}{2714}
\pdfglyphtounicode{tfm:rpzdr/a200}{27AE}
\pdfglyphtounicode{tfm:rpzdr/a201}{27B1}
\pdfglyphtounicode{tfm:rpzdr/a202}{2703}
\pdfglyphtounicode{tfm:rpzdr/a203}{2750}
\pdfglyphtounicode{tfm:rpzdr/a204}{2752}
\pdfglyphtounicode{tfm:rpzdr/a205}{276E}
\pdfglyphtounicode{tfm:rpzdr/a206}{2770}
\pdfglyphtounicode{tfm:rpzdr/a21}{2715}
\pdfglyphtounicode{tfm:rpzdr/a22}{2716}
\pdfglyphtounicode{tfm:rpzdr/a23}{2717}
\pdfglyphtounicode{tfm:rpzdr/a24}{2718}
\pdfglyphtounicode{tfm:rpzdr/a25}{2719}
\pdfglyphtounicode{tfm:rpzdr/a26}{271A}
\pdfglyphtounicode{tfm:rpzdr/a27}{271B}
\pdfglyphtounicode{tfm:rpzdr/a28}{271C}
\pdfglyphtounicode{tfm:rpzdr/a29}{2722}
\pdfglyphtounicode{tfm:rpzdr/a3}{2704}
\pdfglyphtounicode{tfm:rpzdr/a30}{2723}
\pdfglyphtounicode{tfm:rpzdr/a31}{2724}
\pdfglyphtounicode{tfm:rpzdr/a32}{2725}
\pdfglyphtounicode{tfm:rpzdr/a33}{2726}
\pdfglyphtounicode{tfm:rpzdr/a34}{2727}
\pdfglyphtounicode{tfm:rpzdr/a35}{2605}
\pdfglyphtounicode{tfm:rpzdr/a36}{2729}
\pdfglyphtounicode{tfm:rpzdr/a37}{272A}
\pdfglyphtounicode{tfm:rpzdr/a38}{272B}
\pdfglyphtounicode{tfm:rpzdr/a39}{272C}
\pdfglyphtounicode{tfm:rpzdr/a4}{260E}
\pdfglyphtounicode{tfm:rpzdr/a40}{272D}
\pdfglyphtounicode{tfm:rpzdr/a41}{272E}
\pdfglyphtounicode{tfm:rpzdr/a42}{272F}
\pdfglyphtounicode{tfm:rpzdr/a43}{2730}
\pdfglyphtounicode{tfm:rpzdr/a44}{2731}
\pdfglyphtounicode{tfm:rpzdr/a45}{2732}
\pdfglyphtounicode{tfm:rpzdr/a46}{2733}
\pdfglyphtounicode{tfm:rpzdr/a47}{2734}
\pdfglyphtounicode{tfm:rpzdr/a48}{2735}
\pdfglyphtounicode{tfm:rpzdr/a49}{2736}
\pdfglyphtounicode{tfm:rpzdr/a5}{2706}
\pdfglyphtounicode{tfm:rpzdr/a50}{2737}
\pdfglyphtounicode{tfm:rpzdr/a51}{2738}
\pdfglyphtounicode{tfm:rpzdr/a52}{2739}
\pdfglyphtounicode{tfm:rpzdr/a53}{273A}
\pdfglyphtounicode{tfm:rpzdr/a54}{273B}
\pdfglyphtounicode{tfm:rpzdr/a55}{273C}
\pdfglyphtounicode{tfm:rpzdr/a56}{273D}
\pdfglyphtounicode{tfm:rpzdr/a57}{273E}
\pdfglyphtounicode{tfm:rpzdr/a58}{273F}
\pdfglyphtounicode{tfm:rpzdr/a59}{2740}
\pdfglyphtounicode{tfm:rpzdr/a6}{271D}
\pdfglyphtounicode{tfm:rpzdr/a60}{2741}
\pdfglyphtounicode{tfm:rpzdr/a61}{2742}
\pdfglyphtounicode{tfm:rpzdr/a62}{2743}
\pdfglyphtounicode{tfm:rpzdr/a63}{2744}
\pdfglyphtounicode{tfm:rpzdr/a64}{2745}
\pdfglyphtounicode{tfm:rpzdr/a65}{2746}
\pdfglyphtounicode{tfm:rpzdr/a66}{2747}
\pdfglyphtounicode{tfm:rpzdr/a67}{2748}
\pdfglyphtounicode{tfm:rpzdr/a68}{2749}
\pdfglyphtounicode{tfm:rpzdr/a69}{274A}
\pdfglyphtounicode{tfm:rpzdr/a7}{271E}
\pdfglyphtounicode{tfm:rpzdr/a70}{274B}
\pdfglyphtounicode{tfm:rpzdr/a71}{25CF}
\pdfglyphtounicode{tfm:rpzdr/a72}{274D}
\pdfglyphtounicode{tfm:rpzdr/a73}{25A0}
\pdfglyphtounicode{tfm:rpzdr/a74}{274F}
\pdfglyphtounicode{tfm:rpzdr/a75}{2751}
\pdfglyphtounicode{tfm:rpzdr/a76}{25B2}
\pdfglyphtounicode{tfm:rpzdr/a77}{25BC}
\pdfglyphtounicode{tfm:rpzdr/a78}{25C6}
\pdfglyphtounicode{tfm:rpzdr/a79}{2756}
\pdfglyphtounicode{tfm:rpzdr/a8}{271F}
\pdfglyphtounicode{tfm:rpzdr/a81}{25D7}
\pdfglyphtounicode{tfm:rpzdr/a82}{2758}
\pdfglyphtounicode{tfm:rpzdr/a83}{2759}
\pdfglyphtounicode{tfm:rpzdr/a84}{275A}
\pdfglyphtounicode{tfm:rpzdr/a85}{276F}
\pdfglyphtounicode{tfm:rpzdr/a86}{2771}
\pdfglyphtounicode{tfm:rpzdr/a87}{2772}
\pdfglyphtounicode{tfm:rpzdr/a88}{2773}
\pdfglyphtounicode{tfm:rpzdr/a89}{2768}
\pdfglyphtounicode{tfm:rpzdr/a9}{2720}
\pdfglyphtounicode{tfm:rpzdr/a90}{2769}
\pdfglyphtounicode{tfm:rpzdr/a91}{276C}
\pdfglyphtounicode{tfm:rpzdr/a92}{276D}
\pdfglyphtounicode{tfm:rpzdr/a93}{276A}
\pdfglyphtounicode{tfm:rpzdr/a94}{276B}
\pdfglyphtounicode{tfm:rpzdr/a95}{2774}
\pdfglyphtounicode{tfm:rpzdr/a96}{2775}
\pdfglyphtounicode{tfm:rpzdr/a97}{275B}
\pdfglyphtounicode{tfm:rpzdr/a98}{275C}
\pdfglyphtounicode{tfm:rpzdr/a99}{275D}
\pdfglyphtounicode{tfm:rtxbmi/phi}{03D5}
\pdfglyphtounicode{tfm:rtxbmi/phi1}{03C6}
\pdfglyphtounicode{tfm:rtxmi/phi}{03D5}
\pdfglyphtounicode{tfm:rtxmi/phi1}{03C6}
\pdfglyphtounicode{tfm:txbmia/phi}{03D5}
\pdfglyphtounicode{tfm:txbmia/phi1}{03C6}
\pdfglyphtounicode{tfm:txbsy/diamond}{2662}
\pdfglyphtounicode{tfm:txbsy/heart}{2661}
\pdfglyphtounicode{tfm:txbsya/diamond}{2662}
\pdfglyphtounicode{tfm:txmia/phi}{03D5}
\pdfglyphtounicode{tfm:txmia/phi1}{03C6}
\pdfglyphtounicode{tfm:txsy/diamond}{2662}
\pdfglyphtounicode{tfm:txsy/heart}{2661}
\pdfglyphtounicode{tfm:txsya/diamond}{2662}
\pdfglyphtounicode{tfm:zd/a1}{2701}
\pdfglyphtounicode{tfm:zd/a10}{2721}
\pdfglyphtounicode{tfm:zd/a100}{275E}
\pdfglyphtounicode{tfm:zd/a101}{2761}
\pdfglyphtounicode{tfm:zd/a102}{2762}
\pdfglyphtounicode{tfm:zd/a103}{2763}
\pdfglyphtounicode{tfm:zd/a104}{2764}
\pdfglyphtounicode{tfm:zd/a105}{2710}
\pdfglyphtounicode{tfm:zd/a106}{2765}
\pdfglyphtounicode{tfm:zd/a107}{2766}
\pdfglyphtounicode{tfm:zd/a108}{2767}
\pdfglyphtounicode{tfm:zd/a109}{2660}
\pdfglyphtounicode{tfm:zd/a11}{261B}
\pdfglyphtounicode{tfm:zd/a110}{2665}
\pdfglyphtounicode{tfm:zd/a111}{2666}
\pdfglyphtounicode{tfm:zd/a112}{2663}
\pdfglyphtounicode{tfm:zd/a117}{2709}
\pdfglyphtounicode{tfm:zd/a118}{2708}
\pdfglyphtounicode{tfm:zd/a119}{2707}
\pdfglyphtounicode{tfm:zd/a12}{261E}
\pdfglyphtounicode{tfm:zd/a120}{2460}
\pdfglyphtounicode{tfm:zd/a121}{2461}
\pdfglyphtounicode{tfm:zd/a122}{2462}
\pdfglyphtounicode{tfm:zd/a123}{2463}
\pdfglyphtounicode{tfm:zd/a124}{2464}
\pdfglyphtounicode{tfm:zd/a125}{2465}
\pdfglyphtounicode{tfm:zd/a126}{2466}
\pdfglyphtounicode{tfm:zd/a127}{2467}
\pdfglyphtounicode{tfm:zd/a128}{2468}
\pdfglyphtounicode{tfm:zd/a129}{2469}
\pdfglyphtounicode{tfm:zd/a13}{270C}
\pdfglyphtounicode{tfm:zd/a130}{2776}
\pdfglyphtounicode{tfm:zd/a131}{2777}
\pdfglyphtounicode{tfm:zd/a132}{2778}
\pdfglyphtounicode{tfm:zd/a133}{2779}
\pdfglyphtounicode{tfm:zd/a134}{277A}
\pdfglyphtounicode{tfm:zd/a135}{277B}
\pdfglyphtounicode{tfm:zd/a136}{277C}
\pdfglyphtounicode{tfm:zd/a137}{277D}
\pdfglyphtounicode{tfm:zd/a138}{277E}
\pdfglyphtounicode{tfm:zd/a139}{277F}
\pdfglyphtounicode{tfm:zd/a14}{270D}
\pdfglyphtounicode{tfm:zd/a140}{2780}
\pdfglyphtounicode{tfm:zd/a141}{2781}
\pdfglyphtounicode{tfm:zd/a142}{2782}
\pdfglyphtounicode{tfm:zd/a143}{2783}
\pdfglyphtounicode{tfm:zd/a144}{2784}
\pdfglyphtounicode{tfm:zd/a145}{2785}
\pdfglyphtounicode{tfm:zd/a146}{2786}
\pdfglyphtounicode{tfm:zd/a147}{2787}
\pdfglyphtounicode{tfm:zd/a148}{2788}
\pdfglyphtounicode{tfm:zd/a149}{2789}
\pdfglyphtounicode{tfm:zd/a15}{270E}
\pdfglyphtounicode{tfm:zd/a150}{278A}
\pdfglyphtounicode{tfm:zd/a151}{278B}
\pdfglyphtounicode{tfm:zd/a152}{278C}
\pdfglyphtounicode{tfm:zd/a153}{278D}
\pdfglyphtounicode{tfm:zd/a154}{278E}
\pdfglyphtounicode{tfm:zd/a155}{278F}
\pdfglyphtounicode{tfm:zd/a156}{2790}
\pdfglyphtounicode{tfm:zd/a157}{2791}
\pdfglyphtounicode{tfm:zd/a158}{2792}
\pdfglyphtounicode{tfm:zd/a159}{2793}
\pdfglyphtounicode{tfm:zd/a16}{270F}
\pdfglyphtounicode{tfm:zd/a160}{2794}
\pdfglyphtounicode{tfm:zd/a161}{2192}
\pdfglyphtounicode{tfm:zd/a162}{27A3}
\pdfglyphtounicode{tfm:zd/a163}{2194}
\pdfglyphtounicode{tfm:zd/a164}{2195}
\pdfglyphtounicode{tfm:zd/a165}{2799}
\pdfglyphtounicode{tfm:zd/a166}{279B}
\pdfglyphtounicode{tfm:zd/a167}{279C}
\pdfglyphtounicode{tfm:zd/a168}{279D}
\pdfglyphtounicode{tfm:zd/a169}{279E}
\pdfglyphtounicode{tfm:zd/a17}{2711}
\pdfglyphtounicode{tfm:zd/a170}{279F}
\pdfglyphtounicode{tfm:zd/a171}{27A0}
\pdfglyphtounicode{tfm:zd/a172}{27A1}
\pdfglyphtounicode{tfm:zd/a173}{27A2}
\pdfglyphtounicode{tfm:zd/a174}{27A4}
\pdfglyphtounicode{tfm:zd/a175}{27A5}
\pdfglyphtounicode{tfm:zd/a176}{27A6}
\pdfglyphtounicode{tfm:zd/a177}{27A7}
\pdfglyphtounicode{tfm:zd/a178}{27A8}
\pdfglyphtounicode{tfm:zd/a179}{27A9}
\pdfglyphtounicode{tfm:zd/a18}{2712}
\pdfglyphtounicode{tfm:zd/a180}{27AB}
\pdfglyphtounicode{tfm:zd/a181}{27AD}
\pdfglyphtounicode{tfm:zd/a182}{27AF}
\pdfglyphtounicode{tfm:zd/a183}{27B2}
\pdfglyphtounicode{tfm:zd/a184}{27B3}
\pdfglyphtounicode{tfm:zd/a185}{27B5}
\pdfglyphtounicode{tfm:zd/a186}{27B8}
\pdfglyphtounicode{tfm:zd/a187}{27BA}
\pdfglyphtounicode{tfm:zd/a188}{27BB}
\pdfglyphtounicode{tfm:zd/a189}{27BC}
\pdfglyphtounicode{tfm:zd/a19}{2713}
\pdfglyphtounicode{tfm:zd/a190}{27BD}
\pdfglyphtounicode{tfm:zd/a191}{27BE}
\pdfglyphtounicode{tfm:zd/a192}{279A}
\pdfglyphtounicode{tfm:zd/a193}{27AA}
\pdfglyphtounicode{tfm:zd/a194}{27B6}
\pdfglyphtounicode{tfm:zd/a195}{27B9}
\pdfglyphtounicode{tfm:zd/a196}{2798}
\pdfglyphtounicode{tfm:zd/a197}{27B4}
\pdfglyphtounicode{tfm:zd/a198}{27B7}
\pdfglyphtounicode{tfm:zd/a199}{27AC}
\pdfglyphtounicode{tfm:zd/a2}{2702}
\pdfglyphtounicode{tfm:zd/a20}{2714}
\pdfglyphtounicode{tfm:zd/a200}{27AE}
\pdfglyphtounicode{tfm:zd/a201}{27B1}
\pdfglyphtounicode{tfm:zd/a202}{2703}
\pdfglyphtounicode{tfm:zd/a203}{2750}
\pdfglyphtounicode{tfm:zd/a204}{2752}
\pdfglyphtounicode{tfm:zd/a205}{276E}
\pdfglyphtounicode{tfm:zd/a206}{2770}
\pdfglyphtounicode{tfm:zd/a21}{2715}
\pdfglyphtounicode{tfm:zd/a22}{2716}
\pdfglyphtounicode{tfm:zd/a23}{2717}
\pdfglyphtounicode{tfm:zd/a24}{2718}
\pdfglyphtounicode{tfm:zd/a25}{2719}
\pdfglyphtounicode{tfm:zd/a26}{271A}
\pdfglyphtounicode{tfm:zd/a27}{271B}
\pdfglyphtounicode{tfm:zd/a28}{271C}
\pdfglyphtounicode{tfm:zd/a29}{2722}
\pdfglyphtounicode{tfm:zd/a3}{2704}
\pdfglyphtounicode{tfm:zd/a30}{2723}
\pdfglyphtounicode{tfm:zd/a31}{2724}
\pdfglyphtounicode{tfm:zd/a32}{2725}
\pdfglyphtounicode{tfm:zd/a33}{2726}
\pdfglyphtounicode{tfm:zd/a34}{2727}
\pdfglyphtounicode{tfm:zd/a35}{2605}
\pdfglyphtounicode{tfm:zd/a36}{2729}
\pdfglyphtounicode{tfm:zd/a37}{272A}
\pdfglyphtounicode{tfm:zd/a38}{272B}
\pdfglyphtounicode{tfm:zd/a39}{272C}
\pdfglyphtounicode{tfm:zd/a4}{260E}
\pdfglyphtounicode{tfm:zd/a40}{272D}
\pdfglyphtounicode{tfm:zd/a41}{272E}
\pdfglyphtounicode{tfm:zd/a42}{272F}
\pdfglyphtounicode{tfm:zd/a43}{2730}
\pdfglyphtounicode{tfm:zd/a44}{2731}
\pdfglyphtounicode{tfm:zd/a45}{2732}
\pdfglyphtounicode{tfm:zd/a46}{2733}
\pdfglyphtounicode{tfm:zd/a47}{2734}
\pdfglyphtounicode{tfm:zd/a48}{2735}
\pdfglyphtounicode{tfm:zd/a49}{2736}
\pdfglyphtounicode{tfm:zd/a5}{2706}
\pdfglyphtounicode{tfm:zd/a50}{2737}
\pdfglyphtounicode{tfm:zd/a51}{2738}
\pdfglyphtounicode{tfm:zd/a52}{2739}
\pdfglyphtounicode{tfm:zd/a53}{273A}
\pdfglyphtounicode{tfm:zd/a54}{273B}
\pdfglyphtounicode{tfm:zd/a55}{273C}
\pdfglyphtounicode{tfm:zd/a56}{273D}
\pdfglyphtounicode{tfm:zd/a57}{273E}
\pdfglyphtounicode{tfm:zd/a58}{273F}
\pdfglyphtounicode{tfm:zd/a59}{2740}
\pdfglyphtounicode{tfm:zd/a6}{271D}
\pdfglyphtounicode{tfm:zd/a60}{2741}
\pdfglyphtounicode{tfm:zd/a61}{2742}
\pdfglyphtounicode{tfm:zd/a62}{2743}
\pdfglyphtounicode{tfm:zd/a63}{2744}
\pdfglyphtounicode{tfm:zd/a64}{2745}
\pdfglyphtounicode{tfm:zd/a65}{2746}
\pdfglyphtounicode{tfm:zd/a66}{2747}
\pdfglyphtounicode{tfm:zd/a67}{2748}
\pdfglyphtounicode{tfm:zd/a68}{2749}
\pdfglyphtounicode{tfm:zd/a69}{274A}
\pdfglyphtounicode{tfm:zd/a7}{271E}
\pdfglyphtounicode{tfm:zd/a70}{274B}
\pdfglyphtounicode{tfm:zd/a71}{25CF}
\pdfglyphtounicode{tfm:zd/a72}{274D}
\pdfglyphtounicode{tfm:zd/a73}{25A0}
\pdfglyphtounicode{tfm:zd/a74}{274F}
\pdfglyphtounicode{tfm:zd/a75}{2751}
\pdfglyphtounicode{tfm:zd/a76}{25B2}
\pdfglyphtounicode{tfm:zd/a77}{25BC}
\pdfglyphtounicode{tfm:zd/a78}{25C6}
\pdfglyphtounicode{tfm:zd/a79}{2756}
\pdfglyphtounicode{tfm:zd/a8}{271F}
\pdfglyphtounicode{tfm:zd/a81}{25D7}
\pdfglyphtounicode{tfm:zd/a82}{2758}
\pdfglyphtounicode{tfm:zd/a83}{2759}
\pdfglyphtounicode{tfm:zd/a84}{275A}
\pdfglyphtounicode{tfm:zd/a85}{276F}
\pdfglyphtounicode{tfm:zd/a86}{2771}
\pdfglyphtounicode{tfm:zd/a87}{2772}
\pdfglyphtounicode{tfm:zd/a88}{2773}
\pdfglyphtounicode{tfm:zd/a89}{2768}
\pdfglyphtounicode{tfm:zd/a9}{2720}
\pdfglyphtounicode{tfm:zd/a90}{2769}
\pdfglyphtounicode{tfm:zd/a91}{276C}
\pdfglyphtounicode{tfm:zd/a92}{276D}
\pdfglyphtounicode{tfm:zd/a93}{276A}
\pdfglyphtounicode{tfm:zd/a94}{276B}
\pdfglyphtounicode{tfm:zd/a95}{2774}
\pdfglyphtounicode{tfm:zd/a96}{2775}
\pdfglyphtounicode{tfm:zd/a97}{275B}
\pdfglyphtounicode{tfm:zd/a98}{275C}
\pdfglyphtounicode{tfm:zd/a99}{275D}
\pdfglyphtounicode{tfm:zpzdr-reversed/a1}{2701}
\pdfglyphtounicode{tfm:zpzdr-reversed/a10}{2721}
\pdfglyphtounicode{tfm:zpzdr-reversed/a100}{275E}
\pdfglyphtounicode{tfm:zpzdr-reversed/a101}{2761}
\pdfglyphtounicode{tfm:zpzdr-reversed/a102}{2762}
\pdfglyphtounicode{tfm:zpzdr-reversed/a103}{2763}
\pdfglyphtounicode{tfm:zpzdr-reversed/a104}{2764}
\pdfglyphtounicode{tfm:zpzdr-reversed/a105}{2710}
\pdfglyphtounicode{tfm:zpzdr-reversed/a106}{2765}
\pdfglyphtounicode{tfm:zpzdr-reversed/a107}{2766}
\pdfglyphtounicode{tfm:zpzdr-reversed/a108}{2767}
\pdfglyphtounicode{tfm:zpzdr-reversed/a109}{2660}
\pdfglyphtounicode{tfm:zpzdr-reversed/a11}{261B}
\pdfglyphtounicode{tfm:zpzdr-reversed/a110}{2665}
\pdfglyphtounicode{tfm:zpzdr-reversed/a111}{2666}
\pdfglyphtounicode{tfm:zpzdr-reversed/a112}{2663}
\pdfglyphtounicode{tfm:zpzdr-reversed/a117}{2709}
\pdfglyphtounicode{tfm:zpzdr-reversed/a118}{2708}
\pdfglyphtounicode{tfm:zpzdr-reversed/a119}{2707}
\pdfglyphtounicode{tfm:zpzdr-reversed/a12}{261E}
\pdfglyphtounicode{tfm:zpzdr-reversed/a120}{2460}
\pdfglyphtounicode{tfm:zpzdr-reversed/a121}{2461}
\pdfglyphtounicode{tfm:zpzdr-reversed/a122}{2462}
\pdfglyphtounicode{tfm:zpzdr-reversed/a123}{2463}
\pdfglyphtounicode{tfm:zpzdr-reversed/a124}{2464}
\pdfglyphtounicode{tfm:zpzdr-reversed/a125}{2465}
\pdfglyphtounicode{tfm:zpzdr-reversed/a126}{2466}
\pdfglyphtounicode{tfm:zpzdr-reversed/a127}{2467}
\pdfglyphtounicode{tfm:zpzdr-reversed/a128}{2468}
\pdfglyphtounicode{tfm:zpzdr-reversed/a129}{2469}
\pdfglyphtounicode{tfm:zpzdr-reversed/a13}{270C}
\pdfglyphtounicode{tfm:zpzdr-reversed/a130}{2776}
\pdfglyphtounicode{tfm:zpzdr-reversed/a131}{2777}
\pdfglyphtounicode{tfm:zpzdr-reversed/a132}{2778}
\pdfglyphtounicode{tfm:zpzdr-reversed/a133}{2779}
\pdfglyphtounicode{tfm:zpzdr-reversed/a134}{277A}
\pdfglyphtounicode{tfm:zpzdr-reversed/a135}{277B}
\pdfglyphtounicode{tfm:zpzdr-reversed/a136}{277C}
\pdfglyphtounicode{tfm:zpzdr-reversed/a137}{277D}
\pdfglyphtounicode{tfm:zpzdr-reversed/a138}{277E}
\pdfglyphtounicode{tfm:zpzdr-reversed/a139}{277F}
\pdfglyphtounicode{tfm:zpzdr-reversed/a14}{270D}
\pdfglyphtounicode{tfm:zpzdr-reversed/a140}{2780}
\pdfglyphtounicode{tfm:zpzdr-reversed/a141}{2781}
\pdfglyphtounicode{tfm:zpzdr-reversed/a142}{2782}
\pdfglyphtounicode{tfm:zpzdr-reversed/a143}{2783}
\pdfglyphtounicode{tfm:zpzdr-reversed/a144}{2784}
\pdfglyphtounicode{tfm:zpzdr-reversed/a145}{2785}
\pdfglyphtounicode{tfm:zpzdr-reversed/a146}{2786}
\pdfglyphtounicode{tfm:zpzdr-reversed/a147}{2787}
\pdfglyphtounicode{tfm:zpzdr-reversed/a148}{2788}
\pdfglyphtounicode{tfm:zpzdr-reversed/a149}{2789}
\pdfglyphtounicode{tfm:zpzdr-reversed/a15}{270E}
\pdfglyphtounicode{tfm:zpzdr-reversed/a150}{278A}
\pdfglyphtounicode{tfm:zpzdr-reversed/a151}{278B}
\pdfglyphtounicode{tfm:zpzdr-reversed/a152}{278C}
\pdfglyphtounicode{tfm:zpzdr-reversed/a153}{278D}
\pdfglyphtounicode{tfm:zpzdr-reversed/a154}{278E}
\pdfglyphtounicode{tfm:zpzdr-reversed/a155}{278F}
\pdfglyphtounicode{tfm:zpzdr-reversed/a156}{2790}
\pdfglyphtounicode{tfm:zpzdr-reversed/a157}{2791}
\pdfglyphtounicode{tfm:zpzdr-reversed/a158}{2792}
\pdfglyphtounicode{tfm:zpzdr-reversed/a159}{2793}
\pdfglyphtounicode{tfm:zpzdr-reversed/a16}{270F}
\pdfglyphtounicode{tfm:zpzdr-reversed/a160}{2794}
\pdfglyphtounicode{tfm:zpzdr-reversed/a161}{2192}
\pdfglyphtounicode{tfm:zpzdr-reversed/a162}{27A3}
\pdfglyphtounicode{tfm:zpzdr-reversed/a163}{2194}
\pdfglyphtounicode{tfm:zpzdr-reversed/a164}{2195}
\pdfglyphtounicode{tfm:zpzdr-reversed/a165}{2799}
\pdfglyphtounicode{tfm:zpzdr-reversed/a166}{279B}
\pdfglyphtounicode{tfm:zpzdr-reversed/a167}{279C}
\pdfglyphtounicode{tfm:zpzdr-reversed/a168}{279D}
\pdfglyphtounicode{tfm:zpzdr-reversed/a169}{279E}
\pdfglyphtounicode{tfm:zpzdr-reversed/a17}{2711}
\pdfglyphtounicode{tfm:zpzdr-reversed/a170}{279F}
\pdfglyphtounicode{tfm:zpzdr-reversed/a171}{27A0}
\pdfglyphtounicode{tfm:zpzdr-reversed/a172}{27A1}
\pdfglyphtounicode{tfm:zpzdr-reversed/a173}{27A2}
\pdfglyphtounicode{tfm:zpzdr-reversed/a174}{27A4}
\pdfglyphtounicode{tfm:zpzdr-reversed/a175}{27A5}
\pdfglyphtounicode{tfm:zpzdr-reversed/a176}{27A6}
\pdfglyphtounicode{tfm:zpzdr-reversed/a177}{27A7}
\pdfglyphtounicode{tfm:zpzdr-reversed/a178}{27A8}
\pdfglyphtounicode{tfm:zpzdr-reversed/a179}{27A9}
\pdfglyphtounicode{tfm:zpzdr-reversed/a18}{2712}
\pdfglyphtounicode{tfm:zpzdr-reversed/a180}{27AB}
\pdfglyphtounicode{tfm:zpzdr-reversed/a181}{27AD}
\pdfglyphtounicode{tfm:zpzdr-reversed/a182}{27AF}
\pdfglyphtounicode{tfm:zpzdr-reversed/a183}{27B2}
\pdfglyphtounicode{tfm:zpzdr-reversed/a184}{27B3}
\pdfglyphtounicode{tfm:zpzdr-reversed/a185}{27B5}
\pdfglyphtounicode{tfm:zpzdr-reversed/a186}{27B8}
\pdfglyphtounicode{tfm:zpzdr-reversed/a187}{27BA}
\pdfglyphtounicode{tfm:zpzdr-reversed/a188}{27BB}
\pdfglyphtounicode{tfm:zpzdr-reversed/a189}{27BC}
\pdfglyphtounicode{tfm:zpzdr-reversed/a19}{2713}
\pdfglyphtounicode{tfm:zpzdr-reversed/a190}{27BD}
\pdfglyphtounicode{tfm:zpzdr-reversed/a191}{27BE}
\pdfglyphtounicode{tfm:zpzdr-reversed/a192}{279A}
\pdfglyphtounicode{tfm:zpzdr-reversed/a193}{27AA}
\pdfglyphtounicode{tfm:zpzdr-reversed/a194}{27B6}
\pdfglyphtounicode{tfm:zpzdr-reversed/a195}{27B9}
\pdfglyphtounicode{tfm:zpzdr-reversed/a196}{2798}
\pdfglyphtounicode{tfm:zpzdr-reversed/a197}{27B4}
\pdfglyphtounicode{tfm:zpzdr-reversed/a198}{27B7}
\pdfglyphtounicode{tfm:zpzdr-reversed/a199}{27AC}
\pdfglyphtounicode{tfm:zpzdr-reversed/a2}{2702}
\pdfglyphtounicode{tfm:zpzdr-reversed/a20}{2714}
\pdfglyphtounicode{tfm:zpzdr-reversed/a200}{27AE}
\pdfglyphtounicode{tfm:zpzdr-reversed/a201}{27B1}
\pdfglyphtounicode{tfm:zpzdr-reversed/a202}{2703}
\pdfglyphtounicode{tfm:zpzdr-reversed/a203}{2750}
\pdfglyphtounicode{tfm:zpzdr-reversed/a204}{2752}
\pdfglyphtounicode{tfm:zpzdr-reversed/a205}{276E}
\pdfglyphtounicode{tfm:zpzdr-reversed/a206}{2770}
\pdfglyphtounicode{tfm:zpzdr-reversed/a21}{2715}
\pdfglyphtounicode{tfm:zpzdr-reversed/a22}{2716}
\pdfglyphtounicode{tfm:zpzdr-reversed/a23}{2717}
\pdfglyphtounicode{tfm:zpzdr-reversed/a24}{2718}
\pdfglyphtounicode{tfm:zpzdr-reversed/a25}{2719}
\pdfglyphtounicode{tfm:zpzdr-reversed/a26}{271A}
\pdfglyphtounicode{tfm:zpzdr-reversed/a27}{271B}
\pdfglyphtounicode{tfm:zpzdr-reversed/a28}{271C}
\pdfglyphtounicode{tfm:zpzdr-reversed/a29}{2722}
\pdfglyphtounicode{tfm:zpzdr-reversed/a3}{2704}
\pdfglyphtounicode{tfm:zpzdr-reversed/a30}{2723}
\pdfglyphtounicode{tfm:zpzdr-reversed/a31}{2724}
\pdfglyphtounicode{tfm:zpzdr-reversed/a32}{2725}
\pdfglyphtounicode{tfm:zpzdr-reversed/a33}{2726}
\pdfglyphtounicode{tfm:zpzdr-reversed/a34}{2727}
\pdfglyphtounicode{tfm:zpzdr-reversed/a35}{2605}
\pdfglyphtounicode{tfm:zpzdr-reversed/a36}{2729}
\pdfglyphtounicode{tfm:zpzdr-reversed/a37}{272A}
\pdfglyphtounicode{tfm:zpzdr-reversed/a38}{272B}
\pdfglyphtounicode{tfm:zpzdr-reversed/a39}{272C}
\pdfglyphtounicode{tfm:zpzdr-reversed/a4}{260E}
\pdfglyphtounicode{tfm:zpzdr-reversed/a40}{272D}
\pdfglyphtounicode{tfm:zpzdr-reversed/a41}{272E}
\pdfglyphtounicode{tfm:zpzdr-reversed/a42}{272F}
\pdfglyphtounicode{tfm:zpzdr-reversed/a43}{2730}
\pdfglyphtounicode{tfm:zpzdr-reversed/a44}{2731}
\pdfglyphtounicode{tfm:zpzdr-reversed/a45}{2732}
\pdfglyphtounicode{tfm:zpzdr-reversed/a46}{2733}
\pdfglyphtounicode{tfm:zpzdr-reversed/a47}{2734}
\pdfglyphtounicode{tfm:zpzdr-reversed/a48}{2735}
\pdfglyphtounicode{tfm:zpzdr-reversed/a49}{2736}
\pdfglyphtounicode{tfm:zpzdr-reversed/a5}{2706}
\pdfglyphtounicode{tfm:zpzdr-reversed/a50}{2737}
\pdfglyphtounicode{tfm:zpzdr-reversed/a51}{2738}
\pdfglyphtounicode{tfm:zpzdr-reversed/a52}{2739}
\pdfglyphtounicode{tfm:zpzdr-reversed/a53}{273A}
\pdfglyphtounicode{tfm:zpzdr-reversed/a54}{273B}
\pdfglyphtounicode{tfm:zpzdr-reversed/a55}{273C}
\pdfglyphtounicode{tfm:zpzdr-reversed/a56}{273D}
\pdfglyphtounicode{tfm:zpzdr-reversed/a57}{273E}
\pdfglyphtounicode{tfm:zpzdr-reversed/a58}{273F}
\pdfglyphtounicode{tfm:zpzdr-reversed/a59}{2740}
\pdfglyphtounicode{tfm:zpzdr-reversed/a6}{271D}
\pdfglyphtounicode{tfm:zpzdr-reversed/a60}{2741}
\pdfglyphtounicode{tfm:zpzdr-reversed/a61}{2742}
\pdfglyphtounicode{tfm:zpzdr-reversed/a62}{2743}
\pdfglyphtounicode{tfm:zpzdr-reversed/a63}{2744}
\pdfglyphtounicode{tfm:zpzdr-reversed/a64}{2745}
\pdfglyphtounicode{tfm:zpzdr-reversed/a65}{2746}
\pdfglyphtounicode{tfm:zpzdr-reversed/a66}{2747}
\pdfglyphtounicode{tfm:zpzdr-reversed/a67}{2748}
\pdfglyphtounicode{tfm:zpzdr-reversed/a68}{2749}
\pdfglyphtounicode{tfm:zpzdr-reversed/a69}{274A}
\pdfglyphtounicode{tfm:zpzdr-reversed/a7}{271E}
\pdfglyphtounicode{tfm:zpzdr-reversed/a70}{274B}
\pdfglyphtounicode{tfm:zpzdr-reversed/a71}{25CF}
\pdfglyphtounicode{tfm:zpzdr-reversed/a72}{274D}
\pdfglyphtounicode{tfm:zpzdr-reversed/a73}{25A0}
\pdfglyphtounicode{tfm:zpzdr-reversed/a74}{274F}
\pdfglyphtounicode{tfm:zpzdr-reversed/a75}{2751}
\pdfglyphtounicode{tfm:zpzdr-reversed/a76}{25B2}
\pdfglyphtounicode{tfm:zpzdr-reversed/a77}{25BC}
\pdfglyphtounicode{tfm:zpzdr-reversed/a78}{25C6}
\pdfglyphtounicode{tfm:zpzdr-reversed/a79}{2756}
\pdfglyphtounicode{tfm:zpzdr-reversed/a8}{271F}
\pdfglyphtounicode{tfm:zpzdr-reversed/a81}{25D7}
\pdfglyphtounicode{tfm:zpzdr-reversed/a82}{2758}
\pdfglyphtounicode{tfm:zpzdr-reversed/a83}{2759}
\pdfglyphtounicode{tfm:zpzdr-reversed/a84}{275A}
\pdfglyphtounicode{tfm:zpzdr-reversed/a85}{276F}
\pdfglyphtounicode{tfm:zpzdr-reversed/a86}{2771}
\pdfglyphtounicode{tfm:zpzdr-reversed/a87}{2772}
\pdfglyphtounicode{tfm:zpzdr-reversed/a88}{2773}
\pdfglyphtounicode{tfm:zpzdr-reversed/a89}{2768}
\pdfglyphtounicode{tfm:zpzdr-reversed/a9}{2720}
\pdfglyphtounicode{tfm:zpzdr-reversed/a90}{2769}
\pdfglyphtounicode{tfm:zpzdr-reversed/a91}{276C}
\pdfglyphtounicode{tfm:zpzdr-reversed/a92}{276D}
\pdfglyphtounicode{tfm:zpzdr-reversed/a93}{276A}
\pdfglyphtounicode{tfm:zpzdr-reversed/a94}{276B}
\pdfglyphtounicode{tfm:zpzdr-reversed/a95}{2774}
\pdfglyphtounicode{tfm:zpzdr-reversed/a96}{2775}
\pdfglyphtounicode{tfm:zpzdr-reversed/a97}{275B}
\pdfglyphtounicode{tfm:zpzdr-reversed/a98}{275C}
\pdfglyphtounicode{tfm:zpzdr-reversed/a99}{275D}
\pdfglyphtounicode{thabengali}{09A5}
\pdfglyphtounicode{thadeva}{0925}
\pdfglyphtounicode{thagujarati}{0AA5}
\pdfglyphtounicode{thagurmukhi}{0A25}
\pdfglyphtounicode{thalarabic}{0630}
\pdfglyphtounicode{thalfinalarabic}{FEAC}
\pdfglyphtounicode{thanthakhatlowleftthai}{F898}
\pdfglyphtounicode{thanthakhatlowrightthai}{F897}
\pdfglyphtounicode{thanthakhatthai}{0E4C}
\pdfglyphtounicode{thanthakhatupperleftthai}{F896}
\pdfglyphtounicode{theharabic}{062B}
\pdfglyphtounicode{thehfinalarabic}{FE9A}
\pdfglyphtounicode{thehinitialarabic}{FE9B}
\pdfglyphtounicode{thehmedialarabic}{FE9C}
\pdfglyphtounicode{thereexists}{2203}
\pdfglyphtounicode{therefore}{2234}
\pdfglyphtounicode{theta}{03B8}
\pdfglyphtounicode{theta1}{03D1}
\pdfglyphtounicode{thetasymbolgreek}{03D1}
\pdfglyphtounicode{thieuthacirclekorean}{3279}
\pdfglyphtounicode{thieuthaparenkorean}{3219}
\pdfglyphtounicode{thieuthcirclekorean}{326B}
\pdfglyphtounicode{thieuthkorean}{314C}
\pdfglyphtounicode{thieuthparenkorean}{320B}
\pdfglyphtounicode{thirteencircle}{246C}
\pdfglyphtounicode{thirteenparen}{2480}
\pdfglyphtounicode{thirteenperiod}{2494}
\pdfglyphtounicode{thonangmonthothai}{0E11}
\pdfglyphtounicode{thook}{01AD}
\pdfglyphtounicode{thophuthaothai}{0E12}
\pdfglyphtounicode{thorn}{00FE}
\pdfglyphtounicode{thothahanthai}{0E17}
\pdfglyphtounicode{thothanthai}{0E10}
\pdfglyphtounicode{thothongthai}{0E18}
\pdfglyphtounicode{thothungthai}{0E16}
\pdfglyphtounicode{thousandcyrillic}{0482}
\pdfglyphtounicode{thousandsseparatorarabic}{066C}
\pdfglyphtounicode{thousandsseparatorpersian}{066C}
\pdfglyphtounicode{three}{0033}
\pdfglyphtounicode{threearabic}{0663}
\pdfglyphtounicode{threebengali}{09E9}
\pdfglyphtounicode{threecircle}{2462}
\pdfglyphtounicode{threecircleinversesansserif}{278C}
\pdfglyphtounicode{threedeva}{0969}
\pdfglyphtounicode{threeeighths}{215C}
\pdfglyphtounicode{threegujarati}{0AE9}
\pdfglyphtounicode{threegurmukhi}{0A69}
\pdfglyphtounicode{threehackarabic}{0663}
\pdfglyphtounicode{threehangzhou}{3023}
\pdfglyphtounicode{threeideographicparen}{3222}
\pdfglyphtounicode{threeinferior}{2083}
\pdfglyphtounicode{threemonospace}{FF13}
\pdfglyphtounicode{threenumeratorbengali}{09F6}
\pdfglyphtounicode{threeoldstyle}{0033}
\pdfglyphtounicode{threeparen}{2476}
\pdfglyphtounicode{threeperiod}{248A}
\pdfglyphtounicode{threepersian}{06F3}
\pdfglyphtounicode{threequarters}{00BE}
\pdfglyphtounicode{threequartersemdash}{F6DE}
\pdfglyphtounicode{threeroman}{2172}
\pdfglyphtounicode{threesuperior}{00B3}
\pdfglyphtounicode{threethai}{0E53}
\pdfglyphtounicode{thzsquare}{3394}
\pdfglyphtounicode{tihiragana}{3061}
\pdfglyphtounicode{tikatakana}{30C1}
\pdfglyphtounicode{tikatakanahalfwidth}{FF81}
\pdfglyphtounicode{tikeutacirclekorean}{3270}
\pdfglyphtounicode{tikeutaparenkorean}{3210}
\pdfglyphtounicode{tikeutcirclekorean}{3262}
\pdfglyphtounicode{tikeutkorean}{3137}
\pdfglyphtounicode{tikeutparenkorean}{3202}
\pdfglyphtounicode{tilde}{02DC}
\pdfglyphtounicode{tildebelowcmb}{0330}
\pdfglyphtounicode{tildecmb}{0303}
\pdfglyphtounicode{tildecomb}{0303}
\pdfglyphtounicode{tildedoublecmb}{0360}
\pdfglyphtounicode{tildeoperator}{223C}
\pdfglyphtounicode{tildeoverlaycmb}{0334}
\pdfglyphtounicode{tildeverticalcmb}{033E}
\pdfglyphtounicode{timescircle}{2297}
\pdfglyphtounicode{tipehahebrew}{0596}
\pdfglyphtounicode{tipehalefthebrew}{0596}
\pdfglyphtounicode{tippigurmukhi}{0A70}
\pdfglyphtounicode{titlocyrilliccmb}{0483}
\pdfglyphtounicode{tiwnarmenian}{057F}
\pdfglyphtounicode{tlinebelow}{1E6F}
\pdfglyphtounicode{tmonospace}{FF54}
\pdfglyphtounicode{toarmenian}{0569}
\pdfglyphtounicode{tohiragana}{3068}
\pdfglyphtounicode{tokatakana}{30C8}
\pdfglyphtounicode{tokatakanahalfwidth}{FF84}
\pdfglyphtounicode{tonebarextrahighmod}{02E5}
\pdfglyphtounicode{tonebarextralowmod}{02E9}
\pdfglyphtounicode{tonebarhighmod}{02E6}
\pdfglyphtounicode{tonebarlowmod}{02E8}
\pdfglyphtounicode{tonebarmidmod}{02E7}
\pdfglyphtounicode{tonefive}{01BD}
\pdfglyphtounicode{tonesix}{0185}
\pdfglyphtounicode{tonetwo}{01A8}
\pdfglyphtounicode{tonos}{0384}
\pdfglyphtounicode{tonsquare}{3327}
\pdfglyphtounicode{topatakthai}{0E0F}
\pdfglyphtounicode{tortoiseshellbracketleft}{3014}
\pdfglyphtounicode{tortoiseshellbracketleftsmall}{FE5D}
\pdfglyphtounicode{tortoiseshellbracketleftvertical}{FE39}
\pdfglyphtounicode{tortoiseshellbracketright}{3015}
\pdfglyphtounicode{tortoiseshellbracketrightsmall}{FE5E}
\pdfglyphtounicode{tortoiseshellbracketrightvertical}{FE3A}
\pdfglyphtounicode{totaothai}{0E15}
\pdfglyphtounicode{tpalatalhook}{01AB}
\pdfglyphtounicode{tparen}{24AF}
\pdfglyphtounicode{trademark}{2122}
\pdfglyphtounicode{trademarksans}{2122}
\pdfglyphtounicode{trademarkserif}{2122}
\pdfglyphtounicode{tretroflexhook}{0288}
\pdfglyphtounicode{triagdn}{25BC}
\pdfglyphtounicode{triaglf}{25C4}
\pdfglyphtounicode{triagrt}{25BA}
\pdfglyphtounicode{triagup}{25B2}
\pdfglyphtounicode{triangle}{25B3}
\pdfglyphtounicode{triangledownsld}{25BC}
\pdfglyphtounicode{triangleinv}{25BD}
\pdfglyphtounicode{triangleleft}{25C1}
\pdfglyphtounicode{triangleleftequal}{22B4}
\pdfglyphtounicode{triangleleftsld}{25C0}
\pdfglyphtounicode{triangleright}{25B7}
\pdfglyphtounicode{trianglerightequal}{22B5}
\pdfglyphtounicode{trianglerightsld}{25B6}
\pdfglyphtounicode{trianglesolid}{25B2}
\pdfglyphtounicode{ts}{02A6}
\pdfglyphtounicode{tsadi}{05E6}
\pdfglyphtounicode{tsadidagesh}{FB46}
\pdfglyphtounicode{tsadidageshhebrew}{FB46}
\pdfglyphtounicode{tsadihebrew}{05E6}
\pdfglyphtounicode{tsecyrillic}{0446}
\pdfglyphtounicode{tsere}{05B5}
\pdfglyphtounicode{tsere12}{05B5}
\pdfglyphtounicode{tsere1e}{05B5}
\pdfglyphtounicode{tsere2b}{05B5}
\pdfglyphtounicode{tserehebrew}{05B5}
\pdfglyphtounicode{tserenarrowhebrew}{05B5}
\pdfglyphtounicode{tserequarterhebrew}{05B5}
\pdfglyphtounicode{tserewidehebrew}{05B5}
\pdfglyphtounicode{tshecyrillic}{045B}
\pdfglyphtounicode{tsuperior}{0074}
\pdfglyphtounicode{ttabengali}{099F}
\pdfglyphtounicode{ttadeva}{091F}
\pdfglyphtounicode{ttagujarati}{0A9F}
\pdfglyphtounicode{ttagurmukhi}{0A1F}
\pdfglyphtounicode{tteharabic}{0679}
\pdfglyphtounicode{ttehfinalarabic}{FB67}
\pdfglyphtounicode{ttehinitialarabic}{FB68}
\pdfglyphtounicode{ttehmedialarabic}{FB69}
\pdfglyphtounicode{tthabengali}{09A0}
\pdfglyphtounicode{tthadeva}{0920}
\pdfglyphtounicode{tthagujarati}{0AA0}
\pdfglyphtounicode{tthagurmukhi}{0A20}
\pdfglyphtounicode{tturned}{0287}
\pdfglyphtounicode{tuhiragana}{3064}
\pdfglyphtounicode{tukatakana}{30C4}
\pdfglyphtounicode{tukatakanahalfwidth}{FF82}
\pdfglyphtounicode{turnstileleft}{22A2}
\pdfglyphtounicode{turnstileright}{22A3}
\pdfglyphtounicode{tusmallhiragana}{3063}
\pdfglyphtounicode{tusmallkatakana}{30C3}
\pdfglyphtounicode{tusmallkatakanahalfwidth}{FF6F}
\pdfglyphtounicode{twelvecircle}{246B}
\pdfglyphtounicode{twelveparen}{247F}
\pdfglyphtounicode{twelveperiod}{2493}
\pdfglyphtounicode{twelveroman}{217B}
\pdfglyphtounicode{twentycircle}{2473}
\pdfglyphtounicode{twentyhangzhou}{5344}
\pdfglyphtounicode{twentyparen}{2487}
\pdfglyphtounicode{twentyperiod}{249B}
\pdfglyphtounicode{two}{0032}
\pdfglyphtounicode{twoarabic}{0662}
\pdfglyphtounicode{twobengali}{09E8}
\pdfglyphtounicode{twocircle}{2461}
\pdfglyphtounicode{twocircleinversesansserif}{278B}
\pdfglyphtounicode{twodeva}{0968}
\pdfglyphtounicode{twodotenleader}{2025}
\pdfglyphtounicode{twodotleader}{2025}
\pdfglyphtounicode{twodotleadervertical}{FE30}
\pdfglyphtounicode{twogujarati}{0AE8}
\pdfglyphtounicode{twogurmukhi}{0A68}
\pdfglyphtounicode{twohackarabic}{0662}
\pdfglyphtounicode{twohangzhou}{3022}
\pdfglyphtounicode{twoideographicparen}{3221}
\pdfglyphtounicode{twoinferior}{2082}
\pdfglyphtounicode{twomonospace}{FF12}
\pdfglyphtounicode{twonumeratorbengali}{09F5}
\pdfglyphtounicode{twooldstyle}{0032}
\pdfglyphtounicode{twoparen}{2475}
\pdfglyphtounicode{twoperiod}{2489}
\pdfglyphtounicode{twopersian}{06F2}
\pdfglyphtounicode{tworoman}{2171}
\pdfglyphtounicode{twostroke}{01BB}
\pdfglyphtounicode{twosuperior}{00B2}
\pdfglyphtounicode{twothai}{0E52}
\pdfglyphtounicode{twothirds}{2154}
\pdfglyphtounicode{u}{0075}
\pdfglyphtounicode{uacute}{00FA}
\pdfglyphtounicode{ubar}{0289}
\pdfglyphtounicode{ubengali}{0989}
\pdfglyphtounicode{ubopomofo}{3128}
\pdfglyphtounicode{ubreve}{016D}
\pdfglyphtounicode{ucaron}{01D4}
\pdfglyphtounicode{ucircle}{24E4}
\pdfglyphtounicode{ucircumflex}{00FB}
\pdfglyphtounicode{ucircumflexbelow}{1E77}
\pdfglyphtounicode{ucyrillic}{0443}
\pdfglyphtounicode{udattadeva}{0951}
\pdfglyphtounicode{udblacute}{0171}
\pdfglyphtounicode{udblgrave}{0215}
\pdfglyphtounicode{udeva}{0909}
\pdfglyphtounicode{udieresis}{00FC}
\pdfglyphtounicode{udieresisacute}{01D8}
\pdfglyphtounicode{udieresisbelow}{1E73}
\pdfglyphtounicode{udieresiscaron}{01DA}
\pdfglyphtounicode{udieresiscyrillic}{04F1}
\pdfglyphtounicode{udieresisgrave}{01DC}
\pdfglyphtounicode{udieresismacron}{01D6}
\pdfglyphtounicode{udotbelow}{1EE5}
\pdfglyphtounicode{ugrave}{00F9}
\pdfglyphtounicode{ugujarati}{0A89}
\pdfglyphtounicode{ugurmukhi}{0A09}
\pdfglyphtounicode{uhiragana}{3046}
\pdfglyphtounicode{uhookabove}{1EE7}
\pdfglyphtounicode{uhorn}{01B0}
\pdfglyphtounicode{uhornacute}{1EE9}
\pdfglyphtounicode{uhorndotbelow}{1EF1}
\pdfglyphtounicode{uhorngrave}{1EEB}
\pdfglyphtounicode{uhornhookabove}{1EED}
\pdfglyphtounicode{uhorntilde}{1EEF}
\pdfglyphtounicode{uhungarumlaut}{0171}
\pdfglyphtounicode{uhungarumlautcyrillic}{04F3}
\pdfglyphtounicode{uinvertedbreve}{0217}
\pdfglyphtounicode{ukatakana}{30A6}
\pdfglyphtounicode{ukatakanahalfwidth}{FF73}
\pdfglyphtounicode{ukcyrillic}{0479}
\pdfglyphtounicode{ukorean}{315C}
\pdfglyphtounicode{umacron}{016B}
\pdfglyphtounicode{umacroncyrillic}{04EF}
\pdfglyphtounicode{umacrondieresis}{1E7B}
\pdfglyphtounicode{umatragurmukhi}{0A41}
\pdfglyphtounicode{umonospace}{FF55}
\pdfglyphtounicode{underscore}{005F}
\pdfglyphtounicode{underscoredbl}{2017}
\pdfglyphtounicode{underscoremonospace}{FF3F}
\pdfglyphtounicode{underscorevertical}{FE33}
\pdfglyphtounicode{underscorewavy}{FE4F}
\pdfglyphtounicode{union}{222A}
\pdfglyphtounicode{uniondbl}{22D3}
\pdfglyphtounicode{unionmulti}{228E}
\pdfglyphtounicode{unionsq}{2294}
\pdfglyphtounicode{universal}{2200}
\pdfglyphtounicode{uogonek}{0173}
\pdfglyphtounicode{uparen}{24B0}
\pdfglyphtounicode{upblock}{2580}
\pdfglyphtounicode{upperdothebrew}{05C4}
\pdfglyphtounicode{uprise}{22CF}
\pdfglyphtounicode{upsilon}{03C5}
\pdfglyphtounicode{upsilondieresis}{03CB}
\pdfglyphtounicode{upsilondieresistonos}{03B0}
\pdfglyphtounicode{upsilonlatin}{028A}
\pdfglyphtounicode{upsilontonos}{03CD}
\pdfglyphtounicode{upslope}{29F8}
\pdfglyphtounicode{uptackbelowcmb}{031D}
\pdfglyphtounicode{uptackmod}{02D4}
\pdfglyphtounicode{uragurmukhi}{0A73}
\pdfglyphtounicode{uring}{016F}
\pdfglyphtounicode{ushortcyrillic}{045E}
\pdfglyphtounicode{usmallhiragana}{3045}
\pdfglyphtounicode{usmallkatakana}{30A5}
\pdfglyphtounicode{usmallkatakanahalfwidth}{FF69}
\pdfglyphtounicode{ustraightcyrillic}{04AF}
\pdfglyphtounicode{ustraightstrokecyrillic}{04B1}
\pdfglyphtounicode{utilde}{0169}
\pdfglyphtounicode{utildeacute}{1E79}
\pdfglyphtounicode{utildebelow}{1E75}
\pdfglyphtounicode{uubengali}{098A}
\pdfglyphtounicode{uudeva}{090A}
\pdfglyphtounicode{uugujarati}{0A8A}
\pdfglyphtounicode{uugurmukhi}{0A0A}
\pdfglyphtounicode{uumatragurmukhi}{0A42}
\pdfglyphtounicode{uuvowelsignbengali}{09C2}
\pdfglyphtounicode{uuvowelsigndeva}{0942}
\pdfglyphtounicode{uuvowelsigngujarati}{0AC2}
\pdfglyphtounicode{uvowelsignbengali}{09C1}
\pdfglyphtounicode{uvowelsigndeva}{0941}
\pdfglyphtounicode{uvowelsigngujarati}{0AC1}
\pdfglyphtounicode{v}{0076}
\pdfglyphtounicode{vadeva}{0935}
\pdfglyphtounicode{vagujarati}{0AB5}
\pdfglyphtounicode{vagurmukhi}{0A35}
\pdfglyphtounicode{vakatakana}{30F7}
\pdfglyphtounicode{vav}{05D5}
\pdfglyphtounicode{vavdagesh}{FB35}
\pdfglyphtounicode{vavdagesh65}{FB35}
\pdfglyphtounicode{vavdageshhebrew}{FB35}
\pdfglyphtounicode{vavhebrew}{05D5}
\pdfglyphtounicode{vavholam}{FB4B}
\pdfglyphtounicode{vavholamhebrew}{FB4B}
\pdfglyphtounicode{vavvavhebrew}{05F0}
\pdfglyphtounicode{vavyodhebrew}{05F1}
\pdfglyphtounicode{vcircle}{24E5}
\pdfglyphtounicode{vdotbelow}{1E7F}
\pdfglyphtounicode{vector}{20D7}
\pdfglyphtounicode{vecyrillic}{0432}
\pdfglyphtounicode{veharabic}{06A4}
\pdfglyphtounicode{vehfinalarabic}{FB6B}
\pdfglyphtounicode{vehinitialarabic}{FB6C}
\pdfglyphtounicode{vehmedialarabic}{FB6D}
\pdfglyphtounicode{vekatakana}{30F9}
\pdfglyphtounicode{venus}{2640}
\pdfglyphtounicode{verticalbar}{007C}
\pdfglyphtounicode{verticallineabovecmb}{030D}
\pdfglyphtounicode{verticallinebelowcmb}{0329}
\pdfglyphtounicode{verticallinelowmod}{02CC}
\pdfglyphtounicode{verticallinemod}{02C8}
\pdfglyphtounicode{vewarmenian}{057E}
\pdfglyphtounicode{vhook}{028B}
\pdfglyphtounicode{vikatakana}{30F8}
\pdfglyphtounicode{viramabengali}{09CD}
\pdfglyphtounicode{viramadeva}{094D}
\pdfglyphtounicode{viramagujarati}{0ACD}
\pdfglyphtounicode{visargabengali}{0983}
\pdfglyphtounicode{visargadeva}{0903}
\pdfglyphtounicode{visargagujarati}{0A83}
\pdfglyphtounicode{visiblespace}{2423}
\pdfglyphtounicode{visualspace}{2423}
\pdfglyphtounicode{vmonospace}{FF56}
\pdfglyphtounicode{voarmenian}{0578}
\pdfglyphtounicode{voicediterationhiragana}{309E}
\pdfglyphtounicode{voicediterationkatakana}{30FE}
\pdfglyphtounicode{voicedmarkkana}{309B}
\pdfglyphtounicode{voicedmarkkanahalfwidth}{FF9E}
\pdfglyphtounicode{vokatakana}{30FA}
\pdfglyphtounicode{vparen}{24B1}
\pdfglyphtounicode{vtilde}{1E7D}
\pdfglyphtounicode{vturned}{028C}
\pdfglyphtounicode{vuhiragana}{3094}
\pdfglyphtounicode{vukatakana}{30F4}
\pdfglyphtounicode{w}{0077}
\pdfglyphtounicode{wacute}{1E83}
\pdfglyphtounicode{waekorean}{3159}
\pdfglyphtounicode{wahiragana}{308F}
\pdfglyphtounicode{wakatakana}{30EF}
\pdfglyphtounicode{wakatakanahalfwidth}{FF9C}
\pdfglyphtounicode{wakorean}{3158}
\pdfglyphtounicode{wasmallhiragana}{308E}
\pdfglyphtounicode{wasmallkatakana}{30EE}
\pdfglyphtounicode{wattosquare}{3357}
\pdfglyphtounicode{wavedash}{301C}
\pdfglyphtounicode{wavyunderscorevertical}{FE34}
\pdfglyphtounicode{wawarabic}{0648}
\pdfglyphtounicode{wawfinalarabic}{FEEE}
\pdfglyphtounicode{wawhamzaabovearabic}{0624}
\pdfglyphtounicode{wawhamzaabovefinalarabic}{FE86}
\pdfglyphtounicode{wbsquare}{33DD}
\pdfglyphtounicode{wcircle}{24E6}
\pdfglyphtounicode{wcircumflex}{0175}
\pdfglyphtounicode{wdieresis}{1E85}
\pdfglyphtounicode{wdotaccent}{1E87}
\pdfglyphtounicode{wdotbelow}{1E89}
\pdfglyphtounicode{wehiragana}{3091}
\pdfglyphtounicode{weierstrass}{2118}
\pdfglyphtounicode{wekatakana}{30F1}
\pdfglyphtounicode{wekorean}{315E}
\pdfglyphtounicode{weokorean}{315D}
\pdfglyphtounicode{wgrave}{1E81}
\pdfglyphtounicode{whitebullet}{25E6}
\pdfglyphtounicode{whitecircle}{25CB}
\pdfglyphtounicode{whitecircleinverse}{25D9}
\pdfglyphtounicode{whitecornerbracketleft}{300E}
\pdfglyphtounicode{whitecornerbracketleftvertical}{FE43}
\pdfglyphtounicode{whitecornerbracketright}{300F}
\pdfglyphtounicode{whitecornerbracketrightvertical}{FE44}
\pdfglyphtounicode{whitediamond}{25C7}
\pdfglyphtounicode{whitediamondcontainingblacksmalldiamond}{25C8}
\pdfglyphtounicode{whitedownpointingsmalltriangle}{25BF}
\pdfglyphtounicode{whitedownpointingtriangle}{25BD}
\pdfglyphtounicode{whiteleftpointingsmalltriangle}{25C3}
\pdfglyphtounicode{whiteleftpointingtriangle}{25C1}
\pdfglyphtounicode{whitelenticularbracketleft}{3016}
\pdfglyphtounicode{whitelenticularbracketright}{3017}
\pdfglyphtounicode{whiterightpointingsmalltriangle}{25B9}
\pdfglyphtounicode{whiterightpointingtriangle}{25B7}
\pdfglyphtounicode{whitesmallsquare}{25AB}
\pdfglyphtounicode{whitesmilingface}{263A}
\pdfglyphtounicode{whitesquare}{25A1}
\pdfglyphtounicode{whitestar}{2606}
\pdfglyphtounicode{whitetelephone}{260F}
\pdfglyphtounicode{whitetortoiseshellbracketleft}{3018}
\pdfglyphtounicode{whitetortoiseshellbracketright}{3019}
\pdfglyphtounicode{whiteuppointingsmalltriangle}{25B5}
\pdfglyphtounicode{whiteuppointingtriangle}{25B3}
\pdfglyphtounicode{wihiragana}{3090}
\pdfglyphtounicode{wikatakana}{30F0}
\pdfglyphtounicode{wikorean}{315F}
\pdfglyphtounicode{wmonospace}{FF57}
\pdfglyphtounicode{wohiragana}{3092}
\pdfglyphtounicode{wokatakana}{30F2}
\pdfglyphtounicode{wokatakanahalfwidth}{FF66}
\pdfglyphtounicode{won}{20A9}
\pdfglyphtounicode{wonmonospace}{FFE6}
\pdfglyphtounicode{wowaenthai}{0E27}
\pdfglyphtounicode{wparen}{24B2}
\pdfglyphtounicode{wreathproduct}{2240}
\pdfglyphtounicode{wring}{1E98}
\pdfglyphtounicode{wsuperior}{02B7}
\pdfglyphtounicode{wturned}{028D}
\pdfglyphtounicode{wynn}{01BF}
\pdfglyphtounicode{x}{0078}
\pdfglyphtounicode{xabovecmb}{033D}
\pdfglyphtounicode{xbopomofo}{3112}
\pdfglyphtounicode{xcircle}{24E7}
\pdfglyphtounicode{xdieresis}{1E8D}
\pdfglyphtounicode{xdotaccent}{1E8B}
\pdfglyphtounicode{xeharmenian}{056D}
\pdfglyphtounicode{xi}{03BE}
\pdfglyphtounicode{xmonospace}{FF58}
\pdfglyphtounicode{xparen}{24B3}
\pdfglyphtounicode{xsuperior}{02E3}
\pdfglyphtounicode{y}{0079}
\pdfglyphtounicode{yaadosquare}{334E}
\pdfglyphtounicode{yabengali}{09AF}
\pdfglyphtounicode{yacute}{00FD}
\pdfglyphtounicode{yadeva}{092F}
\pdfglyphtounicode{yaekorean}{3152}
\pdfglyphtounicode{yagujarati}{0AAF}
\pdfglyphtounicode{yagurmukhi}{0A2F}
\pdfglyphtounicode{yahiragana}{3084}
\pdfglyphtounicode{yakatakana}{30E4}
\pdfglyphtounicode{yakatakanahalfwidth}{FF94}
\pdfglyphtounicode{yakorean}{3151}
\pdfglyphtounicode{yamakkanthai}{0E4E}
\pdfglyphtounicode{yasmallhiragana}{3083}
\pdfglyphtounicode{yasmallkatakana}{30E3}
\pdfglyphtounicode{yasmallkatakanahalfwidth}{FF6C}
\pdfglyphtounicode{yatcyrillic}{0463}
\pdfglyphtounicode{ycircle}{24E8}
\pdfglyphtounicode{ycircumflex}{0177}
\pdfglyphtounicode{ydieresis}{00FF}
\pdfglyphtounicode{ydotaccent}{1E8F}
\pdfglyphtounicode{ydotbelow}{1EF5}
\pdfglyphtounicode{yeharabic}{064A}
\pdfglyphtounicode{yehbarreearabic}{06D2}
\pdfglyphtounicode{yehbarreefinalarabic}{FBAF}
\pdfglyphtounicode{yehfinalarabic}{FEF2}
\pdfglyphtounicode{yehhamzaabovearabic}{0626}
\pdfglyphtounicode{yehhamzaabovefinalarabic}{FE8A}
\pdfglyphtounicode{yehhamzaaboveinitialarabic}{FE8B}
\pdfglyphtounicode{yehhamzaabovemedialarabic}{FE8C}
\pdfglyphtounicode{yehinitialarabic}{FEF3}
\pdfglyphtounicode{yehmedialarabic}{FEF4}
\pdfglyphtounicode{yehmeeminitialarabic}{FCDD}
\pdfglyphtounicode{yehmeemisolatedarabic}{FC58}
\pdfglyphtounicode{yehnoonfinalarabic}{FC94}
\pdfglyphtounicode{yehthreedotsbelowarabic}{06D1}
\pdfglyphtounicode{yekorean}{3156}
\pdfglyphtounicode{yen}{00A5}
\pdfglyphtounicode{yenmonospace}{FFE5}
\pdfglyphtounicode{yeokorean}{3155}
\pdfglyphtounicode{yeorinhieuhkorean}{3186}
\pdfglyphtounicode{yerahbenyomohebrew}{05AA}
\pdfglyphtounicode{yerahbenyomolefthebrew}{05AA}
\pdfglyphtounicode{yericyrillic}{044B}
\pdfglyphtounicode{yerudieresiscyrillic}{04F9}
\pdfglyphtounicode{yesieungkorean}{3181}
\pdfglyphtounicode{yesieungpansioskorean}{3183}
\pdfglyphtounicode{yesieungsioskorean}{3182}
\pdfglyphtounicode{yetivhebrew}{059A}
\pdfglyphtounicode{ygrave}{1EF3}
\pdfglyphtounicode{yhook}{01B4}
\pdfglyphtounicode{yhookabove}{1EF7}
\pdfglyphtounicode{yiarmenian}{0575}
\pdfglyphtounicode{yicyrillic}{0457}
\pdfglyphtounicode{yikorean}{3162}
\pdfglyphtounicode{yinyang}{262F}
\pdfglyphtounicode{yiwnarmenian}{0582}
\pdfglyphtounicode{ymonospace}{FF59}
\pdfglyphtounicode{yod}{05D9}
\pdfglyphtounicode{yoddagesh}{FB39}
\pdfglyphtounicode{yoddageshhebrew}{FB39}
\pdfglyphtounicode{yodhebrew}{05D9}
\pdfglyphtounicode{yodyodhebrew}{05F2}
\pdfglyphtounicode{yodyodpatahhebrew}{FB1F}
\pdfglyphtounicode{yohiragana}{3088}
\pdfglyphtounicode{yoikorean}{3189}
\pdfglyphtounicode{yokatakana}{30E8}
\pdfglyphtounicode{yokatakanahalfwidth}{FF96}
\pdfglyphtounicode{yokorean}{315B}
\pdfglyphtounicode{yosmallhiragana}{3087}
\pdfglyphtounicode{yosmallkatakana}{30E7}
\pdfglyphtounicode{yosmallkatakanahalfwidth}{FF6E}
\pdfglyphtounicode{yotgreek}{03F3}
\pdfglyphtounicode{yoyaekorean}{3188}
\pdfglyphtounicode{yoyakorean}{3187}
\pdfglyphtounicode{yoyakthai}{0E22}
\pdfglyphtounicode{yoyingthai}{0E0D}
\pdfglyphtounicode{yparen}{24B4}
\pdfglyphtounicode{ypogegrammeni}{037A}
\pdfglyphtounicode{ypogegrammenigreekcmb}{0345}
\pdfglyphtounicode{yr}{01A6}
\pdfglyphtounicode{yring}{1E99}
\pdfglyphtounicode{ysuperior}{02B8}
\pdfglyphtounicode{ytilde}{1EF9}
\pdfglyphtounicode{yturned}{028E}
\pdfglyphtounicode{yuhiragana}{3086}
\pdfglyphtounicode{yuikorean}{318C}
\pdfglyphtounicode{yukatakana}{30E6}
\pdfglyphtounicode{yukatakanahalfwidth}{FF95}
\pdfglyphtounicode{yukorean}{3160}
\pdfglyphtounicode{yusbigcyrillic}{046B}
\pdfglyphtounicode{yusbigiotifiedcyrillic}{046D}
\pdfglyphtounicode{yuslittlecyrillic}{0467}
\pdfglyphtounicode{yuslittleiotifiedcyrillic}{0469}
\pdfglyphtounicode{yusmallhiragana}{3085}
\pdfglyphtounicode{yusmallkatakana}{30E5}
\pdfglyphtounicode{yusmallkatakanahalfwidth}{FF6D}
\pdfglyphtounicode{yuyekorean}{318B}
\pdfglyphtounicode{yuyeokorean}{318A}
\pdfglyphtounicode{yyabengali}{09DF}
\pdfglyphtounicode{yyadeva}{095F}
\pdfglyphtounicode{z}{007A}
\pdfglyphtounicode{zaarmenian}{0566}
\pdfglyphtounicode{zacute}{017A}
\pdfglyphtounicode{zadeva}{095B}
\pdfglyphtounicode{zagurmukhi}{0A5B}
\pdfglyphtounicode{zaharabic}{0638}
\pdfglyphtounicode{zahfinalarabic}{FEC6}
\pdfglyphtounicode{zahinitialarabic}{FEC7}
\pdfglyphtounicode{zahiragana}{3056}
\pdfglyphtounicode{zahmedialarabic}{FEC8}
\pdfglyphtounicode{zainarabic}{0632}
\pdfglyphtounicode{zainfinalarabic}{FEB0}
\pdfglyphtounicode{zakatakana}{30B6}
\pdfglyphtounicode{zaqefgadolhebrew}{0595}
\pdfglyphtounicode{zaqefqatanhebrew}{0594}
\pdfglyphtounicode{zarqahebrew}{0598}
\pdfglyphtounicode{zayin}{05D6}
\pdfglyphtounicode{zayindagesh}{FB36}
\pdfglyphtounicode{zayindageshhebrew}{FB36}
\pdfglyphtounicode{zayinhebrew}{05D6}
\pdfglyphtounicode{zbopomofo}{3117}
\pdfglyphtounicode{zcaron}{017E}
\pdfglyphtounicode{zcircle}{24E9}
\pdfglyphtounicode{zcircumflex}{1E91}
\pdfglyphtounicode{zcurl}{0291}
\pdfglyphtounicode{zdot}{017C}
\pdfglyphtounicode{zdotaccent}{017C}
\pdfglyphtounicode{zdotbelow}{1E93}
\pdfglyphtounicode{zecyrillic}{0437}
\pdfglyphtounicode{zedescendercyrillic}{0499}
\pdfglyphtounicode{zedieresiscyrillic}{04DF}
\pdfglyphtounicode{zehiragana}{305C}
\pdfglyphtounicode{zekatakana}{30BC}
\pdfglyphtounicode{zero}{0030}
\pdfglyphtounicode{zeroarabic}{0660}
\pdfglyphtounicode{zerobengali}{09E6}
\pdfglyphtounicode{zerodeva}{0966}
\pdfglyphtounicode{zerogujarati}{0AE6}
\pdfglyphtounicode{zerogurmukhi}{0A66}
\pdfglyphtounicode{zerohackarabic}{0660}
\pdfglyphtounicode{zeroinferior}{2080}
\pdfglyphtounicode{zeromonospace}{FF10}
\pdfglyphtounicode{zerooldstyle}{0030}
\pdfglyphtounicode{zeropersian}{06F0}
\pdfglyphtounicode{zerosuperior}{2070}
\pdfglyphtounicode{zerothai}{0E50}
\pdfglyphtounicode{zerowidthjoiner}{FEFF}
\pdfglyphtounicode{zerowidthnonjoiner}{200C}
\pdfglyphtounicode{zerowidthspace}{200B}
\pdfglyphtounicode{zeta}{03B6}
\pdfglyphtounicode{zhbopomofo}{3113}
\pdfglyphtounicode{zhearmenian}{056A}
\pdfglyphtounicode{zhebrevecyrillic}{04C2}
\pdfglyphtounicode{zhecyrillic}{0436}
\pdfglyphtounicode{zhedescendercyrillic}{0497}
\pdfglyphtounicode{zhedieresiscyrillic}{04DD}
\pdfglyphtounicode{zihiragana}{3058}
\pdfglyphtounicode{zikatakana}{30B8}
\pdfglyphtounicode{zinorhebrew}{05AE}
\pdfglyphtounicode{zlinebelow}{1E95}
\pdfglyphtounicode{zmonospace}{FF5A}
\pdfglyphtounicode{zohiragana}{305E}
\pdfglyphtounicode{zokatakana}{30BE}
\pdfglyphtounicode{zparen}{24B5}
\pdfglyphtounicode{zretroflexhook}{0290}
\pdfglyphtounicode{zstroke}{01B6}
\pdfglyphtounicode{zuhiragana}{305A}
\pdfglyphtounicode{zukatakana}{30BA}
 \pdfgentounicode=1
\usepackage[spanish,es-noquoting,es-noshorthands]{babel}
\usepackage{newtxtext}
\usepackage[protrusion=true,expansion=true,tracking=true]{microtype}
\usepackage[a4paper,inner=2.72cm,outer=2.34cm,top=2.42cm,bottom=2.68cm,
  headheight=31pt,headsep=18pt,footskip=30pt]{geometry}
\usepackage{amsmath,amsthm,amscd,mathtools,mathrsfs,bm}
\usepackage{newtxmath}
% Las jerarquías de títulos, el índice y las tablas emplean escalas
% intermedias. Se mantiene la familia RSFS y se habilita su escalado continuo,
% de modo que la notación matemática conserve el cuerpo tipográfico exacto
% sin sustituciones silenciosas de tamaño.
\DeclareFontFamily{U}{rsfs}{\skewchar\font127}
\DeclareFontShape{U}{rsfs}{m}{n}{%
  <-6> s*[1.0] rsfs5
  <6-8> s*[1.0] rsfs7
  <8-> s*[1.0] rsfs10
}{}
\usepackage{array,booktabs,longtable,tabularx,multirow}
\usepackage{enumitem,graphicx,float,pdfpages,pdflscape}
\usepackage{subcaption}
\usepackage[table]{xcolor}
\pdfinclusioncopyfonts=1
\usepackage{tikz}
\usepackage{tikz-cd}
\usetikzlibrary{arrows.meta,positioning,fit,calc,matrix,shapes.geometric,
  decorations.pathreplacing,decorations.markings,backgrounds,patterns}
\usepackage[most]{tcolorbox}
\usepackage{setspace,csquotes,xurl,fancyhdr,fancyvrb,listings,tocloft,titlesec,needspace,etoolbox}
\usepackage{xstring}
\usepackage{hyperref}
% Los marcadores PDF son parte de la jerarquía científica de la obra.  Esta
% tabla conserva en texto Unicode la notación que aparece en los títulos, en
% lugar de dejar que hyperref suprima silenciosamente flechas y constantes.
\pdfstringdefDisableCommands{%
  \def\({}\def\){}\def\[{}\def\]{}%
  \def\ensuremath#1{#1}%
  \def\mathrm#1{#1}\def\mathbf#1{#1}\def\mathsf#1{#1}%
  \def\mathbb#1{#1}\def\mathcal#1{#1}\def\mathfrak#1{#1}%
  \def\operatorname#1{#1}\def\text#1{#1}\def\bm#1{#1}%
  \def\allowbreak{}%
  \def\frac#1#2{#1/#2}\def\sqrt#1{raíz(#1)}%
  \def\to{→}\def\longrightarrow{→}\def\mapsto{↦}%
  \def\leftrightarrow{↔}\def\Rightarrow{⇒}\def\Longrightarrow{⇒}%
  \def\pi{π}\def\varphi{φ}\def\phi{φ}\def\alpha{α}%
  \def\varepsilon{ε}\def\epsilon{ε}\def\mu{μ}\def\lambda{λ}%
  \def\Omega{Ω}\def\zeta{ζ}\def\partial{∂}%
  \def\ast{*}\def\cdot{·}\def\natural{natural}\def\log{log}%
}
\usepackage[nameinlink,noabbrev]{cleveref}
\crefname{appendix}{apéndice}{apéndices}
\Crefname{appendix}{Apéndice}{Apéndices}
% Referencias descriptivas a resultados del tratado que se transducen en esta
% monografía bajo una numeración propia.  La sustitución evita remitir al lector
% a la numeración editorial de otra obra sin duplicar los resultados.
\makeatletter
\newcommand{\HMTDeclareDescriptor}[2]{%
  \expandafter\gdef\csname hmt@xref@#1\endcsname{#2}}
\newcommand{\HMTIfDescriptorTF}[3]{%
  \ifcsname hmt@xref@#1\endcsname #2\else #3\fi}

\HMTDeclareDescriptor{app:catalogo-semillas-app-rev9}{capítulo dedicado al dominio de semillas APP y al catálogo canónico de emisiones}
\HMTDeclareDescriptor{chap:83-08}{capítulo dedicado a la descomposición genealógica de las cinco regiones orientadas de \(\pi\)}
\HMTDeclareDescriptor{chap:alpha}{capítulo dedicado a la clausura dodecafásica de \(\alpha\)}
\HMTDeclareDescriptor{chap:banda-particula-virtual}{capítulo dedicado a la banda de partícula, la antisección y la virtualidad}
\HMTDeclareDescriptor{chap:bandera-excepcional-acumulativa}{capítulo dedicado a la incidencia dodecafásica y al corredor Witt--Leech--Monstruo}
\HMTDeclareDescriptor{chap:biografia-pi-autonoma}{capítulo dedicado a la construcción del carácter de clausura y a la biografía estructural de \(\pi\)}
\HMTDeclareDescriptor{chap:cinco-ciclos-pi}{capítulo dedicado a la minimalidad de los cierres de las cinco regiones de \(\pi\)}
\HMTDeclareDescriptor{chap:funtor-dimensional-hmt}{capítulo dedicado a las rectas de magnitud y a la realización dimensional del núcleo}
\HMTDeclareDescriptor{chap:integral-genealogica}{capítulo dedicado al álgebra genealógica de caminos y a su representación integral}
\HMTDeclareDescriptor{chap:numero-medicion}{capítulo dedicado a los sistemas cilíndricos, la supervivencia y la construcción del número}
\HMTDeclareDescriptor{chap:orden-determinantal-accion}{capítulo dedicado al orden determinantal de la recta de acción}
\HMTDeclareDescriptor{chap:pantallas-dualidad}{capítulo dedicado a subespacios, cocientes, reducciones de rango y dualidad}
\HMTDeclareDescriptor{chap:semillas-emisor-54}{capítulo dedicado al emisor de índice \(54\) y a sus semillas orientadas}
\HMTDeclareDescriptor{chap:sistema-inverso-continuidad-genealogica}{capítulo dedicado al sistema inverso y a la continuidad genealógica}
\HMTDeclareDescriptor{chap:tpk-definicion}{capítulo dedicado al sistema de transporte y a la evolución de fase}
\HMTDeclareDescriptor{def:banda-dodecafasica-ruta}{definición de la banda dodecafásica de ruta}
\HMTDeclareDescriptor{def:odometro-9-12}{definición del odómetro nonádico--dodecafásico}
\HMTDeclareDescriptor{eq:G-interno}{identidad constitutiva interna de la gravitación}
\HMTDeclareDescriptor{eq:consistencia-proyectiva-medidas-u024}{identidad de consistencia proyectiva de las medidas cilíndricas}
\HMTDeclareDescriptor{eq:continuo-lector-arquimediano}{identidad de cobertura del lector arquimediano}
\HMTDeclareDescriptor{eq:descomposiciones-cinco-regiones-pi}{descomposición de las cinco regiones de \(\pi\)}
\HMTDeclareDescriptor{eq:frontal-doble-proyeccion}{diagrama frontal de la doble proyección}
\HMTDeclareDescriptor{eq:modular-finite-first-law}{primera identidad modular finita con término de defecto}
\HMTDeclareDescriptor{eq:modular-inverse-channels}{identidad de reconstrucción conjunta de los canales de área}
\HMTDeclareDescriptor{eq:realizacion-circular-ct108}{identidad de realización circular del calendario de \(108\) fases}
\HMTDeclareDescriptor{espiral:chap:realizacion-em}{capítulo dedicado a la realización electromagnética restringida}
\HMTDeclareDescriptor{espiral:eq:cociclos-radiales}{identidad de los cociclos radiales del levantamiento trítico}
\HMTDeclareDescriptor{hap:def:entropia-bicapa-rev6}{definición de la entropía bicapa orientada}
\HMTDeclareDescriptor{hmtctx:p12:chap:orbitas-elevacion-r36}{capítulo dedicado a las órbitas y a la elevación al calibre \(R_{36}\)}
\HMTDeclareDescriptor{mdg:chap:torres}{capítulo dedicado a la jerarquía global de torres dimensionales}
\HMTDeclareDescriptor{mdg:sec:cierre-electron-two-way-rev11}{sección dedicada a la clausura electrónica bidireccional}
\HMTDeclareDescriptor{mdg:sec:ley-general-masa-accion-autonoma}{sección dedicada al operador universal de acción}
\HMTDeclareDescriptor{mdg:sec:selector-gravitatorio-rev11}{sección dedicada al selector gravitatorio natural}
\HMTDeclareDescriptor{p12-83-c19-s1:hmtch:eq-celula-trit-completa}{identidad de la célula trítica completa}
\HMTDeclareDescriptor{prop:cierre-antiseccion-ct108}{proposición de clausura de las formas del calendario de \(108\) fases}
\HMTDeclareDescriptor{prop:proyeccion-fase-longitud}{proposición de proyección fase--longitud}
\HMTDeclareDescriptor{rm:ch:cosmo}{capítulo dedicado a las escalas cosmológicas del calendario de \(108\) fases}
\HMTDeclareDescriptor{rm:ch:metrology}{capítulo dedicado a las constantes cuántico--eléctricas, la escala térmica y el sistema de Planck}
\HMTDeclareDescriptor{sec:tres-generadores-concretos-trit}{sección dedicada a los tres generadores concretos de los regímenes tríticos}
\HMTDeclareDescriptor{sec:w12-w24-elevacion}{sección dedicada al diseño residual dodecafásico y a la elevación de hexadas}
\HMTDeclareDescriptor{thm:censo-atlas-81}{teorema de clasificación de las familias y generaciones del atlas nonádico}
\HMTDeclareDescriptor{thm:corredor-hmt-moonshine}{teorema del corredor de Leech al álgebra de Moonshine}
\HMTDeclareDescriptor{thm:dos-realizaciones-fuente-cuadratica-rev3}{teorema de las dos realizaciones de una presentación cuadrática integral}
\HMTDeclareDescriptor{thm:generacion-monodromica-conjunta-rev7}{teorema de filtración conjunta a profundidad finita}
\HMTDeclareDescriptor{thm:homogeneidad-pch}{teorema de homogeneidad de la acción de Palatini--Cartan--Holst}
\HMTDeclareDescriptor{thm:principio-conjugacion-masa-banda}{teorema de conjugación de masa y carácter de banda}
\HMTDeclareDescriptor{thm:tpk-cociente-predictivo-36}{teorema del cociente mínimo de memoria predictiva}

% Copias semánticas, no simples alias: \cref y \Cref son órdenes robustas
% gestionadas por xparse y un \let conservaría el despachador mutable,
% provocando recursión después de \RenewDocumentCommand.
\NewCommandCopy{\HMTOriginalCref}{\cref}
\NewCommandCopy{\HMTOriginalCrefCapital}{\Cref}
\NewCommandCopy{\HMTOriginalEqref}{\eqref}
% Los capítulos transducidos conservan centenares de referencias semánticas
% cuya numeración pertenecía al tratado integral.  Una referencia que ya ha
% recibido etiqueta local conserva su numeración propia; una referencia que
% sólo existe en el volumen rector se expresa por su clase matemática.  De
% este modo el texto sigue siendo legible y autosuficiente, sin signos ``??''
% ni remisiones ficticias a una numeración ausente.
\newcommand{\HMTFallbackReference}[1]{%
  \IfSubStr{#1}{eq:}{la identidad correspondiente}{%
  \IfSubStr{#1}{thm:}{el teorema correspondiente}{%
  \IfSubStr{#1}{theorem:}{el teorema correspondiente}{%
  \IfSubStr{#1}{prop:}{la proposición correspondiente}{%
  \IfSubStr{#1}{lem:}{el lema correspondiente}{%
  \IfSubStr{#1}{cor:}{el corolario correspondiente}{%
  \IfSubStr{#1}{def:}{la definición correspondiente}{%
  \IfSubStr{#1}{fig:}{la figura correspondiente}{%
  \IfSubStr{#1}{tab:}{la tabla correspondiente}{%
  \IfSubStr{#1}{chap:}{el capítulo correspondiente}{%
  \IfSubStr{#1}{ch:}{el capítulo correspondiente}{%
  \IfSubStr{#1}{sec:}{el apartado correspondiente}{%
  \IfSubStr{#1}{crit:}{el criterio correspondiente}{%
  el resultado correspondiente}}}}}}}}}}}}}}
\newcommand{\HMTFallbackReferenceCapital}[1]{%
  \IfSubStr{#1}{eq:}{La identidad correspondiente}{%
  \IfSubStr{#1}{thm:}{El teorema correspondiente}{%
  \IfSubStr{#1}{theorem:}{El teorema correspondiente}{%
  \IfSubStr{#1}{prop:}{La proposición correspondiente}{%
  \IfSubStr{#1}{lem:}{El lema correspondiente}{%
  \IfSubStr{#1}{cor:}{El corolario correspondiente}{%
  \IfSubStr{#1}{def:}{La definición correspondiente}{%
  \IfSubStr{#1}{fig:}{La figura correspondiente}{%
  \IfSubStr{#1}{tab:}{La tabla correspondiente}{%
  \IfSubStr{#1}{chap:}{El capítulo correspondiente}{%
  \IfSubStr{#1}{ch:}{El capítulo correspondiente}{%
  \IfSubStr{#1}{sec:}{El apartado correspondiente}{%
  \IfSubStr{#1}{crit:}{El criterio correspondiente}{%
  El resultado correspondiente}}}}}}}}}}}}}}
\newcommand{\HMTIfLocalLabelTF}[3]{%
  \ifcsname r@#1\endcsname #2\else #3\fi}
\RenewDocumentCommand{\cref}{s m}{%
  \HMTIfDescriptorTF{#2}{\csname hmt@xref@#2\endcsname}{%
    \HMTIfLocalLabelTF{#2}{%
      \IfBooleanTF{#1}{\HMTOriginalCref*{#2}}{\HMTOriginalCref{#2}}%
    }{\HMTFallbackReference{#2}}}}
\RenewDocumentCommand{\Cref}{s m}{%
  \HMTIfDescriptorTF{#2}{\csname hmt@xref@#2\endcsname}{%
    \HMTIfLocalLabelTF{#2}{%
      \IfBooleanTF{#1}{\HMTOriginalCrefCapital*{#2}}{\HMTOriginalCrefCapital{#2}}%
    }{\HMTFallbackReferenceCapital{#2}}}}
\RenewDocumentCommand{\eqref}{m}{%
  \HMTIfDescriptorTF{#1}{\textup{(\csname hmt@xref@#1\endcsname)}}{%
    \HMTIfLocalLabelTF{#1}{\HMTOriginalEqref{#1}}{\textup{(identidad correspondiente)}}}}
\makeatother
 \usepackage[noautomatic,quiet]{imakeidx}
\makeindex
\allowdisplaybreaks

\definecolor{hmtblue}{RGB}{24,57,89}
\definecolor{hmtgold}{RGB}{145,101,26}
\definecolor{hmtgreen}{RGB}{29,104,78}
\definecolor{hmtred}{RGB}{151,54,45}
\definecolor{hmtviolet}{RGB}{92,64,125}
\definecolor{hmtgray}{RGB}{87,95,102}
\definecolor{hmtpaper}{RGB}{248,247,243}
\definecolor{hmtlightblue}{RGB}{232,240,247}
\definecolor{hmtlightgold}{RGB}{249,242,225}
\definecolor{hmtlightgreen}{RGB}{232,244,238}
\definecolor{hmtlightred}{RGB}{250,235,233}
\definecolor{hmtlightviolet}{RGB}{241,236,248}

% Paleta pública común para los cuadros, el atlas y los inventarios
% integrales de Mecánica Dimensional. Los nombres en versales son alias
% tipográficos; no introducen una segunda identidad visual.
\colorlet{HMTNavy}{hmtblue}
\colorlet{HMTBlue}{hmtblue}
\colorlet{HMTGold}{hmtgold}
\colorlet{HMTLight}{hmtlightblue}
\colorlet{HMTGray}{hmtgray}
\colorlet{HMTLine}{hmtlightblue}
\providecommand{\HMTtablehead}[1]{\color{white}\sffamily\bfseries #1}
\providecommand{\RaggedRight}{\raggedright}

\newcommand{\HMT}{\ensuremath{\mathrm{HMT}}}
\providecommand{\gammaH}{\gamma_{\HMT}}
\providecommand{\ii}{\mathrm i}
\providecommand{\C}{\mathbb C}
\providecommand{\Z}{\mathbb Z}
\providecommand{\U}{\mathcal U}
\providecommand{\arsinh}{\operatorname{arsinh}}
\providecommand{\arcosh}{\operatorname{arcosh}}
\providecommand{\HHM}{K_{\mathrm{HHM}}}
\providecommand{\Dmod}{\Delta_{4}^{\mathrm{mod}}}
\providecommand{\Val}{\operatorname{Val}}
\providecommand{\Ext}{\operatorname{Ext}}
\providecommand{\ev}{\operatorname{ev}}
\providecommand{\End}{\operatorname{End}}
\providecommand{\diam}{\operatorname{diam}}
\providecommand{\im}{\operatorname{im}}
\providecommand{\supp}{\operatorname{supp}}
\providecommand{\tr}{\operatorname{tr}}
\providecommand{\Gnine}{\mathcal G_9}
\providecommand{\NTHW}{\mathfrak N_{\mathrm{HW}}}
\providecommand{\readerquestion}[1]{%
  \begin{quote}\small\itshape\color{hmtblue}#1\end{quote}}
\newcommand{\TRIT}{\ensuremath{\widehat{\mathbb T}}}
\newcommand{\AAPP}{\mathcal A_9}
\newcommand{\AAPPlift}{\mathcal A_9^{\#}}
\newcommand{\GraphAPP}{\Gamma_{\mathrm{APP},9}}
\newcommand{\MonNine}{\mathscr M_9}
\newcommand{\XTPK}{\mathcal X_{\mathrm{TPK}}}
\newcommand{\CellRQ}{\mathcal C_{9}^{\mathrm{rq}}}
\newcommand{\WordMonoid}{\mathfrak W_{9}^{\times}}
\newcommand{\Cdoce}{C_{12}}
\newcommand{\Kseal}{K}
\newcommand{\Usgn}{U_{12}^{\mathrm{sgn}}}
\newcommand{\Etrans}{\mathcal E_{3,4}}
\newcommand{\Rpi}{\mathscr R_{\pi}}
\newcommand{\Reu}{\mathscr R_{e}}
\newcommand{\Rphi}{\mathscr R_{\varphi}}
\newcommand{\RR}{\mathbb R}
\newcommand{\ZZ}{\mathbb Z}
\newcommand{\QQ}{\mathbb Q}
\newcommand{\FF}{\mathbb F}
\newcommand{\CC}{\mathbb C}
\newcommand{\dd}{\,\mathrm d}
\newcommand{\ee}{\mathrm e}
\newcommand{\Aut}{\operatorname{Aut}}
\newcommand{\Hol}{\operatorname{Hol}}
\newcommand{\dr}{\operatorname{dr}_9}
\newcommand{\Surv}{\operatorname{Surv}}
\newcommand{\val}{\operatorname{val}}
\newcommand{\Dnine}{D_9}
\newcommand{\Wnine}{\mathscr W_9}
\newcommand{\Krad}{\mathcal K_{\#}^{+}}
\newcommand{\Hplus}{\mathcal H_{+}}
\newcommand{\Dom}{\operatorname{Dom}}
\newcommand{\Tr}{\operatorname{Tr}}
\newcommand{\Spec}{\operatorname{Spec}}
\providecommand{\Sha}{\operatorname{Sha}}
\newcommand{\Irr}{\mathfrak I}
\newcommand{\NN}{\mathbb N}
\newcommand{\N}{\mathbb N}
\newcommand{\R}{\mathbb R}
\providecommand{\Adeg}{A_{\rm deg}}
\providecommand{\Arad}{A_{\rm rad}}
\providecommand{\Cdeg}{C^{*}_{\rm deg}}
\providecommand{\Crad}{C^{*}_{\rm rad}}
\providecommand{\etaRetrad}{\eta_{\rm ret}^{(\rm rad)}}
\providecommand{\Gtr}{\ensuremath{\Gamma_{\rm tr}}}
\providecommand{\HTr}{\ensuremath{\mathsf{Tr}_{q}}}
\providecommand{\HAdd}{\ensuremath{\mathsf{Add}}}
\providecommand{\HDef}{\ensuremath{\mathsf{Def}}}
\providecommand{\HRes}{\ensuremath{\mathsf{Res}}}
\providecommand{\HLoc}{\ensuremath{\mathsf{Loc}}}

\theoremstyle{plain}
\newtheorem{theorem}{Teorema}[chapter]
\newtheorem{lemma}[theorem]{Lema}
\newtheorem{proposition}[theorem]{Proposición}
\newtheorem{corollary}[theorem]{Corolario}
\theoremstyle{definition}
\newtheorem{definition}[theorem]{Definición}
\newtheorem{construction}[theorem]{Construcción}
\newtheorem{algorithm}[theorem]{Algoritmo}
\newtheorem{ablation}[theorem]{Proposición de necesidad estructural}
\newtheorem{criterion}[theorem]{Criterio de contraste}
\theoremstyle{remark}
\newtheorem{remark}[theorem]{Observación}

\tcbset{enhanced,breakable,boxrule=.45pt,arc=0mm,left=2.5mm,right=2.5mm,top=1.8mm,bottom=1.8mm}
\newtcolorbox{exactbox}[1][]{colback=white,colframe=black,title={Resultado exacto},fonttitle=\bfseries,#1}

\newcolumntype{Y}{>{\raggedright\arraybackslash}X}
\newcolumntype{R}{>{\raggedleft\arraybackslash}X}
\newcolumntype{L}[1]{>{\raggedright\arraybackslash}p{#1}}
\newcolumntype{C}{>{\centering\arraybackslash}p{1.65cm}}

\setlength{\emergencystretch}{3em}
 % Emisión de afirmaciones vivas durante la compilación de REV.9.
% Una marca sólo cuenta si la rama TeX que la contiene se ejecuta.
\newwrite\HMTClaimsStream
\immediate\openout\HMTClaimsStream=\jobname.hmtclaims
\newcommand{\HMTClaim}[1]{%
  \immediate\write\HMTClaimsStream{HMT-CLAIM:#1}%
}
\AtEndDocument{\immediate\closeout\HMTClaimsStream}
 % Compatibilidad tipada de las tres aportaciones incorporadas en la edición
% sucesora de 2026-08-12. Este archivo no redefine objetos de REV.3: aporta
% únicamente nombres, entornos y operadores que los cuadernos autónomos
% necesitan para residir en el grafo único del tratado.

\RequirePackage{xspace}

\definecolor{hmtteal}{RGB}{10,107,112}

\providecommand{\Rhyper}{\mathcal R_{\mathrm{hyp}}}
\providecommand{\OmegaInf}{\Omega_{\infty}}
\providecommand{\Clop}{\operatorname{Clop}}
\providecommand{\Closed}{\operatorname{Closed}}
\providecommand{\Proj}{\operatorname{Proj}}
\providecommand{\id}{\operatorname{id}}
\providecommand{\coker}{\operatorname{coker}}
\providecommand{\ord}{\operatorname{ord}}
\providecommand{\HALT}{\mathsf{HALT}}
\providecommand{\Prob}{\mathcal P}
\providecommand{\status}[1]{\emph{#1}}
\providecommand{\citep}{\cite}
\providecommand{\QMS}{\ensuremath{\mathrm{QMS}}\xspace}
\providecommand{\QLS}{\ensuremath{\mathrm{QLS}}\xspace}
\providecommand{\QPM}{\ensuremath{\mathrm{QPM}}\xspace}
\providecommand{\UCP}{\ensuremath{\mathrm{UCP}}\xspace}
\providecommand{\MUB}{\ensuremath{\mathrm{MUB}}\xspace}
\providecommand{\PVM}{\ensuremath{\mathrm{PVM}}\xspace}
\providecommand{\POVM}{\ensuremath{\mathrm{POVM}}\xspace}
\providecommand{\Mat}{\operatorname{Mat}}
\providecommand{\Her}{\operatorname{Her}}
\providecommand{\mconv}{\operatorname{mconv}}
\providecommand{\Span}{\operatorname{span}}
\providecommand{\Ad}{\operatorname{Ad}}
\providecommand{\diag}{\operatorname{diag}}

\theoremstyle{definition}
\newtheorem{example}[theorem]{Ejemplo}

\newtcolorbox{resultbox}[1][]{%
  colback=hmtlightgreen,colframe=hmtgreen,
  title={Resultado},fonttitle=\bfseries,#1}
\newtcolorbox{controlbox}[1][]{%
  colback=white,colframe=hmtgold!75!black,
  title={Obstrucción estructural},fonttitle=\bfseries,#1}
\newtcolorbox{sourcebox}[1][]{%
  colback=white,colframe=hmtblue!55,
  title={Premisas y derivación},fonttitle=\bfseries,#1}
\newtcolorbox{intuitionbox}[1][]{%
  colback=hmtlightgold,colframe=hmtgold,
  title={Lectura científica},fonttitle=\bfseries,#1}
\newtcolorbox{workbox}[1][]{%
  colback=white,colframe=hmtgray,
  title={Extensiones matemáticas},fonttitle=\bfseries,#1}
\newtcolorbox{importbox}[1][]{%
  colback=white,colframe=hmtblue!55,
  title={Teorema antecedente},fonttitle=\bfseries,#1}
\newtcolorbox{bridgebox}[1][]{%
  colback=hmtlightgold,colframe=hmtgold,
  title={Resultado del puente matricial},fonttitle=\bfseries,#1}
 % Estado material V3 comentario-only: conserva los 219 resultados de V5 y
% registra las dos correcciones editoriales autorizadas del puente matricial.
% Los sidecars 218/219 previos quedan como testigos inactivos e inmutables.
% HMT--MD REV.4: sidecar comentario-only de continuidad material V3/219.
% Encadena V2, conserva 219 IDs y no publica teoremas, títulos ni contenido paginado.
% Las correcciones registrales V3 no alteran enunciados ni fuerzas probatorias.
% HMT_RESULT_BEGIN:R3-001:residence
% {"material_control":"REV3_BASE_GRANULAR+PASS_NO_PERDIDA_REV3","material_state":"CONSERVADO_POR_BASE_GRANULAR_REV3","residences":[{"path":"/Users/ruben/Documents/HMT2/CONTINUIDAD_EDITORIAL/INTEGRACION_REV3_VON_NEUMANN_MATRICIAL_REALIZACION_2026-08-12/REV3_BASE_GRANULAR.json","sha256":"d89983feac6e064f0ce77d6506d60b94dd1eaf00ed174352d3e32bfc5a9b88d4"}],"result_id":"R3-001","statement_sha256":"ec22d0aed84b1a5ccd6d8675f494a2d54bb746ae05ca67325087c42e7697061d"}
% HMT_RESULT_END:R3-001:residence
% HMT_RESULT_BEGIN:R3-002:residence
% {"material_control":"REV3_BASE_GRANULAR+PASS_NO_PERDIDA_REV3","material_state":"CONSERVADO_POR_BASE_GRANULAR_REV3","residences":[{"path":"/Users/ruben/Documents/HMT2/CONTINUIDAD_EDITORIAL/INTEGRACION_REV3_VON_NEUMANN_MATRICIAL_REALIZACION_2026-08-12/REV3_BASE_GRANULAR.json","sha256":"d89983feac6e064f0ce77d6506d60b94dd1eaf00ed174352d3e32bfc5a9b88d4"}],"result_id":"R3-002","statement_sha256":"94b5e013944e5d31216adcca582bf4469c088ecd478d1b65dccc1a3343942490"}
% HMT_RESULT_END:R3-002:residence
% HMT_RESULT_BEGIN:R3-003:residence
% {"material_control":"REV3_BASE_GRANULAR+PASS_NO_PERDIDA_REV3","material_state":"CONSERVADO_POR_BASE_GRANULAR_REV3","residences":[{"path":"/Users/ruben/Documents/HMT2/CONTINUIDAD_EDITORIAL/INTEGRACION_REV3_VON_NEUMANN_MATRICIAL_REALIZACION_2026-08-12/REV3_BASE_GRANULAR.json","sha256":"d89983feac6e064f0ce77d6506d60b94dd1eaf00ed174352d3e32bfc5a9b88d4"}],"result_id":"R3-003","statement_sha256":"1c2e8dc0a9e7ba35279bfc5f48df699c614e4d96af09dcdaf1b76a665374c7cc"}
% HMT_RESULT_END:R3-003:residence
% HMT_RESULT_BEGIN:R3-004:residence
% {"material_control":"REV3_BASE_GRANULAR+PASS_NO_PERDIDA_REV3","material_state":"CONSERVADO_POR_BASE_GRANULAR_REV3","residences":[{"path":"/Users/ruben/Documents/HMT2/CONTINUIDAD_EDITORIAL/INTEGRACION_REV3_VON_NEUMANN_MATRICIAL_REALIZACION_2026-08-12/REV3_BASE_GRANULAR.json","sha256":"d89983feac6e064f0ce77d6506d60b94dd1eaf00ed174352d3e32bfc5a9b88d4"}],"result_id":"R3-004","statement_sha256":"067caf5027a3ca629eaa3f299767eebe36ab8da36d4a594f0490a6fcba2959a3"}
% HMT_RESULT_END:R3-004:residence
% HMT_RESULT_BEGIN:R3-005:residence
% {"material_control":"REV3_BASE_GRANULAR+PASS_NO_PERDIDA_REV3","material_state":"CONSERVADO_POR_BASE_GRANULAR_REV3","residences":[{"path":"/Users/ruben/Documents/HMT2/CONTINUIDAD_EDITORIAL/INTEGRACION_REV3_VON_NEUMANN_MATRICIAL_REALIZACION_2026-08-12/REV3_BASE_GRANULAR.json","sha256":"d89983feac6e064f0ce77d6506d60b94dd1eaf00ed174352d3e32bfc5a9b88d4"}],"result_id":"R3-005","statement_sha256":"7f015f268d0e7772da1a33238cd19903733ee7cfa8947c6372241210ebcd0f01"}
% HMT_RESULT_END:R3-005:residence
% HMT_RESULT_BEGIN:R3-006:residence
% {"material_control":"REV3_BASE_GRANULAR+PASS_NO_PERDIDA_REV3","material_state":"CONSERVADO_POR_BASE_GRANULAR_REV3","residences":[{"path":"/Users/ruben/Documents/HMT2/CONTINUIDAD_EDITORIAL/INTEGRACION_REV3_VON_NEUMANN_MATRICIAL_REALIZACION_2026-08-12/REV3_BASE_GRANULAR.json","sha256":"d89983feac6e064f0ce77d6506d60b94dd1eaf00ed174352d3e32bfc5a9b88d4"}],"result_id":"R3-006","statement_sha256":"9e3f5af6f313c9c4c256d7e6a1f22a0bc0d112d95a7909a4b61b35638a89ea73"}
% HMT_RESULT_END:R3-006:residence
% HMT_RESULT_BEGIN:R3-007:residence
% {"material_control":"REV3_BASE_GRANULAR+PASS_NO_PERDIDA_REV3","material_state":"CONSERVADO_POR_BASE_GRANULAR_REV3","residences":[{"path":"/Users/ruben/Documents/HMT2/CONTINUIDAD_EDITORIAL/INTEGRACION_REV3_VON_NEUMANN_MATRICIAL_REALIZACION_2026-08-12/REV3_BASE_GRANULAR.json","sha256":"d89983feac6e064f0ce77d6506d60b94dd1eaf00ed174352d3e32bfc5a9b88d4"}],"result_id":"R3-007","statement_sha256":"afe6a7f442b4b9b73f09e57dcf993019f691bf05fabf3cf8a866c0cb30e0b20e"}
% HMT_RESULT_END:R3-007:residence
% HMT_RESULT_BEGIN:R3-008:residence
% {"material_control":"REV3_BASE_GRANULAR+PASS_NO_PERDIDA_REV3","material_state":"CONSERVADO_POR_BASE_GRANULAR_REV3","residences":[{"path":"/Users/ruben/Documents/HMT2/CONTINUIDAD_EDITORIAL/INTEGRACION_REV3_VON_NEUMANN_MATRICIAL_REALIZACION_2026-08-12/REV3_BASE_GRANULAR.json","sha256":"d89983feac6e064f0ce77d6506d60b94dd1eaf00ed174352d3e32bfc5a9b88d4"}],"result_id":"R3-008","statement_sha256":"3f0aaf6fa196a66818755e0fe8d0a900ec8170dc90de754fc9b133303d4627b7"}
% HMT_RESULT_END:R3-008:residence
% HMT_RESULT_BEGIN:R3-009:residence
% {"material_control":"REV3_BASE_GRANULAR+PASS_NO_PERDIDA_REV3","material_state":"CONSERVADO_POR_BASE_GRANULAR_REV3","residences":[{"path":"/Users/ruben/Documents/HMT2/CONTINUIDAD_EDITORIAL/INTEGRACION_REV3_VON_NEUMANN_MATRICIAL_REALIZACION_2026-08-12/REV3_BASE_GRANULAR.json","sha256":"d89983feac6e064f0ce77d6506d60b94dd1eaf00ed174352d3e32bfc5a9b88d4"}],"result_id":"R3-009","statement_sha256":"84c9171db979cfcda487379afcdf640f4d78458cf4f7d5697a3ec0b30a2f3096"}
% HMT_RESULT_END:R3-009:residence
% HMT_RESULT_BEGIN:R3-010:residence
% {"material_control":"REV3_BASE_GRANULAR+PASS_NO_PERDIDA_REV3","material_state":"CONSERVADO_POR_BASE_GRANULAR_REV3","residences":[{"path":"/Users/ruben/Documents/HMT2/CONTINUIDAD_EDITORIAL/INTEGRACION_REV3_VON_NEUMANN_MATRICIAL_REALIZACION_2026-08-12/REV3_BASE_GRANULAR.json","sha256":"d89983feac6e064f0ce77d6506d60b94dd1eaf00ed174352d3e32bfc5a9b88d4"}],"result_id":"R3-010","statement_sha256":"3ce9bef7730788d82ad408181bd8942d9cc48d02b7e2b99b36b5f763b3a20968"}
% HMT_RESULT_END:R3-010:residence
% HMT_RESULT_BEGIN:R3-011:residence
% {"material_control":"REV3_BASE_GRANULAR+PASS_NO_PERDIDA_REV3","material_state":"CONSERVADO_POR_BASE_GRANULAR_REV3","residences":[{"path":"/Users/ruben/Documents/HMT2/CONTINUIDAD_EDITORIAL/INTEGRACION_REV3_VON_NEUMANN_MATRICIAL_REALIZACION_2026-08-12/REV3_BASE_GRANULAR.json","sha256":"d89983feac6e064f0ce77d6506d60b94dd1eaf00ed174352d3e32bfc5a9b88d4"}],"result_id":"R3-011","statement_sha256":"5de78d3622665d2d9585dcb93d9f3d762ca07c9959ab042c4c085b8ef5943597"}
% HMT_RESULT_END:R3-011:residence
% HMT_RESULT_BEGIN:R3-012:residence
% {"material_control":"REV3_BASE_GRANULAR+PASS_NO_PERDIDA_REV3","material_state":"CONSERVADO_POR_BASE_GRANULAR_REV3","residences":[{"path":"/Users/ruben/Documents/HMT2/CONTINUIDAD_EDITORIAL/INTEGRACION_REV3_VON_NEUMANN_MATRICIAL_REALIZACION_2026-08-12/REV3_BASE_GRANULAR.json","sha256":"d89983feac6e064f0ce77d6506d60b94dd1eaf00ed174352d3e32bfc5a9b88d4"}],"result_id":"R3-012","statement_sha256":"a873e4e36ef58361dfd662e6f6b08f0484e978fb07f397dab890d355b8ae0df8"}
% HMT_RESULT_END:R3-012:residence
% HMT_RESULT_BEGIN:R3-013:residence
% {"material_control":"REV3_BASE_GRANULAR+PASS_NO_PERDIDA_REV3","material_state":"CONSERVADO_POR_BASE_GRANULAR_REV3","residences":[{"path":"/Users/ruben/Documents/HMT2/CONTINUIDAD_EDITORIAL/INTEGRACION_REV3_VON_NEUMANN_MATRICIAL_REALIZACION_2026-08-12/REV3_BASE_GRANULAR.json","sha256":"d89983feac6e064f0ce77d6506d60b94dd1eaf00ed174352d3e32bfc5a9b88d4"}],"result_id":"R3-013","statement_sha256":"892af66958bb697bdbc9d87aa439e0dcebdecff82281fa756e35aac1088a3c74"}
% HMT_RESULT_END:R3-013:residence
% HMT_RESULT_BEGIN:R3-014:residence
% {"material_control":"REV3_BASE_GRANULAR+PASS_NO_PERDIDA_REV3","material_state":"CONSERVADO_POR_BASE_GRANULAR_REV3","residences":[{"path":"/Users/ruben/Documents/HMT2/CONTINUIDAD_EDITORIAL/INTEGRACION_REV3_VON_NEUMANN_MATRICIAL_REALIZACION_2026-08-12/REV3_BASE_GRANULAR.json","sha256":"d89983feac6e064f0ce77d6506d60b94dd1eaf00ed174352d3e32bfc5a9b88d4"}],"result_id":"R3-014","statement_sha256":"3b135fbe4cbe5d0e292ffcee032f81cef6c993d833f85dde740957bd50df821f"}
% HMT_RESULT_END:R3-014:residence
% HMT_RESULT_BEGIN:R3-015:residence
% {"material_control":"REV3_BASE_GRANULAR+PASS_NO_PERDIDA_REV3","material_state":"CONSERVADO_POR_BASE_GRANULAR_REV3","residences":[{"path":"/Users/ruben/Documents/HMT2/CONTINUIDAD_EDITORIAL/INTEGRACION_REV3_VON_NEUMANN_MATRICIAL_REALIZACION_2026-08-12/REV3_BASE_GRANULAR.json","sha256":"d89983feac6e064f0ce77d6506d60b94dd1eaf00ed174352d3e32bfc5a9b88d4"}],"result_id":"R3-015","statement_sha256":"f94da8f805b5238f940f359590ae265b58be35f6bf4257fac8fa16cf8933c5f8"}
% HMT_RESULT_END:R3-015:residence
% HMT_RESULT_BEGIN:R3-016:residence
% {"material_control":"REV3_BASE_GRANULAR+PASS_NO_PERDIDA_REV3","material_state":"CONSERVADO_POR_BASE_GRANULAR_REV3","residences":[{"path":"/Users/ruben/Documents/HMT2/CONTINUIDAD_EDITORIAL/INTEGRACION_REV3_VON_NEUMANN_MATRICIAL_REALIZACION_2026-08-12/REV3_BASE_GRANULAR.json","sha256":"d89983feac6e064f0ce77d6506d60b94dd1eaf00ed174352d3e32bfc5a9b88d4"}],"result_id":"R3-016","statement_sha256":"660c7b55c80036946d6df58772d5bebbfded71d302316af3ab877a9e3cd1f3b7"}
% HMT_RESULT_END:R3-016:residence
% HMT_RESULT_BEGIN:R3-017:residence
% {"material_control":"REV3_BASE_GRANULAR+PASS_NO_PERDIDA_REV3","material_state":"CONSERVADO_POR_BASE_GRANULAR_REV3","residences":[{"path":"/Users/ruben/Documents/HMT2/CONTINUIDAD_EDITORIAL/INTEGRACION_REV3_VON_NEUMANN_MATRICIAL_REALIZACION_2026-08-12/REV3_BASE_GRANULAR.json","sha256":"d89983feac6e064f0ce77d6506d60b94dd1eaf00ed174352d3e32bfc5a9b88d4"}],"result_id":"R3-017","statement_sha256":"0834862c1f2fd11986bcab78a6a9ad03cd89152459ab8f8a6b83e9d1d4a0d0a8"}
% HMT_RESULT_END:R3-017:residence
% HMT_RESULT_BEGIN:R3-018:residence
% {"material_control":"REV3_BASE_GRANULAR+PASS_NO_PERDIDA_REV3","material_state":"CONSERVADO_POR_BASE_GRANULAR_REV3","residences":[{"path":"/Users/ruben/Documents/HMT2/CONTINUIDAD_EDITORIAL/INTEGRACION_REV3_VON_NEUMANN_MATRICIAL_REALIZACION_2026-08-12/REV3_BASE_GRANULAR.json","sha256":"d89983feac6e064f0ce77d6506d60b94dd1eaf00ed174352d3e32bfc5a9b88d4"}],"result_id":"R3-018","statement_sha256":"7a6ecf5a52d2b3bb92ca8b82bb5e3852c1a1f729a12ca8d39ad7353a61b2259b"}
% HMT_RESULT_END:R3-018:residence
% HMT_RESULT_BEGIN:R3-019:residence
% {"material_control":"REV3_BASE_GRANULAR+PASS_NO_PERDIDA_REV3","material_state":"CONSERVADO_POR_BASE_GRANULAR_REV3","residences":[{"path":"/Users/ruben/Documents/HMT2/CONTINUIDAD_EDITORIAL/INTEGRACION_REV3_VON_NEUMANN_MATRICIAL_REALIZACION_2026-08-12/REV3_BASE_GRANULAR.json","sha256":"d89983feac6e064f0ce77d6506d60b94dd1eaf00ed174352d3e32bfc5a9b88d4"}],"result_id":"R3-019","statement_sha256":"949668449b4331519738227ea1e782aefbfbe4067e371147c0821de5c4e2314c"}
% HMT_RESULT_END:R3-019:residence
% HMT_RESULT_BEGIN:R3-020:residence
% {"material_control":"REV3_BASE_GRANULAR+PASS_NO_PERDIDA_REV3","material_state":"CONSERVADO_POR_BASE_GRANULAR_REV3","residences":[{"path":"/Users/ruben/Documents/HMT2/CONTINUIDAD_EDITORIAL/INTEGRACION_REV3_VON_NEUMANN_MATRICIAL_REALIZACION_2026-08-12/REV3_BASE_GRANULAR.json","sha256":"d89983feac6e064f0ce77d6506d60b94dd1eaf00ed174352d3e32bfc5a9b88d4"}],"result_id":"R3-020","statement_sha256":"e6097d0d922f2e8e30dd23b38e77066a77f46464e19a875d95d0e9c21b9acd23"}
% HMT_RESULT_END:R3-020:residence
% HMT_RESULT_BEGIN:C-001:residence
% {"material_control":"PASS_MATRIZ_ATOMICA_AUTOCONTENIDA_REV3","material_state":"CONTROL_ACTIVO_EN_PLAN_Y_FUENTE","residences":[{"path":"/Users/ruben/Documents/HMT2/CONTINUIDAD_EDITORIAL/INTEGRACION_REV3_VON_NEUMANN_MATRICIAL_REALIZACION_2026-08-12/verificar_matriz_atomica_autocontenida.py","sha256":"a6598c0e1da8e8a8c61acc4150033a66348c039fb2ebf241791b0517dc458472"}],"result_id":"C-001","statement_sha256":"beae79d27d25d820609da41c6347a8561c6d41d85b1ec0525092421ad8721cf3"}
% HMT_RESULT_END:C-001:residence
% HMT_RESULT_BEGIN:C-002:residence
% {"material_control":"PASS_MATRIZ_ATOMICA_AUTOCONTENIDA_REV3","material_state":"CONTROL_ACTIVO_EN_PLAN_Y_FUENTE","residences":[{"path":"/Users/ruben/Documents/HMT2/CONTINUIDAD_EDITORIAL/INTEGRACION_REV3_VON_NEUMANN_MATRICIAL_REALIZACION_2026-08-12/verificar_matriz_atomica_autocontenida.py","sha256":"a6598c0e1da8e8a8c61acc4150033a66348c039fb2ebf241791b0517dc458472"}],"result_id":"C-002","statement_sha256":"5cf2ae9f5be634f470caa6670cf89d72015dc631a63870dc323d018d9f03d0e0"}
% HMT_RESULT_END:C-002:residence
% HMT_RESULT_BEGIN:C-003:residence
% {"material_control":"PASS_MATRIZ_ATOMICA_AUTOCONTENIDA_REV3","material_state":"CONTROL_ACTIVO_EN_PLAN_Y_FUENTE","residences":[{"path":"/Users/ruben/Documents/HMT2/CONTINUIDAD_EDITORIAL/INTEGRACION_REV3_VON_NEUMANN_MATRICIAL_REALIZACION_2026-08-12/verificar_matriz_atomica_autocontenida.py","sha256":"a6598c0e1da8e8a8c61acc4150033a66348c039fb2ebf241791b0517dc458472"}],"result_id":"C-003","statement_sha256":"2cabb42c3d350ac50184c7358426ca161c539e6dfd45b760168861d7461247ad"}
% HMT_RESULT_END:C-003:residence
% HMT_RESULT_BEGIN:C-004:residence
% {"material_control":"PASS_MATRIZ_ATOMICA_AUTOCONTENIDA_REV3","material_state":"CONTROL_ACTIVO_EN_PLAN_Y_FUENTE","residences":[{"path":"/Users/ruben/Documents/HMT2/CONTINUIDAD_EDITORIAL/INTEGRACION_REV3_VON_NEUMANN_MATRICIAL_REALIZACION_2026-08-12/verificar_matriz_atomica_autocontenida.py","sha256":"a6598c0e1da8e8a8c61acc4150033a66348c039fb2ebf241791b0517dc458472"}],"result_id":"C-004","statement_sha256":"9ff96e268873d82ee4ac8616df2a572ca23330c63f928a3e7a67ea10e2d46084"}
% HMT_RESULT_END:C-004:residence
% HMT_RESULT_BEGIN:C-005:residence
% {"material_control":"PASS_MATRIZ_ATOMICA_AUTOCONTENIDA_REV3","material_state":"CONTROL_ACTIVO_EN_PLAN_Y_FUENTE","residences":[{"path":"/Users/ruben/Documents/HMT2/CONTINUIDAD_EDITORIAL/INTEGRACION_REV3_VON_NEUMANN_MATRICIAL_REALIZACION_2026-08-12/verificar_matriz_atomica_autocontenida.py","sha256":"a6598c0e1da8e8a8c61acc4150033a66348c039fb2ebf241791b0517dc458472"}],"result_id":"C-005","statement_sha256":"f28286ac6df6ee110b8bdd8f2de9acd882a13dafee92dd7bc0645accfe4384d5"}
% HMT_RESULT_END:C-005:residence
% HMT_RESULT_BEGIN:C-006:residence
% {"material_control":"PASS_MATRIZ_ATOMICA_AUTOCONTENIDA_REV3","material_state":"CONTROL_ACTIVO_EN_PLAN_Y_FUENTE","residences":[{"path":"/Users/ruben/Documents/HMT2/CONTINUIDAD_EDITORIAL/INTEGRACION_REV3_VON_NEUMANN_MATRICIAL_REALIZACION_2026-08-12/verificar_matriz_atomica_autocontenida.py","sha256":"a6598c0e1da8e8a8c61acc4150033a66348c039fb2ebf241791b0517dc458472"}],"result_id":"C-006","statement_sha256":"61914b1116661f7db384d4e45bba72ec85e7ff331800fc50353e8e93f850dd76"}
% HMT_RESULT_END:C-006:residence
% HMT_RESULT_BEGIN:C-007:residence
% {"material_control":"PASS_MATRIZ_ATOMICA_AUTOCONTENIDA_REV3","material_state":"CONTROL_ACTIVO_EN_PLAN_Y_FUENTE","residences":[{"path":"/Users/ruben/Documents/HMT2/CONTINUIDAD_EDITORIAL/INTEGRACION_REV3_VON_NEUMANN_MATRICIAL_REALIZACION_2026-08-12/verificar_matriz_atomica_autocontenida.py","sha256":"a6598c0e1da8e8a8c61acc4150033a66348c039fb2ebf241791b0517dc458472"}],"result_id":"C-007","statement_sha256":"12e1672ff9d4988000a45bcd1a7b89ef0b97692c2ca5f6756ca06ad9670ede54"}
% HMT_RESULT_END:C-007:residence
% HMT_RESULT_BEGIN:MAS-RECTOR-001:residence
% {"material_control":"PASS_PAQUETE_RECTOR_MASAS_HMT_MD_2026_08_06","material_state":"CONSERVADO_CON_APTO_FOCAL","residences":[{"path":"/Users/ruben/Documents/New project/PUBLICACION_HMT/PAQUETE_RECTOR_LEY_GENERAL_MASAS_HMT_MD_2026-08-06/MATRIZ_RECTORA.tsv","sha256":"d4d58d9e5394804b83f48a2996ef884b275669b0f9468a28c9ed7c0a58affdc5"}],"result_id":"MAS-RECTOR-001","statement_sha256":"0a768984b144ee78bd8e25d36a2945dd57a6f83a5071c4fcfc7ba1a0b68f8d7f"}
% HMT_RESULT_END:MAS-RECTOR-001:residence
% HMT_RESULT_BEGIN:MAS-RECTOR-002:residence
% {"material_control":"PASS_PAQUETE_RECTOR_MASAS_HMT_MD_2026_08_06","material_state":"CONSERVADO_CON_APTO_FOCAL","residences":[{"path":"/Users/ruben/Documents/New project/PUBLICACION_HMT/PAQUETE_RECTOR_LEY_GENERAL_MASAS_HMT_MD_2026-08-06/MATRIZ_RECTORA.tsv","sha256":"d4d58d9e5394804b83f48a2996ef884b275669b0f9468a28c9ed7c0a58affdc5"}],"result_id":"MAS-RECTOR-002","statement_sha256":"89fa9d9d6538d06602255e139245ecc085a236981ab5f973f21ada8ca0f83cd3"}
% HMT_RESULT_END:MAS-RECTOR-002:residence
% HMT_RESULT_BEGIN:MAS-RECTOR-003:residence
% {"material_control":"PASS_PAQUETE_RECTOR_MASAS_HMT_MD_2026_08_06","material_state":"CONSERVADO_CON_APTO_FOCAL","residences":[{"path":"/Users/ruben/Documents/New project/PUBLICACION_HMT/PAQUETE_RECTOR_LEY_GENERAL_MASAS_HMT_MD_2026-08-06/MATRIZ_RECTORA.tsv","sha256":"d4d58d9e5394804b83f48a2996ef884b275669b0f9468a28c9ed7c0a58affdc5"}],"result_id":"MAS-RECTOR-003","statement_sha256":"a859bf6639b8355be0ff5ebcf8ba32630249b5d4842f80ebc5272733f0715109"}
% HMT_RESULT_END:MAS-RECTOR-003:residence
% HMT_RESULT_BEGIN:MAS-RECTOR-004:residence
% {"material_control":"PASS_PAQUETE_RECTOR_MASAS_HMT_MD_2026_08_06","material_state":"CONSERVADO_CON_APTO_FOCAL","residences":[{"path":"/Users/ruben/Documents/New project/PUBLICACION_HMT/PAQUETE_RECTOR_LEY_GENERAL_MASAS_HMT_MD_2026-08-06/MATRIZ_RECTORA.tsv","sha256":"d4d58d9e5394804b83f48a2996ef884b275669b0f9468a28c9ed7c0a58affdc5"}],"result_id":"MAS-RECTOR-004","statement_sha256":"5478fc19a7db005a9c0a02175258a722aa9289c8aae2b40c0f8be1a5539331e6"}
% HMT_RESULT_END:MAS-RECTOR-004:residence
% HMT_RESULT_BEGIN:MAS-RECTOR-005:residence
% {"material_control":"PASS_PAQUETE_RECTOR_MASAS_HMT_MD_2026_08_06","material_state":"CONSERVADO_CON_APTO_FOCAL","residences":[{"path":"/Users/ruben/Documents/New project/PUBLICACION_HMT/PAQUETE_RECTOR_LEY_GENERAL_MASAS_HMT_MD_2026-08-06/MATRIZ_RECTORA.tsv","sha256":"d4d58d9e5394804b83f48a2996ef884b275669b0f9468a28c9ed7c0a58affdc5"}],"result_id":"MAS-RECTOR-005","statement_sha256":"9ac8216130e6ddb43e8032b3eb72c4dd18ab87c05f238ee8fa65f06fdea782d2"}
% HMT_RESULT_END:MAS-RECTOR-005:residence
% HMT_RESULT_BEGIN:MAS-RECTOR-006:residence
% {"material_control":"PASS_PAQUETE_RECTOR_MASAS_HMT_MD_2026_08_06","material_state":"CONSERVADO_CON_APTO_FOCAL","residences":[{"path":"/Users/ruben/Documents/New project/PUBLICACION_HMT/PAQUETE_RECTOR_LEY_GENERAL_MASAS_HMT_MD_2026-08-06/MATRIZ_RECTORA.tsv","sha256":"d4d58d9e5394804b83f48a2996ef884b275669b0f9468a28c9ed7c0a58affdc5"}],"result_id":"MAS-RECTOR-006","statement_sha256":"f64d301daf63969b4ae52cb1259e81411210bb18d9c5e297f815a8f7ea6698e3"}
% HMT_RESULT_END:MAS-RECTOR-006:residence
% HMT_RESULT_BEGIN:MAS-RECTOR-007:residence
% {"material_control":"PASS_PAQUETE_RECTOR_MASAS_HMT_MD_2026_08_06","material_state":"CONSERVADO_CON_APTO_FOCAL","residences":[{"path":"/Users/ruben/Documents/New project/PUBLICACION_HMT/PAQUETE_RECTOR_LEY_GENERAL_MASAS_HMT_MD_2026-08-06/MATRIZ_RECTORA.tsv","sha256":"d4d58d9e5394804b83f48a2996ef884b275669b0f9468a28c9ed7c0a58affdc5"}],"result_id":"MAS-RECTOR-007","statement_sha256":"537017eed6d0ecd896c8a40aed8e3197aeef62bb6895d356b95a7bbc2948643e"}
% HMT_RESULT_END:MAS-RECTOR-007:residence
% HMT_RESULT_BEGIN:MAS-RECTOR-008:residence
% {"material_control":"PASS_PAQUETE_RECTOR_MASAS_HMT_MD_2026_08_06","material_state":"CONSERVADO_CON_APTO_FOCAL","residences":[{"path":"/Users/ruben/Documents/New project/PUBLICACION_HMT/PAQUETE_RECTOR_LEY_GENERAL_MASAS_HMT_MD_2026-08-06/MATRIZ_RECTORA.tsv","sha256":"d4d58d9e5394804b83f48a2996ef884b275669b0f9468a28c9ed7c0a58affdc5"}],"result_id":"MAS-RECTOR-008","statement_sha256":"9563efbb323a0406d27d86e100b82f2b9ca50dd08d2197ca5d02418e4656c614"}
% HMT_RESULT_END:MAS-RECTOR-008:residence
% HMT_RESULT_BEGIN:MAS-RECTOR-009:residence
% {"material_control":"PASS_PAQUETE_RECTOR_MASAS_HMT_MD_2026_08_06","material_state":"CONSERVADO_CON_APTO_FOCAL","residences":[{"path":"/Users/ruben/Documents/New project/PUBLICACION_HMT/PAQUETE_RECTOR_LEY_GENERAL_MASAS_HMT_MD_2026-08-06/MATRIZ_RECTORA.tsv","sha256":"d4d58d9e5394804b83f48a2996ef884b275669b0f9468a28c9ed7c0a58affdc5"}],"result_id":"MAS-RECTOR-009","statement_sha256":"cda0251b6312a403ff360c8c40e29d102a3395cef7e4b3f6ea6b2b72ef982e76"}
% HMT_RESULT_END:MAS-RECTOR-009:residence
% HMT_RESULT_BEGIN:MAS-RECTOR-010:residence
% {"material_control":"PASS_PAQUETE_RECTOR_MASAS_HMT_MD_2026_08_06","material_state":"CONSERVADO_CON_APTO_FOCAL","residences":[{"path":"/Users/ruben/Documents/New project/PUBLICACION_HMT/PAQUETE_RECTOR_LEY_GENERAL_MASAS_HMT_MD_2026-08-06/MATRIZ_RECTORA.tsv","sha256":"d4d58d9e5394804b83f48a2996ef884b275669b0f9468a28c9ed7c0a58affdc5"}],"result_id":"MAS-RECTOR-010","statement_sha256":"5913a73e04b4c752216af319b77e125d31914639200b39b45d1a6b28b536508b"}
% HMT_RESULT_END:MAS-RECTOR-010:residence
% HMT_RESULT_BEGIN:MAS-RECTOR-011:residence
% {"material_control":"PASS_PAQUETE_RECTOR_MASAS_HMT_MD_2026_08_06","material_state":"CONSERVADO_CON_APTO_FOCAL","residences":[{"path":"/Users/ruben/Documents/New project/PUBLICACION_HMT/PAQUETE_RECTOR_LEY_GENERAL_MASAS_HMT_MD_2026-08-06/MATRIZ_RECTORA.tsv","sha256":"d4d58d9e5394804b83f48a2996ef884b275669b0f9468a28c9ed7c0a58affdc5"}],"result_id":"MAS-RECTOR-011","statement_sha256":"b56b1d18f5b984455801b89eace1fbf9143fa965b645cd443b7bbef759b09251"}
% HMT_RESULT_END:MAS-RECTOR-011:residence
% HMT_RESULT_BEGIN:MAS-RECTOR-012:residence
% {"material_control":"PASS_PAQUETE_RECTOR_MASAS_HMT_MD_2026_08_06","material_state":"CONSERVADO_CON_APTO_FOCAL","residences":[{"path":"/Users/ruben/Documents/New project/PUBLICACION_HMT/PAQUETE_RECTOR_LEY_GENERAL_MASAS_HMT_MD_2026-08-06/MATRIZ_RECTORA.tsv","sha256":"d4d58d9e5394804b83f48a2996ef884b275669b0f9468a28c9ed7c0a58affdc5"}],"result_id":"MAS-RECTOR-012","statement_sha256":"aa319b29fb1a3ad47b2d521716fea82da5660479225223eac504798c36b42f24"}
% HMT_RESULT_END:MAS-RECTOR-012:residence
% HMT_RESULT_BEGIN:MAS-RECTOR-013:residence
% {"material_control":"PASS_PAQUETE_RECTOR_MASAS_HMT_MD_2026_08_06","material_state":"CONSERVADO_CON_APTO_FOCAL","residences":[{"path":"/Users/ruben/Documents/New project/PUBLICACION_HMT/PAQUETE_RECTOR_LEY_GENERAL_MASAS_HMT_MD_2026-08-06/MATRIZ_RECTORA.tsv","sha256":"d4d58d9e5394804b83f48a2996ef884b275669b0f9468a28c9ed7c0a58affdc5"}],"result_id":"MAS-RECTOR-013","statement_sha256":"c9f012d4f1b9fa1d9a642670582f65c700f433930615ddc0500e4beef203d5e9"}
% HMT_RESULT_END:MAS-RECTOR-013:residence
% HMT_RESULT_BEGIN:MAS-RECTOR-014:residence
% {"material_control":"PASS_PAQUETE_RECTOR_MASAS_HMT_MD_2026_08_06","material_state":"CONSERVADO_CON_APTO_FOCAL","residences":[{"path":"/Users/ruben/Documents/New project/PUBLICACION_HMT/PAQUETE_RECTOR_LEY_GENERAL_MASAS_HMT_MD_2026-08-06/MATRIZ_RECTORA.tsv","sha256":"d4d58d9e5394804b83f48a2996ef884b275669b0f9468a28c9ed7c0a58affdc5"}],"result_id":"MAS-RECTOR-014","statement_sha256":"12aadab1bbdd24c3cf826e52a333c34f70ced2a37f655b2919eaaaf3281fef04"}
% HMT_RESULT_END:MAS-RECTOR-014:residence
% HMT_RESULT_BEGIN:MAS-RECTOR-015:residence
% {"material_control":"PASS_PAQUETE_RECTOR_MASAS_HMT_MD_2026_08_06","material_state":"CONSERVADO_CON_APTO_FOCAL","residences":[{"path":"/Users/ruben/Documents/New project/PUBLICACION_HMT/PAQUETE_RECTOR_LEY_GENERAL_MASAS_HMT_MD_2026-08-06/MATRIZ_RECTORA.tsv","sha256":"d4d58d9e5394804b83f48a2996ef884b275669b0f9468a28c9ed7c0a58affdc5"}],"result_id":"MAS-RECTOR-015","statement_sha256":"ae01c9c511c63968eed7516b99a2cf7e2a65d0d1c6f56e756638a3d524bb5f9f"}
% HMT_RESULT_END:MAS-RECTOR-015:residence
% HMT_RESULT_BEGIN:MAS-RECTOR-016:residence
% {"material_control":"PASS_PAQUETE_RECTOR_MASAS_HMT_MD_2026_08_06","material_state":"CONSERVADO_CON_APTO_FOCAL","residences":[{"path":"/Users/ruben/Documents/New project/PUBLICACION_HMT/PAQUETE_RECTOR_LEY_GENERAL_MASAS_HMT_MD_2026-08-06/MATRIZ_RECTORA.tsv","sha256":"d4d58d9e5394804b83f48a2996ef884b275669b0f9468a28c9ed7c0a58affdc5"}],"result_id":"MAS-RECTOR-016","statement_sha256":"685df6bdb27ac04026246ca3f09b212f9aabebf22929fb63ddfe1222da2a6e76"}
% HMT_RESULT_END:MAS-RECTOR-016:residence
% HMT_RESULT_BEGIN:MAS-RECTOR-017:residence
% {"material_control":"PASS_PAQUETE_RECTOR_MASAS_HMT_MD_2026_08_06","material_state":"CONSERVADO_CON_APTO_FOCAL","residences":[{"path":"/Users/ruben/Documents/New project/PUBLICACION_HMT/PAQUETE_RECTOR_LEY_GENERAL_MASAS_HMT_MD_2026-08-06/MATRIZ_RECTORA.tsv","sha256":"d4d58d9e5394804b83f48a2996ef884b275669b0f9468a28c9ed7c0a58affdc5"}],"result_id":"MAS-RECTOR-017","statement_sha256":"2d6f3806c75ac84503e6b8ce4c6fe49dbc650e525915991cfef4ecc29e12bb55"}
% HMT_RESULT_END:MAS-RECTOR-017:residence
% HMT_RESULT_BEGIN:MAS-RECTOR-018:residence
% {"material_control":"PASS_PAQUETE_RECTOR_MASAS_HMT_MD_2026_08_06","material_state":"CONSERVADO_CON_APTO_FOCAL","residences":[{"path":"/Users/ruben/Documents/New project/PUBLICACION_HMT/PAQUETE_RECTOR_LEY_GENERAL_MASAS_HMT_MD_2026-08-06/MATRIZ_RECTORA.tsv","sha256":"d4d58d9e5394804b83f48a2996ef884b275669b0f9468a28c9ed7c0a58affdc5"}],"result_id":"MAS-RECTOR-018","statement_sha256":"1fc9dcfc0f77f0e3dcf950b1b3ffb6ee7ee7314e9e8407bb22f75fa7efd5300c"}
% HMT_RESULT_END:MAS-RECTOR-018:residence
% HMT_RESULT_BEGIN:MAS-RECTOR-019:residence
% {"material_control":"PASS_PAQUETE_RECTOR_MASAS_HMT_MD_2026_08_06","material_state":"CONSERVADO_CON_APTO_FOCAL","residences":[{"path":"/Users/ruben/Documents/New project/PUBLICACION_HMT/PAQUETE_RECTOR_LEY_GENERAL_MASAS_HMT_MD_2026-08-06/MATRIZ_RECTORA.tsv","sha256":"d4d58d9e5394804b83f48a2996ef884b275669b0f9468a28c9ed7c0a58affdc5"}],"result_id":"MAS-RECTOR-019","statement_sha256":"ebbe49ae2f4afbc32a86801a2996b2c2d2bf71c0ff4f86d9dc3e58cf6366ae60"}
% HMT_RESULT_END:MAS-RECTOR-019:residence
% HMT_RESULT_BEGIN:MAS-RECTOR-020:residence
% {"material_control":"PASS_PAQUETE_RECTOR_MASAS_HMT_MD_2026_08_06","material_state":"CONSERVADO_CON_APTO_FOCAL","residences":[{"path":"/Users/ruben/Documents/New project/PUBLICACION_HMT/PAQUETE_RECTOR_LEY_GENERAL_MASAS_HMT_MD_2026-08-06/MATRIZ_RECTORA.tsv","sha256":"d4d58d9e5394804b83f48a2996ef884b275669b0f9468a28c9ed7c0a58affdc5"}],"result_id":"MAS-RECTOR-020","statement_sha256":"ea7797a9dc2c4c52d16299bf176fa0b0178158922903c442a8cb1aa85af17879"}
% HMT_RESULT_END:MAS-RECTOR-020:residence
% HMT_RESULT_BEGIN:MAS-RECTOR-021:residence
% {"material_control":"PASS_PAQUETE_RECTOR_MASAS_HMT_MD_2026_08_06","material_state":"CONSERVADO_CON_APTO_FOCAL","residences":[{"path":"/Users/ruben/Documents/New project/PUBLICACION_HMT/PAQUETE_RECTOR_LEY_GENERAL_MASAS_HMT_MD_2026-08-06/MATRIZ_RECTORA.tsv","sha256":"d4d58d9e5394804b83f48a2996ef884b275669b0f9468a28c9ed7c0a58affdc5"}],"result_id":"MAS-RECTOR-021","statement_sha256":"7f875aaae09118d70a3022037a232a45a3e2fe74a8c92aae01e5486aa2ef7f95"}
% HMT_RESULT_END:MAS-RECTOR-021:residence
% HMT_RESULT_BEGIN:MAS-RECTOR-022:residence
% {"material_control":"PASS_PAQUETE_RECTOR_MASAS_HMT_MD_2026_08_06","material_state":"CONSERVADO_CON_APTO_FOCAL","residences":[{"path":"/Users/ruben/Documents/New project/PUBLICACION_HMT/PAQUETE_RECTOR_LEY_GENERAL_MASAS_HMT_MD_2026-08-06/MATRIZ_RECTORA.tsv","sha256":"d4d58d9e5394804b83f48a2996ef884b275669b0f9468a28c9ed7c0a58affdc5"}],"result_id":"MAS-RECTOR-022","statement_sha256":"c6e9560103ad03e63b2460672442d06e18c09c1f866c7d34ee4946f543255d51"}
% HMT_RESULT_END:MAS-RECTOR-022:residence
% HMT_RESULT_BEGIN:MAS-RECTOR-023:residence
% {"material_control":"PASS_PAQUETE_RECTOR_MASAS_HMT_MD_2026_08_06","material_state":"CONSERVADO_CON_APTO_FOCAL","residences":[{"path":"/Users/ruben/Documents/New project/PUBLICACION_HMT/PAQUETE_RECTOR_LEY_GENERAL_MASAS_HMT_MD_2026-08-06/MATRIZ_RECTORA.tsv","sha256":"d4d58d9e5394804b83f48a2996ef884b275669b0f9468a28c9ed7c0a58affdc5"}],"result_id":"MAS-RECTOR-023","statement_sha256":"dd2590f85503d95b6ed6280bb6540f982aadade680f17a9ca0298c3abb7fe3c6"}
% HMT_RESULT_END:MAS-RECTOR-023:residence
% HMT_RESULT_BEGIN:MAS-RECTOR-024:residence
% {"material_control":"PASS_PAQUETE_RECTOR_MASAS_HMT_MD_2026_08_06","material_state":"CONSERVADO_CON_APTO_FOCAL","residences":[{"path":"/Users/ruben/Documents/New project/PUBLICACION_HMT/PAQUETE_RECTOR_LEY_GENERAL_MASAS_HMT_MD_2026-08-06/MATRIZ_RECTORA.tsv","sha256":"d4d58d9e5394804b83f48a2996ef884b275669b0f9468a28c9ed7c0a58affdc5"}],"result_id":"MAS-RECTOR-024","statement_sha256":"4d3e6032d279cb845f8bb3fe2e3bcd153b9585996d56fd93c0e721793c358eae"}
% HMT_RESULT_END:MAS-RECTOR-024:residence
% HMT_RESULT_BEGIN:MAS-RECTOR-025:residence
% {"material_control":"PASS_PAQUETE_RECTOR_MASAS_HMT_MD_2026_08_06","material_state":"CONSERVADO_CON_APTO_FOCAL","residences":[{"path":"/Users/ruben/Documents/New project/PUBLICACION_HMT/PAQUETE_RECTOR_LEY_GENERAL_MASAS_HMT_MD_2026-08-06/MATRIZ_RECTORA.tsv","sha256":"d4d58d9e5394804b83f48a2996ef884b275669b0f9468a28c9ed7c0a58affdc5"}],"result_id":"MAS-RECTOR-025","statement_sha256":"6e3bb65bdf6f87c71551bf92d9f6b8e7c60ba9785857316070c5d88c9a59fb9a"}
% HMT_RESULT_END:MAS-RECTOR-025:residence
% HMT_RESULT_BEGIN:MAS-RECTOR-026:residence
% {"material_control":"PASS_PAQUETE_RECTOR_MASAS_HMT_MD_2026_08_06","material_state":"CONSERVADO_CON_APTO_FOCAL","residences":[{"path":"/Users/ruben/Documents/New project/PUBLICACION_HMT/PAQUETE_RECTOR_LEY_GENERAL_MASAS_HMT_MD_2026-08-06/MATRIZ_RECTORA.tsv","sha256":"d4d58d9e5394804b83f48a2996ef884b275669b0f9468a28c9ed7c0a58affdc5"}],"result_id":"MAS-RECTOR-026","statement_sha256":"33fe0ca572be6ea55887ec4b42e78250c112e42699de77e13eaf5ddfdf608e12"}
% HMT_RESULT_END:MAS-RECTOR-026:residence
% HMT_RESULT_BEGIN:MAS-RECTOR-027:residence
% {"material_control":"PASS_PAQUETE_RECTOR_MASAS_HMT_MD_2026_08_06","material_state":"CONSERVADO_CON_APTO_FOCAL","residences":[{"path":"/Users/ruben/Documents/New project/PUBLICACION_HMT/PAQUETE_RECTOR_LEY_GENERAL_MASAS_HMT_MD_2026-08-06/MATRIZ_RECTORA.tsv","sha256":"d4d58d9e5394804b83f48a2996ef884b275669b0f9468a28c9ed7c0a58affdc5"}],"result_id":"MAS-RECTOR-027","statement_sha256":"6168ae29b3e71bf7885c166f21fc815528ad39b20fa81a628be34700e902949c"}
% HMT_RESULT_END:MAS-RECTOR-027:residence
% HMT_RESULT_BEGIN:MAS-RECTOR-028:residence
% {"material_control":"PASS_PAQUETE_RECTOR_MASAS_HMT_MD_2026_08_06","material_state":"CONSERVADO_CON_APTO_FOCAL","residences":[{"path":"/Users/ruben/Documents/New project/PUBLICACION_HMT/PAQUETE_RECTOR_LEY_GENERAL_MASAS_HMT_MD_2026-08-06/MATRIZ_RECTORA.tsv","sha256":"d4d58d9e5394804b83f48a2996ef884b275669b0f9468a28c9ed7c0a58affdc5"}],"result_id":"MAS-RECTOR-028","statement_sha256":"52f3251b7aec29d108b1b6c06d3cc064189cfe5cad74e31c311b1eab0be4b3d1"}
% HMT_RESULT_END:MAS-RECTOR-028:residence
% HMT_RESULT_BEGIN:MAS-PUBLIC-023:residence
% {"material_control":"PASS_PAQUETE_RECTOR_MASAS_HMT_MD_2026_08_06","material_state":"CONSERVADO_CON_APTO_FOCAL","residences":[{"path":"/Users/ruben/Documents/New project/PUBLICACION_HMT/PAQUETE_RECTOR_LEY_GENERAL_MASAS_HMT_MD_2026-08-06/MATRIZ_RECTORA.tsv","sha256":"d4d58d9e5394804b83f48a2996ef884b275669b0f9468a28c9ed7c0a58affdc5"}],"result_id":"MAS-PUBLIC-023","statement_sha256":"2490d5f9a177b2e5c480be8158dc8b1bccbb5f0785cb1e610be8e85f3a6b8918"}
% HMT_RESULT_END:MAS-PUBLIC-023:residence
% HMT_RESULT_BEGIN:MAS-COND-017:residence
% {"material_control":"PASS_PAQUETE_RECTOR_MASAS_HMT_MD_2026_08_06","material_state":"CONSERVADO_CON_APTO_FOCAL","residences":[{"path":"/Users/ruben/Documents/New project/PUBLICACION_HMT/PAQUETE_RECTOR_LEY_GENERAL_MASAS_HMT_MD_2026-08-06/MATRIZ_RECTORA.tsv","sha256":"d4d58d9e5394804b83f48a2996ef884b275669b0f9468a28c9ed7c0a58affdc5"}],"result_id":"MAS-COND-017","statement_sha256":"d9b0c82973bb442a515af30c9bfe351d9b9ab183eb81db8250f32c615d5f50ba"}
% HMT_RESULT_END:MAS-COND-017:residence
% HMT_RESULT_BEGIN:INF-001:residence
% {"material_control":"GATE_MATERIAL_64+PASS_PUBLICACION_INTEGRAL","material_state":"INTEGRADO_O_REMITIDO_SIN_PERDIDA","residences":[{"path":"manuscrito/sections/hmt/05_tpk_numero_medicion.tex","sha256":"51deffafc13a6d83bc07e15f77892d1ec32ba05f54cc2b01d6417ea13e3599a2"}],"result_id":"INF-001","statement_sha256":"d2d4fc4dc0a639778129fc6046b7f974c7a2cc39dae5b6aede5bd2ce96d91b1d"}
% HMT_RESULT_END:INF-001:residence
% HMT_RESULT_BEGIN:INF-002:residence
% {"material_control":"GATE_MATERIAL_64+PASS_PUBLICACION_INTEGRAL","material_state":"INTEGRADO_O_REMITIDO_SIN_PERDIDA","residences":[{"path":"manuscrito/sections/hmt/05_tpk_numero_medicion.tex","sha256":"51deffafc13a6d83bc07e15f77892d1ec32ba05f54cc2b01d6417ea13e3599a2"}],"result_id":"INF-002","statement_sha256":"8e32b4ab88538fa1ca2b52a628278c2f084dc54f423cf78093a3270700d157a5"}
% HMT_RESULT_END:INF-002:residence
% HMT_RESULT_BEGIN:INF-003:residence
% {"material_control":"GATE_MATERIAL_64+PASS_PUBLICACION_INTEGRAL","material_state":"INTEGRADO_O_REMITIDO_SIN_PERDIDA","residences":[{"path":"manuscrito/sections/hmt/05_tpk_numero_medicion.tex","sha256":"51deffafc13a6d83bc07e15f77892d1ec32ba05f54cc2b01d6417ea13e3599a2"}],"result_id":"INF-003","statement_sha256":"667c50d6578b21e1abcff2bc241db42bbe6e270c7268fef0d64172e29d096836"}
% HMT_RESULT_END:INF-003:residence
% HMT_RESULT_BEGIN:N9-001:residence
% {"material_control":"GATE_MATERIAL_64+PASS_PUBLICACION_INTEGRAL","material_state":"INTEGRADO_O_REMITIDO_SIN_PERDIDA","residences":[{"path":"manuscrito/sections/hmt/06aa_genealogia_filtracion_g9_autonoma.tex","sha256":"ba4816ac627ae2d7eb4141421a19a240bf66022f8728f4ce8bf873bd439efb39"}],"result_id":"N9-001","statement_sha256":"f38e038d91253c1364af870bc92f1e5a243d0a5e5c69c5f85bc4d05cc2ef6fb9"}
% HMT_RESULT_END:N9-001:residence
% HMT_RESULT_BEGIN:N9-002:residence
% {"material_control":"GATE_MATERIAL_64+PASS_PUBLICACION_INTEGRAL","material_state":"INTEGRADO_O_REMITIDO_SIN_PERDIDA","residences":[{"path":"manuscrito/sections/hmt/06aa_genealogia_filtracion_g9_autonoma.tex","sha256":"ba4816ac627ae2d7eb4141421a19a240bf66022f8728f4ce8bf873bd439efb39"}],"result_id":"N9-002","statement_sha256":"e4c29231470c3fe69fe80b927a6249121b8a79287a407ca4318e1889a8ed9fb6"}
% HMT_RESULT_END:N9-002:residence
% HMT_RESULT_BEGIN:N9-003:residence
% {"material_control":"GATE_MATERIAL_64+PASS_PUBLICACION_INTEGRAL","material_state":"INTEGRADO_O_REMITIDO_SIN_PERDIDA","residences":[{"path":"manuscrito/sections/hmt/06b_cinco_ciclos_pi.tex","sha256":"83f8a2be69639ac01847a0e25720d3e58b193ae689abfe5fd567a552fb301012"},{"path":"manuscrito/sections/hmt/06d_biografias_numeros_rev6.tex","sha256":"a01607a1a6f5ee4fadbd482daa384b09b192d315da15048ff661d2baf7d197f6"},{"path":"manuscrito/sections/hmt/06e_certificado_generacion_monodromica_rev7.tex","sha256":"2a0cf8028e12605a97923c53726fcb62c1798c458af8e9f78f12c1129b0e1b94"}],"result_id":"N9-003","statement_sha256":"a3b1e710d9bbef1129d8ddb4da929c36fe0b2eaf7ea7e20bbc91312eff72cb57"}
% HMT_RESULT_END:N9-003:residence
% HMT_RESULT_BEGIN:N9-004:residence
% {"material_control":"GATE_MATERIAL_64+PASS_PUBLICACION_INTEGRAL","material_state":"INTEGRADO_O_REMITIDO_SIN_PERDIDA","residences":[{"path":"manuscrito/sections/hmt/06e_certificado_generacion_monodromica_rev7.tex","sha256":"2a0cf8028e12605a97923c53726fcb62c1798c458af8e9f78f12c1129b0e1b94"}],"result_id":"N9-004","statement_sha256":"5288a67a9b48b60a9a1896b482933707612c1ae1e4b4862ba2b23f9ff393e77f"}
% HMT_RESULT_END:N9-004:residence
% HMT_RESULT_BEGIN:HEL-001:residence
% {"material_control":"GATE_MATERIAL_64+PASS_PUBLICACION_INTEGRAL","material_state":"INTEGRADO_O_REMITIDO_SIN_PERDIDA","residences":[{"path":"manuscrito/sections/hmt/03l_elevacion_cohomologica_dimension_rev2.tex","sha256":"1b0f6b86dc275a3947b64e319b0b19789d62c83bcc3714d5c0b8a2cea2243a8c"}],"result_id":"HEL-001","statement_sha256":"3ed30cd71d3b16fcdaf015298295033782765571ae3a2c7d60f65aa8b0bee52d"}
% HMT_RESULT_END:HEL-001:residence
% HMT_RESULT_BEGIN:OPS-001:residence
% {"material_control":"GATE_MATERIAL_64+PASS_PUBLICACION_INTEGRAL","material_state":"INTEGRADO_O_REMITIDO_SIN_PERDIDA","residences":[{"path":"manuscrito/sections/hmt/03k_genealogia_operaciones_rev2.tex","sha256":"fae7aca94c6fb20754431ed79a5aa2f005a7cf2f5681736fc5826da4c694fb98"}],"result_id":"OPS-001","statement_sha256":"61ed21a147f32626c42baccde2ffe7594b14cf6c1c83912b96d5f3f8776405e5"}
% HMT_RESULT_END:OPS-001:residence
% HMT_RESULT_BEGIN:ZF-001:residence
% {"material_control":"GATE_MATERIAL_64+PASS_PUBLICACION_INTEGRAL","material_state":"INTEGRADO_O_REMITIDO_SIN_PERDIDA","residences":[{"path":"manuscrito/sections/integracion_20260812/von_neumann_64/04_zfc_fundacion_ordinales.tex","sha256":"5d81fd31bf31a950f9d068f278e76ef0ebfb37cb8cf6ee5483575ab87f5fddb0"}],"result_id":"ZF-001","statement_sha256":"1d6985e3c4c3690eb43465daa6bcc3d7fb39cb3d0ae1fe31a3cd6c70db621314"}
% HMT_RESULT_END:ZF-001:residence
% HMT_RESULT_BEGIN:ZF-002:residence
% {"material_control":"GATE_MATERIAL_64+PASS_PUBLICACION_INTEGRAL","material_state":"INTEGRADO_O_REMITIDO_SIN_PERDIDA","residences":[{"path":"manuscrito/sections/integracion_20260812/von_neumann_64/04_zfc_fundacion_ordinales.tex","sha256":"5d81fd31bf31a950f9d068f278e76ef0ebfb37cb8cf6ee5483575ab87f5fddb0"}],"result_id":"ZF-002","statement_sha256":"e78afa0785eb7978d61364af148ab43c40092242ad1edf80db0e5c19f11915cd"}
% HMT_RESULT_END:ZF-002:residence
% HMT_RESULT_BEGIN:EXT-001:residence
% {"material_control":"GATE_MATERIAL_64+PASS_PUBLICACION_INTEGRAL","material_state":"INTEGRADO_O_REMITIDO_SIN_PERDIDA","residences":[{"path":"manuscrito/sections/hmt/14d_ontologia_numeros_rev11.tex","sha256":"c96b3c94af4668f15c6a061260d1c1900c52b69b8a4a5d2442268f6ff1d0abcc"}],"result_id":"EXT-001","statement_sha256":"664eeafa4a1a0cfad63aa9414daaccdf8f0ea8a15f92851a674ad7cd5d26c9a8"}
% HMT_RESULT_END:EXT-001:residence
% HMT_RESULT_BEGIN:FND-001:residence
% {"material_control":"GATE_MATERIAL_64+PASS_PUBLICACION_INTEGRAL","material_state":"INTEGRADO_O_REMITIDO_SIN_PERDIDA","residences":[{"path":"manuscrito/sections/integracion_20260812/von_neumann_64/04_zfc_fundacion_ordinales.tex","sha256":"5d81fd31bf31a950f9d068f278e76ef0ebfb37cb8cf6ee5483575ab87f5fddb0"},{"path":"manuscrito/sections/hmt/00a_continuo_numero_medida_hmt_revfinal.tex","sha256":"1be7648180a45eb337db43d6d18a63624876b9f568ea29998c80b00d13b1a819"}],"result_id":"FND-001","statement_sha256":"6b0867dcae5d6924d26cbfdfb829827b9b60d7461148fab933762bd8ab425ac7"}
% HMT_RESULT_END:FND-001:residence
% HMT_RESULT_BEGIN:ORD-001:residence
% {"material_control":"GATE_MATERIAL_64+PASS_PUBLICACION_INTEGRAL","material_state":"INTEGRADO_O_REMITIDO_SIN_PERDIDA","residences":[{"path":"manuscrito/sections/integracion_20260812/von_neumann_64/04_zfc_fundacion_ordinales.tex","sha256":"5d81fd31bf31a950f9d068f278e76ef0ebfb37cb8cf6ee5483575ab87f5fddb0"},{"path":"manuscrito/sections/hmt/14f_cierre_genealogico_numero.tex","sha256":"4522230d8297a1420ccb586d0d6c1fefcb08b5e8e2deeda85b11cc73afd8ded5"}],"result_id":"ORD-001","statement_sha256":"1c7c37b7877465ee0663efb162f12aa0614e44999754b59066edcfa489f7f6ae"}
% HMT_RESULT_END:ORD-001:residence
% HMT_RESULT_BEGIN:ORD-002:residence
% {"material_control":"GATE_MATERIAL_64+PASS_PUBLICACION_INTEGRAL","material_state":"INTEGRADO_O_REMITIDO_SIN_PERDIDA","residences":[{"path":"manuscrito/sections/integracion_20260812/von_neumann_64/04_zfc_fundacion_ordinales.tex","sha256":"5d81fd31bf31a950f9d068f278e76ef0ebfb37cb8cf6ee5483575ab87f5fddb0"}],"result_id":"ORD-002","statement_sha256":"c8188c7c87197d0287f9d7bde9efd04e7246384d5271159053cefc1b62278bf1"}
% HMT_RESULT_END:ORD-002:residence
% HMT_RESULT_BEGIN:PWR-001:residence
% {"material_control":"GATE_MATERIAL_64+PASS_PUBLICACION_INTEGRAL","material_state":"INTEGRADO_O_REMITIDO_SIN_PERDIDA","residences":[{"path":"manuscrito/sections/integracion_20260812/von_neumann_64/05_potencia_eleccion.tex","sha256":"cb412bdee19cc8e65181f75c0666a966dc20111f1dafcaee563bf1ee0e3d2684"}],"result_id":"PWR-001","statement_sha256":"15e7dd2441fdb8f717dbe8bbfb825a15a1165399b2158782c52216b71f1cb8e4"}
% HMT_RESULT_END:PWR-001:residence
% HMT_RESULT_BEGIN:PWR-002:residence
% {"material_control":"GATE_MATERIAL_64+PASS_PUBLICACION_INTEGRAL","material_state":"INTEGRADO_O_REMITIDO_SIN_PERDIDA","residences":[{"path":"manuscrito/sections/integracion_20260812/von_neumann_64/05_potencia_eleccion.tex","sha256":"cb412bdee19cc8e65181f75c0666a966dc20111f1dafcaee563bf1ee0e3d2684"}],"result_id":"PWR-002","statement_sha256":"667fc09531d0de7134ded48a16fb499d56212ba6d7734d3cbc1e04ec1144ad52"}
% HMT_RESULT_END:PWR-002:residence
% HMT_RESULT_BEGIN:CHO-001:residence
% {"material_control":"GATE_MATERIAL_64+PASS_PUBLICACION_INTEGRAL","material_state":"INTEGRADO_O_REMITIDO_SIN_PERDIDA","residences":[{"path":"manuscrito/sections/integracion_20260812/von_neumann_64/05_potencia_eleccion.tex","sha256":"cb412bdee19cc8e65181f75c0666a966dc20111f1dafcaee563bf1ee0e3d2684"},{"path":"manuscrito/sections/hmt/14f_cierre_genealogico_numero.tex","sha256":"4522230d8297a1420ccb586d0d6c1fefcb08b5e8e2deeda85b11cc73afd8ded5"}],"result_id":"CHO-001","statement_sha256":"8b2d6fc14cefc2d7b601d295b4441c23d96827e08dd639c47877a67cb5b36929"}
% HMT_RESULT_END:CHO-001:residence
% HMT_RESULT_BEGIN:CHO-002:residence
% {"material_control":"GATE_MATERIAL_64+PASS_PUBLICACION_INTEGRAL","material_state":"INTEGRADO_O_REMITIDO_SIN_PERDIDA","residences":[{"path":"manuscrito/sections/integracion_20260812/von_neumann_64/05_potencia_eleccion.tex","sha256":"cb412bdee19cc8e65181f75c0666a966dc20111f1dafcaee563bf1ee0e3d2684"}],"result_id":"CHO-002","statement_sha256":"26fee3b2fa8f6cd1bc1073d172a9682b085c58a85cae5e2a6864e89d45f7ddea"}
% HMT_RESULT_END:CHO-002:residence
% HMT_RESULT_BEGIN:CAT-001:residence
% {"material_control":"GATE_MATERIAL_64+PASS_PUBLICACION_INTEGRAL","material_state":"INTEGRADO_O_REMITIDO_SIN_PERDIDA","residences":[{"path":"manuscrito/sections/integracion_20260812/von_neumann_64/04_zfc_fundacion_ordinales.tex","sha256":"5d81fd31bf31a950f9d068f278e76ef0ebfb37cb8cf6ee5483575ab87f5fddb0"},{"path":"manuscrito/sections/hmt/00b_complejo_discreto_medida_rev2.tex","sha256":"c4bd51e17232daf45c4fe95b5863bfa0f945f2a926d22ec6c9c8eccdfa5f1dd8"},{"path":"manuscrito/sections/hmt/03e_funtor_fase_longitud.tex","sha256":"b1a874c3ecb9a434f01fa3ed05c0cef286ac57a85a9269945a994c08df70e78b"},{"path":"manuscrito/sections/hmt/14f_cierre_genealogico_numero.tex","sha256":"4522230d8297a1420ccb586d0d6c1fefcb08b5e8e2deeda85b11cc73afd8ded5"}],"result_id":"CAT-001","statement_sha256":"cd69268cc61050d66556a6a33bd6fb31ab1fc444f678229ff48d09dfeeb7ec1b"}
% HMT_RESULT_END:CAT-001:residence
% HMT_RESULT_BEGIN:TOP-001:residence
% {"material_control":"GATE_MATERIAL_64+PASS_PUBLICACION_INTEGRAL","material_state":"INTEGRADO_O_REMITIDO_SIN_PERDIDA","residences":[{"path":"manuscrito/sections/integracion_20260812/von_neumann_64/04_zfc_fundacion_ordinales.tex","sha256":"5d81fd31bf31a950f9d068f278e76ef0ebfb37cb8cf6ee5483575ab87f5fddb0"},{"path":"manuscrito/sections/hmt/05_tpk_numero_medicion.tex","sha256":"51deffafc13a6d83bc07e15f77892d1ec32ba05f54cc2b01d6417ea13e3599a2"}],"result_id":"TOP-001","statement_sha256":"cf2c0f501167aac25e8c173e2c9c6cca8aa4cc123f668b06c5a7ca8c42f21eb3"}
% HMT_RESULT_END:TOP-001:residence
% HMT_RESULT_BEGIN:HIL-001:residence
% {"material_control":"GATE_MATERIAL_64+PASS_PUBLICACION_INTEGRAL","material_state":"INTEGRADO_O_REMITIDO_SIN_PERDIDA","residences":[{"path":"manuscrito/sections/integracion_20260812/von_neumann_64/06_hilbert_heisenberg.tex","sha256":"6f35e2e9b56af3cbf33d20c28ae1b5bfac2b92f4a0827f444be39774046e130b"},{"path":"manuscrito/sections/hmt/03l_elevacion_cohomologica_dimension_rev2.tex","sha256":"1b0f6b86dc275a3947b64e319b0b19789d62c83bcc3714d5c0b8a2cea2243a8c"}],"result_id":"HIL-001","statement_sha256":"7529a7dd47eb815a976287f21bd184c40d9cc3dfb0dcb1d2c3f1a2b4b4eeaaa7"}
% HMT_RESULT_END:HIL-001:residence
% HMT_RESULT_BEGIN:UHF-001:residence
% {"material_control":"GATE_MATERIAL_64+PASS_PUBLICACION_INTEGRAL","material_state":"INTEGRADO_O_REMITIDO_SIN_PERDIDA","residences":[{"path":"manuscrito/sections/integracion_20260812/von_neumann_64/07_uhf_factor_ii1.tex","sha256":"a3f413168ac1514605dc7e9bf8a7866e73420380c5a93390fbd5f5d55fa25765"},{"path":"manuscrito/sections/hmt/03l_elevacion_cohomologica_dimension_rev2.tex","sha256":"1b0f6b86dc275a3947b64e319b0b19789d62c83bcc3714d5c0b8a2cea2243a8c"}],"result_id":"UHF-001","statement_sha256":"0cf650bcfd502563a4ff8cab59c8c8c4ae997565ebed1f058034f5cce7e5b378"}
% HMT_RESULT_END:UHF-001:residence
% HMT_RESULT_BEGIN:II1-001:residence
% {"material_control":"GATE_MATERIAL_64+PASS_PUBLICACION_INTEGRAL","material_state":"INTEGRADO_O_REMITIDO_SIN_PERDIDA","residences":[{"path":"manuscrito/sections/integracion_20260812/von_neumann_64/07_uhf_factor_ii1.tex","sha256":"a3f413168ac1514605dc7e9bf8a7866e73420380c5a93390fbd5f5d55fa25765"},{"path":"manuscrito/sections/hmt/03l_elevacion_cohomologica_dimension_rev2.tex","sha256":"1b0f6b86dc275a3947b64e319b0b19789d62c83bcc3714d5c0b8a2cea2243a8c"}],"result_id":"II1-001","statement_sha256":"dc01c84423637eb3704db747bb7a96106332b3cc6f8f36258e1af6cf02c8ead6"}
% HMT_RESULT_END:II1-001:residence
% HMT_RESULT_BEGIN:DIM-001:residence
% {"material_control":"GATE_MATERIAL_64+PASS_PUBLICACION_INTEGRAL","material_state":"INTEGRADO_O_REMITIDO_SIN_PERDIDA","residences":[{"path":"manuscrito/sections/integracion_20260812/von_neumann_64/07_uhf_factor_ii1.tex","sha256":"a3f413168ac1514605dc7e9bf8a7866e73420380c5a93390fbd5f5d55fa25765"},{"path":"manuscrito/sections/hmt/03l_elevacion_cohomologica_dimension_rev2.tex","sha256":"1b0f6b86dc275a3947b64e319b0b19789d62c83bcc3714d5c0b8a2cea2243a8c"}],"result_id":"DIM-001","statement_sha256":"9b8936bc39190cf70a463a94ca6fb311c2d8d8e66262fb2a3d44d831c5653c4b"}
% HMT_RESULT_END:DIM-001:residence
% HMT_RESULT_BEGIN:CLK-001:residence
% {"material_control":"GATE_MATERIAL_64+PASS_PUBLICACION_INTEGRAL","material_state":"INTEGRADO_O_REMITIDO_SIN_PERDIDA","residences":[{"path":"manuscrito/sections/integracion_20260812/von_neumann_64/08_reloj_logica_ergodicidad.tex","sha256":"aeda0eee8e4376ae44e5f6505cf4414ba767a64705a111c66ef2917d7411a0d9"},{"path":"manuscrito/sections/hmt/19_rotaciones_vacancias_criticidad.tex","sha256":"2d46e077e597811a4b0545b2a6ded98b5d56d0d66c65d71f416a1d06c282c363"}],"result_id":"CLK-001","statement_sha256":"b89435d53375a0c540e471c24439fbb5b83ca598038eaa6e60ad19138fe1d106"}
% HMT_RESULT_END:CLK-001:residence
% HMT_RESULT_BEGIN:INT-001:residence
% {"material_control":"GATE_MATERIAL_64+PASS_PUBLICACION_INTEGRAL","material_state":"INTEGRADO_O_REMITIDO_SIN_PERDIDA","residences":[{"path":"manuscrito/sections/integracion_20260812/von_neumann_64/08_reloj_logica_ergodicidad.tex","sha256":"aeda0eee8e4376ae44e5f6505cf4414ba767a64705a111c66ef2917d7411a0d9"}],"result_id":"INT-001","statement_sha256":"aa1ffba490fcc7437280d3a98fe5f9a543b85b037100711af707edb8fdc41527"}
% HMT_RESULT_END:INT-001:residence
% HMT_RESULT_BEGIN:QLG-001:residence
% {"material_control":"GATE_MATERIAL_64+PASS_PUBLICACION_INTEGRAL","material_state":"INTEGRADO_O_REMITIDO_SIN_PERDIDA","residences":[{"path":"manuscrito/sections/integracion_20260812/von_neumann_64/08_reloj_logica_ergodicidad.tex","sha256":"aeda0eee8e4376ae44e5f6505cf4414ba767a64705a111c66ef2917d7411a0d9"},{"path":"manuscrito/sections/integracion_20260812/von_neumann_64/06_hilbert_heisenberg.tex","sha256":"6f35e2e9b56af3cbf33d20c28ae1b5bfac2b92f4a0827f444be39774046e130b"},{"path":"manuscrito/sections/hmt/03_app_ortograma.tex","sha256":"aaf0fd1dec2c270d544cd19562b4461c57ebb009437e051d973be43395428390"}],"result_id":"QLG-001","statement_sha256":"5af886b5ee8e9ef572b0db2f8860b44c9a5685f03b042ad7934083b6157af5d3"}
% HMT_RESULT_END:QLG-001:residence
% HMT_RESULT_BEGIN:ERG-001:residence
% {"material_control":"GATE_MATERIAL_64+PASS_PUBLICACION_INTEGRAL","material_state":"INTEGRADO_O_REMITIDO_SIN_PERDIDA","residences":[{"path":"manuscrito/sections/integracion_20260812/von_neumann_64/08_reloj_logica_ergodicidad.tex","sha256":"aeda0eee8e4376ae44e5f6505cf4414ba767a64705a111c66ef2917d7411a0d9"},{"path":"manuscrito/sections/hmt/19_rotaciones_vacancias_criticidad.tex","sha256":"2d46e077e597811a4b0545b2a6ded98b5d56d0d66c65d71f416a1d06c282c363"}],"result_id":"ERG-001","statement_sha256":"0e5e1bd65b68a7b812b621df955227e4bd7880b25b43edc0e1a9db24e3a80b8e"}
% HMT_RESULT_END:ERG-001:residence
% HMT_RESULT_BEGIN:DEN-001:residence
% {"material_control":"GATE_MATERIAL_64+PASS_PUBLICACION_INTEGRAL","material_state":"INTEGRADO_O_REMITIDO_SIN_PERDIDA","residences":[{"path":"manuscrito/sections/integracion_20260812/von_neumann_64/08_reloj_logica_ergodicidad.tex","sha256":"aeda0eee8e4376ae44e5f6505cf4414ba767a64705a111c66ef2917d7411a0d9"},{"path":"manuscrito/sections/md/04ab_postulados_cuanticos_genealogicos_rev2.tex","sha256":"1895f84bd75b104e873a3d73c5b0e8ae1b5df606bd6feefeb4f9666a92966804"}],"result_id":"DEN-001","statement_sha256":"d3f54d28f9a40e95ef70680ed00ece354f13bef3fa9616bf65771ff984899f55"}
% HMT_RESULT_END:DEN-001:residence
% HMT_RESULT_BEGIN:LAN-001:residence
% {"material_control":"GATE_MATERIAL_64+PASS_PUBLICACION_INTEGRAL","material_state":"INTEGRADO_O_REMITIDO_SIN_PERDIDA","residences":[{"path":"manuscrito/sections/md/04g_informacion_reversible_landauer_rev7.tex","sha256":"3f60a3da983f28f197fe487896eda3667d44247d0ad30f6f3d1a598b368f0e8e"}],"result_id":"LAN-001","statement_sha256":"197b4b23802d24d8e911bcd3787ba07daf3a5694505b103e27acf6a540021b23"}
% HMT_RESULT_END:LAN-001:residence
% HMT_RESULT_BEGIN:LAN-002:residence
% {"material_control":"GATE_MATERIAL_64+PASS_PUBLICACION_INTEGRAL","material_state":"INTEGRADO_O_REMITIDO_SIN_PERDIDA","residences":[{"path":"manuscrito/sections/integracion_20260812/bloque_64_landauer.tex","sha256":"ef8b27ff2b888ef3c5aa3228093c84ed139d63a8e8a504a4b6da8b56f256cbe6"}],"result_id":"LAN-002","statement_sha256":"aa495d4ee76ccc4ec19bbfee8fca847aa39dfb3e7785b0654264e883afa8e9f3"}
% HMT_RESULT_END:LAN-002:residence
% HMT_RESULT_BEGIN:GDL-001:residence
% {"material_control":"GATE_MATERIAL_64+PASS_PUBLICACION_INTEGRAL","material_state":"INTEGRADO_O_REMITIDO_SIN_PERDIDA","residences":[{"path":"manuscrito/sections/integracion_20260812/von_neumann_64/10_godel_turing.tex","sha256":"c17b938966df53455f1927258ae4e4fd2603927ea1827a08620dfdbb386e7863"}],"result_id":"GDL-001","statement_sha256":"e65a071d6f27cf044d0d665fa43c1c77e8b1c800787b0879ee63f693a954b4eb"}
% HMT_RESULT_END:GDL-001:residence
% HMT_RESULT_BEGIN:GDL-002:residence
% {"material_control":"GATE_MATERIAL_64+PASS_PUBLICACION_INTEGRAL","material_state":"INTEGRADO_O_REMITIDO_SIN_PERDIDA","residences":[{"path":"manuscrito/sections/integracion_20260812/von_neumann_64/10_godel_turing.tex","sha256":"c17b938966df53455f1927258ae4e4fd2603927ea1827a08620dfdbb386e7863"}],"result_id":"GDL-002","statement_sha256":"30af7308ca4cbed2f47ea0f0aaf875baaf611eda5be904a33538de7e5ee41864"}
% HMT_RESULT_END:GDL-002:residence
% HMT_RESULT_BEGIN:IND-001:residence
% {"material_control":"GATE_MATERIAL_64+PASS_PUBLICACION_INTEGRAL","material_state":"INTEGRADO_O_REMITIDO_SIN_PERDIDA","residences":[{"path":"manuscrito/sections/integracion_20260812/von_neumann_64/11_forcing_borel_minimax.tex","sha256":"c4b14714d216d63c6db3b930602a3cd7953d40dd61d50a093bd50d840ccfc33c"}],"result_id":"IND-001","statement_sha256":"5e22be0c9938a9dceab95188d043eebbbffab47eced16a62dc2ded6048b7891d"}
% HMT_RESULT_END:IND-001:residence
% HMT_RESULT_BEGIN:FOR-001:residence
% {"material_control":"GATE_MATERIAL_64+PASS_PUBLICACION_INTEGRAL","material_state":"INTEGRADO_O_REMITIDO_SIN_PERDIDA","residences":[{"path":"manuscrito/sections/integracion_20260812/von_neumann_64/11_forcing_borel_minimax.tex","sha256":"c4b14714d216d63c6db3b930602a3cd7953d40dd61d50a093bd50d840ccfc33c"},{"path":"manuscrito/sections/hmt/05_tpk_numero_medicion.tex","sha256":"51deffafc13a6d83bc07e15f77892d1ec32ba05f54cc2b01d6417ea13e3599a2"}],"result_id":"FOR-001","statement_sha256":"1a8d72a91afdb3a1411a9d0431a099066dffdd6a4603d5d3bc5cb570330fa919"}
% HMT_RESULT_END:FOR-001:residence
% HMT_RESULT_BEGIN:FOR-002:residence
% {"material_control":"GATE_MATERIAL_64+PASS_PUBLICACION_INTEGRAL","material_state":"INTEGRADO_O_REMITIDO_SIN_PERDIDA","residences":[{"path":"manuscrito/sections/integracion_20260812/von_neumann_64/11_forcing_borel_minimax.tex","sha256":"c4b14714d216d63c6db3b930602a3cd7953d40dd61d50a093bd50d840ccfc33c"}],"result_id":"FOR-002","statement_sha256":"59f4f58f84a7735d60c6fd2454ba7965e566bd09c58b3d45392f878cb3dec369"}
% HMT_RESULT_END:FOR-002:residence
% HMT_RESULT_BEGIN:BOR-001:residence
% {"material_control":"GATE_MATERIAL_64+PASS_PUBLICACION_INTEGRAL","material_state":"INTEGRADO_O_REMITIDO_SIN_PERDIDA","residences":[{"path":"manuscrito/sections/integracion_20260812/von_neumann_64/11_forcing_borel_minimax.tex","sha256":"c4b14714d216d63c6db3b930602a3cd7953d40dd61d50a093bd50d840ccfc33c"}],"result_id":"BOR-001","statement_sha256":"ac8e973306e90a8e1806a33402b6ef6e087583edf2b49e475fcf64fc7bffc827"}
% HMT_RESULT_END:BOR-001:residence
% HMT_RESULT_BEGIN:MIN-001:residence
% {"material_control":"GATE_MATERIAL_64+PASS_PUBLICACION_INTEGRAL","material_state":"INTEGRADO_O_REMITIDO_SIN_PERDIDA","residences":[{"path":"manuscrito/sections/integracion_20260812/von_neumann_64/11_forcing_borel_minimax.tex","sha256":"c4b14714d216d63c6db3b930602a3cd7953d40dd61d50a093bd50d840ccfc33c"}],"result_id":"MIN-001","statement_sha256":"ed4856dacc07059c530b3404157bcc15541d967d4e205eaba615ed2e7f06bd2e"}
% HMT_RESULT_END:MIN-001:residence
% HMT_RESULT_BEGIN:III-001:residence
% {"material_control":"GATE_MATERIAL_64+PASS_PUBLICACION_INTEGRAL","material_state":"INTEGRADO_O_REMITIDO_SIN_PERDIDA","residences":[{"path":"manuscrito/sections/integracion_20260812/von_neumann_64/08_reloj_logica_ergodicidad.tex","sha256":"aeda0eee8e4376ae44e5f6505cf4414ba767a64705a111c66ef2917d7411a0d9"},{"path":"manuscrito/sections/md/04e_termodinamica_rutas.tex","sha256":"ab1429f9145eee89e5287095b94f327291fa4aa247ae0abcc95e3f5611de87c1"}],"result_id":"III-001","statement_sha256":"ed27269e518d94d2dd09835c4f6d5c6be473aa08906e8ed2d633e2686f5bb4a7"}
% HMT_RESULT_END:III-001:residence
% HMT_RESULT_BEGIN:TRI-001:residence
% {"material_control":"GATE_MATERIAL_64+PASS_PUBLICACION_INTEGRAL","material_state":"INTEGRADO_O_REMITIDO_SIN_PERDIDA","residences":[{"path":"manuscrito/sections/integracion_20260812/von_neumann_64/12_cierre_triple_doble_proyeccion.tex","sha256":"9773623f121e57b194d95202c9b01d24af38a7f1e391fd219beb76ce9dbf2a86"},{"path":"manuscrito/sections/md/10c_puente_retorno_holonomia_cartan.tex","sha256":"b206dc7147f5eb4295dace2d1960e61cb5f6c4b40a0f6828535d1ba4024a7ba9"}],"result_id":"TRI-001","statement_sha256":"7717c12ab00203b9ef7e04dfea02661344e850d6889ac9e04b4906e296f8b383"}
% HMT_RESULT_END:TRI-001:residence
% HMT_RESULT_BEGIN:EXC-001:residence
% {"material_control":"GATE_MATERIAL_64+PASS_PUBLICACION_INTEGRAL","material_state":"INTEGRADO_O_REMITIDO_SIN_PERDIDA","residences":[{"path":"manuscrito/sections/integracion_20260812/von_neumann_64/12_cierre_triple_doble_proyeccion.tex","sha256":"9773623f121e57b194d95202c9b01d24af38a7f1e391fd219beb76ce9dbf2a86"},{"path":"manuscrito/sections/hmt/09_doble_proyeccion.tex","sha256":"ff538add4f3afedca7f4cbd5cea553e3b8132e05117117c40be30e55d33dcdba"},{"path":"manuscrito/sections/hmt/15_bandera_y_rama_excepcional.tex","sha256":"6c42a7012881b036232a5e12c6e2ce382d48964ef9ed12cc0a4d5256126154e5"},{"path":"manuscrito/sections/hmt/15c_cierre_micro_macro_nivel_tres.tex","sha256":"0a5cf2476d9825b9005494bfab1b937deddbba26d5c69f76a7822077be3797c5"},{"path":"manuscrito/sections/hmt/15d_voa_moonshine_orden_dos_rev11.tex","sha256":"6bc88137eb614a0bba834212c9f8f95be09e48e8f570f6c8126ac33dbe007e86"}],"result_id":"EXC-001","statement_sha256":"b9c1e5194c3ddd0cd113377935663265ca76fce3648e3ea5576c0e8b41aa6747"}
% HMT_RESULT_END:EXC-001:residence
% HMT_RESULT_BEGIN:MD-001:residence
% {"material_control":"GATE_MATERIAL_64+PASS_PUBLICACION_INTEGRAL","material_state":"INTEGRADO_O_REMITIDO_SIN_PERDIDA","residences":[{"path":"manuscrito/sections/integracion_20260812/von_neumann_64/12_cierre_triple_doble_proyeccion.tex","sha256":"9773623f121e57b194d95202c9b01d24af38a7f1e391fd219beb76ce9dbf2a86"},{"path":"manuscrito/sections/md/05f_operadores_masa_two_way_rev11.tex","sha256":"e0c576ff54f405c41adea7093ca6b117134e9322749c969f1fc4cca6e8615d35"},{"path":"manuscrito/sections/md/06h_atlas_familias_rev6.tex","sha256":"f6cf1c90f4d42d24598380ca0bff5c9ee66c3269aad819c08d85d32a31066367"}],"result_id":"MD-001","statement_sha256":"3def61d573286af3af73231496757ed80209391421e2f5b759252b1e4f75a582"}
% HMT_RESULT_END:MD-001:residence
% HMT_RESULT_BEGIN:PSC-001:residence
% {"material_control":"GATE_MATERIAL_64+PASS_PUBLICACION_INTEGRAL","material_state":"INTEGRADO_O_REMITIDO_SIN_PERDIDA","residences":[{"path":"manuscrito/sections/integracion_20260812/von_neumann_64/11_forcing_borel_minimax.tex","sha256":"c4b14714d216d63c6db3b930602a3cd7953d40dd61d50a093bd50d840ccfc33c"}],"result_id":"PSC-001","statement_sha256":"aa11f37512474aa438b0ea7a0c6719d2337b671cab88d78f0bfbda69d2d029d4"}
% HMT_RESULT_END:PSC-001:residence
% HMT_RESULT_BEGIN:COV-001:residence
% {"material_control":"GATE_MATERIAL_64+PASS_PUBLICACION_INTEGRAL","material_state":"INTEGRADO_O_REMITIDO_SIN_PERDIDA","residences":[{"path":"/Users/ruben/Documents/HMT2/CONTINUIDAD_EDITORIAL/INTEGRACION_REV3_VON_NEUMANN_MATRICIAL_REALIZACION_2026-08-12/MATRIZ_ATOMICA_AUTOCONTENIDA_REV3_V4.tsv","sha256":"a89bbb749d42ca8e123b0549913409dadda6196187891b4bf7f83b993d2ffa9b"}],"result_id":"COV-001","statement_sha256":"2d2a3767b8f662d7a68bd271a690c4b557d2bc3f3a256077632ca30e518bd4d7"}
% HMT_RESULT_END:COV-001:residence
% HMT_RESULT_BEGIN:IINF-001:residence
% {"material_control":"GATE_MATERIAL_64+PASS_PUBLICACION_INTEGRAL","material_state":"INTEGRADO_O_REMITIDO_SIN_PERDIDA","residences":[{"path":"manuscrito/sections/integracion_20260812/von_neumann_64/07_uhf_factor_ii1.tex","sha256":"a3f413168ac1514605dc7e9bf8a7866e73420380c5a93390fbd5f5d55fa25765"}],"result_id":"IINF-001","statement_sha256":"217587fdc70a6738c237837099cbbba8d41f6d59c40619c866bf48c95e0198d7"}
% HMT_RESULT_END:IINF-001:residence
% HMT_RESULT_BEGIN:DEN-003:residence
% {"material_control":"GATE_MATERIAL_64+PASS_PUBLICACION_INTEGRAL","material_state":"INTEGRADO_O_REMITIDO_SIN_PERDIDA","residences":[{"path":"manuscrito/sections/integracion_20260812/von_neumann_64/08_reloj_logica_ergodicidad.tex","sha256":"aeda0eee8e4376ae44e5f6505cf4414ba767a64705a111c66ef2917d7411a0d9"}],"result_id":"DEN-003","statement_sha256":"921dc3fb0d30a2764747bddc32233eb05b1589bbcbb17f4a98f3aae0ea462a13"}
% HMT_RESULT_END:DEN-003:residence
% HMT_RESULT_BEGIN:CARD-001:residence
% {"material_control":"GATE_MATERIAL_64+PASS_PUBLICACION_INTEGRAL","material_state":"INTEGRADO_O_REMITIDO_SIN_PERDIDA","residences":[{"path":"manuscrito/sections/integracion_20260812/von_neumann_64/04_zfc_fundacion_ordinales.tex","sha256":"5d81fd31bf31a950f9d068f278e76ef0ebfb37cb8cf6ee5483575ab87f5fddb0"},{"path":"manuscrito/sections/hmt/00_construccion_app_gnomon_rev11.tex","sha256":"f027dd8cb64274e9f984a180e01da910cdb1fc6235d856ad6dc736f7477381b9"}],"result_id":"CARD-001","statement_sha256":"657e91ba07505f78793eeac979c673b8bc6ec88f1adc9067345b9c06ef806c89"}
% HMT_RESULT_END:CARD-001:residence
% HMT_RESULT_BEGIN:PWR-003:residence
% {"material_control":"GATE_MATERIAL_64+PASS_PUBLICACION_INTEGRAL","material_state":"INTEGRADO_O_REMITIDO_SIN_PERDIDA","residences":[{"path":"manuscrito/sections/integracion_20260812/von_neumann_64/05_potencia_eleccion.tex","sha256":"cb412bdee19cc8e65181f75c0666a966dc20111f1dafcaee563bf1ee0e3d2684"},{"path":"manuscrito/sections/hmt/09_doble_proyeccion.tex","sha256":"ff538add4f3afedca7f4cbd5cea553e3b8132e05117117c40be30e55d33dcdba"}],"result_id":"PWR-003","statement_sha256":"eaefebdd49a0763208853c604924034d7d4c705c6371b2990f7fec1b29b0032a"}
% HMT_RESULT_END:PWR-003:residence
% HMT_RESULT_BEGIN:WEYL-001:residence
% {"material_control":"GATE_MATERIAL_64+PASS_PUBLICACION_INTEGRAL","material_state":"INTEGRADO_O_REMITIDO_SIN_PERDIDA","residences":[{"path":"manuscrito/sections/integracion_20260812/von_neumann_64/06_hilbert_heisenberg.tex","sha256":"6f35e2e9b56af3cbf33d20c28ae1b5bfac2b92f4a0827f444be39774046e130b"},{"path":"manuscrito/sections/hmt/14_numero_heisenberg_y_d108.tex","sha256":"36c7812d22940e488c5a5434ffb85399d9430a12d32f6a376138c58474b3247b"}],"result_id":"WEYL-001","statement_sha256":"3c9c006467a020b44e1f152943e542f7efd63f1fbafa62c1459d799fb5a7fb8d"}
% HMT_RESULT_END:WEYL-001:residence
% HMT_RESULT_BEGIN:DSC-001:residence
% {"material_control":"GATE_MATERIAL_64+PASS_PUBLICACION_INTEGRAL","material_state":"INTEGRADO_O_REMITIDO_SIN_PERDIDA","residences":[{"path":"manuscrito/sections/md/02_pantalla_accion.tex","sha256":"40796ae0bdab2ff0e8c2cc2f07181c0f5efc77c14aa4cca2d5334e899dc0f1ef"}],"result_id":"DSC-001","statement_sha256":"67da6399420386a7c52ecce6942df44aef3e17b49f2b6c2ea0005b3f95cae4ca"}
% HMT_RESULT_END:DSC-001:residence
% HMT_RESULT_BEGIN:REC-001:residence
% {"material_control":"GATE_MATERIAL_64+PASS_PUBLICACION_INTEGRAL","material_state":"INTEGRADO_O_REMITIDO_SIN_PERDIDA","residences":[{"path":"manuscrito/sections/md/04h_entropias_tipadas_rev2.tex","sha256":"f47cec558bb3f375af069dfde73f82b69bf48d9f168a635630f0e9a25bab3d91"}],"result_id":"REC-001","statement_sha256":"0bd633094af48b759b62a734d6cefa9feb253f605c7052145c392284c4b6a32f"}
% HMT_RESULT_END:REC-001:residence
% HMT_RESULT_BEGIN:LAN-003:residence
% {"material_control":"GATE_MATERIAL_64+PASS_PUBLICACION_INTEGRAL","material_state":"INTEGRADO_O_REMITIDO_SIN_PERDIDA","residences":[{"path":"manuscrito/sections/md/04g_informacion_reversible_landauer_rev7.tex","sha256":"3f60a3da983f28f197fe487896eda3667d44247d0ad30f6f3d1a598b368f0e8e"}],"result_id":"LAN-003","statement_sha256":"2a5b4ef14486d6da04b11ef915d58b57331c845e9d106cf402e883204d561861"}
% HMT_RESULT_END:LAN-003:residence
% HMT_RESULT_BEGIN:FIN-001:residence
% {"material_control":"GATE_MATERIAL_64+PASS_PUBLICACION_INTEGRAL","material_state":"INTEGRADO_O_REMITIDO_SIN_PERDIDA","residences":[{"path":"manuscrito/sections/integracion_20260812/von_neumann_64/10_godel_turing.tex","sha256":"c17b938966df53455f1927258ae4e4fd2603927ea1827a08620dfdbb386e7863"}],"result_id":"FIN-001","statement_sha256":"a73d6ef20cd13e14f35306811ecee9eb3057af5c97ddf492c7ba1f86dd04f1ce"}
% HMT_RESULT_END:FIN-001:residence
% HMT_RESULT_BEGIN:DEN-002:residence
% {"material_control":"GATE_MATERIAL_64+PASS_PUBLICACION_INTEGRAL","material_state":"INTEGRADO_O_REMITIDO_SIN_PERDIDA","residences":[{"path":"manuscrito/sections/integracion_20260812/von_neumann_64/08_reloj_logica_ergodicidad.tex","sha256":"aeda0eee8e4376ae44e5f6505cf4414ba767a64705a111c66ef2917d7411a0d9"}],"result_id":"DEN-002","statement_sha256":"35b2ba3c1b4020b9163f1ef3d443b6bcb565b09c8a8622d543cf8e4f57289896"}
% HMT_RESULT_END:DEN-002:residence
% HMT_RESULT_BEGIN:UNI-001:residence
% {"material_control":"GATE_MATERIAL_64+PASS_PUBLICACION_INTEGRAL","material_state":"INTEGRADO_O_REMITIDO_SIN_PERDIDA","residences":[{"path":"manuscrito/sections/integracion_20260812/von_neumann_64/12_cierre_triple_doble_proyeccion.tex","sha256":"9773623f121e57b194d95202c9b01d24af38a7f1e391fd219beb76ce9dbf2a86"}],"result_id":"UNI-001","statement_sha256":"c937c36037e7f572961ab452004513b8c1c29a3ba0d578a7b0da6590a849a422"}
% HMT_RESULT_END:UNI-001:residence
% HMT_RESULT_BEGIN:UNI-002:residence
% {"material_control":"GATE_MATERIAL_64+PASS_PUBLICACION_INTEGRAL","material_state":"INTEGRADO_O_REMITIDO_SIN_PERDIDA","residences":[{"path":"manuscrito/sections/integracion_20260812/von_neumann_64/12_cierre_triple_doble_proyeccion.tex","sha256":"9773623f121e57b194d95202c9b01d24af38a7f1e391fd219beb76ce9dbf2a86"}],"result_id":"UNI-002","statement_sha256":"6c0f67e5980d3e7df62bb7613076467f3e8c1e71d7bd85ae9b8e5a5a6e6eefac"}
% HMT_RESULT_END:UNI-002:residence
% HMT_RESULT_BEGIN:P01:residence
% {"material_control":"PASS_PUENTE_GENEALOGICO_MATRICIAL","material_state":"INTEGRADO_EN_GRAFO_ACTIVO","residences":[{"path":"manuscrito/sections/integracion_20260812/matricial_52/02_hmt_en_profundidad.tex","sha256":"fe6c23eac09e147199090a091fe2564e5587b7b1155ea85518d2bf47aade4cfc"}],"result_id":"P01","statement_sha256":"53d8dc07ede69f64f0d718b277459101f4824f60da865637ded87662edab5245"}
% HMT_RESULT_END:P01:residence
% HMT_RESULT_BEGIN:P02:residence
% {"material_control":"PASS_PUENTE_GENEALOGICO_MATRICIAL","material_state":"INTEGRADO_EN_GRAFO_ACTIVO","residences":[{"path":"manuscrito/sections/integracion_20260812/matricial_52/02_hmt_en_profundidad.tex","sha256":"fe6c23eac09e147199090a091fe2564e5587b7b1155ea85518d2bf47aade4cfc"}],"result_id":"P02","statement_sha256":"69c00b394120ef47751db68d34eb4b918078692555e4fa5e663b2f0297a8548c"}
% HMT_RESULT_END:P02:residence
% HMT_RESULT_BEGIN:P03:residence
% {"material_control":"PASS_PUENTE_GENEALOGICO_MATRICIAL","material_state":"INTEGRADO_EN_GRAFO_ACTIVO","residences":[{"path":"manuscrito/sections/integracion_20260812/matricial_52/02_hmt_en_profundidad.tex","sha256":"fe6c23eac09e147199090a091fe2564e5587b7b1155ea85518d2bf47aade4cfc"}],"result_id":"P03","statement_sha256":"2fe6041370eb6e955a3d0ba0f2e07811bc0aebd09c6793150785e46a510b9359"}
% HMT_RESULT_END:P03:residence
% HMT_RESULT_BEGIN:P04:residence
% {"material_control":"PASS_PUENTE_GENEALOGICO_MATRICIAL","material_state":"INTEGRADO_EN_GRAFO_ACTIVO","residences":[{"path":"manuscrito/sections/integracion_20260812/matricial_52/02_hmt_en_profundidad.tex","sha256":"fe6c23eac09e147199090a091fe2564e5587b7b1155ea85518d2bf47aade4cfc"}],"result_id":"P04","statement_sha256":"c9c80bd3c7cc1697f5e94b0d7401b0649ab52d4f48fd42943fe4bd715aa71c05"}
% HMT_RESULT_END:P04:residence
% HMT_RESULT_BEGIN:P05:residence
% {"material_control":"PASS_PUENTE_GENEALOGICO_MATRICIAL","material_state":"INTEGRADO_EN_GRAFO_ACTIVO","residences":[{"path":"manuscrito/sections/integracion_20260812/matricial_52/02_hmt_en_profundidad.tex","sha256":"fe6c23eac09e147199090a091fe2564e5587b7b1155ea85518d2bf47aade4cfc"}],"result_id":"P05","statement_sha256":"87558450ea6ada90154488b513914bbcbe56a44e45cb40411baa8054fa0a9b11"}
% HMT_RESULT_END:P05:residence
% HMT_RESULT_BEGIN:P06:residence
% {"material_control":"PASS_PUENTE_GENEALOGICO_MATRICIAL","material_state":"INTEGRADO_EN_GRAFO_ACTIVO","residences":[{"path":"manuscrito/sections/integracion_20260812/matricial_52/02_hmt_en_profundidad.tex","sha256":"fe6c23eac09e147199090a091fe2564e5587b7b1155ea85518d2bf47aade4cfc"}],"result_id":"P06","statement_sha256":"b3609e5998fdaf08a9020b1ff80a65442e473516d2ecf24b5b68474c1f6c94c0"}
% HMT_RESULT_END:P06:residence
% HMT_RESULT_BEGIN:P07:residence
% {"material_control":"PASS_PUENTE_GENEALOGICO_MATRICIAL","material_state":"INTEGRADO_EN_GRAFO_ACTIVO","residences":[{"path":"manuscrito/sections/integracion_20260812/matricial_52/02_hmt_en_profundidad.tex","sha256":"fe6c23eac09e147199090a091fe2564e5587b7b1155ea85518d2bf47aade4cfc"},{"path":"manuscrito/sections/md/04ab_postulados_cuanticos_genealogicos_rev2.tex","sha256":"1895f84bd75b104e873a3d73c5b0e8ae1b5df606bd6feefeb4f9666a92966804"}],"result_id":"P07","statement_sha256":"f8147d53ff5923fc1c43835f30a9f07bf581e5e699eb63dd74ac1414533a629e"}
% HMT_RESULT_END:P07:residence
% HMT_RESULT_BEGIN:P08:residence
% {"material_control":"PASS_PUENTE_GENEALOGICO_MATRICIAL","material_state":"INTEGRADO_EN_GRAFO_ACTIVO","residences":[{"path":"manuscrito/sections/integracion_20260812/matricial_52/02_hmt_en_profundidad.tex","sha256":"fe6c23eac09e147199090a091fe2564e5587b7b1155ea85518d2bf47aade4cfc"}],"result_id":"P08","statement_sha256":"f9a8d4c5ec3b2afef25f422ad348af58e28287d1d33e38fd8b67b7529c84166d"}
% HMT_RESULT_END:P08:residence
% HMT_RESULT_BEGIN:P09:residence
% {"material_control":"PASS_PUENTE_GENEALOGICO_MATRICIAL","material_state":"INTEGRADO_EN_GRAFO_ACTIVO","residences":[{"path":"manuscrito/sections/integracion_20260812/matricial_52/01_objeto_y_diccionario.tex","sha256":"95d254edbb6dd7184d3a5db5ad83006aee1adbbdd4ee0a5f631e89d6300648da"}],"result_id":"P09","statement_sha256":"cb5502d8e34b02758b04dda42755bf0a91a83e8ba03dbc45f8e4453942b6ce77"}
% HMT_RESULT_END:P09:residence
% HMT_RESULT_BEGIN:P10:residence
% {"material_control":"PASS_PUENTE_GENEALOGICO_MATRICIAL","material_state":"INTEGRADO_EN_GRAFO_ACTIVO","residences":[{"path":"manuscrito/sections/integracion_20260812/matricial_52/01_objeto_y_diccionario.tex","sha256":"95d254edbb6dd7184d3a5db5ad83006aee1adbbdd4ee0a5f631e89d6300648da"}],"result_id":"P10","statement_sha256":"42f48b2ff4254875c1d787ae3c1d25e5f5e086049b6b5f8e963f462ceab4b4d9"}
% HMT_RESULT_END:P10:residence
% HMT_RESULT_BEGIN:P11:residence
% {"material_control":"PASS_PUENTE_GENEALOGICO_MATRICIAL","material_state":"INTEGRADO_EN_GRAFO_ACTIVO","residences":[{"path":"manuscrito/sections/integracion_20260812/matricial_52/01_objeto_y_diccionario.tex","sha256":"95d254edbb6dd7184d3a5db5ad83006aee1adbbdd4ee0a5f631e89d6300648da"}],"result_id":"P11","statement_sha256":"bcd1e3fbe8b15b11e4f83caac1d5b8879ef6f7e0ff3881fdd1ffba891a6f32ce"}
% HMT_RESULT_END:P11:residence
% HMT_RESULT_BEGIN:P12:residence
% {"material_control":"PASS_PUENTE_GENEALOGICO_MATRICIAL","material_state":"INTEGRADO_EN_GRAFO_ACTIVO","residences":[{"path":"manuscrito/sections/integracion_20260812/matricial_52/04_cuadrados_cuanticos.tex","sha256":"7131342598b20a77a08f6d9d98e12d41ec47e331372ea6c19754df21a58ab585"}],"result_id":"P12","statement_sha256":"3eca57321c7a0f5f618f0ffad0376d39eb8c8ab5968847fb959fdbb09fc734ab"}
% HMT_RESULT_END:P12:residence
% HMT_RESULT_BEGIN:P13:residence
% {"material_control":"PASS_PUENTE_GENEALOGICO_MATRICIAL","material_state":"INTEGRADO_EN_GRAFO_ACTIVO","residences":[{"path":"manuscrito/sections/integracion_20260812/matricial_52/04_cuadrados_cuanticos.tex","sha256":"7131342598b20a77a08f6d9d98e12d41ec47e331372ea6c19754df21a58ab585"}],"result_id":"P13","statement_sha256":"55db6d54634c091c6187c36369faa76bdf9f508b14a4a3b18f5a348d2c9cf011"}
% HMT_RESULT_END:P13:residence
% HMT_RESULT_BEGIN:P14:residence
% {"material_control":"PASS_PUENTE_GENEALOGICO_MATRICIAL","material_state":"INTEGRADO_EN_GRAFO_ACTIVO","residences":[{"path":"manuscrito/sections/integracion_20260812/matricial_52/03_bigraduacion_ucp.tex","sha256":"c054ec7cb105bc2fda9c3228bfaee808f01db65e38e126a8865a3420ed37864d"}],"result_id":"P14","statement_sha256":"57d9a1d7ea29cb176674b4f73e0033f34f82419195809fa7e8ff1f235d95e410"}
% HMT_RESULT_END:P14:residence
% HMT_RESULT_BEGIN:P15:residence
% {"material_control":"PASS_PUENTE_GENEALOGICO_MATRICIAL","material_state":"INTEGRADO_EN_GRAFO_ACTIVO","residences":[{"path":"manuscrito/sections/integracion_20260812/matricial_52/03_bigraduacion_ucp.tex","sha256":"c054ec7cb105bc2fda9c3228bfaee808f01db65e38e126a8865a3420ed37864d"}],"result_id":"P15","statement_sha256":"83ae0271d5db4482e0f2f1c4ec1fc0a1b3e5bbde57a18fc5a75a65cf03960e6f"}
% HMT_RESULT_END:P15:residence
% HMT_RESULT_BEGIN:P16:residence
% {"material_control":"PASS_PUENTE_GENEALOGICO_MATRICIAL","material_state":"INTEGRADO_EN_GRAFO_ACTIVO","residences":[{"path":"manuscrito/sections/integracion_20260812/matricial_52/03_bigraduacion_ucp.tex","sha256":"c054ec7cb105bc2fda9c3228bfaee808f01db65e38e126a8865a3420ed37864d"}],"result_id":"P16","statement_sha256":"bdf8eff661630a7ba8edf744bf36bded26cbd2417e7ef7e7259b9c6c7505a21e"}
% HMT_RESULT_END:P16:residence
% HMT_RESULT_BEGIN:P17:residence
% {"material_control":"PASS_PUENTE_GENEALOGICO_MATRICIAL","material_state":"INTEGRADO_EN_GRAFO_ACTIVO","residences":[{"path":"manuscrito/sections/integracion_20260812/matricial_52/03_bigraduacion_ucp.tex","sha256":"c054ec7cb105bc2fda9c3228bfaee808f01db65e38e126a8865a3420ed37864d"}],"result_id":"P17","statement_sha256":"b0cce98edf7fa782d1064e6b7ac175d850f4ad264c1b6d8e62288f2568ad45f5"}
% HMT_RESULT_END:P17:residence
% HMT_RESULT_BEGIN:P18:residence
% {"material_control":"PASS_PUENTE_GENEALOGICO_MATRICIAL","material_state":"INTEGRADO_EN_GRAFO_ACTIVO","residences":[{"path":"manuscrito/sections/integracion_20260812/matricial_52/04_cuadrados_cuanticos.tex","sha256":"7131342598b20a77a08f6d9d98e12d41ec47e331372ea6c19754df21a58ab585"}],"result_id":"P18","statement_sha256":"1473f4abe925a8e28298eb8bf1dfd39fb9af561154966420712347199f61b2ce"}
% HMT_RESULT_END:P18:residence
% HMT_RESULT_BEGIN:P19:residence
% {"material_control":"PASS_PUENTE_GENEALOGICO_MATRICIAL","material_state":"INTEGRADO_EN_GRAFO_ACTIVO","residences":[{"path":"manuscrito/sections/integracion_20260812/matricial_52/04_cuadrados_cuanticos.tex","sha256":"7131342598b20a77a08f6d9d98e12d41ec47e331372ea6c19754df21a58ab585"}],"result_id":"P19","statement_sha256":"fd3e9565106ad07b0ffeff64653cdd8f4e25c01bd8125c56b954e4ad58015bce"}
% HMT_RESULT_END:P19:residence
% HMT_RESULT_BEGIN:P20:residence
% {"material_control":"PASS_PUENTE_GENEALOGICO_MATRICIAL","material_state":"INTEGRADO_EN_GRAFO_ACTIVO","residences":[{"path":"manuscrito/sections/integracion_20260812/matricial_52/04_cuadrados_cuanticos.tex","sha256":"7131342598b20a77a08f6d9d98e12d41ec47e331372ea6c19754df21a58ab585"}],"result_id":"P20","statement_sha256":"f0301b81ceca519949c82bef3e2c9c9d0415d1943aadd55e64bafd83dc369e25"}
% HMT_RESULT_END:P20:residence
% HMT_RESULT_BEGIN:P21:residence
% {"material_control":"PASS_PUENTE_GENEALOGICO_MATRICIAL","material_state":"INTEGRADO_EN_GRAFO_ACTIVO","residences":[{"path":"manuscrito/sections/integracion_20260812/matricial_52/08_sistemas_conicos_y_universalidad.tex","sha256":"2b94acee5d5df46c4c6e2c4b9048962de8961703fc585df45d08a42f9f4ae0d2"}],"result_id":"P21","statement_sha256":"9677a2892d35d3020f115f04121dc2d8f34f6bcda52cdb28ac4ad708ca6fcfe2"}
% HMT_RESULT_END:P21:residence
% HMT_RESULT_BEGIN:P22:residence
% {"material_control":"PASS_PUENTE_GENEALOGICO_MATRICIAL","material_state":"INTEGRADO_EN_GRAFO_ACTIVO","residences":[{"path":"manuscrito/sections/integracion_20260812/matricial_52/05_geometria_simplectica_y_lectores.tex","sha256":"087482ebc4a880dba1096a9364dab1a05b0f0bddd3495b6f5b967fb9ccb79ccd"}],"result_id":"P22","statement_sha256":"f79d8ca5659786ae2dfdcb69f4a77431e62f85c8e159a0cb789c16dd900f0e08"}
% HMT_RESULT_END:P22:residence
% HMT_RESULT_BEGIN:P23:residence
% {"material_control":"PASS_PUENTE_GENEALOGICO_MATRICIAL","material_state":"INTEGRADO_EN_GRAFO_ACTIVO","residences":[{"path":"manuscrito/sections/integracion_20260812/matricial_52/05_geometria_simplectica_y_lectores.tex","sha256":"087482ebc4a880dba1096a9364dab1a05b0f0bddd3495b6f5b967fb9ccb79ccd"}],"result_id":"P23","statement_sha256":"4ebed4b0f0e733875915bad691b12c053780dfaa50e1f0f4898837314d3edf9a"}
% HMT_RESULT_END:P23:residence
% HMT_RESULT_BEGIN:P24:residence
% {"material_control":"PASS_PUENTE_GENEALOGICO_MATRICIAL","material_state":"INTEGRADO_EN_GRAFO_ACTIVO","residences":[{"path":"manuscrito/sections/integracion_20260812/matricial_52/06_operatorios_doble_proyeccion.tex","sha256":"0f183ed734248c42a405ba89f43a6ac6085d705a188f911c44ed76e8eb4bf745"}],"result_id":"P24","statement_sha256":"6196ff6b2fa346db4f4fd64dc27219eba3734ae04559393e65a5ad14e3c297e0"}
% HMT_RESULT_END:P24:residence
% HMT_RESULT_BEGIN:P25:residence
% {"material_control":"PASS_PUENTE_GENEALOGICO_MATRICIAL","material_state":"INTEGRADO_EN_GRAFO_ACTIVO","residences":[{"path":"manuscrito/sections/integracion_20260812/matricial_52/06_operatorios_doble_proyeccion.tex","sha256":"0f183ed734248c42a405ba89f43a6ac6085d705a188f911c44ed76e8eb4bf745"}],"result_id":"P25","statement_sha256":"339d4b25f9c31d5aa693c11a2377a405ec7ba754affe3143c351475a2d50885d"}
% HMT_RESULT_END:P25:residence
% HMT_RESULT_BEGIN:P26:residence
% {"material_control":"PASS_PUENTE_GENEALOGICO_MATRICIAL","material_state":"INTEGRADO_EN_GRAFO_ACTIVO","residences":[{"path":"manuscrito/sections/integracion_20260812/matricial_52/06_operatorios_doble_proyeccion.tex","sha256":"0f183ed734248c42a405ba89f43a6ac6085d705a188f911c44ed76e8eb4bf745"}],"result_id":"P26","statement_sha256":"1f0955a242b45acf69644be639ccca7ceac39f2d83009c4aafcdb1172802a601"}
% HMT_RESULT_END:P26:residence
% HMT_RESULT_BEGIN:P27:residence
% {"material_control":"PASS_PUENTE_GENEALOGICO_MATRICIAL","material_state":"INTEGRADO_EN_GRAFO_ACTIVO","residences":[{"path":"manuscrito/sections/integracion_20260812/matricial_52/06_operatorios_doble_proyeccion.tex","sha256":"0f183ed734248c42a405ba89f43a6ac6085d705a188f911c44ed76e8eb4bf745"}],"result_id":"P27","statement_sha256":"60ff5f6df5943b80338a53b54928755aa4cae85b90764d63a153d6ad50dfed24"}
% HMT_RESULT_END:P27:residence
% HMT_RESULT_BEGIN:P28:residence
% {"material_control":"PASS_PUENTE_GENEALOGICO_MATRICIAL","material_state":"INTEGRADO_EN_GRAFO_ACTIVO","residences":[{"path":"manuscrito/sections/integracion_20260812/matricial_52/07_continuo_tensores_y_memoria.tex","sha256":"3df0b8a366a6458f7acf2199859a25f5b40379b3f9039ccaa0b959e2f25e36a5"}],"result_id":"P28","statement_sha256":"de45d3b80b483962a79b544ce428c4298718a36362562637d064386d09251da4"}
% HMT_RESULT_END:P28:residence
% HMT_RESULT_BEGIN:P29:residence
% {"material_control":"PASS_PUENTE_GENEALOGICO_MATRICIAL","material_state":"INTEGRADO_EN_GRAFO_ACTIVO","residences":[{"path":"manuscrito/sections/integracion_20260812/matricial_52/07_continuo_tensores_y_memoria.tex","sha256":"3df0b8a366a6458f7acf2199859a25f5b40379b3f9039ccaa0b959e2f25e36a5"}],"result_id":"P29","statement_sha256":"eed86be2eeabf3a9de1c3150b3fb89bdee7fc9a0e51ea1e9052cfab634322586"}
% HMT_RESULT_END:P29:residence
% HMT_RESULT_BEGIN:P30:residence
% {"material_control":"PASS_PUENTE_GENEALOGICO_MATRICIAL","material_state":"INTEGRADO_EN_GRAFO_ACTIVO","residences":[{"path":"manuscrito/sections/integracion_20260812/matricial_52/09_metrologia_y_programa.tex","sha256":"d5cf28b96e5cde3577601d038089754378d0eb0e6118b134ed2b27e98004a1e7"}],"result_id":"P30","statement_sha256":"16dae690fc342fd1d4bbb3ee5c7067af92476cfdcc61842908df31b55a844797"}
% HMT_RESULT_END:P30:residence
% HMT_RESULT_BEGIN:P31:residence
% {"material_control":"PASS_PUENTE_GENEALOGICO_MATRICIAL","material_state":"INTEGRADO_EN_GRAFO_ACTIVO","residences":[{"path":"manuscrito/sections/integracion_20260812/matricial_52/08_sistemas_conicos_y_universalidad.tex","sha256":"2b94acee5d5df46c4c6e2c4b9048962de8961703fc585df45d08a42f9f4ae0d2"}],"result_id":"P31","statement_sha256":"cabf80d0abd369646b4013a8ac2aad89b7f4c6cd7e0b0092d5f4e07d76e74471"}
% HMT_RESULT_END:P31:residence
% HMT_RESULT_BEGIN:P32:residence
% {"material_control":"PASS_PUENTE_GENEALOGICO_MATRICIAL","material_state":"INTEGRADO_EN_GRAFO_ACTIVO","residences":[{"path":"manuscrito/sections/integracion_20260812/matricial_52/08_sistemas_conicos_y_universalidad.tex","sha256":"2b94acee5d5df46c4c6e2c4b9048962de8961703fc585df45d08a42f9f4ae0d2"}],"result_id":"P32","statement_sha256":"d808a7bd125136593c39a42df20427716bf90ad4b1ac4a033227a3c4e370c64a"}
% HMT_RESULT_END:P32:residence
% HMT_RESULT_BEGIN:QPM-001:residence
% {"material_control":"PASS_PUENTE_GENEALOGICO_MATRICIAL","material_state":"INTEGRADO_EN_GRAFO_ACTIVO","residences":[{"path":"manuscrito/sections/integracion_20260812/matricial_52/01_objeto_y_diccionario.tex","sha256":"95d254edbb6dd7184d3a5db5ad83006aee1adbbdd4ee0a5f631e89d6300648da"}],"result_id":"QPM-001","statement_sha256":"876638c74d1df52e66cf8b186c110f2cbb879465ca4aed51526e0ca3b191405f"}
% HMT_RESULT_END:QPM-001:residence
% HMT_RESULT_BEGIN:QLS-001:residence
% {"material_control":"PASS_PUENTE_GENEALOGICO_MATRICIAL","material_state":"INTEGRADO_EN_GRAFO_ACTIVO","residences":[{"path":"manuscrito/sections/integracion_20260812/matricial_52/01_objeto_y_diccionario.tex","sha256":"95d254edbb6dd7184d3a5db5ad83006aee1adbbdd4ee0a5f631e89d6300648da"}],"result_id":"QLS-001","statement_sha256":"60abad49376bb9a2df1df2851bfc068d7fd0e59b90f09c5812f8ac93138d5308"}
% HMT_RESULT_END:QLS-001:residence
% HMT_RESULT_BEGIN:QLSI-001:residence
% {"material_control":"PASS_PUENTE_GENEALOGICO_MATRICIAL","material_state":"INTEGRADO_EN_GRAFO_ACTIVO","residences":[{"path":"manuscrito/sections/integracion_20260812/matricial_52/01_objeto_y_diccionario.tex","sha256":"95d254edbb6dd7184d3a5db5ad83006aee1adbbdd4ee0a5f631e89d6300648da"}],"result_id":"QLSI-001","statement_sha256":"3cc12a3dcfe1865a34b892c489519c5c02e989e8075c94e9782f1e01e67b2c1d"}
% HMT_RESULT_END:QLSI-001:residence
% HMT_RESULT_BEGIN:QMSI-001:residence
% {"material_control":"PASS_PUENTE_GENEALOGICO_MATRICIAL","material_state":"INTEGRADO_EN_GRAFO_ACTIVO","residences":[{"path":"manuscrito/sections/integracion_20260812/matricial_52/01_objeto_y_diccionario.tex","sha256":"95d254edbb6dd7184d3a5db5ad83006aee1adbbdd4ee0a5f631e89d6300648da"}],"result_id":"QMSI-001","statement_sha256":"641653814fe7d90a81405d016aa63d2840bdf02215f16c828afe4ece7989c37b"}
% HMT_RESULT_END:QMSI-001:residence
% HMT_RESULT_BEGIN:QMSI-002:residence
% {"material_control":"PASS_PUENTE_GENEALOGICO_MATRICIAL","material_state":"INTEGRADO_EN_GRAFO_ACTIVO","residences":[{"path":"manuscrito/sections/integracion_20260812/matricial_52/01_objeto_y_diccionario.tex","sha256":"95d254edbb6dd7184d3a5db5ad83006aee1adbbdd4ee0a5f631e89d6300648da"}],"result_id":"QMSI-002","statement_sha256":"79afee550d7e55837c8ed6a9067de90d22cfee54bbff7c98636ded289843016f"}
% HMT_RESULT_END:QMSI-002:residence
% HMT_RESULT_BEGIN:QMSI-003:residence
% {"material_control":"PASS_PUENTE_GENEALOGICO_MATRICIAL","material_state":"INTEGRADO_EN_GRAFO_ACTIVO","residences":[{"path":"manuscrito/sections/integracion_20260812/matricial_52/01_objeto_y_diccionario.tex","sha256":"95d254edbb6dd7184d3a5db5ad83006aee1adbbdd4ee0a5f631e89d6300648da"}],"result_id":"QMSI-003","statement_sha256":"1bbe978d6a9ba0ec29aa916bb4b2a72df5bf3936b89bf14a66b2916e1d4fd090"}
% HMT_RESULT_END:QMSI-003:residence
% HMT_RESULT_BEGIN:QBR-001:residence
% {"material_control":"PASS_PUENTE_GENEALOGICO_MATRICIAL","material_state":"INTEGRADO_EN_GRAFO_ACTIVO","residences":[{"path":"manuscrito/sections/integracion_20260812/matricial_52/02_hmt_en_profundidad.tex","sha256":"fe6c23eac09e147199090a091fe2564e5587b7b1155ea85518d2bf47aade4cfc"},{"path":"manuscrito/sections/md/04ab_postulados_cuanticos_genealogicos_rev2.tex","sha256":"1895f84bd75b104e873a3d73c5b0e8ae1b5df606bd6feefeb4f9666a92966804"}],"result_id":"QBR-001","statement_sha256":"38ad3305cd9c643e9067d133a541df231165278ab3b01e906f494ec45be28104"}
% HMT_RESULT_END:QBR-001:residence
% HMT_RESULT_BEGIN:QDY-001:residence
% {"material_control":"PASS_PUENTE_GENEALOGICO_MATRICIAL","material_state":"INTEGRADO_EN_GRAFO_ACTIVO","residences":[{"path":"manuscrito/sections/integracion_20260812/matricial_52/02_hmt_en_profundidad.tex","sha256":"fe6c23eac09e147199090a091fe2564e5587b7b1155ea85518d2bf47aade4cfc"},{"path":"manuscrito/sections/md/04ab_postulados_cuanticos_genealogicos_rev2.tex","sha256":"1895f84bd75b104e873a3d73c5b0e8ae1b5df606bd6feefeb4f9666a92966804"}],"result_id":"QDY-001","statement_sha256":"bffa8460b3d4625a598b73816f5f9e8a9218458a12b69fd3cc36c8ba18d69a89"}
% HMT_RESULT_END:QDY-001:residence
% HMT_RESULT_BEGIN:QKR-001:residence
% {"material_control":"PASS_PUENTE_GENEALOGICO_MATRICIAL","material_state":"INTEGRADO_EN_GRAFO_ACTIVO","residences":[{"path":"manuscrito/sections/integracion_20260812/matricial_52/02_hmt_en_profundidad.tex","sha256":"fe6c23eac09e147199090a091fe2564e5587b7b1155ea85518d2bf47aade4cfc"},{"path":"manuscrito/sections/md/01a_medicion_cuantica_genealogica_rev2.tex","sha256":"725bbfb625e4bd74c3f5845dc2d3ad117b173fac63f11d300f7f6b8344fb505a"}],"result_id":"QKR-001","statement_sha256":"e458e06f5341341d7cbbeee4dbcc1fa333c7c1d293266aec038128fe9c334858"}
% HMT_RESULT_END:QKR-001:residence
% HMT_RESULT_BEGIN:QMO-001:residence
% {"material_control":"PASS_PUENTE_GENEALOGICO_MATRICIAL","material_state":"INTEGRADO_EN_GRAFO_ACTIVO","residences":[{"path":"manuscrito/sections/integracion_20260812/matricial_52/02_hmt_en_profundidad.tex","sha256":"fe6c23eac09e147199090a091fe2564e5587b7b1155ea85518d2bf47aade4cfc"},{"path":"manuscrito/sections/md/04ab_postulados_cuanticos_genealogicos_rev2.tex","sha256":"1895f84bd75b104e873a3d73c5b0e8ae1b5df606bd6feefeb4f9666a92966804"}],"result_id":"QMO-001","statement_sha256":"14d863637e09c435b127f3b5fc5dccc9f9d873d837809884a8efc4b2da2953be"}
% HMT_RESULT_END:QMO-001:residence
% HMT_RESULT_BEGIN:DIC-001:residence
% {"material_control":"PASS_PUENTE_GENEALOGICO_MATRICIAL","material_state":"INTEGRADO_EN_GRAFO_ACTIVO","residences":[{"path":"manuscrito/sections/integracion_20260812/matricial_52/01_objeto_y_diccionario.tex","sha256":"95d254edbb6dd7184d3a5db5ad83006aee1adbbdd4ee0a5f631e89d6300648da"}],"result_id":"DIC-001","statement_sha256":"51d602fff6bbd77d754ba0ebe0732f3998ac22c81a82514440408ff0f6545d30"}
% HMT_RESULT_END:DIC-001:residence
% HMT_RESULT_BEGIN:QMC-001:residence
% {"material_control":"PASS_PUENTE_GENEALOGICO_MATRICIAL","material_state":"INTEGRADO_EN_GRAFO_ACTIVO","residences":[{"path":"manuscrito/sections/integracion_20260812/matricial_52/04_cuadrados_cuanticos.tex","sha256":"7131342598b20a77a08f6d9d98e12d41ec47e331372ea6c19754df21a58ab585"}],"result_id":"QMC-001","statement_sha256":"900c10fc899cb62a7793541419f3c6ce103296ab9690d816954062e980960504"}
% HMT_RESULT_END:QMC-001:residence
% HMT_RESULT_BEGIN:QDT-001:residence
% {"material_control":"PASS_PUENTE_GENEALOGICO_MATRICIAL","material_state":"INTEGRADO_EN_GRAFO_ACTIVO","residences":[{"path":"manuscrito/sections/integracion_20260812/matricial_52/05_geometria_simplectica_y_lectores.tex","sha256":"087482ebc4a880dba1096a9364dab1a05b0f0bddd3495b6f5b967fb9ccb79ccd"}],"result_id":"QDT-001","statement_sha256":"a42d2bae56ce5410c45794269f3181c68a38ea3fb2d5861cebb013294580287d"}
% HMT_RESULT_END:QDT-001:residence
% HMT_RESULT_BEGIN:ENV-001:residence
% {"material_control":"PASS_PUENTE_GENEALOGICO_MATRICIAL","material_state":"INTEGRADO_EN_GRAFO_ACTIVO","residences":[{"path":"manuscrito/sections/integracion_20260812/matricial_52/03_bigraduacion_ucp.tex","sha256":"c054ec7cb105bc2fda9c3228bfaee808f01db65e38e126a8865a3420ed37864d"}],"result_id":"ENV-001","statement_sha256":"a2c7c6251ce55a5d5d21e2825a8eecad64f3da571d4130d8119999126371726c"}
% HMT_RESULT_END:ENV-001:residence
% HMT_RESULT_BEGIN:RNG-001:residence
% {"material_control":"PASS_PUENTE_GENEALOGICO_MATRICIAL","material_state":"INTEGRADO_EN_GRAFO_ACTIVO","residences":[{"path":"manuscrito/sections/integracion_20260812/matricial_52/06_operatorios_doble_proyeccion.tex","sha256":"0f183ed734248c42a405ba89f43a6ac6085d705a188f911c44ed76e8eb4bf745"}],"result_id":"RNG-001","statement_sha256":"6f564d18bb3910cde9dd86c3bcc20cf5be7d388abb45c8acebdc31553910900c"}
% HMT_RESULT_END:RNG-001:residence
% HMT_RESULT_BEGIN:DYN-001:residence
% {"material_control":"PASS_PUENTE_GENEALOGICO_MATRICIAL","material_state":"INTEGRADO_EN_GRAFO_ACTIVO","residences":[{"path":"manuscrito/sections/integracion_20260812/matricial_52/07_continuo_tensores_y_memoria.tex","sha256":"3df0b8a366a6458f7acf2199859a25f5b40379b3f9039ccaa0b959e2f25e36a5"}],"result_id":"DYN-001","statement_sha256":"931716fef6fe78eab6afb720f6b8621e7cf08908dfd2b425ac48605bb361219a"}
% HMT_RESULT_END:DYN-001:residence
% HMT_RESULT_BEGIN:TEN-001:residence
% {"material_control":"PASS_PUENTE_GENEALOGICO_MATRICIAL","material_state":"INTEGRADO_EN_GRAFO_ACTIVO","residences":[{"path":"manuscrito/sections/integracion_20260812/matricial_52/07_continuo_tensores_y_memoria.tex","sha256":"3df0b8a366a6458f7acf2199859a25f5b40379b3f9039ccaa0b959e2f25e36a5"}],"result_id":"TEN-001","statement_sha256":"458096f924295967751d528c1f5b73f3ae28b5e1cec81ab6ac3036eda48c54fd"}
% HMT_RESULT_END:TEN-001:residence
% HMT_RESULT_BEGIN:CHP-001:residence
% {"material_control":"PASS_PUENTE_GENEALOGICO_MATRICIAL","material_state":"INTEGRADO_EN_GRAFO_ACTIVO","residences":[{"path":"manuscrito/sections/integracion_20260812/matricial_52/08_sistemas_conicos_y_universalidad.tex","sha256":"2b94acee5d5df46c4c6e2c4b9048962de8961703fc585df45d08a42f9f4ae0d2"}],"result_id":"CHP-001","statement_sha256":"8fce51f0c85e0120694a493a3cd73f2edb220a7f02de5bc0d820e4697456ff1f"}
% HMT_RESULT_END:CHP-001:residence
% HMT_RESULT_BEGIN:SDP-001:residence
% {"material_control":"PASS_PUENTE_GENEALOGICO_MATRICIAL","material_state":"INTEGRADO_EN_GRAFO_ACTIVO","residences":[{"path":"manuscrito/sections/integracion_20260812/matricial_52/08_sistemas_conicos_y_universalidad.tex","sha256":"2b94acee5d5df46c4c6e2c4b9048962de8961703fc585df45d08a42f9f4ae0d2"}],"result_id":"SDP-001","statement_sha256":"0fa05be883df7a0e3a59ceb8d9386d4775fc4456de540e0b4fa29e7d7716aa46"}
% HMT_RESULT_END:SDP-001:residence
% HMT_RESULT_BEGIN:EQV-001:residence
% {"material_control":"PASS_PUENTE_GENEALOGICO_MATRICIAL","material_state":"INTEGRADO_EN_GRAFO_ACTIVO","residences":[{"path":"manuscrito/sections/integracion_20260812/matricial_52/05_geometria_simplectica_y_lectores.tex","sha256":"087482ebc4a880dba1096a9364dab1a05b0f0bddd3495b6f5b967fb9ccb79ccd"}],"result_id":"EQV-001","statement_sha256":"0ef0dde22fd08d35555aa0a88a1eb4fa608e3b298286a47559fb223af824f79f"}
% HMT_RESULT_END:EQV-001:residence
% HMT_RESULT_BEGIN:QMT-001:residence
% {"material_control":"PASS_PUENTE_GENEALOGICO_MATRICIAL","material_state":"INTEGRADO_EN_GRAFO_ACTIVO","residences":[{"path":"manuscrito/sections/integracion_20260812/matricial_52/04_cuadrados_cuanticos.tex","sha256":"7131342598b20a77a08f6d9d98e12d41ec47e331372ea6c19754df21a58ab585"}],"result_id":"QMT-001","statement_sha256":"e40aff576b435e19d95904fbf6c160fedf5d8f5bcbe4634eab9591b8c6960630"}
% HMT_RESULT_END:QMT-001:residence
% HMT_RESULT_BEGIN:TRN-001:residence
% {"material_control":"PASS_PUENTE_GENEALOGICO_MATRICIAL","material_state":"INTEGRADO_EN_GRAFO_ACTIVO","residences":[{"path":"manuscrito/sections/integracion_20260812/matricial_52/09_metrologia_y_programa.tex","sha256":"d5cf28b96e5cde3577601d038089754378d0eb0e6118b134ed2b27e98004a1e7"}],"result_id":"TRN-001","statement_sha256":"076c0e63737e5420b981d0d4223a441dd2be4af3e62fe56af2608779d7c60bb0"}
% HMT_RESULT_END:TRN-001:residence
% HMT_RESULT_BEGIN:MPSH-001:residence
% {"material_control":"PASS_PUENTE_GENEALOGICO_MATRICIAL","material_state":"INTEGRADO_EN_GRAFO_ACTIVO","residences":[{"path":"manuscrito/sections/integracion_20260812/matricial_52/07_continuo_tensores_y_memoria.tex","sha256":"3df0b8a366a6458f7acf2199859a25f5b40379b3f9039ccaa0b959e2f25e36a5"}],"result_id":"MPSH-001","statement_sha256":"e8ee9870050a98d998dcfa91b8f581a5626e54e41b516743f9666140829b1f9e"}
% HMT_RESULT_END:MPSH-001:residence
% HMT_RESULT_BEGIN:UNIIMP-001:residence
% {"material_control":"PASS_PUENTE_GENEALOGICO_MATRICIAL","material_state":"INTEGRADO_EN_GRAFO_ACTIVO","residences":[{"path":"manuscrito/sections/integracion_20260812/matricial_52/08_sistemas_conicos_y_universalidad.tex","sha256":"2b94acee5d5df46c4c6e2c4b9048962de8961703fc585df45d08a42f9f4ae0d2"}],"result_id":"UNIIMP-001","statement_sha256":"342e5b5e49f7e76a154e4376612d523325fc967212ee3b2f4156d578ed7818be"}
% HMT_RESULT_END:UNIIMP-001:residence
% HMT_RESULT_BEGIN:PRG-001:residence
% {"material_control":"PASS_PUENTE_GENEALOGICO_MATRICIAL","material_state":"INTEGRADO_EN_GRAFO_ACTIVO","residences":[{"path":"manuscrito/sections/integracion_20260812/matricial_52/09_metrologia_y_programa.tex","sha256":"d5cf28b96e5cde3577601d038089754378d0eb0e6118b134ed2b27e98004a1e7"}],"result_id":"PRG-001","statement_sha256":"42475f38a74d5f5c373f586a656ec2ec049a0f9fd5b8ed896dfb8bd11bf3aec5"}
% HMT_RESULT_END:PRG-001:residence
% HMT_RESULT_BEGIN:QCT-001:residence
% {"material_control":"HASH_REGENERACION+NON_RESULTS_6+SCOPE_LIMITS_5","material_state":"CONTROL_FOCAL_ACTIVO","residences":[{"path":"manuscrito/sections/integracion_20260812/matricial_52/A_pruebas_y_certificados.tex","sha256":"ba316b48d94c537b548e76a95774f361aab946f725468bb6acc0700fbe9b0cdb"},{"path":"pruebas/integracion_20260812/verificar_puente_genealogico_matricial.py","sha256":"4661187f03768de0b87fd454620287e9bb407a9402fa6d7e43e6d2969dd2dc69"},{"path":"certificados/integracion_20260812/puente_genealogico_matricial.json","sha256":"fd45cc4e1e1d12225874e214e27fb648e2cea586a5162ba49afdabb981216683"}],"result_id":"QCT-001","statement_sha256":"6a04ea7ae52a0c3f2571d9787272c10ddc18456acf94ddfef7adef2764393d8c"}
% HMT_RESULT_END:QCT-001:residence
% HMT_RESULT_BEGIN:PRV-001:residence
% {"material_control":"HASH_REGENERACION+NON_RESULTS_6+SCOPE_LIMITS_5","material_state":"CONTROL_FOCAL_ACTIVO","residences":[{"path":"manuscrito/sections/integracion_20260812/matricial_52/C_fuentes_y_bibliografia.tex","sha256":"ba54c19f8729fa9e16033a263c3409a80a6fbb6788dfddcb037a657c14f598ff"},{"path":"pruebas/integracion_20260812/verificar_puente_genealogico_matricial.py","sha256":"4661187f03768de0b87fd454620287e9bb407a9402fa6d7e43e6d2969dd2dc69"},{"path":"certificados/integracion_20260812/puente_genealogico_matricial.json","sha256":"fd45cc4e1e1d12225874e214e27fb648e2cea586a5162ba49afdabb981216683"}],"result_id":"PRV-001","statement_sha256":"3898b9fa7cd9ca405abd8eb49767227624e28b2704db43c64ed70e7a0061bf68"}
% HMT_RESULT_END:PRV-001:residence
% HMT_RESULT_BEGIN:MOR-001:residence
% {"material_control":"GATE_MATERIAL_68_22_22+AUDITORIA_ESTATICA","material_state":"INTEGRADO_O_REMITIDO_SIN_PERDIDA","residences":[{"path":"manuscrito/sections/hmt/14f_cierre_genealogico_numero.tex","sha256":"4522230d8297a1420ccb586d0d6c1fefcb08b5e8e2deeda85b11cc73afd8ded5"}],"result_id":"MOR-001","statement_sha256":"81ec51bdba281285df6bfdc1ea9bb82c7f59300a8a07d463fa438037c958c27a"}
% HMT_RESULT_END:MOR-001:residence
% HMT_RESULT_BEGIN:MOR-002:residence
% {"material_control":"GATE_MATERIAL_68_22_22+AUDITORIA_ESTATICA","material_state":"INTEGRADO_O_REMITIDO_SIN_PERDIDA","residences":[{"path":"manuscrito/sections/hmt/03_app_ortograma.tex","sha256":"aaf0fd1dec2c270d544cd19562b4461c57ebb009437e051d973be43395428390"}],"result_id":"MOR-002","statement_sha256":"4d2f31444c015ef52d258636bae153f77c798e12f034cb624fdf8d1721eda040"}
% HMT_RESULT_END:MOR-002:residence
% HMT_RESULT_BEGIN:MOR-003:residence
% {"material_control":"GATE_MATERIAL_68_22_22+AUDITORIA_ESTATICA","material_state":"INTEGRADO_O_REMITIDO_SIN_PERDIDA","residences":[{"path":"manuscrito/sections/hmt/02_trit_geometrias.tex","sha256":"aefb5c1ec6732534c810d0021834eff92a662f59072b97c35c175a6ce7952fa6"}],"result_id":"MOR-003","statement_sha256":"56c13700ad07f59cb91c4278006ecbd0864a2c64b2bc08c14219693a1d5c1961"}
% HMT_RESULT_END:MOR-003:residence
% HMT_RESULT_BEGIN:MOR-004:residence
% {"material_control":"GATE_MATERIAL_68_22_22+AUDITORIA_ESTATICA","material_state":"INTEGRADO_O_REMITIDO_SIN_PERDIDA","residences":[{"path":"manuscrito/sections/hmt/03k_genealogia_operaciones_rev2.tex","sha256":"fae7aca94c6fb20754431ed79a5aa2f005a7cf2f5681736fc5826da4c694fb98"}],"result_id":"MOR-004","statement_sha256":"e709d75fe6d60011165caa4daee14d1c7b07aa4b57f2395a2ef27267672e3448"}
% HMT_RESULT_END:MOR-004:residence
% HMT_RESULT_BEGIN:MOR-005:residence
% {"material_control":"GATE_MATERIAL_68_22_22+AUDITORIA_ESTATICA","material_state":"INTEGRADO_O_REMITIDO_SIN_PERDIDA","residences":[{"path":"manuscrito/sections/integracion_20260812/realizacion_68/01_cilindros_kronecker_dirac.tex","sha256":"782450c6fc1e432641df8b3cc605425c604109f1db5830873b0b778c59ae6946"}],"result_id":"MOR-005","statement_sha256":"39a6a2f053bee151e082142609deae0baa2067c95df7e2cc57eabe0e0cdb0b35"}
% HMT_RESULT_END:MOR-005:residence
% HMT_RESULT_BEGIN:MOR-006:residence
% {"material_control":"GATE_MATERIAL_68_22_22+AUDITORIA_ESTATICA","material_state":"INTEGRADO_O_REMITIDO_SIN_PERDIDA","residences":[{"path":"manuscrito/sections/integracion_20260812/realizacion_68/01_cilindros_kronecker_dirac.tex","sha256":"782450c6fc1e432641df8b3cc605425c604109f1db5830873b0b778c59ae6946"}],"result_id":"MOR-006","statement_sha256":"969223c89123c8150a4a16dde30b45f11ccb2c9ba9a773533e000cbb936e1d0c"}
% HMT_RESULT_END:MOR-006:residence
% HMT_RESULT_BEGIN:MOR-007:residence
% {"material_control":"GATE_MATERIAL_68_22_22+AUDITORIA_ESTATICA","material_state":"INTEGRADO_O_REMITIDO_SIN_PERDIDA","residences":[{"path":"manuscrito/sections/integracion_20260812/bloque_64_operatorial.tex","sha256":"46e39285e0cee5f56142e4c31f8aece125754649847a8864878db03c39927c0f"}],"result_id":"MOR-007","statement_sha256":"cd14cec597ee88698599f74945cfe83dee11ed410ee4f5f9c55eea289b84c15d"}
% HMT_RESULT_END:MOR-007:residence
% HMT_RESULT_BEGIN:MOR-008:residence
% {"material_control":"GATE_MATERIAL_68_22_22+AUDITORIA_ESTATICA","material_state":"INTEGRADO_O_REMITIDO_SIN_PERDIDA","residences":[{"path":"manuscrito/sections/md/01a_medicion_cuantica_genealogica_rev2.tex","sha256":"725bbfb625e4bd74c3f5845dc2d3ad117b173fac63f11d300f7f6b8344fb505a"},{"path":"manuscrito/sections/md/04ab_postulados_cuanticos_genealogicos_rev2.tex","sha256":"1895f84bd75b104e873a3d73c5b0e8ae1b5df606bd6feefeb4f9666a92966804"}],"result_id":"MOR-008","statement_sha256":"a220ba7c55fda4a44314abe4540421d4c8d4b4e75f7ffa6f6bc273734217fc40"}
% HMT_RESULT_END:MOR-008:residence
% HMT_RESULT_BEGIN:MOR-009:residence
% {"material_control":"GATE_MATERIAL_68_22_22+AUDITORIA_ESTATICA","material_state":"INTEGRADO_O_REMITIDO_SIN_PERDIDA","residences":[{"path":"manuscrito/sections/md/04ab_postulados_cuanticos_genealogicos_rev2.tex","sha256":"1895f84bd75b104e873a3d73c5b0e8ae1b5df606bd6feefeb4f9666a92966804"},{"path":"manuscrito/sections/md/00b_realizacion_espectral_numero_particula.tex","sha256":"741b8542e90d89c9973f78712ccfb0625689f72a5af27ecedcc81d193f38d8f6"}],"result_id":"MOR-009","statement_sha256":"3b85ec28f733e941f275615780570ab9c2a3d81fc638a62c0e6728f5ae1f5925"}
% HMT_RESULT_END:MOR-009:residence
% HMT_RESULT_BEGIN:MOR-010:residence
% {"material_control":"GATE_MATERIAL_68_22_22+AUDITORIA_ESTATICA","material_state":"INTEGRADO_O_REMITIDO_SIN_PERDIDA","residences":[{"path":"manuscrito/sections/integracion_20260812/realizacion_68/02_realizacion_fisica_discreta_corregida.tex","sha256":"219a1a44dc7a0979b3728cebdc4308038cfbef74342bb0fcba50fe754c30292e"}],"result_id":"MOR-010","statement_sha256":"f2c89228291df78e200ef2085a4f4d478ac1be2e6684f9a715f658123de31342"}
% HMT_RESULT_END:MOR-010:residence
% HMT_RESULT_BEGIN:MOR-011:residence
% {"material_control":"GATE_MATERIAL_68_22_22+AUDITORIA_ESTATICA","material_state":"INTEGRADO_O_REMITIDO_SIN_PERDIDA","residences":[{"path":"manuscrito/sections/integracion_20260812/realizacion_68/05_involucion_constitutiva_bidireccional.tex","sha256":"c5cc64e9d5b8014a44b0c30ef8b2885e64b0379a107a073c6026b15381f76a7c"}],"result_id":"MOR-011","statement_sha256":"1ff2ade4f9f1721656ba26a95ce77ff15b5def7ccc20e9c135257ef0a7b97f0b"}
% HMT_RESULT_END:MOR-011:residence
% HMT_RESULT_BEGIN:MOR-012:residence
% {"material_control":"GATE_MATERIAL_68_22_22+AUDITORIA_ESTATICA","material_state":"INTEGRADO_O_REMITIDO_SIN_PERDIDA","residences":[{"path":"manuscrito/sections/integracion_20260812/realizacion_68/05_involucion_constitutiva_bidireccional.tex","sha256":"c5cc64e9d5b8014a44b0c30ef8b2885e64b0379a107a073c6026b15381f76a7c"}],"result_id":"MOR-012","statement_sha256":"c71dfcd2e0330c46692b8833f1bc5d41d7878bfc9db5a82f4cbfc84d7549ab7c"}
% HMT_RESULT_END:MOR-012:residence
% HMT_RESULT_BEGIN:MOR-013:residence
% {"material_control":"GATE_MATERIAL_68_22_22+AUDITORIA_ESTATICA","material_state":"INTEGRADO_O_REMITIDO_SIN_PERDIDA","residences":[{"path":"manuscrito/sections/md/04c_realizaciones_nula_material.tex","sha256":"d4bc1ef3e11a544359a35157497479312fce619dec624b4e01d33b6df16a84d9"}],"result_id":"MOR-013","statement_sha256":"ffe00c7b9b7d0edf7f35a609adfde95e7d3498e2638a94dfd78b707d810b9890"}
% HMT_RESULT_END:MOR-013:residence
% HMT_RESULT_BEGIN:MOR-014:residence
% {"material_control":"GATE_MATERIAL_68_22_22+AUDITORIA_ESTATICA","material_state":"INTEGRADO_O_REMITIDO_SIN_PERDIDA","residences":[{"path":"manuscrito/sections/md/04c_realizaciones_nula_material.tex","sha256":"d4bc1ef3e11a544359a35157497479312fce619dec624b4e01d33b6df16a84d9"},{"path":"manuscrito/sections/md/06ka_clausura_nulo_material_electron_rev3.tex","sha256":"4e83aee8d1ef22ee3bc6e71b10989c62a7a01c4f41153c30efc054909337c201"}],"result_id":"MOR-014","statement_sha256":"892c56b47c56b0f46d3c0955bfbb8bd5ef3f1c07286a797739e61c565e0dc2a2"}
% HMT_RESULT_END:MOR-014:residence
% HMT_RESULT_BEGIN:MOR-015:residence
% {"material_control":"GATE_MATERIAL_68_22_22+AUDITORIA_ESTATICA","material_state":"INTEGRADO_O_REMITIDO_SIN_PERDIDA","residences":[{"path":"manuscrito/sections/integracion_20260812/realizacion_68/03_hopf_villarceau_memoria_5_2.tex","sha256":"5ff5c6ec5855f74338db3440d9a7ddeaa0d93fec54b1ca5bb5094d6a14bd0fff"}],"result_id":"MOR-015","statement_sha256":"49efabbeafb0ec3ca4300a5684772f9b39fbc9836d768664adc1d443c2be762b"}
% HMT_RESULT_END:MOR-015:residence
% HMT_RESULT_BEGIN:MOR-016:residence
% {"material_control":"GATE_MATERIAL_68_22_22+AUDITORIA_ESTATICA","material_state":"INTEGRADO_O_REMITIDO_SIN_PERDIDA","residences":[{"path":"manuscrito/sections/md/04b_banda_particula_virtual.tex","sha256":"d9b8c3db3bc26d37b5a78251a5fbcf67076eaacf35c194b87bf6552d1edfac08"},{"path":"manuscrito/sections/md/21_persistencia_ruta_mobius_familias_actualizada.tex","sha256":"375c7e897d21777f05f9a506bb2429e5de76351b55f10e15f275913d9c306999"}],"result_id":"MOR-016","statement_sha256":"d549af55e36e03afb5ec1ca58bf41f7c6c0c5bd39e17a4a0320aff3b421fec3b"}
% HMT_RESULT_END:MOR-016:residence
% HMT_RESULT_BEGIN:MOR-017:residence
% {"material_control":"GATE_MATERIAL_68_22_22+AUDITORIA_ESTATICA","material_state":"INTEGRADO_O_REMITIDO_SIN_PERDIDA","residences":[{"path":"manuscrito/sections/integracion_20260812/realizacion_68/04_control_cglmp_exacto_68.tex","sha256":"278ced97b864e5c8a06650fca5a102c1a213716af14d3148e8206da54f91fdd7"}],"result_id":"MOR-017","statement_sha256":"8aceaf43834801dfe28d27e802dea64160a4e53dd788c2579d1ea1828dd1aba1"}
% HMT_RESULT_END:MOR-017:residence
% HMT_RESULT_BEGIN:MOR-018:residence
% {"material_control":"GATE_MATERIAL_68_22_22+AUDITORIA_ESTATICA","material_state":"INTEGRADO_O_REMITIDO_SIN_PERDIDA","residences":[{"path":"manuscrito/sections/md/10_constitutivas_si.tex","sha256":"852439fe31f4004a04c10217e89f3414c049bc8803c5b86e3d5b2097c04da0a1"}],"result_id":"MOR-018","statement_sha256":"1c4bd3e069ce1c87d0a969b52dbb80e067cb34adc73c877aa272e4018e841ed0"}
% HMT_RESULT_END:MOR-018:residence
% HMT_RESULT_BEGIN:MOR-019:residence
% {"material_control":"GATE_MATERIAL_68_22_22+AUDITORIA_ESTATICA","material_state":"INTEGRADO_O_REMITIDO_SIN_PERDIDA","residences":[{"path":"manuscrito/sections/md/03b_escala_accion_determinantal.tex","sha256":"73e012418f7e982e277e3e978a0fb03d76d0aebacc8fe085341c70121b20d832"},{"path":"manuscrito/sections/md/04e_termodinamica_rutas.tex","sha256":"ab1429f9145eee89e5287095b94f327291fa4aa247ae0abcc95e3f5611de87c1"}],"result_id":"MOR-019","statement_sha256":"12f6cc8a7f2131682e716fb00c3f2546dbb1f63249994d5d8c1c14d83510c994"}
% HMT_RESULT_END:MOR-019:residence
% HMT_RESULT_BEGIN:MOR-020:residence
% {"material_control":"GATE_MATERIAL_68_22_22+AUDITORIA_ESTATICA","material_state":"INTEGRADO_O_REMITIDO_SIN_PERDIDA","residences":[{"path":"manuscrito/sections/md/05aa_ley_general_masa_accion_autonoma.tex","sha256":"43dbef95faacb5db35585fa447b16406038f05b5296852f259687292bc0313a2"},{"path":"manuscrito/sections/md/05f_operadores_masa_two_way_rev11.tex","sha256":"e0c576ff54f405c41adea7093ca6b117134e9322749c969f1fc4cca6e8615d35"},{"path":"manuscrito/sections/md/06m_tabla_completa_masas_hmt_sm_revfinal.tex","sha256":"1dd8ca6908a9b94af9d5533d57785a2222d455f88e7ac70aa9acf34667ed35b1"}],"result_id":"MOR-020","statement_sha256":"ef9647c7358991a2b5046f1b253d9497443b3f0c31a95b7f6a9906100d45378e"}
% HMT_RESULT_END:MOR-020:residence
% HMT_RESULT_BEGIN:MOR-021:residence
% {"material_control":"GATE_MATERIAL_68_22_22+AUDITORIA_ESTATICA","material_state":"INTEGRADO_O_REMITIDO_SIN_PERDIDA","residences":[{"path":"manuscrito/sections/md/06k_electron_two_way_rev11.tex","sha256":"e9857301411fd54a3f9d9ad99f56466fc2d7aa72b59d7a4a12f14076fc151e0b"},{"path":"manuscrito/sections/md/06ka_clausura_nulo_material_electron_rev3.tex","sha256":"4e83aee8d1ef22ee3bc6e71b10989c62a7a01c4f41153c30efc054909337c201"}],"result_id":"MOR-021","statement_sha256":"79bbb961eec13af6b1fb10be55c0f3a17f871661f753bc37ab8c5e6ce19c0a6a"}
% HMT_RESULT_END:MOR-021:residence
% HMT_RESULT_BEGIN:MOR-022:residence
% {"material_control":"GATE_MATERIAL_68_22_22+AUDITORIA_ESTATICA","material_state":"INTEGRADO_O_REMITIDO_SIN_PERDIDA","residences":[{"path":"manuscrito/sections/md/03d_identidad_autodual_planck_ct108_rev3.tex","sha256":"522a1ce5c537fb26513227fc0766eadefd1b6fd898382137129677aaa69cfaed"}],"result_id":"MOR-022","statement_sha256":"e741dca3aef92e05581285a7ecd9802c495fc2b37445a4604229e67b619c98da"}
% HMT_RESULT_END:MOR-022:residence
% HMT_RESULT_BEGIN:MOR-023:residence
% {"material_control":"GATE_MATERIAL_68_22_22+AUDITORIA_ESTATICA","material_state":"INTEGRADO_O_REMITIDO_SIN_PERDIDA","residences":[{"path":"manuscrito/sections/integracion_20260812/realizacion_68/02_realizacion_fisica_discreta_corregida.tex","sha256":"219a1a44dc7a0979b3728cebdc4308038cfbef74342bb0fcba50fe754c30292e"}],"result_id":"MOR-023","statement_sha256":"bfaa0b7d864dd992d466327275ed93894d90946d897887db756d6d92268920b5"}
% HMT_RESULT_END:MOR-023:residence
% HMT_RESULT_BEGIN:MOR-024:residence
% {"material_control":"GATE_MATERIAL_68_22_22+AUDITORIA_ESTATICA","material_state":"INTEGRADO_O_REMITIDO_SIN_PERDIDA","residences":[{"path":"manuscrito/sections/md/11c_gravedad_holonomia_rev6.tex","sha256":"87d99cb98c92f5e7a771a9ab750690274e370d49a630e461b9d524e02cf6cf71"},{"path":"manuscrito/sections/md/11n_selector_gravitatorio_rev11.tex","sha256":"14b45e691dfd01916515b4a15b910eba1ff698031f794e68ba13d72da9973c9d"}],"result_id":"MOR-024","statement_sha256":"7736b062aa565d44a6b82fd9985d0d4dd7dd7bf8810791d818212081641834d4"}
% HMT_RESULT_END:MOR-024:residence
% HMT_RESULT_BEGIN:MOR-025:residence
% {"material_control":"GATE_MATERIAL_68_22_22+AUDITORIA_ESTATICA","material_state":"INTEGRADO_O_REMITIDO_SIN_PERDIDA","residences":[{"path":"manuscrito/sections/integracion_20260812/realizacion_68/06_programas_mach_cosmologia_68.tex","sha256":"5cca10348cd22a33e86ac0edaaffbd486e433e7e3d80bdf4c97a9368ab868e89"}],"result_id":"MOR-025","statement_sha256":"560fed3a547431a0a0bb145c63be86a48bd71bbf9d20c9ce42e8660b1aa64996"}
% HMT_RESULT_END:MOR-025:residence
% HMT_RESULT_BEGIN:MOR-026:residence
% {"material_control":"GATE_MATERIAL_68_22_22+AUDITORIA_ESTATICA","material_state":"INTEGRADO_O_REMITIDO_SIN_PERDIDA","residences":[{"path":"manuscrito/sections/hmt/14f_cierre_genealogico_numero.tex","sha256":"4522230d8297a1420ccb586d0d6c1fefcb08b5e8e2deeda85b11cc73afd8ded5"},{"path":"manuscrito/sections/hmt/19b_zeta_vacancias_primos_rev11.tex","sha256":"37b1f8fff4663d03826091250843b7a33c8a83e3fd08659c5fa6613a4c346534"}],"result_id":"MOR-026","statement_sha256":"e27ece362c8e14e5d61ce86444197f785ee77f3c6b2a79b2057a895ebb30db26"}
% HMT_RESULT_END:MOR-026:residence
% HMT_RESULT_BEGIN:MOR-027:residence
% {"material_control":"GATE_MATERIAL_68_22_22+AUDITORIA_ESTATICA","material_state":"INTEGRADO_O_REMITIDO_SIN_PERDIDA","residences":[{"path":"manuscrito/sections/hmt/14d_ontologia_numeros_rev11.tex","sha256":"c96b3c94af4668f15c6a061260d1c1900c52b69b8a4a5d2442268f6ff1d0abcc"},{"path":"manuscrito/sections/hmt/14f_cierre_genealogico_numero.tex","sha256":"4522230d8297a1420ccb586d0d6c1fefcb08b5e8e2deeda85b11cc73afd8ded5"}],"result_id":"MOR-027","statement_sha256":"5b813e0066c4fb64f0e94a22e81f2741e55dc4dec333c9aea2d7d02f31523a5d"}
% HMT_RESULT_END:MOR-027:residence
% HMT_RESULT_BEGIN:MOR-028:residence
% {"material_control":"GATE_MATERIAL_68_22_22+AUDITORIA_ESTATICA","material_state":"INTEGRADO_O_REMITIDO_SIN_PERDIDA","residences":[{"path":"manuscrito/sections/hmt/03f_puente_central_app_emrh_rev2.tex","sha256":"f4088242d198ece37ba686e3314432777582aa93e3b7960e39986a5e9a54358f"}],"result_id":"MOR-028","statement_sha256":"3efd7a28586cf53974d753f4d22e4e7b42a73517f5454c21aefb37550f6b1eae"}
% HMT_RESULT_END:MOR-028:residence
% HMT_RESULT_BEGIN:MOR-029:residence
% {"material_control":"GATE_MATERIAL_68_22_22+AUDITORIA_ESTATICA","material_state":"INTEGRADO_O_REMITIDO_SIN_PERDIDA","residences":[{"path":"manuscrito/sections/hmt/06e_certificado_generacion_monodromica_rev7.tex","sha256":"2a0cf8028e12605a97923c53726fcb62c1798c458af8e9f78f12c1129b0e1b94"}],"result_id":"MOR-029","statement_sha256":"91fbdaa9566d579a9f7f0fc50c56eae1187bb6ce719915dc9f11c1031024480f"}
% HMT_RESULT_END:MOR-029:residence
% HMT_RESULT_BEGIN:MOR-030:residence
% {"material_control":"GATE_MATERIAL_68_22_22+AUDITORIA_ESTATICA","material_state":"INTEGRADO_O_REMITIDO_SIN_PERDIDA","residences":[{"path":"manuscrito/sections/integracion_20260812/realizacion_68/08_interfaces_falsadores_realizacion_68.tex","sha256":"323f27caa4ce2d4e3e96795394eb2843c9549f0c7137bde4a8702bf79561e05e"}],"result_id":"MOR-030","statement_sha256":"ae3db0712fc04e6f7197110ed9ac53af8a4e6d1ce41056c0f1849cbe71519bc2"}
% HMT_RESULT_END:MOR-030:residence
% HMT_RESULT_BEGIN:PLS-001:residence
% {"material_control":"GATE_MATERIAL_68_22_22+AUDITORIA_ESTATICA","material_state":"INTEGRADO_O_REMITIDO_SIN_PERDIDA","residences":[{"path":"manuscrito/sections/md/03d_identidad_autodual_planck_ct108_rev3.tex","sha256":"522a1ce5c537fb26513227fc0766eadefd1b6fd898382137129677aaa69cfaed"}],"result_id":"PLS-001","statement_sha256":"dd9b9e630ef17a042a9d2eb03dfb5bfb4f12afef20d488d3eab6c9d3236bc69a"}
% HMT_RESULT_END:PLS-001:residence
% HMT_RESULT_BEGIN:MAD-001:residence
% {"material_control":"GATE_MATERIAL_68_22_22+AUDITORIA_ESTATICA","material_state":"INTEGRADO_O_REMITIDO_SIN_PERDIDA","residences":[{"path":"manuscrito/sections/md/03d_identidad_autodual_planck_ct108_rev3.tex","sha256":"522a1ce5c537fb26513227fc0766eadefd1b6fd898382137129677aaa69cfaed"}],"result_id":"MAD-001","statement_sha256":"5bce3c37e8568b7428605807933327d2bddbd9668451a4eac13bc22f055af271"}
% HMT_RESULT_END:MAD-001:residence
% HMT_RESULT_BEGIN:AGR-001:residence
% {"material_control":"GATE_MATERIAL_68_22_22+AUDITORIA_ESTATICA","material_state":"INTEGRADO_O_REMITIDO_SIN_PERDIDA","residences":[{"path":"manuscrito/sections/md/03d_identidad_autodual_planck_ct108_rev3.tex","sha256":"522a1ce5c537fb26513227fc0766eadefd1b6fd898382137129677aaa69cfaed"}],"result_id":"AGR-001","statement_sha256":"4fee4acf8c84a7dbbaf710eff510d2cb37e76574c593f4912526198dafaa1747"}
% HMT_RESULT_END:AGR-001:residence
% HMT_RESULT_BEGIN:PLE-001:residence
% {"material_control":"GATE_MATERIAL_68_22_22+AUDITORIA_ESTATICA","material_state":"INTEGRADO_O_REMITIDO_SIN_PERDIDA","residences":[{"path":"manuscrito/sections/md/03d_identidad_autodual_planck_ct108_rev3.tex","sha256":"522a1ce5c537fb26513227fc0766eadefd1b6fd898382137129677aaa69cfaed"},{"path":"manuscrito/sections/md/06ka_clausura_nulo_material_electron_rev3.tex","sha256":"4e83aee8d1ef22ee3bc6e71b10989c62a7a01c4f41153c30efc054909337c201"}],"result_id":"PLE-001","statement_sha256":"354993392866e30255c6503239114f91f02cba03ce48f1bd36427071e5939fc8"}
% HMT_RESULT_END:PLE-001:residence
% HMT_RESULT_BEGIN:ANG-001:residence
% {"material_control":"GATE_MATERIAL_68_22_22+AUDITORIA_ESTATICA","material_state":"INTEGRADO_O_REMITIDO_SIN_PERDIDA","residences":[{"path":"manuscrito/sections/integracion_20260812/realizacion_68/03_hopf_villarceau_memoria_5_2.tex","sha256":"5ff5c6ec5855f74338db3440d9a7ddeaa0d93fec54b1ca5bb5094d6a14bd0fff"}],"result_id":"ANG-001","statement_sha256":"f8f0e5ef3962f4ee1a30115f54f35751e6d96da9a4b54f83017c293e67683683"}
% HMT_RESULT_END:ANG-001:residence
% HMT_RESULT_BEGIN:BDR-001:residence
% {"material_control":"GATE_MATERIAL_68_22_22+AUDITORIA_ESTATICA","material_state":"INTEGRADO_O_REMITIDO_SIN_PERDIDA","residences":[{"path":"manuscrito/sections/integracion_20260812/realizacion_68/05_involucion_constitutiva_bidireccional.tex","sha256":"c5cc64e9d5b8014a44b0c30ef8b2885e64b0379a107a073c6026b15381f76a7c"}],"result_id":"BDR-001","statement_sha256":"6a0929cf24c9401330b5cc9ca4baa6fbfe43f7d9efa70a75883aa7f5168a4682"}
% HMT_RESULT_END:BDR-001:residence
% HMT_RESULT_BEGIN:OCC-001:residence
% {"material_control":"GATE_MATERIAL_68_22_22+AUDITORIA_ESTATICA","material_state":"INTEGRADO_O_REMITIDO_SIN_PERDIDA","residences":[{"path":"manuscrito/sections/md/04e_termodinamica_rutas.tex","sha256":"ab1429f9145eee89e5287095b94f327291fa4aa247ae0abcc95e3f5611de87c1"}],"result_id":"OCC-001","statement_sha256":"ee25ce148dd36c42aeaace51d97662de4acb35aa1ed8f59a27a40f937515c612"}
% HMT_RESULT_END:OCC-001:residence
% HMT_RESULT_BEGIN:BOL-001:residence
% {"material_control":"GATE_MATERIAL_68_22_22+AUDITORIA_ESTATICA","material_state":"INTEGRADO_O_REMITIDO_SIN_PERDIDA","residences":[{"path":"manuscrito/sections/md/04e_termodinamica_rutas.tex","sha256":"ab1429f9145eee89e5287095b94f327291fa4aa247ae0abcc95e3f5611de87c1"}],"result_id":"BOL-001","statement_sha256":"b336e9229e7eae82bc749151dce330d09f3ec616c75b361fac0728e7c87f6c18"}
% HMT_RESULT_END:BOL-001:residence
% HMT_RESULT_BEGIN:COS-001:residence
% {"material_control":"GATE_MATERIAL_68_22_22+AUDITORIA_ESTATICA","material_state":"INTEGRADO_O_REMITIDO_SIN_PERDIDA","residences":[{"path":"manuscrito/sections/integracion_20260812/realizacion_68/06_programas_mach_cosmologia_68.tex","sha256":"5cca10348cd22a33e86ac0edaaffbd486e433e7e3d80bdf4c97a9368ab868e89"}],"result_id":"COS-001","statement_sha256":"bd40c86eaf9fa174e90b2113b1e1001dc7db1a73fb69eb230c228131867c4969"}
% HMT_RESULT_END:COS-001:residence
% HMT_RESULT_BEGIN:GWV-001:residence
% {"material_control":"GATE_MATERIAL_68_22_22+AUDITORIA_ESTATICA","material_state":"INTEGRADO_O_REMITIDO_SIN_PERDIDA","residences":[{"path":"manuscrito/sections/integracion_20260812/realizacion_68/06_programas_mach_cosmologia_68.tex","sha256":"5cca10348cd22a33e86ac0edaaffbd486e433e7e3d80bdf4c97a9368ab868e89"}],"result_id":"GWV-001","statement_sha256":"3cf172b69a92d50f0b499502ca1a16ea845a1c3adb4ce9e5c07cfda16f9c38fa"}
% HMT_RESULT_END:GWV-001:residence
% HMT_RESULT_BEGIN:DRK-001:residence
% {"material_control":"GATE_MATERIAL_68_22_22+AUDITORIA_ESTATICA","material_state":"INTEGRADO_O_REMITIDO_SIN_PERDIDA","residences":[{"path":"manuscrito/sections/integracion_20260812/realizacion_68/06_programas_mach_cosmologia_68.tex","sha256":"5cca10348cd22a33e86ac0edaaffbd486e433e7e3d80bdf4c97a9368ab868e89"}],"result_id":"DRK-001","statement_sha256":"cfbb50fa8ff9e04eb1292280e0f1cac17fb57e837209733564e0026a2be31639"}
% HMT_RESULT_END:DRK-001:residence
% HMT_RESULT_BEGIN:CMK-001:residence
% {"material_control":"GATE_MATERIAL_68_22_22+AUDITORIA_ESTATICA","material_state":"INTEGRADO_O_REMITIDO_SIN_PERDIDA","residences":[{"path":"manuscrito/sections/integracion_20260812/realizacion_68/07_cmk_coestructura_dinamica.tex","sha256":"66c38ef3af1739a961696c1d7c23c719df6530be5afced7b53e2c3de4b79d1fd"}],"result_id":"CMK-001","statement_sha256":"d549072e45a1534a427abeaf657bff1597774b49ed6b4c93ac49b3a4a6c836a6"}
% HMT_RESULT_END:CMK-001:residence
% HMT_RESULT_BEGIN:MRC-001:residence
% {"material_control":"GATE_MATERIAL_68_22_22+AUDITORIA_ESTATICA","material_state":"INTEGRADO_O_REMITIDO_SIN_PERDIDA","residences":[{"path":"manuscrito/sections/integracion_20260812/realizacion_68/09_certificado_y_disposicion_68.tex","sha256":"87dc8390235dd0f1c606ca180c5b56639642fa4818d7400f4ce633c8df913d0c"}],"result_id":"MRC-001","statement_sha256":"80a7a36147cfc1163735bc538f698ff947ea95661565393b3489da4e24a85dd7"}
% HMT_RESULT_END:MRC-001:residence
% HMT_RESULT_BEGIN:PRG68-001:residence
% {"material_control":"GATE_MATERIAL_68_22_22+AUDITORIA_ESTATICA","material_state":"INTEGRADO_O_REMITIDO_SIN_PERDIDA","residences":[{"path":"manuscrito/sections/integracion_20260812/realizacion_68/08_interfaces_falsadores_realizacion_68.tex","sha256":"323f27caa4ce2d4e3e96795394eb2843c9549f0c7137bde4a8702bf79561e05e"}],"result_id":"PRG68-001","statement_sha256":"cf2a64c0deb4ed89400145c71e75b98276573e3fd427b0a93e0ecf04d4114f26"}
% HMT_RESULT_END:PRG68-001:residence
% HMT_RESULT_BEGIN:GTR-001:residence
% {"material_control":"GATE_MATERIAL_68_22_22+AUDITORIA_ESTATICA","material_state":"INTEGRADO_O_REMITIDO_SIN_PERDIDA","residences":[{"path":"manuscrito/sections/md/11c_gravedad_holonomia_rev6.tex","sha256":"87d99cb98c92f5e7a771a9ab750690274e370d49a630e461b9d524e02cf6cf71"},{"path":"manuscrito/sections/md/11n_selector_gravitatorio_rev11.tex","sha256":"14b45e691dfd01916515b4a15b910eba1ff698031f794e68ba13d72da9973c9d"}],"result_id":"GTR-001","statement_sha256":"ddfba4182cba6421426c961c29f8e8fc1053105316f55429bf5db6218eae8a25"}
% HMT_RESULT_END:GTR-001:residence
% HMT_RESULT_BEGIN:AMB-001:residence
% {"material_control":"PASS_AMBIVALENCE_AREAL_VOLUMETRIC+GATE_MATERIAL_V5_219","material_state":"CONSERVADO_Y_ATOMIZADO_EN_DELTA_SUCESOR","residences":[{"path":"manuscrito/sections/hmt/05b_extensiones_dinamica_cociente.tex","sha256":"2c562fa974ec9637a3fe09b0594b57cf0652f1f6813e6eefcae2a14929248882"}],"result_id":"AMB-001","statement_sha256":"a33b07c534544852bd3e4e64ece4fe0c7c8af500d5eb8b8addbb97736170835e"}
% HMT_RESULT_END:AMB-001:residence
 % Sidecar operativo sucesor: conserva los 219 resultados atómicos de REV.4 y
% añade los 29 resultados CHI, separado del estado material testigo.
% Sidecar atómico de la edición integral reforzada: 261 IDs estables.
% PREPARED: no es recibo y no autoriza por sí solo la compilación.
% HMT_RESULT_BEGIN:R3-001:residence
% HMT_ATOMIC_RESULT {"material_state":"CONSERVADO_POR_BASE_GRANULAR_REV3","proof_strength":"EXACTO_INTERNO","provenance":"RESULTADO_RECUPERADO","reinforcement_state":"DEPLOYED_WITHOUT_NEW_RESULT_ID","residence":"portada, resumen, introducción, síntesis","result_id":"R3-001","statement":"Título y subtítulo como tesis: cierre holográfico, estructura discreta del continuo y origen de constantes","status":"PRESENT_IN_SUCCESSOR_SOURCE"}
% HMT_RESULT_END:R3-001:residence
% HMT_RESULT_BEGIN:R3-002:residence
% HMT_ATOMIC_RESULT {"material_state":"CONSERVADO_POR_BASE_GRANULAR_REV3","proof_strength":"EXACTO_INTERNO","provenance":"RESULTADO_RECUPERADO","reinforcement_state":"DEPLOYED_WITHOUT_NEW_RESULT_ID","residence":"`hmt/00_construccion_app_gnomon_rev11`, `03_app_ortograma`, `03f`","result_id":"R3-002","statement":"APP completo: suma, producto, residuo, cociente, gnomón, fibras, acarreo y bloque central","status":"PRESENT_IN_SUCCESSOR_SOURCE"}
% HMT_RESULT_END:R3-002:residence
% HMT_RESULT_BEGIN:R3-003:residence
% HMT_ATOMIC_RESULT {"material_state":"CONSERVADO_POR_BASE_GRANULAR_REV3","proof_strength":"EXACTO_INTERNO","provenance":"RESULTADO_RECUPERADO","reinforcement_state":"DEPLOYED_WITHOUT_NEW_RESULT_ID","residence":"`hmt/02_trit_geometrias`, `03g`","result_id":"R3-003","statement":"TRIT orientado y temporal, con presente torsional y memoria","status":"PRESENT_IN_SUCCESSOR_SOURCE"}
% HMT_RESULT_END:R3-003:residence
% HMT_RESULT_BEGIN:R3-004:residence
% HMT_ATOMIC_RESULT {"material_state":"CONSERVADO_POR_BASE_GRANULAR_REV3","proof_strength":"EXACTO_INTERNO","provenance":"RESULTADO_RECUPERADO","reinforcement_state":"DEPLOYED_WITHOUT_NEW_RESULT_ID","residence":"/Users/ruben/Documents/HMT2/EDICION_AUTONOMA_HMT_MD_2026-08-09_REV_3_EN_TRABAJO/01_FUENTE_REVISADA/manuscrito/sections/hmt/03h_arbol_operatorio_tpk_rev11.tex; /Users/ruben/Documents/HMT2/EDICION_AUTONOMA_HMT_MD_2026-08-09_REV_3_EN_TRABAJO/01_FUENTE_REVISADA/manuscrito/sections/hmt/03i_inventario_operatorio_tpk_rev11.tex; /Users/ruben/Documents/HMT2/EDICION_AUTONOMA_HMT_MD_2026-08-09_REV_3_EN_TRABAJO/01_FUENTE_REVISADA/manuscrito/sections/hmt/03k_genealogia_operaciones_rev2.tex; /Users/ruben/Documents/HMT2/EDICION_AUTONOMA_HMT_MD_2026-08-09_REV_3_EN_TRABAJO/01_FUENTE_REVISADA/manuscrito/sections/hmt/03l_elevacion_cohomologica_dimension_rev2.tex; /Users/ruben/Documents/HMT2/EDICION_AUTONOMA_HMT_MD_2026-08-09_REV_3_EN_TRABAJO/01_FUENTE_REVISADA/manuscrito/sections/hmt/20_arquitectura_operatoria_tpk_actualizada.tex","result_id":"R3-004","statement":"TPK-full: objetos, rutas, memoria, frontera, prolongación, operadores y cálculo genealógico","status":"PRESENT_IN_SUCCESSOR_SOURCE"}
% HMT_RESULT_END:R3-004:residence
% HMT_RESULT_BEGIN:R3-005:residence
% HMT_ATOMIC_RESULT {"material_state":"CONSERVADO_POR_BASE_GRANULAR_REV3","proof_strength":"EXACTO_INTERNO","provenance":"RESULTADO_RECUPERADO","reinforcement_state":"DEPLOYED_WITHOUT_NEW_RESULT_ID","residence":"`hmt/00a`, `00b`, `14f`","result_id":"R3-005","statement":"Cuenta→medida, continuo como secciones compatibles y recta como evaluación","status":"PRESENT_IN_SUCCESSOR_SOURCE"}
% HMT_RESULT_END:R3-005:residence
% HMT_RESULT_BEGIN:R3-006:residence
% HMT_ATOMIC_RESULT {"material_state":"CONSERVADO_POR_BASE_GRANULAR_REV3","proof_strength":"EXACTO_INTERNO","provenance":"RESULTADO_RECUPERADO","reinforcement_state":"DEPLOYED_WITHOUT_NEW_RESULT_ID","residence":"`hmt/06aa`, `06b--06f`","result_id":"R3-006","statement":"Monodromía \\(9\\to1\\) con memoria, generación coinductiva de \\(\\pi,e,\\varphi\\), 396/729/43/387 y cinco regiones de \\(\\pi\\)","status":"PRESENT_IN_SUCCESSOR_SOURCE"}
% HMT_RESULT_END:R3-006:residence
% HMT_RESULT_BEGIN:R3-007:residence
% HMT_ATOMIC_RESULT {"material_state":"CONSERVADO_POR_BASE_GRANULAR_REV3","proof_strength":"EXACTO_INTERNO","provenance":"RESULTADO_RECUPERADO","reinforcement_state":"DEPLOYED_WITHOUT_NEW_RESULT_ID","residence":"capítulos 8, Planck, coordenadas conjugadas y Barbero","result_id":"R3-007","statement":"\\(\\alpha_{\\rm HMT}\\), Planck, Barbero y constantes en su orden causal","status":"PRESENT_IN_SUCCESSOR_SOURCE"}
% HMT_RESULT_END:R3-007:residence
% HMT_RESULT_BEGIN:R3-008:residence
% HMT_ATOMIC_RESULT {"material_state":"CONSERVADO_POR_BASE_GRANULAR_REV3","proof_strength":"EXACTO_INTERNO","provenance":"RESULTADO_RECUPERADO","reinforcement_state":"DEPLOYED_WITHOUT_NEW_RESULT_ID","residence":"`hmt/09`, `09g`, `09e`","result_id":"R3-008","statement":"Doble proyección desde el mismo estado holonómico integral","status":"PRESENT_IN_SUCCESSOR_SOURCE"}
% HMT_RESULT_END:R3-008:residence
% HMT_RESULT_BEGIN:R3-009:residence
% HMT_ATOMIC_RESULT {"material_state":"CONSERVADO_POR_BASE_GRANULAR_REV3","proof_strength":"EXACTO_INTERNO","provenance":"RESULTADO_RECUPERADO","reinforcement_state":"DEPLOYED_WITHOUT_NEW_RESULT_ID","residence":"`hmt/09c`, `15`, `15c`, `15d`","result_id":"R3-009","statement":"Corredor Paley--Witt--Golay--Leech--\\(V^\\natural\\)--Monstruo--Moonshine","status":"PRESENT_IN_SUCCESSOR_SOURCE"}
% HMT_RESULT_END:R3-009:residence
% HMT_RESULT_BEGIN:R3-010:residence
% HMT_ATOMIC_RESULT {"material_state":"CONSERVADO_POR_BASE_GRANULAR_REV3","proof_strength":"EXACTO_INTERNO","provenance":"RESULTADO_RECUPERADO","reinforcement_state":"DEPLOYED_WITHOUT_NEW_RESULT_ID","residence":"`md/00b_realizacion_espectral_numero_particula`","result_id":"R3-010","statement":"Número→partícula: sección, ruta, espectro, persistencia y frontera MD","status":"PRESENT_IN_SUCCESSOR_SOURCE"}
% HMT_RESULT_END:R3-010:residence
% HMT_RESULT_BEGIN:R3-011:residence
% HMT_ATOMIC_RESULT {"material_state":"CONSERVADO_POR_BASE_GRANULAR_REV3","proof_strength":"EXACTO_INTERNO","provenance":"RESULTADO_RECUPERADO","reinforcement_state":"DEPLOYED_WITHOUT_NEW_RESULT_ID","residence":"`md/04ab_postulados_cuanticos_genealogicos_rev2`","result_id":"R3-011","statement":"Postulados cuánticos HMT: Hilbert, rayos/densidades, observables, PVM, Born, Stone--Schrödinger y tensor","status":"PRESENT_IN_SUCCESSOR_SOURCE"}
% HMT_RESULT_END:R3-011:residence
% HMT_RESULT_BEGIN:R3-012:residence
% HMT_ATOMIC_RESULT {"material_state":"CONSERVADO_POR_BASE_GRANULAR_REV3","proof_strength":"EXACTO_INTERNO","provenance":"RESULTADO_RECUPERADO","reinforcement_state":"DEPLOYED_WITHOUT_NEW_RESULT_ID","residence":"`md/01a_medicion_cuantica_genealogica_rev2`","result_id":"R3-012","statement":"Medición: fibra, Lüders, Kraus, CPTP, Stinespring y PVM común","status":"PRESENT_IN_SUCCESSOR_SOURCE"}
% HMT_RESULT_END:R3-012:residence
% HMT_RESULT_BEGIN:R3-013:residence
% HMT_ATOMIC_RESULT {"material_state":"CONSERVADO_POR_BASE_GRANULAR_REV3","proof_strength":"EXACTO_INTERNO","provenance":"RESULTADO_RECUPERADO","reinforcement_state":"DEPLOYED_WITHOUT_NEW_RESULT_ID","residence":"`md/04aa`, capítulos cuánticos","result_id":"R3-013","statement":"Caminos, operador unitario \\(27\\times27\\), límite cilíndrico y reconstrucción cuántica","status":"PRESENT_IN_SUCCESSOR_SOURCE"}
% HMT_RESULT_END:R3-013:residence
% HMT_RESULT_BEGIN:R3-014:residence
% HMT_ATOMIC_RESULT {"material_state":"CONSERVADO_POR_BASE_GRANULAR_REV3","proof_strength":"EXACTO_INTERNO","provenance":"RESULTADO_RECUPERADO","reinforcement_state":"DEPLOYED_WITHOUT_NEW_RESULT_ID","residence":"`md/04d`, `04e`, `04g`, `04h`","result_id":"R3-014","statement":"Landauer finito/informacional y estadística cuántica tipada","status":"PRESENT_IN_SUCCESSOR_SOURCE"}
% HMT_RESULT_END:R3-014:residence
% HMT_RESULT_BEGIN:R3-015:residence
% HMT_ATOMIC_RESULT {"material_state":"CONSERVADO_POR_BASE_GRANULAR_REV3","proof_strength":"EXACTO_INTERNO","provenance":"RESULTADO_RECUPERADO","reinforcement_state":"DEPLOYED_WITHOUT_NEW_RESULT_ID","residence":"`md/03d`, `06k`, `06ka`","result_id":"R3-015","statement":"Planck--CT108--Compton, recta autodual y clausura electrónica de rango uno","status":"PRESENT_IN_SUCCESSOR_SOURCE"}
% HMT_RESULT_END:R3-015:residence
% HMT_RESULT_BEGIN:R3-016:residence
% HMT_ATOMIC_RESULT {"material_state":"CONSERVADO_POR_BASE_GRANULAR_REV3","proof_strength":"EXACTO_INTERNO","provenance":"RESULTADO_RECUPERADO","reinforcement_state":"DEPLOYED_WITHOUT_NEW_RESULT_ID","residence":"`md/05aa`, `05f`, `06m`, `06i`","result_id":"R3-016","statement":"Ley operatorial de masas: \\(10^{-34}\\), \\(K_D/K_\\Omega\\), 81/56/13/324, 23+23 y 17 condicionadas","status":"PRESENT_IN_SUCCESSOR_SOURCE"}
% HMT_RESULT_END:R3-016:residence
% HMT_RESULT_BEGIN:R3-017:residence
% HMT_ATOMIC_RESULT {"material_state":"CONSERVADO_POR_BASE_GRANULAR_REV3","proof_strength":"EXACTO_INTERNO","provenance":"RESULTADO_RECUPERADO","reinforcement_state":"DEPLOYED_WITHOUT_NEW_RESULT_ID","residence":"`md/10c`, `11`, `11c`, `11m`, `11p`","result_id":"R3-017","statement":"Cartan--Holst, gravitación, autoinercia y sus fronteras declaradas","status":"PRESENT_IN_SUCCESSOR_SOURCE"}
% HMT_RESULT_END:R3-017:residence
% HMT_RESULT_BEGIN:R3-018:residence
% HMT_ATOMIC_RESULT {"material_state":"CONSERVADO_POR_BASE_GRANULAR_REV3","proof_strength":"EXACTO_INTERNO","provenance":"RESULTADO_RECUPERADO","reinforcement_state":"DEPLOYED_WITHOUT_NEW_RESULT_ID","residence":"`hmt/17_sintesis`, `md/12_sintesis`","result_id":"R3-018","statement":"Síntesis causal APP→TRIT→TPK→MD, incluyendo defectos 4/2/6/54","status":"PRESENT_IN_SUCCESSOR_SOURCE"}
% HMT_RESULT_END:R3-018:residence
% HMT_RESULT_BEGIN:R3-019:residence
% HMT_ATOMIC_RESULT {"material_state":"CONSERVADO_POR_BASE_GRANULAR_REV3","proof_strength":"EXACTO_INTERNO","provenance":"RESULTADO_RECUPERADO","reinforcement_state":"DEPLOYED_WITHOUT_NEW_RESULT_ID","residence":"activos de REV.3","result_id":"R3-019","statement":"Portada, referencias, tablas y QA ya aptos","status":"PRESENT_IN_SUCCESSOR_SOURCE"}
% HMT_RESULT_END:R3-019:residence
% HMT_RESULT_BEGIN:R3-020:residence
% HMT_ATOMIC_RESULT {"material_state":"CONSERVADO_POR_BASE_GRANULAR_REV3","proof_strength":"EXACTO_INTERNO","provenance":"RESULTADO_RECUPERADO","reinforcement_state":"DEPLOYED_WITHOUT_NEW_RESULT_ID","residence":"grafo completo","result_id":"R3-020","statement":"Todo archivo y toda afirmación de REV.3 no afectados por el delta","status":"PRESENT_IN_SUCCESSOR_SOURCE"}
% HMT_RESULT_END:R3-020:residence
% HMT_RESULT_BEGIN:C-001:residence
% HMT_ATOMIC_RESULT {"material_state":"CONTROL_ACTIVO_EN_PLAN_Y_FUENTE","proof_strength":"EXACTO_INTERNO","provenance":"FORMALIZACION_NUEVA","reinforcement_state":"DEPLOYED_WITHOUT_NEW_RESULT_ID","residence":"gobierno transversal","result_id":"C-001","statement":"La monografía de 68 pp abrevia KMS como \\(F_{A,B}(t+i\\beta)=F_{B,A}(t)\\) bajo una definición que no lo permite","status":"CONTROLLED"}
% HMT_RESULT_END:C-001:residence
% HMT_RESULT_BEGIN:C-002:residence
% HMT_ATOMIC_RESULT {"material_state":"CONTROL_ACTIVO_EN_PLAN_Y_FUENTE","proof_strength":"EXACTO_INTERNO","provenance":"FORMALIZACION_NUEVA","reinforcement_state":"DEPLOYED_WITHOUT_NEW_RESULT_ID","residence":"gobierno transversal","result_id":"C-002","statement":"MOR-010 enumera siete piezas, mientras teorema y certificado tienen cuatro componentes","status":"CONTROLLED"}
% HMT_RESULT_END:C-002:residence
% HMT_RESULT_BEGIN:C-003:residence
% HMT_ATOMIC_RESULT {"material_state":"CONTROL_ACTIVO_EN_PLAN_Y_FUENTE","proof_strength":"EXACTO_INTERNO","provenance":"FORMALIZACION_NUEVA","reinforcement_state":"DEPLOYED_WITHOUT_NEW_RESULT_ID","residence":"gobierno transversal","result_id":"C-003","statement":"El verificador de 68 pp declara parte de la funtorialidad mediante `True`","status":"CONTROLLED"}
% HMT_RESULT_END:C-003:residence
% HMT_RESULT_BEGIN:C-004:residence
% HMT_ATOMIC_RESULT {"material_state":"CONTROL_ACTIVO_EN_PLAN_Y_FUENTE","proof_strength":"EXACTO_INTERNO","provenance":"FORMALIZACION_NUEVA","reinforcement_state":"DEPLOYED_WITHOUT_NEW_RESULT_ID","residence":"gobierno transversal","result_id":"C-004","statement":"QMS no conmutativos del puente son semiclásicos","status":"CONTROLLED"}
% HMT_RESULT_END:C-004:residence
% HMT_RESULT_BEGIN:C-005:residence
% HMT_ATOMIC_RESULT {"material_state":"CONTROL_ACTIVO_EN_PLAN_Y_FUENTE","proof_strength":"EXACTO_INTERNO","provenance":"FORMALIZACION_NUEVA","reinforcement_state":"DEPLOYED_WITHOUT_NEW_RESULT_ID","residence":"gobierno transversal","result_id":"C-005","statement":"El límite inverso HMT y el continuo MPS son objetos diferentes","status":"CONTROLLED"}
% HMT_RESULT_END:C-005:residence
% HMT_RESULT_BEGIN:C-006:residence
% HMT_ATOMIC_RESULT {"material_state":"CONTROL_ACTIVO_EN_PLAN_Y_FUENTE","proof_strength":"EXACTO_INTERNO","provenance":"FORMALIZACION_NUEVA","reinforcement_state":"DEPLOYED_WITHOUT_NEW_RESULT_ID","residence":"gobierno transversal","result_id":"C-006","statement":"`06_teorema_morfismo.tex` presupone un dominio común y hace preservar inversas a magnitudes positivas; además sugerir que la realización positiva se recupera de la orientada pierde variación total","status":"CONTROLLED"}
% HMT_RESULT_END:C-006:residence
% HMT_RESULT_BEGIN:C-007:residence
% HMT_ATOMIC_RESULT {"material_state":"CONTROL_ACTIVO_EN_PLAN_Y_FUENTE","proof_strength":"EXACTO_INTERNO","provenance":"FORMALIZACION_NUEVA","reinforcement_state":"DEPLOYED_WITHOUT_NEW_RESULT_ID","residence":"gobierno transversal","result_id":"C-007","statement":"El apéndice del puente matricial afirma que sus seis no-certificaciones se incluyen en el JSON, pero `scope_limits` contiene cinco controles distintos","status":"CONTROLLED"}
% HMT_RESULT_END:C-007:residence
% HMT_RESULT_BEGIN:MAS-RECTOR-001:residence
% HMT_ATOMIC_RESULT {"material_state":"CONSERVADO_CON_APTO_FOCAL","proof_strength":"EXACTO_INTERNO","provenance":"RESULTADO_RECUPERADO","reinforcement_state":"DEPLOYED_WITHOUT_NEW_RESULT_ID","residence":"/Users/ruben/Documents/HMT2/EDICION_AUTONOMA_HMT_MD_2026-08-09_REV_3_EN_TRABAJO/01_FUENTE_REVISADA/manuscrito/sections/md/06h_atlas_familias_rev6.tex; /Users/ruben/Documents/HMT2/EDICION_AUTONOMA_HMT_MD_2026-08-09_REV_3_EN_TRABAJO/01_FUENTE_REVISADA/manuscrito/sections/md/06i_inventario_humano_23_sectores_rev8.tex","result_id":"MAS-RECTOR-001","statement":"81 celdas; 56 familias; 13 multisecciones; 324 rutas orientadas; sector nulo separado","status":"PRESENT_IN_SUCCESSOR_SOURCE"}
% HMT_RESULT_END:MAS-RECTOR-001:residence
% HMT_RESULT_BEGIN:MAS-RECTOR-002:residence
% HMT_ATOMIC_RESULT {"material_state":"CONSERVADO_CON_APTO_FOCAL","proof_strength":"EXACTO_INTERNO","provenance":"ARQUITECTURA_AUTORAL_PREEXISTENTE","reinforcement_state":"DEPLOYED_WITHOUT_NEW_RESULT_ID","residence":"/Users/ruben/Documents/HMT2/EDICION_AUTONOMA_HMT_MD_2026-08-09_REV_3_EN_TRABAJO/01_FUENTE_REVISADA/manuscrito/sections/md/00b_realizacion_espectral_numero_particula.tex; /Users/ruben/Documents/HMT2/EDICION_AUTONOMA_HMT_MD_2026-08-09_REV_3_EN_TRABAJO/01_FUENTE_REVISADA/manuscrito/sections/md/05aa_ley_general_masa_accion_autonoma.tex","result_id":"MAS-RECTOR-002","statement":"extensión superviviente→ruta→ledger→firma Z5→carácter→torres→recta de masa","status":"PRESENT_IN_SUCCESSOR_SOURCE"}
% HMT_RESULT_END:MAS-RECTOR-002:residence
% HMT_RESULT_BEGIN:MAS-RECTOR-003:residence
% HMT_ATOMIC_RESULT {"material_state":"CONSERVADO_CON_APTO_FOCAL","proof_strength":"EXACTO_INTERNO","provenance":"RESULTADO_RECUPERADO","reinforcement_state":"DEPLOYED_WITHOUT_NEW_RESULT_ID","residence":"/Users/ruben/Documents/HMT2/EDICION_AUTONOMA_HMT_MD_2026-08-09_REV_3_EN_TRABAJO/01_FUENTE_REVISADA/manuscrito/sections/hmt/03_app_ortograma.tex","result_id":"MAS-RECTOR-003","statement":"Bc=9P_++5P_-; salto primordial=4; defecto global posterior 2→6→54","status":"PRESENT_IN_SUCCESSOR_SOURCE"}
% HMT_RESULT_END:MAS-RECTOR-003:residence
% HMT_RESULT_BEGIN:MAS-RECTOR-004:residence
% HMT_ATOMIC_RESULT {"material_state":"CONSERVADO_CON_APTO_FOCAL","proof_strength":"EXACTO_INTERNO","provenance":"RESULTADO_RECUPERADO","reinforcement_state":"DEPLOYED_WITHOUT_NEW_RESULT_ID","residence":"/Users/ruben/Documents/HMT2/EDICION_AUTONOMA_HMT_MD_2026-08-09_REV_3_EN_TRABAJO/01_FUENTE_REVISADA/manuscrito/sections/md/03b_escala_accion_determinantal.tex","result_id":"MAS-RECTOR-004","statement":"SNF(H4)=(1,2,2,4); |coker|=16; Sigma_H=phi alpha^16; ord10=-34","status":"PRESENT_IN_SUCCESSOR_SOURCE"}
% HMT_RESULT_END:MAS-RECTOR-004:residence
% HMT_RESULT_BEGIN:MAS-RECTOR-005:residence
% HMT_ATOMIC_RESULT {"material_state":"CONSERVADO_CON_APTO_FOCAL","proof_strength":"DEMOSTRADO_CON_ESTRUCTURA_DE_PARTIDA_EXPLICITA","provenance":"RESULTADO_RECUPERADO","reinforcement_state":"DEPLOYED_WITHOUT_NEW_RESULT_ID","residence":"/Users/ruben/Documents/HMT2/EDICION_AUTONOMA_HMT_MD_2026-08-09_REV_3_EN_TRABAJO/01_FUENTE_REVISADA/manuscrito/sections/md/03b_escala_accion_determinantal.tex; /Users/ruben/Documents/HMT2/EDICION_AUTONOMA_HMT_MD_2026-08-09_REV_3_EN_TRABAJO/01_FUENTE_REVISADA/manuscrito/sections/md/03c_elipse_accion_holonomia_rev8.tex","result_id":"MAS-RECTOR-005","statement":"delta_2w=H5-etaRet=(90/pi)Cpi(alpha); R_act=H5/etaRet","status":"PRESENT_IN_SUCCESSOR_SOURCE"}
% HMT_RESULT_END:MAS-RECTOR-005:residence
% HMT_RESULT_BEGIN:MAS-RECTOR-006:residence
% HMT_ATOMIC_RESULT {"material_state":"CONSERVADO_CON_APTO_FOCAL","proof_strength":"DEMOSTRADO_CON_ESTRUCTURA_DE_PARTIDA_EXPLICITA","provenance":"FORMALIZACION_NUEVA","reinforcement_state":"DEPLOYED_WITHOUT_NEW_RESULT_ID","residence":"/Users/ruben/Documents/HMT2/EDICION_AUTONOMA_HMT_MD_2026-08-09_REV_3_EN_TRABAJO/01_FUENTE_REVISADA/manuscrito/sections/md/03d_identidad_autodual_planck_ct108_rev3.tex; /Users/ruben/Documents/HMT2/EDICION_AUTONOMA_HMT_MD_2026-08-09_REV_3_EN_TRABAJO/01_FUENTE_REVISADA/manuscrito/sections/md/05aa_ley_general_masa_accion_autonoma.tex","result_id":"MAS-RECTOR-006","statement":"hbar_ret=etaRet 10^-34 U_S; omega0=2pi/(108t0); m_clk=hbar_ret omega0 c^-2","status":"PRESENT_IN_SUCCESSOR_SOURCE"}
% HMT_RESULT_END:MAS-RECTOR-006:residence
% HMT_RESULT_BEGIN:MAS-RECTOR-007:residence
% HMT_ATOMIC_RESULT {"material_state":"CONSERVADO_CON_APTO_FOCAL","proof_strength":"EXACTO_INTERNO","provenance":"RESULTADO_RECUPERADO","reinforcement_state":"DEPLOYED_WITHOUT_NEW_RESULT_ID","residence":"/Users/ruben/Documents/HMT2/EDICION_AUTONOMA_HMT_MD_2026-08-09_REV_3_EN_TRABAJO/01_FUENTE_REVISADA/manuscrito/sections/md/05f_operadores_masa_two_way_rev11.tex","result_id":"MAS-RECTOR-007","statement":"K_D=(D27+D36,D1,(D4-D3)/2); 14 perfiles","status":"PRESENT_IN_SUCCESSOR_SOURCE"}
% HMT_RESULT_END:MAS-RECTOR-007:residence
% HMT_RESULT_BEGIN:MAS-RECTOR-008:residence
% HMT_ATOMIC_RESULT {"material_state":"CONSERVADO_CON_APTO_FOCAL","proof_strength":"EXACTO_INTERNO","provenance":"RESULTADO_RECUPERADO","reinforcement_state":"DEPLOYED_WITHOUT_NEW_RESULT_ID","residence":"/Users/ruben/Documents/HMT2/EDICION_AUTONOMA_HMT_MD_2026-08-09_REV_3_EN_TRABAJO/01_FUENTE_REVISADA/manuscrito/sections/md/05f_operadores_masa_two_way_rev11.tex","result_id":"MAS-RECTOR-008","statement":"K_Omega=(Omega,54,P6); 9 perfiles","status":"PRESENT_IN_SUCCESSOR_SOURCE"}
% HMT_RESULT_END:MAS-RECTOR-008:residence
% HMT_RESULT_BEGIN:MAS-RECTOR-009:residence
% HMT_ATOMIC_RESULT {"material_state":"CONSERVADO_CON_APTO_FOCAL","proof_strength":"EXACTO_INTERNO","provenance":"CERTIFICADO_NUEVO","reinforcement_state":"DEPLOYED_WITHOUT_NEW_RESULT_ID","residence":"/Users/ruben/Documents/HMT2/EDICION_AUTONOMA_HMT_MD_2026-08-09_REV_3_EN_TRABAJO/01_FUENTE_REVISADA/manuscrito/sections/md/05f_operadores_masa_two_way_rev11.tex","result_id":"MAS-RECTOR-009","statement":"40 pares; 18 coincidencias; valor común único (80,54,6)","status":"PRESENT_IN_SUCCESSOR_SOURCE"}
% HMT_RESULT_END:MAS-RECTOR-009:residence
% HMT_RESULT_BEGIN:MAS-RECTOR-010:residence
% HMT_ATOMIC_RESULT {"material_state":"CONSERVADO_CON_APTO_FOCAL","proof_strength":"EXACTO_INTERNO","provenance":"RESULTADO_RECUPERADO","reinforcement_state":"DEPLOYED_WITHOUT_NEW_RESULT_ID","residence":"/Users/ruben/Documents/HMT2/EDICION_AUTONOMA_HMT_MD_2026-08-09_REV_3_EN_TRABAJO/01_FUENTE_REVISADA/manuscrito/sections/md/06f_biografia_electron_rev6.tex","result_id":"MAS-RECTOR-010","statement":"(5,5,++); tau12=+++---+++---; D=(54,56,68,48,32)","status":"PRESENT_IN_SUCCESSOR_SOURCE"}
% HMT_RESULT_END:MAS-RECTOR-010:residence
% HMT_RESULT_BEGIN:MAS-RECTOR-011:residence
% HMT_ATOMIC_RESULT {"material_state":"CONSERVADO_CON_APTO_FOCAL","proof_strength":"DEMOSTRADO_CON_ESTRUCTURA_DE_PARTIDA_EXPLICITA","provenance":"RESULTADO_RECUPERADO","reinforcement_state":"DEPLOYED_WITHOUT_NEW_RESULT_ID","residence":"/Users/ruben/Documents/HMT2/EDICION_AUTONOMA_HMT_MD_2026-08-09_REV_3_EN_TRABAJO/01_FUENTE_REVISADA/manuscrito/sections/md/06f_biografia_electron_rev6.tex; /Users/ruben/Documents/HMT2/EDICION_AUTONOMA_HMT_MD_2026-08-09_REV_3_EN_TRABAJO/01_FUENTE_REVISADA/manuscrito/sections/md/06ka_clausura_nulo_material_electron_rev3.tex","result_id":"MAS-RECTOR-011","statement":"e0=(sqrt3/4)[exp(phi/pi^2)-22alpha^3](1+15Delta4)=0.51099892248806607250987087719134943763; etapa basal, no salida final","status":"PRESENT_IN_SUCCESSOR_SOURCE"}
% HMT_RESULT_END:MAS-RECTOR-011:residence
% HMT_RESULT_BEGIN:MAS-RECTOR-012:residence
% HMT_ATOMIC_RESULT {"material_state":"CONSERVADO_CON_APTO_FOCAL","proof_strength":"EXACTO_INTERNO","provenance":"RESULTADO_RECUPERADO","reinforcement_state":"DEPLOYED_WITHOUT_NEW_RESULT_ID","residence":"/Users/ruben/Documents/HMT2/EDICION_AUTONOMA_HMT_MD_2026-08-09_REV_3_EN_TRABAJO/01_FUENTE_REVISADA/manuscrito/sections/md/06k_electron_two_way_rev11.tex; /Users/ruben/Documents/HMT2/EDICION_AUTONOMA_HMT_MD_2026-08-09_REV_3_EN_TRABAJO/01_FUENTE_REVISADA/manuscrito/sections/md/06ka_clausura_nulo_material_electron_rev3.tex","result_id":"MAS-RECTOR-012","statement":"kappa_e=80-54alpha+6Delta4; m_e^beta=e0 R_act^kappa_e=0.51099895069000080860825716483792339562463486771513350353936201542059113716425185","status":"PRESENT_IN_SUCCESSOR_SOURCE"}
% HMT_RESULT_END:MAS-RECTOR-012:residence
% HMT_RESULT_BEGIN:MAS-RECTOR-013:residence
% HMT_ATOMIC_RESULT {"material_state":"CONSERVADO_CON_APTO_FOCAL","proof_strength":"EXACTO_INTERNO","provenance":"RESULTADO_RECUPERADO","reinforcement_state":"DEPLOYED_WITHOUT_NEW_RESULT_ID","residence":"/Users/ruben/Documents/HMT2/EDICION_AUTONOMA_HMT_MD_2026-08-09_REV_3_EN_TRABAJO/01_FUENTE_REVISADA/manuscrito/sections/md/06f_biografia_electron_rev6.tex; /Users/ruben/Documents/HMT2/EDICION_AUTONOMA_HMT_MD_2026-08-09_REV_3_EN_TRABAJO/01_FUENTE_REVISADA/manuscrito/sections/md/06k_electron_two_way_rev11.tex","result_id":"MAS-RECTOR-013","statement":"firma cruda; firma par CT108; eventos; v_e=0 no se fusionan","status":"PRESENT_IN_SUCCESSOR_SOURCE"}
% HMT_RESULT_END:MAS-RECTOR-013:residence
% HMT_RESULT_BEGIN:MAS-RECTOR-014:residence
% HMT_ATOMIC_RESULT {"material_state":"CONSERVADO_CON_APTO_FOCAL","proof_strength":"EXACTO_INTERNO","provenance":"RESULTADO_RECUPERADO","reinforcement_state":"DEPLOYED_WITHOUT_NEW_RESULT_ID","residence":"/Users/ruben/Documents/HMT2/EDICION_AUTONOMA_HMT_MD_2026-08-09_REV_3_EN_TRABAJO/01_FUENTE_REVISADA/manuscrito/sections/md/05a_jerarquia_espectral_infinita_rev7.tex","result_id":"MAS-RECTOR-014","statement":"s_n=q_-^n-q_+^n; c_n=q_-^n+q_+^n; S=<90,120>={90,120}∪{30r:r≥6}","status":"PRESENT_IN_SUCCESSOR_SOURCE"}
% HMT_RESULT_END:MAS-RECTOR-014:residence
% HMT_RESULT_BEGIN:MAS-RECTOR-015:residence
% HMT_ATOMIC_RESULT {"material_state":"CONSERVADO_CON_APTO_FOCAL","proof_strength":"EXACTO_INTERNO","provenance":"FORMALIZACION_NUEVA","reinforcement_state":"DEPLOYED_WITHOUT_NEW_RESULT_ID","residence":"/Users/ruben/Documents/HMT2/EDICION_AUTONOMA_HMT_MD_2026-08-09_REV_3_EN_TRABAJO/01_FUENTE_REVISADA/manuscrito/sections/md/05a_jerarquia_espectral_infinita_rev7.tex","result_id":"MAS-RECTOR-015","statement":"k[S0]=k[X90,Y120]/(X90^4-Y120^3); dos genealogías a 360 antes del cociente","status":"PRESENT_IN_SUCCESSOR_SOURCE"}
% HMT_RESULT_END:MAS-RECTOR-015:residence
% HMT_RESULT_BEGIN:MAS-RECTOR-016:residence
% HMT_ATOMIC_RESULT {"material_state":"CONSERVADO_CON_APTO_FOCAL","proof_strength":"EXACTO_INTERNO","provenance":"RESULTADO_RECUPERADO","reinforcement_state":"DEPLOYED_WITHOUT_NEW_RESULT_ID","residence":"/Users/ruben/Documents/HMT2/EDICION_AUTONOMA_HMT_MD_2026-08-09_REV_3_EN_TRABAJO/01_FUENTE_REVISADA/manuscrito/sections/md/05f_operadores_masa_two_way_rev11.tex","result_id":"MAS-RECTOR-016","statement":"B_M^D y B_M^Omega diagonales exactos con espectros completos","status":"PRESENT_IN_SUCCESSOR_SOURCE"}
% HMT_RESULT_END:MAS-RECTOR-016:residence
% HMT_RESULT_BEGIN:MAS-RECTOR-017:residence
% HMT_ATOMIC_RESULT {"material_state":"CONSERVADO_CON_APTO_FOCAL","proof_strength":"EXACTO_INTERNO","provenance":"FORMALIZACION_NUEVA","reinforcement_state":"DEPLOYED_WITHOUT_NEW_RESULT_ID","residence":"/Users/ruben/Documents/HMT2/EDICION_AUTONOMA_HMT_MD_2026-08-09_REV_3_EN_TRABAJO/01_FUENTE_REVISADA/manuscrito/sections/md/05f_operadores_masa_two_way_rev11.tex","result_id":"MAS-RECTOR-017","statement":"M_M^{beta,r}=R_act^{B_M^r/2} M_M^(0) R_act^{B_M^r/2}","status":"PRESENT_IN_SUCCESSOR_SOURCE"}
% HMT_RESULT_END:MAS-RECTOR-017:residence
% HMT_RESULT_BEGIN:MAS-RECTOR-018:residence
% HMT_ATOMIC_RESULT {"material_state":"CONSERVADO_CON_APTO_FOCAL","proof_strength":"EXACTO_INTERNO","provenance":"CERTIFICADO_NUEVO","reinforcement_state":"DEPLOYED_WITHOUT_NEW_RESULT_ID","residence":"/Users/ruben/Documents/HMT2/EDICION_AUTONOMA_HMT_MD_2026-08-09_REV_3_EN_TRABAJO/01_FUENTE_REVISADA/manuscrito/sections/md/06m_tabla_completa_masas_hmt_sm_revfinal.tex","result_id":"MAS-RECTOR-018","statement":"Toda clase conserva residencia, operador/salida y estatuto; electrón beta y nulos cerrados; la publicación mantiene 23 filas de salida interna y 23 filas de comparabilidad física, una por clase","status":"PRESENT_IN_SUCCESSOR_SOURCE"}
% HMT_RESULT_END:MAS-RECTOR-018:residence
% HMT_RESULT_BEGIN:MAS-RECTOR-019:residence
% HMT_ATOMIC_RESULT {"material_state":"CONSERVADO_CON_APTO_FOCAL","proof_strength":"EXACTO_INTERNO","provenance":"RESULTADO_RECUPERADO","reinforcement_state":"DEPLOYED_WITHOUT_NEW_RESULT_ID","residence":"/Users/ruben/Documents/HMT2/EDICION_AUTONOMA_HMT_MD_2026-08-09_REV_3_EN_TRABAJO/01_FUENTE_REVISADA/manuscrito/sections/md/06m_tabla_completa_masas_hmt_sm_revfinal.tex","result_id":"MAS-RECTOR-019","statement":"m=0 en carta nula; kappa no aplicable","status":"PRESENT_IN_SUCCESSOR_SOURCE"}
% HMT_RESULT_END:MAS-RECTOR-019:residence
% HMT_RESULT_BEGIN:MAS-RECTOR-020:residence
% HMT_ATOMIC_RESULT {"material_state":"CONSERVADO_CON_APTO_FOCAL","proof_strength":"INTERFAZ_FISICA","provenance":"RESULTADO_RECUPERADO","reinforcement_state":"DEPLOYED_WITHOUT_NEW_RESULT_ID","residence":"/Users/ruben/Documents/HMT2/EDICION_AUTONOMA_HMT_MD_2026-08-09_REV_3_EN_TRABAJO/01_FUENTE_REVISADA/manuscrito/sections/md/06i_inventario_humano_23_sectores_rev8.tex; /Users/ruben/Documents/HMT2/EDICION_AUTONOMA_HMT_MD_2026-08-09_REV_3_EN_TRABAJO/01_FUENTE_REVISADA/manuscrito/sections/md/06m_tabla_completa_masas_hmt_sm_revfinal.tex","result_id":"MAS-RECTOR-020","statement":"Rama up/down, generación, color y carácter fino; scalarización exige join individual y esquema/escala","status":"PRESENT_IN_SUCCESSOR_SOURCE"}
% HMT_RESULT_END:MAS-RECTOR-020:residence
% HMT_RESULT_BEGIN:MAS-RECTOR-021:residence
% HMT_ATOMIC_RESULT {"material_state":"CONSERVADO_CON_APTO_FOCAL","proof_strength":"INTERFAZ_FISICA","provenance":"RESULTADO_RECUPERADO","reinforcement_state":"DEPLOYED_WITHOUT_NEW_RESULT_ID","residence":"/Users/ruben/Documents/HMT2/EDICION_AUTONOMA_HMT_MD_2026-08-09_REV_3_EN_TRABAJO/01_FUENTE_REVISADA/manuscrito/sections/md/06i_inventario_humano_23_sectores_rev8.tex; /Users/ruben/Documents/HMT2/EDICION_AUTONOMA_HMT_MD_2026-08-09_REV_3_EN_TRABAJO/01_FUENTE_REVISADA/manuscrito/sections/md/08b_transporte_ckm_pmns_rev11.tex; /Users/ruben/Documents/HMT2/EDICION_AUTONOMA_HMT_MD_2026-08-09_REV_3_EN_TRABAJO/01_FUENTE_REVISADA/manuscrito/sections/md/05a1_generaciones_quarks_neutrinos_rev2.tex","result_id":"MAS-RECTOR-021","statement":"Tres clases de sabor neutrínico sobre el par operatorial común B_F^D/B_F^Omega; los tres sabores comparten H_nu=directsum_{j=1}^4 H_{nu,G_j} de dimensión 20; f_e,f_mu,f_tau es base de interacción y G_j profundidad estructural; PMNS transporta bases y no crea autovalores, escala ni masa de sabor","status":"PRESENT_IN_SUCCESSOR_SOURCE"}
% HMT_RESULT_END:MAS-RECTOR-021:residence
% HMT_RESULT_BEGIN:MAS-RECTOR-022:residence
% HMT_ATOMIC_RESULT {"material_state":"CONSERVADO_CON_APTO_FOCAL","proof_strength":"RECONOCIMIENTO_EXTERNO","provenance":"RESULTADO_RECUPERADO","reinforcement_state":"DEPLOYED_WITHOUT_NEW_RESULT_ID","residence":"/Users/ruben/Documents/HMT2/EDICION_AUTONOMA_HMT_MD_2026-08-09_REV_3_EN_TRABAJO/01_FUENTE_REVISADA/manuscrito/sections/md/08_ckm_cp.tex; /Users/ruben/Documents/HMT2/EDICION_AUTONOMA_HMT_MD_2026-08-09_REV_3_EN_TRABAJO/01_FUENTE_REVISADA/manuscrito/sections/md/08b_transporte_ckm_pmns_rev11.tex","result_id":"MAS-RECTOR-022","statement":"Cuatro fórmulas angulares, matriz unitaria y J; contraste posterior","status":"PRESENT_IN_SUCCESSOR_SOURCE"}
% HMT_RESULT_END:MAS-RECTOR-022:residence
% HMT_RESULT_BEGIN:MAS-RECTOR-023:residence
% HMT_ATOMIC_RESULT {"material_state":"CONSERVADO_CON_APTO_FOCAL","proof_strength":"PROGRAMA_ABIERTO","provenance":"RESULTADO_RECUPERADO","reinforcement_state":"DEPLOYED_WITHOUT_NEW_RESULT_ID","residence":"/Users/ruben/Documents/HMT2/EDICION_AUTONOMA_HMT_MD_2026-08-09_REV_3_EN_TRABAJO/01_FUENTE_REVISADA/manuscrito/sections/md/06g_composicion_especies_rev6.tex","result_id":"MAS-RECTOR-023","statement":"Rutas compuestas con frontera/enlace; 432 y 1728 expresiones no son especies","status":"PRESENT_IN_SUCCESSOR_SOURCE"}
% HMT_RESULT_END:MAS-RECTOR-023:residence
% HMT_RESULT_BEGIN:MAS-RECTOR-024:residence
% HMT_ATOMIC_RESULT {"material_state":"CONSERVADO_CON_APTO_FOCAL","proof_strength":"EXACTO_INTERNO","provenance":"RESULTADO_RECUPERADO","reinforcement_state":"DEPLOYED_WITHOUT_NEW_RESULT_ID","residence":"/Users/ruben/Documents/HMT2/EDICION_AUTONOMA_HMT_MD_2026-08-09_REV_3_EN_TRABAJO/01_FUENTE_REVISADA/manuscrito/sections/md/06c_sectores_topologicos_no_omision.tex","result_id":"MAS-RECTOR-024","statement":"Anillos, dimensiones y trenzado son codominios propios, no masas escalares automáticas","status":"PRESENT_IN_SUCCESSOR_SOURCE"}
% HMT_RESULT_END:MAS-RECTOR-024:residence
% HMT_RESULT_BEGIN:MAS-RECTOR-025:residence
% HMT_ATOMIC_RESULT {"material_state":"CONSERVADO_CON_APTO_FOCAL","proof_strength":"EXACTO_INTERNO","provenance":"RESULTADO_RECUPERADO","reinforcement_state":"DEPLOYED_WITHOUT_NEW_RESULT_ID","residence":"/Users/ruben/Documents/HMT2/EDICION_AUTONOMA_HMT_MD_2026-08-09_REV_3_EN_TRABAJO/01_FUENTE_REVISADA/manuscrito/sections/md/04b_banda_particula_virtual.tex","result_id":"MAS-RECTOR-025","statement":"Sección de prolongación finita con ledger hasta L; no especie estable ni masa imaginaria","status":"PRESENT_IN_SUCCESSOR_SOURCE"}
% HMT_RESULT_END:MAS-RECTOR-025:residence
% HMT_RESULT_BEGIN:MAS-RECTOR-026:residence
% HMT_ATOMIC_RESULT {"material_state":"CONSERVADO_CON_APTO_FOCAL","proof_strength":"EXACTO_INTERNO","provenance":"CERTIFICADO_NUEVO","reinforcement_state":"DEPLOYED_WITHOUT_NEW_RESULT_ID","residence":"/Users/ruben/Documents/HMT2/EDICION_AUTONOMA_HMT_MD_2026-08-09_REV_3_EN_TRABAJO/01_FUENTE_REVISADA/manuscrito/sections/md/06m_tabla_completa_masas_hmt_sm_revfinal.tex","result_id":"MAS-RECTOR-026","statement":"Salida interna congelada antes de abrir el catálogo externo","status":"PRESENT_IN_SUCCESSOR_SOURCE"}
% HMT_RESULT_END:MAS-RECTOR-026:residence
% HMT_RESULT_BEGIN:MAS-RECTOR-027:residence
% HMT_ATOMIC_RESULT {"material_state":"CONSERVADO_CON_APTO_FOCAL","proof_strength":"EXACTO_INTERNO","provenance":"RESULTADO_RECUPERADO","reinforcement_state":"DEPLOYED_WITHOUT_NEW_RESULT_ID","residence":"/Users/ruben/Documents/HMT2/EDICION_AUTONOMA_HMT_MD_2026-08-09_REV_3_EN_TRABAJO/01_FUENTE_REVISADA/manuscrito/sections/md/06m_tabla_completa_masas_hmt_sm_revfinal.tex","result_id":"MAS-RECTOR-027","statement":"Sólo trazabilidad histórica; no residencia narrativa pública ni dominio","status":"PRESENT_IN_SUCCESSOR_SOURCE"}
% HMT_RESULT_END:MAS-RECTOR-027:residence
% HMT_RESULT_BEGIN:MAS-RECTOR-028:residence
% HMT_ATOMIC_RESULT {"material_state":"CONSERVADO_CON_APTO_FOCAL","proof_strength":"DEMOSTRADO_CON_ESTRUCTURA_DE_PARTIDA_EXPLICITA","provenance":"RESULTADO_RECUPERADO","reinforcement_state":"DEPLOYED_WITHOUT_NEW_RESULT_ID","residence":"/Users/ruben/Documents/HMT2/EDICION_AUTONOMA_HMT_MD_2026-08-09_REV_3_EN_TRABAJO/01_FUENTE_REVISADA/manuscrito/sections/md/06l_vacancias_corriente_rev11.tex","result_id":"MAS-RECTOR-028","statement":"Matriz [[21,22],[6,7]] con det=15 acompaña (-22,+15); uso local, no global","status":"PRESENT_IN_SUCCESSOR_SOURCE"}
% HMT_RESULT_END:MAS-RECTOR-028:residence
% HMT_RESULT_BEGIN:MAS-PUBLIC-023:residence
% HMT_ATOMIC_RESULT {"material_state":"CONSERVADO_CON_APTO_FOCAL","proof_strength":"EXACTO_INTERNO","provenance":"CERTIFICADO_NUEVO","reinforcement_state":"DEPLOYED_WITHOUT_NEW_RESULT_ID","residence":"/Users/ruben/Documents/HMT2/EDICION_AUTONOMA_HMT_MD_2026-08-09_REV_3_EN_TRABAJO/01_FUENTE_REVISADA/manuscrito/sections/md/06m_tabla_completa_masas_hmt_sm_revfinal.tex","result_id":"MAS-PUBLIC-023","statement":"las 23 claves internas reaparecen una vez y en el mismo orden en la capa pública de comparabilidad; las seis clases de codominio no escalar se comparan en su codominio nativo y no reciben una masa universal inventada","status":"PRESENT_IN_SUCCESSOR_SOURCE"}
% HMT_RESULT_END:MAS-PUBLIC-023:residence
% HMT_RESULT_BEGIN:MAS-COND-017:residence
% HMT_ATOMIC_RESULT {"material_state":"CONSERVADO_CON_APTO_FOCAL","proof_strength":"INTERFAZ_FISICA","provenance":"RESULTADO_RECUPERADO","reinforcement_state":"DEPLOYED_WITHOUT_NEW_RESULT_ID","residence":"/Users/ruben/Documents/HMT2/EDICION_AUTONOMA_HMT_MD_2026-08-09_REV_3_EN_TRABAJO/01_FUENTE_REVISADA/manuscrito/sections/md/06m_tabla_completa_masas_hmt_sm_revfinal.tex","result_id":"MAS-COND-017","statement":"17 evaluaciones exactas, 11 elementales y 6 compuestas, residen en una tabla separada y rotulada como no prospectiva; no son salidas cerradas, no seleccionan ruta y sólo admiten residual cuando los escalares, unidad, esquema y escala son comparables","status":"PRESENT_IN_SUCCESSOR_SOURCE"}
% HMT_RESULT_END:MAS-COND-017:residence
% HMT_RESULT_BEGIN:INF-001:residence
% HMT_ATOMIC_RESULT {"material_state":"INTEGRADO_O_REMITIDO_SIN_PERDIDA","proof_strength":"EXACTO_INTERNO","provenance":"RESULTADO_RECUPERADO","reinforcement_state":"DEPLOYED_WITHOUT_NEW_RESULT_ID","residence":"apertura del capítulo del continuo, tras `hmt/00b`","result_id":"INF-001","statement":"El continuo HMT no se forma con valores desnudos: cada casilla transporta bloque, firma, acarreo, fase, etiqueta y orientación. \\(\\xi_m=(B_m,\\sigma_m,c_m,\\omega_m)\\).","status":"PRESENT_IN_SUCCESSOR_SOURCE"}
% HMT_RESULT_END:INF-001:residence
% HMT_RESULT_BEGIN:INF-002:residence
% HMT_ATOMIC_RESULT {"material_state":"INTEGRADO_O_REMITIDO_SIN_PERDIDA","proof_strength":"EXACTO_INTERNO","provenance":"RESULTADO_RECUPERADO","reinforcement_state":"DEPLOYED_WITHOUT_NEW_RESULT_ID","residence":"bloque fundacional, § Estado, sección y valor, tras INF-001","result_id":"INF-002","statement":"Un número HMT es una sección compatible, no su expansión ni un conjunto acabado sin genealogía. \\(\\Omega_\\infty=\\varprojlim_m\\mathcal S_m\\).","status":"PRESENT_IN_SUCCESSOR_SOURCE"}
% HMT_RESULT_END:INF-002:residence
% HMT_RESULT_BEGIN:INF-003:residence
% HMT_ATOMIC_RESULT {"material_state":"INTEGRADO_O_REMITIDO_SIN_PERDIDA","proof_strength":"DEMOSTRADO_CON_ESTRUCTURA_DE_PARTIDA_EXPLICITA","provenance":"RESULTADO_RECUPERADO","reinforcement_state":"DEPLOYED_WITHOUT_NEW_RESULT_ID","residence":"bloque fundacional, § Existencia y unicidad coinductivas, tras INF-002","result_id":"INF-003","statement":"Finitud, no vaciedad y sobreyectividad dan existencia por König; la unicidad exige además fibras unitarias compatibles a toda profundidad.","status":"PRESENT_IN_SUCCESSOR_SOURCE"}
% HMT_RESULT_END:INF-003:residence
% HMT_RESULT_BEGIN:N9-001:residence
% HMT_ATOMIC_RESULT {"material_state":"INTEGRADO_O_REMITIDO_SIN_PERDIDA","proof_strength":"DEMOSTRADO_CON_ESTRUCTURA_DE_PARTIDA_EXPLICITA","provenance":"ARQUITECTURA_AUTORAL_PREEXISTENTE","reinforcement_state":"DEPLOYED_WITHOUT_NEW_RESULT_ID","residence":"`hmt/06aa` y transición inicial","result_id":"N9-001","statement":"La cadena APP--TRIT--TPK abre la frontera coinductiva; el retorno de fase no es retorno de estado.","status":"PRESENT_IN_SUCCESSOR_SOURCE"}
% HMT_RESULT_END:N9-001:residence
% HMT_RESULT_BEGIN:N9-002:residence
% HMT_ATOMIC_RESULT {"material_state":"INTEGRADO_O_REMITIDO_SIN_PERDIDA","proof_strength":"EXACTO_INTERNO","provenance":"RESULTADO_RECUPERADO","reinforcement_state":"DEPLOYED_WITHOUT_NEW_RESULT_ID","residence":"`hmt/06aa`","result_id":"N9-002","statement":"Las tres evaluaciones son lecturas coordinadas de un mismo estado; \\(\\varphi\\) fija el eje y \\(\\pi/e\\) constituyen la fibra especular orientada.","status":"PRESENT_IN_SUCCESSOR_SOURCE"}
% HMT_RESULT_END:N9-002:residence
% HMT_RESULT_BEGIN:N9-003:residence
% HMT_ATOMIC_RESULT {"material_state":"INTEGRADO_O_REMITIDO_SIN_PERDIDA","proof_strength":"EXACTO_INTERNO","provenance":"RESULTADO_RECUPERADO","reinforcement_state":"DEPLOYED_WITHOUT_NEW_RESULT_ID","residence":"`hmt/06b--06e`","result_id":"N9-003","statement":"Una salida visible conserva cinco genealogías regionales. \\(5=3_{++}+1_{--}+1_\\perp\\), \\(432+4\\cdot144=1008\\).","status":"PRESENT_IN_SUCCESSOR_SOURCE"}
% HMT_RESULT_END:N9-003:residence
% HMT_RESULT_BEGIN:N9-004:residence
% HMT_ATOMIC_RESULT {"material_state":"INTEGRADO_O_REMITIDO_SIN_PERDIDA","proof_strength":"EXACTO_INTERNO","provenance":"CERTIFICADO_NUEVO","reinforcement_state":"DEPLOYED_WITHOUT_NEW_RESULT_ID","residence":"`hmt/06e--06f`","result_id":"N9-004","statement":"Las 1.000 cifras son una restricción reproducible, no el techo ni por sí solas la prueba lógica del límite.","status":"PRESENT_IN_SUCCESSOR_SOURCE"}
% HMT_RESULT_END:N9-004:residence
% HMT_RESULT_BEGIN:HEL-001:residence
% HMT_ATOMIC_RESULT {"material_state":"INTEGRADO_O_REMITIDO_SIN_PERDIDA","proof_strength":"EXACTO_INTERNO","provenance":"RESULTADO_RECUPERADO","reinforcement_state":"DEPLOYED_WITHOUT_NEW_RESULT_ID","residence":"`hmt/03l`","result_id":"HEL-001","statement":"«Diez vuelve a uno» significa misma fase y memoria incrementada.","status":"PRESENT_IN_SUCCESSOR_SOURCE"}
% HMT_RESULT_END:HEL-001:residence
% HMT_RESULT_BEGIN:OPS-001:residence
% HMT_ATOMIC_RESULT {"material_state":"INTEGRADO_O_REMITIDO_SIN_PERDIDA","proof_strength":"EXACTO_INTERNO","provenance":"RESULTADO_RECUPERADO","reinforcement_state":"DEPLOYED_WITHOUT_NEW_RESULT_ID","residence":"`hmt/03k`","result_id":"OPS-001","statement":"Suma, producto, potencia, raíz, logaritmo, derivación e integración deben elevarse con memoria de fibra.","status":"PRESENT_IN_SUCCESSOR_SOURCE"}
% HMT_RESULT_END:OPS-001:residence
% HMT_RESULT_BEGIN:ZF-001:residence
% HMT_ATOMIC_RESULT {"material_state":"INTEGRADO_O_REMITIDO_SIN_PERDIDA","proof_strength":"EXACTO_INTERNO","provenance":"RESULTADO_RECUPERADO","reinforcement_state":"DEPLOYED_WITHOUT_NEW_RESULT_ID","residence":"nuevo bloque «Genealogía, extensionalidad y teoría de conjuntos» después de completar APP--TRIT--TPK y `03l`","result_id":"ZF-001","statement":"HMT no refuta ZFC: construye estructuras genealógicas codificables en ZFC y muestra qué información pierde la decategorización conjuntista.","status":"PRESENT_IN_SUCCESSOR_SOURCE"}
% HMT_RESULT_END:ZF-001:residence
% HMT_RESULT_BEGIN:ZF-002:residence
% HMT_ATOMIC_RESULT {"material_state":"INTEGRADO_O_REMITIDO_SIN_PERDIDA","proof_strength":"PROGRAMA_ABIERTO","provenance":"RESULTADO_NUEVO","reinforcement_state":"DEPLOYED_WITHOUT_NEW_RESULT_ID","residence":"bloque lógico, § Alcance axiomático, tras ZF-001","result_id":"ZF-002","statement":"No se afirmará que HMT sea ya un modelo interno completo de ZF, IZF o NBG: Separación, Reemplazo, Potencia plena e Infinito transfinito no se han verificado conjuntamente.","status":"PRESENT_IN_SUCCESSOR_SOURCE"}
% HMT_RESULT_END:ZF-002:residence
% HMT_RESULT_BEGIN:EXT-001:residence
% HMT_ATOMIC_RESULT {"material_state":"INTEGRADO_O_REMITIDO_SIN_PERDIDA","proof_strength":"EXACTO_INTERNO","provenance":"RESULTADO_RECUPERADO","reinforcement_state":"DEPLOYED_WITHOUT_NEW_RESULT_ID","residence":"bloque fundacional, § Extensionalidad tipada, tras DSC-001","result_id":"EXT-001","statement":"Dos secciones son idénticas si coinciden a toda profundidad; igualdad de valor es un cociente salvo lector fiel.","status":"PRESENT_IN_SUCCESSOR_SOURCE"}
% HMT_RESULT_END:EXT-001:residence
% HMT_RESULT_BEGIN:FND-001:residence
% HMT_ATOMIC_RESULT {"material_state":"INTEGRADO_O_REMITIDO_SIN_PERDIDA","proof_strength":"EXACTO_INTERNO","provenance":"RESULTADO_RECUPERADO","reinforcement_state":"DEPLOYED_WITHOUT_NEW_RESULT_ID","residence":"bloque lógico, § Fundación y coinducción, tras TOP-001","result_id":"FND-001","statement":"Los prefijos son bien fundados por rango natural; una sección infinita es coinductiva y no introduce ciclos de pertenencia.","status":"PRESENT_IN_SUCCESSOR_SOURCE"}
% HMT_RESULT_END:FND-001:residence
% HMT_RESULT_BEGIN:ORD-001:residence
% HMT_ATOMIC_RESULT {"material_state":"INTEGRADO_O_REMITIDO_SIN_PERDIDA","proof_strength":"EXACTO_INTERNO","provenance":"FORMALIZACION_NUEVA","reinforcement_state":"DEPLOYED_WITHOUT_NEW_RESULT_ID","residence":"subsección «Ordinal, cardinal y ruta»","result_id":"ORD-001","statement":"El ordinal de von Neumann se eleva conservando el estado recorrido; la proyección recupera el ordinal desnudo.","status":"PRESENT_IN_SUCCESSOR_SOURCE"}
% HMT_RESULT_END:ORD-001:residence
% HMT_RESULT_BEGIN:ORD-002:residence
% HMT_ATOMIC_RESULT {"material_state":"INTEGRADO_O_REMITIDO_SIN_PERDIDA","proof_strength":"PROGRAMA_ABIERTO","provenance":"FORMALIZACION_NUEVA","reinforcement_state":"DEPLOYED_WITHOUT_NEW_RESULT_ID","residence":"bloque lógico, § Jerarquía transfinita, tras ORD-001","result_id":"ORD-002","statement":"El corpus cierra profundidades finitas y secciones de longitud \\(\\omega\\); la extensión a todo ordinal es programa.","status":"PRESENT_IN_SUCCESSOR_SOURCE"}
% HMT_RESULT_END:ORD-002:residence
% HMT_RESULT_BEGIN:PWR-001:residence
% HMT_ATOMIC_RESULT {"material_state":"INTEGRADO_O_REMITIDO_SIN_PERDIDA","proof_strength":"EXACTO_INTERNO","provenance":"FORMALIZACION_NUEVA","reinforcement_state":"DEPLOYED_WITHOUT_NEW_RESULT_ID","residence":"subsección «Potencia genealógica»","result_id":"PWR-001","statement":"Los predicados decidibles a profundidad finita forman los clopen del límite.","status":"PRESENT_IN_SUCCESSOR_SOURCE"}
% HMT_RESULT_END:PWR-001:residence
% HMT_RESULT_BEGIN:PWR-002:residence
% HMT_ATOMIC_RESULT {"material_state":"INTEGRADO_O_REMITIDO_SIN_PERDIDA","proof_strength":"EXACTO_INTERNO","provenance":"FORMALIZACION_NUEVA","reinforcement_state":"DEPLOYED_WITHOUT_NEW_RESULT_ID","residence":"bloque lógico, § Potencia cerrada, tras PWR-001","result_id":"PWR-002","statement":"Familias compatibles de subconjuntos finitos reconstruyen cerrados, no todos los subconjuntos.","status":"PRESENT_IN_SUCCESSOR_SOURCE"}
% HMT_RESULT_END:PWR-002:residence
% HMT_RESULT_BEGIN:CHO-001:residence
% HMT_ATOMIC_RESULT {"material_state":"INTEGRADO_O_REMITIDO_SIN_PERDIDA","proof_strength":"EXACTO_INTERNO","provenance":"RESULTADO_RECUPERADO","reinforcement_state":"DEPLOYED_WITHOUT_NEW_RESULT_ID","residence":"subsección «Elección tipada»","result_id":"CHO-001","statement":"Se separan elección conjuntista, natural, algebraica y efectiva. Toda sección natural, algebraica o efectiva es en particular una sección conjuntista, pero naturalidad, algebraicidad y efectividad son requisitos adicionales generalmente incomparables entre sí.","status":"PRESENT_IN_SUCCESSOR_SOURCE"}
% HMT_RESULT_END:CHO-001:residence
% HMT_RESULT_BEGIN:CHO-002:residence
% HMT_ATOMIC_RESULT {"material_state":"INTEGRADO_O_REMITIDO_SIN_PERDIDA","proof_strength":"EXACTO_INTERNO","provenance":"RESULTADO_RECUPERADO","reinforcement_state":"DEPLOYED_WITHOUT_NEW_RESULT_ID","residence":"bloque lógico, § Elección tipada y obstrucción homomorfa, tras CHO-001","result_id":"CHO-002","statement":"Hay representantes conjuntistas pero no sección homomorfa; eso mide memoria estructural, no fracaso de AC.","status":"PRESENT_IN_SUCCESSOR_SOURCE"}
% HMT_RESULT_END:CHO-002:residence
% HMT_RESULT_BEGIN:CAT-001:residence
% HMT_ATOMIC_RESULT {"material_state":"INTEGRADO_O_REMITIDO_SIN_PERDIDA","proof_strength":"EXACTO_INTERNO","provenance":"RESULTADO_RECUPERADO","reinforcement_state":"DEPLOYED_WITHOUT_NEW_RESULT_ID","residence":"cierre categorial del nuevo bloque","result_id":"CAT-001","statement":"Objetos como torres tipadas; morfismos conmutan restricciones; las secciones son el funtor límite.","status":"PRESENT_IN_SUCCESSOR_SOURCE"}
% HMT_RESULT_END:CAT-001:residence
% HMT_RESULT_BEGIN:TOP-001:residence
% HMT_ATOMIC_RESULT {"material_state":"INTEGRADO_O_REMITIDO_SIN_PERDIDA","proof_strength":"EXACTO_INTERNO","provenance":"FORMALIZACION_NUEVA","reinforcement_state":"DEPLOYED_WITHOUT_NEW_RESULT_ID","residence":"nota técnica subordinada a CAT-001","result_id":"TOP-001","statement":"La categoría pequeña de caminos TPK admite una envolvente de presheaves; no se le atribuyen elección o lógica clásica sin prueba.","status":"PRESENT_IN_SUCCESSOR_SOURCE"}
% HMT_RESULT_END:TOP-001:residence
% HMT_RESULT_BEGIN:HIL-001:residence
% HMT_ATOMIC_RESULT {"material_state":"INTEGRADO_O_REMITIDO_SIN_PERDIDA","proof_strength":"EXACTO_INTERNO","provenance":"RESULTADO_RECUPERADO","reinforcement_state":"DEPLOYED_WITHOUT_NEW_RESULT_ID","residence":"`hmt/03l` y `md/04ab`","result_id":"HIL-001","statement":"Cada profundidad posee una realización de Hilbert y una álgebra matricial.","status":"PRESENT_IN_SUCCESSOR_SOURCE"}
% HMT_RESULT_END:HIL-001:residence
% HMT_RESULT_BEGIN:UHF-001:residence
% HMT_ATOMIC_RESULT {"material_state":"INTEGRADO_O_REMITIDO_SIN_PERDIDA","proof_strength":"EXACTO_INTERNO","provenance":"FORMALIZACION_NUEVA","reinforcement_state":"DEPLOYED_WITHOUT_NEW_RESULT_ID","residence":"nuevo bloque «Torre nonádica de observables» tras `hmt/03l`","result_id":"UHF-001","statement":"Conservar simultáneamente las nueve fibras preserva unidad y traza y genera la rama UHF--II_1; una esquina marcada no lo hace y su corredor de tipo I_infty se conserva separadamente en IINF-001.","status":"PRESENT_IN_SUCCESSOR_SOURCE"}
% HMT_RESULT_END:UHF-001:residence
% HMT_RESULT_BEGIN:II1-001:residence
% HMT_ATOMIC_RESULT {"material_state":"INTEGRADO_O_REMITIDO_SIN_PERDIDA","proof_strength":"EXACTO_INTERNO","provenance":"FORMALIZACION_NUEVA","reinforcement_state":"DEPLOYED_WITHOUT_NEW_RESULT_ID","residence":"torre operatorial, § Cierre tracial hiperinfinito II_1, tras IINF-001","result_id":"II1-001","statement":"El límite unital de todas las puertas produce UHF nonádico y su cierre tracial es el factor hiperinfinito \\(II_1\\).","status":"PRESENT_IN_SUCCESSOR_SOURCE"}
% HMT_RESULT_END:II1-001:residence
% HMT_RESULT_BEGIN:DIM-001:residence
% HMT_ATOMIC_RESULT {"material_state":"INTEGRADO_O_REMITIDO_SIN_PERDIDA","proof_strength":"EXACTO_INTERNO","provenance":"ARQUITECTURA_AUTORAL_PREEXISTENTE","reinforcement_state":"DEPLOYED_WITHOUT_NEW_RESULT_ID","residence":"torre operatorial, § Dimensión continua de proyectores, tras II1-001","result_id":"DIM-001","statement":"Para todo t∈[0,1], m_d=floor(9^d t) permite elegir proyectores anidados cuyo límite fuerte P_t satisface tau(P_t)=t; en el factor finito, P~Q iff tau(P)=tau(Q).","status":"PRESENT_IN_SUCCESSOR_SOURCE"}
% HMT_RESULT_END:DIM-001:residence
% HMT_RESULT_BEGIN:CLK-001:residence
% HMT_ATOMIC_RESULT {"material_state":"INTEGRADO_O_REMITIDO_SIN_PERDIDA","proof_strength":"EXACTO_INTERNO","provenance":"RESULTADO_RECUPERADO","reinforcement_state":"DEPLOYED_WITHOUT_NEW_RESULT_ID","residence":"subsección «Realización por reloj profinito»","result_id":"CLK-001","statement":"Una segunda ruta llega abstractamente al mismo factor mediante producto cruzado.","status":"PRESENT_IN_SUCCESSOR_SOURCE"}
% HMT_RESULT_END:CLK-001:residence
% HMT_RESULT_BEGIN:INT-001:residence
% HMT_ATOMIC_RESULT {"material_state":"INTEGRADO_O_REMITIDO_SIN_PERDIDA","proof_strength":"PROGRAMA_ABIERTO","provenance":"FORMALIZACION_NUEVA","reinforcement_state":"DEPLOYED_WITHOUT_NEW_RESULT_ID","residence":"torre operatorial, § Entrelazador APP--reloj, tras CLK-001","result_id":"INT-001","statement":"La isomorfía abstracta de factores no identifica genealogías ni los antecedentes \\(C^*\\).","status":"PRESENT_IN_SUCCESSOR_SOURCE"}
% HMT_RESULT_END:INT-001:residence
% HMT_RESULT_BEGIN:QLG-001:residence
% HMT_ATOMIC_RESULT {"material_state":"INTEGRADO_O_REMITIDO_SIN_PERDIDA","proof_strength":"EXACTO_INTERNO","provenance":"RESULTADO_RECUPERADO","reinforcement_state":"DEPLOYED_WITHOUT_NEW_RESULT_ID","residence":"`md/04ab` después de PVM/contextos","result_id":"QLG-001","statement":"\\(\\operatorname{Proj}(R_{\\rm hyp})\\) es un retículo ortomodular completo, continuo y no distributivo; cada contexto conmutativo es booleano.","status":"PRESENT_IN_SUCCESSOR_SOURCE"}
% HMT_RESULT_END:QLG-001:residence
% HMT_RESULT_BEGIN:ERG-001:residence
% HMT_ATOMIC_RESULT {"material_state":"INTEGRADO_O_REMITIDO_SIN_PERDIDA","proof_strength":"EXACTO_INTERNO","provenance":"RESULTADO_RECUPERADO","reinforcement_state":"DEPLOYED_WITHOUT_NEW_RESULT_ID","residence":"capítulo de torre/reloj","result_id":"ERG-001","statement":"Los promedios temporales proyectan sobre invariantes; para el reloj ergódico, sobre la media de Haar.","status":"PRESENT_IN_SUCCESSOR_SOURCE"}
% HMT_RESULT_END:ERG-001:residence
% HMT_RESULT_BEGIN:DEN-001:residence
% HMT_ATOMIC_RESULT {"material_state":"INTEGRADO_O_REMITIDO_SIN_PERDIDA","proof_strength":"EXACTO_INTERNO","provenance":"RESULTADO_RECUPERADO","reinforcement_state":"DEPLOYED_WITHOUT_NEW_RESULT_ID","residence":"`md/04ab` y torre \\(II_1\\)","result_id":"DEN-001","statement":"La incertidumbre sobre rutas se representa por estado positivo; no obliga a borrar la genealogía subyacente.","status":"PRESENT_IN_SUCCESSOR_SOURCE"}
% HMT_RESULT_END:DEN-001:residence
% HMT_RESULT_BEGIN:LAN-001:residence
% HMT_ATOMIC_RESULT {"material_state":"INTEGRADO_O_REMITIDO_SIN_PERDIDA","proof_strength":"EXACTO_INTERNO","provenance":"RESULTADO_RECUPERADO","reinforcement_state":"DEPLOYED_WITHOUT_NEW_RESULT_ID","residence":"`md/04g`","result_id":"LAN-001","statement":"Una actualización biyectiva U del estado completo satisface H(X|U(X))=0 y, para el término ideal de borrado lógico, Q_L^{min}[U]=0.","status":"PRESENT_IN_SUCCESSOR_SOURCE"}
% HMT_RESULT_END:LAN-001:residence
% HMT_RESULT_BEGIN:LAN-002:residence
% HMT_ATOMIC_RESULT {"material_state":"INTEGRADO_O_REMITIDO_SIN_PERDIDA","proof_strength":"EXACTO_INTERNO","provenance":"FORMALIZACION_NUEVA","reinforcement_state":"DEPLOYED_WITHOUT_NEW_RESULT_ID","residence":"`md/04g`, después del caso finito","result_id":"LAN-002","statement":"Automorfismos traza-preservantes conservan entropía; una esperanza condicional incrementa la entropía exactamente por entropía relativa.","status":"PRESENT_IN_SUCCESSOR_SOURCE"}
% HMT_RESULT_END:LAN-002:residence
% HMT_RESULT_BEGIN:GDL-001:residence
% HMT_ATOMIC_RESULT {"material_state":"INTEGRADO_O_REMITIDO_SIN_PERDIDA","proof_strength":"EXACTO_INTERNO","provenance":"RESULTADO_RECUPERADO","reinforcement_state":"DEPLOYED_WITHOUT_NEW_RESULT_ID","residence":"nuevo bloque «Límite efectivo de la clausura»","result_id":"GDL-001","statement":"Cada nivel puede calcularse mientras la unicidad de la fibra global codifica HALT.","status":"PRESENT_IN_SUCCESSOR_SOURCE"}
% HMT_RESULT_END:GDL-001:residence
% HMT_RESULT_BEGIN:GDL-002:residence
% HMT_ATOMIC_RESULT {"material_state":"INTEGRADO_O_REMITIDO_SIN_PERDIDA","proof_strength":"EXACTO_INTERNO","provenance":"RESULTADO_RECUPERADO","reinforcement_state":"DEPLOYED_WITHOUT_NEW_RESULT_ID","residence":"bloque metalógico, § Incompletitud semántica Gödel--Turing, tras GDL-001","result_id":"GDL-002","statement":"Toda teoría recursivamente axiomatizable y correcta para la familia \\(U_e\\) deja alguna sentencia sin decidir; no implica inconsistencia de ZFC.","status":"PRESENT_IN_SUCCESSOR_SOURCE"}
% HMT_RESULT_END:GDL-002:residence
% HMT_RESULT_BEGIN:IND-001:residence
% HMT_ATOMIC_RESULT {"material_state":"INTEGRADO_O_REMITIDO_SIN_PERDIDA","proof_strength":"EXACTO_INTERNO","provenance":"FORMALIZACION_NUEVA","reinforcement_state":"DEPLOYED_WITHOUT_NEW_RESULT_ID","residence":"bloque metalógico, § Frontera Gödel--Cohen, tras GDL-002","result_id":"IND-001","statement":"Bajo las hipótesis de consistencia correspondientes, CH es independiente de ZFC y AC es independiente de ZF; HMT enriquece objetos, no decide por sí sola esos enunciados.","status":"PRESENT_IN_SUCCESSOR_SOURCE"}
% HMT_RESULT_END:IND-001:residence
% HMT_RESULT_BEGIN:FOR-001:residence
% HMT_ATOMIC_RESULT {"material_state":"INTEGRADO_O_REMITIDO_SIN_PERDIDA","proof_strength":"EXACTO_INTERNO","provenance":"FORMALIZACION_NUEVA","reinforcement_state":"DEPLOYED_WITHOUT_NEW_RESULT_ID","residence":"subsección de filtros de prefijos","result_id":"FOR-001","statement":"Una sección determina un filtro que atraviesa todos los densos de profundidad; eso no basta para genericidad sobre un modelo.","status":"PRESENT_IN_SUCCESSOR_SOURCE"}
% HMT_RESULT_END:FOR-001:residence
% HMT_RESULT_BEGIN:FOR-002:residence
% HMT_ATOMIC_RESULT {"material_state":"INTEGRADO_O_REMITIDO_SIN_PERDIDA","proof_strength":"INTERFAZ_FISICA","provenance":"FORMALIZACION_NUEVA","reinforcement_state":"DEPLOYED_WITHOUT_NEW_RESULT_ID","residence":"bloque metalógico, § Forcing cilíndrico y filtro genérico, tras FOR-001","result_id":"FOR-002","statement":"Árbol perfecto, numerable y sin átomos da forcing cilíndrico Cohen-equivalente; rama única da forcing trivial.","status":"PRESENT_IN_SUCCESSOR_SOURCE"}
% HMT_RESULT_END:FOR-002:residence
% HMT_RESULT_BEGIN:BOR-001:residence
% HMT_ATOMIC_RESULT {"material_state":"INTEGRADO_O_REMITIDO_SIN_PERDIDA","proof_strength":"EXACTO_INTERNO","provenance":"FORMALIZACION_NUEVA","reinforcement_state":"DEPLOYED_WITHOUT_NEW_RESULT_ID","residence":"subsección final del bloque","result_id":"BOR-001","statement":"Los cilindros generan una clase Borel estable por complemento, unión numerable e imágenes inversas continuas; sus juegos usan determinación boreliana clásica.","status":"PRESENT_IN_SUCCESSOR_SOURCE"}
% HMT_RESULT_END:BOR-001:residence
% HMT_RESULT_BEGIN:MIN-001:residence
% HMT_ATOMIC_RESULT {"material_state":"INTEGRADO_O_REMITIDO_SIN_PERDIDA","proof_strength":"INTERFAZ_FISICA","provenance":"FORMALIZACION_NUEVA","reinforcement_state":"DEPLOYED_WITHOUT_NEW_RESULT_ID","residence":"bloque metalógico, § Minimax finito y proyectivo, tras BOR-001","result_id":"MIN-001","statement":"En cada profundidad finita, von Neumann da valor mixto; un límite global requiere compactitud, continuidad y consistencia proyectiva, y no selecciona por sí solo rama pura.","status":"PRESENT_IN_SUCCESSOR_SOURCE"}
% HMT_RESULT_END:MIN-001:residence
% HMT_RESULT_BEGIN:III-001:residence
% HMT_ATOMIC_RESULT {"material_state":"INTEGRADO_O_REMITIDO_SIN_PERDIDA","proof_strength":"PROGRAMA_ABIERTO","provenance":"RESULTADO_NUEVO","reinforcement_state":"DEPLOYED_WITHOUT_NEW_RESULT_ID","residence":"frontera tras \\(II_1\\)","result_id":"III-001","statement":"La construcción actual produce naturalmente \\(II_1\\); tipo \\(III\\) exige GNS no tracial/KMS o cociclo de Radon--Nikodym no trivial.","status":"PRESENT_IN_SUCCESSOR_SOURCE"}
% HMT_RESULT_END:III-001:residence
% HMT_RESULT_BEGIN:TRI-001:residence
% HMT_ATOMIC_RESULT {"material_state":"INTEGRADO_O_REMITIDO_SIN_PERDIDA","proof_strength":"INTERFAZ_FISICA","provenance":"FORMALIZACION_NUEVA","reinforcement_state":"DEPLOYED_WITHOUT_NEW_RESULT_ID","residence":"después de exponer íntegramente la doble proyección y completar el corredor Paley--Witt--Golay--Leech--\\(V^\\natural\\)--Monstruo--Moonshine","result_id":"TRI-001","statement":"Punto, dimensión y genealogía requieren tres capas coordinadas.","status":"PRESENT_IN_SUCCESSOR_SOURCE"}
% HMT_RESULT_END:TRI-001:residence
% HMT_RESULT_BEGIN:EXC-001:residence
% HMT_ATOMIC_RESULT {"material_state":"INTEGRADO_O_REMITIDO_SIN_PERDIDA","proof_strength":"DEMOSTRADO_CON_ESTRUCTURA_DE_PARTIDA_EXPLICITA","provenance":"RESULTADO_RECUPERADO","reinforcement_state":"DEPLOYED_WITHOUT_NEW_RESULT_ID","residence":"`hmt/09g`, `15`","result_id":"EXC-001","statement":"El mismo estado superviviente alimenta lectura real e incidencia excepcional hasta Leech--\\(V^\\natural\\)--Monstruo/Moonshine; el Monstruo no actúa literalmente sobre dígitos.","status":"PRESENT_IN_SUCCESSOR_SOURCE"}
% HMT_RESULT_END:EXC-001:residence
% HMT_RESULT_BEGIN:MD-001:residence
% HMT_ATOMIC_RESULT {"material_state":"INTEGRADO_O_REMITIDO_SIN_PERDIDA","proof_strength":"PROGRAMA_ABIERTO","provenance":"RESULTADO_RECUPERADO","reinforcement_state":"DEPLOYED_WITHOUT_NEW_RESULT_ID","residence":"`md/00b`, antes de Planck/electrón","result_id":"MD-001","statement":"Una sección compatible no es todavía una partícula: la ruta es su historia observable y la partícula una sección persistente enriquecida antes del morfismo de realización física.","status":"PRESENT_IN_SUCCESSOR_SOURCE"}
% HMT_RESULT_END:MD-001:residence
% HMT_RESULT_BEGIN:PSC-001:residence
% HMT_ATOMIC_RESULT {"material_state":"INTEGRADO_O_REMITIDO_SIN_PERDIDA","proof_strength":"EXACTO_INTERNO","provenance":"ARQUITECTURA_AUTORAL_PREEXISTENTE","reinforcement_state":"DEPLOYED_WITHOUT_NEW_RESULT_ID","residence":"apéndice crítico, fuera del tronco probatorio","result_id":"PSC-001","statement":"Las ramas AD/MM, CH y «cierre fuerte» conservan prioridad intuitiva, no fuerza probatoria vigente.","status":"PRESENT_IN_SUCCESSOR_SOURCE"}
% HMT_RESULT_END:PSC-001:residence
% HMT_RESULT_BEGIN:COV-001:residence
% HMT_ATOMIC_RESULT {"material_state":"INTEGRADO_O_REMITIDO_SIN_PERDIDA","proof_strength":"EXACTO_INTERNO","provenance":"RESULTADO_RECUPERADO","reinforcement_state":"DEPLOYED_WITHOUT_NEW_RESULT_ID","residence":"gate de cobertura final","result_id":"COV-001","statement":"Constantes, electrón, masas, Barbero, CKM/CP, gravedad y corredor excepcional se preservan por remisión, sin versión degradada.","status":"PRESENT_IN_SUCCESSOR_SOURCE"}
% HMT_RESULT_END:COV-001:residence
% HMT_RESULT_BEGIN:IINF-001:residence
% HMT_ATOMIC_RESULT {"material_state":"INTEGRADO_O_REMITIDO_SIN_PERDIDA","proof_strength":"EXACTO_INTERNO","provenance":"RESULTADO_RECUPERADO","reinforcement_state":"DEPLOYED_WITHOUT_NEW_RESULT_ID","residence":"torre operatorial, § Rama marcada de tipo I_infty, tras UHF-001","result_id":"IINF-001","statement":"Las esquinas j_{d,j}(a)=a⊗p_j, implementadas por V_{d,j}aV_{d,j}^*, generan en el límite los operadores de rango finito; su clausura normativa es K(H_infty) y su clausura débil B(H_infty), factor de tipo I_infty.","status":"PRESENT_IN_SUCCESSOR_SOURCE"}
% HMT_RESULT_END:IINF-001:residence
% HMT_RESULT_BEGIN:DEN-003:residence
% HMT_ATOMIC_RESULT {"material_state":"INTEGRADO_O_REMITIDO_SIN_PERDIDA","proof_strength":"EXACTO_INTERNO","provenance":"RESULTADO_RECUPERADO","reinforcement_state":"DEPLOYED_WITHOUT_NEW_RESULT_ID","residence":"torre \\(II_1\\), inmediatamente después de DEN-002","result_id":"DEN-003","statement":"Si h=9^d rho, entonces log h=d log9 I+log rho y S_tau(h)=S_vN(rho)-d log9; bajo rho_{d+1}=rho_d⊗I9/9, S_vN aumenta en log9 y S_tau permanece invariante.","status":"PRESENT_IN_SUCCESSOR_SOURCE"}
% HMT_RESULT_END:DEN-003:residence
% HMT_RESULT_BEGIN:CARD-001:residence
% HMT_ATOMIC_RESULT {"material_state":"INTEGRADO_O_REMITIDO_SIN_PERDIDA","proof_strength":"EXACTO_INTERNO","provenance":"RESULTADO_RECUPERADO","reinforcement_state":"DEPLOYED_WITHOUT_NEW_RESULT_ID","residence":"bloque lógico, § Cardinalidad como funtor no fiel, tras ORD-002","result_id":"CARD-001","statement":"Contar es decategorizar: el perfil (|X_m|)_m no determina restricciones, acarreo, orientación o memoria.","status":"PRESENT_IN_SUCCESSOR_SOURCE"}
% HMT_RESULT_END:CARD-001:residence
% HMT_RESULT_BEGIN:PWR-003:residence
% HMT_ATOMIC_RESULT {"material_state":"INTEGRADO_O_REMITIDO_SIN_PERDIDA","proof_strength":"EXACTO_INTERNO","provenance":"FORMALIZACION_NUEVA","reinforcement_state":"DEPLOYED_WITHOUT_NEW_RESULT_ID","residence":"bloque lógico, § Tres escalas de potencia, tras PWR-002","result_id":"PWR-003","statement":"En un Cantor HMT, clopen, cerrados y potencia plena tienen tamaños distintos.","status":"PRESENT_IN_SUCCESSOR_SOURCE"}
% HMT_RESULT_END:PWR-003:residence
% HMT_RESULT_BEGIN:WEYL-001:residence
% HMT_ATOMIC_RESULT {"material_state":"INTEGRADO_O_REMITIDO_SIN_PERDIDA","proof_strength":"EXACTO_INTERNO","provenance":"RESULTADO_RECUPERADO","reinforcement_state":"DEPLOYED_WITHOUT_NEW_RESULT_ID","residence":"`hmt/14_numero_heisenberg_y_d108.tex`, `md/04ab`","result_id":"WEYL-001","statement":"Fase y desplazamiento forman un par no conmutativo; fijado el carácter central primitivo, la representación irreducible es única hasta unitario.","status":"PRESENT_IN_SUCCESSOR_SOURCE"}
% HMT_RESULT_END:WEYL-001:residence
% HMT_RESULT_BEGIN:DSC-001:residence
% HMT_ATOMIC_RESULT {"material_state":"INTEGRADO_O_REMITIDO_SIN_PERDIDA","proof_strength":"EXACTO_INTERNO","provenance":"RESULTADO_RECUPERADO","reinforcement_state":"DEPLOYED_WITHOUT_NEW_RESULT_ID","residence":"bloque fundacional, § Presentación y descenso por fibras, tras INF-003","result_id":"DSC-001","statement":"Existe un único operador descendido O_bar:V→Y con O=O_bar q si y sólo si O es constante en cada fibra de q.","status":"PRESENT_IN_SUCCESSOR_SOURCE"}
% HMT_RESULT_END:DSC-001:residence
% HMT_RESULT_BEGIN:REC-001:residence
% HMT_ATOMIC_RESULT {"material_state":"INTEGRADO_O_REMITIDO_SIN_PERDIDA","proof_strength":"EXACTO_INTERNO","provenance":"RESULTADO_RECUPERADO","reinforcement_state":"DEPLOYED_WITHOUT_NEW_RESULT_ID","residence":"`md/04g`, antes del teorema de Landauer lógico","result_id":"REC-001","statement":"Si q_hat=(q,M) es inyectivo sobre supp(X), entonces H(X|q(X),M(X))=0; H(X|q(X)) mide la fibra olvidada por el lector visible.","status":"PRESENT_IN_SUCCESSOR_SOURCE"}
% HMT_RESULT_END:REC-001:residence
% HMT_RESULT_BEGIN:LAN-003:residence
% HMT_ATOMIC_RESULT {"material_state":"INTEGRADO_O_REMITIDO_SIN_PERDIDA","proof_strength":"EXACTO_INTERNO","provenance":"RESULTADO_RECUPERADO","reinforcement_state":"DEPLOYED_WITHOUT_NEW_RESULT_ID","residence":"`md/04g`, después de LAN-001","result_id":"LAN-003","statement":"Para un reset no inyectivo de tres estados equiprobables, H(X)=log 3 y el término ideal satisface Q_reset≥k_B Θ log 3.","status":"PRESENT_IN_SUCCESSOR_SOURCE"}
% HMT_RESULT_END:LAN-003:residence
% HMT_RESULT_BEGIN:FIN-001:residence
% HMT_ATOMIC_RESULT {"material_state":"INTEGRADO_O_REMITIDO_SIN_PERDIDA","proof_strength":"EXACTO_INTERNO","provenance":"FORMALIZACION_NUEVA","reinforcement_state":"DEPLOYED_WITHOUT_NEW_RESULT_ID","residence":"nuevo bloque «Límite efectivo de la clausura», antes de GDL-001","result_id":"FIN-001","statement":"Fijado N, con dominios finitos y observables atómicos computables, toda fórmula del lenguaje finito L_N es semánticamente decidible.","status":"PRESENT_IN_SUCCESSOR_SOURCE"}
% HMT_RESULT_END:FIN-001:residence
% HMT_RESULT_BEGIN:DEN-002:residence
% HMT_ATOMIC_RESULT {"material_state":"INTEGRADO_O_REMITIDO_SIN_PERDIDA","proof_strength":"EXACTO_INTERNO","provenance":"RESULTADO_RECUPERADO","reinforcement_state":"DEPLOYED_WITHOUT_NEW_RESULT_ID","residence":"torre \\(II_1\\), después de DEN-001","result_id":"DEN-002","statement":"Para h=9^d rho, S_tau(h)=S_vN(rho)-d log9; la extensión rho_{d+1}=rho_d⊗I9/9 aumenta S_vN en log9 y deja invariante S_tau.","status":"PRESENT_IN_SUCCESSOR_SOURCE"}
% HMT_RESULT_END:DEN-002:residence
% HMT_RESULT_BEGIN:UNI-001:residence
% HMT_ATOMIC_RESULT {"material_state":"INTEGRADO_O_REMITIDO_SIN_PERDIDA","proof_strength":"EXACTO_INTERNO","provenance":"RESULTADO_RECUPERADO","reinforcement_state":"DEPLOYED_WITHOUT_NEW_RESULT_ID","residence":"cierre APP--EMRH","result_id":"UNI-001","statement":"La frontera distingue 999=729+243+27 de 1000=729+243+27+1=729+271; el acarreo c_{m+1}=1 conserva la unidad visible y registra una vuelta adicional.","status":"PRESENT_IN_SUCCESSOR_SOURCE"}
% HMT_RESULT_END:UNI-001:residence
% HMT_RESULT_BEGIN:UNI-002:residence
% HMT_ATOMIC_RESULT {"material_state":"INTEGRADO_O_REMITIDO_SIN_PERDIDA","proof_strength":"PROGRAMA_ABIERTO","provenance":"ARQUITECTURA_AUTORAL_PREEXISTENTE","reinforcement_state":"DEPLOYED_WITHOUT_NEW_RESULT_ID","residence":"transición APP--EMRH hacia monodromía e inclusión unital nonádica","result_id":"UNI-002","statement":"La compatibilidad entre c_{m+1}=1, la monodromía 9→1 y a↦a⊗I9 es una obligación de entrelazamiento: debe construirse un diagrama que transporte la misma unidad con memoria entre los tres dominios.","status":"PRESENT_IN_SUCCESSOR_SOURCE"}
% HMT_RESULT_END:UNI-002:residence
% HMT_RESULT_BEGIN:P01:residence
% HMT_ATOMIC_RESULT {"material_state":"INTEGRADO_EN_GRAFO_ACTIVO","proof_strength":"EXACTO_INTERNO","provenance":"FORMALIZACION_NUEVA","reinforcement_state":"DEPLOYED_WITHOUT_NEW_RESULT_ID","residence":"APP vigente","result_id":"P01","statement":"APP tiene soporte $X_9$, 81 vértices y todos los caminos finitos.","status":"PRESENT_IN_SUCCESSOR_SOURCE"}
% HMT_RESULT_END:P01:residence
% HMT_RESULT_BEGIN:P02:residence
% HMT_ATOMIC_RESULT {"material_state":"INTEGRADO_EN_GRAFO_ACTIVO","proof_strength":"EXACTO_INTERNO","provenance":"FORMALIZACION_NUEVA","reinforcement_state":"DEPLOYED_WITHOUT_NEW_RESULT_ID","residence":"APP vigente","result_id":"P02","statement":"$\\Sigma$ es latina; $\\Pi$ sólo permuta en unidades.","status":"PRESENT_IN_SUCCESSOR_SOURCE"}
% HMT_RESULT_END:P02:residence
% HMT_RESULT_BEGIN:P03:residence
% HMT_ATOMIC_RESULT {"material_state":"INTEGRADO_EN_GRAFO_ACTIVO","proof_strength":"EXACTO_INTERNO","provenance":"FORMALIZACION_NUEVA","reinforcement_state":"DEPLOYED_WITHOUT_NEW_RESULT_ID","residence":"APP/TRIT","result_id":"P03","statement":"La elevación residuo--cociente recompone $1215=54+9\\cdot129$.","status":"PRESENT_IN_SUCCESSOR_SOURCE"}
% HMT_RESULT_END:P03:residence
% HMT_RESULT_BEGIN:P04:residence
% HMT_ATOMIC_RESULT {"material_state":"INTEGRADO_EN_GRAFO_ACTIVO","proof_strength":"EXACTO_INTERNO","provenance":"FORMALIZACION_NUEVA","reinforcement_state":"DEPLOYED_WITHOUT_NEW_RESULT_ID","residence":"`hmt/03f`","result_id":"P04","statement":"$B_c=9P_++5P_-$ y el salto es cuatro.","status":"PRESENT_IN_SUCCESSOR_SOURCE"}
% HMT_RESULT_END:P04:residence
% HMT_RESULT_BEGIN:P05:residence
% HMT_ATOMIC_RESULT {"material_state":"INTEGRADO_EN_GRAFO_ACTIVO","proof_strength":"EXACTO_INTERNO","provenance":"FORMALIZACION_NUEVA","reinforcement_state":"DEPLOYED_WITHOUT_NEW_RESULT_ID","residence":"TRIT","result_id":"P05","statement":"TRIT visible no determina el estado levantado.","status":"PRESENT_IN_SUCCESSOR_SOURCE"}
% HMT_RESULT_END:P05:residence
% HMT_RESULT_BEGIN:P06:residence
% HMT_ATOMIC_RESULT {"material_state":"INTEGRADO_EN_GRAFO_ACTIVO","proof_strength":"EXACTO_INTERNO","provenance":"FORMALIZACION_NUEVA","reinforcement_state":"DEPLOYED_WITHOUT_NEW_RESULT_ID","residence":"TPK","result_id":"P06","statement":"TPK compone rutas con hoja, acarreo y memoria.","status":"PRESENT_IN_SUCCESSOR_SOURCE"}
% HMT_RESULT_END:P06:residence
% HMT_RESULT_BEGIN:P07:residence
% HMT_ATOMIC_RESULT {"material_state":"INTEGRADO_EN_GRAFO_ACTIVO","proof_strength":"EXACTO_INTERNO","provenance":"RESULTADO_RECUPERADO","reinforcement_state":"DEPLOYED_WITHOUT_NEW_RESULT_ID","residence":"`md/04ab`","result_id":"P07","statement":"Los operadores de traslación X_u y fase Z_v satisfacen Z_vX_u=omega^{u·v}X_uZ_v, generan M_{9^d}(C) y, fijado el carácter central primitivo, dan la representación irreducible finita única salvo equivalencia unitaria.","status":"PRESENT_IN_SUCCESSOR_SOURCE"}
% HMT_RESULT_END:P07:residence
% HMT_RESULT_BEGIN:P08:residence
% HMT_ATOMIC_RESULT {"material_state":"INTEGRADO_EN_GRAFO_ACTIVO","proof_strength":"EXACTO_INTERNO","provenance":"RESULTADO_RECUPERADO","reinforcement_state":"DEPLOYED_WITHOUT_NEW_RESULT_ID","residence":"bloque UHF--\\(II_1\\)","result_id":"P08","statement":"El límite inductivo de a↦a⊗I_9 es UHF(9^infty), canónicamente isomorfa a UHF(3^infty); la representación GNS de su traza única cierra en el factor hiperinfinito II_1.","status":"PRESENT_IN_SUCCESSOR_SOURCE"}
% HMT_RESULT_END:P08:residence
% HMT_RESULT_BEGIN:P09:residence
% HMT_ATOMIC_RESULT {"material_state":"INTEGRADO_EN_GRAFO_ACTIVO","proof_strength":"EXACTO_INTERNO","provenance":"FORMALIZACION_NUEVA","reinforcement_state":"DEPLOYED_WITHOUT_NEW_RESULT_ID","residence":"apertura del corredor matricial","result_id":"P09","statement":"Un QMS es una cuadrícula positiva con sumas $I_s$.","status":"PRESENT_IN_SUCCESSOR_SOURCE"}
% HMT_RESULT_END:P09:residence
% HMT_RESULT_BEGIN:P10:residence
% HMT_ATOMIC_RESULT {"material_state":"INTEGRADO_EN_GRAFO_ACTIVO","proof_strength":"EXACTO_INTERNO","provenance":"RESULTADO_RECUPERADO","reinforcement_state":"DEPLOYED_WITHOUT_NEW_RESULT_ID","residence":"corredor matricial, § Semiclasicidad como factorización conmutativa, tras P09/QPM-001","result_id":"P10","statement":"A es semiclásico si y sólo si existe un mapa positivo unital Phi:C(S_n)→M_s tal que A_ij=Phi(f_ij), donde f_ij(sigma)=1_{sigma(i)=j}; por ser conmutativo el dominio, Phi es completamente positivo.","status":"PRESENT_IN_SUCCESSOR_SOURCE"}
% HMT_RESULT_END:P10:residence
% HMT_RESULT_BEGIN:P11:residence
% HMT_ATOMIC_RESULT {"material_state":"INTEGRADO_EN_GRAFO_ACTIVO","proof_strength":"EXACTO_INTERNO","provenance":"RESULTADO_RECUPERADO","reinforcement_state":"DEPLOYED_WITHOUT_NEW_RESULT_ID","residence":"corredor matricial, § QLS fáciles y semiclásicos, tras QLS-001","result_id":"P11","statement":"Un cuadrado latino cuántico es semiclásico si y sólo si es fácil, es decir, procede de un cuadrado latino clásico y una base ortonormal fija.","status":"PRESENT_IN_SUCCESSOR_SOURCE"}
% HMT_RESULT_END:P11:residence
% HMT_RESULT_BEGIN:P12:residence
% HMT_ATOMIC_RESULT {"material_state":"INTEGRADO_EN_GRAFO_ACTIVO","proof_strength":"EXACTO_INTERNO","provenance":"FORMALIZACION_NUEVA","reinforcement_state":"DEPLOYED_WITHOUT_NEW_RESULT_ID","residence":"capítulo qutrit","result_id":"P12","statement":"Cuatro contextos HMT producen cuatro QLS fáciles.","status":"PRESENT_IN_SUCCESSOR_SOURCE"}
% HMT_RESULT_END:P12:residence
% HMT_RESULT_BEGIN:P13:residence
% HMT_ATOMIC_RESULT {"material_state":"INTEGRADO_EN_GRAFO_ACTIVO","proof_strength":"EXACTO_INTERNO","provenance":"RESULTADO_NUEVO","reinforcement_state":"DEPLOYED_WITHOUT_NEW_RESULT_ID","residence":"cuadrados operatoriales, § No-go de QLS cíclicos, tras P12","result_id":"P13","statement":"Dos filas MUB no pueden apilarse en una QLS.","status":"PRESENT_IN_SUCCESSOR_SOURCE"}
% HMT_RESULT_END:P13:residence
% HMT_RESULT_BEGIN:P14:residence
% HMT_ATOMIC_RESULT {"material_state":"INTEGRADO_EN_GRAFO_ACTIVO","proof_strength":"EXACTO_INTERNO","provenance":"FORMALIZACION_NUEVA","reinforcement_state":"DEPLOYED_WITHOUT_NEW_RESULT_ID","residence":"tras TPK y torre UHF","result_id":"P14","statement":"Para n fijo, profundidad genealógica m y nivel matricial s definen una bigraduación; si n varía, la familia es trigraduada por (m,n,s).","status":"PRESENT_IN_SUCCESSOR_SOURCE"}
% HMT_RESULT_END:P14:residence
% HMT_RESULT_BEGIN:P15:residence
% HMT_ATOMIC_RESULT {"material_state":"INTEGRADO_EN_GRAFO_ACTIVO","proof_strength":"EXACTO_INTERNO","provenance":"FORMALIZACION_NUEVA","reinforcement_state":"DEPLOYED_WITHOUT_NEW_RESULT_ID","residence":"sistema inverso UCP, § Compresiones y truncación, tras P14","result_id":"P15","statement":"Si Phi_r∈K_{m+1,s_r} y V_r:C^s→C^{s_r} satisfacen sum_r V_r^*V_r=I_s, entonces la columna V=(V_1,\\ldots,V_k)^T es isométrica y la restricción genealógica conmuta con la combinación matricial: rho_m^s(sum_r V_r^*Phi_r V_r)=sum_r V_r^*(rho_m^{s_r}Phi_r)V_r.","status":"PRESENT_IN_SUCCESSOR_SOURCE"}
% HMT_RESULT_END:P15:residence
% HMT_RESULT_BEGIN:P16:residence
% HMT_ATOMIC_RESULT {"material_state":"INTEGRADO_EN_GRAFO_ACTIVO","proof_strength":"EXACTO_INTERNO","provenance":"FORMALIZACION_NUEVA","reinforcement_state":"DEPLOYED_WITHOUT_NEW_RESULT_ID","residence":"sistema inverso UCP, § Dilatación cilíndrica compatible y PVM común, tras P15","result_id":"P16","statement":"Existe una única Phi_infty:C(X_infty)→M_s que extiende todas las cartas cilíndricas y una dilatación común Phi_infty(f)=V*∫f dP V, con E_x^(m)=V*P([x]_m)V; la dilatación mínima es única hasta unitario y, fijada pi, la PVM P es única.","status":"PRESENT_IN_SUCCESSOR_SOURCE"}
% HMT_RESULT_END:P16:residence
% HMT_RESULT_BEGIN:P17:residence
% HMT_ATOMIC_RESULT {"material_state":"INTEGRADO_EN_GRAFO_ACTIVO","proof_strength":"EXACTO_INTERNO","provenance":"FORMALIZACION_NUEVA","reinforcement_state":"DEPLOYED_WITHOUT_NEW_RESULT_ID","residence":"continuación UHF","result_id":"P17","statement":"Si Psi_m=Psi_{m+1}∘iota_m, existe una única UCP Psi_infty:A_infty→M_s y una dilatación común Psi_m(a)=V* pi(iota_{m,infty}(a))V; pi puede hacerse fiel.","status":"PRESENT_IN_SUCCESSOR_SOURCE"}
% HMT_RESULT_END:P17:residence
% HMT_RESULT_BEGIN:P18:residence
% HMT_ATOMIC_RESULT {"material_state":"INTEGRADO_EN_GRAFO_ACTIVO","proof_strength":"EXACTO_INTERNO","provenance":"FORMALIZACION_NUEVA","reinforcement_state":"DEPLOYED_WITHOUT_NEW_RESULT_ID","residence":"capítulo «Cuadrados cuánticos de APP»","result_id":"P18","statement":"Existe QMS aditivo $(9,3)$.","status":"PRESENT_IN_SUCCESSOR_SOURCE"}
% HMT_RESULT_END:P18:residence
% HMT_RESULT_BEGIN:P19:residence
% HMT_ATOMIC_RESULT {"material_state":"INTEGRADO_EN_GRAFO_ACTIVO","proof_strength":"EXACTO_INTERNO","provenance":"RESULTADO_NUEVO","reinforcement_state":"DEPLOYED_WITHOUT_NEW_RESULT_ID","residence":"cuadrados APP, § Construcción A^Sigma, tras P18","result_id":"P19","statement":"Tres contextos producen QMS $(9,3)$ no conmutativo y semiclásico.","status":"PRESENT_IN_SUCCESSOR_SOURCE"}
% HMT_RESULT_END:P19:residence
% HMT_RESULT_BEGIN:P20:residence
% HMT_ATOMIC_RESULT {"material_state":"INTEGRADO_EN_GRAFO_ACTIVO","proof_strength":"EXACTO_INTERNO","provenance":"RESULTADO_NUEVO","reinforcement_state":"DEPLOYED_WITHOUT_NEW_RESULT_ID","residence":"cuadrados APP, § Construcción A^(9), tras P19","result_id":"P20","statement":"Las doce direcciones producen QMS $(12,3)$ semiclásico.","status":"PRESENT_IN_SUCCESSOR_SOURCE"}
% HMT_RESULT_END:P20:residence
% HMT_RESULT_BEGIN:P21:residence
% HMT_ATOMIC_RESULT {"material_state":"INTEGRADO_EN_GRAFO_ACTIVO","proof_strength":"EXACTO_INTERNO","provenance":"RESULTADO_NUEVO","reinforcement_state":"DEPLOYED_WITHOUT_NEW_RESULT_ID","residence":"cuadrados APP, § Construcción B(12,3) y cubo de tres índices, tras P20","result_id":"P21","statement":"Las mismas \\POVM producen cubos de órdenes 9 y 12.","status":"PRESENT_IN_SUCCESSOR_SOURCE"}
% HMT_RESULT_END:P21:residence
% HMT_RESULT_BEGIN:P22:residence
% HMT_ATOMIC_RESULT {"material_state":"INTEGRADO_EN_GRAFO_ACTIVO","proof_strength":"EXACTO_INTERNO","provenance":"FORMALIZACION_NUEVA","reinforcement_state":"DEPLOYED_WITHOUT_NEW_RESULT_ID","residence":"subsección de fibración finita","result_id":"P22","statement":"P^1(Z/9Z) tiene doce puntos y r tiene cuatro fibras de tres; la fibración es canónica, mientras el orden de cada fibra y la bijección Xi hacia resultados qutrit constituyen un calibre.","status":"PRESENT_IN_SUCCESSOR_SOURCE"}
% HMT_RESULT_END:P22:residence
% HMT_RESULT_BEGIN:P23:residence
% HMT_ATOMIC_RESULT {"material_state":"INTEGRADO_EN_GRAFO_ACTIVO","proof_strength":"EXACTO_INTERNO","provenance":"FORMALIZACION_NUEVA","reinforcement_state":"DEPLOYED_WITHOUT_NEW_RESULT_ID","residence":"geometría qutrit, § Fibración 12→4×3 y POVM calibrada, tras P22","result_id":"P23","statement":"F_k^(pi)=(3/pi)P_k para k=0,1,2 y F_r^(pi)=((1-3/pi)/6)I_3 para r=3,…,8 forman una POVM; para un proyector Q de otro contexto MUB, Tr(QF_k^(pi))=(3/pi)(1/3)=1/pi","status":"PRESENT_IN_SUCCESSOR_SOURCE"}
% HMT_RESULT_END:P23:residence
% HMT_RESULT_BEGIN:P24:residence
% HMT_ATOMIC_RESULT {"material_state":"INTEGRADO_EN_GRAFO_ACTIVO","proof_strength":"EXACTO_INTERNO","provenance":"RESULTADO_NUEVO","reinforcement_state":"DEPLOYED_WITHOUT_NEW_RESULT_ID","residence":"después de separación espectral APP","result_id":"P24","statement":"El polinomio de CC* es t^5(t-1)(t-5/7)(t-1/3)^2 y la descomposición de Halmos da C*(P_{Sigma,2},P_{Pi,2})≅C^4⊕M_2(C)⊕M_2(C), de dimensión 12.","status":"PRESENT_IN_SUCCESSOR_SOURCE"}
% HMT_RESULT_END:P24:residence
% HMT_RESULT_BEGIN:P25:residence
% HMT_ATOMIC_RESULT {"material_state":"INTEGRADO_EN_GRAFO_ACTIVO","proof_strength":"EXACTO_INTERNO","provenance":"RESULTADO_NUEVO","reinforcement_state":"DEPLOYED_WITHOUT_NEW_RESULT_ID","residence":"doble proyección operatorial, § Intervalo libre del bloque B_c, tras P24","result_id":"P25","statement":"Para todo nivel s, W_s(B_c)={Y=Y*:5I_s≼Y≼9I_s}; 7 es centro, 2 radio operatorial y 4 diámetro espectral.","status":"PRESENT_IN_SUCCESSOR_SOURCE"}
% HMT_RESULT_END:P25:residence
% HMT_RESULT_BEGIN:P26:residence
% HMT_ATOMIC_RESULT {"material_state":"INTEGRADO_EN_GRAFO_ACTIVO","proof_strength":"EXACTO_INTERNO","provenance":"FORMALIZACION_NUEVA","reinforcement_state":"DEPLOYED_WITHOUT_NEW_RESULT_ID","residence":"doble proyección, antes de TRI-001","result_id":"P26","statement":"Sus PVM diagonales conmutan y admiten el refinamiento común P_{a,e}=1_{R_m=a}1_{Q_m=e}; si los lectores respetan truncaciones, el refinamiento pasa a una PVM cilíndrica común.","status":"PRESENT_IN_SUCCESSOR_SOURCE"}
% HMT_RESULT_END:P26:residence
% HMT_RESULT_BEGIN:P27:residence
% HMT_ATOMIC_RESULT {"material_state":"INTEGRADO_EN_GRAFO_ACTIVO","proof_strength":"EXACTO_INTERNO","provenance":"FORMALIZACION_NUEVA","reinforcement_state":"DEPLOYED_WITHOUT_NEW_RESULT_ID","residence":"doble proyección operatorial, § Límite de dilatación MUB, tras P26","result_id":"P27","statement":"No admiten una PVM conjunta aguda: sus proyectores tienen solapamiento 1/3 y no conmutan; una realización conjunta requiere ruido, POVM o un ambiente mayor.","status":"PRESENT_IN_SUCCESSOR_SOURCE"}
% HMT_RESULT_END:P27:residence
% HMT_RESULT_BEGIN:P28:residence
% HMT_ATOMIC_RESULT {"material_state":"INTEGRADO_EN_GRAFO_ACTIVO","proof_strength":"EXACTO_INTERNO","provenance":"RESULTADO_RECUPERADO","reinforcement_state":"DEPLOYED_WITHOUT_NEW_RESULT_ID","residence":"frontera del capítulo de continuo operatorial","result_id":"P28","statement":"El teorema importado afirma que la MPS posee límite continuo si y sólo si E_A es infinitamente divisible, equivalentemente E_A=P exp(L), P^2=P y PLP=PL; una potencia E_A^r basta para el límite continuo grueso.","status":"PRESENT_IN_SUCCESSOR_SOURCE"}
% HMT_RESULT_END:P28:residence
% HMT_RESULT_BEGIN:P29:residence
% HMT_ATOMIC_RESULT {"material_state":"INTEGRADO_EN_GRAFO_ACTIVO","proof_strength":"EXACTO_INTERNO","provenance":"FORMALIZACION_NUEVA","reinforcement_state":"DEPLOYED_WITHOUT_NEW_RESULT_ID","residence":"medición/expectativa tracial","result_id":"P29","statement":"E_m es UCP, satisface E_m(a⊗I_9)=a y posee Kraus W_a psi=(1/3)psi⊗e_a y una dilatación cuyo ancilla registra el dígito promediado; en la diagonal, la retracción de rama kappa_{m,a}(f)(x)=f(x,a) es multiplicativa y no equivale a promediar. Esa memoria no se identifica aún con TPK.","status":"PRESENT_IN_SUCCESSOR_SOURCE"}
% HMT_RESULT_END:P29:residence
% HMT_RESULT_BEGIN:P30:residence
% HMT_ATOMIC_RESULT {"material_state":"INTEGRADO_EN_GRAFO_ACTIVO","proof_strength":"EXACTO_INTERNO","provenance":"FORMALIZACION_NUEVA","reinforcement_state":"DEPLOYED_WITHOUT_NEW_RESULT_ID","residence":"APP matricial","result_id":"P30","statement":"El doble centrado inserta una señal APP en un QMS.","status":"PRESENT_IN_SUCCESSOR_SOURCE"}
% HMT_RESULT_END:P30:residence
% HMT_RESULT_BEGIN:P31:residence
% HMT_ATOMIC_RESULT {"material_state":"INTEGRADO_EN_GRAFO_ACTIVO","proof_strength":"PROGRAMA_ABIERTO","provenance":"ARQUITECTURA_AUTORAL_PREEXISTENTE","reinforcement_state":"DEPLOYED_WITHOUT_NEW_RESULT_ID","residence":"perspectivas operatoriales","result_id":"P31","statement":"El programa asigna a cada tipo c un cono de lectores positivos A(c), con transformaciones naturales por niveles; requiere funtor de conos, objeto dualizante, evaluaciones, coevaluaciones y compatibilidad de la composición TPK. La profundidad de separación es el menor m donde los sistemas mínimo y máximo difieren; el nivel matricial de detección es el menor s que exhibe esa diferencia; y el defecto de memoria es la mínima información de ruta necesaria para levantar un punto del sistema máximo al mínimo.","status":"PRESENT_IN_SUCCESSOR_SOURCE"}
% HMT_RESULT_END:P31:residence
% HMT_RESULT_BEGIN:P32:residence
% HMT_ATOMIC_RESULT {"material_state":"INTEGRADO_EN_GRAFO_ACTIVO","proof_strength":"PROGRAMA_ABIERTO","provenance":"ARQUITECTURA_AUTORAL_PREEXISTENTE","reinforcement_state":"DEPLOYED_WITHOUT_NEW_RESULT_ID","residence":"perspectivas, § Universalidad genealógicamente fiel, tras UNIIMP-001","result_id":"P32","statement":"La versión HMT exige un funtor entre categorías de rutas que preserve composición, retorno y memoria, además de los datos del marco target--context importado.","status":"PRESENT_IN_SUCCESSOR_SOURCE"}
% HMT_RESULT_END:P32:residence
% HMT_RESULT_BEGIN:QPM-001:residence
% HMT_ATOMIC_RESULT {"material_state":"INTEGRADO_EN_GRAFO_ACTIVO","proof_strength":"EXACTO_INTERNO","provenance":"FORMALIZACION_NUEVA","reinforcement_state":"DEPLOYED_WITHOUT_NEW_RESULT_ID","residence":"apertura del corredor matricial, tras P09","result_id":"QPM-001","statement":"Una QPM es un QMS cuyas celdas son proyecciones; es conmutativa sólo si todas esas proyecciones conmutan conjuntamente. La definición no se obtiene dilatando cada fila por separado.","status":"PRESENT_IN_SUCCESSOR_SOURCE"}
% HMT_RESULT_END:QPM-001:residence
% HMT_RESULT_BEGIN:QLS-001:residence
% HMT_ATOMIC_RESULT {"material_state":"INTEGRADO_EN_GRAFO_ACTIVO","proof_strength":"EXACTO_INTERNO","provenance":"FORMALIZACION_NUEVA","reinforcement_state":"DEPLOYED_WITHOUT_NEW_RESULT_ID","residence":"apertura del corredor matricial, antes de QLSI-001","result_id":"QLS-001","statement":"Un QLS es una cuadrícula de vectores unitarios v_ij en C^n cuyas filas y columnas son bases ortonormales; las proyecciones P_ij=|v_ij><v_ij| forman un QMS proyectivo. Un QLS fácil se obtiene de un cuadrado latino clásico y una base ortonormal fija.","status":"PRESENT_IN_SUCCESSOR_SOURCE"}
% HMT_RESULT_END:QLS-001:residence
% HMT_RESULT_BEGIN:QLSI-001:residence
% HMT_ATOMIC_RESULT {"material_state":"INTEGRADO_EN_GRAFO_ACTIVO","proof_strength":"EXACTO_INTERNO","provenance":"RESULTADO_RECUPERADO","reinforcement_state":"DEPLOYED_WITHOUT_NEW_RESULT_ID","residence":"apertura del corredor matricial, tras P11","result_id":"QLSI-001","statement":"La envolvente matricial convexa de los QLS contiene estrictamente más que la clase semiclásica.","status":"PRESENT_IN_SUCCESSOR_SOURCE"}
% HMT_RESULT_END:QLSI-001:residence
% HMT_RESULT_BEGIN:QMSI-001:residence
% HMT_ATOMIC_RESULT {"material_state":"INTEGRADO_EN_GRAFO_ACTIVO","proof_strength":"EXACTO_INTERNO","provenance":"RESULTADO_RECUPERADO","reinforcement_state":"DEPLOYED_WITHOUT_NEW_RESULT_ID","residence":"apertura del corredor matricial, tras P10","result_id":"QMSI-001","statement":"Un QMS es semiclásico si y sólo si admite una dilatación simultánea a una QPM conmutativa mediante una única compresión común.","status":"PRESENT_IN_SUCCESSOR_SOURCE"}
% HMT_RESULT_END:QMSI-001:residence
% HMT_RESULT_BEGIN:QMSI-002:residence
% HMT_ATOMIC_RESULT {"material_state":"INTEGRADO_EN_GRAFO_ACTIVO","proof_strength":"EXACTO_INTERNO","provenance":"RESULTADO_RECUPERADO","reinforcement_state":"DEPLOYED_WITHOUT_NEW_RESULT_ID","residence":"apertura matricial, § Existencia de QMS no semiclásicos, tras QMSI-001","result_id":"QMSI-002","statement":"Para n≥3 y s≥2 existen QMS no semiclásicos.","status":"PRESENT_IN_SUCCESSOR_SOURCE"}
% HMT_RESULT_END:QMSI-002:residence
% HMT_RESULT_BEGIN:QMSI-003:residence
% HMT_ATOMIC_RESULT {"material_state":"INTEGRADO_EN_GRAFO_ACTIVO","proof_strength":"EXACTO_INTERNO","provenance":"RESULTADO_RECUPERADO","reinforcement_state":"DEPLOYED_WITHOUT_NEW_RESULT_ID","residence":"apertura del corredor matricial, tras QMSI-002","result_id":"QMSI-003","statement":"Para todo n≥3 existen QMS sobre M_2(C) que no admiten dilatación a ninguna QPM.","status":"PRESENT_IN_SUCCESSOR_SOURCE"}
% HMT_RESULT_END:QMSI-003:residence
% HMT_RESULT_BEGIN:QBR-001:residence
% HMT_ATOMIC_RESULT {"material_state":"INTEGRADO_EN_GRAFO_ACTIVO","proof_strength":"EXACTO_INTERNO","provenance":"RESULTADO_RECUPERADO","reinforcement_state":"DEPLOYED_WITHOUT_NEW_RESULT_ID","residence":"`md/04ab`, después de la realización hilbertiana finita","result_id":"QBR-001","statement":"Los observables se representan por operadores autoadjuntos y sus PVM; para un suceso cilíndrico P y estado normalizado Psi, la probabilidad es ||P Psi||^2.","status":"PRESENT_IN_SUCCESSOR_SOURCE"}
% HMT_RESULT_END:QBR-001:residence
% HMT_RESULT_BEGIN:QDY-001:residence
% HMT_ATOMIC_RESULT {"material_state":"INTEGRADO_EN_GRAFO_ACTIVO","proof_strength":"EXACTO_INTERNO","provenance":"RESULTADO_RECUPERADO","reinforcement_state":"DEPLOYED_WITHOUT_NEW_RESULT_ID","residence":"`md/04ab`, dinámica cuántica","result_id":"QDY-001","statement":"Un grupo unitario fuertemente continuo posee un generador autoadjunto H y satisface la ecuación de Schrödinger en su dominio.","status":"PRESENT_IN_SUCCESSOR_SOURCE"}
% HMT_RESULT_END:QDY-001:residence
% HMT_RESULT_BEGIN:QKR-001:residence
% HMT_ATOMIC_RESULT {"material_state":"INTEGRADO_EN_GRAFO_ACTIVO","proof_strength":"EXACTO_INTERNO","provenance":"RESULTADO_RECUPERADO","reinforcement_state":"DEPLOYED_WITHOUT_NEW_RESULT_ID","residence":"`md/01a`, instrumentos y canales","result_id":"QKR-001","statement":"Una familia de operadores de Kraus con suma K_a^*K_a=I define un canal completamente positivo y traza-preservante; cada resultado conserva su efecto e instrumento tipados.","status":"PRESENT_IN_SUCCESSOR_SOURCE"}
% HMT_RESULT_END:QKR-001:residence
% HMT_RESULT_BEGIN:QMO-001:residence
% HMT_ATOMIC_RESULT {"material_state":"INTEGRADO_EN_GRAFO_ACTIVO","proof_strength":"EXACTO_INTERNO","provenance":"RESULTADO_RECUPERADO","reinforcement_state":"DEPLOYED_WITHOUT_NEW_RESULT_ID","residence":"`md/04ab`, composición de registros","result_id":"QMO-001","statement":"Una realización monoidal fuerte transporta la composición independiente de registros al producto tensorial, con asociatividad, unidad y coherencias declaradas.","status":"PRESENT_IN_SUCCESSOR_SOURCE"}
% HMT_RESULT_END:QMO-001:residence
% HMT_RESULT_BEGIN:DIC-001:residence
% HMT_ATOMIC_RESULT {"material_state":"INTEGRADO_EN_GRAFO_ACTIVO","proof_strength":"EXACTO_INTERNO","provenance":"FORMALIZACION_NUEVA","reinforcement_state":"DEPLOYED_WITHOUT_NEW_RESULT_ID","residence":"apertura del corredor matricial","result_id":"DIC-001","statement":"El diccionario conserva exactamente: tamaño exterior n↔cuadrícula publicada; tamaño interior s↔carta operatorial; profundidad m↔genealogía; POVM↔familia de lectores positivos; PVM↔contexto espectral; UCP↔lector matricial; compresión↔sombra; dilatación↔antecedente operatorio; semiclasicidad↔POVM primaria común reordenada; no semiclasicidad↔inexistencia de una POVM primaria común que reproduzca todas las celdas; sistema de operadores↔espacio de lectores; rango matricial↔sombras a todos los s; límite inductivo↔cierre de refinamientos operatorios; y límite inverso↔sección infinita compatible. La composición genuinamente bidimensional es la interpretación HMT propuesta, no la definición importada.","status":"PRESENT_IN_SUCCESSOR_SOURCE"}
% HMT_RESULT_END:DIC-001:residence
% HMT_RESULT_BEGIN:QMC-001:residence
% HMT_ATOMIC_RESULT {"material_state":"INTEGRADO_EN_GRAFO_ACTIVO","proof_strength":"EXACTO_INTERNO","provenance":"RESULTADO_RECUPERADO","reinforcement_state":"DEPLOYED_WITHOUT_NEW_RESULT_ID","residence":"capítulo de cuadrados cuánticos, antes de las instancias \\((9,3)\\)","result_id":"QMC-001","statement":"A es un QMS (n,s), admite la descomposición semiclásica A=sum_t P_{sigma_t}⊗E_t y toda dilatación Naimark común E_t=V*Q_tV induce un QPM conmutativo simultáneo.","status":"PRESENT_IN_SUCCESSOR_SOURCE"}
% HMT_RESULT_END:QMC-001:residence
% HMT_RESULT_BEGIN:QDT-001:residence
% HMT_ATOMIC_RESULT {"material_state":"INTEGRADO_EN_GRAFO_ACTIVO","proof_strength":"EXACTO_INTERNO","provenance":"FORMALIZACION_NUEVA","reinforcement_state":"DEPLOYED_WITHOUT_NEW_RESULT_ID","residence":"geometría qutrit, antes de la fibración \\(12\\to4\\times3\\)","result_id":"QDT-001","statement":"M_3(C)=CI⊕direct_sum_{a∈P^1(F_3)}D_a^0 ortogonalmente para Hilbert--Schmidt y dim=9=1+4·2.","status":"PRESENT_IN_SUCCESSOR_SOURCE"}
% HMT_RESULT_END:QDT-001:residence
% HMT_RESULT_BEGIN:ENV-001:residence
% HMT_ATOMIC_RESULT {"material_state":"INTEGRADO_EN_GRAFO_ACTIVO","proof_strength":"EXACTO_INTERNO","provenance":"FORMALIZACION_NUEVA","reinforcement_state":"DEPLOYED_WITHOUT_NEW_RESULT_ID","residence":"dilatación cilíndrica, después de P16","result_id":"ENV-001","statement":"Sumar una representación fiel pi_f a una dilatación conserva todas las compresiones finitas y hace fiel la representación ambiental; esto no prueba todavía fidelidad dinámica TPK.","status":"PRESENT_IN_SUCCESSOR_SOURCE"}
% HMT_RESULT_END:ENV-001:residence
% HMT_RESULT_BEGIN:RNG-001:residence
% HMT_ATOMIC_RESULT {"material_state":"INTEGRADO_EN_GRAFO_ACTIVO","proof_strength":"EXACTO_INTERNO","provenance":"RESULTADO_NUEVO","reinforcement_state":"DEPLOYED_WITHOUT_NEW_RESULT_ID","residence":"rango matricial APP, después de P24","result_id":"RNG-001","statement":"(X,Y) pertenece a W_s(APP_2) exactamente cuando existen cuatro efectos escalares E_ab y dos mapas CP Psi_{5/7},Psi_{1/3}:M_2→M_s que satisfacen las ecuaciones completas de normalización y coordenadas.","status":"PRESENT_IN_SUCCESSOR_SOURCE"}
% HMT_RESULT_END:RNG-001:residence
% HMT_RESULT_BEGIN:DYN-001:residence
% HMT_ATOMIC_RESULT {"material_state":"INTEGRADO_EN_GRAFO_ACTIVO","proof_strength":"EXACTO_INTERNO","provenance":"FORMALIZACION_NUEVA","reinforcement_state":"DEPLOYED_WITHOUT_NEW_RESULT_ID","residence":"continuo operatorial, después de Stinespring","result_id":"DYN-001","statement":"La compatibilidad UCP a toda profundidad no implica ancilla finito uniforme; la fidelidad dinámica exige J_{m+1}T_m^TPK=tilde T_m J_m y naturalidad de las compresiones.","status":"PRESENT_IN_SUCCESSOR_SOURCE"}
% HMT_RESULT_END:DYN-001:residence
% HMT_RESULT_BEGIN:TEN-001:residence
% HMT_ATOMIC_RESULT {"material_state":"INTEGRADO_EN_GRAFO_ACTIVO","proof_strength":"PROGRAMA_ABIERTO","provenance":"ARQUITECTURA_AUTORAL_PREEXISTENTE","reinforcement_state":"DEPLOYED_WITHOUT_NEW_RESULT_ID","residence":"perspectivas del continuo tensorial","result_id":"TEN-001","statement":"El programa exige declarar tensor local, grupo de simetría, régimen exacto/aproximado/asintótico y estabilidad de rango; propone una torre coherente de raíces E_{kℓ}^ℓ=E_k equipada con registros TPK.","status":"PRESENT_IN_SUCCESSOR_SOURCE"}
% HMT_RESULT_END:TEN-001:residence
% HMT_RESULT_BEGIN:CHP-001:residence
% HMT_ATOMIC_RESULT {"material_state":"INTEGRADO_EN_GRAFO_ACTIVO","proof_strength":"PROGRAMA_ABIERTO","provenance":"FORMALIZACION_NUEVA","reinforcement_state":"DEPLOYED_WITHOUT_NEW_RESULT_ID","residence":"sistemas cónicos, antes de universalidad","result_id":"CHP-001","statement":"Las condiciones sum_j J(Phi_ij)=J(id) y sum_i J(Phi_ij)=J(id), con J(Phi_ij) positivo, linealizan la normalización mágica; la fidelidad de hoja y ruta es una obligación adicional.","status":"PRESENT_IN_SUCCESSOR_SOURCE"}
% HMT_RESULT_END:CHP-001:residence
% HMT_RESULT_BEGIN:SDP-001:residence
% HMT_ATOMIC_RESULT {"material_state":"INTEGRADO_EN_GRAFO_ACTIVO","proof_strength":"EXACTO_INTERNO","provenance":"FORMALIZACION_NUEVA","reinforcement_state":"DEPLOYED_WITHOUT_NEW_RESULT_ID","residence":"cierre del capítulo qutrit","result_id":"SDP-001","statement":"Se buscan Q_sigma positivos con sum_sigma Q_sigma=I y A_ij=sum_{sigma: sigma(i)=j}Q_sigma; factibilidad certifica semiclasicidad e inviabilidad requiere testigo dual.","status":"PRESENT_IN_SUCCESSOR_SOURCE"}
% HMT_RESULT_END:SDP-001:residence
% HMT_RESULT_BEGIN:EQV-001:residence
% HMT_ATOMIC_RESULT {"material_state":"INTEGRADO_EN_GRAFO_ACTIVO","proof_strength":"EXACTO_INTERNO","provenance":"FORMALIZACION_NUEVA","reinforcement_state":"DEPLOYED_WITHOUT_NEW_RESULT_ID","residence":"geometría simpléctica y lectores, después del calibre \\(3/\\pi\\)","result_id":"EQV-001","statement":"El calibre preliminar falla un test antiunitario y no es aún entrelazador; la promoción exige Xi(g·ell)=U_g Xi(ell)U_g* en todas las fibras, formulada como naturalidad con origen y destino cuando g es una flecha.","status":"PRESENT_IN_SUCCESSOR_SOURCE"}
% HMT_RESULT_END:EQV-001:residence
% HMT_RESULT_BEGIN:QMT-001:residence
% HMT_ATOMIC_RESULT {"material_state":"INTEGRADO_EN_GRAFO_ACTIVO","proof_strength":"PROGRAMA_ABIERTO","provenance":"FORMALIZACION_NUEVA","reinforcement_state":"DEPLOYED_WITHOUT_NEW_RESULT_ID","residence":"cierre del capítulo de cuadrados cuánticos APP","result_id":"QMT-001","statement":"La tabla conserva tres filas distintas: A^Sigma (9,3), que cierra hoja suma, positividad, normalización y semiclasicidad pero no hoja producto, cociente ni rutas; A^(9) (9,3), que cierra tres contextos, no conmutatividad y dilatación común pero no cuarto contexto, incidencia Hensel, gnomon ni memoria TPK; y B (12,3), que cierra el atlas 12→4×3 y su dilatación común pero no equivarianza, dinámica ni monodromía.","status":"PRESENT_IN_SUCCESSOR_SOURCE"}
% HMT_RESULT_END:QMT-001:residence
% HMT_RESULT_BEGIN:TRN-001:residence
% HMT_ATOMIC_RESULT {"material_state":"INTEGRADO_EN_GRAFO_ACTIVO","proof_strength":"PROGRAMA_ABIERTO","provenance":"ARQUITECTURA_AUTORAL_PREEXISTENTE","reinforcement_state":"DEPLOYED_WITHOUT_NEW_RESULT_ID","residence":"cierre metrológico del corredor matricial","result_id":"TRN-001","statement":"Conserva ocho filas completas: (1) monodromía 9→1→retorno del observable con estado ambiental distinto→clasificar dilataciones que preservan holonomía; (2) pi,e,phi,alpha→cuatro lectores compatibles de una genealogía→existencia y minimalidad de un sistema UCP común; (3) doble proyección→dos marginales de un antecedente cilíndrico→antecedente no conmutativo y separación conjunta de fibras; (4) B_c y salto 4→intervalo matricial libre [5I,9I]→extremos y transporte de la jerarquía 4;2;6;54; (5) Planck y CT108→punto autodual y lectores conjugados→realización como simetría de un sistema cónico con memoria; (6) ley de masas→familia positiva sobre fibras y rutas→rango matricial, extremos y realizaciones por sector; (7) CKM/CP→transporte de rango tres y fase no trivial→canal covariante que preserve espectro y orientación; (8) simetría excepcional→segundo codominio de la doble proyección→entrelazador entre atlas nonádico, código y representación.","status":"PRESENT_IN_SUCCESSOR_SOURCE"}
% HMT_RESULT_END:TRN-001:residence
% HMT_RESULT_BEGIN:MPSH-001:residence
% HMT_ATOMIC_RESULT {"material_state":"INTEGRADO_EN_GRAFO_ACTIVO","proof_strength":"PROGRAMA_ABIERTO","provenance":"ARQUITECTURA_AUTORAL_PREEXISTENTE","reinforcement_state":"DEPLOYED_WITHOUT_NEW_RESULT_ID","residence":"inmediatamente después de P28","result_id":"MPSH-001","statement":"Aplicar el teorema MPS a HMT exige construir una celda local uniforme, tensor A, canal E_A, forma irreducible, divisibilidad infinita propia o de una potencia y transporte al límite de hoja, acarreo y monodromía.","status":"PRESENT_IN_SUCCESSOR_SOURCE"}
% HMT_RESULT_END:MPSH-001:residence
% HMT_RESULT_BEGIN:UNIIMP-001:residence
% HMT_ATOMIC_RESULT {"material_state":"INTEGRADO_EN_GRAFO_ACTIVO","proof_strength":"EXACTO_INTERNO","provenance":"RESULTADO_RECUPERADO","reinforcement_state":"DEPLOYED_WITHOUT_NEW_RESULT_ID","residence":"perspectivas operatoriales, antes de P32","result_id":"UNIIMP-001","statement":"Un simulador s:P⊗C→T⊗C es universal cuando existe la reducción laxa prescrita al simulador trivial; un monótono compatible cuyo supremo sobre la imagen funcional quede por debajo de algún target constituye un no-go.","status":"PRESENT_IN_SUCCESSOR_SOURCE"}
% HMT_RESULT_END:UNIIMP-001:residence
% HMT_RESULT_BEGIN:PRG-001:residence
% HMT_ATOMIC_RESULT {"material_state":"INTEGRADO_EN_GRAFO_ACTIVO","proof_strength":"PROGRAMA_ABIERTO","provenance":"ARQUITECTURA_AUTORAL_PREEXISTENTE","reinforcement_state":"DEPLOYED_WITHOUT_NEW_RESULT_ID","residence":"cierre del corredor matricial","result_id":"PRG-001","statement":"Los cierres son: sistema APP filtrado, QMS APP--TPK fiel, frontera semiclásica, canal de continuo, sistema cónico intrínseco y realización metrológica UCP conjunta; ninguno se promueve sin dominio, prueba/certificado, positividad, compatibilidad, falsadores y trazabilidad.","status":"PRESENT_IN_SUCCESSOR_SOURCE"}
% HMT_RESULT_END:PRG-001:residence
% HMT_RESULT_BEGIN:QCT-001:residence
% HMT_ATOMIC_RESULT {"material_state":"CONTROL_FOCAL_ACTIVO","proof_strength":"EXACTO_INTERNO","provenance":"CERTIFICADO_NUEVO","reinforcement_state":"DEPLOYED_WITHOUT_NEW_RESULT_ID","residence":"apéndice digital y gate del puente","result_id":"QCT-001","statement":"Registra exactamente dieciocho controles: PVM qutrit; idempotencia; MUB; no-go QLS; QMS aditivo (9,3); QMS de tres contextos; no conmutatividad; QMS atlas (12,3); descomposición semiclásica; cubos de órdenes 9 y 12; P^1(Z/9)→P^1(F3); truncación--compresión; lector 1/pi; calibre 3/pi; polinomio de Gram; intersecciones de Halmos; y C*-álgebra APP d=2.","status":"CONTROLLED"}
% HMT_RESULT_END:QCT-001:residence
% HMT_RESULT_BEGIN:PRV-001:residence
% HMT_ATOMIC_RESULT {"material_state":"CONTROL_FOCAL_ACTIVO","proof_strength":"EXACTO_INTERNO","provenance":"FORMALIZACION_NUEVA","reinforcement_state":"DEPLOYED_WITHOUT_NEW_RESULT_ID","residence":"registro bibliográfico del corredor matricial","result_id":"PRV-001","statement":"Conserva autores, firmas e identificadores exactamente: De las Cuevas--Drescher--Netzer, Quantum magic squares: dilations and their limitations—DOI 10.1063/5.0022344—QMS/QPM/dilatación; Magic squares: Latin, Semiclassical and Quantum, De las Cuevas--Netzer--Valentiner-Branth—DOI 10.1063/5.0127393—QLS/envolventes; Quantum information theory and free semialgebraic geometry, De las Cuevas--Netzer—arXiv:2102.04240—torres/Choi; Beyond Operator Systems, De les Coves--Mirte van der Eyden--Netzer—DOI 10.1016/j.jmaa.2025.130102—sistemas cónicos; Continuum limits of Matrix Product States, De las Cuevas--Schuch--Pérez-García--Cirac—DOI 10.1103/PhysRevB.98.174303—divisibilidad; Tensor decompositions on simplicial complexes with invariance, De las Cuevas--Hoogsteder Riera--Netzer—DOI 10.1016/j.jsc.2024.102299—simetría; Approximate tensor decompositions, De las Cuevas--Klingler--Netzer—DOI 10.1063/5.0033876—aproximación; A Framework for Universality in Physics, Computer Science, and Beyond, Gonda--Reinhart--Stengele--De les Coves—DOI 10.46298/compositionality-6-3—target--context. Separa literatura simpléctica vecina y fija De las Cuevas/De les Coves.","status":"CONTROLLED"}
% HMT_RESULT_END:PRV-001:residence
% HMT_RESULT_BEGIN:MOR-001:residence
% HMT_ATOMIC_RESULT {"material_state":"INTEGRADO_O_REMITIDO_SIN_PERDIDA","proof_strength":"EXACTO_INTERNO","provenance":"ARQUITECTURA_AUTORAL_PREEXISTENTE","reinforcement_state":"DEPLOYED_WITHOUT_NEW_RESULT_ID","residence":"apertura continuo y `md/00b`","result_id":"MOR-001","statement":"Una cifra es una coordenada finita publicada; un valor es la imagen arquimediana, espectral o física de una sección; un número HMT es la sección completa con profundidad, hoja, orientación, acarreo, ruta, frontera, memoria y supervivencia; una simetría HMT conserva la identidad enriquecida y transporta su memoria.","status":"PRESENT_IN_SUCCESSOR_SOURCE"}
% HMT_RESULT_END:MOR-001:residence
% HMT_RESULT_BEGIN:MOR-002:residence
% HMT_ATOMIC_RESULT {"material_state":"INTEGRADO_O_REMITIDO_SIN_PERDIDA","proof_strength":"EXACTO_INTERNO","provenance":"RESULTADO_RECUPERADO","reinforcement_state":"DEPLOYED_WITHOUT_NEW_RESULT_ID","residence":"APP","result_id":"MOR-002","statement":"suma/producto y memoria de fibras","status":"PRESENT_IN_SUCCESSOR_SOURCE"}
% HMT_RESULT_END:MOR-002:residence
% HMT_RESULT_BEGIN:MOR-003:residence
% HMT_ATOMIC_RESULT {"material_state":"INTEGRADO_O_REMITIDO_SIN_PERDIDA","proof_strength":"EXACTO_INTERNO","provenance":"ARQUITECTURA_AUTORAL_PREEXISTENTE","reinforcement_state":"DEPLOYED_WITHOUT_NEW_RESULT_ID","residence":"TRIT","result_id":"MOR-003","statement":"orientación temporal y acarreo","status":"PRESENT_IN_SUCCESSOR_SOURCE"}
% HMT_RESULT_END:MOR-003:residence
% HMT_RESULT_BEGIN:MOR-004:residence
% HMT_ATOMIC_RESULT {"material_state":"INTEGRADO_O_REMITIDO_SIN_PERDIDA","proof_strength":"EXACTO_INTERNO","provenance":"FORMALIZACION_NUEVA","reinforcement_state":"DEPLOYED_WITHOUT_NEW_RESULT_ID","residence":"TPK","result_id":"MOR-004","statement":"composición con memoria","status":"PRESENT_IN_SUCCESSOR_SOURCE"}
% HMT_RESULT_END:MOR-004:residence
% HMT_RESULT_BEGIN:MOR-005:residence
% HMT_ATOMIC_RESULT {"material_state":"INTEGRADO_O_REMITIDO_SIN_PERDIDA","proof_strength":"EXACTO_INTERNO","provenance":"RESULTADO_RECUPERADO","reinforcement_state":"DEPLOYED_WITHOUT_NEW_RESULT_ID","residence":"continuo, tras `hmt/00b`","result_id":"MOR-005","statement":"9^n 1_C mu -> delta_x","status":"PRESENT_IN_SUCCESSOR_SOURCE"}
% HMT_RESULT_END:MOR-005:residence
% HMT_RESULT_BEGIN:MOR-006:residence
% HMT_ATOMIC_RESULT {"material_state":"INTEGRADO_O_REMITIDO_SIN_PERDIDA","proof_strength":"EXACTO_INTERNO","provenance":"FORMALIZACION_NUEVA","reinforcement_state":"DEPLOYED_WITHOUT_NEW_RESULT_ID","residence":"bisagra continuo--cuántica, § Kronecker--cilindro--delta de Dirac, tras MOR-005","result_id":"MOR-006","statement":"proyectores condicionales y núcleo","status":"PRESENT_IN_SUCCESSOR_SOURCE"}
% HMT_RESULT_END:MOR-006:residence
% HMT_RESULT_BEGIN:MOR-007:residence
% HMT_ATOMIC_RESULT {"material_state":"INTEGRADO_O_REMITIDO_SIN_PERDIDA","proof_strength":"EXACTO_INTERNO","provenance":"RESULTADO_RECUPERADO","reinforcement_state":"DEPLOYED_WITHOUT_NEW_RESULT_ID","residence":"UHF--\\(II_1\\) y `md/04ab`","result_id":"MOR-007","statement":"La inclusión unital nonádica genera UHF(9^infty), canónicamente isomorfa a UHF(3^infty); la traza GNS conduce al factor hiperinfinito II_1. El cuaderno de 64 páginas conserva la prueba propietaria completa.","status":"PRESENT_IN_SUCCESSOR_SOURCE"}
% HMT_RESULT_END:MOR-007:residence
% HMT_RESULT_BEGIN:MOR-008:residence
% HMT_ATOMIC_RESULT {"material_state":"INTEGRADO_O_REMITIDO_SIN_PERDIDA","proof_strength":"EXACTO_INTERNO","provenance":"RESULTADO_RECUPERADO","reinforcement_state":"DEPLOYED_WITHOUT_NEW_RESULT_ID","residence":"`md/01a`, `04ab`","result_id":"MOR-008","statement":"Born--Lüders--Kraus","status":"PRESENT_IN_SUCCESSOR_SOURCE"}
% HMT_RESULT_END:MOR-008:residence
% HMT_RESULT_BEGIN:MOR-009:residence
% HMT_ATOMIC_RESULT {"material_state":"INTEGRADO_O_REMITIDO_SIN_PERDIDA","proof_strength":"INTERFAZ_FISICA","provenance":"RESULTADO_RECUPERADO","reinforcement_state":"DEPLOYED_WITHOUT_NEW_RESULT_ID","residence":"`md/04ab`, `00b`","result_id":"MOR-009","statement":"Q:P_TPK→Hilb_{R,J}; la extensión a flechas inversas se afirma sólo en el subdominio reversible donde el transporte unitario fuente teoremática la demuestra","status":"PRESENT_IN_SUCCESSOR_SOURCE"}
% HMT_RESULT_END:MOR-009:residence
% HMT_RESULT_BEGIN:MOR-010:residence
% HMT_ATOMIC_RESULT {"material_state":"INTEGRADO_O_REMITIDO_SIN_PERDIDA","proof_strength":"PROGRAMA_ABIERTO","provenance":"FORMALIZACION_NUEVA","reinforcement_state":"DEPLOYED_WITHOUT_NEW_RESULT_ID","residence":"después de `md/00b`, antes de la ley dimensional/Planck","result_id":"MOR-010","statement":"R_MD^{disc,+} es el candidato de producto con blanco Hilbert general y sólo queda construido tras demostrar identidad y composición componente a componente; U_reg existe sólo en el subdominio unitario; R_MD^{disc,or} existe únicamente cuando espectro e Hilbert son invertibles, las relaciones localizadas se respetan y la dimensional toma valores en Gr(M_D^+); positiva y orientada son lectoras hermanas; no se presupone ni exige una factorización horizontal entre ellas. El monoide de coeficientes satisface (n,a,b,lambda)(n',a',b',lambda')=(n+n',a+a',b+b',lambda lambda') y sólo después se aplica la completación de Grothendieck.","status":"PRESENT_IN_SUCCESSOR_SOURCE"}
% HMT_RESULT_END:MOR-010:residence
% HMT_RESULT_BEGIN:MOR-011:residence
% HMT_ATOMIC_RESULT {"material_state":"INTEGRADO_O_REMITIDO_SIN_PERDIDA","proof_strength":"EXACTO_INTERNO","provenance":"RESULTADO_RECUPERADO","reinforcement_state":"DEPLOYED_WITHOUT_NEW_RESULT_ID","residence":"`md/11m`","result_id":"MOR-011","statement":"C U(gamma) C^-1=U(gamma^-1)","status":"PRESENT_IN_SUCCESSOR_SOURCE"}
% HMT_RESULT_END:MOR-011:residence
% HMT_RESULT_BEGIN:MOR-012:residence
% HMT_ATOMIC_RESULT {"material_state":"INTEGRADO_O_REMITIDO_SIN_PERDIDA","proof_strength":"INTERFAZ_FISICA","provenance":"ARQUITECTURA_AUTORAL_PREEXISTENTE","reinforcement_state":"DEPLOYED_WITHOUT_NEW_RESULT_ID","residence":"`md/11m`","result_id":"MOR-012","statement":"retardo/adelanto y borde","status":"PRESENT_IN_SUCCESSOR_SOURCE"}
% HMT_RESULT_END:MOR-012:residence
% HMT_RESULT_BEGIN:MOR-013:residence
% HMT_ATOMIC_RESULT {"material_state":"INTEGRADO_O_REMITIDO_SIN_PERDIDA","proof_strength":"EXACTO_INTERNO","provenance":"RESULTADO_RECUPERADO","reinforcement_state":"DEPLOYED_WITHOUT_NEW_RESULT_ID","residence":"`md/04c`","result_id":"MOR-013","statement":"N_sp(uP_++vP_-)=uv","status":"PRESENT_IN_SUCCESSOR_SOURCE"}
% HMT_RESULT_END:MOR-013:residence
% HMT_RESULT_BEGIN:MOR-014:residence
% HMT_ATOMIC_RESULT {"material_state":"INTEGRADO_O_REMITIDO_SIN_PERDIDA","proof_strength":"INTERFAZ_FISICA","provenance":"RESULTADO_RECUPERADO","reinforcement_state":"DEPLOYED_WITHOUT_NEW_RESULT_ID","residence":"`md/04c`, `06ka`","result_id":"MOR-014","statement":"producto positivo/interior","status":"PRESENT_IN_SUCCESSOR_SOURCE"}
% HMT_RESULT_END:MOR-014:residence
% HMT_RESULT_BEGIN:MOR-015:residence
% HMT_ATOMIC_RESULT {"material_state":"INTEGRADO_O_REMITIDO_SIN_PERDIDA","proof_strength":"PROGRAMA_ABIERTO","provenance":"RESULTADO_NUEVO","reinforcement_state":"DEPLOYED_WITHOUT_NEW_RESULT_ID","residence":"`md/04ac` y persistencia","result_id":"MOR-015","statement":"gamma_+ es fibra de Hopf y gamma_- se proyecta con grado dos; sus clases (1,1) y (1,-1) dan las familias de Villarceau. La conjugación compleja total preserva cada familia no orientada y revierte su orientación: no intercambia ambas familias.","status":"PRESENT_IN_SUCCESSOR_SOURCE"}
% HMT_RESULT_END:MOR-015:residence
% HMT_RESULT_BEGIN:MOR-016:residence
% HMT_ATOMIC_RESULT {"material_state":"INTEGRADO_O_REMITIDO_SIN_PERDIDA","proof_strength":"PROGRAMA_ABIERTO","provenance":"RESULTADO_RECUPERADO","reinforcement_state":"DEPLOYED_WITHOUT_NEW_RESULT_ID","residence":"/Users/ruben/Documents/HMT2/EDICION_AUTONOMA_HMT_MD_2026-08-09_REV_3_EN_TRABAJO/01_FUENTE_REVISADA/manuscrito/sections/md/21_persistencia_ruta_mobius_familias_actualizada.tex","result_id":"MOR-016","statement":"El doble retorno 2pi/4pi y el fibrado real de línea L_chi con chi(1)=-1 son exactos; un círculo aislado no es Möbius. Ascender el signo C_M a la holonomía del lazo Villarceau requiere un único entrelazador físico.","status":"PRESENT_IN_SUCCESSOR_SOURCE"}
% HMT_RESULT_END:MOR-016:residence
% HMT_RESULT_BEGIN:MOR-017:residence
% HMT_ATOMIC_RESULT {"material_state":"INTEGRADO_O_REMITIDO_SIN_PERDIDA","proof_strength":"EXACTO_INTERNO","provenance":"FORMALIZACION_NUEVA","reinforcement_state":"DEPLOYED_WITHOUT_NEW_RESULT_ID","residence":"`md/04fa`","result_id":"MOR-017","statement":"En el escenario qutrit declarado, las POVM completas producen una distribución no señalizante y la cota local es I_3≤2. Para el estado máximamente entrelazado, I_3=4(3+2sqrt(3))/9=2.872934…; para |Psi_gamma>=(|00>+gamma|11>+|22>)/sqrt(2+gamma^2), gamma=(sqrt(11)-sqrt(3))/2, se alcanza I_3=1+sqrt(11/3)=2.914854…. Por tanto, máxima entropía no equivale a máxima violación. Derivar el estado y los instrumentos íntegramente desde el registro HMT sigue siendo una obligación separada.","status":"PRESENT_IN_SUCCESSOR_SOURCE"}
% HMT_RESULT_END:MOR-017:residence
% HMT_RESULT_BEGIN:MOR-018:residence
% HMT_ATOMIC_RESULT {"material_state":"INTEGRADO_O_REMITIDO_SIN_PERDIDA","proof_strength":"INTERFAZ_FISICA","provenance":"RESULTADO_RECUPERADO","reinforcement_state":"DEPLOYED_WITHOUT_NEW_RESULT_ID","residence":"`md/10_constitutivas_si`","result_id":"MOR-018","statement":"mu,epsilon,Z,c","status":"PRESENT_IN_SUCCESSOR_SOURCE"}
% HMT_RESULT_END:MOR-018:residence
% HMT_RESULT_BEGIN:MOR-019:residence
% HMT_ATOMIC_RESULT {"material_state":"INTEGRADO_O_REMITIDO_SIN_PERDIDA","proof_strength":"INTERFAZ_FISICA","provenance":"RESULTADO_RECUPERADO","reinforcement_state":"DEPLOYED_WITHOUT_NEW_RESULT_ID","residence":"Planck y paquete de masas","result_id":"MOR-019","statement":"Sigma_H y 10^-34","status":"PRESENT_IN_SUCCESSOR_SOURCE"}
% HMT_RESULT_END:MOR-019:residence
% HMT_RESULT_BEGIN:MOR-020:residence
% HMT_ATOMIC_RESULT {"material_state":"INTEGRADO_O_REMITIDO_SIN_PERDIDA","proof_strength":"INTERFAZ_FISICA","provenance":"RESULTADO_RECUPERADO","reinforcement_state":"DEPLOYED_WITHOUT_NEW_RESULT_ID","residence":"`md/05aa`, `05f`, `06m`","result_id":"MOR-020","statement":"operadores K_D/K_Omega y torres","status":"PRESENT_IN_SUCCESSOR_SOURCE"}
% HMT_RESULT_END:MOR-020:residence
% HMT_RESULT_BEGIN:MOR-021:residence
% HMT_ATOMIC_RESULT {"material_state":"INTEGRADO_O_REMITIDO_SIN_PERDIDA","proof_strength":"INTERFAZ_FISICA","provenance":"RESULTADO_RECUPERADO","reinforcement_state":"DEPLOYED_WITHOUT_NEW_RESULT_ID","residence":"`md/06k`, `06ka`, `06m`","result_id":"MOR-021","statement":"basal->beta (80,54,6)","status":"PRESENT_IN_SUCCESSOR_SOURCE"}
% HMT_RESULT_END:MOR-021:residence
% HMT_RESULT_BEGIN:MOR-022:residence
% HMT_ATOMIC_RESULT {"material_state":"INTEGRADO_O_REMITIDO_SIN_PERDIDA","proof_strength":"EXACTO_INTERNO","provenance":"RESULTADO_RECUPERADO","reinforcement_state":"DEPLOYED_WITHOUT_NEW_RESULT_ID","residence":"`md/03d`, con remisiones a `06k` y `06ka`","result_id":"MOR-022","statement":"Con T_108=108t_0, omega_108=2pi/T_108, R_108=c/omega_108 y m_108=hbar omega_108/c^2, se cumple hbar/(m_108 c)=R_108 y h/(m_108 c)=2pi R_108=C_108.","status":"PRESENT_IN_SUCCESSOR_SOURCE"}
% HMT_RESULT_END:MOR-022:residence
% HMT_RESULT_BEGIN:MOR-023:residence
% HMT_ATOMIC_RESULT {"material_state":"INTEGRADO_O_REMITIDO_SIN_PERDIDA","proof_strength":"INTERFAZ_FISICA","provenance":"RESULTADO_RECUPERADO","reinforcement_state":"DEPLOYED_WITHOUT_NEW_RESULT_ID","residence":"nuevo teorema producto + `md/10c`","result_id":"MOR-023","statement":"holonomía, coframe, conexión","status":"PRESENT_IN_SUCCESSOR_SOURCE"}
% HMT_RESULT_END:MOR-023:residence
% HMT_RESULT_BEGIN:MOR-024:residence
% HMT_ATOMIC_RESULT {"material_state":"INTEGRADO_O_REMITIDO_SIN_PERDIDA","proof_strength":"INTERFAZ_FISICA","provenance":"RESULTADO_RECUPERADO","reinforcement_state":"DEPLOYED_WITHOUT_NEW_RESULT_ID","residence":"`md/11`, `11c`, `11p`","result_id":"MOR-024","statement":"A=(omega,e), F=(R,T)","status":"PRESENT_IN_SUCCESSOR_SOURCE"}
% HMT_RESULT_END:MOR-024:residence
% HMT_RESULT_BEGIN:MOR-025:residence
% HMT_ATOMIC_RESULT {"material_state":"INTEGRADO_O_REMITIDO_SIN_PERDIDA","proof_strength":"INTERFAZ_FISICA","provenance":"RESULTADO_RECUPERADO","reinforcement_state":"DEPLOYED_WITHOUT_NEW_RESULT_ID","residence":"`md/11m`","result_id":"MOR-025","statement":"Si b es el lector de base, todo observable constante en las fibras de b factoriza de manera única a través de b. Una semilla s_0 se transporta consistentemente alrededor de un lazo exactamente cuando s_0 pertenece a Fix(Hol). El dato de fibra asciende al estado TPK completo sólo si la realización empleada es inyectiva.","status":"PRESENT_IN_SUCCESSOR_SOURCE"}
% HMT_RESULT_END:MOR-025:residence
% HMT_RESULT_BEGIN:MOR-026:residence
% HMT_ATOMIC_RESULT {"material_state":"INTEGRADO_O_REMITIDO_SIN_PERDIDA","proof_strength":"EXACTO_INTERNO","provenance":"RESULTADO_RECUPERADO","reinforcement_state":"DEPLOYED_WITHOUT_NEW_RESULT_ID","residence":"Euler/09b/14f","result_id":"MOR-026","statement":"J^2=-I y z->iz","status":"PRESENT_IN_SUCCESSOR_SOURCE"}
% HMT_RESULT_END:MOR-026:residence
% HMT_RESULT_BEGIN:MOR-027:residence
% HMT_ATOMIC_RESULT {"material_state":"INTEGRADO_O_REMITIDO_SIN_PERDIDA","proof_strength":"EXACTO_INTERNO","provenance":"FORMALIZACION_NUEVA","reinforcement_state":"DEPLOYED_WITHOUT_NEW_RESULT_ID","residence":"/Users/ruben/Documents/HMT2/EDICION_AUTONOMA_HMT_MD_2026-08-09_REV_3_EN_TRABAJO/01_FUENTE_REVISADA/manuscrito/sections/hmt/14d_ontologia_numeros_rev11.tex","result_id":"MOR-027","statement":"valor regularizado de zeta(-1)","status":"PRESENT_IN_SUCCESSOR_SOURCE"}
% HMT_RESULT_END:MOR-027:residence
% HMT_RESULT_BEGIN:MOR-028:residence
% HMT_ATOMIC_RESULT {"material_state":"INTEGRADO_O_REMITIDO_SIN_PERDIDA","proof_strength":"EXACTO_INTERNO","provenance":"ARQUITECTURA_AUTORAL_PREEXISTENTE","reinforcement_state":"DEPLOYED_WITHOUT_NEW_RESULT_ID","residence":"`hmt/03f`, `04_cubo_emrh`","result_id":"MOR-028","statement":"doble evaluación del antecedente","status":"PRESENT_IN_SUCCESSOR_SOURCE"}
% HMT_RESULT_END:MOR-028:residence
% HMT_RESULT_BEGIN:MOR-029:residence
% HMT_ATOMIC_RESULT {"material_state":"INTEGRADO_O_REMITIDO_SIN_PERDIDA","proof_strength":"EXACTO_INTERNO","provenance":"RESULTADO_RECUPERADO","reinforcement_state":"DEPLOYED_WITHOUT_NEW_RESULT_ID","residence":"`hmt/06aa--06f`","result_id":"MOR-029","statement":"Para todo K, la misma regla particiona exhaustivamente las 729 prolongaciones, selecciona una única superviviente y conmuta con las restricciones; los cilindros compatibles anidados determinan una sección coinductiva única.","status":"PRESENT_IN_SUCCESSOR_SOURCE"}
% HMT_RESULT_END:MOR-029:residence
% HMT_RESULT_BEGIN:MOR-030:residence
% HMT_ATOMIC_RESULT {"material_state":"INTEGRADO_O_REMITIDO_SIN_PERDIDA","proof_strength":"EXACTO_INTERNO","provenance":"FORMALIZACION_NUEVA","reinforcement_state":"DEPLOYED_WITHOUT_NEW_RESULT_ID","residence":"junto al teorema producto y síntesis final","result_id":"MOR-030","statement":"Los falsadores exactos son: continuo, pérdida de masa unitaria o de convergencia débil; Hilbert, ausencia de PVM compatible o dinámica no unitaria; morfismo discreto, fallo de composición; conjugación, C^2≠I, CJC^{-1}≠-J o CHC^{-1}≠H; absorbedor, falta de operador de onda o de condiciones causales; constitutivas, inversión que no intercambia mu/epsilon o cambia c sin extensión; masa, selección por PDG o promoción condicionada; electrón, imposibilidad de obtener carga, espín o estabilidad del mismo objeto; Planck, otra solución positiva autodual; Cartan, dependencia de subdivisión o confusión con Delta_4; Bell, señalización o instrumentos escogidos tras el blanco; filtración nonádica, consulta de prefijos objetivo o pérdida de una región de pi; Euler/zeta, regularización tratada como suma ordinaria; cosmología, incapacidad de reproducir simultáneamente expansión, lentes y estructura.","status":"PRESENT_IN_SUCCESSOR_SOURCE"}
% HMT_RESULT_END:MOR-030:residence
% HMT_RESULT_BEGIN:PLS-001:residence
% HMT_ATOMIC_RESULT {"material_state":"INTEGRADO_O_REMITIDO_SIN_PERDIDA","proof_strength":"EXACTO_INTERNO","provenance":"RESULTADO_RECUPERADO","reinforcement_state":"DEPLOYED_WITHOUT_NEW_RESULT_ID","residence":"`md/03d`, subsección CT108--Compton","result_id":"PLS-001","statement":"Para G(theta)=theta c_int^5 t_0^2/hbar_int, la clausura circular selecciona de manera única theta_*=(54/pi)^2 y hace coincidir el radio CT108 con las escalas de Compton, gravitatoria y de Planck en la convención declarada. Se distinguen r_g=Gm/c^2 y r_s=2r_g.","status":"PRESENT_IN_SUCCESSOR_SOURCE"}
% HMT_RESULT_END:PLS-001:residence
% HMT_RESULT_BEGIN:MAD-001:residence
% HMT_ATOMIC_RESULT {"material_state":"INTEGRADO_O_REMITIDO_SIN_PERDIDA","proof_strength":"EXACTO_INTERNO","provenance":"RESULTADO_RECUPERADO","reinforcement_state":"DEPLOYED_WITHOUT_NEW_RESULT_ID","residence":"`md/03d`, recta autodual de masas","result_id":"MAD-001","statement":"Con P_±=(I±j)/2, Z(m)=(m_P/m)P_++(m/m_P)P_- satisface N_sp(Z(m))=1; la conjugación implementa m↦m_P^2/m, el punto fijo positivo es m=m_P y x=log(m/m_P) es la rapidez de la recta.","status":"PRESENT_IN_SUCCESSOR_SOURCE"}
% HMT_RESULT_END:MAD-001:residence
% HMT_RESULT_BEGIN:AGR-001:residence
% HMT_ATOMIC_RESULT {"material_state":"INTEGRADO_O_REMITIDO_SIN_PERDIDA","proof_strength":"DEMOSTRADO_CON_ESTRUCTURA_DE_PARTIDA_EXPLICITA","provenance":"RESULTADO_RECUPERADO","reinforcement_state":"DEPLOYED_WITHOUT_NEW_RESULT_ID","residence":"`md/03d`, covariancia pre/retorno","result_id":"AGR-001","statement":"Las ramas pre y retorno transforman recíprocamente acción y acoplamiento gravitatorio, mantienen invariante su producto geométrico y conservan la longitud característica construida.","status":"PRESENT_IN_SUCCESSOR_SOURCE"}
% HMT_RESULT_END:AGR-001:residence
% HMT_RESULT_BEGIN:PLE-001:residence
% HMT_ATOMIC_RESULT {"material_state":"INTEGRADO_O_REMITIDO_SIN_PERDIDA","proof_strength":"INTERFAZ_FISICA","provenance":"RESULTADO_RECUPERADO","reinforcement_state":"DEPLOYED_WITHOUT_NEW_RESULT_ID","residence":"`md/03d`, `06k`, `06ka`","result_id":"PLE-001","statement":"El punto autodual de Planck fija el origen de la coordenada; la genealogía electrónica determina un desplazamiento posterior y la clausura escalar de rango uno. Este cierre discreto no demuestra por sí solo carga, espín, estabilidad, factor de forma, momento magnético ni un Hamiltoniano electromagnético ligado.","status":"PRESENT_IN_SUCCESSOR_SOURCE"}
% HMT_RESULT_END:PLE-001:residence
% HMT_RESULT_BEGIN:ANG-001:residence
% HMT_ATOMIC_RESULT {"material_state":"INTEGRADO_O_REMITIDO_SIN_PERDIDA","proof_strength":"INTERFAZ_FISICA","provenance":"ARQUITECTURA_AUTORAL_PREEXISTENTE","reinforcement_state":"DEPLOYED_WITHOUT_NEW_RESULT_ID","residence":"Hopf/Villarceau, después de MOR-016","result_id":"ANG-001","statement":"Para epsilon=+1, Spec(B_mem^+)={7,3}; resolver R+r=7ell y R-r=3ell da exactamente (R,r)=(5ell,2ell) y sin(beta)=2/5. La lectura como bordes de toro es interfaz declarada; el bloque total conserva por separado la razón 7:2.","status":"PRESENT_IN_SUCCESSOR_SOURCE"}
% HMT_RESULT_END:ANG-001:residence
% HMT_RESULT_BEGIN:BDR-001:residence
% HMT_ATOMIC_RESULT {"material_state":"INTEGRADO_O_REMITIDO_SIN_PERDIDA","proof_strength":"INTERFAZ_FISICA","provenance":"RESULTADO_RECUPERADO","reinforcement_state":"DEPLOYED_WITHOUT_NEW_RESULT_ID","residence":"constitutivas, con remisión desde bidireccionalidad","result_id":"BDR-001","statement":"La inversión conserva c_hat=exp(-E), invierte Z_hat=exp(B) e intercambia mu_hat y epsilon_hat; la realización vigente tiene velocidad común en ambos sentidos e impedancia orientada. Velocidades distintas exigen una extensión anisótropa adicional.","status":"PRESENT_IN_SUCCESSOR_SOURCE"}
% HMT_RESULT_END:BDR-001:residence
% HMT_RESULT_BEGIN:OCC-001:residence
% HMT_ATOMIC_RESULT {"material_state":"INTEGRADO_O_REMITIDO_SIN_PERDIDA","proof_strength":"EXACTO_INTERNO","provenance":"RESULTADO_RECUPERADO","reinforcement_state":"DEPLOYED_WITHOUT_NEW_RESULT_ID","residence":"Gibbs--KMS y ocupación","result_id":"OCC-001","statement":"Gobierna la identidad tipada de REV.3 F_{A,B}(t+i beta)=Tr(rho alpha_t(B)A), no la abreviación errónea del PDF de 68 pp; para 0<C<I, K_1=log((I-C)C^{-1}) y C=(I+e^{K_1})^{-1}, con tipos CAR/CCR separados.","status":"PRESENT_IN_SUCCESSOR_SOURCE"}
% HMT_RESULT_END:OCC-001:residence
% HMT_RESULT_BEGIN:BOL-001:residence
% HMT_ATOMIC_RESULT {"material_state":"INTEGRADO_O_REMITIDO_SIN_PERDIDA","proof_strength":"EXACTO_INTERNO","provenance":"RESULTADO_RECUPERADO","reinforcement_state":"DEPLOYED_WITHOUT_NEW_RESULT_ID","residence":"`md/04e`, `04g` y síntesis, con remisión CT108","result_id":"BOL-001","statement":"Boltzmann actúa como k_B:L_Theta→L_E y, para el reloj CT108, satisface k_B Theta_clk log(3)=hbar_int omega_0 con omega_0=2pi/(108t_0). La identidad fija el producto acción--información; insertar la coordenada SI declarada de k_B sólo evalúa Theta_clk y no constituye una derivación autónoma de ese número SI.","status":"PRESENT_IN_SUCCESSOR_SOURCE"}
% HMT_RESULT_END:BOL-001:residence
% HMT_RESULT_BEGIN:COS-001:residence
% HMT_ATOMIC_RESULT {"material_state":"INTEGRADO_O_REMITIDO_SIN_PERDIDA","proof_strength":"PROGRAMA_ABIERTO","provenance":"ARQUITECTURA_AUTORAL_PREEXISTENTE","reinforcement_state":"DEPLOYED_WITHOUT_NEW_RESULT_ID","residence":"perspectivas cosmológicas","result_id":"COS-001","statement":"La coexistencia de cartas permite un modelo estratificado, pero la aceleración sólo queda explicada tras derivar una ecuación efectiva con signo observado y límites de curvatura.","status":"PRESENT_IN_SUCCESSOR_SOURCE"}
% HMT_RESULT_END:COS-001:residence
% HMT_RESULT_BEGIN:GWV-001:residence
% HMT_ATOMIC_RESULT {"material_state":"INTEGRADO_O_REMITIDO_SIN_PERDIDA","proof_strength":"PROGRAMA_ABIERTO","provenance":"RESULTADO_RECUPERADO","reinforcement_state":"DEPLOYED_WITHOUT_NEW_RESULT_ID","residence":"gravedad sin gravitón y espectro","result_id":"GWV-001","statement":"Linealizar y reducir por calibre puede producir dos grados transversos de una onda gravitatoria colectiva; esto no obliga ni excluye un gravitón elemental, cuya existencia es una pregunta espectral adicional.","status":"PRESENT_IN_SUCCESSOR_SOURCE"}
% HMT_RESULT_END:GWV-001:residence
% HMT_RESULT_BEGIN:DRK-001:residence
% HMT_ATOMIC_RESULT {"material_state":"INTEGRADO_O_REMITIDO_SIN_PERDIDA","proof_strength":"PROGRAMA_ABIERTO","provenance":"ARQUITECTURA_AUTORAL_PREEXISTENTE","reinforcement_state":"DEPLOYED_WITHOUT_NEW_RESULT_ID","residence":"perspectivas cosmológicas","result_id":"DRK-001","statement":"Una alternativa estructural debe reproducir curvas de rotación, lentes, crecimiento, CMB y expansión con los mismos parámetros; Landauer nulo reversible no genera por sí solo un ciclo, que requiere dinámica de frontera acción--información.","status":"PRESENT_IN_SUCCESSOR_SOURCE"}
% HMT_RESULT_END:DRK-001:residence
% HMT_RESULT_BEGIN:CMK-001:residence
% HMT_ATOMIC_RESULT {"material_state":"INTEGRADO_O_REMITIDO_SIN_PERDIDA","proof_strength":"PROGRAMA_ABIERTO","provenance":"ARQUITECTURA_AUTORAL_PREEXISTENTE","reinforcement_state":"DEPLOYED_WITHOUT_NEW_RESULT_ID","residence":"cierre Euler/TPK","result_id":"CMK-001","statement":"El programa debe construir un diagrama de módulos y transformaciones naturales que tipen pares adjuntos o involutivos; una enumeración de operadores no constituye la coestructura.","status":"PRESENT_IN_SUCCESSOR_SOURCE"}
% HMT_RESULT_END:CMK-001:residence
% HMT_RESULT_BEGIN:MRC-001:residence
% HMT_ATOMIC_RESULT {"material_state":"INTEGRADO_O_REMITIDO_SIN_PERDIDA","proof_strength":"EXACTO_INTERNO","provenance":"CERTIFICADO_NUEVO","reinforcement_state":"DEPLOYED_WITHOUT_NEW_RESULT_ID","residence":"apéndice digital y gate de la monografía de realización","result_id":"MRC-001","statement":"Conserva diez familias: cilindros 0--1--infinito; espectros APP/memoria; identidad Perron--Fibonacci; inversión antiunitaria; producto semidirecto de Poincaré; constitutivas; norma partida y clausura material; Planck/autodualidad; CGLMP; e invariantes de filtración nonádica y masas.","status":"PRESENT_IN_SUCCESSOR_SOURCE"}
% HMT_RESULT_END:MRC-001:residence
% HMT_RESULT_BEGIN:PRG68-001:residence
% HMT_ATOMIC_RESULT {"material_state":"INTEGRADO_O_REMITIDO_SIN_PERDIDA","proof_strength":"PROGRAMA_ABIERTO","provenance":"ARQUITECTURA_AUTORAL_PREEXISTENTE","reinforcement_state":"DEPLOYED_WITHOUT_NEW_RESULT_ID","residence":"cierre del capítulo de realización física","result_id":"PRG68-001","statement":"Los seis cierres distintos son: límite suave Poincaré--Cartan independiente de subdivisión; un único operador C global que separe C, P y T; electrón ligado con carga, espín, estabilidad, factor de forma y momento magnético; mapa desde Delta_4 y la clausura electrónica hasta S_matter y la corriente de espín; estado e instrumentos HMT que deriven Bell/CGLMP sin usar el blanco; y selector general theta_G del transductor gravitatorio sin introducir la coordenada externa.","status":"PRESENT_IN_SUCCESSOR_SOURCE"}
% HMT_RESULT_END:PRG68-001:residence
% HMT_RESULT_BEGIN:GTR-001:residence
% HMT_ATOMIC_RESULT {"material_state":"INTEGRADO_O_REMITIDO_SIN_PERDIDA","proof_strength":"EXACTO_INTERNO","provenance":"RESULTADO_RECUPERADO","reinforcement_state":"DEPLOYED_WITHOUT_NEW_RESULT_ID","residence":"`md/11c`, `11n` y síntesis","result_id":"GTR-001","statement":"El mapa G_int(theta_G)=theta_G c_int^5 t_0^2/hbar_int tiene tipo dimensional exacto. En la clausura circular CT108 se especializa con theta_*=(54/pi)^2; una carta declarada evalúa 6.67430×10^-11. El selector general theta_G sigue siendo una obligación distinta y no se obtiene introduciendo esa coordenada externa.","status":"PRESENT_IN_SUCCESSOR_SOURCE"}
% HMT_RESULT_END:GTR-001:residence
% HMT_RESULT_BEGIN:AMB-001:residence
% HMT_ATOMIC_RESULT {"material_state":"CONSERVADO_Y_ATOMIZADO_EN_DELTA_SUCESOR","proof_strength":"DEMOSTRADO_CON_ESTRUCTURA_DE_PARTIDA_EXPLICITA","provenance":"RESULTADO_RECUPERADO","reinforcement_state":"DEPLOYED_WITHOUT_NEW_RESULT_ID","residence":"manuscrito/sections/hmt/05b_extensiones_dinamica_cociente.tex","result_id":"AMB-001","statement":"El cociente de modos del TPK induce la incidencia primitiva F_av; la medida de Parry del envolvente de memoria uno satisface mu_ar/mu_vol=phi_HMT^-2. El lector coinductivo marca una sola vez el retorno de frontera y produce pi_HMT/phi_HMT^2. Para R(x,y)=x/y^2 y J(x,y)=(y,x), las orientaciones conjugadas son pi_HMT/phi_HMT^2 y phi_HMT/pi_HMT^2; son funcionales coordinados y no intercambiables del mismo estado holonómico integral con memoria, sin factorización horizontal.","status":"PRESENT_IN_SUCCESSOR_SOURCE"}
% HMT_RESULT_END:AMB-001:residence
% HMT_RESULT_BEGIN:CHI-II-HDG-01:residence
% HMT_ATOMIC_RESULT {"material_state":"AUDITADO_PENDIENTE_INTEGRACION","proof_strength":"DEMOSTRADO_CON_ESTRUCTURA_DE_PARTIDA_EXPLICITA","provenance":"FORMALIZACION_NUEVA","reinforcement_state":"DEPLOYED_WITHOUT_NEW_RESULT_ID","residence":"manuscrito/sections/sucesora/02_estructura_discreta_continuo_20260817.tex:Hodge_macroobjetos","result_id":"CHI-II-HDG-01","statement":"La realización tipada R_Hdg:Xχ→Perf conserva la historia enriquecida del carácter χ y publica cada estado de Xχ como complejo perfecto, sin reconstruir el antecedente desde su sombra cohomológica.","status":"PRESENT_IN_SUCCESSOR_SOURCE"}
% HMT_RESULT_END:CHI-II-HDG-01:residence
% HMT_RESULT_BEGIN:CHI-II-HDG-02:residence
% HMT_ATOMIC_RESULT {"material_state":"AUDITADO_PENDIENTE_INTEGRACION","proof_strength":"DEMOSTRADO_CON_ESTRUCTURA_DE_PARTIDA_EXPLICITA","provenance":"FORMALIZACION_NUEVA","reinforcement_state":"DEPLOYED_WITHOUT_NEW_RESULT_ID","residence":"manuscrito/sections/sucesora/02_estructura_discreta_continuo_20260817.tex:Hodge_identidad_realizacion","result_id":"CHI-II-HDG-02","statement":"La identidad cl∘chloc∘R_Hdg=ev∞ conmuta exactamente en Xχ: la realización perfecta, su carácter de Chern localizado y la clase de ciclo publican la misma evaluación terminal.","status":"PRESENT_IN_SUCCESSOR_SOURCE"}
% HMT_RESULT_END:CHI-II-HDG-02:residence
% HMT_RESULT_BEGIN:CHI-II-HDG-03:residence
% HMT_ATOMIC_RESULT {"material_state":"INTEGRADO_EN_SUCESORA_PENDIENTE_DE_SELLO","proof_strength":"DEMOSTRADO_CON_ESTRUCTURA_DE_PARTIDA_EXPLICITA","provenance":"FORMALIZACION_NUEVA","reinforcement_state":"DEPLOYED_WITHOUT_NEW_RESULT_ID","residence":"manuscrito/sections/sucesora/02_estructura_discreta_continuo_20260817.tex:eq:continuo-hodge-completitud; manuscrito/sections/sucesora/04_hodge_bsd_fragment_20260817.tex:thm:hodge-suc-generacion","result_id":"CHI-II-HDG-03","statement":"Xχ cubre el dominio Hdg y, para toda clase α de ese dominio, existe ξ_α∈Xχ con β_α=chloc(R_Hdg(ξ_α)); por la identidad de realización, cl(β_α)=α.","status":"PRESENT_IN_SUCCESSOR_SOURCE"}
% HMT_RESULT_END:CHI-II-HDG-03:residence
% HMT_RESULT_BEGIN:CHI-IV-HDG-01:residence
% HMT_ATOMIC_RESULT {"material_state":"AUDITADO_PENDIENTE_INTEGRACION","proof_strength":"DEMOSTRADO_CON_ESTRUCTURA_DE_PARTIDA_EXPLICITA","provenance":"RESULTADO_NUEVO","reinforcement_state":"DEPLOYED_WITHOUT_NEW_RESULT_ID","residence":"manuscrito/sections/sucesora/04_cinco_realizaciones_fundamentales_20260817.tex:Hodge_APP_MV_coend","result_id":"CHI-IV-HDG-01","statement":"La celda APP–Mayer–Vietoris κ materializa suma, producto, diferencia y carry en Perf; el coend sobre historias compatibles totaliza las nueve fases coordinadas y conserva como terminal Cone(I−U_Γ9), con memoria visible.","status":"PRESENT_IN_SUCCESSOR_SOURCE"}
% HMT_RESULT_END:CHI-IV-HDG-01:residence
% HMT_RESULT_BEGIN:CHI-IV-HDG-C01:residence
% HMT_ATOMIC_RESULT {"material_state":"AUDITADO_PENDIENTE_INTEGRACION","proof_strength":"EXACTO_INTERNO","provenance":"RESULTADO_NUEVO","reinforcement_state":"DEPLOYED_WITHOUT_NEW_RESULT_ID","residence":"manuscrito/sections/sucesora/04_cinco_realizaciones_fundamentales_20260817.tex:Hodge_control_no_go_eliptico","result_id":"CHI-IV-HDG-C01","statement":"En una curva elíptica compleja E no existe un frame racional finito que cubra H²(E,Q(1)) y sea invariante bajo todas las traslaciones, ni un selector natural no basado de la clase de un punto; la obstrucción obliga a conservar la rigidificación y su descenso.","status":"PRESENT_IN_SUCCESSOR_SOURCE"}
% HMT_RESULT_END:CHI-IV-HDG-C01:residence
% HMT_RESULT_BEGIN:CHI-II-BSD-01:residence
% HMT_ATOMIC_RESULT {"material_state":"AUDITADO_PENDIENTE_INTEGRACION","proof_strength":"DEMOSTRADO_CON_ESTRUCTURA_DE_PARTIDA_EXPLICITA","provenance":"FORMALIZACION_NUEVA","reinforcement_state":"DEPLOYED_WITHOUT_NEW_RESULT_ID","residence":"manuscrito/sections/sucesora/02_estructura_discreta_continuo_20260817.tex:BSD_Gmn_AW","result_id":"CHI-II-BSD-01","statement":"El complejo G_mn de cadenas normalizadas del nervio de historias TPK supervivientes porta la diagonal Alexander–Whitney; ésta conserva la unidad genealógica y hace group-like cada vértice de estado completo.","status":"PRESENT_IN_SUCCESSOR_SOURCE"}
% HMT_RESULT_END:CHI-II-BSD-01:residence
% HMT_RESULT_BEGIN:CHI-II-BSD-02:residence
% HMT_ATOMIC_RESULT {"material_state":"AUDITADO_PENDIENTE_INTEGRACION","proof_strength":"DEMOSTRADO_CON_ESTRUCTURA_DE_PARTIDA_EXPLICITA","provenance":"RESULTADO_NUEVO","reinforcement_state":"DEPLOYED_WITHOUT_NEW_RESULT_ID","residence":"manuscrito/sections/sucesora/02_estructura_discreta_continuo_20260817.tex:BSD_Manin_producto_fibrado","result_id":"CHI-II-BSD-02","statement":"El blanco conservativo BSD es Manin×_A G: el producto fibrado acopla el complejo Manin basado con G sobre el aumento A y conserva simultáneamente la proyección modular y la historia profunda.","status":"PRESENT_IN_SUCCESSOR_SOURCE"}
% HMT_RESULT_END:CHI-II-BSD-02:residence
% HMT_RESULT_BEGIN:CHI-II-BSD-03:residence
% HMT_ATOMIC_RESULT {"material_state":"AUDITADO_PENDIENTE_INTEGRACION","proof_strength":"DEMOSTRADO_CON_ESTRUCTURA_DE_PARTIDA_EXPLICITA","provenance":"RESULTADO_NUEVO","reinforcement_state":"DEPLOYED_WITHOUT_NEW_RESULT_ID","residence":"manuscrito/sections/sucesora/02_estructura_discreta_continuo_20260817.tex:BSD_publicacion_acoplada_K_PT","result_id":"CHI-II-BSD-03","statement":"Las publicaciones Kato y Poitou–Tate se construyen sobre Manin×_A G y publican la misma unidad genealógica de G; la comparación conserva raíz y complemento sin factorizar por el mero aumento.","status":"PRESENT_IN_SUCCESSOR_SOURCE"}
% HMT_RESULT_END:CHI-II-BSD-03:residence
% HMT_RESULT_BEGIN:CHI-IV-BSD-01:residence
% HMT_ATOMIC_RESULT {"material_state":"AUDITADO_PENDIENTE_INTEGRACION","proof_strength":"DEMOSTRADO_CON_ESTRUCTURA_DE_PARTIDA_EXPLICITA","provenance":"RESULTADO_NUEVO","reinforcement_state":"DEPLOYED_WITHOUT_NEW_RESULT_ID","residence":"manuscrito/sections/sucesora/04_cinco_realizaciones_fundamentales_20260817.tex:BSD_raiz_dualidad_HFS","result_id":"CHI-IV-BSD-01","statement":"La reducción ortogonal y la dualidad chain-level fijan J_red; la orientación λ elimina el twist y fuerza p_root zK=zPT, mientras H_FS construye h* para contraer el complemento finita–singular sin alterar la raíz.","status":"PRESENT_IN_SUCCESSOR_SOURCE"}
% HMT_RESULT_END:CHI-IV-BSD-01:residence
% HMT_RESULT_BEGIN:CHI-IV-BSD-02:residence
% HMT_ATOMIC_RESULT {"material_state":"INTEGRADO_EN_SUCESORA_PENDIENTE_DE_SELLO","proof_strength":"DEMOSTRADO_CON_ESTRUCTURA_DE_PARTIDA_EXPLICITA","provenance":"FORMALIZACION_NUEVA","reinforcement_state":"DEPLOYED_WITHOUT_NEW_RESULT_ID","residence":"manuscrito/sections/sucesora/04_hodge_bsd_fragment_20260817.tex:thm:bsd-suc-bockstein,thm:bsd-suc-local-global,eq:bsd-suc-formula-final","result_id":"CHI-IV-BSD-02","statement":"El regulador derivado R_Boc^der coincide como elemento basado con el término inicial LT_T(S); unido a la igualdad raíz Kato/PT, la altura, Sha, Tamagawa, periodo y torsión ya tipados, publica la fórmula BSD final sin definir un lado por el otro.","status":"PRESENT_IN_SUCCESSOR_SOURCE"}
% HMT_RESULT_END:CHI-IV-BSD-02:residence
% HMT_RESULT_BEGIN:CHI-IV-BSD-C01:residence
% HMT_ATOMIC_RESULT {"material_state":"AUDITADO_PENDIENTE_INTEGRACION","proof_strength":"EXACTO_INTERNO","provenance":"RESULTADO_NUEVO","reinforcement_state":"DEPLOYED_WITHOUT_NEW_RESULT_ID","residence":"manuscrito/sections/sucesora/04_cinco_realizaciones_fundamentales_20260817.tex:BSD_control_twist4","result_id":"CHI-IV-BSD-C01","statement":"La ablación del pullback genealógico o la inserción del twist raíz 4 conserva sombras modulares parciales pero destruye la unidad basada: zK^(4)=4r+3e+b pasa módulo 3 y falla desde módulo 9, por lo que no puede satisfacer p_root zK=zPT.","status":"PRESENT_IN_SUCCESSOR_SOURCE"}
% HMT_RESULT_END:CHI-IV-BSD-C01:residence
% HMT_RESULT_BEGIN:CHI-II-RH-01:residence
% HMT_ATOMIC_RESULT {"material_state":"AUDITADO_PENDIENTE_INTEGRACION","proof_strength":"EXACTO_INTERNO","provenance":"FORMALIZACION_NUEVA","reinforcement_state":"DEPLOYED_WITHOUT_NEW_RESULT_ID","residence":"manuscrito/sections/sucesora/02_estructura_discreta_continuo_20260817.tex:RH_macroobjeto","result_id":"CHI-II-RH-01","statement":"La Parte II construye M_RH=(Xhat_infinito,Gamma9,Rhat=(R,K),Carr_R,S0,M,O9,Xi,P,D_R_cof): Gamma9 conserva la fase y eleva la profundidad; Carr_R compone como cociclo; S0=sum_{q=0}^{N-1}M^{*q}O9^*O9M^q+M^{*N}S0M^N conserva el terminal; y B_W=Xi^*(I-P)Xi publica el complemento de nivel/fibra sin borrar memoria.","status":"PRESENT_IN_SUCCESSOR_SOURCE"}
% HMT_RESULT_END:CHI-II-RH-01:residence
% HMT_RESULT_BEGIN:CHI-IV-RH-01:residence
% HMT_ATOMIC_RESULT {"material_state":"AUDITADO_PENDIENTE_INTEGRACION","proof_strength":"DEMOSTRADO_CON_ESTRUCTURA_DE_PARTIDA_EXPLICITA","provenance":"RESULTADO_RECUPERADO","reinforcement_state":"DEPLOYED_WITHOUT_NEW_RESULT_ID","residence":"manuscrito/sections/sucesora/04_cinco_realizaciones_fundamentales_20260817.tex:RH_teorema","result_id":"CHI-IV-RH-01","statement":"Con q=5/9, M9=Phi A8···A0=qV9 y la identidad finita de memoria dan B_W,R=L_R^*L_R≥0 al extinguirse su terminal; la realización común produce B_W=Xi^*(I-P)Xi≥0 y, sobre D_R_cof con los canales completos, el criterio reversible publica Re(rho)=1/2 para todo cero no trivial.","status":"PRESENT_IN_SUCCESSOR_SOURCE"}
% HMT_RESULT_END:CHI-IV-RH-01:residence
% HMT_RESULT_BEGIN:CHI-IV-RH-C01:residence
% HMT_ATOMIC_RESULT {"material_state":"AUDITADO_PENDIENTE_INTEGRACION","proof_strength":"EXACTO_INTERNO","provenance":"RESULTADO_RECUPERADO","reinforcement_state":"DEPLOYED_WITHOUT_NEW_RESULT_ID","residence":"manuscrito/sections/sucesora/04_cinco_realizaciones_fundamentales_20260817.tex:RH_control_ablacion","result_id":"CHI-IV-RH-C01","statement":"El control retira por separado ortogonalidad, antecedente común, canal primo/gamma/polar, cofinalidad, carry, memoria o terminal; cada ablación rompe una igualdad anterior al criterio y no puede presentarse como una variante del teorema.","status":"PRESENT_IN_SUCCESSOR_SOURCE"}
% HMT_RESULT_END:CHI-IV-RH-C01:residence
% HMT_RESULT_BEGIN:CHI-II-NS-01:residence
% HMT_ATOMIC_RESULT {"material_state":"AUDITADO_PENDIENTE_INTEGRACION","proof_strength":"EXACTO_INTERNO","provenance":"FORMALIZACION_NUEVA","reinforcement_state":"DEPLOYED_WITHOUT_NEW_RESULT_ID","residence":"manuscrito/sections/sucesora/02_estructura_discreta_continuo_20260817.tex:NS_macroobjeto","result_id":"CHI-II-NS-01","statement":"La Parte II construye M_NS,m=(Xi_NS,m,C_{m←n},M_m^+,M_m^-,D_m,H_m,E9,J9,P9,Q_micro,R_m,E_m,ell_m): C_{m←n}=P_mB(u_n,u_n)-B_m(P_m u_n,P_m u_n) conserva el carry; E9J9=I, P9=J9E9 y Q_micro=I-P9 tipan la fibra; E9||Xi_{m+1}||²+E9 ell_m=||Xi_m||² conserva el terminal y archiva memoria positiva.","status":"PRESENT_IN_SUCCESSOR_SOURCE"}
% HMT_RESULT_END:CHI-II-NS-01:residence
% HMT_RESULT_BEGIN:CHI-IV-NS-01:residence
% HMT_ATOMIC_RESULT {"material_state":"AUDITADO_PENDIENTE_INTEGRACION","proof_strength":"EXACTO_INTERNO","provenance":"RESULTADO_RECUPERADO","reinforcement_state":"DEPLOYED_WITHOUT_NEW_RESULT_ID","residence":"manuscrito/sections/sucesora/04_cinco_realizaciones_fundamentales_20260817.tex:NS_teorema","result_id":"CHI-IV-NS-01","statement":"El puerto del macroobjeto satisface [B_{s,m}+(1-eta)nu D_{s,m}]^-≤C_NS ell_m+b_{parcial,m}; al componerlo con la telescopía y la energía se obtiene sup_m||u_m||_{L∞H^s}²+sup_m||u_m||_{L²H^{s+1}}²≤C_T(u0,f,Xi0) para s>5/2, y se publican existencia global, unicidad y suavidad para t>0 en el dominio declarado.","status":"PRESENT_IN_SUCCESSOR_SOURCE"}
% HMT_RESULT_END:CHI-IV-NS-01:residence
% HMT_RESULT_BEGIN:CHI-IV-NS-C01:residence
% HMT_ATOMIC_RESULT {"material_state":"AUDITADO_PENDIENTE_INTEGRACION","proof_strength":"EXACTO_INTERNO","provenance":"RESULTADO_RECUPERADO","reinforcement_state":"DEPLOYED_WITHOUT_NEW_RESULT_ID","residence":"manuscrito/sections/sucesora/04_cinco_realizaciones_fundamentales_20260817.tex:NS_control_ablacion","result_id":"CHI-IV-NS-C01","statement":"El control elimina por separado retorno conjunto, carry de Galerkin, memoria positiva, Q_micro, terminal o uniformidad del puerto; una rama de norma 3^J, C_NS dependiente del corte o b_parcial,m no sumable impiden la telescopía y rechazan la publicación global.","status":"PRESENT_IN_SUCCESSOR_SOURCE"}
% HMT_RESULT_END:CHI-IV-NS-C01:residence
% HMT_RESULT_BEGIN:CHI-II-YM-01:residence
% HMT_ATOMIC_RESULT {"material_state":"AUDITADO_PENDIENTE_INTEGRACION","proof_strength":"EXACTO_INTERNO","provenance":"FORMALIZACION_NUEVA","reinforcement_state":"DEPLOYED_WITHOUT_NEW_RESULT_ID","residence":"manuscrito/sections/sucesora/02_estructura_discreta_continuo_20260817.tex:YM_macroobjeto","result_id":"CHI-II-YM-01","statement":"La Parte II construye M_YM,m=(Lambda_m,mu_m,T_m,J_m,Theta_{nu,m},K_t,D_m^YM,P_Omega_m,Q_m^phys): Lambda_m=a_m Z^4 y K_{a_{m+1}}^9=K_{a_m}; (T_{m+1})^9J_m=J_mT_m y Theta_{nu,m+1}J_m=J_mTheta_{nu,m} conservan nivel/fibra, carry de fusión y memoria; P_Omega_m es el terminal raíz.","status":"PRESENT_IN_SUCCESSOR_SOURCE"}
% HMT_RESULT_END:CHI-II-YM-01:residence
% HMT_RESULT_BEGIN:CHI-IV-YM-01:residence
% HMT_ATOMIC_RESULT {"material_state":"AUDITADO_PENDIENTE_INTEGRACION","proof_strength":"EXACTO_INTERNO","provenance":"RESULTADO_RECUPERADO","reinforcement_state":"DEPLOYED_WITHOUT_NEW_RESULT_ID","residence":"manuscrito/sections/sucesora/04_cinco_realizaciones_fundamentales_20260817.tex:YM_teorema","result_id":"CHI-IV-YM-01","statement":"Sobre el macroobjeto común, omega_m(overline{F(Theta_nu,m U)}F(U))=||B_nu,mF||²≥0 para nu=1,…,4; el generador D_m^YM satisface ker D_m^YM=C Omega_m y D_m^YM≥3kappa_av(I-P_Omega_m); q_{m+1}^9=q_m y -a_m^{-1}log q_m=3kappa_av publican una teoría local, vacío simple y brecha Delta_YM^HMT≥3kappa_av>0 en el sector colectivo gauge-invariante.","status":"PRESENT_IN_SUCCESSOR_SOURCE"}
% HMT_RESULT_END:CHI-IV-YM-01:residence
% HMT_RESULT_BEGIN:CHI-IV-YM-C01:residence
% HMT_ATOMIC_RESULT {"material_state":"AUDITADO_PENDIENTE_INTEGRACION","proof_strength":"EXACTO_INTERNO","provenance":"RESULTADO_RECUPERADO","reinforcement_state":"DEPLOYED_WITHOUT_NEW_RESULT_ID","residence":"manuscrito/sections/sucesora/04_cinco_realizaciones_fundamentales_20260817.tex:YM_control_ablacion","result_id":"CHI-IV-YM-C01","statement":"El control elimina por separado semigrupo común, medida compartida, localidad uniforme, entrelazamiento de refinamiento, carry, memoria, vacío terminal o escala física; cada ablación rompe una identidad previa a la brecha y no puede repararse usando Wilson como selector retrospectivo.","status":"PRESENT_IN_SUCCESSOR_SOURCE"}
% HMT_RESULT_END:CHI-IV-YM-C01:residence
% HMT_RESULT_BEGIN:CHI-I-TPK-01:residence
% HMT_ATOMIC_RESULT {"material_state":"INTEGRADO_EN_LIENZO_SUCESOR_NO_ACTIVADO","proof_strength":"EXACTO_INTERNO","provenance":"FORMALIZACION_NUEVA","reinforcement_state":"DEPLOYED_WITHOUT_NEW_RESULT_ID","residence":"manuscrito/sections/sucesora/01_nucleo_formal_hmt_20260817.tex:conexion_nonadica_conjunta","result_id":"CHI-I-TPK-01","statement":"La holonomía única Γ9=Hol_C9(∇_TPK) conserva el retorno de fase con evolución del estado; en una realización isométrica satisface U_Γ9 J=JT+η, J*η=0 e I−T*T=η*η, y la iteración anti-especular conserva la suma de memoria junto con el terminal finito positivo.","status":"PRESENT_IN_SUCCESSOR_SOURCE"}
% HMT_RESULT_END:CHI-I-TPK-01:residence
% HMT_RESULT_BEGIN:CHI-I-EMRH-01:residence
% HMT_ATOMIC_RESULT {"material_state":"INTEGRADO_EN_LIENZO_SUCESOR_NO_ACTIVADO","proof_strength":"EXACTO_INTERNO","provenance":"RESULTADO_RECUPERADO","reinforcement_state":"DEPLOYED_WITHOUT_NEW_RESULT_ID","residence":"manuscrito/sections/sucesora/01_nucleo_formal_hmt_20260817.tex:completacion_cubica_conservacion_proyectiva","result_id":"CHI-I-EMRH-01","statement":"Para C̄9^(d)=(E9⊔{⋆})^d, el estrato con r coordenadas de frontera tiene cardinal B_r(d)=binom(d,r)9^(d−r), la proyección que olvida la última coordenada satisface (π_d)_*μ_d=μ_{d−1} y la masa total permanece uno; en dimensión tres, 1000=729+243+27+1=729+271.","status":"PRESENT_IN_SUCCESSOR_SOURCE"}
% HMT_RESULT_END:CHI-I-EMRH-01:residence
% HMT_RESULT_BEGIN:CHI-I-FAV-01:residence
% HMT_ATOMIC_RESULT {"material_state":"INTEGRADO_EN_LIENZO_SUCESOR_NO_ACTIVADO","proof_strength":"DEMOSTRADO_CON_ESTRUCTURA_DE_PARTIDA_EXPLICITA","provenance":"FORMALIZACION_NUEVA","reinforcement_state":"DEPLOYED_WITHOUT_NEW_RESULT_ID","residence":"manuscrito/sections/sucesora/01_nucleo_formal_hmt_20260817.tex:cociente_areal_volumetrico","result_id":"CHI-I-FAV-01","statement":"El bloque primitivo (+1,−1,0) induce la órbita de modos (Σ,Π,Π); el lector tipado de hojas transporta su incidencia a F_av=[[0,1],[1,1]], y la repetición del calendario produce N9=3F_av, N27=9F_av y N108=36F_av.","status":"PRESENT_IN_SUCCESSOR_SOURCE"}
% HMT_RESULT_END:CHI-I-FAV-01:residence
% HMT_RESULT_BEGIN:CHI-III-VAC-01:residence
% HMT_ATOMIC_RESULT {"material_state":"AUDITADO_PENDIENTE_INTEGRACION","proof_strength":"EXACTO_INTERNO","provenance":"FORMALIZACION_NUEVA","reinforcement_state":"DEPLOYED_WITHOUT_NEW_RESULT_ID","residence":"manuscrito/sections/sucesora/03_genealogia_constantes_20260817.tex:descomposicion_Fitting_vacancias","result_id":"CHI-III-VAC-01","statement":"En T9=(Z/9Z)^2, M=[[3,4],[6,7]] satisface M²−M=3I; E=M³ es idempotente y N=M−E cumple N²=0 y EN=NE, de modo que T9=im(E)⊕ker(E) y la filtración de imágenes tiene cardinales 81→27→9→9.","status":"PRESENT_IN_SUCCESSOR_SOURCE"}
% HMT_RESULT_END:CHI-III-VAC-01:residence
% HMT_RESULT_BEGIN:CHI-III-MASS-01:residence
% HMT_ATOMIC_RESULT {"material_state":"AUDITADO_PENDIENTE_INTEGRACION","proof_strength":"DEMOSTRADO_CON_ESTRUCTURA_DE_PARTIDA_EXPLICITA","provenance":"FORMALIZACION_NUEVA","reinforcement_state":"DEPLOYED_WITHOUT_NEW_RESULT_ID","residence":"manuscrito/sections/sucesora/03_genealogia_constantes_20260817.tex:dominio_completo_ley_general_masas","result_id":"CHI-III-MASS-01","statement":"El dominio interno vinculante de la Ley General de Masas comprende 81 celdas, 56 familias de ruta, 13 multisecciones, 324 rutas orientadas y un sector nulo separado; cada extensión superviviente publica ruta, registro, firma, carácter y operador de fibra antes de toda carta metrológica. Las antiguas ventanas de diez o doce masas son sólo testigos históricos.","status":"PRESENT_IN_SUCCESSOR_SOURCE"}
% HMT_RESULT_END:CHI-III-MASS-01:residence
% HMT_RESULT_BEGIN:CHI-III-ELECTRON-01:residence
% HMT_ATOMIC_RESULT {"material_state":"AUDITADO_PENDIENTE_INTEGRACION","proof_strength":"DEMOSTRADO_CON_ESTRUCTURA_DE_PARTIDA_EXPLICITA","provenance":"FORMALIZACION_NUEVA","reinforcement_state":"DEPLOYED_WITHOUT_NEW_RESULT_ID","residence":"manuscrito/sections/sucesora/03_genealogia_constantes_20260817.tex:electron_central_two_way","result_id":"CHI-III-ELECTRON-01","statement":"La fibra electrónica ocupa el único centro (5,5) con ruta Gamma_e=(5,5,++); sus lectores independientes coinciden en K_D(Gamma_e)=K_Omega(Gamma_e)=(80,54,6), fijan kappa_e=80−54alpha_HMT+6Delta_4 y, después del basal y de la recta de acción, publican m_e^beta=0.5109989506900008086082571648… MeV. El basal permanece como etapa, no como salida rectora.","status":"PRESENT_IN_SUCCESSOR_SOURCE"}
% HMT_RESULT_END:CHI-III-ELECTRON-01:residence
% HMT_RESULT_BEGIN:CHI-III-BOLTZ-01:residence
% HMT_ATOMIC_RESULT {"material_state":"AUDITADO_PENDIENTE_INTEGRACION","proof_strength":"EXACTO_INTERNO","provenance":"FORMALIZACION_NUEVA","reinforcement_state":"DEPLOYED_WITHOUT_NEW_RESULT_ID","residence":"manuscrito/sections/sucesora/03_genealogia_constantes_20260817.tex:boltzmann_accion_informacion","result_id":"CHI-III-BOLTZ-01","statement":"La constante de Boltzmann se tipa como el isomorfismo positivo k_B:L_Theta→L_E; una decisión TRIT uniforme aporta I_TRIT=log 3 y la periodicidad de orden 108 fija k_B Theta_clk log 3=hbar_int 2pi/(108 t0). La construcción determina internamente la energía por unidad de información; la carta kelvin–joule publica después k_B=1.380649×10^−23 J K^−1.","status":"PRESENT_IN_SUCCESSOR_SOURCE"}
% HMT_RESULT_END:CHI-III-BOLTZ-01:residence
% HMT_RESULT_BEGIN:CHI-III-PLANCK108-01:residence
% HMT_ATOMIC_RESULT {"material_state":"AUDITADO_PENDIENTE_INTEGRACION","proof_strength":"DEMOSTRADO_CON_ESTRUCTURA_DE_PARTIDA_EXPLICITA","provenance":"FORMALIZACION_NUEVA","reinforcement_state":"DEPLOYED_WITHOUT_NEW_RESULT_ID","residence":"manuscrito/sections/sucesora/03_genealogia_constantes_20260817.tex:planck_periodicidad_108","result_id":"CHI-III-PLANCK108-01","statement":"Las cartas de acción hbar_pre y hbar_ret admiten una realización circular de período T_108=108t0, con omega_108=2pi/T_108, R_108=c_int/omega_108 y C_108=108 ell0; para m_108,a=hbar_a omega_108/c_int² se cumplen barlambda_C=R_108 y lambda_C=C_108. La condición de radio gravitatorio reducido selecciona theta_108=(54/pi)² y coordina las unidades internas de Planck en ambas cartas.","status":"PRESENT_IN_SUCCESSOR_SOURCE"}
% HMT_RESULT_END:CHI-III-PLANCK108-01:residence
% HMT_RESULT_BEGIN:CHI-VI-GRAV-ALG-01:residence
% HMT_ATOMIC_RESULT {"material_state":"INTEGRADO_EN_LIENZO_SUCESOR_NO_ACTIVADO","proof_strength":"EXACTO_INTERNO","provenance":"FORMALIZACION_NUEVA","reinforcement_state":"DEPLOYED_WITHOUT_NEW_RESULT_ID","residence":"manuscrito/sections/sucesora/06_mecanica_dimensional_20260817.tex:sector_tensorial_icosaedrico_y_proyeccion_TT","result_id":"CHI-VI-GRAV-ALG-01","statement":"La representación icosaédrica de A5 sobre V12 posee el proyector ortogonal P5 de rango cinco; la igualdad de caracteres identifica V5=im(P5) con Sym²_0(R³) mediante una isometría A5-equivariante única hasta signo, y el proyector transversal sin traza Lambda(n) tiene rango dos. La cadena exacta es 12→5→5→2 y no se confunde con una realización gravitatoria dinámica.","status":"PRESENT_IN_SUCCESSOR_SOURCE"}
% HMT_RESULT_END:CHI-VI-GRAV-ALG-01:residence
% HMT_RESULT_BEGIN:AMP-I-RECTA-01:residence
% HMT_ATOMIC_RESULT {"material_state":"AMPLIFICACION_AUTORIZADA_2026_08_17","proof_strength":"EXACTO_INTERNO","provenance":"FORMALIZACION_NUEVA","reinforcement_state":"DEPLOYED_WITHOUT_NEW_RESULT_ID","residence":"manuscrito/sections/ampliacion_20260817/00_recta_horizonte_completacion.tex","result_id":"AMP-I-RECTA-01","statement":"La sucesión de cartas [0,10]→[0,1000]→[0,10^m] construye un sistema compatible de marcas, intervalos y fronteras cuya evaluación final vive en R; el quinto postulado euclídeo se trata como antecedente histórico de cierre geométrico y la completación se formula mediante mapas y límites, no como metáfora.","status":"PRESENT_IN_AMPLIFIED_SOURCE"}
% HMT_RESULT_END:AMP-I-RECTA-01:residence
% HMT_RESULT_BEGIN:AMP-I-UNIDAD-01:residence
% HMT_ATOMIC_RESULT {"material_state":"AMPLIFICACION_AUTORIZADA_2026_08_17","proof_strength":"EXACTO_INTERNO","provenance":"FORMALIZACION_NUEVA","reinforcement_state":"DEPLOYED_WITHOUT_NEW_RESULT_ID","residence":"manuscrito/sections/ampliacion_20260817/00_recta_horizonte_completacion.tex","result_id":"AMP-I-UNIDAD-01","statement":"La unidad normalizada se conserva bajo refinamiento, elevación dimensional y doble publicación; en dimensión d la estratificación binomial de (E9⊔{★})^d suma 10^d y los push-forward proyectivos conservan masa uno.","status":"PRESENT_IN_AMPLIFIED_SOURCE"}
% HMT_RESULT_END:AMP-I-UNIDAD-01:residence
% HMT_RESULT_BEGIN:AMP-II-EXC-01:residence
% HMT_ATOMIC_RESULT {"material_state":"AMPLIFICACION_AUTORIZADA_2026_08_17","proof_strength":"DEMOSTRADO_CON_ESTRUCTURA_DE_PARTIDA_EXPLICITA","provenance":"RESULTADO_RECUPERADO","reinforcement_state":"DEPLOYED_WITHOUT_NEW_RESULT_ID","residence":"manuscrito/sections/sucesora/02_estructura_discreta_continuo_20260817.tex; manuscrito/sections/sucesora/06_mecanica_dimensional_20260817.tex","result_id":"AMP-II-EXC-01","statement":"La rama excepcional se bifurca tipadamente: C_W→W12→W24 por diseños, y C_W→N(A2^12)→Λ24→V_Λ→V^natural→Monster→Moonshine por retículo y orbifold; no existe una serialización Witt→Golay→A2^12.","status":"PRESENT_IN_AMPLIFIED_SOURCE"}
% HMT_RESULT_END:AMP-II-EXC-01:residence
% HMT_RESULT_BEGIN:AMP-III-TIPOS-01:residence
% HMT_ATOMIC_RESULT {"material_state":"AMPLIFICACION_AUTORIZADA_2026_08_17","proof_strength":"DEMOSTRADO_CON_ESTRUCTURA_DE_PARTIDA_EXPLICITA","provenance":"RESULTADO_RECUPERADO","reinforcement_state":"DEPLOYED_WITHOUT_NEW_RESULT_ID","residence":"manuscrito/sections/sucesora/03_genealogia_constantes_20260817.tex; manuscrito/sections/sucesora/06_mecanica_dimensional_20260817.tex","result_id":"AMP-III-TIPOS-01","statement":"G se expone en tres capas —transductor exacto, selector circular de orden 108 y criterio holonómico general—; Boltzmann deriva internamente k_B Θ_clk log3 y sólo después abre la carta SI; raíz de -1 permanece en Parte II y -1/12 en Parte III.","status":"PRESENT_IN_AMPLIFIED_SOURCE"}
% HMT_RESULT_END:AMP-III-TIPOS-01:residence
% HMT_RESULT_BEGIN:AMP-IV-REAL-01:residence
% HMT_ATOMIC_RESULT {"material_state":"AMPLIFICACION_AUTORIZADA_2026_08_17","proof_strength":"DEMOSTRADO_CON_ESTRUCTURA_DE_PARTIDA_EXPLICITA","provenance":"FORMALIZACION_NUEVA","reinforcement_state":"DEPLOYED_WITHOUT_NEW_RESULT_ID","residence":"manuscrito/sections/ampliacion_20260817/04_principio_realizacion_y_supersesion.tex","result_id":"AMP-IV-REAL-01","statement":"Una familia F_N:S_N^surv→C_N compatible con truncación, orientación, transporte, acarreo, unidad y terminal induce F_∞ por la propiedad universal del límite inverso y transporta conjuntamente observable y memoria.","status":"PRESENT_IN_AMPLIFIED_SOURCE"}
% HMT_RESULT_END:AMP-IV-REAL-01:residence
% HMT_RESULT_BEGIN:AMP-IV-RHNS-01:residence
% HMT_ATOMIC_RESULT {"material_state":"AMPLIFICACION_AUTORIZADA_2026_08_17","proof_strength":"DEMOSTRADO_CON_ESTRUCTURA_DE_PARTIDA_EXPLICITA","provenance":"RESULTADO_RECUPERADO","reinforcement_state":"DEPLOYED_WITHOUT_NEW_RESULT_ID","residence":"manuscrito/sections/ampliacion_20260817/04a_refuerzo_riemann_weil.tex; manuscrito/sections/ampliacion_20260817/04b_refuerzo_navier_stokes.tex","result_id":"AMP-IV-RHNS-01","statement":"RH incorpora el residual finito 729=81+648 y R_m,R=0 antes del criterio; NS incorpora cota conjunta, ∂_t u en L2_loc H^{s-1}, pegado por unicidad y presión normalizada, sin sustituir sus demostraciones vigentes.","status":"PRESENT_IN_AMPLIFIED_SOURCE"}
% HMT_RESULT_END:AMP-IV-RHNS-01:residence
% HMT_RESULT_BEGIN:AMP-IV-YM-01:residence
% HMT_ATOMIC_RESULT {"material_state":"AMPLIFICACION_AUTORIZADA_2026_08_17","proof_strength":"DEMOSTRADO_CON_ESTRUCTURA_DE_PARTIDA_EXPLICITA","provenance":"RESULTADO_RECUPERADO","reinforcement_state":"DEPLOYED_WITHOUT_NEW_RESULT_ID","residence":"manuscrito/sections/ampliacion_20260817/04c_refuerzo_yang_mills.tex","result_id":"AMP-IV-YM-01","statement":"La familia proyectiva común produce cuatro formas OS, un Hamiltoniano común, acción fuerte de E(4), funcionales de Schwinger, publicación en H^{-s}_loc(R4), s>2, lectores clover y F², vacío simple y brecha exacta 3κ_av.","status":"PRESENT_IN_AMPLIFIED_SOURCE"}
% HMT_RESULT_END:AMP-IV-YM-01:residence
% HMT_RESULT_BEGIN:AMP-IV-HDG-01:residence
% HMT_ATOMIC_RESULT {"material_state":"AMPLIFICACION_AUTORIZADA_2026_08_17","proof_strength":"DEMOSTRADO_CON_ESTRUCTURA_DE_PARTIDA_EXPLICITA","provenance":"FORMALIZACION_NUEVA","reinforcement_state":"DEPLOYED_WITHOUT_NEW_RESULT_ID","residence":"manuscrito/sections/ampliacion_20260817/04d_hodge_completitud_desplegada.tex","result_id":"AMP-IV-HDG-01","statement":"La completitud de Xχ(X,p) se despliega como teorema coinductivo de realizaciones finitas compatibles; produce primero ξ_α y sólo después β_α=ch_p^loc(R_Hdg(ξ_α)), preservando el cuadrado de publicación.","status":"PRESENT_IN_AMPLIFIED_SOURCE"}
% HMT_RESULT_END:AMP-IV-HDG-01:residence
% HMT_RESULT_BEGIN:AMP-IV-BSD-01:residence
% HMT_ATOMIC_RESULT {"material_state":"AMPLIFICACION_AUTORIZADA_2026_08_17","proof_strength":"DEMOSTRADO_CON_ESTRUCTURA_DE_PARTIDA_EXPLICITA","provenance":"FORMALIZACION_NUEVA","reinforcement_state":"DEPLOYED_WITHOUT_NEW_RESULT_ID","residence":"manuscrito/sections/ampliacion_20260817/04e_bsd_comparadores_desplegados.tex","result_id":"AMP-IV-BSD-01","statement":"Los comparadores Kato–FERS/Iwasawa, Poitou–Tate/Poincaré–Néron–Tate, Smith–Tamagawa–Shafarevich, torsión y período se construyen y componen sobre una única línea orientada; la igualdad raíz elimina la unidad relativa y publica rango, finitud de Sha y fórmula final.","status":"PRESENT_IN_AMPLIFIED_SOURCE"}
% HMT_RESULT_END:AMP-IV-BSD-01:residence
% HMT_RESULT_BEGIN:AMP-IV-SUP-01:residence
% HMT_ATOMIC_RESULT {"material_state":"AMPLIFICACION_AUTORIZADA_2026_08_17","proof_strength":"EXACTO_INTERNO","provenance":"RESULTADO_RECUPERADO","reinforcement_state":"DEPLOYED_WITHOUT_NEW_RESULT_ID","residence":"manuscrito/sections/ampliacion_20260817/04_principio_realizacion_y_supersesion.tex; manuscrito/appendices_acumulativa_revfinal.tex","result_id":"AMP-IV-SUP-01","statement":"Los rótulos históricos que declaraban abiertos RH, NS, YM, Hodge y BSD quedan supersedidos en esta edición por los cinco teoremas de Parte IV; se conservan únicamente como trazabilidad de formulaciones anteriores y controles tipados, sin autoridad sobre el corte vigente.","status":"PRESENT_IN_AMPLIFIED_SOURCE"}
% HMT_RESULT_END:AMP-IV-SUP-01:residence
% HMT_RESULT_BEGIN:AMP-V-GRAV5-01:residence
% HMT_ATOMIC_RESULT {"material_state":"AMPLIFICACION_AUTORIZADA_2026_08_17","proof_strength":"EXACTO_INTERNO","provenance":"FORMALIZACION_NUEVA","reinforcement_state":"DEPLOYED_WITHOUT_NEW_RESULT_ID","residence":"manuscrito/sections/ampliacion_20260817/05_sector_tensorial_gravitatorio_cinco_componentes.tex","result_id":"AMP-V-GRAV5-01","statement":"La acción icosaédrica dodecafásica induce V12≅V1⊕V3⊕V3'⊕V5; el proyector central P5, el entrelazador explícito Γ5 hacia Sym0²(R3), la base STF y la reducción TT construyen la cadena 12→5→5→2 y separan portador quíntuple, espín dos masivo, GR y clases E(2).","status":"PRESENT_IN_AMPLIFIED_SOURCE"}
% HMT_RESULT_END:AMP-V-GRAV5-01:residence
% HMT_RESULT_BEGIN:AMP-V-GRAVPHY-01:residence
% HMT_ATOMIC_RESULT {"material_state":"AMPLIFICACION_AUTORIZADA_2026_08_17","proof_strength":"INTERFAZ_FISICA","provenance":"FORMALIZACION_NUEVA","reinforcement_state":"DEPLOYED_WITHOUT_NEW_RESULT_ID","residence":"manuscrito/sections/ampliacion_20260817/05_sector_tensorial_gravitatorio_cinco_componentes.tex","result_id":"AMP-V-GRAVPHY-01","statement":"El transductor G, la realización circular de orden 108 y el sector tensorial algebraico conservan sus conclusiones exactas; la realización PCH se expone como cadena tipada que determina el espacio físico mediante ecuaciones, restricciones y calibre, sin atribuir cinco polarizaciones a Einstein, Feynman ni GR.","status":"PRESENT_IN_AMPLIFIED_SOURCE"}
% HMT_RESULT_END:AMP-V-GRAVPHY-01:residence
% HMT_RESULT_BEGIN:AMP-EPI-R-01:residence
% HMT_ATOMIC_RESULT {"material_state":"AMPLIFICACION_AUTORIZADA_2026_08_17","proof_strength":"DEMOSTRADO_CON_ESTRUCTURA_DE_PARTIDA_EXPLICITA","provenance":"FORMALIZACION_NUEVA","reinforcement_state":"DEPLOYED_WITHOUT_NEW_RESULT_ID","residence":"manuscrito/sections/ampliacion_20260817/99_epilogo_retorno_recta_real.tex","result_id":"AMP-EPI-R-01","statement":"El epílogo retorna a R y demuestra que la recta real es la publicación arquimediana de una sección compatible, mientras la publicación excepcional conserva incidencia; ambas proceden del mismo estado y ninguna reconstruye por sí sola su genealogía.","status":"PRESENT_IN_AMPLIFIED_SOURCE"}
% HMT_RESULT_END:AMP-EPI-R-01:residence
 % Sidecar sucesor de la reparación semántica aprobada: 32 residencias
% rectoras, comentario-only y enlazadas al nuevo recibo de precompilación.
% Sidecar probatorio de la reparación semántica 2026-08-19.
% No produce salida tipográfica; cada intervalo liga una residencia rectora.
% HMT_REPAIR_RESULT_BEGIN::HMT-CORE-APP-TRIT-TPK
% RESULT_ID: HMT-CORE-APP-TRIT-TPK
% STATUS: DEMONSTRATED
% RESIDENCE: manuscrito/sections/sucesora/01_nucleo_formal_hmt_20260817.tex#eq:nucleo-app-dos-hojas
% STATEMENT_SHA256: 6657284d9bd2918464c9d964efdaae3608a94c55cd238f3de467fe85c00ad0ac
% STATEMENT: APP construye X_9; TRIT tipa sus dos hojas; después se forman los dominios de semilla, las emisiones y las palabras visibles, y sólo entonces se produce el estado holonómico integral que conserva hoja, orientación, residuo, cociente, acarreo, firma, cilindro, frontera, ruta, memoria, multisección y supervivencia.
% HMT_REPAIR_RESULT_END::HMT-CORE-APP-TRIT-TPK
% HMT_REPAIR_RESULT_BEGIN::HMT-CORE-NONADIC-JOINT-CONTINUITY
% RESULT_ID: HMT-CORE-NONADIC-JOINT-CONTINUITY
% STATUS: DEMONSTRATED
% RESIDENCE: manuscrito/sections/sucesora/01_nucleo_formal_hmt_20260817.tex#thm:nucleo-continuidad-conexion-nonadica-conjunta
% STATEMENT_SHA256: bea2a7fd24c115638fea7eb3682c82f9b684160dc0a4e855abb04024c2e484c4
% STATEMENT: Γ_9 devuelve la fase tras nueve pasos sin reinicializar el estado; conserva memoria, acarreo y terminal, y sus potencias son composiciones relacionales compatibles con truncamientos.
% HMT_REPAIR_RESULT_END::HMT-CORE-NONADIC-JOINT-CONTINUITY
% HMT_REPAIR_RESULT_BEGIN::HMT-RECTOR-CONTINUUM-CLOSURE-CONSTANTS
% RESULT_ID: HMT-RECTOR-CONTINUUM-CLOSURE-CONSTANTS
% STATUS: DEMONSTRATED
% RESIDENCE: manuscrito/sections/sucesora/02_estructura_discreta_continuo_20260817.tex#eq:continuo-teorema-fundamental-discreto
% STATEMENT_SHA256: 34c506d5e3311116e4a5f5d4baf1febf44fb1fcbd22e17a6540197f17479fbd7
% STATEMENT: HMT cierra el problema del continuo en el dominio genealógico: las publicaciones continuas son lectores compatibles de secciones discretas completas; las constantes emergen del mismo antecedente, y los cinco problemas del Milenio son especializaciones concretas y subordinadas del cierre rector.
% HMT_REPAIR_RESULT_END::HMT-RECTOR-CONTINUUM-CLOSURE-CONSTANTS
% HMT_REPAIR_RESULT_BEGIN::HMT-CORE-NO-ABLATION
% RESULT_ID: HMT-CORE-NO-ABLATION
% STATUS: DEMONSTRATED
% RESIDENCE: gestion/reparacion_semantica_20260818/MAPA_RECONSTRUCCION_CAUSAL_FUENTE_20260818.md §7.1
% STATEMENT_SHA256: 6a8421d9f55a806faa959c19578d710ae27e286555775c9a1f34551268c7e592
% STATEMENT: El mapa de olvido U retiene una sombra visible, conjuntista, ordinal, escalar o convencional y elimina genealogía, memoria, orientación, cociente/residuo, acarreo, frontera/ruta, supervivencia o multisección; ninguna objeción a esa sombra se transfiere al objeto completo sin probar que U conserva y refleja la propiedad examinada.
% HMT_REPAIR_RESULT_END::HMT-CORE-NO-ABLATION
% HMT_REPAIR_RESULT_BEGIN::HMT-ZFC-VON-NEUMANN-NONFAITHFUL-CH
% RESULT_ID: HMT-ZFC-VON-NEUMANN-NONFAITHFUL-CH
% STATUS: DEMONSTRATED
% RESIDENCE: manuscrito/sections/integracion_20260812/von_neumann_64/04_zfc_fundacion_ordinales.tex#thm:zfc-perdida-genealogia-olvido
% STATEMENT_SHA256: 9fcf0f78fc08d6066f7bf7f9f8aeb86d0a5e87348f2efb1affef488e061e0b9b
% STATEMENT: Tres resultados distintos quedan tipados: U_60 colapsa dos objetos no isomorfos con la misma Sig_60 y fibras Ext_81 de cardinales 0 y 12, por lo que no refleja isomorfías ni clasifica; V no es fiel porque id_ρ y Γ_9 tienen la misma sombra unipuntual pero grados de acarreo/memoria (0,0) y (1,1); los ordinales finitos con memorias k y k+1 proyectan al mismo ordinal de von Neumann. CH es un blanco distinto.
% HMT_REPAIR_RESULT_END::HMT-ZFC-VON-NEUMANN-NONFAITHFUL-CH
% HMT_REPAIR_RESULT_BEGIN::HMT-CORE-DOUBLE-PUBLICATION
% RESULT_ID: HMT-CORE-DOUBLE-PUBLICATION
% STATUS: DEMONSTRATED
% RESIDENCE: manuscrito/sections/sucesora/02_estructura_discreta_continuo_20260817.tex#eq:continuo-doble-publicacion
% STATEMENT_SHA256: 3051702377a6942a887f59b9323a35081f9d1ba17e9b3e8a72326b7039ab0e43
% STATEMENT: El lector arquimediano contrae cilindros compatibles y el lector incidencial conserva la genealogía de frontera; sus codominios son distintos. La rama excepcional es bifurcada: Hadamard organiza U_12^sgn en paralelo y Paley produce C_W, desde donde parten las ramas Witt–Golay–Mathieu y Niemeier–Leech–V^natural–Monstruo–Moonshine.
% HMT_REPAIR_RESULT_END::HMT-CORE-DOUBLE-PUBLICATION
% HMT_REPAIR_RESULT_BEGIN::HMT-MILLENNIUM-RH-WEIL-CLOSED
% RESULT_ID: HMT-MILLENNIUM-RH-WEIL-CLOSED
% STATUS: DEMONSTRATED
% RESIDENCE: manuscrito/sections/sucesora/04_cinco_realizaciones_fundamentales_20260817.tex#thm:realizaciones-rh
% STATEMENT_SHA256: a4f088786afdbee666309a8fc5034b3722fade65c9f6a259d05da2b5732626b5
% STATEMENT: En cada región R se fija q=5/9 y se construye E_R=j_-Q_RXi_R con AE_R=qE_R y E_R*E_R=B_{W,R}; la telescopía finita conserva el terminal q^(2N)B_{W,R}, cuya extinción sólo se toma al pasar N→∞, y L_R satisface L_R*L_R=B_{W,R}⪰0. El descenso cofinal publica la forma de Weil global como cuadrado positivo y el criterio de Weil reconoce la hipótesis de Riemann.
% HMT_REPAIR_RESULT_END::HMT-MILLENNIUM-RH-WEIL-CLOSED
% HMT_REPAIR_RESULT_BEGIN::HMT-MILLENNIUM-NS-CLOSED
% RESULT_ID: HMT-MILLENNIUM-NS-CLOSED
% STATUS: DEMONSTRATED
% RESIDENCE: manuscrito/sections/sucesora/04_cinco_realizaciones_fundamentales_20260817.tex#thm:realizaciones-ns
% STATEMENT_SHA256: f8e920750843d201de3d006e4d249eb58a18849ab02f1db10372401a049235db
% STATEMENT: La rama propietaria PC–Fock–KYP construye prospectivamente el Gram del terminal y las pérdidas futuras, deriva la identidad Stein y la telescopía con terminal, factoriza directamente el bloque KYP, establece balance, observabilidad y registro de torre, y produce el presupuesto raíz uniforme y la cota H^s global. El comparador G_{M,j}=P_aux*U^term_{M,j}P_aux es sólo un control no gobernante.
% HMT_REPAIR_RESULT_END::HMT-MILLENNIUM-NS-CLOSED
% HMT_REPAIR_RESULT_BEGIN::HMT-MILLENNIUM-YM-CLOSED
% RESULT_ID: HMT-MILLENNIUM-YM-CLOSED
% STATUS: DEMONSTRATED
% RESIDENCE: manuscrito/sections/sucesora/04_cinco_realizaciones_fundamentales_20260817.tex#thm:realizaciones-ym
% STATEMENT_SHA256: 7bae1f46e683874449e24e2e5e76323c00b6953291777330a09a1c2324fc5167
% STATEMENT: La misma familia proyectiva construye medida, red local, generador D_m^YM, semigrupo K_{m,t}=e^(−tD_m^YM), vacío, cuatro formas OS, acción fuerte E(4) y brecha exacta. Se cumplen K_{m,t}P_{Ω_m}=P_{Ω_m} y ||K_{m,t}(I−P_{Ω_m})||≤e^(−3κ_av t); el testigo W_c alcanza la cota. El lector clover/F² reconoce después la curvatura del mismo proceso.
% HMT_REPAIR_RESULT_END::HMT-MILLENNIUM-YM-CLOSED
% HMT_REPAIR_RESULT_BEGIN::HMT-MILLENNIUM-HODGE-CLOSED
% RESULT_ID: HMT-MILLENNIUM-HODGE-CLOSED
% STATUS: DEMONSTRATED
% RESIDENCE: manuscrito/sections/sucesora/04_hodge_bsd_fragment_20260817.tex#thm:hodge-suc-generacion
% STATEMENT_SHA256: 871872c44580d0e280556a6d45a7a549870a5b324223da5310a7a61a0ee0f72b
% STATEMENT: La realización gobernante parte de X_χ=lim_K Surv_K(X,p), conmuta con el carácter local y la publicación de clases y, para cada α, construye ξ_α y β_α con cl(β_α)=α. Los núcleos relativos y la presentación coend son controles de autonomía expositiva, no la cadena gobernante.
% HMT_REPAIR_RESULT_END::HMT-MILLENNIUM-HODGE-CLOSED
% HMT_REPAIR_RESULT_BEGIN::HMT-MILLENNIUM-BSD-CLOSED
% RESULT_ID: HMT-MILLENNIUM-BSD-CLOSED
% STATUS: DEMONSTRATED
% RESIDENCE: manuscrito/sections/sucesora/04_hodge_bsd_fragment_20260817.tex#eq:bsd-suc-formula-final
% STATEMENT_SHA256: 58f17d0b3189c35e80a258c376c281f796aee2740fea34308df90316cbe709c7
% STATEMENT: La coálgebra Alexander–Whitney y su counidad construyen G_{m,n} y el producto fibrado P_E^gen; una misma unidad genera las publicaciones Kato/PT, la igualdad de raíz, el relleno h_*, la bandera de Bockstein, los comparadores y la igualdad basada que cierra BSD. La unipotencia generador a generador es sólo control de despliegue.
% HMT_REPAIR_RESULT_END::HMT-MILLENNIUM-BSD-CLOSED
% HMT_REPAIR_RESULT_BEGIN::HMT-EXCEPTIONAL-CORRIDOR
% RESULT_ID: HMT-EXCEPTIONAL-CORRIDOR
% STATUS: DEMONSTRATED
% RESIDENCE: manuscrito/sections/hmt/09c_genealogia_excepcional_completa_autonoma.tex#sec:genealogia-excepcional-completa-autonoma
% STATEMENT_SHA256: 9e1f0885079d7e94e3bbf4142ac1fef08314e4d1d537ffc3403e23bd7d3bb08b
% STATEMENT: Hadamard organiza U_12^sgn sin convertirse en antecedente serial de Paley; la cadena C_P→A_W→C_W bifurca desde C_W hacia Witt–Golay–Mathieu y hacia Niemeier–Leech–V_Λ–V^natural; después Aut(V^natural)=Monstruo y J=T_{1A} precede a la familia {T_g} de Moonshine.
% HMT_REPAIR_RESULT_END::HMT-EXCEPTIONAL-CORRIDOR
% HMT_REPAIR_RESULT_BEGIN::HMT-INTERNAL-I-QUARTER-TURN
% RESULT_ID: HMT-INTERNAL-I-QUARTER-TURN
% STATUS: DEMONSTRATED
% RESIDENCE: manuscrito/sections/hmt/09a_paley_codigo_witt.tex#eq:AW-raiz-menos-uno
% STATEMENT_SHA256: 89af0b10c8bc5da185fb92a05117661adee2496d45d963a0922f2c33d9cf247e
% STATEMENT: La matriz incidencial A_W satisface A_W²=−I en el sector Paley–Witt; ésta es la raíz finita interna. La estructura J_E²=−I y su lector exponencial publican después la realización compleja arquimediana, sin identificar ambos operadores por nombre.
% HMT_REPAIR_RESULT_END::HMT-INTERNAL-I-QUARTER-TURN
% HMT_REPAIR_RESULT_BEGIN::HMT-INTERNAL-MINUS-ONE-TWELFTH
% RESULT_ID: HMT-INTERNAL-MINUS-ONE-TWELFTH
% STATUS: DEMONSTRATED
% RESIDENCE: manuscrito/sections/sucesora/03_genealogia_constantes_20260817.tex#eq:constantes-menos-un-doce-proyectores
% STATEMENT_SHA256: d04beccf9787ff880f26f73a4ff543adb5aab378144657d59b5d156e5fd65284
% STATEMENT: El valor −1/12 emerge en seis dominios tipados: incidencia 90/120, diferencia de periodos, Gram positivo con pseudoinversa, extractor de escala, exponente eta y traza de proyectores del corrimiento cíclico S_12:Q^12→Q^12, con S_12^12=I, donde D_r=I−S_12^r y los proyectores sobre ker D_3 y ker D_4 tienen rangos 3 y 4. La igualdad ζ(−1)=−1/12 es un reconocimiento analítico separado, no el origen de esas realizaciones.
% HMT_REPAIR_RESULT_END::HMT-INTERNAL-MINUS-ONE-TWELFTH
% HMT_REPAIR_RESULT_BEGIN::HMT-LORENTZ-MINKOWSKI
% RESULT_ID: HMT-LORENTZ-MINKOWSKI
% STATUS: DEMONSTRATED
% RESIDENCE: manuscrito/sections/hmt/03_app_ortograma.tex#prop:app-trit-espectral-markov-lorentz-rev8
% STATEMENT_SHA256: 7da5585d1b7752e0cd76fc580631d618b0aedef2c4bc51e326b4cc8bae1649e5
% STATEMENT: El bloque central B_c de APP produce M_c=B_c/√45∈SO^+(1,1) y satisface M_c^Tη_{1,1}M_c=η_{1,1}; ésta es una realización exacta de Lorentz y del plano de Minkowski 1+1. Toda elevación 3+1 requiere un mapa dimensional tipado y no demueve el resultado 1+1.
% HMT_REPAIR_RESULT_END::HMT-LORENTZ-MINKOWSKI
% HMT_REPAIR_RESULT_BEGIN::HMT-HILBERT-REALIZATION
% RESULT_ID: HMT-HILBERT-REALIZATION
% STATUS: DEMONSTRATED
% RESIDENCE: manuscrito/sections/integracion_20260812/von_neumann_64/06_hilbert_heisenberg.tex
% STATEMENT_SHA256: 5323eb3e082071f29482590ea4ede54b91c92314e2f7d6f865ddf5cde5fea80d
% STATEMENT: El espacio de Hilbert matricial, el límite UHF y su cierre GNS, la dinámica compleja y Born se construyen por ramas tipadas que no se implican automáticamente: C_9^(d) produce H_d,A_d,UHF(9^∞), su traza τ produce la representación GNS (π_τ,H_τ,Ω_τ) y π_τ(A_(9^∞))ʺ=R_hyp; J_E,H∈B(H_τ), con J_E²=−I y H=H*, producen U(t); el lector Born GNS usa P=P*=P²∈A_UHF actuando por π_τ(P) y p_τ(P)=<Ω_τ,π_τ(P)Ω_τ>; una segunda realización finita conserva Pr_ρ(a)=Tr(ρE_a) desde una PVM y un estado normalizado.
% HMT_REPAIR_RESULT_END::HMT-HILBERT-REALIZATION
% HMT_REPAIR_RESULT_BEGIN::HMT-T-DUALITY-M-SCREENS
% RESULT_ID: HMT-T-DUALITY-M-SCREENS
% STATUS: DEMONSTRATED
% RESIDENCE: manuscrito/sections/hmt/16b_pantallas_dualidad_teoria_m.tex#thm:dualidad-visible-memoria
% STATEMENT_SHA256: a6c3111b2029fd7c5e93e357ea6b65a0cd583f22589de276bcbf2a3215fc5d31
% STATEMENT: En D_dual=Z²×R_(>0), T(m,w,R_0)=(w,m,R_0^(−1)) satisface T²=id y preserva E_(R_0)(m,w)=m²/R_0²+w²R_0² mediante E_(R_0)(m,w)=E_(R_0^(−1))(w,m); su realización unitaria convive con las reducciones 11→10 y 26→24. La interfaz física separada R_M:X_HMT^enr→X_11^phys debe producir y compatibilizar g_(MN)∈Γ(Sym² T*X_11), C_(MNP)∈Ω³(X_11), ψ_M∈Γ(S_11⊗T*X_11), su acción y transformaciones, sin contaminar el teorema interno.
% HMT_REPAIR_RESULT_END::HMT-T-DUALITY-M-SCREENS
% HMT_REPAIR_RESULT_BEGIN::HMT-CONSTANT-PI
% RESULT_ID: HMT-CONSTANT-PI
% STATUS: DEMONSTRATED
% RESIDENCE: manuscrito/sections/ampliacion_20260818/02_generacion_infinita_pi_e_phi.tex#eq:generacion-infinita-gaussiana-pi-20260818
% STATEMENT_SHA256: 26b11a14165eacab3aa7d06fcb90f31ff0f2eb0dd328b3f7c11cf183fee1293f
% STATEMENT: La órbita y el carácter 5/4/239 se seleccionan antes de abrir intervalos reales; la identidad de Machin reconoce posteriormente el lector ya construido.
% HMT_REPAIR_RESULT_END::HMT-CONSTANT-PI
% HMT_REPAIR_RESULT_BEGIN::HMT-CONSTANT-E
% RESULT_ID: HMT-CONSTANT-E
% STATUS: DEMONSTRATED
% RESIDENCE: manuscrito/sections/ampliacion_20260818/02_generacion_infinita_pi_e_phi.tex#eq:generacion-infinita-recurrencia-e-20260818
% STATEMENT_SHA256: 2c21447e726bae153a60dc20b39c536b01f8b4dba8836c1176eeb252af969ec0
% STATEMENT: La ley multiplicativa del canal fuerza la recurrencia a_n=1/n! antes de reconocer la serie exponencial y su valor e.
% HMT_REPAIR_RESULT_END::HMT-CONSTANT-E
% HMT_REPAIR_RESULT_BEGIN::HMT-CONSTANT-PHI
% RESULT_ID: HMT-CONSTANT-PHI
% STATUS: DEMONSTRATED
% RESIDENCE: capítulo autónomo de la constante de autoescala, ecuación de prolongación 20260818
% STATEMENT_SHA256: 8151c9f0c4164daffee20672e7f6751ed9452b96c233d2ecd5e81aedca948338
% STATEMENT: La incidencia areal–volumétrica induce F_av y su rayo de Perron fuerza x²−x−1=0; la raíz positiva se reconoce después como φ.
% HMT_REPAIR_RESULT_END::HMT-CONSTANT-PHI
% HMT_REPAIR_RESULT_BEGIN::HMT-CONSTANT-ALPHA
% RESULT_ID: HMT-CONSTANT-ALPHA
% STATUS: DEMONSTRATED
% RESIDENCE: manuscrito/sections/sucesora/03_genealogia_constantes_20260817.tex#thm:constantes-alpha-independencia-principal
% STATEMENT_SHA256: 8eb80629b88f75eaf8f624a0e6992e43812a576d8f34b530137934ff6bf0d7a6
% STATEMENT: En el perfil HMT completo, B_term y U_12^sgn son proyecciones paralelas del mismo estado terminal; R_sgn extrae U, H_12 invertible da K, y el acarreo y la recurrencia producen α sin α, CODATA ni masa como entrada.
% HMT_REPAIR_RESULT_END::HMT-CONSTANT-ALPHA
% HMT_REPAIR_RESULT_BEGIN::HMT-EULER-HOLOGRAPHIC
% RESULT_ID: HMT-EULER-HOLOGRAPHIC
% STATUS: DEMONSTRATED
% RESIDENCE: manuscrito/sections/ampliacion_20260818/03_extension_holografica_euler.tex#thm:euler-holografica-tres-cierres-20260818
% STATEMENT_SHA256: f54720a289461a9177197a7184922777c9bb186d4617980d744cae299400c025
% STATEMENT: A_W²=−I proporciona la raíz finita incidencial; en una rama distinta, J_E²=−I y el lector exponencial producen la identidad arquimediana de Euler, extendida por J_τ a tres ramas. El contraángulo no nace de Euler: sólo después de α_HMT, la microfibra 169, D_A, H_5 y la recurrencia regional generan C*, cuya identidad inversa se reconoce al final.
% HMT_REPAIR_RESULT_END::HMT-EULER-HOLOGRAPHIC
% HMT_REPAIR_RESULT_BEGIN::HMT-EMRH-CUBIC-UNIT
% RESULT_ID: HMT-EMRH-CUBIC-UNIT
% STATUS: DEMONSTRATED
% RESIDENCE: manuscrito/sections/ampliacion_20260818/01_emrh_desarrollo_integral.tex#eq:emrh-cierre-apice-20260818
% STATEMENT_SHA256: 6f630c70e8d9bf2ba536e9a19b4f59ce7f4f90709e8c1f3410e2cf49cfab5724
% STATEMENT: La completación correcta es 10^3=729+271, con 271 como corona E y 729 como bloque alfa; además (1/271+1/729)^−1=271·729/1000 expresa exactamente la razón inversa asociada a la completación volumétrica.
% HMT_REPAIR_RESULT_END::HMT-EMRH-CUBIC-UNIT
% HMT_REPAIR_RESULT_BEGIN::HMT-CONSTANTS-FOUR-STRATA
% RESULT_ID: HMT-CONSTANTS-FOUR-STRATA
% STATUS: DEMONSTRATED
% RESIDENCE: manuscrito/sections/ampliacion_20260818/05_constantes_fisicas_genealogia.tex#thm:amp18-constantes-cuatro-estratos
% STATEMENT_SHA256: f67dfab70770a2b7aed78a5925e7f56252e00166de0b0e1ce805a0be3520e083
% STATEMENT: Toda constante se presenta en el orden constructor discreto/genealógico, caracterización analítica, transducción dimensional equivariante y comparación metrológica; la carta SI no construye el invariante.
% HMT_REPAIR_RESULT_END::HMT-CONSTANTS-FOUR-STRATA
% HMT_REPAIR_RESULT_BEGIN::HMT-CONSTANT-PLANCK
% RESULT_ID: HMT-CONSTANT-PLANCK
% STATUS: DEMONSTRATED
% RESIDENCE: manuscrito/sections/md/03b_escala_accion_determinantal.tex#thm:orden-accion
% STATEMENT_SHA256: 0bc899b4d06add106a366bf912d47203626d8fde1c6911f1073da88cc17eeef0
% STATEMENT: El lector determinantal Σ_H=φ_HMT α_HMT^16 y su cokernel local de orden 16 fijan la recta de acción y la década −34; la mantisa y la carta J·s se distinguen, y cualquier residual permanece visible.
% HMT_REPAIR_RESULT_END::HMT-CONSTANT-PLANCK
% HMT_REPAIR_RESULT_BEGIN::HMT-CONSTANT-BOLTZMANN
% RESULT_ID: HMT-CONSTANT-BOLTZMANN
% STATUS: DEMONSTRATED
% RESIDENCE: manuscrito/sections/ampliacion_20260818/05_constantes_fisicas_genealogia.tex#prop:amp18-constantes-covariancia-termica
% STATEMENT_SHA256: c61e51204132bcfb86f480ea2592cef625c3aa9bdbbce4632ba8c5995cc47828
% STATEMENT: El invariante interno fija k_B Θ_clk log 3=ħ_int 2π/(108t_0); la covariancia Θ→λΘ, k_B→k_B/λ demuestra que el kelvin fija la carta, no el objeto.
% HMT_REPAIR_RESULT_END::HMT-CONSTANT-BOLTZMANN
% HMT_REPAIR_RESULT_BEGIN::HMT-CONSTITUTIVES-MU0-EPSILON0
% RESULT_ID: HMT-CONSTITUTIVES-MU0-EPSILON0
% STATUS: DEMONSTRATED
% RESIDENCE: manuscrito/sections/ampliacion_20260818/05_constantes_fisicas_genealogia.tex#eq:amp18-constantes-identidades-vacio
% STATEMENT_SHA256: 5a67e3f2c22019847d22853c8038fa4258f414e3c63a0d1568a2e08908c456fd
% STATEMENT: A partir de α_HMT y una única carta global de e,h,c se construye la impedancia Z_0 y, mediante Z_0=μ_0c y Z_0c=1/ε_0, se publican μ_0 y ε_0; sus valores SI son comparación posterior.
% HMT_REPAIR_RESULT_END::HMT-CONSTITUTIVES-MU0-EPSILON0
% HMT_REPAIR_RESULT_BEGIN::HMT-CONSTANT-G
% RESULT_ID: HMT-CONSTANT-G
% STATUS: DEMONSTRATED
% RESIDENCE: manuscrito/sections/ampliacion_20260818/05_constantes_fisicas_genealogia.tex#eq:amp18-constantes-selector-108
% STATEMENT_SHA256: 4e8377ac898d655608bc820b018d8760487159e3f6349a41d5aef00743e4fae5
% STATEMENT: El selector circular fija θ_108=(54/π_HMT)^2 y el transductor exacto define G_int(θ_108). La vía holonómica independiente da la plantilla homogénea G_hol=(c_int^3/ħ_int)(δ_θΛ_27)^2, cuya cifra exige construir y normalizar Λ_27; la carta decimal SI se mantiene separada y no construye ninguna de las dos vías.
% HMT_REPAIR_RESULT_END::HMT-CONSTANT-G
% HMT_REPAIR_RESULT_BEGIN::HMT-MASS-LAW-ATLAS
% RESULT_ID: HMT-MASS-LAW-ATLAS
% STATUS: DEMONSTRATED
% RESIDENCE: manuscrito/sections/ampliacion_20260818/04_electron_y_ley_general_masas.tex#eq:amp18-masas-ley-operatorial
% STATEMENT_SHA256: 4666f47d67c5e3e1b827ca0cb3802e083fbee5770e215ceea03a3ab43eefdc84
% STATEMENT: La partícula es sección persistente y la masa es carácter positivo de esa persistencia. El electrón central único (5,5) construye primero el basal adimensional \widetilde{\mathfrak e}_0=(√3/4)[exp(φ_HMT/π_HMT²)−22α_HMT³](1+15Δ_4), después \widetilde{\mathfrak e}_e^β=\widetilde{\mathfrak e}_0R_act^(80−54α_HMT+6Δ_4), y sólo entonces la carta común T_E:=\varsigma_{u_E}:R_{>0}^{adim}→R_{>0}u_E, T_E(x)=xu_E, u_E=1 MeV publica \mathfrak e_e^β=T_E(\widetilde{\mathfrak e}_e^β)=m_e^βc²=0.5109989506900008086082571648… MeV y m_e^β=\mathfrak e_e^β/c²=0.5109989506900008086082571648… MeV/c². La misma carta común gobierna el atlas 81/56/13/324; los 23 sectores son sólo un inventario humano separado.
% HMT_REPAIR_RESULT_END::HMT-MASS-LAW-ATLAS
% HMT_REPAIR_RESULT_BEGIN::HMT-BARBERO-IMMIRZI
% RESULT_ID: HMT-BARBERO-IMMIRZI
% STATUS: DEMONSTRATED
% RESIDENCE: manuscrito/sections/md/07_hexada_octada_barbero.tex#chap:barbero
% STATEMENT_SHA256: 360316071f48bb431d8e0b9a737194c74511243d354eafeaf5e398f932b048e9
% STATEMENT: Las incidencias 90/120, las series de los dos canales, el residuo 1/24 y la minimalidad diofántica fijan los pesos (12,1) y el funcional Γ_6→8; la identificación γ=Γ ocurre después y no define retroactivamente A o C*.
% HMT_REPAIR_RESULT_END::HMT-BARBERO-IMMIRZI
% HMT_REPAIR_RESULT_BEGIN::HMT-CKM-CP
% RESULT_ID: HMT-CKM-CP
% STATUS: DEMONSTRATED
% RESIDENCE: manuscrito/sections/ampliacion_20260818/05_constantes_fisicas_genealogia.tex#eq:amp18-constantes-ckm-cuatro-ecuaciones
% STATEMENT_SHA256: a245142e9d5486d59a0cd73afac1263a4916dd3b49d3437964457a148a0d3878
% STATEMENT: El registro de rutas y el único par angular construyen cuatro ángulos, la matriz unitaria CKM y un invariante de Jarlskog no nulo; la comparación PDG es posterior y CP no se confunde con CPT.
% HMT_REPAIR_RESULT_END::HMT-CKM-CP
% HMT_REPAIR_RESULT_BEGIN::HMT-GOVERNANCE-NO-REGRESSION
% RESULT_ID: HMT-GOVERNANCE-NO-REGRESSION
% STATUS: DEMONSTRATED
% RESIDENCE: gestion/reparacion_semantica_20260818/MAPA_RECONSTRUCCION_CAUSAL_FUENTE_20260818.md §7
% STATEMENT_SHA256: c59a160a4fe22784718235a339c7a3a3b5df6305829c5ef7e443d93c2ca38b7a
% STATEMENT: Todo artefacto obsoleto que invierta APP→TRIT→estado del TPK, haga externas las semillas, convierta los cinco cierres en tareas pendientes o subordine el continuo al marco clásico se retira de la gobernanza o queda como testigo inactivo; no se arrastra a la edición sucesora.
% HMT_REPAIR_RESULT_END::HMT-GOVERNANCE-NO-REGRESSION
 % Entorno editorial para incorporar cuerpos científicos completos sin
% crear capitulos nuevos. El grupo conserva intactos los archivos de origen:
% solo traduce su jerarquia visible un nivel hacia abajo durante el input.
% La procedencia material de cada uso se registra en comentarios del destino
% y en gestion/RESTAURACION_16_CLAVES_SIN_DESTINO_20260824.tsv.
\newenvironment{HMTScientificOwnerSubordinated}{%
  \begingroup
  \let\HMTScientificOwnerSection\section
  \let\HMTScientificOwnerSubsection\subsection
  \let\HMTScientificOwnerSubsubsection\subsubsection
  \let\HMTScientificOwnerParagraph\paragraph
  \let\HMTScientificOwnerSubparagraph\subparagraph
  \let\chapter\HMTScientificOwnerSection
  \let\section\HMTScientificOwnerSubsection
  \let\subsection\HMTScientificOwnerSubsubsection
  \let\subsubsection\HMTScientificOwnerParagraph
  \let\paragraph\HMTScientificOwnerSubparagraph
}{%
  \par
  \endgroup
}
 \colorlet{hmtlight}{hmtpaper}
\theoremstyle{definition}
\newtheorem{principle}[theorem]{Principio}

% Jerarquía tipográfica de monografía científica. Los títulos de capítulo
% comparten línea con su numeración, se componen con cuerpo moderado y dejan
% espacio suficiente para el primer párrafo o la primera ecuación.
\titleformat{\part}[display]
  {\centering\normalfont\bfseries\color{hmtblue}}
  {\Large\scshape\partname\ \thepart}
  {.45ex}
  {\Huge}
\titlespacing*{\part}{0pt}{1.25cm}{1.0cm}
\titleformat{\chapter}[hang]
  {\normalfont\bfseries\color{hmtblue}\LARGE}
  {\normalsize\scshape\chaptertitlename\ \thechapter}
  {1.0em}
  {\raggedright}
\titlespacing*{\chapter}{0pt}{-0.35cm}{0.75cm}
\titleformat{name=\chapter,numberless}[block]
  {\normalfont\bfseries\color{hmtblue}\LARGE}
  {}{0pt}{\raggedright}
\titlespacing*{name=\chapter,numberless}{0pt}{-0.35cm}{.75cm}
\titleformat{\section}
  {\normalfont\bfseries\Large\color{hmtblue}}
  {\thesection}{.65em}{}
\titlespacing*{\section}{0pt}{2.0ex plus .45ex minus .2ex}{.65ex}
\titleformat{\subsection}
  {\normalfont\bfseries\large}
  {\thesubsection}{.65em}{}
\titlespacing*{\subsection}{0pt}{1.65ex plus .4ex minus .2ex}{.55ex}
\titleformat{\subsubsection}
  {\normalfont\itshape\normalsize}
  {\thesubsubsection}{.65em}{}
\titlespacing*{\subsubsection}{0pt}{1.35ex plus .3ex minus .2ex}{.4ex}

% Jerarquía editorial de la edición revisada. Los archivos heredados
% conservan sus rótulos y etiquetas, pero sus capítulos se presentan como
% secciones dentro de capítulos macroscópicos.
\let\HMTBookChapter\chapter
\let\HMTBookSection\section
\let\HMTBookSubsection\subsection
\let\HMTBookSubsubsection\subsubsection
\let\HMTBookParagraph\paragraph
\let\HMTBookSubparagraph\subparagraph
\newcommand{\HMTDemoteHeadings}{%
  \let\chapter\HMTBookSection
  \let\section\HMTBookSubsection
  \let\subsection\HMTBookSubsubsection
  \let\subsubsection\HMTBookParagraph
  \let\paragraph\HMTBookSubparagraph}
\newcommand{\HMTRestoreHeadings}{%
  \let\chapter\HMTBookChapter
  \let\section\HMTBookSection
  \let\subsection\HMTBookSubsection
  \let\subsubsection\HMTBookSubsubsection
  \let\paragraph\HMTBookParagraph
  \let\subparagraph\HMTBookSubparagraph}

% Notación de la frontera con Mecánica Dimensional.
\providecommand{\APP}{\ensuremath{\mathrm{APP}}}
\providecommand{\TPK}{\ensuremath{\mathrm{TPK}}}
\providecommand{\R}{\mathbb R}
\providecommand{\Z}{\mathbb Z}
\providecommand{\N}{\mathbb N}
\providecommand{\F}{\mathbb F}
\providecommand{\M}{\mathbb M}
\providecommand{\one}{\mathbf 1}
\providecommand{\Halg}{\mathfrak H}
\providecommand{\cA}{\mathcal A}
\providecommand{\cC}{\mathcal C}
\providecommand{\cR}{\mathcal R}
\providecommand{\cX}{\mathcal X}
\providecommand{\dd}{\mathrm d}
\providecommand{\e}{\mathrm e}
\providecommand{\ii}{\mathrm i}
\providecommand{\Znine}{\mathbb Z/9\mathbb Z}
\providecommand{\Cstar}{C^{*}}
\providecommand{\Cstarrad}{C^{*}_{\rm rad}}
\providecommand{\hbarHMT}{\hbar_{\mathrm{HMT}}}
\providecommand{\etaRet}{\eta_{\mathrm{ret}}^{(\mathrm{deg})}}
\providecommand{\etaRetRad}{\eta_{\mathrm{ret}}^{(\mathrm{rad})}}
\providecommand{\alphaAngle}{\alpha_{\angle,\mathrm{deg}}}
\providecommand{\alphaHMT}{\alpha}
\providecommand{\piHMT}{\pi}
\providecommand{\phiHMT}{\varphi}
\providecommand{\eexp}{\mathrm e}
\providecommand{\Dfour}{\Delta_{4}}
\providecommand{\zCorr}{\zeta_{270}^{\mathrm{corr}}}
\providecommand{\sSpec}{s_{270}^{\mathrm{spec}}}
\providecommand{\Kfield}{\mathbb K}
\providecommand{\rank}{\operatorname{rank}}
\providecommand{\nint}{\operatorname{nint}}
\providecommand{\MD}{\ensuremath{\mathrm{MD}}}
\providecommand{\CT}{\ensuremath{\mathrm{CT}_{108}}}
\providecommand{\RR}{\mathbb R}
\providecommand{\ZZ}{\mathbb Z}
\providecommand{\PP}{\mathbb P}
\providecommand{\hbarpre}{\hbar_{\mathrm{pre}}}
\providecommand{\hbarret}{\hbar_{\mathrm{ret}}}
\providecommand{\Ract}{R_{\mathrm{act}}}
\providecommand{\ellP}{\ell_{\mathrm P}}
\providecommand{\mP}{m_{\mathrm P}}
\providecommand{\tP}{t_{\mathrm P}}
\providecommand{\EP}{E_{\mathrm P}}
\providecommand{\rg}{r_{\mathrm g}}
\providecommand{\rs}{r_{\mathrm s}}
\providecommand{\lambdabarC}{\bar\lambda_{\mathrm C}}
\providecommand{\lambdaC}{\lambda_{\mathrm C}}
\providecommand{\Nsp}{N_{\mathrm{sp}}}
\providecommand{\thetaStar}{\theta_{\!*}}
\providecommand{\Gammae}{\Gamma_{e}}
\providecommand{\mebeta}{m_{e}^{\beta}}
% En esta monografía las letras universales no designan duplicados de las
% constantes: son las mismas coordenadas producidas por la construcción y
% reconocidas después en la realización arquimediana.
\renewcommand{\alphaHMT}{\alpha}
\renewcommand{\piHMT}{\pi}
\renewcommand{\phiHMT}{\varphi}
\providecommand{\epigraphtext}[2]{%
  \begin{flushright}%
  \begin{minipage}{0.76\textwidth}\itshape
  #1\par\smallskip
  \raggedleft\normalfont--- #2
  \end{minipage}%
  \end{flushright}%
}

\newtcolorbox{historybox}[1][]{%
  colback=hmtlightblue,colframe=hmtblue,
  title={Hecho histórico documentado},fonttitle=\bfseries,#1}
\newtcolorbox{readingbox}[1][]{%
  colback=hmtlightgold,colframe=hmtgold,
  title={Lectura HMT},fonttitle=\bfseries,#1}
\newtcolorbox{statusbox}[1][]{%
  colback=white,colframe=hmtblue!55,
  title={Hipótesis, dominio y conclusión},fonttitle=\bfseries,#1}
\newtcolorbox{warningbox}[1][]{%
  colback=white,colframe=hmtgold!75!black,
  title={Condición de compatibilidad},fonttitle=\bfseries,#1}
\newtcolorbox{thesisbox}[1][]{%
  colback=white,colframe=hmtviolet,
  title={Tesis rectora},fonttitle=\bfseries,#1}
\newtcolorbox{proofstatusbox}[1][]{%
  colback=white,colframe=hmtgreen,
  title={Estatuto probatorio},fonttitle=\bfseries,#1}
\newtcolorbox{falsifierbox}[1][]{%
  colback=white,colframe=hmtred,
  title={Criterio de refutación},fonttitle=\bfseries,#1}
\newtcolorbox{falsadorBox}[1][]{%
  colback=hmtlightred,colframe=hmtred,
  title={Control negativo},fonttitle=\bfseries,#1}
\newtcolorbox{fronteraBox}[1][]{%
  colback=hmtlightgold,colframe=hmtgold,
  title={Frontera de validez},fonttitle=\bfseries,#1}
\newtcolorbox{programbox}[1][]{%
  colback=hmtlightgold,colframe=hmtgold,
  title={Hipótesis y programa de prueba},fonttitle=\bfseries,#1}
\newtcolorbox{takeawaybox}[1][]{%
  colback=white,colframe=hmtblue,
  title={Resultado del apartado},fonttitle=\bfseries,#1}
% Entornos académicos conservados por los fuentes teoremáticas científicos extensos.
% Se tipan aquí con la jerarquía gráfica única del tratado para evitar que cada
% incorporación restablezca una tipografía o un preámbulo autónomos.
\newtcolorbox{resultado}[1][]{%
  colback=white,colframe=hmtgreen,
  title={Resultado},fonttitle=\bfseries,#1}
\newtcolorbox{definicionnarrativa}[1][]{%
  colback=hmtlightblue,colframe=hmtblue,
  title={Definición},fonttitle=\bfseries,#1}
\newtcolorbox{lecturafisica}[1][]{%
  colback=hmtlightgold,colframe=hmtgold,
  title={Realización física},fonttitle=\bfseries,#1}
\newtcolorbox{frontera}[1][]{%
  colback=hmtlightgold,colframe=hmtgold,
  title={Frontera de validez},fonttitle=\bfseries,#1}
\newtheorem{falsifier}[theorem]{Criterio de refutación}
\newtheorem{test}[theorem]{Programa de contraste}
\newenvironment{keyresult}
  {\begin{quote}\color{hmtblue}\bfseries}
  {\end{quote}}

\lstdefinestyle{hmtcode}{
  basicstyle=\ttfamily\scriptsize,
  backgroundcolor=\color{hmtpaper},
  frame=single,
  rulecolor=\color{hmtgray!45},
  breaklines=true,
  showstringspaces=false,
  columns=fullflexible,
  keepspaces=true
}

\hypersetup{
  colorlinks=true,
  bookmarksdepth=3,
  pdflang={es-ES},
  linkcolor=hmtblue,
  citecolor=hmtblue,
  urlcolor=hmtgold,
  pdftitle={Doble proyección holográfica de un cristal temporal aperiódico},
  pdfsubject={Construcción del continuo, generación de constantes y estabilizadores excepcionales de la incidencia fronteriza},
  pdfauthor={Oumar Haidara Fall; Rubén Ramos Balsa},
  pdfkeywords={aritmética nonádica, orientación ternaria, holonomía,
    cristal temporal aperiódico, continuo, constantes fundamentales,
    simetrías excepcionales, acción, información, gravedad}
}

\setstretch{1.035}
\setlength{\parskip}{0.16em}
\setlength{\parindent}{1.25em}
% El índice científico de la edición definitiva hace visibles capítulos y
% secciones rectoras. Los niveles inferiores permanecen en el cuerpo para
% conservar la demostración sin convertir el índice en un inventario.
\setcounter{tocdepth}{3}
\setcounter{secnumdepth}{3}
\renewcommand{\contentsname}{Índice}
\renewcommand{\cftchappagefont}{\fontsize{11.5}{13.4}\selectfont\bfseries}
\renewcommand{\cftchapfont}{\fontsize{11.5}{13.4}\selectfont\bfseries}
\renewcommand{\cftsecfont}{\fontsize{10.75}{12.8}\selectfont\bfseries}
\renewcommand{\cftsecpagefont}{\fontsize{10.75}{12.8}\selectfont\bfseries}
\renewcommand{\cftsubsecfont}{\fontsize{10.25}{12.2}\selectfont}
\renewcommand{\cftsubsecpagefont}{\fontsize{10.25}{12.2}\selectfont}
\renewcommand{\cftsubsubsecfont}{\fontsize{9.5}{11.4}\selectfont\itshape}
\renewcommand{\cftsubsubsecpagefont}{\fontsize{9.5}{11.4}\selectfont\itshape}
\setlength{\cftbeforechapskip}{.38em}
\setlength{\cftbeforesecskip}{.05em}
\setlength{\cftchapnumwidth}{3.25em}
\setlength{\cftsecindent}{3.25em}
\setlength{\cftsecnumwidth}{4.15em}
\setlength{\cftsubsecnumwidth}{5.4em}
\setlist[itemize]{leftmargin=1.65em,itemsep=.25em,topsep=.35em}
\setlist[enumerate]{leftmargin=2em,itemsep=.25em,topsep=.35em}
\emergencystretch=3em
\widowpenalty=10000
\clubpenalty=10000
\displaywidowpenalty=10000
\brokenpenalty=10000
\predisplaypenalty=500
\postdisplaypenalty=500
\raggedbottom
% Cada epígrafe debe permanecer unido a un desarrollo mínimo. Se protegen
% tanto los mandos públicos como sus copias editoriales, pues los capítulos
% heredados se degradan localmente mediante \HMTDemoteHeadings.
\pretocmd{\section}{\Needspace{9\baselineskip}}{}{}
\pretocmd{\subsection}{\Needspace{7\baselineskip}}{}{}
\pretocmd{\subsubsection}{\Needspace{5\baselineskip}}{}{}
\pretocmd{\HMTBookSection}{\Needspace{9\baselineskip}}{}{}
\pretocmd{\HMTBookSubsection}{\Needspace{7\baselineskip}}{}{}
\pretocmd{\HMTBookSubsubsection}{\Needspace{5\baselineskip}}{}{}
\pretocmd{\HMTBookParagraph}{\Needspace{9\baselineskip}}{}{}

\crefname{chapter}{capítulo}{capítulos}
\Crefname{chapter}{Capítulo}{Capítulos}
\crefname{section}{apartado}{apartados}
\Crefname{section}{Apartado}{Apartados}
\crefname{subsection}{apartado}{apartados}
\Crefname{subsection}{Apartado}{Apartados}
\crefname{subsubsection}{apartado}{apartados}
\Crefname{subsubsection}{Apartado}{Apartados}
\crefname{figure}{figura}{figuras}
\Crefname{figure}{Figura}{Figuras}
\crefname{table}{tabla}{tablas}
\Crefname{table}{Tabla}{Tablas}
\crefname{equation}{ecuación}{ecuaciones}
\Crefname{equation}{Ecuación}{Ecuaciones}
\crefname{theorem}{teorema}{teoremas}
\Crefname{theorem}{Teorema}{Teoremas}
\crefname{lemma}{lema}{lemas}
\Crefname{lemma}{Lema}{Lemas}
\crefname{proposition}{proposición}{proposiciones}
\Crefname{proposition}{Proposición}{Proposiciones}
\crefname{corollary}{corolario}{corolarios}
\Crefname{corollary}{Corolario}{Corolarios}
\crefname{definition}{definición}{definiciones}
\Crefname{definition}{Definición}{Definiciones}
\crefname{construction}{construcción}{construcciones}
\Crefname{construction}{Construcción}{Construcciones}
\crefname{algorithm}{algoritmo}{algoritmos}
\Crefname{algorithm}{Algoritmo}{Algoritmos}
\crefname{ablation}{proposición de necesidad estructural}
  {proposiciones de necesidad estructural}
\Crefname{ablation}{Proposición de necesidad estructural}
  {Proposiciones de necesidad estructural}
\crefname{remark}{observación}{observaciones}
\Crefname{remark}{Observación}{Observaciones}
\crefname{principle}{principio}{principios}
\Crefname{principle}{Principio}{Principios}
\crefname{criterion}{criterio de contraste}{criterios de contraste}
\Crefname{criterion}{Criterio de contraste}{Criterios de contraste}
\AtBeginDocument{%
  \renewcommand{\crefrangeconjunction}{ a\nobreakspace}%
  \renewcommand{\crefpairconjunction}{ y\nobreakspace}%
  \renewcommand{\crefmiddleconjunction}{, }%
  \renewcommand{\creflastconjunction}{ y\nobreakspace}%
  \renewcommand{\crefpairgroupconjunction}{ y\nobreakspace}%
  \renewcommand{\crefmiddlegroupconjunction}{, }%
  \renewcommand{\creflastgroupconjunction}{ y\nobreakspace}%
}

\pagestyle{fancy}
\fancyhf{}
\fancyhead[LE,RO]{\small\thepage}
% Las cajas de marca se reservan dentro del 84 % central del encabezado. Así
% los títulos extensos pueden ocupar dos líneas sin invadir el folio exterior.
\fancyhead[LO]{\parbox[b]{.84\headwidth}{\raggedright\footnotesize\nouppercase{\leftmark}}}
\fancyhead[RE]{\parbox[b]{.84\headwidth}{\raggedleft\footnotesize\nouppercase{\leftmark}}}
\renewcommand{\headrulewidth}{0.25pt}
\fancypagestyle{plain}{%
  \fancyhf{}\fancyfoot[C]{\thepage}\renewcommand{\headrulewidth}{0pt}}

% Jerarquía óptica moderada aprobada para el sucesor y mapa declarativo de
% rótulos públicos. El adaptador del radión aporta sólo sus macros locales;
% no importa el estilo autónomo ni modifica la tipografía general.
% Escala tipográfica aprobada por Rubén el 28 de agosto de 2026.
% Este archivo sustituye exclusivamente la escala de encabezamientos; no
% altera los cuerpos científicos ni sus relaciones de numeración.
\titleformat{\part}[display]
  {\centering\normalfont\color{hmtblue}}
  {\fontsize{12.4}{14.2}\selectfont\mdseries\scshape\partname\ \thepart}
  {2.5pt}
  {\fontsize{16.2}{18.4}\selectfont\bfseries}
\titlespacing*{\part}{0pt}{1.1cm}{.82cm}

\titleformat{\chapter}[display]
  {\normalfont\color{hmtblue}}
  {\fontsize{9.4}{11.2}\selectfont\mdseries\scshape\chaptertitlename\ \thechapter}
  {2.4pt}
  {\fontsize{13.7}{15.7}\selectfont\bfseries\raggedright}
\titlespacing*{\chapter}{0pt}{-.22cm}{.62cm}

\titleformat{name=\chapter,numberless}[block]
  {\normalfont\fontsize{13.7}{15.7}\selectfont\bfseries\color{hmtblue}}
  {}{0pt}{\raggedright}
\titlespacing*{name=\chapter,numberless}{0pt}{-.22cm}{.62cm}

\titleformat{\section}
  {\normalfont\fontsize{11.4}{13.4}\selectfont\bfseries\color{hmtblue}}
  {\thesection}{.65em}{}
\titlespacing*{\section}{0pt}{13pt plus 2pt minus 1pt}{3.5pt}

\titleformat{\subsection}
  {\normalfont\fontsize{10.4}{12.3}\selectfont\bfseries}
  {\thesubsection}{.65em}{}
\titlespacing*{\subsection}{0pt}{8pt plus 1.5pt minus 1pt}{3pt}

\titleformat{\subsubsection}
  {\normalfont\fontsize{9.7}{11.6}\selectfont\itshape}
  {\thesubsubsection}{.65em}{}
\titlespacing*{\subsubsection}{0pt}{6pt plus 1.2pt minus .8pt}{2.5pt}

% El índice reproduce la misma jerarquía descendente y evita que un nivel
% subordinado aparezca ópticamente mayor que su antecedente.
\renewcommand{\cftchappagefont}{\fontsize{10.5}{12.4}\selectfont\bfseries}
\renewcommand{\cftchapfont}{\fontsize{10.5}{12.4}\selectfont\bfseries}
\renewcommand{\cftsecfont}{\fontsize{9.9}{11.8}\selectfont\bfseries}
\renewcommand{\cftsecpagefont}{\fontsize{9.9}{11.8}\selectfont\bfseries}
\renewcommand{\cftsubsecfont}{\fontsize{9.35}{11.2}\selectfont}
\renewcommand{\cftsubsecpagefont}{\fontsize{9.35}{11.2}\selectfont}
\renewcommand{\cftsubsubsecfont}{\fontsize{8.9}{10.8}\selectfont\itshape}
\renewcommand{\cftsubsubsecpagefont}{\fontsize{8.9}{10.8}\selectfont\itshape}
% Los folios de cuatro cifras requieren una caja propia: el margen derecho
% evita que el último término del rótulo colisione con la paginación.
\cftsetpnumwidth{3em}
\cftsetrmarg{3.8em}
 % Mapa declarativo de retitulado aprobado. No altera los cuerpos científicos.
% Cada entrada procede de una coincidencia estática uno-a-uno y conserva el nivel.
\makeatletter
\DeclareUnicodeCharacter{00B2}{\ensuremath{^{2}}}
\DeclareUnicodeCharacter{00B3}{\ensuremath{^{3}}}
\DeclareUnicodeCharacter{00B7}{\ensuremath{\cdot}}
\DeclareUnicodeCharacter{00D7}{\ensuremath{\times}}
\DeclareUnicodeCharacter{0393}{\ensuremath{\Gamma}}
\DeclareUnicodeCharacter{03B1}{\ensuremath{\alpha}}
\DeclareUnicodeCharacter{03B6}{\ensuremath{\zeta}}
\DeclareUnicodeCharacter{03BE}{\ensuremath{\xi}}
\DeclareUnicodeCharacter{2013}{--}
\DeclareUnicodeCharacter{2074}{\ensuremath{^{4}}}
\DeclareUnicodeCharacter{2076}{\ensuremath{^{6}}}
\DeclareUnicodeCharacter{207B}{\ensuremath{^{-}}}
\DeclareUnicodeCharacter{2080}{\ensuremath{_{0}}}
\DeclareUnicodeCharacter{2081}{\ensuremath{_{1}}}
\DeclareUnicodeCharacter{2082}{\ensuremath{_{2}}}
\DeclareUnicodeCharacter{2083}{\ensuremath{_{3}}}
\DeclareUnicodeCharacter{2085}{\ensuremath{_{5}}}
\DeclareUnicodeCharacter{2087}{\ensuremath{_{7}}}
\DeclareUnicodeCharacter{2088}{\ensuremath{_{8}}}
\DeclareUnicodeCharacter{2089}{\ensuremath{_{9}}}
\DeclareUnicodeCharacter{210F}{\ensuremath{\hbar}}
\DeclareUnicodeCharacter{2192}{\ensuremath{\to}}
\DeclareUnicodeCharacter{21A6}{\ensuremath{\mapsto}}
\DeclareUnicodeCharacter{2212}{\ensuremath{-}}
\DeclareUnicodeCharacter{2218}{\ensuremath{\circ}}
\DeclareUnicodeCharacter{221A}{\ensuremath{\surd}}
\DeclareUnicodeCharacter{2260}{\ensuremath{\ne}}
\DeclareUnicodeCharacter{2261}{\ensuremath{\equiv}}
\DeclareUnicodeCharacter{27E8}{\ensuremath{\langle}}
\DeclareUnicodeCharacter{27E9}{\ensuremath{\rangle}}
\newcommand{\HMTTitleMapEntry}[2]{%
  \edef\HMTTitleProbe{\detokenize{#1}}%
  \ifx\HMTTitleKey\HMTTitleProbe\def\HMTResolvedTitle{#2}\fi}
\newcommand{\HMTResolvePublicTitle}[1]{%
  \def\HMTResolvedTitle{#1}%
  \edef\HMTTitleKey{\detokenize{#1}}%
  % El prefacio fue aprobado como texto autoral literal; su subtítulo no se
  % retitula ni se normaliza.
  \HMTTitleMapEntry{Interacción compleja entre dos máquinas simples}{Interacción compleja entre dos máquinas simples}% ADP102-0001-LITERAL
  \HMTTitleMapEntry{Introducción general}{Arquitectura causal del problema del continuo y alcance matemático-físico de HMT–MD}% ADP102-0002
  \HMTTitleMapEntry{Objeto inicial y orden causal: escala finita, horizonte de observación y publicación arquimediana}{Horizonte proyectivo: dependencia entre escala, frontera y medida}% ADP102-0003
  \HMTTitleMapEntry{Intervalo orientado, marcas, celdas y arista de cierre}{Construcción de la recta finita mediante orden, escala y conservación de la unidad}% ADP102-0004
  \HMTTitleMapEntry{Refinamiento, supervivencia y umbrales de prolongación compatible}{Relación de límite entre infinitésimo, prolongación infinita y umbral proyectivo}% ADP102-0005
  \HMTTitleMapEntry{Evaluación arquimediana de secciones compatibles y reconstrucción de \(\mathbb R\)}{Límite compatible de horizontes finitos y realización de la recta real}% ADP102-0006
  \HMTTitleMapEntry{Núcleo formal de la Holografía Modular Triádica}{Construcción y arquitectura operatoria del núcleo formal APP–TRIT–TPK: soporte aritmético, orientación ternaria y dinámica genealógica}% ADP102-0007
  \HMTTitleMapEntry{Levantamientos y cociclo de acarreo}{Extensiones residuo–cociente y cociclo nonádico de acarreo}% ADP102-0008
  \HMTTitleMapEntry{Bloque central de \(2\) posiciones}{Bloque central bivariado de APP y descomposición en 2 modos espectrales}% ADP102-0009
  \HMTTitleMapEntry{Curvatura mixta de las \(2\) leyes}{Curvatura discreta de la interacción entre las hojas aditiva y multiplicativa}% ADP102-0010
  \HMTTitleMapEntry{Correspondencia entre \(2\) cartas de una misma clase residual}{Conjugación entre las cartas positiva y residual de una misma clase módulo 9}% ADP102-0011
  \HMTTitleMapEntry{Núcleo interior y complemento gnomónico}{Descomposición del ortograma APP en núcleo interior y complemento gnomónico}% ADP102-0012
  \HMTTitleMapEntry{Segmentos marcados, intersticios y célula nonádica}{Construcción de la célula nonádica a partir de segmentos marcados e intersticios}% ADP102-0013
  \HMTTitleMapEntry{Carta posicional finita y proyección residual decimal}{Proyección residual decimal de la carta posicional finita de APP}% ADP102-0014
  \HMTTitleMapEntry{Extracción de la célula de \(9\) fases}{Construcción de la célula nonádica de 9 fases mediante proyección residual}% ADP102-0015
  \HMTTitleMapEntry{Posición, altura y división euclídea}{Coordenadas de posición y altura inducidas por la división euclídea}% ADP102-0016
  \HMTTitleMapEntry{Dominio inaugural de la célula nonádica y exclusión causal de APP, TRIT, TPK y constantes analíticas}{Dominio aritmético autónomo de la célula nonádica y precedencia causal respecto de APP–TRIT–TPK y las constantes analíticas}% ADP102-0017
  \HMTTitleMapEntry{Ley del defecto a módulo fijo y longitud variable}{Invariancia del defecto ciclotómico a módulo fijo bajo variación de longitud}% ADP102-0018
  \HMTTitleMapEntry{Operador ciclotómico primario}{Operador ciclotómico de APP y acción sobre las 12 fibras de tipo A₂}% ADP102-0019
  \HMTTitleMapEntry{Localización del agregado \texorpdfstring{\(54\)}{54}}{Derivación del defecto suma–producto 459−405=54 en el bloque central de APP}% ADP102-0020
  \HMTTitleMapEntry{Resonancia del operador de incidencia en rango \(6\)}{Rango 6 del operador de incidencia y multiplicidad de su modo resonante}% ADP102-0021
  \HMTTitleMapEntry{TRIT como unidad mínima de información topológica: orientación ternaria, memoria de acarreo y orden temporal}{TRIT como unidad ternaria orientada: geometrías internas, orden temporal y condiciones dinámicas de cristal de tiempo}% ADP102-0022
  \HMTTitleMapEntry{Composición y cociclo de acarreo}{Cociclo de acarreo de la composición ternaria orientada}% ADP102-0023
  \HMTTitleMapEntry{Familia universal de álgebras cuadráticas}{Clasificación unificada de las 3 álgebras cuadráticas del TRIT}% ADP102-0024
  \HMTTitleMapEntry{Topological Phase Kernel: base observable y estado holonómico integral}{TPK: base observable, estado holonómico integral y separación tipada de sus proyecciones}% ADP102-0025
  \HMTTitleMapEntry{La base doble de caminos APP}{Categoría bivariada de caminos APP como base del estado holonómico integral TPK}% ADP102-0026
  \HMTTitleMapEntry{Fase nonádica, sector dodecafásico y calendario}{Extensión no escindida C₁₂→C₁₀₈→C₉ y registro dodecafásico de la fase nonádica}% ADP102-0027
  \HMTTitleMapEntry{Cilindro y frontera}{Compatibilidad cilíndrica y firma de frontera del estado holonómico integral}% ADP102-0028
  \HMTTitleMapEntry{Transporte de fibras}{Operador cardinal \(T_d\) de transporte entre fibras enriquecidas con conservación de ruta y memoria}% ADP102-0029
  \HMTTitleMapEntry{Relaciones de extensión}{Relación tipada \(\operatorname{Ext}_r\) de prolongación entre fibras enriquecidas}% ADP102-0030
  \HMTTitleMapEntry{Proyecciones y pérdidas tipadas}{Cocientes observables del estado holonómico integral y campos genealógicos conservados}% ADP102-0031
  \HMTTitleMapEntry{Árbol operatorio del contenedor}{Inserción de las 6 familias fibradas del estado holonómico integral en la jerarquía operatoria del TPK}% ADP102-0032
  \HMTTitleMapEntry{Álgebra de operadores del Topological Phase Kernel y dinámica causal de actualización del estado holonómico integral}{Jerarquía tipada de las 34 construcciones canónicas heterogéneas y dinámica causal de actualización del TPK}% ADP102-0033
  \HMTTitleMapEntry{Clasificación del estado holonómico integral en \(8\) clases estructurales invariantes}{Jerarquía tipada L₀–L₇ de las 34 construcciones canónicas heterogéneas del TPK}% ADP102-0034
  \HMTTitleMapEntry{Descomposición funcional en \(14\) agrupaciones compatibles con los operadores APP–TRIT–TPK}{Agrupación funcional de las 34 construcciones canónicas en 14 familias de composición APP–TRIT–TPK}% ADP102-0035
  \HMTTitleMapEntry{Dominio y tipado de las \(34\) entradas del operador de actualización}{Clases matemáticas, dominios, codominios y dependencias causales de las 34 construcciones canónicas}% ADP102-0036
  \HMTTitleMapEntry{Composición causal de la actualización del estado holonómico integral}{Actualización del estado holonómico integral mediante \(U_t=\operatorname{Upd}_t\circ\operatorname{Tra}_t\circ\operatorname{Sel}_t\)}% ADP102-0037
  \HMTTitleMapEntry{Separación tipada entre operadores de transporte, memoria, calendario, extracción y retorno}{Separación tipada entre operadores, estados, proyecciones y publicaciones del TPK}% ADP102-0038
  \HMTTitleMapEntry{Descomposición operatorial de los canales de incidencia \texorpdfstring{\(90/120\)}{90/120}: selector temporal, generador nilpotente y componente angular}{Incidencia trítica 90/120 y separación tipada de los operadores temporal, nilpotente y angular}% ADP102-0039
  \HMTTitleMapEntry{Distribución de los índices \(90/120\) entre los \(6\) objetos del sistema de incidencia}{Residencias tipadas de los índices 90/120 en 6 objetos operatorios no equivalentes}% ADP102-0040
  \HMTTitleMapEntry{Defecto normalizado de la incidencia}{Defecto normalizado 1−90/120=1/4 de la incidencia trítica}% ADP102-0041
  \HMTTitleMapEntry{Invariantes de tipado}{Invariantes de dominio, orientación y memoria de las realizaciones 90/120}% ADP102-0042
  \HMTTitleMapEntry{Estado local y realimentación de hoja}{Operador local de transición orientada y actualización de la hoja APP}% ADP102-0043
  \HMTTitleMapEntry{Cilindros de los \(2\) lectores}{Sistemas cilíndricos de los 2 lectores aritméticos y compatibilidad de truncación}% ADP102-0044
  \HMTTitleMapEntry{Descomposición del Topological Phase Kernel en operadores de avance, giro, acarreo y memoria}{Factorización de la transición TPK como \(U_t=\operatorname{Upd}_t\circ\operatorname{Tra}_t\circ\operatorname{Sel}_t\)}% ADP102-0045
  \HMTTitleMapEntry{Soporte ordenado y coordenada cíclica}{Coordenada cíclica sobre el soporte temporal orientado de 12 sectores}% ADP102-0046
  \HMTTitleMapEntry{Extensión cíclica no escindida}{Extensión no escindida 0→C₁₂→C₁₀₈→C₉→0 y cociclo de acarreo}% ADP102-0047
  \HMTTitleMapEntry{Extractor firmado de eventos}{Extracción del estado dodecafásico firmado a partir del registro temporal enriquecido}% ADP102-0048
  \HMTTitleMapEntry{\(2\) palabras dodecafásicas de distinta procedencia}{Distinción genealógica entre 2 palabras dodecafásicas de igual período}% ADP102-0049
  \HMTTitleMapEntry{Refinamiento mínimo por microinstantes}{Subdivisión temporal mínima compatible con la semiconjugación}% ADP102-0050
  \HMTTitleMapEntry{\(3\) operaciones cuaternarias no intercambiables}{Conmutadores de 3 operaciones cuaternarias y separación de sus acciones}% ADP102-0051
  \HMTTitleMapEntry{Lector trítico del estado orientado: orden temporal y transporte de acarreo}{Lector trítico orientado con transporte de cociente y acarreo}% ADP102-0052
  \HMTTitleMapEntry{Cociente mínimo de memoria predictiva}{Cociente de memoria de orden 1 y suficiencia predictiva de la proyección modal}% ADP102-0053
  \HMTTitleMapEntry{Memoria vertical no acotada}{Crecimiento no acotado de la profundidad de memoria bajo retornos de fase}% ADP102-0054
  \HMTTitleMapEntry{Realizaciones geométricas del TRIT: regímenes elíptico, parabólico e hiperbólico y doble orientación temporal}{TRIT como estado ternario orientado: memoria de acarreo, orden temporal y 3 regímenes cuadráticos}% ADP102-0055
  \HMTTitleMapEntry{APP, TRIT y TPK: soporte, orientación y transporte}{Composición APP–TRIT–TPK: soporte aritmético, orientación ternaria y transporte con memoria}% ADP102-0056
  \HMTTitleMapEntry{TPK como sistema de transporte}{TPK como dinámica de selección, transporte, memoria, retorno y despliegue}% ADP102-0057
  \HMTTitleMapEntry{Periodicidad temporal de (108) fases y leyes de retorno a (27), (54) y (108) iteraciones}{Periodicidad temporal de 108 fases y retornos a 27, 54 y 108 iteraciones}% ADP102-0058
  \HMTTitleMapEntry{APP, TRIT y TPK: definiciones rectoras}{Composición generativa APP→TRIT→TPK y construcción del estado holonómico integral}% ADP102-0059
  \HMTTitleMapEntry{Consecuencias y publicaciones del continuo construido}{Conexión nonádica conjunta Γ₉ como correspondencia compatible del sistema inverso}% ADP102-0060
  \HMTTitleMapEntry{Condicionamiento como poda de futuros}{Condicionamiento de la medida mediante restricción del árbol de prolongaciones futuras}% ADP102-0061
  \HMTTitleMapEntry{Holonomía nonádica \texorpdfstring{\(\Gamma_9\)}{Γ₉} del Topological Phase Kernel: fibra \texorpdfstring{\(\mathbb F_3^6\)}{F₃⁶}, prolongación coinductiva y retorno de fase con memoria de estado}{Holonomía nonádica Γ₉ del Topological Phase Kernel: transporte del estado holonómico integral, prolongación coinductiva y retorno 9→1 con avance de memoria}% ADP102-0062
  \HMTTitleMapEntry{Conexión nonádica conjunta, compresión y memoria terminal}{Holonomía conjunta Γ₉: compresión observable T, componente ortogonal η y conservación terminal de memoria}% ADP102-0063
  \HMTTitleMapEntry{Relaciones locales y holonomía}{Relaciones locales de transición y construcción de la holonomía nonádica Γ₉}% ADP102-0064
  \HMTTitleMapEntry{Obstrucciones de la holonomía \(\Gamma_9\): retorno \(9\to1\), índices de composición y conservación de fase, acarreo, frontera, ruta y memoria}{Obstrucciones de la holonomía Γ₉: retorno 9→1, índices de composición y conservación de fase, acarreo, frontera, ruta y memoria}% ADP102-0065
  \HMTTitleMapEntry{Descomposición de la microfibra de clausura de \texorpdfstring{\(\pi\)}{π} en \texorpdfstring{\(5\)}{5} regiones orientadas: multiplicidad \texorpdfstring{\(432+4\cdot144=1008\)}{432+4·144=1008}, eje \texorpdfstring{\(\varphi\)}{φ} e involución especular \texorpdfstring{\(\pi/e\)}{π/e}}{Clasificación orientada de la microfibra de clausura de π: 5 regiones, multiplicidad 432+4·144=1008, eje φ e involución π/e}% ADP102-0066
  \HMTTitleMapEntry{Generación coinductiva de \texorpdfstring{\(\pi\)}{π}, \texorpdfstring{\(e\)}{e} y \texorpdfstring{\(\varphi\)}{φ} mediante refinamiento nonádico compatible}{Generación coinductiva de π, e y φ mediante refinamiento nonádico compatible}% ADP102-0067
  \HMTTitleMapEntry{Elevación visible y cadena de bloques}{Prolongación coinductiva de la emisión exponencial mediante bloques compatibles}% ADP102-0068
  \HMTTitleMapEntry{Firma decimal y cuadrado posicional}{Construcción de la firma decimal mediante el cuadrado posicional ternario–decimal}% ADP102-0069
  \HMTTitleMapEntry{Retorno visible y ruptura de memoria}{Retorno de fase con actualización no trivial de la memoria del estado}% ADP102-0070
  \HMTTitleMapEntry{Certificado adversarial de la generación de \(e\): unicidad orbital, recurrencia factorial, horquillas racionales y prolongación cilíndrica}{Unicidad orbital, recurrencia factorial y convergencia cilíndrica de la generación de e}% ADP102-0071
  \HMTTitleMapEntry{Elevación por bloques y 1.ª frontera coinductiva}{Prolongación por bloques y construcción de la 1.ª frontera coinductiva}% ADP102-0072
  \HMTTitleMapEntry{Del calendario trítico al multigrafo de incidencia}{Construcción del multigrafo de incidencia a partir del calendario trítico}% ADP102-0073
  \HMTTitleMapEntry{Caracterización, horquillas y prolongación}{Caracterización de la autoescala mediante horquillas racionales y prolongación coinductiva}% ADP102-0075
  \HMTTitleMapEntry{Selección diagonal mínima}{Selección diagonal mínima de cilindros compatibles y convergencia a precisión arbitraria}% ADP102-0076
  \HMTTitleMapEntry{Transferencia Markov compatible con todo el límite inverso}{Dinámica de Markov como cociente modal de memoria de orden 1 de la holonomía Γ₉, compatible con el límite inverso}% ADP102-0077
  \HMTTitleMapEntry{2.º sistema de coordenadas: diferencias de pasos \(3\) y \(4\)}{Sistema diferencial de coordenadas inducido por los operadores de paso 3 y 4}% ADP102-0078
  \HMTTitleMapEntry{Obstrucciones de la prolongación afín de \(\alpha\): telescopía, inversión de \(K\), frontera \(662/663\) y aislamiento de datos objetivo}{Obstrucciones de la prolongación afín de α: telescopía, inversión de K, frontera 662/663 y aislamiento de datos objetivo}% ADP102-0079
  \HMTTitleMapEntry{Construcción de \texorpdfstring{\(\sqrt{-1}\)}{√(-1)} mediante \texorpdfstring{\(J^2=-I\)}{J²=-I} y coordinación de las geometrías elíptica, parabólica e hiperbólica}{Construcción de √(−1) y familia exponencial trítica de Euler: operadores J_τ y coordinación de los regímenes elíptico, parabólico e hiperbólico}% ADP102-0080
  \HMTTitleMapEntry{El Núcleo de Transducción Hensel--Witt}{Núcleo de transducción Hensel–Witt: dominio, operación y preservación genealógica}% ADP102-0081
  \HMTTitleMapEntry{Núcleo de Transducción Hensel–Witt: microregiones de \(\pi\), holonomía \(\Gamma_9\), registro dodecafásico y transporte hexada–octada}{Núcleo de Transducción Hensel–Witt: microregiones de π, holonomía Γ₉, registro dodecafásico y transporte hexada–octada}% ADP102-0082
  \HMTTitleMapEntry{Del código ternario al sistema de Witt}{Transducción del código ternario al sistema de Witt}% ADP102-0083
  \HMTTitleMapEntry{Hexada, octada y el corredor \texorpdfstring{\(90/120\)}{90/120}}{Inclusión hexada–octada y transporte de incidencia por los canales 90/120}% ADP102-0084
  \HMTTitleMapEntry{Extractor de escala triádico}{Operador triádico de extracción de escala y su imagen racional}% ADP102-0085
  \HMTTitleMapEntry{Defecto de potencias primas y operador número}{Representación del defecto de potencias primas mediante el operador número}% ADP102-0086
  \HMTTitleMapEntry{Ablaciones del selector y de la continuación}{Condiciones de unicidad del selector regional y de la prolongación analítica}% ADP102-0087
  \HMTTitleMapEntry{Cadena de generación \(\mathcal F_\pi\to b_\pi\to\rho_\pi\to D_A\to C_*^{\rm HMT}\to\Gamma_{6\to8}\): selección regional, prolongación determinantal y coordenada orientada}{Cadena de generación \(\mathcal F_\pi\to b_\pi\to\rho_\pi\to D_A\to C_*^{\mathrm{HMT}}\to\Gamma_{6\to8}\): selección regional, prolongación determinantal y coordenada orientada}% ADP102-0088
  \HMTTitleMapEntry{Descomposición incidencial hexada--octada en los sectores \texorpdfstring{\(q_{90}\)}{q₉₀} y \texorpdfstring{\(q_{120}\)}{q₁₂₀}: entrelazador \texorpdfstring{\(6\to8\)}{6→8} y estructura precursora del espectro de área}{Canales de incidencia 90/120 del Topological Phase Kernel: semicoronas tríticas, factorización de hojas y realización hexada–octada del morfismo 6→8}% ADP102-0089
  \HMTTitleMapEntry{Geometría dodecafásica de mezcla: antecedente angular CKM--CP y lector precursor del electrón}{Operador de mezcla dodecafásico: fase CKM–CP y construcción del lector predimensional del electrón}% ADP102-0090
  \HMTTitleMapEntry{Morfismos arquimediano y excepcional desde un dominio enriquecido común}{Principio de Ambivalencia: involución de los lectores areal y volumétrico sobre el estado holonómico integral}% ADP102-0091
  \HMTTitleMapEntry{La operación local sobre las 729 palabras}{Dominio enriquecido común y codominios inequivalentes de las 2 proyecciones}% ADP102-0092
  \HMTTitleMapEntry{Naturalidad y compatibilidad a toda profundidad}{Separación conjunta de historias por las proyecciones arquimediana y excepcional}% ADP102-0093
  \HMTTitleMapEntry{Acoplamiento postgenerativo con el cierre operatorial de Euler}{Invariante de incidencia común entre el estado holonómico integral y el sistema de Witt}% ADP102-0094
  \HMTTitleMapEntry{Globalización natural de las proyecciones arquimediana y excepcional sobre un sistema inverso común}{Descomposición de la microfibra de π en 5 regiones orientadas: 4 componentes periféricas y 1 eje fijo}% ADP102-0095
  \HMTTitleMapEntry{Sistemas compatibles a profundidad finita}{Prolongación reticular de la incidencia y construcción de pantallas dimensionales}% ADP102-0096
  \HMTTitleMapEntry{Calibre de la publicación excepcional}{Factorización del lector de incidencia a través del estado holonómico integral}% ADP102-0097
  \HMTTitleMapEntry{Naturalidad conjunta de las publicaciones del estado genealógico}{Doble publicación del estado holonómico integral: evaluación arquimediana y conservación de incidencia}% ADP102-0098
  \HMTTitleMapEntry{Separación conjunta de valor e incidencia excepcional}{Unicidad del calibre inicial de la proyección excepcional}% ADP102-0099
  \HMTTitleMapEntry{Antecedente no conmutativo común de la doble proyección}{Álgebra no conmutativa común y factorización de las 2 proyecciones}% ADP102-0100
  \HMTTitleMapEntry{Del espacio de historias a la recta real}{Cociente del espacio de historias compatibles y realización de la recta real}% ADP102-0101
  \HMTTitleMapEntry{Media arquimediana y doble proyección del mismo estado}{Compatibilidad de la evaluación arquimediana media con la doble proyección del estado holonómico integral}% ADP102-0102
  \HMTTitleMapEntry{Calendario unitario, memoria modular y modo estable}{Modo estable del calendario unitario bajo actualización de memoria modular}% ADP102-0103
  \HMTTitleMapEntry{Publicaciones analíticas de historias HMT}{Realizaciones analíticas de historias HMT mediante las cartas arquimediana y bidireccional}% ADP102-0104
  \HMTTitleMapEntry{Consecuencia para una acción positiva}{Positividad de la acción inducida por el entrelazador dodecafásico–Witt}% ADP102-0105
  \HMTTitleMapEntry{Del salto espectral \texorpdfstring{\(4\)}{4} al retículo de Leech: despliegue APP \texorpdfstring{\(4\to2\to6\to54\)}{4→2→6→54} y corredor Paley--Witt--Mathieu--Golay--Niemeier}{Del salto espectral 4 al retículo de Leech: despliegue APP 4→2→6→54 y corredor Paley–Witt–Mathieu–Golay–Niemeier}% ADP102-0106
  \HMTTitleMapEntry{Naturaleza de la estructura obtenida}{Clasificación algebraica de la estructura ternaria autodual obtenida}% ADP102-0107
  \HMTTitleMapEntry{Ley de incidencia \texorpdfstring{\(4+1\)}{4+1} sobre una tétrada}{Descomposición 4+1 de la incidencia de una tétrada en componentes periféricas y axial}% ADP102-0108
  \HMTTitleMapEntry{Módulo y forma discriminantes}{Descomposición de la incidencia de una tétrada en 4 componentes periféricas y 1 componente axial}% ADP102-0109
  \HMTTitleMapEntry{\(12\) factores y aplicación de pegado}{Correspondencia entre estratos de incidencia y órbita libre del triple positivo}% ADP102-0110
  \HMTTitleMapEntry{Carácter y núcleo de congruencia}{Pegado discriminante de 12 factores A₂ y construcción de la extensión integral}% ADP102-0111
  \HMTTitleMapEntry{Curvaturas mixtas de APP}{Sistema ordenado de incidencias y orientación de la cámara radial}% ADP102-0112
  \HMTTitleMapEntry{Representación de Weyl--Heisenberg y fase central}{Construcción del carácter radial y determinación de su núcleo de congruencia}% ADP102-0113
  \HMTTitleMapEntry{Estatuto exacto de la jerarquía \(4\to2\to6\to54\) y de la identificación reticular.}{Estatuto exacto de la jerarquía 4→2→6→54 y de la identificación reticular.}% ADP102-0114
  \HMTTitleMapEntry{Traza graduada del generador ternario}{Traza graduada de la acción del generador ternario sobre \(V^{\natural}\)}% ADP102-0115
  \HMTTitleMapEntry{Palabra de soporte total e isometría fijo-libre}{Isometría sin puntos fijos inducida por la palabra ternaria de soporte total}% ADP102-0116
  \HMTTitleMapEntry{Construcción del espacio de Hilbert a partir de la conexión nonádica y sus operadores de fase}{Realización proyectiva de Bloch a partir de la estructura compleja interna}% ADP102-0117
  \HMTTitleMapEntry{La rama marcada cambia de tipo}{Inmersión por esquinas de una fase marcada: ideal compacto \(\mathcal K(\mathcal H_\infty)\) y factor \(\mathcal B(\mathcal H_\infty)\) de tipo \(\mathrm I_\infty\)}% ADP102-0118
  \HMTTitleMapEntry{Residuo visible y cociente de frontera}{Extensión operatoria de APP: modos visibles y modos de memoria}% ADP102-0119
  \HMTTitleMapEntry{Involución de dualidad T}{Completación 1000=729+243+27+1 y estratificación de la unidad}% ADP102-0120
  \HMTTitleMapEntry{Completación estratificada de la unidad}{Descomposición 243=1+22+90+120+10 y tipado de sus componentes}% ADP102-0121
  \HMTTitleMapEntry{El bloque \texorpdfstring{\(243\)}{243} como frontera y bola ternaria perfecta}{Derivación de los canales de constantes a partir de la descomposición del bloque 243}% ADP102-0122
  \HMTTitleMapEntry{Descomposición interna del bloque \texorpdfstring{\(243\)}{243}}{Construcción de la dimensión 11 mediante lectores residuales compatibles}% ADP102-0123
  \HMTTitleMapEntry{Dimensión \texorpdfstring{\(11\)}{11} y lecturas compatibles}{Clausura del registro genealógico bajo la involución de dualidad T}% ADP102-0124
  \HMTTitleMapEntry{Identidades verificables de la construcción}{Identidades exactas de involución, modularidad y conservación residual}% ADP102-0125
  \HMTTitleMapEntry{Nivel, intercambio lineal de masas y modularidad: \(3\) controles de no fusión}{Separación por nivel, intercambio lineal de masas y modularidad de las 3 realizaciones}% ADP102-0126
  \HMTTitleMapEntry{\(2\) descendientes del código ternario extendido}{Reducción dimensional 26→24→11→10 y coordinación con las pantallas de teoría M}% ADP102-0127
  \HMTTitleMapEntry{La representación icosaédrica de dimensión \(12\)}{Coordinación de las pantallas exactas con la realización física de teoría M}% ADP102-0128
  \HMTTitleMapEntry{Reducción isotrópica de rango \(26\to24\)}{Representación operatoria de APP y terna espectral supersimétrica en dimensión 11}% ADP102-0129
  \HMTTitleMapEntry{Coordinación del corredor Moonshine con las pantallas de teoría M: gradación, retículo lorentziano y doble realización del estado HMT}{Teorema de compatibilidad entre el corredor Moonshine y las pantallas de teoría M mediante gradación y transporte reticular}% ADP102-0130
  \HMTTitleMapEntry{Antecedente común de lectores: refinamiento espectral y límite de dilatación conjunta}{Refinamiento espectral común y límite de la dilatación conjunta de los lectores matriciales}% ADP102-0131
  \HMTTitleMapEntry{Subestructuras, objeto global y proyecciones}{Construcción de 4 cuadrados latinos cuánticos cíclicos en una clase operatoria elemental}% ADP102-0132
  \HMTTitleMapEntry{Correspondencias directas e inversas entre los formalismos}{Realización de 3 contextos operatorios en un cuadrado cuántico (9,3)}% ADP102-0133
  \HMTTitleMapEntry{Extensión de la bigraduación a 3 índices}{Representación del atlas genealógico completo en un cuadrado cuántico (12,3)}% ADP102-0134
  \HMTTitleMapEntry{4 cuadrados latinos cuánticos cíclicos de clase operatorial fácil}{Clasificación cíclica de 4 cuadrados latinos cuánticos en una clase operatoria elemental}% ADP102-0135
  \HMTTitleMapEntry{Lema cíclico de cuadrados mágicos cuánticos asociados a una medida operatorial positiva}{Bigraduación de la familia genealógico–matricial}% ADP102-0136
  \HMTTitleMapEntry{Conos matriciales, rangos y extensiones completamente positivas}{Representación matricial de la composición genealógica APP–TRIT–TPK}% ADP102-0137
  \HMTTitleMapEntry{Construcción explícita de cubos mágicos cuánticos y compatibilidad de sus márgenes}{TPK como categoría de rutas con selección, transporte, memoria y retorno}% ADP102-0138
  \HMTTitleMapEntry{Combinaciones mágicas de mapas CP}{Composiciones completamente positivas compatibles con la bigraduación genealógico–matricial}% ADP102-0139
  \HMTTitleMapEntry{Jerarquía espectral \(4\to2\to6\to54\) del bloque central de APP}{Realización metrológica posterior de la doble proyección del estado holonómico integral}% ADP102-0140
  \HMTTitleMapEntry{Refinamiento nonádico mediante esperanza condicional}{Descomposición simpléctica 12→4×3 de las coordenadas de fase}% ADP102-0141
  \HMTTitleMapEntry{Taxonomía funcional de \(114\) expresiones numéricas y \(106\) lecturas operatoriales}{Clasificación operatoria de 114 expresiones numéricas por dominio, generador, invariante y codominio}% ADP102-0142
  \HMTTitleMapEntry{Transferencia de fase sobre el multigrafo de bloques completos}{Microfibra de clausura de π: 5 regiones orientadas, multiplicidades \(432+4\cdot144=1008\) y proyección común \(w_\pi\)}% ADP102-0143
  \HMTTitleMapEntry{\(5\) ciclos orientados de clausura de la microfibra}{Publicación de memoria reducida de un bloque nonádico completo mediante el multigrafo de transferencia de fase}% ADP102-0144
  \HMTTitleMapEntry{Operador de incidencia y media vuelta orientada}{Holonomías orientadas de clausura de las 5 regiones de π}% ADP102-0145
  \HMTTitleMapEntry{Composición racional y compensador \texorpdfstring{\(239\)}{239}}{Determinación de la pendiente 1/5 por la multiplicidad de la microfibra regional de π}% ADP102-0146
  \HMTTitleMapEntry{Identidad gaussiana y \(1^{\mathrm a}\) media vuelta orientada}{Composición racional de la media vuelta y determinación del término compensador 239}% ADP102-0147
  \HMTTitleMapEntry{Prolongación coinductiva de la sección de \(\pi\) mediante el sistema inverso de historias supervivientes}{Filtración de F₃⁶ por admisibilidad regional y selección de subcilindros compatibles}% ADP102-0148
  \HMTTitleMapEntry{Sistema inverso de historias compatibles}{Unicidad de la sección diagonal mínima en el sistema de horquillas racionales}% ADP102-0149
  \HMTTitleMapEntry{Firma posicional decimal y cuadrado de compatibilidad}{Prolongación de la sección exponencial mediante la cadena de bloques visibles}% ADP102-0150
  \HMTTitleMapEntry{Retorno de fase con cambio de estado y conservación de la memoria recorrida}{Construcción de la firma decimal y del cuadrado posicional de la sección exponencial}% ADP102-0151
  \HMTTitleMapEntry{Semigrupo factorial del carácter exponencial: recurrencia \(a_n=1/n!\)}{Retorno de fase con desigualdad de estados holonómicos integrales: \(g(K+9)=g(K)\) y \(B_{K+9}\ne B_K\)}% ADP102-0152
  \HMTTitleMapEntry{Certificado interno de unicidad orbital, recurrencia factorial, convergencia y prolongación cilíndrica}{Prolongación coinductiva de e mediante la holonomía nonádica Γ₉ del TPK}% ADP102-0153
  \HMTTitleMapEntry{Rayo positivo invariante y unicidad de la autoescala}{Construcción del multigrafo de incidencia modal a partir del calendario trítico}% ADP102-0154
  \HMTTitleMapEntry{Certificado interno de unicidad, compatibilidad y convergencia de la autoescala de Perron}{Caracterización de φ por el rayo de Perron, encaje de horquillas y prolongación coinductiva}% ADP102-0156
  \HMTTitleMapEntry{Derivaciones internas de la constante de estructura fina \(\alpha\): cierre entero dodecafásico reversible y raíz analítica única del núcleo \(E_{21/22}\)}{Clausura dodecafásica de \(\alpha\) y caracterización analítica mediante \(E_{21/22}\): generación reversible, raíz simple y prolongación nonádica de orden 9}% ADP102-0157
  \HMTTitleMapEntry{Estado terminal común y 2 lecturas internas coordinadas}{Evaluaciones terminal y dodecafásica del estado holonómico integral común: bloques \(B_{\mathrm{term}}\) y carta firmada \(U_{12}^{\mathrm{sgn}}\)}% ADP102-0158
  \HMTTitleMapEntry{Vía A: cierre entero dodecafásico}{Generación reversible de \(\alpha\) mediante la clausura dodecafásica entera}% ADP102-0159
  \HMTTitleMapEntry{Dominio Hensel--Witt, mapas de Teichmüller y objetivo racional \(-1/12\)}{Operador Hensel–Witt de transducción entre microregiones de π, holonomía Γ₉, registro dodecafásico e incidencia hexada–octada}% ADP102-0160
  \HMTTitleMapEntry{Aplicación de Teichmüller del código ternario al anillo de Witt}{Construcción del sistema de Witt desde la codificación ternaria compatible}% ADP102-0161
  \HMTTitleMapEntry{Incidencia hexada--octada y canales modulares \(90/120\)}{Realización hexada–octada de los canales 90/120 previamente derivados de la incidencia trítica del TPK}% ADP102-0162
  \HMTTitleMapEntry{4 rutas compatibles hacia \texorpdfstring{\(-1/12\)}{-1/12}}{4 realizaciones operatorias de −1/12 y separación de sus dominios}% ADP102-0163
  \HMTTitleMapEntry{Derivación determinantal de la escala de acción: forma normal de Smith, conúcleo de orden \(16\) y valoración \(10^{-34}\)}{Derivación determinantal de la escala absoluta de acción: forma normal de Smith, conúcleo de orden 16 y valoración 10⁻³⁴}% ADP102-0164
  \HMTTitleMapEntry{Medida de cociente y carácter positivo de la acción: hilbertización canónica, defecto regional--dodecafásico y ley functorial sobre caminos}{Construcción de la medida positiva de acción sobre el cociente de emisiones: hilbertización, defecto regional–dodecafásico y functorialidad de caminos}% ADP102-0165
  \HMTTitleMapEntry{Derivación autodual de la constante de Planck: periodicidad temporal de \(108\) fases, acción y longitud de Compton}{Derivación autodual de la constante de Planck y de su recta de acción: periodicidad temporal de 108 fases, cierre de Compton y covariancia bidireccional}% ADP102-0166
  \HMTTitleMapEntry{Equivalencia de normalizaciones sobre la identidad autodual}{Conmutatividad de las normalizaciones angular, temporal y gravitatoria de la identidad autodual}% ADP102-0167
  \HMTTitleMapEntry{Operador angular predimensional: descomposición en \(27\) sectores, radio espectral y composición efectiva de la doble proyección}{Operador de fase anterior a la transducción dimensional: descomposición en 27 sectores, radio espectral y factorización de la doble proyección}% ADP102-0168
  \HMTTitleMapEntry{Caracteres cuadráticos y constantes espectrales: Catalán, Apéry y torre de Euler--Stieltjes}{Caracteres cuadráticos y constantes espectrales: Catalán, Apéry, Euler–Mascheroni y jerarquía de Stieltjes}% ADP102-0169
  \HMTTitleMapEntry{Criticidad unimodal de la familia dodecafásica: perfil analítico exacto, aproximantes de Feigenbaum y reducción de la universalidad al certificado hiperbólico}{Criticidad unimodal de la familia dodecafásica: perfil analítico exacto, aproximantes de Feigenbaum y certificado hiperbólico de renormalización}% ADP102-0170
  \HMTTitleMapEntry{Aritmética del modo \(30\): identidades de Rogers--Ramanujan, normalizador \(899\) y levantamiento nonádico de Pell}{Aritmética de la familia modular dodecafásica de periodo 30: identidades de Rogers–Ramanujan, normalizador 899 y levantamiento nonádico de Pell}% ADP102-0171
  \HMTTitleMapEntry{Teorema rector de la red radical y la ley potencia--vacancia}{Clasificación de la red radical 3→6→9 y derivación de la ley potencia–vacancia}% ADP102-0172
  \HMTTitleMapEntry{Órbita cíclica \(369\to693\to936\to369\)}{Órbita cíclica 369→693→936→369}% ADP102-0173
  \HMTTitleMapEntry{Núcleo \(7227\), descomposición espectral y enlace con la cruz radical}{Determinación del núcleo aritmético 7227 por incidencia con la cruz radical de APP}% ADP102-0174
  \HMTTitleMapEntry{Codificación mecánica de las vacancias por las potencias nonádicas}{Codificación sturmiana de las vacancias inducida por las potencias nonádicas}% ADP102-0175
  \HMTTitleMapEntry{Raíz interna de \(-1\), estructura exponencial y geometrías elíptica, parabólica e hiperbólica}{Familia exponencial triádica \(J_\tau^2=-\tau I\): realizaciones elíptica, parabólica e hiperbólica de la identidad de Euler}% ADP102-0176
  \HMTTitleMapEntry{Realización continua del tipo cuadrático inducido}{Extensión continua del operador cuadrático común y conservación de su tipo espectral}% ADP102-0177
  \HMTTitleMapEntry{Soporte operatorio común de \(3\) lecturas geométricas}{Conjugación de las realizaciones aritmética, angular y torsional mediante el operador cuadrático común}% ADP102-0178
  \HMTTitleMapEntry{Cadena constitutiva del sector gravitatorio: escala cronológica \(108\), reducción tensorial, holonomía y evaluación metrológica de \(G\)}{Composición causal de los invariantes de fase, acción y longitud en la derivación de G}% ADP102-0179
  \HMTTitleMapEntry{Cardinalidades de incidencia 90 y 120}{Derivación trítica de las cardinalidades 90/120 y realización posterior en la incidencia hexada–octada}% ADP102-0180
  \HMTTitleMapEntry{Separación causal entre el espectro de área, la memoria modular y la realización Cartan--Holst}{Distinción tipada entre Barbero–área, área–Hamiltoniano modular y realización Cartan–Holst}% ADP102-0182
  \HMTTitleMapEntry{Transductor de Hensel no filtrado y desplazamiento completo sobre las fibras ternarias}{Transductor integral de Hensel como operador de desplazamiento completo}% ADP102-0183
  \HMTTitleMapEntry{Sobreyectividad de la carta arquimediana no filtrada sobre \(\mathbb R\)}{Cobertura exacta de la carta real antecedente mediante secciones compatibles}% ADP102-0184
  \HMTTitleMapEntry{2 ejes de clasificación: valor visible y genealogía de transporte}{Clasificación bifactorial del número por valor arquimediano y genealogía de transporte}% ADP102-0185
  \HMTTitleMapEntry{Posición de \texorpdfstring{\(\pi,e,\varphi\) y \(\alpha\)}{pi, e, phi y alpha} en la clasificación}{Posición de π,e,φ y α en la clasificación}% ADP102-0186
  \HMTTitleMapEntry{Carta de Hensel no filtrada, supervivencia y frontera global}{Carta antecedente, supervivencia coinductiva y construcción de la frontera global}% ADP102-0187
  \HMTTitleMapEntry{La familia operatorial de Euler como gramática de cierre}{Factorización del cierre genealógico del número mediante el sistema operatorio de Euler}% ADP102-0188
  \HMTTitleMapEntry{Dependencias matemáticas.}{Dependencia causal entre sección inversa, evaluación arquimediana y conservación excepcional de la incidencia}% ADP102-0189
  \HMTTitleMapEntry{\(4\) exigencias de elección en el sistema proyectivo}{Condiciones de elección necesarias para la prolongación global de secciones compatibles}% ADP102-0190
  \HMTTitleMapEntry{Decidibilidad efectiva de cada ventana finita de multiplicidades proyectivas}{Decidibilidad algorítmica de cada nivel finito del sistema de prefijos}% ADP102-0191
  \HMTTitleMapEntry{Relación exacta con las \(9\) fases del Topological Phase Kernel}{Restricción de la obstrucción uniforme al sistema temporal nonádico de 9 fases}% ADP102-0192
  \HMTTitleMapEntry{De la consistencia constructible a la independencia por forcing}{Consistencia relativa en L e independencia mediante forzamiento}% ADP102-0193
  \HMTTitleMapEntry{Optimización minimax sobre ventanas finitas de prefijos compatibles}{Teorema minimax en cada nivel finito del poset de prefijos compatibles}% ADP102-0194
  \HMTTitleMapEntry{La cadena que el puente recibe demostrada}{Antecedentes teoremáticos del funtor de transducción genealógico–metrológica}% ADP102-0195
  \HMTTitleMapEntry{Derivación bidireccional del exponente electrónico desde la ruta enriquecida}{Banda de realización del electrón y separación entre masa basal y clausura beta}% ADP102-0196
  \HMTTitleMapEntry{\(2\) ramas causales y su confluencia}{Confluencia de las ramas determinantal y dodecafásica de la escala de acción}% ADP102-0197
  \HMTTitleMapEntry{Carta metrológica posterior y evaluación dimensional de los invariantes generados}{Sistemas de transducción entre número genealógico y magnitud física}% ADP102-0198
  \HMTTitleMapEntry{Principio de Ambivalencia: involución de los lectores areal y volumétrico, equivalencia local y orientación temporal conjugada}{Principio de Ambivalencia en Mecánica Dimensional: involución areal–volumétrica, sistemas autoinerciales y orientación temporal conjugada}% ADP102-0199
  \HMTTitleMapEntry{Condiciones de una evaluación prospectiva}{Condiciones de evaluación prospectiva del Principio de Ambivalencia sobre el estado holonómico integral}% ADP102-0200
  \HMTTitleMapEntry{Cristal temporal finito HMT: drive, subarmonía y estabilidad espectral}{Cristal temporal finito HMT: Hamiltoniano periódicamente forzado, respuesta subarmónica y estabilidad espectral}% ADP102-0201
  \HMTTitleMapEntry{Drive construido desde el reloj de 108 estados}{Construcción del forzamiento periódico a partir del sistema temporal de 108 estados}% ADP102-0202
  \HMTTitleMapEntry{Estabilidad espectral finita}{Jet nilpotente de la evolución triádica y prolongación compatible}% ADP102-0203
  \HMTTitleMapEntry{Teorema de confluencia HMT–MD: forma electrogravitatoria, sistema Einstein–Schrödinger y coordinación de acción, materia, campos, geometría e información}{Teorema de confluencia HMT–MD: sistema Einstein–Schrödinger y coordinación electrogravitatoria de acción, materia, campos, geometría e información}% ADP102-0204
  \HMTTitleMapEntry{Independencia horizontal y compatibilidad vertical}{Naturalidad de las realizaciones y conservación vertical de los invariantes del estado holonómico integral}% ADP102-0205
  \HMTTitleMapEntry{1.ª confluencia: acoplamiento, vacío y sector electrodébil}{Acoplamiento constitutivo entre el vacío y el sector electrodébil}% ADP102-0206
  \HMTTitleMapEntry{2.ª confluencia: acción, gravedad y forma electrogravitatoria}{Reciprocidad entre acción y gravedad en la forma electrogravitatoria}% ADP102-0207
  \HMTTitleMapEntry{3.ª confluencia: geometría, información y mezcla}{Confluencia entre geometría, información y mezcla}% ADP102-0208
  \HMTTitleMapEntry{4.ª confluencia: pantalla, torsión y memoria helicoidal}{Transducción de la memoria helicoidal en la realización torsional}% ADP102-0209
  \HMTTitleMapEntry{Del determinante electrodébil al operador constitutivo positivo del vacío y su completación espectral \(3\to4\)}{Del determinante electrodébil al operador constitutivo positivo del vacío y su completación espectral 3→4}% ADP102-0210
  \HMTTitleMapEntry{Valor publicado y memoria de la fibra}{Desintegración espectral de la medida de trayectorias y condicionamiento de futuros}% ADP102-0211
  \HMTTitleMapEntry{Fundamentos algebraicos del cosido de fase y de la automedida}{Estado conjunto observador–observable, instrumento reversible y automedida}% ADP102-0212
  \HMTTitleMapEntry{Información par y transporte geométrico impar en la bicapa}{Descomposición par–impar del estado bicapa: información interhoja y transporte geométrico orientado}% ADP102-0214
  \HMTTitleMapEntry{\texorpdfstring{\(3\)}{3} realizaciones de la memoria genealógica: transporte dodecafásico, función de Green no local con generador local y expansión cosmológica}{Realizaciones dinámica, propagatoria y cosmológica de la memoria genealógica: transporte dodecafásico, función de Green y expansión espacial}% ADP102-0215
  \HMTTitleMapEntry{Cadena constitutiva del sector gravitatorio: escala cronológica \texorpdfstring{\(108\)}{108}, reducción tensorial, holonomía y evaluación metrológica de \texorpdfstring{\(G\)}{G}}{Descenso icosaédrico A₅/C₅ y selección del portador gravitatorio}% ADP102-0216
  \HMTTitleMapEntry{Entrelazador con tensores simétricos sin traza y reducción transversal \texorpdfstring{\(12\to5\to5\to2\)}{12→5→5→2}}{Entrelazador con tensores simétricos sin traza y reducción transversal 12→5→5→2}% ADP102-0217
  \HMTTitleMapEntry{Criterios estructurales del selector gravitatorio: transductor dimensional, holonomía enriquecida y separación entre selectores torsional y circular}{Selector holonómico gravitatorio: transducción dimensional, naturalidad, refinamiento y separación de los regímenes torsional y circular}% ADP102-0218
  \HMTTitleMapEntry{Carta decimal declarada y contraste metrológico}{Reconocimiento metrológico posterior de G e independencia causal de la derivación}% ADP102-0219
  \HMTTitleMapEntry{Tétrada, conexión y \(2\) defectos}{Tétrada, conexión de espín y torsión de la realización Cartan–Holst}% ADP102-0220
  \HMTTitleMapEntry{Acción de 1.er orden}{Acción de Palatini–Holst y ecuaciones variacionales}% ADP102-0221
  \HMTTitleMapEntry{Confluencia horizonte–área–información–entropía–acción}{Teorema de confluencia entre energía de horizonte, área, información, entropía y acción}% ADP102-0222
  \HMTTitleMapEntry{\(2\) orientaciones de una misma propagación}{Propagadores avanzado y retardado sobre rutas orientadas}% ADP102-0223
  \HMTTitleMapEntry{Órbita elíptica, acción y \(3\) invariantes J}{Invariantes de la involución temporal y conjugación partícula–antipartícula}% ADP102-0224
  \HMTTitleMapEntry{Cosmología de periodicidad \texorpdfstring{\(108\)}{108}: inversión de hoja, horizonte, Big Bang, rebote torsional y sector oscuro}{Cosmología torsional de periodicidad 108: expansión, retorno con memoria y rebote de Einstein–Cartan}% ADP102-0225
  \HMTTitleMapEntry{Separación tipológica de las \(2\) realizaciones cosmológicas}{Separación entre aceleración de de Sitter y rebote torsional de Einstein–Cartan}% ADP102-0226
  \HMTTitleMapEntry{\(3\) foliaciones exactas de de Sitter}{Foliaciones cosmológicas compatibles con la inversión de hoja}% ADP102-0227
  \HMTTitleMapEntry{Universo con frontera}{Estado de frontera del horizonte cosmológico y memoria de retorno}% ADP102-0228
  \HMTTitleMapEntry{Datos espectrales de una realización supersimétrica: \((\mathcal H_{\rm phys},H,\mathcal Q)\), supercarga y compactificación en dimensión \(11\)}{Realización de teoría M mediante compactificación, cuerdas, branas y dimensión 11}% ADP102-0229
  \HMTTitleMapEntry{Ley General de Masas sobre el atlas \texorpdfstring{\(81/56/13/324\)}{81/56/13/324}: trayectorias, registro, firma, carácter, torres \texorpdfstring{\(90/120\)}{90/120} y sector nulo}{Ley General de Masas sobre el atlas 81/56/13/324: trayectorias, registro, firma, carácter, semigrupo armónico ⟨90,120⟩ y sector nulo}% ADP102-0230
  \HMTTitleMapEntry{Algoritmo universal de realización material e igualador espectral de los lectores \(K_D\) y \(K_\Omega\)}{Construcción functorial de las familias de masa sobre el atlas}% ADP102-0231
  \HMTTitleMapEntry{Bifurcación nula y positividad}{Bifurcación del sector nulo anterior a la rama positiva de masa}% ADP102-0232
  \HMTTitleMapEntry{Portal del dominio exhaustivo}{Construcción exhaustiva del atlas 81/56/13/324 y conservación del sector nulo}% ADP102-0233
  \HMTTitleMapEntry{Dominio interno y problema de realización}{Teorema de exhaustividad del dominio material y estabilidad del sector nulo bajo refinamiento}% ADP102-0234
  \HMTTitleMapEntry{Producto fibrado entre descriptores y rutas}{Producto fibrado entre invariantes de especie y rutas genealógicas}% ADP102-0235
  \HMTTitleMapEntry{Lectores bidireccionales \texorpdfstring{\(K_D\)}{K D} y \texorpdfstring{\(K_\Omega\)}{K Omega}: perfiles, coincidencias, torres \texorpdfstring{\(90/120\)}{90/120} y operadores por multisección}{Lectores bidireccionales \(K_D\) y \(K_\Omega\): perfiles, coincidencias, semigrupo armónico \(\langle90,120\rangle\) y operadores por multisección}% ADP102-0236
  \HMTTitleMapEntry{Lectores \(K_D\) y \(K_\Omega\): perfiles dodecafásicos, dominio de rutas y codominio espectral}{Lectores \(K_D\) y \(K_\Omega\): perfiles dodecafásicos, dominio de rutas y codominio espectral}% ADP102-0237
  \HMTTitleMapEntry{Realizaciones espectrales, escalares y nulas de la Ley General de Masas: familias materiales, transporte CKM--PMNS y sectores fotónico y gluónico}{Realizaciones espectrales, escalares y nulas de la Ley General de Masas: familias materiales, mezcla quark CKM, mezcla neutrínica PMNS y sectores nulos fotónico y gluónico}% ADP102-0238
  \HMTTitleMapEntry{Color, sabores quark y mezcla CKM}{Operadores de masa por fibra, órbita y multisección del atlas}% ADP102-0239
  \HMTTitleMapEntry{Estados compuestos y ley de enlace}{Correspondencia por dominio y codominio entre el inventario material y las realizaciones del atlas}% ADP102-0240
  \HMTTitleMapEntry{Clasificador de especies materiales por fibra, representación, carga, espín, sabor y generación}{Clasificación de las especies materiales por fibra, representación, carga, espín, sabor y generación}% ADP102-0241
  \HMTTitleMapEntry{Tipado probatorio de persistencia, atlas y familias.}{Condiciones de persistencia y correspondencia entre atlas, familias y representaciones}% ADP102-0242
  \HMTTitleMapEntry{Ontología material, carga, espín, estadística, color y sabor}{Definiciones estructurales de partícula, masa, carga, espín, estadística, color, sabor y generación}% ADP102-0243
  \HMTTitleMapEntry{Teorema de cierre tipado de la realización másica}{Teorema de clausura de la realización material mediante secciones persistentes y representaciones finitas}% ADP102-0244
  \HMTTitleMapEntry{Contraste metrológico y archivo exhaustivo de realizaciones}{Reconocimiento metrológico posterior e inventario exhaustivo de las realizaciones materiales}% ADP102-0245
  \HMTTitleMapEntry{Arquitectura natural común de las realizaciones terminales del continuo genealógico}{Teorema de factorización conjunta de las \(5\) realizaciones terminales a través de la estructura discreta del continuo}% ADP102-0246
  \HMTTitleMapEntry{Publicación terminal del continuo único mediante \(5\) lectores no permutables}{Publicación terminal mediante 5 lectores no permutables y factor común del estado holonómico integral}% ADP102-0247
  \HMTTitleMapEntry{Clausuras dinámicas del continuo: transporte solenoidal y cociente de calibre}{Dinámica y curvatura del continuo genealógico anteriores a las realizaciones clásicas}% ADP102-0248
  \HMTTitleMapEntry{Dinámica y curvatura del continuo genealógico}{Factorización conjunta de las realizaciones terminales desde el mismo estado del continuo}% ADP102-0249
  \HMTTitleMapEntry{Clausuras algebraicas del continuo y límites de la reconstrucción extensional}{Factorización común, consistencia extensional y límites de la reconstrucción terminal}% ADP102-0250
  \HMTTitleMapEntry{El ejemplo de los ordinales de von Neumann}{Representación ordinal de von Neumann como modelo de consistencia extensional}% ADP102-0251
  \HMTTitleMapEntry{Dominio, mapas y compatibilidades del objeto terminal común}{Vista interna \(G_T^{\mathrm{int}}\) del continuo y aplicación \(\operatorname{App}_{T,d_T}\): conservación de fibra, acarreo, memoria y terminal}% ADP102-0252
  \HMTTitleMapEntry{Despliegue interno de las realizaciones correlativas anterior a su identificación con problemas clásicos}{Emergencia correlativa de las 5 acciones terminales sobre el mismo estado del continuo}% ADP102-0253
  \HMTTitleMapEntry{\(5\) lectores terminales no permutables del continuo genealógico}{Teorema de no permutabilidad de los 5 lectores terminales y clasificación de sus codominios}% ADP102-0254
  \HMTTitleMapEntry{Obstrucciones a la proyección ablativa del continuo: pérdida de fase, orientación o memoria genealógica}{Proyectores positivos de la forma de Guinand–Weil y localización espectral}% ADP102-0255
  \HMTTitleMapEntry{Fijación de la recta crítica por la involución funcional \(s\mapsto1-\bar s\)}{Fijación de la recta crítica por la involución funcional \texorpdfstring{\(s\mapsto 1-\bar{s}\)}{s -> 1 - conjugado(s)}}% ADP102-0256
  \HMTTitleMapEntry{Separación tipada de los ceros de \(\zeta\), \(\xi\) y el determinante espectral}{Separación tipada de los ceros de ζ, ξ y el determinante espectral}% ADP102-0257
  \HMTTitleMapEntry{Representante positivo de la clase residual nula: \(n\equiv0\pmod 9\) y \(\rho_9(n)=9\)}{Representante positivo de la clase residual nula: \(n\equiv0\pmod 9\) y \(\rho_9(n)=9\)}% ADP102-0258
  \HMTTitleMapEntry{Isometría de transporte torcido \(\widetilde U_k\) sobre estados cilíndricos, pesos proyectivos y memoria de Weil}{Isometría de transporte torcido \(\widetilde U_k\) sobre estados cilíndricos, pesos proyectivos y memoria de Weil}% ADP102-0259
  \HMTTitleMapEntry{Invariante graduado del corredor espectral \(4\to2\to6\to54\)}{Invariante graduado del corredor espectral 4→2→6→54}% ADP102-0260
  \HMTTitleMapEntry{Traza acoplada del espectro de Weil y la graduación Moonshine}{Construcción independiente del campo espectral sobre el doble círculo}% ADP102-0261
  \HMTTitleMapEntry{Dominio hidrodinámico y publicación Galerkin}{Dominio solenoidal y conservación de la divergencia nula}% ADP102-0262
  \HMTTitleMapEntry{Dato único de partida y núcleos densos}{Datos iniciales, compatibilidad de frontera y propagación regular}% ADP102-0263
  \HMTTitleMapEntry{Realización finita de la desigualdad de Kalman–Yakubovich–Popov y función auxiliar en la clausura energética}{Desigualdad KYP y control coercivo de la evolución solenoidal}% ADP102-0264
  \HMTTitleMapEntry{Extracción material de la cota de Sobolev}{Estimación de Sobolev y exclusión de la concentración singular}% ADP102-0265
  \HMTTitleMapEntry{La fibra de incidencia es \(1\) coordenada física del estado completo}{Curvatura de fibra y coercividad de la acción de calibre}% ADP102-0266
  \HMTTitleMapEntry{Vacío simple y testigo que alcanza la brecha}{Unicidad del vacío invariante de calibre y saturación de la brecha \(3\kappa_{\mathrm{av}}\) por el autovector \(W_c\)}% ADP102-0267
  \HMTTitleMapEntry{\(4\) reconstrucciones de Osterwalder–Schrader de una dinámica común}{Positividad de Osterwalder–Schrader y reconstrucción cuántica}% ADP102-0268
  \HMTTitleMapEntry{Teoría gauge euclídea y brecha de Yang--Mills}{Brecha espectral del sector de calibre y límite proyectivo}% ADP102-0269
  \HMTTitleMapEntry{Presentación generador a generador de la carta nueva}{Generadores cohomológicos y proyectores compatibles}% ADP102-0270
  \HMTTitleMapEntry{Fibra inicial, Mittag--Leffler y publicación afirmativa}{Lefschetz–Murre y descenso algebraico de las clases de Hodge}% ADP102-0271
  \HMTTitleMapEntry{Bocksteins de todas las páginas}{Homomorfismos de Bockstein de la sucesión espectral: compatibilidad entre páginas y determinación del orden en T del complejo universal}% ADP102-0272
  \HMTTitleMapEntry{Tabla completa Kato--FERS antes del determinante}{Determinante de Kato–FERS, rango y término principal de BSD}% ADP102-0273
  \HMTTitleMapEntry{Separación entre estructura y aplicación}{Vista interna \(G_T\) y lector terminal \(P_T\): independencia lógica entre existencia estructural y realización clásica}% ADP102-0274
  \HMTTitleMapEntry{Obstrucciones comunes a las \(5\) publicaciones terminales: pérdida de dominio, naturalidad, certificado o límite compatible}{Obstrucciones comunes de los lectores terminales: pérdida de dominio, naturalidad o límite compatible}% ADP102-0275
  \HMTTitleMapEntry{Formalización asistida del núcleo algebraico}{Formalización verificable del núcleo algebraico y certificación de identidades finitas}% ADP102-0276
  \HMTTitleMapEntry{Correspondencia entre enunciados, pruebas, certificados y fuentes teoremáticas científicos}{Correspondencia entre enunciados, demostraciones, certificados y fuentes científicas primarias}% ADP102-0277
  \HMTTitleMapEntry{Convenciones de tipado y representación de datos}{Convenciones matemáticas y representación de los datos}% ADP102-0278
  \HMTTitleMapEntry{Heterogeneidad tipada y cobertura del catálogo}{Heterogeneidad de los dominios físicos y exhaustividad del catálogo}% ADP102-0279
  \HMTTitleMapEntry{Fichas individuales de las rutas orientadas}{Despliegue íntegro de las 324 rutas orientadas: correspondencia entre 53 campos canónicos, 57 campos enriquecidos y 49 invariantes publicados por sección}% ADP102-0280
  \HMTTitleMapEntry{Inventario universal tipado de 471 masas y 384 anchuras}{Inventario exhaustivo de 471 masas y 384 anchuras: identificación, procedencia y correspondencia con el atlas}% ADP102-0281
  \HMTTitleMapEntry{Fuente, deduplicación y límites de tipo}{Procedencia de los datos, unicidad de los registros y dominios de validez}% ADP102-0282
  \HMTTitleMapEntry{Código de objeto y lector HMT--MD}{Identificación de objetos y aplicación del lector HMT--MD}% ADP102-0283
  % Correspondencias espirales: sustituciones científicas unívocas aprobadas
  % por cotejo focal. Se modifican sólo los rótulos; los cuerpos permanecen
  % íntegros en sus residencias distribuidas.
  \HMTTitleMapEntry{TPK: selección, transporte, actualización y retorno}{Composición dinámica del Topological Phase Kernel sobre el estado holonómico integral: selección trítica, transporte cardinal, actualización de memoria y retorno tipado}% ADP102-0284
  \HMTTitleMapEntry{La transición tipada}{Definición y factorización del operador de transición enriquecida \(U_t\)}% ADP102-0285
  \HMTTitleMapEntry{Composición del acarreo}{Ley de cociclo bimodular del acarreo bajo concatenación de rutas}% ADP102-0286
  \HMTTitleMapEntry{Calendario dodecafásico y no escisión}{Extensión cíclica no escindida \(0\to C_{12}\to C_{108}\to C_9\to0\) y cronología interna del TPK}% ADP102-0287
  \HMTTitleMapEntry{Hélice nonádica de fase y profundidad}{Recubrimiento helicoidal de la fase nonádica y coordenada integral de profundidad de memoria}% ADP102-0288
  \HMTTitleMapEntry{Inmersión geométrica normalizada}{Inmersión \(H_9\) del cociente--residuo nonádico en \(\mathbb R^3\) y cálculo de los invariantes de Frenet}% ADP102-0289
  \HMTTitleMapEntry{Círculos anidados y cilindros}{Secciones circulares y foliación cilíndrica del recubrimiento fase--memoria}% ADP102-0290
  \HMTTitleMapEntry{Hojas orbitales y costura del espacio de suspensión}{Órbitas del producto sesgado y relación de identificación del espacio de suspensión}% ADP102-0291
  \HMTTitleMapEntry{Caso estacionario afín}{Iteración del endomorfismo afín estacionario y clasificación de sus puntos fijos}% ADP102-0292
  \HMTTitleMapEntry{Inversión de rutas}{Antiinvolución del subgrupoide de rutas reversibles y conjugación de la holonomía afín}% ADP102-0293
  \HMTTitleMapEntry{Espejo de constantes y eje fijo}{Involución \(\pi/e\) de la fibra de constantes, invariante \(u^\pi+u^e\) y eje fijo \(\varphi\)}% ADP102-0294
  \HMTTitleMapEntry{La geometría visible como publicación del estado holonómico integral}{Evaluación arquimediana del estado holonómico integral y defecto de fidelidad genealógica}% ADP102-0295
  \HMTTitleMapEntry{La inversión del orden explicativo}{Factorización causal desde la historia enriquecida hasta la evaluación arquimediana}% ADP102-0296
  \HMTTitleMapEntry{Estado, historia y publicación}{Espacio de estados holonómicos integrales, categoría de historias TPK y evaluación arquimediana}% ADP102-0297
  \HMTTitleMapEntry{Circunferencia intrínseca y circunferencia proyectada}{Distinción entre órbita radial intrínseca y proyección circular de una trayectoria helicoidal}% ADP102-0298
  \HMTTitleMapEntry{La célula nonádica y el nacimiento bidimensional}{Construcción del soporte nonádico bidimensional \(X_9=(\mathbb Z/9\mathbb Z)^2\)}% ADP102-0299
  \HMTTitleMapEntry{Antecedentes APP de la clasificación radial}{Invariantes APP que inducen el carácter radial acumulado}% ADP102-0300
  \HMTTitleMapEntry{Estado ternario levantado}{Estado TRIT enriquecido por residuo y acarreo: multiplicidad de los representantes del umbral}% ADP102-0301
  \HMTTitleMapEntry{Álgebras cuadráticas}{Familia \(\mathcal A_\tau=\mathbb R[J]/(J^2+\tau I)\) y clasificación elíptica, parabólica e hiperbólica}% ADP102-0302
  \HMTTitleMapEntry{Doble proyección}{Factorización conjunta de las publicaciones arquimediana y excepcional del estado holonómico integral}% ADP102-0303
  \HMTTitleMapEntry{Representación local de una transición}{Representación afín \(e\mapsto h_e\) de los generadores de la categoría de rutas}% ADP102-0304
  \HMTTitleMapEntry{Holonomías prefijas}{Productos afines ordenados \texorpdfstring{\(H_j=h_{e_j}\cdots h_{e_1}\)}{H_j = h_ej ... h_e1} y conservación de prefijos}% ADP102-0305
  \HMTTitleMapEntry{Categoría levantada y caracteres elementales}{Categoría de rutas enriquecidas por 2 orientaciones TRIT y definición de los caracteres \(p\) y \(a\)}% ADP102-0306
  \HMTTitleMapEntry{Descomposición balanceada}{Descomposición única \(p=r_3(p)+3c_3(p)\) en residuo balanceado y acarreo}% ADP102-0307
  \HMTTitleMapEntry{Involución}{Acción de \((\tau_+,\tau_-)\mapsto(-\tau_-,-\tau_+)\) sobre los caracteres \(p\) y \(a\)}% ADP102-0308
  \HMTTitleMapEntry{Definición y dominio}{Definición de \(L_p\) sobre \(\mathbb R_{>0}\) y dominio de la inversa \(\exp_p\)}% ADP102-0309
  \HMTTitleMapEntry{Adición deformada}{Ley de composición \(L_p(xy)=L_p(x)+L_p(y)+pL_p(x)L_p(y)\)}% ADP102-0310
  \HMTTitleMapEntry{No exhaustividad y alcance exacto}{Dominio clasificatorio de \(p\in\{2,1,0,-1,-2\}\) y extensión afín general}% ADP102-0311
  \HMTTitleMapEntry{Obstrucción al lector reducido}{Defecto de naturalidad del lector reducido \((p,a,\Lambda,C,m,r,\theta)\)}% ADP102-0312
  \HMTTitleMapEntry{Definición del lector completo}{Lector radial enriquecido \(\mathcal P^{\mathrm{enr}}_{\mathrm{rad},N}\): dominio, codominio y conservación de prefijos}% ADP102-0313
  \HMTTitleMapEntry{Conservación como trazabilidad}{Conservación dinámica del acarreo, la frontera y la memoria por el lector radial}% ADP102-0314
  \HMTTitleMapEntry{Subdivisión cuádruple}{Naturalidad del lector bajo el refinamiento cuádruple \(\operatorname{sd}_4\)}% ADP102-0315
  \HMTTitleMapEntry{Clasificación genealógica y no fidelidad de la publicación}{Clasificación por la firma radial plena y defecto de fidelidad de la evaluación arquimediana}% ADP102-0316
  \HMTTitleMapEntry{Clasificación por 3 incrementos visibles}{Clasificación cinemática mediante los incrementos \((\omega,\beta,\mu)\)}% ADP102-0317
  \HMTTitleMapEntry{3 fibras de una misma evaluación visible}{\(3\) clases genealógicas con una misma imagen arquimediana}% ADP102-0318
  \HMTTitleMapEntry{Ley de tornillo}{Identidad espectral \(V^{-1}\mathscr S V=(2+\sqrt3)\mathrm i\mathscr S\) y ecuación de la suspensión logarítmica}% ADP102-0319
  \HMTTitleMapEntry{Modo trítico común}{Modo singular común de las hojas APP en la carta canónica \(1,\ldots,9\)}% ADP102-0320
  \HMTTitleMapEntry{Ledger material}{Registro genealógico de la ruta material: extensión superviviente, firma entera y carácter adimensional}% ADP102-0321
  \HMTTitleMapEntry{Publicación dimensional posterior}{Realización del operador constitutivo del vacío en los campos \((E_\sigma,B_\sigma)\)}% ADP102-0322
  \HMTTitleMapEntry{Curva plana polar}{Curvatura y longitud de arco de una curva polar \(r=r(\theta)\)}% ADP102-0323
  \HMTTitleMapEntry{Especializaciones}{Evaluación de \(\kappa_{\mathrm F}\) en las familias logarítmica, arquimediana, de Fermat, hiperbólica y lituus}% ADP102-0324
  \HMTTitleMapEntry{Suspensión espacial}{Inmersión espacial \(\Gamma(\theta)\) y criterios diferenciales de Frenet--Lancret}% ADP102-0325
  \HMTTitleMapEntry{Tricromía geométrica}{Coordinación de los regímenes \(J^2=-I,0,+I\) con la geometría de Frenet}% ADP102-0326
  \HMTTitleMapEntry{Conclusión}{Teorema de cierre de las correspondencias espirales genealógicas}% ADP102-0327
  \HMTTitleMapEntry{Semiproducto afín y composición de rutas}{Monoide afín \(\operatorname{End}(U_{\mathrm{APP}})\ltimes U_{\mathrm{APP}}\) y composición ordenada de rutas}% ADP102-0328
  \HMTTitleMapEntry{Tabla exhaustiva de la carta local de 2 orientaciones}{Enumeración exhaustiva de la carta local \((\tau_+,\tau_-)\mapsto(p,a,r_3,c_3)\)}% ADP102-0329
  \HMTTitleMapEntry{5 regiones de \texorpdfstring{\(\pi\)}{pi} y no fidelidad de la evaluación}{Descomposición de la microfibra de \(\pi\) en 5 regiones y multiplicidad de la evaluación arquimediana}% ADP102-0330
  \HMTTitleMapEntry{Álgebra potencia--logarítmica y 5 publicaciones canónicas}{Álgebra de \(L_p\) y clasificación de las 5 familias canónicas \(p\in\{2,1,0,-1,-2\}\)}% ADP102-0331
  \HMTTitleMapEntry{Derivación de las 5 curvas canónicas}{Evaluación de \(\exp_p(u_0+\beta m)\) en las 5 familias canónicas}% ADP102-0332
  \HMTTitleMapEntry{Inversión de hojas}{Conjugación \(p\leftrightarrow-p\) de las hojas potencia--logarítmicas}% ADP102-0333
  \HMTTitleMapEntry{Sector afín general}{Clasificación del sector afín por la pareja \((p,[\Lambda,C])\)}% ADP102-0334
  \HMTTitleMapEntry{Naturalidad, equivariancia y no fidelidad}{Naturalidad bajo truncamiento, equivariancia nonádica y defecto de fidelidad arquimediana}% ADP102-0335
  \HMTTitleMapEntry{El lector de prefijos}{Lector de prefijos afines y conservación del registro enriquecido}% ADP102-0336
  \HMTTitleMapEntry{Refinamientos no triviales}{Criterio de compatibilidad para refinamientos generales de rutas}% ADP102-0337
  \HMTTitleMapEntry{Equivarianza, no invariancia}{Equivarianza del cociclo radial bajo la holonomía \(\Gamma_9\)}% ADP102-0338
  \HMTTitleMapEntry{3 demostraciones de no fidelidad}{\(3\) obstrucciones explícitas a la fidelidad de la evaluación arquimediana}% ADP102-0339
  \HMTTitleMapEntry{Mismo extremo, distinto orden afín}{Coincidencia del producto final y dependencia del orden de los factores afines}% ADP102-0340
  \HMTTitleMapEntry{Misma curva, distinta región genealógica}{Coincidencia de la imagen arquimediana entre regiones genealógicas distintas}% ADP102-0341
  \HMTTitleMapEntry{Clasificación por nivel matemático}{Estratificación por dominio: historia, lector, evaluación arquimediana y realización física}% ADP102-0342
  \HMTTitleMapEntry{Clasificación por giro, radio y memoria}{Clasificación por los incrementos \((\omega,\beta,\mu)\)}% ADP102-0343
  \HMTTitleMapEntry{Realizaciones posteriores y controles de tipo}{Realizaciones electromagnética y material con separación de dominios y codominios}% ADP102-0344
  \HMTTitleMapEntry{Centro electrónico y lector posterior}{Compatibilidad del lector radial con el punto fijo electrónico \((5,5)\) y el atlas \(81/56/13/324\)}% ADP102-0345
  \HMTTitleMapEntry{Matriz cerrada de controles negativos}{Matriz de condiciones de validez por nivel causal}% ADP102-0346
  \HMTTitleMapEntry{Primera prolongación y frontera coinductiva}{Prolongación coinductiva inicial desde \(w_6\) hasta la frontera de compatibilidad de orden \(1\)}% ADP102-0347
  \HMTTitleMapEntry{Vía B: núcleo de vacancias \texorpdfstring{\(E_{21/22}\)}{E21/22}}{Caracterización de \(\alpha\) como raíz simple del núcleo analítico \(E_{21/22}\)}% ADP102-0348
  % Purismo residual: rótulos activos localizados en el índice del sucesor.
  % Las sustituciones siguientes no alteran el nivel jerárquico ni el cuerpo.
  \HMTTitleMapEntry{Cronología observable de los dos cursores}{Evolución cronológica del par de cursores y determinación del estado observable}% ADP102-0349
  \HMTTitleMapEntry{Emisor finito de dos cursores}{Emisor finito asociado al par de cursores y generación de estados admisibles}% ADP102-0350
  \HMTTitleMapEntry{De una célula a dos coordenadas: nacimiento de APP}{Construcción de APP mediante la descomposición de 1 célula en 2 coordenadas}% ADP102-0351
  \HMTTitleMapEntry{La rama de \(20\) bloques}{Prolongación de la sección superviviente mediante 20 bloques compatibles}% ADP102-0352
  \HMTTitleMapEntry{Del prefijo a un punto real}{Evaluación arquimediana del prefijo compatible y determinación del punto real}% ADP102-0353
  \HMTTitleMapEntry{Una emisión reproducible de doce tríadas}{Emisión reproducible de 12 tríadas y reconstrucción de su estado holonómico integral}% ADP102-0354
  \HMTTitleMapEntry{Construcción del carácter de clausura y trayectorias genealógicas estructural de \(\pi\)}{Construcción del carácter de clausura y de las trayectorias genealógicas de \(\pi\)}% ADP102-0355
  \HMTTitleMapEntry{La pendiente 1/5 forzada por las \(5\) regiones}{Determinación de la pendiente 1/5 mediante el balance de las 5 regiones orientadas}% ADP102-0356
  \HMTTitleMapEntry{Ventana dodecafásica del estado firmado y clausura reversible del ciclo}{Restricción finita del estado dodecafásico firmado y clausura reversible del ciclo}% ADP102-0357
  \HMTTitleMapEntry{Ventana dodecafásica finita}{Dominio finito de la clausura dodecafásica}% ADP102-0358
  \HMTTitleMapEntry{Prolongación remota \texorpdfstring{\(\mathsf T_{\alpha,\infty}\)}{T alpha infinito} de la vía entera}{Prolongación coinductiva \(\mathsf T_{\alpha,\infty}\) de la clausura dodecafásica entera}% ADP102-0359
  \HMTTitleMapEntry{Teorema de los cuatro canales enteros}{Teorema de clasificación de los 4 canales enteros}% ADP102-0360
  \HMTTitleMapEntry{Primer codificador intrínseco y su límite}{Codificador intrínseco inicial y existencia de su límite}% ADP102-0361
  \HMTTitleMapEntry{Hipótesis de tipado y controles negativos}{Condiciones de tipado y criterios de incompatibilidad de la teoría de vacancias nonádicas}% ADP102-0362
  \HMTTitleMapEntry{Vía cronológica de longitud, acción y espectro}{Composición cronológica de longitud, acción y espectro en la determinación de \(G\)}% ADP102-0363
  \HMTTitleMapEntry{Operador torcido y control negativo determinista}{Operador con torsión y obstrucción determinista a la continuación multiplicativa}% ADP102-0364
  \HMTTitleMapEntry{Estratificación matemática del objeto Majorana en cuatro niveles}{Estratificación matemática del objeto Majorana en 4 niveles}% ADP102-0365
  \HMTTitleMapEntry{Seis caracteres y cuatro torsiones de trenza}{Clasificación de 6 caracteres por 4 clases de torsión de trenza}% ADP102-0366
  \HMTTitleMapEntry{obstrucciones y condiciones de no degeneración específicos}{Obstrucciones específicas y condiciones de no degeneración}% ADP102-0367
  \HMTTitleMapEntry{Quark arriba y antiquark arriba}{Sectores quark \(u\) y antiquark \(\bar u\)}% ADP102-0368
  \HMTTitleMapEntry{Quark abajo y antiquark abajo}{Sectores quark \(d\) y antiquark \(\bar d\)}% ADP102-0369
  \HMTTitleMapEntry{Controles de refutación.}{Condiciones de validez e incompatibilidad de la realización equivariante de especies}% ADP102-0370
  \HMTTitleMapEntry{El cuadrado local de \texorpdfstring{\(e\)}{e} y la ruptura de memoria}{Cuadrado posicional local de \(e\) y diferenciación genealógica de estados con igual bloque visible}% ADP102-0371
  \HMTTitleMapEntry{La ley general antes de sus realizaciones por especie}{Antecedencia causal de la Ley General de Masas respecto de las realizaciones por especie}% ADP102-0372
  \HMTTitleMapEntry{Fundamento unificado de la masa}{Factorización extensión–ruta y conservación del registro genealógico}% ADP102-0373
  \HMTTitleMapEntry{La banda de ruta}{Definición de la banda de ruta mediante la órbita de una sección persistente}% ADP102-0374
  \HMTTitleMapEntry{El canal de incidencia es curvatura local, no un factor auxiliar}{Interpretación del canal de incidencia como curvatura local del sistema de calibre}% ADP102-0375
  % Publicaciones terminales del capítulo 95: se conserva cada desarrollo
  % completo y se sustituye la enumeración narrativa por objeto, operación y resultado.
  \HMTTitleMapEntry{1.ª resolución: Riemann y la positividad como complemento construido}{Publicación espectral–polar del continuo mediante el complemento positivo de Guinand–Weil}% ADP102-0376
  \HMTTitleMapEntry{Existencia, regularidad y unicidad global para Navier–Stokes mediante cota uniforme y compacidad}{Publicación solenoidal del continuo: cota coerciva, compacidad y regularidad global}% ADP102-0377
  \HMTTitleMapEntry{3.ª resolución: Yang–Mills y la brecha del sector de curvatura}{Publicación de calibre del continuo: positividad por reflexión y brecha del sector de curvatura}% ADP102-0378
  \HMTTitleMapEntry{4.ª resolución: Hodge y la construcción efectiva de la preimagen algebraica}{Publicación cohomológica del continuo: levantamiento compatible y construcción efectiva de la preimagen algebraica}% ADP102-0379
  \HMTTitleMapEntry{5.ª resolución: BSD y la coincidencia de las publicaciones analítica y aritmética}{Publicación determinantal del continuo: comparación Kato–Poitou–Tate e identidad del término principal de Birch–Swinnerton–Dyer}% ADP102-0380
  \HMTTitleMapEntry{La identidad de scattering que conserva la potencia cruzada}{Identidad de dispersión y conservación de la potencia cruzada}% ADP102-0381
  \HMTTitleMapEntry{Persistencia del objeto terminal bajo el paso al límite}{Persistencia del estado terminal solenoidal bajo el paso al límite}% ADP102-0382
  \HMTTitleMapEntry{Existencia, regularidad y unicidad global de la solución de Navier–Stokes obtenida por cota uniforme y compacidad}{Teorema de clausura solenoidal global mediante cota uniforme y compacidad}% ADP102-0383
  \HMTTitleMapEntry{El operador universal de extensión}{Operador universal de extensión de clases cohomológicas compatibles}% ADP102-0384
  \HMTTitleMapEntry{Factorización de los entrelazadores terminales RH–NS–YM–Hodge–BSD a través del continuo genealógico común}{Pérdidas de información específicas de las 5 publicaciones clásicas y conservación en el estado holonómico integral}% ADP102-0385
  \HMTTitleMapEntry{Factorización de RH–NS–YM–Hodge–BSD como publicaciones terminales del continuo genealógico común}{Teorema de factorización terminal de RH–NS–YM–Hodge–BSD a través del continuo genealógico común}% ADP102-0386
  \HMTTitleMapEntry{Emergencia correlativa de \texorpdfstring{\(5\)}{5} acciones intrínsecas no exhaustivas}{Descomposición expositiva de la acción correlativa en 5 vistas intrínsecas del mismo estado}% ADP102-0387
  \HMTTitleMapEntry{Una misma frontera: flujo interior y tensión gauge}{Factorización común del operador de frontera en las realizaciones solenoidal y de calibre}% ADP102-0388
  \HMTTitleMapEntry{El estado común de frontera}{Estado holonómico integral de frontera y descomposición paritaria del transporte}% ADP102-0389
  \HMTTitleMapEntry{El operador mínimo y su espectro}{Complejo de borde mínimo y equivalencia espectral de los laplacianos de grados 0 y 1}% ADP102-0390
  \HMTTitleMapEntry{Lectura hidrodinámica}{Realización solenoidal del operador de frontera: contracción del flujo y conservación de energía}% ADP102-0391
  \HMTTitleMapEntry{Lectura Yang--Mills}{Realización de calibre del operador de frontera: coercividad, cociente gauge y brecha espectral}% ADP102-0392
  % Promoción jerárquica de los certificados causales de \(\Gamma_9\) y de
  % la generación trítica de los canales enteros.
  \HMTTitleMapEntry{Filtración monodrómica y cálculo a profundidad prescrita}{Certificación finita de Γ₉: 396 niveles, 2.376 trits, 43 vueltas completas, 387 pares \((K,K+9)\) por canal y fibra ambiente de 729 elementos}% ADP102-0393
  \HMTTitleMapEntry{Censo finito de transiciones y generación decimal}{Censo exhaustivo de 396 transiciones, 2.376 trits, 43 vueltas y 387 pares de igual fase}% ADP102-0394
  \HMTTitleMapEntry{Estructura entera de los lectores de constantes}{Clasificación aritmética de los canales enteros y compatibilidad con la semicorona trítica}% ADP102-0395
}
\let\HMTNativeChapter\chapter
\let\HMTNativeSection\section
\let\HMTNativeSubsection\subsection
\let\HMTNativeSubsubsection\subsubsection
\let\HMTNativeParagraph\paragraph
\let\HMTNativeSubparagraph\subparagraph
\RenewDocumentCommand{\chapter}{s o m}{%
  \HMTResolvePublicTitle{#3}%
  \IfBooleanTF{#1}{\HMTNativeChapter*{\HMTResolvedTitle}}{%
    \IfNoValueTF{#2}{\HMTNativeChapter{\HMTResolvedTitle}}{\HMTNativeChapter[\HMTResolvedTitle]{\HMTResolvedTitle}}}}
\RenewDocumentCommand{\section}{s o m}{%
  \HMTResolvePublicTitle{#3}%
  \IfBooleanTF{#1}{\HMTNativeSection*{\HMTResolvedTitle}}{%
    \IfNoValueTF{#2}{\HMTNativeSection{\HMTResolvedTitle}}{\HMTNativeSection[\HMTResolvedTitle]{\HMTResolvedTitle}}}}
\RenewDocumentCommand{\subsection}{s o m}{%
  \HMTResolvePublicTitle{#3}%
  \IfBooleanTF{#1}{\HMTNativeSubsection*{\HMTResolvedTitle}}{%
    \IfNoValueTF{#2}{\HMTNativeSubsection{\HMTResolvedTitle}}{\HMTNativeSubsection[\HMTResolvedTitle]{\HMTResolvedTitle}}}}
\RenewDocumentCommand{\subsubsection}{s o m}{%
  \HMTResolvePublicTitle{#3}%
  \IfBooleanTF{#1}{\HMTNativeSubsubsection*{\HMTResolvedTitle}}{%
    \IfNoValueTF{#2}{\HMTNativeSubsubsection{\HMTResolvedTitle}}{\HMTNativeSubsubsection[\HMTResolvedTitle]{\HMTResolvedTitle}}}}
\RenewDocumentCommand{\paragraph}{s o m}{%
  \HMTResolvePublicTitle{#3}%
  \IfBooleanTF{#1}{\HMTNativeParagraph*{\HMTResolvedTitle}}{%
    \IfNoValueTF{#2}{\HMTNativeParagraph{\HMTResolvedTitle}}{\HMTNativeParagraph[\HMTResolvedTitle]{\HMTResolvedTitle}}}}
\RenewDocumentCommand{\subparagraph}{s o m}{%
  \HMTResolvePublicTitle{#3}%
  \IfBooleanTF{#1}{\HMTNativeSubparagraph*{\HMTResolvedTitle}}{%
    \IfNoValueTF{#2}{\HMTNativeSubparagraph{\HMTResolvedTitle}}{\HMTNativeSubparagraph[\HMTResolvedTitle]{\HMTResolvedTitle}}}}
% Las utilidades de degradación posteriores deben conservar el mapa de títulos.
\let\HMTBookChapter\chapter
\let\HMTBookSection\section
\let\HMTBookSubsection\subsection
\let\HMTBookSubsubsection\subsubsection
\let\HMTBookParagraph\paragraph
\let\HMTBookSubparagraph\subparagraph
\makeatother
 % Adaptador de compatibilidad del fuente teoremática íntegro del radión.
%
% El tratado rector ya carga tipografía, geometría, hipervínculos, teoremas,
% tablas, TikZ y tcolorbox.  Este archivo aporta únicamente los nombres y
% entornos exclusivos del fuente teoremática, sin importar su preámbulo autónomo ni
% alterar la jerarquía tipográfica general.

\makeatletter
\@ifundefined{color@RadNavy}{\definecolor{RadNavy}{HTML}{102A43}}{}
\@ifundefined{color@RadBlue}{\definecolor{RadBlue}{HTML}{1F5E7A}}{}
\@ifundefined{color@RadTeal}{\definecolor{RadTeal}{HTML}{168C8C}}{}
\@ifundefined{color@RadCopper}{\definecolor{RadCopper}{HTML}{B66A3C}}{}
\@ifundefined{color@RadGold}{\definecolor{RadGold}{HTML}{D4A73A}}{}
\@ifundefined{color@RadInk}{\definecolor{RadInk}{HTML}{1A202C}}{}
\@ifundefined{color@RadMist}{\definecolor{RadMist}{HTML}{EDF4F7}}{}
\@ifundefined{color@RadCream}{\definecolor{RadCream}{HTML}{FAF7F0}}{}
\@ifundefined{color@RadRose}{\definecolor{RadRose}{HTML}{8D3F55}}{}
\@ifundefined{color@RadGreen}{\definecolor{RadGreen}{HTML}{2F7D5D}}{}
\@ifundefined{color@RadGray}{\definecolor{RadGray}{HTML}{66788A}}{}

\providecommand{\TRIT}{\ensuremath{\mathrm{TRIT}}}
\providecommand{\HMT}{\ensuremath{\mathrm{HMT}}}
\providecommand{\Rstate}{\mathcal R_{12}}
\providecommand{\epsR}{\varepsilon_R}
\providecommand{\epsone}{\varepsilon_R^{(1)}}
\providecommand{\pih}{\pi}
\providecommand{\alphah}{\alpha}
\providecommand{\hret}{\hbar_{\mathrm{ret}}}
\providecommand{\hfull}{h_{\mathrm{ret}}}
\providecommand{\diff}{\mathop{}\!\mathrm d}
\providecommand{\e}{\mathrm e}
\providecommand{\R}{\mathbb R}
\providecommand{\C}{\mathbb C}
\providecommand{\F}{\mathbb F}
\providecommand{\Z}{\mathbb Z}
\providecommand{\dd}{\,\mathrm d}
\providecommand{\status}[1]{\textsf{\bfseries #1}}
\providecommand{\sourcepath}[1]{\nolinkurl{#1}}

\@ifundefined{typebox}{%
  \newtcolorbox{typebox}[1][]{enhanced,breakable,colback=white,
    colframe=RadBlue,boxrule=.7pt,arc=1.5mm,left=4mm,right=4mm,
    top=3mm,bottom=3mm,title={Recibo de tipos},fonttitle=\bfseries,
    coltitle=RadNavy,#1}%
}{}
\@ifundefined{narrativebox}{%
  \newtcolorbox{narrativebox}[1][]{enhanced,breakable,colback=white,
    colframe=RadGold,borderline west={2.4pt}{0pt}{RadGold},
    boxrule=.25pt,arc=0mm,left=5mm,right=4mm,top=3mm,bottom=3mm,#1}%
}{}

% Las once portadillas internas del fuente teoremática autónomo se convierten en
% divisiones científicas ordinarias dentro del capítulo 52.  Se evita así
% introducir partes autónomas, gigantismo tipográfico o páginas vacías.
\providecommand{\partopener}[3]{\section{#2}\noindent\emph{#3}\par\medskip}

\providecommand{\BigRadionEquation}{%
  \begin{equation*}
  \boxed{\displaystyle
  \frac{180^{\circ}}{\pih}
  =50^{\circ}+A^{\circ}-\frac{20}{81}\,\Delta_4^{\circ}+\epsR^{\circ}}
  \end{equation*}}
\makeatother
 % Adaptador mínimo para las correspondencias espirales genealógicas.
% El preámbulo autónomo no se importa: sólo se registran símbolos y colores
% que los cuerpos científicos usan de manera efectiva.
\makeatletter
\@ifundefined{color@HMTNavy}{\definecolor{HMTNavy}{HTML}{17324D}}{}
\@ifundefined{color@HMTBlue}{\definecolor{HMTBlue}{HTML}{315E7D}}{}
\@ifundefined{color@HMTRed}{\definecolor{HMTRed}{HTML}{8A3D36}}{}
\@ifundefined{color@HMTGold}{\definecolor{HMTGold}{HTML}{9B7B3E}}{}
\@ifundefined{color@HMTGray}{\definecolor{HMTGray}{HTML}{5F6870}}{}
\@ifundefined{color@HMTLight}{\definecolor{HMTLight}{HTML}{F2F4F5}}{}
\@ifundefined{color@HMTLine}{\definecolor{HMTLine}{HTML}{C9CFD3}}{}
\providecommand{\enr}{\mathrm{enr}}
\providecommand{\AffEnd}{\operatorname{AffEnd}}
\providecommand{\Hist}{\operatorname{Hist}}
\providecommand{\Led}{\operatorname{Led}}
\providecommand{\Carr}{\operatorname{Carr}}
\providecommand{\ev}{\operatorname{ev}}
\providecommand{\sd}{\operatorname{sd}}
\providecommand{\id}{\operatorname{id}}
\providecommand{\im}{\operatorname{im}}
\providecommand{\Spec}{\operatorname{Spec}}
\providecommand{\arcosh}{\operatorname{arcosh}}
\providecommand{\N}{\mathbb N}
\providecommand{\T}{\mathbb T}
\providecommand{\ii}{\mathrm i}
\providecommand{\twoway}{\textit{bidireccional}}
\providecommand{\ph}{\mathrm{ph}}
\providecommand{\mem}{\mathrm{mem}}
\providecommand{\rad}{\mathrm{rad}}
\providecommand{\arq}{\mathrm{arq}}
\providecommand{\e}{\mathrm e}
\providecommand{\F}{\mathbb F}
\providecommand{\Z}{\mathbb Z}
\providecommand{\R}{\mathbb R}
\providecommand{\dd}{\,\mathrm d}

% Los fuentes teoremáticas de correspondencias espirales se distribuyen por
% residencia causal. Sus etiquetas conservan un espacio de nombres único y
% sus divisiones autónomas se subordinan a la jerarquía del capítulo receptor.
\newenvironment{HMTSpiralIntegratedBody}{%
  \begingroup
  % Las copias públicas ya llevan el prefijo material ``espiral:'' en cada
  % etiqueta y referencia. Esto incluye los labels diferidos de amsmath, que
  % no podían asegurarse mediante una redefinición local de \label.
  \let\part\HMTBookSection
  \let\chapter\HMTBookSubsection
  \let\section\HMTBookSubsubsection
  \let\subsection\HMTBookParagraph
  \let\subsubsection\HMTBookSubparagraph
}{\endgroup}

\newenvironment{HMTSpiralChapterBody}{%
  \begingroup
  \let\chapter\HMTBookSection
  \let\section\HMTBookSubsection
  \let\subsection\HMTBookSubsubsection
  \let\subsubsection\HMTBookParagraph
}{\endgroup}

\newenvironment{HMTSpiralAppendixBody}{%
  \begingroup
  \let\part\HMTBookChapter
  \let\chapter\HMTBookSection
  \let\section\HMTBookSubsection
  \let\subsection\HMTBookSubsubsection
  \let\subsubsection\HMTBookParagraph
}{\endgroup}
\makeatother
 
% Identificador estable para builds reproducibles.
\ifdefined\pdfinfoomitdate
  \pdfinfoomitdate=1
\fi
\ifdefined\pdfsuppressptexinfo
  \pdfsuppressptexinfo=15
\fi
\ifdefined\pdftrailerid
  \pdftrailerid{<554E49444144484F4C4F47524146494341>
                <44494D454E53494F4E414C484D54323026>}
\fi

\makeatletter\begin{document}
\frontmatter
\hypersetup{pageanchor=false}
\begin{titlepage}
\thispagestyle{empty}
\centering
\vspace*{1.8cm}
{\color{hmtgold}\rule{0.68\textwidth}{0.8pt}\par}
\vspace{1.35cm}
{\fontsize{31}{37}\selectfont\color{hmtblue}\bfseries
Doble proyección holográfica\par
de un cristal temporal aperiódico\par}
\vspace{1.0cm}
{\Large\bfseries
Construcción del continuo como límite inverso de cilindros compatibles,\par
generación coinductiva de \(\pi,\varphi,e,\alpha\) y estabilizadores\par
excepcionales de la incidencia fronteriza\par}
\vspace{1.1cm}
{\color{hmtgold}\rule{0.68\textwidth}{0.8pt}\par}
\vfill
{\large Oumar Haidara Fall\par}
{\large Rubén Ramos Balsa\par}
\vspace*{1.8cm}
\end{titlepage}
\hypersetup{pageanchor=true}

\chapter*{Resumen}
\addcontentsline{toc}{chapter}{Resumen}

Se construye una aritmética de historias orientadas en la que la unidad se
conserva bajo refinamiento mientras aumentan la resolución, la altura
holonómica y la memoria. El soporte finito de todos los caminos posibles
mantiene simultáneamente las evaluaciones aditiva y multiplicativa, junto con
sus residuos, cocientes y acarreos; una fibra trítica distingue los regímenes
elíptico, parabólico e hiperbólico; y un núcleo de fase topológica selecciona,
transporta y actualiza cada estado sin confundir el retorno de fase con el
retorno del estado. La composición de estas tres estructuras determina un
álgebra asociativa no conmutativa, graduada por memoria, y un sistema de
refinamientos compatibles cuyo espacio de historias es un límite inverso.
Los observables cilíndricos se prolongan en sentido covariante y constituyen
el límite directo dual. Esta doble varianza proporciona la definición
operatoria de la holografía empleada en el artículo.

La dinámica combina una fase nonádica finita, una codificación mecánica de
pendiente $\delta=\log_{10}(10/9)$ y un cociclo de acarreo. El resultado es
un orden topológico-temporal aperiódico de complejidad lineal y entropía
topológica nula, provisto de una holonomía que eleva la memoria en cada retorno
nonádico. Sus relaciones de Weyl separan rigurosamente la periodicidad local
de la aperiodicidad global. Sobre este soporte, la prolongación
coinductiva genera correlativamente $\pi$, $\varphi$, $e$ y $\alpha$: las
tres primeras coordenadas realizan, respectivamente, la clausura elíptica,
la autoescala hiperbólica y el transporte exponencial; la constante de
estructura fina es el invariante de compatibilidad que reúne, por dos
derivaciones internas, el cierre dodecafásico y la prolongación aperiódica.

La misma estructura produce una vacancia unitaria de acarreo multiplicativo
en una fibra aditiva invariante. Su álgebra local es
$\mathrm{Cl}_{1,1}\cong M_2(\R)$ y contiene simultáneamente generadores
elípticos, hiperbólicos y direcciones nulas de Witt. Esta sección determina
el núcleo electrónico y prolonga la genealogía hacia la acción, las escalas
de Planck y gravedad, el parámetro de Barbero--Immirzi, el espectro de área y
la información de frontera. La segunda proyección conserva la incidencia y
conduce, mediante las configuraciones de Paley y Witt, el código de Golay, el
retículo de Leech y el álgebra de vértices $V^\natural$, a los estabilizadores
excepcionales y al grupo Monstruo. Así, la realización arquimediana y la
realización excepcional no constituyen fuentes independientes: son dos
representaciones de una única estructura discreta del continuo.

El lector dodecafásico uniforme determina además una sucesión de primeras raíces de retorno crítico. Se demuestra que su cociente de separaciones converge a la constante de Feigenbaum, conservando la selección original, la orientación y la memoria de la construcción.

\vspace{4mm}
\noindent\textbf{Palabras clave:} holografía modular triádica; límite
inverso; holonomía nonádica; palabra esturmiana; cristal temporal aperiódico;
álgebra de caminos; constante de estructura fina; Clifford escindido;
Barbero--Immirzi; retículo de Leech; grupo Monstruo.

\clearpage

\chapter*{Abstract}
\addcontentsline{toc}{chapter}{Abstract}

We construct an arithmetic of oriented histories in which the unit is
preserved under refinement while resolution, holonomic height and memory
increase.  The finite support of all possible paths retains the additive and
multiplicative evaluations simultaneously, together with their remainders,
quotients and carries; a tritic fibre distinguishes the elliptic, parabolic
and hyperbolic regimes; and a topological phase kernel selects, transports
and updates each state without identifying phase return with state return.
Their composition determines a memory-graded, associative and
noncommutative algebra and a compatible refinement system whose history
space is an inverse limit.  Cylindrical observables extend covariantly and
form the dual direct limit.  This variance duality is the operative notion of
holography used throughout the article.

The dynamics combines a finite nonadic phase, a mechanical coding of slope
$\delta=\log_{10}(10/9)$ and a carry cocycle.  It yields an aperiodic
topological--temporal order with linear complexity and zero topological
entropy, endowed with a holonomy that raises memory at every nonadic return.
On this support the coinductive prolongation generates $\pi$, $\varphi$,
$e$ and $\alpha$ as correlated outputs.  The same structure produces a
unit multiplicative-carry vacancy on an invariant additive fibre, whose
local algebra is $\mathrm{Cl}_{1,1}\cong M_2(\R)$.  The Archimedean
realization and the exceptional boundary-incidence realization are thereby
shown to be two representations of one discrete structure of the continuum.

The uniform dodecaphasic reader also determines a sequence of first roots of critical return. Its separation ratio is proved to converge to the Feigenbaum constant, preserving the original selection, orientation and memory of the construction.

\vspace{4mm}
\noindent\textbf{Keywords:} triadic modular holography; inverse limit;
nonadic holonomy; Sturmian word; aperiodic time crystal; path algebra;
fine-structure constant; split Clifford algebra; Barbero--Immirzi parameter;
Leech lattice; Monster group.
\clearpage
\tableofcontents
\clearpage

\let\cleardoublepage\clearpage
\mainmatter
\part{Aritmética bifoliada, orientación trigeométrica y dinámica de fase}

\chapter{La unidad como sección compatible de una aritmética de historias}
\begingroup
\renewcommand{\chapter}[2][]{}%
\chapter{Tesis y arquitectura matemática}
\label{sec:tesis-inaugural-helicoidal-hmt-md}

El lector dodecafásico uniforme determina además una sucesión de primeras raíces de retorno crítico. Se demuestra que su cociente de separaciones converge a la constante de Feigenbaum, conservando la selección original, la orientación y la memoria de la construcción.


\section{La unidad como invariante de refinamiento}

La cuestión central no consiste en obtener una colección de valores
escalares a partir de una tabla finita. Consiste en determinar la estructura
mínima que conserva una unidad a través de refinamientos ilimitados,
distingue tres clases de curvatura, transporta la memoria de sus recorridos y
genera, mediante representaciones tipadas, números, constantes, secciones,
campos y simetrías. La recta real no se presupone como dominio universal de
esta construcción. Aparece al término de una contracción compatible de
cilindros; conserva el valor escalar de la historia, pero no necesariamente
la totalidad de su genealogía.

Para cada profundidad $n\geq0$, sea $X_n$ el conjunto finito de estados
admisibles. Un estado no es una cifra aislada, sino una fibra estructurada

\begin{equation}
x=\bigl(a,b;\Sigma,\Pi;R^+,q^+,R^\times,q^\times;
\kappa;\tau,o;\theta_9,\theta_{12};h;\partial;\mu;\rho;s\bigr),
\label{eq:estado-integral}
\end{equation}

donde $(a,b)$ localiza la celda; $\Sigma$ y $\Pi$ son las hojas aditiva y
multiplicativa; $(R^+,q^+)$ y $(R^\times,q^\times)$ conservan residuo y
cociente; $\kappa$ es el acarreo; $\tau\in\{+1,0,-1\}$ tipa la curvatura;
$o$ fija la orientación; $\theta_9$ y $\theta_{12}$ son las fases nonádica y
dodecafásica; $h$ mide la altura helicoidal; $\partial$ registra la incidencia
fronteriza; $\mu$ es la memoria; $\rho$ conserva la ruta y $s$ la condición
de supervivencia. Estas coordenadas no forman un producto libre: satisfacen
las relaciones de la aritmética de caminos, el tipado trigeométrico y la
dinámica de transporte que se construirán en los capítulos siguientes.

Los mapas

\begin{equation}
r_{m,n}:X_m\longrightarrow X_n,\qquad m>n,
\label{eq:restriccion}
\end{equation}

olvidan la última capa de resolución, pero preservan toda relación visible a
profundidad $n$. El espacio de historias infinitamente refinables es

\begin{equation}
X_\infty=\varprojlim_n X_n.
\label{eq:limite-inverso}
\end{equation}

Si $e_x$ denota el idempotente característico del cilindro determinado por
$x\in X_n$, la prolongación de observables satisface

\begin{equation}
r_{m,n}^{*}e_x=\sum_{y:\,r_{m,n}(y)=x}e_y.
\label{eq:elevacion-idempotente}
\end{equation}

Por consiguiente,

\begin{equation}
r_{m,n}^{*}\!\left(\sum_{x\in X_n}e_x\right)
=\sum_{y\in X_m}e_y,
\label{eq:conservacion-unidad}
\end{equation}

de modo que la unidad se conserva aunque crezca el número de extensiones. La
evolución no conserva una configuración inmóvil: conserva el aumento, las
relaciones y la posibilidad de reconstruir cada antecedente desde la historia
completa.

\section{Doble varianza y doble proyección}

La profundidad de los estados y la extensión de los observables avanzan en
sentidos categóricos opuestos. A las restricciones de
\eqref{eq:restriccion} corresponden morfismos de prolongación

\begin{equation}
r_{m,n}^{*}:\Halg_n\longrightarrow\Halg_m,
\qquad f\longmapsto f\circ r_{m,n},
\end{equation}

y el álgebra de observables cilíndricos se obtiene como límite directo:

\begin{equation}
\Halg_{\mathrm{cyl}}=\varinjlim_n(\Halg_n,r_{n+1,n}^{*}).
\label{eq:limite-directo}
\end{equation}

La pareja formada por \eqref{eq:limite-inverso} y
\eqref{eq:limite-directo} proporciona la estructura holográfica fundamental:
el estado adquiere resolución por compatibilidad contravariante y el lenguaje
de observación adquiere alcance por prolongación covariante. No se trata de
una semejanza figurativa entre escalas, sino de la naturalidad de la
evaluación

\begin{equation}
\langle r_{m,n}^{*}a,y\rangle_m
=\langle a,r_{m,n}y\rangle_n.
\label{eq:naturalidad-evaluacion}
\end{equation}

La doble proyección del título designa dos representaciones de esta misma
pareja. La primera contrae cilindros compatibles hacia una realización
arquimediana; la segunda retiene la incidencia de frontera y la prolonga
hacia estructuras excepcionales:

\begin{equation}
\rho_{\mathrm{arq}}:\Halg_{\mathrm{cyl}}\longrightarrow\cA_{\R},
\qquad
\rho_{\mathrm{exc}}:\Halg_{\mathrm{cyl}}\longrightarrow\cA_{\mathrm{exc}}.
\label{eq:doble-proyeccion}
\end{equation}

Ambas poseen la misma fuente y conservan invariantes diferentes. La primera
retiene la coordenada escalar de una sección compatible; la segunda retiene
los patrones de incidencia que determinan sus estabilizadores. La palabra
\emph{doble} no duplica la holografía: distingue dos codominios no
equivalentes de un único mecanismo de refinamiento y lectura.

\section{Enunciado central}

\begin{theorem}[Evolución conservativa y producción de simetrías]
Existe un álgebra graduada de historias orientadas
\begin{equation}
\Halg=\left[
\varinjlim_n
\frac{\R\langle e_x,J_x,T_u\rangle}
{\cR_{\mathrm{APP\text{-}TRIT\text{-}TPK}}}
\right]\rtimes_{\Gamma_9}\N,
\label{eq:algebra-rectora}
\end{equation}
que actúa sobre $X_\infty$ y satisface las propiedades siguientes:
\begin{enumerate}[label=\textup{(\roman*)},leftmargin=9mm]
\item la unidad y el aumento se conservan bajo todo refinamiento compatible;
\item el retorno nonádico cierra la fase y eleva la memoria;
\item los regímenes elíptico, parabólico e hiperbólico pertenecen a un único
cálculo funcional trigeométrico;
\item el factor simbólico posee complejidad lineal, entropía topológica nula
y orden aperiódico de largo alcance;
\item la estructura discreta del continuo surge conjuntamente del espacio de
historias y del álgebra de observables;
\item $\pi$, $\varphi$, $e$ y $\alpha$ son salidas correlacionadas de la
misma prolongación coinductiva;
\item una simetría de una realización es el estabilizador de las relaciones y
de la incidencia fronteriza que dicha realización conserva.
\end{enumerate}
\end{theorem}

La demostración se distribuirá a lo largo del artículo. No se utilizará una
constante conocida para seleccionar su propia genealogía, ni un grupo de
simetría para construir retrospectivamente el objeto sobre el que actúa. El
orden lógico será siempre el mismo: aritmética bifoliada, tipado trigeométrico,
transporte holonómico, refinamiento compatible, continuo, constantes,
secciones físicas y estabilizadores de frontera.

\section{Irreducibilidad del núcleo generador}

La estructura posee tres componentes funcionalmente irreducibles. Si se
identifican las dos hojas aritméticas, desaparecen el potencial
$\Delta_{\Sigma\Pi}$ y el conmutador de sus proyectores. Si se identifican
los tres regímenes tríticos, desaparece el espectro

\begin{equation}
J_+^2=-I,\qquad J_0^2=0,\qquad J_-^2=I.
\end{equation}

Si se trivializa el cociclo de transporte, el incremento

\begin{equation}
[D_M,\Gamma_9]=\Gamma_9
\label{eq:conmutador-memoria-intro}
\end{equation}

se anula y la vuelta de fase se vuelve indistinguible del retorno del estado.
Cada cociente de olvido destruye un invariante de tipo diferente que los otros
dos factores no reconstruyen. Esta irreducibilidad explica por qué la
estructura buscada no puede reducirse a una tabla aritmética, a una palabra
aperiódica, a un álgebra de Clifford o a una matriz aislada.

La consecuencia conceptual es precisa. El número es una coordenada escalar
de una representación; la constante es un invariante estable de retorno o de
autoescala; la partícula es una sección persistente; el campo transporta
secciones; la simetría estabiliza las relaciones de su frontera. Estas cinco
clases proceden del mismo álgebra, pero no comparten codominio ni pueden
confundirse entre sí.
 \endgroup

\chapter{Soporte aritmético de todos los caminos posibles}
\label{chap:83-01}

% RESULTADO_RECUPERADO. Owner integro: sections/hmt/03_app_ortograma.tex
% Destino causal: capitulo 1; la jerarquia se subordina sin alterar el owner.
% Los dos bloques científicos del owner se activan en una única residencia.
\begingroup
\def\HMTAPPFormal{1}
\def\HMTAPPDevelopments{1}
\begin{HMTScientificOwnerSubordinated}
\ifdefined\HMTAPPFormal
\chapter{APP: soporte, operaciones y caminos}
\label{chap:app}
\index{APP}

\noindent
La sigla APP procede de \emph{All Possible Paths}. Designa una única
estructura de caminos sobre un soporte aritmético finito: el grafo
subyacente tiene ochenta y un vértices, mientras que el conjunto de caminos
finitos, reunido sobre todas las longitudes, es numerable. Su matriz
aritmética principal es la tabla multiplicativa nonádica. La suma es la ley
interna cuya acumulación genera cada fila y cada columna de esa tabla y actúa
además como observable de comparación sobre el mismo soporte. APP reúne así
suma, producto y
composición de trayectorias antes de efectuar la proyección ternaria o la
evaluación arquimediana. La definición siguiente fija de manera exhaustiva
el universo de caminos considerado.

\section{Dominio, generadores y relaciones preformales del soporte APP}

El soporte de APP es un toro aritmético finito: cada una de sus dos
coordenadas recorre las nueve clases residuales y vuelve a su origen después
de nueve traslaciones unitarias. Sus ochenta y un puntos son, por tanto,
pares de clases módulo nueve. El grafo de Cayley asociado conserva esta
estructura cíclica: un vértice representa una posición, una arista orientada
representa un desplazamiento cardinal unitario y una sucesión finita de
aristas representa un camino. De este modo, \emph{todos los caminos
posibles} designa el conjunto de las palabras cardinales finitas con punto
inicial en el toro.

Cada vértice recibe simultáneamente dos evaluaciones aritméticas. La
\emph{hoja aditiva} publica la suma de sus dos coordenadas y la
\emph{hoja multiplicativa} publica su producto. El término \emph{hoja}
designa aquí una evaluación del mismo punto y del mismo camino; APP conserva
un único soporte y dos funciones sobre él.

La diferencia operativa entre ambas hojas es exacta. En una fila o columna,
la suma actúa por traslación y acumula reiteradamente el mismo incremento.
La multiplicación por una unidad también permuta las clases; la
multiplicación por un divisor de cero pliega varias clases sobre un mismo
residuo. Para conservar la información que ese plegado oculta, APP eleva
cada evaluación a sus coordenadas residuo--cociente:
\[
 i+j=R^+_{ij}+9q^+_{ij},
 \qquad
 ij=R^\times_{ij}+9q^\times_{ij}.
\]
Los pares \((R^+,q^+)\) y \((R^\times,q^\times)\) recuperan los valores
enteros anteriores a la reducción. Al concatenar bloques, el acarreo aditivo
\(\kappa_9\) registra la discrepancia entre sumar representantes enteros y
sumar primero sus residuos; junto con ambos cocientes, hace composicional la
recuperación. Así, la acumulación aditiva construye las trayectorias visibles,
mientras el levantamiento residuo--cociente despliega las fibras creadas por el
plegado multiplicativo.

Para un camino \(\gamma=(x_0,\ldots,x_r)\), la genealogía visible es la
sucesión de posiciones y pares de evaluaciones
\[
 \mathcal G_{\rm vis}(\gamma)
 =\bigl(\gamma;(\Sigma(x_k),\Pi(x_k))_{k=0}^{r}\bigr).
\]
Su genealogía enriquecida añade, en cada paso, los cocientes, el acarreo y
la orientación transportada. La proyección enriquecida--visible conserva el
camino y los residuos y olvida precisamente esas coordenadas de memoria. Esta
distinción será la entrada del estado completo del TRIT y del Topological Phase Kernel.
El dominio completo de (104\,976) condiciones iniciales de doble cursor y
las (468) fibras de emisión que producen se publican íntegramente en el
\cref{app:catalogo-semillas-app-rev9}; allí se distingue cada semilla de la
emisión a la que pertenece y se conserva la multiplicidad de la fibra.

\section{Soporte aritmético, traslaciones y observables}

Sea
\[
\dr(n)=1+((n-1)\bmod9)\in\{1,\ldots,9\}.
\]
Esta aplicación escoge el representante \(9\) para la clase nula. Definimos
\(X_9=(\ZZ/9\ZZ)^2\), sobre el que actúa el grupo de traslaciones
\((\ZZ/9\ZZ)^2\), y fijamos el conjunto simétrico de generadores
\[
 \mathscr D=\{(1,0),(0,1),(-1,0),(0,-1)\}.
\]
El grafo de Cayley
\[
 \GraphAPP=\operatorname{Cay}(X_9,\mathscr D)
\]
posee \(81\) vértices y cuatro aristas orientadas con origen en cada vértice.
Un camino orientado de longitud \(r\) es una sucesión
\(\gamma=(x_0,\ldots,x_r)\) en \(X_9\) tal que
\(x_{k+1}-x_k\in\mathscr D\) para todo \(k<r\). Denotamos por
\(\operatorname{Path}(\GraphAPP)\) el conjunto de todos los caminos finitos
así definidos.

\begin{definition}[Estructura APP]
La estructura aritmética de caminos posibles es la tupla
\[
 \AAPP=(X_9,\GraphAPP,\Sigma,\Pi,
        \operatorname{Path}(\GraphAPP)),
\]
donde los dos observables de vértice son
\[
\Sigma(i,j)=\dr(i+j),
\qquad
\Pi(i,j)=\dr(ij),
\qquad1\le i,j\le9.
\]
\end{definition}

\subsection*{Marco cardinal y palabras de trayectoria}

Para hacer explícita la orientación del toro, usamos coordenadas
\((\text{fila},\text{columna})\), con la fila creciente hacia el sur y la
columna creciente hacia el este, y escribimos
\[
 \delta(N)=(-1,0),\qquad \delta(E)=(0,1),\qquad
 \delta(S)=(1,0),\qquad \delta(O)=(0,-1).
\]
Sea \(\mathscr D_{\rm card}=\{E,N,O,S\}\), donde \(O\) denota el oeste. Dado un punto inicial
\(x_0\in X_9\) y una palabra \(w=d_1\cdots d_r\in
\mathscr D_{\rm card}^{r}\), su trayectoria es
\[
 \gamma(x_0;w)_k
 =x_0+\sum_{j=1}^{k}\delta(d_j)\pmod 9,
 \qquad 0\leq k\leq r.
\]
La correspondencia \((x_0,w)\mapsto\gamma(x_0;w)\) es biyectiva entre
\(X_9\times\mathscr D_{\rm card}^{r}\) y los caminos orientados de
longitud \(r\). En particular, existen exactamente \(81\cdot4^r\) caminos
de esa longitud cuando se permiten retornos y retrocesos.

La palabra \(w\) cierra en el toro precisamente cuando
\[
\#S(w)-\#N(w)\equiv0\pmod9,
\qquad
\#E(w)-\#O(w)\equiv0\pmod9.
\]
Para una palabra cerrada, el desplazamiento entero calculado antes de reducir
módulo nueve define el vector de enrollamiento
\[
 \operatorname{wind}(w)=\frac1{9}
 \bigl(\#S(w)-\#N(w),\#E(w)-\#O(w)\bigr)\in\ZZ^2.
\]
La inversión de la orientación revierte el orden de la palabra, intercambia
\(E\leftrightarrow O\) y \(N\leftrightarrow S\), y transforma
\(\operatorname{wind}(w)\) en su opuesto. Este marco cardinal es el que se
transporta posteriormente en el núcleo de fase y forma parte de la definición
de cada trayectoria.

Las expresiones se escriben aquí en la carta de representantes
\(\{1,\ldots,9\}\). La fórmula
\(\widetilde\Sigma(i,j)=\dr(i+j-1)\) es una carta trasladada. Representa el
mismo toro sólo cuando la traslación se aplica también a índices, condiciones
iniciales, trayectorias posicionales y etiquetas. Todas las demostraciones siguientes emplean
\(\Sigma(i,j)=\dr(i+j)\) y no alternan ambas cartas.

\begin{figure}[H]
\centering
\resizebox{\textwidth}{!}{\begin{tikzpicture}[font=\scriptsize,cell/.style={minimum width=4.2mm,minimum height=4.2mm,inner sep=0pt,draw=hmtgray!35}]
  \begin{scope}[xshift=0cm]
    \node[font=\bfseries\scriptsize,hmtblue,align=center] at (1.72,.92)
      {matriz aritmética APP\\[-.5mm]\(\Pi(i,j)=\dr(ij)\)};
    \foreach \i in {1,...,9}{
      \node[font=\bfseries] at (-.35,{-(\i-1)*.43}) {\i};
      \foreach \j in {1,...,9}{
        \pgfmathtruncatemacro{\v}{mod(\i*\j-1,9)+1}
        \pgfmathtruncatemacro{\border}{(\i==9)||(\j==9)}
        \pgfmathtruncatemacro{\rad}{(mod(\i,3)==0)||(mod(\j,3)==0)}
        \ifnum\border=1
          \node[cell,fill=hmtlightblue,text=hmtblue,font=\bfseries] at ({(\j-1)*.43},{-(\i-1)*.43}) {\v};
        \else
          \ifnum\rad=1
            \node[cell,fill=hmtlightviolet] at ({(\j-1)*.43},{-(\i-1)*.43}) {\v};
          \else
            \node[cell,fill=white] at ({(\j-1)*.43},{-(\i-1)*.43}) {\v};
          \fi
        \fi
      }
    }
    \foreach \j in {1,...,9}{\node[font=\bfseries] at ({(\j-1)*.43},.38) {\j};}
    \draw[hmtgold,very thick,rounded corners=1pt] (1.075,-1.075) rectangle (1.935,-1.935);
    \draw[hmtblue,very thick] (3.72,.22)--(3.72,-3.72)--(-.22,-3.72);
    \node[hmtblue,anchor=north] at (1.75,-4.02) {subconjunto coordenado de la clase nula};
  \end{scope}

  \begin{scope}[xshift=5.25cm]
    \node[draw=hmtblue,rounded corners=2pt,fill=hmtlightblue,
      text width=6.25cm,align=left,inner sep=8pt] at (2.9,-1.35) {%
      \textbf{Ley interna de acumulación}\par\smallskip
      \[
        \Pi(i,j)=\dr(\underbrace{i+\cdots+i}_{j\ {\rm sumandos}})
        =\dr(ij).
      \]
      La suma genera las filas multiplicativas y, en la dirección
      transversal, las columnas. El observable
      \(\Sigma(i,j)=\dr(i+j)\) conserva esa ley elemental para comparar
      ambas lecturas sobre el mismo punto.};
    \node[draw=hmtgold,rounded corners=2pt,fill=hmtlightgold,
      text width=6.25cm,align=center,inner sep=8pt] at (2.9,-3.65) {%
      \(\displaystyle
      \AAPP=(X_9,\GraphAPP,\Pi;\Sigma,
      \operatorname{Path}(\GraphAPP))\)\par
      una sola APP, una matriz principal, una ley aditiva interna y cuatro
      direcciones orientadas \(N,E,S,O\).};
  \end{scope}
\end{tikzpicture}
 }
\caption{La única matriz APP y su ley interna. La tabla multiplicativa
nonádica ocupa el cuerpo principal; cada una de sus filas se genera por
acumulación aditiva. El sombreado violeta señala las filas o columnas no
unitarias; la línea azul identifica el subconjunto determinado por el
representante \(9\) de la clase nula; el marco dorado destaca el bloque
central \(2\times2\).}
\label{fig:app-ortograma}
\end{figure}

Cada punto de \(X_9\) recibe simultáneamente los valores de producto y suma,
pero no pertenece por ello a dos matrices APP distintas. Los caminos de
\(\operatorname{Path}(\GraphAPP)\) recorren un único soporte; \(\Pi\)
publica su tabla principal y \(\Sigma\) conserva la operación interna que
genera las acumulaciones y permite medir su incompatibilidad.

\begin{proposition}[Estructura de las filas y columnas]
\label{prop:app-latina-unidades}
La matriz de \(\Sigma\) es un cuadrado latino de orden nueve. La matriz de
\(\Pi\) no lo es; sin embargo, para cada
\(u\in(\ZZ/9\ZZ)^\times=\{1,2,4,5,7,8\}\), la fila y la columna
correspondientes a \(u\) son permutaciones de las nueve clases residuales.
\end{proposition}

\begin{proof}
Fijado \(i\), la aplicación \(j\mapsto i+j\) es una traslación de
\(\ZZ/9\ZZ\); lo mismo vale al fijar \(j\). Por ello cada representante
aparece exactamente una vez en cada fila y columna de \(\Sigma\). Para
\(\Pi\), la multiplicación por \(u\) permuta las clases si y sólo si \(u\)
es una unidad. Las filas correspondientes a \(3,6,9\) poseen repeticiones,
por lo que la propiedad latina no se extiende a toda la matriz
multiplicativa.
\end{proof}

\section{Estructura multiplicativa del anillo local}

Cada fila de la tabla multiplicativa es una suma repetida. La fila \(i\)
acumula \(i\) a lo largo de una dirección; la columna \(j\) acumula \(j\) a
lo largo de la dirección perpendicular. La celda \((i,j)\) es la intersección
de ambas acumulaciones y devuelve \(ij=ji\). La tabla representa así producto
y conmutatividad mediante dos acumulaciones transversales sobre la misma
retícula.

La segunda fila es
\[
2,4,6,8,1,3,5,7,9.
\]
Antes del retorno modular aparecen los pares y después los impares. La tercera
fila es
\[
3,6,9,3,6,9,3,6,9,
\]
y hace visible el radical. Sea
\[
 R:=\ZZ/9\ZZ,
 \qquad
 \mathfrak m:=(3)=\{\overline0,\overline3,\overline6\}.
\]
La multiplicación por tres define
\[
 M_3:R\longrightarrow R,
 \qquad x\longmapsto3x,
 \qquad
 \ker M_3=\operatorname{im}M_3=\mathfrak m.
\]
Por el primer teorema de isomorfía induce una identificación lineal
\begin{equation}
 \overline M_3:R/\mathfrak m\xrightarrow{\ \sim\ }\mathfrak m,
 \qquad \overline x\longmapsto3x,
 \label{eq:multiplicador-tres-radical-rev3}
\end{equation}
de espacios vectoriales sobre \(\mathbb F_3\). No es un isomorfismo de
anillos: \(\mathfrak m^2=0\), mientras el cociente es un cuerpo.

En representantes positivos, las tres fases de
\eqref{eq:multiplicador-tres-radical-rev3} son
\begin{equation}
 369\longmapsto693\longmapsto936\longmapsto369,
 \label{eq:orbita-radical-369-rev3}
\end{equation}
producidas por el desplazamiento de fase \(j\mapsto j+1\). La fila
correspondiente a \(6\equiv-3\pmod9\) invierte la orientación del mismo
ciclo. Esta órbita radical no debe confundirse con el entero
\(369\,693\,100\), que aparecerá después como número de cifras de
\(9^{9^9}\).

En reducción directa módulo tres,
\[
3,6,9\longmapsto0,0,0.
\]
Si primero extraemos el factor tres y guardamos su coordenada interna,
\[
\eta_{\mathfrak m}(3a)=a\bmod3,
\qquad
3,6,9\longmapsto1,2,0.
\]
La segunda aplicación conserva la coordenada interna del ideal generado por
tres. No debe confundirse con la proyección ternaria orientada que se definirá
en la sección dedicada al TRIT\@.

\section{Subconjunto coordenado de la clase nula}

En la gráfica del observable multiplicativo,
\[
\Pi(9,j)=\Pi(i,9)=9.
\]
La última fila y la última columna forman una unión de diecisiete posiciones,
todas con valor nueve. No constituyen una frontera intrínseca del toro
aritmético: son el subconjunto de la carta en el que al menos una coordenada
representa la clase nula mediante \(9\). Existen además cuatro representantes nulos interiores en
\[
(3,3),(3,6),(6,3),(6,6).
\]
La figura distingue así este subconjunto y los ceros internos del ideal
no unitario. El símbolo adicional empleado más adelante en la completación
cúbica pertenece a un espacio ampliado y no se identifica con ninguno de
ellos.

\section{Extensión residuo--cociente}

La reducción módulo nueve conserva el residuo y elimina el cociente de la
división euclídea. Para estudiar ambas componentes definimos, en cada punto,
\[
i+j=9q^+_{ij}+R^+_{ij},
\qquad
ij=9q^\times_{ij}+R^\times_{ij},
\]
donde \(R^+=\Sigma\), \(R^\times=\Pi\) y los \(q\) conservan el cociente.
La extensión residuo--cociente de APP se denota por
\[
\AAPPlift=(R^+,q^+;R^\times,q^\times).
\]

Sea \(s_9:\ZZ/9\ZZ\to\{1,\ldots,9\}\) la sección de representantes fijada
por \(\dr\). El acarreo aditivo asociado a esta sección es
\[
 \kappa_9(a,b)
 =\frac{s_9(a)+s_9(b)-s_9(a+b)}9\in\ZZ.
\]

\begin{lemma}[Identidad cocíclica del acarreo]
\label{lem:app-cociclo-acarreo}
Para cualesquiera \(a,b,c\in\ZZ/9\ZZ\),
\[
 \kappa_9(a,b)+\kappa_9(a+b,c)
 =\kappa_9(b,c)+\kappa_9(a,b+c).
\]
\end{lemma}

\begin{proof}
Ambos miembros son iguales a
\[
 \frac{s_9(a)+s_9(b)+s_9(c)-s_9(a+b+c)}9.\qedhere
\]
\end{proof}

Este cociclo registra la discrepancia entre sumar representantes enteros y
sumar primero en el cociente. Es la forma composicional de la coordenada de
cociente que reaparecerá en la concatenación base \(1000\).

Los totales brutos son
\[
\sum_{i,j}(i+j)
=9\sum_i i+9\sum_j j
=2\cdot9\cdot45
=810,
\]
y
\[
\sum_{i,j}ij
=\left(\sum_i i\right)\left(\sum_j j\right)
=45^2
=2025.
\]
Las sumas visibles se obtienen de las tablas:
\[
\sum R^+=405,
\qquad
\sum R^\times=459.
\]
Por tanto, las sumas totales de las coordenadas de cociente son
\[
\sum q^+=\frac{810-405}{9}=45,
\qquad
\sum q^\times=\frac{2025-459}{9}=174.
\]

\begin{exactbox}[title={Descomposición exacta de la diferencia integrada}]
\[
\boxed{
2025-810
=1215
=(459-405)+9(174-45)
=54+9\cdot129.}
\]
La diferencia entre las sumas de residuos es \(54\), mientras que la
diferencia entre las sumas de cocientes es \(129\). La igualdad recompone la
diferencia total entre los observables enteros y demuestra que la componente
de cociente contiene información aritmética no determinada por la clase
residual.
\end{exactbox}

\fi

\ifdefined\HMTAPPDevelopments
\chapter{Operadores aritméticos y geometría local del espacio de caminos}
\label{chap:app-operadores}

\section{Localización del exceso en el radical nilpotente}

El grupo de unidades de \(\ZZ/9\ZZ\) es
\[
(\ZZ/9\ZZ)^\times=\{1,2,4,5,7,8\},
\]
y el ideal máximo es
\[
\mathfrak m=(3)=\{0,3,6\}=\{9,3,6\}_{\rm visual},
\qquad\mathfrak m^2=0.
\]
Cada fila unitaria de \(\Pi\) es una permutación de \(1,\ldots,9\) y suma
\(45\). Las filas \(3\) y \(6\) suman \(54\); la fila \(9\), \(81\).
Luego
\[
459=6\cdot45+2\cdot54+81
\]
y
\[
459-405=2(54-45)+(81-45)=54.
\]
Todo el exceso se localiza en el sector no unitario.

La unión de filas y columnas \(3,6,9\) contiene \(45\) celdas y suma
\(297\). Su cuadrado central \(3\times3\) suma \(81=3\cdot27\); los brazos
exteriores, \(216=3\cdot72\). Así,
\[
297=216+81=3(72+27)=3\cdot99.
\]
La interpretación de \(72\) y \(27\) se obtiene más adelante a partir del
bloque central, seleccionado de manera única.

\section{Esperanza condicional respecto de la proyección \(\mathbb Z/9\mathbb Z\to\mathbb Z/3\mathbb Z\)}

La reducción canónica
\(\ZZ/9\ZZ\to\ZZ/3\ZZ\) determina la partición
\[
 C_1=\{1,4,7\},\qquad
 C_2=\{2,5,8\},\qquad
 C_0=\{3,6,9\}.
\]
La esperanza condicional uniforme promedia cada observable sobre los nueve puntos de cada bloque
\(C_a\times C_b\). Las matrices agregadas son
\[
 M_{\Pi}=
 \begin{pmatrix}4&5&6\\5&4&6\\6&6&9\end{pmatrix},
 \qquad
 M_{\Sigma}=
 \begin{pmatrix}5&6&4\\6&4&5\\4&5&6\end{pmatrix}.
\]
Puesto que cada bloque contiene nueve posiciones,
\[
 9\sum_{a,b}(M_\Pi)_{ab}=459,\qquad
 9\sum_{a,b}(M_\Sigma)_{ab}=405.
\]
Así, el promediado conserva los totales de ambos observables y localiza de
nuevo su diferencia:
\[
 \sum M_\Pi-\sum M_\Sigma=51-45=6,\qquad 9\cdot6=54.
\]

\begin{theorem}[Espectro del operador multiplicativo agregado]
\label{thm:app-espectro-agregado}
El operador matricial \(M_\Pi\) satisface
\[
 \det M_\Pi=-9
\]
y su espectro real es
\[
 \operatorname{Spec}(M_\Pi)
 =\{-1,\,9-6\sqrt2,\,9+6\sqrt2\}.
\]
En particular, la forma cuadrática asociada posee firma \((2,1)\).
\end{theorem}

\begin{proof}
La expansión del determinante da \(-9\). El polinomio característico
factoriza como
\[
 (\lambda+1)\bigl((\lambda-9)^2-72\bigr).
\]
Sus raíces son las indicadas y \(9-6\sqrt2>0\); por tanto existen dos
direcciones positivas y una negativa.
\end{proof}

La descomposición espectral admite una forma integral sobre una subretícula
invariante. Para
\[
 v_-=(1,-1,0),\qquad v_+=(1,1,0),\qquad v_0=(0,0,1),
\]
se tiene \(M_\Pi v_-=-v_-\), mientras que la restricción de \(M_\Pi\) al
submódulo generado por \(v_+,v_0\) está representada por
\[
 \begin{pmatrix}9&6\\12&9\end{pmatrix}
 =3H,\qquad
 H=\begin{pmatrix}3&2\\4&3\end{pmatrix}\in SL(2,\ZZ).
\]
Los autovalores de \(H\) son
\[
 3\pm2\sqrt2=(1\pm\sqrt2)^2.
\]
Los tres vectores \(v_-,v_+,v_0\) generan una subretícula de índice dos en
\(\ZZ^3\), pues el determinante de la matriz que los tiene por columnas vale
\(-2\). Por tanto, no se afirma aquí una suma directa integral de toda
\(\ZZ^3\): sobre esa subretícula invariante, el operador agregado separa una
dirección de autovalor \(-1\) y un submódulo hiperbólico gobernado por la
unidad cuadrática \(3+2\sqrt2\); sobre \(\RR^3\), la descomposición espectral
es completa.

\section{Conmutador de los proyectores aritméticos y separación de sus imágenes}

La diferencia entre \(\Sigma\) y \(\Pi\) no se reduce a sus sumas globales.
Para \(d\ge2\), sea \(\mathscr A_d=\{1,\ldots,9\}^d\) y definamos
\[
 \Sigma_d(x)=\dr(x_1+\cdots+x_d),\qquad
 \Pi_d(x)=\dr(x_1\cdots x_d).
\]
En el espacio de Hilbert \(\ell^2(\mathscr A_d)\), provisto de la medida
uniforme, denotemos por \(P_{\Sigma,d}\) y \(P_{\Pi,d}\) las esperanzas
condicionales ortogonales sobre las funciones medibles respecto de
\(\Sigma_d\) y \(\Pi_d\), respectivamente. La tabla conjunta normalizada es
\[
 C^{(d)}_{ab}
 =\frac{N^{(d)}_{ab}}{\sqrt{n_a m_b}},
\quad
 N^{(d)}_{ab}
 =\#\{x:\Sigma_d(x)=a,\ \Pi_d(x)=b\},
\]
donde \(n_a=\sum_bN^{(d)}_{ab}\) y \(m_b=\sum_aN^{(d)}_{ab}\).

En el soporte bidimensional original, la correlación conjunta exacta es
\[
N^{(2)}=
\begin{pmatrix}
0&0&2&0&0&2&3&0&2\\
3&0&2&0&0&2&0&0&2\\
0&2&0&0&2&0&0&2&3\\
0&0&2&3&0&2&0&0&2\\
0&0&2&3&0&2&0&0&2\\
0&2&0&0&2&0&0&2&3\\
3&0&2&0&0&2&0&0&2\\
0&0&2&0&0&2&3&0&2\\
0&2&0&0&2&0&0&2&3
\end{pmatrix}.
\]
Las filas, ordenadas por el valor de \(\Sigma\), suman nueve; las columnas,
ordenadas por el valor de \(\Pi\), suman
\((6,6,12,6,6,12,6,6,21)\). Así queda determinada por enumeración completa
la correlación conjunta de los dos observables antes de introducir cualquier
aplicación de evaluación posterior.

\begin{proposition}[Conmutadores en dimensiones dos y tres]
\label{prop:app-no-conmutatividad}
Los dos pares de proyectores satisfacen
\[
 \|[P_{\Sigma,2},P_{\Pi,2}]\|=\frac{\sqrt2}{3},
 \qquad
 \|[P_{\Sigma,3},P_{\Pi,3}]\|=\frac{\sqrt3}{4}.
\]
En particular, los observables aditivo y multiplicativo no inducen
descomposiciones simultáneamente compatibles.
\end{proposition}

\begin{proof}
Los valores singulares \(s\) de \(C^{(d)}\) son los cosenos de los ángulos
principales entre los dos subespacios. En el bloque bidimensional
correspondiente, la norma del conmutador es \(s\sqrt{1-s^2}\). La enumeración
exacta de las \(9^d\) palabras produce, para \(d=2\), los autovalores no
triviales
\[
 s^2\in\left\{\frac57,\frac13,\frac13\right\}.
\]
El máximo de \(s^2(1-s^2)\) es \(2/9\), de donde resulta
\(\sqrt2/3\). Para \(d=3\), la misma construcción da el máximo
\(s^2(1-s^2)=3/16\), alcanzado en \(s^2=1/4\), y por tanto la norma es
\(\sqrt3/4\). En ambos casos el conmutador es no nulo.
\end{proof}

\begin{proposition}[Intersección de los subespacios observables]
\label{prop:app-interseccion-proyectores}
La enumeración exacta de las palabras de \(\mathscr A_d\), para
\(2\leq d\leq15\), verifica
\[
 \operatorname{Ran}P_{\Sigma,d}\cap
 \operatorname{Ran}P_{\Pi,d}
 =\CC\mathbf 1.
\]
En ese intervalo de dimensiones, las funciones constantes constituyen el
único subespacio común a las dos descomposiciones aritméticas.
\end{proposition}

\begin{proof}
Para dos proyecciones ortogonales, la dimensión de la intersección de sus
rangos coincide con la multiplicidad del valor singular \(1\) en la matriz
de solapamiento \(C^{(d)}\). El cálculo entero de las tablas
\(N^{(d)}\), seguido del cálculo exacto de esa multiplicidad, da valor uno
para cada \(d=2,\ldots,15\). El vector constante pertenece a ambos rangos;
por tanto genera toda la intersección.
\end{proof}

La sucesión de normas no se extrapola a partir de los dos primeros casos.
En particular, para \(d=4\) el mismo cálculo da
\[
 \|[P_{\Sigma,4},P_{\Pi,4}]\|
 =\sqrt{\frac{2618939}{22240656}}
 \simeq0.3431538652,
\]
que no coincide con la prolongación elemental \(2/(d+1)\). Este caso
constituye un contraejemplo para cualquier fórmula inferida sólo de
\(d=2,3\).

\section{Bloque central y descomposición espectral}

La caracterización del bloque central se obtiene mediante una condición de
congruencia. Para dos representantes consecutivos \(k,k+1\), consideremos
\[
B_k=
\begin{pmatrix}
k^2&k(k+1)\\
k(k+1)&(k+1)^2
\end{pmatrix}\pmod9.
\]

\begin{theorem}[Unicidad del bloque diagonal adyacente]
\label{thm:app-bloque-central}
Existe un único bloque \(B_k\) cuyas entradas diagonales coinciden. Está
determinado por \(k=4\) y es
\[
B_c=
\begin{pmatrix}7&2\\2&7\end{pmatrix}.
\]
\end{theorem}

\begin{proof}
La igualdad de las diagonales equivale a
\[
k^2\equiv(k+1)^2\pmod9,
\]
es decir, \(2k+1\equiv0\pmod9\). Como \(2\) es una unidad de
\(\ZZ/9\ZZ\), la congruencia tiene la única solución \(k\equiv4\pmod9\).
\end{proof}

Por tanto, las posiciones \(4,5\) del observable multiplicativo forman
\[
B_c=
\begin{pmatrix}7&2\\2&7\end{pmatrix}
=7I+2R,
\qquad
R=\begin{pmatrix}0&1\\1&0\end{pmatrix},
\qquad R^2=I.
\]
Con \(P_\pm=(I\pm R)/2\),
\[
B_c=9P_++5P_-.
\]
\begin{theorem}[Jerarquía tipada de la separación suma--producto]
\label{thm:app-jerarquia-cuatro-dos-seis-cincuentaycuatro}
El bloque central tiene salto espectral local (9-5=4). Su coeficiente
fuera de la diagonal es (2). La compresión por las tres clases módulo tres
produce la diferencia agregada (51-45=6), y el despliegue sobre las nueve
posiciones produce el defecto global
\[
 9\cdot6=459-405=54.
\]
Los enteros (4), (2), (6) y (54) pertenecen, respectivamente, al
salto espectral, al acoplamiento local, a la compresión modular y a la
integración bidimensional.
\end{theorem}

\begin{proof}
La descomposición \(B_c=7I+2R=9P_++5P_-\) da el salto y el
coeficiente local. Las matrices agregadas \(M_\Pi,M_\Sigma\) construidas
arriba satisfacen
\(\sum M_\Pi=51\) y \(\sum M_\Sigma=45\). Cada entrada agregada representa
nueve posiciones, por lo que el despliegue multiplica la diferencia por
nueve y recupera los totales (459) y (405).
\end{proof}
Sus concatenaciones ordenadas producen los bloques \(72\) y \(27\):
\[
72+27=99,
\qquad
72-27=45=\det B_c.
\]
La concatenación de las entradas diagonales y antidiagonales produce,
respectivamente, \(77\) y \(22\). Por tanto,
\[
 72+27=77+22=99,\qquad
 77-22=55=54+1.
\]
Estas igualdades expresan dos particiones exactas del mismo bloque; no
identifican por sí solas una constante externa.

\begin{proposition}[Relación local--global del sector radical]
\label{prop:app-local-global-radical}
La partición por filas del bloque central reproduce, con factor de escala
tres, la partición del sector radical:
\[
 3\cdot72=216,\qquad
 3\cdot27=81,\qquad
 3\cdot99=297.
\]
\end{proposition}

\begin{proof}
El bloque \(C_0\times C_0\) contiene nueve representantes nulos y suma
\(81\). Los cuatro brazos restantes de las filas y columnas de \(C_0\)
suman \(216\). Las identidades siguen de \(72+27=99\).
\end{proof}

Además,
\[
M_c:=\frac{B_c}{\sqrt{45}}
=\exp\left(\frac12\log\frac95\,R\right)
\in SO^+(1,1).
\]
Si \(\eta=\operatorname{diag}(1,-1)\), la matriz normalizada \(M_c\) satisface
\(M_c^{\mathsf T}\eta M_c=\eta\), \(\det M_c=1\) y pertenece a
\(SO^+(1,1)\). Esta realización finita se limita al bloque \(2\times2\)
construido; su interpretación espacio-temporal requiere el puente dimensional
desarrollado en el tratado de Mecánica Dimensional.

\begin{proposition}[Familia espectral APP--TRIT y lectura
Markov--Lorentz local]
\label{prop:app-trit-espectral-markov-lorentz-rev8}
Para \(\tau\in\{-1,0,+1\}\), sea
\[
 B_\tau=(4+3\tau)I+(5-3\tau)R.
\]
Entonces
\[
 \operatorname{Spec}(B_\tau)=\{9,\,6\tau-1\},
 \qquad
 B_\tau=9P_++(6\tau-1)P_-,
\]
y el modo diferencial recupera la orientación mediante
\[
 \tau=\frac{\lambda_-+1}{6}.
\]
En particular, \(B_{+1}=B_c\). Su normalización estocástica
\[
 P_B:=\frac{1}{9}B_c=P_++\frac59P_-
\]
satisface, para todo \(n\geq0\),
\[
 P_B^n=P_++\left(\frac59\right)^n P_-.
\]
Por tanto, la contracción del modo antisimétrico es
\(\kappa_B=5/9\) y la brecha espectral de esta cadena local es \(4/9\).
Si
\[
 \eta_c:=\frac12\log\frac95,
\]
la normalización lorentziana del mismo bloque cumple
\[
 \frac{B_c}{\sqrt{45}}=\exp(\eta_c R),
 \qquad
 \boxed{\frac59=e^{-2\eta_c}}.
\]
\end{proposition}

\begin{proof}
Como \(R\) actúa con autovalores \(+1\) y \(-1\) sobre
\(\operatorname{im}P_+\) y \(\operatorname{im}P_-\), respectivamente, los
dos autovalores de \(B_\tau\) son
\[
 (4+3\tau)+(5-3\tau)=9,\qquad
 (4+3\tau)-(5-3\tau)=6\tau-1.
\]
La fórmula de las potencias sigue de
\(P_\pm^2=P_\pm\) y \(P_+P_-=0\). Finalmente,
\(e^{-2\eta_c}=e^{-\log(9/5)}=5/9\).
\end{proof}

Esta coincidencia identifica dos lecturas exactas del bloque local: la
memoria diferencial decrece con el mismo cociente que codifica la rapidez
en la realización descompuesta. No convierte \(B_\tau\) en el núcleo de
Markov global de TPK, ni \(5/9\) en una constante física. Tampoco extiende
por sí sola la realización \(SO^+(1,1)\) a un espacio de Minkowski
\(3+1\): esos pasos requieren mapas dimensionales adicionales.

\section{Curvatura discreta inducida por la interacción de las hojas suma y producto y su defecto de conmutación}

En el toro escribimos, ahora en residuos,
\[
S(x,y)=x+y,
\qquad
P(x,y)=xy\pmod9.
\]
Para toda función \(T:X_9\to\ZZ/9\ZZ\), fijamos las diferencias progresivas
\[
 \Delta_xT(x,y)=T(x+1,y)-T(x,y),\qquad
 \Delta_yT(x,y)=T(x,y+1)-T(x,y),
\]
donde las coordenadas se toman módulo nueve. Esta convención fija también la
orientación positiva de cada plaqueta: primero se recorre la dirección
\(x\), después la dirección \(y\), y se regresa por las aristas opuestas.
Es importante tipar correctamente estos dos objetos. \(S\) y \(P\) son
potenciales de vértice; las variables de enlace son sus diferencias
discretas
\[
A_x^T=\Delta_xT,
\qquad
A_y^T=\Delta_yT.
\]
Para una mezcla ordenada de potenciales \((T_x,T_y)\), definimos la
curvatura de plaqueta por
\[
F(T_x,T_y)
=\Delta_x A_y^{T_y}-\Delta_y A_x^{T_x}
=\Delta_x\Delta_yT_y-\Delta_y\Delta_xT_x.
\]
Como
\[
\Delta_xS=\Delta_yS=1,
\qquad
\Delta_xP=y,
\qquad
\Delta_yP=x,
\]
se obtiene en las \(81\) plaquetas:
\[
F(S,S)=F(P,P)=0,
\qquad
F(S,P)=+1,
\qquad
F(P,S)=-1\pmod9.
\]
Las conexiones asociadas a un único observable son planas; la composición
mixta de los dos observables produce curvatura de signo opuesto según el
orden. Para un rectángulo coordenado \(R_{a,b}\) contenido
en la carta, definimos el Wilson aditivo de borde como la suma orientada de
las variables de enlace. La cancelación de las aristas interiores da
\[
 W_{\partial R_{a,b}}(S,P)=ab,\qquad
 W_{\partial R_{a,b}}(P,S)=-ab\pmod9,
\]
que es la identidad de Stokes finita para esta convención.

La distinción entre potencial y enlace es algebraicamente necesaria. Si se sustituyesen
directamente \(A_x,A_y\) por \(S,P\), aparecerían términos dependientes de
\(x,y\) y no se obtendrían los valores constantes anteriores. La formulación
tipada es la que produce el cociclo mixto utilizado más adelante en el cierre
de Heisenberg finito.

\begin{remark}
La diferencia integrada \(54\) y la curvatura mixta \(\pm1\) son invariantes
distintos. La primera compara sumas globales; la segunda depende del orden de
los dos potenciales y se localiza en cada plaqueta. La formulación mediante
operadores autoadjuntos y conmutadores se desarrollará después de construir el
espacio de estados correspondiente.
\end{remark}

\section{Expansión del operador de transición por caminos}

La definición de APP incluye el espacio
\(\operatorname{Path}(\GraphAPP)\). Una vez fijado un operador de
transición \(U\) sobre los vértices de \(\GraphAPP\), suponemos que
\(U_{yx}=0\) siempre que \(x\to y\) no sea una arista del grafo. Bajo esta
hipótesis de soporte, sus potencias admiten la expansión
\begin{proposition}[Expansión finita por caminos]
\label{prop:app-suma-caminos}
Para cada \(R\ge1\) y cada par de vértices \(i,f\),
\[
(U^R)_{fi}
=\sum_{\substack{\gamma:i\to f\\|\gamma|=R}}
\prod_{k=0}^{R-1}U_{x_{k+1},x_k},
\qquad
\gamma=(x_0=i,x_1,\ldots,x_R=f).
\]
\end{proposition}

\begin{proof}
Para \(R=1\) la identidad es la definición de los elementos de matriz de
\(U\). Si vale para \(R\), entonces
\[
 (U^{R+1})_{fi}
 =\sum_j U_{fj}(U^R)_{ji}.
\]
Sustituir la hipótesis inductiva concatena a cada camino de longitud \(R\)
desde \(i\) hasta \(j\) la última arista \(j\to f\); así, cada camino de
longitud~\(R+1\) aparece exactamente una vez.\qedhere
\end{proof}

En consecuencia, APP reúne en una sola estructura el toro aritmético finito,
los observables de suma y producto, la descomposición residuo--cociente, el
radical nilpotente y una curvatura mixta orientada. El representante \(9\) de
la clase nula fija una carta visible del cociente; reemplazarlo por \(0\) sin
transportar simultáneamente la carta altera los registros empleados por la
dinámica posterior.
\fi
 \end{HMTScientificOwnerSubordinated}
\endgroup

% Unidad U003. Redacción autónoma; tablas y censos reconstruidos desde 1063/live.

\section{Matriz aritmética de APP y estructura del anillo local}
\label{p12-83-c1-s1:chap:app-hojas-algebra-local}

La \cref{p12-83-c1-s2:chap:app-completacion-gnomonica} construyó un único soporte
\(X_9\). APP posee una única matriz aritmética principal: la multiplicación
nonádica publicada por raíz digital positiva. La suma no constituye una
segunda tabla APP; es la ley interna de acumulación que genera las filas y
columnas de la matriz multiplicativa y proporciona un observable auxiliar de
comparación. Esta distinción se mantendrá en toda la obra.

\subsection{Matriz multiplicativa nonádica y raíz digital}

Para \(a,b\in\mathcal D_9=\{1,\ldots,9\}\), la entrada de APP se obtiene
directamente por
\begin{equation}
 \boxed{\Pi(a,b)=\operatorname{dr}_9(ab)
 =1+((ab-1)\bmod9).}
 \label{p12-83-c1-s1:eq:app-matriz-multiplicativa-directa}
\end{equation}
La única matriz APP es, por tanto,
\[
\begin{array}{c|ccccccccc}
\Pi&1&2&3&4&5&6&7&8&9\\\hline
1&1&2&3&4&5&6&7&8&9\\
2&2&4&6&8&1&3&5&7&9\\
3&3&6&9&3&6&9&3&6&9\\
4&4&8&3&7&2&6&1&5&9\\
5&5&1&6&2&7&3&8&4&9\\
6&6&3&9&6&3&9&6&3&9\\
7&7&5&3&1&8&6&4&2&9\\
8&8&7&6&5&4&3&2&1&9\\
9&9&9&9&9&9&9&9&9&9
\end{array}
\]
Cada fila puede reconstruirse también por adición reiterada de su índice,
pero esa recurrencia describe la ley interna de la matriz y no una segunda
hoja. Las filas y columnas de índices \(1,2,4,5,7,8\) recorren las nueve
clases; las de índices \(3,6,9\) son exactamente las que contraen el soporte.

En la carta residual, el conjunto de unidades es
\[
 R_9^\times=(\mathbb Z/9\mathbb Z)^\times
   =\{[1],[2],[4],[5],[7],[8]\}.
\]
El ideal maximal es
\[
 \mathfrak m=(3)=\{[0],[3],[6]\}.
\]

\begin{theorem}[Clasificación de filas multiplicativas]
\label{p12-83-c1-s1:thm:app-clasificacion-filas-producto}
Para \(a\in\mathcal D_9\), la aplicación
\[
 m_a:\mathcal D_9\longrightarrow\mathcal D_9,
 \qquad b\longmapsto\Pi(a,b)
\]
es biyectiva si y sólo si \([a]_9\in R_9^\times\). Si
\([a]_9\in\mathfrak m\setminus\{0\}\), su imagen tiene tres elementos. Si
\([a]_9=[0]_9\), su imagen tiene un solo elemento.
\end{theorem}

\begin{proof}
La reducción residual de \(m_a\) es la multiplicación por \([a]_9\). En un
anillo finito, esa multiplicación es biyectiva exactamente cuando \([a]_9\)
es una unidad. Si \([a]_9=[3]\) o \([6]\), su imagen es el ideal
\(\mathfrak m\), de cardinal tres. Si \([a]_9=[0]\), toda entrada se envía
a la clase nula, publicada como \(9\).
\end{proof}

\begin{proposition}[Latinidad de la ley aditiva interna]
\label{p12-83-c1-s1:prop:app-latinidad-aditiva}
El observable auxiliar
\[
 \Sigma(a,b)=\operatorname{dr}_9(a+b)
\]
actúa por traslaciones: para cada \(a,c\in\mathcal D_9\) existe un único
\(b\in\mathcal D_9\) tal que \(\Sigma(a,b)=c\), y análogamente en la otra
coordenada.
\end{proposition}

\begin{proof}
En la carta residual, \([a+b]_9=[c]_9\) tiene la solución única
\([b]_9=[c-a]_9\). La sección positiva transporta esa clase a un único
elemento de \(\mathcal D_9\).
\end{proof}

La diferencia de geometría de fibras queda así tipada sin duplicar la matriz
APP: la ley aditiva interna sólo traslada, mientras la matriz multiplicativa
permuta sobre las unidades y contrae sobre el radical.

\begin{proposition}[Estructura lineal del radical]
\label{p12-83-c1-s1:prop:app-radical-lineal}
La aplicación
\[
 M_3:R_9\longrightarrow R_9,
 \qquad x\longmapsto3x
\]
satisface
\[
 \ker M_3=\operatorname{im}M_3=\mathfrak m.
\]
El mapa inducido
\[
 R_9/\mathfrak m\longrightarrow\mathfrak m,
 \qquad [x]\longmapsto3x
\]
es un isomorfismo de espacios vectoriales sobre \(\mathbb F_3\), pero no un
isomorfismo de anillos.
\end{proposition}

\begin{proof}
La congruencia \(3x\equiv0\pmod9\) equivale a \(x\equiv0\pmod3\), de
donde el núcleo \(\mathfrak m\). La imagen es
\(\{[0],[3],[6]\}=\mathfrak m\). El primer teorema de isomorfía da el
isomorfismo lineal. No es un isomorfismo de anillos porque
\(\mathfrak m^2=0\) en \(R_9\), mientras el cociente
\(R_9/\mathfrak m\cong\mathbb F_3\) tiene unidad no nula.
\end{proof}

La sección positiva publica el ciclo radical no trivial como
\[
 369\longmapsto693\longmapsto936\longmapsto369,
\]
según la posición desde la que se lea la terna. Esta notación conserva el
orden de las coordenadas; no convierte el ideal en tres números enteros
independientes.

\subsection{Levantamientos y cociclo de acarreo}

Sea \(s_0:R_9\to\{0,\ldots,8\}\) la sección residual normalizada. El
acarreo binario de la suma es
\[
 \kappa_0(a,b)
 =\frac{s_0(a)+s_0(b)-s_0(a+b)}9.
\]

\begin{lemma}[Identidad cocíclica]
\label{p12-83-c1-s1:lem:app-identidad-cociclo-acarreo}
Para \(a,b,c\in R_9\),
\[
 \kappa_0(a,b)+\kappa_0(a+b,c)
 =\kappa_0(b,c)+\kappa_0(a,b+c).
\]
\end{lemma}

\begin{proof}
Multiplicar por nueve y desarrollar cualquiera de los dos miembros produce
\[
 s_0(a)+s_0(b)+s_0(c)-s_0(a+b+c).
\]
Por tanto, ambos son iguales.
\end{proof}

Definimos \(\chi_0(a)=1\) si \(a=[0]_9\) y \(0\) en otro caso. Como
\[
 s_+(a)=s_0(a)+9\chi_0(a),
\]
los acarreos de las dos secciones se relacionan mediante
\[
 \kappa_+(a,b)=
 \kappa_0(a,b)+\chi_0(a)+\chi_0(b)-\chi_0(a+b).
\]
El cambio de representante modifica el cociclo por una coborde. Ésta es la
forma exacta de transportar la memoria cuando la clase nula se publica como
\(9\) en vez de \(0\).

\subsection{Censos integrados de residuos y cocientes}

Sumando sobre todos los pares \((i,j)\in\mathcal D_9^2\), se obtiene
\[
 \sum_{i,j}(i+j)=810,
 \qquad
 \sum_{i,j}ij=2025.
\]
Los residuos positivos satisfacen
\[
 \sum_{i,j}R^+_{i,j}=405,
 \qquad
 \sum_{i,j}R^\times_{i,j}=459.
\]
Por las identidades levantadas,
\[
 \sum_{i,j}Q^+_{i,j}=\frac{810-405}{9}=45,
 \qquad
 \sum_{i,j}Q^\times_{i,j}=\frac{2025-459}{9}=174.
\]

\begin{theorem}[Separación exacta del exceso suma--producto]
\label{p12-83-c1-s1:thm:app-separacion-exceso}
La diferencia integrada entre producto y suma se descompone de manera única
en una parte publicada y una parte de cociente:
\[
 \boxed{
  2025-810
  =(459-405)+9(174-45)
  =54+9\cdot129.}
\]
\end{theorem}

\begin{proof}
Restamos, término a término, las identidades
\(ij=R^\times_{i,j}+9Q^\times_{i,j}\) e
\(i+j=R^+_{i,j}+9Q^+_{i,j}\), y sumamos sobre los ochenta y un pares. Los
censos anteriores dan la igualdad. La unicidad procede de la división
euclídea en cada par.
\end{proof}

El valor \(54\) no es la diferencia completa entre dos supuestas matrices.
Es el defecto visible entre la matriz multiplicativa y el acumulador aditivo
en la carta positiva. El término \(9\cdot129\) es la contribución
que permanece en los cocientes. Omitirlo confundiría publicación con
genealogía.

\subsection{Bloque central de \(2\) posiciones}

Para \(k\in\mathbb Z/9\mathbb Z\), consideramos el bloque simétrico
\[
 B_k=
 \begin{pmatrix}
  k^2&k(k+1)\\
  k(k+1)&(k+1)^2
 \end{pmatrix}\pmod9.
\]
Las diagonales coinciden exactamente cuando
\[
 k^2\equiv(k+1)^2\pmod9
 \quad\Longleftrightarrow\quad
 2k+1\equiv0\pmod9.
\]
Como \(2^{-1}=5\pmod9\), la única solución es \(k=4\). Así aparece el
bloque
\[
 B_c=
 \begin{pmatrix}7&2\\2&7\end{pmatrix}.
\]
Sea
\[
 R=\begin{pmatrix}0&1\\1&0\end{pmatrix},
 \qquad
 P_\pm=\frac12(I\pm R).
\]
Entonces
\[
 B_c=7I+2R=9P_++5P_-.
\]

\begin{theorem}[Unicidad y espectro del bloque central]
\label{p12-83-c1-s1:thm:app-bloque-central}
El bloque \(B_c\) es el único bloque adyacente de la familia \(B_k\) con
diagonales iguales. Sus autovalores en los subespacios simétrico y
antisimétrico son, respectivamente, \(9\) y \(5\).
\end{theorem}

\begin{proof}
La congruencia anterior prueba la unicidad. Puesto que \(R\) actúa por
\(+1\) sobre \(\operatorname{im}P_+\) y por \(-1\) sobre
\(\operatorname{im}P_-\), la matriz \(7I+2R\) actúa por \(9\) y \(5\),
respectivamente.
\end{proof}

Las dos lecturas ordenadas de sus cifras dan
\[
 72+27=99,
 \qquad
 72-27=45=\det B_c,
\]
y las dos diagonales dan
\[
 77+22=99,
 \qquad
 77-22=55=54+1.
\]
Estas identidades son controles internos del bloque; no definen por sí solas
el defecto global.

\subsection{Curvatura mixta de las \(2\) leyes}

En la carta residual escribimos
\[
 S(x,y)=x+y,
 \qquad
 P(x,y)=xy.
\]
Para una función \(T:X_9\to R_9\), sean
\[
 \Delta_xT(x,y)=T(x+1,y)-T(x,y),
 \qquad
 \Delta_yT(x,y)=T(x,y+1)-T(x,y).
\]
Definimos la curvatura ordenada de un par de leyes por
\[
 \mathcal F(T_x,T_y)
 :=\Delta_x\Delta_yT_y-\Delta_y\Delta_xT_x.
\]

\begin{proposition}[Orientaciones de la composición mixta]
\label{p12-83-c1-s1:prop:app-curvatura-mixta}
Se cumplen
\[
 \mathcal F(S,S)=0,
 \qquad
 \mathcal F(P,P)=0,
\]
y, para las composiciones mixtas,
\[
 \mathcal F(S,P)=+1,
 \qquad
 \mathcal F(P,S)=-1
 \quad\text{en }R_9.
\]
\end{proposition}

\begin{proof}
Las segundas diferencias mixtas de la ley lineal \(S\) se anulan. Para
\(P(x,y)=xy\), ambas diferencias mixtas puras valen uno y se cancelan al
compararse con el mismo orden. Al combinar una ley lineal y una bilineal,
queda una única unidad cuyo signo cambia al invertir el orden de las hojas.
\end{proof}

La curvatura no crea otra matriz APP. Registra que el orden
\(\Sigma\to\Pi\) y el orden \(\Pi\to\Sigma\) son orientaciones opuestas,
mientras los sectores puros permanecen planos. Esta tríada
\(-1,0,+1\) será tipada en la unidad TRIT.

\subsection{Obstrucciones de la matriz aritmética de APP: traslaciones, radical local, bloque central, curvatura mixta y conservación del acarreo}

\begin{criterion}[Criterios de refutación de la matriz APP]
La unidad queda refutada si:
\begin{enumerate}
 \item la ley aditiva interna deja de actuar por traslaciones;
 \item una fila multiplicativa no unitaria se declara biyectiva;
 \item se identifica \(R_9/\mathfrak m\) con \(\mathfrak m\) como anillo;
 \item el exceso \(1215\) se sustituye por \(54\) sin conservar cocientes;
 \item el bloque central se escoge sin resolver \(2k+1=0\pmod9\);
 \item la inversión del orden mixto no cambia el signo de la curvatura;
 \item el acarreo se elimina al cambiar entre las secciones \(s_0\) y
       \(s_+\).
\end{enumerate}
\end{criterion}

APP está ahora construida como soporte, categoría de trayectorias, matriz
multiplicativa nonádica, ley aditiva interna, radical local, cociclo de
acarreo y orientación mixta. La unidad
siguiente no resumirá estas estructuras: tomará la tríada de orientaciones
que acaba de emerger y construirá el estado completo del TRIT con sus tres regímenes
cuadráticos.


% Unidad U002. Redacción autónoma a partir del generador integral 1063/live.

\section{Completación gnomónica y soporte aritmético de APP}
\label{p12-83-c1-s2:chap:app-completacion-gnomonica}

La presente unidad construye el soporte de APP a partir de la célula
nonádica de la \cref{p12-83-c1-s4:chap:segmentos-marcas-intersticios}. No se presupone un
plano continuo, una retícula infinita ni una constante numérica. El objeto de
partida es finito: dos copias ordenadas de la carta positiva de
\(\mathbb Z/9\mathbb Z\). La completación de sus intersticios interiores
produce exactamente ochenta y un estados, y la ley de traslación residual
convierte ese soporte cuadrado en un toro discreto.

En toda la unidad, APP designa una estructura matemática y no una imagen.
La sigla APP nombra en lo sucesivo esta estructura; las propiedades que
siguen dependen de sus objetos y morfismos, no de una expansión verbal de la
sigla.

\subsection{Correspondencia entre \(2\) cartas de una misma clase residual}

Definimos
\[
 \mathcal D_9:=\{1,2,\ldots,9\},
 \qquad
 \mathcal D_8:=\{1,2,\ldots,8\}.
\]
La aplicación de publicación residual es
\[
 \lambda_9:\mathcal D_9\longrightarrow\mathbb Z/9\mathbb Z,
 \qquad
 \lambda_9(a)=[a]_9.
\]
Por tanto, \(\lambda_9(9)=[0]_9\). Su inversa como aplicación de
conjuntos es la sección positiva
\[
 s_+([a]_9)=
 \begin{cases}
  a,&[a]_9\neq[0]_9,\quad 1\leq a\leq8,\\
  9,&[a]_9=[0]_9.
 \end{cases}
\]
Las dos expresiones \(9\) y \([0]_9\) no son dos fases: son dos
representantes del mismo elemento en cartas diferentes.

Para dos coordenadas escribiremos
\[
 \lambda_9^{\times2}:\mathcal D_9^2
       \longrightarrow X_9:=(\mathbb Z/9\mathbb Z)^2,
 \qquad
 (a,b)\longmapsto([a]_9,[b]_9).
\]
Esta aplicación es una biyección de conjuntos. La carta
\(\mathcal D_9^2\) hace visibles las marcas positivas; el grupo \(X_9\)
hace visibles la suma residual y las traslaciones.

\begin{lemma}[Cambio de carta sin pérdida de estado]
\label{p12-83-c1-s2:lem:app-cambio-carta-sin-perdida}
Las aplicaciones \(\lambda_9^{\times2}\) y
\(s_+^{\times2}\) son inversas. En particular, la escritura positiva y la
escritura residual contienen exactamente los mismos ochenta y un estados.
\end{lemma}

\begin{proof}
En cada coordenada se cumplen
\(\lambda_9\circ s_+=\operatorname{id}_{\mathbb Z/9\mathbb Z}\) y
\(s_+\circ\lambda_9=\operatorname{id}_{\mathcal D_9}\). El producto
cartesiano de estas identidades proporciona las dos identidades requeridas.
\end{proof}

\subsection{Núcleo interior y complemento gnomónico}

El núcleo interior es
\[
 \mathcal N_8:=\mathcal D_8\times\mathcal D_8,
 \qquad |\mathcal N_8|=8^2=64.
\]
En el plano de índices, \(\mathcal N_8\) es el cuadrado que no utiliza la
marca de cierre en ninguna coordenada. Para publicar también la clase nula
se debe completar la novena fila y la novena columna. Definimos el
complemento gnomónico disjunto por
\[
 \mathcal G_9:=
  (\{9\}\times\mathcal D_9)
  \sqcup
  (\mathcal D_8\times\{9\}).
\]
La esquina \((9,9)\) pertenece a la primera pieza; por eso la segunda usa
\(\mathcal D_8\) y no \(\mathcal D_9\). En consecuencia,
\[
 |\mathcal G_9|=9+8=17,
 \qquad
 \mathcal D_9^2=\mathcal N_8\sqcup\mathcal G_9,
 \qquad
 81=64+17.
\]

\begin{proposition}[Minimalidad de la completación]
\label{p12-83-c1-s2:prop:app-minimalidad-completacion}
Sea \(Y\) un subconjunto de \(\mathcal D_9^2\) que contiene
\(\mathcal N_8\) y cuya proyección sobre cada coordenada es todo
\(\mathcal D_9\). Si, además, \(Y\) contiene todos los pares necesarios
para trasladar una coordenada mientras la otra permanece arbitraria,
entonces \(\mathcal G_9\subseteq Y\). Por tanto, la adición de diecisiete
estados es la completación mínima compatible con las dos traslaciones.
\end{proposition}

\begin{proof}
Para trasladar la primera coordenada a través de la clase de cierre mientras
la segunda permanece en un valor arbitrario \(b\in\mathcal D_9\), son
necesarios todos los pares \((9,b)\). Se obtienen nueve estados. Para
trasladar la segunda coordenada a través del cierre mientras la primera toma
un valor arbitrario \(a\in\mathcal D_8\), son necesarios los ocho pares
\((a,9)\) restantes. La esquina \((9,9)\) ya fue contada en la primera
familia. En total son necesarios \(9+8=17\) estados, exactamente los de
\(\mathcal G_9\).
\end{proof}

La hipótesis de compatibilidad con las dos traslaciones es esencial. Bastaría
añadir una sola marca por coordenada para que ambas proyecciones fueran
sobreyectivas, pero ese conjunto no contendría los estados fronterizos
necesarios para mover una coordenada con la otra fija. APP requiere el
producto completo, no dos ejes desconectados.

\subsection{Realización exacta del intervalo intersticial mediante celdas APP y coordenadas de frontera}

Sea \(P_9\) el camino ordenado con nueve marcas y ocho aristas interiores
\[
 E(P_9)=\{e_1,\ldots,e_8\}.
\]
El cierre de los extremos incorpora una arista \(e_9\) y produce el ciclo
\(C_9\):
\[
 E(C_9)=E(P_9)\sqcup\{e_9\}.
\]
Los pares de intersticios interiores forman \(E(P_9)^2\). Los pares de
intersticios del ciclo completo forman \(E(C_9)^2\). La diferencia es
\begin{align}
 E(C_9)^2\setminus E(P_9)^2
   &=(\{e_9\}\times E(C_9))
     \sqcup(E(P_9)\times\{e_9\}).
\label{p12-83-c1-s2:eq:app-complemento-intersticial}
\end{align}
La descomposición es disjunta porque la esquina \((e_9,e_9)\) se reserva a
la primera pieza.

\begin{theorem}[Equivalencia del modelo de marcas y el modelo de
intersticios]
\label{p12-83-c1-s2:thm:app-equivalencia-modelos}
La aplicación
\[
 \beta:E(C_9)^2\longrightarrow\mathcal D_9^2,
 \qquad
 \beta(e_i,e_j)=(i,j),
\]
es biyectiva. Sus restricciones satisfacen
\[
 \beta(E(P_9)^2)=\mathcal N_8,
 \qquad
 \beta(E(C_9)^2\setminus E(P_9)^2)=\mathcal G_9.
\]
Por tanto, la identidad \(81=64+17\) cuenta simultáneamente pares de
marcas publicadas y pares de intersticios cerrados.
\end{theorem}

\begin{proof}
Cada arista de \(C_9\) tiene un índice único en \(\mathcal D_9\). La
aplicación \((i,j)\mapsto(e_i,e_j)\) es la inversa explícita de \(\beta\).
Si \(i,j\leq8\), el par pertenece a \(E(P_9)^2\) y su imagen pertenece a
\(\mathcal D_8^2\). Si al menos uno de los índices es nueve, el par pertenece
al complemento de \cref{p12-83-c1-s2:eq:app-complemento-intersticial}; separar el caso
\(i=9\) del caso \(i\leq8,j=9\) produce exactamente las dos piezas de
\(\mathcal G_9\).
\end{proof}

\subsection{Cierre toroidal y traslaciones}

Sean
\[
 \mathbf e_1=([1]_9,[0]_9),
 \qquad
 \mathbf e_2=([0]_9,[1]_9).
\]
El conjunto de direcciones elementales es
\[
 \mathscr D_{\rm APP}:=
 \{+\mathbf e_1,-\mathbf e_1,+\mathbf e_2,-\mathbf e_2\}.
\]
Definimos el grafo de Cayley orientado
\[
 \Gamma_{\rm APP}:=\operatorname{Cay}(X_9,\mathscr D_{\rm APP}).
\]
Sus vértices son los ochenta y un elementos de \(X_9\). Para cada
\(x\in X_9\) y \(d\in\mathscr D_{\rm APP}\) existe una arista orientada
\(x\xrightarrow{d}x+d\). La arista inversa lleva la etiqueta \(-d\).

En la carta positiva, la traslación unitaria se escribe mediante el sucesor
publicado y su acarreo:
\[
 u_+(a)=
 \begin{cases}
  a+1,&a<9,\\
  1,&a=9,
 \end{cases}
 \qquad
 \kappa_+(a)=
 \begin{cases}
  0,&a<9,\\
  1,&a=9.
 \end{cases}
\]
La identidad levantada es
\begin{equation}
 a+1=u_+(a)+9\kappa_+(a).
\label{p12-83-c1-s2:eq:app-sucesor-publicado-acarreo}
\end{equation}
Al aplicar \(\lambda_9\), el término de acarreo se anula y queda
\[
 \lambda_9(u_+(a))=\lambda_9(a)+[1]_9.
\]

\begin{lemma}[Compatibilidad entre cierre, residuo y elevación]
\label{p12-83-c1-s2:lem:app-compatibilidad-cierre-residuo}
La adición de \(e_9\), la identificación \(9\sim[0]_9\) y la traslación
\(x\mapsto x+\mathbf e_i\) describen tres operaciones compatibles en
dominios distintos. El cociente residual publica el retorno; el acarreo de
\cref{p12-83-c1-s2:eq:app-sucesor-publicado-acarreo} conserva la elevación que el
cociente omite.
\end{lemma}

\begin{proof}
El cierre topológico incorpora la arista que une los extremos del camino. La
carta residual identifica el extremo publicado con la clase nula. Si
\(a<9\), el sucesor tiene acarreo cero; si \(a=9\), la igualdad levantada
es \(10=1+9\). Ambos casos descienden a la misma suma residual, pero el
valor de \(\kappa_+\) los distingue.
\end{proof}

\begin{corollary}[Soporte toroidal]
\label{p12-83-c1-s2:cor:app-soporte-toroidal}
Identificar lados opuestos de la carta \(\mathcal D_9^2\) produce una
realización geométrica del grupo \(X_9\). Cada traslación elemental está
definida en todo vértice y posee inversa.
\end{corollary}

\subsection{Categoría de trayectorias APP: objetos, morfismos de prolongación y composición compatible}

\begin{definition}[Trayectoria orientada]
Una trayectoria de longitud \(n\geq0\) es un par
\[
 \gamma=(x_0;d_1,\ldots,d_n),
 \qquad
 x_0\in X_9,\qquad d_k\in\mathscr D_{\rm APP}.
\]
Sus vértices son
\[
 x_k=x_0+\sum_{r=1}^k d_r,
 \qquad 0\leq k\leq n.
\]
La fuente es \(s(\gamma)=x_0\) y el blanco es
\(t(\gamma)=x_n\). Para \(n=0\) se obtiene la trayectoria identidad en
\(x_0\).
\end{definition}

Si \(t(\gamma)=s(\delta)\), su concatenación es
\[
 \delta\circ\gamma
  =(x_0;d_1,\ldots,d_n,e_1,\ldots,e_m),
\]
donde \(e_1,\ldots,e_m\) son las direcciones de \(\delta\). La
concatenación es asociativa y las trayectorias de longitud cero son
identidades. Obtenemos así la categoría libre
\[
 \mathsf{Path}(\Gamma_{\rm APP}).
\]

\begin{proposition}[Finitud del soporte y numerabilidad de las trayectorias]
\label{p12-83-c1-s2:prop:app-soporte-finito-trayectorias-numerables}
El conjunto de estados de APP es finito, de cardinal \(81\). Para una fuente
fijada existen \(4^n\) palabras de direcciones de longitud \(n\). La unión
de todas las trayectorias finitas es numerable.
\end{proposition}

\begin{proof}
La primera afirmación procede de \(|X_9|=9^2\). En cada uno de los \(n\)
pasos hay cuatro elecciones independientes, de donde \(4^n\). Finalmente,
\(\bigcup_{n\geq0}X_9\times\mathscr D_{\rm APP}^n\) es una unión numerable
de conjuntos finitos.
\end{proof}

Esta proposición separa dos niveles que no se identificarán: APP posee un
soporte finito, pero admite trayectorias de cualquier longitud. La
prolongación de una ruta no crea nuevos vértices del soporte; crea nueva
genealogía sobre ellos.

\subsection{Matriz multiplicativa y ley aditiva interna sobre un único soporte}

Para \(n\in\mathbb Z_{>0}\), la raíz digital positiva es
\[
 \operatorname{dr}_9(n):=1+((n-1)\bmod9).
\]
Sobre la carta positiva definimos la matriz principal y la ley interna
\[
 \Sigma(a,b):=\operatorname{dr}_9(a+b),
 \qquad
 \Pi(a,b):=\operatorname{dr}_9(ab).
\]
Ambas aplicaciones tienen el mismo dominio \(\mathcal D_9^2\) y el mismo
codominio \(\mathcal D_9\), pero no tienen el mismo estatuto: \(\Pi\) publica
la única matriz APP y \(\Sigma\) genera sus acumulaciones. Sus versiones residuales
son
\[
 \bar\Sigma([a]_9,[b]_9)=[a+b]_9,
 \qquad
 \bar\Pi([a]_9,[b]_9)=[ab]_9.
\]
El cambio de carta es compatible con ambas operaciones:
\begin{align*}
 \lambda_9\circ\Sigma
   &=\bar\Sigma\circ\lambda_9^{\times2},\\
 \lambda_9\circ\Pi
   &=\bar\Pi\circ\lambda_9^{\times2}.
\end{align*}

Para no perder el cociente entero se introducen los levantamientos
\begin{align}
 a+b &= R^+_{a,b}+9Q^+_{a,b},
 &1\leq R^+_{a,b}\leq9,
\label{p12-83-c1-s2:eq:app-lift-suma}\\
 ab &= R^\times_{a,b}+9Q^\times_{a,b},
 &1\leq R^\times_{a,b}\leq9.
\label{p12-83-c1-s2:eq:app-lift-producto}
\end{align}
Aquí
\[
 R^+_{a,b}=\Sigma(a,b),
 \qquad
 R^\times_{a,b}=\Pi(a,b),
\]
y los cocientes están determinados unívocamente por
\[
 Q^+_{a,b}=\frac{a+b-R^+_{a,b}}9,
 \qquad
 Q^\times_{a,b}=\frac{ab-R^\times_{a,b}}9.
\]

\begin{definition}[APP genealógica]
\label{p12-83-c1-s2:def:app-genealogica-autonoma}
La estructura APP en esta etapa es la tupla
\[
 \operatorname{APP}:=
 \bigl(
  X_9,
  \Gamma_{\rm APP},
  \mathsf{Path}(\Gamma_{\rm APP}),
  \Sigma,\Pi,
  R^+,Q^+,R^\times,Q^\times
 \bigr).
\]
La genealogía de una trayectoria conserva, para cada vértice visitado, la
posición, la dirección de entrada, la entrada multiplicativa, la acumulación
aditiva y sus respectivos residuos y cocientes.
\end{definition}

\begin{theorem}[Construcción causal del soporte APP]
\label{p12-83-c1-s2:thm:app-cadena-constructiva-autonoma}
La cadena
\[
 \mathcal D_9
 \longrightarrow
 \mathcal D_8^2
 \longrightarrow
 \mathcal D_8^2\sqcup\mathcal G_9
 \longrightarrow
 \mathcal D_9^2
 \xrightarrow{\lambda_9^{\times2}}
 X_9
 \longrightarrow
 \Gamma_{\rm APP}
 \longrightarrow
 \mathsf{Path}(\Gamma_{\rm APP})
 \longrightarrow
 (\Sigma,\Pi)
\]
construye, sin pérdida cardinal, el soporte, el movimiento, las trayectorias,
la matriz multiplicativa de APP y su ley aditiva interna.
\end{theorem}

\begin{proof}
La \cref{p12-83-c1-s2:prop:app-minimalidad-completacion} construye la carta completa a
partir del núcleo. El \cref{p12-83-c1-s2:thm:app-equivalencia-modelos} identifica su
modelo intersticial. El \cref{p12-83-c1-s2:lem:app-cambio-carta-sin-perdida} transporta
los ochenta y un estados a \(X_9\). El grafo de Cayley incorpora las cuatro
traslaciones reversibles sin cambiar los vértices. La categoría libre
incorpora todas las concatenaciones finitas. Finalmente, \(\Sigma\) y
\(\Pi\) son aplicaciones totales distintas sobre el mismo dominio; las
ecuaciones \cref{p12-83-c1-s2:eq:app-lift-suma,p12-83-c1-s2:eq:app-lift-producto} conservan la parte
que la reducción residual no publica.
\end{proof}

\subsection{Descomposición del ortograma APP en núcleo \(8\times8\), gnomon de \(17\) estados y carta \(9\times9\) con lados opuestos identificados}

La \cref{fig:app-ortograma} debe leerse en tres capas: el cuadrado interior
de sesenta y cuatro estados, el gnomon de diecisiete estados y la
identificación de lados opuestos. La carta de nueve por nueve completa así el
núcleo de ocho por ocho mediante el gnomon de cierre. La matriz
multiplicativa y la ley aditiva interna se evalúan en cada vértice del mismo
estado, conservan levantamientos residuo--cociente distintos y no constituyen
dos cuadrados APP.

\subsection{Obstrucciones de fidelidad del soporte APP: cardinales \(81/17\), traslaciones reversibles y separación de las hojas suma y producto}

\begin{criterion}[Criterios de refutación del soporte APP]
La construcción queda refutada si ocurre cualquiera de los hechos
siguientes:
\begin{enumerate}
 \item el soporte completo tiene cardinal distinto de \(81\);
 \item el complemento gnomónico no tiene cardinal \(17\);
 \item un estado de \(X_9\) carece de alguna de las cuatro traslaciones;
 \item se identifican \(\Sigma\) y \(\Pi\) porque coincidan en un valor;
 \item se elimina \(Q^+\) o \(Q^\times\) al retornar por una frontera;
 \item se presentan las trayectorias de distinta longitud como nuevos
       vértices del soporte;
 \item se interpreta \(9\) y \([0]_9\) como fases independientes.
\end{enumerate}
\end{criterion}

La unidad ha construido APP hasta su evaluación genealógica completa. Todavía
no ha orientado la diferencia entre la matriz y su acumulador, no ha
introducido el TRIT
y no ha definido el Topological Phase Kernel. La unidad siguiente estudiará
la estructura algebraica de filas y columnas, el sector de unidades y el
cociclo de acarreo que prepara esa orientación.


% Unidad U005. Redacción autónoma desde 03c/03h/20 de la autoridad 1063/live.

\long\def\HMTDeferredTPKBase{%
\section{Topological Phase Kernel: base observable y estado holonómico integral}
\label{p12-83-c1-s3:chap:tpk-categoria-estado-enriquecido}

APP ha proporcionado un soporte, una matriz multiplicativa y una ley aditiva
interna; TRIT ha
proporcionado una orientación levantada. El Topological Phase Kernel, en
adelante TPK, especifica cómo se transportan conjuntamente esos datos, qué
coordenadas se conservan al cerrar una ruta y qué lectores pueden publicarse
sin identificar la publicación con el estado. Su objeto matemático es una
categoría sobre una base de caminos, provista de una familia explícita de
fibras y relaciones de prolongación.

\subsection{La base doble de caminos APP}

Sea \(\mathsf P_{\rm APP}\) la categoría libre construida en la
\cref{p12-83-c1-s2:chap:app-completacion-gnomonica}. Definimos
\[
 \mathsf P_{\rm APP}^{(2)}
 :=\mathsf P_{\rm APP}\times\mathsf P_{\rm APP}.
\]
Un objeto es un par de posiciones
\((p^\Sigma,p^\Pi)\in X_9^2\). Un morfismo es un par de trayectorias; una
de ellas puede ser la identidad. La primera componente se evalúa por
\(\Sigma\), la segunda por \(\Pi\). No son dos retículas: son dos cursores
tipados sobre el mismo soporte APP.

El conjunto de direcciones cardinales es
\[
 \mathscr D_{\rm card}:=\{N,E,S,O\},
\]
con desplazamientos
\[
 \delta(N)=(-1,0),\quad
 \delta(E)=(0,1),\quad
 \delta(S)=(1,0),\quad
 \delta(O)=(0,-1)
 \quad\text{en }X_9.
\]
La involución de orientación es
\[
 \overline N=S,
 \quad \overline S=N,
 \quad \overline E=O,
 \quad \overline O=E.
\]

\subsection{Fase nonádica, sector dodecafásico y calendario}

El reloj completo de esta etapa es
\[
 \Theta_{108}:=\mathbb Z/108\mathbb Z.
\]
Para \(\theta\in\Theta_{108}\), sea
\(\widetilde\theta\in\{0,\ldots,107\}\) su representante normalizado.
Definimos tres coordenadas distintas:
\begin{align*}
 g_0(\theta)&=[\widetilde\theta+1]_9\in C_9,\\
 r(\theta)&=1+(\widetilde\theta\bmod9)\in\mathcal D_9,\\
 \beta(\theta)&=
 \left[\left\lfloor\frac{\widetilde\theta}{9}\right\rfloor\right]_{12}
 \in C_{12}.
\end{align*}
La primera es la clase de fase, la segunda su etiqueta positiva y la tercera
el sector dodecafásico. La aplicación \(\beta\) es una aplicación de bloques;
no es la reducción canónica \(\theta\bmod12\).

La firma y la hoja activa son
\[
 \tau(\theta):=\varepsilon_9(r(\theta))\in\mathsf{TRIT},
\]
y
\[
 \mu(\theta)=
 \begin{cases}
  \Sigma,&\tau(\theta)=+1,\\
  \Pi,&\tau(\theta)=-1,\\
  \varnothing,&\tau(\theta)=0.
 \end{cases}
\]
El símbolo \(\varnothing\) designa un paso de umbral en el que ningún cursor
se traslada; no designa ausencia de estado.

\begin{definition}[Configuración observable]
\label{p12-83-c1-s3:def:tpk-configuracion-observable}
La base observable es
\[
 \mathcal Q_{\rm obs}
 :=X_9^2\times\mathscr D_{\rm card}^2\times\Theta_{108}.
\]
Un elemento se escribe
\[
 q=(p^\Sigma,p^\Pi,d^\Sigma,d^\Pi,\theta).
\]
Contiene dos posiciones, dos orientaciones cardinales y un índice de
calendario.
\end{definition}

\subsection{Selector de hoja \(M_{\rm ph}:\mathcal D_9\to\mathsf{TRIT}\) y evolución conjunta de posición, orientación y calendario}

La función
\[
 M_{\rm ph}:\mathcal D_9\longrightarrow\mathsf{TRIT},
 \qquad M_{\rm ph}(r)=\varepsilon_9(r)
\]
es el selector interno de hoja. Selecciona qué evaluación actúa; no desplaza
por sí mismo ninguna posición.

Definimos el indicador de inversión de orientación por
\[
 \eta_{27}(\theta)=
 \begin{cases}
  -1,&\widetilde\theta+1\in\{27,54,81,108\},\\
  +1,&\text{en otro caso}.
 \end{cases}
\]
Las nuevas direcciones son
\[
 d^{\nu,+}=
 \begin{cases}
  \overline{d^\nu},&\eta_{27}(\theta)=-1,\\
  d^\nu,&\eta_{27}(\theta)=+1,
 \end{cases}
 \qquad \nu\in\{\Sigma,\Pi\}.
\]
Las posiciones evolucionan mediante
\[
 (p^{\Sigma,+},p^{\Pi,+})=
 \begin{cases}
 (p^\Sigma+\delta(d^\Sigma),p^\Pi),&\tau(\theta)=+1,\\
 (p^\Sigma,p^\Pi+\delta(d^\Pi)),&\tau(\theta)=-1,\\
 (p^\Sigma,p^\Pi),&\tau(\theta)=0.
 \end{cases}
\]

\begin{definition}[Evolución observable]
\label{p12-83-c1-s3:def:tpk-evolucion-observable}
La aplicación
\[
 F_{\rm obs}:\mathcal Q_{\rm obs}\longrightarrow\mathcal Q_{\rm obs}
\]
se define por
\[
 F_{\rm obs}(q)=
 (p^{\Sigma,+},p^{\Pi,+},d^{\Sigma,+},d^{\Pi,+},[\theta+1]_{108}).
\]
La categoría libre del grafo \(q\to F_{\rm obs}(q)\) se denota por
\(\mathsf P_{\rm obs}\).
\end{definition}

Cada flecha de \(\mathsf P_{\rm obs}\) induce un par de trayectorias en
APP, de modo que existe un funtor
\[
 B:\mathsf P_{\rm obs}\longrightarrow
 \mathsf P_{\rm APP}^{(2)}.
\]
Si \(\tau=+1\), la segunda componente de \(B\) es identidad; si
\(\tau=-1\), lo es la primera; si \(\tau=0\), ambas lo son.

\begin{proposition}[Separación entre selección, movimiento y observación]
\label{p12-83-c1-s3:prop:tpk-separacion-seleccion-movimiento}
El selector \(M_{\rm ph}\), las traslaciones
\(T_d(p)=p+\delta(d)\) y cualquier muestreo temporal poseen dominios y
codominios distintos. No son operadores intercambiables.
\end{proposition}

\begin{proof}
\(M_{\rm ph}\) toma una etiqueta de fase y devuelve una orientación trítica.
\(T_d\) es un automorfismo de \(X_9\). Un muestreo actúa sobre una
sucesión o una categoría de flechas. Identificar dos de ellos violaría el
tipado. \(F_{\rm obs}\) los compone en un orden declarado sin borrar esa
distinción.
\end{proof}

\subsection{Descomposición del estado holonómico integral en \(6\) familias fibradas}

Sobre cada \(q\in\mathcal Q_{\rm obs}\) fijamos la fibra total
\[
 \mathcal F_q=
 \mathsf O_q\times
 \mathsf H_q\times
 \mathsf C_q\times
 \mathsf S_q\times
 \mathsf B_q\times
 \Lambda_q.
\]
Cada factor tiene un tipo y una función propios.

\subsubsection{Determinación conjunta de la orientación y la hoja del estado holonómico integral}

\[
 \mathsf O_q\subseteq
 \mathscr D_{\rm card}^2\times\{+1,-1\}.
\]
Las dos direcciones coinciden con las almacenadas en \(q\); la última
coordenada registra la orientación de hoja cuando la realización la requiere.
La etiqueta de hoja no sustituye las direcciones cardinales.

\subsubsection{Descomposición euclídea del estado en residuo y cociente compatibles}

La primera fibra de elevación es
\[
 \mathsf H^{(1)}=\mathcal D_9\times\mathbb Z,
 \qquad
 \operatorname{rec}_9(r,u)=r+9u.
\]
Para cada entero \(n\) existen únicos
\[
 r=\operatorname{dr}_9(n),
 \qquad
 u=\frac{n-r}{9}.
\]
Una profundidad mayor consiste en una torre declarada de estas fibras y sus
mapas de truncamiento. El término ``Hensel'' no sustituye la definición de
esa torre.

\subsubsection{Operador de acarreo y actualización conjunta de residuo, cociente y memoria}

\(\mathsf C_q\) es un grupo abeliano. En los lectores de bloques decimales
se toma \(\mathsf C_q=\mathbb Z\). El incremento no es una coordenada
independiente, sino una cocadena sobre las flechas:
\[
 \kappa(r)\in\mathsf C_{q'}
 \quad\text{para }r:q\to q'.
\]

\subsubsection{Construcción de la firma entera a partir de orientación, acarreo y frontera}

La firma se deriva mediante una aplicación tipada
\[
 \operatorname{Sig}_q:
 \mathsf H_q\times\mathsf C_q\times\Lambda_q
 \longrightarrow\mathsf S_q.
\]
En la realización de frontera nonádica se conservan, de manera individual,
\[
 q_\partial,a_\partial,c_\partial\in\mathbb F_3^6,
 \qquad
 \operatorname{colw}\in\mathbb Z_{\geq0}^{6},
 \qquad
 \rho_W,r_\partial\in\mathbb F_3^3.
\]
Estas coordenadas son margen, combinación axial, condición cúbica, pesos de
columna, clase residual de Witt y retorno ternario. No se condensan en una
única palabra si una prueba posterior utiliza alguna de ellas.

\subsubsection{Cilindro y frontera}

\[
 \mathsf B_q:=\mathsf{Cyl}_q\times\mathsf{Fr}_q.
\]
El cilindro localiza una familia de prefijos; la frontera determina qué
extensiones son compatibles. La supervivencia será una relación entre
estados provistos de ambos datos, no un número añadido a posteriori.

\subsubsection{Registro genealógico de transporte}

\(\Lambda_q\) contiene los registros de transporte cuyo extremo es
\(q\). Una prolongación \(r:q\to q'\) actúa por concatenación:
\[
 \lambda\longmapsto r\circ\lambda.
\]
El registro genealógico conserva la ruta; su extremo observable no la determina.

\subsection{Transporte de fibras}

Para cada flecha \(r:q\to q'\) se declaran aplicaciones
\[
 \mathsf H(r):\mathsf H_q\to\mathsf H_{q'},
 \quad
 \mathsf C(r):\mathsf C_q\to\mathsf C_{q'},
 \quad
 \Lambda(r):\Lambda_q\to\Lambda_{q'}.
\]
Respetan identidades y composición, y \(\mathsf C(r)\) es un homomorfismo.
El acarreo satisface
\begin{equation}
 \kappa(r_2\circ r_1)
 =\kappa(r_2)+\mathsf C(r_2)\bigl(\kappa(r_1)\bigr).
\label{p12-83-c1-s3:eq:tpk-cocadena-transporte}
\end{equation}
Por tanto, el transporte de las coordenadas determinadas es
\begin{align*}
 h'&=\mathsf H(r)(h),\\
 c'&=\mathsf C(r)(c)+\kappa(r),\\
 \lambda'&=\Lambda(r)(\lambda),\\
 \sigma'&=\operatorname{Sig}_{q'}(h',c',\lambda').
\end{align*}

\begin{lemma}[Asociatividad del acarreo transportado]
\label{p12-83-c1-s3:lem:tpk-asociatividad-acarreo}
La ecuación \cref{p12-83-c1-s3:eq:tpk-cocadena-transporte} garantiza que transportar un
estado por \(r_1\) y después por \(r_2\) produce el mismo acarreo que
transportarlo por \(r_2\circ r_1\).
\end{lemma}

\begin{proof}
El transporte sucesivo da
\[
 \mathsf C(r_2)\bigl(\mathsf C(r_1)c+\kappa(r_1)\bigr)+\kappa(r_2).
\]
Por functorialidad de \(\mathsf C\), el primer término es
\(\mathsf C(r_2r_1)c\); la suma de los otros dos es
\(\kappa(r_2r_1)\) por \cref{p12-83-c1-s3:eq:tpk-cocadena-transporte}.
\end{proof}

\subsection{Relaciones de extensión}

La orientación y la frontera no se fuerzan mediante una función universal.
Para cada flecha \(r:q\to q'\), se declara una relación
\[
 \operatorname{Ext}_r\subseteq\mathcal F_q\times\mathcal F_{q'}.
\]
Se exigen las dos condiciones
\[
 \Delta_q\subseteq\operatorname{Ext}_{\operatorname{id}_q},
 \qquad
 \operatorname{Ext}_{r_2}\circ\operatorname{Ext}_{r_1}
 \subseteq\operatorname{Ext}_{r_2\circ r_1}.
\]
La primera conserva identidades; la segunda conserva composición. Una
relación puede admitir varias extensiones locales y dejar una sola tras un
horizonte mayor.

\begin{definition}[Topological Phase Kernel]
\label{p12-83-c1-s3:def:tpk-categoria-autonoma}
Un objeto del TPK es una tupla
\[
 x=(q,o,h,c,\sigma,C,b,\lambda)
\]
tal que
\[
 q\in\mathcal Q_{\rm obs},\quad
 o\in\mathsf O_q,\quad
 h\in\mathsf H_q,\quad
 c\in\mathsf C_q,\quad
 \lambda\in\Lambda_q,
\]
\[
 \sigma=\operatorname{Sig}_q(h,c,\lambda),
 \qquad
 (C,b)\in\mathsf{Cyl}_q\times\mathsf{Fr}_q.
\]
Asociamos a cada objeto su componente fibral
\[
 \mathbf f(x):=(o,h,c,\sigma,(C,b),\lambda)\in\mathcal F_q.
\]
Un morfismo \(x\to x'\) es una terna \((r,x,x')\) en la que
\(r:q\to q'\) pertenece a \(\mathsf P_{\rm obs}\), las coordenadas
transportables satisfacen las leyes anteriores y
\((\mathbf f(x),\mathbf f(x'))\in
\operatorname{Ext}_r\). La composición es
\[
 (r_2,x',x'')\circ(r_1,x,x')=(r_2r_1,x,x''),
\]
y la identidad de \(x\) es \((\operatorname{id}_q,x,x)\).
\end{definition}

\begin{theorem}[Estructura categorial del TPK]
\label{p12-83-c1-s3:thm:tpk-es-categoria}
La \cref{p12-83-c1-s3:def:tpk-categoria-autonoma} define una categoría y la aplicación
\[
 p:\mathsf{TPK}\longrightarrow\mathsf P_{\rm obs},
 \qquad p(x)=q,\qquad p(r,x,x')=r,
\]
es un funtor.
\end{theorem}

\begin{proof}
La composición de rutas es asociativa. El
\cref{p12-83-c1-s3:lem:tpk-asociatividad-acarreo} prueba la compatibilidad de la
coordenada de acarreo; la functorialidad declarada de \(\mathsf H\),
\(\mathsf C\) y \(\Lambda\) prueba la de las demás coordenadas
determinadas. La estabilidad de \(\operatorname{Ext}\) garantiza que el
compuesto sigue siendo admisible, y su condición diagonal proporciona las
identidades. La aplicación \(p\) conserva identidades y composición por
construcción.
\end{proof}

\subsection{Proyecciones y pérdidas tipadas}

Las proyecciones elementales son
\[
 \phi(x)=g_0(\theta),
 \qquad
 \tau(x)=\varepsilon_9(r(\theta)),
 \qquad
 \mu(x)=\mu(\theta).
\]
La proyección al estado TRIT local es
\[
 \pi_{\rm TRIT}(x)
 =\bigl(\tau,\mu,o,h,c,\sigma,\lambda,\phi\bigr).
\]
Olvida el cilindro, la frontera y parte de \(q\); no define otro contenedor.
Dos estados pueden tener la misma imagen por \(\pi_{\rm TRIT}\) y diferir
en la información olvidada.

\subsection{Criterio de supervivencia de rutas mediante conservación de la frontera proyectiva}

Sea \(\mathcal X_{\rm TPK}^{(n)}\) la familia de estados a profundidad
\(n\). Para \(N\geq n\), definimos
\[
 \operatorname{Surv}_n^{(N)}=
 \left\{x_n:
 \begin{array}{l}
 \text{existen }x_{n+1},\ldots,x_N\text{ con}\\
 (x_j,x_{j+1})\in\operatorname{Ext}_{r_j}
 \text{ para }n\leq j<N
 \end{array}\right\}.
\]

\begin{proposition}[Filtración de supervivencia]
\label{p12-83-c1-s3:prop:tpk-filtracion-supervivencia}
Para \(N\geq n\),
\[
 \operatorname{Surv}_n^{(N+1)}
 \subseteq
 \operatorname{Surv}_n^{(N)}.
\]
Si el esquema está definido a toda profundidad, los estados
indefinidamente prolongables son
\[
 \operatorname{Surv}_n
 =\bigcap_{N\geq n}\operatorname{Surv}_n^{(N)}.
\]
\end{proposition}

\begin{proof}
Una prolongación hasta \(N+1\) se trunca a una prolongación hasta \(N\),
lo que prueba la inclusión. La intersección expresa exactamente la
prolongabilidad para todo horizonte finito.
\end{proof}

\subsection{Árbol operatorio del contenedor}

\begin{figure}[htbp]
 \centering
 \begin{tikzpicture}[
  node distance=7mm and 9mm,
  every node/.style={font=\small,align=center},
  root/.style={draw=black,very thick,rounded corners,inner sep=5pt},
  op/.style={draw=black,rounded corners,inner sep=4pt},
  edge/.style={-{Latex[length=2mm]},semithick}
 ]
  \node[root] (tpk) {TPK\\$\mathcal X_{\rm TPK}^{\rm enr}$};
  \node[op,below left=of tpk] (sel) {selección\\$M_{\rm ph}$};
  \node[op,below=of tpk] (mov) {transporte\\$T_d,F_{\rm obs}$};
  \node[op,below right=of tpk] (fib) {fibras\\$\mathsf H,\mathsf C,\mathsf S,
    \mathsf B,\Lambda$};
  \node[op,below=17mm of mov] (cal) {calendario\\$C_{12}\hookrightarrow
    C_{108}\twoheadrightarrow C_9$};
  \node[op,below=of cal] (ext) {extensión, cilindros y supervivencia};
  \draw[edge] (tpk) -- (sel);
  \draw[edge] (tpk) -- (mov);
  \draw[edge] (tpk) -- (fib);
  \draw[edge] (sel) |- (cal);
  \draw[edge] (mov) -- (cal);
  \draw[edge] (fib) |- (cal);
  \draw[edge] (cal) -- (ext);
 \end{tikzpicture}
 \caption{Contención operatoria del Topological Phase Kernel. El selector,
 el transporte, los lectores y las fibras no son contenedores paralelos: son
 mapas o coordenadas del mismo estado holonómico integral.}
 \label{p12-83-c1-s3:fig:tpk-arbol-operatorio-autonomo}
\end{figure}

\subsection{Obstrucciones del estado holonómico integral del TPK: tipado del selector, cocadena de acarreo, extensión compatible y conservación genealógica}

\begin{criterion}[Criterios de refutación del estado holonómico integral]
La definición queda refutada si:
\begin{enumerate}
 \item el selector \(M_{\rm ph}\) se usa como traslación;
 \item se identifica \(\theta\), \(g_0(\theta)\), \(r(\theta)\) y
       \(\beta(\theta)\);
 \item se permite que una firma se elija independientemente de
       \((h,c,\lambda)\);
 \item la composición de acarreos incumple
       \cref{p12-83-c1-s3:eq:tpk-cocadena-transporte};
 \item se sustituye \(\operatorname{Ext}_r\) por una elección funcional no
       demostrada;
 \item se llama estado completo a una proyección que ha olvidado cilindro,
       frontera o registro genealógico;
 \item el retorno de fase borra alguna de las coordenadas de fibra.
\end{enumerate}
\end{criterion}

Esta unidad ha definido el contenedor. La siguiente enumerará sus lectores,
operadores, odómetro y calendarios, y demostrará por qué los retornos de
nueve, veintisiete y ciento ocho pasos no son reinicios del estado.
}%


% Unidad U001. Redacción autónoma; no procede del borrador condensado.

\section{Segmentos marcados, intersticios y célula nonádica}
\label{p12-83-c1-s4:chap:segmentos-marcas-intersticios}

La construcción comienza antes de toda constante y antes de toda elección de
unidad física. El dato inaugural no es un número real ya constituido, sino un
segmento ordenado, una familia de marcas y la relación de adyacencia entre
marcas consecutivas. Esta separación es necesaria porque una marca, un
intervalo, la longitud de ese intervalo y el numeral que publica dicha
longitud son objetos de tipos diferentes.

\subsection{Carta posicional finita y proyección residual decimal}

\begin{definition}[Segmento marcado de orden diez]
Sea
\[
 \mathsf M_{10}:=\{m_0,m_1,\ldots,m_{10}\}
\]
un conjunto totalmente ordenado por
\(m_0<m_1<\cdots<m_{10}\). Su grafo de adyacencia es el camino
\[
 \mathsf P_{10}:\quad
 m_0\longrightarrow m_1\longrightarrow\cdots\longrightarrow m_{10}.
\]
Denotaremos por
\[
 \mathsf I_k:=(m_k,m_{k+1}),\qquad 0\leq k\leq9,
\]
el intersticio abstracto comprendido entre dos marcas consecutivas. La
familia de intersticios es
\(\mathsf E_{10}:=\{\mathsf I_0,\ldots,\mathsf I_9\}\).
\end{definition}

La notación abierta \((m_k,m_{k+1})\) no presupone aún la recta real. Sólo
indica el par ordenado de extremos y la arista que los une. Cuando se elige la
carta arquimediana
\[
 \chi_{10}(m_k)=k,
\]
el intersticio abstracto \(\mathsf I_k\) se publica mediante el intervalo
semiabierto \([k,k+1)\). Por tanto, el índice \(k\), la arista
\(\mathsf I_k\) y el intervalo \([k,k+1)\) no se identificarán en el texto:
el primero es una etiqueta, el segundo es un elemento combinatorio y el
tercero es su realización en una carta.

\begin{lemma}[Censo de marcas e intersticios]
\label{p12-83-c1-s4:lem:censo-marcas-intersticios}
El segmento \(\mathsf P_{10}\) contiene once marcas y diez intersticios. Más
generalmente, un camino con \(n+1\) marcas posee exactamente \(n\) aristas.
\end{lemma}

\begin{proof}
La aplicación
\[
 \{0,\ldots,n-1\}\longrightarrow E(\mathsf P_n),
 \qquad k\longmapsto(m_k,m_{k+1}),
\]
es biyectiva. La afirmación particular se obtiene con \(n=10\).
\end{proof}

El extremo \(m_0\) fija el origen de la carta y \(m_{10}\) fija su umbral.
Ninguno de los dos constituye por sí mismo un intersticio adicional. Esta
observación elemental impide tres confusiones posteriores: contar extremos
como celdas, confundir el cierre con un nuevo estado de fase y tratar el cero
como una décima fase independiente del ciclo nonádico.

\subsection{Extracción de la célula de \(9\) fases}

La carta decimal distingue el anclaje \([0,1)\) de la célula periódica. Las
nueve marcas positivas
\[
 \mathsf D_9:=\{1,2,3,4,5,6,7,8,9\}
\]
se utilizarán como representantes públicos de
\(\mathbb Z/9\mathbb Z\), haciendo corresponder la marca \(9\) con la clase
residual \(\overline 0\). Definimos
\[
 \lambda_9:\mathsf D_9\longrightarrow\mathbb Z/9\mathbb Z,
 \qquad \lambda_9(a)=\overline a.
\]

\begin{proposition}[Carta positiva de las clases nonádicas]
\label{p12-83-c1-s4:prop:carta-positiva-nonadica}
La aplicación \(\lambda_9\) es biyectiva. Su inversa es la sección positiva
\[
 s_+(\overline a)=
 \begin{cases}
 a,&1\leq a\leq8,\\
 9,&\overline a=\overline0.
 \end{cases}
\]
En particular, la escritura pública \(9\) y la escritura residual \(0\)
representan la misma clase, pero pertenecen a cartas distintas.
\end{proposition}

\begin{proof}
Las nueve clases de \(\mathbb Z/9\mathbb Z\) poseen representantes únicos en
\(\{0,1,\ldots,8\}\). Sustituir el representante \(0\) por \(9\) produce
otra transversal completa, por lo que \(\lambda_9\) es biyectiva y \(s_+\)
es su inversa.
\end{proof}

La célula de fase no se obtiene borrando arbitrariamente el origen. El tramo
\([0,1)\) fija la carta decimal desde la que se observa el ciclo; las nueve
fases son las nueve clases que se publican como \(1,\ldots,9\). En la carta
levantada, el paso que cierra la célula se representa por \([9,10)\); en la
carta residual, ese mismo paso retorna de \(\overline0\) a \(\overline1\).
El cierre y el cambio de representante no son la misma operación.

\subsection{Posición, altura y división euclídea}

\begin{definition}[Coordenadas nonádicas]
Para cada entero \(n\geq1\), definimos
\[
 \rho_9(n):=1+((n-1)\bmod9),
 \qquad
 q_9(n):=\left\lfloor\frac{n-1}{9}\right\rfloor.
\]
La primera coordenada se denomina posición de fase; la segunda, altura de
vuelta o cociente de memoria.
\end{definition}

\begin{theorem}[Descomposición nonádica única]
\label{p12-83-c1-s4:thm:descomposicion-nonadica-unica}
Todo entero \(n\geq1\) admite una única escritura
\[
 n=9q+\rho,
 \qquad q\in\mathbb Z_{\geq0},\quad \rho\in\mathsf D_9.
\]
En esa escritura \(q=q_9(n)\) y \(\rho=\rho_9(n)\).
\end{theorem}

\begin{proof}
Aplicamos la división euclídea a \(n-1\): existen únicos
\(q\geq0\) y \(r\in\{0,\ldots,8\}\) tales que
\(n-1=9q+r\). Tomando \(\rho=r+1\) se obtiene
\(n=9q+\rho\), con \(1\leq\rho\leq9\). La unicidad de \(q,r\) implica la
de \(q,\rho\).
\end{proof}

\begin{corollary}[Retorno de fase con incremento de memoria]
\label{p12-83-c1-s4:cor:retorno-fase-incremento-memoria}
Para todo \(n\geq1\),
\[
 \rho_9(n+9)=\rho_9(n),
 \qquad
 q_9(n+9)=q_9(n)+1.
\]
Por consiguiente, una vuelta completa conserva la fase visible y cambia el
estado levantado.
\end{corollary}

\begin{proof}
Si \(n=9q+\rho\), entonces \(n+9=9(q+1)+\rho\); la unicidad del
\cref{p12-83-c1-s4:thm:descomposicion-nonadica-unica} da las dos identidades.
\end{proof}

La evolución por un paso queda escrita sin ambigüedad como
\[
 (q,\rho)\longmapsto
 \begin{cases}
 (q,\rho+1),&1\leq\rho<9,\\
 (q+1,1),&\rho=9.
 \end{cases}
\]
La segunda línea no es un reinicio. Recupera la etiqueta inicial de fase y
aumenta el cociente que registra cuántas vueltas se han realizado.

\subsection{Periodicidad de la proyección circular e inyectividad del levantamiento helicoidal nonádico}

Sea
\[
 \vartheta_9(\rho)=\frac{2\pi(\rho-1)}9.
\]
La proyección de fase sobre el círculo unitario es
\[
 \mathcal R_9(n)=
 \bigl(\cos\vartheta_9(\rho_9(n)),
       \sin\vartheta_9(\rho_9(n))\bigr).
\]
Para conservar simultáneamente la posición y la altura introducimos la
inmersión helicoidal
\[
 \mathcal H_9(n)=
 \left(
  \cos\vartheta_9(\rho_9(n)),
  \sin\vartheta_9(\rho_9(n)),
  q_9(n)+\frac{\rho_9(n)-1}{9}
 \right).
\]

\begin{proposition}[La proyección circular no determina el estado]
\label{p12-83-c1-s4:prop:rueda-no-determina-estado}
La aplicación \(\mathcal R_9\) es periódica de periodo nueve. La aplicación
\(\mathcal H_9\) es inyectiva. En particular,
\[
 \mathcal R_9(n+9)=\mathcal R_9(n),
 \qquad
 \mathcal H_9(n+9)\neq\mathcal H_9(n).
\]
\end{proposition}

\begin{proof}
La primera igualdad se sigue del
\cref{p12-83-c1-s4:cor:retorno-fase-incremento-memoria}. La tercera coordenada de
\(\mathcal H_9(n+9)\) supera en una unidad la de \(\mathcal H_9(n)\), por
lo que ambos puntos son distintos. Finalmente, de la tercera coordenada se
recupera \(n\) multiplicando por nueve y añadiendo uno a la parte de fase;
por tanto, \(\mathcal H_9\) es inyectiva.
\end{proof}

Esta proposición es el modelo mínimo del principio que recorrerá toda la
obra:
\[
 \boxed{\text{retorno de fase}\neq\text{retorno del estado}.}
\]
La conexión nonádica del Topological Phase Kernel conservará más campos que
la hélice elemental ---hoja, orientación, acarreo, frontera, terminal y
memoria de ruta---, pero no alterará esta distinción inaugural.

\subsection{Correspondencia entre cardinalidad discreta y medida proyectiva}

\begin{definition}[Conteo]
El conteo de un conjunto finito \(X\) es su cardinal \(|X|\). No requiere
una orientación, una escala ni un lector adicional.
\end{definition}

\begin{definition}[Dato de medida]
Un dato de medida discreta es una terna
\[
 (x,\mathcal L,\mathcal C),
\]
donde \(x\) es un estado, \(\mathcal L\) es un lector tipado y
\(\mathcal C\) es una carta de publicación compatible con el codominio de
\(\mathcal L\). Su salida visible es
\[
 \mathcal C(\mathcal L(x)).
\]
\end{definition}

El mismo conjunto puede tener un único cardinal y admitir múltiples medidas,
porque pueden variar la cocadena local, la orientación de integración o la
carta de publicación. Inversamente, dos estados distintos pueden publicar la
misma coordenada. Por ello la obra conservará siempre dos niveles:
\[
 \text{estado con genealogía}
 \quad\longrightarrow\quad
 \text{coordenada publicada}.
\]

\begin{criterion}[Criterios de refutación de la unidad inaugural]
La construcción queda refutada si ocurre cualquiera de los siguientes
hechos:
\begin{enumerate}
 \item se cuentan diez fases independientes en la célula residual;
 \item se identifica la marca \(9\) con el entero \(0\) sin declarar el
 cambio de carta;
 \item el paso \(9\to1\) fuerza \(q\mapsto0\);
 \item dos enteros distintos poseen la misma pareja \((q_9,\rho_9)\);
 \item se atribuye una unidad física a la célula antes de definir el
 transductor dimensional.
\end{enumerate}
\end{criterion}

\subsection{Dominio inaugural de la célula nonádica y exclusión causal de APP, TRIT, TPK y constantes analíticas}

La unidad construida hasta aquí contiene únicamente:
\[
 \mathsf P_{10},\quad
 \mathsf D_9,\quad
 \lambda_9,\quad
 (q_9,\rho_9),\quad
 \mathcal R_9,\quad
 \mathcal H_9.
\]
No contiene todavía APP, TRIT, el Topological Phase Kernel ni una constante
analítica. La siguiente unidad tomará dos copias de la célula nonádica,
construirá su completación gnomónica y demostrará cómo aparece el soporte
bidimensional sobre el que actúan la matriz multiplicativa de APP y su ley
aditiva interna.


% Unidad U003B. Esqueleto ciclotómico de APP y completación binomial.
% Redacción autónoma a partir de los objetos, no de la arquitectura del PDF abreviado.

\section{Esqueleto ciclotómico de APP y completación binomial de la unidad}
\label{p12-83-c1-s5:chap:app-esqueleto-completacion}

El defecto entre la matriz APP y su acumulador interno no es un número
aislado. Su ley en
longitud cartesiana, su descomposición en órbitas de orden tres y la
completación de la celda nonádica producen tres objetos diferentes que deben
mantenerse tipados:
\[
 \text{defecto escalar}\quad D_d,
 \qquad
 \text{módulo ciclotómico}\quad W_{\rm APP},
 \qquad
 \text{completación}\quad\overline{\mathcal C}_9.
\]
El primero cuenta; el segundo conserva orientación y rango; el tercero añade
estratos de frontera. Las igualdades entre sus cardinales no autorizan a
identificar sus elementos.

\subsection{Ley del defecto a módulo fijo y longitud variable}

Para cada entero \(d\geq1\), sea
\[
 B_d:=\{1,\ldots,9\}^d.
\]
La raíz digital nonádica se representa por
\[
 \operatorname{dr}_9(n)=
 \begin{cases}
 n\bmod9,&n\not\equiv0\pmod9,\\
 9,&n\equiv0\pmod9.
 \end{cases}
\]
Definimos la acumulación aditiva y la evaluación matricial multiplicativa
sobre el mismo dominio:
\[
 A_d(x_1,\ldots,x_d)
 =\operatorname{dr}_9\!\left(\sum_{j=1}^d x_j\right),
\]
\[
 M_d(x_1,\ldots,x_d)
 =\operatorname{dr}_9\!\left(\prod_{j=1}^d x_j\right).
\]
Sus masas visibles y su defecto son
\[
 S_\Sigma(d):=\sum_{x\in B_d}A_d(x),
 \qquad
 S_\Pi(d):=\sum_{x\in B_d}M_d(x),
\]
\[
 D_d:=S_\Pi(d)-S_\Sigma(d).
\]
Aquí \(d\) cuenta posiciones cartesianas. No varía el módulo \(9\) y no se
ha introducido todavía una dimensión física.

\begin{lemma}[Masa del acumulador aditivo]
\label{p12-83-c1-s5:lem:masa-hoja-aditiva-longitud}
Para todo \(d\geq1\),
\[
 S_\Sigma(d)=45\,9^{d-1}.
\]
\end{lemma}

\begin{proof}
Fijadas las primeras \(d-1\) coordenadas y un residuo publicado
\(r\in\{1,\ldots,9\}\), existe una única última coordenada que hace que la
suma tenga residuo \(r\). Cada valor visible posee así \(9^{d-1}\)
preimágenes. Sumando los pesos \(1+\cdots+9=45\) se obtiene la fórmula.
\end{proof}

\begin{lemma}[Censo de la matriz multiplicativa]
\label{p12-83-c1-s5:lem:censo-hoja-multiplicativa-longitud}
Para cada uno de los seis residuos unitarios hay \(6^{d-1}\) preimágenes
por valor de producto fijado. Los residuos \(3\) y \(6\) tienen
\(d6^{d-1}\) preimágenes cada uno. El residuo nulo, publicado como \(9\),
tiene
\[
 9^d-6^d-2d6^{d-1}
\]
preimágenes. En consecuencia,
\[
 S_\Pi(d)=9^{d+1}-9(d+3)6^{d-1}.
\]
\end{lemma}

\begin{proof}
El grupo de unidades de \(\mathbb Z/9\mathbb Z\) tiene seis elementos. Al
fijar \(d-1\) unidades, la última queda determinada por el producto
unitario prescrito, lo que da \(6^{d-1}\) soluciones para cada una de las
seis clases.

Para obtener producto de valoración ternaria exactamente uno, se elige cuál
de las \(d\) coordenadas pertenece a una de las dos clases \(3,6\), mientras
las demás son unidades. Fijado el producto \(3\) o \(6\), la última unidad
queda determinada; de ahí \(d6^{d-1}\) para cada clase. Todas las tuplas
restantes tienen producto nulo módulo \(9\). Restar de \(9^d\) los dos
censos anteriores da su multiplicidad. La suma ponderada por los
representantes visibles produce la última identidad.
\end{proof}

\begin{theorem}[Ley dimensional interna de APP]
\label{p12-83-c1-s5:thm:ley-dimensional-app-autonoma}
Para todo \(d\geq1\),
\[
 \boxed{D_d=4\,9^d-9(d+3)6^{d-1}.}
\]
En particular,
\[
 D_1=0,
 \qquad
 D_2=54,
\]
y
\[
 \sum_{d\geq1}D_dz^d
 =\frac{54z^2(1-3z)}{(1-9z)(1-6z)^2}.
\]
Además,
\[
 \frac{D_d}{9^d}
 =4-(d+3)\left(\frac23\right)^{d-1}
 \longrightarrow4.
\]
\end{theorem}

\begin{proof}
Se restan las fórmulas de los dos lemas. Para la función generatriz se usan
\[
 \sum_{d\geq1}9^dz^d=\frac{9z}{1-9z},
 \qquad
 \sum_{d\geq1}(d+3)6^{d-1}z^d
 =\frac{z(4-18z)}{(1-6z)^2},
\]
y se reduce el numerador común. La fórmula normalizada resulta de dividir
por \(9^d\); el término exponencial multiplicado por \(d+3\) converge a
cero.
\end{proof}

La igualdad \(D_2=54\) reproduce el censo \(459-405\), pero la fórmula
anterior y la ley de módulo variable
\(\Delta_{3^m}^{\rm pl}=m3^{2m-1}\) son enunciados distintos: la primera
varía \(d\) manteniendo el módulo \(9\); la segunda varía el módulo.

\subsection{Operador ciclotómico primario}

Sea
\[
 G_{A_2}=\begin{pmatrix}2&-1\\-1&2\end{pmatrix},
 \qquad
 C=\begin{pmatrix}0&-1\\1&-1\end{pmatrix}.
\]

\begin{proposition}[Coxeter de tipo \(A_2\)]
\label{p12-83-c1-s5:prop:coxeter-a2-autonomo}
Se verifican
\[
 C^3=I_2,
 \qquad
 C^{\mathsf T}G_{A_2}C=G_{A_2},
 \qquad
 \det(XI_2-C)=X^2+X+1.
\]
Por tanto \(C\) preserva la forma \(A_2\), tiene orden exactamente tres y
no posee vector fijo no nulo.
\end{proposition}

\begin{proof}
La multiplicación da
\[
 C^2=\begin{pmatrix}-1&1\\-1&0\end{pmatrix},
 \qquad C^3=I_2.
\]
La sustitución directa en \(C^{\mathsf T}G_{A_2}C\) devuelve
\(G_{A_2}\). Finalmente, \(1\) no es raíz de \(X^2+X+1\), luego
\(\ker(C-I_2)=0\).
\end{proof}

\subsection{Descomposición de APP en \(12\) fibras de tipo \(A_2\)}

Considérese la acción
\[
 a:\mathbb Z/9\mathbb Z\longrightarrow\mathbb Z/9\mathbb Z,
 \qquad a(r)=4r.
\]
Como \(4^3\equiv1\pmod9\), su restricción a las unidades tiene dos órbitas
libres:
\[
 \mathcal O_1=(1,4,7),
 \qquad
 \mathcal O_2=(2,8,5).
\]
Para una órbita \(\mathcal O=(r,4r,7r)\), definimos el módulo de aumento
\[
 A(\mathcal O):=
 \ker\!\left(\mathbb Z[\mathcal O]
 \xrightarrow{\,\sum\,}\mathbb Z\right).
\]
En la base
\[
 b_1=e_r-e_{4r},
 \qquad b_2=e_{4r}-e_{7r},
\]
la forma euclídea restringida tiene Gram \(G_{A_2}\), y el ciclo
\(r\mapsto4r\mapsto7r\mapsto r\) actúa por \(C\).

Para \(x=(x_1,x_2,x_3)\in(\mathbb Z/9\mathbb Z)^3\), los tres pares no
ordenados son
\[
 \{1,2\},\quad\{2,3\},\quad\{3,1\}.
\]
Sobre cada par se evalúan dos leyes tipadas:
\[
 \ell_{ij}^{\Sigma}(x)=x_i+x_j,
 \qquad
 \ell_{ij}^{\Pi}(x)=x_ix_j.
\]
Como
\[
 \ell_{ij}^{\Sigma}(ax)=a\ell_{ij}^{\Sigma}(x),
 \qquad
 \ell_{ij}^{\Pi}(ax)=a^{-1}\ell_{ij}^{\Pi}(x),
\]
la suma porta la orientación \(C\) y el producto la orientación inversa
\(C^{-1}\).

\begin{definition}[Módulo ciclotómico de APP]
\label{p12-83-c1-s5:def:modulo-ciclotomico-app-autonomo}
Se define
\[
 \boxed{
 W_{\rm APP}:=
 \bigoplus_{\{i,j\}}
 \bigoplus_{\mu\in\{\Sigma,\Pi\}}
 \bigoplus_{\mathcal O\in\{\mathcal O_1,\mathcal O_2\}}
 A(\mathcal O).}
\]
Hay \(3\cdot2\cdot2=12\) sumandos, cada uno de rango dos; por tanto
\[
 W_{\rm APP}\cong A_2^{12},
 \qquad \operatorname{rank}W_{\rm APP}=24.
\]
Las etiquetas par--ley--órbita son parte del objeto.
\end{definition}

Para \(u\in\mathbb Z/9\mathbb Z\), escribimos
\[
 \beta_{\mathcal O}(u)=
 \begin{cases}
 (1,0),&u=r,\\
 (0,1),&u=4r,\\
 (-1,-1),&u=7r,\\
 (0,0),&u\notin\mathcal O,
 \end{cases}
\]
en la base \((b_1,b_2)\). El mapa estatal es
\[
 \Psi(x)_{ij,\mu,\mathcal O}
 :=\beta_{\mathcal O}(\ell_{ij}^{\mu}(x)).
\]

\begin{theorem}[Equivarianza y rango del esqueleto estatal]
\label{p12-83-c1-s5:thm:esqueleto-estatal-app-autonomo}
Sea \(\rho(a)\) la acción que usa \(C\) en las seis fibras aditivas y
\(C^{-1}\) en las seis multiplicativas. Entonces
\[
 \boxed{\Psi(ax)=\rho(a)\Psi(x)}
\]
para los \(729\) estados de \((\mathbb Z/9\mathbb Z)^3\). El submódulo
racional generado por las imágenes tiene rango \(24\).
\end{theorem}

\begin{proof}
En la ley aditiva, multiplicar todas las coordenadas por \(4\) avanza una
posición en cada órbita; en la ley multiplicativa multiplica el valor por
\(4^2\equiv7\equiv4^{-1}\pmod9\), y por ello invierte el ciclo. Esto prueba
la equivarianza componente a componente.

Para que el rango no quede oculto en una afirmación informática, fijamos el
orden de coordenadas
\[
 (12,\Sigma,\mathcal O_1),(12,\Sigma,\mathcal O_2),
 (12,\Pi,\mathcal O_1),(12,\Pi,\mathcal O_2),
\]
seguido del mismo orden para los pares \(23\) y \(31\), y dentro de cada
fibra las dos coordenadas \((b_1,b_2)\). Consideremos los estados
\[
\begin{gathered}
001,002,004,005,010,011,012,013,014,015,016,020,\\
022,023,040,050,101,102,104,105,110,120,140,150.
\end{gathered}
\]
La matriz entera \(M_\Psi\) cuyas filas son sus veinticuatro imágenes tiene
entradas en ({-1,0,1}). La eliminación fraccionariamente libre de
Bareiss produce como menores principales sucesivos
\[
1,1,1,-1,-1,-1,-1,1,-1,1,1,1,1,-2,-1,1,1,1,2,1,4,-6,12,9.
\]
En particular,
\[
 \det M_\Psi=9\neq0.
\]
Estas veinticuatro imágenes son linealmente independientes. Como el
codominio ya tiene rango \(24\), generan \(W_{\rm APP}\otimes\mathbb Q\).
\end{proof}

La prueba exhibe un certificado finito reproducible; no deduce rango pleno
de la mera sobreyectividad de cada proyección componente.

\subsection{Localización del agregado \texorpdfstring{\(54\)}{54}}

Sea
\[
 \mathcal M:=\mathbb Z[\mathbb Z/9\mathbb Z]
\]
y, para \(\mu\in\{\Sigma,\Pi\}\),
\[
 \nu_\mu:=\sum_{x,y\in\mathbb Z/9\mathbb Z}
 e_{\ell^\mu(x,y)}.
\]
El censo de las dos tablas da
\[
 \delta_{\Pi\Sigma}:=\nu_\Pi-\nu_\Sigma
 =12e_0+3e_3+3e_6
 -3\sum_{u\in(\mathbb Z/9\mathbb Z)^\times}e_u.
\]
Sus coeficientes son constantes en las órbitas de \(a\); por tanto
\[
 (I-a)\delta_{\Pi\Sigma}=0.
\]
Definimos la proyección de aumento
\[
 T:\mathcal M\longrightarrow
 A(\mathcal O_1)\oplus A(\mathcal O_2),
\]
\[
 T(\nu)=
 \left(((I-a)\nu)|_{\mathcal O_1},
       ((I-a)\nu)|_{\mathcal O_2}\right).
\]
Entonces
\[
 T(\delta_{\Pi\Sigma})=0.
\]
Sin embargo, el lector visible
\[
 w(e_0)=9,
 \qquad w(e_r)=r\quad(1\leq r\leq8)
\]
satisface
\[
 w(\delta_{\Pi\Sigma})=54.
\]

\begin{corollary}[Obstrucción isotípica]
\label{p12-83-c1-s5:cor:obstruccion-isotipica-app-autonoma}
El escalar \(54\) pertenece al sector \(C_3\)-invariante y no puede
recuperarse mediante un funcional lineal que factorice exclusivamente por
los módulos de aumento. En particular,
\[
 \sum_x\Psi(x)=0
\]
y, para todo funcional lineal \(L\) sobre \(W_{\rm APP}\otimes\mathbb Q\),
\[
 L\!\left(\sum_x\Psi(x)\right)=0\neq54.
\]
\end{corollary}

Esta obstrucción impide presentar \(54\) como una suma lineal retrospectiva
de las doce fibras. El acoplamiento APP--Fock posterior conserva la
orientación y utiliza un objeto adicional; no contradice el corolario.

\subsection{Completación cúbica de la celda nonádica}

Sea \(E_9\) el conjunto subyacente de \(\mathbb Z/9\mathbb Z\), y sea
\(\star\notin E_9\) un estado de frontera. Definimos
\[
 \mathcal C_9:=E_9^3,
 \qquad
 \overline{\mathcal C}_9:=(E_9\sqcup\{\star\})^3.
\]
El símbolo \(\star\) no es la clase cero: indica que una coordenada ha
alcanzado el estrato de cierre.

\begin{theorem}[Estratificación binomial de la unidad cúbica]
\label{p12-83-c1-s5:thm:emrh-binomial-autonoma}
El estrato con exactamente \(r\) coordenadas iguales a \(\star\) contiene
\[
 \binom3r9^{3-r}
\]
elementos. Por tanto,
\[
 \boxed{1000=729+243+27+1=729+271.}
\]
Al retirar sólo el ápice \((\star,\star,\star)\), queda
\[
 999=729+243+27=729+270.
\]
\end{theorem}

\begin{proof}
Se eligen las \(r\) coordenadas de frontera de \(\binom3r\) maneras; cada
coordenada restante admite nueve valores. Los cuatro estratos son disjuntos
y agotan el producto cartesiano. La suma es la expansión de
\((9+1)^3\).
\end{proof}

\begin{figure}[htbp]
\centering
\resizebox{0.96\linewidth}{!}{\begin{tikzpicture}[font=\small]
  \begin{scope}[x={(.95cm,0cm)},y={(.45cm,.33cm)},z={(0cm,.95cm)}]
  \coordinate (O) at (0,0,0);
  \coordinate (X) at (3.4,0,0);
  \coordinate (Y) at (0,3.4,0);
  \coordinate (XY) at (3.4,3.4,0);
  \coordinate (Z) at (0,0,3.4);
  \coordinate (XZ) at (3.4,0,3.4);
  \coordinate (YZ) at (0,3.4,3.4);
  \coordinate (XYZ) at (3.4,3.4,3.4);
  \draw[hmtgray,dashed,semithick] (O)--(Y)--(XY);
  \draw[hmtgray,dashed,semithick] (Y)--(YZ);
  \path[fill=hmtlightblue,fill opacity=.78]
    (O)--(X)--(XZ)--(Z)--cycle;
  \path[fill=hmtlightgold,fill opacity=.62]
    (X)--(XY)--(XYZ)--(XZ)--cycle;
  \path[fill=hmtlightgreen,fill opacity=.72]
    (Z)--(XZ)--(XYZ)--(YZ)--cycle;
  \draw[hmtblue,thick] (O)--(X)--(XZ)--(Z)--cycle;
  \draw[hmtgold,thick] (X)--(XY)--(XYZ)--(XZ);
  \draw[hmtgreen,thick] (Z)--(YZ)--(XYZ);
  \node[font=\bfseries\Large,hmtgreen,fill=white,fill opacity=.78,
    text opacity=1,inner sep=2pt] at (1.7,.55,1.72) {\(9^3=729\)};
  \node[font=\scriptsize,hmtgray,fill=white,fill opacity=.78,
    text opacity=1,inner sep=1.5pt] at (1.7,.55,1.22)
    {interior nonádico};
  \draw[-{Latex[length=2mm]},hmtblue,semithick]
    (O) -- (1.18,0,0)
    node[pos=.78,below=5pt,font=\scriptsize] {\(x=(x_0,x_1)\)};
  \draw[-{Latex[length=2mm]},hmtgreen,semithick]
    (O) -- (0,1.18,0)
    node[pos=1,right=5pt,font=\scriptsize] {\(y=(y_0,y_1)\)};
  \draw[-{Latex[length=2mm]},hmtblue,semithick]
    (O) -- (0,0,1.18)
    node[pos=.78,left=2pt,font=\scriptsize] {\(z=(z_0,z_1)\)};
  \end{scope}
  \node[align=center,text width=5.2cm,font=\scriptsize,hmtgray,
    anchor=north] at (2.55,-.72)
    {coordenadas ternarias pareadas en tres direcciones independientes};
  \draw[decorate,decoration={brace,amplitude=5pt},hmtgold,thick]
    (5.20,.65)--(5.20,4.80);
  \node[anchor=west,align=left,text width=1.25cm,font=\scriptsize,hmtgold]
    at (5.40,3.82) {de lado \(9\) a \(10\)\\corona \(271\)};
  \begin{scope}[xshift=7cm,yshift=.45cm]
    \node[anchor=west,font=\bfseries\large,hmtblue] at (0,4.05)
      {Completación cúbica};
    \draw[fill=hmtlightgreen,draw=hmtgreen] (0,3.15) rectangle (6.1,3.75);
    \node at (3.05,3.45) {interior: \(729\)};
    \draw[fill=hmtlightblue,draw=hmtblue] (0,2.35) rectangle (2.9,2.95);
    \node[align=center,text width=2.75cm,font=\scriptsize] at (1.45,2.65)
      {caras nuevas\\\(243=3\cdot9^2\)};
    \draw[fill=hmtlightgold,draw=hmtgold] (2.9,2.35) rectangle (4.9,2.95);
    \node[align=center,text width=1.8cm,font=\scriptsize] at (3.9,2.65)
      {aristas\\\(27=3\cdot9\)};
    \draw[fill=hmtlightred,draw=hmtred] (4.9,2.35) rectangle (6.1,2.95);
    \node[align=center,font=\scriptsize] at (5.5,2.65)
      {vértice\\\(1\)};
    \node[font=\bfseries\Large,hmtgold] at (3.05,1.65)
      {\(1000=729+243+27+1=729+271\)};
    \node[align=center,hmtgray] at (3.05,.65)
      {frontera interna del cubo de lado \(9\): \(386\) puntos\\
       corona de la completación \(9^3\subset10^3\): \(271\) estados};
  \end{scope}
\end{tikzpicture}
 }
\caption{Completación cúbica de la celda nonádica. El interior de
cardinal \(729\) y los estratos de frontera de cardinales \(243\), \(27\) y
\(1\) son partes disjuntas de un único objeto.}
\label{p12-83-c1-s5:fig:cubo-emrh-autonomo}
\end{figure}

Esta estratificación recibe en HMT el nombre de extrema y media razón
holográfica. Matemáticamente es la completación binomial tipada de la celda
nonádica; el nombre no sustituye su definición.

Para todo \(d\geq1\), sea
\[
 \overline{\mathcal C}^{(d)}_9=(E_9\sqcup\{\star\})^d.
\]
Asignemos peso \(1/10\) a cada símbolo de una coordenada y tomemos la medida
producto \(\mu_d\).

\begin{theorem}[Compatibilidad proyectiva de la medida]
\label{p12-83-c1-s5:thm:medida-proyectiva-completacion}
La proyección que olvida la última coordenada satisface
\[
 (\pi_d)_*\mu_d=\mu_{d-1}.
\]
Además,
\[
 10^d=\sum_{r=0}^{d}\binom dr9^{d-r},
 \qquad
 \mu_d(\overline{\mathcal C}^{(d)}_9)=1.
\]
El número de estados aumenta con \(d\), mientras la masa total normalizada
se conserva.
\end{theorem}

\begin{proof}
Para todo cilindro \(A\times(E_9\sqcup\{\star\})\),
\[
 \mu_d(A\times(E_9\sqcup\{\star\}))
 =\mu_{d-1}(A)\sum_{u\in E_9\sqcup\{\star\}}\frac1{10}
 =\mu_{d-1}(A).
\]
Los cilindros generan el álgebra finita de subconjuntos, lo que determina la
medida imagen. La identidad cardinal vuelve a ser el binomio de Newton.
\end{proof}

\subsection{Morfismo del bloque central de APP hacia su realización euclídea}

El bloque central ya construido es
\[
 B_c=\begin{pmatrix}7&2\\2&7\end{pmatrix}.
\]
La lectura por filas y la lectura diagonal dan
\[
 72+27=77+22=99,
\]
\[
 72-27=45=\det B_c,
 \qquad
 77-22=55=54+1.
\]
La relación con la completación no se obtiene concatenando símbolos, sino
por la división euclídea
\[
 729=10\cdot72+9,
 \qquad
 271=10\cdot27+1.
\]

\begin{theorem}[Unicidad del levantamiento decimal]
\label{p12-83-c1-s5:thm:levantamiento-decimal-emrh-autonomo}
Para \(L_d(n)=10n+d\), con \(d\in\{0,\ldots,9\}\), existe un único par
\((d,d')\) tal que
\[
 L_d(72)=729,
 \qquad
 L_d(72)+L_{d'}(27)=1000,
 \qquad
 d+d'=10.
\]
Es el par
\[
 (d,d')=(9,1).
\]
\end{theorem}

\begin{proof}
La primera igualdad es \(720+d=729\), luego \(d=9\). La segunda queda
\(729+270+d'=1000\), de donde \(d'=1\). La tercera se verifica y la
unicidad de cociente y resto excluye otra solución.
\end{proof}

Así se obtiene el diagrama causal
\[
 B_c\longrightarrow(72,27)
 \longleftarrow((72,9),(27,1))
 \longleftarrow(729,271).
\]
El bloque y su lectura pertenecen a APP; los dividendos y restos pertenecen
a la completación. El diagrama comparte cocientes, no invierte la
causalidad.

\subsection{Incidencia del cubo y separación respecto de hexadas y octadas}

Sea \(M_d\) la matriz de incidencia de las \(2d\) caras de codimensión uno
de un hipercubo sobre sus \(2^d\) vértices. Cada cara contiene
\(2^{d-1}\) vértices; dos caras opuestas son disjuntas y dos caras
transversales comparten \(2^{d-2}\) vértices.

\begin{theorem}[Espectro de incidencia hipercúbica]
\label{p12-83-c1-s5:thm:espectro-incidencia-cubo-autonomo}
\[
 \operatorname{Spec}(M_dM_d^{\mathsf T})
 =\left\{
 d2^{d-1}\,^{(1)},
 2^{d-1}\,^{(d)},
 0^{(d-1)}
 \right\}.
\]
Para \(d=3\),
\[
 \operatorname{Spec}(M_3M_3^{\mathsf T})
 =\{12,4,4,4,0,0\}.
\]
\end{theorem}

\begin{proof}
El espacio de filas se descompone en: la recta constante; las \(d\)
diferencias de cada par de caras opuestas; y las (d-1) combinaciones de
sumas de pares cuya suma total es cero. Las intersecciones descritas antes
hacen que el Gram actúe por \(d2^{d-1}\), \(2^{d-1}\) y \(0\),
respectivamente.
\end{proof}

El cubo tridimensional posee seis caras, ocho vértices y doce aristas. Esos
cardinales anticipan un patrón de incidencia (6/8), pero una cara no es
una hexada y el conjunto de vértices no es una octada. La identificación
excepcional exige un transportador posterior que conserve soportes y
calibre.

\subsection{Resonancia del operador de incidencia en rango \(6\)}

Definamos los estratos que conservan exactamente \(k\) coordenadas
nonádicas en rango \(d\):
\[
 \mathcal I_k(d):=\binom dk9^k.
\]
Los defectos de APP son
\[
 D_{\rm vis}=54,
 \qquad
 D_{\rm tot}=1215,
 \qquad
 D_{\rm coc}=129,
 \qquad
 1215=54+9\cdot129.
\]

\begin{theorem}[Unicidad de la resonancia binomial]
\label{p12-83-c1-s5:thm:resonancia-binomial-autonoma}
El único entero positivo \(d\) que satisface simultáneamente
\[
 \mathcal I_1(d)=54,
 \qquad
 \mathcal I_2(d)=1215
\]
es \(d=6\). En ese rango,
\[
 \frac{\mathcal I_2(6)-\mathcal I_1(6)}9=129.
\]
\end{theorem}

\begin{proof}
La primera igualdad es \(9d=54\), que fuerza \(d=6\). Entonces
\[
 \mathcal I_2(6)=81\binom62=81\cdot15=1215,
\]
y
\[
 \frac{1215-54}{9}=129.
\]
No queda otro rango posible porque la primera igualdad ya determina \(d\).
\end{proof}

\subsection{Obstrucciones del esqueleto ciclotómico de APP: fibras \(A_2^{12}\), completación \(729+271\), resonancia de rango \(6\) e incidencia cúbica}

\begin{criterion}[Falsación tipada de la unidad]
La construcción queda refutada si ocurre cualquiera de los hechos
siguientes:
\begin{enumerate}
 \item la ley de \(D_d\) no reproduce \(D_1=0\) y \(D_2=54\);
 \item se fusiona la variación de longitud con la variación de módulo;
 \item \(C\) deja de preservar \(G_{A_2}\) o adquiere un vector fijo;
 \item las etiquetas par--hoja--órbita no producen doce fibras;
 \item el menor certificado de \(\Psi\) tiene determinante nulo;
 \item un funcional que factoriza por el módulo de aumento recupera \(54\);
 \item el símbolo de frontera \(\star\) se identifica con el residuo cero;
 \item los cuatro estratos no suman \(1000\);
 \item se identifica una cara cúbica con una hexada sin un transportador de
       incidencia.
\end{enumerate}
\end{criterion}

La unidad ha construido, sin reconocerlos retrospectivamente, el módulo
\(A_2^{12}\), la completación \(729+271\), la medida proyectiva y la
resonancia de rango seis. Estas estructuras alimentarán de manera distinta
la generación de constantes y las ramas excepcionales.


% Unidad U006F. Inventario operatorio completo del Topological Phase Kernel.

\long\def\HMTDeferredTPKDevelopment{%
\section{Álgebra de operadores del Topological Phase Kernel y dinámica causal de actualización del estado holonómico integral}
\label{p12-83-c1-s6:chap:inventario-operatorio-completo-tpk}

El Topological Phase Kernel no es una reducción módulo tres ni un reloj de
ciento ocho posiciones. Es una categoría de estados holonómicos integrales equipada
con operadores de selección, transporte, lectura, memoria, periodicidad,
prolongación y publicación. Para impedir que una abreviatura sustituya al
objeto, cada operador se registra por su tipo y por la información que
conserva.

\subsection{Esquema tipado de las 34 construcciones canónicas del TPK: dominio, codominio, acción, invariantes y memoria}

\begin{definition}[Ficha operatoria]
\label{p12-83-c1-s6:def:u006f-ficha-operatoria}
Una realización canónica del TPK es una séxtupla
\begin{equation}
 \mathcal O=
 (D_{\mathcal O},C_{\mathcal O},f_{\mathcal O},
 I_{\mathcal O},L_{\mathcal O},M_{\mathcal O}),
 \label{p12-83-c1-s6:eq:u006f-ficha}
\end{equation}
donde:
\begin{enumerate}
 \item \(f_{\mathcal O}:D_{\mathcal O}\to C_{\mathcal O}\) es su regla;
 \item \(I_{\mathcal O}\) es la familia de relaciones conservadas;
 \item \(L_{\mathcal O}\) declara la información publicada o perdida;
 \item \(M_{\mathcal O}\) enumera la memoria que permanece en el estado
 holonómico integral.
\end{enumerate}
\end{definition}

Dos realizaciones representan el mismo operador sólo cuando existe una
equivalencia tipada que entrelaza reglas e invariantes. Compartir periodo,
cardinal o salida visible no basta.

\subsection{Clasificación del estado holonómico integral en \(8\) clases estructurales invariantes}

El inventario se particiona en:
\begin{enumerate}
 \item[\textbf{C1.}] soportes y estados;
 \item[\textbf{C2.}] selección interna;
 \item[\textbf{C3.}] transporte;
 \item[\textbf{C4.}] lectores y cocientes;
 \item[\textbf{C5.}] memoria y frontera;
 \item[\textbf{C6.}] periodicidad y retorno;
 \item[\textbf{C7.}] prolongación;
 \item[\textbf{C8.}] salidas.
\end{enumerate}
La clase se determina por dominio, codominio y función, no por el nombre de
una cifra asociada.

\Needspace{.92\textheight}
\subsection{Descomposición funcional en \(14\) agrupaciones compatibles con los operadores APP–TRIT–TPK}

\begingroup
\small
\renewcommand{\arraystretch}{1.12}
\begin{longtable}{@{}>{\raggedright\arraybackslash}p{.20\textwidth}>{\raggedright\arraybackslash}p{.29\textwidth}>{\raggedright\arraybackslash}p{.43\textwidth}@{}}
\toprule
Agrupación&Operadores o espacios&Función e invariantes\\
\midrule
\endfirsthead
\toprule
Agrupación&Operadores o espacios&Función e invariantes\\
\midrule
\endhead
Estado y memoria de ruta&
\(\mathcal X_{\rm TPK}^{\rm enr}\), \(\mathcal F_q\)&
Conservan posición, hoja, orientación, residuo, cociente, acarreo, firma,
frontera, cilindro, prefijo, procedencia y memoria.\\

Selección interna&
\(M_{\rm ph}:\{1,\ldots,9\}\to\{-1,0,+1\}\)&
Activa la evaluación aditiva, multiplicativa o el paso de umbral; no mueve
el estado.\\

Transporte espacial&
\(T_N,T_E,T_S,T_O,F_{\rm loc},F_{\rm obs}\)&
Realiza traslaciones cardinales, realimentación de hoja y cronología de dos
cursores.\\

Lectores modulares&
\(\rho_9,q_9,\rho_3^{\rm or},c_3,q_{1000},r_{1000}\)&
Publican fase, orientación y bloque decimal conservando los cocientes de
levantamiento.\\

Bloque y memoria&
\(\mathcal K_{27}=F_{\rm obs}^{27}\), \(\mathcal N_{27}\)&
El primero compone veintisiete transportes; el segundo filtra eventos de
memoria. Igual índice no implica igual operador.\\

Estado dodecafásico&
\(\mathcal R_{12},C_{12},\Theta_{108},\Gamma_{108}\)&
Distingue registro enriquecido, sector, ley cíclica y soporte ordenado.\\

Periodicidad y profundidades&
\(C_{12}\hookrightarrow C_{108}\twoheadrightarrow C_9\),
\(w_{108},w_{1080}\)&
Conserva la extensión no escindida, las palabras de distintas longitudes y
la memoria vertical.\\

Canales \(90/120\)&
\(\mathcal F_6,\mathcal F_8,\mathcal I_{6\to8},
S_{90},S_{120}\)&
Separa incidencias finitas, colas analíticas y sus normalizaciones.\\

Bidireccionalidad&
\(P_\pm,\operatorname{rev},\operatorname{Ext}_\pm,
N_\pm,\beta,C_M\)&
Realiza intercambio, inversión de ruta, prolongación futura y pasada,
valencia, balance y conjugación.\\

Emisión y coinducción&
\(L_0,L_1,R_{36},G_9,\varprojlim\operatorname{Cyl}_m\)&
Eleva la semilla, abre la frontera y organiza prolongaciones finitas y
compatibles.\\

Lectores numéricos&
\(\mathscr C_\pi,\mathscr P_e,F_{\rm av}\)&
Construyen horquillas racionales, propagación factorial y autoescala de
Fibonacci.\\

Proyección excepcional&
\(\Gamma_{A_W}\), peso, soporte proyectivo&
Transporta la misma emisión a incidencia ternaria sin seleccionar cifras
arquimedianas.\\

Atlas y realización&
firma de ruta, atlas \(81/56/13/324\), \(C_M\)&
Clasifica historias y caracteres antes de aplicar una carta dimensional.\\

Validación finita&
censos, pruebas de necesidad estructural, identidades de integridad y
controles de independencia&
Comprueba cobertura en la profundidad ejecutada y separa construcción,
comprobación independiente y datos de contraste.\\
\bottomrule
\end{longtable}
\endgroup

Las ocho clases describen naturaleza matemática; las catorce agrupaciones,
función de ejecución. Son dos gradaciones del mismo inventario.

\Needspace{.92\textheight}
\subsection{Dominio y tipado de las \(34\) entradas del operador de actualización}

\begin{equation}
 \mathfrak O_{\rm TPK}:=(\mathfrak o_1,\ldots,\mathfrak o_{34}).
 \label{p12-83-c1-s6:eq:u006f-registro-operadores}
\end{equation}

\begingroup
\scriptsize
\renewcommand{\arraystretch}{1.12}
\begin{longtable}{@{}r p{.19\textwidth}p{.28\textwidth}p{.42\textwidth}@{}}
\toprule
\#&Símbolo&Dominio y codominio&Regla o función\\
\midrule
\endfirsthead
\toprule
\#&Símbolo&Dominio y codominio&Regla o función\\
\midrule
\endhead
1&\(\mathcal A_9\)&\((\mathbb Z/9)^2\), matriz y ley interna&
Célula aritmética aditivo--multiplicativa.\\
2&\((\rho_9,q_9)\)&\(\mathbb Z_{>0}\to
\{1,\ldots,9\}\times\mathbb Z_{\geq0}\)&
División nonádica con representante positivo y cociente.\\
3&\(\rho_3^{\rm or},c_3\)&\(\mathbb Z\to
\{-1,0,+1\}\times\mathbb Z\)&
Levantamiento ternario orientado \(n=\rho_3^{\rm or}(n)+3c_3(n)\).\\
4&\(\iota_{1000}:=(q_{1000},r_{1000})\)&\(\mathbb Z\to
\mathbb Z\times\{0,\ldots,999\}\)&
División euclídea \(N=1000q_{1000}+r_{1000}\).\\
5&\(T_d\)&\(X_9\to X_9\), \(d\in\{N,E,S,O\}\)&
Traslación cardinal orientada.\\
6&\(F_{\rm loc}\)&\(\mathcal Q_{\rm loc}\times
\mathscr D_{\rm card}\to\mathcal Q_{\rm loc}\)&
Realimentación local de posición y hoja.\\
7&\(T_{\min}\)&\(\Omega_{\rm seed}\to\mathcal U_6\)&
Emisor mínimo de los dos cursores, incluida la firma coordenada.\\
8&\(G\)&\(\operatorname{im}s_\Sigma\times
\operatorname{im}s_\Pi\to\{0,\ldots,999\}^6\)&
Acoplamiento decimal coordenada a coordenada.\\
9&\(\Psi_6\)&\(\mathcal U_6\to\mathbb F_3^6\)&
Reducción ternaria del catálogo.\\
10&\(\mathcal R_3\)&\(\mathcal Q_{\rm loc}\times
\mathscr D_{\rm card}^3\to\{-1,0,+1\}\)&
Factor observable local de tres pasos.\\
11&\(\mathcal K_{27}\)&\(\mathcal Q_{\rm obs}\to\mathcal Q_{\rm obs}\)&
Composición \(F_{\rm obs}^{27}\).\\
12&\(\mathcal X_{\rm TPK}^{\rm enr}\)&categoría de estados holonómicos integrales&
Conserva todas las coordenadas tipadas del núcleo.\\
13&\(\mathcal H\)&relación sobre estados holonómicos integrales&
Transición Hensel que retiene residuo y cociente.\\
14&\(\mathcal C_{3,1000}\)&estados cruzados
\((N,L,D,K,A,B)\)&
Actualización exacta de cilindros ternario--decimales.\\
15&\(\Sigma_{\rm B}\)&memorias de frontera \(\to\) firma conjunta&
Publica la orientación de bloque conservando su registro genealógico.\\
16&\({\rm EF}\text{--}G_9\)&relación de extensión&
Filtra extensiones compatibles y organiza la genealogía nonádica.\\
17&\(\Gamma_9\)&historias enriquecidas \(\longrightarrow\) historias&
Holonomía de una vuelta de fase con memoria de frontera.\\
18&\(X_\chi\)&\(\varprojlim_m\operatorname{Surv}^{(\chi)}_m\)&
Sección global cuando no vaciedad, compatibilidad y unicidad están
demostradas en toda profundidad.\\
19&\(\mathcal P_{\rm APP}^{(2)}\)&categoría bivariada de caminos&
Transporta simultáneamente la evaluación multiplicativa y la acumulación
aditiva.\\
20&\(F_{\rm obs}\)&\(\mathcal Q_{\rm obs}\to\mathcal Q_{\rm obs}\)&
Cronología observable determinista de dos cursores.\\
21&\(\mathcal A_{42}\)&alfabeto de cuarenta y dos tríadas&
Soporte decimal ampliado de la emisión finita.\\
22&\(\operatorname{Ext}_r\)&
\(\operatorname{Ext}_r\subseteq\mathcal F_q\times\mathcal F_{q'}\)&
Relación tipada de prolongación de fibras.\\
23&\(\mathcal K_{\rm ph}\)&\(\mathcal U_6\times C_9\) sobre sí mismo&
Transferencia de continuaciones compatibles con fase explícita.\\
24&\((P,H,Q)\)&\(\mathcal X_\times\to\mathcal X_\times\)&
Avance, media vuelta y cuarto de giro de semillas.\\
25&\((c_{\rm mode},c_{\rm or})\)&caminos \(\to\mathbb Z^2\)&
Cociclos de cambio de modo y orientación.\\
26&\(E_{20}\)&estados finitos \(\to\) veinte bloques&
Prolongación triádica finita con libro de procedencia.\\
27&\(\mathcal R_{12}\)&historia enriquecida \(\to\) registro de doce
sectores&
Sistema temporal orientado con ruta, firma, acarreo y registro genealógico.\\
28&\(R_{\rm sgn}\)&estado terminal \(\to U_{12}^{\rm sgn}
\subseteq\mathbb Z^{12}\)&
Lector del observable armónico firmado.\\
29&\(\mathcal H_{12}\)&\(\mathbb Z^{12}\to\mathbb Z^{12}\)&
Transformación por tres órbitas de paso tres; inversa
\(\mathcal H_{12}/4\) sobre la imagen pertinente.\\
30&\(\mathcal E_{\rm dod}=(D_3,D_4,Q)\)&
\(\mathbb Z^{12}\to\mathbb Z^{25}\)&
Extractor transversal conjunto de aridades tres y cuatro.\\
31&\(\mathcal C_\alpha\)&bloques y acarreos \(\to\) bloques&
Recurrencia dodecafásica de cierre en base mil.\\
32&\((L_{\rm tick},\chi_{\rm mass})\)&historia enriquecida \(\to\)
carácter temporal&
Interfaz dimensional posterior; no participa en el constructor numérico de
esta obra.\\
33&\(\mathcal W_\pm\)&retícula bidireccional de índice doce&
Lectura en ambos sentidos conservando orientación y procedencia.\\
34&\(\mathcal A_{\rm species}\)&familias de extensión \(\to\) atlas&
Interfaz de realización posterior, no un cuarto primitivo.\\
\bottomrule
\end{longtable}
\endgroup

\subsection{Composición causal de la actualización del estado holonómico integral}

\begin{definition}[Estado canónico completo]
\label{p12-83-c1-s6:def:u006f-esquema-estado}
Sea
\[
 \mathcal E:=
 \mathbb Z_{\rm quo}\times\mathbb Z_{\rm car}\times
 \mathsf{Mem}\times\mathsf{Term}\times\mathsf{Cyl}\times
 \mathsf{Fr}\times\mathsf{Pref}\times\Lambda .
\]
El estado holonómico integral canónico en el instante \(t\) es
\begin{equation}
 x_t=(p_t,m_t,\sigma_t,b_t,r_t;
 q_t,c_t,\eta_t,s_t,C_t,\partial_t,N_t,\lambda_t)
 \in
 X_9\times\{\Sigma,\Pi\}\times\mathsf{TRIT}\times
 C_{12}\times C_9\times\mathcal E.
 \label{p12-83-c1-s6:eq:u006f-esquema-estado}
\end{equation}
Las abreviaturas de estado utilizadas en otras unidades son proyecciones de
esta tupla. En particular, no se omiten el cociente \(q_t\), el acarreo
\(c_t\), el terminal \(s_t\), el cilindro \(C_t\), la frontera
\(\partial_t\), el prefijo \(N_t\) ni el libro \(\lambda_t\).
\end{definition}

Sea \(\rho_t\) la flecha observable determinada por \(x_t\). Las leyes de
transporte de U005 y la transición \(\mathcal C_{3,1000}\) determinan un
endomorfismo de memoria
\begin{align}
 \mathbf E_t:\mathcal E&\longrightarrow\mathcal E,\notag\\
 (q,c,\eta,s,C,\partial,N,\lambda)&\longmapsto
 \left(
 \begin{aligned}
 &\mathsf H(\rho_t)q,
  \mathsf C(\rho_t)c+\kappa(\rho_t),
  \mathsf M(\rho_t)\eta,\\
 &\mathsf T(\rho_t)s,
  \mathsf{Cyl}(\rho_t)C,
  \mathsf{Fr}(\rho_t)\partial,\\
 &\mathsf{Pref}(\rho_t)N,
  \Lambda(\rho_t)\lambda
 \end{aligned}
 \right).
 \label{p12-83-c1-s6:eq:u006f-actualizacion-memoria}
\end{align}
Cada componente publicada en \eqref{p12-83-c1-s6:eq:u006f-actualizacion-memoria} posee el
dominio indicado por su fibra; \(\kappa(\rho_t)\) es la cocadena de acarreo.

Definimos tres endomorfismos del producto canónico. Si
\(\sigma'=M_{\rm ph}(g(t))\) y \(m'=\mu(\sigma',m)\), entonces
\begin{align}
 \mathsf{Sel}_t(p,m,\sigma,b,r;e)
  &:=(p,m',\sigma',b,r;e),\label{p12-83-c1-s6:eq:u006f-selector-tipado}\\
 \mathsf{Tra}_t(p,m,\sigma,b,r;e)
  &:=(T_{d(t)}p,m,\sigma,b,r;e),\label{p12-83-c1-s6:eq:u006f-transporte-tipado}\\
 \mathsf{Upd}_t(p,m,\sigma,b,r;e)
  &:=(p,m,\sigma,b',r';\mathbf E_t(e)),
 \label{p12-83-c1-s6:eq:u006f-memoria-tipada}
\end{align}
donde \((b',r')\) es el destino de \(\rho_t:(b,r)\to(b',r')\). La actualización
causal es la composición de mapas, no la composición de un mapa con un valor
trítico:
\begin{equation}
 \boxed{\mathcal U_t:=
 \mathsf{Upd}_t\circ\mathsf{Tra}_t\circ\mathsf{Sel}_t.}
 \label{p12-83-c1-s6:eq:u006f-tuberia}
\end{equation}

\begin{theorem}[Cierre de la actualización causal]
\label{p12-83-c1-s6:thm:u006f-cierre-tuberia}
La aplicación \(\mathcal U_t\) es un endomorfismo del estado canónico
completo. Los lectores, prolongaciones y proyecciones se aplican después de
\(\mathcal U_t\) sólo cuando su dominio ha sido producido.
\end{theorem}

\begin{proof}
\(\mathsf{Sel}_t\) cambia únicamente modo y TRIT;
\(\mathsf{Tra}_t\) cambia únicamente la posición APP; y
\(\mathsf{Upd}_t\) aplica el producto tipado
\eqref{p12-83-c1-s6:eq:u006f-actualizacion-memoria} y la flecha de fase. Los tres
codominios son el mismo producto de la
\cref{p12-83-c1-s6:def:u006f-esquema-estado}; por tanto su composición está definida y
conserva todas las coordenadas del estado.
\end{proof}

\subsection{Separación tipada entre operadores de transporte, memoria, calendario, extracción y retorno}

El inventario impone
\begin{equation}
 \mathcal S_4\neq D_4\neq\operatorname{sd}_4,
 \qquad
 \mathcal R_{12}\neq C_{12}\neq\Theta_{108}\neq\Gamma_{108},
 \qquad
 F_{\rm obs}\neq\mathcal K_{\rm ph},
 \label{p12-83-c1-s6:eq:u006f-no-identificacion-uno}
\end{equation}
y
\begin{equation}
 U_{12}^{\rm cat}\neq U_{12}^{\rm sgn},
 \qquad
 \mathcal E_{\rm dod}\neq D_{\rm raw}^{90,120}
 \neq\mathcal K_{6\to8}.
 \label{p12-83-c1-s6:eq:u006f-no-identificacion-dos}
\end{equation}
Son incompatibilidades de tipo, no meras desigualdades de valor.

\begin{criterion}[Verificación de completitud operatoria]
\label{p12-83-c1-s6:crit:u006f-completitud}
Una ejecución es operatoriamente completa sólo si:
\begin{enumerate}
 \item declara las entradas de \(\mathfrak O_{\rm TPK}\) empleadas;
 \item registra el dominio y el codominio de cada composición;
 \item conserva el orden causal de \eqref{p12-83-c1-s6:eq:u006f-tuberia};
 \item distingue dinámica, transferencia, lectura, prolongación y salida;
 \item permite reconstruir desde \(\lambda_t\) cada cambio de bloque y
 cada acarreo;
 \item no promueve una comprobación finita a una sección infinita sin
 probar no vaciedad, compatibilidad y unicidad a toda profundidad.
\end{enumerate}
Cualquier símbolo usado sin tipo o cualquier fusión prohibida por
\eqref{p12-83-c1-s6:eq:u006f-no-identificacion-uno} y
\eqref{p12-83-c1-s6:eq:u006f-no-identificacion-dos} constituye un fallo verificable.
\end{criterion}


% Unidad U006E. Separación tipada de los canales 90/120 y de los símbolos
% epsilon. La realización de incidencia se construye después de W12 y W24.

\section{Descomposición operatorial de los canales de incidencia \texorpdfstring{\(90/120\)}{90/120}: selector temporal, generador nilpotente y componente angular}
\label{p12-83-c1-s7:chap:canales-90-120-epsilon}

Los enteros \(90\) y \(120\) aparecen en varios estratos de la construcción. Su
coincidencia numérica no establece una identidad de objetos. Del mismo modo,
la letra epsilon designa tres construcciones con dominios incompatibles. La
finalidad de esta unidad es fijar sus tipos antes de emplearlos y señalar
con precisión qué realizaciones requieren estructuras que todavía no se han
construido.

\subsection{Operadores temporal, nilpotente y angular: dominios y codominios}

\subsubsection{Acción del selector temporal trítico sobre \(\{-1,0,+1\}\): apertura, umbral y retorno con memoria}

El selector interno del calendario es
\begin{equation}
 \varepsilon_9:\{1,\ldots,9\}
 \longrightarrow\{-1,0,+1\},
 \qquad
 \varepsilon_9(r)=
 \begin{cases}
  +1,&r\in\{1,4,7\},\\
  -1,&r\in\{2,5,8\},\\
  0,&r\in\{3,6,9\}.
 \end{cases}
 \label{p12-83-c1-s7:eq:u006e-epsilon-temporal}
\end{equation}
Su codominio es el TRIT orientado. Decide qué hoja actúa y no desplaza por
sí mismo ninguna posición.

\subsubsection{Generador parabólico nilpotente \(J_0^2=0\) y evolución en el umbral temporal}

La clase
\begin{equation}
 \varepsilon_{\rm nil}
 \in
 \mathcal Q_{\rm n}:=
 \mathbb F_3[\varepsilon_{\rm nil}]
 /(\varepsilon_{\rm nil}^{\,2})
 \label{p12-83-c1-s7:eq:u006e-epsilon-nilpotente}
\end{equation}
satisface
\begin{equation}
 \varepsilon_{\rm nil}\neq0,
 \qquad
 \varepsilon_{\rm nil}^{\,2}=0.
 \label{p12-83-c1-s7:eq:u006e-epsilon-nilpotente-relacion}
\end{equation}
Es un elemento de un álgebra cuadrática degenerada; no es un selector
temporal.

\subsubsection{Residuo angular inducido después de la selección temporal y conservación de orientación}

El símbolo
\begin{equation}
 \varepsilon_{\rm res}^{\circ}\in\mathbb R
 \label{p12-83-c1-s7:eq:u006e-epsilon-residual-tipo}
\end{equation}
designa, cuando se introduzcan sus parámetros geométricos, el residuo de una
identidad angular. En esta etapa sólo queda reservado su tipo. No se lo
obtiene desde \(\varepsilon_9\) ni desde
\(\varepsilon_{\rm nil}\), y no interviene en la construcción de
\(\pi,e,\varphi\) o \(\alpha\).

\begin{proposition}[No identificación de los símbolos epsilon]
\label{p12-83-c1-s7:prop:u006e-no-identificacion-epsilon}
No existe una igualdad bien tipada entre
\[
 \varepsilon_9,\qquad
 \varepsilon_{\rm nil},\qquad
 \varepsilon_{\rm res}^{\circ}.
\]
\end{proposition}

\begin{proof}
El primer objeto es una función entre conjuntos finitos; el segundo, un
elemento nilpotente de un álgebra sobre \(\mathbb F_3\); el tercero, un
escalar real con unidad angular. Sus dominios, codominios y leyes de
composición son distintos.
\end{proof}

\subsection{Distribución de los índices \(90/120\) entre los \(6\) objetos del sistema de incidencia}

\begin{definition}[Extractor dodecafásico]
\label{p12-83-c1-s7:def:u006e-extractor-dodecafasico}
El extractor integral
\begin{equation}
 \mathcal E_{\rm dod}:=(D_3,D_4,Q):
 \mathbb Z^{12}\longrightarrow\mathbb Z^{25}
 \label{p12-83-c1-s7:eq:u006e-extractor-dodecafasico}
\end{equation}
lee diferencias de aridades tres y cuatro y una coordenada de suma. Los
subíndices no son longitudes angulares.
\end{definition}

\begin{definition}[Diferencia cruda de series]
\label{p12-83-c1-s7:def:u006e-diferencia-cruda}
Para \(|q|<1\), se define
\begin{equation}
 S_m^{\rm raw}(q):=\frac{q^m}{1-q^{3m}},
 \qquad
 D_{\rm raw}^{90,120}(q)
 :=S_{90}^{\rm raw}(q)-S_{120}^{\rm raw}(q).
 \label{p12-83-c1-s7:eq:u006e-series-crudas}
\end{equation}
Este objeto es una función analítica de \(q\); no es un censo finito.
\end{definition}

\begin{definition}[Familias marcadas de incidencia]
\label{p12-83-c1-s7:def:u006e-familias-incidencia}
Una vez construidas una hexada \(H\), su imagen
\(\ell(H)\) en una dodecada binaria \(D\), y la octada elevada
\(O_H\), se definirán
\begin{align}
 \mathcal F_6(H)
 &:=\ell(H)\times\binom{\ell(H)}2,
 \label{p12-83-c1-s7:eq:u006e-familia-seis}\\
 \mathcal F_8(H)
 &:=O_H\times\binom{\ell(H)}2.
 \label{p12-83-c1-s7:eq:u006e-familia-ocho}
\end{align}
Por construcción,
\begin{equation}
 |\mathcal F_6(H)|
 =6\binom62=90,
 \qquad
 |\mathcal F_8(H)|
 =8\binom62=120.
 \label{p12-83-c1-s7:eq:u006e-cardinales-incidencia}
\end{equation}
La inclusión \(\ell(H)\hookrightarrow O_H\) inducirá
\begin{equation}
 \mathcal I_{6\to8}:
 \mathcal F_6(H)\hookrightarrow\mathcal F_8(H).
 \label{p12-83-c1-s7:eq:u006e-inclusion-incidencia}
\end{equation}
\end{definition}

Las \cref{p12-83-c1-s7:eq:u006e-familia-seis,p12-83-c1-s7:eq:u006e-familia-ocho} fijan desde ahora el
tipo. Su existencia efectiva se demostrará después de construir
\(C_W\), \(W_{12}\), el dato binario \(G_{24},D,\ell\) y \(W_{24}\).
No se anticipa aquí esa rama.

\begin{definition}[Combinación normalizada de colas]
\label{p12-83-c1-s7:def:u006e-combinacion-colas}
Para parámetros \(q_-,q_+\in(0,1)\), sean
\begin{equation}
 S_m(q):=\frac{q^m}{1-q^{3m}},
 \qquad
 \Delta S_m:=S_m(q_-)-S_m(q_+).
 \label{p12-83-c1-s7:eq:u006e-colas-normalizadas}
\end{equation}
El funcional
\begin{equation}
 \mathcal K_{6\to8}
 :=12\,\Delta S_{90}-\Delta S_{120}
 \label{p12-83-c1-s7:eq:u006e-funcional-seis-ocho}
\end{equation}
es una combinación analítica de colas. Su coeficiente doce pertenece a su
normalización y no lo convierte en el registro
\(\mathcal R_{12}\).
\end{definition}

\begin{definition}[Semigrupo de grados y momentos]
\label{p12-83-c1-s7:def:u006e-semigrupo-momentos}
El semigrupo aditivo generado por ambos grados es
\begin{equation}
 \mathfrak S_{90,120}
 :=\langle90,120\rangle
 =\{90a+120b:a,b\in\mathbb N_0\}.
 \label{p12-83-c1-s7:eq:u006e-semigrupo}
\end{equation}
Para \(n\in\mathfrak S_{90,120}\), el momento orientado es
\begin{equation}
 s_n:=q_-^n-q_+^n.
 \label{p12-83-c1-s7:eq:u006e-momento}
\end{equation}
\end{definition}

Los momentos son monomios diferenciales indexados por un semigrupo; las
series de \eqref{p12-83-c1-s7:eq:u006e-colas-normalizadas} son sumas infinitas de esos
momentos sobre progresiones específicas.

\subsection{Defecto normalizado de la incidencia}

Cuando la inclusión \(\mathcal I_{6\to8}\) haya sido construida, su
diferencia de cardinalidades será
\[
 |\mathcal F_6(H)|-|\mathcal F_8(H)|=90-120=-30.
\]
Si se declara la normalización
\[
 24\binom62=360,
\]
se obtiene
\begin{equation}
 \frac{90-120}{24\binom62}
 =\frac{6-8}{24}
 =-\frac1{12}.
 \label{p12-83-c1-s7:eq:u006e-defecto-normalizado}
\end{equation}
La inclusión determina \(-30\); el racional \(-1/12\) depende además de
la normalización declarada. No se lo debe presentar como si la cardinalidad
por sí sola fijara el denominador.

\subsection{Invariantes de tipado}

\begin{theorem}[Separación de los seis objetos]
\label{p12-83-c1-s7:thm:u006e-separacion-seis-objetos}
Los objetos
\begin{equation}
 \mathcal E_{\rm dod},\quad
 D_{\rm raw}^{90,120},\quad
 \mathcal I_{6\to8},\quad
 \mathcal K_{6\to8},\quad
 \mathfrak S_{90,120},\quad
 \varepsilon_{\rm res}^{\circ}
 \label{p12-83-c1-s7:eq:u006e-seis-objetos}
\end{equation}
son mutuamente distintos por tipo: son, respectivamente, un lector entero,
una función analítica, un morfismo finito, una combinación de series, un
semigrupo de índices y un escalar real angular.
\end{theorem}

\begin{proof}
La conclusión se obtiene comparando sus dominios y codominios:
\[
 \mathbb Z^{12}\to\mathbb Z^{25},\quad
 \mathbb D\to\mathbb C,\quad
 \mathcal F_6(H)\to\mathcal F_8(H),\quad
 (0,1)^2\to\mathbb R,\quad
 \mathbb N_0^2\to\mathbb N_0,\quad
 \{\ast\}\to\mathbb R.
\]
No hay dos firmas iguales.
\end{proof}

\begin{criterion}[Criterios de refutación]
\label{p12-83-c1-s7:crit:u006e-refutacion}
La clasificación queda refutada si se altera alguno de sus dominios o
codominios, si el censo de
\eqref{p12-83-c1-s7:eq:u006e-cardinales-incidencia} no reproduce \(90\) y \(120\), si
se atribuye \(-1/12\) a la inclusión sin declarar la normalización, o si
se presenta el residual angular como salida del generador de constantes de
esta monografía.
\end{criterion}


% Unidad U006A. Factor local, observacion estroboscopica y transicion exacta
% de cilindros.

\section{Factor local de transporte y transición ternario--decimal}
\label{p12-83-c1-s8:chap:tpk-factor-local-transicion-cilindros}

Esta unidad separa tres operaciones que intervienen en el transporte del
Topological Phase Kernel y que no comparten dominio: la realización cardinal
de un trit, el muestreo estroboscópico de una sucesión y la actualización
exacta de un par de cilindros ternario--decimales. La primera es una dinámica
finita; la segunda es un lector temporal; la tercera conserva los enteros
que permiten decidir compatibilidad sin aritmética de punto flotante.

\subsection{Estado local y realimentación de hoja}

Sea
\[
 X_9=(\mathbb Z/9\mathbb Z)^2,
 \qquad
 \mathscr D_{\rm card}=\{N,E,S,O\},
\]
con
\[
 \delta(N)=(-1,0),\quad \delta(E)=(0,1),\quad
 \delta(S)=(1,0),\quad \delta(O)=(0,-1).
\]
Las evaluaciones visibles de APP son
\[
 L_\Sigma(x,y)=\rho_9((x+1)+(y+1)),
 \qquad
 L_\Pi(x,y)=\rho_9((x+1)(y+1)),
\]
donde \(\rho_9(n)=1+((n-1)\bmod9)\). Definimos además
\[
 \rho_3^{\rm or}([0])=0,
 \qquad
 \rho_3^{\rm or}([1])=+1,
 \qquad
 \rho_3^{\rm or}([2])=-1.
\]

\begin{definition}[Factor local de APP--TRIT]
\label{p12-83-c1-s8:def:u006a-factor-local}
El espacio de estados locales es
\[
 \mathcal Q_{\rm loc}=X_9\times\{\Sigma,\Pi\},
 \qquad |\mathcal Q_{\rm loc}|=162.
\]
Para \(q=(p,m)\), \(d\in\mathscr D_{\rm card}\), ponemos
\[
 p_d=p+\delta(d),
 \qquad
 a_d(q)=\rho_3^{\rm or}([L_m(p_d)]_3),
\]
y
\[
 \mu(a,m)=
 \begin{cases}
  \Sigma,&a=+1,\\
  \Pi,&a=-1,\\
  m,&a=0.
 \end{cases}
\]
La transición local es
\[
 F_d^{\rm loc}(p,m)=(p_d,\mu(a_d(p,m),m)).
\]
\end{definition}

La posición se modifica antes de evaluar la hoja. El trit nulo conserva el
último modo activo; no elimina la hoja ni reinicializa la historia.

Para una palabra cardinal \(u=d_1d_2d_3\), sea
\[
 q_j=F_{d_j}^{\rm loc}(q_{j-1}),\qquad q_0=q,
\]
y denotemos por
\[
 R_3(q,u)\in\{-1,0,+1\}
\]
el trit producido en el tercer desplazamiento.

\begin{definition}[Observación de paso tres]
\label{p12-83-c1-s8:def:u006a-obs3}
Sobre una sucesión orientada \(a=(a_t)_{t\geq1}\) definimos
\[
 \operatorname{Obs}_3(a)=(a_3,a_6,a_9,\ldots).
\]
El mapa \(R_3\) toma un estado y una palabra cardinal de longitud tres;
\(\operatorname{Obs}_3\) toma una sucesión y publica una subsucesión. Son
operadores de tipos distintos.
\end{definition}

\begin{theorem}[Sobreyectividad exacta del factor muestreado]
\label{p12-83-c1-s8:thm:u006a-sobreyectividad-r3}
Para todo \(q\in\mathcal Q_{\rm loc}\) y todo
\(b\in\{-1,0,+1\}\), la fibra
\[
 \mathcal R_3(q,b)
 =\{u\in\mathscr D_{\rm card}^3:R_3(q,u)=b\}
\]
es no vacía y satisface
\[
 \boxed{8\leq |\mathcal R_3(q,b)|\leq40.}
\]
En particular, las \(162\cdot3=486\) fibras estado--salida son no vacías.
\end{theorem}

\begin{proof}
Desde cada uno de los \(162\) estados se enumeran las \(4^3=64\) palabras
cardinales y se aplica literalmente la recurrencia de la
\cref{p12-83-c1-s8:def:u006a-factor-local}. El censo exhaustivo de los \(486\) pares
\((q,b)\) tiene mínimo ocho, máximo cuarenta y ningún valor cero. El
enunciado es, por tanto, una identidad finita sobre el sistema definido, no
una hipótesis de mezcla.
\end{proof}

\begin{corollary}[Realización de una sucesión orientada arbitraria]
\label{p12-83-c1-s8:cor:u006a-realizacion-sucesion}
Para toda sucesión
\((b_k)_{k\geq1}\in\{-1,0,+1\}^{\mathbb N}\) existe una historia cardinal
cuya observación satisface
\[
 (a_3,a_6,a_9,\ldots)=(b_1,b_2,b_3,\ldots).
\]
\end{corollary}

\begin{proof}
Fijado un orden total de \(\mathscr D_{\rm card}^3\), en la etapa \(k\) se
toma el primer elemento de \(\mathcal R_3(q_{k-1},b_k)\). El estado terminal
del bloque vuelve a pertenecer a \(\mathcal Q_{\rm loc}\), de modo que la
elección puede iterarse indefinidamente.
\end{proof}

El corolario demuestra capacidad de realización del factor muestreado. No
selecciona las secciones de \(\pi,e,\varphi\), no fija los dos trits
intermedios y no sustituye la compatibilidad del estado holonómico integral.

\subsection{Cilindros de los \(2\) lectores}

Una palabra ternaria de longitud \(L\), interpretada como el entero \(N\),
determina
\[
 I_3(N,L)=\left[\frac{N}{3^L},\frac{N+1}{3^L}\right).
\]
Un prefijo de \(K\) tríadas decimales, interpretado como el entero \(D\),
determina
\[
 I_{1000}(D,K)=
 \left[\frac{D}{1000^K},\frac{D+1}{1000^K}\right).
\]
Introducimos los residuos cruzados
\[
 A=N1000^K-D3^L,
 \qquad
 B=(N+1)1000^K-D3^L=A+1000^K.
\]

\begin{lemma}[Criterio entero de compatibilidad]
\label{p12-83-c1-s8:lem:u006a-criterio-entero}
Los cilindros \(I_3(N,L)\) e \(I_{1000}(D,K)\) se intersectan si y sólo si
\[
 \boxed{A<3^L,\qquad B>0.}
\]
Equivalentemente,
\[
 N1000^K<(D+1)3^L,
 \qquad
 (N+1)1000^K>D3^L.
\]
\end{lemma}

\begin{proof}
Dos intervalos semiabiertos \([a,b)\) y \([c,d)\) se intersectan
exactamente cuando \(a<d\) y \(c<b\). La multiplicación por
\(3^L1000^K>0\) produce las dos desigualdades. Restar \(D3^L\) y
\(N1000^K\), respectivamente, da la forma en \(A,B\).
\end{proof}

\subsection{Operador de completación \(C_{3,1000}\): extensión ternaria y conservación de la unidad en \(10^3\) posiciones}

Un nuevo bloque de seis trits se representa por
\[
 u\in\{0,\ldots,3^6-1\}=\{0,\ldots,728\}.
\]
La actualización ternaria es
\[
 N'=729N+u,
 \qquad L'=L+6.
\]
La publicación decimal tiene dos casos. Si el cilindro refinado todavía no
determina una nueva tríada, entonces \(K'=K\), \(D'=D\). Si determina
\(\delta\in\{0,\ldots,999\}\), entonces
\[
 K'=K+1,
 \qquad D'=1000D+\delta.
\]

\begin{definition}[Estado cruzado y transición exacta]
\label{p12-83-c1-s8:def:u006a-c31000}
El estado de cilindros es la séxtupla
\[
 \mathfrak c=(N,L,D,K,A,B)
\]
sujeta a
\[
 A=N1000^K-D3^L,
 \qquad B=A+1000^K.
\]
La transición \(\mathcal C_{3,1000}\) adjunta \(u\) y, cuando corresponde,
\(\delta\), y actualiza
\[
 (N,L,D,K,A,B)\longmapsto
 (N',L',D',K',A',B')
\]
mediante
\[
 A'=
 \begin{cases}
  729A+u1000^K,&K'=K,\\[1mm]
  729000A+u1000^{K+1}-\delta\,729\,3^L,&K'=K+1,
 \end{cases}
 \qquad
 B'=A'+1000^{K'}.
\]
La prolongación es admisible exactamente cuando
\[
 A'<3^{L'},\qquad B'>0.
\]
\end{definition}

\begin{theorem}[Exactitud y suficiencia del estado cruzado]
\label{p12-83-c1-s8:thm:u006a-exactitud-c31000}
La transición de la \cref{p12-83-c1-s8:def:u006a-c31000} conserva exactamente las
identidades
\[
 A'=N'1000^{K'}-D'3^{L'},
 \qquad
 B'=(N'+1)1000^{K'}-D'3^{L'}.
\]
Por consiguiente, \((A',B',L',K')\) basta para decidir la compatibilidad del
nuevo par de cilindros, mientras \((N',D')\) conserva su procedencia
reconstructiva.
\end{theorem}

\begin{proof}
Si no aparece tríada decimal,
\[
 N'1000^K-D3^{L+6}
 =(729N+u)1000^K-729D3^L
 =729A+u1000^K.
\]
Si aparece \(\delta\),
\begin{align*}
 N'1000^{K+1}-D'3^{L+6}
 &=(729N+u)1000^{K+1}-(1000D+\delta)7293^L\\
 &=729000A+u1000^{K+1}-\delta\,729\,3^L.
\end{align*}
En ambos casos, aumentar \(N'\) en una unidad añade \(1000^{K'}\), lo que
da \(B'\). El criterio de la \cref{p12-83-c1-s8:lem:u006a-criterio-entero} concluye.
\end{proof}

\begin{proposition}[Compatibilidad con la concatenación]
\label{p12-83-c1-s8:prop:u006a-concatenacion-c31000}
Aplicar sucesivamente \(\mathcal C_{3,1000}\) a dos bloques produce el
mismo par \((N'',D'')\) y los mismos residuos cruzados que adjuntar de una
vez las dos palabras ternarias y las tríadas decimales publicadas en el
mismo orden.
\end{proposition}

\begin{proof}
Las recurrencias \(N\mapsto729N+u\) y
\(D\mapsto1000D+\delta\) son las leyes de concatenación posicional en
bases \(729\) y \(1000\). El
\cref{p12-83-c1-s8:thm:u006a-exactitud-c31000} identifica \(A,B\) con funciones de esos
cuatro enteros; por tanto, su actualización por etapas coincide con la
evaluación posterior a la concatenación.
\end{proof}

\subsection{Descomposición del Topological Phase Kernel en operadores de avance, giro, acarreo y memoria}

La cadena correcta es
\[
 (q,u)\xmapsto{\ R_3\ }b,
 \qquad
 (a_t)_{t\geq1}
 \xmapsto{\ \operatorname{Obs}_3\ }
 (a_{3k})_{k\geq1},
 \qquad
 (\mathfrak c,u,\delta)
 \xmapsto{\ \mathcal C_{3,1000}\ }
 \mathfrak c'.
\]
\(R_3\) prueba realizabilidad local; \(\operatorname{Obs}_3\) publica una
subsucesión; \(\mathcal C_{3,1000}\) transporta compatibilidad entre dos
bases. Ninguno de los tres puede reemplazar a la relación de supervivencia
del TPK completo.

\begin{criterion}[Criterios de refutación]
\label{p12-83-c1-s8:crit:u006a-refutacion}
La construcción queda refutada en su dominio si ocurre cualquiera de los
hechos siguientes:
\begin{enumerate}
 \item alguna de las \(486\) fibras \(\mathcal R_3(q,b)\) es vacía o queda
       fuera del intervalo \([8,40]\);
 \item el trit nulo cambia el modo activo;
 \item se identifica \(R_3\) con \(\operatorname{Obs}_3\);
 \item la recurrencia de \(A'\) no reproduce su definición directa;
 \item una prolongación declarada compatible viola \(A'<3^{L'}\) o
       \(B'>0\);
 \item la decisión depende de una tolerancia numérica no presente en las
       desigualdades enteras.
\end{enumerate}
\end{criterion}


% Unidad U006D. Registro dodecafásico integral y descomposición 1+5+6.

\section{Registro dodecafásico integral, extensión cíclica y descomposición \texorpdfstring{\(1+5+6\)}{1+5+6}}
\label{p12-83-c1-s9:chap:registro-dodecafasico-integral}

El periodo observable de ciento ocho posiciones contiene doce sectores de
nueve posiciones. Esta igualdad de cardinales no basta para describir el
estado temporal: deben distinguirse el soporte ordenado, la ley cíclica, el
índice sectorial y el registro enriquecido de transportes, orientaciones,
acarreos e incidencias. Esta unidad construye los cuatro objetos y demuestra
la reconstrucción integral de las doce coordenadas de memoria.

\subsection{Soporte ordenado y coordenada cíclica}

Para \(m\in\{1,\ldots,12\}\), definimos
\begin{equation}
 E_m:=\{9(m-1)+1,\ldots,9m\}.
 \label{p12-83-c1-s9:eq:u006d-bloques-nonadicos}
\end{equation}
El soporte ordenado y el grupo cíclico son
\begin{equation}
 \Gamma_{108}:=
 E_1\sqcup\cdots\sqcup E_{12},
 \qquad
 \Theta_{108}:=\mathbb Z/108\mathbb Z.
 \label{p12-83-c1-s9:eq:u006d-gamma-theta}
\end{equation}
\(\Gamma_{108}\) conserva posición, orden y pertenencia a bloque;
\(\Theta_{108}\) conserva la ley de traslación. No son el mismo objeto.

Para \(\theta\in\Theta_{108}\), sea
\(\widetilde\theta\in\{0,\ldots,107\}\) su representante. Definimos
\begin{equation}
 \beta(\theta):=
 \left[
 \left\lfloor\frac{\widetilde\theta}{9}\right\rfloor
 \right]_{12},
 \qquad
 r(\theta):=1+(\widetilde\theta\bmod9).
 \label{p12-83-c1-s9:eq:u006d-coordenadas-bloque-residuo}
\end{equation}
La primera coordenada localiza el sector dodecafásico; la segunda publica
la posición positiva \(1,\ldots,9\) dentro del sector.

\subsection{Extensión cíclica no escindida}

Sean \(C_n=\mathbb Z/n\mathbb Z\). Las aplicaciones
\begin{equation}
 \iota:C_{12}\longrightarrow C_{108},
 \quad \iota([b]_{12})=[9b]_{108},
 \qquad
 p:C_{108}\longrightarrow C_9,
 \quad p([\theta]_{108})=[\theta]_9
 \label{p12-83-c1-s9:eq:u006d-mapas-extension}
\end{equation}
forman la sucesión
\begin{equation}
 0\longrightarrow C_{12}
 \xrightarrow{\;\iota\;}C_{108}
 \xrightarrow{\;p\;}C_9
 \longrightarrow0.
 \label{p12-83-c1-s9:eq:u006d-extension}
\end{equation}
Para la sección conjuntista
\[
 s([r]_9):=[\widetilde r]_{108},
 \qquad 0\leq\widetilde r<9,
\]
el defecto de aditividad es
\begin{equation}
 c(r,r'):=
 \left[
 \left\lfloor
 \frac{\widetilde r+\widetilde r'}9
 \right\rfloor
 \right]_{12},
 \qquad
 s(r)+s(r')
 =s(r+r')+\iota(c(r,r')).
 \label{p12-83-c1-s9:eq:u006d-cociclo-extension}
\end{equation}

\begin{theorem}[Exactitud, no escisión y clase cohomológica]
\label{p12-83-c1-s9:thm:u006d-extension-no-escindida}
La sucesión \eqref{p12-83-c1-s9:eq:u006d-extension} es exacta y no se escinde. Para la
acción trivial,
\begin{equation}
 H^2(C_9,C_{12})
 \simeq C_{12}/9C_{12}
 \simeq C_3,
 \label{p12-83-c1-s9:eq:u006d-cohomologia-extension}
\end{equation}
y el cociclo \eqref{p12-83-c1-s9:eq:u006d-cociclo-extension} representa su clase
generadora.
\end{theorem}

\begin{proof}
La imagen de \(\iota\) es el subgrupo de múltiplos de nueve, que coincide
con \(\ker p\); \(p\) es sobreyectiva. Esto prueba la exactitud. Si la
extensión se escindiera, \(C_{108}\) sería isomorfo a
\(C_{12}\times C_9\), pero
\[
 \exp(C_{108})=108,
 \qquad
 \exp(C_{12}\times C_9)=\operatorname{mcm}(12,9)=36.
\]
La identidad de cociclo es la división euclídea de
\(\widetilde r+\widetilde r'\) por nueve. Para un grupo cíclico con acción
trivial,
\(H^2(C_9,C_{12})\simeq C_{12}/9C_{12}\); el acarreo unitario proyecta a
\([1]\), que genera el cociente de orden tres.
\end{proof}

La no escisión tipa el retorno:
\[
 g_0(\theta+9)=g_0(\theta),
 \qquad
 \beta(\theta+9)=\beta(\theta)+1\pmod{12}.
\]
Nueve iteraciones cierran la fase visible y trasladan un sector; ciento
ocho cierran las dos coordenadas del reloj. El estado holonómico integral conserva
además la memoria vertical, el cilindro, la frontera y la genealogía.

\subsection{Registro temporal enriquecido: codificación conjunta de fase, hoja, frontera y memoria}

\begin{definition}[Registro dodecafásico orientado]
\label{p12-83-c1-s9:def:u006d-registro-dodecafasico}
Para una historia enriquecida, sea \(\tau_m\) la subruta ordenada sobre
\(E_m\), \(\sigma_m\) su orientación, \(\kappa_m\) el acarreo que cruza la
frontera del bloque y \(\Lambda_m\) su libro local de incidencias. El
registro completo es
\begin{equation}
 \mathcal R_{12}:=
 \bigl(
 (E_m,\tau_m,\sigma_m,\kappa_m,\Lambda_m)
 \bigr)_{m=1}^{12}.
 \label{p12-83-c1-s9:eq:u006d-registro-r12}
\end{equation}
\end{definition}

Se distinguen así cuatro niveles:
\begin{equation}
 \boxed{
 \mathcal R_{12}
 \neq C_{12}
 \neq\Theta_{108}
 \neq\Gamma_{108}.}
 \label{p12-83-c1-s9:eq:u006d-cuatro-objetos}
\end{equation}
El primero es un registro enriquecido; el segundo, un índice de sector; el
tercero, un grupo de traslaciones; el cuarto, un soporte ordenado. Existen
proyecciones entre ellos, pero no identificaciones.

Sea \(q_{\rm mem}\in\mathbb Z_{\geq0}\) el número no acotado de vueltas
nonádicas. La coordenada dodecafásica es sólo su residuo
\[
 \beta=[q_{\rm mem}]_{12}.
\]
Tras ciento ocho pasos,
\[
 \beta\longmapsto\beta,
 \qquad
 q_{\rm mem}\longmapsto q_{\rm mem}+12.
\]
El calendario retorna; la memoria vertical no se reinicia.

\subsection{Extractor firmado de eventos}

Sea \((d_t)_{t=1}^{108}\) el registro digital sobre
\(\Gamma_{108}\). Definimos
\begin{equation}
 \mathcal D_{\rm evt}:=\{2,4,6,8\},
 \qquad
 \Sigma_{12}:=(+,+,+,-,-,-,+,+,+,-,-,-).
 \label{p12-83-c1-s9:eq:u006d-datos-eventos}
\end{equation}
Para \(m=0,\ldots,11\), sea
\begin{equation}
 e_m:=
 \Sigma_{12}(m+1)\,
 \#\{t\in E_{m+1}:d_t\in\mathcal D_{\rm evt}\}
 \in\mathbb Z.
 \label{p12-83-c1-s9:eq:u006d-eventos-firmados}
\end{equation}
La aplicación conserva el censo local y la orientación sectorial como
factores separados.

Emparejamos las dos semicoronas mediante
\begin{equation}
 p_i:=e_i+e_{i+6},
 \qquad
 a_i:=e_i-e_{i+6},
 \qquad 0\leq i\leq5.
 \label{p12-83-c1-s9:eq:u006d-pares-antipodales}
\end{equation}
Definimos el censo total, las cinco diferencias respecto del sexto par y
las seis antisimetrías:
\begin{equation}
 s_C:=\sum_{i=0}^{5}p_i,
 \qquad
 M_i:=p_i-p_5\quad(0\leq i\leq4),
 \qquad
 A_6:=(a_0,\ldots,a_5).
 \label{p12-83-c1-s9:eq:u006d-coordenadas-uno-cinco-seis}
\end{equation}

\begin{definition}[Lector integral \(1+5+6\)]
\label{p12-83-c1-s9:def:u006d-lector-integral}
El lector lineal es
\begin{equation}
 \mathcal L_{12}:\mathbb Z^{12}\longrightarrow\mathbb Z^{12},
 \qquad
 \mathcal L_{12}(e)
 :=
 (s_C,M_0,\ldots,M_4,a_0,\ldots,a_5).
 \label{p12-83-c1-s9:eq:u006d-lector-l12}
\end{equation}
La distribución de coordenadas es
\begin{equation}
 \underbrace{1}_{s_C}
 +\underbrace{5}_{M_0,\ldots,M_4}
 +\underbrace{6}_{a_0,\ldots,a_5}
 =12.
 \label{p12-83-c1-s9:eq:u006d-uno-cinco-seis}
\end{equation}
\end{definition}

\begin{theorem}[Reconstrucción exacta e imagen integral]
\label{p12-83-c1-s9:thm:u006d-inversa-l12}
El lector \(\mathcal L_{12}\) es inyectivo. Sobre su imagen,
\begin{align}
 p_5&=
 \frac{s_C-\sum_{i=0}^{4}M_i}{6},
 &
 p_i&=p_5+M_i\quad(0\leq i\leq4),
 \label{p12-83-c1-s9:eq:u006d-inversa-p}\\
 e_i&=\frac{p_i+a_i}{2},
 &
 e_{i+6}&=\frac{p_i-a_i}{2}
 \quad(0\leq i\leq5).
 \label{p12-83-c1-s9:eq:u006d-inversa-e}
\end{align}
Una tupla de salida pertenece a \(\operatorname{im}\mathcal L_{12}\) si y
sólo si
\begin{equation}
 s_C-\sum_{i=0}^{4}M_i\equiv0\pmod6
 \label{p12-83-c1-s9:eq:u006d-condicion-divisibilidad}
\end{equation}
y, para los \(p_i\) reconstruidos,
\begin{equation}
 p_i\equiv a_i\pmod2
 \qquad(0\leq i\leq5).
 \label{p12-83-c1-s9:eq:u006d-condicion-paridad}
\end{equation}
\end{theorem}

\begin{proof}
De \(M_i=p_i-p_5\) se obtiene
\[
 s_C=\sum_{i=0}^{4}(p_5+M_i)+p_5
 =6p_5+\sum_{i=0}^{4}M_i,
\]
que da la primera ecuación de \eqref{p12-83-c1-s9:eq:u006d-inversa-p}; las demás son
inmediatas. Sumar y restar
\(p_i=e_i+e_{i+6}\) y \(a_i=e_i-e_{i+6}\) produce
\eqref{p12-83-c1-s9:eq:u006d-inversa-e}. La divisibilidad por seis y las seis
congruencias de paridad son precisamente las condiciones necesarias y
suficientes para que todas las coordenadas reconstruidas sean enteras.
\end{proof}

La división por seis aparece después del censo orientado, no durante él. Las
condiciones integrales de esta unidad son exactamente
\eqref{p12-83-c1-s9:eq:u006d-condicion-divisibilidad} y
\eqref{p12-83-c1-s9:eq:u006d-condicion-paridad}; no se atribuye aquí una forma normal de
Smith a una matriz que no haya sido previamente definida.

\subsection{\(2\) palabras dodecafásicas de distinta procedencia}

Debe conservarse la distinción
\begin{equation}
 U_{12}^{\rm cat}:=U_6\Vert U_6,
 \qquad
 U_{12}^{\rm sgn}:=\mathcal H_{12}(K).
 \label{p12-83-c1-s9:eq:u006d-dos-palabras-doce}
\end{equation}
La primera concatena dos emisiones catalogales de longitud seis. La segunda
es una coordenada firmada ligada reversiblemente a \(K\) por el operador
\(\mathcal H_{12}\). Su longitud común no implica una misma procedencia,
semántica o ley de transporte.

\begin{criterion}[Criterios de refutación]
\label{p12-83-c1-s9:crit:u006d-refutacion}
La construcción queda refutada si la sucesión
\eqref{p12-83-c1-s9:eq:u006d-extension} no es exacta, si su clase cohomológica es
trivial, si una tupla de la imagen viola
\eqref{p12-83-c1-s9:eq:u006d-condicion-divisibilidad} o
\eqref{p12-83-c1-s9:eq:u006d-condicion-paridad}, si las fórmulas inversas no recuperan
los doce enteros iniciales, o si el retorno de \(\beta\) se interpreta como
reinicio de \(q_{\rm mem}\).
\end{criterion}


% Unidad U006B. Funtor ciclotomico y elevacion modular de caminos.

\section{Funtor ciclotómico de fase--longitud y elevación modular de caminos}
\label{p12-83-c1-s10:chap:funtor-ciclotomico-elevacion-caminos}

El cociente odométrico \(\mathcal O_{9,12}\) conserva simultáneamente la
posición nonádica, el registro dodecafásico y la longitud entera de cada
flecha. En esta unidad se realiza su fase sobre fibras \(A_2\) y se prueba
la elevación cardinal que sustituye cada arista por cuatro aristas
orientadas. La representación ciclotómica y la longitud absoluta se
mantendrán como coordenadas distintas.

\subsection{Representación de Coxeter de orden \(3\) sobre las fibras dodecafásicas de tipo \(A_2\)}

Sea
\begin{equation}
 G_{A_2}:=
 \begin{pmatrix}2&-1\\-1&2\end{pmatrix},
 \qquad
 C:=
 \begin{pmatrix}0&-1\\1&-1\end{pmatrix}.
 \label{p12-83-c1-s10:eq:u006b-dato-coxeter}
\end{equation}
Por el \cref{p12-83-c1-s5:prop:coxeter-a2-autonomo},
\begin{equation}
 C^3=I_2,
 \qquad
 C^{\mathsf T}G_{A_2}C=G_{A_2},
 \qquad
 \det(XI_2-C)=X^2+X+1.
 \label{p12-83-c1-s10:eq:u006b-identidades-coxeter}
\end{equation}
Así, \(C\) realiza el carácter ciclotómico de \(C_3\) sobre \(A_2\).

Para cada \(b\in\mathbb Z/12\mathbb Z\), sea \(A_{2,b}\) una copia de
\(A_2\), sea
\[
 \iota_b:A_{2,b}\xrightarrow{\;\sim\;}A_2
\]
una identificación integral y sea
\(P_b\in O(G_{A_2})\) un marco que depende de \(b\), pero no del residuo
\(r\). Definimos
\begin{equation}
 \rho_b:=
 \iota_b^{-1}P_bCP_b^{-1}\iota_b,
 \qquad
 \rho_b^3=\operatorname{id}_{A_{2,b}}.
 \label{p12-83-c1-s10:eq:u006b-rho-fibra}
\end{equation}
La independencia respecto de \(r\) liga el carácter a la fibra
dodecafásica; \(r\) conserva la posición interna de la vuelta nonádica.

\subsection{Categoría odométrica del TPK: cilindros, truncamientos y morfismos de sucesión}

Recordemos que \(\mathcal O_{9,12}\) tiene objetos
\[
 (b,r)\in\mathbb Z/12\mathbb Z\times\{0,\ldots,8\}.
\]
Para \(n\in\mathbb N\), la flecha única de longitud \(n\) es
\begin{equation}
 \gamma_{(b,r),n}:
 (b,r)\longrightarrow
 \left(
 b+\kappa_r(n),[r+n]_9
 \right),
 \qquad
 \kappa_r(n):=
 \left\lfloor\frac{r+n}{9}\right\rfloor.
 \label{p12-83-c1-s10:eq:u006b-flecha-odometro}
\end{equation}
La composición se apoya en
\begin{equation}
 \kappa_r(n+m)
 =\kappa_r(n)+\kappa_{[r+n]_9}(m).
 \label{p12-83-c1-s10:eq:u006b-cociclo-odometro}
\end{equation}
La longitud \(\ell(\gamma_{(b,r),n})=n\) es aditiva.

\begin{definition}[Funtor ciclotómico de fase--longitud]
\label{p12-83-c1-s10:def:u006b-funtor-ciclotomico}
Sea \(\mathbf{Rep}_{\mathbb C}(C_3)\) la categoría de representaciones
complejas de \(C_3\), y sea \(\mathbf B\mathbb N\) la categoría monoidal
de un solo objeto con endomorfismos \((\mathbb N,+)\). Definimos
\begin{equation}
 \mathfrak F_{\rm cyc}:
 \mathcal O_{9,12}
 \longrightarrow
 \mathbf{Rep}_{\mathbb C}(C_3)\times\mathbf B\mathbb N
 \label{p12-83-c1-s10:eq:u006b-signatura-funtor}
\end{equation}
por
\[
 \mathfrak F_{\rm cyc}(b,r)
 =(A_{2,b}\otimes_{\mathbb Z}\mathbb C,\rho_b)
\]
y
\begin{equation}
 \mathfrak F_{\rm cyc}(\gamma)
 :=
 \left(
 \iota_{b'}^{-1}P_{b'}C^{\ell(\gamma)}
 P_b^{-1}\iota_b,
 \ell(\gamma)
 \right)
 \label{p12-83-c1-s10:eq:u006b-flecha-funtor}
\end{equation}
para \(\gamma:(b,r)\to(b',r')\).
\end{definition}

\begin{theorem}[Functorialidad y separación de los registros]
\label{p12-83-c1-s10:thm:u006b-functorialidad}
La asignación de la \cref{p12-83-c1-s10:def:u006b-funtor-ciclotomico} es un funtor. Su
primera componente conserva únicamente la clase de longitud módulo tres; la
segunda conserva el entero completo y, por reducción, determina esa clase.
La componente ciclotómica no permite reconstruir la longitud absoluta.
\end{theorem}

\begin{proof}
Si \(\gamma_1:(b_0,r_0)\to(b_1,r_1)\) y
\(\gamma_2:(b_1,r_1)\to(b_2,r_2)\) tienen longitudes \(n_1,n_2\), la
trivialización intermedia se cancela:
\begin{align}
 &\bigl(
 \iota_{b_2}^{-1}P_{b_2}C^{n_2}P_{b_1}^{-1}\iota_{b_1}
 \bigr)
 \bigl(
 \iota_{b_1}^{-1}P_{b_1}C^{n_1}P_{b_0}^{-1}\iota_{b_0}
 \bigr)\notag\\
 &\qquad=
 \iota_{b_2}^{-1}P_{b_2}C^{n_1+n_2}
 P_{b_0}^{-1}\iota_{b_0}.
 \label{p12-83-c1-s10:eq:u006b-composicion-funtor}
\end{align}
La segunda componente compone por suma. Como \(C^n\) sólo depende de
\([n]_3\), pero \(\ell=n\) conserva el entero, las dos coordenadas no se
identifican.
\end{proof}

Si
\[
 Q:\mathsf{TPK}\longrightarrow\mathcal O_{9,12}
\]
es el funtor odométrico de U006, la realización tipada es
\begin{equation}
 \mathfrak F_{\rm cyc}\circ Q:
 \mathsf{TPK}\longrightarrow
 \mathbf{Rep}_{\mathbb C}(C_3)\times\mathbf B\mathbb N.
 \label{p12-83-c1-s10:eq:u006b-compuesto-tpk-ciclotomico}
\end{equation}

\subsection{Elevación modular de caminos cardinales}

Sea \(\mathcal A_{N,d}\) la categoría libre de caminos cardinales de
\((\mathbb Z/N\mathbb Z)^d\). Sean
\[
 N\ge2,\quad u\ge1,\quad\gcd(u,N)=1,
 \quad k:=\operatorname{ord}_N(u),
 \quad q:=\frac{u^k-1}{N}.
\]

\begin{definition}[Sustitución cardinal por \(u\)]
\label{p12-83-c1-s10:def:u006b-sustitucion-cardinal}
El funtor
\[
 \mathcal S_u:\mathcal A_{N,d}\longrightarrow\mathcal A_{N,d}
\]
envía un objeto \(p\) a \(up\) y cada arista orientada a la concatenación
ordenada de \(u\) copias con la misma dirección.
\end{definition}

\begin{theorem}[Elevación modular a endomorfismo de caminos]
\label{p12-83-c1-s10:thm:u006b-elevacion-modular}
Para todo objeto \(p\), camino \(\gamma\) y
\[
 \kappa_{N,r}(n):=
 \left\lfloor\frac{r+n}{N}\right\rfloor,
\]
se cumple
\begin{align}
 \mathcal S_u^k(p)&=p,
 \label{p12-83-c1-s10:eq:u006b-retorno-objetos}\\
 \ell(\mathcal S_u^k\gamma)&=u^k\ell(\gamma),
 \label{p12-83-c1-s10:eq:u006b-longitud-sustituida}\\
 \kappa_{N,r}(u^kn)-\kappa_{N,r}(n)&=qn.
 \label{p12-83-c1-s10:eq:u006b-diferencia-vueltas}
\end{align}
\end{theorem}

\begin{proof}
La categoría es libre, de modo que las imágenes de objetos y generadores
determinan un funtor único. Cada aplicación multiplica la longitud por
\(u\). Como \(u^k=1+qN\), los objetos retornan módulo \(N\), y
\[
 \left\lfloor
 \frac{r+n+qNn}{N}
 \right\rfloor
 =
 \left\lfloor\frac{r+n}{N}\right\rfloor+qn.
\]
\end{proof}

En el caso nonádico,
\begin{equation}
 (N,u,k,q)=(9,4,3,7),
 \qquad
 4^3=64=1+7\cdot9.
 \label{p12-83-c1-s10:eq:u006b-parametros-nonadicos}
\end{equation}
Por tanto,
\begin{equation}
 \boxed{
 \mathcal S_4^3(e)
 \text{ conserva los extremos y contiene siete vueltas nonádicas
 orientadas por arista inicial}.}
 \label{p12-83-c1-s10:eq:u006b-siete-vueltas}
\end{equation}
Los extremos y la posición residual retornan en tres aplicaciones; el
estado holonómico integral no tiene por qué retornar. Sin una prueba de
invertibilidad, \(\mathcal S_4^3\) es un endomorfismo vertical de caminos,
no una monodromía.

\begin{remark}[Separación respecto de la dilatación del reloj]
\label{p12-83-c1-s10:rem:u006b-s4-no-d4}
\(\mathcal S_4\) sustituye objetos y aristas en una categoría de caminos.
La aplicación \(D_4\) de U006C multiplicará por cuatro la coordenada
\(\Theta_{108}\). Comparten un factor entero, pero no dominio, codominio ni
ley de composición.
\end{remark}

\begin{criterion}[Criterios de refutación]
\label{p12-83-c1-s10:crit:u006b-refutacion}
El bloque queda refutado si falla cualquiera de las identidades
\eqref{p12-83-c1-s10:eq:u006b-identidades-coxeter},
\eqref{p12-83-c1-s10:eq:u006b-cociclo-odometro},
\eqref{p12-83-c1-s10:eq:u006b-composicion-funtor} o
\eqref{p12-83-c1-s10:eq:u006b-parametros-nonadicos}; también queda mal tipado si se
identifican la fase \(C^n\), la longitud \(n\), la sustitución
\(\mathcal S_4\) y la dilatación \(D_4\).
\end{criterion}
}%


% Unidad U004. Redacción autónoma desde los generadors TRIT de 1063/live.

\section{TRIT como unidad mínima de información topológica: orientación ternaria, memoria de acarreo y orden temporal}
\label{p12-83-c1-s11:chap:trit-levantado-regimenes}

La curvatura mixta de APP produjo tres valores orientados: cero en las hojas
puras y signos opuestos al intercambiar el orden suma--producto. El TRIT
tipa esa tríada y conserva, además del residuo visible, el cociente necesario
para reconstruir el entero levantado. Esta unidad separa expresamente cuatro
niveles: clase ternaria, representante balanceado, acarreo y régimen
cuadrático.

\begin{definition}[Unidad de información topológica TRIT]
El \emph{TRIT} es la unidad elemental orientada cuyo espacio de estados es
\(\mathsf{TRIT}=\{-1,0,+1\}\). Su capacidad logarítmica nativa es
\(\log 3\) y su expresión en unidades binarias es \(\log_2 3\). A diferencia
de un bit, sus tres estados no se reducen a una codificación cardinal: tipan,
respectivamente, apertura, umbral y retorno, y conservan el cociente de
acarreo requerido para reconstruir la transición que produjo el residuo
visible.
\end{definition}

\subsection{Sistema ternario orientado y representación balanceada de los estados TRIT}

Definimos el conjunto orientado
\[
 \mathsf{TRIT}:=\{-1,0,+1\}.
\]
La carta balanceada es la biyección de conjuntos
\[
 \operatorname{bal}:\mathbb F_3\longrightarrow\mathsf{TRIT},
 \qquad
 [0]_3\mapsto0,\qquad[1]_3\mapsto+1,\qquad[2]_3\mapsto-1.
\]
No se trata de una igualdad de tipos. \(\mathbb F_3\) porta suma y producto
de cuerpo; \(\mathsf{TRIT}\) porta orden y orientación. Toda operación
transportada entre ambos debe indicar la carta utilizada.

La orientación de las nueve fases es
\[
 \varepsilon_9(r)=
 \begin{cases}
 +1,&r\in\{1,4,7\},\\
 -1,&r\in\{2,5,8\},\\
 0,&r\in\{3,6,9\}.
 \end{cases}
\]
Su tabla completa es
\[
\begin{array}{c|ccccccccc}
r&1&2&3&4&5&6&7&8&9\\\hline
\varepsilon_9(r)&+1&-1&0&+1&-1&0&+1&-1&0
\end{array}.
\]
La periodicidad de tres no elimina la fase nonádica: las fases
\(1,4,7\), por ejemplo, tienen la misma orientación y posiciones distintas
en la vuelta.

\subsection{Algoritmo de división ternaria balanceada y unicidad del par cociente–residuo}

\begin{theorem}[Levantamiento ternario local]
\label{p12-83-c1-s11:thm:trit-levantamiento-local}
Para cada \(n\in\{-4,-3,\ldots,4\}\) existe un único par
\((r,c)\in\mathsf{TRIT}^2\) tal que
\[
 n=r+3c.
\]
\end{theorem}

\begin{proof}
La existencia se verifica en la tabla
\[
\begin{array}{c|rrrrrrrrr}
n&-4&-3&-2&-1&0&1&2&3&4\\\hline
r&-1&0&+1&-1&0&+1&-1&0&+1\\
c&-1&-1&-1&0&0&0&+1&+1&+1
\end{array}.
\]
Si \(r+3c=r'+3c'\), entonces \(r-r'=3(c'-c)\). Como
\(r-r'\in\{-2,-1,0,1,2\}\), el único múltiplo de tres posible es cero;
por tanto \(r=r'\) y \(c=c'\).
\end{proof}

La proyección visible \((r,c)\mapsto r\) posee, sobre el valor cero, tres
estados locales:
\[
 0^-:=(0,-1),
 \qquad
 0^0:=(0,0),
 \qquad
 0^+:=(0,+1).
\]
Publican el mismo residuo, pero reconstruyen \(-3,0,+3\). Esta fibra triple
es la forma mínima de una distinción que reaparecerá en el estado holonómico integral
del Topological Phase Kernel.

\begin{corollary}[Iteración balanceada]
\label{p12-83-c1-s11:cor:trit-iteracion-balanceada}
Todo entero \(N\) admite una expansión ternaria balanceada finita
\[
 N=r_0+3r_1+\cdots+3^mr_m,
 \qquad r_j\in\mathsf{TRIT}.
\]
Se obtiene iterando la división \(N_k=r_k+3N_{k+1}\), donde \(r_k\) es el
representante balanceado de \(N_k\bmod3\).
\end{corollary}

\begin{proof}
La división euclídea módulo tres da un residuo en \(\{0,1,2\}\); sustituir
\(2\) por \(-1\) incrementa el cociente en una unidad y produce
\(r_k\in\mathsf{TRIT}\). El valor absoluto del cociente disminuye fuera de
un conjunto finito, por lo que el proceso termina. La unicidad se obtiene
comparando la primera cifra no coincidente y reduciendo módulo tres.
\end{proof}

\subsection{Composición y cociclo de acarreo}

Para \(r,s\in\mathsf{TRIT}\), definimos \(r\oplus s\in\mathsf{TRIT}\) y
\(\kappa_3(r,s)\in\mathbb Z\) por la igualdad única
\[
 r+s=(r\oplus s)+3\kappa_3(r,s).
\]
La tabla es
\[
\begin{array}{c|ccc}
\oplus&-1&0&+1\\\hline
-1&+1&-1&0\\
0&-1&0&+1\\
+1&0&+1&-1
\end{array},
\qquad
\begin{array}{c|ccc}
\kappa_3&-1&0&+1\\\hline
-1&-1&0&0\\
0&0&0&0\\
+1&0&0&+1
\end{array}.
\]

\begin{proposition}[Identidad cocíclica ternaria]
\label{p12-83-c1-s11:prop:trit-cociclo}
Para \(r,s,t\in\mathsf{TRIT}\),
\[
 \kappa_3(r,s)+\kappa_3(r\oplus s,t)
 =\kappa_3(s,t)+\kappa_3(r,s\oplus t).
\]
\end{proposition}

\begin{proof}
Ambos miembros son el cociente total que aparece al escribir
\(r+s+t\) como un representante balanceado más un múltiplo de tres. La
unicidad del levantamiento prueba la igualdad.
\end{proof}

Si \((r,c)\) y \((s,d)\) son estados levantados, su suma es
\[
 (r,c)\boxplus(s,d)
 :=\bigl(r\oplus s,\ c+d+\kappa_3(r,s)\bigr).
\]
La \cref{p12-83-c1-s11:prop:trit-cociclo} prueba que \(\boxplus\) es asociativa y que la
reconstrucción \(\operatorname{rec}_3(r,c)=r+3c\) satisface
\[
 \operatorname{rec}_3(x\boxplus y)
 =\operatorname{rec}_3(x)+\operatorname{rec}_3(y).
\]

\subsection{Involución de orientación y orden temporal}

Para una palabra trítica \(\tau=(\tau_1,\ldots,\tau_n)\), definimos
\[
 \mathcal C_{\rm rev}(\tau_1,\ldots,\tau_n)
 :=(-\tau_n,\ldots,-\tau_1).
\]

\begin{lemma}[Reversión orientada]
\label{p12-83-c1-s11:lem:trit-reversion}
La aplicación \(\mathcal C_{\rm rev}\) es una involución. Intercambia los
sectores \(+1\) y \(-1\) y fija el sector \(0\).
\end{lemma}

\begin{proof}
Aplicarla dos veces invierte dos veces el orden y cambia dos veces cada
signo. Así se recupera la palabra inicial. Las demás afirmaciones son
inmediatas en cada componente.
\end{proof}

Cuando se fija un parámetro temporal orientado, la convención del TPK es
\[
 +1:\text{retorno con memoria},
 \qquad
 0:\text{umbral presente},
 \qquad
 -1:\text{apertura de prolongación}.
\]
Esta asignación no añade un cuarto estado temporal. La dirección se conserva
en el signo y la posición exacta se conserva en la fase nonádica.

\subsection{Familia universal de álgebras cuadráticas}

Para \(\tau\in\mathsf{TRIT}\), definimos
\[
 \mathbb A_\tau:=\mathbb R[J_\tau]/(J_\tau^2+\tau I).
\]
Los tres casos son:
\[
\begin{array}{c|c|c|c}
\tau&J_\tau^2&\text{régimen}&\exp(tJ_\tau)\\\hline
+1&-I&\text{elíptico}&\cos t\,I+\sin t\,J_{+1}\\
0&0&\text{parabólico}&I+tJ_0\\
-1&+I&\text{hiperbólico}&\cosh t\,I+\sinh t\,J_{-1}
\end{array}.
\]

\begin{theorem}[Exponencial cuadrática uniforme]
\label{p12-83-c1-s11:thm:trit-exponencial-uniforme}
Si \(J_\tau^2=-\tau I\), entonces
\[
 \exp(tJ_\tau)
 =C_\tau(t)I+S_\tau(t)J_\tau,
\]
donde
\[
\begin{array}{c|ccc}
\tau&+1&0&-1\\\hline
C_\tau(t)&\cos t&1&\cosh t\\
S_\tau(t)&\sin t&t&\sinh t
\end{array}.
\]
\end{theorem}

\begin{proof}
Separamos la serie exponencial en potencias pares e impares. Las relaciones
\(J_{+1}^{2m}=(-1)^mI\), \(J_0^2=0\) y
\(J_{-1}^{2m}=I\) producen, respectivamente, las series de coseno y seno,
la truncación lineal y las series hiperbólicas.
\end{proof}

Sobre \(\mathbb F_3\) aparecen los análogos
\[
 \mathbb F_9=\mathbb F_3[\iota]/(\iota^2+1),
 \qquad
 \mathbb D_3=\mathbb F_3[\epsilon]/(\epsilon^2),
 \qquad
 \mathbb S_3=\mathbb F_3[j]/(j^2-1).
\]
El primero es un cuerpo de nueve elementos; el segundo contiene un nilpotente
no nulo; el tercero se descompone mediante los idempotentes
\((1\pm j)/2\) y es isomorfo a \(\mathbb F_3\times\mathbb F_3\). La
coincidencia cardinal no identifica estas tres álgebras.

La fuente integral del régimen elíptico es
\[
 \mathbb Q_{\mathbb Z}:=\mathbb Z[u]/(u^2+1).
\]
La reducción módulo tres envía \(u\) a \(\iota\) y produce
\(\mathbb F_9\). La extensión de escalares a \(\mathbb R\) y una
representación \(u\mapsto J_{+1}\) producen el realizador real elíptico.
No existe por ello un mapa no tipado \(\mathbb F_9\to\mathbb R\): ambas
realizaciones proceden de una fuente integral común.

\subsection{Representaciones matriciales del TRIT seleccionadas por el signo de \(J_\tau^2\)}

Sea
\[
 C=\begin{pmatrix}0&-1\\1&-1\end{pmatrix},
 \qquad C^3=I,
\]
y definamos
\[
 J_e:=\frac{2C+I}{\sqrt3}
 =\frac1{\sqrt3}\begin{pmatrix}1&-2\\2&-1\end{pmatrix}.
\]
Una multiplicación directa da
\[
 (2C+I)^2=-3I,
 \qquad J_e^2=-I.
\]
El realizador parabólico es
\[
 J_p:=N=\begin{pmatrix}0&1\\0&0\end{pmatrix},
 \qquad N^2=0.
\]
Para el régimen hiperbólico tomamos
\[
 F_{\rm av}=\begin{pmatrix}0&1\\1&1\end{pmatrix},
 \qquad
 A:=F_{\rm av}^2=\begin{pmatrix}1&1\\1&2\end{pmatrix},
\]
y
\[
 J_h:=\frac{2A-3I}{\sqrt5}
 =\frac1{\sqrt5}\begin{pmatrix}-1&2\\2&1\end{pmatrix}.
\]
Entonces
\[
 (2A-3I)^2=5I,
 \qquad J_h^2=I.
\]

La forma bilineal
\[
 G=\begin{pmatrix}2&1\\1&-2\end{pmatrix}
\]
satisface
\[
 A^TGA=G.
\]
Además,
\[
 \cosh(2\log\varphi)=\frac32,
 \qquad
 \sinh(2\log\varphi)=\frac{\sqrt5}{2},
\]
por lo que
\[
 A=\exp(2\log\varphi\,J_h)
  =\cosh(2\log\varphi)I+\sinh(2\log\varphi)J_h.
\]
Esta identidad será recuperada causalmente en la biografía de
\(\varphi\); aquí sólo se verifica la representación matricial.

Los tres controles exponenciales son
\[
 \exp(\pi J_e)+I=0,
 \qquad
 \exp(J_p)=I+J_p,
 \qquad
 \exp(2\log\varphi\,J_h)=F_{\rm av}^2.
\]
La presencia de \(\pi\) y \(\varphi\) en estas identidades de verificación
no autoriza a usarlas como constructores de dichas constantes. Sus
constructores aparecerán después desde la genealogía nonádica.

\subsection{Acoplamiento local APP--TRIT}

Con el intercambio
\[
 R=\begin{pmatrix}0&1\\1&0\end{pmatrix}
\]
definimos, para \(\tau\in\mathsf{TRIT}\),
\[
 B_\tau=(4+3\tau)I+(5-3\tau)R.
\]
Puesto que \(P_\pm=(I\pm R)/2\),
\[
 B_\tau=9P_++(6\tau-1)P_-.
\]

\begin{theorem}[Recuperación espectral de la orientación]
\label{p12-83-c1-s11:thm:app-trit-recuperacion-espectral}
El espectro de \(B_\tau\) es
\[
 \operatorname{Spec}(B_\tau)=\{9,6\tau-1\}.
\]
El autovalor diferencial \(\lambda_-\) determina unívocamente el TRIT:
\[
 \tau=\frac{\lambda_-+1}{6}.
\]
Además, el bloque central de APP es \(B_{+1}\).
\end{theorem}

\begin{proof}
Los proyectores \(P_+\) y \(P_-\) diagonalizan \(R\) con autovalores
\(+1\) y \(-1\). La fórmula espectral se obtiene sustituyendo. La segunda
identidad despeja \(\tau\). Para \(\tau=+1\),
\(B_{+1}=7I+2R=\begin{psmallmatrix}7&2\\2&7\end{psmallmatrix}\), que es
el bloque único de la \cref{p12-83-c1-s1:thm:app-bloque-central}.
\end{proof}

Por tanto, el TRIT no se añade como etiqueta retrospectiva: puede recuperarse
del modo diferencial de una familia construida sobre el intercambio interno
de APP.

\subsection{Obstrucciones del TRIT: separación residuo–representante–acarreo, involutividad y recuperación espectral de \(\tau\)}

\begin{criterion}[Criterios de refutación del TRIT]
La unidad queda refutada si:
\begin{enumerate}
 \item se identifican clase residual, representante balanceado y acarreo;
 \item se afirma que \(0^-\), \(0^0\) y \(0^+\) son el mismo estado;
 \item la composición \(\boxplus\) deja de reconstruir la suma entera;
 \item \(\mathcal C_{\rm rev}\) no es involutiva;
 \item se identifican las tres álgebras por tener el mismo cardinal;
 \item algún realizador incumple \(J_\tau^2=-\tau I\);
 \item el valor de \(\tau\) se elige después de observar el resultado y no
       se recupera del espectro de \(B_\tau\).
\end{enumerate}
\end{criterion}

APP proporciona el soporte, la matriz multiplicativa y la ley aditiva
interna; TRIT proporciona orientación,
levantamiento y tipo cuadrático. La unidad siguiente construirá el
Topological Phase Kernel como categoría de transporte de esos datos, sin
confundir el selector interno con el estado completo.

\HMTDeferredTPKBase
\HMTDeferredTPKDevelopment


% Unidad U006C. Dilatación aritmética, cociclo de acarreo y
% semiconjugación cuaternaria.

\section{Dilatación cuaternaria del registro, obstrucción observable y semiconjugación refinada}
\label{p12-83-c1-s12:chap:dilatacion-cuaternaria-semiconjugacion}

La sustitución cardinal \(\mathcal S_4\) construida en la unidad anterior
actúa sobre caminos. En esta unidad se define una segunda operación, de
naturaleza aritmética, sobre el registro
\(\mathbb Z/12\mathbb Z\times\mathbb Z/9\mathbb Z\), y una tercera
operación, categorial, que subdivide cada transición observable en cuatro
microtransiciones. El objetivo no es fusionarlas, sino demostrar la relación
exacta entre sus realizaciones.

\subsection{Dilatación sobre el registro producto}

Sea
\[
 \mathcal Q_{12,9}:=
 \mathbb Z/12\mathbb Z\times\mathbb Z/9\mathbb Z.
\]
Para \(r\in\mathbb Z/9\mathbb Z\), denotamos por
\(\widetilde r\in\{0,\ldots,8\}\) su representante normalizado.

\begin{definition}[Dilatación cuaternaria del registro]
\label{p12-83-c1-s12:def:u006c-dilatacion-registro}
Definimos
\begin{equation}
 D_4(b,r):=
 \left(
 \left[
  4b+\left\lfloor\frac{4\widetilde r}{9}\right\rfloor
 \right]_{12},
 [4r]_9
 \right).
 \label{p12-83-c1-s12:eq:u006c-d4-producto}
\end{equation}
La coordenada total es
\begin{equation}
 \Theta:\mathcal Q_{12,9}\xrightarrow{\;\sim\;}
 \mathbb Z/108\mathbb Z,
 \qquad
 \Theta(b,r):=[9b+\widetilde r]_{108}.
 \label{p12-83-c1-s12:eq:u006c-coordenada-total}
\end{equation}
\end{definition}

\begin{proposition}[Realización por multiplicación modular]
\label{p12-83-c1-s12:prop:u006c-d4-total}
En la coordenada total,
\begin{equation}
 \boxed{
 \Theta\circ D_4\circ\Theta^{-1}(\theta)=[4\theta]_{108}.}
 \label{p12-83-c1-s12:eq:u006c-d4-total}
\end{equation}
En particular,
\begin{equation}
 \operatorname{im}(D_4)
 =4\mathbb Z/108\mathbb Z
 \simeq\mathbb Z/27\mathbb Z,
 \label{p12-83-c1-s12:eq:u006c-imagen-d4}
\end{equation}
y la restricción inducida sobre la imagen tiene orden nueve.
\end{proposition}

\begin{proof}
Para \(\theta=9b+\widetilde r\),
\[
 4\theta
 =9\left(
  4b+\left\lfloor\frac{4\widetilde r}{9}\right\rfloor
 \right)
 +(4\widetilde r\bmod9).
\]
Ésta es exactamente la descomposición cociente--residuo de
\eqref{p12-83-c1-s12:eq:u006c-d4-producto}. La imagen está formada por las clases
múltiplos de cuatro. Tras dividir por cuatro se obtiene
\(\mathbb Z/27\mathbb Z\), y
\[
\begin{aligned}
 4^1&\equiv4,& 4^2&\equiv16,& 4^3&\equiv10,\\
 4^4&\equiv13,& 4^5&\equiv25,& 4^6&\equiv19,\\
 4^7&\equiv22,& 4^8&\equiv7,& 4^9&\equiv1\pmod{27}.
\end{aligned}
\]
Ninguna potencia anterior es uno; por tanto
\(\operatorname{ord}_{27}(4)=9\).
\end{proof}

\subsection{Identidad exacta de acarreo}

Para \(r\in\{0,\ldots,8\}\) y \(n\in\mathbb N\), definimos
\[
 \kappa_r(n):=\left\lfloor\frac{r+n}{9}\right\rfloor,
 \qquad
 r':=[r+n]_9.
\]

\begin{theorem}[Cociclo cuaternario de acarreo]
\label{p12-83-c1-s12:thm:u006c-cociclo-d4}
Se cumple
\begin{equation}
 \boxed{
 \kappa_{[4r]_9}(4n)
 =4\kappa_r(n)
 +\left\lfloor\frac{4\widetilde r'}{9}\right\rfloor
 -\left\lfloor\frac{4\widetilde r}{9}\right\rfloor.}
 \label{p12-83-c1-s12:eq:u006c-cociclo-d4}
\end{equation}
La última diferencia es el coborde producido por la elección de
representantes \(r\mapsto\widetilde r\).
\end{theorem}

\begin{proof}
Escribamos
\[
 r+n=9\kappa_r(n)+\widetilde r'.
\]
Al multiplicar por cuatro,
\[
 4r+4n=36\kappa_r(n)+4\widetilde r'.
\]
Separar cociente y residuo módulo nueve y restar el cociente ya contenido en
\(4\widetilde r\) produce
\eqref{p12-83-c1-s12:eq:u006c-cociclo-d4}.
\end{proof}

\begin{corollary}[Tercera iteración]
\label{p12-83-c1-s12:cor:u006c-tercera-iteracion}
Para todo \((b,r)\in\mathcal Q_{12,9}\),
\begin{equation}
 \boxed{D_4^3(b,r)=([4b+7\widetilde r]_{12},r).}
 \label{p12-83-c1-s12:eq:u006c-d4-cubo}
\end{equation}
\end{corollary}

\begin{proof}
Por la \cref{p12-83-c1-s12:prop:u006c-d4-total},
\[
 \Theta(D_4^3(b,r))
 =[4^3(9b+\widetilde r)]_{108}
 =[64(9b+\widetilde r)]_{108}.
\]
Como \(64=1+7\cdot9\), el residuo módulo nueve vuelve a ser \(r\), mientras
la coordenada de bloque se transforma en
\([4b+7\widetilde r]_{12}\).
\end{proof}

El retorno de \(r\) no es retorno del registro completo. La diferencia
\(7\widetilde r\) conserva la memoria del representante interno atravesado.

\subsection{Obstrucción en el cociente observable}

En la configuración observable de
\cref{p12-83-c1-s3:def:tpk-configuracion-observable}, escribimos
\[
 q=(p^\Sigma,p^\Pi,d^\Sigma,d^\Pi,\theta)
 \in\mathcal Q_{\rm obs}.
\]
La proyección posicional es
\begin{equation}
 \pi_{\rm pos}(q):=(p^\Sigma,p^\Pi)\in X_9^2.
 \label{p12-83-c1-s12:eq:u006c-proyeccion-posicional}
\end{equation}
La firma temporal es
\[
 \tau(q):=\varepsilon_9(r(\theta))\in\{-1,0,+1\}.
\]
Definimos el desplazamiento observable orientado
\begin{equation}
 a(q):=
 \begin{cases}
  (\delta(d^\Sigma),0),&\tau(q)=+1,\\
  (0,\delta(d^\Pi)),&\tau(q)=-1,\\
  (0,0),&\tau(q)=0.
 \end{cases}
 \label{p12-83-c1-s12:eq:u006c-desplazamiento-observable}
\end{equation}
Entonces
\[
 \pi_{\rm pos}(F_{\rm obs}q)=\pi_{\rm pos}(q)+a(q),
\]
con la actualización de direcciones y reloj definida en U005.

\begin{theorem}[Obstrucción al endofuntor observable ordinario]
\label{p12-83-c1-s12:thm:u006c-obstruccion}
No existe un endofuntor del espacio observable ordinario que satisfaga
simultáneamente:
\begin{enumerate}
 \item multiplicar la posición geométrica por cuatro;
 \item sustituir cada transición por cuatro microtransiciones genuinas;
 \item conservar hoja y orientación en cada microtransición;
 \item enviar una fibra dodecafásica completa a una sola fibra
 dodecafásica.
\end{enumerate}
La obstrucción persiste en un paso de umbral con \(a(q)=0\).
\end{theorem}

\begin{proof}
Sea
\[
 \mathcal B_b:=\{(b,r):0\leq r\leq8\}.
\]
Sustituir una iteración elemental por cuatro fuerza
\[
 h(\theta+1)=h(\theta)+4
\]
y, por inducción,
\[
 h(\theta)=4\theta+c\pmod{108}.
\]
Tomemos \(0\leq\widetilde c<108\) y escribamos
\(\widetilde c=9c_b+c_0\), con \(0\leq c_0<9\). Si toda la fibra
\(\mathcal B_b\) fuese enviada a una sola fibra, el índice de bloque de
destino tendría que ser independiente de \(r\); sin embargo contiene
\begin{equation}
 4b+c_b+
 \left\lfloor\frac{4r+c_0}{9}\right\rfloor
 \pmod{12},
 \label{p12-83-c1-s12:eq:u006c-bloque-obstruccion}
\end{equation}
que no es constante en \(r\). Además, el operador de transición neutra
satisface \(P_0^{\rm TPK}=I\), distinto de las trivializaciones
ciclotómicas \(P_b\); el calendario ordinario \((+1,-1,0)\) hace que
cuatro iteraciones elementales visiten ambos modos aritméticos y el modo
neutro.
\end{proof}

\subsection{Refinamiento mínimo por microinstantes}

\begin{definition}[Estado observable subdividido]
\label{p12-83-c1-s12:def:u006c-estado-subdividido}
Sea
\begin{equation}
 \widetilde{\mathcal Q}_4
 :=\mathcal Q_{\rm obs}\times\{0,1,2,3\},
 \qquad
 \widetilde\pi(q,j)
 :=4\pi_{\rm pos}(q)+j\,a(q).
 \label{p12-83-c1-s12:eq:u006c-estado-refinado}
\end{equation}
Definimos
\begin{equation}
 \widetilde F(q,j):=
 \begin{cases}
  (q,j+1),&0\leq j<3,\\
  (F_{\rm obs}q,0),&j=3,
 \end{cases}
 \qquad
 \iota_4(q):=(q,0).
 \label{p12-83-c1-s12:eq:u006c-dinamica-refinada}
\end{equation}
\end{definition}

\begin{theorem}[Semiconjugación a cuatro pasos]
\label{p12-83-c1-s12:thm:u006c-semiconjugacion}
La inclusión \(\iota_4\) satisface
\begin{equation}
 \boxed{
 \widetilde F^4\circ\iota_4
 =\iota_4\circ F_{\rm obs}.}
 \label{p12-83-c1-s12:eq:u006c-semiconjugacion}
\end{equation}
Las posiciones intermedias son exactamente
\[
 4\pi_{\rm pos}(q)+j\,a(q),
 \qquad j=0,1,2,3.
\]
\end{theorem}

\begin{proof}
Desde \((q,0)\), las primeras tres aplicaciones incrementan sólo el índice
de microinstante. La cuarta aplica \(F_{\rm obs}\) y devuelve el índice a
cero. Esto da
\[
 (q,0)\mapsto(q,1)\mapsto(q,2)\mapsto(q,3)
 \mapsto(F_{\rm obs}q,0).
\]
La fórmula de las posiciones es la definición de
\(\widetilde\pi\). El refinamiento hace visible una subdivisión determinada;
no añade una decisión dinámica.
\end{proof}

\subsection{Categoría subdividida y elevación ciclotómica}

Sea \(\operatorname{sd}_4\mathsf P_{\rm obs}\) la categoría obtenida al
reemplazar cada generador
\(\gamma:q\to F_{\rm obs}q\) por la cadena
\begin{equation}
 (q,0)\to(q,1)\to(q,2)\to(q,3)\to(F_{\rm obs}q,0).
 \label{p12-83-c1-s12:eq:u006c-cadena-subdivision}
\end{equation}
El funtor
\[
 \operatorname{sd}_4:
 \mathsf P_{\rm obs}\longrightarrow
 \operatorname{sd}_4\mathsf P_{\rm obs}
\]
envía cada generador a esa cadena. El funtor de colapso
\[
 \operatorname{col}_4:
 \operatorname{sd}_4\mathsf P_{\rm obs}
 \longrightarrow\mathsf P_{\rm obs}
\]
recompone las cuatro microflechas y satisface
\begin{equation}
 \operatorname{col}_4\circ\operatorname{sd}_4
 =\operatorname{id}_{\mathsf P_{\rm obs}}.
 \label{p12-83-c1-s12:eq:u006c-colapso}
\end{equation}

En una cadena que permanece en la fibra \(b\), las tres primeras
microflechas actúan por \(\rho_b\); la cuarta incorpora el transporte
interfibra si existe cruce dodecafásico. Como \(\rho_b^3=I\), el producto
ordenado de las cuatro acciones coincide con el morfismo ciclotómico de la
flecha original.

\begin{proposition}[Invariancia ciclotómica por subdivisión]
\label{p12-83-c1-s12:prop:u006c-invariancia-ciclotomica}
Existe un levantamiento tipado
\[
 \widetilde{\mathfrak F}_{\rm cyc}:
 \operatorname{sd}_4\mathsf P_{\rm obs}
 \longrightarrow
 \mathbf{Rep}_{\mathbb C}(C_3)\times\mathbf B\mathbb N
\]
tal que, para toda flecha \(\gamma\) de
\(\mathsf P_{\rm obs}\),
\begin{equation}
 \boxed{
 \widetilde{\mathfrak F}_{\rm cyc}
 \bigl(\operatorname{sd}_4\gamma\bigr)
 =
 \bigl(\mathfrak F_{\rm cyc}\circ Q\bigr)(\gamma).}
 \label{p12-83-c1-s12:eq:u006c-invariancia-ciclotomica}
\end{equation}
\end{proposition}

\begin{proof}
Se asignan a las microflechas internas las potencias sucesivas de
\(\rho_b\), y a la cuarta el transporte entre los marcos de las fibras de
origen y destino. La composición de los tres primeros factores es la
identidad porque \(\rho_b^3=I\). El factor restante es exactamente el
morfismo que \(\mathfrak F_{\rm cyc}\circ Q\) asigna a la flecha
colapsada. La categoría subdividida es libre sobre esas microflechas, por lo
que la asignación se extiende de manera única.
\end{proof}

\subsection{\(3\) operaciones cuaternarias no intercambiables}

La arquitectura contiene tres operaciones distintas:
\begin{enumerate}
 \item \(\mathcal S_4\), sustitución cardinal de objetos y caminos;
 \item \(D_4\), multiplicación por cuatro en
 \(\Theta_{108}\);
 \item \(\operatorname{sd}_4\), subdivisión functorial de una transición
 observable.
\end{enumerate}
Sus compatibilidades están dadas, respectivamente, por
\[
 4^3=1+7\cdot9,\qquad
 \Theta D_4=4\Theta,\qquad
 \widetilde F^4\iota_4=\iota_4F_{\rm obs}.
\]
Compartir el entero cuatro no convierte ninguna en otra.

\begin{criterion}[Criterios de refutación]
\label{p12-83-c1-s12:crit:u006c-refutacion}
La construcción queda refutada si falla cualquiera de las identidades
\eqref{p12-83-c1-s12:eq:u006c-d4-total},
\eqref{p12-83-c1-s12:eq:u006c-cociclo-d4},
\eqref{p12-83-c1-s12:eq:u006c-d4-cubo},
\eqref{p12-83-c1-s12:eq:u006c-semiconjugacion} o
\eqref{p12-83-c1-s12:eq:u006c-invariancia-ciclotomica}; también si una realización
observable satisface las cuatro condiciones del
\cref{p12-83-c1-s12:thm:u006c-obstruccion} sin conservar un registro equivalente al
microinstante.
\end{criterion}


% Unidad U006. Lectores, odómetro y calendario; redacción autónoma.

\section{Lectores modulares, odómetro y calendario dodecafásico}
\label{p12-83-c1-s13:chap:tpk-lectores-odometro-calendario}

El estado holonómico integral del TPK no se identifica con ninguna de sus lecturas.
Esta unidad define los lectores módulo nueve, módulo tres y módulo mil, el
odómetro fase--sector, el cociente predictivo de treinta y seis estados y la
extensión cíclica de ciento ocho posiciones. Cada construcción declara qué
publica y qué memoria conserva.

\subsection{Esquema tipado de los lectores modulares del TPK: dominio, codominio, publicación y memoria conservada}

\begin{definition}[Ficha operatoria]
\label{p12-83-c1-s13:def:tpk-ficha-operatoria}
Un operador canónico del TPK se registra mediante una séxtupla
\[
 \mathcal O=(D_{\mathcal O},C_{\mathcal O},f_{\mathcal O},
 I_{\mathcal O},L_{\mathcal O},M_{\mathcal O}),
\]
donde \(D_{\mathcal O}\) es el dominio, \(C_{\mathcal O}\) el codominio,
\(f_{\mathcal O}\) la regla, \(I_{\mathcal O}\) los invariantes preservados,
\(L_{\mathcal O}\) la información publicada u olvidada y
\(M_{\mathcal O}\) la memoria que permanece en el estado holonómico integral.
\end{definition}

Dos símbolos con igual cardinal o igual periodo no representan el mismo
operador salvo que exista una equivalencia que entrelace sus dominios,
reglas e invariantes. Esta convención impedirá confundir el cociente
predictivo \(C_{36}^{\rm pred}\), la frontera \(R_{36}\) y un defecto
cardinal de valor treinta y seis.

\subsection{Lector nonádico del estado holonómico integral: fase, transición y retorno con memoria}

Para \(n\geq1\), definimos
\[
 \rho_9(n)=1+((n-1)\bmod9),
 \qquad
 q_9(n)=\left\lfloor\frac{n-1}{9}\right\rfloor.
\]
Entonces
\[
 n=9q_9(n)+\rho_9(n).
\]
La ficha de \(\rho_9\) es
\[
 \bigl(
 \mathbb Z_{>0},\mathcal D_9,\rho_9,
 n\equiv\rho_9(n)\pmod9,
 \text{fase positiva},
 q_9(n)
 \bigr).
\]
La raíz digital positiva coincide numéricamente con \(\rho_9\), pero su
semántica es distinta: una se presenta como sección del cociente; la otra,
como suma iterada de cifras.

\subsection{Lector trítico del estado orientado: orden temporal y transporte de acarreo}

La división balanceada se escribe
\[
 n=r_3(n)+3c_3(n),
 \qquad r_3(n)\in\mathsf{TRIT},\quad c_3(n)\in\mathbb Z.
\]
Su ficha es
\[
 \bigl(
 \mathbb Z,\mathsf{TRIT}\times\mathbb Z,
 n\mapsto(r_3(n),c_3(n)),
 \operatorname{rec}_3(r,c)=r+3c,
 r_3,
 c_3
 \bigr).
\]
La publicación puede mostrar sólo \(r_3\), pero el TPK conserva
\(c_3\) si una composición posterior requiere reconstrucción.

\subsection{Lector posicional por bloques en base \(10^3\)}

La división euclídea define
\[
 N=1000q_{1000}(N)+r_{1000}(N),
 \qquad 0\leq r_{1000}(N)<1000.
\]
Su ficha es
\[
 \bigl(
 \mathbb Z,\mathbb Z\times\{0,\ldots,999\},
 N\mapsto(q_{1000},r_{1000}),
 \operatorname{rec}_{1000}(q,r)=1000q+r,
 r_{1000},
 q_{1000}
 \bigr).
\]

\begin{proposition}[Concatenación y lectura residual]
\label{p12-83-c1-s13:prop:tpk-concatenacion-base-mil}
Si \(u_1,\ldots,u_m\in\{0,\ldots,999\}\) y
\[
 D_m=1000D_{m-1}+u_m,
 \qquad D_0=0,
\]
entonces
\[
 D_m=\sum_{j=1}^m u_j1000^{m-j}.
\]
Además,
\[
 D_m\equiv\sum_{j=1}^m u_j\pmod9
\]
porque \(1000\equiv1\pmod9\).
\end{proposition}

\begin{proof}
La primera igualdad se obtiene por inducción en \(m\). La segunda procede
al reducirla módulo nueve.
\end{proof}

La congruencia no reconstruye el orden de los bloques. El registro genealógico conserva la
secuencia \((u_1,\ldots,u_m)\); el lector residual sólo publica su suma
módulo nueve.

\subsection{Categoría odométrica del TPK: objetos cilíndricos y composición por truncamiento}

\begin{definition}[Odómetro nonádico--dodecafásico]
\label{p12-83-c1-s13:def:tpk-odometro-9-12}
La categoría \(\mathcal O_{9,12}\) tiene objetos
\[
 (b,r)\in C_{12}\times\{0,\ldots,8\}.
\]
Una flecha de longitud \(n\in\mathbb N\) es
\[
 (b,r)\xrightarrow{\ n\ }
 \left(
  b+\left\lfloor\frac{r+n}{9}\right\rfloor,
  r+n\bmod9
 \right).
\]
La composición suma longitudes.
\end{definition}

Sea
\[
 \kappa_r(n)=\left\lfloor\frac{r+n}{9}\right\rfloor.
\]

\begin{lemma}[Cociclo del odómetro]
\label{p12-83-c1-s13:lem:tpk-cociclo-odometro}
Si \(r_1=r+n_1\bmod9\), entonces
\[
 \kappa_r(n_1+n_2)
 =\kappa_r(n_1)+\kappa_{r_1}(n_2).
\]
\end{lemma}

\begin{proof}
Escribimos \(r+n_1=9\kappa_r(n_1)+r_1\), con
\(0\leq r_1<9\). Al añadir \(n_2\) y dividir de nuevo entre nueve se
obtiene la identidad.
\end{proof}

Para un estado \(x\) cuya base contiene \(\theta\), definimos
\[
 Q_0(x)=(\beta(\theta),\widetilde\theta\bmod9).
\]
Una flecha TPK tiene longitud igual al número de generadores observables que
la componen.

\begin{theorem}[Funtor fase--longitud]
\label{p12-83-c1-s13:thm:tpk-funtor-fase-longitud}
La asignación que envía un estado \(x\) a \(Q_0(x)\) y una flecha de
longitud \(n\) a la flecha odométrica de longitud \(n\) define un funtor
\[
 Q:\mathsf{TPK}\longrightarrow\mathcal O_{9,12}.
\]
\end{theorem}

\begin{proof}
Cada generador aumenta \(\theta\) en una unidad. Tras \(n\) pasos, el
residuo de fase aumenta en \(n\) módulo nueve y el sector aumenta en
\(\lfloor(r+n)/9\rfloor\). La identidad tiene longitud cero y el
\cref{p12-83-c1-s13:lem:tpk-cociclo-odometro} prueba la compatibilidad con la composición.
\end{proof}

Después de nueve pasos,
\[
 (b,r)\longmapsto(b+1,r).
\]
Después de ciento ocho,
\[
 (b,r)\longmapsto(b,r).
\]
Estas identidades pertenecen al cociente odométrico. No afirman que hayan
retornado cilindro, frontera, registro genealógico o memoria vertical.

\subsection{Cociente mínimo de memoria predictiva}

Sea
\[
 X_{9,12}=\{(a,m):a\in\mathbb Z/9\mathbb Z, 0\leq m<12\}
\]
con sucesor
\[
 S(a,m)=
 \begin{cases}
  (a,m+1),&m<11,\\
  (a+1\bmod9,0),&m=11.
 \end{cases}
\]
Definimos los observables
\[
 \beta_{\rm pred}(a,m)=4a\pmod{12},
 \qquad
 d_{\rm pred}(a,m)=1+(m+\beta_{\rm pred}(a,m)\pmod{12}).
\]
Para un observable \(q\), la palabra futura inclusiva de horizonte \(h\)
es
\[
 Q_h^q(x)=(q(x),q(Sx),\ldots,q(S^hx)).
\]

\begin{theorem}[Cociente predictivo mínimo]
\label{p12-83-c1-s13:thm:tpk-cociente-predictivo-36}
Los observables \(\beta_{\rm pred}\), \(d_{\rm pred}\) y el par
\((\beta_{\rm pred},d_{\rm pred})\) producen el mismo cociente estable
\[
 \pi_{36}(a,m)=(a\bmod3,m).
\]
Tiene treinta y seis estados y fibras de cardinal tres. Los horizontes
mínimos son \(11\), \(8\) y \(0\), respectivamente.
\end{theorem}

\begin{proof}
Para un sistema finito \((X,T)\), sea
\[
 E_h=\operatorname{Eq}(Q_h^q).
\]
Entonces
\[
 E_{h+1}=E_h\cap(T\times T)^{-1}(E_h).
\]
La cadena decreciente estabiliza y su límite es la mayor equivalencia
\(T\)-compatible contenida en \(\operatorname{Eq}(q)\). En el sistema
declarado, la enumeración exacta produce el número de clases siguiente:
\[
\begin{array}{c|rrrrrrrrrrrr}
h&0&1&2&3&4&5&6&7&8&9&10&11\\\hline
\#(X/E_h^{\beta})&3&6&9&12&15&18&21&24&27&30&33&36
\end{array},
\]
y
\[
\begin{array}{c|rrrrrrrrr}
h&0&1&2&3&4&5&6&7&8\\\hline
\#(X/E_h^{d})&12&15&18&21&24&27&30&33&36.
\end{array}
\]
El par ya distingue las treinta y seis clases en horizonte cero. En todos
los casos, las clases estables conservan \(m\) y \(a\bmod3\); las tres
posibilidades de \(a\) con igual residuo módulo tres forman cada fibra.
\end{proof}

Denotamos este objeto por \(C_{36}^{\rm pred}\). No es la frontera
\(R_{36}\), que contendrá tres canales ordenados de seis trits, ni un censo
de treinta y seis elementos de otra construcción.

\subsection{Extensión central del calendario}

Definimos
\[
 \iota:C_{12}\longrightarrow C_{108},
 \qquad \iota([b]_{12})=[9b]_{108},
\]
y
\[
 p:C_{108}\longrightarrow C_9,
 \qquad p([\theta]_{108})=[\theta]_9.
\]
Se obtiene la sucesión exacta
\begin{equation}
 0\longrightarrow C_{12}
 \xrightarrow{\ \iota\ }C_{108}
 \xrightarrow{\ p\ }C_9
 \longrightarrow0.
\label{p12-83-c1-s13:eq:tpk-extension-12-108-9}
\end{equation}
Una sección de conjuntos \(s([r]_9)=[r]_{108}\),
\(0\leq r<9\), induce el cociclo
\[
 c(r,r')=
 \left[\left\lfloor\frac{r+r'}{9}\right\rfloor\right]_{12}.
\]

\begin{proposition}[La extensión no se escinde]
\label{p12-83-c1-s13:prop:tpk-extension-no-escindida}
La sucesión \cref{p12-83-c1-s13:eq:tpk-extension-12-108-9} no es isomorfa a la extensión
producto \(C_{12}\times C_9\).
\end{proposition}

\begin{proof}
El exponente de \(C_{108}\) es \(108\). El exponente de
\(C_{12}\times C_9\) es
\(\operatorname{mcm}(12,9)=36\). Dos grupos finitos isomorfos tienen el
mismo exponente; por tanto, no son isomorfos.
\end{proof}

\begin{theorem}[Necesidad y periodo exacto de ciento ocho posiciones]
\label{p12-83-c1-s13:thm:tpk-periodo-exacto-108}
La traslación \(\theta\mapsto\theta+1\) sobre \(C_{108}\) tiene orden
exactamente \(108\). En consecuencia, una iteración \(n\) retorna el
calendario completo si y sólo si \(108\mid n\).
\end{theorem}

\begin{proof}
La \(n\)-ésima potencia envía \([\theta]\) a \([\theta+n]_{108}\). Es
la identidad para todo \(\theta\) exactamente cuando \([n]_{108}=0\), es
decir, cuando \(108\mid n\).
\end{proof}

En particular, \(1008\) no es otro retorno del calendario:
\[
 1008\equiv36\pmod{108}.
\]
Tras \(1008\) pasos, la fase módulo nueve retorna, pero la coordenada de
\(C_{108}\) se ha desplazado treinta y seis posiciones. Naturalmente,
\(216,324,\ldots\) también retornan el calendario, pero son repeticiones de
la vuelta fundamental de longitud \(108\).

\begingroup\fontsize{10.2}{11.7}\selectfont
\setlength{\parskip}{2pt}\setlength{\abovedisplayskip}{5pt}\setlength{\belowdisplayskip}{5pt}
\subsection{Memoria vertical no acotada}

Sea \(q_{\rm mem}\in\mathbb Z_{\geq0}\) el número no acotado de vueltas
nonádicas y sea
\[
 \beta=[q_{\rm mem}]_{12}.
\]
Entonces una vuelta de nueve pasos satisface
\[
 q_{\rm mem}\longmapsto q_{\rm mem}+1,
 \qquad
 \beta\longmapsto\beta+1\pmod{12}.
\]
Después de ciento ocho pasos,
\[
 \beta\longmapsto\beta,
 \qquad
 q_{\rm mem}\longmapsto q_{\rm mem}+12.
\]
El calendario retorna; la memoria vertical no.

\begin{corollary}[Jerarquía de retornos]
\label{p12-83-c1-s13:cor:tpk-jerarquia-retornos}
En el cociente de fase,
\(9\) pasos constituyen una vuelta. Los morfismos compuestos
\[
 \mathcal K_{27}=F_{\rm obs}^{27},
 \qquad
 F_{\rm obs}^{54},
 \qquad
 F_{\rm obs}^{108}
\]
contienen, respectivamente, tres, seis y doce vueltas nonádicas. Ninguno
borra por definición el registro genealógico o la frontera.
\end{corollary}

\subsection{Sistema mínimo de tipos del estado holonómico integral y compatibilidades entre sus lectores}

La tabla siguiente responde de manera explícita a qué objetos internos
existen ya en el TPK antes de introducir semillas:
\[
\begin{array}{p{.23\textwidth}p{.27\textwidth}p{.38\textwidth}}
\toprule
Objeto u operador&Dominio y codominio&Función\\\midrule
\(X_9\)&\((\mathbb Z/9)^2\)&soporte APP\\
\(T_N,T_E,T_S,T_O\)&\(X_9\to X_9\)&traslaciones cardinales\\
\(M_{\rm ph}\)&\(\mathcal D_9\to\mathsf{TRIT}\)&selección de hoja\\
\(F_{\rm obs}\)&\(\mathcal Q_{\rm obs}\to\mathcal Q_{\rm obs}\)&sucesor observable\\
\((\rho_9,q_9)\)&\(\mathbb Z_{>0}\to\mathcal D_9\times\mathbb Z_{\ge0}\)&fase y altura\\
\((r_3,c_3)\)&\(\mathbb Z\to\mathsf{TRIT}\times\mathbb Z\)&orientación y carry\\
\((q_{1000},r_{1000})\)&\(\mathbb Z\to\mathbb Z\times\{0,\ldots,999\}\)&bloque decimal y cociente\\
\(Q\)&\(\mathsf{TPK}\to\mathcal O_{9,12}\)&fase y longitud\\
\(C_{36}^{\rm pred}\)&cociente de \(X_{9,12}\)&memoria predictiva mínima\\
\(C_{108}\)&calendario cíclico&extensión no escindida\\
\(q_{\rm mem}\)&\(\mathbb Z_{\ge0}\)&memoria vertical no acotada\\
\(\operatorname{Ext}_r\)&
\(\operatorname{Ext}_r\subseteq\mathcal F_q\times\mathcal F_{q'}\)&
prolongaciones compatibles\\
\bottomrule
\end{array}
\]

\subsection{Obstrucciones de los lectores y del calendario del TPK: cocientes euclídeos, \(C_{36}^{\mathrm{pred}}\), extensión \(C_{108}\) y memoria vertical}

\begin{criterion}[Criterios de refutación de lectores y calendario]
La unidad queda refutada si:
\begin{enumerate}
 \item un lector residual se presenta como transición completa;
 \item se omite el cociente de alguna división euclídea utilizada después;
 \item se identifica \(C_{36}^{\rm pred}\) con \(R_{36}\);
 \item se afirma que la extensión \(C_{108}\) es el producto
       \(C_{12}\times C_9\);
 \item se declara retorno completo para un número no divisible por \(108\);
 \item el retorno de \(\beta\) fuerza \(q_{\rm mem}=0\);
 \item \(\mathcal K_{27}\) se presenta como un reloj distinto y no como
       una potencia del transporte.
\end{enumerate}
\end{criterion}

\section{Estructura operatorial completa de APP y jerarquía aritmético--geométrica}

El desarrollo operatorial de APP fija la jerarquía
\(4\to2\to6\to54\), los proyectores aritméticos, el bloque central y la
geometría local de caminos. Estas consecuencias permanecen subordinadas al
núcleo formal ya construido y no introducen un segundo origen de APP.

La unidad siguiente construirá la subdivisión cuádruple, el funtor
ciclotómico y los canales internos de noventa y ciento veinte posiciones.
Sólo después se introducirá el espacio exhaustivo de semillas.
 

\par\endgroup
\chapter{Arquitectura operatoria de la selección, el transporte y la actualización}
\section{Estructura categorial del TPK: estados holonómicos integrales, morfismos de
transporte y clasificación operatoria}
\label{sec:arbol-operatorio-tpk-rev11}
\index{TPK!clasificación operatoria}
\index{selector de fase!suboperador del TPK}

El TPK es el contenedor categórico que transporta sobre APP los estados
orientados de TRIT. Su definición precede al inventario de realizaciones:
primero se fijan objetos, morfismos, composición y clases operatorias; sólo
después se introducen los operadores que realizan cada clase.

\HMTClaim{HMT-R-TPK-ENRICHED-STATE}
\begin{definition}[TPK como categoría sobre la base observable con familia de fibras]
\label{def:estado-enriquecido-tpk-rev11}
Sea \(\mathsf P_{\rm obs}\) la categoría de rutas observables sobre el grafo
de Cayley de APP. Para cada configuración \(q\) se fija la fibra
\[
 \mathcal F_q=
 \mathsf O_q\times\mathsf H_q\times\mathsf C_q\times
 \mathsf S_q\times\mathsf B_q\times\mathsf\Lambda_q,
\]
cuyos factores conservan, respectivamente, orientación y hoja; residuo y
cociente; acarreo; firma; cilindro y frontera; y registro de transporte. Un
objeto del TPK es un par \(x=(q,f)\), con \(f\in\mathcal F_q\). Un morfismo
\(\Gamma:x\to y\) es una ruta \(\gamma:q\to q'\) de
\(\mathsf P_{\rm obs}\) junto con el transporte de fibra que preserva las
relaciones de extensión declaradas entre \(f\) y \(f'\). La identidad es la
ruta vacía con transporte identidad y la composición es la concatenación de
rutas con composición de los transportes de fibra. La proyección
\[
 p:\mathcal X_{\rm TPK}^{\rm enr}\longrightarrow\mathsf P_{\rm obs}
\]
envía cada estado a su configuración observable y cada morfismo a la ruta que
transporta. Esta definición fija una categoría sobre la base y una familia de
fibras; no presupone la existencia de levantamientos cartesianos para toda
flecha de \(\mathsf P_{\rm obs}\).
\end{definition}

\begin{definition}[Clasificación ontológico--operatoria del TPK]
\label{def:categorias-operatorias-tpk-rev11}
La estructura interna del TPK se organiza en ocho clases diferenciadas por tipo de
dominio, codominio y función:
\begin{enumerate}
 \item \textbf{Soporte y estado.} La base es \(\mathsf P_{\rm obs}\) y el
 portador es \(\mathcal X_{\rm TPK}^{\rm enr}\to\mathsf P_{\rm obs}\). Esta
 clase contiene objetos; las siete siguientes contienen mapas entre objetos
 o lectores de esos objetos.
 \item \textbf{Selección interna.} El selector
 \[
 M_{\rm ph}:\{1,\ldots,9\}\longrightarrow\{+1,-1,0\}
 \]
 activa la evaluación aditiva, la multiplicativa o el paso de umbral. Actúa
 sobre la coordenada de hoja y no modifica por sí mismo la posición en APP.
 \item \textbf{Transporte.} Para cada dirección cardinal \(d\),
 \(T_d:X_9\to X_9\), \(p\mapsto p+\delta(d)\), transporta la posición. La
 evolución \(F_{\rm obs}:\mathcal Q_{\rm obs}\to\mathcal Q_{\rm obs}\)
 compone selección, traslación, orientación y avance de fase.
 \item \textbf{Lectores y cocientes.} \(\rho_9\) publica la clase nonádica;
 \((r_3,c_3)\) conserva residuo trítico y cociente; y
 \((q_{1000},r_{1000})\) separa bloque decimal y acarreo. Son mapas de
 lectura, no transiciones del estado.
 \item \textbf{Memoria y frontera.} Los transportes de
 \(\mathsf H_q,\mathsf C_q,\mathsf S_q,\mathsf B_q\) y
 \(\mathsf\Lambda_q\) actualizan cociente, acarreo, firma, cilindro, frontera
 y registro. Su proyección visible puede coincidir aunque los estados
 holonómicos integrales difieran.
 \item \textbf{Periodicidad y retorno.} La extensión
 \[
 0\longrightarrow C_{12}\longrightarrow C_{108}
 \longrightarrow C_9\longrightarrow0
 \]
 tipa sector, calendario y fase visible. El retorno nonádico restaura la
 fase; la periodicidad de ciento ocho iteraciones restaura el calendario;
 ninguno borra
 por definición la memoria de frontera. El lazo
 \(\mathcal K_{27}=F_{\rm obs}^{27}\) es un morfismo compuesto, no otro
 reloj.
 \item \textbf{Prolongación.} Las relaciones
 \(\operatorname{Ext}_r\subseteq\mathcal F_q\times\mathcal F_{q'}\)
 seleccionan continuaciones compatibles. Los cilindros y su sistema inverso
 publican supervivencia finita y, bajo las hipótesis de compatibilidad
 declaradas, una sección global.
 \item \textbf{Salidas.} Los funtores de evaluación contraen el mismo estado
 holonómico integral hacia la salida arquimediana, la salida de incidencia
 excepcional o los clasificadores del atlas. Estas salidas son codominios de
 lectura; no constituyen nuevos contenedores junto a TPK.
\end{enumerate}
\end{definition}

Cada operador que aparece a continuación queda determinado por la séxtupla
\((D,C,f,I,L,M)\): dominio \(D\), codominio \(C\), regla \(f\), invariantes
preservados \(I\), información publicada o perdida \(L\) y memoria conservada
\(M\). La igualdad de cardinales, periodos o nombres históricos no identifica
dos operadores con séxtuplas distintas. Una vez fijada esta
clasificación, el siguiente diagrama resume cinco agrupaciones funcionales
de realizaciones que refinan las ocho clases operatorias del TPK:
\[
\begin{tikzpicture}[
  node distance=7mm and 9mm,
  every node/.style={font=\small,align=center},
  root/.style={draw=hmtblue,very thick,rounded corners,fill=hmtlightblue,
    inner sep=5pt},
  op/.style={draw=hmtblue,rounded corners,fill=white,inner sep=4pt},
  sub/.style={draw=hmtgray,rounded corners,fill=hmtpaper,inner sep=3pt},
  edge/.style={-{Latex[length=2mm]},semithick,hmtblue}
]
\node[root] (tpk) {TPK\\$\mathcal X_{\rm TPK}^{\rm enr}$};
\node[op,below left=of tpk] (sel) {selección\\de lectura};
\node[op,below=of tpk] (trans) {transporte\\y retorno};
\node[op,below right=of tpk] (lift) {elevación\\y memoria};
\node[sub,below=of sel] (mov) {$M_{\rm ph}$};
\node[sub,below=of trans] (card) {$T_d,\ \mathcal K_{27},\ F_{\rm obs}$};
\node[sub,below=of lift] (fib) {$\mathsf H,\mathsf C,\mathsf S,
  \mathsf B,\mathsf\Lambda$};
\node[op,below=17mm of trans] (clock) {calendario y registro\\
  $C_{12}\hookrightarrow C_{108}\twoheadrightarrow C_9$, $\Theta_{108}$};
\node[op,below=of clock] (ext) {prolongación, cilindros, frontera\\
  y filtración de supervivencia};
\draw[edge] (tpk) -- (sel);
\draw[edge] (tpk) -- (trans);
\draw[edge] (tpk) -- (lift);
\draw[edge] (sel) -- (mov);
\draw[edge] (trans) -- (card);
\draw[edge] (lift) -- (fib);
\draw[edge] (mov) |- (clock);
\draw[edge] (card) -- (clock);
\draw[edge] (fib) |- (clock);
\draw[edge] (clock) -- (ext);
\end{tikzpicture}
\]
El diagrama fija una contención, no una sucesión de teorías. En particular,
el selector interno, las traslaciones cardinales, la observación
estroboscópica y los lectores modulares conservan los tipos fijados en la
\cref{def:categorias-operatorias-tpk-rev11}. El TPK coordina sus composiciones
y conserva su procedencia en un mismo estado holonómico integral.

\HMTClaim{HMT-R-TPK-OPERATOR-SEPARATION}
\begin{proposition}[Separación de selección, movimiento y observación]
\label{prop:tpk-separacion-operadores-rev11}
El selector \(M_{\rm ph}\), el transporte cardinal \(T_d\) y el muestreo
\(\operatorname{Obs}_3\) son morfismos no intercambiables: el primero elige
la evaluación activa; el segundo modifica la posición y el enrollamiento; el
tercero publica una subsucesión temporal. Su composición dentro del TPK
determina una ruta observable, mientras la fibra \(\mathcal F_q\) conserva
los datos necesarios para distinguir prolongaciones con igual extremo.
\end{proposition}

\begin{proof}
Los tres morfismos tienen codominios diferentes: \(M_{\rm ph}\) toma valores
en el conjunto ternario de activación, \(T_d\) es un automorfismo del toro y
\(\operatorname{Obs}_3\) actúa sobre sucesiones. Una identificación entre dos
de ellos violaría el tipado de dominio o codominio. La extensión por fibras
incorpora, además, residuo--cociente, acarreo, frontera y registro, de modo
que la proyección al extremo del camino es un mapa de olvido explícito.
\end{proof}

\begin{statusbox}[title={Nomenclatura canónica}]
La denominación pública del contenedor es \emph{TPK}. Cuando sea necesario
nombrar el dato transportado se empleará \emph{estado holonómico integral del TPK} o
\(\mathcal X_{\rm TPK}^{\rm enr}\). El morfismo \(M_{\rm ph}\) designa
exclusivamente la selección interna de fase.
\end{statusbox}
 \section{Operadores fundamentales del TPK: dominios, acciones, invariantes y
memoria}
\label{sec:inventario-operatorio-tpk-rev11}
\index{TPK!inventario operatorio}

La \cref{def:categorias-operatorias-tpk-rev11} ha fijado primero las ocho
clases ontológico--operatorias: soporte y estado; selección; transporte;
lectores y cocientes; memoria y frontera; periodicidad y retorno;
prolongación; y salidas. Esta sección define sus realizaciones. Para cada
operador se declara la séxtupla \((D,C,f,I,L,M)\) de dominio, codominio, regla,
invariantes, información publicada o perdida y memoria conservada. El censo
nominal posterior agrupa símbolos ya definidos por esos seis datos y no
sustituye sus dominios ni sus reglas.

\begin{definition}[Ficha tipada de un operador del TPK]
Un operador canónico es una séxtupla
\[
 \mathcal O=(D_{\mathcal O},C_{\mathcal O},f_{\mathcal O},
 I_{\mathcal O},L_{\mathcal O},M_{\mathcal O}),
\]
donde \(f_{\mathcal O}:D_{\mathcal O}\to C_{\mathcal O}\) es la regla de
actuación, \(I_{\mathcal O}\) es la familia de relaciones que conserva,
\(L_{\mathcal O}\) declara las coordenadas publicadas y la relación de
equivalencia inducida por el olvido, y \(M_{\mathcal O}\) enumera los datos
que permanecen en el estado holonómico integral después de la publicación. Dos
realizaciones representan el mismo operador únicamente cuando existe una
equivalencia tipada que entrelaza sus reglas y sus invariantes; compartir
periodo, cardinal o salida visible no basta.
\end{definition}

\subsection{Lectores modulares construidos}

La reducción nonádica se fija por el representante publicado
\begin{equation}
 \rho_9(n)=1+((n-1)\bmod 9)
 =n-9\left\lfloor\frac{n-1}{9}\right\rfloor,
 \label{eq:rho9-construida-tpk-rev11}
\end{equation}
de modo que la clase nula se escribe como \(9\). Para \(n\geq1\), la raíz
digital satisface
\[
 \operatorname{dr}_9(n)=\rho_9(n),
\]
pero las dos expresiones conservan semánticas distintas: la primera es una
sección del cociente; la segunda publica la suma iterada de cifras. El levantamiento
trítico escribe de manera única
\begin{equation}
 n=r_3(n)+3c_3(n),
 \qquad r_3(n)\in\{-1,0,+1\},\quad c_3(n)\in\mathbb Z,
 \label{eq:lift-mod3-tpk-rev11}
\end{equation}
y conserva simultáneamente residuo orientado y cociente de acarreo. Por
último, el lector de bloques decimales es la división euclídea
\begin{equation}
 N=1000q_{1000}(N)+r_{1000}(N),
 \qquad 0\leq r_{1000}(N)<1000,
 \label{eq:lector-mod1000-tpk-rev11}
\end{equation}
por lo que base \(1000\), residuo módulo \(1000\) y acarreo \(q_{1000}\)
son tres datos relacionados y no un mismo objeto.

\subsection{Clasificación de las realizaciones operatorias por dominio y función}

Las definiciones anteriores gobiernan la lectura de la tabla. Sus catorce
agrupaciones funcionales refinan las ocho clases operatorias y resumen las
realizaciones que comparten función; una clase puede, por tanto, ocupar más
de una fila. La tabla no opera como definición sustitutiva. Los símbolos que
comparten un cardinal permanecen separados por dominio y codominio.

\begingroup
\footnotesize
\renewcommand{\arraystretch}{1.06}
\begin{longtable}{@{}p{.21\textwidth}p{.31\textwidth}p{.40\textwidth}@{}}
\toprule
Agrupación funcional & Operadores y espacios & Dominio, acción e invariantes \\
\midrule
\endfirsthead
\toprule
Agrupación funcional & Operadores y espacios & Dominio, acción e invariantes \\
\midrule
\endhead
Estado y memoria de ruta &
\(\mathcal X_{\rm TPK}^{\rm enr}\), fibra \(\mathcal F_q\), memoria de ruta &
Conservan posición, hoja, orientación, residuo, cociente, acarreo, firma,
frontera, cilindro, prefijo, genealogía y memoria. \\

Selección de hoja &
\(M_{\rm ph}:\{1,\ldots,9\}\to\{-1,0,+1\}\) &
Activa la evaluación aditiva, la multiplicativa o el paso de umbral; no
mueve por sí mismo el estado. \\

Transporte espacial &
\(T_N,T_E,T_S,T_O\), motor reversible de dos cursores,
\(F_{\rm obs}\) &
Realiza las traslaciones del grafo de Cayley, actualiza hoja y orientación y
conserva cada transición de la ruta. \\

Lectores modulares &
\(\rho_9\), \((r_3,c_3)\), \((q_{1000},r_{1000})\) &
Publican fase nonádica, orientación trítica y bloque decimal sin borrar los
cocientes de levantamiento. \\

Bloque y memoria &
\(\mathcal K_{27}=F_{\rm obs}^{27}\), filtro de memoria
\(\mathcal N_{27}\) &
El primero compone veintisiete transportes; el segundo selecciona los
eventos de memoria. Igual cardinal no implica igual operador. \\

Estado dodecafásico &
\(\mathcal R_{12}\), centros \(\theta_m=30m-10\) grados,
refinamiento \(\mathcal R_{120}\) &
Organiza las doce fases, su oposición y el refinamiento angular en incrementos de
tres grados. \\

Periodicidad y profundidades &
\(C_{12}\hookrightarrow C_{108}\twoheadrightarrow C_9\),
\(\Pi_{108}^{\rm cur}\), \(w_{108}\), \(w_{1080}\),
\(\Gamma_0(1080)\) &
Distingue la periodicidad total de orden \(108\), la proyección de cursores, las
palabras de 108 y 1080 trits y el nivel modular 1080. \\

Canales (90/120) &
máscara \(C_m\), \(S_{90},S_{120}\), gramática orbital,
supertraza e invariantes \(\nu_{120},\nu_{270}\) &
Separa las dos genealogías, sus colas geométricas, sus pesos y las
coordenadas que extraen del estado. \\

Bidireccionalidad &
\(P_\pm\), inversión de ruta, \(\operatorname{Ext}_\pm\),
\(N_\pm\), \(\beta\), \(C_M=-\operatorname{rev}\) &
Realiza, respectivamente, intercambio espectral, ida y retorno espacial,
prolongación futura/pasada, valencia, balance y conjugación de palabra. \\

Emisión y coinducción &
\(L_0,L_1\), \(R_{36}\), \(G_9\), partición \(3^6=729\),
\(\varprojlim\operatorname{Cyl}_m\) &
Eleva la semilla, abre la frontera y enumera exhaustivamente las
prolongaciones finitas. El límite numérico y el levantamiento enriquecido
pertenecen a tipos distintos. \\

Lectores numéricos &
\(\mathscr C_\pi\), \(\mathscr P_e\), \(F_{\rm av}\) &
Construyen, por horquillas racionales, recurrencia factorial y autoescala de
Fibonacci, las secciones publicadas de \(\pi,e,\varphi\). \\

Proyección excepcional &
levantamiento de Witt, enumerador de peso y sombra no selectora &
Transporta el mismo estado a incidencia ternaria sin intervenir en la
selección arquimediana de cifras. \\

Atlas y realización &
firma de ruta, atlas (81/56/13/324), involución \(C_M\), centro
electrónico y operadores de fibra &
Convierte el estado de ruta en clasificadores, caracteres y operadores antes de
aplicar cualquier carta dimensional. \\

Validación finita &
censos exhaustivos, ablaciones e independencia causal &
Comprueba cobertura y unicidad en la profundidad calculada, y verifica que
ningún prefijo objetivo interviene en el mapa generador. \\
\bottomrule
\end{longtable}
\endgroup

\Needspace{46\baselineskip}
\subsection{Índice funcional de los operadores canónicos}

La clasificación anterior organiza las realizaciones por función. La tabla siguiente
reúne los operadores después de la definición categorial y de los lectores
construidos. Cada fila declara dominio, codominio, regla e invariantes; los
símbolos permiten reconocer una misma definición a lo largo de sus
distintas realizaciones sin identificar operadores de tipos diferentes.

\begingroup
\footnotesize
\renewcommand{\arraystretch}{1.06}
\enlargethispage{2\baselineskip}
\begin{longtable}{@{}r >{\raggedright\arraybackslash}p{.42\textwidth}
  >{\raggedright\arraybackslash}p{.22\textwidth}
  >{\raggedright\arraybackslash}p{.22\textwidth}@{}}
\toprule
\# & Operador canónico & Símbolo & Capa tipada \\
\midrule
\endfirsthead
\toprule
\# & Operador canónico & Símbolo & Capa tipada \\
\midrule
\endhead
1 & Célula aritmética aditivo--multiplicativa de orden nueve
  & \(\mathcal A_9\) & soporte aritmético \\
2 & Descomposición residual y cociente nonádicos
  & \((\rho_9,q_9)\) & reducción y elevación \\
3 & Reducción ternaria orientada
  & \(\rho_3^{\rm or}\) & reducción y elevación \\
4 & Lector posicional ternario--decimal de base mil
  & \(\iota_{1000}\) & reducción y elevación \\
5 & Acción cardinal orientada sobre el toro aritmético
  & \(T_d,\ d\in\{N,E,S,O\}\) & transporte cardinal \\
6 & Realimentación local aditivo--multiplicativa
  & \(F_{\rm loc}\) & transporte cardinal \\
7 & Emisor mínimo de dos cursores
  & \(T_{\min}\) & emisión finita \\
8 & Acoplador decimal de firmas aritméticas
  & \(G\) & emisión finita \\
9 & Proyección ternaria del catálogo
  & \(\Psi_6\) & emisión finita \\
10 & Sobreyectividad local del transporte cardinal
  & \(\mathcal R_3\) & transporte cardinal \\
11 & Lazo de transporte de veintisiete iteraciones
  & \(\mathcal K_{27}\) & transporte cardinal \\
12 & Estado holonómico integral del núcleo topológico de fase
  & \(\mathcal X_{\rm TPK}^{\rm enr}\) & estado holonómico integral \\
13 & Transición de estado Hensel
  & \(\mathcal H\) & estado holonómico integral \\
14 & Transición exacta de cilindros ternario--decimales
  & \(\mathcal C_{3,1000}\) & estado holonómico integral \\
15 & Firma conjunta de frontera
  & \(\Sigma_{\rm B}\) & estado holonómico integral \\
16 & Relación de extensión y supervivencia nonádica
  & \({\rm EF}\text{--}G_9\) & prolongación y monodromía \\
17 & Monodromía nonádica con holonomía de frontera
  & \(\mathcal M_9\) & prolongación y monodromía \\
18 & Sección global del sistema inverso de supervivientes
  & \(X_\chi=\varprojlim_m{\rm Surv}^{(\chi)}_m\)
  & prolongación y monodromía \\
19 & Categoría bivariada de caminos aritméticos
  & \(\mathcal P_{\rm APP}^{(2)}\) & transporte cardinal \\
20 & Cronología observable de dos cursores
  & \(F_{\rm obs}\) & transporte cardinal \\
21 & Alfabeto decimal de cuarenta y dos tríadas
  & \(\mathcal A_{42}\) & emisión finita \\
22 & Relación tipada de extensión de fibras
  & \({\rm Ext}_r\) & estado holonómico integral \\
23 & Transferencia nonádica de fibras compatibles
  & \(K_{\rm ph}\) & prolongación y monodromía \\
24 & Refinamiento de semillas por avance y giro
  & \((P,H,Q)\) & transporte cardinal \\
25 & Cociclos de modo y orientación
  & \((c_{\rm mode},c_{\rm or})\) & estado holonómico integral \\
26 & Prolongación triádica de veinte bloques
  & \(E_{20}\) & prolongación y monodromía \\
27 & Sistema dodecafásico orientado
  & \(\mathcal R_{12}\) & lectura dodecafásica \\
28 & Observable armónico dodecafásico firmado
  & \(U_{12}^{\rm sgn}\) & lectura dodecafásica \\
29 & Transformación armónica por órbitas de paso tres
  & \(\mathcal H_{12}\) & lectura dodecafásica \\
30 & Extractor transversal de diferencias cíclicas
  & \(\mathcal E_{3,4}\) & lectura dodecafásica \\
31 & Recurrencia dodecafásica de cierre en base mil
  & \(\mathcal C_\alpha\) & cierre en base mil \\
32 & Transductor de historia enriquecida a carácter angular
  & \((L_{\rm tick},\chi_{\rm mass})\) & interfaz dimensional \\
33 & Retícula bidireccional de índice doce
  & \(\mathcal W_{\pm}\) & interfaz dimensional \\*
34 & Atlas nonádico de familias de extensión
  & \(\mathcal A_{\rm species}\) & interfaz dimensional \\
\bottomrule
\end{longtable}
\endgroup

El símbolo \(L_{\rm tick}\) designa el lector de iteración y retorno. Su
codominio es la memoria temporal del estado; no define una unidad temporal
física independiente.

El selector interno \(M_{\rm ph}\) pertenece al TPK\@. El transporte cardinal
\(T_d\) ocupa otro factor de la transición \(\mathcal U_t\) y realiza el
movimiento sobre APP\@. Esta separación conserva un único contenedor,
distingue selección de fase y transporte, y preserva la arquitectura completa
por tipos.

\Needspace{9\baselineskip}
\begin{resultbox}[title={Profundidad finita y límite enriquecido}]
La enumeración, los censos y la contracción de cilindros numéricos son
resultados exactos en sus dominios. La existencia y unicidad de una sección
enriquecida sometida conjuntamente a todos los predicados de supervivencia
se obtiene cuando las fibras de levantamiento son no vacías, unitarias y
compatibles en toda profundidad. Un cálculo sobre una ventana finita
comprueba las fibras de esa ventana, pero no implica por sí solo el
cuantificador sobre todas las profundidades.
\end{resultbox}

\begin{theorem}[Composición de la dinámica interna]
\label{thm:composicion-dinamica-tpk-rev11}
Una transición observable del TPK es la composición tipada
\[
 \mathcal U_t
 =\operatorname{Rec}_t\circ\operatorname{Mem}_t
 \circ T_{d(t)}\circ M_{\rm ph}(g(t))
\]
sobre \(\mathcal X_{\rm TPK}^{\rm enr}\), seguida, cuando corresponde, por
un lector modular, una extensión cilíndrica o una de las dos proyecciones.
Cada factor actúa sobre un componente distinto del estado; por ello el
selector, el movimiento, el reloj, la memoria y la publicación no pueden
colapsarse en una sola reducción modular.
\end{theorem}

La tabla anterior enumera las familias operatorias. Las equivalencias entre
realizaciones locales deben entrelazar sus dominios, reglas e invariantes.
Los desarrollos posteriores seleccionan composiciones de estas operaciones;
no introducen un cuarto objeto primitivo junto a APP, TRIT y TPK.
 \chapter{Transporte, selección y prolongación en el núcleo topológico de fase}
\label{sec:arquitectura-tpk}

La sigla TPK designa el \emph{Topological Phase Kernel}, o n\'ucleo
topol\'ogico de fase.\footnote{La denominaci\'on conserva asimismo la
dedicatoria autoral a Tan Peng Kian.} No designa un \`unico algoritmo de 108
pasos, una tabla de resultados ni el cat\'alogo de 468 emisiones. Es el
sistema de operadores que hace convivir, sobre la categor\'ia de caminos de
APP, las lecturas aditiva y multiplicativa sin identificar sus genealog\'ias.

La palabra \emph{topol\'ogico} tiene aqu\'i un contenido preciso. El soporte
de APP es el grafo de Cayley del grupo cociente
\[
 X_9=(\mathbb Z/9\mathbb Z)^2
\]
con sus aristas cardinales. Una ruta cerrada admite una clase de enrollamiento
despu\'es de elevarla al recubridor \(\mathbb Z^2\). TPK transporta esa clase,
la orientaci\'on de la ruta y la holonom\'ia de fase. Dos caminos que terminan
en la misma celda pueden, por tanto, representar estados distintos. El n\'ucleo
de fase es precisamente la informaci\'on que permanece invisible al publicar
solamente la celda final o su residuo.

La base dinámica es la categoría libre \(\mathsf P_{\rm APP}\) del grafo
cardinal de APP\@. Puesto que las dos evaluaciones aritméticas no inducen la
misma partición, el transporte vive sobre la categoría bivariada
\[
 \mathsf P_{\rm APP}^{(2)}
 =\mathsf P_{\rm APP}\times\mathsf P_{\rm APP}.
\]
La primera componente se evalúa mediante la hoja aditiva \(\Sigma\) y la
segunda mediante la hoja multiplicativa \(\Pi\). Esta separación es
constitutiva, pero no duplica APP: registra dos lecturas ordenadas de una
misma matriz de caminos. La componente multiplicativa es la tabla principal;
la aditiva conserva la ley de acumulación que genera sus filas y permite
medir qué información elimina una proyección visible.

\ifdefined\HMTTPKRepeatAPP
\section{Carta canónica y arquitectura suma--producto}

El soporte se escribe en coordenadas de cursor
\[
 I_9=\{0,\ldots,8\},\qquad X_9=I_9^2,
\]
pero sus etiquetas aritm\'eticas pertenecen a \(\{1,\ldots,9\}\). La carta
que relaciona ambos niveles es
\[
 \chi:I_9\longrightarrow\{1,\ldots,9\},\qquad \chi(i)=i+1.
\]
Sea
\[
 \rho_9(n)=1+((n-1)\bmod9),\qquad
 q_9(n)=\frac{n-\rho_9(n)}9.
\]
En particular, \(9\) es el representante visible de la clase nula. Los dos
observables de un mismo v\'ertice son
\[
 \Sigma(i,j)=\rho_9\bigl(\chi(i)+\chi(j)\bigr),\qquad
 \Pi(i,j)=\rho_9\bigl(\chi(i)\chi(j)\bigr).
\]
Esta convenci\'on impide mezclar la coordenada de cursor cero con la etiqueta
visible nueve. APP es la tupla
\[
 \mathcal A_9=
 \bigl(X_9,\operatorname{Cay}(X_9),\Sigma,\Pi,
       \operatorname{Path}(X_9)\bigr),
\]
y no la yuxtaposici\'on de dos tablas: cada semilla \((i,j)\) porta a la vez
las dos evaluaciones y cada ruta conserva el orden en que las consulta.

La elevaci\'on que retiene residuo y cociente es
\[
 \mathcal A_9^\#=(R^+,Q^+;R^\times,Q^\times),
\]
con, para \(a=\chi(i)\) y \(b=\chi(j)\),
\[
 a+b=R^+_{ij}+9Q^+_{ij},\qquad
 ab=R^\times_{ij}+9Q^\times_{ij},
\]
\[
 R^+_{ij}=\Sigma(i,j),\qquad R^\times_{ij}=\Pi(i,j).
\]
Los totales exactos son
\[
 \sum R^+=405,\quad \sum R^\times=459,\quad
 \sum Q^+=45,\quad \sum Q^\times=174.
\]
Por tanto,
\[
 2025-810=1215=(459-405)+9(174-45)=54+9\cdot129.
\]
La igualdad no crea TPK por simple conteo: identifica el dato que el
transporte debe conservar. El residuo visible difiere en \(54\) y la memoria
de cociente en \(129\).

La composici\'on de la hoja aditiva posee el cociclo de acarreo
\[
 \kappa_9(a,b)=
 \frac{\rho_9(a)+\rho_9(b)-\rho_9(a+b)}9,
\]
que satisface
\[
 \kappa_9(a,b)+\kappa_9(a+b,c)
 =\kappa_9(b,c)+\kappa_9(a,b+c).
\]
El sector multiplicativo posee adem\'as el radical
\[
 \mathfrak m=(3)=\{3,6,9\},\qquad \mathfrak m^2=0
 \quad\text{en }\mathbb Z/9\mathbb Z.
\]
Las seis filas unitarias suman \(45\), las filas \(3,6\) suman \(54\) y
la fila \(9\) suma \(81\); as\'i
\[
 459=6\cdot45+2\cdot54+81,
\]
y todo el exceso \(459-405=54\) se localiza en el sector no unitario. Esta
localizaci\'on, el cociclo de acarreo y la diferencia de cocientes son tres
datos distintos de la misma arquitectura suma--producto.
\fi

\section{Reducciones residuales y representación decimal}

La arquitectura usa tres reducciones, cada una con un cometido distinto:
\[
 \mathbb Z\xrightarrow{(\rho_9,q_9)}
 \{1,\ldots,9\}\times\mathbb Z,
 \qquad
 \rho_3^{\rm or}:\mathbb Z/3\mathbb Z\longrightarrow
 \{-1,0,+1\},
\]
\[
 (d_1,d_2,\ldots)\in\{0,\ldots,999\}^{\mathbb N}
 \xrightarrow{\iota_{1000}}
 \sum_{m\geq1}d_m1000^{-m}.
\]
La orientaci\'on ternaria es
\[
 \rho_3^{\rm or}([0])=0,\qquad
 \rho_3^{\rm or}([1])=+1,\qquad
 \rho_3^{\rm or}([2])=-1.
\]
La primera separa residuo visible y cociente non\'adico; la segunda asigna el
tipo orientado de la transici\'on; la tercera agrupa la lectura arquimediana
en tr\'iadas decimales, cuyo registro conserva el acarreo. Ninguna puede
sustituir a otra.
En particular, la reducci\'on m\'odulo tres no reconstruye el cociente m\'odulo
nueve, y una tr\'iada decimal no es una palabra de seis trits.
La suma acumula desplazamientos; el producto pliega clases y contrae las
filas no invertibles. La necesidad de TPK nace de que las dos lecturas
comparten soporte, pero no comparten partici\'on, periodo ni cociente.

\Needspace{36\baselineskip}
\section{Correspondencia entre símbolos del TPK, dominios y composiciones}
\label{sec:diccionario-interno-tpk}

Los nombres históricos se fijan aquí por objeto y no por semejanza
numérica. Esta tabla forma parte de la definición pública de TPK.

\begin{longtable}{@{}>{\raggedright\arraybackslash}p{.19\textwidth}
>{\raggedright\arraybackslash}p{.30\textwidth}
>{\raggedright\arraybackslash}p{.43\textwidth}@{}}
\toprule
Nombre&Objeto&Función dentro del núcleo\\\midrule
\endhead
coordenada nonádica&
\(\rho_9(n)\) junto con \(q_9(n)\)&
Separa el representante visible módulo \(9\) del cociente que conserva la
memoria del cruce de frontera. No es una raíz digital.\\
selector temporal ternario&
\(M_{\rm ph}(t)=\varepsilon(t)\in\{+1,-1,0\}\)&
En \(1,4,7\) activa la acumulación aditiva; en \(2,5,8\) activa la
multiplicativa; en \(3,6,9\) mantiene ambos cursores y registra umbral.
No designa las direcciones espaciales.\\
transporte cardinal&
\(T_N,T_E,T_S,T_O\)&
Desplaza la posición sobre el toro APP y conserva la palabra de ruta y su
enrollamiento. Es independiente del selector temporal.\\
lector ternario&
\(\rho_3^{\rm or}\)&
Convierte las tres clases residuales en orientación
\(0,+1,-1\). No reconstruye por sí solo el cociente nonádico.\\
lector en base mil&
\(\iota_{1000}\)&
Agrupa la publicación arquimediana en tríadas decimales y transporta sus
acarreos. Es un lector posicional, no una congruencia ``módulo \(1000\)''.\\
retorno de veintisiete&
\(\mathcal K_{27}=F_{\rm obs}^{27}\)&
Compone veintisiete actualizaciones posicionales y orientadas. Su cuarto
iterado restaura la fase observable, mientras el registro completo conserva
la memoria acumulada.\\
registro dodecafásico&
\(\Lambda\in\mathcal R_{12}\)&
Conserva seis posiciones, cinco diferencias y un total. Es el soporte del
transporte reversible \(1+5+6=12\).\\
lectura bidireccional&
\(q_+,q_-\) y sus pesos conjugados&
Conserva simultáneamente las orientaciones de ida y retorno. En Mecánica
Dimensional, una firma \((n_A,s_C)\) produce primero
\(W_\pm=6n_A\pm s_C\); sólo después aparece
\(s_C/6\) en el carácter angular.\\
\bottomrule
\end{longtable}

El selector temporal, el transporte cardinal y la orientación ternaria son,
por tanto, tres mapas distintos. Su separación conserva la diferencia entre
\emph{cuándo} se activa una lectura, \emph{hacia dónde} avanza la ruta y
\emph{con qué orientación} se registra el paso. TPK contiene los tres y
conserva sus relaciones.

\ifdefined\HMTTPKWordArithmetic
\section{De la célula local al monoide de palabras nonádicas}

La arquitectura suma--producto no termina en la tabla local. Para probar que
el cociente transportado genera aritm\'etica global, hay que construir la
operaci\'on sobre palabras sin introducir como definici\'on la multiplicaci\'on
entera que se pretende recuperar. Sea
\[
 \mathcal D_9=\{1,\ldots,9\},\qquad
 \mathcal W_9=\{\varnothing\}\sqcup
 \bigsqcup_{\ell\geq1}\mathcal D_9^\ell .
\]
Las palabras se escriben desde la posici\'on menos significativa y se eval\'uan
por
\[
 \operatorname{ev}_9(u_0,\ldots,u_{\ell-1})
 =\sum_{j=0}^{\ell-1}u_j9^j,\qquad
 \operatorname{ev}_9(\varnothing)=0.
\]
El s\'imbolo \(9\) conserva el cruce de frontera; no se sustituye por un cero.

\begin{theorem}[Forma normal non\'adica]
La evaluaci\'on \(\operatorname{ev}_9\) es una biyecci\'on entre
\(\mathcal W_9\) y \(\mathbb N_0\).
\end{theorem}
\begin{proof}
Las palabras de longitud \(\ell\) cubren exactamente el intervalo
\[
 \frac{9^\ell-1}{8}
 \leq\operatorname{ev}_9(u)
 \leq9\frac{9^\ell-1}{8}.
\]
El m\'inimo del intervalo siguiente es una unidad mayor que el m\'aximo del
anterior. Para \(n>0\), la divisi\'on de \(n-1\) por nueve determina de manera
\'unica
\[
 d=1+((n-1)\bmod9)\in\mathcal D_9
\]
y un cociente estrictamente menor, al cual se aplica recursivamente la misma
regla. El proceso termina y es reversible.
\end{proof}
Por ejemplo,
\[
 9\leftrightarrow(9),\quad
 10\leftrightarrow(1,1),\quad
 81\leftrightarrow(9,8),\quad
 1000\leftrightarrow(1,3,3,1).
\]
La \`ultima identidad exhibe en una sola forma normal las dos particiones
fundamentales
\[
 1000=729+271=729+243+27+1;
\]
su interpretaci\'on por canales se introduce despu\'es, no en la definici\'on
de la numeraci\'on.

\section{Suma, producto de Horner y conservación del cociente}

La suma nativa \(u\boxplus_\#v\) recorre las palabras por posiciones. En cada
posici\'on aplica la c\'elula aditiva a los dos d\'igitos presentes y, si existe
un acarreo entrante, aplica una segunda c\'elula al residuo provisional. El
residuo se publica y la suma de cocientes pasa a la posici\'on siguiente. El
invariante parcial es
\[
 \sum_{j<k}(u_j+v_j)9^j
 =\sum_{j<k}w_j9^j+c_k9^k,
\]
de donde se obtiene
\[
 \operatorname{ev}_9(u\boxplus_\#v)
 =\operatorname{ev}_9(u)+\operatorname{ev}_9(v).
\]

Para \(d\in\mathcal D_9\), el transductor escalar
\(d\odot_\#u\) usa primero
\[
 u_jd=9q_j+r_j,\qquad
 (r_j,q_j)=\bigl(R^\times(u_j,d),Q^\times(u_j,d)\bigr),
\]
y normaliza la colisi\'on entre \(r_j\) y el cociente entrante mediante la
c\'elula aditiva. As\'i,
\[
 \operatorname{ev}_9(d\odot_\#u)
 =d\operatorname{ev}_9(u).
\]
Si \(v=(v_0,\ldots,v_{m-1})\), se define, de derecha a izquierda,
\[
 a_m=\varnothing,\qquad
 a_j=(9\odot_\#a_{j+1})\boxplus_\#(v_j\odot_\#u),
\]
y
\[
 u\boxtimes_\#v:=a_0.
\]
Esta es una construcci\'on de Horner realizada s\'olo con las dos c\'elulas
locales y sus cocientes.

Esta descripción fija la dependencia operatoria, pero no constituye una
segunda residencia del resultado. El producto global, su prueba de Horner y
la identidad
\(operatorname{ev}_9(u\boxtimes_\#v)
=\operatorname{ev}_9(u)\operatorname{ev}_9(v)\)
se establecen una sola vez en el
\cref{eq:ev-multiplicative}, dentro del capítulo especializado de producto
nativo.
El ejemplo de frontera m\'as corto es
\[
 (9)\boxtimes_\#(9)=(9,8),\qquad 9+8\cdot9=81.
\]
Conservar s\'olo \(R^\times\) destruir\'ia este teorema: \(2\cdot5=10\)
quedar\'ia reducido al residuo visible \(1\). El cociente no es una anotaci\'on
auxiliar, sino la memoria necesaria para que el producto sobreviva al paso por
la frontera.

\section{Coproducto e irreducibilidad}

La apertura multiplicativa se define internamente por
\[
 \Delta_\#e_w
 =\sum_{u\boxtimes_\#v=w}e_u\otimes e_v.
\]
Su versi\'on reducida elimina los factores \((1)\). Si
\(\widetilde\Delta_\#\) denota esa versi\'on y
\[
 A_\#=
 \bigl(\overline{\widetilde\Delta_\#}\bigr)^*
 \overline{\widetilde\Delta_\#},
\]
entonces \(A_\#e_w=d_\#(w)e_w\), donde \(d_\#(w)\) cuenta los pares ordenados
no triviales que producen \(w\). El selector
\[
 P_{\mathrm{Irr},\#}
 =(I-P_{(1)})\,\mathbf1_{\{0\}}(A_\#)
\]
escoge exactamente las palabras distintas de la unidad sin apertura nativa no
trivial. Por tanto, la irreducibilidad se lee en la propia fibra del producto;
no se introduce mediante una tabla externa de primos.

Este resultado es exacto en el monoide de formas normales. Su extensi\'on a
historias completas exige una compatibilidad adicional. Si \(P_n\) proyecta
sobre historias supervivientes de profundidad \(n\), una familia de
inyecciones \(\iota_n\) debe satisfacer
\[
 P_n\iota_n\mu_\#
 =P_n\mu_n(\iota_n\otimes\iota_n).
\]
Los teoremas anteriores construyen \(\mu_\#\) y fijan esta ecuaci\'on; no
afirman por s\'i solos existencia o unicidad de la familia
\((\iota_n)_n\). Esta separaci\'on impide confundir la aritm\'etica global de
palabras, ya demostrada, con su prolongaci\'on monoidal al espacio profundo de
memoria TPK\@.

\fi

\section{Acción cardinal y realimentación aritmética}

Sean
\[
 \nu_N=(-1,0),\quad \nu_E=(0,1),\quad
 \nu_S=(1,0),\quad \nu_O=(0,-1)
\]
los cuatro generadores cardinales. Para \(d\in\{N,E,S,O\}\),
\[
 T_d(i,j)=(i,j)+\nu_d\pmod9.
\]
La letra \(O\) significa oeste. Estas direcciones son desplazamientos sobre
APP y no nombres de los canales \(\pi,e,\varphi\).

Una palabra \(w=d_1\cdots d_r\) determina
\[
 \gamma(x_0;w)_k=x_0+\sum_{j=1}^{k}\nu_{d_j}\pmod9.
\]
Cuando cierra, su elevaci\'on conserva
\[
 \operatorname{wind}(w)=\frac1{9}
 \bigl(\#S-\#N,\#E-\#O\bigr)\in\mathbb Z^2.
\]
La lectura bidireccional conserva palabra, orientaci\'on y hoja. Su
reversi\'on formal invierte el orden e intercambia
\(N\leftrightarrow S\), \(E\leftrightarrow O\). Su eventual lectura como
adelanto--retardo es una interfaz f\'isica posterior, no parte de esta
definici\'on de ruta.

El estado local observable es
\[
 s_t=(i_t,j_t,m_t),
 \qquad m_t\in\{\Sigma,\Pi\}.
\]
Se eval\'ua la hoja activa,
\[
 a_t=\rho_3^{\rm or}\bigl(m_t(i_t,j_t)\bigr),
\]
y el modo se actualiza mediante
\[
 m_{t+1}=
 \begin{cases}
 \Sigma,&a_t=+1,\\
 \Pi,&a_t=-1,\\
 m_t,&a_t=0.
 \end{cases}
\]
As\'i, cada paso cardinal modifica a la vez la posici\'on y la lectura
aritm\'etica futura. La ruta no se superpone a una palabra ya construida: es
la historia posicional que la produce.

\begin{theorem}[Sobreyectividad del factor muestreado de paso tres]
Para cada uno de los \(162=9^2\cdot2\) estados \((i,j,m)\) y cada trit
orientado objetivo \(b\in\{-1,0,+1\}\), existe una ruta
\(u\in\{N,E,S,O\}^3\) cuyo tercer paso emite \(b\). En los \(486\) casos
estado--objetivo, el n\'umero de rutas admisibles est\'a comprendido entre
\(8\) y \(40\). En consecuencia, para toda sucesi\'on objetivo
\((b_k)_{k\geq1}\) existe una historia cardinal cuya salida satisface
\[
 (a_3,a_6,a_9,\ldots)=(b_1,b_2,b_3,\ldots).
\]
Los dos trits intermedios \(a_{3k-2}\) y \(a_{3k-1}\) de cada bloque no
quedan prescritos por esta afirmaci\'on.
\end{theorem}
\begin{proof}
Se enumeran las \(4^3=64\) rutas de tres pasos desde cada estado. El
certificado exhaustivo incluido registra, para cada uno de los
\(162\cdot3\) casos, la ruta ejecutada, el estado final y el n\'umero total de
testigos; sus índices \(0,1,2\) se transportan por
\(\rho_3^{\rm or}\) a \(0,+1,-1\). Todos los conteos son positivos; el
m\'inimo es \(8\) y el m\'aximo es \(40\). Tras cada bloque se aplica el mismo
resultado al estado alcanzado.
La elecci\'on, por ejemplo, del primer testigo en el orden lexicogr\'afico
produce recursivamente una historia para cualquier sucesi\'on objetivo.
\end{proof}

La sobreyectividad se refiere al factor observado cada tres pasos, no a la
palabra emitida a tiempo unitario. Este teorema establece esa capacidad
muestreada del soporte de rutas. La selecci\'on de las secciones
\(\pi,e,\varphi\) pertenece a la relaci\'on enriquecida EF--G9 que se define
m\'as adelante: capacidad de realizaci\'on y selecci\'on geneal\'ogica son dos
operadores distintos.

\section{Cronología observable de los dos cursores}

El calendario se formula sobre
\[
 \Theta_{108}=\mathbb Z/108\mathbb Z,\qquad
 \beta(\theta)=
 \left[\left\lfloor\frac{\widetilde\theta}{9}\right\rfloor\right]_{12},
 \qquad
 r(\theta)=1+(\theta\bmod9),
\]
donde \(\widetilde\theta\in\{0,\ldots,107\}\). La firma activa es
\[
 \varepsilon(r)=
 \begin{cases}
 +1,&r\in\{1,4,7\},\\
 -1,&r\in\{2,5,8\},\\
 0,&r\in\{3,6,9\}.
 \end{cases}
\]
El espacio observable de dos cursores es
\[
 \mathcal Q_{\rm obs}
 =X_9^2\times\{N,E,S,O\}^2\times\Theta_{108}.
\]
Si \(\varepsilon=+1\) avanza el cursor aditivo; si
\(\varepsilon=-1\), el multiplicativo; si \(\varepsilon=0\), ninguno.
Las dos orientaciones se invierten después de los pasos
\(27,54,81,108\). Esta regla define una permutación
\[
 F_{\rm obs}:\mathcal Q_{\rm obs}\longrightarrow\mathcal Q_{\rm obs}.
\]
Las posiciones y orientaciones se restauran tras \(54\) pasos y la fase
completa tras \(108\):
\[
 F_{\rm obs}^{108}=I.
\]
Por ello la cronología de una sola ruta es reversible y no selecciona por sí
misma una distribución sobre las continuaciones; la selección de todas las
prolongaciones compatibles pertenece al operador de transferencia de fibras
definido más abajo.

\section{Emisor finito de dos cursores}

Una semilla de cursor es
\[
 \mathcal S=\mathbb T_9^2\times\{N,E,S,O\},
 \qquad |\mathcal S|=324.
\]
TPK m\'inimo recibe un cursor aditivo \(P\in\mathcal S\) y uno
multiplicativo \(M\in\mathcal S\). Sus firmas de seis ventanas son
\(s(P)\) y \(e(M)\). En cada coordenada se aplica
\[
 G(s,e)=100s+10\bigl((s+e)\bmod10\bigr)
       +\bigl((3s+5e+1)\bmod10\bigr).
\]
La cifra de centenas recupera \(s\); conocida \(s\), la de decenas recupera
\(e\). Por ello \(G\) es inyectiva sobre las firmas efectivas. El emisor
factoriza como
\[
 \mathcal S^2\xrightarrow{s\times e}
 \operatorname{im}s\times\operatorname{im}e
 \xrightarrow{G}\mathcal U_6,
\]
con
\[
 |\mathcal S^2|=104\,976,\qquad
 |\operatorname{im}s|=18,\qquad
 |\operatorname{im}e|=26,\qquad
 |\mathcal U_6|=18\cdot26=468.
\]
La factorizaci\'on \(468=18\cdot26\) es, por tanto, una descomposici\'on del
objeto: cada emisi\'on conserva una firma aditiva y una multiplicativa.

La proyecci\'on
\[
 \Psi_6:\mathcal U_6\longrightarrow\mathbb F_3^6,
 \qquad \Psi_6(U)=-U\pmod3,
\]
posee una imagen de \(243\) palabras. La cadena tipada es
\[
 104\,976\ \text{semillas}
 \longrightarrow468\ \text{emisiones decimales }U_6
 \longrightarrow243\ \text{palabras ternarias }w_6.
\]
El ambiente \(\mathbb F_3^6\) contiene \(729\) palabras; no se identifica
con la imagen efectiva de \(243\). Tampoco \(U_6\) y \(w_6\) son el mismo
tipo.

El emisor m\'inimo repite, en la carta dodecaf\'asica de cat\'alogo,
\[
 U_{12}^{\rm cat}=U_6\Vert U_6.
\]
Su periodo visible seis explica por qu\'e no puede sustituir al estado firmado
\(U_{12}^{\rm sgn}\). Esta diferencia no es una insuficiencia de TPK: separa
la proyecci\'on de cat\'alogo de la historia enriquecida.

Cada bloque decimal emitido pertenece a un alfabeto exacto de
\(7\cdot6=42\) tríadas. Si \(\overline{d_1d_2d_3}\) es una de ellas,
\[
 \boxed{d_3\equiv5d_2-2d_1+1\pmod{10}.}
\]
El barrido de las \(104\,976\) semillas no produce una tríada fuera de esta
ley. Conviene distinguir, por tanto, la cadena completa de cocientes:
\[
 \mathcal S^2
 \longrightarrow\Omega_{\rm seed}
 \longrightarrow\mathcal U_6
 \xrightarrow{\Psi_6}\operatorname{im}\Psi_6\subset\mathbb F_3^6.
\]
Las presentaciones microscópicas, las \(468\) emisiones decimales, las
\(243\) palabras ternarias y el ambiente de \(729\) palabras son cuatro
objetos distintos.

\paragraph{Conmutatividad orbital de los morfismos tipados.}
Sea
\(W_{243}:=\operatorname{im}\Psi_6\subset\mathbb F_3^6\). El embudo
completo y su cociente diedral se representan sin convertir inclusiones en
selecciones:
\[
 \underbrace{\mathcal S^2}_{104\,976\ \mathrm{semillas}}
 \xrightarrow[\text{censo de fibras}]{\text{emisión}}
 \underbrace{\mathcal U_6}_{468\ U_6}
 \xrightarrow[\text{proyección visible}]{\Psi_6}
 \underbrace{W_{243}}_{243\ w_6}
 \hookrightarrow
 \underbrace{\mathbb F_3^6}_{729\ \mathrm{palabras}},
 \qquad
 W_{243}\xrightarrow{\,/D_3^{\rm cat}\,}
 \underbrace{\mathcal O_{43}}_{43\ \text{\'orbitas}}.
 \tag{TPK-embudo}\label{eq:embudo-orbital-tipado-rev8}
\]
La primera flecha agrupa condiciones iniciales por la emisión que producen;
la segunda olvida la presentación decimal; la tercera es una inclusión en
el ambiente algebraico; la última es un cociente por la acción
\(D_3^{\rm cat}\) sobre \(W_{243}\). En particular, no hay un cociente
\(729\to43\) implícito. El censo exacto es
\[
 243=38\cdot6+5\cdot3,
\]
esto es, treinta y ocho órbitas de tamaño seis y cinco de tamaño tres.
Estas \(43\) clases orbitales tampoco son las \(43\) vueltas nonádicas
completas por canal del testigo G9 de mil cifras: el primer \(43\) pertenece
al cociente finito \(W_{243}/D_3^{\rm cat}\); el segundo es un contador de
profundidad de una ejecución certificada.

\section{Configuración del estado holonómico integral del TPK: hoja, orientación, residuo, cociente, acarreo, frontera, ruta y memoria}

La unidad din\'amica completa es una historia, no una palabra aislada. A
profundidad \(m\) escribimos
\[
 v_m=
 \bigl(
 B_m,\boldsymbol\delta_m,\boldsymbol x_m,
 \boldsymbol{\mathcal C}_m,\Sigma_m,
 \omega_m,\ell_m,g(m),\Lambda_m
 \bigr),
\]
donde
\[
 B_m=(u_m^\pi,u_m^e,u_m^\varphi)\in(\mathbb F_3^6)^3.
\]
Sus componentes tienen funciones distintas:
\begin{enumerate}
 \item \(B_m\) es el bloque visible conjunto de los tres canales;
 \item \(\boldsymbol\delta_m\) conserva las nuevas tr\'iadas decimales y el
       acarreo asociado;
 \item \(\boldsymbol x_m\) es el estado Hensel profundo;
 \item \(\boldsymbol{\mathcal C}_m\) contiene los cilindros ternarios y
       decimales y sus residuos enteros;
 \item \(\Sigma_m\) es la firma conjunta de frontera;
 \item \(\omega_m\) y \(\ell_m\) conservan orientaci\'on y hoja bidireccional;
 \item \(g(m)\in\mathbb Z/9\mathbb Z\) es la fase non\'adica;
 \item \(\Lambda_m\) es el registro dodecaf\'asico orientado.
\end{enumerate}

\begin{definition}[Presentaci\'on operatoria y categor\'ia TPK]
El inventario de operadores del estado holonómico integral de TPK es
\[
 \mathfrak T_{\rm HMT}=
 \bigl(
 \mathcal A_9,\TRIT,T_{N,E,S,O},T_{\min},
 H,\mathcal C_{3,1000},\Sigma_B,EF,\mathcal G_9,
 \mathcal R_{12}
 \bigr)
\]
que transporta historias \((v_0,\ldots,v_m)\) y sus truncamientos. Su
portador formal es la categor\'ia \(\mathsf{TPK}\) sobre la base observable:
sus objetos son estados holonómicos integrales y sus morfismos son las prolongaciones
\(\operatorname{Ext}_r\) estables por identidad y concatenaci\'on, definidas
m\'as abajo. El inventario no sustituye esa categor\'ia. La proyecci\'on
visible elimina componentes de \(v_m\); no define una equivalencia con el
estado completo.
\end{definition}

Esta definici\'on contiene exactamente lo que desaparece cuando se reinician
cursores o se conserva s\'olo el residuo: el camino, la hoja, el cociente, el
acarreo, la orientaci\'on y la frontera. TPK es, por ello, el punto en el que
la aritm\'etica modular deja de ser una operaci\'on sin historia.

\section{Cilindros ternario-decimales y estado de extensión}

Sea \(N\) el entero determinado por un prefijo ternario de longitud \(L\), y
sea \(D\) el entero determinado por \(K\) tr\'iadas decimales. Definimos
\[
 A=N1000^K-D3^L,
 \qquad
 B=(N+1)1000^K-D3^L.
\]
Los cilindros se intersectan exactamente cuando
\[
 \boxed{A<3^L,\qquad B>0.}
\]
Al incorporar un bloque ternario \(u\in\{0,\ldots,728\}\),
\[
 N'=729N+u.
\]
Si no aparece una nueva tr\'iada decimal,
\[
 A'=729A+u1000^K.
\]
Si aparece \(\delta\in\{0,\ldots,999\}\),
\[
 D'=1000D+\delta,
\]
\[
 A'=729000A+u1000^{K+1}-\delta\,729\,3^L.
\]
En ambos casos,
\[
 B'=A'+1000^{K+\Delta K}.
\]
Estas recurrencias enteras transportan el dato decimal sin convertirlo en un
objetivo exterior de la actualizaci\'on.

El estado Hensel conserva la informaci\'on que no aparece en el residuo. La
matriz activa es
\[
A_H=\begin{pmatrix}
2&2&2&1&2&1\\2&1&2&2&1&1\\1&0&0&1&0&2\\
0&0&1&1&1&0\\2&2&2&0&2&1\\2&1&2&1&0&0
\end{pmatrix},\qquad \det A_H=1,
\]
y la divisi\'on residuo--cociente es
\[
 y=xA_H,\qquad u=y\bmod3,\qquad x^+=\frac{y-u}{3}.
\]
Para los tres canales se forman adem\'as los cociclos conjuntos
\[
 q=(\pi+e+\varphi)\bmod3,
 \qquad
 c=\frac{\pi+e+\varphi-q}{3},
\]
\[
 a=(\pi+e-\varphi)\bmod3,
 \qquad
 h=\frac{\pi+e-\varphi-a}{3}.
\]
La actualizaci\'on de segundo orden corrige el siguiente estado con \(c\) y
\(h\). Por eso dos ramas con el mismo bloque visible no son necesariamente el
mismo estado: pueden diferir en su extensi\'on Hensel y en sus cilindros.
Como \(\det A_H=1\), la reducci\'on de \(A_H\) pertenece a
\(\operatorname{GL}_6(\mathbb F_3)\). Por tanto, antes de imponer cilindros,
fronteras o supervivencia, la carta residual Hensel cubre exactamente
\(\mathbb F_3^6\). Esta cobertura bruta no implica la cobertura del sistema
filtrado.

Formalmente, sobre cada configuración observable \(q\) se fija una fibra
\[
 \mathcal F_q=
 \mathsf O_q\times\mathsf H_q\times\mathsf C_q\times
 \mathsf S_q\times\mathsf B_q\times\mathsf\Lambda_q,
\]
correspondiente, respectivamente, a orientación y hoja, residuo y cociente,
acarreo, firma, cilindro y frontera, y registro de transporte. Para una
flecha cardinal \(r:q\to q'\), la prolongación admisible no se reduce a una
función sobre residuos; es la relación
\[
 \operatorname{Ext}_r\subseteq\mathcal F_q\times\mathcal F_{q'},
 \qquad
 \operatorname{Ext}_{r_2}\circ\operatorname{Ext}_{r_1}
 \subseteq\operatorname{Ext}_{r_2\circ r_1}.
\]
El incremento de acarreo es un cociclo:
\[
 \kappa(r_2\circ r_1)
 =\kappa(r_2)+\mathsf C(r_2)\bigl(\kappa(r_1)\bigr).
\]
Así queda tipado el mapa de olvido: proyectar una extensión a su ruta
observable elimina componentes de \(\mathcal F_q\), pero no autoriza a
identificar ambos objetos.

\section{Retornos compuestos y calendarios de transporte}

La longitud de una traza no determina su tipo. La presentación académica
separa el retorno compuesto, el calendario, la traza y el atlas por sus
respectivos estados conservados.

El s\'imbolo \(\mathcal K_{27}\) se reserva para la composici\'on de
veintisiete actualizaciones de un lazo de transporte posicional y orientado.
El entero \(27\) es su longitud de composici\'on; no lo convierte en una
traza de 108 iteraciones ni en un emisor infinito. Otra construcci\'on hist\'orica
agrupa cuatro lazos cuyos periodos conjuntos alcanzan 108; ese rotor conjunto
no se denota por \(\mathcal K_{27}\) aislado.

\begin{longtable}{@{}>{\raggedright\arraybackslash}p{.24\textwidth}>{\raggedright\arraybackslash}p{.31\textwidth}>{\raggedright\arraybackslash}p{.37\textwidth}@{}}
\toprule
Objeto&Estado conservado&Significado de 108\\\midrule
\endhead
Contador simplificado&posici\'on y velocidad&su periodo real es 18; es un
control hist\'orico.\\
Rotor conjunto de cuatro lazos&cuatro rutas, lectura, giro y movimiento&
los periodos individuales son \(54,54,36,36\); el estado conjunto tiene
periodo 108.\\
Traza bruta relojada&cuatro posiciones, orientaciones y fase&el contenido
din\'amico se repite a 54.\\
Traza de \(T_{\min}\)&dos posiciones y dos direcciones&el estado retorna a
54 y la salida es \(U_6\Vert U_6\).\\
Pseudoc\'odigo reinicializado&cursores reiniciados cada ventana&colapsa a
\((222,222,222,333,333,333)^2\); no representa el estado holonómico integral de TPK.\\
Rutas condicionadas&camino y palabra prescrita&108 es horizonte, no periodo.\\
Od\'ometro sectorial&\(\mathbb Z/9\mathbb Z\times\{0,\ldots,11\}\)&el
estado completo tiene periodo 108; el observable sectorial, 36.\\
Proyecci\'on bruta&cociente por olvido de \(X_{\rm TPK}^{\rm enr}\)&no posee
por s\'i sola la lectura firmada.\\
Traza enriquecida&todos los campos de \(v_m\) durante 108 actualizaciones&es
la restricci\'on temporal del estado holonómico integral de TPK.\\
Atlas de rutas&81 celdas/56 firmas o 324 semillas orientadas&es una taxonom\'ia,
no una \`orbita de 108 estados.\\
\bottomrule
\end{longtable}

El lazo \(\mathcal K_{27}\) es el operador de transporte finito de
veintisiete actualizaciones. Las ventanas, el rotor conjunto y la proyección
calendárica reciben símbolos distintos, de modo que la longitud compartida
no sustituye el dominio ni la composición de cada objeto.

\section{Transferencia de fibras y memoria de giro}

La factorización \(\mathcal U_6\simeq\mathcal S_+\times\mathcal S_\times\)
permite definir, respecto del conteo uniforme, las esperanzas condicionales
\[
 (E_+f)(s,e)=\frac1{18}\sum_{s'\in\mathcal S_+}f(s',e),
 \qquad
 (E_\times f)(s,e)=\frac1{26}
 \sum_{e'\in\mathcal S_\times}f(s,e').
\]
\(E_+\) renueva la fibra aditiva y conserva la multiplicativa;
\(E_\times\) realiza la operación dual. Son los únicos proyectores de Markov
que marginalizan por completo la coordenada activa, conservan la inactiva y
preservan el conteo uniforme.

Sea \(\tau=(+1,-1,0)^3\) la palabra de fase sobre \(C_9\), y
\(P_{+1}=E_+\), \(P_{-1}=E_\times\), \(P_0=I\). La transferencia nonádica es
\[
 \mathcal K_{\rm ph}\bigl((u,r),(v,r+1)\bigr)
 =P_{\tau(r)}(u,v).
\]
Es doblemente estocástica, irreducible, de periodo \(9\), y posee una única
distribución estacionaria \(1/(9\cdot468)\). Además,
\[
 \mathcal K_{\rm ph}^{\,9}=E_+E_\times,\qquad
 (E_+E_\times)(u,v)=\frac1{468}.
\]
Esta igualdad expresa la renovación simultánea de todas las continuaciones
compatibles, no el sucesor determinista de una sola ruta.

Tres operaciones geométricas sobre las semillas conservan información
adicional: \(P\), avance en la orientación activa; \(H\), media vuelta; y
\(Q\), cuarto de giro. La pareja \((P,H)\) refina las \(26\) firmas producto
en \(28\) clases. En particular, la firma \(555555\) se separa en tres hojas
de ocho representantes, y
\[
 18\cdot28=504=468+36.
\]
El término \(36\) cuenta hojas adicionales del refinamiento y no es el
objeto \(R_{36}\). La pareja \((P,Q)\) produce \(162\) clases de dos
elementos, permutadas transitivamente por
\[
 J(i,j,d)=\bigl(7-i,\,7-j,\,\overline d\bigr)\pmod9.
\]

En una traza de \(108\) pasos aparecen \(36\) signos de cada tipo,
\(72\) cambios efectivos entre las dos hojas y cuatro inversiones de
orientación; el cociclo firmado total es cero. Estos conteos son invariantes
del registro de evoluci\'on y enlazan la dinámica nonádica con el registro dodecafásico.

\section{Prolongación ternaria y frontera conjunta}

Relativamente a las semillas, fronteras y lectores nombrados del perfil
publicado, las palabras visibles iniciales son
\[
 w_6^\pi=010211,\qquad w_6^e=201101,\qquad
 w_6^\varphi=121200.
\]
Tras cinco bloques, los prefijos son
\[
\begin{aligned}
w_{30}^{\pi}&=010211|012222|010211|002111|110221,\\
w_{30}^{e}&=201101|121221|102011|012222|102011,\\
w_{30}^{\varphi}&=121200|112202|121020|010210|010200.
\end{aligned}
\]
La primera frontera profunda es
\[
 R_{36}=
 \begin{pmatrix}222220\\021222\\102011\end{pmatrix}.
\]
Para cada canal se define el prefijo
\(w_{36}^\chi:=w_{30}^\chi|u_5^\chi\); \(R_{36}\) es el triple ordenado de
los tres sextos bloques. Ninguno es un emisor ni un estado terminal.

Para \(B_k=(u_k^\pi,u_k^e,u_k^\varphi)\), las firmas
\[
 q_k=u_k^\pi+u_k^e+u_k^\varphi,\qquad
 a_k=u_k^\pi+u_k^e-u_k^\varphi
\]
determinan
\[
 u_k^\varphi=2(q_k-a_k),\qquad
 u_k^\pi+u_k^e=q_k-u_k^\varphi
 \quad\text{en }\mathbb F_3^6.
\]
As\'i, \(\varphi\) es el eje fijado y la libertad residual es el reparto
interno \(\pi/e\). En la novena fase,
\[
 B_9=(100100|020112|010122)
\]
admite tres salidas locales con el mismo eje \(\varphi\) y el mismo agregado
\(\pi+e\); las tres sobreviven un horizonte y s\'olo
\[
 B_{10}=(101222|112221|112002)
\]
sobrevive al segundo. Este es el descenso ejecutable \(3\to1\) de EF--G9.
La rama publicada contin\'ua hasta veinte bloques por canal: \(w_{30}\) y
\(R_{36}\) son prefijo y frontera, no clausura.

\section{Relación de extensión y monodromía nonádica}
\label{sec:extension-monodromia-tpk}

Sea \(\mathcal H_m\) el conjunto de historias enriquecidas admisibles de
profundidad \(m\), y sea
\[
 p_{n,m}:\mathcal H_n\longrightarrow\mathcal H_m
\]
el truncamiento. Para \(N\geq m\),
\[
 \operatorname{Surv}_m^{(N)}=p_{N,m}(\mathcal H_N),
 \qquad
 \operatorname{Surv}_m=\bigcap_{N\geq m}\operatorname{Surv}_m^{(N)}.
\]
La relaci\'on EF--G9 conserva las extensiones que pertenecen a este n\'ucleo
estable. La fase cumple
\[
 g(m+9)=g(m),
\]
pero el estado transportado satisface, en general,
\[
 v_{m+9}\ne v_m.
\]
El retorno act\'ua sobre la fibra de estados de una misma fase; no reinicia la
historia.

\begin{definition}[N\'umero HMT]
\label{def:numero-hmt-tpk-completo}
Un n\'umero HMT del canal \(\chi\) es únicamente una secci\'on global
compatible
\[
 x=(x_m)_{m\geq0}\in
 X_\chi:=\varprojlim_m\operatorname{Surv}_m^{(\chi)}.
\]
El lector y el valor publicado no forman parte de la secci\'on.
\end{definition}

Una presentaci\'on o medici\'on del n\'umero es un par \((x,L)\), donde
\(L\) es un lector declarado cuyo dominio contiene \(x\). Cuando \(L\) es el
lector arquimediano de cilindros, \(L(x)\) es el valor publicado; \(x\)
conserva adem\'as orientaci\'on, cociente, acarreo y frontera.

\begin{theorem}[Existencia de una prolongaci\'on global]
\label{thm:prolongacion-global-tpk-completo}
Sea \(h\in\mathcal H_m\). Si el \`arbol de prolongaciones de \(h\) es
finitamente ramificado y posee al menos un nodo a toda profundidad, entonces
existe una historia infinita cuyos truncamientos prolongan \(h\). Si, adem\'as,
cada nivel estable es unitario y los \`unicos estados son compatibles, la
secci\'on global es \`unica.
\end{theorem}
\begin{proof}
Los niveles finitos no vac\'ios forman un \`arbol finitamente ramificado. El
lema de K\"onig proporciona una rama infinita. La compatibilidad de los
truncamientos identifica esa rama con un elemento del l\'imite inverso. Si
cada nivel contiene un solo estado, no existe una segunda rama compatible.
\end{proof}

El teorema anterior es abstracto y condicional: por sí solo no afirma la
no-vaciedad ni la unicidad estable de un canal concreto. Para
\(\pi,e,\varphi\), esas hipótesis se descargan después mediante la filtración
racional ejecutable del
\cref{thm:generacion-monodromica-conjunta-rev7}. Esta residencia conserva
aquí únicamente el principio de límite inverso y no anticipa su aplicación.

\section{Composición de los operadores fundamentales}

Los operadores anteriores conducen al estado holonómico terminal sin
aplanar sus coordenadas:
\[
 \begin{aligned}
 \mathrm{APP}&\longrightarrow\mathrm{TRIT}
 \longrightarrow\text{ruta orientada}
 \longrightarrow\text{estado holonómico integral}\\
 &\longrightarrow\text{cilindros y supervivencia}
 \longrightarrow x_{\rm term}.
 \end{aligned}
\]
El cierre dodecafásico recibe entonces dos evaluaciones coordinadas del mismo
estado:
\[
 x_{\rm term}
 \xrightarrow{(\pi_{\rm term},R_{\rm sgn})}
 \bigl(B_{\rm term},U_{12}^{\rm sgn}\bigr),
 \qquad
 U_{12}^{\rm sgn}\xleftrightarrow{\mathcal H_{12}}K,
\]
\[
 (P,E,\Phi,K,c_{13}=0)\longleftrightarrow\alpha_{12}.
\]
No aparece una flecha \(B_{\rm term}\to U_{12}^{\rm sgn}\): ambas
coordenadas proceden de \(x_{\rm term}\). El paso posterior a Mecánica
Dimensional tampoco crea otra máquina; prolonga la genealogía conservada
mediante registro cronológico enriquecido, firma, carácter y torres.

\ifdefined\HMTInternalTraceability
\section{Registro académico y trazabilidad}

El archivo \path{datos/registro_operadores_tpk.json} forma parte del paquete.
Para cada operador declara nombre acad\'emico, alias hist\'orico, dominio,
codominio, transici\'on, invariantes, estatuto y fuente local. Incluye tambi\'en
la separaci\'on de las diez construcciones hist\'oricas asociadas al n\'umero
108. La puerta de publicaci\'on comprueba este registro y bloquea una salida
que vuelva a reducir TPK a su cat\'alogo o a una traza bruta.

En el resto del art\'iculo se usar\'an cuatro abreviaturas, ya tipadas:
\[
 w_6\in\mathbb F_3^6,\qquad
 w_{30}\in\mathbb F_3^{30},\qquad
 R_{36}\in(\mathbb F_3^6)^3,\qquad
 U_{12}^{\rm sgn}\in\mathbb Z^{12}.
\]
La primera es una palabra semilla; la segunda, un prefijo de cinco bloques;
 la tercera, la frontera conjunta de entrada en EF--G9; la cuarta, la
coordenada firmada declarada del estado dodecaf\'asico enriquecido. Ninguna igualdad cardinal autoriza a
identificarlas.
\fi
  \part{Álgebra holonómica de caminos y estructura electrónica local}

\chapter{Presentación universal del álgebra holonómica}
\section{Datos generadores}

Fijada una profundidad \(n\), sea \(X_n\) el conjunto finito de estados
APP--TRIT--TPK que conservan, en el nivel correspondiente, hoja, régimen,
orientación, residuo, cociente, acarreo, frontera, altura, ruta, memoria y
supervivencia. Sea \(\mathcal C_n\) la categoría pequeña generada por los pasos

\[
U_t={\rm Upd}_t\circ{\rm Tra}_t\circ{\rm Sel}_t.
\]

Los objetos de \(\mathcal C_n\) son los estados de \(X_n\); sus flechas son las
historias admisibles. La composición existe cuando el estado terminal de la
primera historia coincide con el estado inicial de la segunda.

A cada \(x\in X_n\) se asocia la fibra trítica

\[
A_x=A_{\tau(x)}
=\mathbb R[J_x]/(J_x^2+\tau(x)e_x),
\qquad \tau(x)\in\{+1,0,-1\}.
\]

El símbolo \(e_x\) designa la identidad local de la fibra. Para cada flecha
\(u:x\to y\), el transporte TPK induce una acción

\[
\rho_u:A_x\longrightarrow A_y
\]

sobre los coeficientes que sobreviven al paso. Cuando una transición cambia
de régimen, \(\rho_u\) se interpreta como el bimódulo tipado de esa
transición, no como una identificación de las dos álgebras cuadráticas.

El registro de memoria de concatenación define un multiplicador

\[
\Omega(v,u)\in Z(A_z)^\times,
\qquad
u:x\to y,\qquad v:y\to z,
\]

que registra el acarreo, la torsión y la contribución de frontera producidos
al componer. Su ley de compatibilidad es

\[
\rho_w(\Omega(v,u))\Omega(w,v\circ u)
=\Omega(w,v)\Omega(w\circ v,u).
\tag{1.1}
\]

La ecuación (1.1) es la identidad de cociclo que impide que la memoria de una
historia dependa de cómo se parentiza su composición.

\section{Presentación por generadores y relaciones}

Sea \(\mathcal F_n\) el álgebra asociativa libre generada por

\[
\{e_x,J_x:x\in X_n\}\cup\{T_u:u\in{\rm Mor}(\mathcal C_n)\}.
\]

Definimos el ideal \(\mathcal I_{\rm HMT}\) mediante las relaciones

\[
e_xe_y=\delta_{xy}e_x,
\qquad
\sum_{x\in X_n}e_x=1_n,
\tag{2.1}
\]

\[
J_x=e_xJ_xe_x,
\qquad
J_x^2=-\tau(x)e_x,
\tag{2.2}
\]

\[
T_u=e_{t(u)}T_ue_{s(u)},
\tag{2.3}
\]

\[
T_ua=\rho_u(a)T_u,
\qquad a\in A_{s(u)},
\tag{2.4}
\]

y

\[
T_vT_u=\Omega(v,u)T_{v\circ u}
\tag{2.5}
\]

cuando \(v\circ u\) existe; el producto es cero cuando las flechas no son
componibles. El álgebra holonómica finita es

\[
\boxed{
\mathfrak H_n=\mathcal F_n/\mathcal I_{\rm HMT}.}
\tag{2.6}
\]

La suma de \(\mathfrak H_n\) superpone rutas posibles. El producto concatena
historias. El factor \(\Omega\) impide que la concatenación borre el acarreo.
Por tanto, suma y producto han adquirido un significado simultáneamente
aritmético y dinámico.

\section{Teorema de asociatividad}

\textbf{Teorema 3.1.} Si \(\rho\) respeta las identidades locales y \(\Omega\)
satisface (1.1), el producto definido por (2.5) es asociativo.

\textbf{Demostración.} Sean \(u:x\to y\), \(v:y\to z\) y \(w:z\to t\). Para
coeficientes \(a,b,c\) situados en las fibras terminales pertinentes,

\[
\begin{aligned}
((cT_w)(bT_v))(aT_u)
={}&c\rho_w(b)\Omega(w,v)
\rho_{w\circ v}(a)\Omega(w\circ v,u)T_{wvu},\\
(cT_w)((bT_v)(aT_u))
={}&c\rho_w(b\rho_v(a)\Omega(v,u))
\Omega(w,v\circ u)T_{wvu}.
\end{aligned}
\]

La covariancia de \(\rho\) iguala los factores que contienen \(a\) y \(b\).
La identidad (1.1) iguala los dos productos de \(\Omega\). De ahí se sigue la
asociatividad sobre monomios, y la bilinealidad la extiende a toda
\(\mathfrak H_n\). \(\square\)

La demostración precisa la función de la memoria. Sin \(\Omega\), la
concatenación conservaría la ruta desnuda, pero no la discrepancia acumulada
entre sus elevaciones. Con \(\Omega\), dos maneras de agrupar la misma
historia producen el mismo estado completo.

\section{Propiedad universal de representación}

\textbf{Teorema 4.1.} Sea \(\{E_x\}_{x\in X_n}\) una familia de módulos. Supóngase
que se dan representaciones locales

\[
\pi_x:A_x\to\operatorname{End}(E_x)
\]

y operadores \(V_u:E_x\to E_y\) que satisfacen

\[
V_u\pi_x(a)=\pi_y(\rho_u(a))V_u,
\tag{4.1}
\]

\[
V_vV_u=\pi_z(\Omega(v,u))V_{v\circ u}.
\tag{4.2}
\]

Existe una única representación

\[
\Pi:\mathfrak H_n\to
\operatorname{End}\!\left(\bigoplus_{x\in X_n}E_x\right)
\]

que realiza \(e_x,J_x,T_u\) como las proyecciones, generadores y transportes
dados.

\textbf{Demostración.} La asignación de los generadores define un homomorfismo
desde el álgebra libre \(\mathcal F_n\). Las ecuaciones (4.1)--(4.2) y las
relaciones locales implican que \(\mathcal I_{\rm HMT}\) pertenece a su
núcleo. El homomorfismo factoriza por el cociente (2.6). La unicidad se sigue
de que las clases de \(e_x,J_x,T_u\) generan el cociente. \(\square\)

La propiedad universal convierte la precedencia causal del generador en una afirmación
algebraica. Un lector posterior es legítimo cuando representa las relaciones
del generador; no necesita reconstruirlas desde su codominio. El continuo,
el electrón, la gravedad y la rama excepcional son realizaciones distintas
de una misma presentación.

\section{El cociclo nonádico elemental}

La sección canónica \(s:C_9\to\{0,\ldots,8\}\subset C_{108}\) produce

\[
c_9(a,b)=\left\lfloor\frac{a+b}{9}\right\rfloor\in C_{12}.
\tag{5.1}
\]

La identidad

\[
c_9(a,b)+c_9(a+b\bmod9,c)
=c_9(b,c)+c_9(a,b+c\bmod9)
\tag{5.2}
\]

es exacta. Si \(X\) es el operador de elevación y \(\lambda\) la imagen del
generador de memoria, entonces

\[
X^9=\lambda I
\]

y

\[
X^aX^b
=\lambda^{c_9(a,b)}X^{a+b\bmod9}.
\tag{5.3}
\]

La ecuación (5.3) es la forma mínima de la hélice. El exponente visible
retorna módulo nueve; el factor \(\lambda\) conserva la vuelta. Si
\(\lambda\ne1\), retorno de fase y retorno de estado son operaciones
distintas.

\section{Refinamiento de estados y prolongación de lectores}

Los conjuntos finitos de estados forman un sistema inverso

\[
\cdots\xrightarrow{r_{n+2,n+1}}X_{n+1}
\xrightarrow{r_{n+1,n}}X_n.
\]

El estado profundo es

\[
X_\infty=\varprojlim_nX_n.
\tag{6.1}
\]

Las funciones y los idempotentes de cilindro se mueven en sentido opuesto:

\[
r_{n+1,n}^*:f\longmapsto f\circ r_{n+1,n}.
\tag{6.2}
\]

Si \(e_x\) es el indicador del cilindro \(x\in X_n\),

\[
r^*e_x=\sum_{y:r(y)=x}e_y.
\tag{6.3}
\]

La familia de álgebras forma, por tanto, un sistema directo:

\[
\mathfrak H_{\rm cyl}=\varinjlim_n\mathfrak H_n.
\tag{6.4}
\]

\textbf{Teorema 6.1.} Las inclusiones (6.2) son unitales; los idempotentes de
cilindro permanecen idempotentes; y todo lector compatible a profundidad
finita determina un único lector sobre \(\mathfrak H_{\rm cyl}\).

\textbf{Demostración.} La preimagen de una partición es una partición. Por ello
\(r^*(fg)=r^*f\,r^*g\), \(r^*1_n=1_{n+1}\) y la suma (6.3) es ortogonal.
La propiedad universal del límite directo extiende toda familia compatible
de lectores. \(\square\)

Este teorema aporta el significado exacto de «fractalizar hacia dentro y
expandir hacia fuera». Los estados se completan por restricción; los
lectores se completan por prolongación. La conservación de la unidad es la
unitalidad de estas inclusiones.

\section{Holonomía y graduación de memoria}

Sea \(d_M\) el grado de memoria. La holonomía satisface

\[
d_M(\Gamma_9x)=d_M(x)+1.
\tag{7.1}
\]

En el dominio reversible, \(\Gamma_9\) induce un automorfismo de
\(\mathfrak H_{\rm cyl}\) y puede formarse

\[
\mathfrak H_{\rm rev}\rtimes_{\Gamma_9}\mathbb Z.
\tag{7.2}
\]

En el dominio general, el avance produce un endomorfismo y el objeto es

\[
\mathfrak H_{\rm cyl}\rtimes_{\Gamma_9}\mathbb N.
\tag{7.3}
\]

La distinción entre (7.2) y (7.3) impide introducir una inversión donde sólo
se ha construido un avance. La inversión helicoidal ya tipada actúa en
(7.2); el espejo \(\pi/e\) actúa sobre las coordenadas del lector y constituye
otra involución.

\section{Álgebra trítica total}

Sea

\[
\mathbb A_{\rm tri}=A_+\oplus A_0\oplus A_-
\]

con idempotentes centrales \(p_+,p_0,p_-\). Los generadores satisfacen

\[
J_+^2=-p_+,
\qquad
J_0^2=0,
\qquad
J_-^2=p_-.
\tag{8.1}
\]

Sea

\[
\mathbb J=p_+J_++p_0J_0+p_-J_-.
\]

La exponencial definida por su serie interna satisface

\[
\boxed{
\begin{aligned}
\operatorname{Exp}_\star(t\mathbb J)
={}&p_+(\cos t+J_+\sin t)\\
&+p_0(1+tJ_0)\\
&+p_-(\cosh t+J_-\sinh t).
\end{aligned}}
\tag{8.2}
\]

\textbf{Demostración.} En la fibra elíptica las potencias pares e impares de
\(J_+\) alternan signos. En la fibra parabólica desaparecen todas las
potencias de orden al menos dos. En la fibra hiperbólica las potencias pares
son la identidad y las impares son \(J_-\). Agrupar las tres series prueba
(8.2). \(\square\)

La identidad conserva simultáneamente las tres geometrías. Ninguna de ellas
se obtiene por continuación informal de otra; las tres son sumandos
ortogonales del mismo objeto.

\section{Euler como proyección elíptica}

Sea \(\pi{}\) el cierre elíptico ya producido por la conexión
nonádica. Entonces

\[
\cos\pi{}=-1,
\qquad
\sin\pi{}=0
\]

en la realización circular, y (8.2) da

\[
p_+\operatorname{Exp}_\star(\pi{}\mathbb J)+p_+=0.
\tag{9.1}
\]

La ecuación (9.1) recupera la identidad de Euler en la cara elíptica. El
elemento total conserva además

\[
p_0\operatorname{Exp}_\star(\pi{}\mathbb J)
=p_0(1+\pi{}J_0),
\]

y

\[
p_-\operatorname{Exp}_\star(\pi{}\mathbb J)
=p_-(\cosh\pi{}+J_-\sinh\pi{}).
\]

La información que la representación compleja omite es precisamente el jet
de umbral y la dilatación. La identidad trítica no reemplaza la identidad de
Euler: identifica qué proyección la produce y qué memoria queda fuera de esa
proyección.

\section{La autoescala hiperbólica y \(\varphi\)}

Sea

\[
F_{\rm av}=\begin{pmatrix}0&1\\1&1\end{pmatrix}.
\]

Su cuadrado tiene determinante uno y traza tres. Definimos

\[
H_\varphi=\frac{2F_{\rm av}^2-3I}{\sqrt5}.
\tag{10.1}
\]

Un cálculo directo da

\[
H_\varphi^2=I.
\tag{10.2}
\]

Como

\[
\cosh(2\log\varphi)=\frac32,
\qquad
\sinh(2\log\varphi)=\frac{\sqrt5}{2},
\]

se obtiene

\[
\boxed{
F_{\rm av}^2
=\operatorname{Exp}(2\log\varphi{}\,H_\varphi).}
\tag{10.3}
\]

El régimen hiperbólico no se limita a admitir una unidad cuadrática: la
matriz de incidencia área--volumen realiza exactamente esa unidad. La razón
\(\varphi^{-2}\) es la dirección estable de una evolución unimodular.

\section{La función estructural de \(e\)}

La serie \(\operatorname{Exp}_\star\) transforma el generador aditivo en
transporte multiplicativo:

\[
\operatorname{Exp}_\star(a+b)
=\operatorname{Exp}_\star(a)\operatorname{Exp}_\star(b)
\]

cuando \(a\) y \(b\) conmutan. Su realización escalar sobre la identidad
define

\[
\operatorname{Exp}_\star(1)=e{}I.
\tag{11.1}
\]

La ecuación (11.1) no sustituye la generación decimal nonádica de \(e\).
Identifica su función algebraica: \(e\) es la escala que convierte la suma
infinitesimal en una ley multiplicativa de transporte. Éste es el puente
intrínseco entre las dos hojas APP en la representación exponencial.

\section{Alpha como funcional de transgresión dodecafásica}

El cociclo (5.1) representa la imposibilidad de escoger una sección
homomórfica \(C_9\to C_{108}\). La vuelta visible produce un elemento del
núcleo dodecafásico. El cierre de \(\alpha\) recibe ese elemento junto con la
palabra firmada y lo determina en una recta escalar.

Conviene distinguir el cociclo de su evaluación. Sea

\[
\mathcal T_\alpha:
Z^2(C_9;C_{12})\times X_{12}^{\rm sgn}
\longrightarrow\mathcal L_\alpha
\tag{12.1}
\]

el funcional que realiza la concatenación con memoria. Entonces

\[
\alpha{}=\mathcal T_\alpha(c_9,K).
\tag{12.2}
\]

La ecuación (12.2) no identifica \(\alpha\) con una clase de cohomología.
Afirma que \(\alpha\) es la realización escalar de esa clase después de
incorporar el estado dodecafásico. La segunda vía determina la misma sección
mediante

\[
P_9(\alpha{})=\delta,
\qquad
\delta=\log_{10}\frac{10}{9},
\tag{12.3}
\]

y ambas se reúnen en el producto fibrado de resolución.

Este tipo algebraico explica la singularidad de \(\alpha\). \(\pi,\varphi,e\)
son coordenadas de cierre, escala y transporte. \(\alpha\) es el funcional
que mide cómo esas coordenadas sobreviven al levantamiento no escindido.

\section{Matriz compacta del sistema correlacionado}

Las cuatro funciones se reúnen sin colapsarse en

\[
\boxed{
\begin{pmatrix}
\operatorname{Exp}_\star(\pi{}J_+) & 0&0&0\\
0&\operatorname{Exp}_\star(I)&0&0\\
0&0&\operatorname{Exp}_\star(2\log\varphi{}H_\varphi)&0\\
0&0&0&\operatorname{Exp}_\star(P_9(\alpha{})J_0)
\end{pmatrix}
=
\begin{pmatrix}
-I&0&0&0\\
0&e{}I&0&0\\
0&0&F_{\rm av}^2&0\\
0&0&0&I+\delta J_0
\end{pmatrix}.}
\tag{13.1}
\]

La matriz (13.1) es una identidad de funciones, no una tabla de valores. Sus
cuatro bloques significan cierre elíptico, multiplicación exponencial,
autoescala hiperbólica y costura parabólica. La compatibilidad dodecafásica
de la última fila conserva la genealogía conjunta.

\section{El rincón electrónico}

Sea

\[
P=\begin{pmatrix}1&0\\0&0\end{pmatrix},
\qquad
Q=\begin{pmatrix}3/4&\sqrt3/4\\\sqrt3/4&1/4\end{pmatrix}.
\]

Definimos

\[
S=2P-I,
\qquad
J=\frac4{\sqrt3}[P,Q].
\]

Entonces

\[
S^2=I,
\qquad
J^2=-I,
\qquad
SJ=-JS.
\tag{14.1}
\]

Por tanto,

\[
\boxed{
e_{\rm vac}\mathfrak H_ne_{\rm vac}\twoheadrightarrow
{\rm Cl}_{1,1}\cong M_2(\mathbb R).}
\tag{14.2}
\]

La flecha (14.2) expresa que el bloque electrónico es un cociente o una
representación de un rincón local de la aritmética holonómica. La
no conmutatividad no se añade desde la mecánica cuántica: ya reside en el
orden de las hojas APP. La raíz interna de \(-1\) es el generador elíptico;
las involuciones escindidas conservan la vacancia y la dilatación.

\section{Lectores y tipos de salida}

Sea

\[
\operatorname{Read}(\mathfrak H)
=\coprod_B\operatorname{Rep}_{\rm comp}(\mathfrak H,B).
\tag{15.1}
\]

Una constante se obtiene como

\[
c=\sigma_B(\rho(x_\infty)),
\tag{15.2}
\]

donde \(\rho\) es una representación compatible y \(\sigma_B\) un invariante
escalar propio de su codominio. Esta definición incluye sin confundir:

\[
\begin{array}{c|c|c}
\text{salida}&B&\sigma_B\\ \hline
\pi&\text{módulo elíptico}&\text{tiempo mínimo de cierre}\\
\varphi&\text{módulo hiperbólico}&\text{radio espectral estable}\\
e&\text{semigrupo exponencial}&\text{evaluación unitaria}\\
\alpha&\text{producto fibrado dodecafásico}&\text{costura escalar}\\
G_{\rm Cat}&\text{módulo de carácter}&\operatorname{Tr}(N^{-2}Q_4)\\
k_B&\text{rectas térmica y energética}&\text{transducción de unidad}\\
\gamma_{\rm BI}&\text{módulo área--volumen}&(12,-1)\text{ sobre }90/120\\
G&\text{módulo gravitatorio}&\text{carácter contragrediente}
\end{array}
\]

Una partícula no es una cifra de (15.2): es una sección persistente de un
módulo. Un campo es el transporte de esas secciones. Una simetría es el
grupo de automorfismos que preserva las relaciones y el borde del módulo.

\section{Teorema de conservación y simetría emergente}

Sea \(1_n=\sum_{x\in X_n}e_x\). Las inclusiones de refinamiento satisfacen

\[
r^*1_n=1_{n+1}.
\tag{16.1}
\]

Una historia individual define además una familia compatible de
idempotentes de cilindro. La primera es la unidad global del álgebra; la
segunda es la unidad genealógica de una sección. Distinguirlas evita
confundir conservación del espacio de posibilidades con conservación de una
ruta particular.

Para una representación \(\rho\) y una frontera \(b_\partial\), definimos

\[
\operatorname{Sym}(\rho,b_\partial)
=\{g\in\operatorname{Aut}(\rho(\mathfrak H)):
g\rho(\mathcal R_{\rm HMT})=\rho(\mathcal R_{\rm HMT}),
\;gb_\partial=b_\partial\}.
\tag{16.2}
\]

La simetría es, por definición, el estabilizador de una conservación ya
construida. En la rama excepcional, (16.2) conduce a los estabilizadores de
incidencia Witt--Golay--Leech--\(V^\natural\)--Monster. En la rama física,
conduce a los grupos que preservan los módulos electrónico, de gauge o
tensorial. No se impone una simetría para seleccionar el estado; se obtiene
la simetría que deja invariante su proyección.

\section{Información, entropía y cocientes}

Si \(U\) es biyectivo sobre el estado completo,

\[
H(X\mid U(X))=0.
\tag{17.1}
\]

Sea \(q\) un cociente que olvida una coordenada. Entonces

\[
H(X)=H(qX)+H(X\mid qX).
\tag{17.2}
\]

Un trit uniforme olvidado aporta

\[
H(X\mid qX)=\log3.
\]

El transductor térmico determina

\[
Q_{\rm reset}\ge k_B\Theta\log3.
\tag{17.3}
\]

Las ecuaciones (17.1)--(17.3) muestran que la entropía no mide el avance de
la holonomía. Mide la información descartada por un lector. La complejidad
esturmiana \(p(m)=m+1\) permite memoria infinita con entropía topológica cero;
el coste térmico aparece únicamente al hacer irreversible una proyección.

\section{Teorema nuclear}

\textbf{Teorema 18.1.} El objeto

\[
\boxed{
\mathfrak H_{\rm HMT}^{\rm alg}
=\left[\varinjlim_n
\bigl(\mathbb R\langle e_x,J_x,T_u\rangle/
\mathcal R_{\rm HMT}\bigr)\right]
\rtimes_{\Gamma_9}\mathbb N}
\tag{18.1}
\]

es una álgebra holonómica trigeométrica, graduada por fase y memoria, con
producto asociativo torcido por el acarreo. Su espacio de estados es el
límite inverso \(X_\infty\). Toda realización compatible de APP, TRIT y TPK
factoriza por (18.1). Sus representaciones tipadas determinan la estructura
discreta conjunta del continuo, el sistema correlacionado, el electrón, la acción,
la gravedad, las constantes y las simetrías sin convertir sus valores
terminales en generadores.

La notación de (18.1) es deliberadamente algebraica: no impone una norma
antes de escoger el codominio. Si
\(\rho:\mathfrak H_{\rm HMT}^{\rm alg}\to B(\mathcal K)\) es una
representación compatible y acotada, su realización completada es

\[
\mathfrak H_{\rm HMT}^{\rho}
=\overline{\rho(\mathfrak H_{\rm HMT}^{\rm alg})}^{\|\cdot\|}.
\tag{18.2}
\]

De este modo cada lector aporta la topología que efectivamente conserva;
ninguna completación convencional se introduce como parte del generador.

\textbf{Demostración.} La asociatividad procede del Teorema 3.1. La universalidad,
del Teorema 4.1. La conservación de unidad y la dualidad de refinamiento,
del Teorema 6.1. La elevación de memoria procede del cociclo nonádico (5.3).
La coexistencia de curvaturas procede de (8.1). Las realizaciones se obtienen
por la propiedad universal y los lectores (15.1)--(15.2). \(\square\)

\section{Certificación finita}

La verificación exhaustiva sobre los dominios finitos demuestra:

\begin{enumerate}
\item las \(9^3\) identidades del cociclo de acarreo;
\item la representación matricial \(X^9=2I\) y su ley torcida;
\item las tres relaciones \(J_+^2=-I\), \(J_0^2=0\), \(J_-^2=I\);
\item las tres ramas de la exponencial;
\item \(H_\varphi^2=I\) y \(\exp(2\log\varphi H_\varphi)=F_{\rm av}^2\);
\item el rincón \({\rm Cl}_{1,1}\);
\item la unitalidad contravariante/covariante del refinamiento;
\item la sincronía \(A_2\) con avance diagonal \(-16\);
\item la diferencia entre una permutación reversible y un cociente trítico;
\item el par nulo de Witt contenido en el rincón electrónico;
\item las relaciones exactas del par de observables de fase y el conmutador de memoria;
\item la irreducibilidad de APP, TRIT y TPK por cocientes de olvido, junto con la inversión del potencial en el lector de secuencias.
\end{enumerate}

Estas identidades constituyen una comprobación finita reproducible del sistema algebraico.

\section{El par nulo de Witt dentro del rincón electrónico}

El rincón electrónico aporta una realización local todavía más completa de
la trigeometría. Sean

\[
S=2P-I,\qquad
J=\frac{[P,Q]}{\sqrt3/4}.
\tag{20.1}
\]

Las relaciones ya demostradas son

\[
S^2=I,\qquad J^2=-I,\qquad SJ=-JS.
\tag{20.2}\label{eq:relaciones-cl11}
\]

Definamos ahora

\[
n_+=\frac{S+J}{2},\qquad
n_-=\frac{S-J}{2}.
\tag{20.3}
\]

Por (20.2),

\[
n_+^2=n_-^2=0,\qquad
n_+n_-+n_-n_+=I.
\tag{20.4}\label{eq:par-nulo-witt}
\]

La demostración es inmediata pero estructuralmente decisiva:
\((S\pm J)^2=S^2+J^2\pm(SJ+JS)=0\), mientras la suma de los
dos productos cruzados vale
\(\tfrac12(S^2-J^2)=I\). Por tanto, el mismo bloque
\({\rm Cl}_{1,1}\cong M_2(\mathbb R)\) contiene:

\[
\begin{array}{c|c|c}
\text{generador local}&\text{cuadrado}&\text{régimen}\\ \hline
J&-I&\text{elíptico}\\
n_\pm&0&\text{parabólico}\\
S&I&\text{hiperbólico}.
\end{array}
\tag{20.5}
\]

Ésta es una descomposición de Witt local exacta. La vacancia electrónica
adquiere así una segunda caracterización equivalente en su rincón: además de
ser la pérdida mínima de una unidad de acarreo multiplicativo bajo las
conservaciones APP, es una dirección nula parabólica entre las dos
orientaciones \(n_+\) y \(n_-\). TRIT tipa globalmente esas tres clases y TPK
las transporta con hoja, fase y memoria. La estructura de incidencia
Paley--Witt posterior conserva su dominio combinatorio propio; la
descomposición (20.3)--(20.5) aporta el modelo local desde el que puede
formularse su realización sin igualar ambos objetos.

\section{Producto cruzado de los observables de fase nonádica y esturmiana}

Sea

\[
\delta=\log_{10}\frac{10}{9},
\qquad
\zeta_9=\text{carácter primitivo de }C_9,
\qquad
\omega_\delta=e^{2\pi i\delta},
\tag{21.1}
\]

y sea \(T\) el paso orientado sobre la elevación reversible. En la base
\(\lvert t,m\rangle\), donde \(t\) registra el paso y \(m\) la altura de
memoria, introducimos los dos caracteres diagonales

\[
V_9\lvert t,m\rangle=\zeta_9^t\lvert t,m\rangle,\qquad
W_\delta\lvert t,m\rangle=\omega_\delta^t\lvert t,m\rangle.
\tag{21.2}
\]

La covariancia del paso produce las relaciones de Weyl

\[
V_9T=\zeta_9\,TV_9,\qquad
W_\delta T=\omega_\delta\,TW_\delta.
\tag{21.3}
\]

Para la vuelta nonádica \(\Gamma_9=T^9\), enriquecida por la actualización
TPK, se obtiene

\[
V_9\Gamma_9=\Gamma_9V_9,\qquad
W_\delta\Gamma_9=\omega_\delta^9\Gamma_9W_\delta.
\tag{21.4}
\]

El observable nonádico cierra; el observable de rotación avanza porque
\(\delta\notin\mathbb Q\).
Si \(D_M\lvert t,m\rangle=m\lvert t,m\rangle\), la ley
\(\Gamma_9\lvert t,m\rangle=\lvert t+9,m+1\rangle\) equivale a

\[
[D_M,\Gamma_9]=\Gamma_9,
\qquad
\Gamma_9^{-1}D_M\Gamma_9=D_M+I
\tag{21.5}
\]

sobre el sector reversible. Ésta es la ecuación operatoria de la hélice:
la clausura de fase es compatible con una elevación unitaria del grado de
memoria. El grupo basal de traslaciones sigue siendo abeliano; el álgebra
de observables y traslaciones deja de serlo por (21.3). La palabra
esturmiana aparece al aplicar el umbral TRIT a la órbita de
\(W_\delta\); por ello su aperiodicidad no decora la fase finita, sino que
constituye el factor simbólico de la rotación irracional.

La presentación operatoria pertenece al producto cruzado

\[
\mathfrak W_{9,\delta}
=C\!\left(C_9\times\mathbb T\right)\rtimes_T\mathbb Z,
\tag{21.6}
\]

enriquecido por la graduación \(D_M\). El espectro de
\(C(C_9\times\mathbb T)\) registra simultáneamente la fase finita y la fase
irracional; el generador de \(\mathbb Z\) implementa su traslación conjunta;
y (21.5) impide que la representación graduada descienda a una órbita
cerrada después de olvidar la memoria. El subshift obtenido por la partición
trítica de la segunda coordenada es el envolvente esturmiano. La suspensión
helicoidal es, por tanto, la extensión graduada de ese producto cruzado, no
la yuxtaposición metafórica de dos relojes.

Los caracteres \(\zeta_9\) y \(\omega_\delta\) pertenecen al generador
holonómico. Su representación elíptica posterior reconoce

\[
\rho_+(\zeta_9)
=\operatorname{Exp}_+\!\left(\frac{2\pi{}}9J_+\right),
\qquad
\rho_+(\omega_\delta)
=\operatorname{Exp}_+\!\left(2\pi{}\delta J_+\right).
\tag{21.7}
\]

La exponencial convencional aparece así como realización del carácter
interno ya construido; no selecciona los observables ni introduce un valor objetivo
en el generador.

Las relaciones (21.3)--(21.5) separan exactamente tres hechos:

\[
\text{período de fase }9,\qquad
\text{ausencia de período irracional},\qquad
\text{avance entero de memoria}.
\tag{21.8}
\]

Su coexistencia define la suspensión helicoidal esturmiana sobre la extensión
holonómica nonádica. El envolvente simbólico posee orden aperiódico de largo
alcance y entropía topológica nula, aunque su módulo espectral contenga
simultáneamente la frecuencia racional nonádica y la frecuencia irracional
\(\delta\). La denominación física «cristal de tiempo» requiere además fijar
la realización dinámica, su respuesta y su estabilidad; el teorema aquí
formulado es el objeto topológico-simbólico que la precede.

\section{La relación polinómica mínima de las tres curvaturas}

La suma ortogonal de los tres generadores TRIT puede escribirse como

\[
\mathbb J=p_+J_++p_0J_0+p_-J_-,
\qquad p_++p_0+p_-=I.
\tag{22.1}
\]

Sus tres cuadrados implican una sola relación:

\[
\boxed{\mathbb J^6=\mathbb J^2.}
\tag{22.2}
\]

Recíprocamente, si \(K=\mathbb J^2\), entonces \(K^3=K\) y los tres
idempotentes de régimen se recuperan polinómicamente:

\[
p_+=\frac{K^2-K}{2},\qquad
p_0=I-K^2,\qquad
p_-=\frac{K^2+K}{2}.
\tag{22.3}
\]

Así, la clasificación elíptica, parabólica e hiperbólica no requiere tres
álgebras yuxtapuestas desde fuera. Está contenida en el cálculo espectral
algebraico de un solo generador trítico. Sustituyendo (22.3) en la
exponencial total se obtiene la identidad compacta

\[
\boxed{
\operatorname{Exp}_\star(t\mathbb J)=
p_+(\cos t+\mathbb J\sin t)+
p_0(I+t\mathbb J)+
p_-(\cosh t+\mathbb J\sinh t).
}
\tag{22.4}
\]

La cara de Euler es la proyección elíptica de (22.4); la propagación
parabólica ocupa el umbral; la autoescala hiperbólica contiene la realización
de \(\varphi\). Al componerla con la holonomía,

\[
\mathcal E_{\rm HMT}(t)
=\Gamma_9\operatorname{Exp}_\star(t\mathbb J),
\qquad [D_M,\Gamma_9]=\Gamma_9,
\tag{22.5}
\]

se reúnen en dos ecuaciones la trigeometría y la elevación helicoidal. APP
aporta los idempotentes y la tensión entre hojas; TRIT aporta
\(\mathbb J\); TPK aporta \(\Gamma_9\), su cociclo y el grado \(D_M\).
Las salidas numéricas y físicas son representaciones tipadas de (22.5).
Esta presentación constituye la espina operatoria mínima del artículo
autónomo, mientras la álgebra universal (18.1) conserva toda la multiplicidad
de estados, rutas y refinamientos necesaria para reconstruir el desarrollo
integral.
 
\chapter{Vacancia de acarreo, álgebra de Clifford escindida y par nulo de Witt}
\begingroup
\renewcommand{\chapter}[2][]{}%
\chapter{Vacancia electrónica, acción, gravedad e información de frontera}

\section{La vacancia como defecto mínimo de acarreo}

La fibra aditiva de nivel diez,

\begin{equation}
 \mathcal F_{10}=\{(i,j)\in D_9^2:i+j=10\},
 \label{eq:fibra-suma-diez}
\end{equation}

contiene el centro $Q_0=(5,5)$ y las deformaciones orientadas
$Q_t=(5+t,5-t)$. Su producto es $25-t^2$. Conservar el residuo
multiplicativo del centro exige

\begin{equation}
 t^2\equiv0\pmod9.
\end{equation}

Para $|t|\leq4$, las únicas soluciones son $t=0,\pm3$. Las dos
deformaciones no centrales son, por tanto,

\begin{equation}
 Q_{+3}=(8,2),\qquad Q_{-3}=(2,8),
 \label{eq:vacancias-orientadas}
\end{equation}

intercambiadas por la inversión celular. En la carta APP,

\begin{equation}
 \mathcal A(5,5)=((1,1),(7,2)),
 \qquad
 \mathcal A(8,2)=((1,1),(7,1)).
 \label{eq:hojas-vacancia}
\end{equation}

Se conservan la evaluación aditiva, su cociente, el residuo multiplicativo y
el régimen trítico; cambia únicamente
$q^\times:2\mapsto1$. La vacancia es, por definición, esta pérdida mínima de
una unidad de acarreo multiplicativo bajo conservación de los restantes
invariantes. No es un cero añadido al estado.

La misma fibra se alcanza desde la reducción nonádica

\begin{equation}
 (369,126)\xmapsto{/9}(41,14)\longmapsto(5,5).
 \label{eq:369-centro}
\end{equation}

Posee dos prolongaciones distintas:

\begin{align}
 (41,14)&\xrightarrow{(+3,0)}(44,14)\longmapsto(8,5),
 \label{eq:rama-22-7}\\
 (41,14)&\xrightarrow{(+3,-3)}(44,11)\longmapsto(8,2).
 \label{eq:rama-vacancia-conservativa}
\end{align}

La primera satisface

\begin{equation}
 \frac{396}{126}=\frac{22}{7};
 \label{eq:aproximante-pi-vacancia}
\end{equation}

la segunda conserva

\begin{equation}
 41+14=44+11=55=54+1.
 \label{eq:unidad-borde-vacancia}
\end{equation}

Las dos ramas separan el lector de clausura del lector de vacancia. El
primero reconoce una razón de período; el segundo conserva la unidad de
frontera mientras reduce exactamente el acarreo.

\begingroup\fontsize{10.2}{11.7}\selectfont
\setlength{\parskip}{2pt}\setlength{\abovedisplayskip}{5pt}\setlength{\belowdisplayskip}{5pt}
\section{Rincón matricial y sección electrónica}

Las relaciones \eqref{eq:relaciones-cl11} y \eqref{eq:par-nulo-witt}
demuestran que la fibra central realiza localmente las tres geometrías. Para
precisar su modo singular sin ambigüedad de normalización, sean
$\Sigma_3,\Pi_3:D_9^3\to D_9$ las evaluaciones aditiva y multiplicativa de
las ternas y definamos su tabla conjunta por

\begin{equation}
 N^{(3)}_{ab}
 :=\#\{x\in D_9^3:\Sigma_3(x)=a,\ \Pi_3(x)=b\},
 \qquad
 C^{(3)}_{ab}:=\frac{N^{(3)}_{ab}}{\sqrt{n_a m_b^{\Pi}}},
 \label{eq:tabla-conjunta-tridimensional}
\end{equation}

donde $n_a=\sum_bN^{(3)}_{ab}$ y
$m_b^{\Pi}=\sum_aN^{(3)}_{ab}$. La enumeración exhaustiva de las
$9^3=729$ ternas da marginales de fila iguales a $81$ y marginales iguales
a $36$ en las seis columnas activas. Con

\begin{equation}
 m_{\mathrm{ph}}=(1,-1,0)^{\oplus3},
 \qquad
 N^{(3)}m_{\mathrm{ph}}
 =N^{(3)\mathsf T}m_{\mathrm{ph}}=-27m_{\mathrm{ph}},
 \label{eq:modo-entero-electronico}
\end{equation}

la normalización $\sqrt{81\cdot36}=54$ produce exactamente

\begin{equation}
 C^{(3)}m_{\mathrm{ph}}
 =C^{(3)*}m_{\mathrm{ph}}=-\frac12m_{\mathrm{ph}},
 \label{eq:modo-singular-electronico}
\end{equation}

y la incompatibilidad de los proyectores de hoja vale

\begin{equation}
 \|[P_{\Sigma,3},P_{\Pi,3}]\|=\frac{\sqrt3}{4}.
 \label{eq:torsion-minima}
\end{equation}

La incidencia areal--volumétrica proporciona

\begin{equation}
 1-\left(\frac12\right)^2=\frac34=\frac{90}{120}.
 \label{eq:incidencia-electronica}
\end{equation}

El punto $(5,5)$ es el único punto fijo de la inversión celular. De este
modo, la sección electrónica no se selecciona mediante una masa conocida:
se individualiza por la coincidencia de centralidad, vacancia mínima,
dirección nula, torsión \eqref{eq:torsion-minima} e incidencia
\eqref{eq:incidencia-electronica}. Su carácter adimensional es

\begin{equation}
 \mathfrak e=
 \frac{\sqrt3}{4}
 \left(e^{\varphi/\pi^2}-22\alpha^3\right)(1+15\Delta_4),
 \label{eq:caracter-electronico}
\end{equation}

donde cada factor procede de una etapa ya construida: la torsión mínima, la
orientación recíproca del lector $R(x,y)=x/y^2$, el complemento de vacancia,
la constante de compatibilidad y el determinante de la fibra. La carta de
energía se aplica después a esta sección.

\section{Polaridad y ecuaciones constitutivas}

Sea $\mathsf C$ la inversión de hojas. El carácter común $\mathcal E$ y el
carácter orientado $\mathcal B$ se definen por

\begin{equation}
 \mathcal E\circ\mathsf C=\mathcal E,
 \qquad
 \mathcal B\circ\mathsf C=-\mathcal B.
 \label{eq:paridad-constitutiva}
\end{equation}

Las dos combinaciones

\begin{equation}
 \widehat\mu=e^{\mathcal E+\mathcal B},
 \qquad
 \widehat\varepsilon=e^{\mathcal E-\mathcal B}
 \label{eq:mu-epsilon}
\end{equation}

satisfacen

\begin{equation}
 \widehat\mu\,\widehat\varepsilon=e^{2\mathcal E},
 \qquad
 \frac{\widehat\mu}{\widehat\varepsilon}=e^{2\mathcal B}.
 \label{eq:constitutivas-comun-orientada}
\end{equation}

El producto fija el cono de propagación y el cociente retiene la orientación
de impedancia. La diferencia de potencial pertenece al borde de la vacancia;
la respuesta transversal pertenece a su transporte conjugado. El sector
neutro es $\ker\mathcal B$, y toda sección fija por $\mathsf C$ pertenece a
él. Esta estructura hace de la polaridad una propiedad del transporte de la
vacancia, no una etiqueta impuesta a posteriori.

% La sección electrónica estructural concluye aquí. Las realizaciones de
% acción, gravitación, área e información se desarrollan posteriormente con
% sus fuentes rectoras completas y no se duplican en este capítulo.
 % 
 \endgroup
 \part{Orden topológico-temporal aperiódico y holonomía nonádica}


\par\endgroup
\chapter{Calendario interno de capacidad, relojes de fase y renormalización diofántica}
\section{La aperiodicidad como calendario interno de la capacidad nonádica}

La palabra esturmiana y la hélice nonádica no constituyen estructuras
superpuestas. Ambas son realizaciones de una sola aritmética de capacidad.
Sea

\[
 \rho:=6\log_{1000}3=\log_{10}9,
 \qquad
 \delta:=1-\rho=\log_{10}\!\left(\frac{10}{9}\right).
\]

El índice de capacidad decimal abierto después del bloque ternario número
\(t\) es

\[
 K_t:=\left\lfloor (t+1)\rho\right\rfloor,
 \qquad t\geq0,
\]

y la palabra mecánica de vacancias es

\[
 z_t:=\lfloor(t+1)\delta\rfloor-\lfloor t\delta\rfloor,
 \qquad t\geq1.
\]

Como \(\delta\) es irracional, \((t+1)\delta\notin\mathbb Z\). Por tanto,

\[
\begin{aligned}
 K_t
 &=\left\lfloor(t+1)(1-\delta)\right\rfloor\\
 &=t+1-\left\lceil(t+1)\delta\right\rceil\\
 &=t-\left\lfloor(t+1)\delta\right\rfloor.
\end{aligned}
\]

Al restar dos términos sucesivos se obtiene el vínculo que faltaba explicitar:

\[
 \boxed{K_t-K_{t-1}=1-z_t.}
\]

Ésta no es una semejanza entre dos conteos. Es una identidad exacta entre el
incremento de capacidad que usa la generación y la palabra de vacancias
producida por la incompatibilidad \(729/1000\). Sus dos casos son:

\[
 \begin{array}{c|c|c}
 z_t & K_t-K_{t-1} & \text{operación visible}\\
 \hline
 0 & 1 & \text{el bloque abre }729\text{ subceldas y entra una nueva poda decimal},\\
 1 & 0 & \text{el bloque abre }729\text{ subceldas sin aumentar la poda decimal}.
 \end{array}
\]

El segundo caso es la reapertura o \emph{stutter}. La palabra aperiódica decide,
por tanto, cuándo la capacidad avanza y cuándo permanece fija. La
aperiodicidad no gobierna desde fuera un generador periódico: es el calendario
interno de su operación abrir--podar.

La telescopía da una forma global del mismo teorema. Para \(0\leq a<b\),

\[
 K_b-K_a=(b-a)-\sum_{t=a+1}^{b}z_t.
\]

Cada intervalo queda descompuesto, sin resto, en incrementos de capacidad y
vacancias. En particular, toda conservación que se formule sobre la pareja

\[
 \bigl(\text{capacidad adquirida},\text{vacancia registrada}\bigr)
\]

debe conservar su suma, que es la longitud del intervalo de bloques. Esta es
la forma aritmética más elemental de la conservación evolutiva de la unidad
en este nivel; todavía no es la medida proyectiva del continuo, pero sí su
modelo local exacto.

\section{Dos tiempos, dos retornos y una retención de fase con avance de memoria}

La identidad anterior obliga a distinguir dos relojes.

El \textbf{tiempo de bloque} \(t\) aumenta en cada aplicación de la dinámica. El
\textbf{tiempo de capacidad} \(K_t\) cuenta las tríadas decimales efectivamente
podadas. La fase nonádica determinada es

\[
 g_t:=[K_t]_9\in C_9.
\]

De \(\Delta K_t=1-z_t\) se sigue

\[
 \boxed{g_t-g_{t-1}\equiv1-z_t\pmod9.}
\]

Cuando \(z_t=0\), el reloj de capacidad y la fase avanzan. Cuando \(z_t=1\),
el bloque temporal avanza pero \(K_t\) y \(g_t\) permanecen. Éste es un
acontecimiento exacto de \textbf{retención de fase con avance de estado}: el nuevo
bloque abre otra fibra de \(729\) posibilidades, cambia la firma y conserva
la historia de supervivencia aunque el residuo de fase sea el mismo. La
afirmación «la fase vuelve y la memoria cambia» posee así dos realizaciones
internas que no deben confundirse:

\begin{enumerate}
\item la retención local \(K_t=K_{t-1}\), causada por una vacancia esturmiana;
\item el retorno holonómico \(K\mapsto K+9\), que recorre un circuito completo de \(C_9\).
\end{enumerate}

En ambos casos coincide la fase determinada y difiere el estado holonómico integral,
pero el primer mecanismo es una pausa de capacidad y el segundo es una vuelta
nonádica completa.

Esto permite escribir el sistema base sin reducir el TPK a una rotación. Sea
\(\theta_t=\{t\delta\}\in\mathbb T\) y sea
\(x_t\in\mathcal X_{\rm TPK}^{\rm enr}\). El paso efectivo tiene la forma

\[
 (g_{t-1},\theta_{t-1},x_{t-1})
 \longmapsto
 \bigl(g_{t-1}+1-z_t,\;\theta_{t-1}+\delta,\;
       U_t(x_{t-1})\bigr),
\]

con

\[
 U_t=\operatorname{Upd}_t\circ\operatorname{Tra}_t
     \circ\operatorname{Sel}_t.
\]

La coordenada \(z_t\) agenda el avance de capacidad; no reemplaza la acción
de \(U_t\). El selector conserva hoja y régimen, el transporte conserva ruta
y orientación, y la actualización incorpora acarreo, frontera y registro.
Por eso la igualdad de \(g\) nunca autoriza la igualdad de \(x\).

\section{El reloj inverso y la suspensión de techo no constante}

La inversión monótona de la capacidad revela una segunda palabra mecánica.
Para \(k\geq1\), sea

\[
 T(k):=\min\{t:K_t=k\}
      =\left\lceil\frac{k}{\rho}\right\rceil-1.
\]

Introduzcamos

\[
 \sigma:=\frac{1}{\rho}-1=\frac{\delta}{\rho}.
\]

Como \(\rho\) es irracional, \(\sigma\) también lo es, y

\[
 T(k)=k+\lceil k\sigma\rceil-1.
\]

El número de bloques necesarios para incrementar una unidad de capacidad es

\[
 r_k:=T(k+1)-T(k)
 =1+\lceil(k+1)\sigma\rceil-\lceil k\sigma\rceil
 \in\{1,2\}.
\]

El valor \(r_k=2\) registra exactamente que entre las capacidades \(k\) y
\(k+1\) se interpuso una reapertura. Las posiciones excepcionales comienzan

\[
 20,41,62,83,104,125,145,166,187,208,\ldots,
\]

y sus separaciones son \(20\) o \(21\). Ésta es la misma dinámica mirada en
el reloj inverso; no debe confundirse con las separaciones \(21/22\) entre
vacancias en tiempo de bloque.

Para un incremento de capacidad de longitud \(m\), el tiempo de retorno es

\[
\begin{aligned}
 R_m(k)
 &:=T(k+m)-T(k)\\
 &=m+\lceil(k+m)\sigma\rceil-\lceil k\sigma\rceil.
\end{aligned}
\]

Por balance mecánico,

\[
 R_m(k)\in
 \{m+\lfloor m\sigma\rfloor,\;m+\lceil m\sigma\rceil\}.
\]

Se obtienen dos consecuencias exactas:

\[
 \boxed{R_9(k)\in\{9,10\}},
 \qquad
 \boxed{R_{108}(k)\in\{113,114\}}.
\]

Una vuelta de fase nonádica necesita nueve o diez bloques; una vuelta del
reloj completo \(C_{108}\) necesita ciento trece o ciento catorce. El soporte
finito conserva exactamente sus fases, pero el tiempo de bloque requerido
para recorrerlo es aperiódico. Esta es una formulación más precisa de la
\textbf{suspensión helicoidal esturmiana}: el reloj finito \(C_{108}\) es la base,
la rotación irracional proporciona el techo \(r_k\), y el cociclo TPK
transporta la fibra enriquecida.

La pareja \(6/7\) reaparece aquí porque

\[
 R_6(k)\in\{6,7\}.
\]

Sin embargo, este \(6/7\) mide bloques consumidos por seis incrementos de
capacidad. Es distinto tanto del número de vacancias en una ventana de
longitud \(132\) como de la separación entre huecos cortos de la palabra
primaria. Los tres fenómenos descienden de la misma pendiente, pero sólo una
conjugación explícita permitiría identificarlos como el mismo operador.

\section{Renormalización diofántica y conservación unimodular de la unidad}

La fracción continua de la densidad de vacancia, calculada después de que
\(\delta\) haya sido generada por \(729/1000\), comienza por

\[
 \delta=
 [0;21,1,5,1,6,2,3,1,1,8,1,2,2,5,\ldots].
\]

Sus primeros convergentes son

\[
 \frac1{21},\quad
 \frac1{22},\quad
 \frac6{131},\quad
 \frac7{153},\quad
 \frac{48}{1049},\quad
 \frac{103}{2251},\quad
 \frac{357}{7802},\quad
 \frac{460}{10053},\quad
 \frac{817}{17855},\quad
 \frac{6996}{152893},\quad
 \frac{7813}{170748},\ldots
\]

Esta lista explica, dentro de una sola jerarquía, los huecos primarios y sus
retornos inducidos. Los denominadores \(21,22\) son los primeros tiempos de
aproximación; los pares \((6,131)\) y \((7,153)\) son los siguientes tiempos
con sus números de eventos. En particular,

\[
 131=14\cdot9+5,
 \qquad
 153=17\cdot9,
\]

de modo que el retorno de longitud secundaria \(6\) deja residuo de fase
\(+5\), mientras el de longitud \(7\) cierra la fase nonádica.

El entero \(5\) no aparece aquí por ajuste. Es el cociente parcial que
produce

\[
 (6,131)=5(1,22)+(1,21),
\]

y el paso siguiente da

\[
 (7,153)=(6,131)+(1,22).
\]

Así se localiza una función exacta del nivel \(5\) dentro de la
renormalización de vacancias. Identificar este cociente parcial con la
fracción continua de Ramanujan de nivel \(5\) exige aún un mapa entre ambos
objetos; la igualdad del entero no lo proporciona por sí sola.

Sea \(p_n/q_n\) el \(n\)-ésimo convergente. Dos convergentes consecutivos
satisfacen

\[
 p_{n+1}q_n-p_nq_{n+1}=(-1)^n.
\]

En los dos primeros escalones esto se ve como

\[
 1\cdot22-1\cdot21=1,
 \qquad
 6\cdot153-7\cdot131=1.
\]

Por tanto, los vectores consecutivos \((p_n,q_n)\) forman una base primitiva
de \(\mathbb Z^2\) y conservan área reticular unitaria. La orientación
alterna, pero el módulo del determinante permanece igual a uno. Esta identidad
demuestra la evolución conservativa de la unidad dimensional: la pendiente real se aproxima a
profundidad creciente mientras cada cambio de escala conserva una celda
fundamental y admite inversa entera.

La recurrencia se escribe mediante

\[
 A(a)=
 \begin{pmatrix}a&1\\1&0\end{pmatrix},
 \qquad \det A(a)=-1.
\]

Los productos ordenados de las matrices \(A(a_n)\) generan los pares
\((p_n,q_n)\). Para \(a\ne b\),

\[
 A(a)A(b)=
 \begin{pmatrix}ab+1&a\\b&1\end{pmatrix}
 \ne
 \begin{pmatrix}ab+1&b\\a&1\end{pmatrix}
 =A(b)A(a).
\]

La rotación por \(\delta\) sigue siendo abeliana. La no conmutatividad
aparece en el cociclo de renormalización que recuerda el orden de los
cocientes parciales, además de la ya demostrada en las hojas APP y en el
cociclo afín. El cristal posee, por consiguiente, una capa abeliana de fase y
al menos dos capas ordenadas no abelianas de memoria.

Los errores orientados

\[
 \varepsilon_n=q_n\delta-p_n
\]

alternan de signo y nunca se anulan. Su signo define un lector trítico de
renormalización con valores \(\pm1\); el régimen \(0\) queda reservado al
umbral exacto, excluido aquí por la irracionalidad. Este lector no sustituye
al TRIT local de APP. Es una proyección a otra escala cuya naturalidad con el
selector trítico de fase debe probarse sobre la fibra enriquecida.

\section{Primera resonancia diofántica del reloj de bloque y separación del reloj de capacidad}

La sucesión de convergentes contiene retornos que cierran exactamente la
coordenada finita y sólo aproximan la coordenada irracional. Entre los
convergentes escritos arriba:

\[
 153\equiv0\pmod9,
 \qquad
 10053\equiv0\pmod9,
 \qquad
 170748\equiv0\pmod{108}.
\]

El último es el primer convergente de esta cadena cuyo denominador cierra el
reloj de bloque —la coordenada cronológica θ indexada por el número de
iteraciones—:

\[
 170748=1581\cdot108.
\]

Al mismo tiempo,

\[
 170748\,\delta-7813
 =-0.0000017458436865028488500088576\ldots
\]

El retorno no es periódico: la coordenada cronológica
\([t]_{108}\in C_{108}^{(t)}\) vuelve exactamente, pero la coordenada
irracional queda desplazada por un residuo no nulo. Sí es una
\textbf{cuasivuelta diofántica privilegiada}, porque procede de un convergente y
no de una búsqueda decimal.

Hay que separar este reloj del cociente de capacidad

\[
 C_{108}^{(K)}:\qquad g_{108}(t)=[K_t]_{108}.
\]

En el origen canónico, después de \(170748\) bloques, la capacidad adquirida
es

\[
 \lfloor170748\rho\rfloor
 =170748-7813
 =162935
 \equiv71\pmod{108}.
\]

Por tanto, ese convergente \textbf{no} cierra \(C_{108}^{(K)}\). Los retornos
\(R_{108}\in\{113,114\}\) de la sección anterior pertenecen precisamente a
este segundo reloj: cuentan cuántos bloques hacen falta para adquirir
\(108\) unidades de capacidad. La coincidencia de nombre \(C_{108}\) no
autoriza a fundir ambos índices. El calendario observable del TPK es el reloj
cronológico; el reloj de capacidad es una construcción inducida por la
compactificación \(729\to1000\).

El sistema distingue así cuatro nociones:

\begin{enumerate}
\item cierre exacto del reloj cronológico finito;
\item no cierre, en general, del reloj finito de capacidad;
\item proximidad óptima de la fase esturmiana;
\item transformación no trivial del estado holonómico integral por el cociclo TPK.
\end{enumerate}

Una repetición completa exigiría la clausura simultánea de los dos relojes,
la anulación del residuo irracional y el retorno de la fibra. Aquí sólo se
cierra el primero. Este desacoplamiento es el contenido matemático del
cristal aperiódico helicoidal: retorno exacto en un factor finito, deriva en
el otro, casi retorno jerarquizado en el factor irracional y evolución
genuina en la fibra.

\section{Espectro de vacancias sobre los cardinales internos de HMT}

La identidad complementaria permite evaluar cualquier longitud interna
\(N\) sin emplear una constante objetivo. Toda ventana de \(N\) pasos posee

\[
 V_N\in\{\lfloor N\delta\rfloor,\lceil N\delta\rceil\},
 \qquad
 P_N=N-V_N,
\]

donde \(V_N\) cuenta vacancias y \(P_N\) incrementos de capacidad. Para
algunos cardinales ya construidos por HMT se obtiene:

\[
\begin{array}{c|c|c}
N&V_N&P_N\\
\hline
108&4\text{ o }5&104\text{ o }103\\
132&6\text{ o }7&126\text{ o }125\\
243&11\text{ o }12&232\text{ o }231\\
396&18\text{ o }19&378\text{ o }377\\
468&21\text{ o }22&447\text{ o }446\\
729&33\text{ o }34&696\text{ o }695\\
1000&45\text{ o }46&955\text{ o }954.
\end{array}
\]

Dos filas requieren especial atención.

En longitud \(132\), la rama de seis vacancias produce exactamente

\[
 132=6+126.
\]

El entero \(126\) reaparece en el censo electrónico, pero los dos objetos no
se identifican todavía: uno es capacidad complementaria en una ventana
temporal y el otro es cardinal de una clase electrónica. El posible puente
debe enviar ventana, fase, orientación y memoria al censo, no sólo conservar
el número.

En longitud \(468\), que coincide con el cardinal del catálogo \(U_6\), el
número de vacancias es \(21\) o \(22\). Aquí la escala primaria reaparece
como multiplicidad de eventos dentro de un cardinal generado anteriormente.
La identidad es exacta; su promoción a autosimilitud del emisor requiere
demostrar que la enumeración de los \(468\) sextetos respeta el orden de la
rotación mecánica.

La ventana decimal completa da una clausura aún más explícita:

\[
\begin{array}{c|c|c|c}
V_{1000}&P_{1000}&P_{1000}-900&V_{1000}+P_{1000}\\
\hline
45&955&55&1000\\
46&954&54&1000.
\end{array}
\]

Así, las parejas \(45/46\) y \(55/54\) son complementos exactos después de
retirar las cien vueltas nonádicas \(900=100\cdot9\). El corredor
\(45\leftrightarrow54\leftrightarrow55\) no se añade mediante una
coincidencia: aparece como balance de poda y vacancia en la propia ventana
base \(1000\). La interpretación ulterior de \(54\) en el corredor modular
conserva su derivación APP propia; lo que se ha construido aquí es el mapa de
coocurrencia que ambas construcciones deberán hacer conmutar.

\section{La ventana \(396\), el borde \(27\) y la factorización radical \(369\)}

El primer intervalo de \(396=3\cdot132\) pasos contiene, en la fase canónica,
tres ventanas consecutivas con seis vacancias cada una. Por tanto,

\[
 396=18+378,
 \qquad
 18=3\cdot6,
 \qquad
 378=3\cdot126=42\cdot9.
\]

Los censos radical y electrónico satisfacen por otra vía

\[
 396=369+27,
 \qquad
 27=3\cdot9.
\]

Las dos descomposiciones se reúnen aritméticamente en

\[
 \boxed{369=378-9=3(132-6)-9}
\]

y

\[
 \boxed{27=18+9.}
\]

Estas dos identidades fijan la factorización cardinal del cuadrado: \(369\)
es la capacidad obtenida al separar una vuelta nonádica de las tres ventanas
bajas de \(132\), mientras \(27\) reúne las dieciocho vacancias y esa vuelta.
La afirmación es estrictamente cardinal en este punto; la realización sobre
la memoria de acarreo se obtiene más adelante mediante el lector entero de la
matriz de tránsito especular. La distinción de tipos impide confundir una
identidad de cardinales con una identificación prematura de fibras.

También fija el alcance de los números certificados \(396\), \(2376\),
\(387\) y \(43\). El certificado ejecuta \(396\) eventos de seis trits, de
donde

\[
 2376=396\cdot6.
\]

Al comparar cada nivel con el situado nueve posiciones después quedan
\(396-9=387\) pares, agrupados en

\[
 387=43\cdot9
\]

vueltas completas. La igualdad adicional \(2376=22\cdot108\) es verdadera,
pero no identifica los trits del certificado con veintidós copias del reloj
\(C_{108}\); el primero es un conteo de posiciones tríticas y el segundo un
calendario no escindido. Sólo un mapa que preserve sector, fase, carry y
memoria podría promover esa factorización a equivalencia de objetos.

\section{Sincronización, orientación y reapertura de los tres canales}

El censo extendido hasta \(t=120\) modifica la imagen de una única clausura en
la puerta novena. La sincronización de cardinalidades reaparece en racimos:

\[
\begin{aligned}
&\{8,9,10,12,13\},\quad
\{16,17,18,19,20\},\quad
\{29\},\\
&\{35,36,37,38,40,42\},\quad
\{49,51,52,54,55,56,57\},\\
&\{63\},\quad\{72,74\},\quad\{77,79,80\},
\quad\{83,84\},\\
&\{101\},\quad\{104,105,107,108\},
\quad\{116,118,119\}.
\end{aligned}
\]

Las primeras reaperturas están en

\[
 21,43,65,87,109,131,152,174,196,\ldots,
\]

con huecos \(21/22\). En la ventana certificada \(0\leq t\leq120\), ningún
\emph{stutter} es una sincronización triple. El dato es finito, no un teorema
asintótico, pero revela una dinámica de respiración:

\[
 \text{apertura amplia}
 \longrightarrow\text{poda}
 \longrightarrow\text{sincronización intermitente}
 \longrightarrow\text{reapertura aperiódica}.
\]

La primera respiración es explícita. Después de la sincronización triple

\[
 (59,59,59)\to(43,43,43)\to(32,32,32),
\]

aparecen nuevas sincronizaciones pequeñas

\[
 (17,17,17),\ (13,13,13),\ (6,6,6),\ (4,4,4),\
 (4,4,4),\ (3,3,3),\ (2,2,2),
\]

y la vacancia \(t=21\), con \(K_{21}=K_{20}=20\), reabre la fibra como

\[
 (611,527,723).
\]

La misma fase de capacidad \([20]_9=2\) permanece, pero las cardinalidades,
la orientación y la firma cambian de forma drástica. Esta es una evidencia
material directa de que el cristal no es una órbita contraída hacia un punto
fijo: alterna contracción y reapertura bajo un calendario aperiódico, mientras
la supervivencia mantiene la coherencia de largo alcance.

La igualdad de cardinalidades tampoco borra la orientación. En \(t=8\) hay
tres orientaciones locales y en \(t=9\) quedan dos; la supervivencia futura
selecciona una. En las puertas de orientación, \(\varphi\) permanece fijo y
\(\pi+e\) permanece constante, mientras el reparto \(\pi/e\) se refleja. La
sincronización debe definirse, por tanto, como igualdad del tamaño de las
fibras de canal, no como identidad de palabras, firmas o estados.

\section{De la sincronización triple al cierre dodecafásico de \(\alpha\)}

La nueva estructura aclara el lugar causal de \(\alpha\). Los canales
\(\pi\), \(e\) y \(\varphi\) no necesitan \(\alpha\) para sincronizar sus
cardinalidades: los racimos anteriores se producen durante la filtración y la
supervivencia. Tampoco la igualdad de cardinalidades basta para producir
\(\alpha\), porque todavía quedan orientación, reparto especular, carry y
condición de frontera.

El cierre dodecafásico recibe precisamente esa información adicional. Desde
el mismo estado holonómico terminal se abren dos lecturas:

\[
 x_{\rm term}
 \longmapsto
 B_{\rm term},
 \qquad
 x_{\rm term}
 \longmapsto
 U_{12}^{\rm sgn}.
\]

La primera determina los tres canales supervivientes; la segunda conserva la
coordenada firmada necesaria para propagar el acarreo. La vía entera produce

\[
 P+E-\Phi-K+C^+-1000C=A,
\]

mientras la vía de vacancias evalúa la misma genealogía mediante

\[
 \mathcal P(x)=\delta,
 \qquad
 \mathcal P(x)=
 2\pi x-\frac74x^2+\frac{x^3}{2\pi}
 +\frac{x^4}{20}-\frac{2x^5}{21}-\frac{x^6}{46}.
\]

La identidad \(\Delta K=1-z\) explica por qué la segunda vía pertenece al
mismo generador: \(\delta\) no es una constante externa colocada junto a los
tres canales, sino la densidad del calendario que decide su capacidad y sus
reaperturas. La sincronización de \((\pi,\varphi,e)\) y la realización de
\(\alpha\) no son dos historias separadas. Son dos niveles del mismo estado:

\[
 \text{sincronización de fibras}
 \longrightarrow
 \text{resolución orientada y memoria de carry}
 \longrightarrow
 \text{cierre dodecafásico de }\alpha.
\]

Esto conserva el candado causal: el sistema correlacionado
\((\pi,\varphi,e,\alpha)\) es una salida correlacionada de la generación
coinductiva. Dentro de esa genealogía común, los tres canales se determinan y
se sincronizan antes de que el sello dodecafásico determine la coordenada
\(\alpha\). Ningún valor convencional de las cuatro constantes interviene
como entrada.

\section{Cociclo multiplicativo de respiración y deriva logarítmica exactamente nula}

La complementariedad entre capacidad y vacancia admite una realización
multiplicativa que precisa la imagen autoral de cilindros que se contraen y se
reabren. Definimos la escala ambiente del cilindro después del bloque \(t\) por

\[
 \Omega_t:=\frac{729^{t+1}}{1000^{K_t}},
 \qquad
 K_t=\lfloor(t+1)\rho\rfloor,
 \qquad
 \rho=\log_{1000}729=\log_{10}9.
\]

La igualdad \(729=1000^\rho\) da inmediatamente

\[
 \boxed{\Omega_t=1000^{\{(t+1)\rho\}}.}
\]

Por tanto,

\[
 1\leq\Omega_t<1000
\]

para todo \(t\), y el conjunto \(\{\Omega_t:t\geq0\}\) es denso en el
intervalo multiplicativo semiabierto \([1,1000)\). La densidad no es una
hipótesis estadística: se sigue de la irracionalidad de \(\rho\). En efecto,
si \(\rho=a/b\in\mathbb Q\), entonces \(9^b=10^a\), contradiciendo la
unicidad de la factorización prima.

Sea

\[
 c_t:=K_t-K_{t-1}=1-z_t\in\{0,1\}
\]

el acarreo de compactificación. El cociente entre dos escalas consecutivas
es exactamente

\[
 \frac{\Omega_t}{\Omega_{t-1}}
 =\frac{729}{1000^{c_t}}
 =
 \begin{cases}
 729/1000,&z_t=0,\\[1mm]
 729,&z_t=1.
 \end{cases}
\]

La primera rama contrae; la segunda reabre. No son dos reglas añadidas al
calendario: son las dos cartas del mismo producto por \(729\) reducido módulo
la escala \(1000\). Antes de la primera reapertura se observa, por ejemplo,

\[
 \Omega_{20}=1.3100205086376203\ldots,
 \qquad
 \Omega_{21}=729\Omega_{20}
 =955.0049507968252\ldots.
\]

La contracción acumulada no conduce a un atractor porque cada vacancia
restaura la escala mediante un salto exacto. En coordenada logarítmica,

\[
 \ell_t:=\log_{1000}\frac{\Omega_t}{\Omega_{t-1}}
 =\rho-c_t=z_t-\delta,
 \qquad \delta=1-\rho.
\]

Además,

\[
 \sum_{j=1}^{n}\ell_j
 =\log_{1000}\Omega_n-\log_{1000}\Omega_0,
\]

de modo que todas las sumas parciales permanecen uniformemente acotadas. En
particular,

\[
 \boxed{\lim_{n\to\infty}\frac1n
 \log\frac{\Omega_n}{\Omega_0}=0.}
\]

Ésta es una afirmación más fuerte que la mera igualdad de una frecuencia
media: el cociclo logarítmico es un coborde acotado. La respiración conserva
la escala a largo plazo con exponente de Lyapunov nulo, aunque cada paso sea
una contracción o una reapertura muy distinta.

La forma geométrica natural del resultado es el círculo multiplicativo

\[
 \mathbb A_{1000}:=\mathbb R_{>0}/1000^{\mathbb Z}.
\]

La aplicación

\[
 R_{729}:[x]\longmapsto[729x]
\]

es conjugada, mediante \(u=\log_{1000}x\), a la rotación irracional
\(u\mapsto u+\rho\) de \(\mathbb R/\mathbb Z\). El entero \(c_t\) es el
índice de la transformación de cubierta que devuelve \(729\Omega_{t-1}\) a
la sección fundamental \([1,1000)\); la vacancia \(z_t=1-c_t\) es,
exactamente, la ausencia mínima de ese acarreo. Esta formulación reúne en un
solo objeto la compactificación, la elasticidad multiplicativa, la torsión de
la cubierta y la memoria de cuántas veces se cruzó su frontera.

La medida invariante inducida en la sección fundamental es logarítmica:

\[
 d\mu_{\Omega}(x)=\frac{dx}{x\log1000},
 \qquad 1\leq x<1000.
\]

Por consiguiente, para \(s\ne0\),

\[
 \int x^s\,d\mu_{\Omega}(x)
 =\frac{1000^s-1}{s\log1000},
\]

y la media geométrica es \(1000^{1/2}=10\sqrt{10}\). Estas magnitudes
describen la distribución matemática del radio multiplicativo; no se
identifican con una longitud física ni con la longitud de Planck sin el
transductor dimensional correspondiente.

El exponente nulo tampoco sustituye al teorema de Landauer del estado
holonómico integral. Aquí se demuestra ausencia de deriva exponencial en una
coordenada de escala; Landauer nulo exige recuperabilidad de la evolución
completa. Son resultados compatibles y tipológicamente distintos.

\textbf{Criterio de refutación.} Un valor de \(\Omega_t\) fuera de \([1,1000)\), un cociente
distinto de \(729/1000\) o \(729\), o una suma logarítmica no acotada.

\section{La fase y la vacancia como coordenadas de una única rotación levantada}

La fase finita y la palabra aperiódica no forman un producto arbitrario. Se
pueden pegar exactamente como dos coordenadas de una sola traslación
irracional. Sea

\[
 u_t:=\{(t+1)\rho\}\in[0,1),
 \qquad
 g_t^{(q)}:=[K_t]_q\in C_q,
\]

para \(q\geq2\). Sobre

\[
 X_q=C_q\times[0,1)
\]

con los extremos de los intervalos pegados por la cubierta, definimos

\[
 c(u):=\lfloor u+\rho\rfloor\in\{0,1\}
\]

y

\[
 F_q(g,u)=
 \bigl(g+c(u),\;u+\rho-c(u)\bigr).
\]

El segundo componente permanece en \([0,1)\); el primero conserva el cruce
de frontera. La aplicación

\[
 \Psi_q:X_q\longrightarrow\mathbb R/q\mathbb Z,
 \qquad
 \Psi_q(g,u)=[g+u]_q,
\]

satisface

\[
 \boxed{\Psi_q\circ F_q=R_\rho\circ\Psi_q,}
\qquad
R_\rho(y)=y+\rho\pmod q.
\]

En la órbita canónica,

\[
 \Psi_q(g_t^{(q)},u_t)=[(t+1)\rho]_q.
\]

Por tanto, la fase \([K_t]_q\), la coordenada irracional \(u_t\) y el
acarreo \(c_t=1-z_t\) no son tres decoraciones yuxtapuestas. Son la parte
entera, la parte fraccionaria y el cociclo de cubierta de un único punto en
el círculo de longitud \(q\).

Esta conjugación proporciona una prueba estructural de las complejidades
observadas. La codificación de \(R_\rho\) por los \(q\) intervalos unitarios
posee, en longitud \(m\), exactamente \(qm\) palabras. Sus fronteras son

\[
 j-k\rho\pmod q,
 \qquad j=0,\ldots,q-1,\quad k=0,\ldots,m-1,
\]

y son todas distintas porque \(\rho\) es irracional. Al añadir \(z_t\), un
bloque de longitud \(m\) reconstruye el símbolo de fase inmediatamente
anterior mediante

\[
 g_{t-1}^{(q)}=g_t^{(q)}-(1-z_t)\pmod q.
\]

Existe, por ello, una biyección entre los bloques decorados de longitud
\(m\) y los bloques de fase de longitud \(m+1\). Se obtiene

\[
 \boxed{p_{g,q}(m)=qm,\qquad
 p_{(g,z),q}(m)=q(m+1).}
\]

Para \(q=3,9,12,108\), estas leyes explican las complejidades lineales
certificadas sin suponer independencia entre un reloj periódico y una
decoración esturmiana. El crecimiento sigue siendo lineal y la entropía
topológica sigue siendo nula, pero el factor finito multiplica exactamente el
lenguaje disponible.

La inversión de la hélice es ahora inequívoca: en la coordenada \(y\) es la
rotación \(y\mapsto y-\rho\). El recorrido descendente no inventa otra
palabra; lee el mismo sistema invertible con el cociclo de cubierta conjugado.
La fase puede retornar, mientras \(u\) y la fibra TPK retienen la historia.

Esta formulación obliga asimismo a distinguir dos objetos de orden \(108\):

\[
 \theta_t=[t]_{108}\in C_{108}^{(t)},
 \qquad
 g_t^{(108)}=[K_t]_{108}\in C_{108}^{(K)}.
\]

El denominador \(170748=1581\cdot108\) cierra el primero. En cambio,

\[
 \lfloor170748\rho\rfloor=162935\equiv71\pmod{108},
\]

por lo que no cierra el segundo. El reloj observable dodecafásico del TPK y
el reloj inducido de capacidad son compatibles, pero no idénticos. Su futura
comparación requiere declarar el mapa de índice, no reutilizar el símbolo
\(C_{108}\) como si borrara el dominio.

\textbf{Criterio de refutación.} Una colisión entre
fronteras \(j-k\rho\), una complejidad distinta de \(qm\) o \(q(m+1)\), o
la afirmación de que \(170748\) cierra la fase de capacidad.

\section{Reloj inducido de vacancias y tipado trítico de los retornos \(6/7\)}

La escala secundaria \(6/7\) no sólo cuenta retornos de huecos cortos. Al
pasar del reloj de bloques al reloj de eventos, organiza exactamente la
duración de los tres regímenes del selector trítico.

Para \(n\geq1\), sea

\[
 p_n:=\left\lfloor\frac{n}{\delta}\right\rfloor
\]

la posición del \(n\)-ésimo evento de vacancia y sea

\[
 \kappa_n:=K_{p_n}
\]

la capacidad retenida en ese evento. Como

\[
 p_n\delta<n<(p_n+1)\delta,
\]

se obtiene

\[
 \boxed{\kappa_n=p_n-n.}
\]

Si

\[
 h_n:=p_{n+1}-p_n\in\{21,22\},
 \qquad
 s_n:=22-h_n\in\{0,1\},
\]

entonces

\[
 \kappa_{n+1}-\kappa_n=h_n-1=21-s_n.
\]

Reduciendo módulo \(3\),

\[
 \boxed{[\kappa_{n+1}]_3=[\kappa_n-s_n]_3.}
\]

Un hueco largo \(22\) conserva el régimen; un hueco corto \(21\) lo rota una
unidad. Como los huecos cortos reaparecen después de seis o siete eventos,
las mesetas del régimen inducido tienen longitudes exactamente \(6\) o \(7\).
Las veinte primeras clases son

\[
 \underbrace{2,\ldots,2}_{6},
 \underbrace{1,\ldots,1}_{7},
 \underbrace{0,\ldots,0}_{7}.
\]

Para relacionarlas con el selector canónico, se determina la clase positiva

\[
 r_n=1+(\kappa_n\bmod9)
\]

y se aplica

\[
 M_{\rm ph}(r)=
 \begin{cases}
 +1,&r\in\{1,4,7\},\\
 -1,&r\in\{2,5,8\},\\
 0,&r\in\{3,6,9\}.
 \end{cases}
\]

Así, las primeras mesetas recorren

\[
 \boxed{0\longrightarrow-1\longrightarrow+1\longrightarrow0}
\]

y cada transición es disparada por un hueco corto. Esta es una formulación
exacta de TRIT como unidad de información topológico-temporal: no reemplaza
la palabra binaria de vacancia, sino que integra sus eventos y determina en qué
régimen opera la siguiente transferencia. La información se conserva porque
la pareja \((\kappa_n,s_n)\) reconstruye el incremento anterior; la
orientación se conserva porque el orden cíclico de las clases no se aplana.

El resultado tampoco convierte los tres valores en tres teorías separadas.
Son tres regímenes internos de un único selector, y la palabra inducida los
visita bajo el mismo cociclo. Su identificación ulterior con curvaturas
elíptica, parabólica e hiperbólica utiliza las álgebras
\(\mathbb R[J_\tau]/(J_\tau^2+\tau I)\) ya construidas; la realización como
curvatura geométrica o campo físico requiere el transductor correspondiente.

\textbf{Criterio de refutación.} Una meseta de longitud distinta de \(6/7\), una transición de
régimen durante un hueco \(22\), o una transición que no sea el descenso
cíclico por un hueco \(21\).

\section{Matriz universal de las dos ramas equilibradas}

El balance esturmiano posee una forma lineal que reúne conservación y
orientación. En una ventana de longitud \(N\), sea

\[
 n_N=\lfloor N\delta\rfloor.
\]

Las dos posibilidades de conteo son \(n_N\) y \(n_N+1\), y sus complementos
son \(N-n_N\) y \(N-n_N-1\). Formamos la matriz de ramas

\[
 W_N:=
 \begin{pmatrix}
 n_N&N-n_N\\
 n_N+1&N-n_N-1
 \end{pmatrix}.
\]

Sus identidades son universales:

\[
 W_N\binom11=N\binom11,
 \qquad
 (1,-1)W_N=-(1,-1),
\]

\[
 \operatorname{tr}W_N=N-1,
 \qquad
 \det W_N=-N,
\]

y, por tanto,

\[
 \boxed{\chi_{W_N}(\lambda)=(\lambda-N)(\lambda+1).}
\]

El modo diagonal conserva la longitud total de la ventana. El funcional
orientado que resta una rama de la otra cambia de signo con valor propio
\(-1\). La unidad orientada no se introduce desde fuera: aparece porque las
dos palabras equilibradas difieren en un solo evento y en un solo acarreo
complementario.

Para las longitudes internas consideradas se obtienen, entre otras,

\[
 W_{132}=\begin{pmatrix}6&126\\7&125\end{pmatrix},
 \quad
 W_{468}=\begin{pmatrix}21&447\\22&446\end{pmatrix},
 \quad
 W_{1000}=\begin{pmatrix}45&955\\46&954\end{pmatrix}.
\]

Cada una posee los valores propios \(N\) y \(-1\). Esta ley separa con
precisión la conservación de la unidad total y la orientación elemental de
la vacancia. El tercer valor trítico \(0\) corresponde al umbral del
funcional orientado, no a una tercera rama de conteo añadida a la palabra
binaria.

La ventana \(1000\) contiene una reducción interna adicional. Después de
separar \(900=100\cdot9\) unidades de capacidad, queda

\[
 \widetilde W_{100}=
 \begin{pmatrix}
 45&55\\
 46&54
 \end{pmatrix},
\]

con

\[
 \operatorname{tr}\widetilde W_{100}=99,
 \qquad
 \det\widetilde W_{100}=-100,
 \qquad
 \operatorname{Spec}(\widetilde W_{100})=\{100,-1\}.
\]

Los semidesbalances de sus filas son

\[
 \frac{55-45}{2}=5,
 \qquad
 \frac{54-46}{2}=4.
\]

Por otra vía, el bloque central de APP

\[
 B_c=\begin{pmatrix}7&2\\2&7\end{pmatrix}
\]

posee espectro \(\{9,5\}\) y salto \(9-5=4\). Por consiguiente, los dos
lectores escalares

\[
 \chi_{\rm vac}(\widetilde W_{100})=(5,4),
 \qquad
 \chi_{\rm APP}(B_c)=(\lambda_{\min},
 \lambda_{\max}-\lambda_{\min})=(5,4)
\]

coinciden exactamente. Esto construye el codominio común del cuadrado
vacancias--APP. No demuestra todavía que \(\widetilde W_{100}\) y \(B_c\)
sean el mismo operador: para esa promoción falta la flecha que transporte
las dos ramas de ventana a los dos autoespacios de APP conservando fase,
hoja y carry. La reserva se localiza en esa flecha y no rebaja ninguna de las
dos derivaciones.

\textbf{Criterio de refutación.} Una
ventana equilibrada cuya matriz no tenga espectro \(\{N,-1\}\), o una
presunta identificación con APP que no transporte las ramas.

\section{Cociente de sincronización \(A_2\) y dirección proyectiva de la primera reapertura}

Los censos de los tres canales admiten un cociente canónico que elimina la
supervivencia común y conserva sólo su desalineamiento. Para

\[
 N=(N_\pi,N_e,N_\varphi)\in\mathbb Z^3
\]

definimos

\[
 \partial N=
 (N_\pi-N_e,\;N_e-N_\varphi,\;N_\varphi-N_\pi).
\]

Su imagen está contenida en

\[
 A_2=\{(x,y,z)\in\mathbb Z^3:x+y+z=0\},
\]

y es todo \(A_2\). El núcleo es la diagonal
\(\mathbb Z(1,1,1)\). En consecuencia,

\[
 \boxed{\mathbb Z^3/\mathbb Z(1,1,1)\simeq A_2.}
\]

La forma cuadrática

\[
 Q(N)=\frac12\|\partial N\|^2
\]

se anula exactamente en las sincronizaciones triples. Los defectos unitarios
de las fases \(11,14,15\) satisfacen \(Q=1\) y son raíces mínimas de
\(A_2\). Las igualdades de dos canales ocupan los tres espejos del grupo de
Weyl \(W(A_2)\simeq S_3\); la orientación a uno u otro lado del espejo queda
en el signo de la raíz.

El corredor inicial de espejos es

\[
\begin{array}{c|c|c|c|c}
t& (N_\pi,N_e,N_\varphi)&\text{espejo}&\partial N&d_t\\
\hline
4&(207,207,122)&\pi=e&(0,85,-85)&85\\
5&(39,151,151)&e=\varphi&(-112,0,112)&112\\
6&(110,110,69)&\pi=e&(0,41,-41)&41\\
7&(81,29,81)&\pi=\varphi&(52,-52,0)&52\\
8&(59,59,59)&S_3&(0,0,0)&0\\
9&(43,43,43)&S_3&(0,0,0)&0.
\end{array}
\]

Las cuatro amplitudes anteriores al primer cierre triple forman la matriz

\[
 D_{\rm syn}:=
 \begin{pmatrix}85&112\\41&52\end{pmatrix}.
\]

Sin introducir \(\alpha\), el electrón ni un valor experimental, se verifican
las identidades

\[
 112-85=27,
 \qquad
 52-41=11,
\]

\[
 85+41=126,
 \qquad
 112+41=153,
 \qquad
 85+52=137,
\]

y

\[
 \boxed{\operatorname{tr}D_{\rm syn}=137,\qquad
 \det D_{\rm syn}=-172=-4\cdot43.}
\]

El mismo corredor contiene así el censo \(126\), el retorno \(153\), la
dirección entera \(137\), el borde \(27\), el rango \(11\), el salto \(4\)
y la segunda sincronización \(43\). Todas estas igualdades son exactas. La
matriz, considerada aisladamente, determina una dirección entera de traza
\(137\), no la constante \(\alpha\). La derivación de \(\alpha\) requiere
además el estado firmado, el cierre dodecafásico y el jet completo, que se
construyen en las secciones correspondientes.

La primera reapertura ofrece una flecha más precisa. En \(t=20\),

\[
 N_{20}=(2,2,2),\qquad \partial N_{20}=0.
\]

En \(t=21\), la capacidad permanece \(K_{21}=K_{20}=20\), la fase determinada
es la misma y el selector se encuentra en la clase umbral
\(r=3\), para la que \(M_{\rm ph}=0\). Sin embargo,

\[
 N_{21}=(611,527,723)
\]

y

\[
 \boxed{\partial N_{21}
 =(84,-196,112)=28(3,-7,4).}
\]

El vector primitivo \((3,-7,4)\) pertenece a \(A_2\). En la carta proyectiva
orientada cuyo denominador es el componente \(\varphi-\pi\),

\[
 \chi_{e\varphi\mid\varphi\pi}(\partial N)
 :=\frac{N_e-N_\varphi}{N_\varphi-N_\pi},
\]

se obtiene

\[
 \boxed{\chi_{e\varphi\mid\varphi\pi}(\partial N_{21})
 =\frac{-196}{112}=-\frac74.}
\]

Éste es exactamente el coeficiente cuadrático del núcleo
\(E_{21/22}\). Ya no se trata de encontrar por separado los enteros \(7\) y
\(4\): la razón orientada completa emerge en la primera reapertura de los
tres canales, después de la sincronización y antes de la resolución de
\(\alpha\). La identidad es interna y no usa el valor objetivo de
\(\alpha\).

\section{Génesis esturmiano-triádica de \(22/7\) y retrato interno del jet \(E_{21/22}\)}

Las primeras siete posiciones efectivas de vacancia son

\[
 21,43,65,87,109,131,152.
\]

Entre la primera y la séptima transcurren \(132\) pasos y aparecen seis
intervalos. En esa primera ventana cerrada,

\[
 132=6\cdot22,
 \qquad
 126=6\cdot21.
\]

La amplificación trítica da

\[
 396=3\cdot132,
\]

y, por tanto,

\[
 \boxed{\frac{396}{126}
 =\frac{3\cdot6\cdot22}{6\cdot21}
 =\frac{22}{7}.}
\]

El sello racional \(22/7\) no se introduce como aproximación histórica de
\(\pi\): se obtiene del cociente entre la amplificación trítica de la
ventana y su supervivencia en el ramal de seis huecos largos. La comparación
posterior con \(\pi\) conserva su estatuto de reconocimiento; no genera la
razón.

La dinámica aporta además cinco pares de extremos invariantes bajo el cambio
de origen de la palabra mecánica:

\[
 \mathcal G_{\rm prim}=\{21,22\},
 \quad
 \mathcal R_{\rm sec}=\{6,7\},
 \quad
 \mathcal G_{\rm inv}=\{20,21\},
\]

\[
 \mathcal V_{1000}=\{45,46\},
 \qquad
 \mathcal H_{100}=\{4,5\}.
\]

Con la orientación ya fijada por el TPK, consideremos la extracción

\[
 r^+=7,\qquad h^-=4,\qquad
 g_{\rm inv}^-=20,\qquad g_{\rm prim}^-=21,\qquad
 v_{1000}^+=46.
\]

Estos cinco invariantes reproducen exactamente los coeficientes no
trascendentes del jet de orden seis:

\[
 -\frac{r^+}{h^-}=-\frac74,
 \qquad
 \frac1{g_{\rm inv}^-}=\frac1{20},
 \qquad
 -\frac2{g_{\rm prim}^-}=-\frac2{21},
 \qquad
 -\frac1{v_{1000}^+}=-\frac1{46}.
\]

Los grados uno y tres forman el par recíproco

\[
 (2\pi)\,(2\pi)^{-1}=1.
\]

Así, el núcleo declarado puede escribirse como

\[
\begin{aligned}
 P_6(x)={}&
 2\pi{}x
 -\frac{r^+}{h^-}x^2
 +(2\pi)^{-1}x^3
 +\frac{x^4}{g_{\rm inv}^-}\\
 &-\frac{2x^5}{g_{\rm prim}^-}
 -\frac{x^6}{v_{1000}^+}.
\end{aligned}
\]

Esta escritura no usa la raíz de \(P_6(x)=\delta\) para escoger los
coeficientes. Todos los números del lado derecho han aparecido antes como
invariantes del calendario de compactificación. El coeficiente \(-7/4\)
posee, además, la segunda factorización interna de la sección anterior como
coordenada proyectiva del primer defecto \(A_2\).

La igualdad anterior es una \emph{factorización genealógica exacta} del
núcleo. Los extremos seleccionados, sus grados y sus signos quedan reunidos
por el lector graduado \(\mathfrak E_9\), definido más adelante sobre la firma
de invariantes del estado holonómico integral; este lector conmuta con los
transportes que conservan origen, orientación, ventanas y memoria.

La prolongación de grados \(7,8,9\) ya posee otro estatuto. Las construcciones de referencia
demuestran

\[
 T_{7:9}(x)
 =-\frac{x^7}{120}+\frac{x^8}{45}+\frac{2x^9}{495},
\]

con

\[
 120=|\mathcal F_8^{\rm ph}|,
 \qquad
 45=\det B_c,
 \qquad
 11=\operatorname{rank}(D_3,D_4),
 \qquad
 495=45\cdot11,
\]

y la resolvente nilpotente \((I-N)^{-1}\). El jet completo

\[
 P_9=P_6+T_{7:9}
\]

queda así estratificado: los grados \(1{:}6\) proceden de la firma de
sincronización y los grados \(7{:}9\) de la incidencia de fase, del
determinante de borde y del rango dodecafásico. La raíz simple de
\(E_{21/22}\) se obtiene después de esa factorización y no interviene en la
selección de los coeficientes.

\section{Cadena causal refinada: de la compactificación a la coordenada \(\alpha\)}

Los resultados de las secciones anteriores permiten sustituir una sucesión
verbal por un diagrama de dominios y cocientes:

\[
\begin{aligned}
&\text{APP--TRIT--TPK y razón interna }729/1000\\
&\quad\longrightarrow
 (\mathbb A_{1000},R_{729},c_t,z_t)\\
&\quad\longrightarrow
 (g_t^{(9)},u_t,x_t^{\rm enr})\\
&\quad\longrightarrow
 (N_\pi,N_e,N_\varphi)\\
&\quad\xrightarrow{\ \partial\ }
 A_2\\
&\quad\longrightarrow
 \bigl(B_{\rm term},U_{12}^{\rm sgn}\bigr)\\
&\quad\longrightarrow
 \begin{cases}
 \text{cierre entero dodecafásico},\\
 \text{raíz simple del jet }E_{21/22},
 \end{cases}\\
&\quad\longrightarrow\alpha.
\end{aligned}
\]

El primer nivel construye la rotación multiplicativa y su cociclo de
acarreo. El segundo conserva simultáneamente fase, parte irracional y estado.
El tercero determina los tres canales correlacionados. El cociente \(A_2\)
mide su desalineamiento sin borrar la supervivencia común. El cierre
dodecafásico recibe después la orientación y la memoria que el cociente de
censos no contiene.

La sincronización, por ello, no genera por sí sola \(\alpha\). Sí proporciona
dos precursores internos concretos:

\begin{enumerate}
\item el rectángulo de defectos, cuya traza es \(137\);
\item la primera reapertura umbral, cuya coordenada proyectiva es \(-7/4\).
\end{enumerate}

Ambos aparecen antes de resolver la ecuación analítica y antes de toda
comparación experimental. El primer precursor localiza una dirección entera;
el segundo localiza un coeficiente del jet. La información restante reside
en el estado firmado, en el acarreo dodecafásico y en los grados superiores
del núcleo.

Esta arquitectura precisa en qué sentido el cristal aperiódico puede ser el
nervio de la construcción, sin convertirlo en una frase omnicomprensiva. El
cristal es el sistema mínimo

\[
 \bigl(\mathbb A_{1000},R_{729},c_t;
 C_9,\Gamma_9;\mathcal X_{\rm TPK}^{\rm enr}\bigr)
\]

que reúne una base irracional de entropía nula, un acarreo entero, una
holonomía finita levantada y una fibra con memoria. Sus realizaciones
\(\pi,\varphi,e,\alpha\) no son entradas del cristal; son caracteres
correlacionados del estado que éste prolonga. Las realizaciones dimensional,
fractal y excepcional se obtienen después mediante transductores propios.
El lector graduado \(\mathfrak E_9\) reúne el defecto \(28(3,-7,4)\), la
matriz \(D_{\rm syn}\) y los extremos de vacancia en la firma de coeficientes
de \(P_9\), de modo que la sincronización y la derivación de \(\alpha\)
quedan enlazadas sin introducir el valor de la constante como dato.

\section{La pendiente de compactificación es trascendente y su renormalización no es estacionaria}

La estructura aperiódica no depende sólo de que la palabra de vacancias
carezca de periodo. Su pendiente posee una propiedad aritmética más fuerte.
Recordemos que

\[
 \delta=\log_{10}\frac{10}{9},
 \qquad
 \rho=1-\delta=\log_{10}9.
\]

Ambas cantidades nacen de la razón interna entre la apertura ternaria de un
bloque de seis trits, \(3^6=729\), y la compactificación decimal
\(10^3=1000\). No se introducen como parámetros de una rotación elegida
externamente.

Primero, \(\delta\) es irracional. Si
\(\delta=p/q\in\mathbb Q\), con \(q>0\), entonces

\[
 10^{p/q}=\frac{10}{9}
 \quad\Longrightarrow\quad
 9^q10^p=10^q.
\]

La igualdad contradice la unicidad de la factorización prima: el miembro
izquierdo contiene el factor \(3\), mientras el derecho sólo contiene
\(2\) y \(5\). Segundo, \(\delta\) no puede ser algebraica irracional. En
ese caso, el teorema de Gelfond--Schneider haría trascendente a
\(10^\delta\), pero

\[
 10^\delta=\frac{10}{9}
\]

es algebraico. Se concluye

\[
 \boxed{\delta\ \text{es trascendente}.}
\]

Este reconocimiento clásico se aplica después de haber construido
\(\delta\) desde APP--TRIT--TPK; no actúa como generador del calendario.
Sus consecuencias estructurales son inmediatas. La fracción continua de
\(\delta\) no es eventualmente periódica, porque una fracción continua
eventualmente periódica representa un irracional cuadrático. Por ello, la
renormalización

\[
 [0;a_1,a_2,a_3,\ldots]
 =[0;21,1,5,1,6,2,3,1,1,8,\ldots]
\]

no cierra sobre una sustitución esturmiana estacionaria. La palabra admite
una descripción \(S\)-ádica dirigida por los cocientes parciales, pero esa
directiva no se repite periódicamente. El cristal es ordenado y recurrente,
sin ser el punto fijo de una única inflación cuadrática.

La misma conclusión puede formularse como ausencia de estabilizador modular
no trivial. Si

\[
 M=\begin{pmatrix}a&b\\c&d\end{pmatrix}\in
 \operatorname{GL}_2(\mathbb Z)
\]

fijara la pendiente por una transformación fraccionaria,

\[
 \frac{a\delta+b}{c\delta+d}=\delta,
\]

entonces

\[
 c\delta^2+(d-a)\delta-b=0.
\]

La trascendencia de \(\delta\) obliga a
\(c=b=0\) y \(d=a\). Como \(\det M=\pm1\), la acción proyectiva resultante
es la identidad. Por tanto,

\[
 \boxed{\operatorname{Stab}_{\operatorname{PGL}_2(\mathbb Z)}(\delta)
 =\{1\}.}
\]

La no repetición de la renormalización no elimina el orden de largo alcance.
La palabra mecánica

\[
 z_n=\mathbf 1_{[1-\delta,1)}
 \bigl(\{n\delta+\beta\}\bigr)
\]

es un conjunto modelo regular: procede de cortar una órbita reticular por
una ventana intervalar cuya frontera tiene medida nula. Su envolvente es
minimal y únicamente ergódico, posee complejidad \(p(m)=m+1\), entropía
topológica nula y espectro dinámico puro. Así queda precisada la expresión
«cristal aperiódico»: orden difractivo y recurrencia uniforme, junto con
ausencia de periodicidad traslacional y de renormalización estacionaria.

\textbf{Criterio de refutación.} Un periodo de la
palabra, una fracción continua eventualmente periódica o una transformación
modular no trivial que fije \(\delta\).

\section{Los tres canales como cortes de una misma ventana sobre \(729\) extensiones}

La sincronización de \(\pi\), \(e\) y \(\varphi\) puede describirse sin
reducirla a igualdad de censos. En el nivel \(t\), sea \(T_{c,t}\) el entero
representado por los primeros \(6t\) dígitos ternarios del canal
\(c\in\{\pi,e,\varphi\}\), y sea \(D_{c,t}\) el entero representado por sus
primeros \(K_t\) bloques decimales. Cada extensión local se identifica con
un entero

\[
 u\in\{0,1,\ldots,728\}.
\]

La celda ternaria correspondiente es

\[
 \left[
 \frac{729T_{c,t}+u}{3^{6(t+1)}},
 \frac{729T_{c,t}+u+1}{3^{6(t+1)}}
 \right),
\]

mientras el cilindro decimal es

\[
 \left[
 \frac{D_{c,t}}{1000^{K_t}},
 \frac{D_{c,t}+1}{1000^{K_t}}
 \right).
\]

Definamos la escala común y el desplazamiento interno

\[
 \Omega_t=\frac{3^{6(t+1)}}{1000^{K_t}},
 \qquad
 \xi_{c,t}=D_{c,t}\Omega_t-729T_{c,t}.
\]

Ambos cilindros se intersecan exactamente cuando

\[
 u<\xi_{c,t}+\Omega_t,
 \qquad
 u+1>\xi_{c,t}.
\]

Por consiguiente, el conjunto de extensiones supervivientes es

\[
 I_{c,t}
 =\{0,\ldots,728\}\cap
 \bigl[\lfloor\xi_{c,t}\rfloor,
       \lceil\xi_{c,t}+\Omega_t\rceil-1\bigr]_{\mathbb Z},
\]

y el censo determinado satisface

\[
 \boxed{N_c(t)=|I_{c,t}|.}
\]

La fórmula proporciona el mecanismo geométrico de la sincronización. Los
tres canales reciben la misma anchura \(\Omega_t\); sólo difieren en la
traslación \(\xi_{c,t}\) de la ventana. La igualdad
\(N_\pi=N_e=N_\varphi\) significa que tres intervalos trasladados cortan el
dominio local de \(729\) celdas con la misma cardinalidad. No implica que
los intervalos, las palabras ni sus historias sean iguales.

En el estado terminal inmediatamente anterior a la primera vacancia se
obtienen los intervalos enteros

\[
\begin{array}{c|c|c}
c&t=20:\ \text{intervalo bruto}&I_{c,20}\\
\hline
\pi&[407,408]&[407,408]\\
e&[172,173]&[172,173]\\
\varphi&[313,314]&[313,314].
\end{array}
\]

Los tres intervalos ocupan lugares diferentes, pero contienen dos celdas.
Ésta es la forma exacta del estado sincronizado
\((2,2,2)\): igualdad de medida discreta con memoria posicional distinta.

La vacancia siguiente conserva \(K_{21}=K_{20}=20\) y multiplica la escala
por \(729\). Los intervalos brutos pasan a ser

\[
\begin{array}{c|c|c|c}
c&t=21:\ \text{intervalo bruto}&I_{c,21}&N_c(21)\\
\hline
\pi&[118,1073]&[118,728]&611\\
e&[-429,526]&[0,526]&527\\
\varphi&[6,961]&[6,728]&723.
\end{array}
\]

Los tres intervalos brutos poseen la misma anchura entera, \(956\), pero
son recortados por bordes opuestos de la fibra finita. Si

\[
 \operatorname{Def}_*
 :=729(1,1,1)-\bigl(N_\pi(21),N_e(21),N_\varphi(21)\bigr),
\]

entonces

\[
 \operatorname{Def}_*=(118,202,6)
 =6(1,1,1)+28(4,7,0).
\]

El término diagonal \(6(1,1,1)\) registra la pérdida común de un bloque de
seis trits. Al cocientar por la supervivencia diagonal queda el divisor
orientado

\[
 [\operatorname{Def}_*]_{A_2}=28(4,7,0).
\]

Equivalentemente,

\[
 \partial N_{21}
 =(84,-196,112)
 =28(3,-7,4).
\]

La pendiente

\[
 \frac{N_e-N_\varphi}{N_\varphi-N_\pi}
 =-\frac{196}{112}
 =-\frac74
\]

es, por tanto, la razón de dos recortes orientados de una ventana común.
No es una fracción añadida al núcleo analítico.

El recálculo íntegro de los \(171\) niveles disponibles en la construcción de referencia
muestra dos unicidades finitas: \(t=21\) es la única vacancia precedida por
una cadena de cinco sincronizaciones consecutivas, y es el único nivel cuyo
defecto proyectivo en esa carta vale \(-7/4\). La afirmación es exacta en el
dominio de definición; su extensión coinductiva a profundidad arbitraria
deberá conservar explícitamente el estado marcado.

\textbf{Criterio de refutación.} Que los censos no
coincidan con las intersecciones de intervalos, que las tres anchuras brutas
en \(t=21\) difieran o que aparezca otra pendiente \(-7/4\) en el dominio
declarado.

\section{Matriz de tránsito especular y cierre conjunto de \(21/22\), \(6/7\), \(137\) y la red \(369\)–\(126\)}

Las cuatro fases que preceden a la primera sincronización triple no deben
leerse como cuatro anécdotas. Sus amplitudes de defecto forman la matriz

\[
 D_{\rm syn}
 =\begin{pmatrix}
 85&112\\
 41&52
 \end{pmatrix}.
\]

El censo sincronizado de la novena fase es \(N_9=43\). Por tanto,

\[
 \det D_{\rm syn}=-172=-4\cdot43
\]

define internamente

\[
 a:=-\frac{\det D_{\rm syn}}{N_9}=4.
\]

Por una rama independiente, la inducción de la palabra de vacancias sobre
sus huecos cortos produce el espectro secundario
\(\mathcal R_{\rm sec}=\{6,7\}\). Sea

\[
 b:=\max\mathcal R_{\rm sec}=7.
\]

La primera reapertura satisface entonces

\[
 t_*=3b=21,
 \qquad
 K_*=3b-1=20,
\]

donde \(3\) es la cardinalidad de los regímenes de TRIT. Además,

\[
 \left\lfloor1000\delta\right\rfloor=45,
 \qquad
 \left\lceil1000\delta\right\rceil=46
 =2(3b)+a.
\]

La descomposición de recorte de la sección anterior adopta ahora la forma
cerrada

\[
 \boxed{
 \operatorname{Def}_*
 =6(1,1,1)+ab(a,b,0).}
\]

En efecto,

\[
 6(1,1,1)+4\cdot7(4,7,0)
 =(118,202,6).
\]

Esta igualdad es un cuadrado de compatibilidad entre dos extracciones
independientes: \(a=4\) procede del determinante de sincronización,
\(b=7\) procede del reloj inducido y el recorte común los vuelve a encontrar
como pesos de borde. La fracción \(-b/a=-7/4\) queda así determinada antes
de resolver ninguna ecuación para \(\alpha\).

La misma matriz reúne la renormalización diofántica. Los dos retornos
secundarios dan

\[
 q_6=22\cdot6-1=131,
 \qquad
 q_7=22\cdot7-1=153,
\]

y satisfacen

\[
 6q_7-7q_6=1.
\]

Dentro de \(D_{\rm syn}\), esos denominadores reaparecen como

\[
 q_6=(85+41)+5=131,
 \qquad
 q_7=112+41=153,
\]

donde \(5\) es la rama superior del semidesbalance \(4/5\) de
\(\widetilde W_{100}\). Finalmente,

\[
 \boxed{\operatorname{tr}D_{\rm syn}=137=q_6+6.}
\]

Así, \(137\) no se obtiene aquí mediante redondeo de
\(\alpha^{-1}\). Es una dirección entera del tránsito de sincronización y,
simultáneamente, el denominador de retorno \(131\) prolongado por un bloque
de seis trits. Esto no sustituye las dos construcciones de \(\alpha\): fija
un precursor entero común que el cierre dodecafásico y el núcleo analítico
deben transportar.

La red radical--electrónica también se reconstruye desde
\(D_{\rm syn}\). Definamos

\[
 A_{\rm surv}:=85+41=126,
 \qquad
 B_{\rm bord}:=112-85=27,
\]

y

\[
 T_{\rm tri}:=3(A_{\rm surv}+6)=396.
\]

El lector matricial

\[
 \mathcal R(D_{\rm syn})
 :=
 \begin{pmatrix}
 T_{\rm tri}-B_{\rm bord}&T_{\rm tri}\\
 T_{\rm tri}-A_{\rm surv}-B_{\rm bord}&
 T_{\rm tri}-A_{\rm surv}
 \end{pmatrix}
\]

produce exactamente

\[
 \boxed{
 \mathcal R(D_{\rm syn})
 =\begin{pmatrix}369&396\\243&270\end{pmatrix}.}
\]

Por consiguiente, los cuatro vértices numéricos del paralelogramo ya no se
recogen como un conjunto independiente: proceden de un lector entero de la
sincronización de canales, del bloque de seis trits y de la amplificación
trítica. La identificación de cada vértice con su objeto radical o
electrónico conserva, no obstante, el dominio y el transductor de sus
dominios de procedencia; la igualdad de la matriz numérica no borra esas fibras.

\textbf{Criterio de refutación.} Que \(a\) no sea entero, que la palabra
inducida no produzca \(b=7\), que falle la ley de recorte o que
\(\mathcal R(D_{\rm syn})\) no reproduzca los cuatro vértices.

\section{Factorización interna de todos los coeficientes del jet \(P_9\)}

Los resultados anteriores permiten reunir, por primera vez dentro de este
desarrollo, la genealogía de todos los coeficientes del jet sin introducir
su raíz. Sea

\[
 \omega:=2\pi,
 \qquad
 m:=3,
 \qquad
 a:=4,
 \qquad
 b:=7,
\]

y definamos

\[
 t_*=mb=21,
 \qquad
 K_*=mb-1=20,
\]

\[
 v_-:=\lfloor1000\delta\rfloor=45,
 \qquad
 v_+:=\lceil1000\delta\rceil=46.
\]

La rama inferior satisface a la vez

\[
 v_-=45=\det
 \begin{pmatrix}7&2\\2&7\end{pmatrix}.
\]

La diferencia interna de la segunda fila de \(D_{\rm syn}\) da

\[
 r_{\rm dod}:=52-41=11,
\]

que coincide con el rango transversal dodecafásico ya demostrado. La
semicorona de fase aporta

\[
 f_{\rm ph}=8\binom62=120,
\]

y

\[
 v_-r_{\rm dod}=45\cdot11=495.
\]

Sobre la firma de invariantes

\[
 \mathcal I_9
 =(\omega,m,a,b,K_*,t_*,v_-,v_+,f_{\rm ph},r_{\rm dod}),
\]

definimos el lector graduado

\[
\begin{aligned}
 \mathfrak E_9(\mathcal I_9)(x)
 :={}&
 \omega x-\frac ba x^2+\omega^{-1}x^3
 +\frac{x^4}{K_*}-\frac{2x^5}{t_*}
 -\frac{x^6}{v_+}\\
 &-\frac{x^7}{f_{\rm ph}}
 +\frac{x^8}{v_-}
 +\frac{2x^9}{v_-r_{\rm dod}}.
\end{aligned}
\]

La sustitución de los invariantes produce

\[
\boxed{
\begin{aligned}
\mathfrak E_9(\mathcal I_9)(x)
={}&2\pi{}x-\frac74x^2
+\frac{x^3}{2\pi}
+\frac{x^4}{20}-\frac{2x^5}{21}-\frac{x^6}{46}\\
&-\frac{x^7}{120}+\frac{x^8}{45}
+\frac{2x^9}{495}
=P_9(x).
\end{aligned}}
\]

La fórmula revela una relación que quedaba oculta al presentar por separado
\(P_6\) y \(T_{7:9}\): el par de vacancias \(45/46\) atraviesa la frontera
entre ambos. La rama superior aparece en
\(-x^6/46\), mientras la rama inferior reaparece en \(+x^8/45\) y,
acoplada al rango \(11\), en \(2x^9/495\). La prolongación no añade una
segunda aritmética; continúa el mismo calendario equilibrado mediante el
resolvente nilpotente.

El diagrama tipado es

\[
 \mathcal X_{\rm TPK}^{\rm enr,*}
 \xrightarrow{\ \mathcal A_*\ }
 \mathcal I_9
 \xrightarrow{\ \mathfrak E_9\ }
 \mathbb Q(\pi)[x]/(x^{10}),
\]

donde el asterisco indica que el objeto conserva el origen coinductivo, el
orden de los canales y la orientación de sus hojas. El extractor
\(\mathcal A_*\) realiza las operaciones demostradas:

\[
 D_{\rm syn}\mapsto a,\qquad
 \mathcal R_{\rm sec}\mapsto b,\qquad
 \operatorname{Def}_*\mapsto(a,b),
\]

\[
 \delta\mapsto(v_-,v_+),\qquad
 B_c\mapsto v_-,\qquad
 (D_3,D_4)\mapsto r_{\rm dod}.
\]

Todo morfismo que preserve el estado marcado, las ventanas de intersección,
la orientación de canales, el reloj de capacidad y la involución two-way
preserva \(\mathcal I_9\), y por ello hace conmutar el lector. Lo que se ha
cerrado es la factorización interna de la firma de coeficientes. La
afirmación más fuerte de que \(\mathfrak E_9\) sea el único lector posible
en una categoría no apuntada exige clasificar todos los automorfismos que
puedan olvidar o permutar el origen y la orientación. Esa reserva es
concreta y no reabre ni la raíz simple de \(E_{21/22}\), ni el resolvente
nilpotente, ni la genealogía de los coeficientes.

\textbf{Criterio de refutación.} Que alguno de los coeficientes requiera el valor
de la raíz, que dos extracciones declaradas produzcan valores distintos o
que un automorfismo permitido cambie \(\mathcal I_9\).

\section{El factor \(28\), la proyección exacta \(36\times28\) y la partición en siete tétradas}

La escala primitiva del primer defecto es

\[
 \lambda_*:=\gcd(84,196,112)=28=ab.
\]

En el catálogo construido de la región de \(\pi\), la microfibra orientada
posee multiplicidad

\[
 M(w_\pi)=432+4\cdot144=1008
\]

y admite las factorizaciones

\[
 1008=36\cdot28
 =\binom92\binom82
 =7\cdot144.
\]

Por tanto,

\[
 \boxed{M(w_\pi)=36\,\lambda_*}
\]

es una compatibilidad escalar exacta entre el defecto de recorte de la
primera reapertura y la multiplicidad de la microfibra de \(\pi\). El
factor \(36\) pertenece además a \(R_{36}\) y al conjunto de pares de una
nónada; \(28\) es el número de pares de un conjunto de ocho elementos.

La primera lectura de esta igualdad como mera compatibilidad escalar era
incompleta. La construcción de referencia  ya construye sobre la microfibra

\[
 \mathcal R_{010211}
 =\{(P,T):w_6(U_6(P,T))=010211\}
\]

una proyección explícita. Para

\[
 P=(i,j,d_P),\qquad T=(k,l,d_T),
\]

se fija el calibre direccional

\[
 \beta(N)=1,\quad\beta(E)=2,\quad
 \beta(S)=3,\quad\beta(O)=4
\]

y se define

\[
 x(P)=5(i-j)\pmod 9,
 \qquad
 h(T)=5\beta(d_T)\pmod9,
\]

\[
 \boxed{
 q_\beta(P,T)=\{x(P)-h(T),x(P)+h(T)\}.}
\]

El mapa

\[
 q_\beta:\mathcal R_{010211}\longrightarrow
 \binom{\mathbb Z/9\mathbb Z}{2}
\]

es sobreyectivo, tiene base de cardinal \(36\) y todas sus fibras poseen
cardinal \(28\). Cada fibra recibe, en el orden de las cinco regiones de
\(\pi\), exactamente

\[
 12+4+4+4+4=28
\]

rutas. No se obtiene dividiendo \(1008\) por \(36\): el certificado
reconstruye las \(104,976\) parejas de semillas, calcula cada imagen y prueba
la uniformidad de las fibras.

La traslación aditiva de \(C_9\) y su reflexión preservan \(U_6\) y hacen a
\(q_\beta\) equivariante respecto de la acción diedral sobre los pares. En el
cursor multiplicativo, dos involuciones generan un \(V_4\) libre que conserva
la firma completa. El grupo

\[
 H_0=C_9\times V_4
\]

El cociente

\[
 \Omega=\mathcal R_{010211}/H_0,
 \qquad |\Omega|=28,
\]

posee una partición intrínseca adicional:

\[
 \Omega=
 \mathcal B_3\sqcup\mathcal B_6\sqcup\mathcal B_9
 \sqcup\mathcal L_1\sqcup\mathcal L_2
 \sqcup\mathcal L_3\sqcup\mathcal L_4,
\]

con siete bloques de cuatro elementos. Los tres bloques
\(\mathcal B_3,\mathcal B_6,\mathcal B_9\) son transversales y están
distinguidos por el residuo multiplicativo \(3,6,9\); los cuatro bloques
\(\mathcal L_j\) proceden de las cuatro regiones laterales. El vector de
censo reducido de esta partición es

\[
 v_\Omega=(3,-7,4)\in A_2.
\]

La primera reapertura de los canales satisface, exactamente,

\[
 \boxed{
 \partial N_{21}=28v_\Omega.}
\]

Así, la dirección \((3,-7,4)\) no sólo contiene la pendiente \(-7/4\): es la
relación de conservación entre tres tétradas transversales, siete tétradas
totales y cuatro tétradas laterales. Además,

\[
 28=7\cdot4,
 \qquad
 1008=36\cdot7\cdot4.
\]

Los enteros \(a=4\) y \(b=7\), obtenidos antes por el determinante de
sincronización y el retorno inducido, reaparecen ahora como tamaño de bloque y
número de bloques de la microfibra.

Esta identificación posee una rigidez local decisiva. El grupo integral de
isometrías de \(A_2\), representado por permutaciones de las tres coordenadas
y cambio global de signo, tiene orden \(12\). Como las coordenadas
\(3,-7,4\) son distintas y el multiconjunto \(\{3,-7,4\}\) no coincide con
su negativo, sólo la identidad fija \(v_\Omega\). Por tanto,

\[
 \boxed{
 \operatorname{Stab}_{\operatorname{Aut}(A_2)}(v_\Omega)=\{1\}.}
\]

Una vez caracterizada la primera reapertura, el propio defecto recupera el
orden de los tres canales y su orientación: no necesita una marca externa.

La frontera octádica queda tipada por las siete tétradas. Tras escoger un
calibre, cartas biyectivas

\[
 \Phi:\Omega\overset\sim\longrightarrow
 \binom{\mathbb F_2^3}{2}
\]

y con ellas una incidencia \(J(8,2)\) exacta. Sin embargo, esas cartas no son
seleccionadas por las relaciones fuente ensayadas. La reflexión que estabiliza
cada fibra actúa libremente con tipo de ciclos \(2^{14}\), mientras que toda
involución de \(\operatorname{Aut}J(8,2)\cong S_8\) fija al menos cuatro
pares. Se sigue el no-go probado:

\[
 \boxed{
 \text{no existe una familia de cartas }J(8,2)
 \text{ equivariante bajo todo }D_9\times V_4.}
\]

El dato orientado \(v_\Omega\), cuyo estabilizador es trivial, aporta la
ruptura de reflexión que exige este teorema de obstrucción. En consecuencia,
las cartas Johnson sobre cada fibra \(F_b=q_\beta^{-1}(b)\) son cartas de
calibre conjugadas y no una selección canónica bajo la acción completa
\(D_9\times V_4\). La proyección \(36\times28\) permanece intrínseca; la
elección de una carta octádica no forma parte de su definición.

También queda corregida una ambigüedad temporal. El entero \(1008\) no es
el periodo del reloj dodecafásico:

\[
 1008\equiv36\pmod{108}.
\]

Tras \(1008\) pasos retorna la fase nonádica, pero el reloj \(C_{108}\)
avanza \(36\) posiciones. En este frente, \(1008\) es una multiplicidad de
microfibra y una incidencia combinatoria; \(108\) es el periodo temporal
finito. La proyección \(q_\beta\) explica el producto \(36\times28\); una
carta octádica de calibre y el mapa de reloj son operaciones distintas.

\textbf{Criterio de refutación.} Que el barrido de semillas cambie una fibra o la partición
\(7\times4\), que una isometría no trivial de \(A_2\) fije \(v_\Omega\),
que exista una carta Johnson completamente \(D_9\times V_4\)-equivariante o
que se use \(1008\) como periodo de \(C_{108}\).

\section{Dos relojes, una memoria integrada y dos módulos espectrales}

Sea

\[
 Z_t:=\sum_{j=0}^{t}z_j
 =\lfloor(t+1)\delta\rfloor.
\]

Como \(\rho=1-\delta\) y \((t+1)\delta\notin\mathbb Z\),

\[
\begin{aligned}
 K_t
 &=\lfloor(t+1)\rho\rfloor\\
 &=\lfloor(t+1)-(t+1)\delta\rfloor\\
 &=t-\lfloor(t+1)\delta\rfloor
 =t-Z_t.
\end{aligned}
\]

Si

\[
 \theta_t=[t]_q
 \quad\text{y}\quad
 g_t=[K_t]_q,
\]

entonces

\[
 \boxed{g_t=\theta_t-[Z_t]_q.}
\]

La fase de capacidad es la fase uniforme corregida por la memoria integrada
de vacancias. Esta identidad expresa exactamente la diferencia entre
retorno de fase y retorno de estado: conocer \(\theta_t\) no determina
\(g_t\) si se ha olvidado \(Z_t\).

La distinción produce dos módulos espectrales legítimos. En el reloj de
bloques, la traslación

\[
 S_q(\theta,u)=(\theta+1,u+\delta)
\]

sobre \(C_q\times\mathbb T\) posee frecuencias aditivas

\[
 \mathcal M_q^{\rm blk}
 =\frac1q\mathbb Z+\delta\mathbb Z
 =\frac1q\mathbb Z+\rho\mathbb Z
 \pmod1.
\]

En particular,

\[
 \mathcal M_9^{\rm blk}
 =\frac19\mathbb Z+\log_{10}9\,\mathbb Z,
\]

y

\[
 \mathcal M_{108}^{\rm blk}
 =\frac1{108}\mathbb Z+\log_{10}9\,\mathbb Z.
\]

Éste es el sentido exacto en el que el logaritmo decimal interno aparece
como módulo espectral: no como entropía positiva, sino como generador
irracional de un grupo denso de frecuencias puras.

En el reloj pegado de capacidad, la dinámica \(F_q\) de la sección 49 es
conjugada a

\[
 y\longmapsto y+\rho
 \quad\text{sobre}\quad
 \mathbb R/q\mathbb Z.
\]

Sus frecuencias son

\[
 \mathcal M_q^{\rm cap}
 =\frac{\rho}{q}\mathbb Z\pmod1.
\]

No hay contradicción entre ambos módulos: pertenecen a dos álgebras de
observables del mismo estado holonómico integral. La primera conserva el índice de
bloque; la segunda absorbe en la coordenada finita el acarreo acumulado.
La relación \(g_t=\theta_t-Z_t\) es el cambio de calibre que las compara.

Esto precisa también dos retornos mencionados anteriormente. El denominador
\(170748=1581\cdot108\) cierra el reloj de bloques, pero

\[
 K_{170748-1}\equiv71\pmod{108},
\]

de modo que no cierra el reloj de capacidad. Asimismo, \(1008\) cierra
\(C_9\), pero desplaza \(C_{108}\) en \(36\). El estado holonómico integral retiene
en ambos casos la altura y la firma, por lo que ninguno de esos retornos
equivale al reinicio de la dinámica.

\textbf{Criterio de refutación.} Que
\(K_t\ne t-Z_t\), que se atribuya el mismo espectro a dos dominios de
observables distintos o que un retorno finito borre la memoria.
 
\chapter{Conexión nonádica conjunta y biografía del retorno con memoria}
\label{chap:83-07}

La conexión se construye a partir de la clasificación orbital completa y de
su primera elevación coinductiva. El desarrollo siguiente fija las tres
historias iniciales, las fronteras orientadas y la cadena
\(w_6\to w_{30}\to R_{36}\) antes de introducir la holonomía nonádica.

\begin{HMTScientificOwnerSubordinated}
% Unidad U008. Clasificación orbital y elevación; redacción autónoma.

\chapter{Clasificación orbital, palabras iniciales y 1.ª frontera}
\label{chap:orbitas-elevacion-r36}

La reducción ternaria del emisor realiza 243 palabras dentro de la fibra
ambiente \(\mathbb F_3^6\). Esta unidad no escoge tres palabras por
comparación decimal. Construye primero la acción de simetría del catálogo,
clasifica sus órbitas mediante invariantes enteros y sólo después orienta
los tres representantes estructurales.

\section{Acción diedral sobre \(6\) coordenadas}

Para \(u=(u_1,u_2,u_3,u_4,u_5,u_6)\in\mathbb F_3^6\), definimos
\begin{align*}
 r(u)&=(u_3,u_1,u_2,u_5,u_6,u_4),\\
 s(u)&=(u_4,u_5,u_6,u_1,u_2,u_3).
\end{align*}
Se verifica
\[
 r^3=s^2=1,
 \qquad srs=r^{-1}.
\]
Por tanto, \(r,s\) generan una acción de \(D_3\).

La misma permutación actúa sobre los seis bloques de cada emisión
\(U=(U_1,\ldots,U_6)\). La proyección \(\Psi\) satisface
\[
 \Psi(gU)=g\Psi(U),
 \qquad g\in D_3.
\]

\begin{theorem}[Descomposición orbital del catálogo]
\label{thm:catalogo-descomposicion-d3}
Las 243 palabras realizadas se descomponen en 43 órbitas:
\[
 38\text{ órbitas de cardinal }6,
 \qquad
 5\text{ órbitas de cardinal }3.
\]
En particular,
\[
 38\cdot6+5\cdot3=243.
\]
\end{theorem}

\begin{proof}
Se aplica a cada palabra el conjunto
\[
 \{u,ru,r^2u,su,sru,sr^2u\}
\]
y se eliminan repeticiones. La enumeración exhaustiva del catálogo produce
el histograma indicado. La equivariancia de \(\Psi\) garantiza que ninguna
órbita abandona la imagen realizada.
\end{proof}

\section{Invariantes de fibra}

Para una palabra realizada \(w\), sean
\[
 \mathcal F_w:=\{U\in\operatorname{im}T_{\min}:\Psi(U)=w\},
\]
\[
 n(w):=|\mathcal F_w|,
 \qquad
 M(w):=\sum_{U\in\mathcal F_w}\operatorname{mult}(U),
\]
y
\[
 c(w)=\bigl(\#\{j:w_j=0\},\#\{j:w_j=1\},\#\{j:w_j=2\}\bigr).
\]
El catálogo enriquecido conserva además, para cada emisión, una coordenada
diagonal \(\operatorname{diag}_c(U)\in\{0,\ldots,8\}\), una orientación
cardinal \(h_+(U)\) y la firma multiplicativa \(E_{\rm sig}(U)\). Todas
ellas se calculan de la ruta de semilla; no se leen de una constante
objetivo. Su algoritmo componente a componente se imprimirá junto al catálogo
completo en el apéndice de certificación, y forma parte del estado de entrada
de los predicados siguientes.

\begin{definition}[Órbita de clausura]
\label{def:orbita-clausura}
Una órbita \(\mathcal O\) es de clausura si tiene cardinal seis y, para
cada \(w\in\mathcal O\),
\[
 n(w)=5,
 \qquad
 M(w)=1008,
\]
\[
 \{\operatorname{mult}(U):U\in\mathcal F_w\}
 =\{432,144,144,144,144\},
\]
y el censo trítico común es \(c(w)=(2,3,1)\).
\end{definition}

\begin{definition}[Órbitas mínimas radicales]
\label{def:orbitas-minimas-radicales}
Una órbita \(\mathcal O\) es mínima radical si tiene cardinal seis, cada
palabra posee una única emisión de multiplicidad \(144\) y
\[
 \{\operatorname{diag}_c(U):\Psi(U)\in\mathcal O\}=\{0,3,6\}.
\]
Entre ellas, se denomina órbita de propagación a la que tiene censo
\((2,3,1)\), y órbita de autoescala a la que tiene censo \((2,2,2)\).
\end{definition}

\begin{theorem}[Unicidad de las tres órbitas estructurales]
\label{thm:unicidad-tres-orbitas-estructurales}
En el catálogo de 468 emisiones existe una única órbita de clausura, una
única órbita mínima radical de propagación y una única órbita mínima radical
de autoescala.
\end{theorem}

\begin{proof}
Se enumeran las 43 órbitas del \cref{thm:catalogo-descomposicion-d3}. Para
cada una se calculan \(n(w)\), \(M(w)\), el multiconjunto de
multiplicidades, el censo \(c(w)\) y el soporte diagonal. Los predicados de
las \cref{def:orbita-clausura,def:orbitas-minimas-radicales} dejan una sola
órbita en cada clase. La prueba es exhaustiva sobre un conjunto finito y no
utiliza cifras decimales de los caracteres analíticos.
\end{proof}

Las tres órbitas son
\begin{align*}
\mathcal O_{\rm cl}
 &=\{001112,010211,100121,112001,121100,211010\},\\
\mathcal O_{\rm prop}
 &=\{011120,012110,101201,110012,120011,201101\},\\
\mathcal O_{\rm aut}
 &=\{002112,020211,112002,121200,200121,211020\}.
\end{align*}

\section{Orientación interna de cada órbita}

Fijamos el calibre interno
\[
 \operatorname{diag}_c(U)=6,
 \qquad h_+(U)=ES.
\]
Una palabra de una órbita es admisible en ese calibre si alguna emisión de
su fibra satisface ambas condiciones.

\begin{proposition}[Representante orientado único]
\label{prop:representante-orientado-unico}
Cada una de las tres órbitas estructurales contiene una única palabra
admisible en el calibre \((6,ES)\). Son
\[
 \boxed{
 w_6^\pi=010211,
 \qquad
 w_6^e=201101,
 \qquad
 w_6^\varphi=121200.}
\]
\end{proposition}

\begin{proof}
La enumeración de las fibras del catálogo cuenta las palabras que contienen
una emisión con \(\operatorname{diag}_c=6\) y \(h_+=ES\). En cada una de
las tres órbitas el número de coincidencias es uno. Los superíndices
\(\pi,e,\varphi\) nombran desde este punto las funciones estructurales de
clausura, propagación y autoescala; no participaron en la selección.
\end{proof}

\section{Microfibra de clausura}

La fibra del representante \(w_6^\pi\) consta de cinco emisiones:
\[
\begin{array}{c|c|c|c|c}
U_6&E_{\rm sig}&\operatorname{mult}&\operatorname{diag}_c&\text{estrato}\\\hline
501|614|498|169|272|272&555555&432&4&\perp\\
810|923|870|169|272|272&339555&144&6&++\\
810|983|810|169|272|272&393555&144&6&++\\
870|923|810|169|272|272&933555&144&6&++\\
870|923|810|418|674|521&933717&144&5&--
\end{array}
\]
Así,
\[
 432+4\cdot144=1008,
 \qquad
 5=1_\perp+3_{++}+1_{--}.
\]
Las emisiones unipuntuales orientadas de los otros dos canales son
\[
 U_e=850|963|890|149|252|212,
\]
\[
 U_\varphi=830|943|830|109|252|252.
\]
La masa \(1008\) de esta microfibra no es el periodo del calendario
\(C_{108}\); ambas cifras pertenecen a dominios diferentes.

\section{Matrices de elevación ternaria}

Sobre vectores fila de \(\mathbb F_3^6\), definimos
\[
L_0=
\begin{pmatrix}
2&2&2&1&2&1\\
2&1&2&2&1&1\\
1&0&0&1&0&2\\
0&0&1&1&1&0\\
2&2&2&0&2&1\\
2&1&2&1&0&0
\end{pmatrix},
\qquad
L_1=
\begin{pmatrix}
2&1&2&0&2&2\\
1&2&1&1&2&0\\
1&2&1&1&1&2\\
0&1&0&0&2&1\\
2&0&2&1&0&1\\
0&2&2&2&1&1
\end{pmatrix}.
\]
El calendario de elevación es \((L_0,L_0,L_1,L_1)\).

\begin{definition}[Recurrencia de bloques y concatenación]
\label{def:recurrencia-bloques-l0-l1}
Para cada canal \(\chi\in\{\pi,e,\varphi\}\), sea
\(b_0^\chi=w_6^\chi\) y
\[
 b_{k+1}^\chi=b_k^\chi L_{\epsilon_k}\pmod3,
 \qquad
 (\epsilon_0,\epsilon_1,\epsilon_2,\epsilon_3)=(0,0,1,1).
\]
La palabra acumulada es
\[
 w_{30}^\chi
 =b_0^\chi\Vert b_1^\chi\Vert b_2^\chi
  \Vert b_3^\chi\Vert b_4^\chi.
\]
\end{definition}

Esta formulación distingue el bloque corriente \(b_k\in\mathbb F_3^6\) de
la palabra concatenada \(w_{6(k+1)}\). Una palabra de longitud \(6k\) no
se multiplica por una matriz \(6\times6\).

\begin{theorem}[Biografías ternarias de treinta posiciones]
\label{thm:palabras-w30}
La recurrencia de la \cref{def:recurrencia-bloques-l0-l1} produce
\begin{align*}
w_{30}^{\pi}
 &=010211|012222|010211|002111|110221,\\
w_{30}^{e}
 &=201101|121221|102011|012222|102011,\\
w_{30}^{\varphi}
 &=121200|112202|121020|010210|010200.
\end{align*}
\end{theorem}

\begin{proof}
Se multiplican sucesivamente los cinco vectores fila por las matrices
indicadas, reduciendo cada coordenada módulo tres después de cada producto.
La concatenación de los cinco bloques calculados da las tres expresiones.
Cada igualdad puede verificarse componente a componente mediante 24
productos escalares por canal.
\end{proof}

Los índices ambientales de los cinco bloques son
\[
\begin{array}{c|rrrrr}
&b_0&b_1&b_2&b_3&b_4\\\hline
\pi&103&161&103&67&349\\
e&523&457&301&161&301\\
\varphi&450&398&438&102&99
\end{array},
\]
donde cada palabra de seis trits se interpreta como un entero en base tres.

\section{1.ª frontera conjunta}

La frontera se obtiene por una relación finita que conserva eje, agregado,
saturación, ocupación de columnas y neutralidad dual. La matriz de incidencia
utilizada en esta relación es
\begin{equation}
 A_W=
 \begin{pmatrix}
 0&1&1&1&1&1\\
 1&0&1&2&2&1\\
 1&1&0&1&2&2\\
 1&2&1&0&1&2\\
 1&2&2&1&0&1\\
 1&1&2&2&1&0
 \end{pmatrix}\in M_6(\mathbb F_3).
 \label{eq:u008-aw-frontera}
\end{equation}
Si \(u_4^\chi\) es el quinto bloque de \(w_{30}^\chi\), el eje y el
agregado previos a separar los dos canales laterales son
\begin{equation}
 u_5^\varphi=u_4^e=102011,
 \qquad
 u_5^\pi+u_5^e=210112,
 \qquad
 u_4^\varphi A_WL_1^3=111101.
 \label{eq:u008-eje-agregado-r36}
\end{equation}
De ellos se calculan
\begin{equation}
 q=012120,\qquad a=111101,\qquad
 qA_W=012000,\qquad aA_W=120210.
 \label{eq:u008-q-a-r36}
\end{equation}
La profundidad seis impone la saturación \(c=111111\); la ocupación mínima
compatible es \(\operatorname{colw}=223232\); y la neutralidad dual exige
\begin{equation}
 (BA_W)\mathbf1=000.
 \label{eq:u008-neutralidad-dual-r36}
\end{equation}

\begin{proposition}[Enumeración de la primera frontera]
\label{prop:u008-enumeracion-r36}
Las condiciones \eqref{eq:u008-eje-agregado-r36}--
\eqref{eq:u008-neutralidad-dual-r36} dejan exactamente dos matrices:
\begin{equation}
 B_-=
 \begin{pmatrix}021222\\222220\\102011\end{pmatrix},
 \qquad
 B_+=
 \begin{pmatrix}222220\\021222\\102011\end{pmatrix}.
 \label{eq:u008-dos-fronteras}
\end{equation}
La involución especular intercambia sus dos primeras filas. Sus cartas
internas son, respectivamente,
\begin{equation}
 B_-\longmapsto(789,048,894),
 \qquad
 B_+\longmapsto(793,045,894).
 \label{eq:u008-cartas-frontera}
\end{equation}
La orientación APP del cilindro semiabierto selecciona \(B_+\).
\end{proposition}

\begin{proof}
El eje fija la tercera fila y el agregado fija la suma de las dos primeras.
La saturación y la ocupación eliminan las distribuciones que no llenan las
seis columnas. La enumeración de las soluciones de las seis ecuaciones
lineales de \eqref{eq:u008-neutralidad-dual-r36} deja las dos matrices de
\eqref{eq:u008-dos-fronteras}; la carta semiabierta y la orientación
distinguen la segunda sin consultar una expansión decimal objetivo.
\end{proof}

Definimos, por tanto,
\begin{equation}
 R_{36}:=B_+=
 \begin{pmatrix}
 222220\\
 021222\\
 102011
 \end{pmatrix},
 \qquad
 \operatorname{ind}_3(R_{36})=(726,215,301).
 \label{eq:u008-r36-construida}
\end{equation}
El objeto \(R_{36}\) contiene tres bloques ordenados de seis trits. No es el
cociente predictivo \(C_{36}^{\rm pred}\) de la
\cref{thm:tpk-cociente-predictivo-36} ni el exceso de 36 clases del calibre
de semillas.

\section{Levantamiento residuo--cociente}

Sea \(A\in M_6(\mathbb Z)\) el levantamiento entero de \(L_0\):
\[
 A=
\begin{pmatrix}
2&2&2&1&2&1\\
2&1&2&2&1&1\\
1&0&0&1&0&2\\
0&0&1&1&1&0\\
2&2&2&0&2&1\\
2&1&2&1&0&0
\end{pmatrix},
\qquad \det A=1.
\]
Para \(x_k\in\mathbb Z^6\), definimos
\[
 y_k=x_kA,
 \qquad
 u_k=y_k\bmod3\in\{0,1,2\}^6,
 \qquad
 x_{k+1}=\frac{y_k-u_k}{3}.
\]
Entonces
\begin{equation}
 x_kA=u_k+3x_{k+1}.
\label{eq:bockstein-hensel-local}
\end{equation}

\begin{theorem}[Reversibilidad del levantamiento]
\label{thm:reversibilidad-bockstein-hensel}
El par \((u_k,x_{k+1})\) queda determinado unívocamente por \(x_k\), y
\[
 x_k=(u_k+3x_{k+1})A^{-1}.
\]
\end{theorem}

\begin{proof}
La división euclídea por tres en cada coordenada de \(x_kA\) determina
residuo y cociente. Como \(\det A=1\), la inversa \(A^{-1}\) tiene
coeficientes enteros. Multiplicar \cref{eq:bockstein-hensel-local} por ella
recupera \(x_k\).
\end{proof}

En el canal \(e\), el bloque visible \(102011\) reaparece dos veces, pero
los preestados módulo 27 son
\[
 (25,11,7,17,8,16)
 \neq
 (7,8,4,23,26,7),
\]
y los cocientes posteriores son
\[
 (6,13,9,2,13,1)
 \neq
 (15,0,3,19,23,25).
\]
Así se verifica de forma explícita que el retorno visible de una palabra no
es el retorno del estado holonómico integral.

\section{Obstrucciones de la selección orbital y de la elevación \(w_6\to w_{30}\to R_{36}\): unicidad, calibre y reversibilidad Hensel}

\begin{criterion}[Criterios de refutación de la selección orbital]
La unidad queda refutada si:
\begin{enumerate}
 \item la acción \(D_3\) abandona el catálogo o cambia multiplicidades;
 \item alguno de los tres predicados estructurales deja más de una órbita;
 \item el calibre \((6,ES)\) selecciona dos palabras en una órbita;
 \item se usa una cifra objetivo para definir los predicados;
 \item se multiplica la palabra acumulada \(w_{6k}\) por \(L_0\) o
       \(L_1\) en lugar del bloque \(b_k\);
 \item \(\det A\neq1\) o falla la inversa del levantamiento;
 \item dos residuos visibles iguales se declaran estados completos iguales.
\end{enumerate}
\end{criterion}

La primera frontera conjunta ya está construida. La unidad siguiente
clasificará las nueve componentes de fase, demostrará el retorno nonádico
sin reinicio y construirá el espejo orientado con su resolución local
\(3\to1\).
 \end{HMTScientificOwnerSubordinated}

% Unidad U009. Conexión nonádica conjunta; redacción autónoma.

\section{Conexión nonádica conjunta, compresión y memoria terminal}
\label{p12-83-c7-s1:chap:conexion-nonadica-conjunta}

Las nueve fases que aparecen después de \(R_{36}\) no son nueve pruebas
independientes. Son las nueve componentes locales de una única conexión
sobre estados holonómicos integrales. Su composición alrededor de una vuelta define
una holonomía. El retorno de la etiqueta de fase conserva la continuidad de
prefijo, acarreo, frontera, sombra de Witt y memoria.

\subsection{Transporte del estado holonómico integral por la holonomía nonádica y actualización de memoria}

En el nivel \(K\), escribimos
\[
 \widehat x_K=
 (B_K,C_K,p_K,c_K,\Sigma_K,W_K,g(K),q_{{\rm mem},K}),
\]
donde:
\begin{enumerate}
 \item \(B_K=(u_K^\pi,u_K^e,u_K^\varphi)\) es el bloque visible de tres
       palabras de seis trits;
 \item \(C_K\) es el cilindro ternario compatible;
 \item \(p_K\) es el prefijo completo;
 \item \(c_K\) es el acarreo acumulado;
 \item \(\Sigma_K\) es la firma conjunta de frontera;
 \item \(W_K\) es la sombra de incidencia de Witt;
 \item \(g(K)=1+((K-1)\bmod9)\) es la etiqueta pública de fase;
 \item \(q_{{\rm mem},K}\in\mathbb Z_{\geq0}\) es la memoria vertical no
       acotada.
\end{enumerate}
La coordenada dodecafásica es derivada:
\[
 \beta_K=[q_{{\rm mem},K}]_{12}.
\]
No se sustituye \(q_{\rm mem}\) por \(\beta\), pues la reducción módulo
doce perdería el número de vueltas.

\subsection{Relaciones locales y holonomía}

Para cada \(K\), sea
\[
 \mathcal R_K\subseteq
 \widehat X_K\times\widehat X_{K+1}
\]
la relación de extensión compatible entre estados completos. La composición
relacional de una vuelta es
\begin{equation}
 \Gamma_{9,K}
 :=\mathcal R_{K+8}\circ\cdots\circ\mathcal R_{K+1}\circ\mathcal R_K
 :\widehat X_K\longrightarrow\widehat X_{K+9}.
\label{p12-83-c7-s1:eq:gamma9-composicion-local}
\end{equation}

\begin{definition}[Holonomía nonádica]
\label{p12-83-c7-s1:def:holonomia-nonadica}
La holonomía de una vuelta es
\[
 \Gamma_9:=\operatorname{Hol}_{C_9}(\nabla_{\rm TPK}),
\]
cuya realización en el nivel \(K\) es
\(\Gamma_{9,K}\). Las nueve relaciones \(\mathcal R_K,\ldots,
\mathcal R_{K+8}\) son cartas de esta única conexión.
\end{definition}

\begin{theorem}[Retorno de fase sin reinicio]
\label{p12-83-c7-s1:thm:retorno-nonadico-sin-reinicio}
Si \((\widehat x_K,\widehat x_{K+9})\in\Gamma_{9,K}\), entonces
\[
 g(K+9)=g(K),
 \qquad
 q_{{\rm mem},K+9}=q_{{\rm mem},K}+1,
\]
y
\[
 \beta_{K+9}=\beta_K+1\pmod{12}.
\]
En general, \(\widehat x_{K+9}\neq\widehat x_K\).
\end{theorem}

\begin{proof}
La primera identidad procede de la definición de \(g\). Una vuelta completa
incrementa por definición el contador de vueltas no acotado. Reducir esa
identidad módulo doce da la tercera. Las restantes coordenadas se
transportan por la composición \cref{p12-83-c7-s1:eq:gamma9-composicion-local}; ninguna
ley del TPK las reinicializa. Por tanto, la igualdad de fase no implica
igualdad del estado completo.
\end{proof}

\subsection{Clasificación exhaustiva de las \(9\) fases del Topological Phase Kernel}

Los bloques visibles y los censos locales de la primera vuelta son
\[
\begin{array}{c|rrrrr|ccc}
K&N_\pi&N_e&N_\varphi&N_{\rm loc}&N_{\rm ext}
 &u^\pi&u^e&u^\varphi\\\hline
1&378&422&349&3&2&012222&121221&112202\\
2&238&368&388&2&5&010211&102011&121020\\
3&184&284&173&5&4&002111&012222&010210\\
4&207&207&122&4&1&110221&102011&010200\\
5&39&151&151&1&1&222220&021222&102011\\
6&110&110&69&1&3&111201&201202&221021\\
7&81&29&81&3&3&212121&222210&200221\\
8&59&59&59&3&2&200121&212212&110110\\
9&43&43&43&2&1&100100&020112&010122
\end{array}.
\]

\begin{theorem}[Funciones de las nueve componentes]
\label{p12-83-c7-s1:thm:funciones-nueve-componentes}
En la vuelta anterior, las nueve componentes desempeñan las funciones
siguientes:
\begin{enumerate}
 \item apertura local de tres orientaciones y primera resolución;
 \item separación binaria con máximo local de cinco extensiones;
 \item máxima multiplicidad local de cinco estados;
 \item contracción de cuatro estados locales a una extensión;
 \item primera saturación fuerte, con una sola continuación;
 \item presente local único y apertura triple del futuro;
 \item fibra especular triple;
 \item sincronización de los tres censos en \(59\);
 \item cierre orientado: dos estados locales de entrada y una relación que,
       tras el horizonte adicional, conserva una prolongación.
\end{enumerate}
\end{theorem}

\begin{proof}
Cada afirmación se lee de los censos \(N_{\rm loc},N_{\rm ext}\), de la
igualdad o desigualdad de \(N_\pi,N_e,N_\varphi\) y de la relación de
extensión explícita en las fases novena, décima y undécima que se demostrará
en la unidad siguiente. La tabla enumera toda la vuelta; no se ha escogido
un subconjunto de fases.
\end{proof}

La quinta componente coincide con la primera frontera conjunta
\[
 R_{36}=(222220|021222|102011)^T.
\]
La novena componente contiene
\[
 B_9=(100100|020112|010122).
\]
La presencia de dos estados locales en la fase novena y de tres extensiones
desde el estado orientado pertenece a capas diferentes; no son censos
contradictorios.

La holonomía actúa sobre el estado holonómico integral completo. Antes de publicar
su cociente modal se construyen, por tanto, la involución que intercambia
las fibras de \(\pi\) y \(e\), el eje fijo de \(\varphi\), los invariantes
ternarios y la resolución local \(3\to1\) determinada por el segundo
horizonte.

\begin{HMTScientificOwnerSubordinated}
% Unidad U010. Espejo orientado y selección 3->1; redacción autónoma.

\chapter{Involución especular y selección local \(3\to1\) de la extensión superviviente}
\label{chap:espejo-seleccion-tres-uno}

La fase novena contiene dos estados locales de entrada. Una vez fijado el
estado orientado \(B_9\), la relación de frontera produce tres extensiones
hacia la fase primera. Las tres comparten un eje de autoescala y un agregado
ternario; difieren en el reparto ordenado de ese agregado entre las ramas de
clausura y propagación. Un segundo horizonte conserva exactamente una.

\section{Coordenadas simétrica y axial}

Sea
\[
 \mathcal B:=(\mathbb F_3^6)^3.
\]
Un bloque visible es
\[
 B=(u^\pi,u^e,u^\varphi)\in\mathcal B.
\]
Definimos
\[
 q(B):=u^\pi+u^e+u^\varphi,
 \qquad
 a(B):=u^\pi+u^e-u^\varphi.
\]
Todas las operaciones se realizan componente a componente en
\(\mathbb F_3^6\).

\begin{theorem}[Reconstrucción del eje y del agregado]
\label{thm:espejo-reconstruccion-eje-agregado}
Las coordenadas \(q,a\) determinan
\[
 \boxed{u^\varphi=2(q-a),}
 \qquad
 \boxed{u^\pi+u^e=q-u^\varphi.}
\]
\end{theorem}

\begin{proof}
Restar las definiciones da
\(q-a=2u^\varphi\). En \(\mathbb F_3\), \(2^{-1}=2\); por tanto
\(u^\varphi=2(q-a)\). Sustituir en la definición de \(q\) recupera el
agregado restante.
\end{proof}

La involución especular es
\[
 \mathfrak c(u^\pi,u^e,u^\varphi)
 =(u^e,u^\pi,u^\varphi).
\]

\begin{proposition}[Invariantes de la involución]
\label{prop:espejo-invariantes}
\(\mathfrak c^2=\operatorname{id}\), y
\[
 q(\mathfrak cB)=q(B),
 \qquad
 a(\mathfrak cB)=a(B).
\]
En particular, \(u^\varphi\) y \(u^\pi+u^e\) son invariantes.
\end{proposition}

\begin{proof}
Intercambiar dos veces las primeras coordenadas recupera \(B\). Tanto
\(q\) como \(a\) dependen de ellas sólo mediante su suma, de modo que el
intercambio las preserva. El \cref{thm:espejo-reconstruccion-eje-agregado}
transfiere esa invariancia al eje y al agregado.
\end{proof}

La involución no afirma que \(\pi=e\) ni que una sea el cociente de la otra.
Afirma que las dos ramas ocupan posiciones conjugadas respecto de una
estructura fija.

\section{Estado orientado en la fase \(9\) y datos de frontera asociados}

La proyección visible del estado seleccionado es
\[
 B_9=(100100\mid020112\mid010122).
\]
La relación inducida sobre bloques visibles se denota
\[
 \mathcal M_9^B\subseteq\mathcal B\times\mathcal B.
\]
Los dos estados locales de entrada de la tabla de fase preceden a esta
relación. Tras seleccionar el estado completo que proyecta a \(B_9\), su
fibra de salida tiene cardinal tres.

\begin{theorem}[Las tres extensiones de la fase novena]
\label{thm:espejo-tres-extensiones}
La imagen relacional de \(B_9\) contiene exactamente
\begin{align*}
B_{10}^{(0)}&=(101222\mid112221\mid112002),\\
B_{10}^{(1)}&=(102221\mid111222\mid112002),\\
B_{10}^{(2)}&=(111221\mid102222\mid112002).
\end{align*}
Las tres satisfacen
\[
 q_{10}=022112,
 \qquad
 a_{10}=101111,
\]
y, en consecuencia,
\[
 u_{10}^\varphi=112002,
 \qquad
 u_{10}^\pi+u_{10}^e=210110.
\]
\end{theorem}

\begin{proof}
La relación de extensión sobre estados completos se enumera en la fibra de
\(B_9\) y produce las tres filas indicadas. Sumar y restar sus tres
componentes da el mismo par \((q_{10},a_{10})\) en los tres casos. El
\cref{thm:espejo-reconstruccion-eje-agregado} produce el eje y el agregado
comunes. Las primeras dos componentes son distintas, por lo que las tres
filas representan distribuciones diferentes del mismo agregado.
\end{proof}

Sus publicaciones enteras en base tres son
\[
 (279,352,365),
 \qquad
 (280,351,365),
 \qquad
 (282,350,365).
\]

\section{Compatibilidad exacta de cilindros}

Una palabra ternaria de longitud \(L\) e índice \(N\) determina el cilindro
semiabierto
\[
 I_3(N,L)=\left[\frac{N}{3^L},\frac{N+1}{3^L}\right).
\]
Un prefijo de \(b\) bloques decimales, de entero asociado \(D\), determina
\[
 I_{1000}(D,b)=
 \left[\frac{D}{1000^b},\frac{D+1}{1000^b}\right).
\]

\begin{lemma}[Test entero de intersección]
\label{lem:espejo-test-entero-cilindros}
Los cilindros \(I_3(N,L)\) e \(I_{1000}(D,b)\) se intersectan si y sólo si
\[
 N1000^b<(D+1)3^L,
 \qquad
 (N+1)1000^b>D3^L.
\]
\end{lemma}

\begin{proof}
Dos intervalos semiabiertos \([a,b)\) y \([c,d)\) se intersectan si y sólo
si \(a<d\) y \(c<b\). Sustituimos los cuatro extremos y multiplicamos por
el entero positivo \(3^L1000^b\).
\end{proof}

El test utiliza sólo aritmética entera. No depende de una tolerancia de punto
flotante.

\section{2.º horizonte y unicidad}

El bloque siguiente es
\[
 B_{11}=(022212\mid110001\mid100021),
\]
y la frontera decimal correspondiente es
\[
 (502,662,638).
\]

\begin{theorem}[Resolución local \(3\to1\)]
\label{thm:espejo-seleccion-tres-uno}
Las tres extensiones del \cref{thm:espejo-tres-extensiones} son compatibles
en profundidad uno. Al exigir simultáneamente \(B_{11}\) y la frontera
\((502,662,638)\), sólo \(B_{10}^{(0)}\) se prolonga como el mismo estado
completo. Por tanto,
\[
 \boxed{
 \#\operatorname{Surv}^{\rm cal}_1(B_9)=3,
 \qquad
 \#\operatorname{Surv}^{\rm cal}_2(B_9)=1.}
\]
\end{theorem}

\begin{proof}
Aplicamos el \cref{lem:espejo-test-entero-cilindros} a cada uno de los tres
canales. En el primer horizonte, las seis desigualdades de cada candidato se
satisfacen. Tras anexar \(B_{11}\), las desigualdades de los tres canales
se mantienen para \(B_{10}^{(0)}\). En \(B_{10}^{(1)}\) y
\(B_{10}^{(2)}\) se mantiene la compatibilidad del canal \(\varphi\), pero
fallan las condiciones de las ramas \(\pi\) y \(e\). La enumeración deja,
por tanto, un único estado a profundidad dos.
\end{proof}

Esta prueba no define la conexión global a partir de un ejemplo local. Es un
testigo explícito de cómo una relación multivaluada en un horizonte puede
volverse univaluada al conservar memoria y exigir una prolongación mayor.

\section{Retorno de fase con cambio de bloque}

La fase primera de la vuelta inicial es
\[
 B_1=(012222\mid121221\mid112202),
\]
mientras la fase primera de la vuelta siguiente es
\[
 B_{10}=(101222\mid112221\mid112002).
\]
Así,
\[
 g(1)=g(10)=1,
 \qquad
 B_1\neq B_{10}.
\]
La fase retorna, pero el eje, el agregado, el reparto y la memoria de
cociente han evolucionado.

\section{Diagrama de la involución especular \(\pi/e\), el eje fijo \(\varphi\) y la selección local \(3\to1\) por prolongabilidad}

\begin{figure}[htbp]
\centering
\resizebox{0.98\textwidth}{!}{%
\begin{tikzpicture}[
 node distance=10mm and 15mm,
 every node/.style={font=\small,align=center},
 state/.style={draw=black,rounded corners,inner sep=5pt},
 inv/.style={draw=black,double,rounded corners,inner sep=5pt},
 arr/.style={-{Latex[length=2mm]},semithick}
]
 \node[state] (pi) {$u^\pi$\\rama de clausura};
 \node[state,right=70mm of pi] (e) {$u^e$\\rama de propagación};
 \coordinate (mirroraxis) at ($(pi)!0.5!(e)$);
 \node[inv,above=9mm of mirroraxis] (phi) {$u^\varphi$\\eje fijo};
 \draw[arr,<->] (pi) -- node[below] {$\mathfrak c$} (e);
 \node[above=8mm of phi] (aggregate)
   {$u^\pi+u^e$ conservado en la fibra};
 \draw[dashed] (phi.north) -- (aggregate.south);
 \node[state,below=20mm of mirroraxis] (b9) {fase $9$: $B_9$\\dos entradas locales};
 \node[state,below left=18mm and 25mm of b9] (b100) {$B_{10}^{(0)}$};
 \node[state,below=18mm of b9] (b101) {$B_{10}^{(1)}$};
 \node[state,below right=18mm and 25mm of b9] (b102) {$B_{10}^{(2)}$};
 \node[state,below=18mm of b101] (b11) {$B_{11}$\\única prolongación};
 \draw[arr] (b9) -- (b100);
 \draw[arr] (b9) -- (b101);
 \draw[arr] (b9) -- (b102);
 \draw[arr] (b100) -- (b11);
 \draw[arr,densely dotted] (b101) -- node[right] {obstruida} (b11);
\draw[arr,densely dotted] (b102) -- node[right] {obstruida} (b11);
\end{tikzpicture}
}
\caption{El espejo conserva el eje \(u^\varphi\) y el agregado
\(u^\pi+u^e\). El retorno de la fase novena abre tres extensiones visibles;
la memoria del segundo horizonte conserva una. La vuelta no reinicia el
estado.}
\label{fig:espejo-tres-uno-arrastre}
\end{figure}

\section{Memoria residuo--cociente durante el espejo}

La recurrencia
\[
 x_kA=u_k+3x_{k+1}
\]
de la \cref{eq:bockstein-hensel-local} se aplica a cada uno de los tres
canales. La parte \(u_k\) publica el bloque visible; \(x_{k+1}\) conserva
el cociente. Como \(A\in\operatorname{GL}_6(\mathbb Z)\), el par recupera
el estado anterior. El espejo puede intercambiar las ramas visibles sin
eliminar sus cocientes de procedencia.

\section{Obstrucciones de la involución especular \(\pi/e\) y de la selección local \(3\to1\): invariantes, horizonte de profundidad \(2\) y memoria Hensel}

\begin{criterion}[Criterios de refutación del espejo]
La construcción queda refutada si:
\begin{enumerate}
 \item la involución modifica \(u^\varphi\) o \(u^\pi+u^e\);
 \item se confunden los dos estados locales de entrada con las tres salidas;
 \item alguna de las tres salidas tiene distinto \((q_{10},a_{10})\);
 \item el test de cilindros utiliza una tolerancia flotante;
 \item más de una salida sobrevive al segundo horizonte declarado;
 \item el retorno de fase se interpreta como igualdad \(B_1=B_{10}\);
 \item el intercambio visible borra los cocientes de Hensel.
\end{enumerate}
\end{criterion}

El espejo y la selección local están ahora construidos. Las unidades
siguientes desarrollarán las vacancias sturmianas y, después, las tres
biografías analíticas completas.
 \end{HMTScientificOwnerSubordinated}

\subsection{Transferencia conjunta de la fibra de 468 emisiones}

La factorización del catálogo es
\[
 \mathcal U_6^{\rm cat}\cong
 \mathfrak S_+\times\mathfrak S_\times,
 \qquad
 |\mathfrak S_+|=18,\qquad |\mathfrak S_\times|=26.
\]
Para una función \(f\) sobre el catálogo, sean
\[
 (E_+f)(s,e)=\frac1{18}\sum_{s'\in\mathfrak S_+}f(s',e),
\]
y
\[
 (E_\times f)(s,e)=\frac1{26}
 \sum_{e'\in\mathfrak S_\times}f(s,e').
\]
Los proyectores \(E_+\) y \(E_\times\) conmutan.

Sea
\[
 \tau=(+1,-1,0,+1,-1,0,+1,-1,0),
 \qquad
 P_{+1}=E_+,quad P_{-1}=E_\times,quad P_0=I.
\]
En \(\mathcal U_6^{\rm cat}\times C_9\) definimos la transferencia de fase
por
\begin{equation}
 \mathcal K_{\rm ph}\bigl((u,r),(v,[r+1]_9)\bigr)
 :=P_{\tau(r)}(u,v),
 \label{p12-83-c7-s1:eq:u009-kph-explicita}
\end{equation}
y la hacemos cero para los demás pares de fases. Así, la permutación cíclica
y el operador aplicado en cada uno de los nueve pasos quedan determinados.

\begin{proposition}[Renovación conjunta tras una vuelta]
\label{p12-83-c7-s1:prop:transferencia-conjunta-nueve}
\[
 \mathcal K_{\rm ph}^{\,9}=E_+E_\times.
\]
En particular, sobre una masa puntual,
\[
 (E_+E_\times)(u,v)=\frac1{18\cdot26}=\frac1{468}.
\]
\end{proposition}

\begin{proof}
El producto ordenado sobre una vuelta es
\[
 P_{\tau(8)}P_{\tau(7)}\cdots P_{\tau(0)}
 =I E_\times E_+ I E_\times E_+ I E_\times E_+
 =E_+^3E_\times^3I^3=E_+E_\times.
\]
La segunda igualdad usa la conmutatividad y la última la idempotencia.
El segundo resultado es la distribución uniforme sobre el producto de las
dos firmas.
\end{proof}

Esta identidad expresa una renovación conjunta, no nueve filtraciones
independientes.

\subsection{Descomposición de la holonomía \(U_{\Gamma_9}\) en compresión visible \(T\) y componente ortogonal \(\eta\): \(I-T^*T=\eta^*\eta\)}

Sean \(\mathcal H_{\rm vis}\) y \(\mathcal H_{\rm full}\) espacios de
Hilbert para la realización visible y completa. Sea
\[
 J:\mathcal H_{\rm vis}\longrightarrow\mathcal H_{\rm full}
\]
una inclusión isométrica y
\(U_{\Gamma_9}\) una realización isométrica de la holonomía. Definimos
\[
 T:=J^*U_{\Gamma_9}J,
 \qquad
 \eta:=(I-JJ^*)U_{\Gamma_9}J.
\]

\begin{theorem}[Descomposición compresión--expresión]
\label{p12-83-c7-s1:thm:conexion-compresion-expresion}
Se cumplen
\[
 \boxed{U_{\Gamma_9}J=JT+\eta,}
\]
\[
 J^*\eta=0,
 \qquad
 \boxed{I-T^*T=\eta^*\eta.}
\]
\end{theorem}

\begin{proof}
Descomponemos \(U_{\Gamma_9}J\) mediante las proyecciones ortogonales
\(JJ^*\) e \(I-JJ^*\):
\[
 U_{\Gamma_9}J
 =JJ^*U_{\Gamma_9}J+(I-JJ^*)U_{\Gamma_9}J
 =JT+\eta.
\]
La ortogonalidad da \(J^*\eta=0\). Como \(J\) y
\(U_{\Gamma_9}\) son isometrías,
\[
 I=J^*U_{\Gamma_9}^*U_{\Gamma_9}J
 =T^*T+\eta^*\eta.
\]
Reordenar prueba la última identidad.
\end{proof}

\(T\) es la compresión visible; \(\eta\) es la expresión complementaria.
Ninguna de las dos debe llamarse contracción estricta. Para esa función se
reserva un operador distinto.

\subsection{Contracción estricta y telescopía terminal}

Sea \(V_9\) una isometría y sea \(q\) un parámetro con \(0<q<1\). Definimos
\[
 M_9=qV_9,
 \qquad
 D_9=(I-M_9^*M_9)^{1/2}.
\]
Entonces \(\|M_9\|=q<1\) y
\[
 I-M_9^*M_9=D_9^*D_9.
\]

\begin{theorem}[Identidad de Stein con término terminal]
\label{p12-83-c7-s1:thm:stein-telescopia-terminal}
Para todo \(L\geq1\),
\[
 \boxed{
 I=
 \sum_{r=0}^{L-1}M_9^{*r}D_9^*D_9M_9^r
 +M_9^{*L}M_9^L.}
\]
Más generalmente, si \(S_0\) satisface
\(S_0-M_9^*S_0M_9=D_9^*D_9\), entonces
\[
 \boxed{
 S_0=
 \sum_{r=0}^{L-1}M_9^{*r}D_9^*D_9M_9^r
 +M_9^{*L}S_0M_9^L.}
\]
\end{theorem}

\begin{proof}
Multiplicamos la identidad de defecto a izquierda por \(M_9^{*r}\) y a
derecha por \(M_9^r\). Al sumar para \(r=0,\ldots,L-1\), los términos
interiores telescopan y permanecen el inicial y el terminal. La misma prueba
se aplica a la ecuación general de Stein.
\end{proof}

El último término no se borra en profundidad finita. Su desaparición en un
límite requiere una prueba de convergencia; no es parte de la identidad
algebraica.

\subsection{Cociclo de acarreo y potencias relacionales}

Para flechas composables \(\gamma_1,\gamma_2\), el acarreo satisface
\[
 \operatorname{Carr}(\gamma_2\gamma_1)
 =\operatorname{Carr}(\gamma_2)H(\gamma_1)
 +Y(\gamma_2)\operatorname{Carr}(\gamma_1).
\]
La fórmula conserva la contribución de la primera etapa después de
transportarla por la segunda.

Las vueltas de mayor longitud se escriben con índices desplazados:
\begin{align*}
 \Gamma_{27,K}
 &=\Gamma_{9,K+18}\circ\Gamma_{9,K+9}\circ\Gamma_{9,K},\\
 \Gamma_{54,K}
 &=\Gamma_{9,K+45}\circ\cdots\circ\Gamma_{9,K},\\
 \Gamma_{108,K}
 &=\Gamma_{9,K+99}\circ\cdots\circ\Gamma_{9,K}.
\end{align*}
Sólo después de declarar un operador global sobre la unión graduada puede
abreviarse esta notación como \(\Gamma_9^3,\Gamma_9^6,\Gamma_9^{12}\).

\subsection{Estado nonádico completo: fase, acarreo, frontera, ruta y memoria tras la prolongación}

La conexión de una vuelta se registra por la quíntupla
\[
 \mathfrak N_9=
 \bigl(\Gamma_9,T,\eta,\operatorname{Carr},S_{\rm term}\bigr),
\]
donde \(S_{\rm term}=M_9^{*L}S_0M_9^L\) en una profundidad finita
\(L\). Un funtor de aplicación que utilice la conexión debe declarar su
compatibilidad con los cinco campos. Conmutar sólo con la fase visible no
constituye naturalidad respecto del estado nonádico completo.

\subsection{Obstrucciones de la holonomía \(\Gamma_9\): retorno \(9\to1\), índices de composición y conservación de fase, acarreo, frontera, ruta y memoria}

\begin{criterion}[Criterios de refutación de la conexión]
La construcción queda refutada si:
\begin{enumerate}
 \item se presentan las nueve fases como filtros independientes;
 \item el retorno \(9\to1\) reinicia \(q_{\rm mem}\), prefijo o carry;
 \item se llama contracción estricta a \(T\);
 \item falla alguna identidad de la
       \cref{p12-83-c7-s1:thm:conexion-compresion-expresion};
 \item se elimina el término terminal a profundidad finita;
 \item \(\Gamma_{27,K}\) se compone sin desplazar los índices;
 \item un funtor de aplicación ignora alguno de los cinco campos de
       \(\mathfrak N_9\).
\end{enumerate}
\end{criterion}

Los desarrollos completos que siguen reconstruyen primero el emisor finito,
su factorización y sus microfibras; después despliegan la involución
especular, las tres extensiones de \(B_9\), la unicidad a dos horizontes y la
inversa Bockstein--Hensel que conserva la memoria durante el retorno de fase.

\begin{HMTScientificOwnerSubordinated}
\chapter[Imagen aritmética finita y factorización \(468=18\cdot26\)]{Imagen aritmética finita y factorización \(468=18\cdot26\)}
\label{chap:alfabeto-42}
\index{código decimal admisible}
\index{dígito de control}

La imagen en base \(1000\) de la aplicación finita que asigna los bloques decimales es un
subconjunto propio de los mil bloques decimales. La ley de evolución definida en el
\cref{chap:tpk-definicion} y la aritmética de \(\mathbb Z/9\mathbb Z\)
determinan una imagen de 42 ternas, organizada en seis clases de siete. La
misma separación de los observables aditivo y multiplicativo factoriza después las 468
salidas en 18 firmas aditivas y 26 firmas multiplicativas.

\section{Primitivas y recurrencia del emisor finito}

\begin{definition}[Datos de partida de \(T_{\min}\)]
\label{pf84:E02-15}
Sea \(D=\{N,E,S,O\}\), con vectores cardinales
\[
 v(N)=(-1,0),\quad v(E)=(0,1),\quad
 v(S)=(1,0),\quad v(O)=(0,-1)
 \qquad\text{en }(\mathbb Z/9\mathbb Z)^2.
\]
La firma temporal y el índice multiplicativo se fijan por
\[
 \operatorname{sgn}_9(t)=
 \begin{cases}
 +1,&t\bmod9\in\{1,4,7\},\\
 -1,&t\bmod9\in\{2,5,8\},\\
 0,&t\bmod9\in\{0,3,6\},
 \end{cases}
\]
y
\[
 \operatorname{ind}_2(1,2,4,8,7,5)=(0,1,2,3,4,5),
\]
con valor cero fuera de las unidades. Una entrada de \(T_{\min}\) es un par
de semillas
\[
 I_9:=\{0,1,\ldots,8\},\qquad
 s_+=(i_+,j_+,d_+),\qquad s_\times=(i_\times,j_\times,d_\times)
 \in I_9^2\times D.
\]
Estas tablas, los lectores \(L_\Sigma,L_\Pi\), los seis bloques de nueve
instantes y la inversión de ambos vectores cuando \(27\mid t\) son todos los
datos de partida. No se introduce \(\pi\), \(e\), \(\varphi\) ni ningún otro
valor objetivo.
\end{definition}

\begin{algorithm}[Recurrencia de 54 pasos]
\label{pf84:E02-16}
Al comenzar, los dos cursores ocupan las posiciones de sus semillas y llevan
los vectores \(v(d_+)\), \(v(d_\times)\). Para
\(t=1,\ldots,54\), en este orden:
\begin{enumerate}
  \item si \(\operatorname{sgn}_9(t)=+1\), se acumula
  \(L_\Sigma(i_+(t),j_+(t))\) y se desplaza el cursor aditivo;
  \item si \(\operatorname{sgn}_9(t)=-1\), se acumula
  \(\operatorname{ind}_2(L_\Pi(i_\times(t),j_\times(t)))\) y se desplaza el
  cursor multiplicativo;
  \item si \(\operatorname{sgn}_9(t)=0\), ambos cursores permanecen quietos;
  \item después de esa actualización, ambos vectores se invierten si
  \(27\mid t\).
\end{enumerate}
Los acumuladores se reinician al comienzo de cada bloque
\(9(r-1)+1\leq t\leq9r\). Si \(S_r,E_r\) son sus valores finales, se emite
\[
 U_r=100\bar S_r+10(\bar S_r+\bar E_r)_{10}
       +(3S_r+5E_r+1)_{10},
 \qquad
 \bar S_r=S_r\bmod10,\quad \bar E_r=E_r\bmod10,
\]
y \(w_r=(-U_r)\bmod3\). Así se obtiene
\[
 T_{\min}(s_+,s_\times)=(U_1,\ldots,U_6)
 \in\mathcal U_6^{\rm cat},
 \qquad
 \Psi(U_6)=(-U_6)\bmod3=:w_6\in\mathbb F_3^6.
\]
\end{algorithm}

\begin{criterion}[Control dual independiente del valor objetivo]
\label{pf84:E02-17}
La recurrencia admite dos formulaciones independientes: factorización de
canales y simulación conjunta. El recorrido exhaustivo de las
\(104\,976\) entradas produce el mismo sexteto por ambas vías para cada
semilla, sin consultar una constante objetivo. Una discrepancia en una sola
semilla refutaría la equivalencia.
\end{criterion}

\section{Dinámica bivariada y observables por bloques}

Trabajamos en el toro \(\mathbb Z_9^2\). Con la convención precedente, una
semilla completa es
\[
\omega=(p_+,p_\times,h_+,h_\times)
\in(\mathbb Z_9^2)^2\times\{N,E,S,O\}^2,
\]
de modo que
\[
|\Omega|=9^4\cdot4^2=104\,976.
\]
En cada ventana de nueve pasos, los tiempos \(1,4,7\) evalúan el observable aditivo,
los tiempos \(2,5,8\) evalúan el observable multiplicativo y los tiempos \(3,6,9\)
incrementan el contador neutral. Si \(S_m,E_m,T_m\) son los acumuladores
de la ventana \(m\), el observable decimal es
\[
d_1\equiv S_m,\qquad
d_2\equiv S_m+E_m,\qquad
d_3\equiv3S_m+5E_m+7T_m\pmod{10},
\]
y la terna decimal es \(U_m=100d_1+10d_2+d_3\).

\begin{lemma}[Contador neutral]
Para toda ventana, \(T_m=3\).
\end{lemma}
\begin{proof}
Cada bloque de nueve tiempos contiene exactamente los tres tiempos neutrales
\(3,6,9\). Por tanto el contador se incrementa tres veces.
\end{proof}

\section{Imagen del acumulador aditivo: \(\operatorname{Im}S_m=\{6,9,12,15,18,21,24\}\)}

Sea \(s=i+j\pmod9\). El observable aditivo visible es
\[
f(s)=\operatorname{dr}_9(s+2)=(2,3,4,5,6,7,8,9,1).
\]
Una traslación cardinal modifica \(s\) en \(+1\) o en \(-1\). En una ventana se
leen tres términos consecutivos; invertir la orientación sólo invierte su
orden. Las nueve sumas cíclicas son
\[
9,12,15,18,21,24,18,12,6.
\]

\begin{lemma}[Imagen del acumulador aditivo]
\[
S_m\in\{6,9,12,15,18,21,24\}.
\]
En particular, \(d_1=S_m\bmod10\) recorre exactamente
\(\{1,2,4,5,6,8,9\}\).
\end{lemma}
\nopagebreak[4]
\begin{proof}
La lista anterior enumera todas las posiciones iniciales de la trayectoria y las
dos orientaciones. Sus residuos decimales son siete y distintos.\qedhere
\end{proof}

\section{Imagen del acumulador multiplicativo: \(\operatorname{Im}E_m=\{0,1,3,5,7,9\}\)}

El grupo de unidades \((\mathbb Z/9\mathbb Z)^\times\) es cíclico de orden
seis. Fijamos la función de coordenada multiplicativa
\[
\ell_\times(1,2,3,4,5,6,7,8,9)=(0,1,0,2,5,0,4,3,0).
\]
Si la fila o columna recorrida posee un factor divisible por tres, cada
evaluación pertenece a \(\{3,6,9\}\) y aporta cero. Si el factor fijo es una
unidad, las sumas de tres términos consecutivos de
\(b\mapsto\ell_\times(\operatorname{dr}_9(ab))\) son:
\[
\begin{array}{c|c}
a&\text{suma posible}\\
\midrule
1&\{1,3,7,9\}\\
2&\{3,5,9\}\\
4&\{1,5,7\}\\
5&\{1,5,7\}\\
7&\{3,5,9\}\\
8&\{1,3,7,9\}.
\end{array}
\]

\begin{lemma}[Imagen del acumulador multiplicativo]
\[
E_m\in\{0,1,3,5,7,9\}.
\]
\end{lemma}
\begin{proof}
El caso no invertible aporta \(0\). La tabla finita de las seis unidades
contiene todas las posibilidades del caso invertible y su unión es
\(\{1,3,5,7,9\}\).
\end{proof}

\section{Congruencia decimal del alfabeto ternario: unicidad del 3.er dígito}

\begin{theorem}[Congruencia de control]
Toda terna \(\overline{d_1d_2d_3}\) emitida por el sistema satisface
\[
\boxed{d_3\equiv5d_2-2d_1+1\pmod{10}.}
\]
\end{theorem}
\begin{proof}
De \(d_2\equiv S_m+E_m\) se obtiene
\(E_m\equiv d_2-d_1\pmod{10}\). Como \(T_m=3\),
\[
d_3\equiv3d_1+5(d_2-d_1)+21
\equiv5d_2-2d_1+1\pmod{10}.\qedhere
\]
\end{proof}

\begin{theorem}[Imagen exacta de la aplicación decimal]
El sistema produce exactamente \(7\cdot6=42\) ternas decimales. Agrupadas por
\(E=d_2-d_1\pmod{10}\), son:
\[
\begin{array}{c|rrrrrrr}
E&\multicolumn{7}{c}{\text{ternas decimales}}\\
\midrule
0&114&227&443&556&669&885&998\\
1&129&232&458&561&674&890&903\\
3&149&252&478&581&694&810&923\\
5&169&272&498&501&614&830&943\\
7&189&292&418&521&634&850&963\\
9&109&212&438&541&654&870&983.
\end{array}
\]
\end{theorem}
\begin{proof}
Hay siete posibilidades para \(d_1\) y seis para \(E\). Para cada par,
\(d_2\equiv d_1+E\) y la congruencia de control fija un único \(d_3\),
de modo que existen a lo sumo 42 ternas. La tabla aplica estas dos fórmulas
a los 42 pares y no contiene repeticiones. La enumeración de las
104\,976 semillas encuentra todas las ternas de la tabla y ninguna fuera de
ella; la cota se alcanza.
\end{proof}

\section{Factorización de los \(2\) observables}

Sea \(\mathcal S=\mathbb Z_9^2\times\{N,E,S,O\}\), con
\(|\mathcal S|=324\). A cada semilla aditiva \(P\in\mathcal S\) se le
asocia la firma de seis ventanas
\[
s(P)=(S_1,\ldots,S_6)\pmod{10},
\]
y a cada semilla multiplicativa \(T\in\mathcal S\),
\[
e(T)=(E_1,\ldots,E_6)\pmod{10}.
\]
Las dos firmas se combinan coordenada a coordenada mediante
\[
G(s,e)_k=
100s_k+10(s_k+e_k)_{10}+(3s_k+5e_k+1)_{10}.
\]
La aplicación \(G\) es inyectiva: desde una terna se recupera
\(s_k=d_1\) y \(e_k=d_2-d_1\pmod{10}\), mientras el tercer dígito verifica
la pertenencia a la imagen.

\begin{lemma}[Clasificación exacta del canal aditivo]
\label{pf84:E02-18}
\[
 |\operatorname{im}s|=18,
\]
y cada firma aditiva tiene exactamente \(18\) preimágenes.
\end{lemma}
\begin{proof}
La firma sólo depende de
\((i+j\bmod9,\varepsilon)\), donde \(\varepsilon=+1\) para \(E,S\) y
\(-1\) para \(N,O\). Los \(9\cdot2\) pares producen firmas distintas; cada
uno posee nueve posiciones en su diagonal y dos direcciones, es decir,
\(18\) preimágenes. Un número distinto de firmas o una clase de tamaño
distinto de \(18\) refutaría la clasificación.
\end{proof}

\begin{lemma}[Clasificación exacta del canal multiplicativo]
\label{pf84:E02-19}
\[
 |\operatorname{im}e|=26.
\]
Hay \(24\) clases de tamaño \(8\), una de tamaño \(24\) y una de tamaño
\(108\).
\end{lemma}
\begin{proof}
La enumeración de las \(9^2\cdot4=324\) semillas mediante la tabla explícita
de \(\ell_\times\) da \(26\) sextetos
\((E_1,\ldots,E_6)\) y el histograma
\[
 24\cdot8+24+108=324.
\]
Una clase adicional, ausente o de tamaño diferente refutaría el censo.
\end{proof}

\begin{theorem}[Factorización exacta e inyectiva del catálogo]
\label{pf84:E02-20}
\[
|\operatorname{im}s|=18,\qquad
|\operatorname{im}e|=26,\qquad
|\operatorname{im}G|=468=18\cdot26.
\]
Cada firma aditiva tiene 18 preimágenes. Las firmas multiplicativas poseen
el histograma
\[
24\text{ clases de tamaño }8,\qquad
1\text{ de tamaño }24,\qquad
1\text{ de tamaño }108.
\]
\end{theorem}
\begin{proof}
Los dos lemas precedentes dan los cardinales de las firmas. En cada componente
emitida, la cifra de centenas es \(\bar S_r\), mientras la de decenas es
\(\bar S_r+\bar E_r\bmod10\). Por tanto,
\[
 \bar S_r=\left\lfloor\frac{U_r}{100}\right\rfloor,\qquad
 \bar E_r=
 \left(\left\lfloor\frac{U_r}{10}\right\rfloor-\bar S_r\right)\bmod10.
\]
La cifra de unidades verifica \(3S_r+5E_r+1\bmod10\); no introduce otra
identificación. Así, \(G\) es inyectiva y su imagen tiene
\(18\cdot26=468\) elementos. Las \(24\) clases multiplicativas de tamaño
\(8\), combinadas con las \(18\) firmas aditivas, producen \(432\) sextetos
de multiplicidad \(144\); las otras dos clases producen \(18\) sextetos de
multiplicidad \(432\) y \(18\) de multiplicidad \(1944\). En particular,
\[
 432\cdot144+18\cdot432+18\cdot1944=104\,976.
\]
Dos pares de firmas con el mismo \(U_6\), o un histograma que no recupere el
dominio completo, refutarían el teorema.
\end{proof}

La composición completa factoriza, por tanto, como
\[
\mathcal S^2\xrightarrow{s\times e}
\operatorname{im}s\times\operatorname{im}e
\xrightarrow{G}\mathcal U_6
\xrightarrow{\rho_3}\mathcal W_6\subset\mathbb F_3^6.
\]
La última proyección posee 243 valores. El primer descenso conserva la
separación suma--producto; el segundo conserva 468 estados decimales; el
tercero elimina la coordenada de fibra y retiene una palabra ternaria.

\begin{corollary}[Censo de la proyección observable]
\label{cor:censo-proyeccion-observable}
\label{pf84:E02-21}
La enumeración completa determina la cadena de cardinales
\[
 104\,976\longrightarrow468\longrightarrow243\subset729.
\]
Los \(468\) elementos intermedios son sextetos de enteros comprendidos entre
cero y \(999\); sólo su imagen final pertenece a \(\mathbb F_3^6\). Las
multiplicidades de las fibras del primer mapa satisfacen
\[
 432\cdot144+18\cdot432+18\cdot1944=104\,976.
\]
\end{corollary}

\begin{proof}
El teorema de factorización da \(18\cdot26=468\), y la enumeración
exhaustiva de las \(104\,976\) semillas produce exactamente las tres
multiplicidades indicadas. La reducción componente a componente recorre las
\(243\) palabras registradas en el catálogo ternario. Los catálogos de
firmas y rutas adjuntos reproducen ambas enumeraciones.
\end{proof}

\begin{proposition}[Microfibra quíntuple de la palabra \(010211\)]
\label{pf84:E02-22}
Para la reducción restringida
\[
 \Psi:\mathcal U_6^{\rm cat}\longrightarrow\mathbb F_3^6,
 \qquad \Psi(U_6)=(-U_6)\bmod3,
\]
la fibra de \(010211\) consta exactamente de los cinco sextetos siguientes:
\begingroup
\scriptsize
\[
\begin{array}{@{}cclc@{}}
\toprule
U_6&\sigma_+&\sigma_\times&\text{multiplicidad}\\
\midrule
(501,614,498,169,272,272)&(5,6,4,1,2,2)&(5,5,5,5,5,5)&432\\
(810,923,870,169,272,272)&(8,9,8,1,2,2)&(3,3,9,5,5,5)&144\\
(810,983,810,169,272,272)&(8,9,8,1,2,2)&(3,9,3,5,5,5)&144\\
(870,923,810,169,272,272)&(8,9,8,1,2,2)&(9,3,3,5,5,5)&144\\
(870,923,810,418,674,521)&(8,9,8,4,6,5)&(9,3,3,7,1,7)&144\\
\bottomrule
\end{array}
\]
\endgroup
Son cinco preimágenes del catálogo \(U_6\), no cinco historias completas
TPK. Conservan firmas y multiplicidades distintas aunque compartan la misma
palabra visible.
\end{proposition}

\begin{proof}
El barrido exhaustivo del catálogo de \(468\) sextetos selecciona exactamente
las cinco filas de la tabla. La función de recuperación de firmas reproduce
los dos sextetos de cada fila y la suma de multiplicidades es
\(432+4\cdot144=1008\). La enumeración comprueba además que cada celda es el
producto cartesiano de sus semillas de canal. Un cardinal distinto de cinco
o la identificación de dos filas pese a sus datos enriquecidos refutaría el
enunciado.
\end{proof}
 \end{HMTScientificOwnerSubordinated}

\begin{HMTScientificOwnerSubordinated}
\chapter{Sistemas inversos, representación de Weyl--Heisenberg y cociente predictivo}
\label{chap:numero-heisenberg-d108}
\index{número HMT}
\index{Heisenberg!grupo finito}

La construcción distingue tres niveles: el catálogo finito de semillas, el
espacio de estados que conserva el acarreo y la orientación, y el límite
inverso de las prolongaciones compatibles. Esta estratificación permite
formalizar la contracción de los intervalos cilíndricos hasta un punto real,
la persistencia de datos de trayectoria en una misma fibra de evaluación y
la representación finita de Weyl--Heisenberg asociada a la forma alternante
del toro nonádico que soporta los observables aditivo y multiplicativo.

Los símbolos \(V_t\), \(X_t\) y \(\rho_{t+1,t}\) empleados aquí son los
mismos conjuntos de estados locales, prefijos admisibles y truncamientos
introducidos en el \cref{chap:numero-medicion}, con el cambio puramente
notacional \(t=m\) y \(\rho_{n,m}=p_{n,m}\). Esta sección explicita sus
coordenadas y no introduce un segundo sistema inverso.

\section{Proyección observable y espacio completo de estados del TPK}

La proyección observable del sistema de transporte de fase (TPK) es la
aplicación finita de codificación
\[
104\,976\longrightarrow468\longrightarrow243.
\]
que clasifica semillas, firmas y palabras \(w_6\). El espacio completo de
estados del TPK conserva además
\[
\begin{gathered}
\text{orientación bidireccional},\quad
\text{estado Hensel},\quad
\text{cilindro y acarreo},\\
\text{firma conjunta de frontera},\quad
\text{fase nonádica visible }g_{\rm vis}(m).
\end{gathered}
\]
Dos candidatos con una misma firma local pueden pertenecer a clases distintas
cuando se conserva la frontera siguiente. No introducimos una segunda tupla
de estado. Para
\[
x_m=(q_m,o_m,h_m,c_m,\sigma_m,C_m,b_m,\lambda_m)\in\XTPK^{(m)}
\]
usamos la carta de coordenadas
\[
\chi_m(x_m)=
\bigl(
B_m,h_m,(C_m,b_m),\sigma_m,
o_m,\mu(q_m),g_{\rm vis}(m),\lambda_m
\bigr),
\]
con
\[
B_m=(u_m^\pi,u_m^e,u_m^\varphi)\in(\mathbb F_3^6)^3.
\]
Aquí \(h_m\) es la extensión de Hensel; \((C_m,b_m)\), el cilindro y su
frontera; \(o_m,\mu(q_m)\), la orientación y el observable bidireccional; y
\[
g_{\rm vis}(m)=1+((m-1)\bmod9)\in\{1,\ldots,9\},
\]
la etiqueta visible de fase nonádica. La coordenada de clase correspondiente
es \(g_0(\theta_m)\in\ZZ/9\ZZ\), con \(\theta_m=[m-1]_{108}\). La firma conjunta
es la coordenada ya tipada por el núcleo de fase,
\[
\sigma_m=\operatorname{Sig}_{q_m}(h_m,c_m,\lambda_m)
=(q_\partial,a_\partial,c_\partial,
\operatorname{colw},\rho_W,r_\partial).
\]
La firma conjunta no se reduce a tres firmas de canal independientes, pues la
selección asociada a la involución \(\pi/e\) y al eje \(\varphi\) depende de
datos correlacionados. Dos nodos con el mismo \(B_m\) representan estados
distintos cuando difieren en extensión, cilindro, acarreo, firma, orientación
o registro de transporte.

\subsection{Compatibilidad exacta de los cilindros bajo truncación y prolongación nonádica}

Sea \(N\) el entero de un prefijo ternario de longitud \(L\), y sea \(D\)
el entero obtenido al concatenar sus \(b\) tríadas decimales. Definimos
\[
A=N1000^b-D3^L,
\qquad
B=(N+1)1000^b-D3^L.
\]
Los cilindros se intersecan si y sólo si
\[
\boxed{A<3^L,\qquad B>0.}
\]
Al añadir un bloque ternario \(u\in\{0,\ldots,728\}\), se tiene
\[
N'=729N+u.
\]
Si no aparece una nueva tríada decimal,
\[
A'=729A+u1000^b.
\]
Si aparece \(\delta\), de modo que \(D'=1000D+\delta\), entonces
\[
A'=729000A+u1000^{b+1}-\delta\,729\,3^L.
\]
Si \(\Delta b\in\{0,1\}\) registra la aparición de esa nueva tríada, en
ambos casos
\[
B'=A'+1000^{b+\Delta b}.
\]
La transición, fijado el par admisible \((u,\delta)\), no consulta un valor
objetivo. La filtración conjunta por intersección de cilindros y retorno
nonádico actúa en un nivel distinto: selecciona los pares de frontera que
admiten prolongación compatible.

\subsection{Registro cíclico de 108 etapas y su proyección reducida}

Denominamos \emph{registro cíclico de 108 etapas} a la sucesión de
\(108\) pasos agrupados en doce sectores de nueve pasos, con la convención de
índice fijada en el \cref{chap:tpk-definicion}.

Sea \(\XTPK^{[108]}\subseteq\XTPK\) el dominio de estados completos cuyo
registro \(\lambda\) contiene los \(108\) pasos declarados. Denotemos por
\(\mathcal C_{108}^{\rm red}\) el conjunto de registros que resultan al
conservar sólo los campos reducidos, y por \(\mathcal U_{12}^{\rm sgn}\) el
conjunto de dodecafases firmadas admisibles. La instancia distingue dos mapas
declarados:
\[
F:\XTPK^{[108]}\longrightarrow\mathcal C_{108}^{\rm red},
\qquad
R:\XTPK^{[108]}\longrightarrow\mathcal U_{12}^{\rm sgn}.
\]
El mapa \(F\) retiene sólo los campos del registro reducido, mientras que
\(R\) evalúa la dodecafase firmada. La existencia de una aplicación \(\Xi\)
tal que \(R=\Xi\circ F\) equivale a la factorización de \(R\) a través de una
proyección que puede eliminar la coordenada del observable, la orientación,
el acarreo o la frontera. Esta factorización constituye una proposición
adicional y no forma parte implícita de la reconstrucción sobre el espacio
completo de estados.

El registro reducido es, por definición, una proyección del estado completo,
no el estado de partida. Los resultados posteriores usan \(R\) sobre el
estado completo y no presuponen una factorización a través de \(F\).

\section{Sistemas inversos de supervivientes}

La transición TPK no tiene por qué ser una función. Para tipar correctamente
este hecho, sea \(V_t\) el conjunto finito de estados completos posibles en
la profundidad \(t\), y sea
\[
\mathcal R_t\subseteq V_t\times V_{t+1}
\]
la relación de transición admisible, con todos los campos del estado
holonómico integral conservados. Definimos entonces \(X_t\), no como el conjunto de
estados terminales aislados, sino como el conjunto de prefijos admisibles, es
decir, de historias admisibles completas:
\[
X_t=
\Bigl\{
(v_0,\ldots,v_t)\in\prod_{j=0}^{t}V_j:
(v_j,v_{j+1})\in\mathcal R_j\ (0\le j<t),
\ \text{y se satisfacen los cierres hasta }t
\Bigr\}.
\]
Sobre estos objetos la truncación sí es una aplicación canónica:
\[
\begin{aligned}
\rho_{t+1,t}:X_{t+1}&\longrightarrow X_t,\\
(v_0,\ldots,v_t,v_{t+1})&\longmapsto(v_0,\ldots,v_t).
\end{aligned}
\]
Está bien definida precisamente porque todo prefijo de una historia
admisible sigue siendo una historia admisible para los cierres ya impuestos.
Si un filtro utilizase información futura y pudiera invalidar
retroactivamente un prefijo, primero habría que reemplazar \(X_t\) por el
subconjunto de prefijos prolongables; con esa convención, la misma truncación
vuelve a estar bien definida.

Una historia HMT global es una familia de prefijos compatible
\[
x=(x_t)_{t\ge0},
\qquad
x_t\in X_t,
\qquad
\rho_{t+1,t}(x_{t+1})=x_t.
\]
Equivale a una sucesión infinita \((v_0,v_1,\ldots)\) cuyas parejas
consecutivas pertenecen a las relaciones \(\mathcal R_t\) y cuyos prefijos
satisfacen todos los cierres. El espacio de historias es el límite inverso
\[
X_\infty=\varprojlim X_t.
\]
Esta formulación evita imponer una transición local unívoca
\(B_t\mapsto B_{t+1}\). La transición puede ser una relación finita, varias
ramas pueden compartir el mismo estado terminal visible y la unicidad puede
aparecer sólo después de imponer compatibilidad futura. Al conservar el
prefijo completo no se pierde, antes de tiempo, la información que permite
distinguirlas.

\begin{theorem}[Sección coinductiva única]
Supongamos que, para todo \(t\), el conjunto de supervivientes \(S_t\subseteq
X_t\) es no vacío, las restricciones
\(\rho_{t+1,t}:S_{t+1}\to S_t\) son compatibles y cada \(S_t\) contiene un
único prefijo. Entonces \(\varprojlim S_t\) contiene una única historia
global.
\end{theorem}

\begin{proof}
Sea \(x_t\) el único prefijo de \(S_t\). La compatibilidad obliga a
\(\rho_{t+1,t}(x_{t+1})=x_t\), de modo que \((x_t)_t\) pertenece al límite
inverso. Cualquier otra historia tendría en cada nivel un miembro de \(S_t\)
y, por unicidad, coincidiría prefijo a prefijo con \((x_t)_t\).
\end{proof}

El teorema formaliza el paso de la compatibilidad nivel a nivel a una historia
global. No convierte una ventana finita en una conclusión coinductiva: su
aplicación a una familia concreta exige verificar en cada profundidad la no
vacuidad, la compatibilidad de las restricciones y la unicidad de los prefijos
supervivientes.

\section{Conexión nonádica de supervivencia}
\label{sec:conexion-nonadica}

La contracción de cilindros no procede bloque a bloque como una sustitución
determinista. En determinadas fronteras existen varios repartos locales de un
mismo dato agregado; la prolongabilidad decide cuál pertenece a la historia
publicada. La estructura resultante es una conexión sobre las nueve clases de
fase: al completar una vuelta regresa la clase de fase, pero el estado de la
fibra ha cambiado. Esta sección construye la ventana finita completa en la que
se observa ese retorno.

El dominio del cálculo es el registro de transporte calibrado adjunto. Sus estados iniciales,
su matriz entera de propagación y sus relaciones de frontera son premisas
publicadas, no salidas de una enumeración autónoma de TPK a profundidad
arbitraria. Todas las afirmaciones de esta sección son exactas relativamente a
esas premisas. La extensión a toda profundidad se rige por el teorema
coinductivo anterior y requiere verificar sus hipótesis en cada nivel.

\subsection{Prefijo de \(30\) trits y 1.ª frontera}

Escribimos \(u_k^\chi\in\FF_3^6\), con
\(\chi\in\{\pi,e,\varphi\}\), para el bloque visible del canal \(\chi\) a
profundidad \(k\). Los cinco primeros bloques forman tres prefijos de treinta
trits:
\begin{align*}
w_{30}^{\pi}&=010211|012222|010211|002111|110221,\\
w_{30}^{e}&=201101|121221|102011|012222|102011,\\
w_{30}^{\varphi}&=121200|112202|121020|010210|010200.
\end{align*}
El sexto bloque es la primera frontera posterior a esos prefijos:
\[
R_{36}=
\begin{pmatrix}
u_5^\pi\\ u_5^e\\ u_5^\varphi
\end{pmatrix}
=
\begin{pmatrix}
222220\\021222\\102011
\end{pmatrix}.
\]
La notación \(R_{36}\) indica que se han expuesto treinta y seis trits por
canal. No designa un estado terminal: el registro que se reconstruye más abajo
contiene veinte bloques, y la recurrencia admite continuaciones ulteriores
cuando se suministran los estados de Hensel correspondientes.

\subsection{Eje de autoescala y fibra residual}

\ifdefined\HMTAxialTheoremDefined
La descomposición axial de la firma y la identificación del eje
\(\varphi\) frente a la fibra residual \(\pi/e\) ya se demostraron en el
\cref{thm:eje-espejo-nonadico}. Aquí se conserva su consecuencia para la
ventana nonádica sin duplicar el teorema ni la prueba.
\else
Para
\[
B_k=(u_k^\pi,u_k^e,u_k^\varphi)\in(\FF_3^6)^3
\]
definimos, coordenada a coordenada,
\[
q_k=u_k^\pi+u_k^e+u_k^\varphi,
\qquad
a_k=u_k^\pi+u_k^e-u_k^\varphi.
\]

\begin{theorem}[Descomposición axial de la firma de frontera]
\label{thm:eje-espejo-nonadico}
El par \((q_k,a_k)\) determina unívocamente el canal axial y el agregado de
los otros dos canales:
\[
u_k^\varphi=2(q_k-a_k),
\qquad
u_k^\pi+u_k^e=q_k-u_k^\varphi.
\]
En consecuencia, una fibra de la aplicación
\[
q_{\rm ax}(u^\pi,u^e,u^\varphi)
=(u^\pi+u^e,u^\varphi)
\]
sólo puede conservar ambigüedad en el reparto interno entre \(u^\pi\) y
\(u^e\).
\end{theorem}

\begin{proof}
La resta de las dos ecuaciones definitorias da
\(q_k-a_k=2u_k^\varphi\). El inverso de \(2\) en \(\FF_3\) vuelve a ser
\(2\), de donde
\(u_k^\varphi=2(q_k-a_k)\). Sustraer este vector de \(q_k\) produce
\(u_k^\pi+u_k^e\). La operación no recupera cada sumando por separado; por
ello la única fibra residual es el conjunto de repartos del agregado fijo.
\end{proof}

En la terminología geométrica del modelo, \(\varphi\) constituye el eje
fijado por la firma y el par \(\pi/e\) la fibra residual. La expresión no
postula una simetría analítica de las constantes reales: nombra la
descomposición lineal exacta que acaba de demostrarse en \(\FF_3^6\).
\fi

\subsection{Órbita completa de las \(9\) fases consecutivas del Topological Phase Kernel}

Fijamos la fase absoluta mediante
\[
g(k)=
\begin{cases}
k\bmod9,&k\not\equiv0\pmod9,\\
9,&k\equiv0\pmod9.
\end{cases}
\]
Como \(w_{30}\) ocupa \(k=0,\ldots,4\), la primera vuelta completa posterior
al prefijo recorre \(k=5,\ldots,13\), con fases
\(5,6,7,8,9,1,2,3,4\). La enumeración de los \(729=3^6\) candidatos de
firma agregada en cada nivel produce la siguiente ventana:
\begin{center}
\small
\begin{tabular}{@{}ccrrlll@{}}
\toprule
\(k\)&\(g(k)\)&locales&seleccionados&\(u_k^\pi\)&\(u_k^e\)&\(u_k^\varphi\)\\
\midrule
5&5&1&1&222220&021222&102011\\
6&6&1&1&111201&201202&221021\\
7&7&3&1&212121&222210&200221\\
8&8&3&1&200121&212212&110110\\
9&9&2&1&100100&020112&010122\\
10&1&1&1&101222&112221&112002\\
11&2&1&1&022212&110001&100021\\
12&3&1&1&012012&202222&000120\\
13&4&1&1&111210&112102&211122\\
\bottomrule
\end{tabular}
\end{center}
La columna «locales» cuenta los estados de entrada admisibles en el nivel
indicado; no cuenta las salidas de la relación de frontera. Esta distinción
es esencial en \(k=9\): los dos estados locales de entrada preceden a una
fibra de tres salidas del estado seleccionado.

\subsection{Retorno 9→1 de la fase nonádica con cambio de estado y avance de memoria}

Sea
\[
 \pi_B:V_k\longrightarrow\mathcal B:=(\FF_3^6)^3
\]
la proyección al triple visible. La relación de transporte sobre estados
completos induce, una vez fijados los datos de fibra del registro, una relación
visible \(\MonNine^B\). El estado seleccionado que entra en la fase novena es
\[
B_9=(100100\mid020112\mid010122).
\]
Su imagen relacional consta exactamente de
\[
\begin{aligned}
B_{10}^{(0)}&=(101222\mid112221\mid112002),\\
B_{10}^{(1)}&=(102221\mid111222\mid112002),\\
B_{10}^{(2)}&=(111221\mid102222\mid112002).
\end{aligned}
\]
Las tres salidas tienen
\[
q_{10}=022112,
\qquad a_{10}=101111,
\]
y, por el \cref{thm:eje-espejo-nonadico}, comparten
\[
u_{10}^\varphi=112002,
\qquad u_{10}^\pi+u_{10}^e=210110.
\]
Difieren exclusivamente en las tres distribuciones publicadas de ese agregado
entre \(\pi\) y \(e\).

\begin{theorem}[Selección por dos niveles de prolongabilidad]
\label{thm:seleccion-novena-fase}
En la ventana calibrada, las tres salidas anteriores admiten una prolongación
de profundidad uno. Al imponer también el bloque siguiente
\[
B_{11}=(022212\mid110001\mid100021)
\]
y la frontera decimal \((502,662,638)\), sólo
\(B_{10}^{(0)}\) admite una prolongación del mismo estado completo. Por tanto
\[
\#\operatorname{Surv}^{\rm cal}_{1}(B_9)=3,
\qquad
\#\operatorname{Surv}^{\rm cal}_{2}(B_9)=1.
\]
\end{theorem}

\begin{proof}
Las ternas decimales asociadas a las tres salidas son, respectivamente,
\[
(279,352,365),\quad(280,351,365),\quad(282,350,365).
\]
Para una palabra ternaria de longitud \(L\), de entero asociado \(N\), y un
prefijo de \(b\) ternas, de entero asociado \(D\), la compatibilidad de los
dos cilindros equivale a las dos desigualdades enteras
\[
N1000^b<(D+1)3^L,
\qquad
(N+1)1000^b>D3^L.
\]
Las seis desigualdades de cada salida se satisfacen en el nivel actual. Al
anexar \(B_{11}\) y \((502,662,638)\), vuelven a satisfacerse para los tres
canales de \(B_{10}^{(0)}\). Para \(B_{10}^{(1)}\) y \(B_{10}^{(2)}\), se
mantiene la desigualdad del canal \(\varphi\), pero fallan las de \(\pi\) y
\(e\). La evaluación se efectúa con enteros en el certificado finito correspondiente; no
interviene tolerancia numérica. En consecuencia, el conjunto de candidatos
se reduce de tres a uno al pasar de profundidad uno a profundidad dos.
\end{proof}

Las otras dos palabras pueden preceder a una misma cola ternaria sólo si se
cambia la extensión profunda de Hensel. No son, por ello, prolongaciones del
mismo estado completo del registro: el residuo visible coincide en parte, pero
la genealogía de la fibra es distinta.

\subsection{Retorno de fase con cambio de estado}

La fase satisface \(g(k+9)=g(k)\). El estado visible, en cambio, no se
reinicializa. En particular,
\[
B_1=(012222\mid121221\mid112202),
\qquad
B_{10}=(101222\mid112221\mid112002),
\]
de modo que \(g(1)=g(10)=1\) pero \(B_1\ne B_{10}\). La comprobación de los
nueve pares \(k\mapsto k+9\), para \(k=1,\ldots,9\), da además
\[
u_{k+9}^\varphi\ne u_k^\varphi,
\qquad
u_{k+9}^\pi+u_{k+9}^e\ne u_k^\pi+u_k^e.
\]
Por consiguiente, la ventana satisface exactamente
\[
\boxed{g(k+9)=g(k),\qquad B_{k+9}\ne B_k
\quad(1\le k\le9).}
\]
Este es un retorno de fase con holonomía visible no trivial. La denominación
\emph{monodromía nonádica} se refiere a esta acción de retorno sobre la
fibra enriquecida; no introduce una décima fase. Para promover la observación
finita a una monodromía global se debe construir la acción a toda profundidad
y demostrar la compatibilidad coinductiva exigida al comienzo de la sección.

\subsection{La rama de \(20\) bloques}

El retorno anterior no clausura la emisión. Los veinte bloques que el
certificado reconstruye desde los tres estados enteros publicados son
\begin{align*}
\pi:\;&010211|012222|010211|002111|110221|222220|111201|\\
&212121|200121|100100|101222|022212|012012|111210|\\
&121011|200220|120210|000101|022010|020111,\\[2pt]
e:\;&201101|121221|102011|012222|102011|021222|201202|\\
&222210|212212|020112|112221|110001|202222|112102|\\
&102010|022020|222002|012010|001112|022212,\\[2pt]
\varphi:\;&121200|112202|121020|010210|010200|102011|221021|\\
&200221|110110|010122|112002|100021|000120|211122|\\
&202200|202110|221000|100102|111122|020010.
\end{align*}
Así, \(w_{30}\) es el prefijo y \(R_{36}\) la primera frontera de una rama
explícitamente continuada; ninguno de los dos objetos sustituye al estado
holonómico integral que transporta esa continuación.

\subsection{Decodificación exacta de \(12\) ternas decimales mediante el lector nonádico}

Sea \(N_{\chi,k}\) el entero representado en base tres por los primeros \(k\)
bloques del canal \(\chi\). Con \(k=13\), la palabra tiene \(78\) trits y
define
\[
D_{\chi,12}
=\left\lfloor\frac{1000^{12}N_{\chi,13}}{3^{78}}\right\rfloor.
\]
La escritura de \(D_{\chi,12}\) en doce bloques de tres cifras es
\[
\begin{array}{c|rrrrrrrrrrrr}
\pi&141&592&653&589&793&238&462&643&383&279&502&884\\
e&718&281&828&459&045&235&360&287&471&352&662&497\\
\varphi&618&033&988&749&894&848&204&586&834&365&638&117.
\end{array}
\]
La identidad es entera: equivale a
\[
D_{\chi,12}3^{78}
\le 1000^{12}N_{\chi,13}
<(D_{\chi,12}+1)3^{78}.
\]
Por tanto, los dos cilindros se intersecan en cada canal y la operación de
lectura produce exactamente las treinta y seis cifras publicadas. Este
resultado es una consecuencia del registro calibrado. Los estados iniciales de
la carta adjunta fueron construidos históricamente con esas constantes como
datos de referencia; por ello, esta identidad no establece la procedencia
independiente de los estados iniciales a partir de las primitivas de APP\@.

\subsection{Calendario de reapertura de cilindros y periodicidad de la prolongación compatible}

La conversión de un bloque de seis trits a ternas decimales tiene pendiente
media
\[
\lambda=6\log_{1000}3\in(0,1),
\qquad k(t)=\lfloor\lambda(t+1)\rfloor.
\]
Un índice \(t\) es una reapertura sin nueva terna cuando
\(k(t+1)=k(t)\). Sea \(\delta=1-\lambda\).

\begin{theorem}[Calendario de Beatty de la lectura ternario--decimal]
\label{thm:calendario-beatty}
El conjunto de reaperturas es
\[
\mathcal T=
\left\{\left\lceil\frac{m}{\delta}\right\rceil-2:m\ge1\right\}.
\]
Sus primeros elementos son
\[
20,42,64,86,108,130,151,173,195,217,\ldots,
\]
y dos reaperturas consecutivas distan \(21\) o \(22\) bloques.
\end{theorem}

\begin{proof}
La cantidad \(\lambda=\log_{1000}729\) es irracional: una igualdad
\(\lambda=p/q\) implicaría \(729^q=1000^p\), imposible por factorización
prima. También \(\delta\) es irracional, por lo que
\[
k(t)=(t+1)-\lceil\delta(t+1)\rceil.
\]
La igualdad \(k(t+1)=k(t)\) ocurre exactamente cuando la función techo del
segundo término aumenta en una unidad. La \(m\)-ésima ocasión se obtiene
resolviendo
\(\delta(t+1)<m<\delta(t+2)\), lo cual da
\(t=\lceil m/\delta\rceil-2\). Las diferencias de una sucesión de Beatty
pertenecen a
\[
\left\{\left\lfloor\delta^{-1}\right\rfloor,
\left\lceil\delta^{-1}\right\rceil\right\}=\{21,22\}.\qedhere
\]
\end{proof}

\subsection{Recurrencia Bockstein--Hensel de residuo y cociente}

La memoria transportada por la rama procede de la sucesión exacta
\[
0\longrightarrow\ZZ^6\xrightarrow{\,\times3\,}\ZZ^6
\longrightarrow\FF_3^6\longrightarrow0
\]
y de la matriz unimodular
\[
A=
\begin{pmatrix}
2&2&2&1&2&1\\
2&1&2&2&1&1\\
1&0&0&1&0&2\\
0&0&1&1&1&0\\
2&2&2&0&2&1\\
2&1&2&1&0&0
\end{pmatrix},
\qquad\det A=1.
\]
Para un estado fila \(x_k^\chi\in\ZZ^6\), se define
\[
y_k^\chi=x_k^\chi A,
\qquad
u_k^\chi=y_k^\chi\bmod3\in\{0,1,2\}^6,
\qquad
x_{k+1}^\chi=\frac{y_k^\chi-u_k^\chi}{3}.
\]

\begin{proposition}[Conservación exacta de la memoria de cociente]
\label{prop:memoria-bockstein-hensel}
Fijados \(x_k^\chi\) y \(A\), el par
\((u_k^\chi,x_{k+1}^\chi)\) es único y satisface
\[
x_k^\chi A=u_k^\chi+3x_{k+1}^\chi.
\]
Recíprocamente, como \(A\in\operatorname{GL}_6(\ZZ)\), el par determina
unívocamente \(x_k^\chi\).
\end{proposition}

\begin{proof}
La división euclídea de cada coordenada de \(x_k^\chi A\) por tres produce
un residuo único en \(\{0,1,2\}\) y un cociente entero único. Esto prueba
la primera afirmación y la identidad. Puesto que \(\det A=1\), la inversa
\(A^{-1}\) tiene entradas enteras; por tanto
\[
x_k^\chi=(u_k^\chi+3x_{k+1}^\chi)A^{-1}
\]
recupera el estado anterior sin ambigüedad.
\end{proof}

El residuo \(u_k^\chi\) es la parte visible y el cociente
\(x_{k+1}^\chi\) la memoria 3-ádica transportada. El certificado reproduce
sesenta pasos coordenados ---veinte por canal---, verifica \(AA^{-1}=I\) con
aritmética entera y coteja todos los bloques con el registro. La denominación
Bockstein--Hensel alude aquí a esta separación residuo--cociente y a su
iteración 3-ádica; no se identifica sin datos adicionales con un operador de
Bockstein cohomológico.

\section{Del prefijo a un punto real}

Asociemos a cada prefijo compatible \(x_t\in X_t\) un intervalo cerrado
\(I_t(x_t)\subset\mathbb R\). La interpretación geométrica es la de un
cilindro: todos los futuros que comparten el mismo prefijo se proyectan dentro
de \(I_t\).

\begin{theorem}[Convergencia arquimediana de intervalos encajados]
Sea \((x_t)_{t\ge0}\) una historia compatible. Si
\[
I_{t+1}(x_{t+1})\subseteq I_t(x_t)
\]
y
\[
\operatorname{diam} I_t(x_t)\longrightarrow0,
\]
entonces existe un único \(r\in\mathbb R\) tal que
\[
\bigcap_{t\ge0}I_t(x_t)=\{r\}.
\]
\end{theorem}

\begin{proof}
Es el principio de intervalos encajados. Los extremos izquierdos forman una
sucesión creciente acotada y los derechos una decreciente acotada. Sus límites
existen; la diferencia entre ellos está acotada por el diámetro, que tiende a
cero. Ambos límites coinciden en un único punto contenido en todos los
intervalos.
\end{proof}

Sea \(X_\infty^{\rm Arch}\subseteq X_\infty\) el dominio de historias para
las que los cilindros están encajados y sus diámetros tienden a cero. Si esas
dos propiedades forman parte de la admisibilidad global, entonces
\(X_\infty^{\rm Arch}=X_\infty\); si no, esta restricción de dominio es
necesaria. Definimos la evaluación
\[
\operatorname{ev}_{\mathbb R}:X_\infty^{\rm Arch}\longrightarrow\mathbb R,
\qquad
\operatorname{ev}_{\mathbb R}(x)=
\text{el único punto de }\bigcap_{t\ge0}I_t(x_t).
\]
La unicidad del valor no obliga a la unicidad de la historia. Dos genealogías
\(x\ne y\) pueden satisfacer
\[
\operatorname{ev}_{\mathbb R}(x)
=\operatorname{ev}_{\mathbb R}(y).
\]
Definimos la equivalencia de lectura
\[
x\sim_{\rm ev}y
\quad\Longleftrightarrow\quad
\operatorname{ev}_{\mathbb R}(x)
=\operatorname{ev}_{\mathbb R}(y).
\]

\begin{theorem}[Cociente de la evaluación arquimediana]
La aplicación
\[
\begin{aligned}
\overline{\operatorname{ev}}_{\mathbb R}:
X_\infty^{\rm Arch}/\!\sim_{\rm ev}
&\longrightarrow
\operatorname{Im}(\operatorname{ev}_{\mathbb R})\subseteq\mathbb R,\\
[x]&\longmapsto\operatorname{ev}_{\mathbb R}(x)
\end{aligned}
\]
es una biyección. En particular, el cociente queda identificado con todo
\(\mathbb R\) si y sólo si la evaluación
\(\operatorname{ev}_{\mathbb R}\) es sobreyectiva.
\end{theorem}

\begin{proof}
La definición de \(\sim_{\rm ev}\) hace que el valor sea independiente del
representante, así que la aplicación está bien definida. Es inyectiva porque
dos clases con el mismo valor son, por definición, una sola clase; es
sobreyectiva sobre la imagen por la definición misma de imagen. La última
afirmación equivale a
\(\operatorname{Im}(\operatorname{ev}_{\mathbb R})=\mathbb R\).
\end{proof}

El teorema de cobertura del lector arquimediano construye precisamente esos
prefijos: para todo intervalo compacto \(K\subset\mathbb R\),
\(\operatorname{ev}_{\mathbb R}(\Omega_{\infty,K})=K\), y los intervalos
\([-m,m]\) forman una familia dirigida cuya unión es \(\mathbb R\). Por
consiguiente, cada real posee al menos una historia HMT compatible y la
afirmación demostrada es
\[
X_\infty^{\rm Arch}/\!\sim_{\rm ev}
\cong\operatorname{Im}(\operatorname{ev}_{\mathbb R})
=\mathbb R.
\]
El valor arquimediano es, por tanto, el cociente que identifica las fibras de
la evaluación; el número HMT enriquecido conserva además la clase de
trayectoria, el acarreo y la orientación. La identificación no se sigue de la
mera formación abstracta del cociente: se sigue del teorema de cobertura ya
demostrado en la construcción discreta del continuo.

El descenso algebraico es otro problema y debe resolverse operación por
operación. Supongamos que se han construido dos operaciones candidatas
\(\oplus\) y \(\odot\) sobre \(X_\infty^{\rm Arch}\), cerradas en ese
dominio. La suma desciende al cociente si y sólo si
\[
x\sim_{\rm ev}x',\ y\sim_{\rm ev}y'
\Longrightarrow
x\oplus y\sim_{\rm ev}x'\oplus y',
\]
mientras que el producto desciende, de manera independiente, si y sólo si
\[
x\sim_{\rm ev}x',\ y\sim_{\rm ev}y'
\Longrightarrow
x\odot y\sim_{\rm ev}x'\odot y'.
\]
La primera congruencia no implica la segunda. Además, para que las operaciones
inducidas sean precisamente la suma y el producto ordinarios del cociente
arquimediano se necesita probar por separado
\[
\operatorname{ev}_{\mathbb R}(x\oplus y)
=\operatorname{ev}_{\mathbb R}(x)+\operatorname{ev}_{\mathbb R}(y),
\]
\[
\operatorname{ev}_{\mathbb R}(x\odot y)
=\operatorname{ev}_{\mathbb R}(x)\operatorname{ev}_{\mathbb R}(y),
\]
y que la imagen de la evaluación sea cerrada bajo la operación correspondiente.
Sólo después de estos certificados el cociente hereda la estructura de
subanillo ---o de subcuerpo, si se construyen también opuestos e inversos---
de \(\mathbb R\). La biyección conjuntista del teorema anterior es formal;
el contenido específico de HMT reside en construir la evaluación, caracterizar sus
fibras, demostrar cobertura y establecer cada ley de descenso.

\section{Representación de Weyl--Heisenberg asociada al toro nonádico}

La forma de área orientada del toro nonádico es
\[
\sigma((a,b),(c,d))=ad-bc\pmod9.
\]
Esta forma alternante es la que se obtiene, con la orientación elegida, al
integrar la curvatura mixta de APP sobre paralelogramos del toro. Sea
\(\omega=\exp(2\pi i/9)\) y
\(\mathcal H_9=\mathbb C^9\), con base \(\{|n\rangle:n\in\mathbb Z_9\}\).
Definimos
\[
X|n\rangle=|n+1\rangle,
\qquad
Z|n\rangle=\omega^n|n\rangle.
\]
Entonces
\[
ZX=\omega XZ,
\]
y, más generalmente,
\[
(X^aZ^b)(X^cZ^d)
=\omega^{bc}X^{a+c}Z^{b+d}.
\]
Por tanto
\[
 c\bigl((a,b),(c,d)\bigr)=bc\pmod9
\]
es el \(2\)-cociclo polarizado que aparece en esta sección de la extensión
central. Definimos explícitamente
\[
 \operatorname{Heis}_9
 =\{\omega^kX^aZ^b:a,b,k\in\ZZ/9\ZZ\}.
\]
Es un grupo de orden \(9^3=729\). Se trata de una extensión central asociada
al grupo de traslaciones del soporte \(X_9\) de APP, no de un subgrupo del
grafo ni de la categoría de caminos de APP\@.
El conmutador de los dos transportes es
\[
(X^aZ^b)(X^cZ^d)(X^aZ^b)^{-1}(X^cZ^d)^{-1}
=\omega^{bc-ad}I.
\]
En consecuencia, su exponente es \(-\sigma((a,b),(c,d))\). El cociclo
polarizado \(c\) y la forma \(\sigma\) no son el mismo objeto: la segunda es,
con este convenio, la opuesta de la antisimetrización
\(c(u,v)-c(v,u)\). La fase de conmutador es así la inversa de la fase de
Wilson para la orientación fijada.

Para construir observables, tomamos
\[
Q_{\rm H}=\operatorname{diag}(0,1,\ldots,8),
\qquad
F_{jk}=\frac13\omega^{jk},
\qquad
P_{\rm H}=F^\dagger Q_{\rm H}F.
\]
Los operadores \(Q_{\rm H},P_{\rm H}\) son autoadjuntos. Sus bases propias son mutuamente
insesgadas porque \(|F_{jk}|^2=1/9\). Para todo vector unitario
\(\psi\in\mathcal H_9\), la desigualdad de Robertson da
\[
\operatorname{Var}_\psi(Q_{\rm H})\,
\operatorname{Var}_\psi(P_{\rm H})
\ge
\frac14\left|\langle\psi,[Q_{\rm H},P_{\rm H}]\psi\rangle\right|^2,
\]
y la cota entrópica para bases mutuamente insesgadas produce
\[
H_{Q_{\rm H}}(\psi)+H_{P_{\rm H}}(\psi)\ge\log9.
\]
Aquí \(H\) es la entropía de Shannon con logaritmo natural de las
distribuciones espectrales correspondientes. Las relaciones construidas se
resumen, sin atribuirles una causalidad física adicional, por
\[
\begin{gathered}
\text{forma alternante }(X_9,\sigma)
\longrightarrow
\operatorname{Heis}_9,\\
\operatorname{Heis}_9\longrightarrow(X,Z)
\longrightarrow(Q_{\rm H},P_{\rm H})
\longrightarrow\text{incertidumbre finita}.
\end{gathered}
\]
Éste es un modelo algebraico finito de operadores; por sí solo no especifica
una dinámica física ni unidades de medida.

\section{Dinámica cíclica sobre \texorpdfstring{\(\mathbb Z_9\times\mathbb Z_{12}\)}{Z9 x Z12} y cociente predictivo mínimo}

El sistema dinámico finito \(\mathfrak C_{9,12}\) permite determinar la memoria
mínima necesaria para que la evolución observada sea determinista. Sea
\[
X=\{(b,m):b\in\mathbb Z/9\mathbb Z, 0\le m<12\}
\]
con sucesor
\[
S(b,m)=
\begin{cases}
(b,m+1),&m<11,\\
(b+1\bmod9,0),&m=11.
\end{cases}
\]
Las etiquetas marginales son
\[
\beta(b,m)=4b\pmod{12},
\qquad
d(b,m)=1+(m+\beta(b,m)\pmod{12}).
\]
Para una aplicación de observación \(q\), la palabra futura inclusiva de horizonte \(h\) es
\[
Q_h^q(x)=\bigl(q(x),q(Sx),\ldots,q(S^hx)\bigr).
\]

\begin{theorem}[Cociente predictivo mínimo]
Sean \(X\) finito, \(T:X\to X\) y \(q:X\to V\). Definamos
\[
\operatorname{Eq}(f):=\{(x,y)\in X^2:f(x)=f(y)\},
\qquad
E_h=\operatorname{Eq}(Q_h^q),
\qquad
E_\infty=\bigcap_{k\ge0}\operatorname{Eq}(q\circ T^k).
\]
Entonces
\(E_{h+1}=E_h\cap(T\times T)^{-1}(E_h)\), la sucesión estabiliza en tiempo
finito y \(X/E_\infty\) es el cociente con menor número de estados que
refina a \(q\) y admite una dinámica determinista inducida.
\end{theorem}

\begin{proof}
\(E_h\) exige igualdad de observaciones en los tiempos \(0,\ldots,h\), mientras
\((T\times T)^{-1}(E_h)\) exige igualdad en \(1,\ldots,h+1\). Su intersección es
\(E_{h+1}\). Como \(X\) es finito, el número de clases no puede crecer
indefinidamente. En el primer nivel estable, \(T\) respeta las clases y
desciende. Toda equivalencia \(T\)-compatible contenida en
\(\operatorname{Eq}(q)\)
identifica sólo estados con el mismo futuro completo, por lo que está
contenida en \(E_\infty\); de ahí la propiedad minimal.
\end{proof}

La evaluación exhaustiva de \(\mathfrak C_{9,12}\) da:
\begin{center}
\small
\begin{tabular}{@{}lrrrr@{}}
\toprule
observable & valores instantáneos & índice \(h\) & observaciones & cociente estable\\
\midrule
\(\beta\) & 3 & 11 & 12 & \(C_{36}\), fibra 3\\
\(d\) & 12 & 8 & 9 & \(C_{36}\), fibra 3\\
\((\beta,d)\) & 36 & 0 & 1 & \(C_{36}\), fibra 3\\
\bottomrule
\end{tabular}
\end{center}
En los dos observables marginales,
\[
\operatorname{Eq}(Q_{11}^{\beta})
=\operatorname{Eq}(Q_8^d)
=\operatorname{Eq}(\pi_{36}),
\qquad
\pi_{36}(b,m)=(b\bmod3,m).
\]
El factor \(36\) no se introduce como objetivo. Aparece como el cociente
predictivo mínimo de cada marginal. La fase \(m\) es una coordenada de estado
necesaria para predecir esos futuros; esta afirmación matemática no presupone
una interpretación causal física. La transformación \(b\mapsto b+3\) es una
simetría invisible para esos observables en todo tiempo.

Como sistema dinámico abstracto, el cociente es un ciclo de treinta y seis
estados. Escribir \(C_{36}\) no elige, sin embargo, un origen distinguido:
para obtener una coordenada \(\mathbb Z/36\mathbb Z\) hay que fijarlo. Tampoco
convierte automáticamente esos estados en los treinta y seis pares de una
nónada. En efecto, \(\operatorname{Aut}J(9,2)\cong S_9\). Si los dos puntos de
un par están en ciclos distintos de longitudes \(r,s\) de una permutación, la
órbita del par tiene longitud \(\operatorname{mcm}(r,s)\), con \(r+s\leq9\),
y por tanto longitud a lo sumo \(20\). Si están en un mismo ciclo de longitud
\(r\leq9\), la órbita tiene longitud \(r\), salvo el caso antipodal, en que
tiene longitud \(r/2\). Ningún automorfismo induce, pues, una órbita de longitud
\(36\) sobre los pares. Por tanto, el paso
\[
C_{36}\longleftrightarrow \binom{[9]}2
\]
permanece como identificación calibrada, no como identificación equivariante natural.
Incluso una ordenación hamiltoniana de los pares probaría conservación local
de incidencia paso a paso, no un automorfismo global ni la canonicidad del
diccionario.

\section{Número HMT como historia compatible en \(X_\infty\): identidad genealógica, estructura de lectura y evaluación arquimediana}

Los resultados precedentes determinan la siguiente instancia de la definición
rectora. En la notación de esta sección, un número HMT es una historia
compatible
\[
x\in X_\infty=\varprojlim X_t.
\]
Un lector \(L\) es estructura adicional: el par \((x,L)\) constituye una
presentación leída o una medición, pero no altera el tipo de identidad de
\(x\). Cuando \(x\) pertenece a \(X_\infty^{\rm Arch}\), su valor
arquimediano es la intersección puntual de los cilindros; el conjunto de
valores así construido es
\(\operatorname{Im}(\operatorname{ev}_{\mathbb R})\subseteq\mathbb R\), y
su identificación con toda la recta exige el teorema adicional de cobertura.
Su memoria mínima relativa a un observable es el cociente predictivo de
futuros; y la forma alternante del soporte suma--producto admite la
representación de Weyl--Heisenberg finita construida arriba. El valor, la clase de trayectoria y la incertidumbre quedan
definidos por tres operaciones sobre dominios especificados de una misma
arquitectura de transporte.
 \end{HMTScientificOwnerSubordinated}
  \part{Estructura discreta del continuo y doble proyección holográfica}

\chapter{Sistema inverso de trayectorias compatibles}
\label{chap:83-03}

\section{Construcción genealógica del continuo}

\subsection{APP, TRIT y TPK: definiciones rectoras}

La construcción parte de tres objetos irreductibles y causalmente
ordenados. El objeto de \emph{All Possible Paths} (en adelante, \APP)
construye el soporte nonádico común, su grafo de todos los caminos finitos y
las dos hojas de evaluación suma--producto. La unidad trítica (en adelante,
\TRIT) es la unidad local mínima de relación topológica, geométrica,
informacional y, una vez ordenada, temporal: su firma visible posee tres
estados y su levantamiento conserva el acarreo que una lectura binaria o
residual perdería. El \emph{Topological Phase Kernel} (en adelante, \TPK)
es la categoría de transporte que compone esas unidades en rutas y conserva,
dentro del mismo estado, hoja, orientación, frontera, registro genealógico, memoria, fase y
supervivencia.

\begin{center}
\begin{tabularx}{.94\textwidth}{@{}p{.10\textwidth}Y Y@{}}
\toprule
Objeto & Construcción & Invariantes transportados\\
\midrule
APP & carta positiva, toro finito, grafo de Cayley, caminos y hojas
\(\Sigma,\Pi\) & marca, celda, vecindad, suma, producto, cociente y cierre\\
TRIT & firma ternaria visible y levantamiento topológico--temporal local &
residuo, carry, hoja, orientación, umbral y reversión\\
TPK & categoría de rutas, fibras y extensiones compatibles & cilindro,
frontera, registro genealógico, memoria, fase, holonomía y supervivencia\\
\bottomrule
\end{tabularx}
\end{center}

La construcción exacta de APP se establece sobre un segmento marcado y una
carta finita; la recta real comparece después como publicación límite. El
intervalo de referencia, sus marcas y el camino orientado que las recorre son:
\begin{equation}
I_{10}=[0,10],\qquad
\mathcal M_{10}=I_{10}\cap\mathbb Z=\{0,1,\ldots,10\},
\qquad
\mathsf P_{10}:0\longrightarrow1\longrightarrow\cdots
\longrightarrow9\longrightarrow10,
\label{p12-83-c3-s1:eq:app-marked-segment}
\end{equation}
Los extremos fijan el borde y las nueve marcas interiores
\(I_{10}^{\circ}\cap\mathbb Z\) forman el alfabeto positivo. El objeto
matemático exacto de partida es la carta
\begin{equation}
\mathcal D_9=\{1,\ldots,9\},
\qquad \mathcal D_8=\{1,\ldots,8\},
\qquad
\lambda_9:\mathcal D_9\xrightarrow{\sim}\mathbb Z/9\mathbb Z,
\qquad \lambda_9(9)=[0]_9.
\label{p12-83-c3-s1:eq:app-positive-chart}
\end{equation}
La marca \(9\) es simultáneamente el último representante positivo y la
publicación de la clase nula. Para todo \(n\geq1\), raíz digital positiva y
cociente se separan mediante
\begin{equation}
\operatorname{dr}_9(n)=1+((n-1)\bmod9),
\qquad q_9(n)=\left\lfloor\frac{n-1}{9}\right\rfloor,
\qquad n=9q_9(n)+\operatorname{dr}_9(n).
\label{p12-83-c3-s1:eq:positive-digital-root-carry}
\end{equation}
Al camino \(1\to2\to\cdots\to9\) se añade la
única arista de cierre \(9\to1\). El sucesor y su memoria se separan:
\begin{equation}
u_+(a)=\begin{cases}a+1,&a<9,\\1,&a=9,\end{cases}
\qquad
\kappa_+(a)=\begin{cases}0,&a<9,\\1,&a=9,\end{cases}
\qquad a+1=u_+(a)+9\kappa_+(a).
\label{p12-83-c3-s1:eq:app-successor-carry}
\end{equation}
Por ello \(9\to1\) cierra la fase, aumenta la memoria y prolonga el estado
genealógico.

El cuadrado de la carta produce el toro finito y su grafo de movimientos:
\begin{equation}
X_9=(\mathbb Z/9\mathbb Z)^2,
\qquad
\Gamma_{\rm APP}=\operatorname{Cay}
\bigl(X_9,\{\pm e_1,\pm e_2\}\bigr),
\qquad
\operatorname{Path}(\Gamma_{\rm APP}).
\label{p12-83-c3-s1:eq:app-cayley-torus}
\end{equation}
La identificación de lados opuestos mantiene el cardinal de la carta y añade
la ley de continuación. El núcleo de \(8^2=64\) intersticios se completa con el gnomon
mínimo que publica la clase nula en las dos coordenadas:
\begin{equation}
\mathcal D_9^2=\mathcal D_8^2\sqcup\mathcal G_9,
\qquad
81=64+(9+8)=64+17.
\label{p12-83-c3-s1:eq:app-gnomon-public}
\end{equation}
La misma ley se prolonga a cartas de mayor profundidad; la escala cambia,
pero el cierre, el borde compartido y la memoria de vuelta conservan su tipo.

La completación volumétrica asociada se expresa mediante la identidad de
estratos
\begin{equation}
10^3=(9+1)^3=729+243+27+1.
\label{p12-83-c3-s1:eq:completion-unit}
\end{equation}
\begin{figure}[tbp]
\centering
\begin{minipage}[t]{.53\textwidth}
\centering
\resizebox{.96\linewidth}{!}{%
\begin{tikzpicture}[
  font=\scriptsize,
  cell/.style={minimum width=4.35mm,minimum height=4.35mm,
    inner sep=0pt,draw=black!28}]
  \node[font=\bfseries] at (1.72,.88)
    {Hoja multiplicativa de la carta APP: $\Pi(i,j)=\operatorname{dr}_9(ij)$};
  \foreach \i in {1,...,9}{
    \node[font=\bfseries] at (-.35,{-(\i-1)*.44}) {\i};
    \foreach \j in {1,...,9}{
      \pgfmathtruncatemacro{\v}{mod(\i*\j-1,9)+1}
      \pgfmathtruncatemacro{\border}{(\i==9)||(\j==9)}
      \ifnum\border=1
        \node[cell,fill=black!16,font=\bfseries]
          at ({(\j-1)*.44},{-(\i-1)*.44}) {\v};
      \else
        \node[cell,fill=white]
          at ({(\j-1)*.44},{-(\i-1)*.44}) {\v};
      \fi
    }
  }
  \foreach \j in {1,...,9}{
    \node[font=\bfseries] at ({(\j-1)*.44},.38) {\j};
  }
  \draw[very thick] (3.74,.22)--(3.74,-3.74)--(-.22,-3.74);
  \draw[densely dashed] (-.22,.22) rectangle (3.30,-3.30);
  \node[align=center,anchor=north,font=\scriptsize] at (1.76,-3.98)
    {$8^2=64$ celdas interiores; gnomon de cierre
     $9+8=17$\\corona positiva de nueves: la marca $9$ publica $[0]_9$};
\end{tikzpicture}}
\end{minipage}\hfill
\begin{minipage}[t]{.44\textwidth}
\centering
\resizebox{.98\linewidth}{!}{%
\begin{tikzpicture}[font=\scriptsize]
  \begin{scope}[x={(.78cm,0cm)},y={(.36cm,.26cm)},z={(0cm,.78cm)}]
    \coordinate (O) at (0,0,0);
    \coordinate (X) at (3.1,0,0);
    \coordinate (Y) at (0,3.1,0);
    \coordinate (XY) at (3.1,3.1,0);
    \coordinate (Z) at (0,0,3.1);
    \coordinate (XZ) at (3.1,0,3.1);
    \coordinate (YZ) at (0,3.1,3.1);
    \coordinate (XYZ) at (3.1,3.1,3.1);
    \path[fill=black!5] (O)--(X)--(XZ)--(Z)--cycle;
    \path[fill=black!12] (X)--(XY)--(XYZ)--(XZ)--cycle;
    \path[fill=black!20] (Z)--(XZ)--(XYZ)--(YZ)--cycle;
    \draw[densely dashed] (O)--(Y)--(XY) (Y)--(YZ);
    \draw[thick] (O)--(X)--(XZ)--(Z)--cycle
      (X)--(XY)--(XYZ)--(XZ) (Z)--(YZ)--(XYZ);
    \node[font=\bfseries\large,fill=white,inner sep=1.5pt]
      at (1.55,.52,1.55) {$9^3=729$};
  \end{scope}
  \node[font=\bfseries] at (3.62,3.68) {Cubo de Extrema y Media Razón Holográfica};
  \draw[fill=black!8] (0.2,-.78) rectangle (3.15,-.05);
  \node[align=center,font=\tiny] at (1.68,-.415)
    {interior\\$9^3=729$};
  \draw[fill=black!16] (3.15,-.78) rectangle (5.15,-.05);
  \node[align=center,font=\tiny] at (4.15,-.415)
    {$3$ caras nuevas\\$3\cdot9^2=243$};
  \draw[fill=black!24] (5.15,-.78) rectangle (6.35,-.05);
  \node[align=center,font=\tiny] at (5.75,-.415)
    {$3$ aristas\\$3\cdot9=27$};
  \draw[fill=black!34] (6.35,-.78) rectangle (7.02,-.05);
  \node[align=center,font=\tiny] at (6.685,-.415)
    {$1$ vértice};
  \node[font=\bfseries] at (3.61,-1.30)
    {$10^3=729+243+27+1=729+271$};
  \node[align=center] at (3.61,-1.88)
    {corona de crecimiento $9^3\subset10^3$: $271$ estados\\
     frontera interna del cubo de lado $9$: $386$ puntos};
\end{tikzpicture}}
\end{minipage}
\caption{Dos publicaciones finitas del soporte. Izquierda: la tabla de
producto con raíz digital positiva y su gnomon de cierre. Derecha: la
completación cúbica EMRH, que distingue interior, caras, aristas y vértice.
La corona de crecimiento $271$ y la frontera interna $386$ son dos
recuentos tipados: el primero mide la ampliación $9^3\subset10^3$ y el
segundo la frontera del cubo de lado $9$.}
\label{fig:app-emrh-mini}
\end{figure}
 El volumen interior, las caras, las aristas y el cierre puntual quedan así
tipados como capas distintas de una misma unidad completada. Sobre las 81
celdas de la carta positiva actúan las evaluaciones
\[
\Sigma,\Pi:\mathcal D_9^2\longrightarrow\mathcal D_9,
\qquad
\Sigma(i,j)=\operatorname{dr}_9(i+j),
\qquad
\Pi(i,j)=\operatorname{dr}_9(ij),
\]
Sus residuos y cocientes levantados se definen por
\[
\begin{aligned}
R^+(i,j)&:=\Sigma(i,j),&
Q^+(i,j)&:=\frac{i+j-R^+(i,j)}9,\\
R^\times(i,j)&:=\Pi(i,j),&
Q^\times(i,j)&:=\frac{ij-R^\times(i,j)}9,
\end{aligned}
\]
y, por tanto,
\[
\begin{aligned}
\widehat\Sigma(i,j)&:=\bigl(R^+(i,j),Q^+(i,j)\bigr),\\
\widehat\Pi(i,j)&:=\bigl(R^\times(i,j),Q^\times(i,j)\bigr),
\end{aligned}
\qquad
\widehat\Sigma,\widehat\Pi:
\mathcal D_9^2\longrightarrow\mathcal D_9\times\mathbb Z_{\geq0}.
\]
La biyección \(\lambda_9^{\times2}:\mathcal D_9^2\to X_9\) transporta
estas operaciones y sus levantamientos al toro. Por tanto, APP es la tupla
\begin{equation}
\APP=\bigl(\mathcal D_9,\lambda_9,X_9,\Gamma_{\rm APP},
\operatorname{Path}(\Gamma_{\rm APP}),
\Sigma,\Pi,\widehat\Sigma,\widehat\Pi\bigr),
\label{p12-83-c3-s1:eq:app-full-definition-public}
\end{equation}
La tabla constituye una representación bidimensional de esta tupla
completa. La hoja aditiva es latina; la multiplicativa conserva
el radical ternario del anillo \(\mathbb Z/9\mathbb Z\)
\begin{equation}
R_9:=\mathbb Z/9\mathbb Z,
\qquad
\mathfrak m:=([3]_9)=\{[0]_9,[3]_9,[6]_9\}\triangleleft R_9,
\qquad \mathfrak m^2=(0).
\label{p12-83-c3-s1:eq:app-radical-369}
\end{equation}
La rotación interna es el automorfismo
\[
\rho:\mathfrak m^3\longrightarrow\mathfrak m^3,
\qquad \rho(x_1,x_2,x_3)=(x_2,x_3,x_1),
\]
cuya órbita distinguida es
\[
([3]_9,[6]_9,[0]_9)\xrightarrow{\rho}
([6]_9,[0]_9,[3]_9)\xrightarrow{\rho}
([0]_9,[3]_9,[6]_9)\xrightarrow{\rho}
([3]_9,[6]_9,[0]_9).
\]
La reducción directa envía \(3,6,9\) a \(0,0,0\). La coordenada interna
del ideal conserva, en cambio, el factor extraído:
\begin{equation}
\eta_{\mathfrak m}:(\mathfrak m,+)
\xrightarrow{\;\sim\;}(\mathbb Z/3\mathbb Z,+),
\qquad
\eta_{\mathfrak m}([3a]_9):=[a]_3,
\qquad
([3]_9,[6]_9,[0]_9)
\xmapsto{\;\eta_{\mathfrak m}^{\times3}\;}
([1]_3,[2]_3,[0]_3),
\label{p12-83-c3-s1:eq:app-369-120}
\end{equation}
de modo que el sector \(3,6,9\) permanece íntegro y su coordenada interna se
lee como \(1,2,0\). Esta aplicación pertenece al radical de APP; la
proyección temporal de TRIT constituye el lector calendárico posterior. Las
tres clases calendáricas
quedan separadas como
\[
\mathcal C_+=\{1,4,7\},\qquad
\mathcal C_-=\{2,5,8\},\qquad
\mathcal C_0=\{3,6,9\}.
\]
La incompatibilidad entre ambas hojas se vuelve un invariante medible. Sobre
\(\mathscr A_3=\mathcal D_9^3\), con medida uniforme
\(\nu_3(x)=9^{-3}\), sea
\(\mathcal H_3=L^2(\mathscr A_3,\nu_3)\). Definimos
\[
\Sigma_3,\Pi_3:\mathscr A_3\longrightarrow\mathcal D_9,
\qquad
\Sigma_3(x)=\operatorname{dr}_9(x_1+x_2+x_3),
\qquad
\Pi_3(x)=\operatorname{dr}_9(x_1x_2x_3),
\]
y los proyectores ortogonales
\[
P_{\Sigma,3}:=\mathbb E_{\nu_3}[\,\cdot\mid\sigma(\Sigma_3)],
\qquad
P_{\Pi,3}:=\mathbb E_{\nu_3}[\,\cdot\mid\sigma(\Pi_3)]
\in\mathcal B(\mathcal H_3).
\]
La enumeración conjunta de las \(9^3\) palabras determina entonces
\[
\bigl\|[P_{\Sigma,3},P_{\Pi,3}]\bigr\|=\frac{\sqrt3}{4},
\qquad
B_c=7I+2R=9P_++5P_-,
\qquad
459-405=54.
\]
El primer número reaparece como prefactor electrónico; el segundo organiza
el bloque central; el tercero registra el defecto global suma--producto.

\subsection{TRIT: unidad local topológica y unidad de relación temporal}

HMT fija primero la firma visible
\begin{equation}
\mathrm{TRIT}_{\rm vis}=\{-1,0,+1\},
\qquad \tau\in\mathrm{TRIT}_{\rm vis}.
\label{p12-83-c3-s1:eq:trit-visible-signature}
\end{equation}
Frente a la unidad binaria, el TRIT conserva tres regímenes locales. En la
distribución uniforme, su información y su entropía se distinguen como
\begin{equation}
I_{\rm TRIT}=\log3\ \text{nats}=\log_2 3\ \text{bits},
\qquad S_{\rm TRIT}=k_BI_{\rm TRIT}.
\label{p12-83-c3-s1:eq:trit-information-unit}
\end{equation}
El bit sigue siendo la unidad binaria; el TRIT es la unidad mínima capaz de
distinguir simultáneamente retorno, umbral y apertura. La igualdad
\eqref{p12-83-c3-s1:eq:trit-information-unit} cuantifica la información de la terna; el
coste \(k_B\Theta\log3\) pertenece al lector térmico de reinicio posterior.
Mecánica Dimensional conserva además el antecedente aritmético de esa firma.
Para \((i,j)\in\mathcal D_9^2\), la celda aritmética levantada asociada es
\begin{equation}
c_{\rm MD}(i,j;\tau)=
\bigl(i,j;R^+,Q^+,R^\times,Q^\times;\tau\bigr).
\label{p12-83-c3-s1:eq:trit-complete-local-unit}
\end{equation}
Los residuos \(R^+,R^\times\) son la publicación visible de las dos hojas;
los cocientes \(Q^+,Q^\times\) conservan la memoria exacta de la división
euclídea; y \(\tau\) orienta la transición. La firma visible ocupa el primer
nivel; la celda que la porta añade los antecedentes aritméticos, y el estado
holonómico integral completa después frontera, ruta, fase y supervivencia.

Conviene separar cuatro estratos. La firma visible es \(\tau\); el par
residuo--acarreo es su levantamiento local; \(c_{\rm MD}\) es la celda
aritmética de APP provista de esa firma; y el estado TPK añade extensión,
frontera, registro genealógico, fase y supervivencia. Esta jerarquía conserva la identidad
propia de la coordenada observada, su levantamiento, la celda portadora y la
historia enriquecida que la transporta.

En la carta ternaria balanceada, el entero local levantado
\(n\in\{-4,-3,\ldots,3,4\}\) se descompone de manera única como
\begin{equation}
\widehat n=(r,\kappa),
\qquad n=r+3\kappa,
\qquad r,\kappa\in\{-1,0,+1\}.
\label{p12-83-c3-s1:eq:trit-balanced-lift}
\end{equation}
La firma visible es \(\tau=r\); \(\kappa\) conserva el carry que la proyección
\(n\mapsto r\) olvidaría.
La fibra sobre el cero visible contiene tres estados distintos,
\[
0^+=(0,+1),\qquad 0^0=(0,0),\qquad 0^-=(0,-1),
\]
porque \(n=-3,0,3\) poseen el mismo residuo visible y carries distintos. Éste
es el levantamiento topológico mínimo de MD: una transición orientada de APP
queda provista de la información necesaria para invertirla y para decidir qué
prolongaciones siguen siendo compatibles.

La firma posee además una regla de memoria de hoja. Si \(m_k\) es el modo
aritmético activo, entonces
\begin{equation}
m_{k+1}=
\begin{cases}
\Sigma,&\tau_k=+1,\\
\Pi,&\tau_k=-1,\\
m_k,&\tau_k=0.
\end{cases}
\label{p12-83-c3-s1:eq:trit-zero-memory}
\end{equation}
El cero codifica retención de la hoja anterior. Por eso la palabra
\((+1,-1,0)\) publica la secuencia \((\Sigma,\Pi,\Pi)\): el tercer estado
conserva el modo multiplicativo activo.

La relación temporal se obtiene al parametrizar ordenadamente estas celdas.
Cada firma porta un álgebra cuadrática
\begin{equation}
\mathbb A_\tau=\mathbb R[X]/(X^2+\tau),
\qquad J_\tau=[X],
\label{p12-83-c3-s1:eq:trit-quadratic-algebra-local}
\end{equation}
de modo que
\[
\begin{array}{c@{\qquad}c@{\qquad}l}
\toprule
\tau & J_\tau^2 & \text{relación temporal ordenada}\\
\midrule
+1 & -I & \text{retorno elíptico, memoria y pasado},\\
0  & 0  & \text{umbral parabólico y presente},\\
-1 & +I & \text{apertura hiperbólica bidireccional y futuro}.\\
\bottomrule
\end{array}
\]
El TRIT precede al tiempo exterior: primero clasifica el tipo algebraico y
topológico del transporte; después, una parametrización ordenada lo publica
como relación temporal. La doble dirección queda tipada por la involución
exacta
\begin{equation}
\mathcal C_{\rm rev}(\tau_1,\ldots,\tau_n)
=(-\tau_n,\ldots,-\tau_1),
\qquad \mathcal C_{\rm rev}^2=I,
\label{p12-83-c3-s1:eq:trit-time-reversal}
\end{equation}
que intercambia retorno y apertura y conserva el umbral. En la hoja
hiperbólica, \(P_\pm=(I\pm J_{-1})/2\) separan las dos direcciones propias y
la reversión intercambia sus orientaciones. Topología y tiempo son, así, las
dos lecturas coordinadas del TRIT.

La involución combinatoria \(\mathcal C_{\rm rev}\) y una conjugación
operatorial \(C\) pertenecen, sin embargo, a dominios distintos. En una
realización dinámica separada, provista de \(H=H^*\),
\(CJ_EC^{-1}=-J_E\) y \(CHC^{-1}=H\), fijamos una escala de acción interna
\(\hbar_{\rm int}>0\) y escribimos
\(h_{\rm int}:=2\pi\hbar_{\rm int}\). Entonces
\begin{equation}
U(t)=\exp(-J_EHt/\hbar_{\rm int}),
\qquad CU(t)C^{-1}=U(-t).
\label{p12-83-c3-s1:eq:trit-bidirectional-propagator}
\end{equation}
El entrelazador que transporta palabras TRIT a la realización dinámica fija
la correspondencia entre \(\mathcal C_{\rm rev}\) y \(C\). La doble
orientación local proporciona el dominio de esa construcción.

\subsection{TPK: categoría de transporte y evolución de fase}

El \emph{Topological Phase Kernel} es la categoría de transporte sobre la
base observable de APP que compone cifras, palabras, operadores y relaciones
de prolongación en un mismo estado holonómico integral. Su arquitectura interna se distribuye en
ocho clases operatorias, diferenciadas por dominio, codominio y función:
\begin{enumerate}[label=\textup{(\roman*)},leftmargin=2.25em]
\item soporte observable y estado holonómico integral;
\item selector interno de hoja \(M_{\rm ph}:\{1,\ldots,9\}\to\{+1,-1,0\}\);
\item traslaciones cardinales y evolución visible \(F_{\rm obs}\);
\item lectores y cocientes nonádico, trítico y de bloque \(1000\);
\item memoria, acarreo, cilindro, frontera y registro genealógico;
\item periodicidad y retorno
\(C_{12}\hookrightarrow C_{108}\twoheadrightarrow C_9\),
\(\mathcal K_{27}\) y canales \(90/120\);
\item prolongaciones \(\operatorname{Ext}_r\), cilindros, supervivencia y
extensión natural bidireccional;
\item funtores de salida arquimediana, excepcional calibrada y dimensional.
\end{enumerate}
El retorno de 27 pasos es el morfismo compuesto
\begin{equation}
\mathcal K_{27}:=F_{\rm obs}^{27},
\qquad \mathcal K_{27}^{4}=F_{\rm obs}^{108}.
\label{p12-83-c3-s1:eq:tpk-kernel-27}
\end{equation}
Su cuarto iterado restaura la fase observable del calendario, mientras la
memoria y el estado TPK completo continúan su trayectoria;
\(\mathcal K_{27}\) es el morfismo de retorno interno de \(C_{108}\).
El calendario \(C_{108}\), la conexión nonádica, la extensión bilateral y
los lectores ocupan lugares tipados dentro de esta arquitectura. Sea
\(\mathsf P_{\rm APP}\) la categoría
libre del grafo dirigido de APP y
\(\mathsf P_{\rm APP}^{(2)}=\mathsf P_{\rm APP}\times
\mathsf P_{\rm APP}\) su lectura simultánea suma--producto. Sea además
\[
\mathcal Q_{\rm obs}
=X_9^2\times\{N,E,S,O\}^2\times\Theta_{108},
\qquad \Theta_{108}=\mathbb Z/108\mathbb Z,
\]
donde, para el representante
\(\widetilde\theta\in\{0,\ldots,107\}\),
\[
r(\theta)=1+(\widetilde\theta\bmod9),
\qquad g_0(\theta)=[\widetilde\theta]_9,
\]
\[
\tau(\theta)=
\begin{cases}
+1,&r(\theta)\in\{1,4,7\},\\
-1,&r(\theta)\in\{2,5,8\},\\
0,&r(\theta)\in\{3,6,9\},
\end{cases}
\qquad
\mu(\theta)=
\begin{cases}
\Sigma,&\tau(\theta)=+1,\\
\Pi,&\tau(\theta)=-1,\\
\varnothing,&\tau(\theta)=0.
\end{cases}
\]
La etiqueta calendárica \(\mu=\varnothing\) señala el umbral, y la memoria
dinámica \(m_k\) de \eqref{p12-83-c3-s1:eq:trit-zero-memory} retiene la última hoja
activa. Una configuración
\(q=(p^\Sigma,p^\Pi,d^\Sigma,d^\Pi,\theta)\) conserva las dos posiciones,
las dos direcciones cardinales y el índice de calendario. La evolución
visible \(F_{\rm obs}\) activa la hoja indicada por TRIT, transporta la
posición correspondiente y avanza \(\theta\). Denotemos por
\(\mathsf P_{\rm obs}\) la categoría libre generada por las flechas
\(q\to F_{\rm obs}(q)\). Cada flecha induce un par de caminos de APP y
define el funtor de base
\begin{equation}
B:\mathsf P_{\rm obs}\longrightarrow\mathsf P_{\rm APP}^{(2)}.
\label{p12-83-c3-s1:eq:tpk-bivariant-base-functor}
\end{equation}

Sobre cada \(q\) se dispone la fibra tipada
\begin{equation}
\mathcal F_q=
\mathsf O_q\times\mathsf H_q\times\mathsf C_q\times\mathsf S_q
\times\mathsf B_q\times\mathsf\Lambda_q,
\label{p12-83-c3-s1:eq:tpk-fiber}
\end{equation}
Sus seis factores conservan, respectivamente, orientación y hoja; residuo y
cociente; acarreo; firma; cilindro y frontera; y registro ordenado de
transporte. Memoria y supervivencia son propiedades de la composición y de
las relaciones de prolongación.

Para una ruta \(r:q\to q'\), la relación
\begin{equation}
\operatorname{Ext}_r\subseteq\mathcal F_q\times\mathcal F_{q'}
\label{p12-83-c3-s1:eq:tpk-extension-relation}
\end{equation}
contiene todas las continuaciones admisibles y satisface
\[
\Delta_q\subseteq\operatorname{Ext}_{\operatorname{id}_q},
\qquad
\operatorname{Ext}_{r_2}\circ\operatorname{Ext}_{r_1}
\subseteq\operatorname{Ext}_{r_2\circ r_1}.
\]
Los transportes respetan identidades y composición:
\begin{equation}
\begin{gathered}
\mathsf H(\operatorname{id}_q)=I,
\quad \mathsf C(\operatorname{id}_q)=I,
\quad \mathsf\Lambda(\operatorname{id}_q)=I,
\quad \kappa(\operatorname{id}_q)=0,\\
\mathsf H(r_2\!\circ r_1)=\mathsf H(r_2)\mathsf H(r_1),
\quad
\mathsf C(r_2\!\circ r_1)=\mathsf C(r_2)\mathsf C(r_1),
\quad
\mathsf\Lambda(r_2\!\circ r_1)=
\mathsf\Lambda(r_2)\mathsf\Lambda(r_1).
\end{gathered}
\label{p12-83-c3-s1:eq:tpk-functorial-transport}
\end{equation}
Junto con estas leyes functoriales, las coordenadas transportadas obedecen
\begin{equation}
\begin{aligned}
h'&=\mathsf H(r)(h),\\
c'&=\mathsf C(r)(c)+\kappa(r),\\
\lambda'&=\mathsf\Lambda(r)(\lambda),\\
\sigma'&=\operatorname{Sig}_{q'}(h',c',\lambda'),
\end{aligned}
\qquad
\kappa(r_2\circ r_1)
=\kappa(r_2)+\mathsf C(r_2)\kappa(r_1).
\label{p12-83-c3-s1:eq:tpk-typed-transport}
\end{equation}
La última identidad es el cociclo de acarreo que hace asociativa la
composición con memoria.

El TPK es entonces la categoría \(\mathsf{TPK}\) cuyos objetos son
\begin{equation}
x=(q,o,h,c,\sigma,C,b,\lambda),
\qquad (o,h,c,\sigma,(C,b),\lambda)\in\mathcal F_q,
\label{p12-83-c3-s1:eq:tpk-enriched-object}
\end{equation}
y cuyos morfismos \(x\to x'\) son las ternas \((r,x,x')\) tales que
las coordenadas satisfacen \eqref{p12-83-c3-s1:eq:tpk-typed-transport} y
\((f_x,f_{x'})\in\operatorname{Ext}_r\), donde
\(f_x\in\mathcal F_q\) y \(f_{x'}\in\mathcal F_{q'}\). La composición es
\begin{equation}
(r_2,x',x'')\circ(r_1,x,x')
=(r_2\circ r_1,x,x''),
\label{p12-83-c3-s1:eq:tpk-morphism-composition}
\end{equation}
\Needspace{7\baselineskip}
la identidad es \((\operatorname{id}_q,x,x)\), y el funtor
\begin{equation}
p:\mathsf{TPK}\longrightarrow\mathsf P_{\rm obs},
\qquad p(x)=q,
\qquad p(r,x,x')=r,
\label{p12-83-c3-s1:eq:tpk-base-functor}
\end{equation}
publica la configuración visible sin identificar estados de la misma fibra.
La definición fija una categoría sobre una base; no presupone que sea una
fibración categórica ni que toda flecha admita un levantamiento cartesiano.
Las proyecciones locales quedan separadas por tipo:
\[
\pi_{\rm TRIT}:\operatorname{Ob}(\mathsf{TPK})
\longrightarrow\mathrm{TRIT}_{\rm vis},
\qquad
\pi_{\rm TRIT}(x)=\tau(\theta),
\]
\[
\pi_{\rm loc}(x)=
\bigl(\pi_{\rm TRIT}(x),\mu(\theta),o,h,c,
\sigma,\lambda,g_0(\theta)\bigr).
\]
La primera recupera únicamente la firma TRIT visible. La segunda conserva
además la etiqueta calendárica de hoja, la orientación, el par
residuo--cociente, el acarreo, la firma de transporte, el registro genealógico y la fase,
pero todavía olvida el cilindro, su etiqueta de borde y otras coordenadas de
\(q\). Ésta es la razón por la que ni el TRIT visible ni su registro local
sustituyen al estado holonómico integral del TPK.

El nombre recoge tres funciones inseparables. \emph{Topological} remite a
la categoría de caminos, la incidencia, el borde y la prolongación por
cilindros; \emph{Phase} remite al índice nonádico, al calendario orientado y
a la holonomía; \emph{Kernel} nombra el núcleo persistente de transporte que
permanece antes de contraer el estado mediante un lector. Formalmente, el
objeto es la categoría anterior: una topología adicional sobre objetos o
morfismos, o su identificación con el núcleo lineal de un mapa, requiere
datos independientes.

En consecuencia, dos rutas con el mismo extremo visible pueden seguir siendo
objetos distintos por hoja, cociente, acarreo, frontera, cilindro, fase o
registro. El TPK no elige una salida escalar: conserva el antecedente común
del que parten las publicaciones arquimediana, excepcional y dimensional.

% Estas consecuencias se publican después de la construcción correlativa y
% del paso al límite de los capítulos 4 y 5. Se definen aquí para preservar
% íntegramente su procedencia, pero no se ejecutan antes de que exista el
% continuo al que se aplican.
\long\def\HMTDeferredContinuumConsequences{%
\section{Consecuencias y publicaciones del continuo construido}

\subsection{Teorema de ensamblaje terminal de las construcciones correlativas generadas por el estado común}

Cada profundidad \(k\) posee un estado holonómico integral \(\mathcal E_k\) con
todas sus emisiones supervivientes y un único constructor correlativo
\begin{equation}
\mathcal O_k:\mathcal E_k\longrightarrow\mathcal Z_k,
\qquad
\widehat x_k\longmapsto
\mathcal O\bigl(\{\mathfrak e_\alpha:
\alpha\in\operatorname{Ch}_k(\widehat x_k)\}\bigr).
\label{p12-83-c3-s1:eq:five-features-single-constructor}
\end{equation}
El valor de \(\mathcal O_k\) conserva simultáneamente las coordenadas
espectral, solenoidal, gauge, coherente y determinantal,
\begin{equation}
(\mathrm{sp},\mathrm{sol},\mathrm{gau},\mathrm{coh},\mathrm{det}),
\label{p12-83-c3-s1:eq:five-features-coordinates}
\end{equation}
como cinco características del mismo estado y del mismo registro genealógico. Los
truncamientos
\[
u^{\mathfrak e}_{\ell,k}:\mathcal E_\ell\to\mathcal E_k,
\qquad
u^{\mathcal O}_{\ell,k}:\mathcal Z_\ell\to\mathcal Z_k
\]
satisfacen la naturalidad única
\begin{equation}
u^{\mathcal O}_{\ell,k}\circ\mathcal O_\ell
=\mathcal O_k\circ u^{\mathfrak e}_{\ell,k},
\qquad \ell\ge k.
\label{p12-83-c3-s1:eq:five-features-naturality}
\end{equation}
Por ello, la acción correlativa pasa al límite como
\begin{equation}
\mathcal O_\infty:
\mathcal C_{\rm cont}^{\rm disc}:=\varprojlim_k\mathcal E_k
\longrightarrow
\mathcal Z_\infty:=\varprojlim_k\mathcal Z_k.
\label{p12-83-c3-s1:eq:five-features-infinite-operator}
\end{equation}
Las cinco aplicaciones terminales son vistas intrínsecas de ese único valor:
\begin{equation}
\boxed{
\mathfrak G_T^{\rm int}:=
\operatorname{View}_T\!\left(
\mathcal O_\infty(\mathcal C_{\rm cont}^{\rm disc})\right),
\quad
T\in\{\mathrm{sp},\mathrm{sol},\mathrm{gau},
\mathrm{coh},\mathrm{det}\}.}
\label{p12-83-c3-s1:eq:five-features-views}
\end{equation}
La simultaneidad precede a las resoluciones especializadas: cada vista expone
la coordenada requerida y el estado completo conserva las otras cuatro.

\subsection{Supervivencia y límite inverso}

Sea \(\operatorname{Surv}^{(\chi)}_N\) el conjunto de historias de
profundidad \(N\) que satisfacen simultáneamente las leyes activas del canal
\(\chi\) y admiten prolongación. Definimos las truncaciones compatibles
\[
\tau_{N+1,N}:\operatorname{Surv}^{(\chi)}_{N+1}
\longrightarrow\operatorname{Surv}^{(\chi)}_N
\]
Estas conservan el prefijo enriquecido. El continuo genealógico es el límite
\begin{equation}
X_\chi=\varprojlim_N
\bigl(\operatorname{Surv}^{(\chi)}_N,\tau_{N+1,N}\bigr),
\qquad
\tau_{N+1,N}(x_{N+1})=x_N.
\label{p12-83-c3-s1:eq:inverse-survival}
\end{equation}
El espacio global de secciones supervivientes y una de sus secciones
enriquecidas se denotan, respectivamente, por
\begin{equation}
X_\infty:=\bigsqcup_\chi X_\chi,
\qquad
\widehat\Omega_\infty\in X_\infty.
\label{p12-83-c3-s1:eq:global-survival-space}
\end{equation}
El lector arquimediano global es
\begin{equation}
P_{\rm arq}:=\operatorname{ev}_{\mathbb R}:
X_\infty\longrightarrow\mathbb R;
\label{p12-83-c3-s1:eq:archimedean-reader}
\end{equation}
su restricción a cada canal evalúa la intersección unitaria de los cilindros
compatibles de la sección correspondiente.
Cuando interviene el lector excepcional, \(X_\infty^{\rm cal}\subseteq
X_\infty\) denota el subespacio en el que se ha fijado el tipo de calibración
de \eqref{p12-83-c3-s1:eq:exceptional-calibration}.
La supervivencia determina la sección por compatibilidad interna: cada
historia debe ser prolongable en el sistema completo. El lector numérico
interviene después para certificar y publicar la intersección del
descendiente ya compatible.

\subsection{Conexión nonádica conjunta, holonomía y memoria vertical}

En profundidad \(K\), la holonomía del retorno nonádico es inicialmente la
correspondencia tipada
\begin{equation}
\Gamma_{9,K}=\operatorname{Hol}_{C_9}(\nabla_{\rm TPK}):
\widehat{\mathcal X}_K\rightrightarrows
\widehat{\mathcal X}_{K+9}.
\label{p12-83-c3-s1:eq:gamma9}
\end{equation}
Sus iterados son composiciones relacionales. Sólo cuando cada estado posee
una prolongación compatible única se restringe a un mapa, y sólo cuando ese
mapa es invertible constituye una monodromía. Esta distinción conserva todas
las extensiones supervivientes antes de seleccionar una sección.
Tras nueve fases retorna la fase local, mientras el estado completo avanza:
\[
g(K+9)=g(K),\qquad q_{{\rm mem},K+9}=q_{{\rm mem},K}+1,
\qquad \widehat x_{K+9}\ne\widehat x_K\ \text{en general}.
\]
Bloque, cilindro, prefijo, acarreo y memoria continúan evolucionando. Ésta
es la diferencia entre periodicidad visible y retorno con biografía.

El estado transportado en el nivel (K) puede exhibirse como
\begin{equation}
\widehat x_K=
(B_K,C_K,p_K,c_K,\Sigma_K,W_K,g(K),q_{{\rm mem},K},s_K),
\label{p12-83-c3-s1:eq:enriched-state-k}
\end{equation}
donde cada coordenada tiene una ley de transporte propia. El retorno de fase
se representa sobre el catálogo
\[
\mathcal U_6^{\rm cat}\cong\mathfrak S_+\times\mathfrak S_\times,
\qquad |\mathfrak S_+|=18,
\qquad |\mathfrak S_\times|=26.
\]
Los promedios condicionales
\[
(E_+f)(s,e)=\frac1{18}\sum_{s'\in\mathfrak S_+}f(s',e),
\qquad
(E_\times f)(s,e)=\frac1{26}\sum_{e'\in\mathfrak S_\times}f(s,e')
\]
conmutan. La palabra de fase
\(\tau=(+1,-1,0,+1,-1,0,+1,-1,0)\) aplica
\(E_+,E_\times,I\) en ese orden, y una vuelta se cierra mediante
\begin{equation}
\mathcal K_{\rm ph}^{\,9}
=(E_+E_\times)\otimes I_{\ell^2(C_9)},
\qquad
\bigl\langle\delta_u,E_+E_\times\delta_v\bigr\rangle
=\frac1{468}\quad(u,v\in\mathcal U_6^{\rm cat}).
\label{p12-83-c3-s1:eq:joint-nine-holonomy}
\end{equation}
Las hojas suma y producto no actúan como filtros independientes: son dos
componentes de una única holonomía que sólo se cierra después de recorrer el
ciclo completo.

La coligación compresión--expresión separa lo que vuelve y lo que permanece
como memoria. Si (J) incrusta el espacio de historias en el espacio de
fase, (T) es la compresión superviviente y \(\eta\) el canal de defecto,
entonces
\begin{equation}
\boxed{
U_{\Gamma_9}J=JT+\eta,
\qquad J^*\eta=0,
\qquad I-T^*T=\eta^*\eta.}
\label{p12-83-c3-s1:eq:compression-expression}
\end{equation}
La primera igualdad descompone el transporte en retorno comprimido y
expresión de memoria; la segunda prueba su ortogonalidad; la tercera mide
exactamente la información que no cabe en la sombra periódica. Así queda
formalizada la afirmación central: vuelve la fase, no el estado completo.

La memoria vertical se organiza mediante la extensión cíclica orientada
\begin{equation}
0\longrightarrow C_{12}\xrightarrow{\,\iota\,}C_{108}
\xrightarrow{\,p\,}C_9\longrightarrow0,
\qquad
c(r,r')=\left[\left\lfloor\frac{r+r'}9\right\rfloor\right]_{12}.
\label{p12-83-c3-s1:eq:c108-extension}
\end{equation}
Aquí \(\iota([b]_{12})=[9b]_{108}\) y
\(p([\theta]_{108})=[\theta]_9\). La extensión no se escinde:
\(C_{108}\) tiene exponente 108, mientras
\(C_{12}\times C_9\) tiene exponente 36.
El cociclo registra el acarreo de cada vuelta de la base \(C_9\) en la
memoria dodecafásica \(C_{12}\). \(C_{108}\) es, por tanto, un calendario
orientado de fase y memoria. Sólo después de la carta dimensional esa
extensión se realiza como reloj cuántico finito de 108 pasos.

El levantamiento local de MD que porta la firma TRIT es la unidad
topológico--temporal de ese cristal genealógico.
Sobre
\(\mathcal H_{\rm clk}=\mathbb C^{108}\), con base
\(\{|j\rangle:j\in\mathbb Z/108\mathbb Z\}\), sean
\[
\zeta=e^{2\pi\ii/108},\qquad
|\widetilde k\rangle=\frac1{\sqrt{108}}
\sum_{j=0}^{107}\zeta^{jk}|j\rangle,
\qquad U_{\rm clk}|j\rangle=|j+1\rangle.
\]
\Needspace{12\baselineskip}
Con \(\omega_0=2\pi/(108t_0)\), el Hamiltoniano
\begin{equation}
H_{\rm clk}=\hbar_{\rm int}\omega_0
\sum_{k=0}^{107}k\,|\widetilde k\rangle\langle\widetilde k|
\label{p12-83-c3-s1:eq:trit-time-crystal-hamiltonian}
\end{equation}
genera exactamente
\begin{equation}
U_{\rm step}=\exp\!\left(
-\frac{\ii H_{\rm clk}t_0}{\hbar_{\rm int}}\right)=U_{\rm clk},
\qquad U_{\rm step}^{108}=I,
\label{p12-83-c3-s1:eq:trit-time-crystal-period}
\end{equation}
Si \(Z_{\rm clk}|j\rangle=\zeta^j|j\rangle\), entonces
\begin{equation}
\langle j_0|U_{\rm step}^{-n}Z_{\rm clk}U_{\rm step}^{n}|j_0\rangle
=\zeta^{j_0+n},
\label{p12-83-c3-s1:eq:trit-time-crystal-observable}
\end{equation}
cuya órbita posee periodo primitivo 108. En el vocabulario HMT, éste es el
cristal temporal genealógico interno: el calendario observable cierra,
\(C_{12}\) registra la vuelta y el estado holonómico integral no tiene por qué
retornar. La realización ulterior como fase física macroscópica exige,
además, protocolo de Floquet, respuesta subarmónica primitiva, robustez,
persistencia y un mecanismo de protección; esa carta no altera el teorema
finito de periodicidad unitaria.

\subsection{Publicación arquimediana, número genealógico y medición}

La publicación arquimediana utiliza cilindros ternarios semiabiertos
\begin{equation}
I_{N,L}=\left[\frac{N}{3^L},\frac{N+1}{3^L}\right),
\qquad \operatorname{diam}I_{N,L}=3^{-L}\longrightarrow0.
\label{p12-83-c3-s1:eq:ternary-cylinders}
\end{equation}
Una sección compatible determina cilindros encajados y, por completitud, un
único punto real. El número HMT es la sección misma,
\begin{equation}
x=(x_N)_N\in X_\chi
=\varprojlim_N\operatorname{Surv}^{(\chi)}_N,
\label{p12-83-c3-s1:eq:hmt-number-as-section}
\end{equation}
no la tupla formada por ella y sus evaluaciones. El valor aparece cuando un
lector olvida parte de la genealogía.
Por tanto, la precisión es profundidad: medir significa descender a un
cilindro más fino mientras el prefijo causal permanece en el antecedente.
La recta real se reconstruye como cociente arquimediano de un espacio de
historias mucho más rico.

\subsection{Realización dimensional del estado HMT: mapa de pantalla, magnitud observable y carta espectral}

La Mecánica Dimensional comienza donde el estado genealógico se acopla a una
frontera de realización. Fijado un cuerpo \(K_n\), sea \(M_n(x)\) el módulo
generado por los canales de prolongación de un estado \(x\), sea
\(\mathcal R_n(x)\) el submódulo de sus relaciones y sea \(\beta_n(x)\) el
dato de borde conservado. La pantalla dimensional es
\begin{equation}
\Sigma_n(x)=\bigl(K_n,M_n(x),\beta_n(x),\mathcal R_n(x)\bigr),
\qquad
d_{\Sigma_n}(x)=\dim_{K_n}\frac{M_n(x)}{\mathcal R_n(x)}.
\label{p12-83-c3-s1:eq:md-screen}
\end{equation}
Si \(M_n(x)\simeq K_n^{g_n}\) y las columnas de \(R_n(x)\) generan las
relaciones, entonces
\begin{equation}
\boxed{d_{\Sigma_n}(x)=g_n-\operatorname{rank}_{K_n}R_n(x).}
\label{p12-83-c3-s1:eq:md-rank}
\end{equation}
La dimensión no se adjunta como etiqueta: es el número de canales
independientes que sobreviven después de imponer las relaciones de frontera.

Una magnitud física es una sección de una recta dimensional \(L_D\); una
unidad es una carta \(\chi_u:L_D\to\mathbb R\) que publica su coordenada.
Por ello
\begin{equation}
\text{estado TPK}\longrightarrow\Sigma_n(x)
\longrightarrow s_D\in\Gamma(L_D)
\xrightarrow{\;\chi_u\;}[s_D]_u\in\mathbb R
\label{p12-83-c3-s1:eq:md-measurement-chain}
\end{equation}
separa de forma exacta estructura, tipo dimensional y número medido. Medir
consiste en acoplar una sección a un aparato reversible, imponer
compatibilidad de borde y aplicar el lector; cambiar de unidades cambia
\(\chi_u\), no la sección ni su genealogía.

\subsection{Doble proyección matemática y realización dimensional}

La rama excepcional conserva datos de calibre que la publicación real no
necesita. Sea \(\mathscr C_{\rm exc}\) el espacio de calibraciones y sea
\begin{equation}
\mathfrak E:
X_\infty^{\rm cal}\times\mathscr C_{\rm exc}
\longrightarrow\mathbf{Exc}
\label{p12-83-c3-s1:eq:exceptional-calibrated-reader}
\end{equation}
el lector de incidencia. Fijamos una vez
\begin{equation}
\begin{aligned}
\chi_{\rm exc}=\bigl(&\text{orden proyectivo},\ \text{calibre de Paley},
 \ \text{dodecada},\\
 &\text{alineamiento de incidencia},\ \text{cámara reticular}\bigr)
 \in\mathscr C_{\rm exc}
\end{aligned}
\label{p12-83-c3-s1:eq:exceptional-calibration}
\end{equation}
y abreviamos
\(\mathfrak E_{\chi_{\rm exc}}(x):=\mathfrak E(x,\chi_{\rm exc})\).
El calibre inicial se transporta por la historia; no vuelve a elegirse en
cada profundidad.

El mismo límite calibrado admite entonces dos lectores con codominios
distintos:
\begin{equation}
\mathbb R
\xleftarrow{\;\operatorname{ev}_{\mathbb R}\;}
X_\infty^{\rm cal}
\xrightarrow{\;\mathfrak E_{\chi_{\rm exc}}\;}
\mathbf{Exc}.
\label{p12-83-c3-s1:eq:double-projection}
\end{equation}
La primera rama contrae cilindros hasta un valor continuo. La segunda
conserva incidencias, orientaciones, calibre y frontera. Ambas descienden en
paralelo del mismo estado fuente. Esta bifurcación explica conjuntamente la
continuidad métrica y la emergencia de simetrías excepcionales. La segunda
rama es una reconstrucción calibrada: es exacta una vez fijado
\(\chi_{\rm exc}\), pero el estado desnudo no selecciona retrospectivamente
los cinco datos que componen ese calibre.

La realización dimensional no serializa esas dos ramas. Es un tercer lector
del mismo antecedente:
\begin{equation}
P_{\rm dim}:X_\infty\longrightarrow
\mathcal R_{\rm MD},
\qquad
\mathcal R_{\rm MD}=
\{\text{acción, masa, campo, geometría, información}\}.
\label{p12-83-c3-s1:eq:dimensional-publication}
\end{equation}
De este modo, un número, una simetría y una magnitud física pueden compartir
genealogía sin confundirse ni reducirse entre sí.

La tesis del continuo queda así formulada con precisión: la realidad
arquimediana es la sombra de radio cero de historias discretas infinitamente
refinables; aquello que la sombra omite reaparece como memoria, incidencia,
geometría y estructura física.

% RESULTADO_RECUPERADO. Owner íntegro de la dinámica del cociente, el retorno
% observable y los cociclos de modo y orientación. Se subordina a la única
% construcción APP--TRIT--TPK sin abrir un origen alternativo.
\begin{HMTScientificOwnerSubordinated}
\chapter{Extensiones supervivientes, monodromía y dinámica del cociente}
\label{chap:extensiones-dinamica-cociente}
\index{supervivencia!extensión global}
\index{monodromía!retorno nonádico}
\index{medida!catálogo de emisiones}

El lenguaje de los caminos necesita una precisión decisiva. Un camino de
APP, una emisión visible, una historia enriquecida de TPK y una sección
global no son cuatro nombres para el mismo objeto. Constituyen cuatro niveles
consecutivos de una misma construcción. El camino conserva el orden
posicional; la emisión conserva una lectura finita; la historia añade firma,
acarreo, orientación, frontera y memoria; la sección global exige que todos
sus prefijos sobrevivan simultáneamente a cualquier profundidad.

Esta jerarquía permite separar dos cuestiones que habían permanecido
entremezcladas. La primera es lógica: bajo qué hipótesis un prefijo finito
pertenece de verdad a una extensión infinita compatible. La segunda es
dinámica: qué significa transportar una sola historia y qué significa
transferir simultáneamente todas sus continuaciones compatibles. La
distinción es esencial en APP\@. Una trayectoria determinada lleva un estado a
otro; una transferencia de caminos lleva una distribución de extensiones a
otra distribución.

En esta sección se resuelve el problema finito de profundidad seis. La
cronología determinista publicada posee retorno visible y no puede seleccionar
una probabilidad única. En cambio, el transporte sesgado por la fase nonádica,
cuando cada canal activo se interpreta como la renovación normalizada de toda
su fibra de continuaciones compatibles, tiene un retorno de nueve pasos
uniforme sobre las (468) emisiones. Los pesos (1/3) dejan entonces de ser
una elección exterior: proceden de las tres apariciones de cada firma
(+1,-1,0) en el ciclo nonádico orientado. La normalización dentro de cada fibra sigue
siendo una premisa matemática precisa ---el conteo uniforme de todos los
caminos compatibles--- y no debe confundirse con el sucesor de una sola
trayectoria. El paso al límite inverso de todas las profundidades conserva un
enunciado propio, formulado al final de la sección.

\section{\(5\) niveles de la noción de camino}

Sea \(\Gamma_{\rm APP}\) la categoría de caminos cardinales del toro APP\@.
Para cada profundidad \(m\), sea \(\mathcal H_m\) el conjunto finito de
historias enriquecidas admisibles
\[
 h_m=(v_0,e_1,v_1,\ldots,e_m,v_m),
\]
donde cada flecha \(e_j\) transporta todos los campos declarados por TPK\@. Si
\(n\geq m\), la aplicación
\[
 p_{n,m}:\mathcal H_n\longrightarrow\mathcal H_m
\]
elimina la cola posterior a la profundidad \(m\). La sombra posicional
\[
 \operatorname{sh}_m:\mathcal H_m\longrightarrow\Gamma_{\rm APP}
\]
olvida la fibra de memoria y conserva el camino APP subyacente.

\begin{definition}[Jerarquía de extensiones]
\label{def:jerarquia-extensiones}
Distinguimos los objetos siguientes.
\begin{enumerate}
  \item Un \emph{camino APP} es un morfismo de
  \(\Gamma_{\rm APP}\).
  \item Una \emph{ruta observable} es una lectura de ese morfismo que
  conserva conjuntamente las hojas de suma y producto.
  \item Una \emph{extensión TPK} es una elevación de la ruta observable a
  una historia \(h_m\in\mathcal H_m\), con sus datos de residuo, cociente,
  firma, orientación, frontera, acarreo y registro.
  \item Una \emph{extensión superviviente} es una historia finita que admite
  prolongaciones compatibles a todo horizonte.
  \item Una \emph{sección global} es una familia compatible
  \(x=(h_m)_{m\geq0}\) de extensiones supervivientes.
\end{enumerate}
En particular, la ruta es la sombra de la extensión; no es un objeto exterior
que la extensión elija después.
\end{definition}

Para \(h\in\mathcal H_m\), su cilindro de prolongaciones es
\[
 Z_m(h)=
 \{x=(h_n)_n\in\varprojlim_n\mathcal H_n:x_m=h\}.
\]
En particular, toda familia de \(Z_m(h)\) satisface además
\(p_{n,\ell}(h_n)=h_\ell\) para \(n\geq\ell\); no se admiten colecciones de
truncamientos incompatibles.
Definamos también
\[
 \mathcal S_m^{(N)}=p_{N,m}(\mathcal H_N),
 \qquad
 \mathcal S_m^\infty=\bigcap_{N\geq m}\mathcal S_m^{(N)}.
\]

\begin{theorem}[Criterio condicional de extensión superviviente]
\label{thm:criterio-extension-superviviente}
Fijemos \(h\in\mathcal H_m\) y, para \(N\geq m\), escribamos
\[
 \mathcal P_N(h)=\{g\in\mathcal H_N:p_{N,m}(g)=h\}.
\]
Supóngase que los truncamientos son coherentes y que cada
\(\mathcal P_N(h)\) es finito. Son equivalentes:
\begin{enumerate}
  \item \(\mathcal P_N(h)\ne\varnothing\) para todo \(N\geq m\), es decir,
  \(h\in\mathcal S_m^\infty\);
  \item el árbol de prolongaciones finitas de \(h\) es finitamente ramificado
  y tiene un nivel no vacío a toda profundidad;
  \item \(Z_m(h)\ne\varnothing\).
\end{enumerate}
En consecuencia, siempre que los truncamientos supervivientes finitos sean no
vacíos, el conjunto
\[
 \Omega_{\rm H}^{\rm surv}
 =\varprojlim_m\mathcal S_m^\infty
\]
es exactamente el espacio de secciones globales supervivientes. El teorema no
afirma que esos truncamientos sean no vacíos para toda entrada APP: esa
propiedad debe demostrarse mediante el protocolo de supervivencia.
\end{theorem}

\begin{proof}
La implicación de la tercera condición a las dos primeras se obtiene por
truncamiento. Supongamos ahora la primera condición y ordenemos los elementos de
\(\bigsqcup_{N\geq m}\mathcal P_N(h)\) por la relación de prefijo. Cada nivel
es finito y no vacío, y la coherencia de los truncamientos hace de esta unión
un árbol finitamente ramificado con nodos a cualquier profundidad; esto da la
segunda condición. Recíprocamente, desde la segunda condición el lema de König
proporciona una rama infinita. Sus truncamientos forman una familia compatible
perteneciente a \(Z_m(h)\), que es la tercera condición. La misma construcción
identifica el límite inverso indicado allí donde sus truncamientos finitos son
no vacíos.
\end{proof}

El teorema resuelve la posible discordancia entre «prolongable hasta cada
horizonte por separado» y «posee una prolongación global»: la finitud de las
ramas es la hipótesis que permite escoger una única cadena compatible. No
afirma que todos los caminos APP pertenezcan a
\(\operatorname{sh}(\Omega_{\rm H}^{\rm surv})\). Esa cobertura exige probar
que los cierres enriquecidos no vacían las fibras correspondientes.

\section{Retorno de fase y relación de monodromía}

El retorno del cociente fase--longitud se formaliza por separado en
\cref{def:odometro-9-12,prop:proyeccion-fase-longitud}:
\[
(b,r)\xrightarrow{\;9\;}(b+1,r),
\qquad
(b,r)\xrightarrow{\;108\;}(b,r).
\]
Estas igualdades pertenecen al odómetro \(9\times12\). No constituyen la
prueba de \(F_{\rm obs}^{108}=I\), que usa el estado observable y aparece más
adelante.

Sea \(\phi\) la coordenada de fase con valores en \(\mathbb Z/9\mathbb Z\).
Un retorno visible satisface
\[
 \phi(h_{k+9})=\phi(h_k).
\]
El registro completo puede, sin embargo, cumplir
\[
 B(h_{k+9})\ne B(h_k).
\]
Esta desigualdad no es un fallo de periodicidad: es la memoria de la vuelta.
Una acción ordinaria de \(C_9\) sobre el estado completo impondría que la
novena potencia fuese la identidad y borraría precisamente esa información.

Para tipar el retorno, sea \(H\) una hexada y
\(P_2(H)\) su conjunto de quince pares. Sean \(\mathcal B_k\) los estados de
registro que sobreviven en la fase \(k\), sea
\(\mathcal T_k\subseteq\mathcal B_k\times\mathcal B_{k+1}\) la relación de
transporte publicada y supongamos dado un codificador de fibra
\[
 c_k:\mathcal B_k\longrightarrow P_2(H).
\]
Sólo después de declarar esos codificadores se define la relación observable
\[
 R_k=(c_k\times c_{k+1})(\mathcal T_k)
 \subseteq P_2(H)\times P_2(H).
\]
La monodromía relacional de una vuelta es
\[
 \mathcal M_k
 =R_{k+8}\circ R_{k+7}\circ\cdots\circ R_k.
\]

\begin{proposition}[Clasificación del retorno nonádico dado el codificador]
\label{prop:clasificacion-retorno-nonadico}
Se dan tres niveles lógicamente distintos.
\begin{enumerate}
  \item Sin una prueba de funcionalidad, el retorno es la relación
  \(\mathcal M_k\).
  \item Si cada \(R_k\) es la gráfica de una aplicación
  \(\mu_k:P_2(H)\to P_2(H)\), el transporte es el producto sesgado
  \[
   T(k,p)=(k+1,\mu_k(p)),
  \]
  y
  \[
   T^9(k,p)=(k,M_k(p)),
   \qquad
   M_k=\mu_{k+8}\circ\cdots\circ\mu_k.
  \]
  \item Si las \(\mu_k\) son biyectivas, \(M_k\) es una holonomía
  automorfa. El transporte desciende a una acción ordinaria de
  \(C_9\) sobre el estado completo si y sólo si \(M_k=I\) en todas las
  fibras.
\end{enumerate}
\end{proposition}

\begin{proof}
La primera afirmación es la definición de imagen y composición relacional una
vez fijados \((c_k)_k\). Bajo
funcionalidad, nueve aplicaciones consecutivas mantienen fija la primera
coordenada y componen sus acciones sobre la segunda. La composición de
biyecciones es una biyección. Finalmente, \(T^9=I\) equivale exactamente a
\(M_k=I\) para cada fase y cada punto de la fibra.
\end{proof}

\begin{corollary}[No retorno puntual bajo transporte acumulado]
\label{pf84:E02-29}
Sean \(q:X\to V\), \(T:X\to X\) y \(x\in X\). Si
\[
 q(T^9x)=q(x),\qquad
 T^9x=\mathcal M_x\,x,\qquad
 \mathcal M_x\,x\ne x,
\]
entonces la proyección retorna y el estado no:
\[
 q(T^9x)=q(x),\qquad T^9x\ne x.
\]
La hipótesis puntual \(\mathcal M_x\,x\ne x\) es indispensable; la desigualdad
global \(\mathcal M_x\ne I\) no excluye que \(x\) sea un punto fijo.
\end{corollary}

\begin{proof}
La primera igualdad es una hipótesis. La segunda y la separación puntual dan
\(T^9x=\mathcal M_x\,x\ne x\). El enunciado quedaría refutado por un caso que
satisfaga las tres premisas y tuviera \(T^9x=x\); no autoriza a sustituir la
tercera premisa por \(\mathcal M_x\ne I\).
\end{proof}

\begin{theorem}[Suficiencia funcional de la tupla enriquecida]
\label{pf84:E02-30}
Sea
\[
 E:\mathscr X_{\rm TPK}^{\rm enr}\longrightarrow\mathscr S_E
\]
la tupla que conserva bloque visible, bloques nuevos y acarreo, estado de
elevación, cilindros, firma de frontera, orientación, observable, fase y
registro dodecafásico. Sean
\[
 T:\mathscr X_{\rm TPK}^{\rm enr}\to
 \mathscr X_{\rm TPK}^{\rm enr},
 \qquad
 O:\mathscr X_{\rm TPK}^{\rm enr}\to Y.
\]
Cuando existen mapas
\[
 \widehat T:\mathscr S_E\to\mathscr X_{\rm TPK}^{\rm enr},
 \qquad
 \widehat O:\mathscr S_E\to Y
\]
tales que
\[
 T=\widehat T\circ E,\qquad O=\widehat O\circ E,
\]
la igualdad \(E(x)=E(y)\) implica \(T(x)=T(y)\), \(O(x)=O(y)\) y, por
determinismo, las mismas continuaciones y salidas posteriores.
\end{theorem}

\begin{proof}
Por factorización,
\[
 T(x)=\widehat T(E(x))=\widehat T(E(y))=T(y),
 \qquad
 O(x)=\widehat O(E(x))=\widehat O(E(y))=O(y).
\]
La primera igualdad permite iterar el argumento. Ésta es una condición
suficiente, no una equivalencia no probada. Dos estados con la misma tupla y
transiciones o salidas diferentes, el acceso a una coordenada ausente o la
consulta de un valor objetivo son falsadores. La microfibra visible de
cardinal cinco del \cref{pf84:E02-22} sólo controla la no fidelidad de la
palabra \(w_6\); no demuestra esta factorización positiva.
\end{proof}

La ventana finita publicada contiene nueve retornos de fase con cambio del
estado holonómico integral y una selección local de tres continuaciones a una al
aumentar dos horizontes. Eso prueba retorno con memoria y refuta la
identificación del registro con un ciclo sin memoria. La extracción global de
los nueve codificadores \(c_k\) y, a través de ellos, de las relaciones
\(R_k\), desde todas las historias supervivientes es una operación distinta.
La proposición anterior clasifica el retorno una vez dado el mapa
registro--pares; no deriva ese mapa desde el estado holonómico integral de TPK.

El episodio concreto de la novena fase, con sus tres continuaciones, el
bloque \(B_{11}\), la frontera \((502,662,638)\) y las desigualdades enteras
de cilindros, se demuestra una sola vez en el
\cref{thm:seleccion-novena-fase}. En esta residencia sólo se usa su
consecuencia tipada:
\[
\#\operatorname{Surv}_{1}(B_9)=3,\qquad
\#\operatorname{Surv}_{2}(B_9)=1.
\]

Las otras dos palabras pueden preceder la misma cola ternaria sólo al cambiar
la extensión profunda de Hensel. Comparten una parte del residuo visible,
pero no la genealogía del estado.

\section{El catálogo de 468 emisiones como producto de canales}

El catálogo finito parte del conjunto
\(\Omega_{\rm seed}\) de \(104\,976\) presentaciones microscópicas. Sea
\(\mathcal U_6\) el conjunto de emisiones séxtuples distintas y sea
\[
 F:\Omega_{\rm seed}\twoheadrightarrow\mathcal U_6.
\]
El censo exacto es
\[
 |\mathcal U_6|=468,
 \qquad
 \#\{U:|F^{-1}(U)|=144,432,1944\}=432,18,18.
\]
Estos conteos pertenecen a \(F\). Compararlos con historias TPK de
profundidad seis requeriría un mapa explícito desde el dominio de historias
hacia \(\Omega_{\rm seed}\); ningún resultado de esta sección presupone ese
mapa ni identifica presentaciones del catálogo con estados holonómicos integrales.
El catálogo finito proporciona, para cada emisión, una firma aditiva
\(s\) y una firma multiplicativa \(e\). La aplicación conjunta es una
biyección
\[
 \mathcal U_6\xrightarrow{\sim}
 \mathcal S_+\times\mathcal S_\times,
 \qquad
 |\mathcal S_+|=18,
 \quad
 |\mathcal S_\times|=26.
\]
Así,
\[
 \boxed{468=18\cdot26}
\]
es una factorización del objeto y no sólo de su cardinal. La primera
coordenada conserva cómo suma la emisión; la segunda, cómo acumula el
producto. Esta separación es la forma finita explícita de la
incompatibilidad suma--producto: ambas lecturas coexisten en el estado, pero
cada proyección olvida la firma complementaria.

La primera coordenada admite una parametrización adicional que será útil al
compararla con la frontera nonádica. Si una semilla aditiva tiene coordenadas
\((i,j)\in(\mathbb Z/9\mathbb Z)^2\) y orientación
\(\varepsilon=+1\) para las direcciones este--sur y \(-1\) para
norte--oeste, su firma completa queda determinada biyectivamente por
\[
 \bigl(i+j\bmod9,\varepsilon\bigr)
 \in(\mathbb Z/9\mathbb Z)\times\{\pm1\}.
\]
Por tanto,
\[
 \boxed{\mathcal S_+\simeq
 (\mathbb Z/9\mathbb Z)\times\{\pm1\}.}
\]
La coordenada producto no posee una uniformidad análoga: sus \(26\) firmas
tienen fibras de semillas de tamaños \(8\), \(24\) y \(108\), con histograma
\(24\times8+1\times24+1\times108=324\). Esta asimetría es estructural. El
factor \(18\) ya contiene fase y orientación; cualquier descenso posterior
del factor producto a una clase de pares debe conservar memoria adicional o
aceptar fibras no uniformes.

Definamos sobre \(\ell^2(\mathcal S_+\times\mathcal S_\times)\) las
esperanzas condicionales
\[
\begin{aligned}
 (E_+f)(s,e)&=\frac1{18}\sum_{s'\in\mathcal S_+}f(s',e),\\
 (E_\times f)(s,e)&=\frac1{26}\sum_{e'\in\mathcal S_\times}f(s,e').
\end{aligned}
\]
El primer operador renueva la firma aditiva y conserva la multiplicativa;
el segundo realiza la operación dual. El tercer canal retiene ambas.

\section{La trayectoria única y la transferencia de caminos}

A diferencia del retorno del odómetro definido en
\cref{def:odometro-9-12,prop:proyeccion-fase-longitud}, la proposición
siguiente demuestra el retorno del operador observable \(F_{\rm obs}\).

La ley observable publicada actualiza un solo cursor en cada paso. En las
fases \(+1\) avanza el cursor aditivo; en las fases \(-1\), el
multiplicativo; en las fases \(0\), ambos permanecen quietos. Después de los
pasos \(27,54,81,108\) se invierten las dos orientaciones. Denotemos esta
evolución determinista por \(F_{\rm obs}\).

\begin{proposition}[Retorno de la cronología observable]
\label{prop:retorno-fobs-identidad}
Para toda elección de las dos posiciones y orientaciones iniciales, las
posiciones y las orientaciones se restauran después de \(54\) pasos. Al
restaurar además la fase se obtiene
\[
 F_{\rm obs}^{108}=I.
\]
Por tanto, el retorno de Poincaré de \(108\) pasos induce la identidad sobre
\(\mathcal U_6\) y no selecciona una medida estacionaria única.
\end{proposition}

\begin{proof}
En cada bloque de \(27\) pasos, cada cursor realiza exactamente nueve
avances en una orientación fija. Como el espacio de posiciones es
\(X_9=(\mathbb Z/9\mathbb Z)^2\), su desplazamiento es \(9d=0\). Al final
del bloque se invierte \(d\). Dos bloques restauran posición y orientación;
cuatro restauran también el índice en \(\Theta_{108}\). El retorno inducido
es, por consiguiente, la identidad. Toda probabilidad es estacionaria para
la identidad, de modo que ninguna queda seleccionada por esa cronología.
\end{proof}

El problema aparece incluso antes del retorno si se intenta olvidar la
memoria de ruta. La firma producto
\[
 e_0=(5,5,5,5,5,5)
\]
tiene veinticuatro representantes. Tras un avance en su propia orientación,
esos representantes se separan, ocho a ocho, entre
\[
 (5,5,5,7,1,7),\qquad
 (5,5,5,7,7,1),\qquad
 (5,5,5,1,7,7).
\]
Así, el avance producto no desciende a una aplicación de las veintiséis
firmas: el representante que se había olvidado vuelve a ser relevante. Este
testigo explica por qué una ruta individual no puede convertirse, por simple
proyección, en el refresco de una fibra completa.

La lectura APP es otra. En lugar de escoger un representante, conserva todas
las continuaciones compatibles. Para formalizarla en el catálogo finito,
consideremos las subálgebras de funciones que sólo dependen de una de las dos
firmas. Respecto del conteo uniforme, \(E_+\) y \(E_\times\) son las
esperanzas condicionales que olvidan exactamente la coordenada activa. Son los
únicos proyectores de Markov que cumplen simultáneamente:
\begin{enumerate}
  \item preservación de la coordenada inactiva;
  \item olvido completo de la coordenada activa;
  \item preservación del estado de conteo uniforme.
\end{enumerate}
En efecto, una fila que ha olvidado la coordenada de partida debe ser
constante sobre los \(18\), respectivamente \(26\), destinos compatibles; la
normalización fuerza los valores \(1/18\) y \(1/26\). Ésta es la canonicidad
relativa de la renovación: no se postula una preferencia entre ramas que APP
declara igualmente posibles.

Sea \(C_9=\mathbb Z/9\mathbb Z\), con la palabra de fase
\[
 \tau=(+1,-1,0,+1,-1,0,+1,-1,0).
\]
Escribamos
\[
 P_{+1}=E_+,
 \qquad P_{-1}=E_\times,
 \qquad P_0=I.
\]

\begin{definition}[Transferencia TPK sesgada por la fase]
\label{def:transferencia-tpk-fase}
Sobre \(\mathcal U_6\times C_9\) definimos
\[
 \mathcal K_{\rm ph}\bigl((u,r),(v,r+1)\bigr)
 =P_{\tau(r)}(u,v),
\]
y hacemos cero las demás entradas. Esta definición es relativa al lector de
todos los caminos compatibles: una fase activa renueva uniformemente la fibra
que esa fase lee y la fase neutra retiene la emisión.
\end{definition}

\begin{theorem}[Retorno nonádico uniforme]
\label{thm:retorno-nonadico-uniforme}
El núcleo \(\mathcal K_{\rm ph}\) es doblemente estocástico e irreducible,
tiene período exactamente nueve y posee un único estado estacionario,
\[
 \pi_{\rm ph}(u,r)=\frac1{9\cdot468}.
\]
En cada fibra de fase, su retorno de nueve pasos es
\[
 \mathcal K_{\rm ph}^{\,9}\big|_r=E_+E_\times,
 \qquad
 (E_+E_\times)(u,v)=\frac1{468}.
\]
Si la fase inicial se distribuye uniformemente, la sombra de un paso sobre
las emisiones es
\[
 \overline K_{\rm ph}=\frac13(I+E_++E_\times).
\]
\end{theorem}

\begin{proof}
Cada \(P_{\tau(r)}\) es doblemente estocástico y el desplazamiento
\(r\mapsto r+1\) es una permutación; su producto sesgado conserva ambas sumas.
Los tres operadores son idempotentes y conmutan. Como cada uno aparece tres
veces en la palabra nonádica,
\[
 \prod_{r\in C_9}P_{\tau(r)}
 =E_+^3E_\times^3I^3=E_+E_\times.
\]
La primera esperanza olvida la firma aditiva y la segunda la multiplicativa,
de modo que su composición asigna a cada uno de los \(18\cdot26=468\)
destinos la probabilidad \(1/468\). Por ello, en nueve pasos se comunican dos
estados cualesquiera de una misma fase; los pasos restantes comunican las
fases, y el núcleo es irreducible. Todo retorno debe tener longitud múltiplo
de nueve, mientras que la diagonal de \(\mathcal K_{\rm ph}^9\) es positiva;
el período es nueve. La doble estocasticidad y la irreducibilidad dan el
estado uniforme único. Finalmente, promediar las nueve fases cuenta tres
copias de cada canal y produce la última fórmula.
\end{proof}

Olvidar la fase no es un cociente de Markov fuerte: desde una misma emisión,
dos fases distintas aplican \(I\), \(E_+\) o \(E_\times\), que tienen filas
distintas. El operador \(\overline K_{\rm ph}\) es una sombra promediada en la
fase estacionaria, no el sucesor cronológico de una emisión sin fase. Esta
distinción conserva a la vez la dinámica y el significado de la medida.

\begin{definition}[Núcleo promediado inducido por la transferencia]
\label{def:nucleo-tritico-refresco}
Denominamos núcleo promediado del catálogo a la sombra de fase
\[
 \boxed{
 K_{\rm cat}=\frac13(I+E_++E_\times).
 }
\]
Los pesos son inducidos por el calendario TPK, que contiene tres ocurrencias
de cada signo. Los operadores de renovación son inducidos relativamente a la
normalización uniforme de todas las continuaciones compatibles de la fibra
activa. Ninguna de las dos afirmaciones identifica \(K_{\rm cat}\) con
\(F_{\rm obs}\).
\end{definition}

\begin{proposition}[Familia convexa de núcleos estocásticos con medida uniforme común]
\label{prop:familia-refrescos-abc}
Para \(a,b,c\geq0\), \(a+b+c=1\), la familia
\[
 K_{a,b,c}=aI+bE_++cE_\times
\]
está formada por núcleos simétricos y doblemente estocásticos. Si \(b,c>0\),
son irreducibles y, por tener diagonal positiva, aperiódicos. Su espectro es
\[
 \operatorname{Spec}(K_{a,b,c})=
 \{1^{[1]},(a+c)^{[17]},(a+b)^{[25]},a^{[425]}\}.
\]
En particular, existe un continuo de dinámicas con la misma factorización y la
misma medida uniforme. La elección
\(K_{\rm cat}=K_{1/3,1/3,1/3}\) no queda caracterizada por el censo solo. La
seleccionan conjuntamente la palabra de fase TPK y la transferencia
normalizada de la \cref{def:transferencia-tpk-fase}.
\end{proposition}

\begin{proof}
Los tres sumandos son núcleos simétricos y estocásticos; por convexidad ocurre
lo mismo con \(K_{a,b,c}\), y la simetría convierte la estocasticidad por filas
en estocasticidad por columnas. Si \(b,c>0\), dos pasos permiten renovar
sucesivamente ambas coordenadas y comunicar cualquier par de estados. Además,
\[
 K_{a,b,c}((s,e),(s,e))=a+\frac b{18}+\frac c{26}>0,
\]
de modo que el núcleo es aperiódico. En la descomposición ortogonal
\[
 \mathbf1\oplus\ell_0^2(\mathcal S_+)
 \oplus\ell_0^2(\mathcal S_\times)
 \oplus\bigl(\ell_0^2(\mathcal S_+)\otimes
              \ell_0^2(\mathcal S_\times)\bigr),
\]
los pares \((E_+,E_\times)\) actúan, respectivamente, como
\((1,1),(0,1),(1,0),(0,0)\). Se obtienen así los cuatro autovalores y las
multiplicidades \(1,17,25,425\). Para \(b,c>0\), la irreducibilidad hace única
la medida estacionaria uniforme, lo que exhibe la no unicidad incluso bajo
esa exigencia.
\end{proof}

\begin{theorem}[Espectro exacto del refresco equiponderado]
\label{thm:dinamica-cociente-468}
El núcleo \(K_{\rm cat}\) es simétrico, doblemente estocástico,
irreducible y aperiódico. Su espectro, con multiplicidades, es
\[
 \boxed{
 \operatorname{Spec}(K_{\rm cat})
 =\left\{1^{[1]},
 \left(\frac23\right)^{[42]},
 \left(\frac13\right)^{[425]}\right\}.
 }
\]
En consecuencia, su brecha espectral es (1/3) y su único estado
estacionario es
\[
 \boxed{
 \tau_{468}(f)=\frac1{468}\sum_{U\in\mathcal U_6}f(U).
 }
\]
\end{theorem}

\begin{proof}
Los operadores \(E_+\) y \(E_\times\) son proyectores ortogonales de
promedio y conmutan. Las filas y columnas de cada uno suman uno; lo mismo
ocurre con su promedio. Desde \((s,e)\) se alcanza \((s',e')\) en dos
renovaciones, y la contribución de \(I\) proporciona una autoarista. Esto
prueba irreducibilidad y aperiodicidad.

El espacio se descompone ortogonalmente en cuatro sectores:
\[
 \mathbf1,
 \qquad
 \ell_0^2(\mathcal S_+),
 \qquad
 \ell_0^2(\mathcal S_\times),
 \qquad
 \ell_0^2(\mathcal S_+)\otimes
 \ell_0^2(\mathcal S_\times).
\]
Sus dimensiones son \(1,17,25,17\cdot25\). En el sector constante actúan
los tres sumandos como la identidad; en cada sector de una sola coordenada
actúan como (I+I+0); en el sector de interacción, como (I+0+0).
Dividir por tres da los autovalores y multiplicidades anunciados. Un núcleo
finito irreducible y aperiódico posee un único estado estacionario; como es
doblemente estocástico, éste es uniforme.
\end{proof}

El teorema es exacto relativamente a la factorización
\(\mathcal U_6\simeq\mathcal S_+\times\mathcal S_\times\), al calendario
TPK y al conteo uniforme de las continuaciones compatibles de cada fibra. La
familia \(K_{a,b,c}\) demuestra que la dinámica no se deduce del censo por sí
solo. La \cref{prop:retorno-fobs-identidad} distingue la transferencia de
todos los caminos de la cronología de uno de ellos.

\section{Elevación compensada construida desde el cociente}

Escribamos \(m(U)=|F^{-1}(U)|\). A partir del núcleo ya elegido en el cociente
definimos el núcleo de presentaciones por
\[
 \boxed{
 K_{\rm seed}(x,y)
 =\frac{K_{\rm cat}(F(x),F(y))}{m(F(y))}
 }
\]
Este núcleo es estocástico. Al sumar sobre una fibra de destino se obtiene
\[
 \sum_{F(y)=V}K_{\rm seed}(x,y)
 =K_{\rm cat}(F(x),V),
\]
independientemente del representante \(x\). Esta identidad muestra que el
elevación construida proyecta exactamente sobre \(K_{\rm cat}\). No demuestra que
una dinámica microscópica preexistente del estado holonómico integral de TPK descienda
al catálogo.

\begin{theorem}[Elevación reversible construida desde el cociente]
\label{thm:elevacion-reversible-cociente}
La probabilidad
\[
 \widetilde\tau(x)
 =\frac1{468\,m(F(x))}
\]
es estacionaria y reversible para \(K_{\rm seed}\), y
\(F_\#\widetilde\tau=\tau_{468}\). Si
\[
 J_F\delta_U
 =m(U)^{-1/2}\sum_{F(x)=U}\delta_x,
 \qquad
 D=\operatorname{diag}(m(U)),
\]
entonces
\[
 K_{\rm seed}J_F
 =J_F\bigl(D^{1/2}K_{\rm cat}D^{-1/2}\bigr).
\]
En particular, el entrelazamiento hilbertiano sin conjugación por \(D\)
ocurre si y sólo si \(K_{\rm cat}\) conmuta con las multiplicidades de las
fibras.
\end{theorem}

\begin{proof}
La normalización se sigue de que cada una de las 468 fibras recibe masa
total (1/468). Usando la simetría de \(K_{\rm cat}\), para
\(F(x)=U\) y \(F(y)=V\) se tiene
\[
 \widetilde\tau(x)K_{\rm seed}(x,y)
 =\frac{K_{\rm cat}(U,V)}{468m(U)m(V)}
 =\widetilde\tau(y)K_{\rm seed}(y,x).
\]
Esto demuestra balance detallado, estacionariedad y la identidad de la
medida imagen.
Aplicar ambos miembros de la última identidad a un vector base
\(\delta_V\) da, sobre la fibra \(U\), el coeficiente
\(K_{\rm cat}(U,V)/\sqrt{m(V)}\); ése es exactamente el coeficiente del
operador conjugado mostrado.
\end{proof}

\begin{theorem}[Elevación de fase fiel a la quietud]
\label{thm:elevacion-fase-semillas}
Sea \(X=\Omega_{\rm seed}\), \(m(u)=|F^{-1}(u)|\), y definamos
\[
\begin{aligned}
 L_0(x,y)&=\mathbf1_{\{x=y\}},\\
 L_+(x,y)&=\frac{E_+(F(x),F(y))}{m(F(y))},\\
 L_\times(x,y)&=\frac{E_\times(F(x),F(y))}{m(F(y))}.
\end{aligned}
\]
Sobre \(X\times C_9\), sea
\[
 \widehat{\mathcal K}_{\rm ph}\bigl((x,r),(y,r+1)\bigr)
 =L_{\tau(r)}(x,y).
\]
Entonces \((F,\mathrm{id}_{C_9})\) satisface la propiedad de
\emph{agregabilidad fuerte} (\emph{strong lumpability}) de
\(\widehat{\mathcal K}_{\rm ph}\) hacia \(\mathcal K_{\rm ph}\). El núcleo
microscópico es irreducible, tiene período nueve y su estado estacionario
único es
\[
 \widetilde\pi_{\rm ph}(x,r)
 =\frac1{9\cdot468\,m(F(x))}.
\]
Su retorno nonádico satisface
\[
 \widehat{\mathcal K}_{\rm ph}^{\,9}
 \bigl((x,r),(y,r)\bigr)
 =\frac1{468\,m(F(y))}.
\]
\end{theorem}

\begin{proof}
La suma de \(L_+\), respectivamente \(L_\times\), sobre una fibra de destino
es \(E_+\), respectivamente \(E_\times\), e \(L_0\) proyecta a la identidad.
Esto prueba la agregabilidad en cada fase. Los factores \(m(F(y))\) cancelan
al componer: primero se selecciona uniformemente la firma activa y después un
representante de la emisión obtenida. Así, \(L_+\) y \(L_\times\) son
proyectores conmutantes y
\[
 L_+L_\times(x,y)=\frac1{468\,m(F(y))}.
\]
La positividad de esta entrada comunica todos los representantes en nueve
pasos. La fase fuerza que las longitudes de retorno sean múltiplos de nueve,
y la entrada diagonal del retorno es positiva. La fórmula estacionaria se
comprueba sumando primero los \(m(u)\) representantes de cada emisión; cada
una de las \(9\cdot468\) clases fase--emisión recibe la misma masa.
\end{proof}

Esta elevación mejora la lectura del núcleo promediado
\(K_{\rm seed}\): en la fase neutra conserva el representante microscópico,
en vez de remezclar silenciosamente su fibra. Al promediar y olvidar la fase,
ambas construcciones tienen la misma sombra \(K_{\rm cat}\), pero no la misma
dinámica oculta. La diferencia es necesaria para respetar que el cero de TRIT
es memoria de quietud.

\section{Congruencias mínimas estables bajo avance, media vuelta y 4.º de giro: \(28\) y \(162\) clases de memoria}

El testigo \(e_0=(5,5,5,5,5,5)\) muestra que la firma producto es demasiado
gruesa para transportar una ruta. Es posible calcular exactamente la memoria
mínima que falta. Sea
\[
 \mathcal X_\times=X_9\times\mathscr D_{\rm card}
\]
el espacio de \(324\) semillas producto. Denotemos por \(P\) el avance de una
casilla en la orientación almacenada, por \(H\) la media vuelta de esa
orientación y por \(Q\) el cuarto de giro
\[
 N\longmapsto E\longmapsto S\longmapsto O\longmapsto N.
\]

\begin{theorem}[Congruencias de avance y giro]
\label{thm:congruencias-avance-giro}
La congruencia mínima que refina las veintiséis firmas producto y es estable
por \(P\) y \(H\) posee \(28\) clases: veintisiete tienen tamaño ocho y una
tiene tamaño \(108\). La única escisión nueva es
\[
 (5,5,5,5,5,5)\longrightarrow
 \mathcal L_0\sqcup\mathcal L_1\sqcup\mathcal L_2,
 \qquad |\mathcal L_j|=8.
\]
Sus órbitas bajo \(P,H\) tienen tamaños
\[
 9,\quad9,\quad9,\quad1.
\]
Al combinar este refinamiento con las dieciocho firmas aditivas se obtiene
\[
 18\cdot28=504=468+36.
\]

La congruencia mínima estable por \(P\) y \(Q\) posee \(162\) clases de dos
semillas. Cada clase es una órbita de la involución
\[
 J(i,j,d)=\bigl(7-i,\,7-j,\,\overline d\bigr)\pmod9.
\]
Los operadores inducidos por \(P,Q\) actúan transitivamente sobre las
\(162\) clases.
\end{theorem}

\begin{proof}
Se parte de la partición por las veintiséis firmas y se refina
iterativamente una clase siempre que dos de sus elementos tengan imágenes en
clases distintas por alguno de los generadores. El proceso termina porque
\(\mathcal X_\times\) es finito. El censo exhaustivo produce, para
\((P,H)\), el histograma \(27\times8+1\times108=324\), las cuatro órbitas
indicadas y tres clases sobre la fibra excepcional de tamaño veinticuatro.
Para \((P,Q)\) produce \(162\times2=324\). La comprobación directa demuestra
que \(J\) preserva cada firma, conmuta con ambos generadores y empareja
exactamente los dos elementos de cada clase; el grafo de clases generado por
\(P,Q\) es conexo.
\end{proof}

El exceso \(36\) y el objeto \(R_{36}\) tienen tipos distintos. La frontera
de profundidad treinta y seis es
\[
 R_{36}=
 \begin{pmatrix}
 222220\\021222\\102011
 \end{pmatrix},
\]
es decir, tres bloques ordenados de seis trits, uno por canal. No es un
conjunto de treinta y seis estados. Por su parte, el término
\(504-468=36\) cuenta memoria adicional de hoja: la media vuelta que hace
funcional el transporte producto añade exactamente dos hojas nuevas a una
sola firma y, al repetirlas sobre las dieciocho firmas aditivas, produce
\(18(3-1)=36\). La coincidencia del entero (36) expresa dos graduaciones
del transporte, profundidad de frontera y multiplicidad de hoja, pero no
autoriza a identificar sus objetos.

El cuarto de giro \(Q\) tampoco estaba escondido en \(F_{\rm obs}\). La
cronología publicada contiene avances y medias vueltas programadas; incorporar
\(Q\) es una ampliación explícita del alfabeto de rutas. Una vez incorporado,
un paseo simétrico perezoso en \(P^{\pm1},Q^{\pm1}\) es irreducible y
aperiódico sobre las \(162\) clases, y su medida uniforme es única. Ésta es la
forma matemática exacta de «contar los giros»: el giro es una coordenada de
transporte y no una corrección decimal posterior.

\section{Cociclos de modo y orientación}
\label{sec:cociclos-modo-orientacion}

La misma regla permite definir los contadores operativos sin introducir la
Mecánica Dimensional. Añadamos al estado el último modo activo
\(m\in\{\Sigma,\Pi\}\). Una firma \(+1\) fija \(m=\Sigma\), una firma
\(-1\) fija \(m=\Pi\), y una firma \(0\) conserva el modo anterior. Para una
flecha \(a\), definamos el cociclo firmado
\[
 \omega_{\rm sw}(a)=
 \begin{cases}
  +1,&\Pi\longrightarrow\Sigma,\\
  -1,&\Sigma\longrightarrow\Pi,\\
  0,&\text{el modo se conserva}.
 \end{cases}
\]
Su variación positiva es \(s(a)=|\omega_{\rm sw}(a)|\). Definamos además
\(g(a)=1\) cuando cambia el par ordenado de orientaciones cardinales y
\(g(a)=0\) cuando se conserva. Al mantener en los extremos el modo y la
orientación terminal, las sumas
\[
 \Omega_{\rm sw}(\gamma)=\sum_{a\subset\gamma}\omega_{\rm sw}(a),
 \qquad
 S(\gamma)=\sum_{a\subset\gamma}s(a),
 \qquad
 G(\gamma)=\sum_{a\subset\gamma}g(a)
\]
son aditivas por concatenación.

En una vuelta nonádica cerrada, la sucesión de modos es
\[
 (\Sigma,\Pi,\Pi,\Sigma,\Pi,\Pi,\Sigma,\Pi,\Pi).
\]
Por tanto,
\[
 \Omega_{\rm sw}(C_9)=0,\qquad S(C_9)=6:
\]
hay tres cambios \(\Pi\to\Sigma\), tres cambios
\(\Sigma\to\Pi\) y tres retenciones. En el ciclo de \(108\) pasos,
\[
 \#(+1)=\#(-1)=\#(0)=36,\qquad
 \Omega_{\rm sw}=0,\qquad S=72,\qquad G=4.
\]
El último número cuenta las cuatro inversiones simultáneas programadas en
\(27,54,81,108\). Estos cociclos proceden directamente del calendario TPK\@.

\subsection{Cociente de modos y correspondencia areal--volumétrica}
\label{sec:cociente-modos-fav}

El signo del calendario, la etiqueta de modo y la hoja geométrica son tres
objetos tipados. En particular, no se identifica el signo de fase con el
observable angular sobre dígitos. Para la correspondencia geométrica fijamos el
lector HMT tipado
\[
 \mathfrak g_{\rm av}(\Sigma)=\mathscr H_{\rm ar},
 \qquad
 \mathfrak g_{\rm av}(\Pi)=\mathscr H_{\rm vol}.
\]
La primera asignación expresa la lectura aditiva o de borde; la segunda, la
lectura multiplicativa o de interior. El mapa \(\mathfrak g_{\rm av}\) es
parte del tipado de hojas. El resultado siguiente demuestra que, una vez
fijado ese tipado, la dinámica TPK fuerza la incidencia entre ellas.

\begin{theorem}[Descenso necesario de la incidencia de hojas]
\label{thm:descenso-necesario-fav}
Sea \(m_r\) el modo obtenido después de aplicar la fase \(r\) del bloque
primitivo \((+1,-1,0)\), con la regla
\[
 +1:\ m\mapsto\Sigma,\qquad
 -1:\ m\mapsto\Pi,\qquad
 0:\ m\mapsto m.
\]
La órbita cíclica de modos es \((\Sigma,\Pi,\Pi)\). Si
\[
 N_3(i,j)
 :=
 \#\{r\in\mathbb Z/3\mathbb Z:(m_r,m_{r+1})=(i,j)\},
\]
entonces, en el orden \((\Sigma,\Pi)\),
\[
 \boxed{
 N_3=
 \begin{pmatrix}0&1\\1&1\end{pmatrix}.}
\]
Por repetición del calendario,
\[
 N_9=3N_3,\qquad N_{27}=9N_3,\qquad N_{108}=36N_3.
\]
Después de aplicar \(\mathfrak g_{\rm av}\), la matriz primitiva de
incidencia es, en el orden
\((\mathscr H_{\rm ar},\mathscr H_{\rm vol})\),
\[
 \boxed{
 F_{\rm av}=
 \begin{pmatrix}0&1\\1&1\end{pmatrix}.}
\]
Así, las dos transiciones cruzadas, la única autorretención volumétrica, la
ausencia de autorretención areal y las multiplicidades primitivas unitarias
son consecuencias del cociente de modos de TPK; no son cuatro postulados
añadidos.
\end{theorem}

\begin{proof}
Los tres pares consecutivos, con cierre cíclico, son
\[
 \Sigma\longrightarrow\Pi,\qquad
 \Pi\longrightarrow\Pi,\qquad
 \Pi\longrightarrow\Sigma.
\]
Cada uno aparece una vez y el par
\(\Sigma\to\Sigma\) no aparece. Esto determina las cuatro entradas de
\(N_3\). Los calendarios de longitudes \(9,27,108\) contienen,
respectivamente, \(3,9,36\) copias del bloque primitivo, por lo que sus
matrices de incidencia son los múltiplos indicados. El máximo común divisor
de las entradas no nulas de cada matriz larga elimina únicamente esa
repetición global y devuelve \(N_3\). Finalmente,
\(\mathfrak g_{\rm av}\) transporta \(\Sigma,\Pi\) a las hojas areal y
volumétrica sin alterar la incidencia.
\end{proof}

La realización hiperbólica real que utiliza esta matriz y prueba su
invariancia lorentziana se formula en
\cref{sec:tres-generadores-concretos-trit}; allí \(F_{\rm av}\) se usa
como antecedente, sin repetir su derivación.

Este teorema debe distinguirse de una afirmación de Markov más fuerte.
Olvidar la fase no es un cociente de Markov fuerte de
\(\mathcal K_{\rm ph}\): el modo visible no determina cuál de los tres
operadores \(I,E_+,E_\times\) actúa en la iteración elemental siguiente. Lo que desciende
sin hipótesis adicional es el multigrafo de aristas de un bloque TPK\@.

Existe, sin embargo, una completación de memoria uno determinada de forma
única por ese descenso. Sea \(X_{\rm av}\) el menor subshift de tipo finito
sobre
\(\{\mathscr H_{\rm ar},\mathscr H_{\rm vol}\}\) que contiene todas las
palabras obtenidas al olvidar la fase y que sólo conserva pares consecutivos.
Todo subshift de un paso con esa propiedad debe admitir exactamente los tres
pares observados; el menor de ellos excluye el único par no observado. Por
tanto su matriz de adyacencia es \(F_{\rm av}\).

\begin{corollary}[Medida de Parry del envolvente de memoria uno]
\label{cor:parry-envolvente-fav}
El subshift \(X_{\rm av}\) es primitivo, tiene entropía topológica
\(\log\varphi\) y posee una única medida de máxima entropía. Sus pesos
estacionarios de vértice satisfacen
\[
 \boxed{
 \frac{\mu_{\rm ar}}{\mu_{\rm vol}}=\varphi^{-2}.}
\]
\end{corollary}

\begin{proof}
La raíz de Perron de \(F_{\rm av}\) es \(\varphi\), y un vector de Perron
positivo es \((1,\varphi)\). Como la matriz es simétrica, los pesos de vértice
de Parry son proporcionales al producto de los vectores izquierdo y derecho,
es decir, a \((1,\varphi^2)\).
\end{proof}

La medida de Parry pertenece al envolvente de caminos de memoria uno. No es
la imagen de \(\tau_{468}\), ni la frecuencia temporal de la órbita periódica
de modos. Esta última es
\[
 \nu_{\rm cyc}(\mathscr H_{\rm ar})=\frac13,\qquad
 \nu_{\rm cyc}(\mathscr H_{\rm vol})=\frac23,
\]
mientras \(\tau_{468}\) es uniforme sobre emisiones y la medida de Parry
pondera historias de renovación de longitud arbitraria. Las tres medidas
responden a preguntas distintas. La razón irracional
\(\varphi^{-2}\) queda así derivada canónicamente del envolvente de caminos,
sin atribuirla a un cociente de cardinales finitos.

\subsubsection{Del conteo cronológico al retorno marcado por \texorpdfstring{\(\pi\)}{π}}
\label{sec:retorno-marcado-pi-phi2}

Las dos fórmulas
\[
 \left(\frac13,\frac23\right)
 \qquad\text{y}\qquad
 \frac{\mu_{\rm ar}}{\mu_{\rm vol}}=\varphi^{-2}
\]
no son aproximaciones una de otra. La primera cuenta los tres instantes del
calendario primitivo \((\Sigma,\Pi,\Pi)\); la segunda pondera todas las
historias admisibles del envolvente de memoria uno. Por eso la frecuencia
cronológica es racional y la razón de Parry es irracional. En lo que sigue
escribimos
\[
 \frac{\pi}{\varphi^2}
 =\pi\varphi^{-2};
\]
no se trata de dividir por \(\varphi^{-2}\).

La aparición de \(\varphi^{-2}\) puede verse directamente en la dinámica
lineal. Si \(F_n\) denota el \(n\)-ésimo número de Fibonacci,
\[
 F_{\rm av}^{\,n}
 =
 \begin{pmatrix}
  F_{n-1}&F_n\\
  F_n&F_{n+1}
 \end{pmatrix}
 \qquad(n\geq1).
\]
Los autovalores de \(F_{\rm av}\) son
\(\varphi\) y \(-\varphi^{-1}\). La rama contraída invierte orientación en
una iteración. Al pasar a la memoria cuadrática, es decir, a
\(\operatorname{Sym}^2(\mathbb R^2)\), esa inversión se registra antes de
desaparecer el signo. En la base \((x^2,2xy,y^2)\), representando por
filas las imágenes de los vectores básicos, el operador inducido es
\[
 \operatorname{Sym}^2(F_{\rm av})=
 \begin{pmatrix}
  0&0&1\\
  0&1&2\\
  1&1&1
 \end{pmatrix},
 \qquad
 \chi(t)=(t+1)(t^2-3t+1).
\]
Sus tres autovalores son
\[
 \varphi^2,\qquad -1,\qquad \varphi^{-2}.
\]
Así, \(\varphi^{-2}\) es el modo estable positivo de la memoria de segundo
orden. Este paso explica por qué la razón áurea aparece aquí como ley de
retorno y no como media estadística añadida.

La clausura \(\pi\) pertenece a otro nivel del mismo estado HMT: es la salida
coinductiva del lector regional ya construido, no un autovalor nuevo de
\(F_{\rm av}\). Sean
\[
 P_{\rm ar}=
 \begin{pmatrix}1&0\\0&0\end{pmatrix},
 \qquad
 P_{\rm vol}=
 \begin{pmatrix}0&0\\0&1\end{pmatrix}.
\]
Entre los operadores positivos que preservan las dos hojas, hay uno y sólo
uno cuya normalización volumétrica es uno y cuya marca areal es la clausura
generada \(\pi\):
\[
 D_\pi={\pi}P_{\rm ar}+P_{\rm vol}.
\]
La unicidad es relativa a estas tres cláusulas explícitas: positividad,
diagonalidad respecto de las hojas y las dos normalizaciones. Con
\(e_{\rm ar}=(1,0)^T\) y \(e_{\rm vol}=(0,1)^T\), el retorno marcado de
longitud \(n\) es
\[
 \Theta_n
 :=
 \frac{\langle e_{\rm ar},
 D_\pi F_{\rm av}^{\,n}e_{\rm ar}\rangle}
 {\langle e_{\rm vol},
 F_{\rm av}^{\,n}e_{\rm vol}\rangle}
 =
 \pi\frac{F_{n-1}}{F_{n+1}}.
\]
Por tanto,
\[
 \boxed{\displaystyle
 \lim_{n\to\infty}\Theta_n
 =\frac{\pi}{\varphi^2}.}
\]
La marca \(\pi\) se aplica una sola vez, en el retorno de frontera;
no se reinyecta en cada transición. La combinatoria finita produce el modo
estable \(\varphi^{-2}\), el lector coinductivo produce la clausura
\(\pi\), y \(D_\pi\) compone ambos datos sin invertir su orden
causal.

La longitud nativa \(108\) ofrece un testigo exacto de esa convergencia:
\[
 (F_{\rm av}^{108})_{\rm ar,ar}
 =F_{107}=10284720757613717413913,
\]
\[
 (F_{\rm av}^{108})_{\rm vol,vol}
 =F_{109}=26925748508234281076009.
\]
En consecuencia,
\[
 \left|
 \frac{F_{107}}{F_{109}}-\varphi^{-2}
 \right|
 =
 6.1684979916614407936\ldots\times10^{-46},
\]
y, después de la única marca de clausura,
\[
 \left|
 \pi\frac{F_{107}}{F_{109}}
 -\frac{\pi}{\varphi^2}
 \right|
 =
 1.9378907974286976068\ldots\times10^{-45}.
\]
Estos errores miden la aproximación finita del retorno; no son ajustes de
\(\pi\) ni de \(\varphi\).

\paragraph{Lector cilíndrico de la razón.}
Sean \(I_n^\pi=[\underline\pi_n,\overline\pi_n]\) e
\(I_n^\varphi=[\underline\varphi_n,\overline\varphi_n]\) los cilindros
positivos y encajados producidos por los dos lectores HMT\@. La operación de
intervalos
\[
 \mathcal Q(I,J)
 :=
 \left[
 \frac{\inf I}{(\sup J)^2},
 \frac{\sup I}{(\inf J)^2}
 \right]
\]
es el menor intervalo que contiene \(x/y^2\) para todo
\((x,y)\in I\times J\). Preserva el encaje: si los cilindros de entrada se
refinan, también se refina \(\mathcal Q(I_n^\pi,I_n^\varphi)\). Como sus
diámetros tienden a cero y el límite de \(I_n^\varphi\) es positivo,
\[
 \bigcap_n\mathcal Q(I_n^\pi,I_n^\varphi)
 =
 \left\{\frac{\pi}{\varphi^2}\right\}.
\]
Con los doce bloques en base mil ya publicados, el intervalo cociente fuerza
treinta y cinco cifras decimales:
\[
 \frac{\pi}{\varphi^2}
 =
 1.19998161486432666111577775331680692\ldots.
\]
Esta construcción conserva la procedencia de cada cifra: el numerador
procede del cilindro de clausura y el denominador de la autoescala.

\paragraph{Separación algebraica entre lectura areal, lectura volumétrica y sus codominios respectivos}
La razón completa no puede salir de la sola álgebra racional de
\(F_{\rm av}\). En efecto, sus funciones racionales espectrales pertenecen a
\(\mathbb Q(\sqrt5)\); allí se encuentra \(\varphi^{-2}\), pero no el número
trascendente \(\pi\). Por consiguiente, toda derivación que pretendiese
obtener \(\pi/\varphi^2\) únicamente de la matriz \(2\times2\) estaría
mal tipada. El dato trascendente debe proceder del lector coinductivo de
clausura, o de un cociclo de frontera equivalente demostrado
independientemente. Este control negativo fija la división de trabajo entre
la correspondencia finita y el lector de número.

\subsubsection{Involución de intercambio del cociente}
\label{sec:espejo-pi-phi-electron-hmt}

La correspondencia admite dos orientaciones exactas de una misma ley. Si
\[
 \mathscr R(x,y)=\frac{x}{y^2},
 \qquad
 \mathsf J(x,y)=(y,x),
\]
entonces \(\mathsf J^2=I\) y
\[
 \mathscr R(\pi,\varphi)=\frac{\pi}{\varphi^2},
 \qquad
 (\mathscr R\circ\mathsf J)(\pi,\varphi)
 =\frac{\varphi}{\pi^2}.
\]
Se verifican además las identidades
\[
 \mathscr R(\pi,\varphi)\mathscr R(\varphi,\pi)
 =\frac1{\pi\varphi},
 \qquad
 \frac{\mathscr R(\pi,\varphi)}
 {\mathscr R(\varphi,\pi)}
 =\left(\frac{\pi}{\varphi}\right)^3.
\]
Escribiendo
\(\mathcal B_{\rm av}^{(\pi)}=\mathscr R(\pi,\varphi)\), se obtiene
también la reparametrización exacta
\[
 \frac{\varphi}{\pi^2}
 =
 \frac1{\varphi^3
 \bigl(\mathcal B_{\rm av}^{(\pi)}\bigr)^2},
 \qquad
 e^{\varphi/\pi^2}
 =
 \exp\!\left[
 \frac1{\varphi^3
 \bigl(\mathcal B_{\rm av}^{(\pi)}\bigr)^2}
 \right].
\]
Así, el cambio de orientación entre cierre y crecimiento transporta la
razón de frontera a su lectura intercambiada:
\[
 \frac{\pi}{\varphi^2}
 \xleftrightarrow{\ \mathsf J\ }
 \frac{\varphi}{\pi^2}.
\]
La primera orientación lee el cierre areal marcado sobre la memoria
volumétrica cuadrática; la segunda lee el crecimiento interior sobre el
cierre cuadrático. Esta correspondencia es un teorema del lector HMT. La
Mecánica Dimensional aplica después la segunda orientación dentro de la
construcción electrónica, sin convertir al electrón en antecedente de
\(\pi\), \(\varphi\) o de la correspondencia.

\subsubsection{Principio de Ambivalencia de los lectores areal y volumétrico}
\label{sec:ambivalencia-grav-inercial-pi-phi2}

El Principio de Ambivalencia afirma que la lectura areal y la lectura
volumétrica son dos funcionales coordinados, pero no intercambiables, de un
mismo estado con memoria. Su coincidencia en la diagonal es
\[
 \mathcal L_{\rm ar}^{\rm diag}(x)
 =\mathcal L_{\rm vol}^{\rm diag}(x).
\]
La igualdad diagonal no borra la ambivalencia fuera de ella.

Sobre la frontera areal--volumétrica, ambos lectores factorizan a través de
una misma amplitud positiva de memoria \(\mathcal M(x)\):
\[
 \mathcal L_{\rm vol}(x)=\mu_{\rm vol}\mathcal M(x),
 \qquad
 \mathcal L_{\rm ar}(x)=
 \pi\mu_{\rm ar}\mathcal M(x).
\]
Entonces el teorema de Parry y la marca de retorno dan
\[
 \boxed{\displaystyle
 \frac{\mathcal L_{\rm ar}(x)}
 {\mathcal L_{\rm vol}(x)}
 =\frac{\pi}{\varphi^2}.}
\]
Con la normalización volumétrica, la diferencia relativa de lectores es
\[
 \frac{\mathcal L_{\rm ar}-\mathcal L_{\rm vol}}
 {\mathcal L_{\rm vol}}
 =
 \frac{\pi}{\varphi^2}-1
 =
 0.1999816148643266611\ldots.
\]
La lectura física de esta correspondencia se realiza más tarde, en el capítulo de
Cartan--Holst, después de construir holonomía, co-tétrada, conexión, torsión
y curvatura. Aquí el resultado es enteramente HMT: una razón exacta entre
dos lectores de la misma memoria de hoja.

Conviene distinguirlos del extractor enriquecido que se utiliza después en
Mecánica Dimensional. Allí una historia completa conserva el conteo de acción
\(n_A\), los doce eventos orientados \(e_0,\ldots,e_{11}\) y la cocadena
torsional a nivel de paso; de ellos se obtienen
\[
 s_C=\sum_{i=0}^{5}(e_i+e_{i+6}),
 \qquad
 W_\pm=6n_A\pm s_C.
\]
Los canales ya generados por HMT satisfacen entonces la identidad exacta
\[
 \exp\!\left(n_AA_{\rm rad}+\frac{s_C}{6}C^\ast_{\rm rad}\right)
 =\left(q_+^{-W_+}q_-^{-W_-}\right)^{1/12}.
\]
Por tanto, el paso ruta--carácter angular no está ausente: se realiza sobre
el registro dodecafásico enriquecido. Lo que este cálculo de fase prueba por sí
solo son \(\Omega_{\rm sw},S,G\); no los sustituye ni permite reconstruir el
registro completo después de haberlo proyectado.

La medida uniforme sobre las \(104\,976\) presentaciones desciende a átomos
de masas \(1/729\), \(1/243\) y \(1/54\), pero no a \(1/468\). En la
microfibra quíntuple de \(\pi\), las dos probabilidades son
\[
 \tau_{468}(\mathscr R_\pi)=\frac5{468},
 \qquad
 F_\#\mu_{\rm seed}(\mathscr R_\pi)=\frac7{729}.
\]
Tanto la elevación promediada como la elevación de fase realizan la primera: no
empujan la medida uniforme sobre semillas, sino que asignan la misma masa
total a cada clase observable. La elevación de fase conserva además la
quietud microscópica en los pasos neutros.

Para relacionar la elevación de fase con las historias enriquecidas de
profundidad seis hace falta un núcleo \(\mathcal K_6\) sobre
\(\mathcal H_6^{\rm cat}\times C_9\). Éste desciende a
\(\widehat{\mathcal K}_{\rm ph}\) a través de
\(\rho_6\times\mathrm{id}_{C_9}\) si y sólo si satisface, en cada fase, la
condición de agregabilidad fuerte
\[
 \sum_{\rho_6(g)=y}\mathcal K_6((h,r),(g,r+1))
 =L_{\tau(r)}(\rho_6(h),y),
\]
independientemente de \(h\) dentro de cada fibra de \(\rho_6\). En ese caso,
la composición con \(F\) proyecta a \(\mathcal K_{\rm ph}\). La construcción
de una familia \((\mathcal K_m)_{m\geq6}\) que satisfaga esta identidad y
conmute con todos los truncamientos es el paso restante hacia una
transferencia del límite inverso.

\begingroup
\setlength{\parskip}{1pt}
\setlength{\abovedisplayskip}{3pt}\setlength{\belowdisplayskip}{3pt}
\setlength{\abovedisplayshortskip}{2pt}\setlength{\belowdisplayshortskip}{2pt}
\section{Cociclo de escala y conservación cilíndrica}
\label{sec:cociclo-parry-escala}

Sea \(r=(1,\varphi)\) un vector de Perron--Frobenius de
\(F_{\rm av}\). La matriz de transición de Parry es
\[
P_{ij}=\frac{(F_{\rm av})_{ij}r_j}{\varphi r_i}
=
\begin{pmatrix}
0&1\\
\varphi^{-2}&\varphi^{-1}
\end{pmatrix}.
\]
Para una ruta admisible \(\gamma:i\to j\) de longitud \(n\),
\[
\mathbb P(\gamma\mid i)=\frac{r_j}{\varphi^nr_i},
\qquad
\boxed{\chi_P(\gamma)=\varphi^n\frac{r_i}{r_j}.}
\]
La telescopía prueba
\[
\chi_P(\gamma_2\circ\gamma_1)
=\chi_P(\gamma_2)\chi_P(\gamma_1).
\]
En una ruta cerrada, \(\chi_P=\varphi^n\); su cuadrado y su inverso
generan las combinaciones
\[
C_n=\frac{\chi_P^2+\chi_P^{-2}}2,\qquad
S_n=\frac{\chi_P^2-\chi_P^{-2}}2,
\]
que satisfacen
\[
x_{n+1}=3x_n-x_{n-1}.
\]

Como \(F_{\rm av}\) es simétrica, sea
\[
\ell=\frac{r}{1+\varphi^2},\qquad \ell^{\mathsf T}r=1.
\]
La medida de un cilindro es
\[
\mu[x_0\cdots x_n]=\ell_{x_0}r_{x_n}\varphi^{-n}.
\]
Por tanto,
\[
\boxed{
V_n(x)
=\frac{\mu[x_0]}{\mu[x_0\cdots x_n]}
=\varphi^n\frac{r_{x_0}}{r_{x_n}}
=\chi_P(x_0\cdots x_n).}
\]
Para una escala de dimensión tipada \(D>0\),
\[
a_n^{(D)}(x)=V_n(x)^{1/D}
\]
cumple la ley de conservación cilíndrica
\[
\boxed{
\bigl(a_n^{(D)}(x)\bigr)^D
\mu[x_0\cdots x_n]=\mu[x_0].}
\]

En rango espacial tres y sobre retornos cerrados, escribimos
\(a_n:=a_n^{(3)}\) y obtenemos
\[
\boxed{
\frac{a_9}{a_0}=\varphi^3,\qquad
\frac{a_{27}}{a_0}=\varphi^9,\qquad
\frac{a_{108}}{a_0}=\varphi^{36}.}
\]
Si \(A_k:=a_{9k}\) para \(k\geq1\), entonces, para \(k\geq2\),
\[
A_{k+1}-2A_k+A_{k-1}
=A_{k-1}(\varphi^3-1)^2>0.
\]
La desigualdad prueba convexidad en profundidad de retorno; una aceleración
física requiere además una carta de duración.

El cálculo primario de
\(\operatorname{Sym}^2(F_{\rm av})\) en
\cref{sec:retorno-marcado-pi-phi2} demuestra que su espectro es
\(\{\varphi^2,-1,\varphi^{-2}\}\). Conservamos aquí únicamente el corolario
para el cociclo de escala: sobre una ruta cerrada con \(V_n=\varphi^n\), el
modo estable obedece
\[
\boxed{M_n=\varphi^{-2n}=V_n^{-2}.}
\]
La identificación de estos tres autoespacios con las tres hojas del TRIT
exige un entrelazador; la coincidencia de cardinalidad no lo sustituye. El
corolario quedaría refutado si la ruta cerrada no satisficiera
\(V_n=\varphi^n\) o si el modo estable no tuviera autovalor
\(\varphi^{-2}\).

\section{Teorema de prolongación dinámica: hipótesis globales y dominio certificado}

Se han construido y separado explícitamente los objetos siguientes del
problema finito:
\[
 F,
 \qquad
 F_{\rm obs},
 \qquad
 \mathcal K_{\rm ph},
 \qquad
 \widehat{\mathcal K}_{\rm ph},
 \qquad
 K_{\rm cat},
 \qquad
 \tau_{468},
 \qquad
 F_{\rm av},
 \qquad
 \operatorname{Sym}^2(F_{\rm av}),
 \qquad
 D_\pi,
 \qquad
 \mathcal Q.
\]
Para \(F_{\rm obs}\) se ha probado el retorno determinista y se ha exhibido un
testigo explícito de no descenso producto. Para \(\mathcal K_{\rm ph}\) se
han probado la doble estocasticidad, la irreducibilidad, el período nueve, el
retorno nonádico uniforme y la unicidad de
\(\pi_{\rm ph}\). Para \(\widehat{\mathcal K}_{\rm ph}\) se han probado la
agregabilidad fuerte, la conservación microscópica de la quietud y la medida
estacionaria compensada. \(K_{\rm cat}\) ya no es un promedio sin
procedencia: es la sombra de un paso al promediar la fase uniforme.
El cociente de modos de TPK induce \(F_{\rm av}\); su elevación cuadrática
separa los modos \(\varphi^2,-1,\varphi^{-2}\). El operador \(D_\pi\)
marca una sola vez el retorno de frontera con la clausura generada, y el
lector intervalar \(\mathcal Q\) conserva el encaje de los cilindros de
\(\pi\) y \(\varphi\). Estos cuatro objetos construyen
\(\pi/\varphi^2\) sin confundir la frecuencia temporal,
la medida uniforme y la medida de Parry.

La fuerza de esta inducción es relativa y concreta. La palabra TPK fuerza los
pesos \(1/3\). La interpretación APP de todas las continuaciones compatibles,
junto con el conteo uniforme, fuerza \(E_+\) y \(E_\times\). La trayectoria
determinista \(F_{\rm obs}\) no fuerza esos operadores y, de hecho, no los
produce. Éste no es un defecto del resultado, sino la diferencia matemática
entre evolución de un camino y transferencia de un espacio de caminos.

\ifdefined\HMTInternalTraceability
\begin{criterion}[Elevación enriquecida de la reversión]
\label{pf84:SYN31-RHO-LIFT}
Se busca un subdominio \(D_\rho\subset X_{\rm TPK}^{\rm enr}\) y un mapa
\[
 \Theta_\rho:D_\rho\longrightarrow X_{\rm TPK}^{\rm enr}
\]
que eleve la reversión entera \(27\mapsto72\) del
\cref{pf84:SYN30-RHO-DESC,pf84:SYN30-CROSS-WITNESS}. Debe conmutar con los
lectores residual, entero y genealógico, cuantificar el cociente y el
acarreo, y preservar orientación, ruta, retorno y, cuando el transporte sea
invertible, monodromía, región y supervivencia.
Sólo si \(D_\rho\) es total e invariante bajo \(\Theta_\rho\) puede exigirse
\(\Theta_\rho^2=I\) y denominarse involución.

No existe todavía esta construcción ni una prueba de involutividad. El
programa queda falsado si no existe el mapa, si falla uno de los cuadrados
de lectura, si se pierde memoria sin cuantificarla o si
\(\Theta_\rho^2\ne I\) en un dominio donde se reclame involución.
\end{criterion}

\begin{criterion}[Elevación enriquecida del cuadrado]
\label{pf84:SYN31-SQ-LIFT}
Se busca un subdominio \(D_s\subset X_{\rm TPK}^{\rm enr}\) y un mapa
\[
 \Theta_s:D_s\longrightarrow X_{\rm TPK}^{\rm enr}
\]
que eleve \(s(n)=n^2\), en particular \(27\mapsto729\), y conmute con los
lectores residual, entero y genealógico de
\cref{pf84:SYN30-SQ-DESC,pf84:SYN30-CROSS-WITNESS}. Debe cuantificar
cociente y acarreo y preservar orientación, ruta, retorno y, cuando el
transporte sea invertible, monodromía, región y supervivencia. Este mapa no
es involutivo y no se le impone
\(\Theta_s^2=I\).

La prueba permanece ausente. La inexistencia de \(\Theta_s\), el fallo de
un cuadrado de lectura o la pérdida no controlada de cualquiera de las
coordenadas enriquecidas refutarían el programa. La ablación
\(81\mapsto72\) al cambiar \(q^\times(9,9)\) es un control independiente:
no es una rama de \(\Theta_\rho\) ni de \(\Theta_s\), y no prueba ninguno
de los dos mapas.
\end{criterion}
\fi

La conclusión finita es completa en su dominio: el retorno uniforme, la
medida estacionaria, la agregabilidad a profundidad seis y la separación
entre evolución de una ruta y transferencia del espacio de rutas quedan
demostrados por los objetos construidos. En particular,
\(\widehat{\mathcal K}_{\rm ph}\) no se identifica con la cronología de
\(F_{\rm obs}\), como muestra \cref{prop:retorno-fobs-identidad}.
\par\endgroup
 \end{HMTScientificOwnerSubordinated}
}
 
\chapter{Construcción correlativa e inseparable del continuo}
\chaptermark{Estructura discreta del continuo}
\label{chap:83-04}

Las cinco construcciones se generan correlativamente sobre un único estado
APP--TRIT--TPK. Constituyen operaciones consustanciales del continuo y no una
enumeración exhaustiva de sus propiedades, cinco ramas independientes ni una
selección retrospectiva determinada por aplicaciones posteriores.

\begin{HMTScientificOwnerSubordinated}
% Corrección exclusivamente editorial del segundo epígrafe de 04b1. El
% fuente teoremática canónico permanece inmutable y se conserva íntegro; aquí se
% traduce la nomenclatura interna «levantamiento TRIT» a una formulación
% pública que nombra la operación y el resultado matemático.
\newcount\HMTContinuoCommonSection
\HMTContinuoCommonSection=0
\let\HMTContinuoCommonOriginalSection\section
\renewcommand{\section}[1]{%
  \advance\HMTContinuoCommonSection by 1
  \ifnum\HMTContinuoCommonSection=2
    \HMTContinuoCommonOriginalSection{Extensión trítica del dato de incidencia: orientación equilibrada y cociclo de acarreo genealógico}%
  \else
    \HMTContinuoCommonOriginalSection{#1}%
  \fi
}
\chapter{Estado común, árbol de prolongaciones y medida genealógica}
\label{chap:continuo-estado-arbol-medida-comun}

\begin{thesisbox}[title={Un solo estado holonómico integral antes de toda discriminación interna}]
El objeto de este capítulo es el estado finito completo producido por
\(\APP\), orientado por TRIT y transportado por \(\TPK\).  A partir de él se
construyen el grupoide de simetrías reversibles, el árbol de todas las
prolongaciones admisibles, la multisección de historias completas y su medida
canónica.  No se introduce ningún dominio especializado: el dominio es el
conjunto total de estados que efectivamente genera el mismo mecanismo
\(\APP\)--TRIT--\(\TPK\).  Tampoco se selecciona una hoja, una continuación o
una lectura posterior.
\end{thesisbox}

\section{Datos finitos producidos por APP}
\label{sec:continuo-comun-app-finito}

Para cada profundidad \(k\geq0\), denotamos por \(\mathsf A_k\) el conjunto
finito de prefijos celulares producidos por APP y por
\begin{equation}
\label{eq:continuo-comun-app-truncamiento}
 \rho^{\APP}_{k+1,k}:\mathsf A_{k+1}\longrightarrow\mathsf A_k
\end{equation}
el borrado del último refinamiento.  Este borrado conserva el prefijo ya
construido.  La célula neutra de APP da una prolongación
\begin{equation}
\label{eq:continuo-comun-app-degeneracion}
 \eta^{\APP}_k:\mathsf A_k\longrightarrow\mathsf A_{k+1},
 \qquad
 \rho^{\APP}_{k+1,k}\eta^{\APP}_k=I_{\mathsf A_k}.
\end{equation}
Escribimos
\begin{equation}
\label{eq:continuo-comun-emision-neutra-app}
 \varepsilon^0_k(a):a\longrightarrow\eta^{\APP}_k(a)
\end{equation}
para la emisión que realiza esa prolongación.  Así se mantienen separados su
destino, \(\eta^{\APP}_k(a)\), y la operación que lo produce,
\(\varepsilon^0_k(a)\).
No se interpreta \(\eta^{\APP}_k\) como ausencia de historia: registra un
refinamiento de longitud uno con emisión neutra y conserva materialmente el
estado anterior.

Cada \(a\in\mathsf A_k\) porta simultáneamente las dos hojas de APP,
\begin{equation}
\label{eq:continuo-comun-dos-hojas}
 \operatorname{Sh}_k(a)
 =\bigl(\Sigma_k(a),\Pi_k(a)\bigr)
 \in\mathsf H_k^{\Sigma}\times\mathsf H_k^{\Pi},
\end{equation}
y no un elemento elegido de la pareja.  Para una emisión elemental
\(\epsilon:a\to a'\), APP produce además dos divisiones euclídeas tipadas,
\begin{equation}
\label{eq:continuo-comun-residuo-cociente-app}
 \delta^{\diamond}(\epsilon)
 =\operatorname{res}^{\diamond}(\epsilon)
  +9\operatorname{quo}^{\diamond}(\epsilon),
 \qquad
 \operatorname{res}^{\diamond}(\epsilon)\in\{0,\ldots,8\},
 \quad
 \operatorname{quo}^{\diamond}(\epsilon)\in\mathbb Z,
\end{equation}
donde \(\diamond\in\{\Sigma,\Pi\}\).  La etiqueta \(\diamond\) sólo señala
qué hoja se está leyendo; ambas divisiones permanecen en el mismo registro.

La fuente y el destino celulares de \(\epsilon\) determinan su frontera
orientada
\begin{equation}
\label{eq:continuo-comun-frontera-elemental}
 \partial\epsilon:=t(\epsilon)-s(\epsilon).
\end{equation}
La emisión neutra satisface
\begin{equation}
\label{eq:continuo-comun-neutral-app-datos}
 \operatorname{res}^{\diamond}(\varepsilon^0_k(a))=0,
 \qquad
 \operatorname{quo}^{\diamond}(\varepsilon^0_k(a))=0,
 \qquad
 \partial\varepsilon^0_k(a)=0.
\end{equation}

\begin{definition}[Grafo dirigido de emisiones APP]
\label{def:continuo-comun-grafo-app}
Sea \(\mathsf E_k(a)\) el conjunto finito de emisiones APP de longitud uno
que parten de \(a\in\mathsf A_k\).  Conservamos las emisiones con
multiplicidad: dos operaciones con el mismo destino pero con distinto
residuo, cociente o procedencia son elementos diferentes de
\(\mathsf E_k(a)\).  Por \eqref{eq:continuo-comun-app-degeneracion},
\begin{equation}
\label{eq:continuo-comun-out-no-vacio-app}
 \varepsilon^0_k(a)\in\mathsf E_k(a),
 \qquad
 1\leq \#\mathsf E_k(a)<\infty.
\end{equation}
\end{definition}

\section{Levantamiento TRIT: orientación y acarreo genealógico producidos}
\label{sec:continuo-comun-trit-lift}

La incidencia firmada de una emisión APP es un entero
\(\Delta(\epsilon)\).  El representante equilibrado módulo tres da de forma
única
\begin{equation}
\label{eq:continuo-comun-trit-euclid}
 \Delta(\epsilon)
 =\operatorname{or}(\epsilon)+3\kappa(\epsilon),
 \qquad
 \operatorname{or}(\epsilon)\in\{-1,0,+1\},
 \quad
 \kappa(\epsilon)\in\mathbb Z.
\end{equation}
Así, la orientación y el carry no son condiciones impuestas a una historia:
son funciones totales del dato APP.

Para hacer visible la composición, sea
\(\mathsf C_{\epsilon}:\mathsf M_{s(\epsilon)}\to
\mathsf M_{t(\epsilon)}\) el transporte de memoria inducido por la emisión.
El carry nativo de una composición cumple
\begin{equation}
\label{eq:continuo-comun-cociclo-carry}
 \kappa(\epsilon_2\epsilon_1)
 =\kappa(\epsilon_2)
  +\mathsf C_{\epsilon_2}\kappa(\epsilon_1),
\end{equation}
siempre que \(t(\epsilon_1)=s(\epsilon_2)\).  Para la identidad y la emisión
neutra,
\begin{equation}
\label{eq:continuo-comun-carry-unidad}
 \kappa(I_a)=0,
 \qquad
 \kappa(\varepsilon^0_k(a))=0,
 \qquad
 \mathsf C_{I_a}=I_{\mathsf M_a}.
\end{equation}

\begin{lemma}[Asociatividad del ledger de carry]
\label{lem:continuo-comun-asociatividad-carry}
Para tres emisiones componibles, la evaluación de
\(\kappa(\epsilon_3\epsilon_2\epsilon_1)\) no depende del parentizado.
\end{lemma}

\begin{proof}
Aplicando dos veces \eqref{eq:continuo-comun-cociclo-carry} al parentizado
izquierdo se obtiene
\[
 \kappa(\epsilon_3)
 +\mathsf C_{\epsilon_3}\kappa(\epsilon_2)
 +\mathsf C_{\epsilon_3}\mathsf C_{\epsilon_2}
  \kappa(\epsilon_1).
\]
El parentizado derecho produce la misma expresión porque el transporte de
memoria es functorial:
\(\mathsf C_{\epsilon_3\epsilon_2}
=\mathsf C_{\epsilon_3}\mathsf C_{\epsilon_2}\).
\end{proof}

La reversión orientada de una emisión reversible se denota
\(\overline\epsilon\).  TRIT actúa por
\begin{equation}
\label{eq:continuo-comun-reversion-trit}
 \operatorname{or}(\overline\epsilon)
 =-\operatorname{or}(\epsilon),
 \qquad
 \partial\overline\epsilon=-\partial\epsilon,
 \qquad
 \mathsf C_{\overline\epsilon}=\mathsf C_\epsilon^{-1}.
\end{equation}
El carry inverso queda determinado, no elegido, por
\begin{equation}
\label{eq:continuo-comun-carry-inverso}
 \kappa(\overline\epsilon)
 =-\mathsf C_\epsilon^{-1}\kappa(\epsilon),
\end{equation}
pues \(\kappa(\overline\epsilon\epsilon)=0\).

\section{Transporte TPK de las \(2\) hojas y de la memoria}
\label{sec:continuo-comun-tpk-transporte}

La cronolog\'ia de TPK forma parte del tipo de cada emisi\'on y no se
reconstruye a posteriori contando letras. Recuperamos aqu\'i la acci\'on del
od\'ometro de \cref{def:odometro-9-12,prop:proyeccion-fase-longitud}. Sea
\begin{equation}
\label{eq:continuo-comun-ciclo-fase-tpk}
 C_9:=\mathbb Z/9\mathbb Z,\qquad
 \widetilde{\mathcal O}_{9}:=\mathbb Z_{\geq0}\times C_9,\qquad
 \mathcal O_{9,12}:=\mathbb Z/12\mathbb Z\times C_9.
\end{equation}
El estado desenrollado de fase y memoria es
\(\mathsf Q_k=(w_k,r_k)\in\widetilde{\mathcal O}_{9}\); su reducción
\((w,r)\mapsto([w]_{12},r)\) recupera
\(\mathcal O_{9,12}\). La semilla canónica del reloj es
\begin{equation}
\label{eq:continuo-comun-semilla-fase-tpk}
 \mathsf Q_0(a_0):=(0,[0]_9)\qquad(a_0\in\mathsf A_0);
\end{equation}
las demás elecciones de origen son sus traslaciones coherentes por \(C_9\)
y no cambian la ley de actualización. Para
\(\bar r\in\{0,\ldots,8\}\), ponemos
\begin{equation}
\label{eq:continuo-comun-carry-fase-tpk}
 \chi_9(r):=\left\lfloor\frac{\bar r+1}{9}\right\rfloor.
\end{equation}
Toda emisión \(\epsilon\) que lleva profundidad \(k\) a profundidad \(k+1\)
actúa en el odómetro por
\begin{equation}
\label{eq:continuo-comun-accion-fase-tpk}
 \sigma_9(w,r):=(w+\chi_9(r),r+[1]_9),\qquad
 \mathsf Q_{k+1}(\epsilon_*\widehat x_k)
 =\sigma_9\mathsf Q_k(\widehat x_k).
\end{equation}
Su componente de fase es, por definición tipada,
\begin{equation}
\label{eq:continuo-comun-fase-emision-tpk}
 \operatorname{ph}(\epsilon):r_k\longmapsto r_k+[1]_9,
 \qquad
 \operatorname{ph}(\epsilon_n\cdots\epsilon_1):
 r_k\longmapsto r_k+[n]_9.
\end{equation}
Escribimos \(g_k:=r_k\in C_9\) para la fase publicada y
\(q^{(9)}_k:=w_k\) para la memoria vertical, y
\(\beta_k:=[w_k]_{12}\) para su reducción dodecafásica. La iteración de
\eqref{eq:continuo-comun-accion-fase-tpk} es
\begin{equation}
\label{eq:continuo-comun-iteracion-fase-tpk}
 \sigma_9^n(w,r)=
 \left(w+\left\lfloor\frac{\bar r+n}{9}\right\rfloor,
       r+[n]_9\right).
\end{equation}
En efecto, los sumandos
\(\chi_9(r),\chi_9(r+[1]_9),\ldots,\chi_9(r+[n-1]_9)\)
cuentan exactamente los cruces del representante \(8\) al representante
\(0\), que son
\(\lfloor(\bar r+n)/9\rfloor\). En particular,
\begin{equation}
\label{eq:continuo-comun-retorno-fase-avance-memoria}
 \sigma_9^9(w,r)=(w+1,r),\qquad
 g_{k+9}=g_k,\qquad q^{(9)}_{k+9}=q^{(9)}_k+1.
\end{equation}
Incluso la emisión neutra de APP realiza
\(g_k\mapsto g_{k+1}\): es neutra en residuo, cociente, carry y frontera,
pero registra que ha ocurrido una extensión. La palabra de fase
\begin{equation}
\label{eq:continuo-comun-palabra-fase-tpk}
 \boldsymbol\tau_k^{(N)}:=(g_k,g_{k+1},\ldots,g_N),
 \qquad g_{j+1}=g_j+[1]_9,
\end{equation}
y la memoria \(q^{(9)}\) tipan por objetos distintos el retorno de fase y
el cambio de estado.

\subsection{Renovación simultánea de la fibra completa}
\label{sec:continuo-comun-kph}

El odómetro anterior conserva la cronología de cada extensión. La misma
vuelta nonádica posee además la realización Markov interna, ya construida
sobre el catálogo APP--TPK, que renueva simultáneamente las dos coordenadas
de la fibra completa. Es esencial conservar ambas realizaciones y no
identificar una arista genérica del árbol con una emisión del catálogo sin
un mapa de tipado.

\begin{proposition}[Realización Markov interna de la conexión nonádica]
\label{prop:continuo-comun-kph-heredada}
\label{pII-full:thm:seleccion-novena-fase}
Sean el dominio de semillas APP, su emisión observable y sus multiplicidades
\[
 T_{\min}:\Omega_{\rm seed}\twoheadrightarrow\mathcal U_6,
 \qquad
 m(u):=\#T_{\min}^{-1}(u),
\]
y la factorización producida por el TPK
\[
 \mathcal U_6\simeq\mathcal S_+\times\mathcal S_\times,
 \qquad |\mathcal S_+|=18,\qquad |\mathcal S_\times|=26.
\]
Sobre funciones en \(\mathcal U_6\), las dos esperanzas condicionales actúan
por
\[
 (E_+f)(s,e)=\frac1{18}\sum_{s'\in\mathcal S_+}f(s',e),
 \qquad
 (E_\times f)(s,e)=\frac1{26}
 \sum_{e'\in\mathcal S_\times}f(s,e').
\]
Para
\[
 \tau=(+1,-1,0,+1,-1,0,+1,-1,0),
 \qquad
 P_{+1}=E_+,\quad P_{-1}=E_\times,\quad P_0=I,
\]
la transferencia de catálogo y su elevación a todas las semillas son
\begin{align}
 \mathcal K_{\rm ph}\bigl((u,r),(v,r+1)\bigr)
 &=P_{\tau(r)}(u,v),
 \label{eq:continuo-comun-kph-catalogo}\\
 L_0(\omega,\omega')&=\mathbf1_{\{\omega=\omega'\}},\nonumber\\
 L_+(\omega,\omega')
 &=\frac{E_+(T_{\min}\omega,T_{\min}\omega')}
         {m(T_{\min}\omega')},\nonumber\\
 L_\times(\omega,\omega')
 &=\frac{E_\times(T_{\min}\omega,T_{\min}\omega')}
         {m(T_{\min}\omega')},\nonumber\\
 \widehat{\mathcal K}_{\rm ph}
 \bigl((\omega,r),(\omega',r+1)\bigr)
 &=L_{\tau(r)}(\omega,\omega').
 \label{eq:continuo-comun-kph-elevada}
\end{align}
Para \(u=T_{\min}\omega\), se cumple fase a fase
\begin{equation}
\label{eq:continuo-comun-kph-lumpability}
 \sum_{\omega':\,T_{\min}\omega'=v}
 \widehat{\mathcal K}_{\rm ph}
 \bigl((\omega,r),(\omega',r+1)\bigr)
 =
 \mathcal K_{\rm ph}\bigl((u,r),(v,r+1)\bigr).
\end{equation}
En particular,
\begin{equation}
\label{eq:continuo-comun-kph-novena-potencia}
 \boxed{\;
 \mathcal K_{\rm ph}^{\,9}\big|_r=E_+E_\times,
 \qquad
 (E_+E_\times)(u,v)=\frac1{468}.
 \;}
\end{equation}
\end{proposition}

\begin{proof}
En una fase \(+\) o \(\times\), sumar
\eqref{eq:continuo-comun-kph-elevada} sobre
\(T_{\min}^{-1}(v)\) cancela exactamente el denominador \(m(v)\); en una
fase cero, \(L_0\) proyecta a la identidad. Esto demuestra
\eqref{eq:continuo-comun-kph-lumpability}. Los operadores \(E_+\),
\(E_\times\) e \(I\) son idempotentes y conmutan, y cada uno aparece tres
veces en la palabra nonádica. Por tanto
\[
 \prod_{r\in C_9}P_{\tau(r)}
 =E_+^3E_\times^3I^3=E_+E_\times.
\]
La composición renueva primero una coordenada de cardinal \(18\) y después
la otra de cardinal \(26\), de modo que cada destino recibe
\(1/(18\cdot26)=1/468\).
\end{proof}

\begin{remark}[Dos realizaciones internas, sin identificación ficticia]
\label{rem:continuo-comun-kph-no-identificacion-aristas}
La transferencia \(\mathcal K_{\rm ph}\) es la realización global del
catálogo completo y \(\sigma_9\) es la acción cronológica sobre cada estado
del árbol. No se sustituye una por otra. Una igualdad arista a arista
requeriría un mapa adicional
\[
 \operatorname{Cat}_{k,x}:\operatorname{Out}(x)\longrightarrow\mathcal U_6
\]
que verificase la correspondiente identidad de suma sobre fibras. Ese mapa
no se usa en este volumen. La monodromía completa
\eqref{eq:continuo-comun-kph-novena-potencia} procede del catálogo
APP--TPK, mientras la genealogía de cada historia permanece en el árbol y
su ledger.
\end{remark}

TPK asigna a cada emisión admisible un transporte tipado de la pareja de
hojas,
\begin{equation}
\label{eq:continuo-comun-transporte-hojas}
 U_\epsilon:
 \mathsf H_{s(\epsilon)}^{\Sigma}\times
 \mathsf H_{s(\epsilon)}^{\Pi}
 \longrightarrow
 \mathsf H_{t(\epsilon)}^{\Sigma}\times
 \mathsf H_{t(\epsilon)}^{\Pi},
\end{equation}
con
\begin{equation}
\label{eq:continuo-comun-transporte-hojas-funtorial}
 U_{I_a}=I,
 \qquad
 U_{\epsilon_2\epsilon_1}=U_{\epsilon_2}U_{\epsilon_1}.
\end{equation}
El codominio de \(U_\epsilon\) sigue siendo la pareja completa.  Ninguna
operación de este capítulo sustituye
\((\Sigma,\Pi)\) por una de sus coordenadas.

Sea \(m\in\mathsf M_{s(\epsilon)}\) una memoria acumulada.  La acción TPK
sobre ella es la transformación afín
\begin{equation}
\label{eq:continuo-comun-accion-memoria}
 \mathsf U_\epsilon(m)
 :=\mathsf C_\epsilon m+\kappa(\epsilon).
\end{equation}

\begin{lemma}[Acción TPK sobre memoria]
\label{lem:continuo-comun-accion-memoria}
Las transformaciones \(\mathsf U_\epsilon\) satisfacen
\begin{equation}
\label{eq:continuo-comun-accion-memoria-composicion}
 \mathsf U_{I_a}=I,
 \qquad
 \mathsf U_{\epsilon_2\epsilon_1}
 =\mathsf U_{\epsilon_2}\mathsf U_{\epsilon_1}.
\end{equation}
\end{lemma}

\begin{proof}
La identidad se sigue de \eqref{eq:continuo-comun-carry-unidad}.  Para la
composición,
\begin{align*}
 \mathsf U_{\epsilon_2}\mathsf U_{\epsilon_1}(m)
 &=\mathsf C_{\epsilon_2}
   \bigl(\mathsf C_{\epsilon_1}m+\kappa(\epsilon_1)\bigr)
   +\kappa(\epsilon_2)\\
 &=\mathsf C_{\epsilon_2\epsilon_1}m
   +\kappa(\epsilon_2)
   +\mathsf C_{\epsilon_2}\kappa(\epsilon_1)\\
 &=\mathsf C_{\epsilon_2\epsilon_1}m
   +\kappa(\epsilon_2\epsilon_1),
\end{align*}
donde la última igualdad es
\eqref{eq:continuo-comun-cociclo-carry}.
\end{proof}

Una ruta finita es una palabra componible
\begin{equation}
\label{eq:continuo-comun-ruta}
 \gamma=\epsilon_1\epsilon_2\cdots\epsilon_k.
\end{equation}
Su frontera, transporte, carry y memoria son
\begin{align}
 \partial\gamma
 &:=\sum_{j=1}^{k}\partial\epsilon_j,
 \label{eq:continuo-comun-frontera-ruta}\\
 U_\gamma
 &:=U_{\epsilon_k}\cdots U_{\epsilon_1},
 \label{eq:continuo-comun-transporte-ruta}\\
 \kappa(\gamma)
 &:=\kappa(\epsilon_k)
   +\sum_{j=1}^{k-1}
    \mathsf C_{\epsilon_k}\cdots\mathsf C_{\epsilon_{j+1}}
    \kappa(\epsilon_j),
 \label{eq:continuo-comun-carry-ruta}\\
 m(\gamma;m_0)
 &:=\mathsf U_{\epsilon_k}\cdots
       \mathsf U_{\epsilon_1}(m_0).
 \label{eq:continuo-comun-memoria-ruta}
\end{align}
En \eqref{eq:continuo-comun-frontera-ruta}, los extremos interiores se
cancelan y queda
\begin{equation}
\label{eq:continuo-comun-frontera-telescopica}
 \partial\gamma=t(\epsilon_k)-s(\epsilon_1).
\end{equation}

\section{El estado común y su dominio total}
\label{sec:continuo-comun-estado-total}

\begin{definition}[Ledger elemental]
\label{def:continuo-comun-ledger-elemental}
Para una emisión que aparece entre dos prefijos enriquecidos consecutivos,
denotamos por
\(\mathsf Q^-(\epsilon)=(w^-(\epsilon),r^-(\epsilon))\) el odómetro del
prefijo fuente y por
\begin{equation}
\label{eq:continuo-comun-ledger-fase-extremos}
 \mathsf Q^+(\epsilon):=\sigma_9\mathsf Q^-(\epsilon)
 =(w^-(\epsilon)+\chi_9(r^-(\epsilon)),r^-(\epsilon)+[1]_9)
\end{equation}
el del prefijo destino. Así no se aplica el odómetro a una célula APP
desnuda. El registro de la emisión es
\begin{equation}
\label{eq:continuo-comun-ledger-elemental}
 \lambda(\epsilon):=\bigl(
 s(\epsilon),t(\epsilon),
 \operatorname{res}^{\Sigma}(\epsilon),
 \operatorname{quo}^{\Sigma}(\epsilon),
 \operatorname{res}^{\Pi}(\epsilon),
 \operatorname{quo}^{\Pi}(\epsilon),
 \operatorname{or}(\epsilon),
 \kappa(\epsilon),
 \partial\epsilon,
 \mathsf Q^-(\epsilon),\mathsf Q^+(\epsilon)\bigr).
\end{equation}
El ledger de \(\gamma\) es la palabra ordenada
\begin{equation}
\label{eq:continuo-comun-ledger-ruta}
 \operatorname{Led}(\gamma)
 :=\bigl(\lambda(\epsilon_1),\ldots,
          \lambda(\epsilon_k)\bigr).
\end{equation}
Conservar esta palabra, y no sólo su suma, distingue dos historias con la
misma salida terminal.
\end{definition}

\subsection{Lector completo del estado TPK}
\label{sec:continuo-comun-lector-estado-tpk}

La profundidad \(K\) de bloques TPK y la profundidad \(k\) del árbol de
extensiones son tipos distintos. En esta subsección \(K\) denota
exclusivamente la profundidad del estado TPK ya generado por el mecanismo
operatorio; no se identifica con \(k\) ni con una arista genérica.

\begin{definition}[Lector compacto de memoria nonádica]
\label{def:continuo-comun-lector-compacto-nonadico}
El estado holonómico integral TPK de profundidad \(K\) es
\begin{equation}
\label{eq:continuo-comun-estado-tpk-completo}
 v_K=
 \bigl(
 B_K,\boldsymbol\delta_K,\boldsymbol x_K,
 \boldsymbol{\mathcal C}_K,\Sigma_K,
 \omega_K,\ell_K,g(K),\Lambda_K
 \bigr),
 \qquad
 B_K=(u_K^\pi,u_K^e,u_K^\varphi).
\end{equation}
Su lector compacto no añade coordenadas: las publica desde ese mismo estado,
\begin{equation}
\label{eq:continuo-comun-lector-compacto-nonadico}
 \mathfrak r_9(v_K)
 :=\bigl(B_K,C_K,p_K,c_K,\Sigma_K,W_K,g(K)\bigr),
\end{equation}
donde
\begin{equation}
\label{eq:continuo-comun-lector-compacto-componentes}
\begin{aligned}
 C_K&:=\boldsymbol{\mathcal C}_K
      =\bigl(C_K^\chi\bigr)_{\chi\in\{\pi,e,\varphi\}},\\
 C_K^\chi&:=(N_K^\chi,L_K^\chi,r_K^\chi,
             D_K^\chi,A_K^\chi,B_K^\chi),\\
 p_K&:=(B_0,\ldots,B_K),\\
 c_K&:=(\boldsymbol\delta_K,\boldsymbol x_K,
       \kappa(\gamma^{\rm TPK}_{0:K})),
\end{aligned}
\end{equation}
y la sombra orientada de Witt se conserva como
\begin{equation}
\label{eq:continuo-comun-sombra-witt}
 W_K:=
 \bigl((u_K^\chi,u_K^\chi A_W)\bigr)_
       {\chi\in\{\pi,e,\varphi\}},
 \qquad
 A_W=
 \begin{pmatrix}
 0&1&1&1&1&1\\
 1&0&1&2&2&1\\
 1&1&0&1&2&2\\
 1&2&1&0&1&2\\
 1&2&2&1&0&1\\
 1&1&2&2&1&0
 \end{pmatrix},
 \qquad A_W^2=-I_6.
\end{equation}
La ruta \(\gamma^{\rm TPK}_{0:K}\) es la palabra de extensiones codificada
por \(\Lambda_K\); su prefijo de longitud \(j\) es el que produce
\(v_j\). La firma es la salida de la primitiva interna explícita
\begin{equation}
\label{eq:continuo-comun-primitiva-firma-tpk}
 \Sigma_B:
 \mathsf B_{\rm TPK}\times\mathsf C_{\rm TPK}
 \times\mathbb Z[\mathsf A]\times\Omega_{12}\times\mathsf L_{\rm TPK}
 \longrightarrow\mathsf S_{\rm TPK},
 \quad
 (B,c,\partial\gamma,\omega,\ell)\longmapsto
 \Sigma_B(B,c,\partial\gamma,\omega,\ell),
\end{equation}
con \(\Sigma_B(g\cdot-)=g\cdot\Sigma_B(-)\) para cada simetría TPK. Su
estructura de partida explícita es
\((B_K,c_K,\partial\gamma^{\rm TPK}_{0:K},\omega_K,\ell_K)\). No se la reemplaza por una
firma exterior ni se infiere retrospectivamente desde un lector numérico.
\end{definition}

La acción del paso \(K\mapsto K+1\) se escribe sin omitir el estado profundo.
Para cada \(\chi\in\{\pi,e,\varphi\}\),
\begin{equation}
\label{eq:continuo-comun-update-hensel-tpk}
 y_K^\chi=x_K^\chi A_H,\qquad
 u_{K+1}^\chi=y_K^\chi\bmod3,\qquad
 x_{K+1}^\chi=\frac{y_K^\chi-u_{K+1}^\chi}{3},
\end{equation}
con
\[
 A_H=
 \begin{pmatrix}
 2&2&2&1&2&1\\
 2&1&2&2&1&1\\
 1&0&0&1&0&2\\
 0&0&1&1&1&0\\
 2&2&2&0&2&1\\
 2&1&2&1&0&0
 \end{pmatrix},
 \qquad \det A_H=1.
\]
Si
\[
 \operatorname{ind}_3(u_1,\ldots,u_6)
 :=\sum_{j=1}^{6}\overline u_j\,3^{6-j}
 \in\{0,\ldots,728\},
\]
la actualización del prefijo y de sus \(729\) subcilindros es
\begin{align}
 N_{K+1}^\chi
 &=729N_K^\chi+\operatorname{ind}_3(u_{K+1}^\chi),
 &L_{K+1}^\chi&=L_K^\chi+6,
 \label{eq:continuo-comun-update-cilindro-indice}\\
 I_{K,u}^\chi
 &:=
 \left[
  \frac{729N_K^\chi+u}{3^{L_K^\chi+6}},
  \frac{729N_K^\chi+u+1}{3^{L_K^\chi+6}}
 \right),
 &0&\leq u<729.
 \label{eq:continuo-comun-update-cilindro-intervalo}
\end{align}
El cilindro decimal y el acarreo asociado se actualizan por los dos casos
enteros del mismo operador \(\mathcal C_{3,1000}\). Si no se publica una
nueva tríada decimal,
\begin{equation}
\label{eq:continuo-comun-update-cilindro-sin-triada}
 r_{K+1}^\chi=r_K^\chi,\qquad D_{K+1}^\chi=D_K^\chi,\qquad
 A_{K+1}^\chi=729A_K^\chi+
 \operatorname{ind}_3(u_{K+1}^\chi)1000^{r_K^\chi},\qquad
 B_{K+1}^\chi=A_{K+1}^\chi+1000^{r_K^\chi}.
\end{equation}
La primitiva \(\mathcal C_{3,1000}\) produce el dato opción
\[
 d_{K+1,\chi}^{\rm dec}\in
 \{\bot\}\sqcup\{0,\ldots,999\};
\]
en el primer caso vale \(d_{K+1,\chi}^{\rm dec}=\bot\). Si aparece una
tríada, escribimos
\(d_{K+1,\chi}^{\rm dec}=\delta_{K+1}^\chi\in\{0,\ldots,999\}\) y
\begin{equation}
\label{eq:continuo-comun-update-cilindro-con-triada}
 \begin{aligned}
 r_{K+1}^\chi&=r_K^\chi+1,\qquad
 D_{K+1}^\chi=1000D_K^\chi+\delta_{K+1}^\chi,\\
 A_{K+1}^\chi&=729000A_K^\chi
 +\operatorname{ind}_3(u_{K+1}^\chi)1000^{r_K^\chi+1}
 -729\delta_{K+1}^\chi3^{L_K^\chi},\\
 B_{K+1}^\chi&=A_{K+1}^\chi+1000^{r_K^\chi+1}.
 \end{aligned}
\end{equation}
El registro \(\boldsymbol\delta_{K+1}\) añade esas tríadas y los dos
cociclos conjuntos
\begin{equation}
\label{eq:continuo-comun-update-triadas-carry-conjunto}
 \begin{gathered}
 q_{K+1}=(u_{K+1}^\pi+u_{K+1}^e+u_{K+1}^\varphi)\bmod3,
 \quad
 c_{K+1}^{\rm joint}=\frac{u_{K+1}^\pi+u_{K+1}^e+u_{K+1}^\varphi-q_{K+1}}3,\\
 a_{K+1}=(u_{K+1}^\pi+u_{K+1}^e-u_{K+1}^\varphi)\bmod3,
 \quad
 h_{K+1}^{\rm joint}=\frac{u_{K+1}^\pi+u_{K+1}^e-u_{K+1}^\varphi-a_{K+1}}3,\\
 \boldsymbol\delta_{K+1}
 =\boldsymbol\delta_K\star
 \bigl((d_{K+1,\chi}^{\rm dec})_\chi,
       c_{K+1}^{\rm joint},h_{K+1}^{\rm joint}\bigr).
 \end{gathered}
\end{equation}
Así \(C_{K+1}\), \(\boldsymbol\delta_{K+1}\) y
\(\boldsymbol x_{K+1}\) tienen acciones materiales, no sólo nombres.
Las demás coordenadas se actualizan simultáneamente por
\begin{align}
 p_{K+1}&=p_K\star B_{K+1},
 \label{eq:continuo-comun-update-prefijo-tpk}\\
 \kappa(\gamma^{\rm TPK}_{0:K+1})
 &=\kappa(\varepsilon^{\rm TPK}_{K+1})
   +\mathsf C_{\varepsilon^{\rm TPK}_{K+1}}
    \kappa(\gamma^{\rm TPK}_{0:K}),
 \label{eq:continuo-comun-update-carry-lector-tpk}\\
 c_{K+1}^{\rm comp}
 &:=(\boldsymbol\delta_{K+1},\boldsymbol x_{K+1},
       \kappa(\gamma^{\rm TPK}_{0:K+1})),
 \label{eq:continuo-comun-update-carry-compacto-tpk}\\
 \Sigma_{K+1}
 &=\Sigma_B\bigl(
 B_{K+1},c_{K+1}^{\rm comp},
 \partial\gamma^{\rm TPK}_{0:K}
 +\partial\varepsilon^{\rm TPK}_{K+1},
 \omega_{K+1},\ell_{K+1}\bigr),
 \label{eq:continuo-comun-update-firma-tpk}\\
 W_{K+1}
 &=\bigl((u_{K+1}^\chi,u_{K+1}^\chi A_W)\bigr)_\chi,
 &g(K+1)&=g(K)+[1]_9.
 \label{eq:continuo-comun-update-witt-fase-tpk}
\end{align}
Como \(A_H\) es invertible módulo tres y \(A_W^2=-I_6\), las actualizaciones
residual y de sombra no colapsan bloques distintos. El truncamiento conserva
las secuencias completas
\[
 (B_0,\ldots,B_K),\ (C_0,\ldots,C_K),\
 (p_0,\ldots,p_K),\ (c_0,\ldots,c_K),\
 (\Sigma_0,\ldots,\Sigma_K),\ (W_0,\ldots,W_K)
\]
y elimina simultáneamente su última entrada. Por ello el lector
\(\mathfrak r_9\) conmuta con el borrado de un bloque.

\begin{proposition}[Alcance universal y alcance certificado del no reinicio]
\label{prop:continuo-comun-no-reinicio-alcance}
En toda vuelta horizontal de nueve pasos se conservan
\[
 g(K+9)=g(K),\qquad
 q_{\rm mem}(K+9)=q_{\rm mem}(K)+1,
\]
y el prefijo se prolonga con nueve emisiones, por lo que el estado no se
reinicia. En la sección conjunta certificada
\(\chi\in\{\pi,e,\varphi\}\), \(0\leq K<387\), se cumple además
\begin{equation}
\label{eq:continuo-comun-no-reinicio-certificado}
 (p_K,C_K)\neq(p_{K+9},C_{K+9}),\qquad
 (B_K,\mathbf N_K,W_K)\neq
 (B_{K+9},\mathbf N_{K+9},W_{K+9}),
 \qquad
 \mathbf N_K:=(N_K^\pi,N_K^e,N_K^\varphi).
\end{equation}
\end{proposition}

\begin{proof}
Las dos identidades universales son la iteración exacta de
\eqref{eq:continuo-comun-accion-fase-tpk}; el prefijo conserva las nueve
entradas añadidas por \eqref{eq:continuo-comun-update-prefijo-tpk}. Para la
sección conjunta, la recurrencia
\eqref{eq:continuo-comun-update-cilindro-indice} compara los \(387\) pares
nonádicos de cada uno de los tres canales y produce
\eqref{eq:continuo-comun-no-reinicio-certificado}; ésta es la afirmación
certificada en
\cref{thm:generacion-monodromica-conjunta-rev7}. Finalmente,
\(u\mapsto(u,uA_W)\) es inyectiva, de modo que el cambio de bloque conserva
también el cambio de su sombra completa. No se extiende esta última
desigualdad a una prolongación neutra arbitraria fuera de la sección
certificada.
\end{proof}

\begin{definition}[Estado APP--TRIT--TPK de profundidad \(k\)]
\label{def:continuo-comun-estado}
Para una ruta \(\gamma=\epsilon_1\cdots\epsilon_k\) que nace en una célula
APP \(a_0\), el estado holonómico integral común es
\begin{equation}
\label{eq:continuo-comun-estado}
 \widehat x_k(\gamma):=\Bigl(
  a_0,\ldots,a_k;
  \operatorname{Sh}_0,\ldots,\operatorname{Sh}_k;
  \mathbf{res}_k,\mathbf{quo}_k;
  \mathbf{or}_k,\mathbf\kappa_k;
  \mathsf Q_k;g_k;q^{(9)}_k;\beta_k;m_k;\partial\gamma;\gamma;
  \operatorname{Led}(\gamma)
 \Bigr),
\end{equation}
donde
\begin{align*}
 \mathbf{res}_k
 &:=\bigl(
  (\operatorname{res}^{\Sigma}(\epsilon_j),
   \operatorname{res}^{\Pi}(\epsilon_j))
  \bigr)_{1\leq j\leq k},\\
 \mathbf{quo}_k
 &:=\bigl(
  (\operatorname{quo}^{\Sigma}(\epsilon_j),
   \operatorname{quo}^{\Pi}(\epsilon_j))
  \bigr)_{1\leq j\leq k},\\
 \mathbf{or}_k
 &:=\bigl(\operatorname{or}(\epsilon_j)\bigr)_{1\leq j\leq k},
 \qquad
 \mathbf\kappa_k
 :=\bigl(\kappa(\epsilon_j)\bigr)_{1\leq j\leq k},\\
 m_k&:=m(\gamma;m_0).
\end{align*}
Las hojas intermedias se obtienen por
\(\operatorname{Sh}_{j}=U_{\epsilon_j}\operatorname{Sh}_{j-1}\).
\end{definition}

La presencia simultánea de la ruta y del ledger impide identificar palabras
distintas que terminan en la misma célula.  El estado no es una lista de
cardinales: cada coordenada tiene las ecuaciones de tipo y composición de las
secciones anteriores.

\begin{definition}[Dominio total de profundidad finita]
\label{def:continuo-comun-dominio-total}
Sea
\begin{equation}
\label{eq:continuo-comun-dominio-total}
 \widehat{\mathcal X}^{\rm tot}_k
 :=\left\{
  \widehat x_k(\epsilon_1\cdots\epsilon_k):
  \begin{array}{l}
   a_0\in\mathsf A_0,\\
   \epsilon_j\in\mathsf E_{j-1}(a_{j-1}),\\
   s(\epsilon_j)=a_{j-1},\ t(\epsilon_j)=a_j
  \end{array}
 \right\}.
\end{equation}
No se añade a \eqref{eq:continuo-comun-dominio-total} ninguna condición de
una publicación posterior.  Ser un elemento del dominio significa
exactamente haber sido producido por las acciones totales
\eqref{eq:continuo-comun-residuo-cociente-app},
\eqref{eq:continuo-comun-trit-euclid} y
\eqref{eq:continuo-comun-accion-fase-tpk}--
\eqref{eq:continuo-comun-accion-memoria}.
\end{definition}

El truncamiento común borra la última emisión y todas y sólo las entradas del
ledger que ésta produjo:
\begin{equation}
\label{eq:continuo-comun-truncamiento-total}
 \rho_{k+1,k}:\widehat{\mathcal X}^{\rm tot}_{k+1}
 \longrightarrow\widehat{\mathcal X}^{\rm tot}_k,
 \qquad
 \rho_{k+1,k}\widehat x_{k+1}
 =\widehat x_k.
\end{equation}
Sea \(\varepsilon^0_k(x)\) la emisión neutra
\eqref{eq:continuo-comun-emision-neutra-app} aplicada al extremo APP de
\(x\).  La prolongación neutra define
\begin{equation}
\label{eq:continuo-comun-seccion-total}
 \eta_k:\widehat{\mathcal X}^{\rm tot}_k
 \longrightarrow\widehat{\mathcal X}^{\rm tot}_{k+1},
 \qquad
 \eta_k(x):=(\varepsilon^0_k(x))_*x,
 \qquad
 \rho_{k+1,k}\eta_k=I.
\end{equation}

\begin{proposition}[Finitud, no vacuidad y sobreyectividad]
\label{prop:continuo-comun-total-finito}
Para todo \(k\), el conjunto
\(\widehat{\mathcal X}^{\rm tot}_k\) es finito y no vacío, y
\(\rho_{k+1,k}\) es sobreyectivo.
\end{proposition}

\begin{proof}
APP produce un conjunto finito no vacío en profundidad cero.  En cada paso,
cada estado tiene un conjunto finito de emisiones por
\eqref{eq:continuo-comun-out-no-vacio-app}; por inducción, sólo hay un número
finito de palabras de longitud \(k\).  La sucesión de prolongaciones neutras
produce al menos una palabra de cada longitud.  Finalmente,
\eqref{eq:continuo-comun-seccion-total} es una inversa por la derecha del
truncamiento, de modo que éste es sobreyectivo.
\end{proof}

\section{Grupoide finito de simetrías reversibles}
\label{sec:continuo-comun-groupoide}

Las operaciones reversibles de APP, TRIT y TPK actúan sobre el estado entero.
Sea \(\mathsf U\) el grupo de familias coherentes
\(g=(g_k)_{k\geq0}\) generadas por esas operaciones, donde
\begin{equation}
\label{eq:continuo-comun-simetria-coherente}
 \rho_{k+1,k}g_{k+1}=g_k\rho_{k+1,k},
 \qquad
 g_{k+1}\eta_k=\eta_kg_k.
\end{equation}
Denotamos por \(\mathsf U_k\) su imagen en
\(\operatorname{Sym}(\widehat{\mathcal X}^{\rm tot}_k)\).  Esta imagen es
finita porque actúa sobre un conjunto finito.  La acción de \(g_k\) sobre
\begin{equation*}
 \widehat x=(a_{\leq k};\operatorname{Sh}_{\leq k};
 \mathbf{res},\mathbf{quo};\mathbf{or},\mathbf\kappa;
 \mathsf Q_k;g_k;q_k^{(9)};\beta_k;m;\partial\gamma;\gamma;
 \operatorname{Led}(\gamma))
\end{equation*}
es
\begin{equation}
\label{eq:continuo-comun-accion-reversible}
\begin{aligned}
 g_k\cdot\widehat x:=\bigl(&
 g_k a_{\leq k};
 U_{g_k}\operatorname{Sh}_{\leq k};
 \mathbf{res}(g_k\gamma),\mathbf{quo}(g_k\gamma);
 \mathbf{or}(g_k\gamma),\mathbf\kappa(g_k\gamma);\\
 &\mathsf Q(g_k\gamma);g(g_k\gamma);q^{(9)}(g_k\gamma);
 \beta(g_k\gamma);
 \mathsf U_{g_k}(m);
 (g_k)_*\partial\gamma;
 g_k\gamma;
 \operatorname{Led}(g_k\gamma)\bigr).
\end{aligned}
\end{equation}
Todos los términos del lado derecho se recalculan mediante las reglas APP,
TRIT y TPK; no se copian como etiquetas independientes.

\begin{definition}[Grupoide de acción común]
\label{def:continuo-comun-groupoide}
El grupoide finito de profundidad \(k\) es
\begin{equation}
\label{eq:continuo-comun-groupoide}
 \mathfrak G_k
 :=\mathsf U_k\ltimes\widehat{\mathcal X}^{\rm tot}_k.
\end{equation}
Sus objetos son los estados comunes.  Una flecha
\((g_k,\widehat x)\) tiene fuente \(\widehat x\), destino
\(g_k\cdot\widehat x\), identidad \((I,\widehat x)\) e inversa
\((g_k^{-1},g_k\cdot\widehat x)\).
\end{definition}

\begin{lemma}[La acción conserva todas las ecuaciones de estado]
\label{lem:continuo-comun-groupoide-conserva}
La acción \eqref{eq:continuo-comun-accion-reversible} conserva la pareja de
hojas, las divisiones APP, la orientación TRIT, el carry, el odómetro de
fase, la memoria vertical, su reducción dodecafásica, la memoria afín, la
frontera, la ruta y el ledger, en sus tipos respectivos.  Además,
\begin{equation}
\label{eq:continuo-comun-groupoide-accion}
 I\cdot\widehat x=\widehat x,
 \qquad
 (g_{2,k}g_{1,k})\cdot\widehat x
 =g_{2,k}\cdot(g_{1,k}\cdot\widehat x).
\end{equation}
\end{lemma}

\begin{proof}
Las divisiones APP son únicas, por lo que la acción sobre una emisión
determina de nuevo sus residuos y cocientes.  La descomposición equilibrada
\eqref{eq:continuo-comun-trit-euclid} es también única.  La composición del
carry es asociativa por
\cref{lem:continuo-comun-asociatividad-carry}; la acción sobre memoria lo es
por \cref{lem:continuo-comun-accion-memoria}; y la frontera es functorial por
\eqref{eq:continuo-comun-frontera-telescopica}.  Finalmente, la acción sobre
la palabra ordenada induce la acción sobre su ledger entrada por entrada.
La coherencia con el odómetro recalcula \(Q\), \(g\) y \(q^{(9)}\) a partir
de la ruta transformada y conserva
\eqref{eq:continuo-comun-accion-fase-tpk}. Estas identidades prueban
\eqref{eq:continuo-comun-groupoide-accion}.
\end{proof}

\section{Árbol finito de todas las prolongaciones}
\label{sec:continuo-comun-arbol}

Para \(x\in\widehat{\mathcal X}^{\rm tot}_k\), sea
\(\operatorname{Out}(x)\) el conjunto de todas las emisiones TPK admisibles
de longitud uno desde el extremo de \(x\).  Una emisión pertenece a
\(\operatorname{Out}(x)\) sólo cuando su acción actualiza conjuntamente todas
las coordenadas de \eqref{eq:continuo-comun-estado}.  En particular,
\begin{equation}
\label{eq:continuo-comun-out-finito}
 1\leq d(x):=\#\operatorname{Out}(x)<\infty,
 \qquad
 \varepsilon^0_k(x)\in\operatorname{Out}(x).
\end{equation}

\begin{definition}[Árbol truncado de prolongaciones]
\label{def:continuo-comun-arbol-finito}
Dados \(N\geq k\) y
\(x\in\widehat{\mathcal X}^{\rm tot}_k\), definimos
\(\mathscr T_{k,N}(x)\) como el árbol enraizado cuyos vértices de altura
\(j\), \(0\leq j\leq N-k\), son las palabras
\begin{equation}
\label{eq:continuo-comun-vertice-arbol}
 h=(\epsilon_{k+1},\ldots,\epsilon_{k+j}),
 \qquad
 \epsilon_{n+1}\in\operatorname{Out}(x_n),
 \qquad
 x_{n+1}=(\epsilon_{n+1})_*x_n,
\end{equation}
con raíz la palabra vacía.  Una arista añade una emisión al final.  Dos
palabras que alcanzan el mismo estado terminal siguen siendo vértices
distintos si sus ledgers difieren.
\end{definition}

La última condición es esencial: el árbol registra genealogías y no el
cociente que sólo recuerda el destino.  Su nivel terminal es
\begin{equation}
\label{eq:continuo-comun-terminal-finito}
 \operatorname{Term}_{k,N}(x)
 :=\{h\in\mathscr T_{k,N}(x):|h|=N-k\}.
\end{equation}
Cada \(h\) conserva el estado terminal completo
\begin{equation}
\label{eq:continuo-comun-registro-terminal}
 \operatorname{ter}_{N}(h)
 :=\bigl(x_N,m_N,\partial(\gamma h),
          \gamma h,\operatorname{Led}(\gamma h)\bigr).
\end{equation}

\begin{proposition}[Finitud y no vacuidad a todo horizonte]
\label{prop:continuo-comun-arbol-finito-no-vacio}
Para todo \(x\), \(k\) y \(N\geq k\), el árbol
\(\mathscr T_{k,N}(x)\) es finito y
\(\operatorname{Term}_{k,N}(x)\neq\varnothing\).
\end{proposition}

\begin{proof}
La raíz tiene un número finito y positivo de sucesores por
\eqref{eq:continuo-comun-out-finito}.  Aplicando el mismo argumento en cada
nivel se obtiene, por inducción, finitud a todo horizonte finito.  La palabra
formada por emisiones neutras
\((\varepsilon^0_k,\varepsilon^0_{k+1},\ldots,
\varepsilon^0_{N-1})\) es un vértice terminal, por lo que
el nivel \eqref{eq:continuo-comun-terminal-finito} no es vacío.
\end{proof}

El borrado de la última emisión define
\begin{equation}
\label{eq:continuo-comun-proyeccion-terminal}
 \pi_{N+1,N}:\operatorname{Term}_{k,N+1}(x)
 \longrightarrow\operatorname{Term}_{k,N}(x).
\end{equation}

\begin{lemma}[Sobreyectividad terminal]
\label{lem:continuo-comun-terminal-sobreyectivo}
La aplicación \(\pi_{N+1,N}\) es sobreyectiva para todo \(N\geq k\).
\end{lemma}

\begin{proof}
Sea \(h\in\operatorname{Term}_{k,N}(x)\), de estado terminal \(v_h\).
Añadir la emisión \(\varepsilon_N^0(v_h)\) produce
\(h'\in\operatorname{Term}_{k,N+1}(x)\) y
\(\pi_{N+1,N}(h')=h\).
\end{proof}

\section{Conexión nonádica: fase, holonomía y memoria}
\label{sec:continuo-comun-gamma9-total}

Una palabra de nueve emisiones no se denomina holonomía sólo por tener
longitud nueve. Primero se conserva la conexión horizontal determinada por
el funtor de fase. Para \(r\in C_9\), definimos
\begin{equation}
\label{eq:continuo-comun-relacion-horizontal-fase}
 \begin{split}
 \mathscr R_{r,k}:=\bigl\{
 (x,\epsilon,\epsilon_*x):\;&
 x\in\widehat{\mathcal X}^{\rm tot}_k,\quad
 \epsilon\in\operatorname{Out}(x),\\
 &\mathsf Q_k(x)=(w,r),\quad
 \mathsf Q_{k+1}(\epsilon_*x)=\sigma_9(w,r)
 \bigr\}.
 \end{split}
\end{equation}
La última igualdad no recorta el dominio: es la acción total
\eqref{eq:continuo-comun-accion-fase-tpk} de toda emisión TPK elemental.
La conexión transporta simultáneamente hoja, memoria y ledger:
\begin{equation}
\label{eq:continuo-comun-conexion-tpk-generador}
 \nabla_{\rm TPK}(\epsilon):=
 \bigl(\operatorname{ph}(\epsilon),U_\epsilon,
       \mathsf U_\epsilon,\lambda(\epsilon)\bigr).
\end{equation}
Su composición se calcula mediante
\eqref{eq:continuo-comun-transporte-hojas-funtorial},
\eqref{eq:continuo-comun-accion-memoria-composicion} y la concatenación del
ledger.

El ápice de la vuelta nonádica es ahora el conjunto de lifts horizontales
fase-compatibles
\begin{equation}
\label{eq:continuo-comun-gamma9-apice}
 \begin{split}
 \mathsf Z^{\Gamma}_{9,k}:=\coprod_{r\in C_9}\bigl\{
 (x,\epsilon_1,\ldots,\epsilon_9):\;&x_0=x,\quad
 x_i=(\epsilon_i)_*x_{i-1},\\
 &(x_{i-1},\epsilon_i,x_i)
 \in\mathscr R_{r+i-1,k+i-1},\quad 1\leq i\leq9
 \bigr\},
 \end{split}
\end{equation}
 donde los subíndices de \(r+i-1\) se leen en \(C_9\). El sumando está
determinado de forma única por \(r=g_k(x)\), de modo que el coproducto no
duplica testigos. Los mapas fuente y
destino son
\begin{equation}
\label{eq:continuo-comun-gamma9-span}
 \widehat{\mathcal X}^{\rm tot}_k
 \xleftarrow{\ s_{9,k}\ }
 \mathsf Z^{\Gamma}_{9,k}
 \xrightarrow{\ t_{9,k}\ }
 \widehat{\mathcal X}^{\rm tot}_{k+9},
 \qquad
 s_{9,k}(x,\boldsymbol\epsilon)=x,\quad
 t_{9,k}(x,\boldsymbol\epsilon)
 =(\epsilon_9\cdots\epsilon_1)_*x.
\end{equation}
La holonomía relacional es
\begin{equation}
\label{eq:continuo-comun-gamma9-relacion}
 \Gamma_{9,k}:=
 \{(s_{9,k}z,t_{9,k}z):z\in\mathsf Z^{\Gamma}_{9,k}\}
 =\operatorname{Hol}_{C_9}(\nabla_{\rm TPK}).
\end{equation}

\begin{lemma}[El nivel nueve del árbol realiza la holonomía]
\label{lem:continuo-comun-gamma9-arbol}
La aplicación que olvida los estados intermedios ya determinados por las
emisiones induce una biyección
\begin{equation}
\label{eq:continuo-comun-gamma9-apice-arbol}
 \mathsf Z^{\Gamma}_{9,k}
 \xrightarrow{\ \cong\ }
 \coprod_{x\in\widehat{\mathcal X}^{\rm tot}_k}
 \operatorname{Term}_{k,k+9}(x).
\end{equation}
\end{lemma}

\begin{proof}
Cada arista de \(\operatorname{Out}(x_i)\) satisface
\eqref{eq:continuo-comun-accion-fase-tpk}; por ello una palabra terminal de
longitud nueve determina los nueve elementos de
\eqref{eq:continuo-comun-relacion-horizontal-fase}. Recíprocamente, un lift
horizontal de \eqref{eq:continuo-comun-gamma9-apice} es una palabra
componible del árbol. Las dos operaciones conservan las emisiones y son
inversas. Así \eqref{eq:continuo-comun-gamma9-apice-arbol} es un teorema
derivado de la conexión, no la definición de ésta.
\end{proof}

\begin{proposition}[Retorno de fase, avance de memoria y no reinicio]
\label{prop:continuo-comun-gamma9-retorno-memoria}
Para todo \(z=(x,\epsilon_1,\ldots,\epsilon_9)\) del ápice,
\begin{align}
 g(t_{9,k}z)&=g(s_{9,k}z),
 \label{eq:continuo-comun-gamma9-retorno-fase}\\
 q^{(9)}(t_{9,k}z)&=q^{(9)}(s_{9,k}z)+1,
 \label{eq:continuo-comun-gamma9-avance-memoria}\\
 \operatorname{Led}(t_{9,k}z)
 &=\operatorname{Led}(s_{9,k}z)\star
   \bigl(\lambda(\epsilon_1),\ldots,\lambda(\epsilon_9)\bigr).
 \label{eq:continuo-comun-gamma9-ledger}
\end{align}
En particular, la identificación de fase
\(g(t_{9,k}z)=g(s_{9,k}z)\) no induce una identificación de estados: tras
transportar ambos registros al mismo tipo de comparación, sus memorias
verticales difieren en una unidad. Retorna la fase publicada, no el estado
holonómico integral.
\end{proposition}

\begin{proof}
Las dos primeras identidades son
\eqref{eq:continuo-comun-retorno-fase-avance-memoria}. La tercera es la
definición de concatenación del ledger. Como la memoria vertical difiere en
una unidad, los estados completos son distintos aunque alguna coordenada
visible adicional pudiera coincidir.
\end{proof}

\begin{proposition}[Compresión isométrica de la misma holonomía]
\label{prop:continuo-comun-gamma9-compresion-isometrica}
Sea \(U_{\Gamma_9}:\mathcal H_{\rm full}\to\mathcal H_{\rm full}\) una
realización isométrica de la holonomía anterior y sea
\(J:\mathcal H_{\rm vis}\to\mathcal H_{\rm full}\) la inclusión isométrica
del registro visible, con \(E:=J^*\) y \(P:=JE\). Definimos
\begin{equation}
\label{eq:continuo-comun-gamma9-compresion-datos}
 T:=EU_{\Gamma_9}J,\qquad
 \eta:=(I-P)U_{\Gamma_9}J.
\end{equation}
Entonces
\begin{equation}
\label{eq:continuo-comun-gamma9-compresion}
 \boxed{\;
 U_{\Gamma_9}J=JT+\eta,\qquad
 I-T^*T=\eta^*\eta.
 \;}
\end{equation}
\end{proposition}

\begin{proof}
Insertando \(I=P+(I-P)\) delante de \(U_{\Gamma_9}J\) se obtiene
\[
 U_{\Gamma_9}J
 =JEU_{\Gamma_9}J+(I-P)U_{\Gamma_9}J
 =JT+\eta.
\]
Los dos sumandos pertenecen a subespacios ortogonales. Como
\(U_{\Gamma_9}\) y \(J\) son isometrías,
\[
 I=(U_{\Gamma_9}J)^*(U_{\Gamma_9}J)
   =T^*T+\eta^*\eta,
\]
que es la segunda identidad.
\end{proof}

\begin{proposition}[Totalidad y equivariancia de la holonomía]
\label{prop:continuo-comun-gamma9-total-natural}
El mapa \(s_{9,k}\) es sobreyectivo. Las nueve prolongaciones neutras forman
la sección
\begin{equation}
\label{eq:continuo-comun-gamma9-seccion-neutra}
 \begin{gathered}
 x_0:=x,\qquad
 \varepsilon^0_{k+i}:=\varepsilon^0_{k+i}(x_i),\qquad
 x_{i+1}:=(\varepsilon^0_{k+i})_*x_i\quad(0\leq i<9),\\
 \eta^{(9)}_k(x):=
 (x,\varepsilon^0_k,\ldots,\varepsilon^0_{k+8})
 \in\mathsf Z^{\Gamma}_{9,k},
 \qquad s_{9,k}\eta^{(9)}_k=I.
 \end{gathered}
\end{equation}
Toda simetría coherente \(a\) cuya acción en el odómetro satisface
\(\mathsf Q\,a=\bar a\,\mathsf Q\) y
\(\bar a\,\sigma_9=\sigma_9\,\bar a\) induce
\[
 a^Z(x,\epsilon_1,\ldots,\epsilon_9)
 :=(a x,a\epsilon_1,\ldots,a\epsilon_9)
\]
y verifica
\begin{equation}
\label{eq:continuo-comun-gamma9-naturalidad}
 s_{9,k}a^Z=a_ks_{9,k},
 \qquad
 t_{9,k}a^Z=a_{k+9}t_{9,k}.
\end{equation}
\end{proposition}

\begin{proof}
La palabra neutra existe en toda fuente y satisface la acción de fase
\eqref{eq:continuo-comun-accion-fase-tpk}; esto prueba la sección y la
sobreyectividad. La condición
\(\bar a\sigma_9=\sigma_9\bar a\) envía relaciones horizontales a relaciones
horizontales. Actuar entrada por entrada conserva composición, carry,
memoria, frontera y ledger; leer el primer o el último estado antes o después
de actuar da \eqref{eq:continuo-comun-gamma9-naturalidad}.
\end{proof}

\section{Supervivencia total y multisección completa}
\label{sec:continuo-comun-supervivencia}

\begin{definition}[Supervivencia producida]
\label{def:continuo-comun-supervivencia}
Un estado \(x\in\widehat{\mathcal X}^{\rm tot}_k\) sobrevive cuando
\begin{equation}
\label{eq:continuo-comun-supervivencia}
 \operatorname{Term}_{k,N}(x)\neq\varnothing
 \qquad\text{para todo }N\geq k.
\end{equation}
Denotamos por \(\operatorname{Surv}_k\) el conjunto de tales estados.
\end{definition}

\begin{theorem}[El dominio total es superviviente]
\label{thm:continuo-comun-total-superviviente}
Para toda profundidad,
\begin{equation}
\label{eq:continuo-comun-surv-igual-total}
 \operatorname{Surv}_k=\widehat{\mathcal X}^{\rm tot}_k.
\end{equation}
En particular, la supervivencia no define un subconjunto supuesto antes de
construir el estado.
\end{theorem}

\begin{proof}
La inclusión \(\operatorname{Surv}_k\subseteq
\widehat{\mathcal X}^{\rm tot}_k\) pertenece a la definición.  La inclusión
opuesta es exactamente la no vacuidad a todo horizonte probada en
\cref{prop:continuo-comun-arbol-finito-no-vacio}.
\end{proof}

\begin{definition}[Multisección de prolongaciones completas]
\label{def:continuo-comun-multiseccion}
La multisección completa sobre \(x\) es el límite inverso
\begin{equation}
\label{eq:continuo-comun-multiseccion}
 \operatorname{Msect}(x)
 :=\varprojlim_{N\geq k}
 \bigl(\operatorname{Term}_{k,N}(x),\pi_{N+1,N}\bigr).
\end{equation}
Sus elementos son todas las historias infinitas compatibles que prolongan
\(x\).  Ninguna de ellas recibe prioridad dentro de
\eqref{eq:continuo-comun-multiseccion}.
\end{definition}

\begin{theorem}[No vacuidad material de la multisección]
\label{thm:continuo-comun-multiseccion-no-vacia}
Para todo \(x\in\widehat{\mathcal X}^{\rm tot}_k\),
\begin{equation}
\label{eq:continuo-comun-multiseccion-no-vacia}
 \operatorname{Msect}(x)\neq\varnothing.
\end{equation}
Además, cada coordenada finita de una historia de
\(\operatorname{Msect}(x)\) conserva las dos hojas, orientación, carry,
memoria, frontera, ruta y ledger del estado que prolonga.
\end{theorem}

\begin{proof}
Los conjuntos terminales son finitos y no vacíos por
\cref{prop:continuo-comun-arbol-finito-no-vacio}, y sus aplicaciones de
enlace son sobreyectivas por
\cref{lem:continuo-comun-terminal-sobreyectivo}.  Se puede construir un
elemento compatible recursivamente: se comienza en la raíz y, en cada paso,
se conserva una prolongación del prefijo anterior.  La prolongación neutra
garantiza que el proceso nunca se detenga.  La compatibilidad de las
coordenadas se sigue de que cada arista del árbol actúa mediante la única
acción de estado definida en
\eqref{eq:continuo-comun-estado}.
\end{proof}

\section{Medida canónica generada por el árbol}
\label{sec:continuo-comun-medida}

La medida no se introduce como un peso exterior sobre las salidas.  Se obtiene
contando las emisiones admisibles en cada vértice antes de elegir una
continuación.

\begin{definition}[Peso local y medida de horizonte]
\label{def:continuo-comun-medida-finita}
Para \(v\in\widehat{\mathcal X}^{\rm tot}_j\) y
\(\epsilon\in\operatorname{Out}(v)\), ponemos
\begin{equation}
\label{eq:continuo-comun-peso-local}
 p_v(\epsilon):=\frac{1}{d(v)}\in\mathbb Q_{>0}.
\end{equation}
Si
\(h=(\epsilon_{k+1},\ldots,\epsilon_N)\in
\operatorname{Term}_{k,N}(x)\) y \(x_j\) es el prefijo que precede a
\(\epsilon_{j+1}\), definimos
\begin{equation}
\label{eq:continuo-comun-medida-horizonte}
 \mu_{x,N}(h)
 :=\prod_{j=k}^{N-1}p_{x_j}(\epsilon_{j+1})
 =\prod_{j=k}^{N-1}\frac{1}{d(x_j)}.
\end{equation}
\end{definition}

\begin{lemma}[Normalización]
\label{lem:continuo-comun-medida-normalizada}
Para todo horizonte \(N\geq k\),
\begin{equation}
\label{eq:continuo-comun-medida-suma-uno}
 \sum_{h\in\operatorname{Term}_{k,N}(x)}\mu_{x,N}(h)=1.
\end{equation}
\end{lemma}

\begin{proof}
En altura uno,
\(\sum_{\epsilon\in\operatorname{Out}(x)}d(x)^{-1}=1\).  Supuesta la
normalización a altura \(N-k\), cada prefijo terminal \(h\), de estado final
\(v_h\), aporta en el nivel siguiente
\[
 \sum_{\epsilon\in\operatorname{Out}(v_h)}
 \mu_{x,N}(h)p_{v_h}(\epsilon)
 =\mu_{x,N}(h).
\]
Sumar sobre todos los prefijos prueba la identidad por inducción.
\end{proof}

\begin{proposition}[Consistencia proyectiva]
\label{prop:continuo-comun-medida-consistente}
Las medidas finitas satisfacen
\begin{equation}
\label{eq:continuo-comun-medida-pushforward}
 (\pi_{N+1,N})_*\mu_{x,N+1}=\mu_{x,N}.
\end{equation}
\end{proposition}

\begin{proof}
Para \(h\in\operatorname{Term}_{k,N}(x)\), la masa de su fibra es
\begin{align*}
 \sum_{\pi_{N+1,N}(h')=h}\mu_{x,N+1}(h')
 &=\mu_{x,N}(h)
   \sum_{\epsilon\in\operatorname{Out}(v_h)}p_{v_h}(\epsilon)\\
 &=\mu_{x,N}(h).
\end{align*}
Ésta es precisamente la igualdad de medidas puntuales que equivale a
\eqref{eq:continuo-comun-medida-pushforward}.
\end{proof}

Los cilindros de la multisección son
\begin{equation}
\label{eq:continuo-comun-cilindro-msect}
 [h]_{x,N}
 :=\{\omega\in\operatorname{Msect}(x):
       \omega|_N=h\}.
\end{equation}

\begin{theorem}[Medida canónica de prolongaciones]
\label{thm:continuo-comun-medida-canonica}
Existe una única medida de cilindros \(\mu_x\) sobre
\(\operatorname{Msect}(x)\) tal que
\begin{equation}
\label{eq:continuo-comun-medida-cilindro}
 \mu_x([h]_{x,N})=\mu_{x,N}(h).
\end{equation}
Está normalizada, tiene soporte en toda la multisección y se obtiene sólo de
los conjuntos finitos de prolongaciones APP--TRIT--TPK.
\end{theorem}

\begin{proof}
La consistencia de
\cref{prop:continuo-comun-medida-consistente} hace que
\eqref{eq:continuo-comun-medida-cilindro} sea independiente del horizonte en
el que se represente un cilindro.  La normalización procede de
\cref{lem:continuo-comun-medida-normalizada}.  Dos cilindros o bien son
disjuntos o bien uno contiene al otro; por tanto, las masas finitas definen
de manera única una medida aditiva en el álgebra de cilindros.  La
multisección es un límite de conjuntos finitos discretos y los cilindros
forman su base compacta; la medida se prolonga de manera única a la
\(\sigma\)-álgebra generada por ellos.  Finalmente,
\(p_v(\epsilon)>0\) para toda emisión admisible, de modo que todo cilindro no
vacío tiene masa positiva y el soporte es la multisección completa.
\end{proof}

\section{Naturalidad bajo simetrías y truncamientos}
\label{sec:continuo-comun-naturalidad}

Una familia coherente \(g=(g_j)_{j\geq0}\in\mathsf U\) transporta una
emisión saliente de profundidad \(j\) mediante
\begin{equation}
\label{eq:continuo-comun-accion-out}
 (g_{j+1},g_j)_*:\operatorname{Out}(x)\longrightarrow
      \operatorname{Out}(g_j\cdot x),
 \qquad
 \epsilon\longmapsto g_{j+1}\epsilon g_j^{-1}.
\end{equation}
Las relaciones TPK hacen de \eqref{eq:continuo-comun-accion-out} una biyección
que conserva multiplicidades.  Actuando en cada entrada se obtiene un
isomorfismo de árboles
\begin{equation}
\label{eq:continuo-comun-isomorfismo-arbol}
 g_{*,N}:\mathscr T_{k,N}(x)
 \xrightarrow{\ \cong\ }
 \mathscr T_{k,N}(g_k\cdot x).
\end{equation}

\begin{theorem}[Naturalidad del árbol, la multisección y la medida]
\label{thm:continuo-comun-naturalidad}
Para toda simetría reversible coherente \(g=(g_j)_{j\geq0}\),
\begin{align}
 \pi_{N+1,N}g_{*,N+1}
 &=g_{*,N}\pi_{N+1,N},
 \label{eq:continuo-comun-naturalidad-terminal}\\
 g_*\operatorname{Msect}(x)
 &=\operatorname{Msect}(g_k\cdot x),
 \label{eq:continuo-comun-naturalidad-msect}\\
 (g_{*,N})_*\mu_{x,N}
 &=\mu_{g_k\cdot x,N},
 \qquad
 (g_*)_*\mu_x=\mu_{g_k\cdot x}.
 \label{eq:continuo-comun-naturalidad-medida}
\end{align}
\end{theorem}

\begin{proof}
La primera igualdad expresa que actuar entrada por entrada y borrar la última
entrada conmutan.  Al pasar al límite inverso se obtiene
\eqref{eq:continuo-comun-naturalidad-msect}.  Como
\eqref{eq:continuo-comun-accion-out} es biyectiva,
\(d(g_j\cdot v)=d(v)\); por
\eqref{eq:continuo-comun-medida-horizonte}, una palabra y su imagen tienen el
mismo peso.  Esto prueba la primera igualdad de medidas de
\eqref{eq:continuo-comun-naturalidad-medida}.  La segunda se deduce en los
cilindros mediante \eqref{eq:continuo-comun-medida-cilindro}, y los cilindros
determinan la medida.
\end{proof}

Los truncamientos del estado y los del árbol forman asimismo el cuadrado
\begin{equation}
\label{eq:continuo-comun-cuadrado-truncamiento}
\begin{CD}
 \operatorname{Term}_{k,N+1}(x)
 @>{\pi_{N+1,N}}>>
 \operatorname{Term}_{k,N}(x)\\
 @V{\operatorname{ter}_{N+1}}VV
 @VV{\operatorname{ter}_{N}}V\\
 \widehat{\mathcal X}^{\rm tot}_{N+1}
 @>{\rho_{N+1,N}}>>
 \widehat{\mathcal X}^{\rm tot}_{N}.
\end{CD}
\end{equation}
La conmutatividad es literal: ambos recorridos borran la última emisión y su
última entrada de ledger.

\section{Límite inverso \(\widehat\Omega\) de historias completas y estado holonómico integral común APP–TRIT–TPK}
\label{sec:continuo-comun-limite}

Las sobreyecciones \eqref{eq:continuo-comun-truncamiento-total} producen el
objeto de historias completas
\begin{equation}
\label{eq:continuo-comun-omega}
 \widehat\Omega
 :=\varprojlim_k
 \bigl(\widehat{\mathcal X}^{\rm tot}_k,\rho_{k+1,k}\bigr).
\end{equation}
Un elemento \(\widehat x=(\widehat x_k)_{k\geq0}\) contiene, en cada
profundidad, la misma pareja de hojas transportada, las orientaciones y
acarreos producidos, la memoria acumulada, las fronteras, la ruta y el ledger
completo.  Su fibra de continuaciones desde cualquier prefijo es
\(\operatorname{Msect}(\widehat x_k)\), equipada con
\(\mu_{\widehat x_k}\).

\begin{theorem}[Estado holonómico integral común APP--TRIT--TPK]
\label{thm:continuo-comun-fuente teoremática}
El sistema
\begin{equation}
\label{eq:continuo-comun-fuente teoremática}
 \mathfrak C^{\rm gen}:=\Biggl(
 \begin{aligned}
 &\{\widehat{\mathcal X}^{\rm tot}_k,\rho_{k+1,k},\eta_k\}_{k\geq0},
   \widehat\Omega,\{\mathfrak G_k\}_{k\geq0},
   \{\mathscr T_{k,N}\}_{N\geq k},\\
 &\operatorname{Msect},\operatorname{Surv},\{\mu_x\},
   \{\mathfrak r_9(v_K)\}_{K\geq0},
   \mathcal K_{\rm ph},\widehat{\mathcal K}_{\rm ph},\\
 &\operatorname{Sh},\operatorname{or},\kappa,
   \mathsf Q,g,q^{(9)},\beta,m,\partial,\gamma,\operatorname{Led},\\
 &\{\mathsf Z^\Gamma_{9,k},s_{9,k},t_{9,k},\Gamma_{9,k}\}_{k\geq0}
 \end{aligned}
 \Biggr)
\end{equation}
está definido para todo estado producido por APP--TRIT--TPK\@.  Es no vacío,
natural bajo las simetrías reversibles y compatible con refinamiento.  En
particular:
\begin{enumerate}
 \item las dos hojas se transportan juntas;
 \item orientación y carry se calculan por la división TRIT;
 \item la memoria obedece la acción afín TPK;
 \item la fase obedece el odómetro TPK, vuelve tras nueve emisiones y la
       memoria vertical avanza una unidad;
 \item la realización Markov de esa vuelta satisface
       \(\mathcal K_{\rm ph}^9=E_+E_\times\) y renueva uniformemente la
       fibra completa de \(468=18\cdot26\) emisiones;
 \item el lector compacto conserva bloque, cilindro, prefijo, carry, firma,
       sombra de Witt y fase, con sus actualizaciones y truncamientos;
 \item frontera, ruta y ledger se componen sin borrar los extremos ni el
       orden de las emisiones;
 \item todo estado sobrevive y posee una multisección no vacía de
       prolongaciones completas;
 \item la medida canónica es proyectiva y tiene soporte en todas esas
       prolongaciones.
\end{enumerate}
\end{theorem}

\begin{proof}
La totalidad y la finitud proceden de
\cref{prop:continuo-comun-total-finito}.  Las identidades de hojas, carry y
memoria están probadas en
\cref{lem:continuo-comun-asociatividad-carry,lem:continuo-comun-accion-memoria}; la fase, la memoria vertical y la
holonomía nonádica, incluida su realización de fibra y su lector de estado,
están probadas en
\cref{prop:continuo-comun-gamma9-retorno-memoria,prop:continuo-comun-gamma9-total-natural,prop:continuo-comun-kph-heredada,prop:continuo-comun-gamma9-compresion-isometrica,prop:continuo-comun-no-reinicio-alcance}; la frontera conserva sus extremos por
\eqref{eq:continuo-comun-frontera-telescopica}.  La supervivencia total y la
no vacuidad de las multisecciones son
\cref{thm:continuo-comun-total-superviviente,thm:continuo-comun-multiseccion-no-vacia}.  La existencia, consistencia y
soporte de la medida son
\cref{prop:continuo-comun-medida-consistente,thm:continuo-comun-medida-canonica}.  Finalmente, la compatibilidad con las
simetrías y con los truncamientos se demuestra en
\cref{thm:continuo-comun-naturalidad} y
\eqref{eq:continuo-comun-cuadrado-truncamiento}.
\end{proof}

\section{Necesidad estructural y procedencia}
\label{sec:continuo-comun-falsadores}

\begin{proposition}[Falsadores del estado holonómico integral común]
\label{prop:continuo-comun-falsadores}
Cualquiera de las operaciones siguientes sustituye
\(\mathfrak C^{\rm gen}\) por otro objeto y bloquea la transferencia de los
teoremas anteriores:
\begin{enumerate}
 \item elegir \(\Sigma\) o \(\Pi\) y descartar la otra hoja;
 \item conservar el residuo y borrar el cociente;
 \item conservar el signo TRIT y borrar su carry;
 \item reemplazar el ledger ordenado por su valor terminal;
 \item identificar dos palabras porque alcanzan la misma célula;
 \item escoger una continuación y sustituir por ella la multisección;
 \item declarar supervivencia sin exhibir prolongaciones a cada horizonte;
 \item normalizar uniformemente sólo el último nivel, ignorando los grados de
       los prefijos y perdiendo la consistencia proyectiva;
 \item usar una lectura posterior para decidir qué estados entran en
       \(\widehat{\mathcal X}^{\rm tot}_k\).
\end{enumerate}
\end{proposition}

\begin{proof}
Los seis primeros puntos eliminan coordenadas o relaciones que aparecen en la
definición de estado, árbol o multisección.  El séptimo elimina la prueba de
\cref{thm:continuo-comun-total-superviviente}.  El octavo puede violar
\eqref{eq:continuo-comun-medida-pushforward} cuando los grados terminales no
son constantes.  El último cambia el dominio total
\eqref{eq:continuo-comun-dominio-total} por una selección retrospectiva.
\end{proof}

Las células APP, la pareja de hojas, el levantamiento TRIT, las extensiones
TPK, el cociclo de acarreo, la memoria, la frontera, la ruta y la conservación
del registro constituyen las estructuras de partida. Sobre ellas se reúnen en
un único dominio el grupoide finito, el árbol de palabras, la multisección y
la medida local normalizada. Los teoremas de supervivencia total, no vacuidad,
consistencia y naturalidad certifican esta construcción conjunta.
 \let\section\HMTContinuoCommonOriginalSection
\chapter{Operaciones intrínsecas correlativas del mismo continuo}
\label{chap:pdfv2-operaciones-intrinsecas-correlativas}

\begin{thesisbox}
Los símbolos \(\mathrm{sp},\mathrm{sol},\mathrm{gau},\mathrm{coh}\) y
\(\mathrm{det}\) que aparecen en este capítulo marcan cinco operaciones
correlativas sobre un único estado APP--TRIT--TPK\@. Cada operación recibe
la misma extensión superviviente, la
misma orientación, el mismo carry, la misma memoria, la misma frontera y la
misma ruta. Los apartados siguientes son etapas matemáticas sucesivas de una
sola construcción y no exposiciones autónomas.
\end{thesisbox}

\section{Arista generadora y estado genealógico común}

No repetimos aquí el estado, el árbol ni la medida construidos en el capítulo
anterior. Tomamos
\(\widehat x=\widehat x_k(\gamma)\in
\widehat{\mathcal X}^{\rm tot}_k\) de
\eqref{eq:continuo-comun-estado} y todas sus emisiones admisibles
\(\alpha\in\operatorname{Out}(\widehat x)\). Si
\(\widehat y=\alpha_*\widehat x\), escribimos
\begin{equation}
\label{eq:pdfv2-cor-hijos-supervivientes}
 \operatorname{Ch}_k(\widehat x):=
 \{(\alpha,\alpha_*\widehat x):
      \alpha\in\operatorname{Out}(\widehat x)\},
 \qquad
 d_k(\widehat x):=\#\operatorname{Ch}_k(\widehat x).
\end{equation}
La finitud y la no vacuidad de este conjunto son ya
\eqref{eq:continuo-comun-out-finito}; la supervivencia a todos los horizontes,
la multisección completa y su medida son
\cref{thm:continuo-comun-total-superviviente,thm:continuo-comun-multiseccion-no-vacia,thm:continuo-comun-medida-canonica}. No se vuelve a suponer ni a probar una
segunda no vacuidad.

Cada par \((\alpha,\widehat y)\) de
\eqref{eq:pdfv2-cor-hijos-supervivientes} conserva el registro único
\begin{equation}
\label{eq:pdfv2-cor-registro-extension}
 \begin{aligned}
 \mathfrak e_\alpha:=\bigl(&
 \alpha,\operatorname{Ext}_{\gamma_\alpha};
 \operatorname{Sh}_{s(\alpha)},\operatorname{Sh}_{t(\alpha)};\\
 &
 \operatorname{res}^{\Sigma,\Pi}(\alpha),
 \operatorname{quo}^{\Sigma,\Pi}(\alpha);\\
 &\operatorname{or}(\alpha),\kappa(\alpha),\mathsf C_\alpha;
 \mathsf Q(s(\alpha)),\mathsf Q(t(\alpha)),g(s(\alpha)),g(t(\alpha)),
 q^{(9)}(s(\alpha)),q^{(9)}(t(\alpha));\\
 &
 m(\gamma\alpha;m_0),\partial\alpha,\gamma\alpha,
 \operatorname{Led}(\gamma\alpha);\\
 &\{\operatorname{Term}_{k+1,N}(\widehat y)\}_{N\geq k+1},
 \operatorname{Msect}(\widehat y),\mu_{\widehat y}\bigr).
 \end{aligned}
\end{equation}
Aquí \(\operatorname{Ext}_{\gamma_\alpha}\) nombra la relación completa que
la emisión admisible ya satisface; no introduce otro filtro. Las operaciones
posteriores son aplicaciones coordinadas de \(\mathfrak e_\alpha\) y no
reconstruyen el registro desde sus salidas.

\section{Acciones intrínsecas correlativas sobre el generador común}

Las marcas internas se fijan sin cambiar de fuente:
\(\mathrm{sp}\) designa transporte y graduación de las dos hojas;
\(\mathrm{sol}\), incremento firmado, balance y memoria;
\(\mathrm{gau}\), incidencia y curvatura de carry;
\(\mathrm{coh}\), puertos, diferencia rectangular y complejo celular; y
\(\mathrm{det}\), complejo basado, relleno correlativo y línea determinante. Las fórmulas
siguientes muestran sus acoplamientos sobre el mismo
\(\mathfrak e_\alpha\); sus acciones comparten literalmente fuente, destino,
ruta y ledger.

La doble evaluación APP proporciona, sobre el mismo soporte, las etiquetas
\(\Sigma\) y \(\Pi\). Para una extensión \(\alpha\), sus multiplicidades
se calculan del ledger completo, sin un contador exterior:
\begin{equation}
\label{eq:pdfv2-cor-multiplicidades-hojas}
 N_\diamond(\alpha):=
 \sum_{\epsilon\ {\rm en}\ \gamma\alpha}
 \mathbf1_{\{\operatorname{res}^{\diamond}(\epsilon)\neq0\}},
 \qquad \diamond\in\{\Sigma,\Pi\}.
\end{equation}
Sea además el coeficiente orientado total
\[
 \Delta_{\gamma_\alpha}:=
 \sum_{\epsilon\ {\rm en}\ \gamma\alpha}\operatorname{or}(\epsilon)
 \in\mathbb Z.
\]
La acción firmada es
\begin{equation}
\label{eq:pdfv2-cor-b-alpha}
 b_\alpha:=\Delta_{\gamma_\alpha}
 \bigl(N_\Sigma(\alpha)-N_\Pi(\alpha)\bigr),
 \qquad
 b_\alpha^\pm:=\frac{|b_\alpha|\pm b_\alpha}{2}.
\end{equation}
Por construcción,
\begin{equation}
\label{eq:pdfv2-cor-b-identidades}
 b_\alpha=b_\alpha^+-b_\alpha^-,\qquad
 |b_\alpha|=b_\alpha^++b_\alpha^-,\qquad
 b_\alpha^+b_\alpha^-=0.
\end{equation}

El mismo conjunto de extensiones produce simultáneamente los puertos
\begin{equation}
\label{eq:pdfv2-cor-puertos-comunes}
 I_k(\widehat x_k):=
 \operatorname{Ch}_k(\widehat x_k)\times\{\Sigma\},
 \qquad
 J_k(\widehat x_k):=
 \operatorname{Ch}_k(\widehat x_k)\times\{\Pi\}.
\end{equation}
No se supone que una hoja visible contenga más prolongaciones que la otra:
cada extensión superviviente se evalúa en las dos hojas del mismo soporte.
Por \eqref{eq:continuo-comun-out-finito}, ambos conjuntos de puertos son no
vacíos.

Para \(A_N=\mathbb Z/3^N\mathbb Z\), la acción celular sobre generadores es
\begin{equation}
\label{eq:pdfv2-cor-kappa-celda}
 \begin{aligned}
 \kappa_{k,N}:A_N^{I_k}\oplus A_N^{J_k}
 &\longrightarrow A_N^{I_k\times J_k},\\
 \kappa_{k,N}(u,v)_{(\alpha,\Sigma),(\beta,\Pi)}
 &:=u_{(\alpha,\Sigma)}-v_{(\beta,\Pi)}.
 \end{aligned}
\end{equation}
Elegidos los primeros puertos según el orden TPK, la diferencia rectangular
\begin{equation}
\label{eq:pdfv2-cor-delta-celda}
 \delta(w)_{ij}:=w_{ij}-w_{ij_0}-w_{i_0j}+w_{i_0j_0}
\end{equation}
descarga la sucesión escindida
\begin{equation}
\label{eq:pdfv2-cor-sucesion-celda}
 0\longrightarrow A_N\xrightarrow{\Delta}
 A_N^{I_k}\oplus A_N^{J_k}
 \xrightarrow{\kappa_{k,N}}A_N^{I_k\times J_k}
 \xrightarrow{\delta}A_N^{(|I_k|-1)(|J_k|-1)}
 \longrightarrow0.
\end{equation}
Ésta es una identidad de la misma fibra APP: no añade una celda ajena al
generador \(\mathfrak e_\alpha\).

La acción celular local usa las dos coordenadas enteras acumuladas por ese
mismo ledger. Para
\(\widehat x=\widehat x_k(\gamma)\) y
\(\widehat y=\alpha_*\widehat x=\widehat x_{k+1}(\gamma\alpha)\), ponemos
\begin{equation}
\label{eq:pdfv2-cor-c-v}
 \begin{aligned}
 c_{\widehat x}&:=\kappa(\gamma), &
 c_\alpha:=c_{\widehat y}&:=\kappa(\gamma\alpha),\\
 v_{\widehat x}&:=\sum_{\epsilon\ {\rm en}\ \gamma}
                  \operatorname{or}(\epsilon), &
 v_{\widehat y}&:=\sum_{\epsilon\ {\rm en}\ \gamma\alpha}
                  \operatorname{or}(\epsilon),\\
 v_\alpha&:=v_{\widehat y}-v_{\widehat x}
           =\operatorname{or}(\alpha).&&
 \end{aligned}
\end{equation}
Sobre los generadores \(f,r,e,b,g\) se define
\begin{equation}
\label{eq:pdfv2-cor-complejo-local}
 \partial f=3c_\alpha e+b,\qquad
 \partial r=0,\qquad
 \partial e=g,\qquad
 \partial b=-3c_\alpha g,\qquad
 \partial g=0.
\end{equation}
La cancelación
\begin{equation}
\label{eq:pdfv2-cor-dos-cero}
 \partial^2f=3c_\alpha g-3c_\alpha g=0
\end{equation}
es literal y usa el mismo carry con los dos signos impuestos por la
orientación.

La semilla de incidencia se obtiene del módulo libre de las dos hojas,
\begin{equation}
\label{eq:pdfv2-cor-v-tpk}
 V_{\rm TPK}:=\mathbf F_3e_\Sigma\oplus\mathbf F_3e_\Pi,
 \qquad
 \mathscr I:=\mathbf P^1(V_{\rm TPK}),
 \qquad |\mathscr I|=4,
\end{equation}
y de sus pares no ordenados
\begin{equation}
\label{eq:pdfv2-cor-incidencias}
 \mathscr E:=\binom{\mathscr I}{2},
 \qquad |\mathscr E|=6,
 \qquad
 \mathscr B:=\mathscr I\times\{+,-\},
 \qquad |\mathscr B|=8.
\end{equation}
La aplicación de incidencia queda determinada en la base por
\begin{equation}
\label{eq:pdfv2-cor-b-incidencia-generadores}
 B[e_L]=\sum_{L'\neq L}e_{\{L,L'\}},
 \qquad L\in\mathscr I,
\end{equation}
y, por linealidad,
\begin{equation}
\label{eq:pdfv2-cor-b-incidencia-coordenadas}
 (Ba)_{\{L,L'\}}=a_L+a_{L'}.
\end{equation}
Si \(s(q)=\sum_{E\in\mathscr E}q_E\), cada \(L\) aparece tres veces y
\begin{equation}
\label{eq:pdfv2-cor-sb-cero}
 s(Ba)=0.
\end{equation}
Por tanto la función
\begin{equation}
\label{eq:pdfv2-cor-wc}
 W_c(q):=\mathbf1_{\{s(q)=0\}}-\frac13
\end{equation}
satisface \(W_c(q+Ba)=W_c(q)\). Las identidades
\eqref{eq:pdfv2-cor-v-tpk}--\eqref{eq:pdfv2-cor-wc} son internas y exactas.
El lift material de \(q\) desde las celdas y los carries comunes se trata en
la capa de ensamblaje; no se confunde esta identidad de incidencia con la
construcción de ese lift.

\section{Transporte unitario, polarización y balance sobre la fibra común}

La no vacuidad de \eqref{eq:pdfv2-cor-hijos-supervivientes} permite fijar,
sin seleccionar una prolongación,
\begin{equation}
\label{eq:pdfv2-cor-peso-conteo}
 p_k(\alpha,\widehat y\mid\widehat x)
 :=\frac1{d_k(\widehat x)}.
\end{equation}
Sea el espacio de las dos hojas del mismo estado
\begin{equation}
\label{eq:pdfv2-cor-hilbert-hojas}
 \mathscr H_k:=
 \bigoplus_{\widehat x\in\widehat{\mathcal X}^{\rm tot}_k}
 \bigl(\mathsf H_{\widehat x}^{\Sigma}
       \oplus\mathsf H_{\widehat x}^{\Pi}\bigr),
\end{equation}
y sea \(\iota_{\widehat x}\) la inclusión de la fibra indicada. El transporte
TPK de \eqref{eq:continuo-comun-transporte-hojas} da la acción exacta
\begin{equation}
\label{eq:pdfv2-cor-u-comun}
 U_k\iota_{\widehat x}v:=
 \sum_{(\alpha,\widehat y)\in\operatorname{Ch}_k(\widehat x)}
 p_k(\alpha,\widehat y\mid\widehat x)^{1/2}
 \iota_{\widehat y}U_\alpha v.
\end{equation}
Los hijos de dos padres distintos son disjuntos porque el truncamiento es una
función. Además, dos testigos \(\alpha\) distintos no colapsan en el mismo
sumando directo: la ruta completa \(\gamma\), incluido su último paso, es una
coordenada de \(\widehat y\). En consecuencia, la identidad de Gram es
\begin{equation}
\label{eq:pdfv2-cor-u-gram}
 \left.U_k^*U_k\right|_{\mathsf H_{\widehat x}^{\Sigma}
       \oplus\mathsf H_{\widehat x}^{\Pi}}
 =
 \sum_{(\alpha,\widehat y)\in\operatorname{Ch}_k(\widehat x)}
 p_k(\alpha,\widehat y\mid\widehat x)\,U_\alpha^*U_\alpha.
\end{equation}
Por tanto \(U_k^*U_k=I\) exactamente cuando los transportes de hoja
\(U_\alpha\) publicados por TPK son isometrías. La normalización de los pesos,
por sí sola, no puede sustituir esa comprobación.
Para \(\ell>k\) se conserva la historia completa mediante
\begin{equation}
\label{eq:pdfv2-cor-u-iterado}
 U_{\ell,k}:=U_{\ell-1}\cdots U_k,
 \qquad U_{k,k}:=I.
\end{equation}

La graduación actúa sobre las dos coordenadas que ya están presentes en cada
fibra de \eqref{eq:pdfv2-cor-hilbert-hojas}:
\begin{equation}
\label{eq:pdfv2-cor-sigma}
 \sigma_k\iota_{\widehat x}(v_\Sigma,v_\Pi)
 :=\iota_{\widehat x}(v_\Sigma,-v_\Pi),
\qquad
 P_k:=\frac{I+\sigma_k}{2},\quad
 Q_k:=\frac{I-\sigma_k}{2}.
\end{equation}
El umbral TRIT pertenece a la orientación y al carry de
\(\widehat x\); no borra ninguna de las dos hojas APP y no requiere un tercer
caso en \eqref{eq:pdfv2-cor-sigma}.

La parte que conserva hoja y la parte que la cambia son dos componentes del
mismo \(U_k\):
\begin{align}
 A_k&:=P_{k+1}U_kP_k+Q_{k+1}U_kQ_k,
 \label{eq:pdfv2-cor-a-comun}\\
 \eta_k&:=Q_{k+1}U_kP_k+P_{k+1}U_kQ_k.
 \label{eq:pdfv2-cor-eta-comun}
\end{align}
Por ortogonalidad de \(P_{k+1}\) y \(Q_{k+1}\),
\begin{equation}
\label{eq:pdfv2-cor-balance-hojas}
 U_k=A_k+\eta_k,\qquad
 A_k^*A_k+\eta_k^*\eta_k=U_k^*U_k.
\end{equation}
El último miembro es \(I\) bajo, y sólo bajo, la condición isométrica
identificada después de \eqref{eq:pdfv2-cor-u-gram}.
Esta identidad y el incremento \eqref{eq:pdfv2-cor-b-alpha} no son dos
transportes: \(A_k,\eta_k\) actúan sobre la base de historias y
\(b_\alpha\) lee la arista de esa misma base.

\subsection{Polarización firmada, telescopía no autónoma y retorno sin reinicialización}
No se introduce aquí una rama espectral independiente. Se restringen al
mismo espacio \(\mathscr H_k\) y al mismo transporte \(U_k\) ya construidos
las notaciones
\[
 \mathcal H_k^{\mathrm{sp}}:=\mathscr H_k,
 \qquad
 \Lambda(\gamma_{\widehat x})
 :=\bigl(\operatorname{or},\operatorname{Carr},\operatorname{Mem},
          \partial,\gamma\bigr)(\widehat x),
\]
y \(\Sigma_{\rm reg}\Lambda(\gamma_{\widehat x})\) designa su familia de
truncamientos compatibles. De este modo, el lector firmado, el defecto y el
terminal siguientes son acciones del constructor correlativo común, no una
selección retrospectiva de historias.
\begingroup
  \let\section\subsection
  % Extracción quirúrgica del fuente teoremática V2 05a: líneas 293--394.
% El resto permanece en este archivo como procedencia, pero no se publica.
\iffalse
\chapter{Macroestructura espectral--polar interna del continuo}
\label{chap:continuo-macro-sp-interna}

\section{Alcance estrictamente interno}

La construcción de este capítulo termina antes de cualquier reconocimiento
externo. Sus únicos datos son el soporte APP de dos hojas, la orientación
TRIT, las extensiones y el carry TPK, el sistema proyectivo de historias, la
medida cilíndrica y el ledger firmado. El resultado es una macroestructura
graduada con transporte, expresión de cambio de hoja, lector firmado y
telescopía no autónoma. No se utiliza una función exterior para definir
ninguno de sus mapas.

\section{Dominio superviviente de la rama}

Para \(x_k\in\mathcal A_k\), sea
\begin{equation}
\label{eq:continuo-sp-msect-n}
 \operatorname{Msect}^{\mathrm{sp},(N)}_k(x_k)
 =\left\{(x_j)_{j=k}^{N}:
 \tau_{j+1,j}x_{j+1}=x_j,
 \ (x_j,x_{j+9})\in\Gamma_{9,j},
 \ \text{y se conserva el haz }\{\Sigma,\Pi\}\right\}.
\end{equation}
Definimos
\begin{equation}
\label{eq:continuo-sp-dominio-superviviente}
 \mathcal C_k^{\mathrm{sp}}
 =\left\{x_k:
 \operatorname{Msect}^{\mathrm{sp},(N)}_k(x_k)\neq\varnothing
 \text{ para todo }N\geq k\right\},
 \qquad
 \mathcal C_{\rm cont}^{\mathrm{sp}}
 =\varprojlim_k\mathcal C_k^{\mathrm{sp}}.
\end{equation}
El umbral TRIT no se excluye de este dominio. El haz conserva las dos hojas
aritméticas aunque una carta visible no publique ninguna de ellas en un
evento concreto.

\section{Paquete de trece coordenadas}

Para \(x_k\in\mathcal C_k^{\mathrm{sp}}\), el paquete dependiente es
\begin{equation}
\label{eq:continuo-sp-d13}
\begin{aligned}
 D_{\mathrm{sp},k}(x_k)=\bigl(&
 \mathcal F_{\mathrm{sp},k}(x_k),
 \operatorname{Ext}_{\Gamma,\mathrm{sp},k}(x_k),
 \operatorname{Rel}_{\mathrm{sp},k}(x_k),
 \mathcal N_{\mathrm{sp},k}(x_k),
 \operatorname{or}_{\mathrm{sp},k}(x_k),\\
 &\operatorname{Carr}_{\mathrm{sp},k}(x_k),
 \operatorname{Mem}_{\mathrm{sp},k}(x_k),
 \partial_{\mathrm{sp},k}(x_k),
 \gamma_{\mathrm{sp},k}(x_k),
 \operatorname{Msect}_{\mathrm{sp},k}(x_k),\\
 &\operatorname{Surv}_{\mathrm{sp},k}(x_k),
 \operatorname{Term}_{\mathrm{sp},k}(x_k),
 \mathcal L^{\rm sgn}_{\mathrm{sp},k}(x_k)\bigr).
\end{aligned}
\end{equation}
Sus componentes se determinan correlativamente como sigue.

La fibra \(\mathcal F_{\mathrm{sp},k}\) contiene el estado TPK y el haz
ordenado \(\{\Sigma,\Pi\}\). La extensión es la restricción completa de
\(\operatorname{Ext}_\Gamma\) a historias que satisfacen
\eqref{eq:continuo-sp-dominio-superviviente}. Las relaciones contienen las
leyes de identidad y composición, el lift residuo--cociente, el cociclo de
carry, la compatibilidad de frontera y la naturalidad del truncamiento. La
coordenada numérica contiene profundidad, fase, residuos, cocientes y firma
interna. La orientación almacena por separado hoja y orientación TRIT. La
memoria contiene el cociente, el ledger y el incremento de vuelta. La
frontera y la ruta permanecen sin contraer. La multisección es el sistema
completo de prolongaciones de \eqref{eq:continuo-sp-msect-n}. La
supervivencia es la pertenencia estable de
\eqref{eq:continuo-sp-dominio-superviviente}. El terminal y el lector se
construyen en las secciones siguientes.

Los truncamientos de todas las coordenadas forman un mapa
\(u^{\mathrm{sp}}_{\ell,k}\) que satisface
\begin{equation}
\label{eq:continuo-sp-d-naturalidad}
 u^{\mathrm{sp}}_{\ell,k}D_{\mathrm{sp},\ell}
 =D_{\mathrm{sp},k}\tau_{\ell,k}.
\end{equation}

\section{Gráfica de la restricción y sección dependiente}

Se define
\begin{equation}
\label{eq:continuo-sp-g-graph}
 \mathcal G_{\mathrm{sp},k}
 =\operatorname{Graph}(D_{\mathrm{sp},k})
 =\{(x_k,D_{\mathrm{sp},k}x_k):
     x_k\in\mathcal C_k^{\mathrm{sp}}\}.
\end{equation}
La restricción total es
\begin{equation}
\label{eq:continuo-sp-res-graph}
 \operatorname{Res}_{\mathrm{sp}}=s_{\mathrm{sp}}:
 \mathcal C_{\rm cont}^{\mathrm{sp}}
 \longrightarrow\mathcal G_{\mathrm{sp}},
 \qquad
 s_{\mathrm{sp}}(x)=(x,D_{\mathrm{sp}}x).
\end{equation}
La proyección \(\pi_1(x,d)=x\) verifica
\begin{equation}
\label{eq:continuo-sp-res-inversa}
 \pi_1s_{\mathrm{sp}}=I,
 \qquad
 s_{\mathrm{sp}}\pi_1=I_{\mathcal G_{\mathrm{sp}}}.
\end{equation}

Sea \(r_{\mathrm{sp},k}(x_k)\) la multisección de hojas y rutas que ya
pertenece a \(D_{\mathrm{sp},k}(x_k)\). Se define
\begin{equation}
\label{eq:continuo-sp-x-dependiente}
 X_{\chi,\mathrm{sp},k}
 =\{(r_{\mathrm{sp},k}x_k,x_k,D_{\mathrm{sp},k}x_k):
     x_k\in\mathcal C_k^{\mathrm{sp}}\},
\end{equation}
\begin{equation}
\label{eq:continuo-sp-pr-x}
 \operatorname{pr}_{X,\mathrm{sp}}(x,D_{\mathrm{sp}}x)
 =(r_{\mathrm{sp}}x,x,D_{\mathrm{sp}}x).
\end{equation}
Olvidar la primera coordenada de
\eqref{eq:continuo-sp-x-dependiente} es la inversa de
\(\operatorname{pr}_{X,\mathrm{sp}}\).

\section{Haz de hojas y medida proyectiva}

Sea
\begin{equation}
\label{eq:continuo-sp-haz-hojas}
 \widetilde{\mathcal C}_k^{\mathrm{sp}}
 =\{(x_k,h):x_k\in\mathcal C_k^{\mathrm{sp}},
              h\in\{\Sigma,\Pi\}\}.
\end{equation}
La medida proyectiva del continuo induce probabilidades
\(\widetilde\mu_k\) sobre este haz. Se restringe al soporte positivo; para
\(z\) en ese soporte,
\begin{equation}
\label{eq:continuo-sp-consistencia-medida}
 \widetilde\mu_k(z)
 =\sum_{\widetilde\tau_{k+1,k}(z')=z}
  \widetilde\mu_{k+1}(z').
\end{equation}
La linealización de nivel es
\begin{equation}
\label{eq:continuo-sp-hilbert-nivel}
 \mathcal H_k^{\mathrm{sp}}
 =\ell^2(\operatorname{supp}\widetilde\mu_k),
\end{equation}
con base ortonormal \(e_z\).

\section{Transporte de todas las extensiones y carry}

Para \(z'\) que prolonga \(z\), definimos el peso condicional positivo
\begin{equation}
\label{eq:continuo-sp-peso-condicional}
 w_k(z'\mid z)
 =\left(\frac{\widetilde\mu_{k+1}(z')}
               {\widetilde\mu_k(z)}\right)^{1/2}.
\end{equation}
El transporte es
\begin{equation}
\label{eq:continuo-sp-u-accion}
 U_ke_z
 =\sum_{\substack{z':\widetilde\tau_{k+1,k}(z')=z\\
            (z,z')\in\operatorname{Ext}_{\Gamma,\mathrm{sp},k}}}
 w_k(z'\mid z)e_{z'}.
\end{equation}
La suma contiene todas las extensiones admisibles. No depende de un selector.

Por \eqref{eq:continuo-sp-consistencia-medida},
\[
 \sum_{z':\widetilde\tau(z')=z}|w_k(z'\mid z)|^2=1.
\]
Los conjuntos de descendientes de dos padres distintos son disjuntos. En
consecuencia,
\begin{equation}
\label{eq:continuo-sp-u-isometria}
 U_k^*U_k=I_{\mathcal H_k^{\mathrm{sp}}}.
\end{equation}

Cada base \(e_{z'}\) incluye el carry de la extensión correspondiente. Para
rutas componibles se conservan
\begin{align}
 \kappa(\gamma_2\gamma_1)
 &=\kappa(\gamma_2)+
   \mathsf C(\gamma_2)\kappa(\gamma_1),
 \label{eq:continuo-sp-carry-nativo}\\
 \operatorname{Carr}(\gamma_2\gamma_1)
 &=\operatorname{Carr}(\gamma_2)H(\gamma_1)
  +Y(\gamma_2)\operatorname{Carr}(\gamma_1).
 \label{eq:continuo-sp-carry-operador}
\end{align}
Los pesos condicionales se multiplican a lo largo de la única cadena de
truncamientos de un descendiente. Las ecuaciones
\eqref{eq:continuo-sp-carry-nativo}--
\eqref{eq:continuo-sp-carry-operador} garantizan que las etiquetas de carry
coincidan. Por tanto,
\begin{equation}
\label{eq:continuo-sp-u-composicion}
 U_{\ell,k}=U_{\ell-1}\cdots U_k.
\end{equation}

\section{Graduación, proyectores y descomposición del transporte}

La graduación de hoja es la involución
\begin{equation}
\label{eq:continuo-sp-sigma}
 \sigma_ke_{(x,\Sigma)}=+e_{(x,\Sigma)},
 \qquad
 \sigma_ke_{(x,\Pi)}=-e_{(x,\Pi)}.
\end{equation}
Así,
\[
 \sigma_k^*=\sigma_k,
 \qquad
 \sigma_k^2=I.
\]
Los proyectores internos son
\begin{equation}
\label{eq:continuo-sp-p-q}
 P_k=\frac{I+\sigma_k}{2},
 \qquad
 Q_k=\frac{I-\sigma_k}{2}.
\end{equation}
Satisfacen
\begin{equation}
\label{eq:continuo-sp-p-q-relaciones}
 P_k^2=P_k,
 \quad Q_k^2=Q_k,
 \quad P_kQ_k=0,
 \quad P_k+Q_k=I.
\end{equation}

Separamos la parte del transporte que conserva hoja de la que la cambia:
\begin{align}
 A_k&=P_{k+1}U_kP_k+Q_{k+1}U_kQ_k,
 \label{eq:continuo-sp-a}\\
 \eta_k&=Q_{k+1}U_kP_k+P_{k+1}U_kQ_k.
 \label{eq:continuo-sp-eta}
\end{align}
Entonces
\begin{equation}
\label{eq:continuo-sp-u-a-eta}
 U_k=A_k+\eta_k,
\end{equation}
y la graduación verifica
\begin{equation}
\label{eq:continuo-sp-paridad-a-eta}
 \sigma_{k+1}A_k=A_k\sigma_k,
 \qquad
 \sigma_{k+1}\eta_k=-\eta_k\sigma_k.
\end{equation}

\begin{proposition}[Balance local de hojas]
\label{prop:continuo-sp-balance-local}
Para cada nivel,
\begin{equation}
\label{eq:continuo-sp-balance-local}
 A_k^*A_k+\eta_k^*\eta_k=I.
\end{equation}
\end{proposition}

\begin{proof}
La ortogonalidad de \(P_{k+1}\) y \(Q_{k+1}\) da
\begin{align*}
 A_k^*A_k
 &=P_kU_k^*P_{k+1}U_kP_k
   +Q_kU_k^*Q_{k+1}U_kQ_k,\\
 \eta_k^*\eta_k
 &=P_kU_k^*Q_{k+1}U_kP_k
   +Q_kU_k^*P_{k+1}U_kQ_k.
\end{align*}
Sumando y usando \(P_{k+1}+Q_{k+1}=I\),
\[
 A_k^*A_k+\eta_k^*\eta_k
 =P_kU_k^*U_kP_k+Q_kU_k^*U_kQ_k
 =P_k+Q_k=I.
\]
\end{proof}

El operador \(\eta_k\) registra el abandono del canal que conserva hoja. No
es el carry entero: las coordenadas
\eqref{eq:continuo-sp-carry-nativo}--
\eqref{eq:continuo-sp-carry-operador} permanecen en el paquete
\eqref{eq:continuo-sp-d13}.

\fi
\section{Lector firmado con registro genealógico}

Para \(v\in\mathcal H_k^{\mathrm{sp}}\), la polarización firmada es
\begin{equation}
\label{eq:continuo-sp-polarizacion-firmada}
 \operatorname{pol}_k(v)
 =\langle v,\sigma_kv\rangle
 =\|P_kv\|^2-\|Q_kv\|^2.
\end{equation}
El lector especializado conserva además la fuente:
\begin{equation}
\label{eq:continuo-sp-lector-firmado}
 \mathcal L^{\rm sgn}_{\mathrm{sp},k}(x_k;v)
 =\bigl(
   \Lambda(\gamma_{x_k}),
   \Sigma_{\rm reg}(\Lambda(\gamma_{x_k})),
   \|P_kv\|^2,
   \|Q_kv\|^2,
   \operatorname{pol}_k(v)\bigr).
\end{equation}
La orientación TRIT se conserva en \(\Lambda(\gamma_{x_k})\); no se
identifica con el signo de hoja. A horizonte \(N\), el registro guarda la
familia de lectores para todos los niveles \(k\leq N\). Truncar significa
eliminar únicamente las entradas posteriores.

\section{Producto no autónomo y terminal}

Se define
\begin{equation}
\label{eq:continuo-sp-f-producto}
 F_{0,0}=I_{\mathcal H_0^{\mathrm{sp}}},
 \qquad
 F_{0,k}=A_{k-1}\cdots A_1A_0
 \quad(k\geq1).
\end{equation}
Por tanto,
\begin{equation}
\label{eq:continuo-sp-f-recursion}
 F_{0,k+1}=A_kF_{0,k}.
\end{equation}
El término publicado al abandonar el canal en el paso \(k\) y el terminal a
horizonte \(N\) son
\begin{equation}
\label{eq:continuo-sp-defecto-terminal}
 \mathcal D_k=F_{0,k}^*\eta_k^*\eta_kF_{0,k},
 \qquad
 \mathcal T_N=F_{0,N}^*F_{0,N}.
\end{equation}

\begin{theorem}[Telescopía espectral--polar no autónoma]
\label{thm:continuo-sp-telescopia-no-autonoma}
Para todo horizonte finito \(N\),
\begin{equation}
\label{eq:continuo-sp-telescopia-no-autonoma}
 \boxed{
 I=\sum_{k=0}^{N-1}
 F_{0,k}^*\eta_k^*\eta_kF_{0,k}
 +F_{0,N}^*F_{0,N}.}
\end{equation}
Además,
\[
 0\preceq\mathcal T_{N+1}
 \preceq\mathcal T_N\preceq I,
\]
de modo que existe el límite fuerte positivo
\[
 \mathcal T_\infty
 =\operatorname{s-lim}_{N\to\infty}\mathcal T_N
\]
y
\begin{equation}
\label{eq:continuo-sp-telescopia-infinita}
 I=\sum_{k=0}^{\infty}
 F_{0,k}^*\eta_k^*\eta_kF_{0,k}
 +\mathcal T_\infty
\end{equation}
en topología fuerte. El teorema no afirma que
\(\mathcal T_\infty\) sea nulo.
\end{theorem}

\begin{proof}
Multiplicando \eqref{eq:pdfv2-cor-balance-hojas} por
\(F_{0,k}^*\) y \(F_{0,k}\), y usando
\eqref{eq:continuo-sp-f-recursion}, se obtiene
\[
 F_{0,k}^*F_{0,k}
 =F_{0,k+1}^*F_{0,k+1}
  +F_{0,k}^*\eta_k^*\eta_kF_{0,k}.
\]
La suma desde \(k=0\) hasta \(N-1\) cancela todos los términos
intermedios y deja precisamente
\eqref{eq:continuo-sp-telescopia-no-autonoma}. Como
\(A_k^*A_k\preceq I\),
\[
 \mathcal T_{N+1}
 =F_{0,N}^*A_N^*A_NF_{0,N}
 \preceq F_{0,N}^*F_{0,N}=\mathcal T_N.
\]
La monotonía y la cota positiva proporcionan el límite fuerte. Pasar al
límite en la identidad finita da
\eqref{eq:continuo-sp-telescopia-infinita}.
\end{proof}

\iffalse
\section{Ensamblador como gráfica}

Para
\[
 z_N=(r_{\mathrm{sp},N}x_N,x_N,D_{\mathrm{sp},N}x_N)
 \in X_{\chi,\mathrm{sp},N},
\]
se define el ensamblador material
\begin{equation}
\label{eq:continuo-sp-a-accion}
\begin{aligned}
 a_{\mathrm{sp},N}(z_N)=\bigl(&
 (\mathcal H_k^{\mathrm{sp}},\sigma_k,P_k,Q_k)_{k\leq N},
 (U_k,A_k,\eta_k)_{k<N},\\
 &(F_{0,k},\mathcal D_k,\mathcal T_k,
   \mathcal L^{\rm sgn}_{\mathrm{sp},k})_{k\leq N}\bigr).
\end{aligned}
\end{equation}
La macroestructura es
\begin{equation}
\label{eq:continuo-sp-m-graph}
 \mathfrak M_{\mathrm{sp},N}
 =\operatorname{Graph}(a_{\mathrm{sp},N}),
 \qquad
 \mathcal A_{\mathrm{sp},N}(z_N)
 =(z_N,a_{\mathrm{sp},N}z_N).
\end{equation}
La proyección a \(z_N\) es la inversa de
\(\mathcal A_{\mathrm{sp},N}\). El truncamiento elimina las entradas cuyo
índice supera el horizonte nuevo y verifica
\begin{equation}
\label{eq:continuo-sp-a-naturalidad}
 m^{\mathrm{sp}}_{\ell,k}a_{\mathrm{sp},\ell}
 =a_{\mathrm{sp},k}\xi^{\mathrm{sp}}_{\ell,k}.
\end{equation}

\section{Publicador como gráfica}

El publicador no contrae la macroestructura a un escalar. Devuelve el
certificado completo
\begin{equation}
\label{eq:continuo-sp-p-accion}
 p_{\mathrm{sp},N}(m_N)
 =\left(
 \mathcal L^{\rm sgn}_{\mathrm{sp},k},
 \mathcal D_k,
 \mathcal T_k,
 I=\sum_{j=0}^{k-1}\mathcal D_j+\mathcal T_k
 \right)_{0\leq k\leq N}.
\end{equation}
Se define
\begin{equation}
\label{eq:continuo-sp-c-graph}
 \mathcal C_{\mathrm{sp},N}
 =\operatorname{Graph}(p_{\mathrm{sp},N}),
 \qquad
 \mathcal P_{\mathrm{sp},N}(m_N)
 =(m_N,p_{\mathrm{sp},N}m_N).
\end{equation}
La proyección a \(m_N\) es su inversa. Truncar el certificado elimina sólo
las identidades posteriores, por lo que
\begin{equation}
\label{eq:continuo-sp-p-naturalidad}
 c^{\mathrm{sp}}_{\ell,k}p_{\mathrm{sp},\ell}
 =p_{\mathrm{sp},k}m^{\mathrm{sp}}_{\ell,k}.
\end{equation}

\section{Teorema de la cadena interna completa}

\begin{theorem}[Macroestructura espectral--polar generada por el continuo]
\label{thm:continuo-sp-cadena-interna}
Las acciones
\eqref{eq:continuo-sp-res-graph},
\eqref{eq:continuo-sp-pr-x},
\eqref{eq:continuo-sp-m-graph} y
\eqref{eq:continuo-sp-c-graph} forman la cadena natural
\begin{equation}
\label{eq:continuo-sp-cadena-interna}
 \mathcal C_{\rm cont}^{\rm disc}
 \supset\mathcal C_{\rm cont}^{\mathrm{sp}}
 \xrightarrow[\cong]{\operatorname{Res}_{\mathrm{sp}}}
 \mathcal G_{\mathrm{sp}}
 \xrightarrow[\cong]{\operatorname{pr}_{X,\mathrm{sp}}}
 X_{\chi,\mathrm{sp}}
 \xrightarrow[\cong]{\mathcal A_{\mathrm{sp}}}
 \mathfrak M_{\mathrm{sp}}
 \xrightarrow[\cong]{\mathcal P_{\mathrm{sp}}}
 \mathcal C_{\mathrm{sp}}.
\end{equation}
Su contenido publicado es el haz de dos hojas, la graduación
\(\sigma_k\), los proyectores \(P_k,Q_k\), el transporte total \(U_k\),
la separación \(U_k=A_k+\eta_k\), el carry y la memoria, el lector firmado
y la telescopía
\eqref{eq:continuo-sp-telescopia-no-autonoma} con terminal explícito. Cada
objeto conserva una proyección canónica hacia todos sus antecedentes.
\end{theorem}

\begin{proof}
Las identidades
\eqref{eq:continuo-sp-res-inversa} prueban la primera equivalencia. La
definición dependiente \eqref{eq:continuo-sp-x-dependiente} prueba la
segunda. Las gráficas
\eqref{eq:continuo-sp-m-graph} y
\eqref{eq:continuo-sp-c-graph} prueban las dos restantes. La naturalidad de
los cuatro pasos es
\eqref{eq:continuo-sp-d-naturalidad},
\eqref{eq:continuo-sp-a-naturalidad} y
\eqref{eq:continuo-sp-p-naturalidad}, junto con la naturalidad de
\(r_{\mathrm{sp},k}\). El contenido operatorial queda probado por
\eqref{eq:continuo-sp-u-isometria}, el
\cref{prop:continuo-sp-balance-local} y el
\cref{thm:continuo-sp-telescopia-no-autonoma}. Todas las extensiones aparecen
en \eqref{eq:continuo-sp-u-accion}; por ello no interviene una elección de
ruta.
\end{proof}

\section{Retorno de fase y evolución del estado}

La rama conserva expresamente
\begin{equation}
\label{eq:continuo-sp-no-reset}
 g(k+9)=g(k),
 \qquad
 x_{k+9}\neq x_k\quad\text{en general}.
\end{equation}
El segundo miembro registra el incremento de memoria y los posibles cambios
de hoja, carry, cilindro y frontera. Ni \(\sigma_k\), ni los proyectores, ni
el lector firmado contraen esta diferencia.

\section{Procedencia}

Las dos hojas APP, el estado trítico holonómico, las fibras y extensiones TPK,
el carry, el retorno sin reinicialización, la supervivencia y el ledger son
\texttt{ARQUITECTURA\_AUTORAL\_PREEXISTENTE}. La obligación de conservar el
terminal finito es \texttt{RESULTADO\_RECUPERADO}. El haz explícito de hojas,
la involución \(\sigma_k\), los proyectores, la linealización que suma todas
las extensiones, la descomposición \(U_k=A_k+\eta_k\), el producto no
autónomo \(F_{0,k}\) y las gráficas dependientes son
\texttt{FORMALIZACION\_NUEVA}. El
\cref{thm:continuo-sp-telescopia-no-autonoma,thm:continuo-sp-cadena-interna}
constituye el \texttt{CERTIFICADO\_NUEVO} de esa formalización.

\section{Falsadores internos}

\begin{proposition}[Falsadores de la macroestructura]
\label{prop:continuo-sp-falsadores}
La conclusión del
\cref{thm:continuo-sp-cadena-interna} no es transferible si ocurre cualquiera
de las siguientes sustituciones:
\begin{enumerate}
 \item elegir una prolongación en
       \eqref{eq:continuo-sp-u-accion} y borrar las restantes;
 \item sustituir \(\mathcal G_{\mathrm{sp}}\) por la imagen exterior de
       \(D_{\mathrm{sp}}\);
 \item borrar \((x,D_{\mathrm{sp}}x)\) al formar
       \(X_{\chi,\mathrm{sp}}\);
 \item identificar la hoja con la orientación TRIT o borrar el umbral;
 \item romper la consistencia
       \eqref{eq:continuo-sp-consistencia-medida};
 \item borrar el carry o incumplir
       \eqref{eq:continuo-sp-carry-nativo}--
       \eqref{eq:continuo-sp-carry-operador};
 \item identificar \(\eta_k\) con todo el carry;
 \item reemplazar la familia \((A_k)_k\) por un solo operador sin demostrar
       estacionariedad;
 \item suprimir \(F_{0,N}^*F_{0,N}\) en un horizonte finito;
 \item afirmar \(\mathcal T_\infty=0\) sin una prueba adicional;
 \item reconstruir una historia desde la polarización firmada;
 \item usar \(g(k+9)=g(k)\) para afirmar \(x_{k+9}=x_k\);
 \item introducir un dominio, una función o un criterio exterior para
       definir cualquiera de las flechas internas;
 \item introducir una salida posterior como antecedente de la rama.
\end{enumerate}
\end{proposition}
\fi
   % Extracción quirúrgica del fuente teoremática V2 05a: líneas 512--523.
% El resto permanece en este archivo como procedencia, pero no se publica.
\iffalse
\chapter{Macroestructura espectral--polar interna del continuo}
\label{chap:continuo-macro-sp-interna}

\section{Alcance estrictamente interno}

La construcción de este capítulo termina antes de cualquier reconocimiento
externo. Sus únicos datos son el soporte APP de dos hojas, la orientación
TRIT, las extensiones y el carry TPK, el sistema proyectivo de historias, la
medida cilíndrica y el ledger firmado. El resultado es una macroestructura
graduada con transporte, expresión de cambio de hoja, lector firmado y
telescopía no autónoma. No se utiliza una función exterior para definir
ninguno de sus mapas.

\section{Dominio superviviente de la rama}

Para \(x_k\in\mathcal A_k\), sea
\begin{equation}
\label{eq:continuo-sp-msect-n}
 \operatorname{Msect}^{\mathrm{sp},(N)}_k(x_k)
 =\left\{(x_j)_{j=k}^{N}:
 \tau_{j+1,j}x_{j+1}=x_j,
 \ (x_j,x_{j+9})\in\Gamma_{9,j},
 \ \text{y se conserva el haz }\{\Sigma,\Pi\}\right\}.
\end{equation}
Definimos
\begin{equation}
\label{eq:continuo-sp-dominio-superviviente}
 \mathcal C_k^{\mathrm{sp}}
 =\left\{x_k:
 \operatorname{Msect}^{\mathrm{sp},(N)}_k(x_k)\neq\varnothing
 \text{ para todo }N\geq k\right\},
 \qquad
 \mathcal C_{\rm cont}^{\mathrm{sp}}
 =\varprojlim_k\mathcal C_k^{\mathrm{sp}}.
\end{equation}
El umbral TRIT no se excluye de este dominio. El haz conserva las dos hojas
aritméticas aunque una carta visible no publique ninguna de ellas en un
evento concreto.

\section{Paquete de trece coordenadas}

Para \(x_k\in\mathcal C_k^{\mathrm{sp}}\), el paquete dependiente es
\begin{equation}
\label{eq:continuo-sp-d13}
\begin{aligned}
 D_{\mathrm{sp},k}(x_k)=\bigl(&
 \mathcal F_{\mathrm{sp},k}(x_k),
 \operatorname{Ext}_{\Gamma,\mathrm{sp},k}(x_k),
 \operatorname{Rel}_{\mathrm{sp},k}(x_k),
 \mathcal N_{\mathrm{sp},k}(x_k),
 \operatorname{or}_{\mathrm{sp},k}(x_k),\\
 &\operatorname{Carr}_{\mathrm{sp},k}(x_k),
 \operatorname{Mem}_{\mathrm{sp},k}(x_k),
 \partial_{\mathrm{sp},k}(x_k),
 \gamma_{\mathrm{sp},k}(x_k),
 \operatorname{Msect}_{\mathrm{sp},k}(x_k),\\
 &\operatorname{Surv}_{\mathrm{sp},k}(x_k),
 \operatorname{Term}_{\mathrm{sp},k}(x_k),
 \mathcal L^{\rm sgn}_{\mathrm{sp},k}(x_k)\bigr).
\end{aligned}
\end{equation}
Sus componentes se determinan correlativamente como sigue.

La fibra \(\mathcal F_{\mathrm{sp},k}\) contiene el estado TPK y el haz
ordenado \(\{\Sigma,\Pi\}\). La extensión es la restricción completa de
\(\operatorname{Ext}_\Gamma\) a historias que satisfacen
\eqref{eq:continuo-sp-dominio-superviviente}. Las relaciones contienen las
leyes de identidad y composición, el lift residuo--cociente, el cociclo de
carry, la compatibilidad de frontera y la naturalidad del truncamiento. La
coordenada numérica contiene profundidad, fase, residuos, cocientes y firma
interna. La orientación almacena por separado hoja y orientación TRIT. La
memoria contiene el cociente, el ledger y el incremento de vuelta. La
frontera y la ruta permanecen sin contraer. La multisección es el sistema
completo de prolongaciones de \eqref{eq:continuo-sp-msect-n}. La
supervivencia es la pertenencia estable de
\eqref{eq:continuo-sp-dominio-superviviente}. El terminal y el lector se
construyen en las secciones siguientes.

Los truncamientos de todas las coordenadas forman un mapa
\(u^{\mathrm{sp}}_{\ell,k}\) que satisface
\begin{equation}
\label{eq:continuo-sp-d-naturalidad}
 u^{\mathrm{sp}}_{\ell,k}D_{\mathrm{sp},\ell}
 =D_{\mathrm{sp},k}\tau_{\ell,k}.
\end{equation}

\section{Gráfica de la restricción y sección dependiente}

Se define
\begin{equation}
\label{eq:continuo-sp-g-graph}
 \mathcal G_{\mathrm{sp},k}
 =\operatorname{Graph}(D_{\mathrm{sp},k})
 =\{(x_k,D_{\mathrm{sp},k}x_k):
     x_k\in\mathcal C_k^{\mathrm{sp}}\}.
\end{equation}
La restricción total es
\begin{equation}
\label{eq:continuo-sp-res-graph}
 \operatorname{Res}_{\mathrm{sp}}=s_{\mathrm{sp}}:
 \mathcal C_{\rm cont}^{\mathrm{sp}}
 \longrightarrow\mathcal G_{\mathrm{sp}},
 \qquad
 s_{\mathrm{sp}}(x)=(x,D_{\mathrm{sp}}x).
\end{equation}
La proyección \(\pi_1(x,d)=x\) verifica
\begin{equation}
\label{eq:continuo-sp-res-inversa}
 \pi_1s_{\mathrm{sp}}=I,
 \qquad
 s_{\mathrm{sp}}\pi_1=I_{\mathcal G_{\mathrm{sp}}}.
\end{equation}

Sea \(r_{\mathrm{sp},k}(x_k)\) la multisección de hojas y rutas que ya
pertenece a \(D_{\mathrm{sp},k}(x_k)\). Se define
\begin{equation}
\label{eq:continuo-sp-x-dependiente}
 X_{\chi,\mathrm{sp},k}
 =\{(r_{\mathrm{sp},k}x_k,x_k,D_{\mathrm{sp},k}x_k):
     x_k\in\mathcal C_k^{\mathrm{sp}}\},
\end{equation}
\begin{equation}
\label{eq:continuo-sp-pr-x}
 \operatorname{pr}_{X,\mathrm{sp}}(x,D_{\mathrm{sp}}x)
 =(r_{\mathrm{sp}}x,x,D_{\mathrm{sp}}x).
\end{equation}
Olvidar la primera coordenada de
\eqref{eq:continuo-sp-x-dependiente} es la inversa de
\(\operatorname{pr}_{X,\mathrm{sp}}\).

\section{Haz de hojas y medida proyectiva}

Sea
\begin{equation}
\label{eq:continuo-sp-haz-hojas}
 \widetilde{\mathcal C}_k^{\mathrm{sp}}
 =\{(x_k,h):x_k\in\mathcal C_k^{\mathrm{sp}},
              h\in\{\Sigma,\Pi\}\}.
\end{equation}
La medida proyectiva del continuo induce probabilidades
\(\widetilde\mu_k\) sobre este haz. Se restringe al soporte positivo; para
\(z\) en ese soporte,
\begin{equation}
\label{eq:continuo-sp-consistencia-medida}
 \widetilde\mu_k(z)
 =\sum_{\widetilde\tau_{k+1,k}(z')=z}
  \widetilde\mu_{k+1}(z').
\end{equation}
La linealización de nivel es
\begin{equation}
\label{eq:continuo-sp-hilbert-nivel}
 \mathcal H_k^{\mathrm{sp}}
 =\ell^2(\operatorname{supp}\widetilde\mu_k),
\end{equation}
con base ortonormal \(e_z\).

\section{Transporte de todas las extensiones y carry}

Para \(z'\) que prolonga \(z\), definimos el peso condicional positivo
\begin{equation}
\label{eq:continuo-sp-peso-condicional}
 w_k(z'\mid z)
 =\left(\frac{\widetilde\mu_{k+1}(z')}
               {\widetilde\mu_k(z)}\right)^{1/2}.
\end{equation}
El transporte es
\begin{equation}
\label{eq:continuo-sp-u-accion}
 U_ke_z
 =\sum_{\substack{z':\widetilde\tau_{k+1,k}(z')=z\\
            (z,z')\in\operatorname{Ext}_{\Gamma,\mathrm{sp},k}}}
 w_k(z'\mid z)e_{z'}.
\end{equation}
La suma contiene todas las extensiones admisibles. No depende de un selector.

Por \eqref{eq:continuo-sp-consistencia-medida},
\[
 \sum_{z':\widetilde\tau(z')=z}|w_k(z'\mid z)|^2=1.
\]
Los conjuntos de descendientes de dos padres distintos son disjuntos. En
consecuencia,
\begin{equation}
\label{eq:continuo-sp-u-isometria}
 U_k^*U_k=I_{\mathcal H_k^{\mathrm{sp}}}.
\end{equation}

Cada base \(e_{z'}\) incluye el carry de la extensión correspondiente. Para
rutas componibles se conservan
\begin{align}
 \kappa(\gamma_2\gamma_1)
 &=\kappa(\gamma_2)+
   \mathsf C(\gamma_2)\kappa(\gamma_1),
 \label{eq:continuo-sp-carry-nativo}\\
 \operatorname{Carr}(\gamma_2\gamma_1)
 &=\operatorname{Carr}(\gamma_2)H(\gamma_1)
  +Y(\gamma_2)\operatorname{Carr}(\gamma_1).
 \label{eq:continuo-sp-carry-operador}
\end{align}
Los pesos condicionales se multiplican a lo largo de la única cadena de
truncamientos de un descendiente. Las ecuaciones
\eqref{eq:continuo-sp-carry-nativo}--
\eqref{eq:continuo-sp-carry-operador} garantizan que las etiquetas de carry
coincidan. Por tanto,
\begin{equation}
\label{eq:continuo-sp-u-composicion}
 U_{\ell,k}=U_{\ell-1}\cdots U_k.
\end{equation}

\section{Graduación, proyectores y descomposición del transporte}

La graduación de hoja es la involución
\begin{equation}
\label{eq:continuo-sp-sigma}
 \sigma_ke_{(x,\Sigma)}=+e_{(x,\Sigma)},
 \qquad
 \sigma_ke_{(x,\Pi)}=-e_{(x,\Pi)}.
\end{equation}
Así,
\[
 \sigma_k^*=\sigma_k,
 \qquad
 \sigma_k^2=I.
\]
Los proyectores internos son
\begin{equation}
\label{eq:continuo-sp-p-q}
 P_k=\frac{I+\sigma_k}{2},
 \qquad
 Q_k=\frac{I-\sigma_k}{2}.
\end{equation}
Satisfacen
\begin{equation}
\label{eq:continuo-sp-p-q-relaciones}
 P_k^2=P_k,
 \quad Q_k^2=Q_k,
 \quad P_kQ_k=0,
 \quad P_k+Q_k=I.
\end{equation}

Separamos la parte del transporte que conserva hoja de la que la cambia:
\begin{align}
 A_k&=P_{k+1}U_kP_k+Q_{k+1}U_kQ_k,
 \label{eq:continuo-sp-a}\\
 \eta_k&=Q_{k+1}U_kP_k+P_{k+1}U_kQ_k.
 \label{eq:continuo-sp-eta}
\end{align}
Entonces
\begin{equation}
\label{eq:continuo-sp-u-a-eta}
 U_k=A_k+\eta_k,
\end{equation}
y la graduación verifica
\begin{equation}
\label{eq:continuo-sp-paridad-a-eta}
 \sigma_{k+1}A_k=A_k\sigma_k,
 \qquad
 \sigma_{k+1}\eta_k=-\eta_k\sigma_k.
\end{equation}

\begin{proposition}[Balance local de hojas]
\label{prop:continuo-sp-balance-local}
Para cada nivel,
\begin{equation}
\label{eq:continuo-sp-balance-local}
 A_k^*A_k+\eta_k^*\eta_k=I.
\end{equation}
\end{proposition}

\begin{proof}
La ortogonalidad de \(P_{k+1}\) y \(Q_{k+1}\) da
\begin{align*}
 A_k^*A_k
 &=P_kU_k^*P_{k+1}U_kP_k
   +Q_kU_k^*Q_{k+1}U_kQ_k,\\
 \eta_k^*\eta_k
 &=P_kU_k^*Q_{k+1}U_kP_k
   +Q_kU_k^*P_{k+1}U_kQ_k.
\end{align*}
Sumando y usando \(P_{k+1}+Q_{k+1}=I\),
\[
 A_k^*A_k+\eta_k^*\eta_k
 =P_kU_k^*U_kP_k+Q_kU_k^*U_kQ_k
 =P_k+Q_k=I.
\]
\end{proof}

El operador \(\eta_k\) registra el abandono del canal que conserva hoja. No
es el carry entero: las coordenadas
\eqref{eq:continuo-sp-carry-nativo}--
\eqref{eq:continuo-sp-carry-operador} permanecen en el paquete
\eqref{eq:continuo-sp-d13}.

\section{Lector firmado con ledger}

Para \(v\in\mathcal H_k^{\mathrm{sp}}\), la polarización firmada es
\begin{equation}
\label{eq:continuo-sp-polarizacion-firmada}
 \operatorname{pol}_k(v)
 =\langle v,\sigma_kv\rangle
 =\|P_kv\|^2-\|Q_kv\|^2.
\end{equation}
El lector especializado conserva además la fuente:
\begin{equation}
\label{eq:continuo-sp-lector-firmado}
 \mathcal L^{\rm sgn}_{\mathrm{sp},k}(x_k;v)
 =\bigl(
   \Lambda(\gamma_{x_k}),
   \Sigma_{\rm reg}(\Lambda(\gamma_{x_k})),
   \|P_kv\|^2,
   \|Q_kv\|^2,
   \operatorname{pol}_k(v)\bigr).
\end{equation}
La orientación TRIT se conserva en \(\Lambda(\gamma_{x_k})\); no se
identifica con el signo de hoja. A horizonte \(N\), el registro guarda la
familia de lectores para todos los niveles \(k\leq N\). Truncar significa
eliminar únicamente las entradas posteriores.

\section{Producto no autónomo y terminal}

Se define
\begin{equation}
\label{eq:continuo-sp-f-producto}
 F_{0,0}=I_{\mathcal H_0^{\mathrm{sp}}},
 \qquad
 F_{0,k}=A_{k-1}\cdots A_1A_0
 \quad(k\geq1).
\end{equation}
Por tanto,
\begin{equation}
\label{eq:continuo-sp-f-recursion}
 F_{0,k+1}=A_kF_{0,k}.
\end{equation}
El término publicado al abandonar el canal en el paso \(k\) y el terminal a
horizonte \(N\) son
\begin{equation}
\label{eq:continuo-sp-defecto-terminal}
 \mathcal D_k=F_{0,k}^*\eta_k^*\eta_kF_{0,k},
 \qquad
 \mathcal T_N=F_{0,N}^*F_{0,N}.
\end{equation}

\begin{theorem}[Telescopía espectral--polar no autónoma]
\label{thm:continuo-sp-telescopia-no-autonoma}
Para todo horizonte finito \(N\),
\begin{equation}
\label{eq:continuo-sp-telescopia-no-autonoma}
 \boxed{
 I=\sum_{k=0}^{N-1}
 F_{0,k}^*\eta_k^*\eta_kF_{0,k}
 +F_{0,N}^*F_{0,N}.}
\end{equation}
Además,
\[
 0\preceq\mathcal T_{N+1}
 \preceq\mathcal T_N\preceq I,
\]
de modo que existe el límite fuerte positivo
\[
 \mathcal T_\infty
 =\operatorname{s-lim}_{N\to\infty}\mathcal T_N
\]
y
\begin{equation}
\label{eq:continuo-sp-telescopia-infinita}
 I=\sum_{k=0}^{\infty}
 F_{0,k}^*\eta_k^*\eta_kF_{0,k}
 +\mathcal T_\infty
\end{equation}
en topología fuerte. El teorema no afirma que
\(\mathcal T_\infty\) sea nulo.
\end{theorem}

\begin{proof}
Multiplicando \eqref{eq:continuo-sp-balance-local} por
\(F_{0,k}^*\) y \(F_{0,k}\), y usando
\eqref{eq:continuo-sp-f-recursion}, se obtiene
\[
 F_{0,k}^*F_{0,k}
 =F_{0,k+1}^*F_{0,k+1}
  +F_{0,k}^*\eta_k^*\eta_kF_{0,k}.
\]
La suma desde \(k=0\) hasta \(N-1\) cancela todos los términos
intermedios y deja precisamente
\eqref{eq:continuo-sp-telescopia-no-autonoma}. Como
\(A_k^*A_k\preceq I\),
\[
 \mathcal T_{N+1}
 =F_{0,N}^*A_N^*A_NF_{0,N}
 \preceq F_{0,N}^*F_{0,N}=\mathcal T_N.
\]
La monotonía y la cota positiva proporcionan el límite fuerte. Pasar al
límite en la identidad finita da
\eqref{eq:continuo-sp-telescopia-infinita}.
\end{proof}

\section{Ensamblador como gráfica}

Para
\[
 z_N=(r_{\mathrm{sp},N}x_N,x_N,D_{\mathrm{sp},N}x_N)
 \in X_{\chi,\mathrm{sp},N},
\]
se define el ensamblador material
\begin{equation}
\label{eq:continuo-sp-a-accion}
\begin{aligned}
 a_{\mathrm{sp},N}(z_N)=\bigl(&
 (\mathcal H_k^{\mathrm{sp}},\sigma_k,P_k,Q_k)_{k\leq N},
 (U_k,A_k,\eta_k)_{k<N},\\
 &(F_{0,k},\mathcal D_k,\mathcal T_k,
   \mathcal L^{\rm sgn}_{\mathrm{sp},k})_{k\leq N}\bigr).
\end{aligned}
\end{equation}
La macroestructura es
\begin{equation}
\label{eq:continuo-sp-m-graph}
 \mathfrak M_{\mathrm{sp},N}
 =\operatorname{Graph}(a_{\mathrm{sp},N}),
 \qquad
 \mathcal A_{\mathrm{sp},N}(z_N)
 =(z_N,a_{\mathrm{sp},N}z_N).
\end{equation}
La proyección a \(z_N\) es la inversa de
\(\mathcal A_{\mathrm{sp},N}\). El truncamiento elimina las entradas cuyo
índice supera el horizonte nuevo y verifica
\begin{equation}
\label{eq:continuo-sp-a-naturalidad}
 m^{\mathrm{sp}}_{\ell,k}a_{\mathrm{sp},\ell}
 =a_{\mathrm{sp},k}\xi^{\mathrm{sp}}_{\ell,k}.
\end{equation}

\section{Publicador como gráfica}

El publicador no contrae la macroestructura a un escalar. Devuelve el
certificado completo
\begin{equation}
\label{eq:continuo-sp-p-accion}
 p_{\mathrm{sp},N}(m_N)
 =\left(
 \mathcal L^{\rm sgn}_{\mathrm{sp},k},
 \mathcal D_k,
 \mathcal T_k,
 I=\sum_{j=0}^{k-1}\mathcal D_j+\mathcal T_k
 \right)_{0\leq k\leq N}.
\end{equation}
Se define
\begin{equation}
\label{eq:continuo-sp-c-graph}
 \mathcal C_{\mathrm{sp},N}
 =\operatorname{Graph}(p_{\mathrm{sp},N}),
 \qquad
 \mathcal P_{\mathrm{sp},N}(m_N)
 =(m_N,p_{\mathrm{sp},N}m_N).
\end{equation}
La proyección a \(m_N\) es su inversa. Truncar el certificado elimina sólo
las identidades posteriores, por lo que
\begin{equation}
\label{eq:continuo-sp-p-naturalidad}
 c^{\mathrm{sp}}_{\ell,k}p_{\mathrm{sp},\ell}
 =p_{\mathrm{sp},k}m^{\mathrm{sp}}_{\ell,k}.
\end{equation}

\section{Teorema de la cadena interna completa}

\begin{theorem}[Macroestructura espectral--polar generada por el continuo]
\label{thm:continuo-sp-cadena-interna}
Las acciones
\eqref{eq:continuo-sp-res-graph},
\eqref{eq:continuo-sp-pr-x},
\eqref{eq:continuo-sp-m-graph} y
\eqref{eq:continuo-sp-c-graph} forman la cadena natural
\begin{equation}
\label{eq:continuo-sp-cadena-interna}
 \mathcal C_{\rm cont}^{\rm disc}
 \supset\mathcal C_{\rm cont}^{\mathrm{sp}}
 \xrightarrow[\cong]{\operatorname{Res}_{\mathrm{sp}}}
 \mathcal G_{\mathrm{sp}}
 \xrightarrow[\cong]{\operatorname{pr}_{X,\mathrm{sp}}}
 X_{\chi,\mathrm{sp}}
 \xrightarrow[\cong]{\mathcal A_{\mathrm{sp}}}
 \mathfrak M_{\mathrm{sp}}
 \xrightarrow[\cong]{\mathcal P_{\mathrm{sp}}}
 \mathcal C_{\mathrm{sp}}.
\end{equation}
Su contenido publicado es el haz de dos hojas, la graduación
\(\sigma_k\), los proyectores \(P_k,Q_k\), el transporte total \(U_k\),
la separación \(U_k=A_k+\eta_k\), el carry y la memoria, el lector firmado
y la telescopía
\eqref{eq:continuo-sp-telescopia-no-autonoma} con terminal explícito. Cada
objeto conserva una proyección canónica hacia todos sus antecedentes.
\end{theorem}

\begin{proof}
Las identidades
\eqref{eq:continuo-sp-res-inversa} prueban la primera equivalencia. La
definición dependiente \eqref{eq:continuo-sp-x-dependiente} prueba la
segunda. Las gráficas
\eqref{eq:continuo-sp-m-graph} y
\eqref{eq:continuo-sp-c-graph} prueban las dos restantes. La naturalidad de
los cuatro pasos es
\eqref{eq:continuo-sp-d-naturalidad},
\eqref{eq:continuo-sp-a-naturalidad} y
\eqref{eq:continuo-sp-p-naturalidad}, junto con la naturalidad de
\(r_{\mathrm{sp},k}\). El contenido operatorial queda probado por
\eqref{eq:continuo-sp-u-isometria}, el
\cref{prop:continuo-sp-balance-local} y el
\cref{thm:continuo-sp-telescopia-no-autonoma}. Todas las extensiones aparecen
en \eqref{eq:continuo-sp-u-accion}; por ello no interviene una elección de
ruta.
\end{proof}

\fi
\section{Retorno de fase y evolución no reinicializada del estado común}

La rama conserva expresamente
\begin{equation}
\label{eq:continuo-sp-no-reset}
 g(k+9)=g(k),
 \qquad
 x_{k+9}\neq x_k\quad\text{en general}.
\end{equation}
El segundo miembro registra el incremento de memoria y los posibles cambios
de hoja, carry, cilindro y frontera. Ni \(\sigma_k\), ni los proyectores, ni
el lector firmado contraen esta diferencia.

\iffalse
\section{Procedencia}

Las dos hojas APP, el estado trítico holonómico, las fibras y extensiones TPK,
el carry, el retorno sin reinicialización, la supervivencia y el ledger son
\texttt{ARQUITECTURA\_AUTORAL\_PREEXISTENTE}. La obligación de conservar el
terminal finito es \texttt{RESULTADO\_RECUPERADO}. El haz explícito de hojas,
la involución \(\sigma_k\), los proyectores, la linealización que suma todas
las extensiones, la descomposición \(U_k=A_k+\eta_k\), el producto no
autónomo \(F_{0,k}\) y las gráficas dependientes son
\texttt{FORMALIZACION\_NUEVA}. El
\cref{thm:continuo-sp-telescopia-no-autonoma,thm:continuo-sp-cadena-interna}
constituye el \texttt{CERTIFICADO\_NUEVO} de esa formalización.

\section{Falsadores internos}

\begin{proposition}[Falsadores de la macroestructura]
\label{prop:continuo-sp-falsadores}
La conclusión del
\cref{thm:continuo-sp-cadena-interna} no es transferible si ocurre cualquiera
de las siguientes sustituciones:
\begin{enumerate}
 \item elegir una prolongación en
       \eqref{eq:continuo-sp-u-accion} y borrar las restantes;
 \item sustituir \(\mathcal G_{\mathrm{sp}}\) por la imagen exterior de
       \(D_{\mathrm{sp}}\);
 \item borrar \((x,D_{\mathrm{sp}}x)\) al formar
       \(X_{\chi,\mathrm{sp}}\);
 \item identificar la hoja con la orientación TRIT o borrar el umbral;
 \item romper la consistencia
       \eqref{eq:continuo-sp-consistencia-medida};
 \item borrar el carry o incumplir
       \eqref{eq:continuo-sp-carry-nativo}--
       \eqref{eq:continuo-sp-carry-operador};
 \item identificar \(\eta_k\) con todo el carry;
 \item reemplazar la familia \((A_k)_k\) por un solo operador sin demostrar
       estacionariedad;
 \item suprimir \(F_{0,N}^*F_{0,N}\) en un horizonte finito;
 \item afirmar \(\mathcal T_\infty=0\) sin una prueba adicional;
 \item reconstruir una historia desde la polarización firmada;
 \item usar \(g(k+9)=g(k)\) para afirmar \(x_{k+9}=x_k\);
 \item introducir un dominio, una función o un criterio exterior para
       definir cualquiera de las flechas internas;
 \item introducir una salida posterior como antecedente de la rama.
\end{enumerate}
\end{proposition}
\fi
 \endgroup

La esperanza condicionada asociada al mismo peso es
\begin{equation}
\label{eq:pdfv2-cor-esperanza-comun}
 (E_kF)(\widehat x):=
 \sum_{(\alpha,\widehat y)\in\operatorname{Ch}_k(\widehat x)}
 p_k(\alpha,\widehat y\mid\widehat x)F(\alpha,\widehat y).
\end{equation}
Para el incremento fino \(b_f\) se ponen
\begin{equation}
\label{eq:pdfv2-cor-kappa-jensen}
 b_c:=E_kb_f,\qquad
 \kappa_k^{\rm can}:=\frac{E_k|b_f|-|b_c|}{2}\geq0.
\end{equation}
Entonces
\begin{equation}
\label{eq:pdfv2-cor-jensen-hojas}
 E_kb_f^+=b_c^++\kappa_k^{\rm can},\qquad
 E_kb_f^-=b_c^-+\kappa_k^{\rm can}.
\end{equation}
El superíndice \({\rm can}\) distingue este carry de cancelación del carry
entero \(\kappa_\alpha\); ambos proceden del mismo transporte, pero tienen
tipos distintos.

Las memorias firmadas no requieren una condición inicial exterior: para
\(\gamma=(\epsilon_1,\ldots,\epsilon_k)\) se definen por suma del propio
ledger,
\begin{equation}
\label{eq:pdfv2-cor-memorias}
 \begin{aligned}
 M_0^\pm&:=0,\\
 M_k^\pm(\gamma)&:=\sum_{j=1}^k b_{\epsilon_j}^\pm,\\
 M_{k+1}^\pm(\gamma\alpha)&:=M_k^\pm(\gamma)+b_\alpha^\pm,\\
 D_k^0&:=M_k^+-M_k^-, &
 H_k^0&:=M_k^++M_k^-.
 \end{aligned}
\end{equation}
Así, el balance firmado, la memoria total y la polarización
\(\langle v,\sigma_kv\rangle\) son tres lecturas coordinadas del mismo
prefijo, no tres estados reconstruidos por separado.

\subsection{Refinamiento nonádico, pérdidas ortogonales y memoria terminal}
El índice \(j\) del desarrollo siguiente es el nivel del mismo árbol de
prolongaciones y su registro de trece coordenadas es la restricción del
registro común de \eqref{eq:pdfv2-cor-registro-extension}. Por tanto se fija
\(D_{{\rm sol},j}:=D_j^{(13)}\) y no se postula un dominio adicional.
\begingroup
  \let\section\subsection
  % Extracción quirúrgica del fuente teoremática V2 05b: líneas 100--359.
% El resto permanece en este archivo como procedencia, pero no se publica.
\iffalse
% Macroestructura solenoidal estrictamente interna del continuo discreto.
% Este módulo no introduce ninguna ecuación clásica ni un problema exterior.

\chapter{Macroestructura interna de balance, memoria y telescopía}
\label{chap:pdfv2-sol-interna}

\begin{statusbox}
La única entrada de este capítulo es la estructura discreta del continuo.
La rama \(\mathrm{sol}\) se construye antes de cualquier evaluación escalar
exterior y antes de cualquier formulación diferencial clásica.  Su cadena
completa es
\[
 \mathcal C_{\rm cont}^{\rm disc}
 \supset\mathcal C_{\rm cont}^{\rm sol}
 \xrightarrow{\operatorname{Res}_{\rm sol}}
 \mathcal G_{\rm sol}
 \xrightarrow{\operatorname{pr}_{X,{\rm sol}}}
 X_{\chi,{\rm sol}}
 \xrightarrow{\mathcal A_{\rm sol}}
 \mathfrak M_{\rm sol}
 \xrightarrow{\mathcal P_{\rm sol}}
 \mathcal C_{\rm sol}^{\rm HMT}.
\]
Cada flecha conserva como coordenadas la historia y el aparato que la
produce.  Ninguna imagen de esta cadena es un cociente que olvide su fuente.
\end{statusbox}

\section{Dominio propio antes de toda evaluación posterior}
\label{sec:pdfv2-sol-dominio}

Sea
\[
 x=(x_0\xrightarrow{e_0}x_1\xrightarrow{e_1}\cdots)
 \in\mathcal C_{\rm cont}^{\rm disc}
\]
una historia superviviente.  Cada arista \(e_j\) conserva sus dos hojas,
orientación, residuo, cociente y carry entero.  Sean
\(N_j^+(e_j)\) y \(N_j^-(e_j)\) las multiplicidades de las dos hojas y sea
\(\Delta_{\gamma_j}\in\mathbb Z\) el coeficiente orientado de la ruta, con el
carry de la composición ya incorporado.  El incremento firmado interno es
\begin{equation}\label{eq:pdfv2-sol-b}
 b_j=b(e_j):=
 \Delta_{\gamma_j}\bigl(N_j^+(e_j)-N_j^-(e_j)\bigr)\in\mathbb Z.
\end{equation}
No se obtiene evaluando una trayectoria exterior: es una coordenada del
ledger de la propia arista.

Sus partes positiva y negativa se construyen por
\begin{equation}\label{eq:pdfv2-sol-bpm}
 b_j^+:=\frac{|b_j|+b_j}{2},
 \qquad
 b_j^-:=\frac{|b_j|-b_j}{2}.
\end{equation}
Por tanto,
\begin{equation}\label{eq:pdfv2-sol-bpm-identidades}
 b_j=b_j^+-b_j^-,
 \qquad
 |b_j|=b_j^++b_j^-,
 \qquad
 b_j^+b_j^-=0.
\end{equation}

Las dos memorias primitivas y sus dos lectores son
\begin{align}
 M_0^+=M_0^-&=0,
 &M_j^+&=\sum_{0\leq i<j}b_i^+,
 &M_j^-&=\sum_{0\leq i<j}b_i^-;
 \label{eq:pdfv2-sol-memorias}\\
 D_j^0&:=M_j^+-M_j^-,
 &H_j^0&:=M_j^++M_j^-.
 \label{eq:pdfv2-sol-dh}
\end{align}
En consecuencia,
\begin{equation}\label{eq:pdfv2-sol-dh-diferencias}
 D_{j+1}^0-D_j^0=b_j,
 \qquad
 H_{j+1}^0-H_j^0=|b_j|.
\end{equation}
La memoria firmada \(D^0\) y la memoria total \(H^0\) permanecen como
objetos distintos; conservar solamente su diferencia final destruiría la
genealogía.

La subestructura \(\mathcal C_{\rm cont}^{\rm sol}\) consta de las historias
para las que, en todo prefijo y todo refinamiento nonádico:
\begin{enumerate}
 \item existen los incrementos \eqref{eq:pdfv2-sol-b} y las dos memorias
       \eqref{eq:pdfv2-sol-memorias};
 \item los cilindros finos se agrupan en fibras de nueve extensiones con la
       misma ley de probabilidad proyectiva;
 \item las proyecciones de profundidad y resolución definidas abajo son
       compatibles con las extensiones de historia;
 \item cada fila de pérdidas y cada columna terminal es finita a horizonte
       finito;
 \item todo prefijo admite al menos una extensión compatible a cualquier
       horizonte posterior.
\end{enumerate}
Estas son condiciones de pertenencia a la rama, no premisas silenciosas
atribuidas a una historia arbitraria.

\fi
\section{Refinamiento nonádico común: inclusión, esperanza y memoria de balance}
\label{sec:pdfv2-sol-e9j9}

Sea \(Y_j(x)\) el conjunto finito de cilindros compatibles con el prefijo
\(x_{\leq j}\), provisto de la medida proyectiva \(\mu_j\).  La truncación
\[
 \tau_{j+1,j}:Y_{j+1}(x)\longrightarrow Y_j(x)
\]
tiene nueve extensiones admisibles sobre cada cilindro.  Se consideran los
espacios finitos
\[
 \mathscr H_j(x):=\ell^2\bigl(Y_j(x),\mu_j\bigr).
\]
La inclusión y la esperanza nonádicas son
\begin{align}
 J_{9,j}:\mathscr H_j&\longrightarrow\mathscr H_{j+1},
 &(J_{9,j}f)(y)&=f(\tau_{j+1,j}y),
 \label{eq:pdfv2-sol-j9}\\
 E_{9,j}:\mathscr H_{j+1}&\longrightarrow\mathscr H_j,
 &(E_{9,j}g)(z)&=\frac1{9}
   \sum_{\tau_{j+1,j}y=z}g(y).
 \label{eq:pdfv2-sol-e9}
\end{align}
La compatibilidad proyectiva de \(\mu_{j+1}\) y \(\mu_j\) da
\begin{equation}\label{eq:pdfv2-sol-ej-projector}
 E_{9,j}J_{9,j}=I_{\mathscr H_j},
 \qquad
 J_{9,j}^*=E_{9,j},
 \qquad
 P_{9,j}:=J_{9,j}E_{9,j}=P_{9,j}^*=P_{9,j}^2.
\end{equation}
El proyector microscópico complementario es
\begin{equation}\label{eq:pdfv2-sol-qmicro-preliminar}
 Q_{9,j}:=I-P_{9,j}.
\end{equation}

Si \(b_f\) es el incremento sobre la fibra fina y
\begin{equation}\label{eq:pdfv2-sol-bcoarse}
 b_c:=E_{9,j}b_f,
\end{equation}
la convexidad elemental de \(|\cdot|\) construye, no postula, el carry de
cancelación
\begin{equation}\label{eq:pdfv2-sol-kappa}
 \kappa_j
 :=\frac{E_{9,j}|b_f|-|b_c|}{2}\geq0.
\end{equation}
En efecto, usando \(t^\pm=(|t|\pm t)/2\), se obtiene
\begin{align}
 E_{9,j}b_f^+
 &=\frac{E_{9,j}|b_f|+E_{9,j}b_f}{2}
   =b_c^++\kappa_j,
 \label{eq:pdfv2-sol-jensen-plus}\\
 E_{9,j}b_f^-
 &=\frac{E_{9,j}|b_f|-E_{9,j}b_f}{2}
   =b_c^-+\kappa_j.
 \label{eq:pdfv2-sol-jensen-minus}
\end{align}
Las dos hojas pierden la misma cantidad al comprimirse; por eso la diferencia
firmada se conserva mientras la memoria total registra la cancelación:
\begin{equation}\label{eq:pdfv2-sol-jensen-dh}
 E_{9,j}(b_f^+-b_f^-)=b_c,
 \qquad
 E_{9,j}(b_f^++b_f^-)=|b_c|+2\kappa_j.
\end{equation}

\section{Factorización ortogonal de los grados de profundidad y de las capas de resolución en \(\mathscr H_j^{\mathrm{depth}}\widehat\otimes\mathscr H_j^{\mathrm{res}}\)}
\label{sec:pdfv2-sol-micro-shell}

Para impedir que una misma pérdida sea contada a la vez como profundidad y
como resolución, el espacio de cada nivel se tipa como
\begin{equation}\label{eq:pdfv2-sol-factor-hilbert}
 \mathscr H_j
 =\mathscr H_j^{\rm depth}\widehat\otimes
  \mathscr H_j^{\rm res}.
\end{equation}
Para \(j\geq1\), el primer factor porta
\[
 P_j^{\rm depth}:=J_{9,j-1}E_{9,j-1}
 \quad\text{sobre }\mathscr H_j^{\rm depth},
\]
y se fija \(P_0^{\rm depth}=I\). El segundo factor lleva la partición
finita de shells heredada de la frontera de cilindros.  Si
\(\sigma_j:Y_j^{\rm res}\to S_j\) es esa partición, sus operadores son
\begin{align}
 J_j^{\rm res}f&=f\circ\sigma_j,
 &E_j^{\rm res}g(s)&=
 \frac1{\#\sigma_j^{-1}(s)}
 \sum_{\sigma_j(y)=s}g(y),
 \label{eq:pdfv2-sol-shell-je}\\
 P_j^{\rm res}&:=J_j^{\rm res}E_j^{\rm res}.
 \label{eq:pdfv2-sol-shell-p}
\end{align}

Se definen los tres proyectores mutuamente ortogonales
\begin{align}
 P_j^{\rm term}
 &:=P_j^{\rm depth}\otimes P_j^{\rm res},
 \label{eq:pdfv2-sol-pterm}\\
 Q_j^{\rm micro}
 &:=(I-P_j^{\rm depth})\otimes I,
 \label{eq:pdfv2-sol-qmicro}\\
 Q_j^{\rm shell}
 &:=P_j^{\rm depth}\otimes(I-P_j^{\rm res}).
 \label{eq:pdfv2-sol-qshell}
\end{align}
Entonces
\begin{equation}\label{eq:pdfv2-sol-ortogonalidad}
 Q_j^{\rm micro}Q_j^{\rm shell}=0,
 \qquad
 I=P_j^{\rm term}+Q_j^{\rm micro}+Q_j^{\rm shell},
\end{equation}
y, para todo \(\xi\in\mathscr H_j\),
\begin{equation}\label{eq:pdfv2-sol-pitagoras}
 \|\xi\|^2
 =\|P_j^{\rm term}\xi\|^2
  +\|Q_j^{\rm micro}\xi\|^2
  +\|Q_j^{\rm shell}\xi\|^2.
\end{equation}
No aparece el término
\((I-P_j^{\rm depth})\otimes(I-P_j^{\rm res})\) por segunda vez: ya está
contenido exactamente una vez en \(Q_j^{\rm micro}\).

\section{Extensiones de cilindro y fila completa de pérdidas}
\label{sec:pdfv2-sol-gamma-loss}

Una extensión de cilindros de nivel \(j\) a nivel \(j+1\) induce el mapa
\begin{equation}\label{eq:pdfv2-sol-gamma}
 \Gamma_j:\mathscr H_j\longrightarrow\mathscr H_{j+1},
 \qquad
 \Gamma_j f=f\circ\tau_{j+1,j},
\end{equation}
con todas las etiquetas de hoja, orientación, carry, memoria y ruta
transportadas en la base de cilindros.  Para \(n>j\), se escribe
\begin{equation}\label{eq:pdfv2-sol-gamma-composition}
 \Gamma_{j:n}:=\Gamma_{n-1}\cdots\Gamma_j,
 \qquad
 \Gamma_{j:j}:=I.
\end{equation}

Sobre \(\mathscr H_j\), los incrementos no negativos definen operadores de
multiplicación
\begin{equation}\label{eq:pdfv2-sol-loss-multipliers}
 B_j^+:=M_{\sqrt{b_j^+}},
 \qquad
 B_j^-:=M_{\sqrt{b_j^-}},
 \qquad
 K_j:=M_{\sqrt{2\kappa_j}}.
\end{equation}
Cada raíz se toma punto a punto sobre el conjunto finito de cilindros; si una
coordenada es nula, su operador es nulo.  El espacio de pérdidas es la suma
ortogonal externa
\[
 \mathscr E_j
 :=\mathscr H_j^{(+)}\oplus\mathscr H_j^{(-)}
   \oplus\mathscr H_j^{(\kappa)}
   \oplus\mathscr H_j^{({\rm micro})}
   \oplus\mathscr H_j^{({\rm shell})},
\]
y la fila completa es
\begin{equation}\label{eq:pdfv2-sol-loss-row}
 L_j:\mathscr H_j\longrightarrow\mathscr E_j,
 \qquad
 L_j\xi
 =\bigl(
 B_j^+\xi,B_j^-\xi,K_j\xi,
 Q_j^{\rm micro}\xi,Q_j^{\rm shell}\xi
 \bigr).
\end{equation}
Por tanto,
\begin{equation}\label{eq:pdfv2-sol-loss-norm}
 \|L_j\xi\|^2
 =\|B_j^+\xi\|^2+\|B_j^-\xi\|^2+\|K_j\xi\|^2
  +\|Q_j^{\rm micro}\xi\|^2
  +\|Q_j^{\rm shell}\xi\|^2.
\end{equation}
Las partes positiva, negativa, carry de cancelación, microestructura y shell
aparecen una vez cada una.  No se fusionan en una constante anónima ni se
repiten en dos sumandos.

\section{Terminal de historia completa}
\label{sec:pdfv2-sol-terminal}

Sea \(\operatorname{Hist}_{J}^{\rm sol}\) el conjunto de historias
supervivientes hasta el horizonte \(J\), cada una con su paquete de trece
coordenadas.  El espacio terminal se define como el espacio libre de esos
registros completos:
\begin{equation}\label{eq:pdfv2-sol-terminal-space}
 \mathscr T_J^{\rm sol}
 :=\ell^2\!\left(
 \left\{(h,D_{{\rm sol},J}(h)):
 h\in\operatorname{Hist}_{J}^{\rm sol}\right\}
 \right).
\end{equation}
La evaluación terminal
\begin{equation}\label{eq:pdfv2-sol-terminal-evaluation}
 E_J^{\rm term}:\mathscr H_J\longrightarrow\mathscr T_J^{\rm sol}
\end{equation}
envía cada vector básico de cilindro al registro completo de la historia que
lo genera y se extiende linealmente.  En particular, el terminal retiene
fibra, relaciones, números, orientación, carry, memoria, frontera, ruta,
multisección, supervivencia y lector.  No se sustituye por el último valor
visible de la ruta.

\section{Columna de observación, Gram y telescopía}
\label{sec:pdfv2-sol-gram}

Para \(0\leq j\leq J\), la columna futura completa es
\begin{equation}\label{eq:pdfv2-sol-g-column}
 \begin{aligned}
 G_{j;J}:\mathscr H_j&\longrightarrow
 \mathscr T_J^{\rm sol}\oplus
 \bigoplus_{n=j}^{J-1}\mathscr E_n,\\
 G_{j;J}\xi
 &:={}
 E_J^{\rm term}\Gamma_{j:J}\xi
 \ \oplus\ 
 \bigoplus_{n=j}^{J-1}L_n\Gamma_{j:n}\xi.
 \end{aligned}
\end{equation}
Para \(j=J\), la suma de pérdidas es vacía y
\(G_{J;J}=E_J^{\rm term}\).  El Gram terminal es
\begin{equation}\label{eq:pdfv2-sol-u-gram}
 U_{j;J}:=G_{j;J}^*G_{j;J}\succeq0.
\end{equation}

Reordenando la suma ortogonal del codominio se obtiene la identidad literal
\begin{equation}\label{eq:pdfv2-sol-g-recursion}
 G_{j;J}\xi
 \simeq
 \bigl(L_j\xi,\,G_{j+1;J}\Gamma_j\xi\bigr).
\end{equation}
En consecuencia,
\begin{equation}\label{eq:pdfv2-sol-stein}
 \boxed{U_{j;J}
 =L_j^*L_j+\Gamma_j^*U_{j+1;J}\Gamma_j}
\end{equation}
y
\begin{equation}\label{eq:pdfv2-sol-stein-difference}
 U_{j;J}-\Gamma_j^*U_{j+1;J}\Gamma_j
 =L_j^*L_j\succeq0.
\end{equation}
No hay una mutación de signo: el término de pérdida está en el lado positivo
de la identidad porque procede de una suma de cuadrados.

Iterando \eqref{eq:pdfv2-sol-stein} se obtiene la telescopía exacta
\begin{equation}\label{eq:pdfv2-sol-telescopia-operatorial}
 U_{j;J}
 =\Gamma_{j:J}^*(E_J^{\rm term})^*E_J^{\rm term}\Gamma_{j:J}
  +\sum_{n=j}^{J-1}
   \Gamma_{j:n}^*L_n^*L_n\Gamma_{j:n},
\end{equation}
o, de forma cuadrática,
\begin{equation}\label{eq:pdfv2-sol-telescopia-normas}
 \langle U_{j;J}\xi,\xi\rangle
 =\|E_J^{\rm term}\Gamma_{j:J}\xi\|^2
  +\sum_{n=j}^{J-1}\|L_n\Gamma_{j:n}\xi\|^2.
\end{equation}
El primer término conserva el futuro terminal completo; la suma conserva,
nivel a nivel, todas las pérdidas y todos los carries que condujeron hasta él.

\iffalse
\section{Las trece coordenadas de la restricción}
\label{sec:pdfv2-sol-trece}

Para un prefijo \(x_{\leq j}\), se define
\begin{equation}\label{eq:pdfv2-sol-d-package}
 \begin{aligned}
 D_{{\rm sol},j}(x):=\bigl(&
 \mathcal F_{{\rm sol},j},
 \operatorname{Ext}_{\Gamma,{\rm sol},j},
 \operatorname{Rel}_{{\rm sol},j},
 \mathcal N_{{\rm sol},j},
 \operatorname{or}_{{\rm sol},j},
 \operatorname{Carr}_{{\rm sol},j},
 \operatorname{Mem}_{{\rm sol},j},\\
 &\partial_{{\rm sol},j},
 \gamma_{{\rm sol},j},
 \operatorname{Msect}_{{\rm sol},j},
 \operatorname{Surv}_{{\rm sol},j},
 \operatorname{Term}_{{\rm sol},j},
 \mathcal L_{{\rm sol},j}\bigr).
 \end{aligned}
\end{equation}
Sus componentes son las siguientes.
\begin{enumerate}
 \item \(\mathcal F_{{\rm sol},j}\) contiene los cilindros \(Y_j\), su medida
       proyectiva, \(\mathscr H_j^{\rm depth}\),
       \(\mathscr H_j^{\rm res}\), sus productos tensoriales y los cinco
       espacios de pérdidas.
 \item \(\operatorname{Ext}_{\Gamma,{\rm sol},j}\) contiene
       \(J_{9,j},E_{9,j},\Gamma_j,\Gamma_{j:n}\), las inclusiones de shell y
       todos los mapas de truncación y refinamiento.
 \item \(\operatorname{Rel}_{{\rm sol},j}\) contiene
       \eqref{eq:pdfv2-sol-bpm-identidades},
       \eqref{eq:pdfv2-sol-ej-projector},
       \eqref{eq:pdfv2-sol-jensen-plus}--
       \eqref{eq:pdfv2-sol-jensen-minus},
       \eqref{eq:pdfv2-sol-ortogonalidad} y
       \eqref{eq:pdfv2-sol-stein}.
 \item \(\mathcal N_{{\rm sol},j}\) conserva
       \(j,J,9,b_j,b_j^+,b_j^-,\kappa_j\), multiplicidades de hoja, tamaños
       de fibras y dimensiones de shells.
 \item \(\operatorname{or}_{{\rm sol},j}\) conserva
       \(\Delta_{\gamma_j}\), el orden de ruta y la acción de inversión
       \(b_j\mapsto-b_j\), que intercambia \(b_j^+\) y \(b_j^-\) sin cambiar
       \(|b_j|\) ni \(\kappa_j\).
 \item \(\operatorname{Carr}_{{\rm sol},j}\) conserva el carry TPK entero,
       el carry de composición y el carry de cancelación \(\kappa_j\) como
       coordenadas diferentes.
 \item \(\operatorname{Mem}_{{\rm sol},j}\) es
       \((M_j^+,M_j^-,D_j^0,H_j^0)\) junto con el ledger ordenado de todos los
       incrementos anteriores.
 \item \(\partial_{{\rm sol},j}\) contiene la frontera orientada de cilindros,
       los incrementos \(b_j\), la fila \(L_j\) y las diferencias de Stein.
 \item \(\gamma_{{\rm sol},j}\) es la ruta completa
       \(e_0e_1\cdots e_{j-1}\), no sólo sus extremos.
 \item \(\operatorname{Msect}_{{\rm sol},j}\) contiene todas las extensiones
       nonádicas y de shell compatibles; no elige una prolongación.
 \item \(\operatorname{Surv}_{{\rm sol},j}\) registra los horizontes
       \(J\geq j\) para los que existen \(G_{j;J}\), \(U_{j;J}\) y
       terminales compatibles.
 \item \(\operatorname{Term}_{{\rm sol},j}\) contiene
       \(E_J^{\rm term}\Gamma_{j:J}\), el registro completo de historia y
       las trece coordenadas en cada horizonte.
 \item \(\mathcal L_{{\rm sol},j}\) publica internamente
       \((b^\pm,M^\pm,D^0,H^0,\kappa,Q^{\rm micro},Q^{\rm shell},
       L,G,U)\) y sus identidades.
\end{enumerate}

Para una extensión \(u:x_j\to x_{j'}\), el transporte coordinado
\(u_{{\rm sol},j',j}\) satisface
\begin{equation}\label{eq:pdfv2-sol-naturality-d}
 u_{{\rm sol},j',j}D_{{\rm sol},j'}(x_{j'})
 =D_{{\rm sol},j}(\tau_{j',j}x_{j'}).
\end{equation}
La igualdad incluye los cinco componentes de \(L\), el terminal completo y
la composición de \(\Gamma\); no es una igualdad sólo de cardinales.

\section{Grafo de restricción y objeto dependiente}
\label{sec:pdfv2-sol-graphs}

Se define
\begin{equation}\label{eq:pdfv2-sol-g-graph}
 \mathcal G_{\rm sol}:=\operatorname{Graph}(D_{\rm sol}),
 \qquad
 \operatorname{Res}_{\rm sol}(x):=(x,D_{\rm sol}x).
\end{equation}
La primera proyección es inversa sobre el grafo:
\begin{equation}\label{eq:pdfv2-sol-res-inverse}
 \pi_1\operatorname{Res}_{\rm sol}=I,
 \qquad
 \operatorname{Res}_{\rm sol}\pi_1=I_{\mathcal G_{\rm sol}}.
\end{equation}

Sea \(\chi_{\rm sol}(x)\) la multisección completa formada por todos los
cilindros, refinamientos, shells, filas de pérdidas y terminales compatibles
con \(x\).  Entonces
\begin{equation}\label{eq:pdfv2-sol-xchi}
 X_{\chi,{\rm sol}}
 :=\{(\chi_{\rm sol}(x),x,D_{\rm sol}x):
 x\in\mathcal C_{\rm cont}^{\rm sol}\},
\end{equation}
y
\begin{equation}\label{eq:pdfv2-sol-prx}
 \operatorname{pr}_{X,{\rm sol}}(x,D_{\rm sol}x)
 :=(\chi_{\rm sol}(x),x,D_{\rm sol}x).
\end{equation}
De nuevo, la proyección que conserva \((x,D_{\rm sol}x)\) es inversa sobre
esta imagen dependiente.

\section{Ensamblador y publicador internos}
\label{sec:pdfv2-sol-assembler-publisher}

El ambiente \(\mathscr M_{\rm sol}^{\rm amb}\) consta de sistemas
compatibles
\[
 \bigl(
 \{b,b^\pm,M^\pm,D^0,H^0,\kappa\},
 \{E_9,J_9,P^{\rm term},Q^{\rm micro},Q^{\rm shell}\},
 \{\Gamma,L,E^{\rm term},G,U\}
 \bigr)
\]
con todas las identidades anteriores.  El ensamblador bruto es
\begin{equation}\label{eq:pdfv2-sol-assembler}
 \begin{aligned}
 a_{\rm sol}:X_{\chi,{\rm sol}}&\longrightarrow
 \mathscr M_{\rm sol}^{\rm amb},\\
 a_{\rm sol}(\chi(x),x,Dx)&:=
 \bigl(
 \{b,b^\pm,M^\pm,D^0,H^0,\kappa\}_x,
 \{E_9,J_9,P,Q\}_x,
 \{\Gamma,L,E^{\rm term},G,U\}_x
 \bigr).
 \end{aligned}
\end{equation}
La macroestructura conserva su antecedente mediante
\begin{equation}\label{eq:pdfv2-sol-macro-graph}
 \mathfrak M_{\rm sol}:=\operatorname{Graph}(a_{\rm sol}),
 \qquad
 \mathcal A_{\rm sol}(z):=(z,a_{\rm sol}z).
\end{equation}

Sea \(\mathscr Y_{\rm sol}^{\rm int}\) el tipo de certificados formados por
las identidades
\begin{equation}\label{eq:pdfv2-sol-certificate-list}
 \begin{gathered}
 b=b^+-b^-,\qquad |b|=b^++b^-,\\
 E_9J_9=I,\qquad J_9E_9=P_9,\\
 E_9b_f^\pm=b_c^\pm+\kappa,\qquad\kappa\geq0,\\
 I=P^{\rm term}+Q^{\rm micro}+Q^{\rm shell},
 \qquad Q^{\rm micro}Q^{\rm shell}=0,\\
 U_{j;J}=L_j^*L_j+\Gamma_j^*U_{j+1;J}\Gamma_j,\\
 \langle U_{j;J}\xi,\xi\rangle
 =\|E_J^{\rm term}\Gamma_{j:J}\xi\|^2
  +\sum_{n=j}^{J-1}\|L_n\Gamma_{j:n}\xi\|^2,
 \end{gathered}
\end{equation}
junto con la naturalidad de todas las flechas.  El publicador bruto
\begin{equation}\label{eq:pdfv2-sol-publisher}
 p_{\rm sol}^{\rm raw}:\mathfrak M_{\rm sol}
 \longrightarrow\mathscr Y_{\rm sol}^{\rm int}
\end{equation}
publica ese certificado sin eliminar el macroobjeto.  Finalmente,
\begin{equation}\label{eq:pdfv2-sol-output-graph}
 \mathcal C_{\rm sol}^{\rm HMT}
 :=\operatorname{Graph}(p_{\rm sol}^{\rm raw}),
 \qquad
 \mathcal P_{\rm sol}(m):=(m,p_{\rm sol}^{\rm raw}m).
\end{equation}

\section{Teorema interno de generación solenoidal}
\label{sec:pdfv2-sol-theorem}

\begin{theorem}[Generación interna de la macroestructura de balance]
\label{thm:pdfv2-sol-internal-generation}
Para todo \(x\in\mathcal C_{\rm cont}^{\rm sol}\), la cadena
\[
 \mathcal C_{\rm cont}^{\rm sol}
 \xrightarrow{\operatorname{Res}_{\rm sol}}
 \mathcal G_{\rm sol}
 \xrightarrow{\operatorname{pr}_{X,{\rm sol}}}
 X_{\chi,{\rm sol}}
 \xrightarrow{\mathcal A_{\rm sol}}
 \mathfrak M_{\rm sol}
 \xrightarrow{\mathcal P_{\rm sol}}
 \mathcal C_{\rm sol}^{\rm HMT}
\]
construye de forma natural una macroestructura que contiene las dos memorias,
el carry de cancelación, la separación ortogonal micro/shell, la fila completa
de pérdidas, la columna futura, el terminal de historia completa, el Gram y la
telescopía.  Las trece coordenadas del continuo pueden recuperarse por
proyecciones sucesivas desde cualquier punto de la cadena.
\end{theorem}

\begin{proof}
Las identidades \eqref{eq:pdfv2-sol-bpm-identidades} y
\eqref{eq:pdfv2-sol-dh-diferencias} se obtienen directamente de la
descomposición positiva y negativa del incremento entero; por ello las
memorias no son datos añadidos.  Las fórmulas
\eqref{eq:pdfv2-sol-j9}--\eqref{eq:pdfv2-sol-e9}, junto con la medida
proyectiva, prueban \(E_9J_9=I\), \(J_9^*=E_9\) y que \(P_9\) es un
proyector ortogonal.

La desigualdad triangular da
\(E_9|b_f|\geq|E_9b_f|=|b_c|\), de modo que
\(\kappa\geq0\).  Sustituir
\(b^\pm=(|b|\pm b)/2\) prueba
\eqref{eq:pdfv2-sol-jensen-plus} y
\eqref{eq:pdfv2-sol-jensen-minus} sin una hipótesis adicional.

Las definiciones tensoriales
\eqref{eq:pdfv2-sol-pterm}--\eqref{eq:pdfv2-sol-qshell} prueban que
\(P^{\rm term}\), \(Q^{\rm micro}\) y \(Q^{\rm shell}\) son ortogonales y
suman la identidad.  En consecuencia, la fila \(L_j\) cuenta cada sector una
sola vez.

La definición \eqref{eq:pdfv2-sol-g-column} separa ortogonalmente el terminal
y las pérdidas de cada nivel.  Separar el primer sumando de pérdidas produce
\eqref{eq:pdfv2-sol-g-recursion}.  Tomar adjunto y componer prueba la identidad
de Stein \eqref{eq:pdfv2-sol-stein}; iterarla prueba
\eqref{eq:pdfv2-sol-telescopia-operatorial} y
\eqref{eq:pdfv2-sol-telescopia-normas}.  Todo término es una norma cuadrada,
por lo que no puede aparecer el signo opuesto.

La compatibilidad de truncación de cilindros, esperanza, proyecciones y
extensiones \(\Gamma\) da la naturalidad
\eqref{eq:pdfv2-sol-naturality-d}.  Finalmente,
\(\mathcal G_{\rm sol}\), \(\mathfrak M_{\rm sol}\) y
\(\mathcal C_{\rm sol}^{\rm HMT}\) son grafos, y
\(X_{\chi,{\rm sol}}\) es un objeto dependiente que conserva explícitamente
su antecedente.  La primera proyección de cada etapa recupera la etapa
anterior.  Así ninguna composición borra historia, memoria, multisección,
terminal ni lector.
\end{proof}

\section{Procedencia probatoria y falsadores}
\label{sec:pdfv2-sol-provenance-falsifiers}

\paragraph{Resultado recuperado.}
Pertenecen a la arquitectura ya construida el incremento firmado con carry,
la separación \(b=b^+-b^-\), las memorias \(M^\pm,D^0,H^0\), el par
\(E_9,J_9\), el carry de Jensen \(\kappa\), la extensión de cilindros, la
fila de pérdidas, el terminal, el Gram, Stein y la telescopía.

\paragraph{Formalización nueva.}
Son formalizaciones explícitas de esa arquitectura el tipado previo a toda
evaluación exterior, la factorización
\(\mathscr H^{\rm depth}\widehat\otimes\mathscr H^{\rm res}\), la separación
ortogonal que impide el doble conteo, el terminal libre sobre historias con
trece coordenadas y los cuatro empaquetados no ablativos por grafos.

\begin{criterion}[Falsadores internos de la rama \(\mathrm{sol}\)]
\label{crit:pdfv2-sol-falsifiers}
La construcción falla si ocurre cualquiera de los hechos siguientes:
\begin{enumerate}
 \item \(b\neq b^+-b^-\) o \(|b|\neq b^++b^-\);
 \item \(E_9J_9\neq I\), la medida no es proyectiva o \(P_9\) no es
       proyector;
 \item \(\kappa<0\) o fallan
       \(E_9b_f^\pm=b_c^\pm+\kappa\);
 \item \(Q^{\rm micro}Q^{\rm shell}\neq0\), o un mismo sector aparece en
       ambos sumandos de la fila \(L\);
 \item una de \(b^+,b^-,\kappa,Q^{\rm micro},Q^{\rm shell}\) se omite o se
       cuenta dos veces;
 \item el terminal conserva sólo el valor visible y no la historia con sus
       trece coordenadas;
 \item falla la identidad de Stein o la telescopía contiene un signo
       contrario a una suma de cuadrados;
 \item una extensión no conmuta con \(E_9,J_9,L,G,U\);
 \item se sustituye la multisección de continuaciones por una prolongación
       escogida retrospectivamente;
 \item se omite cualquiera de las trece coordenadas.
\end{enumerate}
En el último caso el diagnóstico vinculante es: \emph{objeto sustituido;
conclusión no transferible}.
\end{criterion}

\begin{warningbox}
Este capítulo concluye únicamente la construcción interna de
\(\mathcal C_{\rm sol}^{\rm HMT}\).  No contiene una ecuación diferencial
clásica, datos iniciales clásicos, un mapa desde datos exteriores ni una
conclusión sobre un problema exterior.  Esos entrelazadores pertenecen a otro
documento y no gobiernan la existencia de esta macroestructura.
\end{warningbox}
\fi
 \endgroup

\section{Orientación, acarreo y composición simultánea de trayectorias}

La inversión de la ruta actúa una sola vez sobre
\(\mathfrak e_\alpha\). Sea
\(\mathsf J(v_\Sigma,v_\Pi):=(v_\Pi,v_\Sigma)\) el intercambio de hojas,
y sea \(P_{\rm tr}\) la transposición de matrices de puertos. Sus cinco
expresiones locales son
\begin{equation}
\label{eq:pdfv2-cor-orientacion-cruzada}
 \begin{gathered}
 \mathsf J\sigma_k=-\sigma_k\mathsf J,
 \qquad b_{\bar\alpha}=-b_\alpha,
 \qquad (b_{\bar\alpha}^+,b_{\bar\alpha}^-)
       =(b_\alpha^-,b_\alpha^+),\\
 \vartheta_L(K,\epsilon):=
 \begin{cases}
  (L,-\epsilon),&K=L,\\
  (K,\epsilon),&K\neq L,
 \end{cases}
 \qquad
 \kappa_{J,I}(v,u)=-P_{\rm tr}\kappa_{I,J}(u,v),\\
 \kappa(\bar\alpha)
 =-\mathsf C_\alpha^{-1}\kappa(\alpha),
 \qquad v_{\bar\alpha}=-v_\alpha.
 \end{gathered}
\end{equation}
Aquí \(\mathsf J\) intercambia las dos hojas, \(P_{\rm tr}\) transpone la matriz de
incidencia y \(\vartheta_L\) intercambia exactamente las dos caras de la
dirección \(L\), fijando las otras seis. En la base libre de
\(\mathscr B=\mathscr I\times\{+,-\}\), los ocho vectores son distintos,
\(\vartheta_L^2=I\) y
\(\vartheta_L\neq\vartheta_{L'}\) para \(L\neq L'\). La misma orientación
explica todos los signos de
\eqref{eq:pdfv2-cor-orientacion-cruzada}.

Para extensiones componibles \(\alpha:x\to y\) y \(\beta:y\to z\), el
carry nativo satisface
\begin{equation}
\label{eq:pdfv2-cor-carry-nativo}
 \kappa(\beta\alpha)=
 \kappa(\beta)+\mathsf C(\beta)\kappa(\alpha).
\end{equation}
La elevación matricial no se deja como un nombre. En la base ordenada
\((f,r,e,b,g)\), sea
\(\widetilde L_N(\alpha)\in\operatorname{Mat}_5(\mathbb Z)\) el
representante entero centrado de la matriz de
\(\Psi_\alpha\) de \eqref{eq:pdfv2-cor-psi-generadores} módulo \(3^N\), y
defínase
\begin{equation}
\label{eq:pdfv2-cor-carry-matricial-def}
 c_N(\beta,\alpha):=
 \frac{\widetilde L_N(\beta)\widetilde L_N(\alpha)
       -\widetilde L_N(\beta\alpha)}{3^N}
 \in\operatorname{Mat}_5(\mathbb Z).
\end{equation}
La divisibilidad procede de la igualdad de las dos composiciones módulo
\(3^N\). Escribiendo \(L_N=[\widetilde L_N]_{3^N}\), los lifts enteros
satisfacen
\begin{equation}
\label{eq:pdfv2-cor-carry-matricial}
 \widetilde L_N(\beta)\widetilde L_N(\alpha)
 =\widetilde L_N(\beta\alpha)+3^Nc_N(\beta,\alpha)
\end{equation}
y la asociatividad produce el cociclo
\begin{equation}
\label{eq:pdfv2-cor-cociclo-matricial}
 c_N(\delta,\beta\alpha)+\widetilde L_N(\delta)c_N(\beta,\alpha)
 =c_N(\delta\beta,\alpha)
  +c_N(\delta,\beta)\widetilde L_N(\alpha).
\end{equation}

El mismo incremento de carry actúa sobre el complejo local mediante
\begin{equation}
\label{eq:pdfv2-cor-psi-generadores}
 \begin{gathered}
 \Psi_\alpha(r_x)=r_y,\quad
 \Psi_\alpha(e_x)=e_y,\quad
 \Psi_\alpha(g_x)=g_y,\quad
 \Psi_\alpha(f_x)=f_y,\\
 \Psi_\alpha(b_x)=b_y+3(c_y-c_x)e_y.
 \end{gathered}
\end{equation}
El cambio de frontera orientada queda en
\begin{equation}
\label{eq:pdfv2-cor-k-alpha}
 K_\alpha:=(v_y-v_x)f_y,
 \qquad
 K_{\beta\alpha}=K_\beta+\Psi_\beta K_\alpha.
\end{equation}
Las ecuaciones \eqref{eq:pdfv2-cor-carry-nativo},
\eqref{eq:pdfv2-cor-cociclo-matricial} y
\eqref{eq:pdfv2-cor-k-alpha} son tres grados del mismo ledger de composición.
Ninguna autoriza a borrar las otras dos.

\subsection{Nervio de prolongaciones y diagonal de Alexander--Whitney}
Para una historia del registro común se escribe
\(A_n:=\mathbf Z/3^n\mathbf Z\) y \(\mathsf T_m(x)\) para la categoría
finita de \emph{todos} sus prefijos TPK de profundidad \(m\), con los
morfismos de extensión ya construidos. La diagonal y la counidad que siguen
se aplican a ese mismo nervio; no definen una quinta rama autónoma.
\begingroup
  \let\section\subsection
  % Extracción quirúrgica del fuente teoremática V2 05e: líneas 68--110.
% El resto permanece en este archivo como procedencia, pero no se publica.
\iffalse
\chapter[Macroestructura determinantal interna]{Macroestructura determinantal interna generada por el continuo discreto}
\label{chap:detint-macroestructura}

\begin{thesisbox}
La rama \({\rm det}\) parte exclusivamente de historias supervivientes de
\(\mathcal C_{\rm cont}^{\rm disc}\). El nervio de extensiones, la diagonal
Alexander--Whitney, el complejo local, la unidad, las dos publicaciones, el
filler, la reducción terminal y la torre \(3\)-ádica se construyen antes de
cualquier evaluación posterior. Las condiciones de balance y aciclicidad se
incorporan al dominio tipado y nunca se presuponen silenciosamente.
\end{thesisbox}

\section{Dominio determinantal superviviente}

Sea
\begin{equation}\label{eq:detint-an}
 A_n:=\mathbf Z/3^n\mathbf Z,
 \qquad
 \rho_{n+1,n}:A_{n+1}\longrightarrow A_n.
\end{equation}
Para \(x\in\mathcal C_{\rm cont}^{\rm disc}\), sea
\(\mathsf T_m(x)\) la categoría finita de prefijos TPK de profundidad \(m\)
compatibles con \(x\). Cada objeto conserva hoja, orientación, residuo,
cociente, carry entero, memoria, frontera, ruta, multisección, supervivencia
y terminal; sus morfismos son las extensiones admisibles que conservan esas
coordenadas.

De cada objeto \(y\in\mathsf T_m(x)\) se extraen los enteros internos
\begin{equation}\label{eq:detint-c-v-lifts}
 \widetilde c_y,\widetilde v_y\in\mathbf Z,
 \qquad
 c_{y,n}:=\widetilde c_y\bmod3^n,
 \qquad
 v_{y,n}:=\widetilde v_y\bmod3^n.
\end{equation}
La primera coordenada es el carry no orientado y la segunda, el coeficiente
orientado de frontera. La reversión satisface
\begin{equation}\label{eq:detint-or-cv}
 c_{\bar y,n}=c_{y,n},
 \qquad
 v_{\bar y,n}=-v_{y,n}.
\end{equation}

Se escriben
\begin{equation}\label{eq:detint-z3-k3}
 \mathbf Z_3:=\varprojlim_n A_n,
 \qquad
 K_3:=\operatorname{Frac}(\mathbf Z_3).
\end{equation}
La restricción determinantal se define por
\begin{equation}\label{eq:detint-dominio}
\begin{split}
 \mathcal C_{\rm cont}^{\rm det}:=
 \bigl\{x\in\mathcal C_{\rm cont}^{\rm disc}:{}&
 x\text{ sobrevive a toda profundidad};\\
 &\widehat P_x^{\rm red}\text{ es finito y basado};\\
 &\chi(\widehat P_x^{\rm red})=0;\\
 &H_*(\widehat P_x^{\rm red}\otimes_{\mathbf Z_3}K_3)=0;\\
 &\text{la reducción canónica produce un bloque cuadrado}\\
 &\text{estable bajo refinamiento}\bigr\}.
\end{split}
\end{equation}
Los objetos \(\widehat P_x^{\rm red}\) y el bloque cuadrado se construirán
en \cref{sec:detint-terminal}. Por tanto
\eqref{eq:detint-dominio} es una condición de pertenencia comprobable y no
una conclusión insertada en la definición de los mapas.

\fi
\section{Nervio de rutas, diagonal de Alexander--Whitney y counidad}

Se define el complejo normalizado
\begin{equation}\label{eq:detint-gmn}
 G_{m,n}(x):=
 N_*^{\rm norm}\!\left(N\mathsf T_m(x);A_n\right).
\end{equation}
Para un símplice no degenerado
\(\sigma=[y_0\to y_1\to\cdots\to y_q]\), la diagonal es
\begin{equation}\label{eq:detint-aw}
 \Delta_{\rm AW}(\sigma)
 =\sum_{j=0}^q
 [y_0\to\cdots\to y_j]\otimes
 [y_j\to\cdots\to y_q].
\end{equation}
La counidad viene dada por
\begin{equation}\label{eq:detint-counidad-g}
 \epsilon_G([y])=1,
 \qquad
 \epsilon_G(\sigma)=0\quad(q>0).
\end{equation}
En particular,
\begin{equation}\label{eq:detint-vertice-grouplike}
 \Delta_{\rm AW}[y]=[y]\otimes[y],
 \qquad
 \epsilon_G([y])=1.
\end{equation}

Si \(m'\ge m\) y \(n'\ge n\), truncación y reducción de coeficientes
inducen
\begin{equation}\label{eq:detint-rmn}
 R_{m',m}^{n',n}:G_{m',n'}(x)\longrightarrow G_{m,n}(x),
\end{equation}
y se verifican
\begin{equation}\label{eq:detint-aw-natural}
 R\partial=\partial R,
 \qquad
 (R\otimes R)\Delta_{\rm AW}=\Delta_{\rm AW}R,
 \qquad
 \epsilon_GR=\rho_{n',n}\epsilon_G.
\end{equation}
La ruta no se reduce a sus extremos: el símplice entero permanece en
\(G_{m,n}(x)\).

\iffalse
\section{Complejo local interno}

Para \(c\in A_n\), se define
\begin{equation}\label{eq:detint-p0-grados}
 P^0_{n,1}(c)=A_nf,
 \qquad
 P^0_{n,0}(c)=A_nr\oplus A_ne\oplus A_nb,
 \qquad
 P^0_{n,-1}(c)=A_ng,
\end{equation}
con diferencial
\begin{equation}\label{eq:detint-p0-diferencial}
 \partial f=3ce+b,
 \qquad
 \partial r=0,
 \qquad
 \partial e=g,
 \qquad
 \partial b=-3cg.
\end{equation}
La identidad de cadena es literal:
\begin{equation}\label{eq:detint-d2}
 \partial^2f=3c\,\partial e+\partial b=3cg-3cg=0,
 \qquad
 \partial^2e=\partial^2b=\partial^2r=0.
\end{equation}

La aumentación
\begin{equation}\label{eq:detint-epsilon-p}
 \epsilon_P:P_n^0(c)\longrightarrow A_n[0]
\end{equation}
se fija por
\[
 \epsilon_P(r)=1,
 \qquad
 \epsilon_P(e)=\epsilon_P(b)=\epsilon_P(f)=\epsilon_P(g)=0.
\]
El complejo posee la coalgebra diferencial
\begin{equation}\label{eq:detint-coalgebra-p}
 \Delta_P(r)=r\otimes r,
 \qquad
 \Delta_P(q)=q\otimes r+r\otimes q
 \quad(q\in\{e,b,f,g\}).
\end{equation}
Así, \(r\) es group-like y los otros cuatro generadores son primitivos
relativos a \(r\). Las ecuaciones
\eqref{eq:detint-p0-diferencial} prueban que \(\Delta_P\) conmuta con el
diferencial tensorial.

\section{Sistema local de extensiones y ledger de memoria}

Para una extensión \(\alpha:y\to y'\), definimos
\begin{equation}\label{eq:detint-psi}
 \Psi_\alpha:P_n^0(c_y)\longrightarrow P_n^0(c_{y'})
\end{equation}
mediante
\begin{equation}\label{eq:detint-psi-generadores}
\begin{gathered}
 \Psi_\alpha(r_y)=r_{y'},\quad
 \Psi_\alpha(e_y)=e_{y'},\quad
 \Psi_\alpha(g_y)=g_{y'},\quad
 \Psi_\alpha(f_y)=f_{y'},\\
 \Psi_\alpha(b_y)=b_{y'}+3(c_{y'}-c_y)e_{y'}.
\end{gathered}
\end{equation}
Es un mapa de complejos porque
\begin{equation}\label{eq:detint-psi-b-chain}
 \partial\Psi_\alpha(b_y)
 =-3c_{y'}g+3(c_{y'}-c_y)g
 =-3c_yg
 =\Psi_\alpha(\partial b_y),
\end{equation}
y
\begin{equation}\label{eq:detint-psi-f-chain}
 \Psi_\alpha(\partial f_y)
 =3c_ye+b+3(c_{y'}-c_y)e
 =3c_{y'}e+b
 =\partial\Psi_\alpha(f_y).
\end{equation}
Además,
\begin{equation}\label{eq:detint-psi-functorial}
 \Psi_\beta\Psi_\alpha=\Psi_{\beta\alpha}.
\end{equation}

El cambio de la coordenada orientada no se borra. Se registra por
\begin{equation}\label{eq:detint-k-alpha}
 K_\alpha:=(v_{y'}-v_y)f_{y'}.
\end{equation}
Para extensiones componibles,
\begin{equation}\label{eq:detint-k-coherencia}
 K_{\beta\alpha}=K_\beta+\Psi_\beta K_\alpha.
\end{equation}
Esta igualdad es el ledger exacto del incremento de frontera y hace visible
la memoria que una composición puramente nominal perdería.

El diagrama local completo es la construcción de Grothendieck
\begin{equation}\label{eq:detint-grothendieck-local}
 \int_{y\in\mathsf T_m(x)}P_n^0(c_y).
\end{equation}
En ella una flecha sobre \(\alpha:y\to y'\) lleva como componente lineal
\(\Psi_\alpha\) y como coherencia \(K_\alpha\). Cuando se aplica a un
elemento soportado en un vértice, se usa la notación
\begin{equation}\label{eq:detint-psi-soporte-vertice}
 \Psi_\alpha(p,[y]):=(\Psi_\alpha p,[y']).
\end{equation}
La arista \([y\to y']\) permanece en la coordenada de ruta del nervio. Esta
convención será la que se emplee en las fórmulas de naturalidad de
\(z_\pm\) y \(h_*\); no se afirma un mapa global que borre los restantes
símplices de \(G_{m,n}(x)\).

\section{Pullback, unidad, raíz, publicaciones y filler}

El pullback de complejos es
\begin{equation}\label{eq:detint-pgen}
 P^{\rm gen}_{y;m,n}
 :=P_n^0(c_y)\times_{A_n[0]}G_{m,n}(x)
 =\{(p,\gamma):\epsilon_P(p)=\epsilon_G(\gamma)\}.
\end{equation}
Su diferencial es
\begin{equation}\label{eq:detint-pgen-d}
 \partial(p,\gamma)=(\partial p,\partial\gamma).
\end{equation}
La unidad común asociada al vértice \(y\) es
\begin{equation}\label{eq:detint-unidad}
 \mathbf1_y:=(r,[y]).
\end{equation}
Sus dos proyecciones son group-like por
\eqref{eq:detint-vertice-grouplike} y
\eqref{eq:detint-coalgebra-p}; ése es el sentido preciso de unidad
group-like compatible del pullback.

La proyección raíz es el mapa de complejos
\begin{equation}\label{eq:detint-proyeccion-root}
 \varpi_{\rm root}:P^{\rm gen}_{y;m,n}\longrightarrow
 \operatorname{Root}_{y;m,n},
 \qquad
 \varpi_{\rm root}(p,\gamma):=(\epsilon_P(p)r,\gamma).
\end{equation}
Se definen las dos publicaciones internas
\begin{equation}\label{eq:detint-zpm}
 z_-(y):=(r,[y]),
 \qquad
 z_+(y):=(r+3c_yv_ye+v_yb,[y]).
\end{equation}
Ambas son ciclos:
\begin{equation}\label{eq:detint-zpm-ciclos}
 \partial z_-(y)=0,
 \qquad
 \partial z_+(y)=3c_yv_yg-3c_yv_yg=0.
\end{equation}
Sus raíces coinciden:
\begin{equation}\label{eq:detint-raiz-comun}
 \varpi_{\rm root}z_-(y)
 =(r,[y])
 =\varpi_{\rm root}z_+(y).
\end{equation}

El filler se construye directamente:
\begin{equation}\label{eq:detint-hstar}
 h_*(y):=(-v_yf,0)\in P^{\rm gen}_{y;m,n,1}.
\end{equation}
Entonces
\begin{equation}\label{eq:detint-filler-identidad}
 \partial h_*(y)
 =(-3c_yv_ye-v_yb,0)
 =z_-(y)-z_+(y).
\end{equation}
En particular,
\begin{equation}\label{eq:detint-clases-iguales}
 [z_-(y)]=[z_+(y)]
 \quad\text{en}\quad H_0(P^{\rm gen}_{y;m,n}),
\end{equation}
y la igualdad conserva su cadena testigo.

Bajo una extensión \(\alpha:y\to y'\), se cumplen
\begin{equation}\label{eq:detint-z-natural}
 z_-(y')=\Psi_\alpha z_-(y),
 \qquad
 z_+(y')-\Psi_\alpha z_+(y)=\partial K_\alpha,
\end{equation}
y
\begin{equation}\label{eq:detint-h-natural}
 h_*(y')-\Psi_\alpha h_*(y)=-K_\alpha.
\end{equation}
Por tanto, la naturalidad no identifica valores distintos de \(v\): registra
su diferencia mediante la homotopía coherente \(K_\alpha\).

\section{Orientación interna}

Sobre \(P_n^0(c)\) se define
\begin{equation}\label{eq:detint-omega-p}
 \Omega(r)=r,
 \qquad
 \Omega(q)=-q\quad(q\in\{e,b,f,g\}).
\end{equation}
Es un automorfismo de complejos y de la coalgebra relativa. En el nervio,
la inversión orientada es
\begin{equation}\label{eq:detint-omega-g}
 \mathfrak o_G[y_0\to\cdots\to y_q]
 :=(-1)^{q(q+1)/2}
 [\bar y_q\to\cdots\to\bar y_0].
\end{equation}
Junto con \eqref{eq:detint-or-cv}, estas acciones dan
\begin{equation}\label{eq:detint-or-z-h}
 \Omega z_-(y)=z_-(\bar y),
 \qquad
 \Omega z_+(c_y,v_y)=z_+(c_{\bar y},v_{\bar y}),
 \qquad
 \Omega h_*(y)=h_*(\bar y).
\end{equation}
La orientación actúa sobre ruta, complejo, publicaciones, filler y línea
determinantal; no es una etiqueta sin acción.

\section{Reducción terminal y condición de supervivencia}
\label{sec:detint-terminal}

Se elimina únicamente la raíz común:
\begin{equation}\label{eq:detint-p-reducido}
 \overline P_{y;m,n}
 :=P^{\rm gen}_{y;m,n}/A_n\mathbf1_y,
 \qquad
 \widehat P_{y;m}^{\rm red}:=\varprojlim_n\overline P_{y;m,n}.
\end{equation}
La base se ordena primero por el orden TPK de rutas y después por
\begin{equation}\label{eq:detint-orden-base}
 f\prec e\prec b\prec g.
\end{equation}
Se ejecuta el siguiente algoritmo interno.
\begin{enumerate}
 \item Buscar los coeficientes que son unidades de \(\mathbf Z_3\).
 \item Elegir el primero según \eqref{eq:detint-orden-base} y el orden de
 rutas.
 \item Cancelar el par correspondiente y registrar las operaciones
 elementales, su signo y su contribución a la orientación.
 \item Repetir hasta que no quede un pivote unitario cancelable.
\end{enumerate}
La condición de supervivencia de \eqref{eq:detint-dominio} exige que,
cofinalmente en \(m\), el resultado sea un bloque de dos términos
\begin{equation}\label{eq:detint-sx}
 \mathbf Z_3^{r_x}\xrightarrow{S_x}\mathbf Z_3^{r_x},
 \qquad
 \det S_x\ne0\text{ en }K_3,
\end{equation}
y que un refinamiento sólo añada bloques contractibles unitarios y cambios
de base registrados. Si esas condiciones fallan, la historia no pertenece a
\(\mathcal C_{\rm cont}^{\rm det}\); no se fuerza un terminal inexistente.

La ruta identidad aporta un testigo elemental. Después de retirar la raíz,
el complejo local es
\begin{equation}\label{eq:detint-testigo-contractil}
 A_nf\xrightarrow{\binom{3c}{1}}
 A_ne\oplus A_nb\xrightarrow{(1,-3c)}A_ng.
\end{equation}
Es contractible mediante
\begin{equation}\label{eq:detint-contraccion-local}
 s(g)=e,
 \qquad
 s(ae+\beta b)=\beta f,
 \qquad
 s(f)=0.
\end{equation}
Así, las condiciones terminales admiten al menos el testigo identidad
interno y no describen un subdominio formalmente vacío.

\section{Smith, orden \texorpdfstring{\(3\)}{3}-ádico y unidad inicial}

En bases orientadas, una diagonalización por valoración tiene la forma
\begin{equation}\label{eq:detint-smith}
 U_xS_xV_x
 =\operatorname{diag}(3^{a_1}u_1,\ldots,3^{a_r}u_r),
 \qquad
 a_i\in\mathbf N,quad u_i\in\mathbf Z_3^\times.
\end{equation}
Se define
\begin{equation}\label{eq:detint-orden-d}
 d_x:=v_3(\det S_x)=\sum_{i=1}^r a_i
\end{equation}
y la unidad inicial
\begin{equation}\label{eq:detint-leading-unit}
 \lambda_x:=3^{-d_x}\det S_x\in\mathbf Z_3^\times.
\end{equation}
La coordenada de orientación trivializa la línea determinantal. Si una
transición produce
\begin{equation}\label{eq:detint-estabilidad-bloques}
 S_x\oplus I_q=U S_{x'}V,
\end{equation}
la orientación se transporta por el factor inverso de \(\det U\det V\).
La unidad orientada \(\lambda_x^{\rm or}\) es entonces natural y no cambia
al añadir bloques contractibles.

En precisión finita,
\begin{equation}\label{eq:detint-smith-finito}
 S_{x,n}=S_x\bmod3^n,
 \qquad
 a_i^{(n)}=\min(a_i,n).
\end{equation}
El terminal es, por tanto, la familia compatible de todas las precisiones y
no un entero aislado.

\section{Torre de Bockstein y acarreo de orden superior}

Sea
\[
 C_x^1\xrightarrow{S_x}C_x^0
\]
el bloque terminal. Si una clase \([\bar y]\) admite un lift
\(y\in C_x^1\) con \(S_xy\in3^qC_x^0\), definimos
\begin{equation}\label{eq:detint-beta-q}
 \beta_q[\bar y]
 :=\left[3^{-q}S_xy\bmod3\right].
\end{equation}
El cambio de lift altera el representante por una frontera de la página
correspondiente, de modo que \(\beta_q\) está bien definido. El carry
superior es
\begin{equation}\label{eq:detint-carry-superior}
 \operatorname{Carr}_q(y):=3^{-q}S_xy.
\end{equation}
No sustituye el carry TPK de \eqref{eq:detint-c-v-lifts}; registra su
prolongación en la torre de coeficientes.

Para un bloque \(3^{a_i}u_i\), la clase persiste hasta la página \(a_i\),
y en esa página el diferencial no nulo es multiplicación por
\(\bar u_i\in\mathbf F_3^\times\). Por consiguiente,
\begin{equation}\label{eq:detint-bock-longitudes}
 \{\text{longitudes de supervivencia}\}=\{a_1,\ldots,a_r\},
 \qquad
 d_x=\sum_i a_i.
\end{equation}
El producto orientado de las trivializaciones de página satisface
\begin{equation}\label{eq:detint-bock-leading}
 \Theta_{\rm Bock}(x)=\overline{\lambda_x^{\rm or}}
 \in\mathbf F_3^\times.
\end{equation}
La prueba se hace en cada bloque diagonal y se transporta por los cambios
unimodulares registrados. Smith, Bockstein y unidad inicial son así tres
lecturas del mismo terminal.

\section{Las trece coordenadas de la restricción determinantal}

Para cada \((m,n)\), se define
\begin{equation}\label{eq:detint-d-trece}
 D_{{\rm det},m,n}(x):=
 (\mathcal F,\operatorname{Ext}_\Gamma,\operatorname{Rel},
 \mathcal N,\operatorname{or},\operatorname{Carr},\operatorname{Mem},
 \partial,\gamma,\operatorname{Msect},\operatorname{Surv},
 \operatorname{Term},\mathcal L)_{{\rm det},m,n}(x),
\end{equation}
con el contenido siguiente.
\begin{enumerate}
 \item La fibra es
 \[
  \mathcal F=(A_n,\mathsf T_m(x),
  \{P_n^0(c_y)\}_y,\{P^{\rm gen}_{y;m,n}\}_y).
 \]
 \item \(\operatorname{Ext}_\Gamma\) contiene todos los morfismos
 \(\alpha\), los mapas \(\Psi_\alpha\), las homotopías \(K_\alpha\) y
 los mapas \(R_{m',m}^{n',n}\).
 \item \(\operatorname{Rel}\) contiene \(\partial^2=0\), la ecuación de
 pullback, AW, counidad, \(\Psi_\beta\Psi_\alpha=\Psi_{\beta\alpha}\) y
 \(K_{\beta\alpha}=K_\beta+\Psi_\beta K_\alpha\).
 \item
 \(\mathcal N=(m,n,\widetilde c_y,\widetilde v_y,c_{y,n},v_{y,n},
 r_x,a_1,\ldots,a_{r_x},d_x)\).
 \item \(\operatorname{or}\) contiene \(y\mapsto\bar y\),
 \(\mathfrak o_G\), \(\Omega\) y la orientación de la línea
 determinantal.
 \item \(\operatorname{Carr}\) conserva separadamente el carry entero,
 su reducción \(c_{y,n}\) y todos los carries superiores
 \(\operatorname{Carr}_q\).
 \item \(\operatorname{Mem}\) contiene
 \[
  y_0\to\cdots\to y_m,
  \quad\{(\widetilde c_{y_j},\widetilde v_{y_j},\mathbf1_{y_j})\}_{j=0}^m,
  \quad\{K_{\alpha_j}\}.
 \]
 \item
 \(\partial=(\partial_G,\partial_P,\partial_{P^{\rm gen}},
 \partial_{\overline P},h_*)\).
 \item \(\gamma=[y_0\to\cdots\to y_q]\) conserva el símplice orientado
 completo.
 \item \(\operatorname{Msect}\) contiene todas las extensiones
 compatibles, todos los lifts \(3\)-ádicos y todas las presentaciones
 unimodulares del mismo terminal orientado; no elige una a posteriori.
 \item \(\operatorname{Surv}\) registra profundidad arbitraria,
 compatibilidad de coeficientes, balance, aciclicidad sobre \(K_3\),
 estabilidad por bloques unitarios y permanencia de la raíz.
 \item
 \[
  \operatorname{Term}=(\widehat P_x^{\rm red},S_x,
  \{a_i,u_i\},d_x,\lambda_x^{\rm or},
  \{\beta_q\},\Theta_{\rm Bock}).
 \]
 \item El lector interno es
 \[
  \mathcal L_{\rm det}^{\rm int}
  =(\mathbf1_y,
  \varpi_{\rm root}z_-=\varpi_{\rm root}z_+,
  [z_-]=[z_+],d_x,\lambda_x^{\rm or},\Theta_{\rm Bock}).
 \]
\end{enumerate}

Las transiciones satisfacen, coordenada por coordenada,
\begin{equation}\label{eq:detint-naturalidad-trece}
 u^{\rm det}_{(m',n'),(m,n)}D_{{\rm det},m',n'}(x')
 =D_{{\rm det},m,n}(\tau_{m',m}x').
\end{equation}

\section{Ensamblador, publicador y cadena no ablativa}

Se define
\begin{equation}\label{eq:detint-graph-d}
 \mathcal G_{\rm det}:=\operatorname{Graph}(D_{\rm det}),
 \qquad
 \operatorname{Res}_{\rm det}(x):=(x,D_{\rm det}x).
\end{equation}
Sea \(\chi_{\rm det}(x)\) la multisección completa de gérmenes
\[
 (y,\widetilde c_y,\widetilde v_y,\operatorname{or}_y,\gamma_y)
\]
compatibles a toda profundidad. El objeto dependiente es
\begin{equation}\label{eq:detint-x-dependiente}
 X_{\chi,{\rm det}}
 :=\{(\chi_{\rm det}(x),x,D_{\rm det}x):
 x\in\mathcal C_{\rm cont}^{\rm det}\},
\end{equation}
y
\begin{equation}\label{eq:detint-prx}
 \operatorname{pr}_{X,{\rm det}}(x,D_{\rm det}x)
 :=(\chi_{\rm det}(x),x,D_{\rm det}x).
\end{equation}

El ensamblador crudo
\begin{equation}\label{eq:detint-ensamblador-crudo}
 a_{\rm det}:X_{\chi,{\rm det}}\longrightarrow
 \mathscr M_{\rm det}^{\rm amb}
\end{equation}
produce
\begin{equation}\label{eq:detint-ensamblador-salida}
\begin{split}
 a_{\rm det}(z)=\bigl(&
 \{G_{m,n},\Delta_{\rm AW},\epsilon_G\}_{m,n};\\
 &\{P_n^0(c),\Psi_\alpha,K_\alpha,P^{\rm gen},
 \mathbf1,z_-,z_+,h_*\}_{m,n};\\
 &\{\widehat P^{\rm red},S,d,\lambda^{\rm or},
 \beta_q,\Theta_{\rm Bock}\}\bigr).
\end{split}
\end{equation}
El macroobjeto conserva su entrada:
\begin{equation}\label{eq:detint-graph-a}
 \mathfrak M_{\rm det}:=\operatorname{Graph}(a_{\rm det}),
 \qquad
 \mathcal A_{\rm det}(z):=(z,a_{\rm det}z).
\end{equation}

El publicador crudo
\begin{equation}\label{eq:detint-publicador-crudo}
 p_{\rm det}:\mathfrak M_{\rm det}\longrightarrow\mathscr Y_{\rm det}
\end{equation}
publica conjuntamente el certificado
\begin{equation}\label{eq:detint-certificado-publicado}
\begin{gathered}
 \partial^2=0,quad \text{AW y counidad},quad
 \mathbf1_y\text{ group-like compatible},\\
 \partial z_\pm=0,quad
 \varpi_{\rm root}z_-=\varpi_{\rm root}z_+,quad
 \partial h_*=z_--z_+,\\
 \Psi_\beta\Psi_\alpha=\Psi_{\beta\alpha},quad
 K_{\beta\alpha}=K_\beta+\Psi_\beta K_\alpha,\\
 d=v_3(\det S)=\sum_i a_i,quad
 \Theta_{\rm Bock}=\overline{\lambda^{\rm or}},quad
 \text{naturalidad completa bajo refinamiento}.
\end{gathered}
\end{equation}
Finalmente,
\begin{equation}\label{eq:detint-graph-p}
 \mathcal C_{\rm det}^{\rm HMT}:=\operatorname{Graph}(p_{\rm det}),
 \qquad
 \mathcal P_{\rm det}(M):=(M,p_{\rm det}M).
\end{equation}
La cadena íntegra es
\begin{equation}\label{eq:detint-cadena-completa}
 \mathcal C_{\rm cont}^{\rm disc}\supset
 \mathcal C_{\rm cont}^{\rm det}
 \xrightarrow{\operatorname{Res}_{\rm det}}\mathcal G_{\rm det}
 \xrightarrow{\operatorname{pr}_{X,{\rm det}}}X_{\chi,{\rm det}}
 \xrightarrow{\mathcal A_{\rm det}}\mathfrak M_{\rm det}
 \xrightarrow{\mathcal P_{\rm det}}\mathcal C_{\rm det}^{\rm HMT}.
\end{equation}
Cada flecha tiene una inversa sobre su imagen dada por una proyección de
coordenadas. No se sustituye historia por valor, multisección por selector,
complejo por determinante ni prueba por nombre.

\section{Teorema de realización determinantal interna}

\begin{theorem}[Realización determinantal interna completa]
\label{thm:detint-realizacion}
Para todo \(x\in\mathcal C_{\rm cont}^{\rm det}\), la cadena
\eqref{eq:detint-cadena-completa} construye un macroobjeto natural en
profundidad y precisión. Contiene dos ciclos con raíz común, un filler
explícito que los identifica, una unidad group-like compatible y un terminal
\(3\)-ádico cuyo orden, forma de Smith y torre de Bockstein producen el mismo
lector orientado. Las trece coordenadas de \eqref{eq:detint-d-trece}
permanecen recuperables en la salida.
\end{theorem}

\begin{proof}
Las ecuaciones \eqref{eq:detint-p0-diferencial}--
\eqref{eq:detint-d2} prueban que \(P_n^0(c)\) es un complejo. Las dos
aumentaciones son mapas de complejos, por lo que el pullback está tipado.
Alexander--Whitney es natural y hace group-like a cada vértice; junto con
\(r\), esto da la unidad común. El cálculo directo
\eqref{eq:detint-zpm-ciclos}--\eqref{eq:detint-filler-identidad} prueba que
las publicaciones son ciclos, tienen la misma raíz y están unidas por
\(h_*\). Las fórmulas de \(\Psi_\alpha\) prueban functorialidad, y
\(K_\alpha\) prueba la naturalidad homotópica conservando el incremento de
memoria. La definición de \(\mathcal C_{\rm cont}^{\rm det}\) garantiza que
la reducción terminal produce un bloque cuadrado sin postularlo fuera del
dominio. Smith sobre \(\mathbf Z_3\) da \(d=\sum_i a_i\); el cálculo de cada
bloque \(3^{a_i}u_i\) da
\(\Theta_{\rm Bock}=\overline{\lambda^{\rm or}}\). Truncación, reducción y
adición de bloques unitarios conservan estas identidades. Por último, las
tres construcciones por graph retienen la entrada como primera coordenada,
de modo que la composición completa es no ablativa.
\end{proof}

\section{Procedencia y falsadores}

\paragraph{Procedencia.}
El nervio TPK, Alexander--Whitney, la counidad, la unidad común, la forma
normal \(\partial f=3ce+b\), \(\partial e=g\),
\(\partial b=-3cg\), la raíz común, el filler y la arquitectura de torre
tienen estatuto \texttt{ARQUITECTURA\_AUTORAL\_PREEXISTENTE}. El cálculo
literal de los dos ciclos, su raíz y la identidad
\(\partial h_*=z_--z_+\) tiene estatuto \texttt{RESULTADO\_RECUPERADO}. El
sistema local \(\Psi_\alpha/K_\alpha\), el subdominio superviviente
balanceado, el terminal exclusivamente \(3\)-ádico y el empaquetado por
graphs tienen estatuto \texttt{FORMALIZACION\_NUEVA}; formalizan sin cambiar
la procedencia conceptual del aparato.

\begin{ablation}[Falsadores internos de la rama determinantal]
\label{abl:detint-falsadores}
La construcción queda falsada en el nivel correspondiente si ocurre alguno
de los hechos siguientes.
\begin{enumerate}
 \item Falla \(\partial^2=0\); entonces no existe el complejo local.
 \item Alexander--Whitney no conmuta con refinamiento, la counidad falla o
 el vértice deja de ser group-like.
 \item \(z_\pm\) no son ciclos, sus raíces difieren o falla
 \(\partial h_*=z_--z_+\).
 \item \(\Psi_\alpha\) no es mapa de complejos o falla
 \(K_{\beta\alpha}=K_\beta+\Psi_\beta K_\alpha\); en ese caso se ha
 perdido carry o memoria.
 \item El complejo reducido no es balanceado o no es acíclico sobre
 \(K_3\); la historia queda fuera del dominio terminal tipado.
 \item Las longitudes de Bockstein no coinciden con los \(a_i\), o
 \(\Theta_{\rm Bock}\ne\overline{\lambda^{\rm or}}\).
 \item El refinamiento modifica \(d\) o \(\lambda^{\rm or}\) fuera de
 bloques unitarios y cambios de orientación registrados.
 \item Se omite una de las trece coordenadas: \emph{objeto sustituido;
 conclusión no transferible}.
\end{enumerate}
Estos falsadores comprueban el objeto interno construido; no lo reemplazan
por un control posterior.
\end{ablation}
\fi
 \endgroup

\section{Incidencia celular y relleno correlativo de una misma trayectoria}

La sucesión \eqref{eq:pdfv2-cor-sucesion-celda} asigna a cada prefijo el
complejo de dos términos
\begin{equation}
\label{eq:pdfv2-cor-cono-celda}
 C_{k,N}(\widehat x):=
 \left[A_N^{I_k}\oplus A_N^{J_k}
 \xrightarrow{\kappa_{k,N}}A_N^{I_k\times J_k}\right].
\end{equation}
Un refinamiento actúa en los puertos por truncamiento de la extensión y en
los coeficientes por \(A_{N'}\twoheadrightarrow A_N\). Cuando
\(\alpha_u:I_{k'}\to I_k\) y \(\beta_u:J_{k'}\to J_k\) son esos mapas, su
acción en generadores es
\begin{equation}
\label{eq:pdfv2-cor-refinamiento-celda}
 (p,q)\longmapsto(p\circ\alpha_u,q\circ\beta_u),
 \qquad
 w\longmapsto w\circ(\alpha_u\times\beta_u),
\end{equation}
y verifica
\begin{equation}
\label{eq:pdfv2-cor-refinamiento-cadena}
 \kappa_{k',N}\Gamma(u)_1=\Gamma(u)_0\kappa_{k,N}.
\end{equation}

La incidencia no se fabrica etiquetando cuatro copias de una emisión neutra.
Se genera por la completación celular libre del módulo de las dos hojas que
ya porta el mismo estado. Para
\[
 V_{\widehat x}:=\mathbf F_3e_{\Sigma,\widehat x}
 \oplus\mathbf F_3e_{\Pi,\widehat x},
\]
las cuatro rectas internas son
\begin{equation}
\label{eq:pdfv2-cor-cuatro-rectas-tpk}
 \mathscr I_{\widehat x}:=
 \{L_0,L_\infty,L_+,L_-\}
 =\mathbf P^1(V_{\widehat x}),\qquad
 \begin{aligned}
 L_0&=\langle e_\Sigma\rangle,&
 L_\infty&=\langle e_\Pi\rangle,\\
 L_+&=\langle e_\Sigma+e_\Pi\rangle,&
 L_-&=\langle e_\Sigma-e_\Pi\rangle.
 \end{aligned}
\end{equation}
Si \(v_L\) es, respectivamente,
\((1,0)^t,(0,1)^t,(1,1)^t,(1,-1)^t\), la acción de la dirección \(L\)
sobre los generadores de hoja es el proyector
\begin{equation}
\label{eq:pdfv2-cor-proyectores-cuatro-direcciones}
 P_L:=v_L(v_L^tv_L)^{-1}v_L^t,
\end{equation}
es decir, sobre \(\mathbf F_3\),
\[
 P_0=\begin{pmatrix}1&0\\0&0\end{pmatrix},\quad
 P_\infty=\begin{pmatrix}0&0\\0&1\end{pmatrix},\quad
 P_+=\begin{pmatrix}2&2\\2&2\end{pmatrix},\quad
 P_-=\begin{pmatrix}2&1\\1&2\end{pmatrix}.
\]
El cálculo directo da
\begin{equation}
\label{eq:pdfv2-cor-proyectores-no-colapso}
 P_L^2=P_L,\qquad \operatorname{rk}P_L=1,\qquad
 \operatorname{Im}P_L=L,\qquad
 P_L=P_{L'}\Longrightarrow L=L'.
\end{equation}
Por tanto las cuatro direcciones son salidas distintas del módulo
\(V_{\widehat x}\); no son cuatro nombres puestos sobre la misma acción.

\begin{definition}[Completación celular libre de las hojas TPK]
\label{def:pdfv2-cor-completacion-celular-tpk}
Sea \(\operatorname{Cell}_{\rm TPK}(\widehat x)\) el complejo libre
\[
 C_2(\widehat x)\xrightarrow{\partial_2}
 C_1(\widehat x)\xrightarrow{\partial_1}C_0(\widehat x)
\]
con
\begin{align*}
 C_0(\widehat x)
 &:=
 \mathbf F_3[v_\star]\oplus
 \bigoplus_{L\in\mathscr I}\mathbf F_3[v_L]\oplus
 \bigoplus_{(L,\epsilon)\in\mathscr B}
       \mathbf F_3[v_{L,\epsilon}]\oplus
 \bigoplus_{E\in\mathscr E}\mathbf F_3[v_E],\\
 C_1(\widehat x)
 &:=
 \bigoplus_{L\in\mathscr I}\mathbf F_3[u_L]\oplus
 \bigoplus_{(L,\epsilon)\in\mathscr B}
       \mathbf F_3[u_{L,\epsilon}]\oplus
 \bigoplus_{E=\{L,L'\}\in\mathscr E}
 \bigl(\mathbf F_3[u_{L'\mid L}]
       \oplus\mathbf F_3[u_{L\mid L'}]\bigr),\\
 C_2(\widehat x)
 &:=
 \bigoplus_{E\in\mathscr E}\mathbf F_3[\omega_E].
\end{align*}
Su acción en generadores es
\begin{align}
 \partial u_L&=v_L-v_\star,
 &
 \partial u_{L'\mid L}&=v_E-v_L,
 \label{eq:pdfv2-cor-cell-frontera-aristas}\\
 \partial u_{L,\epsilon}&=v_{L,\epsilon}-v_L,
 &
 \partial\omega_{\{L,L'\}}
 &=u_L+u_{L'\mid L}-u_{L'}-u_{L\mid L'}.
 \label{eq:pdfv2-cor-cell-frontera-celda}
\end{align}
\end{definition}

Los ocho generadores \(v_{L,\epsilon}\) realizan materialmente la pantalla
\(\mathscr B\). La involución anterior se extiende al complejo por
\begin{equation}
\label{eq:pdfv2-cor-pantalla-ocho-generadores}
 \vartheta_Lv_{K,\epsilon}:=v_{\vartheta_L(K,\epsilon)},\qquad
 \vartheta_Lu_{K,\epsilon}:=u_{\vartheta_L(K,\epsilon)},
\end{equation}
y fija los demás generadores. Por ser libre la base,
\begin{equation}
\label{eq:pdfv2-cor-pantalla-ocho-no-colapso}
 |\{v_{L,\epsilon}:(L,\epsilon)\in\mathscr B\}|=8,\qquad
 \vartheta_L^2=I,\qquad
 \vartheta_L\vartheta_{L'}=\vartheta_{L'}\vartheta_L.
\end{equation}
Además \(\partial\vartheta_L=\vartheta_L\partial\), pues ambos recorridos
envían \(u_{L,\epsilon}\) a \(v_{L,-\epsilon}-v_L\) y fijan las fronteras
restantes.

Las dos palabras
\[
 p_{L,L'}:=u_{L'\mid L}u_L,\qquad
 p_{L',L}:=u_{L\mid L'}u_{L'}
\]
son caminos diferentes del complejo libre. Su diferencia es la frontera de
\(\omega_{\{L,L'\}}\), y
\begin{equation}
\label{eq:pdfv2-cor-cell-dos-cero}
 \partial^2\omega_E
 =(v_L-v_\star)+(v_E-v_L)
 -(v_{L'}-v_\star)-(v_E-v_{L'})=0.
\end{equation}
La representación de TPK en las hojas se fija por
\[
 \varrho(u_L)=P_L,\qquad
 \varrho(u_{L'\mid L})=P_{L'},\qquad
 \varrho(p_{L,L'})=P_{L'}P_L.
\]
Aunque dos matrices publicadas coincidieran en una especialización, los dos
caminos no colapsan: siguen siendo generadores libres distintos de
\(C_1(\widehat x)\), unidos por una 2--celda explícita.

La coordenada completa de incidencia es el grupo de 2--cocadenas
\begin{equation}
\label{eq:pdfv2-cor-q-incidencia-total}
 Q_{\rm inc}(\widehat x):=
 C^2(\operatorname{Cell}_{\rm TPK}(\widehat x);\mathbf F_3)
 =\bigoplus_{E\in\mathscr E}\mathbf F_3\omega_E^*
 \simeq\mathbf F_3^6.
\end{equation}
La operación interna devuelve todo \(Q_{\rm inc}(\widehat x)\), no un
elemento seleccionado. Cada \(q\) tiene expansión única
\(q=\sum_Eq_E\omega_E^*\). Para
\(a\in\mathbf F_3^{\mathscr I}\), su cambio de potencial es
\begin{equation}
\label{eq:pdfv2-cor-collar-variacion}
 q\longmapsto q+Ba,\qquad (Ba)_{\{L,L'\}}=a_L+a_{L'}.
\end{equation}
Como cada recta pertenece a tres pares,
\[
 s(Ba)=3\sum_{L\in\mathscr I}a_L=0,
\]
y por ello
\begin{equation}
\label{eq:pdfv2-cor-wc-invariante-celular}
 W_c(q+Ba)=W_c(q)\qquad(q\in Q_{\rm inc}(\widehat x)).
\end{equation}
La unidad de la completación sólo ancla el complejo en el estado:
\begin{equation}
\label{eq:pdfv2-cor-cell-unidad}
 \eta^{\rm cell}_{\widehat x}:V_{\widehat x}\longrightarrow
 V_{\widehat x}\otimes C_0(\operatorname{Cell}_{\rm TPK}(\widehat x)),\qquad
 e_{\diamond,\widehat x}\longmapsto
 e_{\diamond,\widehat x}\otimes v_\star,
\end{equation}
y marca el origen \(q=0\); no afirma que la órbita de ese origen sea todo
\(Q_{\rm inc}\).

La construcción es funtorial en el truncamiento común. Si
\(\tau_{\ell,k}\widehat x_\ell=\widehat x_k\), definimos
\begin{align*}
 \tau_V&:V_{\widehat x_\ell}\longrightarrow V_{\widehat x_k},&
 \tau_Ve_{\diamond,\widehat x_\ell}&=e_{\diamond,\widehat x_k},\\
 \tau_a&:\mathbf F_3^{\mathscr I_{\widehat x_\ell}}
          \longrightarrow\mathbf F_3^{\mathscr I_{\widehat x_k}},&
 (\tau_aa)_{\tau_{\mathscr I}L}&=a_L,\\
 \tau_Q&:Q_{\rm inc}(\widehat x_\ell)
          \longrightarrow Q_{\rm inc}(\widehat x_k),&
 (\tau_Qq)_{\tau_{\mathscr E}E}&=q_E,
\end{align*}
donde \(\tau_{\mathscr I}[v]=[\tau_Vv]\) y
\(\tau_{\mathscr E}\{L,L'\}=\{\tau_{\mathscr I}L,
\tau_{\mathscr I}L'\}\). El mapa de cadenas
\(\operatorname{Cell}(\tau):C_\bullet(\widehat x_\ell)
\to C_\bullet(\widehat x_k)\) actúa en generadores por
\begin{equation}
\label{eq:pdfv2-cor-cell-truncamiento-generadores}
 v_\bullet(\widehat x_\ell)\mapsto v_\bullet(\widehat x_k),\quad
 u_\bullet(\widehat x_\ell)\mapsto u_\bullet(\widehat x_k),\quad
 \omega_E(\widehat x_\ell)\mapsto\omega_E(\widehat x_k).
\end{equation}
De \eqref{eq:pdfv2-cor-cell-frontera-aristas}--
\eqref{eq:pdfv2-cor-cell-frontera-celda} se obtiene
\begin{equation}
\label{eq:pdfv2-cor-collar-naturalidad}
\begin{aligned}
 \operatorname{Cell}(\tau)\partial
 &=\partial\operatorname{Cell}(\tau),\\
 \tau_VP_L&=P_{\tau_{\mathscr I}L}\tau_V,\qquad
 \tau_QB=B\tau_a,\qquad W_c\circ\tau_Q=W_c,\\
 \operatorname{Cell}(\tau)\vartheta_L
 &=\vartheta_{\tau_{\mathscr I}L}\operatorname{Cell}(\tau),\\
 (\tau_V\otimes\operatorname{Cell}_0(\tau))
 \eta^{\rm cell}_{\widehat x_\ell}
 &=\eta^{\rm cell}_{\widehat x_k}\tau_V.
\end{aligned}
\end{equation}
Identidad y composición se verifican en cada generador. Así
\(\operatorname{Cell}_{\rm TPK}\) es la completación celular funtorial de
la misma fibra APP--TRIT--TPK, no un portador añadido desde una publicación
posterior.

\subsection{Medida de incidencia, curvatura de acarreo y naturalidad adjunta}
En el desarrollo siguiente se identifican, exclusivamente como abreviaturas
de la completación común,
\[
 I_{\rm gau}:=\mathscr I_{\widehat x},\qquad
 \mathscr B_{\rm gau}:=\mathscr B_{\widehat x},\qquad
 Q_{\rm gau}:=Q_{\rm inc}(\widehat x).
\]
Con el promedio uniforme de conteo sobre \(Q_{\rm gau}\), la función
\(W_c\) de \eqref{eq:pdfv2-cor-wc} satisface
\begin{equation}
\label{eq:pdfv2-cor-wc-momentos}
 \mathbb E_QW_c=0,\qquad
 \mathbb E_QW_c^2=\frac29,\qquad
 \mathbb E_QW_c^3=\frac{2}{27}.
\end{equation}
Estas identidades son propiedades de la misma celda y del mismo acarreo, no
datos de una teoría exterior.
\begingroup
  \let\section\subsection
  % Extracción quirúrgica del fuente teoremática V2 05c: líneas 99--335.
% El resto permanece en este archivo como procedencia, pero no se publica.
\iffalse
\chapter[Macroestructura gauge interna]{Macroestructura gauge interna generada por el continuo discreto}
\label{chap:gauint-macroestructura}

\begin{thesisbox}
La macroestructura de este capítulo se construye únicamente a partir de una
restricción superviviente de \(\mathcal C_{\rm cont}^{\rm disc}\). Los números
\(4\), \(6\) y \(8\), la pantalla, las cuatro involuciones, la medida
proyectiva, el operador de transferencia, el modo de incidencia y el terminal
celular se derivan dentro de esa restricción. Ninguna evaluación posterior
forma parte del dominio ni de las flechas que siguen.
\end{thesisbox}

\section{La semilla proyectiva y la derivación interna de \texorpdfstring{\(4/6/8\)}{4/6/8}}

Sea \(V_{\rm TPK}=\mathbf F_3^2\), con la base ordenada heredada de las dos
hojas del estado TPK. Definimos el conjunto de direcciones
\begin{equation}\label{eq:gauint-cuatro-direcciones}
 I_{\rm gau}:=\mathbf P^1(\mathbf F_3)
 =\{L_\infty,L_0,L_1,L_{-1}\},
\end{equation}
donde
\[
 L_\infty=[1:0],\qquad L_0=[0:1],\qquad
 L_1=[1:1],\qquad L_{-1}=[1:-1].
\]
Por tanto,
\begin{equation}\label{eq:gauint-cardinal-cuatro}
 |I_{\rm gau}|=\frac{3^2-1}{3-1}=4.
\end{equation}
No se introduce un conjunto de cuatro índices: aparece como el cociente de
los ocho vectores no nulos por la acción de \(\mathbf F_3^\times\).

Las incidencias no orientadas son
\begin{equation}\label{eq:gauint-seis-incidencias}
 E_{\rm gau}:=\binom{I_{\rm gau}}2,
 \qquad |E_{\rm gau}|=\binom42=6,
 \qquad Q_{\rm gau}:=\mathbf F_3^{E_{\rm gau}}.
\end{equation}
La pantalla orientada duplica cada dirección, no cada incidencia:
\begin{equation}\label{eq:gauint-ocho-caras}
 \mathscr B_{\rm gau}:=I_{\rm gau}\times\{+,-\},
 \qquad |\mathscr B_{\rm gau}|=2|I_{\rm gau}|=8.
\end{equation}
Así, \(4\), \(6\) y \(8\) proceden respectivamente de direcciones, pares de
direcciones y caras orientadas de una misma semilla.

Para cada \(L\in I_{\rm gau}\), sea \(\vartheta_L\) la involución de la
pantalla que intercambia \((L,+)\) y \((L,-)\) y deja fijas las otras seis
caras. Entonces
\begin{equation}\label{eq:gauint-cuatro-involuciones}
 \vartheta_L^2=I,
 \qquad
 \vartheta_L\vartheta_{L'}=\vartheta_{L'}\vartheta_L
 \quad(L\ne L').
\end{equation}
Las cuatro involuciones se han obtenido antes de definir medida u operador.

\section{Incidencia, modo interno y testigo no constante}

La aplicación de incidencia es
\begin{equation}\label{eq:gauint-incidencia-b}
 B:\mathbf F_3^{I_{\rm gau}}\longrightarrow Q_{\rm gau},
 \qquad
 (Ba)_{\{L,L'\}}:=a_L+a_{L'}.
\end{equation}
Si
\[
 s(q):=\sum_{e\in E_{\rm gau}}q_e\in\mathbf F_3,
\]
cada dirección aparece exactamente tres veces en la suma de las seis
incidencias, y por ello
\begin{equation}\label{eq:gauint-incidencia-invariante}
 s(Ba)=3\sum_{L\in I_{\rm gau}}a_L=0,
 \qquad
 s(q+Ba)=s(q).
\end{equation}

Sobre \(Q_{\rm gau}\) se usa exclusivamente el promedio de conteo
normalizado. El modo de incidencia es la función racional
\begin{equation}\label{eq:gauint-wc-definicion}
 W_c(q):=\mathbf 1_{\{s(q)=0\}}-\frac13.
\end{equation}
Como \(s:Q_{\rm gau}\to\mathbf F_3\) es sobreyectiva, cada fibra tiene
\(3^5\) elementos. De aquí se obtienen, por conteo exacto,
\begin{equation}\label{eq:gauint-wc-momentos}
 \mathbb E_{Q}W_c=0,
 \qquad
 \mathbb E_{Q}W_c^2=\frac29,
 \qquad
 \mathbb E_{Q}W_c^3=\frac{2}{27}.
\end{equation}
Además, \eqref{eq:gauint-incidencia-invariante} da
\begin{equation}\label{eq:gauint-wc-invariancia}
 W_c(q+Ba)=W_c(q).
\end{equation}
Por tanto \(W_c\) es no nulo, centrado e invariante por la incidencia
generada internamente.

\fi
\section{Incidencia, curvatura de acarreo y dominio común de naturalidad}

Para \(x\in\mathcal C_{\rm cont}^{\rm disc}\), denotemos por
\(Y_m(x)\) el conjunto finito de prefijos gauge de profundidad \(m\) que
conservan simultáneamente
\begin{equation*}
\begin{gathered}
 \text{fibra, relaciones de extensión, números, orientación, carry, memoria,}\\
 \text{frontera, ruta, multisección, supervivencia, terminal y lector}.
\end{gathered}
\end{equation*}
Existe una aplicación de collar
\begin{equation}\label{eq:gauint-collar}
 q_m:Y_m(x)\longrightarrow Q_{\rm gau}
\end{equation}
que registra las seis incidencias sin borrar la historia que las produce.
Para \(\ell\ge m\), el truncamiento se escribe
\[
 \tau_{\ell,m}:Y_\ell(x)\longrightarrow Y_m(x).
\]

Para \(y,z\in Y_m(x)\), sea
\begin{equation}\label{eq:gauint-todas-extensiones}
 \operatorname{Ext}^{\rm adm}_{m}(y,z)
\end{equation}
el conjunto de \emph{todas} las extensiones TPK admisibles de \(y\) a \(z\).
Una extensión pertenece a este conjunto sólo si transporta la fibra entera,
las relaciones \(\operatorname{Ext}_\Gamma\), el carry nativo y operatorial,
la orientación, la memoria acumulada, la frontera, la ruta y la condición de
supervivencia. Definimos
\begin{equation}\label{eq:gauint-conteo-kernel}
 a_m(y,z):=\#\operatorname{Ext}^{\rm adm}_m(y,z),
 \qquad
 d_m(y):=\sum_{z\in Y_m(x)}a_m(y,z),
\end{equation}
y, cuando \(d_m(y)>0\),
\begin{equation}\label{eq:gauint-kernel-normalizado}
 k_m(y,z):=\frac{a_m(y,z)}{d_m(y)}.
\end{equation}
Así, \(k_m\) es un conteo normalizado de extensiones completas, no una
función introducida después de proyectar el estado.

La restricción \(\mathcal C_{\rm cont}^{\rm gau}\) está formada por las
historias supervivientes \(x\) para las que se cumplen las condiciones
siguientes a toda profundidad.
\begin{enumerate}
 \item Los conjuntos \(Y_m(x)\) son no vacíos, \(d_m(y)>0\), y los
 truncamientos \(\tau_{\ell,m}\) son sobreyectivos.
 \item Existe una multisección no vacía de pesos racionales estrictamente
 positivos \(\mu_m:Y_m(x)\to\mathbf Q_{>0}\) tal que
 \begin{equation}\label{eq:gauint-medida-proyectiva}
  \sum_{y\in Y_m}\mu_m(y)=1,
  \qquad
  \mu_m(y)=\sum_{\tau_{\ell,m}\widetilde y=y}\mu_\ell(\widetilde y).
 \end{equation}
 \item Cada \(\mu_m\) es estacionaria:
 \begin{equation}\label{eq:gauint-estacionariedad}
  \sum_{y\in Y_m}\mu_m(y)k_m(y,z)=\mu_m(z).
 \end{equation}
 \item La marginal de \((q_m)_*\mu_m\) es el conteo uniforme de
 \(Q_{\rm gau}\). Esta condición hace aplicables literalmente los cálculos
 de \eqref{eq:pdfv2-cor-wc-momentos} al lift \(W_c\circ q_m\).
 \item Se cumple la lumpabilidad directa
 \begin{equation}\label{eq:gauint-lump-directa}
  \sum_{\tau_{\ell,m}\widetilde z=z}
   k_\ell(\widetilde y,\widetilde z)
  =k_m(\tau_{\ell,m}\widetilde y,z)
 \end{equation}
 para todo \(\widetilde y\in Y_\ell\) y \(z\in Y_m\).
 \item Para el kernel adjunto
 \begin{equation}\label{eq:gauint-kernel-adjunto}
  k_m^*(z,y):=\frac{\mu_m(y)k_m(y,z)}{\mu_m(z)},
 \end{equation}
 se cumple la lumpabilidad adjunta
 \begin{equation}\label{eq:gauint-lump-adjunta}
  \sum_{\tau_{\ell,m}\widetilde y=y}
   k_\ell^*(\widetilde z,\widetilde y)
  =k_m^*(\tau_{\ell,m}\widetilde z,y).
 \end{equation}
 \item Las categorías de prefijos poseen la raíz TPK como objeto inicial,
 de modo que su nervio tiene la contracción terminal explícita de
 \eqref{eq:gauint-sdr} más abajo.
\end{enumerate}
Estas condiciones definen el dominio tipado; fuera de él no se publica una
medida o un operador natural por simple declaración.

\section{Medida proyectiva, kernel y naturalidad directa y adjunta}

Fijada una rama de la multisección \(\{\mu_m\}_m\), sea
\begin{equation}\label{eq:gauint-hm}
 H_m:=\operatorname{Fun}(Y_m,\mathbf Q),
 \qquad
 \langle f,g\rangle_m:=\sum_{y\in Y_m}\mu_m(y)f(y)g(y).
\end{equation}
El operador generado por todas las extensiones es
\begin{equation}\label{eq:gauint-km}
 (K_mf)(y):=\sum_{z\in Y_m}k_m(y,z)f(z).
\end{equation}
Las ecuaciones de Markov y estacionariedad dan
\begin{equation}\label{eq:gauint-markov-estacionario}
 K_m\mathbf1=\mathbf1,
 \qquad
 K_m^*\mathbf1=\mathbf1,
 \qquad
 \|K_mf\|_m\le\|f\|_m.
\end{equation}

El lift y la esperanza de truncamiento son
\begin{equation}\label{eq:gauint-lift-esperanza}
 J_{\ell,m}f:=f\circ\tau_{\ell,m},
 \qquad
 E_{m,\ell}:=J_{\ell,m}^*.
\end{equation}
La proyectividad de \(\mu\) implica
\begin{equation}\label{eq:gauint-ej-identidad}
 E_{m,\ell}J_{\ell,m}=I_{H_m}.
\end{equation}
Las dos condiciones de lumpabilidad descargan, respectivamente,
\begin{equation}\label{eq:gauint-naturalidad-directa-adjunta}
 K_\ell J_{\ell,m}=J_{\ell,m}K_m,
 \qquad
 K_\ell^*J_{\ell,m}=J_{\ell,m}K_m^*.
\end{equation}
La segunda igualdad no se infiere de la primera sin la condición adjunta:
ambas quedan registradas en el dominio.

Se define el operador interno positivo
\begin{equation}\label{eq:gauint-transferencia-t}
 T_m:=K_m^*K_m.
\end{equation}
Entonces
\begin{equation}\label{eq:gauint-t-positivo-natural}
 \langle f,T_mf\rangle_m=\|K_mf\|_m^2\ge0,
 \qquad
 T_\ell J_{\ell,m}=J_{\ell,m}T_m.
\end{equation}

\section{Ley común de \(8\) caras y \(4\) cuadrados de Fubini}

Sobre \(Y_m^{\mathscr B_{\rm gau}}\) se define la masa racional
\begin{equation}\label{eq:gauint-ley-ocho-caras}
 \Omega_m^{(8)}(\mathbf y)
 :=\sum_{z\in Y_m}\mu_m(z)
 \prod_{(L,\varepsilon)\in\mathscr B_{\rm gau}}
 k_m(z,y_{L,\varepsilon}).
\end{equation}
Es una probabilidad porque cada factor es Markov. Todas las caras proceden
del mismo estado latente \(z\) y del mismo kernel; no son ocho medidas
independientes añadidas.

Para funciones \(f_{L,+},f_{L,-}\in H_m\), Fubini finito da
\begin{equation}\label{eq:gauint-fubini-ocho}
\begin{aligned}
 &\sum_{\mathbf y}\Omega_m^{(8)}(\mathbf y)
  \prod_{L\in I_{\rm gau}}
  f_{L,+}(y_{L,+})f_{L,-}(y_{L,-})\\
 &\qquad=
 \sum_{z\in Y_m}\mu_m(z)
 \prod_{L\in I_{\rm gau}}
 (K_mf_{L,+})(z)(K_mf_{L,-})(z).
\end{aligned}
\end{equation}
Al tomar \(f_{L,+}=f_{L,-}=f_L\), aparecen simultáneamente los cuatro
cuadrados
\begin{equation}\label{eq:gauint-cuatro-cuadrados}
 \sum_{z\in Y_m}\mu_m(z)
 \prod_{L\in I_{\rm gau}}(K_mf_L(z))^2.
\end{equation}
La marginal de cualquier par opuesto satisface
\begin{equation}\label{eq:gauint-marginal-t}
 \sum_{y_+,y_-}\Omega_{m,L}^{(2)}(y_+,y_-)
 f(y_+)g(y_-)
 =\langle K_mf,K_mg\rangle_m
 =\langle f,T_mg\rangle_m.
\end{equation}
La conmutatividad del producto en \eqref{eq:gauint-ley-ocho-caras} prueba
\begin{equation}\label{eq:gauint-involuciones-ley}
 (\vartheta_L)_*\Omega_m^{(8)}=\Omega_m^{(8)}
 \qquad(L\in I_{\rm gau}).
\end{equation}

\section{Operador interno exacto y testigo propio}

La proyección sobre el sector constante y su complemento son
\begin{equation}\label{eq:gauint-pq}
 P_{\Omega,m}f:=\langle\mathbf1,f\rangle_m\mathbf1,
 \qquad
 Q_m:=I-P_{\Omega,m}.
\end{equation}
El operador celular interno se define por
\begin{equation}\label{eq:gauint-d-hmt}
 D_m^{\rm HMT}:=3Q_m.
\end{equation}
El coeficiente \(3\) es el cardinal del cuerpo de residuos usado en
\eqref{eq:pdfv2-cor-proyectores-cuatro-direcciones}; no representa una
escala añadida.
Para el lift
\[
 W_{c,m}:=W_c\circ q_m
\]
se tiene, por \eqref{eq:pdfv2-cor-wc-momentos},
\begin{equation}\label{eq:gauint-w-eigen}
 P_{\Omega,m}W_{c,m}=0,
 \qquad
 D_m^{\rm HMT}W_{c,m}=3W_{c,m},
 \qquad
 \|W_{c,m}\|_m^2=\frac29.
\end{equation}
Así, el sector complementario posee un testigo no nulo construido por la
incidencia de las seis aristas.

\section{Terminal celular y retracto fuerte}

Sea \(C_*(N\mathsf Y_m;A_n)\) el complejo normalizado del nervio de la
categoría de prefijos, y sea \(y_*\) su raíz inicial. La inclusión y la
aumentación son
\[
 \iota_m:A_n[0]\longrightarrow C_*(N\mathsf Y_m;A_n),
 \quad \iota_m(1)=[y_*],
 \qquad
 \epsilon_m:C_*(N\mathsf Y_m;A_n)\longrightarrow A_n[0].
\]
El objeto inicial proporciona la degeneración extra
\begin{equation}\label{eq:gauint-homotopia-celular}
 H_m[y_0\to\cdots\to y_q]
 :=[y_*\to y_0\to\cdots\to y_q],
\end{equation}
con la convención normalizada cuando aparece una degeneración. Se verifica
\begin{equation}\label{eq:gauint-sdr}
 \epsilon_m\iota_m=I,
 \qquad
 \partial H_m+H_m\partial=I-\iota_m\epsilon_m,
 \qquad
 H_m^2=H_m\iota_m=\epsilon_mH_m=0.
\end{equation}
Este terminal celular es distinto de la proyección constante
\(P_{\Omega,m}\): el primero contrae el nervio y la segunda separa los
sectores del espacio de funciones. Ambos se conservan.

\iffalse
\section{Las trece coordenadas de la restricción gauge}

Para cada \((m,n)\), se define
\begin{equation}\label{eq:gauint-d-trece}
 D_{{\rm gau},m,n}(x):=
 (\mathcal F,\operatorname{Ext}_\Gamma,\operatorname{Rel},
 \mathcal N,\operatorname{or},\operatorname{Carr},\operatorname{Mem},
 \partial,\gamma,\operatorname{Msect},\operatorname{Surv},
 \operatorname{Term},\mathcal L)_{{\rm gau},m,n}(x),
\end{equation}
con el contenido siguiente.
\begin{enumerate}
 \item \(\mathcal F=(V_{\rm TPK},I_{\rm gau},E_{\rm gau},
 Q_{\rm gau},\mathscr B_{\rm gau},Y_m,H_m)\).
 \item \(\operatorname{Ext}_\Gamma\) contiene todos los conjuntos
 \(\operatorname{Ext}^{\rm adm}_m(y,z)\), truncamientos, lifts,
 esperanzas y sus composiciones.
 \item \(\operatorname{Rel}\) contiene \(B\), \(sB=0\), las cuatro
 involuciones, Markov, estacionariedad y las dos lumpabilidades.
 \item \(\mathcal N=(3,4,6,8,m,n,|Y_m|,a_m,d_m)\); esos números no
 sustituyen ninguna de las restantes coordenadas.
 \item \(\operatorname{or}\) conserva las caras \(+/-\), las cuatro
 inversiones \(\vartheta_L\) y la reversión de cada ruta.
 \item \(\operatorname{Carr}\) contiene el carry completo de cada
 extensión contada en \(a_m(y,z)\), incluido el carry ya propagado.
 \item \(\operatorname{Mem}\) conserva el prefijo completo, el estado
 común \(z\) de \(\Omega_m^{(8)}\) y la memoria de sus ocho extensiones.
 \item \(\partial=(q_m,\partial_{N\mathsf Y_m},Q_m,H_m)\) conserva
 frontera de collar, frontera celular y complemento.
 \item \(\gamma\) es el símplice orientado del nervio que registra la ruta
 y no sólo sus extremos.
 \item \(\operatorname{Msect}\) es la familia de todas las medidas
 estacionarias proyectivas compatibles y de todas las extensiones que
 intervienen en los conteos; no se selecciona una rama retrospectiva.
 \item \(\operatorname{Surv}\) registra no vacuidad, profundidad
 arbitraria, proyectividad, las dos lumpabilidades y la persistencia de la
 marginal uniforme de collar.
 \item \(\operatorname{Term}=(\iota_m,\epsilon_m,H_m;
 P_{\Omega,m},Q_m)\) conserva separados los dos terminales.
 \item \(\mathcal L=(\Omega_m^{(8)},K_m,T_m,D_m^{\rm HMT},W_{c,m})\)
 es el lector interno posterior a la construcción de la historia.
\end{enumerate}

Para \(\ell\ge m\) y \(n'\ge n\), todas las coordenadas satisfacen
\begin{equation}\label{eq:gauint-naturalidad-trece}
 u^{\rm gau}_{(\ell,n'),(m,n)}D_{{\rm gau},\ell,n'}(x')
 =D_{{\rm gau},m,n}(\tau_{\ell,m}x').
\end{equation}

\section{Ensamblador, publicador y cadena no ablativa}

Sea
\begin{equation}\label{eq:gauint-graph-d}
 \mathcal G_{\rm gau}:=\operatorname{Graph}(D_{\rm gau}),
 \qquad
 \operatorname{Res}_{\rm gau}(x):=(x,D_{\rm gau}x).
\end{equation}
La multisección dependiente \(\chi_{\rm gau}(x)\) contiene todas las ramas
compatibles de collar, medida y extensión. Definimos
\begin{equation}\label{eq:gauint-x-dependiente}
 X_{\chi,{\rm gau}}
 :=\{(\chi_{\rm gau}(x),x,D_{\rm gau}x):
 x\in\mathcal C_{\rm cont}^{\rm gau}\},
\end{equation}
\begin{equation}\label{eq:gauint-prx}
 \operatorname{pr}_{X,{\rm gau}}(x,D_{\rm gau}x)
 :=(\chi_{\rm gau}(x),x,D_{\rm gau}x).
\end{equation}

El ensamblador crudo
\begin{equation}\label{eq:gauint-ensamblador-crudo}
 a_{\rm gau}:X_{\chi,{\rm gau}}\longrightarrow
 \mathscr M_{\rm gau}^{\rm amb}
\end{equation}
produce la familia completa
\[
 \bigl(I,E,\mathscr B,B,W_c;
 Y_m,\mu_m,k_m,J_{\ell,m},E_{m,\ell};
 K_m,T_m,\Omega_m^{(8)},P_{\Omega,m},Q_m,D_m^{\rm HMT};
 \iota_m,\epsilon_m,H_m\bigr)_{m,n}.
\]
El macroobjeto conserva su entrada mediante
\begin{equation}\label{eq:gauint-graph-a}
 \mathfrak M_{\rm gau}:=\operatorname{Graph}(a_{\rm gau}),
 \qquad
 \mathcal A_{\rm gau}(z):=(z,a_{\rm gau}z).
\end{equation}

El publicador crudo
\begin{equation}\label{eq:gauint-publicador-crudo}
 p_{\rm gau}:\mathfrak M_{\rm gau}\longrightarrow
 \mathscr Y_{\rm gau}
\end{equation}
publica conjuntamente las identidades
\[
 \begin{gathered}
 |I|=4,quad |E|=6,quad |\mathscr B|=8,quad sB=0,
 \quad W_c(q+Ba)=W_c(q),\\
 K_\ell J=JK_m,quad K_\ell^*J=JK_m^*,
 \quad T_m=K_m^*K_m,\\
 (\vartheta_L)_*\Omega_m^{(8)}=\Omega_m^{(8)},
 \quad \text{las cuatro factorizaciones cuadráticas de Fubini},\\
 D_m^{\rm HMT}=3Q_m,quad D_m^{\rm HMT}W_{c,m}=3W_{c,m},
 \quad \partial H_m+H_m\partial=I-\iota_m\epsilon_m.
 \end{gathered}
\]
Finalmente,
\begin{equation}\label{eq:gauint-graph-p}
 \mathcal C_{\rm gau}^{\rm HMT}:=\operatorname{Graph}(p_{\rm gau}),
 \qquad
 \mathcal P_{\rm gau}(M):=(M,p_{\rm gau}M).
\end{equation}
La cadena completa es
\begin{equation}\label{eq:gauint-cadena-completa}
 \mathcal C_{\rm cont}^{\rm disc}\supset
 \mathcal C_{\rm cont}^{\rm gau}
 \xrightarrow{\operatorname{Res}_{\rm gau}}\mathcal G_{\rm gau}
 \xrightarrow{\operatorname{pr}_{X,{\rm gau}}}X_{\chi,{\rm gau}}
 \xrightarrow{\mathcal A_{\rm gau}}\mathfrak M_{\rm gau}
 \xrightarrow{\mathcal P_{\rm gau}}\mathcal C_{\rm gau}^{\rm HMT}.
\end{equation}
Cada flecha tiene una inversa sobre su imagen dada por una proyección de
coordenadas. Ninguna salida reemplaza su historia generadora.

\section{Teorema de realización gauge interna}

\begin{theorem}[Realización interna completa y natural]
\label{thm:gauint-realizacion}
Para todo \(x\in\mathcal C_{\rm cont}^{\rm gau}\), la cadena
\eqref{eq:gauint-cadena-completa} construye, a partir de la estructura
discreta del continuo, una macroestructura natural con cuatro direcciones,
seis incidencias, ocho caras, cuatro involuciones, una medida proyectiva, un
kernel de conteo completo, su transferencia positiva, una ley común de ocho
caras, un testigo de incidencia no constante y un terminal celular. Las
trece coordenadas de \eqref{eq:gauint-d-trece} permanecen recuperables en la
salida.
\end{theorem}

\begin{proof}
Las ecuaciones \eqref{eq:gauint-cuatro-direcciones}--
\eqref{eq:gauint-ocho-caras} derivan \(4/6/8\). La suma de incidencias
\eqref{eq:gauint-incidencia-invariante} prueba la invariancia de \(W_c\), y
el conteo uniforme prueba \eqref{eq:gauint-wc-momentos}. Las extensiones
completas producen \(k_m\) por \eqref{eq:gauint-kernel-normalizado}; las
condiciones explícitas del dominio prueban Markov, estacionariedad y las dos
igualdades de naturalidad. Por tanto \(T_m=K_m^*K_m\) es positivo y natural.
Fubini aplicado una sola vez a \eqref{eq:gauint-ley-ocho-caras} da las cuatro
formas cuadráticas sin cambiar de medida. La marginal uniforme centra
\(W_{c,m}\), de donde \eqref{eq:gauint-w-eigen}. El objeto inicial del nervio
da la degeneración extra y la identidad SDR \eqref{eq:gauint-sdr}. Finalmente,
las tres construcciones por graph retienen sus entradas como primeras
coordenadas, luego la composición no olvida ninguna de las trece entradas.
\end{proof}

\section{Procedencia y falsadores}

\paragraph{Procedencia.}
La semilla TPK, la pantalla orientada, las extensiones supervivientes, la
medida común y la contracción terminal tienen estatuto
\texttt{ARQUITECTURA\_AUTORAL\_PREEXISTENTE}. La transferencia por conteo,
el modo de incidencia y el sector exacto \(3Q\) se presentan como
\texttt{RESULTADO\_RECUPERADO}. La derivación explícita
\(\mathbf P^1(\mathbf F_3)\to4/6/8\), las dos condiciones de lumpabilidad y
el empaquetado no ablativo por graphs tienen estatuto
\texttt{FORMALIZACION\_NUEVA}; formalizan la arquitectura anterior sin
reclamar un origen conceptual nuevo.

\begin{ablation}[Falsadores internos de la rama gauge]
\label{abl:gauint-falsadores}
La construcción queda falsada en el nivel indicado si ocurre cualquiera de
los hechos siguientes.
\begin{enumerate}
 \item Se reemplaza \(I_{\rm gau}\), \(E_{\rm gau}\) o
 \(\mathscr B_{\rm gau}\) por una lista sin los mapas que los derivan.
 \item Falla \(sB=0\) o \(W_c\) deja de ser invariante y no constante.
 \item El conteo \(a_m(y,z)\) omite carry, orientación, memoria, frontera o
 ruta: en tal caso se ha contado otro objeto.
 \item Falla proyectividad, estacionariedad o alguna de las dos
 lumpabilidades; entonces no se puede publicar la naturalidad correspondiente.
 \item Las ocho caras no proceden de una única \(\mu_m\) y un único
 \(k_m\), o alguna \(\vartheta_L\) no preserva \(\Omega_m^{(8)}\).
 \item \(P_{\Omega,m}W_{c,m}\ne0\) o
 \(D_m^{\rm HMT}W_{c,m}\ne3W_{c,m}\).
 \item Falla la identidad SDR \(\partial H+H\partial=I-\iota\epsilon\).
 \item Se omite una de las trece coordenadas: \emph{objeto sustituido;
 conclusión no transferible}.
\end{enumerate}
Estos falsadores controlan pertenencia y naturalidad; no degradan una
historia que satisface todas las identidades.
\end{ablation}
\fi
 \endgroup

La misma regla define su acción sobre una emisión elemental
\(\alpha:x\to y\). Como \(U_\alpha\) transporta la pareja \emph{ordenada}
de hojas, su mapa de etiquetas es
\(\overline U_\alpha e_{\Sigma,x}=e_{\Sigma,y}\) y
\(\overline U_\alpha e_{\Pi,x}=e_{\Pi,y}\). Su proyectivización
\(\alpha_{\mathscr I}[v]:=[\overline U_\alpha v]\) actúa en
\(\mathbf P^1(V_x)\) y, por tanto,
\begin{equation}
\label{eq:pdfv2-cor-cell-accion-emision}
 \begin{gathered}
 \operatorname{Cell}_{\rm TPK}(\alpha)v_\star(x)=v_\star(y),\qquad
 \operatorname{Cell}_{\rm TPK}(\alpha)v_L(x)
   =v_{\alpha_{\mathscr I}L}(y),\\
 \operatorname{Cell}_{\rm TPK}(\alpha)v_{\{L,L'\}}(x)
   =v_{\{\alpha_{\mathscr I}L,\alpha_{\mathscr I}L'\}}(y),\qquad
 \operatorname{Cell}_{\rm TPK}(\alpha)v_{L,\epsilon}(x)
   =v_{\alpha_{\mathscr I}L,\epsilon}(y),\\
 \operatorname{Cell}_{\rm TPK}(\alpha)u_L(x)
   =u_{\alpha_{\mathscr I}L}(y),\qquad
 \operatorname{Cell}_{\rm TPK}(\alpha)u_{L,\epsilon}(x)
   =u_{\alpha_{\mathscr I}L,\epsilon}(y),\\
 \operatorname{Cell}_{\rm TPK}(\alpha)u_{L'\mid L}(x)
   =u_{\alpha_{\mathscr I}L'\mid\alpha_{\mathscr I}L}(y),\qquad
 \operatorname{Cell}_{\rm TPK}(\alpha)\omega_{\{L,L'\}}(x)
   =\omega_{\{\alpha_{\mathscr I}L,
                 \alpha_{\mathscr I}L'\}}(y).
 \end{gathered}
\end{equation}
En la fibra de hojas y en su pantalla se tiene además
\begin{equation}
\label{eq:pdfv2-cor-cell-entrelaza-proyectores-pantalla}
 \overline U_\alpha P_L=P_{\alpha_{\mathscr I}L}\overline U_\alpha,
 \qquad
 \operatorname{Cell}_{\rm TPK}(\alpha)\vartheta_L
 =\vartheta_{\alpha_{\mathscr I}L}
  \operatorname{Cell}_{\rm TPK}(\alpha).
\end{equation}
Las fórmulas de frontera dan
\begin{equation}
\label{eq:pdfv2-cor-cell-accion-emision-natural}
 \partial\operatorname{Cell}_{\rm TPK}(\alpha)
 =\operatorname{Cell}_{\rm TPK}(\alpha)\partial,\qquad
 \operatorname{Cell}_{\rm TPK}(\beta\alpha)
 =\operatorname{Cell}_{\rm TPK}(\beta)
  \operatorname{Cell}_{\rm TPK}(\alpha).
\end{equation}
Así el mapa de complejos usado en el registro terminal es la acción libre
inducida por la misma emisión, no un símbolo añadido.

\subsection{Módulo de rutas, totalización y cono del retorno nonádico}
Sea \(J_d(x)\) la categoría completa de prefijos de profundidad \(d\) del
árbol común, \(\operatorname{Term}_d(x)\) su conjunto terminal y
\(\Gamma_{d,N}(\nu):=C_{d,N}(\nu)\) el complejo celular de
\eqref{eq:pdfv2-cor-cono-celda}. El módulo, la totalización y el cono que
siguen conservan simultáneamente todas las continuaciones; no seleccionan
una historia ni separan una construcción cohomológica del continuo común.
\begingroup
  \let\section\subsection
  % Extracción quirúrgica del fuente teoremática V2 05d: líneas 286--512.
% El resto permanece en este archivo como procedencia, pero no se publica.
\iffalse
% Macroestructura de coherencia estrictamente interna del continuo discreto.
% No usa geometrías, clases ni conclusiones clásicas.

\chapter{Macroestructura interna de celdas, coherencia y completitud}
\label{chap:pdfv2-coh-interna}

\begin{statusbox}
La única entrada de este capítulo es la estructura discreta del continuo. La
rama \(\mathrm{coh}\) se construye sobre los anillos
\(A_N=\mathbb Z/3^N\mathbb Z\), las celdas APP, sus historias y sus
terminales. Su cadena no ablativa es
\[
 \mathcal C_{\rm cont}^{\rm disc}
 \supset\mathcal C_{\rm cont}^{\rm coh}
 \xrightarrow{\operatorname{Res}_{\rm coh}}
 \mathcal G_{\rm coh}
 \xrightarrow{\operatorname{pr}_{X,{\rm coh}}}
 X_{\chi,{\rm coh}}
 \xrightarrow{\mathcal A_{\rm coh}}
 \mathfrak M_{\rm coh}
 \xrightarrow{\mathcal P_{\rm coh}}
 \mathcal C_{\rm coh}^{\rm HMT}.
\]
No aparece como entrada ni como salida ningún objeto clásico. Cada imagen
conserva literalmente su historia generadora.
\end{statusbox}

\section{Dominio tipado y separación de índices}
\label{sec:pdfv2-coh-domain}

Para no identificar profundidad con precisión, se usa el conjunto dirigido
\begin{equation}\label{eq:pdfv2-coh-lambda}
 \Lambda:=\mathbb N_{\rm prof}\times\mathbb N_{\rm prec},
 \qquad
 \lambda=(d,N),
 \qquad
 A_N:=\mathbb Z/3^N\mathbb Z.
\end{equation}
Si \(\lambda\leq\lambda'=(d',N')\), la transición
\begin{equation}\label{eq:pdfv2-coh-rho-lambda}
 \rho_{\lambda',\lambda}:
 \mathcal C_{{\rm cont},\lambda'}^{\rm disc}
 \longrightarrow
 \mathcal C_{{\rm cont},\lambda}^{\rm disc}
\end{equation}
trunca la historia después de la profundidad \(d\) y reduce coeficientes por
\(A_{N'}\twoheadrightarrow A_N\), conservando hoja, orientación, residuo,
cociente, carry, memoria, frontera, ruta, supervivencia y terminal.

La subestructura \(\mathcal C_{\rm cont}^{\rm coh}\) consta de las historias
\(x\in\mathcal C_{\rm cont}^{\rm disc}\) que cumplen, en todo nivel:
\begin{enumerate}
 \item el prefijo de \(x\) porta las celdas orientadas y basadas definidas en
       la sección siguiente;
 \item cada refinamiento induce un mapa del complejo de la celda;
 \item truncación, reducción de coeficientes, orientación y extensión
       nonádica conmutan;
 \item las multisecciones de extensión superviviente son no vacías a toda
       profundidad;
 \item el retorno nonádico conserva la diferencia terminal entre la historia
       y su continuación.
\end{enumerate}
Estas condiciones definen el dominio de la rama. No se atribuyen a una
historia arbitraria sin verificarlas.

\section{Celda APP sobre \texorpdfstring{\(A_N\)}{AN}}
\label{sec:pdfv2-coh-cell}

Una celda orientada es
\[
 \nu=(h_\nu,I_\nu,J_\nu,\epsilon_\nu),
\]
donde \(h_\nu\) es una historia finita, \(I_\nu,J_\nu\) son conjuntos
ordenados no vacíos de puertos,
\[
 a_\nu:=|I_\nu|\geq1,
 \qquad
 b_\nu:=|J_\nu|\geq1,
\]
y \(\epsilon_\nu\in\{+1,-1\}\) es su orientación.

Se define
\begin{equation}\label{eq:pdfv2-coh-kappa}
 \begin{aligned}
 \kappa_{\nu,N}:
 A_N^{I_\nu}\oplus A_N^{J_\nu}
 &\longrightarrow A_N^{I_\nu\times J_\nu},\\
 \kappa_{\nu,N}(u,v)_{ij}&:=u_i-v_j.
 \end{aligned}
\end{equation}
Elegidos los puertos base \(i_0\in I_\nu\), \(j_0\in J_\nu\), sea
\begin{equation}\label{eq:pdfv2-coh-delta}
 \begin{aligned}
 \delta_{\nu,N}:A_N^{I_\nu\times J_\nu}
 &\longrightarrow
 A_N^{(I_\nu\setminus i_0)\times(J_\nu\setminus j_0)},\\
 \delta_{\nu,N}(w)_{ij}
 &:={}
 w_{ij}-w_{i j_0}-w_{i_0j}+w_{i_0j_0}.
 \end{aligned}
\end{equation}

\begin{proposition}[Sucesión exacta de la celda]
\label{prop:pdfv2-coh-cell-exact}
Para toda celda \(\nu\) y toda precisión \(N\), la sucesión
\begin{equation}\label{eq:pdfv2-coh-exact-sequence}
 0\longrightarrow A_N
 \xrightarrow{\ \Delta\ }
 A_N^{I_\nu}\oplus A_N^{J_\nu}
 \xrightarrow{\ \kappa_{\nu,N}\ }
 A_N^{I_\nu\times J_\nu}
 \xrightarrow{\ \delta_{\nu,N}\ }
 A_N^{(a_\nu-1)(b_\nu-1)}
 \longrightarrow0,
\end{equation}
donde
\[
 \Delta(t)=\bigl((t)_{i\in I_\nu},(t)_{j\in J_\nu}\bigr),
\]
es exacta y escindida sobre \(A_N\).
\end{proposition}

\begin{proof}
Si \(\kappa_{\nu,N}(u,v)=0\), entonces \(u_i=v_j\) para todo \(i,j\),
por lo que \((u,v)\) es diagonal. Así,
\(\ker\kappa_{\nu,N}=\operatorname{im}\Delta\). La sustitución directa da
\(\delta_{\nu,N}\kappa_{\nu,N}=0\).

Si \(\delta_{\nu,N}(w)=0\), póngase
\[
 u_i=w_{i j_0},
 \qquad
 v_j=w_{i_0j_0}-w_{i_0j}.
\]
Entonces
\[
 u_i-v_j=w_{i j_0}+w_{i_0j}-w_{i_0j_0}=w_{ij},
\]
de modo que \(\ker\delta_{\nu,N}=\operatorname{im}\kappa_{\nu,N}\). Para
probar la sobreyectividad de \(\delta_{\nu,N}\), se anulan la fila y la
columna base y se colocan arbitrariamente las coordenadas deseadas en las
posiciones restantes. Esas mismas fórmulas proporcionan escisiones.
\end{proof}

La celda compleja, en grados \(1\to0\), es el cono homológico de
\(\kappa_{\nu,N}\) con la convención que coloca su fuente en grado \(1\):
\begin{equation}\label{eq:pdfv2-coh-cell-complex}
 C_{\nu,N}:=\operatorname{Cone}(\kappa_{\nu,N})
 :=\left[
 A_N^{I_\nu}\oplus A_N^{J_\nu}
 \xrightarrow{\ \kappa_{\nu,N}\ }
 A_N^{I_\nu\times J_\nu}
 \right].
\end{equation}
La proposición anterior da
\begin{equation}\label{eq:pdfv2-coh-cell-homology}
 H_1(C_{\nu,N})\simeq A_N,
 \qquad
 H_0(C_{\nu,N})\simeq A_N^{(a_\nu-1)(b_\nu-1)}.
\end{equation}
El defecto de celda queda así construido, no añadido como un cardinal sin
mapa.

\section{Orientación y carry de la celda}
\label{sec:pdfv2-coh-orientation-carry}

Sean
\[
 \tau_1(u,v)=(v,u),
 \qquad
 P(w)_{ji}=w_{ij}.
\]
El intercambio orientado es el mapa de complejos
\begin{equation}\label{eq:pdfv2-coh-theta}
 \Theta_{\nu,N}:=(\tau_1,-P):
 C_{a_\nu,b_\nu,N}\longrightarrow C_{b_\nu,a_\nu,N},
\end{equation}
porque
\begin{equation}\label{eq:pdfv2-coh-orientation}
 \kappa_{b_\nu,a_\nu,N}\tau_1
 =-P\kappa_{a_\nu,b_\nu,N}.
\end{equation}
Además,
\(\Theta_{\nu^{\rm op},N}\Theta_{\nu,N}=I\). La ruta opuesta no se
identifica con la original: cambia el signo de grado cero y conserva el acto
de inversión.

Para \(r\in A_N\), sea
\(\langle r\rangle_N\in\{0,\ldots,3^N-1\}\) su representante canónico. El
borrow/carry local es
\begin{equation}\label{eq:pdfv2-coh-beta}
 \beta_N(r,s):=
 \frac{\langle r-s\rangle_N-\langle r\rangle_N+\langle s\rangle_N}
 {3^N}\in\{0,1\}.
\end{equation}
Así,
\begin{equation}\label{eq:pdfv2-coh-carry-cell}
 \langle u_i-v_j\rangle_N
 =\langle u_i\rangle_N-\langle v_j\rangle_N
  +3^N\beta_N(u_i,v_j).
\end{equation}
La matriz
\[
 \operatorname{Carr}_{\nu,N}(u,v)
 :=\bigl(\beta_N(u_i,v_j)\bigr)_{ij}
\]
se conserva junto a \(\kappa_{\nu,N}(u,v)\): residuo y carry son dos
coordenadas distintas.

Cuando los tamaños de puertos proceden de las dos hojas APP, la admisibilidad
de la celda incluye la identidad entera
\begin{equation}\label{eq:pdfv2-coh-app-defect}
 (a_\nu-1)(b_\nu-1)
 =\rho_{\Pi,\nu}-\rho_{\Sigma,\nu}+1
  +9(c_{\Pi,\nu}-c_{\Sigma,\nu}).
\end{equation}
La igualdad se exige únicamente a una celda con esa genealogía; no se afirma
para enteros desprovistos de hojas APP.

Si \(L_N(u)\) es la elevación entera canónica de la matriz de una flecha
\(u\), se define el carry de composición por
\begin{equation}\label{eq:pdfv2-coh-carry-composition}
 L_N(v)L_N(u)=L_N(vu)+3^N c_N(v,u).
\end{equation}
La asociatividad de tres flechas da el cociclo exacto
\begin{equation}\label{eq:pdfv2-coh-carry-cocycle}
 c_N(w,vu)+L_N(w)c_N(v,u)
 =c_N(wv,u)+c_N(w,v)L_N(u).
\end{equation}

\section{Categoría de historias y functor celular}
\label{sec:pdfv2-coh-category}

Para un prefijo \(x_{\leq d}\), sea \(J_d(x)\) la categoría finita cuyos
objetos son todas sus celdas \(\nu\). Sus flechas están generadas por:
\begin{enumerate}
 \item refinamientos \(u:\nu\to\nu'\) con aplicaciones sobreyectivas de
       puertos
       \[
        \alpha_u:I_{\nu'}\twoheadrightarrow I_\nu,
        \qquad
        \beta_u:J_{\nu'}\twoheadrightarrow J_\nu;
       \]
 \item los intercambios de orientación \(\Theta_\nu\);
 \item las extensiones nonádicas \(\Gamma_9\nu\);
 \item identidad y composición, sometidas a las relaciones de frontera,
       ruta, orientación y carry de la historia.
\end{enumerate}
La finitud se refiere al prefijo; el sistema de todos los prefijos no se
declara finito.

El functor celular es
\begin{equation}\label{eq:pdfv2-coh-gamma-functor}
 \Gamma_{d,N}:J_d(x)\longrightarrow
 \operatorname{Ch}^{[0,1]}(A_N),
 \qquad
 \Gamma_{d,N}(\nu)=C_{\nu,N}.
\end{equation}
Para una flecha de refinamiento \(u:\nu\to\nu'\), actúa mediante
\begin{align}
 \Gamma_{d,N}(u)_1(p,q)
 &:=(p\circ\alpha_u,q\circ\beta_u),
 \label{eq:pdfv2-coh-gamma-one}\\
 \Gamma_{d,N}(u)_0(w)
 &:=w\circ(\alpha_u\times\beta_u).
 \label{eq:pdfv2-coh-gamma-zero}
\end{align}
Se verifica coordenada a coordenada
\begin{equation}\label{eq:pdfv2-coh-gamma-chain-map}
 \kappa_{\nu',N}\Gamma_{d,N}(u)_1
 =\Gamma_{d,N}(u)_0\kappa_{\nu,N}.
\end{equation}
Sobre una inversión de orientación actúa por
\(\Theta_{\nu,N}\), y sobre composiciones actúa por composición, guardando
separadamente \(c_N(v,u)\).

La reducción \(r_{N+1,N}:A_{N+1}\twoheadrightarrow A_N\) satisface
\begin{equation}\label{eq:pdfv2-coh-kappa-reduction}
 r_{N+1,N}^{I\times J}\kappa_{\nu,N+1}
 =\kappa_{\nu,N}
  \bigl(r_{N+1,N}^{I}\oplus r_{N+1,N}^{J}\bigr).
\end{equation}
Por tanto, el diagrama de celdas es natural simultáneamente en historia y
precisión.

\fi
\section{Módulo de rutas, totalización celular y terminal nonádico}
\label{sec:pdfv2-coh-route-module}

Sea \(\operatorname{Term}_d(x)\) el conjunto de terminales de memoria del
prefijo. Para \(\nu\in J_d(x)\), defínase
\begin{equation}\label{eq:pdfv2-coh-w-module}
 W_{d,N}(\nu):=
 A_N\bigl[
 \operatorname{Path}_{J_d(x)}(\nu,\operatorname{Term}_d(x))
 \bigr].
\end{equation}
Es el módulo libre sobre todas las rutas terminales; no selecciona una. Si
\(u:\nu\to\nu'\), la acción contravariante es
\begin{equation}\label{eq:pdfv2-coh-w-action}
 W_{d,N}(u):W_{d,N}(\nu')\longrightarrow W_{d,N}(\nu),
 \qquad
 [\lambda]\longmapsto[\lambda\circ u].
\end{equation}
Cada generador conserva orientación, carry, frontera, ruta y memoria.

\section{Complejo bar bilateral y totalización de la bicadena}
\label{sec:pdfv2-coh-bar}

Para \(q\geq0\) y \(r\in\{0,1\}\), sea
\begin{equation}\label{eq:pdfv2-coh-bar-bigraded}
 B_{q,r}^{d,N}(x):=
 \bigoplus_{\nu_0\xrightarrow{u_1}\cdots\xrightarrow{u_q}\nu_q}
 W_{d,N}(\nu_q)\otimes_{A_N}\Gamma_{d,N}(\nu_0)_r.
\end{equation}
La diferencial bar se define sin abreviación por
\begin{align}
 d_{\rm bar}(w;u_1,\ldots,u_q;c)
 :={}&(w;u_2,\ldots,u_q;\Gamma(u_1)c)
 \nonumber\\
 &+\sum_{i=1}^{q-1}(-1)^i
 (w;u_1,\ldots,u_{i+1}u_i,\ldots,u_q;c)
 \nonumber\\
 &+(-1)^q
 (W(u_q)w;u_1,\ldots,u_{q-1};c).
 \label{eq:pdfv2-coh-bar-differential}
\end{align}
Para \(q=1\), ésta es exactamente la relación
\begin{equation}\label{eq:pdfv2-coh-coend-relation}
 d_{\rm bar}(w;u;c)
 =w\otimes\Gamma(u)c-W(u)w\otimes c.
\end{equation}
Así se identifica una incidencia transportada con la misma incidencia escrita
en la carta refinada; la relación no queda oculta bajo el nombre del
totalizador.

El complejo total es
\begin{equation}\label{eq:pdfv2-coh-total-complex}
 K_{d,N}(x):=\operatorname{Tot}B_{\bullet,\bullet}^{d,N}(x),
 \qquad
 D_{\rm tot}:=d_{\rm bar}+(-1)^q\partial_\Gamma
 \quad\text{sobre }B_{q,r}^{d,N}.
\end{equation}
La functorialidad de \(W\) y \(\Gamma\), junto con
\(\partial_\Gamma^2=0\), da
\begin{equation}\label{eq:pdfv2-coh-total-square-zero}
 D_{\rm tot}^2=0.
\end{equation}
En notación compacta, pero ahora ya descargada por
\eqref{eq:pdfv2-coh-bar-bigraded}--
\eqref{eq:pdfv2-coh-bar-differential},
\begin{equation}\label{eq:pdfv2-coh-derived-coend}
 K_{d,N}(x)
 \simeq
 W_{d,N}\otimes^{\mathbf L}_{A_N[J_d(x)]}\Gamma_{d,N}.
\end{equation}

\section{Colímite de profundidad y terminal nonádico}
\label{sec:pdfv2-coh-terminal}

Los prefijos forman el colímite
\begin{equation}\label{eq:pdfv2-coh-k-infinity}
 K_{\infty,N}(x):=\varinjlim_d K_{d,N}(x).
\end{equation}
La extensión de nueve fases lleva profundidad \(d\) a profundidad \(d+9\).
Sólo después de formar \eqref{eq:pdfv2-coh-k-infinity} induce un
endomorfismo bien tipado
\begin{equation}\label{eq:pdfv2-coh-u-gamma9}
 U_{\Gamma_9,N}:K_{\infty,N}(x)
 \longrightarrow K_{\infty,N}(x).
\end{equation}
El terminal es
\begin{equation}\label{eq:pdfv2-coh-terminal-cone}
 \mathcal T_{{\rm coh},N}(x)
 :=\operatorname{Cone}\bigl(I-U_{\Gamma_9,N}\bigr).
\end{equation}
Con la convención homológica
\[
 \operatorname{Cone}(f)_q
 =K_{\infty,N,q-1}\oplus K_{\infty,N,q},
\]
su diferencial es
\begin{equation}\label{eq:pdfv2-coh-cone-differential}
 \partial_{\rm Cone}(a,b)
 =\bigl(-\partial_Ka,\partial_Kb+f(a)\bigr).
\end{equation}

Aunque la fase visible de \(h\) y \(\Gamma_9h\) coincida, los generadores
\([h]\) y \([\Gamma_9h]\) son diferentes porque sus ledgers y profundidades
son diferentes. El terminal conserva
\begin{equation}\label{eq:pdfv2-coh-terminal-difference}
 (I-U_{\Gamma_9,N})[h]=[h]-[\Gamma_9h].
\end{equation}
Las reducciones de precisión conmutan con el terminal:
\begin{equation}\label{eq:pdfv2-coh-u-reduction}
 r_{N+1,N}U_{\Gamma_9,N+1}
 =U_{\Gamma_9,N}r_{N+1,N}.
\end{equation}

\section{Lector cohomológico terminal: ciclo de incidencia y clase en \(H_0(\operatorname{Cone}(I-U_{\Gamma_9,N}))\)}
\label{sec:pdfv2-coh-reader}

Cada celda de \(x_{\leq d}\) porta un vector de incidencia
\[
 \omega_{\nu,N}(x)
 \in A_N^{I_\nu\times J_\nu}=C_{\nu,N,0}
\]
y la ruta terminal completa \(\lambda_{\nu,x}\). El ciclo bruto es
\begin{equation}\label{eq:pdfv2-coh-z-cycle}
 z_{d,N}(x):=
 \sum_{\nu\in\operatorname{Cell}_d(x)}
 \epsilon_\nu[\lambda_{\nu,x}]
 \otimes\omega_{\nu,N}(x)
 \in B_{0,0}^{d,N}(x).
\end{equation}
Como el complejo total es homológico no negativo,
\begin{equation}\label{eq:pdfv2-coh-z-closed}
 D_{\rm tot}z_{d,N}(x)=0.
\end{equation}
Se conservan tanto el ciclo como su clase
\[
 [z_{d,N}(x)]\in H_0(K_{d,N}(x)).
\]
Tras incluirlo en \(K_{\infty,N}\), su clase terminal es
\[
 [(0,z_{d,N}(x))]
 \in H_0(\mathcal T_{{\rm coh},N}(x)).
\]
El lector completo es
\begin{equation}\label{eq:pdfv2-coh-reader-full}
 \mathcal L_{{\rm coh},d,N}(x)
 :=\bigl(
 z_{d,N}(x),[z_{d,N}(x)],
 U_{\Gamma_9,N}z_{d,N}(x),
 [(0,z_{d,N}(x))]
 \bigr).
\end{equation}
No se publica únicamente una clase cociente: permanecen el representante, la
ruta que lo produjo, su continuación y el terminal.

\section{Fibras supervivientes y correcciones sin selector}
\label{sec:pdfv2-coh-surviving-fibers}

Sea \(\mathscr S_\lambda^{\rm coh}\) el \(A_N\)-módulo libre generado por las
historias finitas admisibles en el nivel \(\lambda=(d,N)\). El lector se
extiende linealmente a
\begin{equation}\label{eq:pdfv2-coh-e-lambda}
 e_\lambda:\mathscr S_\lambda^{\rm coh}
 \longrightarrow\mathscr O_\lambda^{\rm coh},
 \qquad
 \mathscr O_\lambda^{\rm coh}
 :=H_0(K_{d,N})\oplus H_0(\mathcal T_{{\rm coh},N}).
\end{equation}
El codominio superviviente es
\begin{equation}\label{eq:pdfv2-coh-surviving-output}
 \mathscr O_\infty^{\rm surv}
 :=\operatorname{Im}\!\left(
 e_\infty:\mathcal C_{\rm cont}^{\rm coh}
 \longrightarrow\varprojlim_\lambda\mathscr O_\lambda^{\rm coh}
 \right).
\end{equation}
Para \(a=(a_\lambda)_\lambda\in\mathscr O_\infty^{\rm surv}\), defínase
\begin{equation}\label{eq:pdfv2-coh-fiber}
 \mathscr F_\lambda(a):=
 \left\{s_\lambda\in\mathscr S_\lambda^{\rm coh}:
 \begin{array}{l}
 e_\lambda(s_\lambda)=a_\lambda,\\
 \exists y\in\mathcal C_{\rm cont}^{\rm coh}\text{ con }
 y|_\lambda=s_\lambda\text{ y }e_\infty(y)=a
 \end{array}
 \right\}.
\end{equation}
Estas fibras son no vacías por la definición del codominio superviviente. Si
\(\lambda'\geq\lambda\), la restricción es sobreyectiva:
\begin{equation}\label{eq:pdfv2-coh-fiber-surjection}
 \rho_{\lambda',\lambda}:\mathscr F_{\lambda'}(a)
 \twoheadrightarrow\mathscr F_\lambda(a).
\end{equation}
En efecto, una historia completa que testimonia
\(s_\lambda\in\mathscr F_\lambda(a)\) proporciona por truncación un
antecedente en cualquier nivel posterior. No se elige ese antecedente como
parte del objeto.

La multisección completa es
\begin{equation}\label{eq:pdfv2-coh-extension-multisection}
 \Sigma_{\lambda',\lambda}^{\rm coh}(s_\lambda;a)
 :=\{s_{\lambda'}\in\mathscr F_{\lambda'}(a):
 \rho_{\lambda',\lambda}s_{\lambda'}=s_\lambda\}.
\end{equation}
Para una prolongación provisional \(\widetilde s_{\lambda'}\), sea
\begin{equation}\label{eq:pdfv2-coh-discrepancy}
 \delta_{\lambda'}:=
 a_{\lambda'}-e_{\lambda'}(\widetilde s_{\lambda'}).
\end{equation}
La relación de todas las correcciones es
\begin{equation}\label{eq:pdfv2-coh-correction-relation}
 \operatorname{Corr}_{\lambda',\lambda}
 (\widetilde s_{\lambda'},\delta_{\lambda'})
 :=\left\{\eta\in\mathscr S_{\lambda'}^{\rm coh}:
 \begin{array}{l}
 \rho_{\lambda',\lambda}\eta=0,\\
 e_{\lambda'}(\eta)=\delta_{\lambda'},\\
 \eta\text{ está soportada en celdas nuevas},\\
 \eta\text{ respeta orientación, carry, memoria y terminal}
 \end{array}
 \right\}.
\end{equation}
Si dos levantamientos son comparables, su diferencia
\(\eta=s_{\lambda'}-\widetilde s_{\lambda'}\) pertenece a esta relación
exactamente cuando satisface las cuatro condiciones anteriores. No se afirma
que \eqref{eq:pdfv2-coh-correction-relation} sea no vacía para una
discrepancia arbitraria: la no vacuidad gobernante es la de
\eqref{eq:pdfv2-coh-extension-multisection} sobre el dominio superviviente.

\iffalse
\section{Las trece coordenadas de la restricción}
\label{sec:pdfv2-coh-thirteen}

Para \(x_\lambda\in\mathcal C_{{\rm cont},\lambda}^{\rm coh}\), sea
\begin{equation}\label{eq:pdfv2-coh-d-package}
 \begin{aligned}
 D_{{\rm coh},\lambda}(x_\lambda):=\bigl(&
 \mathcal F_{{\rm coh},\lambda},
 \operatorname{Ext}_{\Gamma,{\rm coh},\lambda},
 \operatorname{Rel}_{{\rm coh},\lambda},
 \mathcal N_{{\rm coh},\lambda},
 \operatorname{or}_{{\rm coh},\lambda},
 \operatorname{Carr}_{{\rm coh},\lambda},
 \operatorname{Mem}_{{\rm coh},\lambda},\\
 &\partial_{{\rm coh},\lambda},
 \gamma_{{\rm coh},\lambda},
 \operatorname{Msect}_{{\rm coh},\lambda},
 \operatorname{Surv}_{{\rm coh},\lambda},
 \operatorname{Term}_{{\rm coh},\lambda},
 \mathcal L_{{\rm coh},\lambda}\bigr).
 \end{aligned}
\end{equation}
Sus coordenadas son:
\begin{enumerate}
 \item \(\mathcal F_{{\rm coh},\lambda}\): todos los módulos
       \(A_N^{I_\nu}\), \(A_N^{J_\nu}\),
       \(A_N^{I_\nu\times J_\nu}\), los complejos \(C_{\nu,N}\), los
       módulos \(W_{d,N}(\nu)\) y \(\mathscr S_\lambda^{\rm coh}\).
 \item \(\operatorname{Ext}_{\Gamma,{\rm coh},\lambda}\): todos los mapas
       \(\Gamma(u)\), reducciones de coeficientes, truncaciones, extensiones
       \(\Gamma_9\) y las relaciones \(\Sigma_{\lambda',\lambda}\).
 \item \(\operatorname{Rel}_{{\rm coh},\lambda}\):
       \(\kappa\), \(\delta\), la sucesión exacta, las caras bar, la
       relación de coend, la orientación y las leyes de composición.
 \item \(\mathcal N_{{\rm coh},\lambda}\):
       \[
       (d,N,3^N,a_\nu,b_\nu,(a_\nu-1)(b_\nu-1),
       \rho_{\Sigma,\nu},\rho_{\Pi,\nu},c_{\Sigma,\nu},c_{\Pi,\nu}).
       \]
 \item \(\operatorname{or}_{{\rm coh},\lambda}\): los signos
       \(\epsilon_\nu\), los órdenes de puertos y los mapas
       \(\Theta_{\nu,N}\).
 \item \(\operatorname{Carr}_{{\rm coh},\lambda}\): las matrices
       \(\beta_N\), los cociclos \(c_N(v,u)\), residuos y cocientes,
       conservados por separado.
 \item \(\operatorname{Mem}_{{\rm coh},\lambda}\): el ledger heredado y la
       palabra ordenada de refinamientos, carries, correcciones y avances
       \(\Gamma_9\), no sólo su valor agregado.
 \item \(\partial_{{\rm coh},\lambda}\):
       \(\kappa_{\nu,N}\), \(D_{\rm tot}\),
       \(\partial_{\rm Cone}\) y las fronteras de las celdas.
 \item \(\gamma_{{\rm coh},\lambda}\): todas las cadenas componibles
       \(\nu_0\to\cdots\to\nu_q\) y sus continuaciones terminales.
 \item \(\operatorname{Msect}_{{\rm coh},\lambda}\): todas las celdas,
       rutas terminales, extensiones \(\Sigma\) y correcciones
       \(\operatorname{Corr}\), sin selector.
 \item \(\operatorname{Surv}_{{\rm coh},\lambda}\): el sistema de fibras
       \(\mathscr F_\lambda(a)\) y restricciones sobreyectivas.
 \item \(\operatorname{Term}_{{\rm coh},\lambda}\):
       \(\operatorname{Cone}(I-U_{\Gamma_9,N})\), junto con \(z\), \(Uz\)
       y su diferencia.
 \item \(\mathcal L_{{\rm coh},\lambda}\):
       \((z,[z],Uz,[(0,z)])\) y todas sus compatibilidades de reducción.
\end{enumerate}

Para \(\lambda'\geq\lambda\), el transporte coordinado satisface
\begin{equation}\label{eq:pdfv2-coh-d-naturality}
 u_{\lambda',\lambda}^{\rm coh}
 D_{{\rm coh},\lambda'}(x_{\lambda'})
 =D_{{\rm coh},\lambda}
  (\rho_{\lambda',\lambda}x_{\lambda'}).
\end{equation}

\section{Grafo de restricción y objeto dependiente}
\label{sec:pdfv2-coh-graphs}

La restricción total es el grafo
\begin{equation}\label{eq:pdfv2-coh-g-graph}
 \mathcal G_{\rm coh}:=\operatorname{Graph}(D_{\rm coh}),
 \qquad
 \operatorname{Res}_{\rm coh}(x):=(x,D_{\rm coh}x).
\end{equation}
Así,
\begin{equation}\label{eq:pdfv2-coh-res-inverse}
 \pi_1\operatorname{Res}_{\rm coh}=I,
 \qquad
 \operatorname{Res}_{\rm coh}\pi_1=I_{\mathcal G_{\rm coh}}.
\end{equation}

El extractor completo es
\begin{equation}\label{eq:pdfv2-coh-chi}
 \chi_{\rm coh}(x):=
 \bigl(
 \{C_{\nu,N}\},
 \{z_\lambda,[z_\lambda]\},
 \{\mathscr F_\lambda(e_\infty x)\},
 \{\Sigma_{\lambda',\lambda}\},
 \{\operatorname{Corr}_{\lambda',\lambda}\}
 \bigr).
\end{equation}
Se define
\begin{equation}\label{eq:pdfv2-coh-xchi}
 X_{\chi,{\rm coh}}
 :=\{(\chi_{\rm coh}(x),x,D_{\rm coh}x):
 x\in\mathcal C_{\rm cont}^{\rm coh}\}
\end{equation}
y
\begin{equation}\label{eq:pdfv2-coh-prx}
 \operatorname{pr}_{X,{\rm coh}}(x,D_{\rm coh}x)
 :=(\chi_{\rm coh}(x),x,D_{\rm coh}x).
\end{equation}
La proyección sobre \((x,D_{\rm coh}x)\) es inversa sobre este objeto
dependiente.

\section{Ensamblador y publicador internos}
\label{sec:pdfv2-coh-assembler-publisher}

El ambiente \(\mathscr M_{\rm coh}^{\rm amb}\) consta de sistemas
compatibles de sucesiones exactas de celdas, categorías, funtores, módulos de
ruta, totalizadores, terminales, lectores, fibras y multisecciones. El
ensamblador bruto es
\begin{equation}\label{eq:pdfv2-coh-assembler}
 \begin{aligned}
 a_{\rm coh}:X_{\chi,{\rm coh}}&\longrightarrow
 \mathscr M_{\rm coh}^{\rm amb},\\
 a_{\rm coh}(\chi(x),x,Dx)
 &:={}
 \bigl(
 \{C_{\nu,N},\kappa_{\nu,N},\delta_{\nu,N}\},
 \{J_d,\Gamma,W,K\},\\
 &\hspace{31mm}
 \{U_{\Gamma_9,N},\mathcal T_N\},
 \{e_\lambda,\mathscr F_\lambda,\Sigma,\operatorname{Corr}\}
 \bigr)_x.
 \end{aligned}
\end{equation}
La macroestructura es otro grafo:
\begin{equation}\label{eq:pdfv2-coh-macro-graph}
 \mathfrak M_{\rm coh}:=\operatorname{Graph}(a_{\rm coh}),
 \qquad
 \mathcal A_{\rm coh}(z):=(z,a_{\rm coh}z).
\end{equation}

Sea \(\mathscr Y_{\rm coh}^{\rm int}\) el tipo de certificados que contiene
\begin{equation}\label{eq:pdfv2-coh-certificate-list}
 \begin{gathered}
 \text{la sucesión exacta \eqref{eq:pdfv2-coh-exact-sequence}},\\
 \kappa_{b,a}\tau=-P\kappa_{a,b},
 \qquad D_{\rm tot}^2=0,
 \qquad\partial_{\rm Cone}^2=0,\\
 rU_{\Gamma_9}=U_{\Gamma_9}r,
 \qquad
 e_\lambda\rho_{\lambda',\lambda}
 =r_{\lambda',\lambda}e_{\lambda'},\\
 \rho_{\lambda',\lambda}:\mathscr F_{\lambda'}(a)
 \twoheadrightarrow\mathscr F_\lambda(a),
 \qquad
 \Sigma_{\lambda',\lambda}(s_\lambda;a)\neq\varnothing,
 \end{gathered}
\end{equation}
junto con los ciclos brutos, sus clases y sus terminales. El publicador
\begin{equation}\label{eq:pdfv2-coh-publisher}
 p_{\rm coh}^{\rm raw}:\mathfrak M_{\rm coh}
 \longrightarrow\mathscr Y_{\rm coh}^{\rm int}
\end{equation}
produce ese certificado sin eliminar la macroestructura. Finalmente,
\begin{equation}\label{eq:pdfv2-coh-output-graph}
 \mathcal C_{\rm coh}^{\rm HMT}
 :=\operatorname{Graph}(p_{\rm coh}^{\rm raw}),
 \qquad
 \mathcal P_{\rm coh}(m):=(m,p_{\rm coh}^{\rm raw}m).
\end{equation}

\section{Teorema interno de generación coherente}
\label{sec:pdfv2-coh-theorem}

\begin{theorem}[Generación interna de la macroestructura de coherencia]
\label{thm:pdfv2-coh-internal-generation}
Sobre el dominio explícito \(\mathcal C_{\rm cont}^{\rm coh}\), la cadena
\[
 \mathcal C_{\rm cont}^{\rm coh}
 \xrightarrow{\operatorname{Res}_{\rm coh}}
 \mathcal G_{\rm coh}
 \xrightarrow{\operatorname{pr}_{X,{\rm coh}}}
 X_{\chi,{\rm coh}}
 \xrightarrow{\mathcal A_{\rm coh}}
 \mathfrak M_{\rm coh}
 \xrightarrow{\mathcal P_{\rm coh}}
 \mathcal C_{\rm coh}^{\rm HMT}
\]
construye de forma natural las celdas, sus conos, el diagrama de historias, el
bar de dos lados, la totalización, el terminal nonádico y las fibras
supervivientes. Conserva las trece coordenadas y todas las correcciones como
multisecciones. Cada etapa admite, sobre su imagen, una proyección que recupera
la etapa anterior.
\end{theorem}

\begin{proof}
La \cref{prop:pdfv2-coh-cell-exact} construye cada celda sobre \(A_N\), con
núcleo diagonal y defecto \((a_\nu-1)(b_\nu-1)\). La identidad
\eqref{eq:pdfv2-coh-orientation} hace del intercambio orientado un mapa de
complejos involutivo. Las fórmulas
\eqref{eq:pdfv2-coh-gamma-one}--\eqref{eq:pdfv2-coh-gamma-zero} prueban que
cada refinamiento es un mapa de complejos, mientras que
\eqref{eq:pdfv2-coh-kappa-reduction} prueba la compatibilidad en precisión.

La precomposición de rutas hace contravariante a \(W\). Las identidades
functoriales de \(W\) y \(\Gamma\) implican
\(d_{\rm bar}^2=0\); su conmutación con la diferencial celular da
\(D_{\rm tot}^2=0\). Así, \(K_{d,N}\) queda construido por las fórmulas bar,
no por una invocación nominal del coend.

La extensión \(\Gamma_9\) aumenta la profundidad. Por ello sólo induce el
endomorfismo \(U_{\Gamma_9,N}\) después de pasar al colímite
\(K_{\infty,N}\). El cono
\(\operatorname{Cone}(I-U_{\Gamma_9,N})\) conserva entonces la diferencia
\([h]-[\Gamma_9h]\) y no identifica la repetición de fase con la repetición
de estado.

El ciclo \(z_{d,N}(x)\) se forma con todas las celdas, signos, rutas e
incidencias de la historia. Truncación y reducción conmutan con el lector. Si
\(s_\lambda\in\mathscr F_\lambda(a)\), la historia completa que aparece en
la propia definición de la fibra proporciona un antecedente en todo nivel
\(\lambda'\geq\lambda\); esto prueba la sobreyectividad
\eqref{eq:pdfv2-coh-fiber-surjection}. El objeto conserva el conjunto entero
\(\Sigma_{\lambda',\lambda}\), no una elección. Las diferencias entre
levantamientos se registran mediante
\eqref{eq:pdfv2-coh-correction-relation}, con soporte, orientación, carry,
memoria y terminal explícitos.

Finalmente, \(\mathcal G_{\rm coh}\), \(\mathfrak M_{\rm coh}\) y
\(\mathcal C_{\rm coh}^{\rm HMT}\) son grafos, mientras que
\(X_{\chi,{\rm coh}}\) conserva como coordenadas \(x\) y \(D_{\rm coh}x\).
La primera proyección de cada objeto recupera el anterior. Por tanto, la
composición completa no sustituye historia por clase, multisección por
selector ni terminal por fase visible.
\end{proof}

\section{Procedencia probatoria y falsadores}
\label{sec:pdfv2-coh-provenance-falsifiers}

\paragraph{Resultado recuperado.}
Pertenecen a la arquitectura ya construida la celda \(\kappa_{a,b}\), su
núcleo diagonal y defecto, la orientación
\(\kappa_{b,a}\tau=-P\kappa_{a,b}\), el carry APP, el cono de la celda, la
categoría de historias, la totalización por coend, la acción \(\Gamma_9\), el
terminal \(\operatorname{Cone}(I-U_{\Gamma_9})\) y la arquitectura de
extensión superviviente.

\paragraph{Formalización nueva.}
Son formalizaciones explícitas la sucesión exacta
\eqref{eq:pdfv2-coh-exact-sequence} sobre \(A_N\), la separación entre
profundidad y precisión, las caras completas del bar, la formación del
terminal sólo después del colímite, el empaquetado de las trece coordenadas,
las correcciones como relación sin selector y los cuatro objetos dependientes
o grafos que impiden pérdida de genealogía.

\begin{criterion}[Falsadores internos de la rama \(\mathrm{coh}\)]
\label{crit:pdfv2-coh-falsifiers}
La construcción falla si ocurre cualquiera de los hechos siguientes:
\begin{enumerate}
 \item no es exacta la sucesión
       \eqref{eq:pdfv2-coh-exact-sequence};
 \item falla \(\kappa_{b,a}\tau=-P\kappa_{a,b}\);
 \item algún refinamiento no satisface
       \(\kappa_{\nu'}\Gamma(u)_1=\Gamma(u)_0\kappa_\nu\);
 \item el carry de composición no cumple
       \eqref{eq:pdfv2-coh-carry-cocycle};
 \item \(D_{\rm tot}^2\neq0\);
 \item reducción, lector o terminal no conmutan con \(U_{\Gamma_9}\);
 \item se forma \(\operatorname{Cone}(I-U_{\Gamma_9})\) identificando antes
       las profundidades \(d\) y \(d+9\);
 \item una fibra declarada superviviente es vacía o una restricción no es
       sobreyectiva;
 \item se sustituye \(\Sigma_{\lambda',\lambda}\) por un único
       levantamiento;
 \item se publica sólo \([z]\) y se eliminan \(z\), ruta, carry, memoria o
       terminal;
 \item se omite cualquiera de las trece coordenadas.
\end{enumerate}
En el último caso el diagnóstico vinculante es: \emph{objeto sustituido;
conclusión no transferible}.
\end{criterion}

\begin{warningbox}
Este capítulo concluye únicamente la construcción interna de
\(\mathcal C_{\rm coh}^{\rm HMT}\). No introduce un espacio clásico, una
clase clásica ni una publicación clásica. Cualquier entrelazador posterior
pertenece a otro documento y no gobierna la existencia de esta
macroestructura.
\end{warningbox}
\fi
 \endgroup

El complejo local \eqref{eq:pdfv2-cor-complejo-local} y el vértice
\(y_\alpha\), extremo de esa misma extensión en el nervio, forman el pullback
aumentado. En él se definen
\begin{equation}
\label{eq:pdfv2-cor-zpm-h}
 \mathbf1_\alpha=(r,[y_\alpha]),\qquad
 z_-(\alpha)=(r,[y_\alpha]),
\end{equation}
\begin{equation}
\label{eq:pdfv2-cor-zplus-h}
 z_+(\alpha)=
 (r+3c_\alpha v_\alpha e+v_\alpha b,[y_\alpha]),
 \qquad
 h_*(\alpha)=(-v_\alpha f,0).
\end{equation}
La acción en generadores da, sin una existencia abstracta adicional,
\begin{equation}
\label{eq:pdfv2-cor-filler}
 \partial z_-=\partial z_+=0,\qquad
 \partial h_*=z_--z_+.
\end{equation}
El relleno correlativo, la celda y el collar conservan el mismo
\((\alpha,\operatorname{or}(\alpha),\kappa(\alpha),\partial\alpha,
\gamma\alpha,m(\gamma\alpha;m_0))\); sólo aplican lectores internos distintos
al registro común.

Para el mismo complejo finito basado se define, sin pedir balance ni
aciclicidad, la línea determinante total
\begin{equation}
\label{eq:pdfv2-cor-linea-determinante-total}
 \operatorname{Det}(C_{k,N}):=
 \bigotimes_q
 \left(\bigwedge^{\rm top}C_{k,N,q}\right)^{(-1)^q}.
\end{equation}
El refinamiento induce el morfismo determinantal del mapa de complejos y
transporta a la vez los módulos de homología, las formas normales de Smith y
las páginas de Bockstein como campos del registro. Ni la existencia de
\eqref{eq:pdfv2-cor-linea-determinante-total} ni su transporte exige que la
homología se anule. Si el cono del refinamiento es contractible y orientado,
el morfismo se trivializa canónicamente; en el caso general se conserva el
cono y no se lo sustituye por un isomorfismo ficticio.

\subsection{Forma normal de Smith del operador terminal y descomposición primaria de su conúcleo}
Para cada prefijo \(y\), el acarreo no orientado y el coeficiente orientado
de frontera del registro común admiten los levantamientos
\begin{equation}
\label{eq:detint-c-v-lifts}
 \widetilde c_y,\widetilde v_y\in\mathbf Z,\qquad
 c_{y,n}:=\widetilde c_y\bmod3^n,\qquad
 v_{y,n}:=\widetilde v_y\bmod3^n.
\end{equation}
La salida determinantal no es una entrada ni una rama separada: es el sector
producido por la misma reducción terminal cuando se verifican las condiciones
tipadas
\begin{equation}
\label{eq:detint-dominio}
\begin{split}
 \mathcal C_{\rm cont}^{\rm det}:=\bigl\{x\in\mathcal C_{\rm cont}^{\rm disc}:{}&
 x\text{ sobrevive a toda profundidad};\\
 &\widehat P_x^{\rm red}\text{ es finito y basado};\\
 &\chi(\widehat P_x^{\rm red})=0;\\
 &H_*(\widehat P_x^{\rm red}\otimes_{\mathbf Z_3}K_3)=0;\\
 &\text{la reducción produce un bloque cuadrado estable bajo refinamiento}
 \bigr\}.
\end{split}
\end{equation}
\begingroup
  \let\section\subsection
  % Extracción quirúrgica del fuente teoremática V2 05e: líneas 325--447.
% El resto permanece en este archivo como procedencia, pero no se publica.
\iffalse
\chapter[Macroestructura determinantal interna]{Macroestructura determinantal interna generada por el continuo discreto}
\label{chap:detint-macroestructura}

\begin{thesisbox}
La rama \({\rm det}\) parte exclusivamente de historias supervivientes de
\(\mathcal C_{\rm cont}^{\rm disc}\). El nervio de extensiones, la diagonal
Alexander--Whitney, el complejo local, la unidad, las dos publicaciones, el
filler, la reducción terminal y la torre \(3\)-ádica se construyen antes de
cualquier evaluación posterior. Las condiciones de balance y aciclicidad se
incorporan al dominio tipado y nunca se presuponen silenciosamente.
\end{thesisbox}

\section{Dominio determinantal superviviente}

Sea
\begin{equation}\label{eq:detint-an}
 A_n:=\mathbf Z/3^n\mathbf Z,
 \qquad
 \rho_{n+1,n}:A_{n+1}\longrightarrow A_n.
\end{equation}
Para \(x\in\mathcal C_{\rm cont}^{\rm disc}\), sea
\(\mathsf T_m(x)\) la categoría finita de prefijos TPK de profundidad \(m\)
compatibles con \(x\). Cada objeto conserva hoja, orientación, residuo,
cociente, carry entero, memoria, frontera, ruta, multisección, supervivencia
y terminal; sus morfismos son las extensiones admisibles que conservan esas
coordenadas.

De cada objeto \(y\in\mathsf T_m(x)\) se extraen los enteros internos
\begin{equation}\label{eq:detint-c-v-lifts}
 \widetilde c_y,\widetilde v_y\in\mathbf Z,
 \qquad
 c_{y,n}:=\widetilde c_y\bmod3^n,
 \qquad
 v_{y,n}:=\widetilde v_y\bmod3^n.
\end{equation}
La primera coordenada es el carry no orientado y la segunda, el coeficiente
orientado de frontera. La reversión satisface
\begin{equation}\label{eq:detint-or-cv}
 c_{\bar y,n}=c_{y,n},
 \qquad
 v_{\bar y,n}=-v_{y,n}.
\end{equation}

Se escriben
\begin{equation}\label{eq:detint-z3-k3}
 \mathbf Z_3:=\varprojlim_n A_n,
 \qquad
 K_3:=\operatorname{Frac}(\mathbf Z_3).
\end{equation}
La restricción determinantal se define por
\begin{equation}\label{eq:detint-dominio}
\begin{split}
 \mathcal C_{\rm cont}^{\rm det}:=
 \bigl\{x\in\mathcal C_{\rm cont}^{\rm disc}:{}&
 x\text{ sobrevive a toda profundidad};\\
 &\widehat P_x^{\rm red}\text{ es finito y basado};\\
 &\chi(\widehat P_x^{\rm red})=0;\\
 &H_*(\widehat P_x^{\rm red}\otimes_{\mathbf Z_3}K_3)=0;\\
 &\text{la reducción canónica produce un bloque cuadrado}\\
 &\text{estable bajo refinamiento}\bigr\}.
\end{split}
\end{equation}
Los objetos \(\widehat P_x^{\rm red}\) y el bloque cuadrado se construirán
en \cref{sec:detint-terminal}. Por tanto
\eqref{eq:detint-dominio} es una condición de pertenencia comprobable y no
una conclusión insertada en la definición de los mapas.

\section{Nervio, Alexander--Whitney y counidad}

Se define el complejo normalizado
\begin{equation}\label{eq:detint-gmn}
 G_{m,n}(x):=
 N_*^{\rm norm}\!\left(N\mathsf T_m(x);A_n\right).
\end{equation}
Para un símplice no degenerado
\(\sigma=[y_0\to y_1\to\cdots\to y_q]\), la diagonal es
\begin{equation}\label{eq:detint-aw}
 \Delta_{\rm AW}(\sigma)
 =\sum_{j=0}^q
 [y_0\to\cdots\to y_j]\otimes
 [y_j\to\cdots\to y_q].
\end{equation}
La counidad viene dada por
\begin{equation}\label{eq:detint-counidad-g}
 \epsilon_G([y])=1,
 \qquad
 \epsilon_G(\sigma)=0\quad(q>0).
\end{equation}
En particular,
\begin{equation}\label{eq:detint-vertice-grouplike}
 \Delta_{\rm AW}[y]=[y]\otimes[y],
 \qquad
 \epsilon_G([y])=1.
\end{equation}

Si \(m'\ge m\) y \(n'\ge n\), truncación y reducción de coeficientes
inducen
\begin{equation}\label{eq:detint-rmn}
 R_{m',m}^{n',n}:G_{m',n'}(x)\longrightarrow G_{m,n}(x),
\end{equation}
y se verifican
\begin{equation}\label{eq:detint-aw-natural}
 R\partial=\partial R,
 \qquad
 (R\otimes R)\Delta_{\rm AW}=\Delta_{\rm AW}R,
 \qquad
 \epsilon_GR=\rho_{n',n}\epsilon_G.
\end{equation}
La ruta no se reduce a sus extremos: el símplice entero permanece en
\(G_{m,n}(x)\).

\section{Complejo local interno}

Para \(c\in A_n\), se define
\begin{equation}\label{eq:detint-p0-grados}
 P^0_{n,1}(c)=A_nf,
 \qquad
 P^0_{n,0}(c)=A_nr\oplus A_ne\oplus A_nb,
 \qquad
 P^0_{n,-1}(c)=A_ng,
\end{equation}
con diferencial
\begin{equation}\label{eq:detint-p0-diferencial}
 \partial f=3ce+b,
 \qquad
 \partial r=0,
 \qquad
 \partial e=g,
 \qquad
 \partial b=-3cg.
\end{equation}
La identidad de cadena es literal:
\begin{equation}\label{eq:detint-d2}
 \partial^2f=3c\,\partial e+\partial b=3cg-3cg=0,
 \qquad
 \partial^2e=\partial^2b=\partial^2r=0.
\end{equation}

La aumentación
\begin{equation}\label{eq:detint-epsilon-p}
 \epsilon_P:P_n^0(c)\longrightarrow A_n[0]
\end{equation}
se fija por
\[
 \epsilon_P(r)=1,
 \qquad
 \epsilon_P(e)=\epsilon_P(b)=\epsilon_P(f)=\epsilon_P(g)=0.
\]
El complejo posee la coalgebra diferencial
\begin{equation}\label{eq:detint-coalgebra-p}
 \Delta_P(r)=r\otimes r,
 \qquad
 \Delta_P(q)=q\otimes r+r\otimes q
 \quad(q\in\{e,b,f,g\}).
\end{equation}
Así, \(r\) es group-like y los otros cuatro generadores son primitivos
relativos a \(r\). Las ecuaciones
\eqref{eq:detint-p0-diferencial} prueban que \(\Delta_P\) conmuta con el
diferencial tensorial.

\section{Sistema local de extensiones y ledger de memoria}

Para una extensión \(\alpha:y\to y'\), definimos
\begin{equation}\label{eq:detint-psi}
 \Psi_\alpha:P_n^0(c_y)\longrightarrow P_n^0(c_{y'})
\end{equation}
mediante
\begin{equation}\label{eq:detint-psi-generadores}
\begin{gathered}
 \Psi_\alpha(r_y)=r_{y'},\quad
 \Psi_\alpha(e_y)=e_{y'},\quad
 \Psi_\alpha(g_y)=g_{y'},\quad
 \Psi_\alpha(f_y)=f_{y'},\\
 \Psi_\alpha(b_y)=b_{y'}+3(c_{y'}-c_y)e_{y'}.
\end{gathered}
\end{equation}
Es un mapa de complejos porque
\begin{equation}\label{eq:detint-psi-b-chain}
 \partial\Psi_\alpha(b_y)
 =-3c_{y'}g+3(c_{y'}-c_y)g
 =-3c_yg
 =\Psi_\alpha(\partial b_y),
\end{equation}
y
\begin{equation}\label{eq:detint-psi-f-chain}
 \Psi_\alpha(\partial f_y)
 =3c_ye+b+3(c_{y'}-c_y)e
 =3c_{y'}e+b
 =\partial\Psi_\alpha(f_y).
\end{equation}
Además,
\begin{equation}\label{eq:detint-psi-functorial}
 \Psi_\beta\Psi_\alpha=\Psi_{\beta\alpha}.
\end{equation}

El cambio de la coordenada orientada no se borra. Se registra por
\begin{equation}\label{eq:detint-k-alpha}
 K_\alpha:=(v_{y'}-v_y)f_{y'}.
\end{equation}
Para extensiones componibles,
\begin{equation}\label{eq:detint-k-coherencia}
 K_{\beta\alpha}=K_\beta+\Psi_\beta K_\alpha.
\end{equation}
Esta igualdad es el ledger exacto del incremento de frontera y hace visible
la memoria que una composición puramente nominal perdería.

El diagrama local completo es la construcción de Grothendieck
\begin{equation}\label{eq:detint-grothendieck-local}
 \int_{y\in\mathsf T_m(x)}P_n^0(c_y).
\end{equation}
En ella una flecha sobre \(\alpha:y\to y'\) lleva como componente lineal
\(\Psi_\alpha\) y como coherencia \(K_\alpha\). Cuando se aplica a un
elemento soportado en un vértice, se usa la notación
\begin{equation}\label{eq:detint-psi-soporte-vertice}
 \Psi_\alpha(p,[y]):=(\Psi_\alpha p,[y']).
\end{equation}
La arista \([y\to y']\) permanece en la coordenada de ruta del nervio. Esta
convención será la que se emplee en las fórmulas de naturalidad de
\(z_\pm\) y \(h_*\); no se afirma un mapa global que borre los restantes
símplices de \(G_{m,n}(x)\).

\section{Pullback, unidad, raíz, publicaciones y filler}

El pullback de complejos es
\begin{equation}\label{eq:detint-pgen}
 P^{\rm gen}_{y;m,n}
 :=P_n^0(c_y)\times_{A_n[0]}G_{m,n}(x)
 =\{(p,\gamma):\epsilon_P(p)=\epsilon_G(\gamma)\}.
\end{equation}
Su diferencial es
\begin{equation}\label{eq:detint-pgen-d}
 \partial(p,\gamma)=(\partial p,\partial\gamma).
\end{equation}
La unidad común asociada al vértice \(y\) es
\begin{equation}\label{eq:detint-unidad}
 \mathbf1_y:=(r,[y]).
\end{equation}
Sus dos proyecciones son group-like por
\eqref{eq:detint-vertice-grouplike} y
\eqref{eq:detint-coalgebra-p}; ése es el sentido preciso de unidad
group-like compatible del pullback.

La proyección raíz es el mapa de complejos
\begin{equation}\label{eq:detint-proyeccion-root}
 \varpi_{\rm root}:P^{\rm gen}_{y;m,n}\longrightarrow
 \operatorname{Root}_{y;m,n},
 \qquad
 \varpi_{\rm root}(p,\gamma):=(\epsilon_P(p)r,\gamma).
\end{equation}
Se definen las dos publicaciones internas
\begin{equation}\label{eq:detint-zpm}
 z_-(y):=(r,[y]),
 \qquad
 z_+(y):=(r+3c_yv_ye+v_yb,[y]).
\end{equation}
Ambas son ciclos:
\begin{equation}\label{eq:detint-zpm-ciclos}
 \partial z_-(y)=0,
 \qquad
 \partial z_+(y)=3c_yv_yg-3c_yv_yg=0.
\end{equation}
Sus raíces coinciden:
\begin{equation}\label{eq:detint-raiz-comun}
 \varpi_{\rm root}z_-(y)
 =(r,[y])
 =\varpi_{\rm root}z_+(y).
\end{equation}

El filler se construye directamente:
\begin{equation}\label{eq:detint-hstar}
 h_*(y):=(-v_yf,0)\in P^{\rm gen}_{y;m,n,1}.
\end{equation}
Entonces
\begin{equation}\label{eq:detint-filler-identidad}
 \partial h_*(y)
 =(-3c_yv_ye-v_yb,0)
 =z_-(y)-z_+(y).
\end{equation}
En particular,
\begin{equation}\label{eq:detint-clases-iguales}
 [z_-(y)]=[z_+(y)]
 \quad\text{en}\quad H_0(P^{\rm gen}_{y;m,n}),
\end{equation}
y la igualdad conserva su cadena testigo.

Bajo una extensión \(\alpha:y\to y'\), se cumplen
\begin{equation}\label{eq:detint-z-natural}
 z_-(y')=\Psi_\alpha z_-(y),
 \qquad
 z_+(y')-\Psi_\alpha z_+(y)=\partial K_\alpha,
\end{equation}
y
\begin{equation}\label{eq:detint-h-natural}
 h_*(y')-\Psi_\alpha h_*(y)=-K_\alpha.
\end{equation}
Por tanto, la naturalidad no identifica valores distintos de \(v\): registra
su diferencia mediante la homotopía coherente \(K_\alpha\).

\section{Orientación interna}

Sobre \(P_n^0(c)\) se define
\begin{equation}\label{eq:detint-omega-p}
 \Omega(r)=r,
 \qquad
 \Omega(q)=-q\quad(q\in\{e,b,f,g\}).
\end{equation}
Es un automorfismo de complejos y de la coalgebra relativa. En el nervio,
la inversión orientada es
\begin{equation}\label{eq:detint-omega-g}
 \mathfrak o_G[y_0\to\cdots\to y_q]
 :=(-1)^{q(q+1)/2}
 [\bar y_q\to\cdots\to\bar y_0].
\end{equation}
Junto con \eqref{eq:detint-or-cv}, estas acciones dan
\begin{equation}\label{eq:detint-or-z-h}
 \Omega z_-(y)=z_-(\bar y),
 \qquad
 \Omega z_+(c_y,v_y)=z_+(c_{\bar y},v_{\bar y}),
 \qquad
 \Omega h_*(y)=h_*(\bar y).
\end{equation}
La orientación actúa sobre ruta, complejo, publicaciones, filler y línea
determinantal; no es una etiqueta sin acción.

\fi
\section{Elevación del conúcleo mediante la torre de Bockstein y conservación del acarreo}
\label{sec:detint-terminal}

Se elimina únicamente la raíz común:
\begin{equation}\label{eq:detint-p-reducido}
 \overline P_{y;m,n}
 :=P^{\rm gen}_{y;m,n}/A_n\mathbf1_y,
 \qquad
 \widehat P_{y;m}^{\rm red}:=\varprojlim_n\overline P_{y;m,n}.
\end{equation}
La base se ordena primero por el orden TPK de rutas y después por
\begin{equation}\label{eq:detint-orden-base}
 f\prec e\prec b\prec g.
\end{equation}
Se ejecuta el siguiente algoritmo interno.
\begin{enumerate}
 \item Buscar los coeficientes que son unidades de \(\mathbf Z_3\).
 \item Elegir el primero según \eqref{eq:detint-orden-base} y el orden de
 rutas.
 \item Cancelar el par correspondiente y registrar las operaciones
 elementales, su signo y su contribución a la orientación.
 \item Repetir hasta que no quede un pivote unitario cancelable.
\end{enumerate}
La condición de supervivencia de \eqref{eq:detint-dominio} exige que,
cofinalmente en \(m\), el resultado sea un bloque de dos términos
\begin{equation}\label{eq:detint-sx}
 \mathbf Z_3^{r_x}\xrightarrow{S_x}\mathbf Z_3^{r_x},
 \qquad
 \det S_x\ne0\text{ en }K_3,
\end{equation}
y que un refinamiento sólo añada bloques contractibles unitarios y cambios
de base registrados. Si esas condiciones fallan, la historia no pertenece a
\(\mathcal C_{\rm cont}^{\rm det}\); no se fuerza un terminal inexistente.

La ruta identidad aporta un testigo elemental. Después de retirar la raíz,
el complejo local es
\begin{equation}\label{eq:detint-testigo-contractil}
 A_nf\xrightarrow{\binom{3c}{1}}
 A_ne\oplus A_nb\xrightarrow{(1,-3c)}A_ng.
\end{equation}
Es contractible mediante
\begin{equation}\label{eq:detint-contraccion-local}
 s(g)=e,
 \qquad
 s(ae+\beta b)=\beta f,
 \qquad
 s(f)=0.
\end{equation}
Así, las condiciones terminales admiten al menos el testigo identidad
interno y no describen un subdominio formalmente vacío.

\section{Smith, orden \texorpdfstring{\(3\)}{3}-ádico y unidad inicial}

En bases orientadas, una diagonalización por valoración tiene la forma
\begin{equation}\label{eq:detint-smith}
 U_xS_xV_x
 =\operatorname{diag}(3^{a_1}u_1,\ldots,3^{a_r}u_r),
 \qquad
 a_i\in\mathbf N,quad u_i\in\mathbf Z_3^\times.
\end{equation}
Se define
\begin{equation}\label{eq:detint-orden-d}
 d_x:=v_3(\det S_x)=\sum_{i=1}^r a_i
\end{equation}
y la unidad inicial
\begin{equation}\label{eq:detint-leading-unit}
 \lambda_x:=3^{-d_x}\det S_x\in\mathbf Z_3^\times.
\end{equation}
La coordenada de orientación trivializa la línea determinantal. Si una
transición produce
\begin{equation}\label{eq:detint-estabilidad-bloques}
 S_x\oplus I_q=U S_{x'}V,
\end{equation}
la orientación se transporta por el factor inverso de \(\det U\det V\).
La unidad orientada \(\lambda_x^{\rm or}\) es entonces natural y no cambia
al añadir bloques contractibles.

En precisión finita,
\begin{equation}\label{eq:detint-smith-finito}
 S_{x,n}=S_x\bmod3^n,
 \qquad
 a_i^{(n)}=\min(a_i,n).
\end{equation}
El terminal es, por tanto, la familia compatible de todas las precisiones y
no un entero aislado.

\section{Torre de Bockstein y acarreo de orden superior}

Sea
\[
 C_x^1\xrightarrow{S_x}C_x^0
\]
el bloque terminal. Si una clase \([\bar y]\) admite un lift
\(y\in C_x^1\) con \(S_xy\in3^qC_x^0\), definimos
\begin{equation}\label{eq:detint-beta-q}
 \beta_q[\bar y]
 :=\left[3^{-q}S_xy\bmod3\right].
\end{equation}
El cambio de lift altera el representante por una frontera de la página
correspondiente, de modo que \(\beta_q\) está bien definido. El carry
superior es
\begin{equation}\label{eq:detint-carry-superior}
 \operatorname{Carr}_q(y):=3^{-q}S_xy.
\end{equation}
No sustituye el carry TPK de \eqref{eq:detint-c-v-lifts}; registra su
prolongación en la torre de coeficientes.

Para un bloque \(3^{a_i}u_i\), la clase persiste hasta la página \(a_i\),
y en esa página el diferencial no nulo es multiplicación por
\(\bar u_i\in\mathbf F_3^\times\). Por consiguiente,
\begin{equation}\label{eq:detint-bock-longitudes}
 \{\text{longitudes de supervivencia}\}=\{a_1,\ldots,a_r\},
 \qquad
 d_x=\sum_i a_i.
\end{equation}
El producto orientado de las trivializaciones de página satisface
\begin{equation}\label{eq:detint-bock-leading}
 \Theta_{\rm Bock}(x)=\overline{\lambda_x^{\rm or}}
 \in\mathbf F_3^\times.
\end{equation}
La prueba se hace en cada bloque diagonal y se transporta por los cambios
unimodulares registrados. Smith, Bockstein y unidad inicial son así tres
lecturas del mismo terminal.

\iffalse
\section{Las trece coordenadas de la restricción determinantal}

Para cada \((m,n)\), se define
\begin{equation}\label{eq:detint-d-trece}
 D_{{\rm det},m,n}(x):=
 (\mathcal F,\operatorname{Ext}_\Gamma,\operatorname{Rel},
 \mathcal N,\operatorname{or},\operatorname{Carr},\operatorname{Mem},
 \partial,\gamma,\operatorname{Msect},\operatorname{Surv},
 \operatorname{Term},\mathcal L)_{{\rm det},m,n}(x),
\end{equation}
con el contenido siguiente.
\begin{enumerate}
 \item La fibra es
 \[
  \mathcal F=(A_n,\mathsf T_m(x),
  \{P_n^0(c_y)\}_y,\{P^{\rm gen}_{y;m,n}\}_y).
 \]
 \item \(\operatorname{Ext}_\Gamma\) contiene todos los morfismos
 \(\alpha\), los mapas \(\Psi_\alpha\), las homotopías \(K_\alpha\) y
 los mapas \(R_{m',m}^{n',n}\).
 \item \(\operatorname{Rel}\) contiene \(\partial^2=0\), la ecuación de
 pullback, AW, counidad, \(\Psi_\beta\Psi_\alpha=\Psi_{\beta\alpha}\) y
 \(K_{\beta\alpha}=K_\beta+\Psi_\beta K_\alpha\).
 \item
 \(\mathcal N=(m,n,\widetilde c_y,\widetilde v_y,c_{y,n},v_{y,n},
 r_x,a_1,\ldots,a_{r_x},d_x)\).
 \item \(\operatorname{or}\) contiene \(y\mapsto\bar y\),
 \(\mathfrak o_G\), \(\Omega\) y la orientación de la línea
 determinantal.
 \item \(\operatorname{Carr}\) conserva separadamente el carry entero,
 su reducción \(c_{y,n}\) y todos los carries superiores
 \(\operatorname{Carr}_q\).
 \item \(\operatorname{Mem}\) contiene
 \[
  y_0\to\cdots\to y_m,
  \quad\{(\widetilde c_{y_j},\widetilde v_{y_j},\mathbf1_{y_j})\}_{j=0}^m,
  \quad\{K_{\alpha_j}\}.
 \]
 \item
 \(\partial=(\partial_G,\partial_P,\partial_{P^{\rm gen}},
 \partial_{\overline P},h_*)\).
 \item \(\gamma=[y_0\to\cdots\to y_q]\) conserva el símplice orientado
 completo.
 \item \(\operatorname{Msect}\) contiene todas las extensiones
 compatibles, todos los lifts \(3\)-ádicos y todas las presentaciones
 unimodulares del mismo terminal orientado; no elige una a posteriori.
 \item \(\operatorname{Surv}\) registra profundidad arbitraria,
 compatibilidad de coeficientes, balance, aciclicidad sobre \(K_3\),
 estabilidad por bloques unitarios y permanencia de la raíz.
 \item
 \[
  \operatorname{Term}=(\widehat P_x^{\rm red},S_x,
  \{a_i,u_i\},d_x,\lambda_x^{\rm or},
  \{\beta_q\},\Theta_{\rm Bock}).
 \]
 \item El lector interno es
 \[
  \mathcal L_{\rm det}^{\rm int}
  =(\mathbf1_y,
  \varpi_{\rm root}z_-=\varpi_{\rm root}z_+,
  [z_-]=[z_+],d_x,\lambda_x^{\rm or},\Theta_{\rm Bock}).
 \]
\end{enumerate}

Las transiciones satisfacen, coordenada por coordenada,
\begin{equation}\label{eq:detint-naturalidad-trece}
 u^{\rm det}_{(m',n'),(m,n)}D_{{\rm det},m',n'}(x')
 =D_{{\rm det},m,n}(\tau_{m',m}x').
\end{equation}

\section{Ensamblador, publicador y cadena no ablativa}

Se define
\begin{equation}\label{eq:detint-graph-d}
 \mathcal G_{\rm det}:=\operatorname{Graph}(D_{\rm det}),
 \qquad
 \operatorname{Res}_{\rm det}(x):=(x,D_{\rm det}x).
\end{equation}
Sea \(\chi_{\rm det}(x)\) la multisección completa de gérmenes
\[
 (y,\widetilde c_y,\widetilde v_y,\operatorname{or}_y,\gamma_y)
\]
compatibles a toda profundidad. El objeto dependiente es
\begin{equation}\label{eq:detint-x-dependiente}
 X_{\chi,{\rm det}}
 :=\{(\chi_{\rm det}(x),x,D_{\rm det}x):
 x\in\mathcal C_{\rm cont}^{\rm det}\},
\end{equation}
y
\begin{equation}\label{eq:detint-prx}
 \operatorname{pr}_{X,{\rm det}}(x,D_{\rm det}x)
 :=(\chi_{\rm det}(x),x,D_{\rm det}x).
\end{equation}

El ensamblador crudo
\begin{equation}\label{eq:detint-ensamblador-crudo}
 a_{\rm det}:X_{\chi,{\rm det}}\longrightarrow
 \mathscr M_{\rm det}^{\rm amb}
\end{equation}
produce
\begin{equation}\label{eq:detint-ensamblador-salida}
\begin{split}
 a_{\rm det}(z)=\bigl(&
 \{G_{m,n},\Delta_{\rm AW},\epsilon_G\}_{m,n};\\
 &\{P_n^0(c),\Psi_\alpha,K_\alpha,P^{\rm gen},
 \mathbf1,z_-,z_+,h_*\}_{m,n};\\
 &\{\widehat P^{\rm red},S,d,\lambda^{\rm or},
 \beta_q,\Theta_{\rm Bock}\}\bigr).
\end{split}
\end{equation}
El macroobjeto conserva su entrada:
\begin{equation}\label{eq:detint-graph-a}
 \mathfrak M_{\rm det}:=\operatorname{Graph}(a_{\rm det}),
 \qquad
 \mathcal A_{\rm det}(z):=(z,a_{\rm det}z).
\end{equation}

El publicador crudo
\begin{equation}\label{eq:detint-publicador-crudo}
 p_{\rm det}:\mathfrak M_{\rm det}\longrightarrow\mathscr Y_{\rm det}
\end{equation}
publica conjuntamente el certificado
\begin{equation}\label{eq:detint-certificado-publicado}
\begin{gathered}
 \partial^2=0,quad \text{AW y counidad},quad
 \mathbf1_y\text{ group-like compatible},\\
 \partial z_\pm=0,quad
 \varpi_{\rm root}z_-=\varpi_{\rm root}z_+,quad
 \partial h_*=z_--z_+,\\
 \Psi_\beta\Psi_\alpha=\Psi_{\beta\alpha},quad
 K_{\beta\alpha}=K_\beta+\Psi_\beta K_\alpha,\\
 d=v_3(\det S)=\sum_i a_i,quad
 \Theta_{\rm Bock}=\overline{\lambda^{\rm or}},quad
 \text{naturalidad completa bajo refinamiento}.
\end{gathered}
\end{equation}
Finalmente,
\begin{equation}\label{eq:detint-graph-p}
 \mathcal C_{\rm det}^{\rm HMT}:=\operatorname{Graph}(p_{\rm det}),
 \qquad
 \mathcal P_{\rm det}(M):=(M,p_{\rm det}M).
\end{equation}
La cadena íntegra es
\begin{equation}\label{eq:detint-cadena-completa}
 \mathcal C_{\rm cont}^{\rm disc}\supset
 \mathcal C_{\rm cont}^{\rm det}
 \xrightarrow{\operatorname{Res}_{\rm det}}\mathcal G_{\rm det}
 \xrightarrow{\operatorname{pr}_{X,{\rm det}}}X_{\chi,{\rm det}}
 \xrightarrow{\mathcal A_{\rm det}}\mathfrak M_{\rm det}
 \xrightarrow{\mathcal P_{\rm det}}\mathcal C_{\rm det}^{\rm HMT}.
\end{equation}
Cada flecha tiene una inversa sobre su imagen dada por una proyección de
coordenadas. No se sustituye historia por valor, multisección por selector,
complejo por determinante ni prueba por nombre.

\section{Teorema de realización determinantal interna}

\begin{theorem}[Realización determinantal interna completa]
\label{thm:detint-realizacion}
Para todo \(x\in\mathcal C_{\rm cont}^{\rm det}\), la cadena
\eqref{eq:detint-cadena-completa} construye un macroobjeto natural en
profundidad y precisión. Contiene dos ciclos con raíz común, un filler
explícito que los identifica, una unidad group-like compatible y un terminal
\(3\)-ádico cuyo orden, forma de Smith y torre de Bockstein producen el mismo
lector orientado. Las trece coordenadas de \eqref{eq:detint-d-trece}
permanecen recuperables en la salida.
\end{theorem}

\begin{proof}
Las ecuaciones \eqref{eq:detint-p0-diferencial}--
\eqref{eq:detint-d2} prueban que \(P_n^0(c)\) es un complejo. Las dos
aumentaciones son mapas de complejos, por lo que el pullback está tipado.
Alexander--Whitney es natural y hace group-like a cada vértice; junto con
\(r\), esto da la unidad común. El cálculo directo
\eqref{eq:detint-zpm-ciclos}--\eqref{eq:detint-filler-identidad} prueba que
las publicaciones son ciclos, tienen la misma raíz y están unidas por
\(h_*\). Las fórmulas de \(\Psi_\alpha\) prueban functorialidad, y
\(K_\alpha\) prueba la naturalidad homotópica conservando el incremento de
memoria. La definición de \(\mathcal C_{\rm cont}^{\rm det}\) garantiza que
la reducción terminal produce un bloque cuadrado sin postularlo fuera del
dominio. Smith sobre \(\mathbf Z_3\) da \(d=\sum_i a_i\); el cálculo de cada
bloque \(3^{a_i}u_i\) da
\(\Theta_{\rm Bock}=\overline{\lambda^{\rm or}}\). Truncación, reducción y
adición de bloques unitarios conservan estas identidades. Por último, las
tres construcciones por graph retienen la entrada como primera coordenada,
de modo que la composición completa es no ablativa.
\end{proof}

\section{Procedencia y falsadores}

\paragraph{Procedencia.}
El nervio TPK, Alexander--Whitney, la counidad, la unidad común, la forma
normal \(\partial f=3ce+b\), \(\partial e=g\),
\(\partial b=-3cg\), la raíz común, el filler y la arquitectura de torre
tienen estatuto \texttt{ARQUITECTURA\_AUTORAL\_PREEXISTENTE}. El cálculo
literal de los dos ciclos, su raíz y la identidad
\(\partial h_*=z_--z_+\) tiene estatuto \texttt{RESULTADO\_RECUPERADO}. El
sistema local \(\Psi_\alpha/K_\alpha\), el subdominio superviviente
balanceado, el terminal exclusivamente \(3\)-ádico y el empaquetado por
graphs tienen estatuto \texttt{FORMALIZACION\_NUEVA}; formalizan sin cambiar
la procedencia conceptual del aparato.

\begin{ablation}[Falsadores internos de la rama determinantal]
\label{abl:detint-falsadores}
La construcción queda falsada en el nivel correspondiente si ocurre alguno
de los hechos siguientes.
\begin{enumerate}
 \item Falla \(\partial^2=0\); entonces no existe el complejo local.
 \item Alexander--Whitney no conmuta con refinamiento, la counidad falla o
 el vértice deja de ser group-like.
 \item \(z_\pm\) no son ciclos, sus raíces difieren o falla
 \(\partial h_*=z_--z_+\).
 \item \(\Psi_\alpha\) no es mapa de complejos o falla
 \(K_{\beta\alpha}=K_\beta+\Psi_\beta K_\alpha\); en ese caso se ha
 perdido carry o memoria.
 \item El complejo reducido no es balanceado o no es acíclico sobre
 \(K_3\); la historia queda fuera del dominio terminal tipado.
 \item Las longitudes de Bockstein no coinciden con los \(a_i\), o
 \(\Theta_{\rm Bock}\ne\overline{\lambda^{\rm or}}\).
 \item El refinamiento modifica \(d\) o \(\lambda^{\rm or}\) fuera de
 bloques unitarios y cambios de orientación registrados.
 \item Se omite una de las trece coordenadas: \emph{objeto sustituido;
 conclusión no transferible}.
\end{enumerate}
Estos falsadores comprueban el objeto interno construido; no lo reemplazan
por un control posterior.
\end{ablation}
\fi
 \endgroup

\section{Naturalidad del ensamblaje y compatibilidad bajo refinamiento}

Antes de formar el registro de horizonte, hacemos explícita la acción total
de las cinco operaciones aquí desarrolladas sobre \emph{la misma} arista
\(\mathfrak e_\alpha\). Si \(P_x,Q_x\) denotan las restricciones de
\eqref{eq:pdfv2-cor-sigma} a la fibra de \(x=s(\alpha)\), y análogamente en
\(y=t(\alpha)\), ponemos primero
\begin{equation}
\label{eq:pdfv2-cor-a-eta-emision}
 A_\alpha:=P_yU_\alpha P_x+Q_yU_\alpha Q_x,
 \qquad
 \eta_\alpha:=Q_yU_\alpha P_x+P_yU_\alpha Q_x.
\end{equation}
Las cinco lecturas son familias dependientes indexadas por la propia
\(\mathfrak e_\alpha\). El índice fija fuente, destino, ruta y ledger, pero
no se vuelve a copiar dentro del valor. Sus tipos son los registros
dependientes
\(\operatorname{LedgerEntry}:=\operatorname{cod}(\lambda)\) y
\begin{align*}
 \mathsf Y_{\rm sp}[\alpha]&:=
 \operatorname{Rec}[U:\mathsf H_x\to\mathsf H_y;
 \sigma_x\in\operatorname{End}\mathsf H_x;
 \sigma_y\in\operatorname{End}\mathsf H_y;A,\eta:\mathsf H_x\to\mathsf H_y],\\
 \mathsf Y_{\rm sol}[\alpha]&:=
 \operatorname{Rec}[b,b^+,b^-\in\operatorname{LedgerEntry};
 \operatorname{upd}:\mathsf{Mem}_x\to\mathsf{Mem}_y],\\
 \mathsf Y_{\rm gau}[\alpha]&:=
 \mathcal P_{\rm fin}\operatorname{Rec}[
 \operatorname{Cell}(\alpha);q\in Q_{\rm inc}(y);B;W_c(q)],\\
 \mathsf Y_{\rm coh}[\alpha]&:=
 \operatorname{Rec}[I_y,J_y,\kappa_{y,N},\delta_y,C_{y,N};
 F_\alpha:C_{y,N}\to C_{x,N}],\\
 \mathsf Y_{\rm det}[\alpha]&:=
 \operatorname{Rec}[\Psi_\alpha,K_\alpha,\mathbf1_\alpha,z_\pm,h_*;
 \operatorname{Det},H_*,\operatorname{Smith},\operatorname{Boc}].
\end{align*}
Por tanto el juicio común es
\[
 \mathfrak e_\alpha:\operatorname{Ext}_k\ \vdash\
 \mathcal O_\bullet(\mathfrak e_\alpha)
 \in\mathsf Y_\bullet[\alpha],
 \qquad
 \bullet\in\{\mathrm{sp,sol,gau,coh,det}\}.
\]
\begin{align}
 \mathcal O_{\rm sp}(\mathfrak e_\alpha)
 &:=
 \bigl(U_\alpha,\sigma_x,\sigma_y,
 A_\alpha,\eta_\alpha\bigr),
 \label{eq:pdfv2-cor-accion-sp}\\
 \mathcal O_{\rm sol}(\mathfrak e_\alpha)
 &:=
 \bigl(b_\alpha,b_\alpha^+,b_\alpha^-;\,
 (M^+,M^-,D,H)\mapsto
 (M^++b_\alpha^+,M^-+b_\alpha^-,
  D+b_\alpha,H+|b_\alpha|)\bigr),
 \label{eq:pdfv2-cor-accion-sol}\\
 \mathcal O_{\rm gau}(\mathfrak e_\alpha)
 &:=
 \left\{
 \bigl(\operatorname{Cell}_{\rm TPK}(\alpha),
 q,B,W_c(q)\bigr):
 q\in Q_{\rm inc}(y)
 \right\},
 \label{eq:pdfv2-cor-accion-gau}\\
 \mathcal O_{\rm coh}(\mathfrak e_\alpha)
 &:=
 \bigl(I_y,J_y,\kappa_{y,N},\delta_y,C_{y,N},
 F_\alpha:C_{y,N}\longrightarrow C_{x,N}\bigr),
 \label{eq:pdfv2-cor-accion-coh}\\
 \mathcal O_{\rm det}(\mathfrak e_\alpha)
 &:=
 \bigl(\Psi_\alpha,K_\alpha,\mathbf1_\alpha,z_\pm(\alpha),
 h_*(\alpha),\operatorname{Det}(C_{y,N}),H_*(C_{y,N}),
 \operatorname{Smith}(C_{y,N}),\operatorname{Boc}(C_{y,N})\bigr).
 \label{eq:pdfv2-cor-accion-det}
\end{align}

La lectura arquimediana interna es la primera componente de esta acción
simultánea; no introduce todavía ninguna forma clásica:
\begin{equation}
\label{pII-full:eq:continuo-lector-arquimediano}
 \mathcal L_{\rm arq}^{\rm int}(\mathfrak e_\alpha)
 :=\mathcal O_{\rm sp}(\mathfrak e_\alpha).
\end{equation}
Para cualquier valor de esas familias, incluido cada elemento de la
multisección \(\mathcal O_{\rm gau}\), su soporte es el índice dependiente
\begin{equation}
\label{eq:pdfv2-cor-soporte-dependiente}
 \begin{aligned}
 \mathsf B_\alpha&:=
 \{(s(\alpha),t(\alpha),\gamma\alpha,
          \operatorname{Led}(\gamma\alpha))\},\\
 \operatorname{supp}_\alpha:\mathsf Y_\bullet[\alpha]&\longrightarrow
 \mathsf B_\alpha,\qquad
 o\longmapsto(s(\alpha),t(\alpha),\gamma\alpha,
          \operatorname{Led}(\gamma\alpha)).
 \end{aligned}
\end{equation}
Por tanto quedan tipados, sin una proyección que recupere la arista,
\begin{equation}
\label{eq:pdfv2-cor-soporte-cinco-lecturas}
 \begin{gathered}
 s_\alpha(\mathcal O_\bullet):=s(\alpha),\qquad
 t_\alpha(\mathcal O_\bullet):=t(\alpha),\\
 \gamma_\alpha(\mathcal O_\bullet):=\gamma\alpha,
 \qquad
 \operatorname{Led}_\alpha(\mathcal O_\bullet)
 :=\operatorname{Led}(\gamma\alpha).
 \end{gathered}
\end{equation}
Aquí \(\operatorname{Cell}_{\rm TPK}(\alpha)\) es el mapa de complejos
\eqref{eq:pdfv2-cor-cell-accion-emision} inducido por el transporte de la
misma fibra de hojas, y \(F_\alpha\) es
\eqref{eq:pdfv2-cor-refinamiento-celda} para esa extensión. En particular,
\(\mathcal O_{\rm gau}\) conserva la multisección completa de \(q\)'s; no
elige el origen ni una órbita.

Las cinco fórmulas anteriores constituyen una sola acción dependiente
\begin{equation}
\label{eq:pdfv2-cor-accion-unica-dependiente}
 \mathcal O(\mathfrak e_\alpha)
 :=\left\langle
 \mathcal O_{\rm sp},\mathcal O_{\rm sol},\mathcal O_{\rm gau},
 \mathcal O_{\rm coh},\mathcal O_{\rm det};
 \operatorname{Cpl}_\alpha
 \right\rangle,
\end{equation}
donde \(\operatorname{Cpl}_\alpha\) es el certificado formado por las
identidades
\begin{equation}
\label{eq:pdfv2-cor-acoplamientos-unicos}
 \begin{gathered}
 \operatorname{supp}_\alpha(\mathcal O_{\rm sp})
 =\operatorname{supp}_\alpha(\mathcal O_{\rm sol})
 =\operatorname{supp}_\alpha(\mathcal O_{\rm gau})
 =\operatorname{supp}_\alpha(\mathcal O_{\rm coh})
 =\operatorname{supp}_\alpha(\mathcal O_{\rm det}),\\
 \mathsf Q(t(\alpha))=\sigma_9\mathsf Q(s(\alpha)),\quad
 g(t(\alpha))=g(s(\alpha))+[1]_9,\quad
 q^{(9)}(t(\alpha))=q^{(9)}(s(\alpha))+\chi_9(g(s(\alpha))),\\
 A_\alpha^*A_\alpha+\eta_\alpha^*\eta_\alpha=U_\alpha^*U_\alpha,\qquad
 (M^+,M^-,D,H)'=(M^++b_\alpha^+,M^-+b_\alpha^-,
 D+b_\alpha,H+|b_\alpha|),\\
 Q_{\rm inc}=C^2(\operatorname{Cell}_{\rm TPK};\mathbf F_3),\qquad
 sB=0,\qquad W_c(q+Ba)=W_c(q),\\
 C_{y,N}\text{ es el mismo complejo en }
 \mathcal O_{\rm coh}\text{ y }\mathcal O_{\rm det},\qquad
 \partial h_*=z_--z_+,\qquad
 K_{\beta\alpha}=K_\beta+\Psi_\beta K_\alpha.
 \end{gathered}
\end{equation}
\Needspace{18\baselineskip}
El símbolo \(\bullet\) recorre las cinco lecturas de
\eqref{eq:pdfv2-cor-accion-unica-dependiente}; todas conservan la misma
fuente, destino, ruta y ledger.

Reunimos, sin separar el estado, las operaciones construidas hasta aquí:
\begin{equation}
\label{eq:pdfv2-cor-constructor-unico}
 \begin{aligned}
 \mathfrak D_{k,N}(\widehat x)
 &:=\bigl(\{\mathfrak e_\alpha\}_\alpha;\\
 &\quad U_k,\sigma_k,P_k,Q_k,A_k,\eta_k;\\
 &\quad b_\alpha,b_\alpha^\pm,M_k^\pm,D_k^0,H_k^0,
 \kappa_k^{\rm can};\\
 &\quad V_{\rm TPK},\mathscr I,\mathscr E,\mathscr B,B;\\
 &\quad \operatorname{Cell}_{\rm TPK}(\widehat x),Q_{\rm inc}(\widehat x),W_c;\\
 &\quad I_k,J_k,\kappa_{k,N},\delta,C_{k,N};\\
 &\quad c_\alpha,v_\alpha,\Psi_\alpha,K_\alpha,
 \mathbf1_\alpha,z_\pm(\alpha),h_*(\alpha);\\
 &\quad
 \operatorname{Det}(C_{k,N}),H_*(C_{k,N}),
 \operatorname{Smith}(C_{k,N}),\operatorname{Boc}(C_{k,N});\\
 &\quad \{\mathcal O(\mathfrak e_\alpha)\}_\alpha
 \bigr).
 \end{aligned}
\end{equation}
El punto y coma sólo separa niveles de lectura de una misma tupla dependiente;
no define estructuras independientes. El estado de entrada no se guarda como
primera coordenada ni se recupera mediante una proyección. Cada operación se
demuestra por su acción en los generadores \(\mathfrak e_\alpha\), que
conservan el mismo soporte, la misma ruta y el mismo ledger.

Para \(\ell\geq k\) y \(N'\geq N\), el transporte se define componente a
componente por truncación de historias, reducción de coeficientes y
precomposición de puertos. La afirmación de naturalidad que debe probarse es
una sola:
\begin{equation}
\label{eq:pdfv2-cor-naturalidad-unica}
 u_{(\ell,N'),(k,N)}\mathfrak D_{\ell,N'}(\widehat x_\ell)
 =\mathfrak D_{k,N}(\tau_{\ell,k}\widehat x_\ell).
\end{equation}
La naturalidad se formula sobre la historia generada, no fingiendo que una
arista fina aislada deba seguir siendo arista después de truncar sus dos
extremos. Para
\(\gamma=(\varepsilon_1,\ldots,\varepsilon_\ell)\), definimos
\[
 \mathsf E_{\leq\ell}(\gamma)
 :=(\mathfrak e_{\varepsilon_i})_{1\leq i\leq\ell},
 \qquad
 \mathcal O_{\leq\ell}(\gamma)
 :=(\mathcal O(\mathfrak e_{\varepsilon_i}))_{1\leq i\leq\ell}.
\]
Los dos transportes son los mapas de prefijo explícitos
\begin{align}
 u^{\mathfrak e}_{\ell,k}:
 \mathsf E_{\leq\ell}(\gamma)&\longrightarrow
 \mathsf E_{\leq k}(\gamma),&
 (\mathfrak e_{\varepsilon_i})_{i\leq\ell}
 &\longmapsto(\mathfrak e_{\varepsilon_i})_{i\leq k},
 \label{eq:pdfv2-cor-transporte-arista}\\
 u^{\mathcal O}_{\ell,k}:
 \mathcal O_{\leq\ell}(\gamma)&\longrightarrow
 \mathcal O_{\leq k}(\gamma),&
 (\mathcal O(\mathfrak e_{\varepsilon_i}))_{i\leq\ell}
 &\longmapsto
 (\mathcal O(\mathfrak e_{\varepsilon_i}))_{i\leq k}.
 \label{eq:pdfv2-cor-transporte-accion}
\end{align}
En cada entrada conservada, la componente \(\mathrm{sp}\) entrelaza
\(U,\sigma,P,Q,A,\eta\); la componente \(\mathrm{sol}\) aplica la esperanza
condicional a \(b,b^\pm,\kappa\); la componente \(\mathrm{gau}\) aplica
\(\operatorname{Cell}(\tau)\oplus\tau_Q\); la componente \(\mathrm{coh}\)
precompone los puertos con \(F\); y la componente \(\mathrm{det}\) aplica
\(\Psi,K,H_*,\operatorname{Smith},\operatorname{Boc}\) al mismo mapa de
complejos. El prefijo actúa simultáneamente sobre esas cinco componentes.
Sus componentes ya descargadas son
\begin{equation}
\label{eq:pdfv2-cor-naturalidad-componentes}
 \begin{gathered}
 U_{m,k}=U_{m,\ell}U_{\ell,k}\quad(m\geq\ell\geq k),\qquad
 E_kb_f^\pm=b_c^\pm+\kappa_k^{\rm can},\\
 r_{N',N}\kappa_{k,N'}=
 \kappa_{k,N}(r_{N',N}\oplus r_{N',N}),\\
 \Psi_\beta\Psi_\alpha=\Psi_{\beta\alpha},
 \qquad K_{\beta\alpha}=K_\beta+\Psi_\beta K_\alpha,\\
 u^{\mathcal O}_{\ell,k}\mathcal O_{\leq\ell}(\gamma)
 =\mathcal O_{\leq k}(\gamma)
 =\operatorname{map}(\mathcal O)\bigl(u^{\mathfrak e}_{\ell,k}
                    \mathsf E_{\leq\ell}(\gamma)\bigr).
 \end{gathered}
\end{equation}
La primera igualdad usa la multiplicación de los pesos condicionales sobre el mismo
sistema de cilindros comunes; las restantes son las fórmulas explícitas
anteriores. Ninguna proyección visible
puede usarse para reconstruir o seleccionar \(\widehat x\).

Para completar el último componente, denotemos por
\(F_{(\ell,N'),(k,N)}:C_{\ell,N'}\to C_{k,N}\) el mapa de complejos de
\eqref{eq:pdfv2-cor-refinamiento-celda}. Su cono se conserva como parte de la
transición. La functorialidad de homología y de la sucesión de Bockstein da
\begin{equation}
\label{eq:pdfv2-cor-naturalidad-homologia-bockstein}
 H_*(F_{m,k})=H_*(F_{\ell,k})H_*(F_{m,\ell}),
 \qquad
 \operatorname{Boc}(F_{m,k})
 =\operatorname{Boc}(F_{\ell,k})
  \operatorname{Boc}(F_{m,\ell}).
\end{equation}
La línea determinante se transporta mediante la identidad del triángulo
exacto
\begin{equation}
\label{eq:pdfv2-cor-naturalidad-determinante}
 \operatorname{Det}(C_{k,N})
 \simeq
 \operatorname{Det}(C_{\ell,N'})\otimes
 \operatorname{Det}\!\left(\operatorname{Cone}
 F_{(\ell,N'),(k,N)}\right).
\end{equation}
Esta fórmula registra también los refinamientos que añaden homología. La
lectura de Smith se conserva sin fingir invariancia: a cada transición se le
asigna el certificado total
\begin{equation}
\label{eq:pdfv2-cor-transicion-smith}
 \mathcal S(F):=
 \bigl(\operatorname{Smith}(C_{\ell,N'}),
       \operatorname{Smith}(C_{k,N}),
       \operatorname{Smith}(\operatorname{Cone}F)\bigr),
\end{equation}
que determina literalmente qué sumandos se conservan, nacen o mueren.

La acción de una simetría coherente \(g\) en el árbol común induce
isometrías de permutación \(\widetilde g_k\) en
\(\mathscr H_k\). La biyección de hijos y la conservación de los pesos de
\eqref{eq:continuo-comun-naturalidad-medida} prueban en los vectores básicos
\begin{equation}
\label{eq:pdfv2-cor-naturalidad-transporte-hojas}
 \widetilde g_{k+1}U_k=U_k\widetilde g_k,
 \qquad
 \widetilde g_k\sigma_k=\sigma_k\widetilde g_k,
 \qquad
 \widetilde g_kP_k=P_k\widetilde g_k,
 \qquad
 \widetilde g_kQ_k=Q_k\widetilde g_k.
\end{equation}

\begin{theorem}[Totalidad y naturalidad del constructor correlativo]
\label{thm:pdfv2-cor-totalidad-naturalidad}
Para todo estado total de \eqref{eq:continuo-comun-estado}, todo horizonte
finito y todo nivel de coeficientes, la expresión
\eqref{eq:pdfv2-cor-constructor-unico} está definida. Sus cinco
características son operaciones simultáneas sobre el mismo registro de
extensión y satisfacen la única identidad
\eqref{eq:pdfv2-cor-naturalidad-unica}. El collar, el complejo y la línea
determinante permanecen no vacíos aun cuando sus lectores posteriores no
sean acíclicos ni invertibles.
\end{theorem}

\begin{proof}
La no vacuidad del árbol y de sus puertos es
\cref{prop:continuo-comun-arbol-finito-no-vacio}; la medida y la expectativa
son totales por \cref{thm:continuo-comun-medida-canonica}. Las fórmulas
  \eqref{eq:pdfv2-cor-u-comun}--\eqref{eq:pdfv2-cor-memorias} definen transporte,
polaridad, balance y memoria para todos los generadores. Las cuatro
  direcciones de \eqref{eq:pdfv2-cor-cuatro-rectas-tpk} son distintas por
  \eqref{eq:pdfv2-cor-proyectores-no-colapso}; la completación celular
  \cref{def:pdfv2-cor-completacion-celular-tpk} contiene sus dos órdenes de
  camino y seis celdas, mientras
  \eqref{eq:pdfv2-cor-collar-naturalidad} prueba su naturalidad. La
sucesión celular es exacta por \eqref{eq:pdfv2-cor-sucesion-celda}, el
complejo local cumple \(\partial^2=0\) y el relleno correlativo es explícito por
\eqref{eq:pdfv2-cor-filler}. Toda construcción finita basada posee la línea
\eqref{eq:pdfv2-cor-linea-determinante-total}; homología, Bockstein y el
certificado de Smith se transportan por
\eqref{eq:pdfv2-cor-naturalidad-homologia-bockstein}--
\eqref{eq:pdfv2-cor-transicion-smith}. Finalmente, las identidades se
verifican en cada generador \(\mathfrak e_\alpha\) mediante
\eqref{eq:pdfv2-cor-naturalidad-componentes},
\eqref{eq:pdfv2-cor-collar-naturalidad} y
\eqref{eq:pdfv2-cor-naturalidad-transporte-hojas}; por linealidad y por la
composición functorial de rutas se obtiene
\eqref{eq:pdfv2-cor-naturalidad-unica} para el registro entero.
\end{proof}

\section{Separación causal de las identidades, deducciones y realizaciones del constructor correlativo}

\paragraph{No vacuidad e identidades exactas del constructor correlativo}
La no vacuidad importada de \eqref{eq:continuo-comun-out-finito} y las
identidades
\eqref{eq:pdfv2-cor-b-identidades},
\eqref{eq:pdfv2-cor-sucesion-celda},
\eqref{eq:pdfv2-cor-sb-cero},
\eqref{eq:pdfv2-cor-u-gram},
\eqref{eq:pdfv2-cor-balance-hojas},
\eqref{eq:pdfv2-cor-jensen-hojas},
\eqref{eq:pdfv2-cor-dos-cero},
\eqref{eq:pdfv2-cor-cociclo-matricial} y
\eqref{eq:pdfv2-cor-filler} se deducen de las definiciones y de las relaciones
APP--TRIT--TPK exhibidas. Su estatuto es
\texttt{EXACTO\_INTERNO}.

\paragraph{Hipótesis de partida para la isometría de \(U_k\), la sucesión celular y la composición \(\Psi/K\)}
La construcción de \(U_k\) usa como estructura de partida la finitud de las
fibras, la supervivencia a todo horizonte, el truncamiento funcional y los
transportes de las dos hojas ya producidos por TPK\@. La igualdad
\(U_k^*U_k=I\) exige además verificar \(U_\alpha^*U_\alpha=I\) para cada
emisión; sin esa verificación permanece la identidad de Gram
\eqref{eq:pdfv2-cor-u-gram}, que sí es incondicional. La sucesión celular usa
los dos puertos APP ordenados. La composición \(\Psi/K\) usa carry entero,
orientación y frontera. Bajo esas premisas enumeradas, las conclusiones
indicadas son afirmativas.

\paragraph{Datos adicionales para realizaciones isométricas, retractos balanceados y trivialización de la línea determinante}
El constructor interno queda completo con el Gram positivo de
\eqref{eq:pdfv2-cor-u-gram}, el espacio total \(Q_{\rm inc}\), la sucesión
celular y la línea determinante con su cono de transición. Una publicación
especializada puede exigir datos adicionales: la isometría de cada
\(U_\alpha\), una identificación numérica concreta del defecto APP, un
retracto balanceado y acíclico o la trivialización de una línea
determinante. Esos datos son entrelazadores del codominio publicado; no
forman parte del dominio ni de la existencia de
\(\mathfrak D_{k,N}\). Su ausencia impide exclusivamente la promoción que
los mencione.

\begin{criterion}[Falsadores de la correlación intrínseca]
\label{crit:pdfv2-cor-falsadores}
La construcción conjunta falla en el primer nivel en que ocurra cualquiera de
los hechos siguientes:
\begin{enumerate}
 \item una operación usa una extensión, orientación, carry, memoria, frontera
 o ruta diferente de la registrada en \(\mathfrak e_\alpha\);
 \item se escoge un hijo de
 \(\operatorname{Ch}_k(\widehat x_k)\) y se eliminan los restantes;
 \item el peso de \eqref{eq:pdfv2-cor-peso-conteo} no suma uno, falla
 \eqref{eq:pdfv2-cor-u-gram}, o se declara \(U_k^*U_k=I\) sin probar
 \(U_\alpha^*U_\alpha=I\) para toda emisión;
 \item se identifica \(\kappa_k^{\rm can}\) con el carry entero
 \(\kappa_\alpha\);
 \item falla \(sB=0\), se contrae
 \(Q_{\rm inc}(\widehat x)\) a su origen, se altera la acción
 \(q\mapsto q+Ba\), o se identifican los dos órdenes de camino antes de
 aplicar el lector;
 \item un refinamiento rompe
 \eqref{eq:pdfv2-cor-refinamiento-cadena} o el cociclo
 \eqref{eq:pdfv2-cor-cociclo-matricial};
 \item falla \(\partial^2=0\), \(z_\pm\) no son ciclos o
 \(\partial h_*\neq z_--z_+\);
 \item se afirma un escalar determinantal no nulo sin el retracto, el
 balance y la aciclicidad exigidos por esa publicación;
 \item se presenta cualquiera de las cinco operaciones como un objeto
 separado del constructor \(\mathfrak D_{k,N}\).
\end{enumerate}
En todos esos casos el diagnóstico es: \emph{objeto común sustituido;
conclusión no transferible}.
\end{criterion}
 \end{HMTScientificOwnerSubordinated}
 \label{chap:83-05}

\begin{HMTScientificOwnerSubordinated}
\chapter[Terminal, naturalidad y l\'imite]{Terminal multivaluado, naturalidad \'unica y paso al l\'imite}
\label{chap:pdfv2-terminal-naturalidad-limite}

\begin{thesisbox}
El continuo discreto produce un solo sistema de historias, extensiones,
memorias y terminales. Los discriminantes internos
\(\mathrm{sp},\mathrm{sol},\mathrm{gau},\mathrm{coh},\mathrm{det}\)
nombran cinco caracter\'isticas simult\'aneas de ese mismo sistema: no fijan
entradas distintas ni autorizan a escoger historias diferentes. En este
cap\'itulo el terminal conserva todas las continuaciones, la conexi\'on
non\'adica permanece como correspondencia y la naturalidad se expresa por una
sola identidad.
\end{thesisbox}

\section{Historias supervivientes con testigos de extensi\'on}

Sea
\(\widehat x_k\in\widehat{\mathcal X}^{\rm tot}_k\) el estado total de
\eqref{eq:continuo-comun-dominio-total}. Para \(N\geq k\), definimos sin
otro predicado de pertenencia
\begin{equation}
\label{eq:pdfv2-terminal-historias-finita}
 \mathscr H_{k,N}(\widehat x_k)
 :=\operatorname{Term}_{k,N}(\widehat x_k).
\end{equation}
Cada
\(h=(\varepsilon_{k+1},\ldots,\varepsilon_N)\) de la derecha contiene por
\eqref{eq:continuo-comun-vertice-arbol} sus estados intermedios, los testigos
de \(\operatorname{Ext}\), los carries, las fronteras, las rutas y el ledger.
Cada arista satisface la acción de fase
\eqref{eq:continuo-comun-accion-fase-tpk}; por ello todo subcamino de nueve
pasos determina, con sus nueve emisiones y no sólo con su longitud, un
elemento de \(\mathsf Z^\Gamma_{9,j}\) y un testigo de
\(\Gamma_{9,j}\) por
\eqref{eq:continuo-comun-gamma9-apice}--
\eqref{eq:continuo-comun-gamma9-relacion}. Por tanto, Ext y la holonomía
non\'adica nacen de la misma palabra fase--compatible; no se añade a
\eqref{eq:pdfv2-terminal-historias-finita} un filtro de supervivencia ni una
segunda condici\'on \(\Gamma_9\).

Para \(k\leq\ell\leq N\), sea
\begin{equation}
\label{eq:pdfv2-terminal-prefijos-intermedios}
 \mathscr P_{\ell\mid k}(\widehat x_k)
 :=\{h_{[k,\ell]}:h\in\mathscr H_{k,N}(\widehat x_k)
       \text{ para alg\'un }N\geq\ell\}.
\end{equation}
Para \(p\in\mathscr P_{\ell\mid k}(\widehat x_k)\), usamos la abreviatura
\begin{equation}
\label{eq:pdfv2-terminal-continuaciones-prefijo}
 \mathscr H_{\ell,N}(p)
 :=\{q:p\star q\in\mathscr H_{k,N}(\widehat x_k)\}.
\end{equation}
La prolongaci\'on neutra de \eqref{eq:continuo-comun-seccion-total} extiende
cada prefijo a cualquier horizonte. En consecuencia,
\begin{equation}
\label{eq:pdfv2-terminal-no-vacuidad-historias}
 \mathscr H_{k,N}(\widehat x_k)\neq\varnothing,
 \qquad
 \mathscr P_{\ell\mid k}(\widehat x_k)\neq\varnothing,
\end{equation}
por \cref{prop:continuo-comun-arbol-finito-no-vacio}.

El truncamiento
\begin{equation}
\label{eq:pdfv2-terminal-truncamiento-historias}
 \pi_{N+1,N}^{\mathscr H}:
 \mathscr H_{k,N+1}(\widehat x_k)
 \longrightarrow\mathscr H_{k,N}(\widehat x_k)
\end{equation}
es exactamente \eqref{eq:continuo-comun-proyeccion-terminal} y es
sobreyectivo por \cref{lem:continuo-comun-terminal-sobreyectivo}. La
supervivencia es, por
\eqref{eq:continuo-comun-surv-igual-total}, una salida demostrada del dominio
total y no una selecci\'on previa.

\section{Registro terminal generado y sus \(13\) coordenadas}

Fijemos una historia
\[
 h=(\varepsilon_{j+1}:x_j\longrightarrow x_{j+1})_{j=k}^{N-1},
 \qquad
 \gamma_{k,k}:=I_{x_k},\qquad
 \gamma_{k,j+1}:=\gamma_{k,j}\star\varepsilon_{j+1},
\]
La notación operatorial
\(U_{\gamma_{k,j+1}}
=U_{\varepsilon_{j+1}}U_{\gamma_{k,j}}\) conserva el orden de aplicación,
mientras \(\star\) conserva el orden cronológico de la palabra. Fijamos
además un nivel de coeficientes \(r\geq1\). Para abreviar, escribimos
\[
 C_{x_j,r}:=C_{j,r}(x_j),\qquad
 F_{\varepsilon_{j+1}}:C_{x_{j+1},r}\longrightarrow C_{x_j,r}
\]
para el complejo y el mapa de complejos de
\eqref{eq:pdfv2-cor-cono-celda} y
\eqref{eq:pdfv2-cor-refinamiento-celda}. El tipo del registro de nivel
\(j\), denotado \(\mathfrak R_{j,r}\), es un registro dependiente con los
trece campos siguientes. En cada campo indicamos el tipo, la semilla y la
acción de una arista \(\varepsilon:=\varepsilon_{j+1}\).

El estado \(x_k\) ya contiene la palabra completa
\(\gamma_{0,k}=\varepsilon_1\star\cdots\star\varepsilon_k\). Por ello una
``semilla en \(k\)'' nunca reinicia la genealogía: es el fold de las mismas
actualizaciones desde la raíz hasta \(k\). Denotamos
\begin{equation}
\label{eq:pdfv2-terminal-prefijo-heredado}
 \begin{gathered}
 M_k^\pm:=\sum_{i=1}^{k}b_{\varepsilon_i}^\pm,qquad
 D_k:=\sum_{i=1}^{k}b_{\varepsilon_i},qquad
 H_k:=\sum_{i=1}^{k}|b_{\varepsilon_i}|,\\
 \mathbf c_{\leq k}:=
  \bigl(c_r(\varepsilon_i,\gamma_{0,i-1})\bigr)_{1\leq i\leq k},
 \qquad
 \mathbf k_{\leq k}:=
  \bigl(\kappa_{i-1}^{\rm can}\bigr)_{1\leq i\leq k},\\
 \mathbf F_{\leq k}:=
  \bigl(F_{\varepsilon_i}:C_{x_i,r}\to C_{x_{i-1},r}\bigr)_{1\leq i\leq k}.
 \end{gathered}
\end{equation}
Todas las listas vacías que aparecen abajo se entienden literalmente sólo
en \(k=0\); para \(k>0\) se sustituyen por los prefijos heredados de
\eqref{eq:pdfv2-terminal-prefijo-heredado} y del ledger de
\(\gamma_{0,k}\). Así el registro no conserva una copia opaca de \(x_k\),
pero tampoco pierde su pasado APP--TRIT--TPK.

\paragraph{Fibra \(\mathsf F\): dominio local y conservación bajo prolongación}
La fibra elemental es
\begin{equation}
\label{eq:pdfv2-terminal-tipo-fibra}
 \operatorname{Fib}_r(x):=
 \bigl(\mathsf H_x^\Sigma\oplus\mathsf H_x^\Pi,\,
       \mathsf M_x,\,
       \operatorname{Cell}_{\rm TPK}(x),\,C_{x,r}\bigr).
\end{equation}
El campo \(\mathsf F_{k,j}\) es una palabra en la categoría de estas fibras.
Su semilla y actualización son
\begin{align}
 \mathsf F_{k,k}
 &:=\bigl(\operatorname{Fib}_r(x_0),I\bigr)\star
 \mathop{\star}_{i=1}^{k}
 \bigl(\operatorname{Fib}_r(x_i),U_{\varepsilon_i},
 \mathsf C_{\varepsilon_i},\operatorname{Cell}_{\rm TPK}(\varepsilon_i),
 F_{\varepsilon_i}\bigr),
 \label{eq:pdfv2-terminal-semilla-f}\\
 \operatorname{Upd}^{\mathsf F}_{\varepsilon}\mathsf F_{k,j}
 &:=
 \mathsf F_{k,j}\star
 \bigl(\operatorname{Fib}_r(x_{j+1}),U_\varepsilon,
 \mathsf C_\varepsilon,\operatorname{Cell}_{\rm TPK}(\varepsilon),
 F_\varepsilon\bigr).
 \label{eq:pdfv2-terminal-update-f}
\end{align}
Es no vacía porque las dos hojas, la memoria, la unidad celular
\eqref{eq:pdfv2-cor-cell-unidad} y el complejo de puertos ya están
construidos. La composición de sus cinco mapas prueba la naturalidad.

\paragraph{Extensiones \(\mathsf E\): morfismos de prolongación compatible del estado}
El codominio es la categoría de palabras de relaciones tipadas. Ponemos
\begin{align}
 \mathsf E_{k,k}
 &:=(I_{x_0},\operatorname{Ext}_{I_{x_0}})\star
 \mathop{\star}_{i=1}^{k}
 \bigl(\varepsilon_i,\operatorname{Ext}_{\varepsilon_i},
 \operatorname{Ext}_{\varepsilon_i}\circ
 \operatorname{Ext}_{\gamma_{0,i-1}}
 \subseteq\operatorname{Ext}_{\gamma_{0,i}}\bigr),
 \label{eq:pdfv2-terminal-semilla-e}\\
 \operatorname{Upd}^{\mathsf E}_{\varepsilon}\mathsf E_{k,j}
 &:=
 \mathsf E_{k,j}\star
 \bigl(\varepsilon,\operatorname{Ext}_{\varepsilon},
 \operatorname{Ext}_{\varepsilon}\circ
 \operatorname{Ext}_{\gamma_{0,j}}
 \subseteq\operatorname{Ext}_{\gamma_{0,j+1}}\bigr).
 \label{eq:pdfv2-terminal-update-e}
\end{align}
Identidad y emisión neutra prueban no vacuidad; asociatividad de la
composición relacional prueba naturalidad.

\paragraph{Relaciones \(\mathsf R\): incidencia, orientación y compatibilidad de composición}
Sea \(\operatorname{RelBlock}_{j,r}\) el tipo de bloques de ecuaciones
etiquetados por la arista y el prefijo sobre los que se verifican, y sea
\(\operatorname{Pres}_{j,r}:=\operatorname{List}
(\operatorname{RelBlock}_{j,r})\). El bloque unidad y la semilla son
\begin{equation}
\label{eq:pdfv2-terminal-semilla-r}
 \mathsf B_I:=(I_{x_0};U_I=I,\ \mathsf U_I=I,\ \kappa(I)=0,\ \partial I=0),
 \qquad
 \mathsf R_{k,k}:=(\mathsf B_I)\star
 \mathop{\star}_{i=1}^{k}
 \mathsf B_R(\varepsilon_i;\gamma_{0,i-1}).
\end{equation}
Para una arista \(\varepsilon:x_j\to x_{j+1}\), el bloque nuevo es la
tupla etiquetada
\begin{equation}
\label{eq:pdfv2-terminal-bloque-r}
 \begin{split}
 \mathsf B_R(\varepsilon;\gamma_{0,j}):=\bigl(
 &\varepsilon,\gamma_{0,j};\ 
 \delta^\diamond(\varepsilon)
 =\operatorname{res}^\diamond(\varepsilon)
  +9\operatorname{quo}^\diamond(\varepsilon),\\
 &\Delta(\varepsilon)=\operatorname{or}(\varepsilon)
  +3\kappa(\varepsilon),\quad
 U_{\gamma_{0,j+1}}=U_\varepsilon U_{\gamma_{0,j}},\\
 &\mathsf U_{\gamma_{0,j+1}}
 =\mathsf U_\varepsilon\mathsf U_{\gamma_{0,j}},\quad
 \kappa(\gamma_{0,j+1})
 =\kappa(\varepsilon)+\mathsf C_\varepsilon\kappa(\gamma_{0,j}),\\
 &\partial\gamma_{0,j+1}
 =\partial\varepsilon+\partial\gamma_{0,j},\quad
 \operatorname{Cpl}_\varepsilon
 \text{ de }\eqref{eq:pdfv2-cor-acoplamientos-unicos}
 \bigr).
 \end{split}
\end{equation}
La actualización concatena el bloque sin perder fecha ni multiplicidad:
\begin{equation}
\label{eq:pdfv2-terminal-update-r}
 \operatorname{Upd}^{\mathsf R}_{\varepsilon}\mathsf R_{k,j}
 :=\mathsf R_{k,j}\star\mathsf B_R(\varepsilon;\gamma_{0,j}).
\end{equation}
El truncamiento \(\operatorname{tr}^{\mathsf R}_{j+1,j}\) elimina el último
bloque etiquetado. Además
\(a_*\mathsf B_R(\varepsilon;\gamma)
=\mathsf B_R(a\varepsilon;a\gamma)\), porque cada ecuación es natural en
la misma arista y el mismo prefijo.

\paragraph{Números \(\mathsf N\): publicación del valor con conservación de genealogía}
Para una arista \(\varepsilon:x_j\to x_{j+1}\), definimos el dato numérico
\begin{equation}
\label{eq:pdfv2-terminal-numero-arista}
 \nu(\varepsilon):=
 \bigl(j+1,\mathsf Q_{j+1},g_{j+1},q^{(9)}_{j+1},\beta_{j+1},
 \operatorname{res}^{\Sigma},\operatorname{quo}^{\Sigma},
 \operatorname{res}^{\Pi},\operatorname{quo}^{\Pi},
 \operatorname{or},\kappa,\Delta\bigr)(\varepsilon).
\end{equation}
Su tipo es
\[
 \mathbb N\times\widetilde{\mathcal O}_9\times C_9\times\mathbb N
 \times\mathbb Z/12\mathbb Z
 \times\{0,\ldots,8\}^2\times\mathbb Z^5.
\]
La semilla contiene la lista numérica completa del prefijo \(0\to k\);
la actualización añade exactamente el dato de la arista siguiente:
\begin{equation}
\label{eq:pdfv2-terminal-update-n}
 \mathsf N_{k,k}:=
 ((0,\mathsf Q_0,g_0,q^{(9)}_0,\beta_0;0,\ldots,0))
 \star\mathop{\star}_{i=1}^{k}\nu(\varepsilon_i),\qquad
 \operatorname{Upd}^{\mathsf N}_{\varepsilon}\mathsf N_{k,j}
 :=\mathsf N_{k,j}\star\nu(\varepsilon).
\end{equation}
Totalidad procede de las divisiones APP, la división TRIT y el odómetro TPK.

\paragraph{Orientación \(\mathsf O\): transporte de signo y selección de hoja}
El tipo de un estado de orientación es
\[
 \mathbb Z\times
 (\mathsf H^\Sigma_x\oplus\mathsf H^\Pi_x)
 \times\mathbb Z[\mathsf A],
\]
con semilla y actualización
\begin{align}
 \mathsf O_{k,k}&:=
 \left(\sum_{i=1}^{k}\operatorname{or}(\varepsilon_i),
       \operatorname{Sh}(x_k),\partial\gamma_{0,k}\right),
 \label{eq:pdfv2-terminal-semilla-o}\\
 \operatorname{Upd}^{\mathsf O}_{\varepsilon}
 (\mathfrak o,\operatorname{Sh},\mathfrak b)
 &:=
 \bigl(\mathfrak o+\operatorname{or}(\varepsilon),
 U_\varepsilon\operatorname{Sh},
 \mathfrak b+\partial\varepsilon\bigr).
 \label{eq:pdfv2-terminal-update-o}
\end{align}
Las hojas APP prueban no vacuidad. Functorialidad de \(U\), aditividad de
\(\operatorname{or}\) y telescopía de la frontera prueban naturalidad.

\paragraph{Acarreo \(\mathsf K\): actualización del cociente y memoria residual}
Usamos los lifts enteros tipados en
\eqref{eq:pdfv2-cor-carry-matricial-def}. El campo tiene tipo
\[
 \mathsf M_{x_j}\times
 \operatorname{Hom}(\mathsf M_{x_0},\mathsf M_{x_j})
 \times\operatorname{Mat}_5(\mathbb Z)
 \times\operatorname{List}(\operatorname{Mat}_5(\mathbb Z))
 \times\operatorname{List}(\mathbb Q_{\geq0}).
\]
Para \(\gamma_{k,j}=\varepsilon_j\cdots\varepsilon_{k+1}\), fijamos
explícitamente
\begin{equation}
\label{eq:pdfv2-terminal-lift-entero-ruta}
 \widetilde L_{\gamma_{k,k}}:=I_5,\qquad
 \widetilde L_{\gamma_{k,j}}
 :=\widetilde L_r(\varepsilon_j)\cdots
   \widetilde L_r(\varepsilon_{k+1}).
\end{equation}
Para el prefijo heredado fijamos asimismo
\[
 \widetilde L_{\gamma_{0,k}}
 :=\widetilde L_r(\varepsilon_k)\cdots
   \widetilde L_r(\varepsilon_1).
\]
La semilla y actualización son
\begin{align}
 \mathsf K_{k,k}&:=
 \bigl(\kappa(\gamma_{0,k}),\mathsf C_{\gamma_{0,k}},
       \widetilde L_{\gamma_{0,k}},
       \mathbf c_{\leq k},\mathbf k_{\leq k}\bigr),
 \label{eq:pdfv2-terminal-semilla-k}\\
 \operatorname{Upd}^{\mathsf K}_{\varepsilon}
 (\kappa_\gamma,\mathsf C_\gamma,\widetilde L_\gamma,
   \mathbf c,\mathbf k)
 &:=
 \bigl(\kappa(\varepsilon)+\mathsf C_\varepsilon\kappa_\gamma,\,
 \mathsf C_\varepsilon\mathsf C_\gamma,\,
 \widetilde L_r(\varepsilon)\widetilde L_\gamma,\,
 \mathbf c\star c_r(\varepsilon,\gamma),\,
 \mathbf k\star\kappa_j^{\rm can}\bigr).
 \label{eq:pdfv2-terminal-update-k}
\end{align}
La identidad proporciona la semilla; los dos cociclos
\eqref{eq:continuo-comun-cociclo-carry} y
\eqref{eq:pdfv2-cor-cociclo-matricial} prueban independencia del
parentizado y naturalidad.

\paragraph{Memoria \(\mathsf M\): registro del recorrido y compatibilidad bajo truncamiento}
Su tipo conserva a la vez memoria afín, ledger y memorias firmadas:
\begin{equation}
\label{eq:pdfv2-terminal-tipo-ledger-entry}
 \operatorname{LedgerEntry}:=
 \{\lambda(\varepsilon):\varepsilon\text{ emisión elemental tipada}\}.
\end{equation}
\[
 \mathsf M_x\times\operatorname{List}(\operatorname{LedgerEntry})
 \times\mathbb N^2\times\mathbb Z\times\mathbb N.
\]
Ponemos
\begin{align}
 \mathsf M_{k,k}&:=
 \bigl(m_k,\operatorname{Led}(\gamma_{0,k}),
       M_k^+,M_k^-,D_k,H_k\bigr),
 \label{eq:pdfv2-terminal-semilla-m}\\
 \operatorname{Upd}^{\mathsf M}_{\varepsilon}
 (m,\mathbf{Led},M^+,M^-,D,H)
 &:=
 \bigl(\mathsf C_\varepsilon m+\kappa(\varepsilon),\,
 \mathbf{Led}\star\lambda(\varepsilon),\,
 M^++b_\varepsilon^+,\,M^-+b_\varepsilon^-,\,
 D+b_\varepsilon,\,H+|b_\varepsilon|\bigr).
 \label{eq:pdfv2-terminal-update-m}
\end{align}
La acción afín y \eqref{eq:pdfv2-cor-memorias} prueban totalidad y
naturalidad; el pasado no se recalcula del último valor.

\paragraph{Frontera \(\mathsf B\): datos de empalme y supervivencia de la extensión}
El tipo es una cadena de frontera acompañada por una palabra de complejos
basados y mapas de cadena. Su semilla y actualización son
\[
 \mathsf B_{k,j}\in
 \operatorname{Rec}[\mathfrak b\in\mathbb Z[\mathsf A_j];
 C_{x_j,r}\in\operatorname{Ch}_{\rm bas};
 \mathbf F\in\operatorname{Path}_{\operatorname{Ch}_{\rm bas}}
 (C_{x_j,r},C_{x_0,r})].
\]
\begin{align}
 \mathsf B_{k,k}&:=
 \bigl(\partial\gamma_{0,k},C_{x_k,r},\mathbf F_{\leq k}\bigr),
 \label{eq:pdfv2-terminal-semilla-b}\\
 \operatorname{Upd}^{\mathsf B}_{\varepsilon}
 (\mathfrak b,C_{x_j,r},\mathbf F)
 &:=
 \bigl(\mathfrak b+\partial\varepsilon,\,
 C_{x_{j+1},r},\,\mathbf F\star F_\varepsilon\bigr).
 \label{eq:pdfv2-terminal-update-b}
\end{align}
No vacuidad procede de los puertos comunes; \(\partial^2=0\),
\(F_{\beta\varepsilon}=F_\varepsilon F_\beta\) y
\eqref{eq:continuo-comun-frontera-telescopica} prueban compatibilidad.

\paragraph{Ruta \(\mathsf G\): historia observable de la extensión superviviente}
El campo es la lista de prefijos en la categoría de caminos:
\begin{equation}
\label{eq:pdfv2-terminal-update-g}
 \mathsf G_{k,k}:=(\gamma_{0,i})_{0\leq i\leq k},\qquad
 \operatorname{Upd}^{\mathsf G}_{\varepsilon}\mathsf G_{k,j}
 :=\mathsf G_{k,j}\star\gamma_{0,j+1}.
\end{equation}
Identidad, concatenación asociativa y
\(a(\varepsilon\gamma)=a(\varepsilon)a(\gamma)\) prueban las tres
propiedades requeridas.

\paragraph{Multisección \(\mathsf X\): fibra finita estable y obstrucción a la selección puntual}
Para \(M\geq j\), definimos por acción
\begin{equation}
\label{eq:pdfv2-terminal-multiseccion-interna}
 \mathsf X_{k,j}(M):=
 \{\gamma_{k,j}\star q:
 q\in\operatorname{Term}_{j,M}(x_j)\}.
\end{equation}
Así
\[
 \mathsf X_{k,k}(M)=\operatorname{Term}_{k,M}(x_k),
\]
y
\begin{equation}
\label{eq:pdfv2-terminal-update-x}
 \operatorname{Upd}^{\mathsf X}_{\varepsilon}\mathsf X_{k,j}(M)
 :=
 \{\gamma_{k,j}\star\varepsilon\star q:
 q\in\operatorname{Term}_{j+1,M}(x_{j+1})\}.
\end{equation}
La prolongación neutra prueba no vacuidad. Borrar el último paso y actuar
por una simetría conmutan literalmente con la concatenación.

\paragraph{Supervivencia \(\mathsf S\): persistencia bajo refinamiento y paso al límite}
Aquí no se guarda una proposición, sino testigos. Sea
\begin{equation}
\label{eq:pdfv2-terminal-testigo-neutro}
 s_{j,j}:=I_{x_j},\qquad
 s_{j,M+1}:=s_{j,M}\star
 \varepsilon_M^0(t(s_{j,M})).
\end{equation}
La semilla es
\(\mathsf S_{k,k}:=(s_{k,M})_{M\geq k}\), y
\begin{equation}
\label{eq:pdfv2-terminal-update-s}
 \operatorname{Upd}^{\mathsf S}_{\varepsilon}\mathsf S_{k,j}
 :=(\gamma_{k,j}\star\varepsilon\star s_{j+1,M})_{M\geq j+1}.
\end{equation}
Por construcción
\(\pi_{M+1,M}s_{j,M+1}=s_{j,M}\); la naturalidad de la emisión neutra
prueba la equivariancia.

\paragraph{Terminal \(\mathsf T\): universalidad del objeto límite y factorización de lectores}
El terminal es el sistema inverso, no una cola escogida:
\begin{equation}
\label{eq:pdfv2-terminal-sistema-inverso-campo}
 \mathsf T(x_j):=
 \bigl((\operatorname{Term}_{j,M}(x_j))_{M\geq j},
       (\pi_{M+1,M})_{M\geq j}\bigr).
\end{equation}
Para cada \(M\geq j\), el mapa de prefijo tipado es
\begin{equation}
\label{eq:pdfv2-terminal-update-t}
 \operatorname{pref}_{\gamma_{k,j},M}:
 \operatorname{Term}_{j,M}(x_j)\longrightarrow
 \operatorname{Term}_{k,M}(x_k),\qquad
 q\longmapsto\gamma_{k,j}\star q.
\end{equation}
Definimos el estado del campo y su actualización por
\begin{align}
 \mathsf T_{k,j}
 &:=
 \left(\mathsf T(x_j),
   (\operatorname{pref}_{\gamma_{k,j},M})_{M\geq j}\right),
 \label{eq:pdfv2-terminal-estado-t}\\
 \operatorname{Upd}^{\mathsf T}_{\varepsilon}\mathsf T_{k,j}
 &:=
 \left(\mathsf T(x_{j+1}),
   (\operatorname{pref}_{\gamma_{k,j+1},M})_{M\geq j+1}\right),
 \qquad \gamma_{k,j+1}:=\gamma_{k,j}\star\varepsilon.
 \label{eq:pdfv2-terminal-update-t-explicita}
\end{align}
En particular, \(\mathsf T_{k,k}\) usa
\(\gamma_{k,k}=I_{x_k}\). Los testigos
\eqref{eq:pdfv2-terminal-testigo-neutro} y la sobreyectividad de \(\pi\)
prueban que el sistema es no vacío; nivel por nivel,
\begin{equation}
\label{eq:pdfv2-terminal-prefijo-natural}
 \pi^{\mathscr H}_{M+1,M}
 \operatorname{pref}_{\gamma_{k,j},M+1}
 =\operatorname{pref}_{\gamma_{k,j},M}
  \pi^{\mathscr H}_{M+1,M},
\end{equation}
porque ambos recorridos sólo eliminan la última emisión de \(q\).

\paragraph{13. Lector terminal \(\mathsf L\): evaluación de la extensión compatible y conservación de genealogía}
El lector no contiene ningún objeto clásico. Su tipo es la categoría de
caminos del único constructor dependiente
\(\mathfrak D_{j,r}\) de \eqref{eq:pdfv2-cor-constructor-unico}. Ponemos
\begin{equation}
\label{eq:pdfv2-terminal-transicion-lector-tipado}
 u^{\mathfrak D}_{j+1,j}
 :=u_{(j+1,r),(j,r)}:
 \mathfrak D_{j+1,r}(x_{j+1})
 \longrightarrow\mathfrak D_{j,r}(x_j),
\end{equation}
la componente del transporte común
\eqref{eq:pdfv2-cor-naturalidad-unica} que trunca la última emisión, reduce
los coeficientes y precompone los puertos. Por composición,
\begin{equation}
\label{eq:pdfv2-terminal-transicion-lector-composicion}
 u^{\mathfrak D}_{m,k}
 =u^{\mathfrak D}_{\ell,k}u^{\mathfrak D}_{m,\ell}
 \qquad(m\geq\ell\geq k).
\end{equation}
\begin{align}
 \mathsf L_{k,k}
 &:=(\mathfrak D_{0,r}(x_0))\star
 \mathop{\star}_{i=1}^{k}
 \bigl(\mathfrak D_{i,r}(x_i),u^{\mathfrak D}_{i,i-1},
       \mathcal O(\mathfrak e_{\varepsilon_i})\bigr),
 \label{eq:pdfv2-terminal-semilla-l}\\
 \operatorname{Upd}^{\mathsf L}_{\varepsilon}\mathsf L_{k,j}
 &:=
 \mathsf L_{k,j}\star
 \bigl(\mathfrak D_{j+1,r}(x_{j+1}),
 u^{\mathfrak D}_{j+1,j},\mathcal O(\mathfrak e_\varepsilon)\bigr).
 \label{eq:pdfv2-terminal-update-l}
\end{align}
La totalidad de \(\mathfrak D\) y
\eqref{eq:pdfv2-terminal-transicion-lector-composicion} prueban no vacuidad
y naturalidad.

Las trece semillas anteriores forman un único elemento
\begin{equation}
\label{eq:pdfv2-terminal-semilla-trece}
 \mathfrak r^0_{k,r}:=
 \langle
 \mathsf F_{k,k};\mathsf E_{k,k};\mathsf R_{k,k};\mathsf N_{k,k};
 \mathsf O_{k,k};\mathsf K_{k,k};\mathsf M_{k,k};\mathsf B_{k,k};
 \mathsf G_{k,k};\mathsf X_{k,k};\mathsf S_{k,k};\mathsf T_{k,k};
 \mathsf L_{k,k}
 \rangle\in\mathfrak R_{k,r}.
\end{equation}
La acción de una arista es la aplicación total
\begin{equation}
\label{eq:pdfv2-terminal-actualizacion-unica}
 \operatorname{Upd}_\varepsilon:
 \mathfrak R_{j,r}(\gamma_{0,j})
 \longrightarrow\mathfrak R_{j+1,r}(\gamma_{0,j+1}),\qquad
 \operatorname{Upd}_\varepsilon:=
 \langle\operatorname{Upd}^{\mathsf F}_\varepsilon,\ldots,
 \operatorname{Upd}^{\mathsf L}_\varepsilon\rangle,
\end{equation}
con las trece acciones
\eqref{eq:pdfv2-terminal-update-f}--
\eqref{eq:pdfv2-terminal-update-l}. Finalmente, el registro terminal se
\emph{genera} por el fold
\begin{equation}
\label{eq:pdfv2-terminal-registro-trece}
 \boxed{
 \mathfrak r_{k,N}(h):=
 \operatorname{Upd}_{\varepsilon_N}\circ\cdots\circ
 \operatorname{Upd}_{\varepsilon_{k+1}}(\mathfrak r^0_{k,r}).}
\end{equation}
No es el desempaquetado de una copia intacta del estado.
Sobre la imagen de este fold se define, componente a componente, el
truncamiento
\begin{equation}
\label{eq:pdfv2-terminal-truncamiento-trece-def}
 \operatorname{tr}_{13;j+1,j}
 \mathfrak r_{k,j+1}(h\star\varepsilon)
 :=\mathfrak r_{k,j}(h).
\end{equation}
Está bien definido: los campos \(\mathsf E\) y \(\mathsf G\) conservan la
lista ordenada de emisiones y prefijos, de modo que la igualdad de dos
registros implica la igualdad de sus historias etiquetadas. En
\(\mathsf F,\mathsf E,\mathsf R,\mathsf N,\mathsf G,\mathsf T,\mathsf L\)
se elimina la última entrada o bloque; en \(\mathsf O,\mathsf K,\mathsf M,
\mathsf B\) se aplica la fórmula acumulativa al prefijo recuperado; y en
\(\mathsf X,\mathsf S\) se restringe el cono de colas. No se invierte una
unión idempotente ni se selecciona una continuación.

\begin{proposition}[Totalidad y naturalidad de las trece actualizaciones]
\label{prop:pdfv2-terminal-trece-total-natural}
Todas las semillas de \eqref{eq:pdfv2-terminal-semilla-trece} existen, cada
\(\operatorname{Upd}_\varepsilon\) es total y, para toda simetría coherente
\(a\) y todo truncamiento de horizonte,
\begin{equation}
\label{eq:pdfv2-terminal-trece-naturalidad}
 a_*\operatorname{Upd}_\varepsilon
 =\operatorname{Upd}_{a\varepsilon}a_*,
 \qquad
 \operatorname{tr}_{13;j+1,j}
 \operatorname{Upd}_\varepsilon\mathfrak r_{k,j}(h)
 =\mathfrak r_{k,j}(h).
\end{equation}
\end{proposition}

\begin{proof}
La existencia de cada semilla fue exhibida en los trece párrafos. En
\(\mathsf F,\mathsf E,\mathsf G,\mathsf X,\mathsf S,\mathsf T,\mathsf L\),
la fórmula es concatenación con un mapa ya tipado. En
\(\mathsf R,\mathsf N,\mathsf O,\mathsf K,\mathsf M,\mathsf B\), la
fórmula es, respectivamente, concatenación de un bloque de ecuaciones
etiquetado, adjunción de un dato numérico, transporte de hojas, composición
de cociclos, acción afín y mapa de cadena. Las identidades de naturalidad
citadas después de cada fórmula prueban la primera igualdad de
\eqref{eq:pdfv2-terminal-trece-naturalidad} componente a componente. La
segunda es exactamente
\eqref{eq:pdfv2-terminal-truncamiento-trece-def}; su buena definición usa la
historia ordenada conservada por \(\mathsf E\) y \(\mathsf G\).
\end{proof}

\begin{proposition}[Certificado conjunto de los cinco desarrollos internos]
\label{prop:pdfv2-terminal-cinco-desarrollos-integrales}
Las actualizaciones de las trece coordenadas transportan conjuntamente:
la telescopía espectral--polar y el retorno sin reinicialización de
\cref{thm:continuo-sp-telescopia-no-autonoma,eq:continuo-sp-no-reset};
el balance ortogonal, la columna de observación y la identidad de Stein de
\cref{eq:pdfv2-sol-ortogonalidad,eq:pdfv2-sol-stein}; la naturalidad directa
y adjunta y el retracto celular de
\cref{eq:gauint-naturalidad-directa-adjunta,eq:gauint-sdr}; la totalización,
el cono del retorno nonádico y las fibras completas de
\cref{eq:pdfv2-coh-total-square-zero,eq:pdfv2-coh-terminal-cone,eq:pdfv2-coh-extension-multisection}; y la diagonal de
Alexander--Whitney, la reducción de Smith y la torre de Bockstein de
\cref{eq:detint-aw-natural,eq:detint-smith,eq:detint-bock-leading}.
Todas estas acciones tienen como dominio el mismo registro generado por
\eqref{eq:pdfv2-terminal-registro-trece}; ninguna introduce un subdominio
inicial, una historia privilegiada ni una selección motivada por una
aplicación posterior.
\end{proposition}

\begin{proof}
Cada identidad citada actúa sobre una coordenada del registro común y sus
mapas de truncamiento son restricciones de
\eqref{eq:pdfv2-terminal-truncamiento-trece-def}. La naturalidad componente
a componente demostrada en
\eqref{eq:pdfv2-terminal-trece-naturalidad} permite ensamblarlas en una sola
actualización. La igualdad de dominio y la conservación de la historia
ordenada impiden reemplazar la generación correlativa por cinco
proyecciones retrospectivas.
\end{proof}

\section{Terminal multivaluado sin selector}

El terminal finito generado por \(\widehat x_k\) es la multisecci\'on
\begin{equation}
\label{eq:pdfv2-terminal-multivaluado}
 \mathfrak T_{k,N}(\widehat x_k)
 :=\left\{\mathfrak r_{k,N}(h):
            h\in\mathscr H_{k,N}(\widehat x_k)\right\}.
\end{equation}
La expresi\'on de la derecha devuelve el conjunto entero. No contiene una
regla que distinga una historia privilegiada.

El ambiente terminal de nivel se obtiene reuniendo esas multisecciones sin
identificar registros:
\begin{equation}
\label{eq:pdfv2-terminal-ambiente-nivel}
 \mathfrak T_{k,N}
 :=\bigcup_{\widehat x_k\in\widehat{\mathcal X}^{\rm tot}_k}
   \mathfrak T_{k,N}(\widehat x_k).
\end{equation}

La versi\'on lineal, cuando se necesita sumar contribuciones, se define en
el m\'odulo libre por
\begin{equation}
\label{eq:pdfv2-terminal-evaluacion-todas}
 \begin{aligned}
 \mathsf E_{k,N}^{\rm all}:
 \mathbb Z[\widehat{\mathcal X}^{\rm tot}_k]&\longrightarrow
 \mathbb Z[\mathfrak T_{k,N}],\\
 [\widehat x_k]&\longmapsto
 \sum_{h\in\mathscr H_{k,N}(\widehat x_k)}
 [\mathfrak r_{k,N}(h)].
 \end{aligned}
\end{equation}
Cada testigo aparece con su multiplicidad. Si existe una medida proyectiva
positiva \(\mu\), la versi\'on isom\'etrica usa todas las probabilidades
condicionales:
\begin{equation}
\label{eq:pdfv2-terminal-evaluacion-ponderada}
 \widehat{\mathsf E}_{k,N}^{\mu}e_{\widehat x_k}
 :=\sum_{h\in\mathscr H_{k,N}(\widehat x_k)}
 \mu_{\widehat x_k,N}(h)^{1/2}
 e_{\mathfrak r_{k,N}(h)}.
\end{equation}
La proyectividad da
\begin{equation}
\label{eq:pdfv2-terminal-isometria}
 \left\|\widehat{\mathsf E}_{k,N}^{\mu}e_{\widehat x_k}\right\|^2
 =\sum_{h\in\mathscr H_{k,N}(\widehat x_k)}
   \mu_{\widehat x_k,N}(h)=1.
\end{equation}
La suma no se sustituye por ``la historia que genera el cilindro''.

\section{La conexión nonádica como correspondencia}

La holonom\'ia non\'adica fue generada por el od\'ometro y la conexi\'on TPK,
sin hip\'otesis de pertenencia, en
\eqref{eq:continuo-comun-gamma9-apice}--
\eqref{eq:continuo-comun-gamma9-relacion}. En este cap\'itulo usamos
\begin{equation}
\label{eq:pdfv2-terminal-apice-gamma9}
 \mathsf Z_{9,k}:=\mathsf Z^{\Gamma}_{9,k}.
\end{equation}
Los mapas
\begin{equation}
\label{eq:pdfv2-terminal-span-gamma9}
 \widehat{\mathcal X}^{\rm tot}_k
 \xleftarrow{\ s_{9,k}\ }
 \mathsf Z^{\Gamma}_{9,k}
 \xrightarrow{\ t_{9,k}\ }
 \widehat{\mathcal X}^{\rm tot}_{k+9}
\end{equation}
son los de \eqref{eq:continuo-comun-gamma9-span}; su fuente es sobreyectiva
por \cref{prop:continuo-comun-gamma9-total-natural}. Cada punto del \`apice
conserva las nueve emisiones fase--compatibles. No se reemplaza por el par
de sus extremos.

Para \(N\geq k+9\), definimos el \`apice de horizonte antes de linealizar:
\begin{equation}
\label{eq:pdfv2-terminal-apice-horizonte-gamma9}
 \mathsf Z^\Gamma_{9;k,N}:=
 \coprod_{z=(\widehat x_k,\boldsymbol\omega)
            \in\mathsf Z^\Gamma_{9,k}}
 \operatorname{Term}_{k+9,N}(t_{9,k}z).
\end{equation}
Un elemento se escribe \((z,q)\), donde \(\boldsymbol\omega\) es la palabra
non\'adica de nueve emisiones y \(q\) su cola hasta \(N\). Los mapas sobre
historias usan el codominio etiquetado
\begin{equation}
\label{eq:pdfv2-terminal-gamma9-codominio-colas}
 \mathscr D^\Gamma_{9;k,N}:=
 \coprod_{z\in\mathsf Z^\Gamma_{9,k}}
 \operatorname{Term}_{k+9,N}(t_{9,k}z).
\end{equation}
Así los mapas son
\begin{align}
 S_{k,N}(z,q)&:=\boldsymbol\omega\star q,
 \label{eq:pdfv2-terminal-gamma9-s-horizonte}\\
 T_{k,N}(z,q)&:=(z,q)\in\mathscr D^\Gamma_{9;k,N}.
 \label{eq:pdfv2-terminal-gamma9-t-horizonte}
\end{align}
El primero aterriza en
\(\coprod_x\mathscr H_{k,N}(x)\); el segundo conserva expresamente el
sumando etiquetado por \(z\).

El mapa de enlace del \`apice es la acci\'on expl\'icita
\begin{equation}
\label{eq:pdfv2-terminal-gamma9-enlace-apice}
 \pi^Z_{N+1,N}(z,q\star\varepsilon):=(z,q).
\end{equation}
Es sobreyectivo: si \((z,q)\) est\'a dado, se a\~nade a la cola la emisi\'on
neutra del extremo de \(q\). El codominio de colas tiene el enlace
\begin{equation}
\label{eq:pdfv2-terminal-gamma9-enlace-codominio}
 \pi^{\mathscr D}_{N+1,N}(z,q\star\varepsilon):=(z,q).
\end{equation}
Denotamos por
\(\pi^{\mathscr H,j}_{N+1,N}\) el mapa
\eqref{eq:pdfv2-terminal-truncamiento-historias} con nivel inicial \(j\).
Entonces
\begin{align}
 \pi^{\mathscr H,k}_{N+1,N}S_{k,N+1}
 &=S_{k,N}\pi^Z_{N+1,N},
 \label{eq:pdfv2-terminal-gamma9-cuadrado-fuente}\\
 \pi^{\mathscr D}_{N+1,N}T_{k,N+1}
 &=T_{k,N}\pi^Z_{N+1,N}.
 \label{eq:pdfv2-terminal-gamma9-cuadrado-destino}
\end{align}
En el primer cuadrado ambos recorridos env\'ian
\(\boldsymbol\omega\star q\star\varepsilon\) a
\(\boldsymbol\omega\star q\); en el segundo env\'ian
\((z,q\star\varepsilon)\) a \((z,q)\). Por descomposici\'on \`unica de una palabra
en sus primeras nueve emisiones y su cola, \(S_{k,N}\) es biyectivo. El
segundo cuadrado es cartesiano por la definici\'on coproducto--fibrada de
\eqref{eq:pdfv2-terminal-apice-horizonte-gamma9}. Son los cuadrados de
Beck--Chevalley que conservan testigos y multiplicidades.

El \`apice correspondiente de registros, ahora generado por las trece
actualizaciones, es
\begin{equation}
\label{eq:pdfv2-terminal-apice-terminal-gamma9}
 \begin{split}
 \mathsf Z^{\mathfrak T}_{9;k,N}:=\bigl\{
 (&z,q;
   \mathfrak r_{k,N}(\boldsymbol\omega\star q),
   \mathfrak r_{k+9,N}(q)):\\
 &z=(\widehat x_k,\boldsymbol\omega)\in
       \mathsf Z^\Gamma_{9,k},\quad
 q\in\operatorname{Term}_{k+9,N}(t_{9,k}z)
 \bigr\}.
 \end{split}
\end{equation}
Sus mapas fuente y destino leen, respectivamente, el primer o el segundo
registro, pero nunca eliminan de su dominio el testigo \((z,q)\):
\begin{equation}
\label{eq:pdfv2-terminal-span-terminal}
 \mathfrak T_{k,N}
 \xleftarrow{\ s^{\mathfrak T}_{9;k,N}\ }
 \mathsf Z^{\mathfrak T}_{9;k,N}
 \xrightarrow{\ t^{\mathfrak T}_{9;k,N}\ }
 \mathfrak T_{k+9,N}.
\end{equation}
Si \(\zeta=(z,q;r_s,r_t)\), su acción es
\begin{equation}
\label{eq:pdfv2-terminal-span-terminal-accion}
 s^{\mathfrak T}_{9;k,N}(\zeta):=r_s,
 \qquad
 t^{\mathfrak T}_{9;k,N}(\zeta):=r_t.
\end{equation}
Su enlace se define por
\begin{equation}
\label{eq:pdfv2-terminal-gamma9-enlace-registros}
 \begin{split}
 \pi^{Z,\mathfrak T}_{N+1,N}
 \bigl(&z,q\star\varepsilon;
       \mathfrak r_{k,N+1}
          (\boldsymbol\omega\star q\star\varepsilon),
       \mathfrak r_{k+9,N+1}(q\star\varepsilon)\bigr)\\
 &:=(z,q;\mathfrak r_{k,N}(\boldsymbol\omega\star q),
       \mathfrak r_{k+9,N}(q)).
 \end{split}
\end{equation}
La pareja \((z,q)\) es una coordenada del \`apice, no una representación
elegida de una terna; por ello el enlace está bien definido incluso si dos
registros terminales coinciden. La compatibilidad de
\(\operatorname{Upd}\) con el borrado de la \`ultima arista prueba que
\eqref{eq:pdfv2-terminal-gamma9-enlace-registros} conmuta con ambos mapas de
\eqref{eq:pdfv2-terminal-span-terminal}.

Si se desea una acci\'on sobre m\'odulos libres, se linealiza esta
correspondencia sumando todos los testigos:
\begin{equation}
\label{eq:pdfv2-terminal-linealizacion-span}
 [\Gamma_9]_{k,N}[r]
 :=\sum_{\zeta\in(s^{\mathfrak T}_{9;k,N})^{-1}(r)}
 [t^{\mathfrak T}_{9;k,N}(\zeta)].
\end{equation}
La f\'ormula conserva multiplicidades. La igualdad de fases despu\'es de
nueve pasos no contrae los estados, los testigos ni sus ledgers.

\section{Unicidad de la descomposición por prefijos y naturalidad del sistema de truncamientos}

Para \(k\leq\ell\leq N\), toda historia de
\(\mathscr H_{k,N}(\widehat x_k)\) posee un \'unico prefijo
\(p\in\mathscr P_{\ell\mid k}(\widehat x_k)\). Por tanto hay una descomposici\'on
disjunta por prefijos, expresada sin elegir ninguno:
\begin{equation}
\label{eq:pdfv2-terminal-particion-historias}
 \mathscr H_{k,N}(\widehat x_k)
 =\bigsqcup_{p\in\mathscr P_{\ell\mid k}(\widehat x_k)}
 \left\{p\star q:q\in\mathscr H_{\ell,N}(p)\right\}.
\end{equation}
Sea \(\operatorname{Conc}_p\) la operaci\'on que antepone el registro del
prefijo \(p\) y aplica sucesivamente
\eqref{eq:pdfv2-terminal-actualizacion-unica}. Definimos la transici\'on
com\'un sobre una familia completa de terminales por
\begin{equation}
\label{eq:pdfv2-terminal-transicion-comun}
 \mathfrak u_{\ell,k}
 \left(
   \bigl(\mathfrak T_{\ell,N}(p)\bigr)_{p\in\mathscr P_{\ell\mid k}(\widehat x_k)}
 \right)
 :=\bigcup_{p\in\mathscr P_{\ell\mid k}(\widehat x_k)}
   \operatorname{Conc}_p\mathfrak T_{\ell,N}(p).
\end{equation}

Escribiendo
\(\mathbf G_{k,N}(\widehat x_k):=\mathfrak T_{k,N}(\widehat x_k)\), toda la naturalidad se
concentra en la identidad \'unica
\begin{equation}
\label{eq:pdfv2-terminal-naturalidad-unica}
 \boxed{
 \mathfrak u_{\ell,k}
 \left(
   \bigl(\mathbf G_{\ell,N}(p)\bigr)_{p\in\mathscr P_{\ell\mid k}(\widehat x_k)}
 \right)
 =\mathbf G_{k,N}(\widehat x_k).}
\end{equation}
No hay una ley distinta para cada discriminante interno. Las cinco
caracter\'isticas obedecen
\eqref{eq:pdfv2-terminal-naturalidad-unica} porque cada una se actualiza
mediante el mismo prefijo, la misma extensi\'on y el mismo ledger.

La correspondencia non\'adica es compatible con esta identidad mediante los
mapas ya construidos, no mediante una condici\'on impl\'icita. Para todo
horizonte, \eqref{eq:pdfv2-terminal-gamma9-cuadrado-fuente} y
\eqref{eq:pdfv2-terminal-gamma9-cuadrado-destino} conservan el testigo
non\'adico con su multiplicidad, y
\eqref{eq:pdfv2-terminal-gamma9-enlace-registros} hace lo mismo en las
trece coordenadas. Si \(a\) es una simetr\'ia coherente del estado
com\'un, su acci\'on
\[
 a^Z(z,q):=(a^Zz,aq)
\]
satisface
\begin{equation}
\label{eq:pdfv2-terminal-gamma9-equivarianza-completa}
 S a^Z=aS,\qquad T a^Z=aT,\qquad
 \pi^Za^Z=a^Z\pi^Z,\qquad
 s^{\mathfrak T}a^Z=as^{\mathfrak T},\qquad
 t^{\mathfrak T}a^Z=at^{\mathfrak T}.
\end{equation}
Estas igualdades se verifican entrada por entrada usando
\eqref{eq:continuo-comun-gamma9-naturalidad} y la naturalidad de las trece
actualizaciones. Son la condici\'on material que permite usar
\eqref{eq:pdfv2-terminal-linealizacion-span}; una mera igualdad de valores
visibles no la sustituye.

\section{Paso al l\'imite sin intercambiar imagen y l\'imite}

Truncar una historia y cada uno de los trece campos define
\begin{equation}
\label{eq:pdfv2-terminal-transicion-horizonte}
 \pi^{\mathfrak T}_{N+1,N}:
 \mathfrak T_{k,N+1}(\widehat x_k)\longrightarrow
 \mathfrak T_{k,N}(\widehat x_k),
 \qquad
 \pi^{\mathfrak T}_{N+1,N}\mathfrak r_{k,N+1}(h)
 =\mathfrak r_{k,N}(h_{\leq N}).
\end{equation}
Es sobreyectiva por
\eqref{eq:pdfv2-terminal-truncamiento-historias}. El terminal completo es
entonces
\begin{equation}
\label{eq:pdfv2-terminal-limite-inverso}
 \mathfrak T_{k,\infty}(\widehat x_k)
 :=\left\{(r_N)_{N\geq k}:
 r_N\in\mathfrak T_{k,N}(\widehat x_k),\quad
 \pi^{\mathfrak T}_{N+1,N}r_{N+1}=r_N\right\}.
\end{equation}
Los conjuntos finitos de
\eqref{eq:pdfv2-terminal-multivaluado} son no vac\'ios y sus transiciones
son sobreyectivas; por ello
\begin{equation}
\label{eq:pdfv2-terminal-limite-no-vacio}
 \mathfrak T_{k,\infty}(\widehat x_k)\neq\varnothing.
\end{equation}

\begin{corollary}[Persistencia simultánea de los cinco desarrollos]
\label{cor:pdfv2-terminal-limite-cinco-desarrollos}
En cada elemento de
\(\mathfrak T_{k,\infty}(\widehat x_k)\) persisten simultáneamente el límite
fuerte \(\mathcal T_\infty\) de la telescopía polar, la familia de Gram y
pérdidas, la medida de incidencia con sus dos naturalidades, el cono
\(\operatorname{Cone}(I-U_{\Gamma_9,N})\) y la familia compatible de formas
de Smith y páginas de Bockstein. Ninguno de estos datos se obtiene
reconstruyendo sólo el valor terminal visible.
\end{corollary}

\begin{proof}
La sobreyectividad de los truncamientos conserva los testigos completos. Las
ecuaciones de telescopía, naturalidad y reducción de
\cref{prop:pdfv2-terminal-cinco-desarrollos-integrales} conmutan con esos
truncamientos; por universalidad del límite inverso definen una única familia
compatible sobre cada historia superviviente.
\end{proof}

El \`apice non\'adico en el l\'imite se define, no se infiere de una imagen:
\begin{equation}
\label{eq:pdfv2-terminal-gamma9-apice-limite-raw}
 \mathsf Z^\Gamma_{9;k,\infty}
 :=\varprojlim_{N\geq k+9}
 \bigl(\mathsf Z^\Gamma_{9;k,N},\pi^Z_{N+1,N}\bigr).
\end{equation}
Como los enlaces no modifican el testigo finito \(z\) y s\'olo truncan su
cola, existe el isomorfismo expl\'icito
\begin{equation}
\label{eq:pdfv2-terminal-gamma9-apice-limite-msect}
 \mathsf Z^\Gamma_{9;k,\infty}
 \xrightarrow{\ \cong\ }
 \coprod_{z\in\mathsf Z^\Gamma_{9,k}}
 \operatorname{Msect}(t_{9,k}z),\qquad
 ((z,q_N))_N\longmapsto(z,(q_N)_N).
\end{equation}
La inversa toma un testigo \(z\) y una cola compatible \((q_N)_N\), y forma
\(((z,q_N))_N\). Ambas composiciones son identidades coordenada a
coordenada. Los mapas l\'imite son
\begin{equation}
\label{eq:pdfv2-terminal-gamma9-mapas-limite-raw}
 S_{k,\infty}((z_N)_N):=(S_{k,N}z_N)_N,\qquad
 T_{k,\infty}((z_N)_N):=(T_{k,N}z_N)_N;
\end{equation}
est\'an bien definidos precisamente por
\eqref{eq:pdfv2-terminal-gamma9-cuadrado-fuente} y
\eqref{eq:pdfv2-terminal-gamma9-cuadrado-destino}.

Del mismo modo se define el \`apice de registros
\begin{equation}
\label{eq:pdfv2-terminal-gamma9-apice-registros-limite}
 \mathsf Z^{\mathfrak T}_{9;k,\infty}
 :=\varprojlim_{N\geq k+9}
 \bigl(\mathsf Z^{\mathfrak T}_{9;k,N},
       \pi^{Z,\mathfrak T}_{N+1,N}\bigr).
\end{equation}
Las correspondencias
\eqref{eq:pdfv2-terminal-span-terminal} forman entonces, componente a
componente, el span
\begin{equation}
\label{eq:pdfv2-terminal-span-limite}
 \mathfrak T_{k,\infty}
 \xleftarrow{\ s^{\mathfrak T}_{9;k,\infty}\ }
 \mathsf Z^{\mathfrak T}_{9;k,\infty}
 \xrightarrow{\ t^{\mathfrak T}_{9;k,\infty}\ }
 \mathfrak T_{k+9,\infty}.
\end{equation}
No se ha identificado la imagen de un l\'imite con el l\'imite de unas
im\'agenes: \eqref{eq:pdfv2-terminal-limite-inverso},
\eqref{eq:pdfv2-terminal-gamma9-apice-limite-raw} y
\eqref{eq:pdfv2-terminal-gamma9-apice-registros-limite} construyen
directamente las tres familias compatibles y sus mapas.

\section{Reconstrucción interna del estado holonómico integral a partir de sus \(13\) coordenadas genealógicas}

La recuperaci\'on no desempaqueta nombres: repite el fold que gener\'o el
registro. Para
\(h=(\varepsilon_{k+1},\ldots,\varepsilon_N)\), definimos
\begin{equation}
\label{eq:pdfv2-terminal-recuperacion-trece}
 \boxed{
 \operatorname{Rec}_{13;k,N}(h):=
 \operatorname{Upd}_{\varepsilon_N}\circ\cdots\circ
 \operatorname{Upd}_{\varepsilon_{k+1}}
 (\mathfrak r^0_{k,r}).}
\end{equation}
Por \eqref{eq:pdfv2-terminal-registro-trece},
\(\operatorname{Rec}_{13;k,N}(h)=\mathfrak r_{k,N}(h)\), pero la igualdad
es ahora una identidad entre dos composiciones de acciones, no la
proyecci\'on de una entrada preservada. Si
\(\operatorname{pr}_{\mathsf A}\) denota el proyector del registro al campo
\(\mathsf A\in\{\mathsf F,\mathsf E,\mathsf R,\mathsf N,\mathsf O,
\mathsf K,\mathsf M,\mathsf B,\mathsf G,\mathsf X,\mathsf S,\mathsf T,
\mathsf L\}\), entonces
\begin{equation}
\label{eq:pdfv2-terminal-recuperacion-campo-fold}
 \operatorname{pr}_{\mathsf A}\operatorname{Rec}_{13;k,N}(h)
 =\operatorname{Upd}^{\mathsf A}_{\varepsilon_N}\circ\cdots\circ
  \operatorname{Upd}^{\mathsf A}_{\varepsilon_{k+1}}
  (\mathsf A_{k,k}).
\end{equation}
La prueba es inducci\'on en \(N-k\): para \(N=k\) ambos lados son la
semilla; el paso inductivo aplica la componente \(\mathsf A\) de
\(\operatorname{Upd}_{\varepsilon_N}\).

Adem\'as,
\begin{equation}
\label{eq:pdfv2-terminal-recuperacion-compatible}
 \operatorname{Rec}_{13;k,N}(h_{\leq N})
 =\operatorname{tr}_{13,N+1,N}
   \bigl(\operatorname{Rec}_{13;k,N+1}(h)\bigr),
\end{equation}
donde \(\operatorname{tr}_{13,N+1,N}\) elimina s\'olo la \'ultima
actualizaci\'on y restringe su cono de colas. En consecuencia, un elemento
de \(\mathfrak T_{k,\infty}(\widehat x_k)\) determina una familia compatible de las
trece coordenadas a toda profundidad.

\section{Terminalidad multivaluada, naturalidad y límite compatible del continuo}

\begin{theorem}[Terminal multivaluado, naturalidad y l\'imite del continuo]
\label{thm:pdfv2-terminal-conjunto}
Para todo \(\widehat x_k\in\widehat{\mathcal X}^{\rm tot}_k\), las reglas
\eqref{eq:pdfv2-terminal-historias-finita}--
\eqref{eq:pdfv2-terminal-recuperacion-compatible} construyen, a partir del
\'unico continuo discreto:
\begin{enumerate}
 \item una multisecci\'on terminal finita no vac\'ia en cada horizonte;
 \item una evaluaci\'on que conserva todas las historias y multiplicidades;
 \item la conexi\'on non\'adica como correspondencia con testigos;
 \item la identidad \'unica de naturalidad
       \eqref{eq:pdfv2-terminal-naturalidad-unica};
 \item un terminal inverso no vac\'io y una correspondencia non\'adica en el
       l\'imite;
 \item la recuperaci\'on compatible de fibra, extensiones, relaciones,
       n\'umeros, orientaci\'on, acarreo, memoria, frontera, ruta,
       multisecci\'on, supervivencia, terminal y lector.
\end{enumerate}
Las cinco caracter\'isticas intr\'insecas permanecen correlacionadas dentro
de cada registro y de cada testigo de extensi\'on.
\end{theorem}

\begin{proof}
La no vacuidad de
\(\mathscr H_{k,N}(\widehat x_k)\) es no vac\'ia por
\cref{prop:continuo-comun-arbol-finito-no-vacio}, y su truncamiento es
sobreyectivo por \cref{lem:continuo-comun-terminal-sobreyectivo}. La regla
\eqref{eq:pdfv2-terminal-actualizacion-unica} aplica las identidades APP,
TRIT y TPK sobre cada testigo \(\varepsilon_j\); por inducci\'on en
\(N-k\), produce todos los campos de
\eqref{eq:pdfv2-terminal-registro-trece} y conserva sus relaciones.

Tomar el conjunto completo en
\eqref{eq:pdfv2-terminal-multivaluado}, o la suma completa en
\eqref{eq:pdfv2-terminal-evaluacion-todas}, demuestra que el terminal no
contiene un selector. La versi\'on ponderada es isom\'etrica por la igualdad
proyectiva \eqref{eq:pdfv2-terminal-isometria}.

El testigo \(z=(\widehat x_k,\boldsymbol\omega)\) lleva la acción de fase de cada una
de sus nueve emisiones; por
\cref{prop:continuo-comun-gamma9-retorno-memoria} define los dos mapas de
\eqref{eq:pdfv2-terminal-span-gamma9} sin convertir la relación en una
función ni equiparar fase y estado. Las fórmulas
\eqref{eq:pdfv2-terminal-gamma9-s-horizonte}--
\eqref{eq:pdfv2-terminal-gamma9-enlace-registros} prueban los dos cuadrados
de horizonte y la misma construcción sobre registros. Sumar el espacio
fuente de cada testigo da
\eqref{eq:pdfv2-terminal-linealizacion-span} sin borrar multiplicidades.

La partici\'on
\eqref{eq:pdfv2-terminal-particion-historias} es disjunta porque toda
historia tiene un prefijo de longitud \(\ell\) determinado. Concatenar cada
prefijo con todas sus continuaciones reconstruye una vez y s\'olo una cada
historia de \(\mathscr H_{k,N}(\widehat x_k)\). Esto prueba directamente
\eqref{eq:pdfv2-terminal-naturalidad-unica} para el registro entero; no se
necesitan cinco demostraciones inconexas.

La sobreyectividad de
\eqref{eq:pdfv2-terminal-transicion-horizonte} y la finitud de cada nivel
dan la no vacuidad de
\eqref{eq:pdfv2-terminal-limite-inverso}. Los enlaces explícitos
\eqref{eq:pdfv2-terminal-gamma9-enlace-apice} y
\eqref{eq:pdfv2-terminal-gamma9-enlace-registros}, junto con sus cuadrados,
construyen los límites
\eqref{eq:pdfv2-terminal-gamma9-apice-limite-raw} y
\eqref{eq:pdfv2-terminal-gamma9-apice-registros-limite}; de ellos sale
\eqref{eq:pdfv2-terminal-span-limite}. Finalmente,
\eqref{eq:pdfv2-terminal-recuperacion-trece} vuelve a ejecutar las
actualizaciones sucesivas, y
\eqref{eq:pdfv2-terminal-recuperacion-compatible} demuestra que esas
lecturas pasan conjuntamente al l\'imite.
\end{proof}

\section{Obstrucciones a la terminalidad del estado holonómico integral: pérdida de fibra, orientación, memoria o lector}

\begin{proposition}[Falsadores del terminal conjunto]
\label{prop:pdfv2-terminal-falsadores}
Cualquiera de las operaciones siguientes sustituye el objeto construido y
bloquea la conclusi\'on del
\cref{thm:pdfv2-terminal-conjunto}:
\begin{enumerate}
 \item enviar un cilindro a una sola historia cuando admite varias;
 \item elegir una cola del campo terminal \(\mathsf T\) y eliminar las
       restantes;
 \item reemplazar \(\Gamma_9\) por una funci\'on sin demostrar unicidad de
       cada destino y de su testigo;
 \item linealizar una correspondencia ignorando multiplicidades;
 \item identificar retorno de fase con retorno de estado;
 \item conservar la entrada intacta como primera coordenada y usar su
       recuperabilidad como supuesta generaci\'on;
 \item omitir cualquiera de los trece campos de
       \eqref{eq:pdfv2-terminal-registro-trece};
 \item permitir que dos caracter\'isticas intr\'insecas usen historias,
       carries, fronteras o terminales diferentes;
 \item sustituir
       \eqref{eq:pdfv2-terminal-naturalidad-unica} por cinco afirmaciones
       independientes;
 \item afirmar compatibilidad operatorial de \(\Gamma_9\) sin conservar el
       espacio de testigos y sus multiplicidades;
 \item intercambiar, sin prueba, formaci\'on de im\'agenes y paso al l\'imite;
 \item presentar una lista de trece nombres sin las acciones
       \eqref{eq:pdfv2-terminal-actualizacion-unica} que los producen.
 \item omitir la telescopía polar, la fila de pérdidas, la naturalidad
       adjunta, el módulo de todas las continuaciones o la reducción
       Smith--Bockstein acreditadas en
       \cref{prop:pdfv2-terminal-cinco-desarrollos-integrales};
 \item convertir cualquiera de esos cinco desarrollos en una rama autónoma
       o usar una aplicación clásica posterior para seleccionar su dominio.
\end{enumerate}
\end{proposition}

\begin{proof}
Los cinco primeros actos contraen la multisecci\'on o la correspondencia. Los
cuatro siguientes rompen el soporte y la naturalidad comunes. Los tres
\'ultimos eliminan la prueba generativa o alteran el paso al l\'imite. En
cada caso, el objeto resultante deja de satisfacer al menos una de
\eqref{eq:pdfv2-terminal-multivaluado},
\eqref{eq:pdfv2-terminal-span-terminal},
\eqref{eq:pdfv2-terminal-naturalidad-unica} o
\eqref{eq:pdfv2-terminal-recuperacion-compatible}; por tanto, la conclusi\'on
no se transfiere.
\end{proof}

\begin{statusbox}
La organización del terminal completo, la naturalidad única y el paso al
límite se establecen conjuntamente. La multisección superviviente, los
acarreos, el registro, la frontera, la ruta y la conexión nonádica son sus
antecedentes estructurales.\quad
La linealización sobre módulos libres está demostrada por
\eqref{eq:pdfv2-terminal-linealizacion-span} y sus cuadrados de
Beck--Chevalley. Una promoción posterior a un operador acotado sobre una
completación analítica exige además declarar esa norma y probar la cota; la
correspondencia \eqref{eq:pdfv2-terminal-span-limite} no depende de esa
interfaz.
\end{statusbox}
 \end{HMTScientificOwnerSubordinated}

\HMTDeferredContinuumConsequences
 
\chapter{Medida proyectiva y conservación evolutiva de la unidad}
\label{chap:83-06}

\section{Medida proyectiva y conservación evolutiva de la unidad}
\label{p12-83-c6-s1:chap:medida-proyectiva}

La topología de cilindros permite hablar de proximidad; una teoría de caminos
necesita además pesos. La conservación evolutiva de la unidad adquiere aquí
una forma probabilística exacta: cada cilindro reparte su peso entre sus
extensiones y la suma total permanece igual a uno. La unidad no queda fija en
una celda; se transporta a través del refinamiento.

\subsection{Pesos consistentes sobre niveles finitos}

Sea \(\mu_n\) una probabilidad sobre el conjunto finito
\(\mathcal S_n\). La familia es \emph{proyectivamente consistente} cuando
\begin{equation}
 \mu_n(a)
 =\sum_{\substack{b\in\mathcal S_{n+1}\\
                   \rho_{n+1,n}(b)=a}}
   \mu_{n+1}(b)
 \qquad(a\in\mathcal S_n).
 \label{p12-83-c6-s1:eq:consistencia-pesos}
\end{equation}
Esta igualdad expresa la conservación local de la unidad: el peso de una
historia truncada es la suma de los pesos de todos sus hijos.

\begin{theorem}[Medida genealógica sobre cilindros]
\label{p12-83-c6-s1:thm:medida-genealogica}
Toda familia de probabilidades finitas que satisface
\eqref{p12-83-c6-s1:eq:consistencia-pesos} determina una única probabilidad de Borel
\(\mu_\infty\) sobre \(\Omega_\infty\) tal que
\begin{equation}
 \mu_\infty([a]_n)=\mu_n(a).
 \label{p12-83-c6-s1:eq:medida-cilindro}
\end{equation}
\end{theorem}

\begin{proof}
\[
 \mathscr A
 :=\{\text{uniones finitas de cilindros}\}
\]
es un álgebra: la intersección de dos cilindros es vacía o un cilindro del
nivel más profundo. Todo elemento de \(\mathscr A\) puede refinarse a una
unión disjunta de cilindros de un mismo nivel. Si
\(A=\bigsqcup_{a\in I}[a]_n\), ponemos
\[
 \mu_0(A)=\sum_{a\in I}\mu_n(a).
\]
Iterar \eqref{p12-83-c6-s1:eq:consistencia-pesos} demuestra que el valor no depende del
nivel común elegido.

Si un cilindro \([a]_n\) es vacío, la ramificación finita y el lema de König
\cite{Koenig1936}
implican que existe una profundidad sin descendientes de \(a\). La
consistencia iterada da entonces \(\mu_n(a)=0\). Por ello \(\mu_0\) está bien
definida incluso para representaciones que contienen cilindros vacíos y es
finitamente aditiva.

Los cilindros son abiertos y cerrados. Si una sucesión decreciente de
elementos de \(\mathscr A\) tuviera intersección vacía, la compacidad
implicaría que alguno de ellos ya es vacío; por ello \(\mu_0\) es continua en
el vacío y constituye una premedida numerablemente aditiva. El teorema de extensión de
Carathéodory ---equivalentemente, la extensión de Kolmogórov para este sistema
proyectivo finito--- la prolonga a la \(\sigma\)-álgebra generada por los
cilindros \cite{Caratheodory1914,Kolmogorov1933}. Ésta es la Boreliana de
\(\Omega_\infty\), y la finitud de la
medida garantiza unicidad.
\end{proof}

La medida puede proceder de transiciones locales. Si
\(P_n(b\mid a)\ge0\) satisface
\[
 \sum_{\rho_{n+1,n}(b)=a}P_n(b\mid a)=1,
\]
entonces
\(\mu_{n+1}(b)=\mu_n(a)P_n(b\mid a)\) es consistente. La uniformidad es una
elección posible, no una obligación. Los pesos pueden depender de hoja,
acarreo, orientación, un coste de acción no negativo o un registro siempre
que preserven la suma. Por ejemplo,
\[
 P_n(b\mid a)
 =\frac{e^{-\beta a_n(b\mid a)}}
        {\sum_{b'\mapsto a}e^{-\beta a_n(b'\mid a)}},
 \qquad a_n\geq0,
\]
define pesos positivos normalizados. La fase ortogonal de acción pertenece al
kernel del \cref{chap:integral-genealogica}; no se inserta en estas
probabilidades positivas.

\subsection{Desintegración de la medida genealógica sobre fibras de igual valor arquimediano}

La evaluación arquimediana transporta la medida a \(K\):
\begin{equation}
 \nu:=\Val_*\mu_\infty,
 \qquad
 \nu(B)=\mu_\infty(\Val^{-1}(B)).
 \label{p12-83-c6-s1:eq:pushforward-arquimediano}
\end{equation}
La probabilidad visible de una región suma el peso de todas las genealogías que
la publican. Cuando dos secciones poseen el mismo valor, sus contribuciones se
agregan en \(\nu\), aunque permanezcan separadas en \(\mu_\infty\).

\begin{proposition}[Desintegración por valores]
Si \(K\) es un espacio estándar de Borel, existe una familia de medidas
condicionadas \(\{\mu_y\}_{y\in K}\), definida para \(\nu\)-casi todo \(y\),
tal que \(\mu_y\) está soportada en \(\mathfrak G_y=\Val^{-1}(y)\) y
\[
 \int_{\Omega_\infty}F\dd\mu_\infty
 =\int_K\left(\int_{\mathfrak G_y}F\dd\mu_y\right)\dd\nu(y)
\]
para toda función integrable \(F\).
\end{proposition}

Esta fórmula separa valor y genealogía en el nivel de la medida. La física que
dependa sólo del valor integra primero toda la fibra; una observable de memoria
puede distinguir distribuciones \(\mu_y\) diferentes sobre la misma
coordenada real.

\subsection{Condicionamiento como poda de futuros}

Sea \(A\subseteq\Omega_\infty\) un suceso observable con
\(\mu_\infty(A)>0\). La medida condicionada es
\begin{equation}
 \mu_\infty(B\mid A)
 =\frac{\mu_\infty(B\cap A)}{\mu_\infty(A)}.
 \label{p12-83-c6-s1:eq:condicionamiento-genealogico}
\end{equation}
Cuando \(A\) es la lectura de un registro, su intersección con cada cilindro
elimina las extensiones incompatibles. La normalización no crea nuevas ramas:
redistribuye la unidad entre las que han sobrevivido.

\begin{proposition}[Actualización compatible]
Si \(A\) es unión de cilindros a nivel \(m\), entonces para todo \(n\ge m\)
la distribución condicionada sobre \(\mathcal S_n\) se obtiene suprimiendo los
estados cuyos truncamientos no pertenecen a \(A\) y renormalizando los pesos
restantes. Las distribuciones condicionadas siguen satisfaciendo
\eqref{p12-83-c6-s1:eq:consistencia-pesos}.
\end{proposition}

\begin{proof}
La identidad es la regla de Bayes sobre conjuntos finitos. La suma de los
pesos de los hijos condicionados coincide con el peso condicionado del padre,
porque intersección y unión disjunta conmutan.
\end{proof}

\subsection{Conservación combinatoria y completación dimensional}

La misma ley aparece en la completación binomial:
\[
 1=\sum_{r=0}^{d}\binom dr
 \left(\frac9{10}\right)^{d-r}
 \left(\frac1{10}\right)^r.
\]
En dimensión tres,
\[
 1000=729+243+27+1,
\]
de modo que interior, caras nuevas, aristas extremas y vértice final forman
una partición de la unidad decimal. El paso
\(729+270=999\to729+271=1000\) distingue el estado anterior al cierre de la
unidad completada. En el sistema de caminos, la identidad análoga es
\[
 \mu([a]_n)=\sum_{b\mapsto a}\mu([b]_{n+1}).
\]
La unidad evoluciona al distribuirse por los estratos y se recupera al sumar
la genealogía completa.

\begin{figure}[htbp]
\centering
\begin{tikzpicture}[>=Latex,node distance=12mm and 17mm,
  box/.style={draw=hmtblue,rounded corners,fill=hmtlightblue,align=center,
              minimum width=28mm,minimum height=8mm},
  leaf/.style={draw=hmtgreen,rounded corners,fill=hmtlightgreen,align=center,
               minimum width=22mm,minimum height=7mm}]
\node[box] (a) {$[a]_n$\\peso $\mu_n(a)$};
\node[leaf,below left=of a] (b1) {$[b_1]_{n+1}$\\$\mu_{n+1}(b_1)$};
\node[leaf,below=of a] (b2) {$[b_2]_{n+1}$\\$\mu_{n+1}(b_2)$};
\node[leaf,below right=of a] (b3) {$[b_3]_{n+1}$\\$\mu_{n+1}(b_3)$};
\draw[->] (a)--(b1); \draw[->] (a)--(b2); \draw[->] (a)--(b3);
\node[below=20mm of b2,align=center,text width=.78\textwidth] (eq)
 {$\mu_n(a)=\mu_{n+1}(b_1)+\mu_{n+1}(b_2)+\mu_{n+1}(b_3)$:\\
  la unidad se conserva al refinar la genealogía.};
\end{tikzpicture}
\caption{Conservación proyectiva del peso sobre cilindros.}
\label{p12-83-c6-s1:fig:conservacion-peso-cilindros}
\end{figure}

\begin{statusbox}
La medida proyectiva se construye sobre el sistema inverso HMT.
El teorema de extensión es matemática clásica; la identificación de sus
cilindros con estados TPK es exacta una vez fijados los pesos de transición.
La selección de una medida física concreta pertenece al modelo dinámico y no
se deduce únicamente del censo de ramas.
\end{statusbox}

La medida proyectiva queda caracterizada por positividad, conservación del
peso entre padres e hijos y condicionamiento dentro del suceso observado. Una
realización experimental completa añade la preparación y el protocolo de
muestreo; con esos datos, la misma medida produce probabilidades de lectura
comparables con frecuencias.
 
\chapter{Morfismos arquimediano y excepcional desde un dominio común}
\label{chap:83-17}

\section{Morfismos arquimediano y excepcional desde un dominio enriquecido común}
\label{p12-83-c17-s1:chap:doble}
\index{doble proyección}
\index{Moonshine}

La doble proyección compara dos evaluaciones del dominio enriquecido común
\(X_\infty^{\rm cal}\) ya construido en el
\cref{p12-83-c17-s2:thm:lector-excepcional-global-rev6}. Una produce valores
arquimedianos; la otra, junto con una calibración de su codominio, conserva
la incidencia combinatoria que conduce a códigos, diseños y grupos de
automorfismos. Ambas evaluaciones parten del único estado inicial
\(X_\infty^{\rm cal}\) y preservan la procedencia común en codominios
diferenciados.

\subsection{Dominios de las \(2\) proyecciones}

Sobre \(X_\infty^{\rm cal}\), los cilindros se contraen a un punto y la
frontera admite la prolongación de incidencia. En esta residencia fijamos
\(\mathbf{Arch}:=\mathbb R\). La evaluación arquimediana y la evaluación
excepcional son, respectivamente,
\[
\mathfrak R:X_\infty^{\rm cal}\longrightarrow\mathbf{Arch},
\qquad
\mathfrak E:X_\infty^{\rm cal}\times\mathscr C_{\rm exc}
\longrightarrow\mathbf{Exc}.
\]
Aquí \(\mathscr C_{\rm exc}\) contiene el orden proyectivo, el calibre de
Paley, la dodecada, el alineamiento de incidencia y la cámara reticular
especificados en cada construcción. El primer mapa produce valores; el
segundo, códigos, diseños, retículos y automorfismos. El teorema global
precedente es la fuente de la compatibilidad y permite comparar
invariantes obtenidos por operaciones distintas sin volver a intersectar
dominios definidos por separado. La comparación finita de esta obra se
realiza en la bandera
\(B_K\subset H^-_\alpha\).

\begin{figure}[H]
\centering
\resizebox{\textwidth}{!}{\begin{tikzpicture}[font=\small,
  big/.style={draw,rounded corners=3mm,align=center,minimum width=3.7cm,minimum height=1.25cm},
  arr/.style={-{Stealth[length=2.6mm]},very thick}]
  \node[big,draw=hmtgold,fill=hmtlightgold,font=\bfseries] (x) at (0,0) {estado holonómico integral\\\(\XTPK\)};
  \node[big,draw=hmtblue,fill=hmtlightblue] (arq) at (5,2.6) {evaluación arquimediana\\cilindros \(\to\RR\)};
  \node[big,draw=hmtgreen,fill=hmtlightgreen] (exc) at (5,0) {proyección excepcional\\Witt, 3-vecino y Leech};
  \draw[arr,hmtblue] (x)--node[above left]{\(\mathfrak R\)} (arq);
  \draw[arr,hmtgreen] (x)--node[above]{\(\mathfrak E\)} (exc);
  \node[big,draw=hmtgray,fill=hmtpaper] (value) at (10,2.6) {valor real\\clase de evaluación};
  \node[big,draw=hmtgray,fill=hmtpaper] (sym) at (10,0) {automorfismo\\de la incidencia};
  \draw[arr,hmtblue] (arq)--(value);
  \draw[arr,hmtgreen] (exc)--(sym);
  \draw[dashed,hmtred,thick,<->] (value.south west) to[bend right=18] node[fill=white,inner sep=1pt]{invariantes comunes} (sym.north west);
\end{tikzpicture}
 }
\caption{Esquema tipado de las dos proyecciones. La evaluación arquimediana
identifica estados de una misma fibra; la evaluación excepcional, relativa
a su calibración, conserva estructura de incidencia y automorfismos.}
\label{p12-83-c17-s1:fig:doble-proyeccion}
\end{figure}

\subsection{Topología del espacio de historias}

La formulación mediante cilindros admite una realización estándar como
límite inverso. Sean \(\mathsf S_n\) los conjuntos finitos de estados
admisibles a profundidad \(n\), dotados de la topología discreta, y sean
\[
 \tau_{n+1,n}:\mathsf S_{n+1}\longrightarrow\mathsf S_n
\]
las aplicaciones de truncamiento. Supondremos en esta sección que son
sobreyectivas y definiremos
\[
 \mathcal X=\varprojlim_n\mathsf S_n.
\]
Si \(x\ne y\), sea \(m(x,y)\) la primera profundidad en la que sus
componentes difieren. La función
\[
 d(x,y)=3^{-m(x,y)},\qquad d(x,x)=0,
\]
es una ultramétrica que induce la topología del límite inverso.

\begin{theorem}[Espacio compacto de historias]
\label{p12-83-c17-s1:thm:compacto-historias}
\label{thm:compacto-historias}
El espacio \(\mathcal X\) es compacto y totalmente desconectado. Cuando
toda historia admite bifurcaciones a profundidades arbitrariamente grandes,
\(\mathcal X\) carece de puntos aislados y es homeomorfo al espacio de
Cantor.
\end{theorem}

\begin{proof}
El producto \(\prod_n\mathsf S_n\) es compacto. Las ecuaciones
\(\tau_{n+1,n}(x_{n+1})=x_n\) definen un subconjunto cerrado, de modo que
\(\mathcal X\) es compacto. Los conjuntos de historias con un prefijo fijo
son simultáneamente abiertos y cerrados y forman una base; esto prueba la
desconexión total. Bajo la hipótesis de bifurcación arbitrariamente profunda,
ningún cilindro contiene una historia aislada. La caracterización de los
espacios compactos, métricos, perfectos y totalmente desconectados concluye
la última afirmación.
\end{proof}

Asociemos a cada \(s\in\mathsf S_n\) un intervalo cerrado \(I_n(s)\), de
modo compatible con el truncamiento, y supongamos que
\[
 \operatorname{diam}I_n(s)\leq C\lambda^n,
 \qquad 0<\lambda<1.
\]
La intersección de los intervalos correspondientes a una historia contiene
entonces un único punto.

\begin{theorem}[Evaluación arquimediana]
\label{p12-83-c17-s1:thm:evaluacion-holder}
\label{thm:evaluacion-holder}
La aplicación
\[
 \operatorname{val}:\mathcal X\longrightarrow\mathbb R,
 \qquad
 \{\operatorname{val}(x)\}=\bigcap_n I_n(x_n),
\]
es continua y Hölder respecto de \(d\). Si
\(x\sim_{\mathrm{Arch}}y\) equivale a
\(\operatorname{val}(x)=\operatorname{val}(y)\), la aplicación inducida
es un homeomorfismo
\[
 \mathcal X/{\sim_{\mathrm{Arch}}}
 \;\cong\;\operatorname{val}(\mathcal X).
\]
\end{theorem}

\begin{proof}
Dos historias con un prefijo común de longitud \(n\) tienen valores en un
mismo intervalo de diámetro no mayor que \(C\lambda^n\). Como
\(d(x,y)=3^{-n}\), esta estimación es una cota de Hölder. La aplicación
inducida sobre el cociente es una biyección continua de un compacto sobre
un espacio de Hausdorff y, por tanto, un homeomorfismo.
\end{proof}

El teorema formaliza la proyección arquimediana: el valor retiene la clase
de evaluación y elimina la distinción entre historias situadas en una misma
fibra. Su imagen es \(\operatorname{val}(\mathcal X)\); identificarla con
la totalidad de \(\mathbb R\) exige demostrar además la sobreyectividad.

\subsection{Separación por las \(2\) proyecciones}

\begin{theorem}[Separación conjunta]
\label{p12-83-c17-s1:thm:separacion-conjunta}
\label{thm:separacion-conjunta}
Sean \(P_1:\mathcal X\to Y_1\) y \(P_2:\mathcal X\to Y_2\) aplicaciones
continuas con codominios de Hausdorff. Si
\[
 x\ne y\quad\Longrightarrow\quad
 P_1(x)\ne P_1(y)\ \text{o}\ P_2(x)\ne P_2(y),
\]
la aplicación diagonal
\[
 (P_1,P_2):\mathcal X\longrightarrow Y_1\times Y_2
\]
es un encaje topológico.
\end{theorem}

\begin{proof}
La hipótesis expresa exactamente la inyectividad de la aplicación diagonal.
Una aplicación continua e inyectiva de un compacto a un espacio de Hausdorff
es un homeomorfismo sobre su imagen.
\end{proof}

Fijado un calibre declarado \(c\in\mathscr C_{\rm exc}\), sea
\[
 \mathfrak E_c:X_\infty^{\rm cal}\longrightarrow\mathbf{Exc},
 \qquad
 \mathfrak E_c(x):=\mathfrak E(x,c).
\]
Aplicado a HMT, el teorema toma
\[
 P_1=\mathfrak R,\qquad P_2=\mathfrak E_c
\]
sobre el dominio común, o bien sobre el cociente por los cambios de
coordenadas que se hayan declarado. La afirmación global de separación
requiere verificar que ninguna pareja de historias inequivalentes permanezca
en una misma fibra de ambas aplicaciones.

En un nivel finito \(S\), esta condición es decidible. Dados
\(p:S\to A\), \(q:S\to B\) y una relación de equivalencia de coordenadas,
se enumeran las parejas \(x\not\sim y\) que cumplen simultáneamente
\[
 p(x)=p(y),\qquad q(x)=q(y).
\]
La ausencia de tales parejas equivale a la inyectividad conjunta en ese
nivel. El cálculo debe efectuar las dos aplicaciones por separado: cuando
\(q=f\circ p\), la segunda posee las mismas fibras que la primera y no
recupera información eliminada por \(p\).

En el nivel desarrollado aquí, la igualdad
\[
 B_K=H_e\cap P_\pi
\]
constituye un testigo finito de compatibilidad entre las dos ramas. No
establece por sí sola la separación conjunta a todas las profundidades; esa
afirmación corresponde exactamente a la hipótesis del
\cref{p12-83-c17-s1:thm:separacion-conjunta}.

\subsection{Cociente arquimediano por las fibras de evaluación}

Una sección HMT contiene bloque, orientación, acarreo, etiqueta de observable
y coordenada de extensión. La
evaluación \(\val\) contrae sus cilindros a un punto. Si dos secciones
\(x,y\) satisfacen
\[
\val(x)=\val(y),
\]
la evaluación arquimediana las identifica aunque su transporte sea distinto. En
este sentido preciso,
\[
\boxed{\text{un valor real es una clase de equivalencia de la evaluación}.}
\]
En cualquier subdominio en el que suma, producto y evaluación desciendan de
manera compatible, la imagen hereda esas operaciones. Esta sección no
construye por ese hecho todo \(\RR\), ni prueba la sobreyectividad sobre
\(\RR\).

Una identificación global con \(\RR\) requeriría la existencia del límite, la determinación
de su imagen, el descenso de la suma y el producto y la caracterización de las
fibras. Las ventanas finitas y el teorema condicionado de límite inverso
proporcionan la estructura; la prolongación a toda profundidad requiere las
hipótesis de supervivencia unitaria y compatibilidad.

\subsection{Invariante común con el sistema de Witt}

El vector \(K\) define el soporte superior
\[
B_K=\{m:K_m\ge729\}=\{4,6,9,10\}.
\]
En el atlas de Witt aparecen los soportes
\[
P_\pi=\{2,4,5,6,7,8,9,10,11\},
\]
\[
H_e=\{1,3,4,6,9,10\},
\qquad
H_\varphi=\{1,2,3,4,7,9\}.
\]
El soporte de los acarreos negativos de \(\alpha\) es
\[
H_\alpha^-=\{4,5,6,7,9,10\}.
\]
Se verifican las identidades de soportes
\[
\boxed{
B_K=H_e\cap P_\pi=H_e\cap H_\alpha^-,}
\]
y
\[
\boxed{H_\alpha^-=B_K\sqcup\{5,7\}.}
\]
La polaridad de Witt envía los cuatro puntos de \(B_K\) a los pares
\[
\{1,3\},\quad\{2,11\},\quad\{8,12\},\quad\{5,7\}.
\]
Tres reciben la etiqueta positiva; el último completa el soporte negativo.
Estas identidades son exactas en el calibre declarado, pero son controles
condicionados: el vector, el umbral, las palabras y el calibre de Paley no son
entradas estadísticamente independientes. La igualdad se obtiene por
evaluación directa y no constituye una estimación probabilística.

\subsection{Diseño residual sobre una dodecada y elevación de hexadas}
\label{p12-83-c17-s1:sec:w12-w24-elevacion}

La relación entre los dos diseños de Witt se obtiene dentro del código
binario extendido de Golay y no mediante una identificación de cardinalidades.
Sea
\[
\mathcal G_{24}\subset\mathbb F_2^{24}
\]
el código binario extendido de Golay, y denote \(\mathcal O_{24}\) la familia
de soportes de sus palabras de peso ocho. Es un resultado clásico que
\((\{1,\ldots,24\},\mathcal O_{24})\) es el sistema de Steiner
\(W_{24}=S(5,8,24)\) \cite{MacWilliamsSloane1977,ConwaySloane1999}.
Fijemos una palabra de peso doce y escribamos \(D\) para su soporte. Definimos
\begin{equation}
 \mathcal H(D)
 :=\{O\cap D:O\in\mathcal O_{24},\ |O\cap D|=6\}.
 \label{p12-83-c17-s1:eq:hexadas-residuales-dodecada}
\end{equation}
La familia \(\mathcal H(D)\) es intrínseca al par
\((\mathcal G_{24},D)\): no requiere elegir coordenadas en \(D\).

\begin{theorem}[Diseño residual de una dodecada]
\label{p12-83-c17-s1:thm:dodecada-w12}
El sistema de incidencia
\[
 (D,\mathcal H(D))
\]
es un sistema de Steiner \(S(5,6,12)\). En particular,
\(|\mathcal H(D)|=132\), y todo subconjunto de cinco puntos de \(D\) está
contenido en una única hexada de \(\mathcal H(D)\).
\end{theorem}

\begin{proof}
Sea \(A\subset D\), con \(|A|=5\). Como las octadas forman
\(S(5,8,24)\), existe una única octada \(O_A\) que contiene \(A\). Si
\(j=|O_A\cap D|\), la suma binaria de las palabras características de
\(O_A\) y \(D\) pertenece a \(\mathcal G_{24}\) y tiene peso
\[
 \operatorname{wt}(\mathbf1_{O_A}+\mathbf1_D)=8+12-2j=20-2j.
\]
Los pesos del código binario extendido de Golay son
\(0,8,12,16,24\). Puesto que \(j\leq8\), se sigue que
\(j\in\{2,4,6\}\). La inclusión \(A\subset O_A\cap D\) da \(j\geq5\),
luego \(j=6\). Por tanto \(O_A\cap D\in\mathcal H(D)\) es una hexada que
contiene \(A\).

Si dos miembros de \(\mathcal H(D)\) contuvieran \(A\), sus octadas de
origen contendrían el mismo conjunto de cinco puntos; la propiedad de
Steiner de \(W_{24}\) las haría coincidir. Esto prueba la unicidad. El
número de bloques se obtiene contando incidencias con subconjuntos de cinco
puntos:
\[
 |\mathcal H(D)|
 =\frac{\binom{12}{5}}{\binom65}=132.\qedhere
\]
\end{proof}

\begin{theorem}[Elevación única de una hexada]
\label{p12-83-c17-s1:thm:elevacion-unica-hexada}
Para cada \(H\in\mathcal H(D)\) existe una única octada
\(O_H\in\mathcal O_{24}\) tal que
\[
 O_H\cap D=H.
\]
En consecuencia, la aplicación
\[
 \iota_D:\mathcal H(D)\longrightarrow\mathcal O_{24},
 \qquad H\longmapsto O_H,
\]
está bien definida y es inyectiva.
\end{theorem}

\begin{proof}
La existencia forma parte de la definición de \(\mathcal H(D)\). Si dos
octadas \(O_1,O_2\) tuvieran intersección \(H\) con \(D\), ambas
contendrían cualquier subconjunto de cinco puntos de \(H\). La unicidad de
la octada incidente con esos cinco puntos implica \(O_1=O_2\). La
inyectividad se sigue al intersectar de nuevo con \(D\).
\end{proof}

La elevación tiene una ley local que será utilizada para comparar las cinco
regiones de \(\pi\). En un sistema \(S(t,k,v)\), el número de bloques que
contienen un subconjunto fijo de \(s\leq t\) puntos es
\[
 \lambda_s=\frac{\binom{v-s}{t-s}}{\binom{k-s}{t-s}}.
\]
Por consiguiente,
\begin{equation}
 \lambda_4(W_{12})
 =\frac{\binom81}{\binom21}=4,
 \qquad
 \lambda_4(W_{24})
 =\frac{\binom{20}1}{\binom41}=5.
 \label{p12-83-c17-s1:eq:lambda4-w12-w24}
\end{equation}

\begin{corollary}[Descomposición \(4+1\) sobre una cara]
\label{p12-83-c17-s1:cor:cuatro-mas-uno-w24}
Sea \(B\subset D\), con \(|B|=4\). Las cinco octadas de \(W_{24}\) que
contienen \(B\) se descomponen de manera única en
\[
 \{O_H:H\in\mathcal H(D),\ B\subset H\}
 \;\sqcup\;\{O_B^\perp\}.
\]
El primer conjunto contiene cuatro octadas y satisface
\(|O_H\cap D|=6\); la octada restante satisface
\[
 O_B^\perp\cap D=B.
\]
\end{corollary}

\begin{proof}
La primera igualdad de \eqref{p12-83-c17-s1:eq:lambda4-w12-w24} proporciona cuatro hexadas
residuales que contienen \(B\), y el
\cref{p12-83-c17-s1:thm:elevacion-unica-hexada} proporciona cuatro octadas distintas. La
segunda igualdad de \eqref{p12-83-c17-s1:eq:lambda4-w12-w24} muestra que existe exactamente
una quinta octada \(O_B^\perp\).

Como \(B\subset O_B^\perp\cap D\), el espectro de intersecciones empleado en
la prueba del \cref{p12-83-c17-s1:thm:dodecada-w12} deja sólo las posibilidades cuatro y
seis. Una intersección de tamaño seis sería una quinta hexada residual que
contiene \(B\), contradiciendo \(\lambda_4(W_{12})=4\). Por tanto la
intersección tiene tamaño cuatro y coincide con \(B\).
\end{proof}

\subsubsection{Alineamiento con las \(12\) posiciones HMT}

El diseño residual está definido sobre los puntos de \(D\), mientras que el
diseño obtenido del código ternario \(C_W\) está definido sobre las doce
posiciones dodecafásicas. Para no confundir ambos conjuntos, llamaremos
\emph{alineamiento binario HMT} a una biyección
\begin{equation}
 \ell:\{1,\ldots,12\}_{\rm HMT}\xrightarrow{\sim}D
 \quad\text{tal que}\quad
 \{\ell(H):H\in W_{12}^{\rm HMT}\}=\mathcal H(D).
 \label{p12-83-c17-s1:eq:alineamiento-hmt-w24}
\end{equation}
Aquí \(W_{12}^{\rm HMT}\) denota el diseño de soportes de peso seis
obtenido del código ternario \(C_W\). La unicidad del sistema de Witt
\(S(5,6,12)\) hasta isomorfía garantiza la existencia de \(\ell\). Dos
alineamientos distintos difieren por automorfismos de los diseños de origen y
destino; por ello \(\ell\) es un dato de coordenadas y no un invariante
adicional.
La pareja \((D,\ell)\) es el dato de calibración necesario para transportar
la elevación \(\iota_D\) al sistema de coordenadas HMT\@. Una vez fijada, cada
hexada HMT \(H\) determina sin ambigüedad la octada
\[
 \widehat H:=\iota_D(\ell(H)).
\]
El dodecad y el alineamiento son datos del codominio binario; no se identifican
con una salida decimal del núcleo de transporte.

Para la cara distinguida
\[
 B_K=\{4,6,9,10\},
\]
las cuatro hexadas incidentes son
\[
 B_K\sqcup\{1,3\},\quad
 B_K\sqcup\{2,11\},\quad
 B_K\sqcup\{5,7\},\quad
 B_K\sqcup\{8,12\}.
\]
La tercera coincide con
\(H_\alpha^-=B_K\sqcup\{5,7\}\). El alineamiento \(\ell\) eleva estas
cuatro hexadas a las cuatro octadas residuales sobre \(\ell(B_K)\); el
\cref{p12-83-c17-s1:cor:cuatro-mas-uno-w24} determina además una única octada transversal.

\subsubsection{Correspondencia de incidencia entre la microfibra de \texorpdfstring{\(\pi\)}{pi} y las octadas residuales: multiplicidad \texorpdfstring{\(432+4\cdot144=1008\)}{432+4x144=1008}}

La microfibra exacta de la palabra
\(w_\pi=010211\) contiene cinco regiones decimales distintas:
\[
\begin{array}{c|c|c}
\text{multiplicidad}&U_6& e(U_6)\\
\hline
432&501|614|498|169|272|272&(5,5,5,5,5,5)\\
144&810|923|870|169|272|272&(3,3,9,5,5,5)\\
144&810|983|810|169|272|272&(3,9,3,5,5,5)\\
144&870|923|810|169|272|272&(9,3,3,5,5,5)\\
144&870|923|810|418|674|521&(9,3,3,7,1,7).
\end{array}
\]
Aquí \(e(U_6)\) es la firma multiplicativa recuperada de cada bloque
decimal \(\overline{d_1d_2d_3}\) mediante
\(e=d_2-d_1\pmod {10}\). La primera región es la única de multiplicidad
432 y firma constante. Las cuatro restantes tienen multiplicidad 144; entre
ellas, tres comparten el bloque de retorno \(169|272|272\), mientras que la
cuarta posee el retorno \(418|674|521\). El catálogo distingue las cinco
filas mediante un predicado reproducible: las tres filas de tipo
\(\mathrm{ES}\) forman el triple positivo; entre las dos filas de tipo
\(\mathrm{NO}\), la de multiplicidad máxima y firma constante se denomina
transversal, y la restante se denomina negativa. Relativa a este predicado de
rol, la partición es
\[
 \boxed{1_\perp+3_{++}+1_{--}.}
\]
Los nombres de signo son una convención; lo intrínseco al catálogo es la
partición obtenida de sus columnas, no una polaridad binaria previa.

El alineamiento regional completa el dato \((D,\ell)\) con una biyección
que envía la región transversal a \(O_{\ell(B_K)}^\perp\), la región
negativa a \(\widehat{H_\alpha^-}\) y las tres regiones positivas a las
elevaciones de
\[
 B_K\sqcup\{1,3\},\qquad
 B_K\sqcup\{2,11\},\qquad
 B_K\sqcup\{8,12\}.
\]
La asignación ordenada entre las tres regiones positivas y esas tres hexadas
es una elección de coordenadas. Las dos filas distinguidas y el triple son
reproducibles mediante el predicado de catálogo anterior; su denominación como
negativa, transversal y positivas pertenece al alineamiento. Así, la ley \(4+1\) no identifica por mera
cardinalidad las cinco regiones con cinco octadas: proporciona el espacio de
incidencia, mientras que \((D,\ell)\) y el alineamiento regional especifican
el transporte entre ambos objetos.

\subsection{Prolongación reticular y pantallas dimensionales}

La elevación de hexadas abre dos ramas que no deben confundirse. La primera
conserva la incidencia binaria y prolonga \(W_{12}\) hacia \(W_{24}\). La
segunda conserva el transporte ternario:
\[
C_{12}^{(3)}
\longrightarrow N(A_2^{12})
\longrightarrow\Lambda_{24}
\longrightarrow V^\natural
\longrightarrow\mathbb M
\longrightarrow\text{Moonshine}.
\]
La prueba primaria del pegado, del \(3\)-vecino sin raíces y del corredor
hasta Moonshine reside en el
\cref{chap:bandera-excepcional-acumulativa,thm:corredor-hmt-moonshine}.
No existe una flecha canónica \(W_{24}\to N(A_2^{12})\), ni una acción
directa del Monstruo sobre las palabras TPK: ambas ramas parten del mismo
código calibrado y preservan estructuras distintas.

Las reducciones \(243/11\), \(12\to11\to10\) y \(26\to24\) tampoco son
identidades entre numerales. El \cref{chap:pantallas-dualidad} construye sus
espacios, cocientes y mapas: el cociente perfecto ternario, la pantalla
marcada por \(K\) y la reducción isotrópica del retículo lorentziano. Esta
sección conserva así una sola función: enlazar la incidencia finita ya
demostrada con sus dos prolongaciones, sin repetir las pruebas de referencia.

\subsection{Involución residuo–cociente y conservación de la información del estado}

Para \(n=9Q+r\), definimos la energía de retículo
\[
E_R(r,Q)=\frac{r^2}{R^2}+Q^2R^2.
\]
Entonces
\[
\boxed{E_R(r,Q)=E_{1/R}(Q,r).}
\]
El intercambio \((r,Q,R)\leftrightarrow(Q,r,1/R)\) es una involución exacta
del modelo de retículo: el residuo y el cociente intercambian sus funciones.

\subsection{Compatibilidad finita y prolongación conjunta}

El testigo finito \(B_K\), el patrón de acarreos y el marco de Witt fijan la
compatibilidad local. Su prolongación ya no se deja como un esquema
abstracto: el \cref{p12-83-c17-s2:thm:lector-excepcional-global-rev6} construye bloque a
bloque el lector excepcional con un único calibre inicial transportado, y el
\cref{p12-83-c17-s2:thm:doble-realizacion-global-rev6} demuestra que esa lectura y la
evaluación arquimediana globalizan sobre la misma sección TPK. Truncar la
historia antes de leer equivale en ambos casos a truncar la salida después de
leer.

La evaluación arquimediana identifica historias por su valor límite. La
evaluación excepcional conserva la corriente visible y su incidencia
Paley--Witt; no recupera por ello las memorias ocultas que el emisor visible
no publica. La doble proyección es, por tanto, una construcción conjunta con
dos codominios distintos, no una identificación entre el continuo y la rama
excepcional.

\begin{remark}[Compatibilidad de las dos proyecciones]
El soporte \(B_K\), el patrón de acarreos y el marco de Witt proporcionan los
invariantes comunes. La cadena reticular prolonga esta compatibilidad hasta un
vecino sin raíces y, por clasificación clásica, hasta el retículo de Leech. La
rama binaria hacia \(W_{24}\) permanece separada: su morfismo de incidencia es
\(\iota_D\), definido después de fijar el dodecad y el alineamiento
\((D,\ell)\), y no tiene por codominio un retículo de Niemeier ni el retículo
de Leech.
\end{remark}


\section{Proyección excepcional del estado holonómico integral: incidencia compatible a toda profundidad}
\label{p12-83-c17-s2:chap:lector-excepcional-global-rev6}
\index{lector excepcional}\index{doble proyección}
\index{Paley--Witt}\index{sistema inverso}

El testigo finito de la doble proyección adquiere alcance global cuando la
lectura excepcional se aplica a cada bloque visible y el calibre se
transporta, en lugar de elegirse de nuevo, al aumentar la profundidad. Esta
sección construye esa familia. El dominio sigue siendo la historia TPK
enriquecida; la lectura excepcional es fiel sobre su corriente visible y no
pretende sustituir los campos de memoria que esa corriente no publica.

\subsection{La operación local sobre las 729 palabras}

En la carta de Paley--Witt se fija
\[
A_W=
\begin{pmatrix}
0&1&1&1&1&1\\
1&0&1&2&2&1\\
1&1&0&1&2&2\\
1&2&1&0&1&2\\
1&2&2&1&0&1\\
1&1&2&2&1&0
\end{pmatrix}
\in M_6(\mathbb F_3).
\]
La multiplicación matricial da
\[
A_W^{\mathsf T}=A_W,\qquad A_W^2=-I_6.
\]
Para todo operador conjugado \(A=fA_Wf^{-1}\), definimos
\[
\Gamma_A:\mathbb F_3^6\longrightarrow\mathbb F_3^{12},
\qquad
\Gamma_A(w)=(w,wA).
\]
Su imagen es una copia calibrada \(C_W(A)\) del código ternario de Witt.
La proyección sobre las seis primeras coordenadas es una inversa izquierda:
\[
\operatorname{pr}_1\Gamma_A=\operatorname{id}_{\mathbb F_3^6}.
\]
Por tanto, \(\Gamma_A\) es total e inyectiva sobre las \(729\) palabras y
no necesita distinguir de antemano \(\pi\), \(e\), \(\varphi\) ni ninguna
otra salida.

\subsection{Prolongaciones compatibles y unicidad del calibre inicial}

Sea \(X_n\) el conjunto de prefijos TPK admisibles de profundidad \(n\), con
truncamientos
\[
p_{n,m}:X_n\longrightarrow X_m,\qquad m\leq n.
\]
Fijemos un marco excepcional inicial \(f_0\). Como los cambios \(g_j\) forman
parte del transporte almacenado por la historia, un prefijo calibrado es el
par \((x,f_0)\); escribimos
\[
 X_n^{\rm cal}:=X_n\times\{f_0\},\qquad
 p_{n,m}^{\rm cal}(x,f_0):=(p_{n,m}x,f_0),
\]
y
\[
 X_\infty^{\rm cal}
 :=\varprojlim_n(X_n^{\rm cal},p_{n,m}^{\rm cal}).
\]
En lo que sigue abreviamos \(p_{n,m}^{\rm cal}\) por \(p_{n,m}\).
La emisión visible compatible es
\[
\varepsilon_n:X_n^{\rm cal}\longrightarrow(\mathbb F_3^6)^n,
\qquad
\varepsilon_n(x)=(w_0,\ldots,w_{n-1}),
\]
y satisface
\[
\operatorname{pr}_{n,m}\varepsilon_n
=\varepsilon_m p_{n,m}.
\]
Esta aplicación distingue sus tres tipos:
\[
|\mathcal U_6|=468,\qquad
|\operatorname{im}w_6|=243,\qquad
|\mathbb F_3^6|=729.
\]
La emisión \(w_j\) no se identifica con el estado completo que conserva
hoja, cociente, acarreo, orientación, fase y supervivencia.

El calibre excepcional consta de una carta de Witt, una dodecada, un
alineamiento, un origen reticular y una cámara orientada. Se fija una sola
vez en el bloque inicial. Para
\(x\in X_n^{\rm cal}\), el transporte de la historia suministra
\[
 g_j(x)\in\operatorname{GL}_6(\mathbb F_3),
 \qquad 0\le j<n-1,
\]
y \(g_j(x)\) depende sólo del prefijo
\(p_{n,j+1}(x)\). El marco se transporta por
\[
 f_{j+1}(x)=g_j(x)f_j(x),\qquad
 A_j=f_jA_Wf_j^{-1}.
\]
Así, la familia \((A_j)\) queda determinada por la historia y por el calibre
inicial; no es una sucesión de elecciones independientes.

Sea \(\mathcal F_{\rm cal}\subseteq\operatorname{GL}_6(\mathbb F_3)\) el
conjunto de marcos alcanzables desde \(f_0\). La emisión que conserva el
calibre transportado es
\[
 \widetilde\varepsilon_n:X_n^{\rm cal}
 \longrightarrow(\mathbb F_3^6\times\mathcal F_{\rm cal})^n,
 \qquad
 \widetilde\varepsilon_n(x)
 =\bigl((w_0,f_0),\ldots,(w_{n-1},f_{n-1})\bigr).
\]
La dependencia por prefijos implica
\(\operatorname{pr}_{n,m}\widetilde\varepsilon_n
=\widetilde\varepsilon_m p_{n,m}\). Su proyección sobre las palabras es
\(\varepsilon_n\); el olvido de los marcos no basta para reconstruir la
lectura excepcional.

\subsection{Naturalidad y compatibilidad a toda profundidad}

La operación local es covariante con la convención de vectores fila. Para
todo cambio de coordenadas \(f\in\operatorname{GL}_6(\mathbb F_3)\),
\[
A^f=fAf^{-1}
\quad\Longrightarrow\quad
\Gamma_{A^f}(wf^{-1})
=\Gamma_A(w)(f^{-1}\oplus f^{-1}).
\]
El cambio de marco transporta simultáneamente la palabra, el operador y el
atlas. No se exige que permanezca dentro del estabilizador de una única copia
etiquetada.

Con el conjunto de marcos alcanzables \(\mathcal F_{\rm cal}\) ya definido,
construimos el atlas calibrado fijo
\[
 \mathscr C_W^{\rm cal}
 :=
 \left\{
 \bigl(f,\Gamma_{fA_Wf^{-1}}(w)\bigr):
 f\in\mathcal F_{\rm cal},\ w\in\mathbb F_3^6
 \right\}.
\]
La etiqueta \(f\) conserva la copia calibrada en la que vive cada palabra.
Definimos
\[
Q_n(x)=
\bigl(
\bigl(f_0,\Gamma_{A_0}(w_0)\bigr),\ldots,
\bigl(f_{n-1},\Gamma_{A_{n-1}}(w_{n-1})\bigr)
\bigr)\in\bigl(\mathscr C_W^{\rm cal}\bigr)^n.
\]
Sea
\[
 r_{n,m}:
 \bigl(\mathscr C_W^{\rm cal}\bigr)^n
 \longrightarrow
 \bigl(\mathscr C_W^{\rm cal}\bigr)^m
\]
el truncamiento de los últimos \(n-m\) bloques.

\begin{theorem}[Compatibilidad proyectiva de la incidencia excepcional]
\label{p12-83-c17-s2:thm:lector-excepcional-global-rev6}
\label{thm:lector-excepcional-global-rev6}
Para todo \(m\leq n\),
\[
\boxed{r_{n,m}\circ Q_n=Q_m\circ p_{n,m}.}
\]
En consecuencia existe una única aplicación
\[
\boxed{
Q_\infty:X_\infty^{\rm cal}
\longrightarrow
\bigl(\mathscr C_W^{\rm cal}\bigr)^{\mathbb N}}
\]
cuyas proyecciones finitas son los \(Q_n\). La aplicación es fiel sobre la
corriente visible.
\end{theorem}

\begin{proof}
Escribamos
\[
\varepsilon_n(x)
=(w_0,\ldots,w_{m-1},w_m,\ldots,w_{n-1}).
\]
La compatibilidad de \(\varepsilon\) muestra que la corriente visible de
\(p_{n,m}(x)\) es \((w_0,\ldots,w_{m-1})\). Para \(j<m\), el marco
\(f_j\) depende sólo del calibre inicial y de los cambios anteriores a
\(j\), todos contenidos en ese prefijo. Los primeros \(m\) operadores
\(A_j\) son, por tanto, los mismos antes y después de truncar. Aplicar
\(\Gamma_{A_j}\) bloque a bloque da la identidad del enunciado. La propiedad
universal del límite inverso produce \(Q_\infty\). Finalmente,
\(\operatorname{pr}_1\Gamma_{A_j}=I\) recupera cada \(w_j\), lo que prueba
la fidelidad sobre la corriente visible.
\end{proof}

La última afirmación conserva su tipo. Dos historias con la misma corriente
\((w_j)_j\) y distinta memoria oculta pueden tener la misma imagen
excepcional. El lector global preserva exactamente lo que publica cada bloque
visible y la incidencia que se añade a ese bloque.

La compatibilidad de las emisiones enriquecidas induce una única aplicación
\[
 \widetilde\varepsilon_\infty:X_\infty^{\rm cal}
 \longrightarrow
 (\mathbb F_3^6\times\mathcal F_{\rm cal})^{\mathbb N}.
\]
Sean \(\operatorname{pr}_w\) la proyección sobre la corriente de palabras y
\[
 \widetilde\Gamma_\infty
 \bigl((w_j,f_j)_{j\ge0}\bigr)
 :=\bigl(
   (f_j,\Gamma_{f_jA_Wf_j^{-1}}(w_j))
 \bigr)_{j\ge0}.
\]
Entonces
\(Q_\infty=\widetilde\Gamma_\infty
\circ\widetilde\varepsilon_\infty\), con dominios y codominios declarados.

\subsection{Bifurcación de un prefijo superviviente en valor arquimediano e incidencia excepcional}

Fijado \(m\in\mathbb Z\), denotemos por
\[
 \operatorname{flat}:(\mathbb F_3^6)^{\mathbb N}
 \longrightarrow\mathbb F_3^{\mathbb N}
\]
la concatenación de los bloques en el orden coordenado declarado, y por
\[
 \operatorname{ev}_3((d_k)_{k\ge1})
 :=\sum_{k\ge1}\frac{d_k}{3^k}\in[0,1],
 \qquad
 \tau_m(t):=m+t,
\]
la evaluación ternaria y la traslación por la parte entera. El prefijo visible
\((d_1,\ldots,d_{6n})\) determina el cilindro
\[
I_n(x)=
\left[
m+\sum_{k=1}^{6n}\frac{d_k}{3^k},
m+\sum_{k=1}^{6n}\frac{d_k}{3^k}+3^{-6n}
\right].
\]
Una prolongación restringe el cilindro y
\(\operatorname{diam}I_n=3^{-6n}\to0\). Por tanto, toda historia infinita
compatible determina un único punto real
\(\mathcal P_{\rm arq}(x)\).

\begin{theorem}[Doble realización global]
\label{p12-83-c17-s2:thm:doble-realizacion-global-rev6}
\label{thm:doble-realizacion-global-rev6}
Toda sección compatible \(x=(x_n)_n\in X_\infty^{\rm cal}\) determina,
desde los mismos prefijos, dos salidas compatibles:
\[
\mathcal P_{\rm arq}(x)\in\mathbb R,
\qquad
\mathcal P_{\rm exc}(x)=Q_\infty(x)
\in\bigl(\mathscr C_W^{\rm cal}\bigr)^{\mathbb N}.
\]
Truncar antes de leer equivale a truncar cada lectura después de leer. La
primera salida identifica historias por su valor límite; la segunda conserva
cada bloque visible y adjunta su rotación de Paley--Witt.
\end{theorem}

\begin{proof}
El encaje de los cilindros y la convergencia de sus diámetros prueban la
primera afirmación. El
\cref{p12-83-c17-s2:thm:lector-excepcional-global-rev6} prueba la segunda y la
compatibilidad con truncamientos. Ambas factorizan por la misma emisión:
\[
\mathcal P_{\rm arq}
=\tau_m\circ\operatorname{ev}_3\circ\operatorname{flat}
 \circ\operatorname{pr}_w
 \circ\widetilde\varepsilon_\infty,
\qquad
\mathcal P_{\rm exc}
=\widetilde\Gamma_\infty\circ\widetilde\varepsilon_\infty.
\]
\end{proof}

La doble proyección queda así construida sobre historias admisibles y un
calibre inicial transportado. El continuo es la evaluación arquimediana de
la sección; la rama excepcional es la realización compatible de su corriente
visible. Desde \(C_W\) se abren, con los alineamientos ya declarados, las
ramas hacia \(W_{12}\), \(W_{24}\) y hacia
\(N(A_2^{12})\), Leech, \(V^\natural\) y el Monstruo. Esas ramas no convierten
las cifras en un conjunto sobre el que actúe directamente el Monstruo:
transportan la incidencia que la evaluación por valor deja de ver.

\subsection{Acoplamiento postgenerativo con el cierre operatorial de Euler}
\label{p12-83-c17-s2:sec:acoplamiento-euler-doble-proyeccion-rev3}

Sea \(x_\pi\in X_\infty^{\rm cal}\) la sección superviviente del canal de
\(\pi\). La doble realización del
\cref{p12-83-c17-s2:thm:doble-realizacion-global-rev6} produce
\[
 \mathcal P_{\rm arq}(x_\pi)\in\mathbb R,
 \qquad
 \mathcal P_{\rm exc}(x_\pi)=Q_\infty(x_\pi).
\]
La primera rama publica el ángulo generado; la segunda conserva la corriente
visible y realiza en cada bloque la relación cuadrática de Paley--Witt.

\begin{theorem}[Acoplamiento postgenerativo Euler--doble proyección]
\label{p12-83-c17-s2:thm:acoplamiento-euler-doble-proyeccion-rev3}
Supóngase que la construcción coinductiva del canal de \(\pi\) determina
\[
 \mathcal P_{\rm arq}(x_\pi)=\pi,
\]
y que la rama excepcional utiliza el calibre transportado con
\(A_W^2=-I_6\). Sea \(J_{\mathrm e}\) la realización real de la fuente
integral del \cref{thm:dos-realizaciones-fuente-cuadratica-rev3}. Entonces
\begin{equation}
 \boxed{
 \operatorname{Exp}\!\bigl(
   \mathcal P_{\rm arq}(x_\pi)J_{\mathrm e}
 \bigr)+I=0.}
 \label{p12-83-c17-s2:eq:acoplamiento-euler-doble-proyeccion-rev3}
\end{equation}
\end{theorem}

\begin{proof}
Por el \cref{thm:dos-realizaciones-fuente-cuadratica-rev3}, \(A_W\) y
\(J_{\mathrm e}\) realizan la relación \(u^2+1=0\) en categorías distintas.
En la rama real,
\[
 \operatorname{Exp}(\theta J_{\mathrm e})
 =\cos\theta\,I+\sin\theta\,J_{\mathrm e}.
\]
La sustitución de la salida previamente generada
\(\theta=\mathcal P_{\rm arq}(x_\pi)=\pi\) produce \(-I\), y la
suma de \(I\) da la identidad anunciada.
\end{proof}

\paragraph{Orden causal y ausencia de circularidad.}
La dependencia demostrada es
\[
 \text{emisor}
 \longrightarrow x_\pi
 \longrightarrow\mathcal P_{\rm arq}(x_\pi)
 \longrightarrow\pi
 \longrightarrow
 \operatorname{Exp}({\pi}J_{\mathrm e})+I.
\]
La rama excepcional fija el tipo cuadrático; no selecciona el ángulo ni sus
cifras. La identidad de Euler reconoce después la compatibilidad de las dos
realizaciones y no interviene como entrada del generador de
\(\pi\).


\section{Globalización natural de las proyecciones arquimediana y excepcional sobre un sistema inverso común}
\label{p12-83-c17-s3:chap:globalizacion-doble-proyeccion}

La tesis de doble proyección adquiere contenido matemático cuando las dos
lecturas se definen sobre un mismo sistema de estados y conmutan con la
restricción de profundidad. La evaluación arquimediana contrae una familia de
cilindros a un valor; la evaluación excepcional conserva relaciones de
incidencia que la primera puede olvidar. No se trata de dos expresiones
numéricas yuxtapuestas, sino de dos transformaciones naturales con dominio
común.

\subsection{Sistemas compatibles a profundidad finita}

Sean \((X_n,p_{n+1,n})\), \((A_n,a_{n+1,n})\) y
\((E_n,e_{n+1,n})\) tres sistemas inversos. El primero contiene historias
TPK; el segundo, cilindros o aproximantes arquimedianos; el tercero, datos de
incidencia excepcional. Supongamos dadas aplicaciones
\[
 R_n:X_n\longrightarrow A_n,
 \qquad
 Q_n:X_n\times C_n\longrightarrow E_n,
\]
donde \(C_n\) reúne los datos de coordenadas del lector excepcional. Se exige
que los calibres formen también una familia compatible y que, para todo
\(n\), conmuten los diagramas
\[
 a_{n+1,n}\circ R_{n+1}=R_n\circ p_{n+1,n},
\]
\[
 e_{n+1,n}\circ Q_{n+1}
 =Q_n\circ(p_{n+1,n}\times c_{n+1,n}).
\]

\begin{theorem}[Globalización de las dos evaluaciones]
\label{p12-83-c17-s3:thm:globalizacion-evaluaciones}
Bajo las compatibilidades anteriores existen aplicaciones canónicas
\[
 R_\infty:\varprojlim X_n\longrightarrow\varprojlim A_n,
 \qquad
 Q_\infty:\varprojlim X_n\times\varprojlim C_n
            \longrightarrow\varprojlim E_n.
\]
Su aplicación diagonal
\[
 (R_\infty,Q_\infty):X_\infty\times C_\infty
 \longrightarrow A_\infty\times E_\infty
\]
es única entre las aplicaciones cuyas proyecciones coordenadas son
\(R_n\) y \(Q_n\).
\end{theorem}

\begin{proof}
Sea \(x=(x_n)_n\in\varprojlim X_n\). La primera ecuación de conmutatividad
implica
\[
 a_{n+1,n}(R_{n+1}(x_{n+1}))=R_n(x_n),
\]
por lo que \((R_n(x_n))_n\) pertenece a \(\varprojlim A_n\). Definimos
\(R_\infty(x)\) mediante esa familia. El mismo argumento, aplicado a
\((x_n,c_n)_n\), define \(Q_\infty\). La unicidad se sigue porque un elemento
de un límite inverso queda determinado por todas sus coordenadas.
\end{proof}

El teorema identifica la prueba concreta que requiere una realización HMT:
no basta disponer de un resultado arquimediano en una profundidad y de un
código en otra; los dos deben proceder de estados compatibles y respetar los
mapas de truncamiento.

\subsection{Contracción arquimediana de cilindros compatibles y realización de la recta real}

A cada \(x_n\in X_n\) se asocia un intervalo compacto no vacío
\(I_n(x_n)\subset\mathbb R\). Si
\[
 p_{n+1,n}(x_{n+1})=x_n
 \quad\Longrightarrow\quad
 I_{n+1}(x_{n+1})\subseteq I_n(x_n)
\]
y existe \(\varepsilon_n\to0\) con
\(\operatorname{diam}I_n(x_n)\leq\varepsilon_n\), la intersección
\[
 \bigcap_{n\geq0}I_n(x_n)
\]
contiene un único real. Así se obtiene
\[
 \operatorname{ev}_{\rm arq}:X_\infty^{\rm arq}\longrightarrow\mathbb R.
\]
La frase ``los cilindros tienden a cero'' significa aquí que converge a cero
su diámetro, no que el número evaluado tienda a cero.

La relación
\[
 x\sim_{\rm arq}y
 \quad\Longleftrightarrow\quad
 \operatorname{ev}_{\rm arq}(x)=\operatorname{ev}_{\rm arq}(y)
\]
describe exactamente la información eliminada por el lector. El valor real
es una clase de evaluación; la historia conserva, además, orientación,
acarreo, hoja y firma de frontera.

\subsection{Proyección excepcional por conservación de incidencia y frontera}

La segunda evaluación no toma valores en \(\mathbb R\). Su codominio es una
categoría de configuraciones de incidencia. En la realización finita vigente,
el dato de coordenadas contiene:
\[
 \begin{aligned}
 C_{\rm exc}=(\,&\text{orden proyectivo},\text{calibre de Paley},
 \text{dodecada},\\
 &\text{alineamiento},\text{origen reticular},\text{cámara}).
 \end{aligned}
\]
Después de fijar ese dato, la cadena contiene dos descendencias que deben
permanecer separadas:
\[
 C_W\longrightarrow W_{12}\longrightarrow W_{24},
\]
\[
 C_W\longrightarrow N(A_2^{12})
 \longrightarrow\Lambda_{24}
 \longrightarrow V^\natural
 \longrightarrow\mathbb M.
\]
La primera transporta diseños; la segunda conduce, mediante resultados
clásicos posteriores, al retículo de Leech, el álgebra de operadores de
vértice y el grupo Monstruo. La holonomía central y el corredor
Schläfli--\(E_6\) forman una extracción paralela de frontera; no se obtienen
reemplazando una de las dos descendencias anteriores.

\subsection{Compatibilidad finita del dominio común}

El sello dodecafásico
\[
 K=(234,543,140,729,659,824,621,58,914,794,146,601)
\]
determina el soporte
\[
 B_K=\{m:K_m\geq729\}=\{4,6,9,10\}.
\]
En la carta de Witt publicada,
\[
 P_\pi=\{2,4,5,6,7,8,9,10,11\},
 \qquad
 H_e=\{1,3,4,6,9,10\},
\]
y, por evaluación independiente de los dos soportes,
\[
 \boxed{B_K=H_e\cap P_\pi.}
\]
Además, para el soporte de acarreo negativo
\[
 H_\alpha^-=\{4,5,6,7,9,10\}
\]
se tiene
\[
 B_K=H_e\cap H_\alpha^-,
 \qquad H_\alpha^-=B_K\sqcup\{5,7\}.
\]
Estas identidades constituyen un testigo de compatibilidad finita: el mismo
subconjunto se recupera desde el sello aritmético y desde la incidencia de
Witt. No es una igualdad de cardinalidades, pues identifica posiciones
concretas.

\subsection{Separación tipada de las publicaciones arquimediana y excepcional del estado común}

\begin{theorem}[Criterio de fidelidad conjunta]
Sean \(X\) compacto y \(Y,Z\) espacios de Hausdorff. Si dos aplicaciones
continuas \(R:X\to Y\) y \(Q:X\to Z\) satisfacen
\[
 x\neq y\Longrightarrow R(x)\neq R(y)\ \text{o}\ Q(x)\neq Q(y),
\]
entonces \((R,Q):X\to Y\times Z\) es un encaje topológico.
\end{theorem}

\begin{proof}
La condición es la inyectividad de la aplicación diagonal. Toda aplicación
continua e inyectiva de un compacto en un espacio de Hausdorff es un
homeomorfismo sobre su imagen.
\end{proof}

En cada profundidad finita, la condición es decidible por enumeración de las
parejas \(x\neq y\) que permanecen simultáneamente en una fibra de \(R_n\) y
de \(Q_n\). La globalización exige compatibilidad de esos controles al pasar
de \(n+1\) a \(n\). El testigo \(B_K\) demuestra que las dos lecturas ya
comparten un invariante no trivial; la fidelidad conjunta requiere el control
de todas las fibras y no se sustituye por ese único testigo.

\subsection{Formulación precisa de la tesis sobre la recta real}

El objeto enriquecido asociado a la lectura arquimediana es la flecha
\[
 \mathbb R_{\mathfrak H}:=
 \bigl(X_\infty^{\rm arq}
 \xrightarrow{\operatorname{ev}_{\rm arq}}
 \operatorname{Im}(\operatorname{ev}_{\rm arq})\subseteq\mathbb R\bigr).
\]
El dominio conserva genealogías; la imagen conserva valores. La cobertura de
los cilindros a toda escala y la prolongabilidad compatible se demuestran en
la construcción del lector arquimediano
\eqref{eq:continuo-lector-arquimediano}: para todo compacto
\(K\subset\mathbb R\),
\(\operatorname{ev}_{\mathbb R}(\Omega_{\infty,K})=K\), y la unión dirigida
de esos intervalos cubre \(\mathbb R\). Por tanto,
\[
 \operatorname{Im}(\operatorname{ev}_{\rm arq})=\mathbb R,
 \qquad
 X_\infty^{\rm arq}/\!\sim_{\mathbb R}\ \cong\mathbb R.
\]
El cuerpo real no actúa como punto de partida de la construcción, sino como
cociente arquimediano de un espacio de secciones con memoria. El descenso de
suma y producto se verifica operación por operación y no condiciona esta
sobreyectividad ya demostrada.

La doble lectura queda entonces expresada por
\[
 X_\infty^{\rm com}\times C_\infty
 \xrightarrow{\ (\operatorname{ev}_{\rm arq},Q_\infty)\ }
 \mathbb R\times E_\infty.
\]
Esta fórmula es la traducción matemática de la afirmación conceptual: una
misma historia puede contraerse hasta un valor y, en la lectura transversal
de su frontera, conservar la incidencia que conduce a las simetrías
excepcionales.


% Fuente normativa activa de la obstrucción de factorización.
\section{Dependencia del lector de incidencia respecto del estado holonómico integral}
\label{p12-83-c17-s4:sec:obstruccion-lector-excepcional-minimo}

La compatibilidad finita de las dos evaluaciones no implica que cualquier
lector de incidencia refine la partición producida por el lector visible. El
catálogo regional permite decidir esta cuestión exactamente para la
realización mínima de Paley--Witt. Sea \(\mathcal U_{468}\) el conjunto de
los 468 estados regionales publicados y sea
\[
 \Psi:\mathcal U_{468}\longrightarrow\mathcal W_{243}\subset\mathbb F_3^6,
 \qquad
 \Psi(U_1,\ldots,U_6)=(-U_1,\ldots,-U_6)\pmod 3.
\]
Para \(w\in\mathbb F_3^6\), la aplicación
\[
 c(w)=(w,wA_W)\in C_W\subset\mathbb F_3^{12}
\]
es el codificador ternario fijado en la sección de Paley--Witt. Definimos
el lector de incidencia mínimo por
\[
 Q_{\min}=\operatorname{Supp}\circ c\circ\Psi:
 \mathcal U_{468}\longrightarrow\mathcal P([12]).
\]

Para una aplicación \(f:X\to Y\), denotemos por
\[
 \operatorname{Eq}(f)=\{(x,x')\in X^2:f(x)=f(x')\}
\]
su relación de equivalencia por fibras.

\begin{theorem}[Obstrucción del lector excepcional mínimo]
\label{p12-83-c17-s4:thm:obstruccion-lector-excepcional-minimo}
En el catálogo completo de 468 estados se verifica
\[
 \operatorname{Eq}(\Psi,Q_{\min})=\operatorname{Eq}(\Psi).
\]
En la notación abreviada habitual para particiones por fibras,
\[
 \ker(\Psi,Q_{\min})=\ker\Psi.
\]
Por consiguiente, \((\Psi,Q_{\min})\) no es inyectiva y el lector de
incidencia mínimo no distingue ningún par de estados que ya haya sido
identificado por \(\Psi\).
\end{theorem}

\begin{proof}
La factorización \(Q_{\min}=S\circ\Psi\), con
\(S=\operatorname{Supp}\circ c\), implica
\[
 \Psi(u)=\Psi(v)\Longrightarrow Q_{\min}(u)=Q_{\min}(v).
\]
La implicación recíproca entre las relaciones de equivalencia es inmediata,
pues \(\Psi\) es la primera componente de la aplicación diagonal. La
igualdad de relaciones se obtiene, por tanto, sin hipótesis estadísticas.
La enumeración exhaustiva confirma además que ambas particiones poseen el
mismo histograma:
\[
 132\,[1]+66\,[2]+18\,[3]+3\,[4]+6\,[5]+18\,[6],
\]
donde \(a\,[m]\) significa que existen \(a\) fibras de cardinal \(m\).
\end{proof}

La microfibra de \(w_\pi=010211\) proporciona el testigo distinguido. Sus
cinco regiones son
\[
\begin{split}
 &(501,614,498,169,272,272),\quad
 (810,923,870,169,272,272),\\
 &(810,983,810,169,272,272),\quad
 (870,923,810,169,272,272),\\
 &(870,923,810,418,674,521).
\end{split}
\]
Las cinco tienen la misma imagen visible y el mismo soporte excepcional
\[
 P_\pi=\{2,4,5,6,7,8,9,10,11\}.
\]
En el nivel de las 243 palabras aparecen 213 soportes distintos: 185 tienen
una preimagen, 27 tienen dos y uno tiene cuatro. La toma de soporte introduce,
por tanto, un segundo cociente además del impuesto por \(\Psi\).

El teorema delimita el lector mínimo, no la evaluación excepcional
enriquecida. Un lector capaz de separar las regiones de una misma palabra
debe actuar antes del cociente visible y utilizar al menos uno de los datos
que éste elimina: orientación, hoja, firma, retorno o estado de frontera. El
entrelazador dodecafásico desarrollado en la parte siguiente pertenece a esa
clase enriquecida y no contradice la obstrucción precedente.

La enumeración completa de las fibras produce el histograma anterior y
demuestra que la toma de soporte introduce un segundo cociente.


\section{Publicación arquimediana e incidencial}

Para cada nivel, sea \(\mathcal E_n\) un objeto incidencial y sean
\begin{equation}
 e_n:S_n\longrightarrow\mathcal E_n,
 \qquad
 \sigma_{n,m}:\mathcal E_n\longrightarrow\mathcal E_m
 \label{p12-83-c17-s5:eq:datos-lector-incidencial-u024}
\end{equation}
aplicaciones que satisfacen
\begin{equation}
 \sigma_{n,m}\circ e_n=e_m\circ r_{n,m}.
 \label{p12-83-c17-s5:eq:naturalidad-lector-incidencial-u024}
\end{equation}
La propiedad universal construye
\begin{equation}
 \operatorname{ev}_{\rm inc}:\Omega_\infty
 \longrightarrow\mathcal E_\infty,
 \qquad
 \mathcal E_\infty:=\varprojlim_n\mathcal E_n.
 \label{p12-83-c17-s5:eq:evaluacion-incidencial-limite-u024}
\end{equation}
La doble publicación es
\begin{equation}
 \mathcal P_\infty
 :=(\operatorname{ev}_{\mathbb R},\operatorname{ev}_{\rm inc}):
 \Omega_\infty\longrightarrow\mathbb R\times\mathcal E_\infty.
 \label{p12-83-c17-s5:eq:doble-publicacion-continuo-u024}
\end{equation}

\begin{figure}[H]
\centering
\begin{tikzcd}[column sep=10em,row sep=5em,
  cells={nodes={font=\large}}]
&\Omega_\infty
 \arrow[dl,"\operatorname{ev}_{\mathbb R}"']
 \arrow[dr,"\operatorname{ev}_{\rm inc}"]&\\
\mathbb R
 \arrow[dr,"\operatorname{coord}"']&&
\mathcal E_\infty
 \arrow[dl,"\operatorname{pub}_{\rm inc}"]\\
&\mathcal O&
\end{tikzcd}
\caption{Dos publicaciones de una misma historia. La conmutatividad en
\(\mathcal O\) no identifica los codominios \(\mathbb R\) y
\(\mathcal E_\infty\).}
\label{p12-83-c17-s5:fig:doble-publicacion-continuo-u024}
\end{figure}

Si se demuestra
\begin{equation}
 \operatorname{pub}_{\rm inc}\circ\operatorname{ev}_{\rm inc}
 =\operatorname{coord}\circ\operatorname{ev}_{\mathbb R},
 \label{p12-83-c17-s5:eq:compatibilidad-doble-publicacion-u024}
\end{equation}
se obtiene una comparación tipada de las dos publicaciones. La identidad no
afirma que una rama se construya a partir de la otra. Ambas parten de
\(\Omega_\infty\).


\section{Calibre de la publicación excepcional}

La segunda proyección no es un lector unario desprovisto de datos de
incidencia. Sea
\[
 \mathscr C_{\rm exc}
 =\mathscr C_{\rm ord}\times\mathscr C_{\rm Paley}
  \times\mathscr C_{\rm dod}\times\mathscr C_{\rm inc}
  \times\mathscr C_{\rm ret}
\]
el espacio de calibres excepcionales, cuyas coordenadas fijan,
respectivamente, el orden proyectivo, el calibre de Paley, la dodecada, el
alineamiento de incidencia y la cámara reticular. Se fija
\[
 \chi_{\rm exc}
 =(\chi_{\rm ord},\chi_{\rm Paley},\chi_{\rm dod},
   \chi_{\rm inc},\chi_{\rm ret})\in\mathscr C_{\rm exc}.
\]
El lector tipado tiene dominio
\[
 \mathfrak E:
 X_\infty^{\rm cal}\times\mathscr C_{\rm exc}
 \longrightarrow\mathbf{Exc},
 \qquad
 \mathfrak E_{\chi_{\rm exc}}(u):=\mathfrak E(u,\chi_{\rm exc}).
\]
En lo sucesivo, toda abreviatura \(\mathfrak E(u)\) significa
\(\mathfrak E_{\chi_{\rm exc}}(u)\) para este calibre explícitamente fijado.
La dependencia respecto de \(\chi_{\rm exc}\) no afecta a la existencia del
continuo ni a su publicación arquimediana: tipa la conservación excepcional
de la incidencia una vez construido el estado común.

\begin{falsifier}[Pérdida del calibre excepcional]
La comparación queda incompleta si se cambia el orden proyectivo, la dodecada,
el alineamiento de incidencia o la cámara reticular sin transportar las demás
coordenadas del calibre, o si se usa \(\mathfrak E\) como lector unario antes
de fijar \(\chi_{\rm exc}\). Tal fallo invalida esa comparación excepcional;
no modifica el cociente arquimediano ni regenera el continuo.
\end{falsifier}


La segunda lectura no nace del número real. Sea
\[
 \mathfrak E:
 \mathscr O_{\mathbb R,\leftrightarrow}^{\rm HMT}
 \longrightarrow\mathbf{Exc}
\]
el lector continuo de incidencia, con codominio de Hausdorff, que prolonga
la cadena de Paley--Witt, Golay,
Leech, \(V^\natural\), Monstruo y Moonshine. Ambas aplicaciones parten de
la misma historia:
\begin{equation}
 \mathfrak D(u)=(\mathfrak R(u),\mathfrak E(u)).
 \label{p12-83-c17-s7:hmtch:eq-doble-proyeccion-two-way}
\end{equation}
La primera contrae las fibras compatibles hasta su media arquimediana; la
segunda conserva incidencia, orientación y memoria de borde.

Fijado el origen del calendario dodecafásico, sea
\[
 R_+:=\{k\in\mathbb Z:1+(k\bmod 9)\in\{1,4,7\}\}
\]
el conjunto biinfinito de retornos completos \((\tau=+1)\). Es sindético:
sus huecos son uniformemente acotados. Para
\(u=(m,(x_k)_{k\in\mathbb Z})\) definimos
\[
 \operatorname{Ret}_+(u)
 :=\bigl(m,(x_k)_{k\in R_+}\bigr)
 \in\mathbb Z\times\prod_{k\in R_+}X,
\]
pero cada \(x_k\) es aquí el estado entero
\eqref{p12-83-c19-s1:hmtch:eq-celula-trit-completa} con hoja, carry, ruta e incidencia,
no el signo aislado.

\begin{proposition}[Muestreo sin pérdida del registro genealógico de retorno]
\label{p12-83-c17-s7:hmtch:prop-retorno-reconstruye}
La aplicación \(\operatorname{Ret}_+\), con la topología de subespacio en el
producto, es un homeomorfismo sobre su imagen.
En consecuencia existe una única aplicación continua
\(\mathfrak E_+\) sobre esa imagen tal que
\begin{equation}
 \mathfrak E=\mathfrak E_+\circ\operatorname{Ret}_+.
 \label{p12-83-c17-s7:hmtch:eq-factorizacion-excepcional-retorno}
\end{equation}
\end{proposition}

\begin{proof}
Entre dos retornos consecutivos hay un número uniformemente acotado de pasos
del calendario. El estado completo en un retorno contiene la transición y
el registro genealógico que determinan los pasos intermedios; hacia la otra dirección se
usa la flecha inversa de la completación grupoidal. Así se reconstruye toda
la órbita. La aplicación es continua y, al restringirla a una carta
compacta, una inyección continua en un espacio de Hausdorff es un
homeomorfismo sobre su imagen. La coordenada entera se conserva
explícitamente. La factorización se define transportando \(\mathfrak E\) por
la inversa y es continua carta a carta.
\end{proof}

Por eso el \(+1\) participa en el corredor excepcional como memoria completa
de retorno: el muestreo de estados enteros no pierde la salida excepcional.
Esta proposición demuestra recuperación y factorización, no una generación
del corredor a partir del escalar \(+1\) aislado. La salida arquimediana y la
salida excepcional son dos resultantes coordinadas del mismo antecedente.
 
\chapter{Reconstrucción arquimediana por cilindros semicerrados con memoria}
\label{chap:83-18}

El límite inverso \(\Omega_\infty\) conserva la historia completa. Para
obtener una coordenada real hay que añadir una familia compatible de
intervalos y demostrar tres hechos diferentes: cada historia determina un
único valor; todo valor del compacto declarado es alcanzado; y las medidas
finitas descienden de forma consistente. Ninguno de estos hechos se deduce
del mero crecimiento del número de estados.

La publicación arquimediana conserva asimismo una aproximación de la
identidad por esperanzas condicionales cilíndricas. La concentración de masa
unitaria sobre prefijos compatibles proporciona el paso exacto desde los
cilindros nonádicos hasta la medida de Dirac, sin identificar densidad,
medida del soporte y masa total.

\begin{HMTScientificOwnerSubordinated}
\chapter{Cilindros nonádicos, identidad y núcleo de Dirac}

\section{Sistema inverso de cilindros nonádicos compatibles con memoria genealógica}

Sea
\[
X=(\Znine)^{\mathbb N}
\]
con la topología producto y la medida uniforme \(\mu\). Para
\(x=(x_1,x_2,\ldots)\), el cilindro de profundidad \(n\) es
\[
C_n(x)=\{y\in X:y_j=x_j\text{ para }1\le j\le n\}.
\]
Los cilindros son simultáneamente abiertos y cerrados, forman una base y
satisfacen
\[
\mu(C_n(x))=9^{-n},\qquad
C_{n+1}(x)\subset C_n(x),\qquad
\bigcap_n C_n(x)=\{x\}.
\]
Añadir una cifra no crea un número nuevo: restringe el cilindro de
prolongaciones admisibles. Esta afirmación da contenido topológico a la tesis
genealógica del número.

\section{Concentración cilíndrica de la identidad y construcción del núcleo de Dirac}

\begin{theorem}[Concentración unitaria nonádica]
Defínase
\[
f_{n,x}=9^n\mathbf1_{C_n(x)}.
\]
Entonces
\[
\mu(C_n(x))\to0,\qquad f_{n,x}(x)\to\infty,\qquad
\int_Xf_{n,x}\,\dd\mu=1,
\]
y las medidas \(f_{n,x}\mu\) convergen débil-* a \(\delta_x\).
\end{theorem}

\begin{proof}
La primera identidad es inmediata. La altura de \(f_{n,x}\) sobre su soporte
es \(9^n\), de modo que diverge. La integral vale
\(9^n9^{-n}=1\). Para \(g\in C(X)\), la oscilación de \(g\) en \(C_n(x)\)
tiende a cero por continuidad uniforme; por tanto
\[
\left|\int g f_{n,x}\,\dd\mu-g(x)\right|
\le \sup_{y\in C_n(x)}|g(y)-g(x)|\longrightarrow0.
\]
Ésta es la convergencia débil-* a \(\delta_x\).
\end{proof}

El teorema no afirma \(1/0=\infty\). Exhibe tres evaluaciones de un mismo
objeto: medida del soporte, densidad y masa total. En bloques de tres cifras,
\[
\mu(C_{3m})=729^{-m},\qquad
f_{3m}=729^m\mathbf1_{C_{3m}},\qquad
729^{-m}729^m=1.
\]
La conservación evolutiva de la unidad adopta aquí una forma analítica exacta.

\section{Convergencia del núcleo de Kronecker a la distribución de Dirac}

En el nivel finito \(X_n=(\Znine)^n\), sea
\[
K_n(x,y)=9^n\mathbf1_{\{x=y\}}.
\]
Con la medida uniforme, el operador
\[
(E_nf)(x)=\sum_{y\in X_n}K_n(x,y)f(y)9^{-n}
\]
es la identidad. En el límite cilíndrico la condición \(x=y\) se reemplaza por
«mismo prefijo de longitud \(n\)» y \(E_n\) es la esperanza condicional sobre
el álgebra cilíndrica. Se tiene
\[
E_n^2=E_n,\qquad E_n^*=E_n,\qquad E_nf\to f
\quad\text{en }L^2(X,\mu).
\]

El puente exacto es, por tanto,
\[
\delta_{ij}
\longrightarrow K_n(x,y)
\longrightarrow\delta_x.
\]
La delta de Kronecker es un núcleo matricial; la delta de Dirac es una medida o
distribución. Ninguna es el operador de Dirac
\(D=\gamma^\mu\partial_\mu\). El tránsito hacia espinores y Clifford requiere
un segundo entrelazador y será tratado después.

\section{Ley de escala, dimensión efectiva y densidad de la medida cilíndrica}

Supóngase que una contracción de razón \(\lambda\) actúa de modo independiente
en \(d\) direcciones y que la medida producto escala como \(\lambda^d\). Para
conservar masa uno, la densidad debe escalar como
\[
\rho\longmapsto\lambda^{-d}\rho.
\]
Los exponentes \((-1,-2,-3)\) aparecen entonces como pesos de densidad en línea,
área y volumen. Esta proposición es general y exacta bajo la hipótesis de
producto. No permite identificar automáticamente \(\lambda\) con
\(\alpha\); esa identificación necesita una carta métrica.

\section{Monodromía nonádica y computabilidad de la evaluación arquimediana}

La existencia de una rama infinita en un árbol finitamente ramificado puede
seguirse del lema de König, pero König no produce un selector computable. Hay
árboles binarios recursivos infinitos sin rama recursiva. La conexión nonádica
conjunta no descansa en esa existencia abstracta: posee un selector uniforme
certificado que genera a profundidad finita arbitraria y abre los testigos
decimales sólo después de la generación.

\begin{falsadorBox}[title={Tres confusiones prohibidas}]
No se escribirá \(1/0\); no se representará la delta de Dirac como función
ordinaria con «un uno infinito»; no se usará König como algoritmo. El resultado
fuerte es la compatibilidad exacta entre cilindros, densidad unitaria y límite
débil, más el selector constructivo independiente de la comparación decimal.
\end{falsadorBox}
 \end{HMTScientificOwnerSubordinated}

\section{Sistema de intervalos asociado a los estados finitos}

Sea \((S_n,r_{n,m})\) el sistema inverso de la
\cref{chap:sistema-inverso-continuidad-genealogica}. Para cada
\(a\in S_n\), asignamos un intervalo compacto no vacío
\begin{equation}
 I_n(a)=[\ell_n(a),u_n(a)]\subset\mathbb R.
 \label{p12-83-c18-s1:eq:intervalo-asociado-estado}
\end{equation}
La asignación es \emph{compatible con el refinamiento} cuando
\begin{equation}
 r_{n+1,n}(b)=a
 \quad\Longrightarrow\quad
 I_{n+1}(b)\subseteq I_n(a).
 \label{p12-83-c18-s1:eq:anidamiento-intervalos-estados}
\end{equation}
Es \emph{uniformemente contractiva} cuando existe una sucesión
\(\delta_n\downarrow0\) tal que
\begin{equation}
 \operatorname{diam}I_n(a)\leq\delta_n
 \qquad(a\in S_n).
 \label{p12-83-c18-s1:eq:contraccion-uniforme-intervalos}
\end{equation}

Para una historia \(x=(x_n)_n\), las inclusiones
\eqref{p12-83-c18-s1:eq:anidamiento-intervalos-estados} producen una cadena
\begin{equation}
 I_0(x_0)\supseteq I_1(x_1)\supseteq I_2(x_2)\supseteq\cdots.
 \label{p12-83-c18-s1:eq:cadena-intervalos-historia}
\end{equation}

\begin{theorem}[Evaluación de una historia]
\label{p12-83-c18-s1:thm:evaluacion-unica-historia-u024}
Si la asignación es compatible y uniformemente contractiva, entonces para
cada \(x\in\Omega_\infty\) existe un único número real
\begin{equation}
 \operatorname{ev}_{\mathbb R}(x)
 \in\bigcap_{n\geq0}I_n(x_n).
 \label{p12-83-c18-s1:eq:evaluacion-real-historia-u024}
\end{equation}
La aplicación
\begin{equation}
 \operatorname{ev}_{\mathbb R}:\Omega_\infty\longrightarrow\mathbb R
 \label{p12-83-c18-s1:eq:aplicacion-evaluacion-real-u024}
\end{equation}
es continua.
\end{theorem}

\begin{proof}
Los intervalos de \eqref{p12-83-c18-s1:eq:cadena-intervalos-historia} son compactos,
no vacíos y anidados. El teorema de los intervalos encajados garantiza que
su intersección es no vacía. Si \(s,t\) pertenecen a ella, entonces
\[
 |s-t|\leq\operatorname{diam}I_n(x_n)\leq\delta_n
\]
para todo \(n\). Al hacer \(n\to\infty\), se obtiene \(s=t\).

Para la continuidad, sea \(\varepsilon>0\) y elijamos \(n\) con
\(\delta_n<\varepsilon\). Si \(y\in[x_n]_n\), entonces
\(I_n(y_n)=I_n(x_n)\). Las dos evaluaciones pertenecen a ese intervalo y,
por tanto,
\[
 |\operatorname{ev}_{\mathbb R}(x)
  -\operatorname{ev}_{\mathbb R}(y)|<\varepsilon.
\]
El cilindro \([x_n]_n\) es un entorno abierto de \(x\).
\end{proof}

La construcción no identifica \(x\) con
\(\operatorname{ev}_{\mathbb R}(x)\). Una misma coordenada puede tener
varias historias, y una historia contiene más datos que cualquier punto de
la recta.

\section{Cobertura coherente del compacto visible}

Sea \(K\subset\mathbb R\) un compacto no vacío. Además de anidamiento y
contracción, exigimos
\begin{align}
 K&=\bigcup_{a\in S_0}I_0(a),
 \label{p12-83-c18-s1:eq:cobertura-inicial-compacto-u024}\\
 I_n(a)&=
 \bigcup_{\substack{b\in S_{n+1}\\r_{n+1,n}(b)=a}}
 I_{n+1}(b)
 \qquad(a\in S_n).
 \label{p12-83-c18-s1:eq:cobertura-extensiones-compatible-u024}
\end{align}
La segunda igualdad es una condición geométrica adicional. La
sobreyectividad de \(r_{n+1,n}\) sólo garantiza la existencia de estados
por encima de \(a\); no garantiza que sus intervalos cubran todo
\(I_n(a)\).

\begin{theorem}[Sobreyectividad de la evaluación]
\label{p12-83-c18-s1:thm:sobreyectividad-evaluacion-compacto}
Si se cumplen
\eqref{p12-83-c18-s1:eq:cobertura-inicial-compacto-u024} y
\eqref{p12-83-c18-s1:eq:cobertura-extensiones-compatible-u024}, entonces
\begin{equation}
 \operatorname{ev}_{\mathbb R}(\Omega_\infty)=K.
 \label{p12-83-c18-s1:eq:imagen-evaluacion-compacto}
\end{equation}
\end{theorem}

\begin{proof}
La inclusión de izquierda a derecha sigue de
\(I_n(x_n)\subseteq I_0(x_0)\subseteq K\).

Fijemos \(t\in K\). Definamos
\begin{equation}
 T_n(t):=\{a\in S_n:t\in I_n(a)\}.
 \label{p12-83-c18-s1:eq:estados-que-cubren-punto}
\end{equation}
La cobertura inicial hace no vacío \(T_0(t)\). Si \(a\in T_n(t)\), la
identidad \eqref{p12-83-c18-s1:eq:cobertura-extensiones-compatible-u024} proporciona al
menos un \(b\in T_{n+1}(t)\) con \(r_{n+1,n}(b)=a\). El grafo de estas
extensiones es finitamente ramificado y posee niveles no vacíos a toda
profundidad. Existe, por tanto, una cadena compatible \(x=(x_n)_n\) con
\(t\in I_n(x_n)\) para todo \(n\). La unicidad de la intersección obliga a
\(\operatorname{ev}_{\mathbb R}(x)=t\).
\end{proof}

\section{Cociente genealógico arquimediano y homeomorfismo con \texorpdfstring{\(\mathbb R\)}{R}}

Definimos la relación
\begin{equation}
 x\sim_{\mathbb R}y
 \quad\Longleftrightarrow\quad
 \operatorname{ev}_{\mathbb R}(x)
 =\operatorname{ev}_{\mathbb R}(y).
 \label{p12-83-c18-s1:eq:relacion-equivalencia-arquimediana}
\end{equation}
Sea
\begin{equation}
 q_{\mathbb R}:\Omega_\infty
 \longrightarrow\Omega_\infty/\!\sim_{\mathbb R}
 \label{p12-83-c18-s1:eq:proyeccion-cociente-arquimediano}
\end{equation}
la proyección canónica.

\begin{theorem}[Reconstrucción del compacto como cociente]
\label{p12-83-c18-s1:thm:homeomorfismo-cociente-compacto-u024}
Bajo las hipótesis anteriores, existe un único homeomorfismo
\begin{equation}
 \overline{\operatorname{ev}}_{\mathbb R}:
 \Omega_\infty/\!\sim_{\mathbb R}
 \xrightarrow{\ \cong\ }K
 \label{p12-83-c18-s1:eq:homeomorfismo-cociente-compacto-u024}
\end{equation}
tal que
\begin{equation}
 \operatorname{ev}_{\mathbb R}
 =\overline{\operatorname{ev}}_{\mathbb R}\circ q_{\mathbb R}.
 \label{p12-83-c18-s1:eq:factorizacion-evaluacion-cociente}
\end{equation}
La fibra sobre \(t\in K\) es
\begin{equation}
 \mathfrak G_t
 :=\operatorname{ev}_{\mathbb R}^{-1}(t),
 \label{p12-83-c18-s1:eq:fibra-genealogias-mismo-valor}
\end{equation}
el conjunto completo de genealogías que publican \(t\).
\end{theorem}

\begin{proof}
La aplicación \(\operatorname{ev}_{\mathbb R}\) es continua, sobreyectiva y
constante exactamente en las clases de \(\sim_{\mathbb R}\). Por la
propiedad universal del cociente induce una aplicación continua y biyectiva
\(\overline{\operatorname{ev}}_{\mathbb R}\). El dominio es compacto por
ser imagen continua del compacto \(\Omega_\infty\), y \(K\) es Hausdorff.
Una biyección continua de un compacto a un Hausdorff es un homeomorfismo.
\end{proof}

El teorema prueba dos cosas y las mantiene separadas:
\begin{equation}
 \boxed{
 \begin{array}{c}
 \text{el espacio de historias es profinito y conserva memoria,}\\
 \text{su cociente por igualdad de evaluación es el compacto visible.}
 \end{array}}
 \label{p12-83-c18-s1:eq:dos-niveles-continuo-u024}
\end{equation}
No se afirma que \(\Omega_\infty\) sea homeomorfo a \(K\).

\section{Atlas genealógico de la recta real}

Para \(m\geq1\), sea \(K_m=[-m,m]\). Supongamos dado para cada \(m\) un
sistema \(\Omega_\infty^{(m)}\) que reconstruye \(K_m\), junto con
inclusiones cerradas
\begin{equation}
 j_m:\Omega_\infty^{(m)}\hookrightarrow
 \Omega_\infty^{(m+1)}
 \label{p12-83-c18-s1:eq:inclusion-atlas-genealogico}
\end{equation}
que satisfacen
\begin{equation}
 \operatorname{ev}_{m+1}\circ j_m
 =\iota_m\circ\operatorname{ev}_m,
 \label{p12-83-c18-s1:eq:compatibilidad-atlas-real}
\end{equation}
donde \(\iota_m:K_m\hookrightarrow K_{m+1}\) es la inclusión ordinaria.

\begin{proposition}[Agotamiento de la recta]
\label{p12-83-c18-s1:prop:agotamiento-recta-real-u024}
El cociente arquimediano del colímite estricto
\begin{equation}
 \Omega_\infty^{\mathbb R}
 :=\varinjlim_m\Omega_\infty^{(m)}
 \label{p12-83-c18-s1:eq:colimite-historias-recta}
\end{equation}
es
\begin{equation}
 \varinjlim_m[-m,m]
 =\bigcup_{m\geq1}[-m,m]
 =\mathbb R
 \label{p12-83-c18-s1:eq:colimite-compactos-recta}
\end{equation}
con su topología usual.
\end{proposition}

\begin{proof}
Cada pieza posee cociente \(K_m\) por el
\cref{p12-83-c18-s1:thm:homeomorfismo-cociente-compacto-u024};
\eqref{p12-83-c18-s1:eq:compatibilidad-atlas-real} entrelaza los cocientes. La topología
final del agotamiento por compactos coincide con la topología usual de
\(\mathbb R\): un subconjunto \(U\subseteq\mathbb R\) es abierto si y sólo
si \(U\cap[-m,m]\) es abierto en \([-m,m]\) para todo \(m\).
\end{proof}

El paso de un intervalo acotado a la recta no borra las fibras
\(\mathfrak G_t\). Las inclusiones \(j_m\) conservan la historia y sólo
amplían la carta visible disponible.

\section{Cartas ternaria balanceada y decimal por bloques}

La representación ternaria balanceada de un entero tiene la forma
\begin{equation}
 n=\sum_{j=0}^{J}\epsilon_j3^j,
 \qquad
 \epsilon_j\in\{-1,0,+1\},
 \label{p12-83-c18-s1:eq:carta-ternaria-balanceada-u024}
\end{equation}
con una regla de acarreo fijada. Esta carta conserva el signo local y el
cociente de la reducción; no debe confundirse con la palabra no negativa en
\(\mathbb F_3\).

La carta decimal por bloques utiliza
\begin{equation}
 D_N=1000D_{N-1}+u_N,
 \qquad
 u_N\in\{0,\ldots,999\},
 \qquad D_0=0.
 \label{p12-83-c18-s1:eq:concatenacion-base-mil-u024}
\end{equation}
El bloque \(u_N\), su clase módulo nueve y el acarreo son campos distintos.
Como \(1000\equiv1\pmod9\),
\begin{equation}
 D_N\equiv\sum_{j=1}^{N}u_j\pmod9.
 \label{p12-83-c18-s1:eq:base-mil-compatible-mod9-u024}
\end{equation}
Esta congruencia no recupera el orden de los bloques; la secuencia
\((u_1,\ldots,u_N)\) permanece en el registro.

Para \(M\geq1\), las celdas decimales semiabiertas son
\begin{equation}
 \mathcal D_{M,q}
 :=[q10^{-M},(q+1)10^{-M}),
 \qquad q\in\mathbb Z.
 \label{p12-83-c18-s1:eq:celda-decimal-semiabierta-u024}
\end{equation}
Un intervalo semiabierto \([\ell,h)\) está contenido en una única celda si
y sólo si
\begin{equation}
 \left\lfloor10^M\ell\right\rfloor
 =\left\lceil10^Mh\right\rceil-1.
 \label{p12-83-c18-s1:eq:criterio-exacto-celda-decimal-u024}
\end{equation}
La fórmula con techo en el extremo derecho cubre también el caso en que
\(h\) coincide exactamente con una frontera decimal.

\begin{definition}[Truncamiento y redondeo]
Para \(x\geq0\), definimos
\begin{align}
 \operatorname{Tr}_M(x)
 &:=10^{-M}\left\lfloor10^Mx\right\rfloor,
 \label{p12-83-c18-s1:eq:truncamiento-decimal-u024}\\
 \operatorname{Rd}_M(x)
 &:=10^{-M}\left\lfloor10^Mx+\frac12\right\rfloor.
 \label{p12-83-c18-s1:eq:redondeo-decimal-u024}
\end{align}
\end{definition}
Estas operaciones coinciden salvo cuando la fracción descartada cruza el
umbral de media unidad en la última posición. La distinción es semántica y
aritmética; no puede sustituirse una por otra para forzar compatibilidad de
prefijos.

\section{Descenso de operaciones por las fibras de evaluación}

Sea \(D_\star\subseteq\Omega_\infty\times\Omega_\infty\) el dominio de una
operación parcial continua
\begin{equation}
 \star:D_\star\longrightarrow\Omega_\infty.
 \label{p12-83-c18-s1:eq:operacion-historias-u024}
\end{equation}
Suponemos que \(D_\star\) es saturado por las fibras de
\(\operatorname{ev}_{\mathbb R}\times\operatorname{ev}_{\mathbb R}\).

\begin{theorem}[Criterio exacto de descenso]
\label{p12-83-c18-s1:thm:criterio-descenso-operaciones-u024}
Existe una única operación visible
\begin{equation}
 \overline\star:
 (\operatorname{ev}_{\mathbb R}\times\operatorname{ev}_{\mathbb R})
 (D_\star)\longrightarrow K
 \label{p12-83-c18-s1:eq:operacion-descendida-visible}
\end{equation}
que satisface
\begin{equation}
 \operatorname{ev}_{\mathbb R}(x\star y)
 =\operatorname{ev}_{\mathbb R}(x)
  \ \overline\star\ 
  \operatorname{ev}_{\mathbb R}(y)
 \label{p12-83-c18-s1:eq:entrelazamiento-operacion-descendida}
\end{equation}
si y sólo si
\begin{equation}
 \begin{split}
 &\operatorname{ev}_{\mathbb R}(x)=
  \operatorname{ev}_{\mathbb R}(x'),\\
 &\operatorname{ev}_{\mathbb R}(y)=
  \operatorname{ev}_{\mathbb R}(y')
 \end{split}
 \quad\Longrightarrow\quad
 \operatorname{ev}_{\mathbb R}(x\star y)=
 \operatorname{ev}_{\mathbb R}(x'\star y')
 \label{p12-83-c18-s1:eq:constancia-operacion-sobre-fibras}
\end{equation}
para todos los pares del dominio.
\end{theorem}

\begin{proof}
Si \(\overline\star\) existe, el miembro derecho de
\eqref{p12-83-c18-s1:eq:entrelazamiento-operacion-descendida} depende sólo de las dos
evaluaciones, y se obtiene
\eqref{p12-83-c18-s1:eq:constancia-operacion-sobre-fibras}. Recíprocamente, esta
constancia permite definir
\[
 s\ \overline\star\ t
 :=\operatorname{ev}_{\mathbb R}(x\star y)
\]
para cualesquiera representantes \(x,y\) con evaluaciones \(s,t\). La
condición prueba que la definición no depende de los representantes. La
unicidad es inmediata.
\end{proof}

Las hojas aditiva y multiplicativa de APP suministran operaciones en un
soporte común, pero sus acarreos y sus memorias no se identifican. El
teorema anterior permite que suma y producto desciendan a la carta visible
sin exigir que sus realizaciones genealógicas coincidan.

\section{Correspondencia entre cardinalidad discreta, medida cilíndrica y publicación arquimediana}

En cada profundidad aparecen tres operaciones que deben conservarse
separadas:
\begin{align}
 \operatorname{Count}_n(A)&:=|A|,
 \label{p12-83-c18-s1:eq:operacion-conteo-u024}\\
 \operatorname{Measure}_n(A)&:=\mu_n(A),
 \label{p12-83-c18-s1:eq:operacion-medida-u024}\\
 \operatorname{Publish}_n(a)&:=L_n(a).
 \label{p12-83-c18-s1:eq:operacion-publicacion-u024}
\end{align}
El conteo depende del número de estados; la medida depende de pesos
compatibles; la publicación depende del lector. Dos fibras de igual
cardinal pueden tener pesos distintos, y dos estados de igual publicación
pueden pertenecer a cilindros distintos. La igualdad de una de estas tres
salidas no implica la igualdad de las otras dos.

\section{Obstrucciones de la evaluación arquimediana: anidamiento, continuidad, cociente, medidas proyectivas y descenso de operaciones}

\begin{criterion}[Refutación de la evaluación y de la medida]
La construcción queda refutada si se verifica cualquiera de los hechos
siguientes:
\begin{enumerate}
 \item una extensión compatible recibe un intervalo que no está contenido
       en el intervalo de su restricción;
 \item los diámetros no tienden a cero pero se afirma unicidad de la
       intersección;
 \item la evaluación no es continua en la topología de cilindros;
 \item se afirma que todo punto de \(K\) aparece sin demostrar la cobertura
       coherente \eqref{p12-83-c18-s1:eq:cobertura-extensiones-compatible-u024};
 \item la aplicación inducida por el cociente no es biyectiva;
 \item se identifica el espacio profinito con su cociente visible;
 \item las inclusiones del atlas no entrelazan sus evaluaciones;
 \item se usa \(\lfloor10^Mh\rfloor\) en lugar de
       \(\lceil10^Mh\rceil-1\) para un intervalo semiabierto y se pierde
       el caso de frontera;
 \item se confunden truncamiento y redondeo;
 \item una operación descendida depende del representante elegido en una
       fibra;
 \item una familia \((\mu_n)_n\) viola
       \eqref{eq:consistencia-proyectiva-medidas-u024};
 \item la medida de un cilindro depende del nivel al que se lo refina;
 \item una medida condicionada deja de sumar uno o pierde consistencia;
 \item se identifica la medida visible \(\nu\) con las medidas internas
       \(\mu_t\);
 \item se deduce una flecha entre \(\mathbb R\) y
       \(\mathcal E_\infty\) sólo porque ambas publicaciones comparten
       antecedente;
 \item se usa el conteo de una fibra como si determinara su medida o su
       valor publicado.
\end{enumerate}
\end{criterion}


\section{Del espacio de historias a la recta real}

Sea
\[
 \Psi:\mathbb Z_3^6\xrightarrow{\ \sim\ }
       (\mathbb F_3^6)^{\mathbb N}
\]
la conjugación Hensel entre la memoria inicial y la cinta completa de
bloques de seis trits. Concatenar las seis coordenadas de cada bloque da
una cinta ternaria
\(
 (a_n)_{n\geq 1}\in\{0,1,2\}^{\mathbb N}
\), sobre la que actúa la evaluación
\[
 \operatorname{ev}_3((a_n))
 :=\sum_{n\geq1}\frac{a_n}{3^n}\in[0,1].
\]
Las dos expansiones de un extremo ternario representan el mismo valor y
permanecen diferenciadas antes de aplicar la evaluación. Esto no constituye
una ambigüedad del objeto genealógico: es precisamente una fibra del lector.

Definimos
\[
 X_{\mathrm{raw}}:=\mathbb Z\times\mathbb Z_3^6,
 \qquad
 P(m,x):=m+\operatorname{ev}_3(\operatorname{flat}(\Psi x)).
\]
La coordenada entera proporciona el atlas numerable de la recta y la
coordenada 3-ádica conserva la historia completa de la parte fraccionaria.

\begin{theorem}[Reconstrucción exacta del cociente arquimediano]
\label{p12-83-c18-s2:hmtch:thm-cociente-real}
La aplicación \(P:X_{\mathrm{raw}}\to\mathbb R\) es sobreyectiva. Si
\[
 x\sim_P y
 \quad\Longleftrightarrow\quad
 P(x)=P(y),
\]
entonces
\[
 \overline P:X_{\mathrm{raw}}/\!\sim_P\;\longrightarrow\mathbb R,
 \qquad [x]\longmapsto P(x),
\]
es una biyección. En consecuencia,
\[
 \left|X_{\mathrm{raw}}/\!\sim_P\right|
 =|\mathbb R|
 =2^{\aleph_0}.
\]
\end{theorem}

\begin{proof}
Toda sucesión ternaria se agrupa de manera única en bloques consecutivos de
seis cifras. La conjugación \(\Psi\) asigna a esa cinta una única memoria
3-ádica inicial. Dado \(r\in\mathbb R\), se toma
\(m=\lfloor r\rfloor\) y una expansión ternaria de
\(r-m\in[0,1)\); la reconstrucción anterior produce
\(x\in\mathbb Z_3^6\) con \(P(m,x)=r\). Así, \(P\) es sobreyectiva.

La aplicación \(\overline P\) está bien definida por la definición de
\(\sim_P\), es sobreyectiva porque lo es \(P\), e inyectiva porque dos
clases con la misma imagen pertenecen a una misma fibra. Finalmente,
\[
 |\mathbb Z_3^6|
 =\left|(\mathbb F_3^6)^{\mathbb N}\right|
 =729^{\aleph_0}
 =2^{\aleph_0},
\]
y multiplicar por la coordenada numerable \(\mathbb Z\) no altera ese
cardinal. El cálculo coincide con el del cociente, ya identificado con
\(\mathbb R\).
\end{proof}

\begin{remark}[Versión superviviente]
\label{p12-83-c18-s2:hmtch:rem-version-superviviente}
El teorema anterior es incondicional para la completación bruta. Para un
subespacio superviviente \(X_{\mathrm{surv}}\), el mismo resultado se
obtiene cuando las coberturas de cada nivel son coherentes, los diámetros se
contraen uniformemente y toda celda visible posee una prolongación. Bajo
esas hipótesis, el lema de la rama infinita produce una sección sobre cada
punto y la evaluación induce, compacto a compacto, un homeomorfismo del
cociente con el intervalo publicado. El agotamiento por
\([-m,m]\) reconstruye después toda la recta.
\end{remark}

\section{Media arquimediana y doble proyección del mismo estado}

Sea \(u=(m,\mathbf u)\in
\mathscr O_{\mathbb R,\leftrightarrow}^{\rm HMT}\). Un prefijo
\(\mathbf u_n\) determina en la carta fraccionaria un cilindro orientado
semiabierto
\[
 I_n(\mathbf u_n)=[a_n,b_n),
\]
con la convención complementaria en los extremos de doble expansión. Para
aplicar compacidad se usa su clausura
\(K_n=[a_n,b_n]\). Las compatibilidades dan
\[
 K_{n+1}\subseteq K_n,
 \qquad
 \delta_n:=b_n-a_n\longrightarrow0.
\]
La media arquimediana global del cilindro es
\[
 \mu_n(u)=m+\frac{a_n+b_n}{2}.
\]
Si \(q\geq n\), entonces
\[
 |\mu_q(u)-\mu_n(u)|\leq\delta_n,
\]
de modo que
\begin{equation}
 \mathfrak R(u):=\lim_{n\to\infty}\mu_n(u)
 \label{p12-83-c18-s2:hmtch:eq-media-arquimediana}
\end{equation}
existe. Lo que tiende a cero es el diámetro del cilindro; el valor real es
la resultante del encajamiento, no aquello que se anula.

\begin{theorem}[Cociente two-way y recta real]
\label{p12-83-c18-s2:hmtch:thm-cociente-two-way-real}
Sea
\[
 u\sim_{\mathfrak R}v
 \quad\Longleftrightarrow\quad
 \mathfrak R(u)=\mathfrak R(v).
\]
Bajo la cobertura coherente de las cartas HMT, la evaluación induce un
homeomorfismo
\begin{equation}
 \overline{\mathfrak R}:
 \mathscr O_{\mathbb R,\leftrightarrow}^{\rm HMT}/\!
 \sim_{\mathfrak R}
 \xrightarrow{\ \sim\ }\mathbb R.
 \label{p12-83-c18-s2:hmtch:eq-homeomorfismo-two-way-real}
\end{equation}
En particular, éste es un morfismo HMT continuo: no el resultado de elegir
una historia aislada de cada fibra.
\end{theorem}

\begin{proof}
La cota de cilindros prueba la continuidad de \(\mathfrak R\). La cobertura
produce una historia sobre cada real y, por tanto, sobreyectividad. Al pasar
al cociente por las fibras, la aplicación inducida es biyectiva. En cada
carta compacta, una biyección continua hacia un intervalo de Hausdorff es un
homeomorfismo. La topología global del cociente es la topología final del
pegado de esas cartas; la fibra de evaluación identifica el extremo
\(m+1\) de una carta con el extremo \(0\) trasladado de la carta siguiente.
El recubrimiento resultante es localmente finito, por lo que los
homeomorfismos de carta se pegan en un homeomorfismo global hacia
\(\mathbb R\).
\end{proof}
  \part{Caracteres matemáticos de la prolongación coinductiva}

\chapter{Semillas, completación cúbica y sincronización de los caracteres fundamentales}
\section{Las dos proyecciones mecánicas y la conservación exacta de la unidad de borde}

La afirmación de que la unidad se conserva durante la prolongación admite ya
una formulación aritmética que no depende de una analogía física. Sea

\[
 \delta:=\log_{10}\!\left(\frac{10}{9}\right),
 \qquad 0<\delta<1,
\]

la discrepancia entre la capacidad triádica \(3^6=729\) y la completación
decimal \(10^3=1000\). La irracionalidad de \(\delta\) ya está demostrada en
el construcción por valuaciones primas; su trascendencia ha sido probada en este
cuaderno mediante Gelfond--Schneider. Consideremos las dos codificaciones
unilaterales de la misma recta irracional:

\[
 z_n^-:=\lfloor(n+1)\delta\rfloor-\lfloor n\delta\rfloor,
 \qquad
 z_n^+:=\lceil(n+1)\delta\rceil-\lceil n\delta\rceil,
 \qquad n\geq0.
\]

Ambas toman valores en \(\{0,1\}\). No son dos densidades ni dos dinámicas:
son las dos convenciones de cierre de un mismo corte irracional. Como
\(n\delta\notin\mathbb Z\) para todo \(n\geq1\), se cumple

\[
 z_0^-=0,\qquad z_0^+=1,\qquad z_n^-=z_n^+\quad(n\geq1).
\]

Por telescopía, sus censos en el prefijo de longitud \(N>0\) son

\[
 Z_N^-:=\sum_{n=0}^{N-1}z_n^-=\lfloor N\delta\rfloor,
 \qquad
 Z_N^+:=\sum_{n=0}^{N-1}z_n^+=\lceil N\delta\rceil.
\]

Se obtiene por tanto el teorema de borde

\[
 \boxed{Z_N^+-Z_N^-=1\qquad(N>0).}
\]

La unidad no se crea al final del bloque: estaba localizada en la elección de
qué semiborde pertenece al cilindro. Si se definen las dos polarizaciones

\[
 r_N^-:=N\delta-\lfloor N\delta\rfloor,
 \qquad
 r_N^+:=\lceil N\delta\rceil-N\delta,
\]

entonces

\[
 0<r_N^\pm<1,
 \qquad
 \boxed{r_N^-+r_N^+=1.}
\]

Ésta es la primera forma exacta de la \textbf{evolución conservativa de la unidad}:
las dos proyecciones holográficas reparten una unidad de frontera y su suma la
recupera en toda profundidad finita. No se está afirmando todavía una
conservación de energía. La promoción termodinámica exige además que la
realización física del operador sea reversible. En el estado TPK completo, la
actualización biyectiva ya implica \(H(X\mid U(X))=0\); esa afirmación lógica
es distinta tanto de la entropía topológica nula de la palabra esturmiana como
de una implementación térmica concreta.

La ley es también bidireccional. La lectura inferior fija el semiborde de
salida y la superior fija el de entrada. Invertir la orientación intercambia
las dos convenciones y transporta la unidad, pero no altera la pendiente ni
el orden de largo alcance. Éste es el contenido mínimo que debe conservar
cualquier lector \emph{two-way}: la inversión cambia la residencia de la unidad de
borde, no la elimina.

\textbf{Criterio de refutación.} Un
\(N>0\) para el que \(Z_N^+-Z_N^-\neq1\), o una redacción que convierta esta
identidad aritmética, sin transductor, en una igualdad energética.

\section{Cociclo de concatenación, memoria de acarreo y restricción exacta del TPK}

Para una fase inicial \(k\in\mathbb Z\), el número de vacancias de un bloque
de longitud \(N\) es

\[
 V_N(k):=\lfloor(k+N)\delta\rfloor-\lfloor k\delta\rfloor.
\]

La palabra es balanceada, de modo que

\[
 V_N(k)\in
 \left\{\lfloor N\delta\rfloor,\lceil N\delta\rceil\right\}
\]

para toda fase. Más importante para la dinámica enriquecida es la identidad
de composición

\[
 \boxed{V_{M+N}(k)=V_M(k)+V_N(k+M).}
\]

El segundo sumando no puede evaluarse de nuevo en el origen: recibe la fase
transportada por el primer bloque. Al olvidar esa fase se pierde exactamente
el acarreo

\[
 \kappa_\delta(M,N)
 :=\lfloor(M+N)\delta\rfloor
   -\lfloor M\delta\rfloor-\lfloor N\delta\rfloor
 \in\{0,1\}.
\]

Así,

\[
 \lfloor(M+N)\delta\rfloor
 =\lfloor M\delta\rfloor+\lfloor N\delta\rfloor
  +\kappa_\delta(M,N).
\]

La asociatividad de la suma fuerza la ecuación de cociclo

\[
 \kappa_\delta(L,M)+\kappa_\delta(L+M,N)
 =\kappa_\delta(M,N)+\kappa_\delta(L,M+N).
\]

Por tanto, el acarreo no es una corrección numérica añadida al terminar el
cálculo: es el dato que hace asociativa la concatenación de bloques cuando se
conserva su historia. Ésta es precisamente la estructura que una proyección a
meros cardinales oculta.

La restricción de la dinámica TPK al subsistema de vacancias puede escribirse
sin sustituir la fibra completa. Sobre

\[
 \mathcal Y_{108}:=C_{108}\times\mathbb T\times\mathbb Z
\]

se define

\[
 \mathcal F(r,\theta,m)
 =\bigl(r+1,\{\theta+\delta\},
         m+\lfloor\theta+\delta\rfloor\bigr),
 \qquad 0\leq\theta<1.
\]

El primer componente es el reloj finito; el segundo, la fase irracional; el
tercero, la memoria integrada. El mapa es biyectivo, con inversa

\[
 \mathcal F^{-1}(r,\theta,m)
 =\bigl(r-1,\{\theta-\delta\},
 m-\lfloor\{\theta-\delta\}+\delta\rfloor\bigr).
\]

En el lenguaje operatorio rector, esta fórmula es un \textbf{factor} de
\(U_t=\operatorname{Upd}_t\circ\operatorname{Tra}_t\circ
\operatorname{Sel}_t\): el selector distingue evento y no evento; el
transporte avanza la fase nonádica y la fase irracional; la actualización
incorpora el pulso al registro de memoria. No reemplaza la fibra TPK, que conserva además
hoja, orientación, residuo, cociente, firma, cilindro, frontera, ruta y
supervivencia.

Después de \(108\) iteraciones se tiene

\[
 \mathcal F^{108}(r,\theta,m)
 =\bigl(r,\{\theta+108\delta\},m+V_{108}(\theta)\bigr),
\]

con \(V_{108}(\theta)\in\{4,5\}\). Retorna el primer reloj; no retorna ni la
fase irracional ni la memoria. Esta igualdad tipa de manera exacta la frase
«retorna la fase, no el estado» dentro del subsistema aperiódico.

\textbf{Criterio de refutación.} Que
falle la ecuación de cociclo, que \(\mathcal F\) no sea invertible o que la
proyección del factor omita un campo que luego se use para distinguir estados.

\section{El espectro renormalizado de cardinales y sus dos autorreferencias exactas}

Definamos el lector balanceado de longitud

\[
 \mathcal W_\delta(N)
 :=\{\lfloor N\delta\rfloor,\lceil N\delta\rceil\}.
\]

Este operador asigna a cada escala entera los dos únicos censos permitidos por
el corte irracional. No consulta \(\pi\), \(\varphi\), \(e\), \(\alpha\), una
masa ni una constante experimental. Al aplicarlo a cardinales ya producidos
por APP--TRIT--TPK aparecen las siguientes igualdades exactas:

\[
 \begin{aligned}
 \mathcal W_\delta(54)&=\{2,3\},&
 \mathcal W_\delta(81)&=\{3,4\},\\
 \mathcal W_\delta(90)&=\{4,5\},&
 \mathcal W_\delta(120)&=\{5,6\},\\
 \mathcal W_\delta(137)&=\{6,7\},&
 \mathcal W_\delta(202)&=\{9,10\},\\
 \mathcal W_\delta(210)&=\{9,10\},&
 \mathcal W_\delta(243)&=\{11,12\},\\
 \mathcal W_\delta(271)&=\{12,13\},&
 \mathcal W_\delta(324)&=\{14,15\},\\
 \mathcal W_\delta(369)&=\{16,17\},&
 \mathcal W_\delta(468)&=\{21,22\},\\
 \mathcal W_\delta(729)&=\{33,34\},&
 \mathcal W_\delta(1000)&=\{45,46\}.
 \end{aligned}
\]

Dos filas son autorreferenciales en un sentido preciso. El entero \(137\),
obtenido previamente como traza de la matriz de sincronización, reproduce al
ser leído por \(\mathcal W_\delta\) los retornos secundarios:

\[
 \boxed{\mathcal W_\delta(137)=\{6,7\}.}
\]

El censo demostrado \(468\) reproduce del mismo modo los huecos primarios:

\[
 \boxed{\mathcal W_\delta(468)=\{21,22\}.}
\]

No se trata de redondear \(\alpha^{-1}\) a \(137\). La traza \(137\) se
construyó antes desde la sincronización entera, y el lector de vacancias
recupera después \(6/7\). Tampoco se introduce \(21/22\) para seleccionar
\(468\): el censo pertenece a la prolongación nonádica y el mismo lector lo
devuelve como espectro de ventanas. Estas dos identidades constituyen un
cierre de renormalización dentro de la genealogía, no una definición
retrospectiva.

Hay una tercera identidad estructural particularmente limpia. Puesto que

\[
 210=90+120,
\]

la firma ordenada del intervalo es

\[
 \bigl(V_{90}(0),V_{120}(90)\bigr)=(4,5),
\]

y por tanto

\[
 Z^-_{210}=4+5=9,
 \qquad
 Z^+_{210}=1+4+5=10.
\]

Los canales \(90/120\), generados en el TPK antes de sus evaluadores
parangulares, comprimen así la transición \(9\to10\) mediante la misma unidad
de borde. La igualdad no dice que \(90\) y \(120\) sean «los números» nueve y
diez; dice que el lector de capacidad \(729/1000\), aplicado a su
concatenación causal, produce exactamente esas dos coordenadas.

\textbf{Criterio de refutación.} Una
sola discrepancia en cualquiera de los pisos o techos exactos, o el uso de un
valor reconocido externamente para escoger las longitudes.

\section{La completación cúbica contiene conjuntamente las escalas 34 y 11}

La completación

\[
 10^3=(9+1)^3=729+243+27+1
\]

posee ahora una imagen exacta bajo el cociclo de vacancias. Conservando el
orden de los cuatro sumandos, la palabra mecánica inferior produce la firma

\[
 \boxed{(33,11,1,0),}
\]

porque

\[
 \begin{aligned}
 V_{729}(0)&=33,\\
 V_{243}(729)&=11,\\
 V_{27}(972)&=1,\\
 V_1(999)&=0.
 \end{aligned}
\]

La suma es \(45\), mientras que la lectura superior coloca la unidad de borde
en el primer bloque y da

\[
 \boxed{(34,11,1,0),\qquad 34+11+1+0=46.}
\]

Así se obtiene una forma refinada de la ley \(45/46\): no sólo cuenta eventos
en una ventana de mil posiciones, sino que conserva la procedencia de cada
evento dentro de la completación \(729+243+27+1\). El orden de los sumandos es
parte del dato. Una permutación puede redistribuir el acarreo entre bloques,
aunque la suma total permanezca fija; por eso la firma enriquecida contiene
más información que el cardinal \(45\) o \(46\).

La escala \(729\) se descompone a su vez en tres bloques visibles de longitud
\(243\):

\[
 \bigl(V_{243}(0),V_{243}(243),V_{243}(486)\bigr)
 =(11,11,11).
\]

En consecuencia,

\[
 Z^-_{729}=11+11+11=33,
 \qquad
 \boxed{Z^+_{729}=1+11+11+11=34.}
\]

Este resultado sitúa \(11\) y \(34\) en una única jerarquía aperiódica antes
de toda interpretación dimensional. Además,

\[
 \begin{aligned}
 324&=243+81,& (V_{243}(0),V_{81}(243))&=(11,3),\\
 369&=243+126,&(V_{243}(0),V_{126}(243))&=(11,5),
 \end{aligned}
\]

de donde

\[
 Z^-_{324}=14,\quad Z^+_{324}=15,
 \qquad
 Z^-_{369}=16,\quad Z^+_{369}=17.
\]

Dos contactos con objetos ya demostrados son exactos en el nivel cardinal:

\[
 \boxed{Z^+_{324}=15
 =\det\begin{pmatrix}21&22\\6&7\end{pmatrix},}
\]

y

\[
 \boxed{Z^-_{369}=16=|\det H_4|=|\operatorname{coker}H_4|.}
\]

Combinados con el teorema determinantal activo,

\[
 \operatorname{ord}_{10}
 \bigl(\varphi\alpha^{16}\bigr)=-34,
\]

producen la factorización escalar

\[
 \boxed{
 -\operatorname{ord}_{10}
 \bigl(\varphi\alpha^{Z^-_{369}}\bigr)
 =Z^+_{729}.}
\]

La formulación preliminar de este resultado lo describía como un cuadrado
desclasificado a cardinales. Las Secciones 95--96 cierran su tipado: el dominio
pertinente no es una enumeración arbitraria de dieciséis vacancias, sino la
firma ordenada de la completación cúbica y su lector de profundidad. La rama
de acción factoriza exactamente como

\[
 \boxed{
 \mathcal V^-_{369}
 \xrightarrow{\;Z^-\;}16
 =|\operatorname{coker}(H_4)|
 \xrightarrow{\;\mathcal N_\alpha\;}
 \alpha^{16}
 \xrightarrow{\;\varphi\cdot\;}
 \Sigma_H
 \xrightarrow{\;-\operatorname{ord}_{10}\;}34
 =Z^+_{729}.}
\]

Aquí \(\mathcal N_\alpha(n)=\alpha^{n}\) es el lector normativo
del exponente determinantal. La forma de Smith de \(H_4\) produce el conúcleo
de orden \(16\); la sección \(\alpha\) lo transporta a la norma
\(\alpha^{16}\); la autoescala \(\varphi\) forma
\(\Sigma_H\); y el orden decimal determina la profundidad \(34\). Por tanto, la
igualdad \(Z^-_{369}=|\operatorname{coker}(H_4)|\) no se usa como una biyección
elemento a elemento: actúa como igualdad de invariantes en la composición que
demuestra la profundidad de acción. Éste es precisamente el tipo de
afirmación requerido por el teorema.

Del mismo modo,

\[
 Z^-_{243}=11
\]

es la segunda coordenada del mismo lector cúbico. Su factorización tipada es

\[
 \boxed{
 \mathcal V^-_{243}
 \xrightarrow{\;Z^-\;}11
 =\operatorname{pr}_{243}\!\left(
   \mathcal D_{\rm cub}(\sigma^-,\sigma^+)\right)
 \xrightarrow{\;\operatorname{Pub}_G\;}
 \mathcal D_{10}(G)=11.}
\]

El ensamblaje con la rama de acción da, sin consulta de valores objetivo,

\[
 \boxed{
 \bigl(\mathcal D_{10}(\hbar),\mathcal D_{10}(G)\bigr)=(34,11),
 \qquad
 \mathcal D_{10}(G\otimes\hbar)=34+11=45,
 \qquad 45+1=46.}
\]

La identidad \(3^6\cdot243=3^{11}\) conserva su dominio propio como lector de
dimensión; la coordenada \(11\) de la completación y la profundidad
gravitatoria se conectan mediante \(\mathcal D_{\rm cub}\) y
\(\operatorname{Pub}_G\), no mediante una identificación tácita de sus
elementos. El producto \(G_a\hbar_a\), demostrado en la Sección 95, cierra el
cuadrado dinámico: el retorno two-way invierte recíprocamente ambas ramas y
conserva su producto.

\textbf{Criterio de refutación.} Que las firmas
ordenadas no telescopen, que el determinante o el conúcleo no valgan lo
declarado, que \(\mathcal D_{\rm cub}\) no produzca \((34,11,45,46)\), o que
la realización dimensional dependa de \(\hbar\), \(G\), CODATA o de la
posición decimal que precisamente debe generar.

\section{La escisión electrónica del ambiente y la aparición interna de 9, 10 y 24}

En la primera reapertura conjunta, el canal electrónico contiene \(527\)
extensiones de un ambiente de \(729\). Su déficit es

\[
 202=729-527=6+28\cdot7.
\]

La descomposición ordenada del ambiente

\[
 729=202+527
\]

es compatible de manera exacta con el lector aperiódico:

\[
 \bigl(V_{202}(0),V_{527}(202)\bigr)=(9,24).
\]

Por tanto,

\[
 \boxed{Z^-_{729}=9+24=33,\qquad
        Z^+_{729}=10+24=34.}
\]

El mismo déficit que singulariza la ventana electrónica reproduce, bajo
\(\mathcal W_\delta\), la frontera nonádica--decimal:

\[
 \boxed{\mathcal W_\delta(202)=\{9,10\}.}
\]

La unidad de borde actúa aquí sobre el sector deficitario; el sector
superviviente conserva veinticuatro eventos en la fase canónica. Esta
descomposición refina la igualdad anterior \(34=1+3\cdot11\): ahora la misma
coordenada se lee como \(34=10+24\). Las dos factorizaciones coexisten porque
registran cortes distintos del mismo cilindro, no porque una sustituya a la
otra.

El entero \(24\) coincide con el rango del retículo de Leech, pero la igualdad
cardinal no identifica las \(527\) extensiones electrónicas con vectores del
retículo. Su contenido exacto es un contacto de rango en la segunda
proyección. El entrelazamiento excepcional se construye, con incidencia y
forma cuadrática propias, mediante la cadena Paley--Witt--Golay--Leech;
ninguna acción del grupo Monstruo se atribuye a dígitos ni a vacancias.

\textbf{Criterio de refutación.} Que el censo
electrónico no sea \(527\), que el bloque complementario no contenga
veinticuatro eventos o que ningún transporte preserve la incidencia
excepcional exigida.

\section{Los dos relojes finitos y la suspensión helicoidal aperiódica}

El reloj \(108\) se divide en dos semiciclos de \(54\) pasos. La palabra
inferior tiene firma

\[
 \bigl(V_{54}(0),V_{54}(54)\bigr)=(2,2),
\]

de modo que

\[
 Z^-_{108}=2+2=4,
 \qquad Z^+_{108}=1+2+2=5.
\]

El retorno \(54+54\) cierra la fase finita y conserva una elección de
semiborde. La repetición triple del reloj no repite, sin embargo, la misma
firma:

\[
 \bigl(V_{108}(0),V_{108}(108),V_{108}(216)\bigr)=(4,5,5),
\]

y por ello

\[
 Z^-_{324}=4+5+5=14,
 \qquad Z^+_{324}=1+4+5+5=15.
\]

Ésta es una manifestación directa de orden aperiódico: tres bloques de igual
longitud poseen censos balanceados, pero no idénticos. La secuencia conserva
orden de largo alcance y discrepancia menor que uno, sin adquirir período
traslacional.

El cardinal \(324\) admite dos lecturas internas ya presentes en el construcción:
\(324=81\cdot4\) para las rutas orientadas y \(324=108\cdot3\) para un
levantamiento trítico del reloj. La segunda factorización define el modelo
cíclico mínimo

\[
 0\longrightarrow C_3\longrightarrow C_{324}
 \longrightarrow C_{108}\longrightarrow0.
\]

La extensión no es escindida: \(C_{324}\) contiene un elemento de orden
\(324\), mientras \(C_3\times C_{108}\) tiene exponente \(108\). En el reloj
levantado, una traslación de \(108\) pasos retorna a la misma coordenada de
\(C_{108}\), pero avanza a la siguiente hoja del núcleo \(C_3\); sólo tras
\(324\) pasos retorna la coordenada levantada. Éste es el modelo algebraico
más pequeño de una hélice de tres hojas sobre el reloj \(108\).

Al decorar este reloj por la rotación irracional de pendiente \(\delta\) se
obtiene una suspensión helicoidal esturmiana: la base finita puede retornar,
la hoja levantada puede retornar a mayor periodo y la fase irracional no
retorna jamás exactamente. La memoria acumulada distingue esos tres hechos.
La hélice ascendente usa \(\mathcal F\); la descendente usa
\(\mathcal F^{-1}\). Son lecturas conjugadas de la misma dinámica reversible,
no dos series independientes. El espejo \(\pi/e\) sigue siendo otra
involución: actúa sobre las ventanas de salida y fija \(\varphi\); no invierte
por sí mismo la orientación temporal.

\textbf{Criterio de refutación.} Una sección grupal de \(C_{324}\to C_{108}\), una repetición
\((4,4,4)\) o una conjugación que borre la memoria.

\section{La primera supercelda que cierra el reloj 108}

La fracción continua de \(\delta\) comienza

\[
 [0;21,1,5,1,6,2,3,1,1,8,1,2,\ldots].
\]

El primer convergente cuyo denominador es divisible por \(108\) es

\[
 \boxed{\frac{7813}{170748}.}
\]

No hay un denominador convergente anterior con esa propiedad. Además, el
certificado racional prueba que

\[
 \left\lfloor n\delta\right\rfloor
 =\left\lfloor\frac{7813n}{170748}\right\rfloor
 \qquad(0\leq n<170748),
\]

mientras en el extremo final

\[
 \lfloor170748\delta\rfloor=7812,
 \qquad
 \frac{7813\cdot170748}{170748}=7813.
\]

La palabra aperiódica y su aproximante periódico de Christoffel coinciden en
todo prefijo propio de la supercelda y difieren únicamente en el cierre. El
aproximante periódico no corrige una nube de errores: transporta exactamente
la unidad de borde desde la proyección inferior a la superior.

El denominador posee la factorización

\[
 \boxed{
 170748=108\cdot1581
       =324\cdot527
       =2\cdot54\cdot3\cdot527.}
\]

Aquí \(527\) es exactamente el censo del canal electrónico en la primera
reapertura y \(324\) el cardinal de las rutas orientadas. Como
\(\gcd(324,527)=1\), existe la descomposición canónica de relojes

\[
 C_{170748}\cong C_{324}\times C_{527}.
\]

El teorema chino del resto justifica esta descomposición grupal. El segundo
factor tiene el cardinal de la fibra electrónica y el primero el cardinal del
atlas orientado; esta coincidencia de cardinales no identifica los objetos ni
traslada automáticamente sus acciones.

La pareja de relojes ofrece un segundo dato inesperado. En la rama inferior,
después de \(170748\) bloques,

\[
 7812\equiv36\pmod{108},
 \qquad
 170748-7812=162936\equiv72\pmod{108}.
\]

Por tanto, aunque el reloj de bloques cierre,

\[
 \boxed{(r_{\rm vac},r_{\rm cap})=(36,72),\qquad36+72=108.}
\]

La rama superior transfiere la unidad de borde:

\[
 (r_{\rm vac}^+,r_{\rm cap}^+)=(37,71),
 \qquad37+71=108.
\]

La aparición de \(36\) en la rama inferior coincide con el índice del estadio
\(R_{36}\) de la prolongación coinductiva. Esta congruencia no identifica por
sí sola el residuo de vacancias con la coordenada de \(R_{36}\); demuestra que
la supercelda cierra el reloj de bloque, conserva el reparto de la unidad y
deja una memoria residual no trivial en el reloj de capacidad.

El numerador posee otra factorización interna:

\[
 7813=13\cdot601,
 \qquad
 7812=36\cdot217,
 \qquad217=2\cdot108+1.
\]

De aquí surge el marco unimodular

\[
 \boxed{
 M_\star=
 \begin{pmatrix}
 601&217\\
 36&13
 \end{pmatrix}\in\operatorname{SL}_2(\mathbb Z),
 \qquad
 601\cdot13-217\cdot36=1.}
\]

Los cuatro enteros están localizados en la arquitectura: \(601\) es la
duodécima coordenada del sello dodecafásico \(K\); \(13\) es el número de
multisecciones del atlas; \(36\) es el estadio \(R_{36}\); y \(216\) es el
retorno cinemático extendido. La igualdad es exacta y no utiliza un valor
físico. Define un marco unimodular de sus cardinales; no postula, por esa sola
razón, una equivalencia categórica entre los cuatro objetos.

\textbf{Criterio de refutación.} Un convergente previo con
denominador múltiplo de \(108\), una discrepancia en un prefijo propio, una
factorización falsa o un transporte que no conjugue las acciones.

\section{La segunda derivación de \(\alpha\) como renormalización de una sola palabra}

El polinomio \(P_9\) de la segunda vía de \(\alpha\) contiene los pares
\(21/22\), \(6/7\) y \(45/46\). Hasta ahora se habían reconstruido sus
coeficientes por la matriz de sincronización, los retornos de Beatty y la
frontera \(7{:}9\). El nuevo lector demuestra que los tres pares son además
imágenes de \textbf{la misma} palabra mecánica en tres escalas internas:

\[
 \boxed{
 \mathcal W_\delta(468)=\{21,22\},\qquad
 \mathcal W_\delta(137)=\{6,7\},\qquad
 \mathcal W_\delta(1000)=\{45,46\}.}
\]

La primera longitud es el censo de emisiones decimales, la segunda es la
traza entera de sincronización y la tercera es la completación cúbica. Por
ello, la genealogía de la segunda vía puede escribirse sin introducir sus
coeficientes como una lista independiente:

\[
 \begin{aligned}
 &\mathrm{APP}\longrightarrow\mathrm{TRIT}\longrightarrow\mathrm{TPK}
 \longrightarrow X^{\rm enr},\\
 &729\longrightarrow1000\longrightarrow\delta
 \longrightarrow\mathcal W_\delta,\\
 &(468,137,1000)
 \longmapsto
 \bigl(\{21,22\},\{6,7\},\{45,46\}\bigr),\\
 &\text{firma de sincronización y borde }7{:}9
 \longrightarrow P_9
 \longrightarrow\alpha.
 \end{aligned}
\]

Esta cadena no sustituye la primera vía dodecafásica. Ambas parten del estado
coinductivo que también determina \(\pi\), \(\varphi\) y
\(e\); después conservan lectores propios. La raíz de \(P_9\) se
calcula al final, una vez construida la firma. El valor de \(\alpha\) nunca
selecciona \(468\), \(137\), \(1000\) ni sus ventanas.

El numerador \(7813=13\cdot601\) sitúa además un posible cierre entre la
supercelda aperiódica y el sello dodecafásico: \(601=K_{12}\) pertenece a la
primera vía, mientras \(13\) pertenece al atlas posterior. La identidad
unimodular del apartado anterior permite formular la pregunta correctamente:
hay que determinar si el módulo de memoria sobre el que actúan \(D_3K\),
\(R_{36}\), el retorno \(216\) y las trece multisecciones admite el cambio de
base \(M_\star\). No se debe convertir la coincidencia de factores en prueba,
pero tampoco perderla como una mera curiosidad: está localizada en la primera
supercelda seleccionada por el reloj \(108\), no en un entero buscado después.

\textbf{Criterio de refutación.} Que alguno de los pares no proceda de
\(\mathcal W_\delta\), que \(P_9\) intervenga en la selección de sus propias
entradas o que el módulo propuesto para \(M_\star\) no conserve el registro de memoria.

\section{Profundización de las semillas y separación rigurosa de las escalas físicas}

El espacio de una semilla orientada posee cardinal

\[
 |\mathcal S|=9^2\cdot4=324,
\]

y el barrido de parejas tiene

\[
 |\mathcal S\times\mathcal S|=324^2=104976.
\]

Aplicar el lector aperiódico al cardinal total da

\[
 \mathcal W_\delta(104976)=\{4803,4804\}.
\]

Esta igualdad es un censo de la recta mecánica, no todavía un censo de
semillas. Para convertirla en una propiedad del barrido APP hay que fijar una
ordenación canónica

\[
 \operatorname{ord}_{\mathcal S^2}:
 \mathcal S\times\mathcal S\overset\sim\longrightarrow
 \{0,\ldots,104975\}
\]

compatible con las dos hojas, la acción diedral, el selector TRIT y la
actualización TPK. Sólo entonces puede preguntarse qué semillas ocupan las
vacancias y cómo cambia el censo bajo refinamiento interior o expansión
exterior. La cardinalidad por sí sola no debe seleccionar esa ordenación.

El nuevo análisis, sin embargo, ya determina dónde buscar las escalas
\(10^{-11}\) y \(10^{-34}\). No se obtienen elevando arbitrariamente una
semilla a una potencia decimal. El cristal produce primero los enteros

\[
 Z^-_{243}=11,
 \qquad
 Z^+_{729}=34,
\]

dentro de la misma firma de completación. Después, dos transductores distintos
pueden realizar esos órdenes en rectas dimensionales:

\[
 \begin{aligned}
 \mathscr T_G&:\bigl(X^{\rm enr},\mathcal V^-_{243}\bigr)
 \longrightarrow 10^{-11}\mathcal L_G,\\
 \mathscr T_\hbar&:\bigl(X^{\rm enr},\mathcal V^+_{729},G_H,\alpha,\varphi\bigr)
 \longrightarrow10^{-34}\mathcal L_S.
 \end{aligned}
\]

El segundo transductor posee ya una vía probada mediante
\(|G_H|=16\) y \(\varphi\alpha^{16}\). El primero conserva en la autoridad una
carta decimal y criterios de selección, y debe ser reconciliado con el lector
de vacancias sin usar CODATA. La frase «fractalizar hacia dentro hasta
\(10^{-11}\) y \(10^{-34}\)» adquiere así un significado matemático preciso:
refinar el mismo estado hasta dos cardinales de ventana y sólo después aplicar
la transducción dimensional correspondiente. Las potencias de diez pertenecen
a la carta de la recta dimensional, no a la semilla desnuda.

La conservación de la unidad se mantiene durante ese paso si el transductor
conserva simultáneamente las dos proyecciones y el cociclo. El requisito es un
cuadrado de naturalidad

\[
 \mathscr T\circ\operatorname{Upd}
 =\operatorname{Upd}_{\mathcal L}\circ\mathscr T
\]

en el que el lado izquierdo actualiza primero la memoria discreta y el derecho
transporta primero la sección dimensional. Esta ecuación, y no una igualdad
decimal aislada, es la prueba exigible de que una constante física emerge «en
la escala adecuada».

La estructura discreta del continuo debe recibir todas estas fibras a la vez.
El refinamiento de semillas, la resolubilidad de extensiones, la construcción
de medida, la coherencia y el lector determinantal no constituyen cinco
objetos separables: son cinco aspectos inseparables del mismo paso al límite.
El cristal aperiódico aporta el calendario balanceado; APP conserva suma,
producto, residuo, cociente y acarreo; TRIT tipa la orientación y las tres
curvaturas; TPK conserva la historia durante selección, transporte y
actualización. Sólo su ensamblaje produce la recta real como salida. La recta
\(\mathbb R\) no es, en esta genealogía, el soporte que se presupone para
construir el protocolo: es una de sus realizaciones terminales.

\textbf{Criterio de refutación.} Una ordenación que no sea natural respecto de APP,
TRIT y TPK; una potencia decimal elegida desde el valor físico; o un
transductor que no preserve la unidad de borde y el registro de memoria.

\section{La corona celular del continuo contiene el reloj 108 y la firma TRIT}

La revisión dla construcción de referencia celular aporta una pieza que enlaza la
completación, el reloj y la doble proyección con mayor rigidez que una
coincidencia de cifras. Para el par cubical

\[
 Q_{9,3}\subset Q_{10,3},
\]

el complejo relativo posee el vector de cardinales

\[
 \mathbf c=(c_0,c_1,c_2,c_3)=(271,756,702,217).
\]

Su característica de Euler es nula:

\[
 \chi(\mathbf c)=271-756+702-217=0.
\]

La nulidad, sin embargo, no es una cancelación opaca. Cada coordenada admite
una descomposición en cardinales que ya pertenecen a la arquitectura:

\[
 \begin{aligned}
 271&=243+27+1,\\
 756&=729+27,\\
 702&=729-27,\\
 217&=216+1=2\cdot108+1.
 \end{aligned}
\]

Al emparejar grados extremos e interiores aparecen dos defectos iguales:

\[
 \boxed{c_0-c_3=271-217=54,}
 \qquad
 \boxed{c_1-c_2=756-702=54.}
\]

Por consiguiente,

\[
 \boxed{(c_0-c_3)+(c_1-c_2)=54+54=108.}
\]

El reloj \(108\) no se ha impuesto a la corona desde un calendario exterior:
su doble semiciclo está contenido en la diferencia entre los grados de borde
y en la diferencia entre los grados interiores. La característica nula se
reescribe como

\[
 \chi(\mathbf c)=(c_0-c_3)-(c_1-c_2)=54-54=0.
\]

Esta identidad distingue dos sentidos. El primer \(54\) mide la separación
entre los extremos del complejo; el segundo mide la separación interna. Su
igualdad hace compatible el cierre: lo que sale por un emparejamiento vuelve
por el otro sin borrar el grado celular. La suma \(54+54\) mide la longitud
del recorrido bidireccional; la diferencia \(54-54\) certifica su cierre de
Euler. Son dos operaciones distintas sobre la misma pareja y no deben
aplanarse en el único número \(108\).

Ahora se aplica a cada grado el lector aperiódico

\[
 \mathcal W_\delta(N)=
 \{\lfloor N\delta\rfloor,\lceil N\delta\rceil\}.
\]

Las dos proyecciones de vacancia de la corona son

\[
 \mathbf v^-=(12,34,32,9),
 \qquad
 \mathbf v^+=(13,35,33,10).
\]

Sus complementos de capacidad son

\[
 \mathbf b^+=\mathbf c-\mathbf v^-
 =(259,722,670,208),
\]

\[
 \mathbf b^-=\mathbf c-\mathbf v^+
 =(258,721,669,207).
\]

En ambos cierres unilaterales se verifica, componente a componente,

\[
 \mathbf c=\mathbf v^-+\mathbf b^+
           =\mathbf v^++\mathbf b^-.
\]

La información decisiva aparece al tomar la característica alternada:

\[
 \boxed{\chi(\mathbf v^-)=\chi(\mathbf v^+)=+1,}
\]

\[
 \boxed{\chi(\mathbf b^-)=\chi(\mathbf b^+)=-1.}
\]

Por tanto, la corona de carga total nula se polariza exactamente como

\[
 \boxed{0=(+1)+(-1).}
\]

La independencia respecto de la convención de semiborde es estructural:

\[
 \mathbf v^+-\mathbf v^-=(1,1,1,1),
 \qquad
 \chi(1,1,1,1)=0.
\]

Mover la unidad de borde entre los dos cierres no altera la carga de Euler de
la rama. Se obtiene así una realización interna del alfabeto TRIT:

\[
 \boxed{
 \text{proyección de vacancia}\mapsto+1,
 \quad
 \text{corona completa}\mapsto0,
 \quad
 \text{proyección de capacidad}\mapsto-1.}
\]

Esto no sustituye la definición general del TRIT ni identifica por sí solo
las tres cargas con campos físicos. Prueba algo más elemental y anterior: la
doble proyección de un único objeto celular genera, sin añadir símbolos, las
tres clases orientadas (+1,0,-1). Las álgebras

\[
 \mathbb A_\tau=\mathbb R[J_\tau]/(J_\tau^2+\tau I)
\]

pueden actuar después sobre esas clases. El tipo elíptico, parabólico o
hiperbólico ya no se selecciona desde una geometría continua presupuesta: la
firma que lo indexa ha emergido de la polarización discreta del complejo.

\textbf{Criterio de refutación.} Que cambie cualquiera de los
cuatro cardinales, que uno de los defectos deje de valer \(54\), que las
proyecciones no tengan cargas (+1/-1), o que se use una geometría clásica
para seleccionar retroactivamente el signo.

\section{Marco polarizado, raíz de borde y evolución conservativa de la unidad}

El resultado anterior es un caso de una ley exacta válida para toda longitud.
Sea

\[
 a_N=\lfloor N\delta\rfloor,
 \qquad N>0.
\]

Las dos convenciones de cierre determinan los vectores

\[
 u_N^-=(a_N,N-a_N),
 \qquad
 u_N^+=(a_N+1,N-a_N-1).
\]

Cada fila conserva exactamente el bloque:

\[
 (1,1)u_N^-=(1,1)u_N^+=N.
\]

La diferencia es independiente de (N):

\[
 \boxed{u_N^+-u_N^-=(1,-1).}
\]

La unidad no aparece ni desaparece; se transporta de la coordenada de
capacidad a la de vacancia. El vector ((1,-1)) genera el núcleo de la
augmentación

\[
 \mathbb Z^2\xrightarrow{\ (x,y)\mapsto x+y\ }\mathbb Z,
\]

es decir, la raíz \((A_1)\). La doble proyección posee por tanto una dirección
de frontera canónica. No se necesita elegir una métrica física para definirla.

El marco de las dos polarizaciones tiene determinante

\[
 \boxed{
 \det\begin{pmatrix}
 a_N&N-a_N\\
 a_N+1&N-a_N-1
 \end{pmatrix}=-N.}
\]

La demostración es inmediata:

\[
 a_N(N-a_N-1)-(N-a_N)(a_N+1)=-N.
\]

Así, el traslado de una sola unidad de borde barre un paralelogramo entero de
área orientada (N). Tras normalizar por (N), las dos filas viven en el
símplice (x+y=1), se separan por ((1,-1)/N) y convergen hacia

\[
 (\delta,1-\delta)=(\delta,\rho).
\]

Ésta es una forma precisa de la evolución conservativa de la unidad
dimensional: en cada escala hay una unidad normalizada; el refinamiento no la
multiplica, sino que reduce el tamaño de su cuanto de frontera mientras
conserva la dirección integral que permite recuperarlo.

La trayectoria completa se escribe en el retículo como

\[
 \gamma(N)=
 \bigl(\lfloor N\delta\rfloor,
       N-\lfloor N\delta\rfloor\bigr).
\]

Cada paso incrementa exactamente una coordenada:

\[
 \gamma(N+1)-\gamma(N)\in\{(1,0),(0,1)\}.
\]

Además,

\[
 \left\|\gamma(N)-N(\delta,\rho)\right\|_\infty<1.
\]

La hélice aperiódica es, en su sección reticular, una recta digital de
discrepancia uniformemente acotada. Posee orden de largo alcance porque nunca
se aleja más de una celda de la dirección límite, pero carece de período
traslacional porque \(\delta\) es irracional. La palabra esturmiana no adorna
esa trayectoria: es la sucesión de decisiones que indica cuál de las dos
coordenadas recibe la unidad siguiente.

La misma estructura de augmentación reaparece en los tres canales. El vector
de reapertura

\[
 v_\Omega=(3,-7,4)
\]

satisface (3-7+4=0), por lo que pertenece al retículo

\[
 A_2=\ker\bigl(\mathbb Z^3\xrightarrow{x+y+z}\mathbb Z\bigr).
\]

La raíz \((A_1)\) gobierna la transferencia entre las dos proyecciones; el
defecto \((A_2)\) gobierna la redistribución entre las tres ventanas
correlacionadas. No son el mismo vector ni viven en el mismo rango, pero son
dos grados sucesivos de una única ley de conservación por augmentación. Esto
explica por qué la doble lectura y la orientación ternaria pueden coexistir
sin convertir la teoría en una codificación binaria: la conservación fija la
diagonal, mientras el dato HMT reside en su complemento orientado.

\textbf{Criterio de refutación.} Una escala cuyo marco no tenga determinante (-N), un paso que
modifique ambas coordenadas a la vez, o una identificación de \((A_1)\) y
\((A_2)\) que no declare el mapa de cambio de rango.

\section{Carga topológico-temporal de borde y selección no destructiva de curvatura}

La derivada trítica de la palabra de vacancias,

\[
 \tau_n=z_{n-1}-z_n\in\{+1,0,-1\},
\]

posee una ley de conservación que completa su interpretación como unidad de
información topológico-temporal. Para todo intervalo finito \((a\leq n\leq b)\),

\[
 \boxed{
 Q_{[a,b]}:=\sum_{n=a}^{b}\tau_n
 =z_{a-1}-z_b\in\{-1,0,+1\}.}
\]

La carga orientada del interior telescopa por completo y queda soportada en
los dos extremos. El TRIT visible no acumula una carga extensiva arbitraria:
registra cómo el bloque se abre, persiste o cierra respecto de su frontera.
El registro de memoria del TPK conserva los términos interiores que la suma alternada
oculta, de modo que la telescopía no equivale a pérdida de historia.

Como \((0<\delta<1/2)\), dos vacancias no son consecutivas. Los símbolos no
nulos de \(((\tau_n))\) alternan por ello estrictamente:

\[
 -1,+1,-1,+1,\ldots
\]

cuando se suprimen los ceros intermedios. Una apertura debe cerrarse antes de
que pueda producirse otra apertura. Los ceros no son ausencia de dinámica:
son las mesetas durante las cuales el régimen se conserva mientras el reloj
de bloque continúa avanzando. Al inducir la dinámica sobre los eventos, esas
mesetas poseen longitudes \(6\) o \(7\), tal como ya se demostró.

Cada símbolo selecciona una de las tres exponenciales

\[
 \exp(tJ_\tau)=C_\tau(t)I+S_\tau(t)J_\tau,
 \qquad J_\tau^2=-\tau I.
\]

La palabra ordena así rotaciones elípticas, transportes nilpotentes y
aperturas hiperbólicas. La selección no destruye la aperiodicidad porque la
aplicación diferencial es invertible una vez conservado un bit de frontera:

\[
 z_n=z_0-\sum_{k=1}^{n}\tau_k.
\]

Ésta es la razón exacta por la que el TRIT puede leer tres curvaturas sin
reemplazar la palabra que las genera. La unidad trítica es el cambio visible;
el TPK conserva el dato de frontera, la hoja APP, el acarreo y el orden de los
operadores necesarios para invertir la lectura.

La reversión temporal

\[
 C_{\rm rev}(\tau_1,\ldots,\tau_n)
 =(-\tau_n,\ldots,-\tau_1)
\]

intercambia apertura y cierre, preserva el umbral y satisface
\((C_{\rm rev}^2=I)\). En la fibra matricial invierte también el orden del
producto

\[
 \mathcal A_N=exp(t_NJ_{\tau_N})\cdots
               \exp(t_1J_{\tau_1}).
\]

La doble dirección temporal no es entonces la lectura de una lista al revés.
Es la conjugación de una palabra, sus signos y el orden no conmutativo de sus
transportes, manteniendo la frontera que permite reconstruir el estado.

\textbf{Criterio de refutación.} Dos aperturas consecutivas en la palabra inducida, una carga de
intervalo que no sea de borde o una reconstrucción que prescinda de \((z_0)\).

\section{Torre \(S\)-ádica no estacionaria de los dos relojes}

La fracción continua de la pendiente interna comienza

\[
 \delta=[0;21,1,5,1,6,2,3,1,1,8,1,2,\ldots].
\]

Para cada coeficiente positivo (a), definamos el operador de continuante

\[
 A(a)=\begin{pmatrix}a&1\\1&0\end{pmatrix},
 \qquad\det A(a)=-1.
\]

Si \((p_n/q_n)\) es el (n)-ésimo convergente, entonces

\[
 A(a_1)\cdots A(a_n)
 =\begin{pmatrix}q_n&q_{n-1}\\p_n&p_{n-1}\end{pmatrix}.
\]

Los primeros estadios son

\[
 \begin{aligned}
 A(21)&=\begin{pmatrix}21&1\\1&0\end{pmatrix},\\
 A(21)A(1)&=\begin{pmatrix}22&21\\1&1\end{pmatrix},\\
 A(21)A(1)A(5)&=\begin{pmatrix}131&22\\6&1\end{pmatrix},\\
 A(21)A(1)A(5)A(1)&=\begin{pmatrix}153&131\\7&6\end{pmatrix}.
 \end{aligned}
\]

Esto demuestra que \(21/22\) y \(6/7\) no son dos regularidades puestas una
junto a otra. Los enteros (21,22) son los primeros denominadores del
esqueleto diofántico; (6,7) son los numeradores que aparecen dos
inducciones después:

\[
 \begin{aligned}
 131&=5\cdot22+21,&6&=5\cdot1+1,\\
 153&=1\cdot131+22,&7&=1\cdot6+1.
 \end{aligned}
\]

Por eso los huecos primarios \(21/22\) y los retornos secundarios \(6/7\)
pertenecen a una misma renormalización. Los convergentes siguientes son

\[
 \frac{48}{1049},\quad
 \frac{103}{2251},\quad
 \frac{357}{7802},\quad
 \frac{460}{10053},\quad
 \frac{817}{17855},\quad
 \frac{6996}{152893},\quad
 \frac{7813}{170748}.
\]

Cada par consecutivo conserva una celda reticular primitiva:

\[
 \boxed{p_nq_{n-1}-p_{n-1}q_n=(-1)^{n-1}.}
\]

El primer denominador de esta torre divisible por \(108\) es \(170748\), y
su aparición conserva la recurrencia completa:

\[
 \boxed{
 170748=152893+17855,
 \qquad
 7813=6996+817.}
\]

En ese estadio,

\[
 \begin{pmatrix}
 170748&152893\\
 7813&6996
 \end{pmatrix}
\in GL_2(\mathbb Z),
 \qquad\det=-1.
\]

La torre conserva por ello la unidad reticular incluso cuando el tamaño de
la supercelda crece hasta cerrar el reloj finito. El cierre de \(C_{108}\)
no hace periódica la dirección irracional: el aproximante racional coincide
con la palabra en todo prefijo propio y transfiere una unidad sólo en el
último semiborde.

La renormalización es además genuinamente ordenada. Para \(a\ne b\),

\[
 \boxed{
 A(a)A(b)-A(b)A(a)
 =(a-b)\begin{pmatrix}0&1\\-1&0\end{pmatrix}.}
\]

Por ejemplo, \(a=21\) y \(b=5\) producen el defecto orientado

\[
 16\begin{pmatrix}0&1\\-1&0\end{pmatrix}.
\]

La rotación \((x\mapsto x+\delta)\) sobre el círculo sigue siendo abeliana. La
no conmutatividad aparece al ordenar los operadores de renormalización y, en
la fibra completa, al ordenar las hojas APP y los generadores de curvatura.
Esta distinción elimina dos errores posibles: llamar no abeliana a la
rotación basal o reducir toda la dinámica a una rotación abeliana olvidando
el cociclo.

Como \(\delta\) es trascendente, no es una irracional cuadrática. Por el
teorema de Lagrange, su fracción continua no es eventualmente periódica. La
directiva de continuantes no se estabiliza en una sustitución periódica
única: el cristal es \(S\)-ádico y no estacionario. Conserva, no obstante,
complejidad lineal, entropía nula y determinantes unimodulares. Ésa es una
definición mucho más fuerte de «cristal aperiódico» que la mera ausencia de
repetición: hay renormalización a todas las escalas, pero no hay una celda de
renormalización final que se repita indefinidamente.

\textbf{Criterio de refutación.} Una discrepancia
en los coeficientes, un determinante distinto de \(\pm1\), un denominador
convergente anterior divisible por \(108\), o una directiva eventualmente
periódica.

\section{Genealogía cardinal de la corona, la frontera \(369\) y el sello dodecafásico}

La torre permite resolver los cardinales estructurales sin seleccionarlos por
parecido decimal. Los siguientes pares proceden del mismo lector y poseen
derivaciones independientes:

\[
 \begin{aligned}
 \mathcal W_\delta(217)&=\{9,10\},\\
 \mathcal W_\delta(271)&=\{12,13\},\\
 \mathcal W_\delta(601)&=\{27,28\},\\
 \mathcal W_\delta(973)&=\{44,45\}.
 \end{aligned}
\]

El entero \(217\) es el cardinal de las \(3\)-celdas de la corona y también
\(2\cdot108+1\). Su lectura determina la frontera entre la base nonádica y la
completación decimal. El entero \(271\) es la corona cúbica
\(1000-729\); su lectura determina \(12/13\), precisamente la pareja formada
por la dodecafase y el número de multisecciones del atlas. El entero \(601\)
es la duodécima coordenada del sello \(K\); su lectura determina \(27/28\), la
frontera cúbica y el cardinal de la microfibra. Finalmente,

\[
 973=271+702=756+217
\]

es la suma común de los grados pares e impares de la corona, y

\[
 \mathcal W_\delta(973)=\{44,45\}.
\]

Aquí \(44=396/9\) es el censo APP reducido y \(45\) es el determinante de la
matriz de borde y el censo de la región radical. El balance de Euler del
continuo se proyecta, por tanto, sobre dos cardinales ya situados en APP y en
la rama radical.

Estas igualdades no declaran que una multisección sea literalmente una
vacancia ni que \((K_{12})\) sea un conjunto de veintiocho elementos. El
contenido demostrado es más preciso: un único functor cardinal de corte,
construido antes de esas realizaciones, envía cuatro objetos tipados a
cuatro pares internos con residencia conocida. El siguiente paso no consiste
en buscar más semejanzas, sino en levantar este functor cardinal a las
acciones que cada representación conserva.

La frontera \(369\) admite una genealogía todavía más rígida. La construcción
distingue el cardinal \(386\) de la frontera interna y los \(387\) pares
\((K,K+9)\) por canal. El lector da

\[
 \mathcal W_\delta(386)=\mathcal W_\delta(387)=\{17,18\}.
\]

Usando cierres opuestos,

\[
 \boxed{386=17+369,\qquad387=18+369.}
\]

Así, \(369\) es la capacidad común que sobrevive al traslado de la unidad de
borde entre dos cardinales consecutivos. La transición

\[
 (386,17)\longleftrightarrow(387,18)
\]

incrementa a la vez el ambiente y la vacancia, pero deja fija la capacidad
radical \(369\). Esta es una definición exacta de estabilidad bajo cambio de
semiborde. Se suma, sin sustituirla, a las genealogías ya demostradas

\[
 369=396-27,
 \qquad
 369=3(132-6)-9,
 \qquad
 369+126=495.
\]

La convergencia de cuatro vías independientes —frontera adyacente, cuadrado
APP, retornos inducidos y transferencia conservativa— convierte \(369\) en
un nodo interno de la arquitectura, no en una cadena decimal escogida.

Otras proyecciones cardinales conservan un estatuto distinto:

\[
 \mathcal W_\delta(756)=\{34,35\},
 \qquad
 \mathcal W_\delta(702)=\{32,33\}.
\]

El \(34\) coincide con la profundidad de acción ya derivada y el \(33\) con
el censo inferior de \(729\). La derivación física de la década \(-34\)
pertenece al conúcleo de \(H_4\) y a \(\varphi\alpha^{16}\); la lectura de
las \(756\) aristas no la reemplaza. La misma corona que produce el trit
euleriano contiene así una proyección cardinal adicional de esa profundidad,
sin identificar el lector de aristas con la línea de acción.

El marco unimodular de la supercelda adquiere asimismo una interpretación algebraica adicional
fuerte:

\[
 M_\star=
 \begin{pmatrix}601&217\\36&13\end{pmatrix},
 \qquad\det M_\star=1.
\]

Sus cuatro entradas son ahora: la coordenada dodecafásica \((K_{12}=601)\), el
cardinal de \(3\)-celdas de la corona \(217\), el estadio \(R_{36}\) y las
trece multisecciones. La identidad \(217=216+1\) permanece compatible con
esta lectura. El determinante unitario demuestra la conservación reticular
del marco formado por los cuatro cardinales, sin identificar sus dominios.

\textbf{Criterio de refutación.} Que cualquiera de las
ventanas cambie, que \(369\) no sea el complemento común, o que un mapa
propuesto preserve cardinales pero destruya incidencia, acción o memoria.

\section{El núcleo del cristal y la salida del continuo}

Las piezas anteriores permiten precisar la tesis rectora sin confundir el
resultado con una metáfora cosmológica. APP suministra un soporte finito con
dos evaluaciones inseparables. La hoja aditiva acumula; la multiplicativa
pliega; residuo, cociente y acarreo conservan la diferencia de sus historias.
TRIT extrae la orientación local de esa diferencia y la realiza en tres tipos
cuadráticos. TPK convierte las firmas locales en dinámica mediante selección,
transporte y actualización, conservando fase, hoja, frontera y registro de memoria. La
completación \(729\to1000\) induce entonces una recta digital de pendiente
trascendente; su palabra de vacancias es aperiódica, balanceada y de entropía
nula; su derivada es trítica; su torre de renormalización es \(S\)-ádica; y su
suspensión sobre \(C_{108}\) es helicoidal.

El continuo no entra en esta cadena como \(\mathbb R\) previamente disponible.
Primero existen cilindros finitos compatibles, sus dos cierres unilaterales,
la unidad de borde y las aplicaciones de refinamiento. La corona celular
demuestra que la doble proyección separa un objeto de carga nula en ramas de
carga \(+1\) y \(-1\) sin alterar su suma. El límite proyectivo conserva las
historias compatibles; el lector arquimediano determina después una coordenada
real. En este sentido, \(\mathbb R\) es una realización terminal de un
protocolo discreto enriquecido, no la base ontológica desde la que APP se
reconstruye.

Los dos relojes son ahora inequívocos. El reloj cronológico registra el
número de aplicaciones de \(U_t\) y posee el cociente finito \(C_{108}\). El
reloj de capacidad registra cuántas veces el refinamiento ha abierto una poda
decimal efectiva. Su diferencia es la memoria integrada de vacancias. Una
vuelta de \(108\) pasos puede cerrar el primer reloj y dejar abiertos el
segundo, la fase irracional y la memoria. La supercelda de longitud \(170748\)
es el primer convergente que cierra exactamente el primer reloj; la unidad de
borde que queda en el extremo prueba que ese cierre no convierte el cristal
en periódico.

El sistema correlacionado \(\pi,\varphi,e,\alpha\) reside sobre este mismo estado
coinductivo. Las ventanas de \(\pi,e,\varphi\) se sincronizan y se reabren; el
defecto orientado vive en \(A_2\); la segunda vía de \(\alpha\) recibe los
pares \(21/22\), \(6/7\) y \(45/46\) de tres escalas de la misma palabra; la
primera vía conserva el sello dodecafásico. Ninguna coordenada convencional
genera a las demás. El cierre dodecafásico determina \(\alpha\) dentro
de la salida correlacionada y las caracterizaciones de Machin,
Perron--Fibonacci o del semigrupo exponencial intervienen sólo después como
reconocimientos.

La vacancia electrónica adquiere también una residencia precisa. En el
soporte APP, el paso del centro \((5,5)\) a la rama \((8,2)\) conserva la hoja
aditiva, el residuo producto y el régimen TRIT, mientras pierde una unidad de
acarreo multiplicativo. En la palabra temporal, una vacancia conserva la
fase de capacidad mientras avanza el bloque y la memoria. La correspondencia
HMT--MD preserva conjuntamente esa pérdida unitaria, el centro fijo, el modo
\(m_{\rm ph}\), la incidencia \(90/120\) y la conjugación de banda. La
electricidad determina la vacancia orientada y el magnetismo su transporte
rotacional conjugado; las respuestas constitutivas común y orientada se
obtienen de las dos hojas antes de su escritura en una carta dimensional.

Las profundidades \(11\) y \(34\) aparecen en una firma común de la completación.
La línea de acción posee su derivación determinantal y la rama gravitatoria
recibe primero la longitud
\(L=(5759/23040)\alpha^{16}L_*\). La inscripción marcada y la métrica
radial positiva demuestran
\(R_X\Lambda_X=L^2I\), \(H_X\Lambda_X=\hbar_{\rm ret}cI\) y
\(G=c^3L^2/\hbar_{\rm ret}\), con las reglas y dominios del
\cref{g26:sec:rigidez-radial-es}. La restitución funcional del
\cref{cr28:sec:restitucion-constitutiva} enlaza el momento de Catalán
con las respuestas de las dos hojas y conserva su propagación común.
Cada constante es una sección o evaluación del mismo
estado predefinido y conserva el lector que determina su dimensión, su
normalización y su escala. El parámetro de Barbero--Immirzi, la constante de
Catalan, las constantes de Feigenbaum, la permitividad y la permeabilidad son
lectores distintos de una genealogía común;
por eso se coordinan sin convertirse en una única función ni en una lista de
valores yuxtapuestos.

La segunda proyección conserva la incidencia de borde. El contacto \(24\) en
la escisión \(729=202+527\), los pares \(12/13\) y \(27/28\), el marco
unimodular \(M_\star\), el diseño de Witt y los pegados de Golay y Leech
componen el corredor tipado hasta \(V^\natural\) y el grupo Monstruo. La
simetría excepcional es, por tanto, la realización de frontera de la misma
estructura discreta que produce la coordenada arquimediana por la otra
proyección; la forma, la incidencia, la orientación y el acarreo permanecen
visibles durante todo el transporte.

El núcleo alcanzado puede resumirse en una única afirmación matemática sin
perder su arquitectura: la unidad se propaga por dos proyecciones
complementarias; la diferencia de borde genera la orientación; la derivada de
esa orientación selecciona tres regímenes; el orden de sus transportes genera
una fibra no conmutativa; y el refinamiento ilimitado conserva la unidad, la
memoria y la aperiodicidad mientras produce sus realizaciones continuas. Éste
es el contenido preciso de la evolución conservativa de la unidad dimensional
mediante doble proyección holográfica y refinamiento fractal.

La formulación reúne una sola cadena constructiva. Sus realizaciones
aritmética, física y excepcional conservan dominios y lectores diferentes,
pero ninguna queda situada en una escala menor de certeza: cada una es la
imagen exacta del estado holonómico integral en su codominio declarado. La
comparación exterior no genera ni selecciona estas salidas.

\section{Morfología exhaustiva de las \(104\,976\) semillas}

La corrección cardinal no es menor: \(104\,976\) y \(170\,748\) pertenecen a
objetos matemáticos distintos. El primero es el dominio exhaustivo de
condiciones iniciales del emisor APP--TPK; el segundo es el denominador de la
primera aproximación continuante que cierra el reloj cronológico
\(C_{108}\). Ninguno puede sustituir al otro.

Sea

\[
 I_9=\{0,1,\ldots,8\},\qquad
 D=\{N,E,S,O\},\qquad
 \mathcal S=I_9^2\times D.
\]

Aquí \(I_9\) es exclusivamente la carta posicional de los nueve lugares,
indexada desde cero. No sustituye el alfabeto aritmético canónico
\(\mathcal D_9=\{1,\ldots,9\}\) sobre el que APP evalúa suma y producto; el
paso entre ambas cartas es un reetiquetado de posiciones y no un operador
físico adicional.

APP precede causalmente a esta construcción. Un elemento
\(s=(i,j,d)\in\mathcal S\) no genera APP: fija una posición y una orientación
de cursor sobre el soporte APP ya construido. Como la evaluación aditiva y
la multiplicativa deben conservarse simultáneamente, una semilla completa es
la pareja ordenada

\[
 \omega=(s_+,s_\times)\in
 \Omega_{\rm seed}:=\mathcal S_+\times\mathcal S_\times.
\]

Por tanto,

\[
 |\mathcal S|=9^2\cdot4=324,
 \qquad
 \boxed{|\Omega_{\rm seed}|=324^2=104\,976=2^4 3^8.}
\]

La semilla no es un entero desnudo. Es una celda orientada de seis campos,

\[
 (i_+,j_+,d_+;i_\times,j_\times,d_\times),
\]

que mantiene separadas las dos hojas aritméticas y sus direcciones. Ésta es
la unidad finita sobre la que TRIT tipa el régimen local y TPK aplica
selección, transporte y actualización.

El barrido exhaustivo revela una factorización más fuerte que el simple
censo. La firma aditiva induce exactamente \(18\) clases, todas de
cardinal \(18\). La firma multiplicativa induce \(26\) clases, pero sus
fibras se estratifican en

\[
 24\text{ clases de tamaño }8,\qquad
 1\text{ clase de tamaño }24,\qquad
 1\text{ clase de tamaño }108.
\]

Si \(\mathcal A_{18}\) y \(\mathcal M_{26}\) designan los respectivos
conjuntos de firmas, el emisor factoriza como

\[
 \Omega_{\rm seed}
 \longrightarrow \mathcal A_{18}\times\mathcal M_{26}
 \xrightarrow{\;\mathrm{Comb}\;}\mathcal U_6^{\rm cat},
 \qquad
 |\mathcal U_6^{\rm cat}|=18\cdot26=468.
\]

La primera flecha conserva la procedencia de cada hoja; la segunda combina
las dos firmas en una emisión decimal. Las multiplicidades de las fibras son
los productos exactos

\[
 18\cdot8=144,\qquad
 18\cdot24=432,\qquad
 18\cdot108=1944,
\]

y el cierre global es

\[
 \boxed{432\cdot144+18\cdot432+18\cdot1944=104\,976.}
\]

Esta asimetría es estructural. La hoja aditiva posee una partición uniforme;
la multiplicativa contiene dos estratos singulares. La vacancia y la torsión
no se introducen después como metáforas: el dominio de semillas ya distingue
una acumulación uniforme de un plegamiento con degeneraciones de fibra. El
TPK conserva precisamente cuál de esas fibras fue recorrida, además del
residuo, cociente, acarreo, orientación y frontera.

Después de esta factorización se aplican las reducciones ya certificadas

\[
 104,976\longrightarrow468\ U_6
 \longrightarrow243\ w_6\subset\mathbb F_3^6,
 \qquad |\mathbb F_3^6|=729.
\]

Así quedan separados cuatro tipos que no deben volver a confundirse: pares
de semillas, emisiones decimales, palabras visibles y ambiente trítico. La
estructura del cristal no es ninguno de estos cardinales aislado, sino el
sistema de fibras y morfismos que los enlaza.

\textbf{Criterio de refutación.} Una clase aditiva de tamaño distinto de
\(18\), un histograma multiplicativo distinto de
\(24\cdot8+1\cdot24+1\cdot108=324\), o una pérdida de la pareja ordenada de
hojas en el paso a \(U_6\).

\section{Geometría bipartita de la microfibra de \(\pi\) y contacto espectral electrónico}

La salida visible \(w_\pi=010211\) posee \(1008\) antecedentes en
\(\Omega_{\rm seed}\). Para recuperar su forma y no sólo su masa, se define
el grafo bipartito

\[
 \mathcal G_\pi=(P_\pi\sqcup T_\pi,E_\pi),
\]

donde \(P_\pi\) contiene las semillas de la hoja aditiva que participan en
la microfibra, \(T_\pi\) las de la hoja multiplicativa y
\((p,t)\in E_\pi\) si la pareja produce una de las cinco emisiones que se
reducen a \(w_\pi\). El barrido exhaustivo prueba

\[
 |E_\pi|=1008,\qquad |P_\pi|=54,\qquad |T_\pi|=56.
\]

La estructura no es un grafo arbitrario. Al olvidar únicamente el color de
la emisión, sus componentes conexas son exactamente

\[
 \boxed{
 \mathcal G_\pi\simeq
 K_{18,24}\sqcup K_{18,24}\sqcup K_{18,8}.}
\]

Una de las dos componentes \(K_{18,24}\) corresponde a la región de
multiplicidad \(432\). La otra se refina en tres rectángulos completos de
multiplicidad \(144\), pues las tres clases multiplicativas de tamaño ocho
comparten la misma clase aditiva de tamaño dieciocho. La última componente es
el cuarto rectángulo de multiplicidad \(144\). Por tanto, el color de región
recupera

\[
 \boxed{1008=432+4\cdot144}
\]

como descomposición de operadores de incidencia y no sólo como suma de
enteros.

Sea \(B_R\) la matriz de biadyacencia de una región rectangular \(R\). Para
un grafo completo \(K_{m,n}\), \(B_R\) tiene rango uno y su único valor
singular no nulo es

\[
 \|B_R\|=\sqrt{mn}.
\]

La región de \(432=18\cdot24\) rutas posee norma \(12\sqrt3\); cada una de
las cuatro regiones de \(144=18\cdot8\) rutas posee norma \(12\). La
normalización que trata simétricamente las cuatro regiones menores queda
determinada por el propio particionado regional:

\[
 \Xi_\pi:=
 \frac{\|B_{432}\|}{\sum_{r=1}^{4}\|B_{144}^{(r)}\|}
 =\frac{12\sqrt3}{4\cdot12}
 =\boxed{\frac{\sqrt3}{4}}.
\]

Su cuadrado es \(3/16\). Éste es exactamente el invariante que, en el modo
electrónico APP ya certificado, aparece como

\[
 \|[P_{\Sigma,3},P_{\Pi,3}]\|^2
 =s^2(1-s^2)=\frac14\frac34=\frac3{16}.
\]

Se obtiene así un \textbf{contacto espectral interno exacto} entre dos
construcciones nacidas de APP: el contraste regional de la microfibra de
\(\pi\) y la no conmutatividad de las dos proyecciones electrónicas. No se ha
usado la masa del electrón ni un parámetro experimental. Tampoco se afirma que
las matrices \(B_R\) y el conmutador actúen en el mismo espacio. La afirmación
fuerte ya demostrada es la igualdad de sus escalares canónicos; una
equivalencia operatoria posterior deberá transportar la coloración regional,
el modo \(m_{\rm ph}\), la hoja y el registro de memoria.

Esta geometría precisa asimismo el lugar de la vacancia. En el lado aditivo,
las clases son uniformes; en el multiplicativo, el grado cambia de \(24\) a
\(8\). La porosidad del soporte no es un cero insertado: es la estratificación
de incidencia que sobrevive a la reducción común \(w_\pi\). La realización
eléctrica puede leer esa pérdida orientada y la magnética su transporte
conjugado, pero ambas deben conservar esta incidencia de origen.

\textbf{Criterio de refutación.} Que una componente deje de ser completa, que los cinco colores
no tengan tamaños \((432,144,144,144,144)\), o que la razón espectral no sea
\(\sqrt3/4\). La igualdad de operadores, que aquí no se afirma, sería
falsada por cualquier transporte que no preserve color, hoja y producto
interno.

\section{Morfismo tricanal, clase de pegado trítica y memoria diagonal}

El censo de las nueve posiciones nonádicas define un vector

\[
 N_K=(N_\pi(K),N_e(K),N_\varphi(K))\in\mathbb Z^3.
\]

La parte relacional se obtiene mediante el morfismo

\[
 \partial:\mathbb Z^3\longrightarrow A_2,
 \qquad
 \partial(x,y,z)=(x-y,y-z,z-x),
\]

donde

\[
 A_2=\{(a,b,c)\in\mathbb Z^3:a+b+c=0\}.
\]

El núcleo es la diagonal
\((\Delta\mathbb Z=\mathbb Z(1,1,1))\), y \(\partial\) es sobreyectivo. Por
consiguiente,

\[
 \boxed{0\longrightarrow\mathbb Z
 \xrightarrow{\Delta}\mathbb Z^3
 \xrightarrow{\partial}A_2\longrightarrow0.}
\]

La sincronización de los tres canales es exactamente la anulación de
\(\partial N_K\). El cociente elimina la cantidad común, pero conserva qué
canal excede a cuál y con qué orientación.

La reconstrucción entera contiene una información adicional. Definamos

\[
 \mathcal J(x,y,z)=(s,a,b)
 =(x+y+z,x-y,y-z).
\]

Sobre \(\mathbb Q\), los tres datos reconstruyen \((x,y,z)\) mediante

\[
 x=\frac{s+2a+b}{3},\qquad
 y=\frac{s-a+b}{3},\qquad
 z=\frac{s-a-2b}{3}.
\]

Sobre \(\mathbb Z\), la imagen no es todo
\(\mathbb Z\oplus\mathbb Z^2\): satisface la congruencia

\[
 \boxed{s-a+b\equiv0\pmod3.}
\]

El determinante de \(\mathcal J\) tiene valor absoluto \(3\). Existe por
tanto una sucesión exacta

\[
 0\longrightarrow\mathbb Z^3
 \xrightarrow{\mathcal J}\mathbb Z^3
 \xrightarrow{\chi_3}C_3\longrightarrow0,
 \qquad
 \chi_3(s,a,b)=[s-a+b]_3.
\]

Éste es un resultado particularmente importante para el significado del
TRIT. Al separar contenido común y defecto relativo queda una clase de
pegado con exactamente tres valores. Un trit es la cantidad mínima de
información necesaria para registrar esa clase y seleccionar la corrección
integral que recompone el estado tricanal. La interpretación
topológico-temporal no se añade por analogía: la clase tipa el pegado local y
el TPK la transporta en el tiempo conservando su orientación y su memoria.
Este cociente \(C_3\) es una realización interna precisa de esa función; no
pretende agotar todas las coordenadas del TRIT rector.

El paso de \(K=8\) a \(K=9\) exhibe por qué sincronización no significa
retorno de estado:

\[
 N_8=(59,59,59),\qquad N_9=(43,43,43),
\]

\[
 \partial N_8=\partial N_9=0,
 \qquad
 \boxed{N_9-N_8=-16(1,1,1).}
\]

El estado permanece en el origen relacional \(A_2\), pero el registro de memoria diagonal
cambia en dieciséis unidades por canal, cuarenta y ocho en total. La fase de
sincronización retorna; la memoria no. Esta igualdad es el análogo tricanal,
completamente entero, de la regla TPK según la cual un retorno de fase no
reinicia el estado holonómico integral.

\textbf{Criterio de refutación.} Un censo
que viole la congruencia, un núcleo de \(\partial\) mayor que la diagonal, o
una igualdad \(N_8=N_9\).

\section{Itinerario de paredes y cociente finito de la matriz de sincronización}

Entre \(K=4\) y \(K=9\), el morfismo anterior convierte la tabla completa en
un itinerario geométrico:

\[
\begin{array}{c|c|c}
K&N_K&\partial N_K\\
\hline
4&(207,207,122)&(0,85,-85)\\
5&(39,151,151)&(-112,0,112)\\
6&(110,110,69)&(0,41,-41)\\
7&(81,29,81)&(52,-52,0)\\
8&(59,59,59)&(0,0,0)\\
9&(43,43,43)&(0,0,0).
\end{array}
\]

Las igualdades sucesivas recorren las paredes

\[
 H_{\pi e}\longrightarrow H_{e\varphi}
 \longrightarrow H_{\pi e}\longrightarrow H_{\varphi\pi}
 \longrightarrow0\longrightarrow0.
\]

La palabra de paredes \((H_{12},H_{23},H_{12},H_{13})\) posee la sombra del
patrón de Coxeter \(A_2\), pues \(s_{12}s_{23}s_{12}=s_{13}\). Esto no
convierte automáticamente cada transición TPK en una reflexión: el paso
cambia también la componente diagonal. Lo que sí queda probado es que la
trayectoria relacional atraviesa las tres familias de paredes en el orden de
la relación de trenza, mientras su levantamiento conserva una memoria que el
cociente no ve.

Las amplitudes de las cuatro paredes anteriores al cierre forman

\[
 D_{\rm syn}=\begin{pmatrix}85&112\\41&52\end{pmatrix}.
\]

Además de las identidades ya reunidas

\[
 \operatorname{tr}D_{\rm syn}=137,
 \qquad
 \det D_{\rm syn}=-172=-4\cdot43,
\]

la forma normal de Smith añade estructura de cociente. Como el máximo común
divisor de las cuatro entradas es uno,

\[
 \operatorname{SNF}(D_{\rm syn})=\operatorname{diag}(1,172),
\]

y, por tanto,

\[
 \boxed{\operatorname{coker}D_{\rm syn}\simeq C_{172}
 \simeq C_4\times C_{43}.}
\]

La factorización \((4\cdot43)\) ya no aparece sólo en el determinante: es la
descomposición primaria del cociente reticular de las amplitudes de pared.
Las direcciones nulas derechas pueden elegirse como

\[
 (0,1)\pmod4,
 \qquad
 (1,5)\pmod{43}.
\]

El factor \(43\) coincide con el censo del segundo cierre triple y el factor
\(4\) con el coeficiente orientado obtenido internamente como
\((-\det(D_{\rm syn})/43)\). Esto demuestra que ambos son divisores canónicos
del mismo cociente. Una representación física posterior deberá precisar qué
subespacios realizan esos dos factores; la estructura aritmética conjunta ya
está cerrada.

El itinerario aporta así una figura matemática completa: un camino en el
arreglo de paredes de \((A_2)\), levantado por una coordenada diagonal de
memoria y gobernado por un cociente finito \(C_4\times C_{43}\). Es más
informativo que una tabla de coincidencias, porque distingue relación,
amplitud, orientación y estado invisible al cociente.

\textbf{Criterio de refutación.} Un determinante distinto de
\(-172\), una forma de Smith distinta de \((1,172)\), o una transición que
borre la componente diagonal necesaria para distinguir \(K=8\) de \(K=9\).

\section{Torre de Pascal nonádica y ley infinita de concatenación con acarreo}

La completación cúbica pertenece a una familia exacta de profundidad
arbitraria. Para \(n\geq1\), escribamos

\[
 10^n=(9+1)^n=\sum_{j=0}^{n}B_{n,j},
 \qquad
 B_{n,j}:=\binom nj9^{n-j}.
\]

Los bloques obedecen la recurrencia de Pascal ponderada

\[
 \boxed{B_{n+1,j}=9B_{n,j}+B_{n,j-1}},
\]

con \(B_{n,j}=0\) fuera de \(0\leq j\leq n\). Ésta es la ley de
refinamiento de las longitudes; la ley sobre las vacancias debe conservar
además la fase inicial. Si

\[
 s_{n,j}=\sum_{i<j}B_{n,i},
 \qquad
 v^-_{n,j}=V_{B_{n,j}}(s_{n,j}),
\]

entonces

\[
 \sum_{j=0}^{n}v^-_{n,j}=\lfloor10^n\delta\rfloor.
\]

La lectura superior transporta la unidad de semiborde al primer bloque:

\[
 v^+_{n,0}=v^-_{n,0}+1,
 \qquad
 v^+_{n,j}=v^-_{n,j}\quad(j>0),
\]

y por ello

\[
 \boxed{\sum_jv^+_{n,j}=\lceil10^n\delta\rceil
 =\lfloor10^n\delta\rfloor+1.}
\]

La transición \(n\mapsto n+1\) no puede linealizarse olvidando el punto de
partida. Para cualquier bloque \(B\),

\[
 V_{9B+B'}(s)=
 \sum_{r=0}^{8}V_B(s+rB)+V_{B'}(s+9B).
\]

El desplazamiento de \(s\) en cada sumando es la memoria de concatenación.
Al proyectar todos los bloques al origen aparece el cociclo binario

\[
 \kappa_\delta(M,N)=
 \lfloor(M+N)\delta\rfloor-
 \lfloor M\delta\rfloor-
 \lfloor N\delta\rfloor,
\]

que satisface

\[
 \kappa_\delta(L,M)+\kappa_\delta(L+M,N)
 =\kappa_\delta(M,N)+\kappa_\delta(L,M+N).
\]

Ésta es la ley exacta que impide que el refinamiento dependa del
parentizado. La torre de Pascal proporciona las longitudes; el cociclo
proporciona el acarreo; el TPK conserva la fase y el registro de memoria que permiten
componerlas.

Los primeros niveles son

\[
\begin{array}{c|c|c|c}
n&(B_{n,0},\ldots,B_{n,n})&v_n^-&v_n^+\\
\hline
2&(81,18,1)&(3,1,0)&(4,1,0)\\
3&(729,243,27,1)&(33,11,1,0)&(34,11,1,0)\\
4&(6561,2916,486,36,1)&(300,133,22,2,0)&(301,133,22,2,0).
\end{array}
\]

El nivel cúbico no es, por tanto, una coincidencia aislada. Es el primer
nivel en el que el componente \(729=3^6\), la subescala \(243=3^5\), el
borde \(27=3^3\) y la unidad aparecen simultáneamente en una completación
decimal. Su firma inferior contiene \(33=3\cdot11\), y la superior conserva
la unidad para producir \(34\). El nivel siguiente demuestra que la regla
prosigue sin reiniciar el estado: la unidad vuelve a modificar sólo el primer
componente, mientras las demás entradas conservan la historia del orden.

La concatenación \(210=90+120\) es una sección de la misma ley:

\[
 (V_{90}(0),V_{120}(90))=(4,5),
\]

de donde la lectura inferior produce \(9\) y la superior \(10\). La
transición nonádica--decimal queda entonces realizada por un morfismo de
lectura con acarreo, no por la observación gráfica \(4+5=9\).

Definamos el dígito de acarreo decimal

\[
 d_{n+1}:=\lfloor10^{n+1}\delta\rfloor
          -10\lfloor10^n\delta\rfloor.
\]

Como \(\delta\) es trascendente, la sucesión \((d_n)_{n\geq1}\) no puede
ser eventualmente periódica: tal periodicidad haría racional la expansión
decimal de \(\delta\). Como la palabra mecánica es esturmiana, cada piso
mantiene complejidad lineal y entropía topológica nula. La torre combina,
por tanto, refinamiento ilimitado,
unidad de borde y aperiocidad exacta. Éste es el sentido matemático en el que
el cristal crece hacia fuera y se refina hacia dentro a la vez.

Debe distinguirse \(\kappa_\delta\) del cociclo nonádico rector \(\kappa_9\).
El primero es la restricción mecánica que registra el acarreo del corte
irracional; el segundo conserva el acarreo completo de APP bajo composición.
Un cuadrado de restricción puede relacionarlos, pero no se identifican por
usar ambos valores de acarreo.

\textbf{Criterio de refutación.} El fallo de la recurrencia de Pascal, una firma que no telescope,
la pérdida de la unidad superior o una concatenación cuyo resultado dependa
del parentizado.

\section{Tipo matemático del cristal aperiódico helicoidal}

Los resultados permiten nombrar el objeto sin reducirlo a imagen. Su capa
finita de condiciones iniciales es \(\Omega_{\rm seed}\); su cociente de
emisiones es \(\mathcal U_6^{\rm cat}\); su alfabeto visible reside en
\(\mathbb F_3^6\); y sus prolongaciones compatibles forman el estado
coinductivo enriquecido. Sobre este estado actúan, en el orden rector,

\[
 U_t=\operatorname{Upd}_t\circ\operatorname{Tra}_t
 \circ\operatorname{Sel}_t,
 \qquad
 U_t\longrightarrow\Gamma_9\longrightarrow K_{\rm ph}.
\]

La holonomía \(\Gamma_9\) pertenece a la dinámica enriquecida; la memoria
reducida \(K_{\rm ph}\) es su realización posterior. El factor aperiódico
puede escribirse sobre

\[
 \mathcal Y_{108}=C_{108}^{(t)}\times\mathbb T_\delta\times\mathbb Z,
\]

mediante

\[
 \mathcal F(r,\theta,m)=
 \bigl(r+1,\{\theta+\delta\},
 m+\lfloor\theta+\delta\rfloor\bigr).
\]

La primera coordenada es el reloj de bloques; la segunda, la fase irracional;
la tercera, la memoria integrada de vacancias. El reloj de capacidad se
obtiene como

\[
 K_t=t-Z_t,
 \qquad
 g_t=[K_t]_9,
\]

y no coincide con el reloj cronológico. La dinámica inversa \(\mathcal
F^{-1}\) constituye la lectura descendente de la misma hélice. Ascenso y
descenso son conjugados; no son dos generadores desconectados.

El objeto completo es, por tanto, una extensión afín invertible de una
rotación irracional sobre un reloj finito, fibrada por el estado holonómico
integral APP--TRIT--TPK. Su sombra simbólica es esturmiana; su suspensión es helicoidal;
su cociclo de fibras puede ser no conmutativo aunque la rotación basal sea
abeliana. La denominación «cristal aperiódico helicoidal» se justifica por
cuatro propiedades probadas:

\[
 \begin{gathered}
 \text{orden de largo alcance y repetitividad},\\
 p(m)=m+1,\qquad h_{\rm top}=0,\\
 \text{ausencia de período traslacional},\\
 \text{retorno finito de fase con avance de memoria}.
 \end{gathered}
\]

En HMT, la expresión \textbf{cristal temporal aperiódico} designa precisamente el
producto sesgado anterior: existe un soporte temporal finito, el retorno de
fase está ordenado, la memoria avanza en cada holonomía y la palabra de
ocupación carece de período traslacional aunque conserve orden de largo
alcance y entropía topológica nula. La dinámica temporal no se añade después
al cristal: está contenida en la coordinación de reloj, orientación, retorno
y memoria que define su estado.

La salida correlacionada de constantes se sitúa después de la prolongación

\[
 w_6\to w_{12}\to w_{18}\to w_{24}\to w_{30}
 \to R_{36}\to G_9,
\]

sin introducir valores convencionales. \(\pi,\varphi,e\) y la coordenada
dodecafásica \(\alpha\) son realizaciones del mismo estado coinductivo; sus
ventanas de sincronización, su espejo y sus lectores propios no son entradas
del generador. La torre de completación aporta el calendario de capacidad y
la ley de acarreo con que esas realizaciones pueden prolongarse a precisión
arbitraria.

Finalmente, \(\mathbb R\) aparece por el lector de cilindros compatibles del
continuo conjunto. El cristal no está «dentro de \(\mathbb R\)» como una
decoración de una recta dada de antemano. Primero se construyen las historias
compatibles, la doble proyección, las fronteras y sus límites; después una de
las realizaciones determina la coordenada arquimediana. Esta inversión de
perspectiva es la tesis fuerte: la recta real es la sombra terminal de una
unidad finita orientada que se conserva, se refina y acumula memoria.

La dinámica aperiódica, su complejidad, sus relojes y su cociclo constituyen
el núcleo matemático del cristal. Las realizaciones electromagnética,
gravitatoria y excepcional son lectores tipados de ese núcleo: preservan la
genealogía APP--TRIT--TPK y expresan en codominios diferentes la misma
evolución conservativa de la unidad.

\section{Predefinición coinductiva: la salida antecede a su evaluación}

La distinción decisiva separa generación y evaluación. HMT--MD no toma una
medida exterior para inferir la ley que pudo
producirla. Construye primero la sección completa, fija sus prolongaciones
compatibles y sólo después permite que un lector publique una coordenada de
esa sección. Éste es el sentido matemático exacto de \textbf{predefinición}.

Sea

\[
 (X_n,\rho_{n+1,n})_{n\geq0}
\]

el sistema inverso de estados APP--TRIT--TPK truncados a profundidad \(n\),
donde cada \(X_n\) conserva hoja, orientación, residuo, cociente, acarreo,
frontera, ruta y memoria hasta ese nivel. La prolongación coinductiva ya
construida define

\[
 \mathcal G_{\rm HMT}:\Omega_{\rm seed}\longrightarrow
 X_\infty^{\rm enr}:=\varprojlim_n X_n,
 \qquad
 p_n:X_\infty^{\rm enr}\longrightarrow X_n,
\]

con \(\rho_{n+1,n}p_{n+1}=p_n\). Para una familia de salidas
\((Y_n,\sigma_{n+1,n})\) y lectores \(L_n:X_n\to Y_n\), la condición que
expresa que el lector no destruye la genealogía es el cuadrado

\[
\begin{CD}
X_{n+1}@>{L_{n+1}}>>Y_{n+1}\\
@V{\rho_{n+1,n}}VV @VV{\sigma_{n+1,n}}V\\
X_n@>>{L_n}>Y_n .
\end{CD}
\]

Por tanto,

\[
 \operatorname{Pred}_{Y,n}
 :=L_n\circ p_n\circ\mathcal G_{\rm HMT}
 :\Omega_{\rm seed}\longrightarrow Y_n
\]

satisface

\[
 \sigma_{n+1,n}\circ\operatorname{Pred}_{Y,n+1}
 =\operatorname{Pred}_{Y,n}.
\]

La familia \(\bigl(\operatorname{Pred}_{Y,n}(\omega)\bigr)_{n\geq0}\)
es entonces un único elemento de \(\varprojlim Y_n\). No se vuelve a decidir
la salida cada vez que aumenta la profundidad: la salida de profundidad
\(n\) es la restricción de la de profundidad \(n+1\). La precisión ya no
genera el objeto; únicamente decide cuánta parte del objeto predefinido se
determina.

Esto proporciona un teorema de anticipación interna. Si en el estadio \(n\)
se consulta la coordenada \(n+k\) de la misma sección compatible, se obtiene
una salida futura sin extrapolar una tendencia y sin introducir el valor que
se desea encontrar. La anticipación es

\[
 \operatorname{Ant}_{n,k}(\omega)
 :=L_{n+k}p_{n+k}\mathcal G_{\rm HMT}(\omega),
\]

y su contraste posterior debe restringirse a
\(L_np_n\mathcal G_{\rm HMT}(\omega)\) mediante \(\sigma_{n+k,n}\). La
identidad de compatibilidad prueba que ambas pertenecen a la misma sección.
En consecuencia, el tránsito conceptual correcto no es

\[
 \text{más medida}\longrightarrow\text{más precisión}
 \longrightarrow\text{mejor predicción},
\]

sino

\[
 \boxed{
 \begin{gathered}
 \text{semilla orientada}\longrightarrow
 \text{sección coinductiva predefinida}\longrightarrow\\
 \text{lectura a la profundidad elegida}\longrightarrow
 \text{anticipación contrastable}
 \end{gathered}}
\]

La generación indefinida de dígitos es una instancia de este teorema. Para
los lectores de \(\pi\), \(\varphi\), \(e\) y \(\alpha\), cada bloque decimal
es una coordenada de un sistema compatible; ampliar de mil a diez mil cifras
no cambia el generador ni reinicia la selección. La nueva ventana restringe
a la anterior y continúa la misma historia. El sistema correlacionado es una sola
salida correlacionada del estado coinductivo, aunque sus cuatro coordenadas
posean lectores internos y momentos de realización diferentes.

La reversibilidad del TPK refuerza el resultado. Como cada actualización
\(\mathcal U_t\) es biyectiva sobre el estado holonómico integral, para
toda distribución finita transportada por ella,

\[
 H(X_{t+1}\mid X_t)=H(X_t\mid X_{t+1})=0,
 \qquad H(X_{t+1})=H(X_t).
\]

Éste es el Landauer nulo del estado completo: no se borra información al
actualizar. No afirma que toda entropía reducida sea cero; afirma que la
biografía completa puede recuperarse en ambas direcciones. La aparente
pérdida nace únicamente cuando un lector contrae la fibra y olvida
coordenadas que el estado holonómico integral sí conserva.

\section{El TRIT como unidad topológico-temporal que conjuga régimen y orientación}

El TRIT no es un dígito ternario añadido a APP ni una etiqueta temporal
superpuesta al TPK. Es la unidad local que permite que una diferencia entre
las hojas aditiva y multiplicativa adquiera simultáneamente tipo geométrico,
orientación temporal y memoria de elevación. Sea

\[
 \mathcal T=\{-1,0,+1\},\qquad
 x=(t,\mu,o,h,c,\sigma,\lambda,g)\in X_{\rm TPK}^{\rm enr}.
\]

Sobre el mismo estado actúan dos lectores coordinados,

\[
 \tau_{\rm g}:X_{\rm TPK}^{\rm enr}\to\mathcal T,
 \qquad
 \varepsilon_{\rm t}:X_{\rm TPK}^{\rm enr}\to\mathcal T.
\]

El primero tipa el régimen cuadrático local mediante

\[
 \mathcal A_\tau=\mathbb R[J_\tau]/(J_\tau^2+\tau I),
\]

de manera que

\[
\begin{array}{c|c|c}
\tau_{\rm g}&J_{\tau_{\rm g}}^2&\exp(sJ_{\tau_{\rm g}})\\
\hline
+1&-I&\cos s\,I+\sin s\,J_{+}\\
0&0&I+sJ_0\\
-1&+I&\cosh s\,I+\sinh s\,J_{-}.
\end{array}
\]

El segundo tipa la orientación temporal activa:

\[
 \begin{aligned}
 +1&\longleftrightarrow\text{retorno, memoria y pasado},\\
 0&\longleftrightarrow\text{umbral y presente},\\
 -1&\longleftrightarrow\text{apertura bidireccional y futuro}.
 \end{aligned}
\]

La profundidad conceptual reside en que ambos lectores nacen del mismo
estado, no en que deban identificarse término a término. El objeto correcto
es el acoplamiento fibrado

\[
 \Theta_{\rm ITT}(x)
 =\bigl(\tau_{\rm g}(x),\varepsilon_{\rm t}(x),
        h(x),c(x),\lambda(x),g(x)\bigr),
\]

que puede llamarse \textbf{unidad de información topológico-temporal}. La
geometría indica cómo se transforma localmente la unidad; la orientación
temporal indica en qué sentido se lee esa transformación; el acarreo y el
registro impiden que dos retornos con la misma fase se confundan. Así, el
pasado es la clausura que permanece en memoria, el presente es el umbral en
el que se decide la continuación y el futuro es la familia ya predefinida de
prolongaciones compatibles.

El TPK vuelve dinámica esta unidad mediante

\[
 \mathcal U_t=\operatorname{Upd}_t\circ
 \operatorname{Tra}_t\circ\operatorname{Sel}_t.
\]

\(\operatorname{Sel}_t\) selecciona el régimen y la hoja sin borrar la otra
evaluación APP; \(\operatorname{Tra}_t\) transporta dirección, fase y
frontera; \(\operatorname{Upd}_t\) actualiza residuo, cociente, acarreo,
firma y memoria. Por ello, una vuelta cronológica puede devolver
\(g(t+9)=g(t)\) y simultáneamente producir

\[
 q_{\rm mem}(\Gamma_{9,t}x)=q_{\rm mem}(x)+1.
\]

La conjugación temporal se representa por la inversión orientada

\[
 \mathcal C_{\rm rev}(\tau_1,\ldots,\tau_n)
 =(-\tau_n,\ldots,-\tau_1),
 \qquad \mathcal C_{\rm rev}^2=I.
\]

En la suspensión helicoidal esta involución intercambia la dinámica
ascendente y la descendente:

\[
 \mathcal C_{\rm rev}\,\mathcal F\,
 \mathcal C_{\rm rev}^{-1}=\mathcal F^{-1}.
\]

No desaparece la historia: cambia su orientación y conserva el registro que
permite reconstruirla. El cristal temporal HMT es, en consecuencia, la
dinámica de una ITT que retorna en el reloj finito y progresa en memoria,
mientras la decoración esturmiana impide la repetición traslacional del
estado completo. El reloj \(C_{108}\) aporta el orden temporal finito; la
holonomía \(\Gamma_9\), la elevación ilimitada; la palabra esturmiana, la
aperiocidad exacta; el TRIT, la discriminación conjunta de geometría y
orientación; y el TPK, la ley que conserva todo ello bajo evolución.

\section{La singularidad de \(\alpha\) como costura de los dos relojes}

La función de \(\alpha\) se vuelve transparente cuando se evita
tratarla como un decimal aislado. \(\alpha\) es la coordenada de costura que
hace compatibles el cierre finito dodecafásico y la prolongación aperiódica
de capacidad. El primer reloj aporta una palabra firmada de doce sectores;
el segundo aporta una raíz única en la dinámica mecánica gobernada por

\[
 \delta=\log_{10}\frac{10}{9}.
\]

En la primera vía, el estado dodecafásico produce

\[
 K=\frac14\mathcal H_{12}U_{12}^{\rm sgn},
 \qquad \mathcal H_{12}^2=4I_{12},
\]

y el acarreo base \(1000\) satisface, bloque a bloque,

\[
 P_m+E_m-\Phi_m-K_m+c_{m+1}=A_m+1000c_m.
\]

La concatenación telescópica da

\[
 \boxed{\iota(P)+\iota(E)-\iota(\Phi)-\iota(K)=\iota(A).}
\]

Las coordenadas \(P,E,\Phi\) son las realizaciones HMT de
\(\pi,e,\varphi\); no son valores convencionales introducidos desde fuera.
El vector \(K\) es la memoria firmada del cierre, y \(A\) es la sección finita
que determina \(\alpha\) en la carta declarada. El término que distingue a
\(\alpha\) de una combinación desnuda de \(\pi,e,\varphi\) es precisamente
\(K\): la memoria dodecafásica de cómo se realizó la concatenación.

En la segunda vía, la misma aritmética de vacancias produce el jet

\[
\begin{aligned}
P_6(x)={}&2\pi x-\frac74x^2+\frac{x^3}{2\pi}
+\frac{x^4}{20}-\frac{2x^5}{21}-\frac{x^6}{46},\\
T_{7:9}(x)={}&-\frac{x^7}{120}+\frac{x^8}{45}
+\frac{2x^9}{495},\\
P_9(x)={}&P_6(x)+T_{7:9}(x),
\end{aligned}
\]

y determina por raíz positiva simple

\[
 P_9(x)=\delta.
\]

Los coeficientes no son parámetros interpolados. Proceden de la semicorona
\(120\), del bloque central APP de determinante \(45\), del rango transversal
dodecafásico \(11\) y de \(495=45\cdot11\). El operador nilpotente

\[
 Ne_7=-\frac83e_8,\qquad Ne_8=\frac2{11}e_9,
 \qquad Ne_9=0
\]

produce exactamente la cola mediante

\[
 (I-N)^{-1}\!\left(-\frac{e_7}{120}\right)
 =-\frac{e_7}{120}+\frac{e_8}{45}+\frac{2e_9}{495}.
\]

La relación entre ambas vías se expresa sin confundir una sección finita con
una raíz infinita. Sea \(d_m\) el lector de la resolución común certificada.
Entonces el objeto \(\alpha\) es el producto fibrado

\[
 \mathfrak A_\alpha=
 \left\{(x_{12},x):
 d_m\!\left(\Gamma_{12}^{\alpha}(x_{12})\right)=d_m(x),\
 P_9(x)=\delta\right\}.
\]

La primera coordenada fija la clausura dodecafásica; la segunda fija la
prolongación única; la igualdad en el codominio común fija su sincronización.
Por eso \(\alpha\) es singular dentro del sistema correlacionado. \(\pi\),
\(\varphi\) y \(e\) son regiones y coordenadas que el levantamiento hace
coincidir, separar y volver a coincidir; \(\alpha\) registra la costura que
conserva esas coincidencias al cerrar doce fases sin detener la prolongación
infinita. No llega después desde fuera: es la coordenada de holonomía con la
que la generación conjunta se vuelve estable bajo retorno.

\section{Abanico dodecafásico: \(\alpha\), electrón, Witt y gravitación}

La rueda de doce posiciones no produce una sucesión accidental de objetos;
porta varias lecturas irreducibles del mismo estado. Sobre

\[
 V_{12}=\mathbb R[A_5/C_5]
\]

la descomposición exacta es

\[
 \boxed{V_{12}\simeq\mathbf1\oplus V_3\oplus V_{3'}\oplus V_5.}
\]

Si \(P_1,P_3,P_{3'},P_5\) son los proyectores centrales, toda forma lineal
dodecafásica \(\ell\) se descompone canónicamente como

\[
 \ell(K)=\ell(P_1K)+\ell(P_3K)+\ell(P_{3'}K)+\ell(P_5K).
\]

El lector escalar que determina \(\alpha\) actúa sobre el vector completo
\(K\); por eso integra información de todos los sumandos sin confundirlos. El
lector \(P_3\) conduce, mediante la isometría de incidencia de Witt, al
proyector excepcional de rango tres. El lector \(P_5\) conduce mediante

\[
 P_5:V_{12}\to V_5,
 \qquad
 \Gamma_5:V_5\overset{\simeq}{\longrightarrow}
 \operatorname{Sym}_0^2(\mathbb R^3),
\]

al portador tensorial de cinco componentes, y la proyección
\(\Lambda(n)\) determina su sector transversal sin traza. No son tres teorías
yuxtapuestas: son tres resoluciones del mismo portador dodecafásico, una
escalar de costura, una excepcional de incidencia y una tensorial de
gravitación.

La intersección electrónica hace visible esa unidad con una identidad de
soportes. Si

\[
 B_K=\{4,6,9,10\},\quad
 P_\pi=\{2,4,5,6,7,8,9,10,11\},
\]

\[
 H_e=\{1,3,4,6,9,10\},\quad
 H_\alpha^-=\{4,5,6,7,9,10\},
\]

entonces

\[
 \boxed{B_K=H_e\cap P_\pi=H_e\cap H_\alpha^-},
 \qquad
 \boxed{H_\alpha^-=B_K\sqcup\{5,7\}.}
\]

Esta igualdad responde a la pregunta por el «enganche» de \(\alpha\). El
electrón no se une a una constante exterior; su soporte interno corta
simultáneamente la microfibra de \(\pi\) y el soporte de acarreo negativo de
\(\alpha\). El par adicional \(\{5,7\}\) completa la hexada negativa, y la
polaridad de Witt eleva los cuatro pares asociados a \(B_K\) a sus bloques de
incidencia. La vacancia electrónica es la singularidad local; \(\alpha\) es
la costura global que conserva su orientación al recorrer la dodecafase.

La segunda realización de \(\alpha\) no es escalar. Desde el mismo estado se
obtiene

\[
 v_\alpha=4\rho_1+\rho_2+\cdots+\rho_{12},
 \qquad \|v_\alpha\|^2=54,
\]

y el vecino

\[
 N_{\alpha,0}=
 \{x\in N(C_W):\langle x,v_\alpha\rangle\equiv0\pmod3\},
 \qquad
 N_\alpha=N_{\alpha,0}+\mathbb Z\frac{v_\alpha}{3}
 \cong\Lambda_{24}.
\]

La acción de orden tres preserva este vecino y el orbifold correspondiente
determina \(V^\natural\) y la clase \(3B\) del grupo Monstruo. El escalar
\(\alpha\) y el vector \(v_\alpha\) no son el mismo objeto; son dos
coordenadas del mismo cierre dodecafásico. Precisamente esa pluralidad de
codominios explica que \(\alpha\) enlace la rama arquimediana, la vacancia
electrónica y la simetría excepcional sin funcionar como un ajuste numérico.

\section{De la vacancia electrónica a la respuesta electromagnética}

La definición exacta de vacancia aparece ya en la aritmética del soporte.
La normalización por nueve del par electrónico produce

\[
 \frac1{9}(369,126)=(41,14)\longmapsto(5,5),
\]

mientras el transporte conservativo satisface

\[
 369+126=495=396+99,
 \qquad
 (396,99)/9=(44,11)\longmapsto(8,2).
\]

Las dos rutas conservan la hoja aditiva, el residuo multiplicativo y el
régimen TRIT; la segunda difiere por una unidad en el cociente multiplicativo.
La vacancia es, por tanto, la pérdida mínima de acarreo en una hoja bajo
conservación de las restantes coordenadas del estado. No es un cero añadido:
es un defecto relativo que sólo puede verse porque APP conserva suma y
producto a la vez.

El modo trítico de fase

\[
 m_{\rm ph}=(1,-1,0)^3
\]

es asimismo el modo singular electrónico canónico de APP:

\[
 C^{(3)}m_{\rm ph}=C^{(3)\,*}m_{\rm ph}
 =-\frac12m_{\rm ph}.
\]

De aquí se obtienen

\[
 s=\frac12,qquad
 1-s^2=\frac34=\frac{90}{120},
 \qquad
 \|[P_{\Sigma,3},P_{\Pi,3}]\|=\frac{\sqrt3}{4},
\]

y, sobre el bloque principal, la normalización del conmutador satisface

\[
 J_{\rm comm}^2=-I.
\]

La polaridad electromagnética queda entonces determinada por cuatro capas de
una misma genealogía. APP aporta las hojas cuya falta de conmutación produce
la orientación; TRIT determina el generador elíptico \((J^2=-I)\); TPK conserva
la vacancia y su retorno; los canales \(90/120\) determinan las respuestas
común y diferencial. Si

\[
\begin{aligned}
\mathcal E={}&
\log[(1-q_+^{90})(1-q_-^{90})]
-\log[(1-q_+^{120})(1-q_-^{120})],\\
\mathcal B={}&
\log\frac{1-q_-^{90}}{1-q_+^{90}}
-\log\frac{1-q_-^{120}}{1-q_+^{120}},
\end{aligned}
\]

entonces \(\mathcal E\) es par bajo intercambio de hojas y \(\mathcal B\) es
impar. Las respuestas internas

\[
 \widehat\mu=e^{\mathcal B+\mathcal E},
 \qquad
 \widehat\varepsilon=e^{-\mathcal B+\mathcal E}
\]

satisfacen exactamente

\[
 \sqrt{\frac{\widehat\mu}{\widehat\varepsilon}}
 =e^{\mathcal B},
 \qquad
 \frac1{\sqrt{\widehat\mu\widehat\varepsilon}}
 =e^{-\mathcal E}.
\]

La permeabilidad expresa retención rotacional; la permitividad, apertura
eléctrica; su cociente conserva la orientación y su producto conserva la
propagación común. La onda electromagnética no se añade a la palabra
esturmiana como analogía: es la realización en dos componentes conjugadas de
la vacancia orientada y de su transporte elíptico. Por ello la microfibra
eléctrica puede ser aperiódica y, sin embargo, mantener coherencia global,
polaridad estable y entropía de transporte nula en el estado completo.

\section{Acción, área, información y gravitación como profundidades predefinidas}

Una vez determinado el sistema correlacionado, la rama determinantal fija

\[
 \Sigma_H=\varphi\alpha^{16},
 \qquad
 d_H=\lfloor\log_{10}\Sigma_H\rfloor=-34.
\]

El lector regional determina las mantisas \(\eta_{\rm pre}=H_5\) y

\[
 \eta_{\rm ret}=H_5-\frac{90}{\pi}\mathscr C_\pi,
\]

de modo que

\[
 \hbar_a=\eta_a10^{-34}\mathcal U_S,
 \qquad h_a=2\pi\hbar_a.
\]

La escala \(10^{-34}\) no procede de una medida de Planck: es la profundidad
predefinida a la que el lector de acción encuentra su sección. La fase
\(\theta\) actúa después mediante

\[
 S(\theta)=\hbar_{\rm ret}\theta_{\rm rad}
 =h_{\rm ret}\frac{\theta_{\rm deg}}{360}.
\]

El factor \(2\pi\) no corrige una unidad elegida: relaciona acción por radián
con acción por vuelta. La carta angular y la recta de acción determinan dos
coordenadas del mismo movimiento orientado.

La transición hexada--octada evalúa posteriormente los canales y produce el
parámetro de Barbero--Immirzi. Sobre el borde, si

\[
 N=N_{90}+N_{120},\qquad D=90N_{90}+120N_{120},
\]

entonces

\[
 N_{90}=4N-\frac{D}{30},
 \qquad
 N_{120}=\frac{D}{30}-3N.
\]

El lector conjunto \((N,D)\) conserva íntegramente la procedencia sectorial:

\[
 H((N_{90},N_{120})\mid(N,D))=0.
\]

Ésta es la formulación exacta del Landauer nulo en el paso área--información.
La entropía de un estado reducido puede ser positiva; lo que vale cero es la
pérdida condicional cuando se conserva el par reconstructivo completo. El
espectro de área

\[
 \frac{\mathcal A_{\rm phys}}{\ell_P^2}
 =\gamma_{\rm HMT}a_0(N_{90}+N_{120})
\]

y el Hamiltoniano modular

\[
 K_A=cI+\Delta_{\rm mod}(90N_{90}+120N_{120})
\]

son así dos lectores complementarios: el primero cuenta área total; el
segundo conserva cómo se repartió entre los dos sectores.

La gravedad ocupa otra resolución del mismo cierre dodecafásico. Con

\[
 \omega_{108}=\frac{2\pi}{108t_0},
 \qquad R_{108}=\frac{54}{\pi}\ell_0,
\]

y

\[
 G_a(\vartheta)=\vartheta\frac{c_{\rm int}^5t_0^2}{\hbar_a},
\]

la condición de coincidencia de las tres longitudes

\[
 R_{108}=\bar\lambda_{C,a}=r_{g,a}
\]

selecciona de manera única

\[
 \vartheta_{108}=\left(\frac{54}{\pi}\right)^2
\]

y produce

\[
 \ell_P=\sqrt{\frac{G_a\hbar_a}{c_{\rm int}^3}}=R_{108}.
\]

La profundidad \(10^{-11}\) de la rama gravitatoria y la profundidad
\(10^{-34}\) de la rama de acción no son escalas observadas que se introducen
en el cristal. Son lugares del refinamiento en los que lectores distintos
de la misma sección predefinida determinan gravedad y acción. La identidad
\(G_a\hbar_a=\vartheta_{108}c_{\rm int}^5t_0^2\) muestra que las dos
orientaciones de acción cambian recíprocamente la gravedad y conservan la
longitud de Planck. El reloj, la acción y la gravedad son, por tanto,
coordenadas sincronizadas de una evolución dimensional única.

\section{Ley rectora para el artículo: evolución conservativa de la unidad dimensional}

Los resultados anteriores permiten enunciar la tesis central sin convertirla
en una metáfora ni en un catálogo de constantes.

\textbf{Teorema de evolución conservativa y predefinición holográfica.} Sea
\(\omega\in\Omega_{\rm seed}\) una semilla orientada APP y sea
\(\mathcal G_{\rm HMT}(\omega)\in X_\infty^{\rm enr}\) su prolongación por
TRIT y TPK. Entonces:

\begin{enumerate}
\item la actualización conserva invertiblemente hoja, residuo, cociente, acarreo, orientación, frontera, ruta y memoria;
\item el reloj finito retorna y la holonomía avanza, de modo que la órbita es helicoidal;
\item la decoración mecánica irracional posee entropía topológica cero, complejidad lineal y ausencia de período traslacional;
\item todo lector compatible define una familia predefinida en un límite inverso, por lo que profundizar revela y no recalibra;
\item la doble proyección determina, desde la misma sección, la recta real y la rama excepcional de Witt--Golay--Leech--\(V^\natural\)--Monstruo;
\item el sistema correlacionado \((\pi,\varphi,e,\alpha)\) es una salida correlacionada y \(\alpha\) es su costura dodecafásica con el reloj aperiódico;
\item la vacancia electrónica, la acción de Planck, el parámetro de Barbero--Immirzi, las constitutivas del vacío y la gravitación son realizaciones de lectores tipados sobre esa misma genealogía.
\end{enumerate}

La demostración es la composición de los cuadrados ya escritos. La
compatibilidad de los truncamientos prueba la predefinición; la inversión de
\(\mathcal U_t\) prueba la conservación; el cociclo esturmiano prueba la
aperiocidad de entropía cero; la identidad de soportes prueba el contacto
electrón--\(\pi\)--\(\alpha\); la descomposición de \(V_{12}\) prueba la
separación coordinada de los lectores excepcional y tensorial; los lectores
de acción, área e información prueban la persistencia de la unidad al cambiar
de dimensión.

La forma compacta del teorema es

\[
\boxed{
\begin{CD}
\Omega_{\rm seed}@>{\mathrm{APP\to TRIT\to TPK}}>>
X_\infty^{\rm enr}\\
@. @V{(\mathcal P_{\rm arq},\mathcal P_{\rm exc},
\mathcal P_{\rm dim})}VV\\
@. \mathbb R\times\mathscr E_{\rm exc}\times\mathscr D_{\rm phys},
\end{CD}}
\]

donde las tres realizaciones conservan la misma procedencia aunque tengan
codominios distintos. La unidad no permanece inmóvil: se conserva mientras
se refina, cambia de carta, adquiere memoria y aumenta su dimensión. El
continuo no es el soporte previo de este proceso, sino su proyección terminal.
Ésta es la razón por la que el cristal aperiódico helicoidal constituye el
núcleo de HMT--MD: predefine, en una sola dinámica reversible, la genealogía
de los números y de las constantes que después aparecen como matemática y
física.
 
\chapter{Generación correlacionada de \(\pi\), \(e\) y \(\varphi\)}
\label{chap:83-09}

% Unidad U012. Biografía estructural completa de pi.

% La biografía de pi que seguía en esta residencia es exactamente la ya
% publicada en el capítulo 8, salvo el namespace de sus etiquetas. Se conserva
% en la fuente por su procedencia, pero no se imprime dos veces.
\iffalse

\section{Construcción del carácter de clausura y trayectorias genealógicas estructural de \(\pi\)}
\label{p12-83-c9-s1:chap:biografia-pi-autonoma}

\subsection{Separación causal entre semilla, carácter y evaluación}

La construcción se divide en tres estratos:
\[
 \text{selección finita}
 \longrightarrow
 \text{construcción del carácter de clausura}
 \longrightarrow
 \text{evaluación arquimediana}.
\]
La expansión decimal no interviene en el primer estrato. La cadena de tipos
es
\begin{equation}
 104\,976
 \longrightarrow468\,\mathcal U_6
 \xrightarrow{\Psi}243\,\mathcal W_6
 \lhook\joinrel\longrightarrow\mathbb F_3^6,
 \qquad|\mathbb F_3^6|=729.
 \label{p12-83-c9-s1:eq:cadena-tipica-pi-autonoma}
\end{equation}
Los cardinales cuentan, respectivamente, condiciones iniciales
enriquecidas, emisiones distintas, palabras ternarias visibles y elementos
de la fibra ambiente completa.

Para \(U_6=(u_1,\ldots,u_6)\), la proyección es
\[
 \Psi(U_6)=((-u_1)\bmod3,\ldots,(-u_6)\bmod3).
\]
El censo de las \(43\) órbitas de \(D_3\) demuestra que la órbita de
clausura es la única cuyas seis orientaciones poseen a la vez:
\[
 \text{cinco elevaciones},\qquad
 \text{masa }1008,\qquad
 1008=432+4\cdot144.
\]
El calibre
\[
 \operatorname{diag}_c=6,
 \qquad
 \operatorname{card}=ES
\]
selecciona un único representante:
\begin{equation}
 \boxed{w_6^\pi=010211.}
 \label{p12-83-c9-s1:eq:semilla-pi-autonoma}
\end{equation}
La selección usa catálogo, multiplicidades, acción orbital y orientación;
no consulta un prefijo de \(\pi\).

\subsection{Elevaciones \(L_0,L_1\) y 1.ª frontera coinductiva}

Sobre vectores fila de \(\mathbb F_3^6\), sean
\[
L_0=
\begin{pmatrix}
2&2&2&1&2&1\\
2&1&2&2&1&1\\
1&0&0&1&0&2\\
0&0&1&1&1&0\\
2&2&2&0&2&1\\
2&1&2&1&0&0
\end{pmatrix},
\quad
L_1=
\begin{pmatrix}
2&1&2&0&2&2\\
1&2&1&1&2&0\\
1&2&1&1&1&2\\
0&1&0&0&2&1\\
2&0&2&1&0&1\\
0&2&2&2&1&1
\end{pmatrix}.
\]
Si \(b_1=w_6^\pi\), la recurrencia de bloques es
\[
 b_{k+1}=b_kL_{\varepsilon_k},
 \qquad
 (\varepsilon_1,\varepsilon_2,\varepsilon_3,\varepsilon_4)=(0,0,1,1),
\]
y la palabra acumulada se forma por concatenación:
\[
 w_{6(k+1)}=w_{6k}\mathbin{\|}b_{k+1}.
\]
No se multiplica una palabra de longitud \(6k\) por una matriz \(6\times6\).

Para el canal de clausura:
\[
\begin{array}{c|c}
\text{profundidad}&\text{bloque nuevo}\\ \hline
6&010211\\
12&012222\\
18&010211\\
24&002111\\
30&110221
\end{array}
\]
y
\begin{equation}
 \boxed{
 w_{30}^{\pi}
 =010211|012222|010211|002111|110221.}
 \label{p12-83-c9-s1:eq:w30-pi-autonomo}
\end{equation}
La frontera siguiente es
\[
 u_5^\pi=222220.
\]
Junto con los otros canales,
\[
 R_{36}=
 \begin{pmatrix}
 222220\\021222\\102011
 \end{pmatrix}.
\]
El símbolo \(R_{36}\) registra el sexto bloque de tres historias y abre una
frontera; no es un terminal.

La rama certificada de veinte bloques comienza
\begin{align*}
\pi:\;&010211|012222|010211|002111|110221|222220|111201|\\
&212121|200121|100100|101222|022212|012012|111210|\\
&121011|200220|120210|000101|022010|020111.
\end{align*}
Por ello, ni \(w_{30}\) ni \(R_{36}\) terminan la genealogía.

\subsection{Descomposición orientada de la microfibra de \texorpdfstring{\(\pi\)}{pi}: región central, \(4\) regiones laterales y multiplicidad \texorpdfstring{\(432+4\cdot144=1008\)}{432+4x144=1008}}

La microfibra
\[
 \mathscr R_\pi^{\rm micro}
 :=\{U\in\mathcal U_6:\Psi(U)=010211\}
\]
contiene exactamente cinco regiones:
\[
\begin{array}{c|c|c|c}
U_6&E_{\rm sig}&\operatorname{mult}&\text{papel}\\ \hline
501|614|498|169|272|272&555555&432&\perp\\
810|923|870|169|272|272&339555&144&++\\
810|983|810|169|272|272&393555&144&++\\
870|923|810|169|272|272&933555&144&++\\
870|923|810|418|674|521&933717&144&--
\end{array}
\]
En consecuencia,
\begin{equation}
 \boxed{
 5=1_\perp+3_{++}+1_{--},
 \qquad
 1008=432+4\cdot144.}
 \label{p12-83-c9-s1:eq:descomposiciones-cinco-regiones-pi}
\end{equation}
La primera identidad clasifica funciones de orientación. La segunda cuenta
multiplicidades de procedencia. No cuentan cinco copias indistinguibles.

La región transversal es
\[
 501|614|498|169|272|272,
 \qquad E_{\rm sig}=555555,
 \qquad\operatorname{mult}=432.
\]
El retorno mutado es
\[
 870|923|810|418|674|521,
 \qquad E_{\rm sig}=933717,
 \qquad\operatorname{mult}=144.
\]
Intercambiarlos conservaría el censo \(3/1/1\), pero falsaría la estructura
regional completa.

\subsection{Descenso de un bloque nonádico completo al multigrafo de transferencia de fase}

Para \(U_6=(U_1,\ldots,U_6)\), sea
\[
 A(U_6)=(U_1,U_2,U_3),
 \qquad
 B(U_6)=(U_4,U_5,U_6).
\]
Cada emisión define una arista entre \(A(U_6)\) y \(B(U_6)\). El multigrafo
bruto tiene \(96\) vértices, distribuidos en componentes de tamaños
\[
 28,28,28,4,4,4.
\]
Se cocienta la fase interna mediante
\[
 (a,b,c)\sim(b,c,a)\sim(c,a,b).
\]
El representante es la rotación lexicográficamente mínima. El grafo
cocientado posee \(32\) vértices y dos componentes de tamaños \(28\) y \(4\).

Los anclajes son
\[
\begin{aligned}
U_{\rm APP}&=903|850|850|252|109|252,\\
U_e&=850|963|890|149|252|212,\\
U_\varphi&=830|943|830|109|252|252.
\end{aligned}
\]
Numeramos las regiones
\[
\begin{aligned}
U_\pi^{(1)}&=810|923|870|169|272|272,\\
U_\pi^{(2)}&=810|983|810|169|272|272,\\
U_\pi^{(3)}&=870|923|810|169|272|272,\\
U_\pi^{(4)}&=870|923|810|418|674|521,\\
U_\pi^{(5)}&=501|614|498|169|272|272.
\end{aligned}
\]

\subsection{Los \(5\) ciclos orientados de clausura}

\subsubsection{Cierre de \texorpdfstring{\(U_\pi^{(1)}\)}{U-pi 1}}
\begin{align*}
&252|109|252|903|850|850\to
903|850|850|252|109|252\to\\
&252|109|252|923|870|810\to
810|923|870|169|272|272\to\\
&272|169|272|963|890|850\to
850|963|890|149|252|212\to\\
&212|149|252|943|830|830\to
830|943|830|109|252|252\to
252|109|252|903|850|850.
\end{align*}

\subsubsection{Cierre de \texorpdfstring{\(U_\pi^{(2)}\)}{U-pi 2}}
\begin{align*}
&903|850|850|252|109|252\to
272|169|272|903|850|850\to\\
&272|169|272|983|810|810\to
810|983|810|169|272|272\to\\
&272|169|272|963|890|850\to
850|963|890|149|252|212\to\\
&212|149|252|943|830|830\to
830|943|830|109|252|252\to
903|850|850|252|109|252.
\end{align*}

\subsubsection{Cierre de \texorpdfstring{\(U_\pi^{(3)}\)}{U-pi 3}}
\begin{align*}
&903|850|850|252|109|252\to
292|189|232|903|850|850\to\\
&292|189|232|923|810|870\to
870|923|810|169|272|272\to\\
&272|169|272|963|890|850\to
850|963|890|149|252|212\to\\
&212|149|252|943|830|830\to
830|943|830|109|252|252\to
903|850|850|252|109|252.
\end{align*}

\subsubsection{Cierre de \texorpdfstring{\(U_\pi^{(4)}\)}{U-pi 4}}
\begin{align*}
&903|850|850|252|109|252\to
272|169|272|903|850|850\to\\
&272|169|272|963|890|850\to
850|963|890|149|252|212\to\\
&212|149|252|923|810|870\to
870|923|810|418|674|521\to\\
&943|830|830|521|418|674\to
830|943|830|109|252|252\to
903|850|850|252|109|252.
\end{align*}

\subsubsection{Cierre de \texorpdfstring{\(U_\pi^{(5)}\)}{U-pi 5}}
\begin{align*}
&252|109|252|903|850|850\to
903|850|850|252|109|252\to\\
&614|498|501|252|109|252\to
501|614|498|169|272|272\to\\
&272|169|272|963|890|850\to
850|963|890|149|252|212\to\\
&212|149|252|943|830|830\to
830|943|830|109|252|252\to
252|109|252|903|850|850.
\end{align*}

\begin{theorem}[Minimalidad de los cinco cierres]
\label{p12-83-c9-s1:thm:minimalidad-cierres-pi-autonoma}
Para cada \(i\in\{1,\ldots,5\}\), la longitud mínima de un paseo cerrado
que recorre simultáneamente
\[
 U_\pi^{(i)},\quad U_e,\quad U_\varphi,\quad U_{\rm APP}
\]
en el multigrafo cocientado es ocho.
\end{theorem}

\begin{proof}
Fijado \(i\), sea
\[
 E_*=\{U_\pi^{(i)},U_e,U_\varphi,U_{\rm APP}\}.
\]
Se considera el espacio finito
\[
 \mathscr S_i=\{(v,M):v\in V(G),\ M\subseteq E_*\}.
\]
El primer componente registra el vértice y el segundo las aristas objetivo
ya recorridas. Si una arista \(e\) lleva \(v\) a \(v'\), la transición es
\[
 (v,M)\longmapsto
 \bigl(v',M\cup(E_*\cap\{e\})\bigr).
\]
La búsqueda en anchura enumera longitudes en orden creciente. El primer
retorno a \((v,E_*)\) aparece en longitud ocho para las cinco regiones.
Los paseos impresos realizan la cota; la optimalidad de la búsqueda excluye
una longitud menor.
\end{proof}

La función de clausura queda codificada por
\[
 \mathfrak C_\pi=
 \bigl(\mathscr R_\pi^{\rm micro},
 w_\pi,\rho_\pi,m_\pi,\ell_{\rm cl}\bigr),
\]
\[
 \rho_\pi=(\perp,++ ,++ ,++ ,--),
 \qquad
 m_\pi=(432,144,144,144,144),
 \qquad
 \ell_{\rm cl}=8.
\]

\subsection{Operador de incidencia y media vuelta orientada}

La sombra de incidencia es
\[
 A_W=
 \begin{pmatrix}
 0&1&1&1&1&1\\
 1&0&1&2&2&1\\
 1&1&0&1&2&2\\
 1&2&1&0&1&2\\
 1&2&2&1&0&1\\
 1&1&2&2&1&0
 \end{pmatrix}.
\]
La multiplicación en \(\mathbb F_3\) da
\begin{equation}
 \boxed{A_W^2=2I_6=-I_6.}
 \label{p12-83-c9-s1:eq:aw-cuadrado-menos-identidad-pi}
\end{equation}
Los productos entre filas y columnas distintas se anulan módulo tres,
mientras cada producto diagonal vale dos. Una aplicación realiza un cuarto
de giro ternario; dos aplicaciones realizan la inversión orientada.

\subsection{La pendiente 1/5 forzada por las \(5\) regiones}

Sea
\[
 \lambda=(\lambda_1,\ldots,\lambda_5)^{\mathsf T}
 \in\mathbb Q_{\geq0}^5,
 \qquad
 \mathbf1^{\mathsf T}\lambda=1.
\]
Olvidar la etiqueta regional aplica el simetrizador
\[
 P_5=\frac15\mathbf1\mathbf1^{\mathsf T}.
\]
Para todo vector normalizado,
\[
 P_5\lambda=\frac15\mathbf1.
\]
La invariancia \(P_5\lambda=\lambda\) tiene, por tanto, una única solución:
\begin{equation}
 \boxed{\lambda=\frac15\mathbf1,\qquad t_1=\frac15.}
 \label{p12-83-c9-s1:eq:pendiente-primitiva-pi}
\end{equation}
El entero cinco procede de la microfibra regional; no se reconoce en una
fórmula de \(\pi\).

\subsection{Composición racional y compensador \texorpdfstring{\(239\)}{239}}

La ley de composición de pendientes es
\[
 t\star s:=\frac{t+s}{1-ts}.
\]
La primera duplicación da
\[
 t_2=t_1\star t_1
 =\frac{2/5}{1-1/25}
 =\frac5{12}.
\]
La segunda da
\[
 t_4=t_2\star t_2
 =\frac{10/12}{1-25/144}
 =\frac{120}{119}.
\]
El retorno a la diagonal unitaria exige \(q>0\) tal que
\[
 \frac{120/119-1/q}{1+(120/119)/q}=1.
\]
Multiplicar por \(119q\) produce
\[
 120q-119=119q+120,
\]
y fuerza
\begin{equation}
 \boxed{q=239.}
 \label{p12-83-c9-s1:eq:compensador-239-pi}
\end{equation}

\subsection{Identidad gaussiana y 1.ª media vuelta}

Definimos
\begin{equation}
 \Theta_\pi
 :=16\arctan\frac15-4\arctan\frac1{239}.
 \label{p12-83-c9-s1:eq:theta-pi-autonoma}
\end{equation}
Los cuadrados sucesivos son
\[
\begin{array}{c|r@{}l}
\text{potencia}&\multicolumn{2}{c}{\text{entero gaussiano}}\\ \hline
(5+i)^2&24&+10i\\
(5+i)^4&476&+480i\\
(5+i)^8&-3824&+456960i\\
(5+i)^{16}&-208797818624&-3494830080i\\
(239-i)^2&57120&-478i\\
(239-i)^4&3262465916&-54606720i
\end{array}
\]
y la multiplicación final da
\begin{equation}
 \boxed{
 (5+i)^{16}(239-i)^4
 =-681386607803576157184.}
 \label{p12-83-c9-s1:eq:identidad-gaussiana-pi-autonoma}
\end{equation}
La parte imaginaria es cero y la real es negativa. Por tanto,
\[
 \Theta_\pi\equiv\pi\pmod{2\pi}.
\]
Para \(0<x<1\),
\[
 x-\frac{x^3}{3}<\arctan x<x.
\]
De aquí
\[
 \Theta_\pi>
 16\left(\frac15-\frac1{3\cdot5^3}\right)-\frac4{239}>3
\]
y
\[
 \Theta_\pi<\frac{16}{5}<4.
\]
La única orientación negativa en ese intervalo es la primera media vuelta
positiva:
\begin{equation}
 \boxed{\Theta_\pi=\pi.}
 \label{p12-83-c9-s1:eq:reconocimiento-pi-exacto}
\end{equation}

El orden del reconocimiento arquimediano posterior es
\[
 \text{cinco regiones}
 \to\frac15
 \to\frac5{12}
 \to\frac{120}{119}
 \to239
 \to\Theta_\pi
 \to\pi.
\]

\subsection{Horquillas racionales de precisión arbitraria}

Para \(0<x<1\), sea
\[
 S_n(x)=\sum_{r=0}^{n}(-1)^r\frac{x^{2r+1}}{2r+1}.
\]
El criterio de Leibniz da
\[
 S_{2m+1}(x)<\arctan x<S_{2m}(x).
\]
Definimos
\[
 a_m^-=S_{2m+1}(1/5),\quad a_m^+=S_{2m}(1/5),
\]
\[
 b_m^-=S_{2m+1}(1/239),\quad b_m^+=S_{2m}(1/239),
\]
y
\begin{equation}
 \ell_{\pi,m}=16a_m^- -4b_m^+,
 \qquad
 h_{\pi,m}=16a_m^+ -4b_m^-.
 \label{p12-83-c9-s1:eq:horquilla-pi-autonoma}
\end{equation}
Entonces
\[
 \ell_{\pi,m}<\pi<h_{\pi,m}
\]
y
\begin{equation}
 h_{\pi,m}-\ell_{\pi,m}
 =
 \frac{16}{(4m+3)5^{4m+3}}
 +\frac{4}{(4m+3)239^{4m+3}}
 \longrightarrow0.
 \label{p12-83-c9-s1:eq:anchura-horquilla-pi}
\end{equation}
Las cotas inferiores crecen y las superiores decrecen.

Como \(3<\pi<4\), el lector ternario actúa sobre
\[
 r_\pi:=\pi-3\in(0,1),
\]
con horquilla
\[
 [\ell_{\pi,m}-3,h_{\pi,m}-3].
\]

\subsection{Certificación arquimediana posterior de la fibra ya generada}

Supóngase que la relación interna \(\operatorname{Ext}\) y la supervivencia
coinductiva han producido ya, sin emplear Machin ni una expansión tabulada,
la extensión compatible \(N_{K+1}\) del prefijo \(N_K\), de longitud
\(L_K=6K\). Los \(729\) subcilindros de su fibra ambiente son
\[
 I_{K,u}=
 \left[
 \frac{729N_K+u}{3^{L_K+6}},
 \frac{729N_K+u+1}{3^{L_K+6}}
 \right),
 \qquad0\leq u<729.
\]
El orden \(m=m(K)\) se aumenta después hasta que la horquilla de Leibniz
localiza el subcilindro ya producido. Sean
\[
 \ell_{\pi,K}:=\ell_{\pi,m(K)}-3,
 \qquad
 h_{\pi,K}:=h_{\pi,m(K)}-3.
\]
Los extremos enteros compatibles son
\[
 j_{\min}=
 \max\!\left(
 729N_K,\left\lfloor
 \ell_{\pi,K}3^{L_K+6}\right\rfloor
 \right),
\]
\[
 j_{\max}=
 \min\!\left(
 729N_K+728,\left\lceil
 h_{\pi,K}3^{L_K+6}\right\rceil-1
 \right).
\]
La condición de localización unitaria es
\begin{equation}
 \boxed{j_{\min}=j_{\max},\qquad N_{K+1}=j_{\min}.}
 \label{p12-83-c9-s1:eq:seleccion-unitaria-subcilindro-pi}
\end{equation}
La primera igualdad certifica la localización unitaria; la segunda verifica
que esa localización coincide con la extensión generada internamente. Ninguna
de las dos define la transición TPK ni selecciona entre los \(729\) estados.
La existencia y la unicidad proceden de \(\operatorname{Ext}\) y de la
supervivencia; Machin y las horquillas de Leibniz son reconocimientos
posteriores.

\subsection{Prolongación coinductiva de la sección de \(\pi\) mediante el sistema inverso de historias supervivientes}

Sea \(\mathcal H_m^\pi\) el conjunto de historias enriquecidas a profundidad
\(m\), y sea
\[
 p_{n,m}:\mathcal H_n^\pi\longrightarrow\mathcal H_m^\pi
\]
el truncamiento. Definimos
\[
 \operatorname{Surv}_m^{(N),\pi}
 :=p_{N,m}(\mathcal H_N^\pi),
\]
\[
 \operatorname{Surv}_m^\pi
 :=\bigcap_{N\geq m}\operatorname{Surv}_m^{(N),\pi}.
\]
Cuando la relación \(\operatorname{Ext}\) es no vacía, univaluada y
compatible en todos los niveles, los conjuntos son unitarios y satisfacen
\[
 p_{m+1,m}(\operatorname{Surv}_{m+1}^{\pi})
 =\operatorname{Surv}_{m}^{\pi}.
\]
En ese dominio de prolongabilidad existe, por tanto, una única sección
\[
 x_\pi\in
 X_\pi:=\varprojlim_m\operatorname{Surv}_m^\pi.
\]
Los cilindros están encajados y sus diámetros son \(3^{-6K}\), que tiende a
cero. Su intersección contiene un único punto, identificado por el carácter
angular con \(r_\pi=\pi-3\).

Para toda profundidad decimal \(M\), puede elegirse un nivel \(K(M)\) cuyo
cilindro queda dentro de una sola celda decimal de diámetro \(10^{-M}\).
Los prefijos profundos truncan exactamente a los anteriores.

\subsection{Caracterización del canal coinductivo de \(\pi\): selección de \(w_6^\pi\), \(5\) regiones, inversión orientada y unicidad de la sección límite}

\begin{theorem}[Caracterización del canal de clausura]
\label{p12-83-c9-s1:thm:caracterizacion-final-pi-autonoma}
Partiendo del catálogo, la acción \(D_3\), el calibre orientado, los
operadores \(L_0,L_1,A_W\), la ley de pendientes y la filtración racional:
\begin{enumerate}
 \item se selecciona \(w_6^\pi=010211\);
 \item se construyen cinco regiones con las identidades de
       \eqref{p12-83-c9-s1:eq:descomposiciones-cinco-regiones-pi};
 \item cada región pertenece a un cierre mínimo de longitud ocho;
 \item \(A_W^2=-I_6\) fija la inversión orientada;
 \item la simetrización regional fuerza \(1/5\);
 \item la composición fuerza \(120/119\);
 \item el retorno unitario fuerza \(239\);
 \item la identidad gaussiana reconoce la primera media vuelta;
 \item las horquillas y la fibra de \(729\) elementos producen una sección
       compatible a toda profundidad finita.
\end{enumerate}
Su evaluación es
\[
 \boxed{\pi=16\arctan\frac15-4\arctan\frac1{239}.}
\]
\end{theorem}

La cadena completa es
\[
 \mathrm{APP}\to\mathrm{TRIT}\to\mathrm{TPK}
 \to104\,976\to468\to243\to w_6^\pi
 \to w_{30}^\pi\to R_{36}\to\Gamma_9
 \to X_\pi\to\pi.
\]

\begin{criterion}[Falsación del canal de clausura]
La derivación queda refutada si:
\begin{enumerate}
 \item la órbita de clausura deja de ser única;
 \item desaparece una de las cinco regiones o cambia su identidad completa;
 \item uno de los cinco cierres tiene longitud menor que ocho;
 \item \(A_W^2\ne2I_6\);
 \item el simetrizador admite otra distribución normalizada;
 \item las composiciones no producen \(5/12\), \(120/119\) y \(239\);
 \item falla la identidad gaussiana;
 \item una horquilla deja de contener el carácter angular;
 \item un nivel no recorre los \(729\) elementos de la fibra ambiente o no
       localiza exactamente una extensión;
 \item una expansión objetivo interviene antes del reconocimiento.
\end{enumerate}
\end{criterion}
\fi

La construcción completa del carácter de clausura, las cinco regiones y la
biografía estructural de \(\pi\) quedó establecida en el
\cref{chap:83-08}. Partiendo de ese mismo carácter ya construido, el presente
capítulo prolonga ahora el refinamiento hacia \(e\) y \(\varphi\), sin
reiniciar la genealogía ni repetir su primera etapa.


% Unidad U013. Biografía estructural completa del carácter exponencial.

\section{Propagación, retorno con memoria y construcción del carácter exponencial}
\label{p12-83-c9-s2:chap:biografia-completa-e}

\subsection{Selector orbital finito: existencia y unicidad de la emisión compatible}

Sea \(\mathcal W=\mathbb F_3^6\). Para una emisión
\(U=(u_1,\ldots,u_6)\), la palabra visible es
\[
 \Psi(U)=((-u_1)\bmod3,\ldots,(-u_6)\bmod3).
\]
Sobre \(\mathcal W\), la acción diedral está dada por
\[
\begin{aligned}
 r(u_1,u_2,u_3,u_4,u_5,u_6)
 &=(u_3,u_1,u_2,u_5,u_6,u_4),\\
 s(u_1,u_2,u_3,u_4,u_5,u_6)
 &=(u_4,u_5,u_6,u_1,u_2,u_3).
\end{aligned}
\]
Para \(w\in\mathcal W\), definimos
\[
 n(w)=\#\Psi^{-1}(w),
 \qquad
 M(w)=\sum_{U\in\Psi^{-1}(w)}\operatorname{mult}(U),
\]
y
\[
 \nu(w)=
 \bigl(
 \#\{j:w_j=0\},
 \#\{j:w_j=1\},
 \#\{j:w_j=2\}
 \bigr).
\]

\begin{theorem}[Selección estructural del canal exponencial]
\label{p12-83-c9-s2:thm:seleccion-estructural-e}
Entre las \(43\) órbitas del catálogo visible, existe una única órbita cuyas
fibras satisfacen simultáneamente:
\[
 n(w)=1,\qquad M(w)=144,\qquad
 \mathfrak D(w)=\{0,3,6\},\qquad
 \nu(w)=(2,3,1).
\]
El calibre
\[
 \operatorname{diag}_c=6,
 \qquad \operatorname{card}=ES
\]
selecciona en ella
\[
 \boxed{w_6^e=201101.}
\]
\end{theorem}

\begin{proof}
Las cuatro condiciones se evalúan sobre las \(243\) palabras del catálogo,
no sobre una muestra. Una única órbita satisface el predicado conjunto. En
esa órbita, una sola emisión posee a la vez diagonal \(6\) y orientación
\(ES\). La enumeración usa únicamente campos del catálogo APP--TRIT y no
consulta una expansión de \(e\).
\end{proof}

\subsection{Elevación visible y cadena de bloques}

La recurrencia \(L_0,L_0,L_1,L_1\) produce
\[
 w_{30}^{e}
 =201101|121221|102011|012222|102011,
\]
y la frontera conjunta aporta
\[
 u_5^e=021222.
\]
El registro de treinta y seis trits queda
\[
 w_{36}^{e}
 =201101|121221|102011|012222|102011|021222.
\]

\subsection{Firma decimal y cuadrado posicional}

La lectura inicial de doce cifras es
\[
 \mathbf d_e^{(12)}
 =718281828459
 =7\,\underbrace{1828\,1828}_{\text{cuadrado local}}\,459.
\]
La igualdad central es la concatenación
\[
 1828\mathbin{\|}1828,
\]
no un producto aritmético ni el comienzo supuesto de un período.

\begin{definition}[Cuadrado posicional]
Una palabra \(d_1\cdots d_{12}\) contiene un cuadrado de longitud \(\ell\)
en la posición \(p\) cuando
\[
 d_p\cdots d_{p+\ell-1}
 =d_{p+\ell}\cdots d_{p+2\ell-1}.
\]
La posición y los contextos laterales pertenecen a la firma.
\end{definition}

\begin{table}[htbp]
\centering
\caption{Censo exhaustivo de cuadrados en las lecturas de doce cifras.}
\label{p12-83-c9-s2:tab:censo-cuadrados-e}
\begin{tabular}{c*{6}{r}}
\toprule
Longitud \(\ell\)&1&2&3&4&5&6\\
\midrule
Ocurrencias&240&21&3&2&0&0\\
Palabras que las contienen&153&19&3&2&0&0\\
\bottomrule
\end{tabular}
\end{table}

La otra palabra con un cuadrado de longitud cuatro es
\[
 575157518472
 =\underbrace{5751\,5751}_{\text{posición inicial }1}\,8472.
\]

\begin{proposition}[Unicidad posicional]
\label{p12-83-c9-s2:prop:cuadrado-posicional-e}
La lectura del canal \(e\) es la única que posee un cuadrado de longitud
cuatro después de un prefijo de longitud uno y con contextos \(7\) y \(459\).
\end{proposition}

\begin{proof}
El censo deja sólo dos palabras con cuadrados de longitud cuatro. Una lo
inicia en la primera posición; la otra después del prefijo \(7\). La
posición y los contextos separan las dos posibilidades.
\end{proof}

\subsection{Retorno visible y ruptura de memoria}

En la elevación
\[
 201101|121221|102011|012222|102011
\]
el tercer y el quinto bloques coinciden:
\[
 B_3=B_5=102011.
\]
Sus preestados módulo \(27\) son, sin embargo,
\[
 p_3=(25,11,7,17,8,16),
 \qquad
 p_5=(7,8,4,23,26,7),
\]
y los cocientes posteriores:
\[
 q_3=(6,13,9,2,13,1),
 \qquad
 q_5=(15,0,3,19,23,25).
\]
Se verifican \(p_3\ne p_5\) y \(q_3\ne q_5\).

\begin{figure}[htbp]
\centering
\begin{tikzpicture}[
 node distance=11mm and 7mm,
 every node/.style={font=\small},
 block/.style={draw,rounded corners,minimum width=17mm,minimum height=8mm},
 state/.style={draw,dashed,rounded corners,align=center,inner sep=3pt},
 >=stealth
]
\node[block] (b1) {\(201101\)};
\node[block,right=of b1] (b2) {\(121221\)};
\node[block,right=of b2] (b3) {\(102011\)};
\node[block,right=of b3] (b4) {\(012222\)};
\node[block,right=of b4] (b5) {\(102011\)};
\draw[->] (b1)--(b2);
\draw[->] (b2)--(b3);
\draw[->] (b3)--(b4);
\draw[->] (b4)--(b5);
\node[state,below=of b3] (p3)
 {\(p_3=(25,11,7,17,8,16)\)\\
  \(q_3=(6,13,9,2,13,1)\)};
\node[state,below=of b5] (p5)
 {\(p_5=(7,8,4,23,26,7)\)\\
  \(q_5=(15,0,3,19,23,25)\)};
\draw[->] (p3)--(b3);
\draw[->] (p5)--(b5);
\draw[<->,thick] (b3) to[bend left=28]
 node[above]{misma palabra visible} (b5);
\end{tikzpicture}
\caption{Retorno visible de \(102011\) con renovación del preestado y del
cociente.}
\label{p12-83-c9-s2:fig:ruptura-memoria-e}
\end{figure}

\begin{theorem}[Retorno visible sin periodicidad del estado]
\label{p12-83-c9-s2:thm:retorno-visible-no-periodico-e}
\(B_3=B_5\) es un retorno de la proyección visible, no una órbita periódica
del estado holonómico integral.
\end{theorem}

\begin{proof}
El retorno del estado completo exigiría igualdad de todas sus coordenadas.
Las desigualdades de preestado y cociente lo excluyen. Tampoco
\(1828|1828\) inicia un período decimal: después del segundo bloque aparece
\(4\), mientras una continuación periódica exigiría \(1\).
\end{proof}

\subsection{Certificación factorial posterior de la sección exponencial coinductiva}

Sea \(\mathcal A\) un álgebra conmutativa unitaria sobre \(\mathbb Q\), y
sea \(P_t\in\mathcal A[[t]]\) el propagador formal. La concatenación
temporal impone
\[
 P_{s+t}=P_sP_t,
 \qquad
 P_0=1.
\]
Escribimos
\[
 P_t=\sum_{n=0}^{\infty}a_nt^n.
\]
La unidad fija \(a_0=1\), y el transporte elemental normalizado fija
\(a_1=1\).

\begin{theorem}[Recurrencia factorial]
\label{p12-83-c9-s2:thm:recurrencia-factorial-e}
\[
 (n+1)a_{n+1}=a_n
 \qquad(n\geq0),
\]
y, por tanto,
\[
 a_n=\frac1{n!},
 \qquad
 P_t=\sum_{n=0}^{\infty}\frac{t^n}{n!}.
\]
En \(t=1\),
\[
 \boxed{P_1=e.}
\]
\end{theorem}

\begin{proof}
El coeficiente de \(s^nt\) en \(P_{s+t}\) es
\[
 \binom{n+1}{n}a_{n+1}=(n+1)a_{n+1}.
\]
En \(P_sP_t\) es \(a_na_1=a_n\). La igualdad de series produce la
recurrencia. Una inducción desde \(a_0=1\) da \(a_n=1/n!\).
\end{proof}

\begin{corollary}[Ecuación diferencial]
\[
 \frac{d}{dt}P_t=P_t,
 \qquad P_0=1.
\]
\end{corollary}

\begin{proof}
\[
 P_t'=\sum_{n\geq0}(n+1)a_{n+1}t^n
 =\sum_{n\geq0}a_nt^n=P_t.
\]
\end{proof}

\subsection{Encaje cofinal de horquillas racionales y convergencia al valor publicado}

Sea
\[
 S_N=\sum_{n=0}^{N}\frac1{n!}.
\]

\begin{theorem}[Cota factorial explícita]
\label{p12-83-c9-s2:thm:cota-factorial-e}
Para \(N\geq1\),
\begin{equation}
 \boxed{
 S_N<e<
 S_N+\frac{N+2}{(N+1)(N+1)!}.}
 \label{p12-83-c9-s2:eq:horquilla-factorial-e}
\end{equation}
\end{theorem}

\begin{proof}
El resto es
\[
 R_N=\sum_{k=0}^{\infty}\frac1{(N+1+k)!}>0.
\]
Para \(k\geq0\),
\[
 (N+1+k)!
 \geq(N+1)!(N+2)^k.
\]
Por tanto,
\[
 R_N
 \leq\frac1{(N+1)!}
 \sum_{k=0}^{\infty}\frac1{(N+2)^k}
 =\frac{N+2}{(N+1)(N+1)!}.
\]
La desigualdad es estricta desde \(k=2\).
\end{proof}

Definimos
\[
 J_N^{(e)}=
 \left[
 S_N,\,
 S_N+\frac{N+2}{(N+1)(N+1)!}
 \right]
\]
y, para la parte fraccionaria,
\[
 \widetilde J_N^{(e)}=J_N^{(e)}-2.
\]
La anchura converge a cero. Para cada \(M\), existe \(N\) tal que
\[
 \operatorname{diam}J_N^{(e)}<10^{-M}.
\]
Si los extremos pertenecen a una misma celda decimal, las primeras \(M\)
cifras quedan certificadas mediante aritmética racional.

\subsection{Prolongación coinductiva de la sección exponencial mediante subcilindros ternarios del Topological Phase Kernel}

En el nivel \(K\), TPK y la relación \(\operatorname{Ext}\) han generado ya
un único subcilindro compatible. Se toma el menor \(N=N_e(K)\) para el cual
la horquilla factorial queda contenida en ese subcilindro. \(N_e(K)\) es un
orden de certificación analítica, no un índice generador ni una regla de
selección de la fibra.

Los cilindros TPK son compatibles bajo truncamiento con independencia del
semigrupo. El teorema semigrupal identifica después su evaluación
arquimediana con \(e-2\); la parte entera recupera \(e\). Si fallan la órbita,
\(\operatorname{Ext}\) o el truncamiento, falla la generación; si fallan la
recurrencia o la cota factorial, falla este certificado analítico, no la
genealogía APP--TRIT--TPK.

\subsection{Certificado adversarial de la generación de \(e\): unicidad orbital, recurrencia factorial, horquillas racionales y prolongación cilíndrica}

\begin{criterion}[Falsación del canal exponencial]
La construcción queda refutada si:
\begin{enumerate}
 \item otra órbita satisface el predicado finito completo;
 \item la firma \(7[1828\,1828]459\) no es posicionalmente única;
 \item coinciden los preestados o cocientes de los dos bloques \(102011\);
 \item el semigrupo no fuerza \((n+1)a_{n+1}=a_n\);
 \item la cota \eqref{p12-83-c9-s2:eq:horquilla-factorial-e} falla;
 \item un nivel no localiza de forma unitaria una extensión compatible;
 \item una expansión tabulada interviene antes del reconocimiento.
\end{enumerate}
\end{criterion}


% Unidad U014. Biografía estructural completa del carácter de autoescala.

\section{Incidencia modal, rayo de Perron y construcción del carácter de autoescala}
\label{p12-83-c9-s3:chap:biografia-completa-phi}

\subsection{Selección finita de la semilla de autoescala}

Sea \(\mathcal U_6\subset\{0,\ldots,999\}^6\) el catálogo de las
\(468\) emisiones distintas construido en el
\cref{chap:semillas-emisor-54}, y sea
\[
 \Psi:\mathcal U_6\longrightarrow\mathbb F_3^6,
 \qquad
 \Psi(U)=((-U_j)\bmod3)_{j=1}^{6}.
\]
Para \(w\in\Psi(\mathcal U_6)\), recordamos los invariantes
\[
 n(w)=\#\Psi^{-1}(w),
 \qquad
 M(w)=\sum_{U\in\Psi^{-1}(w)}\operatorname{mult}(U),
\]
y el vector de ocupación
\[
 \nu(w)=
 \bigl(\#\{j:w_j=0\},\#\{j:w_j=1\},\#\{j:w_j=2\}\bigr).
\]

\begin{theorem}[Selección orbital de la autoescala]
\label{p12-83-c9-s3:thm:seleccion-orbital-phi}
La acción diedral sobre \(\Psi(\mathcal U_6)\) contiene una única órbita
que satisface el predicado de autoescala
\[
 n(w)=1,
 \qquad
 M(w)=144,
 \qquad
 \nu(w)=(2,2,2).
\]
El calibre interno
\[
 \operatorname{diag}_c=6,
 \qquad
 \operatorname{card}=ES
\]
selecciona en esa órbita la palabra
\[
 \boxed{w_6^{\varphi}=121200.}
\]
Su única emisión es
\[
 \boxed{U_{\varphi}=(830,943,830,109,252,252),}
\]
de multiplicidad \(144\), firma multiplicativa \(555933\), orientación
aditiva \(ES\), soporte multiplicativo de cardinal seis y tipo residual
\(A\).
\end{theorem}

\begin{proof}
La proyección directa da
\[
 U_{\varphi}\bmod3=(2,1,2,1,0,0),
 \qquad
 -U_{\varphi}\bmod3=(1,2,1,2,0,0).
\]
La enumeración se efectúa sobre las \(243\) palabras visibles y sus
\(43\) órbitas, no sobre una selección de ejemplos. El predicado conjunto
\(n=1\), \(M=144\), \(\nu=(2,2,2)\) deja una órbita. Dentro de ella, la
diagonal y la orientación declaradas dejan una emisión. En esta etapa no
se ha introducido la ecuación \(x^2-x-1=0\), una tabla decimal ni el nombre
convencional de la constante.
\end{proof}

\subsection{Elevación por bloques y 1.ª frontera coinductiva}

Sea \(b_1=w_6^{\varphi}\), y sea
\(\varepsilon=(0,0,1,1)\). Con las matrices \(L_0,L_1\) construidas en el
\cref{hmtctx:p12:chap:orbitas-elevacion-r36}, definimos
\[
 b_{k+1}=b_kL_{\varepsilon_k}\pmod3
 \qquad(1\leq k\leq4),
\]
y concatenamos
\[
 w_{6(k+1)}^{\varphi}=w_{6k}^{\varphi}\mathbin{\|}b_{k+1}.
\]
El cálculo matricial produce, bloque por bloque,
\[
\begin{aligned}
 b_1&=121200,\\
 b_2&=112202,\\
 b_3&=121020,\\
 b_4&=010210,\\
 b_5&=010200,
\end{aligned}
\]
y, por tanto,
\begin{equation}
 \boxed{
 w_{30}^{\varphi}
 =121200|112202|121020|010210|010200.}
 \label{p12-83-c9-s3:eq:w30-phi}
\end{equation}
La neutralidad dual de la frontera añade
\[
 u_5^{\varphi}=102011,
\]
de modo que
\begin{equation}
 \boxed{
 w_{36}^{\varphi}
 =121200|112202|121020|010210|010200|102011.}
 \label{p12-83-c9-s3:eq:w36-phi}
\end{equation}

\begin{remark}[Dos usos distintos de una palabra visible]
La palabra \(102011\) es también visible en la biografía de \(e\). La
igualdad de una proyección no identifica los estados holonómicos integrales: el canal,
el prefijo, el cociente, el acarreo, la fase y la memoria pertenecen al tipo
del objeto. Esta separación es obligatoria antes de construir cualquier
carácter real.
\end{remark}

\subsection{Del calendario trítico al multigrafo de incidencia}

\subsubsection{Incidencia modal primitiva \(F_{\rm av}\) y dinámica de memoria de orden \(1\)}

El calendario trítico primitivo es
\[
 (+1,-1,0).
\]
Fijamos dos modos visibles, \(\Sigma\) y \(\Pi\), y la regla
\begin{equation}
 +1:m\longmapsto\Sigma,
 \qquad
 -1:m\longmapsto\Pi,
 \qquad
 0:m\longmapsto m.
 \label{p12-83-c9-s3:eq:regla-modal-primitiva-phi}
\end{equation}
Con la condición de cierre cíclico, la órbita modal producida por
\eqref{p12-83-c9-s3:eq:regla-modal-primitiva-phi} es
\[
 \boxed{(\Sigma,\Pi,\Pi).}
\]
Los tres pares consecutivos son
\[
 \Sigma\longrightarrow\Pi,
 \qquad
 \Pi\longrightarrow\Pi,
 \qquad
 \Pi\longrightarrow\Sigma.
\]

\begin{definition}[Matriz de incidencia cronológica]
Para \(i,j\in\{\Sigma,\Pi\}\), sea
\[
 N_3(i,j)=
 \#\{r\in\mathbb Z/3\mathbb Z:(m_r,m_{r+1})=(i,j)\}.
\]
El subíndice \(3\) cuenta instantes del calendario primitivo; no denota una
potencia matricial.
\end{definition}

\begin{theorem}[Descenso necesario de la incidencia modal]
\label{p12-83-c9-s3:thm:descenso-incidencia-modal-phi}
En el orden \((\Sigma,\Pi)\),
\begin{equation}
 \boxed{
 N_3=
 \begin{pmatrix}
 0&1\\
 1&1
 \end{pmatrix}.}
 \label{p12-83-c9-s3:eq:N3-phi}
\end{equation}
Por repetición exacta del calendario,
\begin{equation}
 \boxed{
 N_9=3N_3,
 \qquad
 N_{27}=9N_3,
 \qquad
 N_{108}=36N_3.}
 \label{p12-83-c9-s3:eq:incidencias-9-27-108-phi}
\end{equation}
Estas tres igualdades son igualdades de conteos de aristas; no afirman
\(N_9=N_3^3\), \(N_{27}=N_3^9\) ni \(N_{108}=N_3^{36}\).
\end{theorem}

\begin{proof}
No aparece \(\Sigma\to\Sigma\); cada uno de los otros tres pares aparece
una vez. Esto determina las cuatro entradas de \(N_3\). Los calendarios de
longitud \(9\), \(27\) y \(108\) contienen respectivamente \(3\), \(9\)
y \(36\) copias del bloque primitivo. Cada repetición multiplica todos los
conteos por el mismo entero, lo que prueba
\eqref{p12-83-c9-s3:eq:incidencias-9-27-108-phi}.
\end{proof}

\subsubsection{Publicación areal–volumétrica y razón estable \(\varphi^{-2}\) bajo la medida de Parry}

La etiqueta modal y la hoja geométrica son objetos distintos. La carta
de publicación es
\[
 \mathfrak g_{\rm av}(\Sigma)=\mathscr H_{\rm ar},
 \qquad
 \mathfrak g_{\rm av}(\Pi)=\mathscr H_{\rm vol}.
\]
Transportando \eqref{p12-83-c9-s3:eq:N3-phi} por \(\mathfrak g_{\rm av}\), obtenemos
el operador de incidencia
\begin{equation}
 \boxed{
 F_{\rm av}=
 \begin{pmatrix}
 0&1\\
 1&1
 \end{pmatrix}.}
 \label{p12-83-c9-s3:eq:Fav-derivada-phi}
\end{equation}
Así, las dos transiciones cruzadas, la única autorretención volumétrica y la
ausencia de autorretención areal se derivan del calendario; no son cuatro
entradas independientes.

\begin{remark}[Alcance del descenso]
Olvidar la fase elemental no produce por sí solo un cociente de Markov
fuerte del Topological Phase Kernel, porque el modo visible no determina
cuál de los tres operadores internos actuará en el siguiente paso. Lo que
sí desciende sin hipótesis adicional es el multigrafo de aristas de un
bloque completo. Su envolvente mínimo de memoria uno tiene matriz de
adyacencia \(F_{\rm av}\).
\end{remark}

\subsection{Rayo positivo invariante y unicidad de la autoescala}

\begin{theorem}[Rayo de Perron derivado]
\label{p12-83-c9-s3:thm:rayo-perron-phi}
El cono positivo de \(F_{\rm av}\) contiene un único rayo propio. Si lo
normalizamos como \(v=(1,x)^{\mathsf T}\), entonces
\[
 x^2-x-1=0,
 \qquad
 x>0,
\]
y, por tanto,
\begin{equation}
 \boxed{x=\frac{1+\sqrt5}{2}=\varphi.}
 \label{p12-83-c9-s3:eq:phi-perron}
\end{equation}
\end{theorem}

\begin{proof}
La ecuación \(F_{\rm av}v=\lambda v\) da, coordenada a coordenada,
\[
 x=\lambda,
 \qquad
 1+x=\lambda x.
\]
Al eliminar \(\lambda\) resulta \(x^2-x-1=0\). Sus raíces son
\((1\pm\sqrt5)/2\), y sólo la positiva pertenece al interior del cono. La
matriz \(F_{\rm av}\) es primitiva, pues \(F_{\rm av}^2\) tiene todas sus
entradas positivas; el teorema de Perron--Frobenius garantiza que ese rayo
positivo es simple y único. El nombre \(\varphi\) se asigna después de esta
construcción.
\end{proof}

\begin{corollary}[Ecuación escalar de autoescala]
\[
 \varphi=1+\frac1\varphi,
 \qquad
 \varphi^2=\varphi+1,
 \qquad
 \varphi^{-1}=\varphi-1.
\]
Estas identidades son consecuencias de \eqref{p12-83-c9-s3:eq:phi-perron}; no seleccionan
la semilla ni la matriz.
\end{corollary}

\subsection{Recurrencia de Fibonacci del rayo de Perron y encaje diofántico de las horquillas racionales}

Sea \(F_0=0\), \(F_1=1\) y
\[
 F_{n+1}=F_n+F_{n-1}\qquad(n\geq1).
\]

\begin{theorem}[Potencias exactas de la incidencia]
\label{p12-83-c9-s3:thm:potencias-fibonacci-phi}
Para todo \(n\geq1\),
\begin{equation}
 \boxed{
 F_{\rm av}^{\,n}=
 \begin{pmatrix}
 F_{n-1}&F_n\\
 F_n&F_{n+1}
 \end{pmatrix}.}
 \label{p12-83-c9-s3:eq:potencias-Fav-fibonacci}
\end{equation}
\end{theorem}

\begin{proof}
Para \(n=1\), la identidad es la definición de \(F_{\rm av}\). Si vale
para \(n\), la multiplicación a la derecha por \(F_{\rm av}\) da
\[
 \begin{pmatrix}
 F_n&F_{n-1}+F_n\\
 F_{n+1}&F_n+F_{n+1}
 \end{pmatrix}
 =
 \begin{pmatrix}
 F_n&F_{n+1}\\
 F_{n+1}&F_{n+2}
 \end{pmatrix},
\]
que es el caso \(n+1\).
\end{proof}

\begin{theorem}[Identidad de Cassini]
\label{p12-83-c9-s3:thm:cassini-phi}
Para todo \(n\geq1\),
\begin{equation}
 \boxed{F_{n+1}F_{n-1}-F_n^2=(-1)^n.}
 \label{p12-83-c9-s3:eq:cassini-phi}
\end{equation}
\end{theorem}

\begin{proof}
El determinante de \(F_{\rm av}\) es \(-1\). Tomando determinantes en
\eqref{p12-83-c9-s3:eq:potencias-Fav-fibonacci},
\[
 F_{n-1}F_{n+1}-F_n^2
 =\det(F_{\rm av}^{\,n})=(-1)^n.
\]
\end{proof}

Definimos \(r_n=F_{n+1}/F_n\) para \(n\geq1\). De Cassini se sigue
\[
 r_n-r_{n-1}
 =\frac{F_{n+1}F_{n-1}-F_n^2}{F_nF_{n-1}}
 =\frac{(-1)^n}{F_nF_{n-1}}.
\]
Por tanto, las aproximaciones alternan alrededor del rayo de Perron.

\begin{proposition}[Horquillas racionales de Perron]
\label{p12-83-c9-s3:prop:horquillas-perron-phi}
Para \(m\geq1\),
\begin{equation}
 \frac{F_{2m+2}}{F_{2m+1}}
 <\varphi<
 \frac{F_{2m+1}}{F_{2m}},
 \label{p12-83-c9-s3:eq:horquilla-fibonacci-phi}
\end{equation}
y la anchura es
\begin{equation}
 \boxed{
 \frac{F_{2m+1}}{F_{2m}}
 -\frac{F_{2m+2}}{F_{2m+1}}
 =\frac{1}{F_{2m}F_{2m+1}}.}
 \label{p12-83-c9-s3:eq:anchura-fibonacci-phi}
\end{equation}
\end{proposition}

\begin{proof}
La alternancia anterior sitúa los cocientes pares por debajo y los impares
por encima de \(\varphi\). La fórmula de la anchura es
\eqref{p12-83-c9-s3:eq:cassini-phi} con \(n=2m+1\), tras poner ambas fracciones sobre
denominador común. Como \(F_n\to\infty\), la anchura tiende a cero.
\end{proof}

\subsection{Incidencia de memoria de orden \(1\) y medida de Parry del envolvente simbólico}

Sea \(X_{\rm av}\) el menor subshift de tipo finito sobre
\(\{\mathscr H_{\rm ar},\mathscr H_{\rm vol}\}\) que admite exactamente
los tres pares observados en el calendario modal. Su matriz de adyacencia es
\(F_{\rm av}\).

\begin{theorem}[Transición y medida de Parry]
\label{p12-83-c9-s3:thm:parry-phi}
Sea \(r=(1,\varphi)^{\mathsf T}\). La matriz estocástica de Parry es
\begin{equation}
 \boxed{
 P_{ij}=
 \frac{(F_{\rm av})_{ij}r_j}{\varphi r_i}
 =
 \begin{pmatrix}
 0&1\\
 \varphi^{-2}&\varphi^{-1}
 \end{pmatrix}.}
 \label{p12-83-c9-s3:eq:parry-matriz-phi}
\end{equation}
La medida estacionaria de vértices es proporcional a
\((1,\varphi^2)\), luego
\begin{equation}
 \boxed{
 \frac{\mu_{\rm ar}}{\mu_{\rm vol}}=\varphi^{-2}.}
 \label{p12-83-c9-s3:eq:parry-razon-phi}
\end{equation}
\end{theorem}

\begin{proof}
La suma de la fila \(i\) de \(P\) es
\[
 \frac{(F_{\rm av}r)_i}{\varphi r_i}=1.
\]
Como \(F_{\rm av}\) es simétrica, los vectores propios izquierdo y derecho
coinciden. Los pesos estacionarios son proporcionales a sus productos
coordenada a coordenada, es decir, a \((1,\varphi^2)\). La razón entre los
dos pesos es \(\varphi^{-2}\).
\end{proof}

\begin{remark}[Tres medidas que no deben confundirse]
La frecuencia cronológica de la órbita primitiva es
\[
 \nu_{\rm cyc}(\mathscr H_{\rm ar})=\frac13,
 \qquad
 \nu_{\rm cyc}(\mathscr H_{\rm vol})=\frac23.
\]
La medida uniforme \(\tau_{468}\) vive sobre el catálogo de emisiones. La
medida de Parry vive sobre historias admisibles de longitud arbitraria. La
primera cuenta instantes, la segunda cuenta emisiones y la tercera pondera
cilindros dinámicos. Ninguna es una aproximación de las otras.
\end{remark}

\subsection{Carácter multiplicativo y conservación cilíndrica}

Para una ruta admisible \(\gamma:i\to j\) de longitud \(n\), definimos
\begin{equation}
 \chi_P(\gamma)=\varphi^n\frac{r_i}{r_j}.
 \label{p12-83-c9-s3:eq:caracter-parry-phi}
\end{equation}

\begin{proposition}[Multiplicatividad por concatenación]
\label{p12-83-c9-s3:prop:multiplicatividad-parry-phi}
Si \(\gamma_1:i\to j\) y \(\gamma_2:j\to k\), entonces
\[
 \chi_P(\gamma_2\circ\gamma_1)
 =\chi_P(\gamma_2)\chi_P(\gamma_1).
\]
Sobre una ruta cerrada de longitud \(n\),
\[
 \chi_P(\gamma)=\varphi^n.
\]
\end{proposition}

\begin{proof}
La longitud se suma y la coordenada intermedia telescopa:
\[
 \varphi^{n_1+n_2}\frac{r_i}{r_k}
 =\left(\varphi^{n_1}\frac{r_i}{r_j}\right)
  \left(\varphi^{n_2}\frac{r_j}{r_k}\right).
\]
Si \(i=k\), el cociente terminal es uno.
\end{proof}

Normalizamos el vector izquierdo como
\[
 \ell=\frac{r}{1+\varphi^2},
 \qquad
 \ell^{\mathsf T}r=1.
\]
La medida de un cilindro es
\[
 \mu[x_0\cdots x_n]
 =\ell_{x_0}r_{x_n}\varphi^{-n}.
\]
Por tanto,
\begin{equation}
 V_n(x)
 :=\frac{\mu[x_0]}{\mu[x_0\cdots x_n]}
 =\varphi^n\frac{r_{x_0}}{r_{x_n}}
 =\chi_P(x_0\cdots x_n).
 \label{p12-83-c9-s3:eq:volumen-cilindrico-phi}
\end{equation}
Para una dimensión tipada \(D>0\), la escala
\[
 a_n^{(D)}(x)=V_n(x)^{1/D}
\]
satisface la ley exacta
\begin{equation}
 \boxed{
 \bigl(a_n^{(D)}(x)\bigr)^D
 \mu[x_0\cdots x_n]=\mu[x_0].}
 \label{p12-83-c9-s3:eq:conservacion-cilindrica-phi}
\end{equation}

En dimensión tres y sobre retornos cerrados,
\[
 \frac{a_9}{a_0}=\varphi^3,
 \qquad
 \frac{a_{27}}{a_0}=\varphi^9,
 \qquad
 \frac{a_{108}}{a_0}=\varphi^{36}.
\]
Si \(A_k=a_{9k}\), entonces
\begin{equation}
 A_{k+1}-2A_k+A_{k-1}
 =A_{k-1}(\varphi^3-1)^2>0.
 \label{p12-83-c9-s3:eq:convexidad-retornos-phi}
\end{equation}
Ésta es convexidad en profundidad de retorno; no se interpreta como una
magnitud cinemática sin una carta adicional de duración.

\subsection{Memoria cuadrática y rama contraída}

Los autovalores de \(F_{\rm av}\) son
\[
 \varphi,
 \qquad
 -\varphi^{-1}.
\]
La rama contraída cambia de orientación en cada iteración. Para conservar
ese signo antes de elevar al cuadrado, pasamos a
\(\operatorname{Sym}^2(\mathbb R^2)\). En la base
\((x^2,2xy,y^2)\),
\begin{equation}
 \operatorname{Sym}^2(F_{\rm av})=
 \begin{pmatrix}
 0&0&1\\
 0&1&2\\
 1&1&1
 \end{pmatrix}.
 \label{p12-83-c9-s3:eq:sym2-Fav-phi}
\end{equation}

\begin{theorem}[Espectro de la memoria de segundo orden]
\label{p12-83-c9-s3:thm:espectro-sym2-phi}
\[
 \det\bigl(tI-\operatorname{Sym}^2(F_{\rm av})\bigr)
 =(t+1)(t^2-3t+1),
\]
y
\begin{equation}
 \boxed{
 \operatorname{Spec}\bigl(\operatorname{Sym}^2(F_{\rm av})\bigr)
 =\{\varphi^2,-1,\varphi^{-2}\}.}
 \label{p12-83-c9-s3:eq:espectro-sym2-phi}
\end{equation}
El único modo estable positivo es \(\varphi^{-2}\).
\end{theorem}

\begin{proof}
Los autovalores de una potencia simétrica de grado dos son los productos
\(\lambda_1^2\), \(\lambda_1\lambda_2\) y \(\lambda_2^2\). Con
\(\lambda_1=\varphi\) y \(\lambda_2=-\varphi^{-1}\), se obtienen
\(\varphi^2\), \(-1\) y \(\varphi^{-2}\). La primera excede uno, la
segunda tiene módulo uno y la tercera pertenece a \((0,1)\).
\end{proof}

Sobre una ruta cerrada con \(V_n=\varphi^n\), el modo estable satisface
\begin{equation}
 \boxed{M_n=\varphi^{-2n}=V_n^{-2}.}
 \label{p12-83-c9-s3:eq:modo-estable-phi}
\end{equation}

\subsection{Realización hiperbólica normalizada de la autoescala: \(\exp(2\log\varphi\,J_h)=F_{\mathrm{av}}^2\)}

Sea
\[
 A=F_{\rm av}^{\,2}=
 \begin{pmatrix}1&1\\1&2\end{pmatrix},
 \qquad
 J_h=\frac{2A-3I}{\sqrt5}.
\]
Entonces
\[
 J_h^2=I,
 \qquad
 A=\frac32I+\frac{\sqrt5}{2}J_h.
\]
Como
\[
 \cosh(2\log\varphi)=\frac32,
 \qquad
 \sinh(2\log\varphi)=\frac{\sqrt5}{2},
\]
se obtiene
\begin{equation}
 \boxed{
 \exp(2\log\varphi\,J_h)=F_{\rm av}^{\,2}.}
 \label{p12-83-c9-s3:eq:realizacion-hiperbolica-phi}
\end{equation}
La constante de autoescala es, por tanto, simultáneamente el autovalor de
Perron de la incidencia y el parámetro exponencial de su realización
hiperbólica normalizada.

\subsection{Caracterización, horquillas y prolongación}

La construcción finita determina el rayo positivo de \(F_{\rm av}\). La
caracterización analítica posterior es
\[
 x^2-x-1=0,
 \qquad x>0.
\]
Las horquillas \eqref{p12-83-c9-s3:eq:horquilla-fibonacci-phi} tienen extremos racionales
y anchura explícita \eqref{p12-83-c9-s3:eq:anchura-fibonacci-phi}. Para cada nivel
nonádico \(K\), elegimos el menor \(m=m_{\varphi}(K)\) cuya horquilla queda
contenida en un único subcilindro entre los \(729\) elementos de la fibra
ambiente. La minimalidad convierte la localización en una función, no en una
elección tácita.

Los cilindros seleccionados son compatibles con los truncamientos y sus
diámetros tienden a cero. Su intersección contiene un único real; el
\cref{p12-83-c9-s3:thm:rayo-perron-phi} lo identifica con \(\varphi\). Los cocientes de
Fibonacci son, pues, certificados racionales de la localización, no datos
decimales empleados para decidir la semilla.

\subsection{Certificado interno de unicidad, compatibilidad y convergencia de la autoescala}

\begin{criterion}[Refutación de la biografía de autoescala]
La construcción queda refutada si ocurre cualquiera de los hechos
siguientes:
\begin{enumerate}
 \item el predicado finito selecciona más de una órbita o más de una emisión;
 \item la recurrencia \(L_0,L_0,L_1,L_1\) no produce
       \eqref{p12-83-c9-s3:eq:w30-phi};
 \item el calendario no produce exactamente los tres pares usados en
       \eqref{p12-83-c9-s3:eq:N3-phi};
 \item se confunde un múltiplo de conteo \(N_{9},N_{27},N_{108}\) con una
       potencia de \(N_3\);
 \item \(F_{\rm av}\) posee otro rayo propio positivo;
 \item falla \eqref{p12-83-c9-s3:eq:potencias-Fav-fibonacci} o la identidad de Cassini;
 \item las horquillas racionales no son encajadas o no contraen su anchura;
 \item se identifica la frecuencia cronológica, la medida uniforme del
       catálogo o la medida de Parry;
 \item una tabla decimal interviene antes de construir el rayo de Perron;
 \item un cilindro seleccionado no trunca al cilindro anterior.
\end{enumerate}
\end{criterion}


% Unidad U015. Acoplamiento de horquillas racionales, refinamiento TPK y
% certificación a profundidad arbitraria.

\section{Certificación diagonal, límite inverso y generación a precisión arbitraria}
\label{p12-83-c9-s4:chap:precision-arbitraria-constantes}

\subsection{Separación tipada de los \(4\) índices de la prolongación nonádica}

Esta unidad tiene estatuto de certificación posterior. Los prefijos y sus
extensiones compatibles proceden de la relación interna
\(\operatorname{Ext}\), de la supervivencia coinductiva y de las leyes de
transición APP--TRIT--TPK. Las horquillas racionales que siguen no eligen
entre los \(729\) elementos de una fibra: localizan en la carta arquimediana
la extensión ya generada y certifican su precisión a profundidad arbitraria.

Para cada carácter \(\chi\in\{\pi,e,\varphi\}\) intervienen cuatro
parámetros de naturaleza diferente:
\begin{enumerate}
 \item \(n\), orden de la horquilla analítica racional;
 \item \(K\), nivel de refinamiento del Topological Phase Kernel;
 \item \(L_K=6K\), longitud ternaria del prefijo publicado;
 \item \(M\), número de posiciones decimales solicitado.
\end{enumerate}
No se postula \(n=K\), \(n=L_K\), \(K=M\) ni una proporcionalidad fija
entre ellos. El acoplamiento se construye mediante inclusiones de intervalos
decididas por aritmética entera.

Las partes enteras son
\[
 m_\pi=3,
 \qquad
 m_e=2,
 \qquad
 m_\varphi=1,
\]
y los caracteres fraccionarios son
\[
 \rho_\pi=\pi-3,
 \qquad
 \rho_e=e-2,
 \qquad
 \rho_\varphi=\varphi-1.
\]
Las tres familias de horquillas semiabiertas se escriben
\[
 J_{\chi,n}=[\ell_{\chi,n},h_{\chi,n}).
\]
Sus extremos son racionales, contienen a \(\rho_\chi\) y satisfacen
\[
 \ell_{\chi,n}\nearrow\rho_\chi,
 \qquad
 h_{\chi,n}\searrow\rho_\chi,
 \qquad
 h_{\chi,n}-\ell_{\chi,n}\longrightarrow0.
\]
Para \(\pi\), los extremos proceden de las sumas alternadas de la identidad
de Machin demostrada en el \cref{chap:biografia-pi-autonoma}; para \(e\),
de las sumas factoriales y su resto racional del
\cref{p12-83-c9-s2:chap:biografia-completa-e}; para \(\varphi\), de los cocientes de
Fibonacci y Cassini del \cref{p12-83-c9-s3:chap:biografia-completa-phi}.

\subsection{Cilindros ternarios y fibra ambiente inmediata}

Un prefijo ternario de longitud \(L_K=6K\) se representa por un entero
\(N_K\) con
\[
 0\leq N_K<3^{L_K}.
\]
Su cilindro semiabierto es
\begin{equation}
 C(N_K,L_K)=
 \left[
 \frac{N_K}{3^{L_K}},
 \frac{N_K+1}{3^{L_K}}
 \right).
 \label{p12-83-c9-s4:eq:cilindro-nivel-K}
\end{equation}
La siguiente emisión añade seis trits. La fibra ambiente inmediata contiene
\(3^6=729\) elementos, indexados por \(u\in\{0,\ldots,728\}\):
\begin{equation}
 C_{K,u}=
 \left[
 \frac{729N_K+u}{3^{L_K+6}},
 \frac{729N_K+u+1}{3^{L_K+6}}
 \right).
 \label{p12-83-c9-s4:eq:729-subcilindros}
\end{equation}
Los \(C_{K,u}\) son disjuntos y
\[
 C(N_K,L_K)=\bigsqcup_{u=0}^{728}C_{K,u}.
\]
El número \(729\) cuenta coordenadas de esta fibra ambiente; no cuenta
semillas, emisiones distintas, órbitas ni estados sobrevivientes.

\subsubsection{Comparación diofántica de los extremos racionales de la horquilla}

Sea \(J=[\ell,h)\subset C(N_K,L_K)\), con
\(\ell,h\in\mathbb Q\). Escribimos \(S_K=3^{L_K+6}\) y definimos
\begin{align}
 j_{\min}(J,K)
 &=\max\!\left(729N_K,\left\lfloor S_K\ell\right\rfloor\right),
 \label{p12-83-c9-s4:eq:jmin-diagonal}\\
 j_{\max}(J,K)
 &=\min\!\left(729N_K+728,\left\lceil S_Kh\right\rceil-1\right).
 \label{p12-83-c9-s4:eq:jmax-diagonal}
\end{align}

\begin{lemma}[Criterio de localización unitaria]
\label{p12-83-c9-s4:lem:criterio-localizacion-unitaria}
La horquilla \(J\) está contenida en un único elemento de la fibra
\eqref{p12-83-c9-s4:eq:729-subcilindros} si y sólo si
\begin{equation}
 \boxed{j_{\min}(J,K)=j_{\max}(J,K).}
 \label{p12-83-c9-s4:eq:jmin-igual-jmax}
\end{equation}
En ese caso, si \(j\) es el valor común,
\[
 u=j-729N_K,
 \qquad
 N_{K+1}=j.
\]
\end{lemma}

\begin{proof}
Los enteros \(j\) cuyos intervalos elementales de longitud \(S_K^{-1}\)
intersectan \(J\) forman el intervalo entero
\[
 \left[
 \lfloor S_K\ell\rfloor,
 \lceil S_Kh\rceil-1
 \right]\cap
 [729N_K,729N_K+728].
\]
Éste contiene exactamente un elemento si y sólo si sus extremos coinciden.
La igualdad identifica el numerador del subcilindro y, al truncar sus seis
últimos trits, se recupera \(N_K\).
\end{proof}

\subsection{Certificación diagonal mínima}

La generación no fija de antemano un mismo orden analítico para todos los
niveles. Una vez producida la extensión compatible, el orden de la
horquilla se aumenta sólo hasta certificarla de manera unitaria.

\begin{definition}[Recurrencia diagonal de certificación]
\label{p12-83-c9-s4:def:recurrencia-diagonal-constantes}
Fijada por \(\operatorname{Ext}\) una familia de cilindros compatibles
\(C(N_K,L_K)\) cuya publicación contiene a \(\rho_\chi\), se elige un orden
inicial \(n_\chi(K_0-1)\geq1\) y definimos recursivamente
\begin{equation}
 n_\chi(K)=
 \min\left\{
 n\geq n_\chi(K-1):
 j_{\min}(J_{\chi,n},K)=j_{\max}(J_{\chi,n},K)
 \right\},
 \label{p12-83-c9-s4:eq:nchi-K-minimo}
\end{equation}
y verificamos
\begin{equation}
 N_{\chi,K+1}
 =j_{\min}(J_{\chi,n_\chi(K)},K).
 \label{p12-83-c9-s4:eq:Nchi-K-mas-uno}
\end{equation}
El conjunto
\[
 \Delta_\chi=
 \{(n_\chi(K),K):K\geq K_0\}
\]
es la diagonal analítico--nonádica del carácter \(\chi\).
\end{definition}

\begin{theorem}[Existencia y unicidad de la diagonal de certificación]
\label{p12-83-c9-s4:thm:existencia-unicidad-diagonal}
Para \(\chi\in\{\pi,e,\varphi\}\), la recurrencia
\eqref{p12-83-c9-s4:eq:nchi-K-minimo} está definida en todo nivel finito. El índice
\(n_\chi(K)\) y la coordenada \(N_{\chi,K+1}\) son únicos.
\end{theorem}

\begin{proof}
Supongamos generada internamente la extensión cilíndrica de nivel \(K\) que contiene a
\(\rho_\chi\). Como \(\rho_\chi\) es irracional, no pertenece a ninguna
frontera ternaria \(3^{-L}\mathbb Z\). Por ello está en el interior de un
único subcilindro \(C_{K,u_*}\), a distancia positiva de sus dos extremos.
Las horquillas \(J_{\chi,n}\) contienen a \(\rho_\chi\) y su diámetro
tiende a cero; existe, por tanto, un \(n\) para el que la horquilla queda
contenida en \(C_{K,u_*}\). El conjunto del mínimo en
\eqref{p12-83-c9-s4:eq:nchi-K-minimo} es no vacío y, por el buen orden de
\(\mathbb N\), posee un menor elemento. El
\cref{p12-83-c9-s4:lem:criterio-localizacion-unitaria} prueba que la horquilla la
localiza de manera unitaria y que el numerador localizado coincide con el ya
generado. La inducción sobre \(K\) completa el argumento.
\end{proof}

\begin{corollary}[Ausencia de parámetros decimales]
La diagonal se calcula con sumas, productos, potencias, divisiones euclídeas,
pisos y techos de enteros. No recibe una cifra decimal de \(\pi\), \(e\) o
\(\varphi\) como dato de certificación.
\end{corollary}

\subsection{Compatibilidad con la conexión nonádica conjunta}

Sea \(\widehat x_{\chi,K}\) el estado holonómico integral del canal \(\chi\) en
el nivel \(K\). Sus coordenadas incluyen el prefijo, el cilindro, la fase,
el cociente, el acarreo, la memoria vertical no acotada, el término terminal
y el registro de procedencia. La fase residual es
\[
 g_0(K)=[K]_9\in\mathbb Z/9\mathbb Z,
\]
mientras la etiqueta pública es
\[
 g(K)=1+((K-1)\bmod9)\in\{1,\ldots,9\}.
\]
La memoria vertical \(q_{\rm mem}(K)\in\mathbb Z_{\geq0}\) no se reduce
módulo \(12\); sólo su proyección dodecafásica
\[
 \beta(K)=q_{\rm mem}(K)\bmod12
\]
lo hace.

Para cada \(K\), sea
\[
 \Gamma_{9,K}:\widehat X_K\longrightarrow\widehat X_{K+9}
\]
la holonomía de una vuelta de fase. Entonces
\[
 g_0(K+9)=g_0(K),
 \qquad
 q_{\rm mem}(K+9)=q_{\rm mem}(K)+1,
 \qquad
 \beta(K+9)=\beta(K)+1\pmod{12}.
\]

\begin{theorem}[Retorno de fase con conservación genealógica]
\label{p12-83-c9-s4:thm:retorno-fase-conservacion-genealogica}
La igualdad de fase no reinicia el estado:
\[
 g_0(K+9)=g_0(K),
 \qquad
 \widehat x_{\chi,K+9}\ne\widehat x_{\chi,K}.
\]
Las profundidades \(27\), \(54\) y \(108\) son las composiciones tipadas
\begin{align*}
 \Gamma_{27,K}
 &=\Gamma_{9,K+18}\circ\Gamma_{9,K+9}\circ\Gamma_{9,K},\\
 \Gamma_{54,K}
 &=\Gamma_{9,K+45}\circ\cdots\circ\Gamma_{9,K},\\
 \Gamma_{108,K}
 &=\Gamma_{9,K+99}\circ\cdots\circ\Gamma_{9,K}.
\end{align*}
En ningún caso son reinicios del prefijo, del acarreo, del terminal o de la
memoria.
\end{theorem}

\begin{proof}
La primera coordenada se calcula en \(\mathbb Z/9\mathbb Z\), mientras
\(q_{\rm mem}\) pertenece a un monoide no acotado. Cada aplicación de
\(\Gamma_{9,K}\) añade una unidad de memoria, seis bloques por nivel
intermedio y la composición correspondiente de cociclos. Por ello la fase
puede coincidir aunque el estado holonómico integral no lo haga. Las fórmulas de
profundidad muestran los índices de cada mapa y evitan componer como si todos
fueran endomorfismos del mismo espacio.
\end{proof}

\subsubsection{Compresión, expresión y término terminal}

La realización operatorial de una vuelta satisface
\begin{equation}
 U_\Gamma J=JT+\eta,
 \qquad
 J^*\eta=0,
 \qquad
 I-T^*T=\eta^*\eta.
 \label{p12-83-c9-s4:eq:compresion-expresion-precision}
\end{equation}
La contracción visible estricta se reserva para
\[
 M_9=\frac59V_9;
\]
no se identifica con el operador de compresión \(T\). La telescopía de
profundidad \(L\) conserva el término terminal:
\begin{equation}
 S_0=
 \sum_{r=0}^{L-1}M_9^{*r}D_9^*D_9M_9^r
 +M_9^{*L}S_0M_9^L.
 \label{p12-83-c9-s4:eq:telescopia-terminal-precision}
\end{equation}
Omitir el último sumando antes de probar su desaparición borraría
información de frontera y convertiría el retorno de fase en un reinicio
falso.

\subsection{Sistema inverso de historias compatibles}

Sea \(\mathcal H_K^\chi\) el conjunto de historias enriquecidas admisibles
del canal \(\chi\) a profundidad \(K\), y sea
\[
 p_{N,K}:\mathcal H_N^\chi\longrightarrow\mathcal H_K^\chi
 \qquad(N\geq K)
\]
el truncamiento. Definimos
\[
 \operatorname{Surv}_K^{(N)}(\chi)
 =p_{N,K}(\mathcal H_N^\chi)
\]
y
\begin{equation}
 \operatorname{Surv}_K(\chi)
 =\bigcap_{N\geq K}\operatorname{Surv}_K^{(N)}(\chi).
 \label{p12-83-c9-s4:eq:supervivencia-todas-profundidades}
\end{equation}
El objeto coinductivo es
\begin{equation}
 X_\chi=
 \varprojlim_K
 \bigl(\operatorname{Surv}_K(\chi),p_{K+1,K}\bigr).
 \label{p12-83-c9-s4:eq:limite-inverso-caracteres}
\end{equation}

\begin{theorem}[Sección compatible bajo prolongabilidad enriquecida]
\label{p12-83-c9-s4:thm:seccion-compatible-caracteres}
Supóngase además que, sobre cada subcilindro numérico generado, la
relación \(\operatorname{Ext}\) es no vacía, univaluada y compatible con los
truncamientos. Entonces \(\operatorname{Ext}\) y la supervivencia producen una familia
\[
 x_\chi=(\widehat x_{\chi,K})_{K\geq K_0}
\]
que satisface
\[
 p_{K+1,K}(\widehat x_{\chi,K+1})
 =\widehat x_{\chi,K}.
\]
Por tanto, \(x_\chi\in X_\chi\).
\end{theorem}

\begin{proof}
El subcilindro de nivel \(K+1\) tiene numerador
\(N_{\chi,K+1}=729N_{\chi,K}+u\). El truncamiento elimina sus seis últimos
trits y recupera \(N_{\chi,K}\). Las demás coordenadas se obtienen por las
leyes functoriales de acarreo, memoria, terminal y fase. La hipótesis de
univaluación de \(\operatorname{Ext}\) fija el subcilindro y la extensión
enriquecida; la unicidad local del
\cref{p12-83-c9-s4:lem:criterio-localizacion-unitaria} certifica después su
localización arquimediana, y la compatibilidad de \(\operatorname{Ext}\)
proporciona la igualdad de truncamiento.
\end{proof}

\subsection{Celdas decimales y precisión prescrita}

Para \(M\geq1\) y \(q\in\{0,\ldots,10^M-1\}\), sea
\[
 D_{M,q}=
 \left[
 \frac{q}{10^M},
 \frac{q+1}{10^M}
 \right).
\]
Una horquilla \(J=[\ell,h)\) está contenida en una única celda decimal si
y sólo si
\begin{equation}
 \boxed{
 \left\lfloor10^M\ell\right\rfloor
 =\left\lceil10^Mh\right\rceil-1=q.}
 \label{p12-83-c9-s4:eq:criterio-celda-decimal}
\end{equation}

\begin{definition}[Nivel mínimo de certificación decimal]
Para \(\chi\in\{\pi,e,\varphi\}\), definimos \(K_\chi(M)\) como el menor
nivel para el que existen un orden analítico \(n\) y un entero \(q\) tales
que
\begin{enumerate}
 \item \(J_{\chi,n}\) satisface
       \eqref{p12-83-c9-s4:eq:criterio-celda-decimal};
 \item \(J_{\chi,n}\) queda contenido en el cilindro generado por el
       estado \(\widehat x_{\chi,K}\);
 \item la transición al nivel siguiente satisface
       \eqref{p12-83-c9-s4:eq:jmin-igual-jmax}.
\end{enumerate}
El par \((K,n)\) se minimiza lexicográficamente.
\end{definition}

\begin{theorem}[Certificación a precisión decimal arbitraria de las salidas coinductivas]
\label{p12-83-c9-s4:thm:precision-decimal-arbitraria}
Para cada \(\chi\in\{\pi,e,\varphi\}\) y cada \(M\geq1\), existe
\(K_\chi(M)<\infty\). La sección generada determina un único prefijo decimal
de \(M\) posiciones, y la horquilla reconoce y certifica ese prefijo. Si
\(M'<M\), el prefijo de orden \(M\) trunca al de
orden \(M'\), y las celdas correspondientes están encajadas.
\end{theorem}

\begin{proof}
Los tres caracteres son irracionales: \(\varphi\) lo es por su polinomio
cuadrático con discriminante no cuadrado; la irracionalidad de \(e\) y
\(\pi\) se toma de sus teoremas clásicos de caracterización, posteriores a
la construcción finita. Ninguno pertenece, pues, a una frontera decimal o
ternaria racional. Para un \(M\) fijo, la distancia del carácter al conjunto
finito de fronteras decimales relevantes es positiva. Las horquillas
racionales contraen su diámetro hasta quedar en una única celda decimal.
El \cref{p12-83-c9-s4:thm:existencia-unicidad-diagonal} acopla esa horquilla con un
cilindro TPK y con una extensión única. El mínimo lexicográfico existe por
buen orden. La compatibilidad de la sección prueba el encajamiento al
aumentar \(M\).
\end{proof}

\begin{corollary}[Unicidad real]
Las celdas encajadas tienen diámetros \(10^{-M}\to0\). Su intersección es
un singleton, identificado por las caracterizaciones precedentes con
\(\pi\), \(e\) o \(\varphi\), respectivamente.
\end{corollary}

\subsection{Equivalencia de \(2\) realizaciones finitas de la ley uniforme}

Los testigos de \(1,000\) y \(10,000\) posiciones no definen los
caracteres ni limitan la prolongación. Son ejecuciones de la misma recurrencia
uniforme con dos profundidades distintas.

\begin{table}[htbp]
\centering
\caption{Dimensiones exactas de las ejecuciones certificadas, por canal.}
\label{p12-83-c9-s4:tab:dimensiones-ejecuciones-1000-10000}
\begin{tabular}{lrr}
\toprule
Objeto contado & \(1,000\) posiciones & \(10,000\) posiciones\\
\midrule
Particiones exhaustivas de la fibra ambiente & 396 & 3501\\
Trits transportados & 2376 & 21006\\
Pares de igual fase \((K,K+9)\) & 387 & 3492\\
Vueltas nonádicas completas & 43 & 388\\
Elementos examinados por partición & 729 & 729\\
Extensiones compatibles conservadas por partición & 1 & 1\\
\bottomrule
\end{tabular}
\end{table}

Para \(10,000\) posiciones,
\[
 B=3501=389\cdot9,
 \qquad
 6B=21006,
 \qquad
 3^{6B}>10^{10000}.
\]
Los primeros \(396\) bloques de la ejecución larga coinciden con la
ejecución de \(1,000\) posiciones. En cada uno de los \(3492\) pares de
igual fase se preserva la coordenada de fase y se prolongan el prefijo, la
profundidad, el cilindro y la memoria. El certificado registra también las
repeticiones eventuales de proyecciones locales sin confundirlas con retorno
del estado completo.

\subsection{Separación entre generación y verificación}

\begin{definition}[Generador estructural]
El generador recibe el catálogo finito, \(L_0,L_1\), el estado inicial, las
leyes de conexión y la relación interna \(\operatorname{Ext}\). En cada nivel:
\begin{enumerate}
 \item enumera los \(729\) elementos de la fibra ambiente;
 \item aplica \(\operatorname{Ext}\) y conserva la única extensión compatible;
 \item actualiza fase, cociente, acarreo, memoria y terminal;
 \item escribe el bloque ternario, el registro de selección y el registro genealógico.
\end{enumerate}
\end{definition}

\begin{definition}[Comprobación posterior independiente]
La comprobación posterior independiente recibe las salidas cerradas del
generador y las tres familias de horquillas racionales; calcula
\(j_{\min}\) y \(j_{\max}\), certifica su igualdad con la extensión ya
producida, reconstruye las
particiones, comprueba los cuadrados de truncamiento, los pares
\((K,K+9)\), las sumas de partición
\[
 \operatorname{killed}_{\mathrm{izq}}+1+\operatorname{killed}_{\mathrm{der}}=729,
\]
la integridad de las salidas y, al final, la coincidencia con expansiones de referencia.
\end{definition}

\begin{theorem}[Aislamiento de los datos objetivo]
\label{p12-83-c9-s4:thm:aislamiento-datos-objetivo-precision}
Las expansiones decimales de referencia y los datos metrológicos no
intervienen en la selección de una extensión compatible. La comparación se
ejecuta después de escribir los tres prefijos.
\end{theorem}

\begin{proof}
La entrada matemática de cada decisión generativa consiste en el catálogo
finito, \(L_0,L_1\), el estado heredado y la relación interna
\(\operatorname{Ext}\). Los extremos racionales, potencias de \(3\), pisos y
techos pertenecen a la comprobación posterior. El generador termina antes de
que ésta abra las horquillas o ejecute los controles decimales. La dirección
del flujo de datos es
\[
 \text{estructura}\longrightarrow
 \text{prefijos y ledgers}\longrightarrow
 \text{comprobación posterior independiente};
\]
no existe una flecha en sentido contrario.
\end{proof}

\subsection{Obstrucciones diferenciadas de la generación y de su certificación a precisión arbitraria}

\begin{criterion}[Refutación de la generación a precisión arbitraria]
Los fallos de órbita, extensión, truncamiento o retorno refutan la generación;
los fallos de horquilla, orden o control decimal refutan el certificado
analítico correspondiente; y una discrepancia entre ambos refuta la
identificación conjunta de la salida publicada. Con esta separación de tipos,
la afirmación completa queda refutada si:
\begin{enumerate}
 \item se identifica cualquiera de los índices \(n,K,L_K,M\) sin un
       teorema de enlace;
 \item una horquilla deja de contener su carácter o no contrae su diámetro;
 \item una etapa satisface \(j_{\min}\ne j_{\max}\) después del orden
       declarado como mínimo;
 \item una extensión no trunca al cilindro del nivel anterior;
 \item el retorno de fase reinicia prefijo, acarreo, terminal o memoria;
 \item se omite el término terminal en
       \eqref{p12-83-c9-s4:eq:telescopia-terminal-precision} sin demostrar su anulación;
 \item los primeros \(396\) bloques de la ejecución larga difieren de la
       ejecución corta;
 \item un control decimal o metrológico es leído antes de cerrar la salida;
 \item una cifra impresa difiere de la sección matemática certificada;
 \item la prueba depende de que la precisión solicitada sea \(1,000\) o
       \(10,000\), en vez de un \(M\) arbitrario.
\end{enumerate}
\end{criterion}
 
\chapter{Clausura dodecafásica y dos derivaciones internas de \(\alpha\)}
\label{chap:83-10}

% Unidad U016. Registro firmado, lectura terminal y reconstruccion de K.
% Redaccion autonoma desde la autoridad material 1063 y su sucesor V2.

\section{Generación inicial del estado dodecafásico firmado y reconstrucción del sello \(K\)}
\label{p12-83-c10-s1:chap:registro-dodecafasico-hadamard-k}

La construcción del sello entero no comienza con \(\alpha\), con su inversa
ni con una tabla metrológica. Su dominio es el estado holonómico terminal
del Topological Phase Kernel obtenido después de la cadena ya construida
\begin{equation}
 \boxed{
 \mathrm{APP}\longrightarrow\mathrm{TRIT}\longrightarrow\mathrm{TPK}
 \longrightarrow\Omega_{\rm seed}
 \xrightarrow{T_{\min}}\mathcal U_6^{\rm cat}
 \xrightarrow{\rho_3}W_{243}
 \longrightarrow w_{30}\longrightarrow R_{36}
 \longrightarrow G_9\longrightarrow x_{\rm term}.}
 \label{p12-83-c10-s1:eq:cadena-upstream-registro-dodecafasico}
\end{equation}
Las flechas anteriores fueron definidas con sus respectivos dominios:
pares de semillas, firmas, emisiones, palabras ternarias, fronteras,
historias compatibles y holonomía nonádica. Esta unidad no vuelve a escoger
ninguno de esos datos. Construye las cartas enteras del mismo estado
terminal y demuestra su equivalencia.

\subsection{Trazas orientadas sobre ventanas de longitud \(108\)}

Sea
\[
 \Theta_{108}=\mathbb Z/108\mathbb Z
\]
el calendario observable. Sus doce ventanas nonádicas son
\[
 E_m=\{9m+1,\ldots,9m+9\},
 \qquad 0\leq m<12.
\]
La partición es disjunta y exhaustiva:
\[
 \{1,\ldots,108\}=\bigsqcup_{m=0}^{11}E_m.
\]
El conjunto de códigos de dirección que publican un evento se fija como
\[
 D_{\rm evt}=\{2,4,6,8\}.
\]
La orientación de las doce ventanas es
\begin{equation}
 \Sigma_{12}=(+,+,+,-,-,-,+,+,+,-,-,-).
 \label{p12-83-c10-s1:eq:signatura-doce-ventanas}
\end{equation}
Si \(d_t\) es la dirección codificada en el paso \(t\), el evento firmado
de la ventana \(m\) es
\begin{equation}
 e_m=\Sigma_{12}(m)\,
 \#\{t\in E_m:d_t\in D_{\rm evt}\}.
 \label{p12-83-c10-s1:eq:evento-firmado-ventana}
\end{equation}
La fórmula distingue el número no negativo de eventos y el signo de la
orientación. No reemplaza una dirección cardinal por un signo: las dos
coordenadas siguen presentes en el estado TPK.

\begin{definition}[Registro enriquecido dodecafásico]
\label{p12-83-c10-s1:def:registro-enriquecido-dodecafasico}
El registro transportado por la traza es
\[
 \mathcal R_{12}(x)=
 \bigl((E_m,\tau_m,\sigma_m,\kappa_m,\Lambda_m)\bigr)_{m=0}^{11},
\]
donde \(\tau_m\) es la orientación TRIT, \(\sigma_m\) la firma de
frontera, \(\kappa_m\) el acarreo, y \(\Lambda_m\) el segmento de registro genealógico
que conserva la ruta. Existe una cadena de proyecciones
\[
 \mathcal X_{\rm TPK}^{\rm enr}
 \longrightarrow\mathcal R_{12}
 \longrightarrow\Theta_{108},
\]
pero ninguna de las dos flechas es una identificación.
\end{definition}

En particular, \(\Theta_{108}\) sólo conserva la posición del calendario.
El registro \(\mathcal R_{12}\) conserva además orientación, firma,
acarreo y procedencia. El estado completo conserva todavía cilindro,
frontera, cociente, terminal y memoria vertical.

\subsection{Carta integral de la oposición dodecafásica}

Para \(0\leq i<6\), las posiciones opuestas determinan
\begin{equation}
 p_i=e_i+e_{i+6},
 \qquad
 a_i=e_i-e_{i+6}.
 \label{p12-83-c10-s1:eq:pares-opuestos-registro}
\end{equation}
La carga central y los cinco márgenes relativos son
\begin{equation}
 s_C=\sum_{i=0}^{5}p_i,
 \qquad
 M_i=p_i-p_5\quad(0\leq i<5).
 \label{p12-83-c10-s1:eq:carga-margenes-registro}
\end{equation}

\begin{definition}[Carta integral \(1+5+6\)]
\label{p12-83-c10-s1:def:carta-integral-156}
Se define
\begin{equation}
 \mathcal L_{12}(e_0,\ldots,e_{11})
 =\bigl(s_C,M_0,\ldots,M_4,a_0,\ldots,a_5\bigr).
 \label{p12-83-c10-s1:eq:carta-integral-156}
\end{equation}
La partición de coordenadas es
\[
 1+5+6=12.
\]
\end{definition}

\begin{theorem}[Inversión exacta de la carta integral]
\label{p12-83-c10-s1:thm:inversion-carta-integral-156}
Una tupla
\[
 \ell=(s_C,M_0,\ldots,M_4,a_0,\ldots,a_5)\in\mathbb Z^{12}
\]
pertenece a la imagen de \(\mathcal L_{12}\) si y sólo si
\begin{align}
 s_C-\sum_{i=0}^{4}M_i&\equiv0\pmod6,
 \label{p12-83-c10-s1:eq:congruencia-seis-carta}\\
 p_i&\equiv a_i\pmod2
 \quad(0\leq i<6),
 \label{p12-83-c10-s1:eq:congruencia-paridad-carta}
\end{align}
donde
\begin{align}
 p_5&=\frac{s_C-\sum_{i=0}^{4}M_i}{6},
 \label{p12-83-c10-s1:eq:inversa-p5-carta}\\
 p_i&=p_5+M_i\quad(0\leq i<5).
 \label{p12-83-c10-s1:eq:inversa-pi-carta}
\end{align}
En ese caso, la única preimagen entera es
\begin{equation}
 e_i=\frac{p_i+a_i}{2},
 \qquad
 e_{i+6}=\frac{p_i-a_i}{2}.
 \label{p12-83-c10-s1:eq:inversa-eventos-carta}
\end{equation}
\end{theorem}

\begin{proof}
Sumando \(M_i=p_i-p_5\) para \(0\leq i<5\), obtenemos
\[
 \sum_{i=0}^{4}M_i
 =\sum_{i=0}^{4}p_i-5p_5
 =s_C-6p_5,
\]
lo que fuerza \eqref{p12-83-c10-s1:eq:inversa-p5-carta} y la congruencia módulo seis.
Una vez reconstruidos los seis \(p_i\), el sistema
\[
 p_i=e_i+e_{i+6},
 \qquad a_i=e_i-e_{i+6}
\]
tiene la única solución \eqref{p12-83-c10-s1:eq:inversa-eventos-carta}. Esa solución es
entera exactamente cuando \(p_i\) y \(a_i\) tienen la misma paridad.
Las dos familias de congruencias son, por tanto, necesarias y suficientes.
\end{proof}

Este teorema es anterior al cierre numérico de \(\alpha\). Explica cómo la
traza orientada produce una carta integral reversible, y por qué no basta
con conservar el calendario reducido.

\subsection{Selección del estado terminal común}

Las tres secciones compatibles llegan a profundidad terminal con catálogos
de cardinalidades
\[
 |\mathcal T_\pi|=196,
 \qquad
 |\mathcal T_e|=197,
 \qquad
 |\mathcal T_\varphi|=196.
\]
El producto contiene
\begin{equation}
 196\cdot197\cdot196=7\,567\,952
 \label{p12-83-c10-s1:eq:cardinal-producto-terminal}
\end{equation}
ternas ordenadas.

El margen de columnas del lector terminal es
\[
 q_\partial=(1,2,2,1,2,0)\in\mathbb F_3^6,
\]
y la orientación de filas del canal conjunto es
\[
 \sigma=(1,1,-1)=(1,1,2)\in\mathbb F_3^3.
\]
El primer dato actúa sobre columnas; el segundo sobre filas. No se deduce
uno del otro.

Las firmas locales de Gram, norma, peso y distancia reducen
\eqref{p12-83-c10-s1:eq:cardinal-producto-terminal} a \(331\) ternas. El margen
\(q_\partial\) deja dos índices:
\[
 (120,197,157),
 \qquad
 (129,197,148).
\]
Sus matrices son
\[
 B_0=
 \begin{pmatrix}
 0&2&0&2&2&0\\
 0&0&1&2&2&0\\
 1&0&1&0&1&0
 \end{pmatrix},
 \qquad
 B_1=
 \begin{pmatrix}
 0&2&1&0&2&0\\
 0&0&1&2&2&0\\
 1&0&0&2&1&0
 \end{pmatrix}.
\]
Las sumas de filas son
\[
 (0,2,0),
 \qquad
 (2,2,1).
\]
Como
\[
 \langle\sigma\rangle
 =\{(0,0,0),(1,1,2),(2,2,1)\},
\]
sólo \(B_1\) satisface la condición axial.

\begin{theorem}[Unicidad de la publicación terminal visible]
\label{p12-83-c10-s1:thm:unicidad-publicacion-terminal-visible}
La selección conjunta determina
\begin{equation}
 B_{\rm term}=
 \begin{pmatrix}
 021020\\001220\\100210
 \end{pmatrix},
 \label{p12-83-c10-s1:eq:matriz-terminal-visible}
\end{equation}
correspondiente a la terna \((129,197,148)\). Ninguna evaluación decimal de
\(\alpha\) interviene en el censo ni en las dos condiciones de selección.
\end{theorem}

\begin{proof}
El censo y el margen dejan exactamente \(B_0,B_1\). La carga de \(B_0\)
no pertenece a \(\langle\sigma\rangle\), mientras la de \(B_1\) es
\(2\sigma\). La segunda candidata es, por tanto, la única que satisface
simultáneamente ambos requisitos tipados.
\end{proof}

\subsection{Publicaciones arquimediana y excepcional del estado común: codominios inequivalentes y genealogía compartida}

Sea
\[
 x_{\rm term}\in
 \mathcal X_{\rm TPK}^{\rm term,enr}
 \subseteq\mathcal X_{\rm TPK}^{[108]}.
\]
La lectura visible es
\[
 \pi_{\rm term}(x_{\rm term})=B_{\rm term}.
\]
La carta firmada conserva otra coordenada del mismo estado.

\begin{definition}[Subespacio firmado terminal]
\label{p12-83-c10-s1:def:subespacio-firmado-terminal}
Sea \(\mathcal H_{12}\) la transformación integral que se construirá en la
sección siguiente. Definimos
\begin{equation}
 \mathcal X_{\rm TPK}^{[108],\rm sgn}
 :=\{(x_{\rm core},K,U):U=\mathcal H_{12}K\}
 \label{p12-83-c10-s1:eq:subespacio-firmado-terminal}
\end{equation}
y el módulo de publicación
\begin{equation}
 \mathcal U_{12}^{\rm sgn}
 :=\operatorname{im}
 \bigl(\mathcal H_{12}:\mathbb Z^{12}\to\mathbb Z^{12}\bigr).
 \label{p12-83-c10-s1:eq:modulo-publicacion-firmada}
\end{equation}
La lectura firmada es
\begin{equation}
 R_{\rm sgn}:
 \mathcal X_{\rm TPK}^{[108],\rm sgn}
 \longrightarrow\mathcal U_{12}^{\rm sgn},
 \qquad
 R_{\rm sgn}(x_{\rm core},K,U)=U.
 \label{p12-83-c10-s1:eq:lector-firmado-terminal}
\end{equation}
\end{definition}

La escritura conjunta correcta es
\begin{equation}
 x_{\rm term}
 \xrightarrow{(\pi_{\rm term},R_{\rm sgn})}
 (B_{\rm term},U_{12}^{\rm sgn}).
 \label{p12-83-c10-s1:eq:doble-publicacion-terminal}
\end{equation}
No se interpone una flecha
\(B_{\rm term}\to U_{12}^{\rm sgn}\): la matriz visible ha olvidado
precisamente las coordenadas de orientación, oposición, cociente, acarreo y
registro genealógico utilizadas por la carta firmada. La fuente de ambas publicaciones es
el estado holonómico integral producido por \eqref{p12-83-c10-s1:eq:cadena-upstream-registro-dodecafasico}.

La evaluación del registro firmado da
\begin{equation}
 \boxed{
 U=(2378,1406,2479,-452,998,-551,
 -668,-204,-371,-322,-28,-997).}
 \label{p12-83-c10-s1:eq:vector-u-firmado-autonomo}
\end{equation}

\subsection{Naturalidad respecto de la profundidad}

Para \(N'\geq N\geq108\), los truncamientos del sistema firmado son
\begin{equation}
 \tau_{N',N}^{\rm sgn}
 (x_{\rm core}^{[N']},K,U)
 =\bigl(\tau_{N',N}x_{\rm core}^{[N']},K,U\bigr).
 \label{p12-83-c10-s1:eq:truncamiento-preserva-ku-autonomo}
\end{equation}

\begin{theorem}[Cuadrado de naturalidad de la lectura firmada]
\label{p12-83-c10-s1:thm:naturalidad-lector-firmado-autonomo}
Se cumple
\begin{equation}
 \boxed{
 R_{{\rm sgn},N}\circ\tau_{N',N}^{\rm sgn}
 =\operatorname{id}_{\mathcal U_{12}^{\rm sgn}}
  \circ R_{{\rm sgn},N'}.}
 \label{p12-83-c10-s1:eq:cuadrado-naturalidad-lector-firmado}
\end{equation}
\end{theorem}

\begin{proof}
Ambos miembros, aplicados a
\((x_{\rm core}^{[N']},K,U)\), devuelven literalmente \(U\). El
truncamiento modifica la historia ulterior pero conserva la coordenada
terminal ya fijada. La conmutatividad es una igualdad de aplicaciones, no
una analogía entre diagramas.
\end{proof}

\subsection{Órbitas de paso \(3\) y transformada de Hadamard}

El orden cíclico del sello se distribuye en las tres órbitas
\begin{align}
 K^{(1)}&=(K_1,K_4,K_7,K_{10}),
 \label{p12-83-c10-s1:eq:orbita-k-1}\\
 K^{(2)}&=(K_2,K_5,K_8,K_{11}),
 \label{p12-83-c10-s1:eq:orbita-k-2}\\
 K^{(3)}&=(K_3,K_6,K_9,K_{12}).
 \label{p12-83-c10-s1:eq:orbita-k-3}
\end{align}
Sea
\begin{equation}
 H_4=
 \begin{pmatrix}
 1&1&1&1\\
 1&1&-1&-1\\
 1&-1&1&-1\\
 1&-1&-1&1
 \end{pmatrix}.
 \label{p12-83-c10-s1:eq:matriz-hadamard-cuatro-autonoma}
\end{equation}
Si \(P_{\rm orb}\) reordena las doce coordenadas de acuerdo con
\eqref{p12-83-c10-s1:eq:orbita-k-1}--\eqref{p12-83-c10-s1:eq:orbita-k-3}, la transformación global es
\begin{equation}
 \mathcal H_{12}
 =P_{\rm orb}^{-1}
 (H_4\oplus H_4\oplus H_4)P_{\rm orb}.
 \label{p12-83-c10-s1:eq:hadamard-global-doce}
\end{equation}

Los tres bloques de \eqref{p12-83-c10-s1:eq:vector-u-firmado-autonomo}, leídos por
órbitas, son
\begin{align*}
 U^{(1)}&=(2378,-452,-668,-322),\\
 U^{(2)}&=(1406,998,-204,-28),\\
 U^{(3)}&=(2479,-551,-371,-997).
\end{align*}

\begin{lemma}[Cuarto de vuelta algebraico de Hadamard]
\label{p12-83-c10-s1:lem:cuarto-inversa-hadamard}
La matriz \(H_4\) satisface
\begin{equation}
 H_4^2=4I_4,
 \qquad
 H_4^{-1}=\frac14H_4,
 \qquad
 \det H_4=-16.
 \label{p12-83-c10-s1:eq:identidades-hadamard-cuatro}
\end{equation}
\end{lemma}

\begin{proof}
Las filas de \(H_4\) son ortogonales y cada una posee norma cuadrada cuatro.
Como la matriz es simétrica, \(H_4^{\mathsf T}H_4=H_4^2=4I_4\). La
fórmula de la inversa y el determinante se deducen inmediatamente.
\end{proof}

\begin{theorem}[Reconstrucción entera única del sello]
\label{p12-83-c10-s1:thm:reconstruccion-entera-k-hadamard}
La inversión
\begin{equation}
 K^{(r)}=\frac14H_4U^{(r)}
 \qquad(r=1,2,3)
 \label{p12-83-c10-s1:eq:inversion-bloques-hadamard}
\end{equation}
produce
\begin{align*}
 K^{(1)}&=(234,729,621,794),\\
 K^{(2)}&=(543,659,058,146),\\
 K^{(3)}&=(140,824,914,601).
\end{align*}
Al deshacer el orden orbital se obtiene
\begin{equation}
 \boxed{
 K=(234,543,140,729,659,824,621,058,914,794,146,601).}
 \label{p12-83-c10-s1:eq:sello-k-doce-autonomo}
\end{equation}
Es la única preimagen racional de \(U\), y pertenece a
\(\mathbb Z^{12}\).
\end{theorem}

\begin{proof}
Por \eqref{p12-83-c10-s1:eq:identidades-hadamard-cuatro}, la preimagen racional es única.
Las multiplicaciones explícitas son
\begin{align*}
 H_4U^{(1)}&=(936,2916,2484,3176),\\
 H_4U^{(2)}&=(2172,2636,232,584),\\
 H_4U^{(3)}&=(560,3296,3656,2404).
\end{align*}
Todas las coordenadas son divisibles por cuatro. Dividir y reintercalar
produce \eqref{p12-83-c10-s1:eq:sello-k-doce-autonomo}.
\end{proof}

El factor \(1/4\) no se escoge para ajustar una salida: es la inversa
forzada por \(H_4^2=4I_4\). Tampoco es la raíz cuadrada de \(-1\). La
estructura compleja interna \(A_W^2=-I\) aparecerá después de construir la
matriz de Paley; son dos hechos de tipos diferentes.

\subsection{Estructura integral y forma normal de Smith}

\begin{theorem}[Imagen integral de \(H_4\)]
\label{p12-83-c10-s1:thm:smith-hadamard-autonomo}
La forma normal de Smith es
\begin{equation}
 \operatorname{SNF}(H_4)=\operatorname{diag}(1,2,2,4).
 \label{p12-83-c10-s1:eq:snf-hadamard-autonoma}
\end{equation}
En consecuencia,
\begin{equation}
 \operatorname{coker}\mathcal H_{12}
 \cong(\mathbb Z/2\mathbb Z)^6
 \oplus(\mathbb Z/4\mathbb Z)^3,
 \label{p12-83-c10-s1:eq:cociente-hadamard-doce}
\end{equation}
y la imagen de \(\mathcal H_{12}\) posee índice
\[
 16^3=4096
\]
en \(\mathbb Z^{12}\). Para \(u\in\mathbb Z^4\),
\begin{equation}
 u\in H_4\mathbb Z^4
 \quad\Longleftrightarrow\quad
 H_4u\equiv0\pmod4.
 \label{p12-83-c10-s1:eq:criterio-imagen-hadamard}
\end{equation}
\end{theorem}

\begin{proof}
El máximo común divisor de las entradas es uno. Los menores de orden dos son
pares y alguno tiene valor absoluto dos; los de orden tres son múltiplos de
cuatro y alguno tiene valor absoluto cuatro; el determinante tiene módulo
dieciséis. Los divisores determinantes son \(1,2,4,16\), de donde se obtienen
los factores invariantes \(1,2,2,4\). La suma directa triple repite los
factores y multiplica los índices. Finalmente,
\(H_4^{-1}=H_4/4\), por lo que la preimagen es entera exactamente bajo la
congruencia \eqref{p12-83-c10-s1:eq:criterio-imagen-hadamard}.
\end{proof}

\subsection{2.º sistema de coordenadas: diferencias de pasos \(3\) y \(4\)}

Con índices cíclicos módulo doce, definimos
\begin{equation}
 (D_3K)_m=K_m-K_{m+3},
 \qquad
 (D_4K)_m=K_m-K_{m+4},
 \qquad
 Q(K)=\sum_{m=1}^{12}K_m.
 \label{p12-83-c10-s1:eq:extractor-diferencias-k}
\end{equation}
Para \eqref{p12-83-c10-s1:eq:sello-k-doce-autonomo},
\begin{align*}
 D_3K={}&(-495,-116,-684,108,601,-90,\\
         &\qquad-173,-88,313,560,-397,461),\\
 D_4K={}&(-425,-281,-481,671,-255,30,\\
         &\qquad475,-543,680,251,6,-128),\\
 Q(K)={}&6263.
\end{align*}

\begin{theorem}[Completitud del extractor transversal]
\label{p12-83-c10-s1:thm:completitud-extractor-transversal-k}
El operador
\[
 \mathcal E_{\rm dod}=(D_3,D_4,Q):
 \mathbb Z^{12}\longrightarrow\mathbb Z^{25}
\]
tiene rango doce y forma normal de Smith
\begin{equation}
 \operatorname{SNF}(\mathcal E_{\rm dod})
 =\operatorname{diag}(\underbrace{1,\ldots,1}_{11},12).
 \label{p12-83-c10-s1:eq:snf-extractor-transversal-k}
\end{equation}
En particular, \((D_3K,D_4K,Q(K))\) determina un único \(K\).
\end{theorem}

\begin{proof}
Si \(D_3v=D_4v=0\), entonces \(v\) es constante a lo largo de los pasos
tres y cuatro. Como \(\gcd(3,4)=1\), esos pasos generan todo
\(\mathbb Z/12\mathbb Z\), de modo que
\[
 \ker D_3\cap\ker D_4=\langle\mathbf1\rangle.
\]
La carga total no se anula en esa recta, pues \(Q(\mathbf1)=12\). El rango
es doce. Las filas de \((D_3,D_4)\) son incidencias orientadas de un grafo
de Cayley conexo; un árbol generador aporta once factores invariantes
unitarios. Todo menor máximo debe contener la fila \(Q\), y las operaciones
unimodulares reducen su última contribución a doce. Existe un menor máximo de
módulo doce; por tanto el último factor es exactamente doce.
\end{proof}

La reconstrucción de \(U\) desde estas coordenadas no necesita volver a
multiplicar por \(H_4\). Para \(r=1,2,3\), sea
\begin{equation}
 B_r=\sum_{j=0}^{3}(D_4K)_{r+3j}.
 \label{p12-83-c10-s1:eq:br-diferencias-k}
\end{equation}
Entonces
\[
 (B_1,B_2,B_3)=(972,-1073,101).
\]
Las sumas orbitales son
\begin{align}
 A_1&=\frac{Q+2B_1+B_2}{3}=2378,
 \label{p12-83-c10-s1:eq:a1-reconstruccion-u}\\
 A_2&=A_1-B_1=1406,
 \label{p12-83-c10-s1:eq:a2-reconstruccion-u}\\
 A_3&=A_2-B_2=2479.
 \label{p12-83-c10-s1:eq:a3-reconstruccion-u}
\end{align}
Si \(d_{r,j}=(D_3K)_{r+3j}\), el bloque firmado de la órbita \(r\) es
\begin{equation}
 \bigl(A_r,
 d_{r,1}-d_{r,3},
 d_{r,0}+d_{r,2},
 d_{r,0}-d_{r,2}\bigr),
 \label{p12-83-c10-s1:eq:bloque-firmado-desde-diferencias}
\end{equation}
que reproduce los tres bloques \(U^{(r)}\).

\subsection{Tétrada de umbral del sello}

La completación decimal local satisface
\[
 1000=729+271.
\]
El soporte de las coordenadas del sello que alcanzan el umbral \(729\) es
\begin{equation}
 \boxed{
 B_K:=\{m:K_m\geq729\}=\{4,6,9,10\}.}
 \label{p12-83-c10-s1:eq:tetrada-umbral-k}
\end{equation}
Esta tétrada se obtiene después de construir \(K\). En esta unidad no se la
denomina todavía estrella de Witt: tal objeto sólo estará definido después
de construir el código ternario, sus hexadas y el sistema de Steiner.

\subsection{Independencia de la construcción genealógica respecto de datos metrológicos o físicos externos}

\begin{theorem}[Aislamiento de \(U\) y \(K\)]
\label{p12-83-c10-s1:thm:aislamiento-u-k-alpha-codata}
La composición
\begin{equation}
 \Omega_{\rm seed}\longrightarrow x_{\rm term}
 \xrightarrow{R_{\rm sgn}}U
 \xrightarrow{\mathcal H_{12}^{-1}}K
 \label{p12-83-c10-s1:eq:composicion-upstream-u-k}
\end{equation}
no recibe como entrada \(\alpha\), \(\alpha^{-1}\), CODATA ni una cadena
decimal objetivo. El sello \(K\) queda determinado antes de definir la
recurrencia que producirá \(\alpha\).
\end{theorem}

\begin{proof}
La primera parte de la composición utiliza únicamente semillas, emisiones,
lectores ternarios, operadores \(L_0,L_1\), relaciones de supervivencia,
transporte de fase y las cartas del estado TPK. La lectura
\(R_{\rm sgn}\) es una proyección del registro terminal, y la última flecha
es la transformación lineal \(\mathcal H_{12}^{-1}=\mathcal H_{12}/4\)
sobre su imagen integral. Ninguna de las fórmulas contiene una variable
física o una tabla metrológica. La recurrencia en base mil se introduce sólo
en la unidad siguiente, después de haber fijado \(K\).
\end{proof}

\subsection{Obstrucciones del registro dodecafásico y del sello \(K\): naturalidad, integridad de Hadamard y aislamiento metrológico}

\begin{criterion}[Refutación del registro y del sello]
La construcción queda refutada si se verifica cualquiera de los hechos
siguientes:
\begin{enumerate}
 \item las ventanas \(E_m\) no particionan los ciento ocho pasos;
 \item la carta \(\mathcal L_{12}\) incumple su inversa o sus congruencias;
 \item el censo terminal no produce \(331\) ternas, dos matrices y una única
       carga axial;
 \item se intenta obtener \(U\) de \(B_{\rm term}\) después de olvidar las
       coordenadas firmadas;
 \item falla el cuadrado de naturalidad
       \eqref{p12-83-c10-s1:eq:cuadrado-naturalidad-lector-firmado};
 \item \(H_4^2\neq4I_4\), alguna preimagen no es entera o el vector
       reconstruido difiere de \eqref{p12-83-c10-s1:eq:sello-k-doce-autonomo};
 \item la forma normal de Smith difiere de
       \eqref{p12-83-c10-s1:eq:snf-hadamard-autonoma};
 \item el extractor \((D_3,D_4,Q)\) pierde rango o no reproduce \(U\);
 \item \(B_K\) se introduce antes de construir \(K\), o se le atribuye una
       estructura de Witt antes de definir \(W_{12}\);
 \item una cifra de \(\alpha\) o una referencia metrológica interviene en
       cualquiera de las flechas de \eqref{p12-83-c10-s1:eq:composicion-upstream-u-k}.
\end{enumerate}
\end{criterion}


% Unidad U017. Cierre afín, prolongación y semántica de frontera de alpha.
% Redacción autónoma desde la autoridad material 1063 y su sucesor V2.

\section{Construcción inicial de la frontera \texorpdfstring{\(662/663\)}{662/663} mediante el cociclo dodecafásico de acarreo}
\label{p12-83-c10-s2:chap:acarreo-alpha-frontera-662-663}

\subsection{Dominio aritmético del cierre}

La construcción recibe cuatro coordenadas dodecafásicas obtenidas en las
unidades precedentes:
\[
\begin{aligned}
 P={}&(141,592,653,589,793,238,462,643,383,279,502,884),\\
 E={}&(718,281,828,459,045,235,360,287,471,352,662,497),\\
 \Phi={}&(618,033,988,749,894,848,204,586,834,365,638,117),\\
 K={}&(234,543,140,729,659,824,621,058,914,794,146,601).
\end{aligned}
\]
Los tres primeros vectores son las doce primeras tríadas de las lecturas
arquimedianas fraccionarias de \(\pi\), \(e\) y \(\varphi\). El cuarto es
la coordenada dodecafásica del estado holonómico terminal del Topological
Phase Kernel, construida antes de abrir la operación de acarreo.

Fijamos la base
\[
 B:=10^3=1000
\]
y la concatenación
\begin{equation}
 \iota(v_1,\ldots,v_{12})
 :=\sum_{m=1}^{12}v_mB^{12-m}.
 \label{p12-83-c10-s2:eq:concatenacion-base-mil-alpha}
\end{equation}
\Needspace{8\baselineskip}
En particular,
\[
\begin{aligned}
 \iota(P)&=141592653589793238462643383279502884,\\
 \iota(E)&=718281828459045235360287471352662497,\\
 \iota(\Phi)&=618033988749894848204586834365638117,\\
 \iota(K)&=234543140729659824621058914794146601.
\end{aligned}
\]

La operación estudiada en esta unidad es el mapa
\begin{equation}
 \mathfrak C_{12}:
 (P,E,\Phi,K,c_{13})
 \longmapsto(A,C),
 \label{p12-83-c10-s2:eq:mapa-cierre-dodecafasico-alpha}
\end{equation}
donde \(A=(A_1,\ldots,A_{12})\) es la palabra de salida y
\[
 C=(c_1,\ldots,c_{13})
\]
es el sistema completo de trece cocientes euclídeos: doce cocientes
internos y una condición de frontera derecha.

\subsection{División euclídea orientada de derecha a izquierda}

Para \(m=12,11,\ldots,1\), se define
\begin{equation}
\begin{aligned}
 s_m&:=P_m+E_m-\Phi_m-K_m+c_{m+1},\\
 A_m&:=s_m\bmod B,
       \qquad 0\leq A_m<B,\\
 c_m&:=\frac{s_m-A_m}{B}
      =\left\lfloor\frac{s_m}{B}\right\rfloor .
\end{aligned}
\label{p12-83-c10-s2:eq:recurrencia-corta-alpha}
\end{equation}
La segunda igualdad para \(c_m\) es válida también cuando \(s_m<0\),
porque se utiliza división euclídea con resto en
\(\{0,\ldots,B-1\}\). Por ello,
\[
 -431=569-1000,
 \qquad
 -717=283-1000,
 \qquad
 -1200=800-2\cdot1000,
\]
y los préstamos negativos no introducen ambigüedad.

\begin{theorem}[Existencia y unicidad del acarreo orientado]
\label{p12-83-c10-s2:thm:unicidad-acarreo-orientado-alpha}
Fijados \(P,E,\Phi,K\) y el acarreo terminal \(c_{13}\), la recurrencia
\eqref{p12-83-c10-s2:eq:recurrencia-corta-alpha} determina de manera única
\[
 (A_1,\ldots,A_{12};c_1,\ldots,c_{12}).
\]
\end{theorem}

\begin{proof}
La fila \(m=12\) queda determinada por la división euclídea de
\[
 P_{12}+E_{12}-\Phi_{12}-K_{12}+c_{13}
\]
entre \(B\). Conocido \(c_{12}\), la fila \(m=11\) queda determinada por
la misma propiedad. La inducción descendente alcanza \(m=1\). En cada
paso, la unicidad del cociente y del resto euclídeos impide dos salidas
distintas.
\end{proof}

\subsection{Ventana dodecafásica del estado firmado y clausura reversible del ciclo}

La ventana finita se define mediante la condición terminal
\[
 c_{13}^{\rm fin}=0.
\]
La evaluación completa es la siguiente. La tabla se imprime en orden
creciente, aunque debe leerse causalmente desde la última fila hacia la
primera.

\begingroup
\scriptsize
\setlength{\tabcolsep}{3.2pt}
\begin{longtable}{@{}r|rrrr|r|r|r|r@{}}
\caption{Las doce divisiones euclídeas y los trece acarreos de la ventana
dodecafásica finita.}
\label{p12-83-c10-s2:tab:acarreo-alpha-finito-autonomo}\\
\toprule
\(m\)&\(P_m\)&\(E_m\)&\(\Phi_m\)&\(K_m\)&
\(c_{m+1}\)&\(s_m\)&\(A_m\)&\(c_m\)\\
\midrule
\endfirsthead
\toprule
\(m\)&\(P_m\)&\(E_m\)&\(\Phi_m\)&\(K_m\)&
\(c_{m+1}\)&\(s_m\)&\(A_m\)&\(c_m\)\\
\midrule
\endhead
1  &141&718&618&234& \(0\)& \(7\)&007& \(0\)\\
2  &592&281&033&543& \(0\)& \(297\)&297& \(0\)\\
3  &653&828&988&140& \(-1\)& \(352\)&352& \(0\)\\
4  &589&459&749&729& \(-1\)& \(-431\)&569& \(-1\)\\
5  &793&045&894&659& \(-2\)& \(-717\)&283& \(-1\)\\
6  &238&235&848&824& \(-1\)& \(-1200\)&800& \(-2\)\\
7  &462&360&204&621& \(0\)& \(-3\)&997& \(-1\)\\
8  &643&287&586&058& \(-1\)& \(285\)&285& \(0\)\\
9  &383&471&834&914& \(-1\)& \(-895\)&105& \(-1\)\\
10 &279&352&365&794& \(0\)& \(-528\)&472& \(-1\)\\
11 &502&662&638&146& \(0\)& \(380\)&380& \(0\)\\
12 &884&497&117&601& \(0\)& \(663\)&663& \(0\)\\
\bottomrule
\end{longtable}
\endgroup

La palabra completa de trece acarreos es
\[
 \boxed{
 C^{\rm fin}
 =(c_1,\ldots,c_{13})
 =(0,0,0,-1,-1,-2,-1,0,-1,-1,0,0,0).}
\]
La salida dodecafásica es
\[
 \boxed{
 A^{\rm fin}
 =(007,297,352,569,283,800,997,285,105,472,380,663).}
\]
Definimos el racional finito
\begin{equation}
 \boxed{
 \alpha_{12}^{\rm fin}
 :=10^{-36}\iota(A^{\rm fin})
 =\frac{7297352569283800997285105472380663}{10^{36}}.}
 \label{p12-83-c10-s2:eq:alpha-finita-autonoma}
\end{equation}
Por tanto,
\[
 \alpha_{12}^{\rm fin}
 =0.007297352569283800997285105472380663.
\]

\subsection{Reconocimiento metrológico posterior: CODATA 2018}

La comparación metrológica se efectúa sólo después de cerrar la recurrencia,
la frontera derecha y el racional anterior. La identidad entera da
\begin{equation}
 \boxed{
 \alpha_{12}^{\rm fin}
 =10^{-36}
 \left\lfloor\frac{10^{36}}{137.035999084}\right\rfloor.}
 \label{p12-83-c10-s2:eq:alpha-finita-codata-2018}
\end{equation}
En consecuencia,
\begin{equation}
 (\alpha_{12}^{\rm fin})^{-1}
 =137.03599908400000000000000000000001066588\ldots,
 \label{p12-83-c10-s2:eq:alpha-inversa-codata-2018}
\end{equation}
y reproduce exactamente la cadena central publicada por CODATA 2018
\cite{Tiesinga2021CODATA}:
\[
 \boxed{137.035999084.}
\]
La palabra \emph{exactamente} se refiere a la cadena central publicada, no
a una supuesta medición experimental de treinta y seis cifras. Las cifras
posteriores de \eqref{p12-83-c10-s2:eq:alpha-inversa-codata-2018} son consecuencias del
racional HMT ya fijado. CODATA no selecciona el estado, no determina (K),
no decide el sentido del acarreo y no interviene en la tabla dodecafásica.

\subsection{Identidad afín completa}

Cada fila de la tabla satisface exactamente
\begin{equation}
 P_m+E_m-\Phi_m-K_m+c_{m+1}
 =A_m+Bc_m.
 \label{p12-83-c10-s2:eq:identidad-local-acarreo-alpha}
\end{equation}
Introduciendo
\[
 C^-:=(c_1,\ldots,c_{12}),
 \qquad
 C^+:=(c_2,\ldots,c_{13}),
\]
las doce identidades locales se reúnen en
\begin{equation}
 \boxed{
 P+E-\Phi-K+C^+-BC^-=A.}
 \label{p12-83-c10-s2:eq:identidad-afin-vectorial-alpha}
\end{equation}
El término
\[
 \partial_BC:=C^+-BC^-
\]
es el cociclo de acarreo de la cadena orientada. En consecuencia,
\(P+E-\Phi-K=A\) no es una igualdad coordenada en
\(\mathbb Z^{12}\): sólo se vuelve correcta al incorporar
\(\partial_BC\).

\begin{theorem}[Telescopado de los trece acarreos]
\label{p12-83-c10-s2:thm:telescopado-acarreos-alpha}
Para cualquier condición terminal \(c_{13}\), la recurrencia satisface
\begin{equation}
 \boxed{
 \iota(P)+\iota(E)-\iota(\Phi)-\iota(K)-\iota(A)
 =B^{12}c_1-c_{13}.}
 \label{p12-83-c10-s2:eq:telescopado-general-alpha}
\end{equation}
En la ventana finita, \(c_1=c_{13}=0\), y por tanto
\begin{equation}
 \boxed{
 \iota(P)+\iota(E)-\iota(\Phi)-\iota(K)
 =\iota(A^{\rm fin}).}
 \label{p12-83-c10-s2:eq:identidad-entera-alpha-finita}
\end{equation}
\end{theorem}

\begin{proof}
Multiplicamos \eqref{p12-83-c10-s2:eq:identidad-local-acarreo-alpha} por
\(B^{12-m}\) y sumamos sobre \(m=1,\ldots,12\). La contribución de los
acarreos es
\[
 \sum_{m=1}^{12}
 (Bc_m-c_{m+1})B^{12-m}.
\]
Las componentes interiores se cancelan:
\begin{align*}
 \sum_{m=1}^{12}(Bc_m-c_{m+1})B^{12-m}
 &=\sum_{m=1}^{12}c_mB^{13-m}
   -\sum_{m=1}^{12}c_{m+1}B^{12-m}\\
 &=B^{12}c_1-c_{13}.
\end{align*}
Esto prueba \eqref{p12-83-c10-s2:eq:telescopado-general-alpha}. La especialización
\(c_1=c_{13}=0\) da \eqref{p12-83-c10-s2:eq:identidad-entera-alpha-finita}.
\end{proof}

\subsection{Inversión exacta del sello \texorpdfstring{\(K\)}{K}}

La recurrencia local es reversible. Despejando \(K_m\) en
\eqref{p12-83-c10-s2:eq:identidad-local-acarreo-alpha},
\begin{equation}
 \boxed{
 K_m=P_m+E_m-\Phi_m+c_{m+1}-A_m-Bc_m.}
 \label{p12-83-c10-s2:eq:inversion-local-k-desde-alpha}
\end{equation}
Por concatenación,
\begin{equation}
 \boxed{
 \iota(K)
 =\iota(P)+\iota(E)-\iota(\Phi)-\iota(A)
  -B^{12}c_1+c_{13}.}
 \label{p12-83-c10-s2:eq:inversion-concatenada-k-desde-alpha}
\end{equation}
En la frontera finita,
\[
 \iota(K)
 =\iota(P)+\iota(E)-\iota(\Phi)-\iota(A^{\rm fin}).
\]

La reversibilidad algebraica no invierte el orden causal. El orden efectivo
es
\[
 x_{\rm term}
 \longrightarrow K
 \longrightarrow(P,E,\Phi,K,c_{13})
 \longrightarrow(A,C).
\]
La fórmula \eqref{p12-83-c10-s2:eq:inversion-local-k-desde-alpha} es un control de
conmutatividad: demuestra que la salida y los acarreos recuperan el sello
que ya estaba fijado. No constituye una regla para seleccionar
retrospectivamente \(K\) a partir de un valor deseado de \(\alpha\).

\subsection{Prolongación coinductiva del sello dodecafásico y recurrencia inversa del acarreo en \(\mathcal C=\{-2,-1,0,1,2\}\)}

Sean
\[
 (P_k)_{k\geq1},
 \qquad
 (E_k)_{k\geq1},
 \qquad
 (\Phi_k)_{k\geq1}
\]
las secuencias infinitas de tríadas producidas por las tres ramas
coinductivas, y sea
\[
 \kappa(k):=1+((k-1)\bmod12).
\]
La prolongación del sello es periódica:
\[
 K_k^{\rm per}:=K_{\kappa(k)}.
\]
Para una profundidad finita \(N\) y una condición terminal
\(t=c_{N+1}\), la recurrencia es
\begin{equation}
\begin{aligned}
 s_k^{(N,t)}
 &=P_k+E_k-\Phi_k-K_{\kappa(k)}+c_{k+1}^{(N,t)},\\
 A_k^{(N,t)}
 &=s_k^{(N,t)}\bmod1000,\\
 c_k^{(N,t)}
 &=\left\lfloor\frac{s_k^{(N,t)}}{1000}\right\rfloor,
 \qquad k=N,\ldots,1.
\end{aligned}
\label{p12-83-c10-s2:eq:recurrencia-larga-alpha}
\end{equation}

El conjunto de acarreos terminales
\[
 \mathcal C:=\{-2,-1,0,1,2\}
\]
es invariante para las entradas certificadas: para cada tríada y cada
\(c_{k+1}\in\mathcal C\), el cociente \(c_k\) vuelve a pertenecer a
\(\mathcal C\). No se elige una condición terminal favorable. Se propagan
las cinco condiciones y sólo se publica el prefijo común a todas ellas.

La ejecución de profundidad certificada emplea
\[
\begin{array}{c|r}
\text{cantidad}&\text{valor}\\ \hline
\text{bloques ternarios por canal}&3501\\
\text{trits por canal}&3501\cdot6=21006\\
\text{cifras decimales de guarda}&10020\\
\text{tríadas de guarda}&3340\\
\text{condiciones terminales propagadas}&5\\
\text{tríadas comunes}&3339\\
\text{cifras comunes}&10017\\
\text{cifras publicadas}&10000
\end{array}
\]
Las cinco ramas terminales coalescen antes del prefijo publicado. En
particular, ninguna condición terminal extraída de un valor objetivo de
\(\alpha\) interviene en la salida.

\begin{theorem}[Compatibilidad de las salidas estabilizadas]
\label{p12-83-c10-s2:thm:compatibilidad-prolongacion-alpha}
Sean \(A^{[N]}\) y \(A^{[N']}\), con \(N'>N\), dos ejecuciones cuyas
cinco condiciones terminales han coalescido antes de la posición \(N\).
Entonces el truncamiento de \(A^{[N']}\) a sus primeros \(N\) bloques
coincide con \(A^{[N]}\). Los prefijos estabilizados definen una única
sección en el límite inverso de las palabras finitas.
\end{theorem}

\begin{proof}
La recurrencia es determinista de derecha a izquierda una vez conocido el
acarreo que entra en una posición. La coalescencia afirma que todas las
condiciones remotas producen el mismo acarreo antes de la frontera del
prefijo. Desde ese punto, las mismas tríadas \(P_k,E_k,\Phi_k\) y el mismo
sello periódico producen, por unicidad euclídea, los mismos restos y
cocientes en todas las posiciones anteriores. Esto es exactamente la
conmutatividad con el truncamiento.
\end{proof}

\subsection{Frontera procedente de la cola remota}

La prolongación determina en la frontera visible
\[
 c_{13}^{\infty}=-1.
\]
Los doce acarreos interiores siguen siendo los mismos que en la ventana
finita:
\[
 \boxed{
 C_{\leq13}^{\infty}
 =(0,0,0,-1,-1,-2,-1,0,-1,-1,0,0,-1).}
\]
Sólo cambia la información que entra desde la decimotercera tríada. Las
dos divisiones relevantes son
\[
\begin{aligned}
 s_{12}^{\infty}
 &=884+497-117-601-1=662,\\
 A_{12}^{\infty}&=662,
 \qquad c_{12}^{\infty}=0,
\end{aligned}
\]
y
\[
\begin{aligned}
 s_{13}^{\infty}
 &=197+757-720-234-1=-1,\\
 A_{13}^{\infty}&=999,
 \qquad c_{13}^{\infty}=-1.
\end{aligned}
\]
Por tanto, el comienzo de la palabra prolongada es
\[
 \boxed{
 A^{\infty}
 =(007,297,352,569,283,800,997,285,105,472,380,662,999,\ldots).}
\]
La correspondiente expansión comienza
\begin{equation}
 \boxed{
 \alpha_{\infty}
 =0.007297352569283800997285105472380662999564172539642\ldots.}
 \label{p12-83-c10-s2:eq:alpha-infinita-autonoma}
\end{equation}

La diferencia entre ambas fronteras es local y exacta:
\[
 c_{13}^{\rm fin}-c_{13}^{\infty}=1,
 \qquad
 A_{12}^{\rm fin}-A_{12}^{\infty}=1,
\]
mientras
\[
 A_m^{\rm fin}=A_m^\infty,
 \qquad 1\leq m\leq11.
\]

Aplicando el teorema de telescopado a los doce primeros bloques de la
prolongación,
\[
 \iota(P)+\iota(E)-\iota(\Phi)-\iota(K)
 -\iota(A^\infty_{\leq12})
 =-c_{13}^{\infty}=1.
\]
Luego
\begin{equation}
 \boxed{
 \iota(A^{\rm fin})
 =\iota(A^\infty_{\leq12})+1.}
 \label{p12-83-c10-s2:eq:relacion-unidad-fronteras-alpha}
\end{equation}
La terminación \(663\) no contradice la terminación \(662\): la primera
es el representante finito obtenido al clausurar la ventana con
\(c_{13}=0\); la segunda conserva el acarreo \(-1\) procedente de la cola
remota.

\begin{figure}[htbp]
\centering
\begin{tikzpicture}[
  >=Stealth,
  node distance=12mm and 16mm,
  box/.style={draw,rounded corners=1.5mm,align=center,
    inner xsep=3.2mm,inner ysep=2.5mm},
  arr/.style={->,thick}
]
\node[box] (shared) {datos compartidos\\
 \(P,E,\Phi,K\)};
\node[box,below left=of shared] (finite)
 {frontera finita\\\(c_{13}=0\)};
\node[box,below right=of shared] (remote)
 {cola remota\\\(c_{13}=-1\)};
\node[box,below=of finite] (a663)
 {\(A_{12}=663\)\\clausura finita};
\node[box,below=of remote] (a662)
 {\(A_{12}=662\), \(A_{13}=999\)\\prefijo prolongado};
\draw[arr] (shared)--(finite);
\draw[arr] (shared)--(remote);
\draw[arr] (finite)--(a663);
\draw[arr] (remote)--(a662);
\draw[<->,densely dashed] (a663)--node[below,align=center]
 {diferencia entera exacta \(1\)} (a662);
\end{tikzpicture}
\caption{Bifurcación de frontera de una misma recurrencia. No se comparan
dos generadores distintos: cambia únicamente el dato que entra desde la
derecha.}
\label{p12-83-c10-s2:fig:bifurcacion-frontera-alpha}
\end{figure}

\subsection{Truncamiento y redondeo a \(36\) cifras}

Para \(x\geq0\), definimos
\[
 \operatorname{Tr}_{36}(x)
 :=10^{-36}\left\lfloor10^{36}x\right\rfloor
\]
y
\[
 \operatorname{Rd}_{36}(x)
 :=10^{-36}\left\lfloor10^{36}x+\frac12\right\rfloor.
\]
Sea
\[
 a:=7297352569283800997285105472380662.
\]
La expansión \eqref{p12-83-c10-s2:eq:alpha-infinita-autonoma} puede escribirse como
\[
 \alpha_\infty=\frac{a+\rho}{10^{36}},
 \qquad
 \rho=0.999564172539642709\ldots,
\]
por lo que \(0.999<\rho<1\).

\begin{theorem}[Resolución exacta de la frontera \(662/663\)]
\label{p12-83-c10-s2:thm:frontera-alpha-662-663}
La prolongación coinductiva satisface
\[
 \boxed{
 \operatorname{Tr}_{36}(\alpha_\infty)
 =0.007297352569283800997285105472380662,}
\]
mientras su redondeo a treinta y seis cifras satisface
\[
 \boxed{
 \operatorname{Rd}_{36}(\alpha_\infty)
 =0.007297352569283800997285105472380663
 =\alpha_{12}^{\rm fin}.}
\]
\end{theorem}

\begin{proof}
Como \(0\leq\rho<1\),
\[
 \left\lfloor10^{36}\alpha_\infty\right\rfloor=a.
\]
Esto da el truncamiento terminado en \(662\). Como \(\rho>1/2\),
\[
 \left\lfloor10^{36}\alpha_\infty+\frac12\right\rfloor=a+1,
\]
que termina en \(663\). Por \eqref{p12-83-c10-s2:eq:alpha-finita-autonoma},
\((a+1)10^{-36}\) es exactamente \(\alpha_{12}^{\rm fin}\).
\end{proof}

Además,
\[
 0<\alpha_{12}^{\rm fin}-\alpha_\infty
 =\frac{1-\rho}{10^{36}}<10^{-39}.
\]
La diferencia no es un error aritmético de la tabla: cuantifica la
diferencia entre dos operadores de frontera distintos.

\subsection{Aislamiento causal respecto de \texorpdfstring{\(\alpha\)}{alpha} y de los datos metrológicos}

Para profundidad \(N\), el dominio completo de la operación es
\[
 \mathscr D_N
 =\bigl((P_k)_{k=1}^{N},(E_k)_{k=1}^{N},
        (\Phi_k)_{k=1}^{N},K,c_{N+1}\bigr).
\]
El mapa de acarreo
\[
 \mathfrak C_N:\mathscr D_N
 \longrightarrow
 \{0,\ldots,999\}^{N}\times\mathcal C^{N}
\]
está definido exclusivamente por \eqref{p12-83-c10-s2:eq:recurrencia-larga-alpha}. Su
dominio no contiene un valor objetivo de \(\alpha\), su inversa, una tabla
metrológica, una cadena decimal objetivo ni un acarreo terminal elegido por
comparación.

\begin{theorem}[No interferencia del blanco metrológico]
\label{p12-83-c10-s2:thm:no-interferencia-metrologica-alpha}
Sea \(\mathscr Y\) cualquier espacio de datos externos que contenga valores
metrológicos o tablas de comparación. La extensión formal
\[
 \widetilde{\mathfrak C}_N:
 \mathscr D_N\times\mathscr Y
 \longrightarrow
 \{0,\ldots,999\}^{N}\times\mathcal C^{N},
 \qquad
 \widetilde{\mathfrak C}_N(d,y):=\mathfrak C_N(d),
\]
factoriza por la proyección
\[
 \operatorname{pr}_{\mathscr D_N}:
 \mathscr D_N\times\mathscr Y\longrightarrow\mathscr D_N.
\]
Por tanto, para cualesquiera \(y,y'\in\mathscr Y\),
\[
 \widetilde{\mathfrak C}_N(d,y)
 =\widetilde{\mathfrak C}_N(d,y').
\]
Ninguna intervención sobre un valor de \(\alpha\) o sobre una tabla
metrológica altera \(K\), los cocientes, los restos ni el prefijo producido.
\end{theorem}

\begin{proof}
La afirmación resulta de
\[
 \widetilde{\mathfrak C}_N
 =\mathfrak C_N\circ\operatorname{pr}_{\mathscr D_N}.
\]
Todas las operaciones de \(\mathfrak C_N\) son sumas, restas y divisiones
euclídeas de los elementos de \(\mathscr D_N\). Ninguna coordenada de
\(\mathscr Y\) aparece en ellas.
\end{proof}

El orden causal completo es
\[
 \text{ramas enriquecidas de }\pi,e,\varphi
 \longrightarrow(P,E,\Phi),
\]
\[
 x_{\rm term}\longrightarrow K,
\]
\[
 (P,E,\Phi,K,\text{frontera})
 \longrightarrow(A,C)
 \longrightarrow\alpha_\infty,
\]
\[
 \alpha_\infty
 \longrightarrow\text{comparación externa}.
\]
La última flecha no puede invertirse para seleccionar ninguna de las
anteriores.

\subsection{Obstrucciones de la prolongación afín de \(\alpha\): telescopía, inversión de \(K\), frontera \(662/663\) y aislamiento de datos objetivo}

\begin{criterion}[Refutación del cierre afín y de su prolongación]
El bloque queda refutado en su dominio si ocurre cualquiera de los hechos
siguientes:
\begin{enumerate}
 \item alguna fila de la \cref{p12-83-c10-s2:tab:acarreo-alpha-finito-autonomo}
       incumple \(s_m=A_m+1000c_m\);
 \item el vector publicado de trece acarreos no reproduce las doce filas;
 \item falla la identidad telescópica
       \[
       \iota(P)+\iota(E)-\iota(\Phi)-\iota(K)-\iota(A)
       =1000^{12}c_1-c_{13};
       \]
 \item la inversión \eqref{p12-83-c10-s2:eq:inversion-local-k-desde-alpha} no recupera
       las doce coordenadas de \(K\);
 \item la propagación remota produce \(c_{13}\neq-1\) para el prefijo
       certificado;
 \item las cinco condiciones terminales de
       \(\mathcal C=\{-2,-1,0,1,2\}\) no coalescen antes de las cifras
       publicadas;
 \item la tríada siguiente a \(662\) no es \(999\), en cuyo caso deja de
       estar demostrada la identificación del redondeo con \(663\);
 \item el generador abre una expansión objetivo de \(\alpha\), su inversa,
       una tabla metrológica o una condición terminal elegida para forzar
       la coincidencia.
\end{enumerate}
\end{criterion}

Quedan así demostradas dos identidades de frontera distintas:
\[
 \boxed{
 663=\text{clausura finita con }c_{13}=0
     =\text{redondeo a 36 cifras de }\alpha_\infty,}
\]
\[
 \boxed{
 662=\text{truncamiento a 36 cifras de }\alpha_\infty
 \quad\text{con}\quad c_{13}=-1,\ A_{13}=999.}
\]
El resultado metrológico central queda, finalmente, publicado en la misma
residencia que la derivación:
\begin{equation}
 \boxed{
 \begin{gathered}
  (\alpha_{12}^{\rm fin})^{-1}
  =137.03599908400000000000000000000001066588\ldots,\\
  \text{cadena central CODATA 2018:}\qquad137.035999084.
 \end{gathered}}
 \label{p12-83-c10-s2:eq:resultado-final-alpha-codata-2018}
\end{equation}
La implicación expresa comparación posterior con la cadena central
publicada; no introduce esa cadena en el constructor.
 % Insercion editor-ready para el capitulo 31.

\section{Prolongación nonádica graduada del núcleo \texorpdfstring{\(E_{21/22}\)}{E21/22}: operador \texorpdfstring{\(T_{7:9}\)}{T7:9}, genealogía \texorpdfstring{\(120/45/11/495\)}{120/45/11/495} y residuo de orden \texorpdfstring{\(9\)}{9}}
\label{sec:l9s102:alpha-dos-vias}

Esta sección desarrolla dos construcciones internas relativas a la
coordenada \(\alpha\) de la
\hyperref[sec:tesis-inaugural-helicoidal-hmt-md]{tesis inaugural}. La vía A
opera mediante la clausura dodecafásica reversible del estado firmado y
publica \(\alpha_{12}\). La vía B, una vez construidos internamente el
núcleo analítico \(E_{21/22}\) y sus coeficientes, publica en su carta
analítica una raíz positiva simple y única, que denotamos
\(\alpha_E:=x_{21/22}\). Ambas salidas comparten genealogía; su concordancia
demostrada pertenece al codominio decimal tipado que se especifica abajo.
No se identifica por definición \(\alpha_E\) con \(\alpha_{12}\). Sus dominios,
operaciones y certificados se mantienen diferenciados:
\begin{equation}
 X_{12}^{\mathrm{sig}}
 \xrightarrow{\ \Gamma_{12}^{\alpha}\ }
 \alpha_{12},
 \qquad
 E_{21/22}:I_\alpha\longrightarrow\mathbb R,
 \quad E_{21/22}(x_{21/22})=0,
 \quad E_{21/22}'(x_{21/22})\ne0.
 \label{eq:l9s102:alpha-dos-mapas}
\end{equation}
La clausura entera publica \(\alpha_{12}\) mediante acarreo, mientras que la
ecuación analítica determina \(\alpha_E=x_{21/22}\) por existencia y unicidad
de la raíz en el intervalo certificado. La concordancia actualmente
demostrada de ambas salidas reside exactamente en el codominio del morfismo
\begin{equation}
 \operatorname{Tr}_9(y):=10^{-9}\lfloor10^9y\rfloor,
 \qquad
 \operatorname{Tr}_9(x_{21/22}^{-1})
 =137.035999084
 =\operatorname{Tr}_9(\alpha_{12}^{-1}).
 \label{eq:l9s102:concordancia-tipica-alpha}
\end{equation}
La identidad analítica exacta corresponde a
\(E_{21/22}(\alpha_E)=0\), mientras
\eqref{eq:l9s102:concordancia-tipica-alpha} es la igualdad exacta de las dos
publicaciones tras aplicar \(\operatorname{Tr}_9\). La promoción más fuerte
\(\alpha_E=\alpha_{12}\) en \(\mathbb R_{\rm HMT}\) requiere demostrar
\(E_{21/22}(\alpha_{12})=0\), o bien un morfismo inyectivo que levante la
igualdad truncada; el residuo no nulo de la
\cref{tab:l9s102:alpha-residuos} no suministra esa identidad. Ninguna vía
introduce el valor objetivo como dato de entrada; la identificación
metrológica pertenece únicamente al reconocimiento convencional posterior.

\subsection{Núcleo analítico \texorpdfstring{\(P_6\)}{P6} y raíz única de la ecuación \texorpdfstring{\(E_{21/22}\)}{E21/22}}
\label{sec:l9s102:alpha-p6}

La caracterización de vacancias emplea
\begin{equation}
 \begin{aligned}
 P_6(x)={}&2\pi x-\frac74x^2
 +\frac{x^3}{2\pi}
 +\frac{x^4}{20}-\frac{2x^5}{21}-\frac{x^6}{46},\\
 P_6(x)={}&\log_{10}\frac{10}{9}.
 \end{aligned}
 \label{eq:l9s102:p6}
\end{equation}
La raíz positiva certificada de esta ecuación satisface
\begin{equation}
 x_6^{-1}=137.03599908400002702009\ldots .
 \label{eq:l9s102:raíz-p6}
\end{equation}
El polinomio \(P_6\) proporciona el jet inicial de la caracterización
analítica; la holonomía nonádica prolonga ahora sus grados \(7,8,9\).

\subsection{Genealogía APP--TRIT--TPK de los invariantes \texorpdfstring{\(120,45,11\)}{120, 45, 11} y \texorpdfstring{\(495\)}{495}}
\label{sec:l9s102:alpha-genealogia-7-9}

La semicorona de fase del TPK construye
\begin{equation}
 |\mathcal F_8^{\mathrm{ph}}|=120,
 \qquad |L_{\mathrm{ph}}|=45.
 \label{eq:l9s102:alpha-120-45}
\end{equation}
APP aporta el bloque central
\begin{equation}
 B_c=\begin{pmatrix}7&2\\2&7\end{pmatrix},
 \qquad
 \det B_c=45,
 \qquad
 \operatorname{Spec}(B_c)=\{9,5\}.
 \label{eq:l9s102:alpha-bc}
\end{equation}
El extractor transversal dodecafásico proporciona
\begin{equation}
 r_{\mathrm{dod}}=\operatorname{rank}(D_3,D_4)=11,
 \qquad
 (D_3K)_1=-495,
 \qquad
 495=45\cdot11=\det(B_c)\,r_{\mathrm{dod}}.
 \label{eq:l9s102:alpha-495}
\end{equation}
Los coeficientes del jet prolongado proceden, por tanto, de tres
invariantes previos de APP--TRIT--TPK:
\begin{equation}
 -\frac1{120}\longmapsto\frac1{45}
 \longmapsto\frac2{45\cdot11}.
 \label{eq:l9s102:alpha-coeficientes}
\end{equation}

\subsection{Resolvente nilpotente \texorpdfstring{\((I-N)^{-1}\)}{(I-N)^-1} y construcción exacta de la cola \texorpdfstring{\(T_{7:9}\)}{T7:9}}
\label{sec:l9s102:alpha-resolvente}

Sobre la base \(e_7,e_8,e_9\), definimos
\begin{equation}
 Ne_7=-\frac83e_8,
 \qquad
 Ne_8=\frac2{11}e_9,
 \qquad
 Ne_9=0,
 \qquad
 v_7=-\frac1{120}e_7.
 \label{eq:l9s102:alpha-nilpotente}
\end{equation}
Como \(N^3=0\), su resolvente finito actua exactamente por
\begin{equation}
 (I-N)^{-1}v_7
 =v_7+Nv_7+N^2v_7
 =-\frac{e_7}{120}+\frac{e_8}{45}+\frac{2e_9}{495}.
 \label{eq:l9s102:alpha-resolvente-finito}
\end{equation}
La evaluación \(e_n\mapsto x^n\) produce
\begin{equation}
 \boxed{
 T_{7:9}(x)=-\frac{x^7}{120}+\frac{x^8}{45}
 +\frac{2x^9}{495}.}
 \label{eq:l9s102:alpha-t79}
\end{equation}

\subsection{Jet nonádico \texorpdfstring{\(P_9\)}{P9}, raíz prolongada y escala de residuos certificados}
\label{sec:l9s102:alpha-jet-p9}

El jet nonádico de orden \(9\) queda definido por
\begin{equation}
 P_9(x)=P_6(x)+T_{7:9}(x),
 \qquad
 P_9(x)=\log_{10}\frac{10}{9}.
 \label{eq:l9s102:alpha-p9}
\end{equation}
Su raíz positiva satisface
\begin{equation}
 x_9^{-1}
 =137.0359990840000000000000111926291900\ldots .
 \label{eq:l9s102:alpha-raíz-p9}
\end{equation}

\begin{table}[htbp]
\centering
\caption{Descenso certificado del residuo de la caracterización analítica}
\label{tab:l9s102:alpha-residuos}
\begin{tabular}{@{}lc@{}}
\toprule
Núcleo evaluado en \(\alpha_{12}\) & Residuo \\
\midrule
\(P_6\) & \(9.0038886\times10^{-18}\) \\
\(P_6-x^7/120\) & \(-1.7893115\times10^{-19}\) \\
\(P_6-x^7/120+x^8/45\) & \(-2.3708601\times10^{-22}\) \\
\(P_9\) & \(3.7297132\times10^{-27}\) \\
\bottomrule
\end{tabular}
\end{table}

\subsection{Operador dodecafásico alternativo \texorpdfstring{\(A_{\mathrm{dod}}\)}{A_dod} y demostración de la no unicidad de la selección}
\label{sec:l9s102:alpha-control}

El operador de control
\begin{equation}
 A_{\mathrm{dod}}
 =I-\frac14(D_3^*D_3+D_4^*D_4)
 \label{eq:l9s102:alpha-control-alternativo}
\end{equation}
produce una sucesión diferente de coeficientes. Este control delimita la
fuerza selectiva del resolvente
\cref{eq:l9s102:alpha-resolvente-finito,eq:l9s102:alpha-t79}: un operador
dodecafásico genérico no determina por sí solo la prolongación \(T_{7:9}\).
No rebaja la vía de vacancias, cuya condición probatoria es la existencia y
unicidad de la raíz del núcleo declarado. El morfismo de control
\(\Gamma_{E21}^{(9)}\) se reserva exclusivamente para la unicidad canónica de
la prolongación nonádica \(7{:}9\); no certifica ni condiciona el estatuto de
la vía B.

\subsection{Alcance del operador nilpotente \texorpdfstring{\(N\)}{N} sobre el jet nonádico \texorpdfstring{\(P_9\)}{P9} y control de unicidad de la prolongación \texorpdfstring{\(7{:}9\)}{7:9}}
\label{sec:l9s102:alpha-alcance}

La aritmética de \(120,45,11,495\), el operador nilpotente, el resolvente y
la identidad \cref{eq:l9s102:alpha-t79} poseen exactitud interna. La raíz
positiva simple y única del núcleo \(E_{21/22}\) constituye una segunda
construcción interna exacta, \(\alpha_E\), para el dominio y los coeficientes
declarados. Su identificación con la salida dodecafásica está demostrada
exactamente después de \(\operatorname{Tr}_9\); elevarla a igualdad en
\(\mathbb R_{\rm HMT}\) exige la identidad o el levantamiento inyectivo
enunciados al comienzo de la sección. De manera independiente, la unicidad
canónica de la prolongación \(7{:}9\) queda tipada por el morfismo
\begin{equation}
 \Gamma_{E21}^{(9)}:
 X_{\mathrm{APP\text{-}TRIT\text{-}TPK}}
 \longrightarrow \mathbb Q(\pi)[x]/(x^{10}),
 \label{eq:l9s102:alpha-gamma-e21}
\end{equation}
cuya acción determina conjuntamente la semilla orientada, las dos
transferencias, la graduación \(7,8,9\) y el cociente de jets
\(F^7/F^{10}\). El falsador separa, por tanto, dos afirmaciones: invalidaría
la vía B si fallasen la monotonía, el cambio de signo o el aislamiento de la
raíz; invalidaría la promoción canónica de \(T_{7:9}\) si otro lector del
mismo dominio produjese la misma graduación sin factorizar por
\(\Gamma_{E21}^{(9)}\).
 
% Selección íntegra de los cuerpos matemáticos de la Parte III rectora. Las
% cartas dimensionales se reservan para la parte siguiente, de modo que una
% constante física no pueda retroactuar como generador de estos caracteres.
% Punto de montaje redefinible por el compilador único.
% La ruta por defecto es válida cuando el módulo se consume desde
% ``manuscrito/``; el editor único puede redefinirla antes del \input.
\providecommand{\HMTParteIIIRoot}{../colaboracion/parte_iii}
\providecommand{\HMTParteIIISource}{\HMTParteIIIRoot/source/public_final}
\providecommand{\HMTParteIIICapitulos}{\HMTParteIIIRoot/tex/capitulos_finales}
\providecommand{\APP}{\mathrm{APP}}
\providecommand{\TRIT}{\mathrm{TRIT}}
\providecommand{\TPK}{\mathrm{TPK}}
\providecommand{\etaRet}{\eta_{\mathrm{ret}}^{(\mathrm{deg})}}
\providecommand{\etaRetrad}{\eta_{\mathrm{ret}}^{(\mathrm{rad})}}
\providecommand{\alphaAngle}{\alpha_{\angle}}
 
% \newcommand{\HMTMonographChapterInput}[2]{%
%   \chapter{#1}%
%   \begingroup
%   \renewcommand{\chapter}[2][]{}%
%   \input{#2}%
%   \endgroup}

\chapter{Clasificación algebraica y genealógica de las constantes fundamentales}
\begingroup
\renewcommand{\chapter}[2][]{}%
\chapter{Álgebra y clasificación genealógica de las constantes fundamentales}
\label{chap:p3-final:27:algebra_clasificacion}
\label{chap:p3-83-23-algebra_clasificacion}


\label{int:chap:producto-nativo}
\index{producto nativo}

La extensión euclídea residuo--cociente de APP contiene dos aplicaciones
locales: la aplicación aditiva normaliza colisiones y la multiplicativa
produce productos parciales. En esta sección se construye la multiplicación
global de palabras usando exclusivamente ambas aplicaciones. La correspondencia con la multiplicación entera será una
consecuencia del entrelazador, no la definición de la operación.

\section{Coordenadas euclídeas de la célula local}

\begin{definition}[Aplicaciones aditiva y multiplicativa sobre un mismo soporte]
Sea $\Dnine=\{1,\ldots,9\}$. Para $a,b\in\Dnine$ se definen por división
euclídea con resto en $\Dnine$:
\begin{align}
 a+b&=9q^+(a,b)+R^+(a,b),\label{int:eq:local-plus}\\
 ab&=9q^\times(a,b)+R^\times(a,b).\label{int:eq:local-times}
\end{align}
En particular, el resto nulo se representa por $9$ y el cociente se reduce en
una unidad. La cuádrupla
\[
 \CellRQ(a,b)=
 (R^+(a,b),q^+(a,b);R^\times(a,b),q^\times(a,b))
\]
es la extensión aritmética local por residuo y cociente.
\end{definition}

La agregación por las tres clases módulo tres de la tabla visible
\(R^\times\) define una matriz distinta de las aplicaciones locales:
\[
(M_\Pi)_{rs}
=\frac19
\sum_{\substack{a\equiv r\;(3)\\b\equiv s\;(3)}}R^\times(a,b),
\qquad r,s\in\{1,2,0\},
\]
y, en ese orden de clases,
\[
M_\Pi=
\begin{pmatrix}4&5&6\\5&4&6\\6&6&9\end{pmatrix}.
\]
Así, \(q^+\) y \(q^\times\) son coordenadas de cociente de una celda,
mientras que \(M_\Pi\) es un operador agregado sobre clases residuales.

Las $81$ celdas satisfacen exactamente
\[
\begin{array}{c|rrrr}
 &\sum R&\sum q&\sum(9q+R)\\\hline
 +&405&45&810\\
 \times&459&174&2025
\end{array}
\]
y por tanto
\[
 2025-810=(459-405)+9(174-45)=54+9\cdot129=1215.
\]
El cociente $q^\times$ es necesario para la reconstrucción. Por ejemplo,
\[
 2\cdot5=9\cdot1+1.
\]
Si se conserva sólo $R^\times$, el producto $10$ se confunde con el valor
visible $1$.

\subsection{Cociclos de asociatividad y distributividad}

La coherencia del transductor puede decidirse ya en la célula. Para todo
$a,b,c\in\Dnine$ se verifican las tres identidades
\begin{align}
q^+(a,b)+q^+(R^+(a,b),c)
&=q^+(b,c)+q^+(a,R^+(b,c)),\label{int:eq:cocycle-plus}\\
q^\times(R^\times(a,b),c)+cq^\times(a,b)
&=q^\times(a,R^\times(b,c))+aq^\times(b,c),
\label{int:eq:cocycle-times}\\
q^\times(a,R^+(b,c))+aq^+(b,c)
&=q^+(R^\times(a,b),R^\times(a,c))\notag\\
&\quad+q^\times(a,b)+q^\times(a,c).
\label{int:eq:cocycle-distributive}
\end{align}

\begin{theorem}[Coherencia local de la extensión euclídea]
Las identidades \eqref{int:eq:cocycle-plus}--\eqref{int:eq:cocycle-distributive},
junto con las igualdades correspondientes de residuos, expresan
respectivamente asociatividad aditiva, asociatividad multiplicativa y
distributividad en la extensión residuo--cociente. Las tres familias tienen
cero discrepancias en los $9^3=729$ triples.
\end{theorem}

\begin{proof}
Al desarrollar $(a+b)+c$ mediante \eqref{int:eq:local-plus}, el coeficiente de
nueve es el lado izquierdo de \eqref{int:eq:cocycle-plus}; al desarrollar
$a+(b+c)$ se obtiene el derecho. La unicidad de la división con resto en
$\Dnine$ iguala tanto residuos como cocientes. El mismo argumento aplicado a
$(ab)c=a(bc)$ da \eqref{int:eq:cocycle-times}. Finalmente, las dos expansiones de
$a(b+c)=ab+ac$ producen \eqref{int:eq:cocycle-distributive}. La enumeración
directa de los \(729\) casos de cada identidad confirma el cálculo
algebraico.
\end{proof}

\section{Biyectividad de la numeración}

El símbolo $9$ no se identifica con el cero. En $\CellRQ$, $9$ representa un
cruce completo de la frontera nonádica. La sustitución por los dígitos
$0,\ldots,8$ eliminaría el cociente euclídeo que interviene en el
entrelazamiento.

\begin{definition}[Palabras canónicas]
Sea
\[
 \Wnine^+=\bigsqcup_{\ell\geq1}\Dnine^\ell,
 \qquad \Wnine=\{\varnothing\}\sqcup\Wnine^+.
\]
Las palabras se escriben desde el dígito menos significativo. Definimos
\[
 \operatorname{ev}_9(u_0,\ldots,u_{\ell-1})
 =\sum_{j=0}^{\ell-1}u_j9^j,
 \qquad \operatorname{ev}_9(\varnothing)=0.
\]
\end{definition}

\begin{theorem}[Unicidad de la forma normal]
La evaluación $\operatorname{ev}_9$ es una biyección de $\Wnine^+$ sobre
$\N_{\geq1}$, y de $\Wnine$ sobre $\N_0$.
\end{theorem}

\begin{proof}
Las palabras de longitud $\ell$ cubren el intervalo
\[
 \frac{9^\ell-1}{8}
 \leq \operatorname{ev}_9(u)
 \leq 9\frac{9^\ell-1}{8}.
\]
El mínimo del intervalo siguiente es una unidad mayor que el máximo del
anterior. Dentro de cada intervalo, la división de $n-1$ por nueve determina
de manera única el dígito
\[
 d=(n-1)\bmod9+1
\]
y el cociente que se codifica recursivamente. El proceso termina porque el
cociente decrece estrictamente.
\end{proof}

Ejemplos:
\[
 9\longleftrightarrow(9),\qquad
 10\longleftrightarrow(1,1),\qquad
 81\longleftrightarrow(9,8),\qquad
 1000\longleftrightarrow(1,3,3,1).
\]

\section{Fibras de las proyecciones residuales en los levantamientos enteros}

Para este bloque distinguimos dos alfabetos. La forma normal
\(\Wnine\) usa los dígitos \(\{1,\ldots,9\}\), escritos desde el menos
significativo. En cambio, \(\mathscr B_9\) designa las escrituras posicionales
convencionales en base nueve, con dígitos \(\{0,\ldots,8\}\) y el más
significativo a la izquierda. No se identifican ambos tipos.

Sea
\[
 \mathscr W_{10}^{(2)}=\{0,\ldots,9\}^2,\qquad
 \rho_{10}(a,b)=(b,a),\qquad
 \operatorname{ev}_{10}(a,b)=10a+b,
\]
y sea \(q_9:\mathbb N\to\mathbb Z/9\mathbb Z\) la reducción residual.

\begin{lemma}[Descenso residual de la reversión decimal]
\label{int:pf84:SYN30-RHO-DESC}
La reversión sobre palabras decimales de longitud dos es involutiva y su
sombra módulo nueve es la identidad:
\[
 \rho_{10}^2=I,\qquad
 q_9\!\left(\operatorname{ev}_{10}(\rho_{10}(a,b))\right)
 =
 q_9\!\left(\operatorname{ev}_{10}(a,b)\right).
\]
La longitud fija conserva los ceros iniciales.
\end{lemma}

\begin{proof}
La doble permutación recupera \((a,b)\). Como \(10\equiv1\pmod9\),
\[
 10b+a\equiv a+b\equiv10a+b\pmod9.
\]
Una palabra cuya doble reversión no fuese ella misma o cuya suma de cifras
cambiase módulo nueve refutaría el lema.
\end{proof}

\begin{lemma}[Descenso residual del cuadrado]
\label{int:pf84:SYN30-SQ-DESC}
La aplicación \(s(n)=n^2\) induce el endomapa residual
\[
 [n]_9\longmapsto[n]_9^2,
 \qquad q_9(s(n))=q_9(n)^2.
\]
\end{lemma}

\begin{proof}
La compatibilidad de una congruencia con el producto da
\(n\equiv m\pmod9\Rightarrow n^2\equiv m^2\pmod9\). Un entero para el que
\(q_9(n^2)\ne q_9(n)^2\) refutaría el enunciado.
\end{proof}

\begin{proposition}[Fibras residuales de la reversión entera]
\label{int:pf84:SYN30-RHO-NONFAITH}
La sombra residual no reconstruye el levantamiento entero
\(\widehat\rho_{10}=\operatorname{ev}_{10}\circ\rho_{10}\). En efecto,
\[
 18\equiv27\pmod9,\qquad
 \widehat\rho_{10}(1,8)=81\ne72=\widehat\rho_{10}(2,7).
\]
\end{proposition}

\begin{proof}
Ambos enteros pertenecen a la clase residual cero, pero sus reversiones son
distintas. Por tanto no existe
\(\overline\rho:\mathbb Z/9\mathbb Z\to\mathbb N\) que factorice
simultáneamente esos dos valores.
\end{proof}

\begin{proposition}[Fibras residuales del cuadrado entero]
\label{int:pf84:SYN30-SQ-NONFAITH}
La misma sombra tampoco reconstruye el cuadrado entero:
\[
 18\equiv27\pmod9,\qquad
 18^2=324\ne729=27^2.
\]
\end{proposition}

\begin{proof}
Una función \(\overline s:\mathbb Z/9\mathbb Z\to\mathbb N\) que
factorizase \(s\) debería asignar a la misma clase cero los dos valores
\(324\) y \(729\), lo cual es imposible.
\end{proof}

\begin{corollary}[Testigo cruzado \(27\)--\(72\)--\(729\)]
\label{int:pf84:SYN30-CROSS-WITNESS}
\[
 \widehat\rho_{10}(2,7)=72,\qquad
 27^2=729=9^3=3^6,
\]
y, en el alfabeto convencional \(\mathscr B_9\),
\[
 27=30_9,\qquad72=80_9,\qquad729=1000_9.
\]
Los tres enteros tienen sombra residual cero y raíz digital positiva nueve.
\end{corollary}

\begin{proof}
Las dos operaciones se calculan directamente; las divisiones sucesivas por
nueve producen las tres escrituras. El resultado fallaría si se confundiese
\(\mathscr B_9\) con \(\Wnine\), si una de las escrituras fuese incorrecta
o si se afirmase que el descenso residual fiel reconstruye la operación y su
genealogía completa.
\end{proof}

\section{Sumador nativo y propagación del cociente}

\begin{algorithm}[Suma de palabras]
Se recorren dos palabras por posiciones, con acarreo inicial \(c_0=0\). Si en
una posición hay dos dígitos, se aplica \eqref{int:eq:local-plus}; si sólo hay uno,
ese dígito se toma como residuo provisional y el cociente provisional es cero.
Cuando \(c_j\ne0\), se aplica una segunda celda aditiva al residuo provisional
y a \(c_j\). Se publica el nuevo residuo y se transporta la suma de los
cocientes a la posición siguiente. Si al terminar queda un acarreo no nulo, se
publica como último dígito. De este modo las celdas locales nunca reciben el
símbolo ausente \(0\), que no pertenece a \(\Dnine\).
\end{algorithm}

\begin{lemma}[Invariante del sumador]
Si $u\boxplus_\#v$ es la palabra producida por el algoritmo anterior,
\[
 \operatorname{ev}_9(u\boxplus_\#v)
 =\operatorname{ev}_9(u)+\operatorname{ev}_9(v).
\]
El acarreo intermedio está en $\{0,1,2\}$.
\end{lemma}

\begin{proof}
Tras procesar las posiciones $0,\ldots,k-1$, sea $c_k$ el cociente pendiente.
Las identidades locales implican
\[
 \sum_{j<k}(u_j+v_j)9^j
 =\sum_{j<k}w_j9^j+c_k9^k,
\]
incluyendo como cero los dígitos ausentes. Cada nueva celda conserva esta
igualdad. El máximo local es $9+9+2=20$, de modo que $c_{k+1}\leq2$. Al
terminar se publica el último cociente, y queda la identidad global.
\end{proof}

\section{Multiplicación por un dígito}

\begin{algorithm}[Transductor escalar]
Fijado $d\in\Dnine$, cada dígito $u_j$ produce
\[
 u_jd=9q_j+r_j,
 \qquad
 (r_j,q_j)=\bigl(R^\times(u_j,d),q^\times(u_j,d)\bigr).
\]
Sea $c_j\in\{0,\ldots,9\}$ el cociente entrante. Como $0\notin\Dnine$,
la transición se define por ramas:
\[
 (w_j,c_{j+1})=
 \begin{cases}
 (r_j,q_j),&c_j=0,\\[1mm]
 \bigl(R^+(r_j,c_j),q_j+q^+(r_j,c_j)\bigr),&c_j\in\Dnine.
 \end{cases}
\]
Así, ninguna aplicación local se evalúa fuera de su dominio. La primera
rama expresa la ausencia de cociente pendiente; la segunda normaliza la
colisión mediante la aplicación aditiva.
\end{algorithm}

\begin{lemma}[Invariante escalar]
La palabra $d\odot_\#u$ producida por el transductor satisface
\[
 \operatorname{ev}_9(d\odot_\#u)
 =d\operatorname{ev}_9(u).
\]
El cociente multiplicativo pertenece a $\{0,\ldots,9\}$.
\end{lemma}

\begin{proof}
El argumento es el del sumador, sustituyendo cada sumando local por
$u_jd$. Si $c_j=0$, la identidad
$u_jd=9q_j+r_j$ conserva directamente la igualdad parcial. Si
$c_j\in\Dnine$, la segunda división euclídea da
\[
 r_j+c_j=9q^+(r_j,c_j)+R^+(r_j,c_j),
\]
y conserva la misma igualdad después de transportar
$q_j+q^+(r_j,c_j)$. El máximo $9\cdot9+9=90$ y la forma de resto en
$\Dnine$ mantienen el cociente en la cota declarada. El último cociente se
publica como dígito.
\end{proof}

\section{Producto global por Horner}

\begin{definition}[Producto nativo]
Sea $v=(v_0,\ldots,v_{m-1})$. Partiendo de $a_m=\varnothing$, se define para
$j=m-1,\ldots,0$
\[
 a_j=(9\odot_\#a_{j+1})\boxplus_\#(v_j\odot_\#u).
\]
El resultado es
\[
 u\boxtimes_\#v:=a_0.
\]
\end{definition}

Esta definición no contiene
\[
 \operatorname{dec}_9(\operatorname{ev}_9(u)\operatorname{ev}_9(v)).
\]
La evaluación sólo aparecerá en la demostración del entrelazamiento.

\begin{theorem}[Multiplicatividad nativa]
Para todas las palabras $u,v\in\Wnine$,
\begin{equation}
 \label{int:eq:ev-multiplicative}
 \boxed{
 \operatorname{ev}_9(u\boxtimes_\#v)
 =\operatorname{ev}_9(u)\operatorname{ev}_9(v).}
\end{equation}
La palabra vacía es absorbente y $(1)$ es la unidad.
\end{theorem}

\begin{proof}
Por los dos lemas anteriores,
\[
 \operatorname{ev}_9(a_j)
 =9\operatorname{ev}_9(a_{j+1})
  +v_j\operatorname{ev}_9(u).
\]
La iteración de Horner da
\[
 \operatorname{ev}_9(a_0)
 =\operatorname{ev}_9(u)
  \sum_{j=0}^{m-1}v_j9^j.
\]
El último factor es $\operatorname{ev}_9(v)$. La unicidad de la forma normal
identifica la palabra resultante.
\end{proof}

\begin{corollary}[Monoide multiplicativo de palabras]
La estructura
\[
\WordMonoid=(\Wnine,\boxtimes_\#,(1))
\]
es un monoide conmutativo; la palabra vacía es su elemento absorbente. Su
construcción no utiliza la multiplicación global de enteros.
\end{corollary}

El ejemplo de frontera más corto es
\[
 (9)\boxtimes_\#(9)=(9,8),
 \qquad 9+8\cdot9=81.
\]
Las \(17\) celdas locales de la multiplicación producen
\[
 (9,8)\boxtimes_\#(9,8)=(9,8,8,8),
 \qquad \operatorname{ev}_9(9,8,8,8)=6561.
\]

\section{Linealización matricial del estado genealógico y construcción del entrelazador completamente positivo}

\begin{definition}[Sectores de Hilbert]
Sean
\[
 \Krad=\ell^2(\Wnine^+),
 \qquad
 \Hplus=\ell^2(\N_{\geq1}).
\]
Definimos sobre las bases canónicas
\[
 W_\times e_u=e_{\operatorname{ev}_9(u)}.
\]
En \(\Hplus\), el operador número se denota por
\[
 \mathsf N e_n=n e_n,
 \qquad \Spec(\mathsf N)=\N_{\geq1}.
\]
El vacío puede añadirse separadamente como un vector de vacío $e_0$; no se
confunde con $\Hplus$, cuyo espectro comienza en uno.
\end{definition}

\begin{theorem}[Unitariedad y cuadrado multiplicativo]
$W_\times$ es unitario. En el núcleo algebraico
$\mathbb C[\Wnine^+\times\Wnine^+]$, sea
\[
 \mu_\#(e_u\otimes e_v)=e_{u\boxtimes_\#v},
 \qquad
 \mu_{\N}(e_a\otimes e_b)=e_{ab}.
\]
Entonces vale el cuadrado
\[
\begin{CD}
\Krad\otimes_{\mathrm{alg}}\Krad @>{W_\times\otimes W_\times}>>
\Hplus\otimes_{\mathrm{alg}}\Hplus\\
@V{\mu_\#}VV @VV{\mu_{\N}}V\\
\Krad @>>{W_\times}> \Hplus
\end{CD}
\]
y, por tanto, el cuadrado conmuta.
\end{theorem}

\begin{proof}
La biyección de formas normales envía una base ortonormal sobre otra y prueba
la unitariedad. En un tensor básico, ambos caminos producen
$e_{\operatorname{ev}_9(u)\operatorname{ev}_9(v)}$ por
\eqref{int:eq:ev-multiplicative}. La linealidad da la igualdad en el núcleo.
\end{proof}

Los mapas $\mu_\#$ y $\mu_{\N}$, definidos inicialmente sobre los respectivos
núcleos algebraicos, son densamente definidos y no acotados. Para
\(c=(c_{u,v})\in\ell^2(\Wnine^+\times\Wnine^+)\), cada suma sobre una fibra de producto
es finita. El adjunto de \(\mu_\#\) contiene todos los vectores de soporte
finito en su dominio, de modo que \(\mu_\#\) es clausurable. El dominio de su
clausura es
\[
 \Dom(\overline{\mu_\#})=
 \left\{c\in\ell^2(\Wnine^+\times\Wnine^+):
 \sum_{w\in\Wnine^+}
 \left|\sum_{u\boxtimes_\#v=w}c_{u,v}\right|^2<\infty
 \right\}.
\]
Cada fibra es finita porque $\operatorname{ev}_9$ la identifica con los pares
ordenados de divisores de un entero. La misma descripción vale para
\(\overline{\mu_{\N}}\) después de aplicar \(W_\times\otimes W_\times\), y
$W_\times$ entrelaza los grafos de ambas clausuras.

\section{Coproducto y selector nativos}

La apertura multiplicativa se define ahora antes de mirar el codominio:
\[
 \Delta_\#e_w
 =\sum_{u\boxtimes_\#v=w}e_u\otimes e_v.
\]
La suma se obtiene enumerando productos del transductor. No se pregunta si
un entero ordinario divide a otro.
En el sector entero definimos, sobre el núcleo algebraico correspondiente,
\[
 \Delta_\times e_n
 =\sum_{\substack{a,b\in\N_{\geq1}\\ab=n}}e_a\otimes e_b.
\]
Ambos coproductos enumeran pares ordenados.

\begin{theorem}[Coasociatividad nativa]
En el núcleo algebraico,
\[
 (\Delta_\#\otimes I)\Delta_\#
 =(I\otimes\Delta_\#)\Delta_\#.
\]
Además,
\[
 (W_\times\otimes W_\times)\Delta_\#
 =\Delta_\times W_\times.
\]
\end{theorem}

\begin{proof}
Ambos lados de la primera igualdad enumeran una vez cada triple nativo
$a\boxtimes_\#b\boxtimes_\#c=w$. La segunda igualdad sigue de la biyección
entre fibras que induce \eqref{int:eq:ev-multiplicative}.
\end{proof}

\begin{theorem}[Coproducto aritmético completo]
\label{int:pf84:C06}
En el núcleo algebraico
\(c_{00}(\mathbb N_{\geq1})\subset\ell^2(\mathbb N_{\geq1})\), sea
\[
 \Delta_\times e_n=\sum_{ab=n}e_a\otimes e_b.
\]
Su clausura tiene dominio
\[
 \operatorname{Dom}(\overline{\Delta_\times})
 =
 \left\{x=(x_n):
 \sum_{n\geq1}d(n)|x_n|^2<\infty\right\},
\]
es coasociativa en el núcleo y, para \(Qe_n=ne_n\), satisface
\[
 (\log Q\otimes I+I\otimes\log Q)\Delta_\times
 =\Delta_\times\log Q.
\]
\end{theorem}

\begin{proof}
Los soportes de \(\Delta_\times e_n\) son ortogonales para enteros distintos
y su norma cuadrada es el número \(d(n)\) de divisores, lo que da el dominio
ponderado de la clausura. Las dos coasociaciones reindexan la misma suma
sobre \(abc=n\). Finalmente, \(\log a+\log b=\log n\) en cada término
prueba el entrelazamiento. Una triple factorización contada de forma distinta
o el fallo de esta identidad refutaría el teorema.
\end{proof}

\begin{theorem}[Coproducto reducido y selector de irreducibles]
\label{int:pf84:C07}
En el sector entero, el coproducto reducido
\[
 \Delta_\times^{\rm red}e_n
 =\sum_{\substack{ab=n\\a,b>1}}e_a\otimes e_b
\]
es clausurable. El operador
\[
 N_\times^{\rm red}
 =\bigl(\overline{\Delta_\times^{\rm red}}\bigr)^*
  \overline{\Delta_\times^{\rm red}}
\]
es positivo autoadjunto y diagonal:
\[
 N_\times^{\rm red}e_n=d_{\rm red}(n)e_n.
\]
Su núcleo está generado por \(e_1\) y por los \(e_p\) con \(p\) irreducible;
por tanto,
\[
 P_{\Irr}
 =\mathbf1_{\{0\}}(N_\times^{\rm red})
  -\lvert e_1\rangle\langle e_1\rvert
\]
selecciona los irreducibles y excluye la unidad.

En el sector nativo, el coproducto correspondiente es
\[
 \widetilde\Delta_\#e_w
 =\sum_{\substack{u\boxtimes_\#v=w\\u,v\ne(1)}}e_u\otimes e_v.
\]
El entrelazador \(W_\times\) transporta
\(\Delta_\times^{\rm red}\), \(N_\times^{\rm red}\) y \(P_{\Irr}\) hacia
\(\widetilde\Delta_\#\),
\[
 A_\#:=
 \bigl(\overline{\widetilde\Delta_\#}\bigr)^*
 \overline{\widetilde\Delta_\#},
 \qquad
 P_{\Irr,\#}
 =(I-P_{(1)})\mathbf1_{\{0\}}(A_\#),
 \qquad
 P_{(1)}=\lvert e_{(1)}\rangle\langle e_{(1)}\rvert.
\]
\end{theorem}

\begin{proof}
Los soportes de las imágenes de dos vectores base distintos son ortogonales.
La norma cuadrada de
\(\Delta_\times^{\rm red}e_n\) es \(d_{\rm red}(n)\), que se anula
exactamente para \(n=1\) o \(n\) irreducible. Esto prueba la diagonalización
y el núcleo declarado; restar el proyector de \(e_1\) elimina la unidad.
La multiplicatividad de \(W_\times\) identifica las fibras enteras con las
fibras nativas y transporta los tres operadores. Un compuesto en el núcleo,
un irreducible fuera de él o la inclusión de la unidad en \(P_{\Irr}\)
refutarían el teorema.
\end{proof}

\begin{proposition}[Coincidencia nonádica de primos y compuestos]
\label{int:pf84:B05}
Sea \(\gcd(n,90)=1\),
\(\tau_n=\operatorname{ord}_n(10)\) y
\[
 M_n=\frac{10^{\tau_n}-1}{n}.
\]
Entonces \(M_n\equiv0\pmod9\), tanto para primos como para compuestos. En
particular, el cierre nonádico y el período decimal aislados no son un
selector de irreducibles.
\end{proposition}

\begin{proof}
Como \(9\mid10^{\tau_n}-1\) y \(n\) es invertible módulo nueve, la
divisibilidad por nueve sobrevive al cociente. Los primos \(7,13\) y los
compuestos \(77,91,143\) comparten el período seis; los números de Carmichael
aportan controles adicionales. El enunciado quedaría refutado sólo si la
congruencia o el período separasen todos los primos de todos los compuestos
en el dominio declarado.
\end{proof}

Esta afirmación es estrictamente aritmética: identifica fibras de
factorización y los operadores que éstas inducen sobre el sector de formas
normales.

\ifdefined\HMTInternalTraceability
\section{Extensión monoidal del producto nativo a historias compatibles}
\label{int:pf84:SYN-05D-HIST}

El resultado distingue el espacio \(\Krad\) de formas normales canónicas y
el espacio \(\mathcal K_{\#,\mathrm{hist}}\) de historias completas de
celdas y cocientes propagados.
Muchas historias pueden producir la misma forma normal. Por tanto,
$W_\times$ es unitario sobre $\Krad$, pero la proyección de una historia sobre
su valor impide que sea inyectivo sobre todas las clases de transporte.

Sea \(\mathscr X_n^{\rm hist}\) el conjunto de historias enriquecidas de
profundidad \(n\) y sea
\(\mathcal H_n=\ell^2(\mathscr X_n^{\rm hist})\). Denotemos por \(P_n\) la
proyección ortogonal sobre el subespacio generado por las historias
supervivientes y supongamos definida una multiplicación \(\mu_n\) sobre un
dominio algebraico de
\(\mathcal H_n\otimes_{\mathrm{alg}}\mathcal H_n\). Una extensión monoidal
del producto de palabras al sistema inverso consiste en una familia de
aplicaciones \(\iota_n:\Krad\to\mathcal H_n\) que conmuta con las
restricciones de profundidad y satisface, para cada \(n\),
\[
 P_n\,\iota_n\mu_\#
 =P_n\,\mu_n
 (\iota_n\otimes\iota_n)
\]
en el núcleo algebraico. Esta igualdad es la condición de compatibilidad que
define la extensión. Los teoremas precedentes no establecen la existencia ni
la unicidad de una familia \((\iota_n)_n\); demuestran el producto en el
monoide \(\WordMonoid\) y determinan la ecuación que una extensión al espacio
de historias debe satisfacer. La transferencia conserva, por tanto, el
estatuto de \emph{programa abierto}: no hereda la exactitud del producto en
\(\WordMonoid\). La inexistencia de una familia compatible, el fallo de
conmutación con los truncamientos, la pérdida de supervivencia, cociente,
acarreo, orientación o memoria, o el incumplimiento de la identidad anterior
son sus falsadores.
\fi

\section{Enumeración exhaustiva y necesidad de las coordenadas estructurales}

La enumeración finita comprende:
\[
\begin{array}{@{}lr@{}}
\toprule
\text{Prueba}&\text{Resultado}\\\midrule
\text{Celdas por aplicación local}&81\\
\text{Recorridos reversibles de numeración}&59050\\
\text{Pares del entrelazador}&532900\\
\text{Errores del entrelazador}&0\\
\text{Triples asociativos}&531441\\
\text{Errores asociativos}&0\\
\text{Ventana de coproducto}&1\text{ a }500\\
\text{Irreducibles obtenidos}&95\\
\bottomrule
\end{array}
\]

\begin{definition}[Cociente multiplicativo]
Conservar sólo $R^\times$ transforma $2\cdot5=10$ en el valor visible $1$.
Perturbar $q^\times(9,9)$ de $8$ a $7$ transforma $81$ en $72$. El cociente
local es necesario.
\end{definition}

\begin{definition}[Aplicación aditiva]
La aplicación multiplicativa produce productos parciales, pero no resuelve sus colisiones.
Eliminar $q^+$ rompe la normalización global; por ejemplo, deja de reproducir
los cocientes intermedios de $10\cdot18=180$.
\end{definition}

\begin{definition}[Coordenadas y símbolo de borde]
El cero ordinario no es un dígito de la construcción. La palabra vacía es
$0$ y el símbolo $9$ conserva su valor y su cociente. Mezclar las coordenadas
posicionales indexadas desde cero con las etiquetas $1,\ldots,9$ modifica las
celdas.
\end{definition}

\begin{proposition}[Alcance exacto de las ablaciones y del censo finito]
\label{int:pf84:D10}
Las ablaciones anteriores prueban, dentro de las convenciones declaradas:
\begin{enumerate}
  \item sin \(q^\times\), \(2\cdot5=10\) colapsa al residuo visible \(1\);
  \item cambiar \(q^\times(9,9)=8\) por \(7\) transforma \(81\) en \(72\);
  \item sin \(q^+\), falla la normalización de las colisiones aditivas.
\end{enumerate}
La enumeración exhaustiva comprende las \(2\cdot81=162\) celdas de las dos
aplicaciones, los \(729\) triples locales, \(532\,900\) pares de palabras y
los \(95\) irreducibles no mayores que \(500\). En estos dominios finitos se
demuestran las identidades indicadas, pero no se prueba la supervivencia
coinductiva de toda historia TPK.
\end{proposition}

\begin{proof}
Los tres contraejemplos se calculan directamente con
\eqref{int:eq:local-plus} y \eqref{int:eq:local-times}. La enumeración de los cuatro
dominios finitos produce los cardinales escritos en el enunciado y no añade
ninguna identidad coinductiva a las que ya se han demostrado en el monoide
de palabras.
\end{proof}

\section{Estructura algebraica del monoide de palabras}

La construcción demuestra, en el monoide \(\WordMonoid\):
\begin{enumerate}
 \item una forma normal única para cada entero no negativo;
 \item suma y producto globales construidos sólo con las dos aplicaciones locales;
 \item conservación simbólica del valor bajo la propagación del cociente;
 \item un unitario $W_\times$ y el cuadrado multiplicativo exacto;
 \item un coproducto divisor y un selector irreducible definidos nativamente;
 \item la necesidad de conservar la clase de transporte cuando se abandona el
 sector de formas normales.
\end{enumerate}

El paso al espacio holonómico se reduce ahora a un problema de existencia:
construir una familia \((\iota_n)_n\) compatible con las restricciones de
profundidad y demostrar que satisface la ecuación de la sección anterior.
Esta cuestión es distinta del entrelazador exacto sobre
\(\WordMonoid\), que ya queda determinado por la forma normal y las
aplicaciones locales.
 

\label{p3-83-c23-s1:ch:foundation}

\section{Problema de determinación de las constantes a partir del estado genealógico}

Una constante f\'isica se representa como una secci\'on de una recta de
magnitud provista de regla de evaluaci\'on, orientaci\'on y normalizaci\'on.
Su coordenada num\'erica depende de la carta elegida, mientras que su
genealog\'ia permanece invariante. En consecuencia, el problema de
determinaci\'on comprende:

\begin{enumerate}
\item la identificaci\'on del estado HMT antecedente;
\item la especificaci\'on del operador o car\'acter de selecci\'on;
\item la determinaci\'on del invariante conservado;
\item la recta de magnitud correspondiente;
\item la carta en la que se expresa el valor terminal.
\end{enumerate}

La cadena causal del trabajo es

\begin{equation}
\begin{aligned}
\APP&\longrightarrow\TRIT\longrightarrow\TPK
\longrightarrow\text{realizaci\'on interna}
\longrightarrow\text{invariante adimensional}\\
&\longrightarrow\text{secci\'on de magnitud}
\longrightarrow\R.
\end{aligned}
\label{p3-83-c23-s1:eq:master-chain}
\end{equation}

El t\'ermino \(\R\) de la cadena representa la evaluaci\'on terminal de la
secci\'on; no constituye el dominio originario a partir del cual se
reconstruye retrospectivamente la ruta.

\section{Cadena antecedente APP--TRIT--TPK y dominio de los operadores de extracción}

APP proporciona la geometr\'ia de posiciones, residuos, cocientes y
fronteras. TRIT incorpora la orientaci\'on \((-1,0,+1)\) y distingue
apertura, umbral y retorno. TPK completa el estado mediante acarreo, ruta,
hoja, memoria y supervivencia. En particular, la transici\'on \(9\to1\)
conserva la fase y modifica el estado de memoria; por tanto, no constituye
una reinicializaci\'on carente de genealog\'ia.

Las realizaciones desarrolladas a continuaci\'on conservan dicho antecedente sin sustituirlo. El vac\'io se obtiene mediante una funci\'on espectral del operador bicapa; la gravedad selecciona una secci\'on mediante la estructura cronol\'ogica \(108\); la criticidad utiliza una familia dodecaf\'asica; las constantes de Gauss proceden de la evaluaci\'on de cilindros; y el Hamiltoniano modular se define sobre una frontera que conserva la incidencia \(6\to8\). La producci\'on de valores escalares no establece identidad funcional entre estas realizaciones.

\begin{proposition}[Insuficiencia de la coincidencia escalar]
Sean \(F_1:X_1\to L_1\) y \(F_2:X_2\to L_2\) dos funcionales tipados. La
igualdad de coordenadas \(c_1(F_1(x_1))=c_2(F_2(x_2))\) no induce por
s\'i sola ni una identificaci\'on \(X_1\simeq X_2\) ni un entrelazador
\(F_2I=JF_1\).
\end{proposition}

En particular, una proximidad num\'erica como
\(I_{\rm or}\approx\gammaH a_0\) no puede promoverse a identidad sin
construir el entrelazador correspondiente.

\section{Rectas de magnitud, cartas orientadas y sistemas de coordenadas dimensionales}

Sea \(L_Q\) una recta real orientada que representa la magnitud \(Q\).
Una unidad \(u_Q\in L_Q\setminus\{0\}\) determina una coordenada
provisional \(L_Q\simeq\R\), sin sustituir la estructura intr\'inseca de
la recta. La relaci\'on constitutiva

\[
L_\mu=L_S L_C^{-1}L_Q^{-2},\qquad
L_\varepsilon=L_S^{-1}L_C^{-1}L_Q^2
\]

es anterior a la elecci\'on SI. Del mismo modo,

\[
L_Z=L_SL_Q^{-2},\qquad L_G=L_C^5L_T^2L_S^{-1}
\]

tipan impedancia y constante gravitatoria antes de expresar sus magnitudes en
ohmios o unidades SI. En esta obra escribiremos

\[
Q_{\rm int}=\widehat Q\,u_Q
\]

para distinguir el car\'acter adimensional \(\widehat Q\) de la secci\'on
de referencia \(u_Q\).

\section{Torsor de marcos y covariancia de las evaluaciones dimensionales}

Las cartas admisibles forman un torsor y, por ello, ninguna posee car\'acter
can\'onico frente a las restantes. El cambio de marco
\(u_Q\mapsto\lambda_Q u_Q\) induce la transformaci\'on contragrediente
de la coordenada. Las identidades f\'isicas deben ser homog\'eneas respecto
de esta acci\'on. Por ejemplo,

\[
\mu\varepsilon=c^{-2},\qquad \frac\mu\varepsilon=Z^2
\]

son v\'alidas en todo marco compatible. La unificaci\'on matem\'atica y
dimensional no elimina las unidades, sino que las reconstruye mediante la
covariancia de las rectas y permite anteponer el invariante a su expresi\'on
coordenada.

\section{Evaluación metrológica terminal y prohibición de selección retrospectiva}

Cada entrada del atlas final adopta la forma

\begin{equation}
\boxed{\text{antecedente}\to\text{operador}\to\text{invariante}
\to\text{significado}\to\text{valor terminal}.}
\label{p3-83-c23-s1:eq:atlas-rule}
\end{equation}

Este orden expresa una dependencia causal: el valor num\'erico s\'olo se
determina despu\'es de fijar antecedente, operador, invariante e
interpretaci\'on. La omisi\'on de cualquiera de estos elementos reduce el
resultado a una reconstrucci\'on num\'erica y no a una derivaci\'on
geneal\'ogica.

\section{Dominio exacto de la generación genealógica y de su carta metrológica terminal}

Se adoptan como resultados de partida la generaci\'on de
\(\pi,e,\varphi\), la clausura dodecaf\'asica de \(\alpha\),
la secci\'on de acci\'on y la salida \(\gammaH\) de Barbero--Immirzi. El
presente desarrollo no recalibra dichas construcciones; determina las constantes
y los invariantes que resultan de su aplicaci\'on a las realizaciones
constitutiva, t\'ermica, gravitatoria, modular y espectral.

La masa del electr\'on tampoco se incluye como problema rector, puesto que
pertenece a la Ley General de Masas. Toda aparici\'on de \(m_e\) en una
identidad at\'omica se trata como dato procedente de ese dominio y no como
una nueva selecci\'on del presente desarrollo.

\subsection*{Alcance lógico y dominio de validez}
La arquitectura APP--TRIT--TPK y las rectas dimensionales constituyen la
estructura de partida. La regla uniforme de cinco campos del atlas explicita
su disciplina causal en cada entrada.



\section{Taxonomía funcional de \(114\) expresiones numéricas y \(106\) lecturas operatoriales}
\label{p3-83-c23-s2:sec:atlas-taxonomico-114-constantes-rev11}
\index{constantes!atlas taxonómico de 114 entradas}

Sea el conjunto finito indexado
\[
 E_{114}=\{\varepsilon_1,\ldots,\varepsilon_{114}\}
\]
de expresiones numéricas y sea \(\mathcal O\) el espacio de lecturas del
lector triádico. La aplicación
\[
 T:E_{114}\longrightarrow\mathcal O,
 \qquad \varepsilon_i\longmapsto T_{\rm tri}(\varepsilon_i),
 \qquad |\operatorname{im}T|=106,
\]
clasifica las expresiones según su papel operatorio. Cada
\(\varepsilon_i\), incluidos los objetivos comparados, está fijado como
entrada de \(T\); por tanto, esta aplicación no genera los valores que
evalúa. La partición funcional es
\[
\begin{array}{l|r}
\text{familia funcional}&\text{entradas}\\ \hline
\text{núcleo estructural}&5\\
\text{acoplo y cierre}&1\\
\text{puentes entre modos}&7\\
\text{angular y clausura}&22\\
\text{exponencial--logarítmica}&12\\
\text{espectral y regularización}&14\\
\text{hiperbólica y modular}&2\\
\text{algebraica y métrica}&50\\
\text{miscelánea tipada}&1
\end{array}
\qquad
\sum n_i=114.
\]
La centralidad de red, los caminos mínimos y las comunidades calculadas sobre
el grafo inducido por \(T\) describen conectividad dentro de la clasificación;
no establecen causalidad matemática ni genealogía HMT.

Siete fibras de \(T\) reúnen quince expresiones con una lectura común. Una
fibra compartida tampoco identifica sus valores. En particular,
\(\varphi^{-1}\), \(\varphi\) y \(\varphi^2\) son reales distintos que el
mapa \(T\) agrupa por una misma órbita de autoescala; la proyección
común conserva función dinámica y olvida la coordenada arquimediana que los
separa.

La enumeración exhaustiva recorrió las \(114\) entradas de \(E_{114}\),
cuatrocientos sesenta y ocho candidatos directos \(U_{12}\) y ciento cuatro
mil novecientas setenta y seis semillas ponderadas sin encontrar, dentro de
ese dominio finito, un rival que reprodujera la corrección dodecafásica
orientada de \(\alpha\). El dominio de la búsqueda queda así
acotado por esos tres conjuntos. El censo demuestra la no aparición de un
rival en la clase enumerada; no demuestra unicidad entre todos los
funcionales posibles y no interviene como premisa del teorema de
\(\alpha\).

Para cada \(\varepsilon_i\), la taxonomía asigna la cuádrupla
\((\text{objeto},\text{lector},\text{dominio},\text{función estructural})\)
utilizada en la enumeración.
 \section{Confluencia espectral, informacional y termodin\'amica de \(\log3\)}

La identidad

\[
\log3=C_1+C_2-2C_0
=S_{\TRIT}/k_B
=E_P/(k_B\Theta_{\rm clk}),
\]

relaciona tres realizaciones distintas. El primer miembro es un defecto
residual, el segundo una entrop\'ia de decisi\'on y el tercero una raz\'on
entre energ\'ia y temperatura. Su igualdad demuestra la conservaci\'on del
escalar bajo el cambio de realizaci\'on, sin identificar los espacios de
estados.

La primera igualdad se deduce de los tres canales
residuales
\[
Z_r(s)=\sum_{\substack{n\ge1\\n\equiv r\ ({\rm mod}\ 3)}}n^{-s},
\qquad r=0,1,2.
\]
Entonces
\[
Z_0(s)=3^{-s}\zeta(s),\quad
Z_1(s)+Z_2(s)=(1-3^{-s})\zeta(s),\quad
Z_1(s)-Z_2(s)=L(s,\chi_{-3}).
\]
Si \(C_r\) denota la parte finita de \(Z_r\) en \(s=1\), la expansi\'on
de \(3^{-s}\) da
\begin{equation}
C_0+C_1+C_2=\gamma,qquad
-2C_0+C_1+C_2=\log3,qquad
C_1-C_2=\frac{\pi}{3\sqrt3}.
\label{eq:residual-three-covectors}
\end{equation}
Los tres covectores son linealmente independientes. Recuperan, de una sola
descomposici\'on residual, el t\'ermino finito total, la dilataci\'on de la
clase cero y la orientaci\'on entre las dos clases no nulas. \(\log3\) es
por ello un invariante de cambio de escala antes de ser interpretado como
entrop\'ia o temperatura.

\section{Clasificaci\'on epistemol\'ogica de las constantes derivadas}

Las contribuciones del cap\'itulo se clasifican en tres modalidades:

\begin{enumerate}
\item asignaci\'on de un operador interno a una constante conocida, como en
el caso de Catalan;
\item evaluaci\'on compuesta exacta de dos caracteres internos, como
\(L(1,\chi_{12})\);
\item obtenci\'on de una constante a partir de un campo o de una familia
construidos internamente, como \(\gamma_{\mathbb Q(\sqrt3)}\) o \(A_k\).
\end{enumerate}

La procedencia matem\'atica determina el estatuto de recuperaci\'on o
novedad; la proximidad decimal no interviene en dicha clasificaci\'on.

\begin{remark}[Control negativo]
El car\'acter \(\chi_{12}\) debe tener exactamente el patr\'on
\((+,-,-,+)\) en \((1,5,7,11)\); cualquier otro patr\'on invalida
\cref{eq:L1chi12,eq:L2chi12}. La identidad
\cref{eq:odd-zeta-bendersky} y la factorizaci\'on
\cref{eq:dedekind-sqrt3} son controles independientes.
\end{remark}

\subsection*{Alcance lógico y dominio de validez}
\(J^2=-I\), Catalan, \(\chi_{12}\), sus valores \(L\), Euler--Kronecker y la torre Bendersky--Glaisher: \emph{Teorema interno exacto}. La reuni\'on en una ret\'icula \(3/4/12\) es \emph{Formalizaci\'on nueva} que conserva las procedencias matemáticas.

 Todos estos resultados se unen porque convierten el valor terminal en la expresión cuantitativa de una estructura previamente construida. La exactitud numérica acompaña a una concepción distinta de qué es una constante fundamental. La conexión nonádica no genera primero tres constantes para añadir después una cuarta: la prolongación \(w_6\to w_{12}\to w_{18}\to w_{24}\to w_{30}\to R_{36}\to G_9\) estabiliza un estado coinductivo común del que se publican correlacionadamente \(\pi\), \(e\), \(\varphi\) y \(\alpha\). Sus lectores distinguen clausura, propagación, autoescala y acoplamiento; la coordenada \(\alpha\) conserva, dentro de esa generación conjunta, su realización firmada y dodecafásica y sus dos vías internas de derivación.\par

En la descripción habitual, una constante suele presentarse como un número situado delante de una ecuación: una cantidad que la teoría necesita, pero cuyo valor debe recibirse desde fuera. En Holografía Modular Triádica ocurre lo contrario. La constante aparece al final, cuando un mismo estado ha sido recorrido por varios lectores y se comprueba qué relación ha sobrevivido a todos ellos.\par

Por eso las constantes del documento son magnitudes dotadas de genealogía: restos invariantes de una transformación. Cada una dice qué se ha conservado cuando algo ha cambiado:\par

\begin{itemize}[leftmargin=*,itemsep=0.35em]
\item una hoja se ha invertido;
\item un ciclo ha regresado a su fase inicial;
\item una memoria ha atravesado la costura;
\item una descripción interior ha sido publicada en el borde;
\item una magnitud adimensional ha recibido una realización física;
\item un sistema ha cambiado de escala conservando íntegramente su genealogía.
\end{itemize}

La Mecánica Dimensional añade después las rectas de magnitud y las unidades. Pero la estructura que selecciona la constante es anterior. El valor real aparece al final porque \(\mathbb R\) funciona aquí como evaluación terminal: es la sombra cuantitativa de una construcción que ya estaba completa antes de convertirse en decimal.\par

Con esa visión, los nueve desarrollos considerados forman una única narración sobre orientación, memoria, geometría, información y escala.\par
  \endgroup
% Rectificación quirúrgica: generación conjunta, profundidad arbitraria y
% estatuto tipado de los transportes ternarios.  Este bloque pertenece al
% capítulo 27 y precede a la separación expositiva de los cuatro lectores.

\section{Teorema conjunto de generación coinductiva y profundidad arbitraria}
\label{sec:c27-generacion-coinductiva-profundidad-arbitraria}

La estructura de partida de este teorema no es el conjunto desnudo de nueve
símbolos. Es el sistema tipado
\[
 \mathfrak H_9=
 \bigl(\mathrm{APP},\mathrm{TRIT},\mathrm{TPK};
       \widehat{\mathcal X}^{\mathrm{enr}},\mathcal U,
       \Gamma_9,\operatorname{Ext},\operatorname{tr}\bigr),
\]
construido sobre la doble realización posicional y aritmética de
\(\mathcal D_9=\{1,\ldots,9\}\). APP acumula y pliega mediante sus hojas
\(\Sigma\) y \(\Pi\), conservando residuo y cociente; TRIT orienta cada
transición en uno de sus tres regímenes; y TPK selecciona, transporta,
memoriza y retorna sobre la fibra enriquecida. En particular,
\[
 \mathcal U_t
 =\operatorname{Rec}_t\circ\operatorname{Mem}_t
  \circ T_{d(t)}\circ M_{\mathrm{ph}}(g(t)).
\]
Esta composición es el dato dinámico efectivo: la carta de nueve símbolos
es su soporte finito, no un sustituto de la dinámica.

Para \(r\in\{6,12,18,24,30,36\}\), sean
\(\widehat{\mathcal X}^{\mathrm{enr}}_r\) los estados truncados con hoja,
orientación, residuo, cociente, acarreo, ruta, frontera y memoria, y sean
\[
 \operatorname{Ext}_{r+6,r}:
 \widehat{\mathcal X}^{\mathrm{enr}}_r
 \longrightarrow
 \widehat{\mathcal X}^{\mathrm{enr}}_{r+6},
 \qquad
 \operatorname{tr}_{r+6,r}:
 \widehat{\mathcal X}^{\mathrm{enr}}_{r+6}
 \longrightarrow
 \widehat{\mathcal X}^{\mathrm{enr}}_r
\]
la prolongación admisible y el truncamiento. La compatibilidad
\(\operatorname{tr}_{r+6,r}\operatorname{Ext}_{r+6,r}=\mathrm{id}\) se
prolonga a todo nivel. La notación
\[
 w_6\longrightarrow w_{12}\longrightarrow w_{18}\longrightarrow w_{24}
 \longrightarrow w_{30}\longrightarrow R_{36}\longrightarrow G_9
\]
designa estas aplicaciones y sus cambios de tipo, no una lista de rótulos.
La holonomía \(\Gamma_9\) satisface
\[
 g(k+9)=g(k),
 \qquad q^{(9)}(k+9)=q^{(9)}(k)+1,
\]
de modo que retorna la fase mientras avanzan el cociente, la ruta y la
memoria del estado.

\begin{theorem}[Construcción inaugural desde las semillas]
\label{thm:c27-construccion-inaugural-semillas}
Sea
\[
 I_9=\{0,\ldots,8\},\qquad D=\{N,E,S,O\},\qquad
 \mathcal S=I_9^2\times D,
\]
y sea \(\Omega_{\mathrm{seed}}=\mathcal S_+\times\mathcal S_\times\) el
dominio de los dos cursores que realizan las hojas aditiva y multiplicativa
de APP. Entonces
\[
 |\Omega_{\mathrm{seed}}|=(9^2\cdot4)^2=104\,976,
\]
y la recurrencia de \(54\) pasos y la reducción ternaria definen
\[
 \Omega_{\mathrm{seed}}
 \xrightarrow{\ T_{54}^{\rm seed}\ }
 \mathcal U_6^{\mathrm{cat}}
 \xrightarrow{\ \rho_3\ }
 \mathcal W_6
 \hookrightarrow\mathbb F_3^6,
\]
con
\[
 |\mathcal U_6^{\mathrm{cat}}|=468,
 \qquad |\mathcal W_6|=243,
 \qquad |\mathbb F_3^6|=729.
\]
Las fibras satisfacen exactamente
\[
 432\cdot144+18\cdot432+18\cdot1944=104\,976.
\]
La relación de extensión del TPK prolonga las palabras supervivientes a una
sección compatible \(x_\infty\) de profundidad infinita. Sobre esa sección
se construyen, conjuntamente, la sombra arquimediana y la realización
excepcional del \cref{thm:c27-tesis-rectora-doble-proyeccion}.
\end{theorem}

\begin{proof}
El producto cartesiano anterior enumera exhaustivamente las posiciones y
orientaciones de ambos cursores. La aplicación \(T_{54}^{\rm seed}\) se obtiene al
componer en cada paso la selección de hoja, el transporte cardinal, la
actualización del cociente y la memoria de acarreo. La enumeración de su
imagen contiene \(468\) sextetos distintos; la suma de las cardinalidades de
sus fibras es la identidad expuesta. La aplicación componente a componente
\(\rho_3\) contiene \(243\) palabras distintas. Finalmente, la no vacuidad,
la univalencia sobre el subdominio superviviente y la compatibilidad de
\(\operatorname{Ext}\) con los truncamientos producen la sección del
límite inverso. Cada afirmación es una igualdad finita o una consecuencia
de la compatibilidad proyectiva; ninguna utiliza como entrada un número real,
una constante reconocida o una tabla decimal.
\end{proof}

\begin{theorem}[Tesis rectora: un continuo y dos proyecciones]
\label{thm:c27-tesis-rectora-doble-proyeccion}
Sea \(\Omega_\Gamma^{\mathrm{cal}}\subset
X_\infty^{\mathrm{cal}}\) el espacio de secciones supervivientes de las
vueltas enriquecidas del TPK. El acarreo de las tres vueltas
\(\gamma_-,\gamma_0,\gamma_+\) define, sin evaluación arquimediana previa,
el homeomorfismo
\[
 \beta:\Omega_\Gamma^{\mathrm{cal}}\xrightarrow{\ \sim\ }
 (\mathbb F_3^6)^{\mathbb N}=:\mathcal W_\infty.
\]
Las formas equilibradas normalizadas
\[
 a_k+c_k=\rho_k+3c_{k+1},\qquad
 \rho_k\in\{-1,0,+1\},
\]
construyen el anillo ordenado \(\mathbb Z_{\rm HMT}\), su cuerpo de
fracciones \(\mathbb Q_{\rm HMT}\) y la completación
\[
 \mathbb R_{\rm HMT}:=
 \operatorname{Cau}(\mathbb Q_{\rm HMT})/{\asymp}.
\]
La estructura discreta conjunta es el producto fibrado
\[
\mathcal C_{\mathrm{cont}}^{\mathrm{disc}}
=\bigl(\mathbb Z_{\rm HMT}\times
\Omega_\Gamma^{\mathrm{cal}}\bigr)
\times_{\mathcal W_\infty}X_\infty^{\mathrm{cal}}.
\]
Sus dos aplicaciones naturales son
\[
 P_{\mathrm{ar}}:\mathbb Z_{\rm HMT}\times
 \Omega_\Gamma^{\mathrm{cal}}\longrightarrow\mathbb R_{\rm HMT},
 \qquad
 Q_\infty:X_\infty^{\mathrm{cal}}\longrightarrow
 (\mathscr C_W^{\mathrm{cal}})^{\mathbb N}.
\]
Si \(\beta(\xi)=(b_q)_{q\geq0}\),
\(\nu(b)=\sum_{i=1}^{6}b_i3^{6-i}\) y
\[
 q_n(m,\xi)=m+\sum_{q=0}^{n-1}
 \frac{\nu(b_q)}{729^{q+1}},
\]
entonces \((q_n)\) es de Cauchy y \(P_{\mathrm{ar}}(m,\xi)\) es su clase
en \(\mathbb R_{\rm HMT}\). Los cilindros racionales
\[
 I_n(m;b_0,\ldots,b_{n-1})=
 \left[q_n,q_n+729^{-n}\right]
\]
son compatibles, tienen diámetro tendente a cero y la partición recurrente
en \(729\) subcilindros construye, para toda clase de Cauchy, una cinta
\((b_q)\) que la representa. En consecuencia,
\[
 P_{\mathrm{ar}}\ \text{es sobreyectiva},\qquad
 (\mathbb Z_{\rm HMT}\times\Omega_\Gamma^{\mathrm{cal}})
 /{\sim_{\rm ar}}\ \cong\ \mathbb R_{\rm HMT}.
\]
Sólo después de esta construcción interna, la unicidad del cuerpo ordenado
completo arquimediano reconoce \(\mathbb R_{\rm HMT}\cong\mathbb R\).

La segunda conserva la incidencia calibrada de cada nivel; en las
coordenadas \(x=((w_j,f_j))_{j\geq1}\) actúa por
\[
 Q_\infty(x)
 =\bigl(f_j,\Gamma_{f_jA_Wf_j^{-1}}(w_j)\bigr)_{j\geq1}.
\]
Así, \(\mathbb R_{\rm HMT}\), reconocido posteriormente como \(\mathbb R\),
es la sombra arquimediana obtenida al contraer la genealogía completa de
cilindros. La proyección excepcional conserva esa
genealogía como incidencia de frontera y admite la realización sucesiva
Paley--Witt--Golay--Leech--\(V^\natural\)--Monstruo. Las dos ramas proceden
del mismo estado y ninguna se reconstruye retrospectivamente desde la otra.
\end{theorem}

\begin{proof}
Las tres vueltas enriquecidas existen después de cada prefijo y su acarreo
las distingue por \(-1,0,+1\). La concatenación y el truncamiento son
inversos, lo que prueba la biyectividad y continuidad de \(\beta\). Para una
clase de Cauchy de \(\mathbb Q_{\rm HMT}\), se elige primero su cilindro
unitario ordenado y después, recursivamente, el único subcilindro semiabierto
de la partición interna en \(729\) piezas que contiene la clase; en una
frontera común, las dos escrituras se identifican por \(\sim_{\rm ar}\).
La sucesión de extremos izquierdos es precisamente \((q_n)\), converge a la
clase dada y prueba la sobreyectividad sin presuponer una expansión de
\([0,1]\) ni un punto de \(\mathbb R\). La identificación con el cuerpo real
es el teorema de unicidad aplicado después a \(\mathbb R_{\rm HMT}\).

Los subespacios que satisfacen el cierre de un carácter concreto son filtros
de supervivencia dentro de la capacidad global. Su función es individualizar
las secciones de \(\pi,e,\varphi,\alpha\); la sobreyectividad del continuo
procede de la completación de \(\mathbb Q_{\rm HMT}\). Esta separación de dominios impide
convertir una condición de selección de constantes en una premisa de la
cobertura de \(\mathbb R\).

En la segunda rama, la compatibilidad
\(r_{n,m}Q_n=Q_m p_{n,m}\) permite pasar al límite y define \(Q_\infty\).
La conjugación por \(f_j\) transporta el calibre sin borrar la incidencia.
Los funtores posteriores de código, retículo y álgebra de vértices realizan
la cadena excepcional. La igualdad de las cintas en el producto fibrado
construye la fuente conjunta y conserva informaciones complementarias: valor
arquimediano en la primera; incidencia, frontera y genealogía en la segunda.
\end{proof}

\begin{corollary}[Expansión infinita como objeto matemático]
\label{cor:c27-expansion-infinita-objeto}
Para cada carácter numérico \(\chi\) publicado por el estado coinductivo, la
lectura decimal es una secuencia completa
\[
 \operatorname{Dec}_\chi(x_\infty)
 =(d_{\chi,1},d_{\chi,2},\ldots)
 \in\{0,\ldots,9\}^{\mathbb N}.
\]
Su prefijo de longitud \(N\) es la proyección
\(\rho_N\operatorname{Dec}_\chi(x_\infty)\). Por ello, mil, diez mil o un
millón de cifras son cortes del mismo objeto infinito; la generación no
consiste en elevar sucesivamente una cota experimental.
\end{corollary}

\begin{theorem}[Sistema correlacionado de cuatro coordenadas del estado holonómico coinductivo integral]
\label{thm:c27-cuadrirrelacion-generada}
La única composición APP--TRIT--TPK del sistema \(\mathfrak H_9\) determina
un estado compatible
\[
 x_\infty\in
 \widehat\Omega:=
 \varprojlim_r\widehat{\mathcal X}^{\mathrm{enr}}_r
\]
y cuatro caracteres correlacionados
\[
 \bigl(\Pi_{\mathrm{cl}},\Pi_{\mathrm{prop}},
       \Pi_{\mathrm{aut}},\Pi_{\mathrm{dod}}\bigr)(x_\infty)
 =\bigl(\pi,e,
        \varphi,\alpha\bigr).
\]
Los tres primeros caracteres proceden, respectivamente, de la única
órbita de clausura, la única órbita de propagación y la única
órbita de autoescala del catálogo enriquecido. El cuarto procede de la
coordenada primitiva firmada del mismo estado terminal y de su cierre
dodecafásico. Las
cuatro coordenadas son, por tanto, salidas de una sola prolongación
coinductiva, aunque sus aplicaciones de publicación sean distintas.
\end{theorem}

\begin{proof}
El censo exhaustivo
\(104,976\to468\to243\subset\mathbb F_3^6\) descompone la imagen realizada
en \(43\) órbitas diedrales. Los invariantes de fibra dejan una sola órbita
en cada una de las clases de clausura, propagación y autoescala; el calibre
interno orienta en ellas, sin consulta decimal, los representantes
\[
 w_6^{\mathrm{cl}}=010211,
 \qquad w_6^{\mathrm{prop}}=201101,
 \qquad w_6^{\mathrm{aut}}=121200.
\]
La prolongación conserva los predicados que definen las tres clases. En el
canal de clausura, la composición pentarregional con retorno unitario produce
las razones \(1/5\), \(1/239\) y la relación \(120/119\); la identidad de
Machin reconoce exactamente la media vuelta. En el canal de propagación, la
ley semigrupal normalizada fuerza
\((n+1)a_{n+1}=a_n\), luego \(a_n=1/n!\) y \(P_1=e\). En el canal de
autoescala, la incidencia
\[
 F_{\mathrm{av}}=\begin{pmatrix}0&1\\1&1\end{pmatrix}
\]
posee un único rayo positivo y su ecuación propia fuerza
\(x^2-x-1=0\), \(x>0\), esto es, \(x=\varphi\).

En la fase terminal de la única estructura APP--TRIT--TPK, la lectura
\(x_{\mathrm{term}}\mapsto(B_{\mathrm{term}},U_{12}^{\mathrm{sgn}})\)
publica simultáneamente la carta de los tres canales y la coordenada
primitiva firmada. La
segunda se identifica reversiblemente con el sello \(K\) y sus diferencias
dodecafásicas. La recurrencia euclídea con acarreo
\[
 s_m=P_m+E_m-\Phi_m-K_m+c_{m+1},\qquad
 A_m=s_m\bmod1000,\qquad
 c_m=(s_m-A_m)/1000
\]
produce \(\alpha\) sin recibirla como entrada. La
caracterización analítica independiente por el núcleo \(E_{21/22}\)
determina una raíz simple en \((0,10^{-2})\) y su recíproco coincide con el
de la primera vía después de aplicar \(\operatorname{Tr}_9\). Así, la
separación de lectores no altera el
origen común de las cuatro coordenadas.
\end{proof}

\begin{proposition}[Registro dodecafásico del estado completo anterior a \(\alpha\)]
\label{prop:c27-registro-u-k-producido}
El estado holonómico terminal posee la doble publicación
\[
x_{\mathrm{term}}
\xrightarrow{(\pi_{\mathrm{term}},R_{\mathrm{sgn}})}
(B_{\mathrm{term}},U),
\]
donde
\[
U=(2378,1406,2479,-452,998,-551,
-668,-204,-371,-322,-28,-997).
\]
Al ordenar las coordenadas por las tres órbitas de paso tres, la transformada
\[
\mathcal H_{12}
=P_{\mathrm{orb}}^{-1}(H_4\oplus H_4\oplus H_4)P_{\mathrm{orb}},
\qquad H_4^2=4I_4,
\]
posee inversa \(\mathcal H_{12}/4\) sobre su imagen integral y determina
unívocamente
\[
K=(234,543,140,729,659,824,621,058,914,794,146,601).
\]
Por tanto, dentro del estado holonómico integral,
\[
x_{\mathrm{term}}\xrightarrow{R_{\mathrm{sgn}}}U
\xrightarrow{\mathcal H_{12}^{-1}}K
\]
es una composición exacta anterior al acarreo que publica \(\alpha\). La matriz
visible \(B_{\mathrm{term}}\) y el registro firmado \(U\) son dos
coordenadas del mismo estado; la segunda conserva los campos de orientación,
oposición, cociente y registro que la primera contrae. No existe aquí una
segunda APP ni un perfil reducido competidor: APP produce las hojas y el
registro aritmético, TRIT los orienta y TPK acumula en su ledger la ventana
firmada que \(R_{\mathrm{sgn}}\) proyecta. Exigir que APP aislada publique
una coordenada que pertenece al codominio posterior del TPK constituye un
error de tipo. La composición completa no recibe \(U\), \(K\), sus
diferencias, \(\alpha\) ni expansiones objetivo como entradas.
\end{proposition}

\begin{theorem}[Clausura bidireccional posterior del carácter \(\alpha\)]
\label{thm:c27-alpha-segunda-via-bidireccional}
Sea \(\alpha=\alpha\) la coordenada ya publicada por el acarreo
dodecafásico. La evaluación del jet de vacancias de grados \(7{:}9\) en esa
salida determina el funcional
\[
\begin{aligned}
B_9(\alpha):={}&\log_{10}\!\left(\frac{10}{9}\right)
+\frac74\alpha^2-\frac{\alpha^4}{20}
+\frac{2\alpha^5}{21}+\frac{\alpha^6}{46}\\
&+\frac{\alpha^7}{120}-\frac{\alpha^8}{45}
-\frac{2\alpha^9}{495}.
\end{aligned}
\]
Al elevar el giro a la incógnita positiva \(T\), el resolvente equivale a
\[
\boxed{\alpha T^2-B_9(\alpha)T+\alpha^3=0.}
\]
Con
\(\xi=\operatorname{arcosh}(B_9(\alpha)/(2\alpha^2))\), las soluciones son
\[
T_\pm=\alpha e^{\pm\xi},
\qquad T_+T_-=\alpha^2,
\qquad
J_\alpha(T)=\frac{\alpha^2}{T},
\qquad J_\alpha(T_\pm)=T_\mp.
\]
La cámara \(6<T<7\) selecciona \(T_+\); la clausura finita
\(\pi_9=T_+/2\) es compatible, en la resolución declarada, con el canal de
clausura que \(\Gamma_9\) prolonga independientemente hasta
\(\pi\). La identidad
\(\alpha=\sqrt{T_+T_-}\) fija la media geométrica de las ramas conjugadas.
El jet \(7{:}9\), las dos raíces y la involución constituyen una clausura
bidireccional exacta \emph{sobre} el carácter ya generado. Este teorema no
recibe una constante convencional, pero tampoco constituye una segunda
derivación independiente de \(\alpha\): su dominio contiene
\(\alpha\), y su resultado es un certificado funcional posterior.
\end{theorem}

\begin{theorem}[Publicación excepcional correlacionada con el sector \(\alpha\)]
\label{thm:c27-alpha-rama-excepcional}
En la carta incidencial del mismo registro dodecafásico, las raíces radiales
de \(A_2^{12}\) determinan
\[
v_\alpha=4\rho_1+\rho_2+\cdots+\rho_{12},
\qquad v_\alpha^2=54,
\]
y el vecino de índice tres
\[
N_{\alpha,0}
=\{x\in N(C_W):\langle x,v_\alpha\rangle\equiv0\pmod3\},
\qquad
N_\alpha=N_{\alpha,0}+\mathbb Z\frac{v_\alpha}{3}.
\]
El retículo \(N_\alpha\) es par, unimodular, de rango veinticuatro y sin
raíces, de modo que \(N_\alpha\cong\Lambda_{24}\). La cadena
\[
C_W\longrightarrow N(A_2^{12})
\xrightarrow{v_\alpha}\Lambda_{24}
\longrightarrow V^\natural
\xrightarrow{\operatorname{Aut}}\mathbb M
\]
expone la bisagra excepcional. El escalar \(\alpha\) y el
vector \(v_\alpha\) pertenecen a codominios distintos y comparten el estado
dodecafásico de origen. El resultado exacto de esta rama es la construcción
\(N_\alpha\cong\Lambda_{24}\to V^\natural\) y la acción posterior de
\(\operatorname{Aut}(V^\natural)\cong\mathbb M\) sobre \(V^\natural\).
No se identifica el escalar con el vector ni se deduce de la correlación una
acción del Monstruo sobre cifras o rutas; tal promoción exigiría un morfismo
equivariante con dominio y codominio explícitos.
\end{theorem}

\begin{proposition}[Estatuto de los transportes \(L_0,L_1\)]
\label{prop:c27-estatuto-l0-l1}
Los endomorfismos \(L_0,L_1\in\operatorname{End}(\mathbb F_3^6)\) y el
calendario \((L_0,L_0,L_1,L_1)\) pertenecen a la ley de transporte del estado
holonómico integral. Las seis transiciones publicadas de rango completo fijan
cada matriz de manera única por \(L=X^{-1}Y\) en \(\mathbb F_3\), y el
calendario es el único de los \(16\) calendarios binarios que conserva
simultáneamente las tres biografías estructurales. Su función es prolongar
los representantes ya seleccionados; la evaluación arquimediana y la
nomenclatura convencional se aplican después.
\end{proposition}

El dominio de esta proposición es la única estructura \(\mathfrak H_9\).
Sus proyecciones de olvido son lectores posteriores y no constituyen perfiles
generadores alternativos. Esta tipificación impide confundir una sombra
modal con el estado que la produce.

\begin{theorem}[Profundidad arbitraria y determinación de las expansiones completas]
\label{thm:c27-profundidad-arbitraria}
Para cada
\(\chi\in\{\pi,e,
\varphi\}\) y cada \(N\ge1\), existe un único prefijo
decimal \(p_{\chi,N}\) de longitud \(N\), producido por un truncamiento
finito del estado holonómico integral, tal que
\[
 \rho_{N+1,N}(p_{\chi,N+1})=p_{\chi,N}.
\]
La familia compatible es la familia de proyecciones de un único elemento
\[
 p_{\chi,\infty}
 \in\varprojlim_N\{0,\ldots,9\}^{N}
 \cong\{0,\ldots,9\}^{\mathbb N},
\]
y la intersección de sus cilindros arquimedianos consta de un único real.
Por consiguiente, la conexión genera la expansión completa; cada precisión
finita es una proyección de esa salida infinita.
\end{theorem}

\begin{proof}
Para una precisión \(N\), las horquillas racionales de clausura y
propagación y los cocientes consecutivos de Perron se refinan hasta quedar
contenidos en un único cilindro decimal de diámetro \(10^{-N}\). La
aplicación \(\operatorname{Ext}\), el cociclo de acarreo y el truncamiento
TPK hacen que el cilindro de orden \(N+1\) se proyecte sobre el de orden
\(N\). Los cilindros están encajados y sus diámetros tienden a cero; su
límite inverso determina un único punto. Puesto que \(N\) es arbitrario y
la relación de extensión carece de profundidad terminal, la conclusión vale
para toda longitud finita y, conjuntamente, para la expansión infinita.
\end{proof}

\begingroup\fontsize{10}{11.4}\selectfont
\setlength{\parskip}{1.5pt}\setlength{\abovedisplayskip}{4pt}\setlength{\belowdisplayskip}{4pt}
La ejecución de \(396\) niveles, \(2376\) trits, \(43\) vueltas,
\(387\) pares \((K,K+9)\) por canal y \(1000\) cifras verifica una instancia
material extensa del teorema. La ejecución de \(10,000\) cifras verifica
otra instancia. El cuantificador \(\forall N\) y la compatibilidad
\(\rho_{N+1,N}p_{N+1}=p_N\), junto con la existencia del elemento
\(p_\infty\) del límite inverso, demuestran la salida decimal infinita; los
cálculos finitos son certificados reproducibles de sus proyecciones.

La coordenada \(\alpha\) pertenece al mismo estado límite.
Su vía entera se prolonga mediante \(\mathsf T_{\alpha,\infty}\), mientras
que la raíz simple de \(E_{21/22}\) proporciona un segundo lector interno.
La estabilidad comprobada a \(10,000\) cifras certifica la propagación
remota del acarreo; el objeto matemático que ambos lectores publican queda
fijado antes de cualquier comparación metrológica.

\begin{proposition}[Prolongación completa de la vía entera de \(\alpha\)]
\label{prop:c27-alpha-prolongacion-completa}
Sea \(B=1000\), sean \((P_k),(E_k),(\Phi_k)\) las tres sucesiones de bloques
publicadas por el estado coinductivo y sea \((K_k)\) la sucesión
dodecafásica prolongada. Las cuatro series
\[
 p=\sum_{k\geq1}P_kB^{-k},\qquad
 e_0=\sum_{k\geq1}E_kB^{-k},\qquad
 \phi_0=\sum_{k\geq1}\Phi_kB^{-k},\qquad
 \kappa=\sum_{k\geq1}K_kB^{-k}
\]
convergen absolutamente y determinan
\[
 \alpha_A=p+e_0-\phi_0-\kappa.
\]
La recurrencia de acarreo es la escritura en base \(B\) de esta igualdad.
Para cada \(M\), un truncamiento suficientemente profundo determina las
primeras \(M\) cifras de la expansión canónica de \(\alpha_A\); el error de
la cola está acotado por una constante multiplicada por \(B^{-N}\). Por
consiguiente, \(\mathsf T_{\alpha,\infty}\) publica una sucesión decimal
completa, y sus ejecuciones a \(3000\) o \(10\,000\) cifras son proyecciones
finitas de esa salida.
\end{proposition}

\begin{proof}
Cada bloque pertenece a un conjunto entero acotado, de modo que las cuatro
series están dominadas por una serie geométrica. La suma parcial hasta \(N\)
difiere de la suma completa en \(O(B^{-N})\). La división euclídea de
derecha a izquierda transporta exactamente esa cola mediante el acarreo;
por unicidad de la expansión canónica, todo prefijo que queda separado de
una frontera de cilindro se estabiliza al aumentar \(N\), y en una frontera
se aplica la convención canónica no terminal. Esto define simultáneamente
todas las cifras, no una sucesión de objetivos numéricos independientes.
\end{proof}

\begingroup\fontsize{10.2}{11.7}\selectfont
\setlength{\parskip}{2pt}\setlength{\abovedisplayskip}{5pt}\setlength{\belowdisplayskip}{5pt}
\begin{corollary}[Separación entre generación y reconocimiento]
\label{cor:c27-generacion-reconocimiento}
Las identidades de Machin, la serie exponencial, el rayo de Perron y la carta
decimal son realizaciones exactas de caracteres previamente seleccionados
por \(\mathfrak H_9\). Los nombres convencionales \(\pi,e,\varphi,\alpha\)
y las tablas metrológicas reconocen las salidas; no determinan órbitas,
transportes, acarreos ni prefijos.
\end{corollary}
 

\par\endgroup
\par\endgroup
\chapter{Construcción del generador elíptico \texorpdfstring{\(i=\sqrt{-1}\)}{i=sqrt(-1)} y geometría de fase}
\begingroup
\renewcommand{\chapter}[2][]{}%
\chapter{Construcción interna de \(\sqrt{-1}\): rotación de orden \(4\), operador complejo y geometría de fase}
\label{chap:p3-final:32:raiz_menos_uno}
\label{chap:p3-83-26-raiz_menos_uno}

\section*{La raíz interna de \texorpdfstring{\(-1\)}{-1} y la familia exponencial trítica}

La raíz de \(-1\) no se introduce como símbolo formal sin antecedente. Las
seis unidades módulo nueve se coordinatizan sobre
\(\mathbb P^1(\mathbb F_5)\), y el carácter cuadrático de
\(\mathbb F_5\) determina la matriz de conferencia de Paley. Su reducción
ternaria es
\[
 A_W=
 \begin{pmatrix}
 0&1&1&1&1&1\\
 1&0&1&2&2&1\\
 1&1&0&1&2&2\\
 1&2&1&0&1&2\\
 1&2&2&1&0&1\\
 1&1&2&2&1&0
 \end{pmatrix}\in M_6(\mathbb F_3).
\]
La identidad de conferencia reduce a
\begin{equation}\label{p3-83-c26-s1:eq:constantes-raiz-interna-menos-uno}
 \boxed{A_W^2=2I_6=-I_6.}
\end{equation}
Por tanto \(A_W\) es un operador de cuarto de giro en el espacio ternario.
Después de extender escala a \(\mathbb F_9\), se descompone en los espacios
propios conjugados de las raíces de \(X^2+1\). La fuente integral común
\[
 \mathcal Q_{\mathbb Z}=\mathbb Z[u]/(u^2+1)
\]
posee una realización finita \(u\mapsto A_W\) y una realización real
\(u\mapsto J_{\rm e}\). La fuente coordina la relación cuadrática; no
identifica cuerpos de características distintas.

Los tres regímenes de TRIT se realizan mediante operadores
\[
 J_{+1}^2=-I,\qquad J_0^2=0,\qquad J_{-1}^2=I.
\]
Separando las potencias pares e impares de la exponencial se obtiene una
sola familia:
\begin{equation}\label{p3-83-c26-s1:eq:constantes-euler-extendido}
 \boxed{
 \exp(tJ_\tau)=C_\tau(t)I+S_\tau(t)J_\tau,
 \qquad \tau\in\{+1,0,-1\},}
\end{equation}
donde
\[
 (C_{+1},S_{+1})=(\cos,\sin),\qquad
 (C_0,S_0)=(1,t),\qquad
 (C_{-1},S_{-1})=(\cosh,\sinh).
\]
Sus tres especializaciones rectoras son
\begin{equation}\label{p3-83-c26-s1:eq:constantes-euler-tres-ramas}
 \exp(\pi J_{+1})+I=0,\qquad
 \exp(J_0)=I+J_0,\qquad
 \exp\!\bigl(2(\log\varphi)J_{-1}\bigr)=F_{\rm av}^{\,2}.
\end{equation}
La rama elíptica contiene el cierre circular complejo; la parabólica, el
transporte de umbral; la hiperbólica, la autoescala áurea. La ecuación de
Euler aparece así como una especialización de una familia que coordina giro,
frontera y expansión hiperbólica.


\section{Realización elíptica de \(J^2=-I\) mediante la fibración de Hopf y la holonomía de Möbius}

\subsection{Estructura compleja y fibración de Hopf}

Sea \(S^3\subset\mathbb C^2\) y
\[
h(z_0,z_1)=
\bigl(2\Re(z_0\overline z_1),
2\Im(z_0\overline z_1),|z_0|^2-|z_1|^2\bigr)
\in S^2.
\]
En el toro de Hopf
\[
T_\eta=\{(z_0,z_1):|z_0|=\cos\eta,
|z_1|=\sin\eta\},
\]
las curvas
\[
\gamma_+(t)=(\cos\eta\,e^{it},\sin\eta\,e^{it}),
\]
\[
\gamma_-(t)=(\cos\eta\,e^{it},\sin\eta\,e^{-it})
\]
poseen comportamientos distintos: \(\gamma_+\) es una fibra de Hopf y
\(\gamma_-\) se proyecta con grado dos. Bajo proyección estereográfica, las
clases homológicas \((1,1)\) y \((1,-1)\) producen las dos familias diagonales
de Villarceau.

La conjugación compleja total
\[
K(z_0,z_1)=(\overline z_0,\overline z_1)
\]
preserva cada familia no orientada y revierte su orientación. No intercambia
las dos familias. La conjugación parcial puede intercambiarlas, pero no es una
aplicación compleja lineal ni el antiunitario estándar global. Esta corrección
es esencial para no fabricar una simetría partícula--antipartícula a partir del
dibujo toroidal.

\subsection{Doble recubrimiento y retorno spinorial}

La aplicación \(SU(2)\to SO(3)\) es un doble recubrimiento. Una vuelta de
\(2\pi\) en \(SO(3)\) puede levantar a \(-I\) en \(SU(2)\); la vuelta de
\(4\pi\) recupera \(I\). La estructura de signo no convierte un círculo en
banda de Möbius. Para obtener una banda se necesita un fibrado real de línea
cuyo carácter de holonomía sea \(-1\).

Sea \(L_\chi\to S^1\) el fibrado asociado a
\[
\chi:\pi_1(S^1)\longrightarrow\{\pm1\},
\qquad \chi(1)=-1.
\]
Su espacio total es una banda de Möbius. La ruta \(C_M=-\operatorname{rev}\)
aporta el candidato de signo; el círculo Villarceau aporta el lazo; el
entrelazador físico debe probar que el primero es la holonomía del segundo.

\subsection{Realización toroidal del operador \texorpdfstring{\(B_{\mathrm{mem}}^{\varepsilon}\)}{Bmem} y determinación del ángulo de Villarceau}

Con \(B_c=7I+2S\), defínase el operador orientado de memoria
\[
B_{\rm mem}^{\varepsilon}=5I+2\varepsilon S,
\qquad\varepsilon=\pm1.
\]
Para \(\varepsilon=+1\),
\[
\Spec(B_{\rm mem}^{+})=\{7,3\},
\qquad\det B_{\rm mem}^{+}=21.
\]
Interpretar los valores \((7,3)\) como bordes ecuatoriales exterior e interior
en una carta positiva común produce
\[
R+r=7\ell,\qquad R-r=3\ell,
\]
y por tanto
\[
R=5\ell,\qquad r=2\ell,\qquad
\sin\beta=\frac rR=\frac25.
\]
El ángulo Villarceau es
\[
\beta=\arcsin\frac25
=23.5781784782\ldots^\circ.
\]

El teorema algebraico \(\{7,3\}\leftrightarrow(5,2)\) es exacto. La premisa
geométrica que interpreta el espectro como los dos bordes de un toro es una
interfaz declarada. No debe ocultarse ni emplearse el ángulo como blanco.

La realización del bloque completo \(B_c\) conserva otra lectura:
\(R:r=7:2\) y \(\beta=\arcsin(2/7)\). No son valores rivales; pertenecen a
dos funtores: bloque total y modo de memoria. La prioridad física del segundo
debe decidirse por el entrelazador con la dinámica, no por cercanía decimal con
otra constante.

\subsection{Acción del operador \texorpdfstring{\(J^2=-I\)}{J2=-I} sobre las coordenadas angulares conjugadas}

El ángulo nativo de HMT es una fase \(u\in\mathbb R/\mathbb Z\). Los lectores
publican
\[
\theta_{\rm rad}=2\pi u,
\qquad
\theta_{\rm deg}=360u.
\]
Radianes y grados son cartas distintas del mismo elemento circular. El radián
no es más «ontológico» por ser adimensional; su naturalidad procede de la
longitud de arco. Los grados proporcionan otra coordenada absoluta de la
vuelta. Cualquier relación con \(\alpha^{-1}\), \(\varphi\) o CKM debe
formularse primero como cociente o holonomía invariante y sólo después
evaluarse en una carta angular.

La identidad exacta ya disponible
\[
\tanh^2(2\log\varphi)=\frac59
\]
conecta el canal Perron--Fibonacci con el cociente espectral de \(B_c\). Su
complemento es \(4/9\), el salto normalizado. Esta relación es estructural.
No implica que todo ángulo construido con \(\varphi\) sea un ángulo físico.

\subsection{Brouwer, Möbius y grado}

El teorema del punto fijo de Brouwer, las transformaciones fraccionarias de
Möbius y la banda de Möbius son tres objetos distintos. Se relacionan sólo
después de fijar dominio, acción y fibrado. Las órbitas elípticas de adelanto y
retardo pueden modelarse como dos secciones de un fibrado con holonomía de
signo; la transformación de Möbius describe la acción conformal en la carta;
el teorema de Brouwer garantiza un punto fijo para mapas continuos del disco.
Ninguna de estas frases sustituye el diagrama de realización.

\begin{remark}[Resultado y frontera]
Se demuestran el grado doble de la curva Hopf, la orientación de la
conjugación total y la construcción algebraica \(7,3\mapsto5,2\). La lectura
toroidal (5{:}2), la holonomía Möbius y la interpretación
partícula--antipartícula forman una cadena de interfaces que debe compartir un
único entrelazador. El volumen no fuerza ese entrelazador mediante una
coincidencia angular.
\end{remark}
  \endgroup

\chapter{Prolongación Hensel--Witt de \texorpdfstring{\(-1/12\)}{-1/12} y memoria modular}
\begingroup
\renewcommand{\chapter}[2][]{}%
\chapter{Derivación Hensel--Witt de \(-1/12\): prolongación compatible, regularización y memoria modular}
\label{chap:p3-final:33:menos_un_doceavo}
\label{chap:p3-83-27-menos_un_doceavo}

\section{Dominio Hensel--Witt, mapas de Teichmüller y objetivo racional \(-1/12\)}

El NTHW es el paquete de datos y mapas
\[
\NTHW=
\bigl(
\mathscr R_\pi^{\rm micro},
\Gnine,
\Cdoce,
K,
\mathcal W_{12\to24},
\mathcal I_{6\to8}
\bigr),
\]
con las compatibilidades siguientes:
\begin{enumerate}
  \item las regiones Hensel comparten una palabra cociente y conservan
  distinta memoria;
  \item \(\Gnine\) transporta la orientación y produce monodromía;
  \item \(\Cdoce\) organiza el observable firmado y el sello \(K\);
  \item la lectura Witt convierte soportes del mismo borde en hexadas y
  estrellas;
  \item la extensión \(W_{12}\to W_{24}\) produce la interfaz
  hexada--octada;
  \item los lectores arquimediano, excepcional y dimensional reciben
  invariantes del mismo estado.
\end{enumerate}

\begin{figure}[!htbp]
\centering
\resizebox{\textwidth}{!}{\begin{tikzpicture}[font=\small,node distance=1.2cm and 1cm,
  nodebox/.style={draw,rounded corners=2mm,align=center,minimum width=3cm,minimum height=1.05cm},
  arr/.style={-{Stealth[length=2.4mm]},thick}]
  \node[nodebox,draw=hmtblue,fill=hmtlightblue] (hensel) at (-5,2.3) {cilindros henselianos\\cociclo de acarreo y \(\MonNine\)};
  \node[nodebox,draw=hmtblue,fill=hmtlightblue] (regions) at (-5,-.1) {cinco regiones de \(\pi\)\\regiones de \(e,\varphi\)};
  \node[nodebox,draw=hmtgold,fill=hmtlightgold,minimum width=3.8cm,minimum height=2.05cm,font=\bfseries] (core) at (0,1.1) {\(\mathscr T_{\rm HW}\)\\Correspondencias\\Hensel--Witt};
  \node[nodebox,draw=hmtviolet,fill=hmtlightviolet] (c12) at (0,-1.8) {\(\Cdoce\), \(\Usgn\), \(K\)\\diferencias \(D_3,D_4,Q\)};
  \node[nodebox,draw=hmtgreen,fill=hmtlightgreen] (witt) at (5,2.3) {Witt \(W_{12}\)\\soportes y \(M_{12}\)};
  \node[nodebox,draw=hmtgreen,fill=hmtlightgreen] (oct) at (5,-.1) {inclusión de incidencia \(6\to8\)\\canales \(90/120\), normalización \(-1/12\)};
  \node[nodebox,draw=hmtred,fill=hmtlightred] (readers) at (5,-2.5) {lectores coordinados\\valor y simetría};
  \draw[arr,hmtblue] (hensel)--(core);
  \draw[arr,hmtblue] (regions)--(core);
  \draw[arr,hmtgold] (core)--(c12);
  \draw[arr,hmtgreen] (core)--(witt);
  \draw[arr,hmtgreen] (core)--(oct);
  \draw[arr,hmtred] (c12)--(readers);
  \draw[arr,hmtred] (witt.east)--++(1.9,0)|-(readers.east);
  \draw[arr,hmtred] (oct)--(readers);
  \node[align=center,hmtgray,font=\footnotesize] at (0,-3.45) {Compatibilidad de cilindros, dodecafase e incidencias sobre un mismo estado holonómico integral.};
\end{tikzpicture}
 }
\caption{El NTHW como interfaz tipada. La tarea matemática central es hacer
conmutar el diagrama, no acumular coincidencias alrededor de un nombre.}
\label{p3-83-c27-s1:fig:nthw}
\end{figure}

El NTHW da un nombre formal al lugar que antes se describía como «Piedra
Rosetta». La sustitución es conceptual: una piedra traduce por semejanza; una
transducción especifica entrada, salida, dato conservado y pérdida de
información.

\section{Aplicación de Teichmüller del código ternario al anillo de Witt}

La frontera reversible de APP selecciona, con calibres declarados, una matriz
\(A_W\) sobre \(\FF_3\) que satisface
\[
A_WA_W^T=2I\pmod3.
\]
El código gráfico
\[
C_W=\{(q,qA_W):q\in\FF_3^6\}
\]
es el Golay ternario extendido. Su enumerador de pesos es
\[
1+264y^6+440y^9+24y^{12}.
\]
Las \(264\) palabras de peso seis, identificadas por signo, dan \(132\)
soportes. Cada subconjunto de cinco puntos pertenece a una única hexada:
\[
W_{12}=S(5,6,12).
\]

El campo de \(495\) involuciones de estrella genera un grupo de orden
\(95\,040\), reconocido como \(M_{12}\). Una estrella aislada no genera el
grupo. El marco HMT ordenado tiene estabilizador trivial y por ello fija un
calibre dentro de la acción.

\section{Incidencia hexada--octada y canales modulares \(90/120\)}

Sea \(H\) una hexada y \(O_H\) una octada que la contiene. Al marcar un punto
y un par de puntos de la hexada se obtienen
\[
F_6(H)=H\times\binom H2,
\qquad
|F_6|=6\binom62=90.
\]
Si el punto marcado puede recorrer la octada:
\[
F_8(H)=O_H\times\binom H2,
\qquad
|F_8|=8\binom62=120.
\]
La inclusión \(H\hookrightarrow O_H\) induce \(F_6\hookrightarrow F_8\).
Su defecto normalizado es
\[
\boxed{
\frac{90-120}{24\binom62}
=\frac{6-8}{24}
=-\frac1{12}.}
\]
Este teorema combinatorio es anterior a la lectura angular de Barbero. Los
dos objetos comparten la interfaz \(6\to8\), pero no son idénticos.

\section{4 rutas compatibles hacia \texorpdfstring{\(-1/12\)}{-1/12}}

La construcción contiene cuatro realizaciones que deben exponerse juntas y no
contarse como cuatro pruebas independientes de una misma procedencia.

\subsection{Realización angular mediante el incremento elemental dodecafásico}

Un diente de \(\Cdoce\) mide \(30^\circ\). Entonces
\[
\frac{30}{120}-\frac{30}{90}
=\frac14-\frac13
=-\frac1{12}.
\]
Es la sombra angular del patrón \(90/120\).

\subsection{Pseudoinversa orientada del operador de Gram sobre el subespacio superviviente}

En el salto certificado \(60\to81\), dos fibras tienen tamaños \(0\) y
\(12\). El operador de Gram es
\[
K_{60,81}=\operatorname{diag}(0,12).
\]
Sobre el subespacio vivo, su pseudoinverso vale \(1/12\); con la orientación
negativa de borde,
\[
E_{\rm surv}=-K_{60,81}^{+}
\]
produce \(-1/12\).

\subsection{Extracción de escala por 2.ª diferencia triádica}

La realización mediante un extractor de escala triádico es la construcción por
segunda diferencia desarrollada íntegramente en
\cref{sec:c27-segunda-diferencia-triadica}. Esta primera comparecencia fija su
función dentro del inventario de realizaciones; la residencia indicada conserva
las definiciones de \(T_L\), \(a(L)\), \(A_1(L)\), \(C(L)\), la independencia
del residuo respecto de \(L\) y el estatuto exacto de la normalización.

\subsection{Realización modular mediante el cociente \(\eta(\tau)/\eta(3\tau)\) y la función \(t_3\)}

La función eta satisface
\[
\frac{\eta(\tau)}{\eta(3\tau)}
=q^{-1/12}\prod_{n\ge1}(1+q^n+q^{2n})^{-1}.
\]
El exponente procede de \(1/24-3/24=-1/12\). Al tomar doce copias:
\[
t_3(q)=\left(\frac{\eta(\tau)}{\eta(3\tau)}\right)^{12}
=q^{-1}-12+54q-76q^2-243q^3+\cdots.
\]
El coeficiente \(54\) coincide con el defecto APP. El corredor
\[
j=\frac{(t_3+27)(t_3+243)^3}{t_3^3}
\]
conecta después con la rama modular clásica.

\begin{remark}[Qué se ha demostrado en el corredor]
Son exactos el defecto APP \(54\), el exponente eta \(-1/12\), el defecto
hexada--octada y las identidades modulares una vez fijado \(t_3\). La
afirmación fuerte es que todos sean imágenes naturales de un único operador
TPK. Esa naturalidad debe demostrarse mediante el diagrama del NTHW; no se
deduce sólo de repetir los mismos enteros.
\end{remark}


\section{5 realizaciones algebraicas de \texorpdfstring{\(-1/12\)}{-1/12} y separación de sus dominios}

Además del defecto normalizado de la inclusión hexada--octada, aparecen
cinco realizaciones operatoriales del mismo racional. Sus dominios se
especifican por separado; la igualdad de valores no identifica los
operadores sin un entrelazador adicional.

\subsection{Diferencia de periodicidades dodecafásicas}

Un diente de \(\Cdoce\) mide \(30^\circ\). Entonces
\[
\frac{30}{120}-\frac{30}{90}
=\frac14-\frac13
=-\frac1{12}.
\]
Es la diferencia orientada entre los periodos de \(90^\circ\) y
\(120^\circ\) en esta realización.

\subsection{Pseudoinversa del operador de Gram sobre la componente viva}

En el salto exacto \(60\to81\), dos fibras tienen tamaños \(0\) y
\(12\). El operador de Gram es
\[
K_{60,81}=\operatorname{diag}(0,12).
\]
Sobre la componente no nula de su rango, la pseudoinversa vale \(1/12\); con la orientación
negativa de borde,
\[
E_{\rm surv}=-K_{60,81}^{+}
\]
produce \(-1/12\).

\subsection{2.ª diferencia entre escalas triádicas}
\label{sec:c27-segunda-diferencia-triadica}

Definamos
\[
T_L=\sum_{n=1}^{L-1}n(L-n)=\frac{L(L^2-1)}6,
\qquad
a(L)=-\frac{T_L}{L^2}=-\frac L6+\frac1{6L}.
\]
La primera resta de escala es
\[
A_1(L)=a(L)-\frac{a(3L)}3=\frac4{27L},
\]
y la segunda
\[
C(L)=LA_1(L)-\frac{3L}{3}A_1(3L)=\frac8{81}.
\]
El residuo \(8/81\) es independiente de \(L\). Para obtener
\(-1/12\) se aplica la normalización declarada
\[
R(L)=-\frac{27}{32}C(L)=-\frac1{12}.
\]
El cálculo deriva \(8/81\); no deriva por sí solo el factor \(-27/32\).

\subsection{Cadena modular}

La función eta satisface
\[
\frac{\eta(\tau)}{\eta(3\tau)}
=q^{-1/12}\prod_{n\ge1}(1+q^n+q^{2n})^{-1}.
\]
El exponente procede de \(1/24-3/24=-1/12\). Al tomar doce copias:
\[
t_3(q)=\left(\frac{\eta(\tau)}{\eta(3\tau)}\right)^{12}
=q^{-1}-12+54q-76q^2-243q^3+\cdots.
\]
El coeficiente \(54\) coincide con el defecto APP\@. La identidad racional
\[
j=\frac{(t_3+27)(t_3+243)^3}{t_3^3}
\]
establece la conexión con la rama modular clásica.

\Needspace{32\baselineskip}
\subsection{Defecto de los proyectores de diferencia cíclica}

Sea \(S\) el desplazamiento cíclico sobre \(\QQ^{12}\) y definamos
\[
 D_r=I-S^r.
\]
El núcleo de \(D_r\) está formado por los vectores constantes en cada órbita
de \(S^r\); por tanto,
\[
 \dim\ker D_r=\gcd(12,r).
\]
Sean \(\Pi_{\ker D_3}\) y \(\Pi_{\ker D_4}\) los proyectores racionales sobre
\(\ker D_3\) y \(\ker D_4\), y sea
\(\tau(T)=\operatorname{Tr}(T)/12\) la traza normalizada. Entonces
\[
 \boxed{
 \tau(\Pi_{\ker D_3}-\Pi_{\ker D_4})
 =\frac{\operatorname{rank}\Pi_{\ker D_3}
       -\operatorname{rank}\Pi_{\ker D_4}}{12}
 =\frac{3-4}{12}
 =-\frac1{12}.}
\]
Esta realización expresa el racional como defecto transversal de los pasos
tres y cuatro en el ciclo de doce fases.

\begin{remark}[Dominios de las realizaciones]
Son exactos el defecto APP \(54\), el exponente eta \(-1/12\), el defecto
combinatorio de la elevación hexada--octada, y las
identidades modulares una vez fijado \(t_3\). También son exactos el valor
de la pseudoinversa, el extractor normalizado y el defecto de proyectores.
Cada igualdad conserva su dominio y su aplicación correspondiente.
\end{remark}
  \endgroup

\chapter{Caracteres cuadráticos y constantes espectrales: Catalán, Apéry, Euler--Mascheroni y Stieltjes}
\begingroup
\renewcommand{\chapter}[2][]{}%
\chapter{Caracteres cuadráticos y constantes espectrales: Catalán, Apéry y torre de Euler--Stieltjes}
\label{chap:p3-final:41:catalan_apery_euler}
\label{ch:spectral}

\section{Representaci\'on operatoria de los caracteres de Dirichlet}

Las constantes de este cap\'itulo se agrupan por compartir una construcci\'on
espectral independiente de sus expansiones decimales y libre de valores
objetivo. Sea

\[
Ne_n=ne_n,\qquad n\in\N,
\]

en \(\ell^2(\N)\). Un car\'acter residual \(\chi\) define el operador diagonal \(Q_\chi e_n=\chi(n)e_n\); siempre que la potencia sea traza--clase,

\begin{equation}
\Tr(N^{-s}Q_\chi)=L(s,\chi).
\label{eq:L-trace}
\end{equation}

La ecuaci\'on conserva el m\'odulo, la orientaci\'on y la procedencia de cada
car\'acter; por tanto, no identifica objetos residuales distintos.

\section{Construcci\'on HMT de los caracteres residuales}

El operador diagonal anterior no constituye el dato algebraico inicial. La
cadena de selecci\'on es

\[
\begin{aligned}
\APP&\longrightarrow\TRIT\longrightarrow\TPK
\longrightarrow\{\text{residuo m\'odulo }3,\ C_P\}\\
&\longrightarrow\{\chi_{-3},\ J=C_P\bmod3\}
\longrightarrow\{\chi_{-3},\ J^2=-I,\chi_{-4}\}\\
&\longrightarrow Q_3,Q_4.
\end{aligned}
\]

El funcional residual tri\'adico determina la diferencia entre las dos clases
orientadas no nulas. En el sector de incidencia se emplea la matriz de conferencia
sim\'etrica de Paley de orden seis,
\[
C_P^{\mathsf T}=C_P,\qquad C_PC_P^{\mathsf T}=5I_6.
\]
Su reducci\'on \(J=C_P\bmod3\) cumple
\[
J^2=2I_6=-I_6
\quad\text{en }\operatorname{End}_{\mathbb F_3}(\mathbb F_3^6).
\]
El cuarto de torsi\'on determina entonces la fase c\'iclica
\(1,\ii,-1,-\ii\). La aplicaci\'on del teorema chino del resto conserva
ambos datos residuales y produce \(Q_{12}=Q_3Q_4\). La traza
\(\Tr(N^{-s}Q)\) se eval\'ua en la etapa terminal; ni la serie de Dirichlet
ni su expresi\'on decimal intervienen en la selecci\'on del car\'acter.

La construcci\'on se explicita mediante las tablas
\begin{equation}
\chi_{-3}(n)=
\begin{cases}
0,&3\mid n,\\
1,&n\equiv1\pmod3,\\
-1,&n\equiv2\pmod3,
\end{cases}
\qquad
\chi_{-4}(n)=
\begin{cases}
0,&2\mid n,\\
1,&n\equiv1\pmod4,\\
-1,&n\equiv3\pmod4.
\end{cases}
\label{eq:characters34-tables}
\end{equation}
En \(\ell^2(\N)\),
\[
Q_3e_n=\chi_{-3}(n)e_n,
\qquad Q_4e_n=\chi_{-4}(n)e_n.
\]
Como son diagonales, conmutan; el teorema chino del resto da
\begin{equation}
(Q_3Q_4)e_n=\chi_{-3}(n)\chi_{-4}(n)e_n
=\chi_{12}(n)e_n.
\label{eq:q12-product-proof}
\end{equation}
El producto conserva simult\'aneamente el residuo tri\'adico y el cuarto de
torsi\'on. Por consiguiente, el car\'acter m\'odulo doce queda determinado
antes de efectuar cualquier evaluaci\'on de sus valores \(L\).

\begin{remark}[Control negativo]
La cadena queda invalidada si \(Q_3\) no distingue las dos clases orientadas,
si \(J^2\ne-I\), si \(Q_3Q_4\) no tiene la tabla CRT declarada o si el
operador se selecciona a partir del valor de la constante.
\end{remark}

\section{Estructura compleja interna y constante de Catalan}

La estructura compleja interna satisface

\begin{equation}
J^2=-I.
\label{eq:J-square}
\end{equation}

La igualdad \(J^2=-I\) define una representaci\'on
\(\rho:C_4\to\operatorname{Aut}_{\mathbb F_3}(\mathbb F_3^6)\),
\(\rho(j)=J\). Para realizarla espectralmente se toma el car\'acter unitario
\(\upsilon:C_4\to U(1)\), \(\upsilon(j)=\ii\), y el mapa de fase
\(n\mapsto j^n\). Sobre \(\ell^2(\N)\), la representaci\'on diagonal inducida
es
\[
U_4e_n=\upsilon(j^n)e_n=\ii^ne_n.
\]
Su parte orientada es
\[
Q_4=\frac{U_4-U_4^*}{2\ii},\qquad
Q_4e_n=\sin\!\left(\frac{\pi n}{2}\right)e_n
=\chi_{-4}(n)e_n.
\]
Este es el levantamiento expl\'icito entre el generador finito de orden cuatro y el operador
residual; no se identifica \(\mathbb F_9\) con \(\mathbb C\).

Sobre enteros impares, el mismo car\'acter puede escribirse

\[
\chi_J(n)=\ii^{n-1}=\chi_{-4}(n).
\]

Equivalentemente,

\begin{equation}
Q_4=\frac{U_4-U_4^*}{2\ii},
\qquad Q_4e_n=\chi_{-4}(n)e_n.
\label{eq:Q4}
\end{equation}

La constante de Catalan aparece como

\begin{equation}
G=\Tr(N^{-2}Q_4)=L(2,\chi_{-4})
=\sum_{n\ge0}\frac{(-1)^n}{(2n+1)^2}.
\label{eq:catalan}
\end{equation}

El levantamiento unitario del cuarto de torsi\'on genera exactamente el
car\'acter cuyo momento cuadr\'atico es la constante de Catalan.

\begin{proof}[Prueba operatoria]
El operador \(N^{-2}Q_4\) es traza--clase porque
\(\sum_{n\ge1}|\chi_{-4}(n)|n^{-2}<\infty\). Su traza es, por tanto, la
suma de los elementos diagonales. La tabla \cref{eq:characters34-tables}
anula los pares y alterna los impares:
\[
\Tr(N^{-2}Q_4)
=1-\frac1{3^2}+\frac1{5^2}-\frac1{7^2}+\cdots.
\]
La identidad con \(L(2,\chi_{-4})\) es la definici\'on de la serie de
Dirichlet; la igualdad con \(\operatorname{Im}\operatorname{Li}_2(\ii)\)
se obtiene tomando la parte imaginaria de
\(\sum_{n\ge1}\ii^n/n^2\).
\end{proof}

El valor terminal es
\[
G=0.9159655941772190150546035149\ldots.
\]
La estructura determinante es la identidad \(J^2=-I\), que induce
exactamente el signo alternante de las clases impares; la expansi\'on decimal
es su evaluaci\'on terminal.

\section{Composici\'on residual tri\'adico--cuaternaria}

El funcional residual m\'odulo tres proporciona \(\chi_{-3}\). Su producto
CRT con \(\chi_{-4}\) define

\begin{equation}
\chi_{12}=\chi_{-3}\chi_{-4}.
\label{eq:chi12}
\end{equation}

Sobre las unidades m\'odulo doce,

\[
\chi_{12}(1)=1,\quad
\chi_{12}(5)=-1,\quad
\chi_{12}(7)=-1,\quad
\chi_{12}(11)=1,
\]

y se anula fuera de ellas. Dos evaluaciones exactas son

\begin{align}
L(1,\chi_{12})&=\frac{\log(2+\sqrt3)}{\sqrt3},
\label{eq:L1chi12}\\
L(2,\chi_{12})&=\frac{\pi^2}{6\sqrt3}.
\label{eq:L2chi12}
\end{align}

La primera contiene la unidad hiperb\'olica

\begin{equation}
\sqrt3L(1,\chi_{12})=\arcosh2=\log(2+\sqrt3).
\label{eq:pell2}
\end{equation}

Junto con \(\log(3+2\sqrt2)=2\log(1+\sqrt2)\) y
\(\log(30+\sqrt{899})\), esta identidad determina una familia de escalas
de Pell cuyos miembros conservan sus respectivas procedencias algebraicas.

Las evaluaciones terminales certificadas son

\[
L(1,\chi_{12})=0.7603459963009463475310942549\ldots,
\]

\[
L(2,\chi_{12})=0.9497031262940093952634984917\ldots.
\]

\begin{proof}[Evaluaciones dodecaf\'asicas]
Agrupando por clases residuales,
\begin{align*}
L(1,\chi_{12})
&=\sum_{n\ge0}\left(
\frac1{12n+1}-\frac1{12n+5}
-\frac1{12n+7}+\frac1{12n+11}\right)\\
&=\int_0^1
\frac{1-x^4-x^6+x^{10}}{1-x^{12}}\,dx
=\int_0^1\frac{1-x^4}{1+x^6}\,dx.
\end{align*}
Una descomposici\'on real en los factores cuadr\'aticos de \(1+x^6\)
produce
\(\log(2+\sqrt3)/\sqrt3\). Para el segundo momento, si
\(\psi_1\) es la trigamma,
\begin{align*}
L(2,\chi_{12})
&=\frac1{144}\bigl[
\psi_1(1/12)-\psi_1(5/12)
-\psi_1(7/12)+\psi_1(11/12)\bigr]\\
&=\frac{\pi^2}{144}\left(
\csc^2\frac\pi{12}-\csc^2\frac{5\pi}{12}\right)
=\frac{\pi^2}{6\sqrt3},
\end{align*}
donde se ha usado
\(\psi_1(z)+\psi_1(1-z)=\pi^2\csc^2(\pi z)\). Las dos pruebas emplean
la tabla CRT y no un valor decimal objetivo.
\end{proof}

\section{Constantes de Euler y Stieltjes y transformada theta--Mellin}

Sea
\[
\vartheta(t)=\sum_{n\in\Z}e^{-\pi n^2t},\qquad t>0.
\]
Para \(\Re s>1\), el funcional theta--Mellin recupera la identidad convergente

\[
\pi^{-s/2}\Gamma\!\left(\frac{s}{2}\right)\zeta(s)
=\frac12\int_0^\infty
\bigl(\vartheta(t)-1\bigr)t^{s/2-1}\,dt.
\]

La continuaci\'on meromorfa se obtiene al partir la integral en \(t=1\), usar
\(\vartheta(t)=t^{-1/2}\vartheta(1/t)\) y separar los t\'erminos polares. Su expansi\'on en \(s=1\),

\begin{equation}
\zeta(s)=\frac1{s-1}+\sum_{j\ge0}
\frac{(-1)^j\gamma_j}{j!}(s-1)^j,
\label{eq:stieltjes}
\end{equation}

sit\'ua \(\gamma=\gamma_0\) y las constantes de Stieltjes en una misma
familia de momentos renormalizados. Esta organizaci\'on expresa una
procedencia anal\'itica com\'un, sin atribuir una interpretaci\'on f\'isica
id\'entica a todos sus niveles.

 La constante de Catalan suele aparecer como una serie alternada:\par

\[
G
=
1-\frac1{3^2}+\frac1{5^2}-\frac1{7^2}+\cdots.
\]

Esa fórmula da su valor y la genealogía HMT explica por qué existe esa alternancia concreta.\par

En HMT, el origen está en el cuarto de torsión.\par

Existe un operador interno \(J\) tal que\par

\[
J^2=-I.
\]

La raíz de \(-1\) emerge de dos aplicaciones sucesivas de un cuarto de giro. La estructura compleja es la memoria algebraica de una rotación interna.\par

Ese giro genera un carácter de orden cuatro. Cuando se publica sobre los enteros positivos, produce exactamente el patrón residual que alterna entre \(+1\), \(0\), \(-1\), \(0\).\par

La constante de Catalan aparece entonces como traza espectral:\par

\[
G
=
L(2,\chi_{-4})
=
\operatorname{Tr}(N^{-2}Q_4).
\]

Esto cambia su significado. Catalan es simultáneamente la suma de una serie particular y la respuesta espectral de los enteros ante un cuarto de torsión.\par

Cada término de la serie registra cómo el giro interno actúa sobre una clase residual. La alternancia es la sombra aritmética de la rotación.\par

La estructura triádica posee su propio carácter módulo \(3\). Al combinarlo con el carácter de orden cuatro, el teorema chino del resto produce un carácter módulo \(12\):\par

\[
\chi_{12}
=
\chi_{-3}\chi_{-4}.
\]

El módulo \(12\) aparece como el espacio común mínimo en el que pueden coexistir la orientación triádica y el cuarto de torsión.\par

La dodecafase es, en ese sentido, una reconciliación de \(3\) y \(4\).\par

Del carácter combinado emergen\par

\[
L(1,\chi_{12})
=
\frac{\log(2+\sqrt3)}{\sqrt3},
\]

y\par

\[
L(2,\chi_{12})
=
\frac{\pi^2}{6\sqrt3}.
\]

La cantidad\par

\[
2+\sqrt3
\]

es la unidad fundamental del campo real \(\mathbb Q(\sqrt3)\). Su logaritmo representa un desplazamiento hiperbólico. De nuevo, una estructura finita de residuos abre una dinámica multiplicativa infinita.\par

Aquí se encuentran giro, torsión, carácter aritmético, unidad de Pell y desplazamiento hiperbólico.\par

La constante de Euler–Kronecker del campo \(\mathbb Q(\sqrt3)\) prolonga esta estructura. Mide el término regular que permanece tras separar el polo principal de la función zeta del campo. En la genealogía HMT, este término constituye el invariante regular de memoria aritmética fina que subsiste después de retirar la divergencia universal.\par

La familia Bendersky–Glaisher prolonga el mecanismo en una torre. Las derivadas espectrales de zeta generan esa torre. Apéry y otras constantes asociadas a valores impares se organizan como miembros de una misma familia funcional.\par

La aportación novedosa reside en la genealogía que reúne las constantes clásicas:\par

\begin{itemize}[leftmargin=*,itemsep=0.35em]
\item el cuarto de torsión produce la estructura compleja;
\item el residuo triádico produce la estructura módulo \(3\);
\item su composición produce el carácter módulo \(12\);
\item el carácter genera valores \(L\);
\item los valores \(L\) abren el campo cuadrático y sus constantes espectrales.
\end{itemize}

La torsión geométrica se ha convertido en aritmética conservando íntegramente su orientación.\par
  \endgroup

\chapter{Euler--Kronecker, Meissel--Mertens y distribuciones primas en la doble lectura de las hojas}
\begingroup
\renewcommand{\chapter}[2][]{}%
\chapter{Constantes de Euler--Kronecker y distribuciones primas: caracteres dodecafásicos, Meissel--Mertens y series de primos gemelos}
\label{chap:p3-final:42:euler_kronecker_primos}
\section{Constante de Euler--Kronecker del car\'acter dodecaf\'asico}

La factorizaci\'on de Dedekind

\begin{equation}
\zeta_{\mathbb Q(\sqrt3)}(s)=\zeta(s)L(s,\chi_{12})
\label{eq:dedekind-sqrt3}
\end{equation}

da, al comparar los t\'erminos constantes en \(s=1\),

\begin{equation}
\gamma_{\mathbb Q(\sqrt3)}
=\gamma+\frac{L'(1,\chi_{12})}{L(1,\chi_{12})}
=1.0539656082484825800\ldots.
\label{eq:euler-kronecker}
\end{equation}

Esta constante representa el t\'ermino finito normalizado en el polo de la
funci\'on zeta del campo real dodecaf\'asico. No constituye una variante de
\(\gamma\) de Euler, sino la correcci\'on aritm\'etica inducida por
\(\chi_{12}\).

En efecto, escribiendo \(u=s-1\),
\[
\zeta(s)=u^{-1}+\gamma+O(u),
\qquad
L(s,\chi_{12})=L_1+L_1'u+O(u^2),
\]
se obtiene
\begin{equation}
\zeta_{\mathbb Q(\sqrt3)}(s)
=\frac{L_1}{u}+\bigl(\gamma L_1+L_1'\bigr)+O(u).
\label{eq:ek-laurent-derivation}
\end{equation}
La constante de Euler--Kronecker es el cociente entre el t\'ermino finito y
el residuo; dividir \cref{eq:ek-laurent-derivation} por \(L_1\) prueba
\cref{eq:euler-kronecker}. Esta derivaci\'on explica simult\'aneamente el
residuo
\[
\operatorname*{Res}_{s=1}\zeta_{\mathbb Q(\sqrt3)}(s)
=\frac{\log(2+\sqrt3)}{\sqrt3}
\]
y la correcci\'on finita. El valor num\'erico se certifica evaluando
\(L(1,\chi_{12})\) y \(L'(1,\chi_{12})\) con la misma tabla residual y una
cota de cola; no se importa de una tabla de campos cuadr\'aticos.

\section{Constantes aritm\'eticas asociadas a los n\'umeros primos}

Una vez seleccionados los ciclos aritm\'eticos irreducibles, se obtienen dos
familias distintas:

\begin{align}
B_1&=\lim_{x\to\infty}\left(
\sum_{p\le x}\frac1p-\log\log x
\right),
\label{eq:meissel-mertens}\\
C_2&=\prod_{p>2}\frac{p(p-2)}{(p-1)^2}.
\label{eq:twin-prime-constant}
\end{align}

\(B_1\) regulariza una suma local; \(C_2\) es una densidad multiplicativa de pares gemelos. Ambos requieren control de cola. No se deducen el uno del otro por compartir los primos.

  \endgroup

\chapter{Productos regulares y familia Bendersky--Glaisher--Kinkelin}
\begingroup
\renewcommand{\chapter}[2][]{}%
\chapter{Familia Bendersky--Glaisher--Kinkelin: productos regulares, derivadas de la función zeta y jerarquías de normalización}
\label{chap:p3-final:43:bendersky_glaisher}
\section{Generación de la familia Bendersky--Glaisher--Kinkelin mediante regularización zeta}

Sea \(\mathcal Z_3(s):=\zeta(s)\) la realizaci\'on escalar del funcional
zeta tri\'adico en la normalizaci\'on de este volumen,
\(H_k=\sum_{j=1}^k j^{-1}\) (con \(H_0=0\)) y \(B_{k+1}\) el n\'umero de
Bernoulli de convenio \(B_1=-1/2\). Definimos

\begin{equation}
A_k=\exp\!\left(
\frac{H_kB_{k+1}}{k+1}-\mathcal Z_3'(-k)
\right),\qquad k\ge0.
\label{eq:bendersky}
\end{equation}

Los primeros niveles incluyen

\[
A_0=\sqrt{2\pi},
\]

y la constante de Glaisher--Kinkelin en \(A_1\). Para \(n\ge1\),

\begin{equation}
\zeta(2n+1)=(-1)^{n+1}
\frac{2(2\pi)^{2n}}{(2n)!}\log A_{2n}.
\label{eq:odd-zeta-bendersky}
\end{equation}

La constante de Ap\'ery, \(\zeta(3)\), es el primer nivel par no trivial.
Su presencia en el gas fot\'onico del \cref{ch:metrology} constituye una
realizaci\'on f\'isica posterior de este valor espectral.

La primera parte de la torre puede verse de forma sinóptica:
\begin{center}
\small
\begin{tabular}{@{}cll@{}}
\toprule
\(k\) & identidad estructural & valor terminal\\
\midrule
0 & \(A_0=\exp[-\zeta'(0)]=\sqrt{2\pi}\)
  & \(2.5066282746310005024\ldots\)\\
1 & \(\log A_1=1/12-\zeta'(-1)\)
  & \(1.2824271291006226369\ldots\)\\
2 & \(\log A_2=-\zeta'(-2)=\zeta(3)/(4\pi^2)\)
  & \(1.0309167521973921\ldots\)\\
3 & \(\log A_3=-11/720-\zeta'(-3)\)
  & \(0.9795555269428446\ldots\)\\
4 & \(\log A_4=-\zeta'(-4)\)
  & \(0.9920479745250403\ldots\)\\
\bottomrule
\end{tabular}
\end{center}

\begin{proof}[Relaci\'on con los valores impares]
La ecuaci\'on funcional de \(\zeta\), desarrollada en \(s=-2n\), da
\begin{equation}
\zeta'(-2n)=(-1)^n
\frac{(2n)!}{2(2\pi)^{2n}}\zeta(2n+1).
\label{eq:zeta-derivative-even-negative}
\end{equation}
Como \(B_{2n+1}=0\) para \(n\ge1\), la definici\'on
\cref{eq:bendersky} reduce \(\log A_{2n}=-\zeta'(-2n)\). Despejar
\(\zeta(2n+1)\) prueba \cref{eq:odd-zeta-bendersky}. El \`indice \(2n\)
conserva el nivel de derivaci\'on: no es una etiqueta editorial.
\end{proof}

  \endgroup

\chapter{Dinámica de Gauss y constantes de Khinchin, Lévy y Lochs}
\begingroup
\renewcommand{\chapter}[2][]{}%
\chapter{Dinámica de Gauss y constantes de Khinchin, Lévy y Lochs: cofinalidad de cilindros, operador de transferencia y espectro de Gauss--Kuzmin--Wirsing}
\label{chap:p3-final:44:gauss_khinchin_levy_lochs}
\label{ch:gauss}

\section{Cofinalidad entre cilindros HMT y cilindros de fracciones continuas}

Una secci\'on HMT irracional determina una familia anidada de cilindros compactos. Su evaluaci\'on arquimediana es un punto \(x\in(0,1)\setminus\mathbb Q\). La fracci\'on continua de \(x\) determina cilindros cl\'asicos

\[
I(a_1,\ldots,a_n).
\]

\begin{proposition}[Cofinalidad]
Para todo cilindro finito \(I(a_1,\ldots,a_n)\) que contiene la evaluaci\'on \(x\), existe profundidad HMT \(m(n)\) tal que el cilindro HMT de nivel \(m\ge m(n)\) est\'a contenido en \(I(a_1,\ldots,a_n)\).
\end{proposition}

\begin{proof}
Fijado \(n\), el cilindro de fracci\'on continua \(I(a_1,\ldots,a_n)\) es un entorno abierto de la irracional \(x\). Existe entonces \(\varepsilon_n>0\) con \((x-\varepsilon_n,x+\varepsilon_n)\subset I(a_1,\ldots,a_n)\). Los cilindros HMT \(C_m\) contienen \(x\), son anidados y satisfacen \(\operatorname{diam}C_m\to0\). Elija \(m(n)\) con \(\operatorname{diam}C_{m(n)}<\varepsilon_n\). Como \(x\in C_{m(n)}\), todo punto de ese cilindro dista de \(x\) menos de \(\varepsilon_n\), luego \(C_{m(n)}\subset I(a_1,\ldots,a_n)\). El anidamiento da la inclusi\'on para todo nivel posterior.
\end{proof}

La demostraci\'on utiliza compacidad, anidamiento y unicidad de la
evaluaci\'on. El resultado establece el transporte de aproximaciones, pero
no implica la conjugaci\'on del desplazamiento de memoria TPK con el mapa de
Gauss sobre todas las fibras profundas.

La genealog\'ia completa de esta realizaci\'on es
\[
\APP\to\TRIT\to\TPK
\to (C_m,\rho_{m+1,m})_{m\ge0}
\to x=\bigcap_m C_m
\to \bigl(I(a_1,\ldots,a_n)\bigr)_{n\ge1}.
\]
Los cilindros HMT conservan ruta, fase y acarreo; el mapa terminal proyecta
estas coordenadas sobre los prefijos de fracci\'on continua. La cofinalidad
demuestra que dicha proyecci\'on converge al mismo punto. Un entrelazador
din\'amico exigir\'ia adem\'as un levantamiento \(\widetilde T_G\) con
cuadrados conmutativos sobre las fibras de valoraci\'on, propiedad que no se
deduce de la cofinalidad.

\section{Medida de Gauss y observables de Khinchin y L\'evy}

El mapa de Gauss \cite{Khinchin,IosifescuKraaikamp}

\[
T_G(x)=\frac1x-\left\lfloor\frac1x\right\rfloor
\]

preserva

\begin{equation}
d\mu_G(x)=\frac{dx}{\log2(1+x)}.
\label{eq:gauss-measure}
\end{equation}

La constante de Khinchin se define mediante

\begin{equation}
\log K=\int_0^1\log a_1(x)\,d\mu_G(x).
\label{eq:khinchin}
\end{equation}

La distribuci\'on del primer d\'igito se calcula directamente:
\begin{equation}
\mu_G(a_1=k)
=\frac1{\log2}\int_{1/(k+1)}^{1/k}\frac{dx}{1+x}
=\log_2\!\left(\frac{(k+1)^2}{k(k+2)}\right).
\label{eq:gauss-digit-law}
\end{equation}
Por el teorema erg\'odico,
\begin{equation}
K=\prod_{k\ge1}
k^{\,\log_2((k+1)^2/[k(k+2)])}
=2.685452001065306445309714835\ldots
\label{eq:khinchin-product}
\end{equation}
para \(\mu_G\)-casi toda secci\'on. El producto no se aplica a cada punto:
la secci\'on \(\varphi^{-1}\), por ejemplo, tiene todos sus d\'igitos iguales
a uno y es un criterio de refutación inmediato de una lectura punto a punto. Una
certificaci\'on de \cref{eq:khinchin-product} trunca la suma de logaritmos y
acota la cola mediante la estimaci\'on
\[
\mu_G(a_1=k)=O(k^{-2}),
\qquad
\sum_{k>N}\mu_G(a_1=k)\log k
=O\!\left(\frac{\log N+1}{N}\right).
\]

La constante de L\'evy es

\begin{equation}
L=-\int_0^1\log x\,d\mu_G(x)
=\frac{\pi^2}{12\log2}.
\label{eq:levy}
\end{equation}

En efecto,

\[
\frac1{1+x}=\sum_{n\ge0}(-1)^nx^n,\qquad
\int_0^1(-\log x)x^n\,dx=\frac1{(n+1)^2},
\]

por lo que la integral es \(\eta(2)/\log2=\pi^2/(12\log2)\).

La constante de Khinchin determina la media logar\'itmica de los cocientes
parciales, mientras que la constante de L\'evy determina la tasa exponencial
de crecimiento de los denominadores. Ambas son observables distintos de un
mismo sistema medido.

\section{Entrop\'ia m\'etrica, exponente de Lyapunov y constante de Lochs}

Como

\[
\log|T_G'(x)|=-2\log x,
\]

el exponente de Lyapunov y la entrop\'ia m\'etrica son

\begin{equation}
\chi_G=h_G=2L=\frac{\pi^2}{6\log2}.
\label{eq:gauss-entropy}
\end{equation}

La raz\'on de Lochs entre fracciones continuas y una expansi\'on en base \(b\) es

\begin{equation}
\Lambda_b=\frac{\log b}{h_G}
=\frac{6\log2\log b}{\pi^2}.
\label{eq:lochs}
\end{equation}

En particular,

\begin{align*}
\Lambda_{10}&=0.9702701143920339257402560192\ldots,\\
\Lambda_{729}&=\frac{36\log2\log3}{\pi^2}
=2.777618966374305777705782976\ldots.
\end{align*}

La base \(729=3^6\) es relevante para el ambiente \(\mathbb F_3^6\), pero \cref{eq:lochs} sigue siendo un teorema casi seguro de \(\mu_G\); no una igualdad punto a punto para toda ruta HMT.

El valor \(\Lambda_{729}\) representa la tasa asint\'otica t\'ipica de
cocientes parciales obtenidos por bloque completo de seis trits. No identifica el
ambiente de \(729\) palabras con los otros objetos de cardinal \(729\) del
corpus. La igualdad es una tasa de informaci\'on entre particiones, no una
bijecci\'on entre estados.

\section{Operador de transferencia de Gauss--Kuzmin--Wirsing}

El operador de Gauss--Kuzmin--Wirsing act\'ua por

\begin{equation}
(\mathcal L f)(x)=\sum_{n\ge1}\frac1{(x+n)^2}
f\!\left(\frac1{x+n}\right).
\label{eq:gkw-operator}
\end{equation}

Su autovalor dominante corresponde a la densidad invariante. La constante GKW es el m\'odulo del autovalor subdominante sobre el subespacio de media cero. Un esquema HMT compatible usa particiones de \(9^m\) cilindros, matrices de transferencia con intervalos y una cota de cola.

Para promover el c\'alculo se deben certificar: aislamiento del autovalor, invariancia del subespacio, error de truncaci\'on, convergencia de proyectores espectrales y ausencia de otro autovalor de igual m\'odulo. El decimal conocido no puede introducirse como centro del intervalo inicial.

Una deformaci\'on conjunta contiene Khinchin y L\'evy en un solo objeto:
\begin{equation}
(\mathcal L_{s,t}f)(x)
=\sum_{n\ge1}n^t(n+x)^{-2s}
 f\!\left(\frac1{n+x}\right),
\qquad
P(s,t)=\log r(\mathcal L_{s,t}).
\label{eq:gauss-pressure-operator}
\end{equation}
En el punto \((s,t)=(1,0)\), la perturbaci\'on espectral da
\begin{equation}
\partial_tP(1,0)=\log K,
\qquad
-\partial_sP(1,0)=\chi_G=2L.
\label{eq:pressure-derivatives}
\end{equation}
La Hessiana de \(P\) determina las varianzas y covarianzas asint\'oticas de
\(\log a_1\) y \(-2\log x\). Por tanto, las constantes de Khinchin y
L\'evy son dos derivadas transversales de la misma funci\'on de presi\'on.

El reconocimiento num\'erico del autovalor subdominante es
\[
\lambda_2\simeq-0.3036630028987326586,\qquad
\kappa_{\rm GKW}=|\lambda_2|,\qquad
-\log\kappa_{\rm GKW}\simeq1.19183673556055.
\]
En este volumen, dichos decimales constituyen valores de referencia para el
contraste y no un certificado espectral.
La prueba HMT propuesta usa Galerkin o Ulam sobre los cilindros cofinales de
profundidad \(9^m\), una desigualdad de Lasota--Yorke, aritm\'etica de
intervalos para aislar \(\lambda_2\) y una cota rigurosa de la cola. Un
intervalo que contenga dos autovalores del mismo m\'odulo refuta la
promoci\'on, aunque no afecte a Khinchin, L\'evy o Lochs.

 \section{Certificado cofinal del espectro de Gauss--Kuzmin--Wirsing sobre la jerarquía nonádica}
\label{sec:p3-final:cierre-gkw}

Para Gauss--Kuzmin--Wirsing, sea \(P_N\) la familia de truncamientos
certificada sobre el espacio funcional declarado y definamos
\[
 \Pi_m=P_{9^m}.
\]
La jerarquía \((\Pi_m)\) es una subsecuencia cofinal; las cotas uniformes, la
convergencia de proyectores y la ausencia de contaminación espectral se
heredan de la familia \((P_N)\). Una discretización cilíndrica distinta deberá
aportar su entrelazador con \((P_{9^m})\).
  \endgroup

\chapter{Criticidad unimodal, constantes de Feigenbaum y renormalización hiperbólica}
\begingroup
\renewcommand{\chapter}[2][]{}%
\chapter{Criticidad unimodal de la familia dodecafásica: perfil analítico exacto, aproximantes de Feigenbaum y reducción de la universalidad al certificado hiperbólico}
\label{chap:p3-final:45:feigenbaum}
\label{mab:rm:ch:chaos}

\section{Construcci\'on intr\'inseca de la familia dodecaf\'asica}

Las constantes de Feigenbaum caracterizan la universalidad de la
duplicaci\'on de periodo. Su derivaci\'on en HMT exige construir previamente
una familia anal\'itica unimodal, demostrar su criticidad cuadr\'atica y
estudiar posteriormente la acci\'on del operador de renormalizaci\'on.

El mapa inicial no procede de una elecci\'on arm\'onica externa. El retorno
de \(\TPK\) determina la rueda dodecaf\'asica orientada
\(U_{12}^{\rm sign}\). En el sector uniforme \(\eta=0\), la traza
normalizada de sus doce sectores fija \(w_m=1/12\), y el funcional de fase
\begin{equation}
\mathcal P_{\rm crit}:U_{12}^{\rm sign}\longrightarrow\R/2\pi\Z,
\qquad
m\longmapsto \frac{(30m-10)\pi}{180}
=\pi\frac{3m-1}{18}
\label{mab:rm:eq:critical-phase-reader}
\end{equation}
conserva orden y orientaci\'on. La composici\'on causal es
\[
\APP\to\TRIT\to\TPK\to U_{12}^{\rm sign}
\xrightarrow{\mathcal P_{\rm crit}}
\{(\phi_m,w_m)\}_{m=1}^{12}
\to u_0\to f_\lambda.
\]
Toda modificaci\'on de \(\mathcal P_{\rm crit}\), de los pesos o del orden de
las fases altera la familia y constituye un criterio de refutaci\'on. Estas
elecciones deben fijarse con anterioridad al contraste con \(\delta_F\).

La familia dodecaf\'asica resultante parte de

\[
a_m=\frac{3m-1}{18},\qquad
\phi_m=\pi a_m,\qquad m=1,\ldots,12,
\]

y del segundo momento

\[
D_0=\frac1{12}\sum_{m=1}^{12}a_m^2.
\]

La identidad

\begin{equation}
\sum_{m=1}^{12}(3m-1)^2=5394=6\cdot899
\label{mab:rm:eq:chaos-sum2}
\end{equation}

produce

\[
D_0=\frac{899}{648},\qquad
s_0=\frac{36}{\pi\sqrt{899}}.
\]

Por tanto,

\begin{equation}
s_0\phi_m=\frac{2(3m-1)}{\sqrt{899}}.
\label{mab:rm:eq:pi-cancellation}
\end{equation}

El factor \(\pi\) se cancela antes de efectuar las iteraciones. La familia no
depende de \(\delta_F\) ni de \(\alpha_F\) como datos de entrada.

La cancelaci\'on es una propiedad estructural de la normalizaci\'on. Antes de
normalizar,
el segundo momento angular es
\[
\frac1{12}\sum_{m=1}^{12}\phi_m^2
=\pi^2D_0.
\]
La condici\'on de criticidad cuadr\'atica fija
\(s_0^2\pi^2D_0=2\), y por ello cada fase normalizada depende s\'olo de la
aritm\'etica entera de \(3m-1\). El funcional conserva la procedencia angular de
la rueda, pero la iteraci\'on recibe una familia real intr\'inseca. Esto es
precisamente lo que permite distinguir una derivaci\'on de una calibraci\'on
por \(\pi\).

\section{Representaci\'on anal\'itica exacta y expresi\'on trigonom\'etrica cerrada}

Definimos

\begin{equation}
u_0(x)=1-\frac1{12}\sum_{m=1}^{12}
\cos\!\left(\frac{2(3m-1)}{\sqrt{899}}x\right).
\label{mab:rm:eq:chaos-profile}
\end{equation}

El desarrollo en el origen es

\begin{equation}
u_0(x)=x^2-\frac{715769}{2424603}x^4+O(x^6).
\label{mab:rm:eq:chaos-taylor}
\end{equation}

La suma finita admite la siguiente expresi\'on cerrada. Si \(y=x/\sqrt{899}\),

\begin{equation}
u_0(x)=1-\frac{\cos(37y)\sin(36y)}{12\sin(3y)}
=1-\frac{\sin(73y)-\sin y}{24\sin(3y)}.
\label{mab:rm:eq:chaos-closed}
\end{equation}

La singularidad aparente en \(y=0\) es removible y reproduce \cref{mab:rm:eq:chaos-taylor}.

\begin{proof}[Derivaci\'on de la forma cerrada]
Los doce argumentos forman una progresi\'on aritm\'etica:
\[
2(3m-1)y=(6m-2)y,\qquad m=1,\ldots,12.
\]
La suma de cosenos de una progresi\'on satisface
\[
\sum_{m=1}^{N}\cos(a+md)
=\frac{\sin(Nd/2)}{\sin(d/2)}
 \cos\!\left(a+\frac{N+1}{2}d\right).
\]
Tomando \(N=12\), \(d=6y\) y \(a=-2y\) resulta
\[
\sum_{m=1}^{12}\cos((6m-2)y)
=\frac{\sin36y}{\sin3y}\cos37y.
\]
La identidad producto--a--suma
\(2\cos37y\sin36y=\sin73y-\sin y\) da la segunda expresi\'on de
\cref{mab:rm:eq:chaos-closed}.
\end{proof}

El coeficiente cu\'artico tampoco se ajusta. Las sumas enteras
\begin{equation}
\sum_{m=1}^{12}(3m-1)^2=5394,\qquad
\sum_{m=1}^{12}(3m-1)^4=4294614
\label{mab:rm:eq:chaos-moments24}
\end{equation}
insertadas en el desarrollo de \(\cos z\) producen exactamente
\[
\frac1{24\cdot12}\sum_m
\left(\frac{2(3m-1)}{\sqrt{899}}\right)^4
=\frac{715769}{2424603}.
\]
As\'i, tanto la normalizaci\'on cuadr\'atica como la primera correcci\'on no
lineal son momentos finitos de la misma rueda.

\section{Pertenencia a la clase anal\'itica unimodal cuadr\'atica}

Consideremos

\[
f_\lambda(x)=1-\lambda u_0(x),\qquad \lambda>0.
\]

Para \(0<x\le1\), todos los argumentos

\[
\frac{2(3m-1)}{\sqrt{899}}x
\]

pertenecen a \((0,\pi)\). Por ello,

\[
u_0'(x)>0,\qquad u_0'''(x)<0.
\]

Como \(u_0\) es par, \(x=0\) es su único punto cr\'itico y \(u_0''(0)=2\). El Schwarziano de \(f_\lambda\) es

\begin{equation}
S(f_\lambda)
=\frac{u_0'''}{u_0'}
-\frac32\left(\frac{u_0''}{u_0'}\right)^2<0
\qquad(0<|x|\le1).
\label{mab:rm:eq:negative-schwarzian}
\end{equation}

\begin{theorem}[Pertenencia exacta a la clase unimodal]
Para
\[
0<\lambda\le \frac{2}{u_0(1)},
\]
la familia \(f_\lambda=1-\lambda u_0\) es una autoaplicaci\'on real--anal\'itica de \([-1,1]\), par, con un único punto cr\'itico cuadr\'atico en el origen y Schwarziano negativo fuera de ese punto. Por tanto pertenece, sin usar valores de Feigenbaum como valor objetivo, a la clase \(S\)-unimodal cuadr\'atica sobre la que act\'ua el renormalizador de duplicaci\'on de periodo \cite{rm:Feigenbaum1978,rm:Lanford1982}. Para \(\lambda\) fuera de ese intervalo se conserva el germen anal\'itico y su forma local, pero no se afirma autom\'aticamente la autoaplicaci\'on del intervalo.
\end{theorem}

El teorema establece la contribuci\'on intr\'inseca de HMT previa al teorema
de universalidad: el germen anal\'itico y su tipo de criticidad.

\begin{proof}[Detalles del signo y de la autoaplicaci\'on]
Derivando \cref{mab:rm:eq:chaos-profile},
\begin{align*}
u_0'(x)&=\frac1{12}\sum_m k_m\sin(k_mx),\\
u_0'''(x)&=-\frac1{12}\sum_m k_m^3\sin(k_mx),
\qquad k_m=\frac{2(3m-1)}{\sqrt{899}}.
\end{align*}
Para \(0<x\le1\), \(0<k_mx<\pi\), luego todos los sumandos de la primera
l\'inea son positivos y los de la segunda negativos. La paridad extiende la
monoton\'ia a ambos lados del origen. Como \(u_0(0)=0\) y
\(0\le u_0(x)\le u_0(1)\), la cota
\(0<\lambda\le2/u_0(1)\) implica
\(-1\le1-\lambda u_0(x)\le1\). Finalmente,
\(u_0''(0)=1/12\sum k_m^2=2\), de modo que el punto cr\'itico es
cuadr\'atico.
\end{proof}

\section{Sucesiones finitas de par\'ametros superestables}

La sucesi\'on superestable certificada hasta nivel ocho produce

\[
\delta_8=4.66921218915\ldots,
\qquad
\alpha_8=2.50296847988\ldots.
\]

Su estatuto es \emph{Certificado finito libre de valores objetivo}: cada
aproximante se obtiene de la familia previamente definida, sin introducir
como dato el l\'imite universal. La coincidencia decimal proporciona un
indicador de convergencia, pero no sustituye el teorema de hiperbolicidad.

Una extensi\'on num\'erica independiente de la misma sucesi\'on, con paso
adaptado a la separaci\'on entre ra\'ices y control del periodo exacto,
alcanza provisionalmente:

\begin{center}
\small
\begin{tabular}{@{}ccll@{}}
\toprule
Nivel & \(\lambda_n\) & \(\delta_n\) & \(\alpha_n\)\\
\midrule
11 & 1.87403828195439908045 & 4.66920170694424982745 & 2.50290847517165181764\\
12 & 1.87403836411023460451 & 4.66920163036308158848 & 2.50290800379001546915\\
13 & 1.87403838170549870372 & 4.66920161361839219913 & 2.50290790268748193108\\
14 & 1.87403838547386545431 & 4.66920161007472591023 & 2.50290788100969754274\\
\bottomrule
\end{tabular}
\end{center}

Los niveles 11--14 conservan su carácter de evaluación numérica posterior al certificado persistente hasta nivel ocho. La extrapolación de Aitken proporciona $\lambda_\infty^{\HMT}\simeq1.8740383865008918080$ y $\alpha_F\simeq2.50290787509$. El teorema de \ref{sec:p3-final:cierre-feigenbaum} establece el escalado paramétrico de las primeras raíces; la evaluación espacial conserva su alcance propio.

Para evitar un procedimiento num\'erico no explicitado, el m\'etodo se formula como
sigue. Defina
\begin{equation}
F_n(\lambda)=f_\lambda^{\,2^n}(0).
\label{mab:rm:eq:superstable-equation}
\end{equation}
La ra\'iz superestable \(\lambda_n\) es la primera ra\'iz de \(F_n\) a la
derecha de \(\lambda_{n-1}\) que satisface simult\'aneamente
\[
f_{\lambda_n}^{\,2^j}(0)\ne0
\quad(0\le j<n).
\]
Esta desigualdad excluye los periodos divisores. Con intervalos disjuntos
\(I_n\ni\lambda_n\), se calculan
\begin{equation}
\delta_n=
\frac{\lambda_{n-1}-\lambda_{n-2}}
     {\lambda_n-\lambda_{n-1}},
\qquad
\alpha_n=-\frac{d_{n-1}}{d_n},
\label{mab:rm:eq:feigenbaum-approximants}
\end{equation}
donde \(d_n=f_{\lambda_n}^{\,2^{n-1}}(0)\) es la distancia orientada del
\'ultimo punto nuevo al punto cr\'itico. Un certificado persistente debe guardar
los intervalos de ra\'iz, la exclusi\'on de periodos menores y la propagaci\'on
de error de \cref{mab:rm:eq:feigenbaum-approximants}; un simple decimal impreso no
cumple esas tres condiciones.

La regla de discretizaci\'on constituye asimismo un control negativo. Una
cota inferior fija mayor que la separaci\'on
\(|\lambda_n-\lambda_{n-1}|\) puede omitir la primera ra\'iz y producir, no
obstante, una sucesi\'on num\'ericamente plausible. Por ello el incremento
debe escalar con la separaci\'on previa y cada candidato debe verificarse
mediante su periodo exacto.

\section{Condiciones de certificaci\'on del operador de renormalizaci\'on}

Sobre un espacio de funciones pares anal\'iticas, definimos

\begin{equation}
\mathcal R(f)(x)=-a(f)^{-1}f\bigl(f(-a(f)x)\bigr),
\qquad a(f)=-f(1).
\label{mab:rm:eq:renormalization}
\end{equation}

El cierre completo de \(\delta_F\) y \(\alpha_F\) requiere certificar conjuntamente:

\begin{enumerate}
\item existencia de un punto fijo \(g=\mathcal R(g)\) en una bola anal\'itica declarada;
\item unicidad local de ese punto fijo;
\item un único autovalor real inestable de \(D\mathcal R_g\), igual a \(\delta_F\);
\item radio espectral estrictamente menor que uno en el complemento estable;
\item transversalidad de la curva \(\lambda\mapsto f_\lambda\) respecto de la variedad estable de \(g\).
\end{enumerate}

Una implementaci\'on HMT natural emplea truncaciones de Chebyshev de orden
\(9^m\), aritm\'etica de intervalos y polinomios de radios. La prueba debe
proporcionar una bola, una inversa aproximada, una cota de defecto y una cota
de Lipschitz que satisfagan \(Y+Z(r)<r\). De este modo, la profundidad
non\'adica parametriza un certificado de punto fijo y no una mera sucesi\'on
de aproximaciones decimales.

En una base par de Chebyshev,
\[
g(x)=\sum_{j=0}^{N}g_{2j}T_{2j}(x),\qquad N=9^m,
\]
el residuo finito es \(F_N(g)=\Pi_N(\mathcal Rg-g)\). Si \(A_N\) aproxima
\((DF_N(\bar g))^{-1}\), las cantidades que deben certificarse son
\begin{align*}
Y&=\|A_NF_N(\bar g)\|,\\
Z_0&=\|I-A_NDF_N(\bar g)\|,\\
Z_1(r)&=\sup_{\|h\|\le r}
\|A_N(DF_N(\bar g+h)-DF_N(\bar g))\|.
\end{align*}
Una desigualdad
\(Y+(Z_0+Z_1(r))r<r\) encierra un punto fijo \`unico. El bloque espectral
posterior debe aislar una direcci\'on inestable y contraer el complemento. El
uso de \(9^m\) entra as\'i como ley de refinamiento del certificado, no como
decoraci\'on numerol\'ogica.

\begin{corollary}[Confluencia exacta del modo treinta]
El modo \(30\) admite tres realizaciones simult\'aneas y tipol\'ogicamente
distintas:
\begin{align*}
30&\longmapsto(3\cdot30,4\cdot30)=(90,120)
&&\text{en la completaci\'on constitutiva},\\
30&\longmapsto30^2-1=899=29\cdot31
&&\text{en la normalizaci\'on cr\'itica},\\
30&\longmapsto
M_{30}=\begin{pmatrix}30&899\\1&30\end{pmatrix}
&&\text{en la transformaci\'on hiperb\'olica de Pell}.
\end{align*}
Las tres identidades son exactas y se obtienen sin prescribir valores objetivo. El invariante conservado es
el entero central \(30\) junto con su orientaci\'on: extensi\'on \(3\to4\),
defecto cuadr\'atico \(-1\) y unidad de norma uno. Lo que no se afirma es la
igualdad entre el operador de vac\'io y el renormalizador de duplicaci\'on.
\end{corollary}

\begin{remark}[Control negativo]
Un segundo autovalor fuera del disco unidad, una bola de radios sin soluci\'on, p\'erdida de unicidad o tangencia de la familia HMT a la variedad estable refutan la promoci\'on universal. Ninguno refuta el germen exacto ni el teorema de Schwarziano.
\end{remark}

\begin{proposition}[Resultado principal]
La significaci\'on estructural de la constante de duplicaci\'on reside en la
cancelaci\'on dodecaf\'asica de \(\pi\), en la identidad
\(899=30^2-1\), en la forma trigonom\'etrica cerrada y en la pertenencia
demostrada a la clase anal\'itica sobre la que act\'ua el renormalizador. Su
expansi\'on decimal constituye exclusivamente la evaluaci\'on terminal.
\end{proposition}

\subsection*{Alcance lógico y dominio de validez}
Perfil, normalización, unidad de Pell, unimodalidad y Schwarziano conservan las demostraciones precedentes. Los ocho niveles certificados y las extrapolaciones de los niveles 11--14 conservan sus alcances numéricos respectivos. El desarrollo de \ref{sec:p3-final:cierre-feigenbaum} demuestra el escalado paramétrico de la sucesión original de primeras raíces en el sector uniforme, mediante identificación de centros y control del resto. Esta conclusión no atribuye al cociente espacial una demostración que corresponde al cociente paramétrico.

\subsection{Deformación armónica ponderada y selección de la rama superestable}
\label{mab:sec:extension-dodecafasica-ponderada-armonica}

El perfil uniforme posee una extensión paramétrica que conserva las doce
fases y separa la deformación armónica del parámetro dinámico. Para
\(1\le m\le12\), sean

\begin{equation}
 \phi_m=\frac{(30m-10)\pi}{180},
 \qquad
 p(\phi)=12\cos(3\phi)-\cos(4\phi).
 \label{mab:eq:fases-y-deformacion-armonica-armonizada}
\end{equation}

Para todo \(\eta\) real tal que los pesos permanezcan positivos, definimos

\begin{equation}
 w_m(\eta)=\frac{1+\eta p(\phi_m)}{12}.
 \label{mab:eq:pesos-dodecafasicos-armonizados}
\end{equation}

Las sumas de los armónicos tercero y cuarto sobre las doce fases se anulan;
por tanto,

\[
 \sum_{m=1}^{12}w_m(\eta)=1.
\]

Sea

\begin{equation}
 s_\eta^2=\frac{2}{\sum_{m=1}^{12}w_m(\eta)\phi_m^2},
 \qquad
 u_\eta(x)=1-\sum_{m=1}^{12}w_m(\eta)
 \cos(s_\eta\phi_m x).
 \label{mab:eq:familia-ponderada-armonizada}
\end{equation}

La expansión del coseno y la definición de \(s_\eta\) dan

\[
 u_\eta(x)=x^2+O(x^4)\qquad(x\to0).
\]

Por consiguiente, \(u_\eta\) es real--analítica, par y conserva un punto
crítico cuadrático en el origen. Introducimos el parámetro dinámico mediante

\begin{equation}
 f_{\lambda,\eta}(x)=1-\lambda u_\eta(x).
 \label{mab:eq:familia-feigenbaum-ponderada-armonizada}
\end{equation}

Un parámetro superestable de periodo exacto \(2^n\) satisface

\[
 f_{\lambda,\eta}^{\,2^n}(0)=0,
 \qquad
 f_{\lambda,\eta}^{\,2^k}(0)\ne0\quad(0\le k<n).
\]

La ecuación puede poseer varias ramas. Se fija una rama
\(\mathscr B\) recorriendo, inmediatamente después del parámetro del
nivel anterior, una malla adaptativa; se toma el primer intervalo con cambio
de signo y se aplica bisección. La notación
\(\lambda_n^{\mathscr B}(\eta)\) incorpora este selector y no atribuye
unicidad global a la ecuación superestable. El primer nivel se determina
exactamente:

\[
 \lambda_1^{\mathscr B}(\eta)=\frac1{u_\eta(1)},
\]

si \(u_\eta(1)>0\), porque
\(f_{\lambda,\eta}^{\,2}(0)=f_{\lambda,\eta}(1)=0\).
Dados tres niveles consecutivos y

\[
 x_n^{\mathscr B}(\eta)
 :=f_{\lambda_n^{\mathscr B}(\eta),\eta}^{\,2^{n-1}}(0),
\]

se forman

\begin{equation}
 \delta_n^{\mathscr B}(\eta)
 =\frac{\lambda_{n-1}^{\mathscr B}(\eta)-
         \lambda_{n-2}^{\mathscr B}(\eta)}
        {\lambda_n^{\mathscr B}(\eta)-
         \lambda_{n-1}^{\mathscr B}(\eta)},
 \qquad
 \alpha_n^{\mathscr B}(\eta)
 =\frac{|x_{n-1}^{\mathscr B}(\eta)|}
        {|x_n^{\mathscr B}(\eta)|}.
 \label{mab:eq:aproximantes-ponderados-armonizados}
\end{equation}

Para \(\eta=0\) se recupera exactamente el perfil uniforme. La
extensión ponderada introduce un parámetro transversal: permite estudiar
la transversalidad de la familia respecto de la variedad estable del
renormalizador sin convertir una coincidencia finita en una prueba de
universalidad.

\begin{remark}[Control negativo]
La extensión queda refutada si los pesos pierden positividad en el dominio
declarado, si la normalización deja de producir coeficiente cuadrático
unitario, si el selector no excluye periodos divisores o si una rama se elige
usando como entrada los valores universales terminales.
\end{remark}
 \section{Selección de primeras raíces y escalado paramétrico de Feigenbaum}
\label{sec:p3-final:cierre-feigenbaum}

El lector dodecafásico uniforme determina además una sucesión de primeras raíces de retorno crítico. Se demuestra que su cociente de separaciones converge a la constante de Feigenbaum, conservando la selección original, la orientación y la memoria de la construcción.

% Focal insertion; hierarchy inherited from the host.
\begingroup
\let\HMTFGHostSection\section
\edef\HMTFGSavedSecnumdepth{\number\value{secnumdepth}}
\setcounter{secnumdepth}{4}
\let\HMTFGHostSubsection\subsection
\let\HMTFGHostSubsubsection\subsubsection
\let\section\HMTFGHostSubsection
\let\subsection\HMTFGHostSubsubsection
\let\subsubsection\paragraph
% Fragmento maestro: incluir desde el anfitrion, con \HMTFGRoot fijada a la raiz del paquete.
% Requiere amsmath, amssymb/mathrsfs o unicode-math, amsthm, booktabs, longtable, hyperref, needspace.
% El anfitrion conserva sus entornos theorem, proposition, lemma y proof.
\providecommand{\HMTFGRoot}{.}
\begingroup
\providecommand{\ZZ}{}\renewcommand{\ZZ}{\mathbb Z}
\providecommand{\RR}{}\renewcommand{\RR}{\mathbb R}
\providecommand{\FF}{}\renewcommand{\FF}{\mathbb F}
\providecommand{\APP}{}\renewcommand{\APP}{\mathrm{APP}}
\providecommand{\TPK}{}\renewcommand{\TPK}{\mathrm{TPK}}
\providecommand{\TRIT}{}\renewcommand{\TRIT}{\{-1,0,+1\}}
% unicode-math usa lBrack/rBrack; newtxmath ya proporciona llbracket/rrbracket.
\providecommand{\llbracket}{\lBrack}
\providecommand{\rrbracket}{\rBrack}
% Fragmento de inserción; se incluye desde la raíz de este paquete editorial.
% Las fuentes completas y los cortes materiales constan en ANTECEDENTES_Y_DEPENDENCIAS.md.
% Procedencia: XI REV11 ES, sections/nucleo.tex, tpk_desarrollo_integrado.tex,
% ch_tres_vueltas_ext.tex, ch_elevacion_catalogal.tex, ch_cierre_cardinal.tex,
% ch_descenso_por_fibras.tex; perfil integral public_final/feigenbaum/c37_feigenbaum.tex;
% Grassmannianos REV07 ES, sections/jacobi_spectral.tex;
% TRANSVERSALIDAD_Y_ESCALADO_LOCAL_FEIGENBAUM.md, sección 14.
% Ampliación expositiva de demostraciones existentes; no introduce resultados nuevos.
\providecommand{\ZZ}{\mathbb Z}
\providecommand{\RR}{\mathbb R}
\providecommand{\FF}{\mathbb F}
\providecommand{\APP}{\mathrm{APP}}
\providecommand{\TPK}{\mathrm{TPK}}
\providecommand{\TRIT}{\{-1,0,+1\}}

\section{Antecedentes operatorios de la familia de criticidad}
\label{hmtfg:antecedentes-operatorios}

La familia de criticidad se construye a partir de una publicación del
calendario TPK. Las dos hojas aritméticas, la orientación trítica y el
transporte con memoria determinan primero los estados y sus prolongaciones.
La lectura ordenada del continuo permite después evaluar funciones
analíticas, comparar parámetros y seleccionar raíces. En este desarrollo
se conservan el estado genealógico y sus dos publicaciones: la carta
arquimediana proporciona el espacio de evaluación, mientras la publicación
de incidencia conserva las relaciones de frontera. La estructura discreta
conjunta del continuo mantiene sus cinco construcciones consustanciales.
El presente corte expositivo utiliza su lector ordenado y conserva el
antecedente común del que procede.

\subsection{Células aritméticas, orientación y transporte}
\label{hmtfg:app-trit-transporte}

% XI REV11, sections/nucleo.tex: definiciones de APP y TRIT;
% tpk_desarrollo_integrado.tex:22--165, transporte levantado y calendario.
Sean \(D_9=\{1,\ldots,9\}\) y \(X_9=(\ZZ/9\ZZ)^2\). Para
\(n\in\ZZ\), la división con representante positivo es
\[
 \rho_9(n)=1+((n-1)\bmod9),\qquad
 q_9(n)=\frac{n-\rho_9(n)}9,\qquad
 n=\rho_9(n)+9q_9(n).
\]
La célula de APP asociada a \((p,q)\in D_9^2\) es
\[
 a_{p,q}=(p,q;R^\Sigma,q^\Sigma,R^\Pi,q^\Pi),
 \quad
 p+q=R^\Sigma+9q^\Sigma,\qquad
 pq=R^\Pi+9q^\Pi.
\]
Los dos pares \((R^\Sigma,q^\Sigma)\) y \((R^\Pi,q^\Pi)\)
conservan respectivamente la evaluación aditiva y la multiplicativa.
La fórmula de división reconstruye ambas evaluaciones antes de cualquier
reducción posterior. La selección de una hoja activa se registra en una
coordenada del estado y conserva la célula completa.

El representante trítico balanceado \(r_3(n)\in\{-1,0,+1\}\) satisface
\(n\equiv r_3(n)\pmod3\). Su acarreo es
\[
 \chi(a,b)=\frac{a+b-r_3(a+b)}3,\qquad
 a,b\in\{-1,0,+1\}.
\]
La igualdad
\[
 \chi(a,b)+\chi(r_3(a+b),c)
 =\chi(b,c)+\chi(a,r_3(b+c))
\]
se obtiene escribiendo ambos miembros como
\((a+b+c-r_3(a+b+c))/3\). Esta identidad conserva el acarreo
al reagrupar una suma trítica. El selector operativo es
\[
 M_{\rm ph}(r)=
 \begin{cases}
 +1,&r\in\{1,4,7\},\\
 -1,&r\in\{2,5,8\},\\
 0,&r\in\{3,6,9\}.
 \end{cases}
\]
El valor \(+1\) activa la hoja aditiva, el valor \(-1\) la
multiplicativa y el valor \(0\) conserva la posición de ambos cursores
durante el evento neutral. Los tres eventos pertenecen a la cronología
completa y se registran en la ruta.

Sean \(\mathcal D=\{N,E,S,O\}\), con vectores
\(\nu_N=(-1,0)\), \(\nu_E=(0,1)\), \(\nu_S=(1,0)\) y
\(\nu_O=(0,-1)\), y \(\Theta_{108}=\ZZ/108\ZZ\). La configuración
observable es
\[
 q=(x_+,d_+,x_\times,d_\times,\theta)
 \in\mathcal Q_{\rm obs}=(X_9\times\mathcal D)^2\times\Theta_{108}.
\]
Para el representante \(\bar\theta\in\{0,\ldots,107\}\), se definen
\[
 r(\theta)=1+(\bar\theta\bmod9),\qquad
 \beta(\theta)=\left[\left\lfloor\frac{\bar\theta}{9}\right\rfloor\right]_{12},
\]
\[
 a_+(\theta)=\mathbf1_{\{1,4,7\}}(r(\theta)),\qquad
 a_\times(\theta)=\mathbf1_{\{2,5,8\}}(r(\theta)).
\]
La transición de posiciones es
\[
 x_+'=x_++a_+(\theta)\nu_{d_+},\qquad
 x_\times'=x_\times+a_\times(\theta)\nu_{d_\times}.
\]
La involución \(\jmath\) intercambia \(N,S\) y \(E,O\). Con
\[
 h(\theta)=\mathbf1_{\{0,27,54,81\}}\bigl((\bar\theta+1)\bmod108\bigr),
\]
el operador observable completo queda fijado por
\[
 F_{\rm obs}(q)=
 (x_+',\jmath^{h(\theta)}d_+,x_\times',
   \jmath^{h(\theta)}d_\times,\theta+1).
\]

\begin{proposition}[Conservación del transporte levantado y retorno observable]
\label{hmtfg:transporte-calendario}
El transporte conserva el desplazamiento entero mediante el registro de
cruces de frontera. El operador \(F_{\rm obs}\) es invertible y satisface
\(F_{\rm obs}^{108}=I\).
\end{proposition}
\begin{proof}
Sea \(s:X_9\to D_9^2\) la sección de representantes positivos. Para
\(x'=x+a\nu_d\), \(a\in\{0,1\}\), se define
\[
 b_d(x,a)=\frac{s(x)+a\nu_d-s(x')}{9}\in\ZZ^2.
\]
Si \(z'=z+b_d(x,a)\), entonces
\[
 s(x')+9z'=s(x)+9z+a\nu_d.
\]
La suma de esta igualdad sobre una ruta proporciona
\[
 s(x_n)+9z_n=s(x_0)+9z_0+\sum_{t=0}^{n-1}a_t\nu_{d_t}.
\]
La partición de la suma conserva la identidad bajo concatenación.
Para invertir \(F_{\rm obs}\), se resta primero una unidad a la fase;
esa fase determina la inversión de direcciones y el cursor desplazado.
Se recuperan las direcciones previas y se resta el vector correspondiente.
Cada bloque de \(27\) pasos contiene nueve activaciones de cada cursor.
El siguiente bloque invierte las direcciones y cancela sus desplazamientos
enteros. En \(54\) pasos se restauran posiciones y direcciones; en
\(108\) se restaura también la coordenada de calendario.
\end{proof}

El estado completo conserva, además de esta configuración, las dos
evaluaciones aritméticas, los acarreos, la orientación, la trayectoria,
el prefijo, el cilindro compatible, su frontera, el grado de memoria y el
registro de supervivencia. Su actualización se factoriza como
\[
 U_t=\mathsf{Mem}_t\circ\mathsf{Tra}_t\circ\mathsf{Sel}_t.
\]
La holonomía \(\Gamma_9\) compone nueve transiciones completas. El retorno
de la fase y el avance de memoria son coordenadas diferentes de esta
composición. La construcción siguiente fija las transiciones y su
actualización campo por campo.

% Recuperación íntegra del propietario de 855 líneas, con etiquetas prefijadas.
\input{inserciones_20260930/feigenbaum/00a_transiciones_nonadicas.tex}

\subsection{Refinamiento de cilindros y lectura ordenada}
\label{hmtfg:refinamiento-ordenado}

% XI REV11, ch_elevacion_catalogal.tex: bloques de 54 aristas, beta_blk,
% partición efectiva de 729 subcilindros; ch_cierre_cardinal.tex:375--497.
Para \(\mathbf b=(b_1,\ldots,b_6)\in\FF_3^6\), sea
\(\upsilon:\FF_3\to\{-1,0,+1\}\) el representante balanceado.
Las tres prolongaciones construidas arista por arista definen
\[
 \Gamma^{[54]}_{\mathbf b,k}
 =\gamma_{\upsilon(b_6),k+45}\circ\cdots\circ
   \gamma_{\upsilon(b_1),k}.
\]
Partiendo de \(Y_0=\{x_\star\}\), se pone
\[
 Y_{N+1}
 =\{\Gamma^{[54]}_{\mathbf b,1+54N}(x):
          x\in Y_N,\ \mathbf b\in\FF_3^6\}.
\]
El truncamiento \(\rho^{\rm blk}_{N+1,N}\) elimina las últimas
cincuenta y cuatro aristas, recuperando el estado anterior mediante
el registro de transición.

\begin{proposition}[Codificación efectiva de los bloques de prolongación]
\label{hmtfg:bloques-729}
La lectura de las seis parejas tríticas de cada bloque define biyecciones
\[
 \beta_N:Y_N\xrightarrow{\ \cong\ }(\FF_3^6)^N
\]
que conmutan con el truncamiento. Cada estado de \(Y_N\) posee
exactamente \(729\) prolongaciones de bloque en \(Y_{N+1}\).
\end{proposition}
\begin{proof}
En cada vuelta, el registro conserva la pareja trítica y su acarreo.
Las tres parejas producen las tres salidas distintas
\(\chi(\varepsilon,\varepsilon)=\varepsilon\), y el registro permite
recuperar la rama utilizada, incluida la rama de acarreo nulo.
Las seis concatenaciones producen, por tanto, todas las palabras de
\(\FF_3^6\). Dos palabras diferentes se distinguen en la primera pareja
registrada en la que difieren. La recuperación del prefijo prueba la
compatibilidad con el truncamiento. Por inducción sobre \(N\), cada
palabra de longitud \(N\) tiene un único estado antecedente; la siguiente
palabra puede ser cualquiera de las \(3^6=729\) posibilidades.
\end{proof}

Sea \(\Omega_{\Gamma}^{\rm cal}=\varprojlim_NY_N\). Las biyecciones
anteriores inducen
\[
 \Omega_{\Gamma}^{\rm cal}\cong(\FF_3^6)^{\mathbb N}.
\]
La correspondencia y su inversa son continuas: cada coordenada finita
depende de un prefijo finito. Los cilindros son abiertos y cerrados.
Cada cilindro de profundidad \(N\) se particiona en sus \(729\)
subcilindros de profundidad \(N+1\); cada uno contiene una prolongación
infinita, obtenida, por ejemplo, continuando con la rama nula en todos
los bloques posteriores.

% Antecedentes efectivos del continuo conjunto y del universo interno.
% Se imprimen en ambos destinos; la etiqueta de universo no reemplaza su regla.
\input{inserciones_20260930/feigenbaum/00b_continuo_universo.tex}

\subsection{Publicación ordenada de los cilindros de bloque}
\label{hmtfg:publicacion-cilindros-bloque}

Para el representante \([b]_3\in\{0,1,2\}\), se fija
\[
 \nu(\mathbf b)=\sum_{i=1}^6[b_i]_3\,3^{6-i}\in\{0,\ldots,728\}.
\]
En la completación ordenada HMT, un prefijo
\(s=(m;\mathbf b_0,\ldots,\mathbf b_{N-1})\), \(m\in\ZZ\), determina
\[
 a_N(s)=m+\sum_{j=0}^{N-1}\nu(\mathbf b_j)\,729^{-j-1},\qquad
 I_N(s)=[a_N(s),a_N(s)+729^{-N})_{\rm HMT}.
\]
El acarreo conserva el empalme entre los dos representantes de una
frontera. Si \(C_N(s)=\overline{I_N(s)}\), las relaciones anteriores dan
\[
 I_N(s)=\bigsqcup_{\mathbf b\in\FF_3^6}I_{N+1}(s;\mathbf b),
 \qquad \operatorname{diam}C_N(s)=729^{-N}\longrightarrow0.
\]

\begin{proposition}[Cobertura de la lectura arquimediana]
\label{hmtfg:cobertura-arquimediana}
La aplicación
\[
 P_{\rm ar}:\ZZ\times\Omega_\Gamma^{\rm cal}
       \longrightarrow\mathbb R_{\rm HMT},\qquad
 \{P_{\rm ar}(m,\xi)\}
    =\bigcap_N C_N(m;\beta_N(\rho_{\infty,N}\xi))
\]
está bien definida, es sobreyectiva y determina
\[
 (\ZZ\times\Omega_\Gamma^{\rm cal})/\!\sim_{P_{\rm ar}}
       \ \cong\ \mathbb R_{\rm HMT}.
\]
\end{proposition}
\begin{proof}
Los extremos izquierdos forman una sucesión de Cauchy en la completación
HMT, pues los intervalos son anidados y sus diámetros tienden a cero.
Su límite pertenece a todos los intervalos cerrados y es único.
Para cualquier elemento de la completación, se elige primero el intervalo
entero que lo contiene y después el único subintervalo semiabierto de
cada partición que lo contiene. La proposición anterior realiza cada
una de esas elecciones mediante una prolongación TPK efectiva.
Los prefijos compatibles definen \(\xi\), cuya lectura es el elemento
elegido. En las fronteras, el acarreo identifica las escrituras
equivalentes después de su construcción. El cociente por igualdad
de lectura resulta, por definición y sobreyectividad, biyectivo con
la completación.
\end{proof}

La construcción se realiza asimismo en el universo genealógico interno
\(\mathfrak G=\mathfrak G_{\rm HMT}(h,m_\infty)\). Aquí \(h\) codifica
las reglas operatorias y \(m_\infty\) la historia compatible; el universo
se obtiene mediante clausura transitiva inicial, definibilidad en cada
sucesor y unión en los estadios límite:
\[
 G_0=\operatorname{TC}\{\omega,u_{h,m},h,m_\infty\},\qquad
 G_{\beta+1}=\operatorname{Def}_{\Sigma_{\rm HMT}}(G_\beta),\qquad
 G_\lambda=\bigcup_{\beta<\lambda}G_\beta .
\]
Las relaciones de \(\Sigma_{\rm HMT}\) son las de las operaciones ya
construidas y sus códigos. Se mantiene, por tanto, el dominio interno
en todas las afirmaciones de cobertura. Denotemos por
\(\mathbb R_{\rm HMT}^{\mathfrak G}\) su completación y por
\(\mathbb R^{\mathfrak G}\) la realización ordenada dentro del mismo
universo.

\begin{proposition}[Carta ordenada interna]
\label{hmtfg:carta-interna}
Existe un único isomorfismo de cuerpos ordenados arquimedianos completos
\[
 \iota:\mathbb R_{\rm HMT}^{\mathfrak G}
             \xrightarrow{\ \cong\ }\mathbb R^{\mathfrak G}
\]
que fija la unidad. Conserva los extremos racionales, las raíces
cuadradas positivas y los límites de sucesiones convergentes.
\end{proposition}
\begin{proof}
La unidad fija los enteros y los racionales. A cada elemento de la
completación HMT se asocia el corte de los racionales menores que él;
la completitud de la realización ordenada proporciona el único elemento
con ese corte. La densidad racional prueba inyectividad y conservación
estricta del orden. El mismo procedimiento en sentido inverso prueba
sobreyectividad. Las aproximaciones racionales por exceso y por defecto
conservan suma y producto y, con ello, la estructura de cuerpo.
El elemento positivo cuyo cuadrado es \(a>0\) se transporta al elemento
positivo cuyo cuadrado es \(\iota(a)\). Si una sucesión converge, para
cada tolerancia racional positiva sus términos finales están en el
entorno racional del límite; la conservación del orden transporta esa
condición. Todas estas operaciones se efectúan dentro de
\(\mathfrak G\). La unicidad sigue de la conservación de los cortes
racionales.
\end{proof}

\subsection{Descenso por fibras y conservación del selector}
\label{hmtfg:descenso-selector}

% XI REV11, ch_descenso_por_fibras.tex:10--39; prueba completa pertinente.
\begin{proposition}[Criterio de descenso de una operación]
\label{hmtfg:criterio-descenso}
Sea \(q:\Omega\twoheadrightarrow J\) y
\(\star:D\subseteq\Omega^2\to\Omega\), con \(D\) saturado por las
fibras de \(q\times q\). Existe una única operación \(\bar\star\)
sobre \((q\times q)(D)\) tal que
\[
 q(x\star y)=q(x)\,\bar\star\,q(y)
\]
si y sólo si
\[
 q(x)=q(x'),\quad q(y)=q(y')
 \ \Longrightarrow\ q(x\star y)=q(x'\star y')
\]
para los pares del dominio.
\end{proposition}
\begin{proof}
Si existe \(\bar\star\), ambos resultados son la evaluación de la
misma pareja de argumentos. Recíprocamente, se eligen antecedentes
de los argumentos visibles y se define el resultado como la lectura
de su operación. La saturación fija el dominio; la condición sobre
fibras hace independiente el resultado de los antecedentes elegidos.
La sobreyectividad prueba la unicidad.
\end{proof}

En una aplicación mediante \(P_{\rm ar}\), este criterio identifica la
condición precisa para transportar una operación desde estados completos
hasta su lectura. En la carta \(\iota\), la estructura de cuerpo y
los límites ya satisfacen esa condición. A continuación se construye
el perfil analítico a partir del calendario y se utiliza esta
compatibilidad para transportar su iteración y sus raíces.

\subsection{Publicación dodecafásica y normalización cuadrática}
\label{hmtfg:perfil-dodecafasico}

El perfil, sus fases y su normalización son los ya construidos en \eqref{mab:rm:eq:critical-phase-reader}--\eqref{mab:rm:eq:chaos-profile}. Su forma cerrada y sus momentos están demostrados en \eqref{mab:rm:eq:chaos-closed} y \eqref{mab:rm:eq:chaos-moments24}. La suma finita de cosenos determina la prolongación removible en todos los ceros del denominador de la forma cerrada. El momento angular conserva representantes orientados reales: cambiarlos módulo $2\pi$ modifica el momento. El lector y su sector uniforme son la estructura de partida explícita de esta realización.

\subsection{Transporte de iteradas, raíces y mínimos}
\label{hmtfg:transporte-raices}

% Reunión existente: TRANSVERSALIDAD_Y_ESCALADO_LOCAL_FEIGENBAUM.md, sección 14.1.
Dentro de \(\mathbb R_{\rm HMT}^{\mathfrak G}\), la serie convergente
\[
 \cos_{\rm HMT}(t)=\sum_{j=0}^{\infty}
                   \frac{(-1)^j t^{2j}}{(2j)!}
\]
define la realización analítica correspondiente. Sean
\[
 u_{\rm HMT}(x)=\frac1{12}\sum_{m=1}^{12}
  \left[1-\cos_{\rm HMT}
   \left(\frac{2(3m-1)x}{\sqrt{899}_{\rm HMT}}\right)\right],
 \qquad F_a(x)=1-a\,u_{\rm HMT}(x),
\]
\[
 S_n^{\rm HMT}(a)=F_a^{\,2^n}(0),\qquad
 S_n(\lambda)=f_\lambda^{\,2^n}(0).
\]

\begin{proposition}[Conjugación ordenada de la familia y del selector de raíces]
\label{hmtfg:conjugacion-selector}
Para los argumentos internos de las expresiones anteriores,
\[
 \iota\circ F_a=f_{\iota(a)}\circ\iota,\qquad
 \iota(S_n^{\rm HMT}(a))=S_n(\iota(a)).
\]
La carta transporta exactamente las condiciones de retorno, periodo
exacto y orden de los parámetros. Si un conjunto de raíces admisibles
posee mínimo, la imagen de ese mínimo es el mínimo del conjunto imagen.
\end{proposition}
\begin{proof}
La carta fija los racionales, transporta la raíz cuadrada positiva y
conserva los límites de las sumas parciales de coseno. Por ello,
\(\iota(u_{\rm HMT}(x))=u_0(\iota(x))\). La compatibilidad con
suma y producto prueba la primera igualdad; la inducción sobre el
número de iteraciones prueba la segunda. Una raíz se caracteriza por
la igualdad con cero, que se conserva en ambas direcciones. El
periodo exacto \(2^n\) añade las desigualdades
\(F_a^{2^j}(0)\ne0\), \(0\le j<n\), también conservadas por la
inyectividad de \(\iota\). Si se exige que el parámetro supere un
centro precedente, se transporta esa desigualdad por conservación
del orden. Finalmente, \(a\le b\) para todo \(b\) del conjunto
equivale a \(\iota(a)\le\iota(b)\) para todo elemento de su
imagen. La sobreyectividad sobre los conjuntos así definidos
completa la afirmación acerca del mínimo.
\end{proof}

El dominio de este transporte es el continuo interno indicado.
La existencia de cada mínimo y su identificación con una rama dinámica
concreta son resultados de la construcción de criticidad que sigue;
la carta conserva esas propiedades cuando han sido establecidas.

\subsection{Representación espectral finita del lector uniforme}
\label{hmtfg:jacobi-finito}

% Grassmannianos REV07, sections/jacobi_spectral.tex:35--136,
% mecanismo de Hankel, Gram--Schmidt, recurrencia y representación espectral.
% Especialización atómica ya reunida en la sección 14.2 del desarrollo Feigenbaum.
Las frecuencias cuadradas del perfil definen la medida positiva
\[
 s_m=k_m^2=\frac{4(3m-1)^2}{899},\qquad
 \nu_0=\frac1{12}\sum_{m=1}^{12}\delta_{s_m},\qquad
 M_r=\int s^r\,d\nu_0(s).
\]
Sus doce nodos son distintos y positivos. Para \(1\le r\le12\),
la matriz de Hankel \(H_r=(M_{i+j})_{0\le i,j<r}\) es definida
positiva. En efecto, para \(p\) de grado menor que \(r\),
\[
 \mathbf p^{\mathsf T}H_r\mathbf p
   =\frac1{12}\sum_{m=1}^{12}p(s_m)^2>0
\]
salvo que \(p=0\), porque un polinomio de grado a lo sumo once
que se anula en doce nodos distintos es nulo. En grado doce aparece
el polinomio \(\prod_m(s-s_m)\), de norma nula en esta medida.
La representación espectral correspondiente tiene dimensión doce.

Sean \(p_0,\ldots,p_{11}\) los polinomios mónicos ortogonales,
\(h_j=\int p_j^2\,d\nu_0>0\) y
\(\psi_j=p_j/\sqrt{h_j}\). Como \(\nu_0\) es una probabilidad,
\(\psi_0=1\). La multiplicación por \(s\) en \(L^2(\nu_0)\)
es autoadjunta y posee una matriz tridiagonal \(J_{12}\) en esta base.

\begin{proposition}[Recurrencia y entrelazamiento espectral exactos]
\label{hmtfg:jacobi-entrelazamiento}
Con
\[
 a_j=\frac{\int s p_j(s)^2\,d\nu_0(s)}{h_j},\qquad
 \beta_j=\frac{h_j}{h_{j-1}}>0\quad(1\le j\le11),
\]
la matriz \(J_{12}\) tiene diagonal \(a_j\) y subdiagonal
\(\sqrt{\beta_j}\). Si
\[
 Q_{mj}=\frac{\psi_j(s_m)}{\sqrt{12}}
 \quad(1\le m\le12,\ 0\le j\le11),\qquad
 D=\operatorname{diag}(s_1,\ldots,s_{12}),
\]
entonces
\[
 Q^{\mathsf T}Q=I,\qquad DQ=QJ_{12},\qquad
 J_{12}=Q^{\mathsf T}DQ,
\]
y el perfil satisface
\begin{equation}
 u_0(x)=e_0^{\mathsf T}
          \bigl[I-\cos(x\sqrt{J_{12}})\bigr]e_0.
 \label{hmtfg:perfil-espectral}
\end{equation}
\end{proposition}
\begin{proof}
La positividad de los sistemas de Hankel proporciona de forma única
los polinomios mónicos ortogonales. Al desarrollar \(s p_j\) en la
base ortogonal, los coeficientes de \(p_i\), \(i\le j-2\), se
anulan porque
\(\langle sp_j,p_i\rangle=\langle p_j,sp_i\rangle=0\).
El coeficiente de \(p_{j+1}\) es uno por monicidad, el de \(p_j\)
es \(a_j\) y el de \(p_{j-1}\) es \(h_j/h_{j-1}\).
En el último grado, la misma igualdad se entiende en
\(L^2(\nu_0)\), donde el polinomio mónico de grado doce que se
anula en los nodos representa el vector cero. La normalización da
la subdiagonal positiva indicada.

La ortonormalidad equivale a \(Q^{\mathsf T}Q=I\). La igualdad
de la recurrencia en cada nodo da \(DQ=QJ_{12}\).
Como \(Q\) es cuadrada, también \(QQ^{\mathsf T}=I\), y se
obtiene la diagonalización de \(J_{12}\).
Sus autovalores positivos permiten definir su raíz cuadrada
positiva. Para toda función definida en esos doce nodos,
\[
 e_0^{\mathsf T}\Phi(J_{12})e_0
   =\frac1{12}\sum_{m=1}^{12}\Phi(s_m),
\]
pues \(Qe_0=(1,\ldots,1)^{\mathsf T}/\sqrt{12}\).
La elección \(\Phi(s)=1-\cos(x\sqrt{s})\) prueba
\eqref{hmtfg:perfil-espectral}.
\end{proof}

La representación conserva exactamente el funcional uniforme, sus
momentos y su perfil. Sus primeros coeficientes son
\[
 a_0=M_1=2,\qquad
 \beta_1=M_2-M_1^2=\frac{2493348}{808201}.
\]
La genealogía completa permanece en el estado TPK del que procede
el lector; \(J_{12}\) representa su publicación espectral concreta.
La construcción de Gram--Schmidt y la recurrencia de Jacobi se aplican
a esta medida atómica, con su terminación en dimensión doce.

La cadena operatoria utilizada en las secciones siguientes queda así
fijada: las células de APP y la orientación TRIT alimentan la
transición TPK; sus prolongaciones producen los cilindros y su
lectura ordenada; las doce ventanas orientadas y el funcional uniforme
producen \(u_0\); la familia \(f_\lambda\) determina retornos y raíces;
la carta interna conserva el orden de sus selectores. La realización
espectral describe el mismo perfil mediante una matriz autoadjunta
finita y un vector cíclico.

% Fragmento integrable. Preambulo anfitrion: amsmath, amssymb, longtable, hyperref.
\section{Perfil dodecafásico y retorno cuadrático}
\label{hmtfg:fragmento:perfil-dodecafasico-y-retorno-cuadratico}

\subsection{Genealogía y objeto de la demostración}
\label{hmtfg:desarrollo:1}

El objeto inicial es la publicación crítica de la rueda dodecafásica orientada ya construida en el corpus. APP evalúa suma y producto sobre las dos hojas de la retícula de dígitos, conservando cociente y residuo; TRIT determina régimen y orientación; el transporte TPK selecciona, transporta y actualiza el estado con sus registros de acarreo, frontera y memoria. La conexión nonádica prolonga ese estado, con retorno de fase y avance de memoria. Sus lectores actúan sobre la misma estructura discreta del continuo.

El lector focal de la fuente es
\[
\mathcal P_{\mathrm{crit}}:U_{12}^{\mathrm{sign}}\longrightarrow
\mathbb R/2\pi\mathbb Z,\qquad
m\longmapsto\phi_m=\frac{(3m-1)\pi}{18},\quad 1\le m\le12.
\]
Para evaluar el momento angular se conserva el representante orientado de esa carta: sustituirlo por otro representante módulo \(2\pi\) cambiaría el momento. La traza uniforme fija \(w_m=1/12\). La cadena focal es
\[\begin{aligned}
\text{APP–TRIT–TPK}&\longrightarrow
\text{estructura discreta del continuo}\\ &\longrightarrow U_{12}^{\mathrm{sign}}
\xrightarrow{\mathcal P_{\mathrm{crit}}}(\phi_m,w_m)
\\ &\longrightarrow u_0\longrightarrow f_\lambda
\longrightarrow\text{retornos y renormalización}.
\end{aligned}\]

La escritura abreviada conserva como antecedente el estado completo y sus prolongaciones; la familia escalar es una lectura de ese antecedente. Se toma como estructura de partida explícita el lector ya fijado, con su orden, orientación y sector uniforme. Las conclusiones se refieren a ese lector concreto; la unicidad de ese lector entre todas las publicaciones posibles del TPK queda fuera de estas pruebas.

La coordenada \(\pi_{\mathrm{HMT}}\) empleada por la carta pertenece a las salidas internas del corpus. Su factor se cancela exactamente al normalizar el segundo momento. Los valores de Feigenbaum intervienen exclusivamente en el reconocimiento posterior, nunca en la definición de la familia o en la certificación de raíces.


\subsection{Perfil exacto y criticidad cuadrática}
\label{hmtfg:desarrollo:2}

Se conserva la notación $s_m=3m-1$, $\kappa_0=s_0$, $k_m=2s_m/\sqrt{899}$ y $f_{\lambda,0}=f_\lambda$ del perfil ya demostrado. Su prolongación entera tiene la serie racional
\[
u_0(z)=\sum_{j\ge1}(-1)^{j+1}
\frac{4^j\sum_m s_m^{2j}}{12\,899^j(2j)!}z^{2j}
=z^2-\frac{715769}{2424603}z^4
+\frac{1353832978}{32695771455}z^6+\cdots.
\]
Así, el registro finito de doce sectores produce un perfil entero, con criticidad cuadrática y escala fijadas antes de iterar.


\subsection{Robustez estructural con cotas uniformes}
\label{hmtfg:desarrollo:3}

La deformación es la definida en \eqref{mab:eq:pesos-dodecafasicos-armonizados}--\eqref{mab:eq:familia-feigenbaum-ponderada-armonizada}. Escribimos $M_{2,\eta}=\sum_mw_m\phi_m^2$ y $\kappa_\eta=s_\eta=\sqrt{2/M_{2,\eta}}$, conservando los mismos pesos, fases y perfiles.
Las notaciones son $p_m=p(\phi_m)=12\cos(3\phi_m)-\cos(4\phi_m)$ y $k_{m,\eta}=\kappa_\eta\phi_m$; en las fórmulas ponderadas siguientes, $k_m$ abrevia $k_{m,\eta}$, mientras que $k_{m,0}$ corresponde al sector uniforme.

\textbf{Teorema de robustez estructural.} Para \(|\eta|\le1/40\), los pesos son positivos y suman uno. El perfil es par, entero, satisface \(u_\eta''(0)=2\) y es estrictamente creciente en \((0,1]\). Para \(0<\lambda\le2/u_\eta(1)\), el mapa lleva \([-1,1]\) en sí mismo, tiene un único punto crítico en ese intervalo y cumple
\[
S f_{\lambda,\eta}(x)\le-\frac{640}{47647}<0,\qquad0<|x|\le1.
\]

\textbf{Demostración.} Las sumas de los cosenos \(3\phi_m\) y \(4\phi_m\) se anulan por ciclos de longitudes cuatro y tres. Como \(|p_m|\le13\),
\[
\frac9{160}\le w_m\le\frac{53}{480},\qquad
\frac{27}{40}M_{2,0}\le M_{2,\eta}\le\frac{53}{40}M_{2,0}.
\]
Se deduce
\[
k_{\max,\eta}^2\le\frac{196000}{24273}<9<\pi^2,\qquad
k_{\min,\eta}^2\ge\frac{640}{47647}.
\]
Para \(0<x\le1\), todos los senos \(\sin(k_{m,\eta}x)\) son positivos, de donde
\[
u_\eta'(x)=\sum_mw_mk_m\sin(k_mx)>0,\qquad
u_\eta'''(x)=-\sum_mw_mk_m^3\sin(k_mx)<0.
\]
El cociente \(u_\eta'''/u_\eta'\) es el negativo de una media positiva de los \(k_m^2\). La invariancia de la schwarziana bajo cambios afines del valor da
\[
S f_{\lambda,\eta}
=\frac{u_\eta'''}{u_\eta'}-\frac32
\left(\frac{u_\eta''}{u_\eta'}\right)^2
\le-k_{\min,\eta}^2.
\]
La paridad cubre \(x<0\), y la imagen del intervalo es
\([1-\lambda u_\eta(1),1]\). \(\square\)

La anchura \(1/40\) es una elección para obtener cotas uniformes, distinta de una constante fundamental o de un umbral óptimo. El sector original sigue siendo \(\eta=0\).


\subsection{Coordenada holomorfa cuadrática global en una franja}
\label{hmtfg:desarrollo:4}

\textbf{Teorema de factorización holomorfa.} Para \(|\eta|\le1/40\), en
\[
\mathcal S=\{z\in\mathbb C:|\operatorname{Re}z|<\pi/3\}
\]
existe una función holomorfa, impar y univalente \(b_\eta\), con \(b_\eta'(0)=1\), tal que
\[
\boxed{u_\eta(z)=b_\eta(z)^2,\qquad
f_{\lambda,\eta}(z)=1-\lambda b_\eta(z)^2.}
\]

\textbf{Demostración.} Para \(z=x+iy\) en la semibanda derecha,
\[
\operatorname{Re}u_\eta'(z)
=\sum_mw_mk_m\sin(k_mx)\cosh(k_my)>0.
\]
La integral de \(u_\eta'\) sobre el segmento entre dos puntos, dividida por la diferencia de los extremos, tiene parte real positiva. La convexidad de la semibanda prueba allí la inyectividad. La paridad da el resultado correspondiente a la izquierda.

Además,
\[
\operatorname{Im}u_\eta(x+iy)
=\sum_mw_m\sin(k_mx)\sinh(k_my)
\]
tiene el signo de \(y\) para \(x>0\). Sobre el eje real positivo \(u_\eta>0\); sobre el imaginario,
\[
u_\eta(iy)=\sum_mw_m(1-\cosh(k_my))
\]
es estrictamente negativo fuera de cero y decrece estrictamente con \(|y|\). Estas separaciones, la inyectividad de cada semibanda y la paridad prueban que las fibras de \(u_\eta\) en \(\mathcal S\) son exactamente \(\{z,-z\}\), salvo la fibra del origen.

La función \(u_\eta(z)/z^2\), prolongada con valor uno en cero, es holomorfa y carece de ceros en la franja simplemente conexa. Tiene un logaritmo holomorfo \(L_\eta\) con \(L_\eta(0)=0\), que es par por unicidad de la normalización. Definimos
\[
b_\eta(z)=z\exp\!\left(\frac12L_\eta(z)\right).
\]
Entonces \(b_\eta^2=u_\eta\), \(b_\eta\) es impar y \(b_\eta'(0)=1\). La igualdad \(b_\eta(z)=b_\eta(w)\) obliga a \(w=z\) o \(w=-z\). El segundo caso sólo da igualdad en el origen. Esto prueba su univalencia. \(\square\)

Las doce fases determinan así una deformación conforme explícita del pliegue cuadrático. La conjugación dinámica completa conserva esa deformación:
\[
b_\eta\circ f_{\lambda,\eta}\circ b_\eta^{-1}(w)
=b_\eta(1-\lambda w^2),
\]
en sus dominios de composición. La factorización demostrada y una conjugación a toda la familia cuadrática son afirmaciones distintas.


\subsection{Retorno, cambio de escala y memoria}
\label{hmtfg:desarrollo:5}

Sean \(a=-f(1)>0\), \(H_a(x)=-ax\), \(J=H_a([-1,1])\). Cuando \(J\subset[-1,1]\) es admisible para el retorno, \(f^2(J)\subset J\) y las composiciones están definidas,
\[
\mathcal Rf=H_a^{-1}f^2H_a,\qquad
\mathcal Rf(x)=-a^{-1}f(f(-ax)).
\]
Si además \(f(J)\subset(0,1]\), se conserva un único máximo cuadrático central:
\[
(\mathcal Rf)(0)=1,\qquad
(\mathcal Rf)''(0)=-a f'(1)f''(0)<0.
\]
La composición de schwarzianas da, en los puntos regulares,
\[
S(\mathcal Rf)(x)=a^2
\left[(Sf)(f(-ax))f'(-ax)^2+(Sf)(-ax)\right]<0.
\]
Las inclusiones de dominios forman parte de las hipótesis y se comprueban a cada escala.

Sea \(v_\eta=(u_\eta|_{[0,1]})^{-1}\). Las ramas inversas son
\[
\beta_\sigma(y)=\sigma v_\eta((1-y)/\lambda),\quad
\sigma\in\{-1,+1\},
\]
con el registro crítico \(\beta_0(1)=0\). El paso
\[
(x,w)\longmapsto(f(x),w\mathbin{\|}\sigma(x))
\]
se invierte sobre su imagen leyendo la última hoja y aplicando \(\beta_\sigma\). La palabra conserva decisiones de rama; el valor anterior se reconstruye por la ecuación del mapa.

El retorno agrupa dos transiciones. Desde el extremo \(x'\) y sus hojas \((\sigma_0,\sigma_1)\) se recupera
\[
x=H_a^{-1}\beta_{\sigma_0}
\!\left(\beta_{\sigma_1}(H_ax')\right).
\]
También se reconstruyen los puntos intermedios. Los registros APP–TRIT–TPK de hoja, acarreo, ruta, frontera y memoria se concatenan conservando su identidad. El signo de la rama analítica es un registro adicional: no reemplaza esos datos por una etiqueta binaria.

Si \(C_N\) agrupa las historias de longitud \(2N\) en pares, las restricciones satisfacen
\[
\pi^{\mathrm{ret}}_{N,M}C_N
=C_M\pi^f_{2N,2M},\qquad M\le N.
\]
Ambos miembros suprimen los mismos pares terminales y conservan las mismas ramas inversas. Queda demostrado el transporte compatible de historias finitas.


\subsection{Escalas acumuladas y derivada del operador}
\label{hmtfg:desarrollo:6}

Sean \(c_n=f^{2^n}(0)\), \(c_0=1\) y \(H_n(x)=c_nx\). Mientras las composiciones sean admisibles y \(c_n\ne0\),
\[
\boxed{\mathcal R^nf=H_n^{-1}f^{2^n}H_n},\qquad
(\mathcal R^nf)(1)=\frac{c_{n+1}}{c_n}=-a_n.
\]
La inducción se obtiene de \(H_nH_{a_n}=H_{n+1}\). La convención \(a_n>0\) exige alternancia de signo de \(c_n\). En una raíz superestable de período \(2^N\), \(c_N=0\): esa normalización se hace singular y deben utilizarse los niveles anteriores o un límite controlado.

Para \(z=-ax\), \(w=f(z)\) y una perturbación \(h\), la derivada completa es
\[
\boxed{
D\mathcal R_f[h](x)=
-\frac{h(w)+f'(w)h(z)}a+
\frac{h(1)}a
\left[\mathcal Rf(x)-xf'(w)f'(z)\right].
}
\]
Se deriva usando \(\dot a=-h(1)\), \(\dot z=xh(1)\),
\(\dot w=h(z)+xf'(z)h(1)\) y la variación del denominador. Si \(h(0)=0\), resulta \(D\mathcal R_f[h](0)=0\). La dirección paramétrica inicial HMT es \(h=-u_\eta\).

Esta fórmula incluye la variación de escala; congelar \(a\) daría otro operador linealizado.


\subsection{Raíces e itinerarios certificados hasta período 256}
\label{hmtfg:desarrollo:7}

Sea \(F_n(\lambda)=f_{\lambda,0}^{2^n}(0)\). Se calcula
\[
x_0=p_0=0,\qquad x_{j+1}=1-\lambda u_0(x_j),\qquad
p_{j+1}=-u_0(x_j)-\lambda u_0'(x_j)p_j.
\]
Entonces \(F_n=x_{2^n}\) y \(F_n'=p_{2^n}\).

El programa adjunto usa intervalos racionales con redondeo entero hacia afuera, raíz cuadrada encerrada por raíz entera y seno/coseno con resto de Taylor explícito. La cota exacta \(|u_0''|\le2\) permite evaluación centrada con resto cuadrático.

Las aproximaciones históricas sólo proponen cajas de radio \(10^{-42}\). Se certifican signos opuestos en sus extremos, \(F_n'\) separado de cero en toda la caja y exclusión de los retornos de períodos propios divisores de \(2^n\). Valor intermedio y monotonía prueban existencia y unicidad local; la exclusión de divisores prueba período mínimo.

\begingroup\small
\begin{longtable}{p{.12\linewidth}p{.17\linewidth}p{.62\linewidth}}
\hline
Nivel & Período mínimo & Centro de la caja, redondeado \\
\hline
1 & 2 & 1.345913990471832792210949 \\
2 & 4 & 1.763379787846910127290794 \\
3 & 8 & 1.850413796950136464530567 \\
4 & 16 & 1.868979756606798125150574 \\
5 & 32 & 1.872955034435659740004205 \\
6 & 64 & 1.873806367467030119412262 \\
7 & 128 & 1.873988695221365362307169 \\
8 & 256 & 1.874027744154450672897008 \\
\hline\end{longtable}
\endgroup

La primera raíz tiene la expresión exacta \(\lambda_1=1/u_0(1)\). Las cajas completas y sus comprobaciones están en el certificado JSON.

Se han encerrado los 502 estados preclausura de estos ocho niveles, todos separados de cero, y sus palabras completas de signos. En particular,
\[
\operatorname{sign}f_{\lambda_n,0}^{2^j}(0)=(-1)^j,\qquad0\le j<n.
\]
Quedan acreditadas las orientaciones finitas de los reescalados anteriores a la clausura. Las inclusiones de intervalos restrictivos a todas las escalas constituyen una comprobación adicional.

Para \(\delta_{\mathrm F,n}=(\lambda_{n-1}-\lambda_{n-2})/(\lambda_n-\lambda_{n-1})\), un intervalo decimal exterior abreviado es
\[
\boxed{\begin{gathered}
4.66921218915020415455609621588966392804<\delta_{\mathrm F,8},\\
\delta_{\mathrm F,8}<4.66921218915020415455609621588966392863.
\end{gathered}}
\]
La anchura del intervalo completo es inferior a \(5.808\cdot10^{-37}\). Mide el error numérico de la razón finita \(\delta_{\mathrm F,8}\), distinto del error de aproximación al límite universal.


\subsection{Persistencia analítica de las raíces deformadas}
\label{hmtfg:desarrollo:8}

Para cada nivel certificado, \(F_n'(\lambda_n)\ne0\). La función implícita produce una rama analítica local \(\lambda_n(\eta)\) con
\[
f_{\lambda_n(\eta),\eta}^{2^n}(0)=0.
\]
Por continuidad, los retornos propios siguen separados de cero para \(\eta\) suficientemente pequeño; se preservan período mínimo e itinerario. La existencia de esos entornos locales queda demostrada. Su anchura común aún no se identifica con toda la franja estructural \(1/40\).

Con \(M=M_{2,\eta}\), la respuesta del perfil es
\[
v_\eta(x)=\partial_\eta u_\eta(x)
=\sum_m\frac{p_m}{12}(1-\cos(k_{m,\eta}x))
-\frac{M'}{2M}\,x u_\eta'(x).
\]
El segundo término conserva la normalización del segundo momento. Si \(q_j=\partial_\eta x_j\) a \(\lambda\) fija,
\[
q_0=0,\qquad q_{j+1}=-\lambda v_\eta(x_j)-\lambda u_\eta'(x_j)q_j,
\qquad
\boxed{\lambda_n'(\eta)=-q_{2^n}/p_{2^n}}.
\]
Así queda explícita la cadena pesos–escala–respuesta–parámetro superestable. Esta no degeneración finita es distinta de la transversalidad asintótica a la hoja estable del punto fijo universal.


\subsection{Colas analíticas y propagación del error}
\label{hmtfg:desarrollo:9}

Para \(\eta=0\), sean
\[
\mu_{2j}=\frac1{12}\sum_mk_m^{2j},\quad
p_N(z)=\sum_{j=1}^N(-1)^{j+1}\frac{\mu_{2j}}{(2j)!}z^{2j},
\quad\Omega=\frac{70}{\sqrt{899}}.
\]
Si \(q_N=\Omega^2r^2/[(2N+3)(2N+4)]<1\),
\[
\|u_0-p_N\|_r\le
\frac{\mu_{2N+2}r^{2N+2}}{(2N+2)!(1-q_N)}.
\]
La desigualdad \(\mu_{2j+2}\le\Omega^2\mu_{2j}\) domina la cola por una serie geométrica. Para \(0\le s\le2N+2\),
\[
\|(u_0-p_N)^{(s)}\|_r\le
\frac{\mu_{2N+2}r^{2N+2-s}}{(2N+2-s)!}
\left(1-\frac{\Omega^2r^2}{(2N+3-s)(2N+4-s)}\right)^{-1},
\]
cuando el denominador es positivo.

Para transportar errores por un retorno se controlan también sus dominios: si \(\|f-g\|\le\varepsilon\), \(a_f,a_g\ge a_*>0\), las composiciones quedan en un dominio convexo común y allí \(\|f'\|\le L,\ \|g\|\le M\), entonces
\[
\|\mathcal Rf-\mathcal Rg\|_r
\le\varepsilon\left[\frac{1+L+L^2r}{a_*}+\frac{M}{a_*^2}\right].
\]
La prueba suma la perturbación exterior, la interior, el cambio del argumento por \(|a_f-a_g|\le\varepsilon\) y la variación del denominador. Estas cotas evitan reutilizar la suma inicial de cosenos como si todos los perfiles renormalizados conservaran esa misma forma.

% Fragmento integrable. Preambulo anfitrion: amsmath, amssymb, longtable, hyperref.
\section{Profundidad nonádica y persistencia del itinerario}
\label{hmtfg:fragmento:profundidad-nonadica-y-persistencia-del-itinerario}

\subsection{Crecimiento exterior y resolución interior}
\label{hmtfg:prolongacion:2}

La ley de refinamiento del artículo XI, con base entera \(B\ge2\), es
\[
L_c(x,y)=(Bx-c,By+c),\qquad 0\le c<B.
\]
Para una palabra \(c_1\ldots c_m\), su registro es
\[
C_m=\sum_{j=1}^m c_jB^{m-j},\qquad
(x_m,y_m)=(B^mx_0-C_m,B^my_0+C_m).
\]
La suma exterior vale \(B^mM_0\), con \(M_0=x_0+y_0\). La lectura normalizada del registro pertenece al cilindro
\[
I_m(C_m)=\left[\frac{C_m}{B^m},\frac{C_m+1}{B^m}\right],
\qquad \operatorname{diam}I_m=B^{-m}.
\]
La pareja entera y la escala recuperan los cocientes, residuos y ceros iniciales de la palabra. Los extremos con doble expansión conservan su marca de frontera. Así se coordinan crecimiento de soporte y refinamiento de lectura sin identificar historias distintas por su sola imagen límite.

La composición es exacta:
\[
L_D^{[b]}\circ L_C^{[a]}=L_{B^bC+D}^{[a+b]}.
\]
Por asociatividad, agrupar una historia de dieciocho eventos en nueve bloques de dos o en dos bloques de nueve produce el mismo registro. A nivel de retorno \(2^N\), el horizonte común es \(9\,2^N\). La prueba es una igualdad de composición sobre la misma historia ordenada; mantiene sus acarreos y sus subbloques recuperables.

La restricción de estados y la prolongación de sus lectores se expresan mediante los límites siguientes:
\[
X_\infty=\varprojlim X_m,\qquad
r^*a=a\circ r,\qquad
r^*e_x=\sum_{r(y)=x}e_y,
\qquad \mathfrak H_{\rm cyl}=\varinjlim\mathfrak H_m.
\]
Los estados se restringen; los lectores se prolongan. La preimagen preserva productos, unidad e idempotentes. Una familia exactamente compatible induce un lector en el límite directo. Los aproximantes que sólo son compatibles hasta error requieren, además, una estimación de convergencia en su completación. La sección siguiente proporciona esa estimación para el retorno crítico.

\subsubsection{Capacidad y vacancias del registro}

El lector de capacidad de XI compara bloques de \(729\) y \(1000\) posibilidades. Con \(\rho=\log_{1000}729=\log_{10}9\), conserva
\[
\mathcal C_k=\lfloor(k+1)\rho\rfloor,\qquad
\mathcal C_k-\mathcal C_{k-1}=1-z_k,
\]
\[
\mathcal C_b-\mathcal C_a+\sum_{k=a+1}^b z_k=b-a,
\qquad T_{\rm cap}(a)=\lceil a/\rho\rceil-1\quad(a\ge1).
\]
Cada vacancia registra una transición durante la cual se mantiene la capacidad; la longitud cronológica queda conservada en el balance. La fase de capacidad y la fase cronológica son lectores distintos.

Para almacenar un prefijo de \(m\) dígitos de base nueve, una capacidad decimal suficiente es
\[
a(m)=\left\lceil\frac{m\rho}{3}\right\rceil.
\]
El cociente y el residuo permiten convertir el entero \(C_m\) entre bases, manteniendo aparte la profundidad \(m\). Se cumplen
\[
9^m\le1000^{a(m)}\le729^{T_{\rm cap}(a(m))+1}.
\]
Estas desigualdades también se deciden mediante potencias enteras, sin redondear logaritmos. El presupuesto de lectura de la sección siguiente se traduce así a capacidad del registro conservando las unidades. Esta cuenta dimensiona el almacenamiento; las reglas de supervivencia determinan, adicionalmente, qué historias nativas son admisibles.


\subsection{Presupuesto explícito de profundidad para cada retorno}
\label{hmtfg:prolongacion:3}

\textbf{Teorema.} Fijado \(|\eta|\le1/40\), normalícense
\[
t=\lambda u_\eta(1)/2\in[0,1],\qquad z=(x+1)/2,
\]
\[
F_{t,\eta}(z)=1-t\frac{u_\eta(2z-1)}{u_\eta(1)},\qquad z\in[0,1].
\]
Para todo entero \(q\ge1\) y dos parámetros \(t,s\),
\[
\left|F_{t,\eta}^{q}(1/2)-F_{s,\eta}^{q}(1/2)\right|
\le S_q|t-s|,
\qquad S_q=\sum_{j=0}^{q-1}(6\sqrt2)^j<9^q.
\]
En consecuencia, un cilindro de parámetros en base nueve, de profundidad \(m=2^N+r\), encierra la lectura del retorno \(q=2^N\) en un intervalo de diámetro inferior a \(9^{-r}\).

\textbf{Demostración.} La función \((1-\cos v)/v^2\) decrece para \(0<v<\pi\): el signo de su derivada es el de \(v\sin v-2(1-\cos v)\), negativo porque \(\tan(v/2)>v/2\). Por el segundo momento,
\[
u_\eta(1)\ge\frac{1-\cos3}{9}\sum_jw_jk_j^2
=\frac{2(1-\cos3)}9>\frac13.
\]
La última desigualdad se verifica, por ejemplo, acotando \(\cos3\) mediante su suma de Taylor alternada hasta grado ocho, que es menor que \(-1/2\). Además, Cauchy–Schwarz da \(|u_\eta'(x)|\le\sqrt2\) en la recta real. Se deduce
\[
\operatorname{Lip}_zF_{t,\eta}\le6\sqrt2<9,\qquad
\operatorname{Lip}_tF_{t,\eta}\le1.
\]
Si \(d_j\) es la diferencia entre las dos órbitas, entonces \(d_0=0\) y \(d_{j+1}\le6\sqrt2\,d_j+|t-s|\). La inducción prueba la suma geométrica. Para \(|t-s|\le9^{-m}\), el diámetro es menor que \(9^{q-m}\). Sustituir \(q=2^N\) y \(m=q+r\) termina la prueba. \(\square\)

La cota controla la incertidumbre de parámetro con \(\eta\) fijado. Una implementación añade y propaga separadamente el error de evaluación de cosenos y de redondeo; el certificado racional antecedente ofrece esas operaciones. La cota es suficiente y conservadora, con coste que crece con el horizonte.

\textbf{Lector límite.} Realícese \(t\) desde \((x_0,y_0)=(1,0)\), \(B=9\), de modo que
\[
(x_m,y_m)=(9^m-C_m,C_m),\qquad t_m=C_m/9^m.
\]
Los lectores \(E_{q,m}=F_{t_m,\eta}^{q}(1/2)\), llevados a un nivel común, cumplen
\[
\|E_{q,m+a}-r^*E_{q,m}\|_\infty\le S_q9^{-m}.
\]
Convergen uniformemente a un lector continuo \(E_q\) en la completación del álgebra cilíndrica. La versión conjuntista
\[
J_{q,m}(C)=\{F_{t,\eta}^{q}(1/2):t\in I_m(C)\}
\]
es exactamente anidada y sus diámetros tienden a cero. Este resultado da el paso al límite de la \textbf{lectura a horizonte fijo}, con control de profundidad a cualquier horizonte solicitado.


\subsection{Palabra de duplicación y lector nonádico de la cronología}
\label{hmtfg:prolongacion:4}

Los ocho itinerarios certificados en el antecedente obedecen
\[
W_1=+,\qquad W_{n+1}=W_n\,(-1)^n\,W_n.
\]
Su continuación prefijal única es
\[
\tau_j=(-1)^{v_2(j)},\qquad j\ge1;
\qquad \tau_{2j}=-\tau_j,\quad\tau_{2j+1}=+.
\]
La prueba utiliza \(v_2(2^n+r)=v_2(r)\) para \(0<r<2^n\). Es una recurrencia a toda profundidad, y la coincidencia con los ocho retornos calculados constituye su comprobación finita en la familia.

En el orden lexicográfico con paridad de un mapa con máximo crítico, el factor de orientación previo al lugar \(k\) es \(\prod_{j<k}(-\tau_j)\). La palabra \(\tau\) es estrictamente mayor que cualquiera de sus desplazamientos positivos. En efecto, el primer desacuerdo con un desplazamiento impar ocurre en \(k=1\) o \(2\); un desplazamiento par reduce la comparación a la mitad y duplica \(k\). Por tanto el primer desacuerdo está en \(k=2^r\). Allí
\[
\prod_{j<2^r}(-\tau_j)=(-1)^r=\tau_{2^r},
\]
y la diferencia orientada con el signo opuesto es positiva. Esta demostración también prueba aperiodicidad.

El registro cronológico conserva
\[
K=\rho_9(K)+9q_9(K).
\]
Su lector diádico es
\[
\ell_n(B_K)=[\rho_9(K)+9q_9(K)]_{2^n}=[K]_{2^n}.
\]
Sobre la carta de historias en que \(\Gamma_9 B_K=B_{K+9}\),
\[
\ell_n(\Gamma_9B_K)=\ell_n(B_K)+9\pmod{2^n},\qquad
\pi_{n+1,n}\ell_{n+1}=\ell_n.
\]
Como nueve es impar, el avance por nueve visita las \(2^n\) posiciones. Su conjugación con el avance unitario es \(j\mapsto9j\). Para \(0<j<2^n\),
\[
v_2(9j\bmod2^n)=v_2(j),
\]
de modo que el remuestreo nonádico conserva también la palabra orientada en cada ciclo certificado.

Este lector utiliza residuo \textbf{y} cociente. Las fases de \(K=1\) y \(K=10\) coinciden módulo nueve, mientras sus signos diádicos son opuestos. La memoria distingue las dos historias. La propiedad aritmética del remuestreo vale para todo multiplicador impar; el calendario HMT fija aquí la elección de nueve.


\subsection{Memoria bilateral compatible a profundidad infinita}
\label{hmtfg:prolongacion:5}

Sean \(\mathcal H_n=\ell^2(\mathbb Z/2^n\mathbb Z)\),
\[
C_ne_j=e_{j+1},\qquad
J_ne_j=(e_j+e_{j+2^n})/\sqrt2.
\]
Una comprobación sobre la base da \(J_n^*J_n=I\) y \(C_{n+1}J_n=J_nC_n\), incluidos los términos de borde.

El artículo XI aporta los proyectores
\[
P_0=\frac19\begin{pmatrix}8&-\sqrt8\\-\sqrt8&1\end{pmatrix},\qquad
P_1=\frac19\begin{pmatrix}1&\sqrt8\\\sqrt8&8\end{pmatrix}.
\]
Son ortogonales y complementarios. Por tanto
\[
\mathbb U_n=P_0\otimes I+P_1\otimes C_n
=\begin{pmatrix}T_n&D_n\\D_n&R_n\end{pmatrix}
\]
es unitario, con
\[
T_n=(8I+C_n)/9,\quad D_n=\sqrt8(C_n-I)/9,\quad R_n=(I+8C_n)/9.
\]
La naturalidad ya probada implica
\[
\mathbb U_{n+1}(I_2\otimes J_n)=(I_2\otimes J_n)\mathbb U_n,
\quad \mathbb U_n^a=P_0\otimes I+P_1\otimes C_n^a
\quad(a\in\mathbb Z).
\]
El límite de Hilbert se identifica con \(L^2(\mathbb Z_2,\mu_{\rm Haar})\), llevando \(e_j^{(n)}\) a \(2^{n/2}\mathbf1_{j+2^n\mathbb Z_2}\). Las identidades se prolongan por densidad y acotación. Para todo vector \(v\) y todo tiempo entero \(a\),
\[
\boxed{\|T_{\infty,a}v\|^2+\|D_{\infty,a}v\|^2=\|v\|^2,}
\]
\[
\boxed{v=T_{\infty,a}^*(T_{\infty,a}v)+D_{\infty,a}^*(D_{\infty,a}v).}
\]
Aquí \(T_{\infty,a}=(8I+C_\infty^a)/9\). La composición bilateral conserva el archivo entre pasos; componer solamente \(T_n\) describe otra operación.

La permutación \(S_ne_j=e_{9j\bmod2^n}\) satisface
\[
S_{n+1}J_n=J_nS_n,\qquad S_nC_nS_n^{-1}=C_n^9.
\]
Así se obtiene en el límite
\[
S_\infty C_\infty S_\infty^{-1}=C_\infty^9,
\quad (I_2\otimes S_\infty)\mathbb U_\infty
(I_2\otimes S_\infty^{-1})=\mathbb U_\infty^9.
\]
La representación temporal es fiel: \(C_\infty^a=C_\infty^b\) fuerza \(2^n\mid a-b\) para todo \(n\), de donde \(a=b\). Preparaciones particulares pueden conservar simetrías temporales. La identidad operatoria conserva esta distinción.

El punto proyectivo \(([K]_{2^n})_n\) y un vector de \(L^2\) tienen tipos diferentes: una historia puntual induce una Dirac, mientras las amplitudes normalizables pertenecen al espacio de Hilbert. El resultado construye un lector equivarante de la cronología; las restantes coordenadas nativas del TPK acompañan el registro completo.


\subsection{Existencia y separación de las obligaciones asintóticas}
\label{hmtfg:prolongacion:6}

\subsubsection{Realización de la palabra infinita en toda la franja}
\label{hmtfg:prolongacion:6.1}

\textbf{Teorema de existencia.} Para cada \(|\eta|\le1/40\) existe al menos un
\[
\lambda_\infty(\eta)\in
\left[\frac1{2u_\eta(1)},\frac2{u_\eta(1)}\right]
\]
tal que
\[
\boxed{\operatorname{sign}f_{\lambda_\infty(\eta),\eta}^{j}(0)
=(-1)^{v_2(j)}\quad\text{para todo }j\ge1.}
\]
La conclusión realiza el itinerario infinito de duplicación dentro de la familia fijada. El cuantificador de existencia se mantiene separado de una selección única y de la ley de distancias entre parámetros superestables.

\textbf{Demostración.} Fijamos \(\eta\). Escribimos \(K_\lambda\) para la palabra de signos de la órbita crítica; cuando ésta retorna por primera vez a cero, la palabra termina en ese cero. Usamos el orden unimodal con máximo, cuyo factor antes de comparar el símbolo \(j\) es el producto de los signos \(-K_i\), \(i<j\).

En el extremo izquierdo la imagen de todo el intervalo está en \([1/2,1]\); por tanto \(K_-=+++\cdots\). En el extremo derecho la órbita es \(0\mapsto1\mapsto-1\mapsto-1\), de donde \(K_+=+---\cdots\). Comparando los símbolos segundo y tercero, respectivamente, se obtiene
\[
K_-<\tau<K_+.
\]

Si \(K_{\lambda_0}\) es infinita y distinta de \(\tau\), la primera diferencia ocurre en un lugar finito. La continuidad de los correspondientes iterados hace localmente constante el lado de la comparación. Queda examinar un parámetro superestable, de primer período \(p\), con palabra \(W0\), donde \(|W|=p-1\).

Para \(\lambda\) cercano a \(\lambda_0\), sea \(g_\lambda=f_{\lambda,\eta}^p\) y \(a=g_\lambda(0)\). Se tiene \(g_\lambda'(0)=0\), y una cota local uniforme de la segunda derivada da
\[
g_\lambda(x)=a+O(x^2).
\]
Para \(|a|\) suficientemente pequeño y distinto de cero, \(g_\lambda\) lleva \([-2|a|,2|a|]\) en \([a-|a|/2,a+|a|/2]\). Los retornos conservan el signo de \(a\), y las \(p-1\) posiciones intermedias conservan \(W\). Las palabras vecinas son, por ello, \((W+)^\infty\) o \((W-)^\infty\); los parámetros con \(a=0\) conservan \(W0\).

Si \(\tau\) difiere de \(W\) antes del lugar \(p\), esa diferencia fija el mismo lado para las tres palabras. Si comparte \(W\), pongamos
\[
s=\tau_p,\qquad e=\prod_{j<p}(-\tau_j),\qquad A=Ws.
\]
El signo de la comparación de \(\tau\) con \(W0\) y con \((W,-s)^\infty\) es \(es\). Para compararla con \(A^\infty\), sea \(q\) el primer lugar en que \(\tau_{p+q}\ne\tau_q\). Existe por aperiodicidad. Hasta el lugar \(p+q-1\), ambas palabras coinciden, y el signo de orientación se factoriza como \(O_pO_{q-1}\), con \(O_p=-es\). La maximalidad estricta ya demostrada da
\[
O_{q-1}(\tau_q-\tau_{p+q})>0.
\]
En consecuencia, la comparación de \(\tau\) con \(A^\infty\) tiene signo \(-O_p=es\). Las tres posibilidades locales están al mismo lado de \(\tau\).

Si ningún parámetro realizara \(\tau\), los conjuntos \(\{\lambda:K_\lambda<\tau\}\) y \(\{\lambda:K_\lambda>\tau\}\) serían abiertos, disjuntos, no vacíos y cubrirían el intervalo de parámetros. La conexidad del intervalo excluye esa partición. Existe, por tanto, el parámetro anunciado. \(\square\)

Este argumento de continuidad de itinerarios se aplica después de fijar el generador HMT y su familia analítica; sus hipótesis se verifican aquí directamente. La construcción conserva la palabra infinita sin introducir como dato un parámetro de la familia cuadrática ni el valor universal.

\subsubsection{Márgenes que impiden perder el itinerario al tomar límites}
\label{hmtfg:prolongacion:6.2}

Escribamos \(c_p=f_{\lambda,\eta}^p(0)\). Una cota racional uniforme suficiente es
\[
|f'|,|f''|\le16,\qquad
B_p:=16^p\frac{16^p-1}{15}\ge\|(f^p)''\|_{[-1,1]}.
\]
En efecto, \(1-\cos v\ge v^2/2-v^4/24\ge v^2/8\) para \(|v|\le3\), de modo que \(u_\eta(1)\ge1/4\), \(\lambda\le8\) y ambas derivadas se acotan por \(16\), usando los momentos. La regla de la cadena prueba por inducción la cota \(B_p\).

La palabra \(\tau\) satisface \(\tau_{2p}=-\tau_p\) para todo \(p\). Para cualquiera de sus realizaciones, \(c_{2p}\) y \(c_p\) tienen signos opuestos. Como \((f^p)'(0)=0\), Taylor da
\[
|c_{2p}-c_p|\le\tfrac12B_p|c_p|^2.
\]
El lado izquierdo es mayor que \(|c_p|\), y por tanto
\[
\boxed{|c_p|>2/B_p>0.}
\]
Para cada horizonte fijo existe así un margen uniforme que impide que una sucesión de realizaciones pierda su signo al converger. En las coordenadas compactas \((\eta,b)\), \(b=\lambda u_\eta(1)\in[1/2,2]\), el conjunto de realizaciones de \(\tau\) es cerrado y compacto, y su proyección sobre toda la franja de \(\eta\) es sobreyectiva.

El mismo argumento sirve para un árbol de prefijos: si se realizan todos los prefijos de \(\tau\), una subsucesión compacta conserva cada signo, porque los prefijos suficientemente largos incluyen simultáneamente los lugares \(p\) y \(2p\). El teorema de existencia proporciona ahora la no vacuidad que ese argumento necesita.

\subsubsection{Intervalos invariantes de retorno a todos los niveles}
\label{hmtfg:prolongacion:6.3}

\textbf{Teorema de retornos anidados.} Para cualquiera de las realizaciones del teorema de existencia de la sección \ref{hmtfg:prolongacion:6.1}, sean \(c_j=f^j(0)\) y
\[
J_n=\operatorname{conv}\{c_{2^n},c_{2^{n+1}}\},\qquad n\ge1.
\]
Entonces \(0\in\operatorname{int}J_n\), \(J_{n+1}\subset J_n\),
\[
f^{2^n}(J_n)=J_n,
\]
y las \(2^n\) imágenes \(J_n,f(J_n),\ldots,f^{2^n-1}(J_n)\) son disjuntas. Los retornos normalizados son unimodales de máximo, con valor crítico uno y la misma palabra \(\tau\). Estos intervalos son núcleos invariantes de retorno; una normalización que exija extremos periódicos puede ampliar el núcleo al intervalo restrictivo maximal correspondiente y debe especificar ese paso.

\textbf{Demostración del primer retorno.} Sea \(F\) un mapa unimodal de máximo con \(F(0)=1\) y palabra \(\tau\), y \(d_j=F^j(0)\). Sus primeros signos son \(+,-,+,+,+,-,+\). El intervalo \(I=[d_2,1]\) es invariante: \(F(d_2)=d_3>0\), \(F(1)=d_2<0\), y el máximo es uno. Pongamos \(J=[d_2,d_4]\). Como \(F\) decrece en positivos y
\[
F(d_3)=d_4>0> d_6=F(d_5),
\]
se tiene \(d_3<d_5\), y por tanto \(F(J)=[d_3,1]\).

Además \(d_3>d_4\). Para comprobarlo, \(F^2\) es estrictamente creciente entre los dos puntos positivos \(d_3,d_4\), porque sus imágenes \(d_4,d_5\) son positivas. Sus valores son \(F^2(d_3)=d_5>0>d_6=F^2(d_4)\), lo que fuerza ese orden. Así \(J\) y \(F(J)\) quedan separados por un intervalo abierto, y
\[
F^2(J)=F([d_3,1])=[d_2,d_4]=J.
\]
El mapa \(F^2|_J\) tiene un único mínimo crítico: la imagen \(F(J)\) permanece positiva. La conjugación \(h(x)=d_2x\) invierte orientación y normaliza el retorno como máximo con valor uno. Su itinerario es
\[
\operatorname{sign}\frac{d_{2j}}{d_2}
=-\tau_{2j}=\tau_j.
\]
Se puede repetir el mismo argumento a cualquier profundidad.

\textbf{Disjunción inductiva.} Supónganse disjuntas las \(p=2^n\) imágenes del padre \(J_n\), y sea \(g=f^p\). El primer retorno aplicado al mapa normalizado produce \(B=J_{n+1}\), \(A=g(B)\), con \(A\cap B=\varnothing\), \(g(B)=A\), \(g(A)=B\). Para \(0\le j<p\), una intersección de \(f^j(A)\) y \(f^j(B)\), al aplicarle \(f^{p-j}\), produciría una intersección de \(B\) y \(A\). Por tanto esas dos imágenes son disjuntas. Las imágenes correspondientes a distintos \(j\) pertenecen a imágenes disjuntas del padre. Quedan demostradas las \(2p\) imágenes disjuntas y el retorno \(f^{2p}(B)=B\). La fórmula de los extremos resulta de componer las normalizaciones \(x\mapsto c_{2^n}x\). \(\square\)

Los conjuntos compactos
\[
\Lambda_n=\bigcup_{0\le j<2^n}f^j(J_n)
\]
son no vacíos y anidados. Su intersección \(\Lambda_\infty\) tiene un lector continuo sobre el odómetro diádico: a cada punto se le asigna la etiqueta de la única componente que ocupa a cada nivel. Todas las sucesiones compatibles de etiquetas tienen preimagen por compacidad, y el avance por \(f\) suma uno en las etiquetas. Esta aplicación es una semiconjugación sobreyectiva. Para promoverla a conjugación puntual se requiere demostrar que los diámetros de las componentes tienden a cero; esa promoción métrica se mantiene separada.

% Fragmento integrable. Preambulo anfitrion: amsmath, amssymb, longtable, hyperref.
\section{Compatibilidad de lectores y restitución del parámetro}
\label{hmtfg:fragmento:compatibilidad-de-lectores-y-restitucion-del-parametro}

\subsection{Compatibilidad de lectores, mínimos y recuperación de la historia}
\label{hmtfg:escalado:13}

Las operaciones de prolongación y los lectores ordenados conservan la información requerida por la selección de raíces. Las siguientes identidades fijan sus dominios y su composición.

\subsubsection{Regla forward recuperada}
\label{hmtfg:escalado:13.1}

La actualización determinista parte del estado de frontera R36 con carácter \(\chi\) y residuo interno. La actualización es
\begin{equation}
\mathcal U_K^\chi=
\operatorname{Upd}_K^\chi\circ
\operatorname{Tra}_K^\chi\circ
\operatorname{Sel}_K^\chi,
\qquad
N_{K+1}=729N_K+d_K^\chi.
\label{hmtfg:escalado:eq:26}
\end{equation}
El dígito y la actualización dependen del estado precedente. Dos secciones forward con el mismo estado completo R36 coinciden por inducción. La subrelación \(\operatorname{Ext}^{\chi,\rightarrow}\) es el grafo de esa actualización; la relación total \(\operatorname{Ext}\) admite multiplicidad de continuaciones según los datos y caracteres.

Las prolongaciones compatibles transportan los lectores regionales y la firma conjunta sobre los mismos prefijos.

Estas reglas transportan orientación, residuo, cociente, acarreo, frontera y memoria. Su aplicación a la cascada debe conservar además la selección ordenada del parámetro del lector crítico.

\subsubsection{Independencia de las constantes antecedentes respecto del parámetro libre}
\label{hmtfg:escalado:13.2}

Sea \(B\) el estado HMT que proporciona la rueda y sus constantes antecedentes. Escribamos \(\mathcal C(B)\) para el conjunto de esas publicaciones y \(u_B\) para el perfil crítico. En el dominio del lector crítico, la construcción posterior es
\begin{equation}
(B,\lambda)\longmapsto f_{B,\lambda}(x)=1-\lambda u_B(x),
\qquad 0<\lambda\le 2/u_B(1).
\label{hmtfg:escalado:eq:27}
\end{equation}
Para un mismo \(B\), la lectura de las constantes anteriores se factoriza por la proyección al primer componente:
\begin{equation}
\widetilde{\mathcal C}(B,\lambda)
=\mathcal C\circ\operatorname{pr}_B(B,\lambda)
=\mathcal C(B).
\label{hmtfg:escalado:eq:28}
\end{equation}
En consecuencia, dos valores distintos de \(\lambda\) conservan las mismas publicaciones antecedentes. La igualdad es una propiedad de esta construcción causal: el perfil se fija antes de variar \(\lambda\), y el factor \(\pi_{\mathrm{HMT}}\) se cancela en su normalización cuadrática.

\textbf{Lema.} Supóngase que una condición \(A\) se evalúa únicamente sobre esas publicaciones antecedentes. Para todo \(B\) fijo,
\begin{equation}
\{\lambda\in\Lambda:A(\widetilde{\mathcal C}(B,\lambda))\}
=
\begin{cases}
\Lambda,&A(\mathcal C(B))\text{ se cumple},\\
\varnothing,&A(\mathcal C(B))\text{ falla}.
\end{cases}
\label{hmtfg:escalado:eq:29}
\end{equation}
\textbf{Demostración.} Sustituir \eqref{hmtfg:escalado:eq:28} hace independiente de \(\lambda\) el predicado. Por tanto, lo satisfacen todos los parámetros del dominio o ninguno. Esto incluye una comparación de las constantes antecedentes con intervalos metrológicos.

La unicidad de la sección \(B\) y la restricción posterior \(f_{B,\lambda}^{2^n}(0)=0\) producen una fibra de soluciones en \(\lambda\). La regla de primera raíz actúa sobre esta fibra mediante el retorno y su itinerario. La selección y su prolongación se demuestran en los resultados siguientes.

Este resultado mantiene el orden HMT: la metrología contrasta salidas previamente construidas y la selección de la raíz se demuestra en el lector que la produce.

\subsubsection{Lema de transporte ordenado del selector}
\label{hmtfg:escalado:13.3}

Sea \(P_n\) un conjunto ordenado de candidatos genealógicos con mínimo \(p_n^*\). Sea \(Z_n\) el conjunto original de raíces candidatas del corpus, con el orden real, y sea \(L_n:P_n\to Z_n\) un lector monótono. Se dice que su imagen es \textbf{coinicial} cuando
\begin{equation}
\forall z\in Z_n\quad\exists p\in P_n:
L_n(p)\le z.
\label{hmtfg:escalado:eq:30}
\end{equation}
Entonces
\begin{equation}
\boxed{L_n(p_n^*)=\min Z_n
\quad\Longleftrightarrow\quad
L_n(P_n)\text{ es coinicial en }Z_n.}
\label{hmtfg:escalado:eq:31}
\end{equation}
\textbf{Demostración.} Si la imagen del mínimo es el mínimo global, se toma \(p=p_n^*\) en \eqref{hmtfg:escalado:eq:30}. Recíprocamente, dado \(z\in Z_n\), la coinicialidad proporciona \(p\) con \(L_n(p)\le z\). La monotonía da \(L_n(p_n^*)\le L_n(p)\le z\). Como \(L_n(p_n^*)\in Z_n\), es el mínimo global. La sobreyectividad de \(L_n\) constituye una condición suficiente, más fuerte que \eqref{hmtfg:escalado:eq:30}.

El isomorfismo ordenado de la carta del continuo satisface la sobreyectividad del lector. La aplicación concreta al retorno y a su conjunto completo de ceros se demuestra en la sección de transporte ordenado; la clasificación y el aislamiento determinan después su continuación dinámica.

\subsubsection{Selección mínima mediante supervivencia a profundidad arbitraria}
\label{hmtfg:escalado:13.6}

La composición con la supervivencia produce además un selector global preciso. Se mantiene la familia uniforme \(\eta=0\), construida en sección~\ref{hmtfg:escalado:1}, y se utilizan los teoremas de realización del itinerario infinito y conservación de sus signos. Sean
\[
I=\left[\frac1{2u_0(1)},\frac2{u_0(1)}\right],\qquad
c_j(\lambda)=f_\lambda^j(0),\qquad
\tau_j=(-1)^{v_2(j)}.
\]
El conjunto
\[
\Lambda_\tau=\{\lambda\in I:\operatorname{sign}c_j(\lambda)=\tau_j
\text{ para todo }j\ge1\}
\]
es compacto y no vacío. El control de derivadas de ese propietario proporciona
\begin{equation}
B_j=\frac{16^j(16^j-1)}{15},\qquad
|c_j(\lambda)|>\frac2{B_j}\quad(\lambda\in\Lambda_\tau).
\label{hmtfg:escalado:eq:32}
\end{equation}
Definamos
\begin{equation}
\varepsilon_j=\frac1{B_j},\qquad
A_N=\{\lambda\in I:\tau_jc_j(\lambda)\ge\varepsilon_j,\
1\le j\le N\}.
\label{hmtfg:escalado:eq:33}
\end{equation}
Cada \(A_N\) es compacto por continuidad de los iterados y no vacío porque contiene \(\Lambda_\tau\). Además,
\begin{equation}
A_{N+1}\subseteq A_N,\qquad
\bigcap_{N\ge1}A_N=\Lambda_\tau.
\label{hmtfg:escalado:eq:34}
\end{equation}
La inclusión de derecha a izquierda se sigue de \eqref{hmtfg:escalado:eq:32}. La inclusión inversa se sigue de los márgenes estrictamente positivos de \eqref{hmtfg:escalado:eq:33}, que fijan cada signo.

\textbf{Teorema del mínimo superviviente.} Los mínimos \(a_N=\min A_N\) satisfacen
\begin{equation}
\boxed{a_N\uparrow a_*=\min\Lambda_\tau.}
\label{hmtfg:escalado:eq:35}
\end{equation}
\textbf{Demostración.} La inclusión \eqref{hmtfg:escalado:eq:34} hace creciente la sucesión de mínimos, acotada dentro de \(I\); tiene, por tanto, un límite \(a\). Para cada \(r\), todos los términos \(a_N\), \(N\ge r\), pertenecen al cerrado \(A_r\); así \(a\in A_r\). Por \eqref{hmtfg:escalado:eq:34}, \(a\in\Lambda_\tau\). Para cualquier \(\lambda\in\Lambda_\tau\), se tiene \(a_N\le\lambda\) a todo nivel, y al pasar al límite, \(a\le\lambda\). Esto prueba \eqref{hmtfg:escalado:eq:35}.

En el lector cilíndrico HMT, el objeto \(a_*\) posee prefijos compatibles a toda profundidad, con orientación, residuo y memoria conservados. Su elección utiliza el orden de la lectura paramétrica y la supervivencia completa. La compacidad demuestra la existencia y convergencia de estos mínimos.

Los parámetros \(a_N\) resuelven restricciones de signos con margen. Sus fronteras pueden satisfacer \(\tau_jc_j=\varepsilon_j\), mientras las raíces superestables satisfacen \(c_{2^n}=0\). Por ello \eqref{hmtfg:escalado:eq:35} establece un selector global de realización infinita y preserva, como objeto distinto, el selector de primeras raíces del corpus.

\subsubsection{Recuperación del parámetro desde la historia crítica completa}
\label{hmtfg:escalado:13.7}

El segundo término de la historia crítica proporciona un lector inverso exacto:
\begin{equation}
c_1(\lambda)=1,\qquad
c_2(\lambda)=1-\lambda u_0(1),\qquad
\boxed{\lambda=\frac{1-c_2(\lambda)}{u_0(1)}.}
\label{hmtfg:escalado:eq:36}
\end{equation}
La positividad \(u_0(1)>0\), demostrada para el perfil, hace válida la división. En consecuencia, dos historias críticas completas iguales tienen el mismo parámetro. La palabra de signos \(\tau\) constituye una publicación reducida de esa historia; conserva su combinatoria y omite la magnitud de \(c_2\).

Las fórmulas \eqref{hmtfg:escalado:eq:35}–\eqref{hmtfg:escalado:eq:36} unen la selección por supervivencia y la restitución del parámetro desde la historia crítica. La igualdad de su límite con el cruce estable se demuestra por primera salida, y la identificación de los centros por el transporte analítico de sus clases híbridas.

% Fragmento integrable. Preambulo anfitrion: amsmath, amssymb, longtable, hyperref.
\section{Transporte ordenado y representación de Jacobi}
\label{hmtfg:fragmento:transporte-ordenado-y-representacion-de-jacobi}

\subsection{Transporte completo desde el continuo y realización espectral del perfil}
\label{hmtfg:escalado:14}

El transporte ordenado del continuo y la representación espectral del lector se componen antes de la selección paramétrica.

\subsubsection{Cobertura de todas las raíces y transporte de su mínimo}
\label{hmtfg:escalado:14.1}

La construcción ordenada del continuo desarrollada en los antecedentes produce la publicación arquimediana sobreyectiva desde cilindros compatibles y su cociente ordenado. Dentro de su universo declarado \(\mathfrak G\), demuestra el isomorfismo de cuerpos ordenados completos
\begin{equation}
\iota:\mathbb R_{\rm HMT}^{\mathfrak G}
   \xrightarrow{\cong}\mathbb R^{\mathfrak G}.
\label{hmtfg:escalado:eq:37}
\end{equation}
Se cambia aquí el nombre \(\eta\) del propietario por \(\iota\), para distinguirlo del parámetro de ponderación, fijado a cero. La segunda proyección conserva incidencia, frontera y memoria sobre el mismo antecedente; el cociente de evaluación y la historia completa mantienen sus dominios respectivos.

Construyamos en el cuerpo HMT, antes de seleccionar una raíz,
\[
\cos_{\rm H}(t)=\sum_{k=0}^{\infty}\frac{(-1)^kt^{2k}}{(2k)!},
\qquad
u_{\rm H}(x)=\frac1{12}\sum_{m=1}^{12}
\left[1-\cos_{\rm H}\!\left(\frac{2(3m-1)x}{\sqrt[\rm H]{899}}\right)\right],
\]
\begin{equation}
F_a(x)=1-a\,u_{\rm H}(x),\qquad S_n^{\rm H}(a)=F_a^{\,2^n}(0).
\label{hmtfg:escalado:eq:38}
\end{equation}
Los enteros y el segundo momento son los ya derivados del lector dodecafásico. La raíz cuadrada es la positiva del cuerpo ordenado. La serie convergente realiza analíticamente ese lector; conserva el perfil y el coeficiente normalizador de sección~\ref{hmtfg:escalado:1}.

\textbf{Teorema de transporte del retorno.} Para todo parámetro del dominio interno de \eqref{hmtfg:escalado:eq:38},
\begin{equation}
\iota(F_a(x))=f_{\iota(a)}(\iota(x)),\qquad
\iota(S_n^{\rm H}(a))=S_n(\iota(a)).
\label{hmtfg:escalado:eq:39}
\end{equation}
\textbf{Demostración.} El isomorfismo ordenado fija los racionales y preserva la raíz positiva. Preserva los límites convergentes: cada entorno de radio racional positivo se transporta a otro del mismo radio. Aplicado a las sumas parciales de \eqref{hmtfg:escalado:eq:38}, transporta el coseno y \(u_{\rm H}\). La primera igualdad resulta de conservar suma y producto; la segunda, por inducción sobre el número de iteraciones.

Fijado un centro anterior \(a_{n-1}\), definamos \(Z_n^{\rm H}\) por retorno cero, período crítico exacto \(2^n\) y \(a>a_{n-1}\), dentro del dominio paramétrico declarado. Sea \(Z_n\) el conjunto definido por las mismas condiciones para \(f_\lambda\), a la derecha de \(\iota(a_{n-1})\). Entonces
\begin{equation}
\boxed{\iota[Z_n^{\rm H}]=Z_n,\qquad
\iota(\min Z_n^{\rm H})=\min Z_n.}
\label{hmtfg:escalado:eq:40}
\end{equation}
La segunda igualdad se afirma cuando existe el mínimo; la primera preserva también la no existencia de candidatos.

\textbf{Demostración.} \eqref{hmtfg:escalado:eq:39} preserva el retorno cero y los retornos anteriores que determinan el período exacto. El orden conserva la posición respecto del centro previo. La sobreyectividad de \eqref{hmtfg:escalado:eq:37}, aplicada a cualquier parámetro raíz, prueba la cobertura inversa. Si \(a_*\) es el mínimo, cualquier \(z\in Z_n\) tiene la forma \(\iota(a)\), con \(a\in Z_n^{\rm H}\), de donde \(\iota(a_*)\le z\). La afirmación inversa utiliza \(\iota^{-1}\).

La sobreyectividad cubre todos los parámetros de la carta, incluidos aquellos cuyas raíces no se hayan aislado por un cálculo finito. El isomorfismo conserva el universo de interpretación declarado y el orden de selección.

\subsubsection{Representación de Jacobi del perfil dodecafásico completo}
\label{hmtfg:escalado:14.2}

Definamos la medida atómica del lector,
\begin{equation}
s_m=\frac{4(3m-1)^2}{899},\qquad
\nu_0=\frac1{12}\sum_{m=1}^{12}\delta_{s_m},\qquad
M_r=\int s^r\,d\nu_0(s).
\label{hmtfg:escalado:eq:41}
\end{equation}
Los doce nodos son distintos y positivos. Para un polinomio \(p\ne0\) de grado menor que \(r\le12\),
\begin{equation}
\sum_{i,j=0}^{r-1}p_iM_{i+j}p_j
=\frac1{12}\sum_{m=1}^{12}p(s_m)^2>0.
\label{hmtfg:escalado:eq:42}
\end{equation}
Un polinomio de grado menor que doce no se anula en los doce nodos. En grado doce aparece \(\prod_m(s-s_m)\), de norma nula.

La ortogonalización respecto de esta medida produce polinomios ortonormales \(\psi_0,\ldots,\psi_{11}\), con \(\psi_0=1\). Sean
\[
Q_{mj}=\psi_j(s_m)/\sqrt{12},\quad
D=\operatorname{diag}(s_m),\quad J_{12}=Q^{\mathsf T}DQ.
\]
Se cumplen
\begin{equation}
Q^{\mathsf T}Q=I,\quad DQ=QJ_{12},\qquad
\boxed{u_0(x)=e_0^{\mathsf T}
\bigl[I-\cos(x\sqrt{J_{12}})\bigr]e_0.}
\label{hmtfg:escalado:eq:43}
\end{equation}
\textbf{Demostración.} La ortogonalidad procede de \eqref{hmtfg:escalado:eq:42}. La multiplicación por \(s\) produce una recurrencia de tres términos: los elementos separados por más de una posición se anulan por ortogonalidad y grado. Fijar positivos los coeficientes principales hace positivas las subdiagonales. Por cálculo funcional,
\(\cos(x\sqrt{J_{12}})=Q^{\mathsf T}\cos(x\sqrt D)Q\).
El vector \(Qe_0\) es constante, de valor \(1/\sqrt{12}\); sustituirlo recupera la suma de sección~\ref{hmtfg:escalado:1}.

El artículo de Grassmannianos desarrolla este mecanismo para medidas marcadas y refinamientos. Aquí se aplica a \eqref{hmtfg:escalado:eq:41}; el rango doce pertenece a este lector. Las marcas genealógicas se conservan junto al lector, conforme a la fidelidad marcada del propietario. \(J_{12}\) representa el perfil analítico; el estado completo conserva adicionalmente rutas, orientación y memoria.

El verificador de esta ampliación comprueba con fracciones exactas los doce grados de norma positiva, la anulación del grado doce, el entrelazador \(VT=DV\) en base mónica y la autoadjunción ponderada \(HT=T^{\mathsf T}H\). Comprueba los momentos de grado cero a treinta; la identidad para todo grado se sigue del entrelazador. Sus primeros coeficientes son
\begin{equation}
a_0=2,\qquad \beta_1=\frac{2493348}{808201}.
\label{hmtfg:escalado:eq:44}
\end{equation}

\subsubsection{Sensibilidad orientada del retorno}
\label{hmtfg:escalado:14.3}

Para un retorno exacto de período \(p=2^n\), sean
\[
c_j=f_\lambda^j(0),\quad q_j=\partial_\lambda c_j,\quad
D_j=\prod_{k=1}^j f_\lambda'(c_k),\quad D_0=1.
\]
Antes del retorno crítico, los factores de \(D_{p-1}\) son distintos de cero. Diferenciación y telescopaje dan
\begin{equation}
q_{j+1}=-u_0(c_j)+f_\lambda'(c_j)q_j,\qquad
\boxed{\frac{q_p}{D_{p-1}}
=-\sum_{j=1}^{p-1}\frac{u_0(c_j)}{D_j}.}
\label{hmtfg:escalado:eq:45}
\end{equation}
En la palabra canónica, \(\operatorname{sign}c_j=(-1)^{v_2(j)}\). El signo de \(f_\lambda'(c_j)\) es el opuesto; por tanto,
\begin{equation}
\operatorname{sign}D_j=(-1)^{j+\sum_{k=1}^jv_2(k)}
=(-1)^{s_2(j)},
\label{hmtfg:escalado:eq:46}
\end{equation}
usando \(\sum_{k=1}^jv_2(k)=j-s_2(j)\), donde \(s_2(j)\) cuenta los unos binarios. La sensibilidad normalizada es
\begin{equation}
\mathcal T_n=-\sum_{j=1}^{2^n-1}(-1)^{s_2(j)}t_j,\qquad
t_j=\frac{u_0(c_j)}{|D_j|}>0.
\label{hmtfg:escalado:eq:47}
\end{equation}
La positividad nodal de \eqref{hmtfg:escalado:eq:42} determina \(u_0\); \eqref{hmtfg:escalado:eq:47} conserva además los productos orbitales y sus orientaciones.

La representación adicional por momentos exigiría una medida positiva \(\rho\) en \([0,1]\) con \(t_j=\int z^{j-1}\,d\rho\). Bajo esa condición,
\begin{equation}
\mathcal T_n=\int_0^1
\frac{1-\prod_{r=0}^{n-1}(1-z^{2^r})}{z}\,d\rho(z)>0
\label{hmtfg:escalado:eq:48}
\end{equation}
para \(\rho\ne0\), extendiendo el integrando por su valor \(1\) en cero. La identidad se obtiene expandiendo el producto binario.

El control racional en la raíz de período cuatro previamente certificada obtiene
\[
t_1\approx0.40425224267600139730,\quad
t_2\approx0.04925249167432646403,\quad
t_3\approx0.15740870098294833737
\]
y el intervalo exterior riguroso
\begin{equation}
\begin{gathered}
0.296096033367379523967764444238832345360107<\mathcal T_2,\\
\mathcal T_2<0.296096033367379523967764444238832345360112.
\end{gathered}
\label{hmtfg:escalado:eq:49}
\end{equation}
También certifica \(t_3>t_2\). Una medida positiva sobre \([0,1]\) impondría \(t_3\le t_2\); ese levantamiento concreto de los pesos orbitales a momentos positivos queda descartado. El signo positivo de \eqref{hmtfg:escalado:eq:49} permanece certificado. Este control distingue el fallo de una prueba propuesta del fallo de la propiedad que se pretendía probar.

% Fragmento integrable. Preambulo anfitrion: amsmath, amssymb, longtable, hyperref.
\section{Clasificación y continuación del selector de primeras raíces}
\label{hmtfg:fragmento:clasificacion-y-continuacion-del-selector-de-primeras-raices}

\subsection{Objeto y procedencia}
\label{hmtfg:clasificacion:1}

Este apéndice conserva la familia publicada por el lector dodecafásico HMT:

\[
u_0(x)=\frac1{12}\sum_{m=1}^{12}\left(1-\cos\frac{2(3m-1)x}{\sqrt{899}}\right),
\qquad f_\lambda(x)=1-\lambda u_0(x).
\]

La familia procede de los antecedentes operatorios y analíticos expuestos en este desarrollo. La existencia y la compacidad de sus realizaciones del itinerario infinito se han demostrado en las subsecciones de realización de la palabra y conservación de sus signos.

El argumento siguiente es propio y utiliza exclusivamente esas propiedades y las operaciones de retorno de la familia fijada. No utiliza como entrada un parámetro de otra familia ni el valor de una constante universal. Las hojas de los itinerarios aquí utilizados son las dos ramas monótonas del lector crítico; las restantes coordenadas y la memoria HMT conservan su residencia en el estado del que procede el lector.

Escribimos
\[
c_j(\lambda)=f_\lambda^j(0),\qquad
\tau_j=(-1)^{v_2(j)},\qquad
W_n=\tau_1\cdots\tau_{2^n-1}.
\]
El itinerario de una órbita crítica que retorna por primera vez al origen termina en el símbolo cero. El orden de itinerarios es el orden unimodal de máximo: antes de comparar el símbolo de índice \(j\), el factor de orientación es
\[
O_{j-1}(K)=\prod_{i=1}^{j-1}(-K_i).
\]
El orden ordinario de los símbolos es \(-1<0<1\).


\subsection{Clasificación anterior a la palabra infinita}
\label{hmtfg:clasificacion:2}

\textbf{Teorema de clasificación.} Sea \(F:I\to I\), con \(I=[\ell,1]\), \(\ell<0\), un mapa continuo, estrictamente creciente a la izquierda de cero y estrictamente decreciente a su derecha, cuyo máximo satisface \(F(0)=1\). Supóngase que cero tiene período exacto \(2^n\), con \(n\ge1\), y que su itinerario crítico \(K\) satisface \(K<\tau\). Entonces
\[
\boxed{K=W_n0.}
\]
La conclusión se refiere al itinerario, y permite que parámetros distintos de una familia tengan ese mismo itinerario.

\textbf{Demostración.} Denotemos \(d_j=F^j(0)\). Para período dos el itinerario es \(+0=W_10\). Consideremos un período exacto \(p=2^n\ge4\).

Se tiene \(d_2<0\). En efecto, \(d_2=0\) daría período dos. Si \(d_2>0\), la monotonía de la rama positiva da
\[
F([d_2,1])=[d_2,d_3]\subset[d_2,1],
\]
y la órbita crítica permanece en un intervalo estrictamente positivo, incompatible con el retorno a cero.

Los signos iniciales son entonces \(+,-\). El factor de orientación para comparar el tercer signo es \(-1\); por ello \(K<\tau\) obliga a \(d_3>0\). Un cero en ese lugar produciría período tres.

Tampoco puede ocurrir \(d_4<0\). Bajo esta hipótesis,
\[
d_2<d_4<0,
\qquad F([d_2,d_4])=[d_3,d_5]\subset(0,1],
\]
pues \(d_4=F(d_3)>F(1)=d_2\) y \(d_5=F(d_4)>F(d_2)=d_3\). La rama positiva es decreciente, de modo que
\[
F^2([d_2,d_4])=[d_6,d_4]\subset[d_2,d_4].
\]
Aquí \(d_6\ge d_2\), porque \(d_5\le1\), y \(d_6<d_4\), porque \(d_5>d_3\). Este intervalo estrictamente negativo atrapa los retornos pares de la órbita crítica, y excluye el retorno a cero. Por tanto \(d_4\ge0\). Si \(p=4\), se obtiene precisamente \(+-+0=W_20\).

Supongamos ahora \(p\ge8\). Entonces \(d_4>0\). Los factores de orientación para los signos quinto y sexto son, respectivamente, \(-1\) y \(+1\). Como \(\tau_5=+1\) y \(\tau_6=-1\), la desigualdad \(K<\tau\) obliga a
\begin{equation}
\operatorname{sign}(d_1,\ldots,d_6)=(+,-,+,+,+,-).
\label{hmtfg:clasificacion:eq:1}
\end{equation}
Los signos cero en esos lugares están excluidos por el período exacto.

Construimos el intervalo de retorno
\[
J=[d_2,d_4].
\]
La desigualdad \(F(d_3)=d_4>0>d_6=F(d_5)\), junto con \(d_3,d_5>0\), implica \(d_3<d_5\); por tanto
\[
F(J)=[d_3,1].
\]
Además \(d_4<d_3\). Para comprobarlo, la imagen por \(F\) del intervalo de extremos \(d_3,d_4\) es estrictamente positiva, pues sus imágenes extremas son \(d_4,d_5>0\). En ese intervalo \(F^2\) es estrictamente creciente. Sus valores extremos satisfacen
\[
F^2(d_3)=d_5>0>d_6=F^2(d_4),
\]
lo que fuerza \(d_3>d_4\). En consecuencia,
\begin{equation}
J\cap F(J)=\varnothing,
\qquad F^2(J)=F([d_3,1])=[d_2,d_4]=J.
\label{hmtfg:clasificacion:eq:2}
\end{equation}

El mapa \(F^2|_J\) posee un único mínimo en cero. La conjugación por \(h(x)=d_2x\), que invierte la orientación, define
\begin{equation}
G=h^{-1}\circ F^2\circ h:
\left[\frac{d_4}{d_2},1\right]\longrightarrow
\left[\frac{d_4}{d_2},1\right].
\label{hmtfg:clasificacion:eq:3}
\end{equation}
Este mapa satisface las mismas hipótesis de máximo que \(F\), con \(G(0)=1\), y su crítico tiene período exacto \(p/2\). Su itinerario es
\begin{equation}
K'_j=-K_{2j}.
\label{hmtfg:clasificacion:eq:4}
\end{equation}
Todos los símbolos impares de \(K\) son positivos, porque las posiciones pares pertenecen a \(J\) y las impares a \(F(J)\subset(0,1]\).

La transformación \eqref{hmtfg:clasificacion:eq:4} conserva el orden de comparación con \(\tau\). En efecto, \(\tau_{2j-1}=+1\) y \(\tau_{2j}=-\tau_j\). Si \(q\) es el primer lugar en que \(K'\) y \(\tau\) difieren, el primer lugar de diferencia entre \(K\) y \(\tau\) es \(2q\), y
\begin{equation}
O_{2q-1}(K)=-O_{q-1}(K'),
\qquad K_{2q}-\tau_{2q}=-(K'_q-\tau_q).
\label{hmtfg:clasificacion:eq:5}
\end{equation}
El producto que determina el orden es idéntico. Las identidades siguen siendo válidas cuando el primer símbolo distinto es el cero final. Así \(K'<\tau\).

La inducción aplicada a \(G\) da \(K'=W_{n-1}0\). Las posiciones impares positivas y \eqref{hmtfg:clasificacion:eq:4} reconstruyen exactamente \(K=W_n0\). Queda demostrado el teorema. \(\square\)

El teorema demuestra también que cada mapa clasificado posee los retornos restrictivos de duplicación anteriores a su retorno superestable, con los dominios de \eqref{hmtfg:clasificacion:eq:2}–\eqref{hmtfg:clasificacion:eq:3}. No presupone una monotonía del parámetro de una familia.

\textbf{Normalización simétrica de la familia fijada.} Cuando \(F\) es par y está definido sobre \([-1,1]\), las desigualdades anteriores dan además
\[
F(d_4)=d_5>d_3=F(d_2)=F(-d_2).
\]
Como ambos argumentos \(d_4,-d_2\) son positivos y \(F\) decrece en esa rama,
\begin{equation}
0<d_4<-d_2<1.
\label{hmtfg:clasificacion:eq:5a}
\end{equation}
Por tanto \(F^2([d_2,-d_2])=[d_2,d_4]\), y la misma conjugación \(h(x)=d_2x\) define sobre todo \([-1,1]\) un mapa par de máximo cuya imagen es \([d_4/d_2,1]\subset[-1,1]\). Coincide exactamente con el operador \(RF=-F(F(-ax))/a\), \(a=-F(1)\), utilizado en el corpus. Esta ampliación simétrica conserva el retorno crítico, su período y su itinerario.


\subsection{Primer parámetro de itinerario infinito}
\label{hmtfg:clasificacion:3}

Trabajamos sobre el intervalo compacto ya utilizado en la demostración de existencia,
\[
L=\left[\frac1{2u_0(1)},\frac2{u_0(1)}\right].
\]
Sea
\[
\Lambda_\tau=\{\lambda\in L:
\operatorname{sign}c_j(\lambda)=\tau_j\text{ para todo }j\ge1\}.
\]
Por las demostraciones de existencia y separación de signos de la prolongación, este conjunto es no vacío y compacto. Por tanto existe
\begin{equation}
a=\min\Lambda_\tau.
\label{hmtfg:clasificacion:eq:6}
\end{equation}

\textbf{Lema de comparación.} Para todo \(\lambda\in[1/(2u_0(1)),a)\), se tiene \(K_\lambda<\tau\).

\textbf{Demostración.} En el extremo izquierdo, el itinerario es \(+^\infty<\tau\). El argumento de separación de la realización de la palabra infinita demuestra que los conjuntos de comparación estricta con \(\tau\) son abiertos: en un itinerario infinito distinto de \(\tau\), la primera diferencia finita persiste por continuidad; en un parámetro superestable, las dos palabras periódicas vecinas y la palabra con cero final se sitúan en el mismo lado de \(\tau\), por su maximalidad estricta. El intervalo anterior a \(a\) carece, por definición, de parámetros con itinerario \(\tau\). Su conexidad mantiene la comparación inicial en todo ese intervalo. \(\square\)


\subsection{Existencia de cada primera raíz nueva}
\label{hmtfg:clasificacion:4}

Conservamos el selector original. La primera raíz es
\[
\lambda_1=\frac1{u_0(1)},
\]
y, para \(n\ge2\), \(\lambda_n\) es la primera raíz de \(c_{2^n}\) situada a la derecha de \(\lambda_{n-1}\) que satisface
\begin{equation}
c_{2^j}(\lambda_n)\ne0\quad(1\le j<n).
\label{hmtfg:clasificacion:eq:7}
\end{equation}
Como cualquier primer período que divida \(2^n\) es una potencia de dos, \eqref{hmtfg:clasificacion:eq:7} equivale a exigir período crítico exacto \(2^n\).

\textbf{Lema de existencia.} Todas las raíces de este selector existen y cumplen
\begin{equation}
\lambda_1<\lambda_2<\cdots<\lambda_n<a.
\label{hmtfg:clasificacion:eq:8}
\end{equation}

\textbf{Demostración.} El signo \(c_2(a)<0\) implica \(\lambda_1<a\). Supongamos construida \(\lambda_{n-1}<a\), y pongamos \(p=2^{n-1}\). La función analítica \(c_p(\lambda)\) no es idénticamente cero, puesto que \(c_p(a)\ne0\). Su conjunto de ceros en \([\lambda_{n-1},a]\) es finito, no vacío y carece de \(a\). Sea \(z<a\) su último cero.

En \((z,a]\), el signo de \(c_p\) es constante e igual a \(s=\tau_p\). Para \(g_\lambda=f_\lambda^p\), se tiene
\[
g_\lambda(0)=c_p(\lambda),\qquad g'_\lambda(0)=0.
\]
Una cota uniforme local de la segunda derivada da
\begin{equation}
c_{2p}(\lambda)
=g_\lambda(c_p(\lambda))
=c_p(\lambda)+O(c_p(\lambda)^2).
\label{hmtfg:clasificacion:eq:9}
\end{equation}
Al aproximarse \(\lambda\) a \(z\) por la derecha, el término de primer orden domina; por tanto \(c_{2p}(\lambda)\) tiene signo \(s\). En \(a\), su signo es
\[
\tau_{2p}=-\tau_p=-s.
\]
El teorema del valor intermedio proporciona un cero \(\beta\in(z,a)\) de \(c_{2p}\). Como \(c_p\ne0\) en ese intervalo, el período crítico en \(\beta\) es exactamente \(2p\).

La analiticidad y \(c_{2p}(a)\ne0\) implican que los ceros de \(c_{2p}\) en \([\lambda_{n-1},a]\) son finitos. Entre sus ceros de período exacto \(2p\), situados a la derecha de \(\lambda_{n-1}\), existe así un primero. Es el selector original \(\lambda_n\), y satisface \(\lambda_n\le\beta<a\). \(\square\)

La utilización del último cero auxiliar \(z\) pertenece únicamente a la prueba de existencia. El parámetro seleccionado continúa siendo la \textbf{primera raíz nueva} \(\lambda_n\).


\subsection{Convergencia del selector global al primer itinerario infinito}
\label{hmtfg:clasificacion:5}

\textbf{Teorema de continuación global.} Para la familia HMT fijada y el selector original, todos los itinerarios seleccionados son canónicos y
\begin{equation}
\boxed{K_{\lambda_n}=W_n0,\qquad
\lambda_n\nearrow\min\Lambda_\tau.}
\label{hmtfg:clasificacion:eq:10}
\end{equation}

\textbf{Demostración.} Los lemas de comparación y existencia sitúan cada parámetro seleccionado por debajo de \(a\), con itinerario menor que \(\tau\). El teorema de clasificación identifica ese itinerario con \(W_n0\). La sucesión es creciente y acotada por \(a\); sea \(b\le a\) su límite.

Fijemos \(p\ge1\). Para \(n\) suficientemente grande, \(2p<2^n\). Entonces los signos de \(c_p(\lambda_n)\) y \(c_{2p}(\lambda_n)\) coinciden con \(\tau_p\) y \(\tau_{2p}=-\tau_p\), respectivamente. La cota uniforme de las segundas derivadas,
\[
\|(f_\lambda^p)''\|\le B_p,
\qquad B_p=16^p\frac{16^p-1}{15},
\]
y la identidad \((f_\lambda^p)'(0)=0\) implican
\[
|c_{2p}(\lambda_n)-c_p(\lambda_n)|
\le\tfrac12B_p|c_p(\lambda_n)|^2.
\]
Como los dos términos del lado izquierdo tienen signos contrarios,
\begin{equation}
|c_p(\lambda_n)|>2/B_p.
\label{hmtfg:clasificacion:eq:11}
\end{equation}
Por continuidad, \(c_p(b)\) conserva el signo \(\tau_p\) y satisface \(|c_p(b)|\ge2/B_p>0\). Esto vale para cada \(p\), de modo que \(b\in\Lambda_\tau\). La minimalidad de \(a\) da \(b\ge a\), y por tanto \(b=a\). \(\square\)

La prueba completa conserva el selector global y produce sus itinerarios y su límite. La información de los primeros ocho niveles es compatible con el resultado, pero la demostración de \eqref{hmtfg:clasificacion:eq:10} no extrapola esos ocho niveles: utiliza inducción, intervalos invariantes y separación uniforme de cada signo.


\subsection{Controles lógicos y falsadores}
\label{hmtfg:clasificacion:7}

\par\noindent\textbf{1.}\enspace El teorema de clasificación requiere un máximo unimodal estricto y un intervalo invariante. Si alguno de los dos ramos pierde monotonía, los intervalos \eqref{hmtfg:clasificacion:eq:2} pueden dejar de ser restrictivos y la clasificación debe comprobarse de nuevo.
\par\noindent\textbf{2.}\enspace La condición \(K<\tau\) es esencial. Después de la palabra límite pueden existir centros de período \(2^n\) con otras combinatorias; el teorema no los identifica con \(W_n0\).
\par\noindent\textbf{3.}\enspace La analiticidad en el parámetro proporciona aislamiento y finitud de los ceros sobre el compacto. Para una familia solamente continua, la existencia de una primera raíz estrictamente posterior requiere un argumento adicional.
\par\noindent\textbf{4.}\enspace Un límite de raíces críticas puede, en general, perder signos al alcanzar un retorno de período menor. La desigualdad \eqref{hmtfg:clasificacion:eq:11} excluye precisamente ese fenómeno para esta sucesión y cada horizonte fijo.
\par\noindent\textbf{5.}\enspace La convergencia de las primeras raíces al primer parámetro de itinerario \(\tau\) no equivale a la unicidad de todos los parámetros de ese itinerario ni determina por sí misma una tasa métrica. El cruce transversal y la identificación de la rama aportan esas conclusiones posteriores.

% Fragmento integrable. Preambulo anfitrion: amsmath, amssymb, longtable, hyperref.
\section{Entrada transversal y escalado local}
\label{hmtfg:fragmento:entrada-transversal-y-escalado-local}

\subsection{Procedencia y resultado}
\label{hmtfg:escalado:1}

APP evalúa suma y producto sobre las dos hojas de la retícula, conservando cociente y residuo. TRIT determina régimen y orientación. El transporte TPK selecciona, transporta y actualiza el estado con sus registros de acarreo, frontera y memoria. Sus lectores actúan sobre la misma estructura discreta del continuo. La prolongación \texttt{w6→w12→w18→w24→w30→R36→G9} conserva esta procedencia y distingue retorno de fase de incremento de memoria.

El lector crítico de la rueda dodecafásica, construido en los antecedentes operatorios y transportado por el refinamiento compatible, produce
\[
\phi_j=\frac{(3j-1)\pi_{\mathrm{HMT}}}{18},\qquad
w_j=\frac1{12},\qquad
k_j=\sqrt{\frac2{\sum_{\ell=1}^{12}w_\ell\phi_\ell^2}}\,\phi_j
=\frac{2(3j-1)}{\sqrt{899}}.
\]
La coordenada \(\pi_{\mathrm{HMT}}\) conserva su origen APP–TRIT–TPK; su factor común se cancela en este segundo momento. El lector resultante es
\begin{equation}
u_0(x)=\frac1{12}\sum_{j=1}^{12}(1-\cos k_jx),\qquad
f_\lambda(x)=1-\lambda u_0(x).
\label{hmtfg:escalado:eq:1}
\end{equation}
El sector focal es el uniforme, \(\eta=0\). La ampliación antecedente conserva por separado la familia ponderada \(|\eta|\le1/40\). Las nuevas cotas numéricas corresponden a \eqref{hmtfg:escalado:eq:1}.

La familia es par, entera, con \(f_\lambda(0)=1\), crítico cuadrático y segundo momento \(\sum_jw_jk_j^2=2\). El parámetro \(\lambda\) recorre una familia posterior al generador HMT. La constante objetivo de Feigenbaum permanece fuera de las entradas de \eqref{hmtfg:escalado:eq:1}.

\textbf{Resultado principal.} Existe un único
\begin{equation}
\lambda_*\in[1.874038,1.874039]
\label{hmtfg:escalado:eq:2}
\end{equation}
para el que el cuarto retorno normalizado de \eqref{hmtfg:escalado:eq:1} pertenece a la hoja estable local del punto fijo real de duplicación. El cruce de la familia con esa hoja es transversal, con separación cuantificada. Para todo \(n\ge0\),
\begin{equation}
\boxed{\|R^{4+n}f_{\lambda_*}-g\|_{\mathcal A}
<0.001666\,(0.83996)^n.}
\label{hmtfg:escalado:eq:3}
\end{equation}
La cascada superestable local asociada a este cruce posee parámetros \(\mu_j\) y constantes \(C\ne0\), \(0<\theta<1\), tales que
\begin{equation}
\mu_j-\lambda_*=C\delta_{\mathrm F}^{-j}\bigl[1+O(\theta^j)\bigr],
\qquad
\frac{\mu_{j-1}-\mu_{j-2}}{\mu_j-\mu_{j-1}}\longrightarrow\delta_{\mathrm F}.
\label{hmtfg:escalado:eq:4}
\end{equation}
Aquí \(\delta_{\mathrm F}>1\) es el autovalor inestable del operador de duplicación. La selección de \(\mu_j\) se realiza mediante los retornos locales de sección~\ref{hmtfg:escalado:8}. La identificación con el selector global se demuestra en secciones~\ref{hmtfg:escalado:15}--\ref{hmtfg:escalado:17}: todos los términos globales conservan la combinatoria y los centros coinciden para todo nivel suficientemente alto al igualar sus períodos.


\subsection{Carta analítica y retorno efectivo}
\label{hmtfg:escalado:2}

Se emplea la carta de funciones pares
\[
s=x^2,\qquad z=\frac{s-1}{5/2},\qquad
f(x)=1-x^2V(z),\qquad V(z)=\sum_{j\ge0}v_jz^j.
\]
Las coordenadas del espacio de Banach son
\begin{equation}
u=10v_0,\qquad \nu=(v_1,v_2,\ldots),\qquad
\|(\Delta u,\Delta\nu)\|_{\mathcal A}
=|\Delta u|+\sum_{j\ge1}|\Delta v_j|.
\label{hmtfg:escalado:eq:5}
\end{equation}
En particular, el coeficiente constante de \(V\) tiene peso diez. Esta normalización coincide con la carta \(u/10\) de Hertling–Spandl, sección 2 de la fuente.

Sea \(a=-f(1)=v_0-1\). Para los retornos reales admisibles,
\begin{equation}
(Rf)(x)=-\frac{f(f(-ax))}{a}.
\label{hmtfg:escalado:eq:6}
\end{equation}
El programa evalúa una factorización exacta de \eqref{hmtfg:escalado:eq:6}. Definiendo
\[
S=\frac{a^2s-1}{5/2},\qquad b=V(S),\qquad
h=1-a^2sb,\qquad t=\frac{h^2-1}{5/2},
\]
\[
\operatorname{shift}V(z)=\sum_{j\ge0}v_{j+1}z^j,
\]
resulta
\begin{equation}
\boxed{V_R=a\,b\,(h+1)\left[V(t)+\frac25\operatorname{shift}V(t)\right].}
\label{hmtfg:escalado:eq:7}
\end{equation}
Para comprobarla, se usa \(v_0=1+a\),
\(h^2-1=-a^2sb(h+1)\) y
\(V(t)-v_0=t\operatorname{shift}V(t)\). Así,
\(a+1-h^2V(t)=a^2sb(h+1)[V(t)+(2/5)\operatorname{shift}V(t)]\),
que es precisamente \(asV_R\).

Los momentos iniciales son racionales:
\begin{equation}
\frac1{12}\sum_jk_j^{2m}
=\frac1{12}\left(\frac4{899}\right)^m
\sum_{j=1}^{12}(3j-1)^{2m}.
\label{hmtfg:escalado:eq:8}
\end{equation}
Por ello, tanto \eqref{hmtfg:escalado:eq:7} como la construcción inicial utilizan intervalos racionales; las funciones trigonométricas iniciales se representan por sus series con resto exterior.


\subsection{Cotas externas utilizadas como reconocimiento posterior}
\label{hmtfg:escalado:3}

Una vez construida \eqref{hmtfg:escalado:eq:1}, se utiliza el teorema de hiperbolicidad de Lanford como certificado analítico del operador \eqref{hmtfg:escalado:eq:6}. Sean \(g_0\) el polinomio central publicado y \(g\) el punto fijo. Las estimaciones utilizadas son
\[
\|g-g_0\|_{\mathcal A}<\varepsilon_0=4\,10^{-5},
\]
y, en la bola \(\|f-g_0\|_{\mathcal A}<0.01\), los bloques de \(DR\), escritos \(R=(U,V)\), satisfacen
\begin{equation}
U_u\ge a_0=4.521,\quad
\|U_\nu\|\le b_0=0.560,\quad
\|V_u\|\le c_0=0.756,\quad
\|V_\nu\|\le d_0=0.719.
\label{hmtfg:escalado:eq:9}
\end{equation}
El punto fijo posee una única dirección inestable, con autovalor real positivo \(\delta_{\mathrm F}>1\); el resto del espectro queda dentro del disco unidad. Las cifras de \eqref{hmtfg:escalado:eq:9} proceden de las estimaciones publicadas, no de una nueva reproducción de la prueba de Lanford en este expediente. Los coeficientes de \(g_0\) están incorporados explícitamente en el polinomio central de referencia del programa, conforme a la tabla 2 de Hertling–Spandl.

Fuentes: \href{https://repo-archives.ihes.fr/A_GARDER/20180409/Test/I_Prepublications/LANFORD/1968-2013/P_81_17/P_81_17_web.pdf}{Lanford, estimaciones de hiperbolicidad}; \href{https://arxiv.org/pdf/1410.3277}{Hertling–Spandl, carta analítica y tabla 2}.


\subsection{Certificado racional de entrada y dirección}
\label{hmtfg:escalado:4}

El certificado de entrada a la renormalización divide \eqref{hmtfg:escalado:eq:2} en dieciséis intervalos racionales iguales. Propaga cuatro retornos de \eqref{hmtfg:escalado:eq:7} y la derivada exacta respecto a \(\lambda\), usando diferenciación de la composición y redondeo exterior a cien posiciones decimales.

Cada serie consta de un polinomio de grado 32 con coeficientes intervalares y un error analítico arbitrario \(E\), \(\|E\|_1\le\epsilon\). Este error puede modificar también coeficientes bajos. Para \(\|q\|_1<1\),
\begin{equation}
\|E\circ q\|_1\le\epsilon,\qquad
\|E'\circ q\|_1\le\frac{\epsilon}{(1-\|q\|_1)^2},\qquad
\|\operatorname{shift}E\|_1\le\epsilon.
\label{hmtfg:escalado:eq:10}
\end{equation}
Los productos incluyen los términos polinomio–error y error–error. La cola inicial se acota por su primer término absoluto y una razón geométrica, utilizando \(\|1+(5/2)z\|_1=7/2\). Cada argumento compuesto permanece estrictamente dentro de la bola unidad de esta álgebra.

Para \(\Gamma(\lambda)=R^4f_\lambda=(u_\lambda,\nu_\lambda)\), el certificado da las siguientes cotas conservadoras, uniformes en \eqref{hmtfg:escalado:eq:2}:

\begingroup\small
\begin{longtable}{p{0.465\linewidth}p{0.465\linewidth}}
\hline
Cantidad & Cota certificada exterior \\
\hline
\(\|\Gamma(\lambda)-g_0\|_{\mathcal A}\) & \(<0.004993128\) \\
\(\|\nu_\lambda-\nu_{g_0}\|_1\) & \(<0.001396077\) \\
\(u'_\lambda\) & \(>3821.289115\) \\
\(\|\nu'_\lambda\|_1\) & \(<551.757000\) \\
\(\|\nu'_\lambda\|_1/u'_\lambda\) & \(<0.144386<1/4\) \\
Error analítico de la función & \(<8.482\,10^{-19}\) \\
Error analítico de la tangente & \(<4.569\,10^{-15}\) \\
\hline\end{longtable}
\endgroup

Los decimales completos, los dieciséis subintervalos y cada retorno intermedio se conservan en el certificado racional de entrada.

El cono \(\|\dot\nu\|_1\le|\dot u|/4\) es invariante bajo \eqref{hmtfg:escalado:eq:9}:
\[
|\dot u_{\rm nuevo}|\ge(4.521-0.560/4)|\dot u|
=4.381|\dot u|,
\]
\begin{equation}
\|\dot\nu_{\rm nueva}\|_1\le(0.756+0.719/4)|\dot u|
=0.93575|\dot u|<\tfrac14(4.381)|\dot u|.
\label{hmtfg:escalado:eq:11}
\end{equation}
La dirección paramétrica certificada entra, por tanto, en un cono uniformemente expansivo del operador.


\subsection{Construcción cuantitativa de la hoja estable}
\label{hmtfg:escalado:5}

Traslademos \(g\) al origen y escribamos las coordenadas como \((x,y)=(u-u_g,\nu-\nu_g)\). Fijemos
\begin{equation}
L=0.16,\qquad r=0.004,\qquad
A=a_0-Lc_0=4.40004,\qquad q=d_0+c_0L=0.83996.
\label{hmtfg:escalado:eq:12}
\end{equation}
El cilindro \(|x|\le Lr\), \(\|y\|_1\le r\) está dentro de la bola de \eqref{hmtfg:escalado:eq:9}, puesto que
\((1+L)r+\varepsilon_0=0.00468<0.01\).

Considérese el espacio completo de funciones \(\sigma:B_r\to\mathbb R\) con \(\sigma(0)=0\) y constante de Lipschitz a lo sumo \(L\), dotado de la distancia uniforme. Para cada \(y\), se define \((\mathcal G\sigma)(y)\) como la raíz en \([-L\|y\|_1,L\|y\|_1]\) de
\begin{equation}
F_\sigma(x,y)=U(x,y)-\sigma(V(x,y)).
\label{hmtfg:escalado:eq:13}
\end{equation}
En ese intervalo, \(\|V(x,y)\|_1\le q\|y\|_1<r\). Al incrementar \(x\), \eqref{hmtfg:escalado:eq:9} muestra que \(F_\sigma\) crece al menos con pendiente \(A\). En sus extremos,
\[
F_\sigma(L\|y\|_1,y)
\ge[AL-b_0-Ld_0]\|y\|_1,
\]
y la desigualdad opuesta vale en el extremo negativo. El margen es
\begin{equation}
AL-b_0-Ld_0=0.0289664>0.
\label{hmtfg:escalado:eq:14}
\end{equation}
Se obtiene una raíz única. Comparando \eqref{hmtfg:escalado:eq:13} para dos valores de \(y\),
\[
\operatorname{Lip}(\mathcal G\sigma)
\le\frac{b_0+Ld_0}{A}
=\frac{0.67504}{4.40004}<0.16.
\]
Comparando dos gráficos,
\begin{equation}
\|\mathcal G\sigma-\mathcal G\widetilde\sigma\|_\infty
\le A^{-1}\|\sigma-\widetilde\sigma\|_\infty,
\qquad A^{-1}<0.228.
\label{hmtfg:escalado:eq:15}
\end{equation}
El teorema de contracción produce un único gráfico fijo \(x=\sigma_*(y)\). Sobre él,
\begin{equation}
\|y_n\|_1\le q^n\|y_0\|_1,
\qquad
\|(x_n,y_n)\|_{\mathcal A}\le(1+L)q^n\|y_0\|_1.
\label{hmtfg:escalado:eq:16}
\end{equation}
Además, la distancia vertical \(d(x,y)=x-\sigma_*(y)\) cumple
\[
|d(R(x,y))|\ge A|d(x,y)|
\]
mientras la órbita permanece en el cilindro. Toda órbita que permanece allí indefinidamente pertenece al gráfico. Este identifica exactamente la hoja estable local y proporciona la estimación \eqref{hmtfg:escalado:eq:16}. Su regularidad local coincide con la variedad estable analítica del punto fijo hiperbólico.


\subsection{Existencia, unicidad y transversalidad del parámetro}
\label{hmtfg:escalado:6}

Defínase, para \(\lambda\) en \eqref{hmtfg:escalado:eq:2},
\[
H(\lambda)=u_\lambda-u_g-
\sigma_*(\nu_\lambda-\nu_g).
\]
El certificado asegura \(\|\nu_\lambda-\nu_g\|_1<0.001436077<r\). Para \(\lambda_2>\lambda_1\),
\begin{equation}
H(\lambda_2)-H(\lambda_1)
\ge\int_{\lambda_1}^{\lambda_2}
\bigl[u'_t-L\|\nu'_t\|_1\bigr]dt
>3733.007995\,(\lambda_2-\lambda_1).
\label{hmtfg:escalado:eq:17}
\end{equation}
Para decidir los signos extremos se utiliza \(\|g-g_0\|_{\mathcal A}<\varepsilon_0\) y \(|\sigma_*(y)|\le L\|y\|_1\). La evaluación puntual racional produce
\begin{equation}
H(1.874038)<-0.0011856110945,
\qquad H(1.874039)>0.0021866053399.
\label{hmtfg:escalado:eq:18}
\end{equation}
Por continuidad, existe una raíz; \eqref{hmtfg:escalado:eq:17} prueba su unicidad y la transversalidad del cruce. En esa raíz, \eqref{hmtfg:escalado:eq:16} y
\((1+L)(0.001396077+0.00004)<0.001666\)
demuestran \eqref{hmtfg:escalado:eq:3}.

La cota es uniforme en el número de retornos: convierte el avance de profundidad en un error operatorial geométrico explícito.


\subsection{Admisibilidad real de todos los retornos}
\label{hmtfg:escalado:7}

Los cuatro retornos iniciales se acompañan, en cada subintervalo racional, de las desigualdades
\begin{equation}
0<a=-f(1)<1,\qquad b=f(a)>a,\qquad f(b)<a.
\label{hmtfg:escalado:eq:19}
\end{equation}
El programa guarda \(a,b,f(b)\) antes de cada aplicación de \eqref{hmtfg:escalado:eq:7}. La unimodalidad, la paridad y el crítico cuadrático se conservan bajo estos retornos normalizados.

Para las iteraciones siguientes basta una verificación uniforme en la bola de \eqref{hmtfg:escalado:eq:9}. Si \(E=V-V_{g_0}\), la norma \eqref{hmtfg:escalado:eq:5} da, para \(-0.4\le z\le0\),
\[
|E(z)|\le0.004,\qquad |E'(z)|\le0.01.
\]
La segunda desigualdad usa \(\sup_{j\ge1}j(0.4)^{j-1}=1\). La evaluación racional de los coeficientes centrales proporciona
\begin{equation}
\frac65<V(z)<\frac85,\qquad |V'(z)|<\frac12,
\qquad 0.39<a<0.41.
\label{hmtfg:escalado:eq:20}
\end{equation}
En consecuencia,
\[
f'(x)=-2x\left[V(z)+\frac25x^2V'(z)\right],
\]
y el corchete supera uno para \(0\le x\le1\). El mapa es unimodal y preserva \([-1,1]\). Además,
\[
b=f(a)>1-(0.41)^2\frac85=\frac{4569}{6250}>a,
\]
\begin{equation}
f(b)<1-\left(\frac{4569}{6250}\right)^2\frac65
=\frac{35028967}{97656250}<0.39<a.
\label{hmtfg:escalado:eq:21}
\end{equation}
Las órbitas de la hoja estable satisfacen \eqref{hmtfg:escalado:eq:19} a toda profundidad. Así, las iteraciones analíticas del operador corresponden a retornos reales de duplicación, con sus dominios y orientaciones conservados.


\subsection{Escalado de la cascada superestable local}
\label{hmtfg:escalado:8}

La conclusión métrica requiere dos cruces: el de la familia con la hoja estable, demostrado en sección~\ref{hmtfg:escalado:6}, y el de la variedad inestable con el blanco superestable. El segundo está probado por Eckmann–Wittwer para
\(\Sigma_2=\{\psi:\psi(1)=0\}\), correspondiente al ciclo superestable de período dos. La normalización \eqref{hmtfg:escalado:eq:6} coincide con la del operador real par empleado allí.

El teorema general de Collet–Eckmann–Lanford, sección 6 de la fuente citada, proposición 6.1 y teoremas 6.2–6.3, aplica a un mapa de Banach \(C^2\), un punto fijo con un único autovalor inestable \(\delta_{\mathrm F}>1\), una curva transversal a su hoja estable y un blanco transversal a su variedad inestable. Las preimágenes locales del blanco intersectan la curva en puntos que satisfacen
\begin{equation}
\delta_{\mathrm F}^j(\mu_j-\lambda_*)\longrightarrow C\ne0.
\label{hmtfg:escalado:eq:22}
\end{equation}
El índice \(j\) cuenta los retornos locales. Como la curva inicial aquí es \(R^4f_\lambda\), el período físico es \(2^{j+5}\) cuando el blanco se toma en \(\Sigma_2\); cualquier número fijo de pasos empleado para traer el blanco al entorno local se incorpora en un desplazamiento fijo del índice y en \(C\).

Las condiciones reales de sección~\ref{hmtfg:escalado:7} y la rama real de duplicación del blanco fijan el período exacto y el itinerario. La sucesión se selecciona por esta continuación local orientada. Tomando diferencias en \eqref{hmtfg:escalado:eq:22} se obtiene el límite de razones de \eqref{hmtfg:escalado:eq:4}.

Fuentes de los dos resultados utilizados: \href{https://link.springer.com/article/10.1007/BF01013368}{Eckmann–Wittwer, \emph{A complete proof of the Feigenbaum conjectures}}; \href{https://people.math.harvard.edu/~knill/history/lanford/papers/ColletEckmannLanford.pdf}{Collet–Eckmann–Lanford, sección 6}.

\subsubsection{Resto asintótico de potencia}

La analiticidad de nuestra familia permite precisar \eqref{hmtfg:escalado:eq:22}. En la construcción de coordenadas normales de Collet–Eckmann–Lanford se aplanan únicamente la variedad estable, la variedad inestable y la hipersuperficie blanco. Estas tres subvariedades admiten cambios locales \(C^2\) con derivadas acotadas. El corte suave empleado se realiza en la coordenada escalar inestable; la foliación completa se obtiene después mediante el cambio construido.

En esas coordenadas, la corrección se escribe como una serie de incrementos \(w_j\) cuya norma \(C^1\) satisface
\[
\|w_j\|_{C^1}\le C_1\rho^j,\qquad 0<\rho<1.
\]
La regla de la cadena para el mapa cortado, con derivadas de órdenes uno y dos uniformemente acotadas, proporciona
\[
\|w_j\|_{C^2}\le C_2B^j
\]
para algún \(B\ge1\). La norma ponderada de la construcción controla la norma \(C^1\) ordinaria en el dominio considerado. Por interpolación,
\[
[Dw_j]_\beta\le C_3
\left(\rho^{1-\beta}B^\beta\right)^j.
\]
Se elige
\begin{equation}
0<\beta<\min\left\{1,
\frac{-\log\rho}{\log B-\log\rho}\right\}.
\label{hmtfg:escalado:eq:23}
\end{equation}
La serie converge en \(C^{1,\beta}\). Sea \(z\) la coordenada normal resultante, normalizada para que las preimágenes del blanco satisfagan \(z=\delta_{\mathrm F}^{-j}\), absorbiendo el desplazamiento fijo de índice. Como nuestra curva es analítica y transversal,
\[
z(\Gamma(\lambda))=c(\lambda-\lambda_*)
+O(|\lambda-\lambda_*|^{1+\beta}),\qquad c\ne0.
\]
La inversión local da
\begin{equation}
\mu_j-\lambda_*=C\delta_{\mathrm F}^{-j}
+O\left(\delta_{\mathrm F}^{-(1+\beta)j}\right).
\label{hmtfg:escalado:eq:24}
\end{equation}
Esto demuestra el resto de \eqref{hmtfg:escalado:eq:4}, con \(\theta=\delta_{\mathrm F}^{-\beta}\). El enunciado afirma la existencia de las constantes del resto paramétrico. La tasa operatorial explícita \(0.83996\) de \eqref{hmtfg:escalado:eq:3} controla, por su parte, la convergencia de los mapas renormalizados.

\subsubsection{Conversión del resto en error de las razones}

Si \(\mu_j-\lambda_*=C\delta_{\mathrm F}^{-j}(1+e_j)\), con \(|e_j|\le M\theta^j\), se escribe
\[
\mu_{j-1}-\mu_j=C(\delta_{\mathrm F}-1)\delta_{\mathrm F}^{-j}(1+r_j),
\quad
r_j=\frac{\delta_{\mathrm F} e_{j-1}-e_j}{\delta_{\mathrm F}-1}.
\]
Con \(K=M(\delta_{\mathrm F}+\theta)/(\delta_{\mathrm F}-1)\), se tiene \(|r_j|\le K\theta^{j-1}\). Por tanto,
\begin{equation}
\boxed{
\left|\frac{\mu_{j-2}-\mu_{j-1}}{\mu_{j-1}-\mu_j}-\delta_{\mathrm F}\right|
\le\frac{\delta_{\mathrm F} K(1+\theta)\theta^{j-2}}
{1-K\theta^{j-1}}}
\label{hmtfg:escalado:eq:25}
\end{equation}
cuando \(K\theta^{j-1}<1\). Esta fórmula separa la precisión de evaluación de cada raíz del error de truncamiento de la cascada.


\subsection{Continuación certificada de las ocho raíces iniciales}
\label{hmtfg:escalado:9}

El certificado de continuación finita conserva la familia \eqref{hmtfg:escalado:eq:1} y el certificado anterior de ocho raíces. Para \(2\le n\le8\), estudia
\[
S_n(\lambda)=f_\lambda^{2^n}(0)
\]
desde la caja certificada de \(\lambda_{n-1}\) hasta \(1.874039\). La cobertura racional completa contiene exactamente la raíz heredada \(\lambda_{n-1}\) y la raíz nueva \(\lambda_n\). En particular, \(\lambda_n\) es la única raíz nueva en \((\lambda_{n-1},1.874039]\).

Las cajas se excluyen por signo del retorno o por una derivada local de signo certificado. Las cajas ancla conservan existencia y unicidad por signos extremos y monotonía local. Sus extremos cubren exactamente cada intervalo; se verifican las adyacencias. La evaluación usa los momentos \eqref{hmtfg:escalado:eq:8}, Horner de orden 32, colas geométricas exteriores y \(|u_0''|\le2\). La cola uniforme del potencial es menor que \(1.294\,10^{-58}\) en \(|x|\le3/2\).

El certificado final contiene 650 cajas procesadas y 332 hojas de cobertura. El primer nivel tiene la expresión exacta \(\lambda_1=1/u_0(1)\). Las aproximaciones siguientes se muestran únicamente como orientación; sus cajas completas se conservan en los certificados:

\begingroup\small
\begin{longtable}{p{0.465\linewidth}p{0.465\linewidth}}
\hline
Nivel & Parámetro superestable \\
\hline
1 & 1.345913990471832792210949095… \\
2 & 1.763379787846910127290794030… \\
3 & 1.850413796950136464530566778… \\
4 & 1.868979756606798125150574309… \\
5 & 1.872955034435659740004205128… \\
6 & 1.873806367467030119412262311… \\
7 & 1.873988695221365362307168780… \\
8 & 1.874027744154450672897007573… \\
\hline\end{longtable}
\endgroup

% Fragmento integrable. Preambulo anfitrion: amsmath, amssymb, longtable, hyperref.
\section{Aislamiento del límite y escalado del selector global}
\label{hmtfg:fragmento:aislamiento-del-limite-y-escalado-del-selector-global}

\subsection{Continuación indefinida de la primera raíz y exclusión del intervalo intermedio}
\label{hmtfg:escalado:15}

La clasificación completa demostrada anteriormente conserva la familia y el selector de primera raíz nueva. El transporte ordenado conserva los ceros y su mínimo; la palabra de retornos conserva la lectura cronológica nonádica.

\subsubsection{Clasificación y existencia a toda profundidad}
\label{hmtfg:escalado:15.1}

Escribamos \(W_n=\tau_1\cdots\tau_{2^n-1}\), con \(\tau_j=(-1)^{v_2(j)}\), y \(a_*=\min\Lambda_\tau\). La existencia y compacidad de \(\Lambda_\tau\) están demostradas en el desarrollo de prolongación, \ref{hmtfg:prolongacion:6.1}–\ref{hmtfg:prolongacion:6.2}; su mínimo existe por compacidad y orden.

\textbf{Clasificación.} Todo mapa unimodal estricto de máximo, sobre un intervalo invariante, cuyo crítico tiene período exacto \(2^n\) y cuyo itinerario \(K\) satisface \(K<\tau\), tiene itinerario \(W_n0\).

La inducción se realiza mediante el retorno efectivo. Para período al menos ocho, la comparación con \(\tau\), junto con los intervalos que atraparían órbitas de signos incompatibles, fuerza el prefijo \((+,-,+,+,+,-)\). De aquí se obtienen
\[
J_1=[d_2,d_4],\quad F(J_1)=[d_3,1],\quad
d_4<d_3,\quad F^2(J_1)=J_1.
\]
El retorno normalizado por \(h(x)=d_2x\) tiene máximo uno, período mitad e itinerario \(K'_j=-K_{2j}<\tau\). El orden se conserva porque
\[
O_{2q-1}(K)=-O_{q-1}(K'),\qquad
K_{2q}-\tau_{2q}=-(K'_q-\tau_q).
\]
Los casos iniciales son \(+0\) y \(+-+0\); la inducción reconstruye \(W_n0\). Para los mapas pares, también \(0<d_4<-d_2<1\), de modo que el retorno se extiende a \([-1,1]\) y coincide con \eqref{hmtfg:escalado:eq:6}.

Antes de \(a_*\), todos los itinerarios son menores que \(\tau\): los conjuntos de comparación estricta son abiertos por la separación demostrada en \ref{hmtfg:prolongacion:6.1}; el intervalo anterior a \(a_*\) es conexo y comienza con \(+^\infty<\tau\).

\textbf{Existencia de cada primera raíz.} Supóngase construida \(\lambda_{n-1}<a_*\), y sea \(p=2^{n-1}\). Los ceros de \(c_p\) en \([\lambda_{n-1},a_*]\) forman un conjunto finito no vacío: la función es analítica y \(c_p(a_*)\ne0\). Sea \(z<a_*\) su último cero. En \((z,a_*]\), \(c_p\) tiene signo \(\tau_p\). Para \(g_\lambda=f_\lambda^p\),
\[
g_\lambda'(0)=0,\qquad
c_{2p}(\lambda)=c_p(\lambda)+O(c_p(\lambda)^2).
\]
Cerca de \(z\), por la derecha, \(c_{2p}\) tiene signo \(\tau_p\); en \(a_*\), tiene el contrario. Existe un cero de \(c_{2p}\) en \((z,a_*)\), y su período es exactamente \(2p\), pues allí \(c_p\ne0\). La finitud de los ceros proporciona la primera raíz nueva en \((\lambda_{n-1},a_*)\). El último cero auxiliar prueba existencia; la elección efectiva sigue siendo la primera raíz. La base es \(\lambda_1=1/u_0(1)<a_*\).

\subsubsection{Límite de la sucesión original}
\label{hmtfg:escalado:15.2}

\textbf{Teorema.} El selector original existe a todos los niveles y satisface
\begin{equation}
\boxed{K_{\lambda_n}=W_n0,\qquad
\lambda_n\nearrow a_*=\min\Lambda_\tau.}
\label{hmtfg:escalado:eq:50}
\end{equation}
\textbf{Demostración.} La existencia precedente da una sucesión creciente acotada por \(a_*\). La clasificación fija todos sus itinerarios. Sea \(b\le a_*\) su límite. Para cada \(p\), los términos suficientemente avanzados conservan los signos opuestos \(\tau_p,\tau_{2p}\). Taylor y la cota uniforme \(B_p=16^p(16^p-1)/15\) dan
\[
|c_{2p}(\lambda_n)-c_p(\lambda_n)|
\le\tfrac12 B_p|c_p(\lambda_n)|^2,\qquad
|c_p(\lambda_n)|>2/B_p.
\]
Al pasar al límite se conserva \(\operatorname{sign}c_p(b)=\tau_p\), con módulo al menos \(2/B_p>0\). Para todo \(p\), esto implica \(b\in\Lambda_\tau\); la minimalidad da \(b=a_*\).

El resultado resuelve existencia indefinida, conservación de la cascada simbólica y selección de su primer parámetro límite para la regla global original. La completación HMT transporta esta sucesión y su límite conservando orden y lectura.

\subsubsection{Certificado de exclusión del intervalo intermedio}
\label{hmtfg:escalado:15.3}

El programa de puente certifica
\begin{equation}
\Lambda_\tau\cap[\lambda_{8,\rm inf},1.874038]=\varnothing,
\label{hmtfg:escalado:eq:51}
\end{equation}
donde
\[
\lambda_{8,\rm inf}
=1.874027744154450672897007573595092408435133414523352839687903472138975
\]
es el extremo inferior racional del intervalo certificado de \(\lambda_8\).

La cobertura contiene 374 intervalos exactamente adyacentes y ninguna región pendiente: 319 con \(c_{512}>0\), 10 con \(c_{1024}<0\), dos con \(c_{2048}>0\), todos opuestos al signo de \(\tau\), y 43 próximos a centros de períodos 256, 512 o 1024, excluidos por su segundo retorno.

En este último caso, \(g=f_\lambda^p\), \(d=g(0)\) y una cota uniforme \(|g''|\le M\) entre cero y \(d\) satisfacen \(M|d|<2\). Entonces
\[
|g(d)-d|\le\tfrac12M|d|^2<|d|\quad(d\ne0).
\]
Ambos retornos tienen el mismo signo; para \(d=0\), ambos son cero. Las dos situaciones excluyen \(\tau_{2p}=-\tau_p\) con signos estrictos.

La evaluación utiliza momentos racionales, colas geométricas del potencial y de su segunda derivada, y \(|u_0'''|<6\). El centrado paramétrico intersecta la envolvente natural con
\[
c_j(\lambda_{\rm medio})+
\partial_\lambda c_j(I)\,(I-\lambda_{\rm medio});
\]
la derivada se propaga sobre la caja completa. La autoaplicación de \([-1,1]\) se demuestra antes de restringir a ella las envolventes.

El certificado completo conserva la cobertura y comprueba su teselado exacto, junto con cada desigualdad de signo o de Taylor.

\subsubsection{Localización conjunta}
\label{hmtfg:escalado:15.4}

El cruce local \(\lambda_*\) realiza \(\tau\) por sus retornos reales a toda profundidad, de modo que \(a_*\le\lambda_*\). Por \eqref{hmtfg:escalado:eq:50}, \(\lambda_8<a_*\); junto con \eqref{hmtfg:escalado:eq:51},
\begin{equation}
\boxed{1.874038<a_*\le\lambda_*<1.874039.}
\label{hmtfg:escalado:eq:52}
\end{equation}
La localización conserva el itinerario canónico y sitúa su primer parámetro límite en el mismo intervalo que el cruce estable.

\subsection{Aislamiento por primera salida y coincidencia de los límites}
\label{hmtfg:escalado:16}

Esta sección completa el aislamiento identificado en sección~\ref{hmtfg:escalado:15.4}. Conserva la familia HMT \eqref{hmtfg:escalado:eq:1}, la selección mínima \eqref{hmtfg:escalado:eq:50}, la carta \eqref{hmtfg:escalado:eq:5} y la hoja estable de secciones~\ref{hmtfg:escalado:5}--\ref{hmtfg:escalado:7}. La prueba utiliza el orden \(a_*\le\lambda_*\): basta excluir una separación hacia el lado negativo del cono expansivo.

\subsubsection{Cotas superior e inferior del transporte transversal}
\label{hmtfg:escalado:16.1}

Además de \eqref{hmtfg:escalado:eq:9}, el lema 2.1 de \href{https://arxiv.org/pdf/1410.3277}{Hertling–Spandl, sección 2} proporciona \(\|D\Phi\|<0.9\) en la misma bola de radio \(0.01\), donde
\[
\Phi_u(u,\nu)=u-\frac{U(u,\nu)-u}{3.669}.
\]
Por tanto,
\begin{equation}
\left|1-\frac{U_u-1}{3.669}\right|<0.9,\qquad
U_u<1+1.9(3.669)=7.9711.
\label{hmtfg:escalado:eq:53}
\end{equation}
El número \(3.669\) pertenece al precondicionador publicado de este certificado analítico; la familia generada \eqref{hmtfg:escalado:eq:1} conserva sus coeficientes previos.

Tomemos \(\kappa=0.21\). Para dos puntos de la bola unidos por una secante con
\(\Delta u<0\), \(\|\Delta\nu\|_1\le\kappa|\Delta u|\), la integración de los bloques de \(DR\) sobre el segmento da
\begin{equation}
4.4034|\Delta u|
\le|\Delta u'|\le8.0887|\Delta u|,\qquad \Delta u'<0,
\label{hmtfg:escalado:eq:54}
\end{equation}
\begin{equation}
\|\Delta\nu'\|_1
\le(0.756+0.719\kappa)|\Delta u|
=0.90699|\Delta u|
<\kappa(4.4034)|\Delta u|.
\label{hmtfg:escalado:eq:55}
\end{equation}
Aquí \(4.4034=4.521-0.560\kappa\) y \(8.0887=7.9711+0.560\kappa\). Así, el cono de secantes conserva orientación y se expande mientras el segmento permanezca en la bola.

\subsubsection{Reducción de una separación infinita a una región finita}
\label{hmtfg:escalado:16.2}

Supongamos \(a_*<\lambda_*\), y pongamos
\[
f_n=R^{4+n}f_{a_*},\qquad
s_n=R^{4+n}f_{\lambda_*},\qquad
(\Delta u_n,\Delta\nu_n)=f_n-s_n.
\]
La cota positiva de \(u'\) y \(\|\nu'\|_1/u'<0.144386<\kappa\) sobre \(\Gamma(J)\) sitúan la secante inicial en el cono negativo. La hoja estable está escrita respecto de \(g\), con
\begin{equation}
s_n=g+e_n,\quad
\|e_n\|_{\mathcal A}<0.001666(0.83996)^n,\quad
|(e_n)_u|\le0.16\|(e_n)_\nu\|_1.
\label{hmtfg:escalado:eq:56}
\end{equation}
Los controles de entrada implican
\begin{equation}
|\Delta u_0|
<0.004993128+0.00004
 +0.16(0.001396077+0.00004)
=0.00526290032< h,\qquad h=0.006.
\label{hmtfg:escalado:eq:57}
\end{equation}
Mientras \(|\Delta u_n|<h\), \eqref{hmtfg:escalado:eq:55}–\eqref{hmtfg:escalado:eq:56} dan
\begin{equation}
\|f_n-g_0\|_{\mathcal A}
<(1+\kappa)h+0.001666+0.00004
=0.008966<0.01.
\label{hmtfg:escalado:eq:58}
\end{equation}
El punto \(s_n\) y el segmento entre ambos también permanecen en esa bola convexa. La expansión inferior \eqref{hmtfg:escalado:eq:54} obliga a una primera salida finita \(N\), y la cota superior aplicada al paso precedente proporciona
\begin{equation}
0.006\le-\Delta u_N<0.0485322,\qquad
\|\Delta\nu_N\|_1\le0.21|\Delta u_N|.
\label{hmtfg:escalado:eq:59}
\end{equation}
Se utiliza la desigualdad \(8.0887h=0.0485322\); el salto completo, y no únicamente la superficie de primera salida, queda incluido.

En consecuencia, \(f_N\) pertenece a la siguiente región:
\begin{equation}
\begin{split}
\mathcal C_-=\{\,g+(d_u,d_\nu)+e:\;&
-0.0485322\le d_u\le-0.006,\\
&\|d_\nu\|_1\le0.21|d_u|,\quad
|e_u|+\|e_\nu\|_1\le0.001666,\\
&|e_u|\le0.16\|e_\nu\|_1\,\}.
\end{split}
\label{hmtfg:escalado:eq:60}
\end{equation}
Como \(a_*\) realiza \(\tau\), todos sus retornos normalizados son autoaplicaciones unimodales de \([-1,1]\) con el mismo itinerario. Esta propiedad procede de la inducción de retornos, independientemente de la distancia de \(f_N\) a \(g_0\). Por eso el certificado siguiente trata \(\mathcal C_-\) intersectada con esa clase de autoaplicaciones.

\subsubsection{Exclusión certificada de la región de primera salida}
\label{hmtfg:escalado:16.3}

El certificado racional de exclusión por primera salida, reproducido en el suplemento computacional, cubre exactamente el intervalo de \eqref{hmtfg:escalado:eq:59} mediante 64 bandas racionales adyacentes. Todas quedan excluidas; ninguna subdivisión adicional ni región pendiente intervienen en el resultado.

De las 64 bandas, 54 presentan un signo contrario y diez un retorno contractivo de período 16 o 32. La mayor cota de Lipschitz de esos retornos es menor que \(0.475146491\); el horizonte máximo de exclusión es el iterado 64.

La evaluación conserva como variables compartidas los primeros cuarenta coeficientes de la perturbación, junto con una cota explícita de su cola infinita. Para cada banda, el centro es \(g_0\) desplazado en la coordenada \(u\). Si \(a,b\ge0\) acotan las sensibilidades duales en \(u,\nu\), el soporte de la incertidumbre del punto fijo y del residuo estable es
\begin{equation}
0.00004\max(a,b)
+0.001666\max\left(b,\frac{0.16a+b}{1.16}\right).
\label{hmtfg:escalado:eq:61}
\end{equation}
La segunda expresión se obtiene maximizando \(a|e_u|+b\|e_\nu\|_1\) bajo las dos restricciones de \eqref{hmtfg:escalado:eq:56}. De este modo se utiliza la orientación de la hoja estable demostrada previamente, además de su norma.

Sea \(x_j\) la órbita del centro y \(L_j\) su sensibilidad lineal a todos los coeficientes compartidos. El modelo de Taylor tiene la forma
\[
c_j=x_j+L_j(\Delta v)+r_j,\qquad |r_j|\le E_j.
\]
Si \(\rho_j\) acota \(|L_j(\Delta v)|+E_j\), la recurrencia exterior empleada es
\begin{equation}
E_{j+1}\le |f_0'(x_j)|E_j
+\tfrac12\sup|f_0''|\rho_j^2
+\sup|\Delta f'|\rho_j.
\label{hmtfg:escalado:eq:62}
\end{equation}
Las derivadas se acotan sobre el segmento de los dos estados. La órbita central se verifica dentro de \([-1,1]\); las otras órbitas pertenecen a ese intervalo por la autoaplicación incluida en el dominio certificado. Los cálculos utilizan redondeo racional exterior, incluyendo las colas geométricas.

Cada banda se excluye por una de dos desigualdades:

\par\noindent\textbf{1.}\enspace Un valor crítico tiene signo estrictamente opuesto a \(\tau_j\).
\par\noindent\textbf{2.}\enspace Para \(p=2^k\), el retorno \(H=f^p\) tiene constante de Lipschitz \(L_p<1\) sobre el segmento entre \(0\) y \(H(0)\). Entonces
\[
|H(H(0))-H(0)|\le L_p|H(0)|.
\]
Si \(H(0)\ne0\), los dos valores tienen el mismo signo; si \(H(0)=0\), ambos se anulan. En ambos casos queda excluido el itinerario estricto \(\tau_{2p}=-\tau_p\).

La constante \(L_p\) se obtiene propagando el segmento inicial \([-r,r]\), \(r\ge|H(0)|\), junto a las cajas de la órbita crítica y multiplicando las cotas exteriores de las derivadas. Se acotan los mapas de la banda completa, manteniendo \eqref{hmtfg:escalado:eq:61}.

\subsubsection{Teorema de coincidencia del parámetro de acumulación}
\label{hmtfg:escalado:16.4}

\textbf{Teorema.} Para la familia uniforme HMT \eqref{hmtfg:escalado:eq:1}, el límite de la selección original de primeras raíces coincide con el cruce estable certificado:
\begin{equation}
\boxed{\lambda_n\nearrow a_*=\lambda_*,
\qquad 1.874038<\lambda_*<1.874039.}
\label{hmtfg:escalado:eq:63}
\end{equation}

\textbf{Demostración.} \eqref{hmtfg:escalado:eq:50}–\eqref{hmtfg:escalado:eq:52} dan \(\lambda_n\uparrow a_*\le\lambda_*\), con ambos parámetros en \(J\). Si la desigualdad fuera estricta, \eqref{hmtfg:escalado:eq:54}–\eqref{hmtfg:escalado:eq:59} producirían un retorno \(f_N\in\mathcal C_-\) que realiza \(\tau\). La cobertura completa de sección~\ref{hmtfg:escalado:16.3} excluye precisamente esa posibilidad. Luego \(a_*=\lambda_*\).

La prueba controla todas las profundidades mediante una primera salida finita y una cobertura uniforme de su región. La identidad \eqref{hmtfg:escalado:eq:63} enlaza el límite del selector global con el cruce transversal.


\subsection{Identificación de los centros originales y escalado global}
\label{hmtfg:escalado:17}

El último paso conserva la selección de \eqref{hmtfg:escalado:eq:50}. La continuación local y los parámetros originales se comparan por su período exacto y su clase híbrida; la coincidencia se obtiene mediante una unicidad en un entorno paramétrico fijo.

\subsubsection{Realización analítica de grado dos y transversalidad}
\label{hmtfg:escalado:17.1}

El punto fijo \(g\) de sección~\ref{hmtfg:escalado:3} es par, analítico, de crítico cuadrático y combinatoria estacionaria de duplicación. El teorema 2.1 de \href{https://annals.math.princeton.edu/wp-content/uploads/annals-v164-n3-p01.pdf}{de Faria–de Melo–Pinto}, aplicado a esa única combinatoria, da un conjunto límite de un solo elemento, con extensión holomorfa propia de grado dos entre discos topológicos anidados. Su atracción sobre el intervalo identifica ese elemento con \(g\), pues \(Rg=g\).

El lema 3.2 de la misma fuente extiende un número fijo de retornos a una vecindad analítica completa de \(g\), con imágenes de grado dos y módulo positivo. La norma \eqref{hmtfg:escalado:eq:5} controla la norma uniforme sobre una vecindad compleja suficientemente pequeña de \([-1,1]\), porque \(|(x^2-1)/2.5|<1\) allí. La convergencia \eqref{hmtfg:escalado:eq:3} entra en dicha vecindad; la dependencia analítica de los retornos proporciona, para cierto entero fijo \(d\), una familia
\begin{equation}
\mathcal F_\lambda=R^d f_\lambda
\label{hmtfg:escalado:eq:64}
\end{equation}
de mapas de grado dos, analítica en un entorno complejo de \(\lambda_*\). Los discos y el número de retornos se fijan antes de variar \(\lambda\) en ese entorno.

La tangente paramétrica pertenece al cono expansivo, mientras la hoja estable satisface la pendiente de sección~\ref{hmtfg:escalado:5}. Los retornos preservan esa separación. La realización de grado dos identifica localmente la hoja estable con la clase híbrida del mapa infinitamente renormalizable; por ello \eqref{hmtfg:escalado:eq:64} es transversal a esa clase.

El cambio de espacio analítico se justifica detalladamente en \ref{hmtfg:clasificacion:8.2}, dedicada al transporte de la transversalidad. Una dirección tangente a la clase híbrida produciría tangentes renormalizadas decrecientes, por la convergencia uniforme sobre un disco híbrido y la estimación de Cauchy. La dirección certificada satisface, en cambio, \(|\dot f(1)|=|\dot u|/10\) y crece bajo los retornos por el cono expansivo. Esta evaluación común de ambos representantes prueba la transversalidad sin suponer equivalencia de sus normas.

\subsubsection{Unicidad en una vecindad fija}
\label{hmtfg:escalado:17.2}

Para \(n>d\), la clasificación de sección~\ref{hmtfg:escalado:15} da
\begin{equation}
K(\mathcal F_{\lambda_n})=W_{n-d}0.
\label{hmtfg:escalado:eq:65}
\end{equation}
El enderezamiento de un mapa de grado dos conserva la órbita crítica; \eqref{hmtfg:escalado:eq:65} identifica su clase híbrida con la del centro cuadrático canónico de período \(2^{n-d}\). Las cotas a priori de la sucesión real de duplicación son las empleadas en el teorema de universalidad de Lyubich.

El \href{https://www.maths.tcd.ie/EMIS/journals/Annals/149_2/lyubich.pdf\#page=69}{teorema 7.4 de Lyubich, pp. 387–388} da una sola intersección con cada una de esas clases, a niveles suficientemente altos, sobre la transversal analítica y dentro de una vecindad fija del parámetro de acumulación.

Por \eqref{hmtfg:escalado:eq:63}, \(\lambda_n\) entra finalmente en esa vecindad. La continuación local de sección~\ref{hmtfg:escalado:8}, indexada por el mismo período, pertenece a la misma clase híbrida. La unicidad implica
\begin{equation}
\boxed{\lambda_n=\widehat\mu_n\quad(n\ge n_0),}
\label{hmtfg:escalado:eq:66}
\end{equation}
donde \(\widehat\mu_n\) designa el centro local de período físico \(2^n\). Es una reindexación de \(\mu_j\) mediante un desplazamiento fijo determinado por los retornos y el blanco de sección~\ref{hmtfg:escalado:8}. La selección original de \(\lambda_n\) permanece intacta.

El entero \(n_0\) es existencial en esta prueba. La continuación de todos los niveles y la palabra \(W_n0\) ya están demostradas en \eqref{hmtfg:escalado:eq:50}; \eqref{hmtfg:escalado:eq:66} aporta la identidad métrica necesaria para el límite.

\subsubsection{Teorema de escalado para las primeras raíces HMT}
\label{hmtfg:escalado:17.3}

\textbf{Teorema.} Para la familia uniforme \eqref{hmtfg:escalado:eq:1}, sean \(\lambda_n\) las sucesivas primeras raíces nuevas de período crítico exacto \(2^n\). Existen \(C>0\), \(0<\theta<1\) y \(M<\infty\) tales que, para todo \(n\) suficientemente grande,
\begin{equation}
\lambda_*-\lambda_n=C\delta_{\mathrm F}^{-n}(1+e_n),
\qquad |e_n|\le M\theta^n,
\label{hmtfg:escalado:eq:67}
\end{equation}
y
\begin{equation}
\boxed{\lim_{n\to\infty}
\frac{\lambda_{n-1}-\lambda_{n-2}}
{\lambda_n-\lambda_{n-1}}=\delta_{\mathrm F}.}
\label{hmtfg:escalado:eq:68}
\end{equation}

\textbf{Demostración.} \eqref{hmtfg:escalado:eq:66} transporta el escalado local \eqref{hmtfg:escalado:eq:24} a las raíces originales; el desplazamiento finito de índices se absorbe en \(C,M\). Su signo es positivo porque \(\lambda_n<\lambda_*\).

El resto de potencia admite además una justificación directa mediante la holonomía de clases híbridas. El lema 7.3 de Lyubich da regularidad \(C^{1+\beta}\), con derivada no nula, sobre la transversal en el punto de acumulación. La dinámica unidimensional inestable es analítica y tiene multiplicador \(\delta_{\mathrm F}>1\); su coordenada linealizante hace que los centros sucesivos estén a distancia proporcional a \(\delta_{\mathrm F}^{-n}\). Componer con la inversa de la holonomía da un término lineal no nulo y un resto \(O(\delta_{\mathrm F}^{-(1+\beta)n})\). Así se puede tomar \(\theta=\delta_{\mathrm F}^{-\beta}\).

Escribiendo
\[
r_n=\frac{\delta_{\mathrm F} e_{n-1}-e_n}{\delta_{\mathrm F}-1},
\qquad K=\frac{M(\delta_{\mathrm F}+\theta)}{\delta_{\mathrm F}-1},
\]
se obtiene
\[
\lambda_n-\lambda_{n-1}
=C(\delta_{\mathrm F}-1)\delta_{\mathrm F}^{-n}(1+r_n),\qquad
|r_n|\le K\theta^{n-1}.
\]
Por tanto,
\begin{equation}
\left|
\frac{\lambda_{n-1}-\lambda_{n-2}}
{\lambda_n-\lambda_{n-1}}-\delta_{\mathrm F}
\right|
\le
\frac{\delta_{\mathrm F} K(1+\theta)\theta^{n-2}}
{1-K\theta^{n-1}},
\quad K\theta^{n-1}<1.
\label{hmtfg:escalado:eq:69}
\end{equation}
El lado derecho tiende a cero, demostrando \eqref{hmtfg:escalado:eq:68}.

\subsubsection{Encadenamiento y alcance del cierre}
\label{hmtfg:escalado:17.4}

La composición efectiva es: lector dodecafásico APP–TRIT–TPK; perfil \eqref{hmtfg:escalado:eq:1}; transporte ordenado de retornos \eqref{hmtfg:escalado:eq:38}–\eqref{hmtfg:escalado:eq:40}; clasificación y existencia de las primeras raíces \eqref{hmtfg:escalado:eq:50}; localización \eqref{hmtfg:escalado:eq:52}; aislamiento por primera salida \eqref{hmtfg:escalado:eq:63}; identidad eventual de centros \eqref{hmtfg:escalado:eq:66}; escalado \eqref{hmtfg:escalado:eq:67}–\eqref{hmtfg:escalado:eq:69}. La representación de Jacobi \eqref{hmtfg:escalado:eq:43} conserva el perfil completo de este lector y sus marcas permanecen en el antecedente.

La expansión espacial del lector HMT y el refinamiento nonádico aportan la genealogía y la evaluación a profundidad arbitraria. El escalado entre parámetros se demuestra mediante la dinámica del retorno, la transversalidad y el aislamiento que acabamos de componer. Las pruebas analíticas citadas intervienen sobre la familia ya construida; sus valores de comparación no seleccionan sus fases ni sus coeficientes.

El cierre se refiere al selector global original del sector uniforme \(\eta=0\), con las dependencias teoremáticas enumeradas. La ley asintótica conserva sus cuantificadores existenciales. La anchura de una caja de raíces finitas mide precisión de evaluación; la estimación \eqref{hmtfg:escalado:eq:69} expresa la convergencia hacia el límite.


Los programas y certificados del suplemento de reproducción permiten reproducir las partes finitas de la prueba. La continuación a toda profundidad se sustenta en las inducciones, el argumento uniforme de primera salida y los teoremas de identificación, con sus dominios explícitos.

% Fragmento integrable. Preambulo anfitrion: amsmath, amssymb, longtable, hyperref.
\section{Identificación analítica de los centros originales}
\label{hmtfg:fragmento:identificacion-analitica-de-los-centros-originales}

\subsection{Realización de grado dos e identificación métrica de los centros}
\label{hmtfg:clasificacion:8}

La igualdad entre el primer parámetro del itinerario infinito y el cruce estable, demostrada por primera salida, proporciona
\begin{equation}
\boxed{a_*:=\min\Lambda_\tau=\lambda_*.}
\label{hmtfg:clasificacion:eq:H}
\end{equation}

La familia, el perfil dodecafásico, los retornos y el selector se mantienen exactamente como en secciones~\ref{hmtfg:clasificacion:1}--\ref{hmtfg:clasificacion:5}. Los resultados externos empleados a continuación reconocen la dinámica de ese objeto ya construido. Sus valores universales no intervienen en la selección de coeficientes ni de parámetros de la familia HMT.

\textbf{Teorema de transferencia métrica.} La identificación \eqref{hmtfg:clasificacion:eq:H}, junto con las cotas de entrada, admisibilidad y transversalidad anteriormente demostradas, implica que la sucesión original de primeras raíces nuevas satisface, para ciertas constantes \(C>0\) y \(0<\theta<1\),
\begin{equation}
\boxed{\lambda_*-\lambda_n=C\delta_{\mathrm F}^{-n}
\bigl(1+O(\theta^n)\bigr),}
\label{hmtfg:clasificacion:eq:14}
\end{equation}
donde \(\delta_{\mathrm F}>1\) es el autovalor inestable del operador de duplicación en el punto fijo reconocido. En particular,
\begin{equation}
\frac{\lambda_{n-1}-\lambda_{n-2}}
     {\lambda_n-\lambda_{n-1}}
=\delta_{\mathrm F}+O(\theta^n).
\label{hmtfg:clasificacion:eq:15}
\end{equation}
La sucesión coincide eventualmente con los centros de la continuación local, una vez que ambos índices se refieren al mismo período físico. La prueba mantiene el selector de primera raíz de sección~\ref{hmtfg:clasificacion:4}.

\subsubsection{Obtención de un dominio cuadrático semejante común}
\label{hmtfg:clasificacion:8.1}

Escribamos \(g\) para el punto fijo de la renormalización y
\(\Gamma(\lambda)=R^4f_\lambda\). La carta analítica utilizada en el certificado es
\[
f(x)=1-x^2V\!\left(\frac{x^2-1}{2.5}\right),\qquad
V(z)=\frac{u}{10}+\sum_{k\ge1}\nu_kz^k,
\]
con norma de diferencias \(|\Delta u|+\|\Delta\nu\|_1\). Las desigualdades de admisibilidad prueban que \(g\) es par, analítico, unimodal, cuadrático y de tipo de duplicación; además, \(g(x)=\varphi(x^2)\) con \(\varphi'<0\) en \([0,1]\).

El teorema 2.1 de de Faria–de Melo–Pinto, aplicado al conjunto de combinatorias formado únicamente por la duplicación, proporciona un conjunto límite de un solo punto con extensión cuadrática semejante. Como \(Rg=g\), la convergencia real de ese teorema identifica dicho punto con \(g\). El lema 3.2 proporciona, en un entorno analítico de \(g\), una iteración finita de renormalización con dominio cuadrático semejante y módulo uniformemente positivo. Los localizadores son secciones 2.1–2.3, pp. 736–738, y sección 3.1, lema 3.2, p. 745, de \href{https://annals.math.princeton.edu/wp-content/uploads/annals-v164-n3-p01.pdf}{de Faria–de Melo–Pinto, \emph{Global hyperbolicity of renormalization for \(C^r\) unimodal mappings}}.

Para comprobar la pertenencia al entorno de ese lema, sea \(\Omega_\epsilon\) un entorno complejo suficientemente pequeño de \([-1,1]\), con
\[
\rho_\epsilon:=\sup_{x\in\Omega_\epsilon}
\left|\frac{x^2-1}{2.5}\right|<1,
\qquad M_\epsilon:=\sup_{x\in\Omega_\epsilon}|x|^2<\infty.
\]
La restricción de la carta certificada satisface
\begin{equation}
\|\Delta f\|_{\Omega_\epsilon}
\le M_\epsilon\left(\frac{|\Delta u|}{10}
          +\rho_\epsilon\|\Delta\nu\|_1\right)
\le C_\epsilon\|\Delta f\|_{\mathcal A}.
\label{hmtfg:clasificacion:eq:16}
\end{equation}
El certificado da
\[
\|R^m\Gamma(\lambda_*)-g\|_{\mathcal A}
<0.001666(0.83996)^m.
\]
Se elige por ello un entero fijo \(m\) que sitúe esta imagen en el entorno del lema 3.2. La continuidad de la composición finita proporciona un intervalo \(I_*\) alrededor de \(\lambda_*\) donde las mismas operaciones están definidas. Si \(N\) es la iteración suministrada por dicho lema, ponemos
\begin{equation}
d=4+m+N,\qquad
\mathcal F_\lambda=R^df_\lambda,\quad \lambda\in I_*.
\label{hmtfg:clasificacion:eq:17}
\end{equation}
La familia \(\mathcal F\) posee representantes cuadráticos semejantes de grado dos, sobre un dominio común y con módulo positivo uniforme. Su dependencia analítica real admite una complejificación local del parámetro. El número \(d\) es fijo y cuenta las duplicaciones realizadas antes de aplicar el resultado paramétrico.

Esta construcción fija un dominio común y un número finito de retornos \(d\), suficientes para aplicar el teorema paramétrico.

\subsubsection{Transporte de la transversalidad certificada}
\label{hmtfg:clasificacion:8.2}

La dirección de \(\Gamma\) satisface el cono
\[
\|\dot\nu\|_1\le\tfrac14|\dot u|,\qquad \dot u\ne0.
\]
Las cotas diferenciales certificadas implican invariancia de este cono y
\begin{equation}
|\dot u_{k+1}|\ge4.381|\dot u_k|
\label{hmtfg:clasificacion:eq:18}
\end{equation}
a lo largo de la órbita de \(\Gamma(\lambda_*)\). Todos sus puntos permanecen en el dominio certificado. En particular, la dirección transportada por los pasos finitos de \eqref{hmtfg:clasificacion:eq:17} conserva una componente expansiva no nula.

La clase híbrida del punto fijo coincide con su variedad estable en el espacio de gérmenes cuadráticos semejantes: \href{https://www.maths.tcd.ie/EMIS/journals/Annals/149_2/lyubich.pdf}{Lyubich, \emph{Feigenbaum–Coullet–Tresser universality and Milnor's hairiness conjecture}, teorema 6.1, pp. 371–372}. La transversalidad sigue siendo válida al pasar a ese espacio, como puede comprobarse sin identificar normas de Banach diferentes. En efecto, si \(\partial_\lambda\mathcal F_{\lambda_*}\) fuese tangente a la clase híbrida, sería la tangente de un disco analítico contenido en ella. La convergencia exponencial uniforme de la renormalización sobre un subdisco compacto, seguida de la estimación de Cauchy en su parámetro, haría converger a cero la evaluación en \(x=1\) de sus tangentes renormalizadas. Sin embargo, en la carta anterior,
\[
\dot f(1)=-\dot u/10,
\]
y \eqref{hmtfg:clasificacion:eq:18} hace crecer su valor absoluto. Ambas derivadas representan el mismo germen y tienen la misma evaluación real. La contradicción demuestra
\begin{equation}
\partial_\lambda\mathcal F_{\lambda_*}
\notin T_{\mathcal F_{\lambda_*}}\mathcal H_{c_\infty}.
\label{hmtfg:clasificacion:eq:19}
\end{equation}
Así, \(S=\{\mathcal F_\lambda\}\) es una transversal analítica compleja a la clase híbrida del límite de duplicación.

\subsubsection{Identificación eventual de las primeras raíces por su período}
\label{hmtfg:clasificacion:8.3}

Denotemos por \(c_r\) el centro cuadrático canónico de período \(2^r\), obtenido mediante \(r\) duplicaciones sucesivas del centro de período uno. Este símbolo designa el parámetro de la familia cuadrática de reconocimiento; se distingue de las funciones orbitales \(c_j(\lambda)=f_\lambda^j(0)\) de las secciones anteriores. Se escribe \(c_r^{\rm quad}\) para conservar la distinción.

El teorema de continuación global y \eqref{hmtfg:clasificacion:eq:H} dan \(\lambda_n\to\lambda_*\), por lo que \(\lambda_n\in I_*\) para \(n\) suficientemente grande. La clasificación e intervalos restrictivos de sección~\ref{hmtfg:clasificacion:2} muestran que, cuando \(n>d\),
\begin{equation}
\mathcal F_{\lambda_n}=R^df_{\lambda_n}
\quad\text{tiene período crítico exacto }2^{n-d}
\quad\text{e itinerario }W_{n-d}0.
\label{hmtfg:clasificacion:eq:20}
\end{equation}
Su órbita crítica permanece en el intervalo real invariante, contenido en el dominio cuadrático semejante; su conjunto de Julia lleno es, por tanto, conexo. La rectificación conserva la órbita postcrítica y las ramas reales. Al recorrer sus retornos restrictivos se obtienen \(n-d\) duplicaciones canónicas y, al final, un punto crítico fijo. Este último rectifica al centro \(0\); las inversas sucesivas de la rectificación de las copias de duplicación dan el centro único \(c_{n-d}^{\rm quad}\). Esta es la aplicación concreta de la parametrización por combinatoria y del homeomorfismo de rectificación de las copias, descritos en Lyubich, secciones 5.1–5.3, pp. 362–364. Se concluye
\begin{equation}
\mathcal F_{\lambda_n}\in\mathcal H_{c_{n-d}^{\rm quad}}.
\label{hmtfg:clasificacion:eq:21}
\end{equation}

El teorema 7.4 de Lyubich, pp. 387–388, da una única intersección local de \(S\) con cada clase \(\mathcal H_{c_r^{\rm quad}}\) para \(r\) suficientemente grande. Sus hipótesis de cotas complejas a priori se cumplen para la cascada cuadrática real de duplicación, conforme al teorema 5.6, p. 367, del mismo trabajo. Escribamos \(\zeta_r\) para el parámetro de esa intersección. Las ecuaciones \eqref{hmtfg:clasificacion:eq:20}–\eqref{hmtfg:clasificacion:eq:21}, la convergencia al cruce y la unicidad local implican
\begin{equation}
\boxed{\lambda_n=\zeta_{n-d}\quad\text{para todo }n\text{ suficientemente grande}.}
\label{hmtfg:clasificacion:eq:22}
\end{equation}
La unicidad utilizada aquí corresponde a \textbf{cada clase híbrida canónica} en un entorno fijo del cruce. No se infiere de la mera existencia de una cascada local ni de la unicidad del cruce con la hoja estable. La ecuación \eqref{hmtfg:clasificacion:eq:22} identifica el selector original, ya definido globalmente, con las intersecciones locales; no introduce un selector sustitutorio.

En particular, sea \(\mu_j\) la cascada local de la sección \ref{hmtfg:escalado:8}, y sea \(s\) el entero fijo determinado por su período físico exacto:
\[
\operatorname{per}_{\rm crit}(f_{\mu_j})=2^{j+s}.
\]
En la presentación con curva \(R^4f_\lambda\) y blanco de período dos, \(s=5\); los pasos fijos adicionales usados para localizar el blanco modifican \(s\). Por \eqref{hmtfg:clasificacion:eq:21} y la misma unicidad,
\begin{equation}
\mu_j=\zeta_{j+s-d},\qquad
\boxed{\lambda_n=\mu_{n-s}\quad(n\text{ suficientemente grande}).}
\label{hmtfg:clasificacion:eq:23}
\end{equation}
El desplazamiento \(d\) de la preparación cuadrática semejante y el desplazamiento \(s\) de los períodos del blanco son objetos distintos. La identificación se efectúa por el período físico, no igualando esos dos enteros.

\subsubsection{Resto de potencia mediante holonomía transversal}
\label{hmtfg:clasificacion:8.4}

El lema 7.3 de Lyubich, p. 384 y demostración en pp. 384–387, aporta regularidad \(C^{1+\beta}\)-conforme, para algún \(\beta>0\), de la holonomía entre transversales a la clase híbrida del punto de Feigenbaum. Su definición en p. 384 se refiere a los conjuntos de conexidad sobre las transversales; contiene, en particular, los centros considerados aquí. Podemos reducir el exponente para que \(0<\beta\le1\).

Sea \(h\) la holonomía desde \(S\) hacia la variedad inestable y sea \(z\) una coordenada analítica que linealiza la restricción unidimensional de la renormalización:
\begin{equation}
z(g)=0,\qquad z(Rq)=\delta_{\mathrm F} z(q).
\label{hmtfg:clasificacion:eq:24}
\end{equation}
La existencia de esta coordenada se obtiene aplicando la linealización de un germen holomorfo repulsor, de multiplicador \(\delta_{\mathrm F}>1\), a la variedad inestable analítica. La dirección expansiva de \eqref{hmtfg:clasificacion:eq:18}, la invariancia del cono y el retorno estacionario identifican ese multiplicador con el autovalor inestable positivo empleado en el certificado.

Si \(q_r\) es el centro de la clase \(\mathcal H_{c_r^{\rm quad}}\) situado en la variedad inestable, la compatibilidad de rectificación y renormalización da \(Rq_r=q_{r-1}\). Esta identidad es la que se utiliza en la demostración del teorema 7.4, p. 388. Por \eqref{hmtfg:clasificacion:eq:24}, existe \(b\ne0\), que absorbe cualquier índice inicial fijo, tal que
\begin{equation}
z(q_r)=b\delta_{\mathrm F}^{-r}
\label{hmtfg:clasificacion:eq:25}
\end{equation}
para todo \(r\) suficientemente grande.

La regularidad de \(h\) y la transversalidad \eqref{hmtfg:clasificacion:eq:19} dan, sobre el conjunto de conexidad cercano a \(\lambda_*\),
\begin{equation}
H(\lambda):=z\bigl(h(\mathcal F_\lambda)\bigr)
=A(\lambda-\lambda_*)
+O\!\left(|\lambda-\lambda_*|^{1+\beta}\right),
\qquad A\ne0.
\label{hmtfg:clasificacion:eq:26}
\end{equation}
Como \(H(\zeta_r)=b\delta_{\mathrm F}^{-r}\), la estimación inferior
\(|H(\lambda)|\ge |A|\,|\lambda-\lambda_*|/2\) en un entorno suficientemente pequeño implica primero
\(|\zeta_r-\lambda_*|=O(\delta_{\mathrm F}^{-r})\). Sustituyendo esta cota en el resto de \eqref{hmtfg:clasificacion:eq:26}, resulta
\begin{equation}
\zeta_r-\lambda_*=\frac bA\delta_{\mathrm F}^{-r}
+O\!\left(\delta_{\mathrm F}^{-(1+\beta)r}\right).
\label{hmtfg:clasificacion:eq:27}
\end{equation}
Aplicamos \eqref{hmtfg:clasificacion:eq:22} con \(r=n-d\). El desplazamiento finito modifica el coeficiente principal y la constante del resto. Puesto que \(\lambda_n<\lambda_*\), se obtiene \eqref{hmtfg:clasificacion:eq:14}, con
\[
C=-\frac bA\delta_{\mathrm F}^d>0,
\qquad \theta=\delta_{\mathrm F}^{-\beta}\in(0,1).
\]
Finalmente,
\[
\lambda_n-\lambda_{n-1}
=C(\delta_{\mathrm F}-1)\delta_{\mathrm F}^{-n}\bigl(1+O(\theta^n)\bigr),
\]
lo que demuestra \eqref{hmtfg:clasificacion:eq:15}. \(\square\)


\subsubsection{Alcance cuantitativo del teorema}
\label{hmtfg:clasificacion:8.5}

La identificación de límites demostrada por primera salida activa la transferencia métrica para las primeras raíces globales del sector uniforme. La tasa operatorial \(0.83996\) controla los mapas renormalizados sobre la hoja estable. La tasa paramétrica \(\theta=\delta_{\mathrm F}^{-\beta}\) procede de la regularidad de la holonomía transversal. Ambas estimaciones conservan su variable y su dominio.

El teorema acredita la existencia de \(\beta,\theta,C\), de la vecindad \(I_*\), del número finito de retornos \(d\) y de un índice \(n_0\) a partir del cual coinciden los centros. Sus valores numéricos individuales no se asignan en este enunciado asintótico. Esta formulación completa el límite y su resto de potencia; la evaluación numérica de esos testigos constituye una especialización cuantitativa distinta.

\section{Retorno, memoria y universalidad de la cascada dodecafásica}
\label{sec:hmt-feig-interpretacion}

\subsection{La estructura generadora y la lectura crítica}

La construcción de Feigenbaum en Holografía Modular Triádica enlaza una
organización aritmética orientada con una ley asintótica de cambio de escala.
Su punto de partida es la composición efectiva de APP, TRIT y TPK. APP evalúa
suma y producto sobre las dos hojas de la retícula de dígitos y conserva los
cocientes y residuos de esas operaciones. TRIT determina el régimen local y
la orientación de su lectura. TPK selecciona y transporta los estados,
actualiza sus registros y prolonga sus rutas con acarreo, frontera y memoria.
La estructura discreta conjunta del continuo proporciona la residencia común
de esos estados y de los lectores que actúan sobre ellos.

La prolongación
\[
w_6\longrightarrow w_{12}\longrightarrow w_{18}
\longrightarrow w_{24}\longrightarrow w_{30}
\longrightarrow R_{36}\longrightarrow G_9
\]
organiza esta persistencia. La semilla inicial se eleva conservando sus
prefijos; las elevaciones sucesivas cambian el régimen sin perder orientación
ni memoria; la frontera coinductiva permite continuar la construcción, y las
nueve cartas reúnen sus lecturas locales. La holonomía nonádica expresa el
retorno de fase acompañado de un incremento de memoria. Así, una vuelta
visible conserva la posición que ocupa dentro de la historia completa. La
operación de retorno adquiere desde el origen dos componentes inseparables:
reiteración de una regla y conservación de su procedencia.

El lector crítico fijado sobre la rueda dodecafásica orientada convierte esta
estructura en doce fases con una traza uniforme. Sus representantes son
\[
\phi_m=\frac{(3m-1)\pi_{\mathrm{HMT}}}{18},
\qquad w_m=\frac1{12},\qquad 1\le m\le12.
\]
La elección del representante orientado importa porque el segundo momento
utiliza su magnitud. El cociente angular registra la fase, mientras el
representante conserva la información necesaria para normalizarla. La
coordenada \(\pi_{\mathrm{HMT}}\) procede de los lectores coinductivos del mismo
núcleo APP--TRIT--TPK. Dentro de su imagen correlacionada, las coordenadas
\(\pi,\varphi,e\) y el cierre dodecafásico de \(\alpha_{\mathrm{HMT}}\)
conservan su orden constructivo y sus dominios propios. En el lector crítico
aquí considerado se utiliza la coordenada angular ya generada.

La normalización se determina mediante una suma entera:
\[
\sum_{m=1}^{12}(3m-1)^2=5394=6\cdot899,
\qquad k_m=\frac{2(3m-1)}{\sqrt{899}}.
\]
El factor común \(\pi_{\mathrm{HMT}}\) se cancela al fijar el segundo momento.
Esta cancelación conserva la procedencia de las fases y produce una
representación particularmente precisa de su lectura normalizada. Se obtiene
\[
u_0(x)=\frac1{12}\sum_{m=1}^{12}\bigl(1-\cos(k_mx)\bigr),
\qquad f_\lambda(x)=1-\lambda u_0(x),
\qquad \sum_mw_mk_m^2=2.
\]
El perfil es par y entero; satisface \(u_0(0)=u_0'(0)=0\) y
\(u_0''(0)=2\). Para \(0<\lambda\le2/u_0(1)\), la familia lleva
\([-1,1]\) en sí mismo y tiene un máximo crítico cuadrático con
\(f_\lambda''(0)=-2\lambda\). El parámetro \(\lambda\) recorre esta familia
después de fijados el lector, la orientación, los pesos y la escala. La cifra
universal que se reconocerá al final queda así separada de la selección de
los coeficientes.

\subsection{Doce fases, momentos y representación espectral}

El perfil admite una segunda expresión que reúne sus fases en un operador
finito. Los nodos
\[
s_m=k_m^2=\frac{4(3m-1)^2}{899}
\]
determinan la medida positiva \(\mu_0=12^{-1}\sum_m\delta_{s_m}\). Sus
momentos \(M_r=\int s^r\,d\mu_0(s)\) definen formas de Hankel estrictamente
positivas hasta el rango doce. Esta positividad procede de la evaluación de
un polinomio sobre doce nodos distintos: un polinomio de grado menor que
doce conserva al menos una evaluación no nula. En grado doce, el producto
de los doce factores nodales anula la norma y determina exactamente el rango
del lector.

La ortogonalización de estos momentos produce una matriz de Jacobi
\(J_{12}\), real simétrica, con subdiagonales positivas. Su cálculo funcional
recupera el perfil completo:
\[
u_0(x)=e_0^{\mathsf T}
\bigl[I-\cos(x\sqrt{J_{12}})\bigr]e_0.
\]
La suma de cosenos y el operador de Jacobi son dos realizaciones del mismo
lector. Los momentos mantienen todas las frecuencias y sus pesos; las
marcas de orientación, ruta y memoria permanecen en el estado que las
genera. El rango doce identifica la dimensión espectral de esta lectura
concreta. La profundidad de sus prolongaciones pertenece a otra dimensión
de la construcción: la longitud de las historias compatibles y la resolución
con que se evalúan.

La geometría analítica también se hace explícita. Para la deformación
ponderada \(u_\eta\), en la franja \(|\eta|\le1/40\), los pesos siguen siendo
positivos y el segundo momento permanece normalizado. En
\(|\operatorname{Re}z|<\pi/3\) existe una función holomorfa impar y univalente
\(b_\eta\), normalizada por \(b_\eta'(0)=1\), tal que
\(u_\eta=b_\eta^2\). La rueda produce, por tanto, una deformación conforme
del pliegue cuadrático. Su conjugación dinámica conserva esa deformación a
través de \(b_\eta(1-\lambda w^2)\). El teorema de escalado paramétrico que
se desarrolla aquí corresponde al miembro uniforme \(\eta=0\); las
propiedades estructurales de la franja ponderada y la persistencia local de
sus raíces conservan sus cuantificadores específicos.

\subsection{Profundidad de lectura y restitución de las historias}

La profundidad nonádica refina la lectura de un parámetro, mientras la
duplicación dinámica cambia el horizonte de retorno. El refinamiento con
acarreo expresa la primera operación mediante
\[
L_c(x,y)=(Bx-c,By+c),\qquad
C_m=\sum_{j=1}^m c_jB^{m-j},
\]
\[
(x_m,y_m)=(B^mx_0-C_m,B^my_0+C_m).
\]
El soporte exterior crece como \(B^m\), y el cilindro normalizado del registro
tiene diámetro \(B^{-m}\). La pareja entera, la profundidad y la marca de
frontera conservan la reconstrucción de cocientes, residuos y prefijos.
Crecimiento exterior y resolución interior describen así operaciones
compatibles sobre una misma historia.

La composición exacta
\(L_D^{[b]}L_C^{[a]}=L_{B^bC+D}^{[a+b]}\) permite agrupar una historia en
bloques sin alterar sus datos. Un retorno de duplicación reúne dos
transiciones; el calendario nonádico reúne nueve. Ambos agrupamientos se
comparan sobre un horizonte común de \(9\,2^N\) eventos, conservando el
orden de los subbloques y sus acarreos.

La contribución de la memoria puede observarse directamente en las ramas
inversas del perfil. Si \(v=(u_0|_{[0,1]})^{-1}\), las dos ramas son
\[
\beta_\sigma(y)=\sigma v\!\left(\frac{1-y}{\lambda}\right),
\qquad \sigma\in\{-1,+1\}.
\]
El valor final, junto con la última hoja, reconstruye el valor anterior. Una
sucesión de hojas permite recuperar todos los pasos intermedios. La palabra
de ramas acompaña los registros APP--TRIT--TPK de ruta, frontera, acarreo y
memoria; cada registro conserva su función. Esta restitución concreta
explica qué información requiere la inversión de la lectura crítica.

La resolución de cada retorno dispone además de una cota explícita. Con
\(t=\lambda u_\eta(1)/2\), \(z=(x+1)/2\) y
\[
F_{t,\eta}(z)=1-t\frac{u_\eta(2z-1)}{u_\eta(1)},
\]
se obtiene
\[
\left|F_{t,\eta}^{q}(1/2)-F_{s,\eta}^{q}(1/2)\right|
\le S_q|t-s|,\qquad S_q=\sum_{j=0}^{q-1}(6\sqrt2)^j<9^q.
\]
Un cilindro de base nueve y profundidad \(m=2^N+r\) permite, por ello, leer
el retorno de horizonte \(2^N\) con diámetro inferior a \(9^{-r}\). Esta
estimación da contenido cuantitativo a la profundidad arbitraria: para cada
horizonte y precisión solicitados determina una profundidad suficiente. La
evaluación de las funciones analíticas y el redondeo poseen sus propios
restos, que se propagan junto con la incertidumbre del parámetro.

Las vacancias intervienen en el balance de capacidad del registro. La
comparación entre bloques de 729 y 1000 posibilidades distingue capacidad
acumulada y longitud cronológica, conservando los pasos en que la primera
permanece constante. Esta contabilidad determina el espacio requerido para
conservar una lectura y su profundidad. La admisibilidad de las historias
procede, a su vez, de las reglas de supervivencia. Ambas operaciones
colaboran en la construcción y mantienen objetos matemáticos distintos.

\subsection{Cronología nonádica y combinatoria de duplicación}

La palabra crítica de duplicación tiene una expresión a toda profundidad:
\[
\tau_j=(-1)^{v_2(j)},\qquad
\tau_{2j}=-\tau_j,\qquad \tau_{2j+1}=+1.
\]
Sus prefijos satisfacen \(W_{n+1}=W_n(-1)^nW_n\). La regla determina cómo
se orientan las dos copias del retorno y dónde aparece su separación. La
lectura nonádica de la cronología conserva el entero completo
\(K=\rho_9(K)+9q_9(K)\). Al reducirlo módulo \(2^n\), la holonomía que
avanza nueve eventos induce la traslación por nueve sobre ese ciclo. Puesto
que nueve es impar, esta traslación visita todas sus posiciones. La
permutación \(j\mapsto9j\pmod{2^n}\), que la conjuga con el avance unitario,
preserva la valoración diádica de cada posición no nula. El residuo y el cociente
cronológicos sostienen conjuntamente esa compatibilidad.

La memoria bilateral posee asimismo una realización operatoria compatible
con el refinamiento. Sobre los ciclos de longitud \(2^n\), el operador de
avance \(C_n\) se incorpora en
\[
\mathbb U_n=P_0\otimes I+P_1\otimes C_n,
\]
con los proyectores complementarios del lector nonádico. En el límite, sus
componentes visible y complementaria satisfacen
\[
\|T_{\infty,a}v\|^2+\|D_{\infty,a}v\|^2=\|v\|^2,
\qquad
v=T_{\infty,a}^*T_{\infty,a}v+D_{\infty,a}^*D_{\infty,a}v.
\]
Estas identidades describen conservación y restitución dentro del espacio
y la representación declarados. La contracción de los mapas renormalizados,
que aparecerá después, mide otra operación: la aproximación de sus perfiles
a una forma estable. La historia conservada por el levantamiento y la
convergencia de una lectura funcional se articulan en niveles diferentes.

\subsection{La carta ordenada y la selección global de las raíces}

La publicación arquimediana del continuo HMT transporta las operaciones
mediante un isomorfismo de cuerpos ordenados completos, dentro del universo
declarado \(\mathfrak G\):
\[
\iota:\mathbb R_{\mathrm{HMT}}^{\mathfrak G}
\longrightarrow\mathbb R^{\mathfrak G}.
\]
La construcción interna del coseno por su serie, la raíz positiva de 899,
el perfil y sus iterados precede a la selección de ceros. La conservación
de suma, producto, orden y límites convergentes da
\[
\iota\bigl(S_n^{\mathrm H}(a)\bigr)=S_n\bigl(\iota(a)\bigr),
\qquad S_n(\lambda)=f_\lambda^{2^n}(0).
\]
La sobreyectividad cubre todas las raíces del dominio de esa carta; el orden
conserva los retornos anteriores, el período exacto y la condición de
encontrarse a la derecha de la raíz precedente. De este modo se transporta
también el mínimo de cada conjunto de candidatos. La genealogía llega
hasta la selección efectiva del parámetro, con la segunda proyección y sus
marcas conservadas sobre el mismo antecedente.

El selector queda fijado por \(\lambda_1=1/u_0(1)\) y por la primera raíz
nueva de período crítico exacto \(2^n\) situada a la derecha de
\(\lambda_{n-1}\). La clasificación de itinerarios demuestra que cada raíz
seleccionada tiene palabra \(W_n0\). Los retornos restrictivos reducen su
período a la mitad conservando el orden unimodal, y la inducción restituye
la palabra completa. La analiticidad de los retornos no idénticamente nulos
aísla sus ceros en cada compacto; la
comparación de los signos antes del primer parámetro de itinerario infinito
garantiza la existencia de una raíz nueva en todos los niveles.

Sea \(\Lambda_\tau\) el conjunto compacto no vacío de parámetros que realizan
la palabra infinita, y sea \(a_*\) su mínimo. Se demuestra
\[
\lambda_n\nearrow a_*.
\]
El paso al límite conserva cada signo mediante una cota uniforme a horizonte
fijo. Si \(c_p=f_\lambda^p(0)\), los signos opuestos de \(c_p\) y
\(c_{2p}\), junto con \((f_\lambda^p)'(0)=0\), implican
\[
|c_p|>\frac2{B_p},\qquad
B_p=16^p\frac{16^p-1}{15}.
\]
Así, los prefijos de períodos crecientes preservan la palabra infinita al
acumularse. La prueba utiliza la sucesión original de primeras raíces desde
su definición global.

\subsection{La hoja estable y la identificación del límite seleccionado}

El retorno normalizado
\[
(Rf)(x)=-\frac{f(f(-ax))}{a},\qquad a=-f(1)>0,
\]
reúne dos transiciones, cambia la escala y conserva la orientación necesaria
para restituir un máximo de valor uno. Los dominios restrictivos acreditan
su interpretación como retorno real. La derivada de este operador incluye
la variación de la escala \(a\); esa contribución determina la sensibilidad
paramétrica de la transformación completa.

En la carta \(f(x)=1-x^2V((x^2-1)/(5/2))\), con coordenadas \(u=10v_0\)
y \(\nu=(v_1,v_2,\ldots)\), el cuarto retorno de la familia entra en un
entorno cuantificado del punto fijo \(g\). La curva cruza transversalmente
su hoja estable en un único parámetro local
\[
1.874038<\lambda_*<1.874039,
\]
y satisface
\[
\|R^{4+n}f_{\lambda_*}-g\|_{\mathcal A}
<0.001666(0.83996)^n.
\]
La hoja estable es una gráfica de pendiente a lo sumo \(0.16\). Las
secantes paramétricas pertenecen a un cono transversal expansivo. Ambas
propiedades distinguen la disminución de las deformaciones estables y el
crecimiento de la separación entre parámetros.

La identificación de \(a_*\) y \(\lambda_*\) reúne esta geometría local
con el mínimo global. Primero se excluye el intervalo entre la octava raíz
y el extremo inferior de la caja local. Después, el orden
\(a_*\le\lambda_*\) permite estudiar una hipotética separación mediante
el cono negativo. Mientras su coordenada transversal tiene módulo inferior
a \(0.006\), la secante permanece en el dominio analítico y crece por un
factor comprendido entre \(4.4034\) y \(8.0887\). Una separación persistente
produciría, por tanto, una primera salida de amplitud comprendida entre
\(0.006\) y \(0.0485322\).

El residuo de la órbita estable conserva información más precisa que su
norma: sus componentes satisfacen \(|e_u|\le0.16\|e_\nu\|_1\). Esta
relación permite acotar conjuntamente sus efectos sobre la órbita crítica.
La cobertura de toda la región de primera salida muestra en cada caso un
signo incompatible con \(\tau\), o un retorno contractivo que obliga a dos
valores críticos diádicos sucesivos a tener el mismo signo. La palabra
infinita exige signos opuestos. La contradicción identifica los parámetros:
\[
\lambda_n\nearrow a_*=\lambda_*.
\]
La demostración combina una inducción válida a toda profundidad con una
exclusión finita uniforme de la región a la que llegaría cualquier
separación. El alcance infinito descansa en esa composición.

\subsection{Escalado universal y significado de la precisión}

La realización holomorfa de grado dos de los retornos y la transversalidad
acreditada permiten comparar sus clases híbridas. Para cada período
suficientemente alto, la clase canónica posee una única intersección local
con la familia. Las primeras raíces globales ya convergen al cruce y
conservan la combinatoria de duplicación; por ello coinciden finalmente
con los centros locales al igualar sus períodos físicos. La prueba
transporta la selección original hasta la ley métrica universal.

Existen \(C>0\), \(M<\infty\) y \(0<\theta<1\) tales que
\[
\lambda_*-\lambda_n=C\delta_{\mathrm F}^{-n}(1+e_n),
\qquad |e_n|\le M\theta^n,
\]
y, en consecuencia,
\[
\lim_{n\to\infty}
\frac{\lambda_{n-1}-\lambda_{n-2}}
{\lambda_n-\lambda_{n-1}}=\delta_{\mathrm F}.
\]
El factor \(\delta_{\mathrm F}\) es el autovalor expansivo del operador de duplicación
reconocido al final de la construcción. La familia concreta conserva sus
doce fases, sus momentos y su historia de procedencia; el límite de sus
razones paramétricas pertenece a la misma clase universal de criticidad
cuadrática. Esta relación explica conjuntamente la especificidad del
generador y la universalidad de la salida.

La precisión posee aquí tres residencias. Los intervalos de las raíces
controlan su evaluación finita. La estimación
\(0.001666(0.83996)^n\) controla la convergencia operatorial sobre la hoja
estable. El término \(M\theta^n\) controla el traslado de esa geometría a
la escala paramétrica mediante la holonomía transversal. Cada estimación
responde a su variable y a su operación. El octavo cociente queda encerrado
alrededor de \(4.669212189150204154556096215889663928\), con anchura inferior
a \(5.808\cdot10^{-37}\); esa anchura mide la evaluación de \(\delta_{\mathrm F,8}\).
La distancia de \(\delta_{\mathrm F,8}\) al límite requiere cuantificar las constantes
del resto paramétrico, cuya existencia ya queda demostrada.

La notación mantiene separados \(\delta_{\mathrm F}\), factor de escalado paramétrico
de Feigenbaum; \(\alpha_{\mathrm F}\), constante espacial asociada a su
problema de renormalización; y \(\alpha_{\mathrm{HMT}}\), constante de
estructura fina del corpus. El resultado presente identifica \(\delta_{\mathrm F}\)
mediante las primeras raíces del sector uniforme \(\eta=0\). La franja
ponderada conserva sus teoremas estructurales y de existencia, y su extensión
métrica cuantificada constituye un enunciado específico distinto.

La aportación del encadenamiento reside en la continuidad entre generación,
lectura, memoria, selección y límite. La aritmética orientada determina el
perfil; la representación espectral conserva sus doce componentes; los
cilindros permiten resolverlo con precisión arbitraria a cada horizonte;
el transporte de historias restituye sus recorridos; la carta ordenada
conserva todas las raíces y sus mínimos; la clasificación prolonga el
selector; y la expansión transversal, junto con la contracción estable,
determina su ley asintótica. La universalidad aparece como una propiedad
demostrada de esta construcción completa y de su realización analítica.

\endgroup

% Fragmento compartido por VII y CRISTAL; incluir despues del cuerpo Feigenbaum.
% Procedencia: resultados recuperados, con demostraciones elementales reunidas.
% Perfil, Jacobi y escalado: 00_antecedentes_operatorios.tex, 03, 06 y 07
% del presente paquete; no se cambia su selector ni su dominio uniforme.
% Constitutiva: CRISTAL_ES/main_autosuficiente.tex:58738--58840,58864--58894.
% Inversion: CRISTAL_ES/sections/cr28_restitucion_constitutiva.tex:18--59.
% Catalan: el mismo propietario, 277--373; identidad integral, unicidad y serie.
% Corte material: CUMPLIMIENTO_EDITORIAL_20260928/fuentes/CRISTAL_ES.
% No contiene una nueva calibracion ni una nueva campana de calculo.
\section{Articulación de los lectores crítico, espectral y constitutivo}
\label{hmtfg:articulacion:seccion}

La composición APP--TRIT--TPK precede a los lectores de esta sección.
APP conserva en cada célula las evaluaciones aditiva y multiplicativa,
con sus residuos y cocientes; la reducción trítica determina el régimen,
la orientación y el acarreo de la transición. La composición efectiva
de selección, transporte y actualización de memoria produce el estado
enriquecido y sus prolongaciones, desarrollados en
\ref{hmtfg:app-trit-transporte} y
\ref{hmtfg:subsubsec:tres-vueltas-ext-ch}.
La holonomía nonádica compone esas transiciones: el retorno de fase
avanza la memoria y conserva la historia. Los lectores actúan sobre la
misma estructura discreta conjunta del continuo, con sus cinco
construcciones consustanciales y con las dos publicaciones del mismo
antecedente: evaluación ordenada e incidencia de frontera.

La rueda orientada y los canales angulares son salidas de esta
composición. Sus coordenadas HMT ya generadas se utilizan en operaciones
posteriores, con sus marcas y dominios; no son cifras metrológicas
introducidas para seleccionar una respuesta. La publicación ordenada
permite evaluar las fórmulas siguientes, mientras el estado de procedencia
conserva hoja, ruta, residuo, cociente, acarreo, frontera y memoria.
Las restituciones que se prueban recuperan el objeto de su dominio
declarado; la representación evaluada permanece acompañada de esos
registros.

\subsection{La separación crítica como lectura asintótica intrínseca}
\label{hmtfg:articulacion:separacion}

En el sector uniforme de la rueda dodecafásica, los representantes
orientados y sus pesos son
\[
 \phi_m=\frac{(3m-1)\pi_{\mathrm{HMT}}}{18},\qquad
 w_m=\frac1{12},\qquad 1\le m\le12.
\]
El segundo momento se normaliza mediante
\(\sum_{m=1}^{12}(3m-1)^2=5394=6\cdot899\).
De este modo, el factor angular común se cancela y quedan
\begin{equation}
 k_m=\frac{2(3m-1)}{\sqrt{899}},\qquad
 u_0(z)=\frac1{12}\sum_{m=1}^{12}\bigl(1-\cos(k_mz)\bigr),
 \qquad f_\lambda(z)=1-\lambda u_0(z).
 \label{hmtfg:articulacion:perfil}
\end{equation}
La variable \(\lambda\) recorre la familia después de fijados el
lector, sus pesos y la normalización. En
\(0<\lambda\le2/u_0(1)\), el dominio dinámico es \([-1,1]\).
La primera raíz es \(\lambda_1=1/u_0(1)\); para \(n\ge2\),
\(\lambda_n\) es la menor raíz mayor que \(\lambda_{n-1}\)
cuyo punto crítico tiene período exacto \(2^n\).
La clasificación, la existencia de cada mínimo y la identificación
métrica de estos centros se han demostrado en
\ref{hmtfg:escalado:15}--\ref{hmtfg:escalado:17}.

\begin{proposition}[Lectura estructural de la constante de escalado]
\label{hmtfg:articulacion:delta}
Para el lector uniforme \eqref{hmtfg:articulacion:perfil} y su selector
de primeras raíces, la constante de separación crítica es
\begin{equation}
 \delta_{\mathrm{HMT,crit}}
 :=\lim_{n\to\infty}
 \frac{\lambda_n-\lambda_{n-1}}{\lambda_{n+1}-\lambda_n}
 =\delta_{\mathrm F}.
 \label{hmtfg:articulacion:limite}
\end{equation}
Es el invariante asintótico de separación de parámetros de los retornos
críticos de ese lector, con su orden de selección conservado.
\end{proposition}
\begin{proof}
El teorema de escalado \eqref{hmtfg:escalado:eq:67} proporciona
\(\lambda_*-\lambda_n=C\delta_{\mathrm F}^{-n}(1+e_n)\),
con \(C>0\) y \(e_n\to0\). Restando dos términos consecutivos,
\[
 \lambda_n-\lambda_{n-1}
 =C\delta_{\mathrm F}^{-n}
   \bigl[\delta_{\mathrm F}-1+
          \delta_{\mathrm F}e_{n-1}-e_n\bigr].
\]
El cociente de esta igualdad y la del índice \(n+1\) converge a
\(\delta_{\mathrm F}\). La identidad eventual de los centros locales
con las primeras raíces, demostrada en
\eqref{hmtfg:escalado:eq:66}, permite aplicar el escalado precisamente
a la sucesión seleccionada y no a otra cascada elegida después.
\end{proof}

Esta formulación conserva más información que el numeral terminal:
identifica la rueda, su lectura, la familia y el selector cuyo límite
se evalúa. Su interpretación es intrínseca a esa composición HMT;
el reconocimiento de la constante universal utiliza después el teorema
analítico aplicado al lector ya construido.

\subsection{La medida finita y la restitución espectral del perfil}
\label{hmtfg:articulacion:jacobi}

Las frecuencias cuadradas producen la probabilidad atómica
\begin{equation}
 s_m=k_m^2=\frac{4(3m-1)^2}{899},\qquad
 \nu_0=\frac1{12}\sum_{m=1}^{12}\delta_{s_m},\qquad
 M_j=\int s^j\,d\nu_0(s).
 \label{hmtfg:articulacion:medida-finita}
\end{equation}
Aquí \(\delta_{s_m}\) es la masa de Dirac en un nodo, distinta
de la constante \(\delta_{\mathrm F}\). La medida conserva las
doce frecuencias y sus pesos. Para un polinomio no nulo de grado menor
que \(r\le12\),
\[
 \int p(s)^2\,d\nu_0(s)
 =\frac1{12}\sum_{m=1}^{12}p(s_m)^2>0,
\]
pues un polinomio de grado menor que doce no puede anularse en los
doce nodos distintos. Así, las matrices de Hankel
\((M_{i+j})_{0\le i,j<r}\) son positivas definidas.
El polinomio \(\prod_m(s-s_m)\) tiene norma nula; el espacio
\(L^2(\nu_0)\) tiene dimensión doce.

\begin{proposition}[Representación exacta del lector crítico]
\label{hmtfg:articulacion:representacion}
Sean \(\psi_0,\ldots,\psi_{11}\) los polinomios ortonormales
de coeficiente principal positivo para \(\nu_0\), con
\(\psi_0=1\). La multiplicación por \(s\) tiene matriz de
Jacobi \(J_{12}\), real simétrica y de subdiagonal positiva.
Con \(e_0=(1,0,\ldots,0)^{\mathsf T}\),
\begin{equation}
 u_0(z)=\int\bigl(1-\cos(z\sqrt{s})\bigr)\,d\nu_0(s)
 =e_0^{\mathsf T}\bigl[I-\cos(z\sqrt{J_{12}})\bigr]e_0.
 \label{hmtfg:articulacion:perfil-jacobi}
\end{equation}
\end{proposition}
\begin{proof}
La positividad anterior permite ortogonalizar hasta grado once.
Si \(i<j-1\), la autoadjunción de la multiplicación da
\(\langle s\psi_j,\psi_i\rangle
=\langle\psi_j,s\psi_i\rangle=0\).
Por simetría, sólo quedan la diagonal y sus dos vecinas; el cociente
positivo de los coeficientes principales fija la positividad de la
subdiagonal. Definamos
\[
 Q_{mj}=\frac{\psi_j(s_m)}{\sqrt{12}},\qquad
 D=\operatorname{diag}(s_1,\ldots,s_{12}).
\]
La ortonormalidad da \(Q^{\mathsf T}Q=I\), y evaluar la
recurrencia en los nodos da \(DQ=QJ_{12}\). Como \(Q\) es
cuadrada, \(J_{12}=Q^{\mathsf T}DQ\). Sus autovalores son
positivos, por lo que la raíz positiva y el coseno se calculan en
esta diagonalización. Finalmente,
\(Qe_0=(1,\ldots,1)^{\mathsf T}/\sqrt{12}\), lo que prueba
\eqref{hmtfg:articulacion:perfil-jacobi}.
\end{proof}

La realización por \((J_{12},e_0)\) recupera el mismo perfil;
por ello, al conservar la familia y el selector, recupera también
sus retornos y sus primeras raíces. La positividad espectral
justifica esta representación. La transversalidad paramétrica
utiliza además la derivada del retorno y sus productos orbitales,
desarrollados en \ref{hmtfg:escalado:14.3} y
\ref{hmtfg:escalado:6}. Son propiedades de operaciones diferentes:
la demostración del escalado conserva ambas y no sustituye el
control de la tangente por la positividad de los momentos.

\subsection{Los canales constitutivos y su inversión etiquetada}
\label{hmtfg:articulacion:constitutiva}

El lector constitutivo actúa sobre otro objeto tipado del mismo
soporte: el transporte angular de dos hojas. Sus coordenadas ya
generadas son \(x=A_{\rm rad}\) y \(y=C^*_{\rm rad}\), en la
cámara contractiva \(x>|y|\). La involución de hojas
\(R=R^*=R^{-1}\) determina los proyectores marcados
\(P_\pm=(I\pm R)/2\). La representación del transporte es
\begin{equation}
 q_+=\exp[-(x+y)],\qquad q_-=\exp[-(x-y)],\qquad
 T=q_+P_++q_-P_-,\qquad 0<q_\pm<1.
 \label{hmtfg:articulacion:canales}
\end{equation}
La evaluación exponencial realiza esos canales después de su
construcción angular. Las incidencias \(90\) y \(120\) del
transporte fijan la respuesta
\[
 \mathcal V=(I-T^{90})(I-T^{120})^{-1}.
\]
El inverso existe porque \(1-q_\pm^{120}>0\).
Escribiendo \(s=q^{30}\), con \(30=\gcd(90,120)\), se obtiene
\begin{equation}
 f(s)=\frac{1-s^3}{1-s^4}
     =\frac{1+s+s^2}{1+s+s^2+s^3},\qquad
 r_\pm=f(q_\pm^{30}),\qquad
 \mathcal V=r_+P_++r_-P_-.
 \label{hmtfg:articulacion:respuesta}
\end{equation}
En esta subsección \(f\) designa la respuesta constitutiva;
\(f_\lambda\) sigue designando el mapa dinámico de
\eqref{hmtfg:articulacion:perfil}. Sus argumentos y dominios
son distintos.

\begin{proposition}[Restitución de los canales y de la orientación]
\label{hmtfg:articulacion:inversion}
La función constitutiva \(f:(0,1)\to(3/4,1)\) es una biyección
decreciente. Su inversa \(\psi\) asigna a \(r\) la única raíz
\(s\in(0,1)\) de
\begin{equation}
 r s^3+(r-1)(1+s+s^2)=0.
 \label{hmtfg:articulacion:cubica}
\end{equation}
Las lecturas constitutivas normalizadas, con sus hojas etiquetadas,
\[
 \widehat\varepsilon=r_+^2,\qquad
 \widehat\mu=r_-^2,\qquad
 \widehat Z=\frac{r_-}{r_+},\qquad
 \widehat c=\frac1{r_+r_-},
\]
con \(\widehat\varepsilon,\widehat\mu\in(9/16,1)\), permiten recuperar
\begin{equation}
 \begin{aligned}
 r_+&=\sqrt{\widehat\varepsilon},&
 r_-&=\sqrt{\widehat\mu},\\
 s_\pm&=\psi(r_\pm),& q_\pm&=s_\pm^{1/30}.
 \end{aligned}
 \label{hmtfg:articulacion:recuperacion-canales}
\end{equation}
\begin{equation}
 x=-\frac12\log(q_+q_-),\qquad
 y=\frac12\log\frac{q_-}{q_+}.
 \label{hmtfg:articulacion:recuperacion-angular}
\end{equation}
\end{proposition}
\begin{proof}
La derivación del cociente racional da
\begin{equation}
 f'(s)=-\frac{s^2(3+2s+s^2)}{(1+s+s^2+s^3)^2}<0
 \qquad(0<s<1).
 \label{hmtfg:articulacion:monotonia}
\end{equation}
Los límites \(f(0^+)=1\) y \(f(1^-)=3/4\), junto con la
continuidad, prueban la biyección. Multiplicar \(f(s)=r\) por
el denominador positivo produce \eqref{hmtfg:articulacion:cubica}.
La positividad de \(r_\pm\) determina las raíces cuadradas
y la de \(s_\pm\) determina las raíces de orden treinta.
Por \eqref{hmtfg:articulacion:canales},
\(\log q_+=-x-y\) y \(\log q_-=-x+y\);
sumar y restar recupera \(x,y\).
Recíprocamente, las raíces recuperadas pertenecen a \((0,1)\),
de modo que \(x\pm y=-\log q_\pm>0\); la pareja obtenida
permanece en la cámara \(x>|y|\).
\end{proof}

En particular,
\(\widehat\mu\widehat\varepsilon=\widehat c^{-2}\) y
\(\widehat\mu/\widehat\varepsilon=\widehat Z^2\).
El intercambio de las etiquetas permuta \(q_+\) y \(q_-\),
conserva \(x\) y cambia el signo de \(y\). La restitución
requiere esa orientación: el producto aislado \(r_+r_-\)
conserva la propagación común, pero no determina las dos respuestas
etiquetadas. Las raíces positivas y la inversión cúbica recuperan
los canales dentro del lector fijado; el registro completo de su
generación permanece en el estado enriquecido.

\subsection{El momento de Catalán y la restitución funcional}
\label{hmtfg:articulacion:catalan}

Las mismas incidencias determinan una relación entre funciones,
antes de evaluar cualquiera de los canales. Para \(0<s<1\), sea
\(\mathcal D=s\,d/ds\), y definamos el momento orientado
\begin{equation}
 \mathcal C_2(s)
 =\int_0^s\log\frac{s}{t}\,
     \frac{1+t}{1+t+t^2}f(t)\,dt
 =\int_0^s\frac{\log(s/t)}{1+t^2}\,dt.
 \label{hmtfg:articulacion:integral}
\end{equation}
La segunda igualdad procede de la factorización
\(1+t+t^2+t^3=(1+t)(1+t^2)\). El peso continuo y positivo
\(dt/(1+t^2)\) es así producido por la composición racional
del lector constitutivo, y no por una identificación con
\(\nu_0\).

\begin{proposition}[Inversión diferencial, unicidad y lectura extrema]
\label{hmtfg:articulacion:restitucion-catalan}
El momento \(\mathcal C_2\) es la única función
\(F\in C^2((0,1))\) que satisface
\begin{equation}
 \mathcal D^2F(s)=\frac{s(1+s)}{1+s+s^2}f(s),\qquad
 \lim_{s\downarrow0}F(s)=0.
 \label{hmtfg:articulacion:problema-diferencial}
\end{equation}
La respuesta constitutiva se restituye por
\begin{equation}
 \mathcal D^2\mathcal C_2(s)=\frac{s}{1+s^2},\qquad
 f(s)=\frac{1+s+s^2}{s(1+s)}\,
          \mathcal D^2\mathcal C_2(s).
 \label{hmtfg:articulacion:inversa-diferencial}
\end{equation}
Su extensión continua al intervalo cerrado tiene la serie
\begin{equation}
 \mathcal C_2(s)=\sum_{j=0}^{\infty}
       \frac{(-1)^j s^{2j+1}}{(2j+1)^2},\qquad 0\le s\le1,
 \qquad
 \mathcal C_2(1)=\sum_{j=0}^{\infty}\frac{(-1)^j}{(2j+1)^2}
 =G_{\rm Cat}.
 \label{hmtfg:articulacion:serie-catalan}
\end{equation}
\end{proposition}
\begin{proof}
El integrando de \eqref{hmtfg:articulacion:integral} es no negativo
y
\[
 0\le\mathcal C_2(s)\le\int_0^s\log(s/t)\,dt=s.
\]
Esto prueba integrabilidad en cero y el límite requerido.
La identidad \(\log(s/t)=\int_t^s du/u\), seguida del
intercambio de integrales no negativas, proporciona
\[
 \mathcal C_2(s)
 =\int_0^s\frac1u\left(\int_0^u\frac{dt}{1+t^2}\right)du.
\]
En todo intervalo compacto del dominio abierto se aplica el
teorema fundamental del cálculo. Por tanto,
\[
 \mathcal D\mathcal C_2(s)=\int_0^s\frac{dt}{1+t^2},
 \qquad
 \mathcal D^2\mathcal C_2(s)=\frac{s}{1+s^2}.
\]
La factorización racional usada en
\eqref{hmtfg:articulacion:integral} da
\eqref{hmtfg:articulacion:problema-diferencial}; despejar \(f\)
da \eqref{hmtfg:articulacion:inversa-diferencial}, cuyos
denominadores son positivos en \((0,1)\).

Si dos soluciones tienen el límite nulo requerido, su diferencia
\(H\) cumple \(\mathcal D^2H=0\). El cambio \(z=\log s\)
transforma esta ecuación en \(d^2H(e^z)/dz^2=0\), de donde
\(H(s)=a\log s+b\). La existencia del límite nulo obliga a
\(a=b=0\), y prueba la unicidad.

Para \(s<1\), la expansión geométrica de \((1+t^2)^{-1}\)
puede integrarse término a término con el peso logarítmico, pues
\[
 \int_0^s t^{2j}\log(s/t)\,dt
 =\frac{s^{2j+1}}{(2j+1)^2},\qquad
 \sum_{j\ge0}\frac{s^{2j+1}}{(2j+1)^2}<\infty.
\]
Esto prueba la serie en el interior. La cota
\(1/(2j+1)^2\) da convergencia absoluta y uniforme de esa
serie en \([0,1]\). En la integral, el integrando extendido
por cero para \(t>s\) está dominado por \(\log(1/t)\)
en \((0,1)\). La convergencia dominada permite alcanzar
\(s=1\) y coincide con la extensión de la serie. Su valor
extremo es precisamente la serie de Catalán indicada.
\end{proof}

La inversión diferencial utiliza la familia \(\mathcal C_2(s)\),
no solamente su valor extremo \(G_{\rm Cat}\). En particular,
\(\mathcal D^2\mathcal C_2(s)\to1/2\) cuando \(s\uparrow1\).
La serie diferenciada dos veces con \(\mathcal D\) tiene allí
un límite radial de Abel; la serie del propio momento es
absolutamente convergente en la frontera. Esta distinción mantiene
las condiciones precisas de cada paso al límite.

\subsection{Composición de las restituciones y alcance conjunto}
\label{hmtfg:articulacion:conclusion}

Las interrelaciones anteriores tienen operaciones y dominios
explícitos. Los momentos de la medida atómica \(\nu_0\)
producen \(J_{12}\), y su cálculo funcional recupera el perfil
\(u_0\). El perfil y el selector de primeras raíces determinan
la sucesión cuyo escalado publica \(\delta_{\mathrm F}\).
Por otra parte, la función constitutiva \(f\) produce
\(\mathcal C_2\) mediante el momento integral, y
\eqref{hmtfg:articulacion:inversa-diferencial} recupera \(f\)
a partir de esa familia. Sus evaluaciones etiquetadas en
\(s_\pm=q_\pm^{30}\) determinan las respuestas constitutivas;
la inversión cúbica recupera ambos canales y sus coordenadas
angulares. La lectura extrema del momento da \(G_{\rm Cat}\).

La medida \(\nu_0\) es positiva, atómica y de rango doce.
El momento de Catalán integra un peso positivo continuo y posee
una expansión alternada. Los signos de esa expansión son
coeficientes de la serie geométrica; no son pesos negativos de
\(\nu_0\), ni convierten la integral positiva en la misma
representación espectral. Cada lector conserva su soporte,
orientación, normalización y operación de evaluación.

La unidad reside en la estructura APP--TRIT--TPK que genera y
transporta sus antecedentes; las relaciones funcionales se
establecen por las composiciones demostradas. La igualdad de
origen no identifica operadores de dominios diferentes.
En particular, esta articulación no postula una fórmula escalar
que determine \(\delta_{\mathrm F}\) a partir de Catalán y de
otras constantes físicas. Sitúa el escalado crítico junto a
restituciones efectivas entre lectores del mismo soporte,
conservando la información requerida por cada dirección.

\setcounter{secnumdepth}{\HMTFGSavedSecnumdepth}
\endgroup

\endgroup

\chapter{Identidades de Rogers--Ramanujan y levantamiento nonádico de Pell}
\begingroup
\renewcommand{\chapter}[2][]{}%
\chapter{Aritmética del modo \(30\): identidades de Rogers--Ramanujan, normalizador \(899\) y levantamiento nonádico de Pell}
\label{chap:p3-final:46:modo_30_pell}
\section{Normalizador \(899\), unidad de Pell y estructura aritm\'etica del modo \(30\)}

El normalizador posee la factorizaci\'on aritm\'etica

\begin{equation}
899=30^2-1=29\cdot31.
\label{mab:rm:eq:899}
\end{equation}

El mismo modo \(30\) que genera los sectores \(90\) y \(120\) determina
simult\'aneamente el par de primos gemelos \(29,31\). Adem\'as,

\begin{equation}
\varepsilon_{30}=30+\sqrt{899},\qquad
\varepsilon_{30}^{-1}=30-\sqrt{899},\qquad
\log\varepsilon_{30}=\arcosh30.
\label{mab:rm:eq:pell30}
\end{equation}

Esta unidad de Pell determina la escala hiperb\'olica asociada al
normalizador y relaciona la fase dodecaf\'asica, el modo \(30\), los primos
gemelos y la geometr\'ia hiperb\'olica. Esta relaci\'on no identifica el
operador constitutivo del vac\'io con el renormalizador de duplicaci\'on; tal
identificaci\'on requerir\'ia un entrelazador expl\'icito.

La unidad act\'ua de forma integral mediante
\begin{equation}
M_{30}=\begin{pmatrix}30&899\\1&30\end{pmatrix},
\qquad \det M_{30}=1.
\label{mab:rm:eq:pell-matrix30}
\end{equation}
Sobre \(\Z[\sqrt{899}]\), esta matriz representa la multiplicaci\'on por
\(30+\sqrt{899}\). Sus autovalores son \(\varepsilon_{30}^{\pm1}\); por
tanto \(\log\varepsilon_{30}\) es la longitud de traslaci\'on de la transformaci\'on
hiperb\'olica asociada. El mismo entero central \(30\) reaparece en el vac\'io
porque \(90=3\cdot30\) y \(120=4\cdot30\), pero los dos operadores siguen
teniendo dominios distintos. La confluencia demostrada es la del modo y sus
completaciones; la identidad de din\'amicas no se postula.

 \section{Levantamiento non\'adico de la unidad de Pell}

Para \(k\ge1\), sea el anillo cuadr\'atico no ramificado
\begin{equation}
\mathcal R_k=(\Z/3^k\Z)[r]/(r^2-2),
\label{eq:pell-ring}
\end{equation}
y sea

\[
u=3+2r.
\]

La conjugaci\'on \(r\mapsto-r\) da
\[
N_{\mathcal R_k/(\Z/3^k\Z)}(u)
=(3+2r)(3-2r)=9-8=1,
\]
de modo que \(u\) es una unidad de norma uno. Adem\'as,
\[
u^2=17+12r,\qquad
u^4-1=576+408r=3(192+136r).
\]
El segundo factor es unidad m\'odulo tres porque reduce a \(r\ne0\). Por
tanto la valuaci\'on \(3\)-\'adica satisface
\(v_3(u^4-1)=1\). El lema de elevaci\'on para una unidad congruente con uno
m\'odulo tres produce, para todo \(j\ge0\),
\[
v_3\!\left(u^{4\cdot3^j}-1\right)=j+1.
\]
Como \(u\bmod3=2r\) y \((2r)^2=-1\), su orden m\'odulo tres es cuatro.
En consecuencia,

\begin{equation}
\operatorname{ord}(u)=4\cdot3^{k-1},
\label{eq:pell-order}
\end{equation}

y el cierre pro-finito de los grupos c\'iclicos compatibles es
\begin{equation}
\varprojlim_k\langle u\bmod3^k\rangle
\simeq C_4\times\mathbb Z_3.
\label{eq:pell-profinite}
\end{equation}
Su realizaci\'on arquimediana es

\[
3+2\sqrt2=(1+\sqrt2)^2.
\]

Esta construcci\'on demuestra la conservaci\'on de la profundidad non\'adica
por una unidad hiperb\'olica. Las unidades \(3+2\sqrt2\),
\(2+\sqrt3\) y \(30+\sqrt{899}\) forman una familia de Pell con
procedencias distintas; cualquier identificaci\'on entre ellas requiere
mapas expl\'icitos entre las extensiones correspondientes.

 El modo \(30\) es uno de los puntos en los que la arquitectura empieza a mostrar una unidad determinada por la genealogía operatoria y aritmética.\par

Aparece primero en el vacío porque\par

\[
90=3\cdot30,
\qquad
120=4\cdot30.
\]

Es la cadencia elemental que permite comparar los sectores \(90\) y \(120\). El vacío lee esa cadencia como completación \(3\to4\).\par

Pero el mismo \(30\) reaparece en la construcción del germen de criticidad. Allí la identidad central es\par

\[
899=30^2-1=29\cdot31.
\]

La normalización de las doce fases produce \(\sqrt{899}\), y el perfil analítico puede escribirse de forma cerrada. Lo realmente llamativo es que \(\pi\) desaparece antes de comenzar la iteración.\par

Esto debe leerse con atención. La geometría angular de \(\pi\) queda absorbida en la organización de las fases. Una vez normalizada la configuración dodecafásica, el germen dinámico queda gobernado por la aritmética interna de \(899\).\par

Por eso la construcción comienza generando internamente un objeto dinámico propio:\par

\begin{itemize}[leftmargin=*,itemsep=0.35em]
\item es par;
\item es analítico;
\item tiene un crítico cuadrático;
\item pertenece a la clase unimodal correspondiente;
\item posee Schwarziano negativo;
\item puede someterse al operador de renormalización.
\end{itemize}

Las razones superestables empiezan entonces a aproximarse a las constantes universales. El germen ha sido producido primero y la universalidad aparece como comportamiento terminal de su dinámica.\par

Ésta es una diferencia epistemológica profunda. Una coincidencia numérica diría: «nuestro cálculo se parece a Feigenbaum». La construcción HMT dice: «éste es el objeto exacto que entra en la clase de universalidad, y éstos son los cocientes que emergen al iterarlo».\par

El \(30\) reaparece todavía una tercera vez en la memoria modular. Si\par

\[
N=N_{90}+N_{120}
\]

mide la ocupación total y\par

\[
D=90N_{90}+120N_{120}
\]

mide la ocupación ponderada por procedencia, la matriz que transforma \((N_{90},N_{120})\) en \((N,D)\) tiene determinante \(30\).\par

Eso permite invertir el proceso:\par

\[
N_{90}=4N-\frac D{30},
\qquad
N_{120}=\frac D{30}-3N.
\]

De modo que el mismo \(30\) cumple tres funciones:\par

\begin{itemize}[leftmargin=*,itemsep=0.35em]
\item en el vacío, hace conmensurables las propagaciones \(90\) y \(120\);
\item en la criticidad, determina la normalización \(899=30^2-1\);
\item en la memoria modular, es el determinante que permite recuperar el origen de cada ocupación.
\end{itemize}

El operador del vacío, el germen de caos y el Hamiltoniano modular permanecen tipológicamente distintos. La confluencia es más sutil: los tres conservan una partición común de la genealogía.\par

Podríamos decir que el modo \(30\) es un engranaje. En una máquina mide propagación; en otra normaliza criticidad; en otra hace reversible la pérdida aparente de procedencia. El engranaje es el mismo, pero cada máquina realiza una función distinta.\par

La aparición de la unidad de Pell\par

\[
30+\sqrt{899}
\]

refuerza esta interpretación. El \(30\) organiza simultáneamente una suma finita y un crecimiento hiperbólico. La misma aritmética que cierra el germen abre una dinámica de escala.\par

Ésa es quizá la razón profunda por la que el resultado resulta tan bello: el modo que hace posible la completación finita del vacío es también el modo que prepara una universalidad infinita de la criticidad.\par

Finito e infinito se articulan mediante la repetición coherente de una estructura finita que conserva memoria; de ella nace el infinito dinámico.\par

  \endgroup
\begin{HMTSpiralChapterBody}
% Corte mecánico íntegro: parte_ii.tex:323--419.
% Residencia material del ensamblaje sucesor de 102 capítulos.
% Residencia causal: capítulo 46.

\chapter{Tornillo espectral de Pell y suspensión logarítmica}

\section{Ley de tornillo}

El corpus contiene la identidad exacta
\begin{equation}
 V^{-1}\mathscr S V=(2+\sqrt3)\ii\,\mathscr S.
 \label{espiral:eq:tornillo-pell}
\end{equation}
Sea \(\varepsilon=2+\sqrt3\). La órbita transversal
\[
 z_n=z_0(\varepsilon\ii)^n
\]
satisface
\[
 r_n=r_0\varepsilon^n,
 \qquad
 \theta_n=\theta_0+\frac{\pi n}{2}.
\]
Eliminando \(n\),
\begin{equation}
 r(\theta)=r_0
 \exp\!\left[
 \frac{2\log(2+\sqrt3)}{\pi}
 (\theta-\theta_0)
 \right].
 \label{espiral:eq:espiral-log-pell}
\end{equation}
La espiral logarítmica es una salida del operador: no se eligió su ecuación
para ajustar después \(\varepsilon\).

\begin{proof}[Derivación de \cref{espiral:eq:espiral-log-pell}]
De \(\theta_n-\theta_0=n\pi/2\) se sigue
\(n=2(\theta_n-\theta_0)/\pi\). Sustituyendo en
\(r_n=r_0\varepsilon^n\),
\[
 r_n=r_0\exp\left(
 \frac{2\log\varepsilon}{\pi}(\theta_n-\theta_0)
 \right).
\]
La interpolación continua es la ecuación indicada.
\end{proof}

\section{Involución conjugada de contracción}

La extensión bilateral \(n\in\Z\) produce
\[
 n\to+\infty:\ r_n\to\infty,
 \qquad
 n\to-\infty:\ r_n\to0.
\]
La inversión \(n\mapsto-n\) intercambia expansión y contracción. Como
\(\varepsilon^{-1}=2-\sqrt3\), ambas ramas pertenecen a la misma unidad de
Pell. No son dos espirales independientes, sino dos orientaciones del mismo
tornillo espectral.

\begin{figure}[htbp]
\centering
\begin{tikzpicture}[scale=.82]
 \draw[->,HMTGray] (-5.1,0)--(5.1,0) node[right] {$x$};
 \draw[->,HMTGray] (0,-3.1)--(0,3.1) node[above] {$y$};
 \draw[HMTBlue,very thick,domain=-7:5.2,samples=240,smooth,variable=\t]
   plot ({.27*exp(.22*\t)*cos(deg(\t))},{.27*exp(.22*\t)*sin(deg(\t))});
 \draw[HMTRed,thick,domain=-7:5.2,samples=240,smooth,variable=\t]
   plot ({.27*exp(.22*\t)*cos(deg(-\t))},{.27*exp(.22*\t)*sin(deg(-\t))});
 \foreach \k in {-4,-3,...,4}{
   \pgfmathsetmacro{\ang}{1.57079632679*\k}
   \fill[HMTNavy]
     ({.27*exp(.22*\ang)*cos(deg(\ang))},
      {.27*exp(.22*\ang)*sin(deg(\ang))}) circle (.8pt);
 }
 \node[HMTBlue,font=\small\sffamily] at (3.5,2.4) {expansión};
 \node[HMTRed,font=\small\sffamily] at (3.5,-2.4) {orientación conjugada};
\end{tikzpicture}
\caption{Proyección transversal del tornillo de Pell. Las marcas representan cuartos de giro.}
\label{espiral:fig:tornillo-pell}
\end{figure}

\section{Covariancia de escala y conservación energética}

La dilatación \(\varepsilon\) no debe interpretarse como crecimiento
gratuito de la amplitud energética de una onda en vacío. Su realización
coherente actúa sobre la escala espacial o espectral:
\[
 (r,\varphi,\lambda)
 \longmapsto
 (\varepsilon r,\varphi+\pi/2,\varepsilon\lambda).
\]
Esta distinción será esencial en la realización electromagnética de la
\cref{espiral:chap:realizacion-em}.

\begin{resultbox}[title={Compatibilidad exacta entre el tornillo espectral de Pell y la suspensión fase--memoria}]
La elevación fase--memoria mínima es una hélice circular. El tornillo
espectral combina un cuarto de giro y una dilatación de Pell; su proyección
es una espiral logarítmica. Ambos proceden del registro HMT, pero no son la
misma periodicidad ni se obtienen el uno del otro por olvido de tipos.
\end{resultbox}
 \end{HMTSpiralChapterBody}

\chapter{Teoría aritmética de las vacancias nonádicas y torsión de frontera}
\begingroup
\renewcommand{\chapter}[2][]{}%
\chapter{Teoría aritmética de las vacancias nonádicas: defecto \(10^3-3^6\), red radical \(3\to6\to9\), ley potencia--vacancia y torsión de borde}
\label{chap:p3-final:47:vacancias_nonadicas}
\section{Clasificación de los ciclos aritméticos irreducibles de la red radical nonádica}

En la realizaci\'on aritm\'etica, un primo corresponde a un ciclo
irreducible que no se factoriza en ciclos positivos menores. La periodicidad
non\'adica puede seleccionar candidatos y organizar clases residuales; la
primalidad permanece definida por la propiedad multiplicativa exacta

\[
p=ab\Longrightarrow a=1\ \text{o}\ b=1.
\]

Una vez seleccionado el coproducto de ciclos aritm\'eticos primos, se
obtienen \cref{eq:meissel-mertens,eq:twin-prime-constant}. La estructura
helicoidal y la ra\'iz digital parametrizan fase y acarreo, pero no sustituyen
el criterio de factorizaci\'on.

\section{Densidad, t\'ermino regularizado y holonom\'ia de vacancias}

La densidad positiva b\'asica es

\begin{equation}
\varepsilon_{\rm vac}=\log_{10}\frac{10}{9}.
\label{eq:vacancy-density}
\end{equation}

La secuencia se deduce de la completaci\'on
\(1000=729+271\) determina
\begin{equation}
\varepsilon_{\rm vac}
=1-\log_{1000}(729)
=\log_{10}\frac{10}{9},
\label{eq:vacancy-density-two-forms}
\end{equation}
y la palabra mec\'anica exacta
\begin{equation}
z_n=\lfloor(n+1)\varepsilon_{\rm vac}\rfloor
-\lfloor n\varepsilon_{\rm vac}\rfloor,
\qquad n\ge1.
\label{eq:vacancy-mechanical-word}
\end{equation}
Como \(\varepsilon_{\rm vac}\) es irracional, la palabra es Sturmiana:
tiene complejidad \(p(N)=N+1\), balance uno y entrop\'ia topol\'ogica cero.
Esta din\'amica difiere del desplazamiento completo sobre el alfabeto
tri\'adico: la vacancia determina una frontera ordenada de entrop\'ia nula.

La suma telesc\'opica produce
\begin{equation}
\sum_{n=1}^{N}z_n
=\lfloor(N+1)\varepsilon_{\rm vac}\rfloor,
\qquad
\sum_{n=1}^{N}(z_n-\varepsilon_{\rm vac})
=\varepsilon_{\rm vac}
-\{(N+1)\varepsilon_{\rm vac}\}.
\label{eq:vacancy-discrepancy}
\end{equation}
Las sumas parciales de la discrepancia tienen m\'odulo menor que uno. Esta
es la estimaci\'on que sostiene las constantes arm\'onicas posteriores.

La fracci\'on continua comienza
\[
\varepsilon_{\rm vac}
=[0;21,1,5,1,6,2,3,\ldots].
\]
Por ello los unos de \(z_n\) est\'an separados por \(21\) o \(22\) pasos.
Si
\[
\sigma=22-\varepsilon_{\rm vac}^{-1}
=1-T_G(\varepsilon_{\rm vac}),
\]
entonces \(a_1(\sigma)=6\) y los retornos de la separaci\'on corta son de
longitud \(6\) o \(7\). Las parejas \(21/22\) y \(6/7\) son dos escalas de
una misma rotaci\'on, no coincidencias de enteros.

Para la palabra exacta \cref{eq:vacancy-mechanical-word}, la constante regularizada es

\begin{equation}
\gamma_{\rm vac}^{+}
=\lim_{N\to\infty}\left(
\sum_{n\le N}\frac{z_n}{n}
-\varepsilon_{\rm vac}\log N
\right).
\label{eq:vacancy-gamma}
\end{equation}

La holonom\'ia asociada es

\begin{equation}
\rho_\varepsilon(k)=\exp(2\pi\ii k\varepsilon_{\rm vac}).
\label{eq:vacancy-holonomy}
\end{equation}

Las tres cantidades tienen tipos distintos: densidad, t\'ermino finito y fase. La torre de Stieltjes de vacancias prolonga \cref{eq:vacancy-gamma}:

\begin{equation}
\gamma_{\rm vac,j}^{+}
=\varepsilon_{\rm vac}\gamma_j
+\sum_{n\ge1}\frac{(z_n-\varepsilon_{\rm vac})(\log n)^j}{n}.
\label{eq:vacancy-stieltjes}
\end{equation}

Para \(j=0\) se recupera \(\gamma_{\rm vac}^{+}\). La convergencia requiere la cota de discrepancia del desarrollo integral de vacancias.

La convergencia se prueba aqu\'i. Para \(j\ge0\), la sucesi\'on
\(b_n=(\log n)^j/n\) tiende a cero y es decreciente a partir de
\(n>e^j\). La cota uniforme \cref{eq:vacancy-discrepancy} y el criterio de
Dirichlet prueban la convergencia de
\[
\sum_{n\ge1}(z_n-\varepsilon_{\rm vac})b_n.
\]
Para \(j=0\), separar
\(\varepsilon_{\rm vac}\sum_{n\le N}n^{-1}\) da exactamente
\cref{eq:vacancy-gamma}. Una sumaci\'on de Abel acota conservadoramente la
cola posterior a \(N\) por \(2/(N+1)\): un t\'ermino procede del borde y
otro de la variaci\'on total de \(1/n\). La ejecuci\'on exhaustiva declarada usa
\(N=10^8\), enumera las posiciones Beatty de los unos, comprueba los
conjuntos de huecos \(\{21,22\}\) y retornos \(\{6,7\}\), y obtiene como
parte finita
\[
-0.11207107505876366224207105242526397929\ldots.
\]
Al combinar la cota anal\'itica con las guardas de redondeo y frontera
registradas, el certificado finito encierra
\begin{equation}
-0.112071094143613634
<\gamma_{\rm vac}^{+}
<-0.112071055516338732.
\label{eq:vacancy-gamma-interval}
\end{equation}
Este intervalo es un valor terminal con error controlado. No define la
palabra; la palabra y su discrepancia lo preceden.

Una segunda evaluación mediante aritmética multiprecisión y control dirigido
del redondeo proporciona, a dieciocho decimales, el intervalo certificado
\[
-0.112071086058763562
<\gamma_{\rm vac}^{+}
<-0.112071064058763762,
\]
contenido estrictamente en \cref{eq:vacancy-gamma-interval}. El encaje de
ambos intervalos verifica la estabilidad de la parte finita respecto del
truncamiento, de las guardas de frontera y del régimen de precisión. El
resultado depende únicamente de las recurrencias, las cotas de cola y la
aritmética intervalar expuestas en este capítulo.

\section{Contenido operatorio de la representaci\'on helicoidal}

La representaci\'on helicoidal codifica una fase peri\'odica junto con un
incremento de memoria. Posee contenido operatorio cuando existe un cociclo

\[
c(\gamma\eta)=c(\gamma)+c(\eta)
\]

que transporta el acarreo y cuya monodrom\'ia se entrelaza con el operador
aritm\'etico. En ausencia de dicho mapa, la h\'elice constituye s\'olo una
parametrizaci\'on de fase; con el entrelazador, distingue el retorno de fase
del retorno de estado.

\begin{remark}[Control negativo]
La cofinalidad de cilindros no implica conjugaci\'on din\'amica. Un
contraejemplo a la conmutaci\'on entre el desplazamiento TPK y \(T_G\)
impide el entrelazador, pero no afecta a
\cref{eq:levy,eq:gauss-entropy}. En el sector de vacancias, la divergencia
de \cref{eq:vacancy-stieltjes} bajo la cota declarada refuta la familia de
constantes regularizadas.
\end{remark}

\subsection*{Alcance lógico y dominio de validez}
Cofinalidad, identidades de Gauss, L\'evy, entrop\'ia y Lochs: \emph{Teorema interno exacto en el sistema transportado}. GKW y entrelazador din\'amico: \emph{Programa abierto con criterio de refutación}. Densidad, constante y holonom\'ia de vacancia: \emph{Teorema interno exacto con control de discrepancia}.
 \section{Teorema rector de la red radical y la ley potencia--vacancia}

La aparición de \(369\) y \(693\) no se reduce a una semejanza gráfica. Hay dos mecanismos exactos que convergen:

\begin{enumerate}
\item \textbf{mecanismo algebraico nonádico:} \(3,6,9\) son los representantes visuales del radical nilpotente \(\mathfrak m=(3)\subset\mathbb Z/9\mathbb Z\); las palabras \(369\), \(693\) y \(936\) son las tres fases de su órbita multiplicativa;
\item \textbf{mecanismo arquimediano de vacancias:} la diferencia entre la longitud decimal de \(9^n\) y la frontera \(10^n\) está gobernada exactamente por la misma rotación irracional \(\delta_{\rm vac}=\log_{10}(10/9)\) que ya organiza el calendario de vacancias.
\end{enumerate}

Para \(n=9^9\), ambos mecanismos se encuentran en la identidad

\[
9^9-\left\lceil 9^9\log_{10}\frac{10}{9}\right\rceil
=369\,693\,099,
\]

de modo que

\[
\#\operatorname{cifras}\bigl(9^{9^9}\bigr)=369\,693\,100.
\]

El bloque \(369\mid693\) es, por tanto, una salida exacta del lector logarítmico de la potencia nonádica; no es el prefijo decimal de \(9^{9^9}\). La unidad final procede del operador universal que transforma orden decimal en longitud decimal.

Esta identidad abre una lectura estructural de mayor alcance: suma, producto, potencia, resta, división, raíz y logaritmo pueden organizarse como una jerarquía de transportes y cocientes. Cada ascenso de operación compactifica historias; el TPK conserva las coordenadas que permiten reconstruirlas.


\section{Estructura radical de la APP sobre \(\mathbb Z/9\mathbb Z\)}

\subsection{Anillo local nonádico \(\mathbb Z/9\mathbb Z\) y radical nilpotente}

Sea

\[
R=\mathbb Z/9\mathbb Z.
\]

Su grupo de unidades y su ideal maximal son

\[
R^\times=\{1,2,4,5,7,8\},
\qquad
\mathfrak m=(3)=\{0,3,6\}=\{9,3,6\}_{\rm vis}.
\]

Además,

\[
\mathfrak m^2=0.
\]

La igualdad anterior es una igualdad de ideales: todo producto de dos elementos de \(\mathfrak m\) se anula módulo nueve. No debe confundirse con el producto cartesiano \(\mathfrak m\times\mathfrak m\), que contiene nueve posiciones de la APP.

\subsection{Secuencia exacta del multiplicador \(3\)}

La multiplicación por tres,

\[
M_3:R\longrightarrow R,
\qquad x\longmapsto3x,
\]

tiene

\[
\ker M_3=\mathfrak m,
\qquad
\operatorname{im}M_3=\mathfrak m.
\]

Por ello induce un isomorfismo de espacios sobre \(\mathbb F_3\),

\[
\overline M_3:R/\mathfrak m\xrightarrow{\sim}\mathfrak m,
\qquad
\overline x\longmapsto3x.
\]

En representantes visuales,

\[
\overline0\longmapsto9,
\qquad
\overline1\longmapsto3,
\qquad
\overline2\longmapsto6.
\]

Así, la banda \(3,6,9\) no es un adorno decimal. Es una copia interna del cociente ternario alojada en el radical nonádico.

\subsection{Órbita cíclica \(369\to693\to936\to369\)}

La fila de la tabla multiplicativa correspondiente a \(3\) es

\[
3,6,9,3,6,9,3,6,9.
\]

El cambio de fase \(j\mapsto j+1\) produce

\[
369\longmapsto693\longmapsto936\longmapsto369.
\]

Estas tres palabras son las tres lecturas cíclicas de la aplicación

\[
\mathbb F_3\longrightarrow\mathfrak m,
\qquad
\bar j\longmapsto3j.
\]

La fila \(6\), al satisfacer \(6\equiv-3\pmod9\), invierte la orientación del mismo ciclo. La marca \(9\) cumple dos funciones tipadas:

\[
9\bmod9=0,
\qquad
\operatorname{dr}_9(9)=9.
\]

El módulo publica la clase nula; la raíz digital positiva conserva la marca de clausura. Esta diferencia de lectores es precisamente la que permite que \(9\) sea a la vez cero residual y cierre visible.

\subsection{Cruz radical de la APP y reducción coordenada a \(\mathbb F_3^2\)}

Sobre

\[
X_9=R^2,
\]

la unión

\[
\mathscr C_{369}
=(\mathfrak m\times R)\cup(R\times\mathfrak m)
\]

es la cruz formada por las filas y columnas \(3,6,9\). Contiene

\[
3\cdot9+3\cdot9-3\cdot3=45
\]

posiciones. Su intersección central \(\mathfrak m\times\mathfrak m\) contiene nueve posiciones. En el observable multiplicativo, la suma de la cruz es

\[
297=216+81=3(72+27).
\]

La reducción coordenada produce

\[
X_9/(\mathfrak m\times\mathfrak m)\cong\mathbb F_3^2
\]

como conjunto cociente tipado por las clases módulo tres. El lift inverso requiere conservar la coordenada dentro de cada fibra de tres elementos.

\subsection{Núcleo \(7227\), descomposición espectral y enlace con la cruz radical}

El bloque central de la APP es

\[
B_c=
\begin{pmatrix}
7&2\\
2&7
\end{pmatrix}
=7I+2R
=9P_++5P_-.
\]

Su lectura por filas produce \(72\mid27\); su lectura por diagonales produce \(77\mid22\). Se verifican

\[
72+27=77+22=99,
\]

\[
72-27=45=\det B_c,
\]

\[
77-22=55=54+1.
\]

La relación con la cruz radical no es nominal:

\[
3\cdot72=216,
\qquad
3\cdot27=81,
\qquad
3\cdot99=297.
\]

El bloque unitario central y la red no unitaria \(3,6,9\) quedan así enlazados por una identidad local–global exacta. La jerarquía de la diferencia suma–producto conserva sus tipos:

\[
4\quad\text{(salto espectral local)},\qquad
2\quad\text{(acoplamiento fuera de la diagonal)},
\]

\[
6\quad\text{(diferencia comprimida módulo tres)},\qquad
54\quad\text{(despliegue bidimensional)}.
\]


\section{Ley potencia--vacancia para las escalas \texorpdfstring{\(9^n\)}{9 elevado a n} y \texorpdfstring{\(10^n\)}{10 elevado a n}}

\subsection{Funciones de longitud decimal y vacancia acumulada}

Para \(n\ge1\), sea

\[
D_9(n)=\#\operatorname{cifras}(9^n)
=\left\lfloor n\log_{10}9\right\rfloor+1.
\]

La frontera decimal \(10^n\) posee \(n+1\) cifras. Definimos la vacancia acumulada de la potencia nonádica por

\[
V_9(n)=(n+1)-D_9(n).
\]

Introducimos los dos números complementarios

\[
\rho_{\rm vac}=\log_{10}9,
\qquad
\delta_{\rm vac}=1-\rho_{\rm vac}
=\log_{10}\frac{10}{9}.
\]

\subsection{Identidad exacta entre longitud decimal y rotación de vacancias}

\textbf{Teorema.} Para todo entero \(n\ge1\),

\[
\boxed{V_9(n)=\left\lceil n\delta_{\rm vac}\right\rceil},
\]

y, por tanto,

\[
\boxed{D_9(n)=n-\left\lceil n\delta_{\rm vac}\right\rceil+1}.
\]

\textbf{Demostración.} Como \(\delta_{\rm vac}\) es irracional, \(n\delta_{\rm vac}\notin\mathbb Z\). Entonces

\[
\begin{aligned}
D_9(n)
&=\lfloor n(1-\delta_{\rm vac})\rfloor+1\\
&=\lfloor n-n\delta_{\rm vac}\rfloor+1\\
&=n-\lceil n\delta_{\rm vac}\rceil+1.
\end{aligned}
\]

Al restar de \(n+1\) se obtiene la primera identidad. \(\square\)

\subsection{Codificación mecánica de las vacancias por las potencias nonádicas}

Los incrementos de la vacancia son

\[
Z_n=V_9(n+1)-V_9(n)
=\lceil(n+1)\delta_{\rm vac}\rceil
-\lceil n\delta_{\rm vac}\rceil
\in\{0,1\}.
\]

Ésta es exactamente la palabra mecánica de pendiente \(\delta_{\rm vac}\) utilizada por el calendario HMT de vacancias. Por tanto, dicho calendario admite una lectura adicional completamente exacta:

\begin{quote}
\(Z_n=1\) si y sólo si, al pasar de \(9^n\) a \(9^{n+1}\), se pierde un nuevo rango decimal respecto de la frontera \(10^n\to10^{n+1}\).
\end{quote}

La potencia no añade una coincidencia a la teoría de vacancias: proporciona su realización natural como contador de pérdida de rango decimal.

\subsection{Ley \(21/22\) como caso del teorema potencia--vacancia}

Puesto que

\[
\delta_{\rm vac}^{-1}
=21.85434532678283256\ldots,
\]

los índices en los que \(Z_n=1\) están separados por \(21\) o \(22\). Los primeros veintiún exponentes satisfacen

\[
D_9(n)=n,\qquad1\le n\le21,
\]

mientras que

\[
D_9(22)=21.
\]

En particular,

\[
9^9=387\,420\,489
\]

posee exactamente nueve cifras. Esta conservación inicial del rango no es una regla indefinida: el calendario de vacancias determina exactamente cuándo deja de cumplirse.

\subsection{Evaluación exacta en \(n=9^9\) y longitud de \(9^{9^9}\)}

Sea

\[
n_*=9^9=387\,420\,489.
\]

Entonces

\[
n_*\delta_{\rm vac}
=17\,727\,389.3684296412564569\ldots,
\]

y, por ello,

\[
\boxed{V_9(n_*)=17\,727\,390}.
\]

La coordenada conservada después de descontar las vacancias es

\[
n_*-V_9(n_*)
=387\,420\,489-17\,727\,390
=369\,693\,099.
\]

Por tanto,

\[
\boxed{\operatorname{ord}_{10}\bigl(9^{9^9}\bigr)=369\,693\,099},
\]

\[
\boxed{D_9(9^9)=369\,693\,100}.
\]

Se obtienen además las congruencias

\[
9^9\equiv0\pmod9,\qquad
V_9(9^9)\equiv0\pmod9,
\]

\[
369\,693\,099\equiv0\pmod9,\qquad
369\,693\,100\equiv1\pmod9.
\]

Las sumas de cifras correspondientes son

\[
s_{10}(9^9)=45,\qquad
s_{10}(V_9(9^9))=36,
\]

\[
s_{10}(369\,693\,099)=54,\qquad
s_{10}(369\,693\,100)=37.
\]

El retorno \(9\to1\) aparece aquí con tipos precisos:

\begin{enumerate}
\item \(9^9\) y la vacancia acumulada pertenecen a la clase nula módulo nueve;
\item su diferencia, el orden decimal, sigue en la clase nula y posee suma de cifras \(54\);
\item el operador universal \(D=\operatorname{ord}+1\) añade la unidad de clausura y publica la clase positiva \(1\).
\end{enumerate}

\subsection{Dominio de validez de la ley potencia--vacancia}

Queda demostrado:

\begin{itemize}
\item el teorema potencia–vacancia para todo \(n\);
\item la identidad con la palabra mecánica de vacancias;
\item la evaluación exacta en \(n=9^9\);
\item el orden y la longitud decimal de \(9^{9^9}\);
\item las congruencias y sumas de cifras anteriores.
\end{itemize}

No se afirma que \(9\) sea la única base con una propiedad análoga en todo dominio, que cada aparición externa de \(369\) proceda de esta torre ni que la torre sustituya la construcción de las nueve fases o el certificado de generación de \(\pi,e,\varphi\).


\section{Jerarquía operatoria, reversibilidad y memoria genealógica}

\subsection{Traslación, dilatación y exponenciación como niveles operatorios}

Las operaciones elementales pueden verse como acciones:

\[
T_a(x)=x+a,\qquad
M_a(x)=ax,\qquad
P_k(x)=x^k.
\]

Sus inversas parciales son, respectivamente, resta, división y extracción de raíces. El paso de una operación a la siguiente modifica la estructura de las fibras.

\subsection{Reversibilidad estratificada en \(\mathbb Z/9\mathbb Z\)}

\textbf{Proposición.} En el anillo nonádico:

\begin{enumerate}
\item toda traslación \(T_a\) es biyectiva;
\item \(M_a\) es biyectiva si y sólo si \(a\in R^\times\);
\item si \(a\in\mathfrak m\), la multiplicación posee fibras no triviales y su imagen queda dentro de \(\mathfrak m\);
\item para \(x\in\mathfrak m\), \(x^k=0\) para todo \(k\ge2\);
\item sobre \(R^\times\cong C_6\), el mapa \(P_k\) es un automorfismo si y sólo si \(\gcd(k,6)=1\).
\end{enumerate}

\textbf{Demostración.} La primera afirmación usa la estructura de grupo aditivo. La segunda es la caracterización de las unidades. La tercera sigue de \(aR\subseteq\mathfrak m\). La cuarta procede de \(\mathfrak m^2=0\). La quinta es la condición para que la multiplicación por \(k\) sea un automorfismo del grupo cíclico de orden seis. \(\square\)

El resultado produce una jerarquía exacta:

\[
\text{suma: reversibilidad global},
\]

\[
\text{producto: reversibilidad restringida a unidades},
\]

\[
\text{potencia: dinámica cíclica en unidades y colapso nilpotente en }\mathfrak m.
\]

Las raíces son funciones sólo cuando la potencia correspondiente es inyectiva en el dominio elegido; en el dominio total son correspondencias de varias ramas o relaciones vacías.

\subsection{Lectura del Topological Phase Kernel sobre la jerarquía operatoria}

Esta estratificación sugiere una definición funcional:

\begin{quote}
El TPK no añade una operación exterior a suma, producto y potencia. Conserva la hoja, el cociente, el acarreo, la orientación, la ruta y la frontera que cada evaluación ordinaria elimina al identificar varias historias con una misma salida.
\end{quote}

La suma y el producto son las dos primeras evaluaciones sobre la APP. La potencia es composición repetida del producto; la raíz selecciona o reconstruye ramas de esa composición. El registro distingue el valor publicado de la historia operatoria que lo produce.

\subsection{Descenso logarítmico entre niveles consecutivos de la jerarquía}

En un dominio positivo,

\[
\log(xy)=\log x+\log y,\qquad
\log(x^k)=k\log x,\qquad
\log(a^x)=x\log a.
\]

El logaritmo convierte producto en suma, potencia en producto y exponencial en dilatación. El defecto aparece cuando el lector posterior discretiza:

\[
\mathcal L_{10}(a^n)
=\lfloor n\log_{10}a\rfloor+1.
\]

La parte entera introduce el calendario de cruces de frontera. Para \(a=9\), ese calendario es precisamente la palabra de vacancias.

\subsection{Conjugación afín de las traslaciones y dilataciones}

Las traslaciones y dilataciones satisfacen

\[
M_aT_bM_a^{-1}=T_{ab}
\]

cuando \(a\) es invertible. Esta identidad expresa que el producto reorganiza las unidades de suma. Cuando \(a\) deja de ser invertible, la conjugación se rompe porque desaparecen ramas. El radical \(3,6,9\) localiza exactamente ese fallo.


\section{Álgebra de caminos: composición multiplicativa, acción aditiva y superposición}

En un álgebra de caminos, la suma agrega alternativas y el producto concatena trayectorias. Para pesos locales \(w(e)\),

\[
K(x,y)
=\sum_{\gamma:x\to y}
\prod_{e\in\gamma}w(e).
\]

Si se fija una rama logarítmica coherente,

\[
\prod_{e\in\gamma}w(e)
=\exp\!\left(\sum_{e\in\gamma}\log w(e)\right),
\]

y, con la acción discreta apropiada,

\[
K(x,y)=\sum_{\gamma:x\to y}e^{-S(\gamma)/\hbar}.
\]

La ecuación reúne las operaciones:

\begin{itemize}
\item el producto compone pasos sucesivos;
\item el logaritmo acumula la acción local;
\item la exponencial reconstruye el peso de la ruta;
\item la suma agrega todas las alternativas.
\end{itemize}

La APP es el soporte mínimo en el que suma y producto actúan sobre las mismas historias. El TPK añade los datos que impiden que la suma sobre caminos se convierta en una suma sobre valores sin genealogía.

En este marco, una cifra es la salida de un lector; un valor es la proyección arquimediana de una sección; un número HMT conserva la familia compatible de caminos que sostiene ese valor; una partícula conserva esa genealogía bajo una realización espectral persistente.


\section{Raíz interna de \(-1\), estructura exponencial y geometrías elíptica, parabólica e hiperbólica}

\subsection{Raíz interna de \(-1\) en la rama finita}

La matriz de incidencia de Paley proporciona un operador \(A_W\) con

\[
A_W^2=-I.
\]

Esto convierte \(\mathbb F_3^6\) en un espacio tridimensional sobre

\[
\mathbb F_9=\mathbb F_3[\iota]/(\iota^2+1).
\]

La raíz de \(-1\) no se adjunta desde fuera como una notación compleja. Emerge como endomorfismo de un soporte ternario ya construido.

\subsection{Realización continua del tipo cuadrático inducido}

En la rama real elíptica, el generador \(J_{\rm e}\) satisface

\[
J_{\rm e}^2=-I,
\]

y

\[
e^{tJ_{\rm e}}=\cos t\,I+\sin t\,J_{\rm e}.
\]

En particular,

\[
e^{\frac\pi2J_{\rm e}}=J_{\rm e},\qquad
e^{\pi J_{\rm e}}=-I,\qquad
e^{2\pi J_{\rm e}}=I.
\]

Así, \(J_{\rm e}\) es el cuarto de giro interno y la ecuación de Euler es su cierre a media vuelta.

\subsection{Soporte operatorio común de \(3\) lecturas geométricas}

Sea

\[
\mathbb D=\{z\in\mathbb C:|z|<1\}.
\]

Este mismo conjunto contiene el diámetro real \((-1,1)\), el interior complejo, la frontera trigonométrica \(S^1\) y la métrica hiperbólica de Poincaré

\[
ds_{\rm P}^2=\frac{4\ell^2|dz|^2}{(1-|z|^2)^2}.
\]

La aplicación

\[
R_{\pi/2}(z)=iz
\]

preserva \(|z|\) y \(|dz|\). Por consiguiente, es simultáneamente:

\begin{itemize}
\item multiplicación compleja por \(\sqrt{-1}\);
\item rotación trigonométrica de \(90^\circ\);
\item isometría euclídea del disco;
\item isometría hiperbólica de Poincaré;
\item automorfismo de su frontera circular.
\end{itemize}

La coincidencia es literal: el mismo mapa actúa sobre el mismo portador; cambia el lector métrico.

\subsection{Geometrías elíptica, parabólica e hiperbólica de la ecuación de Euler extendida}

Los generadores

\[
J_{+1}^2=-I,\qquad
J_0^2=0,\qquad
J_{-1}^2=I
\]

producen

\[
e^{tJ_{+1}}=\cos t\,I+\sin t\,J_{+1},
\]

\[
e^{tJ_0}=I+tJ_0,
\]

\[
e^{tJ_{-1}}=\cosh t\,I+\sinh t\,J_{-1}.
\]

La misma serie exponencial se separa en tres regímenes según el cuadrado del generador: elíptico, parabólico e hiperbólico. La especialización

\[
e^{2\log\varphi\,J_{-1}}=F_{\rm av}^{\,2}
\]

integra la autoescala de Fibonacci. Euler y Fibonacci aparecen como especializaciones de una familia exponencial cuadrática común.

\subsection{Portador común, operador común y familia exponencial cuadrática}

La fusión geométrica se formula en tres niveles:

\begin{enumerate}
\item \textbf{portador común:} el disco unitario contiene el diámetro real, el interior complejo y el borde circular;
\item \textbf{operador común:} \(z\mapsto iz\) es cuarto de giro, multiplicación compleja e isometría hiperbólica;
\item \textbf{familia común:} la ecuación exponencial trítica reúne las ramas elíptica, parabólica e hiperbólica.
\end{enumerate}

El resultado no identifica las métricas euclídea e hiperbólica. Afirma que ambas son realizaciones distintas sobre un mismo soporte y comparten el automorfismo originado por la raíz interna de \(-I\).


\section{Restricción de la rotación interior, holonomía de borde y torsión mínima}

La construcción contiene dos objetos positivos y diferenciados:

\begin{enumerate}
\item el generador interno \(J_{\rm e}\), con \(J_{\rm e}^2=-I\);
\item el invariante global
\end{enumerate}
   \[
   \Delta_4=
   \frac{\pi-e\log\pi}{270}
   \left(1+\frac{\pi}{729}\right),
   \]
   denominado torsión mínima y utilizado antes de la selección de especie.

La tesis autoral añade que la dinámica rotacional interna, al publicarse en el borde, queda absorbida como torsión mínima. La relación no debe escribirse como una igualdad desnuda \(J_{\rm e}=\Delta_4\): los objetos tienen tipos diferentes. Su forma natural requiere

\[
\text{rotación interior}
\xrightarrow{\operatorname{Res}_{\partial}}
\text{holonomía de borde}
\xrightarrow{\operatorname{Def}}
\text{defecto escalar}.
\]

Sea \(q_\partial\) el lector de frontera, \(\nabla_\partial\) la conexión inducida y \(e_\partial\) el marco discreto. La torsión de borde tiene tipo

\[
T_\partial=d_{\nabla_\partial}e_\partial.
\]

La afirmación autoral adquiere forma probatoria mediante un funcional

\[
\Lambda_\partial:
\operatorname{Tors}(\partial X)
\longrightarrow\mathbb R
\]

tal que

\[
\Lambda_\partial(T_\partial)=\Delta_4
\]

y un teorema que haga proceder \(T_\partial\) de la restricción de \(J_{\rm e}\) a través de los operadores de borde del TPK.

La raíz interna de \(-I\), el cuarto de giro, la ecuación de Euler extendida y la fórmula de \(\Delta_4\) quedan demostrados en sus dominios. El paso de la rotación interior a la torsión mínima de borde exige el morfismo \(\operatorname{Res}_{\partial}\) y el funcional \(\Lambda_{\partial}\) declarados; no se identifica por una igualdad escalar.


\section{Triplicidad generacional, conjugación y mezcla sobre el estado holonómico integral}

La frase de Isidor Isaac Rabi puede servir como referencia histórica, pero no debe gobernar el título ni constituir una sección autónoma. La cuestión científica es:

\begin{quote}
¿Qué mecanismo interno distingue tres órdenes estables de realización fermiónica y conserva esa distinción bajo color, conjugación y mezcla?
\end{quote}

La respuesta HMT no es «porque TRIT tiene tres valores». La estructura separa

\[
\operatorname{flav}(X)
=\bigl(\sigma_{ud}(X),g(X),\chi_{\rm fina}(X)\bigr),
\]

donde \(\sigma_{ud}\) distingue ramas, \(g\) distingue generación física y \(\chi_{\rm fina}\) conserva orientación y memoria dentro de una rama. En quarks, la órbita de color se añade después del sabor y antes del carácter.

La triplicidad se obtiene como consecuencia de la estructura enriquecida; sus efectos sobre color, conjugación y mezcla se desarrollan a continuación:

\begin{enumerate}
\item color, sabor, generación y masa son lectores distintos del mismo estado holonómico integral;
\item la conjugación partícula–antipartícula actúa sobre la orientación sin borrar la clase generacional;
\item CKM y PMNS transportan bases de estados ya tipados; no crean la triplicidad;
\item la clasificación neutrínica debe derivarse de representación, ocupación y mezcla, no de renombrar profundidades \(G0,\ldots,G4\).
\end{enumerate}

Que los neutrinos sean leptones de espín semientero y obediencia fermiónica es una clasificación física conocida. La aportación estructural HMT consiste en obtener internamente por qué el sector neutro ocupa esa representación, por qué admite tres realizaciones físicas y cómo se acopla al operador de mezcla sin leer sus ángulos externos.


\section{Memoria operatoria del TPK: conservación de residuo, cociente, acarreo, hoja, orientación, ruta y frontera bajo proyección}

\subsection{Pérdida de reversibilidad y levantamiento genealógico}

En la APP, distintas historias pueden producir el mismo residuo. El producto pliega más historias que la suma; la potencia pliega más que el producto; la raíz intenta reconstruir ramas ya identificadas. El TPK conserva el dato que hace posible invertir esa pérdida:

\[
\text{residuo},\quad
\text{cociente},\quad
\text{acarreo},\quad
\text{hoja},\quad
\text{orientación},\quad
\text{ruta},\quad
\text{frontera},\quad
\text{memoria}.
\]

La operación algebraica puede caracterizarse no sólo por su valor, sino por el tipo de genealogía que conserva o destruye.

\subsection{Proyección de historias compatibles sobre el continuo arquimediano}

La recta real aparece cuando una familia de cilindros compatibles se contrae bajo la evaluación arquimediana. El valor real es la sombra de la sección; no su genealogía completa. La proyección excepcional conserva otra parte como incidencia y simetría.

La suma sobre caminos expresa el mismo principio:

\[
\text{alternativas}\xrightarrow{\sum}\text{amplitud global},
\qquad
\text{etapas sucesivas}\xrightarrow{\prod}\text{peso de ruta}.
\]

El logaritmo convierte la composición multiplicativa en acción aditiva; la exponencial la devuelve a amplitud. La estructura discreta del continuo puede formularse como conservación genealógica de todas las rutas antes de su proyección sobre \(\mathbb R\).

\subsection{Distinción entre sección numérica y realización persistente de partícula}

Un número y una partícula comparten invariantes ontológicos mínimos: memoria, orientación, persistencia, clausura, firma y posición en una familia de prolongaciones.

El número publica una coordenada arquimediana. La partícula añade realización espectral, ocupación, propagación y estabilidad. El paso HMT \(\to\) MD cambia el lector aplicado a una misma genealogía enriquecida; no cambia arbitrariamente de objeto.


\section{Cierre conjunto de la red radical, la ley potencia--vacancia, la geometría de Euler y la jerarquía operatoria}

\subsection{Teorema de la red radical nonádica}

La red \(3,6,9\) de la APP es la representación visual del radical nilpotente \((3)\subset\mathbb Z/9\mathbb Z\). Sus fases \(369\), \(693\) y \(936\) realizan el isomorfismo \(R/(3)\cong(3)\) inducido por la multiplicación por tres. Las bandas localizan exactamente el fallo de invertibilidad del producto.

\subsection{Teorema potencia--vacancia}

Para todo \(n\ge1\), el número de rangos decimales perdidos por \(9^n\) respecto de \(10^n\) es \(\lceil n\log_{10}(10/9)\rceil\). Sus incrementos forman la palabra mecánica de vacancias. En \(n=9^9\), la coordenada retenida es \(369\,693\,099\) y la longitud resultante es \(369\,693\,100\).

\subsection{Teorema geométrico de la familia exponencial de Euler}

La raíz interna de \(-I\) proporciona un cuarto de giro. En el disco unitario, el mismo mapa \(z\mapsto iz\) es multiplicación compleja, rotación trigonométrica e isometría de Poincaré. La ecuación exponencial trítica prolonga este régimen a los regímenes parabólico e hiperbólico, incorporando en el último la autoescala de Fibonacci.

\subsection{Principio de conservación operatoria entre niveles}

La jerarquía suma–producto–potencia es una jerarquía de fibras y pérdida de reversibilidad. El TPK es el levantamiento que conserva las coordenadas necesarias para distinguir historias identificadas por los lectores ordinarios.


\section{Hipótesis de tipado y controles negativos}

\begin{enumerate}
\item No identificar \(369\,693\,100\) con una palabra TPK sin declarar el lector entre dominios.
\item No confundir el ideal \(\mathfrak m^2=0\) con \(\mathfrak m\times\mathfrak m\).
\item No confundir módulo nueve con raíz digital: \(9\mapsto0\) en el primero y \(9\mapsto9\) en la segunda.
\item No usar la triplicidad de TRIT como demostración suficiente de tres generaciones físicas.
\item No convertir \(J^2=-I\) en \(\Delta_4\) sin el mapa de borde y el funcional de torsión.
\item No presentar \(369\) como prefijo de \(9^{9^9}\): pertenece a su orden y a su longitud decimal.
\item No reducir las diez mil cifras a una prueba de magnitud: certifican truncamientos extensos de un proceso parametrizable cuya tesis es la prolongación indefinida.
\end{enumerate}

\section{Consecuencias conjuntas para la red radical, la memoria operatoria y la torsión de borde}

El hallazgo principal no es que una torre contenga visualmente las cifras \(3,6,9\). La potencia nonádica y el calendario de vacancias son dos presentaciones del mismo defecto logarítmico, mientras la APP localiza \(3,6,9\) como el radical donde la multiplicación pierde invertibilidad. La torre conecta ambos extremos:

\[
\text{radical nonádico}
\longrightarrow
\text{potencia}
\longrightarrow
\text{logaritmo}
\longrightarrow
\text{vacancia}
\longrightarrow
\text{orden decimal}
\longrightarrow
\text{clausura }9\to1.
\]

La ecuación de Euler extendida añade la segunda dirección: la raíz interna de \(-I\) convierte la estructura discreta en dinámica de rotación y reúne, sobre un soporte común, la recta real, el círculo trigonométrico y el disco hiperbólico. El paso ulterior a torsión mínima debe conservar el mapa de borde ya desarrollado en el corpus, no reemplazarlo por una identificación nominal.
  \endgroup

\chapter{Radión, fase angular y equivalencia entre las realizaciones cíclica y lineal}
\begingroup
\renewcommand{\chapter}[2][]{}%
% Integración científica del capítulo del radión. La numeración y los rótulos
% se adaptan a la monografía; los cuerpos matemáticos permanecen íntegros.
% Adaptador de compatibilidad del fuente teoremática íntegro del radión.
%
% El tratado rector ya carga tipografía, geometría, hipervínculos, teoremas,
% tablas, TikZ y tcolorbox.  Este archivo aporta únicamente los nombres y
% entornos exclusivos del fuente teoremática, sin importar su preámbulo autónomo ni
% alterar la jerarquía tipográfica general.

\makeatletter
\@ifundefined{color@RadNavy}{\definecolor{RadNavy}{HTML}{102A43}}{}
\@ifundefined{color@RadBlue}{\definecolor{RadBlue}{HTML}{1F5E7A}}{}
\@ifundefined{color@RadTeal}{\definecolor{RadTeal}{HTML}{168C8C}}{}
\@ifundefined{color@RadCopper}{\definecolor{RadCopper}{HTML}{B66A3C}}{}
\@ifundefined{color@RadGold}{\definecolor{RadGold}{HTML}{D4A73A}}{}
\@ifundefined{color@RadInk}{\definecolor{RadInk}{HTML}{1A202C}}{}
\@ifundefined{color@RadMist}{\definecolor{RadMist}{HTML}{EDF4F7}}{}
\@ifundefined{color@RadCream}{\definecolor{RadCream}{HTML}{FAF7F0}}{}
\@ifundefined{color@RadRose}{\definecolor{RadRose}{HTML}{8D3F55}}{}
\@ifundefined{color@RadGreen}{\definecolor{RadGreen}{HTML}{2F7D5D}}{}
\@ifundefined{color@RadGray}{\definecolor{RadGray}{HTML}{66788A}}{}

\providecommand{\TRIT}{\ensuremath{\mathrm{TRIT}}}
\providecommand{\HMT}{\ensuremath{\mathrm{HMT}}}
\providecommand{\Rstate}{\mathcal R_{12}}
\providecommand{\epsR}{\varepsilon_R}
\providecommand{\epsone}{\varepsilon_R^{(1)}}
\providecommand{\pih}{\pi}
\providecommand{\alphah}{\alpha}
\providecommand{\hret}{\hbar_{\mathrm{ret}}}
\providecommand{\hfull}{h_{\mathrm{ret}}}
\providecommand{\diff}{\mathop{}\!\mathrm d}
\providecommand{\e}{\mathrm e}
\providecommand{\R}{\mathbb R}
\providecommand{\C}{\mathbb C}
\providecommand{\F}{\mathbb F}
\providecommand{\Z}{\mathbb Z}
\providecommand{\dd}{\,\mathrm d}
\providecommand{\status}[1]{\textsf{\bfseries #1}}
\providecommand{\sourcepath}[1]{\nolinkurl{#1}}

\@ifundefined{typebox}{%
  \newtcolorbox{typebox}[1][]{enhanced,breakable,colback=white,
    colframe=RadBlue,boxrule=.7pt,arc=1.5mm,left=4mm,right=4mm,
    top=3mm,bottom=3mm,title={Recibo de tipos},fonttitle=\bfseries,
    coltitle=RadNavy,#1}%
}{}
\@ifundefined{narrativebox}{%
  \newtcolorbox{narrativebox}[1][]{enhanced,breakable,colback=white,
    colframe=RadGold,borderline west={2.4pt}{0pt}{RadGold},
    boxrule=.25pt,arc=0mm,left=5mm,right=4mm,top=3mm,bottom=3mm,#1}%
}{}

% Las once portadillas internas del fuente teoremática autónomo se convierten en
% divisiones científicas ordinarias dentro del capítulo 52.  Se evita así
% introducir partes autónomas, gigantismo tipográfico o páginas vacías.
\providecommand{\partopener}[3]{\section{#2}\noindent\emph{#3}\par\medskip}

\providecommand{\BigRadionEquation}{%
  \begin{equation*}
  \boxed{\displaystyle
  \frac{180^{\circ}}{\pih}
  =50^{\circ}+A^{\circ}-\frac{20}{81}\,\Delta_4^{\circ}+\epsR^{\circ}}
  \end{equation*}}
\makeatother
 
\chapter{El radión angular: empalme con acarreo de las secciones directa y torsional y cierre exacto de la unidad de fase, acción y memoria}
\label{chap:sucesor-102-radion}

\section{Composición APP--TRIT--TPK del estado de fase: generación de la unidad orientada con memoria}
% Fragmento público generado desde fuentes teoremáticas íntegros.
% Residencia: 52.1.
% Fuente: /Users/ruben/Documents/HMT2/EL_CIERRE_HOLOGRAFICO_DEL_INFINITO_SUCESOR_102_CAPITULOS_2026-08-28_EN_TRABAJO/incoming/radion_monografia_20260827/manuscrito/sections/01_resumen_y_guia.tex:1-128
\paragraph{Síntesis matemática del objeto y de su cadena causal}

Esta monografía reconstruye el radián desde la cadena generativa
\[
\APP\longrightarrow\TRIT\longrightarrow\TPK
\longrightarrow X_{\mathrm{enr}}
\longrightarrow\text{estructura discreta del continuo}
\longrightarrow\MD,
\]
sin tomar como entrada ni el valor metrológico de \(180/\pi\), ni el pequeño
residuo angular que se desea explicar. El resultado es una unidad enriquecida,
el \emph{radión}, definida como la unidad orientada que recorre un radián de
fase y barre exactamente una acción \(\hret\) en la elipse positiva asociada
al mismo estado.

El cierre rector es
\BigRadionEquation
con
\[
S_E=2\pih\left(\frac5{36}+\frac{25}{9}\alphah\right),
\qquad
\Delta_4=
\frac{\pih-e\log\pih}{270}
\left(1+\frac{\pih}{729}\right),
\]
\[
\rho_{\mathrm{tw}}=\frac{2\cdot90}{729}=\frac{20}{81},
\qquad
\epsone=\frac{20}{81}\Delta_4-(S_E-1).
\]
La división APP prueba \(1<S_E<2\), de modo que el cociente es exactamente
uno y el resto es \(D_E=S_E-1\). La igualdad
\[
1=S_E-\frac{20}{81}\Delta_4+\epsone
\]
es entonces una identidad interna. La carta angular posterior produce
\[
\epsR^\circ
=5.13830167585752820716851646226459474899085744\ldots
\times10^{-8}\,{}^\circ.
\]

El texto demuestra además que la elipse de acción
\[
Q=a_S\cos\theta,\qquad P=b_S\sin\theta,\qquad
a_Sb_S=2\hret
\]
satisface
\[
\operatorname{Area}(0,\theta)=\hret\theta.
\]
Por tanto, un radión barre \(\hret\) y la vuelta \(2\pi\) encierra
\(h_{\mathrm{ret}}=2\pi\hret\). La misma elipse tiene borde simpléctico
finito e interior de área infinita cuando se equipa con la métrica
Hilbert--Beltrami--Klein. Ésta es la forma exacta de la tesis física que guía
la obra: la superficie visible cierra la acción; la profundidad interior
conserva un desarrollo no agotable por la proyección plana.

La exposición mantiene tipadas cuatro realizaciones del mismo estado:

\begin{enumerate}[label=\textbf{\arabic*.}]
\item el círculo \(S^1\), que publica la fase;
\item la elipse positiva, que publica la acción y la anisotropía;
\item el interior de Klein, que mide la profundidad mediante razón doble;
\item la hélice nonádica, que conserva retorno, memoria y contracción.
\end{enumerate}

No son cuatro geometrías superpuestas sin criterio. Son cuatro lectores con
dominio, codominio y dato conservado diferentes, aplicados a un único estado
holonómico integral. El sector \(90/120\), la dodecafase, el acarreo del radión, la cola
de Lambert, el núcleo de Barbero y las torres globales se presentan asimismo
como objetos emparentados pero no intercambiables.

\begin{thesisbox}
El radión es la unidad orientada que barre \(\hret\) en la elipse de acción.
El círculo publica su fase, Klein mide su profundidad, la hélice conserva su
memoria y el acarreo coinductivo empalma exactamente sus secciones directa y
torsional.
\end{thesisbox}

\paragraph{Orden expositivo y dependencias de lectura}

El orden no es ornamental. Las Partes I y II producen el objeto y su cierre
desde APP--TRIT--TPK. Sólo después se introducen las realizaciones en lenguaje
de Hilbert, LC, Maxwell, Minkowski, Hodge, Lorentz, modularidad y órbitas. Cada
capítulo distingue tres niveles:

\begin{description}[style=nextline,leftmargin=4.2cm,labelwidth=3.9cm]
\item[Resultado interno.] Igualdad o construcción demostrada en el dominio
HMT con sus primitivas declaradas.
\item[Formalización reunida.] Mapa nuevo que articula resultados ya existentes
sin usar el blanco como selector.
\item[Interfaz física.] Traducción dimensional posterior que reconoce una
estructura convencional sin hacerla origen de la selección.
\end{description}

Los recuadros de \emph{estatuto y falsador} indican qué observación concreta
invalidaría cada promoción fuerte. Esta disciplina no debilita la tesis; la
hace auditable. El objetivo de la narración extensa es que ninguna palabra
como \emph{círculo}, \emph{elipse}, \emph{torsión}, \emph{carga} o
\emph{memoria} oculte un cambio de tipo.

\paragraph{Resultados rectores y alcance probatorio}

\begin{enumerate}[label=\textbf{C\arabic*.}]
\item Cierre independiente del valor final del radión por dos publicaciones del mismo estado y
resta con acarreo APP.
\item Derivación del coeficiente \(20/81\) como dos hojas del sector activo
\(90\), normalizadas por la fibra \(3^6=729\).
\item Igualador tipado entre la costura \(\Re s=1/2\), la cara APP de cien
marcas y la fase dodecafásica \(\theta_2=50^\circ\).
\item Teorema de área del radión: \(\operatorname{Area}(\theta)=\hret\theta\).
\item Estudio focal exacto de la elipse y de sus invariantes euclídeos,
hiperbólicos, simplécticos y modulares.
\item Teorema borde--interior: acción total finita \(h\) e interior de Klein
de área hiperbólica infinita.
\item Realización exacta mediante covarianza mínima y oscilador LC, seguida de
la publicación electromagnética \(90/120\).
\item Atlas de no-identificaciones para todos los objetos \(90/120\), incluida
la separación entre \(\epsR\), cola espectral y núcleo Barbero \(12:-1\).
\item Caracterización estructural de \(\pi\) en la carta del radión, distinta
del generador coinductivo primario.
\item Síntesis ontológica: dimensión como rango mínimo de memorias que una
proyección de fase ya no puede conservar.
\end{enumerate}
% Fuente: /Users/ruben/Documents/HMT2/EL_CIERRE_HOLOGRAFICO_DEL_INFINITO_SUCESOR_102_CAPITULOS_2026-08-28_EN_TRABAJO/incoming/radion_monografia_20260827/manuscrito/sections/01b_mapa_resultados.tex:1-133
\paragraph{Correspondencia entre resultados, tipos y residencias causales}

La correspondencia siguiente ordena las afirmaciones rectoras y sus demostraciones,
manteniendo visible la jerarquía que impide que
la riqueza física borre el cierre matemático que la sostiene.

% HMT_RESULT_BEGIN:RDN-01-GENEALOGIA:theorem
\begin{exactbox}[title={Genealogía APP--TRIT--TPK del radión}]
El radión nace de la composición efectiva
\(\APP\to\TRIT\to\TPK\to X_{\mathrm{enr}}\to\MD\); no se obtiene
redefiniendo retrospectivamente el radián convencional.
\end{exactbox}
% HMT_RESULT_END:RDN-01-GENEALOGIA:theorem

% HMT_RESULT_BEGIN:RDN-02-CIERRE:theorem
\begin{exactbox}[title={Cierre exacto e independencia prospectiva}]
El acarreo APP produce internamente
\[
1=S_E-\frac{20}{81}\Delta_4+\varepsilon_R^{(1)},
\qquad
\varepsilon_R^{(1)}=\frac{20}{81}\Delta_4-(S_E-1),
\]
y la carta angular transporta esta misma identidad a
\(180^\circ/\pi\) sin introducir el residuo objetivo.
\end{exactbox}
% HMT_RESULT_END:RDN-02-CIERRE:theorem

% HMT_RESULT_BEGIN:RDN-03-BISAGRA-50:theorem
\begin{exactbox}[title={Coordenada crítica de \(50^\circ\)}]
La coordenada crítica \(\Re s=1/2\), publicada en la cara APP de cien
marcas y en la rueda dodecafásica, selecciona la fase
\(\theta_2=50^\circ\). Es una igualdad de lectores tipados del mismo estado.
\end{exactbox}
% HMT_RESULT_END:RDN-03-BISAGRA-50:theorem

% HMT_RESULT_BEGIN:RDN-04-COEFICIENTE-2081:theorem
\begin{exactbox}[title={Densidad bidireccional nonádica \(20/81\)}]
El coeficiente torsional no se ajusta en la red \(1/81\): se construye como
\(\rho_{\mathrm{tw}}=2N_6/3^6=2\cdot90/729=20/81\), conservando hoja,
orientación y normalización de la fibra nonádica.
\end{exactbox}
% HMT_RESULT_END:RDN-04-COEFICIENTE-2081:theorem

% HMT_RESULT_BEGIN:RDN-05-AREA-ACCION:theorem
\begin{exactbox}[title={Ley areal del radión}]
Para la elipse positiva
\(Q=a_S\cos\theta\), \(P=b_S\sin\theta\),
\(a_Sb_S=2\hbar_{\mathrm{ret}}\), el área orientada barrida es
\(\mathcal A(\theta)=\hbar_{\mathrm{ret}}\theta\). Un radión barre exactamente
\(\hbar_{\mathrm{ret}}\), y la vuelta \(2\pi\) encierra
\(h_{\mathrm{ret}}=2\pi\hbar_{\mathrm{ret}}\).
\end{exactbox}
% HMT_RESULT_END:RDN-05-AREA-ACCION:theorem

% HMT_RESULT_BEGIN:RDN-06-INVARIANTES-FOCALES:theorem
\begin{exactbox}[title={Invariantes focales del operador positivo}]
El par publicado \(A\pm C^\ast\) fija los semiejes, la excentricidad, la
separación focal y el parámetro de compresión simpléctica. En particular,
\((d_1/d_0)^2=C^\ast/2\), identidad adimensional que separa forma de escala.
\end{exactbox}
% HMT_RESULT_END:RDN-06-INVARIANTES-FOCALES:theorem

% HMT_RESULT_BEGIN:RDN-07-BORDE-INTERIOR:theorem
\begin{exactbox}[title={Finitud simpléctica del borde y divergencia del área hiperbólica interior}]
El borde simpléctico de la elipse encierra la acción finita
\(h_{\mathrm{ret}}\), mientras el área Hilbert--Beltrami--Klein del interior
diverge al alcanzar la frontera. Fase visible y profundidad no son la misma
medida.
\end{exactbox}
% HMT_RESULT_END:RDN-07-BORDE-INTERIOR:theorem

% HMT_RESULT_BEGIN:RDN-08-HELICE-MEMORIA:theorem
\begin{exactbox}[title={Elevación helicoidal nonádica con retorno de fase y avance de memoria}]
La monodromía nonádica devuelve la fase sin reiniciar el estado. La hélice
levanta el círculo conservando acarreo y memoria; el cierre \(2\pi/4\pi\)
distingue retorno proyectivo, inversión de orientación y retorno completo.
\end{exactbox}
% HMT_RESULT_END:RDN-08-HELICE-MEMORIA:theorem

% HMT_RESULT_BEGIN:RDN-09-OSCILADOR-EM:development
\begin{narrativebox}[title={Realización oscilatoria y constitutiva electromagnética}]
La misma anisotropía se realiza mediante una covarianza mínima y un oscilador
LC: una coordenada almacena el aspecto eléctrico-elástico y su conjugada el
aspecto magnético-rotacional. La frecuencia y la impedancia quedan tipadas por
separado.
\end{narrativebox}
% HMT_RESULT_END:RDN-09-OSCILADOR-EM:development

% HMT_RESULT_BEGIN:RDN-10-TIPOS-90120:theorem
\begin{exactbox}[title={Separación tipada de los operadores \(90/120\)}]
Extractor dodecafásico, acarreo del radión, cola de Lambert, funcional
hexada--octada, núcleo de Barbero y semigrupo global de torres son objetos
distintos. El peso \(12:-1\) queda forzado sólo al fijar el polo \(1/24\) y
el contratérmino entero mínimo.
\end{exactbox}
% HMT_RESULT_END:RDN-10-TIPOS-90120:theorem

% HMT_RESULT_BEGIN:RDN-11-GEOMETRIAS-CAMPO:development
\begin{narrativebox}[title={Realizaciones compleja, lorentziana y dual de Hodge}]
Los regímenes elíptico, parabólico e hiperbólico del TRIT se publican como
rotación, umbral y transformación hiperbólica. La estrella de Hodge, las constitutivas y la fuerza
de Lorentz actúan después sobre el estado ya orientado; no seleccionan sus
constantes.
\end{narrativebox}
% HMT_RESULT_END:RDN-11-GEOMETRIAS-CAMPO:development

% HMT_RESULT_BEGIN:RDN-12-MODULAR-ORBITAL:development
\begin{narrativebox}[title={Caracteres modular, orbital y aritmético de frontera}]
La involución modular, las órbitas elípticas, el avance-retardo y la aritmética
3--6--9 reconocen distintas memorias del mismo estado. La función modular y
las vacancias se usan como lectores posteriores, nunca como sustitutos del
generador APP--TRIT--TPK.
\end{narrativebox}
% HMT_RESULT_END:RDN-12-MODULAR-ORBITAL:development

% HMT_RESULT_BEGIN:RDN-13-SINTESIS-ONTOLOGICA:conclusion
\begin{thesisbox}[title={Síntesis operatoria de fase, acción, profundidad y memoria}]
El círculo es la proyección de fase; la elipse es la sección de acción; Klein
mide la profundidad; la hélice conserva la historia. El radión es la unidad
orientada que mantiene acopladas esas publicaciones sin borrar el residuo que
las distingue.
\end{thesisbox}
% HMT_RESULT_END:RDN-13-SINTESIS-ONTOLOGICA:conclusion

% HMT_RESULT_BEGIN:RDN-14-CONTROL-ALCANCE:control
\begin{statusbox}[title={Delimitación de dominios y continuidad causal}]
Las leyes completas de masas, el refinamiento electrónico, CKM/CP y el
corredor excepcional conservan sus demostraciones rectoras propias. Esta monografía los
cita sólo donde comparten antecedente o carta; no traslada al radión sus
selectores específicos ni confunde sus residencias probatorias.
\end{statusbox}
% HMT_RESULT_END:RDN-14-CONTROL-ALCANCE:control
 
\subsection{Soporte bivariado de APP, orientación ternaria de TRIT y actualización del estado holonómico integral mediante TPK}
% Fragmento público generado desde fuentes teoremáticas íntegros.
% Residencia: 52.1.1.
% Fuente: /Users/ruben/Documents/HMT2/EL_CIERRE_HOLOGRAFICO_DEL_INFINITO_SUCESOR_102_CAPITULOS_2026-08-28_EN_TRABAJO/incoming/radion_monografia_20260827/manuscrito/sections/02_genealogia.tex:1-180
\noindent La genealogía operatoria antecede a la medida angular y determina las variables conservadas por cada lector.\par\medskip

\subsubsection{Soporte APP--TRIT--TPK y construcción del estado holonómico integral}

\paragraph{Antecedencia del estado holonómico integral respecto de la publicación circular}

La definición escolar de radián comienza con un círculo ya dado, un centro ya
dado y una longitud de arco ya mensurable. Si el radio y el arco tienen la
misma longitud, el ángulo subtendido mide un radián. La definición es correcta
dentro de la geometría euclídea, pero no responde a una pregunta anterior:
¿qué estructura hace posible que una vuelta posea fase, orientación, acción y
memoria, y por qué la unidad natural se relaciona con \(2\pi\)?

HMT invierte el orden causal. Primero construye el estado finito, después su
prolongación coinductiva, luego la vuelta y sólo al final su publicación
angular. El radión no corrige el radián convencional; lo levanta a su dominio
genealógico.

\begin{typebox}
\[
\APP\to\TRIT\to\TPK\to X_{\mathrm{enr}}
\to\mathcal C^{\mathrm{disc}}_{\mathrm{cont}}
\to\bigl(\pih,e,\varphi,\alphah,A,C^*,\Delta_4,\hret\bigr)
\to\mathfrak r_\sigma.
\]
La matemática convencional sólo aparece después como lenguaje de realización:
\(S^1\), elipse simpléctica, métrica de Klein, espacio de Hilbert, circuito
LC, pantalla de Minkowski y ecuaciones de campo.
\end{typebox}

\paragraph{Soporte bivariado de APP: hojas, residuos, cocientes y acarreo}

La Plataforma de Aritmética Posicional (APP) conserva simultáneamente dos
realizaciones de los símbolos \(1,\dots,9\): una hoja aditiva y una hoja
multiplicativa. No se reduce un número a su valor publicado. Junto al valor
permanecen hoja, orientación, cociente, resto, acarreo, frontera, ruta y
supervivencia.

El censo básico tiene tres estratos que jamás deben identificarse por mera
coincidencia cardinal:
\[
104,976\longrightarrow468\;U_6\longrightarrow243\;w_6
\subset\F_3^6,\qquad |\F_3^6|=729.
\]
Las \(468\) emisiones decimales, las \(243\) palabras visibles y el ambiente
de \(729\) palabras son objetos distintos. El ambiente será decisivo para la
normalización \(20/81\).

En base cúbica \(B=1000\), la resta orientada de dos palabras de tríadas se
realiza de derecha a izquierda:
\begin{equation}
u_i=t_i-v_i+c_{i+1},\qquad
r_i=u_i\bmod1000,\qquad
c_i=\frac{u_i-r_i}{1000}.
\label{eq:subacarreo}
\end{equation}
Fijados los operandos y el acarreo remoto, cociente y resto son únicos. Éste es
el operador que, más adelante, publicará las cifras de \(\epsone\). No hay
ningún decimal elegido después de ver el radián.

La completación
\[
10^3=729+243+27+1
\]
no es una curiosidad decimal. Expresa la conservación de la unidad al pasar
del ambiente ternario a una cámara de mil posiciones. La diferencia
\(1000-729=271\) es la corona completante; la identidad correcta es
\(999=729+270\), no \(729+271\).

\paragraph{Orientación ternaria y clasificación de los 3 regímenes operatorios del TRIT}

El TRIT asigna a cada estado local uno de tres regímenes orientados. En la
convención temporal activa,
\[
\begin{aligned}
+1&\longleftrightarrow\text{retorno, memoria, pasado},\\
0&\longleftrightarrow\text{umbral, presente},\\
-1&\longleftrightarrow\text{apertura bidireccional, futuro}.
\end{aligned}
\]
La realización cuadrática distingue:
\begin{equation}
J^2=-I,\qquad N^2=0,\qquad R^2=I.
\label{eq:tres-regimenes}
\end{equation}
Así aparecen, respectivamente, la rotación elíptica, el umbral parabólico y
la separación hiperbólica. No se mezclan tres métricas. Se conservan tres
acciones tipadas sobre un antecedente común.

La raíz interna de \(-1\) no se importa suponiendo de antemano el plano
complejo. En el bloque APP, dos proyecciones \(P,Q\) con razón
\(r=3/4=90/120\) producen
\[
J_{\mathrm{comm}}=\frac{[P,Q]}{\sqrt{r(1-r)}},\qquad J_{\mathrm{comm}}^2=-I,
\]
mientras
\[
R_{\mathrm{comm}}=2P-I,\qquad R_{\mathrm{comm}}^2=I,\qquad
R_{\mathrm{comm}}J_{\mathrm{comm}}=-J_{\mathrm{comm}}R_{\mathrm{comm}}.
\]
La pareja genera la forma mínima de \(\mathrm{Cl}_{1,1}\simeq M_2(\R)\):
rotación y apertura emergen juntas, pero siguen siendo operaciones diferentes.

\paragraph{Composición dinámica del TPK: selección, transporte, memoria y retorno}

El Topological Phase Núcleo no es un nombre para una caja de objetos. Es la
composición que selecciona fase, transporta cardinalidad, levanta hojas,
registra memoria y ejecuta el retorno. La dependencia fundamental es
\[
U_t\longrightarrow\Gamma_9\longrightarrow K_{\mathrm{ph}}.
\]
La holonomía nonádica \(\Gamma_9\) actúa sobre el estado holonómico integral; la
publicación \(K_{\mathrm{ph}}\) es posterior y conserva memoria reducida. La
identidad \(K_{\mathrm{ph}}^9|_r=E_+E_\times\) no convierte a \(K_{\mathrm{ph}}\) en
generador.

El certificado nonádico conserva la cadena completa
\[
w_6\to w_{12}\to w_{18}\to w_{24}\to w_{30}
\to R_{36}\to G_9,
\]
con \(396\) niveles, \(2376\) trits, \(43\) vueltas, \(387\) pares
\((K,K+9)\) por canal y cinco regiones de \(\pi\). La fase cumple
\[
g(k+9)=g(k),
\]
pero el estado completo no retorna:
\[
B_{k+9}\ne B_k.
\]
Ésta es la semilla matemática de la hélice: la proyección angular cierra; la
memoria avanza.

\paragraph{Extensión temporal no escindida \(C_{12}\to C_{108}\to C_9\)}

La memoria de doce pasos se organiza mediante
\begin{equation}
0\longrightarrow C_{12}\longrightarrow C_{108}
\longrightarrow C_9\longrightarrow0.
\label{eq:extension-108}
\end{equation}
El cociclo
\[
c(r,r')=\left[\left\lfloor\frac{r+r'}9\right\rfloor\right]_{12}
\]
mide el acarreo entre retorno de fase y avance de memoria. Nueve pasos cierran
la fase observable \(C_9\); no reinician el registro dodecafásico.

Conviene separar cuatro objetos:
\begin{center}
\begin{tabularx}{.95\textwidth}{>{\bfseries}l X}
\toprule
\(C_{12}\) & grupo cíclico de doce posiciones;\\
\(\Rstate\) & registro enriquecido \((E_m,\tau_m,\sigma_m,\kappa_m,\Lambda_m)_{m=1}^{12}\);\\
\(\Theta_{108}\) & publicación cruda en el calendario de 108;\\
\(\Gamma_{108}\) & historia u holonomía levantada.\\
\bottomrule
\end{tabularx}
\end{center}
Por ello, la expresión vaga \enquote{post-\(C_{12}\)} se sustituirá por
\enquote{posterior al registro dodecafásico \(\Rstate\), en su lectura
sexádica} cuando ése sea el mapa efectivo.

\paragraph{Emergencia correlativa de las 5 construcciones del continuo desde el estado único}

La estructura discreta del continuo no se obtiene ensamblando cinco teorías
independientes. Sus cinco aspectos internos ---espectral, solenoidal, calibre,
cohomológico y determinantal--- emergen correlativamente del mismo estado
APP--TRIT--TPK. Esta obra no vuelve a demostrar el tratado integral, pero
preserva el hecho causal que necesita el radión: la sección coinductiva que
genera \(\pi,e,\varphi\) y \(\alpha\) no vive en una recta real desnuda; vive
en un objeto con fibras, relaciones, orientación, acarreo, memoria, terminal y
lector.

\begin{statusbox}
\textbf{Falsador de genealogía.} Si una fórmula posterior recibe
\(180/\pi\), \(\epsR\), una masa o una tabla metrológica para elegir una fase,
un coeficiente o una rama, deja de ser una generación HMT-first. El contraste
externo puede reconocer una salida; no puede seleccionarla.
\end{statusbox}
 
\subsection{Independencia prospectiva del cierre respecto del valor final}
% Fragmento público generado desde fuentes teoremáticas íntegros.
% Residencia: 52.1.2.
% Fuente: /Users/ruben/Documents/HMT2/EL_CIERRE_HOLOGRAFICO_DEL_INFINITO_SUCESOR_102_CAPITULOS_2026-08-28_EN_TRABAJO/incoming/radion_monografia_20260827/manuscrito/sections/02_genealogia.tex:181-272

\subsubsection{Genealogía de las constantes y construcción del par angular}

\paragraph{Sección coinductiva y unicidad del valor generado de \(\pi\)}

Sea \(x_\infty\) la sección global del sistema inverso. Los lectores
coinductivos producen
\[
\operatorname{ev}_\pi(x_\infty)=\pih,\qquad
\operatorname{ev}_e(x_\infty)=e,\qquad
\operatorname{ev}_\varphi(x_\infty)=\varphi.
\]
Estas evaluaciones determinan las constantes universales mismas, no variantes
dependientes del procedimiento que las construye. En particular,
\[
\pih=3.14159265358979323846264338327950288419716939937510\ldots.
\]
La vuelta positiva es
\begin{equation}
T_+(x_\infty)=2\pih.
\label{eq:vuelta-plus}
\end{equation}
Su publicación angular de una unidad es
\begin{equation}
R_{\HMT}^{\circ}=\frac{360^\circ}{T_+}=\frac{180^\circ}{\pih}.
\label{eq:radian-publicado}
\end{equation}
La ecuación no define todavía el radión enriquecido; define la carta visible
del radián una vez generada la vuelta.

\paragraph{Bidireccionalidad de rutas y prolongaciones conjugadas}

La bidireccionalidad no se enuncia como eslogan. La involución concreta
\begin{equation}
J_\alpha(T)=\frac{\alphah^2}{T},\qquad
T_-=J_\alpha(T_+),\qquad T_+T_-=\alphah^2
\label{eq:bidireccional-turn}
\end{equation}
conserva un producto y revierte la hoja. Expansión y contracción son dos
prolongaciones conjugadas del mismo estado. La raíz finita de un jet
cuadrático puede aproximar \(T_+\), pero no sustituye a la sección
coinductiva: la diferencia certificada es
\[
T_{P_9}-T_+=-5.1110566290335\ldots\times10^{-25}.
\]
El control finito es un testigo; el fundamento es el límite.

\paragraph{Dependencia causal entre \(\alpha\), acción y publicación angular}

El cierre dodecafásico genera \(\alphah\) desde el estado firmado y sus
cartas reversibles. Después, el lector determinantal local fija la década de
acción mediante
\[
\Sigma_H=\varphi\alphah^{16},
\qquad \operatorname{ord}_{10}\Sigma_H=-34.
\]
La sección de retorno produce \(\hret\). Ni \(\alpha\) ni \(\hbar\) se
\enquote{miden en grados}. Se publican posteriormente mediante el par angular
\begin{equation}
P(\alpha,\eta_{\mathrm{ret}})=
\left([10^3\alpha]_{360},[2(\eta_{\mathrm{ret}}+\alpha)]_{360}\right)
=(A,C^*).
\label{eq:parangular}
\end{equation}
En la rama canónica,
\[
A=7.297352569283800997285105472380663\ldots{}^\circ,
\]
\[
C^*=2.1237383389686457242619513803933607937\ldots{}^\circ.
\]
Sólo entonces la carta analítica
\[
x=\frac{\pi A}{180},\qquad y=\frac{\pi C^*}{180}
\]
abre los canales \(q_\pm=e^{-x\mp y}\).

El orden causal completo es
\[
\alpha\to(\eta\parallel-34)\to\hbar
\to P\to\kappa_{360}\to q_\pm
\to\text{realizaciones MD}.
\]
La relación \(C^*/2-\alpha\) puede servir de verificación inversa. No se usa
para generar \(\hret\) retroactivamente.

\begin{narrativebox}
La carta angular no convierte una constante adimensional o una acción en un
ángulo por decreto. Publica dos coordenadas del estado en una rueda finita.
El grado es el idioma de esa publicación; no es la sustancia ontológica de
\(\alpha\) ni de \(\hbar\).
\end{narrativebox}
% Fuente: /Users/ruben/Documents/HMT2/EL_CIERRE_HOLOGRAFICO_DEL_INFINITO_SUCESOR_102_CAPITULOS_2026-08-28_EN_TRABAJO/incoming/radion_monografia_20260827/manuscrito/sections/03_cierre_exacto.tex:254-285
\paragraph{Independencia prospectiva respecto del valor final}

El verificador del operador de profundidad prohíbe expresamente tres entradas:
\[
\frac{180}{\pi},\qquad \epsR,\qquad
\text{el decimal objetivo}.
\]
Primero construye \(S_E\) y \(S_T\), después ejecuta la resta APP, y sólo al
final compara el cierre. Los controles dan
\[
\texttt{target\_epsilon\_used=False},\qquad
\texttt{target\_180\_over\_pi\_used=False}.
\]
Asimismo:
\begin{itemize}
\item una sola hoja da un defecto de orden \(10^{-5}\);
\item sustituir \(\theta_2\) por \(\theta_1\) o \(\theta_3\) desplaza el
resultado en aproximadamente \(\pm\pi/6\);
\item eliminar el canal \(e\) deja un defecto de orden \(10^{-3}\);
\item la variante histórica \(\Delta_4(A,C^*)\) produce otro número;
\item la cola espectral posterior a \(\Rstate\) y el operador armónico de
cuatro términos tampoco coinciden con \(\epsone\).
\end{itemize}

\begin{statusbox}
\textbf{Estatuto.} El cierre es una \status{formalización nueva} demostrada
con estructura de partida explícita y certificado computacional independiente del valor final.
La igualdad es exacta en el límite coinductivo. El perfil corto de 36 cifras
de \(\alpha\) genera una variante que diverge sólo después de muchas cifras;
se conserva como testigo histórico, no como valor canónico.
\end{statusbox}

 
\section{Secciones directa y torsional: acarreo APP y cierre exacto del radión}
\subsection{Construcción independiente de las secciones \(S_E\) y \(S_T\) antes de definir el cociclo \(\varepsilon_R=S_T-S_E\)}
% Fragmento público generado desde fuentes teoremáticas íntegros.
% Residencia: 52.2.1.
% Fuente: /Users/ruben/Documents/HMT2/EL_CIERRE_HOLOGRAFICO_DEL_INFINITO_SUCESOR_102_CAPITULOS_2026-08-28_EN_TRABAJO/incoming/radion_monografia_20260827/manuscrito/sections/03_cierre_exacto.tex:1-50
\noindent El cierre exacto se obtiene mediante dos secciones construidas antes del cociclo que las relaciona.\par\medskip

\subsubsection{Construcción independiente de las secciones directa y torsional}

\paragraph{Sección directa del sistema temporal dodecafásico}

La rueda dodecafásica tiene centros
\begin{equation}
\theta_m=30m-10^\circ,\qquad m=1,\ldots,12.
\label{eq:rueda-dodecafasica}
\end{equation}
La segunda fase es \(\theta_2=50^\circ\). Junto a la publicación
\(A=10^3\alphah\), forma la sección directa
\begin{align}
S_E
&=\frac{T_+}{360}\,(50+10^3\alphah)\\
&=2\pih\left(\frac5{36}+\frac{25}{9}\alphah\right).
\label{eq:seccion-electrica}
\end{align}
El subíndice \(E\) significa aquí \emph{directa o común}; su realización
eléctrica llegará después del lector constitutivo. No se postula
literalmente \(A=E\).

Los cilindros coinductivos prueban
\begin{equation}
1<S_E<2.
\label{eq:SE-intervalo}
\end{equation}
La división APP fuerza entonces
\begin{equation}
q_E=\lfloor S_E\rfloor=1,
\qquad
D_E=\operatorname{rem}_1(S_E)=S_E-1.
\label{eq:DE}
\end{equation}
El entero uno no es el número de vueltas de la hélice. Es el cociente único
de la sección directa dentro de la celda de unidad.

\paragraph{Sección torsional determinada por el operador \(\Delta_4\)}

La segunda sección utiliza una salida coinductiva anterior:
\begin{equation}
\Delta_4=
\frac{\pih-e\log\pih}{270}
\left(1+\frac{\pih}{729}\right).
\label{eq:delta4}
\end{equation}
La palabra \emph{torsión} designa aquí el cuanto global de la carta HMT. No
lo identifica automáticamente con la torsión de Cartan
\(T^I=D_\omega e^I\). El puente hacia Cartan exige otra flecha tipada.
 \subsection{Densidad orientada del sector activo y determinación única de \(20/81\)}
% Fragmento público generado desde fuentes teoremáticas íntegros.
% Residencia: 52.2.2.
% Fuente: /Users/ruben/Documents/HMT2/EL_CIERRE_HOLOGRAFICO_DEL_INFINITO_SUCESOR_102_CAPITULOS_2026-08-28_EN_TRABAJO/incoming/radion_monografia_20260827/manuscrito/sections/03_cierre_exacto.tex:51-103

\paragraph{Densidad orientada \(20/81\) del sector activo \(90/120\)}

El selector trítico de fase separa
\[
H_+=\{1,4,7\},\qquad H_-=\{2,5,8\},\qquad H_0=\{3,6,9\}.
\]
Hay seis fases activas y quince parejas no ordenadas:
\[
H_{\mathrm{act}}=H_+\cup H_-,\qquad
\left|\binom{H_{\mathrm{act}}}2\right|=\binom62=15.
\]
El modo activo y su extensión a ocho posiciones son
\begin{equation}
N_6=6\binom62=90,\qquad
N_8=8\binom62=120.
\label{eq:90-120-incidencia}
\end{equation}
Dos hojas conjugadas del sector activo, normalizadas por la fibra ambiente,
dan
\begin{equation}
\rho_{\mathrm{tw}}
=\frac{2N_6}{|\F_3^6|}
=\frac{2\cdot90}{729}
=\frac{180}{729}
=\boxed{\frac{20}{81}}.
\label{eq:rho-twoway}
\end{equation}

\begin{proposition}[Unicidad relativa de la densidad]
Fijados el sector activo \(N_6=90\), las dos hojas conjugadas y la
normalización por \(|\F_3^6|=729\), la densidad orientada es necesariamente
\(20/81\).
\end{proposition}
\begin{proof}
La cardinalidad transportada es \(2N_6=180\). La normalización por el ambiente
da \(180/729\); dividir numerador y denominador por \(9\) produce \(20/81\).
No interviene ninguna comparación con \(180/\pi\) ni con \(\epsR\).
\end{proof}

La lectura histórica como proyección sobre una red \(1/81\) es un control de
rigidez útil, pero ya no es el fundamento: el coeficiente se obtiene antes
del cierre. La ablación de una hoja produce \(10/81\), y la salida cambia en
orden \(10^{-5}\); por tanto, el factor dos no es decorativo.

La sección torsional es
\begin{equation}
\Phi_T=\rho_{\mathrm{tw}}\Delta_4,
\qquad
S_T=1+\Phi_T.
\label{eq:seccion-torsional}
\end{equation}

 \subsection{El acarreo \(\varepsilon_R\) como cociclo de transición entre secciones}
% Fragmento público generado desde fuentes teoremáticas íntegros.
% Residencia: 52.2.3.
% Fuente: /Users/ruben/Documents/HMT2/EL_CIERRE_HOLOGRAFICO_DEL_INFINITO_SUCESOR_102_CAPITULOS_2026-08-28_EN_TRABAJO/incoming/radion_monografia_20260827/manuscrito/sections/03_cierre_exacto.tex:104-253
\subsubsection{Cociclo de transición entre las secciones directa y torsional}

\paragraph{Definición operatoria del acarreo \(\varepsilon_R\)}

\begin{definition}[Acarreo real del radión]
La salida unitaria del radión es el cociclo de transición
\begin{equation}
\epsone:=S_T-S_E
=\Phi_T-D_E
=\frac{20}{81}\Delta_4-operatorname{rem}_1(S_E).
\label{eq:epsilon-def}
\end{equation}
\end{definition}

La ecuación contiene una resta, pero no es un ajuste al blanco. Sus dos
operandos \(S_T\) y \(S_E\) se generan separadamente. La independencia que
se demuestra es independencia respecto de \(180/\pi\) y de un
\(\varepsilon\) objetivo; no se afirma independencia algebraica entre una
diferencia y sus sumandos.

\begin{theorem}[Cierre exacto unitario]
Con \(S_E\), \(\Phi_T\) y \(\epsone\) definidos por
\eqref{eq:seccion-electrica}, \eqref{eq:delta4}, \eqref{eq:rho-twoway} y
\eqref{eq:epsilon-def}, se cumple exactamente
\begin{equation}
\boxed{1=S_E-\Phi_T+\epsone.}
\label{eq:cierre-unitario}
\end{equation}
\end{theorem}
\begin{proof}
Por \eqref{eq:SE-intervalo}, \(D_E=S_E-1\). Por definición,
\(\epsone=\Phi_T-D_E\). Entonces
\[
S_E-\Phi_T+\epsone
=S_E-\Phi_T+\Phi_T-(S_E-1)=1.
\]
La fuerza de la prueba no está en esta cancelación elemental, sino en que la
genealogía previa fija los dos operandos sin usar el cierre como entrada.
\end{proof}

\paragraph{Publicación del cierre en la carta angular}

La carta de la vuelta multiplica una fracción unitaria por
\(360/T_+=180/\pih\). En particular,
\[
\epsR^\circ=\frac{360^\circ}{T_+}\epsone.
\]
Multiplicar \eqref{eq:cierre-unitario} por \(360^\circ/T_+\) y usar
\[
\frac{360}{T_+}S_E
=360\left(\frac5{36}+\frac{25}{9}\alphah\right)
=50+1000\alphah=50+A
\]
produce el cierre angular.

\begin{theorem}[Ecuación exacta del radión]
\BigRadionEquation
\end{theorem}

La ecuación dice algo más preciso que \enquote{el radián es
\(57.2957\ldots{}^\circ\)}. Descompone esa publicación en cuatro operaciones
con demostración de referencia:
\begin{center}
\begin{tabularx}{.96\textwidth}{>{\bfseries}l X}
\toprule
\(50^\circ\) & fase segunda de la rueda; bisagra crítica en una carta declarada;\\
\(A^\circ\) & coordenada común de la publicación parangular;\\
\(-\frac{20}{81}\Delta_4^\circ\) & transporte orientado de la torsión global;\\
\(+\epsR^\circ\) & acarreo real que empalma los dos lifts.\\
\bottomrule
\end{tabularx}
\end{center}

\paragraph{Convergencia por profundidad del refinamiento}

A profundidad \(N\), el generador entrega cilindros racionales
\(p_N\to\pih\) y \(e_N\to e\). Se definen
\[
S_{E,N}=2p_N\left(\frac5{36}+\frac{25}{9}\alphah\right),
\]
\[
\Delta_{4,N}=\frac{p_N-e_N\log p_N}{270}
\left(1+\frac{p_N}{729}\right),\qquad
S_{T,N}=1+\frac{20}{81}\Delta_{4,N},
\]
y
\[
\mathcal E_N=S_{T,N}-S_{E,N}.
\]
La continuidad del funcional en \(3<p<4\), \(2<e<3\), junto a las
derivadas monótonas, transporta cada par de cilindros a un intervalo
certificado. A profundidad \(45\) bloques, su anchura es menor que
\(4.81\times10^{-130}\).

Cada retorno de nueve niveles contiene seis trits por nivel. Por eso la tasa
de refinamiento es
\begin{equation}
729^{-9}=3^{-54}.
\label{eq:contraccion-54}
\end{equation}
Esta cifra controla la anchura de los intervalos; no es el valor de
\(\epsR\).

La resta de tríadas estabiliza
\begin{center}
\ttfamily
000|000|000|896|802|822|044|562|981|999|050|695|020|651|286|413
\end{center}
y el prefijo de carries
\begin{center}
\ttfamily
0,0,0,0,0,-1,0,-1,-1,-1,0,-1,0.
\end{center}
Éste es el mecanismo concreto de memoria decimal: no una cola elegida, sino
la recurrencia \eqref{eq:subacarreo} aplicada a operandos previamente
construidos.

\paragraph{Evaluación de los valores canónicos}

La salida coinductiva es
\begin{align}
\epsone
&=8.9680282204456298199905069502065128641372736205009\ldots
\times10^{-10},\\
\epsR^\circ
&=5.1383016758575282071685164622645947489908574400054\ldots
\times10^{-8}\,{}^\circ.
\label{eq:epsilon-valores}
\end{align}
La publicación completa resulta
\[
\frac{180^\circ}{\pih}
=57.2957795130823208767981548141051703324054724665643\ldots{}^\circ.
\]

\begin{exactbox}
\begin{align*}
50^\circ+A^\circ
&=57.29735256928380099728510547238066299956\ldots{}^\circ,\\
\frac{20}{81}\Delta_4^\circ
&=0.00157310758449687906223272996065728980465\ldots{}^\circ,\\
50^\circ+A^\circ-\frac{20}{81}\Delta_4^\circ
&=57.29577946169930411822287274242000570976\ldots{}^\circ,\\
\epsR^\circ
&=0.00000005138301675857528207168516462264594749\ldots{}^\circ,\\
\text{suma final}
&=57.29577951308232087679815481410517033241\ldots{}^\circ.
\end{align*}
\end{exactbox}

 
\section{Coordenada crítica y sistema temporal dodecafásico orientado: determinación afín de la fase de \(50^\circ\)}
\subsection{Aplicación afín \(j_{100}\) de la carta APP y unicidad de la fase \(m=2\) en \(j_{100}(1/2)=\theta_m=50^\circ\)}
% Fragmento público generado desde fuentes teoremáticas íntegros.
% Residencia: 52.3.1.
% Fuente: /Users/ruben/Documents/HMT2/EL_CIERRE_HOLOGRAFICO_DEL_INFINITO_SUCESOR_102_CAPITULOS_2026-08-28_EN_TRABAJO/incoming/radion_monografia_20260827/manuscrito/sections/04_cincuenta_critica.tex:1-105
\noindent La coordenada crítica enlaza la carta de Cayley, el registro APP y la fase dodecafásica mediante aplicaciones tipadas.\par\medskip

\subsubsection{Dos realizaciones exactas de la coordenada angular \(50^\circ\)}

\paragraph{Segunda fase del sistema temporal dodecafásico}

Por \eqref{eq:rueda-dodecafasica}, la sucesión de centros comienza
\[
20^\circ,\ 50^\circ,\ 80^\circ,\ 110^\circ,\ldots.
\]
Resolver
\[
30m-10=50
\]
da únicamente \(m=2\). Por tanto, \(50^\circ\) no es una redondeo de \(A\),
una media vuelta ni el centro de un círculo: es una fase tipada del calendario
dodecafásico. La fracción de vuelta es
\[
\frac{50}{360}=\frac5{36},
\]
exactamente el primer sumando de \(S_E/T_+\).

\paragraph{Igualación de la coordenada crítica de Cayley}

La involución funcional
\[
\rho\longmapsto1-\bar\rho
\]
fija
\[
\operatorname{Fix}(\rho\mapsto1-\bar\rho)
=\{\rho:\Re\rho=1/2\}.
\]
El número \(1/2\) es la coordenada real autodual del intervalo funcional
\(0\leq\Re s\leq1\). No es \enquote{la mitad de una recta} en sentido
longitudinal, y no contiene la altura espectral \(\gamma\). Para
\[
\rho_\gamma=\frac12+i\gamma,
\]
la transformación de Cayley se puede escribir
\begin{equation}
\tau_\gamma
=\frac{\rho_\gamma-1}{\rho_\gamma}
=\frac{\gamma+i/2}{\gamma-i/2}
\in S^1,
\qquad
\gamma=\frac12\cot\frac{\theta_\gamma}{2}.
\label{eq:cayley-critica}
\end{equation}
Así, \(1/2\) fija la costura, \(\gamma\) determina el carácter angular y
\(\Gamma_9\) conserva la profundidad de sus iteraciones.

\paragraph{Inyección afín \(j_{100}\) del registro APP en la carta angular}

La completación cúbica dispone una cámara de diez marcas por coordenada. Una
cara contiene \(10^2=100\) posiciones. Declaramos la carta afín
\begin{equation}
j_{100}:[0,1]\longrightarrow[0,100^\circ],
\qquad j_{100}(\sigma)=100\sigma^\circ.
\label{eq:j100}
\end{equation}
Entonces
\[
j_{100}\left(\frac12\right)=50^\circ.
\]
Al exigir simultáneamente \(j_{100}(1/2)=\theta_m\), resulta
\[
100\cdot\frac12=30m-10
\quad\Longleftrightarrow\quad m=2.
\]

\begin{theorem}[Igualador crítico--dodecafásico]
En la carta APP de cien marcas, la coordenada fija de la involución crítica
se publica exactamente en la segunda fase de la rueda:
\begin{equation}
\boxed{j_{100}(1/2)=\theta_2=50^\circ.}
\label{eq:igualador-50}
\end{equation}
La fase \(m=2\) es única.
\end{theorem}

El teorema es exacto una vez declarada \eqref{eq:j100}. Su estatuto es
\status{formalización nueva}: el medio crítico, la completación \(10^3\), la
cara de cien marcas y \(\theta_2\) ya estaban construidos; la elección de esa
cara como carta crítica afín reúne por primera vez sus dominios.

\paragraph{Compatibilidad del valor crítico con la carta de Cayley}

La transformación de Cayley no envía \(1/2\) a \(50^\circ\). Si
\(\gamma=0\), \eqref{eq:cayley-critica} da
\[
\tau_0=-1=e^{i\pi},
\]
es decir, \(180^\circ\). Los \(50^\circ\) pertenecen a la carta APP de cien
marcas y al calendario dodecafásico; el ángulo Cayley depende de \(\gamma\).
Esta distinción impide colapsar todos los ceros de la función zeta en una sola
fase.

\begin{falsifier}
La identificación \(1/2\leftrightarrow50^\circ\) falla si se elimina la
carta \eqref{eq:j100}. La igualdad escalar \(100(1/2)=50\) no basta para
identificar por sí sola dos objetos. El teorema depende de la cara APP
distinguida y de la rueda dodecafásica.
\end{falsifier}

 \subsection{Compatibilidad de la fase crítica con el retorno nonádico y el avance de memoria}
% Fragmento público generado desde fuentes teoremáticas íntegros.
% Residencia: 52.3.2.
% Fuente: /Users/ruben/Documents/HMT2/EL_CIERRE_HOLOGRAFICO_DEL_INFINITO_SUCESOR_102_CAPITULOS_2026-08-28_EN_TRABAJO/incoming/radion_monografia_20260827/manuscrito/sections/04_cincuenta_critica.tex:106-171
\subsubsection{Separación tipada de los 2 objetos de valor \(1/2\)}

APP contiene además un modo singular con
\[
s=\frac12,\qquad s^2=\frac14,\qquad
1-s^2=\frac34=\frac{90}{120},\qquad
s\sqrt{1-s^2}=\frac{\sqrt3}{4}.
\]
Este \(s\) nace de los ángulos principales de las hojas suma y producto. La
costura crítica \(\Re\rho=1/2\) nace de una involución funcional. Que ambas
coordenadas valgan \(1/2\) es una coincidencia escalar con posible contenido,
pero no autoriza a identificarlas sin un cuadrado tipado.

La forma correcta de conservar el hallazgo es un diagrama de igualadores:
\[
\begin{tikzpicture}[baseline=(current bounding box.center),node distance=14mm and 23mm,
 every node/.style={font=\small}]
\node[draw=RadBlue,rounded corners,fill=RadMist] (crit) {costura crítica \(\Re s=1/2\)};
\node[draw=RadTeal,rounded corners,fill=RadCream,right=of crit] (face) {carta \(j_{100}\)};
\node[draw=RadCopper,rounded corners,fill=RadCream,right=of face] (theta) {fase \(\theta_2=50^\circ\)};
\node[draw=RadRose,rounded corners,fill=white,below=of face] (mode) {modo APP \(s=1/2\), \(1-s^2=90/120\)};
\draw[-{Latex},thick,RadBlue] (crit)--node[above]{\(j_{100}\)}(face);
\draw[-{Latex},thick,RadTeal] (face)--node[above]{igualador}(theta);
\draw[-{Latex},dashed,RadRose] (mode)--node[right,align=left]{misma cifra;\\dominio distinto}(face);
\end{tikzpicture}
\]
El diagrama no finge la flecha ausente. Muestra qué está demostrado, qué se
ha formalizado y qué igualdad sigue necesitando un transporte adicional si
se pretende una identificación operatoria global.

\subsubsection{Compatibilidad entre círculo espectral y memoria helicoidal}

El lector Cayley--Pontryagin transporta una fibra espectral mediante
\[
\widetilde U_{\Gamma_9^n}J\psi_\gamma
=\tau_\gamma^nJ\psi_\gamma.
\]
La publicación circular es \(\tau_\gamma^n\in S^1\). El lift
\[
(\tau_\gamma^n,n)\in S^1\times\Z
\]
conserva el número de retornos. En la hoja contractiva aparece
\[
\widetilde M_9^nJ\psi_\gamma
=\left(\frac59\right)^n\tau_\gamma^nJ\psi_\gamma.
\]
Tenemos, por tanto, tres imágenes tipadas:
\[
\text{círculo de fase},\qquad
\text{hélice de memoria},\qquad
\text{espiral contractiva}.
\]

Un cero no es una vuelta, y un nueve no es un cero particular de la función
zeta. El cero fija un carácter \(\tau_\gamma\) que persiste a través de las
vueltas; el retorno nonádico incrementa memoria. El nueve funciona como cero
residual en \(\Z/9\Z\), una operación aritmética distinta de la enumeración
de ceros espectrales.

\begin{narrativebox}
La recta crítica se vuelve círculo cuando olvidamos la altura lineal y
conservamos el carácter unitario. Se vuelve hélice cuando volvemos a insertar
la memoria de cuántas veces el carácter ha sido transportado. La profundidad
no es una tercera coordenada decorativa: es la información que la proyección
circular no puede retener.
\end{narrativebox}
 
\section{Par angular y acción reducida: construcción de la elipse positiva y del área por radión}
\subsection{Ley areal del radión: \(\mathcal A(\theta)=\hbar_{\mathrm{ret}}\theta\) y \(\mathcal A(2\pi)=h_{\mathrm{ret}}\)}
% Fragmento público generado desde fuentes teoremáticas íntegros.
% Residencia: 52.4.1.
% Fuente: /Users/ruben/Documents/HMT2/EL_CIERRE_HOLOGRAFICO_DEL_INFINITO_SUCESOR_102_CAPITULOS_2026-08-28_EN_TRABAJO/incoming/radion_monografia_20260827/manuscrito/sections/05_geometria_accion.tex:1-217
\noindent La geometría espectral determina la forma y la ley areal determina la acción transportada por la fase.\par\medskip

\subsubsection{Equivalencia entre las cartas espectral y geométrica}

\paragraph{Operador espectral positivo y autovalores conjugados}

El par angular produce el operador bicapa
\begin{equation}
B_1=AI+C^*R,\qquad R^2=I,\qquad
P_\pm=\frac{I\pm R}{2}.
\label{eq:B1}
\end{equation}
Sus autovalores son \(A\pm C^*\). Como \(0<C^*<A\), ambos son positivos.
En la carta \(\kappa=\pi/180\), la elipse espectral tiene semiejes
\begin{equation}
a_1^2=\kappa(A+C^*),\qquad
b_1^2=\kappa(A-C^*).
\label{eq:semiejes-espectrales}
\end{equation}
La descomposición es reversible:
\begin{equation}
A=\frac{a_1^2+b_1^2}{2\kappa},\qquad
C^*=\frac{a_1^2-b_1^2}{2\kappa}.
\label{eq:reconstruccion-AC}
\end{equation}
La elipse no se dibuja para ilustrar un número. Reconstruye exactamente el
par que la generó.

Definamos
\begin{equation}
\rho_1=\sqrt{A^2-C^{*2}},\qquad
\eta=\operatorname{artanh}\frac{C^*}{A}
=\frac12\log\frac{A+C^*}{A-C^*}.
\label{eq:eta}
\end{equation}
Entonces
\[
A\pm C^*=\rho_1e^{\pm\eta},
\]
y
\begin{equation}
a_1=\sqrt{\kappa\rho_1}\,e^{\eta/2},
\qquad
b_1=\sqrt{\kappa\rho_1}\,e^{-\eta/2}.
\label{eq:semiejes-eta}
\end{equation}
El producto conserva escala; el cociente conserva forma:
\[
a_1b_1=\kappa\rho_1,
\qquad
\frac{b_1}{a_1}=e^{-\eta}
=\sqrt{\frac{A-C^*}{A+C^*}}.
\]

\paragraph{Invariantes focales del operador positivo}

La semidistancia focal es
\begin{equation}
c_1^2=a_1^2-b_1^2=2\kappa C^*,
\qquad
F_\pm=(\pm c_1,0).
\label{eq:focos}
\end{equation}
La distancia total vale
\begin{equation}
d_1=2c_1=2\sqrt{2\kappa C^*}
=0.544545509325626623406299779866023\ldots.
\label{eq:d1}
\end{equation}
El número aislado no es la tesis. La identidad \(c_1^2=2\kappa C^*\) dice que
\(C^*\) es literalmente la escisión focal cuadrática en la carta
distinguida. Los focos son los extremos conjugados de la apertura espectral;
no son, por sí solos, los polos eléctrico y magnético. Esa lectura requiere
el transductor constitutivo.

La excentricidad es
\begin{equation}
e_f=\frac{c_1}{a_1},\qquad
e_f^2=1-e^{-2\eta}=\frac{2C^*}{A+C^*}.
\label{eq:excentricidad}
\end{equation}
Numéricamente,
\[
\begin{aligned}
\eta&=0.299689683921630799684926562863927\ldots,\\
e_f&=0.671451895550191986916277055864332\ldots.
\end{aligned}
\]

\paragraph{Bloque APP de referencia y transporte normalizado}

El bloque central
\[
B_0=7I+2R=9P_++5P_-
\]
posee
\[
\rho_0=\sqrt{45},\qquad
\eta_0=\frac12\log\frac95,\qquad
e_0=\frac23.
\]
Su distancia focal es
\[
d_0=4\sqrt\kappa
=0.528443639680801468446003835349358\ldots.
\]
La comparación exacta es
\begin{equation}
\boxed{\left(\frac{d_1}{d_0}\right)^2=\frac{C^*}{2}},
\qquad
d_1=d_0\sqrt{\frac{C^*}{2}}.
\label{eq:invariante-focal-relativo}
\end{equation}
Ésta es la genealogía defendible de \(0.544545\ldots\): una escala focal APP
de referencia multiplicada por un invariante relativo. Su cercanía visual a
\(0.54\) no es una identidad con el defecto APP \(54\).

El transporte desde \(B_0\) se separa en escala y anisotropía:
\begin{equation}
T_{\mathrm{rel}}=\frac{\rho_1}{\sqrt{45}}
\exp\bigl((\eta-\eta_0)R\bigr).
\label{eq:transporte-rel}
\end{equation}
Sobre las hojas,
\[
t_+=\frac{A+C^*}{9}=1.0467878786947163\ldots,
\qquad
t_-=\frac{A-C^*}{5}=1.0347228460630311\ldots.
\]
La deformación no es un solo decimal. Su parte isotrópica es
\(\sqrt{t_+t_-}\); su parte hiperbólica es
\(\tfrac12\log(t_+/t_-)\).

\subsubsection{Elipse positiva asociada al par angular}

\paragraph{Construcción mediante semiejes conjugados}

La escala espectral anterior carece de unidades de acción. Mecánica
Dimensional publica la misma forma sobre la recta interna de acción:
\begin{equation}
a_S=\sqrt{2\hret e^\eta},
\qquad
b_S=\sqrt{2\hret e^{-\eta}}.
\label{eq:semiejes-accion}
\end{equation}
Inmediatamente,
\begin{equation}
a_Sb_S=2\hret,\qquad
\frac{a_S}{b_S}=e^\eta.
\label{eq:producto-forma}
\end{equation}
El producto fija la acción; el cociente fija la anisotropía. La elipse se
parametriza por
\begin{equation}
Q(\theta)=a_S\cos\theta,\qquad
P(\theta)=b_S\sin\theta.
\label{eq:elipse-param}
\end{equation}

Las elipses espectral y de acción son homotéticas:
\begin{equation}
\lambda_{\mathrm{act}}^2=\frac{2\hret}{\kappa\rho_1},
\qquad
(a_S,b_S,c_S)=\lambda_{\mathrm{act}}(a_1,b_1,c_1).
\label{eq:homotecia}
\end{equation}
Comparten la forma \(e^{-\eta}\), no la dimensión física.

\paragraph{Ley areal de la acción por radión}

La forma simpléctica canónica es \(\omega=\diff Q\wedge\diff P\). El área
orientada barrida desde \(0\) hasta \(\theta\) es
\begin{align}
\mathcal A(\theta)
&=\frac12\int_0^\theta
\bigl(Q(t)P'(t)-P(t)Q'(t)\bigr)\,\diff t\\
&=\frac12a_Sb_S\int_0^\theta
(\cos^2t+\sin^2t)\,\diff t.
\end{align}
Usando \eqref{eq:producto-forma}, obtenemos el resultado nuclear.

\begin{theorem}[Acción por radión]
Para la elipse \eqref{eq:elipse-param},
\begin{equation}
\boxed{\mathcal A(\theta)=\hret\theta.}
\label{eq:area-theorem}
\end{equation}
En particular,
\begin{equation}
\boxed{\mathcal A(1)=\hret},
\qquad
\boxed{\mathcal A(2\pi)=2\pi\hret=\hfull}.
\label{eq:area-radion-vuelta}
\end{equation}
\end{theorem}

Esta igualdad permite definir el radión sin apelar primero a la longitud de
arco.

\begin{definition}[Radión orientado]
\begin{equation}
\mathfrak r_\sigma=(\theta=1,\sigma),
\qquad \sigma\in\{+1,-1\},
\qquad
\mathcal A(\mathfrak r_\sigma)=\sigma\hret.
\label{eq:radion-oriented}
\end{equation}
\end{definition}

El radián es la proyección que olvida \(\sigma\), la hoja, la memoria y el
área. El radión conserva esa información. De ahí el nombre \emph{rad--ión}:
la orientación precede a la carga física, y la carga aparece más tarde como
\[
Q(x)=n_Q(x)e_{\mathrm{int}},\qquad
e_{\mathrm{int}}=\sqrt{\frac{2\alphah\hfull}{Z_{\mathrm{int}}}}.
\]

 \subsection{Semiejes conjugados, invariantes focales y orientación de la anisotropía}
% Fragmento público generado desde fuentes teoremáticas íntegros.
% Residencia: 52.4.2.
% Fuente: /Users/ruben/Documents/HMT2/EL_CIERRE_HOLOGRAFICO_DEL_INFINITO_SUCESOR_102_CAPITULOS_2026-08-28_EN_TRABAJO/incoming/radion_monografia_20260827/manuscrito/sections/05_geometria_accion.tex:218-290
\paragraph{Correspondencia entre fase, acción, área y orientación}

\begin{figure}[H]
\centering
\begin{tikzpicture}[scale=1.12]
  \def\aa{4.0}
  \def\bb{2.95}
  \def\cc{2.70}
  \draw[->,RadGray] (-4.7,0)--(4.7,0) node[right]{\(Q\)};
  \draw[->,RadGray] (0,-3.5)--(0,3.5) node[above]{\(P\)};
  \fill[RadTeal!8] (0,0) ellipse[x radius=\aa,y radius=\bb];
  \draw[RadTeal,line width=1.35pt] (0,0) ellipse[x radius=\aa,y radius=\bb];
  \fill[RadCopper] (-\cc,0) circle (2pt) node[below=2mm]{\(F_-\)};
  \fill[RadCopper] (\cc,0) circle (2pt) node[below=2mm]{\(F_+\)};
  \draw[RadGold,line width=1.1pt,-{Latex}] (\aa,0)
     arc[start angle=0,end angle=57.3,x radius=\aa,y radius=\bb];
  \node[RadGold] at (3.7,2.3) {un radión};
  \draw[dashed,RadBlue] (0,0)--(\aa,0) node[midway,below]{\(a_S\)};
  \draw[dashed,RadBlue] (0,0)--(0,\bb) node[midway,left]{\(b_S\)};
  \node[align=center,RadNavy] at (0,-3.25)
    {área total \(\pi a_Sb_S=2\pi\hret=\hfull\)};
\end{tikzpicture}
\caption{La elipse de acción. Un arco paramétrico de un radián barre
\(\hret\); una vuelta encierra \(\hfull\). Los focos codifican la apertura
\(C^*\) en la carta distinguida.}
\label{fig:elipse-accion}
\end{figure}

\paragraph{Relación entre la acción reducida, la incertidumbre y el área}

La relación \(\hbar=h/(2\pi)\) suele presentarse como una conveniencia
notacional: al describir oscilaciones, rotaciones y transformadas de Fourier,
los factores \(2\pi\) desaparecen si se usa \(\hbar\). La elipse de acción
revela una lectura más fuerte. La división por \(2\pi\) es exactamente el
paso de la acción de una vuelta a la acción por unidad de fase:
\[
h=\mathcal A(2\pi),\qquad \hbar=\mathcal A(1).
\]
La constante reducida no es sólo \(h\) con una notación más cómoda; es la
densidad areal de acción respecto a la fase.

Heisenberg, Dirac y Jordan establecieron el lenguaje de variables conjugadas,
conmutadores y transformaciones canónicas. No se les atribuye aquí una
\enquote{elipse ontológica} literal sin fuente. La contribución HMT--MD es
identificar la figura positiva, construir sus semiejes desde \((A,C^*)\) y
demostrar \eqref{eq:area-theorem}. La genealogía histórica abre el lenguaje;
el teorema actual fija el objeto.

\begin{narrativebox}
La acción no se añade después al movimiento. En la elipse del radión, la
acción es el área que el movimiento genera. Una vuelta no transporta una
constante prefabricada: la encierra.
\end{narrativebox}

\Needspace{10\baselineskip}
\paragraph{Controles algebraicos de los invariantes focales}

Las semejanzas decimales se someten a control:
\begin{center}
\begin{tabularx}{.96\textwidth}{l r X}
\toprule
Comparación & Residuo & Dictamen\\
\midrule
\(d_1\) frente a \(0.54\) & \(4.5455\times10^{-3}\) & no igualdad; el \(54\) no se deduce;\\
\(d_1^2\) frente a \(8/27\) & \(2.3352\times10^{-4}\) & proximidad sin mapa;\\
\(\eta\) frente a \(3/10\) & \(-3.1032\times10^{-4}\) & no igualdad;\\
\(e_f\) frente a \(2/3\) & \(4.7852\times10^{-3}\) & \(2/3\) pertenece exactamente a \(B_0\);\\
\(e^{-\eta}\) frente a \(20/27\) & \(3.0740\times10^{-4}\) & no igualdad.\\
\bottomrule
\end{tabularx}
\end{center}
Los invariantes sólidos son \eqref{eq:eta}, \eqref{eq:excentricidad} y
\eqref{eq:invariante-focal-relativo}. El criterio HMT no consiste en premiar
un parecido decimal, sino en exigir objeto, operación y conservación.
 
\section{Retorno nonádico e interior proyectivo: disco de Klein y profundidad helicoidal de memoria}
\subsection{Normalización al disco e interior hiperbólico de Hilbert--Beltrami--Klein}
% Fragmento público generado desde fuentes teoremáticas íntegros.
% Residencia: 52.5.1.
% Fuente: /Users/ruben/Documents/HMT2/EL_CIERRE_HOLOGRAFICO_DEL_INFINITO_SUCESOR_102_CAPITULOS_2026-08-28_EN_TRABAJO/incoming/radion_monografia_20260827/manuscrito/sections/06_interior_helice.tex:1-108
\noindent La geometría proyectiva distingue la acción finita de la profundidad hiperbólica y la memoria del retorno.\par\medskip

\subsubsection{Interior hiperbólico de Hilbert--Beltrami--Klein}

\paragraph{Normalización proyectiva en el disco de Klein}

La aplicación afín
\[
u=\frac{Q}{a_S},\qquad v=\frac{P}{b_S}
\]
lleva la elipse de acción al disco unitario
\[
D=\{(u,v):u^2+v^2<1\}.
\]
Sobre cada cuerda, la distancia de Hilbert se define mediante la razón doble
de sus extremos de frontera. La métrica infinitesimal resultante es la forma
Beltrami--Klein
\begin{equation}
\diff s_K^2=
\frac{(1-r^2)(\diff u^2+\diff v^2)+(u\diff u+v\diff v)^2}
{(1-r^2)^2},
\qquad r^2=u^2+v^2.
\label{eq:klein-metric}
\end{equation}
Con esta normalización, la curvatura gaussiana es constante:
\begin{equation}
K=-1.
\label{eq:klein-curvature}
\end{equation}
Las geodésicas son cuerdas euclídeas, pero su longitud hiperbólica diverge al
acercarse a la frontera.

\paragraph{Finitud de la frontera e infinitud de la profundidad hiperbólica}

El determinante métrico de \eqref{eq:klein-metric} produce
\begin{equation}
\diff A_K=\frac{\diff u\,\diff v}{(1-r^2)^{3/2}}.
\label{eq:klein-area-density}
\end{equation}
El área del subdisco \(r<R<1\) es
\begin{align}
A_K(R)
&=\int_0^{2\pi}\int_0^R
\frac{r\,\diff r\,\diff\theta}{(1-r^2)^{3/2}}\\
&=2\pi\left(\frac1{\sqrt{1-R^2}}-1\right).
\label{eq:klein-area-R}
\end{align}
\Needspace{4\baselineskip}
Por tanto,
\[
\lim_{R\to1^-}A_K(R)=+\infty.
\]
Simultáneamente, la misma frontera encierra el área simpléctica
\[
\pi a_Sb_S=2\pi\hret=\hfull<\infty.
\]

\begin{theorem}[Finitud de acción e infinitud de profundidad]
La elipse del radión posee una acción de vuelta finita \(\hfull\), mientras
su interior, equipado con la métrica Hilbert--Beltrami--Klein, tiene área
hiperbólica infinita:
\begin{equation}
\boxed{
\operatorname{Area}_{\mathrm{simp}}(\partial\mathbb E_S)=\hfull<\infty,
\qquad
\operatorname{Area}_{K}(\mathbb E_S)=+\infty.}
\label{eq:finite-infinite}
\end{equation}
\end{theorem}

No hay contradicción: \(\operatorname{Area}_{\mathrm{simp}}\) y
\(\operatorname{Area}_K\) son medidas diferentes. La primera cuenta acción
orientada en fase; la segunda mide profundidad proyectiva interior. La tesis
\enquote{borde finito, interior infinito} es ahora una igualdad y un límite,
no una metáfora.

\paragraph{Coordenadas focales en el interior de Klein}

En el disco normalizado, los focos se sitúan en \((\pm e_f,0)\). La distancia
del centro a un foco es
\begin{equation}
u_f=\operatorname{artanh}e_f
=\operatorname{arcosh}(e^\eta)
=0.813382331175342333760889731088228\ldots.
\label{eq:klein-focal}
\end{equation}
La distancia foco--foco es
\[
d_K(F_-,F_+)=2u_f
=1.62676466235068466752177946217646\ldots.
\]
Se cumplen
\begin{equation}
e_f=\tanh u_f,\qquad
e^{-\eta}=\operatorname{sech}u_f,\qquad
\eta=\log\cosh u_f.
\label{eq:focal-hyperbolic-chain}
\end{equation}
La región hasta el radio focal tiene área
\[
A_K(e_f)=2\pi(e^\eta-1)
=2.19559620884061869541206115980360\ldots.
\]

Todas las elipses abiertas son proyectivamente isométricas como dominios de
Hilbert. Por eso la curvatura \(-1\) sola no recuerda la excentricidad. La
forma HMT reside en la estructura conjunta: carta afín distinguida, forma
simpléctica, etiquetado de hojas y parametrización de frontera.
 \subsection{Diferencia entre retorno de fase y retorno del estado holonómico integral}
% Fragmento público generado desde fuentes teoremáticas íntegros.
% Residencia: 52.5.2.
% Fuente: /Users/ruben/Documents/HMT2/EL_CIERRE_HOLOGRAFICO_DEL_INFINITO_SUCESOR_102_CAPITULOS_2026-08-28_EN_TRABAJO/incoming/radion_monografia_20260827/manuscrito/sections/06_interior_helice.tex:109-201
\subsubsection{Elevación helicoidal del retorno nonádico}

\paragraph{Descomposición nonádica y cociente de memoria}

La operación elemental
\begin{equation}
n=9q_9(n)+\rho_9(n),
\qquad
H_9(n)=(\rho_9(n),q_9(n))
\label{eq:H9}
\end{equation}
satisface
\begin{equation}
H_9(n+9)=H_9(n)+(0,1).
\label{eq:H9-helix}
\end{equation}
La primera coordenada cierra; la segunda asciende. Esta identidad es la
versión aritmética mínima de un tornillo.

Para una fibra espectral,
\[
n\longmapsto(\tau_\gamma^n,n)
\]
es una hélice sobre \(S^1\). Añadir la contracción \((5/9)^n\) produce una
espiral interior. El radión participa en ambas lecturas: una unidad de fase
en la frontera y una unidad orientada de acción en el lift.

\paragraph{Rango de memoria como lector dimensional}

Supongamos que una proyección visible conserva el valor de fase, pero pierde
\(d\) cociclos de memoria linealmente independientes. Cualquier lift fiel
necesita al menos \(d\) coordenadas adicionales para reconstruirlos. Esto
motiva una lectura precisa de Mecánica Dimensional:
\begin{equation}
\begin{aligned}
\text{dimensión emergente}
&=\text{rango mínimo de memoria independiente}\\
&\phantom{=}\text{no publicable en la hoja visible}.
\end{aligned}
\label{eq:dimension-memory}
\end{equation}
No se afirma que toda dimensión física sea un contador. Se afirma que, en el
dominio TPK, cada memoria independiente que debe sobrevivir fuerza una
coordenada del lift.

\begin{narrativebox}
La dimensión aparece cuando la proyección ya no puede conservar toda la
historia. El círculo es suficiente para decir dónde está la fase; no es
suficiente para decir cuántas veces volvió, qué hoja transportó, qué intervalo
se contrajo o qué acarreo sobrevivió.
\end{narrativebox}

\paragraph{Ensamblaje de cilindros, espirales y trayectorias helicoidales}

Los cilindros coinductivos no son tubos físicos añadidos a una espiral. Son
intervalos encajados que conservan prefijos y estrechan el valor. La espiral
describe contracción transversal; la hélice describe memoria longitudinal;
el cilindro describe el conjunto de valores compatibles a una profundidad
finita. En el límite, la intersección de cilindros selecciona una sección
única.

La imagen espacial reúne, sin confundir:
\[
\begin{array}{ccl}
\text{ángulo} &\leftrightarrow& \text{fase visible},\\
\text{altura} &\leftrightarrow& \text{retorno/memoria},\\
\text{radio de cilindro} &\leftrightarrow& \text{incertidumbre de refinamiento},\\
\text{hoja} &\leftrightarrow& \text{orientación bidireccional}.
\end{array}
\]
La tasa \(3^{-54}\) contrae el cilindro por retorno; no se suma al ángulo y
no reemplaza el acarreo real \(\epsR\).

\subsubsection{Cierres de fase de órdenes \(2\pi\) y \(4\pi\)}

En la proyección circular, \(2\pi\) restaura la fase. En un lift orientado
con dos hojas, una vuelta puede cerrar en la hoja conjugada y necesitar una
segunda vuelta para restaurar la orientación:
\[
U_{2\pi}=-I,\qquad U_{4\pi}=I.
\]
La estructura coincide con la periodicidad espinorial cuando se aplica el
funtor a una representación \(SU(2)\); no se identifica el lift TPK con
\(SU(2)\) sin ese mapa.

La lectura areal es compatible:
\[
\mathcal A(2\pi)=\hfull,\qquad
\mathcal A(4\pi)=2\hfull.
\]
En una banda de Möbius, un circuito visible invierte hoja y dos circuitos
restauran la orientación. El \(2\pi/4\pi\) no es una duplicación accidental;
es la diferencia entre retorno de fase y retorno del estado holonómico integral.
 
\section{Elipse orientada y variables conjugadas: incertidumbre mínima y oscilación electromagnética}
\subsection{Variables simplécticas conjugadas y determinación de la incertidumbre mínima}
% Fragmento público generado desde fuentes teoremáticas íntegros.
% Residencia: 52.6.1.
% Fuente: /Users/ruben/Documents/HMT2/EL_CIERRE_HOLOGRAFICO_DEL_INFINITO_SUCESOR_102_CAPITULOS_2026-08-28_EN_TRABAJO/incoming/radion_monografia_20260827/manuscrito/sections/07_incertidumbre_oscilador.tex:1-111
\noindent La forma simpléctica coordina acción, incertidumbre y dinámica oscilatoria sin alterar la genealogía del estado.\par\medskip

\subsubsection{Realización del estado en un espacio de Hilbert}

\paragraph{Matriz de covarianza mínima}

Una vez construida la elipse de acción, puede realizarse en un par de
operadores conjugados
\[
[\widehat Q,\widehat P]=i\hret.
\]
Definimos
\begin{equation}
V_\eta=\frac{\hret}{2}
\begin{pmatrix}
e^\eta&0\\[2pt]0&e^{-\eta}
\end{pmatrix}.
\label{eq:covariance}
\end{equation}
Su determinante es
\begin{equation}
\det V_\eta=\frac{\hret^2}{4}.
\label{eq:cov-det}
\end{equation}
Por consiguiente,
\begin{equation}
(\Delta Q)^2=\frac{\hret}{2}e^\eta,
\qquad
(\Delta P)^2=\frac{\hret}{2}e^{-\eta},
\qquad
\Delta Q\,\Delta P=\frac{\hret}{2}.
\label{eq:uncertainty}
\end{equation}
La desigualdad de Robertson--Schrödinger se satura.

El aniquilador comprimido
\begin{equation}
\widehat a_\eta=
\frac{e^{-\eta/2}\widehat Q+i e^{\eta/2}\widehat P}
{\sqrt{2\hret}}
\label{eq:a-eta}
\end{equation}
satisface \([\widehat a_\eta,\widehat a_\eta^\dagger]=1\). Su vacío tiene
función de onda
\begin{equation}
\psi_\eta(Q)=
(\pi\hret e^\eta)^{-1/4}
\exp\left(-\frac{Q^2}{2\hret e^\eta}\right).
\label{eq:gaussian}
\end{equation}

\paragraph{Elipse como nivel del funcional de Mahalanobis}

El conjunto
\begin{equation}
\mathbb E_{\mathrm{act}}
=\{x\in\R^2:x^{\mathsf T}V_\eta^{-1}x\le4\}
\label{eq:mahalanobis}
\end{equation}
tiene semiejes
\begin{equation}
a_S=2\Delta Q,\qquad b_S=2\Delta P.
\label{eq:two-sigma}
\end{equation}
Su área es
\[
4\pi\sqrt{\det V_\eta}=2\pi\hret=\hfull.
\]
La elipse HMT construida antes coincide exactamente con el contorno de dos
desviaciones del estado gaussiano mínimo.

\begin{proposition}[Cadena única del cociente de forma]
Tras declarar las interfaces espectral, simpléctica y de Hilbert,
\begin{equation}
\boxed{
e^{-\eta}
=\sqrt{\frac{A-C^*}{A+C^*}}
=\frac{b_1}{a_1}
=\frac{b_S}{a_S}
=\frac{\Delta P}{\Delta Q}.}
\label{eq:shape-chain-hilbert}
\end{equation}
\end{proposition}

La elipse de acción corresponde al nivel clásico \(H=\hret\omega\) y al
contorno \(2\sigma\) del vacío comprimido. No debe confundirse con la energía
media del vacío, \(\hret\omega/2\). Esta diferencia de factor dos es un
control físico obligatorio.

\paragraph{Condiciones de invariancia de la elipse como objeto operatorio}

En la formulación convencional, una elipse de incertidumbre puede verse como
un recurso gráfico para representar una matriz de covarianza. Aquí el orden
causal es inverso: la elipse positiva ya existe como publicación de
\((A,C^*,\hret)\); el espacio de Hilbert la reconoce mediante
\eqref{eq:covariance}. La realidad ontológica reclamada no significa que una
partícula clásica recorra físicamente un óvalo en el plano \((Q,P)\). Significa
que el estado conserva una forma positiva, una acción total y un cociente de
orientación antes de que una teoría probabilista los interprete.

La función de onda no crea la forma; la representa. La incertidumbre no crea
\(\hbar\); publica la imposibilidad de comprimir simultáneamente las dos
direcciones sin alterar el producto simpléctico.

\begin{narrativebox}
La lectura probabilista es una realización posible de la elipse, no su causa.
El objeto primario es la conservación conjunta de producto y orientación. La
probabilidad aparece cuando esa geometría se publica como covarianza de un
estado sobre Hilbert.
\end{narrativebox}

 \subsection{Intercambio \(\mu\leftrightarrow\varepsilon\), inversión \(Z\leftrightarrow Z^{-1}\) y conservación de la velocidad causal}
% Fragmento público generado desde fuentes teoremáticas íntegros.
% Residencia: 52.6.2.
% Fuente: /Users/ruben/Documents/HMT2/EL_CIERRE_HOLOGRAFICO_DEL_INFINITO_SUCESOR_102_CAPITULOS_2026-08-28_EN_TRABAJO/incoming/radion_monografia_20260827/manuscrito/sections/07_incertidumbre_oscilador.tex:156-197
\paragraph{Intercambio energético entre los sectores eléctrico y magnético}

Los dos términos de \eqref{eq:HLC} se convierten en
\begin{equation}
U_E(\theta)=\hret\omega\cos^2\theta,
\qquad
U_B(\theta)=\hret\omega\sin^2\theta.
\label{eq:energy-exchange}
\end{equation}
Por tanto,
\[
U_E+U_B=\hret\omega.
\]
La polaridad eléctrica--magnética no consiste en etiquetar arbitrariamente
los focos. Consiste en dos reservas conjugadas que se intercambian durante la
órbita manteniendo constantes frecuencia y acción total.

El cociente común se prolonga:
\begin{equation}
\boxed{
e^{-\eta}
=\frac{b_S}{a_S}
=\frac{\Delta P}{\Delta Q}
=\frac{Z_\eta}{Z_*}
=\sqrt{\frac{A-C^*}{A+C^*}}.}
\label{eq:shape-chain-LC}
\end{equation}
Los productos conservados son
\[
a_Sb_S=2\hret,\qquad
\det V_\eta=\frac{\hret^2}{4},
\qquad
L_\eta C_\eta=\omega^{-2}.
\]

\begin{exactbox}
El oscilador LC es la realización física mínima de la doble lectura: el
producto conserva acción y frecuencia; el cociente conserva orientación e
impedancia. Un ciclo completo intercambia electricidad y magnetismo sin
perder el área \(\oint P\,\diff Q=\hfull\).
\end{exactbox}

 \subsection{Par inductivo--capacitivo \((L_\eta,C_\eta)\): invariancia de \(L_\eta C_\eta=\omega^{-2}\) y orientación de la impedancia}
% Fragmento público generado desde fuentes teoremáticas íntegros.
% Residencia: 52.6.3.
% Fuente: /Users/ruben/Documents/HMT2/EL_CIERRE_HOLOGRAFICO_DEL_INFINITO_SUCESOR_102_CAPITULOS_2026-08-28_EN_TRABAJO/incoming/radion_monografia_20260827/manuscrito/sections/07_incertidumbre_oscilador.tex:112-155
\subsubsection{Oscilador inductivo--capacitivo asociado al par angular}

\paragraph{Construcción independiente de los 2 depósitos conjugados}

Dados \(\omega>0\) y una impedancia de referencia \(Z_*>0\), definimos
\begin{equation}
L_\eta=\frac{Z_*}{\omega}e^{-\eta},
\qquad
C_\eta=\frac{1}{Z_*\omega}e^\eta.
\label{eq:LC-def}
\end{equation}
Entonces
\begin{equation}
L_\eta C_\eta=\omega^{-2},
\qquad
\sqrt{\frac{L_\eta}{C_\eta}}=Z_*e^{-\eta}.
\label{eq:LC-invariants}
\end{equation}
La frecuencia permanece fija; la impedancia conserva la orientación de la
forma.

El Hamiltoniano es
\begin{equation}
H_{LC}=\frac{q^2}{2C_\eta}+\frac{\Phi^2}{2L_\eta}.
\label{eq:HLC}
\end{equation}
Con las variables normalizadas
\begin{equation}
Q=\sqrt{Z_*}\,q,\qquad
P=\frac{\Phi}{\sqrt{Z_*}},
\label{eq:LC-map}
\end{equation}
se obtiene
\begin{equation}
H_{LC}=\frac\omega2
\left(e^{-\eta}Q^2+e^\eta P^2\right).
\label{eq:HLC-QP}
\end{equation}
Sobre \eqref{eq:elipse-param},
\begin{equation}
H_{LC}=\hret\omega.
\label{eq:LC-level}
\end{equation}

% Fuente: /Users/ruben/Documents/HMT2/EL_CIERRE_HOLOGRAFICO_DEL_INFINITO_SUCESOR_102_CAPITULOS_2026-08-28_EN_TRABAJO/incoming/radion_monografia_20260827/manuscrito/sections/07_incertidumbre_oscilador.tex:198-239
\paragraph{Oscilador LC como unidad dinámica elemental}

El oscilador armónico atraviesa casi toda la física porque linealiza el
entorno de un equilibrio, descompone campos en modos, organiza cuantos y
permite leer propagación como superposición de fases. En el radión, esa
universalidad recibe una genealogía concreta: el modo no empieza en una
ecuación diferencial externa, sino en una forma positiva con acción de vuelta
fija. La dinámica convencional aparece al elegir un reloj \(\omega\) y un
transductor \(Z_*\).

La frecuencia no altera la elipse de acción: cambia la energía del nivel
mediante \(E=\hret\omega\). Esta separación es crucial. La geometría fija la
acción; el reloj convierte acción en energía.

\subsubsection{Reversión temporal y propagaciones conjugadas}

Sea \(H=H^*\) y sea \(J_E\) una estructura compleja con \(J_E^2=-I\).
Supongamos un involutivo \(C\) tal que
\[
CJ_EC^{-1}=-J_E,\qquad CHC^{-1}=H.
\]
La evolución
\begin{equation}
U(t)=\exp\left(-\frac{J_EHt}{\hret}\right)
\label{eq:evolution-J}
\end{equation}
satisface exactamente
\begin{equation}
CU(t)C^{-1}=U(-t).
\label{eq:time-reversal}
\end{equation}
Ésta es la forma operatoria mínima de la bidireccionalidad: la conjugación
intercambia las prolongaciones temporalmente opuestas.

La teoría del absorbedor de Wheeler--Feynman y la lectura de antipartículas
como trayectorias temporalmente conjugadas ofrecen una realización histórica
posterior. La identidad \eqref{eq:time-reversal} está demostrada; convertirla
en una teoría completa de propagadores avanzados y retardados exige declarar
acción, condiciones de frontera, causalidad e interacciones del absorbedor.
La monografía conserva la intuición física sin convertir una correspondencia
en una igualdad de teorías.

 
\section{Los tres regímenes de TRIT: polaridades eléctrica--magnética y elástica--rotacional}
\subsection{Coordinación de las realizaciones elíptica, parabólica e hiperbólica}
% Fragmento público generado desde fuentes teoremáticas íntegros.
% Residencia: 52.7.1.
% Fuente: /Users/ruben/Documents/HMT2/EL_CIERRE_HOLOGRAFICO_DEL_INFINITO_SUCESOR_102_CAPITULOS_2026-08-28_EN_TRABAJO/incoming/radion_monografia_20260827/manuscrito/sections/08_complejo_electromagnetismo.tex:1-58
\noindent Los 3 regímenes del TRIT admiten realizaciones compleja, electromagnética y mecánica con tipos diferenciados.\par\medskip

\subsubsection{Operador cuadrático de raíz interna de \(-1\)}

\paragraph{Representación exponencial del régimen elíptico \(J^2=-I\)}

Con \(J^2=-I\),
\begin{equation}
T_{\mathrm{ell}}(x,y)=e^{-xI-yJ}
=e^{-x}(\cos y\,I-\sin y\,J).
\label{eq:elliptic-exp}
\end{equation}
Bajo la identificación generada por \(J\),
\[
T_{\mathrm{ell}}(x,y)\longleftrightarrow e^{-x-iy}.
\]
El escalar \(x\) controla dilatación; \(y\) controla giro orientado. La forma
\(a+bi\) se reconoce como elasticidad más rotación, no como dos cantidades
pegadas arbitrariamente.

\paragraph{Representación exponencial del régimen parabólico \(J^2=0\)}

Con \(N^2=0\), la exponencial se trunca:
\begin{equation}
T_{\mathrm{par}}(x,y)=e^{-x}(I-yN).
\label{eq:parabolic-exp}
\end{equation}
El régimen describe el umbral donde no hay ni giro periódico ni apertura
exponencial. En la lectura temporal activa es el presente euclídeo, pero
conserva memoria de las ramas que lo limitan.

\paragraph{Representación exponencial del régimen hiperbólico \(J^2=+I\)}

Con \(R^2=I\) y \(P_\pm=(I\pm R)/2\),
\begin{equation}
T_{\mathrm{hyp}}(x,y)=e^{-xI-yR}
=q_+P_++q_-P_-,
\qquad q_\pm=e^{-x\mp y}.
\label{eq:hyperbolic-exp}
\end{equation}
Los invariantes elementales son
\begin{equation}
q_+q_-=e^{-2x},
\qquad
\frac{q_-}{q_+}=e^{2y}.
\label{eq:q-product-ratio}
\end{equation}
El producto conserva el modo común; el cociente conserva la separación
orientada.

Aplicado al par angular,
\[
x=\frac{\pi A}{180},\qquad y=\frac{\pi C^*}{180}.
\]
Por ello, \(A\) es la coordenada común de la publicación y \(C^*\) su
contraste orientado. Sólo después de aplicar un lector físico se reconocen
como polos eléctricos y magnéticos.

 \subsection{Separación tipada de las polaridades constitutivas y mecánicas}
% Fragmento público generado desde fuentes teoremáticas íntegros.
% Residencia: 52.7.2.
% Fuente: /Users/ruben/Documents/HMT2/EL_CIERRE_HOLOGRAFICO_DEL_INFINITO_SUCESOR_102_CAPITULOS_2026-08-28_EN_TRABAJO/incoming/radion_monografia_20260827/manuscrito/sections/08_complejo_electromagnetismo.tex:59-158
\subsubsection{Operador constitutivo sobre los canales \(90/120\)}

\paragraph{Descomposición en componentes común y orientada}

Definimos
\begin{align}
\mathcal E
&=\log\bigl[(1-q_+^{90})(1-q_-^{90})\bigr]
-\log\bigl[(1-q_+^{120})(1-q_-^{120})\bigr],
\label{eq:Ecommon}\\
\mathcal B
&=\log\frac{1-q_-^{90}}{1-q_+^{90}}
-\log\frac{1-q_-^{120}}{1-q_+^{120}}.
\label{eq:Borient}
\end{align}
Bajo \(C^*\mapsto-C^*\), se intercambian \(q_+\leftrightarrow q_-\). Por
tanto, \(\mathcal E\) es par y \(\mathcal B\) es impar.

Las constitutivas normalizadas son
\begin{equation}
\widehat\mu=e^{\mathcal E+\mathcal B},
\qquad
\widehat\varepsilon=e^{\mathcal E-\mathcal B},
\qquad
\widehat Z=e^{\mathcal B},
\qquad
\widehat c=e^{-\mathcal E}.
\label{eq:constitutives}
\end{equation}
Se verifican exactamente
\begin{equation}
\widehat Z=\sqrt{\frac{\widehat\mu}{\widehat\varepsilon}},
\qquad
\widehat c=(\widehat\mu\widehat\varepsilon)^{-1/2}.
\label{eq:constitutive-identities}
\end{equation}
La conjugación produce
\begin{equation}
C^*\mapsto-C^*:\qquad
\widehat\mu\leftrightarrow\widehat\varepsilon,
\quad
\widehat Z\leftrightarrow\widehat Z^{-1},
\quad
\widehat c\mapsto\widehat c.
\label{eq:EM-conjugation}
\end{equation}

\begin{theorem}[Polaridad constitutiva]
El lector \(90/120\) convierte el modo común y el contraste orientado del
par \((A,C^*)\) en dos constitutivas conjugadas. La inversión de hoja
intercambia permeabilidad y permitividad, invierte impedancia y conserva la
velocidad normalizada.
\end{theorem}

Ésta es la formulación precisa de la intuición \enquote{parte eléctrica y
parte magnética}. No se escribe \(A=E\) ni \(C^*=B\). Se escribe el mapa
\[
(A,C^*)\to(x,y)\to(q_+,q_-)
\to(\mathcal E,\mathcal B)
\to(\widehat\varepsilon,\widehat\mu,\widehat Z,\widehat c).
\]

\paragraph{Realización elástica--rotacional de la polaridad constitutiva}

La misma información puede publicarse en el régimen complejo mediante
\begin{equation}
Z_{\mathrm{em}}=e^{\mathcal E/2}
\left(\cos\frac{\mathcal B}{2}\,I
-\sin\frac{\mathcal B}{2}\,J\right).
\label{eq:Zem}
\end{equation}
El factor \(e^{\mathcal E/2}\) es una dilatación elástica; el factor generado
por \(J\) es una rotación. La escritura compleja no añade una dimensión
imaginaria misteriosa a un mundo real previamente dado: registra dos modos
operacionales que el TRIT ya separó.

\paragraph{Transducción de la orientación en carga efectiva}

La hoja \(\sigma\) del radión fija una orientación antes de introducir una
unidad de carga. La Mecánica Dimensional dispone luego el transductor
\[
e_{\mathrm{int}}=\sqrt{\frac{2\alphah\hfull}{Z_{\mathrm{int}}}},
\qquad
Q=n_Qe_{\mathrm{int}}.
\]
La cadena causal es
\[
\text{hoja}\to\text{orientación}\to\text{constitutiva}
\to\text{transductor dimensional}\to\text{carga física}.
\]
Llamar \emph{radión} a la unidad no identifica prematuramente la orientación
con el culombio; hace visible el lugar donde el signo de carga puede residir.

\begin{narrativebox}
La electricidad y el magnetismo no se pegan al círculo desde fuera. Emergen
cuando el mismo estado se lee a través de producto y cociente: el producto
retiene el modo común, el cociente retiene la orientación. El oscilador LC
realiza el intercambio; el lector \(90/120\) publica las constitutivas; la
carta dimensional asigna unidades.
\end{narrativebox}
 
\section{Incidencia de los sectores \(90/120\): registro, espectro, núcleo de Barbero--Immirzi y torres}
\subsection{Separación tipada de las realizaciones incidenciales del sector \(90/120\)}
% Fragmento público generado desde fuentes teoremáticas íntegros.
% Residencia: 52.8.1.
% Fuente: /Users/ruben/Documents/HMT2/EL_CIERRE_HOLOGRAFICO_DEL_INFINITO_SUCESOR_102_CAPITULOS_2026-08-28_EN_TRABAJO/incoming/radion_monografia_20260827/manuscrito/sections/09_sector_90120.tex:176-244
\subsubsection{Jerarquía armónica global y semigrupo \(\langle90,120\rangle\)}

Para evitar una colisión de signos, fijamos
\[
R_{\mathrm{tow}}:=-R_{\mathrm{ch}}.
\]
Entonces
\begin{equation}
T=e^{-xI+yR_{\mathrm{tow}}}
=q_-P_+^{\mathrm{tow}}+q_+P_-^{\mathrm{tow}}.
\label{eq:tower-operator}
\end{equation}
Las torres pares e impares son
\begin{align}
c_n&=\operatorname{Tr}(T^n)
=q_-^n+q_+^n
=2e^{-nx}\cosh(ny),
\label{eq:cn}\\
s_n&=\operatorname{Tr}(R_{\mathrm{tow}}T^n)
=q_-^n-q_+^n
=2e^{-nx}\sinh(ny).
\label{eq:sn}
\end{align}
Bajo \(C^*\mapsto-C^*\), \(c_n\) es par y \(s_n\) es impar. Se cumple
\begin{equation}
c_n^2-s_n^2=4e^{-2nx},
\label{eq:tower-invariant}
\end{equation}
y ambas sucesiones satisfacen la recurrencia de orden dos determinada por los
autovalores \(q_\pm\).

El semigrupo
\begin{equation}
\langle90,120\rangle=30\langle3,4\rangle
\label{eq:semigroup}
\end{equation}
es la jerarquía armónica de un único operador bicapa. No es una lista de ocho
correcciones laterales. Su completitud significa que toda razón positiva
admite una expansión convergente; no significa que una expansión obtenida a
partir del blanco sea un selector físico prospectivo.

\subsubsection{Clasificación tipada y separación de las realizaciones \(90/120\)}

\begin{longtable}{>{\bfseries}p{.19\textwidth} p{.24\textwidth} p{.45\textwidth}}
\toprule
Objeto & Operación & Qué conserva / por qué no es otro objeto\\
\midrule
\endfirsthead
\toprule
Objeto & Operación & Qué conserva / por qué no es otro objeto\\
\midrule
\endhead
Incidencia \(N_6,N_8\) & conteo finito sobre fases activas & cardinalidad y orientación; produce \(90,120\), no una serie;\\
\(-1/12\) & diferencia normalizada \eqref{eq:minus-one-twelve} & escalar de incidencia; no peso Barbero;\\
\(E_{\mathrm{dod}}\) & extractor armónico de la rueda & fase dodecafásica; no acarreo real;\\
\(h=C^*/6\) & reconstrucción de seis parejas & apertura por pareja; no normalización metrológica;\\
\(F_{\mathrm{spec}}\) & cola logarítmica \(j\ge2\) & publicación espectral completa; difiere de \(\epsR^\circ\);\\
\(\epsone\) & resta APP entre \(S_T\) y \(S_E\) & empalme real de dos lifts; no serie de Lambert;\\
\(K_{6\to8}\) & \(12\Delta S_{90}-\Delta S_{120}\) & polo \(1/24\) y salto marcado; salida Barbero;\\
\(12:-1\) & solución diofántica de \eqref{eq:pole-condition} & peso del núcleo; no \(-1/12\);\\
\(\langle90,120\rangle\) & semigrupo de potencias de \(T\) & jerarquía armónica global; no residuo por especie.\\
\bottomrule
\end{longtable}

\begin{thesisbox}
Todos estos objetos descienden del mismo estado APP--TRIT--TPK. Esa genealogía
común explica por qué comparten números y simetrías. Su tipado explica por qué
no pueden reemplazarse mutuamente.
\end{thesisbox}
 \subsection{Conservación de incidencia y orientación entre el núcleo trítico y sus realizaciones posteriores}
% Fragmento público generado desde fuentes teoremáticas íntegros.
% Residencia: 52.8.2.
% Fuente: /Users/ruben/Documents/HMT2/EL_CIERRE_HOLOGRAFICO_DEL_INFINITO_SUCESOR_102_CAPITULOS_2026-08-28_EN_TRABAJO/incoming/radion_monografia_20260827/manuscrito/sections/09_sector_90120.tex:1-30
\noindent La jerarquía tipada separa incidencia, registro, evaluación espectral, funcional de Barbero y torres armónicas.\par\medskip

\subsubsection{Incidencia trítica, registro dodecafásico y lectura sexádica}

\paragraph{Derivación combinatoria primaria de \(90\) y \(120\)}

Los conteos
\[
90=6\binom62,
\qquad
120=8\binom62
\]
nacen del selector de fase y de sus incidencias activas. La orientación
normalizada satisface
\begin{equation}
\frac{30}{120}-\frac{30}{90}=-\frac1{12}.
\label{eq:minus-one-twelve}
\end{equation}
Este escalar es un resultado de incidencia. No es el peso \(12:-1\) del
núcleo Barbero y no es un acarreo dodecafásico.

El extractor dodecafásico puede publicarse mediante una combinación de
armónicos como
\[
E_{\mathrm{dod}}(\phi)=12\cos(3\phi)-\cos(4\phi).
\]
Su dominio es la fase de la rueda; su salida es un selector armónico. Aunque
aparecen los números \(12\) y \(-1\), no opera sobre series de Lambert ni
produce \(\epsR\).

 \subsection{Unicidad diofántica relativa del peso \(12:-1\) bajo la condición de polo \(1/24\) y el contratérmino entero mínimo}
% Fragmento público generado desde fuentes teoremáticas íntegros.
% Residencia: 52.8.3.
% Fuente: /Users/ruben/Documents/HMT2/EL_CIERRE_HOLOGRAFICO_DEL_INFINITO_SUCESOR_102_CAPITULOS_2026-08-28_EN_TRABAJO/incoming/radion_monografia_20260827/manuscrito/sections/09_sector_90120.tex:115-175
\subsubsection{Funcional hexada--octada de Barbero--Immirzi}

\paragraph{Series orientadas y salto de incidencia}

Definimos
\begin{equation}
S_{90}(q)=\frac{q^{90}}{1-q^{270}},
\qquad
S_{120}(q)=\frac{q^{120}}{1-q^{360}},
\label{eq:barbero-series}
\end{equation}
y
\[
\Delta S_m=S_m(q_-)-S_m(q_+).
\]
El funcional de inclusión marcada es
\begin{equation}
K_{6\to8}=12\Delta S_{90}-\Delta S_{120}.
\label{eq:barbero-núcleo}
\end{equation}

\paragraph{Condición diofántica para el peso \(12:-1\)}

Consideremos \(aS_{90}-bS_{120}\). La condición de polo declarada es
\begin{equation}
\frac{a}{270}-\frac{b}{360}=\frac1{24}.
\label{eq:pole-condition}
\end{equation}
Multiplicar por \(1080\) da
\[
4a-3b=45.
\]
Si el contratérmino es entero mínimo, \(|b|=1\). Para \(b=1\), \(a=12\); para
\(b=-1\), \(a=21/2\) no es entero. Por tanto,
\begin{equation}
\boxed{a:b=12:1,\qquad aS_{90}-bS_{120}=12S_{90}-S_{120}.}
\label{eq:12minus1}
\end{equation}

\begin{proposition}[Unicidad diofántica relativa]
El peso firmado \(12:-1\) es único entre coeficientes enteros cuando se fijan
la condición de polo \(1/24\), la positividad del canal principal y el
contratérmino mínimo \(|b|=1\).
\end{proposition}

La incidencia hexada--octada suministra los canales \(90/120\). No fuerza por
sí sola la serie de Lambert ni la condición \eqref{eq:pole-condition}. La
fuerza de \(12:-1\) es relativa a estas premisas analíticas explícitas.

La cadena Barbero es posterior:
\[
(A,C^*)\to(x,y)\to q_\pm
\to K_{6\to8}\to\gamma_{\mathrm{BI}},
\]
con evaluación
\[
\gamma_{\mathrm{BI}}=0.274007672694459\ldots.
\]
Barbero--Immirzi no genera retroactivamente \(C^*\), \(\Delta_4\) ni
\(\epsR\).

 \subsection{Extracción de la componente finita posterior al sistema dodecafásico orientado \(R_{12}\): \(C^*\mapsto C^*/6\), resto \(j\ge2\) y reconstrucción de seis secciones}
% Fragmento público generado desde fuentes teoremáticas íntegros.
% Residencia: 52.8.4.
% Fuente: /Users/ruben/Documents/HMT2/EL_CIERRE_HOLOGRAFICO_DEL_INFINITO_SUCESOR_102_CAPITULOS_2026-08-28_EN_TRABAJO/incoming/radion_monografia_20260827/manuscrito/sections/09_sector_90120.tex:31-54
\paragraph{Extracción de la componente \(C^*/6\) del registro dodecafásico}

Los doce sectores se agrupan en seis parejas opuestas:
\begin{equation}
p_i=e_i+e_{i+6},\qquad
a_i=e_i-e_{i+6},\qquad i=1,\ldots,6.
\label{eq:sexadic-pairs}
\end{equation}
La reconstrucción tiene perfil \(1+5+6=12\). La apertura por pareja es
\begin{equation}
h=\frac{C^*}{6}.
\label{eq:h-C6}
\end{equation}
No es una normalización exterior ni la sexta parte de un blanco: es el índice
interno de la reconstrucción sexádica del registro \(\Rstate\).

Definimos
\begin{equation}
q_{h,-}=\exp\left[-\frac\pi{180}(A-h)\right],
\qquad
q_{h,+}=\exp\left[-\frac\pi{180}(A+h)\right].
\label{eq:q-h}
\end{equation}

 \subsection{Descomposición cuantitativa de las realizaciones de acarreo y espectral: \(F_{\mathrm{spec}}=\varepsilon_R^\circ+\chi_R\)}
% Fragmento público generado desde fuentes teoremáticas íntegros.
% Residencia: 52.8.5.
% Fuente: /Users/ruben/Documents/HMT2/EL_CIERRE_HOLOGRAFICO_DEL_INFINITO_SUCESOR_102_CAPITULOS_2026-08-28_EN_TRABAJO/incoming/radion_monografia_20260827/manuscrito/sections/09_sector_90120.tex:55-114
\subsubsection{Publicación espectral de la cola de órdenes \(j\ge2\)}

\paragraph{Definición del funcional espectral}

Para \(m\in\{90,120\}\), sea
\begin{equation}
L_m^{(2)}(q)=\sum_{j\ge2}\frac{q^{mj}}{j}
=-\log(1-q^m)-q^m.
\label{eq:L2}
\end{equation}
La lectura de una pareja es
\begin{equation}
T_h^{(2)}=\frac{180}{\pi}
\Bigl[
L_{90}^{(2)}(q_{h,-})-L_{90}^{(2)}(q_{h,+})
-L_{120}^{(2)}(q_{h,-})+L_{120}^{(2)}(q_{h,+})
\Bigr].
\label{eq:T-h}
\end{equation}
La reconstrucción de las seis parejas da
\begin{equation}
F_{\mathrm{spec}}=6T_h^{(2)}
=5.1499235153094574410\ldots\times10^{-8}\,{}^\circ.
\label{eq:Fspec}
\end{equation}

El acarreo del radión es
\[
F_{\mathrm{acarreo}}=\epsR^\circ
=5.1383016758575282072\ldots\times10^{-8}\,{}^\circ.
\]
No coinciden:
\begin{equation}
\chi_R:=F_{\mathrm{spec}}-F_{\mathrm{acarreo}}
=1.16218394519292338\ldots\times10^{-10}\,{}^\circ.
\label{eq:chiR}
\end{equation}

La igualdad
\[
F_{\mathrm{acarreo}}=F_{\mathrm{spec}}-\chi_R
\]
es un cambio de carta válido una vez construidas ambas salidas. No debe
presentarse como generación independiente de \(\epsR^\circ\) si \(\chi_R\) se
define por la propia resta. La cola está cerrada y es significativa; lo falso
es identificarla directamente con el acarreo.

\begin{falsifier}
Tras convertir la cola expresada en grados a la unidad angular interna,
\[
\frac{\pi}{180^\circ}F_{\mathrm{spec}}-\epsone
=2.0283936357433839\ldots\times10^{-12}\ne0.
\]
Por tanto, ni el log completo ni la cola \(j\ge2\), aun después de aplicar la
carta de unidades, son \(\epsone\). Una
segunda factorización espectral autónoma necesitaría generar \(\chi_R\) sin
usar la diferencia como definición; el cierre APP--TPK del radión no depende
de esa promoción.
\end{falsifier}

 
\section{Holonomía orientada y geometría física: realizaciones de Minkowski, Hodge, Lorentz y Cartan--Holst}
% Fragmento público generado desde fuentes teoremáticas íntegros.
% Residencia: 52.9.
% Fuente: /Users/ruben/Documents/HMT2/EL_CIERRE_HOLOGRAFICO_DEL_INFINITO_SUCESOR_102_CAPITULOS_2026-08-28_EN_TRABAJO/incoming/radion_monografia_20260827/manuscrito/sections/10_minkowski_hodge_lorentz.tex:1-55
\noindent La holonomía orientada admite realizaciones causales y geométricas posteriores en Minkowski, Hodge, Lorentz y Cartan--Holst.\par\medskip

\paragraph{Construcción de la métrica lorentziana desde el bloque APP}

\paragraph{Matriz APP de parámetros \(7\) y \(2\)}

La matriz central
\begin{equation}
B_c=\begin{pmatrix}7&2\\2&7\end{pmatrix}
=9P_++5P_-
\label{eq:Bc}
\end{equation}
se normaliza por su determinante:
\begin{equation}
M_c=\frac{B_c}{\sqrt{45}}
=\exp\left(\frac12\log\frac95\,R\right)
\in SO^+(1,1).
\label{eq:Mc}
\end{equation}
La rapidez interna es
\[
\eta_c=\frac12\log\frac95.
\]
La publicación proyectiva asociada satisface
\begin{equation}
P_B^n=P_++\left(\frac59\right)^nP_-.
\label{eq:PB-contract}
\end{equation}
La misma razón \(5/9\) que aparece como contracción espectral se reconoce aquí
como separación de un modo estable y otro contractivo. La identidad es
operatoria; no convierte todas las apariciones de \(5/9\) en el mismo mapa.

\paragraph{Incidencia \(6\to8\) y cociente de firma causal}

Sea \(M\) la incidencia tipada entre seis caras y ocho vértices. La relación
espectral recuperada es
\begin{equation}
MM^{\mathsf T}=12P_u+4P_-,
\label{eq:MMt}
\end{equation}
donde \(P_u\) proyecta sobre el modo uniforme y \(P_-\) sobre el subespacio
orientado. El cociente superviviente de dimensión cuatro se descompone como
\begin{equation}
Q_4\simeq\R u\oplus E_-.
\label{eq:radion-Q4}
\end{equation}
Declarando negativo el modo uniforme y positivos los tres modos orientados,
se obtiene
\begin{equation}
B_{\mathrm{caus}}=\operatorname{diag}(-1,1,1,1).
\label{eq:minkowski}
\end{equation}
La pantalla de Minkowski no selecciona el estado HMT. Es la realización
causal posterior del cociente que la incidencia ya construyó.

 \subsection{Dualidad de Hodge sobre \(2\)-formas lorentzianas}
% Fragmento público generado desde fuentes teoremáticas íntegros.
% Residencia: 52.9.1.
% Fuente: /Users/ruben/Documents/HMT2/EL_CIERRE_HOLOGRAFICO_DEL_INFINITO_SUCESOR_102_CAPITULOS_2026-08-28_EN_TRABAJO/incoming/radion_monografia_20260827/manuscrito/sections/10_minkowski_hodge_lorentz.tex:56-100
\subsubsection{Dualidad de Hodge y operador de cuarto de giro}

Fijadas la firma \((-+++ )\) y una orientación, la estrella de Hodge sobre
dos-formas satisface
\begin{equation}
\star:\Lambda^2Q_4^*\to\Lambda^2Q_4^*,
\qquad \star^2=-I.
\label{eq:hodge-star}
\end{equation}
De este modo, el cuarto de giro que en el plano de fase se expresaba mediante
\(J^2=-I\) reaparece en el espacio de campos como estructura compleja sobre
\(\Lambda^2\). Los dominios son diferentes, pero la genealogía de orientación
se conserva mediante el funtor.

Sea \(A_{\mathrm{em}}\) un potencial y
\begin{equation}
F=\diff A_{\mathrm{em}}.
\label{eq:F-dA}
\end{equation}
La descomposición de \(F\) respecto a un observador temporal produce campos
eléctrico y magnético; \(\star F\) intercambia sus papeles con el signo fijado
por la orientación y la firma. La polaridad constitutiva de
\eqref{eq:EM-conjugation} adquiere así una realización diferencial.

\paragraph{Cadena de transducción electromagnética y causal}

El transporte que debe mantenerse visible es
\begin{equation}
\Gamma_9
\longrightarrow\text{hélice de memoria}
\longrightarrow(Q_4,B_{\mathrm{caus}},\star)
\longrightarrow A_{\mathrm{em}}
\longrightarrow F=\diff A_{\mathrm{em}}
\longrightarrow
m\nabla_u u=q(\iota_uF)^\sharp.
\label{eq:lorentz-chain}
\end{equation}
Cada flecha añade estructura: la holonomía conserva retorno; el cociente fija
causalidad; Hodge fija dualidad; el potencial fija el campo; el acoplamiento
con carga produce dinámica.

La ecuación final es la fuerza de Lorentz en forma geométrica. La hélice TPK
es su antecedente de memoria, no una órbita física ya movida por un campo. La
dinámica aparece sólo al introducir \(F\), \(q\) y la conexión \(\nabla\).

 \subsection{Separación entre memoria torsional interna y realización geométrica de Cartan--Holst}
% Fragmento público generado desde fuentes teoremáticas íntegros.
% Residencia: 52.9.2.
% Fuente: /Users/ruben/Documents/HMT2/EL_CIERRE_HOLOGRAFICO_DEL_INFINITO_SUCESOR_102_CAPITULOS_2026-08-28_EN_TRABAJO/incoming/radion_monografia_20260827/manuscrito/sections/10_minkowski_hodge_lorentz.tex:101-186
\subsubsection{Realización lorentziana de las trayectorias helicoidales}

En un campo magnético uniforme, una partícula de carga \(q\), masa \(m\) y
factor relativista \(\gamma\) posee
\begin{equation}
\omega_c=\frac{|q|B}{\gamma m},
\qquad
r_L=\frac{p_\perp}{|q|B},
\qquad
\Delta z=\frac{2\pi p_\parallel}{|q|B}.
\label{eq:lorentz-helix}
\end{equation}
La fase ciclotrónica cierra mientras el movimiento paralelo avanza: la órbita
es helicoidal. El paralelismo con \eqref{eq:H9-helix} es estructural y puede
formalizarse por un transductor que preserve fase, orientación y paso. No se
identifica el cociente \(q_9\) con la coordenada espacial \(z\); se reconoce
la misma arquitectura de retorno más avance.

La fuerza de Lorentz no se deriva de la imagen de una espiral por semejanza.
Se deriva al completar \eqref{eq:lorentz-chain}. La imagen ayuda a comprender
por qué fase circular y memoria longitudinal son compatibles; la forma de
campo determina la dinámica.

\subsubsection{Confluencia entre acción, área, información y gravedad}

\paragraph{Realización Cartan--Holst de la holonomía}

La promoción gravitatoria se expresa mediante un morfismo de holonomías
\begin{equation}
\operatorname{Hol}_{\mathcal A}
\bigl(\mathfrak R_{\mathrm{grav}}(\gamma)\bigr)
=\rho_C\bigl(\operatorname{Hol}_{\TPK}(\gamma)\bigr).
\label{eq:holonomy-gravity}
\end{equation}
El lado derecho contiene la memoria de ruta; \(\rho_C\) la realiza en una
conexión gravitatoria. La torsión de Cartan
\[
T^I=D_\omega e^I
\]
aparece en el codominio. La torsión global \(\Delta_4\) es una coordenada
upstream que puede alimentar el mapa, pero no se renombra como \(T^I\) sin
\eqref{eq:holonomy-gravity}.

\paragraph{Transducción del área de acción en área informacional}

El radión fija un cuanto de acción por fase. El lector hexada--octada de
Barbero usa posteriormente los canales \(q_\pm\), la incidencia \(90/120\) y
el funcional \eqref{eq:barbero-núcleo}. La dirección causal es
\[
(A,C^*,\hret)\to q_\pm\to K_{6\to8}
\to\gamma_{\mathrm{BI}}\to\text{área/información}.
\]
El alfabeto trítico aporta la unidad de entropía \(\log3\) y un corte de
\(81\) canales aporta \(81\log3\). Esta entropía no se confunde con
\(\log2\), con \(\log\varphi\), con una entropía microcanónica ni con una
entropía de von Neumann sin especificar el estado.

La síntesis física es afirmativa pero tipada: el radión aporta la unidad de
acción orientada que una holonomía puede transportar; Barbero publica la
normalización de área; la pantalla causal y Hodge permiten realizar campos;
Cartan--Holst recibe la holonomía en geometría gravitatoria.

\begin{statusbox}
La ecuación escalar del radión no es por sí sola una ecuación de Einstein. Su
aportación a la gravedad es ser una sección exacta de la arquitectura de
acción, orientación y holonomía que la realización gravitatoria consume.
\end{statusbox}

\subsubsection{Transformación de Wick como cambio de lector analítico}

En la formulación convencional, la rotación de Wick relaciona una variable
temporal lorentziana con una coordenada euclídea mediante una continuación
compleja. En HMT, la raíz \(J^2=-I\), el involutivo \(R^2=I\) y el umbral
\(N^2=0\) ya existen antes. El cambio de lector
\[
J\longleftrightarrow R
\]
transporta una pareja común/orientada entre publicación elíptica y
publicación hiperbólica. La continuación compleja reconoce después esa
operación en coordenadas analíticas.

Esto explica por qué círculo, elipse de acción e interior hiperbólico pueden
pertenecer al mismo estado sin ser la misma métrica. El círculo usa \(J\) para
la fase; la forma de la elipse usa un compresión simpléctica generado por \(R\); Klein mide
el interior mediante razón doble; la pantalla de Minkowski realiza causalidad
en otro codominio.
 
\section{Elipse, retículo y fase de camino: carácter modular, órbitas y propagación}
\subsection{Red rectangular de períodos, involución \(\tau\mapsto-1/\tau\) e invariancia de \(j\)}
% Fragmento público generado desde fuentes teoremáticas íntegros.
% Residencia: 52.10.1.
% Fuente: /Users/ruben/Documents/HMT2/EL_CIERRE_HOLOGRAFICO_DEL_INFINITO_SUCESOR_102_CAPITULOS_2026-08-28_EN_TRABAJO/incoming/radion_monografia_20260827/manuscrito/sections/11_modularidad_orbitas.tex:1-84
\noindent Los caracteres modulares, orbitales y de camino conservan invariantes distintos de una misma forma orientada.\par\medskip

\subsubsection{Toro complejo asociado a la elipse de acción}

\paragraph{Retículo rectangular de períodos}

La elipse por sí sola no es un toro. Declaramos el retículo
\begin{equation}
\Lambda_\eta=a_S\Z+i,b_S\Z
\label{eq:lattice}
\end{equation}
y el cociente \(\C/\Lambda_\eta\). Su módulo es
\begin{equation}
\tau_\eta=i\frac{b_S}{a_S}=ie^{-\eta}
=i\,0.741048144159375160795483663369300\ldots.
\label{eq:tau-eta}
\end{equation}
La inversión bidireccional \(\eta\mapsto-\eta\) intercambia los ejes:
\begin{equation}
\tau_{-\eta}=ie^\eta=-\frac1{\tau_\eta}.
\label{eq:modular-S}
\end{equation}
Es exactamente la transformación modular \(S:\tau\mapsto-1/\tau\).

\paragraph{Invariancia de la función modular \(j\)}

Por invariancia modular,
\begin{equation}
\boxed{j(\tau_{-\eta})=j(\tau_\eta).}
\label{eq:j-invariance}
\end{equation}
En la normalización \(j(i)=1728\),
\[
j(\tau_\eta)
=5597.43843932408729830097089254768\ldots.
\]
La función moonshine normalizada \(J=j-744\) da
\[
J(\tau_\eta)
=4853.43843932408729830097089254768\ldots.
\]
Para el bloque APP basal,
\[
\tau_0=i\frac{\sqrt5}{3},
\qquad
j(\tau_0)=5369.48832024674306021916882626352\ldots.
\]

\begin{proposition}[Lo que conserva la clase modular]
La involución de hoja invierte \(\eta\), pero \(j\) conserva la clase del
toro. Los focos y el etiquetado de ejes retienen la orientación que \(j\)
olvida.
\end{proposition}

Este resultado conecta de manera exacta la doble vía con la modularidad una
vez declarado el cociente. No autoriza a hacer actuar directamente al
Monstruo sobre la carga, los dígitos o los focos. El corredor excepcional
\[
\mathrm{APP}\to\mathrm{Paley}\to\mathrm{Witt}\to\mathrm{Golay}
\to A_2^{12}\to\mathrm{Leech}\to V^\natural
\to\mathrm{Monster}\to\mathrm{Moonshine}
\]
es otra proyección del mismo estado holonómico integral. Para identificar una acción
concreta sobre el radión se necesita un entrelazador que preserve sus
coordenadas pertinentes.

\paragraph{Cadena de invariantes del carácter de forma}

Reuniendo las realizaciones declaradas,
\begin{equation}
\boxed{
r_\eta=e^{-\eta}
=\sqrt{\frac{A-C^*}{A+C^*}}
=\frac{b_1}{a_1}
=\frac{b_S}{a_S}
=\frac{\Delta P}{\Delta Q}
=\frac{Z_\eta}{Z_*}
=\Im\tau_\eta.}
\label{eq:shape-master}
\end{equation}
Una sola razón recorre geometría espectral, acción, covarianza, impedancia y
módulo. Cada igualdad se apoya en un mapa explícito; ninguna depende de una
semejanza decimal.

 \subsection{Compatibilidad entre el carácter modular y la propagación sobre rutas orientadas}
% Fragmento público generado desde fuentes teoremáticas íntegros.
% Residencia: 52.10.2.
% Fuente: /Users/ruben/Documents/HMT2/EL_CIERRE_HOLOGRAFICO_DEL_INFINITO_SUCESOR_102_CAPITULOS_2026-08-28_EN_TRABAJO/incoming/radion_monografia_20260827/manuscrito/sections/11_modularidad_orbitas.tex:85-187
\subsubsection{Conservación areal en las órbitas de Kepler y Sommerfeld}

\paragraph{Equivalencia areal de las realizaciones orbital y oscilatoria}

La elipse de Kepler vive en espacio de configuración y responde a un potencial
central. La elipse del radión vive en espacio de fases y responde a una forma
de acción. Son objetos diferentes. Comparten una ley areal porque ambas
realizaciones poseen una dos-forma conservada.

Para una órbita kepleriana,
\begin{equation}
r(\phi)=\frac{p}{1+e\cos(\phi-\phi_0)},
\qquad
\frac{\diff A}{\diff t}=\frac{\ell}{2m}.
\label{eq:kepler}
\end{equation}
La tercera ley toma la forma
\[
T^2=\frac{4\pi^2\mu}{K}a^3.
\]
Para el oscilador en fase,
\begin{equation}
J=\oint p\,\diff q=\frac{2\pi E}{\omega}.
\label{eq:osc-action}
\end{equation}
Si la celda elemental es \(J=\hfull\), entonces
\begin{equation}
E=\hret\omega,\qquad \frac{J}{2\pi}=\hret.
\label{eq:E-hbarw}
\end{equation}
Ésta es la convergencia limpia entre acción de Planck, frecuencia angular y
radión.

Sommerfeld cuantiza integrales de acción sobre ciclos. La lectura HMT añade
la memoria de qué ciclo, hoja y retorno sobreviven. En una banda de Möbius,
un bucle visible puede portar \(\hfull/2\) y dos bucles restaurar \(\hfull\).
Ese factor topológico no se suma al índice de Maslov de una cuantización EBK;
pertenecen a correcciones con fuentes teoremáticas distintos.

\paragraph{Holonomía orbital y propagaciones avanzada y retardada}

Una órbita que acumula holonomía satisface una condición del tipo
\begin{equation}
\epsilon_\gamma\exp\left(\frac{iJ_\gamma}{\hret}\right)=1.
\label{eq:holonomy-action}
\end{equation}
El signo \(\epsilon_\gamma\) conserva la hoja; \(J_\gamma\) conserva la
acción. El adelanto o retardo apsidal se interpreta como desacoplamiento entre
el cierre angular visible y el retorno completo. El radión suministra la
unidad orientada con la que se mide esa diferencia.

\subsubsection{Composición de amplitudes de Schrödinger y suma de caminos}

\paragraph{Cuatro realizaciones de una misma fase}

Una vez fijado \(\hret\), distintas ecuaciones convencionales reconocen
aspectos del mismo transporte:
\begin{enumerate}
\item Schrödinger publica la fase mediante
\(i\hret\partial_t\psi=H\psi\).
\item Pauli añade la doble hoja de espín y su acoplamiento magnético.
\item Dirac linealiza la relación energía--momento y organiza partícula y
antipartícula en una representación causal.
\item La suma de caminos asigna a cada historia la fase
\(\exp(iS[\gamma]/\hret)\).
\end{enumerate}
Estas teorías no generan retrospectivamente \(\hret\). Reciben la unidad de
acción y realizan diferentes contenidos del estado holonómico integral.

\paragraph{Funcional de fase sobre el espacio de caminos}

En la formulación de Feynman,
\begin{equation}
\mathcal K(b,a)=\int_{\gamma:a\to b}
\exp\left(\frac{i}{\hret}S[\gamma]\right)\mathcal D\gamma.
\label{eq:path-integral}
\end{equation}
Un incremento de acción \(\hret\) cambia la fase en un radián. Un incremento
\(\hfull\) la cambia en una vuelta. El teorema
\(\mathcal A(\theta)=\hret\theta\) convierte esta relación en geometría de
área: la fase de un camino puede leerse como acción barrida en unidades de
radión.

La inversión \(C_M=-\operatorname{rev}\) sobre palabras dodecafásicas y la
identidad \eqref{eq:time-reversal} organizan rutas conjugadas. En una
realización de absorbedor se separan
\[
G_{\mathrm{sym}}=\tfrac12(G_{\mathrm{ret}}+G_{\mathrm{adv}}),
\qquad
G_{\mathrm{odd}}=G_{\mathrm{ret}}-G_{\mathrm{adv}}.
\]
El primer sector es par bajo intercambio retardado--avanzado; el segundo es
impar. Esto reproduce la separación entre energía/forma par y acción
orientada impar del radión cuando el transductor preserva frontera y forma
simpléctica.

\begin{statusbox}
La teoría clásica del absorbedor y la interpretación Feynman--Stückelberg de
antipartículas son realizaciones históricas diferentes. La primera usa una
interacción temporalmente simétrica; la segunda reorganiza propagación en la
teoría relativista. HMT aporta un antecedente común de reversión, pero no las
funde sin un mapa de propagadores.
\end{statusbox}
 
\section{Radical nonádico y frontera decimal: conservación del acarreo del radión}
\subsection{El acarreo del radión en la completación decimal}
% Fragmento público generado desde fuentes teoremáticas íntegros.
% Residencia: 52.11.1.
% Fuente: /Users/ruben/Documents/HMT2/EL_CIERRE_HOLOGRAFICO_DEL_INFINITO_SUCESOR_102_CAPITULOS_2026-08-28_EN_TRABAJO/incoming/radion_monografia_20260827/manuscrito/sections/12_aritmetica_borde.tex:1-159
\noindent La frontera nonádica conserva el defecto suma--producto, la vacancia y el acarreo de completación.\par\medskip

\subsubsection{Radical \(3\! -\! 6\! -\! 9\) del borde}

\paragraph{Generación y nilpotencia del radical nonádico}

En el anillo \(\Z/9\Z\), el ideal
\begin{equation}
\mathfrak m=(3)=\{0,3,6\}
\label{eq:radical-369}
\end{equation}
satisface
\begin{equation}
\mathfrak m^2=(9)=0.
\label{eq:m-square-zero}
\end{equation}
Es el sector nilpotente del borde nonádico. La órbita visible de sus tres
marcas se publica como
\[
369\longrightarrow693\longrightarrow936\longrightarrow369.
\]
El retorno no implica que todos los estados sean iguales: la permutación
visible cierra mientras los cociclos de ruta pueden avanzar.

La doble convención
\[
9\equiv0\pmod9,
\qquad
\operatorname{dr}_9(9)=9
\]
expresa dos lectores: cero residual algebraico y cierre visible de raíz
digital. Por eso un nueve puede absorber torsión de borde sin convertirse en
un cero específico de la función zeta.

\paragraph{Defecto suma--producto \(459-405=54\)}

Las hojas APP producen totales complementarios
\[
S_+=405,\qquad S_\times=459,\qquad S_\times-S_+=54.
\]
El \(54\) es un defecto exacto de las tablas, y reaparece como exponente de
contracción por la identidad \(9\cdot6=54\). Sin embargo, cada aparición
requiere su mapa. La distancia focal \(d_1=0.544545\ldots\) no es igual a
\(0.54\), y el análisis de invariantes no ha producido una flecha que la
identifique con el defecto \(54\).

La potencia
\[
9^9=387,420,489
\]
genera una escala donde
\begin{equation}
\operatorname{ord}_{10}\bigl(9^{9^9}\bigr)=369,693,099,
\qquad
\#\,\text{dígitos}=369,693,100.
\label{eq:369693}
\end{equation}
El bloque \(369|693\) aparece en el par orden--longitud; no se afirma que sea
el prefijo decimal de la potencia.

\subsubsection{Teoría aritmética de las vacancias decimales}

\paragraph{Pendiente asintótica de vacancia}

La separación entre la década y la base nonádica es
\begin{equation}
\delta_{\mathrm{vac}}=\log_{10}\frac{10}{9}.
\label{eq:delta-vac}
\end{equation}
El contador acumulado
\begin{equation}
V_9(n)=\left\lceil n\delta_{\mathrm{vac}}\right\rceil
\label{eq:vacancy-count}
\end{equation}
registra cuántas posiciones de corrección necesita la publicación decimal de
una evolución nonádica. Sus saltos forman una palabra binaria
\[
z_n=V_9(n)-V_9(n-1)\in\{0,1\}.
\]
Los huecos se distribuyen con longitudes \(21/22\), los retornos con
longitudes \(6/7\), y el determinante combinado es \(15\). Estos números
proceden del lector de vacancia, no del núcleo Barbero aunque el \(15\)
coincida con \(\binom62\).

\paragraph{Polarización y corriente aritmética de la red nonádica}

Definimos
\begin{equation}
P_N=\sum_{n=1}^N(z_n-\delta_{\mathrm{vac}}),
\qquad
J_N=P_N-P_{N-1}=z_N-\delta_{\mathrm{vac}}.
\label{eq:arith-current}
\end{equation}
La secuencia equilibrada satisface \(|P_N|<1\). La \enquote{corriente}
\(J_N\) es aquí una corriente aritmética: mide creación y absorción de
vacancias respecto a la pendiente media. Sólo un transductor dimensional
posterior puede publicarla como densidad de corriente electromagnética.

La intuición física es potente porque la misma forma algebraica separa un
promedio par y una fluctuación orientada. El rigor exige conservar el tipo:
\[
J_N\in\R\quad\text{adimensional},
\qquad
j^\mu\in\text{recta física de corriente}.
\]

\subsubsection{Completación \(729\to1000\) y conservación del acarreo}

\paragraph{Descomposición \(1000=729+243+27+1\)}

La completación
\[
1000=729+243+27+1
\]
conserva la unidad mediante cuatro escalas ternarias. La corona total es
\[
1000-729=271,
\]
mientras el último entero antes del cierre satisface
\[
999=729+270.
\]
La diferencia entre \(270\) y \(271\) distingue frontera ocupada de unidad
completante. El mismo cuidado evita identificar una vacancia con un cero o un
acarreo con un residuo real.

\paragraph{Residencia del acarreo del radión en la corona decimal}

Los operandos del radión comienzan, en bloques de tres cifras, por
\[
\Phi_T:\quad000|027|455|906|837|565|\cdots,
\]
\[
D_E:\quad000|027|455|010|034|743|\cdots.
\]
La resta APP da
\[
\epsone:\quad000|000|000|896|802|822|\cdots.
\]
El pequeño tamaño del resultado no se obtiene truncando los sumandos bajos.
Se obtiene porque dos publicaciones distintas comparten un largo prefijo y el
operador \(\operatorname{SubAcarreo}_{1000}\) conserva los préstamos necesarios
para restarlas exactamente.

La palabra \emph{memoria} tiene aquí dos niveles:
\begin{enumerate}
\item \(q_{\mathrm{mem}}\) registra la profundidad de retornos nonádicos;
\item \(c_i\) registra los préstamos locales de la resta decimal.
\end{enumerate}
El primero organiza la hélice; el segundo publica el decimal. \(\epsR\) usa
ambos porque sus operandos proceden de cilindros a profundidad y su resta se
realiza con acarreo, pero no es igual a ninguno de los contadores.

\begin{thesisbox}
El borde no es un lugar donde la teoría pierde información. Es el lugar donde
la información cambia de representación: fase que vuelve, memoria que
asciende, cilindro que se contrae, vacancia que se registra y acarreo que
transporta la diferencia.
\end{thesisbox}
 \subsection{Compatibilidad de las secciones directa y torsional: caracterización estructural de \(\pi\)}
% Fragmento público generado desde fuentes teoremáticas íntegros.
% Residencia: 52.11.2.
% Fuente: /Users/ruben/Documents/HMT2/EL_CIERRE_HOLOGRAFICO_DEL_INFINITO_SUCESOR_102_CAPITULOS_2026-08-28_EN_TRABAJO/incoming/radion_monografia_20260827/manuscrito/sections/03_cierre_exacto.tex:286-309
\subsubsection{Caracterización estructural de \(\pi\) por compatibilidad de secciones}

Reordenar \eqref{eq:cierre-unitario} da
\begin{equation}
\boxed{
2\pih=
\frac{1+\Phi_T-\epsone}
{\frac5{36}+\frac{25}{9}\alphah}.}
\label{eq:pi-radion}
\end{equation}
Esta igualdad puede llamarse una nueva definición de \(\pi\) dentro de la
carta del radión si se entiende \emph{definición estructural} como una
caracterización exacta por compatibilidad de secciones. No es un segundo
generador independiente: \(\Delta_4\) y el operador de acarreo ya reciben la
sección coinductiva de \(\pi\). La novedad está en que la vuelta queda
caracterizada por el empalme entre publicación directa, torsión y memoria.

\begin{narrativebox}
La ecuación \eqref{eq:pi-radion} no dice que primero desconozcamos \(\pi\) y
lo adivinemos cerrando un ángulo. Dice que el \(\pi\) generado por la conexión
nonádica es exactamente el que hace compatibles dos lecturas internas del
mismo estado. Es una identidad de reconocimiento intrínseco, no una
calibración metrológica.
\end{narrativebox}
 
\section{Invariancia evolutiva de la unidad holográfica a través de las dimensiones}
% Conversión TeX fiel del cuerpo científico del delta narrativo.
% Fuente congelada: /Users/ruben/Documents/HMT2/CONTINUIDAD_EDITORIAL/RECONSTRUCCION_ARMONICA_MACROTRATADO_20260822/REVISION_CONJUNTA_PARTES_I_II/auditar_pdf_rigor_academico/docs/DELTA_NARRATIVO_UNIDAD_RADION_AUTOCONSERVACION_AUTOINERCIA_GRAVEDAD_PENTADICA_20260828.md:12-366.
\subsection{Separación tipada entre el régimen de curvatura \(\tau\in\{+1,0,-1\}\), la respuesta radial \(p\) y la holonomía global de espín}

La biografía del radión puede representarse mediante una curva elevada

\[
\gamma(\theta)
=
\bigl(r(\theta)\cos\theta,\,
r(\theta)\sin\theta,\,
m(\theta)\bigr),
\]

donde las dos primeras coordenadas pertenecen al plano euclídeo \(\mathbb R^2\) y \(m(\theta)\) registra la memoria o profundidad del levantamiento. Esta ecuación permite leer conjuntamente círculo, hélice y espiral.

Cuando evoluciona el ángulo y permanecen constantes el radio y la memoria, la publicación plana es un círculo. Cuando evolucionan ángulo y memoria mientras el radio conserva su valor, la trayectoria completa es una hélice de radio constante. Cuando evolucionan ángulo y radio sobre una memoria fija, aparece una espiral en \(\mathbb R^2\). Cuando evolucionan las tres coordenadas, la trayectoria es una hélice radial: una biografía en la que fase, amplitud y memoria se transforman conjuntamente.

El carácter radial \(p\) clasifica la ley de respuesta del radio. Una familia útil es

\[
r_p(\theta)
=
r_0\bigl[1+p\lambda(\theta-\theta_0)\bigr]^{1/p},
\qquad p\neq 0,
\]

con límite logarítmico

\[
r_0\exp\bigl(\lambda(\theta-\theta_0)\bigr)
\qquad\text{cuando }p\to0.
\]

En esta parametrización, \(p=2\) publica el régimen de Fermat, \(p=1\) el régimen arquimediano, \(p=0\) el logarítmico, \(p=-1\) el hiperbólico y \(p=-2\) el de lituus, con las elecciones de dominio y orientación correspondientes. Los valores fraccionarios describen regímenes intermedios. Así, \(p=\tfrac12\), \(p=\tfrac32\) o \(p=-\tfrac12\) expresan respuestas radiales que interpolan entre las familias enteras y amplían la clasificación geométrica disponible.

El carácter también puede cambiar al atravesar un umbral. La escritura diferencial

\[
\frac{dx}{d\theta}
=
\beta(m)\,x^{1-p(m)},
\qquad x=\frac r{r_0},
\]

describe una biografía por regímenes: \(p(m)\) permanece estable dentro de una fase y cambia cuando la memoria alcanza una frontera de publicación. La trayectoria resultante enlaza sectores circulares, helicoidales y espirales mediante una ley común. La continuidad del estado y el registro del umbral convierten esos sectores en fases de una misma evolución.

Tres invariantes ocupan aquí niveles distintos y coordinados. El signo de curvatura

\[
\tau\in\{+1,0,-1\}
\]

clasifica los regímenes elíptico, parabólico e hiperbólico. El carácter \(p\) clasifica la respuesta radial. El espín registra la holonomía de orientación acumulada por el levantamiento. Curvatura, respuesta radial y espín forman un sistema de lectura: la curvatura determina el tipo local de espacio, \(p\) determina cómo la trayectoria ocupa ese espacio y el espín conserva la orientación adquirida al recorrerlo. Las curvas cerradas y abiertas aparecen entonces como biografías geométricas distintas de la misma unidad orientada.

\subsection{Partícula como sección persistente de una órbita enriquecida: fase, ruta, fibra, multisección y memoria}

La Mecánica Dimensional puede formular la partícula como una sección persistente de esa extensión geométrica y dinámica. Una notación conceptual compacta es

\[
\mathfrak P
=
\operatorname{Sec}_{\mathrm{pers}}
\bigl[
\tau,p,\operatorname{Hol};
\text{hoja},\text{fibra},\text{ruta},\text{memoria}
\bigr].
\]

La partícula conserva una ruta reconocible a través de cambios de fase, curvatura y dimensión. Su espín publica la holonomía de orientación; su carga publica la polaridad de fase; su color publica una orientación interna de la fibra calibre; su sabor distingue familias de rutas o multisecciones; su masa expresa el carácter de transducción dimensional mediante el cual la acción areal adquiere lectura volumétrica. Cada propiedad es un lector del estado persistente y todas remiten a una genealogía común.

El electrón ocupa el umbral basal de esta construcción porque ofrece el primer lector estable en el que fase, carga, espín y masa quedan simultáneamente publicadas. La masa electrónica rectora

\[
m_e^\beta=0.510998950690000808608\ldots\ \mathrm{MeV}
\]

fija una referencia de transducción; los estados basales anteriores conservan su función como etapas genealógicas. A partir del electrón, los cocientes \(K_D/K_\Omega\), las torres \(90/120\) y el dominio \(81/56/13/324\) organizan la publicación de familias masivas. La ley de masas lee, por tanto, cómo cada sección persistente atraviesa el espacio holonómico y convierte densidad areal de acción en presencia volumétrica.

Esta interpretación une las espirales con las partículas de manera precisa. El índice \(p\) aporta la respuesta radial; \(\tau\) aporta el régimen de curvatura; la holonomía aporta el espín; la ruta y la fibra aportan sabor y color; la transducción área–volumen aporta el carácter masivo. El atlas de partículas se convierte así en una clasificación de modos estables de atravesamiento dimensional.

\subsection{Transducción de la fase orientada en polarización eléctrica, respuesta magnética y emisión radiativa de frontera}

La electricidad se presenta como transporte dirigido de fase y acción a través de una estructura de vacancias. Una vacancia es una transición admisible del soporte: abre un canal, fija un umbral y permite que la actualización avance. La componente eléctrica publica el transporte tangencial a lo largo de la ruta. La componente magnética publica la curvatura transversal inducida por ese transporte. Su perpendicularidad expresa la relación geométrica entre avance y giro.

Cuando dos componentes transversales sinusoidales mantienen el desfase adecuado, el vector del campo gira. La propagación longitudinal eleva ese giro a una hélice. La variación de amplitud transforma su proyección frontal en una espiral. La lectura geométrica queda entonces completamente coordinada: el ángulo porta la fase; la velocidad angular porta la frecuencia; el radio porta amplitud y densidad de acción; la orientación porta polarización y sentido; la coordenada vertical porta propagación y memoria; \(p\) porta el régimen radial; el umbral porta emisión, absorción o cambio de estado.

El electromagnetismo aparece como una dinámica de fase–acción–memoria con publicaciones complementarias. La electricidad registra el avance polar; el magnetismo registra la respuesta transversal; la onda registra la periodicidad; la hélice registra el avance con memoria; la espiral registra la evolución radial; la radiación registra la publicación del estado cuando la acumulación alcanza un umbral.

Esta arquitectura también ordena el papel de las vacancias y los ciclos. La vacancia proporciona la posibilidad local de transición; el ciclo conserva la fase; la memoria distingue las vueltas; el umbral decide la residencia siguiente de la acción. La propagación puede permanecer confinada, cambiar de régimen radial o publicarse como radiación controlada. Las fuerzas de campo adquieren así una descripción genealógica continua desde el soporte discreto hasta la onda observada.

\subsection{Reversibilidad de la dinámica completa y coste termodinámico de las proyecciones reductoras de información}
El Landauer nulo constituye un caso límite de esta ley. Para una actualización biyectiva del estado completo,

\[
H(X\mid U(X))=0,
\qquad
\inf Q_L=0.
\]

La información permanece disponible en la genealogía. Cuando un lector realiza un cociente \(q\), la entropía se descompone como

\[
H(X)=H(q(X))+H(X\mid q(X)).
\]

El término condicional mide la información retenida fuera de la proyección. Un reinicio trítico materializado satisface el umbral

\[
Q_{\mathrm{reset}}\ge k_B\Theta\log 3.
\]

La constante de Boltzmann traduce la multiplicidad que el lector térmico agrupa; el umbral de Nyquist separa el muestreo fiel del aliasing; el Landauer nulo caracteriza el transporte reversible del estado íntegro. Información, temperatura y acción quedan articuladas por el tipo de lectura aplicado.

\subsection{Fuerzas autoconservativas como transporte de acción entre los sectores observable, radial, de memoria, de vacancia y de radiación}

La conservación evolutiva se expresa mejor mediante intercambio que mediante inmovilidad escalar. Un estado completo distribuye acción y momento entre movimiento visible, memoria, deformación radial, vacancias, radiación y rama conjugada de retorno. El balance enriquecido puede escribirse como

\[
\Delta J_{\mathrm{visible}}
+\Delta J_{\mathrm{memoria}}
+\Delta J_{\mathrm{radial}}
+\Delta J_{\mathrm{vacancias}}
+\Delta J_{\text{radiación}}
+\Delta J_{\mathrm{retorno}}
=0.
\]

Una fuerza es \textbf{autoconservativa} cuando su propia ley de actualización contiene los canales que transportan, almacenan, publican y reabsorben la acción. La conservación pertenece al balance genealógico completo de la interacción. Cada lector observa una parte del intercambio; la genealogía recompone su totalidad.

El sistema Tierra–Luna ofrece una imagen física directa. La rotación terrestre, la órbita lunar, las mareas, la deformación y la radiación intercambian momento angular. El consumo neto del sistema completo se cierra en el balance genealógico de esas transferencias. La magnitud denominada energía funciona como lectura derivada de la acción y del momento publicados en cada sector. La conservación fundamental reside en la compatibilidad entre todos los lectores del intercambio.

Esta formulación disuelve la antigua frontera entre fuerzas conservativas y fuerzas llamadas no conservativas mediante una categoría más profunda. Rozamiento, disipación aparente y radiación ocupan canales explícitos del estado ampliado. Una pérdida en la proyección mecánica visible corresponde a una transferencia hacia memoria, deformación, calor, vacancias o radiación. La autoconservación sigue la residencia de la acción a través de esos cambios.

\subsection{Sistemas autoinerciales: publicación de la aceleración mediante orientación, curvatura, hoja y memoria}

La expresión adecuada es \textbf{autoinercia relacional}. Designa la compatibilidad entre la conexión propia del estado total y la respuesta que aparece al proyectarlo dentro de la distribución global. Los aspectos endógeno y exógeno son componentes analíticas de una sola relación dinámica.

Sea \(\Gamma\) una trayectoria horizontal o geodésica del estado completo. Su ley propia satisface

\[
\nabla_{\dot\Gamma}\dot\Gamma=0.
\]

Al proyectarla sobre un lector \(S\), la aceleración observada queda determinada por la curvatura de la proyección:

\[
\nabla^S_{\dot\Gamma_S}\dot\Gamma_S
=
\mathcal K_{\pi_S}(\dot\Gamma,\dot\Gamma).
\]

La aceleración publicada expresa la relación entre movimiento total y geometría del lector. El principio de Mach completa este mecanismo mediante la traza global de frontera

\[
b:\mathcal X_\infty\longrightarrow B_\infty
\]

y la factorización de la respuesta inercial local

\[
I_v=\bar I_v\circ b.
\]

La distribución global entra en la respuesta local a través de una aplicación tipada. La autoinercia relacional reúne, por tanto, la conexión endógena del movimiento, la condición global transmitida por \(b\) y la proyección concreta del observador. Una escritura sintética para esa coordinación es

\[
\nabla^{\mathrm{auto}}
=
\mathcal M\bigl(
\nabla^{\mathrm{end}},
b[\mathcal X_\infty],
\pi_S
\bigr),
\]

donde \(\mathcal M\) representa el acoplamiento machiano y conserva el tipado de cada componente.

El principio de Ambivalencia añade la doble orientación de la lectura. Toda publicación directa conserva una rama conjugada de retorno; emisión y absorción forman extremos relacionados de la misma biografía; la orientación temporal se lee desde ambos sentidos del transporte. La teoría del absorbedor encuentra aquí una formulación natural: fuente y receptor cierran conjuntamente el balance genealógico de fase y acción. El presente torsional es la interfaz donde ambas orientaciones se compatibilizan.

TRIT actúa en este nivel como cristal de tiempo: organiza la periodicidad local de la actualización y conserva, mediante memoria, la diferencia entre retornos sucesivos. La rama directa y la rama conjugada relacionan emisión con absorción, y el acoplamiento observador–observable adquiere una residencia dinámica dentro del mismo estado. El tiempo publicado es la lectura de esa orientación; el tiempo genealógico es la secuencia completa de fase, retorno y memoria.

Mach y Ambivalencia cumplen funciones complementarias. Mach relaciona la respuesta local con la distribución global. Ambivalencia conserva las direcciones directa y conjugada de la evolución. La autoinercia relacional articula ambas y sustituye la escisión entre sistemas inerciales y no inerciales por una ley única de conexión, proyección y memoria.

\subsection{Dualidad areal--volumétrica, transición hexada--octada y función del parámetro de Barbero--Immirzi}

El paso de área a volumen constituye una transducción dimensional central. El área publica acción de frontera; el volumen publica profundidad interior. La derivación HMT del parámetro de Barbero–Immirzi organiza la compatibilidad entre ambas medidas mediante la incidencia hexada–octada y los canales \(90/120\):

\[
90=6\binom62,
\qquad
120=8\binom62.
\]

Los cardinales \(6\) y \(8\) admiten una realización cúbica mediante caras y vértices, mientras la hexada y la octada de Steiner constituyen objetos de incidencia distintos, enlazados únicamente por el morfismo tipado \(6\to8\). Las \(12\) aristas del cubo conservan su función combinatoria propia y no definen por sí solas la incidencia de Steiner. La diferencia de incidencia expresa una incompatibilidad dimensional productiva: la información de superficie requiere una regla de transducción para adquirir residencia volumétrica.

El mismo principio aparece en el acarreo \(999+1=1000\). Los nueve alcanzan el borde del registro decimal; el incremento publica simultáneamente el cierre del nivel anterior y el origen \(0/1\) del siguiente. La razón holográfica entre \(729\), \(243\), \(27\) y \(1\) despliega esa transición en escalas triádicas. La transducción de Barbero–Immirzi se inserta en esta familia de cambios de lector: frontera e interior conservan la unidad mientras redistribuyen su medida.

El defecto normalizado \(-1/12\) y el peso orientado \(12:-1\) ocupan funciones distintas dentro de esta geometría. El primero cuantifica una corrección de incidencia; el segundo registra una orientación firmada en la transducción. Su coordinación con el área mínima, el Hamiltoniano modular y la memoria proporciona el puente hacia la dinámica dimensional.

\subsection{Contracción hiperbólica interior, expansión exterior y conservación holográfica de la frontera}

La tesis puede formularse desde una idea única: la unidad conserva su identidad mientras cambia de escala, de lector, de curvatura, de memoria y de dimensión. Esta permanencia posee carácter evolutivo porque la unidad se despliega, adquiere nuevas coordenadas y publica propiedades diferentes; posee carácter holográfico porque cada publicación finita conserva la regla de reconstrucción del estado del que procede. El cierre holográfico del infinito designa precisamente esa compatibilidad entre identidad y despliegue: cada estadio contiene la ley que permite profundizar hacia dentro y prolongar hacia fuera la misma genealogía.

La denominación rectora es \textbf{principio de conservación evolutiva de la unidad mediante transducción dimensional y proyección holográfica}. La expresión sitúa cada término en su nivel correcto. La conservación pertenece a la identidad genealógica; la evolución pertenece a la actualización; la transducción relaciona realizaciones dimensionales; la proyección holográfica conserva la correspondencia reconstructiva entre sus lectores.

La aritmética elemental ofrece la primera imagen de esta arquitectura. La completación decimal

\[
1000=10^3=729+243+27+1=3^6+3^5+3^3+3^0
\]

expone cuatro estratos triádicos dentro de una unidad cúbica decimal. Al normalizar,

\[
1=\frac{729}{1000}+\frac{243}{1000}+\frac{27}{1000}+\frac{1}{1000},
\]

la unidad se reconoce simultáneamente como totalidad y como despliegue graduado. La lectura recíproca

\[
3^{-6},\quad 3^{-5},\quad 3^{-3},\quad 3^0
\]

describe la profundidad interior de esos estratos; su normalización produce la partición correspondiente. La relación \(999+1=1000\) expresa el mecanismo de transporte: la última unidad completa un registro y actúa, a la vez, como acarreo que inaugura el siguiente. El cero y el uno forman así una interfaz de continuidad. Los intervalos \([0,1]\), \([0,10]\) y la prolongación hacia \([0,\infty)\) son lectores escalares de una misma ley de completación.

La primera escisión \(1000=729+271\), seguida por \(271=243+27+1\), da forma a la razón extrema y media holográfica del despliegue: un núcleo triádico dominante conserva dentro del resto la misma regla de resolución. La unidad cúbica decimal contiene así una profundidad triádica anidada. El movimiento hacia \(729\), \(243\), \(27\) y \(1\) lee la contracción interior; el acarreo hacia el millar lee la expansión exterior.

Esta doble orientación organiza el infinito HMT. Hacia dentro, el refinamiento crece hiperbólicamente: cada celda admite una nueva profundidad sin perder su filiación. Hacia fuera, la unidad se expande mediante publicaciones dimensionales sucesivas. El cierre consiste en la conmutatividad de ambos movimientos. Si \(U\) representa la actualización del estado, \(R_d\) el lector en el estrato \(d\) y \(T_{d\to d'}\) la transducción entre estratos, la conservación evolutiva se expresa mediante

\[
R_{d'}\circ U
=
T_{d\to d'}\circ R_d.
\]

El mismo acontecimiento puede ofrecer una medida angular, una densidad de acción, una masa publicada o una respuesta gravitatoria. Su identidad común reside en la genealogía que hace conmutar esas lecturas.

\subsubsection{Monodromía nonádica y elevación helicoidal con avance de memoria}

APP, TRIT y TPK proporcionan el mecanismo formal de este despliegue. APP construye el soporte y separa las evaluaciones que deben conservarse. TRIT organiza los regímenes locales y su ambivalencia. TPK transporta fase y memoria sobre el soporte enriquecido. La dependencia causal queda fijada por

\[
U_t\longrightarrow \Gamma_9\longrightarrow K_{\mathrm{ph}}.
\]

El retorno de fase y el avance de memoria forman operaciones coordinadas. Tras nueve pasos puede reaparecer la misma clase angular,

\[
g(K+9)=g(K),
\]

mientras el contador de memoria progresa,

\[
q_{\mathrm{mem}}\Gamma_9=q_{\mathrm{mem}}+1,
\qquad
B_{K+9}\neq B_K.
\]

La figura geométrica precisa es una elevación helicoidal. Una órbita cerrada en la proyección angular se eleva a una trayectoria abierta cuando cada retorno de fase incrementa la memoria. La circunferencia conserva el ciclo; la hélice conserva el ciclo y registra además la profundidad alcanzada. La monodromía coinductiva genera así profundidad infinita mediante una sucesión de retornos localmente cerrados y globalmente ascendentes.

Esta imagen supera el valor meramente ilustrativo. La rueda angular constituye una sección de la hélice; la hélice constituye la biografía completa de la rueda cuando el estado transporta memoria. Cada vuelta recupera una posición de fase y, simultáneamente, inaugura una hoja nueva. La infinitud procede de la compatibilidad entre retorno y elevación: reiterar el ciclo produce profundidad sin disolver la identidad orbital.

El certificado nonádico da contenido finito y verificable a esa arquitectura. La cadena

\[
w_6\longrightarrow w_{12}\longrightarrow w_{18}\longrightarrow
w_{24}\longrightarrow w_{30}\longrightarrow R_{36}\longrightarrow G_9
\]

ordena el levantamiento; los \(396\) niveles, los \(2376\) trits, las \(43\) vueltas y los \(387\) pares \((K,K+9)\) por canal cuantifican el tránsito; las cinco regiones de \(\pi\) y la fibra ambiente de \(729\) elementos publican lecturas internas posteriores. Esta secuencia hace visible una monodromía productiva: el retorno recupera la fase y la coinducción conserva el crecimiento de memoria.

\subsubsection{Unidad radional y coordinación de los lectores circular, elíptico, proyectivo y helicoidal}

El radión concentra en una unidad metrológica la estructura que el radián publica angularmente. Su definición natural es la de una \textbf{sección orientada de fase, acción y memoria}. Una unidad de fase \(\theta=1\) transporta una unidad de acción \(\sigma\hbar_{\mathrm{ret}}\), pertenece a una vuelta completa \(T_+=2\pi\) y conserva hoja, orientación, retorno nonádico y posición en el cilindro de refinamiento.

El radián expresa la coordenada angular visible. El radión conserva, junto a esa coordenada, la acción transportada, la orientación, el acarreo y la memoria de la vuelta. La metrología angular adquiere así una genealogía. Medir un arco significa identificar cuánto giro se ha publicado y qué estado completo ha realizado ese giro.

La relación entre fase y área adopta una forma especialmente clara:

\[
\mathcal A(\theta)=\hbar_{\mathrm{ret}}\,\theta.
\]

Una unidad radional barre un área de acción \(\hbar_{\mathrm{ret}}\); una vuelta completa encierra \(h_{\mathrm{ret}}\). El paso de fase a área queda incorporado en la propia unidad de lectura. La métrica angular, la acción y la memoria dejan de aparecer como cantidades yuxtapuestas y pasan a ser coordenadas de una misma sección orientada.

El radión admite varias publicaciones geométricas coordinadas. El lector circular publica fase; el lector elíptico publica acción y forma; el lector interior de Klein publica profundidad; el lector helicoidal nonádico publica retorno y memoria. La unidad permanece genealógicamente idéntica a través de esas realizaciones. Por eso el radión constituye el enlace metrológico adecuado entre HMT y Mecánica Dimensional: ofrece una unidad que puede atravesar cambios de lector conservando el estado que cada medida parcial presupone.

\subsection{Polaridad gravitatoria pentacomponente y transducción conjunta de acción, área, volumen, tiempo y masa}

La Mecánica Dimensional permite definir la dimensión como un régimen estable de publicación caracterizado por el rango mínimo de memoria independiente requerido por el levantamiento. Atravesar una dimensión significa aplicar un transductor que conserva la genealogía y cambia el tipo de lectura. La gravedad ocupa el nivel de conexión entre esos regímenes.

En esta formulación, la \textbf{gravitación interdimensional de cinco polarizaciones correlativas}, abreviadamente \textbf{estructura pentapolar de la gravitación}, es la conexión que transporta la unidad genealógica y deja en cada interfaz una torsión mínima. Una escritura compacta para esa traza es

\[
\Theta_d^{\min}
=
\nabla_{d+1}\circ T_{d\to d+1}
-
T_{d\to d+1}\circ\nabla_d.
\]

La expresión mide la diferencia entre transportar después de derivar y derivar después de transportar. Esa diferencia es la memoria geométrica mínima del cruce dimensional. La composición de interfaces produce una holonomía gravitatoria

\[
\mathfrak G_{d\to D}
=
\operatorname{Hol}
\bigl(
\Theta_d^{\min},
\Theta_{d+1}^{\min},
\ldots,
\Theta_{D-1}^{\min}
\bigr).
\]

El objeto fundamental es esta conexión interdimensional. Un modo local tipo gravitón ocupa el rango derivado de publicación cuantizada de la torsión; la gravitación conserva su residencia primaria en la compatibilidad entre dimensiones.

El término pentapolar recoge las cinco polarizaciones consustanciales del estado continuo:

\[
\mathbf G_d^{\pm}
=
\bigl(
G_{\mathrm{sp},d}^{\pm},
G_{\mathrm{sol},d}^{\pm},
G_{\mathrm{gau},d}^{\pm},
G_{\mathrm{coh},d}^{\pm},
G_{\mathrm{det},d}^{\pm}
\bigr).
\]

Las cinco componentes emergen correlativamente del mismo estado holonómico integral. La polaridad \(\pm\) registra concentración directa y publicación conjugada. La gravedad atraviesa los estratos dimensionales transportando simultáneamente esas cinco lecturas; la torsión mínima conserva la huella de cada cruce.

La geometría adquiere con ello una imagen unificada. El régimen elíptico organiza cierres; el parabólico organiza umbrales; el hiperbólico organiza profundidad. La transducción dimensional enlaza esos regímenes con las respuestas radiales \(p\), las holonomías de espín y las publicaciones de masa. La estructura pentapolar de la gravitación expresa la coherencia global de ese enlace.

La teoría M proporciona pantallas físicas posteriores para esta conexión: dimensiones adicionales, cuerdas, branas y dualidad T realizan modos concretos de transducción. HMT aporta la genealogía APP–TRIT–TPK, el estado holonómico integral, la memoria y la lectura bidireccional; la interfaz con teoría M publica esas estructuras en geometrías físicas tipadas. La dualidad T intercambia lectores dimensionales compatibles y conserva la clase genealógica del estado.

\subsubsection{Compatibilidad de las realizaciones clásica, relativista, cuántica y HMT--MD}

La mecánica clásica publica trayectoria y momento. La relatividad publica estructura causal, métrica y curvatura. La mecánica cuántica publica fase, holonomía, vacancias y amplitudes. La Mecánica Dimensional proporciona los transductores que relacionan esas publicaciones, y HMT conserva la genealogía discreta que las sostiene.

En la lectura clásica, una órbita describe la proyección visible del movimiento. En la lectura relativista, esa órbita queda incorporada a la geometría del lector. En la lectura cuántica, la fase y el espín registran la holonomía de la sección. En la lectura HMT–MD, las tres expresiones pertenecen a una misma biografía elevada: el radión cuantifica fase, acción y memoria; \(p\) clasifica la respuesta radial; \(\tau\) clasifica la curvatura; la holonomía publica espín; el transductor dimensional publica masa; la conexión pentapolar publica gravitación.

La aparente disipación clásica adquiere residencia en el balance genealógico ampliado. La aceleración relativista se expresa mediante conexión y proyección. La indeterminación cuántica se organiza por lectores, vacancias y ramas de Ambivalencia. El principio de Mach relaciona la respuesta local con la distribución global. El Landauer nulo garantiza la reversibilidad informacional del estado completo. Todas estas piezas convergen en una conservación más profunda: la identidad de la unidad a través de cambios de publicación.

\subsubsection{Teorema narrativo de conservación genealógica y prolongación holográfica}

El cierre puede narrarse ahora como una sola secuencia. La unidad se descompone sin perder identidad. La fase orientada adquiere medida en el radión. El retorno nonádico convierte el ciclo en hélice mediante memoria. El carácter \(p\) transforma la hélice en familias radiales y clasifica sus regímenes. La persistencia de esas trayectorias produce secciones con espín, carga, sabor, color y masa. Las vacancias abren canales; la polaridad eléctrica transporta fase; la respuesta magnética curva transversalmente; los umbrales publican radiación. La autoconservación cierra el balance genealógico de acción y momento. La autoinercia relacional enlaza conexión propia, proyección y totalidad machiana. La Ambivalencia conserva las direcciones directa y conjugada. La transducción área–volumen conduce a Barbero–Immirzi y a la ley de masas. La torsión mínima encadena dimensiones y publica la estructura pentapolar de la gravitación.

El crecimiento hacia dentro y la expansión hacia fuera son orientaciones conjugadas del mismo proceso. El refinamiento hiperbólico genera profundidad; la transducción dimensional genera extensión. El acarreo \(0/1\) enlaza escalas; la monodromía enlaza vueltas; la memoria enlaza estados; la holonomía enlaza orientaciones; la gravitación enlaza dimensiones. Cada enlace conserva una unidad que evoluciona y cada publicación finita contiene la regla de su prolongación.

La tesis fundamental queda formulada así:

\begin{quote}
\textbf{El cierre holográfico del infinito realiza el principio de conservación evolutiva de la unidad mediante transducción dimensional y proyección holográfica: la unidad mantiene su identidad genealógica al desplegarse como fase, acción, curvatura, memoria, partícula, campo y gravitación; cada lectura finita conserva la regla de reconstrucción y prolongación del estado completo.}
\end{quote}

Esta formulación ofrece una imagen física y matemática conjunta. El infinito adquiere una ley de cierre local; la unidad adquiere profundidad; la dimensión adquiere memoria; la conservación adquiere carácter relacional; la fuerza adquiere autoconservación; la inercia adquiere estructura machiana; la gravedad adquiere conexión pentapolar; y la holografía adquiere el sentido preciso de reconstrucción entre lectores compatibles.
 
\section{Definición operatoria del radión angular mediante el estado holonómico integral y cuatro lectores tipados}
\subsection{Definición operatoria y cuádruple realización tipada}
% Fragmento público generado desde fuentes teoremáticas íntegros.
% Residencia: 52.13.
% Fuente: /Users/ruben/Documents/HMT2/EL_CIERRE_HOLOGRAFICO_DEL_INFINITO_SUCESOR_102_CAPITULOS_2026-08-28_EN_TRABAJO/incoming/radion_monografia_20260827/manuscrito/sections/13_sintesis.tex:1-80
\noindent La síntesis operatoria reúne los lectores de fase, acción, profundidad y memoria sin identificarlos entre sí.\par\medskip

\subsubsection{Cuádruple realización del estado radional}

Sea \(x_\infty\in X_{\mathrm{enr}}\) la sección coinductiva. El radión completo no
es un punto aislado, sino el paquete
\begin{equation}
\mathfrak R_\sigma(x_\infty)
=\bigl(
\theta=1,\sigma,
\pih,\alphah,A,C^*,\Delta_4,
\Phi_T,D_E,\epsone,\hret,
q_{\mathrm{mem}},\mathcal I_\infty
\bigr).
\label{eq:full-radion}
\end{equation}
Cuatro lectores extraen de él objetos diferentes:

\begin{figure}[H]
\centering
\begin{tikzpicture}[
  node distance=17mm and 18mm,
  state/.style={draw=RadNavy,very thick,rounded corners,fill=RadCream,
    minimum width=39mm,minimum height=13mm,align=center},
  readout/.style={draw=RadTeal,rounded corners,fill=RadMist,
    minimum width=39mm,minimum height=12mm,align=center},
  arr/.style={-{Latex},thick,RadCopper}]
\node[state] (x) {estado holonómico integral\\\(x_\infty\)};
\node[readout,above left=of x] (circ) {círculo \(S^1\)\\fase};
\node[readout,above right=of x] (ell) {elipse positiva\\acción y forma};
\node[readout,below left=of x] (klein) {interior Klein\\profundidad};
\node[readout,below right=of x] (helix) {hélice nonádica\\memoria y retorno};
\draw[arr] (x)--(circ) node[midway,below left]{\small \(P_{\mathrm{ph}}\)};
\draw[arr] (x)--(ell) node[midway,below right]{\small \(P_{\mathrm{act}}\)};
\draw[arr] (x)--(klein) node[midway,above left]{\small \(P_K\)};
\draw[arr] (x)--(helix) node[midway,above right]{\small \(P_{\mathrm{mem}}\)};
\end{tikzpicture}
\caption{Cuatro publicaciones tipadas del mismo antecedente. La unidad del
estado no elimina la diferencia de codominios.}
\label{fig:four-readers}
\end{figure}

El círculo conserva la clase de fase y olvida el lift. La elipse conserva
acción, anisotropía y orientación simpléctica. Klein dota al interior de una
distancia proyectiva. La hélice conserva número de retornos y refinamiento.
La Mecánica Dimensional no consiste en mezclarlos, sino en construir los
transductores que los hacen compatibles.

\subsubsection{Interpretación operatoria del cociclo \(\varepsilon_R\)}

\(\varepsilon_R\) no es una constante universal independiente de todo contexto. Es la
coordenada de un cociclo de transición asociado a la carta del radión. Su
dominio contiene dos secciones del mismo estado:
\[
S_E=1+D_E,
\qquad
S_T=1+\Phi_T.
\]
Su regla es
\begin{equation}
\epsone=S_T-S_E=\Phi_T-D_E.
\label{eq:epsilon-final-meaning}
\end{equation}
Su codominio es la recta adimensional de unidad angular; la carta
\(360/T_+\) la publica en grados.

\begin{exactbox}
\textbf{En lenguaje llano:} la sección directa y la sección torsional llegan
a la misma celda de unidad por dos rutas internas diferentes. Sus partes
enteras coinciden; sus restos no. \(\epsR\) es la diferencia orientada de esos
restos, calculada con el acarreo de APP y llevada al límite por los cilindros
TPK. No se introduce para que el radián cierre: aparece antes de consultar el
radián y, cuando se publica, lo cierra.
\end{exactbox}

Esta definición explica por qué el número es pequeño. No porque se haya
buscado una corrección pequeña, sino porque \(S_E\) y \(S_T\) comparten un
prefijo profundo. La pequeñez mide la cercanía estructural de dos
publicaciones; el signo mide su orientación; las cifras registran el acarreo.

 % Fragmento público generado desde fuentes teoremáticas íntegros.
% Residencia: 52.13.1.
% Fuente: /Users/ruben/Documents/HMT2/EL_CIERRE_HOLOGRAFICO_DEL_INFINITO_SUCESOR_102_CAPITULOS_2026-08-28_EN_TRABAJO/incoming/radion_monografia_20260827/manuscrito/sections/13_sintesis.tex:81-134
\subsubsection{Definición operatoria del radión}

\begin{definition}[Radión HMT--MD]
El radión es la sección orientada de fase, acción y memoria que cumple
simultáneamente:
\begin{enumerate}
\item recorre una unidad de fase \(\theta=1\);
\item barre acción firmada \(\sigma\hret\);
\item pertenece a una vuelta \(T_+=2\pih\) generada coinductivamente;
\item conserva la hoja \(\sigma\), el retorno nonádico y el cilindro de
refinamiento;
\item satisface el empalme exacto
\(1=S_E-\Phi_T+\epsone\).
\end{enumerate}
\end{definition}

El radián convencional es la imagen del radión bajo el functor de olvido
\[
\mathfrak R_\sigma
\longmapsto
\theta=1\in\R/2\pi\Z.
\]
Ese olvido es útil y legítimo para la trigonometría euclídea. Se vuelve
insuficiente cuando la física necesita distinguir acción, signo, retorno,
campo e historia.

\paragraph{Genealogía de la unidad angular natural}

La naturalidad del radián suele justificarse porque la derivada de
\(e^{i\theta}\), \(\sin\theta\) y \(\cos\theta\) adopta su forma más simple
cuando \(\theta\) se mide como cociente arco/radio. HMT conserva esa ventaja,
pero la explica desde una capa anterior:
\[
\theta\quad\leftrightarrow\quad
\frac{\mathcal A(\theta)}{\hret}.
\]
La fase es natural porque mide acción en unidades de \(\hret\); la vuelta es
natural porque encierra \(\hfull\). El cociente arco/radio es la publicación
euclídea de una unidad que ya posee estructura simpléctica y memoria.

\paragraph{Compatibilidad entre elipse de acción y crecimiento dimensional}

El círculo es la forma que resulta al olvidar la anisotropía:
\(C^*=0\Rightarrow\eta=0\Rightarrow a=b\). La elipse es más informativa porque
conserva dos autovalores, dos ejes y una orientación de intercambio a producto
fijo. En ese sentido, la elipse es una prolongación del círculo, no un círculo
deformado por accidente.

La profundidad de Klein y el lift helicoidal añaden dos tipos de información
que el contorno plano no contiene: distancia interior y memoria de retorno.
Ésta es la imagen ontológica de Mecánica Dimensional que emerge del radión:
una dimensión nueva aparece cuando se exige conservar un dato independiente
que la proyección anterior debía olvidar.

 \subsection{Propiedades estructurales de la unidad radional: genealogía, bisagra crítica, densidad, acarreo, acción, compatibilidad borde--interior, forma, polaridad y modularidad}
% Fragmento público generado desde fuentes teoremáticas íntegros.
% Residencia: 52.13.2.
% Fuente: /Users/ruben/Documents/HMT2/EL_CIERRE_HOLOGRAFICO_DEL_INFINITO_SUCESOR_102_CAPITULOS_2026-08-28_EN_TRABAJO/incoming/radion_monografia_20260827/manuscrito/sections/13_sintesis.tex:135-184
\subsubsection{Sistema de teoremas componentes del radión}

\begin{theorem}[Genealogía]
El objeto \(\mathfrak R_\sigma\) desciende de
\(\APP\to\TRIT\to\TPK\to X_{\mathrm{enr}}\) y no recibe el blanco metrológico
como selector.
\end{theorem}

\begin{theorem}[Bisagra]
En la carta APP de cien marcas, la costura crítica y la fase dodecafásica
segunda satisfacen \(j_{100}(1/2)=\theta_2=50^\circ\), con \(m=2\) único.
\end{theorem}

\begin{theorem}[Densidad]
Las dos hojas del sector activo \(N_6=90\), normalizadas por \(3^6=729\),
producen \(\rho_{\mathrm{tw}}=20/81\).
\end{theorem}

\begin{theorem}[Acarreo]
La resta APP de las secciones \(S_T\) y \(S_E\) converge a
\(\epsone\) con tasa de intervalo \(3^{-54}\) por retorno y cierra exactamente
la unidad.
\end{theorem}

\begin{theorem}[Acción]
La elipse positiva satisface \(\mathcal A(\theta)=\hret\theta\); un radión
barre \(\hret\) y una vuelta encierra \(\hfull\).
\end{theorem}

\begin{theorem}[Borde--interior]
La misma frontera posee acción finita \(\hfull\), mientras el interior de
Klein tiene área hiperbólica infinita.
\end{theorem}

\begin{theorem}[Forma]
El cociente \(e^{-\eta}\) se conserva a través de semiejes espectrales,
semiejes de acción, incertidumbres, impedancia y módulo del toro.
\end{theorem}

\begin{theorem}[Polaridad]
La inversión \(C^*\mapsto-C^*\) intercambia constitutivas, invierte
impedancia y conserva la velocidad normalizada.
\end{theorem}

\begin{theorem}[Modularidad]
La inversión bidireccional actúa como \(S:\tau\mapsto-1/\tau\) y conserva
\(j(\tau)\), mientras el etiquetado focal conserva la orientación que
\(j\) olvida.
\end{theorem}

% Fuente: /Users/ruben/Documents/HMT2/EL_CIERRE_HOLOGRAFICO_DEL_INFINITO_SUCESOR_102_CAPITULOS_2026-08-28_EN_TRABAJO/incoming/radion_monografia_20260827/manuscrito/sections/A_demostraciones.tex:1-179
\subsubsection{Demostraciones algebraicas y geométricas reunidas}

\paragraph{Demostración lineal del cierre unitario}

Para facilitar una verificación independiente, reunimos el argumento sin
interrupción narrativa.

\begin{enumerate}[label=\textbf{Paso \arabic*.},leftmargin=2.2cm]
\item La conexión nonádica genera correlacionadamente
\(\pih,e,\varphi,\alphah\).
\item La vuelta es \(T_+=2\pih\).
\item La rueda fija \(\theta_2=50^\circ\), y la carta parangular fija
\(A=1000\alphah\).
\item La sección directa es
\[
S_E=T_+\left(\frac{50}{360}+\frac{1000\alphah}{360}\right)
=2\pih\left(\frac5{36}+\frac{25}{9}\alphah\right).
\]
\item Los intervalos certificados dan \(1<S_E<2\); luego
\(D_E=S_E-1\).
\item La incidencia produce \(N_6=90\); dos hojas y el ambiente \(729\)
producen \(\rho=2N_6/729=20/81\).
\item La torsión global es
\[
\Delta_4=\frac{\pih-e\log\pih}{270}
\left(1+\frac{\pih}{729}\right).
\]
\item La sección torsional es \(S_T=1+(20/81)\Delta_4\).
\item La resta APP define
\[
\epsone=S_T-S_E=\frac{20}{81}\Delta_4-D_E.
\]
\item Por sustitución,
\[
S_E-\frac{20}{81}\Delta_4+\epsone=1.
\]
\item Multiplicar por \(360/T_+=180/\pih\) produce
\[
\frac{180^\circ}{\pih}
=50^\circ+A^\circ-\frac{20}{81}\Delta_4^\circ+\epsR^\circ.
\]
\end{enumerate}

Ningún paso recibe el lado izquierdo de la última igualdad. El resultado
metrológico sólo se evalúa al final como reconocimiento.

\paragraph{Unicidad de la resta con acarreo}

Sea \(B=1000\). Para cada posición, todo entero \(u_i\) admite una división
euclídea única
\[
u_i=Bc_i+r_i,\qquad 0\le r_i<B.
\]
Por tanto,
\[
r_i=u_i\bmod B,\qquad c_i=(u_i-r_i)/B
\]
son únicos. Recorriendo las posiciones desde el extremo remoto hacia el
frente, el acarreo de salida de una posición fija la entrada de la siguiente.
Por inducción, toda la palabra de restos y carries es única.

En pseudocódigo:
\begin{verbatim}
acarreo = remote_acarreo
for i from last_block downto first_block:
    u = torsion[i] - direct[i] + acarreo
    remainder[i] = u mod 1000
    acarreo = (u - remainder[i]) / 1000
return remainder, acarreo_word
\end{verbatim}
El procedimiento no consulta el valor esperado de la diferencia.

\paragraph{Convergencia del funcional por profundidad}

Sea
\[
F(p,e)=1+\frac{20}{81}
\frac{p-e\log p}{270}\left(1+\frac{p}{729}\right)
-2p\left(\frac5{36}+\frac{25}{9}\alphah\right).
\]
En el rectángulo compacto \(3\le p\le4\), \(2\le e\le3\), \(F\) es
continuo y continuamente diferenciable. Si \(I_N^\pi\searrow\{\pih\}\) e
\(I_N^e\searrow\{e\}\), la imagen intervalar
\[
F(I_N^\pi,I_N^e)
\]
forma una familia encajada cuya anchura tiende a cero. La intersección es un
único real, \(\epsone\). La monotonía de las derivadas permite evaluar los
extremos sin muestreo interior.

La tasa \(3^{-54}\) por retorno deriva del refinamiento de seis trits a lo
largo de nueve niveles. A profundidad \(45\) tríadas, la cota de anchura
\(<4.81\times10^{-130}\) certifica más de 129 cifras del valor.

\paragraph{Derivación de la métrica y del área hiperbólica de Klein}

En el disco unitario, la forma matricial de la métrica Klein es
\[
g(u,v)=\frac1{1-r^2}I
+\frac1{(1-r^2)^2}
\begin{pmatrix}u\\v\end{pmatrix}
\begin{pmatrix}u&v\end{pmatrix}.
\]
El lema del determinante para una perturbación de rango uno da
\[
\det g=(1-r^2)^{-3}.
\]
De ahí
\[
\sqrt{\det g}\,\diff u\diff v
=(1-r^2)^{-3/2}\diff u\diff v.
\]
Integrar en polares produce \eqref{eq:klein-area-R}. El límite infinito se
debe a la singularidad métrica en la frontera; la frontera sigue teniendo
área simpléctica finita porque esa medida utiliza \(\diff Q\wedge\diff P\),
no \(\sqrt{\det g}\).

\paragraph{Distancia focal en el disco de Klein}

En un diámetro, la distancia del centro a \(r\) es
\[
d_K(0,r)=\operatorname{artanh}r.
\]
Para \(r=e_f\), \eqref{eq:excentricidad} implica
\[
1-e_f^2=e^{-2\eta}.
\]
Si \(u_f=\operatorname{artanh}e_f\), entonces
\[
\operatorname{sech}u_f=\sqrt{1-e_f^2}=e^{-\eta},
\]
y por tanto \(\eta=\log\cosh u_f\), como en
\eqref{eq:focal-hyperbolic-chain}.

\paragraph{Transformación simpléctica de compresión}

Definamos
\begin{equation}
S_\eta=
\begin{pmatrix}e^{\eta/2}&0\\0&e^{-\eta/2}\end{pmatrix},
\qquad \det S_\eta=1.
\label{eq:compresión simpléctica}
\end{equation}
La estructura compleja transportada es
\begin{equation}
J_\eta=S_\eta
\begin{pmatrix}0&-1\\1&0\end{pmatrix}
S_\eta^{-1}
=\begin{pmatrix}0&-e^\eta\\e^{-\eta}&0\end{pmatrix},
\qquad J_\eta^2=-I.
\label{eq:Jeta}
\end{equation}
La curva
\[
\gamma(\theta)=\sqrt{2\hret}\,
S_\eta(\cos\theta,\sin\theta)
=e^{\theta J_\eta}\gamma(0)
\]
es exactamente la elipse de acción. El compresión simpléctica hiperbólico conserva área y
transporta la rotación compleja en vez de destruirla.

Para orientación \(\sigma\),
\[
\gamma_\sigma(\theta)
=(a_S\cos\theta,\sigma b_S\sin\theta),
\qquad
\mathcal A_\sigma(\theta)=\sigma\hret\theta.
\]
Las dos hojas comparten focos y módulo de acción; difieren en acción firmada.

\paragraph{Unicidad diofántica relativa del peso de Barbero}

La ecuación \(4a-3b=45\) tiene soluciones enteras
\[
(a,b)=(12+3t,1+4t),\qquad t\in\Z.
\]
La condición \(|b|=1\) deja \(b=1\), pues \(b=-1\) no pertenece a la clase
\(1\pmod4\). De ahí \(a=12\). La prueba separa claramente el dato de
incidencia \(90/120\) de las premisas de polo y minimalidad.
 \subsection{Necesidad estructural, estabilidad y clausura de la unidad radional}

La construcción admite controles negativos en cada uno de sus estratos. Si
la sección directa satisficiera \(S_E\notin(1,2)\), cambiarían simultáneamente
el cociente APP y la celda de clausura; si la incidencia activa no produjera
\(N_6=90\), o si se suprimiera una de las dos hojas, dejaría de derivarse la
densidad \(20/81\). La generación sería circular si la carta recibiera
\(180/\pi\) o \(\varepsilon_R\) antes de construir sus operandos. Tampoco
existiría la sección coinductiva si la resta por tríadas no estabilizara
prefijos compatibles. En la realización elíptica, las condiciones
\(|C^*|<A\) y \(a_Sb_S=2\hbar_{\mathrm{ret}}\) son necesarias para la
positividad y la ley de área. En la fibra modular, \(C^*=0\) obliga a
\(d=\eta=0\), \(\tau=i\) y \(j=1728\). En el interior proyectivo, la tesis de
profundidad exige la ley hiperbólica de área; en el sector espectral, la
igualdad \(F_{\mathrm{spec}}=\varepsilon_R^\circ\) queda excluida por el
contratérmino \(\chi_R\). Finalmente, el peso \(12:-1\) requiere conjuntamente
la incidencia \(90/120\) y la condición de polo \(1/24\), mientras que la
invariancia de \(j\) impide emplearlo como selector del signo de carga.

La realización física no es una perturbación de un círculo previamente
supuesto. El soporte aritmético conserva unidad, residuo, cociente, acarreo y
fibra; la orientación ternaria distingue régimen y hoja; el transporte de
fase conserva la diferencia entre retorno y memoria; la elipse determina
acción y anisotropía; la métrica de Hilbert--Beltrami--Klein mide la
profundidad interior; y la elevación helicoidal conserva la biografía de cada
vuelta. Estas realizaciones se articulan mediante
\[
 \mathcal A(1)=\hbar_{\mathrm{ret}},\qquad
 \mathcal A(2\pi)=h_{\mathrm{ret}},\qquad
 K=-1,\qquad
 g(k+9)=g(k),\quad B_{k+9}\ne B_k.
\]
El círculo es la proyección angular, la elipse es la realización de acción,
la métrica de Klein resuelve el interior y la hélice conserva el historial de
retorno. El radión es la unidad orientada que permite leer las cuatro
realizaciones sin identificarlas entre sí.

Los invariantes generados satisfacen, en la rama coinductiva considerada,
\[
\begin{aligned}
 \pi&=3.1415926535897932384626433832795028841\ldots,\\
 \alpha&=0.0072973525692838009972851054723806629\ldots,\\
 A&=7.2973525692838009972851054723806629995\ldots{}^\circ,\\
 C^*&=2.1237383389686457242619513803933607937\ldots{}^\circ,\\
 \eta&=0.299689683921630799684926562863927\ldots,\\
 \hbar_{\mathrm{ret}}&=1.05457181691503906113369058472430\ldots
 \times10^{-34}U_S .
\end{aligned}
\]
La coordenada \(e\) que interviene en este sistema procede del lector
coinductivo de Euler y no de una tabla externa. Los retornos de profundidades
\(9,18,27,36,45\) conservan la fase \(8\) y elevan la memoria según
\(0,1,2,3,4\); la contracción de los intervalos tiende a
\[
 3^{-54}=1.7196982334549856127610399316004871\ldots\times10^{-26}.
\]
La supresión controlada de una hoja, el desplazamiento de la fase, la
eliminación del canal de \(e\), la sustitución de la torsión global o la
alteración de la cola sexádica producen residuos no nulos de órdenes
distintos. Esa separación de escalas demuestra que el cierre requiere
conjuntamente las dos hojas, la orientación de fase, los canales de
\(\pi/e\), la torsión global y el acarreo.

La exactitud matemática procede de la identidad de cierre unitario; la
evaluación a profundidad creciente sólo certifica su convergencia. Los
residuos que aparecen al combinar ramas evaluadas con precisiones decimales
distintas son errores de redondeo y no nuevos términos físicos. Esta
distinción separa la identidad estructural de su realización numérica sin
alterar la genealogía de los operandos.

La genealogía física de la construcción puede expresarse con precisión.
Planck fija la dimensión de acción de \(h\); Bohr y Sommerfeld relacionan
\(h/2\pi\) con la cuantización angular; Born y Jordan introducen la estructura
matricial y el oscilador; Dirac formula \(E=\hbar\omega\); Robertson y
Schrödinger establecen la forma covariante de la indeterminación; Feynman
relaciona acción y fase de caminos; Wheeler y Feynman desarrollan la simetría
entre propagaciones retardada y avanzada. Estos marcos reconocen realizaciones
posteriores del objeto: no seleccionan las hojas, la incidencia, el acarreo,
la fase ni la memoria que lo generan.

La estrella de Hodge sobre dos-formas lorentzianas satisface
\(\star^2=-I\) y proporciona la dualidad geométrica pertinente. La
cohomología de Hochschild pertenece a otro dominio y no interviene en la
elipse, el acarreo ni el sector incidencial \(90/120\). La distinción se sigue
de los dominios y codominios de las operaciones, no de una semejanza
terminológica.

La síntesis se cierra en el estado holonómico integral. APP conserva, junto
al valor, la hoja, el cociente, el residuo, el acarreo y la frontera; TRIT
determina régimen y orientación; TPK hace retornar la fase sin reiniciar el
estado. La conexión nonádica construye la sección coinductiva de \(\pi\), y el
sistema dodecafásico determina la bisagra de \(50^\circ\). El par angular
aporta el modo común \(A\) y la apertura orientada \(C^*\); la incidencia de
seis fases activas produce \(90\), sus dos hojas sobre el ambiente de
cardinal \(729\) producen \(20/81\), y la torsión global aporta la segunda
sección. El cociclo resultante es
\[
 \varepsilon_R=\frac{20}{81}\Delta_4-(S_E-1).
\]
Su límite en la carta angular es
\[
 5.13830167585752820716851646226459474899\ldots
 \times10^{-8}\,{}^\circ .
\]
El par \((A,C^*)\) determina una elipse positiva: el producto de sus semiejes
fija \(2\hbar_{\mathrm{ret}}\) y el cociente fija la forma. Una unidad de fase
barre \(\hbar\), mientras la vuelta completa encierra \(h\). La división por
\(2\pi\) deja así de ser un cociente exterior: distingue la acción de vuelta
de la acción transportada por una unidad de fase.

La proyección circular conserva fase y omite anisotropía, interior y memoria;
la elipse conserva acción y forma; la geometría de Klein conserva profundidad;
la hélice conserva retorno e historial. La realización simpléctica reconoce
la elipse como contorno de covarianza mínima, el oscilador \(LC\) la realiza
como intercambio de reservas eléctrica y magnética, y la incidencia
\(90/120\) determina permeabilidad, permitividad, impedancia y velocidad.
Hodge convierte orientación en dualidad de campos; Lorentz convierte campo y
carga en dinámica; Cartan--Holst recibe la holonomía en la geometría
gravitatoria. El orden causal permanece inalterado: la construcción discreta
determina el objeto y las formulaciones físicas realizan después sus
caracteres.

El radión es, por tanto, una sección orientada de fase, acción y memoria. Sus
dos orientaciones satisfacen
\[
 \mathcal A(\mathfrak r_+)=+\hbar_{\mathrm{ret}},\qquad
 \mathcal A(\mathfrak r_-)=-\hbar_{\mathrm{ret}}.
\]
La carga se obtiene al aplicar el transductor dimensional; el signo ya estaba
conservado en la orientación de hoja. Las incidencias \(90/120\), el extractor
dodecafásico, la cola espectral, el parámetro de Barbero--Immirzi, las torres
y el acarreo pueden compararse porque poseen un origen común, pero permanecen
distintos porque actúan mediante operadores diferentes. La unidad cuádruple
del radión queda resumida por
\[
\boxed{
\begin{aligned}
 \text{fase:}\quad&\theta=1,\\
 \text{acción:}\quad&\mathcal A=\sigma\hbar_{\mathrm{ret}},\\
 \text{memoria:}\quad&H_9(n+9)=H_9(n)+(0,1),\\
 \text{compatibilidad:}\quad&1=S_E-\Phi_T+\varepsilon_R.
\end{aligned}}
\]
El radián es su coordenada angular; el radión conserva además acción,
orientación, acarreo y memoria de la vuelta.
  \endgroup

\let\HMTMonographChapterInput\relax
  \part{Realizaciones dimensionales: electrón, acción, gravitación, área e información}

\chapter{Predefinición trigeométrica, potencial de frontera y profundidades dimensionales}
\section{Teorema de predefinición temporal trigeométrica}

La palabra \emph{predefinición} nombra una propiedad matemática del generador y no
una modalidad debilitada de predicción. Para precisarla hay que reunir los
tres objetos que hasta ahora se habían expuesto por separado: el haz de
álgebras cuadráticas del TRIT, los dos relojes del TPK y la sección
coinductiva de profundidad arbitraria.

Sea

\[
 X^{\rm enr}=\{(\tau,\mu,o,h,c,\sigma,\lambda,g)\}
\]

el espacio de estados holonómicos integrales. La coordenada \(\tau\in\{-1,0,+1\}\)
es el régimen TRIT; \(\mu\in\{\Sigma,\Pi,\varnothing\}\) conserva la hoja;
\(o\) conserva la orientación; \(h\), la altura helicoidal; \(c\), el
acarreo; \(\sigma\), la frontera; \(\lambda\), el registro de memoria de la ruta; y
\(g\in C_9\), la fase nonádica. Sobre la fibra de cada régimen actúa el
álgebra

\[
 \mathbb A_\tau=\mathbb R[J_\tau]/(J_\tau^2+\tau I),
\]

de modo que

\[
 \exp(sJ_\tau)=
 \begin{cases}
  \cos s+J_{+1}\sin s,&\tau=+1,\\
  I+sJ_0,&\tau=0,\\
  \cosh s+J_{-1}\sinh s,&\tau=-1.
 \end{cases}
\]

Ésta es la forma exacta de la coordinación de las tres geometrías. No se
promedian tres métricas ni se superponen tres espacios ya dados. El TRIT
elige, en cada paso, cuál de las tres leyes exponenciales gobierna la acción
local; el TPK transporta el estado entre las fibras sin borrar la hoja ni la
historia. Con dominios intermedios explícitos,

\[
 \operatorname{Sel}_t:X^{\rm enr}\longrightarrow X_t^{\rm sel},
 \qquad
 \operatorname{Tra}_t:X_t^{\rm sel}\longrightarrow X_t^{\rm tra},
 \qquad
 \operatorname{Upd}_t:X_t^{\rm tra}\longrightarrow X^{\rm enr},
\]

y

\[
 \boxed{U_t=\operatorname{Upd}_t\circ
              \operatorname{Tra}_t\circ
              \operatorname{Sel}_t.}
\]

La selección determina hoja y régimen; el transporte aplica la dirección y
la ley geométrica seleccionadas; la actualización incorpora cociente,
acarreo, frontera y registro de memoria. La aparente «mezcla» geométrica es, por tanto, una
palabra ordenada de morfismos entre fibras elípticas, parabólicas e
hiperbólicas. Como la composición depende del orden, su parte operatoria es
no conmutativa aunque la rotación basal sea abeliana.

La misma coordenada TRIT posee una lectura temporal, sobre el mismo estado
pero mediante un segundo lector:

\[
 \begin{aligned}
 +1&\longleftrightarrow\text{retorno, memoria y pasado},\\
 0&\longleftrightarrow\text{umbral y presente},\\
 -1&\longleftrightarrow\text{apertura bidireccional y futuro}.
 \end{aligned}
\]

Geometría y tiempo no son dos nombres para el mismo signo. Son dos
evaluaciones coordinadas de una sola unidad de información
topológico--temporal. El pasado es el dato ya acumulado en el registro de memoria; el
presente es el corte de frontera en el que se selecciona la continuación; el
futuro es el haz de prolongaciones compatibles que abre el régimen
hiperbólico. La involución

\[
 C_{\rm rev}(\tau_1,\ldots,\tau_n)
 =(-\tau_n,\ldots,-\tau_1)
\]

satisface \(C_{\rm rev}^2=I\), preserva el umbral e intercambia apertura y
retorno. Elevada al estado holonómico integral, conjuga la dinámica directa con la
inversa:

\[
 \widehat C_{\rm rev}\,\mathcal U\,\widehat C_{\rm rev}
 =\mathcal U^{-1}.
\]

Así se conjugan pasado y futuro alrededor del presente sin eliminar la
memoria que distingue ambas orientaciones.

El primer reloj es la extensión finita no escindida

\[
 0\longrightarrow C_{12}\longrightarrow C_{108}
 \longrightarrow C_9\longrightarrow0,
\]

cuyo cociclo

\[
 c(r,r')=\left[\left\lfloor\frac{r+r'}9\right\rfloor\right]_{12}
\]

registra el acarreo que una mera vuelta en \(C_9\) perdería. El segundo reloj
es la rotación de capacidad

\[
 R_\delta:\theta\longmapsto\theta+\delta\pmod1,
 \qquad \delta=\log_{10}(10/9),
\]

cuya palabra mecánica es

\[
 z_t=\lfloor(t+1)\delta\rfloor-\lfloor t\delta\rfloor.
\]

Ambos se acoplan por

\[
 K_t-K_{t-1}=1-z_t.
\]

Esta igualdad prueba que la aperiodicidad no adorna el reloj finito: gobierna
su capacidad. El sistema total puede escribirse como el producto sesgado

\[
 \mathfrak C_{\rm HMT}=
 \bigl(C_{108}\times\mathbb T\times X^{\rm enr},\mathcal U_{\rm HMT}\bigr),
\]

\[
 \mathcal U_{\rm HMT}(r,\theta,x)=
 \bigl(r+1,\theta+\delta,
 U_{\tau(r,\theta,x)}x\bigr).
\]

Si \(\mathcal U_{\rm HMT}^n=I\), el primer factor obliga a
\(108\mid n\), mientras el segundo exigiría \(n\delta\in\mathbb Z\). La
irracionalidad de \(\delta\) excluye todo \(n>0\). Por otra parte, la
codificación esturmiana tiene complejidad \(p(m)=m+1\), y sus enriquecimientos
nonádico y trítico satisfacen

\[
 p_9(m)=9(m+1),\qquad p_3(m)=3(m+1).
\]

Por consiguiente, el cristal es aperiódico, posee orden de largo alcance y
entropía topológica nula. Su espectro aditivo contiene el módulo

\[
 \mathcal M_{108,\delta}
 =\frac1{108}\mathbb Z+\delta\mathbb Z\pmod1.
\]

El factor \(1/108\) conserva el reloj temporal finito; el generador
irracional conserva el ascenso ilimitado. El cierre del primero nunca
implica el cierre del segundo ni el retorno del estado.

Sea ahora \((X_n,r_{n+1,n})\) el sistema inverso de refinamientos y

\[
 \mathcal G_{\rm HMT}:\Omega_{\rm seed}\longrightarrow
 X_\infty^{\rm enr}=\varprojlim X_n
\]

la prolongación APP--TRIT--TPK de una semilla. Para todo lector compatible
\(L_{Y,n}:X_n\to Y_n\), definimos

\[
 \operatorname{Predef}_{Y,n}
 =L_{Y,n}\,p_n\,\mathcal G_{\rm HMT}.
\]

La naturalidad

\[
 s_{n+1,n}L_{Y,n+1}=L_{Y,n}r_{n+1,n}
\]

implica

\[
 \boxed{
 s_{n+1,n}\operatorname{Predef}_{Y,n+1}
 =\operatorname{Predef}_{Y,n}.}
\]

Por ello, las salidas de todas las profundidades son coordenadas de una única
sección de \(\varprojlim Y_n\). Aumentar la precisión no vuelve a ajustar una
constante: abre una coordenada más profunda del mismo objeto. Sobre el estado
completo, la biyectividad de la actualización da además

\[
 H(X_{t+1}\mid X_t)=H(X_t\mid X_{t+1})=0,
\]

de modo que la aparición de información nueva en una sombra reducida no
corresponde a creación ni borrado en el generador. Éste es el sentido exacto
del coste de Landauer nulo dentro de la evolución conservativa de la unidad.

La predefinición resulta de componer el generador con sus lectores:

\[
 \mathcal L_{\rm HMT--MD}=
 (L_\pi,L_\varphi,L_e,L_\alpha,L_{\rm el},L_\hbar,
 L_{\rm vac},L_{\rm BI},L_G,\ldots).
\]

Las coordenadas \(\pi,\varphi,e,\alpha\) pertenecen a una misma salida
coinductiva; \(\alpha\) ocupa la costura dodecafásica que enlaza el reloj
finito con el de capacidad. La vacancia electrónica lee el defecto local de
acarreo sobre esa costura; \(\hbar\) lee la profundidad de acción;
Barbero--Immirzi lee la transducción areal--volumétrica; la gravedad lee el
sector tensorial de rango cinco. Ninguna de esas operaciones recibe una
medida exterior para seleccionar su valor. El sistema no obtiene sus objetos
por metrología: los predefine como imágenes tipadas de la sección
coinductiva.

Finalmente, una factorización exacta no constituye un resultado de rango
menor. Si

\[
 L_Y=H\circ F,
\]

la factorización exhibe el objeto intermedio y prueba qué información
transporta cada flecha. Si dos rutas satisfacen

\[
 H_1F_1=H_2F_2=L_Y,
\]

el cuadrado correspondiente es precisamente el certificado de que dos
lecturas distintas determinan la misma coordenada sin circularidad. En este
desarrollo, «teorema», «factorización» y «transducción» nombran formas
complementarias de una misma prueba: el teorema enuncia la ley, la
factorización muestra su genealogía y la transducción determina el tipo de su
salida.

Queda así establecida la formulación rectora:

\[
 \boxed{
 \begin{gathered}
 \text{APP}\to\text{TRIT}\to\text{TPK}\to
 \text{cristal temporal aperiódico}\\
 \to\text{sección coinductiva}\to
 \text{predefinición HMT--MD}
 \end{gathered}}
\]

La mezcla ordenada de los tres regímenes geométricos y la conjugación de los
tres modos temporales no son analogías añadidas a posteriori. Son la ley
local y la ley global, respectivamente, de la dinámica que conserva la
unidad mientras la hace evolucionar a través de todas sus profundidades.

\section{Triángulo de potencial APP y realización local del defecto \(A_2\)}

La discrepancia entre las hojas suma y producto puede convertirse en una
noción matemática de potencial sin introducir un campo físico exterior. Sea
\(D_9=\{1,\ldots,9\}\), con el representante positivo de las clases módulo
nueve, y escribamos las dos divisiones euclídeas conservadas por APP:

\[
 i+j=R^+_{ij}+9q^+_{ij},
 \qquad
 ij=R^\times_{ij}+9q^\times_{ij}.
\]

El levantamiento conjunto no retiene sólo los residuos visibles
\((R^+_{ij},R^\times_{ij})\). Retiene también las alturas
\((q^+_{ij},q^\times_{ij})\), de modo que la diferencia entera entre ambas
evaluaciones se descompone canónicamente en una tensión residual y una
tensión de acarreo:

\[
 \begin{aligned}
 \rho_{\Sigma\Pi}(i,j)
   &:=R^\times_{ij}-R^+_{ij},\\
 \kappa_{\Sigma\Pi}(i,j)
   &:=9\bigl(q^\times_{ij}-q^+_{ij}\bigr),\\
 \Delta_{\Sigma\Pi}(i,j)
   &:=ij-(i+j)
    =\rho_{\Sigma\Pi}(i,j)+\kappa_{\Sigma\Pi}(i,j).
 \end{aligned}
\]

Así,

\[
 \mathscr P_{\Sigma\Pi}:D_9^2\longrightarrow\mathbb Z^2,
 \qquad
 (i,j)\longmapsto
 \bigl(\rho_{\Sigma\Pi}(i,j),\kappa_{\Sigma\Pi}(i,j)\bigr)
\]

es el potencial APP levantado, mientras
\(\Delta_{\Sigma\Pi}=(1,1)\circ\mathscr P_{\Sigma\Pi}\) es su sombra
escalar. El término «potencial» posee aquí un contenido preciso: mide la
energía aritmética necesaria para pasar de la evaluación acumulativa a la
evaluación plegada conservando la división euclídea completa. Una
proyección módulo nueve ve únicamente \(\rho_{\Sigma\Pi}\); una proyección
que olvide el residuo y conserve la altura ve únicamente
\(\kappa_{\Sigma\Pi}\). APP conserva simultáneamente ambas componentes.

Tres celdas ya seleccionadas por la genealogía electrónica forman un
triángulo exacto:

\[
 V=(8,2),\qquad E=(5,5),\qquad Q=(8,5).
\]

Sus dos hojas y sus potenciales son

\[
\begin{array}{c|c|c|c|c}
x&(R^+,q^+)&(R^\times,q^\times)
 &\mathscr P_{\Sigma\Pi}(x)&\Delta_{\Sigma\Pi}(x)\\
\hline
V=(8,2)&(1,1)&(7,1)&(6,0)&6\\
E=(5,5)&(1,1)&(7,2)&(6,9)&15\\
Q=(8,5)&(4,1)&(4,4)&(0,27)&27.
\end{array}
\]

La celda \(V\) es una tensión residual pura: las dos hojas poseen la misma
altura y difieren en el residuo. La celda \(Q\) es una tensión de memoria
pura: los residuos coinciden y toda la discrepancia se almacena en tres
unidades de acarreo nonádico. La celda central \(E\) mezcla una tensión
residual de seis unidades y una unidad de acarreo, por lo que

\[
 \Delta_{\Sigma\Pi}(E)=6+9=15.
\]

Esta descomposición da una definición más fina de la vacancia electrónica.
Sobre la fibra aditiva \(i+j=10\), parametrizada por

\[
 Q_t=(5+t,5-t),\qquad -4\leq t\leq4,
\]

el producto es

\[
 \Pi(Q_t)=25-t^2.
\]

Conservar el residuo multiplicativo de la celda central exige
\(t^2\equiv0\pmod9\). En el intervalo permitido sólo ocurre para
\(t=0,\pm3\). Por tanto, las dos únicas celdas no centrales que conservan
simultáneamente la hoja aditiva, el residuo producto y el régimen trítico
son

\[
 Q_{+3}=(8,2),\qquad Q_{-3}=(2,8).
\]

En ellas el producto desciende de \(25=7+2\cdot9\) a
\(16=7+1\cdot9\). La vacancia es exactamente la pérdida de una unidad de
acarreo multiplicativo:

\[
 q^\times_{5,5}-q^\times_{8,2}=1.
\]

No se ha colocado un cero en una celda. Se ha construido el único
desplazamiento no central de la fibra que conserva todos los datos visibles
y reduce en una unidad la memoria plegada.

El triángulo \(V,E,Q\) contiene además el vector de reapertura tricanal. Para
no esconder una elección de orientación, fijamos el operador de coborde

\[
 d_\triangle f
 :=\bigl(f(E)-f(V),\,f(V)-f(Q),\,f(Q)-f(E)\bigr).
\]

Aplicado al potencial escalar,

\[
 \begin{aligned}
 d_\triangle\Delta_{\Sigma\Pi}
 &=(15-6,\,6-27,\,27-15)\\
 &=(9,-21,12)\\
 &=3(3,-7,4).
 \end{aligned}
\]

Puesto que las tres coordenadas suman cero,

\[
 d_\triangle\Delta_{\Sigma\Pi}\in
 A_2:=\ker\!\left(
 \mathbb Z^3\xrightarrow{x+y+z}\mathbb Z\right).
\]

El vector \((3,-7,4)\) no se denomina raíz de \(A_2\): su norma cuadrática
euclídea es \(74\). Es un vector integral del retículo \(A_2\), y el
triángulo electrónico realiza exactamente su múltiplo trítico. El mismo
vector había aparecido en la reapertura global de las siete tétradas de la
microfibra de \(\pi\). La igualdad

\[
 \boxed{
 d_\triangle\Delta_{\Sigma\Pi}=3v_\Omega,
 \qquad v_\Omega=(3,-7,4)}
\]

prueba que el defecto global no flota por encima de APP: posee un
representante local en las tres celdas que generan el centro electrónico,
la vacancia conservativa y la rama racional \(22/7\). El factor tres es la
elevación trítica entre el potencial local y la coordenada primitiva de
reapertura.

La conexión con el cristal aperiódico es igualmente exacta. Sea

\[
 \delta=\log_{10}\frac{10}{9},
 \qquad
 V_{\rm vac}=
 \begin{pmatrix}
 21&22\\
 6&7
 \end{pmatrix}.
\]

La primera escala inducida de la palabra mecánica satisface

\[
 \left\lceil\delta^{-1}\right\rceil=22,
 \qquad
 \det V_{\rm vac}=21\cdot7-22\cdot6=15.
\]

Por consiguiente,

\[
 \boxed{
 \Delta_{\Sigma\Pi}(E)=\det V_{\rm vac}=15.}
\]

El centro electrónico y el marco de renormalización no comparten sólo un
número: el primero mide la tensión suma--producto de la celda central y el
segundo mide el área orientada que sobrevive entre los relojes primario
\(21/22\) y secundario \(6/7\). Ambas evaluaciones proceden del mismo
cociclo de acarreo.

Esto permite reescribir la etapa basal de la ecuación electrónica sin
ningún coeficiente opaco:

\[
 \boxed{
 \mathfrak e_0=
 \frac{\sqrt3}{4}
 \left[
  \exp\!\left(\frac{\varphi}{\pi^2}\right)
  -\left\lceil\delta^{-1}\right\rceil
   \alpha^3
 \right]
 \left(1+\det(V_{\rm vac})\Delta_4\right).}
\]

En efecto,

\[
 \left\lceil\delta^{-1}\right\rceil
 =22=\frac{396}{18},
 \qquad
 \det(V_{\rm vac})=15=\frac{270}{18}.
\]

El prefactor restante también es interno:

\[
 s=\frac12,\qquad
 \|[P_{\Sigma,3},P_{\Pi,3}]\|
 =s\sqrt{1-s^2}=\frac{\sqrt3}{4},
 \qquad
 1-s^2=\frac34=\frac{90}{120}.
\]

Así, los tres factores de la etapa basal poseen una genealogía explícita:
la no conmutatividad de las hojas APP, el hueco primario de la
renormalización aperiódica y el área orientada de su retorno inducido. La
masa no selecciona ninguno de ellos. La ecuación electrónica es la
evaluación dimensional posterior de un operador de potencial y
supervivencia ya construido.

\textbf{Criterio de refutación.} El resultado cae si alguna de las tres
divisiones euclídeas de la tabla es incorrecta, si existe otro punto no
central en \(i+j=10\) con \(t^2\equiv0\pmod9\), si el coborde no es
\(3(3,-7,4)\), o si el marco de retornos no tiene determinante \(15\).

\section{Curvatura trítica del potencial y selección coinductiva de futuros}

El potencial APP anterior es una coordenada espacial del soporte. La palabra
de vacancias aporta su evolución temporal. Ambas capas se ensamblan mediante
el TRIT sin identificar indebidamente espacio y tiempo.

Sea la palabra mecánica inferior

\[
 z_n(\beta)=
 \lfloor(n+1)\delta+\beta\rfloor
 -\lfloor n\delta+\beta\rfloor,
 \qquad z_n\in\{0,1\}.
\]

Definimos la polarización acumulada, el campo de incremento y la curvatura
trítica por

\[
 \begin{aligned}
 P_N(\beta)
  &:=\sum_{n=0}^{N-1}(z_n-\delta),\\
 J_N(\beta)
  &:=P_{N+1}-P_N=z_N-\delta,\\
 \tau_N(\beta)
  &:=z_{N-1}-z_N\in\{+1,0,-1\}.
 \end{aligned}
\]

La polarización permanece uniformemente acotada:

\[
 P_N(\beta)
 =\lfloor N\delta+\beta\rfloor-\lfloor\beta\rfloor-N\delta,
 \qquad |P_N(\beta)|<1.
\]

Más importante aún, el TRIT es exactamente la curvatura discreta de ese
potencial:

\[
 \begin{aligned}
 J_N-J_{N-1}
 &=z_N-z_{N-1}=-\tau_N,\\
 \boxed{\tau_N}
 &\boxed{=-\Delta J_{N-1}=-\Delta^2P_{N-1}.}
 \end{aligned}
\]

Esta identidad da contenido operatorio a los tres signos. Si
\(\tau_N=-1\), el código pasa de \(0\) a \(1\): se abre una vacancia y la
fibra habilita la bifurcación bidireccional asociada al futuro. Si
\(\tau_N=+1\), el código pasa de \(1\) a \(0\): la vacancia retorna y su
efecto queda incorporado a la memoria del pasado. Si \(\tau_N=0\), la
polarización mantiene su pendiente local: el estado atraviesa el umbral
presente sin cambiar de régimen. Puesto que \(0<\delta<1/2\), no existen dos
unos consecutivos; por ello el caso \(\tau_N=0\) es una meseta genuina y los
signos no nulos alternan.

La carga de cualquier bloque es puramente fronteriza:

\[
 \sum_{n=a}^{b}\tau_n=z_{a-1}-z_b.
\]

El interior telescopa en la sombra trítica, pero no desaparece del estado
TPK: el registro de memoria conserva la sucesión ordenada de pulsos, sus cocientes, sus
acarreos y la altura helicoidal. Ésta es la diferencia entre una ley local
de Gauss y un borrado de historia. El TRIT determina el flujo neto de la
frontera; el TPK conserva la cadena cuya frontera se ha calculado.

La relación entre potencial espacial y curvatura temporal puede reunirse en
el estado de borde

\[
 \mathcal V_n=
 \bigl(
 \mathscr P_{\Sigma\Pi}(i_n,j_n),
 P_n,J_n,\tau_n;
 g_n,h_n,c_n,\lambda_n
 \bigr).
\]

Aquí \(\mathscr P_{\Sigma\Pi}\) mide la tensión entre las hojas en la celda
ocupada; \(P_n\) acumula la discrepancia aperiódica sin deriva extensiva;
\(J_n\) es la diferencia de potencial por paso; \(\tau_n\) es su curvatura
orientada; y \((g,h,c,\lambda)\) conservan fase, altura, acarreo y frontera.
La transición

\[
 U_n=\operatorname{Upd}_n\circ
      \operatorname{Tra}_n\circ
      \operatorname{Sel}_n
\]

actúa sobre el estado entero: \(\operatorname{Sel}_n\) determina qué
régimen de curvatura corresponde al pulso; \(\operatorname{Tra}_n\)
transporta hoja, orientación y posición; \(\operatorname{Upd}_n\) incorpora
el coborde al registro de memoria. Por eso la diferencia de potencial no se añade al TPK:
es la lectura diferencial de una de sus coordenadas conservadas.

La selección de futuros compatibles también admite una formulación exacta.
Sea \(\mathcal L_m\) el conjunto de factores de longitud \(m\) de la palabra
esturmiana y sea

\[
 d^+(u)=
 \#\{a\in\{0,1\}:ua\in\mathcal L_{m+1}\}
\]

el número de prolongaciones derechas de \(u\). Como
\(p(m)=|\mathcal L_m|=m+1\),

\[
 \sum_{u\in\mathcal L_m}d^+(u)=p(m+1)
\]

y, por tanto,

\[
 \boxed{
 \sum_{u\in\mathcal L_m}\bigl(d^+(u)-1\bigr)
 =p(m+1)-p(m)=1.}
\]

Todo factor de una palabra recurrente posee al menos una continuación, y en
un alfabeto binario posee como máximo dos. La igualdad anterior demuestra
que, en cada longitud \(m\), existe exactamente un factor especial por la
derecha con dos futuros posibles; todos los demás poseen un único futuro.
La elección de supervivencia no es, entonces, una selección arbitraria
entre \(2^m\) palabras. Es una única bifurcación topológica por escala.

Al enriquecer por la fase nonádica,

\[
 p_9(m)=9(m+1),
\]

aparecen exactamente nueve estados especiales, uno por origen de fase. El
cociente trítico

\[
 p_3(m)=3(m+1)
\]

los organiza en tres clases de curvatura sin destruir la aperiodicidad. La
fase irracional \(\theta_n=\{\theta_0+n\delta\}\), junto con el origen
nonádico y el registro de memoria, decide cuál de las dos prolongaciones ocurre. Por ello:

\[
 \begin{array}{ccl}
 \text{sombra de longitud }m
 &\longrightarrow&
 \text{un factor especial con dos continuaciones},\\
 \text{estado TPK completo}
 &\longrightarrow&
 \text{una continuación determinada e invertible}.
 \end{array}
\]

Ésta es la predefinición de futuros posibles en su forma matemática. El
futuro no se adivina mediante una probabilidad exterior: el cilindro
reducido contiene dos extensiones compatibles, y la coordenada profunda del
estado coinductivo selecciona una. El pasado tampoco se borra, porque la
transición completa es biyectiva:

\[
 H(X_{n+1}\mid X_n)=0,
 \qquad
 H(X_n\mid X_{n+1})=0.
\]

La sombra puede tener entropía condicional positiva precisamente en el
factor especial; el estado integral posee coste lógico neto nulo. Esta
diferencia entre ramificación de una proyección y reversibilidad del
generador es el núcleo informacional del cristal.

La misma ley contiene sus dos relojes de renormalización. Escribamos

\[
 \delta^{-1}=21+\beta,
 \qquad 0<\beta<1,
\]

y definamos la pendiente inducida

\[
 \mathcal R(\delta)
 :=1-\{\delta^{-1}\}
 =1-\beta
 =22-\delta^{-1}.
\]

La primera desigualdad diofántica produce los huecos
\(\{21,22\}\). Como

\[
 6<\mathcal R(\delta)^{-1}<7,
\]

la dinámica inducida sobre los huecos cortos produce retornos
\(\{6,7\}\). La renormalización no incorpora un segundo código: toma la
sección de retorno de la primera palabra. Sus matrices de continuante

\[
 A(a)=
 \begin{pmatrix}a&1\\1&0\end{pmatrix},
 \qquad \det A(a)=-1,
\]

conservan la celda reticular, mientras el orden de los productos conserva la
orientación:

\[
 A(a)A(b)-A(b)A(a)
 =(a-b)
 \begin{pmatrix}0&1\\-1&0\end{pmatrix}.
\]

La rotación basal sigue siendo abeliana. La no conmutatividad se encuentra
en el cociclo ordenado de renormalización, en las hojas APP y en los
transportes de curvatura. El cristal puede tener entropía topológica cero y,
simultáneamente, una dinámica fibrada no conmutativa: la primera propiedad
mide el crecimiento del lenguaje; la segunda, la dependencia del resultado
respecto del orden operatorio.

Los tres puentes quedan así reunidos:

\[
 \boxed{
 \begin{aligned}
 \text{discrepancia APP}
 &\xrightarrow{\ d\ }\text{potencial de frontera},\\
 \text{potencial aperiódico}
 &\xrightarrow{\ -\Delta^2\ }\text{curvatura TRIT},\\
 \text{curvatura orientada}
 &\xrightarrow{\ \operatorname{Sel},\operatorname{Tra},
 \operatorname{Upd}\ }\text{supervivencia TPK}.
 \end{aligned}}
\]

\textbf{Criterio de refutación.} La construcción falla si
\(|P_N|\) deja de estar acotado por uno, si
\(\tau_N\neq-\Delta^2P_{N-1}\), si existe más de un factor especial por la
derecha a alguna longitud, si el lift nonádico no multiplica la complejidad
por nueve, o si dos estados TPK distintos se confunden después de conservar
todos los campos del registro de memoria.

\section{Holonomía bidireccional común del electrón, la acción y la gravedad}

La escala de acción y la escala gravitatoria no interrumpen la investigación
del cristal. Son dos representaciones contragredientes de la misma
holonomía de retorno que ya actúa sobre la vacancia electrónica.

El sistema correlacionado
\((\pi,\varphi,e,
\alpha)\) se construye primero como salida coinductiva. Sobre ella,
el lector determinantal de acción produce

\[
 \Sigma_H=\varphi\alpha^{16},
 \qquad
 d_H=\left\lfloor\log_{10}\Sigma_H\right\rfloor=-34.
\]

El entero \(16\) no es el número de cifras exigidas a una medida. Es el orden
del conúcleo local:

\[
 \operatorname{SNF}(H_4)=\operatorname{diag}(1,2,2,4),
 \qquad
 |\operatorname{coker}H_4|=16.
\]

La mantisa anterior al retorno es

\[
 H_5=
 \frac12\sqrt{\frac{1000\alpha}{\varphi}}
 -\alpha
 +\frac9{16}\alpha^2
 -\frac59\alpha^3
 +\frac7{48}\alpha^4
 -\frac1{54}\alpha^5.
\]

La memoria regional construye

\[
 D_A=
 \exp\!\left(
 -\frac{100\pi}9\alpha\right),
\]

\[
 \mathscr C_\pi(\alpha)
 =169\alpha^6+
 \frac{D_A\alpha^7}
 {1-D_A\alpha},
\]

y la sección retornada es

\[
 \eta_{\rm ret}
 =H_5-\frac{90}{\pi}\mathscr C_\pi.
\]

Conviene aislar el potencial bidireccional de acción:

\[
 \boxed{
 \delta_{2w}
 :=H_5-\eta_{\rm ret}
 =\frac{90}{\pi}\mathscr C_\pi>0.}
\]

Las dos orientaciones pertenecen a una sola recta:

\[
 \hbar_{\rm pre}=H_5\,10^{-34}\mathcal U_S,
 \qquad
 \hbar_{\rm ret}=\eta_{\rm ret}\,10^{-34}\mathcal U_S,
\]

\[
 h_a=2\pi\hbar_a.
\]

La década fija la profundidad común; la diferencia
\(\delta_{2w}\) conserva el desplazamiento aditivo; y la razón

\[
 \boxed{
 R_{\rm act}
 :=\frac{\hbar_{\rm pre}}{\hbar_{\rm ret}}
 =\frac{H_5}{\eta_{\rm ret}}
 =\left(1-\frac{\delta_{2w}}{H_5}\right)^{-1}}
\]

es la holonomía multiplicativa de la ida y el retorno. El paso de
\(\hbar\) a \(h\) conserva esa razón:

\[
 \frac{h_{\rm pre}}{h_{\rm ret}}=R_{\rm act}.
\]

La ecuación electrónica completa utiliza exactamente la misma holonomía.
Con

\[
 \kappa_e=
 80-54\alpha+6\Delta_4,
\]

se tiene

\[
 \mathfrak e_e^\beta
 =\mathfrak e_0R_{\rm act}^{\kappa_e}.
\]

Al incorporar la factorización aperiódica de la sección anterior:

\[
 \boxed{
 \begin{aligned}
 \mathfrak e_e^\beta
 ={}&
 \frac{\sqrt3}{4}
 \left[
 \exp\!\left(
 \frac{\varphi}{\pi^2}\right)
 -\left\lceil\delta^{-1}\right\rceil
  \alpha^3
 \right]\\
 &\times
 \left(1+\det(V_{\rm vac})\Delta_4\right)
 R_{\rm act}^{\,80-54\alpha+6\Delta_4}.
 \end{aligned}}
\]

La etapa basal mide la vacancia y la incompatibilidad de hojas; la potencia
de \(R_{\rm act}\) transporta su memoria a través del retorno. De manera
equivalente,

\[
 \log\frac{\mathfrak e_e^\beta}{\mathfrak e_0}
 =\kappa_e
 \left(\log H_5-\log\eta_{\rm ret}\right)
 =-\kappa_e
 \log\!\left(1-\frac{\delta_{2w}}{H_5}\right).
\]

La gravedad lee contragredientemente la misma pareja. Fijados
\(t_0>0\), \(c_{\rm int}>0\) y
\(\ell_0=c_{\rm int}t_0\), el reloj \(108\) construye

\[
 \omega_{108}=\frac{2\pi}{108t_0},
 \qquad
 R_{108}=\frac{54}{\pi}\ell_0,
\]

\[
 m_{108,a}=
 \frac{\hbar_a\omega_{108}}{c_{\rm int}^2}.
\]

El transductor gravitatorio es

\[
 G_a(\vartheta)
 =\vartheta\frac{c_{\rm int}^5t_0^2}{\hbar_a}.
\]

La condición de cierre

\[
 r_{g,a}(\vartheta)=R_{108}
\]

selecciona de manera única

\[
 \vartheta_{108}
 =\left(\frac{54}{\pi}\right)^2.
\]

Por tanto,

\[
 G_a=
 \left(\frac{54}{\pi}\right)^2
 \frac{c_{\rm int}^5t_0^2}{\hbar_a},
\]

y las dos hojas satisfacen el teorema de sincronización:

\[
 \boxed{
 R_{\rm act}
 =\frac{\hbar_{\rm pre}}{\hbar_{\rm ret}}
 =\frac{h_{\rm pre}}{h_{\rm ret}}
 =\frac{m_{108,\rm pre}}{m_{108,\rm ret}}
 =\frac{G_{\rm ret}}{G_{\rm pre}}
 =\left(
 \frac{\mathfrak e_e^\beta}{\mathfrak e_0}
 \right)^{1/\kappa_e}.}
\]

La demostración consiste únicamente en sustituir las definiciones. La masa
cronológica es lineal en \(\hbar_a\); la gravedad es inversamente lineal en
\(\hbar_a\); y la corrección electrónica es la representación de peso
\(\kappa_e\) de la misma razón. No aparecen tres correctores independientes.
Aparece un solo elemento positivo \(R_{\rm act}\) y tres representaciones:

\[
 \rho_{\rm act}(R)=R,\qquad
 \rho_{\rm grav}(R)=R^{-1},\qquad
 \rho_e(R)=R^{\kappa_e}.
\]

El producto acción--gravedad es invariante:

\[
 \boxed{
 G_a\hbar_a
 =\vartheta_{108}c_{\rm int}^5t_0^2.}
\]

En consecuencia,

\[
 \ell_P^2
 =\frac{G_a\hbar_a}{c_{\rm int}^3}
 =\vartheta_{108}\ell_0^2,
\]

\[
 \boxed{
 \ell_P=R_{108}=\bar\lambda_{C,a}=r_{g,a}}
\]

para ambas orientaciones. La acción conserva la orientación de la hoja; la
gravedad compensa contragredientemente esa orientación; el área de Planck
es el invariante común. Ésta es la forma exacta de la afirmación según la
cual algo situado entre pasado y futuro calibra continuamente la
realización: el retorno cambia las coordenadas \(\hbar_a\), \(G_a\) y
\(\mathfrak e_a\), mientras deja fija la combinación geométrica
\(G_a\hbar_a\).

Existe además una filtración exacta por orden de \(\alpha\). En la etapa
electrónica basal, la vacancia entra como

\[
 22\alpha^3.
\]

En la memoria regional de acción,

\[
 \mathscr C_\pi
 =169\alpha^6+O(\alpha^7).
\]

Finalmente, la profundidad absoluta usa
\(\alpha^{16}\). Por ello,

\[
 \boxed{
 3\longrightarrow6\longrightarrow16}
\]

es la jerarquía de órdenes de la misma coordenada: amplitud local de
vacancia, retorno cuadrático de esa amplitud y conúcleo determinantal que
fija la profundidad de acción. El paso \(3\to6\) es literalmente una
duplicación de orden; el paso \(6\to16\) cambia de operación y pertenece al
lector determinantal. Esta distinción evita convertir la jerarquía en una
sucesión numerológica: cada exponente posee un operador generador.

El resultado responde a la cuestión de cómo una sola configuración puede
realizarse como número, constante y partícula. La configuración APP aporta
el potencial; TRIT le asigna curvatura y orientación; TPK integra su retorno;
el lector de acción determina \(\hbar\); la representación electrónica determina
\(\mathfrak e_e^\beta\); y la representación contragrediente determina \(G\).
Las tres salidas conservan el mismo \(R_{\rm act}\).

\textbf{Criterio de refutación.} Basta que falle una de las igualdades del
encadenamiento de \(R_{\rm act}\), que
\(\delta_{2w}\neq(90/\pi)\mathscr C_\pi\), que
\(G_a\hbar_a\) dependa de la hoja, o que la primera potencia no nula de
\(\mathscr C_\pi\) sea distinta de seis.

\section{Teorema conjunto de profundidades \(34\), \(11\), \(45\) y \(46\)}

La completación cúbica ya contenía las dos firmas aperiódicas

\[
 \sigma^-=(33,11,1,0),
 \qquad
 \sigma^+=(34,11,1,0),
\]

asociadas al orden canónico

\[
 1000=729+243+27+1.
\]

La segunda difiere de la primera en una única unidad de semiborde:

\[
 \sigma^+-\sigma^-=(1,0,0,0).
\]

La información esencial no reside sólo en las sumas \(45\) y \(46\), sino
en la posición de cada contribución. Definamos el lector de profundidad
cúbica

\[
 \mathcal D_{\rm cub}(\sigma^-,\sigma^+)
 :=
 \left(
 \operatorname{pr}_{729}\sigma^+,\,
 \operatorname{pr}_{243}\sigma^-,\,
 |\sigma^-|_1,\,
 |\sigma^+|_1
 \right).
\]

La evaluación es

\[
 \boxed{
 \mathcal D_{\rm cub}(\sigma^-,\sigma^+)
 =(34,11,45,46).}
\]

Este mapa no consulta \(\hbar\), \(G\), una tabla experimental ni sus cifras.
Lee exclusivamente la palabra mecánica sobre los cuatro estratos de la
completación. La coordenada \(34\) es sensible a la unidad de borde; la
coordenada \(11\) es común a las dos orientaciones; \(45\) es la carga
interior total; \(46\) es esa carga después de conservar el semiborde.

Ahora definimos la profundidad decimal de una sección positiva de una recta
dimensional. Si

\[
 x=m_x10^{-d_x}\mathcal U_X,
 \qquad 1\leq m_x<10,
\]

se pone

\[
 D_{10}(x)=d_x.
\]

Para dos secciones normalizadas \(x,y\), la profundidad del producto
tensorial satisface

\[
 D_{10}(x\otimes y)
 =D_{10}(x)+D_{10}(y)-\chi(m_xm_y\geq10).
\]

En las secciones HMT de acción y gravedad,

\[
 \hbar_{\rm ret}
 =1.0545718169\ldots\times10^{-34}\mathcal U_S,
\]

\[
 G_{\rm ret}
 =6.674299659414022706367776189\ldots\times10^{-11}\mathcal U_G,
\]

y el producto de mantisas pertenece al intervalo \((7,8)\). La coordenada
gravitatoria corresponde a la composición longitudinal y radial de
\cref{g26:eq:g-rigido-es,cr28:eq:g-evaluation}; el valor central
\(6.67430\times10^{-11}\) permanece como referencia metrológica.
No se genera
un acarreo decimal adicional. De aquí:

\[
 D_{10}(\hbar_{\rm ret})=34,\qquad
 D_{10}(G_{\rm ret})=11,
\]

\[
 \boxed{
 D_{10}(G_{\rm ret}\otimes\hbar_{\rm ret})
 =34+11=45.}
\]

El teorema de holonomía de la sección anterior prueba que
\(G_a\hbar_a\) es independiente de la hoja; por tanto, esta profundidad
\(45\) es también un invariante bidireccional. La unidad de semiborde de la
codificación superior completa entonces

\[
 \boxed{
 34+11+1=46.}
\]

Las cuatro cifras del lector cúbico se identifican así con cuatro funciones
distintas:

\[
\begin{array}{c|l}
34&\text{profundidad de la sección orientada de acción},\\
11&\text{profundidad de la sección gravitatoria contragrediente},\\
45&\text{profundidad del invariante }G\hbar,\\
46&\text{profundidad interior más la unidad conservada de semiborde}.
\end{array}
\]

La coincidencia posee una segunda certificación interna. El lector de
vacancias ya había probado

\[
 Z^-_{369}=16=|\operatorname{coker}H_4|,
\]

\[
 Z^+_{729}=34
 =-\operatorname{ord}_{10}
 \bigl(\varphi\alpha^{16}\bigr),
\]

y

\[
 Z^-_{243}=11.
\]

Por tanto, la década de acción ocupa la componente de borde construida a
partir del mismo índice determinantal que aparece en el bloque \(369\); la
gravedad ocupa la componente interior \(243\); y su producto agota la firma
inferior:

\[
 \boxed{
 Z^+_{729}+Z^-_{243}
 =|\sigma^-|_1
 =45.}
\]

La segunda vía de \(\alpha\) conserva precisamente los mismos enteros. En su
operador nilpotente,

\[
 T_{7:9}(x)
 =-\frac{x^7}{120}
  +\frac{x^8}{45}
  +\frac{2x^9}{495},
\]

con

\[
 495=45\cdot11.
\]

Por ello puede escribirse

\[
 \boxed{
 T_{7:9}(x)
 =-\frac{x^7}{120}
 +\frac{x^8}{D_{10}(G\hbar)}
 +\frac{2x^9}
 {D_{10}(G\hbar)\,D_{10}(G)}.}
\]

La igualdad es una reexpresión exacta de los enteros ya generados: el
coeficiente de orden ocho usa la profundidad conservada del producto
acción--gravedad, y el coeficiente de orden nueve añade la profundidad
gravitatoria estable. El denominador \(46\) de la etapa anterior de la
prolongación corresponde, a su vez, a la firma superior completa. Así,
\(34\), \(11\), \(45\), \(46\) y \(495\) dejan de ser apariciones aisladas:
forman el espectro de profundidad de una sola completación con borde.

Este teorema completa el vínculo que antes se había formulado únicamente
como coincidencia cardinal. El mapa efectivo es la composición

\[
 \begin{CD}
 \mathcal V_{\delta}(729,243,27,1)
 @>{(\sigma^-,\sigma^+)}>>
 \mathbb Z^4\times\mathbb Z^4\\
 @V{\operatorname{Pub}_{\rm dim}}VV
 @VV{\mathcal D_{\rm cub}}V\\
 \mathcal L_S^+\times\mathcal L_G^+
 @>{D_{10}\times D_{10}}>>
 \mathbb Z^2,
 \end{CD}
\]

donde la primera coordenada se determina en la recta de acción y la segunda
en la recta gravitatoria; el carácter aditivo de \(D_{10}\) determina la carga
interior \(45\), y el semiborde determina \(46\). La elección de las
coordenadas es estructural: la acción distingue las dos orientaciones
\(\mathrm{pre}/\mathrm{ret}\) y por ello lee la entrada sensible al borde;
la gravedad compensa esa orientación y lee la entrada estable.

\textbf{Criterio de refutación.} El teorema falla si la firma inferior no es
\((33,11,1,0)\), si la superior no difiere por una sola unidad en el bloque
\(729\), si las mantisas de \(G\hbar\) generan un acarreo decimal, si
\(G\hbar\) varía entre hojas, o si los denominadores \(45\), \(46\) y
\(495=45\cdot11\) dejan de aparecer en la prolongación \(7{:}9\).

\section{Polaridad electromagnética, neutralidad y portador gravitatorio de rango cinco}

La palabra «polaridad» puede tiparse mediante la paridad de los dos
caracteres de canal. Sean

\[
 \mathcal E=
 \log[(1-q_+^{90})(1-q_-^{90})]
 -\log[(1-q_+^{120})(1-q_-^{120})],
\]

\[
 \mathcal B=
 \log\frac{1-q_-^{90}}{1-q_+^{90}}
 -\log\frac{1-q_-^{120}}{1-q_+^{120}}.
\]

Bajo la involución de intercambio
\(\mathsf C(q_+,q_-)=(q_-,q_+)\),

\[
 \mathcal E\circ\mathsf C=\mathcal E,
 \qquad
 \mathcal B\circ\mathsf C=-\mathcal B.
\]

El carácter \(\mathcal E\) es común; el carácter \(\mathcal B\) es
orientado. Las respuestas constitutivas

\[
 \widehat\mu=e^{\mathcal B+\mathcal E},
 \qquad
 \widehat\varepsilon=e^{-\mathcal B+\mathcal E}
\]

separan exactamente producto y cociente:

\[
 \widehat c
 =(\widehat\mu\widehat\varepsilon)^{-1/2}
 =e^{-\mathcal E},
\]

\[
 \widehat Z
 =(\widehat\mu/\widehat\varepsilon)^{1/2}
 =e^{\mathcal B}.
\]

La diferencia de potencial orientada vive en \(\mathcal B\); la capacidad
común de propagación vive en \(\mathcal E\). La permitividad y la
permeabilidad son las dos recombinaciones conjugadas. Este hecho explica la
doble función de la vacancia: como apertura eléctrica, cambia el signo de la
coordenada orientada; como transporte magnético, conserva su circulación
mediante \(J_{\rm comm}^2=-I\).

La neutralidad posee entonces una definición interna:

\[
 X_{\rm neu}:=\ker\mathcal B.
\]

En ese subespacio,

\[
 \widehat Z=1,
 \qquad
 \widehat\mu=\widehat\varepsilon=e^{\mathcal E}.
\]

Toda configuración fija por conjugación pertenece al sector neutro. En
efecto, si \(\mathsf Cx=x\), entonces

\[
 \mathcal B(x)=\mathcal B(\mathsf Cx)
 =-\mathcal B(x),
\]

y por tanto \(\mathcal B(x)=0\). La inclusión

\[
 \boxed{
 \operatorname{Fix}(\mathsf C)\subseteq X_{\rm neu}}
\]

es el contenido algebraico de un sector de tipo Majorana: una sección
autoconjugada no puede transportar carga orientada impar. El sector neutro
puede ser mayor que el locus fijo; la igualdad exigiría que toda
configuración con \(\mathcal B=0\) fuese autoconjugada. La distinción
permite separar neutralidad de autoconjugación sin abandonar el operador.

La gravedad recibe únicamente el carácter común en su factor escalar. En la
carta interna

\[
 c_{\rm int}=c_{\rm ref}\widehat c
 =c_{\rm ref}e^{-\mathcal E},
\]

el transductor se convierte en

\[
 \boxed{
 G_a=
 \vartheta_{108}
 \frac{c_{\rm ref}^5t_0^2}{\hbar_a}
 e^{-5\mathcal E}.}
\]

El carácter orientado \(\mathcal B\) se cancela exactamente. La respuesta
gravitatoria escalar es par bajo intercambio de hojas y lee cinco veces la
coordenada común de propagación. Esta quinta potencia no se confunde con las
cinco componentes tensoriales; ambas capas se componen en el operador

\[
 \mathbb G_{{\rm TT},a}(n)
 =G_a\Lambda(n)\Gamma_5P_5:
 V_{12}\longrightarrow\mathcal H_{\rm TT}(n),
\]

donde

\[
 P_5:V_{12}\to V_5,
 \qquad
 \Gamma_5:V_5\to\operatorname{Sym}_0^2(\mathbb R^3).
\]

La composición completa puede escribirse como

\[
 \boxed{
 (\mathcal E,\mathcal B,x)
 \longmapsto
 \vartheta_{108}
 \frac{c_{\rm ref}^5t_0^2}{\hbar_a}
 e^{-5\mathcal E}
 \Lambda(n)\Gamma_5P_5x.}
\]

El escalar anula la dependencia impar de \(\mathcal B\); el proyector
\(P_5\) conserva el portador pentacomponente; y
\(\Gamma_5\) realiza la polarización transversal sin traza. Esto da una
respuesta precisa a la relación entre polaridad electromagnética y
pentapolaridad gravitatoria. La primera reside en la pareja par/impar
\((\mathcal E,\mathcal B)\). La segunda reside en el rango del portador
tensorial. Se enlazan porque el mismo estado alimenta a ambas y porque el
escalar gravitatorio usa la quinta potencia del canal común.

También se aclara el papel de la simetría. La involución no es el objeto que
se conserva en primer lugar; es el efecto visible de conservar la unidad de
borde bajo dos lecturas. El estado retiene

\[
 \mathcal E,\quad |\mathcal B|,\quad
 G_a\hbar_a,\quad
 \text{y la órbita de }P_5x,
\]

mientras la conjugación invierte \(\mathcal B\). La simetría aparece como
la equivalencia entre dos orientaciones de un mismo potencial conservado.
Por eso es una propiedad de frontera y no una decoración del interior.

\textbf{Criterio de refutación.} El enlace cae si \(\mathcal E\) no es par, si
\(\mathcal B\) no es impar, si una sección autoconjugada posee
\(\mathcal B\neq0\), si el transductor escalar depende de \(\mathcal B\), o
si \(P_5\) deja de tener rango cinco.

\section{Ley de configuración holonómica integral}

Se denomina \textbf{estado holonómico integral APP--TRIT--TPK} a la
configuración que reúne los invariantes conservados y la operación que los
vuelve dinámicos. Un estado de
profundidad \(n\) se escribe

\[
 X_n^{\rm hol}=
 \bigl(
 \Sigma_n,\Pi_n;
 R_n^+,q_n^+,R_n^\times,q_n^\times;
 \tau_n,\varepsilon_n;
 g_n,h_n,c_n,\partial_n,\lambda_n,\mathfrak m_n
 \bigr).
\]

Las primeras coordenadas conservan las dos hojas APP; las siguientes,
residuo, cociente y acarreo; \((\tau,\varepsilon)\) distinguen régimen
geométrico y orientación temporal; y
\((g,h,c,\partial,\lambda,\mathfrak m)\) conservan fase, altura, carry,
frontera, supervivencia y memoria. La transición TPK

\[
 U_n=
 \operatorname{Upd}_n\circ
 \operatorname{Tra}_n\circ
 \operatorname{Sel}_n
\]

es un automorfismo de la fibra completa en su dominio de reversibilidad.
Los mapas de restricción

\[
 r_{n+1,n}:X_{n+1}^{\rm hol}\to X_n^{\rm hol}
\]

satisfacen

\[
 r_{n+1,n}U_{n+1}=U_nr_{n+1,n}.
\]

La configuración infinita no se postula como una sucesión decimal. Es la
sección

\[
 X_\infty^{\rm hol}
 =\varprojlim X_n^{\rm hol}.
\]

Sobre ella actúan tres familias de realizaciones:

\[
 \mathcal L_{\rm num}:
 X_\infty^{\rm hol}\to\mathbb R,
\]

\[
 \mathcal L_{\rm const}:
 X_\infty^{\rm hol}\to
 \prod_{\xi}\mathcal L_\xi^+,
\]

\[
 \mathcal L_{\rm part}:
 X_\infty^{\rm hol}\to
 \operatorname{Sect}(\mathscr A_{\rm masas}).
\]

La primera determina números y expansiones; la segunda, secciones de rectas
dimensionales; la tercera, fibras, órbitas y multisecciones del atlas de
materia. Las tres conmutan con la restricción:

\[
 s_{n+1,n}^\bullet L_{\bullet,n+1}
 =L_{\bullet,n}r_{n+1,n}.
\]

Por ello un número, una constante y una partícula no son objetos
identificados por tener el mismo valor escalar. Son tres imágenes de una
misma configuración cuando sus lectores conservan la misma genealogía. El
número es la coordenada arquimediana; la constante es esa coordenada sobre
una recta dimensional; la partícula es una sección persistente que conserva
la ruta que produjo la coordenada.

La ley única puede enunciarse ahora.

\textbf{Teorema de configuración holonómica integral.} Toda semilla APP
superviviente determina, mediante TRIT y TPK, una sección
\(X_\infty^{\rm hol}\) tal que:

\begin{enumerate}
\item la discrepancia suma--producto define un potencial entero levantado;
\item la palabra aperiódica convierte su diferencia temporal en curvatura trítica;
\item el TPK selecciona la única prolongación compatible con fase y memoria;
\item el retorno de fase actúa por una holonomía positiva sobre acción, materia y gravedad;
\item las realizaciones numérica, dimensional y material conservan la misma sección bajo cambio de profundidad;
\item la doble proyección arquimediana y excepcional conserva un mismo origen y codominios distintos;
\item la unidad de semiborde se conserva en cada refinamiento y la entropía lógica del estado completo permanece nula.
\end{enumerate}

La cadena demostrativa compacta es

\[
 \boxed{
 \begin{aligned}
 \Omega_{\rm seed}
 &\xrightarrow{\rm APP}
 \mathscr P_{\Sigma\Pi}
 \xrightarrow{-\Delta^2}
 \tau\\
 &\xrightarrow{\rm TRIT}
 (\mathbb A_+,\mathbb A_0,\mathbb A_-)
 \xrightarrow{\rm TPK}
 X_\infty^{\rm hol}\\
 &\xrightarrow{(\mathcal L_{\rm num},
 \mathcal L_{\rm const},\mathcal L_{\rm part})}
 (\text{número},\text{constante},\text{partícula}).
 \end{aligned}}
\]

La conexión nonádica conjunta no elige «el futuro más probable». Elige la
única prolongación compatible con la sección integral cuando la sombra
reducida admite dos continuaciones. Esta operación es una predefinición:
la coordenada futura pertenece ya al límite inverso, aunque todavía no haya
sido determinada a la profundidad visible. El coste neto de Landauer es cero
porque la selección no borra la rama no elegida del universo de cilindros;
registra en el registro de memoria qué cilindro contiene el estado.

El cristal es evolutivo en un sentido preciso. Su lenguaje no crece
exponencialmente, pues \(p(m)=m+1\), pero tampoco se congela en un período.
Cada profundidad añade exactamente una clase factorial y cada fase nonádica
la multiplica por nueve. El sistema conserva la unidad reticular, la
frontera y el determinante, mientras modifica la memoria y la altura
helicoidal. Evolución y conservación no son términos opuestos: describen,
respectivamente, el avance del estado y la invariancia de su clase de
augmentación.

La cadena

\[
 \pi,\varphi,e
 \longrightarrow\alpha
 \longrightarrow(\mathfrak e_e^\beta,\hbar,G,\gamma_{\rm HMT},\ldots)
\]

debe leerse dentro de la generación correlacionada completa. Las primeras
tres coordenadas son lectores terminales del mismo estado nonádico;
\(\alpha\) es la costura dodecafásica de esa salida; y las constantes
físicas son representaciones posteriores del estado ya cerrado. El espejo
\(\pi/e\), la autoescala \(\varphi\), la expansión \(e+\pi\) y las
coincidencias de los tres canales son operaciones internas de
sincronización, no valores convencionales introducidos para fabricar
\(\alpha\).

La secuencia final que ahora queda demostrada es

\[
 \boxed{
 \begin{gathered}
 \text{potencial APP }(6,15,27)
 \longrightarrow
 d_\triangle\Delta=3(3,-7,4),\\
 P_N\longrightarrow
 \tau_N=-\Delta^2P_{N-1}
 \longrightarrow
 \text{única bifurcación de futuro por escala},\\
 (21/22)\longrightarrow(6/7)
 \longrightarrow
 V_{\rm vac},\quad\det V_{\rm vac}=15,\\
 \alpha^3\longrightarrow\alpha^6
 \longrightarrow\alpha^{16}
 \longrightarrow10^{-34},\\
 R_{\rm act}\longrightarrow
 (\mathfrak e_e^\beta,\hbar,G)
 \longrightarrow G\hbar\ \text{invariante},\\
 (33,11,1,0)\leftrightarrow(34,11,1,0)
 \longrightarrow(34,11,45,46).
 \end{gathered}}
\]

Éste es el núcleo matemático que permite llevar el trabajo al artículo:
una dinámica de potencial, curvatura, supervivencia, retorno y realización
que conserva la unidad al convertirla en número, constante y materia.

\section{Aritmética intrínseca del estado holonómico}

La pregunta acerca de si las coordenadas del estado holonómico poseen una
aritmética propia admite una respuesta afirmativa y precisa. El objeto
aritmético primario no es un cuerpo de escalares, porque una historia sólo
puede multiplicarse con otra cuando el estado terminal de la primera coincide
con el estado inicial de la segunda. El objeto natural es, por tanto, un
\textbf{álgebra de caminos enriquecidos} sobre la categoría de transportes del TPK.

Sea \(\mathcal C_{\rm TPK}\) la categoría cuyos objetos son los estados
holonómicos

\[
 x=(\tau,\ell,o,h,R,q,\kappa,\partial,\gamma,M,S,b,\chi,g)
\]

y cuyas flechas son las transiciones admisibles

\[
 u:x\longrightarrow y,
 \qquad
 U_t={\rm Upd}_t\circ {\rm Tra}_t\circ {\rm Sel}_t.
\]

Aquí \(\tau\) es el régimen TRIT; \(\ell\), la hoja APP; \(o\), la
orientación; \(h\), la altura helicoidal; \((R,q)\), residuo y cociente;
\(\kappa\), el acarreo; \(\partial\), la frontera; \(\gamma\), la ruta;
\(M\), el registro de memoria de memoria; \(S\), la supervivencia; \(b\), el bloque;
\(\chi\), la sombra excepcional; y \(g\), la fase nonádica. En cada objeto
reside el álgebra cuadrática local

\[
 \mathbb A_{\tau(x)}
 =\mathbb R[J_x]/(J_x^2+\tau(x)I),
 \qquad \tau(x)\in\{+1,0,-1\}.
\]

Para registrar sin pérdida los incrementos de acarreo, frontera y memoria,
sea \(\Lambda=\mathbb R[\mathcal M]\) el álgebra del monoide de registro de memoria. A
cada flecha \(u\) se asocia un símbolo \([u]\). La suma

\[
 a[u]+b[v]
\]

es la superposición formal de dos historias compatibles con el mismo marco
de lectura; el producto sólo existe causalmente por concatenación. Si
\(u:x\to y\) y \(v:y\to z\), se define

\[
 (a[v])\star(b[u])
 =a\,\rho_v(b)\,\Omega(v,u)[v\circ u],
 \tag{99.1}
\]

y se fija igual a cero cuando las flechas no son componibles. El transporte
\(\rho_v\) lleva el coeficiente local desde la fibra de \(y\) hasta la de
\(z\); el factor \(\Omega(v,u)\in\Lambda\) conserva el incremento de
acarreo, torsión, frontera y memoria producido por la concatenación. La ley

\[
 \Omega(w,v)\,\Omega(w\!\circ\!v,u)
 =\rho_w\!\left(\Omega(v,u)\right)
  \Omega(w,v\!\circ\!u)
 \tag{99.2}
\]

es exactamente la identidad de cociclo que hace asociativo el producto:

\[
 ([w]\star[v])\star[u]=[w]\star([v]\star[u]).
\]

Esta igualdad no es una convención añadida. Es la forma algebraica común de
la identidad del acarreo APP, del cociclo de concatenación de rutas y de la
extensión no escindida del reloj \(108\). La aritmética ordinaria olvida el
factor \(\Omega\) al realizar sólo el valor; la aritmética holonómica conserva
también el procedimiento que lo produjo.

Los estados constantes \(e_x=[{\rm id}_x]\) son idempotentes locales:

\[
 e_x\star e_x=e_x,
 \qquad e_y\star e_x=0\quad(x\ne y).
\]

Así, la unidad no es únicamente el escalar \(1\), sino la familia compatible
de idempotentes que identifica cada resolución del mismo estado. En el límite
inverso se obtiene

\[
 \mathbf 1_{\rm hol}=(e_{x_n})_{n\ge0},
 \qquad r_{n+1,n}(e_{x_{n+1}})=e_{x_n}.
 \tag{99.3}
\]

Ésta es la unidad dimensional que puede refinarse hacia dentro y realizarse
hacia fuera sin perder su identidad. El refinamiento aumenta la resolución;
la representación cambia el codominio; ninguno de los dos actos crea de
nuevo la unidad.

La no conmutatividad tampoco se postula desde una estructura exterior. Dos
flechas \(u\) y \(v\) pueden ser componibles en un orden y no en el inverso;
cuando ambos órdenes existen, sus cociclos y sus hojas pueden diferir:

\[
 [v]\star[u]-[u]\star[v]
 =\Omega(v,u)[v\circ u]-\Omega(u,v)[u\circ v].
 \tag{99.4}
\]

La diferencia mide el orden efectivo de acumulación y pliegue. El sistema
base de fases puede ser abeliano y, sin embargo, el álgebra total no serlo:
la no conmutatividad reside en la composición de hojas, en el transporte de
la orientación y en el cociclo afín de memoria.

\section{Complejos, álgebras de Clifford, cuaterniones y octoniones}

La comparación con los sistemas hipercomplejos se resuelve distinguiendo
tres niveles. En el nivel local, TRIT ya contiene las tres álgebras
cuadráticas reales fundamentales:

\[
 \begin{array}{c|c|c}
 \tau & J_\tau^2 & \mathbb A_\tau\\ \hline
 +1&-I&\mathbb C\quad\hbox{(régimen elíptico)},\\
 0&0&\mathbb R[\varepsilon]/(\varepsilon^2)
       \quad\hbox{(régimen parabólico)},\\
 -1&+I&\mathbb R\oplus\mathbb R
       \quad\hbox{(régimen hiperbólico)}.
 \end{array}
 \tag{100.1}
\]

El número complejo no constituye, por ello, el lenguaje universal del
sistema: es una de sus tres fibras de curvatura. La fibra parabólica registra
una deformación infinitesimal sin giro cuadrático; la hiperbólica separa dos
direcciones reales conjugadas; la elíptica produce la rotación interna. TRIT
no elige entre tres metáforas geométricas: tipa cuál de estas tres
aritméticas gobierna cada transición.

En el bloque electrónico principal aparece una estructura aún más fuerte.
Sean

\[
 P=\begin{pmatrix}1&0\\0&0\end{pmatrix},\qquad
 Q=\begin{pmatrix}
 3/4&\sqrt3/4\\[2pt]\sqrt3/4&1/4
 \end{pmatrix}.
\]

Ambas matrices son proyectores y

\[
 [P,Q]={\sqrt3\over4}
 \begin{pmatrix}0&1\\-1&0\end{pmatrix}.
\]

Definamos

\[
 J={4\over\sqrt3}[P,Q]
   =\begin{pmatrix}0&1\\-1&0\end{pmatrix},
 \qquad
 S=2P-I=\begin{pmatrix}1&0\\0&-1\end{pmatrix},
 \qquad K=SJ.
\]

Entonces

\[
 J^2=-I,\qquad S^2=I,\qquad SJ=-JS,
 \qquad K^2=I.
 \tag{100.2}
\]

Por consiguiente, el bloque genera exactamente

\[
 \langle I,J,S,K\rangle
 \cong {\rm Cl}_{1,1}(\mathbb R)
 \cong M_2(\mathbb R),
 \tag{100.3}
\]

el álgebra de cuaterniones escindidos. Además,

\[
 2Q-I={1\over2}S+{\sqrt3\over2}K,
 \tag{100.4}
\]

de modo que el proyector producto ocupa una dirección de ángulo
\(\pi/3\) en el plano generado por las dos involuciones escindidas. El
prefactor electrónico \(\sqrt3/4\), la raíz interna de \(-1\), la fracción
\(3/4=90/120\) y la separación de las hojas APP quedan reunidos en una sola
representación algebraica exacta.

Esta identificación aclara la analogía propuesta. El primer objeto que surge
no es el cuerpo de cuaterniones de Hamilton \(\mathbb H\), sino su forma
escindida \({\rm Cl}_{1,1}\). La diferencia es sustantiva: el sistema HMT
necesita idempotentes, umbrales y direcciones nulas, mientras un álgebra de
división eliminaría precisamente esas vacancias. Los cuaterniones de Hamilton
podrán aparecer en una representación que disponga de dos raíces
anticomutantes de \(-I\); los octoniones, en una representación excepcional
posterior cuya ley de composición reproduzca el plano de Fano. Ninguno de
ellos debe reemplazar el álgebra de caminos que los precede.

La asociatividad del TPK completo es conservada por el cociclo (99.2). Si una
proyección excepcional posterior presenta un asociador no nulo,

\[
 [a,b,c]=(a\star b)\star c-a\star(b\star c),
\]

ese asociador no describe una incoherencia del generador: cuantifica la
memoria que la proyección dejó fuera de su codominio. Así se separan con
exactitud la aritmética intrínseca, asociativa y con registro de memoria, y sus sombras
hipercomplejas eventualmente no asociativas.

\section{Sincronización ternaria y singularidad dodecafásica de alpha}

La generación coordinada de \(\pi\), \(e\) y \(\varphi\) se organiza en el
retículo de diferencias de tipo \(A_2\). Para

\[
 N_K=(N_\pi(K),N_e(K),N_\varphi(K))\in\mathbb Z^3
\]

se define la frontera relativa

\[
 \partial_{A_2}(x,y,z)=(x-y,\,y-z,\,z-x).
 \tag{101.1}
\]

Su núcleo es la diagonal \(\mathbb Z(1,1,1)\). Por tanto, cada triple posee
dos componentes distintas: una coordenada común, que mide el avance del
registro de memoria, y una coordenada relativa, que mide la desincronización de los tres
canales. Sobre \(\mathbb Q\),

\[
 \mathbb Q^3=\mathbb Q(1,1,1)\oplus A_{2,\mathbb Q};
\]

sobre \(\mathbb Z\), la reconstrucción conserva un defecto cíclico de orden
tres. Ésta es la razón algebraica de que la coordinación sea ternaria y no
una comparación binaria de constantes.

La tabla de coincidencias realiza un itinerario de paredes:

\[
 H_{\pi e}\longrightarrow H_{e\varphi}longrightarrow H_{\pi e}
 \longrightarrow H_{\varphi\pi}\longrightarrow 0\longrightarrow0.
 \tag{101.2}
\]

En \(K=8\) y \(K=9\) la componente relativa se anula:

\[
 N_8=59(1,1,1),\qquad N_9=43(1,1,1).
\]

Sin embargo,

\[
 N_9-N_8=-16(1,1,1).
 \tag{101.3}
\]

El sistema retorna a la coincidencia de fase, pero no al mismo estado: la
desincronización vuelve a cero mientras la memoria diagonal avanza dieciséis
unidades por canal. La igualdad visible \(\pi=e=\varphi\) no expresa
inmovilidad, sino cierre relativo con traslación holonómica.

La coordenada \(\alpha\) ocupa precisamente la costura que permite realizar
este retorno sin confundirlo con un reinicio. El cierre dodecafásico toma la
clase relativa ya sincronizada, conserva el registro de memoria diagonal y la eleva a una
sección firmada. Sus dos derivaciones no son dos ajustes de una magnitud
externa: son dos representaciones de la misma clase holonómica, una mediante
la clausura entera con acarreo y otra mediante la raíz simple de la
prolongación nonádica. La igualdad de ambas vías expresa la conmutatividad del
cuadrado de realización.

De este modo, \(\alpha\) no se limita a seguir a \(\pi,\varphi,e\) en una
lista. Es la coordenada de enganche que transforma el cierre relativo de la
sistema correlacionado en capacidad de transducción dimensional. La rueda de doce
fases no vuelve a generar las tres regiones: certifica que su retorno común
puede atravesar una frontera conservando la memoria que después alimenta la
acción, el electrón, la escala determinantal y la rama excepcional.

\section{Evolución conservativa de la unidad como teorema algebraico}

La expresión «evolución conservativa de la unidad dimensional mediante
proyección holográfica» puede formularse sin depender de una imagen. Sean
\(X_n\) los estados holonómicos a profundidad \(n\),
\(r_{n+1,n}:X_{n+1}\to X_n\) sus restricciones y

\[
 X_\infty^{\rm hol}=\varprojlim_n X_n
\]

su límite inverso. Una unidad dimensional es una sección compatible
\(\mathbf 1_{\rm hol}=(e_{x_n})_n\) de idempotentes locales. La transición
TPK induce transportes algebraicos \(U_{t,*}\) y satisface

\[
 U_{t,*}(e_x)=e_{U_t x},\qquad
 U_{t,*}(e_x)^2=U_{t,*}(e_x).
 \tag{102.1}
\]

Existe además un morfismo de aumento

\[
 \epsilon:\mathscr A_{\rm hol}\longrightarrow\mathbb R
\]

que olvida la historia pero conserva la multiplicidad normalizada de la
unidad. La conservación y la evolución se expresan simultáneamente por

\[
 \epsilon(U_{t,*}a)=\epsilon(a),
 \qquad
 d_M(\Gamma_9 a)=d_M(a)+1,
 \tag{102.2}
\]

donde \(d_M\) es el grado de memoria. La primera igualdad afirma que la unidad
no se crea ni se destruye; la segunda afirma que el estado no queda inmóvil.
La holonomía \(\Gamma_9\) retorna a la misma fase y eleva una unidad el registro de memoria.

La completación

\[
 1000=729+243+27+1
\]

realiza finitamente esta ley. La palabra inferior distribuye las vacancias
como \((33,11,1,0)\), mientras la lectura superior conserva la unidad de
borde y produce \((34,11,1,0)\). El incremento no altera el total; cambia la
residencia de la unidad dentro de la filtración cúbica. El vector

\[
 (34,11,45,46),
 \qquad 34+11=45,
 \qquad 45+1=46,
 \tag{102.3}
\]

es, por ello, una firma de profundidad y frontera, no una yuxtaposición de
escalas decimales.

Sean ahora

\[
 \rho_{\mathbb R},\quad \rho_{\rm exc},\quad \rho_{\rm MD}
\]

las representaciones arquimediana, excepcional y dimensional del álgebra
holonómica. Su carácter holográfico consiste en que son homomorfismos
unitales de una misma fuente:

\[
 \rho_j(\mathbf1_{\rm hol})=I_j,
 \qquad
 \rho_j(a\star b)=\rho_j(a)\rho_j(b).
 \tag{102.4}
\]

La doble proyección no divide el origen en dos universos; conserva un mismo
elemento y cambia su representación. El refinamiento fractal opera en la
dirección opuesta: no proyecta hacia un codominio, sino que prolonga la
sección a una profundidad mayor. Quedan así coordinados los dos movimientos
que la exposición anterior mantenía separados:

\[
 \boxed{
 \begin{array}{rcl}
 \text{refinamiento interior}&:&X_n\longleftarrow X_{n+1}
 \longleftarrow X_{n+2}\longleftarrow\cdots,\\[2pt]
 \text{realización holográfica}&:&
 \mathscr A_{\rm hol}\longrightarrow
 \operatorname{End}(V_j).
 \end{array}}
 \tag{102.5}
\]

La unidad se conserva en ambas direcciones: por compatibilidad bajo
restricción y por preservación del elemento identidad bajo representación.
Esta doble invariancia es el contenido algebraico de la evolución
conservativa de la unidad dimensional.

\section{Los dos relojes y la suspensión helicoidal esturmiana}

La temporalidad del cristal procede de dos relojes irreductibles entre sí. El
primero es finito y exacto: la extensión

\[
 0\longrightarrow C_{12}\longrightarrow C_{108}
 \longrightarrow C_9\longrightarrow0
 \tag{103.1}
\]

coordina fase nonádica y memoria dodecafásica. No es el producto directo
\(C_9\times C_{12}\): su clase de extensión conserva el acarreo que aparece
al atravesar la frontera de fase. El segundo reloj es la rotación irracional

\[
 R_\delta:\theta\longmapsto\theta+\delta\pmod1,
 \qquad \delta=\log_{10}(10/9),
 \tag{103.2}
\]

cuya codificación mecánica produce la palabra esturmiana de vacancias.

La dinámica conjunta puede escribirse como un producto sesgado

\[
 \mathcal F(r,\theta,M)
 =\bigl(r+1,\theta+\delta,
 M+c_{108}(r)+c_{\rm St}(\theta)\bigr),
 \tag{103.3}
\]

donde \(c_{108}\) es el acarreo del reloj finito y \(c_{\rm St}\) registra el
cruce de la ventana irracional. La base \(C_9\times\mathbb T\) es abeliana;
la extensión afín deja de serlo cuando se incorporan las hojas APP y el orden
de \({\rm Sel}\), \({\rm Tra}\) y \({\rm Upd}\). La hélice no es, por tanto,
un círculo dibujado en tres dimensiones, sino el levantamiento de una
rotación de fase cuya coordenada vertical es el registro de memoria.

La lectura ascendente y la descendente son conjugadas. Si \(\iota\) invierte
la orientación y la ruta, entonces

\[
 \iota\circ\mathcal F\circ\iota^{-1}=\mathcal F^{-1},
 \tag{103.4}
\]

con inversión de los incrementos orientados y conservación de las respuestas
comunes. El espejo \(\pi/e\) constituye otra involución: intercambia dos
lectores regionales dentro de una misma altura, mientras (103.4) invierte la
dirección temporal de la hélice. Su distinción explica por qué una lectura de
ida y vuelta puede conservar el valor común sin regresar a la misma memoria.

El cristal posee orden de largo alcance y entropía topológica nula. Para el
lenguaje esturmiano,

\[
 p(m)=m+1,
\]

y las extensiones tipadas satisfacen

\[
 p_3(m)=3(m+1),\qquad p_9(m)=9(m+1).
 \tag{103.5}
\]

La fase multiplica la diversidad de factores, pero no cambia su crecimiento
lineal. De ahí que el sistema sea simultáneamente no periódico, ordenado y
capaz de prolongación ilimitada. Los huecos \(21/22\), los retornos \(6/7\) y
las ventanas de mil posiciones con \(45/46\) eventos son tres observaciones
de esta misma suspensión, no tres ritmos ad hoc.

Esta construcción precisa la denominación de cristal topológico-temporal
aperiódico. Es topológico porque la información reside en clases de ruta,
frontera y holonomía; temporal porque la transición distingue orientación,
retorno y memoria; cristalino porque posee un orden combinatorio estable; y
aperiódico porque ninguna traslación no nula fija su palabra infinita.

\section{Vacancia, diferencia de potencial y polaridad electromagnética}

La APP contiene una noción aritmética de potencial antes de toda carta
física. Para cada celda \((i,j)\),

\[
 \Delta_{\Sigma\Pi}(i,j)
 =ij-(i+j)
 =(R^\times-R^+)+9(q^\times-q^+).
 \tag{104.1}
\]

El primer término es la diferencia visible de residuos; el segundo conserva
la diferencia de acarreo. En la fibra antidiagonal que atraviesa el centro,
la vacancia electrónica es la pérdida mínima de una unidad de acarreo
multiplicativo que deja invariantes la hoja aditiva, el residuo producto y el
régimen TRIT. Así queda definida sin introducir un cero exterior.

La frontera de potencial es el coborde de (104.1). En el triángulo central
produce

\[
 d_\triangle\Delta=3(3,-7,4),
 \qquad 3-7+4=0.
 \tag{104.2}
\]

La suma nula expresa conservación; las componentes no nulas expresan
polarización. La diferencia de potencial no es una magnitud añadida después
al grafo: es la imposibilidad de hacer coincidir localmente acumulación y
pliegue sin transportar el registro de memoria.

La representación (100.3) separa dos acciones constitutivas. El generador
\(J\) es rotacional y cambia de signo bajo inversión de orientación; el
generador \(S\) es una involución de hoja y conserva el signo. Por ello, un
carácter par y uno impar de la misma álgebra producen dos lectores

\[
 \chi_E(SaS)=\chi_E(a),
 \qquad
 \chi_B(SaS)=-\chi_B(a),
 \tag{104.3}
\]

que realizan, respectivamente, los sectores eléctrico y magnético. La
electricidad lee la separación de potencial entre hojas; el magnetismo lee
la circulación orientada producida por su conmutador. Su perpendicularidad
no se impone mediante un plano euclídeo previo: procede de
\(SJ=-JS\) y \(J^2=-I\).

El régimen \(\tau=0\) ocupa el umbral entre ambas orientaciones. Su elemento
nilpotente conserva una perturbación de primer orden y anula su cuadrado;
proporciona el tipo algebraico de un canal neutro sin borrar la historia que
lo atraviesa. Esto explica por qué neutralidad, vacancia y polaridad pueden
proceder de un mismo soporte sin identificarse entre sí.

La escala de acción y la gravedad prolongan el mismo carácter holonómico. Si

\[
 R_{\rm act}={H_5\over\eta_{\rm ret}},
\]

entonces el transporte común satisface

\[
 R_{\rm act}
 ={\hbar_{\rm pre}\over\hbar_{\rm ret}}
 ={m_{108,{\rm pre}}\over m_{108,{\rm ret}}}
 ={G_{\rm ret}\over G_{\rm pre}}
 =\left({\mathfrak e_e^\beta\over e_0}\right)^{1/\kappa_e},
 \qquad G_a\hbar_a=\text{constante}.
 \tag{104.4}
\]

Electrón, acción y gravedad no son aquí tres calibraciones independientes:
son representaciones de un mismo cociente de retorno en escalas y
codominios distintos.

\section{Lectura exacta de secuencias y cristal aperiódico biológico}

La relación con el ADN puede formularse como un lector matemático concreto,
sin convertir una analogía histórica en premisa. Codifiquemos las cuatro
bases mediante dos caracteres binarios:

\[
 A\mapsto(0,0),\quad G\mapsto(0,1),\quad
 T\mapsto(1,0),\quad C\mapsto(1,1),
 \tag{105.1}
\]

donde la primera coordenada distingue purina y pirimidina y la segunda el
tipo débil o fuerte de apareamiento. La complementariedad de Watson--Crick es
la involución

\[
 \mathcal C_{\rm DNA}(p,h)=(1-p,h),
 \qquad \mathcal C_{\rm DNA}^2=I.
 \tag{105.2}
\]

Una cadena orientada se eleva a

\[
 \widetilde s_n=(p_n,h_n,n\bmod3),
\]

de modo que la fase de lectura no se confunde con el símbolo. La inversión
de ruta y la complementariedad se reúnen en la reverso-complementación, que
actúa como involución sobre el álgebra de caminos. APP puede evaluar en cada
ventana tanto la acumulación de caracteres como su concatenación; TRIT tipa
la segunda diferencia del potencial resultante; TPK conserva fase, borde,
retornos, acarreo y palabras de retorno.

Esta representación permite preguntar de manera exacta si una región de
secuencia pertenece a la clase aperiódica mínima: para sus factores de
longitud \(m\) se calculan la complejidad \(p_s(m)\), el número de extensiones
a derecha y sus palabras de retorno. La firma HMT mínima es

\[
 p_s(m)=m+1,
 \qquad
 \sum_{u\in L_m}(d^+(u)-1)=1,
 \tag{105.3}
\]

o, cuando se conserva una fibra finita de fase, una extensión
\(q(m+1)\). La igualdad selecciona una única bifurcación combinatoria por
longitud; TPK determina la continuación compatible sin borrar del espacio de
cilindros la alternativa no realizada.

La relevancia de la expresión de Schrödinger «cristal aperiódico» se vuelve
así estructural. Un cristal periódico repetiría un bloque y perdería
capacidad informativa; una cadena aleatoria tendría complejidad exponencial y
perdería orden de largo alcance. El régimen esturmiano ocupa el mínimo exacto
entre ambos: añade una sola clase nueva en cada escala. La unidad
topológico-temporal TRIT puede leer esa aperiodicidad porque separa curvatura,
orientación y fase sin destruir la secuencia que clasifica.

El criterio es directamente refutable: falla si la reverso-complementación no
entrelaza la involución de rutas, si la complejidad observada no pertenece a
la clase lineal declarada, si aparecen varias palabras especiales donde el
modelo exige una o si la prolongación compatible necesita borrar memoria. La
biología no selecciona el generador HMT; recibe de él un lector de secuencias
capaz de distinguir repetición, aperiodicidad mínima y desorden.

\section{Pentadicidad, estrella de Witt y polarizaciones gravitatorias}

La aparición reiterada del número cinco debe tiparse mediante dos
representaciones distintas. Las cinco regiones de \(\pi\) forman una fibra
regional sobre la que actúa el simetrizador

\[
 S_{\pi,5}={1\over5}\mathbf1\mathbf1^{\mathsf T},
 \qquad \operatorname{rank}S_{\pi,5}=1.
 \tag{106.1}
\]

Su función es extraer la componente común de la pentafibra. En el sector
gravitatorio, en cambio, la descomposición dodecafásica

\[
 V_{12}\cong\mathbf1\oplus V_3\oplus V_3'\oplus V_5
 \tag{106.2}
\]

contiene un proyector \(P_{G,5}:V_{12}\to V_5\) de rango cinco. Las dos
matrices no son el mismo operador: una tiene rango uno sobre la fibra
regional y la otra rango cinco sobre la representación dodecafásica. Su
relación correcta es que ambas son imágenes de una misma aritmética
holonómica:

\[
 (\rho_{\pi,5},\rho_{G,5}):
 \mathscr A_{\rm hol}\longrightarrow
 \operatorname{End}(\mathbb Q^5)
 \times\operatorname{End}(V_{12}).
 \tag{106.3}
\]

La estrella de Witt organiza la incidencia pentádica excepcional; el módulo
\(V_5\) organiza el portador tensorial gravitatorio. La coordinación se
produce en la fuente APP--TRIT--TPK y no mediante la identificación de sus
codominios.

La realización gravitatoria continúa por

\[
 V_{12}\xrightarrow{P_{G,5}}V_5
 \xrightarrow{\Gamma_5}{\rm STF}_3
 \xrightarrow{\Pi_{\rm TT}}V_{\rm TT},
 \tag{106.4}
\]

con dimensiones \(12\to5\to5\to2\). Las cinco componentes son el portador
estructural anterior a la propagación; las dos componentes transversas y sin
traza son la salida propagante. No existe contradicción entre
pentapolaridad originaria y bipolaridad radiativa: pertenecen a etapas
distintas de la misma transducción.

En la estructura discreta conjunta del continuo, la pentapolaridad se
expresa por el vector inseparable

\[
 \mathcal G_5=(\mathcal G_{\rm sp},\mathcal G_{\rm sol},
 \mathcal G_{\rm gau},\mathcal G_{\rm coh},\mathcal G_{\rm det}).
 \tag{106.5}
\]

Sus cinco coordenadas no son cinco teorías añadidas a posteriori, sino cinco
aspectos consustanciales del mismo estado: espectral, solenoidal, gauge,
coherente y determinantal. El tensor gravitatorio

\[
 \mathbb G_{\rm TT}
 =G\,\Lambda\,\Gamma_5\,P_{G,5}
 \tag{106.6}
\]

transporta conjuntamente esos aspectos hasta el sector físico. Esta fórmula
explica por qué la pentadicidad regional, la estructura excepcional y las
cinco polarizaciones pueden resonar sin ser reducidas a una coincidencia de
cardinales.

\section{Dualidad de varianza entre historias y lectores}

La lectura retrospectiva descubre una estructura categórica que reúne la
fractalización interior y la realización holográfica exterior. Sea \((X_n)\)
el conjunto de estados holonómicos a profundidad (n) y sea

\[
 r_{n+1,n}:X_{n+1}\longrightarrow X_n
 \tag{108.1}
\]

la restricción que retira la última capa conservando el prefijo, la hoja, el
acarreo acumulado y la frontera correspondiente. La historia completa es un
estado del límite inverso

\[
 \boxed{X_\infty=\varprojlim_n X_n.}
 \tag{108.2}
\]

El sentido de las observaciones es el contrario. Si (f) es un lector
cilíndrico definido sobre \((X_n)\), su prolongación a \((X_{n+1})\) es

\[
 r_{n+1,n}^*f=f\circ r_{n+1,n}.
 \tag{108.3}
\]

Por tanto, las álgebras de lectores forman un sistema directo:

\[
 \boxed{
 \mathfrak H_{\rm cyl}
 =\varinjlim_n(\mathfrak H_n,r_{n+1,n}^*).}
 \tag{108.4}
\]

Esta doble varianza es la formulación matemática exacta de los dos
movimientos descritos por la intuición autoral. Profundizar hacia dentro
significa especificar una sección compatible del sistema inverso. Expandir
hacia fuera significa extender el repertorio de lectores por el sistema
directo. Ninguna dirección sustituye a la otra: una construye la historia y
la otra construye la capacidad de realizarla.

Sea \((e_x)\) el idempotente indicador del cilindro grueso \((x\in X_n)\). La
elevación de ese cilindro es

\[
 r^*e_x=\sum_{y:r(y)=x}e_y.
 \tag{108.5}
\]

Los sumandos son ortogonales, de modo que \((r^*e_x)\) sigue siendo
idempotente. Además,

\[
 r^*\!\left(\sum_{x\in X_n}e_x\right)
 =\sum_{y\in X_{n+1}}e_y.
 \tag{108.6}
\]

La unidad global se conserva mientras el número de refinamientos aumenta.
Una trayectoria individual determina, a su vez, una familia compatible
\(((e_{x_n})_n)\). La primera es la identidad del álgebra; la segunda es la
unidad genealógica de una biografía. Esta distinción precisa la evolución
conservativa: el espacio total de posibilidades conserva su unidad y cada
ruta conserva su pertenencia mediante restricciones compatibles.

Un lector continuo (L) del estado profundo es equivalente a una familia de
lectores finitos \((L_n)\) que satisface

\[
 L_{n+1}\circ r_{n+1,n}^*=L_n.
 \tag{108.7}
\]

En consecuencia, cada dígito o coordenada determinada a profundidad finita
queda predefinido por un cilindro; aumentar la profundidad no recalibra ese
resultado, sino que determina nuevos cilindros compatibles. La predefinición
es la naturalidad de (108.7), no la inserción anticipada de un valor
experimental.

\textbf{Criterio de refutación.} La construcción falla si una restricción no
recupera el prefijo, si el \emph{pullback} no es multiplicativo, si la unidad no
se conserva o si un lector profundo modifica una coordenada ya fijada en un
cilindro anterior.

\section{Presentación universal del álgebra holonómica}

Sea \((mathcal C_n)\) la categoría de estados a profundidad (n) y de
historias admisibles generadas por

\[
 U_t={\rm Upd}_t\circ{\rm Tra}_t\circ{\rm Sel}_t.
 \tag{109.1}
\]

Para cada objeto (x), sea \((e_x)\) su identidad local y sea \((J_x)\) el
generador de la fibra TRIT:

\[
 J_x^2=-\tau(x)e_x,
 \qquad \tau(x)\in\{+1,0,-1\}.
 \tag{109.2}
\]

Para cada flecha \((u:x\to y)\), sea \((T_u)\) el transporte. El registro de memoria define
un multiplicador central (Omega(v,u)) en cada par componible. La
presentación finita es

\[
 \boxed{
 \mathfrak H_n
 =\mathbb R\langle e_x,J_x,T_u\rangle/\mathcal R_{\rm HMT},}
 \tag{109.3}
\]

donde

\[
 \begin{gathered}
 e_xe_y=\delta_{xy}e_x,
 \qquad
 T_u=e_{t(u)}T_ue_{s(u)},\\
 T_u a=\rho_u(a)T_u,
 \qquad
 T_vT_u=\Omega(v,u)T_{v\circ u}.
 \end{gathered}
 \tag{109.4}
\]

El producto se declara nulo cuando las flechas no son componibles. La
asociatividad equivale a

\[
 \rho_w(\Omega(v,u))\Omega(w,v\circ u)
 =\Omega(w,v)\Omega(w\circ v,u).
 \tag{109.5}
\]

En efecto, los dos productos \((((T_wT_v)T_u))\) y
\((T_w(T_vT_u))\) difieren únicamente en los dos lados de (109.5). La
identidad de cociclo prueba que toda historia posee una memoria independiente
de la parentización elegida.

La suma de (109.3) superpone historias alternativas. Su producto concatena
historias causalmente componibles. Esta distinción convierte la diferencia
suma--producto de APP en una aritmética dinámica: la suma no borra la
alternativa y el producto no borra la procedencia.

La conexión nonádica aparece en la sección elemental del cociente. Para
\((a,b\in\{0,\ldots,8\})\),

\[
 c_9(a,b)=\left\lfloor{a+b\over9}\right\rfloor
 \tag{109.6}
\]

satisface

\[
 c_9(a,b)+c_9(a+b\bmod9,c)
 =c_9(b,c)+c_9(a,b+c\bmod9).
 \tag{109.7}
\]

Si (X) representa un paso y (lambda) la memoria de una vuelta,

\[
 X^9=\lambda I,
 \qquad
 X^aX^b=\lambda^{c_9(a,b)}X^{a+b\bmod9}.
 \tag{109.8}
\]

La hélice queda condensada en (109.8): el exponente retorna; el coeficiente
recuerda cuántas veces se cruzó la frontera.

\section{Propiedad universal y espectro de realizaciones}

Sea \(({E_x})\) una familia de módulos, sean
\((pi_x:A_x\to\operatorname{End}(E_x))\) representaciones de las fibras
TRIT, y sean \((V_u:E_x\to E_y)\) transportes que cumplen

\[
 V_u\pi_x(a)=\pi_y(\rho_u(a))V_u,
 \tag{110.1}
\]

\[
 V_vV_u=\pi_z(\Omega(v,u))V_{v\circ u}.
 \tag{110.2}
\]

La asignación de los generadores define un homomorfismo del álgebra libre.
Las ecuaciones (110.1)--(110.2) hacen que ese homomorfismo anule el ideal
\((mathcal R_{\rm HMT})\). Existe entonces una única representación

\[
 \Pi:\mathfrak H_n\longrightarrow
 \operatorname{End}\!\left(\bigoplus_xE_x\right)
 \tag{110.3}
\]

que realiza los datos prescritos. La unicidad procede de que
\((e_x,J_x,T_u)\) generan el álgebra.

Este teorema proporciona el principio de ultrasimplificación. La física, la
teoría de números, el continuo y la biología no necesitan añadir generadores
al núcleo cuando sus lectores satisfacen las mismas relaciones. Cada uno es
una representación de (109.3) en un codominio diferente.

No todos los lectores son caracteres escalares. Un carácter
\((mathfrak H\to\mathbb R)\) anularía los conmutadores y no podría conservar el
bloque electrónico. Definimos, por ello, el espectro tipado

\[
 \operatorname{Read}(\mathfrak H)
 =\coprod_B\operatorname{Rep}_{\rm comp}(\mathfrak H,B).
 \tag{110.4}
\]

Una salida escalar adopta la forma

\[
 c=\sigma_B(\rho(x_\infty)),
 \tag{110.5}
\]

donde \(\rho\) conserva la genealogía y \(\sigma_B\) es el invariante propio
del codominio. Así quedan tipados cinco conceptos que antes parecían
heterogéneos:

\[
 \begin{array}{c|l}
 \text{número}&\text{coordenada de una representación arquimediana},\\
 \text{constante}&\text{invariante estable de retorno, escala o profundidad},\\
 \text{partícula}&\text{sección persistente de un módulo},\\
 \text{campo}&\text{representación que transporta secciones},\\
 \text{simetría}&\text{estabilizador de borde y relaciones conservadas}.
 \end{array}
 \tag{110.6}
\]

\section{Identidad exponencial de las tres curvaturas}

Introducimos el álgebra total

\[
 \mathbb A_{\rm tri}=A_+\oplus A_0\oplus A_-
 \tag{111.1}
\]

con idempotentes centrales \((p_+,p_0,p_-)\), y el generador

\[
 \mathbb J=p_+J_++p_0J_0+p_-J_-.
 \tag{111.2}
\]

La exponencial se define dentro del álgebra por

\[
 \operatorname{Exp}_\star(X)
 =\sum_{n\ge0}{X^{\star n}\over n!}.
 \tag{111.3}
\]

Las tres relaciones cuadráticas separan la serie en

\[
 \boxed{
 \begin{aligned}
 \operatorname{Exp}_\star(t\mathbb J)
 ={}&p_+(\cos t+J_+\sin t)\\
 &+p_0(1+tJ_0)\\
 &+p_-(\cosh t+J_-\sinh t).
 \end{aligned}}
 \tag{111.4}
\]

La primera fila es rotación; la segunda es transporte infinitesimal de
umbral; la tercera es dilatación. Una única ley de potencias contiene las
tres curvaturas.

La identidad de Euler ocupa la cara elíptica. En el cierre producido por
\((pi_{\rm HMT})\),

\[
 \boxed{
 p_+\operatorname{Exp}_\star(\pi\mathbb J)+p_+=0.}
 \tag{111.5}
\]

La proyección compleja conserva (111.5) y olvida las componentes parabólica e
hiperbólica. HMT conserva conjuntamente el cierre, el jet y la escala. Por
ello la identidad de Euler aparece como una cara exacta de la identidad
trítica, no como modelo exterior del generador.

El número \((e)\) ocupa una función distinta:

\[
 \operatorname{Exp}_\star(I)=e\,I.
 \tag{111.6}
\]

La ecuación define la función que convierte transporte aditivo en transporte
multiplicativo. La generación nonádica de sus cifras conserva su certificado
propio; (111.6) identifica su papel en el álgebra.

La autoescala \(\varphi\) ocupa la cara hiperbólica. Para

\[
 F_{\rm av}=\begin{pmatrix}0&1\\1&1\end{pmatrix},
 \qquad
 H_\varphi={2F_{\rm av}^2-3I\over\sqrt5},
 \tag{111.7}
\]

se verifica

\[
 H_\varphi^2=I,
 \tag{111.8}
\]

y

\[
 \boxed{
 F_{\rm av}^2
 =\operatorname{Exp}_\star(2\log\varphi\,H_\varphi).}
 \tag{111.9}
\]

Así, \(\pi\) mide el cierre elíptico, \(e\) realiza el paso exponencial y
\(\varphi\) mide la dilatación hiperbólica estable. Las igualaciones de los
tres canales son encuentros de estas funciones dentro de la misma órbita, no
igualdades universales entre sus valores reales.

\section{\(\alpha\) como transgresión dodecafásica}

La extensión no escindida produce el cociclo \(c_9\) de (109.6). Una vuelta
de fase determina un elemento del núcleo \(C_{12}\), y la palabra firmada
\(K\) conserva cómo se atravesaron sus sectores. Definimos el funcional de
costura

\[
 \mathcal T_\alpha:
 Z^2(C_9;C_{12})\times X_{12}^{\rm sgn}
 \longrightarrow\mathcal L_\alpha.
 \tag{112.1}
\]

La coordenada determinada es

\[
 \boxed{
 \alpha=\mathcal T_\alpha(c_9,K).}
 \tag{112.2}
\]

La fórmula no identifica el número \(\alpha\) con una clase de cohomología.
Identifica su tipo: es el escalar que resulta de evaluar la obstrucción de
escisión junto con la memoria dodecafásica.

La primera vía realiza (112.2) mediante

\[
 \iota(P)+\iota(E)-\iota(\Phi)-\iota(K)=\iota(A).
 \tag{112.3}
\]

La segunda vía realiza el mismo pegado sobre la rotación aperiódica:

\[
 P_9(\alpha)=\delta,
 \qquad
 \delta=\log_{10}{10\over9}.
 \tag{112.4}
\]

Su compatibilidad se expresa en el producto fibrado

\[
 \mathfrak A_\alpha
 =\{(x_{12},x):
 d_m(\Gamma_{12}^\alpha x_{12})=d_m(x),
 \ P_9(x)=\delta\}.
 \tag{112.5}
\]

Ahora puede formularse una matriz nuclear del sistema correlacionado:

\[
 \boxed{
 \begin{pmatrix}
 \operatorname{Exp}_\star(\pi J_+)&0&0&0\\
 0&\operatorname{Exp}_\star(I)&0&0\\
 0&0&\operatorname{Exp}_\star(2\log\varphi H_\varphi)&0\\
 0&0&0&\operatorname{Exp}_\star(P_9(\alpha)J_0)
 \end{pmatrix}
 =
 \begin{pmatrix}
 -I&0&0&0\\
 0&e\,I&0&0\\
 0&0&F_{\rm av}^2&0\\
 0&0&0&I+\delta J_0
 \end{pmatrix}.}
 \tag{112.6}
\]

La matriz no linealiza la genealogía. Expone cuatro funciones de una sola
álgebra: cierre, exponenciación, autoescala y costura. La memoria firmada
permanece en (112.3)--(112.5).

En el retículo \(A_2\), la misma costura se ve como

\[
 \partial_{A_2}N_8=\partial_{A_2}N_9=0,
 \qquad
 N_9-N_8=-16\mathbf1.
 \tag{112.7}
\]

\(\alpha\) determina el cierre de la parte relativa sin borrar la traslación
diagonal. Éste es el sentido exacto de su singularidad: es la coordenada que
distingue sincronización de reinicio.

\section{El espectro funcional de las constantes}

El lector dodecafásico uniforme determina además una sucesión de primeras raíces de retorno crítico. Se demuestra que su cociente de separaciones converge a la constante de Feigenbaum, conservando la selección original, la orientación y la memoria de la construcción. \eqref{hmtfg:articulacion:limite}.


La estructura anterior reorganiza las constantes ya derivadas por su función
en el álgebra, no por su pertenencia histórica a matemática o física.

La constante de Catalan es el momento de orden dos del carácter elíptico:

\[
 G_{\rm Cat}=\operatorname{Tr}(N^{-2}Q_4)
 =L(2,\chi_{-4}).
 \tag{113.1}
\]

Su residencia funcional conserva toda la familia
\(\mathcal C_2(s)\), no sólo el extremo \(s=1\). En
\cref{cr28:thm:catalan-recovery} se demuestra que la respuesta
constitutiva \(f(s)=(1-s^3)/(1-s^4)\), junto con la condición
\(\mathcal C_2(0)=0\), determina unívocamente el momento mediante
\(\mathcal D^2\mathcal C_2=s(1+s)f(s)/(1+s+s^2)\).
La inversión cúbica recupera, a su vez, los dos argumentos orientados.

El carácter de cuarto de giro y el carácter triádico se reúnen en

\[
 \chi_{12}=\chi_{-3}\chi_{-4}.
 \tag{113.2}
\]

Feigenbaum ocupa el lector no lineal. La rueda dodecafásica genera

\[
 u_0(x)=1-{1\over12}\sum_{m=1}^{12}
 \cos\!\left({2(3m-1)\over\sqrt{899}}x\right),
 \tag{113.3}
\]

con

\[
 899=30^2-1=29\cdot31,
 \qquad
 M_{30}=\begin{pmatrix}30&899\\1&30\end{pmatrix},
 \qquad \det M_{30}=1.
 \tag{113.4}
\]

La familia construida es analítica, par, unimodal, de criticidad cuadrática
y Schwarziano negativo en el dominio demostrado. El renormalizador de
duplicación lee la tasa de acumulación y la reescala espacial. \(\alpha_{\rm F}\)
es una coordenada de ese lector no lineal; \(\alpha\) es la costura
de (112.2). La distinción de subíndices expresa una diferencia de operación.

Boltzmann ocupa el transductor entre las rectas térmica y energética:

\[
 k_B:\mathcal L_\Theta\longrightarrow\mathcal L_E,
 \tag{113.5}
\]

\[
 k_B\Theta_{\rm clk}\log3
 =\hbar_{\rm int}{2\pi\over108t_0}.
 \tag{113.6}
\]

El \(\log 3\) cuantifica una unidad trítica; \(\log\varphi\) cuantifica una
autoescala hiperbólica. No son intercambiables. Para el transporte completo,

\[
 H(X\mid U(X))=0,
 \qquad \inf Q_L=0.
 \tag{113.7}
\]

Para un cociente que borra un trit,

\[
 Q_{\rm reset}\ge k_B\Theta\log3.
 \tag{113.8}
\]

El coste pertenece a la proyección reductora, no al avance de la holonomía.

Barbero--Immirzi ocupa el dual del módulo de incidencia \(90/120\). La
ecuación

\[
 {a\over270}-{b\over360}={1\over24}
 \tag{113.9}
\]

selecciona el covector entero mínimo \((a,-b)=(12,-1)\), y por tanto

\[
 K_{6\to8}=12\Delta S_{90}-\Delta S_{120}.
 \tag{113.10}
\]

\(\alpha\) cierra los relojes; \(\gamma_{\rm BI}\) evalúa la costura sobre el
defecto área--volumen. Esta diferencia funcional explica por qué ambos
objetos son capitales sin desempeñar la misma operación.

\section{Simetría, pentadicidad, continuo y secuencias}

La definición de simetría compatible con el núcleo es

\[
 \operatorname{Sym}(\rho,b_\partial)
 =\{g\in\operatorname{Aut}(\rho(\mathfrak H)):
 g\rho(\mathcal R_{\rm HMT})=\rho(\mathcal R_{\rm HMT}),
 \ gb_\partial=b_\partial\}.
 \tag{114.1}
\]

Las simetrías no son la magnitud primaria conservada. Son los automorfismos
que quedan disponibles después de exigir la conservación de relaciones y
frontera. La rama excepcional realiza (114.1) mediante la incidencia
Witt--Golay--Leech--\(V^\natural\)--Monstruo.

La pentadicidad posee dos representaciones:

\[
 S_{\pi,5}={1\over5}\mathbf1\mathbf1^{\mathsf T},
 \qquad \operatorname{rank}S_{\pi,5}=1,
 \tag{114.2}
\]

y

\[
 P_{G,5}:V_{12}\to V_5,
 \qquad \operatorname{rank}P_{G,5}=5.
 \tag{114.3}
\]

Ambos reciben la misma fuente holonómica, pero uno extrae la componente
común de cinco regiones y el otro conserva el portador gravitatorio de cinco
dimensiones. La cadena

\[
 V_{12}\to V_5\to {\rm STF}_3\to V_{\rm TT},
 \qquad 12\to5\to5\to2,
 \tag{114.4}
\]

demuestra cómo cinco modos estructurales determinan dos modos propagantes.

La doble proyección del continuo se expresa como

\[
 \rho_{\rm arq}:\mathfrak H_{\rm cyl}\to\mathcal A_{\mathbb R},
 \qquad
 \rho_{\rm exc}:\mathfrak H_{\rm cyl}\to\mathcal A_{\rm exc}.
 \tag{114.5}
\]

La primera contrae cilindros compatibles; la segunda conserva incidencia de
frontera. Los cinco aspectos espectral, solenoidal, gauge, coherente y
determinantal pertenecen al mismo dominio de (114.5). \(\mathbb R\) es la
salida de una representación conmutativa; la aritmética holonómica conserva
además los conmutadores y el registro de memoria que esa salida no ve.

La lectura de secuencias constituye otra representación. Con

\[
 A\mapsto(0,0),\quad G\mapsto(0,1),\quad
 T\mapsto(1,0),\quad C\mapsto(1,1),
 \tag{114.6}
\]

y fase \((n\bmod3)\), la reverso-complementación define una involución sobre la
categoría de palabras. APP evalúa acumulación y concatenación; TRIT tipa la
curvatura; TPK conserva retornos y frontera. La condición esturmiana

\[
 p_s(m)=m+1
 \tag{114.7}
\]

proporciona el umbral exacto entre periodicidad y complejidad exponencial.
De este modo, la idea de cristal aperiódico aplicada a secuencias biológicas
se convierte en un lector falsable de factores, extensiones y palabras de
retorno.

\section{Teorema nuclear de la aritmética holonómica HMT}

Las secciones anteriores permiten condensar la genealogía sin reducirla.
Definimos

 \[
 \boxed{
 \mathfrak H_{\rm HMT}^{\rm alg}
 =\left[\varinjlim_n
 \left(
 \mathbb R\langle e_x,J_x,T_u\rangle/
 \mathcal R_{\rm HMT}
 \right)\right]
 \rtimes_{\Gamma_9}\mathbb N.}
 \tag{115.1}
 \]

En el sector reversible, el último factor es \(\mathbb Z\). El espacio de
estados es \(X_\infty=\varprojlim X_n\). Para una representación compatible
y acotada \(\rho\), la completación
\(\mathfrak H_{\rm HMT}^{\rho}
=\overline{\rho(\mathfrak H_{\rm HMT}^{\rm alg})}^{\|\cdot\|}\)
pertenece al lector; el generador universal no presupone una norma.

\begin{quote}
\textbf{Teorema nuclear.} La doble evaluación APP, la orientación trigeométrica
TRIT y la dinámica TPK generan una aritmética asociativa no conmutativa de
historias, torcida por el cociclo de acarreo y graduada por fase y memoria.
El refinamiento conserva la unidad, la holonomía retorna la fase mientras
avanza la memoria y el espectro tipado de representaciones determina la
sistema correlacionado \((\pi,\varphi,e,\alpha)\), la estructura discreta conjunta
del continuo, el electrón, las escalas de acción y gravedad, los parámetros
geométricos, las constantes derivadas y las simetrías de frontera.
\end{quote}

La demostración se obtiene concatenando cinco hechos ya establecidos:

\begin{enumerate}
\item las ecuaciones APP conservan las dos hojas, el residuo y el cociente;
\item \(J_\tau^2=-\tau I\) conserva las tres curvaturas dentro de un solo TRIT;
\item la identidad de cociclo prueba la asociatividad del producto de historias;
\item la extensión no escindida prueba que retorno de fase y retorno de estado son distintos;
\item la propiedad universal prueba que todo lector compatible factoriza por (115.1).
\end{enumerate}

La primera consecuencia es que \(\mathbb R\) no constituye el soporte
fundamental. Es una representación terminal de una estructura que conserva
más información que su coordenada real. La segunda es que una constante no
es una cifra almacenada: es un invariante de un lector. La tercera es que una
partícula no es un número: es una sección persistente de un módulo. La cuarta
es que una simetría no se postula: aparece como estabilizador de una frontera
conservada.

La forma más breve de la evolución conservativa es

\[
 \boxed{
 \underbrace{\varprojlim X_n}_{\text{historia que se profundiza}}
 \quad\longleftrightarrow\quad
 \underbrace{\varinjlim\mathfrak H_n}_{\text{lectores que se expanden}}.}
 \tag{115.2}
\]

El cristal aperiódico helicoidal coordina ambos límites. La identidad
trigeométrica (111.4) aporta su cálculo local; \(\alpha\) aporta la costura de
memoria; el bloque \({\rm Cl}_{1,1}\) aporta su rincón electrónico; el
invariante \((G\hbar)\) aporta su conservación dimensional; Barbero--Immirzi
aporta su covector de frontera; y la doble proyección aporta sus salidas
arquimediana y excepcional. Ésta es la estructura mínima desde la que puede
desplegarse el artículo autónomo sin depender expositivamente de las
(2\,099) páginas y sin perder ninguno de sus operadores esenciales.

\textbf{Criterio de refutación.} El teorema queda refutado si falla la identidad de
cociclo, si el refinamiento deja de ser unital, si la holonomía reinicia la
memoria, si desaparece una fibra TRIT, si una realización no satisface las
relaciones que declara representar o si una salida utiliza su propio valor
convencional para seleccionar la historia que debía producirla.

\section{Irreducibilidad triádica y diferencia de potencial en lectores de secuencia}

La composición APP--TRIT--TPK puede probarse irreducible mediante tres
cocientes de olvido. Sea \(q_{\rm APP}\) el cociente que identifica las
hojas aditiva y multiplicativa. Entonces

\[
q_{\rm APP}(\Delta_{\Sigma\Pi})=0,
\qquad
q_{\rm APP}([P_\Sigma,P_\Pi])=0.
\tag{116.1}
\]

Sea \(q_{\rm TRIT}\) el cociente que identifica los tres idempotentes de
régimen. Se pierde entonces la terna de cuadrados

\[
(J_+^2,J_0^2,J_-^2)=(-I,0,I).
\tag{116.2}
\]

Sea \(q_{\rm TPK}\) el cociente que trivializa el cociclo de acarreo y la
graduación de memoria. En ese cociente,

\[
q_{\rm TPK}\!\left(d_M(\Gamma_9x)-d_M(x)\right)=0,
\tag{116.3}
\]

aunque en el estado holonómico completo la diferencia vale \(1\). Los tres
invariantes (116.1)--(116.3) son de tipos distintos y ninguno puede
reconstruirse a partir de los otros dos después de aplicar su cociente. APP
aporta discrepancia y orden de hojas; TRIT, firma de curvatura; TPK,
composición, holonomía y memoria. Ésta es la forma algebraica exacta del
carácter ternariamente irreducible del núcleo.

La misma separación clarifica la aplicación a secuencias. En el alfabeto
\(\{A,G,T,C\}\), defínanse

\[
\chi_p(A,G,T,C)=(0,0,1,1),
\qquad
\chi_h(A,G,T,C)=(0,1,0,1),
\tag{116.4}
\]

y la involución de complemento
\(\mathcal C_{\rm DNA}(A,G,T,C)=(T,C,A,G)\). Se verifica

\[
\chi_p(\mathcal C_{\rm DNA}b)=1-\chi_p(b),
\qquad
\chi_h(\mathcal C_{\rm DNA}b)=\chi_h(b).
\tag{116.5}
\]

Por consiguiente,

\[
\mathcal V_n=\chi_p(b_{n+1})-\chi_p(b_n)
\tag{116.6}
\]

es una diferencia de potencial simbólica orientada, mientras \(\chi_h\)
conserva la clase transversal de enlace. La elevación

\[
b_n\longmapsto
\bigl(\chi_p(b_n),\chi_h(b_n),n\bmod3,\lambda_n,c_n,m_n\bigr)
\tag{116.7}
\]

añade régimen TRIT, frontera, acarreo y memoria TPK. Así queda formulada la
razón estructural por la que polaridad, potencial y palabra aperiódica
pueden ser leídos dentro de la misma aritmética sin convertir la biología en
generador del formalismo.

\textbf{Criterio de refutación.} La irreducibilidad falla si uno de los tres
invariantes sobrevive al cociente que debe borrarlo o puede reconstruirse de
los otros dos sin reintroducir el dato eliminado. El lector de secuencias
falla si el complemento no invierte \(\chi_p\), no conserva \(\chi_h\) o si
la elevación no entrelaza la fase ternaria con la dinámica TPK.

\section{La vacancia electrónica como par nulo de Witt}

La aritmética del rincón electrónico permite completar localmente los tres
regímenes sin introducir otro álgebra. Con

\[
S=2P-I,\qquad
J=\frac{[P,Q]}{\sqrt3/4},
\tag{117.1}
\]

se tiene

\[
S^2=I,\qquad J^2=-I,\qquad SJ=-JS.
\tag{117.2}
\]

Las combinaciones

\[
n_+=\frac{S+J}{2},\qquad
n_-=\frac{S-J}{2}
\tag{117.3}
\]

satisfacen

\[
n_+^2=n_-^2=0,\qquad
n_+n_-+n_-n_+=I.
\tag{117.4}
\]

Por tanto,

\[
\{J,n_\pm,S\}
\quad\longmapsto\quad
\{\text{elíptico},\text{parabólico},\text{hiperbólico}\}
\tag{117.5}
\]

es una realización local completa del tipado TRIT dentro de
\({\rm Cl}_{1,1}\cong M_2(\mathbb R)\). La vacancia queda caracterizada por
dos descripciones compatibles: en APP es la pérdida unitaria del
cociente-producto que conserva hoja aditiva, residuo producto y régimen; en
el rincón de Clifford es una dirección nula parabólica. La pareja
\((n_+,n_-)\) es el par de Witt local de las dos orientaciones. Esta
descomposición ofrece el dominio algebraico preciso desde el que la
incidencia Paley--Witt posterior puede realizar la misma genealogía en su
propio módulo combinatorio.

La afirmación es más fuerte que la analogía con cuaterniones. Los
cuaterniones de Hamilton poseen tres generadores con cuadrado \(-I\) y forman
una álgebra de división. El bloque (117.2) es escindido: contiene divisores
de cero, idempotentes y direcciones nulas, justamente los objetos requeridos
por hoja, vacancia y umbral. Los octoniones amplían dimensión sacrificando
asociatividad; la aritmética holonómica mantiene asociatividad mediante el
cociclo y amplía la estructura por rutas, memoria y refinamiento.

\textbf{Criterio de refutación.} Basta que falle una de las cuatro identidades
(117.2)--(117.4), que la vacancia no preserve sus datos APP declarados o que
el transporte TPK no conserve el idempotente de régimen.

\section{Relaciones de Weyl del reloj finito y del reloj irracional}

Sea \(T\) el paso de la órbita, sea
\(\delta=\log_{10}(10/9)\), sean \(\zeta_9\) y \(\omega_\delta\)
los caracteres internos de fase nonádica y capacidad, y definamos sobre
\(\lvert t,m\rangle\)

\[
V_9\lvert t,m\rangle=\zeta_9^t\lvert t,m\rangle,
\qquad
W_\delta\lvert t,m\rangle=\omega_\delta^t\lvert t,m\rangle.
\tag{118.1}
\]

El avance temporal satisface

\[
V_9T=\zeta_9TV_9,\qquad
W_\delta T=\omega_\delta TW_\delta.
\tag{118.2}
\]

La vuelta \(\Gamma_9=T^9\) produce

\[
V_9\Gamma_9=\Gamma_9V_9,\qquad
W_\delta\Gamma_9=\omega_\delta^9\Gamma_9W_\delta.
\tag{118.3}
\]

Así se distinguen algebraicamente los dos relojes. El primero posee período
nueve; el segundo conserva una fase irracional. La palabra mecánica de
vacancias es la codificación por umbral de la órbita del segundo reloj. Si
\(D_M\) es el operador de grado de memoria, el lift TPK verifica

\[
[D_M,\Gamma_9]=\Gamma_9.
\tag{118.4}
\]

La representación elíptica posterior reconoce estos caracteres mediante

\[
\rho_+(\zeta_9)
=\operatorname{Exp}_+\!\left(\frac{2\pi}9J_+\right),
\qquad
\rho_+(\omega_\delta)
=\operatorname{Exp}_+\!\left(2\pi\delta J_+\right).
\tag{118.5}
\]

Los caracteres, y no sus expresiones complejas reconocidas, son los datos
operatorios del reloj.

La ecuación (118.4) dice exactamente que \(\Gamma_9\) es un operador de
elevación: la fase vuelve y el estado asciende una unidad en la hélice. La
acción de traslación es abeliana; el producto cruzado de esa acción con los
observables \(V_9,W_\delta,D_M\) es no conmutativo por (118.2)--(118.4).
La estructura temporal no requiere, por ello, elegir entre «reloj finito» y
«cristal aperiódico»: ambos son caracteres conjugados de un único sistema
sesgado.

La entropía topológica nula procede de la rotación irracional y la
complejidad \(p(m)=m+1\); el crecimiento de memoria procede de (118.4).
Entropía nula y biografía infinita coexisten porque la segunda mide
profundidad reconstructible, mientras la primera mide crecimiento
exponencial de factores.

\textbf{Criterio de refutación.} El modelo falla si \(V_9\) no conmuta con
\(\Gamma_9\), si \(W_\delta\) adquiere período finito, si el conmutador de
memoria difiere de \(\Gamma_9\) o si la palabra umbral deja de tener
complejidad esturmiana.

\section{Ecuación mínima del generador trigeométrico}

Sea

\[
\mathbb J=p_+J_++p_0J_0+p_-J_-,
\qquad p_++p_0+p_-=I.
\tag{119.1}
\]

Las tres relaciones de curvatura equivalen a

\[
\boxed{\mathbb J^6=\mathbb J^2.}
\tag{119.2}
\]

En efecto, \(K=\mathbb J^2\) satisface \(K^3=K\), y los tres idempotentes se
recuperan sin datos externos:

\[
p_+=\frac{K^2-K}{2},\qquad
p_0=I-K^2,\qquad
p_-=\frac{K^2+K}{2}.
\tag{119.3}
\]

La identidad exponencial intrínseca es

\[
\operatorname{Exp}_\star(t\mathbb J)=
p_+(\cos t+\mathbb J\sin t)+
p_0(I+t\mathbb J)+
p_-(\cosh t+\mathbb J\sinh t).
\tag{119.4}
\]

Las ecuaciones (118.4) y (119.4) se reúnen en el propagador holonómico

\[
\boxed{
\mathcal E_{\rm HMT}(t)
=\Gamma_9\operatorname{Exp}_\star(t\mathbb J),
\qquad
[D_M,\Gamma_9]=\Gamma_9.}
\tag{119.5}
\]

Ésta es la compresión operatoria buscada. La primera coordenada de (119.5)
gobierna la respuesta local en las tres curvaturas; la segunda gobierna la
elevación de memoria. APP entra mediante los idempotentes de hoja, sus
proyectores y el cociclo que define \(\Gamma_9\); TRIT entra por
\(\mathbb J\); TPK entra por la composición, el lift y \(D_M\). La
sistema correlacionado, el electrón, la acción, la gravedad, Barbero--Immirzi, las
constantes de escala, las simetrías y los lectores de secuencia se obtienen
como representaciones tipadas del álgebra universal que contiene (119.5).

La función de \(\alpha\) se vuelve transparente en esta presentación.
\(\operatorname{Exp}_\star\) contiene la dinámica local y \(\Gamma_9\) la
vuelta con memoria; el cierre dodecafásico evalúa su compatibilidad y determina
\(\alpha\) como costura. Así, \(\alpha\) enlaza la sincronización
de \((\pi,e,\varphi)\) con los lectores dimensionales porque mide una clase
de cierre del propagador holonómico, no porque se agregue a tres constantes
preexistentes.

\textbf{Criterio de refutación.} La ecuación mínima falla si
\(\mathbb J^6\ne\mathbb J^2\), si los polinomios (119.3) dejan de ser
idempotentes ortogonales, si (119.4) no se restringe a las tres
exponenciales certificadas o si \(\alpha\) no coincide en sus dos vías
tipadas dentro del cuadrado de compatibilidad existente.
 
% Realizaciones dimensionales completas tomadas de la autoridad material viva.
% El orden respeta la precedencia electrón estructural -> acción -> constitución
% del vacío -> gravitación -> área e información.
% \newcommand{\HMTMonographChapterInput}[2]{%
%   \chapter{#1}%
%   \begingroup
%   \renewcommand{\chapter}[2][]{}%
%   \input{#2}%
%   \endgroup}

\chapter{El electrón como sección persistente de la vacancia orientada}
\begingroup
\renewcommand{\chapter}[2][]{}%
\chapter{El electrón como 1.ª realización material: etapa basal, clausura beta y estructura atomística de la electricidad}
\label{chap:parte-iv-48}

El estado genealógico adquiere aquí su primera realización material completa. La localización central, la banda electrón--positrón y la lectura bidireccional se construyen antes de toda comparación metrológica; la etapa basal queda distinguida de la clausura beta rectora.

\section{Clausura \(\beta\) de la sección electrónica y publicación dimensional de su energía de reposo}
\label{chap:biografia-electron-rev6}
\index{electrón!construcción holográfica}
\index{positrón!banda conjugada}
\index{masa!ley electrónica}

El electrón aparece en HMT como el punto en el que se ensamblan seis
construcciones anteriores y mutuamente tipadas: la incompatibilidad entre
las dos hojas aritméticas de APP, la orientación inducida por la involución
de la bisagra areal--volumétrica, dos cocientes enteros del censo APP--TPK,
el cierre
dodecafásico de \(\alpha\), el defecto global
\(\Delta_4\) y el centro único del atlas nonádico. La banda de ruta añade
después la conjugación electrón--positrón. La carta en MeV y la comparación
con CODATA son posteriores a todas ellas.

Esta precedencia forma parte del contenido matemático: los factores, sus
tipos y su orden causal se determinan antes de abrir la carta tabulada, que
interviene únicamente en el contraste metrológico posterior.

\subsection{Escalar electrónico \(\mathfrak e\) y grafo causal de sus factores APP–TPK, \(\alpha\), \(\Delta_4\) y atlas nonádico}

La salida interna es el escalar adimensional
\begin{equation}
 \label{eq:electron-holografico-rev6}
 \boxed{
 \mathfrak e
 =
 \frac{\sqrt3}{4}
 \left[
 \exp\!\left(\frac{\varphi}{\pi^2}\right)
 -\frac{396}{18}\,\alpha^3
 \right]
 \left[
 1+\frac{270}{18}\,\Delta_4
 \right].}
\end{equation}
Las identidades
\[
 \frac{396}{18}=22=27-5,
 \qquad
 \frac{270}{18}=15=\binom62=3\cdot5
 \]
permiten recuperar la escritura abreviada
\[
 \mathfrak e
 =
 \frac{\sqrt3}{4}
 \left(e^{\varphi/\pi^2}-22\alpha^3\right)
 (1+15\Delta_4).
\]
Aquí la letra \(e\) dentro de la exponencial es el número de Euler; el
subíndice de \(\mathfrak e\) designa la sección electrónica.

El grafo causal de la ecuación es
\[
\begin{array}{ccl}
\text{proyectores APP en dimensión tres}
&\longrightarrow& \sqrt3/4,\\[2mm]
\text{calendario TPK y hojas areal--volumétrica}
&\longrightarrow&
F_{\mathrm{av}},\ \pi/\varphi^2,\
\varphi/\pi^2,\\[2mm]
\text{cierre dodecafásico}
&\longrightarrow&\alpha,\\[2mm]
\text{censo APP--TPK y familia electrónica}
&\longrightarrow&396,\ 18,\ 270,\ 22,\ 15,\ 5,\\[2mm]
\text{lecturas globales de }\pi,e\text{ y registro}
&\longrightarrow&\Delta_4,\\[2mm]
\text{atlas nonádico y anti-sección}
&\longrightarrow&(5,5),\
\text{banda electrón--positrón}.
\end{array}
\]
Ninguna flecha de este diagrama recibe como argumento la masa electrónica
de referencia.

\subsection{1.er factor: incompatibilidad exacta de las hojas APP}

Para \(d\geq2\), sea
\[
 \mathscr A_d=\{1,\ldots,9\}^d
\]
con medida uniforme, y sean
\[
 \Sigma_d(x)=\operatorname{dr}(x_1+\cdots+x_d),
 \qquad
 \Pi_d(x)=\operatorname{dr}(x_1\cdots x_d)
\]
los observables aditivo y multiplicativo. Sus esperanzas condicionales
ortogonales sobre \(\ell^2(\mathscr A_d)\) son
\(P_{\Sigma,d}\) y \(P_{\Pi,d}\). La
\cref{prop:app-no-conmutatividad} demuestra por enumeración exacta que
\[
 \left\|[P_{\Sigma,3},P_{\Pi,3}]\right\|
 =\frac{\sqrt3}{4}.
\]

\begin{proposition}[Contenido APP del prefactor electrónico]
\label{prop:prefactor-app-electron-rev6}
El prefactor \(\sqrt3/4\) de
\cref{eq:electron-holografico-rev6} es la norma exacta de la
no conmutatividad entre las hojas suma y producto en dimensión tres.
No es un factor geométrico introducido desde la fórmula electrónica.
\end{proposition}

\begin{proof}
La tabla conjunta de \(\Sigma_3\) y \(\Pi_3\) determina los ángulos
principales entre
\(\operatorname{Ran}P_{\Sigma,3}\) y
\(\operatorname{Ran}P_{\Pi,3}\). Si \(s\) es uno de sus valores
singulares, la contribución correspondiente a la norma del conmutador es
\(s\sqrt{1-s^2}\). La enumeración de las \(9^3\) palabras da como máximo
\[
 s^2(1-s^2)=\frac3{16},
 \qquad s^2=\frac14.
\]
Por tanto la norma máxima es \(\sqrt{3/16}=\sqrt3/4\), antes de
introducir \(\alpha\), \(\Delta_4\), una unidad de energía
o una masa de referencia.
\end{proof}

La función física del factor queda así fijada: mide la incompatibilidad
máxima de las dos lecturas aritméticas sobre el mismo soporte ternario. El
electrón no hereda dos números independientes, uno aditivo y otro
multiplicativo; hereda la obstrucción exacta a hacer simultáneamente
compatibles las dos hojas. El prefactor conserva en la ley de masa la
memoria local del desacuerdo APP del que nace la ruta.

\subsection{2.º factor: la involución de intercambio entre \(\pi\) y \(\varphi\)}

El calendario primitivo de modos
\((\Sigma,\Pi,\Pi)\) induce, en las hojas areal y volumétrica, la matriz
\[
 F_{\mathrm{av}}
 =
 \begin{pmatrix}
 0&1\\
 1&1
 \end{pmatrix}.
\]
El \cref{thm:descenso-necesario-fav} la obtiene contando las transiciones
\(\Sigma\to\Pi\), \(\Pi\to\Pi\) y \(\Pi\to\Sigma\), cada una una vez.
No se postula la matriz a partir de la razón áurea: sus autovalores
\(\varphi\) y \(-\varphi^{-1}\) aparecen después de haber construido la
incidencia.

En la memoria de segundo orden,
\(\operatorname{Sym}^2(F_{\mathrm{av}})\) tiene espectro
\[
 \varphi^2,\qquad -1,\qquad\varphi^{-2}.
\]
Así, \(\varphi^{-2}\) es el modo estable positivo de la bisagra. La marca
de clausura \(\pi\), aplicada una sola vez al retorno areal, produce
\[
 \frac{\pi}{\varphi^2}.
\]
Esta razón conserva el orden cierre--autoescala. La ley electrónica usa su
orientación contraria.

\begin{theorem}[Involución exacta de la bisagra electrónica]
\label{thm:espejo-electron-rev6}
Sean
\[
 \mathscr R(x,y)=\frac{x}{y^2},
 \qquad
 \mathsf J(x,y)=(y,x).
\]
Entonces \(\mathsf J^2=I\) y
\[
 \mathscr R(\pi,\varphi)=\frac{\pi}{\varphi^2},
 \qquad
 (\mathscr R\circ\mathsf J)(\pi,\varphi)
 =\frac{\varphi}{\pi^2}.
\]
Por consiguiente,
\[
 \frac{\varphi}{\pi^2}
 =
 \frac{1}{
 \varphi^3\bigl(\pi/\varphi^2\bigr)^2},
 \qquad
 e^{\varphi/\pi^2}
 =
 \exp\!\left[
 \frac{1}{
 \varphi^3\bigl(\pi/\varphi^2\bigr)^2}
 \right].
\]
\end{theorem}

\begin{proof}
La primera pareja de identidades se obtiene aplicando
\(\mathscr R\) antes y después de la involución \(\mathsf J\). Para la
reparametrización,
\[
 \varphi^3
 \left(\frac{\pi}{\varphi^2}\right)^2
 =\frac{\pi^2}{\varphi},
\]
y su inversa es \(\varphi/\pi^2\). La última igualdad se sigue al aplicar
la exponencial.
\end{proof}

La presencia simultánea de \(\pi\) y \(\varphi\) tiene, por tanto, un
significado estructural preciso. \(\pi/\varphi^2\) lee el cierre areal
sobre la memoria volumétrica cuadrática; \(\varphi/\pi^2\) lee el
autocrecimiento interior sobre el cierre cuadrático. No son dos cocientes
parecidos elegidos por conveniencia. Son las dos orientaciones de un mismo
lector, intercambiadas por una involución exacta. El término
\(e^{\varphi/\pi^2}\) publica en la sección electrónica la orientación
interior de esa bisagra.

\subsection{3.er y 4.º factores: los enteros 22, 15 y el centro 5}

El censo de la carta APP empleada por la familia electrónica proporciona
tres sumas enteras:
\[
 \Sigma(\mathrm{APP})=396,
 \qquad
 \Sigma(\mathrm{centro}\ 2\times2)=18,
 \qquad
 \Sigma(\mathrm{canal})=270.
\]
La primera admite la descomposición por anillos
\[
 396=18+72+108+198.
\]
El núcleo \(2\times2\) es, por tanto, una subestructura explícita de la
carta completa, y el entero 270 pertenece al canal global del mismo censo
APP--TPK.
La suma de los anillos pares es
\[
 18+108=126,
 \qquad
 \frac{396}{126}=\frac{22}{7}.
\]
Este último cociente es el sello racional de cierre de la carta, no una
identificación de \(22/7\) con el número \(\pi\) generado por el lector
coinductivo. Sí muestra que el numerador \(22\) ya ocupa una función
interna antes de aparecer en la ley electrónica.

\begin{theorem}[Procedencia interna de los coeficientes electrónicos]
\label{thm:coeficientes-electron-rev6}
En la familia electrónica fijada por la carta APP y su canal global,
\[
 \boxed{
 22=\frac{\Sigma(\mathrm{APP})}
 {\Sigma(\mathrm{centro})}
 =\frac{396}{18}=27-5,}
\]
\[
 \boxed{
 15=\frac{\Sigma(\mathrm{canal})}
 {\Sigma(\mathrm{centro})}
 =\frac{270}{18}=3\cdot5=\binom62.}
\]
En particular, los coeficientes se determinan sin consultar
\(Z_0\), la masa del electrón ni una carta de unidades.
\end{theorem}

\begin{proof}
Las dos primeras igualdades son divisiones exactas del censo finito:
\[
 396=22\cdot18,
 \qquad
 270=15\cdot18.
\]
Al definir
\[
 p:=27-22,
\]
se obtiene \(p=5\), y entonces
\[
 15=3p=3\cdot5=\binom62.
\]
Todas las cantidades pertenecen al dominio entero de la carta y de la
familia antes de la transducción metrológica.
\end{proof}

\subsubsection{Selector de persistencia en el retorno central}

La misma pareja posee una construcción dinámica que fija además su signo sin
consultar una masa. Las cuatro orientaciones de la semilla central
\((5,5)\) regresan por primera vez a esa celda tras veintisiete pasos. El
espacio tipado del retorno es
\[
 \mathcal K_e=\mathbb Z/9\mathbb Z\times\mathbb F_3,
 \qquad |\mathcal K_e|=27.
\]
Por otra parte, la microfibra de \(\pi\) contiene las cinco regiones
genealógicamente distintas
\[
 \mathscr R_\pi=3_{++}+1_{--}+1_{\perp}.
\]
El calibre regional \(\operatorname{diag}_c\) y, dentro de una fase repetida,
el orden interno de \(U_6\), definen la inyección
\[
 \iota:\mathscr R_\pi\hookrightarrow\mathcal K_e
\]
mediante los cinco lugares distintos
\[
 (4,0),\quad(5,0),\quad(6,0),\quad(6,1),\quad(6,2).
\]
Los dos primeros corresponden, respectivamente, a la región transversal y
al retorno mutado; los tres últimos separan las regiones positivas por su
hoja trítica. La palabra visible común no interviene como inversa del mapa.

\begin{theorem}[Selector interno de los coeficientes electrónicos]
\label{thm:selector-persistencia-electron-rev7}
Fijados el retorno central y las cinco regiones tipadas, los coeficientes de
la ecuación electrónica son
\[
 \boxed{
 n_e=-\bigl|\mathcal K_e\setminus\iota(\mathscr R_\pi)\bigr|=-22,
 \qquad
 m_e^{\rm tor}=\bigl|\mathscr R_\pi\times\mathbb F_3\bigr|=+15.}
\]
El primer signo registra el déficit de frontera; el segundo, la memoria
torsional retenida.
\end{theorem}

\begin{proof}
Los cinco lugares anteriores son distintos, de modo que \(\iota\) es
inyectiva y
\[
 |\mathcal K_e\setminus\iota(\mathscr R_\pi)|=27-5=22.
\]
Cada una de las cinco regiones conserva tres hojas tríticas, luego
\[
 |\mathscr R_\pi\times\mathbb F_3|=5\cdot3=15.
\]
La regla de orientación asigna signo negativo a los estados retirados de la
cáscara y positivo a la memoria retenida. La construcción sólo abre el
centro, el retorno de veintisiete estados y las identidades regionales; no
abre una masa observada, una impedancia ni una firma calibrada.
\end{proof}

El número \(22\) mide cuántos núcleos centrales de peso \(18\) recompone la
suma APP \(396\). El número \(15\) mide cuántos de esos mismos núcleos
recompone el canal \(270\), y simultáneamente cuenta los pares de una
hexada. El número \(5\) no es un parámetro adicional:
\[
 5=27-22,\qquad 15=3\cdot5.
\]
Las tres escrituras muestran el mismo cierre desde la profundidad 27,
desde la normalización central y desde la combinatoria hexádica.

\subsubsection{Conservación de los coeficientes 22 y 15 en la realización material del electrón}

La exploración histórica localizó primero \(p=5\) mediante una comparación
finita con \(Z_0\), y de allí escribió
\((27-p,3p)=(22,15)\). Ese procedimiento conserva valor como historia de
descubrimiento y como control externo. No es, sin embargo, la demostración
que gobierna la ley: el recuento posterior
\[
 \frac{396}{18}=22,
 \qquad
 \frac{270}{18}=15
\]
construye ambos enteros dentro de APP--TPK y recupera \(p=5\) como
consecuencia. La historia explica cómo se encontró la pareja; el censo
explica por qué pertenece a la ecuación.

\subsection{5.º factor: el cubo del cierre \(\alpha\)}

El cierre de \(\alpha\) precede a Mecánica Dimensional.
Como se establece en el \cref{chap:alpha}, los doce bloques en base mil
de \(\pi\), \(e\), \(\varphi\) y del sello \(K\) se combinan mediante el
acarreo reversible
\[
 P_m+E_m-\Phi_m-K_m+c_{m+1}
 =A_m+1000c_m,
 \qquad c_{13}=c_1=0.
\]
La salida es
\[
 A_\alpha=
 (007,297,352,569,283,800,997,285,105,472,380,663)
\]
y, por concatenación, la ventana dodecafásica exacta
\[
 \alpha_{12}=\pi_{12}(\alpha)
 =
 \frac{
 7297352569283800997285105472380663
 }{10^{36}}.
\]
El telescopado de los acarreos demuestra la identidad entera completa; la
prolongación compatible construye
\(\alpha=\varprojlim_N A^{(N)}\), y \(\alpha_{12}\) es su
proyección racional de treinta y seis cifras. La
comparación con la constante de estructura fina se realiza después y no
interviene en su construcción.

La ley electrónica toma la tercera potencia
\[
 \alpha^3.
\]
Matemáticamente, se trata de la evaluación multiplicativa de la tercera
potencia simétrica del mismo escalar ya cerrado; no se introducen tres
copias metrológicas ni tres parámetros diferentes. Su contribución es
\[
 22\alpha^3
 =
 0.00000854906599200551727014979807598228\ldots.
\]

\begin{proposition}[Estatuto del canal cúbico electrónico]
\label{prop:alpha-cubica-electron-rev6}
El término \(22\alpha^3\) de la ley electrónica depende
únicamente del cierre previo de \(\alpha\), de la
multiplicación interna y del entero \(22\) del
\cref{thm:coeficientes-electron-rev6}. No depende de una lectura de la
masa electrónica.
\end{proposition}

\begin{proof}
\(\alpha_{12}\) es el racional fijado por la identidad de concatenación y
acarreo, y la familia compatible de ventanas determina la sección
\(\alpha\). La proyección \(\alpha_{12}^3\) certifica las
cifras impresas de \(\alpha^3\).
El coeficiente \(22\) está determinado por \(396/18\). La composición de
ambas operaciones contiene sólo datos construidos con anterioridad a la
carta de energía.
\end{proof}

Esta proposición fija exactamente lo que se usa en el electrón. La potencia
tres es el grado cúbico de este canal en la ley publicada. No es necesario
postular aquí una regla universal según la cual toda aparición
\(\alpha^d\) represente automáticamente una dimensión \(d\); una
afirmación universal de ese tipo exigiría su propio teorema de
multiplicatividad entre grados. La genealogía electrónica queda completa
sin convertir esa generalización en premisa.

\subsection{6.º factor: torsión mínima y defecto global \(\Delta_4\)}
\index{torsión mínima!\(\Delta_4\)}

Las denominaciones \emph{torsión mínima} y \emph{defecto global}
designan aquí un único invariante, no dos correcciones distintas. El objeto
empleado por la ley de masas es
\begin{equation}
 \label{eq:delta4-electron-rev6}
 \boxed{
 \Delta_4
 =
 \frac{\pi-e\log\pi}{270}
 \left(1+\frac{\pi}{729}\right).}
\end{equation}
Sus entradas son las lecturas HMT ya fijadas de \(\pi\) y \(e\), el canal
global \(270\) y la celda interior \(729\). Por consiguiente,
\[
 \Delta_4
 =
 0.00011119642269214005405113478612067909\ldots.
\]
La primera razón mide el defecto firmado
\(\pi-e\log\pi\) en unidades del canal 270; el segundo factor acopla esa
razón a la celda interior mediante \(\pi/729\). Esta lectura no añade
parámetros a la fórmula: separa sus dos operaciones constitutivas.
El \cref{p3-20260826:chap:medida-accion-positiva} muestra que es una entrada matemática
pre-dimensional y la distingue del operador que eventualmente la realiza
como cambio de hoja.

\begin{proposition}[Globalidad del defecto electrónico]
\label{prop:delta4-global-electron-rev6}
\(\Delta_4\) es un invariante global anterior a la especie. En la ley
electrónica entra mediante
\[
 1+15\Delta_4
 =
 1.00166794634038210081076702179181019\ldots.
\]
No es un corrector elegido después de conocer el residuo de la masa del
electrón.
\end{proposition}

\begin{proof}
La \cref{eq:delta4-electron-rev6} no contiene una etiqueta de especie, una
masa ni una unidad. El entero \(15\) ya está construido por \(270/18\).
Por tanto la composición \(1+15\Delta_4\) queda determinada por los datos
globales y el censo interno.
\end{proof}

Debe distinguirse \(\Delta_4\) de los otros bloques del carácter de masa y de
las otras nociones de torsión. En particular, no es \(s_{270}\), ni el
invariante cuadrático \(135\Delta_4^2\), ni el residuo angular
\(\varepsilon_{\rm res}\), ni la torsión geométrica de Cartan--Holst.
Sustituir uno por otro repetiría o cambiaría el nivel de la corrección. En el
electrón, \(\Delta_4\) conserva la memoria global del refinamiento; el factor
\(15\) declara cuántas veces la familia electrónica la lee. El calificativo
\emph{mínima} conserva el nombre histórico de este invariante y no se usa
como sustituto de un teorema variacional de minimalidad.

\subsection{Realización material del punto fijo (5,5) en la clausura electrónica}

Sea
\[
 \mathcal C_{81}=\{1,\ldots,9\}^2
\]
el atlas nonádico y
\[
 \rho(r,c)=(10-r,10-c)
\]
su inversión central. El \cref{thm:censo-atlas-81} prueba que el atlas
contiene una única celda de generación cero y que ésta porta la etiqueta
posicional \(e_{\mathrm{core}}\).

\begin{proposition}[Unicidad del centro electrónico]
\label{prop:centro-electron-rev6}
La involución \(\rho\) posee un único punto fijo:
\[
 \boxed{(r,c)=(5,5).}
\]
Su unicidad es combinatoria y no depende del valor de
\(\mathfrak e\).
\end{proposition}

\begin{proof}
La ecuación \((r,c)=\rho(r,c)\) equivale a
\[
 2r=10,\qquad2c=10.
\]
En \(\{1,\ldots,9\}^2\) tiene la única solución \(r=c=5\). La evaluación
de la masa no aparece en el argumento.
\end{proof}

\begin{corollary}[Universalidad de la sección central]
\label{cor:universalidad-electron-central-rev6}
Todo automorfismo admisible del atlas que entrelace la inversión central
preserva la celda \((5,5)\).
\end{corollary}

\begin{proof}
Si \(f\rho=\rho f\) y \(\rho(5,5)=(5,5)\), entonces
\(\rho f(5,5)=f\rho(5,5)=f(5,5)\). Por tanto \(f(5,5)\) es otro punto fijo
de \(\rho\); la unicidad obliga a \(f(5,5)=(5,5)\).
\end{proof}

Ésta es la razón estructural por la que la clase electrónica es la misma en
todas las realizaciones que preservan el atlas y su inversión: no se repite
por casualidad un decimal local, sino que toda carta admisible debe
transportar el mismo punto fijo y la misma cadena de invariantes. La lectura
física de esta unicidad identifica cada electrón con una realización de esa
clase central; la proposición demuestra la universalidad interna del atlas.

En la familia electrónica fijada por la construcción, este punto es el anclaje
geométrico de la sección. Conviene no confundir cuatro objetos:
\[
 (5,5),\qquad
 e_{\mathrm{core}},\qquad
 v_e=(0,0,0,0,0),\qquad
 \text{la firma cruda de la celda central}.
\]
El primero es una posición; el segundo, su etiqueta interna; el tercero,
el origen elegido en el retículo fino de acción; el cuarto conserva los
datos crudos del atlas y no tiene por qué ser el vector nulo. Su
coordinación explica por qué el electrón ocupa el centro sin borrar la
genealogía de las representaciones que lo describen.

\paragraph{Rutas electrónicas orientadas \(e_{++},e_{--}\), envolvente \(e_{\mathrm{env}}\), origen de acción y firmas de la celda central}
La celda central porta dos historias orientadas de doce eventos. Escribiendo
\(v^2\) para la concatenación de dos copias de una sextupla, la ruta activa
de la orientación \({++}\), su historia opuesta y la envolvente histórica
son, respectivamente,
\[
\begin{aligned}
 e_{\rm act}=e_{++}
   &=(2,2,5,-4,-3,-2)^2,\\
 e_{--}
   &=(4,3,2,-2,-2,-5)^2,\\
 e_{\rm env}
   &=(4,3,5,-4,-3,-5)^2.
\end{aligned}
\tag{E-centro-1}\label{eq:electron-rutas-finas-rev8}
\]
Las dos historias publican la misma palabra
\[
 \tau_{12}=\texttt{+++---+++---},
\]
pero no son la misma ruta. La envolvente toma el máximo absoluto
coordenada a coordenada conservando el signo; registra amplitudes, pero
aplana el orden genealógico y no sustituye a \(e_{\rm act}\).

Los otros lectores del centro viven en codominios distintos:
\[
\begin{array}{rcll}
 v_e&=&(0,0,0,0,0),
   &\text{origen afín del retículo de acción},\\
 R(5,5)&=&(24,0,12,0,-12),
   &\text{firma cruda APP de la celda},\\
 \operatorname{Sig}^{\Gamma_{108}}_{\rm masa}
   &=&(-12,-4,12,-6,12),
   &\text{proyección par de masa},\\
 \operatorname{Sig}^{\Gamma_{108}}_{\rm carga}
   &=&(0,0,0,0),
   &\text{proyección impar de carga}.
\end{array}
\tag{E-centro-2}\label{eq:electron-firmas-tipadas-rev8}
\]
Que la palabra de carga tenga firma impar nula no convierte la ruta en
cero, ni identifica la firma cruda con el origen afín. Posición, ruta
activa, ruta opuesta, envolvente, origen y ambas firmas son objetos
coordinados, no seis notaciones para un mismo vector.

\subsection{Conjugación partícula–antipartícula y persistencia de la banda electrónica}

La identidad electrónica no termina en el punto central ni en el valor de
masa. La conjugación material se realiza sobre palabras dodecafásicas. Sea
\[
 C_M(\mathbf t)=-\operatorname{rev}(\mathbf t)
\]
la anti-sección orientada y sea
\(\chi\in\{-1,+1\}\) la orientación de carga. La banda es la clase
\[
 [(\mathbf t,\chi)]_M
 =
 \{(\mathbf t,\chi),(C_M\mathbf t,-\chi)\}
\]
de la \cref{def:banda-dodecafasica-ruta}. La banda conserva sus componentes
pares y presenta dos secciones con componentes impares opuestas.

\begin{theorem}[Electrón y positrón como secciones de una banda]
\label{thm:electron-positron-biografia-rev6}
Para el carácter de banda del
\cref{thm:principio-conjugacion-masa-banda},
\[
 m(C_M\mathbf t,-\chi)=m(\mathbf t,\chi).
\]
En el canal electrónico,
\[
 (n,\chi)=(-22,+1)
 \quad\longleftrightarrow\quad
 (+22,-1)
\]
y ambos pares producen
\[
 \chi n=-22.
\]
Por tanto electrón y positrón tienen la misma masa y orientaciones de
carga opuestas.
\end{theorem}

\begin{proof}
Si \(E_+\) y \(E_-\) son, respectivamente, las componentes par e impar del
carácter respecto de \(C_M\), entonces
\[
 E_+(C_M\mathbf t)=E_+(\mathbf t),
 \qquad
 E_-(C_M\mathbf t)=-E_-(\mathbf t).
\]
De aquí,
\[
 E_{\mathrm{band}}(C_M\mathbf t,-\chi)
 =
 E_{\mathrm{band}}(\mathbf t,\chi),
\]
y la igualdad de masa se obtiene al aplicar la exponencial. En el término
electrónico,
\[
 (+1)(-22)=(-1)(+22)=-22.
\]
\end{proof}

El cierre finito de la \cref{prop:cierre-antiseccion-ct108} refuerza esta
interpretación: sobre las ochenta y una rutas CT--108, ni el signo puro ni
la reversión pura preservan la taxonomía, mientras
\(-\operatorname{rev}\) sí la preserva, con
\[
 81=54+9+9+9.
\]
La antipartícula no es, por tanto, una masa negativa ni una copia
independiente. Es la otra sección orientada de la misma banda: conserva los
invariantes pares que determinan la masa e invierte los impares que
determinan carga y orientación.

\subsection{Teorema de cierre electrónico}

\begin{theorem}[Composición holográfica de la sección electrónica]
\label{thm:cierre-electron-holografico-rev6}
Fijados los objetos construidos en las secciones anteriores, la sección
electrónica adimensional es
\[
 \mathfrak e
 =
 \left\|[P_{\Sigma,3},P_{\Pi,3}]\right\|
 \left[
 \exp\bigl((\mathscr R\circ\mathsf J)(\pi,\varphi)\bigr)
 -\frac{\Sigma(\mathrm{APP})}
 {\Sigma(\mathrm{centro})}\,
 \alpha^3
 \right]
 \left[
 1+
 \frac{\Sigma(\mathrm{canal})}
 {\Sigma(\mathrm{centro})}\,
 \Delta_4
 \right].
\]
Equivalentemente,
\[
 \boxed{
 \mathfrak e
 =
 \frac{\sqrt3}{4}
 \left(e^{\varphi/\pi^2}
 -22\alpha^3\right)
 (1+15\Delta_4).}
\]
Su evaluación es
\[
 \boxed{
 \mathfrak e
 =
 0.51099892248806607250987087719134943763\ldots.}
\]
\end{theorem}

\begin{proof}
La \cref{prop:prefactor-app-electron-rev6} aporta
\[
 \left\|[P_{\Sigma,3},P_{\Pi,3}]\right\|=\sqrt3/4.
\]
El \cref{thm:espejo-electron-rev6} aporta
\[
 (\mathscr R\circ\mathsf J)(\pi,\varphi)=\varphi/\pi^2.
\]
El \cref{thm:coeficientes-electron-rev6} aporta
\[
 \frac{\Sigma(\mathrm{APP})}{\Sigma(\mathrm{centro})}=22,
 \qquad
 \frac{\Sigma(\mathrm{canal})}{\Sigma(\mathrm{centro})}=15.
\]
El cierre dodecafásico aporta la ventana racional \(\alpha_{12}\) de la
sección coinductiva \(\alpha\), y
\cref{eq:delta4-electron-rev6} aporta el defecto global. La sustitución
produce la fórmula abreviada. La evaluación sucesiva de
\[
\begin{aligned}
 e^{\varphi/\pi^2}
 &=1.17814494259630678081160095680888910\ldots,\\
 22\alpha^3
 &=0.00000854906599200551727014979807598228\ldots,\\
 1+15\Delta_4
 &=1.00166794634038210081076702179181019\ldots
\end{aligned}
\]
da el valor declarado.
\end{proof}

El teorema es holográfico en un sentido verificable: cada factor local
retiene una capa distinta del sistema completo. El conmutador retiene la
dualidad de hojas; la exponencial retiene la orientación interior de la
bisagra \(\pi\)--\(\varphi\); \(22\) y \(15\) retienen el censo global y
su núcleo central; \(\alpha^3\) retiene el cierre fino; y
\(\Delta_4\) retiene el refinamiento global. El centro determina dónde se
ancla la sección y la banda determina cómo sobrevive bajo conjugación.
El valor no reemplaza esta estructura: es su evaluación común.

\subsubsection{Desdoblamiento de acción anterior y posterior al retorno}
\label{sec:electron-desdoblamiento-two-way-rev9}

La realización regional de la coordenada angular conjugada \(C^*\) distingue dos coordenadas sobre la
misma recta interna de acción. La coordenada anterior al retorno es \(H_5\) y
la posterior es
\[
\etaRet
 =H_5-\frac{90}{\pi}\mathscr C_\pi(\alpha).
\]
La notación \(\etaRet\) es la coordenada sexagesimal canónica definida en la
rama de acción por \(\etaRet=\Cstar/2-\alphaAngle\). La igualdad anterior se
evalúa en esa carta; no introduce una segunda coordenada de retorno.
Por tanto, las cartas aditiva y multiplicativa del desdoblamiento orientado
son
\begin{equation}
 \boxed{
 \delta_{2w}:=H_5-\etaRet
 =\frac{90}{\pi}\mathscr C_\pi(\alpha),
 \qquad
 R_{\rm act}:=\frac{\hbar^{\rm pre}_{\rm HMT}}
 {\hbar^{\rm ret}_{\rm HMT}}
 =\frac{H_5}{\etaRet}.}
 \label{eq:electron-desdoblamiento-two-way-rev9}
\end{equation}
Ambas coordenadas proceden de un único transporte pre-retorno/post-retorno.
No son dos constantes de Planck: el factor común
\(10^{-34}\mathcal U_S\) fija la recta dimensional y cancela exactamente en
\(R_{\rm act}\). La elipse de acción realiza después
\(\operatorname{Area}(\theta)=\hbar_{\rm HMT}^{\rm ret}\theta\), de modo que
el desdoblamiento antecede a la lectura de holonomía.

El acoplamiento
\begin{equation}
 \kappa_e^{\beta}
 =80-54\alpha+6\Delta_4,
 \qquad
 \mathfrak e_e^{\beta}
 =\mathfrak eR_{\rm act}^{\kappa_e^{\beta}}
 =0.5109989506900008086\ldots
 \label{eq:electron-beta-two-way-rev9}
\end{equation}
es la clausura cuantitativa bidireccional del electrón. Los dos lectores
independientes del registro electrónico construidos en el capítulo siguiente
seleccionan, sin consultar la masa de contraste,
\[
 K_D(\Gamma_e)=K_\Omega(\Gamma_e)=(80,54,6),
\]
y derivan exactamente \(\kappa_e^\beta\). El entero \(80\) procede de la
primera coordenada de esos lectores; la palabra \(w_e=201101\) conserva su
residencia propia en el canal de la constante de Euler.

\begin{statusbox}[title={Clausura interna del refinamiento bidireccional}]
La identidad \eqref{eq:electron-desdoblamiento-two-way-rev9}, la reducción
\((80,54,6)\) y la ecuación
\eqref{eq:electron-beta-two-way-rev9} constituyen resultados internos exactos
y target--free sobre el registro electrónico. La propagación a fibras de
rango mayor se formula mediante los operadores de multisección.
\end{statusbox}

\subsection{Carta energética y contraste de la clausura bidireccional}

Sea \(\mathcal U_E\) una recta real de dimensiones de energía y sea
\(u_E\) una base elegida. La transducción metrológica es
\[
 \varsigma_{u_E}:\mathbb R\longrightarrow\mathcal U_E,
 \qquad
 \varsigma_{u_E}(x)=x\,u_E.
\]
La carta conserva dos etapas de una misma genealogía. La etapa basal es
\[
 E_e^{(0)}
 =
 \varsigma_{u_E}(\mathfrak e).
\]
Con \(u_E=1\,\mathrm{MeV}\),
\[
 E_e^{(0)}
 =
 0.5109989224880660725\ldots\ \mathrm{MeV}.
\]
La razón de acción y el exponente derivado completan la etapa rectora,
\[
 E_e^\beta=E_e^{(0)}R_{\rm act}^{\kappa_e^\beta}
 =0.5109989506900008086082571648\ldots\ \mathrm{MeV},
 \qquad
 m_e^\beta=E_e^\beta/c^2.
\]
La unidad no aparece en la
\cref{eq:electron-holografico-rev6}; es la base de una recta dimensional
aplicada a una salida ya construida.

La referencia CODATA 2022 para la energía equivalente electrónica es
\cite{CODATA2022}
\[
 E_e^{\mathrm{CODATA}}
 =
 0.51099895069(16)\ \mathrm{MeV}.
\]
En esa misma carta, la salida bidireccional satisface
\[
 E_e^\beta-E_e^{\mathrm{CODATA}}
 =
 8.086082571648\times10^{-16}\ \mathrm{MeV},
\]
con diferencia relativa
\[
 1.582406883758\times10^{-15}.
\]
La comparación se realiza después de congelar la evaluación interna. Las
identidades del selector basal permanecen
\[
 22=\frac{396}{18},
 \qquad
 15=\frac{270}{18},
 \qquad
 p=27-22=5.
\]
y reciben además la demostración combinatoria
\(-22=-|K_e\setminus\iota(R_\pi)|\),
\(+15=|R_\pi\times\mathbb F_3|\). El contraste metrológico verifica la
salida; la ruta, los lectores y la razón de acción la generan.

\subsection{Consecuencias estructurales del cierre electrónico}

La lectura estructural anterior impone varias distinciones:

\begin{enumerate}
\item \(\sqrt3/4\) es la norma del conmutador APP en dimensión tres. Su
coincidencia con otras fórmulas geométricas no sustituye esta procedencia.

\item \(\pi/\varphi^2\) y \(\varphi/\pi^2\) no son intercambiables. La
primera razón es la orientación de cierre; la segunda, obtenida mediante
\(\mathsf J\), es la orientación electrónica de autocrecimiento.

\item \(22\) y \(15\) proceden de \(396/18\) y \(270/18\).
Presentarlos como valores seleccionados lógicamente por \(Z_0\) invertiría
la dependencia demostrada.

\item \(\alpha^3\) usa el cierre previo de
\(\alpha\). La potencia cúbica pertenece a la ley
electrónica; no autoriza por sí sola una regla universal
\(\alpha^d\leftrightarrow d\) sin un teorema adicional.

\item \(\Delta_4\) es global. No debe reemplazarse por
\(s_{270}\), por \(135\Delta_4^2\) ni por un residuo particular de la
especie.

\item La posición \((5,5)\), la etiqueta \(e_{\mathrm{core}}\), el origen
fino \(v_e=0\) y la firma cruda central son objetos coordinados, no
sinónimos.

\item Electrón y positrón no se distinguen por dos leyes de masa: son las
dos secciones orientadas de una misma banda, con masa par y carga impar.

\item \(0.5109989224\ldots\) es primero un escalar adimensional. MeV,
\(\mathrm{MeV}/c^2\), CODATA y \(Z_0\) pertenecen a la comparación
posterior.
\end{enumerate}

Estas distinciones explican por qué la ecuación electrónica es una
miniatura del sistema HMT--MD. En una sola sección reaparecen soporte,
hojas, orientación, censo, cierre fino, refinamiento, centro, persistencia
y conjugación. Éste es el contenido de la palabra \emph{holográfica}: no
que una cifra represente poéticamente al todo, sino que cada factor
conserve una dependencia verificable del todo que la produce.
 
\section{Derivación bidireccional del exponente electrónico desde la ruta enriquecida}
\label{sec:cierre-electron-two-way-rev11}
\index{electrón!lector bidireccional}
\index{calendario de orden 108!exponente electrónico}

La ruta enriquecida del punto central publica dos representaciones
complementarias de su memoria: la palabra trítica completa
\(\gamma_{108}\in\mathbb F_3^{108}\) y el registro dodecafásico firmado
\(\tau_{12}\in\{-1,0,+1\}^{12}\). Los dos lectores definidos a continuación
se calculan exclusivamente desde esas representaciones, el calendario y las
constantes HMT ya construidas.

\subsection{Lector de desplazamientos sobre el calendario cíclico de orden 108}

Sea \(S\) el desplazamiento cíclico en \(\mathbb Z/108\mathbb Z\). Para
\(k\in\mathbb Z/108\mathbb Z\), definimos
\begin{equation}
 D_k(\Gamma):=
 \#\{t\in\mathbb Z/108\mathbb Z:\gamma_t\ne\gamma_{t+k}\}
 =\lVert(I-S^k)\gamma\rVert_0.
 \label{eq:lector-dk-ct108-rev11}
\end{equation}
El triple entero del lector es
\begin{equation}
 K_D(\Gamma)=
 \left(D_{27}+D_{36},\ D_1,\ \frac{D_4-D_3}{2}\right)
 =:(A_D,B_D,C_D),
 \label{eq:KD-rev11}
\end{equation}
y su evaluación escalar es
\begin{equation}
 \kappa_D(\Gamma)=A_D-{\alpha}B_D+\Delta_4 C_D.
 \label{eq:kappaD-rev11}
\end{equation}
Los desplazamientos (27) y (36) son las elevaciones cronológicas de
\(90^\circ\) y \(120^\circ\) porque nueve pasos discretos representan
\(30^\circ\). El término \(D_1\) mide la frontera local inmediata y la
diferencia \(D_4-D_3\) compara las dos genealogías antes de su cociente de
altura. En el dominio de las 324 rutas se cumple
\(\gamma_{t+54}=\gamma_t\); todos los \(D_k\) son pares y \(C_D\) es entero.

\subsection{Lector dodecafásico de diferencias de correlación}

Para la palabra \(\tau=(\tau_r)_{r\in\mathbb Z/12\mathbb Z}\), sean
\[
 C_k^{\rm corr}(\tau)=\sum_{r=0}^{11}\tau_r\tau_{r+k},
 \qquad
 A_\ell(\tau)=\sum_{r=0}^{11}
 \left(\sum_{j=0}^{\ell-1}\tau_{r+j}\right)^2.
\]
La obstrucción genealógica primitiva y la memoria antipodal son
\begin{align}
 \Omega_{90\mid120}(\tau)
 &=4A_3-3A_4
 =-2C_1^{\rm corr}-4C_2^{\rm corr}-6C_3^{\rm corr},
 \label{eq:omega-two-way-rev11}\\
 P_6(\tau)&=\sum_{r=0}^{5}\tau_r\tau_{r+6}.
 \label{eq:p6-two-way-rev11}
\end{align}
El símbolo \(P_6\) evita la colisión con el ciclo corto del corredor de
constantes. El segundo triple y su evaluación son
\begin{equation}
 K_\Omega(\Gamma)=\bigl(\Omega_{90\mid120}(\tau_\Gamma),54,
 P_6(\tau_\Gamma)\bigr),
 \qquad
 \kappa_\Omega=\Omega_{90\mid120}-54\alpha+P_6\Delta_4.
 \label{eq:kappa-omega-rev11}
\end{equation}

\HMTClaim{MD-R-ELECTRON-OMEGA-PRIMITIVE}
\begin{lemma}[Unicidad lineal integral primitiva]
\label{lem:unicidad-omega-rev11}
Sea \(L=uA_3+vA_4\), con \(u,v\in\mathbb Z\), un contraste que se anula
cuando las genealogías descienden a una misma altura,
\(A_3=3a\), \(A_4=4a\). Entonces
\((u,v)=n(4,-3)\). Fijado el signo \(90-120\),
\(\Omega_{90\mid120}=4A_3-3A_4\) es el generador primitivo.
\end{lemma}

\begin{proof}
La anulación para todo \(a\) exige \(3u+4v=0\). Como
\(\gcd(3,4)=1\), las soluciones enteras son
\((u,v)=n(4,-3)\). La primitividad corresponde a \(|n|=1\).
\end{proof}

El lema selecciona \(\Omega\) dentro de la clase de contrastes lineales
integrales de las energías \(A_3,A_4\). Su dominio de unicidad es esa clase
funcional; los invariantes no lineales pertenecen a un problema de selección
distinto.

\subsection{Evaluación de la ruta central}

La ruta electrónica \(\Gamma_e=(5,5,++)\) proporciona, con anterioridad al
contraste metrológico,
\[
 (D_1,D_3,D_4,D_{27},D_{36})=(54,56,68,48,32),
\]
y su registro dodecafásico firmado es
\[
 \tau_e=(+,+,+,-,-,-,+,+,+,-,-,-).
\]
Por tanto,
\[
 K_D(\Gamma_e)=(48+32,54,(68-56)/2)=(80,54,6),
\]
\[
 K_\Omega(\Gamma_e)=(80,54,6).
\]

\HMTClaim{MD-R-ELECTRON-TWOWAY-REDUCTION}
\begin{theorem}[Reducción electrónica común de los lectores bidireccionales]
\label{thm:reduccion-electron-two-way-rev11}
Los lectores \(K_D\) y \(K_\Omega\) coinciden sobre la ruta central y
determinan
\begin{equation}
 \kappa_e=80-54\alpha+6\Delta_4.
 \label{eq:kappa-electron-demostrado-rev11}
\end{equation}
\end{theorem}

\begin{proof}
Las dos evaluaciones anteriores dan el mismo triple entero. La sustitución
de sus componentes en
\(A-{\alpha}B+\Delta_4C\) prueba
\eqref{eq:kappa-electron-demostrado-rev11}. Toda cantidad empleada en esta
reducción pertenece a la ruta, a sus dos registros o a las constantes HMT
previamente construidas.
\end{proof}

\paragraph{Realización posterior en la carta de energía.}
Con \(R_{\rm act}=H_5/\etaRet\), la acción positiva sobre el basal
\(\mathfrak e_0\) produce
\begin{equation}
 \mathfrak e_e^\beta
 =\mathfrak e_0R_{\rm act}^{\kappa_e}.
 \label{eq:masa-electron-adimensional-two-way-rev11}
\end{equation}
Después de aplicar la carta vigente a MeV se obtiene
\begin{equation}
 \varsigma_{u_E}(\mathfrak e_e^\beta)\big|_{u_E=1\,\mathrm{MeV}}
 =0.5109989506900008086082571648\ldots\ \mathrm{MeV}.
 \label{eq:masa-electron-two-way-rev11}
\end{equation}
La primera igualdad pertenece al teorema interno condicionado por el basal y
la recta de acción; la segunda es su realización dimensional. El valor
externo se reserva para el contraste metrológico.

El entero (80) procede de \(D_{27}+D_{36}\) y, en el lector dodecafásico,
de \(\Omega_{90\mid120}\). La coincidencia con otros cardinales internos es
un control de compatibilidad y no sustituye esta procedencia operatoria. La
ruta de orientación \((--)\) es una traslación temporal de la ruta \((++)\)
y conserva el mismo perfil, conforme a la paridad de masa bajo conjugación.

\subsection{Separación de lectores y escalas de retorno}

La coincidencia numérica de componentes internos queda subordinada a su
tipo. En el punto central intervienen, entre otros, los objetos
\[
 R_{\rm atlas}(5,5)=(24,0,12,0,-12),
 \qquad
 L_{\kappa_0}(\Gamma_e^{++})=(24,0,0,0,-12),
\]
\[
 \begin{gathered}
 k_{\Delta_4}^{\rm dod},k_{\Delta_4}^{\rm pro}:
 \mathscr R_{324}\longrightarrow\mathbb Z,\\
 k_{\Delta_4}^{\rm dod}(\Gamma_e^{++})=12,
 \qquad
 k_{\Delta_4}^{\rm pro}(\Gamma_e^{++})=4,\\
 K_D(\Gamma_e^{++})=K_\Omega(\Gamma_e^{++})=(80,54,6).
 \end{gathered}
\]
Aquí \(\mathscr R_{324}\) denota el dominio de las 324 rutas orientadas.
El primero es una firma estática de familia; el segundo, un extractor de
genealogía; el tercer y el cuarto término son lectores torsionales definidos,
respectivamente, por la lectura dodecafásica y por la prolongación prospectiva;
el último objeto es el par de lectores bidireccionales. Ninguna igualdad
parcial identifica esos morfismos ni permite trasladar un coeficiente entre
lectores.

También se distinguen cuatro escalas temporales: \(27\) es el primer retorno
de clase, órbita y celda; \(54=D_1(\Gamma_e)\) es el retorno cinemático
observable dentro del ciclo; \(108\) cierra el calendario; y \(216\) cierra el estado cinemático
orientado en la extensión. Esta tipificación conserva separados retorno
posicional, observación y holonomía orientada.

\begingroup\fontsize{10.2}{11.7}\selectfont
\setlength{\parskip}{2pt}\setlength{\abovedisplayskip}{5pt}\setlength{\belowdisplayskip}{5pt}
\paragraph{Cadena causal del resultado electrónico}
La ruta y sus dos registros determinan \(K_D\), \(K_\Omega\), su reducción
común y la evaluación de \(\kappa_e\). La realización en MeV aplica después
la carta de energía; la masa externa interviene sólo en el contraste. El par
\(\mathbf K=(K_D,K_\Omega)\) es la salida operatoria completa de HMT sobre
cada ruta. En una fibra no central, su representación por un único escalar
físico exige un morfismo adicional entre el operador de fibra y la carta
dimensional; esta exigencia no modifica la ley interna ni autoriza a elegir
una celda por proximidad metrológica.
 
\paragraph{Paso de la realización Cartan–Holst a la confluencia horizonte–área–información mediante conservación de la acción} La clausura electrónica fija una realización material; el paso siguiente identifica la gramática que convierte acción, reloj y masa en magnitudes físicas sin utilizar esa realización como dato de entrada.

\par\endgroup
 \endgroup

\chapter{Derivación determinantal de la escala absoluta de acción y valoración \texorpdfstring{\(10^{-34}\)}{10^-34}}
\begingroup
\renewcommand{\chapter}[2][]{}%
\chapter{Derivación determinantal de la escala de acción: forma normal de Smith, conúcleo de orden \(16\) y valoración \(10^{-34}\)}
\label{chap:p3-final:34:escala_accion}
\label{ch:vacuum}

\section{Derivaci\'on de la escala de acci\'on a partir del estado dodecaf\'asico}

La d\'ecada de acci\'on se determina intr\'insecamente, sin adoptarla como
unidad heredada. El funcional determinantal local se aplica a la salida HMT
previamente construida \(\alpha\) y define
\begin{equation}
\Sigma_H=\varphi\alpha^{16}.
\label{eq:action-determinantal-reader}
\end{equation}
El exponente \(16\) es el orden del cokernel local del cierre firmado; su valoraci\'on decimal fija
\begin{equation}
\operatorname{ord}_{10}(\Sigma_H)=-34.
\label{eq:action-decade}
\end{equation}

El origen del exponente se puede verificar dentro de este volumen. Sobre la
\'orbita local de cuatro hojas act\'ua el bloque de Hadamard
\begin{equation}
H_4=
\begin{pmatrix}
1&1&1&1\\
1&1&-1&-1\\
1&-1&1&-1\\
1&-1&-1&1
\end{pmatrix}.
\label{eq:action-hadamard}
\end{equation}

\begin{lemma}[\'Indice local de la recta de acci\'on]
La forma normal de Smith de \(H_4\) es
\begin{equation}
\operatorname{SNF}(H_4)=\operatorname{diag}(1,2,2,4).
\label{eq:action-smith}
\end{equation}
En consecuencia,
\[
\operatorname{coker}(H_4)\simeq
\Z/2\Z\oplus\Z/2\Z\oplus\Z/4\Z,
\qquad
|\operatorname{coker}(H_4)|=16.
\]
\end{lemma}

\begin{proof}
El m\'aximo com\'un divisor de las entradas es uno; el de los menores
\(2\times2\) no nulos es dos; el de los menores \(3\times3\) es cuatro; y
\(|\det H_4|=16\). Los cocientes sucesivos de estos divisores determinantes
son \(1,2,2,4\), que prueban \cref{eq:action-smith}. El producto de los
invariantes da el orden del con\'ucleo.
\end{proof}

La norma de la secci\'on \(\alpha\) sobre esas diecis\'eis hojas
produce \(\alpha^{16}\); la autoescala \(\varphi\) act\'ua
una sola vez sobre la base, no una vez por hoja. De ah\'i
\cref{eq:action-determinantal-reader}. Las comparaciones certificadas
\begin{equation}
\alpha^{16}<10^{-34}
<\varphi\alpha^{16}<10^{-33}
\label{eq:action-decade-bounds}
\end{equation}
demuestran, sin introducir la unidad SI, que el orden decimal es
exactamente \(-34\). La evaluaci\'on terminal comienza por
\begin{equation}
\Sigma_H
=1.046243838104756580184229208303089\ldots\times10^{-34}.
\label{eq:action-sigma-value}
\end{equation}
La sustituci\'on del con\'ucleo por tres copias con exponente \(16^3\), o la
aplicaci\'on reiterada de \(\varphi\) en cada hoja, modificar\'ia la
genealog\'ia. Ambas operaciones constituyen criterios de refutaci\'on y no
convenciones equivalentes.
La cadena causal que se conserva es, por tanto,
\[
\APP\to\TRIT\to\TPK\to U_{12}^{\rm sign}
\to\alpha\to\Sigma_H
\to10^{-34}\U_S\to
\{\hbar_{\rm pre},\hbar_{\rm ret}\}.
\]
Este volumen no rederiva \(\alpha\), sino que desarrolla la geometr\'ia
metrol\'ogica posterior a la selecci\'on del tipo y de la d\'ecada por
\(\Sigma_H\). El entero \(-34\) no es una mantisa ni un ajuste al Sistema
Internacional, sino la valoraci\'on terminal del funcional local.

\section{Secciones orientadas de la recta de acci\'on}

Las realizaciones anterior y posterior al retorno constituyen dos secciones
de una misma recta de acci\'on:

\begin{equation}
\hbar_{\rm pre}=H_5\,10^{-34}\U_S,\qquad
\hbar_{\rm ret}=\eta_{\rm ret}\,10^{-34}\U_S.
\label{eq:action-twoway}
\end{equation}

La d\'ecada com\'un determina la recta, mientras que la informaci\'on
orientada se expresa mediante

\begin{equation}
R_{\rm act}=\frac{\hbar_{\rm pre}}{\hbar_{\rm ret}}
=\frac{H_5}{\eta_{\rm ret}},
\qquad
\xi_{\rm act}=\frac12\log R_{\rm act}.
\label{eq:action-rapidity}
\end{equation}

La primera cantidad proporciona una coordenada multiplicativa y la segunda
una coordenada aditiva de tipo rapidez. No representan dos constantes de
Planck independientes, sino las dos orientaciones de una misma secci\'on
respecto del retorno.

Si \(h_a=2\pi\hbar_a\), \(a\in\{\mathrm{pre},\mathrm{ret}\}\), una magnitud proporcional a \(h_a^{1/2}\) cambia como \(R_{\rm act}^{1/2}\), una proporcional a \(h_a^{-1/2}\) cambia como \(R_{\rm act}^{-1/2}\), y un cociente donde la acci\'on se cancela es bidireccional invariante.

  \endgroup
% Distribución causal exacta del delta Ley 9; contenido conservado sin síntesis.
\Needspace{28\baselineskip}
\section{Bifurcación causal posterior a \texorpdfstring{\(\alpha\)}{alpha HMT}}
\label{sec:l9s102:bifurcacion-accion-angular}

La genealogía de la acción y del par angular forma un diagrama bifurcado.
La rama regional produce la mantisa reducida \(\eta_{\mathrm{ret}}\); la rama
determinantal produce el orden decimal \(\varepsilon_H=-34\). Ambas
confluyen en la recta de acción. El par angular se construye en una rama
hermana a partir de \(\alpha\) y
\(\eta_{\mathrm{ret}}\):
\begin{equation}
\begin{tikzcd}[column sep=small,row sep=normal]
\alpha,\pi,\varphi
  \arrow[r] \arrow[d]
& H_5,\mathscr C_\pi
  \arrow[d]
&& \alpha,\varphi,H_4
  \arrow[d] & {}\\
x_A,D_A,\rho_\pi
  \arrow[r]
& \eta_{\mathrm{ret}}
  \arrow[r] \arrow[d]
& \hbar_{\mathrm{ret}}\in\mathcal L_S
& \Sigma_H,\varepsilon_H=-34
  \arrow[l] & {}\\
& (A_{\deg},C^*_{\deg})
  \arrow[r]
& (A_{\mathrm{rad}},C^*_{\mathrm{rad}})
  \arrow[r]
& (q_+,q_-)
  \arrow[r]
& V_{90,120}.
\end{tikzcd}
\label{diag:l9s102:dag-accion-angular}
\end{equation}
El orden expositivo situa la derivación de Planck antes de la publicación
angular completa; el diagrama conserva simultaneamente la dependencia
matemática exacta de cada salida.

\section{Conúcleo local y derivación del orden absoluto \texorpdfstring{\(10^{-34}\)}{10^-34}}
\label{sec:l9s102:orden-accion}

La forma normal de Smith del bloque dodecafásico local es
\begin{equation}
 \operatorname{SNF}(H_4)=\operatorname{diag}(1,2,2,4),
 \qquad
 \operatorname{coker}(H_4)
 \cong\mathbb Z/2\mathbb Z\oplus\mathbb Z/2\mathbb Z
 \oplus\mathbb Z/4\mathbb Z.
 \label{eq:l9s102:snf-h4}
\end{equation}
El conúcleo local posee orden \(16\). Para la sección constante
\(s_\alpha(g)=\alpha\), la norma local y una actuación de
\(\varphi\) producen
\begin{equation}
 \boxed{
 \Sigma_H=\varphi\mathcal N_{G_H}(s_\alpha)
 =\varphi\alpha^{16}.}
 \label{eq:l9s102:sigma-h}
\end{equation}
Las desigualdades exactas
\begin{equation}
 \alpha^{16}<10^{-34}
 <\varphi\alpha^{16}<10^{-33}
 \label{eq:l9s102:orden-34}
\end{equation}
determinan
\begin{equation}
 \boxed{
 \varepsilon_H=\left\lfloor\log_{10}\Sigma_H\right\rfloor=-34.}
 \label{eq:l9s102:epsilon-h}
\end{equation}
 
\chapter{Carácter regional de grado cinco y medida positiva de acción}
\begingroup
\renewcommand{\chapter}[2][]{}%
\chapter{Medida de cociente y carácter positivo de la acción: hilbertización canónica, defecto regional--dodecafásico y ley functorial sobre caminos}
\label{chap:p3-final:35:medida_accion_positiva}
\label{p3-20260826:chap:medida-accion-positiva}
\index{medida!catálogo de emisiones}
\index{acción!cociclo positivo}
\index{traza normalizada}

La sección anterior identifica isométricamente un sector positivo de rango tres en
la interfaz orientada \(t=7\), la dodecafase y el módulo de incidencia de Witt. La construcción
siguiente deriva el otro coeficiente de la ley angular, \(5/468\), y
demuestra por qué ambos términos se combinan mediante una suma positiva. La
cuestión exige separar
tres operaciones:
\begin{enumerate}
  \item contar presentaciones microscópicas;
  \item contar emisiones distintas después de formar el cociente TPK;
  \item asignar una acción aditiva a una ruta.
\end{enumerate}
Las dos primeras operaciones producen medidas diferentes. La tercera se
caracteriza por una propiedad universal y no por la proximidad posterior de
un decimal.

Denotaremos por \(\Gtr\) el grafo dirigido de \(1944\) estados del lector de
caminos enriquecidos. Su categoría libre de caminos se escribirá
\(\mathsf P(\Gtr)\). La definición local de sus aristas y de su acción se
desarrolla en la sección correspondiente.

\section{Espacio de semillas, cociente de emisiones y fibras regionales}

Sea
\[
 \Omega_{\rm seed}
 =\bigl((\mathbb Z/9\mathbb Z)^2\times D_4\bigr)^2,
 \qquad
 |\Omega_{\rm seed}|=104\,976,
\]
el conjunto ejecutable de parejas de semillas del emisor finito, y sea
\[
 F:\Omega_{\rm seed}\twoheadrightarrow\mathcal U_6,
 \qquad
 |\mathcal U_6|=468,
\]
la aplicación que conserva la emisión séxtuple distinta. El censo exacto de
fibras es
\[
\begin{array}{c|c|c}
|F^{-1}(U)|&\#\{U\text{ con ese tamaño}\}&\text{semillas aportadas}\\
\hline
144&432&62\,208\\
432&18&7\,776\\
1\,944&18&34\,992.
\end{array}
\]
Las tres contribuciones suman \(104\,976\). Por tanto, \(F\) no es una
cubierta regular: emisiones diferentes poseen números diferentes de
presentaciones microscópicas.

La microfibra de \(\pi\) contiene cinco emisiones distintas. Cuatro tienen
multiplicidad \(144\) y una tiene multiplicidad \(432\). Este dato basta
para demostrar que la medida uniforme sobre semillas y la medida uniforme
sobre emisiones no son la misma.

\section{\(2\) medidas exactas sobre el cociente de emisiones y palabras}

Si cada semilla recibe peso \(1/104\,976\), la medida imagen por \(F\) es
\[
 \nu_{\rm seed}(U)
 =\frac{|F^{-1}(U)|}{104\,976}.
\]
Sus átomos posibles son
\[
 \frac1{729},\qquad\frac1{243},\qquad\frac1{54}.
\]
\Needspace{5\baselineskip}
En particular,
\[
 \nu_{\rm seed}(\mathscr R_\pi^{\rm micro})
 =\frac{4\cdot144+432}{104\,976}
 =\boxed{\frac7{729}}.
\]

Si el estado elemental es, en cambio, una emisión distinta
\(U\in\mathcal U_6\), el álgebra diagonal de observables es
\[
 \mathcal D_{468}=\ell^\infty(\mathcal U_6)
\]
y su estado de conteo normalizado es
\[
 \tau_{468}(f)
 =\frac1{468}\sum_{U\in\mathcal U_6}f(U).
\]
Sea \(\Pi_\pi^{(468)}\) la función característica de las cinco regiones de la
microfibra. Entonces
\[
 \boxed{\tau_{468}(\Pi_\pi^{(468)})=\frac5{468}.}
\]
La diferencia exacta entre ambas lecturas es
\[
 \frac5{468}-\frac7{729}
 =\frac{41}{37\,908}.
\]
No son dos aproximaciones. La primera medida cuenta presentaciones; la
segunda cuenta clases de emisión.

\begin{theorem}[Unicidad de la medida uniforme del catálogo]
\label{p3-20260826:thm:unicidad-medida-468}
La única probabilidad sobre \(\mathcal U_6\) invariante bajo toda
relabelización de sus \(468\) átomos es la medida uniforme.
Equivalentemente, la única traza positiva y unital sobre
\(M_{468}(\mathbb C)\) es
\[
 \tau_{468}(A)=\frac1{468}\operatorname{Tr}A.
\]
\end{theorem}

\begin{proof}
Las transposiciones llevan cualquier átomo a cualquier otro. La invariancia
fuerza que todos tengan la misma masa, y la normalización la fija en
\(1/468\). En la versión matricial, la invariancia unitaria iguala el valor
de todos los proyectores de rango uno. La descomposición espectral y la
linealidad reconstruyen la traza normalizada.
\end{proof}

La hipótesis decisiva está a la vista: el lector trabaja con el
\emph{catálogo cociente}. El teorema no afirma que la medida uniforme de
semillas se convierta en uniforme al olvidar las fibras. La simetría usada
es la naturalidad del objeto finito sin etiquetas, no una propiedad
ergódica que se atribuya sin prueba a la dinámica microscópica.

\section{Hilbertización canónica del cociente}

La relación entre semillas y emisiones puede formularse sin ambigüedad.
Sean
\[
 H_{\rm seed}=\ell^2(\Omega_{\rm seed}),
 \qquad
 H_U=\ell^2(\mathcal U_6).
\]
El emisor determina
\[
 \boxed{
 J_F\delta_U
 =|F^{-1}(U)|^{-1/2}
 \sum_{\omega\in F^{-1}(U)}\delta_\omega.
 }
\]
Las fibras distintas son disjuntas y cada columna posee norma uno; por
tanto,
\[
 J_F^*J_F=I_{H_U}.
\]
Además, para todo observable \(f\) del catálogo,
\[
 M_{f\circ F}J_F=J_FM_f.
\]
La imagen \(J_FH_U\) es exactamente el subespacio de funciones constantes
en cada fibra.

Si \(F^\sharp:H_U\to H_{\rm seed}\) denota la aplicación inducida por
precomposición, sin normalizar,
entonces
\[
 J_F
 =F^\sharp\bigl((F^\sharp)^*F^\sharp\bigr)^{-1/2}.
\]
Es el factor isométrico de la descomposición polar de esa aplicación. Esta
expresión demuestra su naturalidad bajo isomorfismos y relabelizaciones del
emisor. No implica una functorialidad falsa bajo cualquier composición de
sobreyeciones no regulares: la normalización depende de las fibras de la
aplicación concreta.

En lenguaje de medidas, la elevación de la traza de cociente que es
uniforme dentro de cada fibra vale
\[
 \widetilde\tau(\omega)
 =\frac1{468\,|F^{-1}(F\omega)|}.
\]
Cada fibra recibe masa total \(1/468\). La compensación inversa por
multiplicidad, y no una invocación informal a órbitas, es la prueba exacta.

\begin{corollary}[Canonicidad de \(5/468\)]
\label{p3-20260826:cor:canonicidad-cinco-468}
Una vez que el lector regional se define sobre emisiones distintas, la masa
de la microfibra de \(\pi\) queda forzada:
\[
 \tau_{468}(\Pi_\pi^{(468)})
 =\frac{\operatorname{rank}\Pi_\pi^{(468)}}{468}
 =\frac5{468}.
\]
\end{corollary}

\section{Defecto positivo entre las medidas regional y dodecafásica}

Escribamos
\[
 H_{\rm ang}=H_U\otimes V_{11}.
\]
El escalar interno
\[
 \Delta_4
 =\frac{\pi-e\log\pi}{270}
 \left(1+\frac\pi{729}\right)
 =0.0001111964226921402\ldots
\]
se construye con las lecturas HMT ya fijadas de \(\pi\) y \(e\), y con los
enteros \(270\) y \(729\) del registro. Aquí es una entrada matemática
pre-dimensional: no intervienen unidades ni observables físicos.

La construcción interna del escalar no determina por sí sola qué operador
debe multiplicar. Para mantener visible esta premisa, declaramos:
\begin{description}
  \item[\HDef --- defecto de cambio de hoja.] La coordenada dodecafásica del defecto
  de cambio de hoja es \(\Delta_4P_3\).
\end{description}
\HAdd fijará la forma aditiva de la acción, pero no selecciona \HDef: para cualquier
\(c\ge0\), el mismo argumento admitiría el par
\((\Pi_\pi^{(468)},cP_3)\). Demostrar \HDef directamente desde una transición
total del núcleo topológico de fase enriquecido es una condición adicional.

Definimos la suma de Kronecker
\[
 \boxed{
 D_{\rm switch}^{\rm dod}
 =\Pi_\pi^{(468)}\otimes I_{11}
 +\Delta_4 I_{468}\otimes P_3.
 }
\]
Los dos sumandos son positivos, conmutan y actúan en factores diferentes.
Por tanto,
\[
 D_{\rm switch}^{\rm dod}\ge0.
\]
El entrelazador \(U_W\) del
\cref{thm:entrelazador-dodecafase-witt} produce la carta Witt:
\[
\begin{aligned}
 D_{\rm switch}^{\rm Witt}
 &=(I_{468}\otimes U_W)
 D_{\rm switch}^{\rm dod}
 (I_{468}\otimes U_W)^*\\
 &=\Pi_\pi^{(468)}\otimes I_{E_{30}}
 +\Delta_4I_{468}\otimes Q_3.
\end{aligned}
\]
Así, el mismo operador positivo se lee en la dodecafase y en el módulo de
incidencia.

La cadena isométrica orientada identifica además los sectores activos
\[
 I_{L_{\rm or}}\otimes P_G,\qquad P_3,\qquad Q_3.
\]
No se afirma una equivalencia de los complementos completos de los tres
módulos: \(W_{GK}\) es una equivalencia parcial de los rangos
seleccionados. Ésa es exactamente la información necesaria para identificar
la segunda contribución de la traza.

\begin{theorem}[Alcance exacto del entrelazamiento y de la inducción de acción]
\label{p3-20260826:thm:alcance-induccion-accion}
Sean \(W_{GK}\) y \(U_W\) las isometrías parciales del
\cref{thm:identificacion-isometrica-tres-proyectores}. Entonces:
\begin{enumerate}
  \item para todo \(c\geq0\),
  \[
   W_{GK}^{*}(cP_3)W_{GK}
   =c(I_{L_{\rm or}}\otimes P_G),
   \qquad
   U_W(cP_3)U_W^{*}=cQ_3;
  \]
  por tanto, el entrelazamiento transporta el coeficiente, pero no lo
  selecciona;
  \item si \(P_{\rm sgn}\) y \(P_{\rm std}\) son los proyectores centrales
  de la descomposición
  \[
   P_3V_{11}\big|_{H_*}
   \simeq \operatorname{sgn}_{H_*}\oplus\operatorname{std},
   \qquad
   \operatorname{rank}P_{\rm sgn}=1,
   \quad
   \operatorname{rank}P_{\rm std}=2,
  \]
  todo operador positivo autoadjunto soportado en \(P_3V_{11}\) y
  \(H_*\)-equivariante tiene la forma
  \[
   D_{a,b}=aP_{\rm sgn}+bP_{\rm std},
   \qquad a,b\geq0.
  \]
  Incluso imponer la traza de \(\Delta_4P_3\) deja la familia
  \[
   a+2b=3\Delta_4.
  \]
  La isotropía \(a=b\) no se deduce de la simetría residual \(S_3\);
  \item la probabilidad uniforme del catálogo,
  \(\mu_U(U)=1/468\), es la traza canónica del cociente finito y su
  elevación compensada es la determinada por \(J_F\). No es la
  medida imagen de la probabilidad uniforme sobre semillas, cuyos átomos son
  \(1/729,1/243,1/54\). En particular,
  \[
   \mu_U(\mathscr R_\pi^{\rm micro})=\frac5{468},
   \qquad
   F_\#\mu_{\rm seed}(\mathscr R_\pi^{\rm micro})=\frac7{729};
  \]
  \item la ley completa
  \[
   \delta_*=\frac5{468}+\frac3{11}\Delta_4,
   \qquad
   \lambda_*=1-\delta_*,
   \qquad
   a_*(e)=g(e)+\lambda_*s(e)
  \]
  es un teorema relativo a las cinco cláusulas \HTr--\HAdd--\HDef--\HRes--\HLoc. El motor
  ejecutable del núcleo topológico de fase reducido no induce por cubierta dirigida el operador ramificado
  \Gtr: la dinámica asociada retorna a la identidad tras \(54\) iteraciones,
  mientras cada estado de
  \Gtr posee cuatro sucesores distintos y ninguna autoarista.
\end{enumerate}
En consecuencia, los resultados anteriores demuestran la canonicidad de las identificaciones
isométricas, de la traza sobre el catálogo cociente y del carácter aditivo
dentro de sus premisas. La afirmación más fuerte de que una transición total
del núcleo topológico de fase enriquecido induce necesariamente \(\Delta_4P_3\), la medida uniforme y la
ley local requiere todavía el teorema dinámico adicional especificado en el
análisis de inducción.
\end{theorem}

\begin{proof}
Las identidades del primer punto se obtienen multiplicando
\(W_{GK}^{*}W_{GK}=I_{L_{\rm or}}\otimes P_G\),
\(W_{GK}W_{GK}^{*}=P_3\) y \(U_WP_3U_W^{*}=Q_3\) por el escalar \(c\).
Para el segundo punto, la restricción a \(H_*\simeq S_3\) es suma sin
multiplicidades de dos representaciones irreducibles no isomorfas. El lema de
Schur diagonaliza todo operador autoadjunto equivariante en esos dos bloques;
la positividad equivale a \(a,b\geq0\). Su traza normalizada es
\((a+2b)/11\), lo que prueba la libertad residual aun después de fijar la
traza. La acción completa de \(A_5\) sobre el triplete irreducible reduciría
la familia a \(cP_3\), pero tampoco fijaría el valor de \(c\).

El tercer punto sigue del censo de fibras: \(F\) tiene multiplicidades
\(144,432,1944\), de modo que la medida imagen pondera una emisión por su
número de presentaciones. La traza \(\tau_{468}\), en cambio, asigna el
mismo peso a cada clase; \(J_F\) realiza exactamente su elevación
normalizado. El cuarto punto es la concatenación lógica ya demostrada:
\HTr fija la traza de cociente, \HAdd la suma de trazas, \HDef declara
\(\Delta_4P_3\), \HRes define la capacidad residual y \HLoc la lleva a las aristas.
La enumeración finita prueba la obstrucción entre el núcleo reducido y
\Gtr: una cubierta dirigida preservaría la estrella saliente y una
semiconjugación del retorno identidad exigiría autoaristas en la imagen,
contradiciendo ambos censos de \Gtr.
\end{proof}

El espectro completo de \(D_{\rm switch}^{\rm dod}\) es
\[
\begin{array}{c|c}
\text{autovalor}&\text{multiplicidad}\\ \hline
0&(468-5)(11-3)=3704\\
\Delta_4&(468-5)3=1389\\
1&5(11-3)=40\\
1+\Delta_4&5\cdot3=15.
\end{array}
\]
Por consiguiente,
\[
\begin{aligned}
 \tau_{468\cdot11}(D_{\rm switch})
 &=\tau_{468}(\Pi_\pi^{(468)})+\Delta_4\tau_{11}(P_3)\\
 &=\boxed{\frac5{468}+\frac3{11}\Delta_4}.
\end{aligned}
\]
La misma identidad vale en la carta Witt. En la carta del sector orientado, el término
\(3/11\) es la traza de \(P_G\) sobre el espacio homogéneo de once ramas.

\section{Propiedad universal del funcional positivo de acción}

La palabra \emph{acción} se usa en sentido aditivo: la concatenación de dos
rutas recibe la suma de sus incrementos. Para separar este contenido de una
fórmula escogida, definimos primero la categoría en la que la ley será única.

\begin{definition}[Carácter positivo de acción HMT]
\label{p3-20260826:def:caracter-accion-hmt}
Un carácter local de acción sobre el producto
regional--dodecafásico es una aplicación
\[
 \delta:
 M_{468}(\mathbb C)_+\oplus M_{11}(\mathbb C)_+
 \longrightarrow\mathbb R_{\ge0}
\]
que satisface:
\begin{enumerate}
  \item aditividad exacta en el cono producto;
  \item homogeneidad positiva;
  \item invariancia por conjugación unitaria independiente en los dos
  factores;
  \item normalización bifactorial,
  \[
  \delta(I_{468},0)=\delta(0,I_{11})=1.
  \]
\end{enumerate}
\end{definition}

La aditividad traduce la concatenación; la invariancia elimina las
coordenadas; la normalización bifactorial fija una escala en cada sumando.
No se trata de un funcional unital sobre el álgebra suma directa: su unidad
es \((I_{468},I_{11})\) y la aditividad da
\(\delta(I_{468},I_{11})=2\). Ninguna cláusula contiene el valor de un
ángulo.

\begin{theorem}[Unicidad del carácter positivo y capacidad residual relativa]
\label{p3-20260826:thm:ley-accion-hmt}
En la clase de la definición anterior existe un único carácter positivo de
acción HMT\@. Es
\[
 \boxed{\delta(A,B)=\tau_{468}(A)+\tau_{11}(B).}
\]
Aplicado al defecto declarado por \HDef,
\((\Pi_\pi^{(468)},\Delta_4P_3)\), produce
\[
 \delta_*=\frac5{468}+\frac3{11}\Delta_4.
\]
Si, además, se declara \HRes, según la cual la capacidad autodual de cambio de
hoja se normaliza en uno y el defecto la reduce mediante
\(\lambda=1-\delta\), la capacidad residual dentro de \HTr--\HAdd--\HDef--\HRes es
\[
 \boxed{
 \lambda_*
 =1-\delta_*
 =1-\frac5{468}-\frac3{11}\Delta_4.
 }
\]
\end{theorem}

\begin{proof}
La aditividad separa los factores:
\[
 \delta(A,B)=\delta(A,0)+\delta(0,B).
\]
Cada restricción es positiva, homogénea e invariante por conjugación. Para
un proyector de rango uno, la invariancia unitaria da un valor común. La
aditividad sobre una descomposición espectral fuerza que cada restricción
sea un múltiplo de la traza. La normalización bifactorial fija esos múltiplos en
\(1/468\) y \(1/11\). De aquí resulta la suma de trazas. La última fórmula
aplica la normalización declarada de capacidad residual.
\end{proof}

Numéricamente,
\[
 \boxed{\lambda_*=0.9892859130191415\ldots},
 \qquad
 0<\lambda_*<1.
\]
La positividad se deduce ya de
\(0<\Delta_4<10^{-3}\) y \(5/468<1\). No depende de una comparación
metrológica.

\section{Functor de acción positiva sobre la categoría de caminos compatibles}

Recordemos el grafo dirigido \(\Gtr\) y su categoría libre
\(\mathsf P(\Gtr)\). Para una arista
\(e:x\to y\), definimos
\[
 g(e)=\mathbf1_{\{d_y\ne d_x\}},
 \qquad
 s(e)=\mathbf1_{\{m_y\ne m_x\}},
\]
donde \(g\) registra un giro y \(s\) un cambio de hoja. La flecha desde el
escalar \(\lambda_*\) hacia las aristas no se deduce del
teorema de trazas. Se declara mediante la cláusula \HLoc: un giro posee coste
unitario, un cambio de hoja posee coste \(\lambda_*\), y ambos indicadores se
suman cuando ocurren en la misma arista. Bajo \HLoc, la acción local es
\[
 \boxed{a_*(e)=g(e)+\lambda_*s(e).}
\]
Para una ruta \(\gamma=e_1\cdots e_n\),
\[
 \mathcal A_*(\gamma)
 =\sum_{j=1}^na_*(e_j)
 =\#\operatorname{giros}(\gamma)
 +\lambda_*\#\operatorname{cambios\ de\ hoja}(\gamma).
\]
Para la ruta vacía \(\mathbf1_x\), se fija
\(\mathcal A_*(\mathbf1_x)=0\). Entonces
\[
 \boxed{
 \mathcal A_*(\gamma_1\gamma_2)
 =\mathcal A_*(\gamma_1)+\mathcal A_*(\gamma_2).
 }
\]
Por tanto,
\[
 \mathcal A_*:\mathsf P(\Gtr)\longrightarrow(\mathbb R_{\ge0},+)
\]
es un funtor monoidal, o cociclo aditivo de caminos.

Los indicadores \(g\) y \(s\) son insensibles al sentido de recorrido. La
orientación no se pierde: reside en la línea \(L_{\rm or}\) y en la
involución de estrella, mientras el coste positivo reside en los
proyectores. Confundir ambos datos introduciría signos negativos en la
acción.

Para \(\beta>0\), el operador de transferencia asociado es
\[
 (\mathcal L_{\beta,*}f)(y)
 =\sum_{e:x\to y}e^{-\beta a_*(e)}f(x).
\]
Todos sus pesos son estrictamente positivos. Como el potencial depende sólo
de los dos indicadores locales, conmuta con la acción gruesa
\((\mathbb Z/3\mathbb Z)^3\) y respeta la descomposición en \(27\)
sectores de Bloch usada en la sección siguiente.

Una dinámica determinista finita no ponderada actúa por matrices de
permutación y tiene valores singulares iguales a uno. El cociclo
\(\mathcal A_*\) convierte el inventario de caminos en un operador de
transferencia que preserva el cono positivo, es decir, en una matriz con
pesos estrictamente positivos sobre las aristas permitidas. Esa propiedad no
demuestra por sí sola contracción, simplicidad del autovalor dominante ni
separación espectral; estas conclusiones requieren una normalización o un
argumento espectral adicional.

\ifdefined\HMTInternalTraceability
\section{Inducción por el cociente y extensión dinámica}

La frase \emph{inducida por el núcleo topológico de fase enriquecido} posee dos sentidos que deben mantenerse
separados.

\begin{enumerate}
  \item \textbf{Inducción por el objeto cociente.} El emisor ejecutable
  determina \(F\), su factor polar \(J_F\), el espacio \(H_U\) y la
  traza normalizada. El estado holonómico integral aporta
  \(\Pi_\pi^{(468)}\), \(P_3\),
  \(K\), \(B_K\), el marco Witt y \(\Delta_4\). En este sentido
  categorial, \(D_{\rm switch}\) se construye sólo con datos internos
  APP--TRIT--TPK\@.
  \item \textbf{Inducción por una transición total.} Esta lectura exigiría
  una transición finita ejecutable sobre todas las coordenadas de
  \(\mathcal X_{\mathrm{TPK}}^{\mathrm{enr}}\) y una prueba de que su
  medida imagen dinámica
  preserva \(J_FH_U\) y produce el potencial \(a_*\). Esa transición total
  no pertenece al dominio ejecutable construido en esta sección.
\end{enumerate}

El \cref{p3-20260826:thm:ley-accion-hmt} y la construcción de caminos cierran la primera
lectura mediante cinco cláusulas explícitas del lector:
\begin{description}
  \item[\HTr --- traza de cociente.] Cuando la multiplicidad de presentación
  permanece en la fibra olvidada, una emisión distinta es un átomo y el
  estado del catálogo es la traza normalizada de su álgebra finita.
  \item[\HAdd --- cociclo de acción.] El defecto de una ruta es el carácter
  positivo, bifactorialmente normalizado, aditivo e invariante del producto
  regional--dodecafásico.
  \item[\HDef --- defecto de cambio de hoja.] El segundo argumento del carácter es
  \(\Delta_4P_3\).
  \item[\HRes --- capacidad residual.] La capacidad autodual se normaliza en
  uno y el defecto la reduce según \(\lambda=1-\delta\).
  \item[\HLoc --- acción local.] Sobre las aristas de \Gtr se fija
  \(a(e)=g(e)+\lambda s(e)\), y la acción de una ruta es la suma sobre sus
  aristas.
\end{description}
\HTr fuerza \(5/468\). \HAdd fuerza la suma de trazas y excluye términos
mixtos. \HDef fija el escalar que ocupa la coordenada de switch. \HRes transforma
el defecto en capacidad residual y \HLoc transporta esa capacidad a las aristas.
Así, \(\delta_*\) es consecuencia de \HTr--\HAdd--\HDef, \(\lambda_*\) lo es de
\HTr--\HAdd--\HDef--\HRes y \(\mathcal A_*\) lo es de
\HTr--\HAdd--\HDef--\HRes--\HLoc. \HTr y \HAdd poseen caracterizaciones universales; \HDef, \HRes y
\HLoc son identificaciones constitutivas internas; no se deducen de la
transición finita anterior.

\section{Familias compatibles con las simetrías parciales}

\subsection{Simetría de \(1/2\) ciclo y familia compatible de medidas}

El intercambio de las dos mitades de una emisión descompone las \(468\)
emisiones en \(234\) pares. Tanto la medida uniforme del catálogo como las
medidas imagen de semillas son invariantes bajo ese intercambio. También lo
es cualquier asignación constante en cada par. La invariancia del retorno,
por sí sola, no prueba uniformidad.

\subsection{Simetría de \texorpdfstring{\Gtr}{Gamma tr} y familia de costes}

Las \(7776\) aristas de \Gtr se dividen en \(288\) órbitas de tamaño
\(27\) bajo \((\mathbb Z/3\mathbb Z)^3\). Las \(2160\) aristas de
cambio de hoja forman \(80\) órbitas:
\[
 80=5\cdot16.
\]
La simetría permite asignar costes diferentes a esas órbitas. El coste
escalar común procede de \HAdd, no de la cardinalidad.

\subsection{Libertad escalar de la \(2^{\mathrm a}\) coordenada}

Para todo \(c\ge0\), el operador
\[
D_c=\Pi_\pi^{(468)}\otimes I_{11}
+cI_{468}\otimes P_3
\]
es positivo y el carácter único de \HAdd da
\[
\tau(D_c)=\frac5{468}+\frac3{11}c.
\]
La fórmula interna de \(\Delta_4\) construye un candidato distinguido, pero
la propiedad universal de la traza no demuestra por sí sola
\(c=\Delta_4\). Esa igualdad es precisamente el contenido de \HDef.

\subsection{Positividad, linealidad y términos de orden \(2\)}

Las funciones
\[
 1-u-v+cuv,\qquad(1-u)(1-v),\qquad e^{-(u+v)}
\]
pueden compartir normalización y primer orden. La ausencia de \(uv\) se
deduce de la aditividad exacta. Sin \HAdd, la ley lineal no sería única.

\section{Hipótesis de existencia, positividad, naturalidad y extensión dinámica}

El resultado depende de las siguientes condiciones verificables:
\begin{enumerate}
  \item el catálogo no contiene \(468\) emisiones o la microfibra no tiene
  rango cinco;
  \item el censo de fibras no es
  \(432\times144,\ 18\times432,\ 18\times1944\);
  \item \(J_F\) no es isométrico o no entrelaza observables del cociente;
  \item alguno de \(P_G,P_3,Q_3\) deja de ser un proyector de rango tres;
  \item existe otro carácter positivo, bifactorialmente normalizado, aditivo e
  invariante que no sea la suma de trazas;
  \item una transición total del núcleo topológico de fase enriquecido, una vez construida, induce un
  coeficiente de switch distinto de \(\Delta_4\);
  \item \(D_{\rm switch}\) adquiere un autovalor negativo;
  \item la acción acumulada deja de ser aditiva bajo concatenación.
\end{enumerate}

Los certificados reconstruyen las fibras, las dos medidas, el factor polar,
el espectro del defecto, las órbitas de aristas y las ablaciones. El valor de
\(\lambda_*\) aparece al final como consecuencia de esos objetos.
\fi
  \endgroup
% Distribución causal exacta del delta Ley 9; contenido conservado sin síntesis.
\section{Carácter regional de grado \texorpdfstring{\(5\)}{5} y mantisa reducida de acción}
\label{sec:l9s102:eta-ret}

El bloque central
\begin{equation}
 \mathcal B_c=\begin{pmatrix}7&2\\2&7\end{pmatrix}
 \label{eq:l9s102:bc-accion}
\end{equation}
posee autovalores \(9\) y \(5\). La longitud \(4\) de las órbitas de paso
\(3\), el rango \(12\) del ciclo dodecafásico y
\(104976/1944=54\) fijan el carácter
\begin{equation}
 \boxed{
 H_5(\alpha,\varphi)
 =\frac12\sqrt{\frac{1000\alpha}{\varphi}}-\alpha
 +\frac9{16}\alpha^2-\frac59\alpha^3
 +\frac7{48}\alpha^4-\frac1{54}\alpha^5.}
 \label{eq:l9s102:h5}
\end{equation}
La completación nonádica construye la coordenada analítica y el determinante
regional
\begin{equation}
 x_A=\frac{50\pi}{9}\alpha,
 \qquad
 D_A=e^{-2x_A}
 =\exp\!\left(-\frac{100\pi}{9}\alpha\right).
 \label{eq:l9s102:xa-da}
\end{equation}
La microfibra aporta \(\rho_\pi=169\), y su serie absolutamente
convergente es
\begin{equation}
 \mathscr C_\pi(\alpha)
 =169\alpha^6+\frac{D_A\alpha^7}{1-D_A\alpha}.
 \label{eq:l9s102:cpi}
\end{equation}
La mantisa reducida posterior al retorno se construye entonces por
\begin{equation}
 \boxed{
 \eta_{\mathrm{ret}}
 =H_5-\frac{90}{\pi}\mathscr C_\pi(\alpha)
 \in\mathbb R_{>0}.}
 \label{eq:l9s102:eta-ret}
\end{equation}
Su miembro derecho contiene exclusivamente objetos generados con anterioridad
por APP--TRIT--TPK y por la clausura de \(\alpha\).

 
\chapter{Derivación autodual de la constante de Planck y de su recta de acción}
\begingroup
\renewcommand{\chapter}[2][]{}%
\chapter{Derivación autodual de la constante de Planck: periodicidad temporal de \(108\) fases, acción y longitud de Compton}
\label{chap:p3-final:36:planck_autodual}
\label{chap:planck-autodual}

\section{Realización de Compton asociada a la periodicidad temporal de \(108\) fases}

Sobre el ciclo de \cref{eq:realizacion-circular-ct108}, fijemos una sección
positiva de acción \(\hbar\). El cuanto angular asociado es
\begin{equation}
E_{108}=\hbar\omega_{108},
\qquad
m_{108}=\frac{E_{108}}{c^2}
=\frac{\hbar\omega_{108}}{c^2}.
\label{eq:cuanto-ct108}
\end{equation}
No se ha empleado todavía \(G\). La geometría del ciclo y la cuantización
angular bastan para fijar sus dos lectores de Compton.

\begin{theorem}[Identidad entre la periodicidad temporal de \(108\) fases y la longitud de Compton]
\label{thm:ct108-compton}
Para el cuanto \eqref{eq:cuanto-ct108},
\begin{equation}
\lambdabarC(m_{108})=\frac{\hbar}{m_{108}c}=R_{108},
\qquad
\lambdaC(m_{108})=\frac{h}{m_{108}c}=C_{108}.
\label{eq:identidad-ct108-compton}
\end{equation}
Por tanto, el radio del ciclo es su longitud de Compton reducida y la
longitud de una vuelta es su longitud de Compton completa.
\end{theorem}

\begin{proof}
Sustituyendo \(m_{108}=\hbar\omega_{108}/c^2\),
\[
\frac{\hbar}{m_{108}c}
=\frac{\hbar c^2}{\hbar\omega_{108}c}
=\frac{c}{\omega_{108}}=R_{108}.
\]
Como \(h=2\pi\hbar\), la segunda igualdad es \(2\pi R_{108}=C_{108}\).
\end{proof}

La coincidencia no es un ajuste numérico. Es una identidad homogénea: vale
para todo \(\hbar,c,t_0>0\) y depende sólo de que el mismo periodo determine
la frecuencia, la longitud del ciclo y el cuanto de acción.

\section{Transductor gravitatorio y ecuación de selección autodual}

La recta gravitatoria HMT--MD contiene el transductor
\begin{equation}
G(\theta)=\theta\frac{c^5t_0^2}{\hbar},
\qquad \theta>0.
\label{eq:transductor-g-planck-anexo}
\end{equation}
La homogeneidad es exacta, pero no selecciona por sí sola \(\theta\). El
ciclo CT108 aporta una condición interna adicional: el lector gravitatorio
debe devolver la misma escala radial que el lector geométrico y el lector de
Compton.

Definimos el radio gravitatorio reducido
\begin{equation}
\rg(\theta):=\frac{G(\theta)m_{108}}{c^2}.
\label{eq:radio-gravitatorio-reducido}
\end{equation}
Al sustituir \eqref{eq:cuanto-ct108} y
\eqref{eq:transductor-g-planck-anexo},
\begin{equation}
\rg(\theta)
=\theta\frac{\pi}{54}\ell_0.
\label{eq:rg-theta-lineal}
\end{equation}

\begin{theorem}[Selector autodual de la periodicidad temporal de \(108\) fases]
\label{thm:selector-autodual-ct108}
La ecuación
\[
\rg(\theta)=R_{108}
\]
posee una única solución positiva:
\begin{equation}
\boxed{\thetaStar=\left(\frac{54}{\pi}\right)^2.}
\label{eq:theta-star}
\end{equation}
\end{theorem}

\begin{proof}
Por \eqref{eq:realizacion-circular-ct108} y
\eqref{eq:rg-theta-lineal}, la ecuación es
\[
\theta\frac{\pi}{54}\ell_0
=\frac{54}{\pi}\ell_0.
\]
Como \(\ell_0>0\), se simplifica y se obtiene
\(\theta=(54/\pi)^2\). La función
\(\theta\mapsto\rg(\theta)\) es estrictamente creciente en
\(\RR_{>0}\), de modo que no existe otra solución positiva.
\end{proof}

\begin{corollary}[Unidad de Planck interna]
\label{cor:unidad-planck-interna}
Con \(G_*=G(\thetaStar)\), las secciones de Planck satisfacen
\begin{equation}
\ellP:=\sqrt{\frac{\hbar G_*}{c^3}}=R_{108},
\qquad
\tP:=\sqrt{\frac{\hbar G_*}{c^5}}=\frac1{\omega_{108}},
\label{eq:planck-long-time-ct108}
\end{equation}
\begin{equation}
\mP:=\sqrt{\frac{\hbar c}{G_*}}=m_{108},
\qquad
\EP:=\mP c^2=E_{108}.
\label{eq:planck-mass-energy-ct108}
\end{equation}
En particular,
\begin{equation}
\boxed{
R_{108}=\lambdabarC=\rg=\ellP,
\qquad
C_{108}=\lambdaC=2\pi\rg=2\pi\ellP.}
\label{eq:tres-lectores-planck}
\end{equation}
\end{corollary}

\begin{proof}
Se sustituye \(G_*=(54/\pi)^2c^5t_0^2/\hbar\). Entonces
\[
\ellP=\frac{54}{\pi}ct_0=R_{108},
\qquad
\tP=\frac{54}{\pi}t_0=\frac1{\omega_{108}},
\]
y
\[
\mP=\frac{\pi}{54}\frac{\hbar}{c^2t_0}
=\frac{\hbar\omega_{108}}{c^2}=m_{108}.
\]
El resto se sigue del \cref{thm:ct108-compton}.
\end{proof}

\begin{remark}[Tres lectores, una sección]
La geometría circular publica \(R_{108}\) y \(C_{108}\); la cuantización de
Compton publica \(\lambdabarC\) y \(\lambdaC\); la realización gravitatoria
publica \(\rg\) y \(2\pi\rg\). En \(\thetaStar\) esas tres parejas son
coordenadas de una única sección. Esta es la razón precisa por la que la
unidad de Planck deja de aparecer como una combinación dimensional accidental
dentro de esta formalización.
\end{remark}

\section{Equivalencia de normalizaciones sobre la identidad autodual}

El transductor \eqref{eq:transductor-g-planck-anexo} admite dos formas de
fijar la escala elemental, que no deben confundirse.

\begin{enumerate}[label=\textup{(\roman*)}]
\item Si se fija \(\theta=1\), entonces
\(\sqrt{\hbar G/c^3}=\ell_0\). La transición elemental \(ct_0\) se identifica
con la longitud de Planck de esa carta.
\item Si se exige que la \emph{órbita completa} CT108 posea radio de Planck,
la condición selecciona \(\thetaStar=(54/\pi)^2\) y
\(R_{108}=\ellP\).
\end{enumerate}

No son valores rivales de \(G\). Son dos calibres de la misma ecuación
homogénea: uno asigna la unidad de Planck al paso elemental; el otro, al radio
del ciclo completo. El anexo utiliza el segundo porque su tesis concierne a
la identidad de la clausura CT108.

\section{Convenciones gravitatorias: radio reducido y radio de Schwarzschild}

La notación gravitatoria se fija de manera inequívoca:
\[
\rg=\frac{Gm}{c^2},
\qquad
\rs=\frac{2Gm}{c^2}=2\rg.
\]
El teorema anterior emplea \(\rg\), no \(\rs\). En la convención habitual de
masa de Planck, la igualdad \(\hbar/(mc)=Gm/c^2\) se produce en \(m=\mP\),
mientras
\begin{equation}
\frac{\hbar}{mc}=\frac{2Gm}{c^2}
\quad\Longleftrightarrow\quad
m=\frac{\mP}{\sqrt2}.
\label{eq:cruce-schwarzschild-factor-dos}
\end{equation}
Por ello la frase «escala gravitatoria» designa en la tesis rectora el radio
reducido. El radio de horizonte convencional permanece como lector distinto.

\section{Desdoblamiento bidireccional de la recta de acción}

La recta interna de acción posee dos secciones orientadas:
\begin{equation}
\hbarpre=H_5\,10^{-34}\mathcal U_S,
\qquad
\hbarret=\eta_{\mathrm{ret}}\,10^{-34}\mathcal U_S,
\qquad
\Ract=\frac{H_5}{\eta_{\mathrm{ret}}}.
\label{eq:hbar-pre-ret-anexo}
\end{equation}
La década \(10^{-34}\) fija la escala absoluta; el cociente \(\Ract\)
transporta la diferencia orientada. Esta separación permite preguntar si la
autodualidad de Planck sobrevive a ambas coordenadas.

\begin{theorem}[Covariancia recíproca de acción y gravedad]
\label{thm:accion-gravedad-reciproca}
Para \(a\in\{\mathrm{pre},\mathrm{ret}\}\), sea
\[
G_a(\theta)=\theta\frac{c^5t_0^2}{\hbar_a},
\qquad
m_{108,a}=\frac{\hbar_a\omega_{108}}{c^2}.
\]
Entonces
\begin{equation}
\frac{G_a(\theta)m_{108,a}}{c^2}
=\theta\frac{\pi}{54}\ell_0
\label{eq:invariante-pre-ret-planck}
\end{equation}
es independiente de \(a\). Además,
\begin{equation}
\frac{m_{108,\mathrm{pre}}}{m_{108,\mathrm{ret}}}=\Ract,
\qquad
\frac{G_{\mathrm{pre}}}{G_{\mathrm{ret}}}=\Ract^{-1},
\label{eq:reciprocidad-pre-ret}
\end{equation}
y la longitud de Planck
\(\sqrt{\hbar_aG_a/c^3}=\sqrt\theta\,\ell_0\) es la misma en ambas
orientaciones.
\end{theorem}

\begin{proof}
Los factores \(\hbar_a\) se cancelan en el producto \(G_am_{108,a}\).
Las dos razones de \eqref{eq:reciprocidad-pre-ret} se obtienen directamente
de las definiciones. La invariancia de la longitud de Planck se sigue de
\(\hbar_aG_a=\theta c^5t_0^2\).
\end{proof}

Una parametrización simétrica hace visible la geometría:
\begin{equation}
\hbar_{\mathrm{geom}}:=\sqrt{\hbarpre\hbarret},
\qquad
\xi:=\frac12\log\Ract,
\qquad
\hbarpre=\hbar_{\mathrm{geom}}e^\xi,
\quad
\hbarret=\hbar_{\mathrm{geom}}e^{-\xi}.
\label{eq:rapidez-accion-pre-ret}
\end{equation}
La memoria bidireccional se aloja en \(\xi\); la geometría de Planck se
aloja en el producto invariante. Esto permite una diferencia real entre las
dos coordenadas de acción sin exigir una diferencia de velocidad de
propagación.

\begin{remark}[Alcance físico del teorema]
La reciprocidad de \cref{thm:accion-gravedad-reciproca} es exacta dentro del
transductor. Para identificar \(G_{\mathrm{pre}}\) y \(G_{\mathrm{ret}}\) con dos
mediciones gravitatorias orientadas hace falta especificar el protocolo, el
observable y la simetría que relaciona ambos. El teorema prueba la
compatibilidad estructural; no anticipa por sí solo una violación de Lorentz
ni una variación observable de \(G\).
\end{remark}
 \label{mdg:sec:identidad-autodual-planck-ct108-rev3}
\index{Planck!identidad autodual}
\index{calendario de orden 108!realización circular}
\index{Compton!lecturas reducida y completa}

La recta de acción construida en la sección precedente admite una
realización circular compatible con la periodicidad temporal de orden 108.
Esta realización coordina tres objetos que ya poseen tipos propios: la
sección de acción, la geometría de una vuelta y la transducción
gravitatoria. La coordinación resultante proporciona una identidad de
Planck interna y prepara la recta positiva sobre la que actuará después el
operador general de masa.

\subsection{Cartas orientadas de acción y realización circular declarada}

Sean
\[
 \mathcal L_S^+,\quad \mathcal L_T^+,\quad \mathcal L_C^+,
 \quad \mathcal L_L^+,\quad \mathcal L_M^+,\quad \mathcal L_G^+
\]
las semirrectas positivas de acción, tiempo, velocidad, longitud y masa.
Las coordenadas anterior y posterior al retorno se escriben
\begin{equation}
 \hbar_{\rm pre}=H_5\,10^{-34}\mathcal U_S,
 \qquad
 \hbar_{\rm ret}=\eta_{\rm ret}\,10^{-34}\mathcal U_S,
 \qquad
 R_{\rm act}=\frac{\hbar_{\rm pre}}{\hbar_{\rm ret}}
 =\frac{H_5}{\eta_{\rm ret}}.
 \label{mdg:eq:planck-ct108-cartas-pre-ret-rev3}
\end{equation}
El subíndice \(a\in\{\mathrm{pre},\mathrm{ret}\}\) designará cualquiera
de ambas cartas. En cada una se conserva la distinción exacta
\begin{equation}
 h_a=2\pi\hbar_a .
 \label{mdg:eq:planck-ct108-h-hbar-rev3}
\end{equation}

Fijemos una duración elemental \(t_0\in\mathcal L_T^+\), una velocidad
interna \(c_{\rm int}\in\mathcal L_C^+\) y
\(\ell_0=c_{\rm int}t_0\). La realización circular declarada de la
periodicidad temporal de orden 108 queda determinada por
\begin{equation}
 T_{108}=108t_0,
 \qquad
 \omega_{108}=\frac{2\pi}{T_{108}},
 \qquad
 R_{108}=\frac{c_{\rm int}}{\omega_{108}}
 =\frac{54}{\pi}\ell_0,
 \label{mdg:eq:planck-ct108-realizacion-circular-rev3}
\end{equation}
y por la longitud de la vuelta completa
\begin{equation}
 C_{108}=2\pi R_{108}=c_{\rm int}T_{108}=108\ell_0.
 \label{mdg:eq:planck-ct108-perimetro-rev3}
\end{equation}
Así, \(R_{108}\) es el radio de la realización y \(C_{108}\) su
perímetro. La frecuencia angular, la frecuencia cíclica y las dos
constantes de acción se relacionan mediante
\(\omega_{108}=2\pi/T_{108}\) y \(h_a=2\pi\hbar_a\).

\subsection{Correspondencia entre la periodicidad de orden 108 y las lecturas de Compton}

En la carta \(a\), el cuanto circular y su sección de masa son
\begin{equation}
 E_{108,a}=\hbar_a\omega_{108},
 \qquad
 m_{108,a}=\frac{E_{108,a}}{c_{\rm int}^{2}}
 =\frac{\hbar_a\omega_{108}}{c_{\rm int}^{2}}.
 \label{mdg:eq:planck-ct108-cuanto-circular-rev3}
\end{equation}

\begin{theorem}[Identidad entre la periodicidad circular y las lecturas de Compton]
\label{mdg:thm:planck-ct108-compton-rev3}
Para cada carta \(a\in\{\mathrm{pre},\mathrm{ret}\}\), las lecturas de
Compton reducida y completa satisfacen
\begin{equation}
 \boxed{
 \bar\lambda_{\rm C}(m_{108,a})
 =\frac{\hbar_a}{m_{108,a}c_{\rm int}}=R_{108},
 \qquad
 \lambda_{\rm C}(m_{108,a})
 =\frac{h_a}{m_{108,a}c_{\rm int}}=C_{108}.}
 \label{mdg:eq:planck-ct108-identidad-compton-rev3}
\end{equation}
\end{theorem}

\begin{proof}
La sustitución de
\eqref{mdg:eq:planck-ct108-cuanto-circular-rev3} produce
\[
 \frac{\hbar_a}{m_{108,a}c_{\rm int}}
 =\frac{c_{\rm int}}{\omega_{108}}=R_{108}.
\]
La segunda igualdad resulta al aplicar simultáneamente
\eqref{mdg:eq:planck-ct108-h-hbar-rev3} y
\eqref{mdg:eq:planck-ct108-perimetro-rev3}.
\end{proof}

El teorema es homogéneo en
\((\hbar_a,c_{\rm int},t_0)\): la misma periodicidad determina el radio,
el perímetro y las dos cartas de Compton. La realización circular es la
primitiva geométrica declarada; las igualdades precedentes son sus
consecuencias exactas.

\subsection{Selector circular y unidad interna de Planck}

La recta gravitatoria se coordina con la carta \(a\) mediante el transductor
homogéneo
\begin{equation}
 G_a(\theta)=\theta\,
 \frac{c_{\rm int}^{5}t_0^2}{\hbar_a}
 \in\mathcal L_G^+,
 \qquad \theta\in\mathbb R_{>0}.
 \label{mdg:eq:planck-ct108-transductor-rev3}
\end{equation}
Su radio gravitatorio reducido, evaluado sobre el cuanto circular, es
\begin{equation}
 r_{{\rm g},a}(\theta)
 :=\frac{G_a(\theta)m_{108,a}}{c_{\rm int}^{2}}
 =\theta\frac{\pi}{54}\ell_0.
 \label{mdg:eq:planck-ct108-radio-gravitatorio-rev3}
\end{equation}

\begin{proposition}[Selector circular de Planck]
\label{mdg:prop:selector-circular-planck-ct108-rev3}
Dentro de la realización circular
\eqref{mdg:eq:planck-ct108-realizacion-circular-rev3}, la condición
\(r_{{\rm g},a}(\theta)=R_{108}\) selecciona una única coordenada
positiva,
\begin{equation}
 \boxed{\theta_{108}^{\circ}=\left(\frac{54}{\pi}\right)^2.}
 \label{mdg:eq:theta-circular-planck-ct108-rev3}
\end{equation}
La coordenada es independiente de la carta orientada \(a\).
\end{proposition}

\begin{proof}
Las ecuaciones
\eqref{mdg:eq:planck-ct108-realizacion-circular-rev3} y
\eqref{mdg:eq:planck-ct108-radio-gravitatorio-rev3} reducen la condición a
\[
 \theta\frac{\pi}{54}\ell_0=\frac{54}{\pi}\ell_0.
\]
La positividad de \(\ell_0\) y la estricta monotonía de la aplicación
\(\theta\mapsto r_{{\rm g},a}(\theta)\) proporcionan la solución y su
unicidad.
\end{proof}

Sea \(G_{108,a}^{\circ}:=G_a(\theta_{108}^{\circ})\). Las unidades internas
asociadas a esta carta satisfacen
\begin{align}
 \ell_{{\rm P},a}^{\circ}
 &:=\sqrt{\frac{\hbar_a G_{108,a}^{\circ}}{c_{\rm int}^{3}}}
 =R_{108},
 &
 t_{{\rm P},a}^{\circ}
 &:=\sqrt{\frac{\hbar_a G_{108,a}^{\circ}}{c_{\rm int}^{5}}}
 =\omega_{108}^{-1},
 \label{mdg:eq:planck-ct108-longitud-tiempo-rev3}\\
 m_{{\rm P},a}^{\circ}
 &:=\sqrt{\frac{\hbar_a c_{\rm int}}{G_{108,a}^{\circ}}}
 =m_{108,a},
 &
 E_{{\rm P},a}^{\circ}
 &:=m_{{\rm P},a}^{\circ}c_{\rm int}^{2}=E_{108,a}.
 \label{mdg:eq:planck-ct108-masa-energia-rev3}
\end{align}
La igualdad de radios se publica, por tanto, como
\begin{equation}
 R_{108}=\bar\lambda_{\rm C}=r_{\rm g}=\ell_{\rm P}^{\circ},
 \qquad
 C_{108}=\lambda_{\rm C}=2\pi r_{\rm g}=2\pi\ell_{\rm P}^{\circ}.
 \label{mdg:eq:planck-ct108-tres-lectores-rev3}
\end{equation}

La convención de este teorema es el radio gravitatorio reducido
\(r_{\rm g}=Gm/c_{\rm int}^{2}\). El radio de Schwarzschild satisface
\(r_{\rm s}=2r_{\rm g}\); si se adopta esta segunda convención, su cruce con
\(\bar\lambda_{\rm C}\) ocurre en
\(m=m_{{\rm P},a}^{\circ}/\sqrt2\). La declaración del radio forma, por
tanto, parte de la carta y evita trasladar inadvertidamente un factor dos
entre las dos identidades.

La coordenada \(\theta_{108}^{\circ}\) pertenece a la especialización
circular declarada. El selector gravitatorio general, formulado sobre la
holonomía enriquecida y su respuesta espectral, conserva su residencia en
la \cref{mdg:sec:selector-gravitatorio-rev11}. De este modo, la identidad
circular fija una normalización interna de Planck y el criterio general
fija las condiciones que debe satisfacer la realización física de la
sección gravitatoria.

\subsection{Covariancia pre-retorno/post-retorno}

El producto de las secciones recíprocas de acción y gravedad conserva la
geometría circular:
\begin{equation}
 \frac{m_{108,\rm pre}}{m_{108,\rm ret}}=R_{\rm act},
 \qquad
 \frac{G_{108,\rm pre}^{\circ}}{G_{108,\rm ret}^{\circ}}
 =R_{\rm act}^{-1},
 \qquad
 \hbar_a G_{108,a}^{\circ}
 =\theta_{108}^{\circ}c_{\rm int}^{5}t_0^2.
 \label{mdg:eq:planck-ct108-covariancia-pre-ret-rev3}
\end{equation}
En particular, \(\ell_{{\rm P},\rm pre}^{\circ}\) y
\(\ell_{{\rm P},\rm ret}^{\circ}\) coinciden con \(R_{108}\). La memoria
bidireccional reside en el cociente de las masas y en su recíproco
gravitatorio; la escala geométrica reside en el producto invariante.

La reciprocidad admite una parametrización simétrica que separa la escala
común de la orientación del retorno. Defínanse
\begin{equation}
 \begin{aligned}
 \hbar_{\rm geom}
 &:=\sqrt{\hbar_{\rm pre}\hbar_{\rm ret}},
 &
 G_{108,\rm geom}^{\circ}
 &:=\sqrt{G_{108,\rm pre}^{\circ}G_{108,\rm ret}^{\circ}},
 &
 \xi_{\rm act}&:=\frac12\log R_{\rm act}.
 \end{aligned}
 \label{mdg:eq:planck-ct108-parametrizacion-simetrica-rev3}
\end{equation}
Entonces
\[
 \hbar_{\rm pre}=\hbar_{\rm geom}e^{\xi_{\rm act}},
 \qquad
 \hbar_{\rm ret}=\hbar_{\rm geom}e^{-\xi_{\rm act}},
\]
mientras las secciones gravitatorias conjugadas satisfacen
\[
 G_{108,\rm pre}^{\circ}
 =G_{108,\rm geom}^{\circ}e^{-\xi_{\rm act}},
 \qquad
 G_{108,\rm ret}^{\circ}
 =G_{108,\rm geom}^{\circ}e^{\xi_{\rm act}}.
\]
La coordenada \(\xi_{\rm act}\) conserva la orientación
pre-retorno/post-retorno; el producto \(\hbar_aG_{108,a}^{\circ}\)
conserva la escala geométrica. Esta parametrización de las cartas de acción
no introduce por sí sola velocidades distintas de propagación.

\subsection{Recta autodual partida de masas}

Fijada una carta \(a\), sea el álgebra partida
\[
 \mathbb D_{\rm sp}:=\mathbb R[j]/(j^2-1),
 \qquad
 P_\pm:=\frac{I\pm j}{2},
\]
y considérese su descomposición espectral
\[
 \mathbb D_{\rm sp}
 =\mathbb R P_+\oplus\mathbb R P_-,
 \qquad
 P_\pm^2=P_\pm,\quad
 P_+P_-=0,\quad
 P_++P_-=I.
\]
La conjugación partida intercambia los proyectores,
\[
 \overline{uP_++vP_-}=vP_++uP_-,
\]
y su norma es
\[
 N_{\rm sp}(uP_++vP_-):=uv.
\]

Para \(m\in\mathcal L_M^+\), defínase la sección
\begin{equation}
 Z_a(m)
 :=\frac{m_{{\rm P},a}^{\circ}}{m}P_+
   +\frac{m}{m_{{\rm P},a}^{\circ}}P_- .
 \label{mdg:eq:planck-ct108-seccion-partida-masa-rev3}
\end{equation}
La unidad \(m_{{\rm P},a}^{\circ}\) determina asimismo la involución
\begin{equation}
 \iota_{{\rm P},a}(m)
 =\frac{(m_{{\rm P},a}^{\circ})^2}{m}.
 \label{mdg:eq:planck-ct108-involucion-masas-rev3}
\end{equation}

\begin{theorem}[Recta autodual partida de masas]
\label{mdg:thm:planck-ct108-recta-partida-masa-rev3}
Para todo \(m\in\mathcal L_M^+\),
\begin{equation}
 \boxed{
 N_{\rm sp}(Z_a(m))=1,
 \qquad
 \overline{Z_a(m)}
 =Z_a\!\left(\iota_{{\rm P},a}(m)\right).}
 \label{mdg:eq:planck-ct108-norma-conjugacion-rev3}
\end{equation}
La conjugación posee un único punto fijo positivo,
\[
 m=m_{{\rm P},a}^{\circ},
 \qquad
 Z_a(m_{{\rm P},a}^{\circ})=I.
\]
\end{theorem}

\begin{proof}
La norma es el producto de las dos coordenadas recíprocas:
\[
 \frac{m_{{\rm P},a}^{\circ}}m
 \frac m{m_{{\rm P},a}^{\circ}}=1.
\]
La conjugación intercambia ambas coordenadas, lo que equivale a sustituir
\(m\) por \((m_{{\rm P},a}^{\circ})^2/m\). El punto fijo positivo satisface
\(m^2=(m_{{\rm P},a}^{\circ})^2\).
\end{proof}

Con la rapidez
\begin{equation}
 x_a(m)=\log\frac{m}{m_{{\rm P},a}^{\circ}},
 \label{mdg:eq:planck-ct108-rapidez-masa-rev3}
\end{equation}
la sección y su conjugada se escriben
\[
 Z_a(m)=e^{-x_a}P_++e^{x_a}P_-,
 \qquad
 \overline{Z_a(m)}
 =e^{x_a}P_++e^{-x_a}P_-.
\]
Sus lectores par e impar son
\begin{equation}
 Q_a(m):=\frac{e^{x_a}+e^{-x_a}}2=\cosh x_a,
 \qquad
 A_a(m):=\frac{e^{x_a}-e^{-x_a}}2=\sinh x_a.
 \label{mdg:eq:planck-ct108-lectores-par-impar-rev3}
\end{equation}
El primero conserva la distancia al punto autodual y el segundo conserva su
orientación.

Los lectores físicos recíprocos asociados a la misma coordenada son
\[
 r_{{\rm g},a}(m)=\frac{G_{108,a}^{\circ}m}{c_{\rm int}^{2}},
 \qquad
 \bar\lambda_{{\rm C},a}(m)=\frac{\hbar_a}{mc_{\rm int}},
\]
y satisfacen
\begin{equation}
 \frac{r_{{\rm g},a}(m)}{\bar\lambda_{{\rm C},a}(m)}
 =e^{2x_a(m)},
 \qquad
 \frac{\bar\lambda_{{\rm C},a}(m)}{r_{{\rm g},a}(m)}
 =e^{-2x_a(m)}.
 \label{mdg:eq:planck-ct108-lectores-exponenciales-rev3}
\end{equation}
La conjugación intercambia los dos lectores; la identidad de Planck es su
equilibrio unitario.

\begin{remark}[Identidad de Planck y clausura electrónica]
La identidad \(m=m_{{\rm P},a}^{\circ}\) fija el origen autodual de la recta
positiva de masas. La sección electrónica fija una posición distinta de
esa misma recta mediante una genealogía propia. La
\cref{mdg:sec:ley-general-masa-accion-autonoma} construye el operador universal
que transporta escala absoluta, carácter, torres y lectura bidireccional;
la \cref{mdg:sec:cierre-electron-two-way-rev11} realiza después su reducción
central de rango uno. Así, la unidad de Planck proporciona la identidad de
la coordenada y el electrón proporciona la primera transición material
completa sobre ella.
\end{remark}

\begin{proposition}[Realización genealógica del centro electrónico]
\label{mdg:prop:planck-ct108-realizacion-centro-electron-rev3}
Sea
\[
 F_e^+=\mathbb R_{>0}(P_++P_-)
\]
la semirrecta positiva de la fibra central de rango uno, sea
\(\mathcal G_e\) el espacio de genealogías
electrónicas enriquecidas y sea \(\mathcal L_M^+\) la recta positiva de
masas. Un morfismo de realización
\[
 \mathfrak R_e:
 F_e^+\times\mathcal G_e\longrightarrow\mathcal L_M^+
\]
puede enviar la identidad \(I=P_++P_-\) a una masa distinta de
\(m_{{\rm P},a}^{\circ}\) sin alterar el
\cref{mdg:thm:planck-ct108-recta-partida-masa-rev3}: la identidad de Planck fija el
origen autodual de la carta, mientras la genealogía electrónica fija un
desplazamiento sobre esa carta.
\end{proposition}

\begin{proof}
Si \(\chi_e:\mathcal G_e\to\mathbb R_{>0}\) es el carácter positivo construido
por el registro electrónico, la regla
\[
 \mathfrak R_e(\lambda I,g)
 =\lambda\,m_{{\rm P},a}^{\circ}\chi_e(g)
\]
es positiva y multiplicativa en la coordenada genealógica. Para la genealogía
neutra, \(\chi_e=1\), se recupera el origen de Planck; para una genealogía con
\(\chi_e(g)\ne1\), la misma identidad interna ocupa una posición distinta de
la recta. La clausura electrónica concreta conserva su desarrollo antecedente en la
\cref{mdg:sec:cierre-electron-two-way-rev11}.
\end{proof}
 \section{Autodualidad de la unidad interna de Planck bajo inversión de orientación}

El selector CT108 construye

\begin{equation}
\theta_{108}=\left(\frac{54}{\pi}\right)^2,
\qquad
\ell_P=\frac{54}{\pi}\ell_0,
\qquad
t_P=\frac{54}{\pi}t_0,
\label{eq:planck-length-time}
\end{equation}

y

\begin{equation}
E_P=\frac{\hbar_{\rm int}}{t_P}
=\frac{\pi\hbar_{\rm int}}{54t_0}.
\label{eq:planck-energy}
\end{equation}

Comparando \cref{eq:boltzmann-clock,eq:planck-energy} se obtiene

\begin{equation}
\boxed{E_P=k_B\Theta_{\rm clk}\log3.}
\label{eq:planck-trit}
\end{equation}

La relaci\'on no depende de una coincidencia entre unidades: ambos miembros
proceden de la misma estructura cronol\'ogica \(108\).

Cuando la constante gravitatoria del \cref{ch:gravity} se introduce, la familia de Planck satisface

\begin{align}
\ell_P^2&=\frac{G\hbar}{c^3}=\theta_{108}\ell_0^2,
\label{eq:planck-area}\\
F_P&=\frac{c^4}{G},
\qquad
P_P=\frac{c^5}{G},
\label{eq:planck-force-power}\\
\rho_P&=\frac{\hbar}{\theta_{108}^2c^5t_0^4}.
\label{eq:planck-density}
\end{align}

Aunque \(G\) y \(\hbar\) se transformen rec\'iprocamente entre las secciones
anterior y posterior al retorno, su producto en \(\ell_P^2\) conserva la
geometr\'ia. Esta invariancia determina el \`area como interfaz entre las
realizaciones geom\'etrica, informacional y de entrelazamiento.
  \endgroup
% Distribución causal exacta del delta Ley 9; contenido conservado sin síntesis.
\section{Recta de acción y constante de Planck}
\label{sec:l9s102:recta-planck}

Sea \(\mathcal L_S\) la recta real de acción con base positiva
\(\mathcal U_S\). La transducción
\begin{equation}
 \tau_S:\mathbb R_{>0}\longrightarrow\mathcal L_S,
 \qquad u\longmapsto u\,10^{\varepsilon_H}\mathcal U_S
 \label{eq:l9s102:transductor-accion}
\end{equation}
publica las secciones previa y posterior al retorno:
\begin{equation}
 \hbar_{\mathrm{pre}}=H_5\,10^{-34}\mathcal U_S,
 \qquad
 \boxed{
 \hbar_{\mathrm{ret}}=\eta_{\mathrm{ret}}\,10^{-34}\mathcal U_S.}
 \label{eq:l9s102:hbar-pre-ret}
\end{equation}
Las acciones de ciclo completo son
\begin{equation}
 h_{\mathrm{pre}}=2\pi\hbar_{\mathrm{pre}},
 \qquad
 h_{\mathrm{ret}}=2\pi\hbar_{\mathrm{ret}}.
 \label{eq:l9s102:h-pre-ret}
\end{equation}
Para una fase \(\theta\), las cartas angular y de acción conmutan:
\begin{equation}
 S(\theta)
 =\hbar_{\mathrm{ret}}\theta_{\mathrm{rad}}
 =\hbar_{\mathrm{ret}}\frac{\pi}{180}\theta_{\deg}
 =h_{\mathrm{ret}}\frac{\theta_{\deg}}{360}.
 \label{eq:l9s102:accion-fase}
\end{equation}

 
\chapter{Operador de fase, radio espectral y factorización de la doble proyección}
\begingroup
\renewcommand{\chapter}[2][]{}%
\chapter{Operador angular predimensional: descomposición en \(27\) sectores, radio espectral y composición efectiva de la doble proyección}
\label{chap:p3-final:37:operador_angular}
\label{p3-20260826:chap:operador-angular-dictamen}
\index{operador angular}
\index{fase angular!salida espectral}
\index{cierre!pre-dimensional}

Las dos secciones anteriores han reunido tres piezas que antes aparecían
separadas: una identificación isométrica entre proyectores, una medida de
cociente y un cociclo positivo de acción. La presente sección estudia la consecuencia
espectral de esa ley y sitúa el resultado dentro de la arquitectura completa
de HMT\@.

La separación de estatutos es precisa. La identificación triple es un teorema
exacto relativo a la nueva regla variacional interna y a un ancla orientada.
La medida y el carácter positivo son teoremas exactos de \HTr--\HAdd; el defecto
es relativo además a \HDef, la capacidad a \HRes y el coste local a \HLoc. El valor
angular es una salida intervalarmente certificada del operador \Gtr. Su
proximidad a la coordenada angular conjugada canónica es extraordinaria, pero el intervalo de
la diferencia excluye la igualdad.

\section{Descomposición del operador angular en \(27\) sectores de \(72\) estados}

El lector de todos los caminos \Gtr posee \(1944\) estados microscópicos. La
acción de traslación gruesa de \((\mathbb Z/3\mathbb Z)^3\) permite la
descomposición
\[
 1944=27\cdot72.
\]
Cada bloque de \(72\) estados se parametriza por
\[
 (r,s,m,d)
 \in(\mathbb Z/3\mathbb Z)^2\times\{0,1\}\times D_4,
\]
donde \(m\) es la hoja y \(d\) la dirección anterior. Desde cada estado
parten cuatro continuaciones cardinales. El TRIT de la celda de llegada
decide la hoja siguiente.

Con la acción del \cref{p3-20260826:chap:medida-accion-positiva}, el peso de una arista
es
\[
 \exp[-g(e)-\lambda_*s(e)].
\]
La transformada de caracteres reduce el operador a matrices complejas
\(L_{k_a,k_b}\) de tamaño \(72\times72\). El programa enumera las
aristas de cada bloque; no recibe una matriz espectral precalculada. La
iteración de potencias se conserva como exploración, pero el resultado que
sigue ya no descansa en ella: los radios espectrales se aíslan mediante una
certificación racional completa.

Sean \(\rho_{01}\) y \(\rho_{12}\) los radios espectrales de los sectores
orientados \((0,1)\) y \((1,2)\). Definimos la \emph{fase de banda}
\[
 \boxed{y_{\rm band}=\log\rho_{12}-\log\rho_{01}.}
\]
Para
\[
 \lambda_*
 =1-\frac5{468}-\frac3{11}\Delta_4,
\]
la certificación da
\[
\begin{aligned}
1.460655120862011111
&<\rho_{01}<1.460655120862404024,\\
1.515812008270944328
&<\rho_{12}<1.515812008271195416,
\end{aligned}
\]
y, por propagación intervalar,
\[
 0.037066226434427529
 <y_{\rm band}<
 0.037066226434862172,
\]
\[
 \boxed{
 2.123738337168943^\circ
 <C_{\rm band}:=\frac{180}{\pi}y_{\rm band}
 <2.123738337193847^\circ.}
\]
La cantidad primaria es \(y_{\rm band}\), un logaritmo adimensional de cociente
espectral. Los grados aparecen después al aplicar la carta
\(180/\pi\). El operador produce primero una fase interna y sólo después
una representación angular arquimediana.

\subsection{Cálculo exacto del radio espectral del operador angular}

Los bloques son no normales. Por ello, un residuo pequeño
\(\|Lv-\mu v\|\) no basta para demostrar que \(\mu\) sea el autovalor de
mayor módulo. La prueba utilizada aquí consta de cuatro pasos.

Primero se observan doce filas estructuralmente nulas en cada bloque. Al
expandir el polinomio característico por ellas se obtiene un factor
\(z^{12}\) y un menor principal de orden sesenta. En éste aparecen otras
ocho filas estructuralmente nulas. Por tanto,
\[
 \det(zI_{72}-L)=z^{20}\det(zI_{52}-A),
\]
y todo autovalor no nulo queda contenido, con multiplicidad, en un bloque
reducido \(A\) de orden cincuenta y dos.

Segundo, el certificado proporciona matrices gaussianas racionales
\(S,R\), con denominador \(10^{16}\). El verificador calcula exactamente
\(E=I-RS\) y obtiene
\[
 \|E\|_\infty\le8.5664\times10^{-14}
 \quad\hbox{y}\quad
 \|E\|_\infty\le3.7628\times10^{-14}
\]
para los sectores \((0,1)\) y \((1,2)\), respectivamente. La serie de
Neumann demuestra entonces que \(S\) es invertible.

Tercero, si \(A_0\) es el punto medio racional del bloque intervalar, la
identidad
\[
 S^{-1}AS=(I-E)^{-1}RAS
\]
acota la distancia a \(C_0=RA_0S\) por
\[
 1.25125\times10^{-13}
 \quad\hbox{y}\quad
 5.7037\times10^{-14}.
\]
Los discos de Gershgorin de \(C_0\), ampliados por esas cotas, contienen
los discos de la similitud exacta. En cada bloque un disco queda separado
de los otros cincuenta y uno; el segundo teorema de Gershgorin le asigna
exactamente un autovalor. Su módulo inferior supera el módulo superior de
todos los demás discos. Ésa es la prueba de que el autovalor aislado es el
radio espectral.

Cuarto, \(\pi,e,\log\pi,e^{-1},e^{-\lambda_*}\) y \(\sqrt3\) no se
aceptan como decimales de biblioteca. El verificador los encierra mediante
la fórmula de Machin, las series factorial y logarítmica con resto, la
serie alternante de la exponencial y una horquilla racional cuadrática para
\(\sqrt3\). La verificación emplea únicamente enteros y fracciones; el
argumento queda cerrado por la reducción racional, la serie de Neumann y
la separación de Gershgorin expuestas en este apartado. Las expansiones
decimales reproducibles a profundidad diez mil se publican, como testigos
finitos independientes, en el
\hyperref[app:desarrollos-10000-rev9]{apéndice de expansiones}.

\section{Operador efectivo de banda \texorpdfstring{\(T_{\rm band}\)}{Tband} y aislamiento de la fase angular}

Sea
\[
 x=\frac{50\pi}{9}\alpha
 =0.127362829012869966\ldots
\]
y sea
\[
 R=\begin{pmatrix}1&0\\0&-1\end{pmatrix}.
\]
El operador bicapa efectivo asociado exclusivamente a la fase de banda es
\[
 \boxed{T_{\rm band}=\exp(-xI+y_{\rm band}R).}
\]
En la base de sectores orientados,
\[
 T_{\rm band}
 =\begin{pmatrix}
 e^{-x+y_{\rm band}}&0\\
 0&e^{-x-y_{\rm band}}
 \end{pmatrix}.
\]
Se verifican
\[
 \det T_{\rm band}=e^{-2x},
 \qquad
 \frac{(T_{\rm band})_{11}}{(T_{\rm band})_{22}}
 =e^{2y_{\rm band}}.
\]
Por tanto,
\[
 -\frac9{100\pi}\log\det T_{\rm band}=\alpha,
 \qquad
 \frac{90}{\pi}
 \log\frac{(T_{\rm band})_{11}}{(T_{\rm band})_{22}}
 =C_{\rm band}.
\]

Las dos recuperaciones tienen estatutos distintos. La segunda devuelve la
salida producida por la brecha. La primera vuelve a leer el
\(\alpha\) que entró en \(x\); demuestra reversibilidad de la
carta, no una segunda derivación independiente de \(\alpha\).

La forma exponencial separa dos funciones. El factor común \(e^{-x}\)
describe la contracción compartida de las dos hojas. Los factores
\(e^{\pm y_{\rm band}}\) distinguen orientaciones conservando el determinante. Como
\(I\) y \(R\) conmutan, el par
\[
 \left(\det T_{\rm band},
 \frac{(T_{\rm band})_{11}}{(T_{\rm band})_{22}}\right)
\]
reconstruye unívocamente \((x,y_{\rm band})\).

\section{Comparación posterior y escala de precisión}

La sección siguiente construye otra fase, la \emph{fase regional
regularizada} \(y_{\rm reg}\) ---denotada allí también por \(y_*\)---. Para
formular aquí una comparación con ese resultado posterior, anticipamos la
notación
\[
 C_{\rm reg}:=\frac{180}{\pi}y_{\rm reg}
 =2.123738338968646^\circ.
\]
Esta línea no constituye una segunda definición ni una entrada del operador
espectral: cita la conclusión del
\cref{p3-20260826:thm:contraangulo-regional}, cuya construcción se desarrolla en el
apartado siguiente.
No se identifican los dos pares de objetos:
\[
 (y_{\rm band},C_{\rm band})
 \ne
 (y_{\rm reg},C_{\rm reg}).
\]
La diferencia intervalar certificada de la ley lineal es
\[
 -1.799703\times10^{-9}\ {}^\circ
 <C_{\rm band}-C_{\rm reg}
 <-1.774799\times10^{-9}\ {}^\circ,
\]
con diferencia relativa encerrada por
\[
 -8.475\times10^{-10}
 <\frac{C_{\rm band}-C_{\rm reg}}{C_{\rm reg}}
 <-8.356\times10^{-10}.
\]
La ley que produce \(C_{\rm band}\) no recibe \(C_{\rm reg}\), CKM, Barbero,
masas ni datos SI\@. Esta independencia de objetivo hace de la proximidad un
candidato de reconocimiento fuerte y falsable. La certificación espectral
elimina la antigua laguna numérica; la diferencia intervalar, que excluye
cero, fija ahora el límite exacto: \(C_{\rm band}\) es una salida espectral
pre-dimensional distinta de la realización regional
\(C_{\rm reg}\).

El término adicional
\[
 -\frac{\Delta_4^2}{27}
\]
en \(\lambda\) produce
\[
 C_{\rm num}\approx2.12373833894104980{}^\circ,
\]
a \(2.76\times10^{-11}\) grados de \(C_{\rm reg}\). El entero \(27\) es interno,
pues cuenta los caracteres de \((\mathbb Z/3\mathbb Z)^3\), pero ese
hecho no fuerza el coeficiente del segundo momento. La aditividad exacta de
\HAdd excluye términos cuadráticos. La variante se conserva como ablación y no
como ley principal.

\section{Dependencia funcional del radio espectral respecto de los componentes del operador}

Todas las filas siguientes usan el mismo operador; sólo cambia una entrada
de la ley. Son diagnósticos numéricos de ablación; el intervalo certificado
precedente corresponde exclusivamente a la ley completa
\HTr--\HAdd--\HDef--\HRes--\HLoc (y, dentro de ella, el defecto positivo se fija en el
estrato \HTr--\HAdd--\HDef).
\[
\begin{array}{l|c|c}
\text{regla}&\lambda&C_{\rm num}\text{ en grados}\\ \hline
\text{autodual}&1&2.082755794784418\\
\text{sólo microfibra}&1-5/468&2.123621809570142\\
\text{traza }5/468+(3/11)\Delta_4
&0.989285913019141&2.123738337181414\\
4/468\text{ en lugar de }5/468
&0.991422665155894&2.115535228138553\\
3/12\text{ en lugar de }3/11
&0.989288440210566&2.123728626433667\\
5/243\text{ en lugar de }5/468
&0.979393542015659&2.161907060464779\\
5/729\text{ en lugar de }5/468
&0.993110963140488&2.109064222037112\\
7/729\text{, medida de semillas}
&0.990367478915522&2.119584299852700.
\end{array}
\]
Las ablaciones demuestran que la salida distingue la microfibra completa,
la dimensión reducida once y la medida de emisiones. No asignan por sí
solas una probabilidad de azar: tal probabilidad exigiría definir antes un
espacio de modelos alternativos y una medida sobre él.

\section{Composición tipada de los morfismos arquimediano, excepcional y dimensional}

La construcción completa puede escribirse como una sucesión de aplicaciones
con dominios visibles:
\[
\begin{aligned}
\Omega_{\rm seed}
&\xrightarrow{F}\mathcal U_6
\xrightarrow{J_F}H_U,\\
\mathscr R_\pi^{++}
&\xrightarrow{\beta_\pi}P_GH_G,\\
L_{\rm or}\otimes P_GH_G
&\xrightarrow{W_{GK}}P_3V_{11}
\xrightarrow{U_W}Q_3E_{30},\\
(\Pi_\pi^{(468)},\Delta_4P_3)
&\longmapsto D_{\rm switch}
\xrightarrow{\tau}\delta_*
\xrightarrow{\rm \HRes}\lambda_*,\\
\lambda_*
&\xrightarrow{\rm \HLoc}a_*
\xrightarrow{\mathcal L_{1,*}}y_{\rm band}
\xrightarrow{180/\pi}C_{\rm band},\\
(\alpha,y_{\rm band})
&\xrightarrow{\exp(-xI+y_{\rm band}R)}T_{\rm band}.
\end{aligned}
\]
Cada flecha declara lo que conserva:
\begin{itemize}
  \item \(F\) conserva la clase de emisión y deja en la fibra su
  multiplicidad de presentación;
  \item \(J_F\) compensa esa multiplicidad en la norma;
  \item \(\beta_\pi\) conserva la posición excepcional;
  \item \(W_{GK}\) conserva la representación orientada;
  \item \(U_W\) conserva incidencia y métrica;
  \item la traza olvida coordenadas y conserva rangos;
  \item el operador de transferencia suma caminos con el peso de la acción;
  \item la carta \(180/\pi\) representa en grados una fase interna.
\end{itemize}

\section{Diagrama de dependencias del operador angular y sus publicaciones}

\begin{longtable}{>{\raggedright\arraybackslash}p{.31\textwidth}>{\raggedright\arraybackslash}p{.22\textwidth}>{\raggedright\arraybackslash}p{.37\textwidth}}
\toprule
Objeto&Naturaleza matemática&Contenido demostrado\\
\midrule
\endfirsthead
\toprule
Objeto&Naturaleza matemática&Contenido demostrado\\
\midrule
\endhead
\(U_W=I/6:V_{11}\to E_{30}\)
&teorema exacto
&isometría, sobreyectividad, incidencia equivariante y unicidad positiva\\
\(P_3\leftrightarrow Q_3\)
&teorema exacto
&proyectores conjugados, rango tres y traza \(3/11\)\\
estrella \(R_{B_K}\)
&teorema exacto tipado
&espectro \(1^{[7]},(-1)^{[4]}\), índice \(3/11\) y obstrucción de conjugación\\
\(P_G\) en \(t=7\)
&teorema finito exacto
&filtración \(11\to3\to1\), proyector positivo y acción \(S_3\)\\
\(\mathscr R_\pi^{++}\leftrightarrow P_G\)
&teorema exacto
&bijección \(S_3\)-equivariante por posición excepcional\\
identificación isométrica orientada
&teorema relativo nuevo
&selector único por \(B_K,u_K\), Reynolds, polar y reducción \(240\to1\)\\
medida \(1/468\)
&teorema exacto de \HTr
&traza normalizada única y aplicación polar inducida por fibras\\
\(D_{\rm switch}\) y \(\delta_*\)
&teorema exacto relativo a \HTr--\HAdd--\HDef
&positividad, espectro y traza \(5/468+(3/11)\Delta_4\)\\
\(\lambda_*\)
&teorema interno relativo a \HTr--\HAdd--\HDef--\HRes
&capacidad residual obtenida del defecto mediante la regla \HRes\\
\(\mathcal A_*\)
&teorema interno relativo a \HTr--\HAdd--\HDef--\HRes--\HLoc
&cociclo aditivo de caminos con coste local \(g+\lambda_*s\)\\
\(T_{\rm band}\) y \(C_{\rm band}\)
&teorema computacional intervalar
&reducción \(72\to60\to52\), similitud racional, aislamiento de
Gershgorin y salida cuya distancia certificada está entre
\(1.774799\times10^{-9}\) y \(1.799703\times10^{-9}\) grados de
\(C_{\rm reg}\)\\
\(C_{\rm band}=C_{\rm reg}\)
&afirmación refutada para la ley lineal
&el intervalo certificado de la diferencia excluye cero\\
transición total
\(\mathcal X_{\mathrm{TPK}}^{\mathrm{enr}}\to \Gtr\)
&realizador global requerido
&\HTr--\HAdd--\HDef--\HRes--\HLoc define el lector relativo; la transición
total exige además una medida imagen dinámica sobre \(\Gtr\)\\
\bottomrule
\end{longtable}

\section{Acción predimensional en el centro de la doble proyección del estado genealógico}

La doble proyección parte del mismo estado holonómico integral y produce dos
lecturas. En la rama arquimediana, los cilindros compatibles se refinan y
sus diámetros decrecen; la evaluación fija un valor límite y olvida parte
de la genealogía. En la rama excepcional, las incidencias se prolongan hacia
Witt, Golay, Niemeier--Leech, \(V^\natural\), el Monstruo y moonshine.

El nuevo resultado añade un objeto común más fuerte que una coincidencia de
cardinales. El módulo de once modos, su sector positivo de rango tres y la
traza de la acción pasan de dodecafase a Witt mediante \(U_W\). El sector orientado
\(t=7\) se incorpora mediante \(W_{GK}\). La microfibra de
\(\pi\) aporta el proyector regional \(\Pi_\pi^{(468)}\), y la medida uniforme del
catálogo fija su peso.

La acción positiva se sitúa así en el centro de las dos lecturas:
\[
 \text{región}
 \quad\oplus\quad
 \text{polarización orientada}
 \quad\longmapsto\quad
 \text{defecto positivo}.
\]
La evaluación arquimediana y la evaluación excepcional siguen siendo
aplicaciones diferentes; el entrelazador identifica un sector común, no
colapsa sus codominios.

\section{Compatibilidad del operador angular con los cilindros del continuo discreto}

Una historia HMT determina cilindros anidados
\[
 C_1\supset C_2\supset\cdots
\]
con prefijos compatibles. Cuando el lector arquimediano satisface
\[
 \operatorname{diam}(C_n)\longrightarrow0,
\]
la completitud de \(\mathbb R\) fija a lo sumo un valor en la intersección.
La nueva acción añade una estructura antes de ese límite: distingue el coste
de la ruta, la clase de emisión y la multiplicidad de presentaciones que el
valor final ya no muestra.

Éste es el sentido matemático preciso de afirmar que el continuo se lee como
proyección de genealogías discretas. El sistema de cilindros y sus
evaluaciones se prolonga a toda escala; el teorema global de cobertura
demuestra
\(\operatorname{ev}_{\mathbb R}(\Omega_{\infty,K})=K\) para cada compacto
\(K\subset\mathbb R\) y, al dirigir los intervalos, que la imagen es toda
\(\mathbb R\). El teorema finito de este apartado controla la acción angular
en cada nivel y se integra en esa cobertura global; no la reabre ni la
condiciona. El descenso algebraico de cada operación conserva su verificación
propia.

\section{Compatibilidad con el corredor excepcional, Moonshine y las pantallas de teoría M}

El entrelazador \(U_W\) aterriza en el módulo de incidencia de \(W_{12}\).
Desde el código ternario, las dos continuaciones excepcionales permanecen
separadas:
\[
 C_W\longrightarrow W_{12}
 \xrightarrow[\text{relativo a }(D,\ell)]{}W_{24},
\]
y
\[
 C_W\xrightarrow{\ \phi\ } N(A_2^{12})
\xrightarrow{\ v_\alpha\ }\Lambda_{24}
\longrightarrow V_\Lambda
\longrightarrow V^\natural
\longrightarrow\mathbb M
\longrightarrow\text{moonshine}.
\]
La acción positiva aporta a esta rama un sector transportado y una línea de
orientación. Las flechas desde Leech hasta moonshine son las construcciones
clásicas aplicadas al objeto finito obtenido. La acción del Monstruo se
realiza sobre \(V^\natural\); las palabras HMT no se identifican aquí con un
módulo directo para esa acción.

Las pantallas \(12\to11\to10\), \(26\to24\) y la involución T-dual forman
una interfaz matemática exacta con las dimensiones y dualidades usadas en
teoría M. La identificación física completa pertenece a Mecánica
Dimensional y se tipa en la segunda parte mediante el mapa adicional hacia
\(g_{MN}\), \(C_{MNP}\), \(\psi_M\), acción y supersimetría. No se reserva
para otra obra ni se deduce de la acción del Monstruo.

\section{Factorización espectral y aislamiento de los sectores dominantes}

La línea pre-dimensional queda organizada en cuatro estratos.
\begin{enumerate}
  \item \textbf{Núcleo finito.} APP, TRIT, el emisor TPK, el catálogo
  \(104\,976\to468\to243\), las cinco regiones de \(\pi\), la
  dodecafase, \(K\), \(\alpha\), Witt y la rama excepcional poseen
  construcciones explícitas y certificados.
  \item \textbf{Puente operatorial.} La incidencia construye \(U_W\);
  el sector orientado \(t=7\) construye \(P_G\); \(B_K\) y \(u_K\) seleccionan
  la identificación isométrica orientada; los tres proyectores de rango tres
  quedan identificados relativamente a esa regla.
  \item \textbf{Extensión de acción.} \HTr fija la medida \(1/468\), \HAdd
  fija el carácter aditivo y \HDef declara el acoplamiento
  \(\Delta_4P_3\). \HRes transforma el defecto en capacidad y \HLoc lleva la
  capacidad a las aristas. Las cinco cláusulas producen la ley local.
  \item \textbf{Consecuencia espectral.} \Gtr produce
  \(T_{\rm band}\) y una fase angular pre-dimensional certificada, con
  diferencia relativa entre \(8.356\times10^{-10}\) y
  \(8.475\times10^{-10}\) respecto de \(C_{\rm reg}\).
\end{enumerate}

Esto permite un cierre fuerte y tipado de esta ruta: el emisor finito de la
sección coinductiva infinita, la incidencia, la identificación de
proyectores, la medida y la acción quedan cerrados
respecto de las primitivas declaradas, de la regla variacional
\(\mathcal E_B\) y de las cláusulas de acción \HTr--\HAdd--\HDef--\HRes--\HLoc. \HAdd
selecciona la forma lineal; \HDef selecciona su segunda coordenada; \HRes y \HLoc
especifican cómo se lee como coste. El teorema global de
cobertura de todo \(\mathbb R\), la inducción desde una transición total de
\(\mathcal X_{\mathrm{TPK}}^{\mathrm{enr}}\) y la equivalencia física con
teoría M se formulan como extensiones de los resultados anteriores.

\section{Compatibilidad entre las realizaciones matricial, angular y reticular}

Las representaciones construidas satisfacen simultáneamente las condiciones
de compatibilidad siguientes:
\begin{enumerate}
  \item cinco regiones de \(\pi\), con tres positivas, una negativa y una
  transversal;
  \item \(P_G,P_3,Q_3\) como proyectores positivos y \(R_{B_K}\) como
  involución de orientación;
  \item el sector homogéneo orientado \(t=7\), no el cociente heterogéneo de \(t=8\);
  \item la línea \(L_{\rm or}\) en la identificación isométrica triple;
  \item el carácter relativo de la regla variacional que selecciona
  \(H_*,c_*,r_*\);
  \item \(5/468\) sobre emisiones y \(7/729\) sobre semillas, sin
  intercambiarlos;
  \item la linealidad del carácter como consecuencia de \HAdd y el coeficiente
  \(\Delta_4\) como contenido separado de \HDef;
  \item \((y_{\rm band},C_{\rm band})\) como salida espectral certificada,
  separado de \((y_{\rm reg},C_{\rm reg})\), la realización regional
  canónica de la sección siguiente;
  \item \(\alpha\) como entrada de \(x\), de modo que la lectura
  del determinante es reversible y no independiente;
  \item las ramas \(W_{12}\to W_{24}\) y
  \(A_2^{12}\to\Lambda_{24}\) como construcciones separadas;
  \item el lector relativo \HTr--\HAdd--\HDef--\HRes--\HLoc distinguido de
  una transición total, cuya medida imagen dinámica constituye una
  construcción adicional.
\end{enumerate}

Estas condiciones no son notas editoriales. Forman parte de la estructura
lógica del resultado.
  \endgroup

\chapter{Equivalencia de las cartas global y analítica del par angular}
La medida positiva posterior al retorno se obtiene antes del contraángulo. La
rama determinantal determina la valoración de acción y ambas construcciones
confluyen en la sección de la recta dimensional. Sólo entonces se forman el
ángulo, su complemento orientado y los caracteres \(q_{+}\) y \(q_{-}\), que
evalúan las incidencias ya generadas por el núcleo.
% Adaptador único del capítulo 38 del sucesor 102.
% Sustituye exclusivamente la jerarquía de encabezados. Los cuerpos de los
% dos fuentes teoremáticas científicos y del delta Ley 9 se conservan línea a línea,
% reordenados conforme al DAG causal aprobado. Véase MAPA_SEGMENTOS_CAPITULO_38.tsv.

% HMT_SOURCE_SEGMENT_BEGIN id=B00 source=colaboracion/parte_iii/source/public_final/base_83/c30_body.tex body_lines=1-2
\label{chap:p3-83-30-angulo_contraangulo}

% HMT_SOURCE_SEGMENT_END id=B00
% HMT_SOURCE_SEGMENT_BEGIN id=A00 source=colaboracion/parte_iii/source/public_final/ascensos_20260826/16f_realizacion_analitica_contraangulo.tex body_lines=1-4
\label{p3-20260826:chap:realizacion-analitica-contraangulo}
\index{contraángulo}
\index{microfibra de pi@microfibra de \(\pi\)!carácter analítico}
\index{carácter determinantal}
% HMT_SOURCE_SEGMENT_END id=A00
% HMT_SOURCE_SEGMENT_BEGIN id=D00 source=manuscrito/sucesor_102/deltas_ley9/c38_par_angular_cartas_y_canales.tex body_lines=1
% Distribución causal exacta del delta Ley 9; contenido conservado sin síntesis.
% HMT_SOURCE_SEGMENT_END id=D00

\section{Fase angular intrínseca \texorpdfstring{\(\alpha\)}{alpha HMT}, elevación nonádica \texorpdfstring{\(10^3\)}{10^3} y precursor \texorpdfstring{\(x_A\)}{x_A}}
\label{sec:s102:c38:1}

\subsection{Construcción del precursor analítico \texorpdfstring{\(x_A\)}{x_A}, de la contracción \texorpdfstring{\(D_A\)}{D_A} y de la corrección regional \texorpdfstring{\(\mathscr C_\pi(\alpha)\)}{C_pi(alpha)}}
\label{sec:s102:c38:1:1}

% HMT_SOURCE_SEGMENT_BEGIN id=B01 source=colaboracion/parte_iii/source/public_final/base_83/c30_body.tex body_lines=4-32
\label{p3-83-c30-s1:chap:realizacion-analitica-contraangulo}
\index{coordenada angular conjugada \(C^*\)}
\index{microfibra de pi@microfibra de \(\pi\)!carácter analítico}
\index{carácter determinantal}

El catálogo finito y una evaluación analítica responden a preguntas de tipos
distintos. El catálogo determina qué regiones existen, cuáles comparten un
retorno y qué información conserva cada una. Una evaluación analítica debe
decidir, además, cómo se transforma la longitud de una palabra en grado, cómo
se prolonga una rama después de su retorno y con qué carácter escalar se lee
esa prolongación. Este capítulo construye esa segunda operación de manera
explícita.

La construcción parte de cuatro datos ya obtenidos en los capítulos
anteriores: la microfibra de cinco regiones de \(\pi\), la longitud seis de
cada palabra regional, el cierre dodecafásico de \(\alpha\) y el transporte
bicapa asociado a la coordenada común
\(A=10^3\alpha\). Ningún valor de la coordenada angular conjugada \(C_*^{\rm HMT}\), de la componente reducida de
acción, del funcional de Barbero--Immirzi o de una constante metrológica se
emplea en la evaluación. El resultado es un carácter regional convergente,
una recurrencia exacta y una fase orientada cuya carta sexagesimal reproduce
la coordenada angular conjugada publicada.

La separación entre la parte finita y la parte analítica es esencial. El
entero \(169\) será un teorema del catálogo. La serie que lo prolonga será el
teorema de un lector analítico cuyas reglas se declaran antes de evaluarlo.
Así, las premisas adicionales no quedan ocultas dentro de una coincidencia
decimal.

% HMT_SOURCE_SEGMENT_END id=B01
\section{Construcción preangular de \texorpdfstring{\(\eta_{\mathrm{ret}}\)}{eta_ret} mediante \texorpdfstring{\(H_5\)}{H5}, \texorpdfstring{\(D_A\)}{D_A} y \texorpdfstring{\(\mathscr C_\pi(\alpha)\)}{C_pi(alpha)}}
\label{sec:s102:c38:2}

\subsection{Derivación independiente de \texorpdfstring{\(\eta_{\mathrm{ret}}=H_5-(90/\pi)\mathscr C_\pi(\alpha)\)}{eta_ret=H5-(90/pi)C_pi(alpha)} mediante el lector analítico regional}
\label{sec:s102:c38:2:1}

% HMT_SOURCE_SEGMENT_BEGIN id=A01 source=colaboracion/parte_iii/source/public_final/ascensos_20260826/16f_realizacion_analitica_contraangulo.tex body_lines=5-38,40-47

El catálogo finito y una evaluación analítica responden a preguntas de tipos
distintos. El catálogo determina qué regiones existen, cuáles comparten un
retorno y qué información conserva cada una. Una evaluación analítica debe
decidir, además, cómo se transforma la longitud de una palabra en grado, cómo
se prolonga una rama después de su retorno y con qué carácter escalar se lee
esa prolongación. Esta sección construye esa segunda operación de manera
explícita.

La construcción parte de la microfibra de cinco regiones de \(\pi\), la
longitud seis de cada palabra regional, el cierre dodecafásico de \(\alpha\)
y la inserción angular declarada
\[
 \iota_\angle(x):=10^3x\,u_\circ,\qquad
 \Adeg:=10^3\alpha .
\]
Aquí \(u_\circ\) es la base sexagesimal cuya vuelta completa tiene
coordenada \(360\).
El factor \(10^3\) pertenece a esa normalización de carta y no se deduce del
cierre aritmético de \(\alpha\). El transporte bicapa se construye después
a partir de \(\Adeg\). La función que evalúa el lector no recibe como
argumento el contraángulo, la coordenada posterior al retorno, el funcional
hexada--octada ni una constante metrológica. El resultado es un carácter
regional convergente, una recurrencia exacta y una fase orientada en la carta
sexagesimal declarada.

La separación entre la parte finita y la parte analítica es esencial. El
entero \(169\) es un teorema del catálogo. La serie que lo prolonga pertenece
a un lector analítico cuyas reglas se declaran antes de evaluarlo.

La fase se denota por \(y_*\) y, cuando se enfatiza su origen regional, por
\(y_{\rm reg}\); ambos símbolos designan el mismo lector construido a
continuación.


La evaluación recibe el catálogo de \(468\) emisiones, el valor exacto de
\(\alpha\), \(\varphi=(1+\sqrt5)/2\), la fórmula fijada de \(H_5\),
la carta angular y las cinco reglas del lector regional. Con esos datos
calcula \(169\), \(D_A\), la recurrencia y \(C^*\). El teorema es exacto en
esta clase declarada de lectores; no usa \(C^*\) como entrada ni afirma que
las reglas de \(H_5\) sean la única clase natural posible.

% HMT_SOURCE_SEGMENT_END id=A01
\subsection{Persistencia del retorno en las \texorpdfstring{\(5\)}{5} regiones de la microfibra}
\label{sec:s102:c38:2:2}

% HMT_SOURCE_SEGMENT_BEGIN id=B02 source=colaboracion/parte_iii/source/public_final/base_83/c30_body.tex body_lines=34-155

Sea
\[
\mathcal F_\pi=\Psi^{-1}(010211)\subset\mathcal U_6
\]
la microfibra relativa al catálogo. Recordemos sus cinco elementos, ahora
separados en bloque inicial y bloque terminal:
\[
\begin{aligned}
 &(501,614,498\mid169,272,272),\\
 &(810,923,870\mid169,272,272),\\
 &(810,983,810\mid169,272,272),\\
 &(870,923,810\mid169,272,272),\\
 &(870,923,810\mid418,674,521).
\end{aligned}
\]
Definimos la proyección de retorno
\[
p_{\rm ret}:\mathbb Z^6\longrightarrow\mathbb Z^3,
\qquad
p_{\rm ret}(u_1,\ldots,u_6)=(u_4,u_5,u_6).
\]

\begin{proposition}[Bloque terminal persistente]
\label{p3-83-c30-s1:prop:bloque-terminal-persistente}
El multiconjunto \(p_{\rm ret}(\mathcal F_\pi)\) posee un único elemento de
multiplicidad máxima:
\[
b_\pi=(169,272,272),
\qquad
\operatorname{mult}_{\mathcal F_\pi}(b_\pi)=4.
\]
El único bloque terminal restante es
\[
b_\pi^-=(418,674,521)
\]
y tiene multiplicidad uno.
\end{proposition}

\begin{proof}
La afirmación se obtiene por enumeración exacta de las cinco filas del
catálogo que satisfacen \(\Psi(U)=010211\). Cuatro filas poseen bloque
terminal \((169,272,272)\); la quinta, que porta el retorno mutado, posee
\((418,674,521)\). La multiplicidad máxima es cuatro y se alcanza una sola
vez.
\end{proof}

La diferencia orientada entre el retorno mutado y el persistente es
\[
 \omega_\pi=b_\pi^- - b_\pi=(249,402,249),
 \qquad
 \langle(1,1,1),\omega_\pi\rangle=900.
\]
Esta identidad separa dos objetos que no deben confundirse. El entero
\(169\) es una coordenada de un dato terminal, es decir, un dato de grado
cero. El vector \(\omega_\pi\) puede leerse como una cocadena de grado uno
después de orientar la arista que une ambos bloques. Llamarlo cociclo exigiría,
además, fijar un complejo de cocadenas y verificar que su coborde se anula.
La prolongación analítica construida más adelante no presupone esa
identificación.

La selección de \(169\) no exige privilegiar la primera posición del bloque.
El grupo simétrico \(\mathfrak S_3\) actúa permutando sus tres coordenadas.
El estabilizador de \(b_\pi\) intercambia las dos apariciones de \(272\) y
fija la única coordenada de multiplicidad uno. Denotemos por
\(\operatorname{sing}(b)\) esa coordenada cuando el bloque tiene tipo
\((r,s,s)\), con \(r\ne s\). Equivalentemente,
\[
\operatorname{sing}(b)
=\sum_{j=1}^3b_j-2\operatorname{moda}(b).
\]
En particular,
\[
\boxed{\operatorname{sing}(b_\pi)=169.}
\]

\begin{definition}[Selector residual de la microfibra]
\label{p3-83-c30-s1:def:selector-residual-pi}
El selector \(\rho_\pi\) aplica sucesivamente dos operaciones:
\begin{enumerate}
  \item selecciona el bloque terminal de multiplicidad máxima dentro de
  \(p_{\rm ret}(\mathcal F_\pi)\);
  \item retiene la coordenada fija por el estabilizador de ese bloque.
\end{enumerate}
Por la \cref{p3-83-c30-s1:prop:bloque-terminal-persistente},
\[
\boxed{\rho_\pi=169.}
\]
\end{definition}

\begin{theorem}[Unicidad equivariante del retorno]
\label{p3-83-c30-s1:thm:unicidad-retorno-169}
Entre los selectores que
\begin{enumerate}
  \item son invariantes bajo permutaciones de las cinco regiones;
  \item son equivariantes bajo permutaciones de las tres coordenadas
  terminales;
  \item seleccionan el bloque de multiplicidad máxima;
  \item retienen una coordenada fija por su estabilizador,
\end{enumerate}
el valor de retorno está determinado de manera única y es \(169\).
\end{theorem}

\begin{proof}
La invariancia bajo \(\mathfrak S_5\) impide usar el orden de las filas. La
tercera condición selecciona \(b_\pi\), que es el único modo por la
\cref{p3-83-c30-s1:prop:bloque-terminal-persistente}. Su estabilizador en
\(\mathfrak S_3\) posee dos órbitas sobre las coordenadas: la pareja de
valores \(272\) y el valor unipuntual \(169\). La cuarta condición selecciona la
segunda órbita. La equivarianza conserva su valor al desplazar su posición.
\end{proof}

Este argumento mejora la lectura posicional del retorno. El \(169\) no es
«la primera coordenada» de una fila escogida: es un invariante del
multiconjunto regional y de la acción de sus simetrías de reordenación.
Tampoco se lo selecciona por rareza global. En las \(468\) emisiones del
catálogo, \(169\) aparece \(84\) veces, uniformemente distribuido con
\(14\) apariciones en cada una de las seis posiciones. La unicidad del
\cref{p3-83-c30-s1:thm:unicidad-retorno-169} es, por tanto, relacional y regional: reside
en la combinación de microfibra, persistencia modal y estabilizador, no en la
frecuencia aislada del entero.

% HMT_SOURCE_SEGMENT_END id=B02
\subsection{Construcción del lector analítico regional y de su dominio de convergencia}
\label{sec:s102:c38:2:3}

% HMT_SOURCE_SEGMENT_BEGIN id=B05 source=colaboracion/parte_iii/source/public_final/base_83/c30_body.tex body_lines=261-301

La extensión analítica se fija mediante las reglas siguientes.

\begin{definition}[Lector regional determinantal]
\label{p3-83-c30-s1:def:lector-regional-determinantal}
El lector regional determinantal satisface:
\begin{enumerate}
  \item \emph{compatibilidad de grado}: la longitud cilíndrica es el grado
  analítico;
  \item \emph{descomposición de renovación}: el impulso finito de retorno y
  la continuación estricta pertenecen a sumandos aditivos distintos;
  \item \emph{normalización de continuación}: después de registrar una sola
  vez el residuo \(\rho_\pi\), la rama superviviente comienza en la unidad de
  la línea determinante;
  \item \emph{covariancia por concatenación}: cada prolongación elemental
  actúa mediante \(\bigwedge^2\mathcal T=D_A\);
  \item \emph{orientación}: en la convención cardinal fijada para APP, el
  retorno regional pertenece al sector compensador y se sustrae de la fase
  de acción de quinto grado.
\end{enumerate}
\end{definition}

La tercera regla tiene contenido matemático. El valor \(169\) se registra
una vez como amplitud del retorno; no se multiplica de nuevo en cada
prolongación. La continuación estricta cuenta la única línea que prosigue y
la normaliza por su unidad. Esta es la diferencia entre una recurrencia de
renovación y la iteración homogénea del impulso inicial.

La necesidad de declarar esa normalización puede probarse. Para todo
\(\lambda\in\mathbb R\), la familia
\[
F_\lambda(z)
=169z^6+\lambda\frac{D_Az^7}{1-D_Az}
\]
conserva el mismo retorno finito y la misma razón determinantal entre
coeficientes consecutivos. El catálogo y la razón por sí solos no distinguen
\(\lambda\). La unidad de la línea determinante fija
\(\lambda=1\) antes de evaluar \(z=\alpha\). De este modo, la normalización
no se obtiene ajustando el valor final; forma parte de la definición del
carácter.

% HMT_SOURCE_SEGMENT_END id=B05
\subsection{Resolución analítica de la recurrencia de retorno}
\label{sec:s102:c38:2:4}

% HMT_SOURCE_SEGMENT_BEGIN id=B06 source=colaboracion/parte_iii/source/public_final/base_83/c30_body.tex body_lines=303-364

Definimos los coeficientes \(\kappa_n\) por
\[
\kappa_6=169,
\qquad
\kappa_7=D_A,
\qquad
\kappa_{n+1}=D_A\kappa_n\quad(n\ge7).
\]
La primera igualdad es el impulso regional; las dos siguientes describen la
continuación normalizada.

\begin{theorem}[Realización analítica graduada]
\label{p3-83-c30-s1:thm:realizacion-analitica-graduada}
El germen analítico compatible con las reglas de la
\cref{p3-83-c30-s1:def:lector-regional-determinantal} es único y vale
\[
\boxed{
\mathscr C_\pi(z)
=169z^6+\frac{D_Az^7}{1-D_Az}.}
\]
Sus coeficientes satisfacen
\[
\kappa_n=D_A^{n-6}\qquad(n\ge7).
\]
\end{theorem}

\begin{proof}
Escribamos
\[
F(z)=169z^6+\sum_{k\ge1}c_kz^{6+k}.
\]
La normalización de continuación y la covariancia determinantal dan
\[
c_1=D_A,
\qquad
c_{k+1}=D_Ac_k.
\]
Por inducción, \(c_k=D_A^k\). La suma geométrica produce
\[
\sum_{k\ge1}D_A^kz^{6+k}
=\frac{D_Az^7}{1-D_Az}.
\]
No queda ningún coeficiente libre.
\end{proof}

Como \(0<\alpha<1\) y \(0<D_A<1\),
\[
0<D_A\alpha<1.
\]
En consecuencia, \(\mathscr C_\pi(\alpha)\) converge absolutamente y su
radio de convergencia es \(D_A^{-1}>1\). Numéricamente,
\[
D_A=0.7751291192471564301888840964802\ldots,
\]
\[
D_A\alpha=0.00565639046986492673885079095469\ldots.
\]
La rapidez de esta contracción explica por qué la suma completa y su
truncamiento de grado siete poseen las mismas quince primeras cifras en la
carta final.

% HMT_SOURCE_SEGMENT_END id=B06
\subsection{Persistencia del retorno en las \texorpdfstring{\(5\)}{5} regiones y compatibilidad de la componente \texorpdfstring{\(\eta_{\mathrm{ret}}\)}{eta_ret}}
\label{sec:s102:c38:2:5}

% HMT_SOURCE_SEGMENT_BEGIN id=A02 source=colaboracion/parte_iii/source/public_final/ascensos_20260826/16f_realizacion_analitica_contraangulo.tex body_lines=49-170

Sea
\[
\mathcal F_\pi=\Psi^{-1}(010211)\subset\mathcal U_6
\]
la microfibra relativa al catálogo. Recordemos sus cinco elementos, ahora
separados en bloque inicial y bloque terminal:
\[
\begin{aligned}
 &(501,614,498\mid169,272,272),\\
 &(810,923,870\mid169,272,272),\\
 &(810,983,810\mid169,272,272),\\
 &(870,923,810\mid169,272,272),\\
 &(870,923,810\mid418,674,521).
\end{aligned}
\]
Definimos la proyección de retorno
\[
p_{\rm ret}:\mathbb Z^6\longrightarrow\mathbb Z^3,
\qquad
p_{\rm ret}(u_1,\ldots,u_6)=(u_4,u_5,u_6).
\]

\begin{proposition}[Bloque terminal persistente]
\label{p3-20260826:prop:bloque-terminal-persistente}
El multiconjunto \(p_{\rm ret}(\mathcal F_\pi)\) posee un único elemento de
multiplicidad máxima:
\[
b_\pi=(169,272,272),
\qquad
\operatorname{mult}_{\mathcal F_\pi}(b_\pi)=4.
\]
El único bloque terminal restante es
\[
b_\pi^-=(418,674,521)
\]
y tiene multiplicidad uno.
\end{proposition}

\begin{proof}
La afirmación se obtiene por enumeración exacta de las cinco filas del
catálogo que satisfacen \(\Psi(U)=010211\). Cuatro filas poseen bloque
terminal \((169,272,272)\); la quinta, que porta el retorno mutado, posee
\((418,674,521)\). La multiplicidad máxima es cuatro y se alcanza una sola
vez.
\end{proof}

La diferencia orientada entre el retorno mutado y el persistente es
\[
 \omega_\pi=b_\pi^- - b_\pi=(249,402,249),
 \qquad
 \langle(1,1,1),\omega_\pi\rangle=900.
\]
Esta identidad separa dos objetos que no deben confundirse. El entero
\(169\) es una coordenada de un dato terminal, es decir, un dato de grado
cero. El vector \(\omega_\pi\) puede leerse como una cocadena de grado uno
después de orientar la arista que une ambos bloques. Llamarlo cociclo exigiría,
además, fijar un complejo de cocadenas y verificar que su coborde se anula.
La prolongación analítica construida más adelante no presupone esa
identificación.

La selección de \(169\) no exige privilegiar la primera posición del bloque.
El grupo simétrico \(\mathfrak S_3\) actúa permutando sus tres coordenadas.
El estabilizador de \(b_\pi\) intercambia las dos apariciones de \(272\) y
fija la única coordenada de multiplicidad uno. Denotemos por
\(\operatorname{sing}(b)\) esa coordenada cuando el bloque tiene tipo
\((r,s,s)\), con \(r\ne s\). Equivalentemente,
\[
\operatorname{sing}(b)
=\sum_{j=1}^3b_j-2\operatorname{moda}(b).
\]
En particular,
\[
\boxed{\operatorname{sing}(b_\pi)=169.}
\]

\begin{definition}[Selector residual de la microfibra]
\label{p3-20260826:def:selector-residual-pi}
El selector \(\rho_\pi\) aplica sucesivamente dos operaciones:
\begin{enumerate}
  \item selecciona el bloque terminal de multiplicidad máxima dentro de
  \(p_{\rm ret}(\mathcal F_\pi)\);
  \item retiene la coordenada fija por el estabilizador de ese bloque.
\end{enumerate}
Por la \cref{p3-20260826:prop:bloque-terminal-persistente},
\[
\boxed{\rho_\pi=169.}
\]
\end{definition}

\begin{theorem}[Unicidad equivariante del retorno]
\label{p3-20260826:thm:unicidad-retorno-169}
Entre los selectores que
\begin{enumerate}
  \item son invariantes bajo permutaciones de las cinco regiones;
  \item son equivariantes bajo permutaciones de las tres coordenadas
  terminales;
  \item seleccionan el bloque de multiplicidad máxima;
  \item retienen una coordenada fija por su estabilizador,
\end{enumerate}
el valor de retorno está determinado de manera única y es \(169\).
\end{theorem}

\begin{proof}
La invariancia bajo \(\mathfrak S_5\) impide usar el orden de las filas. La
tercera condición selecciona \(b_\pi\), que es el único modo por la
\cref{p3-20260826:prop:bloque-terminal-persistente}. Su estabilizador en
\(\mathfrak S_3\) posee dos órbitas sobre las coordenadas: la pareja de
valores \(272\) y el valor unipuntual \(169\). La cuarta condición selecciona la
segunda órbita. La equivarianza conserva su valor al desplazar su posición.
\end{proof}

Este argumento mejora la lectura posicional del retorno. El \(169\) no es
«la primera coordenada» de una fila escogida: es un invariante del
multiconjunto regional y de la acción de sus simetrías de reordenación.
Tampoco se lo selecciona por rareza global. En las \(468\) emisiones del
catálogo, \(169\) aparece \(84\) veces, uniformemente distribuido con
\(14\) apariciones en cada una de las seis posiciones. La unicidad del
\cref{p3-20260826:thm:unicidad-retorno-169} es, por tanto, relacional y regional: reside
en la combinación de microfibra, persistencia modal y estabilizador, no en la
frecuencia aislada del entero.

% HMT_SOURCE_SEGMENT_END id=A02
\subsection{Aplicación del lector regional al par angular conjugado}
\label{sec:s102:c38:2:6}

% HMT_SOURCE_SEGMENT_BEGIN id=A05 source=colaboracion/parte_iii/source/public_final/ascensos_20260826/16f_realizacion_analitica_contraangulo.tex body_lines=277-317

La extensión analítica se fija mediante las reglas siguientes.

\begin{definition}[Lector regional determinantal]
\label{p3-20260826:def:lector-regional-determinantal}
El lector regional determinantal satisface:
\begin{enumerate}
  \item \emph{compatibilidad de grado}: la longitud cilíndrica es el grado
  analítico;
  \item \emph{descomposición de renovación}: el impulso finito de retorno y
  la continuación estricta pertenecen a sumandos aditivos distintos;
  \item \emph{normalización de continuación}: después de registrar una sola
  vez el residuo \(\rho_\pi\), la rama superviviente comienza en la unidad de
  la línea determinante;
  \item \emph{covariancia por concatenación}: cada prolongación elemental
  actúa mediante \(\bigwedge^2\mathcal T=D_A\);
  \item \emph{orientación}: en la convención cardinal fijada para APP, el
  retorno regional pertenece al sector compensador y se sustrae de la fase
  de acción de quinto grado.
\end{enumerate}
\end{definition}

La tercera regla tiene contenido matemático. El valor \(169\) se registra
una vez como amplitud del retorno; no se multiplica de nuevo en cada
prolongación. La continuación estricta cuenta la única línea que prosigue y
la normaliza por su unidad. Esta es la diferencia entre una recurrencia de
renovación y la iteración homogénea del impulso inicial.

La necesidad de declarar esa normalización puede probarse. Para todo
\(\lambda\in\mathbb R\), la familia
\[
F_\lambda(z)
=169z^6+\lambda\frac{D_Az^7}{1-D_Az}
\]
conserva el mismo retorno finito y la misma razón determinantal entre
coeficientes consecutivos. El catálogo y la razón por sí solos no distinguen
\(\lambda\). La unidad de la línea determinante fija
\(\lambda=1\) antes de evaluar \(z=\alpha\). De este modo, la normalización
no se obtiene ajustando el valor final; forma parte de la definición del
carácter.

% HMT_SOURCE_SEGMENT_END id=A05
\subsection{Suma cerrada y publicación de la componente posterior al retorno}
\label{sec:s102:c38:2:7}

% HMT_SOURCE_SEGMENT_BEGIN id=A06 source=colaboracion/parte_iii/source/public_final/ascensos_20260826/16f_realizacion_analitica_contraangulo.tex body_lines=319-380

Definimos los coeficientes \(\kappa_n\) por
\[
\kappa_6=169,
\qquad
\kappa_7=D_A,
\qquad
\kappa_{n+1}=D_A\kappa_n\quad(n\ge7).
\]
La primera igualdad es el impulso regional; las dos siguientes describen la
continuación normalizada.

\begin{theorem}[Realización analítica graduada]
\label{p3-20260826:thm:realizacion-analitica-graduada}
El germen analítico compatible con las reglas de la
\cref{p3-20260826:def:lector-regional-determinantal} es único y vale
\[
\boxed{
\mathscr C_\pi(z)
=169z^6+\frac{D_Az^7}{1-D_Az}.}
\]
Sus coeficientes satisfacen
\[
\kappa_n=D_A^{n-6}\qquad(n\ge7).
\]
\end{theorem}

\begin{proof}
Escribamos
\[
F(z)=169z^6+\sum_{k\ge1}c_kz^{6+k}.
\]
La normalización de continuación y la covariancia determinantal dan
\[
c_1=D_A,
\qquad
c_{k+1}=D_Ac_k.
\]
Por inducción, \(c_k=D_A^k\). La suma geométrica produce
\[
\sum_{k\ge1}D_A^kz^{6+k}
=\frac{D_Az^7}{1-D_Az}.
\]
No queda ningún coeficiente libre.
\end{proof}

Como \(0<\alpha<1\) y \(0<D_A<1\),
\[
0<D_A\alpha<1.
\]
En consecuencia, \(\mathscr C_\pi(\alpha)\) converge absolutamente y su
radio de convergencia es \(D_A^{-1}>1\). Numéricamente,
\[
D_A=0.7751291192471564301888840964802\ldots,
\]
\[
D_A\alpha=0.00565639046986492673885079095469\ldots.
\]
La rapidez de esta contracción explica por qué la suma completa y su
truncamiento de grado siete poseen las mismas quince primeras cifras en la
carta final.

% HMT_SOURCE_SEGMENT_END id=A06
\section{Confluencia de \texorpdfstring{\(\eta_{\mathrm{ret}}\)}{eta_ret} con la valoración \texorpdfstring{\(\varepsilon_H=-34\)}{epsilon_H=-34} en la magnitud \texorpdfstring{\(\hbar_{\mathrm{ret}}\)}{hbar_ret}}
\label{sec:s102:c38:3}

\subsection{Derivación del carácter determinantal del transporte parangular}
\label{sec:s102:c38:3:1}

% HMT_SOURCE_SEGMENT_BEGIN id=B04 source=colaboracion/parte_iii/source/public_final/base_83/c30_body.tex body_lines=196-259

El cierre dodecafásico proporciona \(\alpha\). La completación de las tres
parejas nonádicas proporciona la escala \(10^3\), de modo que
\[
A=10^3\alpha.
\]
Elegida la carta angular arquimediana, la coordenada común en radianes es
\[
x=\frac{\pi A}{180}=\frac{50\pi}{9}\alpha.
\]
Sea \(R\) una involución real de traza nula en dimensión dos,
\[
R^2=I_2,
\qquad
\operatorname{tr}R=0,
\]
y consideremos el transporte formal bicapa
\[
\mathcal T(x,y)=\exp(-xI_2-yR).
\]
La variable orientada \(y\) todavía no se ha evaluado. Sin embargo, la acción
de \(\mathcal T\) sobre la línea determinante es independiente de ella:
\[
\begin{aligned}
\det\mathcal T(x,y)
&=\exp\!\left(\operatorname{tr}(-xI_2-yR)\right)\\
&=e^{-2x}.
\end{aligned}
\]
Definimos
\[
\boxed{
D_A=e^{-2x}=e^{-100\pi\alpha/9}.}
\]
El subíndice recuerda que se trata del determinante de la contracción común
asociada a \(A\), no de una función de \(C_*^{\rm HMT}\).

\begin{proposition}[Carácter de la línea determinante]
\label{p3-83-c30-s1:prop:caracter-determinante}
La aplicación
\[
\delta_A:(\mathbb N_0,+)\longrightarrow(\mathbb R_{>0},\cdot),
\qquad
\delta_A(k)=D_A^k,
\]
es el único carácter multiplicativo cuyo valor sobre una prolongación
elemental es \(D_A\).
\end{proposition}

\begin{proof}
La multiplicatividad da
\[
\delta_A(k)
=\delta_A(\underbrace{1+\cdots+1}_{k\text{ veces}})
=\delta_A(1)^k=D_A^k.
\]
La misma identidad prueba existencia y unicidad.
\end{proof}

El determinante permanece invariante bajo conjugación de
\(\mathcal T\) y bajo el intercambio de hojas \(R\mapsto-R\). Por tanto, la
cola construida a partir de \(D_A\) pertenece al modo común. La orientación
aparecerá en el signo con el que el carácter regional corrige la fase.

% HMT_SOURCE_SEGMENT_END id=B04
\subsection{Carácter de acción de grado \texorpdfstring{\(5\)}{5} en la construcción regional}
\label{sec:s102:c38:3:2}

% HMT_SOURCE_SEGMENT_BEGIN id=B07 source=colaboracion/parte_iii/source/public_final/base_83/c30_body.tex body_lines=366-405

El bloque central
\[
\mathcal B_c=\begin{pmatrix}7&2\\2&7\end{pmatrix}
\]
posee autovalores \(9\) y \(5\). El ciclo dodecafásico aporta el rango
\(12\); sus órbitas de paso tres poseen longitud cuatro; y el censo completo
aporta el cociente exacto
\[
\frac{104976}{1944}=54.
\]
A partir de estos invariantes se fija el carácter de acción de quinto grado
\[
\boxed{
H_5(\alpha,\varphi)
=\frac12\sqrt{\frac{1000\alpha}{\varphi}}-\alpha
+\frac9{16}\alpha^2-\frac59\alpha^3
+\frac7{48}\alpha^4-\frac1{54}\alpha^5.}
\]
Los cuatro coeficientes racionales pueden escribirse como
\[
\frac9{16}=\frac{\lambda_+}{4^2},
\qquad
\frac59=\frac{\lambda_-}{\lambda_+},
\qquad
\frac7{48}=\frac{\tfrac12\operatorname{tr}\mathcal B_c}{4\cdot12},
\qquad
\frac1{54}=\frac{1944}{104976},
\]
con \((\lambda_+,\lambda_-)=(9,5)\). Estas identidades certifican la
procedencia de los números utilizados. La asignación sucesiva de esos
invariantes a los grados dos, tres, cuatro y cinco forma parte del carácter
analítico; no se deduce de la mera existencia del catálogo finito.

Esta precisión elimina una ambigüedad habitual. Una combinación de
invariantes de la construcción es una definición interna y reproducible; su
naturalidad requiere las reglas del lector que especifican orden, signos y
normalización. Aquí esas reglas son visibles y quedan sometidas a los mismos
controles negativos que el resto de la construcción.

% HMT_SOURCE_SEGMENT_END id=B07
\subsection{Descomposición exacta de la recta de acción en secciones pre-retorno y post-retorno}
\label{sec:s102:c38:3:3}

% HMT_SOURCE_SEGMENT_BEGIN id=B12 source=colaboracion/parte_iii/source/public_final/base_83/c30_body.tex body_lines=571-617

El carácter de quinto grado y la componente de acción posterior al retorno
son dos valores distintos:
\[
H_5
=1.05457181764615446962582182943306\ldots,
\]
\[
\bar h_{\rm HMT}^{\rm ret}
=1.05457181691503906113369058472430\ldots.
\]
El primero es la acción local antes de incorporar el carácter regional; el
segundo es la acción orientada que queda después del retorno. La corrección
que genera la coordenada angular conjugada impide identificarlos.

La comparación metrológica debe respetar esta distinción. El primer valor da
la mantisa
\[
2\pi H_5
=6.62607014999998792600111034989947\ldots,
\]
mientras que la acción posterior da
\[
2\pi\bar h_{\rm HMT}^{\rm ret}
=6.62607014540625433351074984759288\ldots.
\]
En el Sistema Internacional, el valor
\[
h=6.62607015\times10^{-34}\;\mathrm{J\,s}
\]
es exacto por definición \cite{BIPMSI}. Por ello,
\[
2\pi H_5-6.62607015
=-1.2073998889650101\ldots\times10^{-14},
\]
y no es una igualdad exacta. La proximidad del valor anterior al retorno es
un control numérico notable; no autoriza a sustituir la constante definitoria
del SI por la componente posterior ni a confundir las dos acciones HMT\@.

Este resultado determina con precisión la situación matemática. El lector
regional genera \(C_*^{\rm HMT}\) y, desde él, el funcional
\(\Gamma_{6\to8}\). El mismo cálculo separa la aproximación de Planck
producida por \(H_5\) de la acción reducida que interviene en el
la coordenada angular conjugada. La separación no debilita la construcción: elimina una
identificación incompatible con las cifras exactas y convierte cada
comparación en una afirmación falsable de tipo bien definido.

% HMT_SOURCE_SEGMENT_END id=B12
\subsection{Factorización del carácter determinantal en las \texorpdfstring{\(2\)}{2} cartas de fase}
\label{sec:s102:c38:3:4}

% HMT_SOURCE_SEGMENT_BEGIN id=A04 source=colaboracion/parte_iii/source/public_final/ascensos_20260826/16f_realizacion_analitica_contraangulo.tex body_lines=211-275

El cierre dodecafásico proporciona \(\alpha\). La inserción angular declarada
emplea la escala \(10^3\), compatible con la agrupación de tres parejas
nonádicas, de modo que
\[
\Adeg=10^3\alpha\,{}^\circ.
\]
Elegida la carta angular arquimediana, la coordenada común en radianes es
\[
x=: \Arad=\frac{\pi\Adeg}{180}=\frac{50\pi}{9}\alpha.
\]
Sea \(R\) una involución real de traza nula en dimensión dos,
\[
R^2=I_2,
\qquad
\operatorname{tr}R=0,
\]
y consideremos el transporte formal bicapa
\[
\mathcal T(x,y)=\exp(-xI_2-yR).
\]
La variable orientada \(y\) todavía no se ha evaluado. Sin embargo, la acción
de \(\mathcal T\) sobre la línea determinante es independiente de ella:
\[
\begin{aligned}
\det\mathcal T(x,y)
&=\exp\!\left(\operatorname{tr}(-xI_2-yR)\right)\\
&=e^{-2x}.
\end{aligned}
\]
Definimos
\[
\boxed{
D_A=e^{-2x}=e^{-100\pi\alpha/9}.}
\]
El subíndice recuerda que se trata del determinante de la contracción común
asociada a \(A\), no de una función del contraángulo.

\begin{proposition}[Carácter de la línea determinante]
\label{p3-20260826:prop:caracter-determinante}
La aplicación
\[
\delta_A:(\mathbb N_0,+)\longrightarrow(\mathbb R_{>0},\cdot),
\qquad
\delta_A(k)=D_A^k,
\]
es el único carácter multiplicativo cuyo valor sobre una prolongación
elemental es \(D_A\).
\end{proposition}

\begin{proof}
La multiplicatividad da
\[
\delta_A(k)
=\delta_A(\underbrace{1+\cdots+1}_{k\text{ veces}})
=\delta_A(1)^k=D_A^k.
\]
La misma identidad prueba existencia y unicidad.
\end{proof}

El determinante permanece invariante bajo conjugación de
\(\mathcal T\) y bajo el intercambio de hojas \(R\mapsto-R\). Por tanto, la
cola construida a partir de \(D_A\) pertenece al modo común. La orientación
aparecerá en el signo con el que el carácter regional corrige la fase.

% HMT_SOURCE_SEGMENT_END id=A04
\subsection{Publicación del carácter de acción de grado \texorpdfstring{\(5\)}{5} en la carta parangular}
\label{sec:s102:c38:3:5}

% HMT_SOURCE_SEGMENT_BEGIN id=A07 source=colaboracion/parte_iii/source/public_final/ascensos_20260826/16f_realizacion_analitica_contraangulo.tex body_lines=382-418

El bloque central
\[
\mathcal B_c=\begin{pmatrix}7&2\\2&7\end{pmatrix}
\]
posee autovalores \(9\) y \(5\). El ciclo dodecafásico aporta el rango
\(12\); sus órbitas de paso tres poseen longitud cuatro; y el censo completo
aporta el cociente exacto
\[
\frac{104976}{1944}=54.
\]
El lector vigente declara el carácter de acción de quinto grado
\[
\boxed{
H_5(\alpha,\varphi)
=\frac12\sqrt{\frac{1000\alpha}{\varphi}}-\alpha
+\frac9{16}\alpha^2-\frac59\alpha^3
+\frac7{48}\alpha^4-\frac1{54}\alpha^5.}
\]
Los cuatro coeficientes racionales pueden escribirse como
\[
\frac9{16}=\frac{\lambda_+}{4^2},
\qquad
\frac59=\frac{\lambda_-}{\lambda_+},
\qquad
\frac7{48}=\frac{\tfrac12\operatorname{tr}\mathcal B_c}{4\cdot12},
\qquad
\frac1{54}=\frac{1944}{104976},
\]
con \((\lambda_+,\lambda_-)=(9,5)\). Estas identidades fijan la
representabilidad aritmética interna de los coeficientes del carácter. Su
orden, sus signos y su normalización pertenecen a la regla graduada del lector
regional y se aplican antes de toda carta de unidades. Las ablaciones
certifican que cada término interviene efectivamente en la fase resultante.
En particular, ni la mantisa SI de la acción ni el valor sexagesimal del
contraángulo forman parte del dominio de la función generadora.

% HMT_SOURCE_SEGMENT_END id=A07
\section{Aplicación parangular \texorpdfstring{\(P(\alpha,\eta_{\mathrm{ret}})\)}{P(alpha,eta_ret)} y clausura bicapa \texorpdfstring{\(C^*=2(\eta_{\mathrm{ret}}+\alpha)\)}{C*=2(eta_ret+alpha)}}
\label{sec:s102:c38:4}

\subsection{Fase orientada y coordenada angular conjugada}
\label{sec:s102:c38:4:1}

% HMT_SOURCE_SEGMENT_BEGIN id=B08 source=colaboracion/parte_iii/source/public_final/base_83/c30_body.tex body_lines=407-472

La fase de quinto grado anterior al retorno es
\[
y_5=\frac\pi{90}\bigl(H_5+\alpha\bigr).
\]
El carácter regional completo produce la corrección
\(\mathscr C_\pi(\alpha)\). Con la orientación fijada en la
\cref{p3-83-c30-s1:def:lector-regional-determinantal}, definimos
\[
\boxed{
y_*=\frac\pi{90}\bigl(H_5+\alpha\bigr)
-169\alpha^6
-\frac{D_A\alpha^7}{1-D_A\alpha}.}
\]
Ésta es la fórmula primaria. La unidad sexagesimal no interviene hasta aplicar la carta
\[
C_*^{\rm HMT}=\frac{180}{\pi}y_*.
\]

\begin{theorem}[Coordenada angular conjugada del lector regional determinantal]
\label{p3-83-c30-s1:thm:contraangulo-regional}
Fijados el cierre dodecafásico de \(\alpha\),
\(\varphi=(1+\sqrt5)/2\) y las reglas del lector regional determinantal, la
fase \(y_*\) y su coordenada angular conjugada quedan determinados sin introducir el valor
buscado. Su evaluación es
\[
\boxed{
y_*=0.03706622646583826398493779409230638\ldots,}
\]
\[
\boxed{
C_*^{\rm HMT}
=2.12373833896864572426195138039336\ldots{}^\circ.}
\]
Por tanto, a quince cifras decimales,
\[
\boxed{C_*^{\rm HMT}=2.123738338968646^\circ.}
\]
\end{theorem}

\begin{proof}
El \cref{p3-83-c30-s1:thm:unicidad-retorno-169} determina \(169\). El cierre de
\(\alpha\) determina \(A\), \(x\) y
\(D_A=e^{-100\pi\alpha/9}\). El \cref{p3-83-c30-s1:thm:realizacion-analitica-graduada}
determina la serie completa. La fórmula de \(H_5\) contiene únicamente
\(\alpha\), \(\varphi\) e invariantes estructurales declarados. Sustituir
estos datos en la expresión de \(y_*\) da el valor impreso. El valor decimal
de \(C_*^{\rm HMT}\) no aparece en ninguno de los miembros derechos.
\end{proof}

La componente reducida de acción posterior al retorno queda ahora como
salida:
\[
\boxed{
\bar h_{\rm HMT}^{\rm ret}
=\frac{C_*^{\rm HMT}}2-\alpha
=H_5-\frac{90}{\pi}\mathscr C_\pi(\alpha).}
\]
Numéricamente,
\[
\bar h_{\rm HMT}^{\rm ret}
=1.05457181691503906113369058472430\ldots.
\]
La identidad inversa no se ha usado para introducir \(C_*^{\rm HMT}\); se
obtiene después de generar \(y_*\).

% HMT_SOURCE_SEGMENT_END id=B08
\subsection{Coordenada conjugada \texorpdfstring{\(C^*\)}{C*} y reversión de su orientación bajo la involución bicapa}
\label{sec:s102:c38:4:2}

% HMT_SOURCE_SEGMENT_BEGIN id=A08 source=colaboracion/parte_iii/source/public_final/ascensos_20260826/16f_realizacion_analitica_contraangulo.tex body_lines=420-503

La fase de quinto grado anterior al retorno es
\[
y_5=\frac\pi{90}\bigl(H_5+\alpha\bigr).
\]
El carácter regional completo produce la corrección
\(\mathscr C_\pi(\alpha)\). Con la orientación fijada en la
\cref{p3-20260826:def:lector-regional-determinantal}, definimos
\[
\boxed{
\Crad:=y_*=\frac\pi{90}\bigl(H_5+\alpha\bigr)
-169\alpha^6
-\frac{D_A\alpha^7}{1-D_A\alpha}.}
\]
En adelante escribimos también
\[
 \boxed{y_{\rm reg}:=y_*}
\]
cuando sea necesario distinguir esta realización de la fase espectral de
\Gtr. Las dos notaciones designan aquí el mismo objeto regional canónico.
Ésta es la fórmula primaria. La unidad sexagesimal no interviene hasta aplicar la carta
\[
\boxed{\Cdeg:=\frac{180}{\pi}\Crad.}
\]
Los símbolos \(\Crad\) y \(\Cdeg\) no designan dos contraángulos: son las
coordenadas radial y sexagesimal del mismo elemento angular.

\begin{theorem}[Contraángulo del lector regional determinantal]
\label{p3-20260826:thm:contraangulo-regional}
Fijados el cierre dodecafásico de \(\alpha\),
\(\varphi=(1+\sqrt5)/2\), el germen \(H_5\) y las reglas del lector regional
determinantal, la fase \(\Crad\) y su contraángulo quedan determinados sin
introducir el valor buscado en la función de evaluación. Su evaluación es
\[
\boxed{
\Crad=0.03706622646583826398493779409230638\ldots,}
\]
\[
\boxed{
\Cdeg
=2.12373833896864572426195138039336\ldots{}^\circ.}
\]
Por tanto, a quince cifras decimales,
\[
\boxed{\Cdeg=2.123738338968646^\circ.}
\]
\end{theorem}

\begin{proof}
El \cref{p3-20260826:thm:unicidad-retorno-169} determina \(169\). El cierre de
\(\alpha\) determina \(\Adeg\), \(\Arad\) y
\(D_A=e^{-100\pi\alpha/9}\). El \cref{p3-20260826:thm:realizacion-analitica-graduada}
determina la serie completa. La fórmula fijada de \(H_5\) contiene únicamente
\(\alpha\), \(\varphi\) e invariantes estructurales declarados. Sustituir
estos datos en la expresión de \(\Crad\) da el valor impreso. El valor decimal
del contraángulo no aparece en ninguno de los miembros derechos de esta
evaluación. La comparación con una mantisa metrológica de acción se efectúa
únicamente después de obtener la fase y no modifica la función generadora.
\end{proof}

Sea \(u_\circ\) la base sexagesimal cuya vuelta completa tiene coordenada
\(360\), e introduzcamos la coordenada angular
\(\alpha_{\angle,{\rm deg}}:=\alpha\). La salida posterior al
retorno queda entonces tipada por
\[
\boxed{
\etaRet
=\frac{\Cdeg}2-\alpha_{\angle,{\rm deg}}
=H_5-\frac{90}{\pi}\mathscr C_\pi(\alpha).}
\]
Su coordenada radial es
\[
 \etaRetrad=\frac{\pi}{180}\etaRet
 =\frac{\Crad}{2}-\frac{\pi}{180}\alpha.
\]
Numéricamente,
\[
\etaRet
=1.05457181691503906113369058472430\ldots.
\]
La identidad inversa no se ha usado para introducir \(\Cdeg\); se obtiene
después de generar \(\Crad\). El símbolo \(\hbar\) queda reservado para el
elemento de la recta de acción construido en la sección siguiente.

% HMT_SOURCE_SEGMENT_END id=A08
\subsection{Teorema directo \texorpdfstring{\(C^*=2(\eta_{\mathrm{ret}}+\alpha)\)}{C*=2(eta_ret+alpha)} y carácter inverso de la identidad \texorpdfstring{\(\eta_{\mathrm{ret}}=C^*/2-\alpha\)}{eta_ret=C*/2-alpha}}
\label{sec:s102:c38:4:3}

% HMT_SOURCE_SEGMENT_BEGIN id=A13 source=colaboracion/parte_iii/source/public_final/ascensos_20260826/16f_realizacion_analitica_contraangulo.tex body_lines=750-797

El carácter de quinto grado y la coordenada posterior al retorno
son dos valores distintos:
\[
H_5
=1.05457181764615446962582182943306\ldots,
\]
\[
\etaRet
=1.05457181691503906113369058472430\ldots.
\]
El primero es el carácter local antes de incorporar el carácter regional; el
segundo es la coordenada orientada que queda después del retorno. La corrección
que genera el contraángulo impide identificarlos.

La comparación metrológica debe respetar esta distinción. El primer valor da
la mantisa
\[
2\pi H_5
=6.62607014999998792600111034989947\ldots,
\]
mientras que la coordenada posterior da
\[
2\pi\etaRet
=6.62607014540625433351074984759288\ldots.
\]
Como control externo únicamente, retenemos la mantisa \(6.62607015\) del
valor definitorio del Sistema Internacional \cite{BIPMSI}. No se incorpora
todavía su exponente ni su unidad a la recta HMT: la década se deriva en el
\cref{chap:orden-determinantal-accion} y la unidad pertenece a una carta
metrológica posterior. La comparación de mantisas da
\[
2\pi H_5-6.62607015
=-1.2073998889650101\ldots\times10^{-14},
\]
y no es una igualdad exacta. La proximidad del valor anterior al retorno es
un control numérico notable; no autoriza a sustituir la constante definitoria
del SI por la coordenada posterior ni a confundir una coordenada angular con
el elemento de acción que se construirá al incorporar la década.

Este resultado determina con precisión la situación matemática. El lector
regional genera \((\Crad,\Cdeg)\) y, desde ese único elemento angular, el funcional
\(\Gamma_{6\to8}\). El mismo cálculo separa la proximidad externa de la
mantisa producida por \(H_5\) de la coordenada de retorno que interviene en el
contraángulo. La separación no debilita la construcción: elimina una
identificación incompatible con las cifras exactas y convierte cada
comparación en una afirmación falsable de tipo bien definido.

% HMT_SOURCE_SEGMENT_END id=A13
% HMT_SOURCE_SEGMENT_BEGIN id=D01 source=manuscrito/sucesor_102/deltas_ley9/c38_par_angular_cartas_y_canales.tex body_lines=3-20
\label{sec:l9s102:par-angular}

La completación cubica
\begin{equation}
 10^3=(9+1)^3=9^3+3\cdot9^2+3\cdot9+1
 =729+243+27+1
 \label{eq:l9s102:base-1000}
\end{equation}
determina la carta global de la coordenada directa. La clausura bicapa aporta
el factor \(2\) de la coordenada conjugada:
\begin{equation}
 \boxed{
 A_{\deg}=10^3\alpha,
 \qquad
 C^*_{\deg}=2(\eta_{\mathrm{ret}}+\alpha)
 \quad\text{en }\mathbb R/(360\mathbb Z).}
 \label{eq:l9s102:a-cstar-deg}
\end{equation}
% HMT_SOURCE_SEGMENT_END id=D01
% HMT_SOURCE_SEGMENT_BEGIN id=D04 source=manuscrito/sucesor_102/deltas_ley9/c38_par_angular_cartas_y_canales.tex body_lines=60-66
La identidad inversa
\begin{equation}
 \eta_{\mathrm{ret}}=\frac12C^*_{\deg}-\alpha
 \label{eq:l9s102:eta-corolario-inverso}
\end{equation}
queda establecida como corolario algebraico de la construcción directa.

% HMT_SOURCE_SEGMENT_END id=D04
\section{Conmutatividad exacta de las cartas en vueltas, grados y radianes}
\label{sec:s102:c38:5}

\subsection{Derivación de la longitud cilíndrica y del grado analítico del precursor parangular}
\label{sec:s102:c38:5:1}

% HMT_SOURCE_SEGMENT_BEGIN id=B03 source=colaboracion/parte_iii/source/public_final/base_83/c30_body.tex body_lines=157-194

Sea \(\mathcal P\) el módulo libre generado por las palabras regionales y
sea \(F^n\mathcal P\) su filtración por longitud. La graduación asociada
registra el primer nivel en el que aparece cada palabra. Para un parámetro
\(z\), el carácter ordinario de longitud es
\[
\chi_z([w])=z^{|w|}.
\]
La multiplicatividad
\[
\chi_z([uv])=\chi_z([u])\chi_z([v])
\]
expresa que la concatenación suma longitudes. Al evaluar en
\(z=\alpha\), una palabra de longitud \(n\) recibe grado \(\alpha^n\).

\begin{definition}[Compatibilidad de grado]
\label{p3-83-c30-s1:def:compatibilidad-grado}
El lector analítico regional es compatible con la filtración cilíndrica
cuando envía el componente homogéneo de longitud \(n\) al grado analítico
\(n\):
\[
\operatorname{gr}_n\mathcal P\longrightarrow
\alpha^n\mathbb R.
\]
\end{definition}

Las cinco regiones pertenecen a \(\mathcal U_6\); por ello, el dato de
retorno aparece en grado seis:
\[
\rho_\pi[U_6]\longmapsto169\alpha^6.
\]
Este seis es la longitud de la palabra regional. No es la longitud de los
paseos cerrados del \cref{chap:cinco-ciclos-pi}. Aquellos paseos recorren a
la vez los cuatro anclajes
\(U_\pi,U_e,U_\varphi,U_{\rm APP}\) y poseen longitud mínima ocho. Una
palabra, una arista del grafo de bloques y un paseo cerrado son objetos de
tipos distintos; la graduación utilizada aquí pertenece al primero.

% HMT_SOURCE_SEGMENT_END id=B03
\subsection{Publicación del grado analítico en la carta global de fase}
\label{sec:s102:c38:5:2}

% HMT_SOURCE_SEGMENT_BEGIN id=A03 source=colaboracion/parte_iii/source/public_final/ascensos_20260826/16f_realizacion_analitica_contraangulo.tex body_lines=172-209

Sea \(\mathcal P\) el módulo libre generado por las palabras regionales y
sea \(F^n\mathcal P\) su filtración por longitud. La graduación asociada
registra el primer nivel en el que aparece cada palabra. Para un parámetro
\(z\), el carácter ordinario de longitud es
\[
\chi_z([w])=z^{|w|}.
\]
La multiplicatividad
\[
\chi_z([uv])=\chi_z([u])\chi_z([v])
\]
expresa que la concatenación suma longitudes. Al evaluar en
\(z=\alpha\), una palabra de longitud \(n\) recibe grado \(\alpha^n\).

\begin{definition}[Compatibilidad de grado]
\label{p3-20260826:def:compatibilidad-grado}
El lector analítico regional es compatible con la filtración cilíndrica
cuando envía el componente homogéneo de longitud \(n\) al grado analítico
\(n\):
\[
\operatorname{gr}_n\mathcal P\longrightarrow
\alpha^n\mathbb R.
\]
\end{definition}

Las cinco regiones pertenecen a \(\mathcal U_6\); por ello, el dato de
retorno aparece en grado seis:
\[
\rho_\pi[U_6]\longmapsto169\alpha^6.
\]
Este seis es la longitud de la palabra regional. No es la longitud de los
paseos cerrados del \cref{chap:cinco-ciclos-pi}. Aquellos paseos recorren a
la vez los cuatro anclajes
\(U_\pi,U_e,U_\varphi,U_{\rm APP}\) y poseen longitud mínima ocho. Una
palabra, una arista del grafo de bloques y un paseo cerrado son objetos de
tipos distintos; la graduación utilizada aquí pertenece al primero.

% HMT_SOURCE_SEGMENT_END id=A03
% HMT_SOURCE_SEGMENT_BEGIN id=D02 source=manuscrito/sucesor_102/deltas_ley9/c38_par_angular_cartas_y_canales.tex body_lines=21-34
El calendario global exhibe las tres factorizaciones
\begin{equation}
 360=12\cdot30=4\cdot90=3\cdot120.
 \label{eq:l9s102:360-factorizaciones}
\end{equation}
La carta analítica se obtiene mediante
\begin{equation}
 A_{\mathrm{rad}}=\frac\pi{180}A_{\deg}
 =\frac{50\pi}{9}\alpha,
 \qquad
 C^*_{\mathrm{rad}}=\frac\pi{180}C^*_{\deg}
 =\frac\pi{90}(\eta_{\mathrm{ret}}+\alpha).
 \label{eq:l9s102:a-cstar-rad}
\end{equation}
% HMT_SOURCE_SEGMENT_END id=D02
% HMT_SOURCE_SEGMENT_BEGIN id=D03 source=manuscrito/sucesor_102/deltas_ley9/c38_par_angular_cartas_y_canales.tex body_lines=35-59

\begin{theorem}[Conmutatividad de las factorizaciones regional y de acción]
\label{thm:l9s102:conmutatividad-cstar}
Las dos rutas hacia la fase conjugada producen la misma salida:
\begin{equation}
 \boxed{
 \frac\pi{180}C^*_{\deg}
 =\frac\pi{90}(\eta_{\mathrm{ret}}+\alpha)
 =\frac\pi{90}(H_5+\alpha)
 -\mathscr C_\pi(\alpha)=y_*.}
 \label{eq:l9s102:cuadrado-cstar}
\end{equation}
\end{theorem}

\begin{proof}
Sustituyendo \cref{eq:l9s102:eta-ret} en el segundo miembro se obtiene
\[
 \frac\pi{90}
 \left(H_5-\frac{90}{\pi}\mathscr C_\pi+\alpha\right)
 =\frac\pi{90}(H_5+\alpha)-\mathscr C_\pi.
\]
La última expresion es la definicion regional de \(y_*\), mientras la
primera igualdad procede de \cref{eq:l9s102:a-cstar-deg}.
\end{proof}

% HMT_SOURCE_SEGMENT_END id=D03
\section{Descomposición espectral \texorpdfstring{\(q_\pm\)}{q+/-} y cálculo funcional sobre los canales \texorpdfstring{\(90/120\)}{90/120} previamente generados}
\label{sec:s102:c38:6}

\subsection{Reconstrucción espectral del par angular y transición hexada--octada}
\label{sec:s102:c38:6:1}

% HMT_SOURCE_SEGMENT_BEGIN id=B11 source=colaboracion/parte_iii/source/public_final/base_83/c30_body.tex body_lines=523-569

La fase común y la fase orientada determinan los dos canales
\[
q_-=e^{-x+y_*},
\qquad
q_+=e^{-x-y_*}.
\]
Se recuperan exactamente los invariantes
\[
q_-q_+=D_A,
\qquad
\alpha=-\frac9{100\pi}\log D_A,
\qquad
C_*^{\rm HMT}=\frac{90}{\pi}\log\frac{q_-}{q_+}.
\]
La primera igualdad lee la contracción común; la tercera, la separación
orientada de las hojas. El intercambio \(R\mapsto-R\) permuta
\(q_+\) y \(q_-\), conserva \(D_A\) e invierte el signo de \(y_*\).

Para
\[
S_{90}(q)=\frac{q^{90}}{1-q^{270}},
\qquad
S_{120}(q)=\frac{q^{120}}{1-q^{360}},
\]
el funcional de transición hexada--octada se escribe íntegramente en la
carta de fases:
\[
\Gamma_{6\to8}
=\frac{180}{\pi}xy_*
+12\bigl(S_{90}(q_-)-S_{90}(q_+)\bigr)
-\bigl(S_{120}(q_-)-S_{120}(q_+)\bigr).
\]
La evaluación producida por la coordenada \(C_*^{\rm HMT}\) del
\cref{p3-83-c30-s1:thm:contraangulo-regional} es
\[
\boxed{
\Gamma_{6\to8}
=0.274007672694459274747255260774\ldots.}
\]
Así, el valor interno reconocido en Mecánica Dimensional como funcional de
Barbero--Immirzi deja de recibir \(C_*^{\rm HMT}\) como literal independiente:
ambos descienden del mismo par de fases generado en este capítulo. La
interpretación geométrica y física de \(\Gamma_{6\to8}\) pertenece a la obra
de Mecánica Dimensional; su evaluación matemática pertenece ya a la interfaz
previamente construida.

% HMT_SOURCE_SEGMENT_END id=B11
\subsection{Transportador relativo y disjunción espectral de los bloques}
\label{sec:s102:c38:6:2}

% HMT_SOURCE_SEGMENT_BEGIN id=B13 source=colaboracion/parte_iii/source/public_final/base_83/c30_body.tex body_lines=619-675
\label{p3-83-c30-s1:sec:transportador-relativo-split}

En la base ordenada de canales, sea
\[
 R=\begin{pmatrix}0&1\\1&0\end{pmatrix},
 \qquad P_\pm=\frac{I\pm R}{2},
 \qquad
 B_0=7I+2R,
 \qquad B_1=\Adeg I+\Cdeg R.
\]
La descomposición escindida
\[
 xI+yR=\rho(x,y)e^{\eta(x,y)R},
 \quad \rho(x,y)=\sqrt{x^2-y^2},
 \quad \eta(x,y)=\operatorname{artanh}(y/x)
\]
construye el transportador único
\[
 T_{\rm rel}
 =\frac{\sqrt{(\Adeg)^2-(\Cdeg)^2}}{\sqrt{45}}
   \exp\!\left[
   \left(\operatorname{artanh}\frac{\Cdeg}{\Adeg}
   -\operatorname{artanh}\frac27\right)R\right],
 \qquad T_{\rm rel}B_0=B_1.
\]
En la base espectral actúa por
\(9\mapsto\Adeg+\Cdeg\) y
\(5\mapsto\Adeg-\Cdeg\).

\begin{theorem}[Disjunción espectral de las realizaciones local y angular]
\label{p3-83-c30-s1:thm:no-go-entrelazador-bloques-split}
Los espectros
\[
 \operatorname{Spec}(B_0)=\{9,5\},\qquad
 \operatorname{Spec}(B_1)=\{\Adeg+\Cdeg,\Adeg-\Cdeg\}
\]
son disjuntos. En consecuencia,
\[
 \operatorname{Hom}_{B_0,B_1}
 :=\{\Phi:\Phi B_0=B_1\Phi\}=\{0\}.
\]
\end{theorem}

\begin{proof}
Las cotas racionales construidas para las coordenadas angulares dan
\(7.29<\Adeg<7.30\) y \(2.12<\Cdeg<2.13\); por tanto,
\(9.41<\Adeg+\Cdeg<9.43\) y
\(5.16<\Adeg-\Cdeg<5.18\). En la base que diagonaliza \(R\), cada entrada
de un entrelazador queda multiplicada por la diferencia entre un autovalor
de \(B_0\) y uno de \(B_1\); la disjunción obliga a que todas se anulen.
\end{proof}

El transportador relativo compara dos bloques dentro del álgebra escindida;
los entrelazadores dodecafásico--Witt y de incidencia actúan entre
representaciones equivalentes en sus residencias propias. Esta separación
tipa positivamente ambas operaciones y conserva sus dominios.

% HMT_SOURCE_SEGMENT_END id=B13
\subsection{Evaluación espectral \texorpdfstring{\(V_{90,120}\)}{V_90,120} de los canales de incidencia previamente generados por el TPK}
\label{sec:s102:c38:6:3}

% HMT_SOURCE_SEGMENT_BEGIN id=A11 source=colaboracion/parte_iii/source/public_final/ascensos_20260826/16f_realizacion_analitica_contraangulo.tex body_lines=554-572

La fase común y la fase orientada determinan los dos canales
\[
q_-=e^{-\Arad+\Crad},
\qquad
q_+=e^{-\Arad-\Crad}.
\]
Se recuperan exactamente los invariantes
\[
q_-q_+=D_A,
\qquad
\alpha=-\frac9{100\pi}\log D_A,
\qquad
\Cdeg=\frac{90}{\pi}\log\frac{q_-}{q_+}.
\]
La primera igualdad lee la contracción común; la tercera, la separación
orientada de las hojas. El intercambio \(R\mapsto-R\) permuta
\(q_+\) y \(q_-\), conserva \(D_A\) e invierte el signo de \(\Crad\).

% HMT_SOURCE_SEGMENT_END id=A11
% HMT_SOURCE_SEGMENT_BEGIN id=D05 source=manuscrito/sucesor_102/deltas_ley9/c38_par_angular_cartas_y_canales.tex body_lines=68-97
\label{sec:l9s102:qpm-posteriores}

Sea \(H^2=I\) y \(P_\pm=(I\pm H)/2\). El bloque angular
\begin{equation}
 B=A_{\mathrm{rad}}I+C^*_{\mathrm{rad}}H
 =(A_{\mathrm{rad}}+C^*_{\mathrm{rad}})P_+
 +(A_{\mathrm{rad}}-C^*_{\mathrm{rad}})P_-
 \label{eq:l9s102:b-angular}
\end{equation}
se publica multiplicativamente por
\begin{equation}
 e^{-B}=q_+P_++q_-P_-,
 \qquad
 q_\pm=\exp[-(A_{\mathrm{rad}}\pm C^*_{\mathrm{rad}})].
 \label{eq:l9s102:qpm-angular}
\end{equation}
La involución \(C^*\mapsto-C^*\) permuta \(q_+\) y \(q_-\), y conserva
\begin{equation}
 q_+q_-=e^{-2A_{\mathrm{rad}}}=D_A.
 \label{eq:l9s102:det-angular}
\end{equation}
Los cardinales \(90/120\), construidos previamente por la incidencia del
TPK, reciben ahora su evaluación espectral:
\begin{equation}
 V_{90,120}
 =(I-T^{90})(I-T^{120})^{-1}
 =\sum_{\sigma\in\{+,-\}}
 \frac{1-q_\sigma^{90}}{1-q_\sigma^{120}}P_\sigma.
 \label{eq:l9s102:v90120-posterior}
\end{equation}
% HMT_SOURCE_SEGMENT_END id=D05
\subsubsection{Transportador relativo escindido}
\label{sec:s102:c38:6:3:1}

% HMT_SOURCE_SEGMENT_BEGIN id=A12 source=colaboracion/parte_iii/source/public_final/ascensos_20260826/16f_realizacion_analitica_contraangulo.tex body_lines=574-748
\label{p3-20260826:sec:transportador-relativo-split}

La comparación con el bloque central de APP puede efectuarse sin identificar
representaciones distintas. Fijemos la base ordenada de canales y la
involución
\[
 R=\begin{pmatrix}0&1\\1&0\end{pmatrix},
 \qquad R^2=I,
 \qquad
 P_\pm=\frac{I\pm R}{2}.
\]
En el álgebra conmutativa
\[
 \mathscr B_R:=\{xI+yR:x,y\in\mathbb R\}
\]
consideremos, dentro de la carta sexagesimal ya declarada,
\[
 B_0:=7I+2R,
 \qquad
 B_1:=\Adeg I+\Cdeg R.
\]
El cambio ortogonal de base
\[
 S=\frac1{\sqrt2}
 \begin{pmatrix}1&1\\1&-1\end{pmatrix}
\]
diagonaliza simultáneamente la involución y ambos bloques. Para \(x>|y|\),
definamos
\[
 \rho(x,y):=\sqrt{x^2-y^2},
 \qquad
 \eta(x,y):=\operatorname{artanh}\!\left(\frac yx\right).
\]
Entonces
\[
 xI+yR=\rho(x,y)e^{\eta(x,y)R}.
\]
Si
\[
 \rho_0=\sqrt{45},\quad
 \eta_0=\operatorname{artanh}(2/7),\quad
 \rho_1=\rho(\Adeg,\Cdeg),\quad
 \eta_1=\eta(\Adeg,\Cdeg),
\]
la única multiplicación izquierda que lleva \(B_0\) a \(B_1\) es
\[
 \boxed{
 T_{\rm rel}
 :=\frac{\rho_1}{\sqrt{45}}
 e^{(\eta_1-\eta_0)R},
 \qquad
 T_{\rm rel}B_0=B_1.}
\]

\begin{proposition}[Transportador relativo de los dos bloques]
\label{p3-20260826:prop:transportador-relativo-split}
Una vez fijados \(R\), el orden de los canales, la rama positiva y la carta
sexagesimal, \(T_{\rm rel}\) es el único elemento de
\(\mathscr B_R\) que satisface \(TB_0=B_1\). En la base espectral actúa por
\[
 9\longmapsto\Adeg+\Cdeg,
 \qquad
 5\longmapsto\Adeg-\Cdeg.
\]
\end{proposition}

\begin{proof}
La descomposición polar escindida da la fórmula impresa y prueba la existencia.
Como \(\det B_0=45\ne0\), toda solución matricial de \(TB_0=B_1\) satisface
\(T=B_1B_0^{-1}\); por tanto es única. La diagonalización mediante \(S\)
produce los dos cocientes espectrales
\((\Adeg+\Cdeg)/9\) y \((\Adeg-\Cdeg)/5\), equivalentes a la expresión
exponencial de \(T_{\rm rel}\).
\end{proof}

Este transportador no es un entrelazador representacional. En efecto, un
entrelazador \(\Phi\) entre los endomorfismos \(B_0\) y \(B_1\) debería
satisfacer
\[
 \Phi B_0=B_1\Phi.
\]

\begin{theorem}[Obstrucción espectral al entrelazamiento]
\label{p3-20260826:thm:no-go-entrelazador-bloques-split}
Para los valores de \(\Adeg\) y \(\Cdeg\) construidos en esta sección,
\[
 \boxed{\Phi B_0=B_1\Phi\quad\Longrightarrow\quad\Phi=0.}
\]
\end{theorem}

\begin{proof}
Los espectros son
\[
 \operatorname{Spec}(B_0)=\{9,5\},
 \qquad
 \operatorname{Spec}(B_1)
 =\{\Adeg+\Cdeg,\Adeg-\Cdeg\}.
\]
Las cotas racionales previamente demostradas proporcionan los intervalos
disjuntos
\[
 7.29<\Adeg<7.30,
 \qquad
 2.12<\Cdeg<2.13,
\]
y, por consiguiente,
\[
 9.41<\Adeg+\Cdeg<9.43,
 \qquad
 5.16<\Adeg-\Cdeg<5.18.
\]
En la base que diagonaliza \(R\), cada entrada \(\phi_{ij}\) de \(\Phi\)
queda multiplicada por la diferencia entre un autovalor de \(B_0\) y uno de
\(B_1\). Todas esas diferencias son no nulas; luego todas las entradas de
\(\Phi\) se anulan.
\end{proof}

Numéricamente,
\[
 \rho_1=6.9814819335172378048\ldots,
 \qquad
 \frac{\rho_1}{\sqrt{45}}=1.0407378791354140779\ldots,
\]
\[
 \eta_1-\eta_0
 =0.005796351470571295590\ldots.
\]
La fórmula conserva el álgebra escindida, los proyectores \(P_\pm\) y la forma
lorentziana hasta el factor conforme \(\rho_1^2/45\). Depende, sin embargo,
de \(B_1\) ya construido: no constituye una derivación anterior de
\(\Adeg\) o \(\Cdeg\), ni debe confundirse con los entrelazadores
equivariantes entre realizaciones de un mismo módulo.

Para
\[
S_{90}(q)=\frac{q^{90}}{1-q^{270}},
\qquad
S_{120}(q)=\frac{q^{120}}{1-q^{360}},
\]
el funcional de transición hexada--octada se escribe íntegramente en la
carta de fases:
\[
\Gamma_{6\to8}
=\frac{180}{\pi}\Arad\Crad
+12\bigl(S_{90}(q_-)-S_{90}(q_+)\bigr)
-\bigl(S_{120}(q_-)-S_{120}(q_+)\bigr).
\]
La evaluación producida por el contraángulo del
\cref{p3-20260826:thm:contraangulo-regional} es
\[
\boxed{
\Gamma_{6\to8}
=0.274007672694459274747255260774\ldots.}
\]
Las formas equivalentes del funcional matemático interno son
\[
\boxed{
\Gamma_{6\to8}
=\Arad\Cdeg+K_{6\to8}
=\frac{180}{\pi}\Arad\Crad+K_{6\to8}
=\frac{\pi\Adeg\Cdeg}{180}+K_{6\to8}.}
\]
Estas identidades evalúan \(\Gamma_{6\to8}\) después de generar
\((\Crad,\Cdeg)\); no se utiliza en sentido inverso para seleccionar el
contraángulo. La lectura física en la conexión de Ashtekar--Barbero se
representa después mediante
\[
 \Gamma_{6\to8}\dashrightarrow\gamma_{\rm BI}.
\]
La flecha no introduce un tercer funcional ni retroactúa sobre \(C^*\). Las
series \(S_{90}\), \(S_{120}\), los coeficientes \(12\) y \(-1\) y su
asignación al paso \(6\to8\) pertenecen a la construcción matemática
declarada. La identificación física exige, además, la normalización propia
de la conexión.

% HMT_SOURCE_SEGMENT_END id=A12
\section{Naturalidad bajo refinamiento y separación respecto del reconocimiento metrológico}
\label{sec:s102:c38:7}

\subsection{Estimación exacta del resto de la prolongación analítica}
\label{sec:s102:c38:7:1}

% HMT_SOURCE_SEGMENT_BEGIN id=B09 source=colaboracion/parte_iii/source/public_final/base_83/c30_body.tex body_lines=474-499

La expresión hallada inicialmente,
\[
y_7=\frac\pi{90}(H_5+\alpha)-169\alpha^6-D_A\alpha^7,
\]
es el truncamiento de grado siete del carácter completo. Produce
\[
C_7=2.12373833896864600265403827103919\ldots{}^\circ.
\]
El resto omitido no es un término indeterminado; es la cola geométrica
exacta
\[
\mathscr C_\pi(\alpha)
-169\alpha^6-D_A\alpha^7
=\frac{D_A^2\alpha^8}{1-D_A\alpha}.
\]
Por tanto,
\[
\left|C_7-C_*^{\rm HMT}\right|
=\frac{180}{\pi}\frac{D_A^2\alpha^8}{1-D_A\alpha}
=2.7839208689064583\ldots\times10^{-16}\;{}^\circ.
\]
La recurrencia transforma así una aproximación de séptimo grado en una
realización analítica completa y proporciona una cota cerrada para cada
truncamiento posterior.

% HMT_SOURCE_SEGMENT_END id=B09
\subsection{Dependencia funcional del selector regional y de la continuación analítica}
\label{sec:s102:c38:7:2}

% HMT_SOURCE_SEGMENT_BEGIN id=B10 source=colaboracion/parte_iii/source/public_final/base_83/c30_body.tex body_lines=501-521

La construcción cambia al retirar cualquiera de sus componentes activos.
Los siguientes valores se obtienen sin recalibrar los demás términos:
\[
\begin{array}{@{}lc@{}}
\toprule
\text{variante}&\text{coordenada angular conjugada en grados}\\
\midrule
\rho_\pi=168&2.1237383389772976863904487\ldots\\
\rho_\pi=169&2.1237383389686457242619514\ldots\\
\rho_\pi=170&2.1237383389599937621334541\ldots\\
\text{sin cola determinantal}&2.1237383389686949415301675\ldots\\
D_A\text{ sustituido por }1&2.1237383389686313409965826\ldots\\
\bottomrule
\end{array}
\]
Estas ablaciones no constituyen por sí solas una estimación probabilística.
Demuestran que el selector equivariante y el carácter determinantal son
componentes efectivas de la salida y que no pueden suprimirse manteniendo la
misma realización.

% HMT_SOURCE_SEGMENT_END id=B10
\subsection{Interpretación operatoria de la coordenada angular conjugada}
\label{sec:s102:c38:7:3}

% HMT_SOURCE_SEGMENT_BEGIN id=B14 source=colaboracion/parte_iii/source/public_final/base_83/c30_body.tex body_lines=677-715

El procedimiento completo puede resumirse como una cadena de aplicaciones:
\[
\begin{aligned}
\mathcal F_\pi
&\longmapsto b_\pi
\longmapsto\rho_\pi=169,\\
\alpha
&\longmapsto A
\longmapsto x
\longmapsto D_A,\\
(\rho_\pi,D_A)
&\longmapsto\mathscr C_\pi(\alpha),\\
(H_5,\mathscr C_\pi(\alpha))
&\longmapsto y_*
\longmapsto C_*^{\rm HMT},\\
(x,y_*)
&\longmapsto(q_+,q_-)
\longmapsto\Gamma_{6\to8}.
\end{aligned}
\]
La primera línea pertenece al catálogo finito; la segunda, al cierre común de
\(\alpha\); la tercera, a la prolongación determinantal; la cuarta produce la
coordenada orientada; la quinta la transporta hacia la interfaz dimensional.

En esta organización, \(\alpha\) mide la contracción común del transporte y
\(C_*^{\rm HMT}\) mide la separación orientada de sus dos hojas. El retorno
regional no añade una constante externa: selecciona una amplitud finita y la
prolonga mediante el determinante ya fijado por \(\alpha\). La carta en
grados es posterior. El objeto primario es el par de fases \((x,y_*)\).

Una verificación exhaustiva reconstruye las cinco regiones directamente
desde el catálogo, recorre las acciones de
\(\mathfrak S_5\) y \(\mathfrak S_3\), comprueba la recurrencia y suma la
serie. La entrada de \(\alpha\) es la salida exacta de su cierre
dodecafásico, no una cifra metrológica transcrita. La comprobación mantiene
separados los teoremas del catálogo, las reglas del lector y el contraste
externo; por ello ninguna comparación experimental interviene en la función
generadora.
% HMT_SOURCE_SEGMENT_END id=B14
\subsection{Certificación de convergencia del truncamiento parangular}
\label{sec:s102:c38:7:4}

% HMT_SOURCE_SEGMENT_BEGIN id=A09 source=colaboracion/parte_iii/source/public_final/ascensos_20260826/16f_realizacion_analitica_contraangulo.tex body_lines=505-530

La expresión hallada inicialmente,
\[
y_7=\frac\pi{90}(H_5+\alpha)-169\alpha^6-D_A\alpha^7,
\]
es el truncamiento de grado siete del carácter completo. Produce
\[
C_7=2.12373833896864600265403827103919\ldots{}^\circ.
\]
El resto omitido no es un término indeterminado; es la cola geométrica
exacta
\[
\mathscr C_\pi(\alpha)
-169\alpha^6-D_A\alpha^7
=\frac{D_A^2\alpha^8}{1-D_A\alpha}.
\]
Por tanto,
\[
\left|C_7-\Cdeg\right|
=\frac{180}{\pi}\frac{D_A^2\alpha^8}{1-D_A\alpha}
=2.7839208689064583\ldots\times10^{-16}\;{}^\circ.
\]
La recurrencia transforma así una aproximación de séptimo grado en una
realización analítica completa y proporciona una cota cerrada para cada
truncamiento posterior.

% HMT_SOURCE_SEGMENT_END id=A09
\subsection{Naturalidad del selector regional bajo refinamiento y cambio de carta}
\label{sec:s102:c38:7:5}

% HMT_SOURCE_SEGMENT_BEGIN id=A10 source=colaboracion/parte_iii/source/public_final/ascensos_20260826/16f_realizacion_analitica_contraangulo.tex body_lines=532-552

La construcción cambia al retirar cualquiera de sus componentes activos.
Los siguientes valores se obtienen sin recalibrar los demás términos:
\[
\begin{array}{@{}lc@{}}
\toprule
\text{variante}&\text{contraángulo en grados}\\
\midrule
\rho_\pi=168&2.1237383389772976863904487\ldots\\
\rho_\pi=169&2.1237383389686457242619514\ldots\\
\rho_\pi=170&2.1237383389599937621334541\ldots\\
\text{sin cola determinantal}&2.1237383389686949415301675\ldots\\
D_A\text{ sustituido por }1&2.1237383389686313409965826\ldots\\
\bottomrule
\end{array}
\]
Estas ablaciones no constituyen por sí solas una estimación probabilística.
Demuestran que el selector equivariante y el carácter determinantal son
componentes efectivas de la salida y que no pueden suprimirse manteniendo la
misma realización.

% HMT_SOURCE_SEGMENT_END id=A10
\subsection{Independencia metrológica de la construcción parangular}
\label{sec:s102:c38:7:6}

% HMT_SOURCE_SEGMENT_BEGIN id=A14 source=colaboracion/parte_iii/source/public_final/ascensos_20260826/16f_realizacion_analitica_contraangulo.tex body_lines=799-831

El procedimiento completo puede resumirse como una cadena de aplicaciones:
\[
\begin{aligned}
\mathcal F_\pi
&\longmapsto b_\pi
\longmapsto\rho_\pi=169,\\
\alpha
&\longmapsto \Adeg
\longmapsto \Arad
\longmapsto D_A,\\
(\rho_\pi,D_A)
&\longmapsto\mathscr C_\pi(\alpha),\\
(H_5,\mathscr C_\pi(\alpha))
&\longmapsto \Crad
\longmapsto \Cdeg,\\
(\Arad,\Crad)
&\longmapsto(q_+,q_-)
\longmapsto\Gamma_{6\to8}.
\end{aligned}
\]
La primera línea pertenece al catálogo finito; la segunda, al cierre común de
\(\alpha\); la tercera, a la prolongación determinantal; la cuarta produce la
coordenada orientada; la quinta la transporta hacia la interfaz dimensional.

En esta organización, \(\alpha\) mide la contracción común del transporte y
el contraángulo mide la separación orientada de sus dos hojas. En la
evaluación vigente, el retorno regional selecciona una amplitud finita y la
prolonga mediante el determinante ya fijado por \(\alpha\), sin insertar una
constante externa adicional. La carta en grados es posterior. El objeto
primario es el par de fases \((\Arad,\Crad)\). La recurrencia, la suma
regional y las ablaciones dependen del catálogo y de
\(\alpha\), pero no reciben una constante metrológica adicional.
% HMT_SOURCE_SEGMENT_END id=A14

 
\chapter{Operador constitutivo positivo del vacío y derivación conjunta de \texorpdfstring{\(\varepsilon,\mu,Z,c\)}{epsilon, mu, Z, c}}
\begingroup
\renewcommand{\chapter}[2][]{}%
\chapter{Operador constitutivo positivo del vacío: completación \(3\to4\), reconstrucción espectral y derivación conjunta de \(\varepsilon\), \(\mu\), \(Z\) y \(c\)}
\label{chap:p3-final:39:vacio_constitutivo}
\section{Descomposici\'on espectral bicapa del operador angular}

Sea \(R=R^*=R^{-1}\) la involuci\'on activa y

\[
P_\pm=\frac{I\pm R}{2},
\qquad
T=q_+P_+ +q_-P_-,
\qquad 0<q_+<q_-<1.
\]

Con
\[
x=A_{\rm rad}=\frac{\pi A}{180},\qquad
y=C^*_{\rm rad}=\frac{\pi C^*}{180},
\]
los dos canales se construyen antes de evaluar el vac\'io:
\begin{equation}
q_+=e^{-x-y},\qquad q_-=e^{-x+y}.
\label{eq:vacuum-angular-channels}
\end{equation}
El intercambio de hojas invierte \(C^*\), es decir \(y\mapsto-y\), y permuta \(q_+\leftrightarrow q_-\). La pareja es as\'i la lectura espectral de un mismo antecedente angular, no dos par\'ametros ajustados por separado.

\begin{definition}[Respuesta constitutiva]
\begin{equation}
\mathcal V_{90,120}=(I-T^{90})(I-T^{120})^{-1}.
\label{eq:vacuum-operator}
\end{equation}
\end{definition}

Como \(\|T\|<1\), el denominador es positivo e invertible. El c\'alculo funcional da

\begin{equation}
\mathcal V_{90,120}=r_+P_+ +r_-P_-,
\qquad
r_\pm=\frac{1-q_\pm^{90}}{1-q_\pm^{120}}.
\label{eq:vacuum-eigenvalues}
\end{equation}

La afirmaci\'on de invertibilidad merece hacerse expl\'icita. Para cada hoja,
\(0<q_\pm^{120}<1\); por tanto
\[
I-T^{120}=(1-q_+^{120})P_++(1-q_-^{120})P_-
\]
es estrictamente positivo y
\[
(I-T^{120})^{-1}
=\frac{P_+}{1-q_+^{120}}+
 \frac{P_-}{1-q_-^{120}}.
\]
La respuesta \(\mathcal V_{90,120}\) no se obtiene resolviendo cuatro
ecuaciones escalares separadas: es el resultado del c\'alculo funcional de una
sola contracci\'on positiva. Esta observaci\'on impide que permeabilidad,
permitividad, impedancia y propagaci\'on pierdan el antecedente que comparten.

\begin{theorem}[Determinaci\'on espectral conjunta de las cuatro relaciones constitutivas]
Los cuatro caracteres constitutivos
\begin{equation}
\widehat\varepsilon=r_+^2,\qquad
\widehat\mu=r_-^2,\qquad
\widehat Z=\frac{r_-}{r_+},\qquad
\widehat c=\frac1{r_+r_-}
\label{eq:vacuum-four}
\end{equation}
son funciones espectrales de \(\mathcal V_{90,120}\). Satisfacen
\begin{equation}
\widehat\mu\widehat\varepsilon=\widehat c^{-2},
\qquad
\frac{\widehat\mu}{\widehat\varepsilon}=\widehat Z^2.
\label{eq:vacuum-relations}
\end{equation}
\end{theorem}

\begin{proof}
Las identidades se obtienen sustituyendo \cref{eq:vacuum-four}. No interviene ninguna carta dimensional ni ning\'un valor exterior.
\end{proof}

\begin{corollary}[Reconstrucci\'on de las dos hojas]
Los cuatro caracteres no contienen cuatro grados de libertad. Cualquiera de
los pares \((\widehat\mu,\widehat\varepsilon)\),
\((\widehat Z,\widehat c)\) o \((\mathcal E,\mathcal B)\) reconstruye los
dos autovalores positivos:
\begin{align}
r_-&=\sqrt{\widehat\mu}
=\sqrt{\widehat Z/\widehat c},
&r_+&=\sqrt{\widehat\varepsilon}
=\frac1{\sqrt{\widehat Z\widehat c}},
\label{eq:vacuum-reconstruct-r}\\
\log r_-&=\frac{\mathcal E+\mathcal B}{2},
&\log r_+&=\frac{\mathcal E-\mathcal B}{2}.
\label{eq:vacuum-reconstruct-log}
\end{align}
En particular, la orientaci\'on se pierde si se conserva s\'olo el producto,
pero se recupera al a\~nadir el cociente.
\end{corollary}

\begin{proof}
La primera l\'inea es la ra\'iz positiva de las identidades
\cref{eq:vacuum-four,eq:vacuum-relations}. La segunda invierte el sistema
lineal
\(\mathcal E=\log r_-+\log r_+\),
\(\mathcal B=\log r_--\log r_+\).
\end{proof}

Este corolario expresa una propiedad central de HMT: producto y cociente son
funcionales complementarios de un mismo estado. El producto determina la
propagaci\'on com\'un, mientras que el cociente conserva la orientaci\'on
relativa de las hojas interior y exterior.

En el certificado V2,

\begin{align*}
r_+&=0.9999996285482493631786952281\ldots,\\
r_-&=0.9997241371895183641315165853\ldots,\\
\widehat\varepsilon&=0.9999992570966367027604416155\ldots,\\
\widehat\mu&=0.9994483504793269350899815559\ldots,\\
\widehat Z&=0.9997245085389372154555543466\ldots,\\
\widehat c&=1.0002763104861575261063862689\ldots.
\end{align*}

Estos numerales son el reconocimiento terminal de \cref{eq:vacuum-operator}; no definen sus exponentes \(90,120\).

\section{Factorizaci\'on arm\'onica de los sectores \(90\) y \(120\)}

El semigrupo arm\'onico es

\[
\langle90,120\rangle=30\langle3,4\rangle.
\]

Si \(s=q^{30}\), entonces

\begin{equation}
\frac{1-q^{90}}{1-q^{120}}
=\frac{1-s^3}{1-s^4}
=\frac{1+s+s^2}{1+s+s^2+s^3}.
\label{eq:three-four-completion}
\end{equation}

La identidad tambi\'en separa la respuesta de su defecto. Si
\[
F_{3\to4}(s)=\frac{1+s+s^2}{1+s+s^2+s^3},
\qquad 0<s<1,
\]
entonces
\begin{equation}
d_{3\to4}(s):=1-F_{3\to4}(s)
=\frac{s^3}{1+s+s^2+s^3}>0.
\label{eq:vacuum-defect34}
\end{equation}
El t\'ermino adicional se conserva exactamente como defecto
positivo. Adem\'as,
\begin{equation}
F_{3\to4}'(s)
=-\frac{s^2(3+2s+s^2)}{(1+s+s^2+s^3)^2}<0.
\label{eq:vacuum-monotonic34}
\end{equation}
Por consiguiente, el orden de los canales angulares determina un\'ivocamente
el orden de las respuestas constitutivas. La asignaci\'on de
\(\widehat\mu\) y \(\widehat\varepsilon\) a las hojas respectivas se obtiene
anal\'iticamente, sin recurrir a sus expansiones decimales.

Cada hoja del vac\'io constituye la completaci\'on exacta de tres a cuatro
t\'erminos de un mismo modo elemental \(30\). En consecuencia, permeabilidad
y permitividad se interpretan como dos funcionales cuadr\'aticos de una misma
completaci\'on evaluada sobre hojas conjugadas, y no como par\'ametros
constitutivos independientes.

Como la funci\'on \((1-q^{90})/(1-q^{120})\) decrece en \(0<q<1\) y \(q_+<q_-\), se obtiene

\[
0<r_-<r_+<1,\qquad
0<\widehat\mu<\widehat\varepsilon<1,\qquad
0<\widehat Z<1<\widehat c.
\]

Los l\'imites del funcional determinan su geometr\'ia. Cuando \(q\to0^+\),
\(F(q)\to1\): las cuatro lecturas se acercan a la unidad y el defecto de
memoria desaparece. Cuando \(q\to1^-\),
\[
F(q)\longrightarrow\frac{90}{120}=\frac34.
\]
La respuesta toma, por tanto, valores en el intervalo positivo \((3/4,1)\). El
extremo \(3/4\) no es una constante ajustada: es el cociente de las dos
longitudes de canal y la realizaci\'on continua de la completaci\'on
tres--a--cuatro.

\section{Involuci\'on de hojas y conservaci\'on de la orientaci\'on}

El intercambio de hojas \(C^*\mapsto-C^*\) act\'ua por

\begin{equation}
\widehat\mu\longleftrightarrow\widehat\varepsilon,\qquad
\widehat Z\longmapsto\widehat Z^{-1},
\qquad
\widehat c\longmapsto\widehat c.
\label{eq:vacuum-sheet}
\end{equation}

La propagaci\'on com\'un es par; la impedancia conserva la orientaci\'on.
Esta distinci\'on reaparecer\'a en Barbero y CKM: no todos los funcionales
asociados al cambio de hoja tienen la misma paridad.

\section{Caracteres logar\'itmicos y compatibilidad tracial}

Definamos

\[
\mathcal E=\log(r_+r_-),\qquad
\mathcal B=\log(r_-/r_+).
\]

Con la traza normalizada \(\tau=\Tr/2\),

\begin{equation}
2\tau(\log\mathcal V)=\mathcal E,
\qquad
-2\tau(R\log\mathcal V)=\mathcal B.
\label{eq:vacuum-traces}
\end{equation}

La pareja \((\mathcal E,\mathcal B)\) constituye un sistema de coordenadas lineal del logaritmo
operatorio:
\begin{equation}
\log\mathcal V
=\frac{\mathcal E}{2}I-\frac{\mathcal B}{2}R.
\label{eq:vacuum-log-decomposition}
\end{equation}
La componente escalar es invariante bajo el intercambio de hojas; la
componente orientada cambia de signo. Esta descomposici\'on demuestra
directamente la paridad de la propagaci\'on y el car\'acter orientado de la
impedancia.

La amplificaci\'on non\'adica

\[
\mathcal V_m=\mathcal V\otimes I_{9^m}
\]

es compatible con las expectativas \(E_m=\mathrm{id}\otimes\tau_9\):

\[
E_m(\mathcal V_{m+1})=\mathcal V_m.
\]

Por tanto, \(\mathcal E\) y \(\mathcal B\) no son peculiaridades de una matriz \(2\times2\); permanecen en la torre UHF y en su cierre tracial \(II_1\).

M\'as precisamente, para todo polinomio \(p\) y despu\'es, por continuidad,
para toda funci\'on continua en el espectro,
\begin{equation}
E_m\bigl(p(\mathcal V_{m+1})\bigr)=p(\mathcal V_m).
\label{eq:vacuum-functional-refinement}
\end{equation}
La compatibilidad no preserva solamente dos n\'umeros: preserva todo el
\'algebra funcional generada por la respuesta. En particular conserva sus
proyectores espectrales, sus potencias, su logaritmo y las formas cuadr\'aticas
que miden los defectos \cref{eq:vacuum-defect34}.

\section{Realizaci\'on dimensional de las relaciones constitutivas}

La Mec\'anica Dimensional determina

\begin{align}
Z_{\rm vac}&=\widehat Z\,u_Su_Q^{-2},\\
c_{\rm vac}&=\widehat c\,u_C,\\
\mu_{\rm vac}&=\widehat\mu\,u_Su_C^{-1}u_Q^{-2},\\
\varepsilon_{\rm vac}&=\widehat\varepsilon\,u_S^{-1}u_C^{-1}u_Q^2.
\label{eq:vacuum-sections}
\end{align}

Las identidades constitutivas se mantienen en las rectas:

\begin{equation}
\mu_{\rm vac}\varepsilon_{\rm vac}=c_{\rm vac}^{-2},
\qquad
\mu_{\rm vac}/\varepsilon_{\rm vac}=Z_{\rm vac}^2.
\end{equation}

La carta SI es posterior. No es necesario declarar \(c\) primitivo y resolver despu\'es \(\mu\varepsilon=c^{-2}\): los tres objetos descienden conjuntamente de \(\mathcal V\) y de las rectas declaradas.

La separaci\'on de niveles se expresa mediante el diagrama
\[
\begin{array}{c}
(A,C^*)\longrightarrow T\longrightarrow\mathcal V
\longrightarrow
(\widehat\mu,\widehat\varepsilon,\widehat Z,\widehat c)
\\[2mm]
\Big\downarrow\text{ cartas de recta}\hspace{36mm}
\Big\downarrow\text{ evaluaci\'on final}
\\[2mm]
(u_S,u_Q,u_C)\longrightarrow
(\mu_{\rm vac},\varepsilon_{\rm vac},Z_{\rm vac},c_{\rm vac})
\longrightarrow\text{ coordenadas SI}.
\end{array}
\]
La primera l\'inea es adimensional y operatoria; la segunda determina tipos de
magnitud; la tercera selecciona una carta metrol\'ogica. Alterar la carta final
no modifica ni los autovalores ni las identidades constitutivas. En ese
sentido preciso, la unificaci\'on matem\'atica--f\'isica consiste en distinguir
el invariante geneal\'ogico de su sistema de coordenadas, y no en suprimir la
estructura dimensional.

\begin{proposition}[Minimalidad espectral]
Sea \(W\) un operador positivo de dos hojas con espectro simple
\(\{r_+,r_-\}\). Si cuatro caracteres positivos satisfacen
\(\mu\varepsilon=c^{-2}\) y \(\mu/\varepsilon=Z^2\), si la propagaci\'on
com\'un se determina mediante \(c^{-1}=\det W=r_-r_+\), y si su orientaci\'on
fija \(Z=r_-/r_+\), entonces necesariamente
\[
(\mu,\varepsilon,Z,c)
=\left(r_-^2,r_+^2,\frac{r_-}{r_+},\frac1{r_-r_+}\right).
\]
Por tanto \cref{eq:vacuum-four} no es una parametrizaci\'on entre muchas:
es la realizaci\'on positiva m\'inima compatible con producto, cociente y
orientaci\'on.
\end{proposition}

\begin{proof}
De \(Z=r_-/r_+\) y \(c^{-1}=\det W=r_-r_+\) se obtienen de forma \`unica las
ra\'ices positivas. Las dos identidades constitutivas fuerzan despu\'es
\(\mu=r_-^2\) y \(\varepsilon=r_+^2\).
\end{proof}

 La primera inversión conceptual consiste en concebir el vacío como una respuesta relacional positiva.\par

En HMT, el vacío es la respuesta dinámica mínima del sistema cuando sus dos hojas deben coexistir, propagarse y cerrar su diferencia conservando la orientación que las distingue.\par

Las dos hojas constituyen dos lecturas conjugadas del mismo antecedente angular. Una conserva una orientación y la otra conserva la orientación contraria. Por eso aparecen dos canales, \(q_+\) y \(q_-\), como las dos caras espectrales correlacionadas de un mismo estado.\par

El operador\par

\[
\mathcal V_{90,120}
=
(I-T^{90})(I-T^{120})^{-1}
\]

compara dos profundidades de propagación: \(90\) y \(120\). Esas profundidades están determinadas por el mismo modo elemental:\par

\[
90=3\cdot30,
\qquad
120=4\cdot30.
\]

Por eso la respuesta de cada hoja puede reescribirse como\par

\[
\frac{1-q^{90}}{1-q^{120}}
=
\frac{1+s+s^2}{1+s+s^2+s^3},
\qquad
s=q^{30}.
\]

Y aquí es donde la interpretación se vuelve realmente poderosa. El numerador contiene tres estados del modo \(30\); el denominador contiene esos mismos tres estados y uno más. El vacío es, por tanto, una completación exacta \(3\to4\).\par

Ese cuarto término integra y conserva los tres anteriores. Los completa. Lo que se añade queda registrado como un residuo positivo de completación:\par

\[
1-\frac{1+s+s^2}{1+s+s^2+s^3}
=
\frac{s^3}{1+s+s^2+s^3}>0.
\]

La memoria del vacío es precisamente el incremento exacto que permite que tres términos se conviertan en cuatro. Constituye la huella positiva y necesaria de la completación.\par

De los dos autovalores de ese único operador salen las cuatro relaciones constitutivas:\par

\[
\widehat\varepsilon=r_+^2,
\qquad
\widehat\mu=r_-^2,
\qquad
\widehat Z=\frac{r_-}{r_+},
\qquad
\widehat c=\frac1{r_+r_-}.
\]

Esto cambia completamente la interpretación de permeabilidad y permitividad.\par

La permitividad es la respuesta cuadrática de una hoja. La permeabilidad es la respuesta cuadrática de la hoja conjugada. La impedancia conserva la relación orientada entre ambas. La velocidad de propagación depende de su producto común.\par

Así, producto y cociente realizan dos tareas diferentes:\par

\begin{itemize}[leftmargin=*,itemsep=0.35em]
\item el producto olvida cuál era cada hoja y conserva la propagación común;
\item el cociente conserva cuál era la orientación relativa entre ellas.
\end{itemize}

El producto responde a la pregunta: «¿qué comparten las dos hojas?». El cociente responde a la pregunta: «¿cómo se distinguen?».\par

Por eso\par

\[
\widehat\mu\,\widehat\varepsilon=\widehat c^{-2}
\]

y\par

\[
\frac{\widehat\mu}{\widehat\varepsilon}
=
\widehat Z^2
\]

son las dos maneras complementarias de leer el mismo par espectral.\par

Lo holográfico aparece aquí de una forma muy limpia: cuatro magnitudes visibles contienen en realidad sólo dos grados de libertad internos. Cualquiera de ciertos pares adecuados permite reconstruir las dos hojas. La publicación constitutiva conserva la genealogía cuando producto y cociente permanecen conjuntamente disponibles.\par

Ésta es la reparación profunda de la ablación de \(\mu\) y \(\varepsilon\). La reparación reúne sus valores con su origen operatorio común.\par

La carta del Sistema Internacional puede transformar después estas secciones intrínsecas en \(\mu_0\), \(\varepsilon_0\), \(Z_0\) y \(c\). Su función consiste en expresar, dentro de una convención metrológica, una estructura cuya procedencia ya ha sido seleccionada por el operador.\par

El vacío deja así de ser un depósito de cuatro constantes. Se convierte en una operación espectral única con dos memorias conjugadas.\par

Y esa visión tiene una consecuencia filosófica importante: el vacío es una relación de compatibilidad entre dos maneras orientadas de realizar un mismo estado.\par

% BEGIN CR28_CONSTITUTIVE_INVERSE
\section{Recuperación funcional, composición y momento de Catalán}
\label{cr28:sec:restitucion-constitutiva}

La respuesta de las incidencias \(90/120\) conserva más información que
un valor escalar aislado. En la misma notación de
\eqref{eq:three-four-completion}, escribimos
\[
 f(s)=\frac{1+s+s^2}{1+s+s^2+s^3},\qquad
 \mathcal D=s\,\frac{d}{ds},\qquad
 \mathcal V=f(T^{30}).
\]
Las marcas \(P_\pm\) mantienen la orientación de los canales. El
transporte angular es el antecedente de la respuesta; su reconstrucción
funcional recupera esta representación sin sustituir el archivo de memoria
del que procede.

\begin{theorem}[Inversión de la respuesta constitutiva]
\label{cr28:thm:inversion-constitutiva}
La función \(f\) es una biyección decreciente de \((0,1)\) sobre
\((3/4,1)\). Su inversa \(\psi\) asigna a \(r\) la única raíz
\(s\in(0,1)\) de
\begin{equation}
 r s^3+(r-1)(1+s+s^2)=0.
 \label{cr28:eq:cubic}
\end{equation}
Con las dos respuestas etiquetadas \(r_\pm\), quedan determinados
\begin{equation}
 s_\pm=\psi(r_\pm),\quad q_\pm=s_\pm^{1/30},\quad
 x=-\frac12\log(q_+q_-),\quad
 y=\frac12\log\frac{q_-}{q_+}.
 \label{cr28:eq:angular-recovery}
\end{equation}
Para el operador completo \(S=\mathcal V^2\), positivo con espectro
en \((9/16,1)\), el cálculo funcional proporciona
\begin{equation}
 T=\bigl[\psi(S^{1/2})\bigr]^{1/30},\qquad
 \det T=\det\bigl[\psi(S^{1/2})\bigr]^{1/30}.
 \label{cr28:eq:operator-recovery}
\end{equation}
\end{theorem}
\begin{proof}
La derivada negativa de \eqref{eq:vacuum-monotonic34} y los límites
\(f(0^+)=1,\ f(1^-)=3/4\) prueban la biyección. Multiplicar \(f(s)=r\)
por su denominador positivo da \eqref{cr28:eq:cubic}. La raíz positiva
de orden \(30\) recupera \(q_\pm\); la suma y la diferencia de
\(\log q_\pm=-x\mp y\) dan las dos coordenadas angulares. Por
positividad, \(S^{1/2}=\mathcal V\). Aplicar \(\psi\) a cada
proyector espectral recupera \(T^{30}\) y la raíz positiva da \(T\).
Tomar su determinante concluye.
\end{proof}

Esta recuperación del operador completo no se extiende a un determinante
aislado. En efecto,
\[
 T_1=\operatorname{diag}(1/5,1/3),\qquad
 T_2=\operatorname{diag}(1/6,2/5)
\]
tienen determinante \(1/15\), pero poseen espectros distintos en la misma
cámara contractiva. La inyectividad demostrada implica
\(f(T_1^{30})^2\ne f(T_2^{30})^2\). Ambos canales y sus marcas
conservan la recuperación; su reducción a un producto pierde información.
% END CR28_CONSTITUTIVE_INVERSE

% BEGIN CR28_TRANSPORT_COMPOSITION
% Recuperación expositiva de INTERRELACIONES_ESTRUCTURALES_VACIO.md, §§6--8.
% Insertar después de la inversión cúbica y antes de las coordenadas logarítmicas.
\subsection{Composición transportada de las respuestas constitutivas}
\label{cr28:comp-sec}

El transporte angular y las incidencias $90$ y $120$ proceden de la
misma estructura APP--TRIT--TPK: las hojas conservan las evaluaciones y
sus acarreos; el régimen trítico conserva la orientación; el calendario
TPK transporta fase y memoria hasta los canales angulares. La respuesta
constitutiva aplica a esos canales la función $f$ de
\eqref{eq:vacuum-monotonic34}. Su inversión permite expresar, en las coordenadas
de respuesta, la composición de transportes del mismo marco. Se conserva
así el orden constructivo: transporte, composición y evaluación espectral.
La operación obtenida actúa sobre esta representación del estado enriquecido;
su inversa recupera los canales, mientras la historia completa de memoria
permanece en la estructura de procedencia.

Sean $s=q^{30}$ y $\psi=f^{-1}$. El exponente $30$ es el máximo común
divisor de las incidencias $90$ y $120$. La inversión cúbica determina
$\psi:(3/4,1)\longrightarrow(0,1)$. Adjuntando la respuesta de frontera
$f(1)=3/4$, extendemos la biyección a
\[
 \psi:D\longrightarrow(0,1],\qquad D=[3/4,1).
\]
En esa frontera se utiliza la función racional $f$; el cociente
$(1-q^{90})/(1-q^{120})$ tiene allí una singularidad removible.

\begin{proposition}[Monoide de respuestas y coordenada aditiva]
\label{cr28:comp-monoid}
La operación
\begin{equation}
 r\odot t=f\bigl(\psi(r)\psi(t)\bigr),\qquad r,t\in D,
 \label{cr28:comp-law}
\end{equation}
hace de $D$ un monoide conmutativo y cancelativo con elemento neutro $3/4$.
La aplicación
\begin{equation}
 \ell(r)=-\frac1{30}\log\psi(r)
 \label{cr28:comp-character}
\end{equation}
es un isomorfismo de ese monoide con $(\mathbb R_{\geq0},+)$.
El intervalo $(3/4,1)$ es cerrado bajo $\odot$; el neutro pertenece a
la extensión de frontera. El único elemento invertible dentro de $D$
es el neutro.
\end{proposition}
\begin{proof}
El producto de dos elementos de $(0,1]$ pertenece al mismo intervalo y
\[
 \psi(r\odot t)=\psi(r)\psi(t).
\]
Para tres respuestas, ambos órdenes de asociación tienen imagen
$\psi(r)\psi(t)\psi(u)$; la inyectividad de $\psi$ prueba la
asociatividad. La conmutatividad procede del producto escalar. Si
$r\odot t=r\odot u$, la positividad de $\psi(r)$ permite cancelar y
obtener $\psi(t)=\psi(u)$, luego $t=u$. El parámetro $1$ corresponde a
$3/4$ y proporciona la identidad. El logaritmo transforma el producto en
suma, y su rango en $(0,1]$ prueba que $\ell$ es una biyección sobre
$\mathbb R_{\geq0}$. En particular, para la respuesta de un canal $q$,
$\ell(r)=-\log q$. Si ambos parámetros son estrictamente menores que
$1$, su producto también lo es. Finalmente, dos números de $(0,1]$
tienen producto $1$ únicamente si ambos son $1$.
\end{proof}

\begin{proposition}[Cancelación admisible y prolongación inversa]
\label{cr28:comp-inverse}
Para $r,u\in D$, la ecuación $r\odot t=u$ tiene solución $t\in D$
exactamente cuando $u\geq r$. Esa solución es única y viene dada por
\begin{equation}
 t=f\!\left(\frac{\psi(u)}{\psi(r)}\right).
 \label{cr28:comp-cancellation}
\end{equation}
La misma función $f$ es una biyección de $(0,\infty)$ sobre $(0,1)$.
Con esta extensión de $\psi$, la ley~\eqref{cr28:comp-law} transporta el
grupo multiplicativo de los parámetros positivos al intervalo $(0,1)$.
La inversión de grupo satisface
\begin{equation}
 r^{\ominus}=\psi(r)r,\qquad
 r\odot r^{\ominus}=\frac34.
 \label{cr28:comp-group-inverse}
\end{equation}
\end{proposition}
\begin{proof}
La ecuación equivale a
$\psi(t)=\psi(u)/\psi(r)$. Este cociente pertenece a $(0,1]$
exactamente cuando $\psi(u)\leq\psi(r)$, que por monotonía decreciente
equivale a $u\geq r$. La biyección prueba existencia y unicidad.
La derivada de \eqref{eq:vacuum-monotonic34} es negativa para todo $s>0$;
los límites de $f$ en $0$ y en infinito son $1$ y $0$, respectivamente.
Esto establece la biyección extendida. Multiplicando numerador y
denominador por $s^3$ se obtiene
\[
 f(s^{-1})=\frac{s^3+s^2+s}{s^3+s^2+s+1}=s f(s).
\]
La inversa multiplicativa de $s=\psi(r)$ tiene, por tanto, respuesta
$\psi(r)r$. Su composición con $r$ corresponde al parámetro $1$.
\end{proof}

La prolongación a $s>1$, o equivalentemente a $q>1$, es la extensión
algebraica del lector. El sector físico contractivo del capítulo conserva
su dominio declarado. La coordenada $\ell$ se prolonga a un isomorfismo
con $(\mathbb R,+)$ y cambia de signo bajo la inversión de grupo.

\subsubsection{Realización operatoria en un marco de hojas común}

\begin{proposition}[Composición con proyectores marcados]
\label{cr28:comp-operator}
Sean $\mathcal R=\mathcal R^*$, $\mathcal R^2=I$ y
$P_\pm=(I\pm\mathcal R)/2$. Para $j=a,b$, sean
\[
 T_j=q_{j,+}P_++q_{j,-}P_-,\qquad
 \mathcal V_j=f(T_j^{30}),\qquad 0<q_{j,\pm}<1.
\]
La composición $T_{a\circ b}=T_aT_b$ satisface
\begin{equation}
 \mathcal V_{a\circ b}
 =f\bigl(\psi(\mathcal V_a)\psi(\mathcal V_b)\bigr).
 \label{cr28:comp-operator-law}
\end{equation}
Sus dos respuestas son $r_{a,\pm}\odot r_{b,\pm}$. Si
$T_j=\exp[-(x_jI+y_j\mathcal R)]$ con $x_j>|y_j|$, entonces
\begin{equation}
 T_aT_b=
 \exp[-((x_a+x_b)I+(y_a+y_b)\mathcal R)],
 \label{cr28:comp-angular-addition}
\end{equation}
y $x_a+x_b>|y_a+y_b|$.
\end{proposition}
\begin{proof}
La ortogonalidad $P_+P_-=0$ da
\[
 T_aT_b=q_{a,+}q_{b,+}P_++q_{a,-}q_{b,-}P_-.
\]
En particular, $T_a$ y $T_b$ conmutan. El cálculo funcional en estos
proyectores proporciona
$\psi(\mathcal V_j)=T_j^{30}$ y
$(T_aT_b)^{30}=T_a^{30}T_b^{30}$; aplicar $f$ prueba
\eqref{cr28:comp-operator-law}. Para toda función $g$ definida en los
dos autovalores se tiene, más explícitamente,
\[
 P_\pm g(T_j)=g(q_{j,\pm})P_\pm.
\]
Como $q_{j,\pm}=\exp[-(x_j\pm y_j)]$, sus productos prueban
\eqref{cr28:comp-angular-addition}. La desigualdad triangular concluye
$x_a+x_b>|y_a|+|y_b|\geq|y_a+y_b|$.
\end{proof}

Las coordenadas recuperadas pueden escribirse directamente como
\[
 x=\frac{\ell(r_+)+\ell(r_-)}2,\qquad
 y=\frac{\ell(r_+)-\ell(r_-)}2.
\]
La composición suma estas parejas. El intercambio simultáneo de hojas
permuta los dos canales de cada transporte, actúa como
$(x,y)\mapsto(x,-y)$ y es un automorfismo de la ley componente a
componente. La inversión de ambos transportes en la extensión algebraica
actúa como $(x,y)\mapsto(-x,-y)$. Las dos involuciones conmutan y
conservan funciones distintas: una intercambia las marcas de orientación;
la otra invierte el transporte.

Para representar el intercambio mediante conjugación se utiliza un
operador $L$ que satisfaga $L\mathcal R L^{-1}=-\mathcal R$. En la
representación
\[
 \mathcal R=\begin{pmatrix}0&1\\1&0\end{pmatrix},\qquad
 L=\begin{pmatrix}1&0\\0&-1\end{pmatrix},
\]
la multiplicación directa da $LP_\pm L^{-1}=P_\mp$ y, por tanto,
$LT_jL^{-1}=q_{j,-}P_++q_{j,+}P_-$. La conjugación por
$\mathcal R$ conserva sus propios proyectores. Una traza que reúne los
canales tampoco reemplaza a los proyectores en esta identidad funcional:
la información de sus etiquetas interviene en la composición.

La hipótesis de marco común es sustantiva. Por ejemplo, los operadores
positivos definidos
\[
 A_0=\begin{pmatrix}1/2&0\\0&1/3\end{pmatrix},\qquad
 B_0=\begin{pmatrix}1/2&1/6\\1/6&1/2\end{pmatrix}
\]
tienen autovalores positivos, pero
\[
 A_0B_0=\begin{pmatrix}1/4&1/12\\1/18&1/6\end{pmatrix}
\]
carece de autoadjunción. La proposición especifica la composición de
transportes en los proyectores comunes; su dominio no comprende una mezcla
arbitraria de medios materiales. La realización de una concatenación de
rutas HMT conserva además sus datos de memoria: el cierre algebraico del
cono de parámetros describe la representación y no selecciona por sí solo
nuevas rutas primarias.

\subsubsection{Producto de canales y composición de transportes}

La velocidad normalizada de un estado sigue determinada por
\[
 \widehat c=(r_+r_-)^{-1}.
\]
Aquí el producto ordinario reúne las dos respuestas de ese mismo estado.
La ley $\odot$ evalúa otra operación: la composición de dos transportes
en cada hoja. El siguiente cálculo exacto distingue sus argumentos:
\begin{equation}
 f(1/2)=\frac{14}{15},\qquad
 \frac{14}{15}\odot\frac{14}{15}=f(1/4)=\frac{84}{85}
 \ne\frac{196}{225}=\left(\frac{14}{15}\right)^2.
 \label{cr28:comp-rational-example}
\end{equation}
Las fracciones de este ejemplo son parámetros algebraicos de comprobación.
La expresión de $\widehat c$ y las identidades constitutivas anteriores
permanecen inalteradas. El contenido añadido es una ley explícita para
transportar la composición y recuperar su coordenada angular aditiva.

% END CR28_TRANSPORT_COMPOSITION

% BEGIN CR28_CATALAN_RECOVERY
\subsection{Restitución funcional del momento orientado}
\label{cr28:sec:catalan-recovery}

La relación diferencial~\eqref{cr28:eq:catalan-problem} admite una
inversa que conserva la familia completa de momentos. Su dominio es
la función constitutiva producida por los sectores del ciclo, con la
orientación de~\eqref{eq:vacuum-angular-channels}. Las incidencias $90=3\cdot30$ y
$120=4\cdot30$ fijan esa función antes de evaluarla en los canales
$s_\pm=q_\pm^{30}$. La condición en el origen corresponde a la
desaparición del momento cuando se extingue su argumento contractivo.
Esta condición determina también las constantes de integración.

\begin{theorem}[Restitución integral y unicidad]
\label{cr28:thm:catalan-recovery}
Sea $f(s)=(1-s^3)/(1-s^4)$, para $0<s<1$, la respuesta
de~\eqref{eq:three-four-completion}, y sea $\mathcal D=s\,d/ds$.
Existe una única función $F\in C^2((0,1))$ que satisface
\begin{equation}
 \mathcal D^2F(s)=\frac{s(1+s)}{1+s+s^2}f(s),
 \qquad \lim_{s\downarrow0}F(s)=0.
 \label{cr28:eq:catalan-problem}
\end{equation}
Está dada por
\begin{align}
 F(s)&=\int_0^s\log\frac{s}{t}\,
          \frac{1+t}{1+t+t^2}f(t)\,dt
       =\int_0^s\frac{\log(s/t)}{1+t^2}\,dt
       =\mathcal C_2(s),
 \label{cr28:eq:catalan-integral}\\
 f(s)&=\frac{1+s+s^2}{s(1+s)}\mathcal D^2F(s).
 \label{cr28:eq:catalan-inverse}
\end{align}
La extensión continua a $s=1$ recupera $F(1)=G_{\rm Cat}$.
\end{theorem}
\begin{proof}
La factorización $1+t+t^2+t^3=(1+t)(1+t^2)$ da
\[
 \frac{1+t}{1+t+t^2}f(t)=\frac1{1+t^2}.
\]
El integrando de~\eqref{cr28:eq:catalan-integral} es
integrable en cero, porque
\[
 0\leq F(s)\leq\int_0^s\log(s/t)\,dt=s.
\]
Por tanto $F(s)\to0$. La diferenciación en el intervalo abierto,
primero en un intervalo truncado y después por convergencia dominada,
produce
\[
 \mathcal DF(s)=\int_0^s\frac{dt}{1+t^2},
 \qquad \mathcal D^2F(s)=\frac{s}{1+s^2}
     =\frac{s(1+s)}{1+s+s^2}f(s).
\]
En particular, $F$ es de clase $C^2$ y satisface la ecuación.
Si $F_1,F_2$ son dos soluciones, la diferencia $H=F_1-F_2$
satisface $\mathcal D^2H=0$. Mediante $z=\log s$ se obtiene
$d^2H(e^z)/dz^2=0$, luego $H(s)=a\log s+b$.
La existencia del límite nulo cuando $s\downarrow0$ obliga primero
a $a=0$ y después a $b=0$. La condición sobre $F$ determina, pues,
la solución y también el límite $\mathcal DF(s)\to0$.

Para $s<1$, el desarrollo geométrico de $(1+t^2)^{-1}$ y la
integrabilidad absoluta de sus términos ponderados permiten integrar:
\[
 \int_0^s t^{2k}\log(s/t)\,dt
     =\frac{s^{2k+1}}{(2k+1)^2},\qquad
 F(s)=\sum_{k\geq0}\frac{(-1)^ks^{2k+1}}{(2k+1)^2}.
\]
La serie representa la familia integral~\eqref{cr28:eq:catalan-integral}. La convergencia
absoluta y uniforme de esta última serie en $[0,1]$ autoriza el
paso al extremo $s=1$. También lo autoriza en la integral la cota
integrable $\log(1/t)$, extendiendo el integrando por cero para
$t>s$. Finalmente, despejar $f$ en la ecuación diferencial da
\eqref{cr28:eq:catalan-inverse}, con denominadores
positivos en $(0,1)$.
\end{proof}

El límite del momento y el límite de su derivada segunda se toman
en sus respectivas clases de convergencia. En particular,
\[
 \lim_{s\uparrow1}\mathcal D^2\mathcal C_2(s)=\frac12
\]
es el límite radial de Abel de la serie diferenciada. La evaluación
de $\mathcal C_2(1)$ utiliza directamente su serie absolutamente
convergente. Esta distinción conserva el mismo operador de profundidad
en el interior y en su lectura de frontera.

\subsection{Compatibilidad de la restitución con canales y velocidad}
\label{cr28:sec:channels-velocity}

El teorema restituye una función a partir de otra función previamente
construida. La inversión cúbica~\eqref{cr28:eq:cubic} actúa, en cambio,
sobre las evaluaciones $r_\pm$: recupera los argumentos $s_\pm$
dentro de esa familia fijada. La composición de ambas operaciones
conserva la orientación del cuarto de giro, las etiquetas $P_\pm$
y las incidencias del transporte. Así, $G_{\rm Cat}=\mathcal C_2(1)$
es la lectura extrema de un objeto funcional que participa también
en las respuestas interiores $r_\pm=f(s_\pm)$.

Para $\mathcal V_m=\mathcal V\otimes I_{9^m}$ y
$\tau_m=\operatorname{Tr}/(2\cdot9^m)$, la lectura de velocidad
mantiene exactamente la normalización de
\eqref{eq:vacuum-traces}:
\begin{equation}
 \exp[-2\tau_m(\log\mathcal V_m)]
   =\frac1{r_+r_-}
   =(\det\mathcal V)^{-1}=\widehat c.
 \label{cr28:eq:velocity-normalized}
\end{equation}
En efecto, cada autovalor $r_\pm$ se repite $9^m$ veces, de modo
que $2\tau_m(\log\mathcal V_m)=\log r_++\log r_-$.
La identidad utiliza el determinante de la realización de dos hojas
y la traza normalizada de su amplificación. Las secciones
dimensionales ya fijadas dan entonces
\[
 c_{\rm int}=c_{\rm vac}=\widehat c\,u_C,
 \qquad \ell_0=c_{\rm vac}t_0,
\]
con el mismo estado, la misma carta y la misma duración elemental
de~\eqref{cr28:eq:same-velocity}.
La restitución funcional conserva esta sección de velocidad al
componerla con la acción y la longitud; la multiplicidad espectral
queda absorbida por la normalización declarada.

% END CR28_CATALAN_RECOVERY

% BEGIN CR28_CONSTITUTIVE_RADIAL_LINK
\subsection{Una misma sección constitutiva en acción, longitud y gravedad}
\label{cr28:sec:constitutive-radial}
Las relaciones \eqref{eq:vacuum-sections} y la restitución anterior
mantienen una única sección de velocidad:
\begin{equation}
 c_{\rm int}=c_{\rm vac}=\widehat c\,u_C,\qquad
 \ell_0=c_{\rm int}t_0,\qquad
 c_{\rm int}=(\mu_{\rm vac}\varepsilon_{\rm vac})^{-1/2}.
 \label{cr28:eq:same-velocity}
\end{equation}
La última igualdad es una identidad de secciones en bases dimensionales
compatibles. La demostración radial del
\cref{g26:sec:rigidez-radial-es} construye primero
\(L=(5759/23040)\alpha^{16}L_*\) y después determina
\begin{equation}
 G_{\rm ret}=\frac{c_{\rm int}^3L^2}{\hbar_{\rm ret}}
 =\frac{L^2}
 {\hbar_{\rm ret}(\mu_{\rm vac}\varepsilon_{\rm vac})^{3/2}}.
 \label{cr28:eq:constitutive-gravity}
\end{equation}
La segunda expresión resulta de sustituir la identidad constitutiva en
la primera; no introduce una nueva calibración de ninguno de los dos
canales. En coordenadas, si \(c_{\rm int}=c_*\widehat c\), entonces
\(G_{\rm ret}=c_*^3\widehat c^{\,3}L^2/\hbar_{\rm ret}\).
El factor de propagación interviene una sola vez. La recuperación de
Catalán, el transporte compuesto y la respuesta radial conservan así
el mismo origen angular, sus marcas y su normalización.
% END CR28_CONSTITUTIVE_RADIAL_LINK



  \endgroup

\chapter{Carga, resistencia de von Klitzing, conductancia, constante de Josephson y flujo magnético}
\begingroup
\renewcommand{\chapter}[2][]{}%
\chapter{Derivación de la carga y de los cuantos eléctricos: resistencia de von Klitzing, conductancia, constante de Josephson y flujo magnético}
\label{chap:p3-final:40:cuantos_electricos}
\label{ch:metrology}

\section{Derivaci\'on de la secci\'on de carga a partir del vac\'io}

Fijados \(Z_{\rm vac}\), \(h_a=2\pi\hbar_a\) y \(\alpha\), la carga de la hoja \(a\in\{\mathrm{pre},\mathrm{ret}\}\) es

\begin{equation}
e_a=\sqrt{\frac{2\alphah_a}{Z_{\rm vac}}}.
\label{eq:charge}
\end{equation}

La secci\'on \(e_a\) no se introduce mediante una coordenada del Sistema
Internacional destinada a reconstruir \(\alpha\). El orden causal procede de
\(\alpha\), de la acci\'on y del vac\'io hacia la secci\'on de carga.

\section{Constantes de von Klitzing y Josephson, conductancia y flujo magn\'etico}

De \cref{eq:charge} se deducen

\begin{align}
R_K&=\frac{h_a}{e_a^2}=\frac{Z_{\rm vac}}{2\alpha},
\label{eq:rk}\\
G_0&=\frac{2e_a^2}{h_a}=\frac{4\alpha}{Z_{\rm vac}},
\label{eq:g0}\\
\Phi_{0,a}&=\frac{h_a}{2e_a}
=\sqrt{\frac{h_aZ_{\rm vac}}{8\alpha}},
\label{eq:phi0}\\
K_{J,a}&=\frac{2e_a}{h_a}=\Phi_{0,a}^{-1}.
\label{eq:kj}
\end{align}

En consecuencia,

\begin{equation}
R_KG_0=2,\qquad G_0Z_{\rm vac}=4\alpha,
\qquad K_{J,a}\Phi_{0,a}=1.
\label{eq:quantum-electrical-identities}
\end{equation}

La covariancia bidireccional adopta la forma

\[
R_K^{\rm pre}=R_K^{\rm ret},\qquad
G_0^{\rm pre}=G_0^{\rm ret},
\]

\[
\frac{\Phi_{0,\rm pre}}{\Phi_{0,\rm ret}}=R_{\rm act}^{1/2},
\qquad
\frac{K_{J,\rm pre}}{K_{J,\rm ret}}=R_{\rm act}^{-1/2}.
\]

Las relaciones de resistencia y conductancia son independientes de la
orientaci\'on de la acci\'on. El par flujo--frecuencia conserva dicha
orientaci\'on y transforma rec\'iprocamente sus dos componentes.

  \endgroup

\chapter{Constantes térmicas y radiación: Boltzmann, Wien y Stefan--Boltzmann}
\begingroup
\renewcommand{\chapter}[2][]{}%
\chapter{Constantes térmicas y radiación: Boltzmann, Wien, Stefan--Boltzmann y termodinámica espectral del gas fotónico}
\label{chap:p3-final:48:constantes_termicas}
\section{Relaci\'on t\'ermica intr\'inseca de la escala \(108\)}

La estructura cronol\'ogica \(108\) fija

\begin{equation}
k_B\Theta_{\rm clk}\log3
=\frac{h_{\rm int}}{108t_0}
=\frac{2\pi\hbar_{\rm int}}{108t_0}.
\label{eq:boltzmann-clock}
\end{equation}

Esta igualdad no determina por separado una coordenada num\'erica de \(k_B\) y otra de \(\Theta_{\rm clk}\) sin una normalizaci\'on adicional. S\'i determina su producto como secci\'on de energ\'ia. El factor \(\log3\) tiene tres antecedentes compatibles:

\begin{equation}
\log3=C_1+C_2-2C_0
=\frac{S_{\TRIT}}{k_B}
=\frac{E_P}{k_B\Theta_{\rm clk}}.
\label{eq:log3-triple}
\end{equation}

Las tres igualdades corresponden a realizaciones residual, informacional y
termodin\'amica de un mismo escalar. La coincidencia del valor no identifica
sus dominios respectivos.

\section{Formulaciones espectrales de la ley de desplazamiento de Wien}

Para la densidad espectral por longitud de onda, el m\'aximo satisface

\[
5(1-\ee^{-x_5})=x_5,
\qquad
x_5=5+W_0(-5\ee^{-5}).
\]

Con \cref{eq:boltzmann-clock},

\begin{equation}
\lambda_{\max}(\Theta_{\rm clk})
=\frac{108\log3}{x_5}\,\ell_0.
\label{eq:wien-lambda}
\end{equation}

Para la densidad por frecuencia,

\[
3(1-\ee^{-x_3})=x_3,
\qquad
x_3=3+W_0(-3\ee^{-3}),
\]

y

\begin{equation}
\nu_{\max}(\Theta_{\rm clk})
=\frac{x_3}{108t_0\log3}.
\label{eq:wien-frequency}
\end{equation}

Los m\'aximos no son rec\'iprocos, puesto que las densidades espectrales se
transforman con el jacobiano correspondiente. El cociente \(x_5/x_3\)
pertenece a la estructura de la parametrizaci\'on y no representa una
inconsistencia.

\section{Ley de Stefan--Boltzmann y termodin\'amica del gas de fotones}

Evaluar la ley de radiaci\'on en \(\Theta_{\rm clk}\) produce el flujo intr\'inseco

\begin{equation}
j_\star(\Theta_{\rm clk})
=\frac{2\pi^5}{15\,108^4(\log3)^4}\,
\frac{h_{\rm int}}{\ell_0^2t_0^2},
\label{eq:stefan-hmt}
\end{equation}

en la normalizaci\'on declarada. La constante de Stefan--Boltzmann no se inserta como primitiva: resulta de integrar el espectro bos\'onico una vez fijados acci\'on, propagaci\'on y temperatura.
Las leyes espectrales empleadas son las de Planck y Wien \cite{Planck,Wien}; HMT fija aqu\'i la secci\'on sobre la que se eval\'uan, no sus constantes por contraste retrospectivo.

En unidades de la escala de Planck,

\begin{align}
n_\gamma\ell_P^3
&=\frac{2\zeta(3)}{\pi^2\log^3 3},
\label{eq:photon-number}\\
\frac{u_\gamma\ell_P^3}{E_P}
&=\frac{\pi^2}{15\log^4 3},
\label{eq:photon-energy}\\
\frac{s_\gamma\ell_P^3}{k_B}
&=\frac{4\pi^2}{45\log^3 3}.
\label{eq:photon-entropy}
\end{align}

\(\zeta(3)\) representa el momento bos\'onico de orden tres, mientras que
\(\log3\) determina la energ\'ia de decisi\'on TRIT. La ecuaci\'on conserva la
distinci\'on entre ambos invariantes.

La conductancia t\'ermica cu\'antica se expresa como

\begin{equation}
\frac{\kappa_Qt_0}{k_B}=\frac{\pi^2}{324\log3}.
\label{eq:thermal-conductance}
\end{equation}

\section{Confluencia metrológica de los invariantes térmicos y radiativos}

\begin{longtable}{@{}p{0.18\textwidth}p{0.27\textwidth}p{0.45\textwidth}@{}}
\toprule
Objeto & Ecuaci\'on interna & Significado estructural\\
\midrule
\endhead
\(R_K\) & \(Z_{\rm vac}/(2\alpha)\) & resistencia cu\'antica bidireccional invariante\\
\(G_0\) & \(4\alpha/Z_{\rm vac}\) & conductancia complementaria\\
\(\Phi_0\) & \(\sqrt{hZ/(8\alpha)}\) & flujo que conserva orientaci\'on de acci\'on\\
\(K_J\) & \(\Phi_0^{-1}\) & relaci\'on rec\'iproca entre frecuencia y flujo\\
\(E_P\) & \(k_B\Theta\log3\) & energ\'ia del ciclo y contenido informacional TRIT\\
Wien & \(108\log3/x_5\) & extremo espectral en la escala cronol\'ogica\\
Ap\'ery fot\'onica & \(2\zeta(3)/(\pi^2\log^3 3)\) & densidad de n\'umero bos\'onica\\
\bottomrule
\end{longtable}

\begin{remark}[Control negativo]
El bloque queda refutado si un cambio de variable espectral conserva
indebidamente el mismo extremo de Wien, si \(\zeta(3)\) sustituye a
\(\log3\), o si se asignan valores separados a \(k_B\) y \(\Theta\) sin
especificar una carta. Las identidades
\cref{eq:quantum-electrical-identities,eq:planck-trit} proporcionan controles
algebraicos directos.
\end{remark}

\subsection*{Alcance lógico y dominio de validez}
Cuantos el\'ectricos y confluencia Planck--TRIT: \emph{Teorema interno exacto}. Wien, Stefan--Boltzmann y gas fot\'onico: \emph{Composici\'on matem\'atica exacta} con el continuo bos\'onico declarado. La carta SI permanece como \emph{Contraste metrol\'ogico externo}.

  \endgroup

\chapter{Constantes electrodébiles, relación \texorpdfstring{\(W/Z\)}{W/Z} y mezcla angular}
\begingroup
\renewcommand{\chapter}[2][]{}%
\chapter{Constantes electrodébiles y mezcla: determinante angular, relación W/Z, ángulo de Weinberg y reflexión racional CKM}
\label{chap:p3-final:49:sector_electrodebil}
\section{Determinante angular y relaci\'on electrod\'ebil de Weinberg}

El determinante de \(T\) es

\begin{equation}
\det T=q_+q_-=D_A.
\label{eq:detT}
\end{equation}

El sector constitutivo del vac\'io eval\'ua \(F(T)^2\), donde
\(F(z)=(1-z^{90})/(1-z^{120})\), mientras que el sector angular
electrod\'ebil eval\'ua \(\det T\). Se trata de dos funciones de un
antecedente com\'un y no de una factorizaci\'on horizontal entre constantes.
Esta distinci\'on se desarrolla en el \cref{ch:ew}.

\begin{remark}[Control negativo]
El bloque queda refutado si se viola cualquiera de estas identidades exactas:
positividad de \(I-T^{120}\), igualdad
\cref{eq:three-four-completion}, relaciones \cref{eq:vacuum-relations} o
compatibilidad \(E_m(\mathcal V_{m+1})=\mathcal V_m\). Una coincidencia SI
no puede compensar un fallo operatorio.
\end{remark}

\subsection*{Alcance lógico y dominio de validez}
Operador, espectro, completaci\'on \(3\to4\), paridad, amplificaci\'on y secciones constitutivas: \emph{Teorema interno exacto}. La formulaci\'on que los re\'une en un solo operador es \emph{Formalizaci\'on nueva}; los datos angulares y las rectas son \emph{Resultado antecedente demostrado} del desarrollo HMT precedente y se declaran aqu\'i como estructura de partida expl\'icita.

 \label{ch:ew}

\section{Funciones constitutiva y determinantal del operador angular}

Recordemos el operador de dos hojas del \cref{ch:vacuum},
\begin{equation}
\label{eq:ew-angular-operator}
T=q_+P_+ +q_-P_-,
\qquad
q_-=\ee^{-(A_{\rm rad}-C^*_{\rm rad})},
\qquad
q_+=\ee^{-(A_{\rm rad}+C^*_{\rm rad})}.
\end{equation}
El vac\'io electromagn\'etico aplica la funci\'on
\begin{equation}
\label{eq:ew-vacuum-reader}
F(T)^2,
\qquad
F(z)=\frac{1-z^{90}}{1-z^{120}},
\end{equation}
mientras el sector angular electrod\'ebil aplica el determinante
\begin{equation}
\label{eq:ew-determinant-reader}
\boxed{
\det T=q_+q_-=\ee^{-2A_{\rm rad}}=:D_A.}
\end{equation}
Son dos realizaciones funcionales del antecedente com\'un \((A,C^*,T)\).
La generaci\'on parte de \(T\); la recuperaci\'on de esa representaci\'on
desde el operador constitutivo completo es, adem\'as, posible en el
dominio contractivo. Por \cref{cr28:thm:inversion-constitutiva},
\[
 T=\bigl[\psi((F(T)^2)^{1/2})\bigr]^{1/30},
\]
y, por tanto, tambi\'en se recupera \(\det T\). El determinante aislado
no reconstruye los dos canales: el contraejemplo positivo del mismo
teorema conserva su producto y modifica sus respuestas. La unidad
estructural reside en \(T\), sus proyectores y la memoria de procedencia;
la recuperaci\'on funcional conserva ese orden y no identifica entre s\'i
las diferentes lecturas de permeabilidad, permitividad y \(\theta_W\).

En la involuci\'on de hojas \(C^*\mapsto-C^*\), los autovalores
\(q_+,q_-\) se intercambian. Por ello \(D_A\) y la velocidad constitutiva
son pares, mientras \((\widehat\mu,\widehat\varepsilon)\) se intercambia y
\(\widehat Z\) se invierte. La relaci\'on electrod\'ebil que sigue usa el
invariante par.

\section{Relaci\'on exacta entre los sectores \(W\) y \(Z\)}

La coordenada angular de duraci\'on--perduraci\'on es
\begin{equation}
\label{eq:ew-angular-character}
\chi_{\rm dur}(n_A,s_C)
=\exp\!\left(n_AA_{\rm rad}
+\frac{s_C}{6}C^*_{\rm rad}\right).
\end{equation}
Esta coordenada no se introduce como una etiqueta aislada. Su procedencia es
\begin{equation}
\label{eq:ew-signature-provenance}
\begin{aligned}
\APP&\longrightarrow\TRIT\longrightarrow\TPK
\longrightarrow\text{extensión superviviente}
\longrightarrow\text{ruta}\\
&\longrightarrow\text{ledger}
\longrightarrow(n_A,s_C,\ldots)
\longrightarrow\chi_{\rm dur}.
\end{aligned}
\end{equation}
Las firmas de \(W\), \(Z\) y \(H\) son magnitudes tipadas obtenidas mediante
esa cadena; el presente cap\'itulo las adopta sin seleccionarlas a partir de
sus masas terminales. En la realizaci\'on angular establecida, los sectores
cargado y neutro satisfacen
\begin{equation}
\label{eq:ew-WZ-signatures}
(n_A,s_C)_W=(94,-1),
\qquad
(n_A,s_C)_Z=(95,-1).
\end{equation}
Esta afirmaci\'on se refiere a la componente angular com\'un. Las
coordenadas adicionales de refinamiento de la ley general de masas no se
eliminan salvo que el transductor f\'isico demuestre su cancelaci\'on.

\begin{theorem}[Defecto de Weinberg en el esquema de masas f\'isicas de la carta angular]
\label{thm:ew-weinberg}
Si \(W\) y \(Z\) se eval\'uan en una misma secci\'on de escala y comparten
el refinamiento que acompa\~na a \cref{eq:ew-WZ-signatures}, entonces
\begin{equation}
\label{eq:ew-WZ-ratio}
\boxed{
\frac{m_W^{\angle}}{m_Z^{\angle}}
=\ee^{-A_{\rm rad}}=\sqrt{D_A},
\qquad
\left(\frac{m_W^{\angle}}{m_Z^{\angle}}\right)^2=D_A.}
\end{equation}
En el esquema de masas f\'isicas,
\begin{equation}
\label{eq:ew-sin2thetaW}
\boxed{
\sin^2\theta_W^{\rm OS,HMT}=1-D_A
=0.2248708807528435698111159035\ldots .}
\end{equation}
Adem\'as,
\begin{equation}
\label{eq:ew-WZ-terminal-ratio}
\frac{m_W^{\angle}}{m_Z^{\angle}}
=0.8804141748331613670128885459\ldots .
\end{equation}
\end{theorem}

\begin{proof}
Por \cref{eq:ew-angular-character,eq:ew-WZ-signatures},
\[
\frac{\chi_{\rm dur}(94,-1)}{\chi_{\rm dur}(95,-1)}
=\exp(-A_{\rm rad});
\]
la coordenada \(C^*\) se cancela porque ambos sectores comparten
\(s_C=-1\). Al elevar al cuadrado y usar
\cref{eq:ew-determinant-reader} se obtiene \(D_A\). La definici\'on
de masas f\'isicas \(s_W^2=1-m_W^2/m_Z^2\) produce
\cref{eq:ew-sin2thetaW}.
\end{proof}

El contenido principal es la identidad entre dos construcciones
internos: el determinante del operador angular y la diferencia de un paso
en la realizaci\'on \(W/Z\). La promoci\'on a selector f\'isico de calibre exige
un entrelazador que lleve las hojas de \(T\) a la base neutra
\((W^3,B)\), y que demuestre qu\'e refinamientos se conservan o cancelan.

\section{Acoplamientos de calibre en normalizaci\'on racionalizada}

Definimos el acoplamiento adimensional de calibre
\begin{equation}
\label{eq:ew-egauge}
e_g:=\sqrt{4\pi\alpha}
=0.3028221208717526427662442689\ldots .
\end{equation}
No debe confundirse \(e_g\) con una secci\'on dimensional de carga
\(e_{\rm int}\): la primera es un acoplamiento en unidades naturales; la
segunda vive en la recta de carga de la Mec\'anica Dimensional.

Sea
\begin{equation}
\label{eq:ew-sw-cw}
s_W^2=1-D_A,
\qquad
c_W^2=D_A.
\end{equation}
La convenci\'on electrod\'ebil racionalizada a nivel de \`arbol da
\begin{align}
g&=\frac{e_g}{s_W}
=2\sqrt{\frac{\pi\alpha}{1-D_A}}
=0.6385883427334574283647003809\ldots,
\label{eq:ew-g}\\
g'&=\frac{e_g}{c_W}
=2\sqrt{\frac{\pi\alpha}{D_A}}
=0.3439541633108502429362319257\ldots,
\label{eq:ew-gprime}\\
\frac{g'}{g}
&=\sqrt{D_A^{-1}-1}
=0.5386164141966093594996610933\ldots .
\label{eq:ew-gprime-over-g}
\end{align}

\begin{proposition}[Relaci\'on custodial de \`arbol]
\label{prop:ew-rho}
En la misma interfaz,
\begin{equation}
\label{eq:ew-rho}
\rho_{\rm EW}^{(0)}
:=\frac{(m_W^{\angle})^2}
{(m_Z^{\angle})^2c_W^2}=1.
\end{equation}
\end{proposition}

\begin{proof}
Por \cref{eq:ew-WZ-ratio,eq:ew-sw-cw}, tanto el numerador relativo como
\(c_W^2\) son \(D_A\).
\end{proof}

Las ecuaciones \cref{eq:ew-g,eq:ew-gprime} son composiciones exactas una
vez declarados el esquema de masas f\'isicas, la racionalizaci\'on y la escala. No
afirman que un cambio de esquema de renormalizaci\'on deje inalterados sus
valores num\'ericos.

\section{Constante de Fermi en aproximaci\'on de \`arbol}

En unidades naturales, la integraci\'on del propagador cargado a baja
energ\'ia produce
\begin{equation}
\label{eq:ew-GF-def}
\frac{G_F^{(0)}}{\sqrt2}=\frac{g^2}{8m_W^2}.
\end{equation}
Sustituyendo \cref{eq:ew-g} se obtiene la composici\'on
\begin{equation}
\label{eq:ew-GF}
\boxed{
G_F^{(0)}
=\frac{\pi\alpha}
{\sqrt2(1-D_A)m_W^2}.}
\end{equation}
Al adoptar la masa obtenida previamente
\begin{equation}
\label{eq:ew-W-output}
m_W^{\HMT}=80.381592550221\ldots\ \mathrm{GeV},
\end{equation}
la evaluaci\'on terminal es
\begin{equation}
\label{eq:ew-GF-value}
\boxed{
G_F^{(0)}
=1.1157162817659203186621784\times10^{-5}\ \mathrm{GeV}^{-2}.}
\end{equation}

La diferencia respecto de una constante de Fermi renormalizada no se usa
para elegir ning\'un coeficiente. En la convenci\'on
\begin{equation}
\label{eq:ew-Delta-r}
G_F
=\frac{\pi\alpha}
{\sqrt2(1-D_A)m_W^2(1-\Delta r)},
\end{equation}
la obligaci\'on restante es derivar \(\Delta r\) desde el sector interno de
bucles, estados y escala. Abrir el valor observado de \(G_F\) para definir
\(\Delta r\) ser\'ia una reconstrucci\'on calibrada, no una predicci\'on.

\section{Dependencia orientada del acoplamiento de Higgs}

En la carta angular, las rutas pertinentes son
\begin{equation}
\label{eq:ew-HW-signatures}
(n_A,s_C)_H=(97,9),
\qquad
(n_A,s_C)_W=(94,-1).
\end{equation}
Por tanto,
\begin{equation}
\label{eq:ew-HW-ratio}
\boxed{
\frac{m_H^{\angle}}{m_W^{\angle}}
=\frac{\chi_{\rm dur}(97,9)}{\chi_{\rm dur}(94,-1)}
=\exp\!\left(3A_{\rm rad}+\frac53C^*_{\rm rad}\right)
=1.55872087212325548564\ldots .}
\end{equation}
Aqu\'i \(C^*\) no se cancela: el sector escalar conserva la orientaci\'on
que el determinante \(D_A=\ee^{-2A_{\rm rad}}\) no distingue.

Tomemos la convenci\'on del potencial
\begin{equation}
\label{eq:ew-Higgs-potential}
V(H)=-\mu_H^2H^\dagger H
+\lambda_H(H^\dagger H)^2,
\end{equation}
para la cual
\begin{equation}
\label{eq:ew-Higgs-tree-relations}
m_H^2=2\lambda_Hv^2,
\qquad
m_W^2=\frac{g^2v^2}{4}.
\end{equation}

\begin{corollary}[Autoacoplamiento orientado de Higgs]
\label{cor:ew-Higgs-lambda}
En la interfaz de \`arbol declarada,
\begin{equation}
\label{eq:ew-Higgs-lambda}
\boxed{
\lambda_H
=\frac{\pi\alpha}{2(1-D_A)}
\exp\!\left(6A_{\rm rad}+\frac{10}{3}C^*_{\rm rad}\right)
=0.1238479115482466804662\ldots .}
\end{equation}
Adem\'as,
\begin{equation}
\label{eq:ew-GF-Higgs}
G_F^{(0)}m_H^2=\sqrt2\,\lambda_H.
\end{equation}
\end{corollary}

\begin{proof}
Dividir las dos ecuaciones de
\cref{eq:ew-Higgs-tree-relations} da
\(m_H^2/m_W^2=8\lambda_H/g^2\). Con
\(g^2=4\pi\alpha/(1-D_A)\) y el cuadrado de
\cref{eq:ew-HW-ratio} se obtiene \cref{eq:ew-Higgs-lambda}. Por otra parte,
\(G_F^{(0)}=1/(\sqrt2v^2)\) y
\(m_H^2=2\lambda_Hv^2\) implican
\cref{eq:ew-GF-Higgs}.
\end{proof}

El valor de \(\lambda_H\) es anterior a la incorporaci\'on de correcciones
radiativas y depende de
la convenci\'on y la escala declaradas. Su inter\'es estructural es que
separa dos memorias: el sector com\'un de calibre emplea el invariante par \(D_A\),
mientras la ruta de Higgs retiene expl\'icitamente \(C^*\).

\section{Clausura determinantal del sector electrodébil: mezcla, acoplamientos y espectro constitutivo}

La cadena completa de este cap\'itulo es
\begin{equation}
\label{eq:ew-causal-diagram}
\boxed{
\begin{aligned}
(A,C^*)&\longrightarrow T
\longrightarrow
\begin{cases}
F(T)^2&\longrightarrow
(\widehat\mu,\widehat\varepsilon,\widehat Z,\widehat c),\\
\det T=D_A&\longrightarrow
(s_W^2,c_W^2),
\end{cases}\\
(D_A,\alpha)&\longrightarrow(e_g,g,g'),\\
(D_A,\alpha,m_W)&\longrightarrow G_F^{(0)},\\
(A,C^*,D_A,\alpha)&\longrightarrow
\left(\frac{m_H}{m_W},\lambda_H\right).
\end{aligned}}
\end{equation}
La cifra aparece siempre al final. Ni \(\sin^2\theta_W\), ni
\(G_F\), ni \(\lambda_H\) se introducen como valores objetivo para seleccionar
\(A,C^*\) o las firmas de ruta.

\begin{proposition}[Resultado principal]
El mismo operador angular determina el vac\'io mediante una funci\'on racional
de sus hojas y el defecto electrod\'ebil mediante su determinante. La
diferencia de un paso entre \(W\) y \(Z\) convierte ese determinante en
\(\sin^2\theta_W^{\rm OS}=1-D_A\); la ruta de Higgs, en cambio, conserva
el contra\'angulo y conserva la orientaci\'on eliminada por el determinante.
\end{proposition}

\begin{remark}[Control negativo]
La relaci\'on de Weinberg queda refutada en su dominio si las firmas
angulares \(W/Z\) no difieren s\'olo por \(n_A:94\to95\), si una coordenada
de refinamiento no se cancela en la interfaz declarada o si
\((m_W^{\angle}/m_Z^{\angle})^2\neq D_A\). La promoci\'on de calibre queda
bloqueada si no existe un entrelazador con la base \((W^3,B)\). Las cifras
de \(g,g',G_F,\lambda_H\) cambian si se altera esquema o escala sin
registrarlo. Usar \(G_F\) o \(m_H\) observados para seleccionar
\(\Delta r\), las firmas o \(C^*\) introduce retroalimentaci\'on por el
valor objetivo e invalida la derivaci\'on. Confundir \(e_g\) con la carga
dimensional constituye un error de tipo.
\end{remark}

\subsection*{Alcance lógico y dominio de validez}
Determinante angular, cociente \(W/Z\) y defecto complementario de Weinberg:
\emph{Teorema interno exacto en la carta angular condicionada}. Acoplamientos \(g,g'\),
\(\rho_{\rm EW}^{(0)}\), Fermi y Higgs: \emph{Composici\'on exacta en la
interfaz electrod\'ebil de \`arbol}. El entrelazador de base de calibre y la
correcci\'on interna \(\Delta r\): \emph{Frontera abierta con criterio de refutación}.
Fermi y Higgs no son constantes independientes seleccionadas sin valores
objetivo: son composiciones exactas que emplean las firmas, la magnitud
\(m_W\) y la
convenci\'on de \`arbol declarada. Los valores convencionales posteriores
son \emph{Contraste metrol\'ogico externo}.

 \section{Involución de hojas en \(4\) representaciones}

La involuci\'on \(C^*\mapsto-C^*\) organiza cuatro respuestas distintas:
\begin{equation}
\label{eq:modular-sheet-parities}
\begin{aligned}
S_{\rm bi}(A,-C^*)&=S_{\rm bi}(A,C^*),\\
\Gamma_{6\to8}(A,-C^*)&=-\Gamma_{6\to8}(A,C^*),\\
(\widehat\mu,\widehat\varepsilon)&\longleftrightarrow
(\widehat\varepsilon,\widehat\mu),\\
\widehat c&\longmapsto\widehat c,
\qquad D_A\longmapsto D_A.
\end{aligned}
\end{equation}
La informaci\'on bicapa es par; la orientaci\'on areal es impar; las dos
constitutivas forman un doblete, y los determinantes quedan fijos. Esta tabla
de paridades precede a cualquier interpretaci\'on como simetr\'ia f\'isica.

\section{Transformaci\'on racional involutiva en el espacio de par\'ametros CKM}

Sea
\begin{equation}
\label{eq:modular-Mckm}
M_{\rm CKM}=
\begin{pmatrix}
2&-5&28\\
0&7&-25/2\\
1&-20&-2/3
\end{pmatrix},
\qquad
\det M_{\rm CKM}=-\frac{3857}{6}.
\end{equation}
La inversa existe sobre \(\mathbb Q\) y es
\begin{equation}
\label{eq:modular-Mckm-inverse}
M_{\rm CKM}^{-1}=
\begin{pmatrix}
1528/3857&3380/3857&801/3857\\
75/3857&176/3857&-150/3857\\
6/551&-30/551&-12/551
\end{pmatrix}.
\end{equation}
Con
\begin{equation}
\label{eq:modular-ckm-coordinates}
h=\frac{C^*}{6},
\qquad
\Delta=\frac{180}{\pi}\Delta_4,
\qquad
\boldsymbol\theta=M_{\rm CKM}(A,h,\Delta)^T,
\end{equation}
la coordenada impar entra en la direcci\'on
\begin{equation}
\label{eq:modular-ch}
c_h=\begin{pmatrix}-5\\7\\-20\end{pmatrix}.
\end{equation}

Si \(c_A=(2,0,1)^T\) y
\(c_\Delta=(28,-25/2,-2/3)^T\), las tres columnas
\((c_A,c_h,c_\Delta)\) son precisamente la base transportada por
\(M_{\rm CKM}\). La segunda fila de
\cref{eq:modular-Mckm-inverse} es el covector \(\ell_h\) que extrae la
coordenada impar:
\begin{equation}
\label{eq:modular-dual-coordinate-check}
\ell_hc_A=0,
\qquad
\ell_hc_h=1,
\qquad
\ell_hc_\Delta=0.
\end{equation}

\begin{theorem}[Reflexi\'on CKM inducida por la hoja]
\label{thm:modular-ckm-reflection}
Sea
\begin{equation}
\label{eq:modular-Dh-lh}
D_h=\operatorname{diag}(1,-1,1),
\qquad
\ell_h=\frac1{3857}\begin{pmatrix}75&176&-150\end{pmatrix}.
\end{equation}
La involuci\'on inducida en el espacio de \(\boldsymbol\theta\) es
\begin{equation}
\label{eq:modular-Rckm}
\boxed{
\mathcal R_{\rm CKM}
=M_{\rm CKM}D_hM_{\rm CKM}^{-1}
=I-2c_h\ell_h.}
\end{equation}
Satisface
\begin{equation}
\label{eq:modular-Rckm-properties}
\mathcal R_{\rm CKM}^2=I,
\qquad
\det\mathcal R_{\rm CKM}=-1,
\qquad
\mathcal R_{\rm CKM}M_{\rm CKM}=M_{\rm CKM}D_h.
\end{equation}
\end{theorem}

\begin{proof}
La multiplicaci\'on de
\cref{eq:modular-Mckm,eq:modular-Mckm-inverse} da la identidad; en
particular, \cref{eq:modular-dual-coordinate-check} prueba que
\(c_h\ell_h\) es el proyector de rango uno sobre \(\mathbb Qc_h\) a lo
largo de \(\operatorname{span}\{c_A,c_\Delta\}\). El producto exacto
\(\ell_hc_h=1\) implica
\[
(I-2c_h\ell_h)^2
=I-4c_h\ell_h+4c_h(\ell_hc_h)\ell_h=I.
\]
El operador fija punto por punto el hiperplano
\(\ker\ell_h=\operatorname{span}\{c_A,c_\Delta\}\) y cambia el signo de
\(c_h\). Sus autovalores son, por tanto, \(1,1,-1\); de aqu\'i
\(\det\mathcal R_{\rm CKM}=-1\) y
\(\Tr\mathcal R_{\rm CKM}=1\). Finalmente,
\[
\mathcal R_{\rm CKM}M_{\rm CKM}
=M_{\rm CKM}D_hM_{\rm CKM}^{-1}M_{\rm CKM}
=M_{\rm CKM}D_h,
\]
que prueba la identidad de entrelazamiento.
\end{proof}

La matriz expl\'icita es
\begin{equation}
\label{eq:modular-Rckm-matrix}
\mathcal R_{\rm CKM}=
\begin{pmatrix}
4607/3857&1760/3857&-1500/3857\\
-150/551&199/551&300/551\\
3000/3857&7040/3857&-2143/3857
\end{pmatrix}.
\end{equation}
Aunque \(\mathcal R_{\rm CKM}\) no es una reflexi\'on ortogonal para el
producto eucl\'ideo de las tres coordenadas angulares, s\'i lo es para la
m\'etrica racional transportada
\begin{equation}
\label{eq:modular-ckm-metric}
G_{\rm CKM}=(M_{\rm CKM}^{-1})^TM_{\rm CKM}^{-1},
\qquad
\mathcal R_{\rm CKM}^TG_{\rm CKM}\mathcal R_{\rm CKM}=G_{\rm CKM}.
\end{equation}
En efecto, en la base \((c_A,c_h,c_\Delta)\), la m\'etrica es la identidad
y la involuci\'on es \(D_h\). Por eso el t\'ermino ``reflexi\'on'' tiene un
contenido m\'etrico preciso y no describe s\'olo que
\(\mathcal R_{\rm CKM}^2=I\).
La fase \(\delta_{\rm CP}=9A\) permanece fija bajo esta reflexi\'on. Por
ello, \(\mathcal R_{\rm CKM}\) no es la conjugaci\'on CP f\'isica, sino el
transporte de la inversi\'on de hoja HMT al espacio de mezcla.

La inversa racional de \(M_{\rm CKM}\) recupera
\begin{align}
A&=\frac{1528\theta_{12}+3380\theta_{23}+801\theta_{13}}{3857},
\label{eq:modular-A-from-ckm}\\
C^*&=\frac{6(75\theta_{12}+176\theta_{23}-150\theta_{13})}{3857}.
\label{eq:modular-C-from-ckm}
\end{align}
La tercera coordenada tambi\'en se recupera sin p\'erdida:
\begin{equation}
\label{eq:modular-Delta-from-ckm}
\Delta
=\frac{6\theta_{12}-30\theta_{23}-12\theta_{13}}{551}.
\end{equation}
Estas ecuaciones son cambios de coordenadas. No convierten CKM en causa
retrospectiva de \(A,C^*\) o de \(\gammaH\).

Como reconocimiento terminal, el constructor previo determina
\begin{equation}
\label{eq:modular-ckm-values}
(\theta_{12},\theta_{23},\theta_{13},\delta_{\rm CP})
=(13.003313589509^\circ,
2.398056157332^\circ,
0.213977382244^\circ,
65.676173123554^\circ).
\end{equation}
La reflexi\'on racional se demuestra antes de usar estos decimales.

\begin{proposition}[Resultado principal]
La transici\'on \(6\to8\) se formaliza primero mediante la inclusi\'on
\(\iota_{6,8}:\mathcal F_6\hookrightarrow\mathcal F_8\), de la que se
obtienen
\(90,120\) y el defecto \(-1/12\). La palabra de seis trits proporciona
despu\'es la biyecci\'on \(729\leftrightarrow729\), y
\(\mathcal J_{6\to8}\) transporta los dos modos colectivos al borde. Sobre
ese antecedente, el \`area y el Hamiltoniano modular reconstruyen exactamente
las dos familias; la purificaci\'on termocampo doble determina la entrop\'ia
de entrelazamiento, y la primera ley
finita conserva el defecto de entrop\'ia relativa. En CKM, la inversi\'on de
hoja se convierte en una reflexi\'on racional de rango impar uno.
\end{proposition}

\begin{remark}[Control negativo]
La genealog\'ia inicial falla si \(\iota_{6,8}\) no es inyectiva, si los
conteos no son \(90/120\), si \(\delta_{6,8}^{\rm inc}\neq-1/12\), si
\(\beta_{729}\) no es biyectiva o si se la presenta falsamente como
isomorfismo aditivo. El transporte colectivo falla si
\(\mathcal J_{6\to8}D_{\rm inc}\neq
D_{\rm area}\mathcal J_{6\to8}\). La reconstrucci\'on falla si el
determinante de los observables de borde no es treinta o si
\cref{eq:modular-inverse-channels} no recupera las ocupaciones.
La TFD falla si su traza parcial no es \(\rho_A\). La primera ley no se
extiende a cambios finitos sin la entrop\'ia relativa de
\cref{eq:modular-finite-first-law}. El tipo
\(III_{\lambda_\partial}\) queda invalidado si un canal no persiste, si no hay
producto compatible o si cambia el grupo de razones. La reflexi\'on CKM
queda refutada por cualquiera de
\(\mathcal R^2\neq I\), \(\det\mathcal R\neq-1\) o
\(\mathcal RM\neq MD_h\). Una f\'ormula
\(\boldsymbol\theta=f(\gammaH)\) sin entrelazador viola la independencia
tipada de las dos realizaciones funcionales del antecedente com\'un.
\end{remark}

\subsection*{Alcance lógico y dominio de validez}
Inclusi\'on \(H\subset O\), conteos \(90/120\), dos normalizaciones del
defecto \(-1/12\), biyecci\'on de seis trits \(729\leftrightarrow729\),
entrelazador colectivo, escala \(a_0\), operadores finitos, espectro areal,
TFD, defecto orientado, primera ley local y finita y reflexi\'on CKM:
\emph{Teorema interno exacto con sus estructuras de partida expl\'icitas}.
Clasificaci\'on
\(III_{\lambda_\partial}\): \emph{Teorema con hip\'otesis expl\'icitas de
persistencia}. Interpretaci\'on de la copia termocampo doble como exterior f\'isico:
\emph{Realizaci\'on f\'isica condicionada}. La magnitud \(\gammaH\) se adopta
como \emph{Resultado antecedente demostrado}; no se rederiva en este volumen.

 El vínculo electrodébil parte del mismo operador angular que produjo las constitutivas del vacío, pero utiliza otro funcional.\par

Para obtener \(\mu\), \(\varepsilon\), \(Z\) y \(c\) había que conservar el espectro completo. Había que distinguir las dos hojas.\par

El determinante hace lo contrario:\par

\[
D_A=\det T=q_+q_-.
\]

Al multiplicar los autovalores, el determinante integra las dos orientaciones en una magnitud angular común e invariante. Publica la parte par y reserva la identidad de hoja en la memoria orientada.\par

Por eso \(D_A\) es par bajo la inversión\par

\[
C^*\longmapsto-C^*.
\]

El sector \(W/Z\) posee justamente esa estructura. Las rutas de \(W\) y \(Z\) son adyacentes en la coordenada angular principal y comparten la misma coordenada orientada. Al formar el cociente, la contribución de \(C^*\) se cancela.\par

Queda\par

\[
\left(\frac{m_W}{m_Z}\right)^2
=
D_A,
\]

y por tanto\par

\[
\sin^2\theta_W
=
1-D_A.
\]

La identidad revela que el ángulo de Weinberg es el defecto complementario del determinante angular.\par

El determinante mide la parte común que sobrevive al colapsar las dos hojas en su producto. \(1-D_A\) mide la parte complementaria necesaria para la mezcla.\par

Podríamos decirlo así: el vacío conserva toda la respuesta; el sector electrodébil utiliza la compresión par de esa respuesta.\par

El mismo antecedente angular produce dos funciones distintas:\par

\begin{itemize}[leftmargin=*,itemsep=0.35em]
\item el cálculo funcional completo genera las relaciones constitutivas;
\item el determinante genera el invariante de mezcla.
\end{itemize}

Eso explica por qué vacío y ángulo de Weinberg son dos realizaciones funcionales distintas de un mismo antecedente angular.\par

A partir de ahí, los acoplamientos \(g\) y \(g'\) son, en el nivel declarado, las dos proyecciones del mismo acoplamiento electromagnético racionalizado sobre las direcciones seno y coseno de la mezcla:\par

\[
g^2
=
\frac{4\pi\alpha}{1-D_A},
\qquad
g'^2
=
\frac{4\pi\alpha}{D_A}.
\]

La identidad custodial de árbol\par

\[
\rho_{\rm EW}^{(0)}=1
\]

expresa que la misma partición angular organiza coherentemente masa y acoplamiento.\par

La constante de Fermi aparece después como la respuesta efectiva de baja energía del canal cargado:\par

\[
G_F^{(0)}
=
\frac{\pi\alpha}
{\sqrt2(1-D_A)m_W^2}.
\]

La constante de Fermi aparece como una composición de la estructura fina, el defecto angular y la escala del portador \(W\).\par

El sector de Higgs introduce una diferencia decisiva. Su ruta posee una coordenada orientada distinta de la ruta de \(W\), por lo que \(C^*\) persiste. El autoacoplamiento escalar publica la orientación que el determinante reserva.\par

El sector de calibre es par: utiliza \(D_A\).\par
El sector escalar es orientado: retiene \(C^*\).\par

Esta división resulta físicamente sugerente. Los campos de calibre organizan la parte común de la transformación; el sector escalar conserva la memoria de la orientación rota.\par

Las correcciones radiativas pertenecen a un nivel posterior y deben construirse prospectivamente desde los bucles, estados y escalas internos, con \(G_F\) como salida y evaluación terminal. El esqueleto causal ya está presente: vacío, mezcla, acoplamientos y respuesta de Fermi proceden de distintas lecturas del mismo estado angular enriquecido.\par

 \section{Cierre angular electrodébil y condiciones de compatibilidad del transductor de calibre}
\label{sec:p3-final:cierre-electrodebil}

Sean
\[
 c_W=\sqrt{D_A},
 \qquad
 s_W=\sqrt{1-D_A}.
\]
La rotación
\[
 \begin{pmatrix}Z\\A\end{pmatrix}
 =
 \begin{pmatrix}c_W&-s_W\\s_W&c_W\end{pmatrix}
 \begin{pmatrix}W^3\\B\end{pmatrix}
\]
es ortogonal y preserva la orientación. Las identidades
\[
 \left(\frac{m_W}{m_Z}\right)^2=D_A,
 \qquad
 \sin^2\theta_W=1-D_A
\]
quedan cerradas en la carta angular. La promoción gauge requiere que el
transductor preserve la forma cinética, las cargas y las identidades de Ward;
la corrección \(\Delta r\) pertenece a la dinámica radiativa y no se obtiene
ajustándola a \(G_F\).

  \endgroup

\chapter{Derivación intrínseca de la gravitación y reciprocidad contragrediente \texorpdfstring{\(G\leftrightarrow\hbar\)}{G<->hbar}}
\begingroup
\renewcommand{\chapter}[2][]{}%
\chapter{Derivación intrínseca de \(G\) y reciprocidad \(G\leftrightarrow\hbar\): cierre de fase, longitudes de Compton y Planck e invariancia areal}
\label{chap:p3-final:50:gravitacion_accion}
\label{ch:gravity}
% BEGIN CR28_RADIAL_FULL
% Sección incorporable; requiere amsmath y amssymb.
% Fuente principal: RIGIDEZ_SIMULTANEA_DEL_LECTOR_RADIAL_20260926.md.
% Complementos demostrativos: COMPOSICION_METRICA_MEMORIA_Y_ACOPLAMIENTO.md
% y ACCION_ACOPLADA_FASE_Y_MEMORIA_DE_FRONTERA.md.
% La integración, la compilación y la revisión visual corresponden al editor.
\section{Rigidez simultánea del lector longitudinal y del acoplamiento radial}
\label{g26:sec:rigidez-radial-es}

La construcción de Holografía Modular Triádica transmite a la realización
radial un estado enriquecido con sus coordenadas, orientación, acarreo,
memoria e incidencias. La sucesión
\[
 \begin{gathered}
 \mathrm{APP}\longrightarrow\mathrm{TRIT}\longrightarrow\mathrm{TPK}
 \longrightarrow\text{estado enriquecido}\\
 \longrightarrow\text{estructura discreta conjunta del continuo}
 \end{gathered}
\]
conserva la proyección arquimediana y la proyección de incidencia, junto con
las cinco construcciones consustanciales del continuo. La prolongación
\[
 w_6\longrightarrow w_{12}\longrightarrow w_{18}\longrightarrow
 w_{24}\longrightarrow w_{30}\longrightarrow R_{36}\longrightarrow G_9
\]
se articula con el orden operatorio
\(U_t\longrightarrow\Gamma_9\longrightarrow K_{\rm ph}\): la fase retorna
y el registro de memoria se incrementa. Las coordenadas regionales, el
registro seleccionado \(K\), la coordenada \(\alpha\) y la acción de retorno
se reciben con esta procedencia común y con sus lectores previamente
construidos.

El resultado de esta sección es una clasificación de realizaciones que
conservan simultáneamente los datos generados y las reglas longitudinales,
métricas y temporales que se precisan a continuación. Para cada carácter
positivo queda determinada la respuesta radial; el coeficiente que relaciona
esa respuesta con la masa es común a todos los caracteres admisibles.
La interpretación gravitatoria se formula después de construir ambas
lecturas sobre el mismo espacio.

\subsection{Registro de incidencias y restricciones constitutivas}

Se consideran espacios de Hilbert complejos \(V_{\rm in}\) y
\(V_{\rm out}\) de dimensión \(N\), con sus productos escalares fijados.
El cardinal marcado que se especifica en R1 es positivo; en particular,
ambas fibras son no nulas. Los adjuntos se toman respecto de estos productos. El símbolo
\(\mathsf U\) designará el transporte normalizado del archivo y
\(U_t\) conservará su significado de evolución del estado enriquecido.

\paragraph{R1. Registro marcado.}
El conjunto de incidencias es
\begin{equation}
 \Omega=\left\{(j,k,q,u,v):
 \begin{array}{l}
 j,k\in\mathbb Z/12\mathbb Z,\quad k-j\in\{0,3,6,9\},\\
 1\leq q\leq8,\quad u<v,\quad
 u,v\in\{1,2,4,5,7,8\}
 \end{array}\right\}.
 \label{g26:eq:omega-rigidez-es}
\end{equation}
Las doce fases, las cuatro posiciones relativas, las ocho posiciones del
índice \(q\) y las quince parejas del último conjunto dan
\begin{equation}
 N=|\Omega|=12\cdot4\cdot8\binom62
   =12\cdot4\cdot120=5760.
 \label{g26:eq:cardinal-rigidez-es}
\end{equation}
Cada incidencia posee su marca individual. En la base correspondiente
\((e_i)_{i\in\Omega}\) de \(V_{\rm out}\), sean
\(P_i=e_ie_i^*\). El modo común \(\boldsymbol u\) está normalizado por
\(|u_i|^2=1/N\); se escribe \(P=\boldsymbol u\boldsymbol u^*\).
El álgebra marcada es
\(\mathcal A_\Omega=\operatorname{span}\{P_i:i\in\Omega\}\).

\paragraph{R2. Normalización y archivo completo.}
El inversor generado \(B:V_{\rm in}\to V_{\rm out}\) satisface
\begin{equation}
 B^*B=\tfrac14 I_{\rm in},\qquad
 BB^*=\tfrac14 I_{\rm out}.
 \label{g26:eq:normalizacion-b-es}
\end{equation}
Se conservan las dos componentes
\begin{equation}
 C_0=(I-P)B,\qquad M_0=PB,\qquad
 C_0+M_0=B,\qquad \mathsf U=2B.
 \label{g26:eq:archivo-completo-es}
\end{equation}
La primera registra la variación respecto del modo común; la segunda
conserva dicho modo. Sus imágenes son ortogonales y su suma recupera
el transporte completo. Por \eqref{g26:eq:normalizacion-b-es},
\(\mathsf U\) es unitario.

\paragraph{R3. Inscripción individual.}
La aplicación
\(\mathcal E:\operatorname{End}(V_{\rm out})\to\mathcal A_\Omega\)
es una expectativa condicional lineal que fija cada \(P_i\) y satisface
\begin{equation}
 \mathcal E(A_1TA_2)=A_1\mathcal E(T)A_2,
 \qquad A_1,A_2\in\mathcal A_\Omega.
 \label{g26:eq:bimodularidad-es}
\end{equation}
La inscripción del resultado conserva, en un registro complementario,
los datos completos de la construcción. La expectativa determina el
observable inscrito; la conservación del archivo determina su
reconstrucción posterior.

\paragraph{R4. Composición longitudinal.}
Se fijan la coordenada generada \(\alpha>0\) y la sección de referencia
longitudinal \(L_*>0\). El transporte longitudinal se define por
\begin{equation}
 D=L_*\alpha^{16}\mathcal E(C_0C_0^*)\mathsf U:
 V_{\rm in}\longrightarrow V_{\rm out}.
 \label{g26:eq:regla-longitudinal-es}
\end{equation}
Esta regla aplica linealmente el coeficiente inscrito a la sección
longitudinal. Su tipo fija la posición de la potencia \(\alpha^{16}\)
y del coeficiente de retorno dentro de la composición. La lectura de
una raíz de intensidad constituye una operación diferente y posee su
propia regla de transformación.

\paragraph{R5. Métrica radial positiva.}
Sea \(X=X^*>0\) el carácter completo sobre \(V_{\rm in}\). Se define
la lectura radial positiva \(R_X\) sobre \(V_{\rm out}\) mediante
\begin{equation}
 \|R_Xv\|^2=\|XD^*v\|^2
 \qquad(v\in V_{\rm out}).
 \label{g26:eq:metrica-radial-es}
\end{equation}
Esta igualdad fija la métrica en cada vector de la fibra. En notación
polar, \(R_X\) es el módulo de llegada
\(\lvert(DX)^*\rvert=\sqrt{DX^2D^*}\).

\paragraph{R6. Acción y reloj.}
La sección de acción generada \(\hbar_{\rm ret}>0\) y el levantamiento
real de fase satisfacen
\begin{equation}
 S(\theta)=\hbar_{\rm ret}\theta.
 \label{g26:eq:accion-fase-es}
\end{equation}
Se conserva la velocidad constitutiva \(c>0\). Una vez obtenida la
longitud intensiva \(L\) de \eqref{g26:eq:regla-longitudinal-es}, el reloj
conjugado satisface \(\omega L=c\). Para \(\theta(t)=\omega t\),
la tasa de acción determina la escala energética
\begin{equation}
 E_{L}=\frac{dS}{dt}=\hbar_{\rm ret}\omega
   =\frac{\hbar_{\rm ret}c}{L}.
 \label{g26:eq:escala-energetica-es}
\end{equation}
El levantamiento y sus refinamientos permanecen asociados a la memoria
de la ruta.

\paragraph{R7. Lecturas del mismo carácter.}
Sobre \(Y=\mathsf U X\mathsf U^*>0\), se definen
\begin{equation}
 H_X=E_{L}Y,\qquad M_X=\frac{H_X}{c^2},\qquad
 \Lambda_X=LY^{-1}.
 \label{g26:eq:lecturas-conjugadas-es}
\end{equation}
La energía, la masa y la longitud conjugada conservan el mismo carácter
y el mismo transporte. La inversa actúa en el sector positivo; el sector
nulo conserva su tratamiento propio.

\paragraph{R8. Radio reducido y secciones dimensionales.}
La identificación física de esta realización es
\begin{equation}
 R_X=\frac{G}{c^2}M_X.
 \label{g26:eq:radio-reducido-es}
\end{equation}
Así, \(R_X\) representa el radio gravitatorio reducido. El radio de
Schwarzschild correspondiente es \(2R_X\). Se mantienen conjuntamente
\(L_*\), la sección de acción y la sección constitutiva de velocidad;
un cambio de unidades transforma sus coordenadas según sus dimensiones.

\paragraph{Procedencia de las restricciones.}
El transporte de Hadamard, la conservación del archivo, la inscripción,
la acción de fase y la reciprocidad poseen antecedentes en el corpus.
El registro ampliado \eqref{g26:eq:omega-rigidez-es} y las composiciones
\eqref{g26:eq:regla-longitudinal-es}--\eqref{g26:eq:metrica-radial-es} se
incorporan explícitamente en la construcción reunida. R4 y R5 son reglas
constitutivas de esta realización; la demostración siguiente establece
su consecuencia conjunta con los restantes antecedentes. R8 precisa
la identificación física que fija el significado del coeficiente.

\subsection{Teorema de determinación simultánea}

\paragraph{Teorema.}
Para los mismos datos generados \(B\), \((P_i)\), \(P\),
\(\alpha\), \(L_*\), \(\hbar_{\rm ret}\), \(c\) y el mismo carácter
\(X\), las restricciones R1--R8 determinan unívocamente la inscripción,
el transporte longitudinal y la respuesta radial. Para todos los
caracteres positivos admisibles, el coeficiente \(G\) es común. Se cumple
\begin{align}
 \mathcal E&=\Delta_\Omega,
 &\Delta_\Omega(T)&=\sum_{i\in\Omega}P_iTP_i,\label{g26:eq:delta-unica-es}\\
 r_N&=\frac{N-1}{4N}=\frac{5759}{23040},
 &L&=L_*\alpha^{16}r_N,\label{g26:eq:longitud-unica-es}\\
 D&=L\mathsf U,
 &R_X&=L\mathsf U X\mathsf U^*,\label{g26:eq:radial-unica-es}\\
 R_X\Lambda_X&=L^2I,
 &H_X\Lambda_X&=\hbar_{\rm ret}cI.\label{g26:eq:productos-radiales-es}
\end{align}
El acoplamiento queda determinado por
\begin{equation}
 \boxed{G=\frac{c^3L^2}{\hbar_{\rm ret}}
 =\frac{c^3L_*^2}{\hbar_{\rm ret}}
   \left(\frac{5759}{23040}\right)^2\alpha^{32}.}
 \label{g26:eq:g-rigido-es}
\end{equation}

\paragraph{Demostración: unicidad de la inscripción.}
Sean \(E_{ij}=e_ie_j^*\) las unidades matriciales. Por bimodularidad,
\[
 \mathcal E(E_{ij})=P_i\mathcal E(E_{ij})P_j.
\]
Para \(i\ne j\), la pertenencia de \(\mathcal E(E_{ij})\) al álgebra
diagonal implica \(P_i\mathcal E(E_{ij})P_j=0\). Para \(i=j\),
la fijación de \(P_i\) da \(\mathcal E(E_{ii})=P_i\).
La linealidad determina entonces \eqref{g26:eq:delta-unica-es} sobre todas
las matrices. La inscripción marcada queda así fijada por sus
propiedades algebraicas.

\paragraph{Demostración: coeficiente de retorno y archivo.}
La normalización de \(B\) y \(P^2=P=P^*\) proporcionan
\[
 C_0C_0^*=\tfrac14(I-P),\qquad M_0M_0^*=\tfrac14P.
\]
Como \(P_iPP_i=|u_i|^2P_i=P_i/N\), se obtiene
\begin{equation}
 \Delta_\Omega(C_0C_0^*)=\frac{N-1}{4N}I,
 \qquad
 \Delta_\Omega(M_0M_0^*)=\frac{1}{4N}I.
 \label{g26:eq:retorno-y-memoria-es}
\end{equation}
Ambas lecturas suman \(I/4\). La primera es escalar, de modo que
cualquier distribución normalizada \((p_i)\) sobre las incidencias,
con \(p_i\geq0\) y \(\sum_i p_i=1\), produce
\(\sum_i p_i r_N=r_N\). Se preserva, por tanto, el mismo coeficiente
bajo una ponderación diagonal arbitraria.

\paragraph{Demostración: longitud y métrica.}
La sustitución de \eqref{g26:eq:retorno-y-memoria-es} en R4 da
\(D=L\mathsf U\), con \(L\) como en \eqref{g26:eq:longitud-unica-es}.
La unitariedad de \(\mathsf U\) implica \(DD^*=L^2I\).
Por R5, las formas cuadráticas de \(R_X^2\) y \(DX^2D^*\) coinciden.
La identidad de polarización determina
\[
 R_X^2=DX^2D^*=L^2\mathsf U X^2\mathsf U^*.
\]
El operador \(L\mathsf U X\mathsf U^*\) es positivo y su cuadrado
es el operador del segundo miembro. La unicidad de la raíz cuadrada
positiva demuestra \eqref{g26:eq:radial-unica-es}. En particular,
\(R_I=LI\).

\paragraph{Demostración: acción, reciprocidad y acoplamiento.}
R6 y R7 dan
\[
 H_X=\frac{\hbar_{\rm ret}c}{L}Y,
 \qquad M_X=\frac{\hbar_{\rm ret}}{cL}Y,
 \qquad \Lambda_X=LY^{-1}.
\]
La multiplicación produce ambos productos de
\eqref{g26:eq:productos-radiales-es}. A continuación, R8 equivale a
\[
 LY=\frac{G\hbar_{\rm ret}}{c^3L}Y.
\]
La invertibilidad de \(Y\) determina \eqref{g26:eq:g-rigido-es}.
La misma sustitución verifica la igualdad para cada carácter positivo.
Si \(G_1\) y \(G_2\) satisfacen la identidad con los mismos datos,
\((G_1-G_2)Y=0\), y por tanto \(G_1=G_2\).
Esta conclusión distingue la respuesta dependiente del estado
\(R_X\) del coeficiente común \(G\). \hfill\(\square\)

\subsection{Reciprocidad y rigidez de la fase}

El antecedente potencia--logarítmico utiliza, para \(x>0\),
\[
 \ell_p(x)=\frac{x^p-1}{p}\quad(p\ne0),
 \qquad \ell_0(x)=\log x.
\]
La sustitución directa demuestra
\(\ell_{-p}(x^{-1})=-\ell_p(x)\). En el dominio positivo de sus
inversas, esta identidad equivale a
\[
 \exp_p(s)\exp_{-p}(-s)=1.
\]
La composición espectral sobre un carácter positivo transmite este
producto al par \(LY\), \(LY^{-1}\). En la realización actual, R5
puede expresarse equivalentemente, una vez fijadas R1--R4 y R7, por
\(R_X\Lambda_X=L^2I\): multiplicar por \(\Lambda_X^{-1}\) determina
\(R_X=L\mathsf U X\mathsf U^*\), y esta expresión satisface R5.

Los dos productos de \eqref{g26:eq:productos-radiales-es} dan además una
forma directa de la rigidez:
\begin{equation}
 R_X=\frac{L^2}{\hbar_{\rm ret}c}H_X.
 \label{g26:eq:proporcion-operadores-es}
\end{equation}
Para una reescala positiva \(s\), el par
\((sR_X,\Lambda_X/s)\) conserva el producto radial; el producto
\(H_X(\Lambda_X/s)=\hbar_{\rm ret}cI/s\), con energía y acción
fijadas, exige \(s=1\). La conservación conjunta determina la escala.
Asimismo, cualquier operador radial positivo con la misma métrica R5
coincide con \(R_X\) por la unicidad de la raíz positiva.

La acción levantada también posee una determinación exacta. Para
\(\theta\ne0\), la igualdad \(S/\hbar_{\rm ret}=\theta\) fija \(S\).
Incluso cuando se conservan las rotaciones de todas las subdivisiones
\(\theta/9^n\), una reescala real fija \(a\) que satisfaga
\[
 \exp(-ia\theta/9^n)=\exp(-i\theta/9^n)
 \qquad(n\geq0)
\]
verifica \((a-1)\theta/9^n\in2\pi\mathbb Z\). Esta sucesión tiende
a cero; para \(n\) suficientemente grande su único valor posible es
cero. De \(\theta\ne0\) se deduce \(a=1\).

\subsection{Escalas de Planck y especialización del carácter unitario}

Escribimos \(L_\alpha=L\); en la sección de retorno, \(E_{P,\rm ret}=E_L\).

La longitud construida $L_\alpha>0$, la sección de acción $\hbar_a>0$ y la
velocidad $c>0$ determinan las escalas

\[
t_P=\frac{L_\alpha}{c},\qquad
m_{P,a}=\frac{\hbar_a}{cL_\alpha},\qquad
E_{P,a}=\frac{\hbar_a c}{L_\alpha}.
\]

En efecto, sustituir $G_a=c^3L_\alpha^2/\hbar_a$ en las expresiones
$\sqrt{\hbar_aG_a/c^5}$, $\sqrt{\hbar_ac/G_a}$ y
$\sqrt{\hbar_ac^5/G_a}$ da, respectivamente, esas tres identidades.
La positividad fija en cada caso la raíz. La longitud antecede al
reconocimiento de estas escalas; sus expresiones convencionales son
identidades de la realización ya construida.

Para un modo de carácter $y>0$, las lecturas másica y longitudinal son

\[
m(y)=m_{P,a}y,\qquad r_g(y)=L_\alpha y,\qquad
\bar\lambda(y)=\frac{L_\alpha}{y}.
\]

Por tanto, $r_g(y)\bar\lambda(y)=L_\alpha^2$ para todos los modos
positivos. La igualdad de ambas longitudes caracteriza el modo $y=1$,
puesto que $y=y^{-1}$ y $y>0$ implican $y=1$. En este modo, la masa es
$m_{P,a}$ y ambas longitudes coinciden con $L_\alpha$. Así se distingue
la reciprocidad general de su especialización en el carácter unitario.

El paso cronológico y el tiempo de Planck satisfacen
$t_0=(\pi_{\rm HMT}/54)t_P$. El período completo de fase de 108 pasos es
$108t_0=2\pi_{\rm HMT}t_P$ y su frecuencia angular es
$\omega_{108}=1/t_P=c/L_\alpha$. Un horizonte de radio
$R_h=\zeta_hL_\alpha$, con $\zeta_h>0$, conserva su variable geométrica
propia. El radio de Schwarzschild asociado a una masa es $2r_g$.

Si se comparan las dos secciones de acción con $L_\alpha$ y $c$ fijados,
el cociente $\rho=\hbar_b/\hbar_a$ da
$G_b=G_a/\rho$, $m_{P,b}=\rho m_{P,a}$ y
$E_{P,b}=\rho E_{P,a}$; $L_\alpha$ y $t_P$ permanecen iguales.
Esta comparación entre secciones conserva la distinción respecto de una
evolución temporal de las constantes.

\subsection{Covarianza del archivo y transformación de unidades}

Sean \(T:V_{\rm in}\to V'_{\rm in}\) y
\(S:V_{\rm out}\to V'_{\rm out}\) unitarios. El transporte conjunto
de los datos es
\[
 B'=SBT^*,\quad P'=SPS^*,\quad P_i'=SP_iS^*,\quad X'=TXT^*.
\]
La expectativa transportada
\(\mathcal E'(A')=S\mathcal E(S^*A'S)S^*\) fija \(P_i'\)
y es bimodular en el álgebra marcada transformada. Por sustitución,
\begin{align*}
 \mathsf U'&=S\mathsf UT^*,&D'&=SDT^*,\\
 R'_{X'}&=SR_XS^*,&H'_{X'}&=SH_XS^*,&
 \Lambda'_{X'}&=S\Lambda_XS^*.
\end{align*}
La conjugación conserva las igualdades de R1--R8 y el mismo coeficiente
\(G\). Las marcas, orientaciones y registros se transportan con las
bases. Las permutaciones de incidencias son un caso particular.

La incorporación de un espacio auxiliar de memoria \(\mathcal F\)
transporta \(P\) a \(P\otimes I_{\mathcal F}\),
\(\mathsf U\) a \(\mathsf U\otimes I_{\mathcal F}\) y
\(\Delta_\Omega\) a \(\Delta_\Omega\otimes\operatorname{id}\).
La inscripción de \(C_0C_0^*\otimes I_{\mathcal F}\) es
\(r_N I\otimes I_{\mathcal F}\); el coeficiente conserva el cardinal
marcado original.

Para precisar la covarianza dimensional, sean \(\ell\), \(\tau\) y
\(m\) los factores positivos de las nuevas unidades de longitud, tiempo
y masa respecto de las anteriores. Las coordenadas numéricas satisfacen
\[
 L'=L/\ell,\qquad c'=c\tau/\ell,\qquad
 \hbar'_{\rm ret}=\hbar_{\rm ret}\tau/(m\ell^2).
\]
Así,
\[
 \frac{(c')^3(L')^2}{\hbar'_{\rm ret}}
 =G\frac{m\tau^2}{\ell^3},
\]
que es la transformación de una magnitud de dimensión
\(\mathrm{longitud}^3\mathrm{masa}^{-1}\mathrm{tiempo}^{-2}\).
La sección \(L_*\) se transforma como \(L\), mientras \(\alpha\)
y \(r_N\) permanecen adimensionales.

\subsection{Eliminación exacta de memoria acoplada}

Sean \(\mathcal K\) y \(\mathcal H\) espacios de Hilbert finitos,
\(W=W^*>0\) un peso global y
\(\mathcal C:\mathcal K\to\mathcal H\) sobreyectiva. La acción
de los residuos \(r\), con defecto de frontera \(z\), es
\[
 \mathcal S(r)=\mathfrak a\,r^*Wr,
 \qquad \mathcal Cr=z,\qquad \mathfrak a>0.
\]
El peso global conserva sus bloques cruzados entre incidencias. Defínase
\(\mathcal R=\mathcal CW^{-1}\mathcal C^*\). Para \(v\ne0\),
la inyectividad de \(\mathcal C^*\) da
\[
 v^*\mathcal Rv=\|W^{-1/2}\mathcal C^*v\|^2>0.
\]
Por tanto, \(\mathcal R\) es invertible. El residuo de mínima acción es
\[
 r_*(z)=W^{-1}\mathcal C^*\mathcal R^{-1}z.
\]
Se verifica \(\mathcal Cr_*=z\). Cada residuo admisible se escribe
\(r=r_*+k\), con \(\mathcal Ck=0\), y
\(k^*Wr_*=k^*\mathcal C^*\mathcal R^{-1}z=0\). En consecuencia,
\begin{equation}
 \mathcal S(r)=\mathfrak a\,z^*\mathcal R^{-1}z
             +\mathfrak a\,k^*Wk.
 \label{g26:eq:accion-residuo-es}
\end{equation}
La positividad de \(W\) prueba la unicidad del mínimo. El par
\((z,k)\) conserva y reconstruye el residuo completo.

Para una ruta con transportes \(U_j\), los residuos
\(r_j=f_j-U_jf_{j-1}\) se reconstruyen mediante
\(f_j=U_jf_{j-1}+r_j\). La aplicación \(\mathcal C\) reúne los
transportes posteriores de cada residuo hacia el extremo. Si sus
bloques son \(\mathcal C_j\), entonces
\[
 \mathcal R=\sum_{j,k}\mathcal C_j(W^{-1})_{jk}\mathcal C_k^*.
\]
Esta expresión conserva las correlaciones entre tramos de la ruta.
Para un transporte longitudinal invertible \(D\), la frontera
\(y=Dz\) tiene respuesta
\[
 \mathcal R_L=D\mathcal RD^*,\qquad
 \mathcal S_{{\rm ef},L}(y)
 =\mathfrak a\,y^*(D\mathcal RD^*)^{-1}y
 =\mathfrak a\,z^*\mathcal R^{-1}z.
\]
Se obtiene la misma acción al eliminar primero la memoria o transportar
primero la frontera. El prefactor conserva la normalización de la acción
fuente; cuando esa acción lleva \(\mathfrak a=\hbar_{\rm ret}\), la
reducción transmite la misma sección.

\subsection{Reducción conjunta de radio, energía y longitud conjugada}

Escríbase el mismo carácter positivo en una descomposición de frontera
y memoria:
\[
 Y=\begin{pmatrix}A&F\\F^*&C\end{pmatrix},\qquad C>0,
 \qquad Y_{\rm ef}=A-FC^{-1}F^*.
\]
La identidad
\[
 \begin{pmatrix}b\\y\end{pmatrix}^{\!*}Y
 \begin{pmatrix}b\\y\end{pmatrix}
 =b^*Y_{\rm ef}b+\eta^*C\eta,
 \qquad \eta=y+C^{-1}F^*b,
\]
prueba \(Y_{\rm ef}>0\), caracteriza la extensión de mínima energía y
conserva la reconstrucción \(y=\eta-C^{-1}F^*b\). Para \(s>0\),
la sustitución en el complemento de Schur da
\(\operatorname{Schur}(sY)=sY_{\rm ef}\). En consecuencia,
\[
 R_{\rm ef}=LY_{\rm ef},\qquad
 H_{\rm ef}=E_{L}Y_{\rm ef}.
\]
Para calcular la longitud conjugada comprimida, se resuelve
\(Y(x,y)=(f,0)\). La segunda ecuación da
\(y=-C^{-1}F^*x\), y la primera da
\(x=Y_{\rm ef}^{-1}f\). El bloque de frontera de \(Y^{-1}\) es,
por tanto, \(Y_{\rm ef}^{-1}\). Así,
\begin{equation}
 \Lambda_{\rm b}=LY_{\rm ef}^{-1},\qquad
 R_{\rm ef}\Lambda_{\rm b}=L^2I,\qquad
 R_{\rm ef}=\frac{L}{E_{L}}H_{\rm ef}.
 \label{g26:eq:schur-reciproco-es}
\end{equation}
La reducción actúa sobre los mismos bloques en las dos lecturas y
conserva el residuo \(\eta\). El carácter efectivo puede depender de
la memoria de forma matricial y el coeficiente \(L/E_{L}\) permanece fijo.

\subsection{Memoria dinámica y norma temporal inducida}

La eliminación dinámica conserva una dependencia espectral adicional.
En una representación estacionaria, sea
\[
 H=\begin{pmatrix}H_{\rm bb}&V\\V^*&H_{\rm ii}\end{pmatrix},
 \qquad H_{\rm ii}>0.
\]
Para \(z\) en el dominio resolvente pertinente, la eliminación de la
componente interior en \((H-zI)(b,y)=(f,0)\) da
\begin{equation}
 \mathcal K_H(z)=H_{\rm bb}-zI
       -V(H_{\rm ii}-zI)^{-1}V^*.
 \label{g26:eq:resolvente-memoria-es}
\end{equation}
Cuando existen ambos inversos, el bloque de frontera de
\((H-zI)^{-1}\) es \(\mathcal K_H(z)^{-1}\). Para
\(|z|<\lambda_{\min}(H_{\rm ii})\), la serie convergente del
resolvente proporciona
\begin{align}
 \mathcal K_H(z)&=H_{\rm est}-zZ-z^2VH_{\rm ii}^{-3}V^*-\cdots,
 \label{g26:eq:expansion-resolvente-es}\\
 H_{\rm est}&=H_{\rm bb}-VH_{\rm ii}^{-1}V^*,\qquad
 Z=I+VH_{\rm ii}^{-2}V^*\geq I.
 \label{g26:eq:norma-z-es}
\end{align}
La positividad se sigue de
\(Z-I=(H_{\rm ii}^{-1}V^*)^*(H_{\rm ii}^{-1}V^*)\).
La extensión estática \((b,-H_{\rm ii}^{-1}V^*b)\) tiene norma
cuadrada \(b^*Zb\); la energía y la norma temporal proceden de la
misma extensión. El parámetro físico es \(z=\hbar_{\rm ret}\omega\).

Al primer orden temporal, \(T=Z^{1/2}\) y \(\psi=Tb\) transforman
la pareja \((H_{\rm est},Z)\) en
\((T^{-*}H_{\rm est}T^{-1},I)\). La sección
\(\hbar_{\rm ret}\) permanece fijada. Si \(T\) depende del tiempo,
la transformación conserva también el término \(\dot T b\) de
\(\dot\psi=\dot T b+T\dot b\).

La proporcionalidad radial se conserva en el resolvente completo.
Escribiendo \(\chi_{\rm rad}=L/E_{L}\) y \(R=\chi_{\rm rad} H\), la sustitución de
cada bloque demuestra
\begin{equation}
 \mathcal K_R(\chi_{\rm rad} z)=\chi_{\rm rad}\mathcal K_H(z).
 \label{g26:eq:resolvente-proporcional-es}
\end{equation}
Esta identidad transporta simultáneamente el operador y el parámetro
espectral, y conserva cada orden de la memoria dinámica. Para las formas
efectivas de \eqref{g26:eq:schur-reciproco-es}, la normalización canónica
requiere el transporte dual de la longitud inversa:
\begin{align*}
 H_{\rm can}&=T^{-*}H_{\rm ef}T^{-1},&
 R_{\rm can}&=T^{-*}R_{\rm ef}T^{-1},&
 \Lambda_{\rm can}&=T\Lambda_{\rm b}T^*.
\end{align*}
Multiplicar estas expresiones da
\[
 R_{\rm can}=\frac{L}{E_{L}}H_{\rm can},\qquad
 R_{\rm can}\Lambda_{\rm can}=L^2I.
\]
La demostración admite \(T\), \(R_{\rm ef}\) y
\(\Lambda_{\rm b}\) no conmutativos. La norma temporal inducida y
la longitud dual conservan conjuntamente el mismo \(G\).

\subsection{Refinamiento, dominios espectrales y alcance}

Sean \(J\) y \(K\) inscripciones isométricas entre dos niveles que
satisfagan \(X'J=JX\) y \(\mathsf U'J=K\mathsf U\),
con la misma sección intensiva \(L\). Entonces
\[
 R'K=L\mathsf U'X'\mathsf U'^*K
     =KL\mathsf U X\mathsf U^*=KR,
\]
y la misma cuenta demuestra \(H'K=KH\). El entrelazamiento de las
inversas se formula en los dominios correspondientes. Cuando el sistema
compatible posee un límite autoadjunto positivo e inyectivo \(X\),
el cálculo espectral define \(Y=\mathsf U X\mathsf U^*\),
\(R=LY\), \(H=E_{L}Y\) y \(\Lambda=LY^{-1}\). La proporcionalidad
\(R=(L/E_{L})H\) vale en \(\operatorname{Dom}(Y)\), mientras
\(R\Lambda=L^2I\) y \(H\Lambda=\hbar_{\rm ret}cI\) valen en
\(\operatorname{Dom}(Y^{-1})\), con prolongación acotada de sus
productos. Los cortes espectrales
\(\mathbf1_{[\varepsilon,M]}(Y)\) verifican las identidades antes
de pasar al límite en estos dominios. Una longitud de arista refinada
\(L/9\) incorpora su factor de escala propio; la sección intensiva
se conserva según el entrelazamiento declarado.

La clasificación distingue las transformaciones que preservan la
realización de las que modifican una regla constitutiva. Las ponderaciones
diagonales, las equivalencias unitarias, el archivo auxiliar y la reducción
conjunta conservan \(G\). El cambio de resolución marcada modifica R1;
la sustitución de \(r_N\) por \(\sqrt{r_N}\), de
\(\alpha^{16}\) por su cuadrado o de la componente inscrita por la
norma total modifica R4. La reescala aislada del radio cambia R5,
y la reescala aislada de la acción cambia R6. R8 mantiene la distinción
entre radio reducido, radio de Schwarzschild y cambio de unidades.

Se obtiene así la unicidad matemática de la realización R1--R8: los datos
generados y el carácter fijan la respuesta, y la composición común fija
\(G\). Las reglas longitudinal y métrica figuran explícitamente en
esta conclusión como composiciones incorporadas; la identificación
gravitatoria conserva su estatuto físico. La prueba abarca la clase
declarada con sus definiciones constitutivas y sus transformaciones
compatibles. Las comprobaciones racionales del desarrollo verifican
casos finitos de las identidades expuestas; la demostración general
reside en las ecuaciones y argumentos anteriores.
% END CR28_RADIAL_FULL


% BEGIN CR28_RADIAL_CONSTITUTIVE_EVALUATION
\section{Continuidad constitutiva y evaluación de la sección gravitatoria}
\label{cr28:sec:radial-evaluation}

La construcción longitudinal determina \(L\) antes de utilizar la
identificación gravitatoria. La composición de las reglas R1--R8 conserva
la sección de velocidad de \eqref{cr28:eq:same-velocity} y la acción
de retorno previamente generada. Con
\[
 L=\frac{5759}{23040}\alpha^{16}L_*,
 \qquad \hbar_{\rm ret}=\eta_{\rm ret}S_*,
 \qquad c_{\rm int}=c_*\widehat c,
\]
el resultado de \eqref{g26:eq:g-rigido-es} se expresa como
\begin{equation}
 G_{\rm ret}
 =\frac{c_*^3L_*^2}{S_*}
   \left(\frac{5759}{23040}\right)^2
   \frac{\alpha^{32}\widehat c^{\,3}}{\eta_{\rm ret}}
 =\frac{L^2}
 {\hbar_{\rm ret}(\mu_{\rm vac}\varepsilon_{\rm vac})^{3/2}}.
 \label{cr28:eq:g-constitutive-evaluation}
\end{equation}
Los canales son los mismos de \eqref{eq:vacuum-angular-channels}:
\[
 q_\pm=\exp\!\left[-\frac{\pi}{180}(A_{\rm deg}\pm C^*_{\rm deg})\right],
 \qquad r_\pm=\frac{1-q_\pm^{90}}{1-q_\pm^{120}},
 \qquad \widehat c=(r_+r_-)^{-1}.
\]
Las bases constitutivas satisfacen
\(\varepsilon=\varepsilon_*r_+^2\),
\(\mu=\mu_*r_-^2\) y
\(c_*=(\mu_*\varepsilon_*)^{-1/2}\). Su producto prueba
\(c_{\rm int}=(\mu\varepsilon)^{-1/2}\). Si se declara directamente
la coordenada dimensional de \(c_{\rm int}\), su base se recupera como
\(c_*=c_{\rm int}/\widehat c\). El coeficiente constitutivo no vuelve
a multiplicarse: ambas expresiones representan la misma sección.

En la representación \(L_*=1\,\mathrm m\),
\(S_*=10^{-34}\,\mathrm{J\,s}\) y
\(c_{\rm int}=299792458\,\mathrm{m\,s^{-1}}\), la evaluación de los
lectores y de la composición radial da
\begin{equation}
 \boxed{
 G_{\rm ret}
 =6.674299659414022706367776189\ldots\times10^{-11}
 \,\mathrm{m^3\,kg^{-1}\,s^{-2}}.}
 \label{cr28:eq:g-evaluation}
\end{equation}
El valor de contraste
\(G_{\rm CODATA\,2022}=6.67430(15)\times10^{-11}\,
\mathrm{m^3\,kg^{-1}\,s^{-2}}\) conserva su incertidumbre
experimental. En las mismas unidades, la diferencia de la evaluación
respecto del valor central es
\(-3.40585977\ldots\times10^{-18}\), aproximadamente
\(-0.00227057\) incertidumbres estándar. Las cifras de la evaluación
corresponden a los lectores y bases declarados; la incertidumbre pertenece
al resultado medido. Ninguna cifra de la referencia selecciona
\(N\), \(r_N\), la regla longitudinal o la norma radial.

\subsection{Realización cronológica del mismo radio construido}
\label{cr28:sec:radial-clock}
Con la duración elemental
\begin{equation}
 t_0=\frac{\pi L}{54c_{\rm int}},
 \qquad \ell_0=c_{\rm int}t_0=\frac{\pi L}{54},
 \label{cr28:eq:clock-from-length}
\end{equation}
el calendario de \(108\) posiciones satisface
\[
 T_{108}=108t_0=\frac{2\pi L}{c_{\rm int}},\qquad
 \omega_{108}=\frac{c_{\rm int}}{L},\qquad
 R_{108}=\frac{c_{\rm int}}{\omega_{108}}=L.
\]
Por sustitución,
\begin{equation}
 \left(\frac{54}{\pi}\right)^2
 \frac{c_{\rm int}^5t_0^2}{\hbar_a}
 =\frac{c_{\rm int}^3L^2}{\hbar_a}.
 \label{cr28:eq:circular-radial-agreement}
\end{equation}
Esta identidad sitúa el selector circular conservado a continuación en
la misma normalización longitudinal. Su teorema de unicidad conserva la
premisa geométrica que enuncia; la construcción anterior especifica además
la longitud utilizada en esa realización.

El carácter positivo \(Y=I\) da simultáneamente
\(R_X=\Lambda_X=L\) y \(M_X=\hbar_a/(c_{\rm int}L)\).
Para un carácter general \(Y>0\), las longitudes son recíprocas:
\(R_X=LY\), \(\Lambda_X=LY^{-1}\).
Su producto permanece \(L^2I\), sin que sus valores individuales deban
coincidir. Así, la identificación de las cuatro longitudes del caso
circular corresponde a un sector preciso de la prueba radial y no
suprime la dependencia del carácter.

La proporcionalidad \(R_X=(G/c_{\rm int}^{2})M_X\) es común a
todos los caracteres del dominio. Su conservación bajo cambio de base,
refinamiento, minimización con memoria, complemento de Schur y
normalización dinámica se ha demostrado en
\cref{g26:sec:rigidez-radial-es}. El portador pentacomponente y su
reducción a dos helicidades, desarrollados después, describen la
representación tensorial; no reemplazan la demostración del coeficiente
radial común.
% END CR28_RADIAL_CONSTITUTIVE_EVALUATION


\section{Cadena constitutiva del sector gravitatorio: escala cronológica \(108\), reducción tensorial, holonomía y evaluación metrológica de \(G\)}

La gravedad se organiza aqu\'i en cuatro estratos que no deben fundirse:

\begin{equation}
\label{eq:gravity-causal-chain}
\begin{aligned}
\text{inscripci\'on y norma radial R1--R8}
&\longrightarrow L
\longrightarrow G=c_{\rm int}^3L^2/\hbar_a,\\
\text{incidencia icosa\'edrica}
&\longrightarrow V_{12}\xrightarrow{P_5}V_5
\xrightarrow{\Gamma_5}\operatorname{STF}_3
\xrightarrow{\Lambda(n)}\mathcal H_{\rm TT}(n),\\
\text{calendario CT108 con }t_0=\pi L/(54c_{\rm int})
&\longrightarrow R_{108}=L,\quad
\theta_{108}^{\rm circ}=(54/\pi)^2,\\
\text{carta metrol\'ogica}
&\longrightarrow 6.6742996594140227\ldots\times10^{-11}\,
\mathrm{m^3\,kg^{-1}\,s^{-2}}.
\end{aligned}
\end{equation}

La primera l\'inea es la determinaci\'on demostrada en
\cref{g26:sec:rigidez-radial-es}: la longitud precede al acoplamiento.
La segunda construye el portador y su reducci\'on transversal sin traza.
La tercera conserva el calendario como realizaci\'on de esa misma longitud,
seg\'un \eqref{cr28:eq:circular-radial-agreement}. La cuarta eval\'ua
la secci\'on en las bases de \cref{cr28:sec:radial-evaluation}.
Las ramas alternativas de holonom\'ia y el lector decimal se conservan
con sus dominios propios; no reemplazan la prueba radial.

\section{Recta de magnitud gravitatoria y escala cronol\'ogica \(108\)}

Sean \(c_{\rm int}>0\), \(t_0>0\),
\(\ell_0=c_{\rm int}t_0\) y una secci\'on positiva
\(\hbar_a\) de la recta de acci\'on, con
\(a\in\{\mathrm{pre},\mathrm{ret}\}\). El transductor dimensional es

\begin{equation}
\label{eq:gravity-transducer}
G_a(\theta)
=\theta\frac{c_{\rm int}^{5}t_0^2}{\hbar_a}
=\theta\frac{c_{\rm int}^{3}\ell_0^2}{\hbar_a},
\qquad \theta>0.
\end{equation}

La igualdad de las dos expresiones usa s\'olo
\(\ell_0=c_{\rm int}t_0\), y
\([G]=L^3M^{-1}T^{-2}\). La dimensi\'on queda fijada por la recta; el
car\'acter adimensional \(\theta\) distingue las realizaciones de esa
familia. En la composici\'on longitudinal ya demostrada,
\(t_0=\pi L/(54c_{\rm int})\) y
\(\theta=(54/\pi)^2\) reproducen exactamente
\(G_a=c_{\rm int}^3L^2/\hbar_a\).

El calendario CT108 define

\begin{equation}
\label{eq:gravity-clock108}
T_{108}=108t_0,
\qquad
\omega_{108}=\frac{2\pi}{108t_0},
\qquad
R_{108}=\frac{c_{\rm int}}{\omega_{108}}
=\frac{54}{\pi}\ell_0.
\end{equation}

Para cada hoja de acci\'on, sean

\begin{equation}
\label{eq:gravity-clock-energy-mass}
E_{108,a}=\hbar_a\omega_{108},
\qquad
m_{108,a}=\frac{E_{108,a}}{c_{\rm int}^2}.
\end{equation}

Entonces, la longitud de Compton reducida coincide con el radio del ciclo:

\begin{equation}
\label{eq:gravity-compton-clock}
\bar\lambda_{C,a}
=\frac{\hbar_a}{m_{108,a}c_{\rm int}}
=R_{108}.
\end{equation}

\begin{theorem}[Unicidad del selector bajo la condici\'on de periodicidad HMT]
\label{thm:gravity-circular-selector}
Con la convenci\'on reducida
\(r_g=Gm/c_{\rm int}^2\) y la premisa de periodicidad HMT que identifica el radio
gravitatorio del cuanto CT108 con su radio de fase, la condici\'on de cierre
\begin{equation}
\label{eq:gravity-closing-condition}
r_{g,a}(\theta)=R_{108}
\end{equation}
selecciona una \`unica soluci\'on positiva,
\begin{equation}
\label{eq:gravity-theta108}
\boxed{
\theta_{108}^{\rm circ}=\left(\frac{54}{\pi}\right)^2
=295.4525715010569415303525\ldots .}
\end{equation}
En esa realizaci\'on,
\begin{equation}
\label{eq:gravity-four-lengths}
R_{108}=\bar\lambda_{C,a}=r_{g,a}=\ell_P^{\rm circ}.
\end{equation}
\end{theorem}

\begin{proof}
Por \cref{eq:gravity-transducer,eq:gravity-clock-energy-mass},
\[
r_{g,a}(\theta)
=\frac{G_a(\theta)m_{108,a}}{c_{\rm int}^2}
=\theta c_{\rm int}t_0^2\omega_{108}
=\theta\frac{\pi}{54}\ell_0.
\]
Al compararlo con
\(R_{108}=(54/\pi)\ell_0\) se obtiene necesariamente
\(\theta=(54/\pi)^2\). La ecuaci\'on es lineal en \(\theta\) y todos los
factores son positivos, luego la soluci\'on es \`unica.
\end{proof}

La igualdad \(r_g=R_{108}\) no se deduce de la dimensionalidad de
\cref{eq:gravity-transducer}. Es la condición geométrica explícita de esta
realizaci\'on HMT: acci\'on, Compton y retorno peri\'odico deben determinar una misma
longitud. El teorema demuestra que, aceptada esa condición estructural, el
carácter \(\theta\) queda fijado sin usar el valor SI de \(G\).

\begin{remark}[Convenci\'on de radio]
El teorema usa el radio gravitatorio reducido \(Gm/c^2\). Si se sustituye
por el radio de Schwarzschild \(2Gm/c^2\), el selector cambia por un factor
dos. Las dos expresiones corresponden a condiciones geom\'etricas distintas,
por lo que la convenci\'on debe declararse antes del c\'alculo.
\end{remark}

\section{Reciprocidad \(G\leftrightarrow\hbar\) y conservación del área de Planck}

\begin{corollary}[Familia de Planck asociada a la periodicidad temporal de \(108\) fases]
\label{cor:gravity-planck-family}
Para \(G_a=G_a(\theta_{108}^{\rm circ})\),
\begin{align}
\ell_P&=\frac{54}{\pi}\ell_0,
&t_P&=\frac{54}{\pi}t_0,
\label{eq:gravity-planck-length-time}\\
E_{P,a}&=\frac{\hbar_a}{t_P}
=\frac{\pi\hbar_a}{54t_0}
=\frac{h_a}{108t_0},
&\ell_P^2&=\frac{G_a\hbar_a}{c_{\rm int}^3}
=\theta_{108}^{\rm circ}\ell_0^2.
\label{eq:gravity-planck-energy-area}
\end{align}
Adem\'as,
\begin{equation}
\label{eq:gravity-planck-force-power-density}
F_{P,a}=\frac{c_{\rm int}^4}{G_a},
\qquad
P_{P,a}=\frac{c_{\rm int}^5}{G_a},
\qquad
\rho_{P,a}=\frac{\hbar_a}
{(\theta_{108}^{\rm circ})^2c_{\rm int}^5t_0^4}.
\end{equation}
\end{corollary}

\begin{proof}
Las tres primeras identidades proceden de
\(\ell_P=R_{108}\), \(t_P=\ell_P/c_{\rm int}\) y
\(h_a=2\pi\hbar_a\). Para el \`area,
\[
\frac{G_a\hbar_a}{c_{\rm int}^3}
=\theta_{108}^{\rm circ}c_{\rm int}^2t_0^2
=\theta_{108}^{\rm circ}\ell_0^2.
\]
Las expresiones restantes son las composiciones dimensionales de fuerza,
potencia y masa por volumen con la misma secci\'on.
\end{proof}

Sea
\begin{equation}
\label{eq:gravity-Ract}
R_{\rm act}:=\frac{\hbar_{\rm pre}}{\hbar_{\rm ret}}.
\end{equation}
Como la acci\'on aparece en el numerador de la masa cronol\'ogica y en el
denominador de \(G\),
\begin{equation}
\label{eq:gravity-twoway-covariance}
\frac{m_{108,\rm pre}}{m_{108,\rm ret}}=R_{\rm act},
\qquad
\frac{G_{\rm pre}}{G_{\rm ret}}=R_{\rm act}^{-1},
\qquad
G_a\hbar_a
=\theta_{108}^{\rm circ}c_{\rm int}^5t_0^2.
\end{equation}
Por tanto, \(\ell_P^2=G_a\hbar_a/c^3\) no cambia de hoja. La gravedad y
la acci\'on se orientan rec\'iprocamente, mientras el \`area conserva la
geometr\'ia com\'un. Esta cancelaci\'on es la interfaz de confluencia hacia el operador de
\`area del \cref{ch:modular}.
 La segunda relación fundamental aparece cuando se comparan acción y gravedad.\par

En HMT, la acción mide cuánto ocurre y conserva simultáneamente la orientación de la realización. Hay una sección anterior al retorno y otra posterior. La fase puede volver acompañada por una historia de memoria diferente. Esa diferencia queda almacenada en la razón bidireccional de las dos secciones de acción.\par

La gravedad responde de manera recíproca.\par

El transductor gravitatorio tiene la forma\par

\[
G_a(\theta)
=
\theta\frac{c_{\rm int}^{5}t_0^2}{\hbar_a}.
\]

La acción aparece en el denominador. Por tanto, si al cambiar de hoja \(\hbar\) se multiplica por una razón \(R_{\rm act}\), la constante gravitatoria se multiplica por la razón inversa:\par

\[
\frac{\hbar_{\rm pre}}{\hbar_{\rm ret}}
=
R_{\rm act},
\qquad
\frac{G_{\rm pre}}{G_{\rm ret}}
=
R_{\rm act}^{-1}.
\]

La acción y la gravedad forman aquí una pareja contragrediente: cuando una avanza en una dirección, la otra compensa exactamente ese desplazamiento.\par

Lo que permanece invariante es\par

\[
G_a\hbar_a.
\]

Y ésa es precisamente la combinación que determina el área de Planck:\par

\[
\ell_P^2=\frac{G_a\hbar_a}{c_{\rm int}^3}.
\]

La acción recuerda la hoja. La gravedad responde a esa memoria. El área permanece invariante.\par

Ésta es una idea extraordinariamente fértil. El área de Planck emerge como el objeto que sobrevive cuando acción y gravedad se transforman recíprocamente.\par

Podría decirse así: \(\hbar\) conserva la orientación dinámica; \(G\) conserva la compensación geométrica; \(\ell_P^2\) conserva aquello en lo que ambas realizaciones están de acuerdo.\par

El área es, por tanto, una moneda de cambio entre física cuántica y geometría, porque la transformación bidireccional de \(\hbar\) y \(G\) selecciona exactamente esa combinación como invariante.\par

Esta reciprocidad permite comprender de otra forma la relación entre información y gravedad. La información necesita una distinción de estados; la acción mide el coste de transición entre ellos. Pero si esa transición debe adquirir una geometría, la respuesta gravitatoria cambia de manera inversa para mantener estable el cuanto de área.\par

La geometría incorpora la información de fase mediante una compensación exacta.\par

Desde una lectura machiana, éste es ya un resultado muy sugestivo. Una magnitud local de geometría queda determinada relacionalmente por la orientación y el retorno del estado completo. El área invariante restringe esa dependencia relacional y fija su consistencia.\par

La versión HMT del principio relacional consiste en algo más preciso que la máxima «la materia determina la geometría»: la orientación de la acción y la respuesta gravitatoria se codeterminan de tal manera que la geometría elemental permanezca invariante.\par

La acción contiene la historia. La gravedad ajusta la respuesta. El área conserva el acuerdo.\par

 La derivación de \(G\) lleva esa reciprocidad a una condición geométrica mucho más concreta.\par

La Mecánica Dimensional determina primero qué tipo de objeto puede ser \(G\). Fija su recta de magnitud y su homogeneidad dimensional. El cierre cronológico selecciona después su valor: la dimensión determina la clase de magnitud y el cierre determina su coeficiente adimensional.\par

La selección procede del periodicidad temporal de (108) fases.\par

El ciclo \(108\) determina una frecuencia, un radio de fase y una energía. A partir de esa energía aparece una masa cronológica. Y al construir la longitud de Compton reducida de esa masa se obtiene exactamente el mismo radio del ciclo:\par

\[
\bar\lambda_C=R_{108}.
\]

La fase y la materia ya están diciendo la misma longitud antes de introducir \(G\).\par

Entonces entra la gravedad. Se construye el radio gravitatorio reducido\par

\[
r_g=\frac{Gm}{c^2},
\]

y se exige que el cierre gravitatorio identifique el radio de fase, la longitud de Compton y el radio gravitatorio como una sola escala estructural:\par

\[
r_g=R_{108}=\bar\lambda_C.
\]

Esa condición selecciona unívocamente\par

\[
\theta_{108}
=
\left(\frac{54}{\pi}\right)^2.
\]

De este modo,\par

\[
R_{108}
=
\bar\lambda_C
=
r_g
=
\ell_P.
\]

La función estructural de \(\theta_{108}\) determina la significación de su decimal terminal: muestra qué está haciendo \(G\).\par

\(G\) aparece como el coeficiente necesario para hacer coincidir las cuatro longitudes publicadas por estos lenguajes:\par

\begin{itemize}[leftmargin=*,itemsep=0.35em]
\item el lenguaje de la fase produce \(R_{108}\);
\item el lenguaje cuántico produce \(\bar\lambda_C\);
\item el lenguaje gravitatorio produce \(r_g\);
\item el lenguaje de Planck produce \(\ell_P\).
\end{itemize}

La constante gravitatoria es el mediador que impide que esos cuatro lectores se separen.\par

Dicho de una manera más física: \(G\) expresa cuánto debe responder la geometría para que un cuanto cuya masa procede del reloj \(108\) cierre sobre su propia fase.\par

Es la relación que hace posible que materia, fase y geometría sean tres publicaciones compatibles del mismo acontecimiento.\par

Aquí la lectura machiana se vuelve más fuerte. La escala gravitatoria local queda seleccionada por una condición de cierre global del ciclo. El radio local de un cuanto adquiere sentido al integrarse en una periodicidad completa. La inercia, la fase y la geometría se determinan relacionalmente.\par

El alcance demostrado proporciona un mecanismo matemático concreto para una lectura machiana: una propiedad local —el acoplamiento gravitatorio— queda seleccionada por la consistencia del estado global. Las demás formulaciones históricas del principio de Mach conservan su dominio propio.\par

Después aparece la dimensión cinco.\par

El escalar \(G\) fija la intensidad gravitatoria y el portador especifica su estructura representacional. HMT construye un espacio tensorial simétrico de traza nula y dimensión cinco. Ese portador admite cinco pesos helicoidales:\par

\[
+2,\quad +1,\quad 0,\quad -1,\quad -2.
\]

Son los cinco grados de libertad naturales de un tensor simétrico sin traza en tres dimensiones antes de imponer las restricciones dinámicas de una teoría concreta.\par

Eso aclara un punto que suele confundirse. «Cinco componentes» designa el portador completo, mientras las dos helicidades designan la reducción transversal de masa nula en relatividad general.\par

Significa que el portador completo tiene cinco coordenadas irreducibles. Después, la dinámica decide qué parte de ese portador es física:\par

\begin{itemize}[leftmargin=*,itemsep=0.35em]
\item en una realización masiva de espín dos pueden sobrevivir los cinco pesos;
\item en una realización sin masa con invariancia de calibre, la proyección transversal sin traza elimina los sectores \(\pm1\) y \(0\);
\item permanecen las helicidades \(\pm2\).
\end{itemize}

Así, HMT proporciona un espacio previo a la decisión dinámica. Construye el portador completo de cinco componentes y contiene, mediante la proyección correspondiente, su reducción física a dos helicidades.\par

La belleza reside en la coexistencia ordenada de dos afirmaciones:\par

\begin{itemize}[leftmargin=*,itemsep=0.35em]
\item la geometría de espín dos tiene un portador pentacomponente;
\item la radiación gravitatoria sin masa ocupa un subespacio bidimensional.
\end{itemize}

Ambas afirmaciones forman una reducción coherente y precisa.\par

Y esa reducción conserva otra enseñanza central: el escalar \(G\) fija la intensidad común y la información de orientación vive en el portador tensorial. La intensidad reside en la constante, las helicidades residen en la representación y la genealogía completa reúne magnitud y portador.\par

  \endgroup
% Secciones 1--3 del delta gravitatorio íntegro.
\section{Dominio genealógico de la realización gravitatoria}
\label{sec:l9s102:gravedad-dominio-genealogico}

La realización gravitatoria recibe como objeto de partida una sección del
continuo discreto ya construida por la composición
\begin{equation}
 \mathrm{APP}\longrightarrow\mathrm{TRIT}\longrightarrow\mathrm{TPK}
 \longrightarrow X_{\mathrm{TPK}}^{\mathrm{enr}}
 \longrightarrow\mathcal C_{\mathrm{cont}}^{\mathrm{disc}}.
 \label{eq:l9s102:gravedad-cadena-genealogica}
\end{equation}
APP conserva las evaluaciones aditiva y multiplicativa mediante residuo,
cociente y acarreo; TRIT determina el régimen y la orientación; el TPK
incorpora el selector trítico de fase sin eliminar ninguna de las dos hojas
APP, efectúa el transporte cardinal, conserva la memoria y realiza el retorno
de fase sin reiniciar el estado. El transductor
dimensional actúa después sobre la sección superviviente del continuo:
\begin{equation}
 \mathfrak G_{108}:\mathcal C_{\mathrm{cont}}^{\mathrm{disc}}
 \times\mathcal L_S^+
 \longrightarrow\mathcal G_{\mathrm{HMT\text{-}MD}},
 \qquad
 (x,\hbar_a)\longmapsto
 \bigl(G_a,R_{108},m_{108,a},\ell_{P,a}^{,2}\bigr).
 \label{eq:l9s102:gravedad-mapa-realizacion}
\end{equation}
Los cardinales \(90/120\), la periodicidad \(108\), la acción y el par
angular entran con su genealogía conservada. Ninguna identificación
gravitatoria posterior selecciona retroactivamente estos antecedentes.

\section{Transductor gravitatorio y selector cronológico de periodicidad}
\label{sec:l9s102:gravedad-transductor}

Sean \(c_{\mathrm{int}}>0\), \(t_0>0\),
\(\ell_0=c_{\mathrm{int}}t_0\) y
\(\hbar_a\in\mathcal L_S^+\), con
\(a\in\{\mathrm{pre},\mathrm{ret}\}\). El transductor dimensional es
\begin{equation}
 G_a(\theta)
 =\theta\frac{c_{\mathrm{int}}^5t_0^2}{\hbar_a}
 =\theta\frac{c_{\mathrm{int}}^3\ell_0^2}{\hbar_a},
 \qquad \theta>0.
 \label{eq:l9s102:g-transductor}
\end{equation}
La cronologia ciclica de \(108\) estados determina
\begin{equation}
 T_{108}=108t_0,
 \qquad
 \omega_{108}=\frac{2\pi}{108t_0},
 \qquad
 R_{108}=\frac{c_{\mathrm{int}}}{\omega_{108}}
 =\frac{54}{\pi}\ell_0.
 \label{eq:l9s102:g-reloj108}
\end{equation}
El cuanto cronológico y su masa asociada son
\begin{equation}
 E_{108,a}=\hbar_a\omega_{108},
 \qquad
 m_{108,a}=\frac{E_{108,a}}{c_{\mathrm{int}}^2}.
 \label{eq:l9s102:g-cuanto-masa}
\end{equation}
La longitud de Compton reducida satisface
\begin{equation}
 \bar\lambda_{C,a}
 =\frac{\hbar_a}{m_{108,a}c_{\mathrm{int}}}=R_{108}.
 \label{eq:l9s102:g-compton}
\end{equation}

\begin{theorem}[Unicidad del selector gravitatorio de periodicidad]
\label{thm:l9s102:g-selector}
Con la convencion reducida \(r_g=Gm/c_{\mathrm{int}}^2\), la condición
\(r_{g,a}(\theta)=R_{108}\) posee la solución positiva única
\begin{equation}
 \boxed{
 \theta_{108}^{\mathrm{circ}}=\left(\frac{54}{\pi}\right)^2.}
 \label{eq:l9s102:g-theta}
\end{equation}
La realización correspondiente identifica estructuralmente
\begin{equation}
 R_{108}=\bar\lambda_{C,a}=r_{g,a}=\ell_P^{\mathrm{circ}}.
 \label{eq:l9s102:g-cuatro-longitudes}
\end{equation}
\end{theorem}

\begin{proof}
Por \cref{eq:l9s102:g-transductor,eq:l9s102:g-cuanto-masa},
\[
 r_{g,a}(\theta)
 =\frac{G_a(\theta)m_{108,a}}{c_{\mathrm{int}}^2}
 =\theta\frac{\pi}{54}\ell_0.
\]
La igualdad con \(R_{108}=(54/\pi)\ell_0\) impone
\(\theta=(54/\pi)^2\). La linealidad en \(\theta\) y la positividad de
los factores determinan la unicidad.
\end{proof}

En la carta en la que \(c_{\mathrm{int}}=(\mu\varepsilon)^{-1/2}\), el
transductor adopta la forma
\begin{equation}
 \boxed{
 G_a=\left(\frac{54}{\pi}\right)^2
 \frac{t_0^2}{\hbar_a(\mu\varepsilon)^{5/2}}.}
 \label{eq:l9s102:g-mu-epsilon}
\end{equation}
El producto constitutivo \(\mu\varepsilon\) gobierna el cono causal y la
escala gravitatoria; el cociente \(Z^2=\mu/\varepsilon\) conserva la
orientación de impedancia.

\section{Covariancia bidireccional de gravedad y acción}
\label{sec:l9s102:g-hbar}

Definimos
\begin{equation}
 R_{\mathrm{act}}=\frac{\hbar_{\mathrm{pre}}}
 {\hbar_{\mathrm{ret}}}.
 \label{eq:l9s102:r-act}
\end{equation}
La masa cronológica y el acoplamiento gravitatorio transforman
contragredientemente:
\begin{equation}
 \frac{m_{108,\mathrm{pre}}}{m_{108,\mathrm{ret}}}=R_{\mathrm{act}},
 \qquad
 \frac{G_{\mathrm{pre}}}{G_{\mathrm{ret}}}=R_{\mathrm{act}}^{-1}.
 \label{eq:l9s102:g-hbar-reciprocidad}
\end{equation}
Su producto permanece fijo:
\begin{equation}
 \boxed{
 G_a\hbar_a
 =\theta_{108}^{\mathrm{circ}}c_{\mathrm{int}}^5t_0^2,
 \qquad
 \ell_P^2=\frac{G_a\hbar_a}{c_{\mathrm{int}}^3}
 =\theta_{108}^{\mathrm{circ}}\ell_0^2.}
 \label{eq:l9s102:g-hbar-invariante}
\end{equation}
La acción conserva la orientación de la hoja; el área de Planck conserva la
geometría común de ambas orientaciones.

\begin{theorem}[Independencia de carta de las secciones \(h\) y \(G\)]
\label{thm:l9s102:covariancia-unidades-g}
Sean \(u_L,u_T,u_S\) bases positivas de las rectas dimensionales de
longitud, tiempo y acción, y cámbiense por
\[
 u_L'=\lambda_Lu_L,\qquad
 u_T'=\lambda_Tu_T,\qquad
 u_S'=\lambda_Su_S,
 \qquad \lambda_L,\lambda_T,\lambda_S>0.
\]
Las coordenadas de las mismas secciones verifican
\begin{equation}
 c_{\rm int}'=\frac{\lambda_T}{\lambda_L}c_{\rm int},\qquad
 t_0'=\frac{t_0}{\lambda_T},\qquad
 \hbar_a'=\frac{\hbar_a}{\lambda_S},\qquad
 G_a'=\frac{\lambda_S\lambda_T^3}{\lambda_L^5}G_a.
 \label{eq:l9s102:cambio-carta-g}
\end{equation}
Para la base inducida
\(u_G=u_L^5u_T^{-3}u_S^{-1}\) se obtiene
\begin{equation}
 \hbar_a'u_S'=\hbar_au_S,
 \qquad G_a'u_G'=G_au_G.
 \label{eq:l9s102:secciones-h-g-invariantes}
\end{equation}
Por consiguiente,
\(h_a=2\pi\hbar_a\),
\(G_a=\theta_{108}^{\rm circ}c_{\rm int}^5t_0^2/\hbar_a\) y
\(G_a\hbar_a=\theta_{108}^{\rm circ}c_{\rm int}^5t_0^2\)
son identidades de secciones, no igualdades dependientes de una elección de
unidades. La carta metrológica publica sus coordenadas numéricas después de
la construcción HMT--MD.
\end{theorem}

\begin{proof}
Las tres primeras transformaciones de
\eqref{eq:l9s102:cambio-carta-g} son la ley contragrediente de coordenadas
en \(L\otimes T^{-1}\), \(T\) y \(S\). Su sustitución en
\eqref{eq:l9s102:g-transductor} produce la cuarta. La base de la recta
gravitatoria se transforma por el factor inverso, lo que prueba
\eqref{eq:l9s102:secciones-h-g-invariantes}. Los escalares
\(\pi\) y \(\theta_{108}^{\rm circ}\) son adimensionales y
permanecen invariantes.
\end{proof}
 
\chapter{Parámetro de Barbero--Immirzi, incidencia hexádico--octádica y espectro del operador de área}
\begingroup
\renewcommand{\chapter}[2][]{}%
\chapter{Derivación HMT del parámetro de Barbero--Immirzi y del área mínima: incidencia hexada--octada, canales \(90/120\) y espectro del operador de área}
\label{chap:p3-final:51:barbero_area}
\section{Morfismo de incidencia hexada--octada y conservación de los canales orientados \(90/120\)}

Fijados la dodecada y el alineamiento del
\cref{sec:w12-w24-elevacion}, sea \(H\) una hexada en las coordenadas HMT,
sea \(\ell(H)\subset D\) su imagen y sea
\(O_H=\iota_D(\ell(H))\) su elevación única. Al marcar un punto
y un par de puntos de la hexada residual se obtienen
\[
F_6(H)=\ell(H)\times\binom{\ell(H)}2,
\qquad
|F_6|=6\binom62=90.
\]
Si el punto marcado puede recorrer la octada:
\[
F_8(H)=O_H\times\binom{\ell(H)}2,
\qquad
|F_8|=8\binom62=120.
\]
La inclusión \(\ell(H)\hookrightarrow O_H\) induce el morfismo de incidencia
\[
 \mathcal I_{6\to8}:F_6(H)\hookrightarrow F_8(H).
\]
Al normalizar la diferencia por el factor
elegido \(24\binom62\), se obtiene
\[
\boxed{
\frac{90-120}{24\binom62}
=\frac{6-8}{24}
=-\frac1{12}.}
\]
La igualdad es un defecto normalizado de la interfaz combinatoria
\(6\to8\). El factor \(24\binom62\) forma parte de la definición de esta
normalización; la inclusión por sí sola determina la diferencia \(-30\).
El operador \(\mathcal I_{6\to8}\) tiene como dominio configuraciones
marcadas de incidencia; no es el extractor transversal del estado
dodecafásico ni la familia de momentos angulares.
 \label{mdg:chap:barbero}

La transición de una hexada a una octada combina tres estructuras de tipos
distintos: una incidencia finita, dos series de Lambert asociadas a los
canales \(q_\pm\) y una identificación con el parámetro real de la acción de
Holst. Las tres se articulan mediante la incidencia, el funcional angular y
el diccionario areal hacia la geometría de la conexión de
Ashtekar--Barbero desarrollado en esta misma residencia.

\section{Cardinalidades de incidencia 90 y 120}

Sea \(H\) un conjunto de seis puntos y sea \(O\) un conjunto de ocho puntos
que contiene una copia distinguida de \(H\). Consideremos
\[
 \mathcal F_6=H\times\binom H2,
 \qquad
 \mathcal F_8=O\times\binom H2.
\]
La inclusión distinguida \(H\hookrightarrow O\) induce
\[
 \iota_{6,8}:\mathcal F_6\hookrightarrow\mathcal F_8,
 \qquad
 (h,\{h_1,h_2\})\longmapsto(h,\{h_1,h_2\}).
\]
Este morfismo de incidencias es el antecedente combinatorio de los
cardinales \(90/120\); no es el extractor dodecafásico ni una serie angular.

\begin{theorem}[Cardinalidades de la inclusión]
\[
 |\mathcal F_6|=6\binom62=90,
 \qquad
 |\mathcal F_8|=8\binom62=120.
\]
La diferencia normalizada por veinticuatro coordenadas es
\[
 \frac{90-120}{24\binom62}=-\frac1{12}.
\]
\end{theorem}
\begin{proof}
Cada elemento de la primera coordenada puede combinarse con cualquiera de los
quince pares de \(H\). El principio multiplicativo da las dos cardinalidades;
la última identidad resulta de sustituir \(\binom62=15\).
\end{proof}

\begin{figure}[htbp]
\centering
\begin{tikzpicture}[
  font=\small,
  setnode/.style={circle,draw,minimum size=6mm,inner sep=0pt,
    font=\scriptsize\bfseries},
  box/.style={draw,rounded corners=1.5mm,align=center,
    text width=5.0cm,inner xsep=6pt,inner ysep=5pt},
  arr/.style={-{Stealth[length=2.2mm]},line width=.7pt}]

  % Inclusión geométrica. Las flechas parten del contorno y no atraviesan nodos.
  \node[setnode,draw=hmtblue,fill=hmtlightblue] (h1) at (-2.3,1.25) {1};
  \node[setnode,draw=hmtblue,fill=hmtlightblue] (h2) at (-1.1,1.9) {2};
  \node[setnode,draw=hmtblue,fill=hmtlightblue] (h3) at (.1,1.25) {3};
  \node[setnode,draw=hmtblue,fill=hmtlightblue] (h4) at (.1,-.15) {4};
  \node[setnode,draw=hmtblue,fill=hmtlightblue] (h5) at (-1.1,-.8) {5};
  \node[setnode,draw=hmtblue,fill=hmtlightblue] (h6) at (-2.3,-.15) {6};
  \node[setnode,draw=hmtgreen,fill=hmtlightgreen] (o7) at (-3.2,.55) {7};
  \node[setnode,draw=hmtgreen,fill=hmtlightgreen] (o8) at (1.0,.55) {8};
  \node[draw=hmtblue,dashed,rounded corners=4mm,
    fit=(h1)(h2)(h3)(h4)(h5)(h6),inner sep=4mm,
    label={[hmtblue,font=\bfseries]above:{hexada \(H\), \(|H|=6\)}}] (H) {};
  \node[draw=hmtgreen,line width=.8pt,rounded corners=6mm,
    fit=(H)(o7)(o8),inner sep=5mm,
    label={[hmtgreen,font=\bfseries]below:{octada \(O\), \(|O|=8\), \(H\subset O\)}}] (O) {};

  \node[box,draw=hmtblue,fill=hmtlightblue] (f6) at (5.3,1.35)
    {Incidencias definidas sobre \(H\)\\[1mm]
     \(\mathcal F_6=H\times\binom H2\),\quad
     \(|\mathcal F_6|=6\binom62=90\)};
  \node[box,draw=hmtgreen,fill=hmtlightgreen] (f8) at (5.3,-1.0)
    {Extensión de incidencias a \(O\)\\[1mm]
     \(\mathcal F_8=O\times\binom H2\),\quad
     \(|\mathcal F_8|=8\binom62=120\)};
  \draw[arr,hmtblue] (O.east) -- ++(.45,0) |- (f6.west);
  \draw[arr,hmtgreen] (O.east) -- ++(.45,0) |- (f8.west);
  \draw[arr,hmtgold] (f6.south) -- node[right,align=left,font=\footnotesize]
    {inclusión \(\iota_{6,8}\)} (f8.north);

  \node[box,draw=hmtgold,fill=hmtlightgold,text width=10.9cm]
    (def) at (1.45,-3.80)
    {Defecto combinatorio orientado\\[-.2mm]
     con la normalización declarada:\\[1mm]
     \(\displaystyle
       \frac{|\mathcal F_6|-|\mathcal F_8|}{24\binom62}
       =\frac{90-120}{24\cdot15}=-\frac1{12}.\)};
  \draw[arr,hmtgold] (f8.south) -- (f8.south |- def.north);

  \node[align=center,text=hmtgray,font=\footnotesize,text width=11.2cm]
    at (1.45,-5.50)
    {Esta identidad pertenece al mapa de incidencias. No identifica el
     extractor dodecafásico, el semigrupo de torres ni el funcional angular.};
\end{tikzpicture}
 \caption{Inclusión de un subconjunto distinguido de seis elementos en uno de
ocho elementos y conteos de incidencia. Los cardinales \(90\) y \(120\) no fusionan el extractor
dodecafásico, el funcional angular ni el semigrupo de torres.}
\label{mdg:fig:incidencia-seis-ocho-rev9}
\end{figure}

Los índices iniciales de la jerarquía angular quedan así asociados a dos
espacios de incidencia: 90 configuraciones sobre la hexada y 120 sobre la
octada ambiente. El cociente \(-1/12\) es una identidad combinatoria de esta
normalización.

\section{Series de Lambert asociadas a los 2 sectores de incidencia}

Para \(0<q<1\), definimos simultáneamente
\[
 S_{90}(q)=\frac{q^{90}}{1-q^{270}},
 \qquad
 S_{120}(q)=\frac{q^{120}}{1-q^{360}},
 \qquad
 \Delta S_m=S_m(q_-)-S_m(q_+).
\]
Las expansiones geométricas
\[
 \Delta S_{90}=\sum_{j\geq0}s_{90+270j},
 \qquad
 \Delta S_{120}=\sum_{j\geq0}s_{120+360j}
\]
sitúan ambos términos dentro del semigrupo global del
\cref{mdg:chap:torres}.

El funcional angular asociado a la transición \(6\to8\) es
\[
 K_{6\to8}=12\Delta S_{90}-\Delta S_{120},
\]
y no se deduce de las cardinalidades de incidencia por sí solas. Sus pesos
quedan determinados únicamente después de fijar las dos series, el residuo
de polo \(1/24\), la positividad entera y el criterio de minimalidad que se
demuestra a continuación.
La coordenada angular resultante es
\[
\Gamma_{6\to8}
=\frac{\pi A\Cstar}{180}+K_{6\to8}.
\]
En el punto HMT,
\[
\begin{aligned}
 \Delta S_{90}&=2.95169439297457621139847987861\times10^{-4},\\
 \Delta S_{120}&=1.96835112502951720093738706217\times10^{-5},\\
 \frac{\pi A\Cstar}{180}&=0.270485322934140078465586458790\ldots,
\end{aligned}
\]
y, por tanto,
\[
 \boxed{\Gamma_{6\to8}
 =0.27400767269445927474725526077398041940\ldots.}
\]
Las tablas calculadas con la coordenada angular conjugada \(C^*\) redondeada a quince decimales dan
\(0.2740076726944593\) en aritmética de doble precisión; el valor anterior
es la evaluación multiprecisión de la fase regional completa.

\section{Ecuación diofántica de los pesos}

La identidad
\[
 \lim_{q\to1^-}(1-q)\frac{q^m}{1-q^n}=\frac1n
\]
asocia a la combinación
\(a\Delta S_{90}-b\Delta S_{120}\) el residuo normalizado
\[
 \frac a{270}-\frac b{360}.
\]

\begin{theorem}[Pesos enteros con residuo \texorpdfstring{\(1/24\)}{1/24}]
Las soluciones enteras positivas de
\[
 \frac a{270}-\frac b{360}=\frac1{24}
\]
son
\[
 (a,b)=(12+3t,1+4t),\qquad t\in\ZZ_{\geq0}.
\]
El par \((12,1)\) es la única solución mínima coordenada a coordenada y la
única de norma \(\ell^1\) mínima.
\end{theorem}
\begin{proof}
Multiplicar por \(1080\) produce \(4a-3b=45\). Módulo tres se obtiene
\(a\equiv0\pmod3\); escribiendo \(a=3u\), resulta \(4u-b=15\). La
positividad implica \(u=4+t\), \(b=1+4t\), con \(t\geq0\). Por tanto,
\(a+b=13+7t\), lo que demuestra la minimalidad.
\end{proof}

La primitividad \(\gcd(a,b)=1\) no selecciona por sí sola el par mínimo, pues
existen otras soluciones primitivas. El peso \(12:-1\) resulta de combinar la
condición de residuo \(1/24\) con la minimalidad. La elección de ese residuo es
una condición de normalización de la construcción; no se deduce de la
evaluación decimal de \(\Gamma_{6\to8}\).

\newpage
\section{Monotonía del funcional respecto de la coordenada angular conjugada}

Para \(m\in\{90,120\}\),
\[
 S_m'(q)=\frac{m q^{m-1}(1+2q^{3m})}{(1-q^{3m})^2}>0.
\]
Además,
\[
 \frac{\dd q_-}{\dd C}=\frac\pi{180}q_-,
 \qquad
 \frac{\dd q_+}{\dd C}=-\frac\pi{180}q_+.
\]

\begin{proposition}[Monotonía estricta]
Para \(1^\circ\leq C\leq3^\circ\), la función
\(C\mapsto\Gamma_{6\to8}(A,C)\) es estrictamente creciente. En consecuencia,
todo valor de su imagen tiene a lo sumo una preimagen en ese intervalo.
\end{proposition}
\begin{proof}
La derivada de cada diferencia \(\Delta S_m\) es positiva. En el intervalo
indicado se obtiene
\(\Delta S'_{120}<5.35\times10^{-4}\), mientras que
\(\pi A/180>0.1273\). Por tanto,
\(\Gamma'(C)>0.1267\). En \(C=\Cstar\),
\[
 \Gamma'(\Cstar)=0.1328995105618907\ldots.
\]
\end{proof}

La inyectividad local excluye una segunda coordenada angular próxima con la
misma evaluación del funcional. Este control no construye ni selecciona
\(C^*\): la coordenada angular conjugada ya ha sido obtenida en la rama regional anterior.
La proposición sólo cuantifica la sensibilidad de la evaluación descendente
una vez fijados \(A\) y \(C^*\); no autoriza a invertir
\(\Gamma_{6\to8}\) ni un valor físico de \(\gamma_{\mathrm{BI}}\) para
redefinirlos.

 \label{ch:modular}

\section{Inclusi\'on hexada--octada y realizaciones funcionales asociadas}

Fijemos un conjunto \(H\) de seis puntos y un conjunto \(O\) de ocho
puntos con una inclusi\'on distinguida \(j:H\hookrightarrow O\). En la
residencia de Witt, \(H\) es una hexada de un dodecad y \(O\) su octada
elevada; en este cap\'itulo s\'olo se usa la inclusi\'on ya fijada, no el
valor terminal de Barbero--Immirzi. Escribimos
\begin{equation}
\label{eq:modular-pairs-H}
\binom H2:=\{P\subset H:|P|=2\},
\qquad \left|\binom H2\right|=\binom62=15,
\end{equation}
y definimos dos espacios de incidencias con tipos distintos:
\begin{equation}
\label{eq:modular-F6-F8}
\mathcal F_6:=H\times\binom H2,
\qquad
\mathcal F_8:=O\times\binom H2.
\end{equation}
La inclusi\'on de incidencias es
\begin{equation}
\label{eq:modular-iota68}
\iota_{6,8}:\mathcal F_6\hookrightarrow\mathcal F_8,
\qquad
(h,P)\longmapsto(j(h),P).
\end{equation}
Es inyectiva porque \(j\) lo es y conserva el par marcado \(P\).

\begin{theorem}[Conteos de incidencia y defecto orientado]
\label{thm:modular-incidence-90-120}
Los dos espacios de \cref{eq:modular-F6-F8} satisfacen
\begin{equation}
\label{eq:modular-counts-90-120}
|\mathcal F_6|=6\binom62=90,
\qquad
|\mathcal F_8|=8\binom62=120,
\qquad
|\mathcal F_8\setminus\iota_{6,8}(\mathcal F_6)|=30.
\end{equation}
La normalizaci\'on por las veinticuatro coordenadas del par
dodecad--veinticuatro y por los quince pares de \(H\) produce
\begin{equation}
\label{eq:modular-incidence-defect}
\boxed{
\delta_{6,8}^{\rm inc}
=\frac{|\mathcal F_6|-|\mathcal F_8|}
{24\binom62}=-\frac1{12}.}
\end{equation}
La misma fracci\'on se obtiene mediante el funcional rec\'iproco asociado
a la memoria incorporada,
\begin{equation}
\label{eq:modular-reciprocal-defect}
\boxed{
\delta_{6,8}^{\rm rec}
=\frac{30}{120}-\frac{30}{90}=-\frac1{12}.}
\end{equation}
Son dos realizaciones exactas del mismo defecto orientado, pero tienen
dominios distintos y no se identifican con \(\zeta(-1)\) sin un
entrelazador adicional.
\end{theorem}

\begin{proof}
El principio multiplicativo aplicado a
\cref{eq:modular-F6-F8} da \(6\cdot15=90\) y \(8\cdot15=120\).
La imagen de \(\iota_{6,8}\) tiene noventa elementos, de modo que el
complemento tiene treinta. Sustituir en
\cref{eq:modular-incidence-defect} da
\(-30/(24\cdot15)=-1/12\), mientras
\cref{eq:modular-reciprocal-defect} es
\(1/4-1/3=-1/12\).
\end{proof}

La transici\'on hexada--octada admite dos realizaciones funcionales que no
se deducen mutuamente. El diagrama algebraico, formulado sobre los espacios
que justifican los conteos, es
\begin{equation}
\label{eq:modular-68-two-readers}
\begin{aligned}
(H\xhookrightarrow{\,j\,}O)
&\longmapsto
\left(|\mathcal F_6|,|\mathcal F_8|\right)=(90,120)
\longmapsto\Gamma_{6\to8}=\gammaH,\\
(h_6,o_8,t)=(6,8,3)
&\longmapsto-\frac{o_8-h_6}{t}=-\frac23,
\end{aligned}
\end{equation}
donde \(-2/3\) es el coeficiente de la tercera coordenada de la fila
\(13\) del constructor CKM. La primera realizaci\'on transforma la
incidencia en \`area y memoria modular; la segunda transforma la misma
transici\'on en par\'ametros de mezcla angular.
No se postula \(\boldsymbol\theta_{\rm CKM}=f(\gammaH)\): ambas salidas
conservan un antecedente com\'un.

El morfismo que determina los sectores \(90/120\) es
\(\iota_{6,8}\). La magnitud \(\gammaH\) se adopta como teorema previamente
demostrado en la construcci\'on hexada--octada y no se reconstruye a partir
de su representaci\'on decimal. Las
consecuencias operatorias comienzan s\'olo despu\'es de fijar la escala
elemental \(a_0\), que se define a continuaci\'on sin usar entrop\'ia ni el
producto terminal \(\gammaH a_0\).

\section{Correspondencia entre el interior non\'adico y los grados de libertad de borde}

Sea \(\mathbb T_3=\{0,1,2\}\). Todo \(x\in\mathbb Z/9\mathbb Z\) posee
una \`unica escritura de d\'igitos
\(x=x_0+3x_1\), con \(x_i\in\mathbb T_3\), y todo
\(y\in\mathbb Z/27\mathbb Z\) una \`unica escritura
\(y=y_0+3y_1+9y_2\). Por ello
\begin{equation}
\label{eq:modular-bulk-boundary-729}
\mathcal B=(\mathbb Z/9\mathbb Z)^3,
\qquad
\mathcal B_\partial=(\mathbb Z/27\mathbb Z)^2,
\qquad
|\mathcal B|=|\mathcal B_\partial|=3^6=729.
\end{equation}
La igualdad de cardinales se acompa\~na del mapa expl\'icito. Si
\(x_j=r_{2j-1}+3r_{2j}\) para \(j=1,2,3\), definimos
\begin{equation}
\label{eq:modular-beta729}
\beta_{729}(x_1,x_2,x_3)
=\bigl(r_1+3r_2+9r_3,\;r_4+3r_5+9r_6\bigr).
\end{equation}

\begin{proposition}[Bijecci\'on geneal\'ogica \(729\leftrightarrow729\)]
\label{prop:modular-beta729}
El mapa \(\beta_{729}:\mathcal B\to\mathcal B_\partial\) es biyectivo y
preserva la palabra ordenada de seis trits. No es, en general, un
isomorfismo de grupos aditivos: reagrupar los d\'igitos cambia d\'onde se
registra el acarreo.
\end{proposition}

\begin{proof}
Extraer los seis d\'igitos de tres coordenadas non\'adicas y agruparlos en
dos ternas es una permutaci\'on de coordenadas de \(\mathbb T_3^6\). La
descomposici\'on posicional es \`unica en ambos lados, por lo que la
operaci\'on tiene inversa. La advertencia aditiva se comprueba, por ejemplo,
al sumar palabras que generan acarreo entre \(r_3\) y \(r_4\): esa
frontera separa las dos coordenadas de orden veintisiete, pero atraviesa
dos coordenadas distintas en el empaquetado non\'adico.
\end{proof}

En el toro de borde \(\mathcal B_\partial\), tomamos el bloque de v\'ertices
\begin{equation}
\label{eq:modular-block-A}
A=\{0,1,\ldots,8\}\times\{0,1,\ldots,8\}.
\end{equation}
Sus enlaces orientados que cruzan el borde se separan por direcci\'on:
\begin{align}
\partial A_{90}
&=\{((-1,j),(0,j)),((8,j),(9,j)):0\leq j\leq8\},
\label{eq:modular-boundary90}\\
\partial A_{120}
&=\{((i,-1),(i,0)),((i,8),(i,9)):0\leq i\leq8\},
\label{eq:modular-boundary120}
\end{align}
donde todas las coordenadas se leen m\'odulo veintisiete. As\'i
\begin{equation}
\label{eq:modular-boundary-counts}
\partial A=\partial A_{90}\sqcup\partial A_{120},
\qquad
|\partial A_{90}|=|\partial A_{120}|=18.
\end{equation}
Los sub\'indices \(90,120\) etiquetan la procedencia de canal; no dicen que
el borde tenga noventa o ciento veinte enlaces.

\begin{theorem}[Ley areal del corte tri\'adico]
\label{thm:modular-triadic-cut}
En el grafo c\'ubico \(N\times N\times N\) con capacidad unitaria, todo
corte entre dos caras opuestas posee capacidad al menos \(N^2\), y una
capa transversal alcanza esa cota. Para \(N=9\),
\begin{equation}
\label{eq:modular-mincut-entropy}
\mathcal A_{\min}=81,
\qquad
S_{\rm tri}=81\log3.
\end{equation}
\end{theorem}

\begin{proof}
Hay \(N^2\) rectas de aristas paralelas a la direcci\'on normal, disjuntas
dos a dos. Todo corte debe interceptar cada una, de modo que su capacidad
es al menos \(N^2\). Una sola capa transversal contiene exactamente
\(N^2\) aristas y realiza el m\'inimo. Si cada enlace cortado porta un trit
m\'aximamente mezclado, su entrop\'ia es \(\log3\); la aditividad sobre los
ochenta y un enlaces da \cref{eq:modular-mincut-entropy}.
\end{proof}

La entrop\'ia \(S_{\rm tri}\) mide capacidad del corte c\'ubico. La
entrop\'ia modular de los treinta y seis enlaces de
\cref{eq:modular-boundary-counts} mide otro estado y no debe igualarse con
ella.

\section{Entrelazador colectivo de incidencia entre interior y borde}

En la suma
\(\ell^2(\mathcal F_6)\oplus\ell^2(\mathcal F_8)\), definimos
\begin{equation}
\label{eq:modular-u90-u120}
u_{90}=\frac1{\sqrt{90}}\sum_{f\in\mathcal F_6}(\delta_f,0),
\qquad
u_{120}=\frac1{\sqrt{120}}\sum_{f\in\mathcal F_8}(0,\delta_f).
\end{equation}
En
\(\ell^2(\partial A_{90})\oplus\ell^2(\partial A_{120})\), sean
\begin{equation}
\label{eq:modular-v90-v120}
v_{90}=\frac1{\sqrt{18}}\sum_{e\in\partial A_{90}}(\delta_e,0),
\qquad
v_{120}=\frac1{\sqrt{18}}\sum_{e\in\partial A_{120}}(0,\delta_e).
\end{equation}
Las parejas \((u_{90},u_{120})\) y \((v_{90},v_{120})\) son ortonormales.
Por tanto existe una \`unica isometr\'ia entre los subespacios colectivos
que satisface
\begin{equation}
\label{eq:modular-J68}
\mathcal J_{6\to8}u_{90}=v_{90},
\qquad
\mathcal J_{6\to8}u_{120}=v_{120}.
\end{equation}
Definamos
\begin{align}
D_{\rm inc}
&=90|u_{90}\rangle\langle u_{90}|
+120|u_{120}\rangle\langle u_{120}|,
\label{eq:modular-Dinc}\\
D_{\rm area}
&=90|v_{90}\rangle\langle v_{90}|
+120|v_{120}\rangle\langle v_{120}|.
\label{eq:modular-Darea}
\end{align}

\begin{theorem}[Transporte exacto de las dos etiquetas]
\label{thm:modular-J68}
La aplicaci\'on \(\mathcal J_{6\to8}\) es unitaria entre los dos planos
colectivos y entrelaza los operadores de incidencia y \`area:
\begin{equation}
\label{eq:modular-J-intertwining}
\boxed{
\mathcal J_{6\to8}D_{\rm inc}
=D_{\rm area}\mathcal J_{6\to8}.}
\end{equation}
Por tanto, toda combinaci\'on polin\'omica de las etiquetas \(90,120\), en
particular la combinaci\'on primitiva \(12:-1\), se transporta con los
mismos coeficientes.
\end{theorem}

\begin{proof}
La ortonormalidad muestra que \(\mathcal J_{6\to8}\) env\'ia una base
ortonormal en otra. En \(u_{90}\) ambos lados de
\cref{eq:modular-J-intertwining} valen \(90v_{90}\); en \(u_{120}\),
ambos valen \(120v_{120}\). La igualdad en una base prueba la identidad.
El c\'alculo funcional polin\'omico conserva el entrelazamiento.
\end{proof}

Este mapa no pretende inyectar noventa incidencias individuales en
dieciocho enlaces: transporta exactamente los dos modos uniformes y sus
etiquetas espectrales. Extenderlo a fluctuaciones no uniformes exigir\'ia
un nuevo mapa de incidencia; el operador modular definido aqu\'i s\'olo emplea
el subespacio colectivo que acaba de construirse.

\section{Determinación de la escala elemental del espectro del operador de área}

Sea \(C^*\) la coordenada angular conjugada, expresada en grados, y fijemos
\begin{equation}
\label{eq:modular-theta-clock-a0}
\theta_{\rm clk}=\frac{C^*}{6}\frac{\pi}{180},
\qquad
\boxed{
a_0=\frac{1+2\cos(10^\circ)}{3}\,\theta_{\rm clk}.}
\end{equation}
El factor que precede a \(\theta_{\rm clk}\) no es una correcci\'on
decimal. Es el car\'acter real normalizado del triplete de fases
\begin{equation}
\label{eq:modular-c10-character}
\frac13\Tr\operatorname{diag}
\left(1,\ee^{\ii\pi/18},\ee^{-\ii\pi/18}\right)
=\frac{1+2\cos(10^\circ)}3.
\end{equation}
As\'i, \(a_0\) es la media espectral de la estructura angular sobre el trit
orientado. Es positiva, se fija antes del operador de \`area y no se
selecciona a partir de la entrop\'ia. La identificaci\'on posterior
\(\gammaH=\Gamma_{6\to8}\) realiza la salida HMT en la conexi\'on de
Ashtekar--Barbero \cite{Barbero,Immirzi}; no reabre su derivaci\'on.

\begin{remark}[Alcance del s\'imbolo \(a_0\)]
En este cap\'itulo, \(a_0\equiv a_0^{\rm area}\) denota exclusivamente la
normalizaci\'on de \cref{eq:modular-theta-clock-a0}. No es el radio de Bohr
ni el factor de escala cosmol\'ogico inicial, aunque esos objetos se escriban
tradicionalmente con el mismo s\'imbolo. Ninguna identidad de este cap\'itulo
permite sustituir uno por otro.
\end{remark}
 \section{Ley areal del corte triádico: capacidad \(81\log 3\), operador de área, espectro y entrelazador colectivo de los canales \(90/120\)}
\label{hap:sec:area-entrelazamiento-barbero-rev11}
\index{Barbero--Immirzi!operador de área}
\index{entropía modular!Hamiltoniano modular}

La identificación HMT del funcional \(\Gamma_{6\to8}\) con la coordenada
de Barbero--Immirzi se realiza sobre el mismo borde triádico que soporta la
incidencia (90/120). El interior y su hoja de borde poseen igual cardinalidad,
\[
 \mathcal B=(\mathbb Z/9\mathbb Z)^3,
 \quad |\mathcal B|=729,
 \qquad
 \partial\mathcal B=(\mathbb Z/27\mathbb Z)^2,
 \quad |\partial\mathcal B|=729,
\]
mediante el reagrupamiento de seis trits en tres coordenadas nonádicas o dos
coordenadas de orden veintisiete.

\begin{theorem}[Ley areal del corte triádico]
\label{hap:thm:ley-areal-corte-triadico-rev11}
En el grafo cúbico \(N\times N\times N\) con capacidad unitaria, el corte
mínimo entre dos caras opuestas tiene capacidad \(N^2\). Para \(N=9\),
\[
 \mathcal A_{\min}=81,
 \qquad S_{\rm tri}=\mathcal A_{\min}\log3=81\log3.
\]
\end{theorem}
\begin{proof}
Las \(N^2\) rectas paralelas a la dirección normal son caminos disjuntos en
aristas, por lo que todo corte las intercepta. Una capa transversal contiene
exactamente \(N^2\) aristas y alcanza la cota. Cada enlace triádico aporta
\(\log3\) a la entropía de un estado máximamente mezclado.
\end{proof}

En el toro de borde \(27\times27\), sea \(A\) el bloque canónico
\(9\times9\). Su frontera se descompone de manera exacta en
\begin{equation}
 \partial A=\partial A_{90}\sqcup\partial A_{120},
 \qquad |\partial A_{90}|=|\partial A_{120}|=18.
 \label{hap:eq:corte-90-120-rev11}
\end{equation}
Los enlaces horizontales representan el canal (90), y los verticales, el
canal (120). Esta descomposición convierte la pareja de incidencias en una
descomposición operatoria del borde.

La entropía areal \(S_{\rm tri}=81\log3\) pertenece al estado máximamente
mezclado de los enlaces del corte. La entropía bicapa
\(S_{\rm bi}(y)\) de la \cref{hap:def:entropia-bicapa-rev6} pertenece a la
distribución normalizada entre las dos hojas; ambos objetos conservan
dominios distintos.

Sea \(N=\operatorname{diag}(0,1,2)\) el operador número de un enlace trítico.
La realización areal declara la escala positiva
\[
 \theta_{\rm clk}=\frac{C^*}{6}\frac{\pi}{180},
 \qquad
 a_0=\frac{1+2\cos(10^\circ)}{3}\,\theta_{\rm clk},
 \qquad \gamma_{\rm HMT}=\Gamma_{6\to8}.
\]
La fórmula de \(a_0\) constituye la estructura de partida de esta
realización; fijada esa escala, el operador y su espectro se deducen
exactamente.

\begin{definition}[Operador de área HMT]
Sobre la hoja de borde
\(\mathcal H_{\partial}=\bigotimes_{e\in\partial A}\mathbb C^3\), se define
\begin{equation}
 \widehat{\mathcal A}_{\rm HMT}
 =\gamma_{\rm HMT}a_0\sum_{e\in\partial A}N_e.
 \label{hap:eq:operador-area-hmt-rev11}
\end{equation}
Su espectro en el corte canónico es
\[
 \operatorname{Spec}\widehat{\mathcal A}_{\rm HMT}
 =\{n\gamma_{\rm HMT}a_0:0\leq n\leq72\},
\]
con multiplicidades dadas por los coeficientes de
\((1+z+z^2)^{36}\). La prolongación por cortes compatibles produce la
familia protegida \(\mathcal A_n^{\rm prot}=n\gamma_{\rm HMT}a_0\).
\end{definition}

El entrelazador que fija la normalización se construye en el subespacio de
canales colectivos. Sean \(u_{90},u_{120}\) los vectores unitarios uniformes
de las incidencias hexádica y octádica, y sean \(v_{90},v_{120}\) los
vectores unitarios uniformes de los dos conjuntos de enlaces de
\eqref{hap:eq:corte-90-120-rev11}. El mapa
\begin{equation}
 \begin{aligned}
 \mathcal J_{6\to8}:{}
 &\operatorname{span}\{u_{90},u_{120}\}
 \longrightarrow
 \operatorname{span}\{v_{90},v_{120}\},\\
 &\mathcal J_{6\to8}u_{90}=v_{90},
 \qquad
 \mathcal J_{6\to8}u_{120}=v_{120}
 \end{aligned}
 \label{hap:eq:entrelazador-area-barbero-rev11}
\end{equation}
es una isometría entre estos subespacios colectivos. Si
\[
 D_{\rm inc}=90|u_{90}\rangle\langle u_{90}|
 +120|u_{120}\rangle\langle u_{120}|,
 \qquad
 D_{\rm area}=90|v_{90}\rangle\langle v_{90}|
 +120|v_{120}\rangle\langle v_{120}|,
\]
entonces
\[
 \boxed{\mathcal J_{6\to8}D_{\rm inc}
 =D_{\rm area}\mathcal J_{6\to8}.}
\]
En particular, la combinación primitiva \(12:-1\) actúa
con los mismos coeficientes antes y después de \(\mathcal J_{6\to8}\). Esta
igualdad fija el diccionario de normalización que inserta
\(\gamma_{\rm HMT}=\Gamma_{6\to8}\) en la conexión de
Ashtekar--Barbero.

 El paso de hexada a octada es una transformación de la estructura de incidencia.\par

Al pasar de \(6\) a \(8\), cambian los pares disponibles, las relaciones internas y la forma en que la orientación puede conservarse. Esa transición genera un defecto orientado. La variación cardinal realiza un giro en el modo de organización.\par

La misma inversión de hoja se expresa mediante\par

\[
C^*\longmapsto -C^*.
\]

Pero distintos lectores responden de formas distintas a esa inversión.\par

El funcional hexada–octada que conduce al parámetro de Barbero–Immirzi es impar: al invertir la hoja, cambia de signo. La orientación geométrica se invierte.\par

La entropía bicapa, en cambio, es par: atribuye la misma entropía a \(C^*\) y \(-C^*\) cuando se observa la distribución de probabilidades. La información total puede permanecer igual aunque la orientación se haya invertido.\par

El espacio CKM realiza una tercera respuesta. La inversión induce sobre los tres ángulos una reflexión racional exacta:\par

\[
R_{\rm CKM}
=
M_{\rm CKM}
\operatorname{diag}(1,-1,1)
M_{\rm CKM}^{-1},
\]

con\par

\[
R_{\rm CKM}^2=I,
\qquad
\det R_{\rm CKM}=-1.
\]

Esto significa que el cambio de hoja se transporta al espacio de mezcla como una verdadera reflexión. Existe una dirección impar de rango uno y un plano que permanece fijo.\par

El contraángulo \(C^*\) es precisamente la coordenada que cambia. La fase construida desde \(A\) permanece invariante. La transformación separa así dos memorias:\par

\begin{itemize}[leftmargin=*,itemsep=0.35em]
\item una memoria de orientación;
\item una memoria de fase.
\end{itemize}

Esto ofrece una lectura muy profunda de CKM. Los ángulos de mezcla son coordenadas de un espacio sobre el que actúa la misma involución que ya organizaba la transición hexada–octada.\par

El giro \(6\to8\) suministra la operación de orientación que CKM realiza en su propio codominio; la causalidad física pertenece al transductor correspondiente.\par

Tenemos entonces tres publicaciones hermanas del mismo tránsito:\par

\begin{enumerate}[leftmargin=*,itemsep=0.65em]
\item \textbf{Publicación geométrica.}
El parámetro de Barbero–Immirzi orienta el cuanto de área.\par

\item \textbf{Publicación informacional.}
El estado reducido y su Hamiltoniano modular miden qué información de esa orientación permanece en el borde.\par

\item \textbf{Publicación de mezcla.}
La reflexión CKM transforma las coordenadas de mezcla y aísla su componente impar.\par
\end{enumerate}

\(\gamma_{\rm BI}\), la entropía y los ángulos CKM permanecen como realizaciones distintas de una bisagra genealógica común. La unidad reside en el antecedente anterior a la ramificación. Comparten la bisagra, pero cada uno conserva una parte distinta de ella.\par

Ésta es una de las expresiones más bellas de la metodología HMT: la unidad fundamental se recupera remontando las genealogías hasta encontrar la operación común que genera resultados terminales diferentes.\par

La gravedad conserva el giro como orientación del área.\par
La información conserva el giro como memoria del estado.\par
La mezcla conserva el giro como reflexión de coordenadas.\par

La misma transformación «respira» de tres maneras.\par

 \section{Separación causal entre el espectro de área, la memoria modular y la realización Cartan--Holst}
La derivación de $\gamma_{\mathrm{BI}}$, del área mínima y del espectro del operador de área concluye en este capítulo. La reconstrucción conjunta de ocupación y procedencia mediante el Hamiltoniano modular de borde, y la realización Palatini--Cartan--Holst, conservan dominios, operadores y pruebas propios en la Parte IV.
 \endgroup
% Secciones 4--7 del delta gravitatorio íntegro.
\section{Incidencia hexada--octada, parámetro de Barbero--Immirzi y espectro de área}
\label{sec:l9s102:barbero-area}

La derivación areal recibe como antecedentes la inclusión de incidencia
\(6\to8\), los cardinales TPK \(90/120\), el defecto normalizado
\(-1/12\) y los canales espectrales \(q_\pm\). El funcional
hexada--octada determina el parámetro de Barbero--Immirzi y la
normalización del operador de área. Esta residencia contiene la derivación
HMT y el espectro; la realización Palatini--Cartan--Holst ocupa una residencia
posterior en Mecánica Dimensional.

Sean \(H\subset O\), con \(|H|=6\) y \(|O|=8\), y consérvese el par
marcado de la hexada en la inclusión
\begin{equation}
 \iota_{6,8}:\mathcal F_6=H\times\binom H2
 \lhook\joinrel\longrightarrow
 \mathcal F_8=O\times\binom H2.
 \label{eq:l9s102:barbero-inclusion-incidencia}
\end{equation}
La cuenta interna da
\begin{equation}
 |\mathcal F_6|=6\binom62=90,
 \qquad |\mathcal F_8|=8\binom62=120,
 \qquad
 \frac{90-120}{24\binom62}=-\frac1{12}.
 \label{eq:l9s102:barbero-conteos}
\end{equation}
Una vez publicados \(\pi\), el ángulo \(A\) y el
contraángulo único \(C^*\), los dos canales son
\begin{equation}
 q_\pm=\exp\!\left[-\frac{\pi}{180}(A\pm C^*)\right],
 \qquad 0<q_+<q_-<1,
 \label{eq:l9s102:barbero-qpm}
\end{equation}
y sus series de retorno orientado son
\begin{equation}
 S_{90}(q)=\frac{q^{90}}{1-q^{270}},\qquad
 S_{120}(q)=\frac{q^{120}}{1-q^{360}},\qquad
 \Delta S_m=S_m(q_-)-S_m(q_+).
 \label{eq:l9s102:barbero-series}
\end{equation}

\begin{theorem}[Determinación diofántica del funcional de Barbero--Immirzi]
\label{thm:l9s102:barbero-funcional}
Entre las combinaciones positivas
\(a\Delta S_{90}-b\Delta S_{120}\), la condición de residuo \(1/24\)
y minimalidad coordenada a coordenada determina el único par primitivo
\((a,b)=(12,1)\). Por tanto la salida adimensional es
\begin{equation}
 \boxed{
 \gamma_{\rm HMT}=\Gamma_{6\to8}
 =\frac{{\pi}AC^*}{180}
   +12\Delta S_{90}-\Delta S_{120}.}
 \label{eq:l9s102:barbero-gamma}
\end{equation}
Esta fórmula recibe el par angular ya generado y no utiliza como entrada
el valor convencional del parámetro de Barbero--Immirzi.
\end{theorem}

\begin{proof}
El residuo de \(S_{90}\) en \(q=1\) es \(1/270\), y el de
\(S_{120}\), \(1/360\). La condición requerida equivale a
\[
 \frac a{270}-\frac b{360}=\frac1{24}
 \quad\Longleftrightarrow\quad 4a-3b=45.
\]
Sus soluciones enteras positivas son
\((a,b)=(12+3t,1+4t)\), \(t\ge0\); la única solución mínima es
\((12,1)\). Sustituir este par en el funcional de incidencia prueba
\eqref{eq:l9s102:barbero-gamma}.
\end{proof}

La secuencia causal queda fijada sin intercambio de papeles:
\begin{equation}
 (\mathcal F_6\hookrightarrow\mathcal F_8)
 \longrightarrow (90,120,-1/12)
 \longrightarrow (A,C^*,q_\pm)
 \longrightarrow \Gamma_{6\to8}=\gamma_{\rm HMT}
 \longrightarrow \widehat{\mathcal A}_{\rm phys}.
 \label{eq:l9s102:barbero-cadena-completa}
\end{equation}
La realización posterior en la acción de Holst reconoce la misma sección
adimensional; no interviene en su selección.

Para los dos sectores de borde, sean
\begin{equation}
 N_m=\sum_{e\in\partial A_m}N_e,
 \qquad m\in\{90,120\},
 \qquad \operatorname{Spec}(N_e)=\{0,1,2\}.
 \label{eq:l9s102:n90-n120}
\end{equation}
El operador físico de área se escribe
\begin{equation}
 \boxed{
 \frac{\widehat{\mathcal A}_{\mathrm{phys}}}{\ell_P^2}
 =\gamma_{\mathrm{HMT}}a_0(N_{90}+N_{120}).}
 \label{eq:l9s102:area-operador}
\end{equation}
Sobre \(36\) enlaces qutrit,
\begin{equation}
 \operatorname{Spec}(N_{90}+N_{120})=\{0,1,\ldots,72\},
 \label{eq:l9s102:area-numero-spectrum}
\end{equation}
y la multiplicidad del autovalor \(n\) es el coeficiente de \(z^n\) en
\((1+z+z^2)^{36}\). De aqui resulta
\begin{equation}
 \operatorname{Spec}\!\left(
 \frac{\widehat{\mathcal A}_{\mathrm{phys}}}{\ell_P^2}\right)
 =\{n\gamma_{\mathrm{HMT}}a_0:0\le n\le72\}.
 \label{eq:l9s102:area-spectrum}
\end{equation}

\section{Geometría total y memoria de procedencia}
\label{sec:l9s102:area-hamiltoniano}

El Hamiltoniano modular del borde es
\begin{equation}
 \boxed{
 K_A=-\log\rho_A
 =cI+\Delta_{\mathrm{mod}}(90N_{90}+120N_{120}).}
 \label{eq:l9s102:hamiltoniano-modular}
\end{equation}
Sean
\begin{equation}
 N=N_{90}+N_{120},
 \qquad
 D=90N_{90}+120N_{120}.
 \label{eq:l9s102:n-d}
\end{equation}
La matriz de paso
\(\left(\begin{smallmatrix}1&1\\90&120\end{smallmatrix}\right)\)
posee determinante \(30\), y su inversa reconstruye los sectores:
\begin{equation}
 \boxed{
 N_{90}=4N-\frac{D}{30},
 \qquad
 N_{120}=\frac{D}{30}-3N.}
 \label{eq:l9s102:reconstruccion-sectores}
\end{equation}
El área publica el número total de excitaciones; el Hamiltoniano modular
publica su sector de procedencia. El par conjunto reconstruye la ocupación
completa de la frontera.

\section{Estado bicapa e información orientada}
\label{sec:l9s102:estado-bicapa}

Sea \(R=R^*=R^{-1}\) la involución de hojas y
\(y=C^*_{\mathrm{rad}}\). El estado normalizado es
\begin{equation}
 \boxed{
 \rho_y=\frac{e^{-yR}}{2\cosh y}.}
 \label{eq:l9s102:rho-y}
\end{equation}
Su Hamiltoniano modular y su entropía son
\begin{equation}
 K_y=-\log\rho_y=\log(2\cosh y)I+yR,
 \label{eq:l9s102:ky}
\end{equation}
\begin{equation}
 S(\rho_y)=\log(2\cosh y)-y\tanh y.
 \label{eq:l9s102:entropia-y}
\end{equation}
La inversión \(y\mapsto-y\) deja fija la entropía y produce la distancia
informacional orientada
\begin{equation}
 \boxed{
 D(\rho_y\Vert\rho_{-y})=2y\tanh y.}
 \label{eq:l9s102:entropia-relativa-hojas}
\end{equation}
La componente par mide la distribucion espectral; la componente impar mide
el coste informacional de invertir la orientación.

\section{Cuanto informacional de energía y celda gravitatoria de media acción}
\label{sec:l9s102:horizonte-media-accion}

La periodicidad de \(108\) estados determina
\begin{equation}
 t_P=\frac{54}{\pi}t_0,
 \qquad
 E_P=\frac{\hbar}{t_P}=\frac{h}{108t_0}.
 \label{eq:l9s102:ep-tp}
\end{equation}
La realización térmica del reloj tritico satisface
\begin{equation}
 \boxed{
 E_P=k_B\Theta_{\mathrm{clk}}\log3.}
 \label{eq:l9s102:planck-trit}
\end{equation}
Para el horizonte de radio \(R\),
\begin{equation}
 E_\Lambda(R)=\frac{c^4R}{2G},
 \qquad
 E_\Lambda(R)=T_H(R)S_H(R).
 \label{eq:l9s102:energia-horizonte}
\end{equation}
En la celda de Planck,
\begin{equation}
 \boxed{
 E_\Lambda
 =T_HS_H
 =\frac{E_P}{2}
 =\frac{k_B\Theta_{\mathrm{clk}}\log3}{2}.}
 \label{eq:l9s102:confluencia-horizonte}
\end{equation}
Por tanto,
\begin{equation}
 \boxed{
 E_\Lambda t_P=\frac{\hbar}{2},
 \qquad
 E_\Lambda T_{108}=\frac h2.}
 \label{eq:l9s102:media-accion}
\end{equation}
La energía de vacío, la entropía areal, la temperatura cronológica y la
acción comparten una misma celda de confluencia con dominios tipados.

Cuando \(S_H/k_B=\pi\), la multiplicidad efectiva satisface
\begin{equation}
 \Omega_{\mathrm{eff}}=e^{S_H/k_B}=e^\pi.
 \label{eq:l9s102:omega-e-pi}
\end{equation}
El operador hiperbólico de Euler posee
\begin{equation}
 \operatorname{Spec}(e^{\pi H})=\{e^\pi,e^{-\pi}\}.
 \label{eq:l9s102:e-pi-spectrum}
\end{equation}
La igualdad de los escalares establece la confluencia espectral; el
entrelazador entre multiplicidad microscópica de horizonte y dinamica de
Perron pertenece a su residencia física especifica.
 
\chapter{Reversibilidad exacta, coste de borrado y límite nulo de Landauer}
\begingroup
\renewcommand{\chapter}[2][]{}%
\chapter{Actualización reversible, coste de borrado y dominio exacto del límite nulo de Landauer}
\label{chap:parte-iv-57}

El transporte completo conserva la información de fibra y puede tener coste lógico nulo; la proyección o reinicialización que descarta la genealogía queda sujeta al coste de Landauer. Ambos dominios se separan de forma explícita.

\section{Reversibilidad lógica y principio de Landauer en las fibras de TPK}
\label{sec:informacion-fibra-landauer-rev7}
\index{Landauer}\index{información!fibra}

Sea \((\Omega,2^\Omega,\mu)\) un espacio probabilizado finito, sea
\(X:\Omega\to\mathcal X\) una variable aleatoria con valores en el espacio
\(\mathcal X\) de estados completos y sea \(q:\mathcal X\to Y\) un lector.
La regla de la cadena da
\[
 \mathsf H(X)
 =
 \mathsf H(q(X))+\mathsf H(X\mid q(X)).
\]
El segundo sumando mide la información de la fibra que el lector no
publica.

\begin{theorem}[Landauer nulo para el transporte completo y coste de la proyección]
\label{thm:landauer-relativo-lector-rev7}
Si \(U:\mathcal X\to\mathcal X\) es una actualización biyectiva del estado
HMT completo,
entonces
\[
 \boxed{\mathsf H(X\mid U(X))=0.}
\]
Para una salida proyectada \(q(U(X))\), la información descartada es
\[
 \mathsf H(X\mid q(U(X))).
\]
Un reset físico que identifique estados distintos de la memoria debe
contabilizar esa información en el entorno.
\end{theorem}

\begin{proof}
La biyectividad implica
\(X=U^{-1}(U(X))\); por tanto \(X\) queda determinado por \(U(X)\) y la
entropía condicional es cero. Para el lector \(q\), la regla de la cadena
separa la entropía publicada de la entropía de la fibra. Si la operación de
reinicio no es inyectiva, la salida ya no permite reconstruir el estado de
memoria. El
principio termodinámico de Landauer acota el coste de la implementación en
función de esa entropía borrada; la actualización reversible no aporta por
sí sola un término de borrado.
\end{proof}

En nats y bits, respectivamente,
\[
 Q\ge k_B\Theta\,\mathsf H_{\rm nat},
 \qquad
 Q\ge k_B\Theta\log 2\,\mathsf H_{\rm bit}.
\]
Para el transporte biyectivo del estado completo, la entropía de borrado es
cero y, por tanto, el término mínimo de Landauer asociado exclusivamente al
borrado satisface
\[
\boxed{Q_{\rm L}^{\min}[U]=k_B\Theta\,
 \mathsf H(X\mid U(X))=0.}
\]
Para el reinicio de un TRIT uniforme, en cambio,
\[
 \boxed{Q_{\rm reset}\ge k_B\Theta\log3.}
\]
Estas dos ecuaciones localizan el coste en la pérdida de información: el
coste aparece cuando la proyección o el reinicio olvidan la fibra.
La preparación, la proyección y el reinicio pueden producir entropía; el
transporte reversible conserva la información genealógica. El valor nulo es
el límite lógico de Landauer para esa actualización completa, mientras los
costes dinámicos de control, velocidad finita y acoplamiento al entorno
pertenecen a la realización física del dispositivo.

\begin{statusbox}[title={Alcance del Landauer nulo}]
La igualdad \(\mathsf H(X\mid U(X))=0\) es exacta para una actualización
biyectiva del estado holonómico integral. La cota
\(Q_{\rm reset}\ge k_B\Theta\log3\) es la especialización de Landauer al
borrado de una decisión triádica uniforme. La realización de ambas fórmulas
en joules utiliza el transductor de Boltzmann y una temperatura física
declarada.
\end{statusbox}
 
\paragraph{Transducción de la información de borde en capacidad del corte triádico mínimo.} La información conservada por el borde reaparece geométricamente en el corte triádico mínimo y en su capacidad.
 \endgroup

\chapter{Área mínima del corte triádico y capacidad \texorpdfstring{\(81\log 3\)}{81 log 3}}
\begingroup
\renewcommand{\chapter}[2][]{}%
\chapter{Área mínima del corte triádico, capacidad \texorpdfstring{\(81\log 3\)}{81 log 3} y transporte colectivo de incidencia}
\label{chap:parte-iv-58}

La cuantización areal parte del interior nonádico, su borde y el transporte colectivo de incidencia. El corte mínimo contiene \(81\) unidades y su capacidad qutrit es \(81\log 3\); el resultado precede a su realización gravitatoria.

\section{Cuantizaci\'on areal, din\'amica modular y transformaci\'on racional CKM}
\label{rm:ch:modular}

\subsection{Inclusi\'on hexada--octada y realizaciones funcionales asociadas}

Fijemos un conjunto \(H\) de seis puntos y un conjunto \(O\) de ocho
puntos con una inclusi\'on distinguida \(j:H\hookrightarrow O\). En la
residencia de Witt, \(H\) es una hexada de un dodecad y \(O\) su octada
elevada; en este cap\'itulo s\'olo se usa la inclusi\'on ya fijada, no el
valor terminal de Barbero--Immirzi. Escribimos
\begin{equation}
\label{rm:eq:modular-pairs-H}
\binom H2:=\{P\subset H:|P|=2\},
\qquad \left|\binom H2\right|=\binom62=15,
\end{equation}
y definimos dos espacios de incidencias con tipos distintos:
\begin{equation}
\label{rm:eq:modular-F6-F8}
\mathcal F_6:=H\times\binom H2,
\qquad
\mathcal F_8:=O\times\binom H2.
\end{equation}
La inclusi\'on de incidencias es
\begin{equation}
\label{rm:eq:modular-iota68}
\iota_{6,8}:\mathcal F_6\hookrightarrow\mathcal F_8,
\qquad
(h,P)\longmapsto(j(h),P).
\end{equation}
Es inyectiva porque \(j\) lo es y conserva el par marcado \(P\).

\begin{theorem}[Conteos de incidencia y defecto orientado]
\label{rm:thm:modular-incidence-90-120}
Los dos espacios de \cref{rm:eq:modular-F6-F8} satisfacen
\begin{equation}
\label{rm:eq:modular-counts-90-120}
|\mathcal F_6|=6\binom62=90,
\qquad
|\mathcal F_8|=8\binom62=120,
\qquad
|\mathcal F_8\setminus\iota_{6,8}(\mathcal F_6)|=30.
\end{equation}
La normalizaci\'on por las veinticuatro coordenadas del par
dodecad--veinticuatro y por los quince pares de \(H\) produce
\begin{equation}
\label{rm:eq:modular-incidence-defect}
\boxed{
\delta_{6,8}^{\rm inc}
=\frac{|\mathcal F_6|-|\mathcal F_8|}
{24\binom62}=-\frac1{12}.}
\end{equation}
La misma fracci\'on se obtiene mediante el funcional rec\'iproco asociado
a la memoria incorporada,
\begin{equation}
\label{rm:eq:modular-reciprocal-defect}
\boxed{
\delta_{6,8}^{\rm rec}
=\frac{30}{120}-\frac{30}{90}=-\frac1{12}.}
\end{equation}
Son dos realizaciones exactas del mismo defecto orientado, pero tienen
dominios distintos y no se identifican con \(\zeta(-1)\) sin un
entrelazador adicional.
\end{theorem}

\begin{proof}
El principio multiplicativo aplicado a
\cref{rm:eq:modular-F6-F8} da \(6\cdot15=90\) y \(8\cdot15=120\).
La imagen de \(\iota_{6,8}\) tiene noventa elementos, de modo que el
complemento tiene treinta. Sustituir en
\cref{rm:eq:modular-incidence-defect} da
\(-30/(24\cdot15)=-1/12\), mientras
\cref{rm:eq:modular-reciprocal-defect} es
\(1/4-1/3=-1/12\).
\end{proof}

La transici\'on hexada--octada admite dos realizaciones funcionales que no
se deducen mutuamente. El diagrama algebraico, formulado sobre los espacios
que justifican los conteos, es
\begin{equation}
\label{rm:eq:modular-68-two-readers}
\begin{aligned}
(H\xhookrightarrow{\,j\,}O)
&\longmapsto
\left(|\mathcal F_6|,|\mathcal F_8|\right)=(90,120)
\longmapsto\Gamma_{6\to8}=\gammaH,\\
(h_6,o_8,t)=(6,8,3)
&\longmapsto-\frac{o_8-h_6}{t}=-\frac23,
\end{aligned}
\end{equation}
donde \(-2/3\) es el coeficiente de la tercera coordenada de la fila
\(13\) del constructor CKM. La primera realizaci\'on transforma la
incidencia en \`area y memoria modular; la segunda transforma la misma
transici\'on en par\'ametros de mezcla angular.
No se postula \(\boldsymbol\theta_{\rm CKM}=f(\gammaH)\): ambas salidas
conservan un antecedente com\'un.

El morfismo que determina los sectores \(90/120\) es
\(\iota_{6,8}\). La magnitud \(\gammaH\) se adopta como teorema previamente
demostrado en la construcci\'on hexada--octada y no se reconstruye a partir
de su representaci\'on decimal. Las
consecuencias operatorias comienzan s\'olo despu\'es de fijar la escala
elemental \(a_0\), que se define a continuaci\'on sin usar entrop\'ia ni el
producto terminal \(\gammaH a_0\).

\subsection{Correspondencia entre el interior non\'adico y los grados de libertad de borde}

Sea \(\mathbb T_3=\{0,1,2\}\). Todo \(x\in\mathbb Z/9\mathbb Z\) posee
una \`unica escritura de d\'igitos
\(x=x_0+3x_1\), con \(x_i\in\mathbb T_3\), y todo
\(y\in\mathbb Z/27\mathbb Z\) una \`unica escritura
\(y=y_0+3y_1+9y_2\). Por ello
\begin{equation}
\label{rm:eq:modular-bulk-boundary-729}
\mathcal B=(\mathbb Z/9\mathbb Z)^3,
\qquad
\mathcal B_\partial=(\mathbb Z/27\mathbb Z)^2,
\qquad
|\mathcal B|=|\mathcal B_\partial|=3^6=729.
\end{equation}
La igualdad de cardinales se acompa\~na del mapa expl\'icito. Si
\(x_j=r_{2j-1}+3r_{2j}\) para \(j=1,2,3\), definimos
\begin{equation}
\label{rm:eq:modular-beta729}
\beta_{729}(x_1,x_2,x_3)
=\bigl(r_1+3r_2+9r_3,\;r_4+3r_5+9r_6\bigr).
\end{equation}

\begin{proposition}[Bijecci\'on geneal\'ogica \(729\leftrightarrow729\)]
\label{rm:prop:modular-beta729}
El mapa \(\beta_{729}:\mathcal B\to\mathcal B_\partial\) es biyectivo y
preserva la palabra ordenada de seis trits. No es, en general, un
isomorfismo de grupos aditivos: reagrupar los d\'igitos cambia d\'onde se
registra el acarreo.
\end{proposition}

\begin{proof}
Extraer los seis d\'igitos de tres coordenadas non\'adicas y agruparlos en
dos ternas es una permutaci\'on de coordenadas de \(\mathbb T_3^6\). La
descomposici\'on posicional es \`unica en ambos lados, por lo que la
operaci\'on tiene inversa. La advertencia aditiva se comprueba, por ejemplo,
al sumar palabras que generan acarreo entre \(r_3\) y \(r_4\): esa
frontera separa las dos coordenadas de orden veintisiete, pero atraviesa
dos coordenadas distintas en el empaquetado non\'adico.
\end{proof}

En el toro de borde \(\mathcal B_\partial\), tomamos el bloque de v\'ertices
\begin{equation}
\label{rm:eq:modular-block-A}
A=\{0,1,\ldots,8\}\times\{0,1,\ldots,8\}.
\end{equation}
Sus enlaces orientados que cruzan el borde se separan por direcci\'on:
\begin{align}
\partial A_{90}
&=\{((-1,j),(0,j)),((8,j),(9,j)):0\leq j\leq8\},
\label{rm:eq:modular-boundary90}\\
\partial A_{120}
&=\{((i,-1),(i,0)),((i,8),(i,9)):0\leq i\leq8\},
\label{rm:eq:modular-boundary120}
\end{align}
donde todas las coordenadas se leen m\'odulo veintisiete. As\'i
\begin{equation}
\label{rm:eq:modular-boundary-counts}
\partial A=\partial A_{90}\sqcup\partial A_{120},
\qquad
|\partial A_{90}|=|\partial A_{120}|=18.
\end{equation}
Los sub\'indices \(90,120\) etiquetan la procedencia de canal; no dicen que
el borde tenga noventa o ciento veinte enlaces.

\begin{theorem}[Ley areal del corte tri\'adico]
\label{rm:thm:modular-triadic-cut}
En el grafo c\'ubico \(N\times N\times N\) con capacidad unitaria, todo
corte entre dos caras opuestas posee capacidad al menos \(N^2\), y una
capa transversal alcanza esa cota. Para \(N=9\),
\begin{equation}
\label{rm:eq:modular-mincut-entropy}
\mathcal A_{\min}=81,
\qquad
S_{\rm tri}=81\log3.
\end{equation}
\end{theorem}

\begin{proof}
Hay \(N^2\) rectas de aristas paralelas a la direcci\'on normal, disjuntas
dos a dos. Todo corte debe interceptar cada una, de modo que su capacidad
es al menos \(N^2\). Una sola capa transversal contiene exactamente
\(N^2\) aristas y realiza el m\'inimo. Si cada enlace cortado porta un trit
m\'aximamente mezclado, su entrop\'ia es \(\log3\); la aditividad sobre los
ochenta y un enlaces da \cref{rm:eq:modular-mincut-entropy}.
\end{proof}

La entrop\'ia \(S_{\rm tri}\) mide capacidad del corte c\'ubico. La
entrop\'ia modular de los treinta y seis enlaces de
\cref{rm:eq:modular-boundary-counts} mide otro estado y no debe igualarse con
ella.

\subsection{Entrelazador colectivo de incidencia entre interior y borde}

En la suma
\(\ell^2(\mathcal F_6)\oplus\ell^2(\mathcal F_8)\), definimos
\begin{equation}
\label{rm:eq:modular-u90-u120}
u_{90}=\frac1{\sqrt{90}}\sum_{f\in\mathcal F_6}(\delta_f,0),
\qquad
u_{120}=\frac1{\sqrt{120}}\sum_{f\in\mathcal F_8}(0,\delta_f).
\end{equation}
En
\(\ell^2(\partial A_{90})\oplus\ell^2(\partial A_{120})\), sean
\begin{equation}
\label{rm:eq:modular-v90-v120}
v_{90}=\frac1{\sqrt{18}}\sum_{e\in\partial A_{90}}(\delta_e,0),
\qquad
v_{120}=\frac1{\sqrt{18}}\sum_{e\in\partial A_{120}}(0,\delta_e).
\end{equation}
Las parejas \((u_{90},u_{120})\) y \((v_{90},v_{120})\) son ortonormales.
Por tanto existe una \`unica isometr\'ia entre los subespacios colectivos
que satisface
\begin{equation}
\label{rm:eq:modular-J68}
\mathcal J_{6\to8}u_{90}=v_{90},
\qquad
\mathcal J_{6\to8}u_{120}=v_{120}.
\end{equation}
Definamos
\begin{align}
D_{\rm inc}
&=90|u_{90}\rangle\langle u_{90}|
+120|u_{120}\rangle\langle u_{120}|,
\label{rm:eq:modular-Dinc}\\
D_{\rm area}
&=90|v_{90}\rangle\langle v_{90}|
+120|v_{120}\rangle\langle v_{120}|.
\label{rm:eq:modular-Darea}
\end{align}

\begin{theorem}[Transporte exacto de las dos etiquetas]
\label{rm:thm:modular-J68}
La aplicaci\'on \(\mathcal J_{6\to8}\) es unitaria entre los dos planos
colectivos y entrelaza los operadores de incidencia y \`area:
\begin{equation}
\label{rm:eq:modular-J-intertwining}
\boxed{
\mathcal J_{6\to8}D_{\rm inc}
=D_{\rm area}\mathcal J_{6\to8}.}
\end{equation}
Por tanto, toda combinaci\'on polin\'omica de las etiquetas \(90,120\), en
particular la combinaci\'on primitiva \(12:-1\), se transporta con los
mismos coeficientes.
\end{theorem}

\begin{proof}
La ortonormalidad muestra que \(\mathcal J_{6\to8}\) env\'ia una base
ortonormal en otra. En \(u_{90}\) ambos lados de
\cref{rm:eq:modular-J-intertwining} valen \(90v_{90}\); en \(u_{120}\),
ambos valen \(120v_{120}\). La igualdad en una base prueba la identidad.
El c\'alculo funcional polin\'omico conserva el entrelazamiento.
\end{proof}

Este mapa no pretende inyectar noventa incidencias individuales en
dieciocho enlaces: transporta exactamente los dos modos uniformes y sus
etiquetas espectrales. Extenderlo a fluctuaciones no uniformes exigir\'ia
un nuevo mapa de incidencia; el operador modular definido aqu\'i s\'olo emplea
el subespacio colectivo que acaba de construirse.

\subsection{Cuanto elemental y escala mínima del operador de área}

Sea \(C^*\) la coordenada angular conjugada, expresada en grados, y fijemos
\begin{equation}
\label{rm:eq:modular-theta-clock-a0}
\theta_{\rm clk}=\frac{C^*}{6}\frac{\pi}{180},
\qquad
\boxed{
a_0=\frac{1+2\cos(10^\circ)}{3}\,\theta_{\rm clk}.}
\end{equation}
El factor que precede a \(\theta_{\rm clk}\) no es una correcci\'on
decimal. Es el car\'acter real normalizado del triplete de fases
\begin{equation}
\label{rm:eq:modular-c10-character}
\frac13\Tr\operatorname{diag}
\left(1,\ee^{\ii\pi/18},\ee^{-\ii\pi/18}\right)
=\frac{1+2\cos(10^\circ)}3.
\end{equation}
As\'i, \(a_0\) es la media espectral de la estructura angular sobre el trit
orientado. Es positiva, se fija antes del operador de \`area y no se
selecciona a partir de la entrop\'ia. La identificaci\'on posterior
\(\gammaH=\Gamma_{6\to8}\) realiza la salida HMT en la conexi\'on de
Ashtekar--Barbero \cite{Barbero1995,Immirzi1997}; no reabre su derivaci\'on.

\begin{remark}[Alcance del s\'imbolo \(a_0\)]
En este cap\'itulo, \(a_0\equiv a_0^{\rm area}\) denota exclusivamente la
normalizaci\'on de \cref{rm:eq:modular-theta-clock-a0}. No es el radio de Bohr
ni el factor de escala cosmol\'ogico inicial, aunque esos objetos se escriban
tradicionalmente con el mismo s\'imbolo. Ninguna identidad de este cap\'itulo
permite sustituir uno por otro.
\end{remark}
 
\paragraph{Construcción del operador de área a partir del corte triádico y separación de canales.} El corte y su capacidad determinan un operador de área; su espectro debe conservar por separado ocupación total y procedencia de canal.
 \endgroup

\chapter{Espectro del operador de área y reconstrucción de la procedencia de canal}
\begingroup
\renewcommand{\chapter}[2][]{}%
\chapter{Espectro del operador de área y reconstrucción conjunta de ocupación y procedencia}
\label{chap:parte-iv-59}

Los sectores \(90\) y \(120\) definen operadores compatibles de ocupación y memoria. La inversión del sistema lineal reconstruye ambos canales y evita que el área total borre la procedencia.

\section{Operadores areal y modular en los sectores \texorpdfstring{\(90\)}{90} y \texorpdfstring{\(120\)}{120}}

En cada una de las dos familias de
\cref{rm:eq:modular-boundary-counts} existen dieciocho canales qutrit. Para
\(m\in\{90,120\}\), definimos
\begin{equation}
\label{rm:eq:modular-number-operators}
N_m=\sum_{e\in\partial A_m}N_e,
\qquad \Spec(N_e)=\{0,1,2\}.
\end{equation}
La salida de Barbero--Immirzi se usa aqu\'i como resultado previo, sin
rederivarla. Con la normalizaci\'on areal construida en
\cref{rm:eq:modular-theta-clock-a0}, resulta
\begin{equation}
\label{rm:eq:modular-area-operator}
\boxed{
\frac{\widehat{\mathcal A}_{\rm phys}}{\ell_P^2}
=\gammaH a_0(N_{90}+N_{120}),}
\qquad
\gammaH a_0=0.0016755940749224991.
\end{equation}

\begin{theorem}[Espectro finito del cuanto de \`area]
\label{rm:thm:modular-area-spectrum}
Sobre los treinta y seis enlaces del corte,
\begin{equation}
\label{rm:eq:modular-total-number-spectrum}
\Spec(N_{90}+N_{120})=\{0,1,\ldots,72\}.
\end{equation}
La multiplicidad del autovalor \(n\) es el coeficiente de \(z^n\) en
\((1+z+z^2)^{36}\). En consecuencia,
\begin{equation}
\label{rm:eq:modular-area-spectrum}
\boxed{
\Spec\!\left(\frac{\widehat{\mathcal A}_{\rm phys}}{\ell_P^2}\right)
=\{n\gammaH a_0:0\leq n\leq72\}.}
\end{equation}
Los cortes compatibles prolongan la familia
\(\mathcal A_n^{\rm prot}=n\gammaH a_0\ell_P^2\).
\end{theorem}

\begin{proof}
Cada operador \(N_e\) tiene funci\'on generatriz espectral
\(1+z+z^2\). Los treinta y seis factores conmutan y act\'uan sobre factores
tensoriales distintos; por tanto, la funci\'on generatriz del total es
\((1+z+z^2)^{36}\), cuyos grados recorren todos los enteros de cero a
setenta y dos. Multiplicar por el escalar positivo
\(\gammaH a_0\ell_P^2\) da \cref{rm:eq:modular-area-spectrum}.
\end{proof}

Es necesario distinguir dos procedimientos de cuantizaci\'on. La ley de corte
\cref{rm:thm:modular-triadic-cut} cuenta ochenta y una capacidades unitarias y
determina \(81\log3\); el operador
\cref{rm:eq:modular-area-operator} cuantiza ocupaciones de treinta y seis
enlaces de respuesta, hasta setenta y dos trits de n\'umero. Comparten la
misma hoja de borde, pero no el mismo observable.

La torsi\'on modular produce
\begin{equation}
\label{rm:eq:modular-delta-value}
\Dmod
=0.0001111970500599773249414566131933856\ldots,
\end{equation}
y el Hamiltoniano de borde
\begin{equation}
\label{rm:eq:modular-hamiltonian}
\boxed{
\HHM^{(A)}=-\log\rho_A
=cI+\Dmod(90N_{90}+120N_{120}).}
\end{equation}
El calificativo \emph{hologr\'afico} no designa una analog\'ia: la
incidencia \(6\to8\) se transporta a los conjuntos de borde
\(\partial A_{90}\) y \(\partial A_{120}\), y el estado reducido reside
en el \`algebra generada por sus operadores n\'umero.

\begin{theorem}[Reconstrucci\'on conjunta de geometr\'ia y memoria]
\label{rm:thm:modular-area-memory}
Sean
\begin{equation}
\label{rm:eq:modular-ND}
N=N_{90}+N_{120},
\qquad
D=90N_{90}+120N_{120}.
\end{equation}
Entonces
\begin{equation}
\label{rm:eq:modular-inverse-channels}
\boxed{
N_{90}=4N-\frac{D}{30},
\qquad
N_{120}=\frac{D}{30}-3N.}
\end{equation}
En particular,
\([\widehat{\mathcal A}_{\rm phys},\HHM^{(A)}]=0\), y el par
\((\widehat{\mathcal A}_{\rm phys},\HHM^{(A)}-cI)\) reconstruye los dos
sectores de ocupaci\'on.
\end{theorem}

\begin{proof}
Los operadores \(N_{90}\) y \(N_{120}\) act\'uan sobre factores distintos
y conmutan. El mapa lineal que lleva \((N_{90},N_{120})\) a \((N,D)\)
tiene matriz
\[
\begin{pmatrix}1&1\\90&120\end{pmatrix},
\qquad \det=30.
\]
Su inversa es exactamente \cref{rm:eq:modular-inverse-channels}. Como
\cref{rm:eq:modular-area-operator,rm:eq:modular-hamiltonian} son combinaciones
lineales de los mismos operadores n\'umero, el conmutador se anula.
\end{proof}

La interpretaci\'on operatoria es precisa: el observable de \`area conserva
el n\'umero total de excitaciones, mientras que el Hamiltoniano modular
conserva su sector de procedencia. Ninguno de ambos observables retiene por
s\'i solo las dos coordenadas; el par conjunto s\'i las determina.

\section{Confluencia de \`area, informaci\'on y energ\'ia}

La escala angular \(\theta_{\rm clk}\) de
\cref{rm:eq:modular-theta-clock-a0} y la temperatura cronol\'ogica
\(\Theta_{\rm clk}\) son objetos distintos; el uso de may\'uscula es parte
del tipado. La realizaci\'on termodin\'amica del \cref{rm:ch:metrology} demuestra
\begin{equation}
\label{rm:eq:modular-planck-trit-bridge}
E_P=k_B\Theta_{\rm clk}\log3,
\end{equation}
mientras la realizaci\'on areal construye el incremento elemental
\begin{equation}
\label{rm:eq:modular-area-quantum}
\delta\mathcal A
=\gammaH a_0\ell_P^2.
\end{equation}
No se igualan: \cref{rm:eq:modular-planck-trit-bridge} fija el cuanto
informacional de energ\'ia y \cref{rm:eq:modular-area-quantum} el cuanto
geom\'etrico de \`area. Su cociente define la tensi\'on de conversi\'on
tipada
\begin{equation}
\label{rm:eq:modular-information-area-tension}
\boxed{
\mathcal T_{I/A}
=\frac{E_P}{\delta\mathcal A}
=\frac{k_B\Theta_{\rm clk}\log3}
{\gammaH a_0\ell_P^2}.}
\end{equation}
La expresi\'on no introduce una constante adicional: determina las
normalizaciones que deben conservarse en la transformaci\'on de informaci\'on
en geometr\'ia.

\enlargethispage{7\baselineskip}
En el corte m\'inimo de ochenta y un enlaces, la capacidad informacional
f\'isica y la energ\'ia de un trit por enlace son
\begin{equation}
\label{rm:eq:modular-cut-info-energy}
S_{\rm cut}^{\rm phys}=81k_B\log3,
\qquad
E_{\rm cut}^{(1)}=81E_P
=\Theta_{\rm clk}S_{\rm cut}^{\rm phys}.
\end{equation}
La igualdad final es una ecuaci\'on de estado de esta realizaci\'on, no una ley de
equipartici\'on universal. En particular, el estado Gibbs de los treinta y
seis enlaces del borde tiene entrop\'ia distinta y se trata a continuaci\'on.

Esta separaci\'on determina el contenido estructural de la confluencia:
Barbero--Immirzi
orienta el espectro de \`area; \(\log3\) mide la decisi\'on irreducible del
TRIT; \(\Theta_{\rm clk}\) convierte esa informaci\'on en energ\'ia; y
\(\ell_P^2\) conecta el cuanto de \`area con la familia gravitatoria. La
gravedad queda descrita como confluencia de realizaciones funcionales de un
antecedente com\'un, y no como una igualdad impuesta entre sus valores.
 
\paragraph{Transporte de la memoria de canal al Hamiltoniano modular y a la purificación termocampo doble.} La memoria de canal se representa a continuación mediante el Hamiltoniano modular y una purificación termocampo doble.
 \endgroup

\chapter{Hamiltoniano modular de frontera, purificación termocampo doble e información orientada}
\begingroup
\renewcommand{\chapter}[2][]{}%
\chapter{Hamiltoniano modular de borde, purificación termocampo doble y déficit informacional orientado}
\label{chap:parte-iv-60}

El estado qutrit fiel admite una purificación termocampo doble. El Hamiltoniano modular registra la orientación de los canales y el déficit informacional mide la separación finita entre estados sin identificarse con el área.

\section{Estado qutrit, purificaci\'on termocampo doble y d\'eficit informacional}

Para un enlace con peso \(w>0\), sean
\begin{equation}
\label{rm:eq:modular-link-state}
Z(w)=1+\ee^{-w}+\ee^{-2w},
\qquad
\rho(w)=\frac1{Z(w)}\sum_{n=0}^2\ee^{-wn}|n\rangle\langle n|.
\end{equation}
Los pesos de las dos familias son
\begin{align}
w_{90}&=90\Dmod
=0.0100077345053979592447310951874\ldots,
\label{rm:eq:modular-w90}\\
w_{120}&=120\Dmod
=0.0133436460071972789929747935832\ldots,
\qquad \frac{w_{120}}{w_{90}}=\frac43.
\label{rm:eq:modular-w120}
\end{align}
El estado finito de borde es
\begin{equation}
\label{rm:eq:modular-product-state}
\rho_A=\rho(w_{90})^{\otimes18}
\otimes\rho(w_{120})^{\otimes18}.
\end{equation}
Su logaritmo reproduce \cref{rm:eq:modular-hamiltonian}, con
\(c=18\log Z(w_{90})+18\log Z(w_{120})\).

\begin{theorem}[Purificaci\'on termocampo doble]
\label{rm:thm:modular-tfd}
El vector
\begin{equation}
\label{rm:eq:modular-tfd-link}
|\operatorname{TFD}(w)\rangle
=\frac1{\sqrt{Z(w)}}\sum_{n=0}^2
\ee^{-wn/2}|n\rangle_A\otimes|n\rangle_{\bar A}
\end{equation}
purifica \(\rho(w)\). El producto de dieciocho copias para cada uno de los
pesos \(w_{90},w_{120}\) purifica \(\rho_A\); en el sistema duplicado,
\(S(\rho_A)\) es exactamente entrop\'ia de entrelazamiento.
\end{theorem}

\begin{proof}
Al trazar el segundo factor, la ortogonalidad de \(|n\rangle_{\bar A}\)
elimina los t\'erminos cruzados y deja
\[
\operatorname{Tr}_{\bar A}
|\operatorname{TFD}(w)\rangle\langle\operatorname{TFD}(w)|
=\rho(w).
\]
La traza parcial conmuta con productos tensoriales, lo que prueba la
afirmaci\'on para los treinta y seis enlaces.
\end{proof}

Para un enlace,
\begin{equation}
\label{rm:eq:modular-link-entropy}
S(w)=\log Z(w)
+w\frac{\ee^{-w}+2\ee^{-2w}}{Z(w)}.
\end{equation}
En consecuencia,
\begin{equation}
\label{rm:eq:modular-entropy-value}
S(\rho_A)=18S(w_{90})+18S(w_{120})
=39.54837320881807994957319730\ldots\ \text{nats}.
\end{equation}
El estado m\'aximamente mezclado de treinta y seis qutrits tiene entrop\'ia
\(36\log3\). Definimos el defecto orientado
\begin{equation}
\label{rm:eq:modular-oriented-information}
\boxed{
I_{\rm or}=36\log3-S(\rho_A)
=0.001669183233868940655631225913\ldots\ \text{nats}.}
\end{equation}
Su positividad es consecuencia de la maximalidad entr\'opica del estado
uniforme. La cercan\'ia de \(I_{\rm or}\) a \(\gammaH a_0\) no constituye
una identidad y no se usa para seleccionar ninguno de los dos valores.
 
\section{Información par y transporte geométrico impar en la bicapa}

La distribución normalizada entre hojas separa la entropía par del transporte
hexada--octada impar. Esta paridad bicapa conserva la información de borde que
el Hamiltoniano modular registra y evita identificar entropía, área y
orientación.

\begingroup
\let\chapter\section
\let\section\subsection
\let\subsection\subsubsection
\let\subsubsection\paragraph
\let\clearpage\relax
\let\cleardoublepage\relax
\chapter{La bicapa: información par y transporte geométrico impar}
\label{chap:paridad-bicapa-barbero-rev6}
\index{entropía!bicapa}
\index{Barbero--Immirzi!paridad}
\index{hoja!intercambio}

Barbero--Immirzi y la entropía de las dos hojas no son dos cantidades
yuxtapuestas. Nacen de los mismos canales \(q_\pm\), pero conservan datos
distintos. La normalización probabilística olvida la escala común y mide la
distribución entre orientaciones; el funcional hexada--octada conserva el
signo de esa orientación. La primera lectura es par y la segunda impar.

\section{Distribución normalizada entre hojas}

Escribamos
\[
 x=A_{\rm rad},
 \qquad
 y=C^*_{\rm rad},
 \qquad
 q_-=e^{-x+y},
 \qquad
 q_+=e^{-x-y}.
\]
La suma \(q_-+q_+=2e^{-x}\cosh y\) permite definir
\[
 p_-=\frac{q_-}{q_-+q_+}
     =\frac{e^y}{2\cosh y},
 \qquad
 p_+=\frac{q_+}{q_-+q_+}
     =\frac{e^{-y}}{2\cosh y}.
\]
El factor común \(e^{-x}\) desaparece. La distribución conserva la asimetría
entre hojas, pero no la escala angular común.

\begin{definition}[Entropía bicapa]
\label{def:entropia-bicapa-rev6}
La entropía de la distribución orientada es
\[
 S_{\rm bi}(y)
 :=
 -p_-\log p_--p_+\log p_+
 =
 \log(2\cosh y)-y\tanh y.
\]
\end{definition}

\begin{proposition}[Paridad y máximo de la entropía bicapa]
\label{prop:paridad-entropia-bicapa-rev6}
La función \(S_{\rm bi}\) es par y satisface
\[
 S_{\rm bi}(-y)=S_{\rm bi}(y),
 \qquad
 S_{\rm bi}'(y)=-y\,\operatorname{sech}^2y,
 \qquad
 0<S_{\rm bi}(y)\leq\log2.
\]
El máximo \(\log2\) se alcanza únicamente en \(y=0\), cuando ambas hojas
tienen el mismo peso.
\end{proposition}

\begin{proof}
La paridad se sigue de que \(\cosh\) es par y \(y\tanh y\) también lo es.
La derivación directa produce
\[
 \frac{\dd}{\dd y}\log(2\cosh y)=\tanh y,
 \qquad
 \frac{\dd}{\dd y}(y\tanh y)
 =\tanh y+y\operatorname{sech}^2y.
\]
Por tanto \(S_{\rm bi}'(y)=-y\operatorname{sech}^2y\): la función crece en
\((-\infty,0)\), decrece en \((0,\infty)\) y toma en el origen el valor
\(\log2\). La positividad es la de la entropía de una distribución con dos
pesos estrictamente positivos.
\end{proof}

\section{Paridad del transporte hexada--octada}

Sea
\[
 s_n(q_-,q_+)=q_-^n-q_+^n.
\]
El intercambio de hojas es la involución
\[
 \mathsf J_{\rm sh}:(x,y)\longmapsto(x,-y),
\]
que permuta \(q_-\leftrightarrow q_+\). De aquí,
\[
 s_n\circ\mathsf J_{\rm sh}=-s_n.
\]
El funcional construido en el \cref{mdg:chap:barbero} es
\[
 \Gamma_{6\to8}(A,C^*)
 =
 \frac{\pi A C^*}{180}
{}+12\bigl(S_{90}(q_-)-S_{90}(q_+)\bigr)
-\bigl(S_{120}(q_-)-S_{120}(q_+)\bigr).
\]

\begin{theorem}[Teorema de paridad bicapa]
\label{thm:paridad-bicapa-barbero-rev6}
Bajo el intercambio \(C^*\mapsto-C^*\),
\[
 \boxed{S_{\rm bi}(-C^*_{\rm rad})
        =S_{\rm bi}(C^*_{\rm rad}),}
\]
\[
 \boxed{\Gamma_{6\to8}(A,-C^*)
        =-\Gamma_{6\to8}(A,C^*).}
\]
La misma bicapa produce, por tanto, una lectura informacional par y un
transporte geométrico impar.
\end{theorem}

\begin{proof}
La primera igualdad es la
\cref{prop:paridad-entropia-bicapa-rev6}. En la segunda, el término
\(\pi AC^*/180\) cambia de signo y cada diferencia de Lambert cambia de
signo porque se intercambian sus dos argumentos. La combinación completa es
impar.
\end{proof}

\section{Interpretación del operador constitutivo como acoplamiento entre vacío, campo y respuesta}

La entropía bicapa y \(\Gamma_{6\to8}\) no son dos nombres para el mismo
número. La primera cuantifica cuánto peso comparten las dos orientaciones
después de cancelar la escala común. El segundo conserva la orientación y la
incidencia del paso \(6\to8\). En el punto HMT,
\[
 S_{\rm bi}
 =0.69246069960358548\ldots\ {\rm nats},
 \qquad
 \Gamma_{6\to8}
 =0.27400767269445927\ldots.
\]
El vínculo entre información y geometría consiste en que ambas lecturas
factorizan por los mismos canales \(q_\pm\), no en igualar sus valores.

La banda partícula--antipartícula del
\cref{chap:banda-particula-virtual} reproduce la misma distinción de
paridad: la masa y la cantidad de correlación descienden al cociente de
banda, mientras carga, orientación y transporte cambian de signo. En la
acción de Cartan--Holst, \(\Gamma_{6\to8}\) ocupa la coordenada orientada;
en la lectura informacional, \(S_{\rm bi}\) cuenta la distribución entre
hojas. Así se explica por qué una sola estructura puede conservar
información e invertir orientación sin confundir ambas operaciones.
 \endgroup

\paragraph{Variación tangente del estado modular y derivación de la ley areal.} La variación tangente del estado modular produce la primera ley areal; las variaciones finitas conservan el término de entropía relativa.
 \endgroup

\chapter{Primera ley modular areal y extensión finita por entropía relativa}
\begingroup
\renewcommand{\chapter}[2][]{}%
\chapter{1.ª ley modular areal y extensión finita mediante entropía relativa}
\label{chap:parte-iv-61}

La primera ley modular relaciona variaciones de entropía y área dentro del sistema finito declarado. Su extensión no lineal retiene el defecto positivo de entropía relativa y, por ello, no confunde igualdad tangente con identidad global.

\section{1.ª ley modular para variaciones del àrea}

Definamos las dos \`areas normalizadas
\begin{equation}
\label{rm:eq:modular-normalized-area-channels}
\mathcal A_m=\gammaH a_0N_m,
\qquad m\in\{90,120\}.
\end{equation}
Entonces
\begin{equation}
\label{rm:eq:modular-beta-definition}
\HHM^{(A)}-cI
=\beta_{90}\mathcal A_{90}+\beta_{120}\mathcal A_{120},
\qquad
\beta_m=\frac{m\Dmod}{\gammaH a_0}.
\end{equation}
Los valores terminales son
\begin{align}
\beta_{90}&=5.9726485401070929914651798889\ldots,
\label{rm:eq:modular-beta90}\\
\beta_{120}&=7.9635313868094573219535731852\ldots,
\qquad
\frac{\beta_{120}}{\beta_{90}}=\frac43.
\label{rm:eq:modular-beta120}
\end{align}

\begin{theorem}[Primera ley modular en variables areales]
\label{rm:thm:modular-first-law}
Para una familia diferenciable de estados normalizados
\(\rho(\varepsilon)\) que pasa por \(\rho_A\), manteniendo fijos
\(\Dmod,\gammaH,a_0\),
\begin{equation}
\label{rm:eq:modular-first-law}
\boxed{
\delta S
=\beta_{90}\,\delta\langle\mathcal A_{90}\rangle
+\beta_{120}\,\delta\langle\mathcal A_{120}\rangle.}
\end{equation}
\end{theorem}

\begin{proof}
La primera ley de entrelazamiento en el estado de referencia es
\(\delta S=\delta\langle-\log\rho_A\rangle\). El t\'ermino escalar
\(cI\) no contribuye porque \(\delta\Tr\rho=0\). Sustituir
\cref{rm:eq:modular-beta-definition} prueba la identidad.
\end{proof}

La f\'ormula es local en el espacio de estados. Para cambios finitos aparece
la entrop\'ia relativa; omitirla convertir\'ia una identidad tangente en una
afirmaci\'on falsa de linealidad global.

\Needspace{22\baselineskip}
\begin{theorem}[Identidad finita con t\'ermino de defecto]
\label{rm:thm:modular-finite-first-law}
Para todo estado \(\sigma\) cuyo soporte est\'e contenido en el de
\(\rho_A\),
\begin{equation}
\label{rm:eq:modular-finite-first-law}
\boxed{
S(\sigma)-S(\rho_A)
=\beta_{90}\Delta_\sigma\langle\mathcal A_{90}\rangle
+\beta_{120}\Delta_\sigma\langle\mathcal A_{120}\rangle
-D(\sigma\Vert\rho_A),}
\end{equation}
donde
\(D(\sigma\Vert\rho_A)=\Tr\sigma(\log\sigma-\log\rho_A)\geq0\).
En particular, la variaci\'on lineal del
\cref{rm:thm:modular-first-law} es el t\'ermino tangente de esta identidad.
\end{theorem}

\begin{proof}
Por definici\'on y por \(\HHM^{(A)}=-\log\rho_A\),
\[
D(\sigma\Vert\rho_A)
=-S(\sigma)+\Tr(\sigma\HHM^{(A)}).
\]
Para \(\sigma=\rho_A\), el miembro izquierdo es cero y
\(S(\rho_A)=\Tr(\rho_A\HHM^{(A)})\). Restar ambas identidades da
\[
S(\sigma)-S(\rho_A)
=\Delta_\sigma\langle\HHM^{(A)}\rangle
-D(\sigma\Vert\rho_A).
\]
El escalar \(cI\) se cancela entre estados normalizados; sustituir
\cref{rm:eq:modular-beta-definition} produce
\cref{rm:eq:modular-finite-first-law}. La positividad de la entrop\'ia
relativa es la desigualdad de Klein.
\end{proof}

La correcci\'on \(D(\sigma\Vert\rho_A)\) mide la distancia informacional
entre el estado real y el estado modular de referencia. As\'i, la primera
ley no pierde contenido al salir del r\'egimen infinitesimal: se convierte
en una identidad exacta con un defecto positivo controlado.
 
\paragraph{Persistencia inductiva de los pesos modulares y clasificación factorial del límite.} La persistencia inductiva de los pesos modulares conduce al cierre operatorial infinito y a su clasificación factorial.
 \endgroup

\chapter{Persistencia de las incidencias \texorpdfstring{\(90/120\)}{90/120} y límite factorial hiperinfinito}
\begingroup
\renewcommand{\chapter}[2][]{}%
\chapter{Persistencia de los canales \texorpdfstring{\(90/120\)}{90/120} y clasificación del límite como factor hiperinfinito \texorpdfstring{\(III_{\lambda_\partial}\)}{III(lambda parcial)}}
\label{chap:parte-iv-62}

Los productos locales fieles conservan las razones de los canales \(90\) y \(120\) bajo refinamiento. La representación GNS y el límite ITPFI determinan el factor hiperinfinito de tipo III con parámetro de borde explícito.

\section{L\'imite local y clasificaci\'on modular del sistema infinito}

En una UHF infinita no existe, en general, una matriz densidad traza-clase
global dentro del \`algebra. El objeto correcto es la red
\begin{equation}
\label{rm:eq:modular-local-net}
\HHM^{(m)}=-\log\rho_m
\end{equation}
con restricciones de estados compatibles, o el generador modular en la
representaci\'on GNS. La expectativa tracial
\(\mathrm{id}\otimes\tau_9\) no preserva autom\'aticamente un Gibbs no
tracial y no sustituye la compatibilidad de estados.

\begin{theorem}[Tipo modular bajo persistencia fiel de ambos canales]
\label{rm:thm:modular-typeIII}
Sup\'ongase que:
\begin{enumerate}[label=(\roman*)]
\item cada corte es un producto de estados qutrit fieles con pesos
\(w_{90}\) y \(w_{120}\);
\item las inclusiones preservan el estado y la etiqueta de canal;
\item aparecen infinitamente muchos factores de cada tipo, ambos con
densidad asint\'otica positiva;
\item no se introducen cocientes espectrales adicionales y la
representaci\'on producto es factorial.
\end{enumerate}
Entonces el grupo aditivo de logaritmos de cocientes de autovalores es
\begin{equation}
\label{rm:eq:modular-ratio-group}
90\Dmod\Z+120\Dmod\Z=30\Dmod\Z.
\end{equation}
El cierre GNS ITPFI es el factor hiperinfinito
\begin{equation}
\label{rm:eq:modular-typeIII-lambda}
\boxed{
III_{\lambda_\partial},
\qquad
\lambda_\partial=\ee^{-30\Dmod}
=0.9966696464689573786108299769\ldots .}
\end{equation}
\end{theorem}

\begin{proof}
En cada enlace los cocientes de autovalores son potencias de
\(\ee^{-w_m}\). La persistencia infinita de ambos tipos genera el subgrupo
\(w_{90}\Z+w_{120}\Z\). Como \(\gcd(90,120)=30\), se obtiene
\cref{rm:eq:modular-ratio-group}. El criterio de raz\'on asint\'otica de
Araki--Woods para el producto fiel da el factor ITPFI hiperinfinito con
par\'ametro \(\ee^{-30\Dmod}\)
\cite{rm:ArakiWoods1968,rm:Takesaki}.
\end{proof}

El teorema declara sus hip\'otesis reales. Si desaparece uno de los canales,
si el producto no es compatible o si cambia el grupo de razones, debe
recalcularse el tipo.
 
\paragraph{Publicación gravitatoria del cierre modular mediante el portador pentacomponente.} El cierre modular prepara la publicación gravitatoria: primero se construye el portador pentacomponente desde la incidencia icosaédrica.
 \endgroup

\chapter{Descenso icosaédrico y portador gravitatorio pentacomponente}
\begingroup
\renewcommand{\chapter}[2][]{}%
\chapter{Descenso icosaédrico \texorpdfstring{\(A_5/C_5\)}{A₅/C₅} y selección del portador gravitatorio pentacomponente}
\label{chap:parte-iv-63}

La recta gravitatoria conserva la longitud construida y su realización
cronológica de \cref{g26:sec:rigidez-radial-es,cr28:sec:radial-clock};
su portador desciende después al cociente icosaédrico. El proyector
irreducible de rango \(5\) selecciona el portador tensorial sin asumir
todavía una realización masiva o radiativa.

\section{Constante gravitatoria intr\'inseca y representaci\'on tensorial}
\label{rm:ch:gravity}

\subsection{Cadena constitutiva del sector gravitatorio: escala cronológica \texorpdfstring{\(108\)}{108}, reducción tensorial, holonomía y evaluación metrológica de \texorpdfstring{\(G\)}{G}}

La gravedad se organiza aqu\'i en cuatro estratos que no deben fundirse:

\begin{equation}
\label{rm:eq:gravity-causal-chain}
\begin{aligned}
\text{inscripci\'on y norma radial R1--R8}
&\longrightarrow L
\longrightarrow G=c_{\rm int}^3L^2/\hbar_a,\\
\text{incidencia icosa\'edrica}
&\longrightarrow V_{12}\xrightarrow{P_5}V_5
\xrightarrow{\Gamma_5}\operatorname{STF}_3
\xrightarrow{\Lambda(n)}\mathcal H_{\rm TT}(n),\\
\text{calendario CT108 con }t_0=\pi L/(54c_{\rm int})
&\longrightarrow R_{108}=L,\quad
\theta_{108}^{\rm circ}=(54/\pi)^2,\\
\text{carta metrol\'ogica}
&\longrightarrow 6.6742996594140227\ldots\times10^{-11}\,
\mathrm{m^3\,kg^{-1}\,s^{-2}}.
\end{aligned}
\end{equation}

La primera l\'inea conserva la determinaci\'on de la longitud y del
acoplamiento mediante las reglas R1--R8 de
\cref{g26:sec:rigidez-radial-es}. La segunda construye el portador y
su reducci\'on transversal sin traza. La tercera es la realizaci\'on
cronol\'ogica de la longitud ya construida, y la cuarta eval\'ua la
misma secci\'on dimensional. La reducci\'on \(12\to5\to5\to2\)
no altera el coeficiente radial com\'un: intensidad y representaci\'on
conservan su procedencia y sus funciones respectivas.

\subsection{Recta de magnitud gravitatoria y escala cronol\'ogica \texorpdfstring{\(108\)}{108}}

Sean \(c_{\rm int}>0\), \(t_0>0\),
\(\ell_0=c_{\rm int}t_0\) y una secci\'on positiva
\(\hbar_a\) de la recta de acci\'on, con
\(a\in\{\mathrm{pre},\mathrm{ret}\}\). El transductor dimensional es

\begin{equation}
\label{rm:eq:gravity-transducer}
G_a(\theta)
=\theta\frac{c_{\rm int}^{5}t_0^2}{\hbar_a}
=\theta\frac{c_{\rm int}^{3}\ell_0^2}{\hbar_a},
\qquad \theta>0.
\end{equation}

La igualdad de las dos expresiones usa s\'olo
\(\ell_0=c_{\rm int}t_0\), y
\([G]=L^3M^{-1}T^{-2}\). La dimensi\'on queda fijada por la recta; el
car\'acter adimensional \(\theta\) distingue las realizaciones de esa
familia. En la composici\'on longitudinal ya demostrada,
\(t_0=\pi L/(54c_{\rm int})\) y
\(\theta=(54/\pi)^2\) reproducen exactamente
\(G_a=c_{\rm int}^3L^2/\hbar_a\).

El calendario CT108 define

\begin{equation}
\label{rm:eq:gravity-clock108}
T_{108}=108t_0,
\qquad
\omega_{108}=\frac{2\pi}{108t_0},
\qquad
R_{108}=\frac{c_{\rm int}}{\omega_{108}}
=\frac{54}{\pi}\ell_0.
\end{equation}

Para cada hoja de acci\'on, sean

\begin{equation}
\label{rm:eq:gravity-clock-energy-mass}
E_{108,a}=\hbar_a\omega_{108},
\qquad
m_{108,a}=\frac{E_{108,a}}{c_{\rm int}^2}.
\end{equation}

Entonces, la longitud de Compton reducida coincide con el radio del ciclo:

\begin{equation}
\label{rm:eq:gravity-compton-clock}
\bar\lambda_{C,a}
=\frac{\hbar_a}{m_{108,a}c_{\rm int}}
=R_{108}.
\end{equation}

\begin{theorem}[Unicidad del selector bajo la condici\'on de periodicidad HMT]
\label{rm:thm:gravity-circular-selector}
Con la convenci\'on reducida
\(r_g=Gm/c_{\rm int}^2\) y la premisa de periodicidad HMT que identifica el radio
gravitatorio del cuanto CT108 con su radio de fase, la condici\'on de cierre
\begin{equation}
\label{rm:eq:gravity-closing-condition}
r_{g,a}(\theta)=R_{108}
\end{equation}
selecciona una \`unica soluci\'on positiva,
\begin{equation}
\label{rm:eq:gravity-theta108}
\boxed{
\theta_{108}^{\rm circ}=\left(\frac{54}{\pi}\right)^2
=295.4525715010569415303525\ldots .}
\end{equation}
En esa realizaci\'on,
\begin{equation}
\label{rm:eq:gravity-four-lengths}
R_{108}=\bar\lambda_{C,a}=r_{g,a}=\ell_P^{\rm circ}.
\end{equation}
\end{theorem}


\begin{proof}
Por \cref{rm:eq:gravity-transducer,rm:eq:gravity-clock-energy-mass},
\[
r_{g,a}(\theta)
=\frac{G_a(\theta)m_{108,a}}{c_{\rm int}^2}
=\theta c_{\rm int}t_0^2\omega_{108}
=\theta\frac{\pi}{54}\ell_0.
\]
Al compararlo con
\(R_{108}=(54/\pi)\ell_0\) se obtiene necesariamente
\(\theta=(54/\pi)^2\). La ecuaci\'on es lineal en \(\theta\) y todos los
factores son positivos, luego la soluci\'on es \`unica.
\end{proof}

La igualdad \(r_g=R_{108}\) no se deduce de la dimensionalidad de
\cref{rm:eq:gravity-transducer}. Es la condición geométrica explícita de esta
realizaci\'on HMT: acci\'on, Compton y retorno peri\'odico deben determinar una misma
longitud. El teorema demuestra que, aceptada esa condición estructural, el
carácter \(\theta\) queda fijado sin usar el valor SI de \(G\).

\begin{remark}[Convenci\'on de radio]
El teorema usa el radio gravitatorio reducido \(Gm/c^2\). Si se sustituye
por el radio de Schwarzschild \(2Gm/c^2\), el selector cambia por un factor
dos. Las dos expresiones corresponden a condiciones geom\'etricas distintas,
por lo que la convenci\'on debe declararse antes del c\'alculo.
\end{remark}

\subsection{Sistema de Planck y covariancia bidireccional}

\begin{corollary}[Familia de Planck asociada a la estructura cronol\'ogica CT108]
\label{rm:cor:gravity-planck-family}
Para \(G_a=G_a(\theta_{108}^{\rm circ})\),
\begin{align}
\ell_P&=\frac{54}{\pi}\ell_0,
&t_P&=\frac{54}{\pi}t_0,
\label{rm:eq:gravity-planck-length-time}\\
E_{P,a}&=\frac{\hbar_a}{t_P}
=\frac{\pi\hbar_a}{54t_0}
=\frac{h_a}{108t_0},
&\ell_P^2&=\frac{G_a\hbar_a}{c_{\rm int}^3}
=\theta_{108}^{\rm circ}\ell_0^2.
\label{rm:eq:gravity-planck-energy-area}
\end{align}
Adem\'as,
\begin{equation}
\label{rm:eq:gravity-planck-force-power-density}
F_{P,a}=\frac{c_{\rm int}^4}{G_a},
\qquad
P_{P,a}=\frac{c_{\rm int}^5}{G_a},
\qquad
\rho_{P,a}=\frac{\hbar_a}
{(\theta_{108}^{\rm circ})^2c_{\rm int}^5t_0^4}.
\end{equation}
\end{corollary}

\begin{proof}
Las tres primeras identidades proceden de
\(\ell_P=R_{108}\), \(t_P=\ell_P/c_{\rm int}\) y
\(h_a=2\pi\hbar_a\). Para el \`area,
\[
\frac{G_a\hbar_a}{c_{\rm int}^3}
=\theta_{108}^{\rm circ}c_{\rm int}^2t_0^2
=\theta_{108}^{\rm circ}\ell_0^2.
\]
Las expresiones restantes son las composiciones dimensionales de fuerza,
potencia y masa por volumen con la misma secci\'on.
\end{proof}

Sea
\begin{equation}
\label{rm:eq:gravity-Ract}
R_{\rm act}:=\frac{\hbar_{\rm pre}}{\hbar_{\rm ret}}.
\end{equation}
Como la acci\'on aparece en el numerador de la masa cronol\'ogica y en el
denominador de \(G\),
\begin{equation}
\label{rm:eq:gravity-twoway-covariance}
\frac{m_{108,\rm pre}}{m_{108,\rm ret}}=R_{\rm act},
\qquad
\frac{G_{\rm pre}}{G_{\rm ret}}=R_{\rm act}^{-1},
\qquad
G_a\hbar_a
=\theta_{108}^{\rm circ}c_{\rm int}^5t_0^2.
\end{equation}
Por tanto, \(\ell_P^2=G_a\hbar_a/c^3\) no cambia de hoja. La gravedad y
la acci\'on se orientan rec\'iprocamente, mientras el \`area conserva la
geometr\'ia com\'un. Esta cancelaci\'on es la interfaz de confluencia hacia el operador de
\`area del \cref{rm:ch:modular}.

\subsection{Componente icosaédrica de dimensión \(5\)}

El antecedente no es \(\R\) ni una elecci\'on posterior de cinco
polarizaciones. El objeto encargado de transportar la incidencia se define
expl\'icitamente antes de introducir la realizaci\'on excepcional.

Sea \(X_{\rm TPK}^{\rm exc}\) el subespacio del estado \(\TPK\) enriquecido
que conserva, adem\'as de la salida aritm\'etica, la orientaci\'on, la ruta,
el acarreo y la incidencia excepcional. Para explicitar el espacio geom\'etrico de llegada,
sea \(\varphi_g=(1+\sqrt5)/2\), \(\nu=\sqrt{1+\varphi_g^2}\), y definamos
\begin{equation}
\label{rm:eq:gravity-icosahedron-vertices}
\Omega_{12}^{\rm ico}
=\frac1\nu\bigl(
\{0\}\times\{\pm1\}\times\{\pm\varphi_g\}
\;\cup\;
\{\pm1\}\times\{\pm\varphi_g\}\times\{0\}
\;\cup\;
\{\pm\varphi_g\}\times\{0\}\times\{\pm1\}
\bigr)\subset S^2.
\end{equation}
Sus doce elementos son las clases orientadas de v\'ertice de un icosaedro.
La incidencia no se deja verbal: se fija por
\begin{equation}
\label{rm:eq:gravity-icosahedron-incidence}
\mathcal I_{12}^{\rm ico}
=\{(v,w)\in\Omega_{12}^{\rm ico}\times\Omega_{12}^{\rm ico}:
v\neq w,\ \langle v,w\rangle=1/\sqrt5\}.
\end{equation}
Cada v\'ertice tiene cinco vecinos en esta relaci\'on. El puente de lectura
se tipa como una aplicaci\'on
\begin{equation}
\label{rm:eq:gravity-exceptional-reader-map}
\pi_{\rm ico}:X_{\rm TPK}^{\rm exc}\longrightarrow\Omega_{12}^{\rm ico}.
\end{equation}
La cardinalidad doce de la imagen de \(\pi_{\rm ico}\) no es suficiente.
Para definir una realizaci\'on de incidencia, la aplicaci\'on debe satisfacer
simult\'aneamente:
\begin{enumerate}[label=(\roman*)]
\item ser sobreyectiva y constante exactamente en las fibras que descartan
la carta aritm\'etica pero conservan la ruta excepcional;
\item entrelazar la acci\'on de los automorfismos orientados del antecedente
con una acci\'on transitiva sobre \(\Omega_{12}^{\rm ico}\);
\item transportar la incidencia TPK declarada a
\(\mathcal I_{12}^{\rm ico}\), no
s\'olo sus cardinales;
\item conservar la involuci\'on de hoja y el retorno non\'adico en el
estabilizador de la fibra correspondiente.
\end{enumerate}

Estas condiciones separan tres afirmaciones: la existencia del mapa
\(\pi_{\rm ico}\), la identificaci\'on del grupo que act\'ua y la
representaci\'on lineal que se construye despu\'es. La segunda y la tercera
no pueden usarse para definir retrospectivamente la primera.

\begin{theorem}[Descenso de incidencia al cociente icosa\'edrico]
\label{rm:thm:gravity-exceptional-descent}
Sup\'ongase que \(\pi_{\rm ico}\) satisface (i)--(iv) y que la acci\'on
orientada inducida en \(\Omega_{12}^{\rm ico}\) tiene grupo
\(G_{\rm ico}\simeq A_5\). Entonces, para toda clase
\(v_0\in\Omega_{12}^{\rm ico}\),
\begin{equation}
\label{rm:eq:gravity-icosahedral-quotient}
H_{v_0}=\operatorname{Stab}_{G_{\rm ico}}(v_0)\simeq C_5,
\qquad
\Omega_{12}^{\rm ico}\simeq G_{\rm ico}/H_{v_0}\simeq A_5/C_5.
\end{equation}
La aplicaci\'on
\begin{equation}
\label{rm:eq:gravity-coset-map}
\kappa_{v_0}:G_{\rm ico}/H_{v_0}\longrightarrow\Omega_{12}^{\rm ico},
\qquad gH_{v_0}\longmapsto g\cdot v_0,
\end{equation}
es una biyecci\'on equivariante y preserva la incidencia si en el cociente
se declara
\begin{equation}
\label{rm:eq:gravity-coset-incidence}
gH_{v_0}\sim g'H_{v_0}
\quad\Longleftrightarrow\quad
(g\cdot v_0,g'\cdot v_0)\in\mathcal I_{12}^{\rm ico}.
\end{equation}
Por tanto, el compuesto \(\kappa_{v_0}^{-1}\pi_{\rm ico}\) es el mapa
tipado del antecedente TPK al cociente \(A_5/C_5\).
\end{theorem}

\begin{proof}
La lista de \cref{rm:eq:gravity-icosahedron-vertices} tiene doce elementos.
Por transitividad, \(|G_{\rm ico}\cdot v_0|=12\), y el teorema
\`orbita--estabilizador da \(|H_{v_0}|=60/12=5\). Todo grupo de orden
primo cinco es c\'iclico; por tanto, \(H_{v_0}\simeq C_5\). La f\'ormula de
\(\kappa_{v_0}\) no depende del representante: si
\(gH_{v_0}=g'H_{v_0}\), entonces \(g^{-1}g'\in H_{v_0}\) y
\(g'v_0=gv_0\). La sobreyectividad procede de la transitividad, y la
igualdad de cardinales da inyectividad. La equivarianza es inmediata:
\(\kappa_{v_0}(agH_{v_0})=a\kappa_{v_0}(gH_{v_0})\). Finalmente,
\cref{rm:eq:gravity-coset-incidence} es precisamente el transporte de
estructura por esa biyecci\'on; la hip\'otesis (iii) hace conmutar este
transporte con \(\pi_{\rm ico}\).
\end{proof}

El teorema no sustituye la construcci\'on de \(\pi_{\rm ico}\): demuestra
qu\'e se obtiene una vez verificados sus cuatro requisitos. En el perfil
HMT completo, el corredor de incidencia que llega a
Paley--Witt--Golay--Leech--\(V^\natural\)--Monstruo--Moonshine es el
dominio de definici\'on de ese mapa. En este volumen aut\'onomo, la entrada exacta es
el mapa \cref{rm:eq:gravity-exceptional-reader-map} con sus condiciones; todo
el tramo desde \(A_5/C_5\) hasta la reducci\'on TT se prueba a continuaci\'on
sin volver a consultar un valor gravitatorio.

La cadena de procedencia se expresa mediante todos sus morfismos:
\begin{equation}
\label{rm:eq:gravity-exceptional-provenance}
\APP\longrightarrow\TRIT\longrightarrow\TPK
\longrightarrow X_{\rm TPK}^{\rm exc}
\xrightarrow{\ \pi_{\rm ico}\ }\Omega_{12}^{\rm ico}
\xleftarrow[\simeq]{\ \kappa_{v_0}\ }A_5/C_5.
\end{equation}
El corredor excepcional completo se demuestra en su residencia propia. La
hip\'otesis que este cap\'itulo no puede reemplazar por un conteo es la
existencia de \(\pi_{\rm ico}\) con (i)--(iv). Fijada esa entrada, el descenso
al cociente es el \cref{rm:thm:gravity-exceptional-descent}, y se construyen
desde \`el \(P_5\), \(\Gamma_5\) y \(\Lambda(n)\). La realizaci\'on del
portador resultante como campo gravitatorio sigue siendo una interfaz
f\'isica tipada, no una consecuencia de cardinalidad.

Sea \(G=A_5\), sea \(H\simeq C_5\) el estabilizador de un v\'ertice
orientado del icosaedro y consid\'erese
\begin{equation}
\label{rm:eq:gravity-v12}
V_{12}=\R[G/H],
\qquad |G/H|=\frac{60}{5}=12.
\end{equation}
La representaci\'on de permutaci\'on \(\rho_{12}\) se descompone como
\begin{equation}
\label{rm:eq:gravity-v12-decomposition}
V_{12}\simeq\mathbf1\oplus V_3\oplus V_{3'}\oplus V_5.
\end{equation}
El idempotente central del sumando de car\'acter
\(\chi_5=(5,1,-1,0,0)\) es
\begin{equation}
\label{rm:eq:gravity-p5}
\boxed{
P_5=\frac5{60}\sum_{g\in A_5}\chi_5(g^{-1})\rho_{12}(g)
=\frac1{12}\left(
5I+\sum_{g\in2A}\rho_{12}(g)
-\sum_{g\in3A}\rho_{12}(g)
\right).}
\end{equation}

\begin{theorem}[Proyector y acoplamiento gravitatorio de rango cinco]
\label{rm:thm:gravity-p5-coupling}
El operador \(P_5\) satisface
\begin{equation}
\label{rm:eq:gravity-p5-properties}
P_5^2=P_5=P_5^*,
\qquad \Tr P_5=\rank P_5=5.
\end{equation}
Para
\begin{equation}
\label{rm:eq:gravity-G5}
\mathbb G_{5,a}:=G_a(\theta_{108}^{\rm circ})P_5
\end{equation}
se tiene
\begin{equation}
\label{rm:eq:gravity-G5-spectrum}
\Spec(\mathbb G_{5,a})
=\{G_a(\theta_{108}^{\rm circ})\text{ con multiplicidad }5,
\;0\text{ con multiplicidad }7\}.
\end{equation}
\end{theorem}

\begin{proof}
La ortogonalidad de caracteres convierte \cref{rm:eq:gravity-p5} en el
idempotente central que act\'ua como la identidad sobre la \`unica copia de
\(V_5\) y como cero en los otros tres sumandos de
\cref{rm:eq:gravity-v12-decomposition}. La realidad del car\'acter y la
ortogonalidad de \(\rho_{12}\) dan la autoadjunci\'on. Su traza y su rango
son cinco. Multiplicar un proyector de rango cinco por el escalar positivo
\(G_a\) produce inmediatamente el espectro indicado.
\end{proof}

\begingroup\fontsize{10.2}{11.7}\selectfont
\setlength{\parskip}{2pt}\setlength{\abovedisplayskip}{5pt}\setlength{\belowdisplayskip}{5pt}
La expresi\'on \(\mathbb G_{5,a}\) no a\~nade cinco valores diferentes de
la constante gravitatoria. Es una misma secci\'on de \(G\) actuando sobre
el sumando tensorial seleccionado, y anulando los otros siete modos del
m\'odulo de permutaci\'on.
 
\paragraph{Reducción equivariante del portador al sector tensorial sin traza.} El portador se entrelaza con los tensores simétricos sin traza y se reduce de forma equivariante antes de imponer transversalidad.

\par\endgroup
 \endgroup

\chapter{Entrelazador tensorial y reducción transversal \texorpdfstring{\(12\to5\to5\to2\)}{12->5->5->2}}
\begingroup
\renewcommand{\chapter}[2][]{}%
\chapter{Entrelazador con tensores simétricos sin traza y reducción transversal \texorpdfstring{\(12\to5\to5\to2\)}{12→5→5→2}}
\label{chap:parte-iv-64}

El marco icosaédrico construye una isometría equivariante hacia el espacio de tensores simétricos sin traza. La proyección transversal realiza \(12\to5\to5\to2\), conserva el portador pentacomponente y localiza el subespacio radiativo de dimensión \(2\).

\section{Entrelazador con el espacio de tensores sim\'etricos sin traza}

Sea \(\mathcal V\subset S^2\) el conjunto de los doce v\'ertices unitarios
del icosaedro. Para cada \(v\in\mathcal V\), def\'inase
\begin{equation}
\label{rm:eq:gravity-Qv}
Q_v=vv^T-\frac13I_3\in
\operatorname{STF}_3:=\operatorname{Sym}^2_0(\R^3).
\end{equation}
El mapa de marco
\begin{equation}
\label{rm:eq:gravity-Gamma0}
\Gamma_0\!\left(\sum_{v\in\mathcal V}x_ve_v\right)
=\sum_{v\in\mathcal V}x_vQ_v
\end{equation}
es \(A_5\)-equivariante y satisface
\begin{equation}
\label{rm:eq:gravity-tight-frame}
\Gamma_0\Gamma_0^*=\frac85I_{\operatorname{STF}_3},
\qquad
\Gamma_0=\Gamma_0P_5.
\end{equation}

\begin{theorem}[Entrelazador tensorial normalizado]
\label{rm:thm:gravity-gamma5}
La aplicaci\'on
\begin{equation}
\label{rm:eq:gravity-gamma5}
\boxed{
\Gamma_5:=\sqrt{\frac58}\,\Gamma_0|_{V_5}:
V_5\xrightarrow{\ \simeq\ }\operatorname{STF}_3}
\end{equation}
es una isometr\'ia \(A_5\)-equivariante. Su adjunto es
\begin{equation}
\label{rm:eq:gravity-gamma5-adjoint}
\Gamma_5^*Q=\sqrt{\frac58}
\sum_{v\in\mathcal V}(v^TQv)e_v.
\end{equation}
\end{theorem}

\begin{proof}
La equivarianza sigue de \(Q_{gv}=gQ_vg^T\). El operador
\(\Gamma_0\Gamma_0^*\) conmuta con la representaci\'on irreducible de
\(A_5\) sobre \(\operatorname{STF}_3\), por lo que es escalar. Su traza
es
\[
\sum_{v\in\mathcal V}\|Q_v\|_F^2
=12\cdot\frac23=8;
\]
al dividir por la dimensi\'on cinco se obtiene \(8/5\). De aqu\'i,
\(\Gamma_5\Gamma_5^*=I\). Como dominio y codominio tienen dimensi\'on
cinco, tambi\'en \(\Gamma_5^*\Gamma_5=I\). La f\'ormula del adjunto
procede de \(\langle Q,Q_v\rangle_F=v^TQv\).
\end{proof}

La normalizaci\'on mediante el marco icosa\'edrico elimina una coincidencia
puramente cardinal: construye el mapa que transporta el sumando finito al
portador continuo de esp\'in dos.

\section{Proyecci\'on transversal sin traza}

Para \(n\in S^2\), sea \(\Pi(n)=I-nn^T\). Definimos
\begin{equation}
\label{rm:eq:gravity-lambda-tt}
\boxed{
\Lambda(n)Q
=\Pi Q\Pi-\frac12\Tr(\Pi Q\Pi)\Pi.}
\end{equation}

\begin{theorem}[Reducci\'on de cinco componentes a dos]
\label{rm:thm:gravity-tt}
La restricci\'on de \(\Lambda(n)\) a \(\operatorname{STF}_3\) es el
proyector ortogonal sobre
\begin{equation}
\label{rm:eq:gravity-htt}
\mathcal H_{\rm TT}(n)
=\{Q\in\operatorname{STF}_3:Qn=0\}.
\end{equation}
En particular,
\begin{equation}
\label{rm:eq:gravity-lambda-properties}
\Lambda(n)^2=\Lambda(n)=\Lambda(n)^*,
\qquad \rank\Lambda(n)=2.
\end{equation}
Para \(n=e_3\), la imagen est\'a generada por
\begin{equation}
\label{rm:eq:gravity-plus-cross}
E_+=\frac1{\sqrt2}\operatorname{diag}(1,-1,0),
\qquad
E_\times=\frac1{\sqrt2}
\begin{pmatrix}0&1&0\\1&0&0\\0&0&0\end{pmatrix}.
\end{equation}
\end{theorem}

\begin{proof}
Como \(\Pi n=0\), la salida es transversal. Su traza vale
\[
\Tr(\Pi Q\Pi)-\frac12\Tr(\Pi)\Tr(\Pi Q\Pi)=0,
\]
pues \(\Tr\Pi=2\). Si \(Qn=0\) y \(\Tr Q=0\), entonces
\(\Pi Q\Pi=Q\) y \(\Tr(\Pi Q\Pi)=0\), por lo que
\(\Lambda(n)Q=Q\). La f\'ormula es autoadjunta para el producto de
Frobenius. En el marco \(n=e_3\) su matriz en una base STF adaptada es
\(\operatorname{diag}(1,1,0,0,0)\), lo que prueba el rango.
\end{proof}

La cadena completa queda, por tanto,
\begin{equation}
\label{rm:eq:gravity-12-5-5-2}
\boxed{
V_{12}\xrightarrow{P_5}V_5
\xrightarrow{\Gamma_5}\operatorname{STF}_3
\xrightarrow{\Lambda(n)}\mathcal H_{\rm TT}(n),
\qquad 12\longrightarrow5\longrightarrow5\longrightarrow2.}
\end{equation}
El acoplamiento reducido
\begin{equation}
\label{rm:eq:gravity-effective-map}
\mathbb G_{{\rm TT},a}(n)
=G_a(\theta_{108}^{\rm circ})\Lambda(n)\Gamma_5P_5
\end{equation}
tiene rango dos.
 
\paragraph{Descomposición helicoidal del espectro y separación de los modos masivo y no masivos.} La reducción permite distinguir el espectro helicoidal completo de una realización masiva y sus dos helicidades sin masa.
 \endgroup

\chapter{Descomposición helicoidal del portador de espín dos}
\begingroup
\renewcommand{\chapter}[2][]{}%
\chapter{Descomposición helicoidal del portador de espín \texorpdfstring{\(2\)}{2} y distinción entre realizaciones masiva y sin masa}
\label{chap:parte-iv-65}

El portador real de dimensión \(5\) se descompone en \(5\) pesos helicoidales. La realización masiva conserva las \(5\) componentes; la reducción transversal sin masa conserva únicamente las helicidades radiativas \(\pm2\), sin convertir una posibilidad colectiva en partícula fundamental postulada.

\section{Descomposición en \(5\) modos de polarización}

La dimensi\'on cinco adquiere contenido tensorial al elegir una direcci\'on
unitaria de propagaci\'on \(n\) y un marco ortonormal orientado
\((e_1,e_2,n)\). Definimos los cinco tensores reales
\begin{align}
E_{+}&=\frac1{\sqrt2}(e_1e_1^T-e_2e_2^T),
&E_{\times}&=\frac1{\sqrt2}(e_1e_2^T+e_2e_1^T),
\label{rm:eq:gravity-helicity2-real}\\
E_{1}&=\frac1{\sqrt2}(e_1n^T+ne_1^T),
&E_{2}&=\frac1{\sqrt2}(e_2n^T+ne_2^T),
\label{rm:eq:gravity-helicity1-real}\\
E_{0}&=\frac1{\sqrt6}(2nn^T-e_1e_1^T-e_2e_2^T).
\label{rm:eq:gravity-helicity0-real}
\end{align}
Todos son sim\'etricos y sin traza. Los dos primeros son transversales; los
dos siguientes mezclan el plano transversal con la direcci\'on de
propagaci\'on, y el \`ultimo es el modo escalar longitudinal sin traza.

\begin{theorem}[Descomposici\'on helicoidal completa del portador]
\label{rm:thm:gravity-five-helicities}
Los tensores de
\cref{rm:eq:gravity-helicity2-real,rm:eq:gravity-helicity1-real,rm:eq:gravity-helicity0-real}
forman una base ortonormal de \(\operatorname{STF}_3\). En la
complejificaci\'on, las combinaciones
\begin{equation}
\label{rm:eq:gravity-complex-helicities}
E_{\pm2}=\frac1{\sqrt2}(E_+\mp\ii E_\times),
\qquad
E_{\pm1}=\frac1{\sqrt2}(E_1\mp\ii E_2),
\qquad E_0=E_0
\end{equation}
son vectores de peso \(\pm2,\pm1,0\), respectivamente, bajo rotaciones
alrededor de \(n\). Adem\'as,
\begin{equation}
\label{rm:eq:gravity-tt-on-five}
\Lambda(n)E_{\pm2}=E_{\pm2},
\qquad
\Lambda(n)E_{\pm1}=0,
\qquad
\Lambda(n)E_0=0.
\end{equation}
\end{theorem}

\begin{proof}
En el marco \((e_1,e_2,n)\), cada tensor tiene norma de Frobenius uno.
Sus patrones de entradas no diagonales son disjuntos, y las tres diagonales
de \(E_+,E_0\) tienen producto interno cero; por tanto, los cinco tensores
son ortonormales. Como \(\dim\operatorname{STF}_3=6-1=5\), forman una base.

Una rotaci\'on de \`angulo \(\psi\) alrededor de \(n\) env\'ia
\((e_1,e_2)\) a
\((\cos\psi\,e_1+\sin\psi\,e_2,
-\sin\psi\,e_1+\cos\psi\,e_2)\). Sustituyendo en las f\'ormulas se obtiene
una rotaci\'on de \`angulo \(2\psi\) sobre
\(\operatorname{span}\{E_+,E_\times\}\), una de \`angulo \(\psi\) sobre
\(\operatorname{span}\{E_1,E_2\}\), y \(E_0\) fijo. De ah\'i los pesos de
\cref{rm:eq:gravity-complex-helicities}. Finalmente,
\(\Pi E_{+}\Pi=E_+\) y \(\Pi E_\times\Pi=E_\times\), mientras
\(\Pi E_1\Pi=\Pi E_2\Pi=0\). Para \(E_0\),
\(\Pi E_0\Pi=-(e_1e_1^T+e_2e_2^T)/\sqrt6\), que es puramente traza en
el plano, y la sustracci\'on de \cref{rm:eq:gravity-lambda-tt} la anula.
\end{proof}

El teorema precisa el significado de ``cinco polarizaciones'' sin atribuir
a todas el mismo estatuto din\'amico. Antes de imponer ecuaciones de campo,
el portador de esp\'in dos tiene cinco coordenadas irreducibles. En una
realizaci\'on masiva de Fierz--Pauli, las restricciones dejan precisamente
los cinco pesos de \cref{rm:eq:gravity-complex-helicities}; en una realizaci\'on
sin masa con invariancia de calibre, los sectores \(\pm1\) y \(0\) se eliminan y
quedan \(\pm2\). HMT proporciona aqu\'i el portador, la base y el proyector;
la elecci\'on de la din\'amica masiva o sin masa pertenece a la interfaz
f\'isica. Por eso cinco no significa ``cinco ondas observadas'', sino cinco
grados tensoriales disponibles antes de la reducci\'on.

La orientaci\'on tambi\'en queda visible: cambiar
\((e_1,e_2,n)\) por \((e_1,-e_2,-n)\) conserva \(E_+,E_0\) e intercambia
las fases de las parejas helicoidales. Esa informaci\'on no est\'a en el
escalar \(G_a\); reside en el portador sobre el que act\'ua
\(\mathbb G_{5,a}\).

\section{Cadena representacional \(12\to5\to5\to2\): portador \(V_5\), tensor sin traza, polarizaciones de Fierz–Pauli y helicidades transversales}

La dimensi\'on cinco aparece en objetos distintos, y s\'olo
\cref{rm:eq:gravity-12-5-5-2} autoriza el transporte entre ellos:

\begin{enumerate}
\item \(V_5\) es el sumando irreducible de la representaci\'on sobre los
doce v\'ertices. Es un resultado finito de teor\'ia de representaciones.
\item \(\operatorname{STF}_3\) es el portador real de dimensi\'on cinco de
los pesos \(m=-2,-1,0,1,2\) de esp\'in dos.
\item Una teor\'ia masiva de Fierz--Pauli puede realizar esos cinco pesos
como cinco estados f\'isicos, despu\'es de fijar acci\'on y restricciones.
\item La relatividad general sin masa conserva dos helicidades radiativas,
\(\pm2\), representadas por \(E_+\) y \(E_\times\).
\item Las clases \(E(2)\) de ondas m\'etricas, como \(III_5\), clasifican
posibles amplitudes bajo otras hip\'otesis; no son autom\'aticamente
\(V_5\), Fierz--Pauli o relatividad general.
\end{enumerate}

En particular, la frase correcta es: \emph{HMT construye un portador
tensorial de cinco componentes antes de las restricciones; la reducci\'on
TT de la realizaci\'on sin masa posee rango dos}. No se afirma que la
relatividad general tenga cinco polarizaciones propagantes.
 
\paragraph{Condiciones de existencia del modo mediante selección holonómica natural.} La existencia física del modo requiere todavía un selector holonómico natural y falsable.
 \endgroup

\chapter{Selector gravitatorio natural y transductor dimensional}
\section[Criterios estructurales y carta metrológica de \(G\)]{Criterios de
selección de la coordenada gravitatoria: holonomía enriquecida, respuesta
espectral y carta metrológica}
\label{sec:selector-gravitatorio-rev11}
\index{constante gravitatoria!transductor y selector}
\index{gravedad!coordenada metrológica}
\index{holonomía!selector gravitatorio}

La acción Palatini--Cartan--Holst y su Hamiltoniano restringido fijan primero
la dinámica. Sobre ella, la constante gravitatoria se organiza en tres niveles
tipados: la recta dimensional que alberga \(G\), el transductor que convierte
una coordenada adimensional en una sección de esa recta y el selector
estructural de dicha coordenada. La carta decimal y la incertidumbre
experimental se presentan al final como representación y contraste.

\subsection{Recta dimensional y transductor interno}

En la retícula dimensional de la Mecánica Dimensional,
\[
 \mathbf G=-\mathbf S+5\mathbf C+2\mathbf T,
 \qquad [G]=L^3M^{-1}T^{-2}.
\]
Sean
\[
 c_{\rm int}\in\mathcal L_C^+,
 \qquad t_0\in\mathcal L_T^+,
 \qquad \ell_0=c_{\rm int}t_0,
 \qquad \hbar_{\rm int}\in\mathcal L_S^+.
\]

\HMTClaim{MD-R-G-TRANSDUCER}
\begin{definition}[Transductor gravitatorio]
\label{def:transductor-gravitatorio-rev11}
Para \(\theta_G>0\), se define
\begin{equation}
 \boxed{
 G_{\rm int}(\theta_G)
 =\theta_G\frac{c_{\rm int}^{5}t_0^2}{\hbar_{\rm int}}
 =\theta_G\frac{c_{\rm int}^{3}\ell_0^2}{\hbar_{\rm int}}
 \in\mathcal L_G^+.}
 \label{eq:transductor-gravitatorio-rev11}
\end{equation}
\end{definition}

La igualdad de las dos expresiones y el tipo de
\eqref{eq:transductor-gravitatorio-rev11} son exactos. Si
\[
 L_G^2:=\frac{\hbar_{\rm int}G_{\rm int}}{c_{\rm int}^3},
\]
entonces
\begin{equation}
 \theta_G=\left(\frac{L_G}{\ell_0}\right)^2.
 \label{eq:thetaG-razon-longitudes-rev11}
\end{equation}
La coordenada gravitatoria es, por tanto, una razón cuadrática de longitudes
o una respuesta de holonomía normalizada. La década
\(10^{-34}\mathcal U_S\) ya pertenece a \(\hbar_{\rm int}\): fija la sección
absoluta de acción y no constituye un factor relativo añadido a \(G\).

\subsection{Torsión observable y caracteres positivos}

El retorno observable de \(27\) pasos
\[
 \mathcal K_{27}:=F_{\rm obs}^{27}
\]
satisface
\[
 \mathcal K_{27}^{4}=F_{\rm obs}^{108}=I.
\]
En la coordenada \(\theta\in\mathbb Z/108\mathbb Z\), el orden es
exactamente cuatro, pues
\(\mathcal K_{27}(\theta)=\theta+27\),
\(\mathcal K_{27}^{2}(\theta)=\theta+54\neq\theta\) y
\(\mathcal K_{27}^{4}(\theta)=\theta\).

\HMTClaim{MD-R-G-POSITIVE-CHARACTER-NOGO}
\begin{theorem}[Trivialidad del carácter positivo observable]
\label{thm:caracter-positivo-G-rev11}
Todo homomorfismo multiplicativo
\[
 \chi:\langle\mathcal K_{27}\rangle
 \longrightarrow\mathbb R_{>0}^{\times}
\]
cumple \(\chi(\mathcal K_{27})=1\).
\end{theorem}

\begin{proof}
La multiplicatividad implica
\(\chi(\mathcal K_{27})^4=\chi(I)=1\).
El único real positivo cuya cuarta potencia es uno es \(1\).
\end{proof}

El mismo ciclo sí transporta los caracteres unitarios
\[
 \chi_j(\mathcal K_{27})=\exp(2\pi ij/4),
 \qquad j=0,1,2,3.
\]
El cociente observable admite fase, mientras que su torsión elimina una
escala positiva multiplicativa no trivial. Una continuación multiplicativa debe, por ello,
actuar sobre un levantamiento \(\widetilde{\mathcal K}_{27}\) cuya clase en
la abelianización del grupo de isotropía sea no torsional, y debe seleccionar
de manera natural tanto el carácter como su normalización.

\subsection{Memoria cronológica y respuesta espectral}
\label{sec:respuesta-espectral-G-rev11}

La cobertura cronológica dirigida realiza el retorno de \(27\) pasos como la
traslación
\[
 (G_{\rm or},S)\longmapsto(G_{\rm or},S)+(1,18),
\]
de orden semigrupal infinito. Esta vía conserva acumulación, pero su
respuesta espectral requiere interferencia entre historias.

Sea \(\Gamma=(V,E)\) un multigrafo dirigido finito con medida positiva,
pesos \(a_\ast(e)\geq0\) y cociclo \(c(e)\in\mathbb Z^2\). El operador
torcido
\[
 (\mathcal L_{\beta,\vartheta}f)(y)
 =\sum_{e:x\to y}
 e^{-\beta a_\ast(e)}e^{i\langle\vartheta,c(e)\rangle}f(x)
\]
induce el operador positivo
\(\mathcal H_{\beta,\vartheta}
=\mathcal L_{\beta,\vartheta}^{*}\mathcal L_{\beta,\vartheta}\).

\HMTClaim{MD-R-G-PHASE-INTERFERENCE}
\begin{proposition}[Criterio de interferencia para la respuesta de fase]
\label{prop:interferencia-selector-G-rev11}
En una órbita determinista con una única arista entrante por vértice,
\(\mathcal H_{\beta,\vartheta}\) es independiente de \(\vartheta\). Toda
respuesta espectral de fase distinta de cero requiere historias coherentes
múltiples o un diferencial magnético cuya holonomía sobreviva en los ciclos
enriquecidos.
\end{proposition}

\begin{proof}
En la órbita determinista, la única fase de cada término se multiplica por su
conjugada en
\(\mathcal L_{\beta,\vartheta}^{*}\mathcal L_{\beta,\vartheta}\)
y se cancela. La ausencia de términos cruzados elimina toda dependencia de
\(\vartheta\).
\end{proof}

Si una banda simple \(\lambda_0(\vartheta)\) posee un mínimo en el origen,
su Hessiano \(Q\) es semidefinido positivo. Con
\[
 v_{27}=(1,18),
 \qquad
 \mathcal R_{27}=v_{27}^{\mathsf T}Qv_{27},
\]
la condición \(\mathcal R_{27}\neq0\) constituye el control mínimo de
sensibilidad a la memoria acumulada. Una normalización natural
\(\mathfrak N\) permitiría definir
\[
 \Lambda_G=\ell_0\mathfrak N(Q,v_{27}),
 \qquad
 G_{\rm spec}
 =\frac{c_{\rm int}^{3}}{\hbar_{\rm int}}
  \bigl(\delta_\theta\Lambda_G\bigr)^2,
 \qquad
 \delta_\theta=\frac{A_{\rm deg}-C_{\rm deg}^{\ast}}{360}.
\]
La homogeneidad de esta plantilla es exacta. Su realización física queda
caracterizada por cuatro condiciones positivas: multigrafo enriquecido,
respuesta no nula, estabilidad por refinamiento e independencia de reescala,
junto con una normalización \(\mathfrak N\) fijada antes del contraste
metrológico.

\subsection{Funcional geométrico sobre la holonomía enriquecida}

La segunda continuación define directamente un funcional de longitud
\[
 \mathscr L:
 \operatorname{Loop}(\mathcal X_{\rm TPK}^{\rm enr})
 \longrightarrow\mathcal L_L^+\cup\{0\},
\]
invariante por cambio cíclico de punto base, conjugación, reetiquetado
admisible y subdivisión. Para un levantamiento enriquecido se escriben
\[
 \Lambda_{27}:=\mathscr L(\widetilde{\mathcal K}_{27}),
 \qquad
 \rho_{27}:=\frac{\Lambda_{27}}{c_{\rm int}t_0},
 \qquad
 \theta_G^{\rm hol}:=(\delta_\theta\rho_{27})^2.
\]
La transducción produce la plantilla homogénea
\begin{equation}
 \boxed{
 G_{\rm hol}
 =\frac{c_{\rm int}^{3}}{\hbar_{\rm int}}
  \bigl(\delta_\theta\Lambda_{27}\bigr)^2.}
 \label{eq:plantilla-holonomia-G-rev11}
\end{equation}
La ley de concatenación de este funcional sobre el ciclo torsional forma
parte de su definición y de su prueba. La selección física se completa
construyendo \(\widetilde{\mathcal K}_{27}\) y \(\mathscr L\), fijando una
normalización única y verificando la estabilidad de la salida. El valor
\(G_{\rm CODATA\,2022}\) interviene exclusivamente después, en el contraste
metrológico.

Estas condiciones describen esta plantilla de holonomía. La realización
longitudinal y radial ya especificada en
\cref{g26:sec:rigidez-radial-es} utiliza el registro marcado, su
expectativa condicional y la norma positiva; su normalización y el valor
común de \(G\) están demostrados allí. No se identifica
\(\delta_\theta\Lambda_{27}\) con \(L\) por semejanza formal:
la plantilla alternativa conserva la obligación de construir su
funcional antes de efectuar esa identificación.

\begin{proposition}[Separación de los dos selectores gravitatorios]
\label{prop:separacion-selectores-G}
El selector circular
\(\theta_{108}^{\rm circ}=(54/\pi)^2\), construido por el cierre
acción--Compton--gravedad, y el selector holonómico
\(\theta_G^{\rm hol}=(\delta_\theta\rho_{27})^2\) son dos realizaciones
tipadas distintas del transductor \(G_{\rm int}\).  No se identifican entre
sí y la carta decimal \(\mathscr R_G\) no constituye una evaluación de
\(G_{\rm hol}\).
\end{proposition}
\begin{proof}
El primer selector depende únicamente del cierre circular de orden 108. El
segundo contiene, además, la longitud espectral
\(\Lambda_{27}=\mathscr L(\widetilde{\mathcal K}_{27})\).  Una igualdad
entre ambos requeriría calcular \(\Lambda_{27}\) y probar la coincidencia de
normalizaciones; ninguna operación del lector decimal realiza ese cálculo.
Por tanto, sus dominios y datos de entrada los distinguen antes de todo
contraste numérico.
\end{proof}

\subsection{Criterio de selección estructural}

\HMTClaim{MD-R-G-SELECTOR-CRITERION}
\begin{criterion}[Selector estructural de la constante gravitatoria]
\label{crit:selector-G-rev11}
Un selector HMT de \(G\) queda identificado cuando satisface conjuntamente:
\begin{enumerate}
 \item dominio enriquecido, composición y orientación declarados;
 \item descenso a una clase abeliana no torsional, o funcional geométrico
 invariante por punto base, conjugación, reetiquetado y subdivisión;
 \item respuesta positiva no trivial y estable bajo refinamiento;
 \item normalización única e independiente de la escala global del operador;
 \item selección interna de la coordenada y de su posición decimal;
 \item congelación del código y de las entradas antes de abrir el valor de
 contraste;
 \item realización final como sección de \(\mathcal L_G\) en una carta
 dimensional común.
\end{enumerate}
\end{criterion}

Las razones adimensionales
\[
 \frac{\mathcal G_{\rm tor}}{\mathcal I_{\rm tor}}=1,
 \qquad
 \frac{\mathcal G_{\rm av}}{\mathcal I_{\rm av}}
 =\frac{\pi}{\varphi^2}
\]
describen lectores gravitatorio--inerciales distintos y permanecen
compatibles con una sola sección dimensional \(G\). La construcción de esa
sección utiliza además
\((\hbar_{\rm int},c_{\rm int},\ell_0)\) y el selector \(\theta_G\).

\subsection{Evaluación de la carta decimal gravitatoria y contraste metrológico posterior}

La representación decimal conservada en el corpus es
\begin{equation}
 \mathscr R_G([n,k,d]_3;t,e)
 :=10^{-n}\left(k+\frac{t}{10^3}+\frac{e}{10^d}\right).
 \label{eq:lector-decimal-G-rev11}
\end{equation}
Para la carta
\([n,k,d]_3=[11,6,5]_3\), \(t=674\) y \(e=30\), la evaluación racional es
\begin{equation}
 \boxed{
 \mathscr R_G
 =10^{-11}\left(6+\frac{674}{1000}+\frac{30}{10^5}\right)
 =6.67430\times10^{-11}.}
 \label{eq:evaluacion-decimal-G-rev11}
\end{equation}
La terna, los dos bloques y la posición decimal constituyen los datos
declarados de esta carta. Su evaluación es exacta; su selección estructural
queda regida por el \cref{crit:selector-G-rev11}.
Esta descomposición se conserva como testigo decimal histórico. La rama
radial de \eqref{g26:eq:g-rigido-es} obtiene su coordenada mediante
la longitud y la acción, y da
\eqref{cr28:eq:g-evaluation}; no utiliza los bloques \(674\) y \(30\)
para definir el acoplamiento.

Tras elegir la base SI de \(\mathcal L_G\), el contraste vigente es
\begin{equation}
 \boxed{
 G_{\rm CODATA\,2022}=6.67430(15)\times10^{-11}
 \ \mathrm{m^3\,kg^{-1}\,s^{-2}}.}
 \label{eq:G-CODATA-2022-rev11}
\end{equation}

\begin{center}
\begin{tabularx}{0.94\textwidth}{@{}L{0.27\textwidth}Y@{}}
\toprule
Campo & Declaración metrológica\\
\midrule
Valor central &
\(6.67430\times10^{-11}\ \mathrm{m^3\,kg^{-1}\,s^{-2}}\)\\
Incertidumbre estándar &
\(u(G)=0.00015\times10^{-11}
=1.5\times10^{-15}\ \mathrm{m^3\,kg^{-1}\,s^{-2}}\)\\
Incertidumbre relativa & \(u_r(G)\simeq2.2\times10^{-5}\)\\
Fuente y temporalidad &
Valores recomendados CODATA 2022, publicados en 2024
\cite{CODATA2022}\\
Función probatoria &
Contraste externo posterior; la unidad y la incertidumbre no intervienen en
la evaluación racional de \eqref{eq:evaluacion-decimal-G-rev11}.\\
\bottomrule
\end{tabularx}
\end{center}

\paragraph{Reproducibilidad del lector gravitatorio.}
Una enumeración independiente comprueba conjuntamente el tipo dimensional, el
lector decimal y las identidades de monodromía sobre el dominio finito
declarado.

\begin{proposition}[Cadena causal de la constante gravitatoria]
La acción covariante, el Hamiltoniano restringido y la ecuación
Einstein--Schrödinger gobiernan la dinámica gravitatoria. El tipo dimensional
y el transductor \eqref{eq:transductor-gravitatorio-rev11} son exactos; la
evaluación \eqref{eq:evaluacion-decimal-G-rev11} es exacta para su carta
declarada. El valor, la unidad y la incertidumbre de
\eqref{eq:G-CODATA-2022-rev11} constituyen contraste externo CODATA 2022. La
selección autónoma de \(\theta_G\) queda tipada por el levantamiento
enriquecido, la respuesta interferente y el funcional de holonomía sometidos
al \cref{crit:selector-G-rev11}.
\end{proposition}
En la realización radial de reglas R1--R8, la selección está determinada
por \eqref{g26:eq:longitud-unica-es} y \eqref{g26:eq:g-rigido-es};
la compatibilidad con el calendario se demuestra en
\eqref{cr28:eq:circular-radial-agreement}. La plantilla de holonomía
conserva, separadamente, las condiciones de su propio dominio.
 % Integración por delta de la adenda de realizaciones físicas.
% Residencia de procedencia: tipo dimensional de G, selector de holonomía y carta decimal.

\chapter{Transducción dimensional y selección holonómica de la constante
gravitatoria}
\label{chap:transductor-gravitatorio-G}
\index{constante gravitatoria!transductor dimensional}
\index{naturalidad!selector gravitatorio}
\index{lector decimal!control de G}

\section{Recta gravitatoria y forma dimensional}
\label{sec:recta-gravitatoria-interna}

En la retícula del \cref{chap:funtor-dimensional-hmt}, la dimensión de la
constante gravitatoria es
\[
 \mathbf G=-\mathbf S+5\mathbf C+2\mathbf T,
 \qquad
 [G]=L^3M^{-1}T^{-2}.
\]
Sean
\[
 c_{\rm int}\in\mathcal L_C^+,\qquad
 t_0\in\mathcal L_T^+,\qquad
 \ell_0:=c_{\rm int}t_0\in\mathcal L_L^+,
 \qquad
 \hbar_{\rm int}\in\mathcal L_S^+.
\]
La sección \(\hbar_{\rm int}\) se obtiene sólo después del lector
determinantal de década \(-34\) y de la elección de una base de la recta de
acción. No es el carácter angular adimensional anterior al retorno.

Para todo carácter adimensional \(\theta_G>0\), definimos
\[
 \boxed{
 G_{\rm int}(\theta_G)
 =\theta_G
 \frac{c_{\rm int}^{5}t_0^2}{\hbar_{\rm int}}
 =\theta_G
 \frac{c_{\rm int}^{3}\ell_0^2}{\hbar_{\rm int}}
 \in\mathcal L_G^+.}
\]
La igualdad de las dos expresiones y su tipo son
\textsc{exacto interno} dentro de la retícula dimensional. No seleccionan
\(\theta_G\). Si
\[
 L_G^2:=\frac{\hbar_{\rm int}G_{\rm int}}{c_{\rm int}^3},
\]
entonces
\[
 \theta_G=\left(\frac{L_G}{\ell_0}\right)^2.
\]
Esta última identidad localiza el dato selector: la razón de longitudes debe
ser producida por una estructura de holonomía o de respuesta espectral
anterior a la carta metrológica.
La realización del \cref{g26:sec:rigidez-radial-es} satisface esta
exigencia mediante la inscripción y la norma radial, con
\(L_G=L\). Lo que sigue examina otras plantillas de retorno y
holonomía, conservando las obstrucciones propias de sus cocientes.

\section{Obstrucción del carácter positivo sobre el retorno observable}
\label{sec:selector-holonomia-G}

El \cref{sec:jerarquia-retornos-27-54-108} define
\[
 \mathcal K_{27}:=F_{\rm obs}^{27}.
\]
La proposición primaria \cref{prop:retorno-fobs-identidad} demuestra
\[
 F_{\rm obs}^{108}=I;
\]
por tanto, en el cociente observable,
\[
 \mathcal K_{27}^{4}=I.
\]

\begin{theorem}[Obstrucción torsional para el selector positivo observable]
\label{thm:no-go-chi-G-observable}
Todo homomorfismo multiplicativo
\[
 \chi:\langle\mathcal K_{27}\rangle
 \longrightarrow\mathbb R_{>0}^{\times}
\]
es trivial. En particular,
\[
 \chi(\mathcal K_{27})=1.
\]
\end{theorem}

\begin{proof}
La multiplicatividad y \(\mathcal K_{27}^{4}=I\) implican
\[
 \chi(\mathcal K_{27})^4=\chi(I)=1.
\]
El único elemento positivo real cuya cuarta potencia es uno es \(1\).
\end{proof}

El orden observable es exactamente cuatro, y no sólo un divisor de cuatro:
en la coordenada \(\theta\in\mathbb Z/108\mathbb Z\),
\[
 \mathcal K_{27}(\theta)=\theta+27,
 \qquad
 \mathcal K_{27}^{2}(\theta)=\theta+54\neq\theta,
 \qquad
 \mathcal K_{27}^{4}(\theta)=\theta.
\]
El mismo grupo cíclico admite caracteres unitarios
\[
 \chi_k(\mathcal K_{27})
 =\exp\!\left(\frac{2\pi i k}{4}\right),
 \qquad k=0,1,2,3,
\]
y, en particular, los valores de fase \(\pm i\). La conclusión tipada es
que el ciclo observable puede transportar fase, pero no una escala positiva
no trivial. Estos caracteres de fase no se identifican con el operador
\(J_W\) de Paley: satisfacen una relación cuadrática análoga sobre dominios
distintos.

La obstrucción no depende de una comparación metrológica. Afecta a la forma
misma del selector: un carácter positivo sobre el ciclo observable no puede
producir un \(\theta_G\) no trivial. Además, un mero conjunto de clases de
conjugación no constituye, en general, un grupo y no puede declararse dominio
de un carácter sin definir antes su composición.

\subsection{Única continuación multiplicativa admisible}

Una ruta multiplicativa sólo puede reabrirse sobre un objeto enriquecido que
no herede la torsión observable. Para un objeto \(x\) del grupoide de memoria,
el dominio correctamente tipado sería el grupo de isotropía
\[
 H_x:=\operatorname{Aut}_{\mathcal G_{\rm mem}}(x)
\]
y, por multiplicatividad hacia un grupo abeliano, su abelianización
\[
 H_x^{\rm ab}\longrightarrow\mathbb R_{>0}^{\times}.
\]
Todo carácter positivo anula la parte torsional de \(H_x^{\rm ab}\). Por ello,
la propuesta debe construir un levantamiento
\[
 \widetilde{\mathcal K}_{27}\in H_x,
 \qquad
 p(\widetilde{\mathcal K}_{27})=\mathcal K_{27},
\]
probar que su clase abeliana es no torsional y seleccionar de manera natural
un carácter \(\widetilde\chi_G\). La ruta positiva sobre el ciclo observable
queda así cerrada por obstrucción: sin ese dato no torsional no existe un
selector multiplicativo no trivial. Si todo levantamiento admisible
satisficiera
\(\widetilde{\mathcal K}_{27}^{\,4}=I\), la ruta multiplicativa quedaría
descartada también en el dominio enriquecido.

\section{Vía cronológica de longitud, acción y espectro}
\label{sec:via-cronologica-espectral-G}

El \cref{thm:cobertura-acumulativa-dirigida} proporciona una tercera vía,
distinta tanto del carácter observable obstruido como del eventual
levantamiento en un grupo de isotropía. Sobre
\[
 \widetilde{\mathcal O}_{108}^{+}
 =\mathcal O_{108}\times\mathbb N^2,
\]
el retorno de veintisiete pasos suma \((1,18)\) y su cuarta potencia suma
\((4,72)\). Esta acción pertenece a una categoría cronológica dirigida. No es
un levantamiento del grupoide ni un campo que se haya demostrado nativo del
estado holonómico integral completo.

El funcional positivo ya construido,
\[
 \mathcal A_\ast(\gamma)
 =G_{\rm or}(\gamma)+\lambda_\ast S(\gamma),
\]
es aditivo por concatenación de historias dirigidas. No es un carácter del
grupo de isotropía, pues sus indicadores no cambian de signo bajo una
inversión formal. La completación
\(\mathbb N^2\to\mathbb Z^2\) permite introducir caracteres unitarios para
análisis armónico, pero no altera este hecho.

\begin{equation}
\label{eq:fraccion-angular-orientada-G}
 \delta_\theta
 :=\frac{A_{\rm deg}-C_{\rm deg}^{\ast}}{360}.
\end{equation}
Esta fracción angular orientada se define una vez construidos el ángulo y el
contraángulo únicos. Se utilizará en las dos plantillas dimensionales
posteriores y no interviene retroactivamente en la obtención de
\(A_{\rm deg}\) ni de \(C_{\rm deg}^{\ast}\).

\subsection{Operador torcido y control negativo determinista}

Para el control exacto, sea \(\Gamma=(V,E)\) un multigrafo dirigido
\emph{finito} enriquecido, con medida
\(\mu:V\to\mathbb R_{>0}\), pesos \(a_\ast(e)\geq0\), parámetro
\(\beta\geq0\) y cociclo
\(c(e)\in\mathbb Z^2\). Para una normalización fijada, considérese
\[
 (\mathcal L_{\beta,\vartheta}f)(y)
 =
 \sum_{e:x\to y}
 e^{-\beta a_\ast(e)}
 e^{i\langle\vartheta,c(e)\rangle}f(x),
 \qquad
 \vartheta\in\mathbb T^2,
\]
y tómese como dominio todo el espacio de Hilbert finito
\(\ell^2(V,\mu)\). Entonces \(\mathcal L_{\beta,\vartheta}\) es acotado y
\[
 \mathcal H_{\beta,\vartheta}
 :=\mathcal L_{\beta,\vartheta}^{*}
   \mathcal L_{\beta,\vartheta}\geq0
\]
es autoadjunto. Una extensión a un grafo infinito requeriría, además,
operadores acotados o bien operadores cerrados densamente definidos tales
que \(\mathcal L_{\beta,\vartheta}^{*}
\mathcal L_{\beta,\vartheta}\) fuese autoadjunto sobre un dominio común;
esa extensión no se presupone aquí.

\begin{proposition}[Cancelación de fase en la órbita determinista]
\label{prop:control-negativo-espectral-determinista}
Si \(\Gamma\) es únicamente la órbita determinista observable y cada vértice
posee una sola arista entrante, entonces
\(\mathcal H_{\beta,\vartheta}\) es independiente de \(\vartheta\). En
particular, el Hessiano de cualquiera de sus bandas respecto de
\(\vartheta\) es nulo.
\end{proposition}

\begin{proof}
En cada componente, \(\mathcal L_{\beta,\vartheta}\) multiplica el único
término entrante por una fase unitaria. En
\(\mathcal L_{\beta,\vartheta}^{*}\mathcal L_{\beta,\vartheta}\), esa fase
se multiplica por su conjugada y se cancela. No quedan términos cruzados
entre historias distintas.
\end{proof}

Este control impide atribuir rigidez espectral al mero ciclo determinista.
Una respuesta de fase no trivial exige, al menos, una de las construcciones
siguientes:
\begin{enumerate}
 \item un multigrafo enriquecido en el que dos o más historias coherentes
 contribuyan a una misma amplitud; o
 \item un diferencial magnético
 \[
  (d_\vartheta f)(e)
  =
  e^{i\langle\vartheta,c(e)\rangle/2}f(t(e))
  -
  e^{-i\langle\vartheta,c(e)\rangle/2}f(s(e)),
 \]
 con laplaciano
 \(\Delta_\vartheta=d_\vartheta^*d_\vartheta\), cuya fase dependa de la
 holonomía de ciclos enriquecidos.
\end{enumerate}
En ambos casos deben declararse el espacio de Hilbert, el dominio, la medida,
los pesos y la normalización, y ha de demostrarse que la familia no es
unitariamente equivalente a la familia sin fase.

Supongamos que
\(\vartheta\mapsto\mathcal H_{\beta,\vartheta}\) es una familia
autoadjunta analítica sobre un dominio común y que
\(\lambda_0(\vartheta)\) es una banda simple y aislada en un entorno de
\(\vartheta=0\). Si \(0\) es un punto crítico y un mínimo local,
\[
 Q_{ij}
 =
 \left.
 \frac{\partial^2\lambda_0(\vartheta)}
      {\partial\vartheta_i\partial\vartheta_j}
 \right|_{\vartheta=0}
\]
es semidefinida positiva. Sin la hipótesis de mínimo, \(Q\) sólo representa
la respuesta cuadrática. Para
\[
 v_{27}=(1,18),
 \qquad
 \mathcal R_{27}=v_{27}^{\mathsf T}Qv_{27},
\]
la condición \(\mathcal R_{27}\neq0\) es una prueba mínima de sensibilidad
espectral a la memoria acumulada.

\subsection{Normalización tipada de la respuesta espectral y construcción de \texorpdfstring{\(\theta_G^{\mathrm{spec}}\)}{theta G spec}}

Si una normalización natural y única \(\mathfrak N\) convierte la respuesta
anterior en una razón positiva independiente de la escala global del
operador, puede definirse
\[
 \Lambda_G
 :=\ell_0\,\mathfrak N(Q,v_{27}),
 \qquad
 \theta_G^{\rm spec}
 :=
 \left(
 \delta_\theta\frac{\Lambda_G}{\ell_0}
 \right)^2,
\]
y, por el transductor dimensional,
\[
 \boxed{
 G_{\rm spec}
 =
 \frac{c_{\rm int}^{3}}{\hbar_{\rm int}}
 \bigl(\delta_\theta\Lambda_G\bigr)^2.}
\]
La homogeneidad de la plantilla es exacta. La clase de sus evaluaciones queda
caracterizada por las condiciones de selección canónica del operador,
respuesta no nula, naturalidad y estabilidad bajo refinamiento, independencia
de reescala, unicidad de \(\mathfrak N\) y sellado previo a cualquier consulta
del valor metrológico de \(G\). Sin esas condiciones, la invariancia por
reescala demuestra la no unicidad del selector; con ellas, la evaluación queda
determinada en el dominio normalizado declarado.

\section{Funcional no homomórfico sobre rutas enriquecidas}
\label{sec:funcional-longitud-K27-G}

La segunda continuación posible abandona el carácter multiplicativo y utiliza
un funcional geométrico, espectral o de acción. Sea
\[
 \mathscr L:
 \operatorname{Loop}(\mathcal X_{\rm TPK}^{\rm enr})
 \longrightarrow\mathcal L_L^+\cup\{0\}
\]
invariante por cambio cíclico de punto base, reetiquetado admisible,
conjugación y subdivisión. Su ley de concatenación debe declararse
explícitamente ---por ejemplo, como norma subaditiva o funcional espectral---;
no se presume que sea un homomorfismo del grupo torsional. Para un
levantamiento enriquecido candidato escribimos
\[
 \Lambda_{27}:=\mathscr L(\widetilde{\mathcal K}_{27}),
 \qquad
 \rho_{27}:=\frac{\Lambda_{27}}{c_{\rm int}t_0}.
\]
La fracción angular orientada de
\eqref{eq:fraccion-angular-orientada-G} produce el candidato adimensional
\[
 \theta_G^{\rm hol}:=(\delta_\theta\rho_{27})^2.
\]
Su especialización en el transductor da
\[
 \boxed{
 G_{\rm hol}
 =\frac{c_{\rm int}^{3}}{\hbar_{\rm int}}
  \bigl(\delta_\theta\Lambda_{27}\bigr)^2.}
\]
La forma temporal es la misma identidad:
\[
 T_{27}:=\frac{\Lambda_{27}}{c_{\rm int}},
 \qquad
 t_\ast:=\delta_\theta T_{27}.
\]
En particular, no corresponde sustituir \(t_\ast\) por
\(T_{27}/\delta_\theta\).

La homogeneidad de estas fórmulas es exacta. Su estatuto no asciende por ello
a derivación de \(G\) por esta plantilla: su selección queda abierta hasta construir
\(\widetilde{\mathcal K}_{27}\) y \(\mathscr L\) sin consultar el valor
externo, demostrar su naturalidad y fijar una normalización única. El uso
descendente de \(A\) y \(C^\ast\) en \(\delta_\theta\) no autoriza a utilizar
\(G\) para redefinirlos.
La prueba radial de \cref{g26:sec:rigidez-radial-es} constituye una
realización especificada diferente: sus datos, norma, longitud y
acoplamiento se construyen explícitamente y su reducción conserva la
memoria. La limitación de la plantilla anterior no se transfiere a esa
demostración.

\section{Transducción de la carta decimal gravitatoria sobre la recta de magnitud y separación causal del reconocimiento metrológico}
\label{sec:lector-decimal-independiente-G}

Sea
\[
 \mathscr R_G([n,k,d]_3;t,e)
 :=10^{-n}\left(k+\frac{t}{10^3}+\frac{e}{10^d}\right),
\]
donde el subíndice de \([n,k,d]_3\) registra la terna del lector y no denota
un numeral ternario. Para los datos declarados
\[
 [n,k,d]_3=[11,6,5]_3,\qquad t=674,\qquad e=30,
\]
la evaluación racional es
\[
 \boxed{
 \mathscr R_G
 =10^{-11}\left(6+\frac{674}{1000}+\frac{30}{10^5}\right)
 =6.67430\times10^{-11}.}
\]
La igualdad aritmética es exacta una vez declarados los parámetros. La terna,
los bloques \(674\), \(30\) y su posición decimal tienen el estatuto de datos
de la carta, de modo que esta rama se conserva como control parametrizado y
reconocimiento posterior. El número es adimensional hasta elegir una base
\(u_G\in\mathcal L_G^+\); sólo entonces se define la sección
\[
 G_{\rm dec}:=\mathscr R_G\,u_G.
\]
La coincidencia de su coordenada con una publicación metrológica es
\textsc{reconocimiento externo}. Una promoción futura exigiría que los enteros
\(674\), \(30\), la terna \([11,6,5]\) y la colocación decimal quedaran
forzados de manera única por simetrías publicadas y se congelaran antes de
abrir el valor externo.

\section{Jerarquía probatoria de las construcciones gravitatorias}
\label{sec:cuadrado-contraste-G}

La realización longitudinal y radial determina el acoplamiento en el
dominio explícito R1--R8. Las plantillas alternativas y los controles
anteriores conservan sus propios alcances; no constituyen por sí mismos
otras genealogías demostradas de esa realización. La organización es:
\begin{center}
\begin{tabularx}{0.96\textwidth}{@{}L{0.25\textwidth}Y L{0.26\textwidth}@{}}
\toprule
Objeto & Resultado disponible & Veredicto\\
\midrule
Inscripción y norma radial R1--R8 &
\(L=(5759/23040)\alpha^{16}L_*\),
\(R_X\Lambda_X=L^2I\), \(H_X\Lambda_X=\hbar_{\rm ret}cI\),
\(G=c^3L^2/\hbar_{\rm ret}\). &
Determinación única en el dominio declarado; memoria, reducción y
refinamiento conservados.\\
Realización circular CT108 &
\(t_0=\pi L/(54c)\), \(R_{108}=L\) y
\(\theta_{108}^{\rm circ}=(54/\pi)^2\). &
Compatibilidad exacta con la longitud construida.\\
Transductor dimensional &
Tipa exactamente \(G_{\rm int}(\theta_G)\), sin fijar \(\theta_G\). &
Teorema dimensional; no selector.\\
Carácter del observable \(\mathcal K_{27}\) &
\(\mathcal K_{27}^4=I\) fuerza
\(\chi(\mathcal K_{27})=1\). &
Formulación obstruida.\\
Levantamiento enriquecido &
Posible sólo con clase abeliana no torsional, carácter natural y
normalización única. &
Obstrucción de canonicidad demostrada sin esos datos.\\
Cobertura cronológica y respuesta espectral &
La traslación acumulativa tiene orden semigrupal infinito; en la órbita
determinista la fase se cancela. &
Respuesta nula demostrada en la órbita determinista; clase no trivial
caracterizada por interferencia y normalización natural.\\
Funcional de ruta &
La fórmula de \(G_{\rm hol}\) es homogénea una vez fijado
\(\Lambda_{27}\). &
Familia parametrizada exacta; no existe selector unívoco sin especificar
\(\mathscr L\) y su normalización.\\
Lector decimal &
Reproduce exactamente la coordenada declarada, pero recibe sus bloques y
posición. &
Control parametrizado.\\
\bottomrule
\end{tabularx}
\end{center}

\begin{criterion}[Selector de la constante gravitatoria]
\label{crit:selector-G}
Para las plantillas alternativas de levantamiento, respuesta espectral
y lector decimal examinadas en este capítulo, una promoción se rechaza si:
\begin{enumerate}
 \item el levantamiento enriquecido conserva únicamente torsión o no desciende
 a una clase abeliana bien definida;
 \item el operador cronológico no conserva interferencia, produce
 \(\mathcal R_{27}=0\), carece de dominio o medida, o su respuesta cambia bajo
 un refinamiento admisible;
 \item \(\mathscr L\) depende del punto base, del reetiquetado, de la
 orientación elegida, de una subdivisión o de una normalización libre;
 \item la normalización \(\mathfrak N\) depende de la escala global del
 operador, no es única o se fija consultando \(G\);
 \item en el lector decimal, \(674\), \(30\) o la colocación decimal admiten alternativas con el
 mismo derecho estructural;
 \item el candidato consulta el valor metrológico antes de quedar sellado;
 \item el valor no es estable bajo refinamiento o no define una sección de
 \(\mathcal L_G\) en una carta dimensional común.
\end{enumerate}
\end{criterion}

\section{Razones gravitatoria--inercial y constante dimensional}
\label{sec:razones-no-rivales-G}

La normalización local de la carta torsional,
\[
 \frac{\mathcal G_{\rm tor}}{\mathcal I_{\rm tor}}=1,
\]
y la interfaz areal--volumétrica del
\cref{sec:ambivalencia-grav-inercial-pi-phi2},
\[
 \frac{\mathcal G_{\rm av}}{\mathcal I_{\rm av}}
 =\frac{\pi}{\varphi^{2}},
\]
son razones adimensionales de lectores. No son dos valores rivales de la
constante dimensional \(G\). La primera expresa una restricción diagonal
local; la segunda es una interfaz condicionada a la factorización por una
memoria común. La construcción de una sección de
\(\mathcal L_G\) requiere, además, el marco
\((\hbar_{\rm int},c_{\rm int},\ell_0)\) y un selector
\(\theta_G\).
 
\chapter{Dinámica de Perron, memoria torsional y realización Einstein--Cartan}
\begingroup
\renewcommand{\chapter}[2][]{}%
\chapter{Dinámica de Perron, dilución \texorpdfstring{\(a^{-6}\)}{a⁻⁶} de la memoria y realización Einstein–Cartan}
\label{chap:parte-iv-67}

El operador de Perron fija el modo estable de memoria; al publicarse como densidad de espín, su contribución se diluye como la sexta potencia inversa del factor de escala. La realización Einstein--Cartan conserva esta procedencia y sus condiciones de validez.

\section{Din\'amica de Perron y modo estable de memoria}

En una ruta cerrada, el cociclo cil\'indrico de Perron satisface

\begin{equation}
V_n=\varphi^n.
\label{rm:eq:cosmo-perron-V}
\end{equation}

En dimensi\'on espacial \(D=3\), la escala determinada por el cociclo es

\begin{equation}
\frac{a_n}{a_0}=V_n^{1/3}=\varphi^{n/3},
\label{rm:eq:cosmo-perron-a}
\end{equation}

mientras que el modo estable de la acci\'on cuadr\'atica es

\begin{equation}
M_n=\varphi^{-2n}=V_n^{-2}.
\label{rm:eq:cosmo-perron-M}
\end{equation}

Eliminar \(n\) entre \cref{rm:eq:cosmo-perron-a,rm:eq:cosmo-perron-M} da el
resultado central:

\begin{theorem}[Ley cil\'indrica de sexto orden para la memoria]
Para todo nivel cerrado de la realizaci\'on de Perron en \(D=3\),
\begin{equation}
\boxed{
M_n=\left(\frac{a_n}{a_0}\right)^{-6}.}
\label{rm:eq:cosmo-minus-six}
\end{equation}
El exponente \(-6\) no se ajusta: es el cociente entre el exponente estable
\(-2\) y la ra\'iz dimensional \(1/3\).
\end{theorem}

En los retornos distinguidos se obtienen

\[
\frac{a_9}{a_0}=\varphi^3,
\qquad
\frac{a_{27}}{a_0}=\varphi^9,
\qquad
\frac{a_{108}}{a_0}=\varphi^{36},
\]

y los mismos niveles verifican exactamente
\(M_n=(a_n/a_0)^{-6}\). Si se declara el par\'ametro temporal aditivo \(t_n=nt_0\), la tasa
logar\'itmica es constante:

\begin{equation}
H_\varphi
:=\frac{\Delta\log a_n}{\Delta t}
=\frac{\log\varphi}{3t_0},
\qquad
H_\varphi t_0
=0.1604039416865344824992529711\ldots.
\label{rm:eq:cosmo-Hphi}
\end{equation}

El s\'imbolo \(H_\varphi\) nombra la tasa de esta escala discreta. Su lectura
como par\'ametro de Hubble requiere el transductor cosmogr\'afico; la identidad
interna no depende de esa promoci\'on.

\section{Realizaci\'on en el sistema de Einstein--Cartan}

En un fluido de esp\'in de Einstein--Cartan, la densidad cuadr\'atica de esp\'in
obedece

\begin{equation}
s^2\propto a^{-6}.
\label{rm:eq:cosmo-EC-scaling}
\end{equation}

La \cref{rm:eq:cosmo-minus-six} suministra el exponente y el modo estable sin
consultar un rebote. Para hacer la lectura f\'isica se debe declarar una
secci\'on positiva \(s_*^2\) y un mapa

\begin{equation}
\mathcal T_{\rm EC}:M_n\longmapsto
s_n^2:=s_*^2M_n
=s_*^2\left(\frac{a_n}{a_0}\right)^{-6}.
\label{rm:eq:cosmo-EC-transducer}
\end{equation}

El mapa conserva la ley de escala, pero no determina por s\'i solo la amplitud
\(s_*^2\), la corriente de esp\'in, el signo del t\'ermino de contorsi\'on ni la
ecuaci\'on de estado. Con densidad de masa, una ecuaci\'on efectiva tiene la
forma

\begin{equation}
H^2=\frac{8\pi G}{3}
\bigl(\rho_{\rm m}-\rho_{\rm spin,m}\bigr)+\cdots,
\label{rm:eq:cosmo-EC-mass}
\end{equation}

mientras que con densidad de energ\'ia debe escribirse

\begin{equation}
H^2=\frac{8\pi G}{3c^2}
\bigl(\rho_{\rm E}-\rho_{\rm spin,E}\bigr)+\cdots.
\label{rm:eq:cosmo-EC-energy}
\end{equation}

Las dos convenciones no se mezclan. Un candidato a rebote exige, adem\'as de
la igualdad de densidades, regularidad, materia y datos de frontera.

\begin{falsifier}
La transducci\'on Einstein--Cartan falla si el funcional de escala produce un
exponente distinto de \(-6\), si \(V_n\ne\varphi^n\) en la ruta declarada, o
si el mapa f\'isico no conserva positividad y dimensiones. Elegir
\(s_*^2\) retrospectivamente para forzar un umbral de rebote no demuestra el
transductor.
\end{falsifier}

\section{Comparaci\'on estructural entre el ciclo de 108 estados y la din\'amica de Perron}

\begin{center}
\begin{tabularx}{\textwidth}{>{\bfseries}p{29mm}XX}
\toprule
& CT108--de Sitter & Perron--Einstein--Cartan\\
\midrule
Antecedente & retorno circular de orden \(108\) & cociclo cil\'indrico y modo estable\\
Selector & \(r_g=R_{108}\) & \(V_n=\varphi^n\), \(D=3\)\\
Magnitudes determinadas & radio, curvatura, temperatura y entrop\'ia & expansi\'on, tasa y diluci\'on de memoria\\
Invariante & \(G_a\hbar_a\) y radio circular & \(M_na_n^6\)\\
Estatuto & interfaz de Sitter condicionada & ley interna exacta; lectura EC condicionada\\
Criterio de refutación & radio o ecuaci\'on de fondo incompatibles & p\'erdida del exponente \(-6\) o del transductor\\
\bottomrule
\end{tabularx}
\end{center}

\begin{statusbox}
Las identidades CT108, el selector \(\theta_{108}\) y la ley
\(M_n=(a_n/a_0)^{-6}\) son \status{Teorema interno exacto}. La lectura del mismo radio
como de Sitter y la identificaci\'on de \(M_n\) con una densidad de esp\'in son
\status{Realizaciones f\'isicas condicionadas}, cada una con su propio criterio de refutación.
\end{statusbox}
 
\paragraph{Coordinación de torsión y área en la acción Palatini--Cartan--Holst.} La torsión y el área se coordinan después en la acción de primer orden de Palatini--Cartan--Holst.
 \endgroup

\chapter{Realización Palatini--Cartan--Holst del parámetro de Barbero--Immirzi}
\begingroup
\renewcommand{\chapter}[2][]{}%
\chapter{Realización Palatini–Cartan–Holst del parámetro de Barbero–Immirzi y transducción física del espectro de área}
\label{chap:parte-iv-68}

La derivación estructural del parámetro y del espectro de área antecede a este capítulo. Aquí se construyen tétrada, conexión, curvatura y acción de primer orden, y sólo entonces se especializa el parámetro de Holst como realización física del lector areal.

\section{Curvatura, torsión y acción Einstein--Cartan--Holst}
\label{chap:gravedad-cartan-holst}

La estructura de transporte HMT distingue la variable de memoria, la
holonomía y el defecto de cierre. Su realización gravitatoria se formula
mediante una tétrada y una conexión. Esta sección establece la geometría de
primer orden y especifica el tipo del mapa entre ésta y el espacio de estados
TPK\@.

\subsection{Identidad discreta de Gauss–Stokes y transporte de incidencia a la frontera}

Sea \(K\) un complejo celular, sea \(W\) un grupo abeliano y sea
\(\Omega\in C^1(K;W)\) una cocadena. Denotamos por
\(\delta:C^1(K;W)\to C^2(K;W)\) el coborde celular y por
\(\langle\,\cdot\,,\,\cdot\,\rangle\) el emparejamiento entre cocadenas con
valores en \(W\) y cadenas celulares enteras.

\begin{theorem}[Identidad celular universal de Stokes]
Para toda \(2\)-cadena finita \(z\in C_2(K;\mathbb Z)\),
\[
 \boxed{\langle\delta\Omega,z\rangle
 =\langle\Omega,\partial z\rangle.}
\]
\end{theorem}
\begin{proof}
Es la identidad que define el coborde como operador adjunto de la frontera
celular. La finitud de \(z\) garantiza que ambos emparejamientos son sumas
finitas en \(W\).
\end{proof}

\begin{corollary}[Gauss--Stokes sobre una pseudovariedad celular]
Sea \(K\) una pseudovariedad celular regular, bidimensional, finita y compacta,
coherentemente orientada, posiblemente con borde. Si \(z_K\) es la suma de
sus \(2\)-celdas con la orientación coherente, entonces
\[
 \sum_{f\in K^{(2)}}(\delta\Omega)(f)
 =\sum_e[\partial z_K:e]\,\Omega(e).
\]
En particular, la regularidad permite identificar las incidencias no
canceladas de \(\partial z_K\) con las aristas del borde geométrico.
\end{corollary}
\begin{proof}
Se aplica el teorema a \(z_K\). La orientación coherente hace que cada arista
interior aparezca en \(\partial z_K\) dos veces con coeficientes opuestos;
las incidencias no canceladas forman la cadena de borde.
\end{proof}

La versión no abeliana no se obtiene sustituyendo sumas por productos sin más.
En esta obra se utiliza únicamente cuando se han fijado transportes a una base
común, un orden de las holonomías y la conjugación correspondiente. Fuera de
ese subdominio no se afirma una identidad de Stokes no abeliana.

\subsection{Capacidad areal de la celda qutrit: \(81\log 3\) y transporte colectivo de incidencia}

\begin{theorem}[Corte mínimo de un cubo discreto]
En el grafo cúbico \(N\times N\times N\) con capacidad unitaria, todo corte
que separa dos caras opuestas tiene capacidad al menos \(N^2\), y existe uno
de capacidad \(N^2\).
\end{theorem}
\begin{proof}
Hay \(N^2\) caminos disjuntos en aristas que atraviesan el cubo en la dirección
normal a las dos caras. Todo corte debe interceptarlos, luego su capacidad es
al menos \(N^2\). Una capa transversal contiene exactamente \(N^2\) aristas y
realiza la cota.
\end{proof}

Para \(N=9\), el corte vale 81 y
\[
 9^3=729=27^2.
\]
La última igualdad permite reagrupar seis trits como un volumen \(9^3\) o una
cuadrícula de cardinalidad \(27^2\). Es una biyección de conjuntos; no es un
isomorfismo entre
\((\ZZ/9\ZZ)^3\) y \((\ZZ/27\ZZ)^2\), cuyos exponentes son distintos.

\subsection{Tétrada, conexión y \(2\) defectos}

Sean \(e^I\) una tétrada y \(\omega^{IJ}=-\omega^{JI}\) una conexión de
Lorentz. Definimos
\[
 T_{\rm tor}^I=D_\omega e^I
 =\dd e^I+\omega^I{}_J\wedge e^J,
\]
\[
 F^{IJ}=\dd\omega^{IJ}+\omega^I{}_K\wedge\omega^{KJ},
 \qquad
 g=\eta_{IJ}e^I\otimes e^J.
\]
La descomposición
\[
 \omega^{IJ}=\mathring\omega^{IJ}(e)+K^{IJ}
\]
separa la conexión de Levi--Civita y la contorsión. Entonces
\[
 T_{\rm tor}^I=K^I{}_J\wedge e^J.
\]
Curvatura y torsión son defectos diferentes: \(F\) mide el cierre del
transporte de vectores; \(T_{\rm tor}\), el cierre de la tétrada. La memoria
\(T\) del principio de Ambivalencia tampoco es torsión geométrica: sólo un
funtor de realización puede relacionarlas.

\subsection{Acción de 1.er orden}

Sea \(M\) una variedad lisa orientada de dimensión cuatro. Si
\(\partial M\ne\varnothing\), suponemos variaciones de soporte compacto en el
interior o, equivalentemente para el problema variacional considerado, un
término de borde y condiciones de contorno que anulen la contribución
fronteriza. Cuando \(\Psi\) contenga campos espinoriales, suponemos además una
estructura de espín y campos compatibles con ella. Sean
\[
 \Sigma^{IJ}=e^I\wedge e^J,
 \qquad
 (\star X)^{IJ}=\frac12\epsilon^{IJ}{}_{KL}X^{KL},
 \qquad
 \kappa=\frac{8\pi G}{c^4},
 \qquad
 \gamma\in\mathbb R\setminus\{0\}.
\]
La acción Palatini--Cartan--Holst es
\[
\boxed{
 S[e,\omega,\Psi]
 =\frac1{2\kappa c}\int_M
 \Sigma_{IJ}\wedge\left(\star F^{IJ}+\frac1\gamma F^{IJ}\right)
 -\frac\Lambda{\kappa c}\int_M\operatorname{vol}_e
 +S_{\rm m}[e,\omega,\Psi].}
\]
La variación de esta acción determina las correcciones torsionales
\cite{Sciama1964,Kibble1961,Hehl1976,Holst1996}.

Definimos la corriente de espín por
\[
 \delta_\omega S_{\rm m}
 =-\frac12\int_M\delta\omega^{IJ}\wedge\sigma_{IJ}.
\]

\begin{theorem}[Ecuación de Cartan--Holst]
Bajo las hipótesis geométricas y de borde anteriores, la variación respecto
de \(\omega\) da
\[
 D_\omega\!\left[(\star+\gamma^{-1}I)\Sigma^{IJ}\right]
 =\kappa c\,\sigma^{IJ},
\]
con la normalización global fijada por la definición anterior de
\(\sigma^{IJ}\).
\end{theorem}
\begin{proof}
Se usa \(\delta F^{IJ}=D_\omega\delta\omega^{IJ}\), se integra por partes y
se iguala a cero el coeficiente de la variación arbitraria
\(\delta\omega^{IJ}\). Con la convención de signo fijada en la definición de
\(\sigma^{IJ}\), el término de materia pasa al miembro derecho con signo
positivo. El factor \(1/(2\kappa c)\) produce el coeficiente
\(\kappa c\). El término cosmológico no depende de \(\omega\), y las
condiciones de borde eliminan el término generado por la integración por
partes.
\end{proof}

En firma lorentziana \(\star^2=-I\) sobre bivectores y, para
\(\gamma\in\RR\setminus\{0\}\),
\[
 (\star+\gamma^{-1}I)^{-1}
 =\frac{\gamma^2}{1+\gamma^2}(-\star+\gamma^{-1}I).
\]

\begin{corollary}[Desacoplamiento torsional]
Para una co-tétrada no degenerada y
\(\gamma\in\mathbb R\setminus\{0\}\),
\[
 \sigma^{IJ}=0\Longrightarrow T_{\rm tor}^I=0
 \Longrightarrow\omega=\mathring\omega(e).
\]
\end{corollary}

Para hacer explícita la normalización, definimos las fuentes normalizadas por
la acción mediante
\[
 \delta_g S_{\rm m}
 =-\frac12\int_M\operatorname{vol}_e\,
 \mathcal T^{S}_{\mu\nu}\,\delta g^{\mu\nu}.
\]
Cuando las coordenadas de la tétrada tienen dimensión de longitud, el tensor
energía--momento físico convencional es
\(T^{\rm phys}_{\mu\nu}:=c\,\mathcal T^{S}_{\mu\nu}\). Denotamos de igual
modo por \(\mathcal U^{S,\rm spin}\) la contribución normalizada por la
acción y por \(U^{\rm phys,\rm spin}:=c\,\mathcal U^{S,\rm spin}\) su
convención física. Al eliminar algebraicamente la conexión, la ecuación
métrica adopta entonces las dos formas equivalentes
\[
 G_{\mu\nu}+\Lambda g_{\mu\nu}
 =\kappa c\,
 \bigl(\mathcal T_{\mu\nu}^{S,\rm LC}
       +\mathcal U_{\mu\nu}^{S,\rm spin}\bigr)
 =\kappa\,
 \bigl(T_{\mu\nu}^{\rm phys,LC}
       +U_{\mu\nu}^{\rm phys,spin}\bigr),
\]
donde la contribución de espín es cuadrática en la corriente correspondiente
y depende del acoplamiento de materia. El coeficiente de esa contribución no
se congela sin especificar la acción fermiónica y las convenciones de
dualidad.

\subsection{Identidad de Bianchi con torsión y balance de flujo en la conexión de Cartan}

Para toda conexión,
\[
 D_\omega T_{\rm tor}^I=F^I{}_J\wedge e^J,
 \qquad
 D_\omega F^{IJ}=0.
\]
Después de eliminar torsión,
\[
\nabla^\mu(T_{\mu\nu}^{\rm phys,LC}
          +U_{\mu\nu}^{\rm phys,spin})=0.
\]
Estas identidades son pruebas de consistencia obligatorias: cualquier término
gravitatorio HMT debe surgir de una acción compatible o satisfacer el balance
de Noether correspondiente.

\subsection{Mapa de realización y especialización del parámetro de Holst}

La inserción definida en el \cref{mdg:chap:barbero} toma
\[
 \gamma=\Gamma_{6\to8}.
\]
Una vez realizada esta sustitución, todas las consecuencias variacionales
anteriores son teoremas de la acción especializada. La selección física de
esta acción por el transporte HMT se expresa mediante un mapa de realización
\[
 \mathfrak R_{\rm grav}:\mathcal X_{\rm TPK}
 \dashrightarrow\{(e,\omega,\Psi)\}
\]
dotado de tres propiedades: entrelazar el transporte TPK con el de
\(\omega\), identificar las clases de holonomía y suministrar las escalas de
\(G\) y \(\Lambda\). Estas propiedades fijan el dominio preciso de una
realización gravitatoria del formalismo.

Para co-tétrada no degenerada, \(\gamma\) real no nulo y materia sin corriente
de espín, el sector torsional se reduce a la formulación de primer orden de
relatividad general con la misma constante cosmológica, sujeto a las
condiciones de borde ya declaradas. En un fluido de espín, el escalado
\(s^2\propto a^{-6}\) puede producir un rebote Einstein--Cartan; obtener su
densidad y sus condiciones iniciales a partir del estado TPK es una cuestión
distinta de la consistencia variacional de la acción.
 
\section{Retorno nonádico, grupoide de memoria y condición de entrelazamiento de Cartan}

La jerarquía \(27\to54\to108\), la cobertura acumulativa dirigida y el
grupoide de memoria suministran la condición exacta bajo la cual el retorno
del TPK admite una conexión de Cartan compatible. Este puente antecede a la
especialización Palatini--Cartan--Holst y conserva íntegros su dominio y sus
criterios de refutación.

\begingroup
\let\chapter\section
\let\section\subsection
\let\subsection\subsubsection
\let\subsubsection\paragraph
\let\clearpage\relax
\let\cleardoublepage\relax
\chapter{Retorno nonádico, memoria y holonomía de Cartan}
\label{chap:puente-retorno-holonomia-cartan}
\index{retorno!27--54--108}\index{holonomía!TPK--Cartan}
\index{memoria!distinción respecto de torsión}

La realización gravitatoria no comienza identificando una variable discreta
con una magnitud geométrica. Comienza por distinguir tres niveles: el retorno
del transporte observable, la memoria que diferencia historias con igual
proyección y la holonomía de una conexión. El puente que se formula aquí
recibe los retornos y cociclos ya demostrados en HMT; no los reconstruye
desde la acción de Holst, desde una masa ni desde una constante metrológica.

\section{Jerarquía exacta de los retornos}
\label{sec:jerarquia-retornos-27-54-108}

Sea \(F_{\rm obs}\) la cronología observable definida en el núcleo HMT y sea
\[
 \mathcal K_{27}:=F_{\rm obs}^{27}.
\]
El símbolo \(\mathcal K_{27}\) designa exclusivamente la composición de
veintisiete actualizaciones del lazo posicional y orientado; no designa una
ventana genérica, una traza de \(108\) pasos ni el estado holonómico integral
completo. Denotemos por \(p\) la proyección a las posiciones de los dos
cursores y por \(\mathsf J_{\rm or}\) la inversión simultánea de sus
orientaciones. La demostración primaria se encuentra en
\cref{prop:retorno-fobs-identidad}; sus consecuencias tipadas son
\[
 p(\mathcal K_{27}x)=p(x),\qquad
 o(\mathcal K_{27}x)=\mathsf J_{\rm or}\,o(x),
 \qquad \mathsf J_{\rm or}^{\,2}=I,
\]
\[
 F_{\rm obs}^{54}:
 \text{restaura posiciones y orientaciones},
 \qquad
 F_{\rm obs}^{108}=I:
 \text{restaura además la fase observable}.
\]
Por tanto, retorno de posición, retorno de orientación y retorno de fase son
afirmaciones diferentes. Ninguna de ellas implica por sí sola el retorno de
todos los campos de una extensión enriquecida.

El cociente de modos y los cociclos de orientación se establecen en
\cref{sec:cociclos-modo-orientacion,thm:descenso-necesario-fav}. En el orden
areal--volumétrico,
\[
 N_{27}=9F_{\rm av},
 \qquad
 F_{\rm av}=
 \begin{pmatrix}0&1\\1&1\end{pmatrix}.
\]
Si
\[
 \Omega_{\rm sw}(\gamma)=\sum_{a\subset\gamma}\omega_{\rm sw}(a),
 \qquad
 S(\gamma)=\sum_{a\subset\gamma}|\omega_{\rm sw}(a)|,
 \qquad
 G_{\rm or}(\gamma)=\sum_{a\subset\gamma}g(a),
\]
entonces un bloque de veintisiete pasos contiene tres vueltas nonádicas y una
inversión de orientación:
\[
 \Omega_{\rm sw}(\mathcal K_{27})=0,\qquad
 S(\mathcal K_{27})=18,\qquad
 G_{\rm or}(\mathcal K_{27})=1.
\]
En el retorno observable de \(108\) pasos,
\[
 \Omega_{\rm sw}(C_{108})=0,\qquad
 S(C_{108})=72,\qquad
 G_{\rm or}(C_{108})=4.
\]
Así, para el funcional aditivo declarado
\[
 \mathcal A_\ast(\gamma)
 :=G_{\rm or}(\gamma)+\lambda_\ast S(\gamma),
\]
se obtienen, sin introducir una escala dimensional,
\[
 \mathcal A_\ast(\mathcal K_{27})=1+18\lambda_\ast,
 \qquad
 \mathcal A_\ast(C_{108})=4+72\lambda_\ast.
\]
Estos valores son conteos del registro. Su posterior uso térmico o
gravitatorio no altera su procedencia.

\subsection{Sistema dirigido de coberturas, refinamiento y límite de la conexión de Cartan}
\label{sec:cobertura-acumulativa-dirigida}

Sea \(\mathcal O_{108}\) la órbita observable y, para cada arista cronológica
\(e\), sea
\[
 c(e)=\bigl(g(e),s(e)\bigr)\in\mathbb N^2
\]
el par formado por los indicadores de inversión de orientación y de cambio
efectivo de hoja. Para una historia dirigida
\(\gamma=e_1\cdots e_n\), se define
\[
 c(\gamma)=\sum_{j=1}^{n}c(e_j).
\]
La ley de concatenación toma la forma
\[
 c_{m+n}(x)
 =c_m(x)+c_n\!\left(F_{\rm obs}^{m}x\right).
\]
Los conteos anteriores expresan, de manera uniforme sobre la órbita,
\[
 c_{27}(x)=(1,18),
 \qquad
 c_{108}(x)=(4,72).
\]

\begin{theorem}[Cobertura acumulativa de la cronología observable]
\label{thm:cobertura-acumulativa-dirigida}
Sobre
\[
 \widetilde{\mathcal O}_{108}^{+}
 :=\mathcal O_{108}\times\mathbb N^2
\]
defínase
\[
 \widetilde F^{+}(x,m)
 :=\bigl(F_{\rm obs}x,m+c(x\to F_{\rm obs}x)\bigr).
\]
Entonces el transporte de veintisiete pasos satisface
\[
 \widetilde{\mathcal K}_{27}^{+}(x,m)
 =\bigl(\mathcal K_{27}x,m+(1,18)\bigr),
\]
y
\[
 \left(\widetilde{\mathcal K}_{27}^{+}\right)^4(x,m)
 =\bigl(x,m+(4,72)\bigr)\neq(x,m).
\]
Por tanto, su proyección observable tiene orden cuatro, mientras que la
acción acumulativa del monoide cronológico posee orden semigrupal infinito.
\end{theorem}

\begin{proof}
La ley cocíclica y la constancia de \(c_{27}\) dan
\[
 c_{108}(x)
 =\sum_{j=0}^{3}c_{27}(\mathcal K_{27}^{j}x)
 =4(1,18)=(4,72).
\]
La primera coordenada retorna por
\(\mathcal K_{27}^{4}=I\); la segunda no retorna. Tras \(q>0\) iteraciones,
el incremento es \(q(4,72)\), distinto de cero.
\end{proof}

\begin{remark}[Tipo de la extensión acumulativa]
La construcción anterior amplía historias dirigidas mediante contadores
positivos. No es un levantamiento automorfo del grupoide
\(\mathcal G_{\rm mem}\), ni demuestra que la coordenada ilimitada pertenezca
ya al estado \(\mathcal X_{\mathrm{TPK}}^{\mathrm{enr}}\). Su completación de Grothendieck
\(\mathbb Z^2\) puede utilizarse para análisis armónico, pero no transforma
el coste positivo
\(\mathcal A_\ast=G_{\rm or}+\lambda_\ast S\) en un carácter del grupo de
isotropía: al invertir formalmente una ruta, sus indicadores no adquieren
signo opuesto.
\end{remark}

\paragraph{Dominio de validez de la extensión acumulativa.}
La jerarquía \(27\to54\to108\), las incidencias de \(F_{\rm av}\) y los
cociclos anteriores son identidades exactas en la cronología y los
lectores declarados. La afirmación se restringe al estado observable que
esos lectores conservan; extenderla al estado holonómico integral completo requeriría
probar el retorno de cada campo olvidado.

\section{Grupoide de memoria}
\label{sec:grupoide-memoria-tpk}

Sea \(\mathsf P_{\rm TPK}^{\rm enr}\) el grupoide de caminos enriquecidos:
sus objetos son estados TPK supervivientes y sus flechas son rutas admisibles
con origen y destino declarados, composición por concatenación e inversión
por reversión de ruta. Sea \(\sim_{\rm mem}\) una equivalencia que sólo
identifica dos flechas cuando conserva los campos de memoria requeridos por
el lector considerado y que es congruente con identidades, inversión y
concatenación. Definimos
\[
 \mathcal G_{\rm mem}
 :=\mathsf P_{\rm TPK}^{\rm enr}/\!\sim_{\rm mem}.
\]
Esta definición no transforma la memoria en torsión. La memoria pertenece al
estado discreto y distingue historias; la torsión pertenece a una conexión
geométrica y mide \(D_\omega e\). Sólo un mapa entre ambos dominios puede
relacionarlas.

Una realización geométrica exige, como datos adicionales, una longitud
\(\ell_0>0\), una representación del grupoide
\[
 \rho_{\rm C}:
 \mathcal G_{\rm mem}
 \longrightarrow
 \operatorname{Spin}(1,3)\ltimes\RR^{1,3},
\]
y un mapa de campos
\[
 \mathfrak R_{\rm grav}:
 \mathsf P_{\rm TPK}^{\rm enr}
 \dashrightarrow
 \mathsf{Cartan}_{1,3}.
\]
Para definir el segundo mapa deben especificarse su acción sobre objetos y
rutas, su compatibilidad con concatenación e inversión y la regularidad de la
conexión resultante. Ninguna de estas propiedades se deduce del mero conteo
de retornos.

\section{Conexión conjunta y condición de entrelazamiento}
\label{sec:conexion-cartan-conjunta}

Bajo las hipótesis geométricas del \cref{chap:gravedad} ---variedad
orientada de dimensión cuatro, co-tétrada \(e^I\), conexión de Lorentz
\(\omega^{IJ}\) y condiciones de borde allí declaradas---, la conexión de
Cartan conjunta se escribe
\[
 \mathcal A_{\ell_0}
 =
 \begin{pmatrix}
  \omega&\ell_0^{-1}e\\
  0&0
 \end{pmatrix}.
\]
Su curvatura es la identidad algebraica
\[
 \mathcal F_{\ell_0}
 =\dd\mathcal A_{\ell_0}
  +\mathcal A_{\ell_0}\wedge\mathcal A_{\ell_0}
 =
 \begin{pmatrix}
  F_\omega&\ell_0^{-1}T_{\rm tor}\\
  0&0
 \end{pmatrix},
\]
donde
\[
 F_\omega^{IJ}
 =\dd\omega^{IJ}+\omega^I{}_K\wedge\omega^{KJ},
 \qquad
 T_{\rm tor}^{I}=D_\omega e^I.
\]
Curvatura y torsión aparecen como componentes distintas de un mismo defecto
de transporte. La igualdad matricial se sigue algebraicamente de los datos
explícitos \(e,\omega,\ell_0\), pero no construye estas entradas desde TPK.

El contenido fuerte del puente es la conmutatividad de holonomías:
\[
 \boxed{
 \Hol_{\mathcal A_{\ell_0}}
 \bigl(\mathfrak R_{\rm grav}(\gamma)\bigr)
 =
 \rho_{\rm C}\!\left(
 \Hol_{\rm TPK}(\gamma)
 \right)}
 \qquad
 (\gamma\in\mathsf P_{\rm TPK}^{\rm enr}).
\]
La igualdad debe ser compatible con identidades, inversión, concatenación,
subdivisión, cambio de punto base y conjugación. Sólo entonces la memoria
discreta queda representada en una conexión; no queda rebautizada como
torsión ni como curvatura.

\begin{criterion}[Condiciones de existencia del puente TPK--Cartan]
\label{crit:puente-tpk-cartan}
No existe una realización con la conmutatividad y la compatibilidad declaradas
si ocurre cualquiera de los siguientes hechos:
\begin{enumerate}
 \item dos representantes equivalentes de una misma ruta producen
 holonomías geométricas no conjugadas;
 \item la imagen no preserva identidades, inversos o concatenación;
 \item una subdivisión admisible altera la holonomía;
 \item la igualdad de holonomías falla para \(\mathcal K_{27}\) o para sus
 concatenaciones de profundidades \(54\) y \(108\);
 \item el mapa identifica memoria y torsión sin exhibir un morfismo que
 preserve sus tipos;
 \item alguna elección descendente ---la acción de Holst, \(G\), una masa o
 una carta SI--- se utiliza para redefinir \(A\), \(C^\ast\) o el retorno
 TPK de origen.
\end{enumerate}
\end{criterion}

\paragraph{Antecedencia del funcional \(\Gamma_{6\to8}\) respecto de la realización Palatini--Cartan--Holst.}
El \cref{mdg:chap:barbero} fija, una vez construidos \(A,C^\ast,q_\pm\), el
funcional interno \(\Gamma_{6\to8}\). El \cref{chap:gravedad} estudia después
la acción Palatini--Cartan--Holst y sus ecuaciones variacionales. El puente
presente no introduce una ruta inversa entre ambas etapas.
 \endgroup

% Owner íntegro de la realización discreta del retorno, la holonomía y la
% respuesta colectiva. Se sitúa después del puente que declara sus hipótesis.
\begin{HMTScientificOwnerSubordinated}
\chapter{Gravedad: holonomía, torsión y respuesta colectiva}
\label{chap:gravedad}
\index{gravedad}\index{Cartan--Holst}\index{holonomía}
\index{torsión}\index{gravitón!no primitivo}

La gravedad no se introduce aquí mediante una fórmula decimal para \(G\).
Su construcción reúne dos ramas ya tipadas. En la primera, una ruta conserva
orientación y memoria; una conexión compara marcos en puntos distintos; su
holonomía mide el retorno y su curvatura, el defecto de cierre. En la segunda,
la carta dimensional construye \(G_{\rm int}(\theta_G)\) y
\(\kappa_{\rm int}\) a partir de la escala interna de acción, la velocidad y
la duración elemental. Sólo después de disponer de ambas ramas se escribe la
acción Palatini--Cartan--Holst. Ninguna cifra metrológica externa participa en
esa precedencia.

\section{Retornos: posición, orientación, fase y memoria}

Sea \(F_{\rm obs}\) la cronología observable del transporte y
\[
 \mathcal K_{27}=F_{\rm obs}^{27}.
\]
La jerarquía exacta distingue tres retornos:
\[
 \begin{array}{c|c}
 27\ \text{pasos}&
 \text{retorno de posición e inversión de orientación},\\
 54\ \text{pasos}&
 \text{retorno de posición y orientación},\\
 108\ \text{pasos}&
 \text{retorno completo de la fase observable}.
 \end{array}
\]
El retorno observable no borra el registro acumulado.  Si \(g(e)\) cuenta
una inversión de orientación y \(s(e)\) un cambio efectivo de hoja, entonces
\[
 c(\gamma)=\sum_{e\subset\gamma}(g(e),s(e))
\]
satisface la ley cocíclica de concatenación.  En los bloques canónicos,
\[
 c_{27}=(1,18),\qquad c_{108}=(4,72).
\]
Por tanto, la proyección observable de \(\mathcal K_{27}\) tiene orden
cuatro, mientras la coordenada acumulativa sigue creciendo.  Volver a la
misma fase no significa volver a la misma historia.

\begin{theorem}[Cobertura acumulativa del retorno]
\label{thm:cobertura-retorno-grav-rev6}
Sobre \(\mathcal O_{108}\times\mathbb N^2\), la elevación
\[
 \widetilde F(x,m)=
 \bigl(F_{\rm obs}x,m+c(x\to F_{\rm obs}x)\bigr)
\]
satisface
\[
 \widetilde{\mathcal K}_{27}(x,m)
 =\bigl(\mathcal K_{27}x,m+(1,18)\bigr)
\]
y
\[
 \widetilde{\mathcal K}_{27}^{\,4}(x,m)
 =\bigl(x,m+(4,72)\bigr)\neq(x,m).
\]
\end{theorem}

\begin{proof}
La constancia de \(c_{27}\) sobre la órbita y la ley de concatenación dan
\[
 c_{108}=\sum_{j=0}^{3}c_{27}(\mathcal K_{27}^jx)
 =4(1,18).
\]
La primera coordenada retorna; la segunda conserva la memoria acumulada.
\end{proof}

Éste es el dato que una realización geométrica debe preservar.  No se
identifica la memoria discreta con la torsión: se construye una representación
que permita compararlas.

\section{Del grupoide de caminos al grupo de Cartan}

Sea \(\mathcal G_{\rm mem}\) el grupoide de rutas enriquecidas, con
composición por concatenación e inversa por reversión. Fijemos un elemento
\(R_4\in\operatorname{Spin}^+(1,3)\) de orden cuatro y un vector temporal
unitario \(u\in\mathbb R^{1,3}\) fijado por \(R_4\). Si
\[
c(\gamma)=\bigl(c_g(\gamma),c_s(\gamma)\bigr)\in\mathbb Z^2
\]
es el cociclo aditivo de inversión y cambio de hoja, definimos
\[
\rho_{\rm C}^{\rm disc}(\gamma)
=
\left(
R_4^{\,c_g(\gamma)},
\ell_0c_s(\gamma)u
\right)
\in\operatorname{Spin}^+(1,3)\ltimes\mathbb R^{1,3}.
\]

\begin{theorem}[Realización discreta del cociclo de retorno]
\label{thm:realizacion-discreta-cartan-rev6}
\(\rho_{\rm C}^{\rm disc}\) es un funtor del grupoide de rutas en el grupo de
Cartan. En particular, conserva identidad, concatenación e inversión, y
transporta el retorno \(27\to54\to108\) sin borrar la coordenada acumulativa.
\end{theorem}

\begin{proof}
La ley del producto semidirecto es
\[
(R,a)(S,b)=(RS,a+Rb).
\]
Como \(R_4u=u\) y \(c\) es aditivo,
\[
\begin{aligned}
\rho_{\rm C}^{\rm disc}(\gamma_2)
\rho_{\rm C}^{\rm disc}(\gamma_1)
&=
\left(
R_4^{c_g(\gamma_2)+c_g(\gamma_1)},
\ell_0\bigl(c_s(\gamma_2)+c_s(\gamma_1)\bigr)u
\right)\\
&=\rho_{\rm C}^{\rm disc}(\gamma_2\gamma_1).
\end{aligned}
\]
La ruta identidad tiene cociclo nulo y la reversión cambia su signo, por lo
que se preservan identidad e inversa. Las fórmulas de retorno siguen de
\(c_{27}=(1,18)\) y \(c_{108}=4c_{27}\).
\end{proof}

Este teorema construye el puente discreto. Una realización suave posterior
consiste en una representación
\[
 \rho_{\rm C}:\mathcal G_{\rm mem}
 \longrightarrow\operatorname{Spin}(1,3)\ltimes\mathbb R^{1,3}
\]
y en campos \(e^I,\omega^{IJ}\) cuya holonomía entrelace esa representación.
El diagrama que fija el tipo es
\[
 \operatorname{Hol}_{\mathcal A_{\ell_0}}
   \bigl(\mathfrak R_{\rm C}(\gamma)\bigr)
 =
 \rho_{\rm C}\bigl(\operatorname{Hol}_{\rm TPK}(\gamma)\bigr).
\]
Identidades, inversión, concatenación y subdivisión deben conservarse.  Esta
ecuación no rebautiza el acarreo como curvatura: exige que el transporte
discreto ya representado y el geométrico realicen la misma composición.

Fijada una escala positiva \(\ell_0\), la conexión conjunta es
\[
 \mathcal A_{\ell_0}
 =
 \begin{pmatrix}
 \omega&\ell_0^{-1}e\\
 0&0
 \end{pmatrix}.
\]

\begin{proposition}[Curvatura conjunta de Cartan]
\label{prop:curvatura-conjunta-cartan-rev6}
La curvatura de \(\mathcal A_{\ell_0}\) es
\[
 \mathcal F_{\ell_0}
 =\dd\mathcal A_{\ell_0}
  +\mathcal A_{\ell_0}\wedge\mathcal A_{\ell_0}
 =
 \begin{pmatrix}
 F_\omega&\ell_0^{-1}T_{\rm tor}\\
 0&0
 \end{pmatrix},
\]
donde
\[
 F_\omega^{IJ}
 =\dd\omega^{IJ}+\omega^I{}_K\wedge\omega^{KJ},
 \qquad
 T_{\rm tor}^{I}=D_\omega e^I.
\]
\end{proposition}

\begin{proof}
Se multiplican las matrices de formas usando el producto exterior.  El
bloque diagonal produce \(F_\omega\); el bloque de traslación produce
\(\dd e+\omega\wedge e=D_\omega e\).
\end{proof}

Curvatura y torsión son dos componentes de un mismo defecto de transporte,
pero cumplen funciones distintas.  La curvatura mide el retorno de marcos;
la torsión mide el cierre de la co-tétrada.

\section{Gauss--Stokes y capacidad areal}

En un complejo celular finito, para toda \(1\)-cocadena \(\Omega\) y toda
\(2\)-cadena \(z\),
\[
 \boxed{\quad
 \langle\delta\Omega,z\rangle
 =\langle\Omega,\partial z\rangle.
 \quad}
\]
La identidad es la definición del coborde como adjunto de la frontera.
Cuando \(z\) suma las caras coherentemente orientadas de una
pseudovariedad, las aristas interiores se cancelan y sólo permanece el
borde geométrico.  La misma igualdad se leyó ya como
vorticidad--circulación y como curvatura--holonomía.

La escala areal también posee un modelo exacto.  En el grafo cúbico
\(N\times N\times N\), todo corte entre dos caras opuestas contiene al menos
\(N^2\) aristas y una capa transversal alcanza la cota.  Para \(N=9\),
\[
 9^3=729=27^2.
\]
Esta igualdad cardinal relaciona volumen y corte mínimo; no identifica los
grupos \((\mathbb Z/9)^3\) y \((\mathbb Z/27)^2\), que tienen exponentes
distintos.

\section{La acción Palatini--Cartan--Holst}

Sea \(M\) una variedad lisa orientada de dimensión cuatro, con co-tétrada no
degenerada \(e^I\), conexión de Lorentz
\(\omega^{IJ}=-\omega^{JI}\), y condiciones de borde que anulen el término
variacional de frontera.  Escribamos
\[
 \Sigma^{IJ}=e^I\wedge e^J,\qquad
 (\star X)^{IJ}=\frac12\epsilon^{IJ}{}_{KL}X^{KL},
 \qquad
 G_{\rm int}=G_{\rm int}(\theta_G),\qquad
 \kappa_{\rm int}=\frac{8\pi G_{\rm int}}{c_{\rm int}^{4}}.
\]
Estas dos últimas cantidades no se definen mediante la acción: proceden de
la construcción constitutiva de la \cref{eq:G-interno}. El
\cref{thm:homogeneidad-pch} demuestra antes que la combinación siguiente
pertenece a la recta de acción.
La acción de primer orden es
\[
\boxed{
 S[e,\omega,\Psi]
 =\frac1{2\kappa_{\rm int}c_{\rm int}}\int_M
 \Sigma_{IJ}\wedge
 \left(\star F^{IJ}+\frac1\gamma F^{IJ}\right)
 -\frac{\Lambda_{\rm int}}{\kappa_{\rm int}c_{\rm int}}
  \int_M\operatorname{vol}_e
 +S_{\rm m}[e,\omega,\Psi].}
\]
El parámetro interno construido en el capítulo anterior especializa
\[
 \gamma=\Gamma_{6\to8}.
\]
La sustitución se efectúa después de obtener
\(\Gamma_{6\to8}\) desde \(A,C^\ast,q_\pm\) y la inclusión
hexada--octada; la acción no se usa para redefinir esas entradas.

\begin{theorem}[Ecuación de Cartan--Holst]
\label{thm:cartan-holst-rev6}
Si la corriente de espín se normaliza mediante
\[
 \delta_\omega S_{\rm m}
 =-\frac12\int_M\delta\omega^{IJ}\wedge\sigma_{IJ},
\]
la variación respecto de \(\omega\) produce
\[
 \boxed{\quad
 D_\omega
 \left[(\star+\gamma^{-1}I)\Sigma^{IJ}\right]
 =\kappa_{\rm int}c_{\rm int}\,\sigma^{IJ}.
 \quad}
\]
\end{theorem}

\begin{proof}
Se usa \(\delta F^{IJ}=D_\omega\delta\omega^{IJ}\), se integra por partes y
se iguala a cero el coeficiente de la variación arbitraria.  El término
cosmológico no depende de \(\omega\).
\end{proof}

En firma lorentziana, \(\star^2=-I\) sobre bivectores.  Para
\(\gamma\in\mathbb R\setminus\{0\}\),
\[
 (\star+\gamma^{-1}I)^{-1}
 =\frac{\gamma^2}{1+\gamma^2}
  (-\star+\gamma^{-1}I).
\]
Por ello, si la corriente de espín se anula,
\[
 T_{\rm tor}^I=0,\qquad
 \omega=\mathring\omega(e).
\]
Con materia espinorial, la torsión se determina algebraicamente por la
corriente.  No constituye, en la acción mínima, un grado propagante
independiente.

\section{Ecuación métrica, espín y balance geométrico}

Para fijar sin ambigüedad la normalización, definimos las fuentes medidas en
la recta de acción por
\[
 \delta_g S_{\rm m}
 =-\frac12\int_M\operatorname{vol}_e\,
 \mathcal T^{S}_{\mu\nu}\,\delta g^{\mu\nu}.
\]
Cuando las coordenadas de la tétrada tienen dimensión de longitud, escribimos
\[
 T^{\rm phys}_{\mu\nu}:=c_{\rm int}\mathcal T^{S}_{\mu\nu},
 \qquad
 U^{\rm phys,spin}_{\mu\nu}
 :=c_{\rm int}\mathcal U^{S,\rm spin}_{\mu\nu}.
\]
Tras eliminar algebraicamente la conexión, la ecuación métrica adopta las
dos formas equivalentes
\[
 \boxed{
 G_{\mu\nu}+\Lambda_{\rm int}g_{\mu\nu}
 =\kappa_{\rm int}c_{\rm int}
 \bigl(\mathcal T_{\mu\nu}^{S,\rm LC}
       +\mathcal U_{\mu\nu}^{S,\rm spin}\bigr)
 =\kappa_{\rm int}
 \bigl(T_{\mu\nu}^{\rm phys,LC}
       +U_{\mu\nu}^{\rm phys,spin}\bigr).}
\]
La contribución de espín es cuadrática en la corriente correspondiente y
depende del acoplamiento de materia; su coeficiente no se fija sin declarar
la acción fermiónica y las convenciones de dualidad.

Las identidades geométricas
\[
 D_\omega T_{\rm tor}^{I}=F^{I}{}_{J}\wedge e^{J},
 \qquad
 D_\omega F^{IJ}=0
\]
implican, después de eliminar la torsión,
\[
 \nabla^\mu\bigl(T_{\mu\nu}^{\rm phys,LC}
             +U_{\mu\nu}^{\rm phys,spin}\bigr)=0.
\]
Este balance de Bianchi--Noether es una condición de consistencia del
transductor gravitatorio: ningún término adicional puede incorporarse sin
proceder de una acción compatible o satisfacer el mismo balance.

\section{Correspondencia areal–volumétrica \(\pi/\varphi^2\)}

El calendario \((\Sigma,\Pi,\Pi)\) produce la matriz de incidencia
\[
 F_{\rm av}=\begin{pmatrix}0&1\\1&1\end{pmatrix}.
\]
Su medida de Parry da
\(\mu_{\rm ar}/\mu_{\rm vol}=\varphi^{-2}\), y el lector regional marca una
sola vez el cierre \(\pi\).  El retorno marcado satisface
\[
 \Theta_n
 =\pi\frac{F_{n-1}}{F_{n+1}}
 \longrightarrow
 \frac{\pi}{\varphi^2}.
\]
Esta razón no es otra versión de \(G\). Es la correspondencia adimensional entre
lectura areal y lectura volumétrica.  Bajo una misma amplitud positiva
\(\mathcal M\),
\[
 \mathcal I_{\rm av}=\mu_{\rm vol}\mathcal M,\qquad
 \mathcal G_{\rm av}=\pi\mu_{\rm ar}\mathcal M
\]
dan exactamente
\[
 \frac{\mathcal G_{\rm av}}{\mathcal I_{\rm av}}
 =\frac{\pi}{\varphi^2}.
\]
En la hoja plana, la restricción diagonal es
\(\mathcal G_{\rm tor}=\mathcal I_{\rm tor}\).  La equivalencia local y la
ambivalencia de frontera son, por tanto, dos regímenes del lector, no dos
valores incompatibles de una masa.

\section{Tipado de \(G\) como transductor gravitatorio y orden causal de sus lectores}

El tipo y la homogeneidad de \(G\) pertenecen a la construcción
constitutiva previa. La \cref{eq:G-interno} fija una sola vez la familia
covariante \(G_{\rm int}(\theta_G)\), y el
\cref{thm:homogeneidad-pch} demuestra que su inserción en la acción posee
dimensión de acción. Este capítulo no redefine esa familia: utiliza su
sección junto con la holonomía ya construida. La causalidad no es lineal,
sino convergente:
\[
\begin{aligned}
 \mathcal B_{\rm geom}&:
 \text{retorno y cociclo}\to\text{representación de Cartan}
 \to(e,\omega),\\
 \mathcal B_{\rm dim}&:
 (\hbar_{\rm int},c_{\rm int},t_0,\theta_G)
 \to G_{\rm int}\to\kappa_{\rm int},
\end{aligned}
\]
\[
 (\mathcal B_{\rm geom},\mathcal B_{\rm dim})
 \to S_{\rm PCH}
 \to\text{torsión algebraica, ecuación métrica y balance}.
\]
La comparación metrológica pertenece al final y no selecciona ninguna de
las dos ramas.

\Needspace{22\baselineskip}
\section{Origen colectivo del modo gravitatorio en la holonomía de rutas}

La acción anterior tiene como objetos básicos la co-tétrada y la conexión.
La torsión de Cartan es algebraica; la curvatura posee el sector propagante que
contiene las ondas gravitatorias.  Por ello el formalismo no necesita
postular un gravitón elemental para definir la interacción.  Una eventual
cuantización de las perturbaciones puede producir un modo colectivo de
helicidades \(\pm2\); ese modo sería una excitación del transporte
geométrico, no un dato añadido al nivel de APP, TRIT o TPK.

Esta distinción también excluye una confusión frecuente.  La existencia de
ondas clásicas demuestra propagación de curvatura; no demuestra por sí sola
un cuanto fundamental.  Recíprocamente, describir el campo mediante
geometría no elimina sus grados radiativos.  HMT--MD sitúa la pregunta en el
lugar correcto: primero se construye la memoria y su holonomía; después se
estudia el espectro de sus excitaciones.
 \end{HMTScientificOwnerSubordinated}

\paragraph{Comparación tipada del cuanto de área con energía, entropía e información de horizonte.} La acción de primer orden permite comparar el cuanto areal con energía, entropía e información de horizonte sin fundir esos observables.
 \endgroup

\chapter{Confluencia de horizonte, área, información, entropía y acción}
\begingroup
\renewcommand{\chapter}[2][]{}%
\chapter{Confluencia horizonte–área–información–entropía–acción}
\label{chap:parte-iv-69}

El parche de horizonte coordina energía, temperatura, entropía y acción mediante identidades tipadas. La celda de media acción y el cierre cronológico se presentan con sus controles de Schwarzschild y de Sitter, sin transformar una interfaz en definición retroactiva.

\section{Energ\'ia de horizonte y celda de media acci\'on}
\label{rm:chap:horizonte-media-accion}

\subsection{Convenciones geom\'etricas del parche de Sitter}

Consideremos un espacio de Sitter cuadridimensional con radio areal
\(R>0\), velocidad causal \(c\) y constante gravitatoria \(G\). Adoptamos
la convenci\'on
\begin{equation}
 H_R=\frac cR,
 \qquad
 \Lambda_R=\frac3{R^2},
 \qquad
 \ell_P^2=\frac{G\hbar}{c^3}.
 \label{rm:eq:horizon-ds-conventions}
\end{equation}
La densidad de energ\'ia asociada al t\'ermino cosmol\'ogico es
\begin{equation}
 \varepsilon_\Lambda
 :=\frac{\Lambda_Rc^4}{8\pi G}.
 \label{rm:eq:horizon-vacuum-density}
\end{equation}
La secci\'on espacial plana del parche contiene, hasta el radio \(R\), el
volumen
\begin{equation}
 V_R=\frac{4\pi}{3}R^3.
 \label{rm:eq:horizon-flat-volume}
\end{equation}
Estas convenciones fijan positivamente la energ\'ia cuasilocal usada en el
cap\'itulo. La elecci\'on de otro signo para la carga de Killing de Sitter
pertenece a una construcci\'on distinta y debe traducirse antes de comparar
f\'ormulas.

\subsection{Energ\'ia de vac\'io y energ\'ia de Misner--Sharp}

En una geometr\'ia esf\'ericamente sim\'etrica, sea
\begin{equation}
 \dd s^2=g_{ab}(x)\dd x^a\dd x^b
 +\mathcal R(x)^2\dd\Omega_2^2.
 \label{rm:eq:horizon-spherical-metric}
\end{equation}
La energ\'ia de Misner--Sharp \cite{rm:MisnerSharp1964}, en la convenci\'on positiva adoptada, se
define por
\begin{equation}
 E_{\rm MS}(\mathcal R)
 :=\frac{c^4\mathcal R}{2G}
 \left(1-g^{ab}\nabla_a\mathcal R\nabla_b\mathcal R\right).
 \label{rm:eq:horizon-Misner-Sharp}
\end{equation}
Una esfera marginal satisface
\begin{equation}
 g^{ab}\nabla_a\mathcal R\nabla_b\mathcal R=0.
 \label{rm:eq:horizon-marginal-sphere}
\end{equation}

\begin{theorem}[Confluencia volum\'etrica y cuasilocal]
\label{rm:thm:horizon-vacuum-MS}
En el horizonte de Sitter \(\mathcal R=R\), la energ\'ia integrada del
vac\'io y la energ\'ia de Misner--Sharp coinciden exactamente:
\begin{equation}
 \boxed{
 E_\Lambda
 :=\varepsilon_\Lambda V_R
 =E_{\rm MS}(R)
 =\frac{c^4R}{2G}.}
 \label{rm:eq:horizon-E-Lambda}
\end{equation}
\end{theorem}

\begin{proof}
Las ecuaciones
\cref{rm:eq:horizon-ds-conventions,rm:eq:horizon-vacuum-density,rm:eq:horizon-flat-volume}
dan
\[
 \varepsilon_\Lambda V_R
 =\frac{3c^4}{8\pi GR^2}\frac{4\pi R^3}{3}
 =\frac{c^4R}{2G}.
\]
La condici\'on marginal \cref{rm:eq:horizon-marginal-sphere} aplicada a
\cref{rm:eq:horizon-Misner-Sharp} produce el mismo resultado.
\end{proof}

La igualdad re\'une dos publicaciones geom\'etricas. La primera integra una
densidad volum\'etrica; la segunda eval\'ua el defecto de atrapamiento de la
esfera areal. Ambas conservan el radio \(R\) y la compliancia gravitatoria
\(G/c^4\).

\subsection{Temperatura, entrop\'ia y producto de horizonte}

La gravedad superficial de Sitter tiene m\'odulo \(c^2/R\). La temperatura
de Gibbons--Hawking \cite{rm:GibbonsHawking1977} y la entrop\'ia de
Bekenstein--Hawking \cite{rm:Bekenstein1973,rm:Hawking1975} son
\begin{align}
 k_BT_H&=\frac{\hbar c}{2\pi R},
 \label{rm:eq:horizon-temperature}\\
 \frac{S_H}{k_B}
 &=\frac{A_H}{4\ell_P^2}
 =\frac{4\pi R^2}{4\ell_P^2}
 =\frac{\pi R^2}{\ell_P^2}.
 \label{rm:eq:horizon-entropy}
\end{align}

\begin{theorem}[Identidad termodin\'amica de Sitter]
\label{rm:thm:horizon-TS}
Para todo radio \(R\) de la familia declarada,
\begin{equation}
 \boxed{
 T_HS_H=\frac{c^4R}{2G}=E_\Lambda.}
 \label{rm:eq:horizon-TS-E}
\end{equation}
\end{theorem}

\begin{proof}
Multiplicar \cref{rm:eq:horizon-temperature,rm:eq:horizon-entropy} da
\[
 T_HS_H
 =\frac{\hbar c}{2\pi Rk_B}
  \frac{k_B\pi R^2}{\ell_P^2}
 =\frac{\hbar cR}{2\ell_P^2}.
\]
La identidad \(\ell_P^2=G\hbar/c^3\) convierte el \'ultimo miembro en
\(c^4R/(2G)\), igual a \cref{rm:eq:horizon-E-Lambda}.
\end{proof}

El factor \(1/2\) posee una descomposici\'on geom\'etrica transparente. El
factor areal \(4\pi\), la normalizaci\'on hologr\'afica por cuatro y la
periodicidad t\'ermica \((2\pi)^{-1}\) satisfacen
\begin{equation}
 \frac{4\pi}{4}\frac1{2\pi}=\frac12.
 \label{rm:eq:horizon-half-factor}
\end{equation}
La mitad es, por tanto, la contracci\'on conjunta entre cuantificaci\'on
areal y periodicidad eucl\'idea.

\subsection{Cuanto completo del horizonte y fuerza de Planck}

Definimos la tensi\'on o fuerza de Planck y el cuanto completo asociado al
radio por
\begin{equation}
 \mathscr F_P:=\frac{c^4}{G},
 \qquad
 \boxed{
 \mathcal Q_H(R):=2T_HS_H
 =\frac{c^4R}{G}=\mathscr F_P\,R.}
 \label{rm:eq:horizon-full-quantum}
\end{equation}
La energ\'ia cuasilocal ocupa media celda del cuanto completo:
\begin{equation}
 \boxed{E_\Lambda=T_HS_H=\frac{\mathcal Q_H(R)}2.}
 \label{rm:eq:horizon-half-quantum-general}
\end{equation}

\begin{proposition}[Compliancia radial]
\label{rm:prop:horizon-compliance}
La derivada del cuanto completo respecto del radio es la fuerza de Planck,
y la derivada de la energ\'ia cuasilocal es su mitad:
\begin{equation}
 \frac{\dd\mathcal Q_H}{\dd R}=\mathscr F_P,
 \qquad
 \frac{\dd E_\Lambda}{\dd R}=\frac{\mathscr F_P}{2}.
 \label{rm:eq:horizon-radial-derivatives}
\end{equation}
La compliancia gravitatoria \(\mathscr C_P=G/c^4\) satisface
\begin{equation}
 R=\mathscr C_P\mathcal Q_H.
 \label{rm:eq:horizon-compliance-law}
\end{equation}
\end{proposition}

\begin{proof}
Las tres identidades son derivadas o inversiones directas de
\cref{rm:eq:horizon-full-quantum}.
\end{proof}

La ecuaci\'on \cref{rm:eq:horizon-compliance-law} expresa la respuesta radial:
el cuanto completo aplica tensi\'on y la geometr\'ia publica un radio mediante
la compliancia universal. El producto \(T_HS_H\) ocupa la mitad cuasilocal
de esa respuesta.

\subsection{Especializaci\'on al cierre cronol\'ogico de 108 estados}

El cierre cronol\'ogico del \cref{rm:ch:cosmo} determina
\begin{equation}
 R_{108}=\ell_P,
 \qquad
 t_P=\frac{\ell_P}{c},
 \qquad
 T_{108}=2\pi t_P,
 \label{rm:eq:horizon-CT108-scales}
\end{equation}
y el cuanto de Planck
\begin{equation}
 E_P=\frac{\hbar}{t_P}
 =k_B\Theta_{\rm clk}\log3.
 \label{rm:eq:horizon-Planck-trit}
\end{equation}

\begin{theorem}[Celda CT108 de media acci\'on]
\label{rm:thm:horizon-half-action}
En la especializaci\'on \(R=R_{108}=\ell_P\),
\begin{align}
 \mathcal Q_H(\ell_P)
 &=\frac{c^4\ell_P}{G}=E_P,
 \label{rm:eq:horizon-QH-EP}\\
 \boxed{
 E_\Lambda=T_HS_H
 =\frac{E_P}{2}
 =\frac{k_B\Theta_{\rm clk}\log3}{2},}
 \label{rm:eq:horizon-CT108-half-energy}\\
 \boxed{E_\Lambda t_P=\frac{\hbar}{2},}
 \qquad
 \boxed{E_\Lambda T_{108}=\frac h2.}
 \label{rm:eq:horizon-half-actions}
\end{align}
\end{theorem}

\begin{proof}
La primera igualdad se obtiene sustituyendo
\(\ell_P=\sqrt{G\hbar/c^3}\) y
\(E_P=\sqrt{\hbar c^5/G}\). La
\cref{rm:eq:horizon-half-quantum-general} produce
\(E_\Lambda=E_P/2\), y
\cref{rm:eq:horizon-Planck-trit} aporta la carta informacional. Finalmente,
\[
 E_\Lambda t_P=\frac{E_Pt_P}{2}=\frac\hbar2
\]
y, como \(T_{108}=2\pi t_P\) y \(h=2\pi\hbar\),
\[
 E_\Lambda T_{108}
 =\frac{E_P}{2}(2\pi t_P)=\pi\hbar=\frac h2.
\]
\end{proof}

La primera igualdad de \cref{rm:eq:horizon-half-actions} publica media acci\'on
reducida sobre el tiempo de Planck; la segunda publica media acci\'on completa
sobre el periodo CT108. Ambas expresan una misma celda en las cartas angular
y circular.

La temperatura y la entrop\'ia quedan tambi\'en determinadas:
\begin{equation}
 k_BT_H=\frac{E_P}{2\pi},
 \qquad
 \boxed{\frac{S_H}{k_B}=\pi.}
 \label{rm:eq:horizon-CT108-TS}
\end{equation}
La temperatura convierte la frecuencia de horizonte en energ\'ia; la
entrop\'ia convierte el \'area de la esfera de Planck en informaci\'on. Su
producto recupera el mismo medio cuanto.

\subsection{Peso efectivo del horizonte y exponencial hiperb\'olica}

Definimos el peso termodin\'amico efectivo
\begin{equation}
 \Omega_{\rm eff}(R):=\exp\!\left(\frac{S_H(R)}{k_B}\right).
 \label{rm:eq:horizon-effective-weight}
\end{equation}
En el radio CT108,
\begin{equation}
 \boxed{\Omega_{\rm eff}(\ell_P)=\ee^\pi.}
 \label{rm:eq:horizon-effective-e-pi}
\end{equation}
El peso \(\Omega_{\rm eff}\) es una magnitud positiva definida por la
entrop\'ia. Su valor real expresa multiplicidad efectiva; una interpretaci\'on
como cardinal entero de microestados requiere una discretizaci\'on adicional.

El operador de incidencia areal--volum\'etrica es
\begin{equation}
 F_{\rm av}=\begin{pmatrix}0&1\\1&1\end{pmatrix},
 \qquad
 J_\varphi:=\frac{2F_{\rm av}^2-3I}{\sqrt5}.
 \label{rm:eq:horizon-J-phi}
\end{equation}
Un c\'alculo directo da
\begin{equation}
 J_\varphi^2=I,
 \qquad
 \Spec(J_\varphi)=\{+1,-1\}.
 \label{rm:eq:horizon-J-involution}
\end{equation}
El flujo hiperb\'olico en el tiempo de clausura \(\pi\) es
\begin{equation}
 \mathcal E_\pi
 :=\ee^{\pi J_\varphi}
 =\cosh\pi\,I+\sinh\pi\,J_\varphi,
 \label{rm:eq:horizon-hyperbolic-flow}
\end{equation}
y posee espectro
\begin{equation}
 \boxed{
 \Spec(\mathcal E_\pi)=\{\ee^\pi,\ee^{-\pi}\},
 \qquad
 \rho(\mathcal E_\pi)=\ee^\pi.}
 \label{rm:eq:horizon-hyperbolic-spectrum}
\end{equation}

\begin{corollary}[Confluencia escalar horizonte--Perron]
\label{rm:cor:horizon-Perron-scalar}
En el cierre CT108,
\begin{equation}
 \boxed{
 \frac{S_H}{k_B}
 =\log\Omega_{\rm eff}
 =\log\rho(\mathcal E_\pi)
 =\pi.}
 \label{rm:eq:horizon-Perron-log}
\end{equation}
\end{corollary}

\begin{proof}
Las ecuaciones
\cref{rm:eq:horizon-CT108-TS,rm:eq:horizon-effective-e-pi,rm:eq:horizon-hyperbolic-spectrum}
dan los tres miembros.
\end{proof}

La identidad escalar es exacta y conserva dos genealog\'ias distintas. El
horizonte produce \(\ee^\pi\) mediante exponenciaci\'on de la entrop\'ia
areal; la din\'amica de transferencia, construida con anterioridad, reconoce
el mismo escalar como autovalor expansivo de una involuci\'on hiperb\'olica.
La igualdad terminal fija una condici\'on de borde
para el entrelazador y preserva la distinci\'on de dominios.

\subsection{Consistencia del transductor gravitatorio con el horizonte de Schwarzschild}

Un horizonte de Schwarzschild con el mismo radio areal satisface
\begin{equation}
 R_H=\frac{2GE_H}{c^4},
 \qquad
 E_H=\frac{c^4R_H}{2G}.
 \label{rm:eq:horizon-Schwarzschild-energy}
\end{equation}
La entrop\'ia areal coincide con \cref{rm:eq:horizon-entropy}, mientras la
temperatura es
\begin{equation}
 k_BT_{\rm Sch}=\frac{\hbar c}{4\pi R_H}
 =\frac12k_BT_H\big|_{R=R_H}.
 \label{rm:eq:horizon-Schwarzschild-temperature}
\end{equation}
En consecuencia,
\begin{equation}
 E_H=2T_{\rm Sch}S_H,
 \qquad
 T_H^{\rm dS}=2T_{\rm Sch}
 \label{rm:eq:horizon-Schwarzschild-Smarr}
\end{equation}
para un mismo radio. En \(R_H=\ell_P\),
\begin{equation}
 E_H=\frac{E_P}{2},
 \qquad
 T_{\rm Sch}S_H=\frac{E_P}{4}.
 \label{rm:eq:horizon-Schwarzschild-Planck}
\end{equation}
El contraste separa las dos geometr\'ias: comparten energ\'ia cuasilocal,
\'area y entrop\'ia, y sus gravedades superficiales producen temperaturas
en raz\'on dos.

Con \(\beta=(k_BT_{\rm Sch})^{-1}\), la acci\'on eucl\'idea satisface
\begin{equation}
 \frac{I_E}{\hbar}
 =\frac{E_H}{k_BT_{\rm Sch}}-\frac{S_H}{k_B}
 =2\pi-\pi=\pi
 \label{rm:eq:horizon-Schwarzschild-action}
\end{equation}
en el radio de Planck. El mismo \(\pi\) aparece como entrop\'ia reducida,
acci\'on eucl\'idea reducida, tiempo del flujo hiperb\'olico y logaritmo del
peso efectivo.

\subsection{Entrelazador horizonte–Perron: compatibilidad espectral entre energía de frontera y memoria de retorno}

La coincidencia
\(\Omega_{\rm eff}=\rho(\mathcal E_\pi)\) fija un dato espectral y deja
abierta la relaci\'on entre los objetos completos. Denotemos por
\(\mathcal A_H^{(m)}\) el \'algebra de observables de horizonte al nivel
non\'adico \(m\), por \(\omega_H^{(m)}\) su estado, y por
\(\mathcal P_\varphi^{(m)}\) el sistema de transferencia de Perron.

\begin{definition}[Entrelazador termodin\'amico--din\'amico]
\label{rm:def:horizon-Perron-intertwiner}
Un entrelazador admisible es una familia de aplicaciones lineales unitales
completamente positivas (UCP)
\begin{equation}
 \mathfrak I_m:\mathcal A_H^{(m)}\longrightarrow
 \mathcal A_\varphi^{(m)}
 \label{rm:eq:horizon-intertwiner}
\end{equation}
que satisface:
\begin{enumerate}[label=\roman*)]
 \item preservaci\'on del estado, adem\'as de la unidad y de la positividad
       completa ya incorporadas en la condici\'on UCP;
 \item compatibilidad con el refinamiento non\'adico;
 \item transporte de la involuci\'on interior--exterior hacia
       \(J_\varphi\mapsto-J_\varphi\);
 \item conmutaci\'on entre el semigrupo modular de horizonte y el semigrupo
       de transferencia sobre el subespacio declarado;
 \item igualdad del radio espectral en el tiempo \(\pi\), dada por
       \cref{rm:eq:horizon-Perron-log}.
\end{enumerate}
\end{definition}

La quinta condici\'on es necesaria y escalar; las cuatro primeras aportan el
contenido algebraico que una coincidencia num\'erica no contiene. Un
entrelazador construido de este modo identificar\'ia la multiplicidad
termodin\'amica efectiva con una din\'amica de expansi\'on y contracci\'on
conservando orientaci\'on, estado y refinamiento.

\begin{proofstatusbox}
Las identidades
\(E_\Lambda=E_{\rm MS}=T_HS_H=c^4R/(2G)\), la definici\'on
\(\mathcal Q_H=\mathscr F_P\,R\), la especializaci\'on CT108, las dos celdas
de media acci\'on y la igualdad
\(S_H/k_B=\log\rho(\ee^{\pi J_\varphi})=\pi\) son resultados exactos bajo
las convenciones y la realizaci\'on de Sitter declaradas. La identificaci\'on
de \(\Omega_{\rm eff}\) con un cardinal microsc\'opico y el entrelazador
\(\mathfrak I_m\) pertenecen a una interfaz f\'isica delimitada.
\end{proofstatusbox}

\enlargethispage{4\baselineskip}
\begin{falsifierbox}
El cierre geom\'etrico falla si \(\Lambda_R\ne3/R^2\), si la esfera de
radio \(R\) incumple la condici\'on marginal o si la energ\'ia integrada difiere de
\(E_{\rm MS}(R)\). El cierre termodin\'amico falla si la temperatura o la
ley de \'area producen \(T_HS_H\ne c^4R/(2G)\). La especializaci\'on CT108
falla si \(R_{108}\ne\ell_P\), \(T_{108}\ne2\pi t_P\) o
\(E_P\ne k_B\Theta_{\rm clk}\log3\). La interfaz Perron queda refutada por
una obstrucci\'on de positividad, estado, orientaci\'on o refinamiento, aun
cuando el autovalor \(\ee^\pi\) coincida escalarmente.
\end{falsifierbox}
 
\paragraph{Factorización machiana de la memoria de transporte desde el cierre global y la respuesta local.} La relación entre cierre global y respuesta local conduce al principio de Mach formulado sobre la memoria de transporte.
 \endgroup

\let\HMTMonographChapterInput\relax
  \part{Incidencia fronteriza y estabilizadores excepcionales}

\chapter{Entrelazador dodecafásico--Witt}
\section{Construcción del entrelazador de incidencia punto--hexada y transporte de la acción residual de \texorpdfstring{\(S_3\)}{S3}}
\label{chap:entrelazador-accion}
\index{entrelazador!dodecafase--Witt}
\index{Witt!módulo de incidencia}
\index{supervivencia!proyector de rango tres}

La aparición reiterada del cociente \(3/11\) en la pantalla dodecafásica,
en la polaridad de una estrella de Witt y en la filtración nonádica plantea
una pregunta más fuerte que la igualdad aritmética de tres fracciones. La
cuestión correcta es operatorial: determinar si esas tres apariciones son
trazas del mismo sector positivo visto en realizaciones diferentes y, si lo
son, construir los mapas que las entrelazan sin introducir un dato
metrológico posterior.

La respuesta posee cuatro niveles. Primero, la incidencia punto--hexada
construye una isometría canónica entre el módulo dodecafásico de dimensión
once y un autoespacio de dimensión once del diseño de Witt. Segundo, la
estrella de Witt actúa sobre ese mismo módulo, pero lo hace como involución
firmada; su \(3/11\) es un índice, no la traza de un proyector positivo.
Tercero, el sector orientado de supervivencia aparece en \(t=7\), donde un
único espacio de once ramas contiene un sector de rango tres. Cuarto, la
línea de orientación de ese triplete, junto con los testigos internos
\(B_K\) y \(u_K\), selecciona de manera variacional un único normalizador
\(S_3\), fija su orientación y produce por promedio de Reynolds y
descomposición polar el entrelazamiento isométrico sector--dodecafase. El
paso dodecafase--Witt es una segunda isometría de incidencia; no se presupone
que ambas flechas realicen una única representación de \(S_3\).
Los promedios de Reynolds, las descomposiciones en representaciones y el
factor polar utilizados en esta sección son los estándar
\cite{Serre1977Representations,HornJohnson2013}.

\begin{figure}[htbp]
\centering
\begin{tikzpicture}[
  font=\footnotesize,
  box/.style={draw,rounded corners=2mm,align=center,minimum width=2.35cm,minimum height=.98cm},
  arr/.style={-{Stealth[length=2.2mm]},thick},
  link/.style={<->,thick,dashed}
]
  \node[box,draw=hmtblue,fill=hmtlightblue] (seed) at (0,2.2)
    {semillas\\\(104\,976\)};
  \node[box,draw=hmtblue,fill=hmtlightblue] (u) at (3,2.2)
    {emisiones\\\(\mathcal U_6,\ |\mathcal U_6|=468\)};
  \node[box,draw=hmtgold,fill=hmtlightgold] (ppi) at (6,2.2)
    {microfibra\\\(\Pi_\pi^{(468)},\ \operatorname{rank}=5\)};
  \node[box,draw=hmtred,fill=hmtlightred] (d) at (9,2.2)
    {defecto positivo\\\(D_{\rm switch}\)};
  \node[box,draw=hmtviolet,fill=hmtlightviolet] (lam) at (12,2.2)
    {capacidad residual\\\(\lambda_*=1-\delta_*\)};

  \draw[arr,hmtblue] (seed)--node[above=16pt,font=\scriptsize,fill=white,inner sep=1pt]{\(F\)} (u);
  \draw[arr,hmtgold] (u)--node[above=16pt,font=\scriptsize,fill=white,inner sep=1pt]{\(\tau_{468}\)} (ppi);
  \draw[arr,hmtred] (ppi)--node[above=16pt,font=\scriptsize,fill=white,inner sep=1pt]{\(D_{\rm switch}\)} (d);
  \draw[arr,hmtviolet] (d)--node[above=16pt,font=\scriptsize,fill=white,inner sep=1pt]{\(\tau+\mathrm{CR}\)} (lam);

  \node[box,draw=hmtgold,fill=hmtlightgold] (g) at (1.5,-.2)
    {sector orientado \(t=7\)\\\(L_{\rm or}\otimes(H_G,P_G)\)};
  \node[box,draw=hmtgreen,fill=hmtlightgreen] (v) at (5,-.2)
    {dodecafase\\\(V_{11},P_3\)};
  \node[box,draw=hmtgreen,fill=hmtlightgreen] (e) at (8.5,-.2)
    {incidencia Witt\\\(E_{30},Q_3\)};
  \node[box,draw=hmtgray,fill=hmtpaper] (t) at (12,-.2)
    {lector espectral\\\(\mathcal L_{\beta,*}\to T_{\rm band}\)};

  \draw[arr,hmtgold] (g)--node[above=16pt,font=\scriptsize,fill=white,inner sep=1pt]{\(W_{GK}\)} (v);
  \draw[arr,hmtgreen] (v)--node[above=16pt,font=\scriptsize,fill=white,inner sep=1pt]{\(U_W=I/6\)} (e);
  \draw[arr,hmtgray] (lam)--(t);
  \draw[arr,hmtred] (v) to[bend left=20]
    node[below=3pt,font=\scriptsize,fill=white,inner sep=1pt]{\((3/11)\Delta_4\)} (d);
\end{tikzpicture}
 \caption{Arquitectura del cierre pre-dimensional. La fila superior construye
la medida y la acción desde el cociente de emisiones; la inferior transporta
el sector orientado \(t=7\) hacia la dodecafase y, después, la dodecafase hacia
Witt. Las flechas continuas son isometrías de proyectores demostradas. La
primera se fija por la selección variacional del
\cref{thm:identificacion-isometrica-tres-proyectores}; para la segunda no se afirma
equivarianza respecto del \(H_*\) seleccionado.}
\label{fig:entrelazador-accion}
\end{figure}

\section{Tipado de los objetos asociados a \texorpdfstring{\(3/11\)}{3/11}}

Sea
\[
 V_{12}=\mathbb R^{12},
 \qquad
 V_{11}=\mathbf 1^\perp
 =\left\{v\in\mathbb R^{12}:\sum_{p=1}^{12}v_p=0\right\}.
\]
La acción icosaédrica declarada sobre las doce posiciones se descompone como
\[
 V_{12}=V_1\oplus V_3\oplus V_{3'}\oplus V_5.
\]
Denotemos por \(P_3\) el proyector ortogonal sobre el primer sumando
tridimensional. Entonces
\[
 P_3^2=P_3=P_3^*,
 \qquad
 \operatorname{rank}P_3=3,
 \qquad
 \tau_{11}(P_3)
 =\frac1{11}\operatorname{Tr}_{V_{11}}P_3
 =\frac3{11}.
\]

Hay tres objetos próximos y un cuarto objeto de tipo distinto:
\begin{enumerate}
  \item la incidencia punto--hexada transporta \(P_3\) a un proyector
  \(Q_3\) de rango tres sobre un módulo de Witt;
  \item la supervivencia en \(t=7\) distingue tres ramas entre once y
  produce un proyector \(P_G\) con traza normalizada \(3/11\);
  \item un entrelazador orientado identifica
  \(P_3V_{11}\) con la realización orientada de \(P_GH_G\);
  \item una estrella de Witt tiene siete modos pares y cuatro impares, por
  lo que su traza normalizada también vale \(3/11\), pero contiene espectro
  negativo y no es un proyector.
\end{enumerate}
La distinción entre proyector e involución será esencial para que la ley de
acción de la sección siguiente sea positiva.

\section{La matriz de incidencia punto--hexada}

Sea \(\mathcal H\) el conjunto de las \(132\) hexadas del diseño
\[
 W_{12}=S(5,6,12).
\]
Definimos
\[
 I:\mathbb R^{12}\longrightarrow\mathbb R^{\mathcal H},
 \qquad
 (Iv)(H)=\sum_{p\in H}v_p,
\]
o, equivalentemente,
\[
 I_{H,p}=\mathbf 1_{\{p\in H\}}.
\]
Este operador sólo registra incidencia. No contiene \(\alpha\), unidades,
masas, CKM, Barbero ni ningún objetivo físico.

\begin{lemma}[Identidad métrica de la incidencia]
\label{lem:incidencia-metrica-witt}
Se cumple
\[
 \boxed{I^*I=36I_{12}+30J_{12}.}
\]
Por tanto, para todo \(v\in V_{11}\),
\[
 \|Iv\|^2=36\|v\|^2.
\]
\end{lemma}

\begin{proof}
En un sistema \(S(5,6,12)\), cada punto pertenece a
\[
 r=\frac{\binom{11}{4}}{\binom54}=66
\]
hexadas, y cada par de puntos pertenece a
\[
 \lambda_2=\frac{\binom{10}{3}}{\binom43}=30
\]
hexadas. Las entradas diagonales de \(I^*I\) valen \(66\) y las no
diagonales valen \(30\). Luego
\[
 I^*I=(66-30)I_{12}+30J_{12}.
\]
Como \(J_{12}\) se anula sobre \(\mathbf1^\perp\), se obtiene la segunda
identidad.
\end{proof}

Para identificar la rama espectral seleccionada por \(I\), introducimos
\[
 G_{H,H'}=\binom{|H\cap H'|}{4}.
\]

\begin{lemma}[Identidad espectral de la incidencia]
\label{lem:incidencia-espectral-witt}
Con \(J_{132,12}\) la matriz rectangular de unos,
\[
 \boxed{GI=30I+15J_{132,12}.}
\]
En particular, si \(v\in V_{11}\), entonces
\[
 G(Iv)=30Iv.
\]
\end{lemma}

\begin{proof}
La entrada \((H,p)\) de \(GI\) es
\[
 \sum_{H'\ni p}\binom{|H\cap H'|}{4}.
\]
El censo exhaustivo de \(W_{12}\) da \(45\) cuando \(p\in H\) y \(15\)
cuando \(p\notin H\). Por consiguiente, la entrada vale
\(30I_{H,p}+15\). La reconstrucción de las \(132\) hexadas desde el código
de Golay ternario extendido comprueba las \(132\cdot12\) entradas como
enteros.
\end{proof}

Definamos
\[
 E_{30}^{\rm inc}:=I(V_{11})\subseteq\ker(G-30I).
\]
Su dimensión es once por el \cref{lem:incidencia-metrica-witt}. El cálculo
módulo \(1009\) obtiene
\[
 \operatorname{rank}_{\mathbb F_{1009}}(G-30I)=121.
\]
Como la reducción módulo \(1009\) da una cota inferior para el rango
racional, se sigue
\(\dim_{\mathbb Q}\ker(G-30I)\le11\). La inclusión anterior proporciona
once vectores linealmente independientes en ese núcleo. Por tanto,
\[
 E_{30}^{\rm inc}=\ker(G-30I).
\]
En adelante escribiremos \(E_{30}\).

\begin{theorem}[Entrelazador canónico dodecafase--Witt]
\label{thm:entrelazador-dodecafase-witt}
El operador
\[
 \boxed{
 U_W=\frac16 I\big|_{V_{11}}:V_{11}\longrightarrow E_{30}
 }
\]
es una isometría sobreyectiva. Para todo
\(g\in\operatorname{Aut}(W_{12})\cong M_{12}\),
\[
 U_W\rho_{12}(g)=\rho_{\mathcal H}(g)U_W.
\]
Entre los kernels punto--hexada equivariantes, la condición de isometría y
el signo positivo sobre las incidencias determinan \(U_W\) de manera única.
\end{theorem}

\begin{proof}
La isometría se sigue de \(I^*I=36I\) sobre \(V_{11}\). Su imagen tiene
dimensión once y está contenida en \(E_{30}\), también de dimensión once;
por tanto es sobreyectiva. La equivariancia traduce
\[
 p\in H\quad\Longleftrightarrow\quad g(p)\in g(H).
\]

Para la unicidad, \(M_{12}\) es transitivo sobre las banderas
\((p,H)\), con \(p\in H\), y sobre las antibanderas. Todo núcleo
equivariante tiene dos valores, \(a\) sobre banderas y \(b\) sobre
antibanderas:
\[
 T=aI+b(J-I)=(a-b)I+bJ.
\]
El término \(J\) se anula sobre \(V_{11}\), de modo que allí \(T\) es un
múltiplo de \(I\). La isometría fuerza \(|a-b|=1/6\) y la normalización
positiva selecciona \(a-b=1/6\).
\end{proof}

\section{Transporte positivo del proyector trítico}

El entrelazador permite definir
\[
 \boxed{Q_3=U_WP_3U_W^*.}
\]
Se obtiene
\[
 Q_3^2=Q_3=Q_3^*,
 \qquad
 \operatorname{rank}Q_3=3,
\]
y el diagrama
\[
\begin{CD}
V_{11} @>{U_W}>> E_{30}\\
@V{P_3}VV @VV{Q_3}V\\
V_{11} @>{U_W}>> E_{30}
\end{CD}
\]
conmuta. En particular,
\[
 \boxed{\tau_{V_{11}}(P_3)=\tau_{E_{30}}(Q_3)=\frac3{11}.}
\]

La canonicidad de \(U_W\) pertenece al diseño etiquetado y a su acción de
automorfismos. Debe distinguirse de la pantalla \(A_5/C_5\) usada para
definir \(P_3\): en la realización etiquetada construida, sólo la identidad
de esa acción concreta de \(A_5\) preserva el conjunto fijo de hexadas. Por
ello \(U_W\) transporta canónicamente el \(P_3\) ya dado hacia
\(Q_3\), pero no convierte por sí solo la acción \(A_5/C_5\) en una
subacción de \(M_{12}\). El puente con el sector orientado se construirá por separado
desde los datos marcados \(K\) y \(B_K\).

\section{Índice firmado de la estrella de Witt}

Sea \(B\) una tetrada. Las cuatro hexadas que la contienen se escriben
\[
 H_i=B\sqcup P_i,\qquad |P_i|=2,
\]
donde los cuatro pétalos \(P_i\) particionan el complemento de \(B\).
La involución de estrella \(R_B\) fija \(B\) e intercambia los dos puntos
de cada pétalo. Su tipo de ciclos sobre doce puntos es \(1^4\,2^4\). En
\(V_{11}\), y mediante \(U_W\) también en \(E_{30}\), su espectro es
\[
 \operatorname{Spec}(R_B)=\{1^{[7]},(-1)^{[4]}\}.
\]
Por ello,
\[
 \boxed{\tau_{11}(R_B)=\frac{7-4}{11}=\frac3{11}.}
\]

\begin{theorem}[Obstrucción proyector--estrella]
\label{thm:no-go-proyector-estrella}
No existe un isomorfismo \(F:V_{11}\to E_{30}\) tal que
\[
 FP_3=R_BF.
\]
\end{theorem}

\begin{proof}
Si existiera, \(P_3\) y \(R_B\) serían semejantes. Sin embargo,
\[
 \operatorname{Spec}(P_3)=\{1^{[3]},0^{[8]}\},
 \qquad
 \operatorname{Spec}(R_B)=\{1^{[7]},(-1)^{[4]}\}.
\]
Sus polinomios mínimos son \(x(x-1)\) y
\((x-1)(x+1)\), respectivamente. La semejanza es imposible.
\end{proof}

La igualdad \(3/11\) posee así dos lecturas sobre el mismo módulo. Para
\(Q_3\) es la masa de un sector positivo de rango tres. Para \(R_B\) es
un índice de paridad. En una acción positiva debe intervenir \(Q_3\); la
estrella conserva orientación y distingue hojas, pero no reemplaza al
proyector.

\section{Reducción homogénea \texorpdfstring{\(11\to3\to1\)}{11-3-1} en \texorpdfstring{\(t=7\)}{t=7}}

La lectura anterior de \(3/11\) en \(t=8\) comparaba tres raíces locales
con once continuaciones de una sola raíz. Numerador y denominador no
pertenecían al mismo espacio. La tabla detallada localiza el objeto correcto
en \(t=7\): hay once ramas locales; las once sobreviven un horizonte, tres
sobreviven dos y una sobrevive tres,
\[
 \boxed{11\longrightarrow3\longrightarrow1.}
\]

\begin{center}
\small
\begin{tabularx}{\textwidth}{cYcc}
\toprule
índice&triple de palabras&\(h=2\)&\(h=3\)\\
\midrule
0&\texttt{200112|212221|110110}&3&0\\
1&\texttt{200121|212212|110110}&3&1\\
2&\texttt{200211|212122|110110}&3&0\\
3&\texttt{202122|210211|110110}&0&0\\
4&\texttt{202212|210121|110110}&0&0\\
5&\texttt{202221|210112|110110}&0&0\\
6&\texttt{210111|202222|110110}&0&0\\
7&\texttt{210222|202111|110110}&0&0\\
8&\texttt{212112|200221|110110}&0&0\\
9&\texttt{212121|200212|110110}&0&0\\
10&\texttt{212211|200122|110110}&0&0\\
\bottomrule
\end{tabularx}
\end{center}

Sea \(\mathcal B_7\) el conjunto de estas once ramas y
\[
 H_G=\ell^2(\mathcal B_7).
\]
Definimos
\[
 P_G\delta_b=
 \begin{cases}
 \delta_b,&b\text{ posee continuación a horizonte dos},\\
 0,&\text{en otro caso}.
 \end{cases}
\]
Entonces
\[
 P_G^2=P_G=P_G^*,
 \qquad
 \operatorname{rank}P_G=3,
 \qquad
 \boxed{\tau_{H_G}(P_G)=\frac3{11}.}
\]
Ésta es la tercera realización positiva.

\section{La acción residual de \texorpdfstring{\(S_3\)}{S3} y las regiones positivas de \texorpdfstring{\(\pi\)}{pi}}

Una permutación de las tres posiciones finales, aplicada simultáneamente a
las dos primeras palabras de cada rama y dejando fija la tercera, preserva
\(\mathcal B_7\). Esta regla define una acción de \(S_3\). Sus órbitas
tienen tamaños
\[
 3+3+1+1+3=11.
\]
El soporte de \(P_G\), formado por los índices \(0,1,2\), es una órbita
completa, por lo que \(P_G\) conmuta con \(S_3\).

El carácter de la representación sobre las once ramas, leído en identidad,
transposición y ciclo de orden tres, es
\[
 \chi_{\mathcal B_7}=(11,5,2)
 =5\chi_{\rm triv}+3\chi_{\rm std}.
\]
Sobre la imagen de \(P_G\),
\[
 \chi_G=(3,1,0)
 =\chi_{\rm triv}+\chi_{\rm std}.
\]

Las tres regiones positivas de la microfibra de \(\pi\) tienen firmas
\[
 339555,\qquad393555,\qquad933555.
\]
Las colas de la primera palabra de las tres ramas supervivientes son
\[
 112,\qquad121,\qquad211.
\]
La sustitución coordenada \(1\mapsto3\), \(2\mapsto9\) da
\[
\begin{array}{c|c|c}
\text{rama}&\text{cola}&\text{región positiva}\\ \hline
0&112&339555\\
1&121&393555\\
2&211&933555.
\end{array}
\]

\begin{theorem}[Puente horizontal--vertical de la microfibra positiva]
\label{thm:puente-pi-puerta-s3}
La regla anterior es una biyección \(S_3\)-equivariante entre las tres
regiones positivas de \(\pi\) y las tres ramas seleccionadas por \(P_G\).
El mapa conserva la posición del único símbolo excepcional.
\end{theorem}

\begin{proof}
Ambos conjuntos son las tres palabras obtenidas al colocar un único símbolo
excepcional en una de tres posiciones. La sustitución se aplica coordenada a
coordenada y, por tanto, conmuta con toda permutación de posiciones. La
posición excepcional recupera unívocamente el elemento de origen.
\end{proof}

\section{Dependencia de la orientación en el entrelazador dodecafásico--Witt}

Para cualquier normalizador \(H=N_{A_5}(C_3)\simeq S_3\) de un subgrupo de
orden tres en \(A_5\), definimos el carácter abstracto del cociente
\[
 \operatorname{sgn}_H:H\longrightarrow H/C_3\simeq C_2
 \longrightarrow\{+1,-1\},
 \qquad
 \operatorname{sgn}_H(h)=
 \begin{cases}
 +1,&h\in C_3,\\
 -1,&h\in H\setminus C_3.
 \end{cases}
\]
Este carácter no es el signo ambiental de una permutación de cinco letras:
todo elemento de \(A_5\) tiene signo ambiental \(+1\). Es el carácter
intrínseco de la reflexión en el cociente \(H/C_3\). Con esta convención, el
carácter del sector icosaédrico restringido a \(H\) es
\[
 \chi_{P_3V_{11}|H}=(3,-1,0)
 =\chi_{\operatorname{sgn}_H}+\chi_{\rm std}.
\]
En cambio,
\[
 \chi_{\operatorname{im}P_G}=(3,1,0)
 =\chi_{\rm triv}+\chi_{\rm std}.
\]

\begin{corollary}[Obstrucción al entrelazador no orientado]
\label{cor:obstruccion-entrelazador-incidencia-p3}
No existe un isomorfismo \(S_3\)-equivariante
\[
 \operatorname{im}P_G\longrightarrow P_3V_{11}
\]
para las acciones no retorcidas.
\end{corollary}

\begin{proof}
Representaciones isomorfas poseen el mismo carácter. Los valores en una
transposición son \(1\) y \(-1\).
\end{proof}

Sea \((e_0,e_1,e_2)\) la base posicional de
\(\mathbb R^3_{\rm perm}\). Definimos
\[
 L_{\rm or}
 =\bigwedge\nolimits^3\mathbb R^3_{\rm perm},
 \qquad
 \ell_{\rm or}=e_0\wedge e_1\wedge e_2
\]
la línea de orientación del triplete, provista del ancla positiva
\(\ell_{\rm or}\). La acción de \(S_3\) sobre esta línea es precisamente su
carácter de signo abstracto. Al tensorizar por ella,
\[
 \chi_{L_{\rm or}\otimes\operatorname{im}P_G}
 =(3,-1,0),
\]
de modo que la obstrucción de caracteres desaparece. La igualdad de
caracteres demuestra la existencia de una clase de isomorfismos, pero aún no
selecciona uno. La selección se obtiene de dos testigos ya presentes en la
dodecafase.

\section{Selección variacional del normalizador \texorpdfstring{\(S_3\)}{S3}}

Sea
\[
 K=(234,543,140,729,659,824,621,58,914,794,146,601)
\]
la carta decimal del sello. El testigo común mínimo de las dos lecturas es
\[
 B_K=H_e\cap P_\pi
 =\{p:K_p\ge729\}
 =\{4,6,9,10\},
\]
en la numeración matemática \(1,\ldots,12\). En la convención equivalente
con índices desde cero, el mismo conjunto es \(\{3,5,8,9\}\). Definimos dos
vectores del sector \(P_3\):
\[
 b_K=P_3\left(\mathbf1_{B_K}-\frac{|B_K|}{12}\mathbf1\right),
 \qquad
 u_K=P_3\left(K-\overline K\,\mathbf1\right).
\]
El primero satisface exactamente
\[
 \|b_K\|^2=1.
\]

El grupo \(A_5\) contiene diez subgrupos \(C_3\). Para cada uno, sea
\[
 H(C_3)=N_{A_5}(C_3)\simeq S_3
\]
y sea el proyector de signo
\[
 e_{{\rm sgn},H}
 =\frac16\sum_{h\in H}\operatorname{sgn}_H(h)\rho(h).
\]
Aquí \(\operatorname{sgn}_H\) es el carácter abstracto de \(H/C_3\)
definido arriba, no el signo ambiental —trivial— de \(A_5\).
La energía de orientación asociada al testigo común es
\[
 \mathcal E_B(H)=\|e_{{\rm sgn},H}b_K\|^2.
\]
El censo exacto en \(\mathbb Q(\sqrt5)\) produce un máximo único:
\[
 \mathcal E_B(H_*)
 =\frac23+\frac{2}{15}\sqrt5.
\]
El segundo nivel es
\[
 \frac13+\frac{2}{15}\sqrt5,
\]
de modo que la brecha exacta es
\[
 \boxed{\mathcal E_B(H_*)-\mathcal E_B(H_{\rm segundo})=\frac13.}
\]
El mismo \(H_*\) es el máximo único al reemplazar \(b_K\) por \(u_K\).
Por tanto, el selector no depende de una sola codificación del testigo.

En esta etiquetación,
\[
 H_*=N_{A_5}\bigl(\langle c_*\rangle\bigr),
\qquad
 c_*=(2\,3\,4).
\]
La regla es natural bajo relabelización simultánea por \(A_5\): al
transportar \(B_K\), \(K\) y \(P_3\), el máximo se conjuga por el mismo
elemento.

\section{Selección de orientación y presentación de \texorpdfstring{\(S_3\)}{S3}}

Dentro de \(H_*\), consideremos
\[
 s(g)=\langle b_K,\rho(g)u_K\rangle.
\]
El máximo entre los dos elementos de orden tres selecciona unívocamente
\[
 c_*=(0,1,3,4,2),
\]
en notación de imágenes de la permutación de cinco puntos, y el máximo entre
las tres involuciones selecciona
\[
 r_*=(1,0,4,3,2).
\]
Se verifica exactamente
\[
 r_*c_*r_*=c_*^{-1}.
\]
La presentación del sector orientado mediante el ciclo de posiciones
\(c=(0\,1\,2)\) y la reflexión \(r=(0\,2)\), que fija la posición central,
determina entonces un único isomorfismo
\[
 \theta:S_3^{\rm or}\xrightarrow{\sim}H_*,
 \qquad
 \theta(c)=c_*,
 \qquad
 \theta(r)=r_*.
\]
Aquí \(S_3^{\rm or}\) denota el grupo de reetiquetados del sector orientado
\(t=7\), generado por \(c\) y \(r\); no introduce una clase operatoria
adicional.
La única rama que alcanza horizonte tres es precisamente la rama central
\(1\); este dato fija cuál de las tres posiciones queda inmóvil bajo la
reflexión.

\section{Promedio de Reynolds y factor polar}

Sea
\[
 W_{\rm or}=L_{\rm or}\otimes\mathbb R^3_{\rm perm}
\]
el triplete del sector orientado. Con el ancla ya fijada, tomamos
\[
 \boxed{w=\ell_{\rm or}\otimes e_1,}
\]
donde \(e_1\) representa la posición central, la única que alcanza horizonte
tres. Definimos
\[
 A
 =\sum_{g\in S_3}
 \rho_{W_{\rm or}}(g)\,
 |w\rangle\langle b_K|\,
 \rho_{V_{11}}(\theta(g))^{-1}.
\]
Por construcción,
\[
 A\rho_{V_{11}}(\theta(g))
 =\rho_{W_{\rm or}}(g)A,
\qquad
 A=AP_3.
\]
El cálculo exacto en \(\mathbb Q(\sqrt5)\) obtiene
\[
 \operatorname{rank}A=3.
\]
Además, \(G_A=AA^*\) tiene todas sus entradas diagonales iguales a
\[
 3+\frac25\sqrt5
\]
y todas sus entradas no diagonales iguales a
\[
 \frac52+\frac35\sqrt5.
\]
Sus dos autovalores isótipo son
\[
 \lambda_{\rm sgn}=8+\frac85\sqrt5,
 \qquad
 \lambda_{\rm std}=\frac12-\frac15\sqrt5.
\]
Ambos son estrictamente positivos. En particular, \(G_A^{-1/2}\) existe y
la descomposición polar
\[
 \boxed{U_{GK}=G_A^{-1/2}A}
\]
es una isometría equivariante de \(P_3V_{11}\) sobre \(W_{\rm or}\):
\[
 U_{GK}^*U_{GK}=P_3,
 \qquad
 U_{GK}U_{GK}^*=I_{W_{\rm or}}.
\]

Sea
\[
 J_G:W_{\rm or}\longrightarrow L_{\rm or}\otimes H_G
\]
la inclusión isométrica en las tres ramas seleccionadas. Entonces
\[
 J_GJ_G^*=I_{L_{\rm or}}\otimes P_G.
\]
La isometría parcial sector--dodecafase se denota
\[
 \boxed{
 W_{GK}=U_{GK}^*J_G^*
 :L_{\rm or}\otimes H_G\longrightarrow V_{11}
 }
\]
satisface
\[
 W_{GK}^*W_{GK}
 =I_{L_{\rm or}}\otimes P_G,
\qquad
 W_{GK}W_{GK}^*=P_3,
\]
y, por tanto,
\[
 P_3W_{GK}
 =W_{GK}
 =W_{GK}(I_{L_{\rm or}}\otimes P_G).
\]

\begin{theorem}[Identificación isométrica relativa de tres proyectores]
\label{thm:identificacion-isometrica-tres-proyectores}
Relativamente a la acción \(A_5/C_5\), al proyector \(P_3\), al sello
\(K\), al testigo doblemente construido \(B_K\), al sector finito orientado
\(t=7\), al carácter abstracto \(\operatorname{sgn}_{H_*}\), al ancla
\(\ell_{\rm or}\) y a la regla variacional \(\mathcal E_B\), existe una
única isometría parcial normalizada \(W_{GK}\) que identifica \(P_G\) con
\(P_3\). Independientemente, la isometría de incidencia \(U_W\) identifica
\(P_3\) con \(Q_3\). En consecuencia, se obtiene la siguiente cadena de
identificaciones isométricas de proyectores:
\[
\begin{CD}
L_{\rm or}\otimes H_G
 @>{W_{GK}}>>
 V_{11}
 @>{U_W}>>
 E_{30}\\
@V{I_{L_{\rm or}}\otimes P_G}VV
 @VV{P_3}V
 @VV{Q_3}V\\
L_{\rm or}\otimes H_G
 @>{W_{GK}}>>
 V_{11}
 @>{U_W}>>
 E_{30}.
\end{CD}
\]
Las trazas tipadas satisfacen
\[
 \tau_{L_{\rm or}\otimes H_G}(I_{L_{\rm or}}\otimes P_G)
 =\tau_{V_{11}}(P_3)
 =\tau_{E_{30}}(Q_3)
 =\frac3{11}.
\]
El enunciado identifica
isométricamente tres proyectores; no afirma que sus tres espacios sean una
misma representación de un grupo común.
\end{theorem}

\begin{proof}
La energía \(\mathcal E_B\) selecciona un único normalizador \(H_*\);
el criterio de puntuación \(s(g)\) selecciona un único par \((c_*,r_*)\), y con él un
único isomorfismo \(\theta\). El promedio de Reynolds es entonces
equivariante para las acciones de \(S_3^{\rm or}\) y \(H_*\) relacionadas
por \(\theta\), y posee rango tres. La parte positiva de su descomposición
polar fija la normalización sin introducir una base. Las identidades
parciales anteriores demuestran que \(W_{GK}\) identifica
\(P_G\) y \(P_3\). Por otra parte, el
\cref{thm:entrelazador-dodecafase-witt} identifica \(P_3\) y \(Q_3\)
como proyectores mediante \(U_W\). Estas dos identidades operatoriales hacen
conmutar el diagrama, sin extender la equivarianza de la primera flecha a la
segunda.
\end{proof}

\begin{proposition}[Obstrucción a la equivarianza triple en el calibre fijado]
\label{prop:no-go-equivarianza-hstar-witt}
La construcción anterior no demuestra que \(U_W\) sea
\(H_*\)-equivariante. En el etiquetado declarado, la acción fija de
\(A_5/C_5\) sobre los doce puntos preserva el conjunto de hexadas de
\(W_{12}\) únicamente para el elemento identidad. Por tanto, el
entrelazamiento \(H_*\)-equivariante de \(W_{GK}\) no se prolonga, con ese
calibre, a una representación común sobre \(E_{30}\).
\end{proposition}

\begin{proof}
La primera afirmación es una separación de dominios de equivarianza:
\(W_{GK}\) usa \(H_*\subset A_5\), mientras que el teorema de incidencia
prueba la equivarianza de \(U_W\) para
\(\operatorname{Aut}(W_{12})\). El censo exacto reproducible aplica los sesenta
elementos de la acción \(A_5/C_5\) al sistema de hexadas etiquetado; sólo la
identidad conserva ese sistema. Como \(H_*\) contiene seis elementos, su
acción no puede entrar por restricción de esa acción fija en
\(\operatorname{Aut}(W_{12})\).
\end{proof}

Para promover la cadena de proyectores a una representación común es
necesario construir explícitamente un monomorfismo
\[
 \boxed{
 \iota:H_*\hookrightarrow\operatorname{Aut}(W_{12})
 }
\]
tal que, para todo \(h\in H_*\),
\[
 \boxed{
 U_W\,\rho_{V_{11}}(h)
 =\rho_{E_{30}}\!\bigl(\iota(h)\bigr)\,U_W.
 }
\]
Esta ecuación es también una condición suficiente para la equivarianza de la
segunda flecha. Ningún \(\iota\) con esta propiedad ha sido construido en el
calibre vigente; por ello la equivarianza demostrada termina en \(P_3\) y no
se prolonga por simple igualdad de rangos.

Antes de usar \(B_K\), \(u_K\), la orientación y el factor polar había
\(10\cdot6\cdot4=240\) mapas etiquetados, o \(40\) después de identificar
relabelizaciones internas del sector. Las reglas anteriores reducen ese
conjunto a uno. Éste es el contenido verificable de la palabra
\emph{canónico} para \(W_{GK}\), dentro de las reglas declaradas. No
extiende la canonicidad a una acción común de \(H_*\) sobre el módulo de Witt.

La regla variacional \(\mathcal E_B\) utiliza únicamente datos HMT anteriores
y fija la clase en la que se obtiene la unicidad. No identifica la involución
de estrella con \(P_3\), pues sus espectros siguen siendo distintos.

\section{Comparación de las lecturas de \texorpdfstring{\(t=7\) y \(t=8\)}{t=7 y t=8}}

En \(t=8\) hay tres raíces locales y once extensiones a horizonte nueve de
la única raíz que continúa. Los tamaños de las tres fibras son
\[
 (11,0,0).
\]
La matriz natural rama--cadena tiene una fila de once unos y dos filas
nulas, por lo que su rango es uno. En \(t=9\), además, el cociente visible
cambia de \(3/11\) a \(3/6\). Aquella fracción no era la traza de un
proyector estable. La corrección no elimina el sector: desplaza la lectura
al objeto homogéneo \(P_G\) de \(t=7\), donde rango y dimensión
pertenecen al mismo espacio.

\section{Consecuencia para una acción positiva}

La ley de acción puede usar ahora una única clase isométrica de proyectores
positivos en tres espacios de soporte:
\[
 P_G
 \xleftrightarrow{\ W_{GK}\ }
 P_3
 \xleftrightarrow{\ U_W\ }
 Q_3.
\]
La estrella \(R_{B_K}\) queda reservada para paridad y orientación. El
proyector \(P_G\) registra supervivencia; \(P_3\) registra la polarización
dodecafásica; \(Q_3\) registra su realización de incidencia en Witt. Esta
separación de funciones es la premisa matemática de la medida y de la ley
positiva de la sección siguiente. La cadena expresa equivalencia isométrica de
proyectores, no una representación común: esta última requeriría el
monomorfismo \(\iota\) y la identidad de equivarianza formulados arriba.
 
\chapter{Del salto espectral al retículo de Leech}
\chaptermark{Publicación excepcional del continuo}
\label{chap:83-20}

\section{Del salto espectral \texorpdfstring{\(4\to2\to6\to54\)}{4 a 2 a 6 a 54}: acoplamiento local, compresión modular e integración bidimensional}
\label{sec:sucesor83-salto-espectral-4-2-6-54}

Para dos representantes consecutivos \((k,k+1)\), considérese
\[
B_k=
\begin{pmatrix}
k^2&k(k+1)\\
k(k+1)&(k+1)^2
\end{pmatrix}\pmod 9.
\]

\begin{theorem}[Unicidad del bloque diagonal adyacente]
\label{thm:sucesor83-app-bloque-central}
Existe un único bloque \(B_k\) cuyas entradas diagonales coinciden. Está
determinado por \(k=4\) y es
\[
B_c=\begin{pmatrix}7&2\\2&7\end{pmatrix}.
\]
\end{theorem}

\begin{proof}
La igualdad de las diagonales equivale a
\[
k^2\equiv(k+1)^2\pmod 9,
\]
es decir, \(2k+1\equiv0\pmod 9\). Como \(2\) es una unidad de
\(\mathbb Z/9\mathbb Z\), la congruencia posee la única solución
\(k\equiv4\pmod 9\).
\end{proof}

Por tanto,
\[
B_c=7I+2R,
\qquad
R=\begin{pmatrix}0&1\\1&0\end{pmatrix},
\qquad R^2=I.
\]
Con \(P_\pm=(I\pm R)/2\),
\[
B_c=9P_++5P_-.
\]

\begin{theorem}[Jerarquía de la separación entre suma y producto]
\label{thm:sucesor83-app-jerarquia-4-2-6-54}
El bloque central tiene salto espectral local \(9-5=4\). Su coeficiente
fuera de la diagonal es \(2\). La compresión por las \(3\) clases módulo
\(3\) produce la diferencia agregada \(51-45=6\), y el despliegue sobre las
\(9\) posiciones produce el defecto global
\[
9\cdot6=459-405=54.
\]
Los enteros \(4,2,6,54\) pertenecen, respectivamente, al salto espectral,
al acoplamiento local, a la compresión modular y a la integración
bidimensional.
\end{theorem}

\begin{proof}
La descomposición \(B_c=7I+2R=9P_++5P_-\) da el salto y el coeficiente
local. Las matrices agregadas \(M_\Pi,M_\Sigma\) satisfacen
\(\sum M_\Pi=51\) y \(\sum M_\Sigma=45\). Cada entrada agregada representa
\(9\) posiciones, por lo que el despliegue multiplica la diferencia por
\(9\) y recupera los totales \(459\) y \(405\).
\end{proof}

Sus concatenaciones ordenadas producen los bloques \(72\) y \(27\):
\[
72+27=99,
\qquad
72-27=45=\det B_c.
\]
La concatenación de las entradas diagonales y antidiagonales produce,
respectivamente, \(77\) y \(22\). Por tanto,
\[
72+27=77+22=99,
\qquad
77-22=55=54+1.
\]
Estas igualdades expresan dos particiones exactas del mismo bloque; no
identifican por sí solas una constante externa.
 
% Unidad U018. Paley, extension cuadratica finita y codigo ternario.
% Redaccion autonoma desde la autoridad material 1063.

\section{Forma de conferencia de Paley, extensión cuadrática finita y código ternario autodual}
\label{p12-83-c20-s1:chap:paley-extension-cuadratica-codigo-ternario}

La tétrada \(B_K\) construida en la unidad anterior todavía no posee una
estrella de Witt: para definirla es preciso construir primero el sistema de
incidencia que contiene sus hexadas. Esta unidad realiza ese paso desde el
calibre proyectivo de las seis unidades módulo nueve, obtiene una raíz
cuadrática interna de \(-I_6\), y construye el código ternario gráfico de
longitud doce. Ninguna identificación se deduce de una mera igualdad de
cardinales.

\subsection{Calibre proyectivo de las \(6\) unidades módulo \(9\)}

Sea
\[
 \mathcal U_9=(\mathbb Z/9\mathbb Z)^\times
 =\{1,2,4,8,7,5\},
\]
ordenado por las potencias sucesivas de \(2\) módulo \(9\). Fijamos la
biyección de coordenadas
\begin{equation}
 \theta:\mathcal U_9\longrightarrow\mathbb P^1(\mathbb F_5),
 \qquad
 \begin{array}{c|cccccc}
 u&1&2&4&8&7&5\\ \hline
 \theta(u)&\infty&0&1&2&3&4
 \end{array}.
 \label{p12-83-c20-s1:eq:calibre-unidades-paley}
\end{equation}
La aplicación \(\theta\) es un calibre de coordenadas. No se afirma que
sea un homomorfismo entre la multiplicación de \(\mathcal U_9\) y una
operación sobre \(\mathbb P^1(\mathbb F_5)\).

Sea
\[
 \chi:\mathbb F_5\longrightarrow\{-1,0,1\}
\]
el carácter cuadrático:
\[
 \chi(0)=0,
 \qquad
 \chi(1)=\chi(4)=1,
 \qquad
 \chi(2)=\chi(3)=-1.
\]
En el orden \((\infty,0,1,2,3,4)\), definimos
\begin{align}
 (C_P)_{\infty,x}&=(C_P)_{x,\infty}=1,
 &&x\in\mathbb F_5,\label{p12-83-c20-s1:eq:paley-fila-infinito}\\
 (C_P)_{x,y}&=\chi(x-y),
 &&x,y\in\mathbb F_5, x\neq y,\label{p12-83-c20-s1:eq:paley-entradas-finitas}\\
 (C_P)_{x,x}&=0.
 \label{p12-83-c20-s1:eq:paley-diagonal}
\end{align}
Así,
\begin{equation}
 C_P=
 \begin{pmatrix}
 0& 1& 1& 1& 1& 1\\
 1& 0& 1&-1&-1& 1\\
 1& 1& 0& 1&-1&-1\\
 1&-1& 1& 0& 1&-1\\
 1&-1&-1& 1& 0& 1\\
 1& 1&-1&-1& 1& 0
 \end{pmatrix}.
 \label{p12-83-c20-s1:eq:matriz-conferencia-paley-autonoma}
\end{equation}

\begin{lemma}[Suma de correlación cuadrática]
\label{p12-83-c20-s1:lem:correlacion-cuadratica-f5}
Para \(a\neq b\) en \(\mathbb F_5\),
\begin{equation}
 \sum_{t\in\mathbb F_5}
 \chi(t-a)\chi(t-b)=-1.
 \label{p12-83-c20-s1:eq:correlacion-cuadratica-f5}
\end{equation}
\end{lemma}

\begin{proof}
La sustitución \(u=(t-a)/(b-a)\) es una biyección de \(\mathbb F_5\).
Como \(\chi((b-a)^2)=1\), la suma se reduce a
\[
 \sum_{u\in\mathbb F_5}\chi(u(u-1)).
\]
La evaluación de \(u=0,1,2,3,4\) da respectivamente
\(0,0,1,-1,-1\), cuya suma es \(-1\).
\end{proof}

\begin{proposition}[Identidad de conferencia]
\label{p12-83-c20-s1:prop:paley-identidad-conferencia-autonoma}
La matriz \(C_P\) es simétrica y satisface
\begin{equation}
 \boxed{C_PC_P^{\mathsf T}=5I_6.}
 \label{p12-83-c20-s1:eq:paley-identidad-conferencia-autonoma}
\end{equation}
\end{proposition}

\begin{proof}
Cada fila contiene cinco entradas de módulo uno y una entrada nula; su
norma cuadrada es cinco. Para dos filas finitas distintas, la contribución
de la coordenada \(\infty\) es \(+1\), mientras
\cref{p12-83-c20-s1:lem:correlacion-cuadratica-f5} da \(-1\) para las cinco coordenadas
finitas. Su producto escalar es cero. La suma de un carácter no trivial
sobre \(\mathbb F_5\) es cero, por lo que la fila correspondiente a
\(\infty\) también es ortogonal a cada fila finita.
\end{proof}

\subsection{Reducción ternaria y raíz cuadrática interna de \texorpdfstring{\(-I_6\)}{-I₆}}

Reducimos \(C_P\) módulo tres, con \(-1\equiv2\pmod3\):
\begin{equation}
 A_W=
 \begin{pmatrix}
 0&1&1&1&1&1\\
 1&0&1&2&2&1\\
 1&1&0&1&2&2\\
 1&2&1&0&1&2\\
 1&2&2&1&0&1\\
 1&1&2&2&1&0
 \end{pmatrix}
 \in M_6(\mathbb F_3).
 \label{p12-83-c20-s1:eq:matriz-paley-witt-ternaria}
\end{equation}

\begin{theorem}[Raíz cuadrática interna de \(-I_6\)]
\label{p12-83-c20-s1:thm:aw-cuadrado-menos-identidad-autonomo}
Se cumple
\begin{equation}
 \boxed{A_W^2=-I_6.}
 \label{p12-83-c20-s1:eq:aw-cuadrado-menos-identidad-autonomo}
\end{equation}
En particular,
\[
 A_W^{-1}=-A_W.
\]
\end{theorem}

\begin{proof}
La \cref{p12-83-c20-s1:prop:paley-identidad-conferencia-autonoma} da
\[
 A_WA_W^{\mathsf T}
 \equiv5I_6
 \equiv2I_6
 =-I_6
 \pmod3.
\]
Como \(A_W=A_W^{\mathsf T}\), se obtiene \(A_W^2=-I_6\).
\end{proof}

La identidad anterior es la residencia algebraica de la raíz de \(-1\)
en este corredor. No procede del factor \(1/4\) de la inversión Hadamard:
aquél es un escalar racional forzado por \(H_4^2=4I_4\); ésta es una
estructura cuadrática sobre característica tres.

\subsection{Extensión cuadrática a \texorpdfstring{\(\mathbb F_9\)}{F₉}}

Definimos
\begin{equation}
 \mathbb F_9=\mathbb F_3[\iota]/(\iota^2+1),
 \qquad \iota^2=-1.
 \label{p12-83-c20-s1:eq:definicion-f9-autonoma}
\end{equation}
El polinomio \(X^2+1\) no tiene raíces en \(\mathbb F_3\), por lo que el
cociente es un cuerpo de nueve elementos. Para
\(V=\mathbb F_3^6\), sea
\[
 V_9=V\otimes_{\mathbb F_3}\mathbb F_9.
\]
Los operadores
\begin{equation}
 P_\iota=\frac12(I-\iota A_W),
 \qquad
 P_{-\iota}=\frac12(I+\iota A_W)
 \label{p12-83-c20-s1:eq:proyectores-f9-paley}
\end{equation}
están bien definidos porque \(2\) es invertible en \(\mathbb F_3\).

\begin{proposition}[Proyectores conjugados]
\label{p12-83-c20-s1:prop:proyectores-conjugados-paley}
Se cumplen
\[
 P_\iota^2=P_\iota,
 \qquad
 P_{-\iota}^2=P_{-\iota},
 \qquad
 P_\iota P_{-\iota}=0,
 \qquad
 P_\iota+P_{-\iota}=I.
\]
\end{proposition}

\begin{proof}
Usando \(\iota^2=-1\) y \(A_W^2=-I\),
\[
 (\iota A_W)^2=I.
\]
Las cuatro identidades son entonces las identidades elementales de los
proyectores \((I\mp\iota A_W)/2\) de una involución.
\end{proof}

\begin{theorem}[Descomposición espectral conjugada]
\label{p12-83-c20-s1:thm:descomposicion-espectral-f9-paley}
Sea
\[
 E_{\pm\iota}=\ker(A_W\mp\iota I).
\]
Entonces
\begin{equation}
 V_9=E_\iota\oplus E_{-\iota},
 \qquad
 \dim_{\mathbb F_9}E_\iota
 =\dim_{\mathbb F_9}E_{-\iota}=3.
 \label{p12-83-c20-s1:eq:descomposicion-espectral-f9-paley}
\end{equation}
\end{theorem}

\begin{proof}
El polinomio mínimo de \(A_W\) divide \(X^2+1\), que se escinde sobre
\(\mathbb F_9\) con raíces simples \(\pm\iota\). Los proyectores de
\cref{p12-83-c20-s1:prop:proyectores-conjugados-paley} producen la suma directa. Como
\(A_W\) tiene coeficientes en \(\mathbb F_3\), el automorfismo de
Frobenius intercambia los dos autoespacios; sus dimensiones son iguales y
suman seis.
\end{proof}

Esta construcción dota \(\mathbb F_3^6\) de estructura de módulo de rango
tres sobre \(\mathbb F_9\), mediante
\begin{equation}
 (a+b\iota)\cdot v:=av+b(vA_W).
 \label{p12-83-c20-s1:eq:accion-f9-sobre-f3-seis}
\end{equation}
La igualdad \(3^6=9^3\) queda así acompañada por una acción explícita; no
se usa como argumento cardinal.

\subsection{Código gráfico ternario de longitud \(12\)}

Definimos
\begin{equation}
 c:\mathbb F_3^6\longrightarrow\mathbb F_3^{12},
 \qquad
 c(q)=(q,qA_W),
 \label{p12-83-c20-s1:eq:aplicacion-codigo-grafico-witt}
\end{equation}
y
\begin{equation}
 C_W:=\operatorname{im}c.
 \label{p12-83-c20-s1:eq:definicion-codigo-witt-ternario}
\end{equation}
Su matriz generadora es
\[
 G_W=(I_6\mid A_W).
\]

\begin{theorem}[Inyectividad y autodualidad]
\label{p12-83-c20-s1:thm:cw-inyectivo-autodual}
La aplicación \(c\) es inyectiva, \(\dim C_W=6\) y
\begin{equation}
 \boxed{C_W=C_W^\perp.}
 \label{p12-83-c20-s1:eq:cw-autodual}
\end{equation}
\end{theorem}

\begin{proof}
La primera mitad de \(c(q)\) es \(q\), por lo que
\(\operatorname{pr}_1\circ c=\operatorname{id}\); esto prueba la
inyectividad y la dimensión. Además,
\[
 G_WG_W^{\mathsf T}
 =I_6+A_WA_W^{\mathsf T}
 =I_6+2I_6=0
\]
en \(\mathbb F_3\). Por tanto \(C_W\subseteq C_W^\perp\). Ambos espacios
tienen dimensión seis, luego son iguales.
\end{proof}

Para \(r=0,\ldots,12\), sea
\[
 N_r:=\#\{q\in\mathbb F_3^6:
            \operatorname{wt}(q,qA_W)=r\}.
\]
La enumeración de las \(3^6=729\) entradas produce
\begin{equation}
 N_0=1,
 \qquad
 N_6=264,
 \qquad
 N_9=440,
 \qquad
 N_{12}=24,
 \label{p12-83-c20-s1:eq:censo-pesos-cw}
\end{equation}
y \(N_r=0\) en los demás pesos.

\begin{theorem}[Enumerador y distancia mínima]
\label{p12-83-c20-s1:thm:enumerador-cw-autonomo}
El código \(C_W\) posee parámetros
\[
 \boxed{C_W=[12,6,6]_3}
\]
y enumerador de pesos
\begin{equation}
 \boxed{W_{C_W}(y)=1+264y^6+440y^9+24y^{12}.}
 \label{p12-83-c20-s1:eq:enumerador-pesos-cw}
\end{equation}
\end{theorem}

\begin{proof}
La dimensión y la autodualidad ya se demostraron. El censo
\eqref{p12-83-c20-s1:eq:censo-pesos-cw} recorre las \(729\) palabras sin muestreo. Su
menor peso no nulo es seis, que es la distancia mínima de un código lineal.
La suma \(1+264+440+24=729\) controla que ninguna palabra fue omitida.
\end{proof}

\subsection{Naturaleza de la estructura obtenida}

Se han construido tres objetos distintos y tres flechas tipadas:
\[
 \mathcal U_9
 \xrightarrow{\theta}\mathbb P^1(\mathbb F_5)
 \xrightarrow{\chi}C_P
 \xrightarrow{\bmod3}A_W,
\]
\[
 A_W^2=-I_6
 \quad\Longrightarrow\quad
 \mathbb F_3^6\cong\mathbb F_9^3
 \quad\Longrightarrow\quad
 C_W=\{(q,qA_W):q\in\mathbb F_3^6\}.
\]
La primera flecha es un calibre, la segunda usa el carácter cuadrático y la
tercera es reducción de coeficientes. La estructura sobre
\(\mathbb F_9\) es finita; no se identifica con una estructura compleja
real o analítica.

\subsection{Obstrucciones del corredor Paley–\(C_W\): biyección, \(A_W^2=-I_6\), autodualidad y censo de \(729\) palabras}

\begin{criterion}[Refutación del corredor Paley--\(C_W\)]
La construcción queda refutada si ocurre cualquiera de los hechos
siguientes:
\begin{enumerate}
 \item \(\theta\) no es biyectiva o se usa como homomorfismo sin prueba;
 \item la suma \eqref{p12-83-c20-s1:eq:correlacion-cuadratica-f5} difiere de \(-1\);
 \item \(C_PC_P^{\mathsf T}\neq5I_6\);
 \item \(A_W^2\neq-I_6\);
 \item los proyectores \eqref{p12-83-c20-s1:eq:proyectores-f9-paley} no son
       complementarios;
 \item la aplicación \(q\mapsto(q,qA_W)\) pierde inyectividad o
       autodualidad;
 \item el censo no suma \(729\), difiere de
       \eqref{p12-83-c20-s1:eq:censo-pesos-cw}, o presenta una palabra de peso entre uno
       y cinco;
 \item se infiere \(\mathbb F_3^6\cong\mathbb F_9^3\) únicamente de
       \(3^6=9^3\), omitiendo la acción \eqref{p12-83-c20-s1:eq:accion-f9-sobre-f3-seis}.
\end{enumerate}
\end{criterion}


% Unidad U019. Sistema de Witt, campo de estrellas y grupo M12.
% Redaccion autonoma desde la autoridad material 1063.

\section{Sistema de Steiner ternario, involuciones de estrella de Witt y generación de \texorpdfstring{\(M_{12}\)}{M12}}
\label{p12-83-c20-s2:chap:sistema-witt-estrellas-m12}

La unidad precedente construyó el código autodual
\(C_W=[12,6,6]_3\). Ahora se proyectivizan sus palabras de peso seis,
se prueba la propiedad de Steiner, y sólo entonces se define la estrella
asociada a la tétrada \(B_K\). Este orden impide convertir una figura de
incidencia en un objeto añadido por semejanza visual.

\subsection{Soportes proyectivos de peso \(6\)}

Sea
\begin{equation}
 \mathcal H_W
 :=\{\operatorname{supp}(c):c\in C_W,
                              \operatorname{wt}(c)=6\}.
 \label{p12-83-c20-s2:eq:familia-soportes-peso-seis}
\end{equation}
Cada soporte aparece para el par de palabras \(\{c,-c\}\). En efecto, si dos
palabras de peso seis con el mismo soporte no fueran proporcionales, se podría
escalar una de ellas para cancelar una coordenada al restarlas; se obtendría
una palabra no nula de \(C_W\) de peso menor que seis, contra la distancia
mínima. Como tampoco existe una palabra no nula igual a su opuesta en
característica tres, cada órbita proyectiva tiene exactamente dos elementos.
Por tanto,
\begin{equation}
 |\mathcal H_W|=\frac{264}{2}=132.
 \label{p12-83-c20-s2:eq:numero-hexadas-cw}
\end{equation}

Para \(T\in\binom{[12]}5\), definimos
\begin{equation}
 \nu(T):=\#\{H\in\mathcal H_W:T\subseteq H\}.
 \label{p12-83-c20-s2:eq:incidencia-cinco-hexada}
\end{equation}
La enumeración exacta de los
\[
 \binom{12}{5}=792
\]
subconjuntos de cinco puntos produce
\begin{equation}
 \nu(T)=1
 \qquad
 \text{para todo }T\in\binom{[12]}5.
 \label{p12-83-c20-s2:eq:unicidad-cinco-en-hexada}
\end{equation}

\begin{theorem}[Sistema de Witt ternario]
\label{p12-83-c20-s2:thm:w12-autonomo}
El sistema de incidencia
\[
 W_{12}:=([12],\mathcal H_W)
\]
es el sistema de Steiner
\begin{equation}
 \boxed{W_{12}=S(5,6,12).}
 \label{p12-83-c20-s2:eq:w12-steiner-autonomo}
\end{equation}
\end{theorem}

\begin{proof}
La igualdad \eqref{p12-83-c20-s2:eq:unicidad-cinco-en-hexada} es precisamente la
propiedad de Steiner. El número de bloques también se recupera de la doble
cuenta de pares \((T,H)\) con \(T\subset H\):
\[
 |\mathcal H_W|\binom65=\binom{12}{5},
 \qquad
 |\mathcal H_W|=\frac{792}{6}=132.
\]
\end{proof}

\subsection{Residuo de una tétrada y descomposición en \(4\) componentes incidentes}

Sea \(B\in\binom{[12]}4\). En un sistema \(S(5,6,12)\), el número de
hexadas que contienen \(B\) es
\begin{equation}
 \lambda_4
 =\frac{\binom{12-4}{5-4}}{\binom{6-4}{5-4}}
 =\frac82=4.
 \label{p12-83-c20-s2:eq:lambda-cuatro-w12}
\end{equation}
Escribimos estas cuatro hexadas como
\[
 H_i=B\sqcup P_i,
 \qquad |P_i|=2,
 \qquad i=1,\ldots,4.
\]

\begin{lemma}[Partición residual de una tétrada]
\label{p12-83-c20-s2:lem:particion-residual-tetrada-witt}
Los cuatro pares \(P_1,\ldots,P_4\) son disjuntos y
\begin{equation}
 [12]\setminus B
 =P_1\sqcup P_2\sqcup P_3\sqcup P_4.
 \label{p12-83-c20-s2:eq:particion-residual-tetrada-witt}
\end{equation}
\end{lemma}

\begin{proof}
Si dos pares residuales compartieran un punto \(p\), las correspondientes
hexadas contendrían el mismo subconjunto \(B\cup\{p\}\) de cinco puntos,
en contradicción con la unicidad de Steiner. Los cuatro pares son, pues,
disjuntos; contienen ocho puntos y agotan el complemento de \(B\).
\end{proof}

\begin{definition}[Estrella e involución de estrella de Witt]
\label{p12-83-c20-s2:def:estrella-witt-involucion}
La estrella centrada en \(B\) es
\begin{equation}
 \mathscr S_B
 :=\{H\in\mathcal H_W:B\subseteq H\}
 =\{B\sqcup P_1,\ldots,B\sqcup P_4\}.
 \label{p12-83-c20-s2:eq:estrella-witt-general}
\end{equation}
Si \(P_i=\{p_i^-,p_i^+\}\), su involución es
\begin{equation}
 \sigma_B:=\prod_{i=1}^{4}(p_i^-\ p_i^+).
 \label{p12-83-c20-s2:eq:involucion-estrella-witt-general}
\end{equation}
La permutación \(\sigma_B\) fija puntualmente \(B\) e intercambia los
extremos de cada pétalo residual.
\end{definition}

El término matemático empleado en adelante es \emph{estrella de Witt} o
\emph{involución de estrella de Witt}. No se trata de una estrella
euclídea.

\subsection{Estrella seleccionada por la tétrada del sello}

La unidad dodecafásica produjo, sin usar el diseño de Steiner,
\[
 B_K=\{4,6,9,10\}.
\]
Ahora que \(W_{12}\) está construido, su estrella queda bien definida. Las
cuatro hexadas incidentes son
\begin{align}
 H_{13}&=B_K\sqcup\{1,3\}
       =\{1,3,4,6,9,10\},
 \label{p12-83-c20-s2:eq:hexada-petalo-13}\\
 H_{2,11}&=B_K\sqcup\{2,11\}
       =\{2,4,6,9,10,11\},
 \label{p12-83-c20-s2:eq:hexada-petalo-211}\\
 H_{5,7}&=B_K\sqcup\{5,7\}
       =\{4,5,6,7,9,10\}
       =H_\alpha^-,
 \label{p12-83-c20-s2:eq:hexada-petalo-57}\\
 H_{8,12}&=B_K\sqcup\{8,12\}
       =\{4,6,8,9,10,12\}.
 \label{p12-83-c20-s2:eq:hexada-petalo-812}
\end{align}
En consecuencia,
\begin{equation}
 \boxed{
 \sigma_{B_K}=(1\ 3)(2\ 11)(5\ 7)(8\ 12).}
 \label{p12-83-c20-s2:eq:involucion-estrella-witt-bk}
\end{equation}

\begin{figure}[htbp]
\centering
\begin{tikzpicture}[
  >=Stealth,
  font=\small,
  core/.style={draw,double,rounded corners=2mm,align=center,
    minimum width=4.0cm,minimum height=1.3cm,fill=black!4},
  petal/.style={draw,rounded corners=2mm,align=center,
    minimum width=3.55cm,minimum height=1.08cm,fill=black!2},
  negative/.style={petal,line width=1.1pt,double,fill=black!10},
  arr/.style={->,thick}
]
\node[core] (B) at (0,0)
 {\(\displaystyle B_K=\{4,6,9,10\}\)\\
  puntos fijados por \(\sigma_{B_K}\)};
\node[petal] (P13) at (0,3.0)
 {\(P_1=\{1,3\}\), \((1\ 3)\)\\
  \(H_{13}=B_K\sqcup P_1\)};
\node[petal] (P211) at (5.25,0)
 {\(P_2=\{2,11\}\), \((2\ 11)\)\\
  \(H_{2,11}=B_K\sqcup P_2\)};
\node[negative] (P57) at (0,-3.0)
 {\(P_3=\{5,7\}\), \((5\ 7)\)\\
  \(H_{5,7}=H_\alpha^-\)};
\node[petal] (P812) at (-5.25,0)
 {\(P_4=\{8,12\}\), \((8\ 12)\)\\
  \(H_{8,12}=B_K\sqcup P_4\)};
\draw[arr] (B)--(P13);
\draw[arr] (B)--(P211);
\draw[arr] (B)--(P57);
\draw[arr] (B)--(P812);
\node[align=center,font=\footnotesize] at (0,-4.35)
 {\(\mathscr S_{B_K}
   =\{H_{13},H_{2,11},H_{5,7},H_{8,12}\}\).\\
  El diagrama representa incidencia, no distancia geométrica.};
\end{tikzpicture}
\caption{Estrella de Witt centrada en \(B_K\). El centro es la tétrada
fija; los cuatro pétalos son los pares residuales que determinan la
involución.}
\label{p12-83-c20-s2:fig:estrella-witt-bk-autonoma}
\end{figure}

\subsection{Campo global de las 495 estrellas}

Existe una estrella para cada tétrada, por lo que
\begin{equation}
 \#\{\mathscr S_B:B\in\tbinom{[12]}4\}
 =\binom{12}{4}=495.
 \label{p12-83-c20-s2:eq:numero-estrellas-witt}
\end{equation}
El censo exacto verifica
\begin{equation}
 \sigma_B(\mathcal H_W)=\mathcal H_W
 \qquad
 \text{para toda }B\in\binom{[12]}4.
 \label{p12-83-c20-s2:eq:estrellas-conservan-hexadas}
\end{equation}
Definimos
\begin{equation}
 \Gamma_\star
 :=\langle\sigma_B:B\in\tbinom{[12]}4\rangle
 \leq\operatorname{Aut}(W_{12}).
 \label{p12-83-c20-s2:eq:grupo-generado-estrellas-witt}
\end{equation}

Un conjunto generador de seis tétradas es
\[
 1234,
 \quad1235,
 \quad1236,
 \quad1245,
 \quad1345,
 \quad2345.
\]
Los órdenes de los subgrupos obtenidos al añadir sus involuciones
sucesivamente son
\begin{equation}
 1, 2, 6, 18, 360, 7920, 95040.
 \label{p12-83-c20-s2:eq:cadena-ordenes-generadores-m12}
\end{equation}

\begin{theorem}[Generación del grupo de Mathieu \(M_{12}\)]
\label{p12-83-c20-s2:thm:495-estrellas-m12-autonomo}
Se cumple
\begin{equation}
 \boxed{
 \langle\sigma_B:B\in\tbinom{[12]}4\rangle
 =\operatorname{Aut}(W_{12})
 \cong M_{12}.}
 \label{p12-83-c20-s2:eq:estrellas-generan-m12}
\end{equation}
\end{theorem}

\begin{proof}
Cada \(\sigma_B\) conserva las \(132\) hexadas, luego
\(\Gamma_\star\leq\operatorname{Aut}(W_{12})\). El censo de Schreier de
los seis generadores indicados produce orden \(95040\). El grupo de
automorfismos de \(S(5,6,12)\) es \(M_{12}\) y tiene ese mismo orden; por
tanto la inclusión es una igualdad.
\end{proof}

Una estrella aislada genera únicamente
\[
 \langle\sigma_B\rangle\cong C_2.
\]
La conclusión sobre \(M_{12}\) pertenece al campo global de las
\(495\) involuciones, no a \(\sigma_{B_K}\) considerada aisladamente.

\subsection{Acción sobre el módulo de aumentación}

Sea
\[
 V_{11}=\mathbf1^\perp\subset\mathbb R^{12}.
\]
La involución \(\sigma_B\) fija las cuatro coordenadas de \(B\) y actúa
por cuatro transposiciones sobre su complemento. En \(\mathbb R^{12}\)
posee ocho direcciones de autovalor \(+1\) y cuatro de autovalor \(-1\).
La dirección constante pertenece al primer autoespacio. Por tanto,
\begin{equation}
 \operatorname{Spec}(\sigma_B|_{V_{11}})
 =\{1^{[7]},(-1)^{[4]}\},
 \qquad
 \frac1{11}\operatorname{Tr}(\sigma_B|_{V_{11}})=\frac3{11}.
 \label{p12-83-c20-s2:eq:espectro-involucion-estrella-aumentacion}
\end{equation}
Es una involución firmada; no es un proyector positivo de rango tres.

\subsection{Obstrucciones de \(W_{12}\), del campo de \(495\) involuciones de estrella y de la generación de \(M_{12}\)}

\begin{criterion}[Refutación de \(W_{12}\), sus estrellas y \(M_{12}\)]
La construcción queda refutada si ocurre cualquiera de los hechos
siguientes:
\begin{enumerate}
 \item la proyectivización de las \(264\) palabras de peso seis no produce
       \(132\) soportes;
 \item existe un subconjunto de cinco puntos contenido en cero o en más de
       una hexada;
 \item una tétrada no pertenece exactamente a cuatro hexadas;
 \item los cuatro pares residuales no particionan el complemento de la
       tétrada;
 \item los pétalos de \(B_K\) no son
       \(\{1,3\},\{2,11\},\{5,7\},\{8,12\}\);
 \item alguna de las \(495\) involuciones no conserva
       \(\mathcal H_W\);
 \item el grupo generado no tiene orden \(95040\);
 \item se atribuye \(M_{12}\) a una única estrella en vez de al campo
       global de involuciones;
 \item se introduce la estrella antes de construir \(W_{12}\).
\end{enumerate}
\end{criterion}


% Unidad U020. Prolongacion binaria W12--W24 y M24.
% Redaccion autonoma desde la autoridad material 1063.

\section{Extensión binaria relativa a una dodecada, sistema \texorpdfstring{\(W_{24}\)}{W24} y grupo \texorpdfstring{\(M_{24}\)}{M24}}
\label{p12-83-c20-s3:chap:dodecada-w24-m24}

El sistema \(W_{12}\) construido sobre doce fases admite una prolongación
binaria cuando se declara un código de Golay extendido, una dodecada y un
isomorfismo de incidencia. Esos datos son entradas relativas de esta rama:
no se deducen únicamente de \(C_W\), y no intervienen en la rama reticular
ternaria que conducirá al retículo de Leech.

\subsection{Código binario de Golay y octadas}

Sea
\[
 \mathcal G_{24}\subset\mathbb F_2^{24}
\]
el código binario extendido de Golay. Sus pesos son
\begin{equation}
 0, 8, 12, 16, 24.
 \label{p12-83-c20-s3:eq:pesos-golay-binario-extendido}
\end{equation}
Los soportes de las palabras de peso ocho forman el sistema de Steiner
\begin{equation}
 W_{24}=S(5,8,24).
 \label{p12-83-c20-s3:eq:w24-steiner-octadas}
\end{equation}
Sus bloques se denominan \emph{octadas}. Fijemos una dodecada \(D\), es
decir, el soporte de una palabra de peso doce, y sea \(\mathcal O_{24}\)
la familia de octadas.

Definimos
\begin{equation}
 \mathcal H(D)
 :=\{O\cap D:O\in\mathcal O_{24},\ |O\cap D|=6\}.
 \label{p12-83-c20-s3:eq:hexadas-residuales-dodecada}
\end{equation}

\begin{theorem}[Diseño residual de una dodecada]
\label{p12-83-c20-s3:thm:diseno-residual-dodecada-autonomo}
El sistema residual es
\begin{equation}
 \boxed{(D,\mathcal H(D))\cong S(5,6,12).}
 \label{p12-83-c20-s3:eq:diseno-residual-dodecada-w12}
\end{equation}
Además, cada \(H\in\mathcal H(D)\) posee una única octada \(O_H\) tal que
\begin{equation}
 O_H\cap D=H.
 \label{p12-83-c20-s3:eq:elevacion-unica-hexada-octada}
\end{equation}
\end{theorem}

\begin{proof}
Sea \(A\subset D\), con \(|A|=5\). En \(W_{24}=S(5,8,24)\) existe una
única octada \(O_A\) que contiene \(A\). Si \(j=|O_A\cap D|\), la suma
binaria de las palabras características de \(O_A\) y \(D\) pertenece a
\(\mathcal G_{24}\) y tiene peso
\[
 8+12-2j=20-2j.
\]
Como \(A\subseteq O_A\cap D\), se tiene \(5\leq j\leq8\). Los cuatro
pesos posibles \(10,8,6,4\), combinados con
\eqref{p12-83-c20-s3:eq:pesos-golay-binario-extendido}, fuerzan \(j=6\). Así, todo
subconjunto de cinco puntos de \(D\) pertenece a una hexada residual única.

Si dos octadas tuvieran la misma intersección hexádica con \(D\), ambas
contendrían los mismos cinco puntos de esa hexada, contradiciendo la
unicidad en \(S(5,8,24)\). Esto prueba también la unicidad del
levantamiento.
\end{proof}

\subsection{Calibre relativo entre el diseño HMT y la dodecada}

Para relacionar la construcción anterior con el diseño obtenido de
\(C_W\), deben declararse explícitamente
\begin{equation}
 (D,\ell),
 \qquad
 \ell:W_{12}^{\rm HMT}
 \xrightarrow{\sim}(D,\mathcal H(D)),
 \label{p12-83-c20-s3:eq:dato-relativo-dodecada-alineamiento}
\end{equation}
donde \(\ell\) es un isomorfismo de incidencia. Si \(\ell_*\) denota su
acción sobre bloques, definimos
\begin{equation}
 \widehat\iota_{D,\ell}
 :=\iota_D\circ\ell_*:
 W_{12}^{\rm HMT}\longrightarrow W_{24},
 \qquad
 \widehat\iota_{D,\ell}(H)=O_{\ell(H)}.
 \label{p12-83-c20-s3:eq:elevacion-tipificada-w12-w24}
\end{equation}
La composición \(\widehat\iota_{D,\ell}\) evita aplicar
\(\iota_D\), cuyo dominio es \(\mathcal H(D)\), directamente a una hexada
que todavía vive en \([12]_{\rm HMT}\).

\begin{proposition}[Inyectividad de la elevación relativa]
\label{p12-83-c20-s3:prop:inyectividad-elevacion-w12-w24}
La aplicación \(\widehat\iota_{D,\ell}\) es inyectiva y conserva la
incidencia entre puntos y hexadas después de aplicar \(\ell\).
\end{proposition}

\begin{proof}
El isomorfismo \(\ell\) es biyectivo sobre bloques. Por el
\cref{p12-83-c20-s3:thm:diseno-residual-dodecada-autonomo}, cada hexada residual posee una
única octada elevadora. Dos hexadas HMT con la misma imagen tendrían la misma
hexada residual y, por biyectividad de \(\ell\), serían iguales.
\end{proof}

No existe una elección absoluta de \((D,\ell)\) previa a este calibre. El
objeto invariante es la clase de la construcción bajo cambio simultáneo de
dodecada y alineamiento.

\subsection{Ley de incidencia \texorpdfstring{\(4+1\)}{4+1} sobre una tétrada}

Para una tétrada \(B\subset D\), los parámetros derivados de los dos
sistemas de Steiner son
\begin{equation}
 \lambda_4(W_{12})
 =\frac{\binom81}{\binom21}=4,
 \qquad
 \lambda_4(W_{24})
 =\frac{\binom{20}1}{\binom41}=5.
 \label{p12-83-c20-s3:eq:ley-cuatro-mas-uno-witt}
\end{equation}
Por tanto, las cinco octadas que contienen \(B\) se descomponen de manera
única en
\begin{equation}
 \{O_H:H\in\mathcal H(D),\ B\subset H\}
 \sqcup\{O_B^\perp\}.
 \label{p12-83-c20-s3:eq:descomposicion-cinco-octadas}
\end{equation}
Las cuatro primeras satisfacen \(|O_H\cap D|=6\); la quinta satisface
\begin{equation}
 O_B^\perp\cap D=B.
 \label{p12-83-c20-s3:eq:octada-transversal-tetrada}
\end{equation}

Para \(B=B_K\), las cuatro elevaciones residuales corresponden a los
pétalos
\[
 \{1,3\},
 \qquad
 \{2,11\},
 \qquad
 \{5,7\},
 \qquad
 \{8,12\},
\]
y la quinta octada es la transversal \(O_{B_K}^\perp\). Esta ley explica
la transición de cuatro hexadas a cinco octadas sin identificar una cara
de un cubo con una hexada ni un vértice con una octada.

\subsection{Correspondencia de estratos y libertad del triple positivo}

La estructura de cinco regiones del canal \(\pi\) tiene tipos
\[
 1_\perp+1_{--}+3_{++}.
\]
La incidencia sobre \(B_K\) tiene tipos
\[
 1_{\rm transversal}+1_{H^-}+3_{\rm restantes}
\]
después de señalar la hexada negativa. Existe, por tanto, una
correspondencia natural entre \emph{estratos de rol}:
\begin{equation}
 1_\perp\longleftrightarrow O_{B_K}^\perp,
 \qquad
 1_{--}\longleftrightarrow
 \widehat\iota_{D,\ell}(H_\alpha^-),
 \qquad
 3_{++}\longleftrightarrow
 \mathcal O_{B_K}^{+,3}.
 \label{p12-83-c20-s3:eq:correspondencia-estratos-cinco-cinco}
\end{equation}
Aquí \(\mathcal O_{B_K}^{+,3}\) es el conjunto no ordenado de las otras
tres octadas residuales. Una biyección individual entre las tres regiones
positivas y esas tres octadas exige declarar un calibre adicional de
\(S_3\). La coincidencia del patrón \(1+1+3\) no determina por sí sola ese
ordenamiento.

\subsection{Grupo de automorfismos del sistema de octadas}

\begin{theorem}[Grupo de Mathieu \(M_{24}\)]
\label{p12-83-c20-s3:thm:automorfismos-w24-m24}
El grupo de automorfismos del sistema de Steiner de octadas es
\begin{equation}
 \boxed{
 \operatorname{Aut}(W_{24})\cong M_{24},
 \qquad
 |M_{24}|=244823040.}
 \label{p12-83-c20-s3:eq:automorfismos-w24-m24}
\end{equation}
\end{theorem}

El teorema identifica el grupo clásico de automorfismos del dato binario
de entrada. No afirma que \(M_{24}\) se genere a partir de una sola
dodecada, ni que el código ternario \(C_W\) determine por sí solo
\(\mathcal G_{24}\).

\subsection{Separación categórica respecto de la rama reticular}

La rama construida en esta unidad es
\begin{equation}
 (W_{12}^{\rm HMT};\mathcal G_{24},D,\ell)
 \longrightarrow W_{24}
 \longrightarrow\operatorname{Aut}(W_{24})\cong M_{24}.
 \label{p12-83-c20-s3:eq:rama-binaria-w24-m24}
\end{equation}
La rama reticular que se construirá después es
\begin{equation}
 C_W\longrightarrow N(A_2^{12})
 \xrightarrow{\text{vecino de índice tres}}\Lambda_{24}.
 \label{p12-83-c20-s3:eq:rama-ternaria-niemeier-leech-anunciada}
\end{equation}
Sus constructores, dominios y codominios son distintos. En particular,
\begin{equation}
 \boxed{
 \text{no se introduce ninguna flecha }
 W_{24}\longrightarrow\Lambda_{24}.}
 \label{p12-83-c20-s3:eq:no-flecha-w24-leech}
\end{equation}

\begin{figure}[htbp]
\centering
\begin{tikzpicture}[
  >=Stealth,
  node distance=13mm and 16mm,
  box/.style={draw,rounded corners=1.5mm,align=center,
    inner xsep=3mm,inner ysep=2.2mm},
  arr/.style={->,thick}
]
\node[box] (cw) {\(C_W\)};
\node[box,right=of cw] (w12) {\(W_{12}\)};
\node[box,right=of w12] (w24) {\(W_{24}\)};
\node[box,right=of w24] (m24) {\(M_{24}\)};
\node[box,below=of cw] (cw2) {\(C_W\)};
\node[box,right=of cw2] (n) {\(N(A_2^{12})\)};
\node[box,right=of n] (leech) {\(\Lambda_{24}\)};
\draw[arr] (cw)--(w12);
\draw[arr] (w12)--node[above,font=\scriptsize]
 {\((\mathcal G_{24},D,\ell)\)} (w24);
\draw[arr] (w24)--(m24);
\draw[arr] (cw2)--node[below,font=\scriptsize]
 {pegado ternario} (n);
\draw[arr] (n)--node[below,font=\scriptsize]
 {\(3\)-vecino} (leech);
\node[below=5mm of w24,font=\footnotesize,align=center]
 {La alineación vertical no es una flecha.};
\end{tikzpicture}
\caption{Bifurcación tipada de las prolongaciones binaria y ternaria.}
\label{p12-83-c20-s3:fig:bifurcacion-w24-leech-autonoma}
\end{figure}

\subsection{Obstrucciones de la prolongación binaria \(W_{12}\to W_{24}\): código de Golay, octadas, alineamiento y separación de la rama de Leech}

\begin{criterion}[Refutación de la prolongación binaria]
La construcción queda refutada si ocurre cualquiera de los hechos
siguientes:
\begin{enumerate}
 \item el espectro de pesos de \(\mathcal G_{24}\) contiene un peso no
       listado en \eqref{p12-83-c20-s3:eq:pesos-golay-binario-extendido};
 \item los soportes de peso ocho no forman \(S(5,8,24)\);
 \item una intersección residual de una octada con \(D\) que contiene cinco
       puntos de \(D\) tiene cardinal distinto de seis;
 \item una hexada residual admite cero o más de una octada elevadora;
 \item se aplica \(\iota_D\) directamente a una hexada HMT sin el
       alineamiento \(\ell\);
 \item la ley de incidencia sobre una tétrada no es \(4+1\);
 \item se ordena el triple positivo sin declarar el calibre \(S_3\);
 \item \(\operatorname{Aut}(W_{24})\) no tiene orden \(244823040\);
 \item se afirma que \(W_{24}\) surge de \(W_{12}\) sin el dato binario
       \(\mathcal G_{24}\), la dodecada y el alineamiento;
 \item se serializan las dos ramas mediante una flecha
       \(W_{24}\to\Lambda_{24}\).
\end{enumerate}
\end{criterion}


% Unidad U021. Pegado discriminante ternario y Niemeier A2^12.
% Redaccion autonoma desde la autoridad material 1063.

\section{Pegado discriminante ternario y construcción del retículo de Niemeier de tipo \texorpdfstring{\(A_2^{12}\)}{A₂¹²}}
\label{p12-83-c20-s4:chap:pegado-niemeier-a2-doce}

Esta unidad inicia la rama reticular directamente desde el código ternario
autodual \(C_W\). El sistema binario \(W_{24}\), la dodecada y el grupo
\(M_{24}\) no son entradas de esta construcción. La sucesión tipada es
\[
 C_W
 \longrightarrow
 (A_2^*/A_2)^{12}
 \longrightarrow
 N(C_W)
 \cong N(A_2^{12}).
\]

\subsection{Retículo raíz \texorpdfstring{\(A_2\)}{A₂}}

Sea
\[
 A_2=\mathbb Z\alpha_1\oplus\mathbb Z\alpha_2
\]
el retículo raíz cuya matriz de Gram, en la base ordenada
\((\alpha_1,\alpha_2)\), es
\begin{equation}
 G_{A_2}=
 \begin{pmatrix}
 2&-1\\
 -1&2
 \end{pmatrix}.
 \label{p12-83-c20-s4:eq:gram-a2-autonomo}
\end{equation}
Para \(x=x_1\alpha_1+x_2\alpha_2\),
\begin{equation}
 \langle x,x\rangle
 =2(x_1^2-x_1x_2+x_2^2).
 \label{p12-83-c20-s4:eq:norma-a2-autonoma}
\end{equation}
En particular, \(A_2\) es par, positivo definido y
\[
 \det A_2=3.
\]

Definimos
\[
 \rho:=\alpha_1+\alpha_2.
\]
Entonces
\begin{equation}
 \rho^2=2,
 \qquad
 \langle\rho,\alpha_1\rangle
 =\langle\rho,\alpha_2\rangle=1.
 \label{p12-83-c20-s4:eq:identidades-rho-a2}
\end{equation}

\subsection{Módulo y forma discriminantes}

El retículo dual es
\[
 A_2^*
 :=\{x\in A_2\otimes_{\mathbb Z}\mathbb Q:
       \langle x,A_2\rangle\subseteq\mathbb Z\}.
\]
Elegimos los representantes
\begin{equation}
 \mu_0=0,
 \qquad
 \mu_1=\frac{2\alpha_1+\alpha_2}{3},
 \qquad
 \mu_2=\frac{\alpha_1+2\alpha_2}{3}.
 \label{p12-83-c20-s4:eq:representantes-discriminantes-a2}
\end{equation}
Sus emparejamientos con la base simple son
\[
\begin{array}{c|cc}
 &\alpha_1&\alpha_2\\ \hline
 \mu_1&1&0\\
 \mu_2&0&1
\end{array},
\]
por lo que \(\mu_1,\mu_2\in A_2^*\). Además,
\[
 3\mu_1=2\alpha_1+\alpha_2\in A_2,
 \qquad
 3\mu_2=\alpha_1+2\alpha_2\in A_2,
\]
mientras \(\mu_1,\mu_2\notin A_2\). Se obtiene
\begin{equation}
 A_2^*/A_2
 =\{\overline\mu_0,\overline\mu_1,\overline\mu_2\}
 \cong\mathbb Z/3\mathbb Z
 \cong\mathbb F_3.
 \label{p12-83-c20-s4:eq:discriminante-a2-f3}
\end{equation}

La suma de representantes se reduce explícitamente mediante
\[
 \mu_1+\mu_1-\mu_2=\alpha_1,
 \qquad
 \mu_1+\mu_2=\rho,
 \qquad
 \mu_2+\mu_2-\mu_1=\alpha_2.
\]
Por consiguiente,
\begin{equation}
 \mu_r+\mu_s-\mu_{r+s\,({\rm mod}\,3)}\in A_2
 \qquad(r,s\in\mathbb F_3).
 \label{p12-83-c20-s4:eq:adicion-representantes-discriminantes-a2}
\end{equation}

\begin{proposition}[Forma discriminante de \(A_2\)]
\label{p12-83-c20-s4:prop:forma-discriminante-a2}
Bajo \eqref{p12-83-c20-s4:eq:discriminante-a2-f3}, las formas bilineal y cuadrática
discriminantes son
\begin{align}
 b(\overline\mu_r,\overline\mu_s)
 &\equiv\frac{2rs}{3}\pmod{\mathbb Z},
 \label{p12-83-c20-s4:eq:bilineal-discriminante-a2}\\
 q(\overline\mu_r)
 &\equiv\frac{2r^2}{3}\pmod{2\mathbb Z}.
 \label{p12-83-c20-s4:eq:cuadratica-discriminante-a2}
\end{align}
\end{proposition}

\begin{proof}
La fórmula \eqref{p12-83-c20-s4:eq:norma-a2-autonoma} da
\[
 \mu_1^2=\mu_2^2=\frac23,
 \qquad
 \langle\mu_1,\mu_2\rangle=\frac13.
\]
Para \(r,s\in\{0,1,2\}\), estos valores coinciden, módulo
\(\mathbb Z\), con \(2rs/3\). La fórmula cuadrática resulta al tomar
\(r=s\) y recordar que el retículo es par.
\end{proof}

\subsection{\(12\) factores y aplicación de pegado}

Sea
\[
 L_0:=A_2^{12}.
\]
Entonces
\[
 L_0^*=(A_2^*)^{12},
 \qquad
 L_0^*/L_0\cong\mathbb F_3^{12}.
\]
Para \(c=(c_1,\ldots,c_{12})\in\mathbb F_3^{12}\), definimos
\begin{equation}
 \phi(c):=(\mu_{c_1},\ldots,\mu_{c_{12}})\in L_0^*.
 \label{p12-83-c20-s4:eq:aplicacion-pegado-cw}
\end{equation}
La extensión determinada por \(C_W\) es
\begin{equation}
 \begin{aligned}
 N(C_W)
 &:=\bigcup_{c\in C_W}(L_0+\phi(c))\\
 &=\{x\in L_0^*:x\bmod L_0\in C_W\}.
 \end{aligned}
 \label{p12-83-c20-s4:eq:definicion-niemeier-cw}
\end{equation}

\begin{lemma}[Clausura aditiva del pegado]
\label{p12-83-c20-s4:lem:clausura-pegado-cw}
El conjunto \(N(C_W)\) es un subgrupo discreto de \(L_0^*\) de rango
veinticuatro.
\end{lemma}

\begin{proof}
Sean \(x=\ell+\phi(c)\) y \(y=m+\phi(d)\), con \(\ell,m\in L_0\) y
\(c,d\in C_W\). Por \eqref{p12-83-c20-s4:eq:adicion-representantes-discriminantes-a2},
\[
 \phi(c)+\phi(d)-\phi(c+d)\in L_0.
\]
Como \(C_W\) es lineal, \(c+d\in C_W\), y entonces
\(x+y\in L_0+\phi(c+d)\). La misma identidad aplicada al opuesto da
inversos. Las inclusiones
\[
 L_0\subseteq N(C_W)\subseteq L_0^*
\]
prueban discreción y rango veinticuatro.
\end{proof}

\subsection{Integralidad inducida por autodualidad}

La forma bilineal discriminante de \(L_0\) es la suma ortogonal de doce
copias de \eqref{p12-83-c20-s4:eq:bilineal-discriminante-a2}. Para
\(c,d\in\mathbb F_3^{12}\),
\begin{equation}
 b(c,d)
 \equiv\frac23\sum_{i=1}^{12}c_id_i
 \pmod{\mathbb Z}.
 \label{p12-83-c20-s4:eq:bilineal-discriminante-a2-doce}
\end{equation}

\begin{lemma}[Integralidad del pegado]
\label{p12-83-c20-s4:lem:integralidad-pegado-cw}
El retículo \(N(C_W)\) es integral.
\end{lemma}

\begin{proof}
La autodualidad \(C_W=C_W^\perp\) implica
\[
 \sum_{i=1}^{12}c_id_i\equiv0\pmod3
 \qquad(c,d\in C_W).
\]
Por \eqref{p12-83-c20-s4:eq:bilineal-discriminante-a2-doce}, la forma discriminante se
anula sobre \(C_W\times C_W\). Por tanto, el emparejamiento de cualesquiera
dos clases pegadas es entero.
\end{proof}

\subsection{Paridad inducida por los pesos del código}

En una coordenada no nula del código, la forma cuadrática discriminante
vale \(2/3\) módulo \(2\mathbb Z\). Por tanto,
\begin{equation}
 q(c)\equiv\frac23\operatorname{wt}(c)\pmod{2\mathbb Z}.
 \label{p12-83-c20-s4:eq:cuadratica-discriminante-cw}
\end{equation}

\begin{lemma}[Paridad del pegado]
\label{p12-83-c20-s4:lem:paridad-pegado-cw}
El retículo \(N(C_W)\) es par.
\end{lemma}

\begin{proof}
Los únicos pesos de \(C_W\) son \(0,6,9,12\), y
\[
 \frac23\{0,6,9,12\}=\{0,4,6,8\}\subseteq2\mathbb Z.
\]
Toda clase de pegado es isotrópica para la forma cuadrática discriminante,
de modo que todas las normas de \(N(C_W)\) son pares.
\end{proof}

Sean \(g_1,\ldots,g_6\) las filas de \(G_W=(I_6\mid A_W)\) y
\(\gamma_i=\phi(g_i)\). La multiplicación exacta en \(A_2^*\) da
\begin{equation}
 (\langle\gamma_i,\gamma_j\rangle)_{i,j=1}^{6}
 =
 \begin{pmatrix}
 4&2&2&2&2&2\\
 2&4&2&2&2&2\\
 2&2&4&2&2&2\\
 2&2&2&4&2&2\\
 2&2&2&2&4&2\\
 2&2&2&2&2&4
 \end{pmatrix}.
 \label{p12-83-c20-s4:eq:gram-generadores-pegado-cw}
\end{equation}
Esta matriz proporciona un control directo adicional de integralidad y
paridad para los seis generadores de pegado.

\subsection{Índice, determinante y unimodularidad}

Las clases laterales de \(L_0\) que aparecen en
\eqref{p12-83-c20-s4:eq:definicion-niemeier-cw} están indexadas por los
\(|C_W|=729=3^6\) elementos del código y son distintas. En consecuencia,
\begin{equation}
 [N(C_W):L_0]=3^6.
 \label{p12-83-c20-s4:eq:indice-niemeier-sobre-a2-doce}
\end{equation}
Como
\[
 \det L_0=(\det A_2)^{12}=3^{12},
\]
la fórmula del determinante para una extensión finita da
\begin{equation}
 \det N(C_W)
 =\frac{\det L_0}{[N(C_W):L_0]^2}
 =\frac{3^{12}}{(3^6)^2}=1.
 \label{p12-83-c20-s4:eq:determinante-niemeier-cw}
\end{equation}

\subsection{Sistema de raíces inducido por la incidencia de Witt y clasificación de sus órbitas}

En cada una de las dos clases discriminantes no nulas de
\(A_2^*/A_2\), la norma mínima es
\[
 \mu_1^2=\mu_2^2=\frac23.
\]
Si \(c\in C_W\setminus\{0\}\), entonces
\(\operatorname{wt}(c)\geq6\). Por ortogonalidad de los doce factores,
cualquier vector de \(L_0+\phi(c)\) tiene norma al menos
\begin{equation}
 6\cdot\frac23=4.
 \label{p12-83-c20-s4:eq:cota-clases-no-triviales-niemeier}
\end{equation}
Por tanto, ninguna clase de pegado no trivial contiene vectores de norma
dos. Las raíces de \(N(C_W)\) son exactamente las raíces ya presentes en
\(L_0=A_2^{12}\).

\begin{theorem}[Realización del Niemeier de tipo \(A_2^{12}\)]
\label{p12-83-c20-s4:thm:niemeier-a2-doce-autonomo}
El retículo \(N(C_W)\) es positivo definido, par y unimodular, y su sistema
de raíces es \(A_2^{12}\). En consecuencia,
\begin{equation}
 \boxed{N(C_W)\cong N(A_2^{12}).}
 \label{p12-83-c20-s4:eq:identificacion-niemeier-a2-doce}
\end{equation}
\end{theorem}

\begin{proof}
La positividad procede de \(L_0^*\). La
\cref{p12-83-c20-s4:lem:integralidad-pegado-cw,p12-83-c20-s4:lem:paridad-pegado-cw} demuestra
integralidad y paridad; \eqref{p12-83-c20-s4:eq:determinante-niemeier-cw} demuestra
unimodularidad. La cota \eqref{p12-83-c20-s4:eq:cota-clases-no-triviales-niemeier}
identifica el sistema de raíces. La clasificación de Niemeier asocia de
manera única a ese sistema de raíces el retículo indicado.
\end{proof}

\subsection{Obstrucciones del pegado ternario \(C_W\to N(A_2^{12})\): paridad, unimodularidad y sistema de raíces}

\begin{criterion}[Refutación del pegado ternario]
La construcción queda refutada si ocurre cualquiera de los hechos
siguientes:
\begin{enumerate}
 \item alguno de los representantes \(\mu_0,\mu_1,\mu_2\) no pertenece a
       \(A_2^*\), dos de sus clases coinciden, o falla
       \eqref{p12-83-c20-s4:eq:adicion-representantes-discriminantes-a2};
 \item la forma discriminante no es la dada en
       \cref{p12-83-c20-s4:prop:forma-discriminante-a2};
 \item \(N(C_W)\) no es cerrado bajo suma o inversos;
 \item existe \(c,d\in C_W\) con
       \(\frac23\sum_i c_id_i\notin\mathbb Z\);
 \item una palabra de \(C_W\) produce
       \(\frac23\operatorname{wt}(c)\notin2\mathbb Z\);
 \item el índice no es \(3^6\) o el determinante no es uno;
 \item una clase de pegado no trivial contiene un vector de norma dos;
 \item se utiliza \(W_{24}\), una dodecada o un dato binario en cualquiera
       de las operaciones de pegado anteriores.
\end{enumerate}
\end{criterion}


% Unidad U022. Origen incidencial, vecino ternario y Leech.
% Redaccion autonoma desde la autoridad material 1063.

\section{Origen incidencial, problema radial tipado y construcción del vecino unimodular sin raíces}
\label{p12-83-c20-s5:chap:origen-incidencial-vecino-leech}

El retículo \(N(A_2^{12})\) ya está construido. Esta unidad determina un
punto base a partir de un marco ordenado de incidencias, resuelve un
problema variacional en un dominio declarado, verifica las tres condiciones
del vecino de índice tres y demuestra que el resultado es el retículo de
Leech. El vector del vecino se denota \(v_{\mathcal F_*}\): depende del
marco incidencial y no interviene en el acarreo numérico de \(\alpha\).

\subsection{Marco ordenado de incidencias}

Consideremos las cuatro hexadas
\begin{align}
 H^-&=\{4,5,6,7,9,10\},
 \label{p12-83-c20-s5:eq:hexada-marco-negativa}\\
 H_e&=\{1,3,4,6,9,10\},
 \label{p12-83-c20-s5:eq:hexada-marco-e}\\
 H_\varphi&=\{1,2,3,4,7,9\},
 \label{p12-83-c20-s5:eq:hexada-marco-phi}\\
 H_{\rm APP}&=\{2,3,5,10,11,12\}.
 \label{p12-83-c20-s5:eq:hexada-marco-app}
\end{align}
El dato es el marco ordenado
\begin{equation}
 \mathcal F:=(H^-,H_e,H_\varphi,H_{\rm APP}).
 \label{p12-83-c20-s5:eq:marco-incidencial-ordenado}
\end{equation}
Definimos
\begin{equation}
 o(\mathcal F)
 :=(H_e\cap H_\varphi)\setminus(H^-\cup H_{\rm APP}).
 \label{p12-83-c20-s5:eq:operador-origen-incidencial}
\end{equation}

\begin{theorem}[Reconstrucción equivariante del origen]
\label{p12-83-c20-s5:thm:origen-incidencial-equivariante}
Para el marco canónico \(\mathcal F_*\),
\begin{equation}
 \boxed{o(\mathcal F_*)=\{1\}.}
 \label{p12-83-c20-s5:eq:origen-incidencial-singleton}
\end{equation}
Además, para toda permutación simultánea \(g\in M_{12}\),
\begin{equation}
 o(g\mathcal F)=g\,o(\mathcal F).
 \label{p12-83-c20-s5:eq:equivarianza-origen-incidencial}
\end{equation}
\end{theorem}

\begin{proof}
Se tiene
\[
 H_e\cap H_\varphi=\{1,3,4,9\}
\]
y
\[
 (H^-\cup H_{\rm APP})\cap\{1,3,4,9\}=\{3,4,9\}.
\]
La diferencia es \(\{1\}\). Para la equivarianza se aplican las
identidades
\[
 g(A\cap B)=gA\cap gB,
 \quad
 g(A\cup B)=gA\cup gB,
 \quad
 g(A\setminus B)=gA\setminus gB,
\]
válidas para toda biyección.
\end{proof}

La tétrada \(B_K=\{4,6,9,10\}\) no determina por sí sola el origen. El
singleton se obtiene del marco ordenado completo; no es una coordenada
elegida después de observar el resultado.

\subsection{Cámara radial positiva del sistema de raíces y selección del cono dominante}

Sea \(\rho_i\) la copia de \(\rho=\alpha_1+\alpha_2\) en el
\(i\)-ésimo factor de \(A_2^{12}\). Consideramos el dominio
\begin{equation}
 \mathscr R_+
 :=\left\{
 v(a)=\sum_{i=1}^{12}a_i\rho_i:
 a_i=1+3b_i, b_i\in\mathbb Z_{\geq0}
 \right\}.
 \label{p12-83-c20-s5:eq:camara-radial-positiva}
\end{equation}
Como los factores son ortogonales y \(\rho_i^2=2\), si
\(S=\sum_i b_i\), entonces
\begin{align}
 v(a)^2
 &=2\sum_{i=1}^{12}(1+3b_i)^2\notag\\
 &=24+12S+18\sum_{i=1}^{12}b_i^2.
 \label{p12-83-c20-s5:eq:norma-radial-expandida}
\end{align}
Reduciendo módulo dieciocho,
\[
 v(a)^2\equiv6+12S\pmod{18},
\]
y por tanto
\begin{equation}
 v(a)^2\equiv0\pmod{18}
 \quad\Longleftrightarrow\quad
 S\equiv1\pmod3.
 \label{p12-83-c20-s5:eq:congruencia-radial-suma-b}
\end{equation}

\begin{proposition}[Mínimos en la cámara radial declarada]
\label{p12-83-c20-s5:prop:minimos-camara-radial}
En el subconjunto de \(\mathscr R_+\) definido por
\(v(a)^2\equiv0\pmod{18}\), la norma mínima es \(54\). Se alcanza
exactamente en doce vectores: para cada \(j\in\{1,\ldots,12\}\), el
coeficiente \(a_j\) vale cuatro y los once restantes valen uno.
\end{proposition}

\begin{proof}
La condición \eqref{p12-83-c20-s5:eq:congruencia-radial-suma-b} y la no negatividad de
los \(b_i\) implican que el menor valor posible de \(S\) es uno. Esto
obliga a un único índice \(j\) con \(b_j=1\) y a \(b_i=0\) para
\(i\neq j\). Sustituyendo en \eqref{p12-83-c20-s5:eq:norma-radial-expandida},
\[
 v(a)^2=24+12+18=54.
\]
Hay doce elecciones de \(j\). El siguiente valor admisible de \(S\) es
cuatro y produce una norma estrictamente mayor.
\end{proof}

El origen \(o(\mathcal F_*)=\{1\}\) selecciona
\begin{equation}
 \boxed{
 v_{\mathcal F_*}
 =4\rho_1+\rho_2+\cdots+\rho_{12},
 \qquad
 v_{\mathcal F_*}^2=54.}
 \label{p12-83-c20-s5:eq:vector-marco-incidencial}
\end{equation}

\begin{remark}[Alcance exacto de la minimalidad]
La norma \(54\) es mínima en la cámara radial positiva
\eqref{p12-83-c20-s5:eq:camara-radial-positiva}, no en toda la clase residual. El vector
\[
 u=(\alpha_1-2\alpha_2,\rho,\ldots,\rho)
\]
satisface
\[
 \alpha_1-2\alpha_2-\rho=-3\alpha_2,
\]
de modo que
\[
 u\equiv(\rho,\ldots,\rho)
 \equiv v_{\mathcal F_*}\pmod{3A_2^{12}}.
\]
Sin embargo,
\[
 (\alpha_1-2\alpha_2)^2
 =2(1+2+4)=14,
 \qquad
 u^2=14+11\cdot2=36.
\]
Este testigo impide extender la afirmación de minimalidad fuera del dominio
de \cref{p12-83-c20-s5:prop:minimos-camara-radial}.
\end{remark}

\subsection{Carácter y núcleo de congruencia}

Escribimos
\[
 N:=N(A_2^{12}),
 \qquad
 v:=v_{\mathcal F_*}.
\]
Definimos
\begin{equation}
 \chi_v:N\longrightarrow\mathbb Z/3\mathbb Z,
 \qquad
 \chi_v(x)=\langle x,v\rangle\bmod3,
 \label{p12-83-c20-s5:eq:caracter-vecino-leech}
\end{equation}
su núcleo
\begin{equation}
 N_0:=\ker\chi_v
 =\{x\in N:\langle x,v\rangle\equiv0\pmod3\},
 \label{p12-83-c20-s5:eq:nucleo-congruencia-vecino}
\end{equation}
y la extensión
\begin{equation}
 N_v:=N_0+\mathbb Z\frac v3.
 \label{p12-83-c20-s5:eq:definicion-vecino-indice-tres}
\end{equation}

\subsection{Condición I1: sobreyectividad y exclusión de raíces}

Por construcción, \(v\in A_2^{12}\subseteq N\). Las raíces de un factor
\(A_2\) son
\[
 \pm\alpha_1,
 \quad
 \pm\alpha_2,
 \quad
 \pm\rho.
\]
Sus emparejamientos con \(\rho\) son \(\pm1\) o \(\pm2\). Todos los
coeficientes de \(v\) son congruentes con uno módulo tres:
\[
 4\equiv1\pmod3,
 \qquad
 1\equiv1\pmod3.
\]
Por tanto, ninguna raíz de \(A_2^{12}\) tiene emparejamiento nulo con
\(v\) módulo tres. En particular,
\[
 \chi_v(\alpha_{1,1})=1,
\]
de modo que \(\chi_v\) es sobreyectivo y
\begin{equation}
 [N:N_0]=3.
 \label{p12-83-c20-s5:eq:indice-nucleo-vecino}
\end{equation}
Además, ninguna raíz de \(N\) pertenece a \(N_0\).

\subsection{Condición I2: divisibilidad de la norma}

La identidad \eqref{p12-83-c20-s5:eq:vector-marco-incidencial} da
\begin{equation}
 v^2=54\equiv0\pmod{18},
 \qquad
 \left(\frac v3\right)^2=6.
 \label{p12-83-c20-s5:eq:condicion-i2-vecino}
\end{equation}
En particular, \(\langle v,v\rangle\equiv0\pmod3\), luego
\(v\in N_0\).

\subsection{Condición I3: clase dual de orden \(3\)}

Si \(x\in N_0\), existe \(k\in\mathbb Z\) con
\(\langle x,v\rangle=3k\). Por tanto,
\[
 \left\langle x,\frac v3\right\rangle=k\in\mathbb Z,
\]
y
\begin{equation}
 \frac v3\in N_0^*.
 \label{p12-83-c20-s5:eq:v-tercio-en-dual-nucleo}
\end{equation}
Por otra parte,
\[
 \left\langle\frac v3,\alpha_{1,1}\right\rangle=\frac43.
\]
Como \(\alpha_{1,1}\in N\) y \(N\) es integral, esto demuestra
\begin{equation}
 \frac v3\notin N.
 \label{p12-83-c20-s5:eq:v-tercio-no-en-niemeier}
\end{equation}
Finalmente, \(3(v/3)=v\in N_0\). La clase de \(v/3\) en
\(N_0^*/N_0\) es no trivial y tiene orden exactamente tres. En
consecuencia,
\begin{equation}
 [N_v:N_0]=3.
 \label{p12-83-c20-s5:eq:indice-extension-vecino}
\end{equation}

\subsection{Paridad, integralidad y determinante del vecino}

\begin{proposition}[Estructura unimodular del vecino]
\label{p12-83-c20-s5:prop:estructura-unimodular-vecino}
El retículo \(N_v\) es integral, par y unimodular.
\end{proposition}

\begin{proof}
Como \(N\) es unimodular y \([N:N_0]=3\),
\[
 \det N_0=[N:N_0]^2\det N=9.
\]
Como \([N_v:N_0]=3\),
\[
 \det N_v=\frac{\det N_0}{[N_v:N_0]^2}=1.
\]

Sea \(y=x+m(v/3)\), con \(x\in N_0\), y escribamos
\(\langle x,v\rangle=3k\). Entonces
\begin{align*}
 y^2
 &=x^2+2m\left\langle x,\frac v3\right\rangle
       +m^2\left(\frac v3\right)^2\\
 &=x^2+2mk+6m^2.
\end{align*}
Los tres sumandos forman un entero par. La forma polarizada prueba
integralidad, y el determinante uno prueba unimodularidad.
\end{proof}

\subsection{\(2\) clases globales y \(6\) clases locales}

Globalmente, \eqref{p12-83-c20-s5:eq:definicion-vecino-indice-tres} añade sólo dos clases
laterales no triviales:
\[
 N_0+\frac v3,
 \qquad
 N_0-\frac v3.
\]
Localmente, en cada factor \(A_2\), una coordenada de esas clases pertenece
a una de las seis clases
\begin{equation}
 \mathscr C_{j,\pm}
 :=A_2+\mu_j\pm\frac\rho3,
 \qquad j\in\{0,1,2\}.
 \label{p12-83-c20-s5:eq:seis-clases-locales-vecino}
\end{equation}
En la primera coordenada, \(4\rho/3\equiv\rho/3\pmod{A_2}\), por lo que
la misma lista vale para los doce factores.

\begin{table}[htbp]
\centering
\caption{Representantes mínimos de las seis clases laterales locales.}
\label{p12-83-c20-s5:tab:seis-clases-locales-vecino}
\begin{tabular}{cclc}
\toprule
clase&representante mínimo&\((u,v)\)&
\(2(u^2-uv+v^2)/9\)\\
\midrule
\(\mathscr C_{0,-}\)&\(-\rho/3\)&\((-1,-1)\)&\(2/9\)\\
\(\mathscr C_{0,+}\)&\( \rho/3\)&\((1,1)\)&\(2/9\)\\
\(\mathscr C_{1,-}\)&\( \mu_1-\rho/3\)&\((1,0)\)&\(2/9\)\\
\(\mathscr C_{1,+}\)&\( \mu_1+\rho/3-\rho\)&\((0,-1)\)&\(2/9\)\\
\(\mathscr C_{2,-}\)&\( \mu_2-\rho/3\)&\((0,1)\)&\(2/9\)\\
\(\mathscr C_{2,+}\)&\( \mu_2+\rho/3-\rho\)&\((-1,0)\)&\(2/9\)\\
\bottomrule
\end{tabular}
\end{table}

Aquí \((u,v)\) significa el representante
\((u\alpha_1+v\alpha_2)/3\).

\begin{lemma}[Mínimo exacto de las seis clases]
\label{p12-83-c20-s5:lem:minimo-seis-clases-locales}
Cada clase de \eqref{p12-83-c20-s5:eq:seis-clases-locales-vecino} tiene norma mínima
exactamente \(2/9\).
\end{lemma}

\begin{proof}
La norma de \((u\alpha_1+v\alpha_2)/3\) es
\[
 \frac29(u^2-uv+v^2).
\]
La forma entera \(u^2-uv+v^2\) es positiva definida. En una clase residual
no nula no toma el valor cero, luego es al menos uno. Los seis
representantes de la \cref{p12-83-c20-s5:tab:seis-clases-locales-vecino} alcanzan uno.
\end{proof}

\subsection{Criterio de ausencia de raíces en el retículo de Leech y consecuencia para la distancia mínima}

\begin{theorem}[Ausencia de vectores de norma dos]
\label{p12-83-c20-s5:thm:vecino-sin-raices}
El retículo \(N_v\) no contiene vectores de norma dos.
\end{theorem}

\begin{proof}
Todas las raíces de \(N\) son las de \(A_2^{12}\). La condición I1 prueba
que ninguna pertenece a \(N_0\). Por tanto, la clase \(N_0\) no contiene
vectores de norma dos.

En una clase nueva \(N_0\pm v/3\), cada una de las doce coordenadas
pertenece a una de las seis clases locales. Por
\cref{p12-83-c20-s5:lem:minimo-seis-clases-locales}, cada coordenada contribuye al menos
\(2/9\). La ortogonalidad da
\[
 \left\|x\pm\frac v3\right\|^2
 \geq12\cdot\frac29=\frac83>2.
\]
Tampoco las clases nuevas contienen raíces.
\end{proof}

\subsection{Identificación con el retículo de Leech}

\begin{theorem}[Construcción del retículo de Leech]
\label{p12-83-c20-s5:thm:identificacion-vecino-leech}
El retículo
\begin{equation}
 N_v
 =\ker(\langle\,\cdot\,,v_{\mathcal F_*}\rangle\bmod3)
  +\mathbb Z\frac{v_{\mathcal F_*}}3
 \label{p12-83-c20-s5:eq:formula-final-vecino-leech}
\end{equation}
es isométrico al retículo de Leech:
\begin{equation}
 \boxed{N_v\cong\Lambda_{24}.}
 \label{p12-83-c20-s5:eq:identificacion-leech-autonoma}
\end{equation}
\end{theorem}

\begin{proof}
El retículo \(N_v\) tiene rango veinticuatro y es positivo definido. La
\cref{p12-83-c20-s5:prop:estructura-unimodular-vecino} demuestra que es par y unimodular;
el \cref{p12-83-c20-s5:thm:vecino-sin-raices} demuestra que carece de raíces. Por la
clasificación de Niemeier, existe una única clase de isometría de retículos
positivos definidos, pares y unimodulares de rango veinticuatro sin raíces:
el retículo de Leech.
\end{proof}

\subsection{Los 243 síndromes del código ternario perforado}

Sea \(G_{11}\) el código ternario de Golay de parámetros
\([11,6,5]_3\). Su espacio de síndromes tiene cardinal
\begin{equation}
 3^{11-6}=3^5=243.
 \label{p12-83-c20-s5:eq:numero-sindromes-golay-ternario}
\end{equation}
El número de errores de peso a lo sumo dos es
\begin{equation}
 1+2\binom{11}{1}+2^2\binom{11}{2}
 =1+22+220=243.
 \label{p12-83-c20-s5:eq:bola-radio-dos-golay-ternario}
\end{equation}

\begin{proposition}[Correspondencia perfecta error--síndrome]
\label{p12-83-c20-s5:prop:correspondencia-error-sindrome-golay}
Cada síndrome de \(G_{11}\) posee un único representante de error de peso
a lo sumo dos.
\end{proposition}

\begin{proof}
La distancia mínima cinco implica que dos errores distintos de peso a lo
sumo dos no pueden diferir en una palabra del código: su diferencia tendría
peso a lo sumo cuatro. La aplicación error--síndrome es, por tanto,
inyectiva sobre la bola de radio dos. Los dos conjuntos tienen cardinal
\(243\) por \eqref{p12-83-c20-s5:eq:numero-sindromes-golay-ternario} y
\eqref{p12-83-c20-s5:eq:bola-radio-dos-golay-ternario}; la aplicación es biyectiva.
\end{proof}

Estos \(243\) síndromes no son las \(243\) palabras visibles obtenidas del
catálogo TPK. Comparten un cardinal, pero poseen dominios, equivalencias y
funciones distintas.

\subsection{Obstrucciones del vecino ternario \(N_v\cong\Lambda_{24}\): clase de orden \(3\), ausencia de raíces y separación de los cardinales \(243\)}

\begin{criterion}[Refutación del vecino y de la identificación de Leech]
La construcción queda refutada si ocurre cualquiera de los hechos
siguientes:
\begin{enumerate}
 \item el operador \(o(\mathcal F)\) no produce \(\{1\}\) o no es
       equivariante;
 \item existe en la cámara radial un vector admisible de norma menor que
       \(54\), o el número de minimizadores no es doce;
 \item se afirma minimalidad global pese al testigo de norma \(36\);
 \item \(\chi_v\) no es sobreyectivo o una raíz pertenece a \(N_0\);
 \item \(v^2\not\equiv0\pmod{18}\) o \((v/3)^2\neq6\);
 \item \(v/3\notin N_0^*\), \(v/3\in N\), o su clase no tiene orden tres;
 \item el determinante de \(N_v\) no es uno o el vecino deja de ser par;
 \item alguna de las seis clases locales tiene mínimo menor que \(2/9\);
 \item una de las dos clases globales nuevas contiene una raíz;
 \item se confunden las dos clases globales con las seis clases locales;
 \item se utiliza \(W_{24}\) como antecedente del pegado o del vecino;
 \item los \(243\) síndromes se identifican con las \(243\) palabras
       visibles por compartir cardinal.
\end{enumerate}
\end{criterion}


\section{Holonomía finita, configuración de Schläfli y simetría \texorpdfstring{\(E_6\)}{E6}}
\label{p12-83-c20-s6:chap:holonomia-schlafli-e6-nuevo}
\index{holonomía finita}
\index{Schläfli}
\index{E6@\(E_6\)}

La rama de Witt no agota las estructuras excepcionales accesibles desde
APP. La rama paralela de este capítulo parte de las dos polarizaciones
mixtas suma--producto y producto--suma, junto con una orientación de
plaqueta y una sección central declaradas. La curvatura calculada determina
la forma alternante; ésta define la extensión central. No se infiere una
extensión a partir de la sola cardinalidad 27.

\subsection{Curvaturas mixtas de APP}

Sobre \(T_9=(\mathbb Z/9\mathbb Z)^2\), con
\(S(x,y)=x+y\) y \(P(x,y)=xy\), sean
\[
 \Delta_xf=f(x+1,y)-f(x,y),\qquad
 \Delta_yf=f(x,y+1)-f(x,y).
\]
Las conexiones mixtas son
\[
 A^{SP}=(\Delta_xS,\Delta_yP)=(1,x),\qquad
 A^{PS}=(\Delta_xP,\Delta_yS)=(y,1).
\]
Su curvatura discreta es
\[
 F_A=\Delta_xA_y-\Delta_yA_x.
\]

\begin{proposition}[Curvaturas orientadas opuestas]
\label{p12-83-c20-s6:prop:curvaturas-app-e6-nuevo}
\[
 F_{A^{SP}}=1,\qquad F_{A^{PS}}=-1\pmod9.
\]
La circulación sobre un rectángulo orientado de lados \(r,s\) es,
respectivamente, \(rs\) y \(-rs\).
\end{proposition}

\begin{proof}
\(\Delta_xx-\Delta_y1=1\) y
\(\Delta_x1-\Delta_yy=-1\). Un rectángulo contiene \(rs\) plaquetas; al
sumar, las aristas interiores se cancelan y queda la circulación del borde.
\end{proof}

Para desplazamientos \(u=(a,b)\), \(v=(c,d)\), la clase alternante es
\begin{equation}
 \sigma(u,v)=ad-bc\pmod9.
 \label{p12-83-c20-s6:eq:forma-alternante-e6-nueva}
\end{equation}
La orientación fija el signo de \(\sigma\). La sección central empleada
abajo pertenece al calibre de esta realización; cambiarla por un coborde
produce una extensión isomorfa, no un nuevo invariante.

\subsection{Extensión central y reducción ternaria}

Definimos el grupo de Heisenberg nonádico
\begin{equation}
 \mathcal H_9=\{(a,b,z):a,b,z\in\mathbb Z/9\mathbb Z\},
 \quad
 (a,b,z)(c,d,w)=(a+c,b+d,z+w+bc).
 \label{p12-83-c20-s6:eq:heisenberg-9-nuevo}
\end{equation}
El conmutador de dos desplazamientos horizontales es
\[
 [(a,b,0),(c,d,0)]=(0,0,bc-ad),
\]
que registra \(-\sigma(u,v)\). Reduciendo cada coordenada módulo tres se
obtiene
\begin{equation}
 H_3=\{(a,b,z):a,b,z\in\mathbb F_3\},\qquad |H_3|=27,
 \label{p12-83-c20-s6:eq:heisenberg-3-nuevo}
\end{equation}
con la misma ley. El núcleo de
\(\mathcal H_9\to H_3\) es
\((3\mathbb Z/9\mathbb Z)^3\), de cardinal 27.

\subsection{Representación de Weyl--Heisenberg y fase central}

La extensión nonádica admite una realización unitaria explícita. Sea
\[
 \zeta_9:=\exp(2\pi i/9),
 \qquad
 \mathscr H_9^{\rm Sch}:=\mathbb C^9,
\]
con base \(\{|n\rangle:n\in\mathbb Z/9\mathbb Z\}\). Definimos
\begin{equation}
 X|n\rangle:=|n+1\rangle,
 \qquad
 Z|n\rangle:=\zeta_9^n|n\rangle.
 \label{p12-83-c20-s6:eq:weyl-desplazamiento-fase-u029}
\end{equation}
Se tiene
\[
 ZX=\zeta_9XZ
\]
y, por multiplicación,
\begin{equation}
 (X^aZ^b)(X^cZ^d)
 =\zeta_9^{bc}X^{a+c}Z^{b+d}.
 \label{p12-83-c20-s6:eq:cociclo-weyl-nonadico-u029}
\end{equation}

\begin{theorem}[Representación fiel de Weyl--Heisenberg]
\label{p12-83-c20-s6:thm:representacion-weyl-heisenberg-u029}
La aplicación
\begin{equation}
 \rho_{\rm WH}:\mathcal H_9\longrightarrow
 U(\mathscr H_9^{\rm Sch}),
 \qquad
 \rho_{\rm WH}(a,b,z):=\zeta_9^zX^aZ^b,
 \label{p12-83-c20-s6:eq:representacion-weyl-heisenberg-u029}
\end{equation}
es una representación unitaria fiel. Para
\(u=(a,b)\) y \(v=(c,d)\), su conmutador satisface
\begin{equation}
 [\rho_{\rm WH}(u,0),\rho_{\rm WH}(v,0)]
 =\zeta_9^{\,bc-ad}I
 =\zeta_9^{-\sigma(u,v)}I.
 \label{p12-83-c20-s6:eq:conmutador-weyl-holonomia-u029}
\end{equation}
\end{theorem}

\begin{proof}
La identidad \eqref{p12-83-c20-s6:eq:cociclo-weyl-nonadico-u029} reproduce el término
\(bc\) de \eqref{p12-83-c20-s6:eq:heisenberg-9-nuevo}; de ahí la propiedad de
homomorfismo. Si \(\zeta_9^zX^aZ^b=I\), la acción sobre la base obliga
sucesivamente a \(a=0\), \(b=0\) y \(z=0\). Esto prueba fidelidad. Comparar
los productos en los dos órdenes da
\(\zeta_9^{bc-ad}\), que es la fase alternante de
\eqref{p12-83-c20-s6:eq:forma-alternante-e6-nueva}.
\end{proof}

La representación distingue dos observables conjugados. Sean
\[
 Q_{\rm H}:=\operatorname{diag}(0,1,\ldots,8),
 \qquad
 F_{jk}:=\frac13\zeta_9^{jk},
 \qquad
 P_{\rm H}:=F^\dagger Q_{\rm H}F.
\]
Las bases propias de \(Q_{\rm H}\) y \(P_{\rm H}\) son mutuamente
insesgadas, pues \(|F_{jk}|^2=1/9\). Para todo vector unitario
\(\psi\in\mathscr H_9^{\rm Sch}\),
\begin{align}
 \operatorname{Var}_\psi(Q_{\rm H})
 \operatorname{Var}_\psi(P_{\rm H})
 &\ge
 \frac14\left|
 \langle\psi,[Q_{\rm H},P_{\rm H}]\psi\rangle
 \right|^2,
 \label{p12-83-c20-s6:eq:incertidumbre-weyl-u029}\\
 H_{Q_{\rm H}}(\psi)+H_{P_{\rm H}}(\psi)
 &\ge\log9.
 \label{p12-83-c20-s6:eq:incertidumbre-entropica-weyl-u029}
\end{align}
Estos operadores son adimensionales. Una carta física de unidades es una
realización posterior y no interviene en la extensión central.

\subsection{Acción de \(10\) desplazamientos sobre el grafo de Cayley}

Para un funcional lineal
\(\ell(a,b)=pa+qb\), ponemos
\[
 D_\ell(a,b)=\left(a,b,\frac{ab}{2}+\ell(a,b)\right),
 \qquad Z=(0,0,1),
\]
donde \(2^{-1}=2\) en \(\mathbb F_3\). El conjunto inverso-cerrado
\begin{equation}
 S_\ell=
 \{D_\ell(v):v\in\mathbb F_3^2\setminus\{0\}\}
 \cup\{Z,Z^{-1}\}
 \label{p12-83-c20-s6:eq:conjunto-conexion-e6}
\end{equation}
tiene diez elementos y no contiene la identidad.

\begin{theorem}[Independencia del calibre lineal]
\label{p12-83-c20-s6:thm:independencia-calibre-e6}
Los nueve funcionales \(\ell\) producen grafos de Cayley isomorfos. Los
nueve conjuntos \(S_\ell\) forman una sola órbita bajo
\(\operatorname{Aut}(H_3)\).
\end{theorem}

\begin{proof}
La aplicación
\[
 \phi_\ell(a,b,z)=(a,b,z+\ell(a,b))
\]
es un automorfismo central y satisface
\(\phi_\ell(S_0)=S_\ell\). Para escribir la familia completa en las
coordenadas polarizadas de \cref{p12-83-c20-s6:eq:heisenberg-9-nuevo}, sea
\(c(v,w)=b c\) para \(v=(a,b)\), \(w=(c,d)\). Dado
\(M\in\operatorname{GL}(2,3)\), el defecto
\[
 B_M(v,w)=c(Mv,Mw)-\det(M)c(v,w)
\]
es simétrico, porque las partes alternantes de los dos términos coinciden.
Definimos la corrección cuadrática
\begin{equation}
 q_M(v)=2B_M(v,v),
 \qquad
 q_M(v+w)-q_M(v)-q_M(w)=B_M(v,w),
 \label{p12-83-c20-s6:eq:correccion-cuadratica-automorfismos-h3}
\end{equation}
donde \(2=2^{-1}\) en \(\mathbb F_3\). La familia completa de
automorfismos se escribe entonces
\begin{equation}
 (v,t)\longmapsto
 \bigl(Mv,\det(M)t+q_M(v)+\ell(v)\bigr),
 \qquad M\in\operatorname{GL}(2,3),\
 \ell\in(\mathbb F_3^2)^*.
 \label{p12-83-c20-s6:eq:automorfismos-h3-completos}
\end{equation}
La identidad polar de \cref{p12-83-c20-s6:eq:correccion-cuadratica-automorfismos-h3}
prueba directamente que el producto central se conserva. La familia tiene
\(48\cdot9=432\) elementos. La enumeración exhaustiva de los
\(\binom{13}{5}=1287\) subconjuntos inverso-cerrados de cardinal diez
encuentra exactamente esos nueve con los parámetros que se derivan a
continuación.
\end{proof}

Escribamos \(\underline S=\sum_{s\in S}s\) en
\(\mathbb Z[H_3]\). El conteo de productos ordenados da
\begin{equation}
 \underline S^{\,2}
 =10e+\underline S+5(\underline{H_3}-e-\underline S)
 =5\underline{H_3}+5e-4\underline S.
 \label{p12-83-c20-s6:eq:identidad-grupo-schlafli-nueva}
\end{equation}
En efecto, la identidad recibe diez factorizaciones; un elemento de
\(S\) recibe una; cada uno de los otros dieciséis recibe cinco.

\begin{theorem}[Grafo de las veintisiete rectas]
\label{p12-83-c20-s6:thm:grafo-27-rectas-nuevo}
La matriz de adyacencia \(A_\ell\) de
\(\Gamma_\ell=\operatorname{Cay}(H_3,S_\ell)\) satisface
\begin{equation}
 A_\ell^2=5I-4A_\ell+5J.
 \label{p12-83-c20-s6:eq:identidad-schlafli-nueva}
\end{equation}
Por tanto,
\[
 \Gamma_\ell=\operatorname{srg}(27,10,1,5),
 \qquad
 \operatorname{Spec}(A_\ell)=\{10^1,1^{20},(-5)^6\}.
\]
Es isomorfo al grafo de intersección de las veintisiete rectas de una
superficie cúbica lisa.
\end{theorem}

\begin{proof}
La representación regular de
\cref{p12-83-c20-s6:eq:identidad-grupo-schlafli-nueva} produce
\cref{p12-83-c20-s6:eq:identidad-schlafli-nueva}. Sus entradas dicen que todo vértice
tiene diez vecinos, que un par adyacente tiene un vecino común y un par no
adyacente tiene cinco. Sobre la recta constante, \(A_\ell\) actúa por 10;
en su complemento, \(J=0\) y
\((A_\ell-I)(A_\ell+5I)=0\). Traza cero fija las multiplicidades 20 y 6.

En el modelo de una cúbica obtenida al soplar seis puntos de
\(\mathbb P^2\), las 27 clases son
\[
 E_i,\qquad L_{ij}=H-E_i-E_j,\qquad
 Q_i=2H-\sum_{j\ne i}E_j.
\]
Con \(H^2=1\), \(E_i^2=-1\), \(H\cdot E_i=0\) y
\(E_i\cdot E_j=0\) para \(i\ne j\), su matriz de intersección tiene los
mismos parámetros. La identificación con la configuración clásica usa el
teorema de unicidad de esta realización de las 27 rectas
\cite{Manivel2005}; no se atribuye a un certificado inexistente.
\end{proof}

La afirmación de isomorfía puede hacerse completamente efectiva mediante
una biyección calculada. El etiquetado no es absoluto: otra sección lineal
o una automorfía de la superficie produce otra tabla de la misma clase.
\begin{table}[!htbp]
\centering\fontsize{9.4}{10.8}\selectfont
\setlength{\tabcolsep}{4pt}\renewcommand{\arraystretch}{1.04}
\caption{Bijección explícita entre \(H_3\) y las veintisiete clases de
rectas.}
\label{p12-83-c20-s6:tab:h3-27-rectas-u029}
\begin{tabular}{@{}c@{\,}c@{\,}c@{\qquad}l@{}}

\toprule
\(a\)&\(b\)&\(z\)&clase de recta\\
\midrule

0&0&0&\(E_1\)\\
0&0&1&\(L_{12}\)\\
0&0&2&\(Q_2\)\\
0&1&0&\(L_{13}\)\\
0&1&1&\(Q_1\)\\
0&1&2&\(E_3\)\\
0&2&0&\(Q_3\)\\
0&2&1&\(E_2\)\\
0&2&2&\(L_{23}\)\\
1&0&0&\(L_{14}\)\\
1&0&1&\(L_{35}\)\\
1&0&2&\(L_{26}\)\\
1&1&0&\(L_{36}\)\\
1&1&1&\(L_{24}\)\\
1&1&2&\(L_{15}\)\\
1&2&0&\(L_{25}\)\\
1&2&1&\(L_{16}\)\\
1&2&2&\(L_{34}\)\\
2&0&0&\(Q_4\)\\
2&0&1&\(L_{46}\)\\
2&0&2&\(E_6\)\\
2&1&0&\(E_5\)\\
2&1&1&\(Q_6\)\\
2&1&2&\(L_{56}\)\\
2&2&0&\(L_{45}\)\\
2&2&1&\(E_4\)\\
2&2&2&\(Q_5\)\\
\bottomrule
\end{tabular}
\end{table}

Como control local, los diez vecinos de \(e=(0,0,0)\) son los elementos de
\(S_0\). La tabla los transforma en las diez clases que cortan a \(E_1\):
las cinco \(L_{1j}\) y las cinco \(Q_j\), con \(2\le j\le6\). La
verificación completa compara las 729 entradas de ambas matrices de
adyacencia.

El complemento \(B_\ell=J-I-A_\ell\) satisface
\begin{equation}
 B_\ell^2=8I+2B_\ell+8J,
 \qquad
 \operatorname{srg}(27,16,10,8),
 \qquad
 \operatorname{Spec}(B_\ell)=\{16^1,4^6,(-2)^{20}\}.
 \label{p12-83-c20-s6:eq:complemento-schlafli}
\end{equation}
Éste es el grafo de Schläfli en la convención de adyacencia complementaria;
\(A_\ell\) es el grafo de intersección de parámetros
\((27,10,1,5)\).

Cada arista pertenece a un triángulo único, pues \(\lambda=1\). El número
de triángulos es
\[
 \frac{27\cdot10}{6}=45.
\]
El grupo completo de automorfismos tiene orden
\[
 27\cdot1920=51840
\]
y, por su acción sobre la configuración de rectas, se identifica con
\(W(E_6)\), no sólo por igualdad de órdenes \cite{Manivel2005}.

\subsection{El retículo de raíces que realiza \texorpdfstring{\(E_6\)}{E6}}

Para hacer explícito el objeto sobre el que actúa ese grupo, sea
\begin{equation}
 \operatorname{Pic}(X)
 =\mathbb ZH\oplus\bigoplus_{i=1}^6\mathbb ZE_i,
 \qquad
 H^2=1,\qquad E_i^2=-1,\qquad H\cdot E_i=E_i\cdot E_j=0\ (i\ne j),
 \label{p12-83-c20-s6:eq:picard-cubica-e6}
\end{equation}
el retículo de Picard de la explosión de seis puntos generales de
\(\mathbb P^2\). Su clase canónica es
\[
 K_X=-3H+E_1+\cdots+E_6.
\]
El complemento ortogonal
\begin{equation}
 L_{E_6}=K_X^\perp\subset\operatorname{Pic}(X),
 \qquad (x,y)_{E_6}:=-x\cdot y,
 \label{p12-83-c20-s6:eq:reticulo-raices-e6}
\end{equation}
es positivo definido de rango seis. Una base de raíces simples es
\begin{equation}
 \begin{aligned}
 \alpha_1&=E_1-E_2,& \alpha_2&=E_2-E_3,&
 \alpha_3&=E_3-E_4,\\
 \alpha_4&=E_4-E_5,& \alpha_5&=E_5-E_6,&
 \alpha_6&=H-E_1-E_2-E_3.
 \end{aligned}
 \label{p12-83-c20-s6:eq:raices-simples-e6}
\end{equation}
Cada \(\alpha_i\) es ortogonal a \(K_X\), tiene norma dos y sus productos
son \(-1\) exactamente para los pares adyacentes del diagrama
\(E_6\): la cadena \(1-2-3-4-5\), con el nodo \(6\) unido al nodo \(3\).
Por tanto, su matriz de Gram es la matriz de Cartan de \(E_6\).

\begin{proposition}[Las setenta y dos raíces y sus reflexiones]
\label{p12-83-c20-s6:prop:setenta-dos-raices-e6}
Las raíces de \(L_{E_6}\) son
\begin{equation}
 \begin{aligned}
 &E_i-E_j &&(i\ne j),\\
 &\pm(H-E_i-E_j-E_k) &&(1\le i<j<k\le6),\\
 &\pm(2H-E_1-\cdots-E_6),
 \end{aligned}
 \label{p12-83-c20-s6:eq:enumeracion-raices-e6}
\end{equation}
en total \(30+40+2=72\). Para cada raíz \(\alpha\),
\begin{equation}
 s_\alpha(x)=x+(x\cdot\alpha)\alpha
 \label{p12-83-c20-s6:eq:reflexion-raiz-e6}
\end{equation}
preserva la intersección y \(K_X\). Las seis reflexiones simples generan
\(W(E_6)\) y permutan fielmente las 27 clases
\(E_i,L_{ij},Q_i\) de \cref{p12-83-c20-s6:thm:grafo-27-rectas-nuevo}.
\end{proposition}

\begin{proof}
Sustituir cada clase de \cref{p12-83-c20-s6:eq:enumeracion-raices-e6} en
\cref{p12-83-c20-s6:eq:picard-cubica-e6} da \(K_X\cdot\alpha=0\) y
\(\alpha^2=-2\). La cuenta de clases es la indicada. La fórmula de reflexión
es la reflexión ortogonal para la forma \(-(\cdot)\); preserva \(K_X\) porque
\(K_X\cdot\alpha=0\). El sistema simple de
\cref{p12-83-c20-s6:eq:raices-simples-e6} genera todas las raíces y su grupo de Coxeter es
\(W(E_6)\). Su acción sobre las 27 clases de curvas excepcionales conserva la
intersección, de modo que coincide con la acción sobre el grafo ya
construido \cite{Manivel2005}.
\end{proof}

\subsection{Estratificación de las 729 palabras}

Escribimos una palabra de seis trits como
\[
 w=(\ell_1,h_1,\ell_2,h_2,\ell_3,h_3),
\]
y definimos dos elementos de \(H_3\),
\[
 L(w)=(\ell_1,\ell_2,\ell_3),\qquad
 H(w)=(h_1,h_2,h_3),
\]
usando la ley central en sus productos. Según
\(L(w)^{-1}H(w)\), las 729 palabras se dividen en
\begin{align*}
 \mathcal D_\ell&=\{w:L(w)^{-1}H(w)=e\},\\
 \mathcal I_\ell&=\{w:L(w)^{-1}H(w)\in S_\ell\},\\
 \mathcal S_\ell&=\{w:L(w)^{-1}H(w)\notin\{e\}\cup S_\ell\}.
\end{align*}

\begin{theorem}[Partición Hensel--Schläfli]
\label{p12-83-c20-s6:thm:particion-hensel-schlafli-nueva}
\begin{equation}
 729=27+270+432,
 \label{p12-83-c20-s6:eq:particion-729-e6-nueva}
\end{equation}
correspondientes a diagonal, incidencia y alabeamiento.
\end{theorem}

\begin{proof}
Para cada uno de los 27 valores de \(L(w)\), hay un único valor de
\(H(w)\) con diferencia identidad, diez con diferencia en \(S_\ell\) y
dieciséis restantes. Se multiplica \(27\) por \(1,10,16\).
\end{proof}

El cardinal \(432\) del estrato alabeado coincide numéricamente con la
multiplicidad de la región transversal de \(\pi\) en
\cref{eq:descomposiciones-cinco-regiones-pi}. No se identifica por ello ambos objetos: aquí
se cuentan palabras clasificadas por \(L(w)^{-1}H(w)\), mientras allí se
cuentan preimágenes de un registro \(U_6\). Sin un entrelazador explícito y
equivariante entre esos dominios, \(432=432\) es sólo un control cardinal.

\subsection{Necesidad de la fase central}

La enumeración de todos los subconjuntos inverso-cerrados de cardinal diez
en los cinco grupos de orden 27 da
\[
\begin{array}{c|c|c}
\text{grupo}&\text{abeliano}&
\#\operatorname{srg}(27,10,1,5)\\ \hline
C_{27}&\text{sí}&0\\
C_9\times C_3&\text{sí}&0\\
C_3^3&\text{sí}&0\\
C_9\rtimes C_3&\text{no}&9\\
H_3&\text{no}&9.
\end{array}
\]
Así, conservar 27 etiquetas pero abelianizar la composición destruye la
configuración. No bastan la cardinalidad ni el grado.

Si \(\mathcal T\) es el conjunto de 45 triángulos, el cúbico
\[
 \mathcal C_\Gamma(x_1,\ldots,x_{27})
 =\sum_{\{i,j,k\}\in\mathcal T}x_ix_jx_k
\]
recupera la adyacencia, porque cada arista está en un triángulo único. Su
estabilizador permutacional es exactamente
\[
 \operatorname{Aut}(\Gamma_\ell)\cong W(E_6).
\]

\begin{figure}[htbp]
\centering
\[
 \text{curvaturas APP }\pm1
 \longrightarrow\sigma
 \longrightarrow\mathcal H_9
 \longrightarrow H_3
 \longrightarrow\Gamma_{27}
 \longrightarrow\mathcal T_{45}
 \longrightarrow W(E_6).
\]
\caption{Cadena causal de la rama de holonomía mixta. Es paralela a
\(C_W\to W_{12}\to M_{12}\) y no un eslabón de ella.}
\label{p12-83-c20-s6:fig:cadena-holonomia-e6-u029}
\end{figure}

\begin{criterion}[Criterios de refutación de la rama \(E_6\)]
\label{p12-83-c20-s6:crit:refutacion-rama-e6-u029}
La rama queda refutada si las curvaturas no son opuestas, si el cociclo
central se borra, si la corrección cuadrática
\eqref{p12-83-c20-s6:eq:correccion-cuadratica-automorfismos-h3} no conserva el producto,
si \eqref{p12-83-c20-s6:eq:identidad-grupo-schlafli-nueva} no produce los coeficientes
\((10,1,5)\), si un grupo abeliano produce la misma configuración, si la
partición no suma 729, si la tabla
\ref{p12-83-c20-s6:tab:h3-27-rectas-u029} no entrelaza las matrices de adyacencia o si
\(W(E_6)\) se identifica sólo por su orden y no por su acción sobre las
veintisiete clases y las 72 raíces.
\end{criterion}


\section{Acoplamiento orientado APP–Fock y rigidez transversal de \(12\) fibras}
\label{p12-83-c20-s7:chap:acoplamiento-app-fock}
\index{APP--Fock!acoplamiento orientado}
\index{Fock!grado dos}
\index{orientación de hoja}

El determinante ciclotómico ordinario no distingue \(C\) de \(C^{-1}\).
La información suma--producto se conserva introduciendo explícitamente la
paridad de hoja y una recta de orientación. El acoplamiento resultante
produce el índice \(54\) desde el agregado APP. Una vez fijadas las cuatro
condiciones del acoplamiento, \(54\), \(J\) y \(196884\) no son entradas de
la evaluación. Esta independencia evaluativa no afirma unicidad universal
del ansatz ni independencia histórica de su diseño.

\subsection{Marco triangular en cada órbita}

Sea
\[
M=\mathbb Z[\mathbb Z/9\mathbb Z]
\]
y sea \(M_{\mathrm{bal}}^{C_3}\) el submódulo de los elementos invariantes
\(\nu\) cuyos seis coeficientes unitarios coinciden:
\[
\nu_u=c(\nu)
\qquad
\bigl(u\in(\mathbb Z/9\mathbb Z)^\times\bigr).
\]
Los coeficientes de \(e_0,e_3,e_6\) permanecen libres.

Para una órbita libre
\(\mathcal O=\{u,4u,7u\}\), pongamos
\[
b_v=e_v-e_{4v}\in A(\mathcal O).
\]
Si \(b_v^\flat\) es el covector definido por la forma de \(A_2\), las tres
raíces del marco triangular satisfacen
\[
\boxed{
\sum_{v\in\mathcal O}b_v\otimes b_v^\flat
=3I_{A(\mathcal O)}.}
\]
En consecuencia,
\[
\sum_{v\in\mathcal O}\nu_v\,b_v\otimes b_v^\flat
=3c(\nu)I.
\]
El factor tres procede de la resolución de la identidad por el marco de
raíces.

\subsection{Acoplamiento orientado \(\mathcal H_{\Sigma,\Pi}^{(m)}\) entre el balance \(C_3\), la paridad de hoja y el grado \(2\) del espacio de Fock}

Para \(m\in2\mathbb N_{>0}\), sea
\[
\begin{aligned}
\mathcal V_m&=\mathcal V_{\Sigma,m}\oplus\mathcal V_{\Pi,m},\\
\mathcal V_{\Sigma,m}&=(A_2\otimes\mathbb C)^{\oplus m/2},\qquad
\mathcal V_{\Pi,m}=(A_2\otimes\mathbb C)^{\oplus m/2},\\
g_m&=C^{\oplus m/2}\oplus(C^{-1})^{\oplus m/2},\\
\mathsf E_m&=I_{\mathcal V_{\Sigma,m}}
\oplus(-I_{\mathcal V_{\Pi,m}}).
\end{aligned}
\]
La especialización \(m=12\), con las etiquetas de par, hoja y órbita,
es \(\mathcal V_{12}=W_{\mathrm{APP}}\otimes\mathbb C\). Esta notación
evita confundir el espacio representacional con el sistema de Witt
\(W_{12}=S(5,6,12)\). El grado dos del espacio de
Fock y la acción inducida son
\[
\mathcal F_2(\mathcal V_m)
=\mathcal V_m[2]\oplus\operatorname{Sym}^2(\mathcal V_m[1]),
\]
\[
\Gamma_2(g_m)
=g_m|_{\mathcal V_m[2]}
\oplus\operatorname{Sym}^2(g_m|_{\mathcal V_m[1]}).
\]

Sea \(o_{\Sigma\Pi}\) la recta de orientación con signo
\(\Pi-\Sigma\), y sea \(M_{\mathrm{bal}}^{C_3,-}\) la copia orientada en
la que el intercambio de hojas actúa por \(\nu\mapsto-\nu\). Definimos
\[
\boxed{
\begin{aligned}
\mathcal H_{\Sigma,\Pi}^{(m)}&:
M_{\mathrm{bal}}^{C_3,-}\longrightarrow
\operatorname{End}(\mathcal F_2(\mathcal V_m))\otimes o_{\Sigma\Pi},\\
\mathcal H_{\Sigma,\Pi}^{(m)}(\nu)
&=3c(\nu)
\left(0_{\mathcal V_m[2]}\oplus\operatorname{Sym}^2\mathsf E_m\right)
\otimes o_{\Sigma\Pi}.
\end{aligned}}
\]
La transformación es lineal en \(c\), está soportada en
\(\operatorname{Sym}^2(\mathcal V_m[1])\), utiliza la paridad de hoja y adopta la
normalización del marco triangular.

Sea \(\sigma\) el intercambio de hojas. Entonces
\[
\nu\longmapsto-\nu,\qquad
o_{\Sigma\Pi}\longmapsto-o_{\Sigma\Pi},\qquad
\mathsf E_m\longmapsto-\mathsf E_m,
\]
y
\[
\operatorname{Sym}^2(-\mathsf E_m)
=\operatorname{Sym}^2(\mathsf E_m).
\]
Por tanto,
\[
\mathcal H_{\Sigma,\Pi}^{(m)}(\sigma\nu)
=\sigma\mathcal H_{\Sigma,\Pi}^{(m)}(\nu).
\]
Al invertir la orientación, el covector dual cambia también de signo; así
\((-o_{\Sigma\Pi}^{\vee})(-o_{\Sigma\Pi})=1\), y la evaluación orientada
permanece normalizada.

Fijado
\(o_{\Sigma\Pi}^{\vee}(o_{\Sigma\Pi})=1\), la traza orientada es
\[
\operatorname{Tr}_{o,\mathcal F_2}
(T\otimes o_{\Sigma\Pi})
=\operatorname{Tr}_{\mathcal F_2}(T).
\]

\begin{theorem}[Índice orientado APP--Fock]
\label{p12-83-c20-s7:thm:indice-orientado-app-fock}
Supóngase que \(m\) fibras \(A_2\), con \(m\) par, se reparten por igual
entre las orientaciones \(C\) y \(C^{-1}\). Entonces
\[
\boxed{
\operatorname{Tr}_{o,\mathcal F_2}
\left(
\mathcal H_{\Sigma,\Pi}^{(m)}(\nu)\Gamma_2(g_m)
\right)
=-\frac{3m\,c(\nu)}2.}
\]
Para el agregado producido por APP,
\[
\delta=\nu_\Pi-\nu_\Sigma
=12e_0+3e_3+3e_6
-3\sum_{u\in(\mathbb Z/9\mathbb Z)^\times}e_u,
\]
se tiene \(c(\delta)=-3\) y
\[
m=3\ \text{pares}\cdot2\ \text{hojas}\cdot2\ \text{órbitas}=12.
\]
Por tanto,
\[
\boxed{
\operatorname{Tr}_{o,\mathcal F_2}
\left(
\mathcal H_{\Sigma,\Pi}^{(12)}(\delta)\Gamma_2(g_{12})
\right)=54.}
\]
\end{theorem}

\begin{proof}
Sea \(A=\mathsf E_mg_m\). Las contribuciones de las dos hojas se cancelan:
\(\operatorname{tr}A=0\). Como \(\mathsf E_m\) conmuta con \(g_m\) y
\(\mathsf E_m^2=I\),
\[
\operatorname{tr}(A^2)=\operatorname{tr}(g_m^2)=-m.
\]
La identidad
\[
\operatorname{tr}(\operatorname{Sym}^2A)
=\frac{(\operatorname{tr}A)^2+\operatorname{tr}(A^2)}2
\]
da \(-m/2\); el marco aporta \(3c(\nu)\).
\end{proof}

\begin{corollary}[Minimalidad del grado bosónico detector]
\label{p12-83-c20-s7:cor:minimalidad-grado-dos-app-fock-u030}
En la clase de acoplamientos declarada, el grado uno no detecta el agregado
orientado y el grado dos es el primer grado bosónico positivo que lo
detecta:
\[
 \operatorname{tr}(\mathsf E_mg_m)=0,
 \qquad
 \operatorname{tr}\!\left(
 \operatorname{Sym}^2(\mathsf E_mg_m)
 \right)=-\frac m2\neq0.
\]
\end{corollary}

\begin{proof}
La primera igualdad es la cancelación de las dos hojas. La segunda se
obtuvo en la prueba del \cref{p12-83-c20-s7:thm:indice-orientado-app-fock}. Como el grado
cero sólo contiene la unidad y el grado uno da traza nula, el grado dos es
el primero con respuesta no nula.
\end{proof}

\subsection{Rigidez transversal del número de fibras}

Para un entero \(m\ge1\), definimos la serie formal
\begin{equation}
 T_m(q):=
 q^{-1}\prod_{n\ge1}(1+q^n+q^{2n})^{-m}
 \in q^{-1}\mathbb Z[[q]]
 \label{p12-83-c20-s7:eq:familia-fock-m-fibras-u030}
\end{equation}
y, en la suma ortogonal \(A_2^m\), el vector radial
\begin{equation}
 v_m:=4\rho_1+\rho_2+\cdots+\rho_m,
 \qquad
 \|\rho_j\|^2=2.
 \label{p12-83-c20-s7:eq:familia-radial-m-fibras-u030}
\end{equation}
Construimos dos funciones enteras sin usar previamente \(m=12\):
\[
 F(m):=[q]T_m(q),
 \qquad
 R(m):=\|v_m\|^2.
\]

\begin{proposition}[Coeficiente de Fock y norma radial]
\label{p12-83-c20-s7:prop:dos-funciones-rigidez-doce-u030}
Para todo \(m\ge1\),
\begin{equation}
 F(m)=\frac{m(m-3)}2,
 \qquad
 R(m)=2(m+15).
 \label{p12-83-c20-s7:eq:funciones-rigidez-doce-u030}
\end{equation}
\end{proposition}

\begin{proof}
Hasta orden dos,
\[
 (1+q+q^2)^{-m}
 =1-mq+\frac{m(m-1)}2q^2+O(q^3),
\]
\[
 (1+q^2+q^4)^{-m}=1-mq^2+O(q^4).
\]
Los factores con \(n\ge3\) no contribuyen al coeficiente de \(q^2\).
Como el factor exterior de \(T_m\) es \(q^{-1}\),
\[
 F(m)=\frac{m(m-1)}2-m=\frac{m(m-3)}2.
\]
Por ortogonalidad de las \(m\) copias de \(A_2\),
\[
 R(m)=4^2\|\rho_1\|^2+
 \sum_{j=2}^{m}\|\rho_j\|^2
 =32+2(m-1)=2(m+15).
\]
\end{proof}

\begin{theorem}[Rigidez ciclotómica--radial]
\label{p12-83-c20-s7:thm:rigidez-transversal-doce-u030}
La igualdad \(F(m)=R(m)\) equivale a
\begin{equation}
 (m-12)(m+5)=0.
 \label{p12-83-c20-s7:eq:ecuacion-rigidez-doce-u030}
\end{equation}
En el dominio \(m\in\mathbb Z_{>0}\), la única solución es
\[
 \boxed{m=12,\qquad F(12)=R(12)=54.}
\]
\end{theorem}

\begin{proof}
Igualar las dos expresiones de
\eqref{p12-83-c20-s7:eq:funciones-rigidez-doce-u030} produce
\[
 m(m-3)=4(m+15),
\]
es decir,
\[
 m^2-7m-60=(m-12)(m+5)=0.
\]
La segunda raíz es negativa y queda fuera del dominio.
\end{proof}

Esta rigidez pertenece a la familia ciclotómica--radial declarada. No usa el
valor \(54\) para escoger \(m\): construye primero \(F(m)\) y \(R(m)\), y
resuelve después su igualdad.

\subsection{Factorización del acoplamiento APP–Fock: hipótesis, invariantes y condiciones necesarias de validez y obstrucciones}

El acoplamiento factoriza por
\[
M_{\mathrm{bal}}^{C_3,-}/\ker c
\cong\mathbb Z\otimes o_{\Sigma\Pi},
\]
y el mapa inducido es inyectivo porque
\(\operatorname{Sym}^2\mathsf E_m\neq0\).
La fórmula queda seleccionada por las cuatro condiciones declaradas:
linealidad en \(c\), soporte en
\(\operatorname{Sym}^2(\mathcal V_m[1])\), paridad
de hoja y normalización por el marco triangular. No se afirma unicidad
entre todos los acoplamientos naturales ni independencia histórica del
ansatz. Una vez fijadas esas condiciones, \(54\), \(J\) y \(196884\) no
son entradas de la evaluación.

Tres controles muestran que el resultado no está fijado por una
normalización retrospectiva:
\[
(m,c)=(12,-2)\longmapsto36,\qquad
(m,c)=(10,-3)\longmapsto45,
\]
y
\[
\delta_k=\delta+ke_0
\]
deja invariante \(\mathcal H_{\Sigma,\Pi}^{(12)}\) mientras el peso digital
cambia de \(54\) a \(54+9k\). En particular, para \(k=1\),
\[
\operatorname{Tr}_{o,\mathcal F_2}
\left(
\mathcal H_{\Sigma,\Pi}^{(12)}(\delta+e_0)\Gamma_2(g_{12})
\right)=54,
\qquad
w(\delta+e_0)=63.
\]

\begin{proposition}[Necesidad del dato de orientación]
\label{p12-83-c20-s7:prop:orientacion-app-fock}
Si se olvida \(o_{\Sigma\Pi}\), todo escalar natural lineal en
\(\nu_\Pi-\nu_\Sigma\) y dependiente sólo de la clase de la representación
se anula.
\end{proposition}

\begin{proof}
\(C\) y \(C^{-1}\) son conjugados, por lo que la clase representacional es
par bajo el intercambio de hojas. El agregado
\(\nu_\Pi-\nu_\Sigma\) es impar. Sin la recta de orientación, una
transformación simultáneamente par e impar vale cero.
\end{proof}

El teorema establece así el transporte exacto del agregado al grado dos de
Fock. El defecto local estado a estado conserva la obstrucción de descenso
del \cref{p12-83-c1-s5:cor:obstruccion-isotipica-app-autonoma}; ambos enunciados pertenecen a
dominios distintos y son compatibles.


% Unidad U031. Alineamiento APP--Witt, automorfia de orden tres,
% construccion VOA y Moonshine.
 \section{Del salto espectral 4 a las simetrías del Monstruo: despliegue APP
\texorpdfstring{$2\to6\to54$}{2--6--54}, retículo de Leech,
\texorpdfstring{$V^\natural$}{V natural} y corredor Moonshine}
\label{sec:cierre-equivalencias-corredor-moonshine-20260825}

El corredor no es una sucesión de coincidencias numéricas. Parte del bloque
central de APP, conserva la orientación suma--producto, se ramifica en los
acoplamientos APP--Fock y APP--Witt, y vuelve a coordinarse mediante la
gradación. El orden causal vinculante es
\[\begin{aligned}
 B_c=9P_++5P_-
 &\longrightarrow 4\longrightarrow2\longrightarrow6\longrightarrow54\\
 &\longrightarrow
 \left\{
 \begin{array}{l}
  \text{grado dos de Fock y }[q]t_3=54,\\
  \text{alineamiento APP--Witt y }N_\alpha\cong\Lambda_{24}
 \end{array}
 \right.\\
 &\longrightarrow V^\natural
 \longrightarrow
 \bigl(\operatorname{Aut},\operatorname{ch},g\mapsto T_g\bigr)
 \longrightarrow\text{Moonshine}.
\end{aligned}\]
Las dos ramas intermedias poseen codominios distintos. La llave común \(54\)
es un invariante graduado transportado por el corredor; no es una aplicación
lineal desde un escalar hacia un álgebra de operadores de vértice.

\subsection{Jerarquía tipada y \(2\) realizaciones del índice \(54\)}

\begin{theorem}[Jerarquía espectral y despliegue APP]
\label{thm:cierre-vtm-jerarquia-app}
Sea \(B_c=7I+2R=9P_++5P_-\) el bloque central de APP.\@ Si
\(M_\Pi,M_\Sigma\) son las compresiones modulares definidas, entonces
\[
 9-5=4,
 \qquad (B_c)_{i\ne j}=2,
 \qquad \sum M_\Pi-\sum M_\Sigma=51-45=6,
\]
y el despliegue sobre las nueve posiciones da
\[
 9\cdot6=459-405=54.
\]
Los cuatro enteros pertenecen, respectivamente, al salto espectral, al
acoplamiento local, a la compresión modular y a la integración bidimensional.
\end{theorem}

\begin{proof}
La descomposición espectral de \(B_c\) da \(9\) y \(5\), mientras que
\(7I+2R\) fija el coeficiente no diagonal. Las dos sumas agregadas son \(51\)
y \(45\). Cada entrada comprimida representa nueve posiciones, de modo que el
despliegue multiplica la diferencia por nueve. Ninguno de los cuatro enteros
se obtiene identificando retrospectivamente sus tipos.
\end{proof}

Fijadas la recta de orientación \(o_{\Sigma\Pi}\), la paridad de hoja y la
normalización triangular, el acoplamiento orientado
\[
 \mathcal H_{\Sigma,\Pi}^{(m)}:
 M_{\mathrm{bal}}^{C_3,-}
 \longrightarrow
 \operatorname{End}\!\bigl(\mathcal F_2(\mathcal V_m)\bigr)
 \otimes o_{\Sigma\Pi}
\]
satisface
\[
 \operatorname{Tr}_{o,\mathcal F_2}
 \left(\mathcal H_{\Sigma,\Pi}^{(m)}(\nu)\Gamma_2(g_m)\right)
 =-\frac{3m\,c(\nu)}2.
\]
Para \(m=12\) y \(c(\delta)=-3\), la traza vale \(54\). Por otro lado,
\[
 t_3(q)=q^{-1}\prod_{n\ge1}(1+q^n+q^{2n})^{-12}
       =q^{-1}-12+54q-76q^2-243q^3+\cdots,
\]
de donde \([q]t_3=54\).

\begin{theorem}[Cono de compatibilidad del índice común]
\label{thm:cierre-vtm-indice-comun-54}
Con las estructuras de partida de cada rama conservadas, se cumple
\[
 \boxed{
 D_{\mathrm{APP}}
 =\operatorname{Tr}_{o,\mathcal F_2}
   \left(\mathcal H_{\Sigma,\Pi}^{(12)}(\delta)\Gamma_2(g_{12})\right)
 =[q]t_3
 =v_\alpha^2
 =\frac{196884}{5\cdot729+1}
 =54.}
\]
Esta igualdad constituye un cono de evaluaciones hacia el mismo entero. No
identifica \(W_{\mathrm{APP}}\), el espacio de Fock, el retículo vecino y
\(V^\natural\).
\end{theorem}

\begin{proof}
El censo APP da \(459-405=54\); la traza orientada da \(54\); la expansión
del producto ciclotómico da \([q]t_3=54\); y el vector radial del vecino
satisface \(v_\alpha^2=4^2\cdot2+11\cdot2=54\). Finalmente,
\(5\cdot729+1=3646\) y \(54\cdot3646=196884\). Las igualdades se comparan
después de construir cada objeto y por eso no son una selección numérica del
corredor.
\end{proof}

\paragraph{Estatuto exacto de la jerarquía \(4\to2\to6\to54\) y de la identificación reticular.}
La jerarquía \(4\to2\to6\to54\), la traza APP--Fock y la extracción
\([q]t_3\) son identidades internas exactas. La identificación del vecino con
Leech usa la cámara, el vector de vecindad y la clasificación reticular
declaradas.
%
 
\chapter{Álgebra de operadores de vértice y Moonshine monstruoso}
\section{Alineamiento ciclotómico, orbifold de orden \(3\) y Moonshine monstruoso}
\label{p12-83-c20-s8:chap:alineamiento-orbifold-moonshine}

El retículo de Leech construido en la
\cref{p12-83-c20-s5:thm:identificacion-vecino-leech} permite pasar de la incidencia
finita a una estructura graduada infinita. Esta transición no se efectuará
mediante una cadena de nombres. Se construirán, en orden tipado:
\begin{enumerate}
 \item el entrelazamiento de las doce fibras \(A_2\);
 \item una isometría reticular fijo-libre de orden tres;
 \item su traza graduada y su elevación de tipo cero;
 \item el orbifold cíclico y su automorfismo dual de clase \(3B\);
 \item el álgebra de operadores de vértice lunar;
 \item por publicaciones distintas, su grupo de automorfismos, su carácter
       graduado y las series de McKay--Thompson.
\end{enumerate}
Moonshine será el teorema que coordina la acción y las series. No se usará
el atajo no tipado
\(V^\natural\to\mathbb M\to\text{Moonshine}\).

\subsection{Alineamiento de las \(12\) fibras APP y Witt}

Las doce fibras de \(W_{\mathrm{APP}}\cong A_2^{12}\) llevan tres
etiquetas independientes:
\[
 i\in\mathbb Z/3\mathbb Z
 \quad\text{(par de coordenadas)},\qquad
 \mu\in\{0,1\}
 \quad\text{(hoja)},\qquad
 \varepsilon\in\{0,1\}
 \quad\text{(órbita ciclotómica)}.
\]
Fijado el origen dodecafásico construido en el TPK, definimos
\begin{equation}
 \lambda(i,\mu,\varepsilon)
 :=4i+2\mu+\varepsilon\pmod{12}.
 \label{p12-83-c20-s8:eq:alineamiento-doce-fibras-u031}
\end{equation}

\begin{lemma}[Biyectividad del alineamiento]
\label{p12-83-c20-s8:lem:biyectividad-alineamiento-u031}
La aplicación \(\lambda\) de
\eqref{p12-83-c20-s8:eq:alineamiento-doce-fibras-u031} es una biyección sobre
\(\mathbb Z/12\mathbb Z\).
\end{lemma}

\begin{proof}
Para cada \(i\), los cuatro valores \(4i+2\mu+\varepsilon\) forman el bloque
\(\{4i,4i+1,4i+2,4i+3\}\). Los tres bloques para \(i=0,1,2\) son disjuntos
y cubren las doce clases.
\end{proof}

Sea \(C\) el Coxeter de \(A_2\) y sea \(\mathsf R\) la reflexión que
satisface
\begin{equation}
 \mathsf R C\mathsf R=C^{-1}.
 \label{p12-83-c20-s8:eq:reflexion-conjuga-coxeter-u031}
\end{equation}
Fijadas las trivializaciones \(\iota_b\) y los marcos orientados \(P_b\),
ponemos
\begin{align}
 g_b&:=\iota_b^{-1}P_bCP_b^{-1}\iota_b,
 \label{p12-83-c20-s8:eq:coxeter-calibrado-fibra-u031}\\
 T_{i\mu\varepsilon}
 &:=\iota_{\lambda(i,\mu,\varepsilon)}^{-1}
 P_{\lambda(i,\mu,\varepsilon)}\mathsf R^\mu.
 \label{p12-83-c20-s8:eq:entrelazador-fibra-u031}
\end{align}

\begin{theorem}[Entrelazamiento APP--Witt]
\label{p12-83-c20-s8:thm:alineamiento-app-witt-u031}
Para cada terna \((i,\mu,\varepsilon)\),
\begin{equation}
 g_{\lambda(i,\mu,\varepsilon)}T_{i\mu\varepsilon}
 =T_{i\mu\varepsilon}C^{(-1)^\mu}.
 \label{p12-83-c20-s8:eq:identidad-entrelazamiento-app-witt-u031}
\end{equation}
La suma directa define un isomorfismo \(C_3\)-equivariante
\begin{equation}
 T_\lambda:
 W_{\mathrm{APP}}\otimes\mathbb C
 \xrightarrow{\;\cong\;}
 W_c:=\bigoplus_{b=0}^{11}
 (A_{2,b}\otimes\mathbb C).
 \label{p12-83-c20-s8:eq:isomorfismo-app-witt-u031}
\end{equation}
\end{theorem}

\begin{proof}
Para \(\mu=0\), la definición de \(g_b\) da directamente la conjugación
por \(C\). Para \(\mu=1\), se inserta
\(\mathsf RC\mathsf R=C^{-1}\). El
\cref{p12-83-c20-s8:lem:biyectividad-alineamiento-u031} garantiza que la suma directa
usa exactamente una vez cada posición Witt y conserva par, hoja y órbita.
\end{proof}

El teorema es relativo al transporte simultáneo del orden de pares, hojas y
órbitas, del origen Witt, de los marcos y de la orientación. Cambiar una
sola de estas coordenadas no define el mismo funtor.

\subsection{Traza graduada del generador ternario}

Sea
\[
 \rho_W:=\bigoplus_{b=0}^{11}g_b
\]
y consideremos la envolvente bosónica
\begin{equation}
 \mathcal H_{\rm bos}(W_c)
 :=
 \operatorname{Sym}\!\left(
 \bigoplus_{n\ge1}W_c[n]
 \right).
 \label{p12-83-c20-s8:eq:fock-bosonico-doce-fibras-u031}
\end{equation}
Si \(N\) es el operador de grado, definimos
\begin{equation}
 Z_{\rm Fock}(q)
 :=
 q^{-1}\operatorname{Tr}_{\mathcal H_{\rm bos}(W_c)}
 \left(\Gamma_{\rm bos}(\rho_W)q^N\right).
 \label{p12-83-c20-s8:eq:traza-graduada-fock-u031}
\end{equation}

\begin{proposition}[Producto ciclotómico de nivel tres]
\label{p12-83-c20-s8:prop:producto-tres-doce-fibras-u031}
Como serie formal de Laurent,
\begin{equation}
 \boxed{
 Z_{\rm Fock}(q)
 =q^{-1}\prod_{n\ge1}(1+q^n+q^{2n})^{-12}
 =:t_3(q).}
 \label{p12-83-c20-s8:eq:t3-producto-u031}
\end{equation}
Su expansión comienza por
\begin{equation}
 t_3(q)
 =q^{-1}-12+54q-76q^2-243q^3+1188q^4+O(q^5).
 \label{p12-83-c20-s8:eq:t3-expansion-u031}
\end{equation}
\end{proposition}

\begin{proof}
Cada \(g_b\) es conjugado a \(C\), por lo que
\[
 \det(I-g_bq^n)=1+q^n+q^{2n}.
\]
La identidad
\[
 \sum_{k\ge0}
 \operatorname{Tr}(\operatorname{Sym}^k\rho_W)t^k
 =\det(I-\rho_Wt)^{-1}
\]
aplicada en cada grado produce el producto. La multiplicación formal hasta
grado cinco da la expansión. En particular, el coeficiente de \(q\) se
obtiene como
\[
 [q^2]\,
 (1+q+q^2)^{-12}(1+q^2+q^4)^{-12}
 =66-12=54.
\]
\end{proof}

El valor \(54\) aparece aquí como coeficiente de una traza construida desde
doce fibras. La unidad U030 lo obtuvo también como índice orientado del
agregado APP y probó que doce es la única solución positiva de la rigidez
ciclotómica--radial. La igualdad de ambos valores es una compatibilidad de
dos realizaciones; ninguna usa el coeficiente de \(J\) como entrada.

\subsection{Palabra de soporte total e isometría fijo-libre}

Sea
\[
 q_*=(1,1,1,1,1,1)\in\mathbb F_3^6.
\]
La multiplicación por la matriz de Paley--Witt da
\begin{equation}
 q_*A_W=(2,1,1,1,1,1).
 \label{p12-83-c20-s8:eq:producto-qstar-aw-u031}
\end{equation}
Por tanto,
\begin{equation}
 c_*:=(q_*,q_*A_W)
 =(1,1,1,1,1,1,2,1,1,1,1,1)
 \in C_W
 \label{p12-83-c20-s8:eq:palabra-soporte-total-u031}
\end{equation}
y todas sus coordenadas son no nulas.

Para cada posición \(b\), fijamos
\begin{equation}
 P_b(c_*)=
 \begin{cases}
  \mathsf R,&(c_*)_b=1,\\
  I,&(c_*)_b=2,
 \end{cases}
 \qquad
 g_c:=\bigoplus_{b=0}^{11}P_b(c_*)CP_b(c_*)^{-1}.
 \label{p12-83-c20-s8:eq:isometria-orden-tres-u031}
\end{equation}

\begin{proposition}[Orden, ausencia de puntos fijos y pegado]
\label{p12-83-c20-s8:prop:isometria-fijo-libre-pegado-u031}
El operador \(g_c\) tiene orden tres, carece de vectores fijos no nulos y
preserva el retículo de Niemeier \(N(A_2^{12})\) construido en U021.
\end{proposition}

\begin{proof}
Cada bloque es conjugado al Coxeter \(C\), cuyo polinomio mínimo es
\(X^2+X+1\). Luego tiene orden tres y no posee autovalor uno. La suma
directa conserva ambas propiedades. El Coxeter actúa trivialmente sobre
\(A_2^*/A_2\); por ello cada bloque conserva la clase discriminante y la
unión de clases laterales que define el pegado.
\end{proof}

\subsection{Preservación del vecino de Leech}

Escribamos
\[
 v:=v_{\mathcal F_*}
 =4\rho_1+\rho_2+\cdots+\rho_{12},
 \qquad
 v^2=54,
\]
para el vector seleccionado por el origen incidencial de
\eqref{p12-83-c20-s5:eq:vector-marco-incidencial}. Sean
\[
 N:=N(A_2^{12}),
 \qquad
 N_0:=\ker(\langle\,\cdot\,,v\rangle\bmod3),
 \qquad
 N_v:=N_0+\mathbb Z\frac v3\cong\Lambda_{24}.
\]

\begin{theorem}[Preservación orientada del vecino]
\label{p12-83-c20-s8:thm:preservacion-vecino-orden-tres-u031}
Los desplazamientos
\begin{equation}
 y_*:=\frac{g_cv-v}{3},
 \qquad
 z_*:=\frac{g_c^{-1}v-v}{3}
 \label{p12-83-c20-s8:eq:desplazamientos-vecino-u031}
\end{equation}
pertenecen a \(N_0\). En consecuencia,
\[
 g_cN_0=N_0,
 \qquad
 g_cN_v=N_v.
\]
\end{theorem}

\begin{proof}
Las palabras discriminantes de \(y_*\) y \(z_*\) son \(c_*\) y
\(-c_*\), ambas en \(C_W\); de ahí \(y_*,z_*\in N\). El cálculo en cada
factor \(A_2\) da
\begin{equation}
 \langle y_*,v\rangle
 =\langle z_*,v\rangle
 =-(4^2+11)=-27\equiv0\pmod3.
 \label{p12-83-c20-s8:eq:emparejamientos-desplazamientos-u031}
\end{equation}
Luego ambos están en \(N_0\). Para \(x\in N_0\),
\[
 \langle g_cx,v\rangle
 =\langle x,g_c^{-1}v\rangle
 =\langle x,v\rangle+3\langle x,z_*\rangle
 \equiv0\pmod3.
\]
La misma cuenta con \(g_c^{-1}\) da igualdad del núcleo. Finalmente,
\[
 g_c(v/3)=v/3+y_*,
\]
por lo que se preserva la clase que genera el vecino.
\end{proof}

\subsection{\(24\) orientaciones y naturalidad reticular}

El enumerador de pesos de \eqref{p12-83-c20-s1:eq:enumerador-pesos-cw} contiene
\(24y^{12}\). Por tanto, el conjunto
\[
 C_W^\times:=
 \{u\in C_W:u_b\in\{1,2\}\text{ para todo }b\}
\]
posee 24 elementos. Para cada \(u\in C_W^\times\), definimos
\[
 P_b(u)=
 \begin{cases}
  \mathsf R,&u_b=1,\\
  I,&u_b=2,
 \end{cases}
 \qquad
 g_u:=\bigoplus_{b=0}^{11}P_b(u)CP_b(u)^{-1}.
\]

\begin{proposition}[Robustez de las orientaciones de soporte total]
\label{p12-83-c20-s8:prop:robustez-orientaciones-u031}
Para toda \(u\in C_W^\times\),
\[
 \frac{g_uv-v}{3},
 \quad
 \frac{g_u^{-1}v-v}{3}
 \in N_0.
\]
En particular, cada \(g_u\) preserva el vecino construido con la
orientación transportada.
\end{proposition}

\begin{proof}
En coordenadas de raíces simples,
\[
 C\rho-\rho=(-2,-1),
 \qquad
 C^{-1}\rho-\rho=(-1,-2).
\]
Las clases discriminantes de los dos desplazamientos son \(u\) y \(-u\).
Para el coeficiente radial \(a_b\in\{4,1\}\),
\[
 \left\langle
 \frac{(C^{\pm1}-I)a_b\rho}{3},a_b\rho
 \right\rangle=-a_b^2.
\]
La suma sobre los doce factores vuelve a dar \(-27\). La prueba del
\cref{p12-83-c20-s8:thm:preservacion-vecino-orden-tres-u031} se aplica componente a
componente.
\end{proof}

Para tipar la recalibración, sea
\(h\in\operatorname{Aut}(C_W)\) monomial: una permutación \(\sigma\) de
las doce coordenadas seguida de multiplicadores
\(\epsilon_b\in\{1,2\}\). Definimos su elevación reticular
\begin{equation}
 \widetilde h
 :=P_\sigma\circ
 \bigoplus_{b=0}^{11}\mathsf R^{\nu_b},
 \qquad
 \nu_b=
 \begin{cases}
 0,&\epsilon_b=1,\\
 1,&\epsilon_b=2.
 \end{cases}
 \label{p12-83-c20-s8:eq:elevacion-monomial-u031}
\end{equation}
Su acción en el discriminante es exactamente \(h\). Si el cambio transporta
simultáneamente código, bandera, origen incidencial, vector \(v\), vecino y
marcos, entonces
\begin{equation}
 g_{hu}=\widetilde h\,g_u\,\widetilde h^{-1}.
 \label{p12-83-c20-s8:eq:naturalidad-isometrias-u031}
\end{equation}
La matriz ternaria \(h\) no actúa sin elevación sobre \(A_2^{12}\) ni sobre
Leech.

\subsection{Traza reticular y cálculo del tipo \(0\)}

Sea \(\widehat g_c\) la elevación estándar de \(g_c\) al VOA reticular
\[
 V_{N_v}=M(1)\otimes\mathbb C_\varepsilon[N_v].
\]
Para \(j=1,2\), la descomposición oscilador--retículo produce
\begin{equation}
 \begin{split}
 Z_{0,j}(\tau)
 &:=
 \operatorname{Tr}_{V_{N_v}}
 \left(\widehat g_c^{\,j}q^{L_0-1}\right)\\
 &=q^{-1}
 \left(
 \sum_{\lambda\in N_v^{g_c^j}}
 \varepsilon_j(\lambda)
 q^{\langle\lambda,\lambda\rangle/2}
 \right)
 \prod_{n\ge1}\det(I-g_c^jq^n)^{-1}.
 \end{split}
 \label{p12-83-c20-s8:eq:factorizacion-traza-reticular-u031}
\end{equation}

\begin{proposition}[Traza de los sectores no triviales]
\label{p12-83-c20-s8:prop:traza-reticular-t3-u031}
Se cumple
\[
 Z_{0,1}(\tau)=Z_{0,2}(\tau)=t_3(\tau).
\]
\end{proposition}

\begin{proof}
La ausencia de puntos fijos implica
\[
 N_v^{g_c}=N_v^{g_c^2}=\{0\},
\]
de modo que la suma reticular de
\eqref{p12-83-c20-s8:eq:factorizacion-traza-reticular-u031} vale uno. En cada factor,
\(g_c^j\) tiene autovalores \(\zeta_3,\zeta_3^2\); por tanto,
\[
 \det(I-g_c^jq^n)=1+q^n+q^{2n}.
\]
Los doce factores reproducen \eqref{p12-83-c20-s8:eq:t3-producto-u031}.
\end{proof}

\begin{theorem}[Peso torcido y tipo de la elevación]
\label{p12-83-c20-s8:thm:peso-torcido-tipo-cero-u031}
La elevación estándar tiene orden tres, peso torcido
\begin{equation}
 \rho_{g_c}=\frac43
 \label{p12-83-c20-s8:eq:peso-torcido-u031}
\end{equation}
y tipo cero.
\end{theorem}

\begin{proof}
En
\(\mathfrak h=N_v\otimes_{\mathbb Z}\mathbb C\), cada uno de los doce
factores aporta un autovector de autovalor \(\zeta_3\) y otro de autovalor
\(\zeta_3^2\). Por tanto,
\[
 \dim\mathfrak h_{\zeta_3}
 =\dim\mathfrak h_{\zeta_3^2}=12.
\]
La fórmula del peso del vacío torcido de orden tres da
\[
 \rho_{g_c}
 =\frac1{4\cdot3^2}
 \bigl(1\cdot2\cdot12+2\cdot1\cdot12\bigr)
 =\frac43.
\]
Para una isometría reticular de orden impar, la elevación estándar conserva
el orden. El tipo es la clase de
\[
 3^2\rho_{g_c}=12\equiv0\pmod3.
\]
Así, las dos hipótesis inicialmente independientes quedan calculadas aquí.
\end{proof}

\subsection{Orbifold cíclico y clase dual \texorpdfstring{\(3B\)}{3B}}

\begin{theorem}[Orbifold triádico del vecino de Leech]
\label{p12-83-c20-s8:thm:orbifold-triadico-vnatural-u031}
Los teoremas clásicos de orbifold cíclico, aplicados a la elevación estándar
del \cref{p12-83-c20-s8:thm:peso-torcido-tipo-cero-u031}, dan
\begin{equation}
 \operatorname{Orb}_{\mathbb Z/3}
 (V_{N_v},\widehat g_c)
 \cong V^\natural.
 \label{p12-83-c20-s8:eq:orbifold-triadico-vnatural-u031}
\end{equation}
El automorfismo dual pertenece a la clase \(3B\) y su serie es
\begin{equation}
 T_{3B}(\tau)=t_3(\tau)+12.
 \label{p12-83-c20-s8:eq:serie-3b-u031}
\end{equation}
\end{theorem}

\begin{proof}
La isometría es fijo-libre de orden tres; el peso torcido y el tipo cero se
calcularon, no se supusieron. Los teoremas de orbifold cíclico identifican
la extensión holomorfa con el módulo lunar y clasifican el automorfismo dual
como \(3B\)
\cite{ChenLamShimakura2016,AbeLamYamada2017,Carnahan2017}. La
identificación usa estos teoremas después de construir el objeto reticular;
no se deduce sólo de una igualdad de series.
\end{proof}

Si \(V^{(i)}\) es el módulo \(\widehat g_c^{\,i}\)-torcido y
\[
 Z_{i,j}(\tau):=
 \operatorname{Tr}_{V^{(i)}}\!
 \left(
 \phi_i(\widehat g_c)^j q^{L_0-1}
 \right),
\]
el carácter del orbifold es el promedio de los nueve sectores:
\begin{equation}
 \operatorname{ch}_{V^\natural}(\tau)
 =\frac13\sum_{i,j\in\mathbb Z/3}Z_{i,j}(\tau)
 =J(\tau).
 \label{p12-83-c20-s8:eq:promedio-nueve-sectores-u031}
\end{equation}
La identidad racional de nivel tres es
\begin{equation}
 J
 =t_3+12+\frac{10\,3^9}{t_3}
 +\frac{4\,3^{14}}{t_3^2}
 +\frac{3^{18}}{t_3^3}.
 \label{p12-83-c20-s8:eq:j-desde-t3-u031}
\end{equation}
Como \(t_3^{-1}=q+O(q^2)\),
\begin{equation}
 [q]J=54+10\cdot3^9=196884
 =54(5\cdot729+1).
 \label{p12-83-c20-s8:eq:coeficiente-196884-u031}
\end{equation}
Esta es una compatibilidad graduada dentro del corredor construido, no un
morfismo escalar \(54\to V^\natural\).

\subsection{Construcción reticular de orden \(2\)}

La ruta de orden tres anterior y la construcción clásica de
Frenkel--Lepowsky--Meurman producen dos presentaciones del mismo VOA lunar.
Para declarar la segunda, fijamos la extensión central
\begin{equation}
 1\longrightarrow\{\pm1\}
 \longrightarrow\widehat\Lambda_{24}
 \longrightarrow\Lambda_{24}
 \longrightarrow1
 \label{p12-83-c20-s8:eq:extension-central-leech-u031}
\end{equation}
con conmutador
\((-1)^{\langle\lambda,\mu\rangle}\), y un cociclo
\(\varepsilon(\lambda,\mu)\) que la representa. El álgebra de grupo
torcida es \(\mathbb C_\varepsilon[\Lambda_{24}]\), y
\begin{equation}
 V_\Lambda
 :=M(1)\otimes\mathbb C_\varepsilon[\Lambda_{24}]
 \label{p12-83-c20-s8:eq:voa-reticular-leech-u031}
\end{equation}
es el VOA reticular de carga central 24.

La negación \(\lambda\mapsto-\lambda\) posee una elevación
\(\theta\) compatible con \eqref{p12-83-c20-s8:eq:extension-central-leech-u031}. Su traza
es
\begin{equation}
 t_2(\tau)
 :=
 \operatorname{Tr}_{V_\Lambda}
 (\theta q^{L_0-1})
 =q^{-1}\prod_{n\ge1}(1+q^n)^{-24}.
 \label{p12-83-c20-s8:eq:t2-leech-u031}
\end{equation}
En efecto, la negación empareja \(e^\lambda\) con \(e^{-\lambda}\), y la
ausencia de torsión deja sólo \(\lambda=0\) en la traza reticular; cada uno
de los 24 osciladores aporta \((1+q^n)^{-1}\).

Sea \(V_\Lambda^T\) el módulo irreducible torcido por \(\theta\), y
\((V_\Lambda^T)^+\) su sector de signo positivo. La construcción FLM es
\begin{equation}
 V^\natural
 =(V_\Lambda)^+\oplus(V_\Lambda^T)^+.
 \label{p12-83-c20-s8:eq:vnatural-flm-u031}
\end{equation}
La ruta \(\mathbb Z/2\) de
\eqref{p12-83-c20-s8:eq:vnatural-flm-u031} y la ruta \(\mathbb Z/3\) de
\eqref{p12-83-c20-s8:eq:orbifold-triadico-vnatural-u031} son construcciones paralelas
identificadas por teoremas clásicos; ninguna sustituye a la otra.

\subsection{Grupo de automorfismos, carácter y series graduadas}

Una vez construido \(V^\natural\), aparecen dos publicaciones distintas:
\begin{align}
 \operatorname{Aut}(V^\natural)&\cong\mathbb M,
 \label{p12-83-c20-s8:eq:aut-vnatural-monstruo-u031}\\
 \operatorname{ch}_{V^\natural}(\tau)
 &=J(\tau)=j(\tau)-744
 =q^{-1}+196884q+21493760q^2+\cdots.
 \label{p12-83-c20-s8:eq:caracter-vnatural-j-u031}
\end{align}
La primera identifica el grupo de simetrías; la segunda cuenta dimensiones
graduadas. La línea de vacío es
\[
 (V^\natural)_0=\mathbb C\mathbf1,
\]
y la ausencia de raíces de \(\Lambda_{24}\) da
\[
 (V^\natural)_1=0.
\]
En grado dos,
\[
 196884=1+196883=196560+324=54(5\cdot729+1).
\]
Sólo la primera igualdad es una descomposición en dimensiones de
representaciones del Monstruo; las demás son controles aritméticos.

Para \(g\in\mathbb M\), se define la serie de McKay--Thompson
\begin{equation}
 T_g(\tau):=
 \operatorname{Tr}_{V^\natural}
 \left(gq^{L_0-1}\right).
 \label{p12-83-c20-s8:eq:serie-mckay-thompson-u031}
\end{equation}

\begin{theorem}[Moonshine monstruoso]
\label{p12-83-c20-s8:thm:moonshine-monstruoso-u031}
La acción graduada de \(\mathbb M\) sobre \(V^\natural\) produce la familia
\(\{T_g\}_{g\in\mathbb M}\). Los teoremas de Moonshine identifican estas
series con funciones modulares de género cero y prueban sus identidades de
replicabilidad
\cite{ConwayNorton1979,Borcherds1992,FLM1988}.
\end{theorem}

Este teorema tiene como entrada el par
\[
 \bigl(V^\natural,\operatorname{Aut}(V^\natural)\bigr)
\]
y su graduación. No es la consecuencia de una flecha cuyo dominio sea el
grupo abstracto \(\mathbb M\).

Para \(J=T_{1A}\), sean \(P_n\) los polinomios de Faber normalizados por
\[
 P_n(J)=q^{-n}+O(q).
\]
La identidad de Hecke--Faber es
\begin{equation}
 P_n(J)=n\,T_nJ.
 \label{p12-83-c20-s8:eq:hecke-faber-u031}
\end{equation}
Controla todos los grados; el coeficiente 196884 es sólo la primera
manifestación no trivial.

\subsection{Grafo de dependencias del corredor Paley–Witt–Golay–Leech–Moonshine}

\begin{figure}[htbp]
\centering
\begin{tikzcd}[column sep=huge,row sep=large]
 &C_W
   \arrow[dl,"W_{12}"']
   \arrow[d,"N(A_2^{12})"',font=\scriptsize]
   \arrow[dr,phantom,"\substack{\text{ramas}\\\text{distintas}}",
          description,pos=.72,font=\scriptsize]&\\
 M_{12}
 &\Lambda_{24}
   \arrow[d,"V_\Lambda"]
 &(\mathcal G_{24},D,\ell)
   \arrow[d,"W_{24}"]\\
 &V^\natural
   \arrow[dl,"\operatorname{Aut}"']
   \arrow[d,"\operatorname{ch}"]
   \arrow[dr,"g\mapsto T_g"]
 &M_{24}\\
 \mathbb M
 &J
 &\{T_g\}_{g\in\mathbb M}
\end{tikzcd}
\caption{Bifurcación excepcional y tres publicaciones de
\(V^\natural\). La entrada binaria \((\mathcal G_{24},D,\ell)\) es
independiente; no existe una flecha \(W_{24}\to\Lambda_{24}\).}
\label{p12-83-c20-s8:fig:grafo-excepcional-moonshine-u031}
\end{figure}

\begin{criterion}[Criterios de refutación del corredor lunar]
\label{p12-83-c20-s8:crit:refutacion-corredor-lunar-u031}
La construcción queda refutada si ocurre alguno de los hechos siguientes:
\begin{enumerate}
 \item \(\lambda\) no es biyectiva o
       \eqref{p12-83-c20-s8:eq:identidad-entrelazamiento-app-witt-u031} falla;
 \item \(c_*\notin C_W\), \(g_c\) posee puntos fijos o no preserva el
       vecino;
 \item los desplazamientos de
       \eqref{p12-83-c20-s8:eq:desplazamientos-vecino-u031} no pertenecen a \(N_0\);
 \item la elevación reticular \(\widetilde h\) no realiza el automorfismo
       monomial \(h\) en el discriminante;
 \item el peso torcido no es \(4/3\) o la elevación no es de tipo cero;
 \item el orbifold se identifica con \(V^\natural\) sin verificar las
       hipótesis del teorema clásico;
 \item se usa \(W_{24}\) como constructor de Leech;
 \item una igualdad de dimensiones se presenta como construcción de
       \(\mathbb M\);
 \item se serializan indebidamente las publicaciones
       \(\operatorname{Aut}\), \(\operatorname{ch}\) y \(g\mapsto T_g\).
\end{enumerate}
\end{criterion}
 \section{Control theta del retículo de Leech: identidad \texorpdfstring{\(\Theta_{\Lambda}=E_4^3-720\Delta\)}{Theta Lambda = E4^3 - 720 Delta} y separación entre coincidencia aritmética y morfismo estructural}
\label{sec:control-theta-leech}

Desde \(\Lambda_{24}\), la cadena
\[
\Lambda_{24}
\longrightarrow V^\natural
\longrightarrow\mathbb M
\longrightarrow\text{Moonshine}
\]
es matemática clásica: álgebra de operadores de vértice, grupo Monstruo y
funciones McKay--Thompson. La contribución verificada aquí está en la entrada
finita, la selección del borde y el \(3\)-vecino sin raíces, no en atribuirse de
nuevo esos teoremas.

Tres identidades numéricas acompañan esta cadena:
\[
\Theta_{\Lambda}=E_4^3-720\Delta,
\]
\[
196884=196560+324,
\]
\[
196884=54(3645+1)=54(5\cdot729+1).
\]
Son identidades aritméticas exactas, pero sin un morfismo o una descomposición
de caracteres no constituyen por sí mismas señales matemáticas de
compatibilidad. Ninguna sustituye la construcción de \(V^\natural\) ni
proporciona una acción directa del Monstruo sobre palabras decimales HMT.
 \section{Uniformización de orden \texorpdfstring{\(2\)}{2} en Moonshine: funciones \texorpdfstring{\(T_{2A}\)}{T2A} y \texorpdfstring{\(j\)}{j}, estructura modular y control del corredor micro--macro}
\label{sec:voa-moonshine-orden-dos-rev11}
\index{VOA!retículo de Leech}
\index{Moonshine!clase 2A}

La rama de incidencia construye primero el retículo de Leech
\(\Lambda_{24}\). A partir de ese dato, la teoría clásica de álgebras de
operadores de vértice proporciona una prolongación canónica. Si
\(\mathfrak h=\Lambda_{24}\otimes_{\mathbb Z}\mathbb C\) y \(M(1)\) es
el módulo de Fock de su álgebra de Heisenberg, el VOA reticular es
\begin{equation}
 V_{\Lambda}=M(1)\otimes\mathbb C[\Lambda_{24}],
 \label{eq:voa-reticular-leech-rev11}
\end{equation}
con carga central \(24\). Para \(\alpha\in\Lambda_{24}\), sus operadores
de vértice se escriben
\[
 Y(e^\alpha,z)=E^-(-\alpha,z)E^+(-\alpha,z)e^\alpha z^{\alpha_0}.
\]
La ausencia de raíces en \(\Lambda_{24}\) implica la anulación del sector
de peso uno en el orbifold de Moonshine.

La involución
\[
 \theta:\Lambda_{24}\longrightarrow\Lambda_{24},
 \qquad \theta(\lambda)=-\lambda,
\]
es intrínseca y se prolonga a \(V_{\Lambda}\). Su traza graduada se obtiene
por conteo directo de los modos osciladores.

\begin{theorem}[Traza de la involución de negación]
\label{thm:traza-theta-leech-rev11}
\begin{equation}
 t_2(\tau)
 :=\operatorname{Tr}_{V_\Lambda}
 \bigl(\theta q^{L_0-1}\bigr)
 =q^{-1}\prod_{n\geq1}(1+q^n)^{-24}.
 \label{eq:t2-producto-leech-rev11}
\end{equation}
\end{theorem}
\begin{proof}
La involución empareja \(e^\lambda\) con \(e^{-\lambda}\); como el retículo
es libre de torsión, sólo \(\lambda=0\) contribuye a la traza del sector
reticular. Cada uno de los veinticuatro osciladores de energía \(n\) aporta
\(\sum_{k\geq0}(-1)^kq^{nk}=(1+q^n)^{-1}\). El desplazamiento del vacío
produce el factor \(q^{-1}\).
\end{proof}

La completación de Fricke asociada a esta traza es
\begin{equation}
 T_{2A}(\tau)=t_2(\tau)+24+\frac{2^{12}}{t_2(\tau)}
 =q^{-1}+4372q+96256q^2+1240002q^3+\cdots.
 \label{eq:T2A-hmt-rev11}
\end{equation}
El término \(24\) cancela el término constante de \(t_2\); el factor
\(2^{12}\) es la degeneración del sector torcido de veinticuatro bosones.
La misma serie de borde determina el invariante modular mediante
\begin{equation}
 \boxed{j(\tau)=\frac{(t_2(\tau)+256)^3}{t_2(\tau)^2}}.
 \label{eq:j-desde-t2-rev11}
\end{equation}
Las ecuaciones \eqref{eq:t2-producto-leech-rev11}--
\eqref{eq:j-desde-t2-rev11} construyen la rama canónica de orden dos sin
parámetros continuos.

El orbifold de Frenkel--Lepowsky--Meurman se define por
\[
 V^\natural=(V_\Lambda)^+\oplus(V_\Lambda^T)^+.
\]
Es un VOA holomorfo de carga central \(24\), satisface
\((V^\natural)_1=0\), posee carácter
\[
 \operatorname{ch}_{V^\natural}(\tau)=j(\tau)-744
\]
y su grupo de automorfismos es el grupo Monstruo
\cite{FLM1988,ConwayNorton1979,Borcherds1992}. En esta cadena HMT fija el
retículo de Leech y la involución \(\theta\); los teoremas clásicos de VOA y
FLM identifican la prolongación resultante con \(V^\natural\) y su grupo de
automorfismos. La afirmación no requiere una acción del Monstruo sobre las
cifras: la evaluación arquimediana y la rama excepcional descienden de una
preimagen HMT común y conservan invariantes diferentes.

La función vecina
\(T_{3A}=t_3+12+729/t_3\) pertenece a la teoría clásica de Hauptmoduln de
orden tres. La realización HMT activa en el corredor trítico es la clase
\(3B\); una identificación con \(3A\) requiere declarar el levantamiento específico
que cambie de clase.
 \subsection{Entrelazamiento APP--Witt, vecino de Leech y orbifold de orden \(3\)}

Las doce fibras llevan las etiquetas
\[
 (i,\mu,\varepsilon)\in
 \mathbb Z/3\mathbb Z\times\{0,1\}\times\{0,1\},
 \qquad
 \lambda(i,\mu,\varepsilon)=4i+2\mu+\varepsilon\pmod{12}.
\]
La aplicación \(\lambda\) es biyectiva y, fijados trivializaciones y marcos,
los entrelazadores satisfacen
\[
 g_{\lambda(i,\mu,\varepsilon)}T_{i\mu\varepsilon}
 =T_{i\mu\varepsilon}C^{(-1)^\mu}.
\]
Su suma directa es un isomorfismo \(C_3\)-equivariante
\[
 T_\lambda:W_{\mathrm{APP}}\otimes\mathbb C
 \xrightarrow{\ \cong\ }
 \bigoplus_{b=0}^{11}(A_{2,b}\otimes\mathbb C).
\]

Sea \(N=N(A_2^{12})\), sea \(v=v_\alpha\) el vector radial de norma \(54\),
y definamos
\[
 N_0=\ker\bigl(\langle\,\cdot\,,v\rangle\bmod3\bigr),
 \qquad N_v=N_0+\mathbb Z\frac v3.
\]
El cálculo reticular verifica que \(N_v\cong\Lambda_{24}\), que la
isometría \(g_c\) tiene orden tres y no posee puntos fijos no nulos, y que
\[
 \frac{g_cv-v}{3},\quad\frac{g_c^{-1}v-v}{3}\in N_0.
\]
Por ello \(g_cN_v=N_v\). Su elevación estándar al VOA reticular tiene peso
torcido \(4/3\) y tipo cero. Los teoremas clásicos de orbifold cíclico dan
\[
 V_{(3)}:=\operatorname{Orb}_{\mathbb Z/3}
       (V_{N_v},\widehat g_c)\cong V^\natural,
 \qquad T_{3B}=t_3+12.
\]
El carácter se recupera por el promedio de los nueve sectores y por la
uniformización de nivel tres:
\[
 \operatorname{ch}_{V_{(3)}}
 =\frac13\sum_{i,j\in\mathbb Z/3}Z_{i,j}=J,
\]
\[
 J=t_3+12+\frac{10\,3^9}{t_3}
       +\frac{4\,3^{14}}{t_3^2}+\frac{3^{18}}{t_3^3}.
\]
En particular,
\[
 [q]J=54+10\cdot3^9=196884=54(5\cdot729+1).
\]

\subsection{Construcciones de órdenes \(2\) y \(3\) del álgebra de operadores de vértice \texorpdfstring{\(V^\natural\)}{V natural}}

Fijemos la extensión central y el cociclo de Leech
\[
 1\longrightarrow\{\pm1\}\longrightarrow\widehat\Lambda_{24}
 \longrightarrow\Lambda_{24}\longrightarrow1,
 \qquad
 V_\Lambda=M(1)\otimes\mathbb C_\varepsilon[\Lambda_{24}].
\]
La negación \(\lambda\mapsto-\lambda\) se eleva a \(\theta\), y
\[
 t_2(\tau)=\operatorname{Tr}_{V_\Lambda}
  \left(\theta q^{L_0-1}\right)
 =q^{-1}\prod_{n\ge1}(1+q^n)^{-24}.
\]
La presentación FLM es
\[
 V_{(2)}:=(V_\Lambda)^+\oplus(V_\Lambda^T)^+\cong V^\natural.
\]

\begin{theorem}[Equivalencia tipada de las presentaciones de órdenes \(3\) y \(2\)]
\label{thm:cierre-vtm-equivalencia-voa-23}
Sean
\[
 \phi_3:V_{(3)}\xrightarrow{\ \cong\ }V^\natural,
 \qquad
 \phi_2:V_{(2)}\xrightarrow{\ \cong\ }V^\natural
\]
los isomorfismos graduados proporcionados por los teoremas de orbifold
aplicados a los datos anteriores. Entonces
\[
 \boxed{\Phi_{23}:=\phi_2^{-1}\circ\phi_3:
 V_{(3)}\xrightarrow{\ \cong\ }V_{(2)}}
\]
es un isomorfismo graduado de álgebras de operadores de vértice. Además,
\[
 a\longmapsto\Phi_{23}a\Phi_{23}^{-1}
\]
identifica sus grupos de automorfismos y conserva las trazas graduadas.
\end{theorem}

\begin{proof}
La composición de isomorfismos de VOA es un isomorfismo de VOA y preserva el
vacío, el vector conforme, los operadores de vértice y la gradación por
\(L_0\). La conjugación transporta automorfismos. Para cada automorfismo
\(a\), la invariancia de la traza bajo conjugación da la igualdad de las
series graduadas correspondientes.
\end{proof}

El isomorfismo \(\Phi_{23}\) depende de las identificaciones \(\phi_2\) y
\(\phi_3\); no se afirma canonicidad absoluta. La construcción de orden dos
conserva además el contenido exclusivo
\[
 T_{2A}=t_2+24+\frac{2^{12}}{t_2},
 \qquad
 \boxed{j=\frac{(t_2+256)^3}{t_2^2}},
\]
que completa la uniformización modular de orden dos.

Una vez construido \(V^\natural\), las tres publicaciones se enuncian sin
serializarlas indebidamente:
\[
 \operatorname{Aut}(V^\natural)\cong\mathbb M,
 \qquad
 \operatorname{ch}_{V^\natural}=J,
 \qquad
 g\longmapsto T_g(\tau)
 =\operatorname{Tr}_{V^\natural}\!\left(gq^{L_0-1}\right).
\]
El teorema de Moonshine coordina la acción graduada con las funciones de
género cero y sus identidades de replicabilidad. La acción del Monstruo vive
en \(V^\natural\); que no sea una permutación puntual de cifras es un control
de tipos final, no la tesis del corredor.

 \section{Cadena probatoria continua de la publicación excepcional: del despliegue APP a Paley--Witt,
\texorpdfstring{\(V^\natural\)}{V natural} y Moonshine}
\label{sec:prueba-corredor-excepcional}

La publicación excepcional parte del lector de incidencia
\(\mathfrak E_{\rm exc}\) de \eqref{eq:frontal-doble-proyeccion}. No toma como
entrada un código, un retículo o un grupo excepcional ya dados: toma la
frontera, la orientación, la hoja y la memoria de una historia enriquecida y
las expresa en el calibre proyectivo fijado por la holonomía nonádica. Sólo
después se aplican los teoremas clásicos de códigos, retículos y álgebras de
operadores de vértice a los objetos producidos.

\subsection{\texorpdfstring{Despliegue APP y cono graduado del índice \(54\)}{Despliegue APP y cono graduado del índice 54}}

Sea \(R\) la involución que intercambia las hojas APP y sean
\(P_+=(I+R)/2\), \(P_-=(I-R)/2\) sus proyectores. El bloque central
\begin{equation}
 B_c=7I+2R=9P_++5P_-
 \label{eq:exc-bloque-central-app}
\end{equation}
produce cuatro objetos de tipos distintos:
\begin{equation}
 9-5=4,\qquad (B_c)_{i\ne j}=2,\qquad
 \sum M_\Pi-\sum M_\Sigma=51-45=6,\qquad
 9\cdot6=459-405=54.
 \label{eq:exc-despliegue-42654}
\end{equation}
El \(4\) es el salto espectral entre hojas, el \(2\) el acoplamiento local,
el \(6\) la diferencia de las compresiones modulares y el \(54\) la
integración de esa diferencia sobre las nueve posiciones. Por tanto,
\(4\to2\to6\to54\) designa una jerarquía de morfismos y evaluaciones, no
una cadena de igualdades entre objetos.

Fijadas la orientación \(o_{\Sigma\Pi}\), la paridad de hoja y la
normalización triangular, la rama de Fock posee el acoplamiento
\begin{equation}
 \mathcal H_{\Sigma,\Pi}^{(m)}:
 M_{\rm bal}^{C_3,-}\longrightarrow
 \operatorname{End}\!\bigl(\mathcal F_2(\mathcal V_m)\bigr)
 \otimes o_{\Sigma\Pi},
 \label{eq:exc-acoplamiento-fock}
\end{equation}
y su traza orientada satisface
\begin{equation}
 \operatorname{Tr}_{o,\mathcal F_2}
 \bigl(\mathcal H_{\Sigma,\Pi}^{(m)}(\nu)\Gamma_2(g_m)\bigr)
 =-\frac{3m\,c(\nu)}2.
 \label{eq:exc-traza-fock}
\end{equation}
Para \(m=12\) y \(c(\delta)=-3\), vale \(54\). Independientemente, el
producto ciclotómico
\begin{equation}
 t_3(q)=q^{-1}\prod_{n\ge1}(1+q^n+q^{2n})^{-12}
 =q^{-1}-12+54q-76q^2-243q^3+\cdots
 \label{eq:exc-t3}
\end{equation}
da \([q]t_3=54\).

\begin{theorem}[Cono tipado del índice común]
\label{thm:exc-cono-54}
Una vez construidos por separado el censo APP, la rama de Fock, el producto
ciclotómico, el vecino reticular de
\cref{eq:exc-vecino-indice-tres} y el módulo lunar de
\cref{eq:exc-flm}, sus lectores verifican
\begin{equation}
 D_{\rm APP}
 =\operatorname{Tr}_{o,\mathcal F_2}
  \bigl(\mathcal H_{\Sigma,\Pi}^{(12)}(\delta)\Gamma_2(g_{12})\bigr)
 =[q]t_3=v^2=\frac{196884}{5\cdot729+1}=54.
 \label{eq:exc-cono-54}
\end{equation}
La igualdad reúne evaluaciones en un cono hacia \(\mathbb Z\); no
identifica sus dominios ni usa \(196884\) para seleccionar el corredor.
\end{theorem}

\begin{proof}
Las tres primeras igualdades son
\eqref{eq:exc-despliegue-42654}, \eqref{eq:exc-traza-fock} y
\eqref{eq:exc-t3}. Para el vector radial
\(v=4\rho_1+\rho_2+\cdots+\rho_{12}\) de
\eqref{eq:exc-vector-vecino}, la ortogonalidad de los factores \(A_2\) da
\(v^2=4^2\cdot2+11\cdot2=54\). Finalmente,
\(5\cdot729+1=3646\) y \(54\cdot3646=196884\). Cada identidad se prueba en
su propio codominio antes de compararlas.
\end{proof}

\subsection{\texorpdfstring{Forma de Paley y raíz cuadrática interna de \(-I_6\)}{Forma de Paley y raíz cuadrática interna de -I6}}

Sea
\[
 \mathcal U_9=(\mathbb Z/9\mathbb Z)^\times=\{1,2,4,8,7,5\},
\]
en el orden de las potencias de \(2\), y sea
\begin{equation}
 \theta:\mathcal U_9\xrightarrow{\sim}\mathbb P^1(\mathbb F_5),
 \qquad
 \begin{array}{c|cccccc}
 u&1&2&4&8&7&5\\ \hline
 \theta(u)&\infty&0&1&2&3&4
 \end{array}.
 \label{eq:exc-calibre-proyectivo}
\end{equation}
Esta biyección es un calibre de coordenadas del lector; no se usa como
homomorfismo. Para el carácter cuadrático
\(\chi:\mathbb F_5\to\{-1,0,1\}\), defínase
\begin{equation}
 C_P=
 \begin{pmatrix}
 0& 1& 1& 1& 1& 1\\
 1& 0& 1&-1&-1& 1\\
 1& 1& 0& 1&-1&-1\\
 1&-1& 1& 0& 1&-1\\
 1&-1&-1& 1& 0& 1\\
 1& 1&-1&-1& 1& 0
 \end{pmatrix},
 \label{eq:exc-matriz-paley}
\end{equation}
donde las entradas finitas fuera de la diagonal son \(\chi(x-y)\) y las
entradas incidentes con \(\infty\) son uno.

\begin{lemma}[Correlación cuadrática]
\label{lem:exc-correlacion-cuadratica}
Si \(a\ne b\) en \(\mathbb F_5\), entonces
\[
 \sum_{t\in\mathbb F_5}\chi(t-a)\chi(t-b)=-1.
\]
\end{lemma}

\begin{proof}
La sustitución \(u=(t-a)/(b-a)\) reduce la suma a
\(\sum_u\chi(u(u-1))\). Para \(u=0,1,2,3,4\), los sumandos son
\(0,0,1,-1,-1\).
\end{proof}

\begin{proposition}[Identidad de conferencia y reducción ternaria]
\label{prop:exc-paley-witt}
La matriz \(C_P\) es simétrica y satisface
\[
 C_PC_P^{\mathsf T}=5I_6.
\]
Su reducción \(A_W=C_P\bmod3\) verifica
\begin{equation}
 \boxed{A_W^2=-I_6\quad\text{en }M_6(\mathbb F_3).}
 \label{eq:exc-aw-cuadrado}
\end{equation}
\end{proposition}

\begin{proof}
Cada fila de \(C_P\) tiene norma cuadrada cinco. Dos filas finitas
distintas tienen producto escalar \(1-1=0\) por el
\cref{lem:exc-correlacion-cuadratica}; la suma de un carácter no trivial
prueba también la ortogonalidad con la fila de \(\infty\). Al reducir módulo
tres se obtiene
\(A_WA_W^{\mathsf T}=5I_6=2I_6=-I_6\). Como \(A_W\) es simétrica,
\eqref{eq:exc-aw-cuadrado} sigue.
\end{proof}

La identidad \eqref{eq:exc-aw-cuadrado} construye la estructura compleja
finita del corredor. En efecto,
\[
 \mathbb F_9=\mathbb F_3[\iota]/(\iota^2+1),
 \qquad
 P_{\pm\iota}=\frac12(I\mp\iota A_W)
\]
son proyectores complementarios de rango tres y dotan a
\(\mathbb F_3^6\) de estructura de \(\mathbb F_9\)-módulo de rango tres.
La acción es explícita:
\begin{equation}
 (a+b\iota)\cdot v=av+b(vA_W).
 \label{eq:exc-accion-f9}
\end{equation}
Por tanto, \(3^6=9^3\) no se usa como sustituto cardinal de una acción.

\subsection{Código ternario autodual y sistema de Witt}

Defínase
\begin{equation}
 c:\mathbb F_3^6\longrightarrow\mathbb F_3^{12},
 \qquad c(q)=(q,qA_W),
 \qquad C_W:=\operatorname{im}c.
 \label{eq:exc-codigo-grafico}
\end{equation}

\begin{theorem}[Código ternario producido por la incidencia]
\label{thm:exc-codigo-ternario}
La aplicación \(c\) es inyectiva y
\begin{equation}
 C_W=C_W^\perp,
 \qquad
 W_{C_W}(y)=1+264y^6+440y^9+24y^{12}.
 \label{eq:exc-enumerador-cw}
\end{equation}
En particular, \(C_W\) tiene parámetros \([12,6,6]_3\).
\end{theorem}

\begin{proof}
La primera proyección de \(c(q)\) es \(q\), luego \(c\) es inyectiva.
Con \(G_W=(I_6\mid A_W)\),
\[
 G_WG_W^{\mathsf T}=I_6+A_WA_W^{\mathsf T}=0
\]
en \(\mathbb F_3\). Así \(C_W\subseteq C_W^\perp\), y la igualdad de
dimensiones da la autodualidad. El enumerador es la suma finita exacta
\[
 \sum_{q\in\mathbb F_3^6}
 y^{\operatorname{wt}(q)+\operatorname{wt}(qA_W)}.
\]
Evaluar las \(3^6=729\) filas de la matriz explícita
\eqref{eq:exc-matriz-paley} da, respectivamente, los censos
\(1,264,440,24\) en pesos \(0,6,9,12\), y cero en los restantes. La suma
de los cuatro censos es \(729\), por lo que no queda ninguna palabra fuera
del certificado finito.
\end{proof}

Los \(264\) vectores de peso seis aparecen en pares \(\{x,-x\}\); sus
\(132\) soportes forman el sistema de Steiner
\begin{equation}
 W_{12}=S(5,6,12).
 \label{eq:exc-w12}
\end{equation}
Éste es el reconocimiento, después de construir \(C_W\), del código ternario
extendido de Golay y de su diseño de Witt
\citep{MacWilliamsSloane1977,ConwaySloane1999}. No se introduce el diseño
para seleccionar la matriz \(A_W\).

La prolongación binaria se mantiene separada. Si
\(\mathcal G_{24}\subset\mathbb F_2^{24}\) es el código binario extendido
de Golay, \(D\) una dodecada y
\(\ell:\{1,\ldots,12\}_{\rm HMT}\to D\) un isomorfismo de diseños, los
soportes de peso ocho forman \(W_{24}=S(5,8,24)\). Para una tétrada fija,
\[
 \lambda_4(W_{12})=4,
 \qquad
 \lambda_4(W_{24})=5,
\]
de modo que cuatro octadas elevan las cuatro hexadas incidentes y una quinta
es transversal. El dato \((\mathcal G_{24},D,\ell)\) pertenece al codominio
binario; no existe una flecha canónica \(W_{24}\to\Lambda_{24}\).

\subsection{Pegado ternario y vecino de Leech}

Sea \(Q=A_2^{12}\), con \(\det Q=3^{12}\). El discriminante de cada
factor \(A_2\) es \(\mathbb F_3\), y el pegado definido por
\eqref{eq:exc-codigo-grafico} es
\begin{equation}
 N:=\{x\in Q^*:x\bmod Q\in C_W\}.
 \label{eq:exc-niemeier-pegado}
\end{equation}

\begin{theorem}[Retículo de Niemeier producido por el código]
\label{thm:exc-niemeier}
El retículo \(N\) es positivo definido, par y unimodular de rango
veinticuatro, y su sistema de raíces es \(A_2^{12}\). Por consiguiente,
\(N\cong N(A_2^{12})\).
\end{theorem}

\begin{proof}
La autodualidad de \(C_W\) hace integral el pegado. Todos sus pesos son
múltiplos de tres, lo que hace par la extensión. Además,
\([N:Q]=|C_W|=3^6\), y por ello
\[
 \det N=\frac{\det Q}{[N:Q]^2}
 =\frac{3^{12}}{3^{12}}=1.
\]
Una clase discriminante no nula de \(A_2\) tiene norma mínima \(2/3\).
Como toda palabra no nula de \(C_W\) tiene peso al menos seis, todo vector
de pegado exterior a \(Q\) tiene norma al menos cuatro. Los vectores de
norma dos son exactamente las raíces de los doce factores \(A_2\). La
clasificación de los retículos pares unimodulares de rango veinticuatro
identifica el tipo indicado \citep{ConwaySloane1999}.
\end{proof}

Sea \(\rho_i\) la raíz máxima del factor \(A_2^{(i)}\), normalizada por
\(\rho_i^2=2\), y defínase
\begin{equation}
 v=4\rho_1+\rho_2+\cdots+\rho_{12},
 \qquad v^2=54.
 \label{eq:exc-vector-vecino}
\end{equation}
El marco ordenado del lector excepcional fija cuál de las doce fibras ocupa
la primera posición; ningún valor arquimediano interviene en esta selección.
Póngase
\begin{equation}
 \chi_v(x)=\langle x,v\rangle\bmod3,
 \quad N_0=\ker\chi_v,
 \quad N_v=N_0+\mathbb Z\frac v3.
 \label{eq:exc-vecino-indice-tres}
\end{equation}

\begin{lemma}[Estructura del vecino]
\label{lem:exc-estructura-vecino}
El carácter \(\chi_v\) es sobreyectivo, \([N:N_0]=3\),
\(v/3\in N_0^*\), y \(N_v\) es par y unimodular.
\end{lemma}

\begin{proof}
Para toda raíz de un factor \(A_2\), el producto con la raíz máxima es
\(\pm1\) o \(\pm2\); los coeficientes de \(v\) son congruentes con uno
módulo tres. Por tanto, ninguna raíz pertenece a \(N_0\) y \(\chi_v\) es
no trivial, luego sobreyectivo. La definición del núcleo da índice tres y
\(\det N_0=9\). Para \(x\in N_0\),
\(\langle x,v/3\rangle\in\mathbb Z\), mientras
\((v/3)^2=6\). Así la extensión por \(v/3\) es integral y par. Como
\([N_v:N_0]=3\),
\(\det N_v=9/3^2=1\).
\end{proof}

\begin{lemma}[Mínimo de las clases nuevas]
\label{lem:exc-minimo-clases}
Cada coordenada de una clase local que aparece en
\(N_0\pm v/3\) tiene norma mínima \(2/9\). En consecuencia,
\[
 \|x\pm v/3\|^2\ge12\cdot\frac29=\frac83
 \qquad(x\in N_0).
\]
\end{lemma}

\begin{proof}
Escribiendo una coordenada de \(A_2^*\) en la base de raíces simples, las
seis clases orientadas relevantes son los desplazamientos de
\(\pm\rho/3\) por las tres clases discriminantes. Completar cuadrados en la
forma de Gram
\(\bigl(\begin{smallmatrix}2&-1\\-1&2\end{smallmatrix}\bigr)\)
da mínimo \(2/9\) en las seis clases. La suma ortogonal de las doce
coordenadas da la desigualdad.
\end{proof}

\begin{theorem}[Construcción del retículo de Leech]
\label{thm:exc-leech}
El vecino \(N_v\) de \eqref{eq:exc-vecino-indice-tres} es isométrico al
retículo de Leech:
\begin{equation}
 \boxed{N_v\cong\Lambda_{24}.}
 \label{eq:exc-leech}
\end{equation}
\end{theorem}

\begin{proof}
Las raíces de \(N\) no pertenecen a \(N_0\), y las dos clases nuevas no
contienen vectores de norma dos por el
\cref{lem:exc-minimo-clases}. El \cref{lem:exc-estructura-vecino} prueba que
\(N_v\) es positivo definido, par, unimodular y de rango veinticuatro. La
unicidad del retículo con esas propiedades y sin raíces lo identifica con
\(\Lambda_{24}\) \citep{ConwaySloane1999}.
\end{proof}

\subsection{Isometría de orden tres, orbifold y módulo lunar}

La rotación de Coxeter de cada factor \(A_2\) induce una isometría
\(g_c\) de orden tres. La estabilidad del código \(C_W\) prueba que
\(g_cN=N\). Para el marco orientado que fija \(v\), los desplazamientos
\begin{equation}
 y_*:=\frac{g_cv-v}{3},
 \qquad
 z_*:=\frac{g_c^{-1}v-v}{3}
 \label{eq:exc-desplazamientos-vecino}
\end{equation}
pertenecen a \(N_0\): sus palabras discriminantes son un par
\(c_*,-c_*\in C_W\), y
\(\langle y_*,v\rangle=\langle z_*,v\rangle=-27\). Por tanto,
\begin{equation}
 g_cN_v=N_v.
 \label{eq:exc-preservacion-vecino}
\end{equation}
La acción no tiene puntos fijos en \(N_v\otimes\mathbb R\), pues cada
factor \(A_2\) posee los dos autovalores primitivos de orden tres.

Sea \(\widehat g_c\) la elevación estándar a la VOA reticular
\(V_{\Lambda_{24}}\). Sus dos autoespacios no triviales tienen dimensión
doce. La energía de vacío del módulo torcido es, por tanto,
\begin{equation}
 \rho_{g_c}
 =\frac1{4\cdot3^2}
 \sum_{k=1}^{2}k(3-k)\dim\mathfrak h_{(k)}
 =\frac43.
 \label{eq:exc-peso-torcido}
\end{equation}
La congruencia \(3^2\rho_{g_c}=12\equiv0\pmod3\) prueba que la elevación
es de tipo cero.

\begin{theorem}[Orbifold triádico]
\label{thm:exc-orbifold}
El orbifold cíclico de la elevación anterior satisface
\begin{equation}
 \operatorname{Orb}_{\mathbb Z/3}
 (V_{\Lambda_{24}},\widehat g_c)\cong V^\natural.
 \label{eq:exc-orbifold}
\end{equation}
El automorfismo dual pertenece a la clase \(3B\).
\end{theorem}

\begin{proof}
Se han verificado las tres hipótesis que no pueden sustituirse por una
igualdad de series: isometría fijo-libre de orden tres, preservación del
vecino y tipo cero. Los teoremas de orbifold cíclico identifican entonces
la extensión holomorfa con el módulo lunar y clasifican el automorfismo dual
\citep{ChenLamShimakura2016,AbeLamYamada2017,Carnahan2017}.
\end{proof}

La construcción de orden dos es una realización paralela. Para la elevación
\(\theta\) de la negación en \(\Lambda_{24}\),
\begin{equation}
 V^\natural=(V_{\Lambda_{24}})^+
 \oplus(V_{\Lambda_{24}}^T)^+.
 \label{eq:exc-flm}
\end{equation}
Las rutas de orden dos y tres se identifican por los teoremas de VOA; no se
serializan causalmente una dentro de la otra \citep{FLM1988}.

\begin{theorem}[Publicaciones del módulo lunar]
\label{thm:exc-moonshine}
Una vez construido \(V^\natural\), se obtienen
\begin{align}
 \operatorname{Aut}(V^\natural)&\cong\mathbb M,
 \label{eq:exc-aut-monster}\\
 \operatorname{ch}_{V^\natural}(\tau)
 &=J(\tau)=j(\tau)-744
 =q^{-1}+196884q+21493760q^2+\cdots,
 \label{eq:exc-caracter-j}\\
 T_g(\tau)&=\operatorname{Tr}_{V^\natural}
 \bigl(gq^{L_0-1}\bigr),\qquad g\in\mathbb M.
 \label{eq:exc-mckay-thompson}
\end{align}
Las series \(T_g\) son las funciones modulares de género cero del teorema
de Moonshine y satisfacen las identidades de replicabilidad.
\end{theorem}

\begin{proof}
La primera identificación y la construcción graduada proceden de la
realización FLM. La coordinación entre la acción graduada, las series de
McKay--Thompson y sus propiedades modulares es el teorema de Moonshine de
Conway--Norton y Borcherds
\citep{FLM1988,ConwayNorton1979,Borcherds1992}. En particular,
\((V^\natural)_0=\mathbb C\mathbf1\) y
\((V^\natural)_1=0\); el coeficiente \(196884=1+196883\) es la primera
descomposición no trivial en dimensiones de representaciones del Monstruo.
\end{proof}

\begin{falsifier}[Corredor excepcional]
\label{fals:exc-corredor}
La publicación queda refutada si falla cualquiera de los pasos efectivos:
\eqref{eq:exc-despliegue-42654}, la traza
\eqref{eq:exc-traza-fock}, el coeficiente \([q]t_3=54\),
\(C_PC_P^{\mathsf T}=5I_6\), \(A_W^2=-I_6\), la autodualidad o el
enumerador de \(C_W\), la paridad o unimodularidad del pegado, la ausencia
de raíces del vecino, su preservación por \(g_c\), el peso torcido
\(4/3\) o el tipo cero. También queda refutada una narración que utilice
\(W_{24}\) para construir \(\Lambda_{24}\), que confunda los dos conjuntos
de \(243\) por su cardinal, o que haga actuar al Monstruo directamente
sobre cifras y rutas. La acción de \(\mathbb M\) reside en \(V^\natural\).
\end{falsifier}
 
\chapter{Imagen espectral--polar y reducción exacta de Riemann--Weil}
\begingroup
\renewcommand{\chapter}[2][]{}%
\chapter{Imagen espectral--polar y reducción exacta de Riemann--Weil}

\section{Punto de partida: la estructura espectral--polar completa}

La entrada de la resolución no es la forma de Weil aislada ni una lista de
ceros. Es la estructura espectral--polar completa reproducida como dato en este
volumen. La cadena causal propia de este capítulo es
\begin{equation}\label{eq:pdf2-rh-cadena-procedencia}
 \mathfrak G_{\rm sp}
 \xrightarrow{\ \operatorname{pr}_{X,{\rm sp}}\ }X_{\chi_{\rm sp}}
 \xrightarrow{\ \mathcal A_{\rm sp}\ }\mathfrak M_{\rm sp}
 \xrightarrow{\ \mathcal P_{\rm sp}\ }\mathcal C_{\rm GW}.
\end{equation}
El codominio clásico aparece sólo en la última flecha; no selecciona
retrospectivamente el estado, las cartas, la proyección ni el terminal.

La estructura de partida contiene el objeto de trece coordenadas
\begin{equation}\label{eq:pdf2-rh-res-sp-completa}
\begin{aligned}
 \mathfrak G_{\rm sp}^{(13)}=\bigl(&
 \mathcal F_{\rm sp},
 \operatorname{Ext}_{\Gamma,{\rm sp}},
 \operatorname{Rel}_{\rm sp},
 \mathcal N_{\rm sp},
 \operatorname{or}_{\rm sp},
 \operatorname{Carr}_{\rm sp},
 \operatorname{Mem}_{\rm sp},\\
 &\partial_{\rm sp},
 \gamma_{\rm sp},
 \Sigma_{\rm sp}^{\rm mult},
 \operatorname{Surv}_{\rm sp},
 \tau_{\rm sp}^{\rm term},
 \mathcal L_{\rm sp}\bigr).
\end{aligned}
\end{equation}
Cada coordenada interviene materialmente:
\begin{enumerate}
 \item \(\mathcal F_{\rm sp}\) conserva las dos hojas, su orientación, el
 par residuo--cociente y la firma de la publicación doble.
 \item \(\operatorname{Ext}_{\Gamma,{\rm sp}}\) transporta esa fibra por la
 clase dirigida de cortes y niveles nonádicos.
    \item \(\operatorname{Rel}_{\rm sp}\) contiene el cociclo de acarreo, la
    ortogonalidad \(P=P^*=P^2\), las dos publicaciones y el complemento
    estructural \(B_{\rm sp}^{\rm str}:=\Xi^*(I-P)\Xi\). La identificación
    de este complemento con la forma de Guinand--Weil es una flecha distinta
    y se audita por separado en este capítulo.
 \item \(\mathcal N_{\rm sp}\) registra el nivel \(m\), la ventana \(R\),
 el cociente \(q=5/9\), los momentos de ambas hojas y los índices de los
 canales primo, gamma y polar.
 \item \(\operatorname{or}_{\rm sp}=\widehat R=(R,K)\) mantiene distintas
 las dos cartas orientadas.
 \item \(\operatorname{Carr}_{\rm sp}=\operatorname{Carr}_R\) es el
 acarreo operatorial y no se anula al publicar una traslación.
 \item \(\operatorname{Mem}_{\rm sp}=(S_0,M,\mathcal O_9)\) conserva la
 identidad finita de Stein y su terminal.
 \item \(\partial_{\rm sp}=(I-P)\Xi\) es la frontera que no aparece en la
 compresión \(P\Xi\) y cuya energía define el complemento estructural. Su
 identificación con la forma de Weil no se incorpora a esta definición.
 \item \(\gamma_{\rm sp}\) es la ruta cofinal de ventanas, con un único
 regulador para todos los canales.
 \item \(\Sigma_{\rm sp}^{\rm mult}\) es la multisección de las dos
 publicaciones y sus componentes primo--gamma--polar; no es la firma de un
 único estado.
 \item \(\operatorname{Surv}_{\rm sp}\) conserva la familia bajo
 \(m\mapsto m+1\) y \(\mathcal D_R\hookrightarrow\mathcal D_{R'}\).
 \item \(\tau_{\rm sp}^{\rm term}\) conserva el terminal
 \(M^{*N}S_0M^N\), que en la carta espejo se publica mediante \(E_R\).
    \item \(\mathcal L_{\rm sp}=(\Xi,P,B_{\rm sp}^{\rm str})\) es el lector final; se aplica
 después de antecedente, cartas, memoria, ruta y terminal.
\end{enumerate}

El ensamblador y el publicador son
\begin{equation}\label{eq:pdf2-rh-ensamblador-sp}
\begin{aligned}
 \mathcal A_{\rm sp}(\operatorname{pr}_{X,{\rm sp}}\widehat x)
  ={}&(\widehat X_\infty,\Gamma_9,\widehat R,
      \operatorname{Carr}_R,S_0,M,\mathcal O_9,
      \Xi,P,\mathcal D_R^{\rm cof}),\\
 \mathfrak M_{\rm sp}={}&\operatorname{Im}\mathcal A_{\rm sp},\\
 \mathcal P_{\rm sp}(\mathfrak M_{\rm sp})
  ={}&(B_{\rm sp}^{\rm str},\Xi^*(I-P)\Xi,
       \mathscr I_{\rm GW},\Delta_{\rm bind},
       \text{familia cofinal primo--gamma--polar}),\\
 \Delta_{\rm bind}:={}&B_{\rm sp}^{\rm str}-\mathscr I_{\rm GW}.
\end{aligned}
\end{equation}
La memoria finita del objeto de partida satisface
\begin{equation}\label{eq:pdf2-rh-stein-finito}
 S_0=\sum_{r=0}^{N-1}M^{*r}\mathcal O_9^*\mathcal O_9M^r
     +M^{*N}S_0M^N.
\end{equation}
El último sumando permanece visible hasta que se demuestra su extinción. Por
tanto, reemplazar \(\mathfrak G_{\rm sp}\) por la sombra publicada
\(B_{\rm sp}^{\rm str}\) cambia el objeto de partida; no es una abreviación
del mismo teorema. El residual \(\Delta_{\rm bind}\) se conserva hasta que la
coligación source--first demuestra materialmente su anulación en
\eqref{eq:pdf2-rh-binding-source-first}.

\section{Dominio cofinal y forma objetivo}

Para $R>0$ se fija
\[
 \mathcal D_R=C_c^\infty((-R/2,R/2)),\qquad
 \widehat\psi(t)=\int_{\mathbb R}\psi(x)e^{-itx}\,dx,
\]
y, para $\psi,\phi\in\mathcal D_R$,
\[
 h_{\psi,\phi}(z)=\widehat\psi(z)
 \overline{\widehat\phi(\overline z)},\qquad
 \check h_{\psi,\phi}(a)=\langle\psi,T_a\phi\rangle.
\]
La unión $\bigcup_R\mathcal D_R=C_c^\infty(\mathbb R)$ es exacta. Para todo
conjunto finito de puntos complejos distintos, las evaluaciones
$\psi\mapsto(\widehat\psi(z_1),\ldots,\widehat\psi(z_n))$ son conjuntamente
sobreyectivas cuando $R$ es suficientemente grande. Por tanto, esta clase cofinal
separa las evaluaciones requeridas por el criterio de Weil.

\begin{lemma}[Prueba de la separación conjunta]
\label{lem:pdf2-rh-separacion}
La aplicación
\[
 \mathcal D_R\longrightarrow\mathbb C^n,
 \qquad
 \psi\longmapsto
 (\widehat\psi(z_1),\ldots,\widehat\psi(z_n))
\]
es sobreyectiva para todo conjunto de puntos complejos distintos
\(z_1,\ldots,z_n\), una vez que \(R\) contiene un intervalo no vacío.
\end{lemma}

\begin{proof}
Si los funcionales de evaluación fueran linealmente dependientes, existirían
\(c_1,\ldots,c_n\), no todos nulos, tales que
\[
 \int_{\mathbb R}\psi(x)
 \left(\sum_{j=1}^nc_je^{-iz_jx}\right)\,dx=0
 \qquad(\psi\in\mathcal D_R).
\]
La función analítica \(\sum_jc_je^{-iz_jx}\) se anularía en
\((-R/2,R/2)\), luego en todo el plano. La independencia de exponenciales
con exponentes distintos fuerza \(c_j=0\) para todo \(j\), contradicción.
Los funcionales son independientes; como el codominio es finito-dimensional,
la aplicación tiene rango \(\mathbb C^n\).
\end{proof}

\section{Las dos columnas y los tres canales calculados antes del signo}

Sea $\vartheta_L$ un regulador común. Con
\[
 \ell_{p,k}=k\log p,\qquad w_{p,k}=(\log p)p^{-k/2},\qquad
 \mu_\Gamma(a)=\frac{e^{-a/2}}{1-e^{-2a}},
\]
se usan la inclusión uniforme $j_m$ y el canal de costura
$\widehat\delta_{a,m}$, que satisfacen
\[
 j_m^*j_m=I,\qquad
 \widehat\delta_{a,m}^*\widehat\delta_{b,m}
 =(I-T_a)^*(I-T_b).
\]
Las columnas reguladas son
\begin{align*}
 J_{+,m,R,L}\psi={}&
 \Bigl(
  (\sqrt{\vartheta_L(\ell_{p,k})w_{p,k}}\,
       \widehat\delta_{\ell_{p,k},m}\psi)_{\ell_{p,k}\leq R},\\
 &\hspace{9mm}a\mapsto
  \sqrt{\vartheta_L(a)\mu_\Gamma(a)}\,
       \widehat\delta_{a,m}\psi,\ \zeta_+b_+(\psi)
 \Bigr),\\
 J_{-,m,R,L}\psi={}&
 \Bigl(
  (\sqrt{2\vartheta_L(\ell_{p,k})w_{p,k}}\,j_m\psi)_{\ell_{p,k}\leq R},
  \sqrt{-c_\Gamma}\,j_m\psi,\ \zeta_-b_-(\psi)
 \Bigr),
\end{align*}
donde $c_\Gamma=\psi_\Gamma(1/4)-\log\pi<0$ y
$b_\pm=(\widehat\psi(i/2)\pm\widehat\psi(-i/2))/\sqrt2$.

Una expansión término a término, sin usar ceros ni el signo buscado, da
\begin{align*}
 &\langle J_{+,m,R,L}\psi,J_{+,m,R,L}\phi\rangle
 -\langle J_{-,m,R,L}\psi,J_{-,m,R,L}\phi\rangle
 =\mathscr I_{{\rm GW},L}(\psi,\phi),\\
 \mathscr I_{{\rm GW},L}(\psi,\phi)
 ={}&h_{\psi,\phi}(i/2)+h_{\psi,\phi}(-i/2)
 +c_\Gamma\check h_{\psi,\phi}(0)\\
 &+\int_0^\infty\vartheta_L(a)\mu_\Gamma(a)
 [2\check h(0)-\check h(a)-\check h(-a)]\,da\\
 &-\sum_{p,k}\vartheta_L(\ell_{p,k})w_{p,k}
 [\check h(\ell_{p,k})+\check h(-\ell_{p,k})].
\end{align*}
El truncamiento de ambas columnas en \(\ell_{p,k}\leq R\) es obligatorio, no
notacional. Si \(\psi,\phi\in\mathcal D_R\), entonces
\(\check h_{\psi,\phi}(\pm\ell)=0\) para \(\ell>R\); la contribución de cada
par de coordenadas omitido a la \emph{diferencia} de Grams es por tanto cero.
La columna negativa no truncada, en cambio, tendría norma infinita porque
\(\sum_p(\log p)p^{-1/2}=\infty\). En el conjunto finito
\(\ell_{p,k}\leq R\) las coordenadas primas convergen individualmente cuando
\(L\to\infty\). En el canal gamma, el integrando es \(O(a)\) en cero y
\(\mu_\Gamma(a)=O(e^{-a/2})\) en infinito, de modo que la convergencia también
es en norma. Definimos, por tanto,
\begin{equation}\label{eq:pdf2-rh-columnas-reducidas}
 J_{\pm,m,R}:=\lim_{L\to\infty}J_{\pm,m,R,L}
 \quad\text{en sus Hilbert de coordenadas truncadas},
\end{equation}
y suprimimos \(m\) cuando permanece fijo. Por convergencia dominada,
\[
 \mathscr I_{{\rm GW},L}\longrightarrow\mathscr I_{\rm GW}
 \qquad(L\to\infty),
\]
y las identidades son naturales bajo $m\mapsto m+1$ y $R\le R'$.

Las columnas reducidas \(J_{+,R}\) y \(J_{-,R}\) son ahora operadores
materiales de \(\mathcal D_R\) a Hilbert. La expansión explícita prueba, sin
consultar ceros ni signo,
\begin{equation}\label{eq:pdf2-rh-gram-difference-gw}
 \langle J_{+,R}f,J_{+,R}g\rangle
 -\langle J_{-,R}f,J_{-,R}g\rangle
 =\mathscr I_{\rm pr}(f,g)+\mathscr I_\Gamma(f,g)
  +\mathscr I_{\rm pol}(f,g)
 =\mathscr I_{\rm GW}(f,g).
\end{equation}
Los relojes \(\ell_{p,k}\) producen el sumando primo; la integral con
\(\mu_\Gamma\) y \(c_\Gamma\check h(0)\), el sumando gamma; y
\[
 \langle b_+(f),b_+(g)\rangle-
 \langle b_-(f),b_-(g)\rangle
 =h_{f,g}(i/2)+h_{f,g}(-i/2)
\]
produce el sumando polar. Esta parte del cálculo es material e independiente
del puente estructural.

\section{Coligación source--first de la imagen espectral--polar}

Fijamos primero un nivel nonádico \(m\); este índice se suprime en
\(J_{\pm,R}\), \(E_R\) y \(\Sigma_R\) hasta construir el codificador. La
descomposición de los codominios explícitos de
\eqref{eq:pdf2-rh-columnas-reducidas} es
\[
 I_+=\{\mathrm{pr},\Gamma,\mathrm{pol}\},\qquad
 I_-(R)=\{\Gamma0,\mathrm{pol}\}
 \sqcup\{(p,k):\ell_{p,k}\leq R\},
\]
\[
 \mathscr H_{+,R}^{\rm amb}=\bigoplus_{i\in I_+}H_{+,i,R},
 \qquad
 \mathscr H_{-,R}^{\rm amb}=\bigoplus_{j\in I_-(R)}H_{-,j,R}.
\]
No identificamos un subespacio alcanzable con una suma de sus proyecciones
coordenadas. Definimos, con las inclusiones ambientales entendidas,
\[
 H_{+,R}:=\overline{\operatorname{Ran}J_{+,R}},\qquad
 H_{-,R}:=\overline{\operatorname{Ran}J_{-,R}},
\]
como subespacios cerrados de los dos Hilbert ambientales. La estructura
\(\mathfrak G_{\rm sp}\) contiene, como parte de su dato matemático y
antes de la publicación clásica, un espacio de estado antiespecular
\(\mathcal H_R^{\rm sh}\) y la fila
isométrica source--first
\begin{equation}\label{eq:pdf2-rh-coligacion-source-first}
 \Sigma_R=
 \begin{pmatrix}\mathcal K_R^{\rm src}\\ \mathcal F_R^{\rm src}\end{pmatrix}:
 H_{+,R}\longrightarrow
 H_{-,R}\oplus\mathcal H_R^{\rm sh},
 \qquad
 \Sigma_R^*\Sigma_R
 = (\mathcal K_R^{\rm src})^*\mathcal K_R^{\rm src}
  +(\mathcal F_R^{\rm src})^*\mathcal F_R^{\rm src}=I_{H_{+,R}}.
\end{equation}
Su acción alcanzable se fija sobre la columna positiva completa por
las relaciones de las dos publicaciones que forman parte de
\(\operatorname{Rel}_{\rm sp}\). Para evitar cualquier sustitución de lector,
denotamos por
\((\mathcal X_R^{\rm str},\Xi_R^{\rm str},P_R^{\rm str})\) la restricción a
\(\mathcal D_R\) del lector estructural de
\eqref{eq:pdf2-rh-res-sp-completa}. Sus dos momentos son, literalmente,
\begin{equation}\label{eq:pdf2-rh-momentos-lector-estructural}
\begin{aligned}
 \langle\Xi_R^{\rm str}f,\Xi_R^{\rm str}g\rangle
  &=\langle J_{+,R}f,J_{+,R}g\rangle,\\
 \langle P_R^{\rm str}\Xi_R^{\rm str}f,
          P_R^{\rm str}\Xi_R^{\rm str}g\rangle
  &=\langle J_{-,R}f,J_{-,R}g\rangle,
 \qquad (P_R^{\rm str})^*= (P_R^{\rm str})^2=P_R^{\rm str}.
\end{aligned}
\end{equation}
La coordenada de frontera viene con la isometría estructural
\[
 \jmath_R^{\rm sh}:
 \overline{\operatorname{Ran}((I-P_R^{\rm str})\Xi_R^{\rm str})}
 \longrightarrow\mathcal H_R^{\rm sh}.
\]
Por definición de la fila source--first, no por una identificación posterior,
\begin{equation}\label{eq:pdf2-rh-source-state-output}
 E_R:=\mathcal F_R^{\rm src}J_{+,R}
     =\jmath_R^{\rm sh}(I-P_R^{\rm str})\Xi_R^{\rm str}:
 \mathcal D_R\longrightarrow\mathcal H_R^{\rm sh},
 \qquad
 \Omega_{0,R}:=\mathcal K_R^{\rm src}J_{+,R}-J_{-,R}.
\end{equation}
No hay una coligación por canal: la misma fila recibe simultáneamente los
relojes primo--potencia, la integral gamma, el modo escalar y las dos líneas
polares. En particular, conserva todos los productos cruzados que aparecen al
polarizar \eqref{eq:pdf2-rh-gram-difference-gw}.
Con las proyecciones ambientales \(\pi_j^-:\mathscr H_{-,R}^{\rm amb}
\to H_{-,j,R}\), las filas de salida correctamente tipadas son
\begin{equation}\label{eq:pdf2-rh-bloques-coligacion}
 \mathcal K_{j,R}^{\rm src}
 :=\pi_j^-\mathcal K_R^{\rm src}:
 H_{+,R}\longrightarrow H_{-,j,R},
 \qquad
 \mathcal F_R^{\rm src}:H_{+,R}\longrightarrow\mathcal H_R^{\rm sh}.
\end{equation}
Así no se aplica \(\mathcal K_R^{\rm src}\) a una coordenada positiva que
pudiera no pertenecer por sí sola a \(H_{+,R}\). Todos los cruces entre
canales permanecen dentro de
\((\mathcal K_R^{\rm src})^*\mathcal K_R^{\rm src}\) y
\((\mathcal F_R^{\rm src})^*\mathcal F_R^{\rm src}\); no se ortogonalizan
antes de la identidad de energía.

La coordenada de supervivencia de
\(\mathfrak G_{\rm sp}\) proporciona espacios residuales
\(\mathcal R_{n,R}^{\rm sh}\), continuaciones
\(\Omega_{n,R}:\mathcal D_R\to\mathcal R_{n,R}^{\rm sh}\) e isometrías
antiespeculares
\(V_{n,R}:\mathcal R_{n+1,R}^{\rm sh}\to\mathcal R_{n,R}^{\rm sh}\). Su
primer espacio y su primer operador son, por definición tipada,
\[
 \mathcal R_{0,R}^{\rm sh}:=H_{-,R},\qquad
 \Omega_{0,R}:=\mathcal K_R^{\rm src}J_{+,R}-J_{-,R}.
\]
Su
recurrencia de una vuelta completa es
\begin{equation}\label{eq:pdf2-rh-recurrencia-residual}
 q=\frac59,\qquad
 \Omega_{n,R}=qV_{n,R}\Omega_{n+1,R},
 \qquad V_{n,R}^*V_{n,R}=I,
 \qquad
 \sup_{n\geq0}\|\Omega_{n,R}f\|<\infty
 \quad(f\in\mathcal D_R).
\end{equation}
Esta recurrencia pertenece a la coligación enriquecida: transporta la
diferencia de salida, no la forma de Weil ni su signo.

\begin{proposition}[Anulación coinductiva de la diferencia de salida]
\label{prop:pdf2-rh-omega-cero}
Para todo \(R>0\) y todo \(f\in\mathcal D_R\),
\begin{equation}\label{eq:pdf2-rh-output-binding-source}
 \Omega_{n,R}f=0\quad(n\geq0),
 \qquad
 \boxed{\mathcal K_R^{\rm src}J_{+,R}=J_{-,R}}.
\end{equation}
\end{proposition}

\begin{proof}
Iterar \eqref{eq:pdf2-rh-recurrencia-residual} \(N\) veces y usar que cada
\(V_{n,R}\) es isométrica da
\[
 \|\Omega_{n,R}f\|
 =q^N\|\Omega_{n+N,R}f\|
 \leq q^N\sup_{j\geq n}\|\Omega_{j,R}f\|.
\]
El miembro derecho converge a cero porque \(0<q=5/9<1\). Por tanto
\(\Omega_{n,R}f=0\); el caso \(n=0\) y
\eqref{eq:pdf2-rh-source-state-output} producen la identidad enmarcada.
\end{proof}

La acción fila a fila de la salida ya no queda implícita. Las proyecciones son
las restricciones a \(H_{-,R}\) de las proyecciones del Hilbert ambiental, y
\begin{equation}\label{eq:pdf2-rh-salidas-generador-a-generador}
\begin{aligned}
 \mathcal K_{\Gamma0,R}^{\rm src}J_{+,R}f
  &=\sqrt{-c_\Gamma}\,j_mf,\\
 \mathcal K_{(p,k),R}^{\rm src}J_{+,R}f
  &=\sqrt{2w_{p,k}}\,j_mf
  &&(\ell_{p,k}\leq R),\\
 \mathcal K_{{\rm pol},R}^{\rm src}J_{+,R}f
  &=\zeta_-b_-(f).
\end{aligned}
\end{equation}
Estas filas agotan \(J_{-,R}\), mientras
\begin{equation}\label{eq:pdf2-rh-frontera-generador-a-generador}
 \mathcal F_R^{\rm src}J_{+,R}f=E_Rf
\end{equation}
conserva, con todos sus cruces, la coordenada de frontera.
Evaluar \eqref{eq:pdf2-rh-coligacion-source-first} en \(J_{+,R}f\) y
\(J_{+,R}g\), y usar \eqref{eq:pdf2-rh-output-binding-source}, produce la
identidad de energía anterior al signo
\begin{equation}\label{eq:pdf2-rh-energia-source-first}
 \boxed{
 \langle J_{+,R}f,J_{+,R}g\rangle
 =\langle J_{-,R}f,J_{-,R}g\rangle
  +\langle E_Rf,E_Rg\rangle.}
\end{equation}
Al compararla con el cálculo independiente
\eqref{eq:pdf2-rh-gram-difference-gw} se obtiene, sin introducir una raíz de
la forma objetivo,
\begin{equation}\label{eq:pdf2-rh-binding-source-first}
 \boxed{E_R^*E_R=\mathscr I_{\rm GW}\big|_{\mathcal D_R}.}
\end{equation}
En particular,
\(\Delta_{{\rm bind},R}
:=E_R^*E_R-\mathscr I_{\rm GW}|_{\mathcal D_R}=0\); esta anulación es la
salida de \eqref{eq:pdf2-rh-coligacion-source-first}--
\eqref{eq:pdf2-rh-energia-source-first}, no una premisa insertada en el
publicador.

\section{Complemento estructural, hoja espejo y telescopía terminal}

El estado \(E_R=\mathcal F_R^{\rm src}J_{+,R}\) construido por la fila
source--first es la coordenada de frontera de
\(\mathfrak G_{\rm sp}\). Definimos su forma estructural por
\begin{equation}\label{eq:pdf2-rh-bw-estructural}
 \mathcal B_{{\rm sp},R}^{\rm str}
 :=E_R^*E_R
 =\mathscr I_{\rm GW}\big|_{\mathcal D_R},
\end{equation}
donde la segunda igualdad es
\eqref{eq:pdf2-rh-binding-source-first}, ya deducida de la coligación y de
la expansión generador a generador.

Sean \(P_+^{\rm sh}\) y \(P_-^{\rm sh}\) los proyectores ortogonales
complementarios de las dos hojas. Se fijan
\begin{equation}\label{eq:pdf2-rh-q-a-d}
 q=\frac59,\qquad
 A=P_+^{\rm sh}+qP_-^{\rm sh},\qquad
 D=\sqrt{1-q^2}\,P_-^{\rm sh}.
\end{equation}
La ortogonalidad da la identidad exacta
\begin{equation}\label{eq:pdf2-rh-conservacion-local}
 A^*A+D^*D
 =P_+^{\rm sh}+q^2P_-^{\rm sh}
 +(1-q^2)P_-^{\rm sh}=I.
\end{equation}
Las nueve fases coordinan una sola vuelta con memoria:
\begin{equation}\label{eq:pdf2-rh-monodromia}
 M_9=\Phi A_8\cdots A_0=qV_9,
 \qquad V_9^*V_9=I.
\end{equation}
La contracción por vuelta es \(q\), no \(q^9\); las fases internas conservan
el orden y el acarreo.

La coligación tipa
\(E_R:\mathcal D_R\to\operatorname{Ran}P_-^{\rm sh}\). Por tanto,
\begin{equation}\label{eq:pdf2-rh-binding-espejo}
 P_+^{\rm sh}E_R=0,
 \qquad AE_R=qE_R,
 \qquad
 E_R^*E_R=\mathcal B_{{\rm sp},R}^{\rm str}
 =\mathscr I_{\rm GW}\big|_{\mathcal D_R}.
\end{equation}

\begin{proposition}[Telescopía finita con terminal conservado]
\label{prop:pdf2-rh-telescopia}
Para cada \(N\geq1\),
\begin{equation}\label{eq:pdf2-rh-telescopia-finita}
 \mathcal B_{{\rm sp},R}^{\rm str}
 =\sum_{n=0}^{N-1}(DA^nE_R)^*(DA^nE_R)
 +(A^NE_R)^*(A^NE_R).
\end{equation}
Equivalentemente, evaluada en \(f,g\in\mathcal D_R\), la identidad finita de
la coligación es
\begin{equation}\label{eq:pdf2-rh-balance-finito-columnas}
\boxed{
\begin{aligned}
 &\langle J_{+,R}f,J_{+,R}g\rangle
  -\langle J_{-,R}f,J_{-,R}g\rangle\\
 &\quad=\sum_{n=0}^{N-1}
  \langle DA^nE_Rf,DA^nE_Rg\rangle
  +\langle A^NE_Rf,A^NE_Rg\rangle.
\end{aligned}}
\end{equation}
El último sumando satisface
\begin{equation}\label{eq:pdf2-rh-terminal-extincion}
 (A^NE_R)^*(A^NE_R)
 =q^{2N}\mathcal B_{{\rm sp},R}^{\rm str}
 \xrightarrow[N\to\infty]{}0.
\end{equation}
\end{proposition}

\begin{proof}
Aplicando \eqref{eq:pdf2-rh-conservacion-local} a \(A^nE_R\),
\[
 (A^nE_R)^*(A^nE_R)
 =(DA^nE_R)^*(DA^nE_R)
 +(A^{n+1}E_R)^*(A^{n+1}E_R).
\]
La suma desde \(n=0\) hasta \(N-1\) cancela sólo los términos intermedios y
conserva \((A^NE_R)^*(A^NE_R)\). Por
\eqref{eq:pdf2-rh-binding-espejo}, \(A^NE_R=q^NE_R\); por tanto el terminal
es \(q^{2N}E_R^*E_R=q^{2N}\mathcal B_{{\rm sp},R}^{\rm str}\). Como
\(0<q=5/9<1\), converge geométricamente a cero.
\end{proof}

Sólo después de demostrar \eqref{eq:pdf2-rh-terminal-extincion} se define
\begin{equation}\label{eq:pdf2-rh-l-r}
 \mathscr L_Rf=(DA^nE_Rf)_{n\geq0}.
\end{equation}
El límite de la telescopía finita da
\begin{equation}\label{eq:pdf2-rh-lstar-l}
 \boxed{
 \mathscr L_R^*\mathscr L_R
 =\mathcal B_{{\rm sp},R}^{\rm str}
 =\mathscr I_{\rm GW}\big|_{\mathcal D_R}\succeq0.}
\end{equation}

\section{Codificador común, proyección y doble publicación explícitas}

Fijemos \(m\geq1\) y dos vectores ortonormales
\(u_m,\eta_m\in\mathbb C^{9^m}\), donde \(u_m\) es el modo uniforme y
\(\eta_m\) el modo de acarreo. El espacio y la proyección del codificador son
\begin{equation}\label{eq:pdf2-rh-hilbert-comun}
 \mathcal X_{m,R}^{\rm sp}
 :=\mathbb C^{9^m}\widehat\otimes
   (H_{-,R}\oplus\mathcal H_R^{\rm sh}),
 \qquad
 P_{m,R}:=|u_m\rangle\langle u_m|\otimes
 I_{H_{-,R}\oplus\mathcal H_R^{\rm sh}},
 \qquad Q_{m,R}:=I-P_{m,R}.
\end{equation}
La acción completa del codificador de los cinco canales es
\begin{equation}\label{eq:pdf2-rh-xi-comun-explicita}
 \boxed{
 \begin{aligned}
  \mathfrak C_{m,R}f
  &:=u_m\otimes(J_{-,R}f,0)
    +\eta_m\otimes(0,E_Rf),\\
  \Xi_{m,R}f&:=\mathfrak C_{m,R}f.
 \end{aligned}}
\end{equation}
No se toma una raíz de \(\mathscr I_{\rm GW}\): las dos entradas son la
salida y el estado de la coligación source--first ya definidos en
\eqref{eq:pdf2-rh-source-state-output}.

\begin{proposition}[Los dos momentos se obtienen de la acción del codificador]
\label{prop:pdf2-rh-vmas-isometria}
Para todo \(f,g\in\mathcal D_R\),
\begin{equation}\label{eq:pdf2-rh-dos-momentos-construidos}
\begin{aligned}
 \langle\mathfrak C_{m,R}f,\mathfrak C_{m,R}g\rangle
  &=\langle J_{+,R}f,J_{+,R}g\rangle,\\
 \langle P_{m,R}\mathfrak C_{m,R}f,
         P_{m,R}\mathfrak C_{m,R}g\rangle
  &=\langle J_{-,R}f,J_{-,R}g\rangle,\\
 \langle Q_{m,R}\Xi_{m,R}f,Q_{m,R}\Xi_{m,R}g\rangle
  &=\mathscr I_{\rm GW}(f,g).
\end{aligned}
\end{equation}
\end{proposition}

\begin{proof}
La ortogonalidad \(u_m\perp\eta_m\) y
\eqref{eq:pdf2-rh-energia-source-first} dan
\[
 \langle\mathfrak C_{m,R}f,\mathfrak C_{m,R}g\rangle
 =\langle J_{-,R}f,J_{-,R}g\rangle
  +\langle E_Rf,E_Rg\rangle
 =\langle J_{+,R}f,J_{+,R}g\rangle.
\]
La proyección \(P_{m,R}\) conserva el sumando uniforme y elimina el de
acarreo; \(Q_{m,R}\) hace lo contrario. Las otras dos identidades siguen de
\eqref{eq:pdf2-rh-binding-source-first}.
\end{proof}

El codificador recién escrito es un representante explícito del lector
estructural inicial, no un lector sustituto. En los subespacios alcanzables se
define
\begin{equation}\label{eq:pdf2-rh-isomorfismo-lector-estructural}
\begin{aligned}
 \mathcal U_{m,R}^{\rm str}
  (P_R^{\rm str}\Xi_R^{\rm str}f)
   &:=u_m\otimes(J_{-,R}f,0),\\
 \mathcal U_{m,R}^{\rm str}
  ((I-P_R^{\rm str})\Xi_R^{\rm str}f)
   &:=\eta_m\otimes(0,E_Rf).
\end{aligned}
\end{equation}
Las dos fuentes son ortogonales. La primera igualdad de Gram de cada fuente
es \eqref{eq:pdf2-rh-momentos-lector-estructural}; para la segunda se usa
\(E_R=\jmath_R^{\rm sh}(I-P_R^{\rm str})\Xi_R^{\rm str}\). Por ello las dos
reglas están bien definidas, son isométricas y se extienden a la suma ortogonal
de sus cierres. En ese cierre alcanzable satisfacen
\begin{equation}\label{eq:pdf2-rh-entrelazamiento-lector-estructural}
 \mathcal U_{m,R}^{\rm str}\Xi_R^{\rm str}=\Xi_{m,R},
 \qquad
 \mathcal U_{m,R}^{\rm str}P_R^{\rm str}
 =P_{m,R}\mathcal U_{m,R}^{\rm str},
 \qquad
 (\Xi_R^{\rm str})^*(I-P_R^{\rm str})\Xi_R^{\rm str}=E_R^*E_R.
\end{equation}
Así, \(B_{{\rm sp},R}^{\rm str}\), el complemento de la estructura de
partida, es exactamente el complemento del codificador construido; no se ha
cambiado de objeto al pasar de \eqref{eq:pdf2-rh-res-sp-completa} a
\eqref{eq:pdf2-rh-xi-comun-explicita}.

La regla
\begin{equation}\label{eq:pdf2-rh-vmas-rango}
 V_{+,m,R}^{0}(J_{+,R}f):=\Xi_{m,R}f
\end{equation}
está bien definida y es isométrica por la primera igualdad de
\eqref{eq:pdf2-rh-dos-momentos-construidos}; se extiende de manera única a
una isometría
\(V_{+,m,R}:H_{+,R}\to\mathcal X_{m,R}^{\rm sp}\). La publicación negativa
es la isometría
\begin{equation}\label{eq:pdf2-rh-vmenos}
 V_{-,m,R}:H_{-,R}\longrightarrow\mathcal X_{m,R}^{\rm sp},
 \qquad V_{-,m,R}y=u_m\otimes(y,0).
\end{equation}
Por construcción,
\begin{equation}\label{eq:pdf2-rh-doble-publicacion-material}
 \boxed{
 V_{+,m,R}J_{+,R}=\Xi_{m,R},
 \qquad
 V_{-,m,R}J_{-,R}=P_{m,R}\Xi_{m,R}.}
\end{equation}
El complemento conserva exactamente el estado antiespecular y su columna de
observación:
\begin{equation}\label{eq:pdf2-rh-complemento-comun}
 \boxed{
 \Xi_{m,R}^*Q_{m,R}\Xi_{m,R}
 =E_R^*E_R
 =\mathscr L_R^*\mathscr L_R
 =\mathscr I_{\rm GW}\big|_{\mathcal D_R}.}
\end{equation}

\begin{proposition}[Conector inducido y coincidencia source--first]
\label{prop:pdf2-rh-conector}
El operador
\begin{equation}\label{eq:pdf2-rh-k-r}
 \mathcal K_R
 :=V_{-,m,R}^*P_{m,R}V_{+,m,R}:
 H_{+,R}\longrightarrow H_{-,R}
\end{equation}
es contractivo, coincide con \(\mathcal K_R^{\rm src}\) en todo
\(H_{+,R}\) y satisface
\begin{equation}\label{eq:pdf2-rh-k-binding}
 \|\mathcal K_R\|\leq1,
 \qquad
 \mathcal K_RJ_{+,R}=J_{-,R}.
\end{equation}
\end{proposition}

\begin{proof}
Las dos cartas son isometrías y \(P_{m,R}\) es una proyección ortogonal,
luego \(\|\mathcal K_R\|\leq1\). Sobre el rango denso de \(J_{+,R}\),
\[
 \mathcal K_RJ_{+,R}f
 =V_{-,m,R}^*P_{m,R}\Xi_{m,R}f
 =V_{-,m,R}^*V_{-,m,R}J_{-,R}f
 =J_{-,R}f.
\]
La misma identidad vale para \(\mathcal K_R^{\rm src}\) por
\eqref{eq:pdf2-rh-output-binding-source}; la continuidad y la densidad
prueban que ambos operadores coinciden.
\end{proof}

La diferencia de Gram queda así publicada en todas sus realizaciones
coincidentes:
\begin{equation}\label{eq:pdf2-rh-tres-publicaciones}
\boxed{
\begin{aligned}
 J_{+,R}^*J_{+,R}-J_{-,R}^*J_{-,R}
 &=\Xi_{m,R}^*Q_{m,R}\Xi_{m,R}\\
 &=E_R^*E_R\\
 &=\mathscr L_R^*\mathscr L_R\\
 &=\mathcal B_{{\rm sp},R}^{\rm str}=\mathscr I_{\rm GW}.
\end{aligned}}
\end{equation}

\section{Naturalidad cofinal y criterio reversible de Weil}

Para \(R\leq R'\), sea
\(i_{R,R'}:\mathcal D_R\hookrightarrow\mathcal D_{R'}\). Las columnas
individuales no forman un sistema isométrico: al aparecer una longitud nueva
\(R<\ell_{p,k}\leq R'\), las dos adquieren la misma energía positiva. Para
\(f,g\in\mathcal D_R\), la separación de soportes da exactamente
\begin{equation}\label{eq:pdf2-rh-naturalidad}
 \left\langle(I-T_{\ell_{p,k}})f,
                    (I-T_{\ell_{p,k}})g\right\rangle
 =2\langle f,g\rangle
 =\left\langle\sqrt2j_mf,\sqrt2j_mg\right\rangle .
\end{equation}
Por tanto el incremento primo de
\(J_{+,R'}^*J_{+,R'}\) coincide con el de
\(J_{-,R'}^*J_{-,R'}\) y se cancela en la diferencia, pero no se declara una
isometría falsa entre las columnas. La compatibilidad cofinal correcta es la
de las formas:
\begin{equation}\label{eq:pdf2-rh-naturalidad-p}
 \boxed{
 \mathscr I_{{\rm GW},R}
 =i_{R,R'}^*\mathscr I_{{\rm GW},R'}i_{R,R'},
 \qquad
 \mathcal B_{{\rm sp},R}^{\rm str}
 =i_{R,R'}^*\mathcal B_{{\rm sp},R'}^{\rm str}i_{R,R'}.}
\end{equation}
Como \(E_R^*E_R=\mathscr L_R^*\mathscr L_R=\mathscr I_{{\rm GW},R}\), las
realizaciones mínimas de Kolmogórov sí admiten isometrías únicas sobre sus
rangos alcanzables, definidas por
\begin{equation}\label{eq:pdf2-rh-naturalidad-bw}
 \iota_{R,R'}^E(E_Rf):=E_{R'}f,
 \qquad
 \iota_{R,R'}^L(\mathscr L_Rf):=\mathscr L_{R'}f
 \quad(f\in\mathcal D_R).
\end{equation}
Estas dos reglas están bien definidas precisamente por
\eqref{eq:pdf2-rh-naturalidad-p}; no se extienden por decreto a
\(J_{\pm,R}\), \(\Xi_{m,R}\) o \(\mathcal K_R\). El refinamiento nonádico
\(m\mapsto m+1\) permanece, a ventana fija, en las isometrías estructurales
que transportan \(j_m\), \(\widehat\delta_{a,m}\), \(u_m\) y \(\eta_m\).

Para no comprimir el último paso, fijamos también sus objetos y
normalizaciones. Sea
\[
 \xi(s):=\frac12s(s-1)\pi^{-s/2}\Gamma(s/2)\zeta(s)
\]
y sea \(\mathscr Z\) el multiconjunto de sus ceros, cada cero repetido según
su multiplicidad. La ecuación funcional y la realidad de los coeficientes
preservan \(\mathscr Z\) por
\(\rho\mapsto1-\rho\), \(\rho\mapsto\overline\rho\) y
\(\rho\mapsto1-\overline\rho\). Con nuestra convención de Fourier se define
\begin{equation}\label{eq:pdf2-rh-transformada-weil}
 \Phi_f(s):=\widehat f\!\left(i(s-\tfrac12)\right),
 \qquad
 \widetilde f(x):=\overline{f(-x)}.
\end{equation}
La fórmula explícita completa calculada en
\eqref{eq:pdf2-rh-gram-difference-gw}, incluidos los términos gamma y polar
que retiran polos y ceros triviales, tiene la forma espectral
\begin{equation}\label{eq:pdf2-rh-formula-espectral-weil}
 \mathscr I_{\rm GW}(f,g)
 =\sum_{\rho\in\mathscr Z}
   \Phi_f(\rho)\,
   \overline{\Phi_g(1-\overline\rho)}.
\end{equation}
La suma cuenta multiplicidades. Converge absolutamente: para
\(f,g\in C_c^\infty(\mathbb R)\), Paley--Wiener da decaimiento más rápido
que cualquier potencia en bandas verticales, mientras el número de ceros con
\(|\operatorname{Im}\rho|\leq T\) es \(O(T\log T)\).

\begin{proposition}[Criterio reversible de Weil con cuantificadores completos]
\label{prop:pdf2-rh-criterio-weil-desplegado}
Bajo la normalización
\eqref{eq:pdf2-rh-transformada-weil} y contando todos los ceros no triviales
con su multiplicidad,
\begin{equation}\label{eq:pdf2-rh-criterio-weil}
 \left[
  \mathscr I_{\rm GW}(f,f)\geq0
  \text{ para todo }f\in C_c^\infty(\mathbb R;\mathbb C)
 \right]
 \Longleftrightarrow
 \left[
  \operatorname{Re}\rho=\frac12
  \text{ para todo }\rho\in\mathscr Z
 \right].
\end{equation}
\end{proposition}

La implicación recíproca no se obtiene interpolando un número creciente de
ceros: esa operación no daría control uniforme sobre la cola. Se usa el
criterio clásico completo de Weil en la formulación de Bombieri, cuyo espacio
de prueba y cuya normalización se trasladan aquí sin abreviarlos.\footnote{%
E.~Bombieri, \emph{Remarks on Weil's quadratic functional in the theory of
prime numbers, I}, Rendiconti Lincei, Matematica e Applicazioni, ser.~9,
vol.~11 (2000), pp.~183--233. El criterio allí fijado afirma que el funcional
cuadrático de la fórmula explícita es semidefinido positivo sobre toda
\(C_c^\infty(0,\infty)\) si y sólo si vale la hipótesis de Riemann.}

\begin{proof}
Escribimos materialmente el cambio de coordenadas que identifica ambos
criterios. Para \(f\in C_c^\infty(\mathbb R)\), sea
\begin{equation}\label{eq:pdf2-rh-cambio-mellin}
 (\mathcal M f)(x):=x^{-1/2}f(\log x),
 \qquad x>0.
\end{equation}
El mapa
\[
 \mathcal M:C_c^\infty(\mathbb R)
 \longrightarrow C_c^\infty(0,\infty)
\]
es biyectivo, con inversa
\begin{equation}\label{eq:pdf2-rh-cambio-mellin-inverso}
 (\mathcal M^{-1}g)(u)=e^{u/2}g(e^u).
\end{equation}
Si
\(\widehat g_{\rm Mel}(s):=\int_0^\infty g(x)x^{s-1}\,dx\), el cambio
\(x=e^u\) da, sin factor residual,
\begin{equation}\label{eq:pdf2-rh-mellin-fourier}
 \begin{aligned}
 \widehat{\mathcal M f}_{\rm Mel}(s)
 &=\int_{\mathbb R}f(u)e^{(s-1/2)u}\,du\\
 &=\widehat f\!\left(i(s-\tfrac12)\right)
 =\Phi_f(s).
 \end{aligned}
\end{equation}
Además,
\begin{equation}\label{eq:pdf2-rh-mellin-involucion}
 \widehat{\overline{\mathcal M f}}_{\rm Mel}(1-\rho)
 =\overline{
   \widehat{\mathcal M f}_{\rm Mel}(1-\overline\rho)}
 =\overline{\Phi_f(1-\overline\rho)}.
\end{equation}
Por tanto el funcional del criterio de Weil--Bombieri es exactamente
\begin{equation}\label{eq:pdf2-rh-bombieri-identificado}
 \sum_{\rho\in\mathscr Z}
 \widehat{\mathcal M f}_{\rm Mel}(\rho)\,
 \widehat{\overline{\mathcal M f}}_{\rm Mel}(1-\rho)
 =
 \sum_{\rho\in\mathscr Z}
 \Phi_f(\rho)\overline{\Phi_f(1-\overline\rho)}
 =\mathscr I_{\rm GW}(f,f).
\end{equation}
No se cambia la clase de pruebas, no se pierde ninguna multiplicidad y no se
trunca el conjunto de ceros, porque \(\mathcal M\) es una biyección. El
criterio clásico citado aplicado a
\eqref{eq:pdf2-rh-bombieri-identificado} da precisamente la equivalencia
\eqref{eq:pdf2-rh-criterio-weil}. En la dirección directa, si
\(\operatorname{Re}\rho=1/2\), la misma identidad se reduce explícitamente a
\[
 \mathscr I_{\rm GW}(f,f)
 =\sum_{\rho\in\mathscr Z}|\Phi_f(\rho)|^2\geq0.
\]
\end{proof}

La igualdad \(\bigcup_R\mathcal D_R=C_c^\infty(\mathbb R)\) es exacta; cada
función empleada en el argumento pertenece a alguna ventana y la naturalidad
cofinal transporta allí la positividad ya demostrada.

\begin{theorem}[Resolución Riemann--Weil desde la imagen espectral--polar]
\label{thm:pdf2-rh-resolucion}
Fijada la estructura completa
\(\mathfrak G_{\rm sp}\), la forma
completa de Guinand--Weil es positiva en la clase cofinal y separadora. En
consecuencia, todo cero no trivial \(\rho\) de la función zeta de Riemann
satisface
\[
 \boxed{\operatorname{Re}\rho=\frac12.}
\]
\end{theorem}

\begin{proof}
La expansión primo--gamma--polar identifica la diferencia de las dos columnas
con \(\mathscr I_{\rm GW}\) usando un único regulador y lo retira por
convergencia dominada. La recurrencia antiespecular
\eqref{eq:pdf2-rh-recurrencia-residual} anula la diferencia de salida antes de
formar una desigualdad; por tanto la isometría \(\Sigma_R\) da
\eqref{eq:pdf2-rh-energia-source-first}. Comparar esa identidad con la
expansión calculada produce
\(E_R^*E_R=\mathscr I_{\rm GW}\). La identidad local
\(A^*A+D^*D=I\) se itera sin borrar el terminal y da
\eqref{eq:pdf2-rh-telescopia-finita}; el terminal es exactamente
\(q^{2N}E_R^*E_R\) y se extingue sólo en el límite. De ahí
\(\mathscr L_R^*\mathscr L_R=\mathscr I_{\rm GW}\succeq0\).
La acción
\eqref{eq:pdf2-rh-xi-comun-explicita} construye entonces el codificador, sus
dos momentos, la proyección y las dos publicaciones, y
\eqref{eq:pdf2-rh-tres-publicaciones} identifica todas las realizaciones del
mismo cuadrado. La naturalidad transporta esa identidad a toda la unión
cofinal. Finalmente,
\eqref{eq:pdf2-rh-criterio-weil} concluye la afirmación.
\end{proof}

\section{Control Redheffer local posterior y no gobernante}

El siguiente despliegue verifica que la pasividad local puede obtenerse desde
el acarreo sin usar la forma de Weil. Es un control de no-circularidad y no una
premisa del \cref{thm:pdf2-rh-resolucion}. Esta transferencia truncada no es
la coligación completa \(\Sigma_R\) de
\eqref{eq:pdf2-rh-coligacion-source-first}.

Para \(N=9^m\), sea \(U_{a,N}\) el odómetro unitario
\[
 U_{a,N}(e_j\otimes f)=e_{j+1}\otimes f\ (j<N-1),
 \qquad U_{a,N}(e_{N-1}\otimes f)=e_0\otimes T_af.
\]
Con \(P_N=|u_N\rangle\langle u_N|\otimes I\), \(Q_N=I-P_N\),
\(A_{a,N}=P_NU_{a,N}P_N\) y \(C_{a,N}=Q_NU_{a,N}P_N\),
\begin{equation}\label{eq:pdf2-rh-source-energia}
 A_{a,N}^*A_{a,N}+C_{a,N}^*C_{a,N}=P_N,
 \qquad
 C_{a,N}^*C_{a,N}
 =\frac{N-1}{N^2}(2I-T_a-T_a^*).
\end{equation}
Para \(0<s<1\), la transferencia
\begin{equation}\label{eq:pdf2-rh-source-redheffer}
 \mathscr R_{a,N}(s)=(sI-A_{a,N})(I-sA_{a,N})^{-1}
\end{equation}
satisface
\begin{equation}\label{eq:pdf2-rh-source-defecto}
\begin{aligned}
 I-\mathscr R_{a,N}(s)^*\mathscr R_{a,N}(s)
 ={}&(1-s^2)(I-sA_{a,N}^*)^{-1}\\
 &\quad\cdot C_{a,N}^*C_{a,N}(I-sA_{a,N})^{-1}\succeq0.
\end{aligned}
\end{equation}
Con
\[
 s_{p,N}=\frac{Np^{-1/2}}{1+(N-1)p^{-1/2}},
 \qquad s_\Gamma(a)=e^{-a/2},
\]
la suma sobre relojes primos y la integral directa gamma conservan norma a lo
sumo uno. De forma explícita,
\begin{equation}\label{eq:pdf2-rh-source-relojes}
 \mathscr R_{N,R}
 =\bigoplus_{\log p\leq R}
   \mathscr R_{\log p,N}(s_{p,N})
 \ \oplus\!
 \int_0^\infty{}^\oplus
   \mathscr R_{a,N}(s_\Gamma(a))\,da,
 \qquad \|\mathscr R_{N,R}\|\leq1.
\end{equation}
La carta finita de roles
\[
 \mathcal H_+=\operatorname{span}\{f_3,f_6,f_9\},\quad
 \mathcal H_-=\operatorname{span}\{f_4,f_8\},\quad
 C_\angle=|f_4\rangle\langle f_3|+|f_8\rangle\langle f_9|
\]
satisface \(I-C_\angle^*C_\angle=P_6\). Con las cartas energéticas
two-way se fija
\begin{equation}\label{eq:pdf2-rh-source-cartas}
\begin{aligned}
 \mathcal K_q&=K_2\otimes(\mathbb C^9)^{\otimes q},\\
 \mathsf B_{\pm,q}:\mathcal M_{\pm,q,R}
 &\longrightarrow
 \mathcal K_q\widehat\otimes\mathcal H_\pm
 \widehat\otimes\mathcal H_{\mathfrak C_R},\\
 \mathsf B_{\pm,q}^*\mathsf B_{\pm,q}&=\mathsf G_{\pm,q},
 \qquad
 \mathsf B_{\pm,q}^{\sharp}
 =\mathsf G_{\pm,q}^{-1}\mathsf B_{\pm,q}^*.
\end{aligned}
\end{equation}
La transferencia
\begin{equation}\label{eq:pdf2-rh-source-transferencia}
 \mathscr C_{q,R}^{\rm src}
 =\mathsf B_{-,q}^{\sharp}
  (I\otimes C_\angle\otimes\mathscr R_{N,R})\mathsf B_{+,q}
\end{equation}
tiene defecto
\begin{align}
 \Delta_{q,R}^{\rm src}
 ={}&I\otimes\bigl[P_6\otimes I
 +(P_3+P_9)\otimes
 (I-\mathscr R_{N,R}^*\mathscr R_{N,R})\bigr]\succeq0,
 \label{eq:pdf2-rh-source-delta}\\
 I-(\mathscr C_{q,R}^{\rm src})^\sharp\mathscr C_{q,R}^{\rm src}
 ={}&\mathsf B_{+,q}^{\sharp}\Delta_{q,R}^{\rm src}
 \mathsf B_{+,q}\succeq0.
 \label{eq:pdf2-rh-source-contractividad}
\end{align}
Este control conserva ruta, reloj, rol, acarreo y los canales antes de formar
una diferencia. No se identifica con \(\mathcal K_R\) y no sustituye la
publicación común explícita de
\eqref{eq:pdf2-rh-doble-publicacion-material}; por ello su alcance local no
reabre la conclusión cofinal.

\section{Falsadores de una sombra ablacionada}

\begin{falsifier}[Sustituciones que cambian el objeto de partida]
El argumento anterior no se transfiere a una sombra que elimine una de las
trece coordenadas de \eqref{eq:pdf2-rh-res-sp-completa}, construya \(J_+\) y
\(J_-\) desde antecedentes distintos, reemplace una proyección ortogonal por
un idempotente no autoadjunto, omita un canal, use reguladores incompatibles,
escoja \(P_R\) después del signo, borre el terminal antes del límite o reduzca
la clase a una familia no cofinal.
\end{falsifier}

\begin{proof}
Cada operación rompe, respectivamente, el dominio de
\(\mathfrak G_{\rm sp}\), la buena definición de
\eqref{eq:pdf2-rh-vmas-rango}, el cuadrado ortogonal, la identidad funcional,
la dirección causal, la telescopía finita o el criterio reversible. El
falsador afecta a la sombra sustituida; no demueve la imagen estructural
completa empleada en el \cref{thm:pdf2-rh-resolucion}.
\end{proof}

\paragraph{Procedencia y fuerza probatoria.}
La restricción espectral--polar, la coligación source--first, su complemento y
la hoja espejo son \texttt{ARQUITECTURA\_AUTORAL\_PREEXISTENTE}. La
anulación coinductiva de \(\Omega\), la identidad finita con terminal y el
codificador de acción
\eqref{eq:pdf2-rh-xi-comun-explicita} son
\texttt{FORMALIZACION\_NUEVA} de esa estructura completa. La identificación
Guinand--Weil y el cierre cofinal son \texttt{RESULTADO\_RECUPERADO}. Las
ecuaciones
\eqref{eq:pdf2-rh-coligacion-source-first}--
\eqref{eq:pdf2-rh-tres-publicaciones} exhiben el objeto, sus dominios, su
acción, el terminal y el límite sin introducir una raíz abstracta de la forma;
su fuerza es \emph{demostrado con estructura de partida explícita}.
 \endgroup
 \part{Síntesis teoremática}

\chapter{La aritmética anterior a la realización arquimediana}
\section{VI. La estructura discreta del continuo}

\subsection{El continuo es una realización de cilindros compatibles}

El límite inverso \(X_\infty\) conserva todas las historias compatibles. Sus
cilindros semicerrados determinan una topología; la refinación determina la
forma; el cociclo determina la holonomía; la orientación determina los
bordes; y el álgebra inductiva determina los lectores. Los cinco aspectos
internos espectral, solenoidal, gauge, coherente y determinantal emergen
conjuntamente de esta construcción. Separarlos destruiría el objeto que les
da naturalidad común.

La doble proyección es una pareja de representaciones del mismo álgebra:

\[
\rho_{\rm arq}:\mathfrak H_{\rm cyl}\to\mathcal A_{\mathbb R},
\qquad
\rho_{\rm exc}:\mathfrak H_{\rm cyl}\to\mathcal A_{\rm exc}.
\]

La primera contrae cilindros compatibles hacia la lectura arquimediana. La
segunda conserva incidencia de frontera hacia
Witt--Golay--Leech--\(V^\natural\)--Monster. \(\mathbb R\) y la rama
excepcional no son dos universos de origen: son dos cocientes representativos
de una fuente más rica.

La abelianización olvida necesariamente el conmutador electrónico y parte
del registro de memoria. Por eso la recta real puede ser completa como continuo escalar y
seguir siendo insuficiente como descripción de la genealogía. La aritmética
holonómica precede a \(\mathbb R\) del mismo modo en que una historia precede
a una de sus coordenadas.

\subsection{Predefinición de futuros compatibles}

Sea

\[
\mathcal S_n(x)=\{y\in X_{n+1}:r_{n+1,n}(y)=x
\text{ y }y\text{ es admisible}\}.
\]

El árbol coinductivo contiene las prolongaciones posibles. La ley esturmiana
garantiza una única palabra especial a derecha por longitud; TRIT tipa su
curvatura y TPK selecciona la continuación conservando la rama
complementaria en la frontera. La predefinición no significa que una cifra
exterior haya sido almacenada. Significa que cada salida finita está fijada
por un cilindro y que la familia de cilindros es compatible a toda
profundidad.

Así se precisa la expresión «elección de futuros de supervivencia». El futuro
no se adivina desde una medida: se selecciona dentro del conjunto ya tipado
de prolongaciones que conservan las relaciones del estado. Pasado, presente
y futuro son funciones del operador de retorno, del umbral y de la apertura,
no simples etiquetas verbales.

\section{VII. Teorema nuclear y consecuencias}

\subsection{Enunciado mínimo}

Definimos el objeto nuclear

\[
\boxed{
\mathfrak H_{\rm HMT}^{\rm alg}
=\left[\varinjlim_n
\left(
\mathbb R\langle e_x,J_x,T_u\rangle/
\mathcal R_{\rm HMT}
\right)\right]
\rtimes_{\Gamma_9}\mathbb N.}
\]

En el sector reversible se sustituye \(\mathbb N\) por \(\mathbb Z\). Su
espacio de estados infinitos es

\[
X_\infty=\varprojlim_nX_n.
\]

La expresión es algebraica. Una representación acotada \(\rho\) determina
su propia completación
\(\mathfrak H_{\rm HMT}^{\rho}
=\overline{\rho(\mathfrak H_{\rm HMT}^{\rm alg})}^{\|\cdot\|}\);
por tanto, la topología del codominio pertenece al lector y no se antepone
al generador.

\textbf{Teorema nuclear de evolución conservativa de la unidad.} La aritmética
APP, orientada por las fibras \(J_\tau^2=-\tau I\) y transportada por el
cociclo TPK, genera un álgebra asociativa no conmutativa cuyo refinamiento
preserva la unidad y cuyos retornos elevan la memoria. El espectro tipado de
representaciones de esta álgebra determina correlativamente la
sistema correlacionado \((\pi,\varphi,e,\alpha)\), la estructura conjunta del
continuo, la vacancia electrónica, las escalas de acción y gravedad, los
parámetros geométricos, las constantes del catálogo y las simetrías de
frontera. Cada realización conserva su dominio, su lector y su grado.

La demostración se divide en seis identidades contiguas:

\begin{enumerate}
\item las ecuaciones de APP prueban la conservación de residuo y acarreo;
\item las relaciones tríticas prueban la coexistencia de las tres curvaturas;
\item la identidad de cociclo prueba la asociatividad de la concatenación;
\item la no escisión de \(C_{108}\) prueba que el retorno de fase eleva memoria;
\item la pareja límite inverso/límite directo prueba la conservación de unidad bajo refinamiento y lectura;
\item la propiedad universal prueba que cada salida tipada es una representación de la misma fuente y no una entrada convencional encubierta.
\end{enumerate}

\subsection{Consecuencia conceptual: la realidad escalar es una sombra fiel pero incompleta}

La representación arquimediana puede ser fiel respecto del valor determinado
y, sin embargo, no ser fiel respecto de la genealogía. Dos historias pueden
tener la misma coordenada real y distinto acarreo, altura o memoria. La
estructura discreta del continuo resuelve esta multiplicidad conservando las
fibras antes de contraerlas. La frase «un punto llega a ser dos» se traduce
así: una coordenada escalar puede admitir dos elevaciones holonómicas
distintas, y la bidireccionalidad vive en la fibra, no en una duplicación
arbitraria del punto.

\subsection{Consecuencia física: la unidad no conserva primero simetrías; las produce}

Lo que se conserva de manera primaria es la posibilidad de reconstrucción:
hoja, residuo, cociente, acarreo, orientación, frontera y memoria. Una
simetría aparece cuando una representación posee automorfismos que preservan
esas relaciones. De este modo, las simetrías excepcionales, la polaridad
electromagnética y la pentapolaridad gravitatoria son formas diferentes de
estabilizar una frontera ya generada. La conservación de simetría es una
consecuencia de la conservación genealógica.

\subsection{Consecuencia informacional: crecimiento infinito con coste interno nulo}

La complejidad lineal \(p(m)=m+1\) permite producir una información no
periódica ilimitada sin crecimiento exponencial de entropía. El refinamiento
añade memoria, pero la actualización integral es reversible. El coste aparece
cuando un lector decide borrar la fibra que permitiría reconstruir el paso.
Esta combinación —aperiodicidad, entropía topológica cero, memoria creciente
y unidad conservada— es la definición matemática operativa del cristal
topológico-temporal HMT.

\section{VIII. Aportaciones autónomas surgidas de la lectura retrospectiva}

\subsection{Primera aportación: dualidad límite inverso--límite directo}

La distinción entre el límite inverso de estados y el límite directo de
observables no estaba formulada de manera nuclear en los desarrollos
anteriores. Explica por qué profundización y proyección son operaciones
conjugadas sin ser inversas triviales; corrige la identificación entre unidad
global y sección individual; y proporciona el marco natural para la
predefinición.

\subsection{Segunda aportación: álgebra universal de historias}

La presentación por generadores y relaciones reúne las fibras TRIT, el orden
APP, los transportes TPK y el cociclo de registro de memoria. Su propiedad universal
demuestra que los lectores posteriores no necesitan nuevos generadores si
preservan las relaciones.

\subsection{Tercera aportación: identidad exponencial trigeométrica}

La exponencial común de \(J_+,J_0,J_-\) muestra que la rotación compleja, el
jet nilpotente y la dilatación hiperbólica son tres componentes de una sola
ley. La identidad de Euler es su proyección elíptica; \(\varphi\) es su escala
hiperbólica; \(e\) es la ley exponencial; y \(\alpha\) es la costura que pega
el jet a la holonomía dodecafásica.

\subsection{Cuarta aportación: constantes como espectro tipado de representaciones}

La noción de carácter escalar se amplía a un espectro de representaciones y
funcionales. Esto permite reunir constantes, partículas y campos sin
confundir sus codominios. La misma definición localiza Catalan como momento,
Feigenbaum como reescala no lineal, Boltzmann como transductor, Barbero--
Immirzi como covector de incidencia, Planck como sección de acción y gravedad
como carácter contragrediente.

La recuperación funcional explicita una relación adicional: la respuesta
constitutiva de incidencias \(90/120\), con su regularidad en el origen,
determina el momento de Catalán completo; su inversión recupera los
canales angulares etiquetados. La composición de transportes se expresa
entonces mediante la ley \(r\odot t=f(\psi(r)\psi(t))\), conservando
producto, orientación y traza normalizada. Sobre esa misma sección de
velocidad, la inscripción y la norma radial determinan primero la longitud
\(L\), después los productos
\(R_X\Lambda_X=L^2I\) y \(H_X\Lambda_X=\hbar_{\rm ret}cI\), y finalmente
el coeficiente común \(G=c^3L^2/\hbar_{\rm ret}\).
Las pruebas de \cref{cr28:sec:restitucion-constitutiva,g26:sec:rigidez-radial-es}
conservan la memoria y las hipótesis constitutivas de cada realización.

La constante de Feigenbaum adquiere una definición estructural precisa: es el invariante asintótico de separación de los parámetros seleccionados por los retornos críticos del lector dodecafásico uniforme. La construcción conserva la procedencia APP--TRIT--TPK, el continuo conjunto y el orden de primeras raíces. La identificación de los centros seleccionados con los centros analíticos permite aplicar el escalado a esa misma sucesión; las pruebas de \ref{hmtfg:escalado:15}--\ref{hmtfg:escalado:17} y \eqref{hmtfg:articulacion:limite} establecen el resultado paramétrico. La representación de Jacobi restituye el perfil desde su medida de doce nodos, mientras la inversión constitutiva recupera los canales orientados y el momento cuyo extremo es Catalán. Estas operaciones muestran relaciones entre lectores del mismo antecedente, conservando el dominio y la información propios de cada una.

\subsection{Quinta aportación: simetría como estabilizador de reconstrucción}

La intuición según la cual las simetrías resultan de la conservación queda
convertida en una definición de estabilizador de borde y relaciones. Esta
formulación enlaza la doble proyección, el corredor excepcional y las
representaciones físicas sin postular la simetría como fuente.

\subsection{Sexta aportación: irreducibilidad triádica del núcleo}

La imagen autoral del nudo borromeo admite una formulación exacta mediante
tres cocientes de olvido. Si se identifican las dos hojas APP, desaparecen
simultáneamente el potencial
\(\Delta_{\Sigma\Pi}\) y el conmutador
\([P_\Sigma,P_\Pi]\). Si se identifican los tres regímenes TRIT, desaparece
el espectro de signos

\[
J_+^2=-I,\qquad J_0^2=0,\qquad J_-^2=I.
\]

Si se trivializa el cociclo TPK, la relación
\(d_M(\Gamma_9x)-d_M(x)=1\) se convierte en cero y el retorno de fase queda
indistinguible del retorno de estado. Cada cociente destruye un invariante
que los otros dos factores no pueden reconstruir. Por ello APP, TRIT y TPK
son tres generadores funcionalmente irreducibles del objeto nuclear: APP
aporta la discrepancia y el orden de hojas; TRIT aporta la firma local de
curvatura; TPK aporta composición, holonomía y memoria.

Esta irreducibilidad explica por qué el núcleo no puede reducirse ni a una
tabla aritmética, ni a un álgebra cuadrática, ni a una dinámica de fases.
La estructura aparece únicamente en su composición.

\begingroup\fontsize{10.2}{11.7}\selectfont
\setlength{\parskip}{2pt}\setlength{\abovedisplayskip}{5pt}\setlength{\belowdisplayskip}{5pt}
\fontsize{9.5}{10.7}\selectfont
\setlength{\parskip}{1.5pt}\setlength{\abovedisplayskip}{2pt}\setlength{\belowdisplayskip}{2pt}
\subsection{Séptima aportación: el electrón contiene el par nulo de Witt}

La comparación retrospectiva del bloque electrónico con los tres regímenes
TRIT revela una clausura local que había permanecido implícita. A partir de

\[
S^2=I,\qquad J^2=-I,\qquad SJ=-JS
\]

se forman

\[
n_\pm=\frac{S\pm J}{2}.
\]

Entonces

\[
n_\pm^2=0,\qquad n_+n_-+n_-n_+=I.
\]

El rincón \({\rm Cl}_{1,1}\cong M_2(\mathbb R)\) contiene, por tanto, las
tres clases locales: \(J\) es elíptico, \(S\) es hiperbólico y \(n_\pm\) son
parabólicos. La vacancia electrónica es simultáneamente un defecto mínimo de
acarreo en APP y una dirección nula del álgebra escindida. TRIT no incorpora
el umbral como una corrección externa al electrón: reconoce su tipo y lo
conserva bajo el transporte TPK. La denominación de Witt es aquí la de la
descomposición nula local; la estructura de incidencia Paley--Witt se obtiene
después como representación combinatoria de frontera de la misma genealogía.

\subsection{Octava aportación: las relaciones de Weyl de los dos relojes}

La frase «la hélice contiene dos relojes» admite una presentación operatoria
exacta. Sea \(T\) el paso, \(V_9\) el carácter de fase nonádica,
\(W_\delta\) el carácter de capacidad irracional y
\(\delta=\log_{10}(10/9)\). Si \(\zeta_9\) y
\(\omega_\delta\) designan sus caracteres internos, se cumple

\[
V_9T=\zeta_9TV_9,\qquad
W_\delta T=\omega_\delta TW_\delta.
\]

Para \(\Gamma_9=T^9\),

\[
V_9\Gamma_9=\Gamma_9V_9,\qquad
W_\delta\Gamma_9=\omega_\delta^9\Gamma_9W_\delta.
\]

El reloj de fase cierra y el reloj de capacidad continúa. Si \(D_M\) mide
la altura helicoidal,

\[
[D_M,\Gamma_9]=\Gamma_9.
\]

Esta última identidad es la forma algebraica más breve de la diferencia
entre retorno de fase y retorno de estado. También localiza con precisión
la no conmutatividad: la traslación basal pertenece a un grupo abeliano,
pero su producto cruzado con los observables de fase y capacidad satisface
relaciones de Weyl no conmutativas. El cristal aperiódico HMT es así una
suspensión ordenada de un reloj finito por un carácter irracional, elevada
por un operador de memoria.

En la representación elíptica posterior,
\(\zeta_9\) y \(\omega_\delta\) se reconocen como
\(\operatorname{Exp}_+(2\pi{}J_+/9)\) y
\(\operatorname{Exp}_+(2\pi{}\delta J_+)\). El carácter HMT precede
a esa escritura compleja.

\subsection{Novena aportación: una ecuación polinómica para las tres geometrías}

Los tres generadores locales pueden reunirse en

\[
\mathbb J=p_+J_++p_0J_0+p_-J_-.
\]

La identidad

\[
\boxed{\mathbb J^6=\mathbb J^2}
\]

contiene exactamente el espectro de cuadrados \(\{-I,0,I\}\). Si
\(K=\mathbb J^2\), los regímenes se recuperan desde el propio operador:

\[
p_+=\frac{K^2-K}{2},\qquad
p_0=I-K^2,\qquad
p_-=\frac{K^2+K}{2}.
\]

Esta fórmula elimina la apariencia de que TRIT es sólo una etiqueta ternaria.
El régimen es una descomposición espectral intrínseca. La correspondiente
identidad exponencial

\[
\operatorname{Exp}_\star(t\mathbb J)=
p_+(\cos t+\mathbb J\sin t)+
p_0(I+t\mathbb J)+
p_-(\cosh t+\mathbb J\sinh t)
\]

contiene como caras la clausura elíptica de Euler, el umbral nilpotente y la
autoescala hiperbólica. Su composición con \(\Gamma_9\), unida a
\([D_M,\Gamma_9]=\Gamma_9\), constituye la identidad operatoria mínima desde
la que puede exponerse la evolución conservativa de la unidad dimensional.
  

\par\endgroup
\appendix
\label{app:desarrollos-10000-rev9}
% Copia mecánica íntegra de apendices.tex:1--761.
% Residencia material del ensamblaje sucesor de 102 capítulos.
\part{Apéndices demostrativos y tablas de clasificación: certificados algebraicos, aritméticos y coinductivos}

\chapter{Cálculo afín ordenado de las hojas APP}

\section{Semiproducto afín y composición de rutas}

Sea \(U_{\APP}\) el módulo radial anterior a su evaluación
arquimediana. La envolvente correcta es el monoide
\[
 \AffEnd(U_{\APP})
 =\operatorname{End}(U_{\APP})\ltimes U_{\APP},
\]
pues las hojas de umbral o plegamiento pueden contener endomorfismos no
invertibles. La acción de \((\Lambda,C)\) sobre \(u\in U_{\APP}\) es
\[
 (\Lambda,C)\cdot u=\Lambda u+C.
\]
La composición se calcula sin conmutar factores:
\begin{align}
 (\Lambda_2,C_2)(\Lambda_1,C_1)\cdot u
 &=\Lambda_2(\Lambda_1u+C_1)+C_2 \notag\\
 &=(\Lambda_2\Lambda_1)u+(C_2+\Lambda_2C_1).
 \label{espiral:eq:ap-comp-afin}
\end{align}
Por tanto,
\[
 (\Lambda_2,C_2)(\Lambda_1,C_1)
 =(\Lambda_2\Lambda_1,C_2+\Lambda_2C_1).
\]
La segunda coordenada no es una suma ordinaria de desplazamientos: es un
cociclo torcido por la acción lineal acumulada.

\begin{proposition}[Fórmula de los prefijos]
Para una ruta de \(N\) factores \(h_j=(\lambda_j,C_j)\), ordenados desde
\(j=0\) hasta \(j=N-1\), la holonomía prefija es
\[
 H_N=h_{N-1}\cdots h_1h_0=(\Lambda_N,\mathcal C_N),
\]
donde
\[
 \Lambda_N=\prod_{j=0}^{N-1}\lambda_j
\]
y
\[
 \mathcal C_N
 =C_{N-1}+\lambda_{N-1}C_{N-2}
 +\lambda_{N-1}\lambda_{N-2}C_{N-3}
 +\cdots+
 \lambda_{N-1}\cdots\lambda_1C_0.
\]
\end{proposition}
\begin{proof}
Para \(N=1\) la fórmula es inmediata. Si vale para \(N\), entonces
\[
 h_NH_N
 =(\lambda_N\Lambda_N,C_N+\lambda_N\mathcal C_N),
\]
que produce exactamente la expresión de longitud \(N+1\). La inducción
preserva el orden de la ruta y no identifica recorridos con los mismos
extremos.
\end{proof}

\section{Diferencia suma--producto como curvatura de orden}

Sean \(A_a=(I,a)\) y \(M_\lambda=(\lambda,0)\). En el sector invertible,
\[
 A_aM_\lambda A_a^{-1}M_\lambda^{-1}
 =A_{a(1-\lambda)}.
\]
La demostración directa sobre \(u\) es
\[
 u\xmapsto{M_\lambda^{-1}}\lambda^{-1}u
 \xmapsto{A_a^{-1}}\lambda^{-1}u-a
 \xmapsto{M_\lambda}u-\lambda a
 \xmapsto{A_a}u+a(1-\lambda).
\]
Así, la diferencia entre acumular y plegar queda publicada como una
traslación. Cuando \(\lambda=1\), el conmutador se anula; cuando
\(\lambda\neq1\), la ruta conserva una holonomía que no puede reconstruirse
a partir del punto final solamente.

\begin{example}[Dos rutas con el mismo dato visible]
Fijemos \(u_0=0\), \(a=1\), \(\lambda=2\). Las rutas
\[
 A_1M_2:u_0\longmapsto1,
 \qquad
 M_2A_{1/2}:u_0\longmapsto1
\]
publican el mismo valor final. Sin embargo, sus ledgers contienen factores,
desplazamientos y prefijos distintos. La evaluación puntual identifica las
rutas; la holonomía afín enriquecida las separa.
\end{example}

\section{Tabla exhaustiva de la carta local de 2 orientaciones}

La tabla siguiente no postula que todo grado radial global proceda de un
solo par. Registra la carta elemental obtenida al componer dos orientaciones
TRIT. El grado de una ruta general es el cociclo entero de
\cref{espiral:eq:cociclos-radiales}; esta carta proporciona sus nueve generadores
locales.

\begin{longtable}{rrrrrrp{.27\textwidth}}
\caption{Los \(9\) pares orientados de la carta TRIT local.}
\label{espiral:tab:ap-nueve-pares}\\
\toprule
\(\tau_+\)&\(\tau_-\)&\(p\)&\(a\)&\(r_3\)&\(c_3\)&Lectura local\\
\midrule
\endfirsthead
\toprule
\(\tau_+\)&\(\tau_-\)&\(p\)&\(a\)&\(r_3\)&\(c_3\)&Lectura local\\
\midrule
\endhead
\( 1\)&\( 1\)&\( 2\)&\( 0\)&\(-1\)&\( 1\)&apertura radial doble con carry positivo\\
\( 1\)&\( 0\)&\( 1\)&\( 1\)&\( 1\)&\( 0\)&apertura y umbral\\
\( 1\)&\(-1\)&\( 0\)&\( 2\)&\( 0\)&\( 0\)&equilibrio con orientación positiva\\
\( 0\)&\( 1\)&\( 1\)&\(-1\)&\( 1\)&\( 0\)&umbral y apertura\\
\( 0\)&\( 0\)&\( 0\)&\( 0\)&\( 0\)&\( 0\)&umbral doble\\
\( 0\)&\(-1\)&\(-1\)&\( 1\)&\(-1\)&\( 0\)&umbral y contracción\\
\(-1\)&\( 1\)&\( 0\)&\(-2\)&\( 0\)&\( 0\)&equilibrio con orientación negativa\\
\(-1\)&\( 0\)&\(-1\)&\(-1\)&\(-1\)&\( 0\)&contracción y umbral\\
\(-1\)&\(-1\)&\(-2\)&\( 0\)&\( 1\)&\(-1\)&contracción doble con carry negativo\\
\bottomrule
\end{longtable}

Aquí \(p=r_3+3c_3\), con residuo balanceado
\(r_3\in\{-1,0,1\}\). El carry distingue \(p=2\) de su residuo \(-1\) y
\(p=-2\) de su residuo \(1\). El par
\((p,a)\) distingue asimismo las tres genealogías locales de
grado \(0\).

\chapter{Certificado nonádico y geometría de la memoria}

\section{Distinción tipada entre catálogo decimal, palabras ternarias y fibra ambiente}

\begin{longtable}{p{.18\textwidth}p{.18\textwidth}p{.52\textwidth}}
\caption{Estratos del certificado de la conexión nonádica conjunta.}
\label{espiral:tab:certificado-nonadico}\\
\toprule
Objeto&Cardinalidad&Función matemática\\
\midrule
\endfirsthead
\toprule
Objeto&Cardinalidad&Función matemática\\
\midrule
\endhead
Catálogo de ternas de origen&\(104\,976\)&Dominio finito previo a la
selección conjunta; no es una tabla decimal.\\
Emisiones decimales \(U_6\)&\(468\)&Emisiones compatibles de longitud \(6\)
en la publicación decimal.\\
Palabras visibles \(w_6\)&\(243\)&Palabras ternarias publicadas tras el
tipado; no son el ambiente completo.\\
Ambiente \(\F_3^6\)&\(729\)&Espacio de todas las palabras ternarias de
longitud \(6\); contiene a las \(243\) visibles.\\
Niveles certificados&\(396\)&Prolongación finita hasta las fronteras
declaradas.\\
Trits recorridos&\(2\,376\)&Longitud total \(396\cdot6\).\\
Vueltas completas&\(43\)&Ciclos nonádicos completos en el tramo
certificado.\\
Pares \((K,K+9)\)&\(387\)&Pares de igual fase por canal; el ledger puede
diferir.\\
\bottomrule
\end{longtable}

La cadena de prolongación se conserva en el orden
\[
 w_6\longrightarrow w_{12}\longrightarrow w_{18}
 \longrightarrow w_{24}\longrightarrow w_{30}
 \longrightarrow R_{36}\longrightarrow G_9.
\]
Ninguno de estos objetos puede sustituirse por el cociente modal posterior.
La holonomía \(\Gamma_9\) actúa sobre la dinámica enriquecida; la
publicación de memoria reducida conserva sólo la información especificada
por su codominio.

\section{Retorno de fase, avance de memoria y estado completo}

Con
\[
 \rho_9(N)=1+((N-1)\bmod9),
 \qquad
 q_9(N)=\left\lfloor\frac{N-1}{9}\right\rfloor,
\]
se tiene
\[
 \rho_9(N+9)=\rho_9(N),
 \qquad q_9(N+9)=q_9(N)+1.
\]
Por ello,
\[
 H_9(N+9)=H_9(N)+(0,0,1).
\]
La igualdad de la primera coordenada de fase no implica igualdad del estado:
\[
 g(K+9)=g(K),
 \qquad B_{K+9}\neq B_K
\]
en general. La fórmula helicoidal es precisamente la representación mínima
de esta distinción.

\section{Extensión dodecafásica no escindida}

La sucesión
\[
 0\longrightarrow C_{12}\longrightarrow C_{108}
 \longrightarrow C_9\longrightarrow0
\]
acopla el calendario dodecafásico con el retorno nonádico. Si existiera una
sección homomorfa global \(s:C_9\to C_{108}\), el generador de \(C_9\)
debería elevarse a un elemento de orden divisor de \(9\) que proyectase en el
generador. El lift temporal activo conserva, sin embargo, el desplazamiento
dodecafásico acumulado y no identifica el estado tras una vuelta de fase.
La costura se registra como carry; geométricamente se manifiesta como
traslación axial de la hélice.

\section{5 regiones de \texorpdfstring{\(\pi\)}{pi} y no fidelidad de la evaluación}

La microfibra de \(\pi\) contiene cinco regiones orientadas con
multiplicidad total
\[
 432+4\cdot144=1008.
\]
Las regiones publican el mismo valor real, pero conservan genealogías
distintas. Este resultado proporciona el precedente interno de la
clasificación espiral: una curva visible no determina por sí sola la región,
la hoja, el carry ni la memoria que la produjo.

\chapter{Álgebra potencia--logarítmica y 5 publicaciones canónicas}

\section[Coordenada potencia--logarítmica e inversa]{La coordenada \(L_p\) y su inversa}

Para \(x>0\), se define
\[
 L_p(x)=
 \begin{cases}
 (x^p-1)/p,&p\neq0,\\
 \log x,&p=0.
 \end{cases}
\]
Su inversa es
\[
 \exp_p(u)=
 \begin{cases}
 (1+pu)^{1/p},&p\neq0,\ 1+pu>0,\\
 \e^u,&p=0.
 \end{cases}
\]
La derivada
\[
 \frac{\dd}{\dd x}L_p(x)=x^{p-1}>0
\]
prueba que \(L_p\) es estrictamente creciente en su dominio y, por tanto,
invertible.

\begin{proposition}[Ley suma--producto]
Para \(x,y>0\),
\[
 L_p(xy)=L_p(x)+L_p(y)+pL_p(x)L_p(y).
\]
\end{proposition}
\begin{proof}
Si \(p\neq0\),
\begin{align*}
 L_p(x)+L_p(y)+pL_p(x)L_p(y)
 &=\frac{x^p-1}{p}+\frac{y^p-1}{p}
   +\frac{(x^p-1)(y^p-1)}p\\
 &=\frac{x^py^p-1}{p}=L_p(xy).
\end{align*}
Para \(p=0\), la igualdad se reduce a
\(\log(xy)=\log x+\log y\).
\end{proof}

\begin{corollary}[Grupo deformado en el dominio positivo]
La operación
\[
 u\boxplus_pv=u+v+puv
\]
satisface
\[
 1+p(u\boxplus_pv)=(1+pu)(1+pv).
\]
Su neutro es \(0\) y el inverso de \(u\) es
\[
 \ominus_pu=-\frac{u}{1+pu},
\]
si \(1+pu\neq0\).
\end{corollary}

\section{Derivación de las 5 curvas canónicas}

Sea \(m\) la profundidad publicada y supongamos
\[
 L_p(r/r_0)=\beta m,
 \qquad
 \theta=\theta_0+\omega m,
 \qquad\omega\neq0.
\]
Eliminando \(m=(\theta-\theta_0)/\omega\),
\[
 r(\theta)=r_0
 \left[1+p\frac{\beta}{\omega}
 (\theta-\theta_0)\right]^{1/p},
 \qquad p\neq0,
\]
y
\[
 r(\theta)=r_0
 \exp\!\left[\frac{\beta}{\omega}
 (\theta-\theta_0)\right]
\]
para \(p=0\). Las reoriginaciones y reescalados positivos producen las
formas convencionales:

\begin{longtable}{rllp{.31\textwidth}}
\caption{Publicaciones canónicas del sector traslacional.}
\label{espiral:tab:cinco-publicaciones-detalle}\\
\toprule
\(p\)&Ecuación normal&Nombre convencional&Dominio y rasgo\\
\midrule
\endfirsthead
\toprule
\(p\)&Ecuación normal&Nombre convencional&Dominio y rasgo\\
\midrule
\endhead
\(2\)&\(r^2=a^2\theta\)&Fermat&\(\theta\ge0\); incremento lineal del área
radial cuadrática.\\
\(1\)&\(r=a\theta\)&Arquímedes&\(\theta\ge0\); incremento radial
constante por unidad angular.\\
\(0\)&\(r=r_0\e^{b\theta}\)&logarítmica&\(\theta\in\R\); dilatación
constante por unidad angular.\\
\(-1\)&\(r=a/\theta\)&hiperbólica&\(\theta>0\); producto \(r\theta\)
constante.\\
\(-2\)&\(r^2=a^2/\theta\)&lituus&\(\theta>0\); producto
\(r^2\theta\) constante.\\
\bottomrule
\end{longtable}

\section{Inversión de hojas}

Para \(p\neq0\),
\[
 L_{-p}(x^{-1})
 =\frac{x^p-1}{-p}=-L_p(x).
\]
La identidad también vale para \(p=0\). Por ello, la inversión radial
intercambia \(p\) con \(-p\) y cambia el signo de la coordenada
linealizada. Esta relación agrupa Fermat con el lituus y Arquímedes con la
hiperbólica como orientaciones conjugadas del lector, sin identificarlas
como curvas iguales.

\section{Sector afín general}

Si la coordenada radial satisface
\[
 u_{n+1}=\Lambda u_n+C,
\]
la iteración interna exacta es
\[
 u_n=\Lambda^n u_0+\sum_{k=0}^{n-1}\Lambda^kC.
\]
Si \(I-\Lambda\) es invertible,
\(u_*=(I-\Lambda)^{-1}C\) y
\(u_n-u_*=\Lambda^n(u_0-u_*)\). Sólo en una realización escalar positiva,
\(\chi\Lambda=\lambda\chi\) con \(\lambda>0\), se define
\(u_*=C/(1-\lambda)\). Con \(\theta_n=\theta_0+n\omega\),
\(\omega\neq0\), la interpolación publicada en esa realización es
\[
 r(\theta)=r_0\exp_p\!\left[
 u_*+\lambda^{(\theta-\theta_0)/\omega}(u_0-u_*)
 \right].
\]
La taxonomía completa no es, pues, un listado de cinco nombres. Contiene al
menos el grado \(p\), la orientación \(a\), la clase afín
\([\Lambda,C]\), el cociclo de carry, la frontera, el incremento angular y
el ledger de memoria.

\chapter{Naturalidad, equivariancia y no fidelidad}

\section{El lector de prefijos}

Para una historia truncada
\[
 \widehat\gamma_N=\widehat e_N\cdots\widehat e_1,
\]
el lector enriquecido registra
\[
 \mathcal P_{\rad,N}^{\enr}(\widehat\gamma_N)
 =
 \left(
 (H_0,\Led_0);
 \bigl(p(\widehat e_j),a(\widehat e_j),P_j,A_j,
 h_{e_j},H_j,\Led_j\bigr)_{j=1}^{N}
 \right).
\]
El término \(H_j=h_{e_j}\cdots h_{e_1}\) conserva los prefijos. El ledger
\(\Led_j\) conserva, según la ruta, hoja, residuo, cociente, carry,
frontera, orientación, cilindro, firma, supervivencia y memoria. Las
coordenadas reales se obtienen después:
\[
 \mathcal P_{\arq}^{\mathrm{esp}}
 =\ev_{\R^2}\circ\mathcal P_{\rad}^{\enr}.
\]

\section{Prueba completa de naturalidad bajo truncamiento}

Sea \(\pi_{N,M}\) el borrado del sufijo de longitud \(N-M\), con \(M<N\).
Sea \(\operatorname{pr}_{N,M}\) la proyección que conserva el registro
inicial y los primeros \(M\) registros de arista. Entonces
\[
 \operatorname{pr}_{N,M}\,
 \mathcal P_{\rad,N}^{\enr}(\widehat\gamma_N)
 =
 \left(
 (H_0,\Led_0);
 \bigl(p(\widehat e_j),a(\widehat e_j),P_j,A_j,
 h_{e_j},H_j,\Led_j\bigr)_{j=1}^{M}
 \right).
\]
Por definición de truncamiento,
\[
 \pi_{N,M}(\widehat\gamma_N)=\widehat e_M\cdots\widehat e_1,
\]
y al volver a aplicar el lector se obtiene la misma secuencia. Por tanto,
\[
 \operatorname{pr}_{N,M}\circ
 \mathcal P_{\rad,N}^{\enr}
 =
 \mathcal P_{\rad,M}^{\enr}\circ\pi_{N,M}.
\]
La demostración no usa el valor de la curva. Usa la identidad de prefijos,
que es la información que el lector reducido de siete coordenadas habría
perdido.

\section{Refinamientos no triviales}

Para un refinamiento \(r:\gamma\to\gamma'\), la naturalidad requiere un
mapa \(\widehat r\) sobre ledgers y la igualdad local
\[
 h_e
 =
 h_{e'_k}\cdots h_{e'_1}
\]
para cada arista \(e\) sustituida por la cadena refinada
\((e'_1,\ldots,e'_k)\). Si además el carry, la frontera y la memoria de la
cadena refinada componen los registros del paso original, entonces
\[
 \widehat r\circ\mathcal P_{\rad}^{\enr}
 =
 \mathcal P_{\rad}^{\enr}\circ r.
\]
Esta condición es un cuadrado verificable; no debe suponerse para un
refinamiento arbitrario sin exhibir los factores.

\section{Equivarianza, no invariancia}

La acción de \(\Gamma_9\) devuelve la fase y prolonga el estado. En el
codominio del lector actúa
\(\widehat\Gamma_9^{\rad}\), que desplaza los registros y actualiza el
último prefijo. La identidad correcta es
\[
 \mathcal P_{\rad}^{\enr}(\Gamma_9\gamma)
 =
 \widehat\Gamma_9^{\rad}
 \mathcal P_{\rad}^{\enr}(\gamma).
\]
Exigir invariancia borraría precisamente la memoria que la conexión
conserva. En el sector reversible, la inversión de rutas satisface
\[
 H_{C_{\mathrm{rev}}\gamma}=H_\gamma^{-1}
\]
cuando todos los factores pertenecen al subgrupo afín invertible. Esta
afirmación tipada no se extiende automáticamente a una involución global
sobre cualquier cociente posterior.

\section{3 demostraciones de no fidelidad}

\begin{proposition}[Circunferencia intrínseca y circunferencia proyectada]
Existen dos historias con la misma curva visible
\[
 c(\theta)=(\cos\theta,\sin\theta)
\]
y ledgers distintos.
\end{proposition}
\begin{proof}
La primera historia posee memoria axial constante \(m=0\). La segunda es
la hélice
\[
 h(\theta)=(\cos\theta,\sin\theta,\theta/2\pi),
\]
cuyo ledger registra \(m(\theta)=\theta/2\pi\). Ambas satisfacen
\(\pi_{xy}\circ h=c\). La proyección es, por tanto, no inyectiva.
\end{proof}

\begin{proposition}[Mismo extremo, distinto orden afín]
Si \(a(1-\lambda)\neq0\), los caminos \(A_aM_\lambda\) y
\(M_\lambda A_a\) no son equivalentes en el lector enriquecido, aunque
puedan calibrarse para publicar un mismo extremo.
\end{proposition}
\begin{proof}
Su conmutador es \(A_{a(1-\lambda)}\neq I\). El dato del conmutador
pertenece al ledger de prefijos y desaparece al conservar sólo el extremo.
\end{proof}

\begin{proposition}[Misma curva, distinta región genealógica]
La igualdad del valor real publicado no identifica las cinco regiones de
la microfibra de \(\pi\); por consiguiente, tampoco basta para identificar
la clase espiral cuando el lector conserva región, frontera o carry.
\end{proposition}

\chapter{Taxonomía estratificada de correspondencias espirales}

\section{Clasificación por nivel matemático}

\begin{longtable}{p{.20\textwidth}p{.30\textwidth}p{.39\textwidth}}
\caption{Niveles que no deben colapsarse en una sola noción de curvatura.}
\label{espiral:tab:niveles-curvatura}\\
\toprule
Nivel&Invariante&Significado\\
\midrule
\endfirsthead
\toprule
Nivel&Invariante&Significado\\
\midrule
\endhead
APP&\(F(S,P)=1,\ F(P,S)=-1\)&Curvatura discreta del orden de acumulación
y plegamiento.\\
TRIT&\(\tau\in\{-1,0,1\}\), \(J_\tau^2=-\tau I\)&Régimen local elíptico,
parabólico o hiperbólico del transporte.\\
TPK&\(\Gamma_9\), carry, frontera, ruta y ledger&Holonomía de fase y
memoria sobre el estado holonómico integral.\\
Radial&\(p,a,[\Lambda,C]\)&Caracteres que determinan la ley
radial antes de reconocer una curva.\\
Arquimediano&\(r(\theta)\)&Curva publicada en \(\R^2\); puede olvidar
genealogía.\\
Frenet&\(\kappa_{\mathrm{F}},\mathcal T_{\mathrm{F}}\)&Curvatura y torsión
diferenciales de la curva o suspensión visible.\\
Dimensional&campos, escalas y representaciones&Realización física posterior,
con premisas propias.\\
\bottomrule
\end{longtable}

\section{Clasificación por giro, radio y memoria}

\begin{longtable}{ccccl}
\caption{Taxonomía cinemática completa de los casos degenerados mínimos.}
\label{espiral:tab:taxonomia-cinematica-ap}\\
\toprule
\(\Delta\theta\)&\(\Delta u_{\rad}\)&\(\Delta m\)&Salida visible&
Información genealógica\\
\midrule
\endfirsthead
\toprule
\(\Delta\theta\)&\(\Delta u_{\rad}\)&\(\Delta m\)&Salida visible&
Información genealógica\\
\midrule
\endhead
\(0\)&\(0\)&\(0\)&punto&estado estacionario publicado\\
\(0\)&\(\neq0\)&\(0\)&recta radial&cambio de escala\\
\(0\)&\(0\)&\(\neq0\)&fibra axial&memoria sin desplazamiento transversal\\
\(0\)&\(\neq0\)&\(\neq0\)&hoja radial suspendida&escala y memoria\\
\(\neq0\)&\(0\)&\(0\)&circunferencia&cierre angular intrínseco\\
\(\neq0\)&\(0\)&\(\neq0\)&hélice cilíndrica&fase cerrada, memoria elevada\\
\(\neq0\)&\(\neq0\)&\(0\)&espiral plana&carácter radial visible\\
\(\neq0\)&\(\neq0\)&\(\neq0\)&suspensión helicoidal radial&giro, radio y
profundidad conservados\\
\bottomrule
\end{longtable}

\section{Clasificación por holonomía afín}

La tabla siguiente pertenece a una realización escalar real de la parte
lineal, \(\chi\Lambda=\lambda\chi\); para aligerar su lectura se escribe
\(\Lambda\) en la columna de condiciones. No clasifica endomorfismos
abstractos por desigualdades numéricas.

\begin{longtable}{p{.18\textwidth}p{.22\textwidth}p{.48\textwidth}}
\caption{Tipos de holonomía sobre la coordenada radial.}
\label{espiral:tab:tipos-holonomia}\\
\toprule
Condición&Órbita de \(u\)&Interpretación\\
\midrule
\endfirsthead
\toprule
Condición&Órbita de \(u\)&Interpretación\\
\midrule
\endhead
\(\Lambda=1,\ C=0\)&\(u_n=u_0\)&radio constante; la memoria puede ser
nula u oculta por proyección.\\
\(\Lambda=1,\ C\neq0\)&\(u_n=u_0+nC\)&sector traslacional que publica la
familia \(L_p\).\\
\(0<\Lambda<1\)&\(u_n\to u_*\)&contracción hacia una sección fija.\\
\(\Lambda>1\)&\(|u_n-u_*|\) crece&expansión desde una sección fija.\\
\(\Lambda<0\)&alternancia de orientación&requiere controlar el dominio de
\(\exp_p\) en cada paso.\\
\(\Lambda=0\)&\(u_{n+1}=C\)&plegamiento no invertible; pertenece al
monoide, no al grupo afín.\\
\bottomrule
\end{longtable}

\section{Criterio de equivalencia genealógica}

Dos rutas no se identifican por publicar la misma curva. Se exige un
isomorfismo tipado que preserve término a término:
\[
 \bigl(
 (x_0,H_0,\Led_0);
 (\widehat e_j,p(\widehat e_j),a(\widehat e_j),P_j,A_j,
 h_{e_j},H_j,\Led_j)_{j=1}^{N}
 \bigr)
\]
y que conmute con truncamientos y refinamientos admitidos. La clasificación
resultante es más fina que la taxonomía convencional porque incorpora la
historia que precede a la aparición de \(\R^2\).

\chapter{Realizaciones posteriores y controles de tipo}

\section{Modo electromagnético circular}

Sea
\[
 \zeta=\frac{z-c_{\mathrm{vac}}t}{\lambda},
\]
y
\[
 \mathbf E_\sigma
 =E_0(\mathbf e_1\cos2\pi\zeta
 +\sigma\mathbf e_2\sin2\pi\zeta),
 \qquad
 \mathbf B_\sigma
 =c_{\mathrm{vac}}^{-1}\widehat{\mathbf k}\times\mathbf E_\sigma.
\]
Entonces
\[
 \mathbf E_\sigma\perp\mathbf B_\sigma\perp\widehat{\mathbf k},
 \quad |E|=c_{\mathrm{vac}}|B|,
 \quad\omega=c_{\mathrm{vac}}k.
\]
A tiempo fijo, el extremo normalizado del campo en el fibrado
fase--posición describe
\[
 (\cos2\pi\zeta,\sin2\pi\zeta,\zeta).
\]
Esta afirmación no dice que toda línea espacial de campo sea una hélice:
describe la trayectoria del extremo del vector polarización junto con el
evento marcado.

\section{Compatibilidad loxodrómica}

Para
\[
 \Phi(r,\varphi,z,t)
 =\ell\varphi+\nu\log(r/r_0)+k_zz-\omega t,
\]
la transformación
\[
 (r,\varphi)\longmapsto
 ((2+\sqrt3)r,\varphi+\pi/2)
\]
preserva la fase módulo \(2\pi\) exactamente cuando
\[
 \nu\log(2+\sqrt3)+\ell\frac{\pi}{2}\in2\pi\Z.
\]
La dilatación se interpreta como cambio de escala entre soluciones; no
como producción espontánea de energía.

\section{Centro electrónico y lector posterior}

El electrón ocupa el punto fijo único \((5,5)\) del atlas material. Esta
centralidad no se deriva otra vez en el presente trabajo. La compatibilidad
operatoria previa incluye
\[
 C^{(3)}M_{\ph}=-\frac12M_{\ph},
 \qquad
 1-\left(\frac12\right)^2=\frac34=\frac{90}{120},
\]
y
\[
 \|[P_{\Sigma,3},P_{\Pi,3}]\|=\frac{\sqrt3}{4}.
\]
La teoría de correspondencias espirales añade un lector de rutas; no
reemplaza la Ley General de Masas ni usa el electrón para seleccionar
retrospectivamente \(p\), \(\Lambda\) o \(C\).

\section{Matriz cerrada de controles negativos}

\begin{longtable}{p{.38\textwidth}p{.50\textwidth}}
\caption{Falsadores y errores de tipo excluidos por la construcción.}
\label{espiral:tab:falsadores-ap}\\
\toprule
Fallo&Consecuencia\\
\midrule
\endfirsthead
\toprule
Fallo&Consecuencia\\
\midrule
\endhead
La curva clásica selecciona \(p,\Lambda,C\).&Circularidad: el lector deja
de ser constructivo--generativo.\\
Se conserva sólo la holonomía final.&Pérdida de prefijos, carry y frontera;
no se acredita naturalidad.\\
Se exige invariancia bajo \(\Gamma_9\).&Se borra el avance de memoria; la
relación correcta es equivarianza.\\
Se confunden \(U_6\), \(w_6\) y \(\F_3^6\).&Colapso de objetos con
cardinalidades \(468\), \(243\) y \(729\).\\
Se identifica retorno de fase con retorno de estado.&Contradicción con
\(B_{K+9}\neq B_K\) en general.\\
Se identifica \(\tau\), \(p\) y \(\kappa_{\mathrm{F}}\).&Colapso de régimen,
carácter radial y curvatura visible.\\
Se llama hélice general a cualquier radio variable.&Afirmación diferencial
sin verificar la condición de Lancret.\\
Se identifica el eje \(\varphi\) con el electrón.&Mezcla de la fibra de
constantes y la realización material.\\
Se interpreta Pell como amplificación energética.&Violación de la lectura
de escala y de la conservación física.\\
Se identifica helicidad electromagnética con espín \(1/2\).&Confusión entre
representaciones distintas.\\
Se impone una celda puntual a toda especie.&Contradicción con el carácter
de fibra u multisección de las especies no centrales.\\
Se separan las cinco construcciones del continuo.&Ruptura de la genealogía
única anterior a cualquier publicación geométrica.\\
\bottomrule
\end{longtable}

\chapter{Glosario matemático}

\begin{longtable}{p{.22\textwidth}p{.66\textwidth}}
\caption{Símbolos y términos rectores.}
\label{espiral:tab:glosario}\\
\toprule
Símbolo o término&Definición\\
\midrule
\endfirsthead
\toprule
Símbolo o término&Definición\\
\midrule
\endhead
\(\Sigma,\Pi\)&Evaluaciones aditiva y multiplicativa sobre el soporte
común APP.\\
\(\tau\)&Régimen local de orientación TRIT; no es el grado radial ni la
torsión de Frenet.\\
\(M_{\ph}\)&Selector trítico de fase operativa.\\
\(\Gamma_9\)&Holonomía nonádica de la dinámica enriquecida.\\
\(H_N\)&Holonomía afín acumulada hasta el prefijo \(N\).\\
\(\Lambda_N,\mathcal C_N\)&Parte lineal y cociclo traslacional de \(H_N\).\\
\(p\)&Cociclo entero del carácter radial.\\
\(a\)&Carácter orientado local que separa genealogías del mismo
grado local.\\
\(L_p,\exp_p\)&Coordenada potencia--logarítmica y su inversa positiva.\\
\(\boxplus_p\)&Ley deformada que publica la interacción suma--producto.\\
\(\Led_N\)&Ledger enriquecido de hoja, carry, frontera, ruta y memoria.\\
\(\mathcal P_{\rad}^{\enr}\)&Lector radial enriquecido anterior a la curva.\\
\(\mathcal P_{\arq}^{\mathrm{esp}}\)&Evaluación posterior en una curva de
\(\R^2\).\\
\(\kappa_{\mathrm{F}},\mathcal T_{\mathrm{F}}\)&Curvatura y torsión de Frenet de la
curva visible.\\
Suspensión helicoidal radial&Superficie u hoja de radio variable que
conserva fase, profundidad y recorrido.\\
Publicación no fiel&Evaluación que puede identificar historias distintas.\\
Constructivo--generativo&Método que construye el objeto y su genealogía
antes de reconocer su apariencia convencional.\\
\bottomrule
\end{longtable}
  \chapter{Diccionario de dominios, operadores y representaciones}
\chapter*{Notación, tipos y desambiguación de símbolos}
\addtocontents{toc}{\protect\enlargethispage{2\baselineskip}}
\addcontentsline{toc}{chapter}{Notación, tipos y desambiguación de símbolos}

\begin{longtable}{@{}p{.22\textwidth}p{.70\textwidth}@{}}
\toprule
Símbolo & Objeto \\ \midrule
\endhead
\(\mathcal A_9\)&Toro aritmético nonádico de caminos con dos evaluaciones
internas.\\
\(\mathrm{TRIT}\)&Coordenada ternaria orientada del transporte; la salida
visible no agota el levantamiento enriquecido con acarreo.\\
\(\mathrm{TPK}\)&Núcleo topológico de fase con estado holonómico integral,
transiciones, rutas, prolongaciones y memoria.\\
\(M_{\rm ph}\)&Selector ternario de fase operativa, distinto del transporte
cardinal y del muestreo estroboscópico.\\
\(\tau_3\)&Bloque primitivo de fase \((+1,-1,0)\).\\
\(\mathcal N,\mathcal E,\mathcal S,\mathcal O\)&Generadores del transporte
cardinal.\\
\(\operatorname{Obs}_3\)&Muestreo estroboscópico de cada tercer paso.\\
\(\mathcal R_{12}\)&Sistema temporal dodecafásico orientado interno a TPK.\\
\(\Theta_{108}=\mathbb Z/108\mathbb Z\)&Calendario observable
\(12\times9\); no es el estado TPK completo.\\
\(C_{12}\)&Rotación de los doce sectores temporales.\\
\(\Gamma_{108}\)&Órbita orientada de transporte de \(108\) pasos.\\
\(\mathcal G_9\)&Poda unitaria dentro de la relación de extensión y
frontera; no constituye por sí sola la monodromía.\\
\(\mathfrak M_9\)&Monodromía nonádica con memoria: retorna la fase y cambia
el estado.\\
\(\pi_{\rm term}\)&Proyección del estado holonómico terminal a los tres
bloques visibles seleccionados.\\
\(R_{\rm sgn}\)&Lectura de la coordenada entera firmada del mismo estado
holonómico terminal.\\
\(U_{12}^{\rm cat}\)&Concatenación visible \(U_6\Vert U_6\).\\
\(U_{12}^{\rm sgn}\)&Coordenada firmada del estado holonómico integral, reversible
con \(K\) mediante \(\mathcal H_{12}\); no se deduce de
\(U_{12}^{\rm cat}\), de la proyección calendárica ni de los tres bloques terminales
aislados.\\
\(\mathcal H_{12}\)&Transformación dodecafásica reversible entre
\(U_{12}^{\rm sgn}\) y \(K\).\\
\(\mathcal E_{\rm dod}\)&Extractor integral
\((D_3,D_4,Q):\mathbb Z^{12}\to\mathbb Z^{25}\) que reconstruye \(K\).\\
\(\overline{\mathcal C}_9^{\,3}\)&Completación nonádica cúbica.\\
\(A,C^*\)&Coordenadas angular y contraangular intrínsecas.\\
\(q_\pm\)&Canales orientados del lector bicapa.\\
\(D_{\rm raw}^{90,120}\)&Diferencia diagnóstica cruda de series en una
carta fijada; no es un extractor canónico.\\
\(\iota_{6,8}\)&Inclusión
\(\mathcal F_6\hookrightarrow\mathcal F_8\) de incidencias hexada--octada.\\
\(K_{6\to8}\)&Funcional angular de las dos subcolas de Lambert, posterior a
\(q_\pm\), las series, el polo, la positividad y la minimalidad.\\
\(\langle90,120\rangle\)&Semigrupo espectral que gobierna la jerarquía de
torres.\\
\(\varepsilon_{\rm res}^{\circ}\)&Residuo de la identidad angular exacta;
no es una extracción finita de los objetos 90/120 anteriores.\\
\(\Delta_4\)&Torsión mínima o invariante torsional global anterior a la
especie; no es la torsión de Cartan.\\
\(\mathsf T_{\alpha,\infty}\)&Prolongación remota del acarreo de la vía
entera de \(\alpha\), tipada separadamente del generador orbital
\(T_1^{\rm orb}=L_1\), del contador neutral y del calendario.\\
\(N_1^{\rm neu}\)&Número de posiciones neutrales en la primera ventana del
ejemplo de transporte; vale $3$ y no es un operador.\\
\(T_1^{\rm orb}:=L_1\)&Segundo generador en el calibre orbital local; en su
sección de referencia se abrevia \(T_1\).\\
\(\varepsilon_{\rm clk}\)&Cuanto energético del reloj,
\(\hbar_{\rm int}2\pi/(108t_0)\); no es una base energética canónica sin
declarar el marco temporal. El símbolo \(\varepsilon_0\) queda reservado a
la permitividad del vacío en la interfaz constitutiva SI.\\
\(\mathrm{CLK}1\text{--}\mathrm{CLK}6\)&Etiquetas de las seis ecuaciones del
reloj finito, desde el Hamiltoniano hasta Schrödinger.\\
\(\mathcal K_{27}\)&Retorno compuesto de veintisiete actualizaciones.\\
\(X_\infty^{\rm cal}\)&Dominio enriquecido común de las evaluaciones
globales.\\
\bottomrule
\end{longtable}

\section*{Cardinales repetidos con fuentes diferentes}

La igualdad de un cardinal no transporta automáticamente dominio,
operación, memoria ni fuerza probatoria. La tabla siguiente fija las
colisiones que más afectan a la lectura.

\begin{longtable}{@{}p{.12\textwidth}p{.27\textwidth}p{.51\textwidth}@{}}
\toprule
Índice & Objeto tipado & Función y separación \\ \midrule
\endhead
\(108\)&\(\Theta_{108}\)&Calendario \(12\times9\); cierra el reloj
observable.\\
\(108\)&\(\Pi_{108}^{\rm cal}\)&Proyección calendárica de un cursor con \(324\)
semillas; olvida segundo cursor, orientación, acarreo, frontera, registro
cronológico enriquecido y
supervivencia.\\
\(108\)&\(w_{108}\)&Palabra de \(108=18\times6\) trits; no es el calendario.\\
\(104,976\)&\(X_{\rm TPK}^{\rm seed}\)&Motor de dos cursores,
\(104,976=324^2\); conserva un dominio distinto del calendario.\\
\(1080\)&\(w_{1080}\)&Longitud de una palabra.\\
\(1080\)&\(\Gamma_0(1080)\)&Nivel de un subgrupo modular; requiere su
realización propia.\\
\(729\)&\(\mathbb F_3^6\)&Ambiente de palabras ternarias.\\
\(729\)&\(9^3\)&Célula cúbica/interior EMRH.\\
\(729\)&\(\operatorname{Ext}(X_K)\)&Número de subcilindros examinados en cada
etapa de prolongación.\\
\(729\)&bulk EMRH&Estrato interior de la completación base mil.\\
\(27\)&\(\mathcal K_{27}\)&Operador de retorno compuesto; no nombre de una
teoría autónoma.\\
\(27\)&estrato cúbico&Contribución de dos coordenadas de frontera en
\(999=729+243+27\).\\
\(27\)&módulo o clasificador&Cada aparición exige su mapa antes de
identificarse con otra.\\
\bottomrule
\end{longtable}

\section*{Seis fuentes distintos de 90/120}

\begin{longtable}{@{}
>{\raggedright\arraybackslash}p{.21\textwidth}
>{\raggedright\arraybackslash}p{.42\textwidth}
>{\raggedright\arraybackslash}p{.27\textwidth}@{}}
\toprule
Objeto & Dominio y operación & Salida y prohibición de tipo \\ \midrule
\endhead
\(\mathcal E_{\rm dod}\)&\(K\in\mathbb Z^{12}\); diferencias por pasos
\(3,4\) y suma \(Q\)&Coordenadas transversales dodecafásicas; no incidencia
ni serie angular.\\
\(D_{\rm raw}^{90,120}\)&Una carta \(q\); diferencia cruda
\(S_{90}(q)-S_{120}(q)\)&Diagnóstico de escala, distinto del funcional de
Barbero--Immirzi y de los operadores de incidencia.\\
\(\iota_{6,8}\)&\(H\subset O\); inclusión de
\(H\times\binom{H}{2}\) en \(O\times\binom{H}{2}\)&Cardinales \(90,120\) y
defecto \(-1/12\); no extractor ni serie.\\
\(K_{6\to8}\)&Canales \(q_\pm\) y subcolas de Lambert; combinación
\(12\Delta S_{90}-\Delta S_{120}\)&Funcional angular relativo a sus
premisas; no se transporta como bloque de masa.\\
\(\mathfrak S=\langle90,120\rangle\)&Alturas de momentos
\(s_n=q_-^n-q_+^n\)&Semigrupo global de torres; no termina en ocho alturas.\\
\bottomrule
\end{longtable}

\section*{Cuatro realizaciones de la reversibilidad bidireccional}

\begin{longtable}{@{}
  >{\raggedright\arraybackslash}p{.26\textwidth}
  >{\raggedright\arraybackslash}p{.26\textwidth}
  >{\raggedright\arraybackslash}p{.38\textwidth}@{}}
\toprule
Operación & Dominio & Acción \\ \midrule
\endhead
\(P_\pm=(I\pm R)/2\)&Álgebra escindida&Proyecta los dos sectores
espectrales.\\
Inversión de rumbo&Ruta y calendario compuesto&Invierte la orientación del
transporte sin borrar su palabra.\\
\(\operatorname{Ext}^{\pm}\)&Árbol de prolongaciones&Compara valencias
futuras y pasadas; su balance puede leerse mediante
\(\log(\#\operatorname{Ext}^+/\#\operatorname{Ext}^-)\).\\
\(C_M(\tau)=-\operatorname{rev}(\tau)\)&Palabra dodecafásica&Conjuga orden
y signo; no es la involución celular \(\rho_c\).\\
\bottomrule
\end{longtable}

La función común de estas cuatro realizaciones es conservar orientación
reversible. No autoriza a identificar retorno de fase, retorno de ruta,
retorno de palabra y retorno de estado.
  
\backmatter
\begingroup
\sloppy
\small
\setlength{\parskip}{0pt}
\setlength{\emergencystretch}{3em}
\bibliographystyle{alphaurl}
\begin{thebibliography}{HvdHKN76}

\bibitem[ALY17]{AbeLamYamada2017}
Toshiyuki Abe, Ching~Hung Lam, and Hiromichi Yamada.
\newblock A remark on {$\mathbb Z_p$}-orbifold constructions of the moonshine
  vertex operator algebra.
\newblock {\em arXiv preprint arXiv:1705.09022}, 2017.
\newblock \href {https://arxiv.org/abs/1705.09022} {\path{arXiv:1705.09022}}.

\bibitem[AW68]{rm:ArakiWoods1968}
H.~Araki and E.~J. Woods.
\newblock A classification of factors.
\newblock {\em Publ. Res. Inst. Math. Sci. Ser. A}, 4:51--130, 1968.
\newblock URL: \url{https://ems.press/content/serial-article-files/41616}.

\bibitem[Bek73]{rm:Bekenstein1973}
J.~D. Bekenstein.
\newblock Black holes and entropy.
\newblock {\em Physical Review D}, 7:2333--2346, 1973.

\bibitem[BG95a]{Barbero}
J.~Fernando Barbero~G.
\newblock Real {Ashtekar} variables for {Lorentzian} signature space-times.
\newblock {\em Physical Review D}, 51:5507--5510, 1995.

\bibitem[BG95b]{Barbero1995}
J.~Fernando Barbero~G.
\newblock Real {Ashtekar} variables for {Lorentzian} signature space-times.
\newblock {\em Physical Review D}, 51:5507--5510, 1995.

\bibitem[Bor92]{Borcherds1992}
Richard~E. Borcherds.
\newblock Monstrous moonshine and monstrous lie superalgebras.
\newblock {\em Inventiones Mathematicae}, 109:405--444, 1992.

\bibitem[{Bur}19]{BIPMSI}
{Bureau International des Poids et Mesures}.
\newblock {\em The International System of Units ({SI} Brochure)}.
\newblock BIPM, S\`evres, 9 edition, 2019.
\newblock Updated online edition.
\newblock \href {https://doi.org/10.59161/AUEZ1291}
  {\path{doi:10.59161/AUEZ1291}}.

\bibitem[Car14]{Caratheodory1914}
Constantin Carath{\'e}odory.
\newblock {"U}ber das lineare ma{\ss} von punktmengen---eine verallgemeinerung
  des l{"a}ngenbegriffs.
\newblock {\em Nachrichten von der Gesellschaft der Wissenschaften zu
  G{"o}ttingen, Mathematisch-Physikalische Klasse}, pages 404--426, 1914.

\bibitem[Car17]{Carnahan2017}
Scott Carnahan.
\newblock 51 constructions of the moonshine module.
\newblock {\em arXiv preprint arXiv:1707.02954}, 2017.
\newblock \href {https://arxiv.org/abs/1707.02954} {\path{arXiv:1707.02954}}.

\bibitem[CLS16]{ChenLamShimakura2016}
Hsian-Yang Chen, Ching~Hung Lam, and Hiroki Shimakura.
\newblock A {$\mathbb Z_3$}-orbifold construction of the moonshine vertex
  operator algebra and some maximal 3-local subgroups of the monster.
\newblock {\em arXiv preprint arXiv:1606.05961}, 2016.
\newblock \href {https://arxiv.org/abs/1606.05961} {\path{arXiv:1606.05961}}.

\bibitem[CN79]{ConwayNorton1979}
John~H. Conway and Simon~P. Norton.
\newblock Monstrous moonshine.
\newblock {\em Bulletin of the London Mathematical Society}, 11:308--339, 1979.

\bibitem[{COD}24]{CODATA2022}
{CODATA Task Group on Fundamental Constants}.
\newblock The 2022 {CODATA} recommended values of the fundamental physical
  constants, 2024.
\newblock URL: \url{https://physics.nist.gov/constants}.

\bibitem[CS99]{ConwaySloane1999}
John~H. Conway and Neil J.~A. Sloane.
\newblock {\em Sphere Packings, Lattices and Groups}.
\newblock Springer, 3 edition, 1999.

\bibitem[Fei78]{rm:Feigenbaum1978}
Mitchell~J. Feigenbaum.
\newblock Quantitative universality for a class of nonlinear transformations.
\newblock {\em Journal of Statistical Physics}, 19(1):25--52, 1978.
\newblock \href {https://doi.org/10.1007/BF01020332}
  {\path{doi:10.1007/BF01020332}}.

\bibitem[FLM88]{FLM1988}
Igor Frenkel, James Lepowsky, and Arne Meurman.
\newblock {\em Vertex Operator Algebras and the Monster}.
\newblock Academic Press, 1988.

\bibitem[GH77]{rm:GibbonsHawking1977}
G.~W. Gibbons and S.~W. Hawking.
\newblock Cosmological event horizons, thermodynamics, and particle creation.
\newblock {\em Physical Review D}, 15:2738--2751, 1977.

\bibitem[Haw75]{rm:Hawking1975}
S.~W. Hawking.
\newblock Particle creation by black holes.
\newblock {\em Communications in Mathematical Physics}, 43:199--220, 1975.

\bibitem[HJ13]{HornJohnson2013}
Roger~A. Horn and Charles~R. Johnson.
\newblock {\em Matrix Analysis}.
\newblock Cambridge University Press, 2 edition, 2013.

\bibitem[Hol96]{Holst1996}
Sören Holst.
\newblock Barbero's hamiltonian derived from a generalized hilbert--palatini
  action.
\newblock {\em Physical Review D}, 53:5966--5969, 1996.

\bibitem[HvdHKN76]{Hehl1976}
Friedrich~W. Hehl, Paul von~der Heyde, G.~David Kerlick, and James~M. Nester.
\newblock General relativity with spin and torsion: Foundations and prospects.
\newblock {\em Reviews of Modern Physics}, 48:393--416, 1976.

\bibitem[IK02]{IosifescuKraaikamp}
M.~Iosifescu and C.~Kraaikamp.
\newblock {\em Metrical Theory of Continued Fractions}.
\newblock Kluwer, 2002.

\bibitem[Imm97a]{Immirzi}
Giorgio Immirzi.
\newblock Real and complex connections for canonical gravity.
\newblock {\em Classical and Quantum Gravity}, 14:L177--L181, 1997.

\bibitem[Imm97b]{Immirzi1997}
Giorgio Immirzi.
\newblock Real and complex connections for canonical gravity.
\newblock {\em Classical and Quantum Gravity}, 14:L177--L181, 1997.

\bibitem[Khi64]{Khinchin}
A.~Ya. Khinchin.
\newblock {\em Continued Fractions}.
\newblock University of Chicago Press, 1964.

\bibitem[Kib61]{Kibble1961}
T.~W.~B. Kibble.
\newblock Lorentz invariance and the gravitational field.
\newblock {\em Journal of Mathematical Physics}, 2:212--221, 1961.

\bibitem[K{"o}36]{Koenig1936}
D{\'e}nes K{"o}nig.
\newblock {\em Theorie der endlichen und unendlichen Graphen}.
\newblock Akademische Verlagsgesellschaft, Leipzig, 1936.

\bibitem[Kol33]{Kolmogorov1933}
Andrei~N. Kolmogorov.
\newblock {\em Grundbegriffe der Wahrscheinlichkeitsrechnung}, volume 2, Heft 3
  of {\em Ergebnisse der Mathematik und ihrer Grenzgebiete}.
\newblock Julius Springer, Berlin, 1933.

\bibitem[Lan82]{rm:Lanford1982}
III Lanford, O.~E.
\newblock A computer-assisted proof of the {Feigenbaum} conjectures.
\newblock {\em Bulletin of the American Mathematical Society}, 6:427--434,
  1982.

\bibitem[Man05]{Manivel2005}
Laurent Manivel.
\newblock Configurations of lines and models of lie algebras.
\newblock arXiv:math/0507118, 2005.
\newblock \href {https://arxiv.org/abs/math/0507118}
  {\path{arXiv:math/0507118}}.

\bibitem[MS64]{rm:MisnerSharp1964}
C.~W. Misner and D.~H. Sharp.
\newblock Relativistic equations for adiabatic, spherically symmetric
  gravitational collapse.
\newblock {\em Physical Review}, 136:B571--B576, 1964.

\bibitem[MS77]{MacWilliamsSloane1977}
Florence~Jessie MacWilliams and Neil J.~A. Sloane.
\newblock {\em The Theory of Error-Correcting Codes}.
\newblock North-Holland, 1977.

\bibitem[Pla01]{Planck}
M.~Planck.
\newblock Ueber das gesetz der energieverteilung im normalspectrum.
\newblock {\em Annalen der Physik}, 309:553--563, 1901.

\bibitem[Sci64]{Sciama1964}
Dennis~W. Sciama.
\newblock The physical structure of general relativity.
\newblock {\em Reviews of Modern Physics}, 36:463--469, 1964.

\bibitem[Ser77]{Serre1977Representations}
Jean-Pierre Serre.
\newblock {\em Linear Representations of Finite Groups}.
\newblock Springer, 1977.

\bibitem[Tak03]{rm:Takesaki}
M.~Takesaki.
\newblock {\em Theory of Operator Algebras II}.
\newblock Springer, 2003.

\bibitem[TMNT21]{Tiesinga2021CODATA}
Eite Tiesinga, Peter~J. Mohr, David~B. Newell, and Barry~N. Taylor.
\newblock {CODATA} recommended values of the fundamental physical constants:
  2018.
\newblock {\em Journal of Physical and Chemical Reference Data}, 50(3):033105,
  2021.
\newblock \href {https://doi.org/10.1063/5.0064853}
  {\path{doi:10.1063/5.0064853}}.

\bibitem[Wie93]{Wien}
W.~Wien.
\newblock Eine neue beziehung der strahlung schwarzer k{\"o}rper zum zweiten
  hauptsatz der w{\"a}rmetheorie.
\newblock {\em Sitzungsberichte der K{\"o}niglich Preussischen Akademie der
  Wissenschaften}, pages 55--62, 1893.

\end{thebibliography}
 
\endgroup
\begingroup
\setlength{\columnsep}{2.2em}\fontsize{9}{10.4}\selectfont
\renewcommand{\indexspace}{\par\vspace{3pt}}
\printindex
\endgroup
\end{document}
